\documentclass{book}
\usepackage{amsmath}
\usepackage{amssymb}
\usepackage{amsthm}
\usepackage{amssymb}

% For dealing with references we use the comment environment
\usepackage{verbatim}
\newenvironment{reference}{\comment}{\endcomment}
%\newenvironment{reference}{}{}
\newenvironment{slogan}{\comment}{\endcomment}
\newenvironment{history}{\comment}{\endcomment}

% For commutative diagrams we use Xy-pic
\usepackage[all]{xy}

% We use 2cell for 2-commutative diagrams.
\xyoption{2cell}
\UseAllTwocells

% We use multicol for the list of chapters between chapters

% This is generally recommended for better output
\usepackage{lmodern}
\usepackage[T1]{fontenc}

% For cross-file-references
% We prefer the snapshot tmp/xref/#2.aux,
% which latexmkrc replaces atomically, so
% that documents compiled simultaneously
% never read a half-written .aux file.
% The last argument keeps links to #2.pdf.

% Package for hypertext links:
\usepackage{hyperref}

% For any local file, say "hello.tex" you want to link to please
% use \externaldocument[hello-]{hello}

% Heap Project documents
variations-of-hodge-structures-]{%
variations-of-hodge-structures}
homological-mirror-symmetry-]{%
homological-mirror-symmetry}
geometric-prequantization-]{%
geometric-prequantization}


16th-alga-meeting-2026-]{%
16th-alga-meeting-2026}

iv-jornadas-lefschetz-2026-]{%
iv-jornadas-lefschetz-2026}






% Theorem environments.
%
\theoremstyle{plain}
\newtheorem{theorem}[subsection]{Theorem}
\newtheorem{proposition}[subsection]{Proposition}
\newtheorem{lemma}[subsection]{Lemma}

\theoremstyle{definition}
\newtheorem{definition}[subsection]{Definition}
\newtheorem{example}[subsection]{Example}
\newtheorem{exercise}[subsection]{Exercise}
\newtheorem{situation}[subsection]{Situation}

\theoremstyle{remark}
\newtheorem{remark}[subsection]{Remark}
\newtheorem{remarks}[subsection]{Remarks}

\numberwithin{equation}{subsection}

% Macros
%
\def\lim{\mathop{\mathrm{lim}}\nolimits}
\def\colim{\mathop{\mathrm{colim}}\nolimits}
\def\Spec{\mathop{\mathrm{Spec}}}
\def\Hom{\mathop{\mathrm{Hom}}\nolimits}
\def\Ext{\mathop{\mathrm{Ext}}\nolimits}
\def\SheafHom{\mathop{\mathcal{H}\!\mathit{om}}\nolimits}
\def\SheafExt{\mathop{\mathcal{E}\!\mathit{xt}}\nolimits}
\def\Sch{\mathit{Sch}}
\def\Mor{\mathop{\mathrm{Mor}}\nolimits}
\def\Ob{\mathop{\mathrm{Ob}}\nolimits}
\def\Sh{\mathop{\mathit{Sh}}\nolimits}
\def\NL{\mathop{N\!L}\nolimits}
\def\CH{\mathop{\mathrm{CH}}\nolimits}
\def\proetale{{pro\text{-}\acute{e}tale}}
\def\etale{{\acute{e}tale}}
\def\QCoh{\mathit{QCoh}}
\def\Ker{\mathop{\mathrm{Ker}}}
\def\Im{\mathop{\mathrm{Im}}}
\def\Coker{\mathop{\mathrm{Coker}}}
\def\Coim{\mathop{\mathrm{Coim}}}

% Boxtimes
%
\DeclareMathSymbol{\boxtimes}{\mathbin}{AMSa}{"02}

%
% Macros for moduli stacks/spaces
%
\def\QCohstack{\mathcal{QC}\!\mathit{oh}}
\def\Cohstack{\mathcal{C}\!\mathit{oh}}
\def\Spacesstack{\mathcal{S}\!\mathit{paces}}
\def\Quotfunctor{\mathrm{Quot}}
\def\Hilbfunctor{\mathrm{Hilb}}
\def\Curvesstack{\mathcal{C}\!\mathit{urves}}
\def\Polarizedstack{\mathcal{P}\!\mathit{olarized}}
\def\Complexesstack{\mathcal{C}\!\mathit{omplexes}}
% \Pic is the operator that assigns to X its picard group, usage \Pic(X)
% \Picardstack_{X/B} denotes the Picard stack of X over B
% \Picardfunctor_{X/B} denotes the Picard functor of X over B
\def\Pic{\mathop{\mathrm{Pic}}\nolimits}
\def\Picardstack{\mathcal{P}\!\mathit{ic}}
\def\Picardfunctor{\mathrm{Pic}}
\def\Deformationcategory{\mathcal{D}\!\mathit{ef}}
\begin{document}

\tableofcontents

\part{Analysis}
\label{part-analysis}

% --- Begin real-analysis.tex ---
\chapter{Real Analysis}
\label{chapter-real-analysis}





\section{Uniform convergence of series}
\label{real-analysis-section-uniform-convergence-series}

\begin{theorem}[Weierstrass M-test]
\label{real-analysis-theorem-m-test}
Suppose that $(f_n)$ is a sequence of real or
complex-valued functions defined on a set
$A$, and that there is a sequence of
non-negative numbers $(M_n)$ satisfying
\begin{itemize}
\item $|f_n(x)|<M_n$ for all $n\geq1$ and
all $x\in A$, and
\item $\sum_{n=1}^\infty M_n$ converges.
\end{itemize}
Then the series
$\sum_{n=1}^\infty f_n(x)$ converges
absolutely and uniformly on $A$.
\end{theorem}

\begin{proof}
Let $S_n$ be the sequence of partial sums.
For any $x\in A$ and $m>n>N$,
$
|S_m(x)-S_n(x)|
\leq
\sum_{k=n+1}^m |f_k(x)|
<
\sum_{k=n+1}^m M_k.
$
Since $\sum_k M_k$ converges, its partial
sums are Cauchy. Thus the right hand side can
be made smaller than any $\varepsilon>0$,
independently of $x$. Hence $(S_n)$ is
uniformly Cauchy, and therefore converges
uniformly. The same estimate gives absolute
convergence.
\end{proof}

\section{Differentiation of series}
\label{real-analysis-section-differentiation-series}

The following is the open-domain form of
\cite[Section 16.3.5, Corollary 5,
pp.~389--390]{zo2}.

\begin{theorem}[Differentiation of series]
\label{real-analysis-theorem-differentiation-series}
Let $E$ be a real or complex normed vector
space and let $U\subset E$ be nonempty,
open, bounded and convex. Suppose that
$f_n:U\to\mathbb{C}$ are differentiable,
that $\sum_{n=1}^{\infty}f_n(x_0)$
converges at some $x_0\in U$, and that
$\sum_{n=1}^{\infty}Df_n$ converges
uniformly on $U$ in the operator norm.
Then $\sum_{n=1}^{\infty}f_n$ converges
uniformly on $U$ to a differentiable
function $f$, and
$$
Df(x)=\sum_{n=1}^{\infty}Df_n(x).
$$
Here $Df_n(x)$ is the derivative as a
bounded linear map. Its operator norm is
$$
\|Df_n(x)\|
=\sup_{\|v\|\leq1}|Df_n(x)v|.
$$
\end{theorem}

\begin{proof}
Write $S_N=\sum_{n=1}^N f_n$ and let $G$
be the uniform limit of $DS_N$. Choose
$R>0$ with $\|x-x_0\|\leq R$ for every
$x\in U$. The mean value inequality along
the segment from $x_0$ to $x$ gives
$$
\begin{aligned}
|S_N(x)-S_M(x)|
&\leq |S_N(x_0)-S_M(x_0)|\\
&\quad+R\sup_{y\in U}
\|DS_N(y)-DS_M(y)\|.
\end{aligned}
$$
Thus $(S_N)$ is uniformly Cauchy and has a
uniform limit $f$.

Fix $x\in U$ and $\varepsilon>0$. Choose
$N$ so that, for every $M\geq N$,
$$
\sup_{y\in U}\|DS_M(y)-DS_N(y)\|
\leq\varepsilon,
\qquad
\|G(x)-DS_N(x)\|\leq\varepsilon.
$$
Apply the same mean value inequality to
$S_M-S_N$ between $x$ and $x+h$ in $U$.
Letting $M\to\infty$ gives
$$
|(f-S_N)(x+h)-(f-S_N)(x)|
\leq\varepsilon\|h\|.
$$
Consequently, for $h\ne0$,
$$
\begin{aligned}
\frac{|f(x+h)-f(x)-G(x)h|}{\|h\|}
&\leq 2\varepsilon\\
&\quad+
\frac{|S_N(x+h)-S_N(x)-DS_N(x)h|}
{\|h\|}.
\end{aligned}
$$
The last quotient tends to zero as
$h\to0$, since $S_N$ is differentiable.
As $\varepsilon$ is arbitrary, $Df(x)=G(x)$.
\end{proof}

\begin{remark}
\label{real-analysis-remark-continuous-derivatives-series}
Continuity of the derivatives is not needed
in the theorem. If each $f_n$ is $C^1$,
then $Df$ is continuous as a uniform limit
of continuous functions, so $f$ is $C^1$.
\end{remark}

\begin{lemma}[Corollary]
\label{real-analysis-lemma-differentiation-series-interval}
Let $a<b$ and $f_n\in C^1([a,b],\mathbb{C})$.
If $\sum_n f_n$ converges at one point of
$[a,b]$ and $\sum_n f_n'$ converges uniformly
on $[a,b]$, then $\sum_n f_n$ converges
uniformly to a $C^1$ function $f$, with
$$
f'=\sum_{n=1}^{\infty}f_n'.
$$
Derivatives at the endpoints are one-sided.
\end{lemma}

\begin{proof}
Let $S_N=\sum_{n=1}^N f_n$, let $g$ be the
uniform limit of $S_N'$, and choose $x_0$
where $S_N(x_0)$ converges to $c$. Since
$$
S_N(x)=S_N(x_0)+\int_{x_0}^x S_N'(t)\,dt,
$$
uniform convergence of the derivatives
gives uniform convergence to
$$
f(x)=c+\int_{x_0}^x g(t)\,dt.
$$
The function $g$ is continuous, so the
fundamental theorem of calculus gives
$f\in C^1([a,b])$ and $f'=g$.
\end{proof}

\begin{remark}
\label{real-analysis-remark-smooth-series}
For smooth summands on $[a,b]$, if the
original series converges at one point and
the series of $k$th derivatives converges
uniformly for every $k\geq1$, repeated
application of the corollary gives a smooth
sum with every derivative computed termwise.
\end{remark}

\section{Density of continuous functions in $L^p$}
\label{real-analysis-section-density-continuous-lp}

The following is \cite[Proposition~7.9]{folland-real-analysis}.

\begin{proposition}
\label{real-analysis-proposition-cc-dense-lp}
Let $X$ be a locally compact Hausdorff space and let $\mu$ be a
Radon measure on $X$. For every $1\leq p<\infty$,
\[
C_c(X)\text{ is dense in }L^p(X,\mu).
\]
\end{proposition}

\begin{remark}
\label{real-analysis-remark-continuous-dense-lp-compact}
If $X$ is compact, then $C_c(X)=C(X)$. Therefore
\[
\overline{C(X)}^{\,L^p}=L^p(X,\mu),
\qquad 1\leq p<\infty.
\]
In particular, for Lebesgue measure on $S^1$,
\[
\overline{C(S^1,\mathbb C)}^{\,L^2}=L^2(S^1).
\]
\end{remark}

\section{Fréchet derivative}
\label{real-analysis-section-Fr-chet-derivative}

\begin{definition}
\label{real-analysis-definition-Frechet-derivative}
\cite{zo1}, Section 8.2.1, Definition 1. A
function $f:E \to \mathbb{R}^n$
defined on a set $E \subset \mathcal{R}^m$ is
{\it differentiable} at the point
$x \in E$, which is a limit point of $E$, if
$$
f(x+h)-f(x)=L(x)h+\alpha(x;h)
$$
where $L(x):\mathbb{R}^m\to \mathbb{R}^n$ is
a function that is linear in $h$
and $\alpha(x;h)=o(h)$ as $h\to 0$, $x+h\in
E$.
\end{definition}

This means that
$$
\lim_{\|h\|\to 0}
\frac{\|f(x+h)-f(x)-Ah\|}{\|h\|}=0
$$
\begin{lemma}
\label{real-analysis-lemma-differentiating-bilinear-form}
Let $V$ be a finite dimensional real vector
space and let $B:V\times V\to \mathbb{R}$ be
a bilinear form. If $a(t),b(t)$ are
differentiable
curves in $V$, then
$$
\frac{d}{dt}B(a(t),b(t))
=
B(a'(t),b(t))+B(a(t),b'(t)).
$$
\end{lemma}

\begin{proof}
Choose a basis $e_1,\ldots,e_n$ of $V$. Write
$$
a(t)=\sum_i a_i(t)e_i,\qquad
b(t)=\sum_j b_j(t)e_j.
$$
If $B_{ij}=B(e_i,e_j)$, then by bilinearity
$$
B(a(t),b(t))=\sum_{i,j}B_{ij}a_i(t)b_j(t).
$$
Therefore the ordinary product rule gives
$$
\frac{d}{dt}B(a(t),b(t))
=
\sum_{i,j}B_{ij}a_i'(t)b_j(t)
+
\sum_{i,j}B_{ij}a_i(t)b_j'(t).
$$
Packing the two sums back into
coordinate-free notation, this is
$$
B(a'(t),b(t))+B(a(t),b'(t)).
$$
\end{proof}

\section{Arzelà-Ascoli Theorem}
\label{real-analysis-section-Arzela-Ascoli-theorem}
\begin{definition}[Equicontinuous family]
\label{real-analysis-definition-equicontinuous}
From Wikipedia. Let $X$ be a compact
Hausdorff space and let  $C(X)$ be the
space of real-valued continuous functions on
$X$. A subset $\mathcal{F} \subset
C(X)$ is said to be {\it equicontinuous} if
for every $x in X$ and every
$\varepsilon>0$, $x$ has a neighbourhood
$U_x$ such that
$$
\forall y \in U_x, \forall f \in \mathcal{F}:
\qquad |f(y)-f(x)|<\varepsilon
$$
\end{definition}

\begin{definition}[Pointwise bounded]
\label{real-analysis-definition-pointwise-bounded}
A set $\mathcal{F}\subset C(X,\mathbb{R})$ is
said to be {\it pointwise
bounded} if for every $x \in X$,
$$
\text{sup}\{|f(x)|:f \in \mathcal{F}\}<\infty
$$
\end{definition}

\begin{theorem}[Arzelà-Ascoli]
\label{real-analysis-theorem-Arzela-Ascoli}
Let $X$ be a compact Hausdorff space. Then a
subset $\mathcal{F}$ of $C(X)$ is
relatively compact in the topology induced by
the uniform norm if and only if it
is equicontinuous and pointwise bounded.
\end{theorem}


% --- End real-analysis.tex ---

% --- Begin complex-analysis.tex ---
\chapter{Complex Analysis}
\label{chapter-complex-analysis}





\section{Introduction}
\label{complex-analysis-section-Introduction}

The basic result of complex analysis,
Cauchy's integral formula
\ref{complex-analysis-equation-cauchy-integral-formula}
is very well explained:
if follows from Cauchy integral theorem
\ref{complex-analysis-theorem-cauchy},
which is nothing else than an application of
Stokes theorem,
(in this case also Green's theorem),
which I do not prove in this document.
After Cauchy's integral formula
I prove that a function is holomorphic
if and only if it is analytic (i.e. has
convergent
power series expansion in a neighbourhood of
the point)
in Theorem
\ref{complex-analysis-theorem-holomorphic-iff-analytic}.
I use this to prove Riemann removable
singularity theorem
\ref{complex-analysis-theorem-removable-singularity},
from which I can explain how to ``factor
zeroes'' of a
holomorphic function as in Equations
\ref{complex-analysis-equation-zero}
and \ref{complex-analysis-equation-factor-zeroes}.

The so-called Argument principle,
which is presented as Exercise
\ref{complex-analysis-exercise-argument-principle},
follows easily from the above development.
Also Morera's and Liouville's
theorems \ref{complex-analysis-theorem-Morera} and
\ref{complex-analysis-theorem-Liouville} follow easily
and their proofs are outlined.
Other important results not stated so far
also follow,
in particular: the
\href{https://en.wikipedia.org/wiki/%
Identity_theorem}{
Identity theorem},
which says that if two functions coincide in
a
set with an accumulation point they will
coincide in their domains,
the open mapping theorem, and the maximum
principle.

Next in Ahlfors book is a generalization of
Cauchy's formula
for cycles. And after that is the Residue
theorem.
Next are harmonic functions, in particular
the mean value theorem and Schwarz theorem.

And after all that elliptic functions.

\medskip\noindent
Then multivariable complex analysis,
with a beautifully worked out proof of
Weierstrass preparation theorem
\ref{complex-analysis-theorem-weier-prepa-theor},
culminating in the proof that the ring of
holomorphic
functions $\mathcal{O}_n$ is a UFD via the
Gauss lemma
from Commutative Algebra
\ref{commutative-algebra-lemma-gauss}.

\section{Holomorphic functions}
\label{complex-analysis-section-holomorphic-functions}

\begin{definition}
\label{complex-analysis-definition-holomorphic-function}
A function $f:\Omega\subset \mathbb{C} \to
\mathbb{C}$
is {\it holomorphic} at $z_0 \in \Omega$ if
$$
\lim_{z\to z_0} \frac{f(z)-f(z_0)}{z-z_0}
$$
exists.
\end{definition}

\noindent
This is equivalent to
\begin{equation}
\label{complex-analysis-equation-holom-funct-secon-defin}
\lim_{h\to 0} \frac{f(z_0+h)-f(z_0)}{h}
\end{equation}
where $h$ is a complex parameter.


\section{Cauchy-Riemann equations}
\label{complex-analysis-sectioncauchy-riemann-equations}

\begin{theorem}
\label{complex-analysis-theorem-cauchy-riemann}
Let $f:\Omega\subset\mathbb{C}\to \mathbb{C}$
be a function. Let $z=x+iy$ coordinates
in $\mathbb{C}$. Define $f:=u+iv$ and
$z_0=x_0+iy_0 \in \Omega$. $f$ is holomorphic
at $z_0$ if and only if
$$
\frac{\partial u}{\partial x}=\frac{\partial
v}{\partial y},\qquad
\frac{\partial v}{\partial x}=-\frac{\partial
u}{\partial y}.
$$
\end{theorem}

\begin{proof}
The idea is to use Equation
\ref{complex-analysis-equation-holomorphic-second-definition}
as definition of holomorphic function.
Let's differentiate along the $x$ axis.
Let $h \in \mathbb{R}$
(this time $h$ {\it is} just a real number,
while in the definition of holomorphic
function
it is a complex parameter).
We see that
$$
\begin{aligned}
f(z_0+h)-f(z_0)&=\Big(u(x_0+h,y_0)+iv(x_0+h,
y_0)\Big)-\Big(u(x_0,y_0)+iv(x_0,y_0)\Big)\\
&=\Big(u(x_0+h,y_0)-u(x_0,y_0)\Big)+
i\Big(v(x_0+h,y_0)-v(x_0,y_0)\Big)
\end{aligned}
$$

\noindent
Takin limit as $h \to 0$ we obtain
\begin{equation}
\label{complex-analysis-equation-crproof1}
\frac{\partial u}{\partial x}+i\frac{\partial
v}{\partial x}.
\end{equation}
If we differentiate along the $y$ axis we
shall obtain
\begin{equation}
\label{complex-analysis-equation-crproof2}
\frac{\partial v}{\partial y}-i\frac{\partial
u}{\partial y}.
\end{equation}
Indeed, we let $h$ be again a real number,
but now we multiply by $i$:
$$
\begin{aligned}
f(z_0+ih)-f(z_0)&=\Big(u(x_0,y_0+ih)+iv(x_0,
y_0+ih)\Big)-\Big(u(x_0,y_0)+iv(x_0,y_0)\Big)
\\
&=\Big(u(x_0,y_0+ih)-u(x_0,y_0)\Big)+
i\Big((v(x_0,y_0+ih)-v(x_0,y_0)\Big)\ldots
\end{aligned}
$$

\noindent
To agree with our definition we must take
limit as $ih \to 0$.
Thus we obtain
$$
\begin{aligned}
\lim_{ih \to 0}\frac{f(z+ih)-f(z)}{ih}
&=\frac{1}{i}\left(\frac{\partial u}{\partial
y}+i\frac{\partial v}{\partial y}\right)
=\frac{\partial v}{\partial
y}-i\frac{\partial u}{\partial y}
\end{aligned}
$$

\noindent
To conclude observe that since $f$ is
holomorphic,
The expressions \ref{complex-analysis-equation-crproof1} and
\ref{complex-analysis-equation-crproof2}
must both equal the derivative of $f$ at
$z_0$.
Matching real and imaginary parts gives the
result.

\medskip\noindent
For the converse first observe that if $u$
and $v$
are two real-valued functions of two real
variables
satisfying Cauchy-Riemann equations
\ref{complex-analysis-equation-cauchy-riemann},
then they are harmonic, i.e.
$$
\begin{aligned}
\Delta u&=\frac{\partial^2 u}{\partial x^2}+
\frac{\partial^2 u}{\partial y^2}=0\\
\Delta v=\frac{\partial^2 v}{\partial x^2}+
\frac{\partial^2v}{\partial y^2}=0.
\end{aligned}
$$

\noindent
Now if $u$ and $v$ have continuous
first-order partial
derivatives, then we know (?) that
$$
\begin{aligned}
u(x+h,y+k)-u(x,y)
&=\frac{\partial u}{\partial
x}h+\frac{\partial u}{\partial
y}k+\varepsilon_1\\
v(x+h,y+k)-v(x,y)
&=\frac{\partial v}{\partial
x}h+\frac{\partial v}{\partial
y}k+\varepsilon_2.
\end{aligned}
$$

\noindent
For real numbers $h,k,\varepsilon_1$ and
$\varepsilon_2$,
where $\varepsilon_1,\varepsilon_2$ tend to
zero
more rapidly than $h+ik$ in the sense that
$\varepsilon_1/(h+ik) \to 0$ and
$\varepsilon_2/(h+ik)\to 0$
for $h+ik\to 0$. Then we compute that
$$
\begin{aligned}
f(z+h+ik)-f(z)
&=u(z+h+ik)+iv(z+h+ik)-(u(z)+iv(z))\\
&=u(x+h,y+k)+iv(x+h,y+k)-u(x,y)-iv(x,y)\\
&=u_xh+u_yk+\varepsilon_1+i(v_xh+v_yk+
\varepsilon_2)
&=h(u_x+iv_x)+k(u_y+iv_y)+\varepsilon_1+
i\varepsilon_2\\
&=h(u_x+iv_x)+k(-v_x+iu_x)+\varepsilon_1+
i\varepsilon_2
\qquad \text{by Cauchy-Riemann}\\
&=h(u_x+iv_x)+ik(iv_x+u_x)+\varepsilon_1+
i\varepsilon_2.
\end{aligned}
$$

\noindent
Then we see that
$$
\lim_{h+ik\to 0}\frac{f(z+h+ik)-f(z)}{h+ik}
=\frac{\partial u}{\partial
x}+i\frac{\partial v}{\partial x},
$$
so that the limit exists and thus $f$ is
holomorphic.
\end{proof}

\section{Integration}
\label{complex-analysis-section-integration}

\noindent
The complex integral may first be defined for
a complex-valued function $f$
defined on a real integral as
$$
\int_a^bf=\int_a^b u + i \int_a^bv.
$$
And then for curves as
$$
\int_\gamma f:=\int_a^b
f(\gamma(t))\gamma'(t)dt.
$$
\section{Cauchy integral theorem}
\label{complex-analysis-section-cauchy-integral}

\noindent
The upshot about complex integration
(and complex analysis, really)
is that the integral $\int_\Delta fdz$
vanishes
by Stokes theorem when $f$ is holomorphic.
This leads to Cauchy integral formula and
several other results.
Beware: this works only when the domain is
simply connected.

\begin{theorem}[Cauchy]
\label{complex-analysis-theorem-cauchy}
If $f(z)$ is analytic in an open disk
$\Delta$, then
\begin{equation}
\label{complex-analysis-equationcauchy-theorem}
\int_\gamma f(z)dz=0
\end{equation}
for every closed curve $\gamma$ in $\Delta$.
\end{theorem}

\begin{proof}
The following proof comes from
\href{https://en.wikipedia.org/wiki/Cauchy
% 27s_integral_theorem}{wiki}.
First write
$$
\int_\gamma f(z)dz=\int_\gamma(u+iv)(dx+idy)
=\int_\gamma (udx-vdy)+i\int_\gamma(vdx+udy).
$$
Now apply Stokes theorem. For example, the
second integral is
$$
\int_{\partial\Delta}(vdx+udy)=\int_\Delta
d(vdx+udy)
=\int_\Delta(\partial_xu-\partial_yv)
\text{Vol}.
$$
There's Cauchy-Riemann equations!
\end{proof}

\noindent
The following result is essential for proving
Cauchy integral formula.
We are interested in taking a point away from
the domain
and seeing if we can extend the function
holomorphically.
The condition we need for this to work is
given by the following theorem:

\begin{theorem}
\label{complex-analysis-theorem-cauchy-removable}
Let $f$ be analytic in the region $\Delta'$
by omitting a finite number of
points $z_j$ from an open disk $\Delta$. If
$f$ satisfies that
\begin{equation}
\label{complex-analysis-equation-removable-singularity}
\lim_{z\to z_j} (z-z_j)f(z)=0
\end{equation}
for all $j$, then $\int_\gamma fdz=0$ for any
closed curve
$\gamma$ in $\Delta'$.
\end{theorem}

\begin{proof}[Proof]
So I think it might be using logarithmic
derivative. The limit above allows to
write
$$
|f(z)|\leq \frac{\varepsilon}{|z-z_j|}
$$
Then integrate. On the left we can bound the
integral of $f$ after applying
integral triangle inequality, and on the
right we have logarithmic derivative of
$z-z_j$. This will be a fixed number (the
index, see Definition
\ref{complex-analysis-definition-index}), so that we have
effectively bounded the integral.
\end{proof}

\section{Cauchy integral formula}
\label{complex-analysis-sectioncauchy-integral-formula}

\noindent
The way to Cauchy's integral formula is
basically:
holomorphic functions satisfy Cauchy-Riemann
equations,
which give Cauchy's integral theorem
(which is just a consequence Stokes)
and then apply this to the function
$\frac{f(z)-f(z_0)}{z-z_0}$.

The theorem says that we can compute the
value of
a holomorphic function at a point as an
integral around the point.
The crucial point here is that the function
inside the integral
has a singularity at the point because of the
fraction $\frac{1}{z-z_0}$.
This is to be understood as the fact
that the domain of holomorphy is not simply
connected, and thus
(we cannot apply Stokes) the integral does
not vanish ---
it gives the value of the function at the
point!

To prove a general version we shall use the
notion of index of a curve
about a point, which tells us how many times
the curve winds about the point.

\begin{lemma}
\label{complex-analysis-lemma-index-is-multiple-of-2pii}
\begin{reference}
\cite[Section 2.2, Lemma 1]{ahl}
\end{reference}
If the piecewise differentiable closed curve
$\gamma$ does not pass through the
point $a$, then the value of the integral
$$
\int_\gamma\frac{dz}{z-a}
$$
is a multple of $2\pi i$.
\end{lemma}

An idea (though not at proof) for a proof is
the following:
We are tempted to simply write the integrand
as the logarithmic derivative of
the function $f(z)=z-a$. But this isn't quite
right, we must be careful with the
domain of the logarithm. But it is
instructive to see the computation:
$$
\int_\gamma\frac{dz}{z-a}=\int_\gamma d
\text{log}(z-a)=
\int_\gamma d\text{log}|z-a|+i\int_\gamma d
\text{arg}(z-a)
$$
If  $\gamma$ is closed then $\text{log}|z-a|$
would return to its initial value
and $\text{arg}(z-a)$ increases or decreases
by a multiple of $2\pi$. The actual
proof by Ahlfors is different.

\begin{proof}[Proof for circle]
Still, the proof for a circle about a point
$z_0$ is also trivial: this is just
the case of the integral
$\int_{\partial\Delta}\frac{dz}{z-z_0}$, but
the curve
here is  $\gamma(t)=e^{it}+z_0$ so that when
substituting in the parametrization
we obtain
$\int_0^{2\pi}\frac{ie^{it}}{(e^{it}+z_0)-
z_0}=2\pi
i$.
\end{proof}

\begin{definition}
\label{complex-analysis-definition-index}
The {\it index} of the point $a$ with respect
to the closed curve $\gamma$ is
\begin{equation}
\label{complex-analysis-equation-index}
n(\gamma,a):=\frac{1}{2\pi
i}\int_\gamma\frac{dz}{z-a}
\end{equation}
\end{definition}

\begin{theorem}[Cauchy Integral Formula]
\label{complex-analysis-theorem-cauchy-integral-formula}
\begin{reference}
\cite[Section 2.2, Theorem 6]{ahl}
\end{reference}
Suppose that $f(z)$ is analytic in an open
disk $\Delta$, and let $\gamma$ be a
closed curve in $\Delta$. For any point not
on $\gamma$,
\begin{equation}
\label{complex-analysis-equation-cauchy-integral-formula}
n(\gamma,a)f(a)=\frac{1}{2\pi
i}\int_\gamma\frac{f(z)dz}{z-a}
\end{equation}
where $n(\gamma,a)$ is the index of $a$ with
respect to $\gamma$.
\end{theorem}

\begin{proof}
This is a simple application of Theorem
\ref{complex-analysis-theorem-cauchy-removable} for the
function
$$
F(z):=\frac{f(z)-f(a)}{z-a}
$$
Notice the limit condition holds, so that the
integral vanishes!
\end{proof}

\noindent
Notice that we can differentiate Cauchy
integral formula to obtain
\begin{equation}
\label{complex-analysis-equation-first-deriv-formu}
f'(z)=\frac{1}{2\pi
i}\int_\gamma\frac{f(\zeta)}{(\zeta-z)^2}
d\zeta
\end{equation}
and more generally
$$
f^{(n)}(z)=\frac{n!}{2\pi i}\int_\gamma
\frac{f(\zeta)}{(\zeta-z)^{n+1}}
$$
Though we can just say that integrals can be
differentiating by splitting into
real and imaginary part, and using the real
derivative, a technical lemma is
used by Ahlfors to confirm that indeed we can
differentiate
under the integral sign.

\begin{lemma}
\label{complex-analysis-lemma-differentiation-of-integral}
Suppose that $\varphi(\zeta)$ is continuous
on the arc $\gamma$.
Then the function
$$
F_n(z)=\int_\gamma
\frac{\varphi(\zeta)}{(\zeta-z)^n}d\zeta
$$
is analytic in each of the regions determined
by $\gamma$,
and its derivative is $F_n'(z)=nF_{n+1}(z)$.
\end{lemma}

\noindent
The importance of this result is that it
shows that a holomorphic function has
derivatives of all degrees.

Then we state two interesting results:

\begin{theorem}[Morera]
\label{complex-analysis-theorem-Morera}
\begin{reference}
\cite[p. 122]{ahl}
\end{reference}
If $f(z)$ is defined and continuous in a
region $\Omega$, and if
$\int_\gamma fdz=0$ for all closed curves
$\gamma$ in $\Omega$, then $f(z)$ is
analytic in $\Omega$.
\end{theorem}

\noindent
This is what Lee calls a {\it conservative
covector field}, i.e. a form whose
integral on closed curves vanishes. The proof
of Morera's theorem then reduces
to the fact that a covector field is
conservative if and only if it is exact
\cite[Theorem 11.42]{les}. The reverse
implication is the fundamental theorem of
calculus. The forward implication is not very
straightforward in \cite{les}.

As a final remark, I add that
(\cite[Propoistion 11.40]{les}) a smooth
covector
field is conservative if and only if its line
integrals are path-independent, in
the sense that integrals coincide along two
{\it piecewise} smooth curve
segments with the same starting and ending
points. That is, form is exact
implies integral independent of homotopy
class?

We finish with

\begin{theorem}[Liouville]
\label{complex-analysis-theorem-Liouville}
A bounded holomorphic function defined on all
of $\mathbb{C}$ must be constant.
\end{theorem}

\begin{proof}
Let $\gamma$ be a circle of radius $r$ about
$z$, then by the first derivative
of Cauchy formula (Eq.
\ref{complex-analysis-equation-first-deriv-formu}),
$$
|f'(z)|\leq \frac{M}{2\pi i}\int_\gamma
\frac{d\zeta}{|\zeta-z|^2}
$$
and compute the integral, which is
$\frac{2\pi i}{r}$.
\end{proof}

\noindent
A fun application of this is proving the
fundamental theorem of algebra. If a
complex polynomial had no zeroes, $1/P(z)$
would be analytic in all of
$\mathbb{C}$, and since $P(z)$ tends to
$\infty$ as $z$ tends to $\infty$,
$1/P(z)$ is bounded (use Riemann sphere
argument).

\section{Taylor series}
\label{complex-analysis-section-taylor}

\noindent
In this section I put
\href{https://en.wikipedia.org/wiki/%
Analyticity_of_holomorphic_functions}
{wiki's proof}
that a function is holomorphic if and only if
it is analytic.
This gives for free the Taylor expansion.
See Remark \ref{complex-analysis-remark-taylor-before} to
understand why this is here
and not elsewhere in this notes.

\medskip\noindent
Before proving the statement I need the
following tool.

We will use the Weierstrass M-test in the form
proved in Real Analysis,
Theorem~\ref{real-analysis-theorem-m-test}.

\begin{definition}
\label{complex-analysis-definition-analytic}
A function $f:\Omega \subset \mathbb{C} \to
\mathbb{C}$ is
{\it analytic} at $z_0 \in \Omega$ if in some
open disk centered at $z_0$
it can be expanded as a convergent power
series
$$
f(z)=\sum_{n=0}^\infty c_n(z-z_0)^n.
$$
\end{definition}

\begin{theorem}
\label{complex-analysis-theorem-holomorphic-iff-analytic}
A function $f:\Omega \subset \mathbb{C} \to
\mathbb{C}$ is holomorphic
in the sense of Definition
\ref{complex-analysis-definition-holomorphic-function}
if and only if it is analytic.
\end{theorem}

\begin{proof}
The backward implication is easy.
Suppose that $f(z)=\sum_{n=0}^\infty
c_n(z-z_0)^n$.
Then
$$
\begin{aligned}
\lim_{z \to z_0}\frac{f(z)-f(z_0)}{z-z_0} &=
\lim_{z\to z_0}\frac{\sum_{n=0}^\infty
c_n(z-z_0)^n-c_0}{z-z_0}\\
&=\lim_{z\to z_0}\frac{\sum_{n=1}^\infty
c_n(z-z_0)^n}{z-z_0}\\
&=\lim_{z\to z_0}\sum_{n=1}^\infty
c_n(z-z_0)^{n-1}=c_1.
\end{aligned}
$$

\noindent
The forward implication uses Cauchy integral
formula
\ref{complex-analysis-equation-cauchy-integral-formula} and
the geometric series.
Let us prove the geometric series, namely
that
$\lim_{n\to\infty} r^n=\frac{1}{1-r}$ for any
complex number $r$ of
norm strictly less than 1:
$$
\xymatrix@C=.2em@R=1em{
S_n & & 1
\ar@{}[r]&+&r\ar@{}[r]&+&r^2\ar@{}[r]&+&
\cdots\ar@{}[r]&+
&r^{n-2}\ar@{}[r]&+&r^{n-1}\\
rS_n & & & &
r\ar@{}[r]&+&r^2\ar@{}[r]&+&\cdots\ar@{}[r]&+
&r^{n-2}\ar@{}[r]&+&r^{n-1}\ar@{}[r]&+&r^n\\
S_n-rS_n& & 1 & & & & & & & & & & & - & r^n
}
$$
Thus $S_n(1-r)=1-r^n$, so that
$S_n=\frac{1-r^n}{1-r}$.
Taking limit as $n\to \infty$ gives the
result.

\medskip\noindent
To prove $f$ is analytic at $z_0\in \Omega$
and suppose that $\gamma$ is a circle
contained in $\Omega$
centered in $z_0$.
We need to construct a series expansion for
every $z$ inside the interior of $\gamma$.
By Cauchy integral formula
\ref{complex-analysis-equation-cauchy-integral-formula},
$$
\begin{aligned}
f(z)&=\frac{1}{2\pi\sqrt{-1}}\int_\gamma
\frac{f(\zeta)}{\zeta-z}d\zeta\\
&=\frac{1}{2\pi\sqrt{-1}}
\int_\gamma\frac{f(\zeta)}{(\zeta-z_0)-(z-
z_0)}d\zeta\\
&=\frac{1}{2\pi\sqrt{-1}}
\int_\gamma \frac{1}{\zeta-z_0}
\sum_{n=0}^\infty\left(\frac{z-z_0}{\zeta-
z_0}\right)^n
f(\zeta)d\zeta\\
&=\sum_{n=0}^\infty \frac{1}{2\pi\sqrt{-1}}
\int_\gamma\frac{(z-z_0)^n}{(\zeta-z_0)^{n+1}
}f(\zeta)d\zeta.
\end{aligned}
$$

\noindent
The critical step in the above computation
is interchanging the integral sign and the
series in the last step.
To make sure we can do this all we need to do
is check
that the last series converges.
This is done via Weierstrass M-test, i.e.
Theorem \ref{complex-analysis-theorem-m-test}.
We need to find uniform bounds for every term
in the series.
This is clear since $f$ is bounded on
$\gamma$ (by compactness of $\gamma$),
and the modulus of the quotient
$(z-z_0)/(\zeta-z_0)$ is strictly smaller
than $1$
(since $\zeta$ is farther from $z_0$ than
$z$), giving
for every $n$
$$
\begin{aligned}
\left| \frac{1}{2\pi\sqrt{-1}}
\int_\gamma\frac{(z-z_0)^n}{(\zeta-z_0)^{n+1}
}f(z)d\zeta\right|
&=\frac{1}{2\pi}\left|
\int_\gamma\left(\frac{z-z_0}{\zeta-z_0}
\right)^n\frac{f(z)}{\zeta-z_0}
d\zeta\right|\\
&< \frac{1}{2\pi}\int_\gamma |K|^n
\left|\frac{M}{r}\right|\\
&=|M|\;|K|^n.
\end{aligned}
$$

\noindent
where $M=\text{max}_{\zeta \in
\gamma}f(\zeta)$,
$\frac{z-z_0}{\zeta-z_0}=K<1$ and $r$ is the
radius of $\gamma$.
\end{proof}
\section{Riemann removable singularity
theorem}
\label{complex-analysis-section-removable-singularity}

\noindent
The upshot about the condition that
$\lim_{z \to z_0} (z-z_0)f(z)=0$ is that this
implies
that $\int_\gamma f=0$. But this actually
doesn't
work in the following proof, since we don't
get much from that
integral vanishing --- that's just to say
that
\cite[Chapter 4, Theorem 7]{Ahlfors} is
confusing.
Instead, I'll put here this great
\href{https://en.wikipedia.org/wiki/%
Removable_singularity}{Wikipedia})
proof:

\begin{theorem}[Riemann's removable
singularity theorem]
\label{complex-analysis-theorem-removable-singularity}
Let $f(z)$ be analytic in a region $\Omega'$
obtained by omitting a point $z_0$
from a region $\Omega$. There exists an
analytic function  defined on all
$\Omega$ that coincides with $f$ on $\Omega'$
if and only if
$\lim_{z \to z_0} (z-z_0)f(z)=0$. Such a
function is unique.
\end{theorem}

\begin{proof}
The direct implication is obvious since we
can replace $f$ by
its holomorphic extension in the limit, which
will converge to a fixed
value, so that the limit is zero. Uniqueness
also is trivial
by continuity of any holomorphic extension.

For the converse define
$$
h(z)=\begin{cases}
(z-z_0)^2f(z)\qquad &z \neq z_0 \\
0\qquad &z=z_0.
\end{cases}
$$
Clearly $h$ is holomorphic on
$\Omega\setminus\{z_0\}$,
and at $z_0$ it's also
complex-differentiable:
$$
h'(z_0)=\lim_{z \to z_0}
\frac{(z-z_0)^2f(z)-0}{z-z_0}
=\lim_{z\to z_0}(z-z_0)f(z)=0.
$$
by hypothesis. That is, $h$ is holomorphic
and thus has a Taylor series about $z_0$
(justify!
\href{https://en.wikipedia.org/wiki/%
Analyticity_of_holomorphic_functions}{wiki}
):
$$
h(z)=c_0+c_1(z-z_0)+c_2(z-z_0)^2+c_3(z-z_0)
^3+\ldots
$$

We have $c_0=h(z_0)=0$ and $c_1=h'(z_0)=0$,
and thus
$$
h(z)=c_2(z-z_0)^2+c_3(z-z_0)^3+\ldots.
$$
Hence, where $z\neq z_0$ we have
$$
f(z)=\frac{h(z)}{(z-z_0)^2}=c_2+c_3(z-z_0)+
\ldots.
$$
However
$$
g(z)=c_2+c_3(z-z_0)+\ldots
$$
is holomorphic on all $\Omega$, and is thus
an extension of $f$.
\end{proof}

\section{Zeroes of a holomorphic function}
\label{complex-analysis-section-zeroes}

\noindent
Perhaps the best way to remember this section
is by the formula
\begin{equation}
\label{complex-analysis-equation-zero}
f(z)=f_n(z)(z-z_0)^n,
\end{equation}
which holds for a holomorphic function with a
{\it zero of order $k$} on
$z_0$, meaning that $f$ and its derivatives
$f^{(k)}(z_0)$ vanish
for $k<n$ (see \cite[p. 126]{Ahlfors}).
Here $f_n$ is a holomorphic function which
does not vanish at $z_0$ .
This fact shows that the zero $z_0$ is
isolated,
since $f_n(z)$ is continuous and nonzero in
$z_0$,
and $(z-z_0)^n$ is also nonzero in a pointed
neighbourhood of $z_0$.

\medskip\noindent
Let's explain why Equation
\ref{complex-analysis-equation-zero} holds.
Recall the function that we used to prove
Cauchy integral formula
(Theorem
\ref{complex-analysis-theorem-cauchy-integral-formula}):
\begin{equation}
\label{complex-analysis-equation-F}
F(z)=\frac{f(z)-f(z_0)}{z-z_0},
\end{equation}
which has $\lim_{z\to z_0} (z-z_0)F(z)=0$.
Then we apply
Riemann's removable singularity theorem
\ref{complex-analysis-theorem-removable-singularity}
to find a holomorphic function $f_1$
extending $F$.
Taking the limit as $z \to z_0$, we see that
in fact
$f_1(z_0)=f'(z_0)$.
Substituting $F$ by $f_1$
in Equation \ref{complex-analysis-equation-F} we obtain
\begin{equation}
\label{complex-analysis-equation-zero2}
f(z)=f(z_0)+f_1(z)(z-z_0).
\end{equation}
If $f(z_0)$ and $f'(z_0)\neq 0$, then we say
that $f$ has
a zero of order 0 at $z_0$, which is one case
of
Equation \ref{complex-analysis-equation-zero}. If  $f'(z_0)$
also vanishes,
we repeat process for $f_1$, and for $f_2$ if
necessary, and so on,
until we arrive at Equation
\ref{complex-analysis-equation-zero}.

Let us now stop assuming that $z_0$ is a zero
of $f$
and just apply the argument to $f_1$. We
obtain
$$
f_1(z)=f_1(z_0)+f_2(z)(z-z_0)
$$
for some function $f_2$ such that
$f_2(z_0)=f''(z_0)$.
Substituing back in Equation
\ref{complex-analysis-equation-zero2} we obtain
$$
f(z)=f(z_0)+f_1(z)(z-z_0)+f_2(z)(z-z_0)^2.
$$

Repeating this process yields:

\begin{lemma}[Finite Taylor expansion]
\label{complex-analysis-lemma-finite-taylor-expansion}
\begin{reference}
\cite[Chapter 4, Section 3, Subsection 2,
Theorem 8]{ahl}
\end{reference}
If $f$ is analytic in a connected open set
$\Omega\subset\mathbb{C}$ and $z_0 \in
\Omega$,
then we can write
\begin{equation}
\label{complex-analysis-equation-finite-taylor-expansion}
f(z)=f(z_0)+f'(z_0)(z-z_0)+\frac{f''(z_0)}
{2!}(z-z_0)^2+\ldots
+\frac{f^{(n-1)}(z_0)}{(n-1)!}(z-z_0)^{n-1}+
f_n(z)(z-z_0)^n
\end{equation}
for some function $f_n$ analytic on $\Omega$.
\end{lemma}

\noindent
The coefficients $\frac{f^{(n)}(z_0)}{n!}$
may be found
by induction differentiating
Equation
\ref{complex-analysis-equation-finite-taylor-expansion} $n$
times
and evaluating at $z_0$:
any term with
factor $(z-z_0)$ will vanish at $z_0$, and we
are left with
$f^{(n)}(z_0)=n!f_n(z_0)$.

\medskip\noindent
\begin{remark}
\label{complex-analysis-remark-taylor-before}
Now I will briefly discuss \cite[Section 1.2,
Chapter 5]{Ahlfors}.
Since the proof presented by Ahlfors of the
Riemann removable singularity theorem
\ref{complex-analysis-theorem-removable-singularity}
did not convince me,
I went to
\href{https://en.wikipedia.org/wiki/%
Removable_singularity}{wiki}
and this forced me to do Taylor series before
Riemann removable singularity.
Of course if I ever understand Ahlfors' idea
I could put Taylor after
Riemann removable singularity, and after
zeroes of holomorphic functions
which seems much more natural
\end{remark}

\medskip\noindent
Actually this part is quite later in Ahlfors
book because it requires
the theory of convergence of series.
To pass from the finite Taylor expansion to
the
infinite version we notice that the
term $f_{n+1}$ in Equation
\ref{complex-analysis-equation-finite-taylor-expansion}
is of the form
$$
f_{n+1}(z)=\frac{1}{2\pi i}\int_\gamma
\frac{f(\zeta)d\zeta}
{(\zeta-z_0)^{n+1}(\zeta-z)}
$$
where $\gamma$ is a curve of radius $\rho$.
Then ``we obtain at once''
$$
|f_{n+1}(z)(z-z_0)^{n+1}|\leq
\frac{M|z-z_0|^{n+1}}{\rho^n(\rho-|z-z_0|}.
$$
Then we must observe that the bound converges
uniformly
to zero in every disk of radius smaller than
$\rho$,
and by Weierstrass theorem there should be a
limit
function identical to zero.

\medskip\noindent
Here's an important result from
\href{https://en.wikipedia.org/wiki/%
Analyticity_of_holomorphic_functions}{wiki}:
a complex-valued function is holomorphic if
and only if it
has a Taylor series expansion at every point
(raidus of convergence is finite).


\section{Maximum principle}
\label{complex-analysis-section-maximum-principle}

\begin{theorem}[Maximum principle]
\label{complex-analysis-theorem-maximum-principle}
\begin{reference}
\cite[Theorem 12, p.134]{ahl}
\end{reference}
If $f$ is analytic and non constant in a
region $\Omega$,
then its absolute value $|f|$ has no maximum
in $\Omega$.
\end{theorem}

Now look at \cite[p. 166]{ahl}:
``the maximum principle for analytic
functions
follows from the maximum principle for
harmonic functions
by applying the latter to $\text{log}|f|$
which is harmonic
when $f(z)\neq 0$''.

\begin{theorem}[Maxmimum principle for
harmonic functions]
\label{complex-analysis-theorem-maximum-principle-harmonic}
A nonconstant harmonic function has
neither a maximum nor a minimum in its region
of definition.
Consequently, the maximum and the minimum
on a closed bounded set $E$ are taken on the
boundary of $E$.
\end{theorem}

\noindent
Which basically says the same thing but now
for harmonic functions
and now also for the minimum.
``It applies also to the minimum
for the reason that $-u$ is harmonic together
with $u$''.


\section{Argument Principle}
\label{complex-analysis-section-argument-principle}

\noindent
A simple version of the Argument Principle
does not need Residue theorem:

\begin{exercise}
\label{complex-analysis-exercise-argument-principle}
Let $f$ be a holomorphic function on a disk,
non-zero on $\partial \Delta$, and
let
$S_k(f):=\frac{1}{2\pi\sqrt{-1}}
\int_{\partial\Delta}\frac{f'}{f}z^kdz$.
Prove the $S_k(f)=\sum d_i\alpha_i^k$ where
$\alpha_i$ are all zeroes of $f$ and
$d_i$ their multiplicities.
\end{exercise}

\begin{proof}
We use Equation \ref{complex-analysis-equation-zero} to write
$f(z)=(z-z_1)h_1(z)$ where $z_1$ is a zero of
$f$.
Applying that to $h_1$ for another zero
$z_2$ of $f$,
and continuing in this way, we obtain
\begin{equation}
\label{complex-analysis-equation-factor-zeroes}
f(z)=(z-z_1)(z-z_2)\ldots(z-z_n)h(z).
\end{equation}

\noindent
Computing the quotient we obtain
$$
\frac{f'(z)}{f(z)}=\frac{\text{derivative of
$(z-z_1)\ldots(z-z_n)$}}{(z-z_1)\ldots(z-z_n)
}+\frac{h'(z)}{h(z)}
$$
The right hand term will vanish upon
integration since it is a holomorphic
function (with no poles because $h(z)\neq
0$). The left-hand term will give a
sum of $\frac{1}{z-z_i}$ upon
differentiation. As it is, the result of
integration is the sum of the orders of each
zero i.e. the number of zeroes
counted without multiplicity
(because the integrals are all 1, and there's
one for each zero without
multiplicity).

Multiplying by $z^k$ yields the desired
result by Cauchy formula
\ref{complex-analysis-equation-cauchy-integral-formula}. More
exactly,
multiplying $z^k$ in the formula gives
instead of a sum
of integrals $\int \frac{1}{z-z_i}$, a sum of
integrals
$\int\frac{z^k}{z-z_i}$ which by Cauchy
formula give $z_i^k$
(disregarding the factor $2\pi\sqrt{-1}$).
\end{proof}

\begin{theorem}[Argument Principle]
\label{complex-analysis-theorem-argum-princ-rouch-theor}
\begin{reference}
\cite[Chapter 5, Theorem 18]{ahl}
\end{reference}
If $f$ is meromorphic in $\Omega$ with zeros
$a_j$ and poles $b_k$, then
\begin{equation}
\label{complex-analysis-equation-argument-principle}
\frac{1}{2\pi
i}\int_\gamma\frac{f'(\zeta)}{f(\zeta)}d\zeta
=\sum_{i}n(\gamma,a_i)-\sum_{k}n(\gamma,b_k)
\end{equation}
\end{theorem}

\begin{lemma}[Rouché theorem]
\label{complex-analysis-lemma-Rouche-theorem}
\begin{reference}
\cite[Chapter 5, Corollary, p. 153]{ahl}
\end{reference}
Let $\gamma$ be homologous to zero in
$\Omega$ and such that $n(\gamma,z)$ is
either 0 or 1 for any point $z$ not on
$\gamma$. Suppose that $f$ and $g$ are
analytic in $\Omega$ and satisfy that
$|f-g|<|f|$ on $\gamma$. Then $f$ and $g$
have the same number of zeros enclosed by
$\gamma$.
\end{lemma}

Compare with Misha's version

\begin{theorem}[Rouché theorem]
\label{complex-analysis-theorem-Rouche-theorem-Mishas-version}
Let $f_t$ be a family of holomorphic
functions on a disk  $\Delta$, continuously
depending on a parameter $t\in \mathbb{R}$
and non-zero everywhere on its
boundary $\partial\Delta$. Prove that the
number of zeros of $f_t$ in $\Delta$
is constant.
\end{theorem}

\begin{proof}
Consider the map $t\mapsto f_t\mapsto
S(f_t)$, which is continuous by hypothesis
and because $S$, being an integral, is
continuous. Then we obtain a continuous
map $\mathbb{R}\to\mathbb{R}$ with integer
values, meaning it must be constant.
\end{proof}

A similar argument may be used to prove that
a holomorphic function $F$ defined
on $\Delta\times\Delta$ gives a holomorphic
map
\begin{equation}
\label{complex-analysis-equation-zeros-on-polydisk}
y_0\mapsto
\int_{\partial\Delta}\frac{F'(x,y_0)}{F(x,
y_0)}\phi(x)dx
=\sum_id_i\phi(\alpha_i)
\end{equation}
for any holomorphic function
$\phi:\Delta\to\mathbb{C}$.

\section{Holomorphic functions in several
variables}
\label{complex-analysis-section-holom-funct-sever-varia}

\begin{lemma}
\label{complex-analysis-lemma-holom-funct-chara}
\cite{lec}, Theorem 1.21. Let
$U\subseteq\mathbb{C}^n$ be open and $f:U\to
\mathbb{C}$. The following are equivalent:
\begin{enumerate}
\item $f$ is holomorphic (i.e. it is
continuous and has a complex partial
derivative with respect to each variable at
each point of $U$)
\item $f$ is smooth and satisfies the
following Cauchy-Riemann equations:
\begin{equation}
\label{complex-analysis-equationcauchy-equat-riema-sever}
\frac{\partial u}{\partial
x^j}=\frac{\partial v}{\partial y^j},\qquad
\frac{\partial u}{\partial
y^j}=-\frac{\partial v}{\partial x^j}
\end{equation}
where $z^j=x^j+\sqrt{-1}y^j$ and
$f(s)=u(z)+\sqrt{-1}v(x)$.
\item For each $p=(p^1,\ldots,p^n)\in U$
there exists a neighbourhood of $p$ in
$U$ on which $f$ is equal to the sum of an
absolutely convergent power series of
the form
\begin{equation}
\label{complex-analysis-equation-taylo-serie-sever-varia}
f(z)=\sum_{i_1,\ldots,i_n}a_{i_1,\ldots,i_n}
(z^1-p^1)\ldots(z^n-p^n)
\end{equation}
\end{enumerate}
\end{lemma}

\begin{proof}
I will prove that if $f$ is holomorphic then
it has a Taylor series for $n=2$.
First apply Cauchy integral formula on each
variable to obtain
$$
f(z^1,z^2)=\frac{1}{(2\pi\sqrt{-1})^2}
\int_{\substack{|z^1-w^1|=r \\ |z^2-w^2|=r}}
\frac{f(w^1,w^2)}{(w^1-z^1)(w^2-z^2)}dw^1dw^2
$$
Now observe:
\begin{equation}
\label{complex-analysis-equation-multivariablecauchy}
\frac{1}{w^1-z^1}=\frac{1}{w^1-p^1+p^1-z^1}
=\frac{1}{w^1-p^1}\frac{1}{1-\frac{p^1-z^1}
{w^1-p^1}}
\end{equation}
And on the right-hand-side we have a
geometric series so that we may write
$$
\frac{1}{w^1-z^1}
=\frac{1}{w^1-p^1}\sum_{k=0}
^\infty\left(\frac{p^1-z^1}{w^1-p^1}\right)^k
$$
finally substituting this into
(\ref{complex-analysis-equation-multivariablecauchy}) we may
take
the products $(p^1-z^1)^{k_1}(p^2-z^2)^{k_2}$
out of the integral and define the
remaining term as $a_{k_1k_2}$.
\end{proof}

\section{Identity theorem in several
variables}
\label{complex-analysis-section-ident-theor-sever-varia}

\begin{theorem}[Identity theorem]
\label{complex-analysis-theorem-identity-several-variables}
\begin{reference}
\cite[Proposition 1.28]{lec}
\end{reference}
If two holomorphic functions
$f,g:\Omega\subset \mathbb{C}^n \to
\mathbb{C}$ coincide
in an open subset of the connected open set
$\Omega$,
then they coincide in all of $\Omega$.
\end{theorem}

\begin{proof}
Let $U$ be the set where the function
$h:=f-g$ and all its partial derivatives
vanish. Then $U$ is open since for every
point in $U$ there is a neighbourhood
where $h$ is expressed as a Taylor series,
whose coefficients must be given by
the partial derivatives of $h$. By definition
of $U$, we see that $h$ must be
zero in such a neighbourhood. $U$ is also
closed by continuity of partial
derivatives of all orders.
\end{proof}

\section{Hartog's phenomenon}
\label{complex-analysis-section-hartog}

\noindent
In particular, Hartog's phenomom says
that any isolated singularity is removable in
dimension $\geq 2$.

More precisely but still roughly,
it says that in dimension at least two,
a holomorphic function defined on an open set
minus a compact set
can be extended to the whole open set
as long as the open minus the compact is
connected.
This doesn't happen in one variable:
$\frac{1}{z}$
is holomorphic away from the origin, but
cannot be extended.

Hopefully we will see how this theorem is
used
in Kodaira embedding theorem and in showing
that
the sections of $\mathcal{O}(d)$ are
precisely degree $d$-polynomials.


\begin{theorem}[Hartog's phenomenon]
\label{complex-analysis-theorem-hartog}
\end{theorem}

\section{Germs of holomorphic functions}
\label{complex-analysis-section-germs-holom-funct}

\begin{definition}
\label{complex-analysis-definition-germ-holom-funct}
Let $U,U' \subset \mathbb{C}^n$ be
neighbourhoods of $0$ and $f \in
\mathcal{O}_U$, $f'\in\mathcal{O}_{U'}$
holomorphic functions. We say that $f$
and $f'$ have the same germ, $f\sim f'$, if
$f|_{U\cap U'}=f'|_{U\cap U'}$.
Clearly (?), $\sim$ gives an equivalence
relation on the set of pairs $(U\ni 0,
f\in \mathcal{O}_U)$. An equivalence class is
called {\bf germ of a holomorphic
function}. The space of germs in $0$ of
holomorphic functions on $\mathbb{C}^n$
is denoted $\mathcal{O}_n$.
\end{definition}

\begin{exercise}
\label{complex-analysis-exercise-stalk-not-finit-gener-over}
Prova that the ring $\mathcal{O}_n$ of germs
of holomorphic functions is not
finitely generated over $\mathbb{C}$ for any
$n>0$.
\end{exercise}

\begin{proof}
I think the existence of $e^z$ as a
holomorphic function satisfying the
differential equation $f'=f$ is enough to
show that the coefficients of its
Taylor polynomial are all nonzero. This
argument works for several variables as
well.
\end{proof}

\begin{definition}
\label{complex-analysis-definition-formal-power-series}
A {\it formal power series} in the variables
$t_1,\ldots,t_n$ is a sum
$\sum_{i=0}^{\infty}P_i(t_1,\ldots,t_n)$
where $P_i$ are homogeneous polynomials
of degree $i$. Addition of power series is
defined componentwise, and
multiplication is defined via
$$
\left(\sum_{i=0}^\infty
P_i(t_1,\ldots,t_n)\right)
\left(\sum_{i=0}^\infty
Q_i(t_1,\ldots,t_n)\right)
=\sum_{i=0}^\infty R_i(t_1,\ldots,t_n)
$$
where
$R_d(t_1,\ldots,t_n)=\sum_{i+j=d}P_i(t_1,
\ldots,t_n)Q_j(t_1,\ldots,t_n)$.
\end{definition}

We can think of germs of functions in
$\mathcal{O}_n$ as elements in the ring of
power series
$\mathbb{C}[\![t_1,\ldots,t_n]\!]$. I think
there is no problem to
prove this statement, nor the fact that power
series is actually a ring with
units the nonzero constants and zero the zero
constant.

\begin{exercise}
\label{complex-analysis-exercise-stalk-has-no-zero-divisors}
Prove that $\mathcal{O}_n$ has no zero
divisors.
\end{exercise}

\begin{proof}
Suppose that $PQ=0$ but neither of $P$ nor
$Q$ are zero. Then $P_i \neq 0$ for
some $i$ and $Q_j \neq 0$ for some $j$. Then
we can write
$$
P_iQ_j=-\sum_{\substack{p+q=i+j \\ p\neq
i,q\neq j}}P_pQ_q
$$
And then what. Other way is by induction. For
$n=0$ we are the complex numbers
so no problem. Suppose that $\mathcal{O}_n$
has no zero divisors. Then it looks
like we can deal with degrees smaller than
$n+1$, but the $d$-th term is not a
simple product of $P_iQ_j$ with $i+j=d$, but
a sum. So not sure too.
\end{proof}

\begin{definition}
\label{complex-analysis-definition-local-ring}
A ring $R$ is called {\it local} if it
contains an ideal $I\subset R$ such that
all elements $r \not \in I$ are invertible.
\end{definition}

It is an easy exercise to show that this
definition is equivalent to having a
unique maximal ideal.

\begin{exercise}
\label{complex-analysis-exercise-power-serie-not-finit-gene}
Prove that the ring $\mathbb{C}[\![
t_1,\ldots,t_n]\!]$ is not
finitely generated over $\mathcal{O}_n
\subset
\mathbb{C}[\![ t_1,\ldots,t_n]\!]$.
\end{exercise}

\begin{proof}
I thought that $\mathcal{O}_n$ would be the
same as
$\mathbb{C}[\![ t_1,..,t_n]\!]$…
the natural map is surely injective
but why not surjective? There are power
series that are not holomorphic
functions? Maybe because of radius of
convergence? No, because germs of
holomorphic functions can be defined very
near the origin… Every power series
has a nonzero radius of convergence, right?
\end{proof}

\noindent
Now I will discuss zeroes of holomorphic
functions of several variables.

\begin{definition}
\label{complex-analysis-definition-zero-holom-funct-sever}
Let  $f\in \mathcal{O}_n$ be a germ of
holomorphic function on $\mathbb{C}^n$.
Write its Taylor series
$f(z)=\sum_{i=0}^\infty P_i(t_1,\ldots,t_n)$,
where
$P_i$ are homogeneous polynomials of degree
$i$. We say that $f$ has a {\it
zero of order (or multiplicity) $k$ in $0$}
if $P_0=\ldots=P_{k-1}=0$. In this
situation {\it principal part} of the
function $f$ is the homogeneous polynomial
$P_k$.
\end{definition}

The following exercise shows that a
holomorphic change of coordinates will
preserve the order of a zero, and the
principal part of the new function will be
determined by the differential of the change
of coordinate map. We will need to
change coordinates when we do Weierstrass
Preparation theorem
\ref{complex-analysis-theorem-weierstrass-preparation}.

\begin{exercise}
\label{complex-analysis-exercise-princ-part-under-coord}
Let
$\Phi(t_1,\ldots,t_n)=(F_1(t_1,\ldots,t_n),
\ldots,F_n(t_1,\ldots,t_n))$
be
the holomorphic coordinate change around zero
(?), with $F_i(0,\ldots,0)=0$ and
$A:=\left(\frac{\partial F_i}{\partial
t_j}\right)_0$ its differential (I suppose
$\det A \neq 0$. Prove that
\begin{enumerate}
\item For any germ $f\in \mathcal{O}_n$ which
has multiplicity $k$, the function
$\Phi^*(f)$ has zero of the same
multiplicity.
\item The principal part of $\Phi^*(f)$ is
obtained from the principal part of
$f$ by action of $A$.
\end{enumerate}
\end{exercise}

\begin{proof}[Sketch of proof]
\begin{enumerate}
\item The condition that $\det A \neq 0$
implies that all $F_i$ must have linear
term---otherwise their partial derivatives
would vanish when evaluated at zero.
When substituting
$f(z_1,\ldots,z_n)=\sum_{|\alpha|\geq
k}\alpha z^\alpha$ with
$\Phi$ we find that there must be a term of
order $k$.
\item The case for $n=1$ is clear. The
general case follows, I think, from the
observation that the derivative $A$ recovers
the linear terms, which, as shown
in the previous item, correspond to the
principal part of $\Phi^*f$.
\end{enumerate}
\end{proof}

\begin{exercise}
\label{complex-analysis-exercise-zero-locus-homog-polyn}
Let $Q$ be a non-zero homogeneous polynomial
on $t_0,\ldots,t_n$, and $V(Q)$ its
zero set, which we consider as a subset in
$\mathbb{C}P^{n}$.
\begin{enumerate}
\item Prove that $\mathbb{C}P^{n}\setminus
V(Q)$ is non-empty.
\item Prove that $V(Q)\subset\mathbb{C}P^{n}$
is a set of measure 0.
\end{enumerate}
\end{exercise}

\begin{proof}
\begin{enumerate}
\item This follows from Identity Theorem
\ref{complex-analysis-theorem-identity-several-variables}.
Indeed, if $V(Q)$ was all of
$\mathbb{C}^{n}$, it would vanish on an open
set, implying that $Q$ is
identically zero.
\item $V(Q)$ may be decomposed in the sets of
regular and singular points.
Regular points have submanifold charts, while
singular points have measure zero
by Sard's theorem.
\end{enumerate}
\end{proof}

\begin{exercise}
\label{complex-analysis-exercise-principal-parts}
Let $f_1,f_2,\ldots\in \mathcal{O}_n$ be a
countable collection of germs, which
vanish with multiplicity $k_1,k_2,\ldots$.
Prove that there exists a coordinate
system $z_1,\ldots,z_n$ such that
$\lim_{z_n\to 0}
\frac{f_i(0,z_n)}{z_n^{k_i}}\neq 0$ for all
$i$.
\end{exercise}

\begin{proof}[Sketch of proof]
I don't understand this exercise: evaluating
any germ $f_i$ on $(0,z_n)$ will
make all terms that have any variable other
than $z_n$ vanish. Thus the first
term will be a homogeneous polynomial in
$z_n$, which is the principal part of
$f_i$ evaluated in $(0,z_n)$. But of course
taking quotient by $z_n$ all terms
with powers of $z_n$ higher than 1 will
vanish, barely leaving the principal
part of $f_i$ evaluated at $(0,1)$. But this
works regardless of the coordinate
system.
\end{proof}

\section{Elementary symmetric polynomials and
Newton formula}
\label{complex-analysis-section-eleme-symme-polyn-newto}

\noindent
The main result in this section
is to show that the elementary symmetric
polynomials
can be given in terms of the Newton
polynomials.
More exactly, as elements of the polynomial
ring with rational coefficients and Newton
polynomials
as indeterminates.
In turn, this says that the elementary
symmetric polynomials
are Weierstrass polynomials as long as the
Newton polynomials
are holomorphic (which is easy to prove using
the Argument Principle,
Exercise \ref{complex-analysis-exercise-argument-principle}).
The elementary symmetric polynomials, in the
form of a product,
are exactly what we get when we factor the
zeroes
of a holomorphic function as in Equation
\ref{complex-analysis-equation-factor-zeroes}.
So essentially this paragraph says the proof
of
Weierstrass Preparation Theorem
\ref{complex-analysis-theorem-weier-prepa-theor}.

\medskip\noindent
The $\alpha_i$ will eventually play the roles
of the zeroes
of the holomorphic function, but in this
section they are just
indeterminates in the ring
$\mathbb{Z}[\alpha_1,\ldots,\alpha_n]$.
We start by defining three types of
polynomials in this ring.

The definition of symmetric polynomial
will not be used in the following,
but we leave include it anyways:

\begin{definition}
\label{complex-analysis-definition-symmetric-polynomial}
{\it Symmetric polynomial} 
$P(z_1,…,z_n) \in \mathbb{C}[z_1,…,z_n]$
is a polynomial which is invariant
with respect to the symmetric group
$\Sigma_n$ acting on $\mathbb{C}[z_1,…,z_n]$
on the usual way.
\end{definition}

\begin{definition}
\label{complex-analysis-definition-elemen-symmetric-poly}
Let $e_i \in
\mathbb{Z}[\alpha_1,\ldots,\alpha_n]$ be the
coefficients of the
polynomial
\begin{equation}
\label{complex-analysis-equation-eleme-symme-polyn}
t^n+e_1t^{n-1}+\ldots+e_{n-1}t+e_n:=\prod_{i=
1}^n(t+\alpha_i).
\end{equation}
Then $e_i$ are called {\it elementary
symmetric polynomials} on $\alpha_i$.
\end{definition}

\noindent
Notice that the polynomial in Equation
\ref{complex-analysis-equation-eleme-symme-polyn} is an
element of
the polynomial ring
on the indeterminates
$t,\alpha_1,\ldots,\alpha_n$.

There are explicit descriptions
for the elementary symmetric polynomials.
In fact, $e_1=\sum_i\alpha_i$
(which is the first Newton polynomial $p_1$
defined below).
Also $e_n=\prod_{i}\alpha_i$, and in general
$e_k=\sum_{1\leq i_1<\ldots<i_k\leq
n}\alpha_{i_1}\ldots\alpha_{i_k}$.

\begin{definition}
\label{complex-analysis-definition-newton-polynomial}
A {\it Newton polynomial} is
$p_j:=\sum_{i=1}^n\alpha_i^j$.
\end{definition}

\begin{definition}
\label{complex-analysis-definition-compl-homog-symme-polyn}
A {\it complete homogeneous symmetric
polynomial} of degree $k$ is $h_k$
obtained as a sum of all homogeneous
monomials of degree $k$, that is,
$\alpha_1^k+\ldots+\alpha_n^k+\alpha_1^{k-1}
\alpha_2+\ldots$.
\end{definition}

\noindent
Upshot: a generating function is a
clothesline on which we hang up a sequence of
numbers
for display
(\cite[p. 1]{generatingfunctionology})

\noindent
Corresponding to the above definitions we
have the generating functions
$E(t):=\sum_{i=0}^ne_it^i$,
$P(t):=\sum_{i=1}^\infty p_it^i$ and
$H(t):=\sum_{i=0}^\infty h_it^i$ which are
formal series in
$\mathbb{Z}[\alpha_1,\ldots,\alpha_n][\![t]
\!]$.

\begin{exercise}
\label{complex-analysis-exercise-H}
Prove that
$H(t)=\prod_{i=1}^n\frac{1}{1-t\alpha_i}$.
\end{exercise}

\begin{proof}
Let us write (possibly as a formal
definition)
$$
\prod_{i=1}^n\frac{1}{1-t\alpha_i}=\prod_{i=
1}^n\sum_{k=0}^\infty\alpha_i^kt^k.
$$

We also write
$f_i\overset{\text{ops}}{\longleftrightarrow}
\{\alpha_i^k\}_{k=1}^\infty$
to mean
that $f_i$ is the power series associated to
the sequence
$\{\alpha_i^k\}_{k=1}^\infty$. Then the
equation above is the product of the
$f_i$. Then all we have to prove is that
$$
\prod_{i=1}^nf_i\overset{\text{ops}}
{\longleftrightarrow}
\left\{\sum_{i_1+\ldots+i_n=k}\alpha_1^{i_1}
\ldots\alpha_n^{i_n}\right\}
_{k=1}^\infty=\{h_k\}_{k=0}^\infty
$$
This is just a generalization of the formula
for product of power series for a
product of $n$ power series.
\end{proof}

\begin{exercise}
\label{complex-analysis-exercise-E}
Prove that $E(t)=\prod_{i=1}^n(1+t\alpha_i)$.
\end{exercise}

\begin{proof}
This is just writing
$\prod_{i=1}^n(1+t\alpha_i)
=\prod_{i=1}^nt(\frac{1}{t}+\alpha_i)$ and
continue until we get to $E(t)$.
\end{proof}

\noindent
It follows from the two previous exercises
that $H(t)E(-t)=1$. Using Exercise
\ref{complex-analysis-exercise-E} and applying logarithm we
obtain that
$\frac{E'(-t)}{E(-t)}=-\sum_{i=1}^n
\frac{\alpha_i}{1-t\alpha_i}$. Expanding
this formula as geometric power series we
obtain that
\begin{equation}
\label{complex-analysis-equation-P-in-terms-of-E}
P(t)=-\frac{E'(-t)}{E(-t)}
\end{equation}

\begin{exercise}
\label{complex-analysis-exercise-p-polynomials-of-e}
Prove that $p_i$ can be expressed as
polynomials of $e_i$ (with integer
coefficients).
\end{exercise}

\begin{proof}
Using Eq. \ref{complex-analysis-equation-P-in-terms-of-E}, and
$H(t)E(-t)=1$, we have
$$
P(t)=E'(-t)H(t).
$$
Expanding the power series we obtain that the
$k$-th term is
$$
p_k=\sum_{i=0}^k(-1)^{i+1}ie_ih_{k-i}
$$
\end{proof}

\begin{exercise}
\label{complex-analysis-exercise-h-and-e}
Prove that $h_i$ can be expressed as
polynomials of $e_i$ with integer
coefficients. Prove that $e_i$ can be
expressed as polynomials of $h_i$ with
integer coefficients.
\end{exercise}

\begin{proof}[Sketch of proof]
By $H(t)E(-t)=1$ we see that
$\frac{E'(-t)}{E(-t)}=\frac{H(t)}{H'(t)}$.
Both
denominators are expressed as power series in
$\alpha_i$. Multiplying by $E(-t)$
as in Exercise \ref{complex-analysis-exercise-E}
would let $E'(-t)$ be expressed as a power
series in $h_i$ and
$\alpha_i$, while multiplying by $H(t)$ as in
Exercise \ref{complex-analysis-exercise-H} would
let $H'(t)$ expressed in terms of
$e_i$ and $\alpha_i$.
\end{proof}

\begin{exercise}[Newton formula]
\label{complex-analysis-exercise-newton-formula}
Prove that
$ke_k=\sum_{i=1}^{k-1}(-1)^ie_{k-i}p_i$.
\end{exercise}

\begin{proof}
Note that
$$
P(t)E(-t)=\left(\sum_{k=0}^{\infty}
p_it^i\right)
\left(\sum_{i=0}^n(-1)^ie_it^i\right)
=\sum_{k=0}^\infty \sum_{i=0}^k
(-1)^{k-i}e_{k-i}p_it^k
$$
using the product formula of power series,
where we define $E(-t)$ as a power
series by letting $e_i=0$ for $i>n$. By Eq.
\ref{complex-analysis-equation-P-in-terms-of-E}, this
equals $E'(-t)$, which we may also see as a
power series. Comparing the $k$-th
term yields the result modulo a minus sign.
\end{proof}

\noindent
Finally,

\begin{exercise}
\label{complex-analysis-exercise-symme-polyn-terms-newto}
Prove that $e_i$ are expressed as polynomials
on $p_i$ with rational
coefficients.
\end{exercise}

\begin{proof}

\end{proof}

\section{Weierstrass preparation theorem}
\label{complex-analysis-section-weier-prepa-theor}

\noindent
The upshot about Weierstrass preparation
theorem is that
we would like it if $\mathcal{O}_n$ was an
Euclidean domain,
i.e. that there was a division algorithm,
as in Commutative Algebra Definition
\ref{commutative-algebra-definition-%
euclidean-domain}.
Probably that's not true because
holomorphic functions are infinite series,
so we cannot put an Euclidean function like
the degree.
Instead, we have Weierstrass preparation
theorem
(which gives a Weierstrass {\it division}
theorem…?)

\medskip\noindent
Before starting let's say the proof.
Start with a germ of holomorphic function
$f$.
Split the domain as $(z',z_n)$ and regard $f$
as a function of $z_n$.
Then factor its zeroes
as in Equation \ref{complex-analysis-equation-factor-zeroes}.
This gives, for every $z'$,
a product of $z_n$ minus the zeroes times a
holomorphic function
without zeros, namely $u(z',z_n)\prod
(z_n-w_i(z'))$.
But hey, that product is a polynomial
function in the ring
$\mathbb{C}[z_n,w_1(z'),\ldots,w_n(z')]$;
i.e.
we shall be done once prove it is in fact a
Weierstrass polynomial!
By definition, its coefficients are the
elementary symmetric polynomials on the
$w_i(z')$.
But we showed that these can be expressed as
polynomials
with indeterminates being the Newton
polynomials.
And hey, the Newton polynomials are
holomorphic functions
by the Argument principle exercise
\ref{complex-analysis-exercise-argument-principle}.
That says that that product of $z_n$ minus
the roots
is a Weierstrass polynomial, as desired.


\begin{definition}
\label{complex-analysis-definition-weierstrass-polynomial}
A {\it Weierstrass polynomial} is a function
$f \in \mathcal{O}_{n-1}[z_n]$,
that is, a function which is polynomial in
the last variables with coefficients
that are analytic functions on the first
$n-1$ variables.
\end{definition}

\noindent
The following formulation is found in
\cite{Demailly}:

\begin{theorem}[Weierstrass preparation
theorem]
\label{complex-analysis-theorem-weier-prepa-theor}
Let $f$ be holomorphic on a neighourhood of
$0$ in $\mathbb{C}^n$,
such that $f(0,z_n)/z_n^s$ has a nonzero
finite limit at $z_n=0$.
With the ``above'' choice of $r'$ and $r_n$,
one can write
$$
f(z)=u(z)P(z',z_n)
$$
where $u$ is an invertible holomorphic
function
in a neighbourhood of the polydisk
$\overline{\Delta}(r',r_n)$
and $P$ is a Weierstrass polynomial with
coefficient
functions defined for $|z'|\leq r'$ in
$\mathbb{C}^{n-1}$.
\end{theorem}

\begin{proof}
Step 0. First notice that since I will use
this theorem
to prove that $\mathcal{O}_n$ is a UFD, I
will start with
an arbitrary function $f \in \mathcal{O}_n$.
(Right now it's not clear why we need that
$f(0)=0$.)
But then it's straightforward to change
coordinates so that $f$
is not identically zero on the  $z_n$-axis:
just complete the point where $f$ is nonzero
to a basis $\{b_i\}_{i=1}^{n-1}$
and precompose with a linear map sending to
$(0,\ldots,0,1)$
to that point, and the rest of the canonical
basis
to the corresponding basis vectors $b_i$.
The composition of this linear map with $f$
remains holomorphic
and is nonvanishing in the $z_n$-axis.

(See Exercise \ref{complex-analysis-exercise-principal-parts}
for a proof that we can suppose
that the coordinate change preserves the
order of the zero,
but again it's not clear why I need that.)

\medskip\noindent
Step 1. Now we define the correct domain.
This is necessary for the next step.

Then $f(0,z_n)$ is a holomorphic function
with a zero on $0$. By
the Taylor expansion we know that this zero
must be isolated, so that there is a
number $r_n>0$ such that $f(0,z_n)\neq 0$ for
$0<|z_n|\leq r_n$.

Since the circle
$\partial\Delta_0(r_n)=\{(0,z_n):|z_n|=r_n\}$
is compact
and $f(0,z_n)\neq 0$ on this circle, we can
construct
an open neighbourhood of the circle where $f$
is nonzero as follows.
For any point of the circle $|z_n|=r_n$
pick a pointed open neighbourhood where $f$
is nonzero (using continuity).
This gives a collection of open
neighbourhoods of $\mathbb{C}^n$
covering the circle
$\partial\Delta_{0}(r_n)$,
from which we take a finite subcover.
Let $r'$ be the minimum radius of these
balls.
We have that $f(z',z_n)\neq 0$ for $z'<r'$
and  $|z_n|<r_n+\varepsilon$
for some small $\varepsilon$.

This means that in the definition of $S_k$ in
the next
step we are confident that we can integrate
the function
$\frac{1}{f(z',z_n)}$ along the circle
$\partial\Delta_{z'}r_n$
for every $z'$.

\medskip\noindent
Step 2. Now we apply Argument principle.
Consider the function
$$
S_k(z'):=\frac{1}{2\pi\sqrt{-1}}
\int_{\partial\Delta(r_n)}
\frac{f'(z',z_n)}{f(z',z_n)}z_n^kdz_n.
$$
In Exercise \ref{complex-analysis-exercise-argument-principle}
we showed that $S_k(z')$ is the sum of the
$k$-th powers of
the zeroes of  $f(z',-)$ without
multiplicities.
In particular
$S_0(z')$ is an integer: the number of
zeroes of $f(z',-)$ on $\Delta_{z'}(r_n)$
(for fixed
$z' \in \Delta(r')=\{z' \in
\mathbb{C}^{n-1}:|z'|<r'\}$).

The map $S_k$ is holomorphic as a function of
$z'$ because
it is holomorphic in each of its $n-1$
complex variables.
This can be checked by fixing all but one of
such variables
and splitting the expression in real and
imaginary parts.
By interchanging partial derivatives and
integrals as real functions
we conclude that Cauchy-Riemann equations are
satisfied
by holomorphicity of the integrands.
The proof for differentiation of integrals of
real functions may be found
\href{https://en.wikipedia.org/wiki/%
Leibniz_integral_rule#Proof_of_basic_form}
{here}.

Moreover, since
$S_0(z')$ is an integer,
it's constant (as a function of $z'$) by
continuity.
Thus we have exactly $s$ functions
$w_1(z'),\ldots,w_s(z')$ which give the $s$
zeroes of $f(z',-)$
at every $z' \in \Delta(r')$.
Now we apply argument principle with
multiplication by
$z^k$ to obtain $S_k(z')=\sum w_i(z')^k$
(without multiplicities).
Again, these maps are holomorphic on $z'$ by
Theorem
\ref{complex-analysis-theorem-holomorphic-iff-analytic} since
they are defined as integrals
and have series expansions.
Notice these are by definition the Newton
polynomials on $w_i(z')$.

\medskip\noindent
Step 3.
By the Newton formula (Exercise
\ref{complex-analysis-exercise-newton-formula})
we know that the elementary symmetric
polynomials $e_i$ can be expressed
as polynomials of the Newton polynomials
$p_i$.
We conclude that the $e_i$ are also
holomorphic functions of $z'$.
This shows that
$$
\begin{aligned}
P(z',z_n):&=\prod_{i=1}^k(z_n-w_i(z'))\\
&=t^n+e_1t^{n-1}+\ldots+e_{n-1}t+e_n,
\end{aligned}
$$

\noindent
where the $e_i$ are the elementary symmetric
polynomials,
defined exactly by the above relation (see
Equation \ref{complex-analysis-equation-eleme-symme-polyn}),
is a Weiestrass polynomial
(because its coefficients are holomorphic
functions).

Factoring the zeroes of $f(z',-)$ as in
Equation \ref{complex-analysis-equation-factor-zeroes}
we obtain
$$
f(z',z_n)=\left(\prod_{i}z_n-w_i(z')\right)
h(z',z_n)
$$
for a function $h$ that doesn't vanish in
$\Delta(r')\times \Delta_0(r_n+\varepsilon)$.

To conclude we need to confirm that $h$ is
holomorphic.
By our factorization we know that for every
fixed $z'$
$h(z',-)$ is holomorphic in $z_n$.
Then we can use Cauchy integral formula
\ref{complex-analysis-equation-cauchy-integral-formula}
to get
$$
h(z',z_n)=\frac{1}{2\pi\sqrt{-1}}
\int_{\partial\Delta(r_n)}\frac{f(z',w_n)}
{w_n-z_n}dw_n.
$$
This is a holomorphic function of $z'$ since
for any $z_0' \in \Delta(r')$ the limit
$$
\lim_{z'\to z_0'}
\frac{
\int_{\partial\Delta_{z'}(r_n)}
\frac{f(z',w_n)}{w_n-z_n}wd_n
-
\int_{\Delta_{z_0'(r_n)}}
\frac{f(z_0',w_n)}{w_n-z_n}dw_n}
{z'-z_0}
$$
exists.


% The classical way is to notice that the
% function $f/P$ is holomorphic in $z_n$
% because $f$ and $P$ have the same zeroes,
% and we write the integral formula to
% complete to a holomorphic function (but I
% think it should be enough with my
% argument).
\end{proof}

\noindent
Recall from Commutative Algebra Lemma
\ref{commutative-algebra-lemma-gauss}
that if $R$ is a UFD, then $R[x]$ is a UFD.

\begin{lemma}
\label{complex-analysis-lemma-stalk-is-ufd}
The stalk
$\mathcal{O}_n:=\mathcal{O}_{\mathbb{C}^n,0}$
is a UFD.
\end{lemma}

\begin{proof}
By induction on $n$. For $n=0$ it is trivial.
Suppose $\mathcal{O}_{n-1}$ is a
UFD. Then by Gauss' Lemma
\ref{commutative-algebra-lemma-gauss},
$\mathcal{O}_{n-1}[w]$ is a UFD
too. Thus we may express any Weierstrass
polynomial $g$ as a product of
irreducible elements (uniquely up to
multiplication by units).

Let $f\in \mathcal{O}_n$. We want to express
$f$ as a product of (unique up to
multiplication by units) of irreducible
elements. By Weierstrass Preparation
Theorem \ref{complex-analysis-theorem-weierstrass-preparation}
there is a Weierstrass polynomial
$g\in\mathcal{O}_n[w]$ and a holomorphic
function not vanishing on $0$ (i.e. a
unit of $\mathcal{O}_n$) such that $f=gh$. By
the previous remark $g$ is
factored uniquely up to multiplication by
units as $g=g_1\ldots g_m$. This shows
existence of the factorization.

To prove uniqueness suppose that $f=f_1\ldots
f_k$ for some irreducible
$f_1,\ldots,f_k\in\mathcal{O}_n$. Since $f$
does not vanish in the $w$ axis,
neither can each $f_i$, so that we may
decompose each of them as  $f_i=g_i'h_i$
by Weierstrass Preparation Theorem. Since
$f_i$ is irreducible, it follows that
$g_i'$ is irreducible. Then we have that $$
f=gh=\prod g_i'\prod h_i $$ so by
uniqueness in Weierstrass Preparation Theorem
we conclude that $g=\prod g_i'$,
and by uniqueness from the fact that
$\mathcal{O}_n[w]$ is a UFD we conclude
that $g$ coincides with $\prod g_i'$ up to
multiplication by units.
\end{proof}

\medskip\noindent
We finish with an exercise from Misha's
handouts
that helped on the way.

We denote by $B_r(z_1,\ldots,z_{n-1})$ the
ball of radius $r$ in
$\mathbb{C}^{n-1}\subset\mathbb{C}^n$.

\begin{exercise}
\label{complex-analysis-exercise-polydisks}
Let $F$ be a holomorphic function on a
neighbourhood of $0$ in $\mathbb{C}^n$,
such that $\lim_{z_n\to 0}
\frac{F(0,z_n)}{z_n^k}\neq 0,\infty$.
Consider the projection map
$\Pi:\mathbb{C}^n\to\mathbb{C}^{n-1}$
$(z_1,\ldots,z_n)\mapsto
(z_1,\ldots,z_{n-1})$.
\begin{enumerate}
\item Prove that for an appropriate pair
$r,r'$, the resitriction of $F$ to the
polydisk
$\Delta(n-1,1):=B_r(z_1,\ldots,z_{n-1})
\times\Delta_{r'}(z_n)$
nowhere
vanishes on the set
$\Pi^{-1}(\partial\Delta_{r'}(z_n))$.
\item Prove that in this case the
resitriction of $F$ to this polydisk has
precisely $k$ zeroes
$\alpha_1,\ldots,\alpha_k$ on each fiber of
$\Pi$.
\item Prove that $\sum_{i=1}^k \alpha_i^d$ is
a holomorphic function on
$B_r(z_1,\ldots,z_{n-1})$.
\item Prove that any elementary symmetric
polynomial on $\alpha_i$ gives a
holomorphic function on
$B_r(z_1,\ldots,z_{n-1})$.
\end{enumerate}
\end{exercise}

\begin{proof}
\begin{enumerate}
\item Suppose that for every $(r,r')$ there
are points
$z_1\in B_r(z_1,\ldots,z_{n-1})$ and $z_n\in
\partial\Delta_{r'}(z_n)$
where $F$ vanishes. Then we obtain a
convergent sequence where $F$
vanishes, and by Identity Principle
\ref{complex-analysis-theorem-identity-several-variables} we
conclude that $F$ mus be identically zero.
\item This is just applying Argument
principle, i.e. the integral
$\frac{1}{2\pi\sqrt{-1}}\int_{\partial
\Delta}\frac{F'(y_0,z)}{F(y_0,z)}dz$
equals the number of zeroes. Since it is a
continuous (holomorphic?) function on
$\Delta_y$ and integer valued, it must be
constant.

The condition that
$\lim_{z_n\to0}\frac{F(0,z_n)}{z_n^k}\neq
0,\infty$ means
that the fiber at 0 has a zero of order $k$:
$$
\frac{F(0,z_n)}{z_n^k}=\frac{F(0)}{z_n^k}+
F'(0)\frac{z_n}{z_n^k}+\ldots
+\frac{F^{(k)}(0)}{k!}\frac{z_n^k}{z_n^k}+
\ldots
$$
Thus, on other fibers we can have
more zeroes, but without multiplicities they
are always $k$.
\item
\end{enumerate}
\end{proof}

\begin{exercise}[Weierstrass preparation
theorem]
\label{complex-analysis-exercise-weier-prepa-theor}
Let $F$ be an analytic function in a
neighbourhood of $0$ in $\mathbb{C}^n$,
such that $\lim_{z_n\to 0}\neq 0,\infty$.
Consider the projection map
$\Pi:\mathbb{C}^n\to\mathbb{C}^{n-1}$,
$(z_1,\ldots,z_n)\mapsto(z_1,\ldots,z_{n-1})
$,
and let
$P(z_n)\in \mathcal{O}_{n_-1}[z_n]$ be the
Weierstrass polynomial given by
$P(z_n)=\sum_{i=0}^ke_iz_n^i$, where $e_i$
are the elementary symmetric
polynomials on the zeros
$\alpha_1,\ldots,\alpha_k$ defined in the
previous
exercise. Prove that $F=P(z_n)u$, where $u$
is a germ of an invertible
holomorphic function.
\end{exercise}

\begin{proof}
For every $y_0$ we can write
$$
F(y_0,z_n)=(z_n-\alpha_1)\ldots(z_n-\alpha_n)
h(z_n)
$$
where implicitly all $\alpha_i$ and $h$
depend on $y_0$. The product of
$(z_n-\alpha_i)$ is the definition of $P$.
Varying $y_0$ we obtain the result.
\end{proof}

\section{Weierstrass division theorem}
\label{complex-analysis-section-weierstrass-division-theorem}

\begin{theorem}[Weierstrass division theorem]
\label{complex-analysis-theorem-weierstrass-division-theorem}
For any $f \in \mathcal{O}_{\mathbb{C}^n,0}$
and
a Weierstrass polynomial $g \in
\mathcal{O}_{\mathbb{C}^{n-1},0}[z_n]$
there exist a unique Weierstrass polynomial
$r \in \mathcal{O}_{\mathbb{C}^{n-1},0}[z_n]$
and a unique germ $h \in
\mathcal{O}_{\mathbb{C}^n,0}$
such that
$$
f=gh+r,\qquad \text{deg }r < \text{deg }g.
$$
\end{theorem}

\begin{proposition}
\label{complex-analysis-proposition-stalk-is-noetherian}
$\mathcal{O}_{\mathbb{C}^n,0}$ is Noetherian.
\end{proposition}

\begin{proof}
Let $I \subset \mathcal{O}_{\mathbb{C}^n,0}$
be any ideal.
We want to show that $I$ is finitely
generated.

First notice that the ring of Weierstrass
polynomials is
Noetherian by induction on Hilbert basis
theorem:
$\mathcal{O}_{\mathbb{C}^0,0}[z_1]=\mathbb{C}
[z_1]$
(because $\mathbb{C}$-valued functions on
$\mathbb{C}^0=pt$
are just constants, i.e. $\mathbb{C}$)
is Noetherian because $\mathbb{C}$ is,
and so on.

This means that $I \cap
\mathcal{O}_{\mathbb{C}^{n-1},0}[z_n]$
is finitely generated, say $I=(g_1,…,g_k)$
for $g_i$ Weierstrass polynomials.

Now pick any Weierstrass polynomial
$g \in I \cap
\mathcal{O}_{\mathbb{C}^{n-1},0}[z_n]$,
which is granted to exist by Weierstrass
Preparation theorem
applied to any $f \in \mathcal{O}_{C_n,0}$:
we must have $f=gh$ for $g \in
\mathcal{O}_{\mathbb{C}^{n-1},0}[z_n]$
and $h$ unit (i.e. holomorphic germ not
vanishing at $0$),
which means that $g=fh^{-1}$ is in $I$,
and pick also any element $\tilde{f} \in I$.
Apply Wierestrass Division theorem to obtain
$$
\tilde{f}=gh+r,\qquad \text{deg }r <
\text{deg }g.
$$
Since $g$ and $\tilde{f}$ are in $I$, so is
$r$.
Moreover, $r$ is a Weierstrass polynomial and
is thus generated
by the $g_i$. Since $g$ does not depend on
$\tilde{f}$,
we see that $I=(g,g_1,…,g_k)$.
\end{proof}

\medskip\noindent
Not knowing exactly where to put this, here
goes:

\medskip\noindent
Now we discuss the analytic version of
Hilbert's Nullstellensatz.
as explained by Demailly.

\noindent
Usually I think of Hilbert's Nullstellensatz
as the statement that $I(V(I))=\sqrt{I}$.
Following are two simple initial
observations.

It is immediate that
$\sqrt{\mathfrak{p}}=\mathfrak{p}$
if $\mathfrak{p}$ is a prime ideal (of any
ring):
if  $f^k \in \mathfrak{p}$ then
either $f \in \mathfrak{p}$ or $f^{k-1}\in
\mathfrak{p}$,
etc.

It is also easy to see that
for any ideal $I$ of functions or germs of
functions
(i.e. either  $I \subset \mathcal{O}_X(U)$
or $I \subset \mathcal{O}_n$)
we have
$\sqrt{I}\subseteq I(V(I))$.
Indeed, if $f^k \in I$,
then $f^k(x)=0$ for all $x \in V(I)$,
and since $\mathbb{C}$ has no nonzero
divisors
we see that $f(x)=0$,
that is $f \in I(V(I))$.

It remains to show that $I(V(I)) \subseteq
\sqrt{I}$.
We shall first suppose that  $I$ is prime
and prove that $I(V(I)) \subseteq I$.
(The complete statement for the radical
follows
from a purely algebraic process,
see \cite[Theorem 4.22]{Demailly}.)

We will need a few preparatory results.
First recall that $S$ is an integral ring
over $R$
if every element of $S$ is expressed as
a root of a monic polynomial with
coefficients in $R$.
(See Stacks Project
\href{https://stacks.math.columbia.edu/tag/%
00GI}{00GI}.)

It will be handy to keep in mind the
following lemma
(Stacks Project
\href{https://stacks.math.columbia.edu/tag/%
00GH}{00GH}):

\begin{lemma}
\label{complex-analysis-lemma-chara-finit-terms-integ}
Let $R \to S$ be a ring map. The following
are equivalent
\begin{enumerate}
\item $R \to S$ is finite,
\item $R \to S$ is integral and of finite
type, and
\item there exist $x_1, \ldots, x_n \in S$
which generate $S$ as an
algebra over $R$ such that each $x_i$ is
integral over $R$.
\end{enumerate}
\end{lemma}

\noindent
This essentially says that an integral
extension
$R \to S$ gives $S$ the structure of a
finitely generated $R$-module
(I think experts call this a ``finite
$R$-module'')
whose generators are integral.
This just means that $S=R[x_1,\ldots,x_k]$
and $x_i \in S$ is the root of a monic
polynomial in $R$.

\begin{proposition}
\label{complex-analysis-proposition-integral-extension}
$\mathcal{O}_n/\mathcal{I}$ is a
finite integral extension of $\mathcal{O}_d$.
\end{proposition}

\begin{proof}
For this it's enough to check that
the map $\mathcal{O}_d \to
\mathcal{O}_n/\mathcal{I}$
is finite.
That is, that every element of
$\mathcal{O}_n/\mathcal{I}$
can be expressed as a tuple of elements
of $\mathcal{O}_d$.
(We can think of this tuple
as a linear combination, or cycle.)
Step 1: produce some Weierstrass
polynomials…
\end{proof}

\begin{lemma}
\label{complex-analysis-lemma-nullstellensatz}
\begin{reference}
\cite[Lemma 4.21]{Demailly}
\end{reference}
If $\mathcal{I} \subset \mathcal{O}_n$ is
prime and
$A = V(\mathcal{I})$, then
$\mathcal{I}_{A,0}=\mathcal{I}$.
\end{lemma}

\begin{proof}
Proposition
\ref{complex-analysis-proposition-integral-extension}
says: take a polynomial $f \in \mathcal{O}_n$
and consider its equivalence class
$\tilde{f}\in\mathcal{O}_n/\mathcal{I}$.
Then $\tilde{f}$ is the root of some monic
polynomial
$t^r+b_{r-1}t^{r-1}+\ldots+b_1t+b_0$
with $b_i \in \mathcal{O}_d$.
That is, putting $f$ instead of $t$ we get an
equation
$$
f^r+b_1f^{r-1}+\ldots+b_r=0\qquad \text{mod
}\mathcal{I}.
$$
In other words, the holomorphic function
$f^r+b_1f^{r-1}+\ldots+b_r$ is in
$\mathcal{I}$.
Since $f$ vanishes on $(A,0)$,
this says that $b_r$ vanishes in $(A,0)$ too.
Now suppose that this implied that the class
$\tilde{b_r}=0$.
We are supposed to use that Proposition
\ref{complex-analysis-proposition-integral-extension}
also tells us is that the above equation is
irreducible,
which I think means that $F=\sum b_it^i \in
\mathcal{O}_d[t]$
is irreducible.
In turn this gives that $r=1$.
\end{proof}


\medskip\noindent
Let's just store what was learned today:
\begin{enumerate}
\item {\bf Noether normalization lemma.}
See
\href{https://www.youtube.com/watch?v=%
VWWoG_4XVvs}
{this video}.
Let $k$ be a (infinite) field,
$A \neq 0$ a finitely generated $k$-algebra.
Then there exists elements $y_1,\ldots,y_r
\in A$
which are algebraically independent over $k$
and such that $A$ is integral over
$k[y_1,\ldots,y_r]$.
That is,
$$
\xymatrix{
A\ar@{-}[d]^{\substack{\text{integral} \\
\text{ext.}}}\\
k[y_1,\ldots,y_r]\ar@{-}[d]
^{\substack{\text{transc.}
\\ \text{ext.}}}\\
k.
}
$$

\item This can be used to prove the weak
Nullstellensatz, i.e.
let $k$ be an algberaically closed field, and
$I \subseteq k[x_1,\ldots,x_n]$
an ideal such that $V(I)=\emptyset$. Then
$I=k[x_1,\ldots,x_n]$.

To prove it using Noether's normalization
lemma,
suppose that the $I\neq k[x_1,\ldots,x_n]$.
Let $\mathfrak{m}$ be a maximal ideal such
that $I \subseteq \mathfrak{m}$.
Then
$A=k[x_1,\ldots,x_n]/\mathfrak{m}=k[y_1,
\ldots,y_n]$,
where the first string is algebraically
independent
while the last string is integral. But
because the
whole extension is algebraic, then
$y_1,\ldots,y_n$
are integral over $k$.
Then we can map $y_i \mapsto \alpha_i$, a
root
of its defining polynomial over $k$.
This means we do have a point where all
polynomials
in $\mathfrak{m}$ vanish.

{\bf Note.} I'm not sure what we mean with
"the first string"
and the "second string".

\item Also reviewed Misha's proof of the
Nullstellensatz.
Looks accessible, though it is not
holomorphic!

\item Perhaps start in Demailly just after
Prop. 4.13.
\end{enumerate}


% --- End complex-analysis.tex ---

% --- Begin functional-analysis.tex ---
\chapter{Functional Analysis}
\label{chapter-functional-analysis}


\label{functional-analysis-section-phantom-functional-analysis}


\section{Measure spaces}
\label{functional-analysis-section-measure-spaces}

\noindent These notes merge the home notes on
measure theory and functional analysis with
the class notes from Bruno Braga's course.
The point is to keep the material in the
style of the stack: definitions, lemmas,
theorems, examples, exercises, and
cross-references by labels.

\begin{definition}
\label{functional-analysis-definition-sigma-algebra}
Let $X$ be a set. An {\it algebra of sets} on
$X$ is a nonempty collection $\mathcal A$ of
subsets of $X$ closed under finite unions and
complements. A {\it $\sigma$-algebra} on $X$
is an algebra of sets closed under countable
unions.
\end{definition}

\begin{definition}
\label{functional-analysis-definition-generated-sigma-algebra}
Let $\mathcal E \subset \mathcal P(X)$. The
{\it $\sigma$-algebra generated by $\mathcal
E$} is the intersection of all
$\sigma$-algebras containing $\mathcal E$.
\end{definition}

\begin{definition}
\label{functional-analysis-definition-borel-algebra}
If $X$ is a topological space, the
$\sigma$-algebra generated by the open sets
of $X$ is called the {\it Borel
$\sigma$-algebra} of $X$ and is denoted
$\mathcal B_X$.
\end{definition}

\begin{definition}
\label{functional-analysis-definition-measure}
Let $\mathcal M$ be a $\sigma$-algebra on
$X$. A {\it measure} on $(X,\mathcal M)$ is a
map
$$
\mu : \mathcal M \longrightarrow [0,\infty]
$$
such that $\mu(\emptyset)=0$ and such that
whenever $E_1,E_2,\ldots$ are pairwise
disjoint measurable sets, one has
$$
\mu\left(\bigcup_{n=1}^\infty E_n\right) =
\sum_{n=1}^\infty \mu(E_n).
$$
The triple $(X,\mathcal M,\mu)$ is called a
{\it measure space}.
\end{definition}

\begin{lemma}
\label{functional-analysis-lemma-basic-properties-measures}
Let $(X,\mathcal M,\mu)$ be a measure space.
\begin{enumerate}
\item If $E \subset F$, then $\mu(E) \leq
\mu(F)$.
\item If $E_1,E_2,\ldots$ are measurable,
then
$$
\mu\left(\bigcup_{n=1}^\infty E_n\right) \leq
\sum_{n=1}^\infty \mu(E_n).
$$
\item If $E_1\subset E_2\subset\ldots$, then
$$
\mu\left(\bigcup_{n=1}^\infty
E_n\right)=\lim_n \mu(E_n).
$$
\item If $E_1\supset E_2\supset\ldots$ and
$\mu(E_1)<\infty$, then
$$
\mu\left(\bigcap_{n=1}^\infty
E_n\right)=\lim_n\mu(E_n).
$$
\end{enumerate}
\end{lemma}

\begin{definition}
\label{functional-analysis-definition-null-set-complete-measure}
A measurable set $N$ is a {\it null set} if
$\mu(N)=0$. A measure is {\it complete} if
every subset of a null set is measurable.
\end{definition}

\begin{theorem}
\label{functional-analysis-theorem-completion-measure}
Every measure has a completion. More
precisely, if $(X,\mathcal M,\mu)$ is a
measure space and
$$
\mathcal N = \{N\in \mathcal M : \mu(N)=0\},
$$
then the collection
$$
\overline{\mathcal M} = \{E\cup A :
E\in\mathcal M,\ A\subset N \text{ for some
}N\in\mathcal N\}
$$
is a $\sigma$-algebra, and $\mu$ extends
uniquely to a complete measure on
$\overline{\mathcal M}$.
\end{theorem}

\begin{definition}
\label{functional-analysis-definition-outer-measure}
An {\it outer measure} on $X$ is a function
$\mu^*:\mathcal P(X)\to[0,\infty]$ such that
\begin{enumerate}
\item $\mu^*(\emptyset)=0$;
\item $A\subset B$ implies
$\mu^*(A)\leq\mu^*(B)$;
\item $\mu^*(\bigcup A_n)\leq \sum
\mu^*(A_n)$.
\end{enumerate}
\end{definition}

\begin{definition}
\label{functional-analysis-definition-carath-measurable}
Let $\mu^*$ be an outer measure on $X$. A set
$A\subset X$ is {\it $\mu^*$-measurable} if
for every $E\subset X$ one has
$$
\mu^*(E)=\mu^*(E\cap A)+\mu^*(E\cap A^c).
$$
\end{definition}

\begin{theorem}[Caratheodory]
\label{functional-analysis-theorem-caratheodory}
If $\mu^*$ is an outer measure on $X$, then
the collection of $\mu^*$-measurable subsets
of $X$ is a $\sigma$-algebra and the
restriction of $\mu^*$ to it is a complete
measure.
\end{theorem}

\begin{definition}
\label{functional-analysis-definition-premeasure}
Let $\mathcal A$ be an algebra of subsets of
$X$. A function $\mu_0:\mathcal
A\to[0,\infty]$ is a {\it premeasure} if
$\mu_0(\emptyset)=0$ and if countable
additivity holds whenever the countable
disjoint union still belongs to $\mathcal A$.
\end{definition}

\begin{theorem}
\label{functional-analysis-theorem-extension-premeasure}
Let $\mu_0$ be a premeasure on an algebra
$\mathcal A$, and let $\mathcal M$ be the
$\sigma$-algebra generated by $\mathcal A$.
Then $\mu_0$ extends to a measure on
$\mathcal M$. If $\mu_0$ is $\sigma$-finite,
this extension is unique.
\end{theorem}

\begin{definition}
\label{functional-analysis-definition-borel-measure-r}
A {\it Borel measure on $\mathbb{R}$} is a
measure on $\mathbb{R}$ whose domain is the
Borel $\sigma$-algebra $\mathcal
B_{\mathbb{R}}$.
\end{definition}

\begin{proposition}
\label{functional-analysis-proposition-premeasure-half-open}
Let $F:\mathbb{R}\to\mathbb{R}$ be an
increasing and right-continuous function. The
collection $\mathcal A$ of finite disjoint
unions of half-open intervals $(a,b]$
(including $(a,\infty)$ and $(-\infty,b]$) is
an algebra. The function $\mu_0$ defined by
$$
\mu_0\left(\bigcup_{j=1}^n (a_j, b_j]\right)
= \sum_{j=1}^n (F(b_j) - F(a_j))
$$
is a premeasure on $\mathcal A$.
\end{proposition}

\begin{theorem}
\label{functional-analysis-theorem-lebesgue-stieltjes}
If $F:\mathbb{R}\to\mathbb{R}$ is an
increasing and right-continuous function,
there is a unique Borel measure $\mu_F$ on
$\mathbb{R}$ such that $\mu_F((a,b]) =
F(b)-F(a)$ for all $a,b\in\mathbb{R}$. The
completion of $\mu_F$ is called the {\it
Lebesgue--Stieltjes measure} associated with
$F$.
\end{theorem}

\begin{example}
\label{functional-analysis-example-lebesgue-measure}
The {\it Lebesgue measure} $m$ on
$\mathbb{R}$ is the Lebesgue--Stieltjes
measure associated with the function
$F(x)=x$. It is the unique measure on
$(\mathbb{R}, \mathcal B_{\mathbb{R}})$ that
assigns to each interval its length.
\end{example}

\section
{Measurable functions and integration}
\label
{section-measurable-functions-integration}

\begin{definition}
\label{functional-analysis-definition-measurable-map}
Let $(X,\mathcal M)$ and $(Y,\mathcal N)$ be
measurable spaces. A map $f:X\to Y$ is {\it
measurable} if $f^{-1}(E)\in\mathcal M$ for
every $E\in\mathcal N$.
\end{definition}

\begin{lemma}
\label{functional-analysis-lemma-measurable-generated}
If $\mathcal N$ is generated by a collection
$\mathcal E$, then $f:X\to Y$ is measurable
if and only if $f^{-1}(E)\in\mathcal M$ for
all $E\in\mathcal E$.
\end{lemma}

\begin{example}
\label{functional-analysis-example-continuous-borel-measurable}
If $X,Y$ are topological spaces and $f:X\to
Y$ is continuous, then $f$ is $(\mathcal
B_X,\mathcal B_Y)$-measurable.
\end{example}

\begin{definition}
\label{functional-analysis-definition-simple-function}
A complex-valued function $\varphi$ on $X$ is
{\it simple} if it is a finite linear
combination of characteristic functions of
measurable sets. Equivalently, $\varphi$ is
measurable and has finite image.
\end{definition}

\begin{theorem}
\label
{theorem-approximation-simple-functions}
Let $(X,\mathcal M)$ be a measurable space.
\begin{enumerate}
\item If $f:X\to[0,\infty]$ is measurable,
then there is an increasing sequence of
nonnegative simple functions $\varphi_n$ with
$\varphi_n\to f$ pointwise.
\item If $f:X\to\mathbb{C}$ is measurable,
then there is a sequence of simple functions
$\varphi_n$ with $\varphi_n\to f$ pointwise
and $|\varphi_n|\leq |f|$.
\end{enumerate}
\end{theorem}

\begin{proof}
For (1), for each $n\in\mathbb{N}$ and $0\leq
k\leq 2^{2n}-1$, let
$$
E_n^k = f^{-1}((k2^{-n}, (k+1)2^{-n}]) \quad
\text{and} \quad F_n = f^{-1}((2^n, \infty]).
$$
Define
$$
\varphi_n = \sum_{k=0}^{2^{2n}-1}
k2^{-n}\chi_{E_n^k} + 2^n\chi_{F_n}.
$$
One checks that $\varphi_n \leq
\varphi_{n+1}$ and $\varphi_n \to f$
pointwise. For (2), decompose $f = f_1 - f_2
+ i(f_3 - f_4)$ where $f_j \geq 0$ are the
positive and negative parts of the real and
imaginary parts of $f$. Apply part (1) to
each $f_j$.
\end{proof}

\begin{definition}
\label
{definition-integral-nonnegative-simple}
Let $\varphi=\sum_{j=1}^n a_j\chi_{E_j}$ be a
nonnegative simple function written in
standard form. Define
$$
\int \varphi\,d\mu = \sum_{j=1}^n
a_j\mu(E_j),
$$
with the convention $0\cdot\infty=0$.
\end{definition}

\begin{definition}
\label{functional-analysis-definition-integral-nonnegative}
For $f:X\to[0,\infty]$ measurable, define
$$
\int f\,d\mu = \sup\left\{\int \varphi\,d\mu:
0\leq \varphi\leq f,\ \varphi\text{
simple}\right\}.
$$
\end{definition}

\begin{theorem}[Monotone convergence]
\label{functional-analysis-theorem-monotone-convergence}
If $0\leq f_1\leq f_2\leq\ldots$ are
measurable and $f_n\to f$ pointwise, then
$$
\int f\,d\mu = \lim_n \int f_n\,d\mu.
$$
\end{theorem}

\begin{lemma}[Fatou]
\label{functional-analysis-lemma-fatou}
If $f_n\geq 0$ are measurable, then
$$
\int \liminf_n f_n\,d\mu \leq \liminf_n \int
f_n\,d\mu.
$$
\end{lemma}

\begin{theorem}[Dominated convergence]
\label{functional-analysis-theorem-dominated-convergence}
Let $f_n$ and $f$ be measurable functions
such that $f_n\to f$ almost everywhere.
Suppose that there is an integrable function
$g$ such that, for every $n$,
$$
|f_n|\leq g
$$
almost everywhere. Then $f$ is integrable
and
$$
\lim_n\int |f_n-f|\,d\mu=0.
$$
In particular, $\int f_n\,d\mu\to\int
f\,d\mu$.
\end{theorem}

\begin{definition}
\label{functional-analysis-definition-counting-measure}
The {\it counting measure} on
$\mathbb{N}$ is the measure $\#$ defined by
$$
\#(A)=|A|.
$$
If $A$ is infinite, this value is $\infty$.
For a nonnegative sequence $(a_n)$,
$$
\int_{\mathbb{N}} a_n\,d\#
=
\sum_{n=1}^{\infty}a_n.
$$
For a complex sequence the same identity
holds whenever it is absolutely summable.
\end{definition}

\begin{theorem}[Differentiating a series]
\label{functional-analysis-theorem-differentiating-series}
Let $I=[a,b]$ and let $f_n\in C^1(I)$.
Assume that
$$
\sum_{n=1}^{\infty}f_n(x)
$$
converges for every $x\in I$, and that
$$
\sum_{n=1}^{\infty}f_n'
$$
converges uniformly on $I$. Then the first
series converges uniformly on $I$ to a
function $f\in C^1(I)$, and
$$
f'=\sum_{n=1}^{\infty}f_n'.
$$
In fact, it is enough that the first series
converges at one point of $I$.
\end{theorem}

\begin{proof}
Let
$$
S_N=\sum_{n=1}^N f_n
$$
and let $g$ be the uniform limit of $S_N'$.
The function $g$ is continuous. Fix
$x_0\in I$ where $S_N(x_0)$ converges. By
the fundamental theorem of calculus,
$$
S_N(x)=S_N(x_0)
+\int_{x_0}^x S_N'(t)\,dt.
$$
Hence
$$
|S_N(x)-S_M(x)|
\leq |S_N(x_0)-S_M(x_0)|
+(b-a)\|S_N'-S_M'\|_\infty.
$$
Thus $S_N$ is uniformly Cauchy on $I$.
Let $f$ be its uniform limit. Passing to the
limit in the preceding integral identity
gives
$$
f(x)=f(x_0)+\int_{x_0}^x g(t)\,dt.
$$
Therefore $f\in C^1(I)$ and $f'=g$.
\end{proof}

\begin{remark}
\label{functional-analysis-remark-differentiating-series-dct}
There is a stronger sufficient hypothesis
under which dominated convergence justifies
the derivative identity directly. Suppose
there are
numbers $M_n\geq0$ such that
$$
|f_n'(x)|\leq M_n
$$
for every $x\in I$, and
$$
\sum_{n=1}^{\infty}M_n<\infty.
$$
These assumptions imply uniform convergence
of the derivative series by the Weierstrass
criterion, but they are stronger than mere
uniform convergence of that series.

Fix $x\in(a,b)$. For small nonzero $h$ set,
as a function on $\mathbb{N}$,
$$
q_h(n)=
\frac{f_n(x+h)-f_n(x)}{h}.
$$
The fundamental theorem of calculus gives
$$
|q_h(n)|\leq M_n.
$$
For every $n$, one has
$$
q_h(n)\longrightarrow f_n'(x).
$$
If $h_k\to0$, dominated convergence for
counting measure therefore gives
$$
\lim_{k\to\infty}
\sum_{n=1}^{\infty}q_{h_k}(n)
=
\sum_{n=1}^{\infty}f_n'(x).
$$
Since the original series converges at
$x$ and at $x+h_k$,
$$
\sum_{n=1}^{\infty}q_{h_k}(n)
=
\frac{f(x+h_k)-f(x)}{h_k}.
$$
Hence $f'(x)$ is the derivative series.
This is the counting-measure version of the
difference-quotient argument in Folland,
{\it Real Analysis: Modern Techniques and
Their Applications}, second edition,
Theorem 2.27(b), p.~56.
\end{remark}

\begin{definition}
\label{functional-analysis-definition-lp-space}
Let $1\leq p<\infty$. The space $L^p(X,\mu)$
is the vector space of measurable functions
$f$ with
$$
\|f\|_p = \left(\int
|f|^p\,d\mu\right)^{1/p}<\infty,
$$
modulo equality almost everywhere. Also
$L^\infty(X,\mu)$ is the space of essentially
bounded measurable functions with the
essential supremum norm.
\end{definition}

\begin{theorem}[Holder inequality]
\label{functional-analysis-theorem-holder-inequality}
Let $1<p,q<\infty$ and $1/p+1/q=1$. If $f\in
L^p$ and $g\in L^q$, then $fg\in L^1$ and
$$
\|fg\|_1\leq \|f\|_p\|g\|_q.
$$
\end{theorem}

\begin{theorem}[Minkowski inequality]
\label{functional-analysis-theorem-minkowski-inequality}
For $1\leq p\leq\infty$, $L^p(X,\mu)$ is a
normed vector space. Equivalently, if $f,g\in
L^p$, then
$$
\|f+g\|_p\leq \|f\|_p+\|g\|_p.
$$
\end{theorem}

\begin{theorem}
\label{functional-analysis-theorem-lp-banach}
For $1\leq p\leq\infty$, the space
$L^p(X,\mu)$ is a Banach space.
\end{theorem}

\begin{example}
\label{functional-analysis-example-sequence-lp}
For $1\leq p<\infty$,
$$
\ell^p = \left\{x=(x_n):\sum_{n=1}^\infty
|x_n|^p<\infty\right\}
$$
with $\|x\|_p=(\sum |x_n|^p)^{1/p}$ is a
Banach space. Also $\ell^\infty$ is a Banach
space with $\|x\|_\infty=\sup_n |x_n|$, and
$c_0\subset\ell^\infty$ is the closed
subspace of sequences converging to zero.
\end{example}

\begin{definition}
\label{functional-analysis-definition-convergence-in-measure}
Let $(X,\mathcal M,\mu)$ be a measure space.
A sequence $f_n$ converges to $f$ {\it in
measure} if for every $\varepsilon>0$,
$$
\mu\{x: |f_n(x)-f(x)|>\varepsilon\}\to 0.
$$
\end{definition}

\begin{theorem}
\label
{theorem-l1-convergence-implies-measure}
If $f_n\to f$ in $L^1$, then $f_n\to f$ in
measure. Moreover, some subsequence converges
to $f$ almost everywhere.
\end{theorem}

\begin{theorem}[Egoroff]
\label{functional-analysis-theorem-egoroff}
Assume $\mu(X)<\infty$. If $f_n\to f$ almost
everywhere, then for every $\varepsilon>0$
there is a measurable set $E$ with
$\mu(E)<\varepsilon$ such that $f_n\to f$
uniformly on $X\setminus E$.
\end{theorem}

\begin{definition}
\label{functional-analysis-definition-almost-uniform-convergence}
A sequence of measurable functions $f_n$ is
said to converge to $f$ {\it almost
uniformly} if for every $\varepsilon>0$ there
exists a measurable set $E$ with
$\mu(E)<\varepsilon$ such that $f_n \to f$
uniformly on $E^c$.
\end{definition}

\begin{remark}
\label{functional-analysis-remark-egoroff-almost-uniform}
Egoroff's theorem can be restated as: on a
finite measure space, almost everywhere
convergence implies almost uniform
convergence.
\end{remark}

\section{Product measures}
\label{functional-analysis-section-product-measures}

\begin{definition}
\label{functional-analysis-definition-product-sigma-algebra}
Let $(X,\mathcal M)$ and $(Y,\mathcal N)$ be
measurable spaces. The {\it product
$\sigma$-algebra} $\mathcal M\otimes\mathcal
N$ is the $\sigma$-algebra on $X\times Y$
generated by the measurable rectangles
$A\times B$, where $A\in\mathcal M$ and
$B\in\mathcal N$.
\end{definition}

\begin{proposition}
\label{functional-analysis-proposition-product-borel}
Let $X_1, \ldots, X_n$ be metric spaces and
$X = \prod X_i$ with the product metric. Then
$\bigotimes \mathcal B_{X_i} \subset \mathcal
B_X$. If each $X_i$ is separable, then
$\bigotimes \mathcal B_{X_i} = \mathcal B_X$.
\end{proposition}

\begin{theorem}
\label{functional-analysis-theorem-product-measure}
Let $(X,\mathcal M,\mu)$ and $(Y,\mathcal
N,\nu)$ be $\sigma$-finite measure spaces.
There is a unique measure $\mu\times\nu$ on
$\mathcal M\otimes\mathcal N$ such that
$$
(\mu\times\nu)(A\times B)=\mu(A)\nu(B)
$$
for all measurable rectangles $A\times B$.
\end{theorem}

\begin{definition}
\label{functional-analysis-definition-sections}
If $E\subset X\times Y$, define
$$
E_x = \{y\in Y:(x,y)\in E\}, \quad E^y =
\{x\in X:(x,y)\in E\}.
$$
If $f:X\times Y\to\mathbb{C}$, define
$f_x(y)=f^y(x)=f(x,y)$.
\end{definition}

\begin{theorem}[Tonelli]
\label{functional-analysis-theorem-tonelli}
Let $(X,\mathcal M,\mu)$ and $(Y,\mathcal
N,\nu)$ be $\sigma$-finite measure spaces. If
$f:X\times Y\to[0,\infty]$ is measurable,
then $x\mapsto\int f_x\,d\nu$ and
$y\mapsto\int f^y\,d\mu$ are measurable, and
$$
\int_{X\times Y} f\,d(\mu\times\nu) =
\int_X\left(\int_Y
f(x,y)\,d\nu(y)\right)d\mu(x) =
\int_Y\left(\int_X
f(x,y)\,d\mu(x)\right)d\nu(y).
$$
\end{theorem}

\begin{theorem}[Fubini]
\label{functional-analysis-theorem-fubini}
Under the hypotheses of Theorem
\ref{functional-analysis-theorem-tonelli}, if $f\in L^1(X\times
Y)$, then almost all sections are integrable
and the same iterated-integral formula holds.
\end{theorem}

\section{Compactness and spaces of functions}
\label{functional-analysis-section-compactness-function-spaces}

\begin{definition}
\label{functional-analysis-definition-uniform-norm}
Let $X$ be a set. The space $B(X)$ of bounded
scalar-valued functions on $X$ is normed by
$$
\|f\|_u = \sup_{x\in X} |f(x)|.
$$
\end{definition}

\begin{proposition}
\label{functional-analysis-proposition-bx-banach}
The space $B(X)$ is a Banach space. If $X$ is
a topological space, then $BC(X)$, the space
of bounded continuous functions, is a closed
subspace of $B(X)$ and is therefore Banach.
\end{proposition}

\begin{definition}
\label{functional-analysis-definition-compact}
A topological space $X$ is {\it compact} if
every open cover of $X$ has a finite
subcover. A subset is {\it precompact} if its
closure is compact.
\end{definition}

\begin{lemma}
\label{functional-analysis-lemma-compact-closed-fip}
A topological space $X$ is compact if and
only if every family of closed sets with the
finite intersection property has nonempty
intersection.
\end{lemma}

\begin{theorem}
\label
{theorem-compact-metric-characterization}
Let $(X,d)$ be a metric space and $E\subset
X$. Then the following are equivalent.
\begin{enumerate}
\item $E$ is compact.
\item Every sequence in $E$ has a subsequence
converging to a point of $E$.
\item $E$ is complete and totally bounded.
\end{enumerate}
\end{theorem}

\begin{theorem}[Urysohn]
\label{functional-analysis-theorem-urysohn}
Let $X$ be a normal topological space. If $A$
and $B$ are disjoint closed subsets of $X$,
then there is a continuous function
$f:X\to[0,1]$ such that $f|_A=0$ and
$f|_B=1$.
\end{theorem}

\begin{theorem}[Tietze]
\label{functional-analysis-theorem-tietze}
Let $X$ be a normal topological space, let
$A\subset X$ be closed, and let $f:A\to[a,b]$
be continuous. Then there is a continuous
$F:X\to[a,b]$ such that $F|_A=f$.
\end{theorem}

\begin{definition}
\label{functional-analysis-definition-equicontinuous}
Let $X$ be a topological space. A family
$\mathcal F \subset C(X)$ is {\it
equicontinuous} if for every $x\in X$ and
every $\varepsilon>0$, there is a
neighborhood $U$ of $x$ such that
$|f(x)-f(y)|<\varepsilon$ for all $y\in U$
and all $f\in\mathcal F$.
\end{definition}

\begin{theorem}[Arzelà--Ascoli]
\label{functional-analysis-theorem-arzela-ascoli}
Let $X$ be a compact Hausdorff space. A
subset $\mathcal F \subset C(X)$ is
precompact in the uniform norm if and only if
it is equicontinuous and pointwise bounded.
\end{theorem}

\begin{theorem}[Stone--Weierstrass]
\label{functional-analysis-theorem-stone-weierstrass}
Let $X$ be a compact Hausdorff space. If
$\mathcal A$ is a subalgebra of $C(X,
\mathbb{R})$ that separates points and
contains the constant functions, then
$\mathcal A$ is dense in $C(X, \mathbb{R})$
for the uniform norm. If $\mathcal A$ is a
subalgebra of $C(X, \mathbb{C})$ that
separates points, contains the constant
functions, and is closed under complex
conjugation, then $\mathcal A$ is dense in
$C(X, \mathbb{C})$.
\end{theorem}

\begin{definition}
\label{functional-analysis-definition-cc-c0}
Let $X$ be a locally compact Hausdorff space.
\begin{enumerate}
\item $C_c(X)$ is the space of continuous
functions with compact support.
\item $C_0(X)$ is the space of continuous
functions that vanish at infinity, i.e., for
every $\varepsilon>0$, the set $\{x\in X :
|f(x)|\geq\varepsilon\}$ is compact.
\end{enumerate}
\end{definition}

\begin{proposition}
\label{functional-analysis-proposition-c0-completion-cc}
If $X$ is a locally compact Hausdorff space,
$C_0(X)$ is the completion of $C_c(X)$ in the
uniform norm.
\end{proposition}

\section{Normed spaces and Banach spaces}
\label{functional-analysis-section-normed-banach-spaces}

\begin{definition}
\label{functional-analysis-definition-norm}
Let $X$ be a vector space over a field
$\mathbb{K} \in \{\mathbb{R}, \mathbb{C}\}$.
A function $\|\cdot\|: X \to [0, \infty)$ is
a {\it norm} if:
\begin{enumerate}
\item $\|x\| = 0$ if and only if $x = 0$;
\item $\|\lambda x\| = |\lambda|\|x\|$ for
all $\lambda \in \mathbb{K}$ and $x \in X$;
\item $\|x+y\| \leq \|x\| + \|y\|$ for all
$x, y \in X$.
\end{enumerate}
The pair $(X, \|\cdot\|)$ is called a {\it
normed space}. A normed space is a metric
space with the metric $d(x,y) = \|x-y\|$. If
$(X, \|\cdot\|)$ is complete as a metric
space, it is called a {\it Banach space}.
\end{definition}

\begin{example}
\label{functional-analysis-example-banach-spaces}
The following are important examples of
Banach spaces:
\begin{enumerate}
\item {\bf Sequence spaces.} For $1 \leq p <
\infty$, the space
$$
\ell_p = \left\{ (x_n) \in
\mathbb{K}^\mathbb{N} : \sum_{n=1}^\infty
|x_n|^p < \infty \right\}
$$
with the norm $\|(x_n)\|_p = (\sum
|x_n|^p)^{1/p}$ is a Banach space.
\item {\bf Bounded sequences.} The space
$$
\ell_\infty = \{ (x_n) \in
\mathbb{K}^\mathbb{N} : \sup_n |x_n| < \infty
\}
$$
with the norm $\|(x_n)\|_\infty = \sup_n
|x_n|$ is a Banach space.
\item {\bf Null sequences.} The space $c_0 =
\{ (x_n) \in \ell_\infty : \lim_n x_n = 0 \}$
is a closed subspace of $\ell_\infty$ and
thus a Banach space.
\item {\bf Continuous functions.} If $K$ is a
compact Hausdorff space, $C(K)$ with the
supremum norm $\|f\| = \sup_{x \in K} |f(x)|$
is a Banach space.
\item {\bf $L^p$ spaces.} Let $(X,
\mathcal{A}, \mu)$ be a measure space. For $1
\leq p < \infty$, the space $L^p(\mu)$ of
equivalence classes of measurable functions
such that $\int |f|^p d\mu < \infty$ is a
Banach space.
\end{enumerate}
\end{example}

\begin{lemma}[Young's inequality]
\label{functional-analysis-lemma-young-inequality}
Let $a, b \geq 0$ and $p, q > 1$ such that
$1/p + 1/q = 1$. Then
$$
ab \leq \frac{a^p}{p} + \frac{b^q}{q}.
$$
\end{lemma}

\begin{proof}
Consider the function $f(x) = x^{p-1}$. The
area under $f$ from $0$ to $a$ is $\int_0^a
x^{p-1} dx = a^p/p$. The inverse function is
$g(y) = y^{1/(p-1)} = y^{q-1}$. The area
under $g$ from $0$ to $b$ is $\int_0^b
y^{q-1} dy = b^q/q$. By drawing a figure, it
is clear that for any $a, b > 0$, the sum of
these areas is greater than or equal to the
area of the rectangle with sides $a$ and $b$.
\end{proof}

\begin{theorem}[Hölder's inequality]
\label{functional-analysis-theorem-holder-inequality-duplicate-2}
Let $p, q \in (1, \infty)$ with $1/p + 1/q =
1$. If $(x_n) \in \ell_p$ and $(y_n) \in
\ell_q$, then $(x_n y_n) \in \ell_1$ and
$$
\|(x_n y_n)\|_1 \leq \|(x_n)\|_p \|(y_n)\|_q.
$$
\end{theorem}

\begin{proof}
If $\|(x_n)\|_p = 0$ or $\|(y_n)\|_q = 0$,
the inequality is trivial. Assume
$\|(x_n)\|_p = \|(y_n)\|_q = 1$. By Young's
inequality, for each $n$:
$$
|x_n y_n| \leq \frac{|x_n|^p}{p} +
\frac{|y_n|^q}{q}.
$$
Summing over $n$:
$$
\sum |x_n y_n| \leq \frac{1}{p} \sum |x_n|^p
+ \frac{1}{q} \sum |y_n|^q = \frac{1}{p} +
\frac{1}{q} = 1.
$$
For the general case, normalize $\tilde{x}_n
= x_n / \|x\|_p$ and $\tilde{y}_n = y_n /
\|y\|_q$.
\end{proof}

\begin{theorem}[Minkowski's inequality]
\label{functional-analysis-theorem-minkowski-inequality-duplicate-2}
If $1 \leq p < \infty$ and $x, y \in \ell_p$,
then $x+y \in \ell_p$ and $\|x+y\|_p \leq
\|x\|_p + \|y\|_p$.
\end{theorem}

\begin{proof}
For $p=1$ it is the triangle inequality. For
$p > 1$:
$$
\begin{aligned}
\|x+y\|_p^p &= \sum |x_n+y_n|^p \leq \sum
|x_n||x_n+y_n|^{p-1} + \sum
|y_n||x_n+y_n|^{p-1} \\ &\leq \|x\|_p
\|(|x_n+y_n|^{p-1})\|_q + \|y\|_p
\|(|x_n+y_n|^{p-1})\|_q \\
&\quad \text{(Hölder)}
\end{aligned}
$$

\noindent
Note that $(p-1)q = p$. Thus
$\|(|x_n+y_n|^{p-1})\|_q = (\sum
|x_n+y_n|^p)^{1/q} = \|x+y\|_p^{p/q}$.
Dividing by $\|x+y\|_p^{p/q}$ (assuming $x+y
\neq 0$) and using $p - p/q = 1$, we get the
result.
\end{proof}

\begin{theorem}
\label{functional-analysis-theorem-lp-completeness}
The space $\ell_p$ is a Banach space for $1
\leq p \leq \infty$.
\end{theorem}

\begin{proof}
We treat the case $1 \leq p < \infty$. Let
$(x^{(n)})$ be a Cauchy sequence in $\ell_p$,
where $x^{(n)} = (x_k^{(n)})_k$. For each
fixed $k$, $|x_k^{(n)} - x_k^{(m)}| \leq
\|x^{(n)} - x^{(m)}\|_p$, so $(x_k^{(n)})_n$
is a Cauchy sequence in $\mathbb{K}$. Let
$x_k = \lim_n x_k^{(n)}$ and $x = (x_k)_k$.
Since $(x^{(n)})$ is Cauchy, it is bounded:
$\|x^{(n)}\|_p \leq M$ for all $n$. For any
$N \in \mathbb{N}$:
$$
\sum_{k=1}^N |x_k^{(n)}|^p \leq M^p.
$$
Taking $n \to \infty$, $\sum_{k=1}^N |x_k|^p
\leq M^p$. As $N$ is arbitrary, $x \in
\ell_p$. Now show $x^{(n)} \to x$ in norm.
Given $\varepsilon > 0$, $\exists n_0$ s.t.
$m, n > n_0 \implies \|x^{(n)} - x^{(m)}\|_p
< \varepsilon$. For any $N$:
$$
\sum_{k=1}^N |x_k^{(n)} - x_k^{(m)}|^p <
\varepsilon^p.
$$
Taking $m \to \infty$: $\sum_{k=1}^N
|x_k^{(n)} - x_k|^p \leq \varepsilon^p$. As
$N$ is arbitrary, $\|x^{(n)} - x\|_p \leq
\varepsilon$.
\end{proof}

\begin{remark}
\label{functional-analysis-remark-lp-inclusions}
\begin{enumerate}
\item If $p \leq q$, then $\ell_p \subseteq
\ell_q$.
\item If $\mu(X) < \infty$ and $p \leq q$,
then $L^q(\mu) \subseteq L^p(\mu)$.
\item The space $c_{00}$ of sequences with
finite support is dense in $\ell_p$ for $p <
\infty$, but not in $\ell_\infty$.
\end{enumerate}
\end{remark}


\section{Quotient spaces}
\label{functional-analysis-section-quotient-spaces}

Given a normed space $X$ and a linear
subspace $Y$, we can consider the quotient
vector space $X/Y$ consisting of cosets $[x]
= x + Y$.

\begin{definition}
\label{functional-analysis-definition-quotient-norm}
Let $X$ be a normed space and $Y \subseteq X$
be a subspace. The {\it quotient seminorm} on
$X/Y$ is defined by
$$
\|x+Y\|_{X/Y} = \inf \{ \|x-y\| : y \in Y \}
= d(x, Y).
$$
\end{definition}

\begin{proposition}
\label{functional-analysis-proposition-quotient-norm-properties}
Let $X$ be a normed space and $Y$ a subspace.
\begin{enumerate}
\item The function $\|\cdot\|_{X/Y}$ is a
seminorm on $X/Y$.
\item It is a norm if and only if $Y$ is
closed.
\item If $X$ is a Banach space and $Y$ is
closed, then $X/Y$ is a Banach space.
\end{enumerate}
\end{proposition}

\begin{proof}
(1) and (2) are straightforward. For (3),
assume $X$ is Banach and $Y$ is closed. Let
$([x_n])$ be a Cauchy sequence in $X/Y$.
Passing to a subsequence, we may assume
$\|[x_n] - [x_{n+1}]\| < 2^{-n}$. We find a
sequence $(y_n)$ in $X$ such that $y_n \in
[x_n]$ and $\|y_n - y_{n+1}\| < 2^{-n}$ by
choosing $w_n \in Y$ such that $\|x_n -
x_{n+1} - w_n\| < 2^{-n}$ and setting $y_n$
appropriately. Then $(y_n)$ is Cauchy in $X$,
hence converges to some $y \in X$. Then
$[x_n] \to [y]$ in $X/Y$.
\end{proof}

\begin{definition}
\label{functional-analysis-definition-canonical-projection}
The map $Q : X \to X/Y$ defined by $Q(x) =
x+Y$ is called the {\it canonical
projection}.
\end{definition}

\begin{proposition}
\label{functional-analysis-proposition-canonical-projection-norm}
The canonical projection $Q : X \to X/Y$ is a
linear operator with $\|Q\| \leq 1$. If $Y$
is a proper closed subspace, then $\|Q\| =
1$.
\end{proposition}

\begin{theorem}[First isomorphism theorem]
\label{functional-analysis-theorem-isomorphism-theorem-operators}
Let $T : X \to Y$ be a bounded linear
operator. There exists a unique bounded
linear operator $\widehat{T} : X/\ker T \to
Y$ such that $T = \widehat{T} \circ Q$.
Moreover, $\|\widehat{T}\| = \|T\|$.
\end{theorem}

\begin{proof}
The algebraic existence of $\widehat{T}$ is
standard. For the norm:
$$
\begin{aligned}
\|\widehat{T}\| &= \sup \{ \|T(x)\| : x+\ker
T \in B_{X/\ker T} \} \\ &= \sup \{
\|T(x-y)\| : y \in \ker T, \|x+\ker T\| \leq
1 \}.
\end{aligned}
$$

\noindent
Given $\varepsilon > 0$, for each $x$ with
$\|x+\ker T\| \leq 1$, we can find $y \in
\ker T$ such that $\|x-y\| < 1+\varepsilon$.
Thus $\|T(x-y)\| \leq \|T\|(1+\varepsilon)$.
This shows $\|\widehat{T}\| \leq \|T\|$. The
reverse inequality $\|T\| \leq
\|\widehat{T}\|\|Q\| \leq \|\widehat{T}\|$ is
immediate.
\end{proof}

\section{Linear maps and bounded operators}
\label{functional-analysis-section-bounded-operators}

\begin{definition}
\label{functional-analysis-definition-bounded-linear-map}
Let $X, Y$ be normed spaces. A linear map $T
: X \to Y$ is {\it bounded} if there is $C
\geq 0$ such that $\|Tx\| \leq C\|x\|$ for
all $x \in X$.
\end{definition}

\begin{proposition}
\label{functional-analysis-proposition-bounded-iff-continuous}
Let $T : X \to Y$ be linear. The following
are equivalent:
\begin{enumerate}
\item $T$ is continuous at $0$.
\item $T$ is continuous.
\item $T$ is bounded.
\item $T$ is uniformly continuous.
\end{enumerate}
\end{proposition}

\begin{proof}
The equivalence $(1) \iff (2) \iff (4)$ is
standard for linear maps. For $(2) \implies
(3)$, if $T$ is continuous at $0$, then for
$\varepsilon=1$ there exists $\delta>0$ such
that $\|x\| < \delta \implies \|Tx\| < 1$.
For any $x \neq 0$, $z =
\frac{\delta}{2\|x\|}x$ has norm $\delta/2 <
\delta$, so $\|Tz\| < 1$, which means
$\frac{\delta}{2\|x\|}\|Tx\| < 1$, i.e.,
$\|Tx\| \leq \frac{2}{\delta}\|x\|$.
\end{proof}

\begin{definition}
\label{functional-analysis-definition-operator-norm}
The vector space of bounded linear maps from
$X$ to $Y$ is denoted $\mathcal{L}(X, Y)$.
The {\it operator norm} on $\mathcal{L}(X,
Y)$ is defined by
$$
\|T\| = \sup_{x \in B_X} \|Tx\| =
\sup_{\|x\|=1} \|Tx\| = \inf \{ C \geq 0 :
\|Tx\| \leq C\|x\| \text{ for all } x \in X
\}.
$$
\end{definition}

\begin{proposition}
\label{functional-analysis-proposition-operator-space-banach}
If $Y$ is a Banach space, then
$\mathcal{L}(X, Y)$ is a Banach space.
\end{proposition}

\begin{proof}
Let $(T_n)$ be a Cauchy sequence in
$\mathcal{L}(X, Y)$. For each $x \in X$,
$\|T_n x - T_m x\| \leq \|T_n - T_m\| \|x\|$,
so $(T_n x)$ is Cauchy in $Y$. Since $Y$ is
complete, let $Tx = \lim_n T_n x$. Linearity
of $T$ is easy. Since $\|T_n\|$ is bounded
(say by $M$), $\|Tx\| = \lim \|T_n x\| \leq
M\|x\|$, so $T$ is bounded. To show $\|T_n -
T\| \to 0$, for fixed $\varepsilon > 0$,
choose $N$ such that $n, m \geq N \implies
\|T_n - T_m\| < \varepsilon$. For $\|x\| \leq
1$, $\|T_n x - T_m x\| < \varepsilon$.
Letting $m \to \infty$, $\|T_n x - Tx\| \leq
\varepsilon$, so $\|T_n - T\| \leq
\varepsilon$.
\end{proof}

\begin{lemma}[Riesz]
\label{functional-analysis-lemma-riesz}
Let $X$ be a normed space, $Y \subset X$ a
proper closed subspace, and $a \in (0, 1)$.
Then there exists $x \in X$ with $\|x\|=1$
such that $d(x, Y) \geq a$.
\end{lemma}

\begin{theorem}
\label
{theorem-finite-dimensional-operators-%
continuous}
If $\dim X < \infty$, then every linear map
$T : X \to Y$ is bounded.
\end{theorem}

\begin{theorem}
\label{functional-analysis-theorem-unit-ball-compact-finite-dim}
A normed space $X$ is finite-dimensional if
and only if $B_X$ is compact.
\end{theorem}



\section{Hahn--Banach theorems}
\label{functional-analysis-section-hahn-banach}

\begin{definition}
\label{functional-analysis-definition-sublinear}
Let $E$ be a real vector space. A function
$p:E\to\mathbb{R}$ is {\it sublinear} if
$p(\lambda x)=\lambda p(x)$ for
$\lambda\geq0$ and $p(x+y)\leq p(x)+p(y)$ for
all $x,y \in E$.
\end{definition}

\begin{theorem}
[Hahn--Banach, real analytic form]
\label{functional-analysis-theorem-hahn-banach-real}
Let $E$ be a real vector space, let
$p:E\to\mathbb{R}$ be sublinear, let
$G\subset E$ be a linear subspace, and let
$g:G\to\mathbb{R}$ be linear with $g(x)\leq
p(x)$ for all $x\in G$. Then there is a
linear functional $f:E\to\mathbb{R}$
extending $g$ such that $f(x)\leq p(x)$ for
all $x\in E$.
\end{theorem}

\begin{theorem}
[Hahn--Banach, complex analytic form]
\label{functional-analysis-theorem-hahn-banach-complex}
Let $E$ be a complex vector space, let
$p:E\to\mathbb{R}$ be a seminorm, let
$G\subset E$ be a subspace, and let
$g:G\to\mathbb{C}$ be linear with $|g(x)|\leq
p(x)$ for all $x\in G$. Then there is a
linear functional $f:E\to\mathbb{C}$
extending $g$ such that $|f(x)|\leq p(x)$ for
all $x\in E$.
\end{theorem}

\begin{proof}
Let $g_1 = \operatorname{Re} g$. Then $g_1$
is a real linear functional on $G$ and
$g_1(x) \leq |g(x)| \leq p(x)$. By the real
Hahn-Banach theorem, there is a real linear
extension $f_1$ of $g_1$ to $E$ such that
$f_1(x) \leq p(x)$ for all $x \in E$. Define
$f(x) = f_1(x) - i f_1(ix)$. One checks that
$f$ is complex linear, extends $g$, and
satisfies $|f(x)| \leq p(x)$.
\end{proof}

\begin{proposition}
\label
{proposition-hahn-banach-norm-preserving-%
extension}
Let $X$ be a normed space, let $Y\subset X$
be a linear subspace, and let $f\in Y^*$.
Then there is $F\in X^*$ such that $F|_Y=f$
and $\|F\|=\|f\|$.
\end{proposition}

\begin{proposition}
\label{functional-analysis-proposition-norming-functional}
Let $X$ be a normed space and let $x_0 \in X
\setminus \{0\}$. There exists $f \in X^*$
such that $\|f\|=1$ and $f(x_0) = \|x_0\|$.
\end{proposition}

\begin{proof}
Let $Y = \operatorname{span}\{x_0\}$ and
define $g(\lambda x_0) = \lambda \|x_0\|$.
Then $\|g\|=1$ on $Y$. Apply Proposition
\ref{proposition-hahn-banach-norm-preserving-%
extension}.
\end{proof}

\begin{proposition}
\label{functional-analysis-proposition-dual-separates-points}
The dual space $X^*$ separates points of $X$:
if $x \neq y$, there exists $f \in X^*$ such
that $f(x) \neq f(y)$.
\end{proposition}

\begin{theorem}
[Hahn--Banach, first geometric form]
\label{functional-analysis-theorem-hahn-banach-geometric-1}
Let $X$ be a normed real vector space and let
$A, B \subset X$ be nonempty disjoint convex
sets. If $A$ is open, there exist $f \in X^*
\setminus \{0\}$ and $\alpha \in \mathbb{R}$
such that $f(a) < \alpha \leq f(b)$ for all
$a \in A, b \in B$.
\end{theorem}

\begin{theorem}
[Hahn--Banach, second geometric form]
\label{functional-analysis-theorem-hahn-banach-geometric-2}
Let $X$ be a normed real vector space and let
$A, B \subset X$ be nonempty disjoint convex
sets. If $A$ is closed and $B$ is compact,
there exist $f \in X^* \setminus \{0\}$ and
$\alpha_1, \alpha_2 \in \mathbb{R}$ such that
$$
f(a) \leq \alpha_1 < \alpha_2 \leq f(b)
$$
for all $a \in A, b \in B$.
\end{theorem}


\section{Weak and weak star topologies}
\label{functional-analysis-section-weak-topologies}

\begin{definition}
\label{functional-analysis-definition-weak-topology}
Let $X$ be a normed space. The {\it weak
topology} on $X$, denoted $\sigma(X,X^*)$, is
the coarsest topology on $X$ such that every
$f \in X^*$ is continuous.
\end{definition}

\begin{definition}
\label{functional-analysis-definition-weak-star-topology}
Let $X$ be a normed space. The {\it weak star
topology} on $X^*$, denoted $\sigma(X^*,X)$,
is the coarsest topology on $X^*$ such that
for every $x \in X$, the evaluation map
$\text{ev}_x : X^* \to \mathbb{K}$ given by
$\text{ev}_x(f) = f(x)$ is continuous.
\end{definition}

\begin{proposition}
\label{functional-analysis-proposition-weak-vs-norm-topology}
The weak topology is coarser than the norm
topology. They coincide if and only if $\dim
X < \infty$.
\end{proposition}

\begin{theorem}[Banach--Alaoglu]
\label{functional-analysis-theorem-banach-alaoglu}
The closed unit ball $B_{X^*}$ of the dual
space $X^*$ is compact in the weak star
topology $\sigma(X^*,X)$.
\end{theorem}

\begin{proof}
For each $x \in X$, let $D_x = \{ \lambda \in
\mathbb{K} : |\lambda| \leq \|x\| \}$. By
Tychonoff's theorem, the product space $D =
\prod_{x \in X} D_x$ is compact in the
product topology. We can view $B_{X^*}$ as a
subset of $D$ via the map $\Phi : f \mapsto
(f(x))_{x \in X}$. The topology on $D$
restricted to $\Phi(B_{X^*})$ is precisely
the weak star topology. One checks that
$\Phi(B_{X^*})$ is a closed subset of $D$ by
verifying that the limit of a net of linear
functionals is linear. Thus $B_{X^*}$ is
compact.
\end{proof}

\begin{theorem}[Goldstine]
\label{functional-analysis-theorem-goldstine}
Let $X$ be a normed space and $J : X \to
X^{**}$ the canonical embedding. Then
$J(B_X)$ is weak star dense in $B_{X^{**}}$.
\end{theorem}

\begin{proof}
Suppose $J(B_X)$ is not weak star dense.
Since it is convex, by the Hahn-Banach
separation theorem (geometric form), there
exists a weak star continuous linear
functional on $X^{**}$ that separates a point
in $B_{X^{**}}$ from $J(B_X)$. Any weak star
continuous functional on $X^{**}$ is of the
form $\text{ev}_f$ for some $f \in X^*$. Thus
there exists $f \in X^*$ and $\alpha \in
\mathbb{R}$ such that $\operatorname{Re}
\xi(f) > \alpha$ for some $\xi \in
B_{X^{**}}$ and $\operatorname{Re} Jx(f) \leq
\alpha$ for all $x \in B_X$. But $\sup_{x \in
B_X} \operatorname{Re} f(x) = \|f\|$, so
$\alpha \geq \|f\|$. On the other hand,
$|\xi(f)| \leq \|\xi\| \|f\| \leq \|f\|$, so
$\alpha < \|f\|$, a contradiction.
\end{proof}

\begin{definition}
\label{functional-analysis-definition-reflexive-space}
A Banach space $X$ is {\it reflexive} if the
canonical embedding $J : X \to X^{**}$ is
surjective.
\end{definition}

\begin{theorem}[Kakutani]
\label{functional-analysis-theorem-kakutani-reflexivity}
A Banach space $X$ is reflexive if and only
if its closed unit ball $B_X$ is compact in
the weak topology.
\end{theorem}

\begin{theorem}[Eberlein--\v{S}mulyan]
\label{functional-analysis-theorem-eberlein-smulyan}
Let $A$ be a subset of a Banach space $X$.
The following are equivalent:
\begin{enumerate}
\item $A$ is relatively weakly compact.
\item $A$ is relatively weakly sequentially
compact.
\end{enumerate}
\end{theorem}


\section
{Baire category and Banach--Steinhaus}
\label{functional-analysis-section-baire-banach-steinhaus}

\begin{theorem}[Baire category theorem]
\label{functional-analysis-theorem-baire}
Let $M$ be a non-empty complete metric space.
\begin{enumerate}
\item If $\{U_n\}_{n=1}^\infty$ is a sequence
of open dense subsets of $M$, then
$\bigcap_{n=1}^\infty U_n$ is dense in $M$.
\item $M$ is not a countable union of nowhere
dense sets.
\end{enumerate}
\end{theorem}

\begin{proof}
(1) Let $W \subseteq M$ be a non-empty open
set. Since $U_1$ is dense, there exists $x_1
\in W \cap U_1$ and $r_1 \in (0,1)$ such that
$\overline{B(x_1, r_1)} \subseteq W \cap
U_1$. Proceeding inductively, given $x_n,
r_n$, there exists $x_{n+1}, r_{n+1}$ such
that $\overline{B(x_{n+1}, r_{n+1})}
\subseteq B(x_n, r_n) \cap U_{n+1}$ and
$r_{n+1} < r_n / 2$. The sequence $(x_n)$ is
Cauchy, and its limit $x$ belongs to $\bigcap
\overline{B(x_n, r_n)}$, which is contained
in $W \cap (\bigcap U_n)$.
\end{proof}

\begin{theorem}
[Banach--Steinhaus / Uniform boundedness
principle]
\label{functional-analysis-theorem-banach-steinhaus}
Let $X$ be a Banach space, $Y$ a normed
space, and let $\mathcal F \subset \mathcal
L(X,Y)$. If $\sup_{T \in \mathcal F} \|Tx\| <
\infty$ for each $x \in X$, then $\sup_{T \in
\mathcal F} \|T\| < \infty$.
\end{theorem}

\begin{proof}
For each $m \in \mathbb{N}$, let $X_m = \{ x
\in X : \|Tx\| \leq m \text{ for all } T \in
\mathcal F \}$. Each $X_m = \bigcap_{T \in
\mathcal F} \{x : \|Tx\| \leq m\}$ is closed.
By assumption, $X = \bigcup_{m=1}^\infty
X_m$. Since $X$ is Banach, by Baire's
theorem, some $X_{m_0}$ has non-empty
interior. Let $B(y, r) \subseteq X_{m_0}$.
For any $x \in X$ with $\|x\| \leq 1$, $y+rx
\in X_{m_0}$, so $\|T(y+rx)\| \leq m_0$. Then
$\|Trx\| \leq \|T(y+rx)\| + \|Ty\| \leq
2m_0$, so $\|Tx\| \leq 2m_0/r$. Thus $\|T\|
\leq 2m_0/r$ for all $T \in \mathcal F$.
\end{proof}

\begin{example}
\label{functional-analysis-example-completeness-necessary-bs}
Let $X = c_{00}$ with the sup-norm. Define
$f_n \in X^*$ by $f_n(x) = n x_n$. For each
$x \in X$, $f_n(x) = 0$ for large $n$, so
$\sup_n |f_n(x)| < \infty$. However, $\|f_n\|
= n \to \infty$. Thus completeness of $X$ is
essential.
\end{example}

\begin{lemma}
\label{functional-analysis-lemma-pointwise-limit-linear}
Let $X$ be a Banach space and $Y$ a normed
space. If $(T_n) \subset \mathcal L(X,Y)$ is
a sequence such that $Tx = \lim_n T_n x$
exists for all $x \in X$, then $T \in
\mathcal L(X,Y)$ and $\|T\| \leq \liminf_n
\|T_n\|$.
\end{lemma}

\section{Krein--Milman theorem}
\label{functional-analysis-section-krein-milman}

\begin{definition}
\label{functional-analysis-definition-extreme-point}
Let $K$ be a convex subset of a vector space.
\begin{enumerate}
\item A non-empty subset $L \subseteq K$ is
an {\it extreme set} of $K$ if for any $x, y
\in K$ and $\lambda \in (0,1)$, $\lambda x +
(1-\lambda)y \in L$ implies $x, y \in L$.
\item A point $x \in K$ is an {\it extreme
point} if $\{x\}$ is an extreme set. The set
of extreme points is denoted $E(K)$.
\end{enumerate}
\end{definition}

\begin{example}
\label{functional-analysis-example-extreme-points}
\begin{enumerate}
\item $E(B_{c_0}) = \emptyset$.
\item $E(B_{\ell_\infty}) = \{ (x_n) : x_n
\in \{-1, 1\} \text{ for all } n \}$.
\item $E(B_{\ell_1}) = \{ \pm e_n : n \in
\mathbb{N} \}$.
\item For $1 < p < \infty$, $E(B_{\ell_p}) =
S_{\ell_p}$ (the unit sphere).
\end{enumerate}
\end{example}

\begin{theorem}[Krein--Milman]
\label{functional-analysis-theorem-krein-milman}
Let $X$ be a locally convex topological
vector space and $K \subseteq X$ a non-empty
compact convex set. Then $K$ is the closed
convex hull of its extreme points: $K =
\overline{\operatorname{conv}}(E(K))$.
\end{theorem}

\begin{lemma}
\label{functional-analysis-lemma-c0-not-dual}
Since $E(B_{c_0}) = \emptyset$, $c_0$ is not
isometrically isomorphic to the dual of any
Banach space (if it were, $B_{c_0}$ would be
weak* compact and would have extreme points
by Krein-Milman).
\end{lemma}

\section
{Open mapping and closed graph theorems}
\label{functional-analysis-section-open-mapping-closed-graph}

\begin{theorem}[Open mapping theorem]
\label{functional-analysis-theorem-open-mapping}
Let $X, Y$ be Banach spaces and let $T \in
\mathcal L(X,Y)$ be surjective. Then $T$ is
an open map.
\end{theorem}

\begin{proof}[Proof sketch]
One shows that $0 \in
\operatorname{int}(\overline{T(B_X)})$ using
Baire's theorem (since $Y = \bigcup_n T(n
B_X)$). Then, using the completeness of $X$,
one proves $\overline{T(B_X)} \subseteq T(2
B_X)$ by constructing a convergent series in
$X$ for any $y \in \overline{T(B_X)}$.
\end{proof}

\begin{proposition}[Bounded inverse theorem]
\label{functional-analysis-proposition-bounded-inverse}
If $T \in \mathcal L(X,Y)$ is a bijection
between Banach spaces, then $T^{-1}$ is
bounded.
\end{proposition}

\begin{theorem}[Closed graph theorem]
\label{functional-analysis-theorem-closed-graph}
Let $X, Y$ be Banach spaces and let $T : X
\to Y$ be a linear map. If the graph
$\Gamma(T) = \{ (x, Tx) : x \in X \}$ is
closed in $X \times Y$, then $T$ is bounded.
\end{theorem}

\begin{proof}
Consider the projection $\pi_1 : \Gamma(T)
\to X$ given by $(x, Tx) \mapsto x$.
$\Gamma(T)$ is a Banach space as it is a
closed subspace of $X \times Y$. $\pi_1$ is a
bounded linear bijection. By the Bounded
Inverse Theorem, $\pi_1^{-1}$ is bounded.
Since $T = \pi_2 \circ \pi_1^{-1}$ (where
$\pi_2$ is the projection on $Y$), $T$ is
bounded.
\end{proof}

\section{Complemented subspaces}
\label{functional-analysis-section-complemented-subspaces}

\begin{definition}
\label{functional-analysis-definition-complemented-subspace}
A closed subspace $Y \subseteq X$ is {\it
complemented} if there exists a closed
subspace $Z \subseteq X$ such that $X = Y
\oplus Z$.
\end{definition}

\begin{proposition}
\label{functional-analysis-proposition-finite-codim-complemented}
If $Y$ is a closed subspace of a Banach space
$X$ and $\operatorname{codim} Y < \infty$,
then $Y$ is complemented.
\end{proposition}

\begin{theorem}[Phillips]
\label{functional-analysis-theorem-c0-not-complemented}
The subspace $c_0$ is not complemented in
$\ell_\infty$.
\end{theorem}


\begin{definition}
\label{functional-analysis-definition-compact-operator}
Let $X, Y$ be normed spaces. $T \in \mathcal
L(X,Y)$ is {\it compact} if $T(B_X)$ has
compact closure in $Y$. The set of compact
operators is denoted $\mathcal K(X,Y)$.
\end{definition}

\begin{proposition}
\label{functional-analysis-proposition-compact-operator-ideal}
$\mathcal K(X)$ is a closed two-sided ideal
in $\mathcal L(X)$.
\end{proposition}

\begin{theorem}[Schauder]
\label{functional-analysis-theorem-schauder-compact-adjoint}
An operator $T \in \mathcal L(X,Y)$ is
compact if and only if its adjoint $T^* \in
\mathcal L(Y^*, X^*)$ is compact.
\end{theorem}

\begin{definition}
\label{functional-analysis-definition-fredholm-operator}
$T \in \mathcal L(X,Y)$ is a {\it Fredholm
operator} if $\dim \ker T < \infty$ and
$\operatorname{codim} \operatorname{Im} T <
\infty$. The {\it index} of $T$ is
$$
\operatorname{ind}(T) = \dim \ker T -
\operatorname{codim} \operatorname{Im} T.
$$
\end{definition}

\begin{theorem}[Atkinson]
\label{functional-analysis-theorem-atkinson}
Let $X$ be an infinite-dimensional Banach
space. $T \in \mathcal L(X)$ is Fredholm if
and only if its image in the Calkin algebra
$\mathcal L(X) / \mathcal K(X)$ is
invertible.
\end{theorem}

\begin{definition}
\label{functional-analysis-definition-adjoint-banach}
For $T \in \mathcal L(X,Y)$, the {\it
adjoint} $T^* : Y^* \to X^*$ is defined by
$(T^*f)(x) = f(Tx)$.
\end{definition}

\begin{proposition}
\label{functional-analysis-proposition-adjoint-properties}
The map $T \mapsto T^*$ is a linear isometry
from $\mathcal L(X,Y)$ to $\mathcal L(Y^*,
X^*)$. Moreover, $(ST)^* = T^*S^*$.
\end{proposition}


\section{Universal Banach spaces}
\label{functional-analysis-section-universal-banach-spaces}

\begin{theorem}
\label{functional-analysis-theorem-separable-embeds-linfty}
Every separable normed space $X$ is
isometrically isomorphic to a subspace of
$\ell_\infty$.
\end{theorem}

\begin{proof}
Let $\{x_n\}_{n=1}^\infty$ be a dense subset
of the unit sphere $S_X$. By Hahn-Banach,
there exist $f_n \in S_{X^*}$ such that
$f_n(x_n) = 1$. Define $T : X \to
\ell_\infty$ by $Tx = (f_n(x))_{n=1}^\infty$.
Since $|f_n(x)| \leq \|x\|$, $T$ is
well-defined and $\|Tx\| \leq \|x\|$. For any
$\varepsilon > 0$, choose $n$ such that $\|x
- \|x\|x_n\| < \varepsilon$. Then $\|x\| =
f_n(\|x\|x_n) < f_n(x) + \varepsilon \leq
\|Tx\| + \varepsilon$. Thus $\|Tx\| = \|x\|$.
\end{proof}

\begin{theorem}[Banach--Mazur]
\label{functional-analysis-theorem-banach-mazur}
Every separable Banach space is isometrically
isomorphic to a subspace of $C[0,1]$.
\end{theorem}

\begin{theorem}
[Pe\l{}czy\'{n}ski decomposition]
\label{functional-analysis-theorem-pelczynski-decomposition}
Let $X, Y$ be Banach spaces such that each is
isomorphic to a complemented subspace of the
other.
\begin{enumerate}
\item If $X \cong X \oplus X$ and $Y \cong Y
\oplus Y$, then $X \cong Y$.
\item If $X \cong \ell_p(X)$ for some $1 \leq
p \leq \infty$, then $X \cong Y$.
\end{enumerate}
\end{theorem}

\begin{proof}
(1) Let $X \cong Y \oplus E$ and $Y \cong X
\oplus F$. Then $X \cong Y \oplus E \cong (Y
\oplus Y) \oplus E \cong Y \oplus (Y \oplus
E) \cong Y \oplus X$. Symmetrically, $Y \cong
X \oplus Y$. Thus $X \cong Y$.
\end{proof}


\section{Hilbert spaces}
\label{functional-analysis-section-hilbert-spaces}

\begin{definition}
\label{functional-analysis-definition-inner-product-space}
An {\it inner product space} over
$\mathbb{K}$ is a vector space $H$ with a
function $\langle\cdot,\cdot\rangle:H\times
H\to\mathbb{K}$ which is linear in the first
variable, conjugate-linear in the second
variable, positive definite, and Hermitian:
$$
\langle x,y\rangle=\overline{\langle
y,x\rangle}.
$$
The associated norm is $\|x\|=\langle
x,x\rangle^{1/2}$.
\end{definition}

\begin{proposition}[Pythagoras]
\label{functional-analysis-proposition-pythagoras}
If $x\perp y$, then
$$
\|x+y\|^2=\|x\|^2+\|y\|^2.
$$
\end{proposition}

\begin{theorem}[Cauchy--Schwarz]
\label{functional-analysis-theorem-cauchy-schwarz}
For all $x,y\in H$,
$$
|\langle x,y\rangle|\leq \|x\|\|y\|.
$$
Equality holds if and only if $x,y$ are
linearly dependent.
\end{theorem}

\begin{definition}
\label{functional-analysis-definition-hilbert-space}
A {\it Hilbert space} is an inner product
space which is Banach for the norm induced by
the inner product.
\end{definition}

\begin{theorem}[Parallelogram law]
\label{functional-analysis-theorem-parallelogram-law}
A norm on a vector space comes from an inner
product if and only if
$$
\|x+y\|^2+\|x-y\|^2=2\|x\|^2+2\|y\|^2
$$
for all $x,y$.
\end{theorem}

\begin{definition}
\label{functional-analysis-definition-orthonormal-set}
A subset $S\subset H$ is {\it orthonormal} if
$\|s\|=1$ for all $s\in S$ and $\langle
s,t\rangle=0$ for $s\neq t$.
\end{definition}

\begin{theorem}[Bessel inequality]
\label{functional-analysis-theorem-bessel}
If $S\subset H$ is orthonormal, then for
every $x\in H$,
$
\sum_{s\in S}|\langle x,s\rangle|^2\leq
\|x\|^2.
$
\end{theorem}

\begin{proof}
It is enough to prove that for every finite
subset $F\subset S$,
$
\sum_{s\in F}|\langle x,s\rangle|^2
\leq \|x\|^2.
$
Indeed, for a nonnegative family
$(a_s)_{s\in S}$, if every finite partial sum
is bounded above by $A$, define
$
a=\sup_{F\subset S,\ F\text{ finite}}
\sum_{s\in F}a_s\leq A.
$
Given $\varepsilon>0$, choose a finite
$F_0\subset S$ such that
$
a-\varepsilon<
\sum_{s\in F_0}a_s\leq a.
$
If $F\supset F_0$ is finite, then
$
a-\varepsilon<
\sum_{s\in F_0}a_s
\leq
\sum_{s\in F}a_s
\leq a,
$
so the finite partial sums converge to $a$.

Now fix a finite $F\subset S$. Since norms
are nonnegative,
$
0\leq
\left\|
x-\sum_{s\in F}\langle x,s\rangle s
\right\|^2.
$
Expanding the inner product and using
orthonormality,
\begin{align*}
\left\|
x-\sum_{s\in F}\langle x,s\rangle s
\right\|^2
&=
\|x\|^2
-2\sum_{s\in F}|\langle x,s\rangle|^2
+\sum_{s\in F}|\langle x,s\rangle|^2\\
&=
\|x\|^2-
\sum_{s\in F}|\langle x,s\rangle|^2.
\end{align*}
Therefore
$
\sum_{s\in F}|\langle x,s\rangle|^2
\leq \|x\|^2
$
for every finite $F\subset S$, and hence
$
\sum_{s\in S}|\langle x,s\rangle|^2
\leq \|x\|^2.
$
\end{proof}

\begin{theorem}[Parseval]
\label{functional-analysis-theorem-parseval}
Let $H$ be a Hilbert space and let $S\subset
H$ be orthonormal. The following are
equivalent.
\begin{enumerate}
\item $S$ is maximal orthonormal.
\item The closed linear span of $S$ is $H$.
\item For every $x\in H$,
$
x=\sum_{s\in S}\langle x,s\rangle s.
$
\item For every $x\in H$,
$
\|x\|^2=\sum_{s\in S}|\langle x,s\rangle|^2.
$
\end{enumerate}
\end{theorem}

\begin{proof}
First, the series
$
\sum_{s\in S}\langle x,s\rangle s
$
exists for every $x\in H$. By the Cauchy
criterion, it is enough to show that for
every $\varepsilon>0$ there is a finite
$S_0\subset S$ such that for every finite
$S_1\subset S\setminus S_0$,
$
\left\|\sum_{s\in S_1}
\langle x,s\rangle s\right\|<\varepsilon.
$
By Pythagoras,
$
\left\|\sum_{s\in S_1}
\langle x,s\rangle s\right\|^2
=
\sum_{s\in S_1}|\langle x,s\rangle|^2.
$
By Bessel's inequality, the family
$(|\langle x,s\rangle|^2)_{s\in S}$ is
summable, so the Cauchy criterion holds.

Suppose now that $S$ is maximal. We have
$
x-\sum_{s\in S}\langle x,s\rangle s
\perp S.
$
Hence maximality gives
$
x=\sum_{s\in S}\langle x,s\rangle s.
$

If $S$ is not maximal, there exists
$y_0\in H$ with $\|y_0\|=1$ and
$y_0\perp S$. Hence
$
y_0\perp\overline{\operatorname{span}}(S),
$
so
$y_0\notin\overline{\operatorname{span}}(S)$.
This proves the equivalence with the
closed-span condition. It also gives
$
1=\|y_0\|^2
\neq
\sum_{s\in S}|\langle y_0,s\rangle|^2=0,
$
so Parseval's equality fails when $S$ is not
maximal.

Finally, if
$
x=\sum_{s\in S}\langle x,s\rangle s,
$
then Pythagoras gives
$
\|x\|^2=
\sum_{s\in S}|\langle x,s\rangle|^2.
$
\end{proof}



\begin{proposition}
\label{functional-analysis-proposition-hilbert-is-l2-set}
Every Hilbert space is isometrically
isomorphic to $\ell^2(S)$ for some set $S$.
\end{proposition}

\begin{theorem}[Riesz representation]
\label{functional-analysis-theorem-riesz-representation-hilbert}
Let $H$ be a Hilbert space. For every $f\in
H^*$ there is a unique $y\in H$ such that
$$
f(x)=\langle x,y\rangle
$$
for all $x\in H$. Moreover $\|f\|=\|y\|$.
\end{theorem}

\begin{theorem}[Projection theorem]
\label{functional-analysis-theorem-projection-hilbert}
Let $H$ be a Hilbert space and let $M\subset
H$ be a closed linear subspace. For every
$x\in H$ there is a unique $m\in M$ such that
$$
\|x-m\|=d(x,M).
$$
Moreover $x-m\perp M$.
\end{theorem}

\begin{definition}
\label{functional-analysis-definition-orthogonal-projection}
The map $P_M:H\to M$ sending $x$ to the
unique nearest point of $M$ is called the
{\it orthogonal projection} onto $M$.
\end{definition}

\begin{proposition}
\label{functional-analysis-proposition-orthogonal-projection}
The orthogonal projection $P_M$ is linear,
bounded, satisfies $P_M^2=P_M$, and has norm
$1$ if $M\neq0$.
\end{proposition}

\begin{proposition}
\label{functional-analysis-example-fourier-orthonormal-system}
\begin{enumerate}
\item In $L_2(-\pi,\pi)$, the collection (or any subset thereof)
\[x_n=\frac{1}{\sqrt{2\pi}}e^{int},\qquad n=0,\pm1,\ldots\]
is an orthonormal set of vectors.
\begin{proof}
For any $n\in\mathbb{Z}$,
\[\int_{-\pi}^\pi x_n\overline{x}_ndt=\frac{1}{2\pi}\int_{-\pi}^\pi e^{int}\overline{e^{int}}dt=\frac{1}{2\pi}\int_{-\pi}^\pi \Vert e^{int}\Vert^2dt=1,\]
and if $m$ is another integer,
\[\int_{-\pi}^\pi x_n\overline{x}_mdt=\frac{1}{2\pi}\int_{-\pi}^\pi e^{int}\overline{e^{imt}}dt=\frac{1}{2\pi}\int_{-\pi}^\pi e^{(n-m)it}dt=\frac{1}{2\pi}\left[\frac{e^{(n-m)it}}{(n-m)i}\right]_{-\pi}^\pi=0.\]
\end{proof}
\item If we restric out attention to only real-valued functions that are square-integrable on the interval $[-\pi,\pi]$, then the collection (or any subset thereof)
\begin{align*}
\frac{1}{\sqrt{2\pi}},&\frac{1}{\sqrt{\pi}}\cos t,\frac{1}{\sqrt{\pi}}\cos 2t,\ldots\\
&\frac{1}{\sqrt{\pi}}\sin t,\frac{1}{\sqrt{\pi}}\sin 2t,\ldots
\end{align*}
is an orthonormal set.
\end{enumerate}
\end{proposition}

\begin{proposition}
\label{functional-analysis-proposition-fourier-system-maximal}
The orthonormal system
$
S=
\left\{
\frac{e^{int}}{\sqrt{2\pi}}:n\in\mathbb Z
\right\}
$
is maximal in $L_2[0,2\pi]$.
\end{proposition}

\begin{proof}
We only prove maximality. Let
$f\in L_2[0,2\pi]$ satisfy
$
f\perp S.
$
We show that $f=0$.

First we replace $f$ by a continuous
function. Define
$
G(x)=\int_0^x f\,d\mu.
$
Then $G$ is absolutely continuous and
$G'=f$ almost everywhere. Hence, for every
$k\in\mathbb C$,
$
(G+k)'\perp S.
$
For $n\neq0$, integration by parts gives
$
\begin{aligned}
0
&=\int_0^{2\pi}(G+k)'e^{int}\,d\mu\\
&=G(2\pi)-G(0)
 -\int_0^{2\pi}(G+k)(in)e^{int}\,d\mu\\
&=-\int_0^{2\pi}(G+k)(in)e^{int}\,d\mu,
\end{aligned}
$
where $G(2\pi)-G(0)=\int_0^{2\pi}f\,d\mu=0$
because $f$ is orthogonal to the constant
function. Therefore
$
\int_0^{2\pi}(G+k)e^{int}\,d\mu=0
\qquad
\text{for every }n\in\mathbb Z\setminus
\{0\}.
$

Choose $k\in\mathbb C$ so that
$
H:=G+k
$
is also orthogonal to the constant function.
Then $H\perp S$.

Since $H$ is continuous, the
Stone--Weierstrass theorem gives, for every
$\varepsilon>0$, a trigonometric polynomial
$T\in\operatorname{span}(S)$ such that
$
\sup_{t\in[0,2\pi]}|H(t)-T(t)|
<\varepsilon.
$
Thus
$
H\in\overline{\operatorname{span}(S)}.
$
But $H\perp S$, hence
$H\perp\overline{\operatorname{span}(S)}$.
Therefore $H=0$. Since
$
f=G'=H'
$
almost everywhere, we conclude that $f=0$.
Thus $S$ is maximal.
\end{proof}

\section{Banach algebras}
\label{functional-analysis-section-banach-algebras}

\begin{definition}
\label{functional-analysis-definition-banach-algebra}
A {\it Banach algebra} $A$ is a Banach space
equipped with an associative multiplication
such that $\|ab\| \leq \|a\| \|b\|$ for all
$a, b \in A$. If there exists $1 \in A$ such
that $a1 = 1a = a$, $A$ is a {\it unital
Banach algebra}.
\end{definition}

\begin{example}
\label{functional-analysis-example-banach-algebras}
\begin{enumerate}
\item $C(K)$ for a compact Hausdorff space
$K$ with point-wise multiplication.
\item $\mathcal{L}(X)$ for a Banach space $X$
with composition.
\item $L^1(\mathbb{R})$ with the convolution
product $(f \ast g)(t) = \int f(t-s)g(s) ds$.
\item $K(X)$, the space of compact operators,
which is a closed ideal in $\mathcal{L}(X)$.
\end{enumerate}
\end{example}

\begin{proposition}
\label{functional-analysis-proposition-algebra-quotient}
If $I \subseteq A$ is a closed two-sided
ideal, then $A/I$ is a Banach algebra with
the quotient norm.
\end{proposition}

\begin{definition}
\label{functional-analysis-definition-spectrum-algebra}
Let $A$ be a unital Banach algebra and $a \in
A$. The {\it spectrum} of $a$ is
$$
\sigma(a) = \{ \lambda \in \mathbb{C} : a -
\lambda 1 \text{ is not invertible in } A \}.
$$
The {\it spectral radius} of $a$ is $r(a) =
\sup \{ |\lambda| : \lambda \in \sigma(a)
\}$.
\end{definition}

\begin{theorem}[Gelfand]
\label{functional-analysis-theorem-spectrum-nonempty}
For any $a \in A$, $\sigma(a)$ is a non-empty
compact subset of $\mathbb{C}$. Moreover,
$\sigma(a) \subseteq \{ z \in \mathbb{C} :
|z| \leq \|a\| \}$.
\end{theorem}

\begin{theorem}[Spectral radius formula]
\label{functional-analysis-theorem-spectral-radius-formula}
For any $a \in A$, $r(a) = \lim_{n \to
\infty} \|a^n\|^{1/n}$.
\end{theorem}

\section
{Fredholm operators and the Calkin algebra}
\label{functional-analysis-section-fredholm-theory}

\begin{definition}
\label{functional-analysis-definition-fredholm-operator-duplicate-2}
A bounded linear operator $T \in
\mathcal{L}(X, Y)$ is a {\it Fredholm
operator} if $\dim \ker T < \infty$ and
$\operatorname{codim} \operatorname{Im} T <
\infty$. The {\it index} of $T$ is
$$
\operatorname{ind}(T) = \dim \ker T -
\operatorname{codim} \operatorname{Im} T.
$$
\end{definition}

\begin{theorem}[Atkinson]
\label{functional-analysis-theorem-atkinson-calkin}
Let $X$ be an infinite-dimensional Banach
space. Let $\pi : \mathcal{L}(X) \to
\mathcal{L}(X)/K(X)$ be the natural
projection into the Calkin algebra. Then $T
\in \mathcal{L}(X)$ is Fredholm if and only
if $\pi(T)$ is invertible in
$\mathcal{L}(X)/K(X)$.
\end{theorem}

\begin{proposition}
\label{functional-analysis-proposition-index-stability}
\begin{enumerate}
\item The set of Fredholm operators $\Phi(X)$
is open in $\mathcal{L}(X)$.
\item The index map $\operatorname{ind} :
\Phi(X) \to \mathbb{Z}$ is continuous (and
thus locally constant).
\item If $T$ is Fredholm and $S$ is compact,
then $T+S$ is Fredholm and
$\operatorname{ind}(T+S) =
\operatorname{ind}(T)$.
\end{enumerate}
\end{proposition}

\section
{Spectral theory for compact operators}
\label{functional-analysis-section-spectral-compact}

\begin{theorem}
\label{functional-analysis-theorem-spectral-compact-ops}
Let $X$ be a Banach space and $T \in
\mathcal{K}(X)$ a compact operator.
\begin{enumerate}
\item $\sigma(T)$ is at most countable and
has no limit point except possibly $0$.
\item Every $\lambda \in \sigma(T) \setminus
\{0\}$ is an eigenvalue of finite
multiplicity.
\item $\operatorname{ind}(T - \lambda I) = 0$
for all $\lambda \neq 0$.
\end{enumerate}
\end{theorem}

\begin{theorem}
[Spectral theorem for compact self-adjoint
operators]
\label{functional-analysis-theorem-spectral-compact-selfadjoint}
Let $H$ be a Hilbert space and $T \in
\mathcal{K}(H)$ be self-adjoint ($T = T^*$).
There exists a maximal orthonormal set $\{
e_n \}$ of eigenvectors of $T$ with real
eigenvalues $\{ \lambda_n \}$. For any $x \in
H$:
$$
Tx = \sum_n \lambda_n \langle x, e_n \rangle
e_n.
$$
\end{theorem}

\section
{Special topics and advanced examples}
\label{functional-analysis-section-advanced-topics}

\begin{example}[Fourier Transform]
\label{functional-analysis-example-fourier-l1}
For $f \in L^1(\mathbb{R})$, the Fourier
transform is $\widehat{f}(\xi) =
\int_\mathbb{R} e^{-2\pi i \xi x} f(x) dx$.
The map $\mathcal{F} : L^1(\mathbb{R}) \to
C_0(\mathbb{R})$ is a bounded linear operator
with $\|\mathcal{F}\| \leq 1$. It is
injective but not surjective.
\end{example}

\begin{example}
[Weak Derivatives and Sobolev Spaces]
\label{functional-analysis-example-sobolev-w1p}
Let $f \in L^1_{\operatorname{loc}}(0,1)$. A
function $g$ is the {\it weak derivative} of
$f$ if
$$
\int_0^1 f(x)\varphi'(x) dx = -\int_0^1
g(x)\varphi(x) dx \quad \forall \varphi \in
C_c^1(0,1).
$$
The Sobolev space $W^{1,p}(0,1) = \{ f \in
L^p(0,1) : f' \in L^p(0,1) \}$ is a Banach
space with norm $\|f\| = \|f\|_p + \|f'\|_p$.
It is reflexive for $1 < p < \infty$.
\end{example}

\begin{example}[Orlicz Spaces]
\label{functional-analysis-example-orlicz-space}
Let $\varphi$ be a continuous, convex,
increasing function on $[0, \infty)$ with
$\varphi(0)=0$. The Orlicz space
$L^\varphi(\mathbb{R})$ consists of
measurable functions such that $\int
\varphi(|f|/C) dx < \infty$ for some $C > 0$.
The Luxemburg norm is $\|f\|_\varphi = \inf
\{ C > 0 : \int \varphi(|f|/C) dx \leq 1 \}$.
\end{example}

\section{Further exercises}
\label{functional-analysis-section-functional-analysis-exercises}

\begin{exercise}
\label{functional-analysis-exercise-dual-c0-l1}
Prove that $c_0^* \cong \ell_1$ and $\ell_1^*
\cong \ell_\infty$.
\end{exercise}

\begin{exercise}
\label{functional-analysis-exercise-weak-vs-strong}
Show that in an infinite-dimensional space,
the weak topology is strictly coarser than
the norm topology. Find a sequence that
converges weakly but not in norm.
\end{exercise}

\begin{exercise}
\label{functional-analysis-exercise-hilbert-projection}
Let $M$ be a closed subspace of a Hilbert
space $H$. Show that $H = M \oplus M^\perp$
and that the orthogonal projection $P_M$ is a
self-adjoint operator of norm 1.
\end{exercise}

% --- End functional-analysis.tex ---

% --- Begin pde.tex ---
\chapter{Partial Differential Equations}
\label{chapter-pde}





\section{Introduction}
\label{pde-section-introduction}

\noindent
These notes develop the theory of partial
differential equations.

Given the material elsewhere in the Heap
Project, the basic things to learn are:

\begin{enumerate}
\item The model elliptic, parabolic, and
hyperbolic equations: the Laplace, heat,
and wave equations.

\item How characteristic curves explain
first-order equations, transport, and the
formation of singularities.

\item The fundamental solutions and Green
functions of the model equations, together
with the maximum principle and mean-value
property.

\item Fourier-series and Fourier-transform
methods, separation of variables, and the
role of distributions.

\item Sobolev spaces, weak derivatives,
embeddings, compactness, and trace
theorems.

\item Weak formulations and the main
existence mechanisms: energy estimates,
the Lax--Milgram theorem, compactness, and
the calculus of variations.

\item Elliptic regularity: when a weak
solution becomes smooth, and how the
regularity of the data controls that of
the solution.

\item Energy and semigroup methods for
evolution equations, including existence,
uniqueness, stability, and long-time
behavior.

\item The distinction between classical,
weak, distributional, and viscosity
solutions, and why different equations
require different notions of solution.

\item How differential operators fit the
geometric language already developed in
the project: symbols, ellipticity, the
Laplace--Beltrami operator, and de Rham
and Hodge Laplacians.

\item The first nonlinear phenomena:
conservation laws, shocks, blow-up,
scaling, and critical exponents.
\end{enumerate}

\section{Some Problems in Hamiltonian
Geometry of PDEs}
\label{pde-section-hamiltonian-geometry-pdes}

\medskip\noindent
{\bf Speaker.}

Raffaele Vitolo,
Universit\`a del Salento.

\medskip\noindent
{\bf Seminar information.}

Semin\'ario das Sextas, PUC-Rio,
September 27, 2024. The talk was also
presented in the Geometric Structures on
Manifolds Seminar at IMPA on
December 5, 2024.

\medskip\noindent
{\bf Abstract.}

The Hamiltonian formulation of partial
differential equations is one of the
cornerstones of the theory of integrable
systems. It is modeled on finite-dimensional
Hamiltonian systems and has an intrinsic
geometric nature. The talk reviews
differential-geometric aspects of the
Hamiltonian theory of PDEs and discusses
projective-geometric properties of
integrable systems that have emerged in
recent years.

\medskip\noindent
{\bf Hamiltonian evolutionary equations.}

Consider an evolutionary system
$$
u_t^i=f^i(t,x,u^j,u_x^j,u_{xx}^j,\ldots).
$$
A local functional has the form
$$
\mathcal{H}=\int h\,dx.
$$
Its variational derivative is
$$
\frac{\delta\mathcal{H}}{\delta u^i}
=
\sum_\sigma(-D)_\sigma
\frac{\partial h}{\partial u^i_\sigma}.
$$

\begin{definition}
\label{pde-definition-hamiltonian-pde}
The system has a \textit{Hamiltonian
formulation} if there is a differential
operator $A$ and a local functional
$\mathcal{H}$ such that
$$
u_t^i=A^{ij}
\frac{\delta\mathcal{H}}{\delta u^j}.
$$
\end{definition}

The operator $A$ plays the role of a
Poisson tensor. It sends the differential
of a Hamiltonian functional to the
corresponding evolutionary vector field.

\begin{definition}
\label{pde-definition-hamiltonian-operator}
A differential operator $A$ is a
\textit{Hamiltonian operator} if
$$
\{\mathcal{F},\mathcal{G}\}_A
=
\int
\frac{\delta\mathcal{F}}{\delta u^i}
A^{ij}
\frac{\delta\mathcal{G}}{\delta u^j}
\,dx
$$
defines a Poisson bracket on local
functionals.
\end{definition}

Thus $A$ must be formally skew-adjoint and
must satisfy the differential analogue of
the Jacobi identity.

There is also a purely algebraic definition
of differential operators. Let $R$ be a
ring, let $P,Q$ be $R$-modules, and let
$\Delta:P\to Q$ be additive. For $a\in R$,
write
$$
\delta_a\Delta=a\Delta-\Delta a.
$$
The operator $\Delta$ has order at most
$k$ if
$$
\delta_{a_0}\delta_{a_1}\cdots
\delta_{a_k}\Delta=0
$$
for all $a_0,\ldots,a_k\in R$. This is the
Grothendieck definition of a differential
operator.

\medskip\noindent
{\bf The KdV equation.}

\begin{example}
\label{pde-example-kdv-bihamiltonian}
Consider the Korteweg--de Vries equation
$$
u_t=uu_x+u_{xxx}.
$$
It admits two Hamiltonian descriptions.
Set
$$
A_1=D_x
$$
and
$$
A_2=D_x^3+\frac{2}{3}uD_x
+\frac{1}{3}u_x.
$$
For the Hamiltonians
$$
\mathcal{H}_1=
\int\left(
\frac{u^3}{6}-\frac{u_x^2}{2}
\right)dx
$$
and
$$
\mathcal{H}_2=\int\frac{u^2}{2}\,dx,
$$
one has
$$
u_t=A_1
\frac{\delta\mathcal{H}_1}{\delta u}
=A_2
\frac{\delta\mathcal{H}_2}{\delta u}.
$$
\end{example}

This is the model for a bi-Hamiltonian
system.

\begin{definition}
\label{pde-definition-compatible-hamiltonian}
Two Hamiltonian operators $A_1,A_2$ are
\textit{compatible} if every linear
combination of them is again Hamiltonian.
A system admitting such a pair is
\textit{bi-Hamiltonian}.
\end{definition}

\medskip\noindent
{\bf The Magri scheme.}

A Hamiltonian operator sends conservation
laws to symmetries. Given compatible
Hamiltonian operators, the Magri recursion
seeks functionals $\mathcal{H}_r$ satisfying
$$
A_1
\frac{\delta\mathcal{H}_{r+1}}{\delta u}
=
A_2
\frac{\delta\mathcal{H}_r}{\delta u}.
$$
Under suitable hypotheses, this produces
an infinite hierarchy of commuting flows
and conserved Hamiltonians in involution:
$$
\{\mathcal{H}_r,\mathcal{H}_s\}=0.
$$

This motivates one notion of integrability
for PDEs: the existence of infinitely many
commuting symmetries or conservation laws.
There is no direct infinite-dimensional
version of the Liouville theorem that
settles the theory, but the hierarchy may
lead to suitable coordinates or to an
inverse-scattering construction.

\medskip\noindent
{\bf Dubrovin--Novikov operators.}

A first-order homogeneous operator has the
form
$$
A^{ij}=g^{ij}(u)D_x
+b^{ij}_k(u)u_x^k.
$$
The conditions for $A$ to be Hamiltonian
have a differential-geometric meaning.
Formal skew-adjointness forces $g^{ij}$ to
be symmetric and relates $b^{ij}_k$ to a
metric connection. The Jacobi identity
then forces that connection to be
torsion-free and the metric to be flat,
when the metric is nondegenerate.

Thus the classification of these
Hamiltonian operators becomes a problem
in the geometry of flat metrics. The
classification of compatible pairs is
more rigid and encodes whole hierarchies
of conservation laws.

\medskip\noindent
{\bf Projective classification.}

The talk described a classification program
using projective reciprocal
transformations. Such transformations can
preserve the order and Hamiltonian nature
of an operator while changing its
coefficients. One first fixes a
higher-order operator and then classifies
compatible lower-order operators modulo
this group.

In the cases discussed in the talk,
second-order homogeneous Hamiltonian
operators are related to linear line
congruences, while third-order operators
and compatible first-order operators lead
to line complexes. Pl\"ucker coordinates
place these objects inside
$$
\mathbb{P}(\Lambda^2V).
$$
Compatibility of Hamiltonian operators is
thereby translated into compatibility of
projective varieties.

The broad point is that many familiar
integrable systems carry algebraic
varieties that organize their Hamiltonian
structures. Conversely, suitable
projective data can be used to construct
integrable systems. The Kaup--Broer system
was mentioned as an example arising from
this mechanism.

\section{PDE Without Tears for Geometry
Students}
\label{pde-section-pde-without-tears}

\medskip\noindent
{\bf Speaker.}

Raffaele Vitolo,
Universit\`a del Salento.

\medskip\noindent
{\bf Seminar information.}

Semin\'ario das Sextas, PUC-Rio,
November 29, 2024.

The geometric study of differential
equations goes back to Lie and was put
into its modern form by Ehresmann. Its
basic idea is to regard derivatives as
coordinates on a larger space.

Let $N$ have dimension $n$ and let $M$
have dimension $m$. A field is a map
$$
u:U\longrightarrow M,
\qquad U\subset N.
$$
Locally it has independent coordinates
$x^1,\ldots,x^n$ and dependent
coordinates $u^1,\ldots,u^m$. It can
equivalently be regarded as the section
$$
s:N\longrightarrow N\times M,
\qquad s(x)=(x,u(x)).
$$

\medskip\noindent
{\bf Contact and jet spaces.}

\begin{definition}
\label{pde-definition-contact-order}
Two local sections $s_1$ and $s_2$ have
\textit{contact of order $r$} at $x$ if
their derivatives at $x$ agree through
order $r$.
\end{definition}

Thus contact of order zero means that the
sections have the same value. Contact of
order one means in addition that their
graphs have the same tangent space.

More generally, let $E$ be a manifold of
dimension $n+m$, and let $L_1,L_2\subset E$
be submanifolds of dimension $n$ through
$z$. After expressing both as graphs in a
common chart, one can compare their
derivatives. The resulting notion of
contact is independent of the chart.

\begin{definition}
\label{pde-definition-jet-space}
The \textit{$r$-jet} of $L$ at $z$ is its
equivalence class under contact of order
$r$, denoted by $[L]^r_z$. The
\textit{jet space} is
$$
J^r(E,n)=
\bigcup_{z\in E}
\left\{[L]^r_z\right\}.
$$
\end{definition}

There are natural projections
$$
\cdots\longrightarrow J^2(E,n)
\longrightarrow J^1(E,n)
\longrightarrow E.
$$
The first jet space records tangent
$n$-planes and is therefore a
Grassmannian bundle. Higher jet spaces
record higher-order Taylor data.

Suppose now that $E\to N$ is a bundle. A
local section has coordinates
$$
(x^\lambda,u^i),
$$
and its $r$-jet has coordinates
$$
(x^\lambda,u^i,u^i_\sigma),
\qquad |\sigma|\leq r.
$$
Here $\sigma$ is a multi-index and
$u^i_\sigma$ represents the corresponding
partial derivative of $u^i$.

\begin{definition}
\label{pde-definition-pde-jet-space}
A \textit{partial differential equation of
order $r$} is a submanifold
$$
\mathcal{E}\subset J^r(E,n).
$$
\end{definition}

A system given locally by equations
$$
F^a(x^\lambda,u^i,u^i_\sigma)=0
$$
defines such a submanifold whenever the
usual regularity conditions hold.

\begin{definition}
\label{pde-definition-prolongation-section}
The \textit{$r$-th prolongation} of a
section $s$ is
$$
j^rs:N\longrightarrow J^r(E,n),
\qquad x\longmapsto[s(N)]^r_{s(x)}.
$$
The section is a \textit{solution} of
$\mathcal{E}$ if
$$
j^rs(N)\subset\mathcal{E}.
$$
\end{definition}

\begin{example}
\label{pde-example-first-derivative-one}
For $n=m=1$, the equation
$$
u_x=1
$$
is a surface in $J^1(\mathbb{R}^2,1)$.
Its solutions are
$$
u(x)=x+c.
$$
Their first prolongations lie in that
surface.
\end{example}

\medskip\noindent
{\bf The inviscid Burgers equation.}

\begin{example}
\label{pde-example-inviscid-burgers}
The Riemann, Hopf, or inviscid Burgers
equation is
$$
u_t+uu_x=0.
$$
Along a characteristic curve one has
$$
\frac{d}{dt}u(t,x(t))=0,
\qquad
\frac{dx}{dt}=u.
$$
Consequently, if the initial condition is
$$
u(0,x)=f(x),
$$
then the solution is determined implicitly
by
$$
u(t,x)=f(x-tu(t,x)).
$$

Set
$$
\xi=x-tu,
\qquad
v(t,\xi)=u(t,x).
$$
Away from points where the coordinate
change is singular, Burgers' equation
becomes
$$
v_t=0.
$$
The singularity occurs when
$$
1+tf'(\xi)=0.
$$
Thus characteristics can intersect and a
smooth solution can develop a shock in
finite time.
\end{example}

This example suggests a general strategy:
find symmetries and conservation laws, and
then seek a change of variables that
simplifies or linearizes the equation.

\medskip\noindent
{\bf The Cartan distribution.}

The images of prolonged sections determine
distinguished tangent directions on a jet
space. These directions form the geometric
structure that must be preserved by any
transformation of differential equations.

\begin{definition}
\label{pde-definition-cartan-distribution}
The \textit{Cartan distribution} on
$J^r(E,n)$ is the distribution spanned by
tangent spaces to prolonged submanifolds,
together with the vertical directions of
the projection to $J^{r-1}(E,n)$.
\end{definition}

For a bundle $E\to N$, introduce the total
derivatives
$$
D_\lambda=
\frac{\partial}{\partial x^\lambda}
+
\sum_{i,|\sigma|<r}
u^i_{\sigma+\lambda}
\frac{\partial}{\partial u^i_\sigma}.
$$
The Cartan distribution is locally spanned
by the $D_\lambda$ and by
$$
\frac{\partial}{\partial u^i_\tau},
\qquad |\tau|=r.
$$
The total derivatives commute:
$$
[D_\lambda,D_\mu]=0.
$$
They differentiate functions on jet space
as if all jet coordinates arose from an
actual section.

\medskip\noindent
{\bf Contact transformations and
symmetries.}

\begin{definition}
\label{pde-definition-contact-transformation}
A \textit{contact transformation} of a jet
space is a local diffeomorphism preserving
the Cartan distribution.
\end{definition}

When there is more than one dependent
variable, the Lie--B\"acklund theorem says
that every contact transformation of a
finite-order jet space is the prolongation
of a transformation of $E$. With one
dependent variable, genuine contact
transformations already occur on
$J^1(E,n)$.

Contact transformations can sometimes
linearize nonlinear equations. A model
case is the Monge--Amp\`ere equation
$$
u_{xx}u_{yy}-u_{xy}^2=1,
$$
which may be studied using a Legendre
transformation.

\begin{definition}
\label{pde-definition-symmetry-pde}
A \textit{symmetry} of a differential
equation $\mathcal{E}$ is a contact
transformation preserving $\mathcal{E}$.
An infinitesimal symmetry is a vector field
tangent to $\mathcal{E}$ whose flow
preserves the restricted Cartan
distribution.
\end{definition}

The bracket of two infinitesimal symmetries
is again an infinitesimal symmetry. For an
ordinary differential equation, the
Lie--Bianchi theorem says that a sufficiently
large solvable Lie algebra of symmetries
allows the equation to be solved by
quadratures. For PDEs, symmetries can reduce
the number of variables or turn a PDE into
an ODE.

\medskip\noindent
{\bf Prolongations and linearization.}

Let
$$
X=\xi^\lambda
\frac{\partial}{\partial x^\lambda}
+\phi^i\frac{\partial}{\partial u^i}
$$
be a vector field on $E$. Its characteristic
is
$$
Q^i=\phi^i-u^i_\lambda\xi^\lambda.
$$
Its prolongation to the infinite jet space
can be written as
$$
X^{(\infty)}=
\xi^\lambda D_\lambda
+
\sum_{i,\sigma}D_\sigma(Q^i)
\frac{\partial}{\partial u^i_\sigma}.
$$
This separates the motion along the base
from the evolutionary, or vertical, part.

Suppose the equation is defined by
$$
F^a(x^\lambda,u^i,u^i_\sigma)=0.
$$
Its linearization is the differential
operator
$$
\ell_F(Q)^a=
\sum_{i,\sigma}
\frac{\partial F^a}
{\partial u^i_\sigma}D_\sigma(Q^i).
$$
A generalized symmetry has characteristic
$Q$ satisfying
$$
\ell_F(Q)=0
$$
on the equation and all its differential
consequences.

\begin{example}
\label{pde-example-kdv-linearization}
For the Korteweg--de Vries equation
$$
F=u_t+uu_x+u_{xxx}=0,
$$
the linearization is
$$
\ell_F(Q)=
D_tQ+uD_xQ+u_xQ+D_x^3Q.
$$
Time translation is a symmetry, represented
by the characteristic $Q=u_t$. Verifying
this means substituting the equation and
its total derivatives into
$\ell_F(u_t)$.
\end{example}

\medskip\noindent
{\bf Conservation laws.}

\begin{definition}
\label{pde-definition-conservation-law}
For an equation in variables $t$ and $x$,
a \textit{conservation law} is a pair of
differential functions $\rho$ and $J$ such
that
$$
D_t\rho+D_xJ=0
$$
on every solution. The function $\rho$ is
the \textit{density}, and $J$ is the
\textit{flux}.
\end{definition}

Under suitable boundary conditions,
integration gives
$$
\frac{d}{dt}\int_a^b\rho\,dx
=-J(t,b)+J(t,a).
$$
If the boundary flux vanishes, the integral
of the density is constant in time.

Geometrically, conservation laws are
horizontal differential forms that become
closed after restriction to the equation.
They are the PDE analogue of first
integrals in ordinary differential
equations and are closely tied to
symmetries and Hamiltonian structures.

\section{Regularity Phenomena in
Nonlocal Equations}
\label{pde-section-regularity-nonlocal}

\medskip\noindent
{\bf Conference information.}

\href{%
https://sites.google.com/view/%
60thanniversary-dmat-puc-rio/Home
}{International Conference Celebrating
the 60th Anniversary of the Department
of Mathematics at PUC-Rio}

August 12--14, 2026,
RDC Auditorium, PUC-Rio,
Rio de Janeiro.

\medskip\noindent
{\bf Speaker.}

Makson Santos,
Universidade de Lisboa.

\medskip\noindent
{\bf Date and time.}

August 13, 2026, 10:30--11:30.

\medskip\noindent
{\bf Title.}

Regularity Phenomena in Nonlocal
Equations.

\medskip\noindent
{\bf Abstract.}

The talk studies regularity for solutions
of fully nonlinear nonlocal equations in
elliptic and degenerate settings. In the
elliptic case, the goal is optimal
regularity for equations whose right-hand
side lies in $L^p$, with the regularity
space determined by $p$. In the degenerate
case, one obtains at least one viscosity
solution in
$C^{1,\alpha}_{\mathrm{loc}}$ for some
$\alpha\in(0,1)$. When the order of the
operator is sufficiently close to $2$,
one also obtains H\"older estimates for
the gradient of every viscosity solution.

\medskip\noindent
{\bf Notes.} Once we have a (weak) solution
of a PDE, we study what differentiability 
properties does the function have only
because it solves the given PDE.



% --- End pde.tex ---

% --- Begin probability.tex ---
\chapter{Probability}
\label{chapter-probability}





\section{Introduction}
\label{probability-section-introduction}

\noindent
These notes develop probability theory.

Given the material elsewhere in the Heap
Project, the basic things to learn are:

\begin{enumerate}
\item Probability spaces, random variables,
distributions, expectation, independence,
and conditional probability.

\item The main discrete and continuous
distributions, their moments, transforms,
and relations to one another.

\item Modes of convergence of random
variables: almost sure, in probability,
in distribution, and in $L^p$.

\item The Borel--Cantelli lemmas, laws of
large numbers, and the central limit
theorem.

\item Conditional expectation as an
orthogonal projection and as the basic
tool for updating information.

\item Martingales, stopping times, optional
stopping, and the principal martingale
convergence theorems.

\item Concentration inequalities, from
Markov and Chebyshev inequalities to
exponential bounds and sub-Gaussian
behavior.

\item Characteristic functions and the
Fourier-analytic proofs of convergence
theorems.

\item Product measures, infinite sequences
of random variables, and the construction
of stochastic processes.

\item Markov chains, invariant measures,
recurrence, mixing, and ergodic theorems.

\item Brownian motion, the It\^o integral,
stochastic differential equations, and
their relation to the heat equation.

\item The statistical viewpoint:
estimation, confidence intervals,
likelihood, minimax risk, robustness, and
asymptotic efficiency.
\end{enumerate}

\section{A Primer for Lattice Random Walks}
\label{probability-section-primer-lattice-random-walks}

This section gives the probabilistic
background needed to read
\emph{Can one hear the shape of a lattice
random walk?} by Pieter Belmans, Sergey
Galkin, and Swarnava Mukhopadhyay. The
paper also uses Laurent polynomials,
convex polytopes, and spectral theory.
Here the emphasis is on the probability
dictionary connecting those subjects.

\medskip\noindent
{\bf Probability laws.}

\begin{definition}
\label{probability-definition-probability-space}
A \emph{probability space} is a triple
$(\Omega,\mathcal F,\mathbb P)$ consisting
of a set, a sigma algebra of events, and a
nonnegative, countably additive measure
satisfying
$$
\mathbb P(\Omega)=1.
$$
A \emph{random variable} with values in a
measurable space $E$ is a measurable map
$$
X:\Omega\longrightarrow E.
$$
\end{definition}

\begin{definition}
\label{probability-definition-conditional-probability}
Let $A$ and $B$ be events with
$\mathbb P(B)>0$. The \emph{conditional
probability} of $A$ given $B$ is
$$
\mathbb P(A\mid B)
=\frac{\mathbb P(A\cap B)}{\mathbb P(B)}.
$$
\end{definition}

Only discrete random variables are needed
for the random-walk part of the paper. In
that setting the distribution of $X$ is
specified by its probability mass function
$$
p:E\longrightarrow[0,1].
$$
For every $v\in E$, its value is
$$
p(v)=\mathbb P(X=v).
$$
Thus
$$
p(v)\geq 0,
\qquad
\sum_{v\in E}p(v)=1.
$$
Conversely, any function satisfying these
conditions defines a random variable on a
suitable probability space.

\begin{definition}
\label{probability-definition-support-finite-range}
The \emph{support} of a discrete law $p$
is
$$
\operatorname{Supp}(p)
=\{v\in E:p(v)>0\}.
$$
The law has \emph{finite range} if its
support is finite.
\end{definition}

The paper sometimes uses the word
distribution for an arbitrary finitely
supported coefficient function. Such a
function is a probability distribution
only when its coefficients are
nonnegative and sum to one.

For a function $f$ on $E$, its expectation
is
$$
\mathbb E[f(X)]
=\sum_{v\in E}f(v)p(v),
$$
whenever the sum is defined. Every such
expectation is finite when $p$ has finite
support and $f$ is finite on that support.

\medskip\noindent
{\bf Independent increments and walks.}

\begin{definition}
\label
{definition-independent-random-variables}
Discrete random variables
$X_1,\ldots,X_m$ are \emph{independent} if
$$
\mathbb P(X_1=v_1,\ldots,X_m=v_m)
=\prod_{j=1}^m\mathbb P(X_j=v_j)
$$
for every choice of values. They are
\emph{identically distributed} if they
have the same probability law.
\end{definition}

The abbreviation i.i.d. means independent
and identically distributed.

\begin{definition}
\label{probability-definition-lattice-random-walk}
Let $L$ be a lattice. A \emph{step law}
on $L$ is a function
$$
p:L\longrightarrow[0,1]
$$
satisfying
$$
\sum_{v\in L}p(v)=1.
$$
Let $(\Omega,\mathcal F,\mathbb P)$ be a
probability space carrying i.i.d. random
variables
$$
X_j:\Omega\longrightarrow L
$$
whose common probability mass function is
$p$. Explicitly, for every $j\geq 1$ and
$v\in L$,
$$
\mathbb P(X_j=v)=p(v).
$$
Fix $s\in L$. The \emph{lattice random
walk} with step law $p$, started at $s$,
is the sequence of random variables
$$
S_m:\Omega\longrightarrow L
$$
defined, for every $\omega\in\Omega$, by
$$
S_0(\omega)=s,
\qquad
S_m(\omega)
=s+\sum_{j=1}^mX_j(\omega).
$$
\end{definition}

Thus $\Omega$ is the sample space and $L$
is the state space. The point $s$ is a
fixed element of $L$, not an element of
$\Omega$. Each $X_j$ is the random
$j$th increment, each $S_m$ is the random
position at time $m$, and the entire
sequence $(S_m)_{m\geq 0}$ is the random
walk. Such a sequence of random variables
is also called a discrete-time stochastic
process.

The next position depends on the past only
through the present position. More
precisely,
$$
\mathbb P(S_{m+1}=x+v\mid S_m=x)=p(v).
$$
This is the Markov property, and $p(v)$ is
the transition probability from $x$ to
$x+v$. Translation invariance makes this
probability independent of $x$ and $m$.

\begin{definition}
\label{probability-definition-convolution-discrete-laws}
The \emph{convolution} of functions $p$
and $q$ on $L$ is
$$
(p*q)(x)=\sum_{v\in L}p(v)q(x-v).
$$
The $m$-fold convolution of $p$ is denoted
by $p^{*m}$, with $p^{*0}$ concentrated at
the origin.
\end{definition}

Independence gives
$$
\mathbb P(S_m=s+x)=p^{*m}(x).
$$
In particular, the probability of return
to the starting point after $m$ steps is
$$
a_m=\mathbb P(S_m=S_0)=p^{*m}(0).
$$
Since the starting point is usually fixed
at zero, the paper also writes this as a
conditional probability given $S_0=0$.

\medskip\noindent
{\bf Reachability and periodicity.}

\begin{definition}
\label{probability-definition-irreducible-lattice-walk}
A lattice walk is \emph{irreducible} if
every $x\in L$ can be reached from the
origin with positive probability. Thus,
for every $x\in L$, there is an $m$ such
that
$$
p^{*m}(x)>0.
$$
\end{definition}

Equivalently, the additive monoid generated
by $\operatorname{Supp}(p)$ is all of $L$.
The subgroup generated by the support can
be smaller than the initially chosen
lattice. Restricting the walk to that
subgroup removes the first ambiguity
discussed in the paper. When zero lies in
the interior of the convex hull of the
support, group generation also gives the
required two-sided reachability.

\begin{definition}
\label{probability-definition-period-lattice-walk}
For an irreducible walk, its \emph{period}
is
$$
h=\gcd\{m\geq 1:a_m>0\}.
$$
The walk is \emph{aperiodic} if $h=1$.
\end{definition}

This notion of period should not be
confused with the paper's \emph{period
sequence}, which is the entire sequence
$(a_m)_{m\geq 0}$. An irreducible walk can
still be periodic.

\medskip\noindent
{\bf Drift and covariance.}

Suppose now that $L=\mathbb Z^d$. Vectors
in another rank-$d$ lattice can be handled
after choosing a basis.

\begin{definition}
\label{probability-definition-drift-lattice-walk}
The \emph{mean} or \emph{drift} of one
step is
$$
\mu=\mathbb E[X]
=\sum_{v\in\mathbb Z^d}v p(v).
$$
The walk is \emph{mean-zero} or
\emph{centered} if $\mu=0$.
\end{definition}

\begin{definition}
\label{probability-definition-covariance-lattice-walk}
The \emph{covariance matrix} of one step
is
$$
\Sigma
=\mathbb E[(X-\mu)(X-\mu)^{\mathsf T}].
$$
Equivalently, its entries are
$$
\Sigma_{ij}
=\mathbb E[(X_i-\mu_i)(X_j-\mu_j)].
$$
The law is \emph{non-degenerate} if
$\Sigma$ is invertible. It is
\emph{isotropic} if, for some
$\sigma^2>0$,
$$
\Sigma=\sigma^2 I_d.
$$
\end{definition}

For every $r\in\mathbb R^d$,
$$
r^{\mathsf T}\Sigma r
=\mathbb E[(r^{\mathsf T}(X-\mu))^2]
\geq 0.
$$
Thus covariance is always positive
semidefinite. It is positive definite
precisely when the support of $p$ is not
contained in a proper affine hyperplane.
This is the paper's geometric meaning of
non-degeneracy.

Independence yields
$$
\mathbb E[S_m]=S_0+m\mu,
\qquad
\operatorname{Cov}(S_m)=m\Sigma.
$$
Mean zero removes the first-order drift,
while isotropic covariance makes the
second-order spreading the same in every
direction.

If a non-degenerate finite law has mean
zero, then zero lies in the interior of
the convex hull of its support. Indeed, a
separating linear functional at a boundary
point would be nonnegative on the support.
Its expectation could vanish only if the
support lay in a hyperplane.

\medskip\noindent
{\bf Three generating functions.}

For $v=(v_1,\ldots,v_d)$ and
$z=(z_1,\ldots,z_d)$, put
$$
z^v=\prod_{j=1}^d z_j^{v_j}.
$$

\begin{definition}
\label{probability-definition-laurent-step-function}
The \emph{Laurent generating function} of
the step law is
$$
W(z)=\sum_{v\in\mathbb Z^d}p(v)z^v.
$$
\end{definition}

Finite range makes $W$ a Laurent
polynomial. Its coefficients are exactly
the step probabilities, and
$$
W(1,\ldots,1)=1.
$$
The paper also calls $W$ a factorial
moment-generating function. In this setting
that phrase refers to this same Laurent
generating function, not to an additional
random object.

\begin{definition}
\label{probability-definition-characteristic-function}
The \emph{characteristic function} of one
step is
$$
\phi(\theta)
=\mathbb E
 [\exp(i\left\langle{}
 X,\theta\right\rangle{})]
=\sum_v p(v)
 \exp(i\left\langle{}
 v,\theta\right\rangle{}).
$$
\end{definition}

Here each coordinate of $\theta$ is
understood modulo $2\pi$. Consequently,
$\phi$ is a function on the torus. It is
the restriction of $W$ to the unit torus:
$$
\phi(\theta)
=W(e^{i\theta_1},\ldots,e^{i\theta_d}).
$$

\begin{definition}
\label{probability-definition-moment-generating-function}
The \emph{moment-generating function} of
one step is
$$
M(u)
=\mathbb E
 [\exp(\left\langle{}
 X,u\right\rangle{})]
=\sum_v p(v)
 \exp(\left\langle{}
 v,u\right\rangle{}).
$$
\end{definition}

It is the restriction of $W$ to the
positive real torus:
$$
M(u)=W(e^{u_1},\ldots,e^{u_d}).
$$
After analytic continuation, the two
transforms satisfy
$$
M(u)=\phi(-iu).
$$
Since the law has finite range, $M(u)$ is
finite for every $u\in\mathbb R^d$. The
three functions $W$, $\phi$, and $M$ are
different restrictions of the same finite
sum.

Differentiation at the origin recovers the
moments:
$$
\nabla M(0)=\mu,
\qquad
\nabla^2M(0)=\mathbb E[XX^{\mathsf T}].
$$
The logarithm gives the centered moments:
$$
\nabla\log M(0)=\mu,
\qquad
\nabla^2\log M(0)=\Sigma.
$$
Thus, in the mean-zero case used in the
paper, the Hessian of $M$ at zero is the
covariance matrix. Mean zero says that
zero is a critical point, and
non-degeneracy says that this critical
point is Morse.

More generally, the Hessian of $\log M$ at
any point is a covariance matrix for a
tilted law. It is positive definite when
the support has full affine span. Hence
$\log M$, and also $M$, are strictly convex
in the non-degenerate setting. A critical
point is therefore the unique global
minimum. This is the probabilistic content
of the paper's conifold-point condition.

\medskip\noindent
{\bf Normalization and graph potentials.}

Suppose instead that a Laurent polynomial
has nonnegative coefficients $c_v$ and
$$
C=\sum_v c_v.
$$
The corresponding step probabilities are
$$
p(v)=\frac{c_v}{C},
$$
so its probabilistic generating function is
the original polynomial divided by $C$.
This distinction is essential in the
paper. Its unnormalized graph potential is
denoted by $\widetilde W_\gamma$, whereas
the random-walk generating function is
$$
W_\gamma
=\frac{1}{8(g-1)}
 \widetilde W_\gamma.
$$
There are $8g-8$ allowed steps, each with
probability $1/(8g-8)$.

\medskip\noindent
{\bf Exponential tilting.}

For $\tau\in\mathbb R^d$, define
$$
p_\tau(v)
=\frac{
 \exp(\left\langle{}
 \tau,v\right\rangle{})p(v)}
 {M(\tau)}.
$$
This is the \emph{exponential tilt} of
$p$. The denominator is necessary because
the unnormalized weights sum to $M(\tau)$.
The tilted moment-generating function is
$$
M_\tau(u)=\frac{M(u+\tau)}{M(\tau)}.
$$
Its mean and covariance are
$$
\mu_\tau=\nabla\log M(\tau),
\qquad
\Sigma_\tau=\nabla^2\log M(\tau).
$$

Translation of the argument of $M$ thus
corresponds to exponential tilting of the
step law. If the division by $M(\tau)$ is
omitted, the new coefficients define a
probability law only when $M(\tau)=1$.
This is the second probabilistic ambiguity
described in the introduction of the
paper.

Tilting does not change the support or the
Newton polytope. For return probabilities,
the exponential factors cancel along every
closed path, and hence
$$
a_m(p_\tau)=\frac{a_m(p)}{M(\tau)^m}.
$$
Thus two tilts lying on the same level set
of $M$ have the same return probabilities.

If zero is in the interior of the step
polytope, $M$ has a unique minimizer
$\tau_0$. It gives a centered tilted law,
since
$$
\nabla\log M(\tau_0)=0.
$$

\medskip\noindent
{\bf Return probabilities and constant
terms.}

Multiplication of generating functions
corresponds to convolution. Namely,
$$
W(z)^m
=\sum_{x\in\mathbb Z^d}p^{*m}(x)z^x.
$$
Equivalently, independence gives
$$
\mathbb E
 [\exp(i\left\langle{}
 S_m-S_0,\theta\right\rangle{})]
=\phi(\theta)^m.
$$
If $[z^x]$ denotes the coefficient of
$z^x$, then
$$
[z^x]W(z)^m
=\mathbb P(S_m-S_0=x).
$$
In particular,
$$
a_m=[z^0]W(z)^m.
$$

Constant coefficients can also be
extracted by Fourier integration on the
torus. For every Laurent polynomial $F$,
$$
[z^0]F(z)
=\frac{1}{(2\pi)^d}
 \int_{[-\pi,\pi]^d}
 F(e^{i\theta_1},\ldots,
 e^{i\theta_d})\,d\theta.
$$
Consequently,
$$
a_m
=\frac{1}{(2\pi)^d}
 \int_{[-\pi,\pi]^d}
 \phi(\theta)^m\,d\theta.
$$
Although the integrand can be complex,
Fourier orthogonality makes the integral a
real, nonnegative return probability.

For an unnormalized polynomial
$\widetilde W=CW$, one instead has
$$
[z^0]\widetilde W(z)^m=C^m a_m.
$$
The constant term then counts weighted
closed paths; division by $C^m$ turns that
count into a probability.

\begin{definition}
\label{probability-definition-period-sequence}
The \emph{period sequence} of a finite
range walk is
$$
\Pi_W=(a_m)_{m\geq 0}.
$$
\end{definition}

Knowing all $a_m$ is weaker than knowing
all $m$-step distributions. It retains only
the coefficient at the origin in every
convolution power. This loss of information
is what makes the paper's non-reconstruction
result possible.

\medskip\noindent
{\bf The geometry of a step law.}

\begin{definition}
\label{probability-definition-newton-polytope-step-law}
The \emph{Newton polytope} of $W$, or the
\emph{step polytope} of $p$, is
$$
\operatorname{NP}(W)
=\operatorname{conv}
 (\operatorname{Supp}(p))
\subset\mathbb R^d.
$$
\end{definition}

It records the convex geometry of the
allowed steps but forgets their
probabilities. In the examples in the
paper all allowed steps have equal weight,
so the support and its convex hull carry
most of the visible input data.

The paper says that two walks have the
same shape when a linear isomorphism of
their ambient real vector spaces maps one
Newton polytope to the other. Equal shape
is therefore much stronger than equality
of the period sequences in the examples
under consideration.

\medskip\noindent
{\bf The Markov operator and the discrete
Laplacian.}

Let $q:L\to\mathbb C$ be a function. The
one-step Markov operator is
$$
(Pq)(x)
=\mathbb E[q(x+X)]
=\sum_{v\in L}p(v)q(x+v).
$$
The associated difference operator is
$$
\Delta_p=P-I,
$$
or explicitly
$$
(\Delta_p q)(x)
=\sum_{v\in L}p(v)
 (q(x+v)-q(x)).
$$
This is the operator used in the
introduction of the paper.

For a smooth function on $\mathbb R^d$, a
Taylor expansion with lattice spacing
$\hbar$ gives
$$
\begin{aligned}
&\mathbb E[q(x+\hbar X)]-q(x) \\
&\quad=\hbar\sum_i\mu_i\partial_iq(x) \\
&\qquad
+\frac{\hbar^2}{2}
 \sum_{i,j}\mathbb E[X_iX_j]
 \partial_i\partial_jq(x)
+O(\hbar^3).
\end{aligned}
$$
Mean zero cancels the first-order term. In
that case $\mathbb E[X_iX_j]=\Sigma_{ij}$.
If the covariance is isotropic, the
quadratic term is a scalar multiple of the
ordinary Laplacian:
$$
\frac{\sigma^2\hbar^2}{2}
\sum_{i=1}^d\partial_i^2q(x).
$$
This explains the mean-zero and isotropic
covariance assumptions behind the paper's
discretization of the standard Laplacian.

Fourier characters diagonalize the Markov
operator. For
$$
q_\theta(x)
=\exp(i\left\langle{}
 x,\theta\right\rangle{}),
$$
one has
$$
Pq_\theta=\phi(\theta)q_\theta,
\qquad
\Delta_pq_\theta
=(\phi(\theta)-1)q_\theta.
$$
Thus the same characteristic function that
encodes the step law is the Fourier symbol
of the translation-invariant operator.
The identity
$$
a_m
=\frac{1}{(2\pi)^d}
 \int_{[-\pi,\pi]^d}
 \phi(\theta)^m\,d\theta
$$
says that return probabilities are moments
of this symbol. This is the bridge to the
paper's spectral language.

When $p(v)=p(-v)$, the operator is
self-adjoint and $\phi$ is real. Pushing
normalized Haar measure on the torus
forward by $\phi$ then gives a spectral
probability measure whose $m$th moment is
$a_m$. The paper uses related
density-of-states constructions even when
the presentation is not phrased in this
symmetric special case.

\medskip\noindent
{\bf Push-forward measures and integral
transforms.}

\begin{definition}
\label{probability-definition-push-forward-measure}
Let $f:X\to Y$ be measurable and let
$\nu$ be a measure on $X$. The
\emph{push-forward measure} $f_*\nu$ on
$Y$ is defined by
$$
(f_*\nu)(A)=\nu(f^{-1}(A)).
$$
\end{definition}

The defining change-of-variables identity
is
$$
\int_Y g(y)\,d(f_*\nu)(y)
=\int_X g(f(x))\,d\nu(x).
$$
If the push-forward has a density, that
density describes how much volume in the
domain is carried to each range of values
of $f$.

The paper pushes Lebesgue measure on
$\mathbb R^d$ forward by $M$. Its two
integral invariants are
$$
K_W(t)
=\int_{\mathbb R^d}
 \exp(-tM(u))\,du
$$
and
$$
\zeta_W(s)
=\int_{\mathbb R^d}M(u)^{-s}\,du.
$$
These formulas are considered only for
parameters for which the integrals
converge.
They are respectively Laplace-type and
Mellin-type transforms of the push-forward
measure. This push-forward is not the law
of $M(X)$. The variable $u$ ranges over
the parameter space of exponential
moments, and Lebesgue measure is being
pushed forward.

A volume-preserving change of variables
$F$ satisfying
$$
M_\gamma(F(u))=M_{\gamma'}(u)
$$
makes these two integrals equal. This is
the measure-theoretic content of the
change-of-variables argument in Section 5
of the paper.

\medskip\noindent
{\bf Limit theorems mentioned in the
paper.}

The main reconstruction argument does not
require the following results, but the
introduction mentions them as part of the
random-walk and mirror-symmetry context.

\begin{theorem}[Central limit theorem]
\label{probability-theorem-central-limit-lattice-walk}
Let the step law have mean $\mu$ and
non-degenerate covariance $\Sigma$. Then
$$
\frac{S_m-S_0-m\mu}{\sqrt m}
\Longrightarrow
\mathcal N(0,\Sigma).
$$
\end{theorem}

The arrow denotes convergence in
distribution. Thus averages of bounded
continuous functions converge to their
averages under the indicated Gaussian
law. The central limit theorem describes
probabilities on regions of scale
$\sqrt m$.

The local central limit theorem refines
this to individual lattice points. In the
centered, non-degenerate, irreducible, and
aperiodic case, it gives
$$
a_m
=\frac{1}{
 (2\pi m)^{d/2}\sqrt{\det\Sigma}}
+o(m^{-d/2}).
$$
For a walk of period $h$, a periodic
correction is required and the formula is
read only along admissible step counts.
In particular, centered return
probabilities decay polynomially. With
nonzero drift, returning to the origin is
a large-deviation event and, under the
finite-range hypotheses in the paper, the
decay is exponential up to a polynomial
factor.

Finite range also gives exponential
concentration. For example, fix
$r\in\mathbb R^d$ and suppose
$$
|\left\langle{}
r,X-\mu\right\rangle{}|\leq B.
$$
Hoeffding's inequality gives
$$
\mathbb P
 (|\left\langle{}
 r,S_m-S_0-m\mu\right\rangle{}|
 \geq x)
\leq
2\exp\left(-\frac{x^2}{2mB^2}\right).
$$
Taking $x$ proportional to $m$ produces a
bound exponentially small in $m$. This is
the meaning of exponential concentration
in the introduction.

\medskip\noindent
{\bf A one-dimensional model.}

Consider the simple symmetric walk on
$\mathbb Z$ with
$$
p(1)=p(-1)=\frac12.
$$
Its three generating functions are
$$
W(z)=\frac{z+z^{-1}}2,
\qquad
\phi(\theta)=\cos\theta,
\qquad
M(u)=\cosh u.
$$
The walk is mean-zero, has covariance one,
and is irreducible. It nevertheless has
period two. Its return probabilities are
$$
a_{2m}
=\frac{1}{2^{2m}}\binom{2m}{m},
\qquad
a_{2m+1}=0.
$$
Stirling's formula gives
$$
a_{2m}\sim\frac{1}{\sqrt{\pi m}}.
$$
This example displays the polynomial decay
and the periodic correction at once.

For comparison, let
$$
p(1)=q,
\qquad
p(-1)=1-q.
$$
Then
$$
M(u)=qe^u+(1-q)e^{-u},
\qquad
\mu=2q-1,
$$
and
$$
a_{2m}
=\binom{2m}{m}q^m(1-q)^m.
$$
Unless $q=1/2$, the factor
$4q(1-q)$ is smaller than one and return
probabilities decay exponentially after
the polynomial correction. The minimizer
of $M$ is
$$
\tau_0
=\frac12\log\left(\frac{1-q}{q}\right),
$$
and exponential tilting by $\tau_0$ turns
the biased walk into the symmetric one.

\medskip\noindent
{\bf Reading the main construction.}

For a trivalent graph $\gamma$ of genus
$g$, the paper constructs $8g-8$ lattice
vectors. A single step $X_\gamma$ is
uniformly distributed on these vectors.
The following dictionary summarizes the
probabilistic argument:

\begin{enumerate}
\item The cancellation of the step vectors
gives
$$
\mathbb E[X_\gamma]=0.
$$
Equivalently, the logarithmic
moment-generating function has a critical
point at the origin.

\item The support spans the lattice and
zero lies in the interior of its Newton
polytope. Hence the walk is irreducible and
non-degenerate.

\item The normalized graph potential is
the Laurent generating function of the
step law.

\item The period number
$$
a_m(\gamma)
=[z^0]W_\gamma(z)^m
$$
is the probability of return after $m$
steps.

\item Birational changes of variables
between graph potentials imply that graphs
of the same genus have the same entire
period sequence.

\item The Newton polytope is the convex
hull of the possible steps. The
reconstruction theorem recovers the graph
from this polytope.
\end{enumerate}

Consequently, non-isomorphic graphs give
walks with different shapes but identical
return probabilities at every time. No
additional probability theory is hidden in
that conclusion: its force comes from
combining the coefficient identity for
$a_m$ with the combinatorial reconstruction
of the Newton polytope.

\section{Robust and Efficient Mean
Estimation via Soft Shrinkage}
\label{probability-section-robust-efficient-mean}

\medskip\noindent
{\bf Conference information.}

\href{%
https://sites.google.com/view/%
60thanniversary-dmat-puc-rio/Home
}{International Conference Celebrating
the 60th Anniversary of the Department
of Mathematics at PUC-Rio}

August 12--14, 2026,
RDC Auditorium, PUC-Rio,
Rio de Janeiro.

\medskip\noindent
{\bf Speaker.}

Paulo Orenstein, IMPA.

\medskip\noindent
{\bf Date and time.}

August 12, 2026, 11:15--12:15.

\medskip\noindent
{\bf Title.}

Robust and Efficient Mean Estimation
via Soft Shrinkage.

\medskip\noindent
{\bf Abstract.}

How should we weight observations when
estimating a mean? The sample mean is
efficient and minimax but fragile under
heavy tails; robust estimators resist
heavy tails but typically sacrifice
efficiency. We show a soft-shrinkage
calibration improves on both fronts:
observations are smoothly downweighted
by their distance from a pilot estimate,
with the shrinkage scale calibrated so
that the total rejected weight equals a
prescribed budget. Any reasonable pilot
is upgraded to sub-Gaussian concentration
with the minimax constant $\sqrt{2}$ under
finite variance alone, and the budget
doubles as an exact breakdown point. The
mechanism trades variance for bias, and
the bias can be quantified exactly. This
yields a second-order theory establishing
an MSE expansion that identifies when
shrinkage beats the sample mean,
Berry--Esseen and moderate deviation
approximations, and asymptotically exact
confidence intervals. Beyond the theory,
the results give concrete guidance for
designing mean estimators and bridge
modern sub-Gaussian estimation with
classical robust statistics. This is
joint work with Ant\^onio Cat\~ao.

\medskip\noindent
{\bf Notes.}
The talk starts introducing the mean
$\mu$ of a sample $X_i$, which is 
the sum of all the variables divided
by the size of the sample. This is
an efficient estimate, but not robust.
Robustness can be explained with atypical
data: suppose that you have a very atypical
datum in the sample. This will move the
mean $\mu$ quite a lot, but this isn't very
good since it was just one extraordinary
datum. Then we introduce the trimmed mean,
which uses a new variable to trim, or
remove, the tails of the data, i.e.
the most extraordinary data points
to avoid the case we just explained.
The trimmed mean is robust, but looses
efficiency since we need to chop off a 
piece of our data, which isn't ideal.

The talk aims at introducing a new 
estimate which lies in between the 
usual mean and the trimmed mean.
The idea is to introduce weights to every
data using a variable  $\alpha$ (which
depends on the whole data).
Such a new estimate is both robust
and efficient.

Developing these kinds of statistical
gadgets has applications in mathematics,
science, computer science, artificial
intelligence, etc. In particular, this can
be thought of as a more sophisticated method
than simple gradient descent used in machine
learning methods, the one used in the AI
course at EGAXII.


\section{Recovering Structure from Binary
Graphical Model Dynamics}
\label{probability-section-recovering-binary-graphical}

\medskip\noindent
{\bf Conference information.}

\href{%
https://sites.google.com/view/%
60thanniversary-dmat-puc-rio/Home
}{International Conference Celebrating
the 60th Anniversary of the Department
of Mathematics at PUC-Rio}

August 12--14, 2026,
RDC Auditorium, PUC-Rio,
Rio de Janeiro.

\medskip\noindent
{\bf Speaker.}

Roberto Imbuzeiro, IMPA.

\medskip\noindent
{\bf Date and time.}

August 12, 2026, 16:00--17:00.

\medskip\noindent
{\bf Title.}

Recovering structure from binary
graphical model dynamics.

\medskip\noindent
{\bf Abstract.}

We report on work in progress with
Guilherme Ost (UFRJ) and Guilherme Reis
(UFF) on a class of toy models for
neuronal dynamics. These models consist
of sets of units taking $0/1$ values,
which interact pairwise via a matrix of
weights. The main question we address is
how much of the weight structure can be
recovered from the dynamics. We will
focus on ``mean field'' cases where
individual pairwise interactions may be
quite weak. Interestingly, this property
turns out to be helpful.

\medskip\noindent
{\bf Notes.}
This talk made me remember Carina Curto's
neural networks toy models. It looks
like they worked in (oversimplified)
models describing neural networks.

(AI paragraph)
The two approaches are complementary.
Curto starts with a directed connectivity
graph and studies the resulting dynamics,
such as fixed points and attractors.
Imbuzeiro starts with observed binary
dynamics and asks what can be recovered
about the unknown interaction graph. Thus
Curto studies structure to dynamics, while
Imbuzeiro studies dynamics to structure.


A paper mentioned in the talk is
\href{%
https://doi.org/10.1214/25-AOS2592
}{Inferring the dependence graph density
of binary graphical models in high
dimension}, by Julien Chevallier,
Eva L\"ocherbach, and Guilherme Ost. It
appeared in \emph{The Annals of
Statistics}, volume 54, number 2 (2026),
pages 861--881.

\section{Detecção de Comunidades para
Modelos Gráficos Binários de Alta Dimensão}
\label{probability-section-ost-colmea-2026}

\noindent
Guilherme Ost (IMPA).
September 23, 2026.
COLMEA, IMPA, 15:40--17:00.

\medskip\noindent
{\bf Abstract.}
Consideramos um sistema de $N$ unidades
binárias que interagem aos pares por meio
de uma matriz de pesos. Esta matriz é
possivelmente aleatória e possui uma
estrutura em blocos, codificando conjuntos
de unidades que desempenham um papel
semelhante no sistema, chamados de
comunidades. Tanto a matriz de pesos quanto
essas comunidades são desconhecidas. Nesta
palestra, focarei na seguinte questão
estatística: ao observar apenas os estados
das $N$ unidades ao longo de $T$ passos de
tempo, conseguimos recuperar as diferentes
comunidades do sistema? Nosso objetivo é
apresentar garantias teóricas de alguns
métodos estatísticos simples que nos
permitem recuperar essas comunidades a
partir dos dados observados, desenvolvidos
recentemente com J. Chevallier (Grenoble,
França).
\medskip\noindent
{\bf Notes.}
The classical version of this problem is
``community detection in networks.'' The
simplest case is described by a symmetric
stochastic block model.

Let the $N$ vertices be divided into two
communities
$
\mathcal{C}_1,\mathcal{C}_2
$
of equal size. Every pair $(i,j)$ of
vertices is connected independently with
probability $p$ if $i\sim j$, meaning that
$i$ and $j$ belong to the same community,
and with probability $q$ if $i\not\sim j$.
We suppose
$
p>q>0.
$

The resulting connectivity is represented
by an adjacency matrix. The problem is:
given a permuted version of the adjacency
matrix, say $\theta$, for an unknown
permutation of the vertices, recover the
communities $\mathcal{C}_1$ and
$\mathcal{C}_2$.

There is a solution to this classical
problem using spectral methods.

Here $p$ is the probability of connecting
two vertices in the same community, while
$q$ is the probability of connecting two
vertices in different communities.

The basic quantitative question is how
separated $p$ and $q$ must be in order to
recover the communities. Equivalently, how
small can $p-q$ be before the community
structure becomes statistically invisible?

\medskip\noindent
{\bf Some classical thresholds.}
Yun and Prouti\`ere (2014) use the sparse
parametrization
$
p=\frac{a_N}{N},
\qquad
q=\frac{b_N}{N}.
$
Take a vertex $v\in\mathcal{C}_1$.
Since $|\mathcal{C}_1|=N/2$, the expected
number of its neighbours inside
$\mathcal{C}_1$ is approximately
$
|\mathcal{C}_1|p
=\frac{N}{2}\frac{a_N}{N}
=\frac{a_N}{2}.
$
Likewise, the expected number of
neighbours of $v$ in $\mathcal{C}_2$ is
$
|\mathcal{C}_2|q=\frac{b_N}{2}.
$
Thus $a_N-b_N$ measures the separation
between the two expected degree profiles.

Abbe, Bandeira and Hall (2016) consider
$
p=\frac{a\log N}{N},
\qquad
q=\frac{b\log N}{N}.
$
They obtain a sharp threshold for exact
recovery. Above
$
\frac{a+b}{2}-\sqrt{ab}>1
$
exact recovery is possible with
probability tending to one; below the
corresponding strict inequality it is
impossible. Equivalently, the threshold is
$
(\sqrt a-\sqrt b)^2=2.
$

\medskip\noindent
{\bf From observing the graph to observing
binary data.}
The problem of interest in this talk is
harder: the graph itself is not observed.
The motivating example comes from
neuroscience and spike-train data. A
``spike'' is an action potential, i.e. an
event in which a neuron fires.

Discretize time into windows and record
$1$ if neuron $i$ fires during a window
and $0$ otherwise. At each time $t$ this
gives a binary state
$
X_t\in\{0,1\}^N.
$
Over $T$ time windows one therefore
observes a binary data matrix. The hidden
question is whether one can reconstruct
the excitatory and inhibitory communities
from these data alone. This is the
community detection problem from binary
data.

\medskip\noindent
{\bf Ingredients.}
\begin{itemize}
\item
An unobserved graph $\theta$, possibly
random, directed, and with block
structure.

\item
A probability law $\mathbb{P}_\theta$ on
trajectories of binary network states.

\item
An observed sample
$
(X_1,\ldots,X_T)\sim\mathbb{P}_\theta,
$
from which we want to recover the
communities.
\end{itemize}

\medskip\noindent
{\bf IA-remark.}
Morally, a trajectory is just the movie of
the network through time. Each frame
$X_t\in\{0,1\}^N$ records which of the
$N$ neurons fired in the $t$th time
window. Thus
$
(X_1,\ldots,X_T)
$
is one complete observed history of the
network. A probability law on trajectories
says how likely the different possible
movies are.

\medskip\noindent
{\bf Independent benchmark.}
In an i.i.d. model one would have
$
\mathbb{P}_\theta
=\bigotimes_{t=1}^T P_\theta,
$
where $P_\theta$ is a probability measure
on $\{0,1\}^N$. A famous example is the
Ising blockmodel of Berthet, Rigollet and
Srivastava (2019), where independent
binary observations are used to recover
an underlying block structure.

\medskip\noindent
{\bf The speaker's model.}
The model goes beyond the i.i.d.
assumption. The vectors
$
X_1,\ldots,X_T\in\{0,1\}^N
$
are dependent under
$\mathbb{P}_\theta$. The random graph has
directed interactions and a block
structure. A connection emanating from an
excitatory neuron has positive weight.

Under $\mathbb{P}_\theta$, the process
$
(X_t)_{1\leq t\leq T}
$
is a stationary Markov chain.

The statistical problem is: from the
sample $X_1,\ldots,X_T$, recover the
partition
$
(\mathcal{P}_+,\mathcal{P}_-).
$
The model parameters $\mu$, $\lambda$ and
$\theta$ are unknown. The community sizes
have the form
$
|\mathcal{P}_+|=Nr_+,
\qquad
|\mathcal{P}_-|=Nr_-,
$
where
$
0<r_+,r_-<1,
\qquad
r_++r_-=1,
$
and $r_+$ and $r_-$ are also unknown.

\medskip\noindent
{\bf AI-reference check.}
Chevallier and Ost's paper
``Community detection for binary graphical
models in high dimension'' studies this
dependent-data problem. Their exact
recovery probability tends to one in the
regime
$
\frac{N}{\sqrt{T}}\log(NT)\longrightarrow0.
$

\medskip\noindent
{\bf AI-speaker profile.}
Guilherme Ost is a professor at UFRJ who
works in probability and stochastic
processes. His research repeatedly studies
how hidden structure can be inferred from
dependent stochastic observations. Topics
include neuronal interaction graphs from
spike trains, high-dimensional binary
graphical models, and random walks on
random graphs. In 2026 he is also teaching
High-Dimensional Statistics at IMPA.

% --- End probability.tex ---

% --- Begin optimal-transport.tex ---
\chapter{Optimal Transport}
\label{chapter-optimal-transport}





\section{Uma perspectiva do transporte
ótimo sobre a Teoria do Funcional da
Densidade: entropia, métodos numéricos e o
regime de elétrons estritamente
correlacionados}
\label{optimal-transport-section-gerolin-colmea-2026}

\noindent
Augusto Gerolin (IMPA).
September 23, 2026.
COLMEA, IMPA, 14:00--15:20.

\medskip\noindent
{\bf Abstract.}
Nesta palestra, apresentaremos alguns
aspectos matemáticos da Teoria do Funcional
da Densidade sob a perspectiva do transporte
ótimo. A abordagem baseia-se no estudo da
regularização entrópica de problemas de
transporte ótimo multi-marginal com funções
de custo repulsivas, como o custo de Coulomb
e o negativo do logaritmo. Serão destacados
tanto os aspectos probabilísticos e
computacionais --- incluindo o algoritmo de
Sinkhorn --- quanto novas perspectivas e
aplicações em química teórica. A palestra
requer poucos pré-requisitos e não apresenta
contraindicações. Assim, estudantes de
mestrado e doutorado em física, química
teórica e matemática são especialmente
convidados a participar.

\medskip\noindent
{\bf Notes.}
\begin{itemize}
\item
See ``The strong-interaction limit of
density functional theory'': it tells the
story of the main characters of density
functional theory and optimal transport.

\item
Optimal transport, in its modern form,
started in the 1940s with Kantorovich.

\item
In 1964, Hohenberg and Kohn introduced the
density-functional viewpoint. In the
strong-interaction regime, the corresponding
density functional can be formulated as an
optimal transport problem. This is the
viewpoint of the talk.
\item
See ``Pushing the frontiers of density
functionals by solving the fractional
electron problem'' (Kirkpatrick et al.,
2021).

\item
See ``Accurate and scalable
exchange-correlation with deep learning''
(Luise et al., 2025).
\item
New branch: quantum optimal transport.
\end{itemize}

\medskip\noindent
{\bf AI explanation: ground state energy.}
The ground state of a quantum system is its
lowest-energy state. The corresponding
{\it ground state energy} is therefore the
smallest value that the energy can take
among all admissible quantum states. In
variational language, one minimizes the
energy functional over the allowed states.
For the electron systems considered here,
density functional theory tries to express
this minimum using the electron density
rather than the full many-electron
wavefunction.

% --- End optimal-transport.tex ---

\part{Algebra and Categories}
\label{part-algebra-and-categories}

% --- Begin linear-algebra.tex ---
\chapter{Linear Algebra}
\label{chapter-linear-algebra}





\section{Projections}
\label{linear-algebra-section-projections}

\begin{definition}
\label{linear-algebra-definition-projection}
Let $V$ be a vector space. A linear map
$P:V\to V$ is a {\it projection} if
$
P^2=P.
$
Its image is the subspace onto which
it projects, and its kernel specifies
the direction of the projection.
\end{definition}

\begin{lemma}
\label{linear-algebra-lemma-projection-decomposition}
If $P$ is a projection, then
$
V=\operatorname{Im}P\oplus\ker P.
$
The restriction of $P$ to its image is
the identity. Conversely, a decomposition
$V=U\oplus W$ defines a projection onto
$U$ with kernel $W$ by $P(u+w)=u$.
\end{lemma}

\begin{proof}
If $u=P(v)$, then $P(u)=P^2(v)=u$.
For any $v\in V$, write
$
v=P(v)+(v-P(v)).
$
The first term lies in the image, and
the second lies in the kernel because
$
P(v-P(v))=P(v)-P^2(v)=0.
$
If $w$ lies in both subspaces, then
$w=P(w)=0$. This proves the direct sum.

Conversely, the decomposition $v=u+w$
is unique. Thus $P(u+w)=u$ is well
defined and linear. Applying it twice
still gives $u$, so $P^2=P$.
Its image is $U$ and its kernel is $W$.
\end{proof}

\begin{remark}
\label{linear-algebra-remark-projection-choice}
A subspace can have different complements.
For example, on $\mathbb{R}^2$ the maps
$
P(x,y)=(x,0),\qquad Q(x,y)=(x-y,0)
$
are projections onto the same horizontal
axis. Their kernels are respectively
the vertical axis and the line $x=y$.

With an inner product, an orthogonal
projection uses the orthogonal complement
as its kernel. See
the section ``Hilbert spaces'' in the
Functional Analysis notebook.
The averaging projection in the
Gr\"unbaum notebook only requires the
algebraic definition above.
\end{remark}

\section{Signature}
\label{linear-algebra-section-signature}

The signature of a form quadratic form can be
equivalently computed as the
maximal dimension of an isotropic subspace,
or as the maximal dimension of the
subspace where the form is negative.

\begin{remark}
\label{linear-algebra-remark-nondegenerate-vs-anisotropic}
Let $B:V \times V \to k$ be a bilinear form.
The {\it radical} of $B$ is
$$
\operatorname{rad}(B)=
\{v \in V : B(v,w)=0\text{ for all }w
\in V\}.
$$
The form is {\it nondegenerate} if
$\operatorname{rad}(B)=0$. In coordinates,
if $A$ is the Gram matrix of $B$, this is
equivalent to $\det(A)\neq 0$.

This should not be confused with the
condition that there are no nonzero isotropic
vectors. A vector $v\neq 0$ is {\it
isotropic} if $B(v,v)=0$, and a form with no
nonzero isotropic vectors is called
{\it anisotropic}. Anisotropic implies
nondegenerate, since a nonzero vector in the
radical would be isotropic. The converse is
false; for example the form with matrix
$$
\begin{pmatrix}
1&0\\
0&-1
\end{pmatrix}
$$
is nondegenerate, but $(1,1)$ is isotropic.
\end{remark}

\section{Adjoint}
\label{linear-algebra-section-adjoint}

\begin{definition}
\label{linear-algebra-definition-adjoint}
Let  $A:(V,\left\langle{}\cdot,\cdot\right\rangle{}_V) \to
(W,\left\langle{}\cdot,\cdot\right\rangle{}_W)$ be a
linear transformation of inner product vector
spaces, its {\it adjoint} is the
map
$$
A^{\dag}:W\to V
$$
satisfying
\begin{equation}
\label{linear-algebra-equation-adjoint}
\left\langle{}Av,w\right\rangle{}_W=\left\langle{}v,A^\dag w\right\rangle{}_V
\end{equation}
\end{definition}

Let $G$ be a matrix representation for
$\left\langle{}\cdot,\cdot\right\rangle{}_V$ and $H$ for
$\left\langle{}\cdot,\cdot\right\rangle{}_W$. Then Eq.
\ref{linear-algebra-equation-adjoint} is
$$
v^{\mathbf{T}}A^{\mathbf{T}}Hw=(Av)
^{\mathbf{T}}Hw=v^{\mathbf{T}}GA^\dag
w
$$
So
\begin{equation}
\label{linear-algebra-equation-adjoint-transpose}
A^{\mathbf{T}}H=GA^\dag \iff
A^\dag:=G^{-1}A^{\mathbf{T}}H
\end{equation}
which gives the relationship between the
usual transpose matrix and the adjoint.

\begin{definition}
\label{linear-algebra-definition-Frobenius-norm}
In the space of matrices we define {\it
Frobenius norm} as
\begin{equation}
\label{linear-algebra-equation-Frobenius-norm}
\|A\|^2=\text{tr}A^\dag
A=\text{tr}(G^{-1}A^{\mathbf{T}}HA
\end{equation}
\end{definition}


% --- End linear-algebra.tex ---

% --- Begin groups.tex ---
\chapter{Groups}
\label{chapter-groups}





\section{Group extensions and semidirect
products}
\label{groups-section-group-extensions}

\begin{definition}
\label{groups-definition-group-extension}
A {\it group extension} of a group $Q$ by a
group $N$ is a short exact sequence
$$
1\longrightarrow N
\overset{i}{\longrightarrow}G
\overset{\pi}{\longrightarrow}Q
\longrightarrow1.
$$
It is {\it split} if there is a group
homomorphism
$$
s:Q\longrightarrow G
$$
such that
$$
\pi\circ s=\operatorname{id}_Q.
$$
Such a map $s$ is called a {\it splitting}
or a {\it section}.
\end{definition}

\noindent
Exactness says in particular that
$i(N)=\operatorname{Ker}\pi$, so $i(N)$ is
normal in $G$. We usually identify $N$ with
its image in $G$.

\begin{proposition}
\label{groups-proposition-split-extension-semidirect}
If
$$
1\longrightarrow N
\longrightarrow G
\longrightarrow Q
\longrightarrow1
$$
splits, then
$$
G\cong N\rtimes_\varphi Q,
$$
where
$$
\varphi:Q\longrightarrow
\operatorname{Aut}(N)
$$
is defined by
$$
\varphi(q)(n)=s(q)ns(q)^{-1}.
$$
\end{proposition}

\begin{proof}
Let $g\in G$ and put $q=\pi(g)$. Then
$$
\pi(gs(q)^{-1})=1,
$$
so there is $n\in N$ such that
$$
g=ns(q).
$$
This expression is unique: if
$$
ns(q)=n's(q'),
$$
applying $\pi$ gives $q=q'$, and then
$n=n'$.

Thus the map
$$
N\times Q\longrightarrow G,
\qquad
(n,q)\longmapsto ns(q)
$$
is a bijection. Multiplication in these
coordinates is
$$
(n,q)(n',q')
=
(n\varphi(q)(n'),qq'),
$$
which is the multiplication in the
semidirect product.
\end{proof}

\begin{remark}
\label{groups-remark-split-not-direct}
A split extension need not be a direct
product. It is a direct product with respect
to the chosen splitting precisely when
$$
\varphi(q)=\operatorname{id}_N
$$
for every $q\in Q$, equivalently when
$s(Q)$ centralizes $N$.
\end{remark}

\begin{example}
\label{groups-example-s3-semidirect}
The sequence
$$
1\longrightarrow C_3
\longrightarrow S_3
\longrightarrow C_2
\longrightarrow1
$$
splits, so
$$
S_3\cong C_3\rtimes C_2.
$$
The nontrivial element of $C_2$ acts on
$C_3$ by inversion. Hence the action is not
trivial, and
$$
S_3\not\cong C_3\times C_2.
$$
\end{example}

% --- End groups.tex ---

% --- Begin categories.tex ---
\chapter{Categories}
\label{chapter-categories}





\section{Definitions}
\label{categories-section-definition-categories}

\noindent
We recall the definitions, partly to fix
notation.

\begin{definition}
\label{categories-definition-category}
A {\it category} $\mathcal{C}$ consists of
the following data:
\begin{enumerate}
\item A set of objects $\Ob(\mathcal{C})$.
\item For each pair $x, y \in
\Ob(\mathcal{C})$ a set of morphisms
$\Mor_\mathcal{C}(x, y)$.
\item For each triple $x, y, z\in
\Ob(\mathcal{C})$ a composition
map $ \Mor_\mathcal{C}(y, z) \times
\Mor_\mathcal{C}(x, y)
\to \Mor_\mathcal{C}(x, z) $, denoted $(\phi,
\psi) \mapsto
\phi \circ \psi$.
\end{enumerate}
These data are to satisfy the following
rules:
\begin{enumerate}
\item For every element $x\in
\Ob(\mathcal{C})$ there exists a
morphism $\text{id}_x\in \Mor_\mathcal{C}(x,
x)$ such that
$\text{id}_x \circ \phi = \phi$ and $\psi
\circ \text{id}_x = \psi $ whenever
these compositions make sense.
\item Composition is associative, i.e.,
$(\phi \circ \psi) \circ \chi =
\phi \circ ( \psi \circ \chi)$ whenever these
compositions make sense.
\end{enumerate}
\end{definition}

\begin{definition}
\label{categories-definition-functor}
A {\it functor} $F : \mathcal{A} \to
\mathcal{B}$
between two categories $\mathcal{A},
\mathcal{B}$ is given by the
following data:
\begin{enumerate}
\item A map $F : \Ob(\mathcal{A}) \to
\Ob(\mathcal{B})$.
\item For every $x, y \in \Ob(\mathcal{A})$ a
map
$F : \Mor_\mathcal{A}(x, y) \to
\Mor_\mathcal{B}(F(x), F(y))$,
denoted $\phi \mapsto F(\phi)$.
\end{enumerate}
These data should be compatible with
composition and identity morphisms
in the following manner: $F(\phi \circ \psi)
=
F(\phi) \circ F(\psi)$ for a composable pair
$(\phi, \psi)$ of
morphisms of $\mathcal{A}$ and
$F(\text{id}_x) = \text{id}_{F(x)}$.
\end{definition}

\noindent
Note that every category $\mathcal{A}$ has an
{\it identity} functor
$\text{id}_\mathcal{A}$.
In addition, given a functor $G : \mathcal{B}
\to \mathcal{C}$
and a functor $F : \mathcal{A} \to
\mathcal{B}$ there is
a {\it composition} functor $G \circ F :
\mathcal{A} \to \mathcal{C}$
defined in an obvious manner.

\begin{definition}
\label{categories-definition-faithful}
Let $F : \mathcal{A} \to \mathcal{B}$ be a
functor.
\begin{enumerate}
\item We say $F$ is {\it faithful} if
for any objects $x, y \in \Ob(\mathcal{A})$
the map
$$
F : \Mor_\mathcal{A}(x, y) \to
\Mor_\mathcal{B}(F(x), F(y))
$$
is injective.
\item If these maps are all bijective then
$F$ is called
{\it fully faithful}.
\item
The functor $F$ is called {\it essentially
surjective} if for any
object $y \in \Ob(\mathcal{B})$ there exists
an object
$x \in \Ob(\mathcal{A})$ such that $F(x)$ is
isomorphic to $y$ in
$\mathcal{B}$.
\end{enumerate}
\end{definition}

\begin{definition}
\label{categories-definition-subcategory}
A {\it subcategory} of a category
$\mathcal{B}$ is a category $\mathcal{A}$
whose objects and arrows form subsets of the
objects and arrows of $\mathcal{B}$
and such that source, target and composition
in $\mathcal{A}$ agree with those
of $\mathcal{B}$ and such that the identity
morphism of an object of
$\mathcal{A}$ matches the one in
$\mathcal{B}$. We say $\mathcal{A}$ is a
{\it full subcategory} of $\mathcal{B}$ if
$\Mor_\mathcal{A}(x, y)
= \Mor_\mathcal{B}(x, y)$ for all $x, y \in
\Ob(\mathcal{A})$.
We say $\mathcal{A}$ is a {\it strictly full}
subcategory of $\mathcal{B}$
if it is a full subcategory and given $x \in
\Ob(\mathcal{A})$ any
object of $\mathcal{B}$ which is isomorphic
to $x$ is also in $\mathcal{A}$.
\end{definition}

\noindent
If $\mathcal{A} \subset \mathcal{B}$ is a
subcategory then the
identity map is a functor from $\mathcal{A}$
to $\mathcal{B}$.
Furthermore a subcategory $\mathcal{A}
\subset \mathcal{B}$
is full if and only if the inclusion functor
is fully faithful.
Note that given a category $\mathcal{B}$ the
set of full subcategories
of $\mathcal{B}$ is the same as the set of
subsets of
$\Ob(\mathcal{B})$.

\begin{example}
\label{categories-example-category-over-X}
Given a category $\mathcal{C}$ and an object
$X\in \Ob(\mathcal{C})$
we define the {\it category of objects over
$X$},
denoted $\mathcal{C}/X$ as follows.
The objects of $\mathcal{C}/X$ are morphisms
$Y\to X$ for
some $Y\in \Ob(\mathcal{C})$. Morphisms
between objects
$Y\to X$ and $Y'\to X$ are morphisms $Y\to
Y'$ in $\mathcal{C}$ that
make the obvious diagram commute. Note that
there is a functor
$p_X : \mathcal{C}/X\to \mathcal{C}$ which
simply forgets the
morphism. Moreover given a morphism $f :
X'\to X$ in
$\mathcal{C}$ there is an induced functor
$F : \mathcal{C}/X' \to \mathcal{C}/X$
obtained by composition with $f$,
and $p_X\circ F = p_{X'}$.
\end{example}

\begin{example}
\label{categories-example-category-under-X}
Given a category $\mathcal{C}$ and an object
$X\in \Ob(\mathcal{C})$
we define the {\it category of objects under
$X$},
denoted $X/\mathcal{C}$ as follows.
The objects of $X/\mathcal{C}$ are morphisms
$X\to Y$ for
some $Y\in \Ob(\mathcal{C})$. Morphisms
between objects
$X\to Y$ and $X\to Y'$ are morphisms $Y\to
Y'$ in $\mathcal{C}$ that
make the obvious diagram commute. Note that
there is a functor
$p_X : X/\mathcal{C}\to \mathcal{C}$ which
simply forgets the
morphism. Moreover given a morphism $f :
X'\to X$ in
$\mathcal{C}$ there is an induced functor
$F : X/\mathcal{C} \to X'/\mathcal{C}$
obtained by composition with $f$,
and $p_{X'}\circ F = p_X$.
\end{example}

\begin{definition}
\label{categories-definition-equivalence-of-categories}
Two categories are {\it equivalent} if there
exist functors
from one to the other and backwards which
composed with
each other are the corresponding identities.
\end{definition}

\begin{lemma}
\label{categories-lemma-equivalence-categories}
A functor is an equivalence of categories if
and only if it is both fully
faithful and essentially surjective.
\end{lemma}

\begin{proof}
Let $F : \mathcal{A} \to \mathcal{B}$ be
essentially surjective and fully
faithful. As by convention all categories are
small and as $F$ is essentially
surjective we can, using the axiom of choice,
choose for every
$X \in \Ob(\mathcal{B})$ an object $j(X)$ of
$\mathcal{A}$ and an
isomorphism $i_X : X \to F(j(X))$. Then we
apply
Lemma \ref{categories-lemma-construct-quasi-inverse}
using that $F$ is fully faithful.
\end{proof}

\section{Monomorphisms}
\label{categories-section-monom-epimo}

\begin{definition}
\label{categories-definition-mono-epi}
Let $\mathcal{C}$ be a category and let $f :
X \to Y$ be
a morphism of $\mathcal{C}$.
\begin{enumerate}
\item We say that $f$ is a {\it monomorphism}
if for every object
$W$ and every pair of morphisms $a, b : W \to
X$ such that
$f \circ a = f \circ b$ we have $a = b$.
\item We say that $f$ is an {\it epimorphism}
if for every object
$W$ and every pair of morphisms $a, b : Y \to
W$ such that
$a \circ f = b \circ f$ we have $a = b$.
\end{enumerate}
\end{definition}

\begin{definition}
\label{categories-definition-presh-injec-surje}
Let $\mathcal{C}$ be a category, and let
$\varphi : \mathcal{F}
\to \mathcal{G}$ be a map of presheaves of
sets.
\begin{enumerate}
\item We say that $\varphi$ is {\it
injective} if for every object
$U$ of $\mathcal{C}$ the map $\varphi_U :
\mathcal{F}(U)
\to \mathcal{G}(U)$ is injective.
\item We say that $\varphi$ is {\it
surjective} if for every object
$U$ of $\mathcal{C}$ the map $\varphi_U :
\mathcal{F}(U)
\to \mathcal{G}(U)$ is surjective.
\end{enumerate}
\end{definition}

\begin{lemma}
\label{categories-lemma-mono-epi}
The injective (resp.\ surjective) maps
defined above
are exactly the monomorphisms (resp.\
epimorphisms) of
$\textit{PSh}(\mathcal{C})$. A map is an
isomorphism
if and only if it is both injective and
surjective.
\end{lemma}

\section{Presheaves}
\label{categories-section-presheaves}

\begin{definition}
\label{categories-definition-presheaves-sets}
A {\it presheaf of sets} on $\mathcal{C}$ is
a contravariant
functor from $\mathcal{C}$ to
$\textit{Sets}$. {\it Morphisms
of presheaves} are natural transformations of
functors.
The category
of presheaves of sets is denoted
$\textit{PSh}(\mathcal{C})$
or $\hat{\mathcal{C}}$.
\end{definition}

\section{Yoneda lemma}
\label{categories-section-Yoneda-lemma}

The Yoneda lemma says that the sections over
$a$
of a presheaf $X$ can be described completely
as natural transformations between the
presheaf $\Hom(-,a)$ and $X$.

\begin{definition}
\label{categories-definition-Yoneda-embedding}
\begin{reference}
\cite[Definition 1.1.3]{Cisinsky}
\end{reference}
Let $A$ be a category.
The {\it Yoneda embedding} is the functor
$$
h:A \to \hat{A}
$$
whose value at an object $a$ of $A$ is the
presheaf
$$
h_a=\Hom_A(-,a).
$$
\end{definition}

In other words, the evaluation of the
presheaf
$h_a$ at an object $c$ of $A$ is the set of
maps from $c$ to $a$.

\begin{theorem}[Yoneda lemma]
\label{categories-theorem-Yoneda-lemma}
\begin{reference}
\cite[Theorem 1.1.4]{Cisinsky}
\end{reference}
For any presheaf $X$ over $A$, there is a
natural bijection
(in $\text{Sets}$ I think!!)
$$
\begin{aligned}
\Hom_{\widehat{A}}(h_a,X)
&\xrightarrow{\simeq} X_a  \\
(h_a \xrightarrow{u}X) &\longmapsto u_a(1_a)
\end{aligned}
$$

\end{theorem}

\section{Internal Hom}
\label{categories-section-internal-hom}

\noindent
{\bf Upshot.} Internal Hom is when the
Hom set of two objects in some category
is in also an object of the category.
Down-to-earth, that for two sheaves
$\mathcal{F},\mathcal{G}$,
$U\mapsto \Hom(\mathcal{F}(U),\mathcal{G}(U)$
is also a sheaf, called $\SheafHom$.

\medskip\noindent
I start with Stacks Project approach.

Let $(X, \mathcal{O}_X)$ be a ringed space.
Let $\mathcal{F}$, $\mathcal{G}$ be
$\mathcal{O}_X$-modules.
Consider the rule
$$
U \longmapsto
\Hom_{\mathcal{O}_X|_U}(\mathcal{F}|_U,
\mathcal{G}|_U).
$$
It follows from the discussion in Sheaves,
Section
\ref{categories-sheaves-section-glueing-sheaves} that
this is a sheaf of
abelian groups. In addition, given an element
$\varphi \in
\Hom_{\mathcal{O}_X|_U}(\mathcal{F}|_U,
\mathcal{G}|_U)$
and a section $f \in \mathcal{O}_X(U)$ then
we can define
$f\varphi \in
\Hom_{\mathcal{O}_X|_U}(\mathcal{F}|_U,
\mathcal{G}|_U)$
by either precomposing with multiplication by
$f$ on $\mathcal{F}|_U$
or postcomposing with multiplication by $f$
on $\mathcal{G}|_U$ (it gives
the same result). Hence we in fact get a
sheaf of $\mathcal{O}_X$-modules.
We will denote this sheaf
$\SheafHom_{\mathcal{O}_X}(\mathcal{F},
\mathcal{G})$.
There is a canonical ``evaluation'' morphism
$$
\mathcal{F}
\otimes_{\mathcal{O}_X}
\SheafHom_{\mathcal{O}_X}(\mathcal{F},
\mathcal{G})
\longrightarrow
\mathcal{G}.
$$
For every $x \in X$ there is also a canonical
morphism
$$
\SheafHom_{\mathcal{O}_X}(\mathcal{F},
\mathcal{G})_x
\to
\Hom_{\mathcal{O}_{X, x}}(\mathcal{F}_x,
\mathcal{G}_x)
$$
which is rarely an isomorphism.

\medskip\noindent
{\bf Cartesian closed category} In the
category of sets there is a bijection
$\Hom(X\times Y, Z)\cong \Hom(X,
\text{Hom}(Y, Z))$ that depends
naturally on $X$, $Y$ and $Z$. The notions
related to this bijection are
“Cartesian closed category”, “currying” and
“internal Hom”.

\begin{definition}
\label{categories-definition-Cartesian-closed-category}
A category $\mathcal{C}$ is {\it Cartesian
closed} if:
\begin{enumerate}
\item $\mathcal{C}$ has all finite products
(Caveat: some require that $\mathcal{C}$ has
all finite limits)
\item For any object $Y$ the functor $-
\times Y$ has a right adjoint,
which we will denote by $\text{Map}(Y,-)$ or
by $-^Y$ .
\end{enumerate}
\end{definition}

\begin{remark}
\label{categories-remark-currying}
By section 3
\href{https://ncatlab.org/nlab/show/internal%
+hom
}{here}, the
second property above implies that we get a
functor
$\text{Map}(-,-) :\mathcal{C}^{\text{op}}
\times \mathcal{C} \to \mathcal{C}$,
and moreover we get natural isomorphisms
$\Hom(X, \text{Map}(Y, Z)) \cong \Hom(X
\times Y, Z)$
and $\text{Map}(X, \text{Map}(Y, Z))\cong
\text{Map}(X \times Y, Z)$.
\end{remark}

\section{Product}
\label{categories-section-product}

\begin{definition}
\label{categories-definition-product}
A {\it product} in a category $\mathcal{C}$
is
$$
\xymatrix{
P\ar[d]_b \ar[r]^a&A\\
B
}
$$
such that for every other
$$
\xymatrix{
P'\ar[d]_{b'} \ar[r]^{a'}&A\\
B
}
$$
there exists a unique map $p:P'\to P$ such
that
$ap=a'$ and  $bp=b'$.

\end{definition}

For completeness here's Stacks Project
formulation:

\begin{definition}
\label{categories-definition-products}
Let $x, y\in \Ob(\mathcal{C})$.
A {\it product} of $x$ and $y$ is
an object $x \times y \in \Ob(\mathcal{C})$
together with morphisms
$p\in \Mor_{\mathcal C}(x \times y, x)$ and
$q\in\Mor_{\mathcal C}(x \times y, y)$ such
that the following universal property holds:
for
any $w\in \Ob(\mathcal{C})$ and morphisms
$\alpha \in \Mor_{\mathcal C}(w, x)$ and
$\beta \in \Mor_\mathcal{C}(w, y)$
there is a unique
$\gamma\in \Mor_{\mathcal C}(w, x \times y)$
making
the diagram
$$
\xymatrix{
w \ar[rrrd]^\beta \ar@{-->}[rrd]_\gamma
\ar[rrdd]_\alpha & & \\
& & x \times y \ar[d]_p \ar[r]_q & y \\
& & x &
}
$$
commute.
\end{definition}

And some nice piece of information and a
definition also from Stacks Project.

\medskip\noindent
If a product exists it is unique up to unique
isomorphism. This follows from the Yoneda
lemma as
the definition requires $x \times y$ to be an
object
of $\mathcal{C}$ such that
$$
h_{x \times y}(w) = h_x(w) \times h_y(w)
$$
functorially in $w$. In other words the
product $x \times y$
is an object representing the functor
$w \mapsto h_x(w) \times h_y(w)$.

\begin{definition}
\label{categories-definition-has-products-of-pairs}
We say the category $\mathcal{C}$ {\it has
products of pairs
of objects} if a product $x \times y$
exists for any $x, y \in \Ob(\mathcal{C})$.
\end{definition}

\section{Coproducts of pairs}
\label{categories-section-coproducts-pairs}

\begin{definition}
\label{categories-definition-coproducts}
Let $x, y \in \Ob(\mathcal{C})$.
A {\it coproduct}, or {\it amalgamated sum}
of $x$ and $y$ is
an object $x \amalg y \in \Ob(\mathcal{C})$
together with morphisms
$i \in \Mor_{\mathcal C}(x, x \amalg y)$ and
$j \in \Mor_{\mathcal C}(y, x \amalg y)$ such
that the following universal property holds:
for
any $w \in \Ob(\mathcal{C})$ and morphisms
$\alpha \in \Mor_{\mathcal C}(x, w)$ and
$\beta \in \Mor_\mathcal{C}(y, w)$
there is a unique
$\gamma \in \Mor_{\mathcal C}(x \amalg y, w)$
making
the diagram
$$
\xymatrix{
& y \ar[d]^j \ar[rrdd]^\beta \\
x \ar[r]^i \ar[rrrd]_\alpha & x \amalg y
\ar@{-->}[rrd]^\gamma \\
& & & w
}
$$
commute.
\end{definition}

\noindent
If a coproduct exists it is unique up to
unique
isomorphism. This follows from the Yoneda
lemma (applied to the
opposite category) as
the definition requires $x \amalg y$ to be an
object
of $\mathcal{C}$ such that
$$
\Mor_\mathcal{C}(x \amalg y, w) =
\Mor_\mathcal{C}(x, w) \times
\Mor_\mathcal{C}(y, w)
$$
functorially in $w$.

\begin{definition}
\label{categories-definition-has-coproducts-of-pairs}
We say the category $\mathcal{C}$ {\it has
coproducts of pairs
of objects} if a coproduct $x \amalg y$
exists for any $x, y \in \Ob(\mathcal{C})$.
\end{definition}
\section{Fibre products}
\label{categories-section-fibre-products}

\begin{definition}
\label{categories-definition-fibre-products}
Let $x, y, z\in \Ob(\mathcal{C})$,
$f\in \Mor_\mathcal{C}(x, y)$
and $g\in \Mor_{\mathcal C}(z, y)$.
A {\it fibre product} of $f$ and $g$ is
an object $x \times_y z\in \Ob(\mathcal{C})$
together with morphisms
$p \in \Mor_{\mathcal C}(x \times_y z, x)$
and
$q \in \Mor_{\mathcal C}(x \times_y z, z)$
making the diagram
$$
\xymatrix{
x \times_y z \ar[r]_q \ar[d]_p & z \ar[d]^g
\\
x \ar[r]^f & y
}
$$
commute, and such that the following
universal property holds: for
any $w\in \Ob(\mathcal{C})$ and morphisms
$\alpha \in \Mor_{\mathcal C}(w, x)$ and
$\beta \in \Mor_\mathcal{C}(w, z)$ with
$f \circ \alpha = g \circ \beta$
there is a unique
$\gamma \in \Mor_{\mathcal C}(w, x \times_y
z)$ making
the diagram
$$
\xymatrix{
w \ar[rrrd]^\beta \ar@{-->}[rrd]_\gamma
\ar[rrdd]_\alpha & & \\
& & x \times_y z \ar[d]^p \ar[r]_q & z
\ar[d]^g \\
& & x \ar[r]^f & y
}
$$
commute.
\end{definition}

\begin{remark}
\label{categories-remark-product-is-fibre-product}
I think that a product is a fibre product
when we put $y$ to be the terminal object!
\end{remark}


\noindent
If a fibre product exists it is unique up to
unique
isomorphism. This follows from the Yoneda
lemma as
the definition requires $x \times_y z$ to be
an object
of $\mathcal{C}$ such that
$$
h_{x \times_y z}(w) = h_x(w) \times_{h_y(w)}
h_z(w)
$$
functorially in $w$. In other words the fibre
product $x \times_y z$
is an object representing the functor
$w \mapsto h_x(w) \times_{h_y(w)} h_z(w)$.

\begin{definition}
\label{categories-definition-cartesian}
We say a commutative diagram
$$
\xymatrix{
w \ar[r] \ar[d] &
z \ar[d] \\
x \ar[r] &
y
}
$$
in a category is {\it cartesian} if $w$ and
the morphisms $w \to x$ and
$w \to z$ form a fibre product of the
morphisms $x \to y$ and $z \to y$.
\end{definition}

\begin{definition}
\label{categories-definition-has-fibre-products}
We say the category $\mathcal{C}$ {\it has
fibre products} if
the fibre product exists for any $f\in
\Mor_{\mathcal C}(x, y)$
and $g\in \Mor_{\mathcal C}(z, y)$.
\end{definition}

\begin{definition}
\label{categories-definition-representable-morphism}
A morphism $f : x \to y$ of a category
$\mathcal{C}$ is said to be
{\it representable} if for every morphism $z
\to y$
in $\mathcal{C}$ the fibre product $x
\times_y z$ exists.
\end{definition}

\begin{lemma}
\label{categories-lemma-composition-representable}
Let $\mathcal{C}$ be a category.
Let $f : x \to y$, and $g : y \to z$ be
representable.
Then $g \circ f : x \to z$ is representable.
\end{lemma}

\begin{proof}
Let $t \in \Ob(\mathcal C)$ and $ \varphi \in
\Mor_{\mathcal C}(t,z) $.
As $g$ and $f$ are representable, we obtain
commutative diagrams
$$
\xymatrix{
y \times_z t \ar[r]_q \ar[d]_p &
t \ar[d]^{\varphi} \\
y \ar[r]^{g} & z
}
\quad\quad
\xymatrix{
x \times_y (y\! \times_z\! t) \ar[r]_{q'}
\ar[d]_{p'} &
y \times_z t \ar[d]^p \\
x \ar[r]^f & y
}
$$
with the universal property of Definition
\ref{categories-definition-fibre-products}.
We claim that $x \times_z t = x \times_y (y
\times_z t)$
with morphisms $q \circ q' : x \times_z t \to
t$ and $p' : x \times_z t \to x$
is a fibre product. First, it follows from
the commutativity of the
diagrams above that $\varphi \circ q \circ q'
= g \circ f \circ p'$.
To verify the universal property, let
$w \in \Ob(\mathcal C)$ and suppose $\alpha :
w \to x$ and
$\beta : w \to y$ are morphisms with
$\varphi \circ \beta = g \circ f \circ
\alpha$.
By definition of the fibre product, there
are unique morphisms $\delta$ and $\gamma$
such that
$$
\xymatrix{
w \ar[rrrd]^\beta \ar@{-->}[rrd]_\delta
\ar[rrdd]_{f\circ\alpha} & & \\
& & y \times_z t \ar[d]_p \ar[r]_q & t
\ar[d]^{\varphi} \\
& & y \ar[r]^{g} & z
}
$$
and
$$
\xymatrix{
w \ar[rrrd]^\delta \ar@{-->}[rrd]_\gamma
\ar[rrdd]_{\alpha} & & \\
& & x \times_y (y\!\times_z\! t) \ar[d]_{p'}
\ar[r]_{q'} &
y \times_z t \ar[d]^{p} \\
& & x \ar[r]^{f} & y
}
$$
commute. Then, $\gamma$ makes the diagram
$$
\xymatrix{
w \ar[rrrd]^\beta \ar@{-->}[rrd]_\gamma
\ar[rrdd]_{\alpha} & & \\
& & x \times_z t \ar[d]_{p'} \ar[r]_{q\circ
q'} & t \ar[d]^{\varphi} \\
& & x \ar[r]^{g\circ f} & z
}
$$
commute. To show its uniqueness, let
$\gamma'$ verify $ q\circ q'\circ\gamma' =
\beta $ and $ p'\circ \gamma' = \alpha $.
Because $\gamma$ is unique, we just
need to prove that $ q'\circ\gamma' = \delta
$ and $ p'\circ\gamma' = \alpha $
to conclude. We supposed the second equality.
For the first one, we also need
to use the uniqueness of delta. Notice that
$\delta$ is the only morphism
verifying $ q\circ\delta = \beta $ and $
p\circ\delta = f\circ\alpha $. We
already supposed that $ q\circ
(q'\circ\gamma') = \beta $. Furthermore, by
definition of the fibre product, we know that
$ f\circ p' = p\circ q' $.
Therefore:
$$
p\circ (q'\circ\gamma') =
(p\circ q')\circ\gamma' =
(f\circ p')\circ\gamma' =
f\circ (p'\circ\gamma') = f\circ\alpha.
$$
Then $ q'\circ\gamma' = \delta $, which
concludes the proof.
\end{proof}

\begin{lemma}
\label{categories-lemma-base-change-representable}
Let $\mathcal{C}$ be a category.
Let $f : x \to y$ be representable.
Let $y' \to y$ be a morphism of
$\mathcal{C}$.
Then the morphism $x' := x \times_y y' \to
y'$ is representable also.
\end{lemma}

\begin{proof}
Let $z \to y'$ be a morphism. The fibre
product
$x' \times_{y'} z$ is supposed to represent
the
functor
\begin{eqnarray*}
w & \mapsto & h_{x'}(w)\times_{h_{y'}(w)}
h_z(w) \\
& = & (h_x(w) \times_{h_y(w)} h_{y'}(w))
\times_{h_{y'}(w)} h_z(w) \\
& = & h_x(w) \times_{h_y(w)} h_z(w)
\end{eqnarray*}
which is representable by assumption.
\end{proof}

\section{Examples of fibre products}
\label{categories-section-example-fibre-products}

\noindent
In this section we list examples of fibre
products and
we describe them.

\medskip\noindent
As a really trivial first example we observe
that the category of sets has fibre products
and hence every
morphism is representable. Namely, if $f : X
\to Y$
and $g : Z \to Y$ are maps of sets then we
define
$X \times_Y Z$ as the subset of $X \times Z$
consisting
of pairs $(x, z)$ such that $f(x) = g(z)$.
The morphisms
$p : X \times_Y Z \to X$ and $q : X \times_Y
Z \to Z$ are
the projection maps $(x, z) \mapsto x$, and
$(x, z) \mapsto z$.
Finally, if $\alpha : W \to X$ and $\beta : W
\to Z$
are morphisms such that $f \circ \alpha = g
\circ \beta$
then the map $W \to X \times Z$, $w\mapsto
(\alpha(w), \beta(w))$
obviously ends up in $X \times_Y Z$ as
desired.

\medskip\noindent
In many categories whose objects are sets
endowed with certain types of
algebraic structures the fibre product of the
underlying sets also
provides the fibre product in the category.
For example, suppose
that $X$, $Y$ and $Z$ above are groups and
that $f$, $g$ are
homomorphisms of groups. Then the
set-theoretic fibre product
$X \times_Y Z$ inherits the structure of a
group, simply by
defining the product of two pairs by the
formula
$(x, z) \cdot (x', z') = (xx', zz')$. Here we
list those categories
for which a similar reasoning works.
\begin{enumerate}
\item The category $\textit{Groups}$ of
groups.
\item The category $G\textit{-Sets}$ of sets
endowed with a left $G$-action for some fixed
group $G$.
\item The category of rings.
\item The category of $R$-modules given a
ring $R$.
\end{enumerate}

\section{Pushouts}
\label{categories-section-pushouts}

\noindent
The dual notion to fibre products is that of
pushouts.

\begin{definition}
\label{categories-definition-pushouts}
Let $x, y, z\in \Ob(\mathcal{C})$,
$f\in \Mor_\mathcal{C}(y, x)$
and $g\in \Mor_{\mathcal C}(y, z)$.
A {\it pushout} of $f$ and $g$ is
an object $x\amalg_y z\in \Ob(\mathcal{C})$
together with morphisms
$p\in \Mor_{\mathcal C}(x, x\amalg_y z)$ and
$q\in\Mor_{\mathcal C}(z, x\amalg_y z)$
making the diagram
$$
\xymatrix{
y \ar[r]_g \ar[d]_f & z \ar[d]^q \\
x \ar[r]^p & x\amalg_y z
}
$$
commute, and such that the following
universal property holds:
For any $w\in \Ob(\mathcal{C})$ and morphisms
$\alpha \in \Mor_{\mathcal C}(x, w)$ and
$\beta \in \Mor_\mathcal{C}(z, w)$ with
$\alpha \circ f = \beta \circ g$ there is a
unique
$\gamma\in \Mor_{\mathcal C}(x\amalg_y z, w)$
making
the diagram
$$
\xymatrix{
y \ar[r]_g \ar[d]_f & z \ar[d]^q
\ar[rrdd]^\beta & & \\
x \ar[r]^p \ar[rrrd]^\alpha & x \amalg_y z
\ar@{-->}[rrd]^\gamma & & \\
& & & w
}
$$
commute.
\end{definition}

\noindent
It is possible and straightforward to prove
the uniqueness of the triple
$(x\amalg_y z, p, q)$ up to unique
isomorphism (if it exists) by direct
arguments. Another possibility is to think of
the pushout as the
fibre product in the opposite category,
thereby getting this uniqueness for
free from the discussion in Section
\ref{categories-section-fibre-products}.

\begin{definition}
\label{categories-definition-cocartesian}
We say a commutative diagram
$$
\xymatrix{
y \ar[r] \ar[d] & z \ar[d] \\
x \ar[r] & w
}
$$
in a category is {\it cocartesian} if $w$ and
the morphisms $x \to w$ and
$z \to w$ form a pushout of the morphisms $y
\to x$ and $y \to z$.
\end{definition}

\medskip\noindent
According to Sergei Burkin, another way
to define a pushout is as follows.
The diagram
$$
\xymatrix{
A\ar[r]\ar[d]\ar@{}[dr]|-{\ulcorner}&B\ar[d]
\\
C\ar[r]&D
}
$$
is a {\it pushout} if the diagram
$$
\xymatrix{
\mathit{Map}(D,X)\ar[r]\ar[d]\ar@{}[dr]|-
{\lrcorner}&
\mathit{Map}(B,X)\ar[d]\\
\mathit{Map}(C,D)\ar[r]&\mathit{Map}(D,X)
}
$$
is a pullback.

In fact, spheres are (homotopy?) pushouts:

$$
\xymatrix{
S^{n-1}\ar[r]\ar[d]\ar@{}[dr]|-{\ulcorner}&
\text{pt}\ar[d]\\
\text{pt}\ar[r]&S^n
}
$$
which is homotopically the same as
$$
\xymatrix{
S^{n-1}\ar[r]\ar[d]\ar@{}[dr]|-{\ulcorner}&
D^n\ar[d]\\
D^n\ar[r]&S^n,
}
$$
the familiar construction of glueing two
disks along their boundary to obtain a
sphere.


\section{Limits and colimits}
\label{categories-section-limits-and-colimits}

\noindent
A colimit is:
take an abstract diagram $\mathcal{I}$.
Realize it in the category that interests you
(for the stack this are open sets which
contain the point,
and are a poset by inclusion).
The colimit is the unique object that
there is an arrow from all the points of the
diagram
and satisfies some universal property that
any other such object there's a morphism bla
bla.
(In the case of the stack it is the unique
open set contained in all the other open
sets.
Which actually it's interesting that such a
thing
exists.)

\noindent
A product is actually a particular case of a
limit.

\medskip\noindent
Let $\mathcal{C}$ be a category. A {\it
diagram} in $\mathcal{C}$ is
simply a functor $M : \mathcal{I} \to
\mathcal{C}$. We say that
$\mathcal{I}$ is the {\it index category} or
that $M$ is an
$\mathcal{I}$-diagram. We will use the
notation $M_i$ to denote the
image of the object
$i$ of $\mathcal{I}$. Hence for $\phi : i \to
i'$ a morphism
in $\mathcal{I}$ we have $M(\phi) : M_i \to
M_{i'}$.

\begin{definition}
\label{categories-definition-limit}
A {\it limit} of the $\mathcal{I}$-diagram
$M$ in the category
$\mathcal{C}$ is given by an object
$\lim_\mathcal{I} M$ in $\mathcal{C}$
together with morphisms $p_i :
\lim_\mathcal{I} M \to M_i$ such that
\begin{enumerate}
\item for $\phi : i \to i'$ a morphism
in $\mathcal{I}$ we have $p_{i'} = M(\phi)
\circ p_i$, and
\item for any object $W$ in $\mathcal{C}$ and
any family of
morphisms $q_i : W \to M_i$ (indexed by $i
\in \Ob(\mathcal{I})$)
such that for all $\phi : i \to i'$
in $\mathcal{I}$ we have $q_{i'} = M(\phi)
\circ q_i$ there
exists a unique morphism $q : W \to
\lim_\mathcal{I} M$ such that
$q_i = p_i \circ q$ for every object $i$ of
$\mathcal{I}$.
\end{enumerate}
\end{definition}

\noindent
Limits $(\lim_\mathcal{I} M, (p_i)_{i\in
\Ob(\mathcal{I})})$ are
(if they exist)
unique up to unique isomorphism by the
uniqueness requirement
in the definition. Products of pairs, fibre
products, and equalizers are
examples of limits. The limit over the empty
diagram is a final object
of $\mathcal{C}$.
In the category of sets all limits exist.
The dual notion is that of colimits.

\begin{definition}
\label{categories-definition-colimit}
A {\it colimit} of the $\mathcal{I}$-diagram
$M$ in the category
$\mathcal{C}$ is given by an object
$\colim_\mathcal{I} M$ in $\mathcal{C}$
together with morphisms $s_i : M_i \to
\colim_\mathcal{I} M$ such that
\begin{enumerate}
\item for $\phi : i \to i'$ a morphism
in $\mathcal{I}$ we have $s_i = s_{i'} \circ
M(\phi)$, and
\item for any object $W$ in $\mathcal{C}$ and
any family of
morphisms $t_i : M_i \to W$ (indexed by $i
\in \Ob(\mathcal{I})$)
such that for all $\phi : i \to i'$
in $\mathcal{I}$ we have $t_i = t_{i'} \circ
M(\phi)$ there
exists a unique morphism $t :
\colim_\mathcal{I} M \to W$ such that
$t_i = t \circ s_i$ for every object $i$ of
$\mathcal{I}$.
\end{enumerate}
\end{definition}

\begin{definition}
\label{categories-definition-product-as-a-limit}
Suppose that $I$ is a set, and suppose given
for every $i \in I$ an
object $M_i$ of the category $\mathcal{C}$. A
{\it product}
$\prod_{i\in I} M_i$ is by definition
$\lim_\mathcal{I} M$
(if it exists)
where $\mathcal{I}$ is the category having
only identities as
morphisms and having the elements of $I$ as
objects.
\end{definition}

\noindent
An important special case is where $I =
\emptyset$ in which case the
product is a final object of the category.
The morphisms $p_i : \prod M_i \to M_i$ are
called the
{\it projection morphisms}.

\begin{definition}
\label{categories-definition-coproduct}
Suppose that $I$ is a set, and suppose given
for every $i \in I$ an
object $M_i$ of the category $\mathcal{C}$. A
{\it coproduct}
$\coprod_{i\in I} M_i$ is by definition
$\colim_\mathcal{I} M$
(if it exists) where $\mathcal{I}$ is the
category having only
identities as morphisms and having the
elements of $I$ as objects.
\end{definition}

\noindent
An important special case is where $I =
\emptyset$ in which case the
coproduct is an initial object of the
category.
Note that the coproduct comes equipped with
morphisms
$M_i \to \coprod M_i$. These are sometimes
called the
{\it coprojections}.

\medskip\noindent
There is also so-called statement
of ``commutativity of limits''.

\section{Categories fibred in groupoids}
\label{categories-section-fibred-groupoids}

\noindent
General fibration: every map on the base has
a lift that is $p$-cartesian.
(Discrete fibration: unique?)
Fibered in groupoids: all maps in the Étale
space are $p$-cartesian.

\medskip\noindent
And to my surprise, these are none other than
categories fibered in groupoids according to
Stacks Project
\href{https://stacks.math.columbia.edu/tag/%
003S}{003S}
section
on Categories fibered in groupoids.

\medskip\noindent
Let $p:\mathcal{E} \to \mathcal{B}$ be any
functor.
The idea of fibrations is to lift maps from
$\mathcal{B}$ to $\mathcal{E}$.
We will fix the target; the condition is that
for maps on $\mathcal{B}$
with target $p(y)$ there exist
(unique? that's the condition for discrete
fibrations)
maps on $\mathcal{E}$ corresponding via $p$.
Why fix the target? Because it's like a
base-change.
But not only that: we need that the map that
exists
is good; it has to be $p$-Cartesian:

\begin{definition}
\label{categories-definition-p-cartesian}
Let $p:\mathcal{E} \to \mathcal{B}$. A
morphism $f$ in  $\mathcal{E}$
is {\it $p$-cartesian} if for all $g$ and for
all $h'$
such that $p(g)=p(f) \circ h'$
there exists a unique $h$ such that
 $p(h)=h'$ and $g=f \circ h$
$$
\xymatrix{
Z \ar@{-->}[d]_h \ar[rd]^g\\
X \ar[r]_f& Y
}
$$
given that
$$
\xymatrix{
p(Z)\ar[d]_{h'}\ar[rd]^{p(g)}\\
p(X)\ar[r]_{p(f)}&p(Y).
}
$$
\end{definition}

\begin{definition}
\label{categories-definition-fibration}
$p:\mathcal{E} \to \mathcal{B}$ is a {\it
fibration} if
for all $Y$ and all $f':X' \to p(Y)$ there
exists $X$ and
$f:X \to Y$ such that
\begin{enumerate}
\item $p(f)=f'$, and
\label{categories-item-fibration1}

 \item $f$ is $p$-Cartesian.
\label{categories-item-fibration2}
\end{enumerate}
\end{definition}

\noindent
Recall Definition
\ref{categories-definition-discrete-fibration} of
discrete fibration:
$p:\mathcal{E} \to \mathcal{B}$ is a discrete
fibration if for all
$Y \in \mathcal{E}$
and for all $f':X' \to p(Y)$ there exist
{\bf unique} ($X$ and  $f:X \to Y$) such that
$p(f)=f'$.

\begin{lemma}
\label{categories-lemma-discrete-fibration-is-fibration}
If $p:\mathcal{E} \to \mathcal{B}$
is a discrete fibration then it is a
fibration,
i.e.  $f$ is also $p$-Cartesian.
\end{lemma}

\begin{proof}
Uniqueness ensures source and target coincide
in all lifts.
\end{proof}

\begin{lemma}
\label{categories-lemma-composition-p-cartesian}
Composition of $p$-Cartesian morphisms is
$p$-Cartesian.
\end{lemma}

\begin{proof}
Write the diagrams and don't forget to prove
uniqueness.
\end{proof}

\begin{lemma}
\label{categories-lemma-2-for-3-p-cartesian}
2-for-3 for $p$-Cartesian.
\end{lemma}

\begin{proof}
Write the diagrams and don't forget to prove
uniqueness.
\end{proof}

\medskip\noindent
The notion of pseudofunctor
$F:\mathcal{B}^{\text{op}} \to \mathsf{Cat}$
is like that of a functor but we ask that
$F(\text{id}_A) \cong \text{id}_{F(A)}$
and $F(f) \circ F(g) \cong F(f \circ g)$
in some weaker sense than equality.
(This is very much like it happens in
$\infty$-categories.)
It turns out that
$$
\text{Fibrations over
$\mathcal{B}$}\leftrightarrow
\text{pseudo functors
}F:\mathcal{B}^{\text{op}}\to \mathsf{Cat}.
$$

\begin{definition}
\label{categories-definition-fibred-in-groupoids}
A functor $p:\mathcal{E} \to \mathcal{B}$ is
called a
{\it fibration in groupoids}
if it is a fibration and every fiber is a
groupoid.
\end{definition}

\noindent
As I said in the beginning of this section,
I think this condition is equivalent to
saying that
every morphism in the Étale space is
cartesian
(this should be an easy exercise).

\begin{lemma}
\label{categories-lemma-fibration-in-groupoids-iff}
$p:\mathcal{E} \to \mathcal{B}$ is a
fibration in groupoids if and only if
every morphism in $\mathcal{E}$ is
$p$-cartesian and any lifting problem
$$
\xymatrix{
& e\ar[d]\\
b\ar[r]& p(e)
}
$$
has a solution.
\end{lemma}

\begin{proof}
The backward implication is easy.
From the hypotheses we get immediately that
$p$ is a fibration.
To show that the fibers are groupoids we lift
an arbitrary morphism $f$
and do lifting problem against itself.
\end{proof}

\begin{definition}
\label{categories-definition-right-fibration-simplicial}

\end{definition}

\begin{lemma}
\label{categories-lemma-fibration-nerves}
Let $X, Y$ be nerves of 1-categories.
Then $p:X \to Y$ is a right fibration if and
only if
it is a fibration in groupoids.
\end{lemma}

\begin{proof}
The condition for right fibration for $n=0$
gives existence,
for $n=1$ gives commutativty,
for $n=2$ gives uniqueness.
\end{proof}

\medskip\noindent
Now I put two other results:

\begin{lemma}
\label{categories-lemma-kan-fibration-left-right}
Kan fibration if and only if both left and
right fibration.
\end{lemma}

\begin{lemma}
\label{categories-lemma-kan-fibration-infty-groupoids}
$X$ is $\infty$-groupoid if and only if $X
\to pt$ is Kan fibration,
if and only if it is left (or right)
fibration.
\end{lemma}










% --- End categories.tex ---

% --- Begin homological-algebra.tex ---
\chapter{Homological algebra}
\label{chapter-homological-algebra}





\section{Categories of G-modules and
R-modules}
\label{homological-algebra-section-g-and-r-modules}

\noindent
Fix a group $G$ and a ring $R$. We write
$\mathit{Mod}_R$ for the category of left
$R$-modules and $R$-linear maps.

\begin{definition}
\label{homological-algebra-definition-g-module}
A {\it $G$-module} is an abelian group $A$
with a left action of $G$ by group
automorphisms. Morphisms of $G$-modules are
$G$-equivariant group homomorphisms. Their
category is denoted $\mathit{Mod}_G$.
\end{definition}

\noindent
Equivalently, a $G$-module is a left
$\mathbb{Z}[G]$-module. The convention in
Definition \ref{homological-algebra-definition-g-module} agrees
with the Stacks Project,
\href{https://stacks.math.columbia.edu/tag/%
04JP}{Tag 04JP}.

\begin{definition}
\label{homological-algebra-definition-invariants-coinvariants}
Let $A$ be a $G$-module. The {\it
invariants} and {\it coinvariants} of $A$
are, respectively,
$$
A^G
=
\{a\in A\mid ga=a\quad\forall g\in G\}
$$
and
$$
A_G
=
A\Big/\sum_{g\in G}(g-1)A.
$$
\end{definition}

\noindent
The invariants in Definition
\ref{homological-algebra-definition-invariants-coinvariants}
are the elements already fixed by $G$.
Taking coinvariants forces $ga$ and $a$
to become equal. See the
Stacks Project,
\href{https://stacks.math.columbia.edu/tag/%
0A2H}{Tag 0A2H} and
\href{https://stacks.math.columbia.edu/tag/%
09WT}{Tag 09WT}.

\begin{example}
\label{homological-algebra-example-invariants-quotient}
Let $G=\{1,s\}$ be the group of order two.
Let $A=\mathbb{Z}$, with $sa=-a$, and let
$B=2\mathbb{Z}$. Then
$$
A^G=B^G=0,
$$
while $G$ acts trivially on
$A/B\cong\mathbb{Z}/2\mathbb{Z}$. Hence
$$
(A/B)^G\ne A^G/B^G.
$$
\end{example}

\noindent
Example \ref{homological-algebra-example-invariants-quotient}
makes the following comment from the
lecture concrete.

\begin{quote}
``Understanding natural operations
(lie $(-)^G$), applied to natural
constructions
(like quotient)
and the failure of the most naive
guess (namely $(A/B)^G=A^G/B^G$)
is kind of what homological algebra
is about.''
\end{quote}

\section{Preadditive and additive categories}
\label{homological-algebra-section-additive-categories}

\noindent
{\bf Note.} Apart from the standard
definition, I highlight from Jethro's
lecture that he proved 
Lemma \ref{homological-algebra-lemma-preadditive-direct-sum}.


\medskip\noindent
Here is the definition of a preadditive
category.

\begin{definition}
\label{homological-algebra-definition-preadditive}
A category $\mathcal{A}$ is called {\it
preadditive} if each
morphism set $\Mor_\mathcal{A}(x, y)$ is
endowed
with the structure of an abelian group such
that the
compositions
$$
\Mor(x, y) \times \Mor(y, z)
\longrightarrow
\Mor(x, z)
$$
are bilinear. A functor $F : \mathcal{A} \to
\mathcal{B}$ of
preadditive categories is called {\it
additive} if and only
if $F : \Mor(x, y) \to \Mor(F(x), F(y))$
is a homomorphism of abelian groups for all
$x, y \in \Ob(\mathcal{A})$.
\end{definition}

\noindent
In particular for every $x, y$ there exists
at least
one morphism $x \to y$, namely the zero map.

\begin{lemma}
\label{homological-algebra-lemma-preadditive-zero}
Let $\mathcal{A}$ be a preadditive category.
Let $x$ be an object of $\mathcal{A}$.
The following are equivalent
\begin{enumerate}
\item $x$ is an initial object,
\item $x$ is a final object, and
\item $\text{id}_x = 0$ in
$\Mor_\mathcal{A}(x, x)$.
\end{enumerate}
Furthermore, if such an object $0$ exists,
then a morphism
$\alpha : y \to z$ factors through $0$ if and
only if $\alpha = 0$.
\end{lemma}

\begin{proof}
First assume that $x$ is either (1) initial
or (2) final.
In both cases, it follows that $\Mor(x,x)$ is
a trivial abelian group
containing $\text{id}_x$, thus $\text{id}_x =
0$ in
$\Mor(x, x)$, which shows that each of (1)
and (2) implies (3).

\medskip\noindent
Now assume that $\text{id}_x = 0$ in
$\Mor(x,x)$. Let $y$
be an arbitrary object of $\mathcal{A}$ and
let $f \in \Mor(x ,y)$.
Denote $C : \Mor(x,x) \times \Mor(x,y) \to
\Mor(x,y)$ the composition map.
Then $f = C(0, f)$ and since $C$ is bilinear
we have $C(0, f) = 0$.
Thus $f = 0$. Hence $x$ is initial in
$\mathcal{A}$.
A similar argument for $f \in \Mor(y, x)$ can
be used to show that
$x$ is also final. Thus (3) implies both (1)
and (2).
\end{proof}

\begin{definition}
\label{homological-algebra-definition-zero-object}
In a preadditive category $\mathcal{A}$, we
call an object that is both
final and initial (as in Lemma
\ref{homological-algebra-lemma-preadditive-zero}) a
{\it zero object} and denote it $0$.
\end{definition}

\begin{lemma}
\label{homological-algebra-lemma-preadditive-direct-sum}
Let $\mathcal{A}$ be a preadditive category.
Let $x, y \in \Ob(\mathcal{A})$.
If the product $x \times y$ exists, then so
does
the coproduct $x \amalg y$.
If the coproduct $x \amalg y$ exists, then so
does
the product $x \times y$. In this case
also $x \amalg y \cong x \times y$.
\end{lemma}

\begin{proof}
Suppose that $z = x \times y$ with
projections
$p : z \to x$ and $q : z \to y$. Denote $i :
x \to z$
the morphism corresponding to $(1, 0)$.
Denote $j : y \to z$
the morphism corresponding to $(0, 1)$. Thus
we have the
commutative diagram
$$
\xymatrix{
x \ar[rr]^1 \ar[rd]^i & & x \\
& z \ar[ru]^p \ar[rd]^q & \\
y \ar[rr]^1 \ar[ru]^j & & y
}
$$
where the diagonal compositions are zero. It
follows that
$i \circ p + j \circ q : z \to z$ is the
identity since
it is a morphism which upon composing with
$p$ gives $p$
and upon composing with $q$ gives $q$.
Suppose given morphisms $a : x \to w$ and $b
: y \to w$.
Then we can form the map $a \circ p + b \circ
q : z \to w$.
In this way we get a bijection $\Mor(z, w)
= \Mor(x, w) \times \Mor(y, w)$ which
show that $z = x \amalg y$.

\medskip\noindent
We leave it to the reader to construct the
morphisms
$p, q$ given a coproduct $x \amalg y$ instead
of a
product.
\end{proof}

\begin{definition}
\label{homological-algebra-definition-direct-sum}
Given a pair of objects $x, y$ in a
preadditive category $\mathcal{A}$,
the {\it direct sum} $x \oplus y$ of $x$ and
$y$ is the direct
product $x \times y$ endowed with the
morphisms
$i, j, p, q$ as in Lemma
\ref{homological-algebra-lemma-preadditive-direct-sum} above.
\end{definition}

\begin{remark}
\label{homological-algebra-remark-direct-sum}
Note that the proof of Lemma
\ref{homological-algebra-lemma-preadditive-direct-sum}
shows that given $p$ and $q$ the morphisms
$i$, $j$ are uniquely
determined by the rules $p \circ i =
\text{id}_x$,
$q \circ j = \text{id}_y$, $p \circ j = 0$,
$q \circ i = 0$.
Moreover, we automatically have
$i \circ p + j \circ q = \text{id}_{x \oplus
y}$.
Similarly, given $i$, $j$ the morphisms $p$
and $q$ are uniquely determined.
Finally, given objects $x, y, z$ and
morphisms
$i : x \to z$, $j : y \to z$, $p : z \to x$
and
$q : z \to y$ such that $p \circ i =
\text{id}_x$,
$q \circ j = \text{id}_y$, $p \circ j = 0$,
$q \circ i = 0$
and $i \circ p + j \circ q = \text{id}_z$,
then $z$
is the direct sum of $x$ and $y$ with the
four morphisms
equal to $i, j, p, q$.
\end{remark}

\begin{lemma}
\label{homological-algebra-lemma-additive-additive}
Let $\mathcal{A}$, $\mathcal{B}$ be
preadditive categories.
Let $F : \mathcal{A} \to \mathcal{B}$ be an
additive functor.
Then $F$ transforms direct sums to direct
sums and zero to zero.
\end{lemma}

\begin{proof}
Suppose $F$ is additive. A direct sum $z$
of $x$ and $y$ is characterized by having
morphisms
$i : x \to z$, $j : y \to z$, $p : z \to x$
and
$q : z \to y$ such that $p \circ i =
\text{id}_x$,
$q \circ j = \text{id}_y$, $p \circ j = 0$,
$q \circ i = 0$
and $i \circ p + j \circ q = \text{id}_z$,
according
to Remark \ref{homological-algebra-remark-direct-sum}. Clearly
$F(x), F(y), F(z)$
and the morphisms $F(i), F(j), F(p), F(q)$
satisfy exactly the
same relations (by additivity) and we see
that $F(z)$ is
a direct sum of $F(x)$ and $F(y)$.
Hence, $F$ transforms direct sums to direct
sums.

\medskip\noindent
To see that $F$ transforms zero to zero, use
the
characterization (3) of the zero object in
Lemma \ref{homological-algebra-lemma-preadditive-zero}.
\end{proof}

\begin{definition}
\label{homological-algebra-definition-additive-category}
A category $\mathcal{A}$ is called {\it
additive}
if it is preadditive and finite products
exist, in other
words it has a zero object and direct sums.
\end{definition}

\noindent
Namely the empty product is a finite product
and
if it exists, then it is a final object.

\begin{definition}
\label{homological-algebra-definition-kernel}
Let $\mathcal{A}$ be a preadditive category.
Let $f : x \to y$ be a morphism.
\begin{enumerate}
\item A {\it kernel} of $f$ is a morphism
$i : z \to x$ such that (a) $f \circ i = 0$
and (b)
for any $i' : z' \to x$ such that $f \circ i'
= 0$ there
exists a unique morphism $g : z' \to z$ such
that
$i' = i \circ g$.
\item If the kernel of $f$ exists, then we
denote
this $\Ker(f) \to x$.
\item A {\it cokernel} of $f$ is a morphism
$p : y \to z$ such that (a) $p \circ f = 0$
and (b)
for any $p' : y \to z'$ such that $p' \circ f
= 0$ there
exists a unique morphism $g : z \to z'$ such
that
$p' = g \circ p$.
\item If a cokernel of $f$ exists we denote
this
$y \to \Coker(f)$.
\item If a kernel of $f$ exists, then a {\it
coimage
of $f$} is a cokernel for the morphism
$\Ker(f) \to x$.
\item If a kernel and coimage exist then we
denote this
$x \to \Coim(f)$.
\item If a cokernel of $f$ exists, then the
{\it image of
$f$} is a kernel of the morphism $y \to
\Coker(f)$.
\item If a cokernel and image of $f$ exist
then we denote
this $\Im(f) \to y$.
\end{enumerate}
\end{definition}

\noindent
In the above definition, we have spoken of
``the kernel'' and
``the cokernel'', tacitly using their
uniqueness
up to unique isomorphism. This follows from
the Yoneda lemma
(Categories, Section
\ref{categories-section-opposite})
because the kernel of $f : x \to y$
represents the functor
sending an object $z$ to the set
$\Ker(\Mor_\mathcal{A}(z, x) \to
\Mor_\mathcal{A}(z, y))$.
The case of cokernels is dual.

\medskip\noindent
We first relate the direct sum to kernels as
follows.

\begin{lemma}
\label{homological-algebra-lemma-additive-cat-biproduct-kernel}
Let $\mathcal{C}$ be a preadditive category.
Let $x \oplus y$ with morphisms $i, j, p, q$
as in
Lemma \ref{homological-algebra-lemma-preadditive-direct-sum}
be a direct sum in $\mathcal{C}$. Then $i : x
\to x \oplus y$
is a kernel of $q : x \oplus y \rightarrow
y$. Dually, $p$ is
a cokernel for $j$.
\end{lemma}

\begin{proof}
Let $f : z' \to x \oplus y$ be a morphism
such that $q \circ f = 0$.
We have to show that there exists a unique
morphism $g : z' \to x$
such that $f = i \circ g$. Since $i \circ p +
j \circ q$ is the identity on
$x \oplus y$ we see that
$$
f = (i \circ p + j \circ q) \circ f = i \circ
p \circ f
$$
and hence $g = p \circ f$ works. Uniqueness
holds because $p \circ i$
is the identity on $x$. The proof of the
second statement is dual.
\end{proof}

\begin{lemma}
\label{homological-algebra-lemma-kernel-mono}
Let $\mathcal{C}$ be a preadditive category.
Let $f : x \to y$ be a morphism in
$\mathcal{C}$.
\begin{enumerate}
\item If a kernel of $f$ exists, then
this kernel is a monomorphism.
\item If a cokernel of $f$ exists, then
this cokernel is an epimorphism.
\item If a kernel and coimage of $f$ exist,
then
the coimage is an epimorphism.
\item If a cokernel and image of $f$ exist,
then
the image is a monomorphism.
\end{enumerate}
\end{lemma}

\begin{proof}
Part (1) follows easily from the uniqueness
required in the
definition of a kernel. The proof of (2) is
dual.
Part (3) follows from (2), since the coimage
is a cokernel.
Similarly, (4) follows from (1).
\end{proof}

\begin{lemma}
\label{homological-algebra-lemma-coim-im-map}
Let $f : x \to y$ be a morphism in a
preadditive category
such that the kernel, cokernel, image and
coimage all exist.
Then $f$ can be factored uniquely as
$x \to \Coim(f) \to \Im(f) \to y$.
\end{lemma}

\begin{proof}
There is a canonical morphism $\Coim(f) \to
y$
because $\Ker(f) \to x \to y$ is zero.
The composition $\Coim(f) \to y \to
\Coker(f)$
is zero, because it is the unique morphism
which gives
rise to the morphism $x \to y \to \Coker(f)$
which
is zero
(the uniqueness follows from
Lemma \ref{homological-algebra-lemma-kernel-mono} (3)).
Hence $\Coim(f) \to y$ factors uniquely
through
$\Im(f) \to y$, which gives us the desired
map.
\end{proof}

\begin{example}
\label{homological-algebra-example-not-abelian}
Let $k$ be a field.
Consider the category
of filtered vector spaces over $k$.
(See Definition \ref{homological-algebra-definition-filtered}.)
Consider the filtered vector spaces $(V, F)$
and $(W, F)$ with
$V = W = k$ and
$$
F^iV
=
\left\{
\begin{matrix}
V & \text{if} & i < 0 \\
0 & \text{if} & i \geq 0
\end{matrix}
\right.
\text{ and }
F^iW
=
\left\{
\begin{matrix}
W & \text{if} & i \leq 0 \\
0 & \text{if} & i > 0
\end{matrix}
\right.
$$
The map $f : V \to W$ corresponding to
$\text{id}_k$ on the underlying
vector spaces has trivial kernel and cokernel
but is not
an isomorphism. Note also that $\Coim(f) = V$
and $\Im(f) = W$.
This means that the category of filtered
vector spaces over $k$
is not abelian.
\end{example}





\section{Abelian categories}
\label{homological-algebra-section-abelian-categories}

\noindent
An abelian category is a category satisfying
just enough axioms so the
snake lemma holds. An axiom (that is
sometimes forgotten)
is that the canonical map $\Coim(f) \to
\Im(f)$
of Lemma \ref{homological-algebra-lemma-coim-im-map} is always an
isomorphism.
Example \ref{homological-algebra-example-not-abelian} shows that
it is necessary.

Also see Remark \ref{homological-algebra-remark-snake-lemma-les}
on how the Snake Lemma is used
to prove the existence of the
induced long exact sequence
in homology.

\begin{definition}
\label{homological-algebra-definition-abelian-category}
A category $\mathcal{A}$ is {\it abelian} if
it is additive, if all kernels and cokernels
exist,
and if the natural map $\Coim(f) \to \Im(f)$
is an isomorphism for all morphisms $f$ of
$\mathcal{A}$.
\end{definition}

\begin{lemma}
\label{homological-algebra-lemma-abelian-opposite}
Let $\mathcal{A}$ be a preadditive category.
The additions on sets of morphisms make
$\mathcal{A}^{opp}$ into a preadditive
category.
Furthermore, $\mathcal{A}$ is additive if and
only if $\mathcal{A}^{opp}$
is additive, and
$\mathcal{A}$ is abelian if and only if
$\mathcal{A}^{opp}$ is abelian.
\end{lemma}

\begin{proof}
The first statement is straightforward.
To see that $\mathcal{A}$ is additive if and
only if $\mathcal{A}^{opp}$
is additive, recall that additivity can be
characterized by
the existence of a zero object and direct
sums, which are both
preserved when passing to the opposite
category.
Finally, to see that
$\mathcal{A}$ is abelian if and only if
$\mathcal{A}^{opp}$ is abelian,
observes that kernels, cokernels, images and
coimages in
$\mathcal{A}^{opp}$ correspond to
cokernels, kernels, coimages and images in
$\mathcal{A}$,
respectively.
\end{proof}

\begin{definition}
\label{homological-algebra-definition-injective-surjective}
Let $f : x \to y$ be a morphism in an abelian
category.
\begin{enumerate}
\item We say $f$ is {\it injective} if
$\Ker(f) = 0$.
\item We say $f$ is {\it surjective} if
$\Coker(f) = 0$.
\end{enumerate}
If $x \to y$ is injective, then we say that
$x$ is a {\it subobject}
of $y$ and we use the notation $x \subset y$.
If $x \to y$ is
surjective, then we say that $y$ is a {\it
quotient} of $x$.
\end{definition}

\begin{lemma}
\label{homological-algebra-lemma-characterize-injective}
Let $f : x \to y$ be a morphism in an abelian
category $\mathcal{A}$. Then
\begin{enumerate}
\item $f$ is injective if and only if $f$ is
a monomorphism,
\item $f$ is surjective if and only if $f$ is
an epimorphism, and
\item $f$ is an isomorphism if and only if
$f$ is injective and surjective.
\end{enumerate}
\end{lemma}

\begin{proof}
Proof of (1). Recall that $\Ker(f)$ is an
object representing the
functor sending $z$ to
$\Ker(\Mor_\mathcal{A}(z, x) \to
\Mor_\mathcal{A}(z, y))$, see
Definition \ref{homological-algebra-definition-kernel}.
Thus $\Ker(f)$ is $0$ if and only if
$\Mor_\mathcal{A}(z, x) \to
\Mor_\mathcal{A}(z, y)$
is injective for all $z$ if and only if $f$
is a monomorphism.
The proof of (2) is similar.
For the proof of (3) note that an isomorphism
is both a monomorphism
and epimorphism, which by (1), (2) proves $f$
is injective
and surjective. If $f$ is both injective and
surjective, then
$x = \Coim(f)$ and $y = \Im(f)$ whence $f$ is
an isomorphism.
\end{proof}

\noindent
In an abelian category, if $x \subset y$ is a
subobject,
then we denote
$$
y/x = \Coker(x \to y).
$$

\begin{lemma}
\label{homological-algebra-lemma-colimit-abelian-category}
Let $\mathcal{A}$ be an abelian category.
All finite limits and finite colimits exist
in $\mathcal{A}$.
\end{lemma}

\begin{proof}
To show that finite limits exist it suffices
to show
that finite products and equalizers exist,
see
Categories, Lemma
\ref{categories-lemma-finite-limits-exist}.
Finite products exist
by definition and the equalizer of $a, b : x
\to y$ is
the kernel of $a - b$. The argument for
finite colimits
is similar but dual to this.
\end{proof}

\begin{example}
\label{homological-algebra-example-fibre-product-pushouts}
Let $\mathcal{A}$ be an abelian category.
Pushouts and fibre products in $\mathcal{A}$
have the following
simple descriptions:
\begin{enumerate}
\item If $a : x \to y$, $b : z \to y$ are
morphisms in $\mathcal{A}$, then
we have the fibre product:
$x \times_y z = \Ker((a, -b) : x \oplus z \to
y)$.
\item If $a : y \to x$, $b : y \to z$ are
morphisms in $\mathcal{A}$, then
we have the pushout:
$x \amalg_y z = \Coker((a, -b) : y \to x
\oplus z)$.
\end{enumerate}
\end{example}

\begin{definition}
\label{homological-algebra-definition-exact}
Let $\mathcal{A}$ be an additive category.
Consider a sequence of morphisms
$$
\ldots \to x \to y \to z \to \ldots
\quad\text{or}\quad
x_1 \to x_2 \to \ldots \to x_n
$$
in $\mathcal{A}$. We say such a sequence is a
{\it complex} if the
composition of any two consecutive (drawn)
arrows is zero.
If $\mathcal{A}$ is abelian then we say a
complex of the first
type above is {\it exact at $y$} if $\Im(x
\to y) = \Ker(y \to z)$
and we say a complex of the second kind is
{\it exact at $x_i$}
where $1 < i < n$ if
$\Im(x_{i - 1} \to x_i) = \Ker(x_i \to x_{i +
1})$. We a
sequence as above is {\it exact} or is an
{\it exact sequence} or is an
{\it exact complex} if it is a complex and
exact at every object (in
the first case) or exact at $x_i$ for all $1
< i < n$ (in the second case).
There are variants of these notions for
sequences of the form
$$
\ldots \to x_{-3} \to x_{-2} \to x_{-1}
\quad\text{and}\quad
x_1 \to x_2 \to x_3 \to \ldots
$$
A {\it short exact sequence} is an exact
complex of the form
$$
0 \to A  \to B \to C \to 0.
$$
\end{definition}

\noindent
In the following lemma we assume the reader
knows what it means
for a sequence of abelian groups to be exact.

\begin{lemma}
\label{homological-algebra-lemma-check-exactness}
Let $\mathcal{A}$ be an abelian category.
Let $0 \to M_1 \to M_2 \to M_3 \to 0$ be a
complex of $\mathcal{A}$.
\begin{enumerate}
\item $M_1 \to M_2 \to M_3 \to 0$ is exact if
and only if
$$
0 \to \Hom_\mathcal{A}(M_3, N) \to
\Hom_\mathcal{A}(M_2, N) \to
\Hom_\mathcal{A}(M_1, N)
$$
is an exact sequence of abelian groups for
all objects $N$ of
$\mathcal{A}$, and
\item $0 \to M_1 \to M_2 \to M_3$ is exact if
and only if
$$
0 \to \Hom_\mathcal{A}(N, M_1) \to
\Hom_\mathcal{A}(N, M_2) \to
\Hom_\mathcal{A}(N, M_3)
$$
is an exact sequence of abelian groups for
all objects $N$ of $\mathcal{A}$.
\end{enumerate}
\end{lemma}

\begin{proof}
Omitted. Hint: See
Algebra, Lemma \ref{algebra-lemma-hom-exact}.
\end{proof}

\begin{definition}
\label{homological-algebra-definition-ses-split}
Let $\mathcal{A}$ be an abelian category.
Let $i : A \to B$ and $q : B \to C$ be
morphisms
of $\mathcal{A}$ such that
$0 \to A \to B \to C \to 0$ is a short
exact sequence. We say the short exact
sequence is {\it split} if there exist
morphisms $j : C \to B$ and $p : B \to A$
such
that $(B, i, j, p, q)$ is the direct sum of
$A$ and $C$.
\end{definition}

\begin{lemma}
\label{homological-algebra-lemma-ses-split}
Let $\mathcal{A}$ be an abelian category.
Let $0 \to A \to B \to C \to 0$
be a short exact sequence.
\begin{enumerate}
\item Given a morphism $s : C \to B$ right
inverse to
$B \to C$, there exists a unique $\pi : B \to
A$
such that $(s, \pi)$ splits the short exact
sequence
as in Definition \ref{homological-algebra-definition-ses-split}.
\item Given a morphism $\pi : B \to A$ left
inverse to
$A \to B$, there exists a unique $s : C \to
B$
such that $(s, \pi)$ splits the short exact
sequence
as in Definition \ref{homological-algebra-definition-ses-split}.
\end{enumerate}
\end{lemma}

\begin{proof}
Omitted.
\end{proof}

\begin{exercise}[Splitting lemma]
\label{homological-algebra-exercise-split-short-exact-sequence-modules}
Let
$$
0\longrightarrow A\xrightarrow{i}B\xrightarrow{q}C\longrightarrow0
$$
be a short exact sequence of $R$-modules.
\begin{reference}
\cite[Algebraic Topology, p.~147, Splitting Lemma]{Hatcher-algebraic-topology}
\end{reference}
Prove directly that the following are equivalent:
\begin{enumerate}
\item there is a retraction $p:B\to A$, that is,
$$
p\circ i=\operatorname{id}_A;
$$
\item there is a section $s:C\to B$, that is,
$$
q\circ s=\operatorname{id}_C;
$$
\item there is an isomorphism
$$
B\cong A\oplus C
$$
under which $i$ becomes $a\mapsto(a,0)$ and $q$ becomes $(a,c)\mapsto c$.
\end{enumerate}
When these equivalent conditions hold, the short exact sequence is called \emph{split}.

For $(1)\Rightarrow(3)$, try
$$
b\longmapsto\bigl(p(b),q(b)\bigr).
$$
For $(2)\Rightarrow(3)$, try
$$
(a,c)\longmapsto i(a)+s(c).
$$
\end{exercise}

\begin{lemma}
\label{homological-algebra-lemma-characterize-cartesian}
Let $\mathcal{A}$ be an abelian category. Let
$$
\xymatrix{
w\ar[r]^f\ar[d]_g
& y\ar[d]^h\\
x\ar[r]^k
& z
}
$$
be a commutative diagram.
\begin{enumerate}
\item The diagram is cartesian if and only if
$$
0 \to w \xrightarrow{(g, f)} x \oplus y
\xrightarrow{(k, -h)} z
$$
is exact.
\item The diagram is cocartesian if and only
if
$$
w \xrightarrow{(g, -f)} x \oplus y
\xrightarrow{(k, h)} z \to 0
$$
is exact.
\end{enumerate}
\end{lemma}

\begin{proof}
Let $u = (g, f) : w \to x \oplus y$ and $v =
(k, -h) : x \oplus y \to z$.
Let $p : x \oplus y \to x$ and $q : x \oplus
y \to y$ be the canonical
projections. Let $i : \Ker(v) \to x \oplus y$
be the canonical
injection. By Example
\ref{homological-algebra-example-fibre-product-pushouts}, the
diagram is
cartesian if and only if there exists an
isomorphism
$r : \Ker(v) \to w$ with $f \circ r = q \circ
i$ and
$g \circ r = p \circ i$. The sequence
$0 \to w \overset{u} \to x \oplus y
\overset{v} \to z$ is exact if and
only if there exists an isomorphism $r :
\Ker(v) \to w$ with
$u \circ r = i$. But given $r : \Ker(v) \to
w$, we have
$f \circ r = q \circ i$ and $g \circ r = p
\circ i$ if and
only if $q \circ u \circ r= f \circ r = q
\circ i$ and
$p \circ u \circ r = g \circ r = p \circ i$,
hence if and only if
$u \circ r = i$. This proves (1), and then
(2) follows by duality.
\end{proof}

\begin{lemma}
\label{homological-algebra-lemma-cartesian-kernel}
Let $\mathcal{A}$ be an abelian category. Let
$$
\xymatrix{
w\ar[r]^f\ar[d]_g
& y\ar[d]^h\\
x\ar[r]^k
& z
}
$$
be a commutative diagram.
\begin{enumerate}
\item If the diagram is cartesian, then the
morphism
$\Ker(f)\to\Ker(k)$ induced by $g$ is an
isomorphism.
\item If the diagram is cocartesian, then the
morphism
$\Coker(f)\to\Coker(k)$ induced by $h$ is an
isomorphism.
\end{enumerate}
\end{lemma}

\begin{proof}
Suppose the diagram is cartesian. Let
$e:\Ker(f)\to\Ker(k)$ be induced by $g$. Let
$i:\Ker(f)\to w$ and $j:\Ker(k)\to x$ be the
canonical
injections. There exists $t:\Ker(k)\to w$
with $f\circ t=0$
and $g\circ t=j$. Hence, there exists
$u:\Ker(k)\to\Ker(f)$
with $i\circ u=t$. It follows
$g\circ i\circ u\circ e=g\circ t\circ
e=j\circ e=g\circ i$ and
$f\circ i\circ u\circ e=0=f\circ i$, hence
$i\circ u\circ e=i$. Since
$i$ is a monomorphism this implies $u\circ
e=\text{id}_{\Ker(f)}$.
Furthermore, we have $j\circ e\circ u=g\circ
i\circ u=g\circ t=j$.
Since $j$ is a monomorphism this implies
$e\circ u=\text{id}_{\Ker(k)}$.
This proves (1). Now, (2) follows by duality.
\end{proof}

\begin{lemma}
\label{homological-algebra-lemma-cartesian-cocartesian}
Let $\mathcal{A}$ be an abelian category. Let
$$
\xymatrix{
w\ar[r]^f\ar[d]_g
& y\ar[d]^h\\
x\ar[r]^k
& z
}
$$
be a commutative diagram.
\begin{enumerate}
\item If the diagram is cartesian and $k$ is
an epimorphism,
then the diagram is cocartesian and $f$ is an
epimorphism.
\item If the diagram is cocartesian and $g$
is a monomorphism,
then the diagram is cartesian and $h$ is a
monomorphism.
\end{enumerate}
\end{lemma}

\begin{proof}
Suppose the diagram is cartesian and $k$ is
an epimorphism.
Let $u = (g, f) : w \to x \oplus y$ and let
$v = (k, -h) : x \oplus y \to z$.
As $k$ is an epimorphism, $v$ is an
epimorphism, too. Therefore
and by Lemma
\ref{homological-algebra-lemma-characterize-cartesian}, the
sequence
$0\to w\overset{u}\to x\oplus y\overset{v}\to
z\to 0$ is exact. Thus, the
diagram is cocartesian by Lemma
\ref{homological-algebra-lemma-characterize-cartesian}. Finally,
$f$ is an epimorphism by Lemma
\ref{homological-algebra-lemma-cartesian-kernel} and
Lemma \ref{homological-algebra-lemma-characterize-injective}.
This proves (1), and (2)
follows by duality.
\end{proof}

\begin{lemma}
\label{homological-algebra-lemma-epimo-unive-abeli-categ}
Let $\mathcal{A}$ be an abelian category.
\begin{enumerate}
\item If $x \to y$ is surjective, then for
every $z \to y$ the
projection $x \times_y z \to z$ is
surjective.
\item If $x \to y$ is injective, then for
every $x \to z$ the
morphism $z \to z \amalg_x y$ is injective.
\end{enumerate}
\end{lemma}

\begin{proof}
Immediately from Lemma
\ref{homological-algebra-lemma-characterize-injective} and
Lemma \ref{homological-algebra-lemma-cartesian-cocartesian}.
\end{proof}

\begin{lemma}
\label{homological-algebra-lemma-check-exactness-fibre-product}
Let $\mathcal{A}$ be an abelian category. Let
$f:x\to y$ and $g:y\to z$
be morphisms with $g\circ f=0$. Then, the
following statements are equivalent:
\begin{enumerate}
\item The sequence $x\overset{f}\to
y\overset{g}\to z$ is exact.
\item For every $h:w\to y$ with $g\circ h=0$
there exist an object $v$,
an epimorphism $k:v\to w$ and a morphism
$l:v\to x$ with $h\circ k=f\circ l$.
\end{enumerate}
\end{lemma}

\begin{proof}
Let $i:\Ker(g)\to y$ be the canonical
injection. Let
$p:x\to\Coim(f)$ be the canonical projection.
Let
$j:\Im(f)\to\Ker(g)$ be the canonical
injection.

\medskip\noindent
Suppose (1) holds. Let $h:w\to y$ with
$g\circ h=0$. There exists
$c:w\to\Ker(g)$ with $i\circ c=h$.
Let $v=x\times_{\Ker(g)}w$ with canonical
projections
$k:v\to w$ and $l:v\to x$, so that $c\circ
k=j\circ p\circ l$.
Then, $h\circ k=i\circ c\circ k=i\circ j\circ
p\circ l=f\circ l$.
As $j\circ p$ is an epimorphism by
hypothesis, $k$ is an
epimorphism by Lemma
\ref{homological-algebra-lemma-cartesian-cocartesian}. This
implies (2).

\medskip\noindent
Suppose (2) holds. Then, $g\circ i=0$. So,
there are an object
$w$, an epimorphism $k:w\to\Ker(g)$ and a
morphism
$l:w\to x$ with $f\circ l=i\circ k$. It
follows
$i\circ j\circ p\circ l=f\circ l=i\circ k$.
Since $i$ is a
monomorphism we see that $j\circ p\circ l=k$
is an epimorphism.
So, $j$ is an epimorphisms and thus an
isomorphism. This implies (1).
\end{proof}

\begin{lemma}
\label{homological-algebra-lemma-exact-kernel-sequence}
Let $\mathcal{A}$ be an abelian category. Let
$$
\xymatrix{
x \ar[r]^f \ar[d]^\alpha &
y \ar[r]^g \ar[d]^\beta &
z \ar[d]^\gamma\\
u \ar[r]^k & v \ar[r]^l & w
}
$$
be a commutative diagram.
\begin{enumerate}
\item If the first row is exact and $k$ is a
monomorphism, then the induced
sequence $\Ker(\alpha) \to \Ker(\beta) \to
\Ker(\gamma)$
is exact.
\item If the second row is exact and $g$ is
an epimorphism, then the induced
sequence
$\Coker(\alpha) \to \Coker(\beta) \to
\Coker(\gamma)$
is exact.
\end{enumerate}
\end{lemma}

\begin{proof}
Suppose the first row is exact and $k$ is a
monomorphism. Let
$a:\Ker(\alpha)\to\Ker(\beta)$ and
$b:\Ker(\beta)\to\Ker(\gamma)$ be the induced
morphisms.
Let $h:\Ker(\alpha)\to x$, $i:\Ker(\beta)\to
y$ and
$j:\Ker(\gamma)\to z$ be the canonical
injections. As $j$ is
a monomorphism we have $b\circ a=0$. Let
$c:s\to\Ker(\beta)$
with $b\circ c=0$. Then, $g\circ i\circ
c=j\circ b\circ c=0$. By
Lemma
\ref{homological-algebra-lemma-check-exactness-fibre-product}
there are an object $t$, an
epimorphism $d:t\to s$ and a morphism $e:t\to
x$ with
$i\circ c\circ d=f\circ e$. Then,
$k\circ \alpha\circ e=\beta\circ f\circ
e=\beta\circ i\circ c\circ d=0$.
As $k$ is a monomorphism we get $\alpha\circ
e=0$. So, there exists
$m:t\to\Ker(\alpha)$ with $h\circ m=e$. It
follows
$i\circ a\circ m=f\circ h\circ m=f\circ
e=i\circ c\circ d$.
As $i$ is a monomorphism we get $a\circ
m=c\circ d$. Thus,
Lemma
\ref{homological-algebra-lemma-check-exactness-fibre-product}
implies (1), and then
(2) follows by duality.
\end{proof}

\noindent
(From Jethro's lecture.)
The point of the snake lemma is the
existence of the connecting morphism.
In particular, this allows us to
construct the long exact sequence in
cohomology from a short exact sequence;
that is, not only do we get induced
morphisms but also the weird connecting
morphism.

\begin{lemma}[Snake lemma, Jethro's version]
\label{homological-algebra-lemma-snake-jethro}
In an abelian category $\mathcal{C}$,
let the following diagram be given,
with exact rows:
$$
\xymatrix{
& A \ar[r]^f \ar[d]^\alpha &
B \ar[r]^g \ar[d]^\beta &
C \ar[r] \ar[d]^\gamma &
0 \\
0 \ar[r] & A' \ar[r]^k & B' \ar[r]^l & C'
}
$$
then there is an exact sequence
$$
\xymatrix{
\ker(f)\ar[r]
&\ker(\alpha) \ar[r]
&\ker(\beta)\ar[r]
&\ker(\gamma)\ar[lld]\\
&\operatorname{coker}(\alpha)\ar[r]
&\operatorname{coker}(\beta)\ar[r]
&\operatorname{coker}(\gamma)\ar[r]
&\operatorname{coker}(g').
}
$$
\end{lemma}

\begin{proof}
Mostly routine, but exhausting.
For instance, one can check from
universal properties, why there
exists a morphism […]

The most interesting part is the 
morphism
$$
\ker(\gamma)\xrightarrow{\delta}
\operatorname{coker}(\alpha).
$$
Let's construct it. […]
\end{proof}

\begin{lemma}[Snake lemma, Stacks Project]
\label{homological-algebra-lemma-snake}
Let $\mathcal{A}$ be an abelian category. Let
$$
\xymatrix{
& x \ar[r]^f \ar[d]^\alpha &
y \ar[r]^g \ar[d]^\beta &
z \ar[r] \ar[d]^\gamma &
0 \\
0 \ar[r] & u \ar[r]^k & v \ar[r]^l & w
}
$$
be a commutative diagram with exact rows.
\begin{enumerate}
\item There exists a unique morphism
$\delta : \Ker(\gamma) \to \Coker(\alpha)$
such that the diagram
$$
\xymatrix{
y \ar[d]_\beta &
y \times_z \Ker(\gamma) \ar[l]_{\pi'}
\ar[r]^{\pi} &
\Ker(\gamma) \ar[d]^\delta \\
v \ar[r]^{\iota'} & \Coker(\alpha) \amalg_u v
&
\Coker(\alpha) \ar[l]_\iota
}
$$
commutes, where $\pi$ and $\pi'$ are the
canonical projections
and $\iota$ and $\iota'$ are the canonical
coprojections.
\item The induced sequence
$$
\Ker(\alpha) \xrightarrow{f'} \Ker(\beta)
\xrightarrow{g'}
\Ker(\gamma) \xrightarrow{\delta}
\Coker(\alpha) \xrightarrow{k'}
\Coker(\beta) \xrightarrow{l'} \Coker(\gamma)
$$
is exact. If $f$ is injective then so is
$f'$, and if $l$ is
surjective then so is $l'$.
\end{enumerate}
\end{lemma}

\begin{proof}
As $\pi$ is an epimorphism and $\iota$ is a
monomorphism by
Lemma \ref{homological-algebra-lemma-cartesian-cocartesian},
uniqueness of $\delta$ is clear.
Let $p = y \times_z \Ker(\gamma)$ and $q =
\Coker(\alpha) \amalg_u v$.
Let $h : \Ker(\beta) \to y$, $i :
\Ker(\gamma) \to z$ and
$j : \Ker(\pi) \to p$ be the canonical
injections.
Let $\pi'' : u \to \Coker(\alpha)$ be the
canonical projection.
Keeping in mind Lemma
\ref{homological-algebra-lemma-cartesian-cocartesian} we get a
commutative
diagram with exact rows
$$
\xymatrix{
0 \ar[r] &
\Ker(\pi) \ar[r]^j &
p \ar[r]^{\pi} \ar[d]_{\pi'} &
\Ker(\gamma) \ar[d]_i \ar[r] & 0 \\
& x \ar[r]^f \ar[d]_\alpha & y \ar[r]^g
\ar[d]_\beta &
z \ar[d]_\gamma \ar[r] & 0 \\
0 \ar[r] & u \ar[r]^k \ar[d]_{\pi''} &
v \ar[r]^l \ar[d]_{\iota'} & w & \\
0 \ar[r] & \Coker(\alpha) \ar[r]^\iota & q &
&
}
$$
As $l \circ \beta \circ \pi' = \gamma \circ i
\circ \pi = 0$ and as the third
row of the diagram above is exact, there is
an $a : p \to u$
with $k \circ a = \beta \circ \pi'$. As the
upper right quadrangle of the
diagram above is cartesian, Lemma
\ref{homological-algebra-lemma-cartesian-kernel} yields an
epimorphism $b : x \to \Ker(\pi)$ with $\pi'
\circ j \circ b = f$.
It follows
$k \circ a \circ j \circ b = \beta \circ \pi'
\circ j \circ b =
\beta \circ f = k \circ \alpha$.
As $k$ is a monomorphism this implies $a
\circ j \circ b = \alpha$. It follows
$\pi'' \circ a \circ j \circ b = \pi'' \circ
\alpha = 0$. As $b$ is an
epimorphism this
implies $\pi'' \circ a \circ j = 0$.
Therefore, as the top row of the diagram
above is exact, there exists
$\delta : \Ker(\gamma) \to \Coker(\alpha)$
with
$\delta \circ \pi = \pi'' \circ a$. It
follows
$\iota \circ \delta \circ \pi = \iota \circ
\pi'' \circ a =
\iota' \circ k \circ a = \iota' \circ \beta
\circ \pi'$
as desired.

\medskip\noindent
As the upper right quadrangle in the diagram
above is cartesian there
is a $c : \Ker(\beta) \to p$ with $\pi' \circ
c = h$ and $\pi \circ c = g'$.
It follows
$\iota \circ \delta \circ g' = \iota \circ
\delta \circ \pi \circ c =
\iota' \circ \beta \circ \pi' \circ c =
\iota' \circ \beta \circ h = 0$.
As $\iota$ is a monomorphism this implies
$\delta \circ g' = 0$.

\medskip\noindent
Next, let $d : r \to \Ker(\gamma)$ with
$\delta \circ d = 0$. Applying
Lemma
\ref{homological-algebra-lemma-check-exactness-fibre-product}
to the exact sequence
$p \xrightarrow{\pi} \Ker(\gamma) \to 0$ and
$d$ yields an object $s$,
an epimorphism $m : s \to r$ and a morphism
$n : s \to p$ with
$\pi \circ n = d \circ m$. As
$\pi'' \circ a \circ n = \delta \circ d \circ
m = 0$,
applying Lemma
\ref{homological-algebra-lemma-check-exactness-fibre-product} to
the exact sequence
$x \xrightarrow{\alpha} u \xrightarrow{p}
\Coker(\alpha)$ and
$a \circ n$ yields an object $t$, an
epimorphism $\varepsilon : t \to s$ and
a morphism $\zeta : t \to x$ with
$a \circ n \circ \varepsilon = \alpha \circ
\zeta$.
It holds
$\beta \circ \pi' \circ n \circ \varepsilon =
k \circ \alpha \circ \zeta = \beta \circ f
\circ \zeta$.
Let $\eta = \pi' \circ n \circ \varepsilon -
f \circ \zeta : t \to y$. Then,
$\beta \circ \eta = 0$. It follows that there
is a
$\vartheta : t \to \Ker(\beta)$ with $\eta =
h \circ \vartheta$. It holds
$i \circ g' \circ \vartheta = g \circ h \circ
\vartheta =
g \circ \pi' \circ n \circ \varepsilon - g
\circ f \circ \zeta =
i \circ \pi \circ n \circ \varepsilon = i
\circ d \circ m \circ \varepsilon$.
As $i$ is a monomorphism we get
$g' \circ \vartheta = d \circ m \circ
\varepsilon$.
Thus, as $m \circ \varepsilon$ is an
epimorphism,
Lemma
\ref{homological-algebra-lemma-check-exactness-fibre-product}
implies that
$\Ker(\beta) \xrightarrow{g'} \Ker(\gamma)
\xrightarrow{\delta} \Coker(\alpha)$
is exact. Then, the claim follows by Lemma
\ref{homological-algebra-lemma-exact-kernel-sequence}
and duality.
\end{proof}

\begin{lemma}
\label{homological-algebra-lemma-snake-natural}
Let $\mathcal{A}$ be an abelian category. Let
$$
\xymatrix{
& & & x\ar[ld]\ar[rr]\ar[dd]^(.4)\alpha
& & y\ar[ld]\ar[rr]\ar[dd]^(.4)\beta
& & z\ar[ld]\ar[rr]\ar[dd]^(.4)\gamma
& & 0\\
& & x'\ar[rr]\ar[dd]^(.4){\alpha'}
& & y'\ar[rr]\ar[dd]^(.4){\beta'}
& & z'\ar[rr]\ar[dd]^(.4){\gamma'}
& & 0
& \\
& 0\ar[rr]
& & u\ar[ld]\ar[rr]
& & v\ar[ld]\ar[rr]
& & w\ar[ld]
& & \\
0\ar[rr]
& & u'\ar[rr]
& & v'\ar[rr]
& & w'
& & &
}
$$
be a commutative diagram with exact rows.
Then, the induced diagram
$$
\xymatrix@C=15pt{
\Ker(\alpha) \ar[r] \ar[d] &
\Ker(\beta) \ar[r] \ar[d] &
\Ker(\gamma) \ar[r]^(.45){\delta} \ar[d] &
\Coker(\alpha) \ar[r] \ar[d] &
\Coker(\beta) \ar[r] \ar[d] &
\Coker(\gamma) \ar[d] \\
\Ker(\alpha') \ar[r] &
\Ker(\beta') \ar[r] &
\Ker(\gamma') \ar[r]^(.45){\delta'} &
\Coker(\alpha') \ar[r] &
\Coker(\beta') \ar[r] &
\Coker(\gamma')
}
$$
commutes.
\end{lemma}

\begin{proof}
Omitted.
\end{proof}

\begin{lemma}
\label{homological-algebra-lemma-four-lemma}
Let $\mathcal{A}$ be an abelian category. Let
$$
\xymatrix{
w \ar[r] \ar[d]^\alpha & x \ar[r]
\ar[d]^\beta & y \ar[r] \ar[d]^\gamma &
z \ar[d]^\delta \\
w' \ar[r] & x' \ar[r] & y' \ar[r] & z'
}
$$
be a commutative diagram with exact rows.
\begin{enumerate}
\item If $\alpha, \gamma$ are surjective and
$\delta$ is injective, then
$\beta$ is surjective.
\item If $\beta, \delta$ are injective and
$\alpha$ is surjective, then
$\gamma$ is injective.
\end{enumerate}
\end{lemma}

\begin{proof}
Assume $\alpha, \gamma$ are surjective and
$\delta$ is injective.
We may replace $w'$ by $\Im(w' \to x')$,
i.e., we may assume
that $w' \to x'$ is injective.
We may replace $z$ by $\Im(y \to z)$, i.e.,
we may assume that
$y \to z$ is surjective. Then we may apply
Lemma \ref{homological-algebra-lemma-snake}
to
$$
\xymatrix{
& \Ker(y \to z) \ar[r] \ar[d] & y \ar[r]
\ar[d] & z \ar[r] \ar[d] & 0 \\
0 \ar[r] & \Ker(y' \to z') \ar[r] & y' \ar[r]
& z'
}
$$
to conclude that $\Ker(y \to z) \to \Ker(y'
\to z')$ is
surjective. Finally, we apply
Lemma \ref{homological-algebra-lemma-snake}
to
$$
\xymatrix{
& w \ar[r] \ar[d] & x \ar[r] \ar[d] & \Ker(y
\to z) \ar[r] \ar[d] & 0 \\
0 \ar[r] & w' \ar[r] & x' \ar[r] & \Ker(y'
\to z')
}
$$
to conclude that $x \to x'$ is surjective.
This proves (1). The proof
of (2) is dual to this.
\end{proof}

\begin{lemma}
\label{homological-algebra-lemma-five-lemma}
\begin{reference}
\cite[Lemma 4.5 page 16]{Eilenberg-Steenrod}
\end{reference}
Let $\mathcal{A}$ be an abelian category. Let
$$
\xymatrix{
v \ar[r] \ar[d]^\alpha &
w \ar[r] \ar[d]^\beta &
x \ar[r] \ar[d]^\gamma &
y \ar[r] \ar[d]^\delta &
z \ar[d]^\epsilon \\
v' \ar[r] & w' \ar[r] & x' \ar[r] & y' \ar[r]
& z'
}
$$
be a commutative diagram with exact rows. If
$\beta, \delta$
are isomorphisms, $\epsilon$ is injective,
and $\alpha$ is surjective
then $\gamma$ is an isomorphism.
\end{lemma}

\begin{proof}
Immediate consequence of
Lemma \ref{homological-algebra-lemma-four-lemma}.
\end{proof}

\medskip\noindent
From Jethro's lecture:

\begin{example}[Non-abelian categories]
\label{homological-algebra-example-non-abelian-categories}
\begin{enumerate}
\item Fix $M$ a smooth manifold,
and consider the category of finite
dimensional vector bundles.
This is not abelian.
Let $f$ be a map of the total spaces
of two vector bundles which is linear
on every fiber, i.e.
$$
\xymatrix{
E_1\ar[rr]^{f}\ar[dr]_{p_1}
&&E_2\ar[dl]^{p_2}\\
&M
}
$$
then ker and coker don't exist in general,
because the dimension of the fibers
might jump.

\item $M=\mathbb{R}$, $E=\mathbb{R}\times
V$,
$V=\mathbb{C}$ (so one-dimensional).

Let
$$
\xymatrix{
E \ar[dr]\ar[rr]&& E \ar[dl]\\
& \mathbb{R}
}
$$
with $f(v)=zv$, $z $ coordinate on
$\mathbb{R}$.
Then $\ker(f)=(\mathbb{R}\setminus 0)
\times \coprod 0 \times \mathbb{C}$.

The notion of coherent sheaf
is motivated by this example.

\item Other silly examples like the category
of all free abelian groups:
$\mathbb{Z} \xrightarrow{3}\mathbb{Z}$
has no cokernel.
\end{enumerate}
\end{example}

\medskip\noindent
We finish with the example of filtered
vector spaces mentioned in the lecture
(this was done by Jethro but the following
notes are AI-written).

\begin{definition}
\label{homological-algebra-definition-filtered}
Let $k$ be a field. A {\it decreasing
filtration} on a vector space $V$ over $k$
is a collection of vector subspaces $F^pV$,
indexed by $p \in \mathbb{Z}$, such that
$$
F^{p + 1}V \subset F^pV
$$
for every $p$. A {\it filtered vector space}
is a pair $(V, F)$ consisting of a vector
space and a decreasing filtration on it.
A {\it morphism of filtered vector spaces}
from $(V, F)$ to $(W, G)$ is a linear map
$f : V \to W$ such that
$$
f(F^pV) \subset G^pW
$$
for every $p$. Such a morphism is called
{\it strict} if
$$
f(F^pV)=f(V)\cap G^pW
$$
for every $p$. We denote the resulting
category by $\mathit{Fil}(\mathit{Vec}_k)$.
No boundedness, exhaustiveness, or
separatedness condition is imposed on the
filtrations.
\end{definition}

\begin{lemma}
\label{homological-algebra-lemma-filtered-vector-spaces}
The category $\mathit{Fil}(\mathit{Vec}_k)$
is additive and has all kernels and
cokernels. More precisely, let
$f : (V, F) \to (W, G)$ be a morphism, and
let $K = \Ker(f)$ and $Q = \Coker(f)$ be
the kernel and cokernel of the underlying
linear map. Their filtrations are
$$
F^pK=K\cap F^pV
$$
and
$$
G^pQ=(G^pW+f(V))/f(V).
$$
The canonical morphism
$$
\Coim(f)\longrightarrow\Im(f)
$$
is an isomorphism if and only if $f$ is
strict. Consequently,
$\mathit{Fil}(\mathit{Vec}_k)$ is not an
abelian category.
\end{lemma}

\begin{proof}
Addition of morphisms is inherited from
linear maps, and composition is bilinear.
The zero vector space with its unique
filtration is a zero object. Direct sums
are formed with the filtration
$$
H^p(V\oplus W)=F^pV\oplus G^pW.
$$
Hence the category is additive. The usual
kernel map is
filtration preserving for the filtration on
$K$ displayed above. If a morphism of
filtered vector spaces has zero composition
with $f$, its factorization through $K$ is
filtration preserving. Thus this is the
kernel in $\mathit{Fil}(\mathit{Vec}_k)$.

The filtration on $Q$ is the image of
$G^pW$ under the quotient map $W \to Q$.
The usual factorization through $Q$ of a
filtration preserving map that annihilates
$f(V)$ is again filtration preserving. Thus
this is the cokernel.

The filtration on the coimage corresponds,
under the usual vector space isomorphism
from the coimage to $f(V)$, to the subspaces
$f(F^pV)$. The filtration on the image is
given by the subspaces
$$
f(V)\cap G^pW.
$$
Hence the canonical morphism from the
coimage to the image is an isomorphism
exactly when $f$ is strict. The morphism in
Example \ref{homological-algebra-example-not-abelian} is not
strict, so the category is not abelian.
\end{proof}

\section{Tensor product}
\label{homological-algebra-section-tensor-product}

If the direct sum $A \oplus B$
'categorifies' $+$, surely what
categorifies $\cdot$ is $\otimes$,
the tensor product. Recall
that for vector spaces $U, V$
with bases $\{u_i\}, \{v_j\}$,
then $U \otimes V$ has basis
$\{u_i \otimes v_j\}$
(and has dimension $\dim(U \otimes V)
=\dim(U)\dim(V)$).

Universal property? The idea is that
a morphism $A \otimes B \to X$
is the same as a bilinear map
$A \times B \to X$.

$A \otimes B$ is caracterized by
the existence of a bijection
$\tilde{f} \leftrightarrow f$
of $\tilde{f}:A \otimes B \to X$,
and $f \in \operatorname{Bil}_{A,B}(X)$
such that… (exercise: such that what?).

\medskip\noindent
Another perspective is: for $T=A \otimes B$
fixed, we want the functos
$$
\Hom_{\underline{\operatorname{Ab}}}
(T,-):\underline{\operatorname{Ab}}
\to \underline{\operatorname{Set}}
$$
and
$$
\operatorname{Bil}_{A,B}(-):
\underline{\operatorname{Ab}}
\to \underline{\operatorname{Set}}
$$
to be isomorphic, i.e. that there exists
an natural isomorphism (i.e. invertible
natural transformation) between them.
That is:

\begin{definition}
\label{homological-algebra-definition-tensor-product-categorical}
A tensor product of $A,B \in \underline{\operatorname{Ab}}$
is an object $T \in \underline{\operatorname{Ab}}$
and a natural isomorphism
$\theta:\Hom_{\operatorname{Ab}}(T,-)
\to \operatorname{Bil}_{A,B}(-)$
of functors $\underline{\operatorname{Ab}}
\to \underline{\operatorname{Set}}$.
\end{definition}

\noindent
If you decode the fancy categorical
language, it's the usual universal
property of $A \otimes B$.

\begin{lemma}
\label{homological-algebra-lemma-tensor-product-uniqueness}
Two tensor products are naturally
isomorphic.
\end{lemma}

\noindent
Does a $\otimes$ exist in
$\underline{\operatorname{Ab}}$? Yes.
Let $A,B \in \underline{\operatorname{Ab}}$.
Consider $\mathbb{Z}[A \times B]$,
the free Abelian group on Cartesian
product. Let $J \subset \mathbb{Z}[A\times
B]$ be the submodule generated
by [distributivity of operations],
then $A \otimes B=\mathbb{Z}[A\times B]/J$
is a tensor product of $A$ and $B$.

\section{Complexes}
\label{homological-algebra-section-complexes}

\noindent
The upshot here is:
complexes on an additive category
$\mathcal{A}$ form a category called
the {\it category of chain complexes of
$\mathcal{A}$},
denoted $\text{Ch}(\mathcal{A})$.
If $\mathcal{A}$ is abelian, then
$\operatorname{Ch}(\mathcal{A})$
is abelian too
(Lemma \ref{homological-algebra-lemma-cat-chain-abelian},
by snake lemma).

For any $i \in \mathbf{Z}$ the $i$-th
homology group is a functor
$$
H_i : \text{Ch}(\mathcal{A}) \longrightarrow
\mathcal{A},
$$
which assigns to a chain complex $A_\bullet$
in an abelian category
its  $i$th {\it homology group}, defined by
the following formula
$$
H_i(A_\bullet) = \Ker(d_i)/\Im(d_{i + 1}).
$$
If $f : A_\bullet \to B_\bullet$ is a
morphism of chain
complexes of $\mathcal{A}$ then we get an
induced
morphism $H_i(f) : H_i(A_\bullet) \to
H_i(B_\bullet)$
because clearly
$f_i(\Ker(d_i : A_i \to A_{i - 1})) \subset
\Ker(d_i : B_i \to B_{i - 1})$, and similarly
for $\Im(d_{i + 1})$.
This proves that $H_i$ is indeed a functor.

Similarly,
for any $i \in \mathbf{Z}$ the $i$th {\it
cohomology group}
of a cochain complex $A^\bullet$ is defined
by
the following formula
$$
H^i(A^\bullet) = \Ker(d^i)/\Im(d^{i - 1}).
$$
If $f : A^\bullet \to B^\bullet$ is a
morphism of cochain
complexes of $\mathcal{A}$ then we get an
induced
morphism $H^i(f) : H^i(A^\bullet) \to
H^i(B^\bullet)$
because clearly
$f^i(\Ker(d^i : A^i \to A^{i + 1})) \subset
\Ker(d^i : B^i \to B^{i + 1})$, and similarly
for $\Im(d^{i - 1})$.
Thus we obtain a functor
$$
H^i : \text{CoCh}(\mathcal{A})
\longrightarrow \mathcal{A}.
$$




\medskip\noindent
Of course the notions of a chain complex and
a cochain complex
are dual and you only have to read one of the
two parts of
this section. So pick the one you like.
(Actually, this doesn't
quite work right since the conventions on
numbering things
are not adapted to an easy transition between
chain and cochain
complexes.)

\medskip\noindent
A {\it chain complex $A_\bullet$} in an
additive category $\mathcal{A}$
is a complex
$$
\ldots \to
A_{n + 1} \xrightarrow{d_{n + 1}}
A_n \xrightarrow{d_n}
A_{n - 1} \to
\ldots
$$
of $\mathcal{A}$. In other words, we are
given an object $A_i$ of
$\mathcal{A}$ for all $i \in \mathbf{Z}$ and
for
all $i \in \mathbf{Z}$ a morphism $d_i : A_i
\to A_{i - 1}$ such that
$d_{i - 1} \circ d_i = 0$ for all $i$. A {\it
morphism of chain
complexes $f : A_\bullet \to B_\bullet$} is
given by a
family of morphisms $f_i : A_i \to B_i$ such
that all
the diagrams
$$
\xymatrix{
A_i \ar[r]_{d_i} \ar[d]_{f_i} & A_{i - 1}
\ar[d]^{f_{i - 1}} \\
B_i \ar[r]^{d_i} & B_{i - 1}
}
$$
commute. The {\it category of chain complexes
of $\mathcal{A}$}
is denoted $\text{Ch}(\mathcal{A})$. The full
subcategory consisting
of objects of the form
$$
\ldots \to A_2 \to A_1 \to A_0 \to 0 \to 0
\to \ldots
$$
is denoted $\text{Ch}_{\geq 0}(\mathcal{A})$.
In other words, a chain complex $A_\bullet$
belongs to
$\text{Ch}_{\geq 0}(\mathcal{A})$ if and only
if
$A_i = 0$ for all $i < 0$.

\medskip\noindent
Given an additive category $\mathcal{A}$ we
identify $\mathcal{A}$
with the full subcategory of
$\text{Ch}(\mathcal{A})$ consisting
of chain complexes zero except in degree $0$
by the functor
$$
\mathcal{A} \longrightarrow
\text{Ch}(\mathcal{A}),\quad
A \longmapsto (\ldots \to 0 \to A \to 0 \to
\ldots)
$$
By abuse of notation we often denote the
object on the right hand side
simply $A$. If we want to stress that we are
viewing $A$ as a chain
complex we may sometimes use the notation
$A[0]$, see
Section \ref{homological-algebra-section-homotopy-shift}.

\medskip\noindent
A {\it homotopy $h$} between a pair of
morphisms
of chain complexes $f, g : A_\bullet \to
B_\bullet$
is a collection of morphisms $h_i : A_i \to
B_{i + 1}$
such that we have
$$
f_i - g_i = d_{i + 1} \circ h_i + h_{i - 1}
\circ d_i
$$
for all $i$.
Two morphisms $f, g : A_\bullet \to
B_\bullet$ are
said to be {\it homotopic} if a homotopy
between $f$
and $g$ exists.
Clearly, the notions of chain complex,
morphism of
chain complexes, and homotopies between
morphisms of chain complexes
make sense even in a preadditive category.

\begin{lemma}
\label{homological-algebra-lemma-compose-homotopy}
\begin{slogan}
Hom functors of $\text{Ch}(\mathcal{A})$
respect the homotopy relation.
\end{slogan}
Let $\mathcal{A}$ be an additive category.
Let $f, g : B_\bullet \to C_\bullet$ be
morphisms
of chain complexes. Suppose given morphisms
of chain
complexes $a : A_\bullet \to B_\bullet$, and
$c : C_\bullet \to D_\bullet$.
If $\{h_i : B_i \to C_{i + 1}\}$ defines a
homotopy
between $f$ and $g$, then $\{c_{i + 1} \circ
h_i \circ a_i\}$
defines a homotopy between $c \circ f \circ
a$ and
$c \circ g \circ a$.
\end{lemma}

\begin{proof}
Omitted.
\end{proof}

\noindent
In particular this means that it makes sense
to define
the category of chain complexes with maps up
to homotopy.
We'll return to this later.

\begin{definition}
\label{homological-algebra-definition-homotopy-equivalent}
Let $\mathcal{A}$ be an additive category.
We say a morphism $a : A_\bullet \to
B_\bullet$
is a {\it homotopy equivalence} if there
exists
a morphism $b : B_\bullet \to A_\bullet$
such that there exists a homotopy between
$a \circ b$ and $\text{id}_A$
and there exists a homotopy between $b \circ
a$ and $\text{id}_B$.
If there exists such a morphism between
$A_\bullet$ and $B_\bullet$, then
we say that $A_\bullet$ and $B_\bullet$ are
{\it homotopy equivalent}.
\end{definition}

\noindent
In other words, two complexes are homotopy
equivalent if they become
isomorphic in the category of complexes up to
homotopy.

\begin{lemma}
\label{homological-algebra-lemma-cat-chain-abelian}
Let $\mathcal{A}$ be an abelian category.
\begin{enumerate}
\item The category of chain complexes in
$\mathcal{A}$ is
abelian.
\item A morphism of complexes
$f : A_\bullet \to B_\bullet$ is injective
if and only if each $f_n : A_n \to B_n$ is
injective.
\item A morphism of complexes
$f : A_\bullet \to B_\bullet$ is surjective
if and only if each $f_n : A_n \to B_n$ is
surjective.
\item A sequence of chain complexes
$$
A_\bullet \xrightarrow{f} B_\bullet
\xrightarrow{g} C_\bullet
$$
is exact at $B_\bullet$ if and only if each
sequence
$$
A_i \xrightarrow{f_i} B_i \xrightarrow{g_i}
C_i
$$
is exact at $B_i$.
\end{enumerate}
\end{lemma}

\begin{proof}
Omitted.
\end{proof}

\noindent
For any $i \in \mathbf{Z}$ the $i$th {\it
homology group}
of a chain complex $A_\bullet$ in an abelian
category is defined by
the following formula
$$
H_i(A_\bullet) = \Ker(d_i)/\Im(d_{i + 1}).
$$
If $f : A_\bullet \to B_\bullet$ is a
morphism of chain
complexes of $\mathcal{A}$ then we get an
induced
morphism $H_i(f) : H_i(A_\bullet) \to
H_i(B_\bullet)$
because clearly
$f_i(\Ker(d_i : A_i \to A_{i - 1})) \subset
\Ker(d_i : B_i \to B_{i - 1})$, and similarly
for $\Im(d_{i + 1})$.
Thus we obtain a functor
$$
H_i : \text{Ch}(\mathcal{A}) \longrightarrow
\mathcal{A}.
$$

\begin{definition}
\label{homological-algebra-definition-quasi-isomorphism}
Let $\mathcal{A}$ be an abelian category.
\begin{enumerate}
\item A morphism of chain complexes $f :
A_\bullet \to B_\bullet$
is called a {\it quasi-isomorphism} if the
induced
map $H_i(f) : H_i(A_\bullet) \to
H_i(B_\bullet)$
is an isomorphism for all $i \in \mathbf{Z}$.
\item A chain complex $A_\bullet$ is called
{\it acyclic} if all of its homology objects
$H_i(A_\bullet)$ are zero.
\end{enumerate}
\end{definition}


\begin{lemma}
\label{homological-algebra-lemma-map-homology-homotopy}
Let $\mathcal{A}$ be an abelian category.
\begin{enumerate}
\item If the maps $f, g : A_\bullet \to
B_\bullet$ are
homotopic, then the induced maps $H_i(f)$ and
$H_i(g)$
are equal.
\item If the map $f : A_\bullet \to
B_\bullet$ is a homotopy
equivalence, then $f$ is a quasi-isomorphism.
\end{enumerate}
\end{lemma}

\begin{proof}
Omitted.
\end{proof}

\begin{lemma}
\label{homological-algebra-lemma-long-exact-sequence-chain}
Let $\mathcal{A}$ be an abelian category.
Suppose that
$$
0 \to
A_\bullet \to
B_\bullet \to
C_\bullet \to
0
$$
is a short exact sequence of chain complexes
of $\mathcal{A}$.
Then there is a canonical long exact homology
sequence
$$
\xymatrix{
\ldots & \ldots & \ldots \ar[lld] \\
H_i(A_\bullet) \ar[r] & H_i(B_\bullet) \ar[r]
& H_i(C_\bullet) \ar[lld] \\
H_{i - 1}(A_\bullet) \ar[r] &
H_{i - 1}(B_\bullet) \ar[r] &
H_{i - 1}(C_\bullet) \ar[lld] \\
\ldots & \ldots & \ldots \\
}
$$
\end{lemma}

\begin{proof}
Omitted. The maps come from the Snake Lemma
\ref{homological-algebra-lemma-snake}
applied to the diagrams
$$
\xymatrix{
&
A_i/\Im(d_{A, i + 1}) \ar[r] \ar[d]^{d_{A,
i}} &
B_i/\Im(d_{B, i + 1}) \ar[r] \ar[d]^{d_{B,
i}} &
C_i/\Im(d_{C, i + 1}) \ar[r] \ar[d]^{d_{C,
i}} &
0 \\
0 \ar[r] &
\Ker(d_{A, i - 1}) \ar[r] &
\Ker(d_{B, i - 1}) \ar[r] &
\Ker(d_{C, i - 1}) &
}
$$
\end{proof}

\begin{remark}
\label{homological-algebra-remark-snake-lemma-les}
The snake lemma gives us the following:
let 
$$
\xymatrix{
0\ar[r]
&M'\ar[r]
&M\ar[r]
&M''\ar[r]
&0\\
& (P')^\bullet\ar[u]
&& (P'')^\bullet \ar[u]
}
$$
be a short exact sequence of $R$-modules
with $(P')^\bullet$ and
$(P'')^\bullet$ projective
resolutions. Then there is a
projective resolution $P^\bullet \to M$
such that
$$
\xymatrix{
0\ar[r]
&(P')^\bullet\ar[r]
&P^\bullet\ar[r]
&(P'')^\bullet\ar[r]
&0
}
$$
is exact in the corresponding category.
\end{remark}

\noindent
A {\it cochain complex $A^\bullet$} in an
additive category $\mathcal{A}$
is a complex
$$
\ldots \to
A^{n - 1} \xrightarrow{d^{n - 1}}
A^n \xrightarrow{d^n}
A^{n + 1} \to
\ldots
$$
of $\mathcal{A}$. In other words, we are
given an object $A^i$ of
$\mathcal{A}$ for all $i \in \mathbf{Z}$ and
for
all $i \in \mathbf{Z}$ a morphism $d^i : A^i
\to A^{i + 1}$ such that
$d^{i + 1} \circ d^i = 0$ for all $i$. A {\it
morphism of cochain
complexes $f : A^\bullet \to B^\bullet$} is
given by a
family of morphisms $f^i : A^i \to B^i$ such
that all
the diagrams
$$
\xymatrix{
A^i \ar[r]_{d^i} \ar[d]_{f^i} & A^{i + 1}
\ar[d]^{f^{i + 1}} \\
B^i \ar[r]^{d^i} & B^{i + 1}
}
$$
commute. The {\it category of cochain
complexes of $\mathcal{A}$}
is denoted $\text{CoCh}(\mathcal{A})$. The
full subcategory consisting
of objects of the form
$$
\ldots \to 0 \to 0 \to A^0 \to A^1 \to A^2
\to \ldots
$$
is denoted $\text{CoCh}_{\geq
0}(\mathcal{A})$.
In other words, a cochain complex $A^\bullet$
belongs to the subcategory
$\text{CoCh}_{\geq 0}(\mathcal{A})$ if and
only if
$A^i = 0$ for all $i < 0$.

\medskip\noindent
Given an additive category $\mathcal{A}$ we
identify $\mathcal{A}$
with the full subcategory of
$\text{CoCh}(\mathcal{A})$ consisting
of cochain complexes zero except in degree
$0$ by the functor
$$
\mathcal{A} \longrightarrow
\text{CoCh}(\mathcal{A}),\quad
A \longmapsto (\ldots \to 0 \to A \to 0 \to
\ldots)
$$
By abuse of notation we often denote the
object on the right hand side
simply $A$. If we want to stress that we are
viewing $A$ as a cochain
complex we may sometimes use the notation
$A[0]$, see
Section \ref{homological-algebra-section-homotopy-shift}.

\medskip\noindent
A {\it homotopy $h$} between a pair of
morphisms
of cochain complexes $f, g : A^\bullet \to
B^\bullet$
is a collection of morphisms $h^i : A^i \to
B^{i - 1}$
such that we have
$$
f^i - g^i = d^{i - 1} \circ h^i + h^{i + 1}
\circ d^i
$$
for all $i$.
Two morphisms $f, g : A^\bullet \to
B^\bullet$ are
said to be {\it homotopic} if a homotopy
between $f$
and $g$ exists.
Clearly, the notions of cochain complex,
morphism of
cochain complexes, and homotopies between
morphisms of cochain complexes
make sense even in a preadditive category.

\begin{lemma}
\label{homological-algebra-lemma-compose-homotopy-cochain}
Let $\mathcal{A}$ be an additive category.
Let $f, g : B^\bullet \to C^\bullet$ be
morphisms
of cochain complexes. Suppose given morphisms
of cochain
complexes $a : A^\bullet \to B^\bullet$, and
$c : C^\bullet \to D^\bullet$.
If $\{h^i : B^i \to C^{i - 1}\}$ defines a
homotopy
between $f$ and $g$, then $\{c^{i - 1} \circ
h^i \circ a^i\}$
defines a homotopy between $c \circ f \circ
a$ and
$c \circ g \circ a$.
\end{lemma}

\begin{proof}
Omitted.
\end{proof}

\noindent
In particular this means that it makes sense
to define
the category of cochain complexes with maps
up to homotopy.
We'll return to this later.

\begin{definition}
\label{homological-algebra-definition-homot-equiv-cocha}
Let $\mathcal{A}$ be an additive category.
We say a morphism $a : A^\bullet \to
B^\bullet$
is a {\it homotopy equivalence} if there
exists
a morphism $b : B^\bullet \to A^\bullet$
such that there exists a homotopy between
$a \circ b$ and $\text{id}_A$
and there exists a homotopy between $b \circ
a$ and $\text{id}_B$.
If there exists such a morphism between
$A^\bullet$ and $B^\bullet$, then
we say that $A^\bullet$ and $B^\bullet$ are
{\it homotopy equivalent}.
\end{definition}

\noindent
In other words, two complexes are homotopy
equivalent if they become
isomorphic in the category of complexes up to
homotopy.

\begin{lemma}
\label{homological-algebra-lemma-cat-cochain-abelian}
Let $\mathcal{A}$ be an abelian category.
\begin{enumerate}
\item The category of cochain complexes in
$\mathcal{A}$ is
abelian.
\item A morphism of cochain complexes
$f : A^\bullet \to B^\bullet$ is injective
if and only if each $f^n : A^n \to B^n$ is
injective.
\item A morphism of cochain complexes
$f : A^\bullet \to B^\bullet$ is surjective
if and only if each $f^n : A^n \to B^n$ is
surjective.
\item A sequence of cochain complexes
$$
A^\bullet \xrightarrow{f} B^\bullet
\xrightarrow{g} C^\bullet
$$
is exact at $B^\bullet$ if and only if each
sequence
$$
A^i \xrightarrow{f^i} B^i \xrightarrow{g^i}
C^i
$$
is exact at $B^i$.
\end{enumerate}
\end{lemma}

\begin{proof}
Omitted.
\end{proof}

\noindent
For any $i \in \mathbf{Z}$ the $i$th {\it
cohomology group}
of a cochain complex $A^\bullet$ is defined
by
the following formula
$$
H^i(A^\bullet) = \Ker(d^i)/\Im(d^{i - 1}).
$$
If $f : A^\bullet \to B^\bullet$ is a
morphism of cochain
complexes of $\mathcal{A}$ then we get an
induced
morphism $H^i(f) : H^i(A^\bullet) \to
H^i(B^\bullet)$
because clearly
$f^i(\Ker(d^i : A^i \to A^{i + 1})) \subset
\Ker(d^i : B^i \to B^{i + 1})$, and similarly
for $\Im(d^{i - 1})$.
Thus we obtain a functor
$$
H^i : \text{CoCh}(\mathcal{A})
\longrightarrow \mathcal{A}.
$$

\begin{definition}
\label{homological-algebra-definition-quasi-isomorphism-cochain}
Let $\mathcal{A}$ be an abelian category.
\begin{enumerate}
\item A morphism of cochain complexes $f :
A^\bullet \to B^\bullet$
of $\mathcal{A}$ is called a {\it
quasi-isomorphism} if the induced
maps $H^i(f) : H^i(A^\bullet) \to
H^i(B^\bullet)$
is an isomorphism for all $i \in \mathbf{Z}$.
\item A cochain complex $A^\bullet$ is called
{\it acyclic} if all of its cohomology
objects
$H^i(A^\bullet)$ are zero.
\end{enumerate}
\end{definition}

\begin{lemma}
\label{homological-algebra-lemma-map-cohomology-homotopy-cochain}
Let $\mathcal{A}$ be an abelian category.
\begin{enumerate}
\item If the maps $f, g : A^\bullet \to
B^\bullet$ are
homotopic, then the induced maps $H^i(f)$ and
$H^i(g)$
are equal.
\item If $f : A^\bullet \to B^\bullet$ is a
homotopy equivalence,
then $f$ is a quasi-isomorphism.
\end{enumerate}
\end{lemma}

\begin{proof}
Omitted.
\end{proof}

\begin{lemma}
\label{homological-algebra-lemma-long-exact-sequence-cochain}
\begin{slogan}
Short exact sequences of complexes give rise
to long exact sequences
of (co)homology.
\end{slogan}
Let $\mathcal{A}$ be an abelian category.
Suppose that
$$
0 \to
A^\bullet \to
B^\bullet \to
C^\bullet \to
0
$$
is a short exact sequence of cochain
complexes of $\mathcal{A}$.
Then there is a long exact cohomology
sequence
$$
\xymatrix{
\ldots & \ldots & \ldots \ar[lld] \\
H^i(A^\bullet) \ar[r] &
H^i(B^\bullet) \ar[r] &
H^i(C^\bullet) \ar[lld] \\
H^{i + 1}(A^\bullet) \ar[r] &
H^{i + 1}(B^\bullet) \ar[r] &
H^{i + 1}(C^\bullet) \ar[lld] \\
\ldots & \ldots & \ldots \\
}
$$
The construction produces long exact
cohomology sequences
which are functorial in the short exact
sequence and compatible with shifts as in
Definition \ref{homological-algebra-definition-cohomology-shift}.
\end{lemma}

\begin{proof}
For the horizontal maps $H^i(A^\bullet) \to
H^i(B^\bullet)$ and
$H^i(B^\bullet) \to H^i(C^\bullet)$ we use
the fact that $H^i$
is a functor, see above. For the ``boundary
map''
$H^i(C^\bullet) \to H^{i + 1}(A^\bullet)$ we
use the map $\delta$
of the Snake Lemma \ref{homological-algebra-lemma-snake}
applied to the diagram
$$
\xymatrix{
&
A^i/\Im(d_A^{i - 1}) \ar[r] \ar[d]^{d_A^i} &
B^i/\Im(d_B^{i - 1}) \ar[r] \ar[d]^{d_B^i} &
C^i/\Im(d_C^{i - 1}) \ar[r] \ar[d]^{d_C^i} &
0 \\
0 \ar[r] &
\Ker(d_A^{i + 1}) \ar[r] &
\Ker(d_B^{i + 1}) \ar[r] &
\Ker(d_C^{i + 1}) &
}
$$
This works as the kernel of the right
vertical map is equal to
$H^i(C^\bullet)$ and the cokernel of the left
vertical map is
$H^{i + 1}(A^\bullet)$. The exactness of the
long sequence is
the exactnesss in part (2) of Lemma
\ref{homological-algebra-lemma-snake}.
The functoriality is Lemma
\ref{homological-algebra-lemma-snake-natural}. Compatibility
with shifts is immediate from the
definitions.
\end{proof}


\section{Injectives}
\label{homological-algebra-section-injectives}

\noindent
The upshot of injectives is
that whenever there is a map
defined on a subobject $A \subset B$,
then we can extend the map to the whole $B$.
Note that this defining property is not
universal: the extension may not be unique.

Also note that we discussed this concept,
and also the projective concept,
in Commutative Algebra.

\begin{definition}
\label{homological-algebra-definition-injective}
Let $\mathcal{A}$ be an abelian category.
An object $J \in \Ob(\mathcal{A})$ is
called {\it injective} if for every
injection $A \hookrightarrow B$ and every
morphism $A \to J$ there exists a morphism
$B \to J$ making the following diagram
commute:
$$
\xymatrix{
A \ar[r] \ar[d] &
B \ar@{-->}[ld] \\
J &
}
$$
\end{definition}

\noindent
Here is the obligatory characterization of
injective objects.

\begin{lemma}
\label{homological-algebra-lemma-characterize-injectives}
Let $\mathcal{A}$ be an abelian category,
and let $I$ be an object of $\mathcal{A}$.
The following are equivalent:
\begin{enumerate}
\item The object $I$ is injective.
\item The functor
$B \mapsto \Hom_\mathcal{A}(B,I)$ is exact.
\item Every short exact sequence
$$
0 \to I \to A \to B \to 0
$$
in $\mathcal{A}$ is split.
\item For every
$B \in \Ob(\mathcal{A})$, we have
$$
\Ext_\mathcal{A}(B,I)=0.
$$
\end{enumerate}
\end{lemma}

\begin{proof}
Omitted.
\end{proof}

\begin{lemma}
\label{homological-algebra-lemma-product-injectives}
Let $\mathcal{A}$ be an abelian category.
Suppose that $I_\omega$, for
$\omega \in \Omega$, is a set of injective
objects of $\mathcal{A}$. If
$$
\prod_{\omega \in \Omega} I_\omega
$$
exists, then it is injective.
\end{lemma}

\begin{proof}
Omitted.
\end{proof}

\begin{definition}
\label{homological-algebra-definition-enough-injectives}
Let $\mathcal{A}$ be an abelian category.
We say that $\mathcal{A}$ has {\it enough
injectives} if every object $A$ has an
injective morphism $A \to J$ into an
injective object $J$.
\end{definition}

\begin{definition}
\label{homological-algebra-definition-injective-resolution}
Let $\mathcal{A}$ be an abelian category.
Let $A \in \Ob(\mathcal{A})$.
An {\it injective resolution of $A$} is a
complex $I^\bullet$ together with a map
$A \to I^0$ such that:
\begin{enumerate}
\item We have $I^n = 0$ for $n < 0$.
\item Each $I^n$ is an injective object of
$\mathcal{A}$.
\item The map $A \to I^0$ is an isomorphism
onto $\Ker(d^0)$.
\item We have $H^i(I^\bullet) = 0$ for
$i > 0$.
\end{enumerate}
Hence $A[0] \to I^\bullet$ is a
quasi-isomorphism.
In other words the complex
$$
\ldots \to 0 \to A \to I^0 \to I^1
\to \ldots
$$
is acyclic.
\end{definition}

\begin{lemma}
\label{homological-algebra-lemma-injective-resolutions-exist}
Let $\mathcal{A}$ be an abelian category.
Assume $\mathcal{A}$ has enough injectives.
\begin{enumerate}
\item Any object of $\mathcal{A}$ has an
injective resolution.
\end{enumerate}
\end{lemma}

\begin{proof}
Proof of (1). First choose an injection
$A \to I^0$ of $A$ into an injective object
of $\mathcal{A}$. Next, choose an injection
$I_0/A \to I^1$ into an injective object of
$\mathcal{A}$. Denote $d^0$ the induced map
$I^0 \to I^1$.
Next, choose an injection
$I^1/\Im(d^0) \to I^2$ into an injective
object of $\mathcal{A}$. Denote $d^1$ the
induced map $I^1 \to I^2$. And so on.
\end{proof}

\begin{definition}
\label{homological-algebra-definition-functorial-injectives}
Let $\mathcal{A}$ be an abelian category.
We say that $\mathcal{A}$ has {\it
functorial injective embeddings} if there
exists a functor $J$ from $\mathcal{A}$ to
its arrow category such that:
\begin{enumerate}
\item the source of $J(A)$ is $A$;
\item for every object
$A \in \Ob(\mathcal{A})$, the morphism
$J(A)$ is injective;
\item for every object
$A \in \Ob(\mathcal{A})$, the target of
$J(A)$ is an injective object of
$\mathcal{A}$.
\end{enumerate}
We denote such a functor by
$A \mapsto (A \to J(A))$.
\end{definition}








\section{Derived functors}
\label{homological-algebra-section-derived-functos}

\noindent
Basic sort of question in algebra is,
say some object $C=B/A$ defined
as a quotient,
and we have some process $F$
(a functor), does $F(C)=F(B)/F(A)$,
and if not, how to correct this naive guess.

More abstractly, if
$$
\xymatrix{
0\ar[r]
&A\ar[r]^i
&B\ar[r]^p
&C\ar[r]
&0
}
$$
is a short exact sequence, then does
$$
\xymatrix{
0\ar[r]
&F(A)\ar[r]
&F(B)\ar[r]
&F(C)\ar[r]
&0
}
$$

\begin{remark}
\label{homological-algebra-remark-exact-functors}
In homological algebra,
exact functors are kind of rare,
but it is very common that a functor be
either left exact or right exact.
\end{remark}

\begin{definition}
\label{homological-algebra-definition-exact-functors}
Let $\mathcal{A}$ and $\mathcal{B}$ be
abelian categories, and let
$F : \mathcal{A} \to \mathcal{B}$ be a
functor. For every short exact sequence
$$
0 \to A \to B \to C \to 0
$$
in $\mathcal{A}$, consider the induced
sequence in $\mathcal{B}$.
\begin{enumerate}
\item The functor $F$ is {\it left exact}
if
$$
0 \to F(A) \to F(B) \to F(C)
$$
is exact.
\item The functor $F$ is {\it right exact}
if
$$
F(A) \to F(B) \to F(C) \to 0
$$
is exact.
\item The functor $F$ is {\it exact} if
$$
0 \to F(A) \to F(B) \to F(C) \to 0
$$
is exact.
\end{enumerate}
\end{definition}

\begin{definition}
\label{homological-algebra-definition-right-derived-functors}
Let $\mathcal{A}$ have enough injectives.
Let $F: \mathcal{A} to \mathcal{B}$ be a
left exact functor.
Define the {\it right derived functors}
of $F$ to be
$$
R^nF(A) = H^n(F(I^\bullet)),
$$
when $0 \to I^0 \to I^1 \to I^2 \to…$
is an injective resolution of $\mathcal{A}$.
\end{definition}

\begin{theorem}
\label{homological-algebra-theorem-derived-functors}
\begin{enumerate}
\item The value of $R^nF(A) \in \mathcal{B}$
does not depend on the choice
of $A \to I^\bullet$ (up to isomorphism).

\item A morphism $f:A \to A'$
in $\mathcal{A}$ induces
$R^nF(A) \to R^nF(A')$
well-defined, and $R^nF$,
$n \geq  0$ are all functors
$\mathcal{A} \to \mathcal{B}$.

\item $R^0F(A) = F(A)$.

\item If
$0 \to A \to B \to C \to 0$
is a short exact sequence, then we have

$$
\xymatrix{
&R^2F(A) \ar[r] & \cdots\\
&R^1F(A) \ar[r] & R^1F(B)\ar[r]&R^1F(C)
\ar[ull]\\
0\ar[r] &
F(A)\ar[r]&F(B)\ar[r]&F(C)\ar[ull]^\delta
}
$$
is exact.
\end{enumerate}
\end{theorem}


\medskip\noindent
{\bf Divisible groups.}

\begin{definition}
\label{homological-algebra-definition-divisible-group}
An abelian group $I$ is called
{\it divisible} if $\forall i \in I$
and $n \in \mathbb{Z}\setminus\{0\}$
there exists $i' \in I$ such that
$ni'=i$.
\end{definition}

\noindent
Example: $I = \mathbb{Q}$.

\section{Projectives}
\label{homological-algebra-section-projectives}

\noindent
Projective objects have a lifting property
dual to the extension property of injective
objects. A morphism from a projective
object to a quotient lifts to the object
being quotiented. The lift need not be
unique.

\begin{definition}
\label{homological-algebra-definition-projective}
Let $\mathcal{A}$ be an abelian category.
An object $P$ of $\mathcal{A}$ is called
{\it projective} if, for every epimorphism
$f:A\to B$ and every morphism
$\alpha:P\to B$, there is a morphism
$\beta:P\to A$ such that
$$
\alpha=f\circ\beta.
$$
Thus there is a commutative diagram
$$
\xymatrix{
A \ar[r]^f & B \ar[r] & 0 \\
P \ar[u]^{\beta} \ar[ur]_{\alpha} &
}
$$
for every such $f$ and $\alpha$.
\end{definition}

\begin{definition}
\label{homological-algebra-definition-projection-module}
Let $R$ be a ring and let $F$ be an
$R$-module. A {\it projection on $F$} is
an endomorphism $\pi:F\to F$ satisfying
$$
\pi^2=\pi.
$$
\end{definition}

\begin{lemma}
\label{homological-algebra-lemma-projection-decomposition}
Let $\pi:F\to F$ be a projection on an
$R$-module. Then
$$
F\cong\Ker(\pi)\oplus\Im(\pi).
$$
\end{lemma}

\begin{proof}
Set $K=\Ker(\pi)$ and $I=\Im(\pi)$.
If $x\in K\cap I$, then $x=\pi(y)$ for
some $y\in F$, and hence
$$
x=\pi(y)=\pi^2(y)=\pi(x)=0.
$$
On the other hand, every $x\in F$ has the
decomposition
$$
x=\pi(x)+(x-\pi(x)),
$$
where the first summand belongs to $I$ and
the second belongs to $K$.
\end{proof}

\begin{definition}
\label{homological-algebra-definition-projective-direct-summand}
Let $R$ be a ring. A {\it projective
$R$-module} is an $R$-module $P$ for which
there are an $R$-module $C$, a free
$R$-module $F$, and an isomorphism
$$
F\cong P\oplus C.
$$
\end{definition}

\begin{lemma}[Corollary]
\label{homological-algebra-lemma-projection-projective}
The kernel and image of a projection on a
free $R$-module are projective
$R$-modules.
\end{lemma}

\begin{proof}
This follows from Lemma
\ref{homological-algebra-lemma-projection-decomposition}.
\end{proof}

\begin{example}
\label{homological-algebra-example-moebius-projective}
Consider the ring
$$
R=
\mathbb{R}[x,y]/(x^2+y^2-1)
$$
and the endomorphism of $R^2$ given by
$$
\pi
=
\frac{1}{2}
\begin{pmatrix}
1+x & y \\
y & 1-x
\end{pmatrix}.
$$
The relation in $R$ gives $\pi^2=\pi$.
Consequently, $\Im(\pi)$ is a projective
$R$-module.

Geometrically, $R$ is the ring of
polynomial functions on the real unit
circle. At the point
$$
p(\theta)
=
(\cos(\theta),\sin(\theta)),
$$
the image of $\pi$ is the line spanned by
$$
(\cos(\theta/2),\sin(\theta/2)).
$$
As the point goes once around the circle,
this line makes a half turn. Thus the
module describes the sections of the
Möbius line bundle.
\end{example}

\begin{lemma}
\label{homological-algebra-lemma-projective-module-object}
An object of $\mathit{Mod}_R$ is projective
in the sense of Definition
\ref{homological-algebra-definition-projective} if and only if
it is a projective $R$-module in the sense
of Definition
\ref{homological-algebra-definition-projective-direct-summand}.
\end{lemma}

\begin{proof}
Suppose first that $P$ is projective as an
object of $\mathit{Mod}_R$. Choose a free
$R$-module $F$ with an epimorphism
$f:F\to P$. The identity of $P$ lifts to
a morphism $s:P\to F$ satisfying
$$
f\circ s=\operatorname{id}_P.
$$
In particular, $s$ is injective. If
$K=\Ker(f)$, every $x\in F$ can be written
uniquely as
$$
x=(x-s(f(x)))+s(f(x)),
$$
with the first summand in $K$ and the
second in $s(P)$. Moreover,
$K\cap s(P)=0$. Thus
$$
F\cong K\oplus P.
$$

Conversely, suppose that
$F\cong P\oplus K$ for a free module $F$.
Let $f:A\to B$ be an epimorphism and let
$\alpha:P\to B$ be a morphism. Extend
$\alpha$ to a morphism $F\to B$ by making
it zero on $K$. A map from a free module
lifts across every epimorphism: choose a
basis of $F$ and a preimage in $A$ of the
image of each basis element. Restricting
such a lift $F\to A$ to $P$ gives the
required lift of $\alpha$.
\end{proof}

\begin{definition}
\label{homological-algebra-definition-enough-projectives}
An abelian category $\mathcal{A}$ has
{\it enough projectives} if every object
$A$ admits an epimorphism $P\to A$ from a
projective object $P$.
\end{definition}

\begin{lemma}
\label{homological-algebra-lemma-modules-enough-projectives}
For every ring $R$, the category
$\mathit{Mod}_R$ has enough projectives.
\end{lemma}

\begin{proof}
Every $R$-module is the quotient of a free
$R$-module. Every free $R$-module is
projective by Lemma
\ref{homological-algebra-lemma-projective-module-object}.
\end{proof}

\begin{definition}
\label{homological-algebra-definition-projective-resolution}
Let $\mathcal{A}$ be an abelian category
and let $A$ be an object of $\mathcal{A}$.
A {\it projective resolution of $A$} is a
chain complex $P_\bullet$ together with an
augmentation $P_0\to A$ such that:
\begin{enumerate}
\item We have $P_n=0$ for $n<0$.
\item Every $P_n$ is projective.
\item The augmented complex
$$
\ldots\to P_2\to P_1\to P_0
\to A\to 0
$$
is exact.
\end{enumerate}
Equivalently, $P_\bullet\to A[0]$ is a
quasi-isomorphism.
\end{definition}

\begin{lemma}
\label{homological-algebra-lemma-projective-resolutions-exist}
If an abelian category $\mathcal{A}$ has
enough projectives, then every object of
$\mathcal{A}$ has a projective resolution.
In particular, every $R$-module has a
projective resolution.
\end{lemma}

\begin{proof}
Choose an epimorphism $P_0\to A$ with
$P_0$ projective. Next choose an
epimorphism
$$
P_1\to\Ker(P_0\to A)
$$
with $P_1$ projective, and compose it with
the inclusion into $P_0$. Continuing with
the kernel at each stage produces an exact
augmented complex
$$
\ldots\to P_2\to P_1\to P_0
\to A\to 0.
$$
The last assertion follows from Lemma
\ref{homological-algebra-lemma-modules-enough-projectives}.
\end{proof}

\noindent
Projective resolutions are the dual input
to injective resolutions. If a right exact
functor admits left derived functors, they
are computed by applying the functor to a
projective resolution and taking homology.

\section{Extensions}
\label{homological-algebra-section-extensions}

\noindent
For our purposes, the point of extensions is to understand concretely why the first derived functor $\operatorname{Ext}^1$ classifies short exact sequences.

\noindent
Let $R$ be a ring and $M$ an $R$-module.
The functor
$
\operatorname{Hom}_R(M,-):
R\operatorname{-mod}
\to \underline{\operatorname{Ab}}
$
is left exact, and

$[R^n\Hom_R(M,-)](A)$
is called $\operatorname{Ext}^n_R(M,A)$.

\begin{remark}
\label{homological-algebra-remark-hom-functor-contravariant}
Fixing an $R$-module $N$,
the functor
$$
\Hom_R(-,N):R\operatorname{-mod}
\to \underline{\operatorname{Ab}}
$$
is a left exact (and contravariant)
functor.
\end{remark}

\noindent
So if $R$-mod has enough projectives
(which it does) we can define left
derived functors
$$
[L^n\Hom_R(-,N)](M).
$$

\begin{theorem}[Balancing]
\label{homological-algebra-theorem-balancing}
$$
[R^n\Hom_R(M,-)](N)
\simeq
[L^n\Hom_R(-,N)](M)
$$
and both equal
$$
\operatorname{Ext}^n_R(M,N).
$$
\end{theorem}

\noindent
We shall prove this theorem later.

\begin{definition}
\label{homological-algebra-definition-extension}
Let $A$ and $B$ be $R$-modules. An \emph{extension of $B$ by $A$} is a short exact sequence
$$
0\longrightarrow A\longrightarrow E\longrightarrow B\longrightarrow0.
$$
Two extensions are \emph{equivalent} if there is an isomorphism between their middle terms making the diagram
$$
\xymatrix{
0\ar[r]&A\ar[r]\ar[d]^{\mathrm{id}}&E\ar[r]\ar[d]^{\sim}&B\ar[r]\ar[d]^{\mathrm{id}}&0\\
0\ar[r]&A\ar[r]&E'\ar[r]&B\ar[r]&0
}
$$
commute.
\end{definition}

\begin{definition}
\label{homological-algebra-definition-split-extension}
An extension is \emph{split} if the surjection $E\to B$ admits a section $B\to E$. Equivalently,
$$
E\cong A\oplus B
$$
compatibly with the extension maps.
\end{definition}

\medskip\noindent
\textbf{From an extension to a derived $\operatorname{Ext}^1$ class.}
Choose a projective resolution of $B$ and look at its beginning:
$$
\cdots\longrightarrow P_2\xrightarrow{d_2}P_1\xrightarrow{d_1}P_0\xrightarrow{p}B\longrightarrow0.
$$
In particular,
$$
P_1\xrightarrow{d_1}P_0\xrightarrow{p}B\longrightarrow0
$$
is a projective presentation of $B$. Such a presentation always exists for $R$-modules: Lemma
\ref{homological-algebra-lemma-modules-enough-projectives} proves that $R$-modules have enough projectives, and Lemma
\ref{homological-algebra-lemma-projective-resolutions-exist} then proves that every $R$-module admits a projective resolution.

Now let
$$
0\longrightarrow A\longrightarrow E\xrightarrow{q}B\longrightarrow0
$$
be an extension. Since $P_0$ is projective and $q$ is surjective, the map $p:P_0\to B$ lifts to a map
$$
\widetilde p:P_0\longrightarrow E
$$
such that
$$
q\widetilde p=p.
$$
This is exactly the lifting property in Definition
\ref{commutative-algebra-definition-projective-module}
in \emph{Commutative Algebra}.

Consider the composite
$$
P_1\xrightarrow{d_1}P_0\xrightarrow{\widetilde p}E.
$$
Its image lies in $A$. Indeed, after composing with $q:E\to B$ we obtain
$$
q\widetilde p d_1=pd_1=0,
$$
because the projective resolution is a complex. Since
$$
A=\ker(q),
$$
the composite factors uniquely through a map
$$
\phi:P_1\longrightarrow A.
$$

If we choose a different lift $\widetilde p$, the resulting $\phi$ changes by an element coming from
$$
\operatorname{Hom}_R(P_0,A)\longrightarrow\operatorname{Hom}_R(P_1,A).
$$
Moreover, since $d_1d_2=0$, the map $\phi$ lies in the kernel of
$$
\operatorname{Hom}_R(P_1,A)\longrightarrow\operatorname{Hom}_R(P_2,A).
$$
Thus the extension determines a well-defined class
$$
[\phi]\in
\frac{\ker\bigl(\operatorname{Hom}_R(P_1,A)\to\operatorname{Hom}_R(P_2,A)\bigr)}
{\operatorname{im}\bigl(\operatorname{Hom}_R(P_0,A)\to\operatorname{Hom}_R(P_1,A)\bigr)}.
$$
But this quotient is exactly
$$
H^1\bigl(\operatorname{Hom}_R(P_\bullet,A)\bigr),
$$
which by definition of the derived functor is
$$
\operatorname{Ext}^1_R(B,A).
$$

\begin{theorem}
\label{homological-algebra-theorem-ext1-classifies-extensions}
The construction above gives a natural bijection
$$
\left\{
\begin{array}{c}
\text{equivalence classes of extensions}\\
0\to A\to E\to B\to0
\end{array}
\right\}
\cong
\operatorname{Ext}^1_R(B,A).
$$
Under this bijection the split extension corresponds to $0$.
\end{theorem}

\noindent
The extension classes themselves form an abelian group. The zero element is the split extension
$$
0\longrightarrow A\longrightarrow A\oplus B\longrightarrow B\longrightarrow0.
$$
The sum of two extension classes is the \emph{Baer sum}: one combines the two short exact sequences by the direct sum, pulls back along the diagonal map $B\to B\oplus B$, and pushes out along addition $A\oplus A\to A$. Under the theorem this group law agrees with the usual addition on
$$
\operatorname{Ext}^1_R(B,A).
$$
Thus saying $\operatorname{Ext}^1_R(B,A)=0$ means precisely that every extension of $B$ by $A$ is equivalent to the split extension.

\section{Group cohomology}
\label{homological-algebra-section-group-cohomology}

\noindent
I did not attend this lecture and have not yet read the contents of this section.

\noindent
We have considered the category
$G$-mod $\sim$ $\mathbb{Z}G$-mod and the
functor of $G$-invariants
$$
(-)^G : G\operatorname{-mod}
\to
\underline{\operatorname{Ab}}.
$$
this functor is left exact, and
$\mathbb{Z}G$-mod has enough injectives,
so we can construct right derived functors
$$
R^n(-)^G: G\operatorname{-mod}
\to \underline{\operatorname{Ab}}.
$$
\medskip\noindent
Applying Ext to group cohomology,
$$
[R^n(-)^G](M)
=
\operatorname{Ext}^n_{\mathbb{Z} G}
(\mathbb{Z},M).
$$
We can compute it by forming an injective
resolution of $M$… OR (!)
a projective resolution of $\mathbb{Z}$
(which is better).



\medskip\noindent
{\bf The bar complex.}
The projective resolution of $\mathbb{Z}$
will be
$$
… \xrightarrow{d}P_2
\xrightarrow{d}P_1
\xrightarrow{d}P_0
\xrightarrow{\varepsilon}
\mathbb{Z}\to 0
$$
with $P_n=\mathbb{Z}G^{n+1}$
as abelian group (free $\mathbb{Z}$-module).

$G$-action on $\mathbb{Z}G^{n+1}$ is
$$
g(g_0,g_1,…,g_n)
=
(gg_0, g g_1,…,g g_n),
$$
extended $\mathbb{Z}$-linearly.

The differential
$$
\mathbb{Z}G^{n+1}
\xrightarrow{d}
\mathbb{Z}G^n,
$$
$$
d((g_0,g_1,…,g_n))
=
\sum_{i=0}^n
(-1)^i(g_0,g_1,…,\widehat{g_i},…,g_n).
$$
So for example
$$
\begin{aligned}
\mathbb{Z}G^2 &\longrightarrow \mathbb{Z}G \\
(g_0,g_1) &\longmapsto g_1-g_0,
\end{aligned}
$$
and
$$
\begin{aligned}
\mathbb{Z}G^3 &\longrightarrow \mathbb{Z}G \\
(g_0,g_1,g_2) &\longmapsto 
(g_1,g_2)-(g_0,g_2)+(g_0,g_1).
\end{aligned}
$$
The ``accumulation'' is
$$
\varepsilon:\mathbb{Z}G \to \mathbb{Z},
\qquad 
\varepsilon(g)=1 \quad \forall g \in G.
$$
\begin{exercise}
\label{homological-algebra-exercise-bar-complex}
$d \circ d=0$.
\end{exercise}

\begin{remark}
\label{homological-algebra-remark-bar}
Every $P_n$ is a free $\mathbb{Z}G$-module,
and free modules are projective.
Why is $\tilde{P}_\bullet$
acyclic? We'll see soon.
\end{remark}

\noindent
This leads to the following
definition/theorem.

\begin{definition}
\label{homological-algebra-definition-group-cohomology}
Let $A$ be a $G$-module.
We define
$$
C^n(G,A)
=
\{\varphi:G^n\to A\}
$$
(no restrictions whatsoever, the $\varphi$
are just functions of sets).
$C^n$ is an abelian group (by the addition
of $A$). Let $C^n = C^n(G,A)$.
Differentials
$$
0\to C^0
\xrightarrow{d} C^1
\xrightarrow{d} C^2
\xrightarrow{d}…
$$
If $\varphi \in C^n$,
$$
[d\varphi](g_0,g_1,…,g_n)
=
g_0\varphi(g_1,g_2,…,g_n)
-
\sum_{i=a}^{n-1}
(-1)^i
\varphi(g_0|g_1|…g_{i-1}|g_ig_{i+1}|
g_{i+2}|…g_n)
+(-1)^{n+1}\varphi(g_0|g_1|…|g_{n-1}).
$$
Let $Z^n=\ker d$,
$B^n=\operatorname{Im}(d)$, since
$d^2=0$,  $B^n \subset Z^n$.
The group cohomology is
$$
H^n(G,A) = Z^n/B^n.
$$
\end{definition}

\begin{exercise}
\label{homological-algebra-exercise-hom-group-cohom}
$$
\Hom_{\mathbb{Z}G}(P_n,A)
=
C^n(G,A).
$$
\end{exercise}

\begin{theorem}
\label{homological-algebra-theorem-ext-group-cohom}
$$
\operatorname{Ext}^n_{\mathbb{Z}A}(\mathbb{Z},A)
=
H^n(G,A).
$$
\end{theorem}

\begin{proof}
\begin{enumerate}
\item Resolve $\mathbb{Z}$ by 
$P_{\bullet}$ (balancing says this is
legit).

\item Exercise says
$$
\operatorname{Ext}^n_{\mathbb{Z}G}(\mathbb{Z},A)
=
[L^n\Hom_{\mathbb{Z}G}(Z,-)](A)
=
\Hom_{\mathbb{Z}G}(P_n,A)
=
C^n(G,A).
$$
\item Remains to prove
$$
\tilde{P}_\bullet=P_\bullet
\to \mathbb{Z}
\xrightarrow{\varepsilon}0
$$
is acyclic. [Proof continued.]
\end{enumerate}
\end{proof}

\noindent
Unwrapping definitions we see that
$$
H^0(G,A) = A^G,
$$
$$
H^1(G,A)
=
\frac{\{\varphi:G \to A:
g\varphi(h)-\varphi(gh)
+\varphi(g)=0\}}{
\varphi \operatorname{da}
\operatorname{forma}
\varphi(g)=ga-a}.
$$
The bar resolution is very big, its
advantage is it works for all $G$.
If I'm intersete in a specific $G$,
maybe there exists a smaller resolution…

\begin{example}
\label{homological-algebra-example-group-cohom-cyclic}
Let $G = C_m \simeq \mathbb{Z}/m$
be the cylcic group…
\end{example}


\noindent
Anything else in this lecture?


\section{Background for group cohomology}
\label{homological-algebra-section-ai-background-sept-four}

\noindent
I did not attend this lecture and have not yet read the contents of this section.

\medskip\noindent
The September 4 reconstruction uses
Francisca Aguayo Moreno's notes,
\emph{Homological Algebra}, pp.~13--17,
and Jethro van Ekeren's
\href{https://w3.impa.br/~jethro/2026-2/%
ah-current.pdf}{course notes}.
The following background is retained
from the AI revision to support its
proofs and references. It is not a
replacement for the earlier lecture
notes.

\begin{proposition}
\label{homological-algebra-proposition-divisible-injective}
An abelian group is injective if and only
if it is divisible.
\end{proposition}

\begin{proof}
If $I$ is injective, extend the map
$n\mathbb{Z}\to I$ sending $n$ to $a$
to a map $f:\mathbb{Z}\to I$. Then
$nf(1)=a$, proving divisibility.

Conversely, suppose $I$ is divisible and
$f:A\to I$ is defined on a subgroup
$A\subset B$. Order its extensions to
subgroups of $B$ by inclusion. Unions of
chains give extensions, so Zorn's lemma
gives a maximal extension $f':A'\to I$.
If $x\in B\setminus A'$, consider the
subgroup of integers $k$ with $kx\in A'$.
If it is zero, extend to $A'+\mathbb{Z}x$
by sending $x$ to zero. Otherwise it is
$n\mathbb{Z}$ for some $n>0$.
Choose $y\in I$ with $ny=f'(nx)$ and set
$$
f''(a+kx)=f'(a)+ky.
$$
The choice of $y$ makes this well-defined.
Both cases contradict maximality. Thus
$A'=B$, proving injectivity.
\end{proof}

\begin{proposition}[Averaging]
\label{homological-algebra-proposition-invariants-averaging}
Let $G$ be finite and $F$ a field of
characteristic zero. The functor
$$
(-)^G:\operatorname{Rep}_F(G)
\to\operatorname{Vec}_F
$$
is exact. Here representations may have
arbitrary dimension.
\end{proposition}

\begin{proof}
Invariants are left exact. Given an
equivariant surjection $p:B\to C$ and
$c\in C^G$, choose $b$ with $p(b)=c$.
Then
$$
\overline b=\frac{1}{|G|}\sum_{g\in G}gb
$$
is invariant and $p(\overline b)=c$.
Thus $B^G\to C^G$ is surjective as well.
\end{proof}

\begin{example}[Cyclic groups]
\label{homological-algebra-example-ai-cyclic-sept-four}
Let $G=C_m=\langle\sigma\rangle$ and put
$$
N=1+\sigma+\cdots+\sigma^{m-1}
\in\mathbb{Z}G.
$$
A smaller free resolution is
$$
\cdots\xrightarrow{\sigma-1}\mathbb{Z}G
\xrightarrow{N}\mathbb{Z}G
\xrightarrow{\sigma-1}\mathbb{Z}G
\xrightarrow{\varepsilon}\mathbb{Z}\to0.
$$
Indeed, the kernel of $\sigma-1$ consists
of elements whose coefficients are all
equal, hence is $\mathbb{Z}N$.
Multiplication by $N$ sends $x$ to
$\varepsilon(x)N$, and the augmentation
kernel is $(\sigma-1)\mathbb{Z}G$.
These identities prove exactness.

Applying Hom gives
$$
0\to A\xrightarrow{\sigma-1}A
\xrightarrow{N}A
\xrightarrow{\sigma-1}A
\xrightarrow{N}\cdots.
$$
Thus, for every integer $r\geq0$,
$$
\begin{aligned}
H^0(G,A)&=A^G,\\
H^{2r+1}(G,A)&=\ker(N)/(\sigma-1)A,\\
H^{2r+2}(G,A)&=A^G/NA.
\end{aligned}
$$
For trivial action, $\sigma-1=0$ and
$N$ is multiplication by $m$. If
$A[m]=\{a\in A\mid ma=0\}$, this gives
$$
H^{2r+1}(G,A)=A[m],
\qquad H^{2r+2}(G,A)=A/mA.
$$
In particular, for trivial coefficients,
$$
\begin{aligned}
H^{2r+1}(G,\mathbb{Z})&=0,\\
H^{2r+2}(G,\mathbb{Z})&=\mathbb{Z}/m,\\
H^{2r+1}(G,\mathbb{C}^{\times})&=\mu_m,\\
H^{2r+2}(G,\mathbb{C}^{\times})&=0.
\end{aligned}
$$
Here $\mu_m$ is the group of $m$th roots
of unity. For the additive group
$\mathbb{C}$, all positive-degree
cohomology vanishes, since multiplication
by $m$ is bijective.
\end{example}

\section{Group cohomology}
\label{homological-algebra-section-group-cohomology-4-sept}

\noindent
Francisca's record of this lecture starts
on p.~13. We first finish the argument
begun on September 2. Unless otherwise
specified, $G$ is finite and $F$ has
characteristic zero.

\medskip\noindent
Exactness of $(-)^G$ on representations
implies that its derived functors in
$\operatorname{Rep}_F(G)$ vanish in
positive degrees. But $H^n(G,V)$ was
defined using injectives in $G$-modules.
We must answer two questions: does
$\operatorname{Rep}_F(G)$ have enough
injectives, and can we choose a resolution
whose terms are injective in $G$-modules
as well?

\begin{definition}
\label{homological-algebra-definition-coinduction-4-sept}
For an $F$-vector space $W$, its
\emph{coinduced representation} is
$$
\operatorname{Coind}(W)
=\operatorname{Hom}_F(FG,W),
$$
with $(g\varphi)(h)=\varphi(hg)$.
The formula on $h\in G$ extends linearly
to $FG$. Equivalently, this is the space
of functions $G\to W$ with right
translation on the argument.
\end{definition}

\noindent
The action law follows from
$$
(g_1(g_2\varphi))(h)=\varphi(hg_1g_2).
$$
The alternative convention
$(g\varphi)(h)=\varphi(g^{-1}h)$ is
isomorphic to this one via
$\varphi(h)\mapsto\varphi(h^{-1})$.

\begin{proposition}
\label{homological-algebra-proposition-forgetful-4-sept}
Coinduction is right adjoint to the
forgetful functor
$$
U:\operatorname{Rep}_F(G)\to
\operatorname{Vec}_F.
$$
\end{proposition}

\begin{proof}
There is a natural bijection
$$
\begin{aligned}
&\operatorname{Hom}_{FG}
(M,\operatorname{Coind}(W))\\
&\qquad\cong\operatorname{Hom}_F(U(M),W).
\end{aligned}
$$
It sends $f$ to $m\mapsto f(m)(e)$.
Conversely, a linear map $a:U(M)\to W$
gives $f_a(m)(h)=a(hm)$.
This is equivariant because
$$
f_a(gm)(h)=a(hgm)=(g f_a(m))(h).
$$
The two constructions are inverse:
equivariance recovers $f(m)(h)$ from
$f(hm)(e)$. They commute with composition
in both variables, proving naturality.
\end{proof}

\begin{proposition}
\label{homological-algebra-proposition-adjoint-exact-4-sept}
An additive right adjoint between abelian
categories is left exact.
\end{proposition}

\begin{proof}
Right adjoints preserve kernels and finite
products by their universal properties.
Thus they preserve left exact sequences.
\end{proof}

\begin{proposition}
\label{homological-algebra-proposition-monos-4-sept}
A left exact additive functor preserves
monomorphisms.
\end{proposition}

\begin{proof}
Apply it to $0\to A\to B$.
\end{proof}

\begin{proposition}
\label{homological-algebra-proposition-adjoint-injective-4-sept}
The right adjoint to an exact additive
functor sends injective objects to
injective objects.
\end{proposition}

\begin{proof}
Write $L:\mathcal A\to\mathcal B$ for
the exact functor and $R$ for its right
adjoint. If $J$ is injective, then
$$
\operatorname{Hom}_{\mathcal A}(-,R(J))
\cong\operatorname{Hom}_{\mathcal B}(L(-),J)
$$
is exact. Hence $R(J)$ is injective.
\end{proof}

\noindent
These three proofs were left as exercises
on the board; the short justifications
above fill them in.

\begin{lemma}[Corollary]
\label{homological-algebra-lemma-adjoint-injectives-4-sept}
Let $L:\mathcal A\to\mathcal B$ be
faithful and exact, with right adjoint
$R$. If $\mathcal B$ has enough
injectives, then so does $\mathcal A$.
\end{lemma}

\begin{proof}
Choose a monomorphism $j:L(X)\to J$
with $J$ injective. Its adjoint is
$i:X\to R(J)$, whose target is injective
by the preceding proposition.
If $a,b:Y\to X$ satisfy $ia=ib$,
naturality of the adjunction gives
$jL(a)=jL(b)$. Since $j$ is monic and
$L$ is faithful, $a=b$. Thus $i$ is monic.

In the notation of the lecture, the unit
$\eta_X:X\to RL(X)$ corresponds to
$\operatorname{id}_{L(X)}$, and
$i=R(j)\eta_X$. The same argument, with
$j$ the identity, proves that $\eta_X$
is monic.
\end{proof}

\begin{remark}
\label{homological-algebra-remark-faithfulness-4-sept}
The faithfulness hypothesis is missing
from the recorded board statement and
the previous transcription. It is needed
in the proof and is included in
Corollary 3.4.9 of Jethro's current notes.
See also the Stacks Project,
\href{https://stacks.math.columbia.edu/tag/%
015Y}{Section 12.29, Lemma 12.29.3}.
Forgetting a group action is faithful,
so the application here satisfies it.
\end{remark}

\noindent
Every vector space is injective in
$\operatorname{Vec}_F$, since a linear
map on a subspace extends after choosing
a complementary subspace. Hence
$\operatorname{Coind}(W)$ is injective
in $\operatorname{Rep}_F(G)$.
The adjunction gives an embedding
$$
\begin{aligned}
\eta_V:V&\to\operatorname{Coind}(U(V)),\\
\eta_V(v)(h)&=hv.
\end{aligned}
$$
Evaluation at $e$ recovers $v$, proving
injectivity directly. This answers the
first question.

\begin{proposition}
\label{homological-algebra-proposition-coind-injective-4-sept}
If $F$ has characteristic zero, then
$\operatorname{Coind}(W)$ is also
injective as a $G$-module.
\end{proposition}

\begin{proof}
For every $G$-module $M$, evaluation at
$e$ gives a natural isomorphism
$$
\operatorname{Hom}_{\mathbb{Z}G}
(M,\operatorname{Coind}(W))
\cong\operatorname{Hom}_{\mathbb{Z}}(M,W).
$$
Its inverse sends $a:M\to W$ to the map
whose value at $m$ is the $F$-linear
extension of $h\mapsto a(hm)$ from the
basis $G$ of $FG$. The verification is
the same as for the first adjunction.
As an abelian group, $W$ is divisible:
for $n\ne0$, division by $n$ is possible
because $F$ has characteristic zero.
Thus $W$ is injective in abelian groups
by Proposition
\ref{homological-algebra-proposition-divisible-injective}.
The functor on the right is exact,
which proves the claim.
\end{proof}

\begin{theorem}
\label{homological-algebra-theorem-group-vanishing-4-sept}
Let $G$ be finite and $V$ an $F$-linear
representation, with $F$ of characteristic
zero. Then $H^n(G,V)=0$ for every $n>0$.
\end{theorem}

\begin{proof}
Embed $V$ into $\operatorname{Coind}(U(V))$,
take the quotient, and repeat. This gives
a resolution by representations which are
injective in both categories. Applying
invariants therefore computes both the
derived functors in representations and
ordinary group cohomology.
By Proposition
\ref{homological-algebra-proposition-invariants-averaging},
the invariants functor on representations
is exact, so the positive cohomology
of this complex is zero.

Equivalently, if $0\to V\to I\to Q\to0$
is the first step, the long exact sequence
and surjectivity of $I^G\to Q^G$ give
$H^1(G,V)=0$. The higher terms give
$H^{n+1}(G,V)\cong H^n(G,Q)$ for $n>0$;
iteration proves all the remaining cases.
\end{proof}

\section{Hilbert 90}
\label{homological-algebra-section-hilberts-theorem-90-4-sept}

\noindent
I did not attend this lecture and have not yet read the contents of this section.

\noindent
The remainder of the lecture is recorded
on pp.~15--17 of Francisca's notes.
Let $L/K$ be a finite field extension.
Its degree is $[L:K]=\dim_K L$, and
$$
G=\operatorname{Gal}(L/K)
=\{g\in\operatorname{Aut}(L)\mid
g|_K=\operatorname{id}_K\}.
$$
For $H\subset G$, its fixed field is
$$
L^H=\{a\in L\mid h(a)=a
\quad\forall h\in H\}.
$$
We always have $K\subset L^H$.
The finite extension is Galois precisely
when $L^G=K$; then $|G|=[L:K]$.
From now on assume $L/K$ is Galois.

\begin{definition}
\label{homological-algebra-definition-galois-norm-4-sept}
For $\alpha\in L$, let
$m_\alpha\in\operatorname{End}_K(L)$
be multiplication by $\alpha$.
The \emph{norm} is
$$
N_{L/K}(\alpha)=\det(m_\alpha)\in K.
$$
\end{definition}

\noindent
For a finite Galois extension,
$$
N_{L/K}(\alpha)=\prod_{g\in G}g(\alpha).
$$
This formula was stated in the lecture.

\begin{theorem}[Noether]
\label{homological-algebra-theorem-noether-hilbert-90-4-sept}
For a finite Galois extension $L/K$,
$$
H^1(G,L^{\times})=0.
$$
Here $G$ acts on $L^{\times}$ by field
automorphisms.
\end{theorem}

\begin{theorem}[Hilbert 90, cyclic form]
\label{homological-algebra-theorem-hilbert-90-cyclic-4-sept}
If $G=\langle\sigma\rangle$ is cyclic,
then for every $\alpha\in L^{\times}$
with $N_{L/K}(\alpha)=1$ there exists
$\beta\in L^{\times}$ such that
$$
\alpha=\sigma(\beta)/\beta.
$$
\end{theorem}

\noindent
The elementary proof of the cyclic form
was explicitly omitted on the board.
Instead, the lecture asked us to recover
it from the periodic complex for $C_m$.

\begin{exercise}
\label{homological-algebra-exercise-hilbert-90-cyclic-4-sept}
Write the maps in Example
\ref{homological-algebra-example-ai-cyclic-sept-four}
multiplicatively for $A=L^{\times}$.
Check that they are
$$
\beta\longmapsto\sigma(\beta)/\beta,
\qquad \alpha\longmapsto N_{L/K}(\alpha).
$$
Consequently,
$$
H^1(G,L^{\times})
=\frac{\ker(N_{L/K})}
{\{\sigma(\beta)/\beta
\mid\beta\in L^{\times}\}}.
$$
This identifies the cyclic theorem with
the cyclic case of Noether's theorem.
\end{exercise}

\begin{lemma}[Artin's lemma]
\label{homological-algebra-lemma-artin-4-sept}
The automorphisms $g\in G$ are linearly
independent over $L$ as functions $L\to L$.
In other words, if $c_g\in L$ and
$$
\sum_{g\in G}c_g g(x)=0
$$
for every $x\in L$, then all $c_g$ vanish.
\end{lemma}

\noindent
The lemma is stated on the board and
used in the following proof. It concerns
linear independence in
$\operatorname{End}_K(L)$, where $L$
acts by multiplication on the output.

\begin{proof}[Proof of Noether's theorem]
Let $\alpha:G\to L^{\times}$ be a
$1$-cocycle. In multiplicative notation,
$$
\alpha(gh)=\alpha(g)g(\alpha(h)).
$$
By Artin's lemma the function
$$
x\longmapsto\sum_{h\in G}\alpha(h)h(x)
$$
is not zero. Choose $\theta\in L$ so that
$$
b=\sum_{h\in G}\alpha(h)h(\theta)\ne0.
$$
The cocycle identity gives
$$
\begin{aligned}
g(b)&=\sum_{h\in G}g(\alpha(h))gh(\theta)\\
&=\alpha(g)^{-1}
\sum_{h\in G}\alpha(gh)gh(\theta)\\
&=\alpha(g)^{-1}b.
\end{aligned}
$$
Thus $\alpha(g)=b/g(b)$. With
$\beta=b^{-1}$ this is
$\alpha(g)=g(\beta)/\beta$, a coboundary
in our convention. The expressions
$b/g(b)$ and $g(\beta)/\beta$ describe
the same subgroup of coboundaries.
Every cocycle therefore represents zero.
\end{proof}

\begin{remark}[Closing statements]
\label{homological-algebra-remark-galois-closing-4-sept}
The final board also states
$$
H^n(G,L)=0
$$
for $n>0$, with the additive action on
$L$, and mentions $H^2(G,L^{\times})$
as the Brauer group of the extension.
No proof of the additive statement or
construction of the Brauer group is
recorded there.

More precisely, the latter is the
\emph{relative Brauer group} of $L/K$.
For this terminology see
\href{https://www.jmilne.org/math/%
CourseNotes/CFT.pdf}{Milne's notes},
the opening paragraph of Chapter III.
The board also recalls the vanishing for
characteristic-zero representations proved
earlier in the lecture.
\end{remark}


\section{Split sequences and Lie algebra
cohomology}
\label{homological-algebra-section-g-modules}

\noindent
Let $\mathfrak{g}$ be a Lie algebra.
Recall that a $\mathfrak{g}$-module
is a representation of $\mathfrak{g}$.

\begin{proposition}
\label{homological-algebra-proposition-g-mod-abelian-category}
$\mathfrak{g}$-mod is an abelian category.
\end{proposition}

\noindent
Let $B$ be a $\mathfrak{g}$-module and
$A \subset B$ a submodule.
Then $C = B/A$ is the quotient module, and
we have
\begin{equation}
\label{homological-algebra-equation-ses-gmod}
\xymatrix{
0\ar[r]
&A\ar[r]^i
&B\ar[r]^p
&C\ar[r]\ar@/^.5pc/[l]^s
&0,
}
\end{equation}

\noindent
a short exact sequence in
$\mathfrak{g}$-mod.

\begin{definition}
\label{homological-algebra-definition-split}
We say Equation \ref{homological-algebra-equation-ses-gmod}
is a {\it split short exact sequence} 
if there exists the section $s$,
i.e. a map $s:C \to B$ such that $p \circ s
= \operatorname{id}_C$.
\end{definition}

\noindent
Note: if Equation \ref{homological-algebra-equation-ses-gmod}
is split, then $s$ is injective.
Call $\widetilde{C} \subset B$
the image of $s$.
Also call $\tilde{A} =
\operatorname{Im}(i)$.
One can check that
\begin{enumerate}
\item $\widetilde{A}$ and
$\widetilde{C}$ are submodules of $B$.

\item $\tilde{A} \cap \widetilde{C}=0$.

\item $\widetilde{A}+\widetilde{C}=B$.
\end{enumerate}

\noindent
Hence
$$
B=A \oplus C.
$$

\medskip\noindent
If $U \subset V$ is a subrepresentation,
it has a complement $W$ if and only if

$$
\xymatrix{
0\ar[r]
&U\ar[r]^i
&V\ar[r]^p
&V/U\ar[r]\ar@/^.5pc/[l]^s
&0,
}
$$
is split. Indeed, if there exists
a complement $W \subset V$, then
$$
\xymatrix{
W \ar@{^{(}->}[d]\ar[dr]^{\simeq}\\
V \ar[r] & V/U.
}
$$
\medskip\noindent
{\bf Some homological algebra, finally.} 
Let Equation
\ref{homological-algebra-equation-ses-gmod}
be a short exact sequence in
$\mathfrak{g}$-mod. Choose a splitting $s$
at the level of vector spaces.
$p \circ s = \operatorname{id}_C$,
but for $x \in \mathfrak{g}$,
$s(xc) \neq  xs(c)$.
Look at difference $x s(c) - s(xc)$:
$$
\begin{aligned}
p(xs(c)-s(xc)) &= x \psi(c)-\psi(xc)\\
&=xc-xc = 0.
\end{aligned}
$$
So $$
xs(c) - s(xc) \in \ker(p)
= \operatorname{Im}(i) \simeq A.
$$
So we have
$$
\begin{aligned}
\mathfrak{g}\times C &\longrightarrow A \\
x,c &\longmapsto a.
\end{aligned}
$$
This map is bilinear.

So we can think of this function as an
element of
$$
\Hom(\mathfrak{g} \otimes C,A),
$$
or as an element
$$
\xi \in
\Hom(\mathfrak{g},\Hom(C,A)).
$$
$\xi$ depends on the random choice of $s$.

\medskip\noindent
Let $M$ be a $\mathfrak{g}$-module.
Consider
$$
\begin{aligned}
C^0 = M, \qquad 
C^1 = \Hom_{\mathbb{F}}(\mathfrak{g},M),\qquad 
C^2 = \Hom(\Lambda^2\mathfrak{g},M).
\end{aligned}
$$
Let
$$
\begin{aligned}
d: C^0 &\longrightarrow C^1 \\
m &\longmapsto (x \mapsto xm),
\end{aligned}
$$
and
$$
\begin{aligned}
d: C^1 &\longrightarrow C^2 \\
\varphi &\longmapsto 
(x \wedge y \mapsto 
x \varphi(y) - y \varphi(x)-
\varphi([x,y]).
\end{aligned}
$$

Check: $d \circ d =0$.


\begin{definition}
\label{homological-algebra-definition-lie-algebra-cohomology}
The {\it Chevalley-Eilenberg complex} is
$$
C^p(\mathfrak{g},M)
=
\Hom_{\mathbb{F}}
(\Lambda^p\mathfrak{g},M)
$$
for $p=0,1,\dots$, and
$$
\begin{aligned}
d: C^p &\longrightarrow C^{p+1} \\
[d\varphi](x_1 \wedge\dots\wedge x_{p+1} &=
\sum_{i=1}^{p+1}(-1)^{i+1}
x_i \varphi(x_1 \wedge\dots \widehat{x_i}
\wedge\dots \wedge x_{p+1}
\\ &+ \sum_{i<j}(-1)^{i+j}
\varphi([x_i,x_j]\wedge x_1 \wedge\dots
\wedge \widehat{x_i} \wedge\dots
\wedge \widehat{x_j}\wedge
\dots \wedge x_{p+1}).
\end{aligned}
$$

\noindent
{\it Lie algebra cohomology} is
$$
H^n(\mathfrak{g},M)
=
\frac{\ker(d)}{\operatorname{Im}d}.
$$
\end{definition}

\begin{lemma}
\label{homological-algebra-lemma-xi}
$\xi:\mathfrak{g} \to
\Hom_{\mathbb{F}}(C,A)$ is in $\ker(d)$,
where $d$ is defined using some
$\mathfrak{g}$-module structure on
$M=\Hom(C,A)$.
\end{lemma}

\noindent
For $f \in \Hom(C,A)$ and
$x \in \mathfrak{g}$,
$$
\begin{aligned}
xf: C &\longrightarrow A \\
[xg](c) &= xf(c) - f(xc).
\end{aligned}
$$

\begin{exercise}
\label{homological-algebra-exercise-g-mod-structure}
Check this defines a $\mathfrak{g}$-mod
structure and check that $\xi \in \ker(d)$.
\end{exercise}

\begin{exercise}
\label{homological-algebra-exercise-does-not-depend}
Show Lie algebra cohomology does not
depend on the choice of $s$.
\end{exercise}

\begin{exercise}
\label{homological-algebra-exercise-homca}
Check that $\Hom(C,A)$…
\end{exercise}

\begin{remark}
\label{homological-algebra-remark-tensor-g-modules}
If $M,N$ are $\mathfrak{g}$-modules
so is $M \otimes_{\mathbb{F}} N$
and $M^*$.
\end{remark}

\noindent
Now, in the setting of
Equation \ref{homological-algebra-equation-ses-gmod},
$$
\Hom_{\mathbb{F}}(C,A)
\simeq
C^* \otimes A,
$$
check (exercise) that with this
$\mathfrak{g}$-mod structure on $\Hom(C,A)=M
$, $\xi \in
\Hom(\mathfrak{g},M)=C^1(\mathfrak{g},M)$
lies in $Z^1(\mathfrak{g},M)=\ker(d)$.

The point is that given an extension (split
sequence), we find a cohomology class!
That is, there is a corresponcence
$$
\begin{aligned}
\{\operatorname{extensions}
\operatorname{of}C
\operatorname{by} A\}
&\leftrightarrow
H^1(\mathfrak{g},\Hom(C,A))\\
0\to A\to B \to C \to 0
&\mapsto [\xi] \in Z^1/B^1.
\end{aligned}
$$
\noindent
Let $s'$ be another splitting, i.e.,
$p \circ s = p \circ s' =
\operatorname{id}_c$.
Define
$$
\xi'(x) = (c \mapsto xs'c-s'xc).
$$

Consider
$$
\xi' - \xi \in
Z^1(\mathfrak{g},\Hom(C,A))
$$
We want to show
$\xi'-\xi \in B^1$,i.e. $\exists m \in M$
such that $\xi'-\xi = dm$ (exercise, hints
were provided).

Recall that $\operatorname{Ext}(C,A)$ is
the set (category) of extensions modulo 
isomorphism. The above argument shows that
$$
\operatorname{Ext}(C,A)
\to H^1(\mathfrak{g},\Hom(C,A))
$$
is well-defined.

\begin{theorem}
\label{homological-algebra-theorem-extensions-cohomology}
$\psi$ is a bijection
$\operatorname{Ext}(C,A) \to
H^1(\mathfrak{g},\Hom(C,A))$.
\end{theorem}

\begin{proof}
(Start.) If $E: 0 \to A \to B \to C \to 0$
is an extension such that $\psi(E) = 0$,
then $E$ was split.
\end{proof}

\noindent
Now we recall a few facts from Lie theory.

\begin{theorem}[Whitehead]
\label{homological-algebra-theorem-whitehead}
Let $\mathfrak{g}$ be semisimple over
$\mathbb{C}$, and $M$ a finite dimensional
$\mathfrak{g}$-module. Then
$$
H^1(\mathfrak{g},M)=0.
$$
\end{theorem}

\begin{lemma}[Corollary]
\label{homological-algebra-lemma-g-modules-ses-split}
If $\mathfrak{g}$ is semisimple, any short
exact sequence 
$0 \to A \to B \to C \to 0$
of finite-dimensional
$\mathfrak{g}$-modules is split. That is,
if $A \subset B$ is a
$\mathfrak{g}$-submodule,
there exists a complement $C \subset B$
such that $A \cap C = 0$,
$B = A + C$.
\end{lemma}

\section{Spectral sequences}
\label{homological-algebra-section-spectral-sequences}

\noindent
[Large detour on Lie theory and universal
enveloping algebra.]
The following problem shall motivate
our discussion on spectral sequences.
We define
$$
V_n
=
U(\mathfrak{g}) \otimes
\Lambda^n\mathfrak{g},
$$
\begin{equation}
\label{homological-algebra-equation-differential-spectral}
\begin{aligned}
d(u \otimes \xi_1 \wedge \dots \wedge
\xi_n)&=
\sum_{i=1}^n(-1)^{i+1}
(u\xi_i) \otimes \xi_i
\wedge\dots\wedge \widehat{\xi_i}
\wedge\dots \widehat{\xi_j}\wedge
\dots\wedge\xi_n\\
&+\sum_{i<j}(-1)^{i+1}
u \otimes[\xi_i,\xi_j]
\wedge\xi_1\wedge\dots\wedge
\widehat{\xi_i}\wedge\dots\wedge
\widehat{\xi_j}\wedge\dots\wedge\xi_n.
\end{aligned}
\end{equation}


\begin{lemma}
\label{homological-algebra-lemma-chevalley-eilenberg-acyclic}
The (augmented) complex
$$
\xymatrix{
0&\ar[l]\mathbb{F}&
\ar[l]_{\varepsilon}V_0
&\ar[l]_{d}V_1
&\ar[l]_d V_2&
\ar[l]\cdots
}
$$
is acyclic.
\end{lemma}

\noindent
To prove this we shall use filtrations.
The PBW filtration is:
$$
F_pV_n
=\dots
$$

A spectral sequence $\{E^r_{p,n}\}$
is a collection of pages indexed
by $r \in \mathbb{Z}_{\geq 0}$
$$
E^0_{\bullet,\bullet},
E^1_{\bullet,\bullet},\dots
$$
Here $p,n \in \mathbb{Z}$,
and each $E_{p,n}^r$ is an
abelian group, $R$-module, vector space
(this is our current case), etc.

Our filtration $F_pV_n$
is compatible with differentials, meaning
$$
d(F_pV_n) \subset F_pV_{n-1}.
$$

\medskip\noindent
{\bf Claim.} A complex with a compatible filtration gives
a spectral sequence.

Page 0 is:
$$
E^0_{p,n}
=
\operatorname{gr}_pV_n
=
\frac{F_pV_n}{F_{p-1}V_n}
$$
plus some differentials
$$
d:F_pV_n \to F_pV_{n-1},
$$
which, since $F_{p-1}V_n \subset F_pV_n$
and $F_{p-1}V_{n-1}\subset F_p V_{n-1}$,
descends to
$$
d:\operatorname{gr}_pV_n
\to \operatorname{gr}_pV_{n-1}.
$$

Now
$$
E^1_{p,n}
=
H^\bullet(E^0_{p,n})
$$
(or perhaps homology instead of
cohomology…).
The {\bf first nontrivial construction} 
is to define the differentials in 
$E^1_{\bullet,\bullet}$.
This uses the long exact sequence in
homology. It looks like a diagram-chasing
type of argument.

Now
$$
E^2_{p,n}
=
H^\bullet(E^1_{p,n}),
$$
$E^2_{\bullet,\bullet}$ acquires a
differential too, and it goes like
[diagram with diagonal arrows, similar as
the above argument].

This continues forever,
$$
E^r_{p,n}
=
H^\bullet(E^{r-1}_{p,n}),
$$
acquires a differential
of bidigree
$$
(\alpha_r,\beta_r)
=
(\alpha_0,\beta_0)
+r(\alpha,\beta).
$$
So in $E_0$ the bidegree is $(0,-1)$,
which means $d$ looks like vertical arrows
in the $E^0$ page. Then for $E^1$
it has bidegree $(-1,-1)$,
and $E^r$ has bidigree $(-r,-1)$.


\begin{theorem}
\label{homological-algebra-theorem-spectral-sequence}
$H_n(V_\bullet)$ is computed by
$$
\bigoplus_{k \in \mathbb{Z}}
E^{\infty}_{(0,n)+k(\alpha,\beta)}
$$
\end{theorem}

\begin{remark}
\label{homological-algebra-remark-spectral-sequence}
\begin{enumerate}
\item You never really go to $E^\infty$.
In practice, usually $d$ becomes
$d=0$ on $E^1$ or $E^2$ already.
So then $E^N_{p,n}=E^{N+1}_{p,n}
=E^{N+2}_{p,n}=\dots=E^{\infty}_{p,n}$.
We say the spectral sequence 
``collapses'' at page $N$.

\item The precise statement of th etheorem.
Note that $F_pV_n \subset V_n$
induces a subspace
$$
\frac{Z_n \cap F_p}{B_n \cap F_p}
\subset
\frac{Z_n}{B_n} = H_n(V_\bullet).
$$
Denote this
$$
G_p H_n(V_\bullet)
$$
(to not confuse with $F_p$).
$$
G_0H^0 \subset G_1H^0 \subset\dots
$$
so we have $\operatorname{gr}_p H_n$.
The theorem really says
$$
\operatorname{gr}_pH_n(V_\bullet)
=
E^\infty_{p,n}.
$$

\item The power of spectral sequence
usually comes from observing
$$
E^2_{p,n} = 0 \qquad \forall n\geq 1,
$$
or something similar (the point is that,
in practice, most entries in $E^2$ are
zero), and deducing
$$
H_n=0\qquad \forall n \geq n
$$
from this.
\end{enumerate}
\end{remark}

\section{Koszul complex}
\label{homological-algebra-section-koszul-complex}

\noindent
Let us return to our Lie algebra
question which motivated the previous
section, namely
Lemma
\ref{homological-algebra-lemma-chevalley-eilenberg-acyclic}.

We have, for Chevalley-Eilenberg,
an $E^0$ page
$$
E^0_{p,n}
=
\operatorname{gr}_pV_n = 0
$$
whenever $n>p$ because…

In the computation, at some point of
computing the entries to obtain polynomial
rings:
$$
\operatorname{gr}_pV_0
=
\operatorname{gr}_pU(\mathfrak{g})
=
S_p(\mathfrak{g})
\simeq
\mathbb{C}[x^1,\dots,x^N],
$$
where $x^1,\dots,x^N$ is a basis of
$\mathfrak{g}$.
Then we compute the differential
as defined in
Equation
\ref{homological-algebra-equation-differential-spectral}.
We note that the last summand, i.e.
the sum indexed by $i<j$,
lies in $F_{p-1}$ and thus is killed in the
quotient.

In fact,
$$
E^1_{p,q}
=
S_p(\mathfrak{g})
\otimes \Lambda^n\mathfrak{g}
$$
and
$$
d(u \otimes \xi_1\wedge\dots\wedge \xi_n)
=
\sum_{i=1}^n(-1)^{i=1}
(u\xi_i) \otimes \xi_1 \wedge\dots\wedge
\widehat{\xi_i}\wedge\dots\wedge \xi_n.
$$
In particular, on column \#1,
$\mathfrak{g} \to \mathfrak{g}$, it is
\begin{equation}
\label{homological-algebra-equation-differential-koszul}
\begin{aligned}
1 \otimes\xi &\longrightarrow \xi \otimes 1 \\
\mathfrak{g}\otimes\mathbb{F} &\longmapsto 
\mathbb{F}\otimes \mathfrak{g}.
\end{aligned}
\end{equation}

In row \#1 to row \#0,
$$
u \otimes \xi \to u\xi \otimes 1
$$
for $u \in S_p(\mathfrak{g})$ for all
$p \geq 1$.

If we sum all the columns together, we get
the {\it Koszul complex}
$$
K_n:= S(\mathfrak{g}) \otimes
\Lambda^n\mathfrak{g},
$$
with differential given by
Equation \ref{homological-algebra-equation-differential-koszul}.

\medskip\noindent
Now we explain the Koszul complex.
Let $R$ be a commutative ring
($R = S(\mathfrak{g})$),
and $\{x_1,\dots,x_N\}$
a set of elements in $R$
(a basis of $\mathfrak{g} \subset
S(\mathfrak{g})$
in our case).
Let $L$ be the free $R$-module
on basis
$\{x_1,\dots,x_N\}$ 
(write generators as $dx_i$
to distinguish them).
Set
$$
K_n = \Lambda^n_R L
$$
and
$$
\begin{aligned}
d: K_n &\longrightarrow K_{n-1} \\
d(u \otimes dx_{i_1}
\wedge\dots dx_{i_n} &=
\sum_{i=1}^n
(-1)^i (ux_i)
\otimes dx_1 \dots
\widehat{dx_i}\wedge\dots dx_n.
\end{aligned}
$$
\begin{definition}
\label{homological-algebra-definition-regular-sequence}
A sequence $x_1,\dots,x_N$
in $R$ is called {\it regular} 
if $x_i$ is not a zero divisor in
$$
\frac{R}{(x_1,\dots,x_{i-1}}
$$
for all $i$.
\end{definition}

\begin{theorem}
\label{homological-algebra-theorem-koszul-homology}
If $x_1,\dots,x_N$ is a regular sequence
in $R$. Then
$$
H_n(K_\bullet)
=
\begin{cases}
R/(x_1,\dots,x_N)\qquad &n=0 \\
0\qquad &n \geq 1
\end{cases}
$$
\end{theorem}

\begin{proof}
Proved in lecture for the case we need only.
\end{proof}

\noindent
For $R = S(\mathfrak{g}) =
\mathbb{C}[x_1,\dots x_N]$
and $x_1,\dots,x_N$ a basis of
$\mathfrak{g}$. Certainly $x_1,\dots,x_N$
is regular.

We have shown that
$$
E^1_{p,n}
=
S_p(\mathfrak{g}) \otimes
\Lambda^n\mathfrak{g},
$$
$$
\bigoplus_p E^n_{p,n}
=
S(\mathfrak{g}) \otimes
\Lambda^n\mathfrak{g},
$$
and 
$$
d = d_{\operatorname{Koszul}}.
$$
Thus, the $E^2_{p,q}$ of our problem is
zero except in the $(0,0)$ entry, where it
is $\mathbb{F}$, and we conclude that
$$
H_\bullet(V_\bullet)
=
\begin{cases}
\mathbb{F}\qquad &n=0 \\
0\qquad &n\neq 0
\end{cases}
$$
as claimed in
Lemma
\ref{homological-algebra-lemma-chevalley-eilenberg-acyclic}.

\medskip\noindent
Let $C_\bullet$ and $D_\bullet$

be two complexes, say,
$$
d:C_n \to C_{n-1},
\qquad 
d:D_n \to D_{n-1}.
$$
We define $C_\bullet \otimes D_\bullet$
to be the complex
$$
(C_\bullet \otimes D_\bullet)_n
=
\bigoplus_{p+q=n}
C_p \otimes D_q.
$$
If $c \in C_p$, $e \in D_q$,
$$
d(c \otimes e)
=
d c \otimes e +
(-1)^p c \otimes de.
$$
This called ``Koszul rule of signs''.
Then $d \circ d=0$.

\begin{proposition}
\label{homological-algebra-proposition-homology-tensor-koszul}
$H_n(C_\bullet \otimes D_\bullet)
\simeq
(H_\bullet(C_\bullet)
\otimes H_\bullet(D_\bullet))_n$.
\end{proposition}

\begin{proposition}
\label{homological-algebra-proposition-koszul-polynomials}
$K_\bullet(\mathbb{F}[x_1,\dots,x_N])
\simeq
K(\mathbb{F}[x_1]) \otimes
K(\mathbb{F}[x_2]) \otimes\dots
\otimes K(\mathbb{F}[x_N])$.
\end{proposition}

\begin{lemma}[Corollary]
\label{homological-algebra-lemma-homology-koszul-polynomials}
$H_n(K(\mathbb{F}[x_1,\dots,x_N])
=
\begin{cases}
F\qquad &n=0 \\
0\qquad &n \neq 0.
\end{cases}$
\end{lemma}


\section{The cone construction}
\label{homological-algebra-section-cone-construction}

\noindent
(Motivated by algebraic topology).

\begin{definition}
\label{homological-algebra-definition-cone}
Let $A^\bullet,B^\bullet \in
\operatorname{CoCh}$ (cochain complexes
of abelian groups, or whatever abelian
category) and $f: A^\bullet \to B^\bullet$ 
a morphism. The {\it cone of $f$} 
$C(f)$ is the cochain complex
$$
C(f)^n=A^{n+1}\oplus B^n,
$$
$$
d(a,b) = (-d_A(s),d_B(b)+f(a)).
$$
\end{definition}

\noindent
{\bf Notation.} 
\begin{itemize}
\item $A^\bullet[1]$ means the complex.
\item $A[1]^n = A^{n+1}$,
$d_{A[1]}=(-1)^kd_A$,
\item $A[k]^n=A^{n+k}$,
$d_{A[k]}=(-1)^kd_A$.
\end{itemize}

\noindent
Observe that
$$
B^\bullet \hookrightarrow C(f)^\bullet.
$$
Indeed $C(f)^n=A^{n+1}\oplus B^n$,
so $B^n \hookrightarrow C(f)^n$
and the differential
$$
d_{C(f)}:B^n \to B^n.
$$
(There is {\bf not} an embedding
$A^\bullet \hookrightarrow C(f)$
unless $f=0$.)

Consider $C(f)^\bullet/B^\bullet$.
We have $C(f)^\bullet/B^\bullet \simeq
A[1]$.
So we have proved

\begin{proposition}
\label{homological-algebra-proposition-ses-complexes-cone}
There is a short exact sequence of complexes
$$
\xymatrix{
0\ar[r]
&B^\bullet\ar[r]
&C(f)\ar[r]
&A[1]\ar[r]
&0
}
$$
\end{proposition}

\medskip\noindent
{\bf Proof of balancing.}
Recall Theorem \ref{homological-algebra-theorem-balancing}.
One of the uses of balancing is that
we may compute $\operatorname{Ext}^n(A,B)$
``by resolving either $A$ projectively,
or $B$ injectively'' (whatever's best for
the in-practice situation). That is,
$$
H^n\Hom(P^\bullet,B)
\cong
\operatorname{Ext}^n(A,B)
\cong
H^n\Hom(A,I^\bullet).
$$

Let's prove balancing for Ext:

\begin{proof}
(Proof of balancing for Ext.)
Fix $A,B \in R$-mod. Choose an
injective resolution
$0 \to B \to I^\bullet$
and projective resolution
$P^\bullet \to A \to 0$.

To compare 
$$
R^n\Hom_R(A,-)(B)
=
H^n(\Hom_R(A,I^\bullet))
$$
with
$$
L^n\Hom_R(-,B)(A)
=
H^n(\Hom_R(P^\bullet,B)),
$$
let's consider
$$
\Hom_R(P^m,I^n)
:=
C^{m,n} \in \underline{\operatorname{Ab}}.
$$
These naturally form a bicomplex
$$
\operatorname{[diagram]}
$$
$\widetilde{C}$ is the ``augmented'' version
of $C$, using the map $P^0
\xrightarrow{\varepsilon} A$.

What's a bicomplex? It is a collection of
$C^{m,n} \in \underline{\operatorname{Ab}}$,
$m,n \in \mathbb{Z}$, morphisms 
$$
d^h:C^{m,n}\to C^{m+1,n}
$$
$$
d^v:C^{m,n} \to C^{m,n+1}
$$
such that
$$
d^h \circ d^h =0,
\quad 
d^b \circ d^v =0
\quad 
d^h \circ d^v = -d^v \circ d^h.
$$
The {\it total
complex} $\operatorname{Tot}(C)$ of a bicomplex $C$
is
$$
\operatorname{Tot}(C)_n
=
\bigoplus_{p+q=n}C^{p,q}
$$
with
$$
d=d^h + d^v.
$$

\noindent
Exercise:
$\operatorname{Tot}(\widetilde{C})$
is a complex.

(I think the proof now pauses because we
need the following lemma.)
\end{proof}

\begin{lemma}[Acyclic assembly lemma]
\label{homological-algebra-lemma-assembly}
$\operatorname{Tot}(\widetilde{C})$
is acyclic.
\end{lemma}

\noindent
How does it help prove balancing?
We interpret $\operatorname{Tot}(\tilde{C})$
as a cone. The morphism
$P^\bullet.$
$\Hom(A,I^\bullet)$ is a complex,
and $\operatorname{Tot}(C)$is a complex,
and
$$
\omega:\Hom(A,I)
\to \operatorname{Tot}(C)
$$
of complexes.

$$
[\operatorname{diagram}]
$$

We convince ourselves that
$$
C(\omega)
\cong
\operatorname{Tot}(\widetilde{C}),
$$
so we have a short exact sequence
$$
\xymatrix{
0\ar[r]
&\operatorname{Tot}(C)\ar[r]
&\operatorname{Tot}(\widetilde{C})\ar[r]
&\Hom(A,I^\bullet)[1]\ar[r]
&0
}
$$
which yields a long exact sequence in
cohomology. [Diagram.]

Consequently,
$$
[R^n\Hom_R(A,-)
\simeq
H^n(\operatorname{Tot}(C)).
$$
We could, instead, use
$$
B \xrightarrow{\eta} I^0
$$
to add one extra {\bf row} 
to the bottom of $C$,
and do a similar argument to conclude that
$$
H^n(\operatorname{Tot}(C))
\simeq
[L^n\Hom_R(-,B)](A).
$$
That's how acyclic assembly
(the two augmented $\widetilde{C}$
bicomplexes are both
acyclic) implies balancing.

\begin{lemma}[Acyclic assembly lemma,
homological version]
\label{homological-algebra-lemma-assembly-homological}
If $C$ is a bicomplex with
$d^h:C_{p,q} \to C_{p-1,q}$,
$d^v:C_{p,q}\to C_{p,q-1}$,
and $C_{p,q}=0$ for $p<0$ or $q<0$.
If the columns of $C$ are acyclic,
then $\operatorname{Tot}(C)$ is acyclic.
\end{lemma}

\begin{proof}
(Done in lecture.)
\end{proof}

\begin{remark}
\label{homological-algebra-remark-rows-acyclic}
By symmetry, acyclic rows also implies
$\operatorname{Tot}(C)$ acyclic.
\end{remark}

\begin{exercise}
\label{homological-algebra-exercise-balancing-cohomology}
Do the cohomology version.
\end{exercise}

\noindent
Finally, why does the lemma apply to our
$\widetilde{C}$? [Explanation including
diagram.]

\section{Resolution by acyclic objects}
\label{homological-algebra-section-resolution-by-acyclic}

\noindent Upshot: sometimes it can be
painful to compute injective resolutions.
But if you care about (co)homology, it
suffices to find an acyclic resolution.
(This is how to prove de Rham/Dolbeault
theorems.)

\medskip\noindent
For $F$ left exact, we defined $R^nF(A)$ to 
be $H^n(F(I^\bullet))$ where
$$
\xymatrix{
0 \ar[r] & I^0 \ar[r]& I^1 \ar[r]& I^2
\ar[r]&\cdots\\
&A \ar[u]
}
$$
In particular, if $I$ is injective, then
$$
R^nF(I)
=
\begin{cases}
F(I)\qquad &\operatorname{if} n=0 \\
0\qquad &\operatorname{if} n>0.
\end{cases}
$$
\begin{definition}
\label{homological-algebra-definition-f-acyclic}
$X \in \mathcal{A}$ is {\it $F$-acyclic} 
if
$$
R^nF(X)
=
\begin{cases}
F(X)\qquad &\operatorname{if}n=0 \\
0\qquad &\operatorname{if} n>0.
\end{cases}
$$
\end{definition}

\begin{proposition}[4.32, Voisin]
\label{homological-algebra-proposition-432-acyclic}
Let $F:\mathcal{A} \to \mathcal{B}$
be left exact, $\mathcal{A}$ with enough
injectives and $0 \to X \to A^\bullet$
a resolution of $X$ by $A^n$,
$F$-acyclic objects. Then
$$
R^nF(X) \simeq H^nF(A^\bullet)).
$$
\end{proposition}

\begin{proof}
Let $K^n=\ker(d:A^n \to A^{n+1})
\subset A^n$.
We obtain short exact sequences
$$
\xymatrix{
0\ar[r]
&X\ar[r]
&A^0\ar[r]
&K^1\ar[r]
&0
}
$$
and
$$
\xymatrix{
0\ar[r]
&K^n\ar[r]
&A^n\ar[r]^d
&K^{n+1}\ar[r]
&0.
}
$$
Now look at the corresponding long
exact sequences:
$$
\xymatrix{
&R^2F(X)\ar[r]&0\ar[r]&R^2F(K^1)\\
&R^1F(X)\ar[r]&0\ar[r]&R^1F(K^1)
\ar[llu]^{\simeq}\\
0\ar[r]&F(X)\ar[r]&F(A^0)\ar[r]&F(K^1)
\ar[llu]
}
$$
and
$$
\xymatrix{
&R^2F(K^2)\ar[r]&0\ar[r]&R^2F(K^{i+2})\\
&R^1F(K^1)\ar[r]&0\ar[r]
&R^1F(K^{i+1})\ar[llu]^{\simeq}\\
0\ar[r]&F(K^2)\ar[r]&F(A^2)\ar[r]&
F(K^{i+1})\ar[ull]
}
$$
for all $i$.
Thus
$$
R^nF(X) \simeq R^{n-1}F(K^1)
\simeq R^{n-2}F(K^2) \simeq\dots
\simeq R^1F(K^{n-1}),
$$
$$
R^1F(K^{n-1}) \xrightarrow{\simeq}
F(K^n)/\operatorname{Im}(F(A^{n-1})
\xrightarrow{F(d)}F(K^n))
\xrightarrow{\simeq}
\frac{\ker(F(A^n) \to F(A^{n+1})}{
\operatorname{Im}(F(A^{n-1}) \to
F(A^n)}
$$
But $F$ is left exact, so the second arrow
in
$$
0 \to F(K^n) \to F(A^n) \to F(A^{n+1})
$$
is injective, so
$$
\ker(F(A^n)\to F(A^{n+1}))=F(K^n).
$$
\end{proof}

\section{Simplicial homology}
\label{homological-algebra-section-simplicial-homology}

\begin{definition}
\label{homological-algebra-definition-n-simplex}
The {\it $n$-simplex} $\Delta^1$ is the set
of points in $\mathbb{R}^{n+1}$ whose
coordinates add to 1.
A nonempty subset of vertices determines
(by taking convex hull) a {\it face} 
of $\Delta^n$. Often by ``face''
we mean only the $(n-1)$-dimensional
faces; there are $n+1$ of these.
\end{definition}

\noindent
The idea is that we decompose a topological
space as a union of simplices ``glued''
along faces.

\begin{definition}
\label{homological-algebra-definition-abstract-simplicial-cx}
An {\it abstract simplicial complex (ASC)} 
$K$ is a set $\operatorname{Vert}(K)$
and a set $K$ of subsets of
$\operatorname{Vert}(K)$ such that
if $A \in K$ and $B \subset A$,
then $B \in K$.
\end{definition}

\noindent
Think of $\operatorname{Vert}(K)$
as th vertices, and $A \in K$ means
the simplex with its vertices is part of our
space. We write $K_n$ for
$\{A \in K|\# A = n+1\}$
and call these the $n$-simplices of $K$.

The ASCs form a category
$\underline{\operatorname{ASC}}$,
where a morphism $f:K \to K'$
is a function
$f:\operatorname{Vert}(K) \to
\operatorname{Vert}(K')$
such that $f(A) \in K'$ for all
$A \in K'$.

\begin{example}
\label{homological-algebra-example-triangulate-s1}
Triangulate $S^1$ by 1-simplices.
[Figure of circle with three points…]
\end{example}

\begin{definition}
\label{homological-algebra-definition-geometric-realization}
Let $K$ be an ASC.
The {\it geometric realisation} of $K$
is the topological space
$$
\coprod_n K_n \times \Delta^n/ \sim
$$
where, for $A \in K_n$ and
$B \in K_r$, a face of $A$
$(B,x) \sim (A,y)$
means identify the vertex set of $\Delta^r$
with $B$, and of $\Delta^n$
with $A$.
If $x \in \Delta^r$ and we are given a map
$B \to A$,
then $x$ is mapped to a point $y \in
\Delta^n$.
Then $(B,x) \sim (A,y)$.
\end{definition}

\noindent At least morally, the above
definition (which is modified from the book
definition in that we removed the choice of
order in the vertices), we use the
induced topology of $\Delta^n$ as
the geometric $n$-simplex, and the
identification is just gluing the
common sub-simplices of pairs of faces,
like two 2-faces sharing an edge.

\begin{definition}
\label{homological-algebra-definition-asc-n-simplex}
The {\it  $n$-simplex as an object of 
$\underline{\operatorname{ASC}}$}
is an $(n+1)$-element set with all
$2^{n+1}$ of its subsets as simplices.
\end{definition}

\medskip\noindent
{\bf Homology complex.} Let $K \in
\underline{\operatorname{ASC}}$.
Choose an order on $\operatorname{Vert}(K)$.
Now for $\alpha \in K_n$,
denote by $\partial_i \alpha \in K_{n-1}$
the subset obtained by deleting the $i$-th
element of $\alpha$.

Define 
$$
C_n(K) = \mathbb{Z}K_n
= \bigoplus_{\alpha \in
K_n}\mathbb{Z}[\alpha],
$$
$$
\begin{aligned}
d: C_n &\longrightarrow C_{n-1} \\
d[\alpha] &= \sum_{i=0}^n
(-1)^i[\partial_i\alpha].
\end{aligned}
$$
\begin{exercise}
\label{homological-algebra-exercise-d2-0}
$d \circ d =0$.
\end{exercise}

\begin{definition}
\label{homological-algebra-definition-simplicial-homology}
The {\it homology} of $K$ is
$$
H_n(K) = \frac{Z_n}{B_n}
$$
where $Z_n = \ker d \subset C_n$
and $B_n = \operatorname{Im}(d) \subset
C_n$.
\end{definition}

\medskip\noindent
{\bf Delta sets.}
Recall from Infinity Categories
that the simplex category, now denoted
$\hat{\Delta}$, has for objects the
ordered sets
$[n] = \{0,1,\dots,n\}$,
$n \in \mathbb{Z}_{\geq 0/}$ and morphisms are
strictly
increasing functions.

\begin{definition}
\label{homological-algebra-definition-delta-set}
A {\it delta set} is a functor
$$
S : \hat{\Delta}^{\mathrm{op}}
\to \underline{\operatorname{Set}}.
$$
\end{definition}

\noindent
More explicitly, such $S$ consists
of:
\begin{itemize}
\item for each $n \geq 0$,
a set $S_n$, i.e. $S([n])$,
\item functions
$$
\partial_n S_n \to S_{n-1}
$$
for $i = 0,\dots,n$, i.e. $\partial_i =
S(\varepsilon_i)$. [We introduced $\varepsilon$
in an example I didn't copy.]
\end{itemize}

\begin{remark}
\label{homological-algebra-remark-delta-sets-vs-simplicial-sets}
This is closely related to the simplicial sets
introduced in Infinity Categories, but slightly
more particular. There we used the full simplex
category $\Delta$, with all order-preserving
maps, so a simplicial set

$$
S:\Delta^{\operatorname{op}}
\to \underline{\operatorname{Set}}
$$

has both face maps and degeneracy maps.

Here we only use the subcategory
$\Delta_{\mathrm{inj}}\subset\Delta$ with the
same objects and only strictly increasing maps.
Thus a delta set

$$
S:\Delta_{\mathrm{inj}}^{\operatorname{op}}
\to \underline{\operatorname{Set}}
$$

has face maps but no degeneracy maps. An
ordered abstract simplicial complex gives such
an object by taking $S_n=K_n$.

The point of the more particular notion is
simplicity: face maps are already enough to
define

$$
d=\sum_i(-1)^id_i
$$

and hence simplicial homology. Passing to full
simplicial sets adds degeneracies, which are
needed for the richer categorical and
homotopical constructions used in Infinity
Categories, such as nerves, Kan complexes and
$\infty$-categories.
\end{remark}


\noindent
Now our more general triangulations
can be modelled as $\Delta$-sets.


\begin{definition}
\label{homological-algebra-definition-homology-complex}
If $S: \hat{\Delta}^{\mathrm{op}}
\to \underline{\operatorname{Set}}$
is a $\Delta$-set, then its {\it homology
complex} is
$$
C_n(S) = \mathbb{Z} S_n,
$$
$$
\begin{aligned}
\partial: C_n &\longrightarrow C_{n-1} \\
\partial[\alpha] &=
\sum_{i=0}^n (-1)^i[\partial_i\alpha].
\end{aligned}
$$
\end{definition}

\noindent
For every $K \in \underline{\operatorname{ASC}}$
and choice of order of
$\operatorname{Vert}(K)$, it's possible
to construct a $\Delta$-set, via
$$
S_n=K_n,
$$
and for $A = \{v_{j_0},v_{j_1},\dots,
v_{j_n}\}$ and $\partial_iA 
=
\{v_{j_0},\dots,\widehat{v_{j_i}},\dots\}$.

This is not a functor.
Indeed, the morphism […]

There {\it is} a way to go from
$\underline{\operatorname{ASC}}$
to chain complexes without choosing
an order on $\operatorname{Vert}(K)$. The
following definition is just a clever way to
do this (no order, but same data):

\begin{definition}
\label{homological-algebra-definition-chain-cx-asc}
Let $K \in \underline{\operatorname{ASC}}$.
For $p \geq 0$ define 
$$
C^{\operatorname{ori}}_p(K)
$$
to be the free abelian group
on the symbols $\langle
v_0,v_1,\dots,v_p\rangle$,
where $v_i$ are elements of a single
$p$-simplex, modulo relations
$$
\langle v_{\sigma(0)},
v_{\sigma(1)},\dots,
v_{\sigma(p)}\rangle
=
\varepsilon(\sigma)
\langle v_0,v_1,\dots,v_p\rangle,
$$
for $\sigma \in S_{p+1}$,
and
$$
\delta \langle v_0,\dots,v_p\rangle
=
\sum_{i=0}^p (-1)^i
\langle
v_0,\dots,\widehat{v_i},\dots,v_p\rangle.
$$
\end{definition}

Where ori stands for oriented.

In fact, $K \mapsto
C_\bullet^{\operatorname{ori}}(K)$
defines a functor
$$
\underline{\operatorname{ASC}} 
\to \operatorname{Ch}_{\geq 0}.
$$

\begin{definition}
\label{homological-algebra-definition-big-delta-set}
The {\it (big) $\Delta$-set} associated
with an ASC is $S^k$ defined by
$$
S_n^k
$$
which is defined as the set of
$f:[n] \to \operatorname{Vert}(K)$
such that the image of $f$ lies in
a single simplex of $K$.

Notice that $S^k$ is indeed big,
$$
S_n^k \neq  \{\}
$$
for all $n$. There are many non-injective
functions $f : [n] \to
\operatorname{Vert}(K)$,
these are called {\it degenerate simplices}.

\end{definition}


\begin{remark}
\label{homological-algebra-remark-big-delta-set-simplicial-set}
The big $\Delta$-set is very close to the
simplicial sets from Infinity Categories.
Its non-injective maps $[n]\to
\operatorname{Vert}(K)$ are precisely the
degenerate simplices. If we let all
order-preserving maps in the full simplex
category $\Delta$ act by precomposition, then
we obtain an actual simplicial set

$$
S^K:\Delta^{\operatorname{op}}
\to \underline{\operatorname{Set}}.
$$

\end{remark}



\section{Simplicial approximation and homotopy}
\label{homological-algebra-section-simplicial-approximation-homotopy}

\noindent
The following reconstructs the lectures from the
blackboard notes and the notes shared by colleagues.

\begin{definition}
\label{homological-algebra-definition-simplicial-map}
Let $K,L$ be abstract simplicial complexes.
A map
\[
f:\operatorname{Vert}(K)\longrightarrow
\operatorname{Vert}(L)
\]
is a {\it simplicial map} if the image of every
simplex of $K$ is contained in a simplex of $L$.
It induces a continuous map
\[
|f|:|K|\longrightarrow |L|
\]
which is affine on each simplex.
\end{definition}

\begin{definition}
\label{homological-algebra-definition-star-of-simplex}
For a simplex $\sigma$ of $K$, its (open) star is
the union of the interiors of simplices containing
$\sigma$.
\end{definition}

\begin{definition}
\label{homological-algebra-definition-simplicial-approximation}
Let $h:|K|\to |L|$ be continuous. A simplicial map
$f:K\to L$ is a {\it simplicial approximation} to
$h$ if, for every vertex $v$ of $K$,
\[
h(\operatorname{St}(v))
\subset
\operatorname{St}(f(v)).
\]
\end{definition}

\begin{theorem}[Simplicial approximation theorem]
\label{homological-algebra-theorem-simplicial-approximation}
After sufficiently many subdivisions of $K$, every
continuous map
\[
h:|K|\longrightarrow |L|
\]
admits a simplicial approximation.
\end{theorem}

\noindent
The point is that continuous maps can therefore be
replaced, up to homotopy, by combinatorial maps.

\begin{definition}
\label{homological-algebra-definition-contiguous-simplicial-maps}
Two simplicial maps $f,g:K\to L$ are
{\it contiguous} if, for every simplex
$\sigma\subset K$, the union
\[
f(\sigma)\cup g(\sigma)
\]
is contained in a simplex of $L$.
\end{definition}

\begin{proposition}
\label{homological-algebra-proposition-contiguous-chain-homotopic}
Contiguous simplicial maps induce chain-homotopic
maps on simplicial chains, and hence the same map
on homology.
\end{proposition}

\section{Acyclic carriers}
\label{homological-algebra-section-acyclic-carriers}

\begin{definition}
\label{homological-algebra-definition-acyclic-chain-complex}
A chain complex $C_\bullet$ is {\it acyclic} if
\[
H_n(C_\bullet)=0
\]
in positive degrees (with the usual augmentation
understood in degree $0$ when working with reduced
homology).
\end{definition}

\begin{definition}
\label{homological-algebra-definition-acyclic-carrier}
Let $C_\bullet,D_\bullet$ be chain complexes with
chosen bases. An {\it acyclic carrier} $\Phi$
assigns to each basis element $\sigma$ of $C$ an
acyclic subcomplex
\[
\Phi(\sigma)\subset D_\bullet
\]
compatibly with faces: if $\tau$ is a face of
$\sigma$, then
\[
\Phi(\tau)\subset \Phi(\sigma).
\]
A chain map $f:C_\bullet\to D_\bullet$ is
{\it carried by $\Phi$} if
\[
f(\sigma)\in \Phi(\sigma)
\]
for every basis element $\sigma$.
\end{definition}

\begin{theorem}[Acyclic carrier theorem]
\label{homological-algebra-theorem-acyclic-carrier}
Two augmentation-preserving chain maps carried by
the same acyclic carrier are chain homotopic.
\end{theorem}

\noindent
This is the general mechanism behind the preceding
statement about contiguous simplicial maps: both
maps land, simplex by simplex, in the chain complex
of one common simplex, and that complex is acyclic.

\section{Singular homology}
\label{homological-algebra-section-singular-homology}

\begin{definition}
\label{homological-algebra-definition-singular-simplex}
Let $X$ be a topological space. A {\it singular
$n$-simplex} in $X$ is a continuous map
\[
\sigma:\Delta^n\longrightarrow X.
\]
The singular chain group is the free abelian group
\[
C_n^{\operatorname{sing}}(X)
=
\mathbb Z\{\sigma:\Delta^n\to X\}.
\]
\end{definition}

\noindent
If $\delta_i:\Delta^{n-1}\to\Delta^n$ is the
inclusion of the $i$-th face, define
\[
\partial\sigma
=
\sum_{i=0}^n(-1)^i\,\sigma\circ\delta_i.
\]
Then $\partial^2=0$, so
$C_\bullet^{\operatorname{sing}}(X)$ is a chain
complex.

\begin{definition}
\label{homological-algebra-definition-singular-homology}
The {\it singular homology} of $X$ is
\[
H_n^{\operatorname{sing}}(X)
=
H_n(C_\bullet^{\operatorname{sing}}(X)).
\]
\end{definition}

\begin{proposition}
\label{homological-algebra-proposition-singular-homology-functor}
A continuous map $f:X\to Y$ induces a chain map
\[
f_\#:C_\bullet^{\operatorname{sing}}(X)
\longrightarrow
C_\bullet^{\operatorname{sing}}(Y),
\qquad
f_\#(\sigma)=f\circ\sigma,
\]
and hence maps
\[
f_*:H_n^{\operatorname{sing}}(X)
\longrightarrow H_n^{\operatorname{sing}}(Y).
\]
Thus singular homology is functorial.
\end{proposition}

\section{Homotopy invariance of singular homology}
\label{homological-algebra-section-homotopy-invariance-singular}

\begin{definition}
\label{homological-algebra-definition-homotopy-maps}
Two continuous maps $f,g:X\to Y$ are homotopic if
there exists
\[
H:X\times I\longrightarrow Y
\]
with $H(-,0)=f$ and $H(-,1)=g$.
\end{definition}

\begin{theorem}[Homotopy invariance]
\label{homological-algebra-theorem-homotopy-invariance-singular}
If $f\simeq g:X\to Y$, then
\[
f_*=g_*:
H_n^{\operatorname{sing}}(X)
\longrightarrow
H_n^{\operatorname{sing}}(Y)
\]
for every $n$.
\end{theorem}

\noindent
The proof constructs a {\it prism operator}
\[
P:C_n^{\operatorname{sing}}(X)
\longrightarrow
C_{n+1}^{\operatorname{sing}}(Y)
\]
from a triangulation of
$\Delta^n\times I$ and checks
\[
\partial P+P\partial=g_\#-f_\#.
\]
Thus $f_\#$ and $g_\#$ are chain homotopic.

\section{Simplicial versus singular homology}
\label{homological-algebra-section-simplicial-versus-singular}

\begin{theorem}
\label{homological-algebra-theorem-simplicial-singular-agree}
If $K$ is an abstract simplicial complex, then
\[
H_n^{\operatorname{simp}}(K)
\cong
H_n^{\operatorname{sing}}(|K|)
\]
naturally in $K$.
\end{theorem}

\noindent
The lecture proof compares the simplicial and
singular chain complexes using simplicial
approximation together with the acyclic-carrier
argument. The important point is that simplicial
homology depends on a triangulation in its
definition, while singular homology is defined
directly from the topological space, and the
theorem says that they compute the same invariant
for a triangulated space.

\section{Homology and cohomology with coefficients}
\label{homological-algebra-section-homology-cohomology-coefficients}

\begin{definition}
\label{homological-algebra-definition-singular-homology-coefficients}
For an abelian group (or ring) $R$, define
\[
C_n^{\operatorname{sing}}(X;R)
=
C_n^{\operatorname{sing}}(X)\otimes_{\mathbb Z}R.
\]
Its homology is denoted
\[
H_n^{\operatorname{sing}}(X;R).
\]
\end{definition}

\begin{definition}
\label{homological-algebra-definition-singular-cochains}
The singular cochain complex with coefficients in
$R$ is
\[
C^n_{\operatorname{sing}}(X;R)
=
\operatorname{Hom}_{R}
(C_n^{\operatorname{sing}}(X;R),R),
\]
with coboundary induced by the boundary on chains.
Its cohomology is
\[
H^n_{\operatorname{sing}}(X;R).
\]
\end{definition}

\section{Universal coefficient theorem}
\label{homological-algebra-section-universal-coefficient-theorem}

\noindent
Write
\[
Z_n=\ker(\partial:C_n\to C_{n-1}),
\qquad
B_n=\operatorname{Im}(\partial:C_{n+1}\to C_n).
\]
Then
\[
0\longrightarrow B_n\longrightarrow Z_n
\longrightarrow H_n\longrightarrow0
\]
is exact. Since the singular chain groups are
free, the lecture chooses a free resolution of
$H_n$ and applies $\operatorname{Hom}_R(-,R)$.
The resulting long exact sequence identifies the
two pieces which occur in cohomology.

\begin{theorem}[Universal coefficient theorem for cohomology]
\label{homological-algebra-theorem-universal-coefficient-cohomology}
For singular chains with coefficients in $R$,
there is a natural short exact sequence
\[
0\longrightarrow
\operatorname{Ext}^1_R(H_{n-1}(X;R),R)
\longrightarrow
H^n(X;R)
\longrightarrow
\operatorname{Hom}_R(H_n(X;R),R)
\longrightarrow0.
\]
\end{theorem}

\noindent
The blackboard proof is precisely the
homological-algebra construction behind this
statement: use the exact sequences containing
$B_n,Z_n,C_n$, apply $\operatorname{Hom}_R(-,R)$,
and identify the kernel and cokernel of the
connecting maps.

\section{Sheaves}
\label{homological-algebra-section-sheaves-pointer}

\noindent
The sheaf-theoretic material from the course is
kept in the Ringed Spaces notebook, beginning with
Section \ref{homological-algebra-section-presheaves-stalks-sheaves}.
That section is a verbatim transcription of
Chapter 5 of Jethro van Ekeren's course notes and
includes stalks and sheafification, morphisms of
sheaves, restriction and local systems, the
singular-cochain resolution of the constant sheaf,
and direct and inverse image functors.

\section{Flabby sheaves and cohomology}
\label{homological-algebra-section-flabby-sheaves}

\noindent
{\bf Upshot.} Flabby sheaves have enough
extendable sections to eliminate higher
sheaf cohomology. Resolving a sheaf by them
therefore computes its cohomology, by
Section \ref{homological-algebra-section-resolution-by-acyclic}.

\medskip\noindent
{\bf Lecture reconstruction, 14 October
2026.} The corollary recorded from the
board matches the argument in Jethro van
Ekeren's
\href{https://w3.impa.br/~jethro/%
2026-2/ah-current.pdf}{course notes},
Section 6.1, Definition 6.1.4 and
Proposition 6.1.4(2), pp.~117--118.
This identifies the topic, but does not
verify the full contents of today's class.
Sheaf basics remain in Ringed Spaces.

\begin{definition}
\label{homological-algebra-definition-flabby}
A sheaf $\mathcal{F}$ of abelian groups
on $X$ is {\it flabby}, or {\it flasque},
if every restriction map
$
\mathcal{F}(U)\longrightarrow\mathcal{F}(V)
$
is surjective for open $V\subset U$.
\end{definition}

\noindent
For example, arbitrary integer-valued
functions form a flabby sheaf: extend by
zero. Continuous functions need not be
flabby: $1/x$ on $(0,1)$ cannot extend
continuously to $\mathbb{R}$.

\begin{lemma}
\label{homological-algebra-lemma-flabby-kernel-sections}
Suppose
$
0\to\mathcal{A}\to\mathcal{B}
\to\mathcal{C}\to0
$
is exact and $\mathcal{A}$ is flabby.
For every open $U$, the map
$
\mathcal{B}(U)\to\mathcal{C}(U)
$
is surjective.
\end{lemma}

\begin{proof}
A section of $\mathcal{C}$ lifts locally.
Two lifts differ on their overlap by an
$\mathcal{A}$-section. Extend this difference
using flabbiness, subtract it from one
lift, and glue. Zorn's lemma gives a
maximal lift; the same correction enlarges
its domain unless it already covers $U$.
See the proof of Stacks Project,
\href{https://stacks.math.columbia.edu/%
tag/09SV}{Lemma 20.12.3}.
\end{proof}

\begin{lemma}[Corollary]
\label{homological-algebra-lemma-flabby-quotient}
In the preceding exact sequence, if
$\mathcal{A}$ and $\mathcal{B}$ are flabby,
then $\mathcal{C}$ is flabby.
\end{lemma}

\begin{proof}
Given $V\subset U$ and
$c\in\mathcal{C}(V)$, lift $c$ to
$b\in\mathcal{B}(V)$ using
Lemma \ref{homological-algebra-lemma-flabby-kernel-sections}.
Extend $b$ to $U$ using flabbiness of
$\mathcal{B}$. Its image extends $c$.
\end{proof}

\begin{proposition}
\label{homological-algebra-proposition-flabby-acyclic}
If $\mathcal{F}$ is flabby, then
$
H^q(X,\mathcal{F})=0
$
for every $q>0$.
\end{proposition}

\begin{proof}
Embed $\mathcal{F}$ into an injective
$\mathcal{I}$ and write
$\mathcal{Q}=\mathcal{I}/\mathcal{F}$.
Injective sheaves are flabby, so the
corollary makes $\mathcal{Q}$ flabby.
The preceding lemma and the long exact
cohomology sequence give
$
H^1(X,\mathcal{F})=0,\qquad
H^{q+1}(X,\mathcal{F})
\cong H^q(X,\mathcal{Q})
$
for $q\geq1$. Induction proves the claim.
References: \cite[III, Lemma 2.4 and
Proposition 2.5]{hart}; Stacks Project,
\href{https://stacks.math.columbia.edu/%
tag/09SV}{Lemmas 20.12.2--20.12.3}.
\end{proof}

\medskip\noindent
{\bf Why the board corollary matters.}
It keeps the quotient flabby, allowing
dimension shifting to continue.
Consequently a flabby resolution
$
0\to\mathcal{F}\to\mathcal{G}^0
\to\mathcal{G}^1\to\cdots
$
computes
$
H^q(X,\mathcal{F})
\cong H^q(\Gamma(X,\mathcal{G}^{\bullet})).
$
This applies the acyclic-resolution
principle already in these notes.


\subsection{Godement envelopes and enough injectives}
\label{homological-algebra-subsection-godement-envelopes}
For a sheaf $\mathcal F$ on $X$, the Godement envelope is
$G(\mathcal F)(U)=\prod_{x\in U}\mathcal F_x$.
The map $\mathcal F\to G(\mathcal F)$ sending sections to germs
is a monomorphism.  The envelope is flabby, since families of
germs can be extended by zero.  For sheaves of abelian groups,
choose an embedding $\mathcal F_x\hookrightarrow I_x$ into an
injective abelian group at each point and define
$\mathcal I(U)=\prod_{x\in U}I_x$.
This is injective as a sheaf: maps into $\mathcal I$ correspond
to maps into $I_x$ on each stalk, and those extend across
monomorphisms by injectivity of $I_x$.  Thus the category of
sheaves of abelian groups has enough injectives.
Reference: Jethro van Ekeren, \emph{Homological Algebra},
Section 6.1, Theorem 6.1.1, pp.~115--116.


\section{\v{C}ech cohomology}
\label{homological-algebra-section-cech-cohomology-jethro-2026}

\noindent
Lecture notes, 16 October 2026. Reference: Jethro van Ekeren,
\href{https://w3.impa.br/~jethro/2026-2/ah-current.pdf}
{Homological Algebra}, Section 6.2.1, p.~118.
The following is the student's first chunk, preserved verbatim
in wording (with mathematical notation typeset in LaTeX).

\subsection{Definition and elementary properties}
\label{homological-algebra-subsection-cech-definition-elementary}

\medskip
Convenient tool for computing sheaf cohomology. Uses a choice of open cover
$\mathcal U = \{U_i\}_{i\in I}$ of $X$.

\begin{lemma}
If $\mathcal I$ in $\operatorname{Sh}(X)$ is an injective sheaf, then the restriction
$\mathcal I|_U$ to any open $U\subset X$ is injective in $\operatorname{Sh}(U)$.
\end{lemma}

Pf. Exercise.

\begin{definition}
(locally finite cover).
\end{definition}

Note. The definition of Chech cohomology we give next does not require that the cover be locally finite, but many theorems that follow do.

\medskip
\noindent
The idea is: grab an open cover, for any finite set of open sets, consider the intersection (which is open because their are finitely many). Now for any sheaf $\mathcal F$ consider $\mathcal F$ of those intersections, and then the $p$-th chain group shall be the product of all those $\mathcal F(U_s)$.

\medskip
Let $\mathcal U=\{U_i\}_{i\in I}$ be a cover of $X$. For any $S=\{i_0,\ldots,i_p\}\subset I$, denote
\[
U_s=U_{i_0}\cap\cdots\cap U_{i_p}.
\]

Let $\mathcal F$ be a sheaf of abelian groups on $X$. Define
\[
\check C^p(\mathcal U,\mathcal F)
=\prod_{S\subset I,\ \#S=p+1}\mathcal F(U_s).
\]

For instance $\check C^0(\mathcal U,\mathcal F)=\prod_{i\in I}\mathcal F(U_i)$.

To define $d$, we fix an order on $I$, now $\alpha$ in $\check C^p(\mathcal U,\mathcal F)$ consists of a section

$\alpha_{i0\ldots}$


\subsection*{From derived global sections to singular cohomology}

\noindent
{\bf Upshot.} On a topological manifold $X$,
sheaf cohomology of the constant sheaf recovers singular cohomology:
\[
H^n(X,\underline{\mathbb Z})
=R^n\Gamma(X,\underline{\mathbb Z})
\cong H^n_{\mathrm{sing}}(X;\mathbb Z).
\]
More generally, the comparison holds for sufficiently well-behaved
spaces, e.g.\ paracompact locally contractible spaces.
This is the assertion of Jethro van Ekeren's
\href{https://w3.impa.br/~jethro/2026-2/ah-current.pdf}
{notes}, Theorem 6.2.16 and surrounding discussion in
Section 6.2.7 (stated there with real coefficients).

\noindent
The mechanism is to resolve the constant sheaf by sheaves of
singular cochains, as constructed in Ringed Spaces,
Definition~\ref{ringed-spaces-definition-singular-cochain-sheaf}.
Acyclicity of the terms permits computation of derived global
sections by an ordinary cochain complex; comparison with the
usual singular cochain complex then identifies the resulting
cohomology. These are logically distinct steps: a resolution
alone does not yet prove the global comparison.



\subsection{The Leray theorem}
\label{homological-algebra-subsection-cech-leray}
Recalling the definition and elementary properties of the previous subsection, choose an ordering
of the cover indices and define
\[
(d\alpha)_{i_0\ldots i_{p+1}}=
\sum_{j=0}^{p+1}(-1)^j
\alpha_{i_0\ldots\widehat{i_j}\ldots i_{p+1}}
\big|_{U_{i_0}\cap\cdots\cap U_{i_{p+1}}}.
\]
The alternating signs imply $d^2=0$, and the cohomology
of the resulting complex is $\check H^p(\mathcal U,\mathcal F)$.
\begin{definition}
A cover $\mathcal U$ is $\mathcal F$-acyclic if, for every
nonempty finite intersection $U_S$ of members of the cover,
$H^q(U_S,\mathcal F|_{U_S})=0$ for all $q>0$.
\end{definition}
\begin{theorem}[Leray]
If $\mathcal U$ is $\mathcal F$-acyclic then the canonical map
$\check H^p(\mathcal U,\mathcal F)\to H^p(X,\mathcal F)$
is an isomorphism for every $p\geq0$.
\end{theorem}
\begin{proof}
Take an injective resolution $\mathcal F\to\mathcal I^\bullet$
and the double complex
$K^{p,q}=\check C^p(\mathcal U,\mathcal I^q)$.
Vertical-first cohomology has
$E_1^{p,q}=\prod_{|S|=p+1}H^q(U_S,\mathcal F)$,
because restrictions of injective sheaves to opens remain
injective.  By the hypothesis only the $q=0$ row remains,
and this is the \v{C}ech complex of $\mathcal F$.
Horizontally first, each $\mathcal I^q$ is flabby, so
higher \v{C}ech cohomology vanishes and the remaining
column is $\Gamma(X,\mathcal I^\bullet)$.
Both compute the cohomology of the same total complex.
\end{proof}
Reference: Jethro van Ekeren, Section 6.2.2,
Theorem 6.2.7, pp.~121--122; \cite[III, Section 4]{hart}.



\subsection{Example: Sheaf cohomology on the elliptic curve}
\label{homological-algebra-subsection-cech-elliptic-curve}
Let $E$ be a smooth complex elliptic curve, and choose
distinct points $p,q\in E$. In the algebraic setting,
$U=E\setminus\{p\}$, $V=E\setminus\{q\}$ and their
intersection are affine. They are $\mathcal O_E$-acyclic,
so Leray identifies sheaf cohomology with the two-term
\v{C}ech complex
\[
0\to\mathcal O_E(U)\oplus\mathcal O_E(V)
\xrightarrow{\delta}\mathcal O_E(U\cap V)\to0,
\qquad\delta(f,g)=g|_{U\cap V}-f|_{U\cap V}.
\]
The kernel is the space of global regular functions,
hence $H^0(E,\mathcal O_E)=\mathbb C$.
The cokernel is one-dimensional, and there are no
higher terms:
\[
H^1(E,\mathcal O_E)\cong\mathbb C,\qquad
H^k(E,\mathcal O_E)=0\quad(k>1).
\]
For example the dimension of $H^1$ also follows from
Riemann--Roch for genus one.
Jethro's explicit calculation in Section 6.2.3,
pp.~123--124, uses the Weierstrass $\wp$-function
and the difference
$\zeta(z-p)-\zeta(z-q)$ as an elliptic meromorphic
function with opposite residues at $p$ and $q$.
A caution about that presentation: holomorphic
functions on punctured compact curves can have
essential singularities, so the claimed spanning
sets should be interpreted in the meromorphic/
algebraic context, not as spanning all analytic
holomorphic sections on the punctured curve.


\subsection{Fine sheaves and the de Rham theorem}
\label{homological-algebra-subsection-fine-sheaves-and-de-rham}
\noindent
{\bf Upshot.} On a paracompact smooth manifold, partitions of
unity make the sheaves of differential forms acyclic. Together
with the Poincar\'e lemma this identifies de Rham cohomology
with sheaf cohomology of the constant sheaf.

A sheaf is \emph{fine} if locally finite open covers admit
endomorphisms supported in the corresponding open sets which
sum to the identity. A sheaf of modules over the sheaf
$C^\infty_X$ is fine: choose a subordinate partition of unity
$\{\rho_i\}$ and multiply sections by $\rho_i$.
On a paracompact space fine sheaves are acyclic.

Write $\mathcal A_X^p$ for the sheaf of smooth real
$p$-forms on a smooth manifold $X$. The Poincar\'e lemma
gives an exact sequence of sheaves
\[
0\longrightarrow\underline{\mathbb R}
\longrightarrow\mathcal A_X^0\xrightarrow{d}
\mathcal A_X^1\xrightarrow{d}\mathcal A_X^2
\xrightarrow{d}\cdots.
\]
Since every $\mathcal A_X^p$ is fine and acyclic,
the acyclic-resolution theorem gives
\[
H^k(X,\underline{\mathbb R})
\cong H^k(\Gamma(X,\mathcal A_X^\bullet))
=H_{\mathrm{dR}}^k(X).
\]
Combining with the comparison to singular cohomology,
$H^k_{\mathrm{dR}}(X)\cong H^k_{\mathrm{sing}}(X;\mathbb R)$.
Reference: Jethro van Ekeren, Sections 6.2.4--6.2.7,
pp.~125--129 (including Proposition 6.2.10).


% --- End homological-algebra.tex ---

% --- Begin infty-categories.tex ---
\chapter{Infty-categories}
\label{chapter-infty-categories}





\section{Discrete fibrations}
\label{infty-categories-section-discrete-fibrations}

\begin{definition}
\label{infty-categories-definition-discrete-fibration}
\begin{reference}
\url{https://ncatlab.org/nlab/show/discrete+%
fibration}
\end{reference}
A functor $F:C \to B$ is a {\it discrete
fibration} if for every object
$c$ in $C$ and every morphism of the form $g:
b \to F(c)$ in $B$
there is a unique morphism $h:d \to c$ in $C$
such that $F(h)=g$.
\end{definition}

\begin{exercise}
\label{infty-categories-exercise-discr-fibra-presh}
Prove that discrete fibrations over a
category $C$ correspond to presheaves over
$C$.
\end{exercise}

\begin{proof}
First suppose that we are given a presheaf
$X$ on $C$. Define a discrete
fibration $F:C/X \to C$ by $(a,s)\mapsto a$
on objects and
mapping a morphism $f:(a,s)\to (b,t)$ in
$C/X$ to the
corresponding morphism $a \to b$ in $C$.
To show $F$ is a discrete fibration
let $g:b\to a$ be a morphism in $C$.
Consider $g^*=X(g):X_a\to X_b$,
and the section $g^*s$ of $X_b$.
Then the
morphism $h:(b,g^*s)\to (a,s)$ is the only
one mapping to $g$ under $h$.

For the converse let $F:B \to C$ be a
discrete fibration over $C$.
To define a presheaf $X:C^{\text{op}}\to
\text{Sets}$ let $c \in \Ob C$.
We assign the set (for now I won't justify
why this is a set)
of objects in $B$ mapped to $c$ under $F$.
To define the correspondence on morphisms,
consider a map $f:c \to d$ in
$C^{\text{op}}$.
In other words, we have a map in  $C$ of the
form $f^{\text{op}}:d \to c$.
Then to any object in $b$ such that $F(b)=c$,
by definition of discrete fibration,
we have a unique morphism of $B$ of the form
$h:r \to b$ such that $F(h)=f^{\text{op}}$.
In particular this means that $F(r)=d$.
This gives a function from $X(c)$ to $X(d)$.
This situation is described in the following
diagram:

$$
\xymatrix{
X(d)\ni&r
\ar@{.>}[d]_{\exists!h}\ar@{|->}[r]& d=F(r)
\ar[d]^{f^{\text{op}}=F(h)}\\
X(c)\ni& b\ar@{|->}[r]_F&c=F(b)
}
$$




To check functoriality of $X$ defined in the
previous paragraph
suppose that $f:c\to d$ and $g:d\to e$ are
two morphisms in $C^{\text{op}}$.
Like before, we have maps $f^{\text{op}}:d\to
c$ and $g^{\text{op}}:e \to d$.
$$
\xymatrix{
X(e)\ni& q\ar@{.>}[d]_{\exists
!j}\ar@{|->}[r]& e=F(q)
\ar[d]^{g^{\text{op}}=F(j)}\\
X(d)\ni&r
\ar@{.>}[d]_{\exists!h}\ar@{|->}[r]& d=F(r)
\ar[d]^{f^{\text{op}}=F(h)}\\
X(c)\ni& b\ar@{|->}[r]_F&c=F(b)
}
$$
on the other hand, $gf:c \to e$ gives by the
same construction a unique map $k:\hat{q}
\to b$
such that $F(k)=f^{\text{op}}g^{\text{op}}$.
To check that $\hat{q}=q$, observe that
by functoriality of $F$, we have
$F(hj)=F(h)F(j)=f^{\text{op}}g^{\text{op}}=
F(k)$.
By uniqueness of $k$, we conclude that
$k=hj$ and thus $q=\hat{q}$.
\end{proof}

\section{Simplicial sets}
\label{infty-categories-section-simplicial-sets}

\begin{definition}
\label{infty-categories-definition-simplex}
The {\it simplex category} is the category of
ordinals, i.e.
non-empty finite ordered sets
$$
[n]=\{0<1<\ldots<n\}, \quad n=0,1,\ldots
$$
with order preserving maps of sets.
\end{definition}

Objects of the simplicial category $\Delta$
are not called simplices.
Instead, simplices are type of simplicial
set:

\begin{definition}
\label{infty-categories-definition-simplicial-set}
A {\it simplicial set} is a presheaf on
$\Delta$, i.e.
an element of
$\mathsf{Fun}(\Delta^{\text{op}},\text{Sets})
$.
\end{definition}

\begin{definition}
\label{infty-categories-definition-simplex-duplicate-2}
The {\it $n$-simplex} is the
representable simplicial set
$\Hom(-,[n]):=\Delta^n$.
\end{definition}

So far the upshot for me is that simplicial
sets
are like a generalization of a triangulated
topological space. The following construction
shows how to associate to $S \in
\mathsf{sSet}$
a topological space $|S|\in \mathsf{Top}$.
In fact, the weak equivalences in
$\mathsf{sSet}$
can be defined using weak equivalences
in $\mathsf{Top}$ just like in Hatcher.

\medskip\noindent
Here's some copy-paste from \cite[Chapter
1]{rune}
transcribed by ChatGPT:

The category $\Delta$ is generated by
\begin{itemize}
\item the \emph{face maps} $d_i : [n-1]
\hookrightarrow [n]$ that skip $i \in [n]$,
\item the \emph{degeneracy maps} $s_i : [n+1]
\xrightarrow{\text{surj.}} [n]$
that repeat $i \in [n]$,
\end{itemize}
subject to certain relations.

\begin{definition}
\label{infty-categories-definition-topological-n-simplex}
\begin{reference}
\cite[Definition 1.4.2]{rune}
\end{reference}
The {\it topological $n$-simplex}
$|\Delta^n|$ is the topological space
\[
|\Delta^n| := \{(x_0,\ldots,x_n) \in
\mathbb{R}^{n+1} : \sum x_i = 1,\; 0 \leq x_i
\leq 1\}
\]
(with the subspace topology from
$\mathbb{R}^{n+1}$).
For $\varphi:[n]\to[m]$ we can define a
continuous map
$\varphi_* : |\Delta^n| \to |\Delta^m|$ by
$$
\varphi_*(x_0,\ldots,x_n)_i =
\sum_{j:\varphi(j)=i} x_j.
$$
This gives a functor $|\Delta^\bullet| :
\Delta \to \mathsf{Top}$.
\end{definition}

We can then define the \emph{singular
simplicial set functor}
\[
\mathrm{Sing}: \mathbf{Top} \to
\mathbf{Set}_\Delta
\]
as
\[
\mathrm{Sing}(X) =
\mathrm{Hom}_{\mathbf{Top}}(|\Delta^\bullet|,
X).
\]

This has a left adjoint $|-|:
\mathbf{Set}_\Delta \to \mathbf{Top}$, called
the
\emph{geometric realization functor}, which
is the unique colimit-preserving functor
that extends $|\Delta^\bullet|$ via the
Yoneda embedding. More concretely,
we can define $|S|$ for a simplicial set $S$
as the quotient of
$\coprod_n S_n \times |\Delta^n|$ where we
identify $(\sigma,\varphi_*p)$ with
$(\varphi^*\sigma,p)$ for
$\varphi:[n]\to[m]$, $\sigma \in S_n$ and $p
\in |\Delta^m|$.
Informally, we build the topological space
$|S|$ out of simplices according to
the ``blueprint'' $S$.

If we say that a morphism $S \to T$ in
$\mathbf{Set}_\Delta$ is a weak equivalence
if
$|S| \to |T|$ is a weak homotopy equivalence,
then the relative category consisting
of $\mathbf{Set}_\Delta$ with these weak
equivalences describes the same homotopy
theory as that of topological spaces; for
example, the counit map $|\mathrm{Sing}\,X|
\to X$
for a topological space $X$ is always a weak
homotopy equivalence.
We can also describe the weak equivalences of
simplicial sets as homotopy equivalences
(or describe them via homotopy groups) if we
restrict to a class of nice objects,
which we will introduce next.

\begin{exercise}
\label{infty-categories-exercise-simpl-n-simpl-topol-n-simp}
It should be possible to show that
the geometric realization of $\Delta^n$ is in
fact
the topological $n$-simplex in
$\mathbb{R}^n$, right?
\end{exercise}

\noindent
I wonder what is the relationship between
$\mathsf{sSet}$ and $\mathsf{Top}$.
How much information do simplicial sets give
us about topological spaces?
Most likely we will be able to see any
tringulated
space, say, a CW complex as a simplicial set.
And also most likely, there are pathological
topological spaces that cannot be retrieved
by
the geometric realization functor.
But probably we just don't care, and forget
$\mathsf{Top}$
(again).

\medskip\noindent
Now I put some exercises.

\begin{exercise}[Observation 1.4.7]
\label{infty-categories-exercise-simpl-sets-categ-has-inter}
Show that the simplicial sets category
$\mathsf{Set}_\Delta$ has internal
Hom $S^T$ for simplicial sets $S$ and $T$,
given by
$$
(S^T)_n:=\Hom_{\mathsf{Set}_\Delta}
(T\times\Delta^n,S)
$$
\end{exercise}

\begin{proof}
We need to show, that $S^T$ is the internal
Hom in the category
$\mathsf{Set}_\Delta$. Different notations
for the internal Hom
are $\text{Map}(-,-)$,
$\underline{\Hom}(-,-)$.
It must be right adjoint to
the functor $U \times -$.
That is,
$$
\Hom_{\text{Set}_\Delta}(U \times
S,T)\cong\Hom_{\mathsf{Set}_\Delta}(U,S^T)
$$
(I think) I understand the statement
correctly
but I don't understand how to apply
\cite[Theorem 1.1.10 (Kan)]{Cisinsky} nor
\cite[Remark 1.1.11]{Cisinsky}
to prove it.


\end{proof}

\begin{exercise}[1.1]
\label{infty-categories-exercise-relat-equiv-relat}
If $S$ is a Kan complex, then the relation
defining
$\pi_0S$ is an equivalence relation.
\end{exercise}

\begin{proof}
\begin{enumerate}
\item (Reflexivity.) Let $a \in S_0$.
Consider the composition
$$
\xymatrix@R=0.5em{
[0]\ar[r]^{d_0}& [1]\ar[r]^{f_0}& [0]\\
0\ar@{|->}[r]& 1\ar@{|->}[r]&0
}
$$
since this gives the identity we must have
$$
\xymatrix@R=0.5em{
S_0\ar[r]^{S(f_0)}& S_1\ar[r]^{S(d_0)}& [0]\\
a\ar@{|->}[r]& S(f_0)(a)\ar@{|->}[r]&a
}
$$
but we can replace $d_0$ by $d_1$ and we'd
still get the identity,
so that $S(d_1)$ also maps $S(f_0)(a)$ to
$a$. In other words,
for any  $a \in S_0$ the 1-simplex
$S(f_0)(a)$ is the desired one.
\item (Symmetry.) Let $a,b \in S_0$…
\end{enumerate}
\end{proof}

Rather informally, I understand a category
$C$ to be
an {\it enriched category over $D$} if for
any objects
$c,d$ in $C$, $\Hom(c,d)$ is an object of
$D$.
Compositions of morphisms exist and are
associative,
and there is an identity morphism for every
object $c$ in $C$.
(See
\url{https://ncatlab.org/nlab/show/enriched+%
category}
for a formal definition.)

\begin{exercise}[1.2]
\label{infty-categories-exercise-enric-simpl-set-categ-subc}
Show that $\mathsf{Cat}_\Delta$ can be
described as the full
subcategory of
$\text{Fun}(\Delta^{\text{op}},\mathsf{Cat})$
containing the functors whose simplicial sets
of objects
are constant.
\end{exercise}

\begin{remark}
\label{infty-categories-remark-induced-functor-on-objects}
The phrase ``simplicial sets of objects are
constants''
means the following. Consider the functor
$\mathsf{Cat} \to \mathsf{Set}$ that maps a
category to its set
of objects (I suppose we may take
$\mathsf{Set}$ to be a universe),
which induces for every functor in
$\mathsf{Fun}(\Delta^{\text{op}},\mathsf{Cat}
)$
a functor in
$\mathsf{Fun}(\Delta^{\text{op}},\mathsf{Set}
)$.
We mean to say the latter map is constant.
\end{remark}

\begin{proof}
We need to construct a fully faithful functor
$$
F:\mathsf{Cat}_\Delta \to
\mathsf{Fun}(\Delta^{\text{op}},\mathsf{Cat})
$$
whose image is the subcategory of functors
whose simplicial sets of objects are
constant.

To a $\mathsf{Set}_\Delta$-enriched category
$C$
associate the functor $F(C)$ which maps $[n]$
to the category $C_n$, which is defined as
follows.
The objects of $C_n$ are the objects of $C$
for all $n$.
(Notice that once we define the functor
completely,
this property will make it indeed a functor
whose simplicial
sets of objects are constant.)
For $a,b$ in $C$, the morphisms of $C_n$
are $\Hom(a,b)_n$. To a map $f:[n] \to [m]$
in $\Delta^{\text{op}}$,
define $F(C)$ to give the functor of $C_m$ to
$C_n$ that fixes
all objects and maps a map in $\Hom(a,b)_m$
to the induced map $\Hom(a,b)_n$ by the
presheaf
$\Hom(a,b)$.

Now let's define how $F$ acts on morphisms.
(This definition is just what it should be,
but let's go over it.)
Choose two
$\mathsf{Set}_\Delta$-enriched categories
$C,D$
and consider their corresponding functors
$F(C),F(D)$.
Fix a morphism $\varphi \in
\Hom_{\mathsf{Cat}_\Delta}(C,D)$.
Define a morphism (of presheaves of
categories)
$F(\varphi):F(C) \to F(D)$
defined as a collection of maps $F(C)_n \to
F(D)_n$
given on objects by $\varphi$ and on
morphisms
also given by $\varphi$, using that $\varphi$
is a morphism
of $\mathsf{Cat}_\Delta$ to ensure
naturality.

Functoriality of $F$ follows from
functoriality of each $\varphi$ as in the
previous paragraph.

Now let's confirm that $F$ is faithful, that
is, it induces injections
on the Hom sets. Suppose $\varphi,\psi \in
\Hom_{\mathsf{Cat}_\Delta}(C,D)$
are such that $F(\varphi)=F(\psi)$. By
definition of $F(\varphi)$ and
 $F(\psi)$, it is immediate that $\varphi$
 and $\psi$ coincide on objects.
In fact, it is also immediate that they
coincide on morphism and as
simplicial sets by definition.

To prove $F$ is fully faithful we only need
to check surjectivity of
the induced maps in Hom sets. Pick a morphism
of presheaves
of categories, denote it
$F(\varphi)$, between two presheaves of
categories $F(C)$ and  $F(D)$,
both of whose simplicial sets of objects are
constant, namely
two sets $C$ and $D$.
Then we can define two
$\mathsf{Set}_\Delta$-enriched categories,
which we also denote by $C$ and $D$,
by defining their objects to be the sets $C$
and $D$, and their morphisms to be
the collections of all the induced morphisms
by $F(C)$ and $F(D)$ coming from
morphisms of $\Delta^{\text{op}}$.
Then it is immediate that the set $\Hom(C,D)$
is indeed a simplicial set.
Thus $C,D \in
\mathsf{Cat}_{\mathsf{Set}_\Delta}$. Further,
we
can define a morphism $\varphi \in
\Hom_{\mathsf{Cat}_\Delta}(C,D)$
which maps on objects as any of the induced
maps by the morphism of presheaves
of categories we started with (since both of
the simplicial sets of objects
of the corresponding categories are
constant!)
and on morphisms as well (any morphism of $C$
was defined as the
induced map by $F(C)$ coming from a map of
$\Delta^{\text{op}}$).
It is clear that this morphism is mapped to
$F(\varphi)$ under $F$.
\end{proof}

\begin{exercise}[1.3]
\label{infty-categories-exercise-N-is-fully-faithful}
Show that
$N:\mathsf{Cat}\to\mathsf{Set}_\Delta$ is
fully faithful.
\end{exercise}

\begin{proof}
We need to show that for any categories
$A,B$,
$\Hom(A,B)=\text{Fun}(A,B)$ is in
``bijection'' with
$\Hom_{\mathsf{Set}_\Delta}(NA,NB)$.
Recall that
$NA$ is the presheaf that maps $[n]$ to
the set of composible sequence of $n$
morphisms in $A$.
Then to a functor $F:A \to B$ we associate
the
map that sends a sequence of $n$ morphisms in
$A$
to the respective sequence of $n$ morphisms
in $B$
after applying $F$ to each object and map.

Conversely, given a morphism in
$\mathsf{Set}_\Delta$
from $NA$ to $NB$
we can reconstruct a functor from $A$ to $B$
 by interpreting objects of $A$ as $NA_0$ and
maps as $NA_1$.
\end{proof}

\section{infty-groupoids}
\label{infty-categories-section-infty-groupoids}

Apparently the philosophy is that
we will not formally construct
infty-categories
(nor infty-groupoids) but
barely start using them.
So we admit ``facts'' such as
``there are objects called
{\it infty-groupoids}''.
We shall admit that although there are points
(and paths),
we cannot distinguish between points if there
is
a path joining two points. Thus we don't
really have points
but a set $\pi_0X$ of path components.

There are also homotopies between
paths, and homotopies between homotopies, and
so on.

Also, there are {\it maps} or {\it morphisms}
between groupoids,
homotopies between morphisms, homotopies
between homotopies,
and so on. In fact, all those things form an
infty-groupoid
we denote by $\mathsf{Map}(X,Y)$.

We can {\it compose} maps, and there is
an {\it identity morphism} for every
infty-groupoid $X$.
Composition is unital and associative in the
only
way that makes sense (?), which is up to
homotopy.

\begin{definition}
\label{infty-categories-definition-equivalence}
An {\it equivalence} of groupoids is a pair
of maps
that may be composed not be the identities,
but homotopical to the identities.
\end{definition}

Dani: it looks like the main idea
is to care about anything only up to
homotopy.

\medskip\noindent
Sets are groupoids and for any set
$$
\Hom_{\mathsf{Set}}(\pi_0X,S)
\xrightarrow{\simeq
}\mathsf{Map}(X,S)
$$
is an equivalence.

Then the 1-point set ends up being
the terminal infty-groupoid. The empty set
is also an infty-groupoid, and it is
the initial one.

\medskip\noindent
Given morphisms $X \to Z$ and $Y \to Z$,
there exists a pullback square
$$
\xymatrix{
X\times_Z
Y\ar[r]\ar[d]\ar@{}[dr]|-{\lrcorner}&X\ar[d]
\\
Y\ar[r]&Z
}
$$
where ``square'' means not that it is
commutative,
but that there exists an homotopy between
the compositions (commutative up to
homotopy).
Looks like it basically a fibre product
(see Definition
\ref{infty-categories-definition-fibre-product})
up to homotopy.

The {\it fibre} $f^{-1}(b)$ at $b$ of a map
$f:E\to B$ is defined as the pullback
$$
\xymatrix{
f^{-1}(b)\ar[r]\ar[d]\ar@{}[dr]|-{\lrcorner}&
E\ar[d]\\
\{b\}\ar[r]&B.
}
$$

The {\it product} $X\times Y$ of two
infty-groupoids
$X$ and  $Y$ is defined as the pullback
$$
\xymatrix{
X\times
Y\ar[r]\ar[d]\ar@{}[dr]|-{\lrcorner}&X\ar[d]
\\
Y\ar[r]& \ *.
}
$$

\medskip\noindent
Composition of squares is another square,
and the composition of two pullback squares
is another pullback square.

\begin{definition}
\label{infty-categories-definition-path-space}
For points $x,y \in X$, the {\it path space}
$X(x,y)$
is the pullback
$$
\xymatrix{
X(x,y)\ar[r]\ar[d]\ar@{}[dr]|-{\lrcorner}&
\{x\}\ar[d]\\
\{y\}\ar[r]&X.
}
$$
\end{definition}

\begin{exercise}[2.1.1]
\label{infty-categories-exercise-2-1-1}
Assuming that pushouts exist and have the
expected
universal property,
show that $\Omega^n_xX \simeq
\mathsf{Map}_*(S^n,X)$,
where the $n$-th sphere is the pushout
$$
S^n:=* \amalg_{S^{n-1}}*,
$$
and the space of pointed maps is the pullback
$$
\xymatrix{
\mathsf{Map}_*(S^n,X)\ar[r]\ar[d]\ar@{}[dr]|-
{\lrcorner}
&\mathsf{Map}(S^n,X)\ar[d]\\
\{x\}\ar[r]&\mathsf{Map}(*,X).
}
$$
\end{exercise}

\begin{proof}
Since both
$\Omega^1_xX$ and $\mathsf{Map}_*(S^1,X)$
are pullbacks (and pullbacks
are unique up to homotopy),
it's enough to show that they are
the pullback of the same diagram
(up to homotopy).

First notice that $\mathsf{Map}(*,X)\simeq X$
in an obvious way:
we identify a map $* \to X$
with the image of $*$.
To identify $\mathsf{Map}(S^n,X)$
with $\{x\} \simeq \mathsf{Map}(*,X)$
pick a map $* \to X$.
Now consider the universal
property of pushouts:
$$
\xymatrix{
S^{n-1} \ar[r] \ar[d] & \ * \ar[d] \ar[rrdd]
& & \\
\ * \ar[r] \ar[rrrd] & \ S^n=*
\amalg_{S^{n-1}} \, * \ar@{-->}[rrd]^{\exists
!}& & \\
& & & X
}
$$
\end{proof}

\begin{definition}
\label{infty-categories-definition-homotopy-groups}
The {\it $n$-th homotopy group} of a groupoid
is
$\pi_0\Omega^n_x$.
\end{definition}

\noindent
$\pi_1(X,x)$ is a group, and $\pi_n(X,x)$ is
an abelian group for $n>1$.

Homotopy group detect equivalences:
a map between groupoids is an equivalence
if and only if all homotopy groups
are isomorphic (as groups for $n>1$ or
as sets for $n=0$).

\begin{lemma}
\label{infty-categories-lemma-equivalences-satisfy-3-for-2}
Equivalences of infty-groupoids
satisfy the 3-for-2 property:
if  $f$ and $g$ are composable and two
out of $f,g,f \circ g$ are equivalences,
so is the third.
\end{lemma}

\begin{lemma}
\label{infty-categories-lemma-equiv-iff-equiv-path-space}
A morphism $f:X\to Y$ is an equivalence if
and only if
$\pi_0 X \to \pi_0Y$ is surjective
and $X(x,x') \to Y(fx,fx')$ is an equivalence
for all $x,x' \in X$.
\end{lemma}

\begin{lemma}
\label{infty-categories-lemma-equiv-iff-path-space-contr}
For an  infty-groupoid $X$,
the map $X \to \pi_0X$ is an equivalence if
and only if
$X(x,x')$ is either empty of contractible for
all $x,x' \in X$.
\end{lemma}

\begin{lemma}
\label{infty-categories-lemma-group-contr-iff-path-space}
A groupoid $X$ is contracible iff $X(x,x')$
is either empty or
contractible for all $x,x' \in X$.
\end{lemma}

\medskip\noindent
We also state as a fact that
for a map $f:E \to B$, a point $b \in B$ and
a point $e \in f^{-1}(b)$,
there is a long exact sequence of homotopy
groups
$$
\xymatrix{
\cdots\ar[r]&\pi_n(f^{-1}(b),e)\ar[r]&
\pi_n(E,e)\ar[r]&\pi_n(B,b)\ar[r]&
\pi_{n-1}(f^{-1}(b),e)\ar[r]&\cdots
\ar[r]&\pi_0(E)\ar[r]&\pi_0(B),
}
$$
with appropriate interpretation near the end
since
$\pi_0$ is only a set while $\pi_1$ is a
group.

\begin{proposition}
\label{infty-categories-proposition-map-equiv-iff-fiber}
A map $f:E \to B$ is an equivalence if and
only if
all the fibers $f^{-1}(b)$ for $b \in B$ are
contractible.
\end{proposition}

\begin{proof}
This is immediate from homotopy long exact
sequence
and the fact that homotopy groups detect
equivalences.
\end{proof}

\begin{exercise}[2.2]
\label{infty-categories-exercise-2-2}
Use the 5-lemma to show that given a
commutative triangle
$$
\xymatrix{
X \ar[rd]_{p} \ar[rr]_{f} && Y \ar[dl]^{q} \\
& B
}
$$
the morphism $f$ is an equivalence if and
only if
the induced maps on the fibres
$p^{-1}(b)\to q^{-1}(b)$ are equivalences for
all $b \in B$.
\end{exercise}

\begin{proof}
For the converse implication,
$$
\xymatrix@C=.7em{
\cdots
\ar[r]&\pi_{n+1}(B,b)\ar[r]\ar[d]^{\simeq}
&\pi_n(p^{-1}(b),x)\ar[r]\ar[d]^{\simeq}
&\pi_n(X,x)\ar[r]\ar@{.>}[d]^{\simeq }
&\pi_n(B,b)\ar[r]\ar[d]^{\simeq}
&\pi_{n-1}(p^{-1}(b),x)\ar[d]^{\simeq}\ar[r]
&\cdots\\
\cdots\ar[r]
&\pi_{n+1}(B,b)\ar[r]
&\pi_n(q^{-1}(b),f(x))\ar[r]
&\pi_n(Y,f(x))\ar[r]
&\pi_n(B,b)\ar[r]
&\pi_{n-1}(q^{-1}(b),f(x))\ar[r]
&\cdots
}
$$
and for the forward implication just do the
same with
the map $\pi_n(p^{-1}(b),x)\to
\pi_n(q^{-1},f(x))$ in the center.
\end{proof}

\begin{lemma}
\label{infty-categories-lemma-pullb-iff-fiber-maps-equiv}
A commutative square
$$
\xymatrix{
X'\ar[r]\ar[d]_{f'}&Y'\ar[d]^{f}\\
X\ar[r]_{g}&Y
}
$$
is a pullback if and only if for every $x \in
X$,
the induced map on fibres is an equivalence.
\end{lemma}

\begin{proof}
We need to check that $X' \simeq Y'\times_Y
X$.
I think we use the universal property of
pullbacks
to get a map $X' \to Y'\times_YX$.
To check this is an equivalence we use
Exercise \ref{infty-categories-exercise-2-2} as in
$$
\xymatrix{
X' \ar[rd]_{f'} \ar[rr] && Y'\times_YX
\ar[dl]^p \\
& X.
}
$$
We see that the map on fibers is non other
than $f'^{-1}(x) \to f^{-1}(g(x))$… but why?
This ammounts to showing that
$p^{-1}(x)\simeq f^{-1}(g(x))$.
$$
\xymatrix{
f^{-1}(x)\ar[d]\ar@{-->}[rr]
& & p^{-1}(x)\ar[dd]\ar[r]^{\simeq?}
&f^{-1}(g(x))\ar[dd]\\
X' \ar[rrrd] \ar[rrd] \ar[rrdd]_{f'} & & \\
& & X \times_YY' \ar[d]_p \ar[r] & Y'
\ar[d]^f \\
& & X \ar[r]_g & Y
}
$$
\end{proof}

\begin{exercise}[2.3]
\label{infty-categories-exercise-2-3}
Consider a commutative square…

\end{exercise}

\begin{exercise}[2.4]
\label{infty-categories-exercise-2-4}
Suppose we have a commutative diagram
$$
\xymatrix{
X\ar[r]\ar[d]&X'\ar[d]\ar[r]& X''\ar[d]\\
Y\ar[r]&Y'\ar[r]&Y''.
}
$$
\begin{enumerate}
\item If the right and composite squares are
both pullbacks,
then so is the left-hand square.
\item If $\pi_0Y \to \pi_0Y'$ is surjective
and the left and composite squares
are both pullbacks,
then so is the right-hand square.
\end{enumerate}
\end{exercise}

\begin{proof}
\begin{enumerate}
\item To use Corollary 2.1.23 we look at the
fibers and apply
 Corollary 2.1.17: the 3-for-2 property.
\item I'm not sure why is that condition on
$\pi_0Y \to \pi_0Y'$… looks like the same
argument should work.
\end{enumerate}
\end{proof}

\section{Monomorphisms of infty-groupoids}
\label{infty-categories-section-monomorphisms}

\noindent
Monomorphisms can be easily recalled
as the maps whose fibers are either empty or
contractible,
and, equivalently, as the maps that are
injective on $\pi_0$ and the square
$$
\xymatrix{
X\ar[r]\ar[d]&Y\ar[d]\\
\pi_0X\ar[r]&\pi_0Y
}
$$
is a pullback.

In class I picked up that a map is
monomorphism if and only if it's injective on
$\pi_0$
and isomorphism in higher homotopy groups
(after fixing a base point).

Monomorphisms of infty-groupoids are…

\begin{exercise}[2.6]
\label{infty-categories-exercise-2-6}
Show that monomorphisms are closed
under base change.
\end{exercise}

This means that in a pullback square,
if the vertical arrow
is a monomorphism, then so is the
other vertical arrow.

\begin{remark}
\label{infty-categories-remark-monomorphism-equivalence}
Using Exercise \ref{infty-categories-exercise-2-2}, it follows
that $f$ is a monomorphism
if and only if for all $x,x' \in X$, the
induced
morphism on path spaces $X(x,x') \to
Y(fx,fx')$ is
an equivalence.
\end{remark}

\noindent
The following lemma is very similar to
Proposition
\ref{infty-categories-proposition-map-equiv-iff-fiber}.
The difference is that in that proposition we
require all the fibers to be contractible,
and in the following we allow some of them to
be empty.

\begin{lemma}
\label{infty-categories-lemma-monom-iff-fibre-empty-contr}
A morphism of infty-groupoids $f:X \to Y$
is
a monomorphism if and only if the fibres are
all
either empty of contractible.
\end{lemma}

\begin{proof}
Using Lemma
\ref{infty-categories-lemma-pullb-iff-fiber-maps-equiv},
which says a map is a pullback iff induced
maps on fibers are equivalences.
\end{proof}

\begin{remark}
\label{infty-categories-remark-lemma-2-2-3}
There is a caveat in Lemma 2.2.3 on the claim
that the vertical map on the left in the
following
diagram is ``diagonal map''
$$
\xymatrix{
\bullet\ar[r]\ar[d]&X\ar[d]\ar[r]&Y\ar[d]\\
\bullet\ar[r]&X \times X\ar[r]Y \times Y&
}
$$
This should be proved.
\end{remark}

\begin{proposition}
\label{infty-categories-proposition-monom-then-monom-sets}
If $X \to Y$ is a monomorphism, then $\pi_0X
\to \pi_0Y$
is a monomorphism of sets, and the
commutative square
$$
\xymatrix{
X\ar[r]\ar[d]&Y\ar[d]\\
\pi_0\ar[r]&\pi_0Y
}
$$
is a pullback.
\end{proposition}

\begin{exercise}
\label{infty-categories-exercise-2-7}
Let $X$ be an infty-groupoid and consider
a subset $S\subseteq \pi_0X$.
Show that if we form the pullback
$$
\xymatrix{
Y\ar[r]\ar[d]\ar@{}[dr]|-{\lrcorner}&X\ar[d]
\\
S\ar@{^{(}->}[r]&\pi_0X
}
$$
then the induced map $\pi_0Y\to S$ is an
isomorphism.
\end{exercise}

\begin{proof}
Might be kind of nonsense but I thought this:
think of $\pi_0$ as a functor from
infty-groupoids
to $\mathsf{Sets}$. Note that this functor
preserves pullbacks almost tautologically:
in fact, morphisms of groupoids are defined
to be morphisms of the $\pi_0$'s along
with some homotopies and so on. Then
$\pi_0Y$ is the pullback of the square above
after applying the functor $\pi_0$.
But so is $S$ by trivial reasons.
Then by universal property of the pullback
in $\mathsf{Sets}$ there is a unique morphism
$S \to \pi_0Y$ which inverts $\pi_0Y \to S$,
i.e.
$$
\xymatrix{
S \ar[rrrd] \ar@{-->}[rrd] \ar[rrdd] & & \\
& & \pi_0Y \ar[d] \ar[r] & \pi_0X \ar[d] \\
& & S \ar[r] & \pi_0X
}
$$
\end{proof}

\noindent
See remark 2.2.8. Looks like there
is a factorization structure involved.

\begin{exercise}[2.8]
\label{infty-categories-exercise-2-8}
Given a commutative triangle
$$
\xymatrix{
X \ar[rd]_{p} \ar[rr]_{f} && Y \ar[dl]^{q} \\
& B
}
$$
of infty-groupoids, the morphism $f$ is a
monomorphism
if and only if for all $b \in B$, the
induced map on fibres $X_b \to Y_b$ is a
monomorphism.
\end{exercise}

\begin{proof}
By Lemma 2.2.3 it's enough to show that all
the fibres of
$f$ are either empty or contractible.
Let $y \in Y$. The map $X_{q(y)}\to Y_{q(y)}$
induced from $f$
is a monomorphism by hypothesis, so that its
fibers
are empty or contractible. Since
$f(q(y))=p^{-1}(q)$,
the fibers of the fiber map coincide with the
fiber of the
original map $f$.
\end{proof}

\begin{exercise}[2.9]
\label{infty-categories-exercise-2-9}
Show that if $X$ is an infty-groupoid,
then so is $\mathsf{Fun}(\mathcal{C},X)$
for any infty-category $\mathcal{C}$.
\end{exercise}

\begin{proof}
We should use that
infty-categories are Cartesian closed.
\end{proof}

\section{infty-categories}
\label{infty-categories-section-infty-categories}

\noindent
There are objects called {\it
infty-categories}. …

\medskip\noindent
Given an infty-category $\mathcal{C}$,
there exists
a {\it localization} $\|\mathcal{C}\|$ to an
infty-groupoid, with a canonical map
$\mathcal{C} \to \|\mathcal{C}\|$,
(which is not equivalent to $\mathcal{C}$,
but)
such that for any infty-groupoid $X$
the induced map
$$
\mathsf{Map}(\|\mathcal{C}\|,X) \to
\mathsf{Map}(\mathcal{C},X)
$$
is an equivalence.
Moreover  $\|[1]\|\simeq *$.

It turns out that this concept of
localization
generalizes the notion of geometric
realization
for simplices.

\medskip\noindent
We say that $[0]$ and $[1]$ ``detect
equivalences''
in the following sense:
a functor $F: \mathcal{C} \to \mathcal{D}$ is
an equivalence if and only if the maps
$$
\mathcal{C}^{\simeq} \to
\mathcal{D}^{\simeq},\qquad
\mathsf{Map}([1],\mathcal{C})\to\mathsf{Map}
([1],\mathcal{D})
$$
are equivalences of infty-groupoids.
(Because $\mathcal{C}^\simeq \simeq
\mathsf{Map}([0], \mathcal{C})$.)
To understand this you can think that
$\mathsf{Map}([1],\mathcal{C})$ is like
the space of all $\Hom$s, which is a
groupoid,
whose arrows are natural equivalences.

\begin{lemma}
\label{infty-categories-lemma-equiv-infty-group}
The following are equivalent for an
infty-category $\mathcal{C}$:
\begin{enumerate}
\item $\mathcal{C}$ is an infty-groupoid.

\item $\mathcal{C}^\simeq \to \mathcal{C}$ is
an equivalence.

\item The map
$$
\mathcal{C}^\simeq \simeq
\mathsf{Map}([0],\mathcal{C})
\to \mathsf{Map}([1],\mathcal{C})
$$
induced by $[1] \to [0]$ is an equivalence
of infty-groupoids.

\item The functor
$$
\mathcal{C}\to \mathsf{Fun}([1],\mathcal{C})
$$
induced by $[1]\to [0]$ is an equivalence
of infty-categories.
\end{enumerate}
\end{lemma}

\begin{proof}
The first two conditions are equivalent by
definition
of $\mathcal{C}^{\simeq}$. The third
condition is implied
by the second one using the object
$\|\cdot\|$ postulated above
(which applies when $\mathcal{C}$ is a
groupoid).
{\bf Caveat:} I think the fiber of this map
is always a point!

$(3\implies 2)$ actually makes sense:
we can construct a commutative diagram with
left and top arrows
equivalences (why the left one is an
equivalence?).
By hypothesis, the one on the right is also
an equivalence.
Then the one on the bottom is too, and since
$[0]$ and $[1]$
detect equivalences, we are done.

($3 \implies 4$) We need to apply
$(-)^\simeq$ and
$\mathsf{Map}([1],-)=\mathsf{Fun}([1],-)
^\simeq$…
\end{proof}

\section{Segal spaces}
\label{infty-categories-section-segal}


\begin{remark}
\label{infty-categories-remark-Segal-spaces}
Think of a Segal space as satisfying
$$
\mathsf{Map}(\text{Spine}^n,X) \simeq
\mathsf{Map}(\Delta^n,X)
$$
where $\text{Spine}^n$ is a simplicial
construction.
In the case of a tetrahedon, $\text{Spine}$
is defined as the
chain of edges joining the four vertices
0,1,2,3.
\end{remark}


The {\it Kan condition} is
$$
\xymatrix{
\Lambda^n_i\ar@{^{(}->}[d]^i \ar[r]&  X\\
\Delta^n\ar@{.>}[ur]_{\exists }
}
$$
for all $i$,
and is called strict/weak whether the
existence
is unique or not. The {\it inner Kan
condition}
is if we require for $i \neq 0$ and $i\neq
n$.

Why is this so important?

\begin{definition}
\label{infty-categories-definition-nerve}
The {\it nerve} functor $N:\mathsf{Cat} \to
\mathsf{Set}_\Delta$
is defined by
$$
\mathcal{C} \mapsto
\mathsf{Hom}_{\mathsf{Cat}}([\bullet],
\mathcal{C}),
$$
so that $N\mathcal{C}_n$ is the set
of all composable sequences of $n$ morphisms.
\end{definition}

Notice that $N\mathcal{C}_0$ is the set of
all objects of $\mathcal{C}$
and $N\mathcal{C}_1$ is the set of all
morphisms.
This is crucial for the following exercise:

\begin{exercise}
\label{infty-categories-exercise}
$N:\mathsf{Cat}\to\mathsf{Set}_\Delta$ is
fully faithful.
\end{exercise}

\begin{proof}
To a functor $F \in
\mathsf{Hom}(\mathcal{C},\mathcal{D})$
we assign a simplicial set
$N(F)_n:N(\mathcal{C})_n \to
N(\mathcal{D})_n$
defined in the obvious way: we map a sequence
$$
\bullet \xrightarrow{f_1}\bullet\to \cdots
\bullet \xrightarrow{f_n}\bullet
$$
to
$$
\bullet \xrightarrow{F(f_1)}\bullet \to
\cdots
\bullet\xrightarrow{F(f_n)}\bullet
$$
This functor is well defined as a functor
of simplicial sets by functoriality of $F$.
To check fully faithfullness we use the fact
that we can reconstruct a category from
$N(F)_0$ and $N(F)_1$.
\end{proof}

Thus, the Kan condition allows to see
categories
as simplicial sets. Here are possible
generalizations
of this:

\begin{center}
\begin{tabular}{c | c | c}
& strict &weak\\
\hline
all $i$ & groupoids & infty-groupoids\\
inner & categories & infty-categories.
\end{tabular}
\end{center}

\section{Segal condition}
\label{infty-categories-section-segal-condition}

The {\it Segal condition} is

$$
\mathsf{Map}([n],X) \xrightarrow{\simeq}
\mathsf{Map}([1],X)
\times_{\mathsf{Map}([0],X)}\ldots
\times_{\mathsf{Map}([0],X)}\mathsf{Map}([1],
X)
$$

\section{Lifting properties}
\label{infty-categories-section-lifting-properties}

The key to proving anything related
to left orthogonality is not to use
the usual lifting diagram but
the equivalent definition involving
$\mathsf{Map}$.

\begin{definition}
\label{infty-categories-definition-orthogonal}
For morphisms $\ell:A \to B$ and $r:X \to Y$
we say that $r$ is {\it right orthogonal}
to $\ell$ (and dually that $\ell$ is
{\it left orthogonal} to  $r$ if
for any commutative square
$$
\xymatrix{
A\ar[r]\ar[d]_{\ell}&X\ar[d]^{r}\\
B\ar@{.>}[ur]\ar[r]&Y
}
$$
the space of diagonal lifts $B \to X$ is
contractible.
(I think of this as: there exists a unique
lift.)
\end{definition}

\noindent
The key observation is that
a choice of horizontal arrows making the
former
diagram commute is the same as a choice of
element in
$\mathsf{Map}(B,Y)\times_{\mathsf{Map}(A,Y)}
\mathsf{Map}(A,X)$.
Indeed, given a map $A\to X$,
we may poscompose with $r$ to obtain
a map $A \to Y$, and similarly precomposing
with
$\ell$ any map $B \to Y$.
Intuitively, the pullback
of these two maps is the set of pairs
of maps $A \to X$ and $B \to Y$
that give the same map  $A \to Y$,
i.e. that the square commutes.

But a lift $B \to X$ also determines
a map $A \to X$ and
a map $B \to Y$.
All this to say that
a choice of horizontal arrows,
i.e. an element of the pullback
$\mathsf{Map}(B,Y)
\times_{\mathsf{Map}(A,Y)}\mathsf{Map}(A,X)$,
that comes from a choice of map
$B \to X$ is the same as
a fibre of the map
$$
\mathsf{Map}(B,X) \to
\mathsf{Map}(B,Y)
\times_{\mathsf{Map}(A,Y)}\mathsf{Map}(A,X).
$$
Then the definition is equivalent
to this fibre being contractible.
In turn, this is equivalent to the
map being an equivalence.
By putting $\mathsf{Map}(B,X)$ instead
of the pullback, this is also equivalent
to the commutative square
$$
\xymatrix{
\mathsf{Map}(B,X)\ar[r]^{r_*}\ar[d]_{\ell^*}&
\mathsf{Map}(B,Y)\ar[d]^{\ell^*}\\
\mathsf{Map}(A,X)\ar[r]_{r_*}&\mathsf{Map}(A,
Y)
}
$$
being a pullback. So that's the
right way to look at things, apparently.

\begin{example}
\label{infty-categories-example-monom-iff-right-ortho-pusho}
It's the same for a map $f:X \to Y$ to be
a monomorphism and to be
right orthogonal to $* \amalg * \to *$.
To see it, draw the diagram of right
orthogonality
and then put it in the $\mathsf{Map}$
pullback form.
Using that $\mathsf{Map}(* \amalg *,X)\simeq
X\times X$,
which follows from writing the
definition of $* \amalg *$ as the pushout of
the empty space, we obtain the definition of
monomorphism.
\end{example}


\begin{lemma}
\label{infty-categories-lemma-epimo-left-ortho-monom}
Epimorphisms in $\mathit{Gpd}_{\infty}$ are
left orthogonal to
monomorphisms.
\end{lemma}

\begin{exercise}
\label{infty-categories-exercise-2-11}
Show that a map is left orthogonal
to itself if and only if it is an
equivalence.
\end{exercise}

\begin{proof}
We use Proposition
\ref{infty-categories-proposition-map-equiv-iff-fiber}.
Consider the diagram
$$
\xymatrix{
X\ar[r]^{\text{id}}\ar[d]_{f}&X\ar[d]^{f}\\
Y\ar[r]_{\text{id}}\ar@{.>}[ur]&Y.
}
$$
If $f$ is an equivalence then
by the existence of the lift
every fiber consists of only one point and
thus it is contractible.
Conversely, if every fiber is contractible
we define the unique lift as the fiber of
every point.
\end{proof}

\begin{lemma}
\label{infty-categories-lemma-orthogonality-and-composition}
Suppose $f:A \to B$ is left orthogonal
to a map $r:X \to Y$.
Then a map $g:B \to C$ is left
orthogonal to $r$ if and only if $g \circ f$
is.
\end{lemma}

\begin{proof}
By 3-for-2 for pullbacks, i.e. Exercise
\ref{infty-categories-exercise-2-4}.
\end{proof}

It is important to remember that
being a monomorphism is the same
as square with $\pi_0$ is pullback.
That is the key to Lemma 2.4.3.
For one implication we need
to put surjectivity on $\pi_0$
to apply 3-for-2.

Look at the case for sets,
it's also true that epimorphisms
is orthogonal to monomorphisms.

\begin{lemma}
\label{infty-categories-lemma-ortho-pusho-diagr}
Suppose we have a commutative diagram
\begin{equation}
\label{infty-categories-equation-orthogonal-lemma}
\xymatrix{
X\ar[d]^f& Y \ar[l]\ar[d]_g \ar[r]&
Z\ar[d]_h\\
X'&Y' \ar[l]\ar[r]&Z,
}
\end{equation}
such that each of the morphisms $f$, $g$ and
$h$
is left orthogonal to a morphism $r:U \to V$.
Then the induced morphism on pushouts
$X \amalg_Y Z \to X' \amalg_{Y'}Z'$ is
also left orthogonal to $r$.
\end{lemma}

\begin{proof}
The induced map on pushouts is
$$
\xymatrix{
Y \ar[r] \ar[d] & Z \ar[d] \ar[rrd] & & \\
X \ar[r] \ar[rd]& X \amalg_Y Z
\ar@{-->}[rrd]^F & & Z'\ar[d] \\
&X'\ar[rr]& & X' \amalg_{Y'}Z'
}
$$
To construct a lift of the form
$$
\xymatrix{
X \amalg_Y Z\ar[r]\ar[d]&U\ar[d]^r\\
X' \amalg_{Y'}Z'\ar[r]\ar@{.>}[ur]&V
}
$$
we use the universal product of pushout for
$X' \amalg_{Y'}Z'$,
that is, we shall be done once we construct a
digram
$$
\xymatrix{
Y'\ar[r]\ar[d]&Z'\ar[d]\\
X'\ar[r]&U.
}
$$
Notice that this is indeed the case:
the universal property of the pushout
in homotopy sense means that the space of
maps
is contractible.

Such a map is obviously constructed by
composing $Y' \to X'$ as in Diagram
\ref{infty-categories-equation-orthogonal-lemma}
followed by the lift of $f$. Likewise
we construct a map  $Y'\to V$.
And in fact
{\bf it is not needed that $g$ is orthogonal
to $r$}.
\end{proof}

\begin{proof}
$$
\xymatrix{
X\ar[d]&Y \ar[l]\ar[d]\ar[r]&Z\ar[d]\\
X'\ar[d]&  Y \ar[l]\ar[d]\ar[r]&Z\ar[d]\\
X' \ar[d]&  Y \ar[l]\ar[d]\ar[r]& Z'\ar[d]\\
X'\amalg_Y Y'\ar[d]&  Y'\ar[l]\ar[d]\ar[r]&
Z'\ar[d]\\
X'& Y'\ar[l]\ar[r]&Z'
}
$$
In each row of vertical arrows there is one
that
is not trivial. We will be done once we show…
using
the next lemma.

Every triple of horizontal arrows gives a
pushout.
Every triple of vertical arrows gives a
morphism of pushouts.
In fact, we have that for every triple
of horizontal arrows, the corresponding
pushout
morphism is a pushout of the corresponding
nontrivial
vertical arrow, e.g. on the top part we get

$$
\xymatrix{
X\ar[r]\ar[d]\ar@{}[dr]|-{\ulcorner}&X
\amalg_YZ\ar[d]\\
X'\ar[r]&X'\amalg_YZ
}
$$

\end{proof}

\begin{lemma}
\label{infty-categories-lemma-orthogonal-lemma2}
Suppose we have a pushout square
$$
\xymatrix{
A\ar[r]\ar[d]_f\ar@{}[dr]|-{\ulcorner}&
B\ar[d]^{f'}\\
C\ar[r]&B \amalg_A C
}
$$
where $f$ is left orthogonal to a morphism
$r$.
Then $f'$ is also left orthogonal to $r$.
\end{lemma}

\begin{proof}
Suppose there is a square
$$
\xymatrix{
B\ar[r]\ar[d]_{f'}&U\ar[d]^{r}\\
C\ar[r]&B \amalg_A C.
}
$$
To find a unique up to homotopy lift of $f'$,
all we need to do is construct a square
$$
\xymatrix{
A\ar[r]^?\ar[d]&B\ar[d]\\
C\ar[r]_?&U.
}
$$
Which we can construct easily by
$$
\xymatrix{
A \ar[d]_f\ar[r]& B\ar[r]& U\ar[d]^r\\
C\ar@{.>}[urr]\ar[r]&  B \amalg_A C \ar[r]&V.
}
$$
\end{proof}

A more formal approach is:

\begin{proof}
First apply $\mathsf{Map}-,U)$.
\end{proof}

\begin{definition}
\label{infty-categories-definition-retract}
Recall that an object $X$ is a {\it retract}
of $Y$
if there are maps $X \to Y \to X$ and
a homotopy between the composite and the
identity of $X$.
Similarly,
we say that a morphism of $f'$ is
a {\it retract} of $f$ if there is a
commutative diagram
…
\end{definition}

\begin{lemma}
\label{infty-categories-lemma-retract-orthogonality}
Suppose $f'$ is a retract of $f'$.
If $f$ is left orthogonal to a morphism
$r$, then so is $f'$.
\end{lemma}

\begin{proof}
Suppose we have a lifting problem
$$
\xymatrix{
X'\ar[r]\ar[d]&U\ar[d]^r\\
Y'\ar[r]&V.
}
$$
Then we construct
$$
\xymatrix{
X\ar[r]\ar[d]&X'\ar[r]\ar[d]&U\ar[d]^{r}\\
Y\ar[r]\ar@{.>}[urr]&Y'\ar[r]&V
}
$$
thus obtaining a lift by precomposing the
bottom arrow with $Y' \to Y$
given in the definition of retract.
\end{proof}

\begin{exercise}
\label{infty-categories-exercise-retra-equiv-equiv}
Show that any retract of an equivalence is an
equivalence.
\end{exercise}

\begin{proof}
Recall that the definition of
equivalence is that it has inverses such
that the compositions are homotopic to the
identity.

$$
\xymatrix{
\bullet\ar[r]\ar[d]^f&\bullet\ar[d]^{f'}&
\bullet\ar[d]^{f}\\
\bullet\ar[r]&\bullet\ar@{.>}[u]\ar[r]&
\bullet
}
$$
In a more formal way we obtain a cube diagram
upon application of $\mathsf{Map}$.
\end{proof}


\section{Conservative functors and mapping
spaces}
\label{infty-categories-section-conse-funct-mappi-space}

\noindent
The fibers of conservative functors are
groupoids.
Fibrations are conservative (Exercise
\ref{infty-categories-exercise-3-1})
and so are fully faithful functors (Theorem
\ref{infty-categories-theorem-fully-faith-impli-conse})

\begin{definition}
\label{infty-categories-definition-conservative-functor}
A functor $\mathcal{C} \to \mathcal{D}$ is
{\it conservative}
if it is right orthogonal to $s_0:[1]\to
[0]$.
\end{definition}

\begin{exercise}
\label{infty-categories-exercise-equiv-conse-funct}
The following are equivalent for a functor
$F:\mathcal{C} \to \mathcal{D}$:
\begin{enumerate}
\item $F$ is conservative.
\label{infty-categories-item-conservative1}

\item The fibres of $F$ are all
infty-groupoids.
\label{infty-categories-item-conservative2}

\item The commutative square
\label{infty-categories-item-conservative3}
$$
\xymatrix{
\mathcal{C}^{\simeq}\ar[r]\ar[d]&\mathcal{C}
\ar[d]\\
\mathcal{D}^{\simeq}\ar[r]&\mathcal{D}
}
$$
is a pullback.
\end{enumerate}
\end{exercise}

\begin{proof}
Maybe the easiest is to check that
\ref{infty-categories-item-conservative3} $\implies$
\ref{infty-categories-item-conservative2}.
This is just because the square being a
pullback
means that the fibres are equivalent,
and the fibres of the left vertical arrow are
groupoids, while the fibres of the right
vertical
arrow are the fibres of $F$.

In a similar way we can prove
\ref{infty-categories-item-conservative1} $\iff$
\ref{infty-categories-item-conservative2}.
Use that a map is an equivalence if and only
if it is
an equivalence upon groupoidification and
$\mathsf{Map}([1],-)$
on the fibres of diagram in
\ref{infty-categories-item-conservative3}.
Applying groupoidification gives
trivial equivalences in the fibres.
Applying $\mathsf{Map}([1],-)$ and getting
equivalences on the fibres is equivalent to
the square
$$
\xymatrix{
\mathcal{C}\simeq\mathsf{Map}([0],\mathcal{C}
)
\ar[r]\ar[d]&
\mathsf{Map}([1],\mathcal{C})\ar[d]\\
\mathcal{D}\simeq
\mathsf{Map}([0],\mathcal{D})\ar[r]&
\mathsf{Map}([1],\mathcal{D})
}
$$
being a pullback, which is the definition of
conservative.

if and only if the maps on fibres
are equivalences. The fiber of the left
vertical arrow
is the groupoidification of the fiber of
groupoidification of $F$,
and the fiber of the right vertical arrow
is the $\mathsf{Map}([1],-)$.

Also, by applying $\mathsf{Map}([1],-)$ to
the square on \ref{infty-categories-item-conservative3}
we immediately obtain the $\mathsf{Map}$
condition
for being conservative, i.e. condition
\ref{infty-categories-item-conservative1}.

Finally to check that
\ref{infty-categories-item-conservative2} $\implies$
\ref{infty-categories-item-conservative3}
we apply $\mathsf{Map}([0],-)$ to obtain
$$
\xymatrix{
\mathsf{Map}([0],\mathcal{C}^\simeq)\ar[r]
\ar[d]
&\mathsf{Map}([0],\mathcal{C})\ar[d]\\
\mathsf{Map}([0],\mathcal{D}^{\simeq})\ar[r]&
\mathsf{Map}([0],\mathcal{D})
}
$$
To check this is a pullback we only need to
check that the maps induced in fibers are
equivalences
by  Lemma
\ref{infty-categories-lemma-pullb-iff-fiber-maps-equiv}.
But since the fibres of $F$ are
infty-groupoids,
they must equal the fiber of the map on the
associated groupoids
(it's a fact! Fibre and groupoidification
commute).
We do the same for $[1]$ and use that $[0]$
and $[1]$ detect
equivalences.
\end{proof}

Condition \ref{infty-categories-item-conservative3}
can be interpreted as:
if a map gives an equivalence in
$\mathcal{D}$,
then it already was an equivalence in
$\mathcal{C}$.
For that we use
$$
\xymatrix{
\pi_0\mathcal{C}^\simeq\ar[r]\ar[d]\ar@{}[dr]
|-{\lrcorner}&
\pi_0\mathcal{D}^\simeq\ar[d]\\
\pi_0\mathsf{Map}([1],\mathcal{C})\ar[r]&
\pi_0\mathsf{Map}([1],\mathcal{D}).
}
$$

Key tip: by Lemma (?) it is equivalent that
a functor $F: \mathcal{C} \to \mathcal{D}$ is
conservative and
that the square [square with $\pi_0$] is a
pullback.

\begin{definition}
\label{infty-categories-definition-arrow-category}
The {\it arrow category} of $\mathcal{C}$ is
$\mathsf{Ar}(\mathcal{C})=\mathsf{Fun}([1],
\mathcal{C})$.
\end{definition}

\begin{lemma}
\label{infty-categories-lemma-1}
$$
\xymatrix{
B\ar[r]\ar[d]^f\ar@{}[dr]|-{\ulcorner}&B
\amalg_YZ\ar[d]^
{f\amalg_{\text{id}}\text{id}}\\
D\ar[r]&D \amalg_YZ
}
$$
\end{lemma}

\begin{proof}
$$
\xymatrix{
Y\ar[r]\ar[d]\ar@{}[dr]|-{\ulcorner}&Z\ar[d]
\\
B\ar[r]&B\amalg_YZ\ar[d]\\
D\ar[r]&D\amalg_YZ
}
$$
\end{proof}

Now observe that the folowing
is a particular case:
$$
\xymatrix{
[1]\ar[r]\ar[d]^{s_0}\ar@{}[dr]|-{\ulcorner}&
[1]\amalg_{[0]}[1]=[2]\ar[d]
_{\text{id}\amalg_{\text{id}}s_0}\\
[0]\ar[r]&[1]\amalg_{[0]}[0]=[1]
}
$$
where the latter equality,
namely $[1]\amalg_{[0]}[0]=[1]$,
holds because it must be a fact that
when one of the horizontal arrows of a
pushout
is an equivalence then the other one must be
too.

\begin{lemma}
\label{infty-categories-lemma-2}
In
$$
\xymatrix{
A\ar[r]\ar[d]\ar@{}[dr]|-{\ulcorner}&V\ar[d]
\ar[r]&
B\amalg_YZ\ar[d]_{f\amalg_{\text{id}}
\text{id}}\\
C\ar[r]&D\ar[r]&D\amalg_YZ
}
$$
We can put identity left or right and it will
still be a pushout
\end{lemma}

We conclude that
$$
\xymatrix{
[1]\amalg [1]\ar[r]\ar[d]^{s_0
\amalg_{\emptyset}\text{id}}
\ar@{}[dr]|-{\ulcorner}&[2]\amalg_{[1]}[2]
\ar[d]\\
[0]\amalg[1]\ar[d]^{\text{id}\amalg
s_0}\ar[r]&[1]\amalg_{[1]}[2]
\ar[d]_{\text{id} \amalg_{\text{id}}s_0}\\
[0]\amalg[0]\ar[r]&[1]\amalg_{[1]}[1]=[1]
}
$$

Why do all this? A map that is right
orthogonal
to $[1]\to[0]$ is also right orthogonal
to $s_0\times \text{id}:[1]\times[1]\to[1]$.

\begin{proposition}
\label{infty-categories-proposition-chara-conse}
It is equivalent that $F:\mathcal{C} \to
\mathcal{D}$
is conservative and that the square
$$
\xymatrix{
\mathsf{Fun}([0],\mathcal{C})=\mathcal{C}
\ar[r]\ar[d]&
\mathsf{Fun}([0],\mathcal{D})=\mathcal{D}
\ar[d]\\
\mathsf{Fun}([1],\mathcal{C})=\mathsf{Ar}
(\mathcal{C})\ar[r]
&\mathsf{Fun}([1],\mathcal{D})=\mathsf{Ar}
(\mathcal{D})
}
$$
is a pullback.
\end{proposition}

\begin{proof}
One implication is easy: just take
$(-)^\simeq$ to obtain the $\mathsf{Map}$
condition
for conservative functors.

For the other direction we can get to
the exercise diagram by looking at
$$
\xymatrix{
[1]&\emptyset\ar[l]\ar[r]\ar[d]& [1]\ar[d]\\
[0]&\emptyset\ar[l]\ar[r]&\emptyset
}
$$
Then we apply $\mathsf{Map}$ (I think) and
get to
$$
\xymatrix{
\mathsf{Map}([1],\mathcal{C})\ar[r]\ar[d]
\ar@{}[dr]|-{\lrcorner}&
\mathsf{Map}([1],\mathcal{D})\ar[d]\\
\mathsf{Map}([1],\mathsf{Fun}([1],\mathcal{C}
)\ar[r]\ar[d]^{=}&
\mathsf{Map}([1],\mathcal{D})\ar[d]^=\\
\mathsf{Map}([1]\times[1],\mathcal{C})\ar[r]&
\mathsf{Map}([1]\times[1],\mathcal{D})
}
$$
Where have used currying.
The latter diagram gives the pullback
after applying $\mathsf{Map}([1],-)$.
The pullback after applying $(-)^\simeq$
gives just the $\mathsf{Map}$ condition
for conservativity.
\end{proof}

\begin{proposition}
\label{infty-categories-proposition-arrow-funct-conse}
$\mathsf{Ar}(\mathcal{C})\to \mathcal{C}
\times \mathcal{C}$
is conservative.
\end{proposition}

\begin{proof}
We use Proposition
\ref{infty-categories-proposition-chara-conse}
putting
$\mathcal{C}=\mathsf{Ar}(\mathcal{C})$
and
$\mathcal{D}=\mathcal{C} \times \mathcal{C}$.
That is, we know that
$\mathsf{Ar}(\mathcal{C}) \to \mathcal{C}
\times \mathcal{C}$
is conservative if and only if
$$
\xymatrix{
\mathsf{Fun}([1],\mathcal{C})=\mathsf{Ar}
(\mathcal{C})\ar[r]\ar[d]&
\mathsf{Fun}([0]\amalg
[0],\mathcal{C})=\mathcal{C}\times\mathcal{C}
\ar[d]\\
\mathsf{Fun}([1],\mathsf{Fun}([1],\mathcal{C}
))
=\mathsf{Fun}([1]\times[1],\mathcal{C}\ar[r]&
\mathsf{Fun}[1],\mathcal{C}
\times\mathcal{C})
=\mathsf{Fun}([1]\amalg[1],\mathcal{C})
}
$$
where we have used
$$
\begin{aligned}
\mathcal{C}\times\mathcal{C}&=\mathsf{Fun}
([0]\amalg[0],\mathcal{C})\\
\mathsf{Fun}([1],\mathcal{C}\times\mathcal{C}
)
&=\mathsf{Fun}([1]\times([0]\amalg[0]),
\mathcal{C})\\
&=\mathsf{Fun}([1]\amalg[1],\mathcal{C})
\end{aligned}
$$

\noindent
where the first equality is given
just by applying
$\mathsf{Fun}(-,\mathcal{C})$
to the basic pushout diagram,
the second one by currying,
and third one by
using that these are also $1$-categories.

The point is that once
we look at this the right way
we realise that it is nothing but
the exercise diagram. So it's a pullback.
\end{proof}

\begin{definition}
\label{infty-categories-definition-mapping-space}
The {\it mapping space} is the fiber
$$
\xymatrix{
\mathcal{C}(x,y)\ar[r]\ar[d]\ar@{}[dr]|-
{\lrcorner}
&\mathsf{Ar}(\mathcal{C})\ar[d]\\
\{(x,y)\}\ar[r]&\mathcal{C}\times\mathcal{C}
}
$$
\end{definition}

\begin{exercise}
\label{infty-categories-exercise-compositions}
Use the pushout decomposition of $[2]$
to define composition maps
$$
\mathcal{C}(x,y)\times\mathcal{C}(y,z)
\to\mathcal{C}(x,z).
$$
(For extra credit, use the decomposition
of $[3]$ to show this is associative
up to a specified homotopy.)
\end{exercise}

\begin{proof}
$$
\xymatrix{
\text{fiber}\ar[d]\ar@{.>}[rr]^{\simeq}&&
\mathcal{C}(x,y)\times
\mathcal{C}(y,z)\ar[d]\\
\mathsf{Fun}([2],\mathcal{C})\ar[r]^-{\simeq}
\ar[d]&
\mathsf{Fun}([1]\amalg_{[0]}[1],\mathcal{C})
\ar[r]^-{\simeq}\ar[d]
&\mathsf{Fun}([1],\mathcal{C})
\times_{\mathsf{Fun}([0],\mathcal{C})}
\mathsf{Fun}([1],\mathcal{C})\ar[ld]\\
\mathcal{C}\times\mathcal{C}\times\mathcal{C}
\ar[r]^-{\simeq}&
(\mathcal{C} \times \mathcal{C})
\times_{\mathcal{C}}(\mathcal{C}
\times\mathcal{C})
}
$$
We shall define the fiber to be the
composition $\mathcal{C}(x,y,z)$.
This makes sense since
the fiber on the rightmost part
is indeed the product of the fibers,
as is displayed in the diagram.
This is because in general
$$
\xymatrix{
fib(F)\times fib(G)\ar[r]&
\mathcal{A}\times_Z \mathcal{B}
\ar[r]^{F\times_{\text{id}}G}&
\mathcal{C}\times_Z \mathcal{D}
}
$$
\end{proof}


\section{Fully faithful functors}
\label{infty-categories-section-fully-faithful-functors}

\begin{definition}
\label{infty-categories-definition-fully-faithful}
A functor of infty-categories
$F:\mathcal{C} \to \mathcal{D}$
is {\it fully faithful} if it is
right orthogonal to $\partial[1] \to [1]$,
that is, if the commutative square
$$
\xymatrix{
\mathsf{Map}([1],\mathcal{C})\ar[r]\ar[d]&
\mathsf{Map}([1],\mathcal{D})\ar[d]\\
(\mathcal{C}^\simeq)^{\times 2}\ar[r]&
(\mathcal{D}^\simeq)^{\times 2}
}
$$
is a pullback.
\end{definition}

By Exercise \ref{infty-categories-exercise-equiv-conse-funct},
since
$\mathsf{Ar}(\mathcal{C})\to
\mathcal{C}\times\mathcal{C}$
is conservative, we get that the fibres
are groupoids, so we know that we get
equivalences
on the fibres when we get a pullback, so
really we have
equivalences
$$
\mathcal{C}(x,y)\xrightarrow{\simeq}
\mathcal{D}(Fx,Fy)
$$
\begin{proposition}
\label{infty-categories-proposition-fully-faith-iff-pullb}
$F:\mathcal{C} \to \mathcal{D}$ is fully
faithful
if and only if the commutative square
$$
\xymatrix{
\mathsf{Ar}(\mathcal{C})\ar[r]\ar[d]&
\mathsf{Ar}(\mathcal{D})\ar[d]\\
\mathcal{C}\times\mathcal{C}\ar[r]&
\mathcal{D}\times\mathcal{D}
}
$$
\end{proposition}

\begin{proof}
At least one implication is easy:
if it's a pullback,
maps on fibres are equivalences.
By the remark on the human-readble
definition of fully faithfulness, we are
done.
\end{proof}

\begin{lemma}
\label{infty-categories-lemma-fully-faith-prese-fun}
If $F:\mathcal{C} \to \mathcal{D}$ is fully
faithful,
so is
$F_*:
\mathsf{Fun}(\mathcal{A},\mathcal{C})
\to\mathsf{Fun}(\mathcal{A},\mathcal{D})$
for any infty-category $\mathcal{A}$.
\end{lemma}

\begin{exercise}
\label{infty-categories-exercise-chara-pullb-fully-faith}
Show that the following are equivalent for a
commutative square
$$
\xymatrix{
\mathcal{C}\ar[r]^F\ar[d]&\mathcal{C}'\ar[d]
\\
\mathcal{D}\ar[r]_G&\mathcal{D}'
}
$$
where $F$ and $G$ are fully faithful:
\begin{enumerate}
\item The square is a pullback of
infty-categories.
\label{infty-categories-item-fully-faithful1}

\item The square gives a pullback of
infty-groupoids on cores.
\label{infty-categories-item-fully-faithful2}

\item The square gives a pullback of sets on
$\pi_0(-)^\simeq$.
\label{infty-categories-item-fully-faithful3}
\end{enumerate}
\end{exercise}

\begin{proof}
\ref{infty-categories-item-fully-faithful1} $\implies$
\ref{infty-categories-item-fully-faithful2}
is immediate since pullbacks are detected on
maps from $[0]$
(and $[1]$).

For \ref{infty-categories-item-fully-faithful2} $\implies$
\ref{infty-categories-item-fully-faithful1}
first apply $\mathsf{Map}([1],-)$, then we
get

For \ref{infty-categories-item-fully-faithful2} $\implies$
\ref{infty-categories-item-fully-faithful3}
we first need to make the diagram in the
exercise
after applying $(-)^{\simeq}$, then
the arrows of $F$ and $G$ become
monomorphisms.
If we can prove that, then
we apply some past exercise to see that
the diagram is pullback iff the diagram on
$\pi_0$.
In fact, to prove those maps are
monomorphisms
we will use the theorem proved next,
Theorem
\ref{infty-categories-theorem-fully-faith-impli-conse}.
\end{proof}

\begin{theorem}
\label{infty-categories-theorem-fully-faith-impli-conse}
$F$ is fully faithfull if and only if
$F$ is conservative.
\end{theorem}

\begin{proof}
We want to prove that $F$
is right orthogonal to $\partial[1]\to [1]$
then it is right orthogonal to $[1]\to [0]$.
First we apply some sort of 3-for-2
property (some past exercise)
we look at
$$
\xymatrix{
\partial[1]=[0]\amalg[0]\ar[r]
^{\substack{\text{left
ort.} \\ \text{to $F$}}}
& [1]\ar[r]^{\substack{\text{want to} \\
\text{prove r.o.}}}
&[0]
}
$$
we realise that it's enough
to prove right-orthogonality of the
composition of the two maps above.
$$
\xymatrix{
[1]\amalg[1]\ar[r]\ar[d]\ar@{}[dr]|-
{\ulcorner}&[3]\ar[d]\\
[0]\amalg[0]\ar[r]&E\simeq \ast
}
$$
where $E \simeq *$ by a fact.

Consider
$$
K=\{0<1<2\}\amalg_{\{1<2\}}\{1<2<3\}
$$
which is nothing more than
$[2]\amalg_{[1]}[2]$.
Putting $[2]\simeq[1]\amalg_{[0]}[1]$
we see that
$$
K\simeq[1]\amalg_{[0]}[1]\amalg_{[1]}[1]
\amalg_{[0]}[1]
\simeq[1]\amalg_{[0]}[1]\amalg_{[0]}[1]
\simeq[3]
$$
Now we do
$$
\xymatrix{
\{0<2\}\amalg\{1<3\}K\ar[r]^{\simeq}\ar[d]&
[3]\ar[d]\\
\ast \amalg \ast\ar[r]&E' \ar[r]&E
}
$$
where we already saw the outer square is
a pushouts by definition of $E$,
and $E'$ is defined as the pushout of
whatever it should be.
Then, since $K \simeq [3]$,
we apply the property of pushouts
that if the map on the ``total spaces''
is equivalence then so is the
map on ``base spaces''
(this is opposite to the property in
pullback,
which is just what we call
``pullback of equivalences is equivalence'',
and then apply $\mathsf{Fun}$
to get to the property for pushouts).

Now back to our objective,
by the 3-for-2-ish for left-orthogonality,
now we aim to prove that
$$
\xymatrix{
\{1<2\}\ar[r]& K\ar[r]&E'\ar[r]&E.
}
$$
Now let $H$ be the pushout
$$
\xymatrix{
\{0<2\} \ar[r]\ar[d]\ar@{}[dr]|-{\ulcorner}&
\ast \ar[d]\\
[2] \ar[r]& H…
}
$$
\end{proof}

Recall from Remark
\ref{infty-categories-remark-monomorphism-equiavlence}
that a map of infty-groupoids $f:X \to Y$
is a monomorphism if and only if
the induced map on mapping spaces
$X(x,x') \to Y(fx,fx')$ is an equivalence.
Which is strange: that's what we wanted for
{\it fully faithful} functors, not
monomorphisms.
(But then again: that was before we defined
infty-categories.)

\begin{lemma}
\label{infty-categories-lemma-full-faith-impli-under-group}
If $F: \mathcal{C} \to \mathcal{D}$ is fully
faithful,
then its underlying morphism of
infty-groupoids is a monomorphism
of spaces.
\end{lemma}

\begin{proof}
Fully faithful says $\partial[1] \to [1]$ is
left orthogonal to $F$. But by Theorem
\ref{infty-categories}
we also know it is conservative, so it's
right orthogonal
to $[1] \to [0]$. By 3-for-2, $\partial[1]
\to [1] \to [0]$
is left orthogonal to $F$. But we know that
this is equivalent
to being an equivalence on underlying
groupoids
\end{proof}

\begin{definition}
\label{infty-categories-definition-essentially-surjective}
$F:\mathcal{C} \to \mathcal{D}$ is {\it
essentially surjective} if
$\pi_0F: \pi_0\mathcal{C}^\simeq \to
\pi_0\mathcal{D}^\simeq$ is surjective.
\end{definition}

\begin{remark}
\label{infty-categories-remark-equiv-cores-impli-fully-fait}
Suppose $F:\mathcal{C} \to \mathcal{D}$ gives
an equivalence
on cores. Let's show that $F$ is fully
faithful
if and only if $\mathsf{Map}([1],\mathcal{C})
\to \mathsf{Map}([1],\mathcal{D})$
is an equivalence. Indeed: by definition of
full faithfullness
we get equivalences on all mapping spaces
$\mathcal{C}(x,x') \simeq
\mathcal{D}(fx,fx')$.
Applying $(-)^\simeq$ we obtain a similar
diagram on groupoids,
namely
$$
\xymatrix{
\mathsf{Map}([1],\mathcal{C})\ar[r]\ar[d]&
\mathsf{Map}([1],\mathcal{D})\ar[d]\\
\mathcal{C}^\simeq \times
\mathcal{C}^\simeq\ar[r]&
\mathcal{D}^\simeq\times \mathcal{D}^\simeq.
}
$$
The lower arrow is an equivalence by
hypothesis
and by Exercise \ref{infty-categories-exercise-2-3} so is the
upper arrow.

In conclusion, provided $F$ is an equivalence
on the cores,
it is an equivalence of infty-categories
if and only if it
is fully faithful.
\end{remark}

\begin{lemma}
\label{infty-categories-lemma-equiv-iff-fully-faith-essen}
A functor of infty-categories is an
equivalence
if and only if it is fully faithful and
essentially surjective.
\end{lemma}

\begin{proof}
Apply Lemma
\ref{infty-categories-lemma-full-faith-impli-under-group}
and then Remark
\ref{infty-categories-remark-equiv-cores-impli-fully-fait}.
\end{proof}

\section{Full subcategories}
\label{infty-categories-section-full-subcactegories}

\noindent
This is the construction of a full
subcategory
from a subgroupoid of the underlying groupoid
of the given category.

Given an infty-category $\mathcal{C}$
and a monomorphism $i:X \hookrightarrow
\mathcal{C}^\simeq$ of infty-groupoids,
then there exist an infty-category
$i^*\mathcal{C}$
and a functor of infty-categories
$\overline{i}:i^*\mathcal{C} \to \mathcal{C}$
such that $(\overline{i})^\simeq \simeq i$
(and $(i^*\mathcal{C})\simeq X$ should also
hold)
and such that the following diagram
$$
\xymatrix{
\mathsf{Map}(\mathcal{D},i^*\mathcal{C})
\ar[r]\ar[d]&
\mathsf{Map}(\mathcal{D},\mathcal{C})\ar[d]\\
\mathsf{Map}(\mathcal{D}^\simeq,X)\ar[r]&
\mathsf{Map}(\mathcal{D}^\simeq,\mathcal{C}
^\simeq).
}
$$
is a pullback.

Note that putting $\mathcal{D}=[1]$
we get that $\mathcal{D}^\simeq=\partial[1]$,
meaning that $\overline{i}$ is fully
faithful.

In the other direction, it's also true that
any fully faithful
functor is of that form:

\begin{lemma}
\label{infty-categories-lemma-fully-faithful-subcategories}
If $j:\mathcal{D} \to \mathcal{C}$ is fully
faithful
then $j \simeq \overline{j^\simeq}$.
\end{lemma}

\begin{proof}
By Lemma
\ref{infty-categories-lemma-full-faith-impli-under-group}
we know that fully faithful gives
monomorphism of underlying groupoids,
which means that we can apply the
construction above.
Now, using the diagram above we can factor
$$
\xymatrix{
\mathcal{D}\ar[r]^f\ar@/_1em/[rr]_j
&(j^\simeq)^*\mathcal{C}
\ar[r]^{\overline{i}}
&\mathcal{C}
}
$$
(Essentially by definition of pullback.)
Further, by 3-for-2 in for fully faithful
morphisms we may
conclude that $f$ is an equivalence.
\end{proof}

\begin{lemma}
\label{infty-categories-lemma-essen-surje-ortho}
Essentially surjective functors are left
orthogonal to fully faithful ones.
\end{lemma}

\medskip\noindent
We are almost done constructing a
factorization system!
Consider the functor  $F:\mathcal{C} \to
\mathcal{D}$ such that
$F:\mathcal{C}^\simeq X \hookrightarrow
\mathcal{D}^\simeq.$
This gives

\section{Equivalences in an
infinity-category}
\label{infty-categories-section-equivalences}

For any infty-category we always have
$\mathcal{C} \to \mathsf{Ar}(\mathcal{C})$
induced
from $[1] \to [0]$ by
$\mathsf{Fun}(-,\mathcal{C})$.
This is a fully faithful functor (picture).

Then we apply that
$$
\xymatrix{
\mathcal{C}\ar[r]
^{\substack{\text{essentially}
\\ \text{surjective}}}&
i^*\mathcal{D}\ar[r]^{\substack{\text{fully}
\\
\text{faithful}}} & ?
}
$$
for the case
$$
C\xrightarrow{\simeq}\mathsf{Ar}_{\text{eq}}
(\mathcal{C})
\hookrightarrow \mathsf{Ar}(\mathcal{C}).
$$
Recall
$$
\xymatrix{
\text{Spaces with a $G$-bundle} \ar[rd]
\ar[rr] && \text{Spaces} \\
& \text{Spaces admitting a $G$-bundle}\ar[ur]
}
$$
So th
Then we get
$$
\text{triangle diagram}.
$$
\begin{lemma}
\label{infty-categories-lemma-equivalence-equivalences}
Suppose $f: x \to y$ is a morphism in an
infty-category.
Then the following are equivalent:

\begin{enumerate}
\item $f$ is an equivalence.
\label{infty-categories-item-equiv1}
\item
\label{infty-categories-item-equiv2}
For all $c \in \mathcal{C}$, the morphism
$f_*\mathcal{C}(c,x) \to \mathcal{C}(c,y)$ is
an equivalence of infty-groupoids.
\item
\label{infty-categories-item-equiv3}
For all $c \in \mathcal{C}$, the morphism
$f^*\mathcal{C}(y,c)\to \mathcal{C}(x,c)$
is an equivalence of infty-groupoids.
\end{enumerate}
\end{lemma}

\begin{proof}
2 $\implies$ 1. Put $c=x$ and then $c=y$. We
obtain that
$\mathcal{C}(x,x) \simeq \mathcal{C}(x,y)$
and $\mathcal{C}(y,x) \simeq
\mathcal{C}(y,y)$.

{\bf Exercise.} Check that $- *$ is
``functorial''.
Consider
$$
\xymatrix{
C(x,y)\ar[r]\ar[d]&\mathcal{C}(c,y)\times
\mathcal{C}(x,y)\ar[d]&
\mathcal{C}(c,x)\ar[d]\\
pt\ar[r]_{f}&\mathcal{C}(x,y)\ar[r]&pt
}
$$
where the three squares are pullbacks.
\end{proof}

\begin{lemma}
\label{infty-categories-lemma-group-iff-every-morph-equiv}
An infty-category $\mathcal{C}$ is an
infty-groupoid if
and only if every morphism in $\mathcal{C}$
is an aquicalence.
\end{lemma}

\begin{lemma}
\label{infty-categories-lemma}
A morphism $f \in
\mathsf{Fun}(\mathcal{C},\mathcal{D})$ is an
equivalence
if and only if its
\end{lemma}

\begin{exercise}
\label{infty-categories-exercise-2-18}
The statement is equivalent to the fact that
essentially surjective is left orthogonal to
conservative! [Picture]
\end{exercise}

\section{Monomorphisms of infinity
categories}
\label{infty-categories-section-monom-infty-categ}

\begin{exercise}
\label{infty-categories-exercise-2-19}
Show that the following are equivalent for
$F:\mathcal{C} \to \mathcal{D}$:
\begin{enumerate}
\item $F$ is a monomorphism.
\item $F([i]):\mathcal{C}([i]) \to
\mathcal{C}([i])$
is a monomorphism of infty-groupoids for
$i=0,1$.
\item $F$ is right orthogonal to
$[i]\amalg[i]\to[i]$ for $i=0,1$.
\end{enumerate}
\end{exercise}

\begin{proof}
For $(1) \implies (2)$ just apply definition
that pullbacks are detected on maps from
$[0]$ and $[1]$.

Now $(2) \implies (3)$ is definition of
orthogonality on Maps.
\end{proof}

\begin{remark}
\label{infty-categories-remark-only-on-1}
By Observation 2.10.3, it really only matters
what happens on $[1]$.
\end{remark}

\begin{exercise}
\label{infty-categories-exercise-2-20}
Use exercise dsadas
\end{exercise}

\begin{proof}
Recall we already proved that
$$
\xymatrix{
E \ar[rd] \ar[rr] && E' \ar[dl] \\
& B
}
$$
then $E \to W'$ is mono iff $F_b \to F'_b$ is
mono for all $b \in B$.

Then we have a square with bases
$\mathcal{C}[0] \times \mathcal{C}[0]$
and $\mathcal{D}[0]\times \mathcal{D}[0]$.
But those are different bases!
Then we can construct the pullback
$$
\xymatrix{
pt\ar[r]\ar[d]\ar@{}[dr]|-{\lrcorner}&
\mathcal{C}[0]\times\mathcal{C}[0]\ar[d]\\
pt\ar[r]&\mathcal{D}[0]\times\mathcal{D}[0],
}
$$
indeed, it's a pullback {\it when it exists},
since
monomorphisms have fibers that are either
empty of points,
but in this case, it cannot be a point
(why?!).
\end{proof}

\noindent
Now some intuition on what is a monomorphism
in 1-categories:
if $f\in \mathcal{C}$ is an isomorphism, then
$F(f)$ is an isomorphism too.

\begin{lemma}
\label{infty-categories-lemma-monomorphisms-are-conservative}
Monomorphisms are conservative.
\end{lemma}

\begin{proof}
consider
$$
[1] \amalg [1] \to [1] \perp F \implies  [1]
\to [0] \perp F.
$$
We want to argue that $[1] \to [0]$ is a
retract of
$[0] \amalg [1]\to [0]$.

$$
\xymatrix{
[1]\ar[d]\ar[r]^f&[0]\ar[d]\\
[1] \amalg [1]\ar[r]^{f \amalg
\text{id}}\ar[d]^{\text{id} \amalg
\text{id}}&
[0] \amalg[1]\ar[d]^{\text{left. o. to }F}\\
[1]\ar[r]& [0]
}
$$
Then the upper square is a pushout.
The outer square it's all identities, so it's
also a pushout.
This gives that the two vertical lower arrows
are orthogonal to each other.
That's equivalent to the diagram in
Observation 2.10.3.
\end{proof}

\begin{lemma}
\label{infty-categories-lemma-fully-faithful-are-conservative}
Fully faithful functors are conservative.
\end{lemma}

\begin{proof}
Use that monomorphism iff monomorphism on
groupoids
and on mapping spaces.
\end{proof}

\section{Subcategories}
\label{infty-categories-section-subcategories}

\noindent
A similar construction to that of full
subcategories
yields the notion of subcategory.

\section{Left and right fibrations}
\label{infty-categories-section-left-and-right-fibrations}

\begin{definition}
\label{infty-categories-definition-left-right-fibrations}
A functor of infty-categories
$p:\mathcal{E} \to \mathcal{B}$
is a {\it left fibration} if it is right
orthogonal to the
inclusion $\{0\}\to [1]$, and a {\it right
fibration}
if it is orthogonal to $\{1\}\to [1]$.
Equivalently, $p$ is a left or right
fibration if
the corresponding commutative square
$$
\xymatrix{
\mathsf{Map}([1],\mathcal{E})\ar[r]\ar[d]&
\mathsf{Map}([1],\mathcal{B}\ar[d]\\
\mathcal{E}^\simeq\ar[r]&\mathcal{B}^\simeq
}
$$
is a pullback.
\end{definition}

\begin{exercise}
\label{infty-categories-exercise-3-1}
Every left/right fibration is conservative.
\end{exercise}

\begin{proof}
Put orthogonal diagram in $\mathsf{Map}$
form,
but it doesn't match exactly.
The same happens when we look at
Exercise \ref{infty-categories-exercise-equiv-conse-funct}.
By the way: shouldn't we use $\mathsf{Fun}$
instead of $\mathsf{Map}$
when we work with infty-categories instead
of infty-groupoids?
[This was solved in class, there must be a
picture.]
\end{proof}

\begin{lemma}
\label{infty-categories-lemma-fibra-etale-space-group}
If $\mathcal{B}$ is an infty-groupoid,
then a functor $p:\mathcal{E} \to
\mathcal{B}$
is a left/right fibration if and only if
$\mathcal{E}$
is an infty-groupoid. In particular,
any morphism of infty-groupoids
is both a left and right fibration.
\end{lemma}

\begin{proof}
The proof looks good: I think the vertical
maps
are induced from maps $[0] \to [1]$ and $[1]
\to [0]$.
But how is the former defined? We can just
take
any of the two maps $[0]\to[1]$:
each of these concerns the left/right
fibration case.
We see that either choice implies that
$\mathcal{E}$ is
an infty-groupoid, and that $\mathcal{E}$
being an
infty-groupoid implies both.
\end{proof}

\noindent
Next we have a characterization:

\begin{proposition}
\label{infty-categories-proposition-fibra-chara}
$p$ is a left/right fibration if and only if
the commutative square
$$
\xymatrix{
\mathsf{Ar}(\mathcal{E})\ar[r]\ar[d]&
\mathsf{Ar}(\mathcal{B})\ar[d]\\
\mathcal{E}\ar[r]&\mathcal{B}
}
$$
is a pullback.
\end{proposition}

\begin{proof}
This is complicated.
\end{proof}

\noindent
But it gives a nice corollary:

\begin{lemma}
\label{infty-categories-lemma-induced-map-fibration}
If $p:\mathcal{E}\to\mathcal{B}$ is a
left/right fibration,
then so is $p_*:
\mathsf{Fun}(\mathcal{C},\mathcal{E})
\to\mathsf{Fun}(\mathcal{C},\mathcal{B})$
for any category $\mathcal{C}$.
\end{lemma}

\begin{proof}
It follows easily from the characterization.
\end{proof}

\noindent
Now we characterize fibrations that are also
equivalences:

\begin{proposition}
\label{infty-categories-proposition-fibration-equivalence}
The following are equivalent for a left (or
right)
fibration $p:\mathcal{E} \to \mathcal{B}$:

\begin{enumerate}
\item $p$ is an equivalence.
\item The map on cores
$p^\simeq:\mathcal{E}^\simeq\to\mathcal{B}
^\simeq$
is an equivalence of infty-groupoids.
\item The fibre $\mathcal{E}_b$ is
contractible for every $b \in \mathcal{B}$.
\end{enumerate}
\end{proposition}

\begin{proof}
Easy, uses previous exercises.
\end{proof}

\begin{remark}[3-for-2]
\label{infty-categories-remark-3-for-2}
Given a commutative triangle
$$
\xymatrix{
\mathcal{E} \ar[rd]_{p} \ar[rr]_{F} &&
\mathcal{F} \ar[dl]^{q} \\
& \mathcal{B}
}
$$
where $p$ and $q$ are both left (or both
right)
fibrations, then so if $F$.
This follows easily by 3-for-2 for pullback
squares.
\end{remark}

\noindent
The following proposition looks very similar
to
Proposition
\ref{infty-categories-proposition-fibration-equivalence}.
The only difference is on the third
condition:
this time we look at the fibers of the base
and total space over $\mathcal{B}$
(last time we asked that the fibers
of the fibration itself).

\begin{lemma}
\label{infty-categories-lemma-fibration-equivalence2}
Suppose given a commutative triangle
$$
\xymatrix{
\mathcal{E} \ar[rd]_{p} \ar[rr]_{F} &&
\mathcal{F} \ar[dl]^{q} \\
& \mathcal{B}
}
$$
where $p$ and $q$ are both left (or both
right) fibrations.
Then the following are equivalent.
\begin{enumerate}
\item $F$ is an equivalence.
\item The underlying map of
infty-groupoids
$\mathcal{E}^\simeq\to\mathcal{F}^\simeq$ is
an equivalence.
\item The functor on fibres
$\mathcal{E}_b\to\mathcal{F}_b$
is an equivalence for all $b \in
\mathcal{B}$.
\end{enumerate}
\end{lemma}

\begin{proof}
The first two conditions are given by
Proposition
\ref{infty-categories-proposition-fibration-equivalence}
and Remark \ref{infty-categories-remark-3-for-2}.
For the last one we apply Exercise
\ref{infty-categories-exercise-2-2}.
\end{proof}

\begin{exercise}
\label{infty-categories-exercise-pullb-iff-fiber-maps-equiv}
Show that a commutative square
$$
\xymatrix{
\mathcal{E}\ar[r]\ar[d]_{p}&\mathcal{F}\ar[d]
^{q}\\
\mathcal{B}\ar[r]_{F}&\mathcal{C}
}
$$
\end{exercise}

\begin{proof}
I think this would be proved
essentially as Lemma
\ref{infty-categories-lemma-pullb-iff-fiber-maps-equiv},
using the recently proved Lemma
\ref{infty-categories-lemma-fibration-equivalence2}.
(Notice that the fact that the fibres are
groupoids
follows from Exercise \ref{infty-categories-exercise-3-1} and
Exercise
\ref{infty-categories-exercise-equiv-conse-funct}.
It's true, with a single nuance: that we need
to prove that the pullback $p$
is a right/left fibration when $q$ is.
To prove this we need to use a cube diagram,
see picture.
\end{proof}



\section{Slice categories}
\label{infty-categories-section-slice}

Let $\mathcal{C}$ be a 1-category and $x \in
\Ob \mathcal{C}$.
Then the {\it slice category} $\mathcal{C}/x$
is the 1-category whose objects are morphisms
$y \to x$ in $\mathcal{C}$, and morphisms are
$$
\xymatrix{
y\ar[dd]\ar[dr]\\
& x\\
y'\ar[ur]
}
$$

It turns out that the slice category is the
pullback
$$
\xymatrix{
\mathcal{C}/x\ar[r]\ar[d]\ar@{}[dr]|-
{\lrcorner}&Ar(\mathcal{C})\ar[d]^{ev_1}\\
\{x\}\ar[r]&\mathcal{C}
}
$$
We claim that the map $ev_1$ is a fibration.
(Shown in class.)

\begin{definition}
\label{infty-categories-definition-cone}
For an infty-category $\mathcal{K}$,
the {\it left cone}
$\mathcal{K}^{\triangleleft}$
and {\it right cone}
$\mathcal{K}^{\triangleright}$ on
$\mathcal{K}$
are defined by the pushouts
$$
\begin{aligned}
\mathcal{K}^{\triangleleft}&
=\{-\infty\}\amalg_{\mathcal{K}\times\{0\}}
\mathcal{K}\times[1]\\
\mathcal{K}^{\triangleright}&
=\mathcal{K}\times[1]\amalg_{\mathcal{K}
\times\{1\}}\{\infty\}.
\end{aligned}
$$
The {\it cone point} is the object $\pm
\infty$.
\end{definition}

It can be shown (or maybe we take as a fact)
that cores and mapping spaces in the left
cone $\mathcal{K}^{\triangleleft}$
are given by
\begin{align*}
(\mathcal{K}^{\triangleleft})^\simeq&
\simeq \{-\infty\}\amalg\mathcal{K}^\simeq,\\
\mathcal{K}^{\triangleleft}(x,y)&=
\begin{cases}
\emptyset\qquad &x \in \mathcal{K}, y \simeq
-\infty \\
\ast,\qquad &x \simeq -\infty,\\
\mathcal{K}(x,y),\qquad &x,y \in \mathcal{K}.
\end{cases}
\end{align*}

\noindent
It looks like: there are no maps from any
object to $-\infty$,
there is a single map coming from $-\infty$
to any object
(i.e., the space of maps from $-\infty$ to
any object is contractible).

\begin{lemma}
\label{infty-categories-lemma-pullback-cone}
For infty-categories $\mathcal{K}$ and
$\mathcal{C}$
and an object $x$ in $\mathcal{C}$,
there are natural equivalences
$$
\begin{aligned}
\mathsf{Fun}(\mathcal{K},\mathcal{C}_{/x})&
\simeq
\mathsf{Fun}(\mathcal{K}^{\triangleright},
\mathcal{C})
\times_{\mathsf{Fun}(\{\infty\},\mathcal{C})}
\{x\},\\
\mathsf{Fun}(\mathcal{K},\mathcal{C}_{x/}&
\simeq
\mathsf{Fun}(\mathcal{K}^{\triangleleft},
\mathcal{C})
\times_{\mathsf{Fun}(\{-\infty\},\mathcal{C})
}\{x\}.
\end{aligned}
$$
\end{lemma}

\noindent
I think of this as describing the fiber of
$x$ from the
projection $\mathcal{K}^{\triangleright} \to
\{\infty\}$:
$$
\xymatrix{
\bullet\ar[r]\ar[d]\ar@{}[dr]|-{\lrcorner}
&\mathsf{Fun}(\mathcal{K}^{\triangleright},
\mathcal{C})\ar[d]\\
\{x\}\ar[r]&\mathsf{Fun}(\{\infty\},
\mathcal{C}),
}
$$
i.e., the space of maps from infty to $x$.

\medskip\noindent
Looking back at the definition of cone
with $\mathcal{K}=[1]$,
if we can show that
$$
\xymatrix{
\mathcal{K}\times\{1\}\simeq[1]\times\{1\}
\ar[r]\ar[d]
&\mathcal{K}\times[1]=[1]\times[1]\ar[d]\\
\{\infty\}\simeq[0]\ar[r]&[2]
}
$$
is a pushout, we can conclude that
$[1]^{\triangleright}\simeq[2]$.
Consider the following diagram,
where we want to check that the double lower
square is a pushout:
$$
\xymatrix{
&[1]\ar[r]\ar[d]^{\substack{0\mapsto 0 \\
1\mapsto 2}}
&[2]\ar[d]\\
[1]\ar[d]\ar[r]^{\substack{0 \mapsto 1 \\ 1
\mapsto 2}}
&[2]\ar[r]\ar[d]^{\substack{0\mapsto 0 \\ 1
\mapsto 1\\2\mapsto 1}}
&[2]\amalg_{[1]}[2]\simeq[1]\times[1]\ar[d]\\
[0]\ar[r]_{0 \mapsto
1}&[1]\ar[r]_-{\substack{0 \mapsto 0 \\ 1
\mapsto 2}}&[2]\simeq[1]\amalg_{[0]}[1].
}
$$
First observe that the lower left square is a
pushout.
This is by 2-for-3 since the upper right
square is a pushout
by pushout decomposition of $[1]\times[1]$
and because the two composite vertical arrows
are the identity.
To conclude it's enough to show that the
lower left square
is also a pushout.

\medskip\noindent
Here's a summary of what else is in the
chapter.
Proposition 3.3.7 on two pullback squares.
The exercise that the left and right cones of
$[n]$ are both $[n+1]$.
Corollaruy: $C_{/c}\to \mathcal{C}$ is right
fibration
if and only if $C_{c/}\to \mathcal{C}$ is
left fibration.
Corollary: for a morphism  $f:x \to y$ in an
infty-category $\mathcal{C}$,
we have
$(C_{/y})_{/f}\xrightarrow{\sim}\mathcal{C}
_{/x}$
and
$(\mathcal{C}_{x/})_{f/}\xrightarrow{\sim}
\mathcal{C}_{y/}$.

\begin{lemma}
\label{infty-categories-lemma-3-3-9}
For a morphism $f:x \to y$ in an
infty-category
we have equivalences $(C_{/y})_{f/}/\simeq
\mathcal{C}{/x} $.
\end{lemma}


\medskip\noindent
This one I like:

\begin{lemma}
\label{infty-categories-lemma-left-fibration-iff-induced}
A functor $p:\mathcal{E} \to \mathcal{B}$ is
a
right fibration of and only if
$\mathcal{E}_{/k}\to\mathcal{B}_{/p(x)}$
is an equivalence for all $x$, and a
left fibration
if and only if for every $x \in \mathcal{E}$,
the induced functor $\mathcal{E}_{x/}\to
\mathcal{B}_{p(x)/}$
is an equivalence.
\end{lemma}

\begin{proof}
Did in class. Maybe this is already somewhere
else in this document?
\end{proof}






\section{Straightening for left and right
fibrations}
\label{infty-categories-section-straightening}

\noindent
Suppose $p:\mathcal{E} \to \mathcal{B}$ is a
left fibration.
``Then $p$ is supposed to determine a functor
from $\mathcal{B}\text{}$ to
infty-groupoids.''

First we note that a morphism $f:b \to b'$ in
$\mathcal{B}$
we get maps on fibers $f_!:\mathcal{E}_b \to
\mathcal{E}_{b'}$.
$$
\xymatrix{
\mathcal{E}[0]\ar[d]&\mathcal{E}[1]\ar[l]
^s\ar[d]\ar[r]^t&
\mathcal{E}[0]\ar[d]\\
\mathcal{B}[0]&\mathcal{B}[1]\ar[l]_s\ar[r]
_t&\mathcal{B}[0]
}
$$
where the left square is a pullback by
definition of
$p$ being a left fibration, which in turn
gives
equivalences on maps on fibers, so that we
may invert the
left morphism in
$$
\xymatrix{
\mathcal{E}_b & \mathcal{E}[1]_f\ar[l]^\simeq
\ar[r]& \mathcal{E}_{b'}.
}
$$
Further, we see that $(\text{id}_b)_!$ is
homotopic to the identity of $\mathcal{E}_b$
and that $(gf)_! \simeq g_!f_!$.
Dually, a right fibration gives the
basic data for a contravariant functor from
$\mathcal{B}^{\text{op}}$ to
infty-groupoids.

\medskip\noindent
We introduce the fact that there is an
infty-category $\mathsf{Cat}_{\infty}$
of infty-categories. We denote
$\mathsf{Gpd}_{\infty}$
the full subcategory spanned by the
infty-groupoids.
For an infty-category $\mathcal{B}$,
we denote $\mathsf{LFib}(\mathcal{B})$ and
$\mathsf{RFib}(\mathcal{B})$
the subcategories of
$\mathsf{Cat}_{\infty/\mathcal{B}}$
of left and right fibrations.

\begin{theorem}[Lurie]
\label{infty-categories-theorem-lurie}
For an infty-category $\mathcal{B}$, there
is an equivalence
$$
\text{Str}_{\mathcal{B}}^L:\mathsf{LFib}
(\mathcal{B})
\overset{\sim}{\to}\mathsf{Fun}(\mathcal{B},
\mathsf{Gpd}_{\infty})
$$
called the {\it straightning equivalence} for
left fibrations.
This is moreover natural in $\mathcal{B}$
with respect
to precomposition of functors and pullback of
fibrations.
\end{theorem}

\section{Cartesian and cocartesian
fibrations}
\label{infty-categories-section-cartesian-fibrations}

\noindent
The idea here is to generalize the notion of
right fibration
to some object whose fibers are not
necessarily infty-groupoids
(recall that a fibration is conservative,
which implies its fibers
are infty-groupoids).

\begin{exercise}
\label{infty-categories-exercise-3-5}
Given a functor  $p: \mathcal{E} \to
\mathcal{B}$ and a morphism
$f: x \to y$ in $\mathcal{E}$, show that for
an object
$z \in \mathcal{E}$ we get commutative
squares
$$
\xymatrix{
\mathcal{E}(y,z)\ar[r]^{f^*}\ar[d]&
\mathcal{E}(x,z)\ar[d]\\
\mathcal{B}(py,pz)\ar[r]_{p(f)^*}&\mathcal{B}
(px,pz)
}
\qquad
\xymatrix{
\mathcal{E}(z,x)\ar[r]^{f_*}\ar[d]&
\mathcal{E}(z,y)\ar[d]\\
\mathcal{B}(pz,px)\ar[r]_{p(f)_*}&\mathcal{B}
(pz,py)
}
$$
by taking appropriate fibers in a cube with
top face
$$
\xymatrix{
\mathcal{E}[2]\ar[r]^{d_1}\ar[d]&\mathcal{E}
[1]\ar[d]\\
\mathcal{B}[2]\ar[r]_{d_1}&\mathcal{B}[1].
}
$$
Similarly we have the diagrams
$$
\xymatrix{
\mathcal{E}_{y/}\ar[r]^{f^*}\ar[d]&
\mathcal{E}_{x/}\ar[d]\\
\mathcal{B}_{py/}\ar[r]_{p(f)^*}&\mathcal{B}
_{px/}
}\qquad
\xymatrix{
\mathcal{E}_{/x}\ar[r]^{f_*}\ar[d]&
\mathcal{E}_{/y}\ar[d]\\
\mathcal{B}_{/px}\ar[r]_{p(f)_*}&\mathcal{B}
_{/py}.
}
$$
\end{exercise}

\begin{proof}
We use the proof of Lemma 3.3.9
to construct the following complicated
diagram:

$$
\xymatrix{
&C_{/y}\ar[r]&\mathsf{Ar}(\mathcal{C})\ar[r]
^s&\mathcal{C}\\
(\mathcal{C}_{/y})_{/f}\ar@{-->}[ur]\ar[r]
\ar[d]
&\mathsf{Ar}(\mathcal{C}_{/y})\ar[d]^t\ar[r]
\ar[u]^s&
\mathsf{Fun}([2],\mathcal{C})\ar[d]\ar[r]
&\mathsf{Fun}([1],\mathcal{C})=\mathsf{Ar}
(\mathcal{C})\ar[d]^t\ar[u]^s\\
pt\ar[r]&\mathcal{C}_{/y}\ar[d]\ar[r]
&\mathsf{Fun}([1],\mathcal{C})\ar[r]^s
\ar[d]&\mathcal{C}\\
&pt\ar[r]&\mathcal{C}
}
$$

This then leads to
$$
\xymatrix{
\mathcal{E}_{/x}\ar[rd]\ar[rr]^{\text{right
fib.}}\ar[ddd]
& & \mathcal{E}_{/y}\ar[ddd]\ar[dl]\\
& \mathcal{E}\ar[d]\\
&\mathcal{B}\\
\mathcal{B}_{px}\ar[ur]\ar[rr]_{\text{right
fib.}}&&
\mathcal{B}_{py}\ar[lu]
}
$$
and in fact the commutativity of the first
diagram
in the exercise is given by taking the fibres
of this last diagram by the diagonal arrows.
\end{proof}
\noindent
From this we can define $p$-cartesian as in
1-categories:

\begin{definition}
\label{infty-categories-definition-p-cocartesian}
Given a functor $p:\mathcal{E} \to
\mathcal{B}$,
we say a morphism $f:x\to y$
in $\mathcal{E}$ is {\it $p$-cocartesian}
if for every object $z \in \mathcal{E}$, the
commutative square
$$
\xymatrix{
\mathcal{E}(y,z)\ar[r]^{f^*}\ar[d]&
\mathcal{E}(x,y)\ar[d]\\
B(py,pz)\ar[r]_{p(f)^*}&\mathcal{B}(px,pz)
}
$$
is a pullback. Dually, we say that $f$ is
{\it $p$-cartesian}
if all squares
$$
\xymatrix{
\mathcal{E}(z,x)\ar[r]^{f_*}\ar[d]&
\mathcal{E}(z,y)\ar[d]\\
\mathcal{B}(pz,px)\ar[r]_{p(f)_*}&\mathcal{B}
(pz,py)
}
$$
are pullbacks.
\end{definition}

\noindent
This is basically the definition we had for
cartesian (and cocartesian, by inverting the
arrows)
for 1-categories. Namely, for cocartesian,
that given the solid diagram
$$
\xymatrix{
&z\ar@{-->}[d]^{\exists !}\\
x\ar[ur]^g\ar[r]_f&y
}
\qquad \qquad
\xymatrix{
&p(z)\ar[d]^\phi\\
p(x)\ar[ur]^{p(g)}\ar[r]_{p(f)}&p(y)
}
$$
there is a unique dashed arrow.

\begin{exercise}
\label{infty-categories-exercise-3-6}
A morphism is $p$-cocartesian if and only if
$$
\xymatrix{
\mathcal{E}_{y/}\ar[r]^{f^*}\ar[d]&
\mathcal{E}_{x/}\ar[d]\\
\mathcal{B}_{py/}\ar[r]_{p(f)^*}&\mathcal{B}
_{px/}
}
$$
is a pullback.
\end{exercise}

\begin{proof}
We have constructed a diagram of two cubes on
top of each other.
In fact, this can be thought of as shifts in
derived category.
\end{proof}

\begin{lemma}
\label{infty-categories-lemma-3-5-3}
This looks like a 2-for-3 for cocartesian
maps.
\end{lemma}

\begin{definition}
\label{infty-categories-definition-cocartesian-lift}
Given a morphism $f:a \to b$ in $\mathcal{B}$
and an object $x \in \mathcal{E}$
\end{definition}

\begin{lemma}
\label{infty-categories-lemma-3-5-5}
Lifts of $p$-cartesian if and only if …
equivalence.
\end{lemma}

\begin{lemma}
\label{infty-categories-lemma-3-5-6}
Complicated monomorphism statement.
\end{lemma}

\begin{proof}
We start by presenting a general situation.
We start with a $2\times 2$ square with
homotopy
fibers. We claim that it is well-defined, and
proved it.

Now we imagine another situation:
a cube with the bottom and top faces that are
pullbacks.
Then we take homotopy fibre and produce a
cube on top of
the original cube. We claim that the top face
is also a pullback. We take homotopy fibres
producing
a cube to one of the sides, then take
homotopy
fibre of the homotopy fibres,
which gives an equivalence since the maps
induced
on homotopy fibres are equivalences.

As for another equivalence for fibration,
we have that the diagram
$$
\xymatrix{
\mathsf{Ar}_{\text{cart}}(\mathcal{E})\ar[r]
\ar[d]_{\text{ev}_1}&
\mathsf{Ar}(\mathcal{B})\ar[d]^{\text{ev}_1}
\\
\mathcal{E}\ar[r]&\mathcal{B}
}
$$
is a pullback. That is, that
$$
\mathsf{Ar}_{\text{cart}}(\mathcal{E})
\simeq
\mathcal{E}\times_{\mathcal{B}}\mathsf{Ar}
(\mathcal{B}).
$$
This can also be represented as
$$
\xymatrix@C=2em@R=2em{
\mathsf{Ar}_{\text{cart}}(\mathcal{E})
\ar@{^{(}->}[r]
&\mathsf{Ar}(\mathcal{E}) \ar[rrrd] \ar[rrd]
\ar[rrdd]_{\text{ev}_1} & & \\
& & &
\mathcal{E}\times_{\mathcal{B}}\mathsf{Ar}
(\mathcal{B})
\ar[d] \ar[r]
& \mathsf{Ar}(\mathcal{B}) \ar[d] \\
& & & \mathcal{E} \ar[r] & \mathcal{B}
}
$$
Applying $\mathsf{Map}([1],-)$ we obtain
$$
\xymatrix@C=2em@R=2em{
\mathcal{E}[1]_{\text{cart}}\ar@/_{1pc}/
[rrrd]\ar@{^{(}->}[r]
&\mathcal{E}[1] \ar[rrrd] \ar[rrd] \ar[rrdd]
& & \\
& & &
\mathcal{E}^{\simeq}\times_{\mathcal{B}
^{\simeq}}\mathcal{B}[1]
 \ar[d] \ar[r]
& \mathcal{B}[1] \ar[d] \\
& & & \mathcal{E}^{\simeq} \ar[r] &
\mathcal{B}^{\simeq}
}
$$
where the arrow from $\mathcal{E}[1]$ to the
pullback
comes from the universal property of
pullback.


\end{proof}

\begin{proposition}
\label{infty-categories-proposition-left-fibra-cocar}
Left fibrations are $p$-cocartesian and right
fibrations are $p$-cartesian.
\end{proposition}

\noindent
The following lemma shows that for
cocartesian fibrations
right fibrations are equivalent
fibrations in infty-groupoids.

\begin{lemma}
\label{infty-categories-lemma-cocar-fibra-equiv}
Suppose $p:\mathcal{E} \to \mathcal{B}$ a
cocartesian fibration.
Then the following are equivalent:
\begin{enumerate}
\item $p$ is a left fibration (which is the
same as saying it's
a fibration in infty-groupoids).
\label{infty-categories-item-cocartesian1}

\item Every morphism in $\mathcal{E}$ is
$p$-cocartesian.
\label{infty-categories-item-cocartesian2}

\item $p$ is conservative.
\label{infty-categories-item-cocartesian3}

\item The fibers of $p$ are
infty-groupoids.
\label{infty-categories-item-cocartesian4}
\end{enumerate}
\end{lemma}

\begin{proof}
We already proved (\ref{infty-categories-item-cocartesian1})
implies (\ref{infty-categories-item-cocartesian2})
(?). For (\ref{infty-categories-item-cocartesian2})implies
(\ref{infty-categories-item-cocartesian1}) we consider
$$
\xymatrix{
\mathcal{E}[1]_{\text{cartesian}}
\ar[rd]_{\simeq} \ar[rr]^{\simeq} &&
\mathcal{B}[1]\times_{\mathcal{B}}\mathcal{E}
^{\simeq}
\\
& \mathcal{E}[1]\ar[ur]_{\simeq}
}
$$
We already know that
(\ref{infty-categories-item-cocartesian3}) is equivalent to
(\ref{infty-categories-item-cocartesian4}).

To prove that (\ref{infty-categories-item-3}) implies
(\ref{infty-categories-item-2})
we use the property that conservative
functors
lift equivalences to equivalences (?).


\end{proof}

\begin{proposition}
\label{infty-categories-proposition-mophi-cocar-fibra-equiv}
If
$$
\xymatrix{
\mathcal{E} \ar[rd]_{p} \ar[rr]_{F} &&
\mathcal{F} \ar[dl]^{q} \\
& \mathcal{B}
}
$$
is a morphism of cocartesian fibrations over
$\mathcal{B}$
then $F$ is an equivalence if and only if the
induced
map on fibres is an equivalence for all $b
\in \mathcal{B}$.
\end{proposition}

\begin{proposition}
\label{infty-categories-proposition-3-5-15}
Pullback of fibration is fibration.
\end{proposition}

\begin{exercise}
\label{infty-categories-exercise-pullb-iff-equiv-fiber-inft}
Suppose we have a morphism of (co-)cartesian
fibrations $F$
$$
\xymatrix{
\mathcal{E}\ar[r]^{F}\ar[d]_{p}&\mathcal{E}
'\ar[d]^{p'}\\
\mathcal{B}\ar[r]_{G}&\mathcal{B}'.
}
$$
\end{exercise}

\begin{proof}
Factor the above square as
$$
\xymatrix{
\mathcal{E}\ar[r]\ar[d]&\mathcal{E}
'\times_{\mathcal{B}'}\mathcal{B}\ar[d]\ar[r]
&\mathcal{E}'\ar[d]\\
\mathcal{B}\ar[r]_{\text{id}}&\mathcal{B}
\ar[r]&\mathcal{B}'
}
$$
and apply the last two propositions.
\end{proof}

\begin{exercise}
\label{infty-categories-exercise-pullback-of-three-fibrations}
Pullback of three cocartesian fibrations…
\end{exercise}






\section{Arrow infty-categories and
fibration}
\label{infty-categories-section-arrow-infty-categories}

\begin{proposition}
\label{infty-categories-proposition-pullb-arrow-categ}
The mapping space of two morphisms
$f:x \to y$ and $g:x'\to y'$ in an
infty-category is the pullback
$$
\xymatrix{
\mathsf{Ar}(\mathcal{C})(f,g)\ar[r]\ar[d]
\ar@{}[dr]|-{\lrcorner}
&\mathcal{C}(x,x')\ar[d]\\
\mathcal{C}(y,y')\ar[r]_{f^*}&
\mathcal{C}(x,y').
}
$$
\end{proposition}

\noindent
Upshot: morphisms of morphisms are
a map of sources and map of targets such that
composing with $f^*$ or $g_*$ accordingly
give the same map.

More precisely,
the idea is as in Section
\ref{infty-categories-section-lifting-properties},
where we defined lifting as a pullback.
In the current situation: a choice of
a map $x \to x'$ and a map $y \to y'$
such that $x \to x' \xrightarrow{g_*}y'$
coincides with
$x \xrightarrow{f} y \to y'$,
is just an element of
$\mathcal{C}(x,x')\times_{\mathcal{C}(x,y')}
\mathcal{C}(y,y')$.
The proposition then says that morphisms from
$f$ to $g$
are the same as elements in such pullback.

\begin{proof}
Rune starts by saying that we want to find
the fibre at $(f,g)$ of the map
$$
\mathsf{Fun}([1],\mathsf{Fun}([1],\mathcal{C}
))=
\mathsf{Map}([1]\times[1],\mathcal{C})
\to\mathcal{C}[1]^{\times
2}.
$$
Why do we want that? Because that's just the
definition
of the mapping space $\mathsf{Ar}(x,y)$
(see Definition
\ref{infty-categories-definition-mapping-space}).

To do it we notice that by the fact that in a
general situation
$$
\xymatrix{
\bullet\ar@{-->}[d]\ar[r]&\mathcal{E}_b\ar[d]
\ar[rr]&&
\mathcal{F}_b\ar[d]\\
\bullet \ar[r]&\mathcal{E} \ar[rd] \ar[rr] &&
\mathcal{F} \ar[dl] \\
&& \mathcal{B}
}
$$
the dashed arrow is
an equivalence (because the square is a
pullback
by 2-for-3:
$$
\xymatrix{
\mathcal{E}_\ell\ar[r]\ar[d]&\mathcal{F}
_\ell\ar[d]\ar[r]&pt\ar[d]\\
\mathcal{E}\ar[r]&\mathcal{F}\ar[r]&
\mathcal{B}),
}
$$
so it's enough to find the fibre of the fibre
map.

Since we want to compute $\mathsf{Ar}(x,y)$,
which is by definition the fibre of
$\mathsf{Ar}(\mathcal{C})\times\mathsf{Ar}
(\mathcal{C})
\to \mathcal{C}\times\mathcal{C}$, a
conservative map,
which means that its fibres are groupoids,
then the following diagram should be enough:
$$
\xymatrix{
\bullet\ar@{-->}[d]\ar[r]
&(\mathcal{C}(x,y)\times\mathcal{C}(y,y'))
\times
_{\mathbb{C}(x,y')}
(\mathcal{C}(x,x'),\mathcal{C}(z,y'))
\ar[d]\ar[rr]
&& \mathcal{C}(x,y)\times
\mathcal{C}(x',y')\ar[d]\\
\bullet \ar[r]
&\mathsf{Map}([1]\times[1],\mathcal{C})
=\mathsf{Ar}(\mathsf{Ar}(\mathcal{C}))^\simeq
\ar[rd] \ar[rr]
&&\mathcal{C}[1]^{\times 2}\ar[dl] \\
&& \mathcal{B}
}
$$
We find that the fibre we are seeking is
\begin{equation}
\label{infty-categories-equation-fibre}
\{f\}\times\mathcal{C}(y,y')
\times_{\mathcal{C}(x,y')}
(\mathcal{C}(x,x')\times \{g\}).
\end{equation}

\noindent
The reason is given by the three diagrams
$$
\xymatrix{
\mathcal{C}(x,y')\ar[r]
&\mathcal{C}(x,y')\ar[r]
&pt\\
\{f\}\times\mathcal{C}(y,y')\ar[r]
&\mathcal{C}(x,y)\times\mathcal{C}(y,y')
\ar[r]
&\mathcal{C}(x,y)\ni f\\
\mathcal{C}(x,x')\times\{y\}\ar[r]
&\mathcal{C}(x,x')\times\mathcal{C}(x',y')
\ar[r]
&\mathcal{C}(x',y')\ni g.
}
$$
\end{proof}

\noindent
Next we observe that morphisms in
$\mathsf{Ar}(\mathcal{C})$
are given by commutative squares
\begin{equation}
\label{infty-categories-equation-morphisms-in-arrow}
\xymatrix{
y\ar[r]^{e}\ar[d]_{f}&y'\ar[d]^{g}\\
x\ar[r]_{h}&x'
}
\end{equation}

\noindent
which in fact are elements of the fibre we
found
in the last proposition, namely Equation
\ref{infty-categories-equation-fibre}.
This just says that ``the mapping spaces in
$\mathsf{Ar}(\mathcal{C})$
are what we expected them to be''!
(namely those commutative squares).

\begin{exercise}
\label{infty-categories-exercise-morphism-of-morphisms}
Extend this description to show that for a
morphism
$\alpha:f \to g$ in
$\mathsf{Ar}(\mathcal{C})$,
given by a commutative square
$$
\xymatrix{
x\ar[r]^{e}\ar[d]_{f}&z\ar[d]^{g}\\
y\ar[r]_{h}&w,
}
$$
precomposition with $\alpha$ fits in a
commutative cube
(where the left and right faces are
pullbacks)
$$
\xymatrix@C=1em@R=1em{
\mathsf{Ar}(\mathcal{C})(g,q)
\ar[rr]^{\alpha^*} \ar[dd] \ar[rd] & &
\mathsf{Ar}(\mathcal{C})(f,q) \ar'[d][dd]
\ar[rd] \\
 & \mathcal{C}(z,u) \ar[rr]^(0.3){e^*}
 \ar[dd] & & \mathcal{C}(x,u) \ar[dd] \\
\mathcal{C}(w,v) \ar'[r][rr]^{h^*} \ar[rd] &
& \mathcal{C}(y,v) \ar[rd] \\
 & \mathcal{C}(z,v) \ar[rr]^{e^*} & &
 \mathcal{C}(x,v)
}
$$
for $q:u \to v$.
\end{exercise}

\begin{proof}
By Proposition
\ref{infty-categories-proposition-pullb-arrow-categ}
we know that the left and right faces are
pullbacks
(and in particular comute).
The front and back squares are
commutative by associativity of composition
(which in fact is not immediate to prove, but
holds).
The hardest part are the top and bottom
faces.
They commute by a careful application of
Exercise \ref{infty-categories-exercise-3-5}.
\end{proof}

\noindent
The next proposition looks like a technical
statement
justifying Exercise
\ref{infty-categories-exercise-morphism-of-morphisms},
and to be used for proving the next lemma.

\begin{proposition}
\label{infty-categories-proposition-evaluation-and-arrow}
For any infty-category $\mathcal{C}$,
the functor
$$
\text{ev}_1:\mathsf{Ar}(\mathcal{C})
\to\mathcal{C}
$$
is a cocartesian fibration.
A morphism in $\mathsf{Ar}(\mathcal{C})$
given by a square
\ref{infty-categories-equation-morphisms-in-arrow}
is $\text{ev}_1$ cocartesian if and only if
$e$
is an equivalence.
\end{proposition}

\begin{lemma}
\label{infty-categories-lemma-p-cocartesian-characterization}
Given a functor
$p:\mathcal{E}\to\mathcal{B}$,
the following are equivalent for a morphism
$\overline{f}:x\to y$ in $\mathcal{E}$
over $f:a \to b$ in $\mathcal{B}$:
\begin{enumerate}
\item The morphism $\overline{f}$ is
$p$-cocartesian.
\item For every morphism  …
\end{enumerate}
\end{lemma}

\noindent
The next proposition describes a situation
in which we have a subcategory $S$ of arrows
and we construct a functor $p:\mathcal{E} \to
\mathcal{B}$
such that $S$ is the $p$-cocartesian maps.

\begin{proposition}
\label{infty-categories-proposition-subcategory-of-arrows}
For a functor $p:\mathcal{E} \to \mathcal{B}$
and a monomorphism
$S\hookrightarrow \mathcal{E}[1]$, let
$\mathsf{Ar}_S(\mathcal{E})\subseteq
\mathsf{Ar}(\mathcal{E})$
denote the full subcategory spanned by
morphisms in $S$.
Then the following are equivalent:
\begin{enumerate}
\item $p$ is a cocartesian fibration and  $S$
is the
infty-groupoid of $p$-cocartesian
morphisms.
\item The commutative square
$$
\xymatrix{
\mathsf{Ar}_S(\mathcal{E})\ar[r]^{\text{ev}
_0}\ar[d]&\mathcal{E}\ar[d]^{p}\\
\mathsf{Ar}(\mathcal{B})\ar[r]_{\text{ev}_0}&
\mathcal{B}
}
$$
is a pullback.
\end{enumerate}
\end{proposition}

We obtain a corollary that seems to establish
that $\mathsf{Fun}(\mathcal{K},-)$ preserves
cocartesian functors.

\begin{lemma}
\label{infty-categories-lemma-covar-fun-prese-cocar}
Suppose $p:\mathcal{E}\to \mathcal{B}$ is a
cocartesian fibration.
Then for any infty-category $\mathcal{K}$,
the functor
$$
p_*:\mathsf{Fun}(\mathcal{K},\mathcal{E})
\to\mathsf{Fun}(\mathcal{K},\mathcal{B})
$$
given by composition with $p$ is also a
cocartesian fibration.
A natural transformation
$\phi:\mathcal{K}\times[1]\to\mathcal{E}$
is $p_*$-cocartesian if and only if its
components
$\phi_a:[1]\to\mathcal{E}$are $p$-cocartesian
for all $a \in \mathcal{K}$.
\end{lemma}

\noindent
Then there's the part on the cocartesian
transport functor.
I think it goes like this: we use Proposition
\ref{infty-categories-proposition-subcategory-of-arrows} to
identify
$\mathcal{E}\times_{\mathcal{B}}\mathsf{Ar}
(\mathcal{B})
\simeq
\mathsf{Ar}_{\text{coct}}(\mathcal{E})$.
We have learnt to read elements in the
pullback
as pairs in the product that project nicely;
namely,
elements of
$\mathcal{E}\times_{\mathcal{B}}\mathsf{Ar}
(\mathcal{B})$
are nothing but pairs $(x,f:p(x)\to b)$
according to the diagram
$$
\xymatrix{
\mathcal{E}\times_{\mathcal{B}}\mathsf{Ar}
(\mathcal{B})
\ar[r]\ar[d]\ar@{}[dr]|-{\lrcorner}&
\mathsf{Ar}(\mathcal{B})\ar[d]^{s}\\
\mathcal{E}\ar[r]&\mathcal{B}.
}
$$
Take any such pair, and consider a
cocartesian lift.
This gives a map on $\mathcal{E}$ with source
$x$.
Denote the target by $f_!x$.
This map,
$$
\mathcal{E}\times_{\mathcal{B}}\mathsf{Ar}
(\mathcal{B})
\xrightarrow{\sim}\mathsf{Ar}_{\text{coct}}
(\mathcal{E})
\xrightarrow{\text{ev}_1}\mathcal{E},
$$
is called the {\it cocartesian transport
functor}.

\medskip\noindent
The upshot about this part of the book
is that we defined what the induced maps on
mapping spaces
associated to some straightening functor
that we didn't define yet should be.

More explicitly, we want to construct
$$
\text{CoCart}(\mathcal{B})
\xrightarrow{\simeq}
\mathsf{Fun}(\mathcal{B},\mathsf{Cat}
_{\infty}),
$$
which means that after fixing a cocartestian
fibration
over $\mathcal{B}$, say
$p:\mathcal{E} \to \mathcal{B}$,
we need to construct a functor $\mathcal{B}
\to \mathsf{Cat}_{\infty}$.
For this we need to specify
the correspondence of morphisms in
$\mathcal{B}$,
namely
$$
\begin{aligned}
\mathcal{B}(b,b') &\longrightarrow
\mathsf{Map}(\mathcal{E}_{b'}, \mathcal{E}_b)
\\
f &\longmapsto f_!.
\end{aligned}
$$

\noindent
This $f_!$ was constructed using the results
of this chapter
(and of the previous!).
This is explained in \cite[Section
3.7]{rune}.
Notice that this is not detailed
enough to be a proof of the equivalence
of infty-categories that we call {\it
straightning}.

This section ends explaining that we may take
op,
and get an isomorphism with cartesian
morphisms
(not cocartesian).

\section{Representability and the Yoneda
embedding}
\label{infty-categories-section-representability}

\noindent
An object $x \in \mathcal{C}$ is {\it initial
} if $\mathcal{C}(x,c)$
is contractible for all $c \in \mathcal{C}$
and
{\it terminal} if $\mathcal{C}(c,x)$ is
contractible for all $c$.
A right fibration $p:\mathcal{E} \to
\mathcal{B}$ is
{\it representable} if $\mathcal{E}$ has a
terminal object
and a left fibration $p:\mathcal{E} \to
\mathcal{B}$ is
{\it corepresentable} if $\mathcal{E}$ has an
initial object.
We say that $p$ is {\it represented} by $b
\in \mathcal{B}$
is $\mathcal{E}$ has a terminal object that
lies over $b$,
and that $x \in \mathcal{E}$ {\it exhibits}
$p$
as represented by $b$ if $x$ is terminal
and $p(x)\simeq b$.

\begin{lemma}
\label{infty-categories-lemma-equivalence-iff-initial}
The object $x$ is initial if and only if
$\mathcal{C}_{x/}\to \mathcal{C}$
is an equivalence, and terminal if and only
if $\mathcal{C}_{/x}\to \mathcal{C}$
is an equivalence.
\end{lemma}

\begin{proof}
Since $C_{x/}\to \mathcal{C}$ is a left
fibration (why?)
it's enough to show that the induced maps on
fibers of
the diagram
$$
\xymatrix{
\mathcal{C}(x,y)\ar[r]\ar[d]&C_{x/}\ar[d]\\
\{y\}\ar[r]&\mathcal{C}
}
$$
is a puback. [See picture.]
\end{proof}

\medskip\noindent
It would be natural to define a left
fibration $p:\mathcal{E} \to \mathcal{B}$
to be co-representable if its corresponding
straightening
$\mathcal{B} \to \text{Grp}_{\infty}$

Recall from
Exercise \ref{infty-categories-exercise-discr-fibra-presh}
that in 1-categories we have an equivalence
between presheaves
$F:\mathcal{B}^{\text{op}}\to \mathsf{Sets}$
and discrete fibrations $p:\mathcal{E} \to
\mathcal{B}$.
We define $F$ to be {\it representable} if
$\mathcal{E}$
has terminal object.

\begin{exercise}
\label{infty-categories-exercise-duplicate-2}
Show that this definition is equivalent to
having $\Hom(-,b_0) \cong F$.
\end{exercise}

\begin{proof}
Recall from Exercise
\ref{infty-categories-exercise-discr-fibra-presh}
that $\mathcal{E}$ is defined so that the
fiber of $b \in \mathcal{B}$
is the set $F(b)$. More exactly, the objects
of $\mathcal{E}$
are pairs $(e,b)$ with $b \in \mathcal{B}$
and $e \in F(b)$.
The morphisms are pairs $(f,F(f))$ for $f:b'
\to b$ in $\mathcal{B}$.

Suppose that $\underline{e}=(e,p(e))$ is a
terminal object of $\mathcal{E}$.
We claim that $\Hom(-,p(e))\cong F$.
Let $b \in \mathcal{B}$. Let's check that
$\Hom(b,p(e))\simeq F(b)$ as sets.
By definition of discrete
fibration for every map $f:b\to p(e)$
there is a unique map $g:e' \to e$ such that
$p(g)=f$.
Then the element $e'$ is in $F(b)$.
To see that this correspondence is natural in
$b$
pick a map $h:b_1 \to b$.
We have an induced map $F(h):F(b) \to F(b_1)$
that for any map $f \in \Hom(b,p(e))$
puts the element we found $e'$
to an element $e_1'$ in the set $b_1$,
which is the $p$-fiber of $b_1$ by definition
of $p$.
$$
\xymatrix{
f\ar@{|->}[d]&\Hom(b,p(e))\ar[r]\ar[d]_{h^*}&
F(b)\ar[d]^{F(h)}&e'\ar@{|->}[d]\\
fh&\Hom(b_1,p(e))\ar[r]&F(b_1)& e_1'.
}
$$
To confirm that the diagram commutes we need
to check
that this element $e_1'$ is the source of the
unique
lift of the map $fh$. Since
$\underline{e}=(e,p(e))$ is terminal,
there is only one map in $\mathcal{E}$
with target $\underline{e}$ and thus
the source of the lift of $fh$ is uniquely
determined.

Conversely, suppose that $\Hom(-,b_0)\cong F$
for some $b_0 \in \mathcal{B}$.
To find a terminal object in $\mathcal{E}$ we
look
for a distinguished element in the fiber of
$b_0$
(which is the set $F(b_0)$ by construction).
We claim that the element corresponding to
the identity
in $\Hom(b_0,b_0)\cong F(b_0)=p^{-1}(b_0)$
is terminal. Let us denote this element by
$\underline{e}=(e,p(e)) \in \mathcal{E}$.
Let $f:\underline{e}' \to \underline{e}$ be
any morphism in $\mathcal{E}$.
This gives a map on $\mathcal{B}$ by
$p(f):p(e')\to p(e)=b_0$.
Since $p$ is a discrete fibration, there
exists a unique morphism
in $\mathcal{E}$ lifting $p(f)$, which must
be $f$.
\end{proof}

\medskip\noindent
The following explanation is quite
illuminating.
Let's informally define a Lie groupoid to be
a groupoid object
in the category of differentiable manifolds
(see nLab \href{url}{Lie groupoid} and
\href{https://ncatlab.org/nlab/show/Lie+%
groupoid}{groupoid
object}).
Actually, this just means that it is a
representable
presheaf of groupoids (perhaps with some
other conditions).
That is, a Lie groupoid is a representable
functor
$\text{Diff}^{\text{op}}\to \text{Grpd}$.

\begin{definition}
\label{infty-categories-definition-representable-presheaf}
We say a presheaf
$\phi:\mathcal{B}^{\text{op}}\to
\mathsf{Gpd}_{\infty}$
is {\it representable} if the corresponding
right fibration
is representable, and that a copresheaf
$\mathcal{B}\to \mathsf{Gpd}_{\infty}$ is {\it
correpresentable}
if the corresponding left fibration is
corepresentable.
We say that $\phi$ is {\it represented} by $b
\in \mathcal{B}$
and that a point $x \in \phi(b)$ {\it
exhibits} $\phi$
as represented by $b$.
\end{definition}

\medskip\noindent
Let's do a quick digression for Yoneda lemma.
First recall the usual Yoneda lemma.
The easiest way to remember this is that
the yoneda map $a \to h_a:=\Hom(-,a)$ is an
embedding $\mathcal{C} \to \text{Psh}$.
But that's not very precise.
What is precise is to say that
$$
\text{Nat}(\Hom(-,b),X) \cong X(b).
$$
This is Yoneda lemma (see Categories, Lemma
\ref{infty-categories-lemma-yoneda}).
There are distinguished objects:
$$
\begin{aligned}
\text{Nat}(\Hom(-,b),X)&\cong X(b)\\
\text{id}_b \mapsto \eta(\text{id}_b)&
\eta(\text{id}_b)\\
(f:c \to b) \mapsto f^*(\eta(\text{id}_b))&
\end{aligned}
$$

\noindent
But we are interested in another version,
let me put it as a lemma:

\begin{lemma}
\label{infty-categories-lemma-yoned-equiv-presh-categ}
Yoneda lemma is equivalent to
$(\Hom(-,b),\text{id}_b)$ is
the initial object in $\mathcal{C}$, where
$\mathcal{C}$ is defined as the
category whose objects are presheaves $X$
over a given category $\mathcal{B}$ together
with a
choice of object $x \in X(b)$,
and morphisms are morphisms of presheaves
that preserve the
chosen object.
\end{lemma}

\noindent
Via straightening, there is more:

\begin{lemma}
\label{infty-categories-lemma-yoneda-equivalence-fibrations}
Yoneda lemma is equivalent
to the category of discrete fibrations
$\mathcal{E} \to \mathcal{B}$
with a choice of object $e \in \mathcal{E}$
over $b$
having an initial object.
(and such an initial object is a discrete
fibration
with a terminal object in $\mathcal{E}$).
\end{lemma}

Left as an exercise.


\medskip\noindent
Now recall Lemma
\ref{infty-categories-lemma-left-fibration-iff-induced},
it says that the top arrow in the following
diagram
is an equivalence
$$
\xymatrix{
\mathcal{E}_x\ar[r]^\sim\ar[d]&\mathcal{B}_{/
b}\ar[d]\\
\mathcal{E}\ar[r]^p&\mathcal{B}.
}
$$
Invert the equivalence to obtain a map over
$\mathcal{B}$:
$$
\xymatrix{
\mathcal{B}_{/b} \ar[rd] \ar[rr] &&
\mathcal{E} \ar[dl]^{p} \\
& \mathcal{B}.
}
$$
which takes $\text{id}_b$ to $x$.
This leads to a section
$s_x:\mathcal{B}_{/b}\to \mathcal{E}$.

\medskip\noindent
A proposition I did not understand properly,
but seems not very hard to prove follows. And
then,

\begin{proposition}
\label{infty-categories-proposition-maps-over-x-repre}
$\text{id}_x$ is initial in
$\mathcal{C}_{x/}$
and terminal in $\mathcal{C}_{/x}$.
In particular, $\mathcal{C}_{/x}\to
\mathcal{C}$
is a representable right fibration
and $C_{x/}\to \mathcal{C}$ is a
corepresentable left fibration.
\end{proposition}

\begin{proof}
Missing.
\end{proof}

\noindent
This leads to a corollary:

\begin{lemma}
\label{infty-categories-lemma-fibra-equiv-iff-maps-over}
A right fibration over $\mathcal{B}$ is
representable
if and only if it is equivalent to
$\mathcal{B}_{/b}\to \mathcal{B}$
for some $b \in \mathcal{B}$,
while a left fibration is corepresentable if
and only if
it is equivalent to $\mathcal{B}_{b/}\to
\mathcal{B}$
for some $b \in \mathcal{B}$.
\end{lemma}

\medskip\noindent
We proved \cite[Corollary 4.2.2]{rune}:
\begin{lemma}
label{lemma-corollary-4.2.2}
Suppose $p:\mathcal{E}\to\mathcal{C}$ is a
right fibration,
then for $x \in \mathcal{C}$, the morphism
$$
Map_{/\mathcal{C}}(\mathcal{C}_{x/},
\mathcal{E})\to
Map_{/\mathcal{C}}(\{x\},\mathcal{E}) \simeq
\mathcal{E}_x
$$
is an equivalence.
Dually, if $p$ is a left fibration we obtain
a corresponding
equivalence.
\end{lemma}

\begin{exercise}
\label{infty-categories-exercise-4-1}
Show that, for a right fibration
$p:\mathcal{E} \to \mathcal{B}$,
the maps $s_x:\mathcal{B}_{/p(x)} \to
\mathcal{E}$
for $x \in \mathcal{E}$ are natural in $x$,
and use this to produce a section of the map
 $$
\mathsf{Map}_{/\mathcal{C}}(\mathcal{C}_{/x},
\mathcal{E})
\to
\mathsf{Map}_{/\mathcal{C}}(\{x\},\mathcal{E}
).
$$

\end{exercise}

\section{Bifibrations}
\label{infty-categories-section-bifibrations}

\noindent
Now consider
 $$
\mathcal{C}(-,-):\mathcal{C}^{\text{op}}
\times\mathcal{C}
\to
\mathsf{\text{Grpd}}_{\infty}
$$
We proved Corollary 4.3.5.
Using this we figured out that we have a map
$$
\begin{aligned}
\mathcal{B}&\longrightarrow
\text{RFib}(\mathcal{A}) \\
b&\longmapsto (\mathcal{E}_b \to \mathcal{A})
\end{aligned}
$$
And we can do it also for (probably)
$$
\begin{aligned}
\mathcal{A}^{\text{op}} &\longrightarrow
\text{LFib}(\mathcal{B}) \\
a &\longmapsto (\mathcal{E}_a \to
\mathcal{A})
\end{aligned}
$$

\noindent
Now we observe that by proposition 3.6.2,
the well-known map
$\mathsf{Ar}(\mathcal{C})\to
\mathcal{C}\times \mathcal{C}$
is a bifibration.
Then we apply the previous observation to get
the {\it Yoneda embedding} statement that
 $$
\mathsf{y}_{\mathcal{C}}\to
\mathsf{RFib}(\mathcal{C})\simeq
\text{Psh}(\mathcal{C})
$$

\begin{lemma}[Pullback of bifibrations is
fibifration]
\label{infty-categories-lemma-pullback-bifibration}
\end{lemma}


% --- End infty-categories.tex ---

% --- Begin commutative-algebra.tex ---
\chapter{Commutative Algebra}
\label{chapter-commutative-algebra}





\section{Rings}
\label{commutative-algebra-section-rings}

\begin{definition}
\label{commutative-algebra-definition-domain}
A ring $R$ is called a {\it domain} if it has
no proper zero divisors,
that is, if $ab=0$ then $a=0$ or $b=0$.
\end{definition}

\begin{theorem}
\label{commutative-algebra-theorem-ideal-quoti-corre}
Let $R$ be a ring and $I \subset R$ an ideal.
There is a correspondence between ideals of
$R$ containing $I$
and ideals of $R/I$.
\end{theorem}

\begin{proof}
We show that there are two maps between the
set of ideals of $R$
containing $I$ and the ideals of $R/I$ whose
compositions
are the identities in the corresponding set.
One map is the projection
to the quotient, and the other map is taking
the elements
whose equivalence class lies in a given
ideal.

Suppose that $\tilde{J}$ is an ideal of
$R/I$.
Define $J=\{j \in R : j+I \in \tilde{J}\}$,
the set of elements that project to
$\tilde{J}$.
Then $I \subset J$, because $i + I = I$ is
the zero of $R/I$
which is contained in any ideal of $R/I$.
Also notice that $J$ is an ideal of $R$ since
for any $r \in R$,
$rj+I \in \tilde{J}$ by $\tilde{J}$ being an
ideal.
Notice that projecting $J$ back to $R/I$
gives $J$
by definition.

Conversely, we project to the quotient.
Let $J$ be any ideal of $R$ containing $I$
Let $\tilde{J}$ be the projection to $R/I$.
Then $\tilde{J}$ is an ideal of $R/I$ by $J$
being an ideal of $R$.
Consider the set $J'$ of all elements whose
projection lies in $\tilde{J}$.

Then $J'=J$ by definition.
\end{proof}

\noindent
This is the essential tool in proving that
$R/I$ is a field
when $I$ is a maximal ideal

\begin{lemma}
\label{commutative-algebra-lemma-quoti-ring-maxim-ideal-field}
Let $R$ be a ring and $I$ a maximal ideal of
$R$.
Then $R/I$ is a field.
\end{lemma}

\begin{proof}
It's easy to prove that a ring is a field if
and only if its
only ideals are $\{0\}$ and $R$. By Theorem
\ref{commutative-algebra-theorem-ideal-quoti-corre},
the ideals of $R/I$ are in correspondence
with ideals of $R$ containing $I$,
but only $R$ is such.
\end{proof}

\section{Polynomial rings}
\label{commutative-algebra-section-polynomial-rings}

\begin{lemma}
\label{commutative-algebra-lemma-ring-polyn-one-indet-over}
The ring of polynomials in one indeterminante
over
a domain $R$ is a domain.
\end{lemma}

\begin{proof}
Suppose that $fg=0$ for $f,g \in R[x]$. Then
the leading coefficient of $fg$ is $a_nb_m$
where $f=\sum_{i=1}^n a_ix^i$ and
$g=\sum_{i=1}^mb_ix^i$.
Since $a_nb_m=0$, then $a_n=0$ or $b_m=0$,
but this is a contradiction since $R$ is a
domain and $a_n \neq 0$, $b_m \neq 0$.
\end{proof}


\begin{exercise}
\label{commutative-algebra-exercise-ring-of-plane-curve-is-pid}
Let $P(x,y)$ be an irreducible polynomial
and $R = \frac{\mathbb{C}[x,y]}{(P)}$.
Prove that all ideals in $R$ are
principal, or find a counterexample.
\end{exercise}

\begin{proof}
This is false. Consider
$R=\mathbb{C}[x,y]/(y-x^2)$
and the ideal $(x,y) \subset
\mathbb{C}[x,y]$.
This ideal contains $y^3-x^2$,
so it's an ideal of $R$,
but it is not principal!
Indeed, even after substituting $y^3=x^2$,
there's no way to obtain $y$ from $x$
or vice versa. [Proof can be improved]
\end{proof}

\begin{exercise}
\label{commutative-algebra-exercise-smooth-irreducible-plane}
Prove that for any $d \in  \mathbb{Z}^{>0}$
there exists an irreducible polynomial
$P \in \mathbb{C}[x,y]$ such that
$P(x,y)=0$ defines a smooth submanifold
in $\mathbb{C}^2$.
\end{exercise}

\begin{proof}
$x^n+y$ since $\partial_y =1$.
This is an irreducible polynomial
because $y$ is irreducible
More precisely,
if $x^n+y = fg$, then either of $f,g$
must have no $y$ variable, say $f \in
\mathbb{C}[x]$,
and the other one must be degree 1 in $y$.
Suppose $g = y g' + g''$ for $g',g'' \in
\mathbb{C}[x]$. But then
$x^n+y = fg = yfg' + f g''$,
which means that $f$ cannot have positive
degree in $x$, but $f \in \mathbb{C}[x]$,
so it is a unit.
\end{proof}


\section{Greatest common divisor and least
common multiple}
\label{commutative-algebra-section-gcd-and-lcm}

Following is the universal definition of lcm
and gdc:

\begin{definition}
\label{commutative-algebra-definition-lcm-and-gdc-universal}
A {\it least common multiple} of $a,b \in R$
is an element $d$ such that if for
every $c\in R$, if $a,b|c$ then $d|c$.

A {\it greatest common divisor} of $a,b \in
R$ is an element $d$ such that
for every  $c \in R$,  if $c|a,b$ then
$c|d$.
\end{definition}

But it looks like it shall be easier to
define them
as follows:

\begin{definition}
\label{commutative-algebra-definition-lcm-and-gdc-polynomial}
A {\it least common multiple} of $f,g \in
R[x]$ is the monic polynomial $d$ of
smallest degree such that $f,g|d$.

A {\it greatest common divisor} of $f,g\in
R[x]$ is the monic polynomial $d$ of
greatest degree such that $d|f,g$.
\end{definition}

\begin{lemma}
\label{commutative-algebra-lemma-existence-of-lcm-and-gcd}
If $R$ is a UFD, in $R[x]$ there exist lcm
and mcd, and are unique up to
multiplication by units.
\end{lemma}

\begin{proof}
We can order by degree, and those of the same
degree order by the leading
term… but this requires an order in $R$.
\end{proof}

\section{Euclidean domains}
\label{commutative-algebra-section-euclidean-domains}

\begin{definition}
\label{commutative-algebra-definition-euclidean-domain}
A domain $R$ is an {\it Euclidean domain} if
there
is a function  $\varphi:R\to \mathbb{Z}_+$,
the positive integers,
such that
\begin{enumerate}
\item $\varphi(ab)\geq \varphi(a)$ for all
$a,b \in R$,
\label{commutative-algebra-item-valuation}
\item For all $a,b \in R$ with $\varphi(a)
\geq \varphi(b)$,
there are $q,r \in R$ such that
$$
a=bq+r\qquad \text{and}
\qquad\varphi(r)<\varphi(b).
$$
\end{enumerate}
\end{definition}

\begin{remark}
\label{commutative-algebra-remark-equivalence-for-valuation}
Item
\ref{commutative-algebra-item-valuation}
is equivalent to
asking that if $a$ divides $b$ then
$\varphi(a) \leq \varphi (b)$.
\end{remark}

\begin{theorem}[Euclidean algorithm for
polynomial rings]
\label{commutative-algebra-theorem-eucli-algor-polyn-rings}
Let $R$ be a domain and $f,g \in R[x]$
with the leading coefficient of $g$ a unit
and $\text{deg}(f)\geq \text{deg}(g)$. Then
there are polynomials $q,r \in R[x]$ such
that
$f=qg+r$ with $\text{deg}(r)<\text{deg}g$.
\end{theorem}

\begin{proof}
Since $\text{deg}(r)$ must be smaller than
$\text{deg}q$
and the leading coefficient of a product is
the
product of the leading coefficients of the
factors,
we readily see that the leading coefficient
of $q$ must
be $L(f)/L(g)$ and of degree $n-m$, where $L$
is the leading
coefficient function and $\text{deg}(f)=n$,
$\text{deg}(g)=m$.

We are lead to the following natural
strategy:
write $g=L(g)x^m+g_1$ for some $g_1$ of
degree
strictly smaller than $m$. Likewise, write
$f=L(f)x^n+f_1$.

Then multiply $g$ with the leading term of
$q$ to obtain
$$
\begin{aligned}
g
\left(\frac{L(f)}{L(g)}x^{n-m}\right)&=(L(g)
x^m+g_1)
\left(\frac{L(f)}{L(g)}x^{n-m}\right)\\
&=L(f)x^n+g_1\left(\frac{L(f)}{L(g)}x^{n-m}
\right).
\end{aligned}
$$
Add $f_1$ on both sides to obtain $f$:
$$
\begin{aligned}
g\left(\frac{L(f)}{L(g)}x^{n-m}\right)+f_1&=
L(f)x^n+f_1+
g_1\left(\frac{L(f)}{L(g)}x^{n-m}\right)\\
\iff
f&=g\left(\frac{L(f)}{L(g)}x^{n-m}\right)+f_1
-g_1\left(\frac{L(f)}{L(g)}x^{n-m}\right)
\end{aligned}
$$
We have great candidates for $q$ and $r$,
except that put this way
it is not necessarily true that the remainder
will have degree
strictly smaller than $g$. However, it is
true that it
will have degree strictly smaller than the
degree of $f$.
Then we can just repeat
the process until we arrive at a reminder
with degree
strictly smaller than the degree of $g$.
\end{proof}

\begin{exercise}
\label{commutative-algebra-exercise-gaussian-integers-euclidean}
Prove that the ring of Gaussian integers
$\mathbb{Z}[\sqrt{-1}]$ is Euclidean.
\end{exercise}

\begin{proof}
Take two Gaussian integers $p,q$ and
take their quotient in $\mathbb{C}$.
By simple geometrical observation,
the nearest Gaussian integer $s$
is at distance at most $1/\sqrt{2}$
from $p/q$. That is,
$$
\left|\frac{p}{q}-s\right|<\frac{1}{\sqrt{2}}
.
$$
Multiplying by $|q|$ we obtain
$$
|p-qs|\leq \frac{|q|}{\sqrt{2}}.
$$
Indeed --- we can view the squared norms
of Gaussian integers as $x \bar{x}$, so that
$$
|q|^2\left|\frac{p}{q}-s\right|^2
=q\bar{q}\left(
\frac{p}{q}-s\right)
\overline{\left(\frac{p}{q}-s\right)}
=\left(q\frac{p}{q}-qs\right)
\overline{\left(q\frac{p}{q}-qs\right)}
=|p-qs|^2.
$$
Defining $r=p-qs$ yields the result.
\end{proof}

\noindent
The same argument shows that
$\mathbb{Z}$,
$\mathbb{Z}[\sqrt{-2}]$ and
$\mathbb{Z}
\left[\frac{1+\sqrt{-3}}{2}\right]$
are also Euclidean rings.
But we can't use it for
$\mathbb{Z}[\sqrt{-3}]$
or $\mathbb{Z}[\sqrt{5}]$
--- is any of these Euclidean?
Neither, since they are not
integrally closed, while we have:

\begin{lemma}
\label{commutative-algebra-lemma-euclidan-is-ufd}
A PID is integrally closed.
\end{lemma}

\begin{proof}
?
\end{proof}

\section{Principal ideal domains}
\label{commutative-algebra-section-pid}

\begin{lemma}
\label{commutative-algebra-lemma-pid-free-module}
Let $R$ be a PID, $M$ a a free
$R$-module of finite rank $n$ and $N$ a
submodule of $M$. Then $N$ is a free module
of rank $m\leq n$ and there exists a basis
$y_1,y_2,\ldots,y_n$ of $M$ such that
$a_1y_1,a_2y_2,\ldots,a_my_m$ is a basis of
$N$, where $a_1,a_2,\ldots,a_m$ are nonzero
elements of $R$ with the divisibility
relations
$$
a_1|a_2|\ldots|a_m.
$$
\end{lemma}

\section{Unique factorization domains and
the
Gauss lemma}
\label{commutative-algebra-section-UFD}

The point of the Gauss lemma is that a
polynomial $f$ over the field of fractions
$F$ of a UFD $R$ can be expressed as a
rational number called the content of $f$
times a polynomial over $R$ called the
primitive part. The point of the Gauss
lemma is that it allows us to make sure that
the content will in fact be an
element of $R$.

\medskip\noindent
The following lemma shows that in UFDs
irreducibles are primes, but this is not
true in general. (Counterexample in
$\mathbb{Z}[\sqrt{-5}]$.)

\begin{lemma}
\label{commutative-algebra-lemma-irreducible-is-prime}
If $R$ is a UFD then irreducible elements are
prime.
\end{lemma}

\begin{proof}
If $f$ is irreducible and $f|ab$ then $f$
must be an element of the unique
factorization of $a$ or $b$.
\end{proof}

In fact, showing that irreducibles are prime
would be enough for our purpose of
showing that $F[x]$ is a UFD when $F$ is a
field. We shall be able to prove this
(without the assumption of UFD as above) via
Bachet-Bézout theorem; which we
state first in its version for integer
numbers:

\begin{theorem}[Bachet-Bézout for integers]
\label{commutative-algebra-theorem-Bachet-Bezout-for-integers}
\begin{reference}
\cite[Theorem 1.7]{MMST}
\end{reference}
Let $a,b \in \mathbb{Z}$ (what is the
essential property of $\mathbb{Z}$ that
makes this true?). Then
there exist $x,y \in \mathbb{Z}$ with
 $$
ax+by=\text{gcd}(a,b)
$$
Therefore, if $c \in \mathbb{Z}$ divides both
$a$ and $b$ then $c$ divides
$\text{gcd}(a,b)$ as well.
\end{theorem}

\begin{lemma}
\label{commutative-algebra-lemma-relat-prime-divid-impli-divid}
Let $a,b,c \in \mathbb{Z}$. If $a|bc$ and
$(a,b)=1$, then $a|c$.
\end{lemma}

\begin{proof}
By Theorem
\ref{theorem-Bachet-Bezout-%
for-integers},
there exist elements $p,q$ such
that $ap+bq=1\implies apc+bcq=c$, so that $a$
divides both summands and thus
divides $c$.
\end{proof}

\begin{theorem}[Division algorithm for
polynomials]
\label{commutative-algebra-theorem-divis-algor-polyn}
\begin{reference}
\cite[Theorem 1.48]{MMST}
\end{reference}
Let $F$ be a field. Given two polynomials
$f,g \in F[x]$ there are two unique
$q,r \in F[x]$ such that
$$
f=qg+r,\qquad \text{deg}r<\text{deg}g.
$$
{\bf Remark.} We use the hypothesis that $F$
is a field to take a fraction with
denominator the leading coefficient of $g$;
that is, the theorem is valid if the
leading coefficient of $g$ is a unit.
\end{theorem}

\begin{proof}
We may suppose that $\text{deg}f \geq
\text{deg}g$ for if not we just pick
$q=0$ and $r=f$. Now we use induction.
If $\text{deg}f=\text{deg}g=0$ then $f$ and
$g$ are numbers and the result is
obvious since $F$ is a field.

If $\text{deg}f>0$ we do as follows. Suppose
$\text{deg}f=n$ and
$\text{deg}g=m$. Write $f(x)=a_nx^n+f_1(x)$
and $g(x)=b_mx^m+g_1(x)$. Notice
that
$$
\frac{a_nx^n}{b_mx^m}g(x)=a_nx^n+\frac{a_n}
{b_m}x^{n-m}g_1(x)
$$
so that
$f(x)-\frac{a_n}{b_m}x^{n-m}g(x)=f_1(x)-
\frac{a_n}{b_m}x^{n-m}g_1(x)$.
This is a polynomial of strictly less degree
than that of $f$, so that we can
use induction hypothesis to write it as
$qg+r$ with $\text{deg}r<m$. Then
$$
f(x)=g(x)\left(q(x)+\frac{a_n}{b_m}x^{n-m}
\right)+r(x).
$$
To prove uniqueness suppose that
$f=q_1g+r_1=q_2g+r_2$.
Then
$$
(q_1-q_2)g=r_1-r_2
$$
Taking degrees we see that the degree of this
polynomial should be,
by the right-hand-side less than the degree
of $g$,
and by the left hand side more than the
degree of $g$.
Then the degree must be zero and $q_1=q_2$,
$r_1=r_2$.
\end{proof}

Recall that the greatest common divisor of
two polynomials $f$ and $g$
is the {\it monic}
polynomial of least degree that divides both
$f$ and $g$.

\begin{theorem}[Bachet-Bezout for
polynomials
over field]
\label{commutative-algebra-theorem-bache-bezou-polyn-over-fiel}
\begin{reference}
\cite[Theorem 1.53]{MMST}
\end{reference}
Let $F$ be a field and $f,g\in F[x]$.
There exist elements  $p,q \in F[x]$ such
that $fp+gq=\text{gcd}(f,g)$.
\end{theorem}

\begin{proof}
Consider the set
$$
I:=\{fx+gy:x,y \in F[x]\}.
$$
Let $d=fp+gq$ be the monic polynomial of
least degree
of the form $fx+gy$, for $x,y \in F[x]$.
Notice that the existence of $d$ monic comes
from $F$ being a field.
We claim that $d=\text{gcd}(f,g)$.
Indeed, we first check by division algorithm
that $d$ divides both $f$ and $g$:
suppose that $f=dq+r$ with
$\text{deg}r<\text{deg}d$.
Then, if $r=f-dq=f(1-xq)-yg$ is not zero,
dividing by its leading coefficient gives an
element in $I$
of degree strictly smaller than the degree of
$d$, a contradiction.
Thus $r$ is zero. The same argument shows
that $d|g$, and we conclude that
Thus $\text{deg}d \leq
\text{deg}\text{gcd}(f,g)$.

Now, since
$\text{gcd}(f,g)$ divides both $f$ and $g$,
it divides $d=xf+yg$. (Indeed, if
$c:=\text{gcd}(f,g)$ is such that $ca=f,cb=g$
then $c(ap+bq)=cap+cbq=fp+gq=d$.)
Thus $\text{deg}d=\text{deg}\text{gcd}(f,g)$,
and since both are monic we are
done.
\end{proof}

Finally we obtain

\begin{lemma}
\label{commutative-algebra-lemma-irreducible-is-prime-in-F-x}
Let $F$ be a field. Irreducible elements of
$F[x]$ are prime. That is, if $f|gh$
then $f|g$ or $f|h$ for any $f,g,h \in F[x]$.
\end{lemma}

\begin{proof}
Let $d=\text{gcd}(f,g)$. Then $dp=f$ and
$dq=g$. If $d$ is not $1$
we get that $p$ is a unit. Then $gp$ is a
multiple by a unit of $g$ which is
divided by $f$: $gp=dpq=fq$. That's good
enough.

Now suppose that $d=1$. Then by
Bachet-Bezout
Theorem
\ref{theorem-bache-bezou-%
polyn-over-fiel}
there are $p,q \in F[x]$ such that $fp+gq=1$.
Then $fph+gqh=h$. Since by
hypothesis $f|gh$ we conclude that $f$
divides both summands, so $f|h$.
\end{proof}

Applying this last lemma over and over again
we get that if $f|f_1\ldots f_k$
then $f$ must divide $f_i$ for some $i$.

\begin{lemma}
\label{commutative-algebra-lemma-polynomials-over-fields-are-UFD}
If $F$ is a field then $F[x]$ is a UFD.
\end{lemma}

\begin{proof}
The existence of the factorization is an easy
induction on the degree of $f\in
R[x]$. If $f$ has degree zero, then we are
done. If $f$ has positive degree and
is irreducible, we are done. If $f$ has
positive degree and it is reducible,
i.e. $f=gh$, then both $g$ and $h$ have
strictly smaller degree and we may apply
induction hypothesis to obtain a
decomposition of each factor into irreducible
elements.

To prove uniqueness suppose that there exists
an element in $F[x]$ that admits
more than one factorization in irreducible
elements
up to multiplication by units. Choose $f$ so
that
$\text{deg}f$ is minimal among all such
elements.

Write two factorizations of $f$ in
irreducible elements as
$f_1\ldots f_k=f_1'\ldots f_\ell'$ . Suppose
that the factors are in
ascending order of degree.

We have that $f_1|f_1'\ldots f_{k'}$, so
that $f_1$ must divide some $f_{i'}'$
by Lemma
\ref{lemma-irreducible-is-%
prime-in-F[x]}.
Since $f'_{i'}$ is irreducible and $f_1$ is
not a unit, we conclude that $f_{i'}=f_1 u$
for some unit $u$.

Then we apply cancellation law using that
$R[x]$ is a domain
to obtain $uf_2\ldots f_k=f_2'\ldots f'_{k'}$
is a
polynomial with two distinct factorizations.
Since $f$ is minimal among all
elements that admit more than one
factorization, we conclude that $k=k'$ and
$f_i=f'_i$ for all $i$.
\end{proof}

\begin{remark}
\label{commutative-algebra-remark-not-the-same-as-in-Z}
Notice that the proof given for integers does
not work identically,
for in this case we cannot conclude that
$f_1=f_1'$ directly from putting the
factors in order of degree (we can conclude
that they have the same degree, but
not that they arr the same polynomial!).
\end{remark}

\medskip\noindent
The definition of content can be done in
three steps:

\begin{definition}
\label{commutative-algebra-definition-content}
\begin{reference}
\cite[p. 181]{Lang-algebra}
\end{reference}
Let $R$ be a factorial ring and
$F$ its fraction field.

\begin{enumerate}
\item (Order of a fraction $a$ at a prime
$p$.)
Let $a \in F$. If $p \in R$ is irreducible,
the {\it order} of $a$ at $p$
to be the number $r$ such that $a=p^rb$ for
some fraction $b$ that does not have
$p$ as factor in the numerator nor in the
denominator. (To understand the point
of the Gauss lemma, notice that the order
can be a fraction, i.e. if $p$ is a factor of
the denominator.)

\item (Order of a polynomial $f \in F[x]$ at
a prime $p$.)
The {\it  order of $f\in F[x]$ at $p$} is the
minimium order of
$p$ in any of the coefficients of $f$.
(Defining this number as the {\bf
minimum} ensures that the primitive part of
$f$ (see below) has all coefficients
in $R$. It also follows that the gcd of the
coefficients of the primitive part is 1.)

\item (Content of $f \in F[x]$.) The {\it
content} of $f$ is the product
$$
\prod p^{\operatorname{ord}_p(f)}
$$
over all primes $p$ such that
$p^{\operatorname{ord}_p(f)}$ is nonzero.

\item (Primitive polynomial.) A polynomial is
called {\it primitive} if it has content 1.
\end{enumerate}
\end{definition}

As explained in the beginning of this
section, the point of the Gauss
lemma is that it will allow us to make sure
that the content of $f$ is an
element of $R$, i.e. not a fraction.

Also, in the definition of order of a
fraction $a$ at a prime $p$, notice
 the order well-defined, i.e. unique, by the
 uniqueness of the factorization
in $R$.

\begin{lemma}[True Gauss lemma]
\label{commutative-algebra-lemma-true-Gauss}
\begin{reference}
\cite{Chapter 4, Theorem 2.1}{Lang-algebra}
\end{reference}
Let $R$ be a UFD and $F$ its fraction field.
Let $f,g\in F[x]$. Then the
$\text{cont}(fg)=\text{cont}(f)\text{cont}(g)
$.
\end{lemma}

This clearly implies that the product of two
primitive polynomials being
primitive, and the converse I think is true
as well.

\begin{proof}
Since $\text{cont}(bf)=b\text{cont}(f)$ it's
enough to suppose that  $f$ and $g$
are primitive. Indeed, taking content of $fg$
we obtain on one hand
$\text{cont}(fg)$ and on the other hand
$\text{cont}(\text{cont}(f)\text{pp}(f)
\text{cont}(g)\text{pp}(g))
=\text{cont}(f)\text{cont}(g)\text{cont}
(\text{pp}(f)\text{pp}(g))$.
Thus we would finish if we show that the
product of two primitive polynomials is
primitive.

Put another way, it's enough to show will
show that if $f,g$ are primitive (as
are the primitive parts of $f$ and $g$ in the
preceding paragraph) then $fg$ is
primitive. This means that I will show that
any $p$ is not a factor of all the
coefficients in $fg$. A coefficient of $fg$
looks as
$c_j=\sum_{i=0}^ja_ib_{j-i}$.  Since $f,g$
are primitive then $p$ cannot divide
all $a_i$ and all $b_i$.  Thus we can pick
the maximum $a_s$ and the maximum
$b_t$ that $p$ does not divide. Then the
$c_j$ containing the product $a_sb_t$
will not be divided by $p$, for $p$ cannot
divide $a_sb_t$ because it is prime,
but will divide the remaining terms in the
sum $c_j$; thus if $p$ divided $c_j$
it would also divide $a_sb_t$lmail.
\end{proof}

\begin{lemma}[Gauss]
\label{commutative-algebra-lemma-R-UFD-implies-R-x-UFD}
\begin{reference}
\cite[Chapter 4, Theorem 2.3]{Lang-algebra}
\end{reference}
If $R$ is a UFD then $R[x]$ is a UFD.
\end{lemma}

\begin{proof}
Let $f\in R[x]$ be distinct from 0. Since
$F[x]$ is UFD we can write
$f=f_1\ldots f_k$. Notice that the primitive
parts of the $f_i$ coincide
with the primitive part of the product
$f_1\ldots f_k$ by
the Gauss Lemma:
$$
\begin{aligned}
f&=c(f_1\ldots f_k)\text{pp}(f_1\ldots
f_k),\\
f&=c(f_1)\text{pp}(f_1)\ldots
c(f_k)\text{pp}(f_k)
\end{aligned}
$$
and we cancel the content part dividing
(using that we are in a field) by the
Gauss Lemma.

We also apply the Gauss Lemma to conclude
that
$c(f)=c(f_1)\ldots c(f_k)$, being the content
of a polynomial with coefficients
in $R$, must be in $R$ as well, so that we
can decompose it as a product of
irreducibles since $R$ is a UFD.

To obtain a factorization in irreducibles we
only need to show that each of the
$\text{pp}(f_i)$ is irreducible in $R[x]$.
To see this suppose that
$\text{pp}(f_i)=gh$.
$c(f_i)\text{pp}(f_i)=f_i$ is irreducible
in $F[x]$, so either $c(f_i)g$ or $h$ is a
unit of $F$. Then
either $g$ or $h$ are in $R$. But the gcd of
the coefficients of $\text{pp}(f_i)$ is 1,
so, since we have factored
an element of $R$ from $\text{pp}(f_i)$, it
must be 1.

The uniqueness of the factorization follows
from the uniqueness of the
factorization in $F[x]$, for another
factorization in $R[x]$ would be a
factorization in $F[x]$
Indeed, to see this suppose that $f=gh$ with
$g,h \in F[x]$.
Then take contents. The content of  $f$ must
be a unit
for otherwise the principal part of $f$ would
have to be
trivial, and $f$ is not a constant.
By the Gauss Lemma the contents of $g$ and
$h$ are units.
Then we are left with
$f=\text{pp}(g)\text{pp}(h)$, an equation
in $R[x]$ that makes either of the principal
parts a unit of
$R[x]$.

Then any factorization in $R[x]$ has the
same number of factors,
and any two factorizations are equal up
to a fraction $u$,
which is a unit of $F[x]$. We just need to
make sure $u$ is a unit of $R$.
Write uniqueness of the factorization in
$F[x]$ as:
$f_1=uf_1'$ and $f_i=f_i'$ for $i\neq 1$.
It is clear that $u \in R$, since otherwise
$f_1$ couldn't
be a polynomial in $R[x]$. But since $f_1$ is
irreducible in $R[x]$
and $f_1'$ is not a unit of $R[x]$,
we conclude that $u$ is a unit of $R[x]$.
\end{proof}

\begin{lemma}
\label{commutative-algebra-lemma-ufd-irred-eleme-gener-prime}
In a UFD, irreducible elements generate prime
ideals.
\end{lemma}

\begin{proof}
Let $f$ be irreducible. Suppose that
$pq\in(f)$. Then $pq=fg$ for some  $g\in
R$. Then $p_1\ldots p_kq_1\ldots
q_\ell=fg_1\ldots g_m$. But since $f$ is
irreducible and factorization is unique we
conclude that $f$ must be one of the
 $p_i$ or one of the $q_i$. Then $f$ divides
 $p$ or $q$, i.e. either $p\in(f)$
or $q\in(f)$.
\end{proof}

\begin{exercise}
\label{commutative-algebra-exercise-rational-functions-irreducible}
Let $k=\mathbb{C}(t)$ be the field of
rational functions on $\mathbb{C}$.
Prove that for any $n$ there exists
an irreducible polynomial
$P(t) \in k[t]$ of degree $n$.
\end{exercise}

\begin{proof}
Let $P(x) \in k[x]=\mathbb{C}(t)[x]$
be $P(x) = x^n-t$.
This polynomial has no root in
$k=\mathbb{C}(t)$,
for if $(f/g)^n=t$ then 
$n(\operatorname{deg}f-
\operatorname{deg}g)=1$,
which is impossible for $n>1$.
But unfortunately this doesn't prove
that it's irreducible: in $\mathbb{R}[x]$,
the polynomial
$(x^2 + 1)(x^2+2)=x^4+3x^2+2$
is reducible and has no roots.

To prove $x^n + t$ is irreducible we shall
use the Gauss lemma. Forget the above
notation and suppose that
$fg = x^n + t$, where
$f,g \in \mathbb{C}(t)[x]$.
We want to clear the denominators in
the coefficients of $f$ and $g$,
so we multiply by the product $D$ of all
such denominators and obtain
$$
f_1g_1 = (x^n + t)D,
$$
where $D \in \mathbb{C}[t]$ and
$f_1,g_1 \in \mathbb{C}[t][x]$.
Now we apply Gauss lemma, that is, the
content of the product is the product
of the contents. Since $x^n +t$ is primitive
(i.e. has content $1$), then the
content of $f_1g_1$ equals the content
of $D$, but $D$ is a constant
(constants here are polynomials in
$\mathbb{C}[t]$).
We obtain
$$
FG = x^n + t
$$
where $F$ and $G$ are the primitive parts
of $f_1$ and $g_1$. This holds up to
multiplication by a complex number.

Now we reduce modulo $t$ and obtain
$$
[F][G] = [x^n],
$$
and conclude that, as classes,
$$
[F]=[cx^m],\qquad 
[G]=[c^{-1}x^{n-m}].
$$
Thus we see that each of $F,G$
lost their constant term when we reduced,
i.e. that their constant terms are divisible
by $t$. This means that $t^2$ divides
the constant term of the product $FG$.
But this is impossible, since the constant
term of $FG$ is $t$. 

The only way
this makes sense is if $m=0$ or $m=n$.
To conclude we note that reduction modulo
$t$ preserves the degree of both $F$, and 
$G$ because their leading coefficients
multiply to $1$. Thus $m=0$ or $m=n$
imply that $F$ or $G$ is a unit, which
in turn implies so for $f$ or $g$.
\end{proof}


\begin{remark}
\label{commutative-algebra-remark-alternative-approach}
At some point of this notes,
there was another write-up of the proof
of the Gauss lemma. This may be found in
the Git repository: look for a section
called ``Unique factorization
domains (alternative approach)''.
\end{remark}

\section{Modules}
\label{commutative-algebra-section-modules}

It looks like most books (at least
\cite{eis} and \cite{Samuel-Zariski-Vol1})
define a module to be an abelian group
along with a certain operation with
elements of the ring.

Which secretly just says that

\begin{definition}
\label{commutative-algebra-definition-module}
Let $R$ be a ring.
An {\it (left) $R$-module} $M$ is an abelian
group
such that there exists a (left) ring morphism
$$
R \to \text{End}(M).
$$
\end{definition}

\noindent
Indeed, the three usual requirements
correspond to:

\begin{enumerate}
\item The endomorphism corresponding to $a
\in R$
respect the group structure of the group $M$:
$$
a(x+y)=ax+ay,
$$

\item The representation map is a map of
rings
$$
(a+b)x=ax+by,\qquad (ab)x=a(bx).
$$
\end{enumerate}

\section{Invariant factors theorem}
\label{commutative-algebra-section-invariant-factors-theorem}
 
\begin{exercise}
\label{commutative-algebra-exercise-invariant-factors-theorem}
Let $R$ be a principal ideal ring. Prove that
any finitely-generated $R$-module is a direct
sum of cyclic $R$-modules. Use this result to
deduce the Jordan normal form theorem, and to
classify the finitely-generated abelian
groups.
\end{exercise}

I follow \cite[Section 12.1]{dummit} 
to prove the
invariant factors theorem in
the following form:

\begin{theorem}[Invariant factors]
\label{commutative-algebra-theorem-invariant-factors}
Let $R$ be a principal ideal domain and
let $M$ be a finitely generated $R$-module.
Then
$$
M\cong R^r\oplus R/(a_1)\oplus R/(a_2)
\oplus \ldots\oplus R/(a_m)
$$
for some integer $r\geq 0$ and nonzero
elements $a_1,a_2,\ldots,a_m$ of $R$ which
are not units in $R$ and which satisfy the
divisibility relation
$$
a_1|a_2|\ldots|a_m.
$$
\end{theorem}

\begin{proof}
This strongly relies on 
Lemma \ref{commutative-algebra-lemma-pid-free-module}.

Choose a basis $\{x_i\}_{i=1}^n$ of $M$.
Consider the free module $R^n$ along with a
basis $\{b_i\}_{i=1}^n$. The homomorphism
$\pi:R^n\to M$, $b_i\mapsto x_i$ is
surjective, so we get
$R^n/\ker \pi\cong M$.

Apply the theorem for the module $R^n$ and
its submodule $\ker \pi$. We get a basis
$\{y_i\}_{i=1}^n$ of $R^n$ such that
$\{a_iy_i\}_{i=1}^m$ is a basis of
$\ker \pi$ for $a_i\in R$ for
$i=1,\ldots,m$ such that
$a_1|\ldots|a_m$. Then
$$
M\cong R^n/\ker \pi
=
(Ry_1\oplus \ldots\oplus Ry_n)/
(Ra_1y_1\oplus \ldots\oplus Ra_my_m).
$$
It remains to show that the quotient on the
right-hand side of last equation is in fact
the desired decomposition of $M$.

Consider the map
$$
Ry_1\oplus \ldots\oplus Ry_n
\longrightarrow
R/(a_1)\oplus \ldots\oplus R/(a_m)
\oplus R^{n-m}
$$
given by
$$
(\alpha_1y_1,\ldots,\alpha_ny_n)
\longmapsto
(\alpha_1+(a_1),\ldots,\alpha_m+(a_m),
\alpha_{m+1},\ldots,a_n).
$$
The kernel of this map is the set of
elements which $\alpha_i\in(a_i)$ for all
$i=1,\ldots,m$. Thus we can write the
kernel as
$$
Ra_1y_1\oplus \ldots\oplus Ra_my_m.
$$
\end{proof}

\noindent
Looks like I didn't use the divisibility
condition $a_1|\ldots|a_m$.

\medskip\noindent
Now let's deduce the Jordan normal form
theorem using again \cite{dummit}. 

\begin{definition}
\label{commutative-algebra-definition-jordan-form}
A $k\times k$ matrix with $\lambda$
along the main diagonal
and 1 along the first
superdiagonal, that is,
$$
\begin{pmatrix} 
\lambda & 1\\
& \lambda & \ddots\\
& & \ddots\\
& & & \lambda & 1\\
& & &   & \lambda
\end{pmatrix} 
$$


A matrix is said to be in {\it Jordan
canonical form} if it is a block
diagonal matrix with Jordan blocks 
along the diagonal.
\end{definition}

\begin{theorem}[Jordan canonical form]
\label{commutative-algebra-theorem-jordan-form}
\begin{reference}
\cite[Theorem 22, Section 12.3]{dummit}
\end{reference}
Let $A$ be an $n \times n$
matrix over the field $F$
and assume $F$ contains
all the eigenvalues of $A$.
\begin{enumerate}
\item The matrix $A$ is similar
to a matrix in Jordan canonical form,
i.e. there is an invertibnle $n \times n$
matrix $P$ over $F$ such that
$P^{-1}AP$ is a block diagonal matrix
whose diagonal blocks are the 
Jordan blocks for the elementary
divisors of $A$.
\item The Jordan canonical form
of $A$ is unique up to permutation
of the Jordan blocks along the diagonal.
\end{enumerate}
\end{theorem}

\begin{proof}
Fix a vector space $V$ over a field $F$
and a linear transformation $T$. This makes
$V$ into an $F[x]$-module by substitution of
the variable $x$ by $T$.
Note that if $V$ is finite-dimensional
then there is no free part,
that is $F[x]$ is not contained
in $V$ since $F[x]$ is
infinite-dimensional as a vector space
over $F$ --- so we only have the
torsion part in the invariant factor
decomposition.

The ``invariant factors'' are the elements
$a_i$, which in our
present case can be shown to be monic
polynomials $a_i(x)$ satisfying
$a_1(x)|\ldots|a_n(x)$. 
AI-facts: $a_n$ is the minimal polynomial
of $T$, and $a_1a_2… a_n$ is the
characteristic polynomial of $T$.
So the canonical example should be
$a_i(x)=(x-\lambda)^k$
where $\lambda$ is an eigenvalue of $T$.

Further, these
elements are associated to the so-called
{\it elementary divisors}, which concern
another similar formulation of the Invariant
factors theorem. The elementary divisors are
powers of the irreducible components of the
$a_i(x)$.

If we assume that the $a_i(x)$ factor into
linear polynomials, 
i.e. that the field contains all the
eigenvalues of $T$,
the elementary divisors
can be written as $(x-\lambda)^k$, where
$\lambda$ is one of the eigenvalues of $T$
(we consider only one of the eigenvalues at
this point since we are constructing
{\it one} of Jordan blocks of $T$).

Using the Invariant factors theorem for
elementary divisors, we see that $V$ is the
direct sum of finitely many cyclic modules of
the form $F[x]/(x-\lambda)^k$. Such quotients
have basis
$$
\bar{x}^{k-1},\bar{x}^{k-2},\ldots,
\bar{x},1.
$$
Now we observe that the polynomials
$$
(\bar{x}-\lambda)^{k-1},
(\bar{x}-\lambda)^{k-2},\ldots,
\bar{x}-\lambda,1
$$
are also a basis. This follows since
expanding the latter in terms of the
$\bar{x}_i$ gives a triangular matrix, which
makes it invertible, so it's a valid change
of coordinates.

Next we observe that action of multiplication
by $x$, which is the same as applying $T$,
maps these basis elements in the quotient as
follows:
$$
1\mapsto 1\bar{x}
=
\lambda\cdot 1+\bar{x}-\lambda
$$
$$
(\bar{x}-\lambda)
\mapsto
\bar{x}^2-\lambda\bar{x}
=
x\lambda\cdot(\bar{x}-\lambda)
+
(\bar{x}-\lambda)^2
$$
$$
\vdots
$$
$$
(\bar{x}-\lambda)^{k-2}
\mapsto
\lambda\cdot(\bar{x}-\lambda)^{k-2}
+
(\bar{x}-\lambda)^{k-1}
$$
$$
(\bar{x}-\lambda)^{k-1}
\mapsto
\lambda\cdot(\bar{x}-\lambda)^{k-1}
+
(\bar{x}-\lambda)^k
$$
(where I'm not sure how to compute the last
two). The last expression simplifies further
since $(\bar{x}-\lambda)^k=0$ in the
quotient, so that multiplication of $x$ by
the basis element $(\bar{x}-\lambda)^{k-1}$
is barely multiplication by $\lambda$. In sum,
we have constructed the {\it Jordan block}
$$
J_{\lambda}
=
\begin{pmatrix}
\lambda &1\\
&\lambda\\
&&\ddots&1\\
&&&\lambda\\
&&&&1\\
&&&&\lambda
\end{pmatrix}
$$
applying this procedure to the rest of the
cyclic components of $V$ we obtain the
matrix representation
$$
\begin{pmatrix}
J_1\\
&J_2\\
&&\ddots\\
&&&J_t
\end{pmatrix}
$$
of $T$.
\end{proof}

\noindent
In the context of representation theory,
$T$ is diagonalizable if and only if
$$
V=\bigoplus_\lambda E_\lambda,
$$
and every eigenspace $E_\lambda$
is invariant under the $F[x]$ action
on $V$ induced by $T$. That is,
the representation is semisimple
if and only if $T$ is diagonalizable.

More demanding is that the representation
be irreducible: this says that
the invariant factor decomposition is
$$
V=\frac{F[x]}{(a_1)}.
$$
Here's a more precise AI-generated note:

\begin{example}
\label{commutative-algebra-example-irreducible-vs-semisimple-z}
Let $V$ be a finite-dimensional complex
vector space and let
$$
T:V\to V
$$
be a linear map. We regard $V$ as a
$\mathbb{C}[x]$-module by setting
(on the generator $x$,
this is just as above)
$$
xv=T(v).
$$
Then a submodule is exactly a
$T$-invariant subspace.

Therefore $V$ is irreducible as a
$\mathbb{C}[x]$-module if and only if $V$
has no nonzero proper $T$-invariant
subspaces. But every linear operator on a
nonzero finite-dimensional complex vector
space has an eigenvector. If $\dim V>1$ and
$$
Tv=\lambda v,
$$
then the line $\mathbb{C}v$ is a nonzero
proper $T$-invariant subspace. Hence $V$
cannot be irreducible.

Thus the irreducible finite-dimensional
$\mathbb{C}[x]$-modules are exactly the
one-dimensional ones.
(An irreducible representation over
a complex vector space is one dimensional!
Regardless the group!)
Equivalently, they are
the modules
$$
\mathbb{C}[x]/(x-\lambda).
$$

This should be compared with semisimplicity.
For example, if
$$
T=
\begin{pmatrix}
1&0\\
0&2
\end{pmatrix},
$$
then $T$ is diagonalizable and
$$
V=
\mathbb{C}e_1\oplus\mathbb{C}e_2.
$$
Each summand is irreducible, but $V$ itself
is not irreducible. It is semisimple.

So irreducible means one simple piece,
whereas semisimple means a direct sum of
simple pieces.

For representations of $\mathbb{Z}$ over
$\mathbb{C}$, the generator acts by an
invertible operator $T$. Hence the
irreducible representations are the
one-dimensional representations where the
generator acts by some
$\lambda\in\mathbb{C}^{*}$, and a
finite-dimensional representation is
semisimple if and only if $T$ is
diagonalizable.
\end{example}


\noindent
Finally, the classification of finitely
generated abelian groups follows by direct
application of the Invariant factors theorem
choosing $R=\mathbb{Z}$, since abelian
groups are in correspondence with
$\mathbb{Z}$-modules.



\section{Algebras}
\label{commutative-algebra-section-algebras}

\begin{definition}
\label{commutative-algebra-definition-algebra}
\begin{reference}
\cite{Eisenbud}
\end{reference}
If $R$ is a commutative ring,
then a {\it commutative algebra} over $R$
is a commutative ring $S$ together
with a ring morphism $R \to S$.
\end{definition}

\noindent
But actually I was expecting that
$S$ would be defined as
a ring that is also an $R$-module.
So let us note that a morphism of rings
$R \to S$ gives a representation
$R \to \text{End}(S)$ via left
multiplication.
But the other way around,
given an endomorphism associated
to some element in $R$,
how do we assign an element of $S$
so as to produce a map $R \to S$?

\medskip\noindent
Now we briefly introduce finitely-generated
algebras and modules.
Upshot: a finitely-generated $R$-algebra
$S$ is such that $S \simeq
R[x_1,\ldots,x_n]/I$ for some ideal $I
\subset R$. Another way of saying this is
that there is a surjective morphism
$R[x_1,…,x_n]\to S$. Finitely-presentedness
is when $I = (s_1,\ldots,s_k)$.

As an aside, recall that finitely presented
$k$-algebras with no nilpotents are in
correspondence with
affine algebraic varieties.

\begin{definition}
\label{commutative-algebra-definition-finite-type}
Let $R \to S$ be a ring map.
\begin{enumerate}
\item We say $R \to S$ is of {\it finite
type},
or that {\it $S$ is a finite
type $R$-algebra} if there exist an $n \in
\mathbf{N}$
and a surjection
of $R$-algebras $R[x_1, \ldots, x_n] \to S$.
\item We say $R \to S$ is of
{\it finite presentation} if there
exist integers $n, m \in \mathbf{N}$ and
polynomials
$f_1, \ldots, f_m \in R[x_1, \ldots, x_n]$
and an isomorphism of $R$-algebras
$R[x_1, \ldots, x_n]/(f_1, \ldots, f_m) \cong
S$.
\end{enumerate}
\end{definition}

\noindent
For modules, the definition of
finitely generated $R$-module is that there
exists a surjective morphism $R^n
\to\!\!\!\!\!\to S$.
So: finitely-generated
algebra means elements are polynomials on
the generators with coefficients in the
ring, and finitely-generated
module means elements are
linear combinations on the generators with
coefficients in the ring.


\section{Hilbert's Nullstellensatz}
\label{commutative-algebra-section-nullstellensatz}

\section{Spectrum of a ring}
\label{commutative-algebra-section-spectrum}

\noindent
From Stacks Project:

\begin{definition}
\label{commutative-algebra-definition-spectrum-ring}
Let $R$ be a ring.
\begin{enumerate}

\item The {\it spectrum} of $R$ is the set of
prime ideals of $R$.
It is usually denoted $\Spec(R)$.

\item Given a subset $T \subset R$ we let
$V(T) \subset \Spec(R)$ be the set of primes
containing $T$, i.e., $V(T) = \{ \mathfrak p
\in \Spec(R) \mid \forall f\in T, f\in
\mathfrak p\}$.

\item Given an element $f \in R$ we let $D(f)
\subset \Spec(R)$ be the set of primes not
containing $f$.

\item The {\it maximal spectrum} if a ring $R$
is the set of maximal ideals with the same
topology.
\end{enumerate}
\end{definition}

\begin{exercise}
\label{commutative-algebra-exercise-maximal-spectrum}
Let $R=\mathbb{R}[t]$ and $X$ be its
maximal spectrum. Prove that $X$ with its
Zariski topology is homeomorphic
to the maximal spectrum of $\mathbb{C}[t]$.
\end{exercise}

\begin{proof}
It's clear that
the maximal spectrums of $\mathbb{R}[t]$
and $\mathbb{C}[t]$ have the same
cardinality. Choose any bijection between
them. We claim this is a homeomorphism
by being a closed map. This follows
since closed sets in any maximal spectrum
are the finite sets (and the
image of a finite set under
a bijection is finite).

To check a closed set $F$ is finite,
suppose $\mathfrak{m} \in F$.
The basis of closed sets is given
by the complements of $U(f)$,
which are defined by
$\mathfrak{m}' \in U(f) \iff
f \not\in \mathfrak{m}'$.
Thus for the closed set $F=U(f)^c$
we have that $\mathfrak{m} \in F
\iff f \in \mathfrak{m}$.
Now since $K[t]$ is a PID,
we may assume $\mathfrak{m}=(g)$,
so $f \in \mathfrak{m}$ becomes
$g$ is a factor of $f$.
But $f$ can have only finitely many factors.

To check any finite set is closed, we
only use the maximal ideals
$t-a$ for any $a \in \mathbb{R}$ or
$\mathbb{C}$.
\end{proof}

\noindent
The following proof is wrong: the map
is not surjective.

\begin{proof}
They key observations are: $k[t]$ is an
Euclidean domain for any field, which
implies that it's a PID and a UFD.
The maximal ideals of $\mathbb{C}[t]$
are generated by linear polynomials
and those of $\mathbb{R}[t]$ are generated
either by linear polynomials
$t-a$ with $a \in \mathbb{R}$
or quadratic polynomials
$t^2+at+b$ for $a,b \in \mathbb{R}$ which
have no real roots
(which is obvious because if it had
a real root then it would split;
and using the fundamental theorem
of algebra and the fact that a number is
a root of a
polynomial with real coefficients
if and only if the conjugate is also
a root we can prove that any irreducible
polynomial with real coefficients
must be at most quadratic).

With these conventions, the map is clear:
$$
\begin{aligned}
\varphi: \operatorname{MSpec}(\mathbb{R}[t])
&\longrightarrow \\
\operatorname{MSpec}(\mathbb{C}[t])
(t-a) &\longmapsto(t-a)\\
t^2+at+b& \longmapsto (t-(a+ib)).
\end{aligned}
$$
Being clearly bijective, it remains
to show it is either open or closed.

It's easier to show it is closed.
Let $D(f)$ be open in
$\operatorname{MSpec}(\mathbb{R}[t])$,
so that its complement,
$$
\operatorname{MSpec}(\mathbb{R}[t])\setminus
D(f)
=\{\mathfrak{m} \subset \mathbb{R}[t]
f \in \mathfrak{m}\}
$$
is a basic closed set. 
Since maximal
ideals are principal and generated
by irreducibles, $f \in \mathfrak{m}=(g)$
means that $\mathfrak{m}=(f)$,
so the basic closed sets are sets
consisting of a single ideal
generated by an irreducible element.
Under $\varphi$
we obtain
Then $\varphi(f)=(f)$ is also closed
in $\operatorname{MSpec}(\mathbb{C}[t])$
by the same argument.
\end{proof}

\section{Category of affine varieties}
\label{commutative-algebra-section-category-affine-varieties}

\section{Localization}
\label{commutative-algebra-section-localization}

\noindent
The general recipe for localization is that
$\mathcal{O}_U=\mathcal{O}_X[f^{-1}]$ where
$U=X\setminus V(f)$.
That is, localizing gives the sheaf of the
variety without the zero set.

\begin{definition}
\label{commutative-algebra-definition-localization}
Let $R$ be a ring and $S \subset R$
a multiplicative subset of $R$ (containing
$1$).
The {\it localization} of $R$ at $S$
is the ring $RS^{-1}$
(where $S^{-1}$ is just a set with
the same elements as $S$ but with a $-1$
exponent)
with the operations of
$(rs^{-1})(r_1s_1^{-1})=rr_1s^{-1}s_1^{-1}$,
$rs^{-1}+r_1s_1^{-1}=(rs_1+r_1s)(s^{-1}s^{-1}
_1)$
and the identification given by
$rs^{-1}=r_1s_1^{-1} \iff rs_1=r_1s$.
(Check this definition if $R$ has zero
divisors!)
\end{definition}

\noindent
In a simple case, consider a ring $R$ and $f
\in R$.
The {\it localization of $R$ at $f$}
is $R[t]/(ft-1):=R[f^{-1}]=R(f)$.
The quotient basically says that $t$ is the
inverse of $f$.

\begin{example}
\label{commutative-algebra-example-localization}
\begin{itemize}
\item Dyadic numbers: $\mathbb{Z}[2^{-1}]$,
which are fractions of integers with
denominators which
are powers of $2$.

\item Laurent polynomials:
$\mathbb{C}[t,t^{-1}]$.
\end{itemize}
\end{example}

\begin{lemma}
\label{commutative-algebra-lemma-localization-nilpotent}
If $R[f^{-1}]=0$ then $f$ is nilpotent.
\end{lemma}

\begin{proof}
If $R[f^{-1}]=0$ then $1 \in (ft-1)$,
which means that $1=(1-ft)p(t)$ for some
$p(t) \in k[t]$.
Then…
\end{proof}

\noindent
The localization a ring is not always a local
ring.
But if $S$ is a prime ideal, then it is.


\begin{exercise}
\label{commutative-algebra-exercise-local-noeth-noeth}
Prove that localization of a Noetherian ring
is Noetherian.
\end{exercise}

\begin{proof}
Let $R$ be a Noetherian ring and $S \subset
R$.
Let $I_1\subset I_2 \subset…$
be a chain of ideals in the localization
$R[S^{-1}]$.
Let $\pi:R \to R[S^{-1}]$ the canonical
projection $r \mapsto r$.

Let $I$ be any of the $I_i$.
Consider $\pi^{-1}(I)$. This gives a chain of
ideals in $R$,
which must stabilize.
To go back to the localization simply recall
that
$\pi^{-1}(I)[S^{-1}]=I$.
\end{proof}






\noindent
\section{Noetherian and Artinian rings}
\label{commutative-algebra-section-noetherian-and-artinian-rings}

\noindent 
Noetherianness is exactly the fact
that there are only finitely many
subvarieties each properly put inside the
previous one. That's it.

\medskip\noindent
In this section we explain
why can any affine algebraic variety
be expressed as a union of irreducible
affine algebraic varieties.
Roughly, the proof is as follows.
Hilbert's basis theorem
assures that the ring of polynomials
in $n$ variables over a field is
Noetherian, which guarantees
that there are only finitely many
prime ideals containing any given ideal.
Since the radical of an ideal $I$ is
the intersection of all primes
$\mathfrak{p}_i$ containing $I$,
and the variety of the radical is
the same as the variety of the ideal,
we can express $V(\sqrt{I})
= V(\cap\mathfrak{p}_i)
=\cup V(\mathfrak{p}_i)$.

\begin{definition}
\label{commutative-algebra-definition-noetherian-artinian-rings}
A ring is called {\it Noetherian}
if any increasing chain of
ideals $I_1 \subset I_2 \subset…$
stabilizes. It is called {\it Artinian}
if any descending chain of
ideals $I_1 \supset I_2\supset…$
stabilizes.
\end{definition}

\begin{exercise}
\label{commutative-algebra-exercise-noeth-iff-ideal-finit-gene}
A ring is Noetherian if and only if
all its ideals are finitely generated.
\end{exercise}

\begin{proof}
First we prove that if not every ideal is
finitely generated
then $R$ cannot be Noetherian.
Indeed, if $I$ is not finitely generated
we construct an ascending chain by of ideals
by starting with
any generator of $I$ and adding other
generators one by one.

For the converse suppose that all ideals of
$R$ are finitely generated
and pick an ascending chain $I_1\subset
I_2\subset…$.
Consider $I=\bigcup_i I_i$,
which is an ideal since for any $a \in I$
there is $i$ such that $a \in I_i$
so that $ra \in I_i$ for any $r \in R$.
However, since it is finitely generated,
all the generators have to be in some $I_j$,
and thus the chain stabilizes.
\end{proof}

\begin{theorem}[Hilbert basis theorem]
\label{commutative-algebra-theorem-hilbert-basis}
If $R$ is Noetherian then so is $R[x]$.
\end{theorem}

\begin{exercise}
\label{commutative-algebra-exercise-hilbe-basis-forma-power}
If $A$ is Noetherian then so is $A[[t]]$.
\end{exercise}

\begin{proof}
Pass from an ideal in $A[[t]]$ to $A$ by
taking
the coefficients of the monomials of minimal
degree
of each element in the ideal.
Then take a finite set of generators.
Then go back to $A[[t]]$ by taking the
polynomials
where the generators where taken.
Then…
\end{proof}

\begin{exercise}
\label{commutative-algebra-exercise-radical-intersection-primes}
The radical of an ideal is the
intersection of all primes containing it,
that is,
$$
\sqrt{I}
=\bigcap_
{\mathfrak{p} \supset I}\mathfrak{p}_i.
$$
\end{exercise}

\noindent
Since in a Noetherian ring there
are only finitely many primes
containing any ideal, we
obtain the irreducible decomposition
of affine algebraic varieties.

\begin{exercise}
\label{commutative-algebra-exercise-intersection-submodules-noet}
Let $M$ be an $A$-module, where $A$
is not necessarily Noetherian,
and $N_1,N_2 \subset M$ submodules
which satisfy $N_1 \cap N_2 = 0$.
Assume that the modules $M/N_1$
and $M/N_2$ are Noetherian.
Prove that $M$ is Noetherian.
\end{exercise}

\begin{proof}[Failed proof]
Let $N$ be a submodule (=ideal) of $M$.
Projecting $N$ to $M/N_1$
we obtain a finitely generated module.
Suppose the generators are
$x_i + N_1$, where $x_i \in N$.
The fact that they are generators means
that for any $x \in N$ we have
$x+N_1=\sum a_ix_i + N_1$,
that is, $x-\sum a_ix_i \in N_1$.
The same goes for $N_2$, say,
$x - \sum b_i x_i' \in N_2$ for any
$x \in N$. But then
$\sum a_i x_i - \sum b_i x_i' \in N_1\cap
N_2$, a contradiction (no…)
\end{proof}

\medskip\noindent
Extra exercise:

\begin{exercise}
\label{commutative-algebra-exercise-descending-chains-is-field}
Let $A$ be a ring without zero divisors.
Suppose that any descending chain of
ideals in $A$ stabilizes. Prove
that $A$ is a field.
\end{exercise}

\begin{proof}
Suppose that $A$ is not a field
and let's show that there's a descending
chain of ideals which does not stabilize.
Let $a$ be not invertible (since $A$
is not a field).
Then 
$$
(a)\subset (a^2) \subset \subset (a^3)
\subset…
$$
does not stabilize. Indeed,
if $(a^n)=(a^{n+1})$ we obtain
$a^n=ba^{n+1}\iff a^n(1-ba)=0
\implies  1=ba$, since $A$ is a domain,
which is a contradiction because
$a$ is not invertible.
\end{proof}

\section{Tensor product}
\label{commutative-algebra-section-tensor-product}

\begin{definition}
\label{commutative-algebra-definition-tensor-product}

\end{definition}

\begin{theorem}[Universal property of tensor
product]
\label{commutative-algebra-theorem-unive-prope-tenso-produ}
Let $M_1,M_2,M$ be modules over a ring $R$.
For any bilinear map $B:M_1\times M_2 \to M$
there exists
a unique homomorphism $b:M_1 \otimes M_2 \to
M$
making the following diagram commute
$$
\xymatrix{
M_1\times M_2 \ar[r]^\pi\ar[dr]_B &  M_1
\otimes M_2 \ar[d]^b\\
& M
}
$$
where $\pi(m_1,m_2)=m_1 \otimes m_2$.
\end{theorem}

\begin{remark}
\label{commutative-algebra-remark-whitneys-definition}
For $R=k$ the above property shows that
$\text{Bil}(M_1\times M_2,k)=(M_1 \otimes
M_2)^*$.
If we assume that the vector spaces are
finite dimensional,
we can take duals to obtain
$$
M_1 \otimes M_2 = (M_1 \otimes M_2)^ {*
*}=\text{Bil}(M_1 \times M_2,k)^*,
$$
which is how Whitney first defined tensor
products.
\end{remark}

\noindent
Here's Misha's way to show that the universal
property
of tensor product implies its uniqueness.
Consider the category $\mathcal{C}$ of pairs
of an $R$-module $M$,
and a bilinear map $M_1 \times M_2 \to M$.
Then we prove that the tensor product $M_1
\otimes M_2$
is an initial object in such category,
and it is easy to show that initial objects
are unique.

\begin{remark}
\label{commutative-algebra-remark-universal-property}
Misha's alternative way of stating the
universal property of tensor product is
that there is a bijective correspondence
$$
\Hom_R(M_1 \otimes_R M_2, M)
\leftrightarrow
\text{Bil}(M_1\times M_2,M).
$$
The key observation is that the LHS are {\it
linear} maps from a module to $M$
while the RHS are {\it bilinear} maps.
\end{remark}

\begin{exercise}[Currying]
\label{commutative-algebra-exercise-currying}
$$
\Hom(M_1\otimes_R M_2, M)
=
\Hom(M_1,\Hom(M_2,M))
$$
\end{exercise}

\begin{proof}
This is implied by the universal property
since the RHS is just $\text{Bil}(M_1\times
M_2,M)$.
\end{proof}

\begin{exercise}
\label{commutative-algebra-exercise-zero-tensor}
Find non-zero $R$-modules $A$ and $B$
such that
$$
A \otimes_R B=0
$$
in each of the following cases.
\begin{enumerate}
\item $R=\mathbb{Z}$.
\item $R=C^\infty(M)$, where $M$ is a
smooth manifold with at least two points.
\item $R=\mathbb{C}[t]$.
\end{enumerate}
\end{exercise}

\begin{proof}
For $R=\mathbb{Z}$, take
$A=\mathbb{R}$ and $B=\mathbb{Z}/2$.
For every $r\in \mathbb{R}$ we can write
$r=2s$ for some $s\in \mathbb{R}$.
Thus
$$
r\otimes \bar n=s\otimes 2\bar n=0.
$$
Every tensor is a finite sum of pure tensors,
so $A\otimes_R B=0$.

For $R=C^\infty(M)$, choose distinct points
$p,q\in M$ and take
$A=\mathcal{O}_{M,p}$ and
$B=\mathcal{O}_{M,q}$.
Let $a\otimes b$ be a pure tensor.
Choose $f\in C^\infty(M)$ which is zero
near $p$ and equal to $1$ near $q$.
Then $fa=0$ in $A$ and $fb=b$ in $B$.
Thus
$$
a\otimes b=a\otimes fb=fa\otimes b=0.
$$
Again $A\otimes_R B=0$.

For $R=\mathbb{C}[t]$, take
$A=R/(t)$ and $B=R/(t-1)$.
In $A$ we have $t=0$, and in $B$ we have
$t=1$.
For every pure tensor $a\otimes b$,
$$
a\otimes b=a\otimes tb=ta\otimes b=0.
$$
Hence $A\otimes_R B=0$.
\end{proof}


\begin{exercise}
\label{commutative-algebra-exercise-tensor-product-fg-or-not}
Let $R_1, R_2$ be rings over $\mathbb{C}$
such that $R_1 \otimes_{\mathbb{C}}R_2$
is finitely generated. Prove that
$R_1,R_2$ are finitely generated,
or find a counterexample.
\end{exercise}

\begin{proof}
Suppose 
$R_1 \otimes_{\mathbb{C}}R_2
= \mathbb{C}[z_1,…,z_n]$
and
$z_j=\sum_{k}a_{jk} \otimes b_{jk}$.
Let $A\subset R_1$ be the algebra
generated by the $a_{jk}$.
It's clear that the generators
$z_j$ are in $A\otimes_{\mathbb{C}}R_2$.
Thus, any polynomial on the $z_j$
lies in $A \otimes_{\mathbb{C}}R_2$.
That is, $R_1 \otimes_{\mathbb{C}}R_2
=A \otimes_{\mathbb{C}}R_2$.

Let $\lambda:R_2 \to \mathbb{C}$
be a linear functional with $\lambda(1)=1$.
Consider the operator
$\operatorname{id} \otimes \lambda:
A \otimes R_2 \to A$.

Now take $r_1 \in R_1$.
Then $r_1 \otimes 1 \in 
R_1 \otimes_{\mathbb{C}}R_2
=A \otimes_{\mathbb{C}}R_2$.
Therefore
$r_1=\operatorname{id}\otimes \lambda
(r_1 \otimes 1) \in A$,
and $A$ is finitely generated.
\end{proof}

\begin{exercise}
\label{commutative-algebra-exercise-left-hom-sends-right-exact}
Let
$$
\xymatrix{
M_1\ar[r]^{f_1}
&M_2\ar[r]^{f_2}
&M_3\ar[r]
&0
}
$$
be exact.
Show that
$$
\xymatrix{
0\ar[r]
&\Hom(M_3,N)\ar[r]^{-\circ f_2}
&\Hom(M_2,N)\ar[r]^{- \circ f_1}
&\Hom(M_1,N)
}
$$
is exact.
\end{exercise}

\begin{proof}
Injectivity of the first arrow is easy.
For the middle term pick $f \in
\Hom_R(M_2,N)$
in the kernel of $- \circ f_1$.
To show that it is in the image of $-\circ
f_2$
we look for $g \in \Hom_R(M_3,N)$
such that $f=g\circ f_2$.

Since $M_1 \to M_2 \to M_3 \to 0$
is exact, we have an isomorphism
$\bar{f}_2:M_2/\text{img}f_1
\xrightarrow{\simeq} M_3$.
Since $f \circ f_1=0$, we see that $f$
induces a map $\bar{f}$ on the quotient since
Indeed, if $m_2-m_2' \in \text{img}f_1$,
we have $0=f(m_2-m_2')=f(m_2)-f(m_2')$.

Then for any $m_3 \in M_3$
we have a unique class $m_2 + M_1 \in
M_2/\text{img}f_1$.
Thus we can define $g = \bar{f} \circ
\bar{f}_2^{-1}$.

Note that this choice satisfies $f = g \circ
f_2$.
Indeed,
$$
g\circ f_2(m_2)=\bar{f} \circ
\bar{f}^{-1}_2(m_3)
= \bar{f} \circ (m_2+M)
=f(m_2).
$$
\end{proof}

\begin{exercise}
\label{commutative-algebra-exercise-count-right-hom-not-right}
Let $R = \mathbb{Z}$.
Find an example of an injective homomorphism
$M_1 \to M_2$ such that $\Hom(M_2,N) \to
\Hom(M_1,N)$
is not surjective.
\end{exercise}

\begin{proof}
Consider
$$
\xymatrix{
0\ar[r]
&2\mathbb{Z}\ar[r]^i\ar[d]_\varphi
&\mathbb{Z}\ar@{-->}[dl]^{\not\exists}\\
&\mathbb{Z}/2\mathbb{Z}
}
$$
where
$$
\begin{aligned}
\varphi: 2\mathbb{Z} &\longrightarrow
\mathbb{Z}/2\mathbb{Z} \\
2n &\longmapsto [n].
\end{aligned}
$$

\noindent
Then there is no extension
$\bar{\varphi}:\mathbb{Z} \to
\mathbb{Z}/2\mathbb{Z}$
making $\varphi=\bar{\varphi} \circ i$.
Indeed, note that $\varphi(2)=[1]$,
while $\bar{\varphi}(2\cdot
1)=2\bar{\varphi}(1)=[0]$.
\end{proof}

\noindent
The idea is that any element in the image
must be divisible
by any integer. This is why $\mathbb{Q}$ is
an injective module
according to the following definition.
(A precise proof of why $\mathbb{Q}$ is
injective
can be found at
\href{https://stacks.math.columbia.edu/tag/%
01D6}{01D6}.)

\begin{definition}
\label{commutative-algebra-definition-injective-module}
An $R$-module is {\it injective} if the
functor $\Hom(-,M)$
is exact.
That is, if the dotted arrow in the following
diagram always exists:
$$
\xymatrix{
0\ar[r]
&M_1\ar@{^{(}->}[r]\ar[d]
&M_2\ar@{-->}[dl]^{\exists}\\
&N
}
$$
\end{definition}

\begin{exercise}
\label{commutative-algebra-exercise-left-hom-is-left-exact}
Let
$$
\xymatrix{
0\ar[r]
&M_1\ar[r]^{f_1}
&M_2\ar[r]^{f_2}
&M_3\ar[r]
&0
}
$$
be an exact sequence of $R$-modules.
Prove that
$$
\xymatrix{
0\ar[r]
&\Hom_R(N,M_1)\ar[r]^{f_1 \circ -}
&\Hom_R(N,M_2)\ar[r]^{f_2 \circ -}
&\Hom_R(N,M_3)
}
$$
is exact
\end{exercise}

\begin{proof}
Injectivity is simple.
For the middle term let $f:N \to M_2$ be in
the kernel
of $f_2 \circ -$,
that is, $f(N)$ is killed by $f_2$.

We want to define a map $g:N \to M_1$
such that $f_1 \circ g = f$.
But observe that the image of $f$ already
lies in $M_1 = \Ker f_2$.
Thus $f$ is already a map $f:N \to M_1$,
save a more precise notation.
\end{proof}

\begin{exercise}
\label{commutative-algebra-exercise-count-left-hom-not-right}
Let $R=\mathbb{Z}$. Find an epimorphism
(surjective homomorphism)
of $R$-modules $M_2 \to M_3$ such that the
corresponding map
$$
\Hom_R(N,M_2) \to \Hom_R(N,M_3)
$$
is not surjective for some $R$-module $N$.
\end{exercise}

\begin{proof}
Now consider
$$
\xymatrix{
&\mathbb{Z}
\ar[d]^\pi\ar@{-->}[dl]_{\not\exists}\\
2\mathbb{Z}\ar[r]_-\varphi&
\mathbb{Z}/2\mathbb{Z}\ar[r]& 0
}
$$
for the same $\varphi$ as in Exercise
\ref{commutative-algebra-exercise-count-right-hom-not-right}.
Then $\pi(1)=[1]$, while any extension
mapping $1 \to 2n$
gives $\varphi(2n)=2\varphi(n)=[0]$.
\end{proof}

\begin{definition}
\label{commutative-algebra-definition-projective-module}
An $R$-module $N$ is {\it projective} if
the functor $\Hom_R(N,-)$ is exact.
That is, if the dotted
arrow in the following diagram always exists:
$$
\xymatrix{
&N \ar[d]\ar@{-->}[dl]_{\exists}\\
M_2\ar@{->>}[r]&M_3\ar[r]& 0
}
$$

\end{definition}

\begin{exercise}
\label{commutative-algebra-exercise-finit-gener-proje-iff-fact}
A finitely generated $R$-module $N$ is
projective
if and only if
it is a direct summand of a free $R$-module
$R^n$,
i.e. there is a direct sum decomposition
$R^n=N \oplus N'$.
\end{exercise}

\begin{proof}
Consider $R^n \simeq \bigoplus_{i=1}^n a_iR$
where $N$ is generated by $a_1,…,a_n$.
Apply projectivity to obtain
$$
\xymatrix{
&N \ar[d]^{\text{id}}\ar@{-->}[dl]_{\exists
\psi}\\
R^n\ar[r]_-\varphi& N \ar[r]& 0
}
$$
where
$$
\begin{aligned}
\varphi: R^n &\longrightarrow N \\
(r_1a_1,…,r_na_n) &\longmapsto \sum r_ia_i.
\end{aligned}
$$

\noindent
This is a section of $\varphi$ in the exact
sequence
$$
\xymatrix{
0\ar[r]
&\Ker \varphi\ar[r]
&R^n\ar[r]^{\varphi}
&N\ar[r]\ar@/_-.5pc/[l]^\psi
&0
}
$$
which gives decomposition
$$
R^n \simeq \Ker\varphi \oplus \text{Im}\psi.
$$
More explicitly, the intersection of those
modules is clearly trivial
while any element of $R^n$ may be written as
$$
x=\underbrace{(x-\psi \varphi(x))}_{\in \Ker
\varphi}
+\underbrace{\psi\varphi(x)}_{\in
\text{Im}\psi}.
$$
For the converse we start from the diagram
given in the definition
of projective module, add $N'$ to every term,
arriving
at $R^n\simeq N\oplus N'$. We observe that
free modules
are projective by constructing an extension
using a basis of $R^n$,
and then back to $N$ by taking the action of
the extension restricted to $N$.
\end{proof}

\noindent
The upshot about projective modules
is precisely that they can't be
direct summands of a free module.
Thus $\mathbb{Z}_2$ is not projective,
since it has torsion, and
so it cannot be a summand of $\mathbb{Z}^n$.

Thus, we create a non-projective module
by quotienting the ring by any ideal
- an element of the ideal is nilpotent,
i.e. is a torsion element.

Further, it looks like projective modules
are flat (confirm this!)

\medskip\noindent
In fact, currying (Exercise
\ref{commutative-algebra-exercise-currying})
is the essential ingredient
for the proof that tensor product is left
exact.
Such result follows from two lemmas.

\begin{lemma}
\label{commutative-algebra-lemma-tenso-produ-funct-left-exact}
Let
$$
\xymatrix{
0\ar[r]
&M_1\ar[r]
&M_2\ar[r]
&M_3\ar[r]
&0
}
$$
be an exact sequence of $R$-modules.
Then for any $R$-modules $N$ and $N'$,
the sequence
$$
\xymatrix{
0\ar[r]
&\Hom_R(M_3 \otimes N',N)\ar[r]
&\Hom_R(M_2 \otimes N',N)\ar[r]
&\Hom_R(M_1 \otimes N',N)
}
$$
is exact.
\end{lemma}

\begin{proof}
Apply $\Hom_R(-,N)$,
then $\Hom_R(N',-)$
and then currying.
\end{proof}

\noindent
There is an (almost) converse
to Exercise
\ref{commutative-algebra-exercise-left-hom-sends-right-exact}:

\begin{lemma}
\label{commutative-algebra-lemma-converse-hom-is-left-exact}
Let
$$
\xymatrix{
E:&
M_1\ar[r]^\mu
&M_2\ar[r]^\rho
&M_3\ar[r]
&0
}
$$
be a complex,
i.e. $\rho\mu=0$,
such that
$$
\xymatrix{
0\ar[r]
&\Hom_R(M_3,N)\ar[r]^{\rho_N}
&\Hom_R(M_2,N)\ar[r]^{\mu_N}
&\Hom_R(M_1,N)
}
$$
is exact for any $R$-module $N$. Then $E$ is
also exact.
\end{lemma}

\noindent
And this is what implies right-exactness of
tensor product:

\begin{theorem}
\label{commutative-algebra-theorem-tensor-product-is-right-exact}
Let
$$
\xymatrix{
M_1\ar[r]
&M_2\ar[r]
&M_3\ar[r]
&0
}
$$
be exact. Then
$$
\xymatrix{
M_1 \otimes_R M\ar[r]
&M_2 \otimes_R M\ar[r]
&M_3 \otimes_R M\ar[r]
&0
}
$$
is exact.
\end{theorem}

\begin{proof}
This is just Lemma
\ref{commutative-algebra-lemma-tenso-produ-funct-left-exact}
and Lemma
\ref{commutative-algebra-lemma-converse-hom-is-left-exact}!!!
\end{proof}

\begin{exercise}
\label{commutative-algebra-exercise-not-flat}
Let $R=\mathbb{C}[t]$.
Find an exact sequence of $R$-modules
$$
\xymatrix{
0\ar[r]
&M_1\ar[r]
&M_2\ar[r]
&M_3\ar[r]
&0
}
$$
and an $R$-module $N$
such that the corresponding map
$$
M_1 \otimes_R N \to M_2 \otimes_R N
$$
is not injective.
\end{exercise}

\begin{proof}
Suppose the given exact sequence is just the
ideal
exact sequence of a point in $\mathbb{A}^1$.
Then tensoring by $M$ means thickening the
point.
Not having injectivity after tensoring
means that tensoring the ideal by $M$
does not give the ideal of the resulting
variety.
\end{proof}

And from this it follows as corollary

\begin{proposition}
\label{commutative-algebra-proposition-tensor-ideals}
Let $I \subset R$ be an ideal of a ring.
Then $M \otimes_R (R/I)=M/IM$.
\end{proposition}

\begin{proof}
Take the exact sequence $0 \to I \to R \to
R/I$ and tensor
by $M$.
\end{proof}

\begin{exercise}
\label{commutative-algebra-exercise-tensor-product-maximal-ideal}
Let $I,I'$ be distinct maximal ideals in a
ring $R$.
Then $R/I \otimes_R R/I'=0$.
\end{exercise}

\begin{proof}
They key observation is that since $I$ and
$I'$
are maximal we must have $I+I'=M$.
This means that $1=f+f'$ for some $f \in I$
and $f' \in I'$.
Then
$1(a+I)\otimes(b+I')
=(fa+I) \otimes (b + I')
+(a+I)\otimes(fb+I')=0$.
\end{proof}

\medskip\noindent
Now we discuss tensor product in the category
of rings over a ring.

\begin{definition}
\label{commutative-algebra-definition-tensor-product-of-rings}
Let $A,B$ and $C$ be rings, and
$C \to A$ and $C \to B$ homomorphisms.
Considering $A$ and $B$ as $C$-modules,
we can define a product on $A \otimes_C B$ by
$$
a \otimes b \cdot a' \otimes b':=a a' \otimes
b b'.
$$
By the universal property this indeed extends
to a product.
\end{definition}

\begin{example}
\label{commutative-algebra-example-tensor-product-of-fg-algebras}
$$
\mathbb{C}[t_1,…,t_k] \otimes_\mathbb{C}
\mathbb{C}[z_1,…,z_n].
$$
\end{example}

\begin{example}
\label{commutative-algebra-example-tenso-produ-rings-ring}
For any homomorphism $\varphi:C \to A$,
the ring $A \otimes_C(C/I)$ is isomorphic to
$\frac{A}{\varphi(I)A}$
by Proposition
\ref{commutative-algebra-proposition-tensor-ideals}.
\end{example}

\noindent
Recall that the spectrum of a ring $\Spec R$
is the corresponding algebraic variety.

\begin{proposition}
\label{commutative-algebra-proposition-preimage-of-point}
Let $f:X \to Y$ be a morphism of affine
varieties
and $f^*:\mathcal{O}_Y \to \mathcal{O}_X$ the
corresponding morphism.
Let $y \in Y$ and $\mathfrak{m}_y$ the
corresponding maximal ideal.
Then
$$
\Spec\left(
\left(\mathcal{O}_X \otimes_{\mathcal{O}_Y}
\frac{\mathcal{O}_Y}{\mathfrak{m}_y}\right)
\Big/\text{nilradical}\right)
\cong f^{-1}(y).
$$
\end{proposition}

\begin{remark}
\label{commutative-algebra-remark-preimage-as-fibered-product}
Notice that we can see $y \in Y$ as a
subvariety.
In fact, it is, like $X$, a variety {\it over
$Y$}.
Thus this expression for $f^{-1}(y)$
matches that of the fibered product
in Proposition
\ref{commutative-algebra-proposition-fibered-product}.
\end{remark}

\begin{lemma}
\label{commutative-algebra-lemma-produ-affin-varie-tenso-algeb}
Let $A,B$ be finitely generated reduced
$\mathbb{C}$-algebras.
If $R:=A \otimes_{\mathbb{C}}B$, and $N
\subset R$
nilradical (the set of nilpotent elements of
$R$).
Then
$$
\Spec(R/N) = \Spec(A) \times \Spec(B).
$$

\end{lemma}

\medskip\noindent
Now we discuss…

\begin{lemma}
\label{commutative-algebra-lemma-inter-maxim-ideal-nilra}
Let $A$ be a finitely-generated
$\mathbb{C}$-algebra.
The intersection of all the prime ideals of
$A$
is its nilradical.
\end{lemma}

\begin{exercise}
\label{commutative-algebra-exercise-tenso-produ-may-conta-nilp}
Let $A$ and $B$ be finitely generated
$\mathbb{C}$-algebras,
$B \to A$ a homomorphism and $\mathfrak{m}
\subset B$ a maximal ideal.
Then $A\otimes_B (B/\mathfrak{m})$ can
contain nilpotents.
\end{exercise}

\begin{proof}
The geometric idea is to project a parabola
to an axis.
At the point of tangency we obtain a fat
point,
whose algebra is not reduced.
\end{proof}

However, over $\mathbb{C}$ the situation is
different:

\begin{theorem}
\label{commutative-algebra-theorem-tenso-produ-over-c-reduc}
Let $A,B$ be finitely generated reduced
$\mathbb{C}$-algebras.
Then $A \otimes_\mathbb{C} B$ is reduced.
\end{theorem}

It follows that

\begin{theorem}
\label{commutative-algebra-theorem-tenso-produ-affin-varie}
$$
\Spec(A)\times \Spec(B) = \Spec(A
\otimes_\mathbb{C}B).
$$
\end{theorem}

\medskip\noindent
Now we can take preimages!

\begin{lemma}
\label{commutative-algebra-lemma-preimage-of-affine-subvariety}
Let $f:X \to Y$ be a morphism of affine
algebraic varieties,
$f^*:\mathcal{O}_Y \to \mathcal{O}_X$ the
corresponding ring homomorphism,
$Z \subset Y $ a subvariety and $I_Z$ its
ideal.
Denote by $R_1$ the quotient ring
$$
(\mathcal{O}_X
\otimes_{\mathcal{O}_Y}(\mathcal{O}_Y/IZ)/
\text{nilradical}
= (\mathcal{O}_X/f^*
(I_Z))/\text{nilradical}.
$$
Then
$$
\Spec(R_1)=f^{-1}(Z).
$$
\end{lemma}

\begin{proof}
Analogous to preimage of a point.
\end{proof}

\noindent
The slogan is:
adding polynomials to an ideal is the same as
taking tensor product.
This operation is given by Exercise
\ref{commutative-algebra-exercise-tensor-product-maximal-ideal}.

We can also consider the diagonal of a
product:

\begin{lemma}
\label{commutative-algebra-lemma-diagonal-of-product}
Let $M$ be an affine algebraic variety,
its diagonal $\Delta \subset M\times M$
and $I \subset \mathcal{O}_M
\otimes_\mathbb{C} \mathcal{O}_M$
the ideal generated by $r \otimes 1 - 1
\otimes r$
for all $r \in \mathcal{O}_M$. Then
$$
\mathcal{O}_M \otimes_\mathbb{C}
\mathcal{O}_M/I.
$$
\end{lemma}

\noindent
And now we introduce fibered product!

\begin{proposition}
\label{commutative-algebra-proposition-fibered-product}
Let $X \to M$ and $Y \to M$ be morphisms of
algebraic varieties,
$R= \mathcal{O}_X
\otimes_{\mathcal{O}_M}\mathcal{O}_Y$,
and $R_1$ the quotient of $R$ by its
nilradical.
Then
$$
\Spec(R_1)=X \times_M Y.
$$
\end{proposition}

\begin{proof}
$$
X\times Y \xrightarrow{f \times g} Z \times
Z,
$$
and compute $(f\times g)^{-1}(\Delta)$.
By Lemma \cite{lemma-diagonal} we must have
$$
A \otimes_\mathbb{C} B \big/(r \otimes 1-1
\otimes r)
=
A \otimes_C B
$$
where $A=\mathcal{O}_X$, $B=\mathcal{O}_Y$
and $C=\mathcal{O}_Z$.
\end{proof}

\begin{exercise}
\label{commutative-algebra-exercise-coproduct-in-rings}
Coproduct in the category of
reduced finitely generated algebras
is tensor product.
\end{exercise}

\section{Flat modules}
\label{commutative-algebra-section-flat-modules}

\noindent
Let $S=A\setminus \mathfrak p$. Then
$$
M_{\mathfrak p}=0
$$
means that for every $m\in M$,
there exists $s\in S$ such that
$$
sm=0.
$$
In other words, every element of $M$
is killed by one of the denominators
allowed in the localization.

If $A$ is a domain, this is the usual
notion of torsion: a nonzero element
$m\in M$ is torsion if there exists
$0\neq a\in A$ such that
$$
am=0.
$$
This is the module analogue of a zero
divisor in a ring.

Flatness implies torsion-freeness over
a domain. Indeed, let $0\neq a\in A$.
Then multiplication by $a$ gives an
injective map
$$
A \xrightarrow{\cdot a} A.
$$
If $M$ is flat, tensoring with $M$
preserves injections, so we get an
injective map
$$
A\otimes_A M
\longrightarrow
A\otimes_A M.
$$
Using the identification
$$
A\otimes_A M \cong M,
\qquad
x\otimes m \mapsto xm,
$$
this map is exactly
$$
M \xrightarrow{\cdot a} M.
$$
Therefore
$$
am=0
\quad\Longrightarrow\quad
m=0.
$$
Thus $M$ has no torsion elements.

The converse is false in general:
torsion-free does not imply flat.
For example, over
$$
A=k[x,y],
$$
the ideal
$$
M=(x,y)
$$
is torsion-free, since it is contained
in the domain $A$, but it is not flat.

\begin{definition}
\label{commutative-algebra-definition-flat-module}
Let $R$ be a ring.
An {\it $R$-module} $M$ is called
{\it flat}
if the functor
$$
X \mapsto X \otimes_R M
$$
is exact.
\end{definition}

\begin{lemma}
\label{commutative-algebra-lemma-free-tensor-copies}
For every $R$-module $N$,
$$
N\otimes_R R^{\oplus r}\cong N^{\oplus r}.
$$
Consequently finite free modules are flat.
\end{lemma}

\begin{proof}
If $e_1,\ldots,e_r$ is the standard
basis, the maps are
$$
\begin{aligned}
n\otimes(b_1,\ldots,b_r)
&\longmapsto(b_1n,\ldots,b_rn),\\
(n_1,\ldots,n_r)
&\longmapsto\sum_i n_i\otimes e_i.
\end{aligned}
$$
They are inverse by bilinearity.
Tensoring with $R^{\oplus r}$ therefore
makes $r$ copies of a short exact
sequence, which remains exact.
The right-hand side is a direct sum,
not a tensor power of $N$.
\end{proof}

Tensoring is always right exact; flatness
is what also preserves the first
injection. For instance, tensor
$0\to R\xrightarrow{x}R\to R/(x)\to0$
over $R=\mathbb{C}[x]$ with $R/(x)$.
The first map becomes multiplication by
$x$ on $R/(x)$, hence the zero map.
It is no longer injective.
Tensoring individual maps is always
defined, regardless of flatness.

\begin{definition}
\label{commutative-algebra-definition-torsion-module}
Let $A$ be a ring without zero divisors,
and $k(A)$ its fraction field.
The {\it torsion} of an $A$-module $M$
is the kernel of the natural map
$$
M \to M \otimes_A k(A).
$$
\end{definition}

\begin{exercise}
\label{commutative-algebra-exercise-flat-modules-torsion-free}
Let $M$ be a flat $R$-module,
where $R$ is a ring without zero divisors.
Prove that $M$ is torsion-free.
\end{exercise}

\begin{exercise}
\label{commutative-algebra-exercise-union-flat-modules-flat}
Let
$
M_1 \subset M_2 \subset \ldots
$
be a sequence of embedded $R$-modules.
Assume that all $M_i$ are flat.
Prove that
$
\bigcup_i M_i
$
is also flat.
\end{exercise}

\begin{exercise}
\label{commutative-algebra-exercise-flat-finite-local-free}
Prove that a finitely generated module
over a Noetherian local ring is flat
if and only if it is free.
\end{exercise}

\begin{exercise}
\label{commutative-algebra-exercise-local-flat-implies-flat}
Let $A$ be a ring without zero divisors,
and $M$ an $A$-module
which is flat over all local rings
$A_\mathfrak{p} \supset A$
obtained by localizing $A$
at prime ideals.
Prove that $M$ is flat.
\end{exercise}

\begin{proof}
Let
$$
\xymatrix{
0\ar[r]
&M_1\ar[r]
&M_2\ar[r]
&M_3\ar[r]
&0
}
$$
be exact. Then consider
$$
\xymatrix{
0\ar[r]
&M_1 \otimes_A M\ar[r]^{\varphi}
&M_2 \otimes_A M\ar[r]
&M_3 \otimes_A M\ar[r]
&0.
}
$$
We need to show that the map
$\varphi: M_1 \otimes_A M 
\to M_2 \otimes_A M$ is injective.
So basically the claim is:
if that map is injective
over all localizations $A_{\mathfrak{p}}$
then it is injective over $A$.

So it's basically two steps:
(1) show that
\begin{equation}
\label{commutative-algebra-equation-localized-varphi}
(\Ker \varphi)_{\mathfrak{p}}
=
\Ker \varphi^{\mathfrak{p}}
\end{equation}
where $\varphi^{\mathfrak{p}}$
is the induced map over $A_{\mathfrak{p}}$
for any $\mathfrak{p} \subset A$
prime. More precisely:
$$
\xymatrix{
0\ar[r]
&(M_1)_{\mathfrak{p}} 
\otimes_{A_{\mathfrak{p}}}M_{\mathfrak{p}}
\ar[r]^{
\varphi^{\mathfrak{p}}}
&
(M_2)_{\mathfrak{p}} 
\otimes_{A_{\mathfrak{p}}}
M_{\mathfrak{p}}\ar[r]
&(M_3)_{\mathfrak{p}}
\otimes_{A_{\mathfrak{p}}}
M_{\mathfrak{p}}\ar[r]
&0
}
$$
where the localizations of each module
are given by tensoring with $A_{\mathfrak{p}}$
over $A$, e.g. $M_{\mathfrak{p}} =
A_{\mathfrak{p}} \otimes_A M$ by definition.
Thus the above sequence is exactly
$$
\xymatrix{
0\ar[r]
&(M_1 \otimes_A M)_{\mathfrak{p}}\ar[r]
^{\varphi^{\mathfrak{p}}}
&(M_2 \otimes_A M)_{\mathfrak{p}}\ar[r]
&(M_3 \otimes_A M)_{\mathfrak{p}}\ar[r]
&0
}
$$
and we see that
Equation \ref{commutative-algebra-equation-localized-varphi}
follows by definition.
Indeed, both sets are barely
$$
\left\{
\frac{a}{b}x:\varphi(x)=0,
\frac{a}{b}\in A_{\mathfrak{p}}\right\}.
$$


By the exercise's hypothesis, we know that
$\Ker \varphi^{\mathfrak{p}}=0$.

Then, step (2) is:
show that if any $A$-module $K$
is zero as an $A_{\mathfrak{p}}$-module
for any $\mathfrak{p}\subset A$,
then it is zero as an $A$-module.
Here's AI's proof: consider the annihilator
of an element $m \in M$, it's the
set of scalars that kill it, i.e.
$\operatorname{Ann}(m) =
\{a \in A: am=0\}$. It's nonempty by
hypothesis (yes, the fact that $m$ is
zero over $A_{\mathfrak{p}}$ literally
means that there is $a \not\in \mathfrak{p}$
such that $a m =0$).
This just says that $\operatorname{Ann}(m)
\not\subset \mathfrak{p}$ for any
$\mathfrak{p}$.
But that's not true: any ideal is
contained in some maximal (and thus prime)
ideal. The reason for that is Zorn's lemma:
the set of ideals containing
$\operatorname{Ann}(m)$ is inductive
(i.e. every chain is bounded, in this case
by the union of the ideals), so it has
a maximal element, and maximal ideals
are prime because the quotient is a field,
which is in particular a domain.
 
\end{proof}

\begin{definition}
\label{commutative-algebra-definition-boolean-ring}
A {\it boolean ring} is a ring such
that all its elements are idempotents,
that is, satisfy $a^2=a$.
\end{definition}

\begin{exercise}
\label{commutative-algebra-exercise-boolean-ring-flat}
Prove that any finitely generated module
over a Boolean ring is flat.
\end{exercise}

\begin{proof}
The claim is false.
Consider the boolean ring
$\mathbb{Z}/2\mathbb{Z} \oplus
\mathbb{Z}/2\mathbb{Z}$,
equipped with the usual ring structure
for direct sums which is given
by applying the operations entry by entry.

This module has torsion:
$$
([1],[0])([0],[1])=([0],[0]).
$$
Ah! But this ring is not a domain. Forget it.
\end{proof}

\section{Artinian algebras}
\label{commutative-algebra-section-artinian-algebras}

\noindent
The aim of this section
is to prove that a commutative
algebra over a field
which is finite-dimensional as a
vector space and has no nilpotents
(i.e. is semisimple)
is a direct sum of fields;
namely Theorem
\ref{commutative-algebra-theorem-semisimple-decomposition}.
The key idea is that in a
finite-dimensional algebra
we can use Chinese Reminder Theorem
applied to the irreducible
decomposition of the minimal
polynomial of a non-invertible
element to construct an idempotent
element.
This equivalent to the claim that
if there are no idempotents
the algebra is a field.
Then the decomposition into fields
is just a decomposition
of unity into indecomposable
idempotents, each of which
gives a summand $eR$ with no
idempotents, i.e. a field.

This decomposition is used to prove
that the preimage of a point by a
finite map of
affine algebraic varieties
is a finite set of points,
Theorem
\ref{commutative-algebra-theorem-finit-map-has-finit-preim}.
Moreover, it also used quite a lot in the
development of Galois theory in
Section \ref{commutative-algebra-section-galois-theory}.


\medskip\noindent
In this section algebras are not
assumed to have unity.

\begin{definition}
\label{commutative-algebra-definition-minimal-polynomial}
The {\it minimal polynomial} of an element
$v$
in a finitely-generated $k$-algebra $R$
a monic polynomial $P \in R[t]$
of smallest possible degree having $v$ as a
root.
\end{definition}

\begin{lemma}
\label{commutative-algebra-lemma-subal-minim-polyn}
The subalgebra generated by $v$ is
isomorphic to $k[t]/(P)$.
\end{lemma}

\begin{proof}
Consider the map $k[t] \to Rv$,
whose kernel is $(P)$ since $P$ divides
any polynomial vanishing in $v$.
\end{proof}

\begin{example}
\label{commutative-algebra-example-subalgebras}
$$
\frac{\mathbb{R}[t]}{t^2+1}\simeq \mathbb{C}
$$
while
$$
\frac{\mathbb{R}[t]}{t^2=1}\simeq \mathbb{R}
\oplus \mathbb{R}.
$$


\end{example}

\medskip\noindent
Recall that $k[t]$ is a UFD.
Also recall that $P$ is an irreducible
element of $k[t]$ if and only if
$k[t]/(P)$ is a field.
This follows since $k[t]/(P)$ cannot
have zerodivisors if $P$ is prime,
which is equivalent to irreducible
in UFDs.
These extensions are called
{\it primitive}.
[They are extensions of $k$
obtained by adding a root of $P$.]

\begin{definition}
\label{commutative-algebra-definition-artinian-algebra}
A commutative, associative $k$-algebra $R$
not necessarily with unit
is called an {\it Artinian algebra} if it is
a finite-dimensional vector space over $k$.
It is called {\it semisimple}
if it has no nonzero nilpotents.

A commutative algebra with unity over a field
is called an {\it Artinian ring} if it is
finite-dimensional
as a vector space over $k$.
\end{definition}

\noindent
We will show that any semisimple
Artinian algebra
is a direct sum (direct sum in
the category of algebras is
the direct sum of the rings
equipped with natural operations)
of fields, and this decomposition
is uniquely defined.

An idempotent element $e$ in $R$,
i.e. $e^2=e$,
gives us a canonical decomposition
$$
eR \oplus (1-e)R
$$
since $1-e$ is also idempotent.

\begin{exercise}
\label{commutative-algebra-exercise-minimal-polynomial}
Let $v \in R$ be an element of a
finite-dimensional algebra $R$ over $k$.
Let $k[v]$ be the subalgebra generated by $1,
v, v^2,…$.
Prove that $P(v)=0$ for a unique monic
polynomial $P$ of degree
the dimension of $k[v]$.
\end{exercise}

\begin{proof}
Division algorithm.
\end{proof}

\noindent
The polynomial above is called the {\it
minimal polynomial} of $v$.

\begin{exercise}
\label{commutative-algebra-exercise-minim-polyn-resid}
Let $v \in R$ be an element of an Artinian
ring over $k$,
and $P(t)$ its minimal polynomial.
Prove that $k[v]$ is isomorphic to
$k[t]/(P)$.
\end{exercise}

\begin{proof}
Evaluation map.
\end{proof}

\begin{exercise}
\label{commutative-algebra-exercise-chinese-reminder}
Let $R=k[t]/P$ where $P \in k[t]$ is a
polynomial
decomposing as a product $P=P_1…P_n$ of
coprime polynomials.
Prove that there exists an isomorphism
$R \to \bigoplus_ik[t]/P_i$.
\end{exercise}

\begin{proof}
The natural map $[F]_P \mapsto
([F]_{P_1},…,[F]_{P_n})$
is well defined.
It is surjective and injective by the Chinese
reminder theorem.
If $[F]_P = [t]_P$, i.e.
$F-t=PQ=P_1…P_nQ$, it's immediate that
$[F]_{P_i}=[t]_{P_i}$ for all $i$.
\end{proof}

\begin{exercise}
\label{commutative-algebra-exercise-artin-witho-idemp-field}
Let $R$ be a semisimple Artinian ring with
all idempotents
equal to 1 or 0. Prove it is a field.
\end{exercise}

\begin{proof}
Suppose $R$ is not a field, thus
there is an element $x$ which
is not invertible. Consider the
algebra $k[x]$ generated by $x$.
By a Exercise \ref{exercise-minimal-
polynomial-and-residues}
we know that $k[x] \simeq k[t]/P$
where $P$ is the minimal polynomial
of $x$.
Now decompose $P$ in its unique
irreducible factorization.
Since there are no nilpotents in $R$,
there are no nilpotents in
$k[x] \simeq k[t]/P$,
which implies that the irreducible
factorization of $P$
is square-free, which means that
the factors of $P$
are coprime and we can apply the
Chinese Reminder theorem.
That is, by Exercise
\ref{commutative-algebra-exercise-chinese-reminder},
we know that $k[t]/P$ is a direct
sum of fields
$\bigoplus k[t]/P_i$.
There is an idempotent element
here:
$(1,0,…,0)$.
\end{proof}

\noindent
Two idempotents $e_1,e_2 \in R$ are
called {\it orthogonal}
if $e_1e_2=0$. If $e_2,e_3$ are
orthogonal idempotents,
then $e_1=e_2+e_3$ is also an
idempotent
satisfying $e_2,e_3 \in e_1R$ and
$e_1R=r_2R \oplus e_3R$.
In characteristic $\neq 2$, for any
three idempotents
$e_1,e_2,e_3$, if $e_1=e_2+e_3$
then $e_2,e_3$ are orthogonal.
An idempotent $e \in R$ is called
{\it indecomposable} if there are
no non-zero
orthogonal idempotents such that
$e = e_2+e_2$.

\begin{exercise}
\label{commutative-algebra-exercise-subal-gener-indec-field}
Let $R$ be a semismple Artinian algebra,
and $e \in R$ a non-decomposable idempotent.
Then $eR$ is a field.
\end{exercise}

\begin{proof}
Using Exercise \ref{exercise-artinian
-without-idempotents-is-field}.
If there is an idempotent $er$
we can decompose $e$ as a sum
of orthogonal idempotents,
$e = (e-er)+er$.
\end{proof}

\noindent
The next exercise show that
there is an
indecomposable-orthogonal-
decomposition
of unity in an Artinian ring:

\begin{exercise}
\label{commutative-algebra-exercise-decom-unity-ortho-idemp}
Let $R$ be a semisimple Artinian ring
over a field $k$, $\text{char }k\neq 2$.
Prove that 1 can be decomposed to a
sum of indecomposable orthogonal
idempotents and such decomposition
is unique.
\end{exercise}

\begin{proof}
If 1 is indecomposable we are done.
If it is decomposable, say
$1=e_1+e_2$ then
we can write
$R=e_1R\oplus e_2R$.
This implies that $e_1R$ and $e_2R$
are vector subspaces of $R$
of strictly less dimension,
and thus we can repeat the process
and finish in finitely many steps.
For uniqueness we write
$1=e_1+…+e_k=e_1'+…+e'_{k'}$.
Multiplying by $e_1$ we obtain
$e_1=e_1e_1'+…+e_1e'_{k'}$
which, if the decompositions
do not coincide, is a decomposition
of $e_1$ into orthogonal idempotents,
a contradiction.
\end{proof}

\begin{exercise}
\label{commutative-algebra-exercise-artin-ring-direc-sum-field}
Let $R$ be an Artinian ring over a
field $k$, $\text{char }k\neq 2$.
Then $R$ is a direct sum of fields
and this decomposition is unique.
\end{exercise}

\begin{proof}
This is just
Exercises
\ref{commutative-algebra-exercise-decom-unity-ortho-idemp}
and \ref{commutative-algebra-exercise-artin-witho-idemp-field}.
\end{proof}

\begin{exercise}
\label{commutative-algebra-exercise-decomposition-char2}
Does it work in characteristic 2?
\end{exercise}

\begin{proof}
ChatGPT: should work but not this
technique.
Me: although we can decompose
$R = eR \oplus(1-e)R$,
we cannot show this is a field
unless we apply Chinese Reminder
Theroem! Perhaps this was a
mistake in the lectures,
or else the whole orthogonality
business is useless.
\end{proof}

\begin{exercise}
\label{commutative-algebra-exercise-n-prime-ideals}
Let $R$ be an Artinian ring over
a field $k$, $\text{char }k \neq 2$
and $1=e_1+…+e_n$ a decomposition
of 1 to a sum of indecomposable
orthogonal ideompotents.
Prove that $R$ has precisely $n$
prime ideals.
\end{exercise}

\begin{proof}
My idea is that the given
decomposition of unity gives a
decomposition of the algebra in
fields, which we may just think as an
affine algebraic variety which is
just $n$ points.
Any prime ideal in the algebra
is an irreducible subvariety
of the $n$ points,
which in turn can only be a point,
which is already one of the
obvious $n$ prime ideals.
\end{proof}

\medskip\noindent
The following result was proved
in lectures:

\begin{theorem}
\label{commutative-algebra-theorem-reduc-artin-conta-idemp}
An Artinian algebra without
nilpotents contains an idempotent.
\end{theorem}

\noindent
The proof of the theorem constructs
a maximal ideal $I \subset A$
which contains a unit, and thus is
a field $I$.

Thus we obtain

\begin{theorem}
\label{commutative-algebra-theorem-semisimple-decomposition}
An Artinian algebra with unity with
no nilpotents
is a direct sum of fields.
\end{theorem}

\begin{proof}
We find an idempotent element $e$,
so that we can decompose
$$
A=eA\oplus (1-e)A
$$
where $e$ is a unity on the first
factor and so is $1-e$ on the second
factor, thus every component
is a field [I think this might be
wrong!!! In the home assignment
we needed Chinese Reminder Theorem
to get the fieldness;
see Exercise
\ref{commutative-algebra-exercise-artin-witho-idemp-field}]
We repeat this process until we
exhaust the dimension of $A$.
\end{proof}

\noindent
We can also prove this theorem for
algebras without unity.
We consider the kernel of the map
$u_0:a \mapsto au$ and put
$$
A=u A \oplus \Ker u_0A.
$$
\begin{lemma}
\label{commutative-algebra-lemma-unique-decomposition-algebras}
Any decomposition of an algebra into
a direct sum of fields
is unique up to permutation of the
factors.
\end{lemma}


\medskip\noindent
We devote the end of this section
to explain the relationship
between semisimplicity and the trace form.
Namely,
Theorem
\ref{commutative-algebra-theorem-nonde-trace-semis}
states that a semisimple Artinian
algebra has nondegenerate trace form.
This is reminiscent of Cartan's
theorem which says that a
finite-dimensional Lie algebra
is semisimple in the sense that can be
decomposed as a direct sum of simple
Lie algebras (i.e. without nontrivial
proper ideals) if and only if the
Killing form (i.e. trace of adjoint
operators) is nondegenerate.

\begin{definition}
\label{commutative-algebra-definition-invariant-bilinear-form}
Let $R$ be a $k$-algebra, and $g:R \times R
\to k$
a symmetric $k$-bilinear form on $R$.
$g$ is called {\it invariant} if
$$
g(x,yz)=g(xy,z).
$$
\end{definition}

If $R$ has a unity we have
$$
g(x,y)=g(xy,1),
$$
so that in fact $g$ is determined by a
linear functional $a \mapsto g(a,1)$.

\begin{example}
\label{commutative-algebra-example-invariant-nondegenerate-form}
In the ring
$\mathbb{C}[x,y]/(x^{n+1},y^{n+1})$
the form $\varepsilon(xy)=a_nn$
is nondegenerate.
\end{example}

\begin{definition}
\label{commutative-algebra-definition-trace-form}
Let $R$ be an Artinian ring over $k$.
The {\it trace form} is the bilinear form
$a,b \mapsto \text{tr}(ab)$
where $\text{tr}(ab)$ is the trace
of the left-multiplication
by $ab$ endomorphism, i.e.
$L_{ab} \in \text{End}_kR$,
$x \mapsto abx$.
\end{definition}

\begin{exercise}
\label{commutative-algebra-exercise-field-extension-trace}
Let $[K:k]$ be a finite field extension
in characteristic 0.
Prove that the trace form is always
non-degenerate.
\end{exercise}

\begin{proof}
Since $[K:k]$ is finite,
we find that $K \simeq k^n$
and we can choose an orthonormal
base $e_i$ with respect to usual
Euclidean norm on $k^n$, i.e. dot product.
Then for any $a,b \in K$,
$$
\text{tr}(ab)=
\sum_i\left\langle{}abe_i,e_i\right\rangle{}
$$
by definition. Thus, since $K$ is a field
there exist inverses and a unity, so that
for any nonzero $x$,
$\text{tr}_k(x,x^{-1})=n \neq 0$.
\end{proof}

\begin{definition}
\label{commutative-algebra-definition-separable-field-extension}
A field extension is {\it separable}
if the trace form is nonzero.
\end{definition}

\begin{theorem}
\label{commutative-algebra-theorem-nonde-trace-semis}
Let $R$ be an Artinian algebra over $k$
with nondegenerate trace form.
Then $R$ is semisimple.
Conversely, for $\text{char }k=0$,
the trace form is nondegenerate
for any semisimple Artinian ring.
\end{theorem}

\begin{proof}
Suppose that $a \in R$ is a nilpotent
element. We shall show that this prevents
the trace form from being nondegenerate.
Indeed, we observe that for any $x \in R$,
the endomorphism $ax$ has no eigenvalues.
Suppose $ax\cdot y=\lambda y$ for some
$y \in R$ and $\lambda \in k$.
Then $(ax)^n\cdot y = \lambda^n y$
vanishes and thus $\lambda=0$.
Thus the trace of $ax$ is zero.

For the converse we use that any semismple
Artinian ring is a direct sum of fields.
Then the endomorphism $ab$ acts on
a vector $x \in R$ by
$(a_1b_1x_1,…,a_nb_nx_n)$
by definition of direct sum of algebras.
Thus the trace is the sum of
the traces in each of the fields that compose
$R$. By Exercise
\ref{commutative-algebra-exercise-field-extension-trace}
we know that each trace form is non
degenerate.
\end{proof}

\noindent
Note that the proof above uses
commutativity of the product in $R$.
The result also holds in noncommutative
algebras, but the proof must be reformulated
using the Jacobson ideal,
see Definition
\ref{commutative-algebra-definition-jacobson-ideal}.
\begin{lemma}
\label{commutative-algebra-lemma-trace-and-tensor-product}
The trace in the tensor product
is the tensor product in of the traces.
\end{lemma}

\section{Finite morphisms}
\label{commutative-algebra-section-finite-morphisms}

\begin{remark}
\label{commutative-algebra-remark-finit-gener-tenso-produ}
If $M$ is a finitely generated
$R$-module and $R \to R'$ is a ring
morphism, then $M\otimes_R R'$ is a
finitely-generated $R'$-module
generated by the same generators of
$M$.
\end{remark}

\begin{definition}
\label{commutative-algebra-definition-finite-morphism}
A morphism of affine
varieties $X \to Y$
is called {\it finite} if
$\mathcal{O}_X$ is a finitely
generated $\mathcal{O}_Y$-module. In
this case, $\mathcal{O}_X$ is called
an {\it integral extension} of
$\mathcal{O}_Y$.
\end{definition}

\begin{theorem}
\label{commutative-algebra-theorem-finit-map-has-finit-preim}
Let $f:X \to Y$ be a finite morphism of
affine varieties.
Then for any point $y \in Y$,
the preimage $f^{-1}(y)$ is finite.
\end{theorem}

\begin{proof}
Let $R$ be the $\mathcal{O}_Y$-module
$R:= \mathcal{O}_X
\otimes_{\mathcal{O}_Y}\mathcal{O}_Y/
\mathfrak{m}_y$,
which
is the coordinate ring of $f^{-1}(y)$
by Proposition
\ref{commutative-algebra-proposition-preimage-of-point}
or Lemma
\ref{commutative-algebra-lemma-preimage-of-affine-subvariety}.

Further, it is finitely generated
because $\mathcal{O}_X$ is finitely generated
over $\mathcal{O}_Y$
(Remark
\ref{commutative-algebra-remark-finit-gener-tenso-produ}).

Now observe that
$R$ is an Artinian algebra:
elements of $\mathcal{O}_X$
are finite $\mathcal{O}_Y$-linear
combinations of certain elements
$f_1,…,f_n \in \mathcal{O}_X$,
but in $R$ we have tensored by
$\mathcal{O}_Y/\mathfrak{m}_y$,
so by definition of tensor product
the $\mathcal{O}_Y$-coefficients
pass to the $\mathcal{O}_Y/\mathfrak{m}_y$
slot, where they collapse to
complex numbers; thus $R$
is a finite-dimensional vector
space over $\mathbb{C}$.
By Theorem
\ref{commutative-algebra-theorem-semisimple-decomposition},
$\Spec(R)$ is a finite set!
\end{proof}

\section{Field extensions}
\label{commutative-algebra-section-field-extensions}

\noindent
We collect the basic language of field
extensions here. It will be used later both
for Galois theory and for the function fields
which appear in Noether normalization.

\begin{definition}
\label{commutative-algebra-definition-field-extension}
A {\it field extension}
of the field $k$ is another field $K$ which
cointains $k$.
This is denoted $[K:k]$.
\end{definition}

\begin{definition}
\label{commutative-algebra-definition-algebraic-element}
Let $[K:k]$ be a field extension.
An element $x \in K$ is {\it algebraic} if it
is a root
of a nonzero polynomial with coefficients in
$k$,
and otherwise it is {\it trascendental}.
\end{definition}

\begin{definition}
\label{commutative-algebra-definition-finite-extension}
Let $[K:k]$ be a field extension.
We say that $[K:k]$ is a {\it finite
extension}
if $K$ is finite-dimensional as a vector
space
over $k$.
The dimension of this vector space is denoted
$[K:k]$ and is called the {\it degree} of the
extension.
\end{definition}

\begin{definition}
\label{commutative-algebra-definition-algebraic-extension}
Let $[K:k]$ be a field extension.
$[K:k]$ is {\it algebraic} if every element
of $K$ is algebraic over $k$.
\end{definition}

\noindent
Finite extensions are algebraic, but
algebraic extensions need not be finite.

\begin{definition}
\label{commutative-algebra-definition-field-generated-by-elements}
Let $[K:k]$ be a field extension, and let
$a_1,\ldots,a_m \in K$.
We write $k(a_1,\ldots,a_m)$ for the smallest
subfield of $K$ containing $k$ and
$a_1,\ldots,a_m$.
Equivalently, it is the fraction field of the
$k$-subalgebra $k[a_1,\ldots,a_m]\subset K$.
\end{definition}

\begin{definition}
\label{commutative-algebra-definition-finitely-generated-field-extension}
Let $[K:k]$ be a field extension.
We say that $[K:k]$ is a {\it finitely
generated field extension} if there exist
elements $a_1,\ldots,a_m \in K$ such that
$K=k(a_1,\ldots,a_m)$.
\end{definition}

\noindent
The simple upshot is that algebraic
independence is the analogue of being
transcendental, but for several elements at
once.

\begin{definition}
\label{commutative-algebra-definition-algebraically-independent}
Let $[K:k]$ be any field extension.
$t_1,…,t_n \in K$ are called
{\it algebraically independent} 
if there is no nonzero polynomial
$P \in k[z_1,…,z_n]$
such that $P(t_1,…,t_n)=0$.
\end{definition}

\begin{example}
\label{commutative-algebra-example-numbers-algebraically-dependent}
The numbers $1$ and $2$ are algebraically
dependent over $\mathbb{Q}$.
Indeed, the nonzero polynomial
$P(z_1,z_2)=z_1-1$ satisfies $P(1,2)=0$.
More generally, any two elements of
$\mathbb{Q}$ are algebraically dependent
over $\mathbb{Q}$.
\end{example}

\begin{example}
\label{commutative-algebra-example-transcendental-number-algebraically-independent}
A single transcendental number, such as
$\pi \in \mathbb{R}$, is algebraically
independent over $\mathbb{Q}$.
Indeed, algebraic independence of one element
just says that there is no nonzero polynomial
$P \in \mathbb{Q}[z]$ such that $P(\pi)=0$.
For two actual numbers, things get hard fast:
for example $e$ and $e^{\sqrt{2}}$ are
algebraically independent over $\mathbb{Q}$,
but this uses the Lindemann-Weierstrass
theorem.
\end{example}

\begin{example}
\label{commutative-algebra-example-coordinate-functions-algebraically-independent}
The elements $x,y \in \mathbb{C}(x,y)$
are algebraically independent over
$\mathbb{C}$.
Indeed, there is no nonzero polynomial
$P \in \mathbb{C}[z_1,z_2]$ such that
$P(x,y)=0$ as a rational function.
In contrast, the elements $x,x^2 \in
\mathbb{C}(x)$ are algebraically dependent,
since $z_2-z_1^2$ vanishes after setting
$z_1=x$ and $z_2=x^2$.
\end{example}

\begin{definition}
\label{commutative-algebra-definition-transcendence-basis}
Let $[K:k]$ be a field extension.
A set $t_1,\ldots,t_n \in K$ is called
a {\it transcendence basis} of $K$ over $k$
if $t_1,\ldots,t_n$ are algebraically
independent
over $k$ and $K$ is algebraic over
$k(t_1,\ldots,t_n)$.
\end{definition}

\begin{exercise}
\label{commutative-algebra-exercise-frac-fg-algebra-fg-field}
Let $A$ be a finitely generated domain over
a field $k$.
Prove that $\text{Frac}(A)$ is a finitely
generated field extension of $k$.
\end{exercise}

\begin{exercise}
\label{commutative-algebra-exercise-fg-algebraic-field-extension-finite}
Let $[K:k]$ be a finitely generated field
extension.
Assume that $K$ is algebraic over $k$.
Prove that $[K:k]$ is finite.
\end{exercise}

\begin{proof}
Suppose that $K=k(a_1,\ldots,a_n)$.
Let $K_0=k$ and
$K_i=K_{i-1}[a_i]$.
Since $K$ is algebraic over $k$, each $a_i$
is algebraic over $k$, hence also over the
larger field $K_{i-1}$.
By Lemma
\ref{commutative-algebra-lemma-integral-element-fg-submodule},
$K_i$ is a finitely generated
$K_{i-1}$-module.
Moreover, $K_i$ is a field: if
$0\neq h(a_i)\in K_{i-1}[a_i]$, let $m_i(t)$
be the minimal polynomial of $a_i$ over
$K_{i-1}$.
Since $h(a_i)\neq 0$, the polynomial $m_i$
does not divide $h$.
Thus $\text{gcd}(h,m_i)=1$, and by
Bachet-Bézout for polynomials over a field,
Theorem
\ref{theorem-bache-bezou-%
polyn-over-fiel},
there exist $u(t),v(t)\in K_{i-1}[t]$ such
that
$$
u(t)h(t)+v(t)m_i(t)=1.
$$
Evaluating at $a_i$ gives
$u(a_i)h(a_i)=1$, so $h(a_i)$ has inverse
$u(a_i)$ in $K_{i-1}[a_i]$.
Thus $K_i=K_{i-1}(a_i)$ is a finite extension
of $K_{i-1}$.
Inducting on $i$ and multiplying degrees,
$K=K_n$ is finite over $k$.
\end{proof}

\noindent
A field is called {\it algebraically closed}
is all algebraic extensions are trivial.
This is equivalent to the field containing
the roots of all polynomials.

\begin{lemma}
\label{commutative-algebra-lemma-algeb-close-iff-facto-polyn}
A field $k$ is algebraically closed if and
only if
any polynomial $P \in k[t]$ is a product of
linear factors.
\end{lemma}

\begin{proof}
If $P$ does not decompose,
there is a polynomial of degree >1 diving
$P$,
and the quotient $k[t]/$such polynomial is an
algebraic field extension.
\end{proof}

\noindent
There is a good notion of algebraic closure
of a field
obtained by adding roots and applying Zorn's
lemma
to ensure this process finishes. And it is
unique.

\section{Galois theory}
\label{commutative-algebra-section-galois-theory}

\noindent
Now we study finite extensions from the point
of view of automorphisms. Geometrically, one
should think of passing from a variety to a
cover and then studying automorphisms of the
corresponding extension of function fields.

\medskip\noindent
In this section we aim to prove:
let $[K:k]$ be a finite extension over with
$\text{char} k =0$.
then there exists an element $l \in k$
generating $k$.

\begin{theorem}
\label{commutative-algebra-theorem-separ-exten-algeb-closu}
Let $[\bar{k}:k]$ be the algebraic closure of
$k$,
and $[K:k]$ a separable finite extension.
Then
$$
K \otimes_k \bar{k} = \bigoplus \bar{k}.
$$
\end{theorem}

\noindent
Recall that by Galois theory we can see
that there only finitely many field
extensions.
Namely Galois theory says that field
extensions
are in correspondence with subgroups of the
Galois group, which is finite.
But this a simpler case:

\begin{lemma}
\label{commutative-algebra-lemma-finitely-many-k-algebras}
Let $k$ be a field,
and $A = \bigoplus_{i=1}^nk$. Then $A$
contains
only finitely many different $k$-algebras.
\end{lemma}

\begin{proof}

\end{proof}

\noindent
Now we can prove (without Galois theory):

\begin{theorem}
\label{commutative-algebra-theorem-finit-exten-char-0-gener}
let $[k:k]$ be a finite extension over with
$\text{char} k =0$.
then there exists an element $l \in k$
generating $k$.
\end{theorem}

\begin{proof}
Using finiteness of the contained algebras…
And the following fun claim:
if $V_i$ are positive codimension subspaces
of $k^n$,
then $\bigcup V_i \neq  k^n$.
\end{proof}

\begin{definition}
\label{commutative-algebra-definition-galois-extension}
Let $[K:k]$ be a finite extension.
It is called a {\it Galois extension} if the
algebra
$K \otimes_k K$ is isomorphic to a direct sum
of several copies of $K$.
\end{definition}

\begin{example}
\label{commutative-algebra-example-galois-extension}
Let $K=k[t]/(P)$
be a primitive, separable extension with
$\text{deg} P=n$.
Suppose that $P$ has $n$ roots in $K$.
Then $[K:k]$ is Galois.
Indeed, this is the Chinese reminder exercise
once again:
$$
K \otimes_k K = \frac{K[t]}{(P)}
=\bigoplus_i \frac{K[t]}{(t-a_i)}.
$$
\end{example}

\begin{exercise}
\label{commutative-algebra-exercise-galois-extension}
$K=k[t]/P$.
$[K:k]$ is Galois if and only if
$P$ has $n$ roots in $K[t]$.
\end{exercise}

\begin{definition}
\label{commutative-algebra-definition-splitting-field}
Let $P \in k[t]$
be an irreducible polynomial.
{\it Splitting field} is the minimal
extension
such that $P$ factors into linear factors.

Equivalently, pick $P \in k[t]$ of degree $n$
without multiple roots,
and $k$ of characteristic 0.
Decompose $P$ into irreducible factors and
construct finite extensions
$K_1\subset … \subset K_\ell$ by adding the
roots of each factor.
This process ends (prove this) and the
resulting field is
the {\it splitting field}.
\end{definition}

\noindent
Splitting field is Galois by Exercise
\ref{commutative-algebra-exercise-galois-extension}.
We can think of the splitting field as
embedded in $\bar{k}$.
There is not a unique embedding but all
embeddings are isomorphic.

\begin{definition}
\label{commutative-algebra-definition-galois-group}
If $[K:k]$ is a Galois extension,
$\text{Aut}_k K$ is called {\it Galois
group}.
\end{definition}

\noindent
Note that $K^*$ acts on $K \otimes_k K$ by
multiplication
on the right and on the left.

\begin{lemma}
\label{commutative-algebra-lemma-galois-group-and-extensions}
Let $[K:k]$ be a Galois extension.
There is a bijection between the following
sets:
\begin{enumerate}
\item $\text{Aut}_k K$.
\item Prime ideals in $K \otimes_k K$.
\item Homomorphisms $K \otimes_k K \to K$
which are linear with respect to the left
action of $K$.
\end{enumerate}
\end{lemma}

\begin{proof}
(2 $\implies$ 3) is simple.
\end{proof}

\begin{lemma}
\label{commutative-algebra-lemma-galois-extension-and-direct-sum}
Let $[K:k]$ be a Galois extension,
and $K\otimes_k K = \bigoplus_{\nu \in
\text{Aut}_k K}K_\nu$
the decomposition of $K \otimes_k K$
enumerated by elements of $\text{Aut}_k K$.
Then the left $\mu_l$ and right $\mu_r$
actions of $K^*$ on
$K\otimes_k K$ are such that
$\mu_l(a)e_\nu = \mu_r(\nu(a))e_\nu$.
\end{lemma}

\begin{lemma}
\label{commutative-algebra-lemma-gal-invariant}
Let $[K:k]$ be a Galois extension,
and $a \in K$ a $\text{Aut}_k K$-invariant
element.
Then $a \in k$.
\end{lemma}

\begin{lemma}
\label{commutative-algebra-lemma-galois-and-g-invariants}
Let $[K:k]$ be a Galois extension, and $K'$
an intermediate field
(which, as we will see in the proof, will
also be Galois)
Then
$$
K' = K^{G'}
$$
where $G' \subset \text{Aut}_{K'}K$ is the
group
of $K'$-linear automorphisms of $K$
and $K^{G'}$ is the space of $G'$-invariants.
\end{lemma}

\begin{theorem}[Main theorem of Galois
theory]
\label{commutative-algebra-theorem-main-theorem-of-galois-theory}
Let $[K:k]$ be a Galois extension.
The subgroups $H \subset \text{Aut}_k K$ are
in bijective correspondence
with intermediate subfields
$k \subset K^H \subset K$.
\end{theorem}


\section{Dominant
and birational morphisms}
\label{commutative-algebra-section-dominant-birational-morphisms}

\noindent
Recall that $\Spec(R)$ is the set of all
prime ideals of the ring $R$. We introduce a
topology whose closed sets are the prime
ideals contained in a given ideal.

\begin{definition}
\label{commutative-algebra-definition-dominant-morphism}
A morphism of algebraic varieties $f:X \to Y$
is called {\it dominant}
if its image is dense,
i.e. $Y$ is the Zariski closure of $f(X)$.
\end{definition}

\begin{proposition}
\label{commutative-algebra-proposition-domin-iff-induc-map}
A morphism of varieties $f:X \to Y$
is dominant if and only if
$f^*:\mathcal{O}_Y \to \mathcal{O}_X$ is
injective.
\end{proposition}

\begin{proof}
Suppose $f^*$ is not injective and
consider its kernel $\Ker f^*$.
Then the image of $f$ is in the zero set of
the kernel,
$V(\Ker f^*)$.
Here we use that points of a variety are
homomorphisms from the variety to
$\mathbb{C}$,
which follows from the Nullstellensatz…
Dominantness means that the image is not
contained in
the zero set of an ideal.
\end{proof}

\begin{lemma}
\label{commutative-algebra-lemma-dominant-morphism-irreducible}
If $f:X \to Y$ is dominant and $X$ is
irreducible, then $Y$ is irreducible.
Moreover, $f^*:\mathcal{O}_Y \to
\mathcal{O}_X$ can be extended to the fields
of fractions $k(Y) \to k(X)$.
\end{lemma}

\begin{proof}
Since $f$ is dominant, $f^*$ is injective,
i.e. an embedding,
irreducibility transforms to not having
zerodivisors.
Descending to the field of fraction also
follows from not having zerodivisors.
\end{proof}

\begin{remark}
\label{commutative-algebra-remark-divisor}
All of a sudden we consider a divisor on the
affine variety $X$
given by $f \in \mathcal{O}_X$.
Then $X\setminus D$ is also an affine variety
given by the ideal
$I+\left\langle{}ft-1\right\rangle{} \subset
\mathbb{C}[x_1,…,x_n,t]$.
Recall that the localization of a ring by an
element is
$\mathbb{C}/\left\langle{}1-ft\right\rangle{}$, where $t$ is
the inverse of $f$.
\end{remark}

\begin{definition}
\label{commutative-algebra-definition-birational-morphism}
A morphism of irreducible algebraic varieties
$f:X \to Y$ is called a
{\it birational morphism} if
the induced map on fraction fields
$k(Y) \to k(X)$ is an isomorphism.
\end{definition}

\noindent
We further show:

\begin{proposition}
\label{commutative-algebra-proposition-birat-then-isomo-away}
Let $f:X \to Y$ be a birational morphism.
Then there exists a divisor $Z \subset Y$
such that $f:(X \setminus f^{-1}(Z)) \to Y
\setminus Z$
is an isomorphism.
\end{proposition}

\medskip\noindent
Now we introduce normal varieties.

Here's some notes from the past that may pop
up in memory: The upshot is: let $\varphi:R
\to S$ be a ring map. An element $s \in S$ is
integral over $R$ if there is a monic
polynomial with coefficients in $R$ that
gives zero when you put $s$ instead of the
variable.

\section{Integrally closed rings
and normal varieties}
\label{commutative-algebra-section-integrally-closed}

Integral closure is for rings like 
algebraic closure for fields,
but it's not exactly the same construction.
For a field, the algebraic closure is a
larger field in which all algebraic equations
split.
For a ring, integral closure is relative to
an ambient ring, usually the fraction field,
and only adds elements satisfying monic
polynomial equations.

\begin{definition}
\label{commutative-algebra-definition-integral-element}
Let $A \subset B$ be rings. An element $b \in
B$ is called {\it integral over $A$} if the
subring
$A[b]=A\cdot \left\langle{}1,b,b^2,\ldots\right\rangle{}$
is generated by $b$ and $A$, is finitely
generated as an $A$-module.
\end{definition}

\noindent
Note that if $A$ is finitely generated, $A
\left\langle{}1,b,b^2,\ldots\right\rangle{}$ is obviously
finitely
generated as a ring, but not as a module,
which is what we care about.

\begin{lemma}
\label{commutative-algebra-lemma-integral-element-fg-submodule}
An element $x \in B$ is integral over $A
\subset B$
if and only if
the chain of submodules
$$
A \subset A\cdot \left\langle{}1,x\right\rangle{}
\subset A \cdot\left\langle{}1,x,x^2\right\rangle{}
\subset…
$$
terminates.
Equivalently, $x$ is the root of a
monic polynomial with coefficients in $A$.
\end{lemma}

\begin{proof}
The chain statement is almost tautological
As for the second, express $x^n$ as an
element of the finitely-generated module.
Conversely, if $x$ is a root of a monic
polynomial
$$
x^n+a_{n-1}x^{n-1}+\cdots+a_0=0,
$$
then $x^n$ is an $A$-linear combination of
$1,x,\ldots,x^{n-1}$.
Multiplying by powers of $x$ and arguing by
induction, every higher power of $x$ is also
an $A$-linear combination of
$1,x,\ldots,x^{n-1}$.
Thus $A[x]$ is generated as an $A$-module by
$1,x,\ldots,x^{n-1}$.
\end{proof}

\begin{proposition}
\label{commutative-algebra-prop-localization-integrally-closed}
Let $R$ be an integrally closed domain and
let $S \subset R$ be a multiplicative
subset with $0 \notin S$.
Then $RS^{-1}$ is integrally
closed.
\end{proposition}

\begin{proof}
Let $x \in \text{Frac}(RS^{-1})$ be
integral over $RS^{-1}$.
Then there is a monic polynomial
\[
x^n+c_{n-1}x^{n-1}+\cdots+c_0=0
\]
with $c_i \in RS^{-1}$.
Choose $s \in S$ so that
$sc_i \in R$ for all $i$.
Multiplying by the right power of $s$,
we get
\[
(sx)^n+s c_{n-1}(sx)^{n-1}
+\cdots+s^n c_0=0.
\]
This is a monic polynomial over $R$
vanishing at $sx$.
Since $R$ is integrally closed,
we have $sx \in R$.
Therefore $x=(sx)/s \in RS^{-1}$.
\end{proof}

\begin{lemma}
\label{commutative-algebra-lemma-sums-of-integral-elements}
In Noetherian rings,
sums and products of integral
elements is integral.
\end{lemma}

\begin{definition}
\label{commutative-algebra-definition-normal-variety}
An irreducible affine algebraic variety is
called {\it normal} if its coordinate ring is
integrally closed, that is, it coincides with
its its integral closure, which in turn is
the set of all elements that are integral
over its field of fractions.
\end{definition}

\noindent
Recall from \cite[Exercise 3.17, Chapter
1]{har}
that for affine varieties this is equivalent
to all
the local rings are all integrally closed.

\begin{proposition}
\label{commutative-algebra-proposition-normal-iff-birational}
An affine variety $X$ is normal if and only
if
for any finite, birational morphism
$Y \to X$ is an isomorphism.
\end{proposition}

\begin{proof}
They have the same fraction field, and
$\mathcal{O}_X \subset \mathcal{O}_Y$ by
dominant (birational maps are dominant),
which ways that $\mathcal{O}_Y$ is finite
over $\mathcal{O}_X$, which says that
elements of $\mathcal{O}_Y$ are integral.
Since $X$ is normal then the elements of
$\mathcal{O}_Y$ belong to $\mathcal{O}_X$.
\end{proof}

\begin{proposition}
\label{commutative-algebra-proposition-ufd-integrally-closed}
A UFD is integrally closed.
\end{proposition}

\begin{proof}
Let $u,v$ be elements of the UFD $A$
and suppose the $u/v \in k(A)$ is integral,
i.e. there is a monic polynomial
$P \in A[t]$ of degree $n$ killing $u/v$.
Then $u^n$ is divisible by $v$.
Write $v=p_1…p_k$.
Then $p_i$ divides $u^n$,
which implies that it divides $u$.
But then $p_i$ divides both $u$ and $v$,
which we may take to be coprime, making
$p_i$ a unit for all $i$.
\end{proof}


\noindent
As a corollary, we see that taking product of
an affine variety, whose coordinate ring is
an UFD, with affine space, $X \times
\mathbb{C}^n$ preserves normality since the
resulting coordinate ring is
$\mathcal{O}_X[x_1,…,x_n]=\mathcal{O}_X
\otimes \mathbb{C}[x_1,…,x_n]$.

\section{Nakayama's lemma}
\label{commutative-algebra-section-nakayama}

\noindent
``By Nakayama lemma,
maximal ideal generates everything:
if $x_1,…,x_n$ generate $M/\mathfrak{m} M$
then they generate $M$''

\medskip\noindent
We would like to have an analaogue of
the inverse function theorem,
i.e. a function that is surjective
on the tangent bundle is locally invertible.
However, this is not true.
Consider $B$ a locally free
$\mathcal{O}_M$-module.

The idea is to find local trivializations
of projective modules,
just as we have local trivializations
of vector bundles.

\medskip\noindent
Nakayama's lemma is sometimes
also called Krull lemma.
We motivate with the following question:
for an ideal $I \subset A$,
when is $\cap I^n=0$?
Answer: when
$$
I\cdot \bigcap I^n = \bigcap I^n.
$$

The setting is:
$\mathfrak{a}\subset A$ an ideal,
$M$ finitely generated module…


\begin{definition}
\label{commutative-algebra-definition-torsion}
Let $A$ be a domain an
$M$ be an $A$-module.
A {\it torsion} is the set of all $m \in M$
such that there exists $0 \neq a \in A$
such that $a m=0$.
\end{definition}

\begin{definition}
\label{commutative-algebra-definition-jacobson-ideal}
If $R$ is a ring,
its {\it Jacobson ideal} is
the intersection of all maximal ideals.
\end{definition}

\begin{lemma}
\label{commutative-algebra-lemma-jacobson-ideal}
$a \in R$ belongs to $J(R)$
if and only if for all $r \in R$,
$1+ar$ is invertible.
\end{lemma}

\begin{definition}
\label{commutative-algebra-definition-quasi-regular-ideal}
An ideal $\mathfrak{a}\subset R$
is {\it quasi-regular} if $1+a$
is invertible for any $a \in \mathfrak{a}$.
\end{definition}

\begin{lemma}
\label{commutative-algebra-lemma-quasi-regular-jacobson}
An ideal is quasi-regular if and only if
it belongs to the Jacobson ideal.
\end{lemma}

\begin{proof}
Straightforward.
\end{proof}

\noindent
Recall that the intuition of localization
is ``to pass to a smaller set'' -
at least geometrically.

\begin{definition}
\label{commutative-algebra-definition-localization-module}
Let $S \subset A$ be a multiplicative
system,
and $A(S^{-1})$ the localization of $A$
in $A$.
The {\it localization} of $M$ in $S$
is $M(S^{-1}):= M \otimes_A A(S^{-1})$.
\end{definition}

\begin{lemma}
\label{commutative-algebra-lemma-kernel-localization}
The natural map $\tau:M\to M(S^{-1})$
has kernel
$\{m \in M:ms=0\text{ for some }
s \in S \}$.
\end{lemma}

\begin{lemma}[Nakayama, first version,
AKA Krull lemma]
\label{commutative-algebra-lemma-nakayama-first}
If $\mathfrak{a}\subset A$ is a quasiregular
ideal and $M$ is a finitely generated
$A$-module. Then
$$
M=\mathfrak{a}M \implies M=0.
$$
\end{lemma}

\begin{lemma}[Nakayama]
\label{commutative-algebra-lemma-nakayama}
Let $A$ be a domain and $M$
a finitely generated $A$-module.
If $M = \mathfrak{a}M$ for some nonzero ideal
$\mathfrak{a}\subset A$
then $(1+a)M=0$
for some $a \in \mathfrak{a}$.
\end{lemma}

\begin{lemma}
\label{commutative-algebra-lemma-nakayama-equivalence}
Lemmas \ref{commutative-algebra-lemma-nakayama-first}
and \ref{commutative-algebra-lemma-nakayama} are equivalent.
\end{lemma}

\begin{proof}
Converse is immediate,
for forward implication uses
Lemma \ref{commutative-algebra-lemma-kernel-localization}.
\end{proof}

\noindent
Recall that we call a module
{\it cyclic} when it is generated by
a single element, i.e.
$M=\left\langle{}x\right\rangle{}$.
This gives a surjective map
$A \to M$, $a \mapsto ax$
and so $M=A/\Ker$.
So cyclic modules are quotients of
the ring by an ideal.

\begin{proof}
First prove it for cyclic $M$.
The definition of quasiregular ideal and
cyclicness of $M$ give us
$(1-a)x=0$ which means that $z=0$
since $1-a$ is invertible.

If $M$ is generated by $z_1…,z_n$,
we have a matrix given by the equation
$\mathfrak{a}M=M$.
Namely
$$
z_i=\sum a_{ij}z_j
$$
then put
$$
(1-a_{11})z_1=\sum_{i=2}^na_{ij}z_j.
$$
Dividing by $1-a_{11}$ since it is
invertible, we can ``get rid of $z_1$'',
and construct a submodule $M_1$, namely
the submodule generated by the remaining
generators, and in fact it satisfies
$\mathfrak{a}M_1=M_1$. Apply induction
hypothesis.

\end{proof}

\noindent
We won't prove the following theorem:

\begin{theorem}[Cayley-Hamilton]
\label{commutative-algebra-theorem-cayley-hamilton}
Let $M$ be a finitely-generated $A$-module.
An $A$-linear map $\varphi:M \to M$
can be represented by a matrix,
which has a characteristic polynomial $P$.
Then $P(A)=0$.
\end{theorem}

The idea is to use this theorem
to express the identity of $M$
as a matrix with coefficients in $A$:
$$
\xymatrix{
A^n\ar[r]^{\tilde{\varphi}}\ar[d]_{\pi}
&A^n\ar[d]^{\pi}\\
M\ar[r]_{\varphi=\text{Id}_M}
&M
}
$$
In fact, this almost concludes the proof
of Nakayama's lemma from Atiyah-Macdonald.

\medskip\noindent
Now we turn to applications.
Recall that localization by a prime ideal
is a local ring. This says that elements
outside the maximal ideal are invertible.
The converse is also true and
is used along with Nakayama's lemma
to prove the following result.
The intuition behind this is that
if a map of modules induces a surjection
on local quotients, then it is invertible
locally (inverse function thorem!).

In the following lemma
I think of $M_1,M_2$ as their
corresponding affine varieties
over the local algebra $(A,\mathfrak{m})$,
which I can think as the stalk of a point.
Then $M_1/\mathfrak{m}M_1
=M_1 \otimes_A A/\mathfrak{m}$
by Exercise
\ref{commutative-algebra-exercise-tensor-product-maximal-ideal}.
But we also know
by Proposition
\ref{commutative-algebra-proposition-preimage-of-point}
that this is the ring of the fiber of
the point corresponding to the maximal ideal.
So the following lemma says:
if the map is surjective on fibers
of the point, then it is surjective
on the stalk.
It looks like for a global statement, i.e.
when $A$ is not local,
we could check this for every
maximal ideal and conclude that
$\varphi$ is a surjective morphism.

\begin{lemma}
\label{commutative-algebra-lemma-localization-nakayama}
Let $(A,\mathfrak{m})$ be a local ring
and $M_1,M_2$ be finitely generated
$A$-modules. Let $\varphi:M_1 \to M_2$
and suppose it induces a surjective map
$M_1/\mathfrak{m}M_1
\to M_2/\mathfrak{m}M_2$.
Then $\varphi$ is surjective.
\end{lemma}

\noindent
This easily implies

\begin{lemma}
\label{commutative-algebra-lemma-local-modules}
Let $M$ be a finitely generated $A$-module,
$A$ local. Let $x_1,…,x_n \in M$
be generators of $M$ mod $\mathfrak{m}$.
Then $x_1,…,x_n$ generate $M$.
\end{lemma}

\begin{theorem}
\label{commutative-algebra-theorem-finite-dominant-surjective}
A finite dominant map is surjective.
\end{theorem}

\begin{proof}
$A = \mathcal{O}_Y$, $B=\mathcal{O}_X$,
$\mathfrak{m}_y \in A$ maximal ideal of $y$.
We need dominance to suppose that
$B \subset A$.

Then $\mathfrak{m}_yB \neq B$ by
Nakayama's lemma.

Now consider $f^{-1}(y)
=\Spec(B \otimes_A A/\mathfrak{m}_y)
=\Spec(B /\mathfrak{m}_yB)\neq 0$.
\end{proof}

\medskip\noindent
Recall that projective modules
are summands of free modules.
Vector bundles over affine varieties
are the same as projective modules
over the rings of functions!
That's a result by Serre, but we won't
prove it here.

\begin{theorem}
\label{commutative-algebra-theorem-projective-locally-free}
A finitely generated projective module
over a local ring is free.
\end{theorem}

\begin{proof}
Consider a map
$$
\begin{aligned}
A^n &\longrightarrow P \\
(a_1,…,a_n) &\longmapsto \sum a_ip_i.
\end{aligned}
$$

It is surjective and thus has a section
since $P$ is projective.

In writing
$A^n = P \oplus Q$
we kill $Q$ when taking quotient
by the maximal ideal
since $\mathfrak{m} Q = Q$
by Nakayama.
\end{proof}

\medskip\noindent
As a final application consider the
Groethendick group $K_0(R)$
where the subindex 0 denotes
that we consider only exact sequences
of projective modules.

\begin{proposition}
\label{commutative-algebra-proposition-grothendieck-projective}
Let $(A,\mathfrak{m})$
be a local ring.
Then $K_0(A) = \mathbb{Z}$.
\end{proposition}

\begin{proof}
Map $P \in K_0(Q)$ to
$\dim_{A/\mathfrak{m}}P/\mathfrak{m}P$.
Since projective modules are free,
this is a bijection.
\end{proof}


\section{Noether normalization lemma}
\label{commutative-algebra-section-noether-normalization}

\noindent
Normalization means taking integral closure
in the field of rational functions.
Noether normalization lemma is the key
ingredient to prove the Finiteness of
Integral Closure theorem,
Theorem
\ref{commutative-algebra-theorem-integrl-clsure-fntly-genertd},
which says that the integral closure
is a finitely generated algebra.
In turn, this is the key fact needed
to define normalization of a variety:
it's the spectrum of the integral closure!

\medskip\noindent
We start discussing finite quotients.

\begin{theorem}[Noether theorem for
G-invariants]
\label{commutative-algebra-theorem-noether-g-invariants}
Let $R$ be a finitely generated
ring over $\mathbb{C}$ and $G$ a finite
group acting on $R$ by automorphisms.
The ring $R^G$ of $G$-invariants
is finitely generated.
\end{theorem}

\begin{proof}
A bit involved - past lecture.
\end{proof}

\begin{lemma}
\label{commutative-algebra-lemma-g-invariants-surjective}
Let $R$ be a Noetherian ring without
zero divisors, $G$ a finite group
acting by automorphisms on $R$,
and $R^G$ the ring of $G$-invariants.
Then $\varphi:\Spec R \to \Spec R^G$
is a finite, dominant morphism.
\end{lemma}

\begin{proof}
Dominance is immediate because the map
on rings
is clearly injective, which is our
definition of dominant map.
We need to show that $R$ is finitely
generated
as a $R^G$-module.
Recall that to show that the ring
is finitely generated is equivalent
to finding minimal polynomials…
\end{proof}

\begin{example}
\label{commutative-algebra-example-quotient-singularity}
$\mathbb{C}^2=\Spec\mathbb{C}[x,y]$
under the antipodal action of $\mathbb{Z}_2$.
We obtain $\mathbb{C}^2/\pm 1
\simeq \mathbb{C}[a,b,c]/(c^2-ab)$
because invariant polynomials are
even polynomials.

Another example is $\mathbb{C}/\mathbb{Z}_n$,
which gives $\Spec \mathbb{C}[t^n]$
which is actually $\mathbb{C}$ again.
\end{example}

\begin{theorem}
\label{commutative-algebra-theorem-orbit-space}
The fibers of the natural map
$\varphi:\Spec R \to \Spec R^G$
are the orbits of $G$.
\end{theorem}

\begin{proof}
We prove that the orbit is a finite
set of points by showing the corresponding
algebra is a semisimple Artinian algebra.
\end{proof}

\medskip\noindent
Now we proceed
on our way to proving Noether's
normalization lemma.

\begin{lemma}
\label{commutative-algebra-lemma-field-extension-finite}
Let $[k:\mathbb{C}]$ be an extension of
$\mathbb{C}$. Then either $z$ is
trascendental and $K$ is isomorphic to
the field of rational functions,
or $z$ is algebraic and the extension
is finite.
\end{lemma}

\begin{proof}
A polynomial vanishing on the point
gives us finiteness once again…
\end{proof}

\noindent
We will use the terminology of field
extensions from Section
\ref{commutative-algebra-section-field-extensions}.
In particular,
a finitely generated field extension does not
mean that $K=k(t_1,\ldots,t_n)$ for some
transcendence basis $t_1,\ldots,t_n$.
Relations among generators are allowed. This
explains why, if
$X \subset \mathbb{C}^n$ is an affine
variety,
then $k(X)=\text{Frac}(\mathcal{O}_X)$ is a
finitely generated field extension of
$\mathbb{C}$: the coordinate classes
$\bar z_1,\ldots,\bar z_n$ generate $k(X)$ as
a field, even though they may satisfy the
relations coming from the ideal of $X$.
See Exercise
\ref{commutative-algebra-exercise-frac-fg-algebra-fg-field}.

\begin{theorem}
\label{commutative-algebra-theorem-finitely-gen-extension}
Let $[K:k]$ be a finitely generated field
extension, e.g. $[k(X):\mathbb{C}]$
where $X$ is an affine algebraic variety.
Then $K$ is a finite extension
of $k(t_1,…,t_n)$ where $t_1,…,t_n$
is a transcendence basis.
\end{theorem}

\begin{proof}
Choose generators $a_1,\ldots,a_m$ such that
$K=k(a_1,\ldots,a_m)$.
Among them choose a maximal algebraically
independent subset, say $t_1,\ldots,t_n$.
By maximality, every remaining generator is
algebraic over $k(t_1,\ldots,t_n)$.
Since there are only finitely many remaining
generators, adjoining them gives an algebraic
finitely generated extension of
$k(t_1,\ldots,t_n)$, hence a finite extension
by Exercise
\ref{exercise-fg-algebraic-field-extension-%
finite}.
\end{proof}

\noindent
Let us spell out the notation in geometric
terms. If $X\subset \mathbb{C}^n$, then
$$
\mathcal{O}_X
=\mathbb{C}[z_1,\ldots,z_n]/I(X),
$$
and the symbols $t_i$ below are really the
classes of the ambient coordinate functions
$z_i+I(X)$ in this quotient ring, viewed also
as elements of the function field $k(X)$.
If some of these elements are algebraically
independent, then they can serve as genuine
independent coordinates on a dense part of
$X$; geometrically this will say that the
projection to those coordinates is dominant.

\begin{proposition}
\label{commutative-algebra-proposition-dominant-transcendence}
Let $X \subset \mathbb{C}^n$ be an 
affine irreducible variety.
Let $\Pi_k:X \to \mathbb{C}^k$
be the projection to the first $k$
coordinates.
The following are equivalent:
\begin{enumerate}
\item $\Pi_k$ is dominant
and the extension $[k(X):k(t_1,…,t_k)]$
is finite.
\item $t_1,…,t_k$ is a transcendence basis
in $k(X)$.
\end{enumerate}
\end{proposition}

\begin{proof}
Dominance of $\Pi_k$ means that
$\mathbb{C}[t_1,\ldots,t_k]\to
\mathcal{O}_X$ is injective.
This is equivalent to saying that
$t_1,\ldots,t_k$ satisfy no nonzero
polynomial relation, i.e. they are
algebraically independent.
Passing to fraction fields gives
$\mathbb{C}(t_1,\ldots,t_k)\subset k(X)$,
the field generated by these independent
coordinates inside $k(X)$.

The remaining condition that
$[k(X):\mathbb{C}(t_1,\ldots,t_k)]$ is finite
is precisely the statement that $k(X)$ is
algebraic over the field generated by these
elements; since $k(X)$ is finitely generated
because $\mathcal{O}_X$ is a finitely
generated algebra, algebraic here implies
finite. Thus
$t_1,\ldots,t_k$ is a transcendence basis
inside $k(X)$.
\end{proof}

\noindent
Here is the density point more concretely.
The image of $\Pi_k$ is not dense exactly
when it is contained in the zero set of some
nonzero polynomial
$P \in \mathbb{C}[z_1,\ldots,z_k]$.
But being contained in that zero set means
$P(t_1,\ldots,t_k)=0$ on $X$, i.e. a
polynomial relation among the classes
$t_i=z_i+I(X)$.
Thus no polynomial relation among the $t_i$
means the projection has dense image.

We are looking for conditions where
this projection is a finite map.
The projection $\Pi_{n-1}$ is finite
if and only if $P(z_n)=0$ in $\mathcal{O}_X$
for some monic polynomial
$P(t) \in \mathbb{C}[z_1,…,z_{n-1}][t]$.

\begin{lemma}
\label{commutative-algebra-lemma-coord-change-transcen-basis}
Let $X\subset \mathbb{C}^n$ be an irreducible
affine variety and suppose that the
coordinate
classes $z_1,\ldots,z_k \in k(X)$ form a
transcendence basis.
Then there exists a linear coordinate change
$z_i':=z_i+\lambda_iz_n$ such that $z_n$
is finite over $z_1',…,z_k'$,
that is, $z_n$ is a root of a monic
polynomial with coefficients in
$\mathbb{C}[z_1',…,z_k']$.
\end{lemma}

\begin{proof}
Since $z_1,\ldots,z_k$ is a transcendence
basis, $z_n$ is algebraic over
$\mathbb{C}(z_1,\ldots,z_k)$.
Thus there is a polynomial relation involving
$z_n$ and $z_1,\ldots,z_k$.
After replacing
$z_i$ by $z_i'=z_i+\lambda_i z_n$, choose the
constants $\lambda_i$ generically so that the
highest-degree part of this relation does not
vanish in the $z_n$ direction.
Then the relation can be viewed as a
polynomial in $z_n$ whose leading coefficient
is a nonzero constant.
Dividing by this leading coefficient gives a
monic polynomial for $z_n$ with coefficients
in $\mathbb{C}[z_1',\ldots,z_k']$.
This is exactly the statement that $z_n$ is
finite over $z_1',\ldots,z_k'$.
\end{proof}

\noindent
Finally we can prove Noether's normalization
lemma. The finite map obtained in this
theorem is the key ingredient
to prove the finiteness of the integral
closure of an algebra.

\begin{lemma}[Noether's normalization]
\label{commutative-algebra-lemma-noether-normalization}
Let $X\subset \mathbb{C}^n$
be an irreducible affine variety.
Then there exists a linear coordinate
change such that the projection of $X$
to the first $k$ coordinates
gives a finite map $X \to \mathbb{C}^k$.
\end{lemma}

\noindent
As a corollary we may obtain a dominant map:

\begin{lemma}
\label{commutative-algebra-lemma-noether-nrmlztin-dominant}
Let $X \subset \mathbb{C}^n$
be an irreducible affine variety,
$z_i$ coordinates on $\mathbb{C}^n$,
and $z_1,…,z_k$ transcendence basis
on $k(X)$.
Then there exists a linear change of 
coordinates $z_i'=z_i+\sum_{j=1}^{n-k}
\lambda_{j+k}z_{j+k}$
such that the projection to the
first $k$ variables is a finite,
dominant morphism.
\end{lemma}

\begin{proof}
The fact that $z_1,\ldots,z_k$ is a
transcendence basis says, by Proposition
\ref{commutative-algebra-proposition-dominant-transcendence},
that the projection to the first $k$
coordinates is dominant and generically
finite.
Apply Lemma
\ref{commutative-algebra-lemma-coord-change-transcen-basis}
successively to the remaining coordinates
$z_{k+1},\ldots,z_n$.
After a suitable linear change, each
remaining
coordinate becomes a root of a monic
polynomial with coefficients in
$\mathbb{C}[z_1',\ldots,z_k']$.
Thus all coordinate functions of $X$ are
integral over
$\mathbb{C}[z_1',\ldots,z_k']$.
Since $\mathcal{O}_X$ is generated by these
coordinate functions, this makes
$\mathcal{O}_X$ finite over
$\mathbb{C}[z_1',\ldots,z_k']$, which is
exactly the finiteness of the projection.
\end{proof}

\noindent
Now we introduce a version that
applies to non-irreducible varieties
--- at the cost of losing dominance.

\begin{proposition}
\label{commutative-algebra-proposition-noether-reducible-vars}
Let $X \subset \mathbb{C}^n$ be an affine
subvariety and $X_i$ its irreducible
components. Denote by $k$ the maximal
transcendence degree for $k(X_i)$.
Then there exists a linear coordinate
change $z_i'$ such that the projection
to the first $k$ variables is finite.
\end{proposition}

\noindent
Recall that the integral closure of $A$ 
in $K$ is the set of all elements
such that $A[x]$ is a finitely
generated $A$-module.
The integral closure of $A$ is its integral
closure in its field of fractions $k(A)$.

\begin{theorem}[Finiteness of Integral
Closure]
\label{commutative-algebra-theorem-integrl-clsure-fntly-genertd}
Let $A$ be an integrally closed Noetherian
ring without zero divisors and
$[K:k(A)]$ a finite extension of its
field of fractions, and $B$ the
integral closure of $A$ in $K$.
Then $B$ is finitely generated
as an $A$-module.
\end{theorem}

\begin{proof}
Here is the expected strategy.
\begin{enumerate}
\item Use Noether normalization to find a
polynomial subring
$R=\mathbb{C}[x_1,\ldots,x_d]\subset A$
such that $A$ is finite over $R$.
\item Reduce the problem from $A$ to the
polynomial ring $R$.
\item Observe that $K$ is a finite extension
of $\text{Frac}(R)$, since it is finite over
$k(A)=\text{Frac}(A)$ and $A$ is finite over
$R$.
\item Prove that the integral closure of
$R$ in $K$ is finite over $R$.
This is the serious algebra step, usually
proved using trace.
\item Conclude that the integral closure of
$A$ in $K$ is contained in this finite
$R$-module, hence is finite over $R$ and
therefore finitely generated over $A$.
\end{enumerate}
\end{proof}

\begin{lemma}
\label{commutative-algebra-lemma-normalization}
Let $B$ be a domain over $\mathbb{C}$.
Assume that there is an injective
ring morphism
$$
A=\mathbb{C}[x_1,\ldots,x_k]\to B
$$
such that $B$ is finitely generated
as an $A$-module.
Then the integral closure $\tilde{B}$
of $B$ is finitely generated
as an $A$-module.
In particular, $\tilde{B}$ is a
finitely generated ring.
\end{lemma}

\begin{proof}
The ring $A$ is factorial by Gauss'
lemma, hence integrally closed.
Since $B$ is finite over $A$,
the field extension
$\text{Frac}(B)/\text{Frac}(A)$
is finite.
Theorem
\ref{commutative-algebra-theorem-integrl-clsure-fntly-genertd}
implies that the integral closure
of $A$ in $\text{Frac}(B)$
is a finitely generated $A$-module.
The integral closure $\tilde{B}$
is contained in this ring.
Since $A$ is Noetherian, $\tilde{B}$
is a finite $A$-module as well.
\end{proof}

\section{Normalization}
\label{commutative-algebra-section-normalization}

\begin{definition}
\label{commutative-algebra-definition-normalization}
Let $X$ be an irreducible affine variety,
and let $A=\mathcal{O}_X$.
Let $\tilde{A}$ be the integral closure
of $A$ in its field of fractions.
Assume that $\tilde{A}$ is a finitely
generated $\mathbb{C}$-algebra.
Then
$$
\tilde{X}=\Spec \tilde{A}
$$
is called the {\it normalization} of $X$.
The inclusion $A \subset \tilde{A}$
gives a morphism
$\tilde{X}\to X$.
\end{definition}

\noindent
The normalization map is birational
to $X$, and it is finite.

\begin{example}
Let $C \subset \mathbb{A}^2$ be the curve
defined by
$$
x^2=y^3.
$$
Let's check if $C$ is normal.

First we need a parametrization of $C$.
This means that we look for functions
$x(t)$ and $y(t)$ such that
$$
x(t)^2=y(t)^3.
$$
In this example it is natural to try powers
of $t$, say
$$
x(t)=t^a,\qquad y(t)=t^b.
$$
Then we need
$$
t^{2a}=t^{3b},
$$
hence
$$
2a=3b.
$$
The smallest positive solution is
$$
a=3,\qquad b=2.
$$
Thus we take
$$
x(t)=t^3,\qquad y(t)=t^2.
$$
Thus we have a parametrization
$$
\varphi:\mathbb{C}\to C,\qquad
t\longmapsto (t^3,t^2).
$$
On coordinate rings this gives
$$
\mathbb{C}[x,y]/(x^2-y^3)
\cong \mathbb{C}[t^2,t^3]
\subset \mathbb{C}[t].
$$
Thus $\mathcal{O}_C=\mathbb{C}[t^2,t^3]$
and its fraction field is $\mathbb{C}(t)$.
The element $t$ lies in this fraction field,
since $t=x/y$, and it is integral over
$\mathcal{O}_C$ because it satisfies
$$
t^2-y=0.
$$
But $t\notin \mathbb{C}[t^2,t^3]$,
so $C$ is not normal.

Moreover $\mathbb{C}[t]$ is integral over
$\mathbb{C}[t^2,t^3]$, because it is
generated by the integral element $t$.
It is integrally closed, since it is a
polynomial ring in one variable.
Therefore the integral closure of
$\mathcal{O}_C$ in $\mathbb{C}(t)$
is $\mathbb{C}[t]$.
Thus the normalization of $C$ is
$$
\Spec \mathbb{C}[t]=\mathbb{A}^1,
$$
with normalization map $\varphi$.
\end{example}


\begin{exercise}
\label{commutative-algebra-exercise-normalization}
\begin{enumerate}
\item Parametrize the curve
$$
C = V(x^3-y^4) \subset \mathbb{A}^2.
$$
\item What is $\mathcal{O}_C$?
\label{commutative-algebra-item-2-exer-normalization}
\item What is its fraction field?
\item What is its normalization?
\end{enumerate}
\end{exercise}

\begin{proof}
We want functions of the form
$$
x(t)=? \quad  y(t)=?
$$
giving points in $C$.
We guess that this will work:
$$
x(t)=t^a \quad  y(t)=t^b.
$$
For those points to be in the curve we need
$$
t^{3a} = t^{4b},
$$
so
$$
3a=4b.
$$
Obviously we pick
$$
a=4,\qquad b=3.
$$
Alright, the candidate parametrization is
$$
\begin{aligned}
\varphi: \mathbb{C} &\longrightarrow
\mathbb{C}^2 \\
t &\longmapsto (t^4,t^3).
\end{aligned}
$$

\noindent
But is it surjective on $C$?
Actually we don't need surjectivity.
To prove this parametrization is good
we look at the induced map on rings,
and confirm that the coordinate algebra
$\mathcal{O}_C$ of $C$ appears.
Look at the pullback
map on regular functions:
$$
\begin{aligned}
\varphi^*: \mathbb{C}[x,y] 
&\longrightarrow \mathbb{C}[t] \\
x &\longmapsto t^4\\
y &\longmapsto t^3.
\end{aligned}
$$

\noindent
What's the kernel? Obviously
$x^3-y^4$ is in the kernel --- that's
the whole point. To prove that's all
the kernel let $f \in \Ker \varphi^*$.

Think of $\mathbb{C}[x,y]$
as
$$
\mathbb{C}[y][x]
$$
that is, as polynomials in $x$
with coefficients in $\mathbb{C}[y]$.
This has a division algorithm
for polynomials whose leading
coefficient is a unit
because $\mathbb{C}[y]$ is a domain
(see Theorem
\ref{commutative-algebra-theorem-eucli-algor-polyn-rings})
which is due to $\mathbb{C}$
being a domain
(see Lemma
\ref{commutative-algebra-lemma-ring-polyn-one-indet-over}).
So we can divide $f$ by
the monic polynomial $x^3 - y^4 
\in \mathbb{C}[y][x]$:
$$
f = q(x)(x^3-y^4) + r(x)
\qquad \text{deg }r < 3,
$$
where $q,r$ have coefficients in 
$\mathbb{C}[y]$. Since the degree
in $x$ of $r$ is at most 2, we can
write:
$$
f
=
q(x,y)(x^3-y^4)
+
r_0(y)+r_1(y)x+r_2(y)x^2,
$$
where
$$
q(x,y)\in \mathbb C[x,y],
\qquad
r_0,r_1,r_2\in \mathbb C[y].
$$

Now use the fact that $\varphi^*f=0$:
$$
0=\varphi^*f=f(t^4,t^3)
=q(t^4,t^3)(t^{12}-t^{12}) 
+ r_0(t^3)
+r_1(t^3)t^4
+r_2(t^3)t^8
$$
Hence
$$
r_0(t^3)+r_1(t^3)t^4+r_2(t^3)t^8=0.
$$
Now the powers of $t$ appearing in
$r_0(t^3)$ are congruent to $0$ modulo $3$.
The powers of $t$ appearing in
$r_1(t^3)t^4$ are congruent to $1$ modulo $3$.
The powers of $t$ appearing in
$r_2(t^3)t^8$ are congruent to $2$ modulo $3$.
Thus these three terms cannot cancel each
other. Therefore
$$
r_0(t^3)=0,\qquad
r_1(t^3)t^4=0,\qquad
r_2(t^3)t^8=0.
$$
So
$$
r_0(t^3)=r_1(t^3)=r_2(t^3)=0.
$$
But the map
$$
\mathbb C[y]\longrightarrow \mathbb C[t],
\qquad
y\longmapsto t^3
$$
is injective. Hence
$$
r_0=r_1=r_2=0.
$$
Therefore
$$
f=q(x,y)(x^3-y^4),
$$
so
$$
f\in (x^3-y^4).
$$
This proves
$$
\Ker\varphi^*=(x^3-y^4).
$$


Then we have an isomorphism
$$
\mathcal{O}_C=
\frac{\mathbb{C}[x,y]}{(x^3-y^4)}
\cong \text{Im} \varphi^*
=\mathbb{C}[t^4,t^3]
$$
This answers Item
(\ref{commutative-algebra-item-2-exer-normalization}).
Writing the coordinate ring of $C$
in this form is what allows us
to check if it's normal or not, as show
next.

Observe that $C$ is not normal.
The fraction $t^4/t^3 = t \in \mathbb{C}(C)
=\text{Frac}\mathbb{C}[t^4,t^3]$.
Further, $t$ is integral over
$\mathcal{O}_C$ since there is a 
monic polynomial 
$P(z)\in \mathcal{O}_C[z]$
vanishing at $t$:
$P(z)=z^3-t^3$. Also we can put
$Q(z)=z^4-t^4$ but not
$R(z)=z-t$ because 
$t \not \in \mathcal{O}_C$.
(This may seem stupid to the reader
but when we face the problem without
the open book it might be not so obvious.)
Precisely because $t \not \in \mathcal{O}_C$
we conclude that $\mathcal{O}_C$
is not integrally closed, and so $C$ cannot
be normal.

To find the normalization we seek
the integral closure of 
$\mathbb{C}[t^4,t^3]$.
As observed above,
such integral closure must contain $t$,
which makes it a ring which conatins
$\mathbb{C}[t]$,
which is already integrally closed
(because it's a UFD,
c.f. Proposition
\ref{commutative-algebra-proposition-ufd-integrally-closed}),
and therefore we conclude that the
normalization of $C$ is
$\Spec\mathbb{C}[t]=\mathbb{C}$,
that is, $C$ is rational
with normalization map
$$
\mathbb A^1\longrightarrow C,
\qquad
t\longmapsto (t^4,t^3).
$$
\end{proof}

\begin{lemma}
\label{commutative-algebra-lemma-universal-property-normalization}
Let $\nu:\tilde{X}\to X$ be the
normalization of an irreducible affine
variety.
Then $\nu$ is finite and birational.
For every finite birational morphism
$\mu:Y\to X$, there is a unique morphism
$\tilde{X}\to Y$ over $X$.
In particular, if $Y$ is normal, then
$Y$ is isomorphic to $\tilde{X}$ over $X$.
\end{lemma}

\begin{proof}
Let $A=\mathcal{O}_X$ and
$\tilde{A}$ be its integral closure in
$k(X)$.
Since $\tilde{A}$ is finite over $A$,
the map $\nu$ is finite.
Since $A$ and $\tilde{A}$ have the same
fraction field, $\nu$ is birational.

Let $Y$ be finite and birational over $X$.
Then $\mathcal{O}_Y$ is a finite
$A$-module inside $k(X)$.
Thus every element of $\mathcal{O}_Y$
is integral over $A$, so
$\mathcal{O}_Y\subset \tilde{A}$.
This inclusion gives the unique morphism
$\tilde{X}\to Y$ over $X$.

If $Y$ is normal, then $\mathcal{O}_Y$
is integrally closed in $k(X)$ and contains
$A$.
Hence it contains $\tilde{A}$, and so
$$
\mathcal{O}_Y=\tilde{A},
$$
which gives the desired isomorphism.
\end{proof}

\noindent
Equivalently, normalization is the initial
object in the category whose objects are
finite birational morphisms
$\mu:X'\to X$, and whose morphisms are maps
over $X$.
If one restricts to normal sources $X'$,
the preceding lemma says that all objects
are canonically isomorphic to the
normalization.

\begin{exercise}
\label{commutative-algebra-exercise-normal-surface}
Prove that the ring 
$\mathbb{C}[x,y,z]/(y^2-xz)$
is integrally closed.
\end{exercise}

\begin{proof}
Consider a parametrization of the
corresponding affine surface:
$$
(t,s) \longmapsto (t^2,t s,s^2)
$$
This gives an isomorphism
$$
\frac{k[x,y,z]}{(y^2-xz)}
\cong
k[t^2,t s,s^2] \subset k[t,s].
$$
Let $A=\mathbb{C}[t^2,ts,s^2]$.
To prove $A$ is integrally closed
notice that $A$ is exactly the subring
of $\mathbb{C}[t,s]$ of polynomials which
remain invariant under
$\varphi:(t,s)\mapsto (-t,-s)$. Indeed,
it is clear that elements $A$ are 
$\varphi$-invariant, and conversely,
any $\varphi$-invariant polynomial
must have even degree in all variables
and thus it must be in $A$.

Now let $u \in \operatorname{Frac}A$
be the root of a monic polynomial
with coefficients in $A$.
Since $A \subset \mathbb{C}[t,s]$,
and $\mathbb{C}[t,s]$ is integrally closed,
we conclude $u \in \mathbb{C}[t,s]$.
But since $u$ is a quotient of elements
in $A$, it is $\varphi$-invariant.
Therefore $u \in A$.
\end{proof}

\begin{exercise}
\label{commutative-algebra-exercise-quotient-not-smooth}
Let $\zeta_n \in \mathbb{C}$
be a primitive root of unity
of degree $n$.
Define the action of the cyclic group
$G = \mathbb{Z}/n\mathbb{Z}$ on $\mathbb{C}^2$
with coordinates $x,y$
in such a way that the generator
$t$ maps $x$ to $\zeta_nx$ and
$y$ to $\zeta_n^{-1}y$.
Prove that the quotient $\mathbb{C}^2/G$
is singular.
\end{exercise}

\begin{proof}
The ring of invariants is
$\mathbb{C}[x^n,y^n,xy]$.
Consider the ideal 
$\mathfrak{m}=(x^n,y^n,xy)$.
It is clear that, for $n\geq 2$,
the cotangent space
$\mathfrak{m}/\mathfrak{m}^2$
has three linearly independent vectors,
namely, the classes of
$x^n,y^n$ and $xy$, and thus it can't
be regular.
\end{proof}
 

\section{Algebraic dimension}
\label{commutative-algebra-section-algebraic-dimension}

\begin{definition}
\label{commutative-algebra-definition-algebraic-dimension}
Let $M$ be an irreducible affine variety.
The {\it algebraic dimension} $a(M)$
is the transcendence degree of the
field extension
$$
[k(M):\mathbb{C}].
$$
\end{definition}

\noindent
Equivalently, if $k(M)$ is finite over
$\mathbb{C}(t_1,\ldots,t_n)$,
then $a(M)=n$.

\medskip\noindent
We shall prove that this definition agrees
with Krull dimension.
The idea is to cut $M$ by a principal
divisor, then cut this divisor by a
principal divisor on it, and continue in
this way until we get a point.
The Krull dimension shall be the number
of such operations.

\begin{definition}
\label{commutative-algebra-definition-principal-divisor}
Let $X$ be an affine variety, and let
$f\in \mathcal{O}_X$ be a regular function
which does not vanish on any irreducible
component of $X$.
The zero set of $f$ is called a
{\it principal divisor} on $X$.
Its irreducible components are called
{\it divisors} on $X$.
\end{definition}

\begin{lemma}
\label{commutative-algebra-lemma-divisors-affine-space}
Any divisor in $\mathbb{A}^n$ is principal.
\end{lemma}

\begin{proof}
Let $I=(f)\subset \mathbb{C}[t_1,\ldots,t_n]$
be a principal ideal.
Write
$$
f=\prod_i f_i^{\alpha_i}
$$
as the prime decomposition of $f$.
Since $\mathbb{C}[t_1,\ldots,t_n]$
is factorial, the ideals $(f_i)$
are prime.
Moreover
$$
V(f)=\bigcup_i V(f_i).
$$
Since all $(f_i)$ are prime, this is the
irreducible decomposition of $V(f)$.
Thus each irreducible component of $V(f)$
is the zero set of a single equation,
hence is principal.
\end{proof}

\begin{lemma}
\label{commutative-algebra-lemma-algebraic-dimension-divisor}
Let $D$ be an irreducible divisor in
$\mathbb{A}^n$.
Then
$$
a(D)=n-1.
$$
\end{lemma}

\begin{proof}
By Lemma
\ref{commutative-algebra-lemma-divisors-affine-space},
we can write $D=V(f)$ for an irreducible
polynomial
$f\in \mathbb{C}[t_1,\ldots,t_n]$.
After reordering the coordinates, suppose
that $f$ depends nontrivially on $t_n$.

The equation $f=0$ gives one algebraic
relation among the coordinate functions
on $D$, so
$$
a(D)\leq n-1.
$$
On the other hand, the classes of
$t_1,\ldots,t_{n-1}$ in $\mathcal{O}_D$
are algebraically independent.
Indeed, if a nonzero polynomial in
$\mathbb{C}[t_1,\ldots,t_{n-1}]$
vanished on $D$, it would be divisible
by $f$, which is impossible because
$f$ depends on $t_n$.

Finally, the class of $t_n$ in $k(D)$
is a root of the polynomial $f$,
viewed as a polynomial in $t_n$
with coefficients in
$\mathbb{C}[t_1,\ldots,t_{n-1}]$.
Thus $k(D)$ is finite over
$\mathbb{C}(t_1,\ldots,t_{n-1})$,
and hence $a(D)=n-1$.
\end{proof}

\noindent
The point is that $t_n$ satisfies one
polynomial equation over
$\mathbb{C}(t_1,\ldots,t_{n-1})$.
Thus, after passing to rational functions,
adjoining $t_n$ is an algebraic extension.
Since the field is generated by this one
algebraic element over
$\mathbb{C}(t_1,\ldots,t_{n-1})$,
the extension is finite.

\medskip\noindent
The length of a chain
$$
A_1\subset A_2\subset \cdots \subset A_n
$$
is $n-1$, namely the number of inclusions.

\begin{definition}
\label{commutative-algebra-definition-krull-dimension-ring}
The {\it Krull dimension} of a ring $A$
is the maximal possible length of a chain
of prime ideals
$$
0\neq \mathfrak{p}_1
\subset \mathfrak{p}_2
\subset \cdots
\subset \mathfrak{p}_n
\subset A.
$$
\end{definition}

\begin{definition}
\label{commutative-algebra-definition-krull-dimension-variety}
The {\it Krull dimension} of a variety $X$
is the maximal possible length of a chain
of nonempty, irreducible, distinct
subvarieties
$$
X_1\subset X_2\subset \cdots \subset X_n.
$$
\end{definition}

\begin{theorem}
\label{commutative-algebra-theorem-krull-algebraic-geometric}
Let $X$ be an affine variety.
Then the Krull dimension of $X$ as a
variety is equal to the Krull dimension
of its coordinate ring $\mathcal{O}_X$.
\end{theorem}

\begin{proof}
Irreducible subvarieties of $X$
correspond to prime ideals of
$\mathcal{O}_X$ by taking vanishing ideals.
This correspondence reverses inclusions:
if $Y\subset Z$, then
$$
I(Z)\subset I(Y).
$$
Thus a chain of irreducible subvarieties
of $X$ is the same thing as a chain of
prime ideals in $\mathcal{O}_X$, with the
order reversed.
The two maximal possible lengths are
therefore equal.
\end{proof}

\noindent
We shall prove below that for any affine
variety, the algebraic dimension and the
Krull dimension coincide.

\begin{lemma}[Nakayama for local rings]
\label{commutative-algebra-lemma-nakayama-local-rings}
Let $A$ be a Noetherian local ring,
and let $\mathfrak{p}$ be its maximal ideal.
Let $M$ be a finitely generated $A$-module.
If
$$
\mathfrak{p}M=M,
$$
then $M=0$.
\end{lemma}

\begin{definition}
\label{commutative-algebra-definition-finite-ring-extension}
Let $A\subset B$ be rings.
We say that $B$ is a {\it finite extension}
of $A$ if $B$ is finitely generated
as an $A$-module.
\end{definition}

\noindent
If $A\subset B$ and $\mathfrak{q}\subset B$
is an ideal, the {\it contraction} of
$\mathfrak{q}$ to $A$ is the ideal
$$
\mathfrak{q}\cap A\subset A.
$$
It is just the set of elements of $A$
which already lie in $\mathfrak{q}$.

\begin{lemma}[Going down]
\label{commutative-algebra-lemma-contraction-nonzero-prime}
Let $A\subset B$ be domains, and assume
that $B$ is a finite extension of $A$.
Let $\mathfrak{q}\subset B$ be a nonzero
prime ideal.
Then
$$
\mathfrak{p}:=\mathfrak{q}\cap A
$$
is a nonzero prime ideal of $A$.
\end{lemma}

\begin{proof}
The contraction of a prime ideal is prime,
so it remains to prove that
$\mathfrak{p}\neq 0$.
Choose $0\neq b\in \mathfrak{q}$.
Since $B$ is finite over $A$,
the element $b$ is integral over $A$.
Thus $b$ satisfies a monic equation
$$
b^n+a_{n-1}b^{n-1}+\cdots+a_0=0
$$
with $a_i\in A$.
Choose such an equation with $n$ minimal.
Then $a_0\neq 0$, otherwise we could divide
by $b$ because $B$ is a domain.
But
$$
a_0=-b(b^{n-1}+a_{n-1}b^{n-2}
+\cdots+a_1)
$$
lies in $\mathfrak{q}$.
Hence $0\neq a_0\in \mathfrak{q}\cap A$.
\end{proof}

\begin{proof}[Alternative proof]
Let
$$
S=A\setminus \{0\}.
$$
Suppose that $\mathfrak{q}\cap A=0$.
Then $S\cap \mathfrak{q}=\emptyset$,
so $\mathfrak{q}S^{-1}B$ is a nonzero
proper ideal in $S^{-1}B$.

On the other hand,
$S^{-1}A=\text{Frac}(A)$ is a field.
Since $B$ is finite over $A$,
the ring $S^{-1}B$ is a finite-dimensional
algebra over $\text{Frac}(A)$.
In particular it is Artinian.
It is also a domain, because $B$ is a
domain.
By Theorem
\ref{commutative-algebra-theorem-semisimple-decomposition},
an Artinian algebra without nilpotents
is a direct sum of fields.
Since $S^{-1}B$ is a domain, this direct
sum has only one summand.
Thus $S^{-1}B$ is a field.

But a field has no nonzero proper ideals,
contradicting the existence of
$\mathfrak{q}S^{-1}B$.
Therefore $\mathfrak{q}\cap A\neq 0$.
\end{proof}

\noindent
The proof of the next lemma first localizes
at $\mathfrak{p}$.
Nakayama shows that the fiber over
$\mathfrak{p}$ is nonzero.
Prime ideals in this fiber correspond
precisely to prime ideals of $B$ lying over
$\mathfrak{p}$.
The fiber is a finite-dimensional algebra
over a field, hence Artinian.
Thus the Artinian algebra classification
shows that it has only finitely many prime
ideals.

\begin{lemma}[Going up]
\label{commutative-algebra-lemma-going-up}
Let $A\subset B$ be rings, with $A$
Noetherian, and assume that $B$ is a
finite extension of $A$.
Let $\mathfrak{p}\subset A$ be a prime
ideal.
Then there exist finitely many prime ideals
$\mathfrak{q}\subset B$ such that
$$
\mathfrak{q}\cap A=\mathfrak{p}.
$$
\end{lemma}

\begin{proof}
Let
$$
S=A\setminus \mathfrak{p}.
$$
Localize:
$$
A_{\mathfrak{p}}=S^{-1}A,\qquad
B_{\mathfrak{p}}=S^{-1}B.
$$
The ring $A_{\mathfrak{p}}$ is local,
with maximal ideal
$\mathfrak{p}A_{\mathfrak{p}}$.
Also $B_{\mathfrak{p}}$ is finite as an
$A_{\mathfrak{p}}$-module.

By Nakayama's lemma,
Lemma
\ref{commutative-algebra-lemma-nakayama-local-rings},
we have
$$
B_{\mathfrak{p}}\neq
\mathfrak{p}B_{\mathfrak{p}}.
$$
Hence
$$
C:=B_{\mathfrak{p}}/
\mathfrak{p}B_{\mathfrak{p}}
$$
is a nonzero finite-dimensional algebra
over the field
$A_{\mathfrak{p}}/
\mathfrak{p}A_{\mathfrak{p}}$.
In particular, $C$ is Artinian.
Every prime ideal of $C$ contains its
nilradical.
After quotienting by the nilradical,
the Artinian algebra decomposition,
Theorem
\ref{commutative-algebra-theorem-semisimple-decomposition},
identifies the reduced quotient with a
finite direct sum of fields.
Thus there are only finitely many prime
ideals in $C$.

Let $\overline{\mathfrak{q}}\subset C$
be a prime ideal, and let
$\tilde{\mathfrak{q}}\subset
B_{\mathfrak{p}}$
be its preimage.
Then $\tilde{\mathfrak{q}}$ is prime.
Put
$$
\mathfrak{q}:=
\tilde{\mathfrak{q}}\cap B.
$$
Then $\mathfrak{q}$ is prime in $B$.
Moreover
$\tilde{\mathfrak{q}}\cap
A_{\mathfrak{p}}
=\mathfrak{p}A_{\mathfrak{p}}$,
because the quotient
$A_{\mathfrak{p}}/
\mathfrak{p}A_{\mathfrak{p}}$
is a field.
Contracting back to $A$ gives
$\mathfrak{q}\cap A=\mathfrak{p}$.

Conversely, every prime of $B$ lying over
$\mathfrak{p}$ arises from a prime ideal
of $C$ in this way.
Thus there are only finitely many such
prime ideals $\mathfrak{q}$.
\end{proof}

\noindent
The Cohen-Seidenberg theorem is essentially
the union of the going-up and going-down
lemmas.

\begin{theorem}[Cohen-Seidenberg]
\label{commutative-algebra-theorem-cohen-seidenberg}
Let $A\subset B$ be a finite extension
of Noetherian rings.
Let
$$
\mathfrak{q}_1\subset \mathfrak{q}_2
\subset \cdots \subset \mathfrak{q}_n
\subset B
$$
be a chain of prime ideals, and put
$$
\mathfrak{p}_i=\mathfrak{q}_i\cap A.
$$
Then
$$
\mathfrak{p}_1\subset \mathfrak{p}_2
\subset \cdots \subset \mathfrak{p}_n
\subset A
$$
is a chain of distinct prime ideals.

Conversely, any chain of prime ideals
$$
\mathfrak{p}_1\subset \mathfrak{p}_2
\subset \cdots \subset \mathfrak{p}_n
\subset A
$$
is obtained in this way from a chain of
prime ideals in $B$.
\end{theorem}

\begin{proof}
We prove (i).
Suppose that for some $i$ we have
$$
\mathfrak{q}_{i-1}\subset \mathfrak{q}_i
$$
but
$$
\mathfrak{p}_{i-1}=\mathfrak{p}_i.
$$
Replacing $A$ by $A/\mathfrak{p}_{i-1}$
and $B$ by $B/\mathfrak{q}_{i-1}$,
we reduce to the case where $A\subset B$
are domains and
$0\neq \mathfrak{q}_i\subset B$
is a prime ideal.
The Going down lemma,
Lemma
\ref{commutative-algebra-lemma-contraction-nonzero-prime},
then gives
$\mathfrak{q}_i\cap A\neq 0$,
a contradiction.
Thus the contractions are distinct.

We prove (ii).
Now let
$$
\mathfrak{p}_1\subset \cdots
\subset \mathfrak{p}_n
$$
be a chain of prime ideals in $A$.
By the Going up lemma,
Lemma
\ref{commutative-algebra-lemma-going-up},
there is a prime
$\mathfrak{q}_1\subset B$ lying over
$\mathfrak{p}_1$.
Assume by induction that
$$
\mathfrak{q}_1\subset \cdots
\subset \mathfrak{q}_r
$$
have been chosen over
$$
\mathfrak{p}_1\subset \cdots
\subset \mathfrak{p}_r.
$$
Apply the Going up lemma to the finite
extension
$$
A/\mathfrak{p}_r
\subset B/\mathfrak{q}_r
$$
and to the prime
$\mathfrak{p}_{r+1}/\mathfrak{p}_r$.
This gives a prime ideal
$\mathfrak{q}_{r+1}/\mathfrak{q}_r$
of $B/\mathfrak{q}_r$ lying over it.
Equivalently, we get a prime
$\mathfrak{q}_{r+1}\subset B$
with
$$
\mathfrak{q}_r\subset \mathfrak{q}_{r+1}
\quad\text{and}\quad
\mathfrak{q}_{r+1}\cap A
=\mathfrak{p}_{r+1}.
$$
This completes the induction.
\end{proof}

\begin{lemma}
\label{commutative-algebra-lemma-krull-dimension-finite-dominant}
Let $f:X\to Y$ be a finite dominant
morphism of irreducible affine varieties.
Then the Krull dimension of $X$ is equal
to the Krull dimension of $Y$.
\end{lemma}

\begin{proof}
The morphism $f$ gives an injective finite
ring extension
$$
\mathcal{O}_Y\subset \mathcal{O}_X.
$$
Any chain of prime ideals in
$\mathcal{O}_Y$ lifts to a chain in
$\mathcal{O}_X$ by Theorem
\ref{commutative-algebra-theorem-cohen-seidenberg}.
Conversely, any chain of distinct prime
ideals in $\mathcal{O}_X$ contracts to a
chain of distinct prime ideals in
$\mathcal{O}_Y$, again by Theorem
\ref{commutative-algebra-theorem-cohen-seidenberg}.
Therefore the maximal possible lengths of
chains are equal.
\end{proof}

\begin{theorem}
\label{commutative-algebra-theorem-algebraic-dimension-krull}
For any affine variety $X$, its algebraic
dimension $a(X)$ is equal to its Krull
dimension.
\end{theorem}

\begin{proof}
Let $d=a(X)$.
By Noether normalization,
Lemma
\ref{commutative-algebra-lemma-noether-normalization},
there exists a finite dominant morphism
$$
X\to \mathbb{A}^d.
$$
By Lemma
\ref{commutative-algebra-lemma-krull-dimension-finite-dominant},
$X$ and $\mathbb{A}^d$ have the same
Krull dimension.
Thus it remains to prove that
$$
\dim \mathbb{A}^d=d.
$$

We prove this by induction on $d$.
For $d=0$, affine space is a point.
Assume the result for affine varieties
of algebraic dimension $<d$.
Let
$$
0=\mathfrak{p}_0\subset \mathfrak{p}_1
\subset \cdots
$$
be a maximal chain of prime ideals in
$$
\mathbb{C}[t_1,\ldots,t_d].
$$
The prime ideal $\mathfrak{p}_1$ contains
a nonzero polynomial $g$.
Writing $g$ as a product of irreducibles,
and using that $\mathfrak{p}_1$ is prime,
we find an irreducible polynomial
$f\in \mathfrak{p}_1$.
Since $\mathbb{C}[t_1,\ldots,t_d]$
is a UFD, $(f)$ is a prime ideal.
Thus
$$
0\subset (f)\subset \mathfrak{p}_1.
$$
By maximality of the chain,
$\mathfrak{p}_1=(f)$.

Let $D=V(f)\subset \mathbb{A}^d$.
Then $D$ is an irreducible divisor.
Prime ideals containing $(f)$ correspond
to prime ideals of
$$
\mathbb{C}[t_1,\ldots,t_d]/(f)
=\mathcal{O}_D.
$$
Hence
$$
\dim \mathbb{A}^d=1+\dim D.
$$
By Lemma
\ref{commutative-algebra-lemma-algebraic-dimension-divisor},
we have $a(D)=d-1$.
By the induction hypothesis,
$$
\dim D=d-1.
$$
Therefore
$$
\dim \mathbb{A}^d=1+(d-1)=d.
$$
This proves the theorem.
\end{proof}

\medskip\noindent
We finish with Hilbert's Nullstellensatz.
The proof below is algebraic: it uses
Noether normalization instead of the usual
complex-analytic trick.
This is more appropriate here, because the
analytic proof is special to
$\mathbb{C}$, while the algebraic proof
works over any algebraically closed field.

\begin{theorem}[Hilbert's Nullstellensatz]
\label{commutative-algebra-theorem-hilbert-nullstellensatz}
Let $k$ be an algebraically closed field.
Then every maximal ideal of
$k[x_1,\ldots,x_n]$ is of the form
$$
(x_1-a_1,\ldots,x_n-a_n)
$$
for some point $a=(a_1,\ldots,a_n)\in k^n$.
Moreover, for every ideal
$I\subset k[x_1,\ldots,x_n]$,
we have
$$
I(V(I))=\sqrt{I}.
$$
\end{theorem}

\begin{proof}
Let $\mathfrak{m}\subset k[x_1,\ldots,x_n]$
be a maximal ideal, and put
$$
A=k[x_1,\ldots,x_n]/\mathfrak{m}.
$$
Then $A$ is a field, finitely generated
as a $k$-algebra.
By Noether normalization, there is a
polynomial subring
$$
k[t_1,\ldots,t_d]\subset A
$$
such that $A$ is finite over it.
Since $A$ is a field and is integral over
$k[t_1,\ldots,t_d]$, the ring
$k[t_1,\ldots,t_d]$ is a field.
This forces $d=0$.
Therefore $A$ is a finite algebraic
extension of $k$.
Since $k$ is algebraically closed, $A=k$.
Thus the quotient map
$$
k[x_1,\ldots,x_n]\to A=k
$$
is evaluation at a point
$a=(a_1,\ldots,a_n)$, and
$$
\mathfrak{m}
=(x_1-a_1,\ldots,x_n-a_n).
$$

Now let $I\subset k[x_1,\ldots,x_n]$.
The inclusion
$$
\sqrt{I}\subset I(V(I))
$$
is immediate.
Conversely, let $f\notin \sqrt{I}$.
Then the ideal
$$
J=(I,1-tf)\subset k[x_1,\ldots,x_n,t]
$$
is proper.
Indeed, if $1\in J$, then after setting
$t=1/f$ we would get that a power of $f$
lies in $I$.
By the first part, $J$ is contained in a
maximal ideal corresponding to a point
$(a,b)\in k^{n+1}$.
Then $a\in V(I)$ and
$$
1-bf(a)=0,
$$
so $f(a)\neq 0$.
Thus $f$ does not vanish on $V(I)$, and
$f\notin I(V(I))$.
Therefore
$$
I(V(I))\subset \sqrt{I},
$$
and the theorem follows.
\end{proof}

\begin{exercise}
\label{commutative-algebra-exercise-trans-deg-hol-functions}
Let $A$ be the ring of complex-analytic
functions on $\mathbb{C}$,
and $k(A)$ its fraction field.
Prove that the transcendence degree
of $A$ over $\mathbb{C}$ is infinite.
\end{exercise}

\begin{proof}

\end{proof}

\section{Fractional ideal}
\label{commutative-algebra-section-fractional-ideals}

\noindent
We shall use the relationship
between fractional ideals and
projective modules in
\ref{commutative-algebra-lemma-invertible-maximal-ideal-dvr},
in the next section.

\medskip\noindent
{\bf AI-generated outline of the section.}
This section develops two related ideas. First,
it studies fractional ideals as submodules of
the fraction field $k(R)$ and shows that their
basic operations, such as products, inverses,
and endomorphism rings, can be described
concretely inside $k(R)$. Then it uses
$\operatorname{Hom}$ to connect that picture
with module theory, leading to the Dual Basis
Theorem and the characterization of finitely
generated projective modules.

\medskip\noindent
In short:
\begin{itemize}
\item It defines a fractional ideal as an
$R$-submodule of the fraction field $k(R)$
that becomes contained in $R$ after
multiplying by some nonzero element of $R$.
\item It shows that products, intersections,
inverses, and endomorphism rings of
fractional ideals are again fractional
ideals.
\item It identifies these objects with
$\operatorname{Hom}$ groups:
\begin{itemize}
\item $I^{-1} = \operatorname{Hom}_R(I,R)$
\item $R(I) = \operatorname{Hom}_R(I,I)$
\end{itemize}
\item The point is that maps between
fractional ideals are just multiplication by
fractions in $k(R)$.
\item Finally, it uses that perspective to
motivate the Dual Basis Theorem for modules,
which characterizes finitely generated
projective modules.
\end{itemize}

\noindent
So the section is really building a bridge
from concrete ideal arithmetic in $k(R)$ to
abstract module theory via $\operatorname{Hom}$.

\medskip\noindent
The link with projectivity is that a
fractional ideal is being treated as an
$R$-module, and projectivity is the condition
that this module has a finite dual
description.

\medskip\noindent
Concretely:
\begin{itemize}
\item For a fractional ideal $I$, the lecture
identifies $I^{-1}=\operatorname{Hom}_R(I,R)$.
\item The Dual Basis Theorem says a finitely
generated module is projective exactly when
the identity map can be written as a finite
sum of maps from $\operatorname{Hom}_R(I,R)$
paired with elements of $I$.
\item So projectivity is the condition that
$I$ has enough dual functionals to split off
from a free module.
\end{itemize}

\medskip\noindent
For fractional ideals, this becomes especially
concrete because the dual module is again a
fractional ideal living in $k(R)$. That is why
$\operatorname{Hom}$ is useful here: it turns
an abstract projectivity question into a
statement about multiplication by fractions.

\medskip\noindent
The usual takeaway is that a finitely
generated fractional ideal is projective
exactly when it behaves like an invertible
ideal, equivalently when it has a dual
fractional ideal that pairs with it the way a
basis and dual basis do.

\medskip\noindent
So the technical bridge is: fractional ideals
give concrete modules inside the fraction
field, and projectivity is the condition that
those modules admit a finite dual basis, i.e.
are locally free-like rather than just
arbitrary submodules of $k(R)$.

\medskip\noindent
A bit of history: Kummer: interested in $
\mathbb{Z}[q]$ with $q^p=1$.
He noted it has no prime decomposition.
Dedekind realised that
ideals are just subsets.



\begin{definition}
\label{commutative-algebra-definition-fractional-ideal}
Let $R$ be a ring without zero divisors,
and $k(R)$ its fraction field. A non-zero
$R$-submodule $I\subset k(R)$ is called a
{\it fractional ideal} of $R$ if for some
$a\in R$ one has $aI\subset R$.
\end{definition}

\begin{lemma}
\label{commutative-algebra-lemma-fractional-ideal-noetherian}
Let $R$ be a Noetherian ring, and
$I\subset k(R)$ an $R$-submodule. Then $I$
is a fractional ideal if and only if $I$ is
finitely generated.
\end{lemma}

\begin{proof}
If $I$ is finitely generated by elements
$a_1/b_1,\ldots,a_n/b_n\in k(R)$ with
$a_i,b_i\in R$, then
$$
(b_1\cdots b_n)I\subset R,
$$
so $I$ is fractional.

Conversely, if $aI\subset R$ for some
$a\in R$, then $aI$ is a submodule of the
Noetherian ring $R$, hence it is finitely
generated. It follows that $I$ is finitely
generated.
\end{proof}

\begin{remark}
\label{commutative-algebra-remark-fractional-product}
If $I_1$ and $I_2$ are fractional ideals,
then the set
$$
I_1I_2:=\left\{\sum_j x_jy_j \,\middle|\,
x_j\in I_1,\ y_j\in I_2\right\}
$$
is a fractional ideal, since if
$aI_1\subset R$ and $bI_2\subset R$, then
$ab(I_1I_2)\subset R$.
Also, $I_1\cap I_2$ is a fractional ideal,
since if $aI_1\subset R$ and $bI_2\subset R$,
then $ab(I_1\cap I_2)\subset R$.
\end{remark}

\begin{lemma}
\label{commutative-algebra-lemma-fractional-hom-ideal}
$$
J:=\{x\in k(R)\mid xI_1\subset I_2\}
$$
is a fractional ideal.
\end{lemma}

\begin{proof}
Step 1. Let $a,b\in R\setminus\{0\}$ be
elements such that $aI_1\subset R$ and
$bI_2\subset R$. Then
$$
abI_1\subset bR\subset I_2,
$$
hence $J$ is non-empty.

Step 2. For any non-zero elements
$c,d\in R$ such that $cI_1\subset R$ and
$dI_2\subset R$, and any $x\in J$, we have
$$
cdx=d(xc)\in dI_2\subset R,
$$
hence $cdx\in R$.
\end{proof}

\begin{lemma}
\label{commutative-algebra-lemma-fractional-inverse-r}
Let $I\subset R$ be a fractional ideal.
Then
$$
I^{-1}:=\{x\in k(R)\mid xI\subset R\}
$$
and
$$
R(I):=\{x\in k(R)\mid xI\subset I\}
$$
are fractional ideals.
Also, $R\subset R(I)$.
If, in addition, $I\subset R$, then
$R\subset I^{-1}$.
\end{lemma}

\begin{proof}
Let $a\in R\setminus\{0\}$ satisfy
$aI\subset R$. If $x\in I^{-1}$, then
$xI\subset R$, so $a x I\subset R$; hence
$ax\in I^{-1}$ and $I^{-1}$ is fractional.

Also, if $x\in R(I)$, then
$a(xI)=x(aI)\subset R$, so $ax\in R(I)$;
thus $R(I)$ is fractional.

Since $I$ is an $R$-module, one has
$R\subset R(I)$.

Finally, if $I\subset R$, then any $x\in R$
satisfies $xI\subset R$, hence $x\in I^{-1}$.
Thus $R\subset I^{-1}$.
\end{proof}

\noindent
Below are some useful ways
of writing the modules introduced above.
The point of viewing $J$ as
$\operatorname{Hom}_R(I_1,I_2)$ is that it
turns a condition about maps between
fractional ideals into a concrete
multiplication description inside $k(R)$.
That makes it easy to recognize familiar
objects: $I^{-1}$ becomes
$\operatorname{Hom}_R(I,R)$ and $R(I)$
becomes $\operatorname{Hom}_R(I,I)$, so one
can study these ideals using standard
module-theoretic ideas like finite generation
and endomorphism rings.

\begin{lemma}
\label{commutative-algebra-lemma-fractional-hom-iso}
$J\simeq \operatorname{Hom}_R(I_1,I_2)$.
\end{lemma}

\begin{proof}
Any $R$-linear map $\phi:I_1\to I_2$ becomes,
after tensoring with $k(R)$, an element of
$$
\operatorname{Hom}_{k(R)}(k(R),k(R))=k(R).
$$
So $\phi$ is multiplication by some
$z\in k(R)$.
The condition that $\phi(I_1)\subset I_2$
is exactly $zI_1\subset I_2$, so $z\in J$.
Conversely, any $z\in J$ defines an
$R$-linear map $I_1\to I_2$ by
multiplication.
\end{proof}

\begin{remark}
\label{commutative-algebra-remark-fractional-hom}
$$
I^{-1}=\operatorname{Hom}_R(I,R)
\quad\text{and}\quad
R(I)=\operatorname{Hom}_R(I,I).
$$
Indeed, any $R$-linear map
$\phi\in\operatorname{Hom}_R(I_1,I_2)$, after
tensoring with $k(R)$, becomes an element of
$$
\operatorname{Hom}_{k(R)}(k(R),k(R))=k(R).
$$
Therefore, $\phi$ is multiplication by some
$z\in k(R)$ which satisfies $zI_1=I_2$.
\end{remark}

\begin{remark}
\label{commutative-algebra-remark-fractional-integrally-closed}
If $R$ is Noetherian and integrally
closed, then $R(I)=R$. Indeed,
$R(I)=\operatorname{Hom}_R(I,I)$ is a finitely
generated $R$-module in $k(R)$.
\end{remark}

\noindent 
Recall the definition of projective
module, Definition
\ref{commutative-algebra-definition-projective-module}. The Dual Basis
Theorem below says that a module is finitely
generated projective exactly when its
identity map can be written as a finite sum
of rank-one pieces coming from $\Hom_R(M,R)$
and $M$. Equivalently, the natural map
$\Hom_R(M,R)\otimes_R M \to \Hom_R(M,M)$ is
an isomorphism precisely in that case.

\begin{theorem}[Dual basis theorem]
\label{commutative-algebra-theorem-dual-basis}
Let $M$ be an $R$-module. Consider the
natural map
$$
\psi:\Hom_R(M,R)\otimes_R M\to \Hom_R(M,M).
$$
Then the following are equivalent:
\begin{enumerate}
\item $\psi$ is an isomorphism.
\item $\mathrm{id}_M \in \text{Im}\psi$.
\item $M$ is projective and finitely
generated.
\end{enumerate}
\end{theorem}

\begin{proof}
Step 1. The map $\psi$ is an isomorphism for a free
finitely generated module, and therefore for
a direct summand of such a module. Hence
(iii) $\Rightarrow$ (i) $\Rightarrow$ (ii).

Step 2. Condition (ii) means there exist
maps $f_i:M\to R$ and elements $m_i\in M$
such that for any $m\in M$,
$$
\sum_i f_i(m)m_i=m.
$$
In particular, $M$ is finitely generated.

Step 3. Let $F$ be the free module generated
by the $m_i$, and define
$f:M\to F$ by
$$
f(m)=\sum_i f_i(m)m_i.
$$
Then $f$ is a section of the natural
surjection $F\to M$, so $M$ is projective.
Thus (ii) $\Rightarrow$ (iii).
\end{proof}

\noindent
The following proposition is very
much like Exercise
\ref{commutative-algebra-exercise-finit-gener-proje-iff-fact},
but this time with a local ring.

\begin{proposition}
\label{commutative-algebra-proposition-projective-local-free}
Let $P$ be a finitely generated projective
module over a local ring $R$. Then $P$ is a
free $R$-module.
\end{proposition}

\begin{proof}
Denote by $\mathfrak m$ the maximal ideal of
$R$. Let $e_1,\ldots,e_n$ be elements of
$P/\mathfrak m P$ generating $P/\mathfrak m P$.
Choose the $e_i$ so that their images in
$P/\mathfrak m P$ are linearly independent
over the field $R/\mathfrak m$.
By Nakayama's lemma, the $e_i$ generate $P$.

Since $P$ is projective, the projection from
the free module
$$
F:=\langle e_1,\ldots,e_n\rangle
\to P
$$
has a section $P\to F$.
Since its image modulo $\mathfrak m F$
generates $F/\mathfrak m F$, we have
$F=\operatorname{im}(\text{section})$,
and $P$ is a free $R$-module.
\end{proof}

\begin{definition}
\label{commutative-algebra-definition-invertible-fractional-ideal}
A fractional ideal $I\subset k(R)$ is
called {\it invertible} if
$II^{-1}=R$.
\end{definition}

\begin{theorem}
\label{commutative-algebra-theorem-invertible-fractional-ideal}
A fractional ideal is invertible if and only
if it is projective, and it is then finitely
generated.
\end{theorem}

\begin{proof}
Step 1. Clearly,
$I^{-1}\subset \operatorname{Hom}_R(I,R)$.
Therefore, $II^{-1}=R$ implies
$$
I\otimes_R I^{-1}\subset
\operatorname{Hom}_R(I,R)\otimes_R I,
$$
which implies projectivity by the Dual Basis
Theorem.

Step 2. Consider an $R$-linear map
$\phi:I_1\to I_2$. If we tensor $\phi$ with
$k(R)$, we obtain a $k(R)$-linear map
$$
I_1\otimes_R k(R)\to I_2\otimes_R k(R),
$$
with $I_i\otimes_R k(R)=k(R)$.
Clearly,
$\operatorname{Hom}_{k(R)}(k(R),k(R))=k(R)$.
Therefore, $\phi$ is expressed as
$v\mapsto \alpha v$ for some
$\alpha\in k(R)$.
This implies that the natural map
$I^{-1}\to \operatorname{Hom}_R(I,R)$ is an
isomorphism.

Step 3. If $I$ is projective, then
$$
\text{id}_I \in
\operatorname{Hom}_R(I,R)\otimes_R I,
$$
hence $II^{-1}\ni 1$, and $I^{-1}I=R$.
\end{proof}

\section{Discrete valuations}
\label{commutative-algebra-section-discrete-valuations}

\noindent
Discrete valuation rings are
the local rings of normal curves!!
The key proposition that
shows this is
Proposition 
\ref{commutative-algebra-proposition-integrally-closed-dvr}.

\medskip\noindent
The prototypical example of a DVR
is the $p$-adic numbers, described below.
The notion of discrete valuation is more
abstract, but we introduce it first.

\begin{definition}
\label{commutative-algebra-definition-discrete-valuation}
Let $K$ be a field. A {\it discrete valuation}
on
$K$ is a surjective map
$$
\nu:K^\times\to \mathbb{Z}
$$
such that
$$
\nu(xy)=\nu(x)+\nu(y)
$$
and
$$
\nu(x+y)\geq \min(\nu(x),\nu(y)).
$$
The {\it valuation ring} of $\nu$ is
$$
\{x\in K\mid \nu(x)\geq 0\}.
$$
\end{definition}

\begin{example}
\label{commutative-algebra-example-p-adic-valuation}
Let $R$ be a factorial ring, and
$p\in R$ a prime. Given $x\in k(R)$, write
the prime decomposition
$$
x=\prod_i p_i^{\alpha_i}
$$
where $\alpha_i\in \mathbb{Z}$ may be
negative, since $x$ belongs to the fraction
field. Define
$$
\nu_{p_1}(x):=\alpha_1.
$$
This function is called the {\it $p$-adic
valuation}. Its valuation ring consists of
fractions $a/b$ where $a,b\in R$ are
coprime and $b$ is not divisible by $p$.
\end{example}

\begin{exercise}
\label{commutative-algebra-exercise-p-adic-metric}
Show that for the $p$-adic valuation, the
function
$$
d(x,y):=p^{-\nu_p(x-y)}
$$
is a metric.
\end{exercise}

\noindent
This metric has two properties:
\begin{itemize}
\item $|x-y|\leq \text{max}(
|x-z|,|z-y|)$, (shouldn't it be min…?)
\item $|xy|=|x||y|$.
\end{itemize}

\noindent
We can extend this norm to the rational
numbers using the product property.
Now let $q \in \mathbb{Q}$
with $q=\frac{m}{n}$, $m,n$ coprime.
Then
$$
|q|\leq  1 
\iff n \not| p.
$$

\begin{definition}
\label{commutative-algebra-definition-p-adic-numbers}
The {\it $p$-adic numbers} is the ring
$\mathbb{Z}_p$, the completion  of $\mathbb{Z}$
with respect to this norm as a metric space.
Similarly, $\mathbb{Q}_p$
is the completion of $\mathbb{Q}$
with respect to the extension
of the norm described above.
\end{definition}

\begin{remark}
\label{commutative-algebra-remark-binary-decimals-etc}
$p$-adic numbers are for $p$
as binary numbers are for 2
and decimals for 10. Look for this!
\end{remark}

\begin{example}
\label{commutative-algebra-example-laurent-series-valuation}
Consider the ring of Laurent power series
$$
\mathbb{C}[[t]][t^{-1}],
$$
which is the fraction field of
$\mathbb{C}[[t]]$, namely
$k(\mathbb{C}[[t]])$.
If
$$
F=t^k+a_1t^{k+1}+a_2t^{k+2}+\cdots,
$$
define
$$
\nu(F):=k,
$$
the minimum exponent.
This makes sense because it's finite
seres in the negative exponents by
definition.

Then $\nu$ is a discrete valuation, with
$\nu(t)=1$ and $\nu(FG)=\nu(F)+\nu(G)$.
\end{example}

\begin{definition}
\label{commutative-algebra-definition-discrete-valuation-ring}
A ring $A$ without zero divisors is called
{\it a discrete valuation ring} if
$k(A)$ admits a discrete valuation such that
$A$ is its valuation ring.
\end{definition}

\begin{lemma}
\label{commutative-algebra-lemma-dvr-principal-maximal-ideal}
Let $R$ be a Noetherian local ring without
zero divisors, and let $\mathfrak{m}$ be its maximal
ideal. Then $\mathfrak{m}$ is principal if and only if
$R$ is a discrete valuation ring.
\end{lemma}

\begin{proof}
Let $R$ be a DVR, and choose an element $p$
with $\nu(p)=1$. For any $x\in R$ with
$\nu(x)=k$, we have $x=up^k$, where $u$ is
invertible. Hence every ideal of $R$ is
principal, so in particular $\mathfrak{m}$ is
principal.

Conversely, suppose that $\mathfrak{m}=(p)$
is principal. Then any element $x\in R$ which
is not divisible by $p$ is invertible. Hence
every nonzero $x\in k(R)$ can be written as
$x=up^k$, where $u$ is invertible and
$k\in \mathbb{Z}$. This defines a discrete
valuation on $k(R)$ with valuation ring $R$.
\end{proof}

\begin{lemma}[Corollary]
\label{commutative-algebra-lemma-dvr-basic-properties}
A discrete valuation ring is local,
Noetherian, integrally closed, all its ideals
are principal, and the Krull dimension of $R$
is $1$.
\end{lemma}

\begin{lemma}
\label{commutative-algebra-lemma-radical-ideal-local-dim1}
Let $R$ be a local ring of Krull dimension
$1$. Then any non-zero radical ideal of $R$
is its maximal ideal.
\end{lemma}

\begin{proof}
Let $I\subset R$ be a non-zero radical ideal.
Since $I$ is radical, the quotient ring
$R/I$ has nilradical $0$, by
Lemma~\ref{commutative-algebra-lemma-inter-maxim-ideal-nilra}.
Pulling this back to $R$, we see that
$I$ is the intersection of all prime ideals
of $R$ containing $I$.

Since $R$ has Krull dimension $1$, any chain
of prime ideals has length at most $2$. Because
$I\neq 0$, the only prime ideal of $R$
containing $I$ is the maximal ideal
$\mathfrak{m}$. Hence $I=\mathfrak{m}$.
\end{proof}

\begin{lemma}[Corollary]
\label{commutative-algebra-lemma-power-in-maximal-ideal}
Let $R$ be a Noetherian local ring of Krull
dimension $1$, and let $\mathfrak{m}$ be its
maximal ideal. Then for any $t\in\mathfrak{m}$
there exists $n$ such that
$\mathfrak{m}^n\subset (t)$.
\end{lemma}

\begin{proof}
The radical of $(t)$ is $\mathfrak{m}$, by the
previous lemma.
Spelling out the definition, this says
that $\forall x \in R$, if
$t|x^n$ for some $n$, 
then $x \in \mathfrak{m}$.
Suppose that $\mathfrak{m}$ is generated
by $a_1,…,a_k$ using that $R$
is Noetherian.
We need to show that there is $n$
such that for any
$\left(\sum r_ia_i\right)^n$
has a factor $t$.
[Details missing… perhaps there's an error]

A standard Noetherian fact is that if
$\sqrt{(t)}=\mathfrak{m}$, then some power of the
radical is contained in $(t)$. Therefore there
exists $n>0$ such that
$\mathfrak{m}^n\subset (t)$.
\end{proof}

\noindent
Recall that we showed before
that a projective module over a local
ring is free --- this is a variation
for local rings of the fact
that projective modules
are direct summands of free modules.
I think that uses Nakayama somehow.

\begin{proposition}
\label{commutative-algebra-proposition-projective-maximal-ideal-dvr}
Let $R$ be a Noetherian local ring, and let
$\mathfrak{m}$ be its maximal ideal. Assume
that $\mathfrak{m}$ is projective as an
$R$-module. Then $R$ is a DVR.
\end{proposition}

\begin{proof}
Since $R$ is Noetherian and local,
$\mathfrak{m}$ is free. Choose elements
$p_1,\dots,p_n$ freely generating
$\mathfrak{m}$. Then
$$
\mathfrak{m}=\bigoplus_i (p_i),
$$
which is impossible unless $n=1$, because
$p_1p_2\in (p_1)\cap(p_2)$. Therefore,
$\mathfrak{m}$ is a principal ideal, and
$R$ is a DVR by Lemma~\ref{commutative-algebra-lemma-dvr-principal-maximal-ideal}.
\end{proof}

\noindent
Here we use the fractional ideal-projective
module construction from the past section!

\begin{lemma}[Corollary]
\label{commutative-algebra-lemma-invertible-maximal-ideal-dvr}
Let $R$ be a Noetherian local ring, and let
$\mathfrak{m}$ be its maximal ideal. Assume
that $\mathfrak{m}$ is invertible as a
fractional ideal. Then $R$ is a DVR.
\end{lemma}

\begin{proof}
Invertible fractional ideals are projective.
\end{proof}

\noindent
The following proposition is the
key step in proving that
normality in dimension 1 is smoothness!!!
A remark: ``The converse statement
[of the following proposition]
is implied by Corollary 1''…
but what's that one?

\begin{proposition}
\label{commutative-algebra-proposition-integrally-closed-dvr}
Let $R$ be an integrally closed Noetherian
local ring of Krull dimension $1$. Then $R$
is a discrete valuation ring.
\end{proposition}

\begin{proof}
\textbf{Step 1.} Let $I\subset k(R)$ be a
fractional ideal, and $R(I)\subset k(R)$ the
ring of all $x\in k(R)$ such that $xI\subset I$.
Since $R(I)\subset \operatorname{Hom}_R(I, I)$,
it is finitely generated as an $R$-module,
hence $R(I)=R$.

\textbf{Step 2.} Let $\mathfrak{m}$ be the
maximal ideal of $R$. Then
$R\subset \mathfrak{m}\mathfrak{m}^{-1}\subset
\mathfrak{m}$. Since any element of
$R\setminus \mathfrak{m}$ is invertible, we
have either $\mathfrak{m}\mathfrak{m}^{-1}=R$
or $\mathfrak{m}\mathfrak{m}^{-1}=\mathfrak{m}$.
In the second case,
$\mathfrak{m}^{-1}=R$ by Step 1, because in
this case $\mathfrak{m}^{-1}\subset R(\mathfrak{m})$.
In the first case, the fractional ideal
$\mathfrak{m}$ is invertible.

\textbf{Step 3.} By Krull lemma
(which is a version
of Nakayama lemma, 
namely Lemma~\ref{commutative-algebra-lemma-nakayama-first}),
$\mathfrak{m}\supset \mathfrak{m}^2\supset
\cdots \supset \mathfrak{m}^n$.
Let $t\in \mathfrak{m}\setminus \mathfrak{m}^2$.
By Corollary~\ref{commutative-algebra-lemma-power-in-maximal-ideal},
$\mathfrak{m}^n\subset (t)$ for some $n>0$;
choose the minimal $n$ with this property, and
let $x\in \mathfrak{m}^{n-1}\setminus (t)$.
Let
$$
y:=\frac{x}{t}.
$$
Since $x\notin (t)$, we have $y\notin R$.
Clearly, $x\mathfrak{m}\subset \mathfrak{m}^n
\subset (t)$. This gives
$y\mathfrak{m}\subset R$, hence
$\mathfrak{m}^{-1}\neq R$. Then Step 2 implies
that $\mathfrak{m}^{-1}\mathfrak{m}=R$, and
Corollary~\ref{commutative-algebra-lemma-invertible-maximal-ideal-dvr}
implies that $R$ is a DVR.
\end{proof}

\begin{lemma}
\label{commutative-algebra-lemma-dvr-quotient-powers}
Let $R$ be a discrete valuation ring, and let
$\mathfrak{m}$ be its maximal ideal. Set
$k:=R/\mathfrak{m}$. Then
$\mathfrak{m}^d/\mathfrak{m}^{d+1}$ is
1-dimensional as a $k$-vector space for all
$d\in \mathbb{Z}_{\geq 0}$.
\end{lemma}

\begin{proof}
Let $p$ be the generator of $\mathfrak{m}$,
and let $L_{p^d}(x):=p^d x$. Then
$L_{p^d}:R\to \mathfrak{m}^d$ is an isomorphism
which maps $\mathfrak{m}$ to $\mathfrak{m}^{d+1}$.
Therefore
$$
\mathfrak{m}^d/\mathfrak{m}^{d+1}\cong R/\mathfrak{m}=k,
$$
so it is 1-dimensional over $k$.
\end{proof}

\begin{theorem}
\label{commutative-algebra-theorem-dvr-characterizations}
Let $R$ be a Noetherian local ring without
zero divisors, $\mathfrak{m}$ its maximal
ideal, and $k:=R/\mathfrak{m}$ its residue
field. Assume that $R$ has Krull dimension
$1$. Then the following are equivalent.
\begin{enumerate}
\item[$(i)$] $R$ is a discrete valuation ring.
\item[$(ii)$] $R$ is integrally closed.
\item[$(iii)$] $\mathfrak{m}$ is a principal ideal.
\item[$(iv)$] $\dim_k \mathfrak{m}/\mathfrak{m}^2 = 1$.
\item[$(v)$] Every non-zero ideal of $R$ 
is a power of $\mathfrak{m}$.
\item[$(vi)$]
All non-zero ideals are principal, 
generated by
some $p \in R$. More precisely,
there exists $p\in R$ such that every 
non-zero ideal of $R$ is generated by $p^k$, 
for some $k\in \mathbb{Z}_{>0}$.
\end{enumerate}
\end{theorem}

\begin{proof}
\textbf{Step 1.} $(i)$ implies $(ii)$,
$(iii)$, $(v)$ and $(vi)$ as follows from
Proposition~
\ref{commutative-algebra-proposition-integrally-closed-dvr}.
$(i)$ implies $(iv)$ by
Lemma~\ref{commutative-algebra-lemma-dvr-quotient-powers}.
Therefore, $(i)$ implies $(ii)$--$(vi)$.

\textbf{Step 2.} $(iii)\Leftrightarrow (i)$
by Lemma~
\ref{commutative-algebra-lemma-dvr-principal-maximal-ideal}.
$(ii)$ implies $(i)$ by
Proposition~
\ref{commutative-algebra-proposition-integrally-closed-dvr}.
Also, $(v)$ implies $(vi)$ which implies that
$R$ is factorial which implies $(ii)$.
We proved that all $(i)$--$(vi)$ are
equivalent, except maybe $(iv)$.

\textbf{Step 3.} Now we prove that $(iv)$
implies $(iii)$. Let $p$ be a generator of
$\mathfrak{m}$ modulo $\mathfrak{m}^2$.
Then the natural map
$$
(p)/(p)\mathfrak{m}
\to \mathfrak{m}/\mathfrak{m}^2
$$
is surjective. By Nakayama lemma, for any
morphism $\varphi:M\to N$ of finitely
generated modules over a local ring,
surjectivity of the map
$\varphi:M/\mathfrak{m}M\to N/\mathfrak{m}N$
implies the surjectivity of $\varphi$,
hence $\mathfrak{m}$ is a principal ideal.
\end{proof}

\section{p-adic numbers}
\label{commutative-algebra-section-p-adic-numbers}

\begin{definition}
\label{commutative-algebra-definition-norm-on-ring}
A {\it norm} on a ring $A$ is a
function
$$
A \to \mathbb R_{\geq 0},
$$
denoted $r \mapsto ||r||$, such that
$$
||x+y|| \leq ||x||+||y||,
$$
$$
||xy||=||x||||y||
$$
and
$$
||x||=0 \Leftrightarrow x=0.
$$
\end{definition}

\begin{remark}
\label{commutative-algebra-remark-usual-absolute-value-norm}
The usual absolute values on
$\mathbb R$ and $\mathbb C$ are
examples of norms.
\end{remark}

\begin{definition}
\label{commutative-algebra-definition-discrete-valuation-duplicate-2}
A {\it discrete valuation} on a ring
$A$ is a function
$$
\nu:A \to \mathbb Z \cup \{\infty\}
$$
such that
$$
\nu(ab)=\nu(a)+\nu(b),
$$
$$
\nu(a+b)\geq \min(\nu(a),\nu(b)),
$$
and
$$
\nu(x)=\infty \Leftrightarrow x=0.
$$
\end{definition}

\begin{remark}
\label{commutative-algebra-remark-additive-valuation-convention}
We use the additive convention for
valuations. Thus products become sums.
\end{remark}

\begin{exercise}
\label{commutative-algebra-exercise-norm-from-discrete-valuation}
Let $\nu$ be a discrete valuation on
a ring. Prove that
$$
x \mapsto u^{-\nu(x)}
$$
is a norm, for any
$u \in \mathbb R_{>1}$.
\end{exercise}

\begin{proof}
We use the convention $u^{-\infty}=0$.
Define
$$
||x||=u^{-\nu(x)}.
$$
Then $||x||=0$ if and only if
$\nu(x)=\infty$, which happens if and
only if $x=0$.

For the product property, we use
additivity of the valuation:
$$
\begin{aligned}
||xy||
&=u^{-\nu(xy)}\\
&=u^{-\nu(x)-\nu(y)}\\
&=u^{-\nu(x)}u^{-\nu(y)}\\
&=||x||||y||.
\end{aligned}
$$

For the triangle inequality, the
valuation inequality gives
$$
\nu(x+y)\geq \min(\nu(x),\nu(y)).
$$
Since $u>1$, this implies
$$
u^{-\nu(x+y)}
\leq
u^{-\min(\nu(x),\nu(y))}.
$$
Therefore
$$
||x+y||
\leq
u^{-\min(\nu(x),\nu(y))}.
$$
Now the minimum is one of the two
numbers $\nu(x),\nu(y)$. If
$$
\min(\nu(x),\nu(y))=\nu(x),
$$
then
$$
u^{-\min(\nu(x),\nu(y))}
=
u^{-\nu(x)}
\leq
u^{-\nu(x)}+u^{-\nu(y)}.
$$
The other case is identical. Hence
$$
||x+y|| \leq ||x||+||y||.
$$
So $x \mapsto u^{-\nu(x)}$ is a norm.
\end{proof}

\begin{definition}
\label{commutative-algebra-definition-p-adic-valuation}
Let $n/m \in \mathbb Q$, where $n,m$
are coprime integers. The {\it
$p$-adic valuation} $\nu_p(n/m)$ is
equal to $\alpha$ if
$$
n=p^\alpha n',
$$
where $\alpha>0$ and $n'$ is coprime
with $p$; to $-\alpha$ if
$$
m=p^\alpha m',
$$
where $\alpha>0$ and $m'$ is coprime
with $p$; to $0$ if $m$ and $n$ are
coprime with $p$; and to $\infty$ if
$n=0$.
\end{definition}

\begin{exercise}
\label{commutative-algebra-exercise-p-adic-valuation-discrete}
Prove that the $p$-adic valuation is
a discrete valuation.
\end{exercise}

\begin{proof}
Let's consider cases.
To prove the product rule, i.e.
that $\nu\left(\frac{n}{m}
\frac{n'}{m'}\right)
=\nu\left(\frac{n}{m}\right)+
\nu\frac{n'}{m'}$, consider:
\begin{itemize}
\item ($\nu
\left(\frac{n}{m}\right)=\nu(n),
\nu\left(\frac{n'}{m'}\right)=\nu(n')$.)
Then $\nu$(product)=$\nu(n)+\nu(n')$.
\item Similar cases give the result
with no problems.
\end{itemize}

\noindent
Now, the inequality property requires
an observation. In the cases
$\nu\left(\frac{n}{m}\right)=\nu(n),
\nu\left(\frac{n'}{m'}\right)=\nu(n')$
and $\nu\left(\frac{n}{m}\right)=-\nu(m),
\nu\left(\frac{n'}{m'}\right)=-\nu(m')$
we obtain equality
(WRONG! if the valuations
are the same we could 
get something not equality 
--- but the required inequality
would hold still)
$$
\nu\left(\frac{n}{m}
+\frac{n'}{m'}\right)
=\text{min}\left\{
\nu\left(\frac{n}{m}\right),
\nu\left(\frac{n'}{m'}\right)\right\}.
$$
This is because the
exponent of $p$ factors through the minimum.
Pending: look more carefully at the
case where the two valuations coincide.

Now consider the case
$\nu\left(\frac{n}{m}\right)=\nu(n),
\nu\left(\frac{n'}{m'}\right)=-\nu(m')$.
Do simple sum:
$$
p^{\nu(n)}+p^{-\nu(m')}
p^{\nu(n)}
\frac{p^{\nu(m')}}{p^{\nu(m')}}
+\frac{1}{p^{\nu(m')}}
=
\frac{p^{\nu(n)+\nu(m')}+1}
{p^{\nu(m')}},
$$
which shows that the valuation is
$-\nu(m')$ since the numerator
has no $p$-factor.
Since both $\nu(n)$ and $\nu(m')$ 
are positive number, the inequality follows.

The last property of valuations 
is verified trivially.


\medskip\noindent
AI-proof for inequality part:
Assume $\nu_p(x)\leq \nu_p(y)$.
Write
$$
x=p^r a,\qquad y=p^s b,
$$
where $r=\nu_p(x)$, $s=\nu_p(y)$,
and $\nu_p(a)=\nu_p(b)=0$.
Then
$$
x+y=p^r(a+p^{s-r}b).
$$
Since $s-r\geq 0$, the second factor
has nonnegative valuation. Hence
$$
\nu_p(x+y)\geq r
=
\min(\nu_p(x),\nu_p(y)).
$$
\end{proof}

\begin{remark}
\label{commutative-algebra-remark-norm-defines-metric}
The norm defines a metric on $R$ by
$$
d(x,y):=||x-y||.
$$
\end{remark}

\begin{exercise}
\label{commutative-algebra-exercise-ring-structure-completion}
Let $R$ be a ring equipped with a norm.
Construct the ring structure on its
metric completion $\overline R$,
compatible with the ring structure on
$R$. Prove that $\overline R$ is
equipped with a norm compatible with
the norm on $R$.
\end{exercise}

\begin{definition}
\label{commutative-algebra-definition-completion-normed-ring}
The normed ring $\overline R$ is called
the {\it completion} of $R$ with
respect to the norm.
\end{definition}

\begin{exercise}
\label{commutative-algebra-exercise-completion-of-field}
Let $R$ be a normed ring, and let
$\overline R$ be its completion.
Assume that $R$ is a field. Prove that
$\overline R$ is also a field.
\end{exercise}

\begin{definition}
\label{commutative-algebra-definition-completion-valuation}
Let $\nu$ be a valuation on a ring
$A$. The {\it completion of $A$ with
respect to $\nu$} is its completion
with respect to the norm
$$
x \mapsto u^{-\nu(x)}
$$
for some $u>1$.
\end{definition}

\begin{exercise}
\label{commutative-algebra-exercise-completion-independent-u}
Prove that the completion is
independent of
$u \in \mathbb R_{>1}$.
\end{exercise}

\begin{definition}
\label{commutative-algebra-definition-p-adic-integers-rationals}
The completion of $\mathbb Z$ with
respect to $\nu_p$ is called the
{\it ring of $p$-adic integers},
denoted $\mathbb Z_p$. The completion
of $\mathbb Q$ with respect to $\nu_p$
is called the {\it field of $p$-adic
rationals}, denoted $\mathbb Q_p$.
\end{definition}

\begin{exercise}
\label{commutative-algebra-exercise-summable-distances-cauchy}
Let $(X,d)$ be a metric space, and let
$\{a_i\}$ be a sequence of points in
$X$. Suppose that the series
$$
\sum_i d(a_i,a_{i-1})
$$
converges. Prove that $\{a_i\}$ is a
Cauchy sequence. Is the converse true?
\end{exercise}

\begin{exercise}
\label{commutative-algebra-exercise-series-converges-zp}
Prove that for any sequence of integers
$a_k$, the series
$$
\sum_k a_kp^k
$$
converges in $\mathbb Z_p$.
\end{exercise}

\begin{exercise}
\label{commutative-algebra-exercise-geometric-series-zp}
Prove that
$$
(1-p)\left(\sum_{k=0}^{\infty}p^k\right)=1
$$
in $\mathbb Z_p$.
\end{exercise}

\begin{exercise}
\label{commutative-algebra-exercise-prime-to-p-invertible-zp}
Prove that every integer not divisible
by $p$ is invertible in $\mathbb Z_p$.
\end{exercise}

\begin{exercise}
\label{commutative-algebra-exercise-qp-element-fraction-zp}
Let $x \in \mathbb Q_p$. Prove that
$$
x=\frac{x'}{p^k},
$$
where $x' \in \mathbb Z_p$.
\end{exercise}

\begin{exercise}
\label{commutative-algebra-exercise-root-n-limit}
Prove that
$$
\lim_{n \to \infty}\sqrt[n]{n}=1.
$$
Here the limit is taken in $\mathbb R$
with the usual metric.
\end{exercise}

\begin{definition}
\label{commutative-algebra-definition-non-archimedean-norm}
A norm $\nu$ on a ring $R$ is called
{\it non-Archimedean} if
$$
\nu(x+y)\leq \max(\nu(x),\nu(y))
$$
for all $x,y$. Otherwise, the norm is
called {\it Archimedean}.
\end{definition}

\begin{exercise}
\label{commutative-algebra-exercise-non-archimedean-unit-ball}
Let $\nu$ be a norm on $\mathbb Q$.
Prove that $\nu$ is non-Archimedean if
and only if $\mathbb Z$ is contained in
the unit ball.
\end{exercise}

\begin{exercise}
\label{commutative-algebra-exercise-non-archimedean-prime-ideal}
Let $\nu$ be a non-Archimedean norm on
$\mathbb Z$. Consider the set
$$
\mathfrak m \subset \mathbb Z
$$
consisting of all integers $n$ such
that $\nu(n)<1$. Deduce from the
non-Archimedean property that
$\mathfrak m$ is an ideal in
$\mathbb Z$. Prove that
$\mathfrak m$ is prime.
\end{exercise}

\begin{exercise}
\label{commutative-algebra-exercise-non-archimedean-powers}
Let $\nu$ be a non-Archimedean norm on
$\mathbb Q$, and let
$$
\mathfrak m=(p)
=
\{0,\pm p,\pm 2p,\pm 3p,\ldots\}
$$
be the ideal constructed above. Prove
that there exists a real number
$\lambda>1$ such that
$$
\nu(n)=\lambda^{-k}
$$
for every
$$
n=p^kr,
$$
where $r$ is coprime with $p$.
\end{exercise}

\begin{exercise}
\label{commutative-algebra-exercise-norm-bound-by-log}
Let $\nu$ be a norm on $\mathbb Q$ such
that $\nu(2)\leq 1$. Prove that
$$
\nu(a)<\log_2(a)+1
$$
for every integer $a>0$.
\end{exercise}

\begin{exercise}
\label{commutative-algebra-exercise-norm-two-less-than-one}
Let $\nu$ be a norm on $\mathbb Q$ such
that $\nu(2)\leq 1$. Prove that
$$
\nu(a)\leq 1
$$
for every integer $a>0$. In particular,
$\nu$ is non-Archimedean.
\end{exercise}

\begin{exercise}
\label{commutative-algebra-exercise-archimedean-cauchy-sequence}
Let $a_i$ be a sequence of rational
numbers of the form
$$
\frac{x}{2^n},
$$
which is Cauchy with respect to the
usual Euclidean norm on $\mathbb R$.
Suppose that the norm $\nu$ on
$\mathbb Q$ is Archimedean. Prove that
$\nu(a_i)$ is a Cauchy sequence.
\end{exercise}

\begin{exercise}
\label{commutative-algebra-exercise-archimedean-norm-extension}
Show that any Archimedean norm $\nu$ on
$\mathbb Q$ extends to a continuous
function on $\mathbb R$ satisfying
$$
\nu(xy)=\nu(x)\nu(y).
$$
Prove that $\nu$ is given by
$$
x \mapsto |x|^\lambda
$$
for some constant $\lambda>0$. Express
$\lambda$ in terms of $\nu(2)$.
\end{exercise}

\begin{remark}
\label{commutative-algebra-remark-ostrowski-theorem}
The previous exercises give a complete
classification of norms on
$\mathbb Q$: every norm is either a
power of a $p$-adic norm or a power of
the usual absolute value. This
classification is called
{\it Ostrowski's theorem}.
\end{remark}

\begin{definition}
\label{commutative-algebra-definition-inverse-limit}
Let $\{C_i\}$ be a set of vector spaces
or abelian groups associated with the
vertices of a commutative diagram $C$.
Let
$$
V \subset \prod_i C_i
$$
be the set of all collections
$v_i \in C_i$ such that for any arrow
$$
\varphi_{ij}:C_i \to C_j
$$
in the diagram, one has
$$
\varphi_{ij}(v_i)=v_j.
$$
The space $V$ is called the
{\it inverse limit} of the diagram
$\{C_i\}$. The same notion is also
called a {\it limit} or
{\it projective limit}. It is denoted
$$
\varprojlim C_i.
$$
\end{definition}

\begin{exercise}
\label{commutative-algebra-exercise-inverse-limit-intersection}
Let
$$
\cdots \to C_3 \to C_2 \to C_1
$$
be a diagram with all arrows injective.
Prove that $\varprojlim C_i$ is an
intersection of all $C_i$.
\end{exercise}

\begin{exercise}
\label{commutative-algebra-exercise-inverse-limit-initial-object}
Let
$$
C_1 \to C_2 \to C_3 \to \cdots \to C_n
$$
be a diagram. Prove that
$$
\varprojlim C_i=C_1.
$$
\end{exercise}

\begin{exercise}
\label{commutative-algebra-exercise-nonzero-diagram-zero-limit}
Find an example of a diagram
$$
\cdots \to C_1 \to C_2 \to C_3 \to
\cdots
$$
where all spaces $C_i$ are non-zero,
but the limit $\varprojlim C_i$
vanishes.
\end{exercise}

\begin{definition}
\label{commutative-algebra-definition-cofiltered-diagram}
A diagram $C$ is called
{\it cofiltered} if for any two
vertices $C_i,C_j$, there exists a
third vertex $C_k$ and sequences of
arrows leading from $C_k$ to $C_i$
and to $C_j$.
\end{definition}

\begin{exercise}
\label{commutative-algebra-exercise-inverse-limit-ring}
Let $C$ be a commutative diagram of
vector spaces $C_i$, with all $C_i$
equipped with a ring structure, and all
arrows ring homomorphisms. Suppose that
the diagram $C$ is cofiltered. Prove
that $\varprojlim C_i$ admits a
structure of a ring, such that the
natural maps
$$
\varprojlim C_i \to C_i
$$
are ring homomorphisms.
\end{exercise}

\begin{definition}
\label{commutative-algebra-definition-completion-topological-group}
Let $G$ be a topological group. A
sequence $\{g_i\}\subset G$ is called
a {\it Cauchy sequence} if for any
neighborhood $U$ of the identity there
exists $g \in G$ such that $gU$
contains almost all elements of
$\{g_i\}$. Cauchy sequences
$\{a_i\},\{b_i\}\subset G$ are called
{\it equivalent} if the sequence
$$
a_1,b_1,a_2,b_2,\ldots
$$
is Cauchy. The {\it completion} of
$G$ is the set of all equivalence
classes of Cauchy sequences.
\end{definition}

\begin{exercise}
\label{commutative-algebra-exercise-topological-group-completion}
Define the following topology on the
completion $\overline G$ of a
topological group $G$. The base of open
sets in $\overline G$ is obtained from
the open sets $U \subset G$, with
$\overline U \subset \overline G$
defined as the set of all sequences
$\{a_i\}$ such that almost all $a_i$
belong to $U$. Define the
multiplicative structure on
$\overline G$ by
$$
\{a_i\}\{b_i\}=\{a_ib_i\}.
$$
Prove that the completion
$\overline G$ of a topological group
$G$ is a topological group.
\end{exercise}

\begin{exercise}
\label{commutative-algebra-exercise-dense-map-to-completion}
Consider the map from a topological
group $G$ to its completion
$\overline G$, taking $g \in G$ to the
constant sequence
$$
g,g,g,g,\ldots .
$$
Prove that this defines a homomorphism
$$
G \to \overline G
$$
of topological groups, and that $G$ is
dense in $\overline G$.
\end{exercise}

\begin{definition}
\label{commutative-algebra-definition-i-adic-topology}
Let $I \subset R$ be an ideal in a
ring. The {\it $I$-adic topology} on
$R$ is the topology with base of open
sets given by
$$
a+I^k \subset R
$$
for all $a \in R$ and
$k \in \mathbb Z_{>0}$. The
{\it $I$-adic completion} of $R$ is
the completion of
$$
R/\bigcap_{k>0}I^k
$$
in this topology.
\end{definition}

\begin{exercise}
\label{commutative-algebra-exercise-i-adic-completion-ring}
Prove that the $I$-adic completion
$\overline R$ of $R$ is equipped with
a ring structure in such a way that
the natural map
$$
R \to \overline R
$$
is a homomorphism. Prove that
$\overline R$ is isomorphic to the
inverse limit
$$
\varprojlim R/I^k.
$$
\end{exercise}

\begin{exercise}
\label{commutative-algebra-exercise-p-adic-topology-completion}
Define the $p$-adic norm on
$\mathbb Z$ as the map
$$
\nu_p:\mathbb Z \to \mathbb Q
$$
taking $x \in \mathbb Z$ to $p^{-n}$,
where $n$ is the maximal number such
that $x$ is divisible by $p^n$ in
$\mathbb Z$. Consider the topology
associated with the ideal $(p)$, called
the $p$-adic topology. Prove that the
completion of $\mathbb Z$ with respect
to the norm $\nu_p$ is equal to the
completion of $\mathbb Z$ with respect
to the $p$-adic topology.
\end{exercise}

\begin{exercise}
\label{commutative-algebra-exercise-zp-as-inverse-limit}
Let $p$ be a prime number. Define the
ring $\mathbb Z_p$ of $p$-adic numbers
as
$$
\varprojlim \mathbb Z/p^k\mathbb Z.
$$

\begin{enumerate}
\item Prove that any power series
$$
\sum_{i=0}^{\infty}a_ip^i
$$
converges in $\mathbb Z_p$, for any
$a_i \in \mathbb Z$. Prove that any
element in $\mathbb Z_p$ can be
obtained as a sum of such a series.

\item Prove that $\mathbb Z_p$ is a
local ring, and every
$x \in \mathbb Z_p$ which is not
divisible by $p$ is invertible in
$\mathbb Z_p$.
\end{enumerate}
\end{exercise}

\begin{definition}
\label{commutative-algebra-definition-profinite-topology}
Let $G$ be a group, and let $S$ be the
set of all normal subgroups of finite
index. The {\it profinite topology} on
$G$ is the topology with base of open
sets given by all
$$
gN,
$$
where $g \in G$ and $N \in S$. The
{\it profinite completion} of $G$ is
the completion of
$$
G/\bigcap_{N\in S}N
$$
with respect to this topology.
\end{definition}

\begin{exercise}
\label{commutative-algebra-exercise-profinite-completion-limit}
Prove that the profinite completion of
$G$ is
$$
\varprojlim G/N,
$$
where the limit is taken over all
normal subgroups $N \subset G$ of
finite index.
\end{exercise}

\begin{exercise}
\label{commutative-algebra-exercise-profinite-completion-compact}
Prove that the profinite completion of
a group is compact and Hausdorff.
\end{exercise}

\begin{exercise}
\label{commutative-algebra-exercise-profinite-completion-z}
Prove that the profinite completion of
$\mathbb Z$ is
$$
\prod_p \mathbb Z_p,
$$
where the product is taken over all
prime numbers.
\end{exercise}

\begin{exercise}
\label{commutative-algebra-exercise-equivalence-p-adic-definitions}
Prove that this definition of
$p$-adics is equivalent to the one
given using metric completion.
\end{exercise}

\section{Dedekind rings}
\label{commutative-algebra-section-dedekind-rings}

\noindent
Dedekind rings are the coordinate algebras
of normal curves, which, by the 
previous section, are smooth.

\medskip\noindent
Historically, Dedekind and Kummer were
trying to understand factorization in
number rings where unique factorization
of elements fails.
Their idea was to replace elements by
ideals, because ideals factor uniquely in
the Dedekind setting, and this gives a
manageable theory of arithmetic in
number fields.
Krull later abstracted this ideal-theoretic
picture and studied it in dimension $1$,
where the geometry is especially simple.

\begin{definition}
\label{commutative-algebra-definition-dedekind-ring}
Let $R$ be a ring without zero
divisors.
Assume that all ideals $I\subset R$
are invertible as fractional ideals.
Then $R$ is called a Dedekind ring,
or a Dedekind domain.
\end{definition}

\begin{lemma}
\label{commutative-algebra-lemma-free-ideal-principal}
Let $\mathfrak{p}$ be an ideal in a
ring $R$.
Assume that $\mathfrak{p}$ is free as
an $R$-module.
Then $\mathfrak{p}$ is principal.
\end{lemma}

\begin{proof}
Let $e_1,\ldots,e_n$ be a free basis of
$\mathfrak{p}$.
Then
$$
\mathfrak{p}
=\bigoplus_i R e_i.
$$
However, $e_1e_2\in Re_1\cap Re_2$,
which is impossible because
$$
Re_1\cap Re_2=0.
$$
Hence $n=1$, so $\mathfrak{p}$ is
principal.
\end{proof}

\begin{proposition}
\label{commutative-algebra-proposition-dedekind-local-dvr}
Let $\mathfrak{m}$ be a maximal ideal of a
Dedekind ring $R$, and let $R_{\mathfrak{m}}$
be its localization.
Then $R_{\mathfrak{m}}$ is a discrete valuation
ring.
\end{proposition}

\begin{proof}
Every ideal of $R_{\mathfrak{m}}$ is the
localization of an ideal of $R$.
Localization of a projective module is
projective, hence every ideal of
$R_{\mathfrak{m}}$ is projective.
Since $R_{\mathfrak{m}}$ is local, every
projective module over it is free.
Therefore, every ideal of $R_{\mathfrak{m}}$
is free.

By Lemma~\ref{commutative-algebra-lemma-free-ideal-principal},
every free ideal is principal.
Hence $R_{\mathfrak{m}}$ is a principal ideal
ring, so it is a discrete valuation ring.
\end{proof}

\begin{lemma}[Corollary]
\label{commutative-algebra-lemma-dedekind-krull-dim-one}
A Dedekind ring $R$ has Krull dimension $1$.
\end{lemma}

\begin{proof}
Let $\mathfrak{m}$ be a maximal ideal of
$R$.
By Proposition
\ref{commutative-algebra-proposition-dedekind-local-dvr},
$R_{\mathfrak{m}}$ is a discrete valuation
ring.
Hence $R_{\mathfrak{m}}$ has only two prime
ideals.
Therefore, every chain of prime ideals in $R$
has length at most $1$, so $R$ has Krull
dimension $1$.
\end{proof}

\noindent
The upshot about the following lemma is:
prime ideals correspond to points
of the scheme $\Spec(R)$.
Localizations at each prime are
the functions on the scheme which
are regular near the point.
The functions which are regular
near every point are just regular functions,
which are the elements of $R$.

\begin{lemma}
\label{commutative-algebra-lemma-intersection-localizations}
Let $R$ be a ring without zero divisors,
and let $\mathfrak{p}$ range over the
maximal ideals of $R$.
Then
$$
\bigcap_{\mathfrak{p}} R_{\mathfrak{p}} = R.
$$
\end{lemma}

\begin{proof}
Let $x$ be an element of
$\bigcap_{\mathfrak{p}} R_{\mathfrak{p}}$.
Let $I$ be the set of all $y\in R$ such
that $yx\in R$.
We will prove that $I=R$, which implies
that $x\in R$.

Assume that $I\neq R$.
Then there exists a maximal ideal
$\mathfrak{m}\supset I$.
For any $a\notin\mathfrak{m}$, we have
$ax\notin R$ by the definition of $I$.
On the other hand, since
$x\in R_{\mathfrak{m}}$, there exists
$a\notin\mathfrak{m}$ such that
$ax\in R$, a contradiction.
Therefore, $I=R$, and so $x\in R$.
\end{proof}

\noindent
In the following lemma,
Noetherianity follows from fractionality
of the local rings
and integral-closedness from
DVR-ness
of the local rings.

\begin{lemma}
\label{commutative-algebra-lemma-dedekind-noetherian-integrally-closed}
Any Dedekind ring $R$ is Noetherian and
integrally closed.
\end{lemma}

\begin{proof}
Step 1. Any invertible ideal is fractional,
hence finitely generated.
Therefore, $R$ is Noetherian.

Step 2. Let $\mathfrak{p}$ be a prime ideal
of $R$.
By Proposition
\ref{commutative-algebra-proposition-dedekind-local-dvr},
$R_{\mathfrak{p}}$ is a discrete valuation
ring, hence integrally closed.
Now let $x\in k(R)$ be integral over $R$.
Then $x\in R_{\mathfrak{p}}$ for every
prime ideal $\mathfrak{p}$, because
$R_{\mathfrak{p}}$ is integrally closed.
By Lemma
\ref{commutative-algebra-lemma-intersection-localizations},
$x\in R$.
\end{proof}

\noindent
Conversely, we have:

\begin{theorem}
\label{commutative-algebra-theorem-noetherian-integrally-closed-dim1-dedekind}
Any Noetherian integrally closed ring
$R$ of Krull dimension $1$ is
Dedekind.
\end{theorem}

\begin{proof}
Step 1. Let $I\subset R$ be an ideal.
If $II^{-1}\subset R$, then there
exists a prime ideal $\mathfrak{p}$
containing $II^{-1}$.
The localization of $I$ at
$\mathfrak{p}$ is a principal ideal,
generated by some $x\in R$, because
$R_{\mathfrak{p}}$ is a DVR.
Let $b_1,\dots,b_n$ be generators of
$I$.
Then $b_i=xr_i/s_i$, where
$s_i\notin\mathfrak{p}$.

Step 2. Let $s=\prod s_i$.
Then $sx^{-1}\in I^{-1}$, which gives
$s\in II^{-1}$.
This is a contradiction, because each
$s_i\notin\mathfrak{p}$.
Hence $II^{-1}=R$, so $I$ is
invertible.
\end{proof}

\begin{remark}[AI-generated]
\label{commutative-algebra-remark-dedekind-iff}
If $R$ is a Dedekind ring which is not a
field, then $R$ has Krull dimension $1$.
A field is also Dedekind, but has Krull
dimension $0$.
\end{remark}

\medskip\noindent
In fact,
Kummer and Dedekind only cared about
the following examples:

\begin{definition}
\label{commutative-algebra-definition-number-field}
Let $[K:\mathbb{Q}]$ be finite.
Then $K$ is called an algebraic number
field.
\end{definition}

\begin{example}
\label{commutative-algebra-example-number-field-canonical}
The ring of integers in an algebraic
number field is Dedekind.
$\mathbb{Z}[\sqrt{-1}],
\mathbb{Z}\left[\frac{
\sqrt{-3}+1}{2}\right],
\mathbb{Z}[\sqrt{-2}]$.
The algebraic closure of $\mathbb{Z}$
in a number field $K$,
denoted $\mathcal{O}_K$,
is Dedekind.
\end{example}

\begin{theorem}
\label{commutative-algebra-theorem-dedekind-equivalent-conditions}
Let $R$ be a ring without zero divisors.
Then the following are equivalent:
\begin{itemize}
\item[(i)] $R$ is Dedekind.
\item[(ii)] $R$ is Noetherian and any
localization of $R$ at a prime ideal is
a DVR.
\item[(iii)] Every non-zero ideal of $R$
is equal to a unique product of maximal
ideals.
\end{itemize}
\end{theorem}

\begin{proof}
Step 1: $(i)$ implies $(ii)$ by
Proposition~\ref{commutative-algebra-proposition-dedekind-local-dvr},
and $(ii)$ implies $(i)$ because
$R=\bigcap_{\mathfrak{p}}R_{\mathfrak{p}}$,
hence $(ii)$ implies that $R$ is
integrally closed.

Step 2: To deduce $(iii)$ from $(i)$, take
any ideal $I\subset R$. If $I$ is maximal,
we are done. If $I$ is not maximal, let
$\mathfrak{m}$ be a maximal ideal strictly
containing $I$. Since
$\mathfrak{m}^{-1}\mathfrak{m}=R$ and
$I\supset \mathfrak{m}$, the product
$J:=\mathfrak{m}^{-1}I\subset R$ is an
ideal. Then $I=\mathfrak{m}J$, and
$J\supset I$. Repeating this procedure,
we obtain an increasing chain of ideals
$$
I\subset \mathfrak{m}_1^{-1}I
\subset \mathfrak{m}_1^{-1}\mathfrak{m}_2^{-1}I
\subset \cdots
$$
which terminates because $R$ is
Noetherian. Then
$\mathfrak{m}_1^{-1}\mathfrak{m}_2^{-1}
\cdots \mathfrak{m}_n^{-1}I=\mathfrak{m}$
is maximal, giving
$I=\mathfrak{m}_1\mathfrak{m}_2\cdots
\mathfrak{m}_n\mathfrak{m}$.

Step 3: Let $R_{\mathfrak{p}}$ be a
localization of $R$ in $\mathfrak{p}$.
Then every non-zero ideal of
$R_{\mathfrak{p}}$ is a product of
non-zero maximal ideals, hence the
maximal ideal of $R_{\mathfrak{p}}$ is
principal, and $R_{\mathfrak{p}}$ is DVR.
We proved that $(iii)$ implies $(ii)$.
\end{proof}

\begin{lemma}[Corollary]
\label{commutative-algebra-lemma-dedekind-factorial-iff-principal}
A Dedekind ring $R$ is factorial if and
only if all ideals in $R$ are principal.
\end{lemma}

\begin{proof}
Step 1. A principal ideal ring is
factorial if it is also Noetherian,
because in a Noetherian ring every
element can be represented as a product
of irreducible ones.

Step 2. Conversely, let $R$ be a
factorial Dedekind ring,
$\mathfrak{p}\subset R$ a prime ideal,
and $a\in \mathfrak{p}$ a non-zero
element.
Then $\mathfrak{p}$ contains an
irreducible factor $t$ of $a$.
Since $R$ is factorial, the ideal $(t)$
is prime.
Since $R$ has Krull dimension $1$, the
chain
$0\subset \mathfrak{p}\subset (t)$
has length $2$, unless
$\mathfrak{p}=(t)$.
Since every ideal is a product of
maximal ideals, every ideal in $R$ is
principal.
\end{proof}

\medskip\noindent
Now we make a slight detour to prove the
Chinese Reminder Theorem.

\begin{definition}
\label{commutative-algebra-definition-comaximal-ideals}
Ideals $P_1,\dots,P_n\subset R$ are
called {\it pairwise coprime}, or
{\it comaximal},
if $P_i+P_j=R$ whenever $i\neq j$.
\end{definition}

\begin{lemma}
\label{commutative-algebra-lemma-coprime-ideals-intersection}
If $P,Q\subset R$ are coprime ideals,
then $PQ=P\cap Q$.
\end{lemma}

\begin{proof}
If $x\in P\cap Q$ and $P+Q=R$, then
$1=p+q$ with $p\in P$ and $q\in Q$.
Hence
$$
x=x(p+q)\in PQ+QP\subset PQ,
$$
which shows that $P\cap Q\subset PQ$.
The reverse inclusion is clear.
\end{proof}

\medskip\noindent
Now we make a slight detour to prove the
Chinese Reminder Theorem.

\begin{definition}
\label{commutative-algebra-definition-comaximal-ideals-duplicate-2}
Ideals $P_1,\dots,P_n\subset R$ are
called {\it pairwise coprime}, or
{\it comaximal},
if $P_i+P_j=R$ whenever $i\neq j$.
\end{definition}

\begin{lemma}
\label{commutative-algebra-lemma-coprime-ideals-intersection-duplicate-2}
If $P,Q\subset R$ are coprime ideals,
then $PQ=P\cap Q$.
\end{lemma}

\begin{proof}
If $x\in P\cap Q$ and $P+Q=R$, then
$1=p+q$ with $p\in P$ and $q\in Q$.
Hence
$$
x=x(p+q)\in PQ+QP\subset PQ,
$$
which shows that $P\cap Q\subset PQ$.
The reverse inclusion is clear.
\end{proof}

\medskip\noindent
The following theorem is the induction
base for the proof of the Chinese Reminder
Theorem.

\begin{theorem}[Chinese reminder theorem]
\label{commutative-algebra-theorem-chinese-remainder}
Let $P_1,\dots,P_n$ be comaximal ideals
in a ring $R$. Then
$$
\prod_i P_i=
\bigcap_i P_i=P_1P_2\cdots P_n.
$$
Moreover, the natural map
$$
\chi:R\Big/\bigcap_i P_i\to
\bigoplus_i R/P_i
$$
is an isomorphism.
\end{theorem}

\begin{proof}
We already proved the base case of
induction for $n=2$.
By the induction assumption,
$$
\bigcap_{i=1}^{n-1}P_i=
P_1P_2\cdots P_{n-1}
$$
and
$$
R\Big/\bigcap_{i=1}^{n-1}P_i=
\bigoplus_{i=1}^{n-1}R/P_i.
$$
We start by showing that the ideals
\[
\bigcap_{i=1}^{n-1}P_i
\quad\text{and}\quad P_n
\]
are comaximal. Indeed, one has
$$
1=\prod_{i=1}^{n-1}(P_i+P_n)
\in \prod_{i=1}^{n-1}P_i+P_n
\subset \bigcap_{i=1}^{n-1}P_i+P_n.
$$
Therefore,
\[
\bigcap_{i=1}^{n-1}P_i
\text{ and }P_n
\]
are comaximal, and
by the base case $n=2$ we get
$$
P_n\bigcap_{i=1}^{n-1}P_i=
P_n\cap\bigcap_{i=1}^{n-1}P_i=
\bigcap_{i=1}^nP_i.
$$
Also,
$$
R\Big/\Bigl(P_n\cap\bigcap_{i=1}^{n-1}P_i\Bigr)
=R/P_n\oplus
R\Big/\bigcap_{i=1}^{n-1}P_i
=\bigoplus_i R/P_i.
$$
\end{proof}

\medskip\noindent
Back to Dedekind rings,
we have the following result.
According to AI:
Geometrically, this says that on a
Dedekind scheme every ideal sheaf is
generated by two global sections.
Equivalently, in dimension $1$ every
closed subscheme is cut out by two
functions.
This does not mean that the scheme is
two-dimensional.
The point is that the statement is
global: an ideal sheaf on a curve is
locally principal, but a single global
equation may fail to exist because of
twisting.
The theorem says that two global
functions always suffice to generate
the ideal. In other words, the ideal
sheaf of any set of points with
multiplicities is generated by two
global functions.

\begin{theorem}
\label{commutative-algebra-theorem-dedekind-ideals-two-generators}
Let $R$ be a Dedekind ring, and
$I\subset R$ an ideal.
Then for any ideal $J\supset I$, there
exists $b\in R$ such that $J=I+(b)$.
Moreover, $I$ can be generated by two
elements of $R$.
\end{theorem}

\begin{proof}
Step 1. Let $\mathfrak{p}_i$ be non-zero
prime ideals, and
$I=\prod_i \mathfrak{p}_i^{\alpha_i}$.
Since $R$ has Krull dimension $1$, the
ideals $\mathfrak{p}_1,\mathfrak{p}_2,\dots$
are pairwise comaximal.
Since $\sqrt{\mathfrak{p}_i^n}=\mathfrak{p}_i$,
one can never have
$\mathfrak{p}_i^n\subset \mathfrak{p}_j$.
This implies that the ideals
$\mathfrak{p}_i^{\alpha_n}$ are also
pairwise comaximal.

Step 2. Chinese Reminder Theorem implies
that
$$
R/I=\bigoplus_i R/\mathfrak{p}_i^{\alpha_i}.
$$
Each of the factors
$R/\mathfrak{p}_i^{\alpha_i}$ is a quotient
of a DVR, hence any ideal in
$R/\mathfrak{p}_i^{\alpha_i}$ and in $R/I$
is principal.
This implies that $J=I+(b)$.

Step 3. If $a\in I$ is non-zero, we apply
Step 2 to the ideals $(a)\subset I$, and
find that $I=(a)+(b)$.
\end{proof}

\medskip\noindent
This is exactly the point where Dedekind
and Kummer were heading: elements may not
factor uniquely, but ideals do.
So the right replacement for unique
factorization in arithmetic rings is
unique factorization of prime ideals.
The idea of the proof is to compare two
prime decompositions of the same ideal
and strip off any common prime factor.
This reduces the number of factors, and
eventually forces a contradiction unless
the two decompositions agree up to order.
The Chinese Reminder Theorem enters when
we use pairwise comaximality to pass from
the equality of the two decompositions to
an equality of products of the
corresponding ideals.

\begin{theorem}
\label{commutative-algebra-theorem-dedekind-prime-factorization-unique}
Let $R$ be a Dedekind ring, $I\subset R$
an ideal, $\mathfrak{p}_i$ non-zero prime
ideals, and
$I=\prod_i \mathfrak{p}_i^{\alpha_i}$
its prime decomposition.
Then this decomposition is unique, up to
renumbering the prime factors.
\end{theorem}

\begin{proof}
Let
$I=\prod_j \mathfrak{q}_j^{\beta_j}$
be another such decomposition.
If there is a common factor, say,
$\mathfrak{p}$, between
$\mathfrak{p}_i$ and $\mathfrak{q}_j$,
we can replace $I$ by
$I\mathfrak{p}^{-1}$, and have two
non-equivalent decompositions with fewer
factors.
Repeating this procedure, we obtain
$$
I=\prod_i \mathfrak{p}_i^{\alpha_i}
=\prod_j \mathfrak{q}_j^{\beta_j}.
$$
However, the ideals
$\mathfrak{q}_i^{\alpha_i}$ and
$\mathfrak{q}_j^{\beta_j}$ are pairwise
comaximal, which implies by the proof of
the Chinese Reminder Theorem that
$$
I=\prod_i \mathfrak{p}_i^{\alpha_i}
\cap \prod_j \mathfrak{q}_j^{\beta_j}
=\prod_i \mathfrak{p}_i^{\alpha_i}
\prod_j \mathfrak{q}_j^{\beta_j}
=I^2,
$$
which is impossible by
Lemma~\ref{commutative-algebra-lemma-nakayama-first},
also called Krull lemma.
\end{proof}

\medskip\noindent
This extends
Theorem~
\ref{commutative-algebra-theorem-dedekind-equivalent-conditions}:

\begin{theorem}
\label{commutative-algebra-theorem-dedekind-equivalent-conditions-full}
Let $R$ be a ring without zero divisors.
Then the following are equivalent:
\begin{itemize}
\item[(i)] $R$ is Dedekind.
\item[(ii)] $R$ is Noetherian and any
localization of $R$ at a prime ideal is
a DVR.
\item[(iii)] Every non-zero ideal in $R$
is equal to a unique product of prime
ideals.
\item[(iv)] $R$ is integrally closed,
Noetherian ring of Krull dimension $1$.
\item[(v)] Every ideal in $R$ is
projective as an $R$-module.
\end{itemize}
\end{theorem}

\begin{proof}
Step 1. The equivalence $(v)\Rightarrow(i)$
is clear because finitely generated
projective fractional ideals are the same
as invertible ones.
The implication $(i)\Rightarrow(ii)$ is
also clear because invertibility of the
ideal implies that it is finitely
generated projective, and after passing
to localizations, every finitely generated
projective module is free, hence the
localization is a principal ideal ring,
therefore a DVR.

Step 2. The implications $(ii)\Rightarrow(i)
\Rightarrow(iii)$ and the equivalence
$(iv)\Leftrightarrow(ii)$ were proven
today.

Step 3. It remains to show that
$(iii)\Rightarrow(i)$.
Theorem~\ref{commutative-algebra-theorem-dedekind-prime-factorization-unique}
implies that any ring where each non-zero
ideal is a product of maximal ideals is
Dedekind.
From the uniqueness of decomposition we
obtain that $\mathfrak{p}\subset \mathfrak{q}$
implies $\mathfrak{p}=\mathfrak{q}$ when
both ideals are prime; indeed,
$\mathfrak{p}\mathfrak{q}=\mathfrak{p}$.
(Oops! There is a mistake here!
It's probably $\mathfrak{p}^2$…
or smaller than $\mathfrak{p}^2$.)
Hence $R$ has Krull dimension $1$, and we
can apply Theorem~
\ref{commutative-algebra-theorem-noetherian-integrally-closed-dim1-dedekind}.
\end{proof}

\section{Hilbert polynomial}
\label{commutative-algebra-section-hilbert-polynomial}

\noindent
Let $A$ be a local ring.
Think of it as the ring of
germs of holomorphic functions at a point.
There is an associated graded:
$$
\bigoplus \frac{\mathfrak{p}^k}{
\mathfrak{p}^{k+1}}
$$
In this section we will show that the map
$$
k \mapsto \dim \frac{\mathfrak{p}^k}{
\mathfrak{p}^{k+1}}
$$
is polynomial.

\begin{definition}
\label{commutative-algebra-definition-filtered-ring}
Let $A$ be a ring.
A {\it multiplicative filtration} on $A$ is
a sequence
$$
A=F_0\supset F_1\supset F_2\supset \cdots
$$
such that each $F_i$ is closed under
multiplication and
$$
F_iF_j\subset F_{i+j}.
$$
A ring equipped with a multiplicative
filtration is called {\it filtered}.
\end{definition}

\begin{lemma}
\label{commutative-algebra-lemma-associated-graded-ring}
Let $A=F_0\supset F_1\supset\cdots$
be a filtered ring, and let
$a_1,a_2\in F_i$ and $b_1,b_2\in F_j$.
Assume that
$$
a_1=a_2 \mod F_{i+1}
\qquad\text{and}\qquad
b_1=b_2 \mod F_{j+1}.
$$
Then
$$
a_1b_1=a_2b_2 \mod
F_{i+1}F_j+F_iF_{j+1}.
$$
\end{lemma}

\begin{proof}
We have
$$
\begin{aligned}
a_1b_1-a_2b_2
&=a_1(b_1-b_2)+(a_1-a_2)b_2.
\end{aligned}
$$
The first term lies in $F_iF_{j+1}$ and the
second lies in $F_{i+1}F_j$.
\end{proof}

\begin{remark}
\label{commutative-algebra-remark-associated-graded-ring}
Since
$$
F_{i+1}F_j+F_iF_{j+1}\subset F_{i+j+1},
$$
the product of classes in
$F_i/F_{i+1}$ and $F_j/F_{j+1}$ is well
defined in $F_{i+j}/F_{i+j+1}$.
Therefore the associated graded object is a
graded ring.
\end{remark}

\begin{lemma}
\label{commutative-algebra-lemma-graded-vs-cstar}
The category of graded vector spaces over
$\mathbb{C}$ is equivalent to the category of
$\mathbb{C}^*$ representations.
\end{lemma}

\begin{proof}
Let $V=\bigoplus_{n\in \mathbb{Z}}V_n$ be a
graded $\mathbb{C}$-vector space.
Define an action of $\mathbb{C}^*$ on $V$
by letting $t\in \mathbb{C}^*$ act as
$$
t\cdot v=t^n v
$$
for $v\in V_n$.
This preserves direct sums and graded maps,
so we get a functor.

Conversely, let $V$ be a
$\mathbb{C}^*$-representation.
Since $\mathbb{C}^*$ is diagonalizable, $V$
decomposes as a direct sum of its weight
spaces
$$
V=\bigoplus_{n\in \mathbb{Z}}V_n,
$$
where $t\in \mathbb{C}^*$ acts on $V_n$ by
the character $t\mapsto t^n$.
This gives a graded vector space, and the two
constructions are inverse to one another.
\end{proof}

\begin{definition}
\label{commutative-algebra-definition-ideal-filtration}
Let $I\subset A$ be an ideal.
Then
$$
A\supset I\supset I^2\supset \cdots
$$
is a multiplicative filtration.
The corresponding associated graded ring is
$$
\bigoplus_{k=0}^{\infty} \frac{I^k}{I^{k+1}}.
$$
\end{definition}

\noindent
Recall from our work in Noether rings
that the associated graded ring
of a Noetherian local ring is finitely
generated.

\begin{definition}
\label{commutative-algebra-definition-graded-module}
Let
$$
A^\ast=\bigoplus_{i=0}^{\infty} A_i
$$
be a graded ring.
A {\it graded module} over $A^\ast$ is a
module
$$
M^\ast=\bigoplus_{i=0}^{\infty} M_i
$$
such that
$$
A_iM_j\subset M_{i+j}.
$$
\end{definition}
\noindent
In the following definition
it is crucial that we use the definition
of graded module starting from $i=0$.
Indeed, we have problems of dimension
with rings like
$\mathbb{C}[z,z^{-1}] \otimes
\mathbb{C}[t,t^{-1}]
=\mathbb{C}[z^i t^j]$,
which has infinite-dimensional components.

\begin{lemma}
\label{commutative-algebra-lemma-graded-module-finite-dim}
Let $M^\ast$ be a finitely generated graded
module over a graded ring $A^\ast$, with
each $A_i$ finite-dimensional over a field
$A_0=k$.
Then
$$
\dim_k M_i<\infty
$$
for all $i$.
\end{lemma}

\begin{proof}
Choose homogeneous generators of $M^\ast$.
Each graded piece $M_i$ is then spanned by
finitely many products of generators with
elements of the finite-dimensional spaces
$A_j$ whose degrees add up to $i$.
Thus each $M_i$ is finite-dimensional over
$k$.
\end{proof}

\begin{exercise}[Difference derivative]
\label{commutative-algebra-exercise-hilbert-function-polynomial}
Let $f:\mathbb{Z}_{\geq 0}\to \mathbb{Z}$
be a function.
Assume that
$$
D(n):=f(n+1)-f(n)
$$
is polynomial for $n\gg 0$.
Prove that $f(n)$ is polynomial for
$n\gg 0$.
\end{exercise}

\begin{proof}
The idea is that this thing acts like a
derivative. If $D(n)$ is a constant,
then $f$ grows arithmetically. If
$D(n)$ is of degree 1, then $f$ grows
quadratically, and so on.
To show this suppose $f(t)=t^k$.
Then compute to find
$$
f(t) - f(t-1) = \sum \binom{1}{n}t^n
$$
and apply the case where $D(n)$ is constant.
\end{proof}

\begin{definition}
\label{commutative-algebra-definition-hilbert-function}
Let $M^\ast$ be a finitely generated graded
module over a graded ring $A^\ast$, with
$A_0=k$ a field.
The {\it Hilbert function} of $M^\ast$ is
$$
h_M(n)=\dim_k M_n.
$$
\end{definition}

\begin{theorem}
\label{commutative-algebra-theorem-hilbert-function-polynomial}
Let $M^\ast$ be a finitely generated graded
module over a finitely generated graded ring
$A^\ast$, with $A_0=k$ a field.
Then the Hilbert function $h_M(n)$ is
polynomial for $n\gg 0$.
\end{theorem}

\begin{proof}
Step 1. Let
$$
0\to M_1^\ast\to M_2^\ast\to M_3^\ast\to 0
$$
be an exact sequence of finitely generated
graded $A^\ast$-modules.
Then
$$
h_{M_2}(n)=h_{M_1}(n)+h_{M_3}(n).
$$
Therefore the Hilbert function is polynomial
for $M_2^\ast$ if it is polynomial for
$M_1^\ast$ and $M_3^\ast$.

Step 2. Let $A^\ast$ be a finitely
generated graded ring, $M^\ast$ a torsion-
free $A^\ast$-module, and $a\in A_k$.
Denote by
$$
L_a:M^\ast\to M^\ast
$$
the multiplication map $m\mapsto ma$.
Consider the exact sequence
$$
0\to \ker L_a\to M^\ast
\xrightarrow{L_a} M^\ast
\to M^\ast+aM^\ast\to 0.
$$
Let
$$
U:=M^\ast+aM^\ast.
$$
Suppose that the Hilbert function $h_U(n)$
and $h_{\ker L_a}(n)$ are polynomial for
$n\gg 0$.
Since
$$
h_M(n+1)-h_M(n)=h_U(n)-h_{\ker L_a}(n),
$$
this implies that $h_M(n)$ is polynomial for
$n\gg 0$.
Exercise
\ref{commutative-algebra-exercise-hilbert-function-polynomial}.

Step 3. For any graded, finitely generated
$k[t]$-module, the Hilbert function is
constant for $n\gg 0$.
This proves the theorem for $A^\ast=k[t]$.

Step 4. Let $A^\ast$ be a graded ring
generated by $d$ elements. We prove the
Hilbert function $h_M(n)$ using induction in
$d$. For $d=1$ it follows from Step 3.
Using induction in $d$, we assume that for
any graded ring $B^\ast$ generated by
$d$ elements, and any finitely-generated
$B^\ast$-module $U^\ast$, the Hilbert
function $h_U(n)$ is polynomial for
$n\gg 0$.

Step 5. Let $a\in A_k$ be a generator of
$A^\ast$, and
$$
0\to \ker L_a\to M^\ast
\xrightarrow{L_a} M^\ast
\to M^\ast/aM^\ast\to 0
$$
the exact sequence of Step 2.
By the induction assumption, the Hilbert
function of $A^\ast/aA^\ast$-modules is
polynomial for $n\gg 0$. Since $\ker L_a$
is a module over $A^\ast/aA^\ast$, the
Hilbert function $h_{\ker L_a}(n)$ is
polynomial. By Step 2, the Hilbert function
$h_M(n)$ is also polynomial.
\end{proof}

\begin{lemma}[Corollary]
\label{commutative-algebra-lemma-hilbert-degree-drop}
Let $M^\ast$ be a graded $A^\ast$-module,
and let $a\in A_k$ be such that the
multiplication map
$L_a:M^\ast\to M^\ast$ is injective.
Denote by $N$ the module
$$
N=M^\ast/aM^\ast.
$$
Then
$$
\deg h_N(n)=\deg h_M(n)-1.
$$
\end{lemma}

\medskip\noindent
The following result is the main result of
the section: that is, the key fact
is that the degree of the Hilbert
function is the Krull dimension of the
underlying ring.

Actually, this shows that the Hilbert
scheme of points parametrizes 0-dimensional
subschemes: it is defined by restricting
the functor of flat maps to families
with constant (degree 0!) Hilbert polynomial!

\begin{theorem}
\label{commutative-algebra-theorem-hilbert-polynomial-slide-7}
Let $R$ be a Noetherian local ring without
zero divisors, $m$ its maximal ideal,
$$
A:=\bigoplus_{i=0}^{\infty} \frac{m^i}{
m^{i+1}}
$$
its associated graded ring, and $h_A(n):=
\dim A_n$.
Then the Hilbert function $h_A(n)$ is a
polynomial for $n\gg 0$.
Moreover, its degree is equal to the Krull
dimension of $R$.
\end{theorem}

\begin{proof}
Step 1. Let
  $$
  0\to M_1^\ast\to M_2^\ast\to M_3^\ast\to 0
  $$
  be an exact sequence of finitely generated
  graded $A^\ast$-modules. Then
  $$
  h_{M_2}(n)=h_{M_1}(n)+h_{M_3}(n).
  $$
  So the Hilbert function is polynomial for
  $M_2^\ast$ whenever it is polynomial for
  $M_1^\ast$ and $M_3^\ast$.

  Step 2. Let $A^\ast$ be a finitely
  generated
  graded ring, $M^\ast$ a torsion-free
  $A^\ast$-module, and $a\in A_k$. Consider
  multiplication by $a$:
  $$
  L_a:M^\ast\to M^\ast.
  $$
  There is an exact sequence
  $$
  0\to \ker L_a\to M^\ast
  \xrightarrow{L_a} M^\ast
  \to M^\ast+aM^\ast\to 0.
  $$
  Set
  $$
  U:=M^\ast+aM^\ast.
  $$
  If $h_U(n)$ and $h_{\ker L_a}(n)$ are
  polynomial for $n\gg 0$, then
  $$
  h_M(n+1)-h_M(n)=h_U(n)-h_{\ker L_a}(n),
  $$
  so $h_M(n)$ is polynomial for $n\gg 0$
.
Step 2: Let $\tilde{R}$ be an integral
closure of $R$. By Cohen-Seidenberg, the
Krull dimension of $\tilde{R}$ is equal to
the Krull dimension of $R$. Denote by
$\tilde{A}$ the associated graded ring of
$\tilde{R}$. Then deg $h_A(n)$ = deg
$h_{\tilde{A}}(n)$, because $\tilde{A}$ is a
finite extension of $A$. Replacing $R$ by
$\tilde{R}$ if necessary, we may assume that
$R$ is normal.

Step 3: Let
$$
0\subset p_1\subset p_2\subset \cdots
\subset p_n\subset R
$$
be a chain of prime ideals of maximal
possible length. Since $R_{p_1}$ is a local
integrally closed ring of Krull dimension
1, it is a discrete valuation ring, and its
maximal ideal is principal. Therefore, in
the ring $R_{p_1}$ we have $p_1=(p)$, where
$p\in p_1$.

Step 4: The localization of $R/(p)$ in
$p_1$ is a field, hence $p_1$ is a minimal
prime ideal in $R/(p)$. Therefore the chain
of prime ideals
$$
p_1\subset p_2\subset \cdots\subset p_n
\subset R/(p)
$$
is maximal, and $R/(p)$ has Krull
dimension $\dim R-1$.

Step 5: Let $a\in A_1=m^1/m^2$ be the class
represented by $p$. Using induction in
$\dim R$, we may assume that
$$
\deg h_{A/(a)}(n)=\dim R/(p)=\dim R-1.
$$
By Theorem 1, Step 2, $\deg h_{A/(a)}(n)=
\deg h_A(n)-1$. This gives
$$
\deg h_A(n)=\deg h_{A/(a)}(n)+1=\dim
R/(p)+1=\dim R.
$$
\end{proof}

\medskip\noindent
Now we include the exercises corresponding
with this lecture. They are intended
to show that rational functions
can be used instead of polynomials
to derive the same results.

\begin{exercise}
\label{commutative-algebra-exercise-poincare-10-1}
Let $n\mapsto H(n)$ be a polynomial of
degree $d$, and $k$ a fixed integer number.
Prove that
$n\mapsto H(n)-H(n-k)$ is a polynomial of
degree $d-1$.
\end{exercise}

\begin{exercise}
\label{commutative-algebra-exercise-poincare-10-2}
Let $p(t)=\sum a_i t^i$, and
$q(t)=\sum t^{k_i}\cdot \sum a_i t^i
=\sum b_i t^i$, for some fixed integer
$k>0$.
Assume that $a_i=h(i)$, where $h(i)$ for
$i\gg 0$ is a polynomial of degree $d$.
Prove that $h_1(i):=b_i$, is a polynomial of
degree $d+1$ when $i\gg 0$.
\end{exercise}

\begin{proof}
Hint. Use the previous exercise.
\end{proof}

\begin{exercise}
\label{commutative-algebra-exercise-poincare-10-3}
Let $f(t)$ be a polynomial, $k_1, \dots,
k_n\in \mathbb{Z}_{>0}$, and $p(t)$ a
rational function
$$
p(t)=f(t)\prod_{i=1}^n(1-t^{k_i})^{-1}.
$$
Consider the Taylor series
$$
p(t)=\sum a_i t^i.
$$
Prove that there is a polynomial $H(i)$
such that $a_i=H(i)$ for $i\gg 0$.
\end{exercise}

\begin{proof}
Hint. Use the previous exercise
\end{proof}

\begin{definition}
\label{commutative-algebra-definition-graded-modules-additive}
Let $C$ be the category of finite generated
graded $R$-modules, where $R$ is a graded
ring.
A function $P:C\to \mathbb{Z}[[t]]$ is
called {\it additive} if for any exact
sequence
$$
0\to M_1\to M_2\to M_3\to 0,
$$
we have
$$
P(M_1)+P(M_2)=P(M_3).
$$
\end{definition}

\begin{exercise}
\label{commutative-algebra-exercise-poincare-10-4}
Let $P$ be an additive function on the
category $C$ of finite generated graded
modules over polynomials
$R=\mathbb{C}[z_1,\dots,z_n]$.
Assume that $P(M)$ is polynomial for any
$R$-module of finite dimension over
$\mathbb{C}$, and $P(M_d)=t^dP(M)$, where
$M_d$ is obtained from $M$ by shifting the
grading by $+d$.
Prove that $P(M)$ is a rational function of
form
$$
g(t)+f(t)\prod_{i=1}^n(1-t^{k_i})^{-1}
$$
for some polynomials $g(t),f(t)$.
\end{exercise}

\begin{proof}
Hint. Use the previous exercise.
\end{proof}

\begin{definition}
\label{commutative-algebra-definition-poincare-series}
The {\it Poincar\'e series} of a graded
module $M^\ast$ over polynomials is the
function
$$
P(t):=\sum_{i=0}^{\infty} h_i t^i,
$$
where
$$
h_i:=\dim M_i.
$$
\end{definition}

\begin{exercise}
\label{commutative-algebra-exercise-poincare-10-5}
Prove that the Poincar\'e series of
$\mathbb{C}[x,y]$ is
$$
\frac{1}{(1-t)^2}.
$$
\end{exercise}

\begin{exercise}
\label{commutative-algebra-exercise-poincare-10-6}
Prove that the Poincar\'e series of
$\mathbb{C}[x,y]/(x^2)$ is
$$
\frac{t+1}{1-t}.
$$
\end{exercise}

\begin{exercise}
\label{commutative-algebra-exercise-poincare-10-7}
Prove that the Poincar\'e series of the
ideal $(x^2,y^3)\subset \mathbb{C}[x,y]$ is
$$
\frac{t^2+t^3-t^5}{(1-t)^2}.
$$
\end{exercise}

\begin{exercise}
\label{commutative-algebra-exercise-poincare-10-8}
(*). Let
$$
P(t):=\sum_{i=0}^{\infty} h_i t^i,
$$
where $i\mapsto h_i$ is a polynomial
function.
Prove that $P(t)$ is rational, or find a
counterexample.
\end{exercise}

\section{Regular rings}
\label{commutative-algebra-section-regular-rings}

\noindent
Let $p$ be a prime integer.
The associated graded to $(p) \subset \mathbb{Z}$
is
$$
\bigoplus \frac{(p^n)}{(p^{n+1)})}
=
\mathbb{Z}/p[t].
$$


\begin{theorem}
\label{commutative-algebra-theorem-regular-rings}
Let $R$ be a Noetherian local ring of
Krull dimension $d$, $m$ its maximal ideal,
and $k:=R/m$ its residue field.
Then
$$
\dim_k(m/m^2)\geq d.
$$
Moreover, the following are equivalent:
\begin{enumerate}
\item The associated graded ring of $R$ is
isomorphic to $k[t_1,\dots,t_d]$.
\item $\dim_k(m/m^2)=d$.
\item The ideal $m$ can be generated by
$d$ elements.
\end{enumerate}
\end{theorem}

\begin{proof}
Step 1. The implication $(i)\Rightarrow(ii)$
is obvious, and $(ii)\Rightarrow(iii)$
follows from Nakayama lemma, which also
implies that the associated graded ring
$A^\ast$ of $R$ is generated by $A_1$.
Since $A^\ast$ is a quotient of
$k[A_1]$, and
$\dim_k k[A_1]=\dim_k A_1$, we obtain
$$
\dim_k(m/m^2)
=\dim_k A_1
\geq \dim_k A^\ast
=\dim R
=d.
$$
It remains to prove
$(iii)\Rightarrow(i)$.

Step 2. Let
$\mu:k[A_1]\to A^\ast$ be a natural
surjective map, and $I:=\ker \mu$.
Let $t_1,\dots,t_d$ be the basis in $A_1$;
then $k[A_1]=k[t_1,\dots,t_d]$.
If $I$ is non-empty, there is an algebraic
relation between $t_i\in A_1$, which implies
$$
d=\dim \operatorname{tr} A^\ast<
\dim \operatorname{tr} k[A_1]
=\dim kA_1=d.
$$
Therefore, $I=0$ and $\mu$ is an
isomorphism.
\end{proof}

\begin{definition}
\label{commutative-algebra-definition-regular-ring}
{\it A Noetherian local ring} is called
regular if $(i)$--$(iii)$ hold.
\end{definition}

\begin{definition}
\label{commutative-algebra-definition-smooth-affine-variety}
An affine variety $X=\operatorname{Spec}(A)$
is called smooth if all local rings of
maximal ideals in $A$ are regular.
\end{definition}

\begin{lemma}
\label{commutative-algebra-lemma-dvr-regular}
Discrete valuation rings are regular.
\end{lemma}

\begin{proof}
Indeed, the maximal ideal of the discrete
valuation ring $R$ is principal, hence
$m/m^2=R/m=k$.
\end{proof}

\begin{lemma}[Corollary]
\label{commutative-algebra-lemma-dim-1-dedekind-smooth}
Let $R$ be a ring of Krull dimension $1$
and without zero divisors. Then $R$ is
Dedekind if and only if it is smooth.
\end{lemma}

\noindent
The idea of the next lemma is that
``denominators stabilize''.

\begin{lemma}
\label{commutative-algebra-lemma-integral-over-fraction-field}
Let $R$ be a Noetherian ring without zero
divisors, and $x\in k(R)$ an element in its
fraction field. Then $x$ is integral over
$R$ if and only if there is $y\in R$ such
that $yx^n\in R$ for all $n>0$.
\end{lemma}

\begin{proof}
Let $S=R[x]\subset k(X)$ be the subring
generated by $x$. If $x=a/b$, with
$a,b\in R$, we have $b^nS\subset R$, for
some $n\in\mathbb{Z}_{>0}$. Then
$b^nx^m\in R$ for all $m\in\mathbb{Z}_{\geq 0}$.
Conversely, if $yx^n\in R$, then
$aS\subset R$; the $R$-module $aS$ is
finitely generated, because $R$ is
Noetherian.
\end{proof}



\begin{proposition}
\label{commutative-algebra-proposition-regular-rings-integrally-closed}
Let $R$ be a Noetherian local ring, $m$
its maximal ideal, and $\operatorname{Gr}(R)$
the associated graded ring. Suppose that
$\operatorname{Gr}(R)$ contains no zero
divisors. Then $R$ has no zero divisors.
Moreover, if $\operatorname{Gr}(R)$ is
integrally closed, then $R$ is also
integrally closed.
\end{proposition}

\begin{proof}
Step 1. Further on, for any
$a\in m^k\setminus m^{k+1}$, we denote the
class of $a$ modulo $m^{k+1}$ by $a\in
\operatorname{Gr}(R)$. Krull's theorem
implies that $m^k\setminus m^{k+1}=R$,
hence $a\in\operatorname{Gr}(R)$ is
non-zero for any non-zero $a\in R$. If $ab$
vanishes in $R$, then
$a\in m^k\setminus m^{k+1}$ and
$b\in m^\ell\setminus m^{\ell+1}$ vanishes
in $\operatorname{Gr}(R)$, hence
$\operatorname{Gr}(R)$ also has zero
divisors. This proves the first statement.

Step 2. Let $x=a/b$ be integral over $R$,
with $a,b\in R$. Let $R_1:=R/(b)$.
By Nakayama lemma,
$$
\bigcap_k m^kR_1=0,
$$
hence
$$
\bigcap_k(m^k+(b))=(b).
$$
To show that $a\in(b)$ it would suffice to
prove that $a\in m^k+(b)$ for all $k$.

Step 3. We prove $a\in m^k+(b)$ using
induction in $k$. Assume that
$a\in m^{k-1}+(b)$, and $a=d+cb$, where
$d\in m^{k-1}\setminus m^k$.
For some $y\in R$, and for all
$n\in\mathbb{Z}_{>0}$ we have
$yd^n/b^n\in R$ (Lemma 1).
This brings $yd^n\in b^nR$ and
$yd^n\in b^n\operatorname{Gr}(R)$.
Since $\operatorname{Gr}(R)$ is integrally
closed, this implies that $d/b\in
\operatorname{Gr}(R)$, hence
$d=bz\mod m^k$, and $a\in m^k+(b)$.
\end{proof}

\begin{lemma}[Corollary]
\label{commutative-algebra-lemma-regular-rings-integrally-closed}
Regular rings are integrally closed.
\end{lemma}

\begin{definition}
\label{commutative-algebra-definition-derivation}
Let $R$ be a ring over a field $k$, and
$V$ an $R$-module. A $k$-linear map
$D:R\to V$ is called a derivation if it
satisfies the Leibnitz identity
$$
D(ab)=aD(b)+bD(a).
$$
The space of derivations from $R$ to $V$
is denoted $\operatorname{Der}_k(R,V)$.
\end{definition}

\begin{remark}
\label{commutative-algebra-remark-derivations-module}
$\operatorname{Der}_k(R,V)$ is an
$R$-module, with a natural $R$-action.
\end{remark}

\section{Gorenstein properties}
\label{commutative-algebra-section-gorenstein-properties}

The goal is to understand what the adjective
``arithmetically Gorenstein'' says about a
projective variety. The route is through
regular sequences, Cohen--Macaulay rings,
and duality. The general structural results
below are stated without proof; elementary
examples and consequences are worked out.
All rings in this section are commutative
and Noetherian.

\medskip\noindent
{\bf Regular sequences and depth.}
Recall the regular local rings of Section
\ref{commutative-algebra-section-regular-rings}.
A regular sequence is a different notion:
it describes equations that remain
non-zero-divisors when imposed successively.

\begin{definition}
\label{commutative-algebra-definition-regular-sequence}
A sequence $f_1,\ldots,f_r$ in a ring $R$
is a {\it regular sequence} if
$R/(f_1,\ldots,f_r)\ne0$ and multiplication
by $f_i$ is injective on
$R/(f_1,\ldots,f_{i-1})$ for each $i$.
\end{definition}

For example, $x,y$ is regular in $k[x,y]$.
The sequence $x,xy$ is not: after setting
$x=0$, the second element is zero.
Thus counting equations alone does not
recognize a complete intersection.

\begin{definition}
\label{commutative-algebra-definition-homogeneous-ci}
Let $S=k[x_0,\ldots,x_N]$ and let $I$ be
a proper homogeneous ideal.
The ideal $I$, and the algebra $S/I$, are
called a {\it complete intersection} if
$I$ is generated by a homogeneous regular
sequence. Equivalently, $I$ can be
generated by exactly $c$ elements, where
$$
c=\operatorname{ht}(I)=N+1-\dim(S/I).
$$
\end{definition}

Thus the number of equations must equal
the codimension, not merely the number
of generators in a chosen presentation.
Each equation in the regular sequence
cuts the dimension by exactly one.

For example, in $k[x,y,z]$ the ideal
$(x^2,y^2)$ is a complete intersection
of codimension two. Its quotient has
dimension one, although it is not reduced.
In contrast, the presentation $(x,xy)$
has two generators but defines $(x)$,
a complete intersection of codimension one.

\begin{definition}
\label{commutative-algebra-definition-depth-cm}
For a local ring $(R,\mathfrak m)$, its
{\it depth} is the maximum length of a
regular sequence contained in $\mathfrak m$.
It is {\it Cohen--Macaulay} if
$$
\operatorname{depth}R=\dim R.
$$
A ring is Cohen--Macaulay if its
localizations at all prime ideals are.
\end{definition}

Depth is at most Krull dimension.
Geometrically, the Cohen--Macaulay condition
allows a full set of parameters to cut down
dimension without creating zero divisors at
any intermediate step. Every Artinian local
ring is Cohen--Macaulay: both numbers are
zero. Thus Cohen--Macaulay does not mean
smooth, reduced, or irreducible.

\medskip\noindent
{\bf A concrete local Gorenstein criterion.}
A {\it system of parameters} of a local
ring of dimension $d$ is a sequence of $d$
elements of its maximal ideal whose quotient
is Artinian. In a Cohen--Macaulay local ring,
every such sequence is regular.

\begin{definition}
\label{commutative-algebra-definition-socle}
For an Artinian local ring
$(B,\mathfrak n,\kappa)$, its {\it socle} is
$$
\operatorname{Soc}(B)
=\{b\in B:\mathfrak n b=0\}.
$$
This is a vector space over the residue
field $\kappa=B/\mathfrak n$.
\end{definition}

\begin{definition}
\label{commutative-algebra-definition-gorenstein-local}
A local ring $R$ is {\it Gorenstein} if it
is Cohen--Macaulay and, after quotienting by
a system of parameters, the resulting
Artinian local ring has one-dimensional
socle over its residue field.
A ring is {\it Gorenstein} if all its
localizations at prime ideals are.
\end{definition}

This criterion is independent of the chosen
parameters. It is equivalent to the usual
homological definition that $R$ has finite
injective dimension as an $R$-module.
We use the socle criterion as our entry
point, without developing injective
resolutions here.

\begin{example}
\label{commutative-algebra-example-gorenstein-socles}
For $B=k[x,y]/(x^2,y^2)$ a basis is
$1,x,y,xy$. Its maximal ideal is $(x,y)$,
and its socle is $kxy$, so $B$ is
Gorenstein. In contrast, for
$C=k[x,y]/(x,y)^2$ a basis is $1,x,y$,
and the socle is $kx\oplus ky$.
Thus $C$ is Cohen--Macaulay but not
Gorenstein.
\end{example}

The next two results have different roles:
the first supplies a starting ring;
the second explains what happens when
one equation is imposed.
Together they show that complete
intersections are Gorenstein.

\begin{theorem}
\label{commutative-algebra-theorem-regular-local-gorenstein}
Source: Stacks Project, Lemma 47.21.3,
\href{https://stacks.math.columbia.edu/tag/0AWX}
{Tag 0AWX}.
\par\smallskip
Every regular local ring is Gorenstein.
\end{theorem}

\begin{theorem}
\label{commutative-algebra-theorem-gorenstein-hypersurface}
Source: Stacks Project, Lemma 47.21.7,
\href{https://stacks.math.columbia.edu/tag/0BJJ}
{Tag 0BJJ}.
\par\smallskip
Let $(R,\mathfrak m)$ be a Noetherian
local ring, and let $f\in\mathfrak m$
be a non-zero-divisor. Then
$R$ is Gorenstein if and only if
$R/(f)$ is Gorenstein.
\end{theorem}

No regularity assumption on $R$ is made
in this second theorem. Geometrically,
if $R$ is the local ring of $X$ at $p$,
then $R/(f)$ is the local ring of the
hypersurface $f=0$ at $p$.
The germ $f$ must be a non-zero-divisor,
not merely nonzero. The theorem concerns
Gorensteinness on both sides, not
regularity of the hypersurface.

\begin{lemma}[Corollary]
\label{commutative-algebra-theorem-gorenstein-ci}
Let $(S,\mathfrak n)$ be a regular local
ring and let $f_1,\ldots,f_c$ be a regular
sequence contained in $\mathfrak n$.
Then $S/(f_1,\ldots,f_c)$ is Gorenstein.
\end{lemma}

\begin{proof}
By Theorem
\ref{commutative-algebra-theorem-regular-local-gorenstein},
$S$ is Gorenstein.
By the definition of a regular sequence,
$f_i$ is a non-zero-divisor in
$S/(f_1,\ldots,f_{i-1})$.
Apply Theorem
\ref{commutative-algebra-theorem-gorenstein-hypersurface}
successively for $i=1,\ldots,c$.
Every resulting quotient is Gorenstein.
\end{proof}

Such quotients of regular local rings
are complete intersection local rings.
Regularity is the starting assumption;
Gorensteinness survives the cuts.
A regular sequence need not produce
a regular quotient.

The implications are: regular implies
complete intersection, which implies
Gorenstein, which implies Cohen--Macaulay.
They do not reverse in general.
For example, the local ring of the nodal
curve $xy=0$ at the origin is a hypersurface,
hence Gorenstein, but is not regular.

\medskip\noindent
{\bf Grading and the canonical module.}
Let $S=k[x_0,\ldots,x_N]$, with each
variable of degree one, and let $R=S/I$
for a proper homogeneous ideal $I$.
Put $d=\dim R$ and $c=N+1-d$.
Here $c$ is the codimension of the ring.
The irrelevant maximal ideal is
$\mathfrak m=R_{>0}$; in this setting the
Cohen--Macaulay and Gorenstein properties
can be tested at $\mathfrak m$.

Here is how the three notions fit together
in this graded setting. The ring $R$ is
Cohen--Macaulay precisely when it admits
a regular sequence $g_1,\ldots,g_d$ of
homogeneous positive-degree elements.
Its length equals $\dim R$, and
$$
B=R/(g_1,\ldots,g_d)
$$
is finite-dimensional over $k$.
The extra condition for $R$ to be
Gorenstein is
$$
\dim_k\operatorname{Soc}(B)=1,
\qquad
\operatorname{Soc}(B)
=\{b\in B:B_{>0}b=0\}.
$$
Thus, up to scalar, there is exactly one
nonzero element killed by every
positive-degree element of $B$.

There are two different regular sequences
here. For a complete intersection, the
$c$ generators of $I$ live in $S$ and
define $R$. The $d$ parameters live in
$R$ and reduce it to $B$.
Theorem \ref{commutative-algebra-theorem-gorenstein-ci}
explains the implication: polynomial
rings are Gorenstein, and quotienting by
a regular sequence preserves this property.
Hence every complete intersection is
Gorenstein and therefore Cohen--Macaulay;
neither converse holds in general.

The shift convention is
$$
M(a)_j=M_{a+j}.
$$
In particular, the generator of $S(-b)$
has degree $b$.

\begin{definition}
\label{commutative-algebra-definition-canonical-module-graded}
If $R$ is Cohen--Macaulay, its graded
{\it canonical module} is
$$
\omega_R
=\operatorname{Ext}^{c}_S(R,S(-N-1)).
$$
\end{definition}

Concretely, take a graded free resolution of
$R$, apply $\operatorname{Hom}_S(-,S(-N-1))$,
and take cohomology in degree $c$.
This describes how to compute the displayed
Ext module. The canonical module packages
algebraic duality; its associated sheaf on
$\operatorname{Proj}(R)$ is the dualizing
sheaf. On a smooth $n$-dimensional variety,
the latter is the canonical line bundle
$\bigwedge^n\Omega^1_X$.

\begin{theorem}
\label{commutative-algebra-theorem-graded-gorenstein}
A Cohen--Macaulay standard graded
$k$-algebra $R$ is Gorenstein if and only if
there is an integer $a$ such that
$$
\omega_R\cong R(a)
$$
as graded modules.
This integer is its {\it $a$-invariant}.
\end{theorem}

The extra condition is therefore that the
canonical module is free of rank one, with
one possible change of grading. It does not
say that the shift is zero.

\medskip\noindent
{\bf What ``arithmetically'' means.}

\begin{definition}
\label{commutative-algebra-definition-arithmetically-gorenstein}
Let $X\subset\mathbb P^N_k$ be a nonempty
projective scheme, with saturated homogeneous
ideal $I_X$ and coordinate ring $R=S/I_X$.
It is {\it arithmetically Cohen--Macaulay}
if $R$ is Cohen--Macaulay, and
{\it arithmetically Gorenstein} if $R$ is
Gorenstein.
\end{definition}

These are properties of the embedding:
$R$ is the ring of the affine cone over $X$,
not the ring of global regular functions
on $X$. In particular,
$$
\dim R=\dim X+1.
$$
The cone vertex matters. Smoothness of $X$
makes its local rings Gorenstein, but does
not by itself make its homogeneous
coordinate ring Gorenstein.

\begin{theorem}
\label{commutative-algebra-theorem-ag-subcanonical}
Let $X\subset\mathbb P^N_k$ have positive
dimension. Then $X$ is arithmetically
Gorenstein if and only if it is
arithmetically Cohen--Macaulay and
$$
\omega_X\cong\mathcal O_X(a)
$$
for some integer $a$.
The latter condition is called
{\it subcanonical}.
\end{theorem}

Here $\omega_X$ is the dualizing sheaf,
so the statement also applies to singular
schemes. For an arithmetically
Cohen--Macaulay scheme of positive dimension,
one recovers the canonical module from all
twisted sections of this sheaf; this explains
why the sheaf condition gives the module
condition above.

For a positive-dimensional $X$, the ACM
condition is equivalent to the following
two requirements: the restriction maps
$$
H^0(\mathbb P^N,\mathcal O(t))
\longrightarrow H^0(X,\mathcal O_X(t))
$$
are surjective for every integer $t$, and
$H^i(X,\mathcal O_X(t))=0$ for every $t$
and every $0<i<\dim X$.
Thus being subcanonical alone is not enough.

\begin{example}
\label{commutative-algebra-example-ag-complete-intersection}
Suppose $I=(f_1,\ldots,f_c)$ is generated
by a homogeneous regular sequence with
$\deg f_i=e_i$. The Koszul resolution gives
$$
\omega_{S/I}\cong
(S/I)(e_1+\cdots+e_c-N-1).
$$
Hence a projective complete intersection is
arithmetically Gorenstein. Its canonical
sheaf is trivial when the degrees sum to
$N+1$. Examples are a quintic in
$\mathbb P^4$, an intersection of degrees
$(2,2,3)$ in $\mathbb P^6$, and four
quadrics in $\mathbb P^7$.
Smoothness still has to be checked.
\end{example}

\medskip\noindent
{\bf The Calabi--Yau connection.}
For a smooth projective complex threefold
with $\omega_X\cong\mathcal O_X$ and
$H^1(X,\mathcal O_X)=0$, Serre duality also
gives $H^2(X,\mathcal O_X)=0$.
For positive twists, Kodaira vanishing
kills higher cohomology; for negative twists,
Serre duality gives the required vanishing
of $H^1$ and $H^2$.
Consequently, if its embedding is
projectively normal, it is ACM and hence
arithmetically Gorenstein with $a=0$.
Here projective normality means that the
homogeneous coordinate ring is integrally
closed; for a smooth variety this amounts
to surjectivity of the restriction maps
above in nonnegative degrees.

An arbitrary Gorenstein ring need not have
$a=0$, and an arithmetically Gorenstein
scheme need not be smooth. Neither word
alone asserts that we have a Calabi--Yau
manifold.

\medskip\noindent
{\bf Reading the free resolution.}
For a Cohen--Macaulay quotient $R=S/I$,
its minimal graded free resolution has
length $c$:
$$
0\longrightarrow F_c\longrightarrow\cdots
\longrightarrow F_1\longrightarrow S
\longrightarrow R\longrightarrow0.
$$
Minimal means that the differential matrices
have entries in $(x_0,\ldots,x_N)$.
The summands $S(-b)$ of $F_1$ record the
degrees of minimal generators of $I$;
$F_2$ records their minimal relations,
also called syzygies.

\begin{theorem}
\label{commutative-algebra-theorem-gorenstein-resolution}
For such a Cohen--Macaulay quotient,
$R$ is Gorenstein if and only if $F_c$ has
rank one. If $\omega_R=R(a)$ and
$b=N+1+a$, then
$$
F_c=S(-b),\qquad
F_i\cong
\operatorname{Hom}_S(F_{c-i},S(-b)).
$$
Thus the resolution is self-dual, with a
shift. In particular its ranks are symmetric.
\end{theorem}

Symmetric ranks alone are not a proof of
Gorensteinness: exactness, minimality and
the Cohen--Macaulay hypothesis matter.
For an aG Calabi--Yau threefold in
$\mathbb P^7$, we have $c=4$ and $b=8$.
The familiar rank pattern
$(1,16,30,16,1)$ fits this duality.
The numbers alone do not establish a
smoothing or determine Hodge numbers.

In codimension two, Gorenstein ideals in a
polynomial ring are complete intersections.
In codimension three, the
Buchsbaum--Eisenbud theorem describes them
by maximal Pfaffians of an odd-sized
alternating matrix. A Pfaffian is the
polynomial whose square is the determinant
of an even-sized alternating matrix;
``maximal'' here means deleting one row and
the corresponding column first.
Codimension four is the next substantially
harder case, directly relevant to threefolds
in $\mathbb P^7$.

\medskip\noindent
{\bf Artinian reduction and Lefschetz.}
Over an infinite field, a standard graded
Cohen--Macaulay ring $R$ of dimension $d$
admits $d$ linear parameters forming a
regular sequence. Set
$$
A=R/(\ell_1,\ldots,\ell_d).
$$
Then $A$ is finite-dimensional, and $R$ is
Gorenstein if and only if $A$ is Gorenstein.
If $\omega_R=R(a)$, its socle is concentrated
in degree $s=a+d$.
In this case multiplication gives perfect
pairings
$$
A_i\times A_{s-i}\longrightarrow A_s
\cong k.
$$
In particular, $\dim_k A_i=\dim_k A_{s-i}$.

To see the nondegeneracy, take a nonzero
homogeneous $u\in A_i$. The ideal $uA$
meets the socle: multiply until the next
multiplication by the maximal ideal is
zero. Since the entire socle is $A_s$,
some $v\in A_{s-i}$ has $uv\ne0$.

This is reminiscent of Poincar\'e duality.
The Lefschetz property asks for more:
isomorphisms given by powers of one fixed
linear element. Perfect multiplication
pairings do not supply such an element.
For the strong Lefschetz condition in this
Gorenstein setting, one asks for
$\ell\in A_1$
so that
$$
\ell^{s-2i}:A_i\longrightarrow A_{s-i}
$$
is an isomorphism for every
$0\leq i\leq s/2$.

For an aG Calabi--Yau threefold, $d=4$
and $a=0$, so the Artinian reduction has
socle degree four. This is why quartic
Artinian Gorenstein algebras enter the
study of these projective threefolds.
The Hilbert series satisfy
$$
H_R(t)=\frac{H_A(t)}{(1-t)^d},\qquad
H_A(t)=\sum_i(\dim_k A_i)t^i.
$$
In particular $\deg X=\dim_k A$.
Passing to $A$ preserves this numerical
information, but recovering a smooth
projective variety from $A$ is a separate
problem.

\medskip\noindent
{\bf Why this is useful for Kapustka.}
Coughlan, Go\l\k{e}biowski, G.~Kapustka and
M.~Kapustka study aG Calabi--Yau threefolds
in $\mathbb P^7$. Their 2016 paper gives
examples and a conjectural complete list,
not a proof of completeness.
It uses bilinkage and unprojection among
its construction methods. The point for
us is that equations and their syzygies
encode the canonical sheaf, and allow one
to seek Calabi--Yau varieties beyond
complete intersections.

In the Stanley--Reisner direction, the
face ring of a triangulated sphere is
Gorenstein with $a=0$. Thus the sphere
condition supplies algebraic duality before
any smoothing is constructed. For a
triangulated $3$-sphere this gives a
projective three-dimensional scheme with
trivial dualizing sheaf. It does not prove
that this singular scheme is smoothable.

\medskip\noindent
{\bf Cox--Gorenstein algebras.}
The same socle idea works with several
degrees. We use the connected, positive
setup: $G$ is a finitely generated abelian
group with a translation-invariant partial
order, $A_0=k$, and $A_g=0$ unless $g\geq0$.
Assume $A$ is finite-dimensional over $k$.
Then
$$
\mathfrak m=\bigoplus_{g\ne0}A_g
$$
is its unique maximal ideal.

\begin{definition}
\label{commutative-algebra-definition-cox-gorenstein}
Such a $G$-graded algebra is
{\it Cox--Gorenstein} if, for some
$\rho\in G$,
$$
\operatorname{Soc}(A)=A_\rho\cong k.
$$
The element $\rho$ is its
{\it socle degree}.
\end{definition}

This is the Artinian Gorenstein condition
with the group grading retained.
It is equivalent to perfect multiplication
pairings
$$
A_g\times A_{\rho-g}\longrightarrow A_\rho.
$$
Here a zero piece has a zero partner.
Every nonzero degree satisfies $g\leq\rho$:
$\rho$ is a greatest degree, not merely
one of several incomparable maximal degrees.

Indeed, any nonzero homogeneous principal
ideal meets the socle, by nilpotence of
$\mathfrak m$. Thus for $0\ne u\in A_g$
there is $v\in A_{\rho-g}$ with $uv\ne0$.
This proves nondegeneracy in both factors.
Conversely, such perfect pairings rule out
socle elements in any other degree, since
those elements annihilate their partners.

The Cox terminology comes from toric
homogeneous coordinates: for a complete
normal toric variety $Y$, its Cox ring is
polynomial, with one variable for each fan
ray, graded by the divisor class group
$\operatorname{Cl}(Y)$.
This group consists of Weil divisors modulo
principal divisors. A variable has the class
of its corresponding invariant divisor.

\begin{example}
\label{commutative-algebra-example-cox-bigraded}
For $Y=\mathbb P^1\times\mathbb P^1$,
$$
S=k[x_0,x_1,y_0,y_1],
$$
with $\deg x_i=(1,0)$ and
$\deg y_i=(0,1)$.
The piece $S_{(a,b)}$ consists of
polynomials homogeneous of degree $a$ in
$x_0,x_1$ and degree $b$ in $y_0,y_1$.
It describes sections of
$\mathcal O_Y(a,b)$.

Consider the homogeneous quotient
$$
A=S/(x_0^2,x_1^2,y_0^2,y_1^2).
$$
Its basis consists of the sixteen squarefree
monomials. A monomial is killed by all four
variables exactly when it contains all four.
Hence
$$
\operatorname{Soc}(A)
=kx_0x_1y_0y_1=A_{(2,2)}.
$$
So $A$ is Cox--Gorenstein with
$\rho=(2,2)$. The complementary monomial
pairs each basis monomial with the socle
monomial. Its bigraded Hilbert series is
$$
H_A(u,v)=(1+u)^2(1+v)^2.
$$
Thus its symmetry reflects
$(a,b)\mapsto(2-a,2-b)$.
\end{example}

This example also separates the objects:
$Y$ is the toric variety, $S$ its Cox ring,
and $A$ an Artinian quotient of that ring.
The ring $S$ is itself Gorenstein because
it is polynomial, but it is not Artinian.
The quotient $A$ is not automatically the
cohomology ring of $Y$: here their total
vector-space dimensions are sixteen and
four, respectively.

Nor does $A$ define a nonempty subscheme of
$Y$ by these equations: all four homogeneous
coordinates would have to vanish.
Artinian quotients can encode duality
without being coordinate rings of nonempty
projective subvarieties.

\medskip\noindent
{\bf Multiplication in different directions.}
For homogeneous $L\in A_\lambda$, powers
of $L$ give maps
$$
L^r:A_g\longrightarrow A_{g+r\lambda}.
$$
To reach the dual degree requires
$g+r\lambda=\rho-g$.
Unlike the integer-graded case, several
independent degree directions can occur.
For example, in the quotient above,
$x_0+x_1$ is homogeneous of degree $(1,0)$,
whereas $x_0+y_0$ is not homogeneous for
the bigrading. Forcing total degree loses
this distinction.

In characteristic zero, multiplication by
$(x_0+x_1)^2=2x_0x_1$ identifies
$A_{(0,b)}$ with $A_{(2,b)}$ for
$b=0,1,2$. This is a concrete Lefschetz-type
map in one degree direction. The socle
pairing alone does not assert maximal rank
for every proposed multiplication map.

For the precise toric Lefschetz definitions,
see Definition 5.1 of
\href{https://arxiv.org/abs/2407.17811}%
{Bruzzo--Gondim--Holanda--Montoya,
\textit{Cox--Gorenstein algebras}}.
The socle definition and duality equivalence
are Definition 4.3 and Theorem 4.6 in
version 4. This is the multigraded bridge
from the local algebra above to the toric
lecture; it is distinct from imposing an
arithmetically Gorenstein embedding in
ordinary projective space.

\medskip\noindent
{\bf References and next steps.}
For regular sequences and the local theory,
see Bruns--Herzog,
{\it Cohen--Macaulay Rings},
Chapters 1 and 3, and the Stacks Project,
Tags \href{https://stacks.math.columbia.edu/%
tag/00N7}{00N7} and
\href{https://stacks.math.columbia.edu/%
tag/0DW6}{0DW6}.
For sphere face rings, see Stanley,
{\it Combinatorics and Commutative Algebra},
the chapter on Stanley--Reisner rings.
For the geometric motivation, start with
Sections 1--2 of
\href{https://arxiv.org/abs/1609.01195}%
{Coughlan--Go\l\k{e}biowski--Kapustka--%
Kapustka,
\textit{Arithmetically Gorenstein
Calabi--Yau threefolds in projective
seven-space}}.
A useful algebraic follow-up is Section 1.1
of \href{https://arxiv.org/abs/2011.10871}%
{\textit{Calabi--Yau threefolds in projective
space and Gorenstein rings}}.

A first exercise is to calculate the socles
in the two examples above, then derive the
canonical shift of a complete intersection
from its Koszul resolution. These are the
smallest concrete computations behind the
terminology in the geometry papers.

\section{Minimal resolutions}
\label{commutative-algebra-section-minimal-resolutions}

\noindent
In this section I explain what is a minimal
resolution and what are syzygies. The
importance of this is that our main tool to
obtain flatness shall be changesiArtin's
criterion for flatness (\cite[Theorem
3.15]{Graffeo-Alessi}), which requires
computation of the first syzygy module.

\medskip\noindent
If an $S$-module $M$ is free, then it won't
have any syzygies. If $M$ is not free, then
its first syzygy module is the module
generated by the relations which define  $M$.
More formally: consider a set $A$ of
generators of $M$ and let $S^A$ be a free
$S$-module generated by $A$. This module maps
onto $M$, and its kernel is the {\it first
syzygy module}:
$$
\xymatrix{
\text{Syz}^1(M)\ar[r]
&S^A\ar[r]
&M\ar[r]
&0.
}
$$
If $\text{Syz}^1(M)$ is free, it won't have any syzygies. But if it is not free, it will have syzygies of itself (because by Hilbert Basis Theorem it is finitely generated), giving
$$
\xymatrix{
\text{Syz}^2(M)\ar[r]
&S^{A_1}\ar[r]
&\text{Syz}(M)
&0
}
$$
where $S^{A_1}$ is the free $S$-module
generated by a set $A_1$ of generators of $\text{Syz}^1(M)$.

The {\it minimal resolution} of $M$ is
$$
\xymatrix{
&S^{A_k}\ar[r]
&\cdots\ar[r]
&S^{A_2}\ar[r]
&S^{A_1}\ar[r]
&M\ar[r]
&0.
}
$$
This means that the syzygy modules
{\bf do not} appear in the minimal resolution!
They are modules lying in-between each of the arrows
joining the free modules.
The minimal resolution is independent
(up to isomorphism) of the choices of generating
sets of each of the syzygy modules.

\medskip\noindent
Consider
$$
\xymatrix{
0\ar[r]
&S(-1)\ar[r]^{\cdot x_0}
&S\ar[r]
&S/(x_0)\ar[r]
&0
}
$$

In $S(-1)$ constans have degree 1, linear
polynomials have degree 2, and so on. This
means that the map $\cdot x_0$ maps degree-1
objects of $S(-1)$ to degree-1 objects in
$S$; i.e., it is degree-preserving. Its image
is $(x_0)$, so its cokernel is $S/\text{img}
= S/(x_0)$.

\medskip\noindent
{\bf Explanation of M2 notation}

\noindent
Usually we see in Macaulay something like
\[
\begin{array}{cccccc}
      & S      & \xleftarrow{} S^{16} & \xleftarrow{} S^{30} & 
	\xleftarrow{} S^{16} & \xleftarrow{} S \\
      &        &        &        &        &    \\
      & 0      & 1      & 2      & 3      & 4      \\
      & 1      & 16     & 30     & 16     & 1
\end{array}
\]

Here the module is again $S/I$, but it
doesn't appear in Macaulay's output. Instead,
$S/I$ is the cokernel of the first map
$S^{16}\to S$, and the map we had before $S
\to M$ is the composition $S^{16}\to S \to
S/I$.



\section{Hilbert--Burch and
almost Cohen--Macaulay ideals}
\label{commutative-algebra-section-hilbert-burch}

\begin{remark}[Minimality and grading]
\label{commutative-algebra-remark-resolution-conventions}
The preceding section is the original
writeup. For precision, minimality requires
minimal generating sets at every step.
A free module has no relations when its
chosen generators are a basis.
For $M=S/I$, the augmentation $S\to S/I$
is the quotient map; the composite
$S^{16}\to S\to S/I$ is zero.

Degree-preserving maps let us take each
degree of a resolution separately and
obtain an exact sequence of vector spaces.
For a finite graded free resolution,
$$
\dim_k M_q
=\sum_i(-1)^i\dim_k(F_i)_q.
$$
Thus the shifts determine the Hilbert
function; the free-module ranks alone
do not. Multiplication by a degree-$d$
polynomial $f$ is recorded as
$$
S(-d)_q=S_{q-d}
\xrightarrow{\cdot f}S_q.
$$
\end{remark}

Let $R=k[x_1,\ldots,x_n]$, with its
standard grading, and let $I$ be a
nonzero proper homogeneous ideal.
Recall depth and Cohen--Macaulayness from
Section \ref{commutative-algebra-section-gorenstein-properties}.
The Auslander--Buchsbaum formula gives
$$
\operatorname{pd}_R(R/I)
+\operatorname{depth}(R/I)=n.
$$
Depth is taken at the homogeneous maximal
ideal. In this polynomial-ring setting,
$I$ is called {\it perfect} when
$$
\operatorname{pd}_R(R/I)
=\operatorname{ht}(I).
$$
Equivalently, $R/I$ is Cohen--Macaulay.
Thus a perfect height-two ideal has
$$
\operatorname{pd}_R(R/I)=2,
\qquad\operatorname{pd}_R I=1.
$$
Its minimal resolution has $F_i=0$ for
$i\geq3$. This length statement comes from
Auslander--Buchsbaum; Hilbert--Burch also
recovers the equations from a matrix.

\begin{theorem}[Hilbert--Burch]
\label{commutative-algebra-theorem-hilbert-burch}
Suppose that $I$ has height two, $R/I$ is
Cohen--Macaulay, and $I$ has $r$ minimal
generators. Omitting grading shifts, its
minimal free resolution has the form
$$
0\longrightarrow R^{r-1}
\xrightarrow{\Phi}R^r
\longrightarrow R\longrightarrow R/I
\longrightarrow0.
$$
The ideal $I$ is generated by the maximal
minors of the $r\times(r-1)$ matrix
$\Phi$. If $\Phi_i$ deletes row $i$, put
$$
\Delta_i=(-1)^{i+1}\det\Phi_i.
$$
Up to a common nonzero scalar, the row
giving $R^r\to R$ is
$$
(\Delta_1,\ldots,\Delta_r).
$$
Conversely, take a homogeneous
$r\times(r-1)$ matrix whose maximal
minors generate an ideal of height two.
Its quotient is Cohen--Macaulay, with a
free resolution of the displayed form.
This resolution is minimal
when the matrix entries lie in
$(x_1,\ldots,x_n)$.
\end{theorem}

See \cite[Theorem 5.1]{migliore-cocoa},
available in
\href{https://academicweb.nd.edu/%
~jmiglior/CoCoANotes.pdf}{Migliore's notes}.
The columns of $\Phi$ form a free basis
of the relations among the generators
of $I$. The relation module is free,
not zero: it is the next kernel that
vanishes.
The Gr\"unbaum example above has
codimension four, so its longer resolution
does not contradict this theorem.

\medskip\noindent
{\bf The first non-perfect case.}
Now assume $n\geq3$ and that $J$ has
height two and three minimal generators
of the same degree $d\geq2$.
The almost Cohen--Macaulay case here means
$$
\operatorname{depth}(R/J)
=\dim(R/J)-1=n-3.
$$
Equivalently,
$$
\operatorname{pd}_R(R/J)=3,
\qquad\operatorname{pd}_R J=2.
$$
Ignoring shifts, the minimal resolution is
$$
0\longrightarrow R^{m-2}
\longrightarrow R^m
\longrightarrow R^3
\longrightarrow R\longrightarrow R/J
\longrightarrow0,
\qquad m\geq3.
$$
The next theorem provides a determinantal
description for precisely this class.

\medskip\noindent
Here are the data needed to state it.
Choose integers $m\geq3$,
$1\leq\delta_1\leq\cdots\leq\delta_m$
and $\epsilon_1,\ldots,\epsilon_{m-2}\geq1$
satisfying
$$
\begin{aligned}
\delta_3+\epsilon_1
&\leq\cdots\leq\delta_m+\epsilon_{m-2},\\
\delta_1+\delta_2
&=d+\sum_{i=1}^{m-2}\epsilon_i,\\
\delta_i+\delta_j&\geq d+1
\quad(1\leq i<j\leq m),\\
\delta_3&\leq d.
\end{aligned}
$$
These are the paper's {\it latent data}.
A corresponding {\it level matrix} is
$$
\eta=\begin{pmatrix}A\\B\end{pmatrix},
$$
where $A$ has three rows, $B$ has $m-2$
rows, and both have $m$ columns. Their
entries have the prescribed degrees
$$
a_{ij}\in R_{d-\delta_j},
\qquad
b_{ij}\in
R_{\delta_{i+2}+\epsilon_i-\delta_j}.
$$
Here $R_q=0$ for $q<0$; additionally,
$b_{ij}=0$ whenever its prescribed degree
is zero. Writing $I_s(C)$ for the ideal
of all $s\times s$ minors of $C$, require
$$
\operatorname{ht}I_m(\eta)=2,
\qquad
\operatorname{ht}I_{m-2}(B)=3.
$$

\begin{theorem}[Burity--Fiel--Ramos--Simis,
2026]
\label{commutative-algebra-theorem-bf-rs}
Under the assumptions on $R$ and $J$ above,
$R/J$ is almost Cohen--Macaulay if and
only if there are latent data and a
corresponding level matrix $\eta$ such that
$$
J=(\det\eta_{\widehat1},
   \det\eta_{\widehat2},
   \det\eta_{\widehat3}).
$$
Here $\eta_{\widehat i}$ deletes the $i$th
row of $\eta$: each determinant retains
all rows of $B$ and two of the three rows
of $A$.
\end{theorem}

This is \cite[Theorem 1.3, Setup 1.1 and
Definition 1.2]{bf-rs-2026}, version 4.
Thus Hilbert--Burch uses all maximal
minors; this theorem uses the three that
retain the lower block. It does not cover
arbitrary non-perfect ideals.


\section{Colon ideals and saturation}
\label{commutative-algebra-section-colon-ideals-saturation}

Let $R$ be a commutative Noetherian ring,
and let $I,J\subseteq R$ be ideals.

\begin{definition}
\label{commutative-algebra-definition-colon-ideal}
The \textit{colon ideal}, or
\textit{ideal quotient}, is
$$
I:J=\{r\in R:rJ\subseteq I\}.
$$
For $J=(f)$, write $I:f$.
This is an ideal of $R$, whereas $R/I$
is a quotient ring.
\end{definition}

If $J=(g_1,\ldots,g_s)$, then $r\in I:J$
means that every product $rg_i$ lies in
$I$. In particular,
$$
I:J=\bigcap_{i=1}^s(I:g_i).
$$

\begin{example}
\label{commutative-algebra-example-colon-two-lines}
In $k[x,y]$,
$$
(xy):(x)=(y).
$$
Indeed, $fx$ is divisible by $xy$ exactly
when $f$ is divisible by $y$.
The union $V(xy)$ consists of two axes;
taking this colon ideal leaves the axis
residual to $V(x)$.
\end{example}

\begin{definition}
\label{commutative-algebra-definition-ideal-saturation}
The \textit{saturation of $I$ with
respect to $J$} is
$$
I:J^\infty=\bigcup_{m\geq1}(I:J^m).
$$
We say that $I$ is \textit{$J$-saturated}
if $I:J^\infty=I$.
\end{definition}

Thus membership allows multiplication
by a sufficiently high power of the
whole ideal $J$. The ascending chain
$$
I:J\subseteq I:J^2\subseteq\cdots
$$
stabilizes because $R$ is Noetherian.
For instance,
$$
(x^2y):(x)=(xy),
\qquad
(x^2y):(x)^\infty=(y).
$$
A single colon removes one factor of
$x$ here; saturation removes all of them.
These operations therefore need not agree.

For ideals in a polynomial ring over an
algebraically closed field, the geometric
interpretation on underlying closed sets
is
$$
V(I:J^\infty)
=\overline{V(I)\setminus V(J)}.
$$
Saturation takes the closure of the part
outside $V(J)$; it need not make the
remaining scheme reduced.

\begin{definition}
\label{commutative-algebra-definition-projective-saturation}
For a homogeneous ideal
$I\subseteq S=k[x_0,\ldots,x_n]$, put
$\mathfrak m=(x_0,\ldots,x_n)$.
Its \textit{projective saturation} is
$$
I^{\mathrm{sat}}=I:\mathfrak m^\infty.
$$
\end{definition}

The ideals $I$ and $I^{\mathrm{sat}}$
define the same closed subscheme of
$\mathbb{P}^n$. This saturation removes
information supported only at the vertex
of the affine cone. In linkage, the
residual ideal is instead obtained by
taking a colon with the ideal of the
chosen component.

\section{Linkage}
\label{commutative-algebra-section-linkage}

\begin{definition}[Direct linkage]
\label{commutative-algebra-definition-direct-linkage}
For reduced $U,V$ without common
components in an arithmetically
Gorenstein $G\subset\mathbb{P}^n$,
\cite[Section 2.2]{luk2} states:
\begin{quote}
``We say that $U$ and $V$ are
\textit{directly linked} if their union
is a complete intersection $F$ in $G$''
\end{quote}
Scheme-theoretically, using ideals in $G$:
$$
\begin{aligned}
\mathcal{I}_F:\mathcal{I}_U
&=\mathcal{I}_V,\\
\mathcal{I}_F:\mathcal{I}_V
&=\mathcal{I}_U.
\end{aligned}
$$
Locally, $J:K$ consists of functions $f$
satisfying $fK\subseteq J$.
Thus each scheme is the residual
to the other inside $F$.
Two successive links give a
\textit{bilinkage}.
\end{definition}

\medskip\noindent
\textbf{Liaison and even liaison.}
Recall that \emph{liaison} means the
equivalence relation generated by chains
of direct links:
$$
X=X_0\longleftrightarrow X_1
\longleftrightarrow\cdots
\longleftrightarrow X_r=Y.
$$
One arrow is a direct link; an arbitrary
finite chain gives a liaison, and a chain
of even length gives an \emph{even liaison}.
The union of the two endpoints need not
itself be a complete intersection.

The ambient scheme and permitted links
must be specified. For \emph{CI-liaison}
in projective space, every linking scheme
is a complete intersection there.
For \emph{Gorenstein liaison}, it may be
any arithmetically Gorenstein scheme;
the two residual-ideal identities still
define the link. A complete intersection
inside an arithmetically Gorenstein
ambient scheme gives a Gorenstein link
in projective space, but need not give a
CI-link there.

\begin{theorem}[Rao for space curves]
\label{commutative-algebra-theorem-rao-space-curves}
Let $C,C'\subset\mathbb{P}^3_k$ be pure
one-dimensional locally Cohen--Macaulay
closed subschemes over an algebraically
closed field. Define their \emph{Rao
modules} by
$$
M_C=\bigoplus_{t\in\mathbb{Z}}
H^1(\mathbb{P}^3,\mathcal{I}_C(t)).
$$
Multiplication by homogeneous polynomials
defines the graded module structure.
Then $C$ and $C'$ belong to the same
even CI-liaison class if and only if
their Rao modules are isomorphic up to
grading shift:
$$
M_C\cong M_{C'}[e],
\qquad N[e]_t=N_{t+e}.
$$
\end{theorem}

\noindent
See the beginning of Section 3 of
Hartshorne,
\href{%
https://arxiv.org/abs/math/0302078%
}{\emph{On Rao's Theorems and the
Lazarsfeld--Rao Property}}.
In degree $t$, the module measures the
failure of ambient sections to restrict
onto all sections on the curve:
$$
(M_C)_t\cong
\operatorname{coker}\bigl(
H^0(\mathbb{P}^3,\mathcal{O}(t))
\longrightarrow
H^0(C,\mathcal{O}_C(t))\bigr).
$$
In particular, $M_C=0$ precisely when
$C$ is arithmetically Cohen--Macaulay.

\begin{remark}[The codimension-four case]
\label{commutative-algebra-remark-rao-codimension-four}
The curve theorem does not directly apply
to a threefold in $\mathbb{P}^7$.
For the arithmetically Gorenstein
Gr\"unbaum Stanley--Reisner threefold
$X_0$, the intermediate deficiency
modules already vanish:
$$
H^i(\mathbb{P}^7,\mathcal{I}_{X_0}(t))=0,
\qquad
1\leq i\leq3,\quad t\in\mathbb{Z}.
$$
This vanishing alone supplies neither a
particular linkage nor its behaviour
under deformation.

There is a sheaf-valued Rao correspondence
for codimension-two subschemes of an
integral arithmetically Gorenstein
ambient scheme: the invariant uses
suitable resolution sheaves, up to
twist and stable equivalence.
Stable equivalence allows adding direct
sums of line bundles to both sheaves.
See Theorem 2.16 of
\href{%
https://arxiv.org/abs/math/0307378%
}{Casanellas--Drozd--Hartshorne,
\emph{Gorenstein Liaison and ACM Sheaves}}.
\end{remark}

\begin{theorem}[Casanellas--Drozd--%
Hartshorne]
\label{commutative-algebra-theorem-ag-glicci}
Every arithmetically Gorenstein subscheme
of projective space of codimension at
least two belongs to the Gorenstein
liaison class of a complete intersection.
\end{theorem}

\noindent
This is Theorem 7.1 of the paper above.
The property is called \emph{glicci}.
It applies to $X_0$.
The links need not be complete-intersection
links, and the theorem does not supply
a rational partner or a deformation of
the linkage along a chosen smoothing.

Its proof puts $X_0$ inside a suitable
complete-intersection fivefold
$W\subset\mathbb{P}^7$, making $X_0$
codimension two in $W$, and uses
$$
0\longrightarrow\mathcal{O}_W(-a)
\longrightarrow\mathcal{E}
\longrightarrow\mathcal{I}_{X_0/W}
\longrightarrow0.
$$
Here $\mathcal{E}$ is a rank-two ACM
sheaf, not necessarily a vector bundle.
This offers a concrete sheaf-theoretic
starting point for the linkage question.

\begin{example}[From projective space
to a Calabi--Yau threefold]
\label{commutative-algebra-example-bilinkage-degree-seventeen}
The construction in
\cite[Section 4.4]{luk2} starts with
a linear $\Pi\cong\mathbb{P}^3$:
$$
\Pi\subset Y=V(q_1,q_2,q_3)
\subset\mathbb{P}^7,
\qquad \deg Y=2^3=8.
$$
Choose general quadrics $q_i$ through
$\Pi$, then a general hyperplane
$H=V(h)$ through $\Pi$.
The fourfold $Y$ has dimension $7-3=4$.
Intersecting with $H$ lowers the dimension
to three, allowing $\Pi$ to be an
irreducible component.

Its residual threefold $D$ satisfies
$$
\begin{aligned}
Y\cap H&=\Pi\cup D,\\
I_D&=(q_1,q_2,q_3,h):I_\Pi,\\
\deg D&=8-1=7.
\end{aligned}
$$
The resolution of $I_D$ supplies
a cubic $c$ containing $D$ but not $\Pi$.
Take the residual again:
$$
\begin{aligned}
Y\cap V(c)&=D\cup X,\\
I_X&=(q_1,q_2,q_3,c):I_D,\\
\deg X&=3\cdot8-7=17.
\end{aligned}
$$
The paper computes a resolution proving
that $X$ is arithmetically Gorenstein,
with
$$
\omega_X\cong\mathcal{O}_X,
\qquad H^1(X,\mathcal{O}_X)=0.
$$
It verifies smoothness for an example;
general choices therefore yield
Calabi--Yau threefolds.
Lemma 4.3 establishes that this
construction reaches general members
of their deformation family.
\end{example}

\noindent
This constructs Calabi--Yau threefolds
from the rational variety $\mathbb{P}^3$.
Linkage does not imply birational
equivalence.

\section{Kähler differentials}
\label{commutative-algebra-section-kahler-differentials}

\noindent
The idea behind Kähler differentials
is that they are inital objects
in some category…

They were invented by Kähler in his
attempt to rewrite algebraic geometry
from zero. This was one of the insights
he had.

\begin{definition}
\label{commutative-algebra-definition-kahler-differentials}
Let $R$ be a ring over a field $k$.
For every $a\in R$, introduce a formal
symbol $da$. At this point $da$ is
literally just a symbol. Take the free
$R$-module generated by all these symbols,
and quotient by the relations expressing
additivity, $k$-linearity, and Leibniz:
$$
d(a+b)=da+db,\qquad
d(\lambda a)=\lambda\,da,\qquad
d(ab)=a\,db+b\,da,
$$
for $a,b\in R$ and $\lambda\in k$.
Equivalently, we quotient by the
submodule generated by the differences
between the two sides of these equations.
The resulting $R$-module is
$\Omega^1_{kR}$, the module of
K\"ahler differentials.
\end{definition}

\noindent
Thus $\Omega^1_{kR}$ is not $R$ again:
it is a new $R$-module made from formal
symbols $da$, with precisely the relations
that make $a\mapsto da$ behave like
differentiation.

\noindent
You can also think of this as generated
by $D(a)$ for $a \in R$ modulo the relation
of the Leibniz identity.

\begin{exercise}[Universal property
of Kähler differentials]
\label{commutative-algebra-exercise-kahler-differentials-derivation}
Prove that the map
$d:R\to\Omega^1_{kR}$ mapping $a$ to $da$
is a derivation.
\end{exercise}

\begin{lemma}
\label{commutative-algebra-lemma-universal-property-kahler-differentials}
Let $V$ be an $R$-module, and
$D\in\operatorname{Der}_k(R,V)$ a
derivation. Then there exists a unique
$R$-module homomorphism
$\varphi_D:\Omega^1_{kR}\to V$ mapping
$da$ to $D(a)$.
\end{lemma}
$$
\xymatrix{
R \ar[r]^d \ar[dr]_{\delta}&  
\Omega^1\ar@{-->}[d]\\
&V
}
$$


\begin{remark}
\label{commutative-algebra-remark-universal-property-kahler-differentials}
Consider a category $\mathcal C$ of
$R$-modules equipped with a derivation
$(V,D:R\to V)$, and define morphisms in
$\mathcal C$ as morphisms of $R$-modules
which commute with the derivation map.
Then $\Omega^1_{kR}$ is an initial object
in this category. This is called the
universal property of the module of
K\"ahler differentials.
\end{remark}

\begin{lemma}
\label{commutative-algebra-lemma-derivations-hom}
$\operatorname{Der}_k(R,V)
=\operatorname{Hom}_R(\Omega^1_{kR},V)$.
\end{lemma}

\begin{proof}
A composition of a derivation
$R\xrightarrow{d}\Omega^1_{kR}$ and an
$R$-module homomorphism
$\Omega^1_{kR}\to V$ lies in
$\operatorname{Der}_k(R,V)$. On the other
hand, any derivation
$\xi\in\operatorname{Der}_k(R,V)$ is
obtained this way, by the universal
property.
\end{proof}

\begin{lemma}[Corollary]
\label{commutative-algebra-lemma-derivations-dual}
$\operatorname{Der}_k(R)
=(\Omega^1_{kR})^\ast$, where
$V^\ast:=\operatorname{Hom}(V,R)$.
\end{lemma}

\noindent
Now let's se how this turns out
for polynomials.

\begin{lemma}
\label{commutative-algebra-lemma-derivations-polynomial-ring}
Let $R=k[t_1,\dots,t_n]$.
Then $\operatorname{Der}_k(R)$ is a free
$R$-module generated by
$\frac{d}{dt_1},\dots,\frac{d}{dt_n}$.
\end{lemma}

\begin{proof}
Any derivation is determined by its values
on $t_1,\dots,t_n$. Indeed, if
$\delta\in\operatorname{Der}_k(R)$ and
$\delta(t_i)=a_i\in R$, then for every
$P\in R$ we have
$$
\delta(P)
=\sum_{i=1}^n a_i
\frac{dP}{dt_i}.
$$
Hence the natural map from the free module
$\langle \frac{d}{dt_1},\dots,\frac{d}{dt_n}
\rangle$ to $\operatorname{Der}_k(R)$ is
surjective.

To see that it is injective, suppose
$$
\sum_{i=1}^n a_i\frac{d}{dt_i}=0
$$
as a derivation. Evaluating on $t_j$ gives
$a_j=0$ for each $j$, so the expression was
already zero in the free module. Thus
\[
\operatorname{Der}_k(R)
\cong
\langle \frac{d}{dt_1},\dots,\frac{d}{dt_n}
\rangle.
\]
\end{proof}

\noindent
The following result is precisely what
fails for smooth manifolds.

\begin{proposition}
\label{commutative-algebra-proposition-kahler-differentials-polynomial-ring}
Let $R=k[t_1,\dots,t_n]$.
Then $\Omega^1_R$ is a free $R$-module
generated by $dt_1,\dots,dt_n$.
\end{proposition}

\begin{proof}
By the Leibniz rule, $\Omega^1_R$ is
generated by $dt_1,\dots,dt_n$.
Let
$$
\Psi:\langle dt_1,\dots,dt_n\rangle
\to \Omega^1_R
$$
be the map defined this way. It is
surjective, so it remains to show that
$\Psi$ is injective.

For any $R$-module $M$, write
$M^\ast:=\operatorname{Hom}_R(M,R)$.
By the universal property,
$(\Omega^1_R)^\ast=\operatorname{Der}_k(R)$.
Since $\operatorname{Der}_k(R)$ is a free
$R$-module, $(\operatorname{Der}_k(R))^\ast$
is also free with basis
$
\left(\frac{d}{dt_1}\right)^\ast,\dots,
\left(\frac{d}{dt_n}\right)^\ast.
$
The natural map
$\Omega^1_R\to (\Omega^1_R)^{\ast\ast}$
takes $dt_i$ to
$\left(\frac{d}{dt_i}\right)^\ast$, hence
$\Psi$ becomes injective after applying
$(-)^\ast$. Therefore $\Psi$ itself is
injective.
\end{proof}

\begin{exercise}
\label{commutative-algebra-exercise-kahler-differentials-polynomial-ring}
Prove that
$\Omega^1_{\mathbb{C}[t_1,\dots,t_n]}$
is a free $\mathbb{C}[t_1,\dots,t_n]$-module
generated by $dt_1,\dots,dt_n$.
\end{exercise}

\begin{proof}
It's clear that $\Omega^1$ is generated
by the $dt_i$.
Now we know that we can pair forms
and vector fields:
$$
\Omega^1(\mathbb{C}[t_1,…,t_n]
=
\text{Der}(\mathbb{C}[t_1,…,t_n])
=\langle \frac{d}{dt}\rangle
$$
but we can pair them! By the pairing,
both are free.
\end{proof}

\noindent
The following lemma says $\Spec(B)$
is a subvariety of $\Spec(A)$ (surjectivity).
Then $I/I^2$ corresponds
with the conormal bundle.
The idea is:
$$
\xymatrix{
0\ar[r]
&N^\vee\ar[r]
&\Lambda^1M|_X\ar[r]
&\Lambda^1X\ar[r]
&0
}
$$

\begin{lemma}
\label{commutative-algebra-lemma-exact-sequence-kahler-differentials}
Let $\varphi:A\to B$ be a surjective ring
homomorphism, and $I$ its kernel. Then the
following sequence of $B$-modules is exact:
$$
I/I^2\xrightarrow{\delta}
\Omega^1_A\otimes_A B
\xrightarrow{\varphi}
\Omega^1_B\to 0.
$$
\end{lemma}

\begin{proof}
Let $\delta:I\to\Omega^1_A$ take
$x\mapsto dx$. Since
$d(xy)=yd(x)+xd(y)$, we have
$\delta(I^2)=0$ modulo $I\Omega^1_A$.
Therefore, the map
$I/I^2\xrightarrow{\delta}
\Omega^1_A\otimes_A B$ is well defined.

The composition $\delta\circ\varphi$
clearly vanishes, and
$\Omega^1_A\otimes_A B\to\Omega^1_B$ is
clearly surjective. The kernel of
$\Omega^1_A\to\Omega^1_B$ is generated by
$I\Omega^1_A$ together with $d(I)$, hence
the kernel of
$\Omega^1_A\otimes_A B\to\Omega^1_B$
is equal to $d(I)$.
\end{proof}

\noindent
Now, if we assume there exists a section
we also get injectivity. The idea
is to define a new map
$\rho$ and prove it's a derivation;
then apply the universal property
of Kähler differentials.

\begin{lemma}
\label{commutative-algebra-lemma-exact-sequence-kahler-differentials-section}
Let $\varphi:A\to B$ be a surjective ring
homomorphism, and $I$ its kernel. Assume
that $\varphi$ admits a section
$\sigma:B\to A$. Then the sequence
$$
0\to I/I^2\xrightarrow{\delta}
\Omega^1_A\otimes_A B
\xrightarrow{\varphi}
\Omega^1_B\to 0
$$
is exact.
\end{lemma}

\begin{proof}
Consider the map $\rho:A\to I/I^2$ taking
$x$ to $x-\sigma(\varphi(x))$. Since
$(x-\sigma(\varphi(x)))(y-\sigma(\varphi(y)))\in I^2$,
we have
$$
xy+\sigma(\varphi(x))\sigma(\varphi(y))
= x\sigma(\varphi(y))
+y\sigma(\varphi(x)).
$$
Then
$$
x(y-\sigma(\varphi(y)))
+y(x-\sigma(\varphi(x)))
= 2xy-xy-\sigma(\varphi(x))\sigma(\varphi(y))
=xy-\sigma(\varphi(x))\sigma(\varphi(y)).
$$
Therefore, $\rho$ is a derivation.
Since $\rho$ takes an element of $I$ to
itself, the corresponding universal map
$\Omega^1_A\otimes_A B\to I/I^2$ is by
construction inverse to $\delta$.
This implies that
$I/I^2\xrightarrow{\delta}
\Omega^1_A\otimes_A B$ is injective.
\end{proof}

\section{Zariski cotangent space}
\label{commutative-algebra-section-zariski-cotangent-space}

\begin{definition}
\label{commutative-algebra-definition-zariski-cotangent-space}
Let $m \subset R$ be a maximal ideal,
associated with $x \in X = \operatorname{Spec} R$.
The Zariski cotangent space $T_x^\ast X$
is $m/m^2$.
\end{definition}

\noindent
Recall that bundles over varieties
are in bijective correspondence with
projective modules over the ring of
functions.

The following proposition says that
the cotangent space is the fiber of
the cotangent bundle at a closed point!

The upshot is that we can construct
a section as in the last lemma
from the last section because the
residue field is $k$.

\begin{proposition}
\label{commutative-algebra-proposition-zariski-cotangent-space-kahler}
Let $R$ be a ring over $k$, and let
$m \subset R$ be a maximal ideal such that
$R/m = k$. Then the Zariski cotangent space
$T_m^\ast \operatorname{Spec} R$ is equal to
$\Omega^1_{kR}\otimes_R R/m$.
\end{proposition}

\begin{proof}
Let $R_1 := R/m$. Since $R/m = k$, the map
$R \to R_1$ has a section. Then the exact
sequence
$$
0 \to m/m^2 \xrightarrow{\delta}
\Omega^1_{kR}\otimes_R R_1
\to \Omega^1_{kR_1} = 0
$$
implies that
$m/m^2 = \Omega^1_{kR}\otimes_R R_1$.
\end{proof}

\begin{lemma}[Corollary]
\label{commutative-algebra-lemma-zariski-cotangent-space-manifold}
If $X$ is a smooth manifold, $x \in X$, and
$m_x$ is the ideal of all smooth functions
vanishing at $x$, then $m_x^2$ is the ideal
of all functions $f$ such that
$f(x)=0$ and $df|_x=0$. In particular, the
natural map
$m_x/m_x^2 \to T_x^\ast X$ mapping
$f$ to $df|_x$ is an isomorphism.
\end{lemma}

This explains the term ``Zariski cotangent
space''.

\medskip\noindent
The following proposition is another
incarnation of the Jacobian criterion!

\begin{proposition}
\label{commutative-algebra-proposition-regular-points-affine-variety}
Let $X \subset \mathbf{C}^n$ be an affine
variety defined by a system of polynomial
equations
$f_1(z)=f_2(z)=\cdots=f_k(z)=0$, and let
$B=\mathcal O_X$ be its coordinate ring.
For a point $x \in X$, let
$W \subset T_x^\ast \mathbf{C}^n$
be the subspace generated by
$df_1,\dots,df_k$.
Then
$T_x^\ast X = T_x^\ast \mathbf{C}^n/W$.
Moreover, $B_x$ is regular if and only if
$\dim X = n-\dim W$
(Theorem \ref{commutative-algebra-theorem-regular-rings}).
\end{proposition}

\begin{proof}
Let $A=\mathbf{C}[t_1,\dots,t_n]$, and let
$I$ be the kernel of the natural surjection
$A \to B$. Consider the exact sequence
$$
I/I^2 \xrightarrow{\delta}
\Omega^1_A \otimes_A B
\xrightarrow{\varphi}
\Omega^1_B \to 0.
$$
Then
$$
T_x^\ast B
=\Omega^1_B \otimes_A A/m_x
=\Omega^1_A \otimes_A A/m_x
\big/ \delta(I/I^2)
$$
and
$$
\delta(I/I^2)
=\langle dt_1,\dots,dt_n\rangle
\big/\langle df_1,\dots,df_k\rangle
\otimes_A A/m_x.
$$
Therefore, the Zariski cotangent space
$T_x^\ast X
=\Omega^1_B\otimes_A A/m_x
=T_x^\ast \mathbf{C}^n/W$
is $(n-\dim_k W)$-dimensional. On the other
hand, $B_x$ is regular if and only if
$\dim T_x^\ast X = \dim X$
(Theorem \ref{commutative-algebra-theorem-regular-rings}).
\end{proof}

\begin{lemma}[Corollary]
\label{commutative-algebra-lemma-regular-values-affine-variety}
Let $X \subset \mathbf{C}^n$ be an affine
variety defined by a system of polynomial
equations
$f_1(z)=f_2(z)=\cdots=f_k(z)=0$.
Suppose that $0$ is a regular value of the
map
$F(z)=\bigl(f_1(z),f_2(z),\dots,f_k(z)\bigr)$.
Then $T_x^\ast X = \dim X$ for any
$x \in X$.
In particular, all localizations of $X$
in maximal ideals are regular.
\end{lemma}

\begin{proof}
Since $X$ is given by $k$ equations,
the dimension of $X$ is at least $n-k$.
On the other hand,
$\dim T_x^\ast X = n-\dim W = n-k$
as shown above. By Theorem
\ref{commutative-algebra-theorem-regular-rings},
$\dim T_x^\ast X \geq \dim X$, and the
equality is realized precisely when $x$ is
regular.
\end{proof}

\begin{definition}
\label{commutative-algebra-definition-smooth-point-variety}
A point on an algebraic variety is called
{\it smooth} if the local ring associated
with the corresponding maximal ideal is
regular. A variety is called {\it smooth}
if all its points are regular.
\end{definition}

\section{Conormal exact sequence}
\label{commutative-algebra-section-conormal-exact-sequence}

\noindent
It was an idea of Grothendieck
that one should define things not for
a variety alone, but for a variety
``over'' another variety.
Surprisingly, this simplifies computations
a lot --- particularly that of the
Riemann-Roch-Grothendieck theorem.

In this case, we can think of the
relative Kähler differentials
as the vertical cotangent bundle:
think that $X \to Y$ is a fibration,
so we look at the cotangent spaces
of the fibers.

\begin{definition}
\label{commutative-algebra-definition-relative-kahler-differentials}
Let $f:X \to Y$ be a morphism of complex
affine varieties associated with a ring
homomorphism $f^\ast:B\to A$. The \emph{module of relative K\"ahler differentials}
$\Omega^1_{X/Y}$ is the module of
$B$-linear K\"ahler differentials on $X$,
that is, the module generated by $da$ for
$a \in A$ with relations
$d(uv)=u\,dv+v\,du$ and $d(b)=0$ for all
$b\in B$.
\end{definition}

\begin{remark}
\label{commutative-algebra-remark-relative-kahler-differentials}
The last condition is equivalent to
$B$-linearity:
$d(b)=0 \iff d(ab)=b\,d(a)$.
That is, everything is linear
over the smaller ring.
\end{remark}

\begin{proposition}[Conormal exact sequence]
\label{commutative-algebra-proposition-conormal-exact-sequence}
Let $X \xrightarrow{f} Y \xrightarrow{g} Z$
be a morphism of affine varieties.
Then there is a natural exact sequence
$$
\Omega^1_{Y/Z}\otimes_{\mathcal O_Y}\mathcal O_X
\xrightarrow{\psi}
\Omega^1_{X/Z}
\xrightarrow{\phi}
\Omega^1_{X/Y}\to 0.
$$
\end{proposition}

\begin{proof}
Surjectivity in the last term is clear. An
image of $\Omega^1_{Y/Z}\otimes_{\mathcal O_Y}
\mathcal O_X$ vanishes in $\Omega^1_{X/Y}$
because $dy=0$ in $\Omega^1_{X/Y}$ for any
$y\in \mathcal O_Y$. It remains to prove
that this sequence is exact in the middle
term.
\medskip\noindent
Step 2: Let
$Q:=\Omega^1_{X/Z}/\operatorname{im}\psi$,
and let $\delta(x)$ be the image of $dx$ in
$Q$ for any $x\in \mathcal O_X$. This is
clearly an $\mathcal O_Y$-linear derivation,
hence defines an element of
$\Omega^1_{X/Y}$. Therefore,
$\operatorname{im}\psi \supset \ker\phi$.
\end{proof}

\begin{lemma}[Corollary]
\label{commutative-algebra-lemma-conormal-exact-sequence-ideal}
Let $X\hookrightarrow Y$ be a closed
embedding defined by an ideal $I$. Then
there is a natural exact sequence
$$
I/I^2\longrightarrow
\Omega^1_Y\otimes_{\mathcal O_Y}\mathcal O_X
\longrightarrow \Omega^1_X\longrightarrow0.
$$
This is the conormal exact sequence for a
closed embedding.
\end{lemma}

\begin{proof}
On affine coordinate rings this is exactly
Lemma \ref{commutative-algebra-lemma-exact-sequence-kahler-differentials}:
if $A\twoheadrightarrow B=A/I$, then
$$
I/I^2\longrightarrow
\Omega^1_A\otimes_A B
\longrightarrow\Omega^1_B\longrightarrow0.
$$
\end{proof}

\begin{lemma}[Corollary]
\label{commutative-algebra-lemma-cotangent-space-point}
Let $y$ be a point on $Y$.
Then $\Omega^1_{\{y\}/Y}=T_y^\ast Y$.
\end{lemma}

\section{Smooth morphisms}
\label{commutative-algebra-section-smooth-morphisms}

\noindent
{\bf Upshot.}
A smooth morphism is the algebraic
analogue of a submersion: a family whose
fibres are nonsingular and vary flatly.
The model is the projection
$
Y\times\mathbb{A}^d\longrightarrow Y.
$
Smoothness concerns the map, not merely
whether its source and target are smooth.

\begin{definition}
\label{commutative-algebra-definition-smooth-morphism}
A morphism $f:X\to Y$ is {\it smooth}
if it is locally of finite presentation,
flat, and has geometrically regular
fibres. The last condition means that
each fibre remains regular after every
extension of its residue field.
\end{definition}

\noindent
For morphisms of complex varieties, this
amounts to flatness and nonsingular
fibres. If both varieties are smooth,
one can instead check that
$
df_x:T_xX\longrightarrow T_{f(x)}Y
$
is surjective at every point.

For a smooth morphism of relative
dimension $d$, the relative differentials
$\Omega^1_{X/Y}$ form a vector bundle
of rank $d$: they measure cotangent
directions along the fibres.

\medskip\noindent
{\bf Example.}
The map
$
\mathbb{A}^2_{\mathbb{C}}
\longrightarrow\mathbb{A}^1_{\mathbb{C}},
\qquad (x,y)\longmapsto xy
$
is smooth away from the origin.
At the origin its differential vanishes,
and the fibre $xy=0$ has a node.

\medskip\noindent
Reference: Stacks Project,
\href{https://stacks.math.columbia.edu/%
tag/01V4}{Smooth morphisms},
Lemmas 29.35.12 and 29.35.14.

\section{Étale maps}
\label{commutative-algebra-section-etale-maps}

\noindent
Upshot: Étale is finite cover in
the complex analytic category.

\begin{definition}
\label{commutative-algebra-definition-etale-map}
Let $f:X \to Y$ be a morphism of complex
varieties. We say that $f$ is {\it unramified}
if $\Omega^1_{X/Y}=0$, and {\it \'etale} if
it is unramified and finite.
\end{definition}

\begin{remark}
\label{commutative-algebra-remark-etale-map-flat}
Caveat emptor: when the base field is not
algebraically closed, one needs to ask that
$f$ is {\it flat} as well.
\end{remark}

\begin{theorem}
\label{commutative-algebra-theorem-etale-cotangent-isomorphism}
Let $g:Y \to Z$ be an \'etale map. Then
$dg^\ast$ induces an isomorphism
$T^\ast_{g(y)}Z \to T^\ast_yY$ for any
$y\in Y$.
\end{theorem}

\begin{proof}
Let $X:=\{y\}$, and let $m$ be its maximal
ideal. Consider the conormal exact sequence
associated with
$X \to Y \to Z$:
$$
\Omega^1_{Y/Z}\otimes_{\mathcal O_Y}
\mathcal O_Y/m
\xrightarrow{\psi}
\Omega^1_{X/Z}
\xrightarrow{\phi}
\Omega^1_{X/Y}\to 0.
$$
Vanishing of the relative differential
implies an isomorphism
$\Omega^1_{X/Z}=\Omega^1_{X/Y}$.
These spaces are Zariski cotangent spaces to
$Y$ and $Z$ because $X=\{x\}$ is a point.
\end{proof}

\begin{lemma}[Corollary]
\label{commutative-algebra-lemma-etale-smooth-points}
Let $f:X \to Y$ be an \'etale map. Then
$x\in X$ is smooth if and only if $f(x)$ is
smooth.
\end{lemma}

\begin{proof}
An \'etale map preserves the Krull dimension
(``by these pull back and forth ideas''),
and by the previous result also the 
dimension of $m/m^2$ is that number.
\end{proof}

\noindent
Grothendieck's idea behind the following was:
replace
$$
U_i\cap U_j \cap U_k
$$
on a Chech cover by
$$
\tilde{U}_i x \tilde{U}_j \times \tilde{U}_k.
$$
And suprisingly, in defining cohomology
just as in the usual Chech cohomology,
we obtain the same result.


\begin{exercise}
\label{commutative-algebra-exercise-derivations-finite-field-extension}
Let $K/k$ be a finite extension of fields.
Prove that $\operatorname{Der}_k(K)=0$.
\end{exercise}

\begin{definition}
\label{commutative-algebra-definition-support}
The {\it support} of an $R$-module $M$ is
$$
\text{supp}(M)
=
\{x \in \Spec R
:R \text{ localized at $x$ is not zero}\}.
$$
\end{definition}

\noindent
The ramification locus is then
the support of $\Omega_{X/Y}^1$.
More precisely:

\begin{definition}
\label{commutative-algebra-definition-ramification-locus}
Let $X \to Y$ be a finite map. Its
ramification locus is the set of maximal
ideals $m \in Y$ such that
$\Omega^1_{X/Y}$
localized at $m$ is non-zero.
\end{definition}

\noindent
The idea is to show that every
variety has a smooth point.
By Noether's lemma, we know that
there is a projection to $\mathbb{C}^n$.
Then we check that the ramification
locus is not everything.

The idea of the proof is to show
that the complement of the support
(supports are Zariski closed sets)
is not dense, a contradiction:

\begin{proposition}
\label{commutative-algebra-proposition-ramification-locus-not-dense}
The ramification locus is not Zariski dense.
\end{proposition}

\begin{proof}
The intersection of localizations of
$\mathcal O_Y$ is $\mathcal O_Y$. If the
ramification locus is Zariski dense, we
obtain a non-zero differential (and
therefore, non-zero derivation) of the
corresponding fraction fields. However, the
module $\operatorname{Der}_{k(X)}(k(Y))$
vanishes because
$[k(Y):k(X)]$ is a finite extension.
\end{proof}

\begin{lemma}[Corollary]
\label{commutative-algebra-lemma-finite-map-etale-outside-closed}
Any finite map is \'etale outside of a
Zariski closed set.
\end{lemma}

\begin{lemma}[Corollary]
\label{commutative-algebra-lemma-smooth-points-zariski-dense}
The set of smooth points on an algebraic
variety $X$ is Zariski dense.
\end{lemma}

\begin{proof}
The Noether normalization lemma is used to
construct a finite map $X \to \mathbf{C}^n$,
and the previous corollary then yields
smooth points on $X$.
\end{proof}


% --- End commutative-algebra.tex ---

% --- Begin lie-algebras.tex ---
\chapter{Lie algebras}
\label{chapter-lie-algebras}





\section{Introduction}
\label{lie-algebras-section-introduction}

\noindent
This introduction is more like a global
upshot.

\medskip\noindent
Semisimple
finite-dimensional Lie algebras can be
classified using root space decomposition.
We choose a Cartan subalgebra
$\mathfrak{h}\subset\mathfrak{g}$ (see
Definition
\ref{lie-algebras-definition-Cartan-subalgebra}) and
consider the generalized weight
decomposition with respect to the adjoint
representation (which is by convention
called the {\it root} space decomposition)
$$ \mathfrak{g}=\bigoplus_{\alpha \in
\mathfrak{h}^*}\mathfrak{g}_\alpha \qquad
\text{with }\mathfrak{g}_0=\mathfrak{h}, $$
where
$[\mathfrak{g}_\alpha,\mathfrak{g}_\beta]
\subset\mathfrak{g}_{\alpha+\beta}$.

``For arbitrary Lie algebras, Cartan
subalgebras can have unwieldy stucture, and
an unpredicatble action under the adjoint
representation. We see from the previous
(\cite[Lecture 12]{KLAL}) result, however,
that with the assumption of semisimplicity
the situation is much more transparent. We
can rewrite the root space decomposition as
$$
\mathfrak{g}=\mathfrak{h}
\oplus\bigoplus_{\alpha\in\Delta}\mathfrak{g}
_\alpha
$$ where $\Delta$ is the set of all
$\alpha\neq 0$ such that
$\mathfrak{g}_\alpha\neq 0$. We call
$\alpha \in \Delta$ a {\it root} of
$\mathfrak{g}$, and $\mathfrak{g}_\alpha$
the corresponding {\it root space}.''

In Lecture 17 we have: ``Previously, given
a semisimple Lie algebra $\mathfrak{g}$ we
constructed its associtated root system
$(V,\Delta)$. (The construction depends on
choosing a Cartan subalgebra, but by
Chevalley's theorem, the root systems
constructed from the same $\mathfrak{g}$
are isomorphic.) Next, given a root system
we'll construct a {\it Cartan matrix} $A$,
and from this we'll eventually see how to
reconstruct $\mathfrak{g}$.

We'll see that to every root system
corresponds a semisimple Lie algebra so
it's important to know all the root
systems. Last time we saw the four series
$A_r,B_r,C_r$ and $D_r$, and the three
exceptions $E_6,E_7$ and $E_8$. The
remaining two exceptions are $F_4$ and
$G_2$.''

Suppose $(V,\Delta)$ is a root system and
$f:V \to \mathbb{R}$ is a linear map such
that $f(\alpha)\neq 0$ for all $\alpha \in
\Delta$. Define \cite[Definition
17.1]{KLAL}
\begin{enumerate}
\item $\alpha \in \Delta$ is {\it positive}
if $f(\alpha)>0$
and  {\it negative} if $f(\alpha)<0$.
\item A positive root is {\it simple} if it
cannot be written
as the sum of two positive roots.
\item A {\it highest root} $\theta \in
\Delta$ is
a root where $f$ is maximal;
that is $f(\theta)\geq f(\alpha)$ for all
$\alpha \in \Delta$.
\end{enumerate}

Finally, on Lecture 21,
``Let $V$ be a finite dimensional real vector
space
with a positive definite bilinear form
$(\cdot,\cdot)$.
Let $a \in V$ be a nonzero vector and denote
by $r_a$ the orthogonal reflection relative
to $a$,
i.e. $r_a(a)=-a$, $r_a(v)=v$ if $(a,v)=0$,
$r_a(v)=v-2\frac{(a,v)}{(a,a)}a$.

Let $(V,\Delta)$ be a root system.
Let $W$ be the subgroup of $O(V)$ generated
by $r_\alpha, \alpha \in \Delta$.
The group $W$ is called the {\it Weyl group}
of $(V,\Delta)$
(and of the corresponding semisimple Lie
algebra $\mathfrak{g}$.)''

\section{Basic definitions and examples}
\label{lie-algebras-section-basic-defin-examp}

\begin{definition}
\label{lie-algebras-definition-algebra}
An {\it algebra} is a vector space over a
field $\mathbb{F}$ endowed with a
binary operation that is bilinear
$$
\begin{aligned}
a(\lambda b+\mu c)&=\lambda ab+\mu ac\\
(\lambda b + \mu c) a &= \lambda b a + \mu c
a
\end{aligned}
$$
\end{definition}

\begin{definition}
\label{lie-algebras-definition-Lie-algebra}
A {\it Lie algebra} is an algebra
$\mathfrak{g}$
whose product $[\cdot,\cdot]$ is called
{\it bracket} and satisfies
\begin{enumerate}
\item $[x,x]=0\qquad \forall x\in
\mathfrak{g}$,
\item (Jacobi identity.)
\end{enumerate}
\end{definition}

\begin{definition}
\label{lie-algebras-definition-simple-Lie-algebra}
A {\it simple Lie algebra} is a
Lie algebra that has no
nontrivial proper ideals.
\end{definition}

As a rule of thumb I keep in mind the
following silly computation:
$$
\text{log}\det A=\text{log}
\prod_{i}\lambda_i=\sum \text{log}\lambda_i
=``\text{tr}\text{log}A''
$$
And recall that exponent map goes from $T_eG=
\mathfrak{g} \to G$, so that
logarithm would go from $G\to \mathfrak{g}$.
This is why I remember that the
condition on a classical Lie group of
{\it having determinant 1} goes to {\it
having vanishing trace} in the
Lie algebra. (Because $\text{log}1=0$.)

\begin{example}
\label{lie-algebras-example-Lie-algebras}
\begin{enumerate}
\item The {\it special linear Lie algebra}
$$
\mathfrak{sl}_n=\{\text{Mat}_n|Tr(A)=0\}
$$
which is just obvious from the slogan above.
\item The {\it special orthogonal Lie
algebra}
$$
\mathfrak{so}_n=\{A\in
\text{Mat}_n|A+A^{\mathbf{T}}=0\}
$$
which is obvious from:
$\text{SO}(n)=$isometries, so
$\left\langle{}v,v\right\rangle{}=\left\langle{}Av,Av\right\rangle{}=\left\langle{}v,
A^{\mathbf{T}}Av\right\rangle{}$,
so
$\text{SO}(n)=\{A \in
\text{Mat}_n:A^{-1}=A^{\mathbf{T}}\}$, and
then
$0=
\text{log}1=\text{log}(A
A^{\mathbf{T}})=\text{log}A+\text{log}
A^{\mathbf{T}}$.

\item The {\it symplectic Lie algebra}
$$
\mathfrak{sp}_{2n}=\{A \in
\text{Mat}_{2n}:\Omega
A+A^{\mathbf{T}}\Omega=0\}
$$
where $\Omega=\begin{pmatrix}
0&\text{Id}_n\\
-\text{Id}_n & 0
\end{pmatrix}$.

{\bf Pending.} Why this makes sense?
\end{enumerate}
\end{example}

\begin{example}
\label{lie-algebras-example-assoc-algeb-gives-lie-algeb}
If $A$ is an associative Lie algebra,
then $A$ with the bracket  $[a,b]=ab-ba$
is a Lie algebra, denoted by  $A_-$.
(It is an exercise to verify Jacobi's
identity.) This gives
 $\mathfrak{gl}_V:=\text{End}(V)_-$.
\end{example}

\begin{definition}
\label{lie-algebras-definition-Lie-subalgebra}
A {\it Lie subalgebra} is a vector
subspace $\mathfrak{h} \subset \mathfrak{g}$
such that $[\mathfrak{h},\mathfrak{h}]
\subset \mathfrak{h}$.
\end{definition}

\begin{definition}
\label{lie-algebras-definition-ideal}
A subspace $\mathfrak{h}\subset \mathfrak{g}$
is called an
{\it ideal} if
$[\mathfrak{h},\mathfrak{g}]
\subset\mathfrak{h}$.
\end{definition}

\medskip\noindent
Recall that a bilinear form
$(\cdot,\cdot):\mathfrak{g}\times\mathfrak{g}
\to
\mathbb{C}$ is {\it invariant}
if $([a,b],c])=(a,[b,c])$ for all  $a,b,c \in
\mathfrak{g}$.

\begin{definition}
\label{lie-algebras-definition-ad}
For any finite dimensional Lie algebra
$\mathfrak{g}$
and $x \in \mathfrak{g}$
we write
$$
\text{ad}_a=[a,-]
$$
\end{definition}

\begin{exercise}
\label{lie-algebras-exercise-adjoint-representation}
Show that the map $\text{ad}:\mathfrak{g}\to
\text{End}(\mathfrak{g})$,
$a \mapsto \text{ad}_a$ is a representation,
i.e. a morphism of Lie algebras.
\end{exercise}

\begin{definition}
\label{lie-algebras-definition-Killing-form}
The {\it Killing form} is
$$
\kappa(x,y)=\text{Tr}_{\mathfrak{g}}\text{ad}
x\text{ad}y
$$
\end{definition}

\begin{exercise}
\label{lie-algebras-exercise-invar-bilin-form-simpl}
Prove that an invariant bilinear form on a
simple Lie algebra must in fact be
symmetric.
\end{exercise}

\begin{proof}
(David.) It's enough to show that
$\mathfrak{g}$ is
{\it perfect}, i.e. that
$[\mathfrak{g},\mathfrak{g}]=\mathfrak{g}$.
In this
case, let $a,b \in \mathfrak{g}$ and suppose
that $b=[x,y]$. Then
$$
\begin{aligned}
(a,b)&=(a,[x,y])=(a,-[y,x])=(-[a,y],x)=([y,a]
,x)\\
&=(y,[a,x])=(y,-[x,a])=(-[y,x],a)=([x,y],a)=
(b,a)
\end{aligned}
$$
To confirm that $\mathfrak{g}$ is perfect
just observe that
$[\mathfrak{g},\mathfrak{g}]$ is a nontrivial
ideal of $\mathfrak{g}$.
\end{proof}

\begin{definition}
\label{lie-algebras-definition-semisimple-Lie-algebra}
A {\it semisimple Lie algebra}
is a direct sum of simple ones
$$
\mathfrak{g}=\mathfrak{g}
_1\oplus\ldots\oplus\mathfrak{g}_s
$$
\end{definition}

\begin{theorem}[Cartan]
\label{lie-algebras-theorem-cartan-semisimple}
The finite dimensional Lie algebra
$\mathfrak{g}$ is semisimple if and only if
its Killing form is nondegenerate.
\end{theorem}

\noindent
There are other equivalent definitions
semisimple Lie algebras:

\begin{theorem}[Levi decomposition]
\label{lie-algebras-theorem-levi-decomposition}
\begin{reference}
\cite[Theorem 11.1]{KLAL}
\end{reference}
If $\mathfrak{g}$ is a finite dimensional Lie
algebra
over a field $\mathbb{F}$ of characteristic
0,
then there exists a semisimple subalgebra
$\mathfrak{s}\subset\mathfrak{g}$,
complementary
to the radical $R(\mathfrak{g})$,
such that $\mathfrak{g}=\mathfrak{s} \oplus
R(\mathfrak{g})$
as vector spaces.
\end{theorem}

Here the {\it radical} $R(\mathfrak{g})$
is a solvable ideal of maximal possible
dimension
(\cite[Definition 11.1]{KLAL}.)

It is equivalent that $\mathfrak{g}$ is
semisimple to

\begin{itemize}
\item

\item
\end{itemize}

\section{Semisimple, nilpotent and solvable
Lie algebras}
\label{lie-algebras-section-nilpo-solva-lie-algeb}

Let $\mathfrak{g}$ be a Lie algebra. Define
$$
\mathfrak{g}^1:=[\mathfrak{g},\mathfrak{g}],
\quad
\mathfrak{g}^2:=[\mathfrak{g},\mathfrak{g}^1]
,\quad
\ldots,
\mathfrak{g}^n:=[\mathfrak{g},\mathfrak{g}
^{n-1}]
$$
and
$$
\mathfrak{g}^{(1)}:=[\mathfrak{g},
\mathfrak{g}],\quad
\mathfrak{g}^{(2)}:=[\mathfrak{g}^{(1)},
\mathfrak{g}^{(1)}],
\ldots,
\mathfrak{g}^{(n)}:=[\mathfrak{g}^{(n-1)},
\mathfrak{g}^{(n-1)}].
$$

\begin{definition}
\label{lie-algebras-definition-nilpo-solva-lie-algeb}
A Lie algebra $\mathfrak{g}$ is
\begin{itemize}
\item {\it nilpotent} if there exists $n$
such that $\mathfrak{g}^n=0$,
\item {\it solvable} is there exists $n$ such
that $\mathfrak{g}^{(n)}=0$.
\end{itemize}
\end{definition}

\section{Lie's theorem}
\label{lie-algebras-section-Lie-theorem}

Lie's theorem says that solvable Lie algebras
over closed fields of
characterstic not zero have weights.

\begin{definition}
\label{lie-algebras-definition-weight-space}
Let $\mathfrak{h}$ be a Lie algebra,
$\pi:\mathfrak{h}\to\mathfrak{gl}_V$ a
representation of $\mathfrak{h}$
and $\lambda \in \mathfrak{h}^*$. The {\it
weight space} of $\mathfrak{h}$
attached to $\lambda$ is
$$
V_\lambda^{\mathfrak{h}}:=\{v \in
V|\pi(h)v=\lambda(h)v,\;
\forall h \in \mathfrak{h}\}.
$$
If $V_\lambda^{\mathfrak{h}}\neq 0$, we say
that $\lambda$ is a {\it weight} for
$\pi$.
\end{definition}

\begin{theorem}[Lie's theorem]
\label{lie-algebras-theorem-Lie}
Let $\mathfrak{g}$ be a solvable Lie algebra
and $\pi$ a representation of
$\mathfrak{g}$ on a finite dimensional vector
space $V \neq 0$, over an
algebraically closed field $\mathbb{F}$ of
characterstic 0.
Then there exists a weight $\lambda \in
\mathfrak{g}^*$ for $\pi$
(that is, $V_\lambda^\mathfrak{g} \neq
\{0\}$).
\end{theorem}

\begin{exercise}
\label{lie-algebras-exercise-corollaries-of-Lie-theorem}
Show the following two corollaries of Lie's
theorem:
\begin{itemize}
\item for all representations $\pi$ of a
solvable Lie algebra $\mathfrak{g}$ on
a finite dimensional vector space
$\mathcal{V}$ over an algebraically closed
field $\mathbb{F}$,
$\text{char}\mathbb{F}=0$,
there exists a basis for $V$ for which the
matrices
of $\pi(\mathfrak{g})$ are upper triangular;
\item a solvable subalgebra $\mathfrak{g}
\subset \mathfrak{gl}_V$,
(where $V$ is finite-dimensional over an
algebraically closed field
$\mathbb{F}$, $\text{char}\mathbb{F}=0$),
is contained in the subalgebra of upper
triangular matrices
over $\mathbb{F}$ for some basis of $V$.
\end{itemize}
\end{exercise}

\section{Generalized Eigenspaces and
Generalized Weight spaces}
\label{lie-algebras-section-gener-eigen-gener-weigh}

\begin{definition}
\label{lie-algebras-definition-generalized-eigenspace}
A {\it generalized eigenspace of $A \in
\text{End}(V)$
with eigenvalue $\lambda$} is
$$
V_\lambda=\{v \in V:(A-\lambda \text{Id})^N=0
\text{ for some positive integer }N\}.
$$
\end{definition}

It turns out that any linear operator on an
algebraically closed field gives a
decomposition into generalized eigenspaces
via Jordan canonical form:

\begin{proposition}
\label{lie-algebras-proposition-gener-eigen-decom}
Let $A$ be a linear operator on a
finite-dimensional
vector space $V$ over an algebraically closed
field $\mathbb{F}$, and let
$\lambda_1,\ldots,\lambda_s$
be all eigenvalues of $A$, and
$n_1,\ldots,n_s$ their
multiplicities.
Then one has the generalized eigenspace
decomposition:
$$
V=\bigoplus_{i=1}^s V_{\lambda_i},\qquad \dim
V_{\lambda_i}=n_i.
$$
\end{proposition}

In particular, for a Lie algebra
$\mathfrak{g}$ with a
representation
(i.e. Lie algebra morphism)
$\pi:\mathfrak{g} \to \text{End}(V)$ we have
$$
V=\bigoplus_{\lambda \in
\mathbb{F}}V^a_\lambda,\qquad
V^a_\lambda=\{v \in V:(\pi(a)v-\lambda
\text{Id})^Nv=0\text{ for some }
N \in \mathbb{N}\}
$$
And even more particularly, for the adjoint
representation we have
$$
V=\bigoplus_{\lambda \in
\mathbb{F}}\mathfrak{g}^a_\lambda,\qquad
\mathfrak{g}_\alpha^a=\{b \in
\mathfrak{g}:([a,\cdot]-\alpha
\text{Id})^Nb=0,
\text{ for some }N \in \mathbb{N}\}
$$
And a little more generally we have
\begin{definition}
\label{lie-algebras-definition-generalized-weight-space}
Let $\mathfrak{h}$ be a Lie algebra with a
representation $\pi$ on a vector
space $V$, and $\lambda \in \mathfrak{h}^*$.
A {\it generalized weight space of
$\mathfrak{h}$ in $V$ attached to $\lambda$}
is
$$
V_\lambda^\mathfrak{h}=
\left\{v \in V:
(\pi(a)-\lambda(a)\text{Id})^Nv=0,
\substack{ \text{ for some }N \in
\mathbb{N}\\
\text{ depending on }a \in \mathfrak{h}\\
\text{for all }a \in \mathfrak{h}}\right\}
$$
\end{definition}

Under the right conditions, a nilpotent
subalgebra
$\mathfrak{h} \subset \mathfrak{g}$ permits
decomposing $V$ as a direct sum of
generalized weight
spaces of $\mathfrak{h}$. Namely,

\begin{theorem}
\label{lie-algebras-theorem-decom-into-gener-weigh-spac}
Let $\mathfrak{g}$ be a finite-dimensional
Lie algebra
and $\pi$ its representation on a
finite-dimensional vector
space $V$, over an algebraically closed field
$\mathbb{F}$ of characteristic 0.
Let $\mathfrak{h}$ be a nilpotent subalgebra
of $\mathfrak{g}$.
Then the following equalities hold:
$$
V=\bigoplus_{\lambda \in
\mathfrak{h}^*}V_\lambda^\mathfrak{h}
$$
\begin{equation}
\label{lie-algebras-equation-inclu-gener-repre}
\pi\left(\mathfrak{g}^\mathfrak{h}
_\alpha\right)V_\lambda^\mathfrak{h}
\subseteq V_{\lambda+\alpha}^\mathfrak{h}
\end{equation}
\end{theorem}

Which in the case of the adjoint
representation, looks like:

\begin{definition}
\label{lie-algebras-definition-gener-root-space-decom}
The {\it generalized root space
decomposition} of a Lie algebra
$\mathfrak{g}$ over an algebraically closed
field $\mathbb{F}$
of characteristic zero with respect to a
nilpotent subalgebra $\mathfrak{h}$
is the generalized weight space decomposition
with respect to the
adjoint representation. That is,
$$
\mathfrak{g}=\bigoplus_{\alpha \in
\mathfrak{h}^*}
\mathfrak{g}_\alpha^\mathfrak{h}
$$
and it has the property that
\begin{equation}
\label{lie-algebras-equation-inclu-adjoi-repre}
[\mathfrak{g}_\alpha^\mathfrak{h},
\mathfrak{g}_\beta^\mathfrak{h}]
\subseteq\mathfrak{g}_{\alpha+\beta}
^\mathfrak{h}
\end{equation}
\end{definition}

It is important to make the distinction
between the generalized weight
space decomposition and the generalized root
space decomposition;
``we will see its convenience in later
lectures, as
we try to better understand the functionals
$\alpha$ appearing
in the decomposition''.

\begin{exercise}
\label{lie-algebras-exercise-gener-weigh-space-decom}
Take
$\mathfrak{g}=\mathfrak{gl}_n(\mathbb{F})$
and
$\mathfrak{h}=\{\text{diagonal matrices}\}$.
Find the generalized weight space
decomposition in both the
tautological and the adjoint representations,
and check the inclusions
\ref{lie-algebras-equation-inclu-gener-repre} and
\ref{lie-algebras-equation-inclu-adjoi-repre}.
\end{exercise}

\section{Zarisky Topology and Regular
Elements}
\label{lie-algebras-section-zaris-topol-regul-eleme}

Let $\mathfrak{g}$ be a finite-dimensional
Lie algebra
of dimension $d$. For every element  $a \in
\mathfrak{g}$
the characeristic polynomial
of $\text{ad}_a$ must be of the form
$$
\det_{\mathfrak{g}}(\text{ad}_a-
\lambda\text{Id})
=(-\lambda)^d+c_{d-1}(-\lambda)^{d-1}+\ldots
+\det(\text{ad}_a).
$$
[to be honest I don't see why the constant
term must be de determinant of $\text{ad}_a$,
but OK…]
since $[a,a]=0$, $\det(\text{ad}_a)=0$, i.e.,
the
constant term vanishes.

\begin{exercise}
\label{lie-algebras-exercise-chara-polyn-ada}
Show that $c_j$ is a homogeneous polynomial
on $\mathfrak{g}$
of degree $d-j$.
\end{exercise}

\begin{proof}
The term ``polynomial on $\mathfrak{g}$'' is
not clear.
But we follow \cite{KLAL} to interpret this
is only as
being a polynomial on the $a^i$ where
$a=a^iX_i$
for some basis $X_i$ of $\mathfrak{g}$.

Then we just notice that denoting
$[X_i,X_j]:=L_{ij}$ we have
$$
[a,X_j]=[a^iX_i,X_j]=a^i[X_i,X_j]=a^iL_{ij}
$$
which is a linear combination of the $a^i$.
Then the matrix
representation of $\text{ad}_a$ in terms of
the basis $X_i$ is
$$
[\text{ad}_a]=\begin{pmatrix}
a^iL_{i1}^1&\cdots &a^iL_{in}^1\\
\vdots &&\vdots \\
a^iL_{i 1}^n&  \cdots &  a^iL_{in}^n
\end{pmatrix}
$$
substracting $\lambda \text{Id}$ and taking
determinant
we obtain that a term with $\lambda^k$ would
have
for coefficient a product of $k$ of the
linear
combinations $a^iL_{i\ell}^j$ for varying $j$
and $\ell$.
Such a product is understood to be a
homogeneous polynomial
of degree $k$ in $\mathfrak{g}$.
\end{proof}

\begin{definition}
\label{lie-algebras-definition-rank}
I think this goes for any Lie algebra
$\mathfrak{g}$:
\begin{itemize}
\item The {\it rank} of $\mathfrak{g}$ is
the smallest integer $r$ such that $c_r(a)$
is not the
zero polynomial on $\mathfrak{g}$.
\item An element $a \in \mathfrak{g}$ is
called {\it regular}
if $c_r(a)\neq 0$.
\item The {\it discriminant} of
$\mathfrak{g}$ is
the nonzero polynomial $c_r(a)$ of degree
$d-r$, what?
\end{itemize}
\end{definition}

Explanation: we compute the polynomial
characteristic w.r.t.
$\text{ad}_a$ for every $a \in \mathfrak{g}$.
Express this polynomial as a polynomial with
coefficients in $\mathfrak{g}[t]$,
(here's the question: this polynomial seems
to me to be a polynomial
in $\mathbb{F}$.)
This polynomial does not have a constant
term.
But what is the next smallest-degree
monomial? 1? 2?
That's the rank of the Lie algebra.

\section{Cartan subalgebra}
\label{lie-algebras-section-Cartan-subalgebra}

Recall that $\mathfrak{h}\subset\mathfrak{g}$
is an ideal
if $[a,\mathfrak{h}]\subset\mathfrak{h}$ for
all $a \in \mathfrak{g}$.
Maybe $\mathfrak{h}$ is not an ideal, but we
can consider the largest subalgebra of
$\mathfrak{g}$
where $\mathfrak{h}$ is an ideal.
This is called the normalizer.

\begin{definition}
\label{lie-algebras-definition-normalizer}
Let $\mathfrak{h}$ be a subalgebra of a Lie
algebra $\mathfrak{g}$.
The {\it normalizer} of $\mathfrak{h}$
is
$N_\mathfrak{g}(\mathfrak{h}):=\{a \in
\mathfrak{g}
|[a,\mathfrak{h}]\subset\mathfrak{h}\}$
\end{definition}

I think behind the Cartan subalgebra is that,
roughly,
taking product with
anything not in $\mathfrak{h}$ leaves
$\mathfrak{h}$.

\begin{definition}
\label{lie-algebras-definition-Cartan-subalgebra}
A {\it Cartan subalgebra} of a Lie algebra
$\mathfrak{g}$ is a subalgebra
$\mathfrak{h}$, satisfying the following two
conditions:
\begin{enumerate}
\item $\mathfrak{h}$ is a nilpotent Lie
algebra,
\item
$N_\mathfrak{g}(\mathfrak{h})=\mathfrak{h}$.
\end{enumerate}
\end{definition}

\begin{proposition}
\label{lie-algebras-proposition-diago-matri-carta-subal}
Let $\mathfrak{g}\subset
\mathfrak{gl}_n(\mathbb{F})$ be a subalgebra
containing a diagonal matrix
$a=\text{diag}(a_1,\ldots,a_n)$
with distinct $a_i$, and let $\mathfrak{h}$
be the subspace
of all diagonal matrices in $\mathfrak{g}$.
Then $\mathfrak{h}$
is a Cartan subalgebra (of $\mathfrak{g}$ I
guess).
\end{proposition}

\begin{proof}[Idea of proof]
Consider the following illustrative example:
\begin{align*}
&\begin{pmatrix}
a&b\\
c&d
\end{pmatrix}\begin{pmatrix}
\lambda_1&0\\
0&\lambda_2
\end{pmatrix}-\begin{pmatrix}
\lambda_1&0\\
0&\lambda_2
\end{pmatrix}\begin{pmatrix}
a&b\\
c&d
\end{pmatrix}\\
&=\begin{pmatrix}
a \lambda_1&b \lambda_2\\
c \lambda_3&d \lambda_2
\end{pmatrix}-
\begin{pmatrix}
\lambda_1a&\lambda_1b\\
\lambda_2c&\lambda_2d
\end{pmatrix}\\
&=\begin{pmatrix}
0&b\lambda_2-\lambda_1b\\
\lambda_3a-\lambda_2c&0
\end{pmatrix}\qquad
\end{align*}
which says that any non-diagonal matrix would
escape $\mathfrak{h}$.
\end{proof}

\begin{theorem}[Cartan]
\label{lie-algebras-theorem-finit-dimen-lie-algeb-over}
Let $\mathfrak{g}$ be a finite-dimensional
Lie algebra over an
algebraically closed field $\mathbb{F}$. Let
$a \in \mathfrak{g}$
be a regular element (which exists since
$\mathbb{F}$ is infinite),
and let $\mathfrak{g}=\bigoplus_{\lambda \in
\mathbb{F}}\mathfrak{g}^a_\lambda$
be the generalized eigenspace decomposition
of $\mathfrak{g}$
with respect to $\text{ad}_a$. Then
$\mathfrak{g}_0^a$
is a Cartan subalgebra.
\end{theorem}

\begin{proof}[Idea of proof]
First recall what is
$\mathfrak{h}=\mathfrak{g}_0^a$, the set of
elements
$b\in\mathfrak{g}$ such that
$\text{ad}_a^N(b)=0$.
\end{proof}

\begin{remark}
\label{lie-algebras-remark-rank-dimen-subal}
The rank of a Lie algebra is the dimension of
$\mathfrak{g}_0^a=\mathfrak{h}$.
\end{remark}

\begin{proposition}
\label{lie-algebras-proposition-g0-equal-h-any-carta}
Let $\mathfrak{g}$ be a finite-dimensional
Lie algebra over
an algebraically closed field $\mathbb{F}$ of
characteristic zero
and let $\mathfrak{h}\subset\mathfrak{g}$ be
a Cartan subalgebra.
Then $\mathfrak{g}_0=\mathfrak{h}$ in the
generalized
weight space decomposition.
\end{proposition}

\begin{proof}[Idea of proof]
Engel's theorem…
\end{proof}



\section{Semisimple Lie algebras}
\label{lie-algebras-section-semisimple-Lie-algebras}

\noindent
Recall that the root space decomposition
works best for semisimple finite-dimensional
Lie algebras!

\begin{exercise}
\label{lie-algebras-exercise-semis-iff-killi-nonde}
\begin{reference}
\cite[Exercise 12.1]{KLAL}
\end{reference}
Let $\mathfrak{g}$ denote a Lie algebra over
a field of characteristic 0.
Show that $\mathfrak{g}$ is semisimple if and
only if
the Killing form on $\mathfrak{g}$ is
nondegenerate.
\end{exercise}

\begin{definition}
\label{lie-algebras-definition-radical}
A {\it radical} $R(\mathfrak{g})$ of a
finite-dimensional Lie algebra
$\mathfrak{g}$ is a solvable ideal of
$\mathfrak{g}$ of maximal possible
dimension.
\end{definition}

\begin{proposition}
\label{lie-algebras-proposition-radic-ideal-conta-any}
The radical ideal of $\mathfrak{g}$ contains
any solvable ideal of
$\mathfrak{g}$ and is unique.
\end{proposition}

If $\mathfrak{g}$ is a finite dimensional
solvable Lie algebra,
then $R(\mathfrak{g})=\mathfrak{g}$.
The opposite case if when
$R(\mathfrak{g})=0$.

\begin{definition}
\label{lie-algebras-definition-semisimple-Lie-algebra-duplicate-2}
A finite-dimensional Lie algebra
$\mathfrak{g}$ is called {\it semisimple} if
$R(\mathfrak{g})=0$.
\end{definition}

In \cite[Lecture 11]{KLAL} we have the tools
needed for a characterization
of semisimple Lie algebras in terms of the
Killing form:

\begin{theorem}[Cartan]
\label{lie-algebras-theorem-Cartan}
Let $\mathfrak{g}$ be a finite-dimensional
Lie algebra
over a field of characteristic 0.
Then the Killing form on $\mathfrak{g}$ is
non-degenerate is and only if
$\mathfrak{g}$ is semisimple.
Moreover, if $\mathfrak{g}$ is semisimple and
$\mathfrak{a}\subset\mathfrak{g}$
is an ideal,
then the restriction of the Killing form to
$\mathfrak{a}$,
$K|_{\mathfrak{a}\times\mathfrak{a}}$,
is also nondegenerate and coincides with the
Killing form of $\mathfrak{a}$.
\end{theorem}

I think we need semisimplicity, i.e.
nondegeneracy of the
Killing form for the following definition:

\begin{definition}
\label{lie-algebras-definition-coroot}
Given a set of simple roots
$\Pi=\{\alpha_1,\ldots,\alpha_\ell\}$,
we define the {\it coroot}
\begin{equation}
\label{lie-algebras-equation-coroot}
\alpha_i^\vee=\frac{2}{(\alpha_i,\alpha_i)}
\nu^{-1}(\alpha_i)
\end{equation}
where $\nu:\mathfrak{h}\to\mathfrak{h}^*$ is
the pairing
induced by the Killing form.
\end{definition}

\section{Borel subalgebras}
\label{lie-algebras-section-Borel-subalgebras}

\noindent
AI section!

\medskip\noindent
Let $\mathfrak{g}$ be a complex semisimple
Lie algebra and let $\mathfrak{h}\subset
\mathfrak{g}$ be a Cartan subalgebra. After
choosing a system of positive roots
$\Delta_+\subset \Delta$, we define
$$
\mathfrak{n}_+
=
\bigoplus_{\alpha\in\Delta_+}\mathfrak{g}_\alpha,
\qquad
\mathfrak{n}_-
=
\bigoplus_{\alpha\in\Delta_+}\mathfrak{g}_{-\alpha}.
$$
Thus
$$
\mathfrak{g}
=
\mathfrak{n}_-
\oplus
\mathfrak{h}
\oplus
\mathfrak{n}_+.
$$

\begin{definition}
\label{lie-algebras-definition-Borel-subalgebra}
A {\it Borel subalgebra} of $\mathfrak{g}$ is
a maximal solvable Lie subalgebra.
\end{definition}

The Borel subalgebra associated to the choice
of positive roots is
$$
\mathfrak{b}
=
\mathfrak{h}\oplus \mathfrak{n}_+.
$$

\begin{remark}
The choice of $\Delta_+$ is exactly the
choice of what we want to call ``raising''
operators. This is why highest weight vectors
are killed by $\mathfrak{n}_+$.
\end{remark}

\begin{proposition}
\label{lie-algebras-proposition-standard-Borel}
The Lie algebra
$\mathfrak{b}=\mathfrak{h}\oplus\mathfrak{n}_+$
is solvable. In fact it is a Borel
subalgebra.
\end{proposition}

\begin{proof}[Idea]
Since
$$
[\mathfrak{h},\mathfrak{h}]=0,
\qquad
[\mathfrak{h},\mathfrak{n}_+]\subset
\mathfrak{n}_+,
\qquad
[\mathfrak{n}_+,\mathfrak{n}_+]\subset
\mathfrak{n}_+,
$$
we have
$$
[\mathfrak{b},\mathfrak{b}]
\subset
\mathfrak{n}_+.
$$
Moreover $\mathfrak{n}_+$ is nilpotent,
because brackets of positive root spaces move
to higher positive roots, and this process
must stop. Therefore $\mathfrak{b}$ is
solvable.

The nontrivial fact is maximality: no strictly
larger solvable subalgebra contains
$\mathfrak{b}$. Thus $\mathfrak{b}$ is a
Borel subalgebra.
\end{proof}

\begin{remark}
For an algebraic group $G$, a Borel subgroup
is a maximal connected solvable closed
subgroup. Its Lie algebra is a Borel
subalgebra.
\end{remark}
\section{g-modules}
\label{lie-algebras-section-g-modules}

\begin{definition}
\label{lie-algebras-definition-g-module}
Let $\mathfrak{g}$ be a Lie algebra. A {\it
$\mathfrak{g}$-module}
(also called a {\it representation of
$\mathfrak{g}$}
is a vector space $V$ and a homomorphism
$$
\rho:\mathfrak{g}\to \text{End}(V)
$$
of Lie algebras, i.e.,
$$
\rho([x,y])=\rho(x)\rho(y)-\rho(y)\rho(x).
$$
A vector space $W \subset V$ such that
$vW\subset W$ for all
$x \in \mathfrak{g}$ is called a
 {\it $\mathfrak{g}$-submodule}.
\end{definition}

\begin{definition}
\label{lie-algebras-definition-irreducible-g-module}
An {\it irreducible $\mathfrak{g}$-module} is
one that has
no $\mathfrak{g}$-submodules other than the
trivial ones (namely, 0 and $V$ itself).
\end{definition}

It is possible to classify and describe the
irreducible finite-dimensional
$\mathfrak{g}$-modules. This is in contrast
with all the modules of a Lie algebra
$\mathfrak{g}$,
which is ``impossible in general''.

\section{Verma module}
\label{lie-algebras-section-Verma-module}

Let $M$ be a $\mathfrak{g}$-module. Recall
Definition \ref{lie-algebras-definition-weight-space}:
let $\mathfrak{h}$ be a Lie algebra,
$\pi:\mathfrak{h}\to\mathfrak{gl}_V$ a
representation of $\mathfrak{h}$
and $\lambda \in \mathfrak{h}^*$. The {\it
weight space} of $\mathfrak{h}$
attached to $\lambda$ is
$$
V_\lambda^{\mathfrak{h}}:=\{v \in
V|\pi(h)v=\lambda(h)v,\;
\forall h \in \mathfrak{h}\}.
$$
We have
that a weight space in $M$ is a nonzero
subspace
$M_\lambda \subset M$ such that
$$
hv=\lambda(h)v \qquad \forall h \in
\mathfrak{h}\subset\mathfrak{g}
\text{ and }v \in M,
$$
and in this case we call $\lambda$ a weight.
If $M$ is a direct sum of weight spaces, we
call it a {\it weight module}.

The {\it support} of a weight module $M$ is
\begin{equation}
\label{lie-algebras-equation-support-of-weight-module}
\text{supp}(M):=\{\lambda \in
\mathfrak{h}^*:M_\lambda\neq 0\}
\end{equation}
We call elements of $M_\lambda$ {\it vectors
of weight $\lambda$}.

\begin{definition}
\label{lie-algebras-definition-highest-weight}
$\Lambda \in \mathfrak{h}^*$ is a
{\it highest weight} of $M$
if $\Lambda \in \text{supp}(M)$
and for all positive roots $\alpha \in
\Delta_+$
we have that $\Lambda+\alpha
\not\in\text{supp}(M)$.
\end{definition}

\begin{exercise}
\label{lie-algebras-exercise-highe-weigh-irred-g-modul}
The highest weight of an irreducible
$\mathfrak{g}$-module
is unique.
\end{exercise}

\begin{proof}

\end{proof}

\begin{definition}
\label{lie-algebras-definition-Verma-module}
\begin{reference}
\cite[Definition 4.1]{jethro-rt}
\end{reference}
Let $\Lambda \in \mathfrak{h}^*$
(a highest weight, I suppose).
The {\it Verma module} $M(\Lambda)$
of highest weight $\Lambda$ is
$$
M(\Lambda)=U(\mathfrak{g})
\otimes_{U(\mathfrak{h}\oplus
\mathfrak{n}_+}
\mathbb{C}_\Lambda,
$$
where
$\mathbb{C}_\Lambda=\mathbb{C}v_\Lambda$
is the $1$-dimensional $(\mathfrak{h} \oplus
\mathfrak{n}_+)$-module
defined by $\mathfrak{n}_+v_\Lambda=0$
and
$hv_\Lambda=\left\langle{}\Lambda,h\right\rangle{}v_\Lambda$
for all $h \in \mathfrak{h}$.
\end{definition}

Long story short, by the PBW Theorem we
have an isomorphism $U(\mathfrak{n}_-) \to
M(\Lambda)$
of $U(\mathfrak{n}_-)$ modules.

\section{Weyl character formula}
\label{lie-algebras-section-Weyl-character-formula}

There are at least two Weyl character
formulas.
This one is \cite[Theorem 10.14]{Hall}:

\begin{theorem}[Weyl Character Formula]
\label{lie-algebras-theorem-Weyl-character-formula-Hall}
If $(\pi,V_\mu)$ is an irreducible
representation of
$\mathfrak{g}$ with highest weight $\mu$,
then
\begin{equation}
\label{lie-algebras-equation-Weyl-character-formula-Hall}
\chi_\pi(H)=\frac{\sum_{w \in
W}\det(w)e^{\left\langle{}w\cdot(\mu+\delta),
H\right\rangle{}}}
{\sum_{w \in W}\det(w)e^{\left\langle{}w\cdot
\delta,H\right\rangle{}}}
\end{equation}
for all $H \in \mathfrak{h}$ for which the
denominator is nonzero.
\end{theorem}

And this one is \cite[Theorem 25.3]{KLAL}:

\begin{theorem}[Weyl Character Formula]
\label{lie-algebras-theorem-Weyl-character-formula-KLAL}
Let $R=\prod_{\alpha \in
\Delta_+}(1-e^{-\alpha}$.
If  $\Lambda \in P_+$, then
\begin{equation}
\label{lie-algebras-equation-Weyl-character-formula-KLAL}
e^{\rho}R \text{ch}L(\Lambda)=\sum_{w \in
W}\det w e^{w(\Lambda+\rho)}
\end{equation}
\end{theorem}

\begin{exercise}
\label{lie-algebras-exercise-weyl-chara-formu-sl3}
Let $\mathfrak{g}=\mathfrak{sl}_3$. Compute
$\text{ch}_{L(\lambda)}$
for some $\lambda$,
\begin{enumerate}
\item $\lambda=0$,
\item $\lambda=\omega_1$,
\item $\lambda=\omega_1+\omega_2$,
\end{enumerate}
etc.
\end{exercise}

\begin{proof}
No matter what formula I use, I must compute
the determinants of the Weyl reflections
and the $\rho$ (which I think is $\delta$ for
Hall),
which is just the sum of all the fundamental
weights,
which in turn are the dual of the coroot
vectors.
Finally applying the reflections to $\rho$
and
to $\lambda+\rho$ would yield the result.
\begin{enumerate}
\item (Find roots of $\mathfrak{sl}_3$.) We
look for $\alpha \in \mathfrak{h}^*
= \{\text{diagonal matrices}\}$ such that
$$
[H,E_{ij}]=\alpha(H)E_{ij}\qquad \forall H
\in \mathfrak{h}
$$
where $E_{ij}$ is the matrix that has zero in
every entry but
in the $(i,j)$-th where it has a 1. One
obtains that
$[H,E_{ij}]=(h_i-h_j)E_{ij}$, so that the
roots are
$\alpha_{ij}(h)=h_i-h_j$.

\item (Compute the Killing form.) Recall that
by definition
$\kappa(H,H')=\text{Tr}(\text{ad}_H\text{ad}
_{H'})$.
But in this case, we have
$\kappa(H,H')=\text{Tr}(HH')$
because
$\text{ad}_H\text{ad}_{H'}=[H,[H',A]]$

The first thing Jethro did was to consider a
basis
$$
H_1=\begin{pmatrix}
1&0&0\\
0&-1&0\\
0&0&0
\end{pmatrix},
\qquad
H_2=\begin{pmatrix}
0&0&0\\
0&1&0\\
0&0&-1
\end{pmatrix}
$$
and then compute the Killing form via
$$
\begin{aligned}
H_1H_1&=2,\qquad H_1H_2=-1,\qquad
H_2H_1=-1,\qquad H_2H_2=2.
\end{aligned}
$$

\item (Find the dual basis vectors.) Now we
use the musical
isomorphism $\nu:\mathfrak{h}\to
\mathfrak{h}^*$,
$H\mapsto (-,H)$ to find
$$
\nu(H_1)=
$$
\item (Find fundamental weights.) It looks
like $\alpha_1=2 \omega_1-\omega_2$
and $\alpha_2=2\omega_2-\omega_1$. This says
$\omega_1=2\omega_2-\alpha_2=2(2\omega_1-
\alpha_1)-\alpha_2$
so
$3\omega_1=2\alpha_1+\alpha_2$. And then
$\omega_2=2(\frac{2}{3}\alpha_1+\frac{1}{3}
\alpha_2)-\alpha_1
=\frac{1}{3}(\alpha_1+2\alpha_2)$.

\item (Find $\rho$.) Then
$\rho=\omega_1+\omega_2=\alpha_1+\alpha_2$.

\item (Find Weyl reflections.) To find Weyl
reflections I need to compute
first the numbers $(\alpha_i, \alpha_j)$. I
think this is still not
very clear --- have to use Killing form? But
OK, I know we must have
$(\alpha_i,\alpha_i)=2$, and the both of the
crossed ones give $-1$.
This should give by bilinearity the values on
$\omega_i$ also.

\item (Find determinants of Weyl
reflections.) I guess for this
I have to evaluate the reflections in each of
the w

\item (Evaluate $\rho$ and $\lambda+\rho$ on
Weyl reflections.)

\item (Find the Weyl character.)
\end{enumerate}

What's funny in this whole story is that
in the end we can just find the Weyl
character geometrically.
For every $\lambda=m \omega_1+n \omega_2$
I want to compute its orbit under the Weyl
group $W$.
This means


I would start by finding the determinant
of the Weyl reflections…
\end{proof}


% --- End lie-algebras.tex ---

% --- Begin representation-theory.tex ---
\chapter{Representation Theory}
\label{chapter-representation-theory}





\section{Introduction}
\label{representation-theory-section-introduction}

\noindent
Basic roadmap 
of what I should know
and currently don't know,
suggested by AI:

\begin{enumerate}
\item Representations as $k[G]$-modules.

\item Subrepresentations, quotient
representations, and direct sums.

\item Irreducible representations, simple
modules, and Schur's lemma.

\item Semisimple representations and
Maschke's theorem.

\item Characters:
$$
\chi_V(g)=\operatorname{tr}(\rho(g)).
$$

\item Orthogonality of characters and the
fact that irreducible characters form a
basis of class functions.

\item Regular representation:
$$
\mathbb{C}[G]
\cong
\bigoplus_i V_i^{\oplus \dim V_i}.
$$

\item The formula
$$
|G|=\sum_i(\dim V_i)^2.
$$

\item Character tables of basic examples:
cyclic groups, finite abelian groups,
$S_3$, dihedral groups, and $Q_8$.

\item Restriction, induction, and Frobenius
reciprocity.

\item Tensor products of representations and
the tensor-category viewpoint.

\item Modular and infinite-group warning:
Maschke can fail, and non-semisimple
phenomena such as Jordan blocks can appear.
\end{enumerate}

\section{Representations of Lie algebras}
\label{representation-theory-section-lie-algebra-reps}

\begin{definition}
\label{representation-theory-definition-representation-lie-algebra}
A {\it representation of Lie algebra}
$\mathfrak{g}$ is a vector space $V$
and homomorphism
$$
\rho:\mathfrak{g} \to \text{End}(V)
$$
of Lie algebras, i.e.,
$$
\rho([x,y])=\rho(x)\rho(y)-\rho(y)\rho(x).
$$
\end{definition}

\begin{remark}
\label{representation-theory-remark-representations-are-modules}
Another name for a representation of a Lie
algebra $\mathfrak{g}$ is a
{\it $\mathfrak{g}$-module}. This is like the
abstract definition of an algebra
using a morphism. It's as if the morphism
lets us define a sort of product of
elements in the codomain by elements in the
domain.
\end{remark}

It is not possible to classify all
representations of a Lie algebras
$\mathfrak{g}$. But there is a theorem by
Weyl that says that if $\mathfrak{g}$
is finite-dimensional and (semi)simple, then
every finite-dimensional
representation of $\mathfrak{g}$
(i.e. a representation where $V$ is
finite-dimensional)
is isomorphic to a direct sum of irreducible
representations

\section{Representations of groups}
\label{representation-theory-section-representations-of-groups}

\noindent
Notes from
\href{http://verbit.ru/IMPA/CA-2026/slides-%
CA-2026-05.pdf}{
Misha's lecture}.

\begin{definition}
\label{representation-theory-definition-representation-group}
A {\it representation of a group}
is a vector space $V$ and a morphism
$\rho:G \to \text{GL}(V)$.
A {\it morphism} of $G$-representations
is a map of the underlying vector spaces
commuting
with the $G$-action, i.e.
$\varphi(gx)=g\varphi(x)$.
\end{definition}

\begin{definition}
\label{representation-theory-definition-irreducible-representation}
A group representation $G \to \text{GL}(V)$
is {\it irreducible} if $V$ has no
$G$-invariant subspaces,
i.e. if there are no subrepresentations.
It is called {\it semisimple} if it is a
direct sum
of irreducible representations.
\end{definition}

\begin{remark}
\label{representation-theory-remark-semis-iff-any-short-exact}
All representations are semismple if and only
if
any short exact sequence of
$G$-representations splits.
\end{remark}

\begin{proposition}
\label{representation-theory-proposition-finit-group-have-semis}
Let $\text{Rep}_k(G)$ be the category of
representations
of a finite group $G$ over $k$,
with $\text{char}(k)$ coprime with $|G|$.
Then any short exact sequence of
$G$-representations splits.
\end{proposition}

\begin{proof}
The idea is to construct a section of the
right-hand map
on a short exact sequence by averaging over
the group.
We need $\text{char}(k)$ coprime with $|G|$
so that
we can divide by $|G|$.
\end{proof}

This means that any finite-dimensional
representation
of finite group with $\text{char}(k)$ coprime
with $|G|$
is semisimple.

\begin{exercise}
\label{representation-theory-exercise-nonsemisimple-borel}
Consider the group $G$ of non-degenerate
upper triangular matrices $n \times n$.
Construct a non-semisimple representation of
$G$.
\end{exercise}

\begin{proof}
Assume $n\geq 2$ and let $G$ act on
$V=k^n$ by the standard representation.
Let $e_1,\ldots,e_n$ be the standard basis.
The line
$$
U=\langle e_1\rangle
$$
is $G$-invariant, since every upper
triangular matrix sends $e_1$ to a scalar
multiple of $e_1$.

We claim that $U$ has no $G$-invariant
complement. Suppose that such a complement
$W$ exists. Since $V=U\oplus W$, there is a
unique vector $w\in W$ whose image in
$V/U$ is the image of $e_2$. Thus
$$
w=e_2+ce_1
$$
for some $c\in k$.

Now consider the unipotent upper triangular
matrix
$$
u_t=
\begin{pmatrix}
1&t&0&\cdots&0\\
0&1&0&\cdots&0\\
0&0&1&\cdots&0\\
\vdots&\vdots&\vdots&\ddots&\vdots\\
0&0&0&\cdots&1
\end{pmatrix}.
$$
Since $W$ is assumed invariant, $u_t w$ must
belong to $W$. But
$$
u_t w=e_2+(c+t)e_1.
$$
This vector has the same image as $e_2$ in
$V/U$. By uniqueness of $w$, we must have
$u_t w=w$, hence $c+t=c$. Taking
$t\neq 0$ gives a contradiction.

Therefore the invariant subspace $U$ has no
$G$-invariant complement. Hence the standard
representation of $G$ on $k^n$ is not
semisimple.
\end{proof}

\begin{exercise}
\label{representation-theory-exercise-finite-field-nonsemisimple}
Construct a finite group $G$ and a
non-semisimple representation
$$
G \longrightarrow \text{GL}(V),
$$
where $V$ is a finite-dimensional vector
space over a finite field.
\end{exercise}

\begin{proof}
We may use the above example for any
nontrivial finite field - the group
of upper triangular matrices is finite
and we get a non-semisimple representation
in the exact same way.
\end{proof}

\begin{definition}
\label{representation-theory-definition-real-complex-quaternionic}
Let $V$ be a finite-dimensional irreducible
group representation over $\mathbb{R}$.
It is called {\it real} if
$\text{End}_G(V)=\mathbb{R}$, {\it complex}
if $\text{End}_G(V)=\mathbb{C}$, and
{\it quaternionic} if $\text{End}_G(V)$ is
the algebra of quaternions.
\end{definition}

\begin{remark}
\label{representation-theory-remark-real-complex-quaternionic}
This definition is not about the image of
the homomorphism
$G \to \text{GL}(V)$. It is about the
endomorphisms of $V$ which commute with the
whole $G$-action:
$$
\text{End}_G(V)=
\{T\in \text{End}_{\mathbb{R}}(V):
T\rho(g)=\rho(g)T\text{ for all }g\in G\}.
$$
Thus $\text{End}_G(V)$ is an algebra of
operators on $V$. For an irreducible real
representation, Schur's lemma says that this
is a division algebra over $\mathbb{R}$.
By Frobenius' theorem, it is isomorphic to
one of
$$
\mathbb{R},\qquad \mathbb{C},\qquad
\mathbb{H}.
$$
These are the real, complex, and
quaternionic cases.

In other words, $\text{End}_G(V)$ asks:
after $G$ acts, what linear operations on
$V$ are still invisible to $G$, because they
commute with everything $G$ does? If these
invisible operations form $\mathbb{H}$, the
representation is quaternionic.
\end{remark}

\begin{exercise}
\label{representation-theory-exercise-quaternionic-finite-group}
Find an example of a quaternionic
representation of a finite group.
\end{exercise}

\begin{proof}
Let $G=Q_8=\{\pm 1,\pm i,\pm j,\pm k\}$ be
the quaternion group, and let $V=\mathbb{H}$
be regarded as a real vector space. Let
$G$ act by left multiplication:
$$
\rho(q)(x)=qx.
$$
This is a real representation of $G$.

We first check that it is irreducible. Let
$0\neq W\subset V$ be a $G$-invariant real
subspace. Since $W$ is stable under left
multiplication by $i,j,k$, it is stable
under left multiplication by every
quaternion.
Choose $0\neq x\in W$. Since $x$ is
invertible in $\mathbb{H}$, left
multiplication by $x^{-1}$ shows that
$1\in W$. Hence $W=\mathbb{H}$.

Now compute the commuting endomorphisms.
For every $a\in \mathbb{H}$, right
multiplication
$$
R_a(x)=xa
$$
commutes with left multiplication by $G$.
Conversely, let $T\in \text{End}_G(V)$ and
put $a=T(1)$. Since $T$ commutes with left
multiplication by $i,j,k$, it commutes with
left multiplication by every quaternion.
Thus, for all $x\in\mathbb{H}$,
$$
T(x)=T(x\cdot 1)=xT(1)=xa.
$$
So every $G$-equivariant endomorphism is
some $R_a$. Therefore
$$
\text{End}_G(V)\cong \mathbb{H}^{op}
\cong \mathbb{H}.
$$
Thus $V$ is quaternionic.
\end{proof}

\begin{exercise}
\label{representation-theory-exercise-quaternionic-cyclic-17}
Find a quaternionic representation of a
cyclic group of order $17$, or prove that
such a representation cannot exist.
\end{exercise}

\begin{exercise}
\label{representation-theory-exercise-quaternionic-tetrahedron}
Find a quaternionic representation of the
group $G$ of isometries of a regular
tetrahedron, or prove that such a
representation cannot exist.
\end{exercise}

\begin{proof}
The full group of isometries of a regular
tetrahedron is isomorphic to $S_4$. We show
that it has no quaternionic irreducible real
representation.

We use the Frobenius-Schur indicator. If
$\chi$ is an irreducible complex character
of a finite group, then
$$
\nu(\chi)=
\frac{1}{|G|}\sum_{g\in G}\chi(g^2)
$$
is $1,0,-1$ according as the corresponding
real representation is of real, complex, or
quaternionic type.

Let the conjugacy classes of $S_4$ be
$$
C_1,\ C_2,\ C_{22},\ C_3,\ C_4
$$
with cycle types
$$
1,\quad (12),\quad (12)(34),\quad
(123),\quad (1234).
$$
Their sizes are $1,6,3,8,6$. Squaring sends
$C_1,C_2,C_{22}$ to $C_1$, sends $C_3$ to
$C_3$, and sends $C_4$ to $C_{22}$. Hence
$$
\nu(\chi)=
\frac{10\chi(C_1)+8\chi(C_3)
+6\chi(C_{22})}{24}.
$$

The irreducible characters of $S_4$ are
$$
\begin{array}{c|ccccc}
&C_1&C_2&C_{22}&C_3&C_4\\
\hline
\chi_1&1&1&1&1&1\\
\chi_{\mathrm{sgn}}&1&-1&1&1&-1\\
\chi_2&2&0&2&-1&0\\
\chi_3&3&1&-1&0&-1\\
\chi_3'&3&-1&-1&0&1
\end{array}
$$
Substituting these rows into the formula for
$\nu(\chi)$ gives $\nu(\chi)=1$ in every
case. Thus every irreducible real
representation of $S_4$ is of real type.
In particular, none is quaternionic.
\end{proof}

\begin{exercise}
\label{representation-theory-exercise-cyclic-rational-2d}
Let $V=\mathbb{Q}^2$ be an irreducible
$2$-dimensional representation of the cyclic
group $G=\mathbb{Z}/n$ over $\mathbb{Q}$.
\begin{enumerate}
\item Prove that the algebra
$\text{End}_G(V)$ of endomorphisms of this
representation is commutative and
$2$-dimensional.
\item Prove that $\text{End}_G(V)$ is a
number field isomorphic to
$\mathbb{Q}[\sqrt{a}]$ for some
$a \in \mathbb{Z}$.
\end{enumerate}
\end{exercise}

\begin{proof}
Let $g$ be a generator of $G$, and write
$$
A=\rho(g)\in \text{GL}_2(\mathbb{Q}).
$$
Geometrically, one may picture this as a
finite order rotation. Then an endomorphism
commuting with the representation should be
something which has the same symmetry: a
homothety, or the rotation itself, or a
linear combination of these two operations.
Over $\mathbb{Q}$ we should say this
algebraically rather than in terms of angles.

Since $G$ is cyclic, an endomorphism
$T\in\text{End}_{\mathbb{Q}}(V)$ is
$G$-equivariant if and only if it commutes
with $A$:
$$
\text{End}_G(V)=
\{T\in\text{End}_{\mathbb{Q}}(V):TA=AT\}.
$$
Thus the problem is to compute the
centralizer of the single matrix $A$.

The scalar matrices give the homotheties,
and $A$ itself is the algebraic replacement
for the rotation. Since the representation is
irreducible, $A$ cannot be scalar; otherwise
every line in $V$ would be $G$-invariant.
Therefore $I$ and $A$ are linearly
independent over $\mathbb{Q}$.

By Cayley-Hamilton, every polynomial in $A$
is of the form
$$
aI+bA,\qquad a,b\in\mathbb{Q}.
$$
So the algebra $\mathbb{Q}[A]$ has dimension
$2$ and is commutative. The remaining point
is to show that every matrix commuting with
$A$ already lies in $\mathbb{Q}[A]$.
\end{proof}


\section{Weight decomposition for
algebraic tori}
\label{representation-theory-section-torus-weights}

\noindent
Let $T=(\mathbb{C}^{\times})^n$.
A finite-dimensional representation
$\rho:T\to\operatorname{GL}(V)$ is
algebraic if its matrix entries are
regular functions on $T$, that is,
Laurent polynomials in $t_1,\ldots,t_n$.
This condition matters: we are not
considering arbitrary representations
of the underlying abstract group.

For characters and their eigenspaces,
see \cite[Section 4g and
Chapter 12, Sections a--d]{milne-ag-2022}.
The algebraic characters of $T$ are
precisely the homomorphisms
$$
\chi_\alpha:T\longrightarrow
\mathbb{C}^{\times},
\qquad
\chi_\alpha(t)=t^\alpha
=\prod_{i=1}^n t_i^{\alpha_i},
\qquad \alpha\in\mathbb{Z}^n.
$$
Negative exponents are allowed because
each $t_i$ is nonzero.

\begin{definition}
\label{representation-theory-definition-torus-weight-space}
The \emph{weight space} of weight
$\alpha$ in $V$ is
$$
V_\alpha=
\{v\in V:\rho(t)v=t^\alpha v
\ \forall t\in T\}.
$$
\end{definition}

\noindent
Thus its vectors are simultaneous
eigenvectors for every torus element,
with eigenvalues specified by $\alpha$.
The definition alone does not assert
that these spaces span $V$.

\begin{theorem}
\label{representation-theory-theorem-torus-weights}
\begin{reference}
\cite[Theorem 12.12]{milne-ag-2022}.
\end{reference}
Every finite-dimensional algebraic
representation of $T$ decomposes as
\begin{equation}
\label{representation-theory-equation-torus-weights}
V=\bigoplus_{\alpha\in\mathbb{Z}^n}
V_\alpha.
\end{equation}
Only finitely many summands are nonzero.
\end{theorem}

\noindent
This is the split-torus case of
Milne's \cite[Theorem 12.12,
pp.~234--235]{milne-ag-2022}.
We accept it for now; the proof is
available there and is postponed here.
The same decomposition holds for a
locally finite algebraic representation:
every vector belongs to a
finite-dimensional invariant subspace
on which the action is algebraic.
There may then be infinitely many
weights, but each vector has only
finitely many nonzero weight components.
See \cite[Corollary 4.8 and
Theorem 12.12]{milne-ag-2022}.

The torus is not a semisimple group.
Its algebraic representations are
semisimple, with one-dimensional
irreducible representations given by
characters. Commutativity alone does
not imply diagonalizability.

To connect with Lie algebra weights,
recall Definition~\ref
{lie-algebras-definition-weight-space}.
The Lie algebra of $T$ is the abelian
algebra $\mathfrak{t}=\mathbb{C}^n$.
Let $\pi$ be the differentiated action.
For $h=(h_1,\ldots,h_n)$, put
$$
t(s)=(e^{s h_1},\ldots,e^{s h_n}).
$$
For $v\in V_\alpha$ we have
$$
\rho(t(s))v
=e^{s\sum_i\alpha_i h_i}v.
$$
Differentiating at $s=0$ gives
$$
\pi(h)v=\lambda_\alpha(h)v,
\qquad
\lambda_\alpha(h)=\sum_i\alpha_i h_i.
$$
These are the familiar Lie algebra
weight spaces. Conversely, this
infinitesimal condition recovers the
torus condition by exponentiation,
since every element of $T$ is an
exponential. In this setting the
spaces are genuine eigenspaces,
not merely generalized weight spaces.

For example, the action on polynomials
given by $x_i\mapsto t_i x_i$ makes
each monomial a weight vector:
$$
t\cdot x_1^{a_1}\cdots x_n^{a_n}
=t_1^{a_1}\cdots t_n^{a_n}
x_1^{a_1}\cdots x_n^{a_n}.
$$
Its weight is $(a_1,\ldots,a_n)$.
The application to cotangent cohomology
is developed in
Section~\ref
{deformations-section-torus-action-ac-multidegrees}
of \emph{Deformation Theory}.
The existence of the decomposition is
standard representation theory.
The contribution of
Altmann--Christophersen \cite{jan3}
is the combinatorial computation of
the cotangent cohomology pieces from
the simplicial complex.

\section{Group algebra}
\label{representation-theory-section-group-algebra}

\noindent
The group algebra is a gadget
that allows us to see group representations
as modules over the some ring --- the
group algebra. This way we can use what
we know about modules over commutative
rings in the realm of group representations
(the problem is that there is no such thing
as a $G$-module in the same
way as an $R$-module).

This is based in \cite[Section
3.4]{fulton-harris-rep}.

\begin{definition}
\label{representation-theory-definition-group-algebra}
Let $G$ be a finite group.
The {\it group algebra} of $G$
is the vector space $\mathbb{C}G$
generated by the vectors $\{e_g\}$
for $g \in G$ equipped with the product
$$
e_g\cdot e_h = e_{gh}.
$$
\end{definition}

\noindent
The basic property we need is, roughly:
There is an equivalence of categories
between representations of groups
and representations of the group algebras.

The AI-formulation of that is:
\begin{theorem}
\label{representation-theory-theorem-representations-group-algebra}
Let $G$ be a group. The category of complex
representations of $G$ is equivalent to the
category of left $\mathbb{C}[G]$-modules.

Under this equivalence, a representation
$$
\rho:G\to \operatorname{GL}(V)
$$
corresponds to the $\mathbb{C}[G]$-module
structure on $V$ defined by
$$
\left(\sum_{g\in G} a_g e_g\right)v
=
\sum_{g\in G} a_g \rho(g)v.
$$
Conversely, a left $\mathbb{C}[G]$-module
$V$ gives a representation of $G$ by letting
$g$ act as multiplication by $e_g$.
\end{theorem}

\section{Representations of SL(2,C)}
\label{representation-theory-section-sl2c}

\noindent
Klein classified nontrivial finite subgroups
of $\text{SL}_2(\mathbb{C})$
up to conjugacy:
\begin{itemize}
\item Cyclic groups $\mathcal{C}_m$, $m \geq
2$. $|\mathcal{C}_m|=m$,
$$
\mathcal{C}_m=\left\{\begin{pmatrix}
\zeta&0\\
0&\zeta^{-1}\end{pmatrix}:\zeta^m=1
\right\}.
$$
\item Dihedral groups $\mathcal{D}_m$, $m\geq
2$,
$$
\mathcal{D}_m=\left\langle{}\mathcal{C}_{2m},
\begin{pmatrix}
0&-1\\
1&0
\end{pmatrix}\right\rangle{}.
$$
For example, $\mathcal{D}_2$ is the {\it
quaternion group}
$\left\langle{}\underbrace{\begin{pmatrix}
i&0\\
0&-i
\end{pmatrix}}_{i},
\underbrace{\begin{pmatrix}
0&-1\\
1&0
\end{pmatrix}}_{j}\right\rangle{}$.

\item Binary tetrahedral group $\mathcal{T}$,
with $|\mathcal{T}|=24$.
\item Binary octahedral group, $\mathcal{O}$,
with $|\mathcal{O}|=48$.
\item Binary icosahedral group $\mathcal{I}$
with $|I|=120$.

These three are related to the rotation
groups of the
Platonic solids by the 2-to-1 covering
$\text{SU}_2 \to \text{SO}_3(\mathbb{R})$.
To classify the finite subgroups of
$\text{SL}_2(\mathbb{C})$:
first show that any such subgroup is
conjugate to a subgroup
of $\text{SU}_2$ so that its image is a
finite subgroup of
$\text{SO}_3(\mathbb{R})$. We can classify
those by considering their
action on the sphere and showing that they
have to be the
rotation group of one of the Platonic solids
or a dihedral group
or a cyclic group (related to planar
polygons).
\end{itemize}

\noindent
At some point of the XX century it was
realised that this
classification is an instance of an ADE
classification.

Now we explain the (historically second
explanation of this) McKay
correspondence.
If $\Gamma$ is a nontrivial finite subgroup
of $\text{SL}_2(\mathbb{C})$,
define the  {\it McKay graph}
$\tilde{\Delta}_\Gamma$ as follows:
the vertices are $0,1,2,\ldots,\ell$ which
index the
irreducible representations
$R_0,R_1,\ldots,R_\ell$ of $\Gamma$.

Let $Q$ be the 2-dimensional $\Gamma$ coming
from
$\Gamma<\text{SL}_2(\mathbb{C})$.
(Note that $Q$ is one of the $R_i$'s unless
$\Gamma=\mathcal{C}_m$,
in which case $Q$ is reducible.)
The number of edges between $i$ and $j$ in
$\tilde{\Delta}_\Gamma$
is the multiplicity of $R_j$ in $Q \otimes
R_i$.
So: take the representation $R_i$, tensor it
with the
2-dimensional representation $Q$,
express it as a direct sum of irreducible
representations
and perhaps $R_j$ will appear with nonzero
multiplicity,
in which case you put $R_j$ with that number
of edges.
So in general it's a multigraph, though in
most
cases its a simple graph.
Also note that this multiplicity coincides
with the multiplicity of $R_i$ in $Q \otimes
R_j$
since $Q^* \cong Q$; that is the graph is not
directed.

The surprising thing is that when we do this
we obtain
the affine Dynkin diagrams. McKay's famous
observation says:
\begin{enumerate}
\item $\tilde{\Delta}_\Gamma$ is a
simply-laced affine Dynkin
diagram.

\item This is a bijection:
$$
\left\{\substack{\text{nontrvial finite} \\
\Gamma \subset
\text{SL}_2(\mathbb{C})}\right\}\Big/
\text{conj.}
\leftrightarrow
\left\{\substack{\text{simply-laced} \\
\text{affine Dynkin diagrams}}\right\}.
$$
\end{enumerate}

\noindent
The correspondance reads
$$
\begin{aligned}
\mathcal{C}_m &\leftrightarrow
\widetilde{A}_{m-1}\\
\mathcal{D}_m &\leftrightarrow
\widetilde{D}_{m+2}\\
\mathcal{T}&\leftrightarrow\tilde{E}
_6\\
\mathcal{O}&\leftrightarrow\tilde{E}
_7\\
\mathcal{I}&\leftrightarrow\tilde{E}
_8.
\end{aligned}
$$

\medskip\noindent
Now we explain the geometric approach.
Given a finite $\Gamma \subset
\text{SL}_2(\mathbb{C})$, consider
the quotient $Y_\Gamma=\Gamma\setminus Q$ as
an algebraic variety,
that is, (since $Q$ is just a representation
of $\Gamma$, i.e.
a 2-dimensional affine space with an action
of $\Gamma$)
we can take its quotient in the sense
described before:
$Y_\Gamma$ is the affine variety with
coordinate algebra
$\mathbb{C}[x,y]^\Gamma$.

\begin{lemma}
\label{representation-theory-lemma-action-is-free-outside-origin}
The action of $\Gamma$ is free on
$Q\setminus\{0\}$ is free.
\end{lemma}


\medskip\noindent
{\bf Unit quaternions and rotations}
Notes from Conway's book on quaternions.

\begin{lemma}
\label{representation-theory-lemma-so3-simple-rotations}
$\text{SO}(3)$ is generated by simple
rotations, i.e. rotations which fix a
vector.
\end{lemma}

\begin{lemma}
\label{representation-theory-lemma-unit-quaternions-cover-so3}
Every element of $\text{SO}(3)$ has the form
$x \mapsto q^{-1}xq$ for some unit
quaternion $q$. The map
$$
q \longmapsto [q],
\qquad
[q](x)=q^{-1}xq
$$
is a $2$-to-$1$ homomorphism from the group
of unit quaternions to $\text{SO}(3)$.
\end{lemma}

\noindent
The imaginary quaternions
$$
\operatorname{Im}\mathbb{H}
=
\mathbb{R}i\oplus\mathbb{R}j
\oplus\mathbb{R}k
$$
are naturally $\mathbb{R}^3$. The unit
quaternions act on them by conjugation:
$$
\rho(\alpha)(\beta)
=
\alpha\beta\overline{\alpha}.
$$
Thus we get a homomorphism
$$
\rho:S^3 \longrightarrow
\text{SO}_3(\mathbb{R})
$$
with kernel $\{\pm 1\}$. This identifies the
unit quaternions with $\text{Spin}(3)$, and
also with $\text{SU}(2)$. The covering is
$2$-to-$1$ because opposite unit quaternions
give the same rotation.

\begin{definition}
\label{representation-theory-definition-binary-icosahedral}
The {\it binary icosahedral group} is the
preimage of the icosahedral rotation group
$I\cong A_5$ under the covering
$$
\text{SU}(2) \longrightarrow
\text{SO}_3(\mathbb{R}).
$$
\end{definition}

\begin{remark}
\label{representation-theory-remark-quaternions-binary-groups}
This is the bridge between the rotation
groups of the Platonic solids and the finite
subgroups of $\text{SU}(2)$, or equivalently
of $\text{SL}_2(\mathbb{C})$. The word
``binary'' means ``take the inverse image
under the double cover''. For example, the
quaternion group $Q_8$ is the binary
dihedral group $\mathcal{D}_2$.
\end{remark}




\medskip\noindent
{\bf The quaternionic double cover of $SO(4)$}

The preceding covering describes rotations of
$\operatorname{Im}\mathbb H\cong\mathbb R^3$.
For rotations of $\mathbb R^4\cong\mathbb H$,
we use \emph{two} unit quaternions rather
than one.

\begin{proposition}
\label{representation-theory-proposition-quaternionic-cover-so4}
The map
$
\begin{aligned}
S^3\times S^3&\longrightarrow SO(4),\\
(u,v)&\longmapsto
\bigl(x\longmapsto uxv^{-1}\bigr)
\end{aligned}
$
is a surjective group homomorphism with
kernel $\{(1,1),(-1,-1)\}$. In particular,
$
SO(4)\cong
(S^3\times S^3)/\{\pm(1,1)\}.
$
\end{proposition}

\begin{proof}
Left and right multiplication by unit
quaternions are real-linear norm-preserving
maps of $\mathbb H$, and each has determinant
$1$ by connectedness. The formula is
multiplicative. If $uxv^{-1}=x$ for every
$x$, setting $x=1$ gives $u=v$; commuting
with every quaternion forces $u=v=\pm1$.
The differential at $(1,1)$ sends imaginary
quaternions $(a,b)$ to $x\mapsto ax-xb$.
It is injective: evaluating at $1$ gives
$a=b$, and then $a$ commutes with all
quaternions only if $a=0$. Domain and
codomain have dimension six, so the image
is an open subgroup of connected $SO(4)$,
hence all of it.
\end{proof}

\begin{remark}
\label{representation-theory-remark-so3-diagonal-inside-so4}
The subgroup fixing $1\in\mathbb H$ is
represented by pairs $(u,u)$: on imaginary
quaternions it acts as
$x\mapsto uxu^{-1}$. Thus the earlier
$S^3\to SO(3)$ covering reappears as the
\emph{diagonal} inside the $SO(4)$
covering. If a base vertex is another
unit quaternion $w$, changing coordinates
by an orthogonal transformation taking
$w$ to $1$ gives the same picture.
\end{remark}

\begin{remark}[Application to the PETCOX trace restriction]
\label{representation-theory-remark-petcox-quaternionic-trace}
In PETCOX, let
$\Gamma=\langle S_1,S_2,S_3\rangle
\leq SO(4)$ and
$J=\langle S_2,S_3\rangle$ be the
base-vertex stabilizer. Assuming a genuine
finite extension, Lemma 4.1 of
\texttt{petcox-extensions-v2} says that
$J$ is tetrahedral, octahedral, or
icosahedral. By the preceding remark,
its full lift is diagonal and both
coordinate projections of the lift of
$\Gamma$ contain the corresponding binary
polyhedral group.

The two projected groups are finite
subgroups of $S^3$. Since they contain a
noncyclic, nondihedral binary polyhedral
group, their images in $SO(3)$ are among
the tetrahedral, octahedral and
icosahedral rotation groups. Their
elements therefore have real parts in
the finite set
$
\mathcal C=
\left\{0,\ \pm1,\ \pm\frac12,\
\pm\frac{\sqrt2}{2},\
\pm\frac{\varphi}{2},\
\pm\frac{1}{2\varphi}\right\},
\qquad \varphi=\frac{1+\sqrt5}{2}.
$
Directly taking traces of the real-linear
map $x\mapsto uxv^{-1}$ gives
$
\operatorname{tr}_{\mathbb R^4}
(x\mapsto uxv^{-1})
=4\operatorname{Re}(u)\operatorname{Re}(v).
$
Consequently all traces in $\Gamma$ lie
in the finite set
$\{4ab:a,b\in\mathcal C\}$.
This is the key input to the polynomial
sieve in Section 4 of that manuscript.
The argument is conditional on
$\Gamma$ being finite; it does not prove
that an arbitrary PETCOX parameter
produces an extension.
\end{remark}

\section{Representations of finite groups}
\label{representation-theory-section-finite-group-reps}

\noindent
Exercises from a home assignment
by Misha on his Commutative Algebra
2026 course.
Current state of this:
to solve these problems
we need to know lots of basic
notions in representation theory
of groups that I do not know yet.

\begin{definition}[AI --- should
look in Fulton for this]
\label{representation-theory-definition-finitely-generated-rep}
Let $G$ be a group and let $k$ be a field.
A representation
$$
\rho:G\to \operatorname{GL}(V)
$$
is called {\it finitely generated} if $V$ is
finitely generated as a $k[G]$-module.

Equivalently, there exist finitely many
vectors
$$
v_1,\ldots,v_n\in V
$$
such that every vector of $V$ is a finite
$k$-linear combination of vectors of the form
$$
\rho(g)v_i,
\qquad g\in G,\quad i=1,\ldots,n.
$$
\end{definition}

\begin{exercise}
\label{representation-theory-exercise-fg-c-t-inv-modules-z-reps}
Prove that the category of all finitely
generated $\mathbb{C}[t,t^{-1}]$-modules is
equivalent to the category of finitely
generated representations of $\mathbb{Z}$.
\end{exercise}

\begin{exercise}
\label{representation-theory-exercise-z2-reps-fd-vs-all}
Let $C_1$ be the category of all
finite-dimensional representations of
$\mathbb{Z}/2\mathbb{Z}$ over $\mathbb{C}$,
and $C_2$ the category of all
representations (finite and
infinite-dimensional). Are the categories
$C_1$ and $C_2$ equivalent?
\end{exercise}

\begin{proof}
The relevant
finite-dimensional representations
are just a choice of $n>0$ and
an idempotent matrix $A \in
\text{GL}(\mathbb{C},n)$.
Perhaps the choice of $A$
is just $-\text{Id}$ up
to conjugation,
so maybe it's literally just
a choice of $n>1$ and a matrix in
$\text{GL}(\mathbb{C},n)$.

Infinite-dimensional representations
are just… well, an infinite-dimensional
complex vector space and an involution.

I need to find a way to deal
with equivalence of categories…
Something like: suppose they were
equivalent. Then there exists fully
faithful essentially surjective functor
between them.
Then for every infinite-dimensional
vector space there is a finite-dimensional
one which corresponds…
but the category of finite-dimensional
ones is just natural numbers,
and the infinite-dimensional one
is…
\end{proof}

\noindent
The upshot is that 
we can define irreducibility
in categorical language:
a representation $A$
is irreducible if and only if for every
monomorphism $f:B \to A$,
$f$ is an isomorphism or $B \simeq 0$.
Note that $0$ is both an initial
and terminal object in the category
of vector spaces.

\begin{exercise}
\label{representation-theory-exercise-z2-reps-vs-vect}
Let $C_1$ be the category of all
finite-dimensional representations of
$\mathbb{Z}/2\mathbb{Z}$ over $\mathbb{C}$,
and $C_2$ the category of all finite-
dimensional vector spaces over
$\mathbb{C}$. Are the categories $C_1$ and
$C_2$ equivalent?
\end{exercise}

\begin{exercise}
\label{representation-theory-exercise-z2-vs-s3-reps}
Let $C_1$ be the category of all
finite-dimensional representations of
$\mathbb{Z}/2\mathbb{Z}$ over $\mathbb{C}$,
and $C_2$ the category of all
finite-dimensional representations of the
symmetric group $S_3$ over $\mathbb{C}$.
Are the categories $C_1$ and $C_2$
equivalent?
\end{exercise}

\begin{proof}
Recall that the symmetric group
is never simple: the kernel
of the parity map $S_n \to \mathbb{Z}/2\mathbb{Z}$
is a normal subgroup. For us,
$$
\xymatrix{
0\ar[r]
&\mathbb{Z}/3\ar[r]
&S_3\ar[r]
&\mathbb{Z}/2\ar[r]
&0
}
$$

\end{proof}

\noindent
Consider doing the following
extra exercise. Pick
a group $G$ and 
consider the group algebra
$\mathbb{C}[G]$. Prove
that $\mathbb{C}[G]$ contains
all the representations of $G$,
where the multiplicity of a representation
$R$ is in fact $\dim R$.
And $\dim \mathbb{C}[R]$
is the order of the group.
See Serre's book on finite-dimensional
group representations.

Meta exercise: which
properties of representations can
be defined categorically? Just
like irreducibility above.

\begin{exercise}
\label{representation-theory-exercise-z2x-z2-vs-z4-reps}
Let $C_1$ be the category of all
finite-dimensional representations of
$\mathbb{Z}/2\mathbb{Z}\times\mathbb{Z}/2\mathbb{Z}$
over $\mathbb{C}$, and $C_2$ the category of
all finite-dimensional representations of
$\mathbb{Z}/4\mathbb{Z}$ over $\mathbb{C}$.
Are the categories $C_1$ and $C_2$
equivalent?
\end{exercise}

\begin{proof}
The idea of the proof is to mimic
the following argument. The
category of all finite-dimensional
vector spaces is equivalent
to the category of finite-dimensional
vector spaces over $\mathbb{C}$.
The functor giving this equivalence
is just given by choosing
a basis of an arbitrary finite-dimensional
vector space and mapping it
to the corresponding $\mathbb{C}^n$.

Somehow, we are supposed to do the same in
this exercise. Here's the AI-answer,
which in fact coincides with Misha's and
Vanya's:

Both categories are equivalent to the product
category
$$
\operatorname{Vect}_{fd}(\mathbb{C})^4.
$$
For $\mathbb{Z}/2\mathbb{Z}\times
\mathbb{Z}/2\mathbb{Z}$, a representation is
given by two commuting involutions $A,B$.
They are simultaneously diagonalizable, so
the representation decomposes as
$$
V=
V_{++}\oplus V_{+-}\oplus V_{-+}
\oplus V_{--}.
$$
A morphism of representations preserves
these four summands.

For $\mathbb{Z}/4\mathbb{Z}$, a
representation is given by an operator $T$
with $T^4=I$. Since $x^4-1$ has four distinct
roots over $\mathbb{C}$, the operator $T$ is
diagonalizable and
$$
V=
V_1\oplus V_{-1}\oplus V_i\oplus V_{-i}.
$$
Again, a morphism of representations
preserves these four summands.

Thus both categories are equivalent to the
category of quadruples of finite-dimensional
complex vector spaces. Hence they are
equivalent as categories.
\end{proof}

\begin{exercise}
\label{representation-theory-exercise-r-vs-c-vect}
Let $C_1,C_2$ be the categories of
finite-dimensional vector spaces over
$\mathbb{R}$ and $\mathbb{C}$. Are the
categories $C_1$ and $C_2$ equivalent?
\end{exercise}

\begin{exercise}
\label{representation-theory-exercise-finite-groups-vs-abelian}
Prove that the category of finite groups is
not equivalent to the category of finite
abelian groups.
\end{exercise}

\begin{proof}
The key idea is that product and coproduct
are the same object in the category of
abelian groups! That's essentially by
the universal property of coproduct
of maps factoring through $A \coprod B$,
by abelianess, the direct sum is
also the coproduct, but this is not
true for non-abelian groups.
In the latter cases the coproduct
is called {\it amalgam product}.

An alternative strategy
is to observe that simple objects in
the category of abelian groups are
$\mathbb{Z}/p$. Then we can find a simple
group $A_5$, in the category
of noncommutative groups,
which has an automorphism group
that cannot be
the automorphism group a $\mathbb{Z}/p$.
\end{proof}

\section{Borel-Weil-Bott theorem}
\label{representation-theory-section-Borel-Weil-Bott-thoerem}

[These are {\it very} sketchy notes after a
conversation with Lyalya Guseva on
Borel-Weil-Bott theorem]

Borel-Weil-Bott theorem is a device for
computing cohomologies of sheaves.

Take the Grassmanian $\text{Gr}(n,k)$, which
may be obtained as a quotient $G$
of a group $G$ by a parabolic subgroup $P$.
We would like it if $P$ was a
semisimple group since those are classified,
but unfortunately it is not.
So instead we use so-called Levi quotient
denoted by $L$ which is semisimple
and allows us to understand the cohomologies
of the variety $G/P$.

Fortunately the representation theory of
$\text{Gr}(k,n)$ is well known, in fact
we get $\text{GL}_k \times \text{GL}_{n-k}$.

There are functors associated to $L$, which
are described by a sequence of
numbers $a_1,\ldots,a_n$. Along with other
sequences of numbers
$k_1,\ldots,k_\ell$, and $\rho$ (the latter
is a concept in representation
theory but we may ultimately think of it as
another sequence of numbers) we may
construct an action of the symmetric group
$S^n$ acting on these sequences and
obtain a result concerning the cohomology
$H^{\ell(\sigma)}(L)$, and in
particular we find that if two entries in our
list of numbers coincide, the
cohomology will vanish.



% --- End representation-theory.tex ---

% --- Begin vertex-algebras.tex ---
\chapter{Vertex Algebras}
\label{chapter-vertex-algebras}





\section{Cartan subalgebra, Cartan matrix and
Serre relations}
\label{vertex-algebras-section-Cartan-subalgebra}

Kac-Moody algebras are Lie algebras, whose
definition is motivated by the
structure of finite-dimensional simple Lie
algebras over $\mathbb{C}$.

Let $\mathfrak{g}$ be a finite-dimensional
semisimple Lie algebra over
$\mathbb{C}$. Then $\mathfrak{g}$ has a {\it
Cartan subalgebra}
$\mathfrak{h}\subset \mathfrak{g}$ (abelian +
…).
Fixing $\mathfrak{h}\subset\mathfrak{g}$
gives a
{\it root space decomposition}
$$
\mathfrak{g}=\mathfrak{h}\oplus
\bigoplus_{\alpha \in
\Delta}\mathfrak{g}_\alpha
$$
where $\Delta \subset \mathfrak{h}^*$ linear
dual, and, by definition
$$
\mathfrak{g}_\alpha=\{X \in \mathfrak{g}|
[H,X]=\alpha(H)X\; \forall H \in
\mathfrak{h}\}
$$
Turns out the $\mathfrak{g}_\alpha$ are all
1-dimensional,
though this property is lost when we go to
Kac-Moody algebras.
$$
[\mathfrak{g}_\alpha,\mathfrak{g}_\beta]
\subset \mathfrak{g}_{\alpha+\beta}
$$
The Killing form  $\kappa:\mathfrak{g} \times
\mathfrak{g} \to \mathbb{C}$,
 $\kappa(x,y)=\text{Tr}_\mathfrak{g}\text{ad}
 (x)\text{ad}(y)$
 is nondegenerate.
``This is kind of the definition of
semisimple.''
(Think of $\mathfrak{h}$ as $\mathfrak{g}_0$,
btw.)

$\kappa |_{\mathfrak{g}_\alpha \times
\mathfrak{g}_\beta}\neq 0$ only when
$\beta=-\alpha$. $\kappa
|_{\mathfrak{h}\times \mathfrak{h}}$ is
non-degenerate.
 This gives a linear isomorphism
$\mathfrak{h} \xrightarrow{\nu} \mathfrak{h}$
via
$\nu(H)(H') = \kappa(H,H')$.

\medskip\noindent
So, $\mathfrak{h}^*$ comes with a
non-degenerate bilinear form.

The {\it reflection}
$r_\alpha:\mathfrak{h}\to \mathfrak{h}^*$ in
$\alpha \in \mathfrak{h}^*$ (usually a root)
is
$r_\alpha(\lambda)=\lambda-
2\frac{(\lambda,\alpha)}{(\alpha,\alpha)}
\cdot \alpha.$

``Classify root systems […] classify
semisimple Lie algebras''
It is a fact that $r_\alpha (\Delta)=\Delta$
for all $\alpha \in \Delta$,
which motivates the definition of {\it root
system}
and permits classification.

\begin{example}
\label{vertex-algebras-example-root-system-sl2}
$\mathfrak{g}=\mathfrak{sl}_2$,
$\mathfrak{h}=$ diagonal matrices
$$
\begin{pmatrix}
1&&\\
&-1&\\
& & 0
\end{pmatrix}
\qquad
\begin{pmatrix}
0&&\\
&1&\\
& & -1
\end{pmatrix}
$$
is a basis of $\mathfrak{h}$. There are 6
roots vectors
$$
E_{12}=
\begin{pmatrix}
0 & 1 & 0\\
0 & 0 & 0\\
0 & 0 & 0
\end{pmatrix}
$$
$E_{23},E_{13}$, etc.
\end{example}

\begin{exercise}
\label{vertex-algebras-exercise-basis-of-h}
$[H_1,E_{12}]=2E_{12}$,
$[H_2,E_{12}]=-E_{12}$, $\alpha_{12}=(2,-1)$.
\end{exercise}

[Drawing of roots]

Notions of {\it positive roots} and {\it
simple roots}
(set of $\text{rank}\mathfrak{g}$ simple
roots has $\ell$ elements,
 where $\ell=\dim(\mathfrak{h}^*)$.
This will also fail for Kac-Moody algebras
more generally).
Next write the {\it Cartan matrix}
$$
A=(a_{ij}),\qquad
a_{ij}=2\frac{(\alpha_i,\alpha_j)}{(\alpha_i,
\alpha_i)}
$$
for $1\leq i,j \leq \ell$.

\begin{example}
\label{vertex-algebras-example-Cartan-matrix-sl3}
$\mathfrak{sl}_3$. [Picture, hexagonal
pattern].
$(\alpha_1,\alpha_1)=(\alpha_2,\alpha_2)=2,
(\alpha_1,\alpha_2)=-1$,
so
$$
A=\begin{pmatrix}
2&-1\\
-1&2
\end{pmatrix}
$$
\end{example}

\begin{example}
\label{vertex-algebras-example-Cartan-matrix-sl5}
$\mathfrak{sl}_5$. [Picture, square pattern].
$|\alpha_2|=1$, $|\alpha_1|=2$,
$(\alpha_1,\alpha_2)=-2$, so
$$
A=\begin{pmatrix}
2&-1\\
-2&2
\end{pmatrix}
$$
\end{example}

\begin{remark}
\label{vertex-algebras-remark}
Since $\mathfrak{g}_\alpha$ is 1-dimensional,
set
$\mathfrak{g}_\alpha=\mathbb{C}E_\alpha$ and
$E_i=E_{\alpha_i}$, $i=1,2,\ldots,\ell$
(simple root vectors).
It turns out that
$$
\text{ad}(E_i)^{1-a_{ij}}E_j=0.
$$
This is called a {\it Serre relation}.
\end{remark}

\section{Some infinite dimensional Lie
algebras}
\label{vertex-algebras-section-some-infin-dimen-lie-algeb}
Let $\mathfrak{g}$ be a finite-dimensional
semisimple Lie algebra, and define
the {\it loop algebra}
$$
\begin{aligned}
L\mathfrak{g}&=\mathfrak{g}[t,t^{-1}],\text{
(with basis
$at^m|\substack{\text{$a \in $ a basis of
$\mathfrak{g}$} \\ m \in \mathbb{Z}}$
)}\\
&=\mathfrak{g} \otimes
_{\mathbb{C}}\mathbb{C}[t,t^{-1}]
\end{aligned}
$$
with the Lie bracket
$$
[at^m,bt^n]=[a,b]t^{m+n}.
$$
``This construction is absurdely general ---
we don't need $\mathfrak{g}$ to be
semisimple […]''

\medskip\noindent
Take $\mathfrak{g}=\mathfrak{sl}_2$. Recall
that
$$
E=\begin{pmatrix}
0&1\\
0&0
\end{pmatrix},\quad
H=\begin{pmatrix}
1&0\\
0&-1
\end{pmatrix},\quad
F=\begin{pmatrix}
0&0\\
1&0
\end{pmatrix}
$$
[Picture with $F,H,E,Ft,Ht,Et,Et^2$…]
$E$ was a root vector, corresponding to the
unique root in $\mathfrak{sl}_2$,
call it $\alpha_1$. We seem to have a second
simple root $\alpha_0$,
 corresponding to $Ft$.

This looks like it wants to have a Cartan
matrix
 $$
A=\begin{pmatrix}
2&-2\\
-2&2
\end{pmatrix}
$$
We will indeed recover (a variant of)
$L\mathfrak{g}$ as a Lie algebra
``built from'' $A=\begin{pmatrix}
2&-2\\
-2&2
\end{pmatrix}$,
 a Kac-Moody algebra.
But note first, $\mathfrak{h}=\mathbb{C}H$ is
too small.
``Problem with $\alpha_0$ and $\alpha_1$
being linearly independent  …''

\begin{exercise}
\label{vertex-algebras-exercise-Lie-algebra-d}
Consider $L\mathfrak{g}\oplus\mathbb{C}d$,
and set
$[d,at^m]=mat^m$, $[d,d]=0$. Check this
defines a Lie algebra.
\end{exercise}

\begin{proof}
Skew-commutativity, i.e. for all $x \in
L\mathfrak{g}\oplus\mathbb{C}d$,
\begin{equation}
\label{vertex-algebras-equation-skew-commutativity}
[x,x]=0,
\end{equation}
is immediate from
skew commutativity in $L\mathfrak{g}$ and
the hypothesis that $[d,d]=0$.

To confirm Jacobi identity, i.e. that for all
$x,y,z \in L\mathfrak{g}\oplus\mathbb{C}d$
\begin{equation}
\label{vertex-algebras-equation-Jacobi-algebra-d-identity}
[x,[y,z]]+[y,[z,x]]+[z,[x,y]]=0,
\end{equation}
notice that since this is a cyclic sum on
$x,y,z$
we only need to consider three elements in
$L\mathfrak{g} \oplus \mathbb{C}d$ up to
cyclic permutation.
The cases in which the three elements are
either in $L\mathfrak{g}$ or in
$\mathbb{C}d$ are obvious, so that there
are only two interesting possibilities:
\begin{equation}
\label{vertex-algebras-equation-Jacobi-algebra-d-case-1}
x=d,\qquad y=at^m,\qquad z=bt^n
\end{equation}
and
\begin{equation}
\label{vertex-algebras-equation-Jacobi-algebra-d-case-2}
x=d,\qquad y=d,\qquad z=at^n.
\end{equation}
Case \ref{vertex-algebras-equation-Jacobi-algebra-d-case-1}
gives
$$
\begin{aligned}
&
[d,[at^m,bt^n]]+[at^m,[bt^n,d]]+[bt^n,[d,
at^m]]\\
&=[d,[a,b]t^{m+n}]+[at^m,-nbt^n]+[bt^n,mat^m]
\\
&=(m+n)[a,b]t^{m+n}-n[a,b]t^{m+n}+m[b,a]t^{m+
n}\\
&=(m+n)[a,b]t^{m+n}-n[a,b]t^{m+n}-m[a,b]t^{m+
n}\\
&=(m+n)[a,b]t^{m+n}-(m+n)[a,b]t^{m+n}=0.
\end{aligned}
$$
Case \ref{vertex-algebras-equation-Jacobi-algebra-d-case-2}
gives
$$
\begin{aligned}
&[d,[d,at^m]]+[d,[at^m,d]]+[at^m,[d,d]]\\
&=[d,mat^m]+[d,-mat^m]=0.
\end{aligned}
$$

\end{proof}

The Kac-Moody algebra turns out to be, not
quite this, but slightly larger
still.

Recall that an {\it invariant bilinear form}
$(\cdot,\cdot)$ on a Lie algebra
$\mathfrak{g}$ is a bilinear form such that
\begin{equation}
\label{vertex-algebras-equation-invariant-bilinear-form}
([a,b],c)=(a,[b,c])\qquad \forall a,b,c \in
\mathfrak{g}.
\end{equation}

\begin{exercise}
\label{vertex-algebras-exercise-invar-bilin-form-simpl}
Prove that an invariant bilinear form on a
simple Lie algebra must in fact be
symmetric.
\end{exercise}

\begin{proof}
It's enough to show that $\mathfrak{g}$ is
{\it perfect}, i.e. that
$[\mathfrak{g},\mathfrak{g}]=\mathfrak{g}$.
In this
case, let $a,b \in \mathfrak{g}$ and suppose
that $b=[x,y]$. Then
$$
\begin{aligned}
(a,b)&=(a,[x,y])=(a,-[y,x])=(-[a,y],x)=([y,a]
,x)\\
&=(y,[a,x])=(y,-[x,a])=(-[y,x],a)=([x,y],a)=
(b,a)
\end{aligned}
$$
To confirm that $\mathfrak{g}$ is perfect
just observe that
$[\mathfrak{g},\mathfrak{g}]$ is a nontrivial
ideal of $\mathfrak{g}$.
\end{proof}

\begin{definition}
\label{vertex-algebras-definition-affine-Lie-algebra}
Given $\mathfrak{g}$ simple, with
$(\cdot,\cdot):\mathfrak{g} \times
\mathfrak{g} \to \mathbb{C}$
invariant bilinear form, the {\it affine Lie
algebra} is
$$
\hat{\mathfrak{g}}=L\mathfrak{g} \oplus
\mathbb{C}K,
$$
with $[K,\hat{\mathfrak{g}}]=0$, and
$[at^m,bt^n]=[a,b]t^{m+n}+m(a,b)\delta_{m,-n}
K$.
\end{definition}

``For the construction to work it doesn't
actually have to be nondegenerate.''

\begin{exercise}
\label{vertex-algebras-exercise-affin-lie-algeb-lie-algeb}
Check that the affine Lie algebra
$\hat{\mathfrak{g}}$ is a Lie algebra.
\end{exercise}

\begin{proof}
(Skew-commutativity.)
Since $[K,\hat{\mathfrak{g}}]=0$ and $K \in
\hat{\mathfrak{g}}$,
it is immediate that $[K,K]=0$.
For the case of an element in
$L\mathfrak{g}$, we see that
$[at^m,at^m]=0$ by skew-commutativity of the
bracket in $\mathfrak{g}$
and the Kronecker delta.

(Jacobi identity.) As in Exercise
\ref{vertex-algebras-exercise-Lie-algebra-d},
any choice of $x,y,z$ involving $K$ is
immediate by $[K,\hat{\mathfrak{g}}]=0$.
Thus the only interesting case is
for Jacobi identity consider the cases
$$
\begin{aligned}
&
[at^m,[bt^n,ct^\ell]]+[bt^n,[ct^\ell,at^m]]+
[ct^\ell,[at^m,bt^n]]\\
&=[at^m,[b,c]t^{n+\ell}+n(b,c)\delta_{n,-
\ell}K]\\
&+[bt^n,[c,a]t^{\ell+m}+\ell(c,a)
\delta_{\ell,-m}K]\\
&+[ct^\ell,[a,b]t^{m+n}+m(a,b)\delta_{m,-n}K]
\\
&=[at^m,[b,c]t^{n+\ell}]+[at^m,n(b,c)
\delta_{n,-\ell}K]\\
&+[bt^n,[c,a]t^{\ell+m}]+[bt^n,\ell(c,a)
\delta_{\ell,-m}K]\\
&+[ct^\ell,[a,b]t^{m+n}]+[ct^\ell,m(a,b)
\delta_{m,-n}K]\\
&=[a,[b,c]]t^{m+(n+\ell)}+m(a,[b,c])
\delta_{m,-(n+\ell)}K\\
&+[b,[c,a]]t^{n+(\ell+m)}+n(b,[c,a])
\delta_{n,-(\ell+m)}K\\
&+[c,[a,b]]t^{\ell+(m+n)}+\ell(c,[a,b])
\delta_{\ell,-(m+n)}K=0
\end{aligned}
$$
It is clear that we obtain a Jacobi equation
on $\mathfrak{g}$. To see that the
remaining terms vanish, notice that the
condition on
the Kronecker delta in its three appearances
is the same,
namely, $m+n+\ell=0$. In this case, we only
need to check
that $(a,[b,c])=(b,[c,a])=(c,[a,b])$ to
conclude.
This follows from the invariance of
$(\cdot,\cdot)$
and the fact that $\mathfrak{g}$ simple using
Exercise
\ref{vertex-algebras-exercise-invar-bilin-form-simpl}.
\end{proof}

We also have
\begin{definition}
\label{vertex-algebras-definition-exten-affin-lie-algeb}
The {\it extended affine Lie algebra} is
$$
\tilde{\mathfrak{g}}=L\mathfrak{g} \oplus
\mathbb{C}K \oplus \mathbb{C}d,
$$
with $[d,at^m]=m a t^m$ as before, and
$[K,d]=0$.
\end{definition}

The extended affine Lie algebra is an example
of a Kac-Moody algebra.

\begin{exercise}[For those who like geometry]
\label{vertex-algebras-exercise-for-those-who-like-geometry}
Let $R=\mathbb{C}[t,t^{-1}]$.
If $D \in \text{Der}(R)$, then
$L\mathfrak{g} \oplus \mathbb{C}d$ is a Lie
algebra with
$[d,a \otimes r]=a\otimes D(r)$.
Is $(\mathfrak{g}\otimes R)\oplus
\text{Der}(R)$ a Lie algebra? (The Lie alegra
$L\mathfrak{g} \oplus \mathbb{C}d$ from
Exercise \ref{vertex-algebras-exercise-Lie-algebra-d} is
a particular case, for $D=t\frac{d}{dt}$.)
\end{exercise}

\begin{proof}
Checking that $L\mathfrak{g} \oplus
\mathbb{C}d$ is a Lie algebra with
$[d,a\otimes r]=a \otimes D(r)$ is similar to
Exercise
\ref{vertex-algebras-exercise-Lie-algebra-d}:
skew-commutativity is immediate from
skew-commutativity in each of the components,
while
Jacobi identity is verified in two cases. For
$x=y=d$ and $z=a \otimes r$
we quickly obtain
$$
\begin{aligned}
&[x,[y,z]]+[y,[z,x]]+[z,[x,y]]\\
&=[d,[d,a \otimes r]]+[d,[a \otimes r,d]]+[a
\otimes r,[d,d]]\\
&=[d,a \otimes D(r)]+[d,-a \otimes D(r)]=0.
\end{aligned}
$$

And for $x=d$, $y=a \otimes r$ and $z=b
\otimes s$, we get
\begin{equation}
\label{vertex-algebras-equation-Jacobi-derivation}
\begin{aligned}
&[x,[y,z]]+[y,[z,x]]+[z,[x,y]]\\
&=[d,[a\otimes r,b \otimes s]]+[a\otimes r,[b
\otimes s,d]]
+[b \otimes s,[d,a\otimes r]]\\
&=[d,[a,b]\otimes rs]+[a\otimes r,-b\otimes
D(s)]
+[b\otimes s,a\otimes D(r)]\\
&=[a,b]\otimes D(rs)-[a,b]\otimes
rD(s)+[b,a]\otimes sD(r)=0.
\end{aligned}
\end{equation}

\medskip\noindent
To check whether $(\mathfrak{g}\otimes
R)\oplus\text{Der}(R)$ is a Lie algebra
first put the Lie bracket on $\text{Der}(R)$
as
$[D,D_1]=DD_1-D_1D$. It is clear that this
bracket is skew-commutative.
Jacobi identity reads
$$
\begin{aligned}
&[D,[D_1,D_2]]+[D_1,[D_2,D]]+[D_2,[D,D_1]]\\
&=[D,D_1D_2-D_2D_1]+[D_1,D_2D-DD_2]+[D_2,
DD_1-D_1D]\\
&=D(D_1D_2-D_2D_1)-(D_1D_2-D_2D_1)D+D_1(D_2D-
DD_2)\\
&-(D_2D-DD_2)D_1+D_2(DD_1-D_1D)-(DD_1-D_1D)
D_2\\
&=DD_1D_2-DD_2D_1-D_1D_2D+D_2D_1D+D_1D_2D-
D_1DD_2\\
&-D_2DD_1+DD_2D_1+D_2DD_1-D_2D_1D-DD_1D_2+
D_1DD_2=0.
\end{aligned}
$$
Now put the bracket on $(\mathfrak{g} \otimes
R)\oplus\text{Der}(R)$ as
$[D,a \otimes  r]=a \otimes D(r)$.
Skew-commutativity is immediate.
Jacobi identity for $x=D,y=a\otimes r$ and
$z=b\otimes s$ is identical to the
computation \ref{vertex-algebras-equation-Jacobi-derivation}.
In the case
$x=D$, $y=D_1$ and $z=a\otimes r$, we get
$$
\begin{aligned}
&[D,[D_1,a\otimes r]]+[D_1,[a\otimes
r,D]]+[a\otimes r,[D,D_1]]\\
&=[D,a\otimes D_1(r)]+[D_1,-a\otimes
D(r)]+[a\otimes r,[D,D_1]]\\
&=a\otimes DD_1(r)-a\otimes D_1D(r)-a\otimes
[D,D_1](r)=0
\end{aligned}
$$
\end{proof}

\section{Kac-Moody algebras}
\label{vertex-algebras-section-Kac-Moody-algebras}

Recall the notion of the free Lie algebra on
a
vector space $V$ of generators
(or a set $X$, think of $V$ as a vector space
with basis $X$):

\begin{definition}
\label{vertex-algebras-definition-free-Lie-algebra}
The {\it free Lie algebra} on $V$ is
characterized by the universal property
$$
\xymatrix{
V\ar[rr]^f\ar[dr]_i&&\mathfrak{g}\\
& F(V)\ar[ur]_{\exists !\tilde{f}}
}
$$
That is, for any linear map $f:V \to
\mathfrak{g}$ with $\mathfrak{g}$ Lie
algebra, there exists a unique $\tilde{f}$
homomorphism of Lie algebras
$F(V)\to \mathfrak{g}$ such that $\tilde{f}
\circ i=f$.
$$
\Hom_{\text{Lie}}(F(V),\mathfrak{g})=
\Hom_{\text{Vec}}(V,\mathfrak{g})
$$
naturally.
\end{definition}

That is, $F$ and the
forgetful functor
$G:\underline{\text{Lie}}\to
\underline{\text{Vec}}$
are adjoint:
$$
\Hom_{\underline{\text{Lie}}}(F(V),
\mathfrak{g})
\xrightarrow{\simeq }
\Hom_{\underline{\text{Vec}}}(V,
G(\mathfrak{g}))
$$

\medskip\noindent
{\bf A realisation of $F(V)$.} Let
$$
T(V)=\mathbb{C} \oplus V \oplus V^{\otimes
2}\oplus V^{\otimes 3}\oplus\ldots
$$
be the tensor algebra of $V$.

Then inside $T(V)$ consider $F(V)$ the span
of iterated commutators of elements
of  $V$.

\begin{proposition}
\label{vertex-algebras-proposition-reali-free-lie-algeb}
This realises the free Lie algebra.
\end{proposition}

\begin{proof}
In online notes.
\end{proof}

\medskip\noindent
In the finite dimensional simple case, we had
$$
a_{ij}=\frac{2(\alpha_i,\alpha_j)}{(\alpha_i,
\alpha_i)},
$$
which we think also as $\alpha_i,\alpha_j \in
\mathfrak{h}^*$,
and
$\alpha_i^\vee=\frac{2}{(\alpha_i,\alpha_i)}
\nu^{-1}(\alpha_i)
\in \mathfrak{h}.$

Clearly, $\alpha_{ii}=2$ for all $i$.
$a_{ij}$ might not equal $a_{ji}$, but
certainly $a_{ij}=0 \iff a_{ji}=0$.
And $\forall  i \neq j$, $a_{ij} \leq 0$.

\medskip\noindent
{\bf One further property.} Set
$$
\varepsilon_i=\frac{2}{(\alpha_i,\alpha_i)},
\quad
\text{and}\quad
D=\substack{\text{diagonal matrix} \\
\text{with entries }\varepsilon_i}
$$
Then $A=DB$, where $B=((\alpha_i,\alpha_j))$
is symmetric.
If a matrix $A$ is equal to
$(\text{diag})(\text{symm})$, we call it
 {\it symmetrizable}.

\begin{definition}
\label{vertex-algebras-definition-generalized-Cartan-matrix}
A {\it generalized Cartan matrix} is an
integer matrix $A=(a_{ij})$
which is
\begin{itemize}
\item symmetrizable,
\item $a_{ii}=2$ for all $i$,
\item $a_{ij}=0 \iff a_{ji}=0$,
\item $a_{ij}\leq 0$ for $i \neq j$.
\end{itemize}
\end{definition}

\begin{definition}
\label{vertex-algebras-definition-realisation}
A {\it realisation} of a generalized Cartan
matrix is a complex vector space
 $\mathfrak{h}$, and two sets
$$
\begin{aligned}
\Pi^\vee&=\{\alpha_1^\vee,
\alpha_2^\vee,\ldots,\alpha_n^\vee\},\quad
\text{and},\\
\Pi&=\{\alpha_1,\alpha_2,\ldots,\alpha_n\}
\end{aligned}
$$
such that
$\left\langle{}\alpha_i^\vee,\alpha_j\right\rangle{}=a_{ij}$,
$1\leq i,j\leq n$.
\end{definition}

\begin{exercise}
\label{vertex-algebras-exercise-realisation}
$\dim(\mathfrak{h})\geq 2n-\text{rank}(A)$.
\end{exercise}

\begin{proof}
For $A=\begin{pmatrix}
2&-2\\
-2&2
\end{pmatrix}$, a realisation is given by
$$
\Pi^\vee=\{H_1,H_0\},\qquad
\Pi=\{\alpha_0,\alpha_1\}
$$
$$
\begin{aligned}
\mathfrak{h}&=\mathbb{C}H\oplus\mathbb{C}
d\oplus\mathbb{C}K,\\
\mathfrak{h}^*&=\mathbb{C}
\alpha_1\oplus\mathbb{C}\delta
\oplus\mathbb{C}\Lambda_0
\end{aligned}
$$
(Canonical dual, $\left\langle{}\alpha_1,H\right\rangle{}=2$,
$\left\langle{}\delta,d\right\rangle{}=1=\left\langle{}\Lambda_0,
K\right\rangle{}$,
every other pairing $0$.)

Then
$$
\begin{cases}
\alpha_1=\alpha_1\\
\alpha_0=\delta-\alpha_1
\end{cases}\qquad
\begin{cases}
\alpha_1^\vee=H \\
\alpha_0^\vee = K-H
\end{cases}
$$
So we obtain
$$
\begin{aligned}
\left\langle{}\alpha_0^\vee,\alpha_1\right\rangle{}&=\left\langle{}K-
H,\alpha_1\right\rangle{}=2\\
\left\langle{}\alpha_1^\vee,\alpha_0\right\rangle{}&=\left\langle{}H,
\delta-\alpha_1\right\rangle{}=-2\\
\left\langle{}\alpha_0^\vee,\alpha_0\right\rangle{}&=\left\langle{}K-
H,\delta-\alpha_1\right\rangle{}=+2
\end{aligned}
$$

\end{proof}

\medskip\noindent
Finally let's define Kac-Moody algebras.

Let $A$ be a generalized Cartan matrix. Let
$$
\tilde{\mathfrak{n}}_+=F(e_1,\ldots,e_n),
$$
the free Lie algebra on $n$ generators, and
similarly
$$
\tilde{\mathfrak{n}}_-=F(f_1,\ldots,f_n).
$$
Let $\mathfrak{h}$ be a realisation of $A$.
Set $\tilde{\mathfrak{g}}(A)
=\tilde{\mathfrak{n}}_- \oplus \mathfrak{h}
\oplus \tilde{\mathfrak{n}}_+$.

Make $\tilde{\mathfrak{g}}(A)$ a Lie algebra
by defining
\begin{itemize}
\item $[\mathfrak{h},\mathfrak{h}]=0$,
\item $\forall  H \in \mathfrak{h}$,
$[H,e_i]=\left\langle{}\alpha_i,H\right\rangle{}e_i
=\alpha_i(H)e_i$. And similarly,
$[H,f_i]=-\alpha_i(H)f_i$.
\item $[e_i,f_j]=\delta_{ij}\alpha_i^\vee$.
\end{itemize}
Then $\tilde{\mathfrak{g}}(A)$ is a Lie
algebra (though not yet the Kac-Moody
algebra). See Kac, \cite[Thorem 1.2]{IDLA}.

\begin{remark}
\label{vertex-algebras-remark-lattice}
In $\mathfrak{h}$ we have a lattice
$$
\begin{aligned}
Q^\vee&=\mathbb{Z}\alpha_1^\vee+\ldots+
\mathbb{Z}\alpha_n^\vee,\quad
\text{and}\\
Q&=\mathbb{Z}\alpha_1+\ldots+\mathbb{Z}
\alpha_n\text{
in }\mathfrak{h}^*
\end{aligned}
$$
(root and coroot lattices).
$\tilde{\mathfrak{g}}(A)$ is naturally
$Q$-graded,
with
$$
\tilde{\mathfrak{g}}(A)_\beta=\text{span}\{
\text{commutators of $e_i$ with $\sum
\alpha_i=\beta$}\}.
$$
$\tilde{g}(A)=\mathfrak{h}$.

\begin{theorem}[Gabber-Kac]
\label{vertex-algebras-theorem-Gabber-Kac}
Denote by $I \subset \tilde{\mathfrak{g}}(A)$
the maximal $Q$-graded ideal, such
that $I \cap\mathfrak{h}=\{0\}$.
Then $I$ is generated by the Serre relations
$$
\text{ad}(e_i)^{1-a_{ij}}e_j\qquad
\text{and}\qquad
\text{ad}(f_i)^{1-a_{ij}}f_j,\; i\neq j.
$$
\end{theorem}

\begin{proof}
\cite[Theorem 9.11]{IDLA}.
\end{proof}

(The existence of the ideal $I$ does not need
the theorem;
the importance of the theorem is providing an
expression for the generators.)

\begin{definition}
\label{vertex-algebras-definition-Kac-Moody-algebra}
The {\it Kac-Moody algebra} $\mathfrak{g}(A)$
is $\tilde{\mathfrak{g}}(A)/I$.
\end{definition}

\end{remark}

\section{Affine Kac-Moody algebras}
\label{vertex-algebras-section-affine-Kac-Moody-algebras}

Let $\mathfrak{g}$ be a finite-dimensional
simple Lie algebra,
with $(\cdot,\cdot):\mathfrak{g} \times
\mathfrak{g} \to \mathbb{C}$
invariant bilinear form,
$$
([x,y],z)=(z,[y,z])\qquad \forall x,y,z \in
\mathfrak{g}
$$
(Eg. the Killing form
$\kappa(x,y)=\text{Tr}_\mathfrak{g}\text{ad}
(x)\text{ad}(y)$
is invariant.)

Typically we normalise $(\cdot,\cdot)$ so
that $(\alpha,\alpha)=2$ for the long
roots of $\mathfrak{g}$.

Then
$\hat{\mathfrak{g}}=\mathfrak{g}[t,t^{-1}]
\oplus\mathbb{C}K$
(affine Lie algebra),
$$
[at^m,bt^n]=[a,b]t^{m+n}+m\delta_{m,-n}(a,b)
K,\qquad
[K,\hat{\mathfrak{g}}]=0
$$
and
$\tilde{\mathfrak{g}}=\hat{\mathfrak{g}}
\oplus\mathbb{C}d$,
$[d,K]=0$,  $[d,at^m]=mat^m$,
(affine Kac-Moody algebra or ``extended
affine Lie algebra'')

\begin{theorem}
\label{vertex-algebras-theorem-gtilde-is-a-Kac-Moody-algebra}
$\tilde{\mathfrak{g}}$ is a Kac-Moody
algebra.
\end{theorem}

Let $\mathfrak{g}=\mathfrak{h}\oplus
\bigoplus_{\alpha \in
\Delta}\mathfrak{g}_\alpha$,
($\mathfrak{g}_\alpha=\mathbb{C}E_\alpha$.)

{\bf The simple roots and coroots.}
$\tilde{\mathfrak{h}}=\mathfrak{h}
\oplus\mathbb{C}K\oplus\mathbb{C}d$.
We identify $\tilde{\mathfrak{h}}^*$ with
 $\mathfrak{h}^*\oplus\mathbb{C}
 \Lambda_0\oplus\mathbb{C}\delta$
 where
$$
\begin{aligned}
\Lambda_0(\mathfrak{h})&=\delta(\mathfrak{h})
=0\\
\Lambda_0(d)&=\delta(K)=0\\
\Lambda_0(K)&=\delta(d)=1
\end{aligned}
$$
The {\it real coroots} are
$$
\hat{\Delta}^{V,re}=\{E_\alpha t^m|\alpha \in
\Delta, m \in \mathbb{Z}\}
$$
and there are also imaginary roots and
coroots
$$
\hat{\Delta}^{V,im}=\{Ht^m|H \in
\mathfrak{h}, m \in
\mathbb{Z}\setminus\{0\}\}
$$
Roots:
$$
\begin{aligned}
\hat{\Delta}^{re}&=\{\alpha+m\delta|\alpha
\in \Delta, m \in \mathbb{Z}\}\\
\hat{\Delta}^{im}&=\{m\delta|m \neq 0\}
\end{aligned}
$$
$Xt^m$:
$$
\begin{aligned}
[H,Xt^m]&=[H,x]t^m,\qquad H \in
\mathfrak{h}\\
[K,xt^m]&=0\\
[d,xt^m]&=m x t^m
\end{aligned}
$$
so it $x \in \mathfrak{g}_\alpha$, $x t^m \in
\tilde{\mathfrak{g}}_{\alpha+m\delta}$.

\medskip\noindent
The invariant bilinear form $(\cdot,\cdot)$
from $\mathfrak{g} \times
\mathfrak{g}$ extends uniquely to
$(\cdot,\cdot):\tilde{\mathfrak{g}}
\times\tilde{\mathfrak{g}}\to
\mathbb{C}$.

$(d,d)=(K,K)=0$, $(d,K)=1$ and
$(d,\mathfrak{h})=(K,\mathfrak{h})=0$.

\medskip\noindent
So, in $\tilde{\mathfrak{h}}^*$:
$$
\begin{aligned}
(\Lambda_0,\Lambda_0)&=(\delta,\delta)=0\\
(\Lambda_0,\mathfrak{h}^*)&=(\delta,
\mathfrak{h}^*)=0\\
(\Lambda_0,\delta)&=1.
\end{aligned}
$$
Hence, $|\alpha+m \delta|^2=|\alpha|^2$, $|m
\delta|^2=0$.

\begin{example}
\label{vertex-algebras-example-sl2hat}
$\widetilde{\mathfrak{sl}_2},
\tilde{\mathfrak{h}}^*
=\text{span}\{\alpha,\Lambda_0,\delta\}$ with
Gram matrix …
\end{example}

We can make a choice of positive roots,
$$
\hat{\Delta}_+=\{\alpha+m \delta |\alpha \in
\Delta, m>0\}
\cup \{m \delta|m>0\} \cup \Delta_+
$$
Obviously, if $\alpha \in \Delta_+$ is
simple,
$\alpha \in \hat{\Delta}_+$ is simple.

\medskip\noindent
{\bf Notation.} Let $\theta \in \Delta_+$ be
a the highest root. ($\not \exists
\alpha \in \Delta_+$ such that $\alpha-\theta
\in \mathbb{Z}_+ \Delta_+$.) and
 $\alpha=\delta-\theta$.

Then $\alpha_0 \in \hat{\Delta}_+$ is simple
and the set of simple roots is
$\hat{\Pi}=\{\alpha_0,\underbrace{\alpha_1,
\ldots,\alpha_\ell}
_{\substack{\text{the finite} \\ \text{simple
roots}}} \}.$
where $\ell=\text{rank}(\mathfrak{g})$.

The coroot corresponding to $\alpha_0$ is
$$
\alpha^\vee_0=K-\theta^\vee, \qquad
\theta^\vee=\frac{2}{(\theta,\theta)}\nu^{-1}
(\theta)
\in \mathfrak{h}
$$
$$
\text{ and }\qquad E_{\alpha_0}=E_{-\theta}t.
$$

\section{Weyl group}
\label{vertex-algebras-section-Weyl-group}

\medskip\noindent
{\bf Upshot.} The Weyl group is a semidirect
product of
pseudoreflections and translations.

\medskip\noindent
In any Kac-Moody algebra, we have
$$
\begin{aligned}
\text{roots }\qquad
\Pi&=\{\alpha_1,\ldots,\alpha_\ell\}
\subset\mathfrak{h}^*
\\
\text{coroots}\qquad
\Pi^\vee&=\{\alpha_1^\vee,\ldots,
\alpha_\ell^\vee\}\subset
\mathfrak{h},
\end{aligned}
$$
and {\it reflections} $r_i \in
\text{GL}(\mathfrak{h}^*)$,
defined by
$$
r_i(\lambda)=\lambda-\left\langle{}\lambda,
\alpha_i^\vee\right\rangle{}\alpha_i.
$$

One can check that
$$
(r_i\lambda,r_i\mu)=(\lambda,\mu)\qquad
\forall \lambda,\mu \in \mathfrak{h}^*
$$
The {\it Weyl group} $W$ is
$\left\langle{}r_i|i=1,\ldots,\ell\right\rangle{}\subset
\text{GL}(\mathfrak{h}^*)$.

\begin{example}
\label{vertex-algebras-example-Weyl-group-sl2tilde}
For $\widetilde{\mathfrak{sl}_2}$, $r_1$ is
easy,
$$
\begin{aligned}
r_1(\alpha)&=-\alpha \qquad \text{(as in
$\mathfrak{sl}_2$)}\\
r_1(\delta)&=\delta,\qquad
r_1(\Lambda_0)=\Lambda_0.
\end{aligned}
$$
To compute $r_0$ take an arbitrary element $m
\alpha_1+k\Lambda_0+f\delta$ and
do:
$$
\begin{aligned}
r_0(m
\alpha_1+k\Lambda_0+f\delta)&=m\alpha_1+
k\Lambda_0+f\delta
-\left\langle{}\alpha_0^\vee,m\alpha_1+k\Lambda_0+
f\delta\right\rangle{}\alpha_0\\
\alpha_0=\delta-\alpha_1,\qquad &
\alpha_0^\vee=K-\alpha^\vee
\end{aligned}
$$
so we obtain
$$
=m\alpha_1+k\Lambda_0+f\delta-(k-2m)(\delta-
\alpha_1)
$$
$$
(k-m)\alpha_1+k\Lambda_0+(f-k+2m)\delta.
$$
Relative to basis
$\{\alpha_1,\Lambda_0,\delta\}$.
$$
r_1=\begin{pmatrix}
-1&0&0\\
0&1&0\\
0&0&1
\end{pmatrix},
r_0=\begin{pmatrix}
-1&1&0\\
0&1&0\\
2&-1&1
\end{pmatrix}
\begin{pmatrix}
m\\
k\\
f
\end{pmatrix}
$$
$$
t=r_1r_0=\begin{pmatrix}
1&-1&0\\
0&1&0\\
2&-1&1
\end{pmatrix}
$$
Notice that $\delta$ is fixed by all $r_i$.
Also
$m \alpha+k \Lambda_0 +f\delta$, the {\it
coefficient} of $\Lambda_0$ is fixed
by all $r_i$.

Then
$$
\begin{aligned}
t(m\alpha_1+&k\Lambda_0+f\delta)=(m-k)
\alpha_1+k\Lambda_0+(f-k+2m)\delta.
\end{aligned}
$$
Think of $t$ as a translation.

The number $k$ in
 $$
\mathfrak{h}^* \ni
\hat{\lambda}=\lambda+k\Lambda_0+f\delta
$$
is called the {\it level} of $\hat{\lambda}$.

$\hat{\mathfrak{h}}=$ union of (hyper)planes
of constant level which are stable
under $W$. The roots $\alpha$ are all of
level $0$.

[Picture]
``$r_1$ changes the sign of the finite
path''. And $t=r_1r_0$ is a sort of
translation. Indeed, in general we can
consider
$t_{\alpha_i}=r_{\alpha_i}\circ r_0 \in W$,

$$
t_\alpha(\beta+m\delta)=\beta+(m+(\beta,
\alpha_i))\delta
$$
\end{example}
One can describe the action of $t_\alpha$ on
$\hat{\lambda}$ in general
(e.g. see \cite[Chapter 6]{IDLA})

\begin{proposition}
\label{vertex-algebras-proposition-weyl-group-affin-kac}
For the affine Kac-Moody algebra
$\hat{\mathfrak{g}}$
 with
 $\hat{W}=\left\langle{}r_0,r_1,\ldots,r_\ell\right\rangle{}$
 its Weyl group
(and $W=\left\langle{}r_1,\ldots,r_\ell\right\rangle{}\subset
\hat{W}$ the Weyl group of
$\mathfrak{g}$), then $\hat{W} \simeq W\times
t_{Q^\vee}$ (where it should be
semidirect product instead of $\times$…)
where $Q^\vee$ is the
coroot lattice of $\mathfrak{g}$.
\end{proposition}

\begin{remark}
\label{vertex-algebras-remark-weyl-group-large-other-kac}
For general Kac-Moody algebras, the Weyl
groups are much larger, hyperbolic
reflection groups.
\end{remark}

In the affine case, $\hat{W}$ fixes level
$k$, and $|\hat{\lambda}|$. One gets,
in the intersection, paraboloids [Picture of
section of hyperboloid that is a
parabola].

\section{Weyl character formula}
\label{vertex-algebras-section-Weyl-character-formula}
Highest weight representations of Kac-Moody
algebras.
Let $\lambda \in \mathfrak{h}^*$, where
$\mathfrak{g}(A)=\mathfrak{n}_-
\oplus\mathfrak{h}\oplus
\mathfrak{n}_+$
is a Kac-Moody algebra. We define a {\it
Verma module}
$$
M(\Lambda)
=U(\mathfrak{g})
\otimes_{U(\mathfrak{h}+\mathfrak{n}_+)}
\mathbb{C}v_\Lambda
$$
where $\mathfrak{h}+\mathfrak{n}_+$ acts on
$V_\Lambda$ by:
\begin{align*}
\begin{aligned}
Xv_\Lambda&=0,\qquad &\forall x &\in
\mathfrak{n}_+,\\
Hv_\Lambda&=\Lambda(H)v_\Lambda,\qquad
&\forall H &\in \mathfrak{h}
\end{aligned}
\end{align*}
So $\mathbb{C} v_\Lambda$ is a
$U(\mathfrak{h}+\mathfrak{n}_+)$-module,
$$
\xymatrix{
U(\mathfrak{h}+\mathfrak{n}_+)\ar[d]\\
U(\mathfrak{g})
}
$$
By the PBW theorem, $M(\Lambda)$ has a linear
$\mathbb{C}$-basis.

Let $\{F_{\alpha,i}:i = 1,\ldots, \dim
\mathfrak{g}_\alpha\}$ be a basis of
$\mathfrak{g}_{-\alpha}$, $\forall \alpha \in
\Delta_+$.
Also choose a total order on $\Delta_+$.
(Some sort of lexicographical order
that takes longer to write than to say.)
$$
F_{\alpha_1,i_1},F_{\alpha_2,i_2},\ldots,
F_{\alpha_s,i_s},v_\Lambda
$$
$$
\alpha_1\leq \alpha_2\leq \ldots \leq
\alpha_2\text{ and if }
\alpha_p=\alpha_{p+1}, i_p \leq  i_{p+1}
$$
We have
$M(\Lambda)_\lambda=\{m|Hm=\lambda(H)m\}$
weight spaces.
$$
M(\Lambda)=\bigoplus_{\lambda \in
\mathfrak{h}^*}M(\Lambda)_\lambda
$$
The vector $v_{\Lambda}$ is in
$M(\Lambda)_\Lambda$ by definition,
$$
\begin{aligned}
F_{\alpha,i}V_\Lambda &\in
M(\Lambda)_{\Lambda-\alpha}\\
H(Fv_\Lambda)&=\underbrace{[H,F]v_\Lambda}_{=
-\alpha(H)Fv_\Lambda}
+\underbrace{FHv_\Lambda}_{=\Lambda(H)
FV_\Lambda}
\end{aligned}
$$
So $\chi_{M(\Lambda)}=\sum_{\lambda \in
\mathfrak{h}^*}\dim M(\Lambda)_\lambda
e^{\lambda}$
is computed by counting monomials $y$ with
fixed $\sum_{i}\alpha_i$.

\begin{equation}
\label{vertex-algebras-equation-character-of-Verma-module}
\chi_{M(\Lambda)}=e^{\Lambda}\prod_{\alpha
\in \Delta_+}
\frac{1}{(1-e^{-\alpha})^{\dim
\mathfrak{g}_\alpha}}.
\end{equation}
The product on Eq.
\ref{vertex-algebras-equation-character-of-Verma-module} is
called
{\it Weyl denominator}.
\begin{exercise}
\label{vertex-algebras-exercise-convince-yourself}
Convince yourself of this.
\end{exercise}

\begin{example}
\label{vertex-algebras-example-sl2-character}
$\mathfrak{g}=\mathfrak{sl}_2$, [Picture]
$$
\begin{aligned}
\chi_{M(\Lambda)}&=e^{\Lambda}+e^{\Lambda-
\alpha}+e^{\Lambda-2\alpha}+\ldots\\
&=e^{\Lambda}(1+e^{-\alpha}+e^{-2\alpha}+
\ldots\\
&=e^{\Lambda}\frac{1}{1-e^{-\alpha}}.
\end{aligned}
$$
\end{example}

\medskip\noindent
For certain $\Lambda$, $M(\Lambda)$ is {\it
reducible} (i.e. there exists
a submodule $0 \neq N \subset M(\Lambda)$
(with proper contention).

\begin{lemma}
\label{vertex-algebras-lemma}
For any submodule $N$,
$$
N=\bigoplus_{\mu \in \mathfrak{h}^*}N \cap
M(\Lambda)_\mu.
$$
\end{lemma}

{\bf Corollary.} The sum of all
\underline{proper} submodules
of  $M(\Lambda)$ is proper, in particular
there is a maximal proper submodule.

{\bf Notation.}
$L(\Lambda)=M(\Lambda)\Big/\left(\substack{
\text{max. proper} \\
\text{submodule}}\right)$

\begin{example}
\label{vertex-algebras-example}
$\mathfrak{sl}_2$. $\Lambda=3\omega$
($\omega$: fundamental weight,
$\alpha=2\omega$.)
$L(3\omega)=\mathbb{C}\left\langle{}e^{3\omega},e^{-
\omega},e^{-3\omega}\right\rangle{}$.

[Picture]
\end{example}

\begin{definition}
\label{vertex-algebras-definition-integrable-g-module}
A $\mathfrak{g}$-module is {\it integrable}
if
\begin{itemize}
\item $V=\bigoplus_{\mu \in
\mathfrak{h}^*}V_\mu$ (weight module).
\item For all simple roots $\alpha_i$; $e_i$
and $f_i$ are
locally nilpotent on $V$ (i.e. for all $v \in
V$
there exists $N$ such that
$e_i^Nv=f_i^Nv=0$.)
\end{itemize}
\end{definition}

\begin{remark}
\label{vertex-algebras-remark-integrability}
\begin{itemize}
\item Vermas are not integrable.
\item $\dim V< \infty \implies V$ integrable.
\item $\mathfrak{g}$ itself (Kac-Moody) is
integrable.
\end{itemize}
\end{remark}

\medskip\noindent
{\bf Dominant integrable weights.}
Let
$\{\alpha_1^\vee,\ldots,\alpha_\ell^\vee\}
\subset\mathfrak{h}$
be the simple
coroots.

\begin{definition}
\label{vertex-algebras-definition-dominant-integral-weights}
The {\it dominant integral weights} are the
weights that pair with the coroots
to give integers:
$$
P_+=\{\lambda \in \mathfrak{h}^* :
\left\langle{}\lambda,
\alpha_i^\vee\right\rangle{}\in\mathbb{Z}_{\geq
0,1},
i=1,\ldots,\ell\}.
$$
\end{definition}

For $L(\Lambda)$ to be integrable, it is
necessary that $\Lambda \in P_+$.

Indeed, suppose $L(\Lambda)$ is integrable.
Then $f_i^Nv_\Lambda=0$ in $L(\Lambda)$,
or rather
$$
\underbrace{e_if_i^{N+1}v_\Lambda}
_{Kf_i^N=0} \in M(\Lambda),
$$
and $K$ can only be zero if
$\left\langle{}\Lambda,\alpha_i^\vee\right\rangle{}
\in \mathbb{Z}_{\geq 0}$. Applying for all
$i$, we find
$\Lambda \in P_+$ is necessary.

\begin{proposition}
\label{vertex-algebras-proposition-integ-iff-domin-integ}
$L(\Lambda)$ is integrable if and only if
$\Lambda \in P_+$.
\end{proposition}

\begin{proof}
For the converse, use induction and Serre
relations
(we know the result for the highest weight,
and want
to prove for others).
\end{proof}

\begin{example}
\label{vertex-algebras-example-sl3-integrability}
$\mathfrak{sl}_3$. [Picture]
\end{example}

\begin{example}
\label{vertex-algebras-example-sl-2hat-integrability}
$\widehat{\mathfrak{sl}_2}$. [Picture, $P_+$
looks like diagonal lines.]
$$
\alpha_0^\vee=K-H,\qquad
\alpha_1^\vee=H\in\mathfrak{sl}_2,
\qquad
\left\langle{}\delta,\alpha_i^\vee\right\rangle{}=0,\quad
i=0,1
$$
\end{example}

\begin{remark}
\label{vertex-algebras-remark-verma-affin-kac-moody-algeb}
For affine Kac-Moody algebras, almost nothing
about the structure of
$M(\Lambda)$ depends on the coefficient of
$\delta$ in $\Lambda$.
So it's common to consider
$$
M(\Lambda)=M_k(\lambda)=M(k\Lambda_0+\lambda)
,\qquad
\lambda \in \mathfrak{h}^*
$$
where $k$, the level of $\Lambda$, is super
important.
\end{remark}

Then
$$
\underbrace{\hat{P}_+}_{
\substack{\delta\text{-coef.} \\ =0}}
=\bigcup_{k \in \mathbb{Z}_{\geq 0}}
\{k\Lambda_0+\lambda|\lambda \in P_+^k\}
\{k\Lambda_0+\lambda|\lambda \in P_+^k\}.
$$
$$
P_+^k=\{\lambda \in P_+|
\left\langle{}\lambda,\theta\right\rangle{}\leq k\}\subset
P_+ \text{ for }\mathfrak{g}.
$$

\medskip\noindent
Consider $V= \bigoplus_{\mu \in
\mathfrak{h}^*}V_\mu$ integrable,
and $V_\lambda \neq 0$.
Let $i \in \{1,\ldots,\ell\}$.
Consider $U=\bigoplus_{n \in
\mathbb{Z}}V_{\lambda+n\alpha_i}\subset V$,
and the action of $e_i,f_i$ and
$h_i=[e_i,f_i]$.

The action of $\mathfrak{sl}_2$ on $U$ is
locally integrable. By structure of
$\mathfrak{sl}_2$-represenations, $U$ must be
finite-dimensional
with ``symmetrical'' weight space
multiplicities, i.e.,
$$
\begin{aligned}
\{\lambda+n\alpha_i|n\in\mathbb{Z}\}\cap
\{\text{weights of }V\}
&=\{\lambda+n\alpha_i|-p\leq n\leq q\}.
\end{aligned}
$$
and
$$
\left\langle{}\lambda-p\alpha_i,h_i\right\rangle{}=-
\left\langle{}\lambda+q\alpha_i,h_i\right\rangle{}.
$$
\medskip\noindent
Consequently, the reflection
$r_i(\lambda):=\lambda-\left\langle{}\lambda,
\alpha_i^\vee\right\rangle{}\alpha_i$
has the same multiplicity as $\lambda$.

[Picture]

Now consider $M(\lambda)$ and $L(\Lambda)$
for
$\Lambda \in P_+$.
Actually, for general $\Lambda$,
$M(\Lambda)$, while not
necessarilly irreducible, has an
($\Omega$-)composition series$^*$ by
irreducibles $L(\lambda)$.
$$
M(\Lambda)=V_0 \supset V_1 \supset V_2
\supset \ldots \supset V_n=0,
$$
such that $V_i/V_{i+1}$ is $\simeq
L(\lambda_i)$ (i.e. an irreducible highest
weight module) for some
$\lambda_i=\Lambda-\beta_i$, $\beta_i \in Q$.
For
Kac-Moody algebras we also consider the case
that $(V_i/V_{i+1})=0$ for all
$\mu \in \Omega + Q_+$, (which is the
$\Omega$-composition series).

[Picture.]

For an $\Omega$-composition series, as above,
we have
$$
\text{ch}_{M(\Lambda)}-\sum_{\lambda \geq
\Omega}
\underbrace{[M(\Lambda):L(\lambda)]}
_{\substack{\text{\# of times $L(\lambda)$}
\\ \text{
appears in the compos. series}}}
\text{ch}_{L(\lambda)}
\in \left\langle{}e^{\mu}:\mu \not \geq \Omega\right\rangle{}
$$
Sending  ``$\Omega \to -\infty$'', the
identity
$$
\text{ch}_{M(\Lambda)}=\sum_{\lambda \leq
\Lambda}
[M(\Lambda):L(\lambda)]\text{ch}_{L(\lambda)}
$$
makes sense.

{\bf Notation.}
$b_{\Lambda,\lambda}=[M(\Lambda):L(\lambda)]
$.

\begin{remark}
\label{vertex-algebras-remark-total-order-on-hstar}
Recall the partial order on weights that
$\lambda \leq \Lambda$ if
$\Lambda-\lambda \in Q= \sum
\mathbb{Z}_+\alpha_i$.
$b_{\Lambda,\lambda}=1$ if $\lambda =
\Lambda$ and
$b_{\Lambda,\lambda}=0$ if not ($\lambda \leq
\Lambda$).

If we choose a total order on
$\mathfrak{h}^*$, compatible with $\leq$.
Then $\{b_{\Lambda,\lambda}\}$ is a lower
triangular matrix with 1
on the diagonal.
We an define $\{m_{\Lambda,\lambda}\}$ the
{\it inverse matrix}.
It's again lower triangular, 1 on the
diagonal, and all
$m_{\Lambda,\lambda}$ are
\underline{integers} (maybe negative now).
And we have
\begin{equation}
\label{vertex-algebras-equation-character}
\text{ch}_{L(\Lambda)}=\sum_{\lambda \leq
\Lambda}m_{\Lambda,\lambda}
\text{ch}_{M(\lambda)}.
\end{equation}
\end{remark}

\begin{example}
\label{vertex-algebras-example-sl2-characters}
$\mathfrak{sl}_2$.
$M(3)\underset{L(3)}{\supset}M(-5)
\underset{L(-5)}{\supset}0$. Since $M(-5)$ is
already irreducible.
[Missing…]
\end{example}

We want to discover $m_{\Lambda,\lambda}$.
In general massively difficult.
For $\Lambda \in P_+$, $m_{\Lambda,\lambda}$
easy.
Multiply Eq. \ref{vertex-algebras-equation-character} by  $R$
$$
\begin{aligned}
R\text{ch}_{L(\Lambda)}&=\sum_{\lambda \geq
\Lambda} m_{\Lambda,\lambda}
e^{\lambda}.
e^{\rho}R\text{ch}_{L(\Lambda)}=\sum_{\lambda
\leq  \Lambda}
m_{\Lambda,\lambda}e^{\lambda+\rho}.
\end{aligned}
$$
What's $\rho$? It's $\rho \in \mathfrak{h}^*$
chosen so that
$\left\langle{}\rho,\alpha_i^\vee\right\rangle{}=1$, and it's
called
the {\it Weyl vector}.

\begin{remark}
\label{vertex-algebras-remark-finite-dimension}
For $\mathfrak{g}$ finite-dimensional,
$\rho=\sum_{i=1}^\ell \omega_i$
necessarily. (And equals
$\frac{1}{2}\sum_{\alpha \in
\Delta_+}\alpha$.

For $\mathfrak{g}$ finite dimensional and the
affine $\hat{\mathfrak{g}}$,
$\hat{\rho}=h^\vee\Lambda_0+\rho$ works.
\end{remark}

For $w \in
W=\left\langle{}r_i|i=1,\ldots,\ell\right\rangle{}$,
define $w(\text{ch}_V)=\sum \dim V_\mu
e^{W(\mu)}$.
We saw $w(\text{ch}_V)=\text{ch}_V$ if $V$
integrable.
In particular
$w(\text{ch}_{L(\Lambda)}=\text{ch}
_{L(\Lambda)}$,
$\Lambda \in P_+$.

\begin{lemma}
\label{vertex-algebras-lemma-for-Weyl-formula}
$m_{\Lambda,\lambda}=0$ unless
$\lambda+\rho=w(\Lambda+\rho)$ for some $w
\in W$
\end{lemma}

\medskip\noindent
{\bf Claim.} $r_i(e^{\rho} R)=-e^\rho R$.
So $w(e^\rho R)=\det(w)e^\rho R$ for all $w
\in W$.

\begin{proof}
$$
R=\prod_{\alpha \in
\Delta_+}(1-e^{-\alpha})^{\text{mult}(\alpha)
}.
$$
Note that
\begin{enumerate}
\item
$\text{mult}(r_i(\alpha))=\text{mult}(\alpha)
$
for all $\alpha \in \Delta$. (Since
$\mathfrak{g}$ is integrable!)
\item $r_i(\Delta_+)=\{-\alpha_i\}\cup
(\Delta_+\setminus\{\alpha_i\}$.
\end{enumerate}
Any $\alpha \in \Delta_+$ is of the form
$\alpha = \sum_{i=1}^\ell k_i \alpha_i$.
If $\alpha \neq \alpha_i$, some $k_{j_0}\neq
0$, $j_0 \neq  i$ and
$$
\begin{aligned}
r_i(\alpha)&=\sum_{j}k_j \alpha_j
\left\langle{}\alpha,\alpha_i^\vee\right\rangle{}\alpha_i\\
&=\sum k_j' \alpha_j\\
&=e^\rho(e^{-\alpha_i}-1)
\left(\prod_{\alpha \in \Delta_+ \setminus
\alpha_i}\right)\\
&=-e^\rho R.
\end{aligned}
$$
for $k_{j_0}'=k_{j_0}>0$. Can't have a
mixture of signs,
so $r_i(\alpha) \in \Delta_{\perp}$.

So
$$
\begin{aligned}
r_i(e^\rho R)&=\prod_{\alpha \in
\Delta_+\setminus \alpha_i}
(1-e^{-\alpha})^{\text{mult}(\alpha)}
\cdot (1-e^{t \alpha_i}\cdot
e^{\rho-\alpha_i}
\end{aligned}
$$
$$
r_i=\rho-\left\langle{}\alpha_i^\vee,
\rho\right\rangle{}\alpha_i=\rho-\alpha_i
$$
Hence
$$
\sum_{\lambda \leq
\Lambda}m_{\Lambda,\lambda}e^{\lambda+\rho}
$$
is $W$-skew-invariant. (Lemma
\ref{vertex-algebras-lemma-for-Weyl-formula}.)

Hence
$$
\sum_{\lambda \leq
\Lambda}m_{\Lambda,\lambda}e^{\lambda+\rho}
=\sum_{w \in W}\det(w)\cdot
e^{w(\Lambda+\rho)}
$$
In conclusion, the {\it Weyl character
formula}
$$
\text{ch}_{L(\Lambda)}=\sum_{w \in W}\det(w)
\cdot \frac{e^{w(\Lambda+\rho)-\rho}}{R}
$$
\end{proof}

\medskip\noindent
{\bf Corollary (Weyl denominator formula).}
$$
e^{\rho}R=\sum_{w \in W}\det(w)e^{w(\rho)}.
$$
Next time: Affine case, $\theta$-functions.
Modular forms (Poisson summation.)

\section{Characters of integrable highest
weight for affine Kac-Moody algebras}
\label{vertex-algebras-section-chara-integ-highe-weigh}

Can be calculated using the Weyl character
formula.

Recall:
$\hat{\mathfrak{g}}=\mathfrak{g}[t,t^{-1}]
\oplus \mathbb{C}K
\oplus \mathbb{C}d$.

Dual Cartan
$\hat{\mathfrak{h}}^*=\mathfrak{h}
^*\oplus\mathbb{C}\Lambda_0
\oplus\mathbb{C}\delta$.

We consider weights
$\Lambda=k\Lambda_0+\lambda$, $\lambda \in
\mathfrak{h}^*$,
$k \in \mathbb{Z}_+$, $\lambda \in P_+^k$.

Simple roots: $\alpha_0=\delta-\theta$,
$\{\alpha_1,\ldots,\alpha_\ell\}
\subset \mathfrak{h}^*$.

The affine Weyl group is (the {\it
semidirect} product!)
$$
\widehat{W}=\left\langle{}r_0,\ldots,r_\ell\right\rangle{}
\cong W \times T
$$
where $T=\{t_\alpha | \alpha \in Q\}$, where
$t_\alpha$ is the translation
$$
t_\alpha(\Lambda)=\Lambda+k\alpha-
\left((\lambda,a)+k
\frac{|\alpha|^2}{2}\right)\delta
$$
$\Lambda=k\Lambda_0+\lambda$.

\medskip\noindent
Let's compute $\chi_{L(\Lambda_0)}$ using the
Weyl character formula
$$
\chi_{L(\Lambda)}=\frac{\sum_{w \in
\tilde{W}}\varepsilon(w)
e^{W(\Lambda+\rho)-\rho}}{R}
$$
Firstly $R=\prod_{\alpha \in
\hat{\Delta}}(1-e^{-\alpha})
^{\text{mult}(\alpha)}$

Recall
$$
\begin{aligned}
\hat{\Delta}_+&=\{m\delta | m \in
\mathbb{Z}_{\geq 1}\}
\cup \{\alpha_1+m\delta|m \in
\mathbb{Z}_{\geq 0}\}
\cup \{-\alpha_1+m \delta |m \in
\mathbb{Z}_{\geq 1}\}
\end{aligned}
$$
So let's write $e^{m \delta}=q^m$ (i.e. $q =
e^{-\delta}$, is a symbol)
and $y=e^{\alpha_1}$.

So
$$
R=\prod_{n=1}^{\infty}\underbrace{(1-y^{-1}
q^{n-1})
}_{\alpha+(n-1)\delta
}\underbrace{(1-q^n)
}_{n\delta}\underbrace{(1-yq^n)
}_{-\alpha+n\delta}
$$
Now let's express numerator also in terms of
$q$ and $y$

$$
\hat{\rho}=h^\vee\Lambda_0+\rho
$$
$h^\vee$ dual Coxeter number for
$\mathfrak{g}$.
For $\mathfrak{g}=\mathfrak{sl}_2$,
$h^\vee=2$.
(For $\mathfrak{g}=\mathfrak{sl}_n$,
$h^\vee=n$ and $\mathfrak{g}=E_8$,
$h^\vee=30$.)

For $\mathfrak{sl}_2$,
$\rho=\frac{1}{2}\alpha_1=\omega_1$.

$\Lambda=\Lambda_0$,
$\Lambda+\rho=3\Lambda_0+\omega_1$.

Using the formula, we find
$$
\begin{aligned}
&t_{m
\alpha_1}(3\Lambda_0+\omega_1)-(\Lambda_0+
\omega_1)=\\
&\ldots,\Lambda_0-3\alpha,-2\delta,\Lambda_0,
\Lambda_0+3\alpha_1-4\delta,\ldots
\end{aligned}
$$

[Picture]

$\hat{W}=T \cup  T\sigma,$ $\sigma=r_1$
finite reflection.

\medskip\noindent
{\bf Notation.} $w(\Lambda+\rho):=w \circ
\Lambda.$

One finds
$$
\begin{aligned}
\sum_{w \in \widehat{W}}\varepsilon(w)e^{w
\circ \Lambda_0}&=
e^{\Lambda_0}\left(
\underbrace{1+y^3q^4+y^{-3}q^2+\ldots}
_{\text{+
signs because $w \in T$}}
\underbrace{-y^{-1}-y^2q^2-\ldots}_{\text{$-$
signs because $w \in T_\sigma$}}
\right)
\end{aligned}
$$

We find explicitly
$$
\begin{aligned}
\chi_{L(\Lambda_0)}&=e^{\Lambda_0}
\frac{\sum_{m \in \mathbb{Z}}
y^{3m}q^{3m^2+m}-\sum_{m \in \mathbb{Z}}
y^{3m-1}q^{3m^2-m}}
{\prod_{n=1}^\infty(1-y^{-1}q^{n-1})(1-q^n)
(1-yq^n)}
\end{aligned}
$$

\begin{exercise}
\label{vertex-algebras-exercise-mathematica}
Put this formula in mathematica and confirm
that [Picture]
$$
\begin{aligned}
\chi_{L(\Lambda_0)}&=e^{\Lambda_0}
\left(1+q(y^{-1}+1+y)+q^2(y^{-1}+2+y)
+(q^3(2y^{-1}+3+2y)+\ldots\right)
\end{aligned}
$$
\end{exercise}

Appears that the central column here are {\it
partitions}, i.e.
$p(n)=$ \# of partitions of $n$. To see why
recall the generating function
$\sum_{n=0}^\infty
q^np(n)=\prod_{k=1}^\infty\frac{1}{1-q^k}$,
and expand.

It would appear that
$$
\chi_{L(\Lambda_0)}=
\prod_{k=1}^\infty\frac{1}{1-q^k}\cdot
\sum_{m \in \mathbb{Z}}y^mq^{m^2}
$$
This identity is true, and we will see it has
a vertex-algebra interpretation.

\begin{remark}
\label{vertex-algebras-remark-Jacobi-triple-product-identity}
If we compute $L(0)=1$ using the formula, we
obtain
$$
\begin{aligned}
\prod_{n=1}^\infty(1-yq^{n-1})(1-q^n)(1-y^{-
1}q^n)=\sum(\text{exercise})
\end{aligned}
$$
This identity is called the {\it Jacobi
triple product identity}.
\end{remark}

These functions
$$
\sum_{n \in \mathbb{Z}}q^{n^2}, \sum_{n \in
\mathbb{Z}}y^n q^{n^2},
\sum_{n \in
\mathbb{Z}}y^{3n}q^{3n^2+n},\text{ etc…}
$$
are all examples of {\it $\theta$-functions}.

\begin{remark}
\label{vertex-algebras-remark-duplicate-2}
In the formula for
$\chi_{L(\Lambda_0)}(y,q)$,
we could put $y=1$ to get
$$
\begin{aligned}
\chi_{L(\Lambda_0)}(q)&=1+3q+4q^2+7q^3+
\ldots\\
&=\prod\frac{1}{1-q^k}\cdot \sum_{m \in
\mathbb{Z}}q^{m^2}
\end{aligned}
$$
If one looks at $L(\Lambda_0+\omega_1)$ (the
other $\Lambda \in P_+^1$),
$$
\chi_{L(\Lambda+\omega_1)}=\prod
\frac{1}{1-q^k}
\sum_{m \in \mathbb{Z}}q^{(m+1/2)^2-1/4}
$$
\end{remark}

\section{Theta-functions}
\label{vertex-algebras-section-theta-functions}

\noindent
Let's consider
$$
\theta(\tau)=\sum_{n \in
\mathbb{Z}}q^{n^2/2},\qquad q=e^{2\pi i \tau}
$$
This converges (absolutely on compact regions
in)
the domain $\text{Im}(\tau)>0$.

Consider the Fourier transform
$$
\hat{g}(y)=\int_{-\infty}^\infty g(x)e^{2\pi
ixy}dx
$$

\begin{theorem}[Poisson summation]
\label{vertex-algebras-theorem-Poisson-summation}
 $$
\sum_{n \in \mathbb{Z}}g(n)=\sum_{n \in
\mathbb{Z}}\hat{g}(n)
$$
\end{theorem}

Let's take $g(x,t)=e^{-\pi tx^2}$.
Then $\theta(it)=\sum_{n\in
\mathbb{Z}}g(n,t)$.

In this case
$$
\hat{g}(y)=\sqrt{t}e^{-\pi y^2/t}
$$
(integral of Gaussian).

So we conclude that
$$
\theta\left(-\frac{1}{\tau}\right)=
\sqrt{\frac{\tau}{i}}\theta(\tau)
$$
Note also that
$$
\theta(\tau+2)=\theta(\tau)
$$
because $q^{\frac{1}{2}}=e^{\pi i\tau}$.

So $\theta(\tau)$ is an example of a modular
for.

\medskip\noindent
What is a modular form? Let
$$
\text{SL}_2(\mathbb{Z})=\left\{
\begin{pmatrix}
a&b\\
c&d
\end{pmatrix}|a,b,c,d \in \mathbb{Z}, \det =1
\right\},
$$
which acts on $\mathbb{H}=\{\tau \in
\mathbb{C}|\text{Im}(\tau)>0\}$ by
 $$
\begin{pmatrix}
a&b\\
c&d
\end{pmatrix}\cdot\tau=\frac{a\tau+b}{c\tau+
d}.
$$
\begin{definition}
\label{vertex-algebras-definition-weak-modular-form}
Let $\Gamma=\text{SL}_2(\mathbb{Z})$ or some
subgroup.
A {\it (weak) modular form (of weight $k$)}
is $f:\mathbb{H}\to \mathbb{C}$
holomorphic such that
\begin{equation}
\label{vertex-algebras-equation-weak-modular-form}
f\left(\frac{a\tau+b}{c\tau+d}\right)=(c\tau+
d)^kf(\tau).
\end{equation}
If we demand that $\Gamma$
 contains $\begin{pmatrix}
1&N\\
0&1
\end{pmatrix}$ for some $N\geq 1$,
(so $f(\tau+N)=f(\tau)$, and so $f(\tau)
=\sum_{n=-\infty}^\infty a(n)q^{2\pi in/2}$),
for $f:\mathbb{H}\to
\mathbb{C}$ satisfying Eq.
\ref{vertex-algebras-equation-weak-modular-form} and such
that
 $f(\tau)$ is ``meromorphic at cusps'', i.e.
$f(\tau)=\sum_{n \geq N_0}a(n)q^{n/N}$, etc.

Finally we would add a factor
$$
f\left(\frac{a\tau+b}{c\tau+d}\right)=
\varepsilon(A)(c\tau+d)^kf(\tau).
$$
\end{definition}

So in particular a weight-0 modular form is a
modular function.

Denote $S=\begin{pmatrix}
0&-1\\
1&0
\end{pmatrix}$. So $S\tau=-\frac{1}{\tau}$,
and denote
$T=\begin{pmatrix}
1&1\\
0&1
\end{pmatrix}$, so $T\tau = \tau+1$.

Then $\Gamma=\left\langle{}S,T^2\right\rangle{}\subset
\text{SL}_2(\mathbb{Z})$.

We saw
$$
\begin{aligned}
\theta(T^2\tau)&=\theta(\tau)\\
\theta(S\tau)&=\sqrt{\frac{\tau}{i}}
\theta(\tau)=i^{-1/2}
(c\tau+d)^{1/2}\theta(\tau).
\end{aligned}
$$
So $\theta(\tau)$ is a modular form of weight
$1/2$, for the group
$\Gamma=\left\langle{}S,T^2\right\rangle{}$ with multiplier
system $\varepsilon:\Gamma \to
\mathbb{C}^\times$ defined by
$$
\begin{aligned}
\varepsilon(S)&=i^{-1/2}\\
\varepsilon(T^2)&=1
\end{aligned}
$$

\begin{theorem}
\label{vertex-algebras-theorem-dedekind-eta-function}
The {\it Dedekind eta function}
$$
\eta(\tau)=q^{1/24}\prod_{n=1}^\infty
(1-q^n),\qquad q=e^{2\pi i \tau}
$$
is also a modular form of weight $1/2$ for
$\text{SL}_2(\mathbb{Z})=\left\langle{}S,T\right\rangle{}$.
$$
\begin{aligned}
\eta\left(-\frac{1}{\tau}\right)&=
\sqrt{\frac{\tau}{i}}\eta(\tau)\\
\eta(T\tau)&=e^{2\pi i/24}\eta(\tau).
\end{aligned}
$$

\end{theorem}

\section{Vertex algebras}
\label{vertex-algebras-section-vertex-algebras}

\begin{enumerate}
\item Recall
$\hat{\mathfrak{a}}=\mathbb{C}K\oplus
\bigoplus_{n \in \mathbb{Z}}\mathbb{C}h_n$
the {\it Heisenberg Lie algebra},
or {\it oscillator Lie algebra},
$$
[h_m,h_n]=m\delta_{m,-n}K,\qquad
[K,\hat{\mathfrak{a}}]=0.
$$
\begin{remark}
\label{vertex-algebras-remark-Heisenberg-is-affine}
``It's the simplest case of an affine Lie
algebra: a 1-dimensional
Lie algebra with a bilinear form''. It's an
example of the affine
Lie algebra construction
$\mathfrak{g} \rightsquigarrow
\hat{\mathfrak{g}}$,
where now $\mathfrak{a}=\mathbb{C}$,
and
 $$
\begin{aligned}
(\cdot,\cdot): \mathfrak{a}\times\mathfrak{a}
&\longrightarrow \mathbb{C} \\
(1,1) &\longmapsto 1
\end{aligned}
$$
\end{remark}

\item Witt Lie algebra
$$
W=\bigoplus_{n \in \mathbb{Z}}\mathbb{C}D_n,
$$
$$
[D_m,D_n]=(m-n)D_{m+n}.
$$
\end{enumerate}

\medskip\noindent
The oscillator algebra $\hat{\mathfrak{a}}$
has a representation which we have seen
already:
$$
H=\mathbb{C}[x_1,x_2,\ldots]
$$

\begin{align*}
h_n&\mapsto \begin{cases}
n \frac{\partial }{\partial x_n}\qquad
&\text{if }n>0\\
x_{-n}\qquad &\text{if }n<0\\
0\qquad &\text{if }n=0
\end{cases},
\qquad \qquad  K\mapsto \text{Id}
\end{align*}
A picture of $H$ [Picture].

Here I am introducing a \underline{grading}
of $H$:
$$
H=\bigoplus_{n \in \mathbb{Z}_{\geq 0}}H_n,
\qquad
H_n=\text{span}\{x_{m_1},\ldots,x_{m_s}|
\sum m_s=n\}.
$$

\begin{remark}
\label{vertex-algebras-remark-dimen-numbe-parti}
$\dim(H_n)=$\# $\{\text{integer partitions of
}n\}=p(n)$
and $\sum_{n=0}^\infty
\dim(H_n)q^n=\prod_{k=1}^\infty
\frac{1}{1-q^k}$.
\end{remark}

\begin{remark}[Verma module style]
\label{vertex-algebras-remark-Verma-module-style}
$H$ can be presented alternatively as
$$
H=U(\hat{\mathfrak{a}})
\otimes_{U(\hat{\mathfrak{a}}_+)}\mathbb{C}1
$$
where $\mathbb{C}1$ is a 1-dimensional
representation of
 $$
\hat{\mathfrak{a}}_+=\bigoplus_{n \geq
0}\mathbb{C}h_n \oplus\mathbb{C}K.
$$
$$
h_n \cdot 1=0\qquad \forall n \geq 0,\qquad
K\cdot 1=1
$$
\end{remark}

\begin{exercise}
\label{vertex-algebras-exercise-H-is-irreducible}
$H$ is irreducible.
\end{exercise}

\begin{proof}
We need to prove that there is no
vector subspace $V \subset
\mathbb{C}[x_1,x_2,\ldots]=H$
such that $xV \subset V$ for all $x \in
\hat{\mathfrak{a}}$.
And then the proof is basically noticing
that the orbit of the action of
$\hat{\mathfrak{a}}$ on $H$
is all of $H$. That is, if we had proper
subspace
of $H$, we can always find an element of
$\hat{\mathfrak{a}}$
that takes some element in $V$ to the one of
the
elements that is not in $V$. Indeed, by
applying
$h_n$ for  different values of positive  $n$
we can take
any element of $V$ to a constant. Then we
apply
$h_n$ for different values of negative $n$ to
obtain
any monomial. Then we add these monomials and
obtain any
polynomial in $H$.
\end{proof}

\noindent
{\bf Notation.} Instead of $1$ let's write
$|0\rangle$.

\begin{remark}
\label{vertex-algebras-remark-generalise-H}
We can easily generalise $H$ to
$$
H^\mu=U(\hat{\mathfrak{a}})
\otimes_{U(\hat{\mathfrak{a}}_+)} \mathbb{C}
|\mu\rangle,
$$
($\mu \in \mathbb{C}$), where
$\hat{\mathfrak{a}}_+$ is the same, but now
\begin{align*}
h_n |\mu\rangle&=\begin{cases}
0\qquad &\text{if }n>0 \\
\mu \cdot|\mu\rangle\qquad& \text{if }n=0
\end{cases}\qquad \qquad
K|\mu\rangle=|\mu\rangle.
\end{align*}
$H^\mu$ is again a (irreducible)
$\hat{\mathfrak{a}}$-module.
\end{remark}

\medskip\noindent
Generalise even more: Let $\mathfrak{h}$ be a
finite-dimensional
vector space, and
$(\cdot,\cdot): \mathfrak{h} \times
\mathfrak{h} \to \mathbb{C}$
a symmetric bilinear form.
$\hat{\mathfrak{h}}=\mathfrak{h}[t,t^{-1}]
\oplus
\mathbb{C}K$,
$$
[at^m,bt^n]=m(a,b)\delta_{m,-n}K,\qquad
[K,\hat{\mathfrak{h}}]=0.
$$
Let $\mu\in \mathfrak{h}$. Define
$H^\mu=U(\hat{\mathfrak{h}})
\otimes_{U(\hat{\mathfrak{h}})}\mathbb{C}
|\mu\rangle$
,
$\hat{\mathfrak{h}}_+=\mathfrak{h}[t]
\oplus\mathbb{C}K$,
and
$$
at^m\cdot |\mu\rangle=\begin{cases}
(\mu,a)|\mu\rangle\qquad &\text{if } m=0\\
0\qquad &\text{if }m>0.
\end{cases}
$$

\medskip\noindent
Returning to the action of
$\hat{\mathfrak{a}}$ on $H=H^0$. Introduce
$$
L_m=\frac{1}{2}\sum_{k \in
\mathbb{Z}}h_{m-k}h_k.
$$
If $m \neq 0$, this sum is well-defined in
$\text{End}(H)$.

Indeed, $h_{m-k}$ and $h_k$ commute ($m \neq
0$), and
$\forall  p(\underline{x})\in H$, there
exists $N$
such that
$p(\underline{x})=p(x_1,\ldots,x_{N-1})$,
then $h_k\cdot p(x)=0$ $\forall  k \geq N$,
and $h_{m-k}\cdot p(x)=0$ $\forall m-k \geq
N$.

So $(\sum h_{m-k}h_k)p(\underline{x})$ is a
\underline{finite} sum.
Of course, the number of terms in the sum
depends on $p(x)$ and can be
arbitrarily large.

\begin{exercise}
\label{vertex-algebras-exercise-hard}
If  $m,n$ and $m+n$ are not zero, then
$$
[L_m,L_n]=(m-n)L_{m+n}
$$
holds in $\text{End}(H)$.

(Trying to be a representation of $W$
 $$
D_n \mapsto L_n \in \text{End}(H).)
$$
\end{exercise}

But $\sum_{k \in \mathbb{Z}}h_{-k}h_k$ is not
well-defined.
Indeed,
$$
\begin{aligned}
\left(\sum_{k \in
\mathbb{Z}}h_{-k}h_k\right)1&=
\sum_{k \geq 1}k\frac{\partial }{\partial
x_k}(x_k 1)\\
&=(1+2+3+4+\ldots)1
\end{aligned}
$$
!

\medskip\noindent
{\bf Normal ordering idea (``that physicist
do'').} Let's cheat
and redefine the product as
$$
\,:\!h_{-k}h_k\!:\,
=\begin{cases}
h_{-k}h_k\qquad & \text{if }k\geq 0\\
h_k h_{-k}\qquad &\text{if }k<0.
\end{cases}
$$
Now $\sum_{k \in \mathbb{Z}}:\!h_{-k}h_k\!:$
is well-defined!

\medskip\noindent
{\bf Notation.} Let us consider series of the
form
$$
a(z)=\sum_{n \in
\mathbb{Z}}a_{(n)}z^{-n-1},\qquad a_{(n)}\in
\text{End}(V),
$$
where $V$ is a vector space.

\begin{definition}[Most important in the
course]
\label{vertex-algebras-definition-quantum-field}
We say $a(z)$ is a {\it quantum field} on $V$
if,
for every $v \in V$ there exists $N \in
\mathbb{Z}$
such that $a_{(n)}v=0$ $\forall  n\geq N$.
\end{definition}

\begin{example}
\label{vertex-algebras-example-quantum-field}
$V=H=\mathbb{C}[x_1,x_2,\ldots]$,
$$
a(z)=h(z)=\sum_{n \in \mathbb{Z}}h_n
z^{-n-1}.
$$
$$
h(z)=\ldots+3z^{-4}\frac{\partial }{\partial
x_3}
+2z^{-3}\frac{\partial }{\partial x_1}+x_1
+x_2z+x_3z^2+\ldots
$$
Any fixed $v \in H$ is
$v=p(x_1,\ldots,x_{N-1})$ for some $N$.
Then $h_Nv=N \frac{\partial }{\partial
x_N}v=0$.

So $h(z)$ is a quantum field on $H$.
\end{example}

\begin{definition}
\label{vertex-algebras-definition-creat-annih-opera}
For a quantum field
(or any series in fact) $a(z)=\sum_{ n \in
\mathbb{Z}}a_{(n)}z^{-n-1}$,
define the {\it creation} and {\it
annihilating operators}
$$
a(z)_+=\sum_{n \leq -1}a_{(n)}z^{-n-1},\qquad
a(z)_-=\sum_{n \geq 0}a_{(n)}z^{-n-1}.
$$
\end{definition}

The quantum field condition solves the
problem of the infinite series…

The normally order product of quantum fields
$a(z)=\sum_{n \in
\mathbb{Z}}a_{(n)}z^{-n-1}$, and
$b(z)=\sum_{n \in \mathbb{Z}}b_{(n)}z^{-n-1}$
is.
$$
:\!a(z)b(z)\!:\,=a(z)_+b(z)+b(z)a(z)_-.
$$

\begin{exercise}[Most important of the
course]
\label{vertex-algebras-exercise-most-important}
Show that if $a(z)$ and $b(z)$ are quantum
fields,
then $:\!a(z)b(z)\!:\,$ is well-defined, and
is a quantum field.
\end{exercise}

\medskip\noindent
Next idea: in our example of $h(z)$ acting on
$H$,
the coefficients $h_n$ were coming from a Lie
algebra,
so we had relations
 $[h_m,h_n]=(\ldots)$.
(Indeed, $[h_m,h_n]=m
\delta_{m,-n}\text{Id}_H$ in this case.)
We would like to interpret such relations at
the level of $h(z)$.

$$
\begin{aligned}
[h(z),h(z)]&=h(z)h(z)-h(z)h(z)\\
&=\sum_{p}\sum_{m+n=p}h_mh_nz^{-(m+n)-2}
-\sum_{m+n=p}h_nh_mz^{-(m+n)-2}\\
&=\sum_p
z^{-p-2}\underbrace{\left(\sum_{m+n=p}[h_m,
h_n]\right)}
_{\text{infinite?}}
\end{aligned}
$$
which is a bad idea, since that sum can be
infinite.
Better idea: change a variable:
$$
[h(z),h(w)]=\sum_{m,n \in
\mathbb{Z}}[h_m,h_n]z^{-m-1}w^{-n-1}.
$$
In this example,
$$
\begin{aligned}
[h(z),h(w)]&=\sum_{m,n \in \mathbb{Z}}m
\delta_{m,-n}z^{-m-1}
w^{-n-1}I_H\\
&=\sum_{m \in \mathbb{Z}}z^{-m-1}w^{m-1}I_H.
\end{aligned}
$$

Observe (geometric series):
$$
\begin{aligned}
\frac{1}{z-w}&=\sum_{k \geq
0}z^{-k-1}w^k\qquad
\text{(convergent for $|z|>|w|$)}\\
\frac{1}{z-w}&=-\sum_{k \geq 0}w^{-k-1}z^k\\
&=-\sum_{k<0}z^{-k-1}w^k\qquad
\text{(convergent for $|z|<|w|$)}
\end{aligned}
$$
So, in a sense
$$
\sum_{k \in
\mathbb{Z}}z^{-k-1}w^k\text{``$=$''}\frac{1}
{z-w}
-\frac{1}{z-w}.
$$
This motivates us to introduce the expression
\begin{definition}
\label{vertex-algebras-definition-delta-function}
The {\it formal delta function} is
$$
\delta(z,w)=\sum_{k \in
\mathbb{Z}}z^{-k-1}w^k \in
\mathbb{C}[\![z,z^{-1},w,w^{-1}]\!].
$$
\end{definition}
Why delta function? It behaves like Dirac
delta.
Recall that $\delta(x)$ is the
{\it distribution} on $\mathbb{R}$ defined by
$\delta[f(x)]=f(0)$ for all test functions
$f(x)$.
Every function $k(x)$ gives a distribution
$D_k$.
$$
D_k[f]=\int_{-\infty}^\infty k(x)f(x)dx,
$$
$f \in C_c^\infty(\mathbb{R})$.

If $\delta$ were of the form $D_k$ (it's not)
it would have to look like
[Picture, positive part of $y$ axis, all $x$
axis.

One can show that (it's a theorem by Plamelj)
$$
\delta(x)=\frac{1}{2\pi i}\lim_{\varepsilon
\to 0_+}
\left(\frac{1}{x-i \varepsilon}-\frac{1}{x-i
\varepsilon}\right)
$$
as distributions (i.e. limit taken in
``distributional sense''.)

(So that's why we called the delta function
that way.)

Denote by $i_{z,w}$ the ``expansion in
positive powers of $w$''.
E.g.
$$
i_{z,w}\frac{1}{z-w}=\sum_{k\geq
0}z^{-k-1}w^k,
$$
similarly,
$$
i_{z,w}\frac{1}{(z-w)^2}=\sum_{k \geq
0}kz^{-k-1}w^{k-1},
$$
$$
i_{z,w}(w^{-2})=w^{-2}.
$$
\begin{exercise}
\label{vertex-algebras-exercise-derivative-wrt-w}
So, in fact,
$$
\begin{aligned}
[h(z),h(w)]&=i_{z,w}\frac{1}{(z-w)^2}-i_{w,z}
\frac{1}{(z-w)^2}
=\partial_w \delta(z,w).
\end{aligned}
$$
\end{exercise}

What exactly is $i_{z,w}$?

\medskip\noindent
{\bf Notation.}
$$
\begin{aligned}
\mathbb{C}[z]&\qquad \text{ring of
polynomials}\\
\mathbb{C}[z,z^{-1}]&\qquad \text{ring of
Laurent polynomials}\\
\mathbb{C}[\![z ]\!]&\qquad \text{ring of
power series}\\
\mathbb{C}[\![z,z^{-1}]\!]&  \qquad
\text{vector space (not ring!) of formal
distributions}\\
\mathbb{C}(\!(z)\!)&\qquad \text{field of
Laurent series}
\end{aligned}
$$
where the last is $\sum_{n \geq N}fa_n z^n$
for some $N$.

Since $\mathbb{C}(\!(z)\!)$ is a field,
$\mathbb{C}(\!(z)\!)(\!(w)\!)$ is
also a field.

There are natural inclusions
$$
\mathbb{C}[z,w] \to
\mathbb{C}(\!(z)\!)(\!(w)\!) \to
\mathbb{C}[\![z^{\pm1},w^{\pm 1}]\!].
$$
Let $\mathbb{C}(z,w)$ denote the fraction
field
of the domain $\mathbb{C}[z,w]$.
By the property of $\mathbb{C}(z,w)$, there
exists an embedding
$$
\xymatrix{
\mathbb{C}[z,w]\ar[rr]\ar[dr]&  &
\mathbb{C}(\!(z)\!)(\!(w)\!)\\
&\mathbb{C}(z,w)\ar[ur]^{i_{z,w}}
}
$$
The following diagram does not commute:
[diagram]

We computed
$$
[h(z),h(w)]=\partial_w \delta(z,w)
$$
We can similarly compute
$$
\begin{aligned}
h(z)h(w)&=\,:\!h(z)h(w)\!:\,+i_{z,w}\frac{1}
{(z-w)^2}\\
h(w)h(z)&=\,:\!h(z)h(w)\!:\,+i_{w,z}\frac{1}
{(z-w)^2}.
\end{aligned}
$$

\medskip\noindent
{\bf Notation.} We write
$\partial_w=\frac{\partial }{\partial w}$ and
$\partial^{(j)}_w=\frac{1}{j!}\partial^j_w$.

\begin{lemma}
\label{vertex-algebras-lemma-elementary}
\begin{enumerate}
\item If we multiply the delta function with
$(z-w)$ we get zero, that is,
$$
(z-w)\delta(z,w)=0.
$$
\item
$$
i_{z,w}\frac{1}{(z-w)^{j+1}}-i_{w,z}\frac{1}
{(z-w)^{j+1}}
=\partial_w^{(j)}\delta(z,w).
$$
\item $\forall j\geq 0$,
$$
(z-w)^{j+1}\partial_w^{(j)}\delta(z,w)=0.
$$
\item Whenever $m \leq j$,
$$
(z-w)^m\partial^{(j)}(z,w)=\partial_w^{(j-m)}
\delta(z,w).
$$
\end{enumerate}
\end{lemma}

\begin{proof}

\end{proof}

\begin{remark}
\label{vertex-algebras-remark-like-usual-Dirac-delta}
All this is completely parallel
to the Dirac $\delta$-distribution
$\delta(x)$.
$$
x\delta(x)=0,\qquad
x\delta'(x)=\delta(x),\text{ etc.},
x^2\delta'(x)=0.
$$
\end{remark}

\section{The residue pairing}
\label{vertex-algebras-section-residue-pairing}

Let $f(z) \in \mathbb{C}[z,z^{-1}]$, which is
a vector space
with basis $\{z^n:n\in\mathbb{Z}\}$.

An element of the dual vector space is a
formal linear combination
$$
\sum_{n \in \mathbb{Z}}c_n\varphi_n,\qquad
\text{where }
\varphi_n(z^k)=\begin{cases}
0\qquad &\text{if }k \neq n \\
1\qquad &\text{if }k=n
\end{cases}
$$
(no restriction on the $c_n \in \mathbb{C}$).

Identify $\sum c_n \varphi_n$ with
$$
c(z)=\sum c_nz^{-n-1}\in \mathbb{C}[z^{\pm 1}
$$
so that $\varphi$ acts on $f(z)=\sum f_nz^n
\in \mathbb{C}[z^{\pm 1}]$,
as
$$
\varphi(f)=\sum_{n}c_nf_n=\text{Res}_zc(z)
f(z)dz,
$$
where
\begin{definition}
\label{vertex-algebras-definition-residue}
If  $U$ is a vector space and
$a(z)=\sum_{n}a_nz^n \in U[\![z^{\pm 1}]\!]$,
$$
\text{Res}_za(z)dz=a_{-1}
$$
\end{definition}

Let's record a few properties of
$\text{Res}_z(\cdot)dz:=\text{Res}_z(\cdot)$.

\begin{lemma}
\label{vertex-algebras-lemma-properties-of-residues}
\begin{enumerate}
\item $\text{Res}_zf(z)\delta(z,w)=f(w)$.
\item
$\text{Res}_zf(z)(\partial_zg(z))=-\text{Res}
_(\partial_zf(z))fg(z)$.
\end{enumerate}
\end{lemma}

\begin{proof}
\begin{enumerate}
\item Straightforward computation.
\item Product rule, and uses that derivatives
have no residues.
\end{enumerate}
\end{proof}

So, in our example, $H,h(z)$ we see that
$$
(z-w)^2[h(z),h(w)]=0,
$$
which is a nontrivial example: we really need
to
take the square, and also the bracket is not
zero.

\medskip\noindent
\begin{definition}
\label{vertex-algebras-definition-mutua-local-quant-field}
Two quantum fields $a(z)$ and $b(z)$ on a
vector space
$V$ are {\it mutually local} if there exists
$N$
such that
$$
(z-w)^N[a(z),b(w)]=0.
$$
\end{definition}

\begin{remark}
\label{vertex-algebras-remark-motivation}
See  \cite[Chapter 1]{KVA} for physics
motivation.
More generally, if $D=\sum_{m,n}D_{m,n}z^mw^n
\in U[\![z^{\pm 1}, w^{\pm
1}]\!]$,
we call $D$ {\it local} if $(z-w)^ND=0$ for
some $N$.
\end{remark}

\begin{proposition}
\label{vertex-algebras-proposition}
Let $a(z),b(z)$ be a local pair of quantum
fields on $V$.
Then
\begin{equation}
\label{vertex-algebras-equation-1}
a(z)b(w)=
:\!a(z)b(z)\!:+\sum_{j=0}^{N-1}c_j(w)
i_{z,w}\frac{1}{(z-w)^{j+1}}
\end{equation}
and
\begin{equation}
\label{vertex-algebras-equation-2}
b(w)a(z)=:\!a(z)b(z)\!:+\sum_{j=0}^{N-1}
c_j(w)
i_{w,z}\frac{1}{(z-w)^{j+1}}
\end{equation}
In particular,
\begin{equation}
\label{vertex-algebras-equation-3}
[a(z),b(w)]=\sum_{j=0}^{N-1}c_j(w)
\partial_w^{(j)}\delta(z,w).
\end{equation}
Furthermore
\begin{equation}
\label{vertex-algebras-equation-4}
c_j(w)=\text{Res}_z(z-w)^j[a(z),b(w)].
\end{equation}
\end{proposition}

\begin{proof}
First we prove \ref{vertex-algebras-equation-3} using
\ref{vertex-algebras-equation-4}.
\end{proof}

\begin{exercise}
\label{vertex-algebras-exercise-deduce-1-and-2-from-3}
Deduce \ref{vertex-algebras-equation-1} and  \ref{vertex-algebras-equation-2}
from \ref{vertex-algebras-equation-3}.
\end{exercise}

\medskip\noindent
{\bf Notation.} If $a(w)$ and $b(w)$ are
quantum fields on $V$,
we denote
$$
a(w)_{(j)}b(w)=\text{Res}_z(z-w)^j[a(z),b(w)]
$$
for $j \in \mathbb{Z}_{\geq 0}$.

\begin{example}
\label{vertex-algebras-example-duplicate-2}
On $H$,
$[h(z),h(w)]=\partial_w\delta(z,w)I_H$.
So $\text{Res}_z\partial_w\delta(z,w)I_H=0$
and
$\text{Res}_z(z-w)\partial_w\delta(z,w)I_H=
\text{Res}_z\delta(z,w)I_H=I_H$.
\end{example}

Note that $I_H$ is a (very simple) quantum
field.

The following definition generalizes the
$j$-product to any
integer, possible negative.

\begin{definition}
\label{vertex-algebras-definition-nth-product}
For $n \in \mathbb{Z}$, and $a(w)$, $b(w)$
quantum fields,
we define the {\it $\text{n}^{\text{th}}$
product} $a(w)_{(n)}b(w)$ as
$$
a(w)_{(n)}b(w)=\text{Res}_z
[i_{z,w}(z-w)^na(z)b(w)-i_{w,z}(z-w)^nb(w)
a(z)].
$$
\end{definition}

\begin{remark}
\label{vertex-algebras-remark-nth-product}
\begin{enumerate}
\item We recover the prior definition: if $n
\geq 0$,
$$
i_{z,w}(z-w)^n=i_{w,z}(z-w)^n=\sum_{r=0}
^n\binom{n}{r}z^{n-r}(-w)^r,
$$
a finite sum.
\item (Exercise.) If $n=1$, then
$$
\begin{aligned}
a(w)_{(-1)}b(w)&=a(w)_+b(w)+b(w)a(w)_-\\
&=:\! a(w)b(w)\!:
\end{aligned}
$$

\item (Exercise.)
The more negative products are not something
new:
for $k \geq 0$,
$$
a(w)_{(-k-1)}b(w)=:\!(\partial_w^{(k)}a(w))
b(w)\!:
$$
\end{enumerate}
\end{remark}

We have actually already used the following
proposition:

\begin{proposition}
\label{vertex-algebras-proposition-nth-produ-quant-field}
If $a(w),b(w)$ are quantum fields, then
$a(w)_{(n)}b(w)$
is also a quantum field for all $n \in
\mathbb{Z}$.
\end{proposition}

\begin{definition}
\label{vertex-algebras-definition-vertex-algebra}
A {\it vertex algebra} consists of a vector
space $V$,
a set $\mathcal{F}$ of quantum fields on $V$,
a nonzero vector $|0\rangle\in V$,
and a linear map $T:V \to V$,
such that
\begin{enumerate}
\item $T|0\rangle=0$, and
$[T,a(z)]=\partial_za(z)$
$\forall a(z) \in \mathcal{F}$.
\item $V$ is spanned by
$a_{(n_1)}^{i_1},\ldots,a_{(n_s)}^{i_s}
|0\rangle$,
where $a^{i_j}(z) \in \mathcal{F}$ (and
$s(z)=\sum a_{(n)}z^{-n-1}$ always).
\item All pairs $a(z),b(z) \in \mathcal{F}$
are mutually local.
(We saw (ref?) that this is equivalent to
$[a(z),b(w)]=\sum_{j=0}^{N-1}c_j(w)
\partial_w^{(j)}\delta(z,w)$
and
described the coefficients $c_j$, …)
\end{enumerate}
\end{definition}

We call $|0\rangle$ the {\it vacuum vector}
and $T$ the {\it translation
operator}.

In fact, our Heisenberg example that we've
been discussing
so far is an example:

\begin{example}
\label{vertex-algebras-example-Heisenberg-is-vertex-algebra}
$V=H=\mathbb{C}[x_1,x_2,\ldots]$,
$\mathcal{F}=\{h(z)\}$, $|0\rangle=1$, $T=?$,
is a vertex algebra.

\noindent
{\bf Answer 1.} Recall
$L(z)=\frac{1}{2}:\!h(z)h(z)\!:=
\frac{1}{2}h(z)_{(-1)}h(z)$, $L(z)=\sum_{m
\in \mathbb{Z}}L_mz^{-m-z}$.
In particular, $L_{-1}=\frac{1}{2}\sum_{k \in
Z}h_kh_{-1-k}$.
Well, $T=L_{-1}$.

\noindent
{\bf Answer 2.} Given that $T$ must satisfy
$T1=0$, and $[T,h(z)]=\partial_zh(z)$, i.e.
$$
[T,h_n]=-nh_{n-1},
$$
and $H$ is generated from $1$ by $\{h_n|n\leq
-1\}$,
the action of $T$ on $H$ (if well-defined) is
completely determined.
\end{example}

\begin{exercise}
\label{vertex-algebras-exercise-write-a-formula}
Write a formula for
$T(x_1^{m_1}x_2^{m_2}\ldots x_s^{m_2})$.
\end{exercise}

\section{A second definition of vertex
algebra}
\label{vertex-algebras-section-secon-defin-verte-algeb}

\noindent
Now we aim for a second definition of vertex
algebras.

Let $(V,|0\rangle,T,\mathcal{F})$ be a vertex
algebra.
Define
$$
\begin{aligned}
\tilde{\mathcal{F}}&=\{\text{quantum fields
$a(z)$ on $V$}
|[T,a(Z)]=\partial_za(z)\}\\
\overline{\mathcal{F}}&=\bigcup_{k\geq
0}\mathcal{F}_k,\\
\mathcal{F}_0&=\{I_V\},\\
\mathcal{F}_k&=\{a(z)_{(n)}b(z)|a(z)\in
\mathcal{F},b(z),\mathcal{F}_{k-1}\}\\
\mathcal{F}'&=\{b(z) \in
\tilde{\mathcal{F}}|a(z),b(z)\text{ is
local},
\forall  a(z) \in \mathcal{F}\}
\end{aligned}
$$
Where $I_V$ is the identity of $V$ considered
as a quantum field.
The idea is that this is like taking a
subalgebra and taking
the commutator over and over again.

\begin{lemma}
\label{vertex-algebras-lemma-simple-lemma1}
$a(z)_{(-1)}I_V=a(z)$.
\end{lemma}

\begin{proof}
Direct calculation.
\end{proof}

\begin{lemma}
\label{vertex-algebras-lemma-simple-lemma2}
$a(z)_{(-n-1)}I_V=\partial_z^{(n)}a(z)$.
\end{lemma}

\begin{proof}
Similar.
\end{proof}

\begin{lemma}[Dong]
\label{vertex-algebras-lemma-dong}
Suppose $a(z),b(z)$ and $c(z)$ are mutually
local in pairs.
Let $n \in \mathbb{Z}$. Then $a(z)_{(n)}b(z)$
and $c(z)$
is a local pair.
\end{lemma}

\begin{proof}
Done in lecture.
\end{proof}

Therefore
$$
\mathcal{F} \subset
\overline{\mathcal{F}}\overset{\text{Dong}}
{\subset}
\mathcal{F}'\subset \tilde{\mathcal{F}}
$$
\begin{remark}
\label{vertex-algebras-remark-Fbar-subset-Ftilde}
To be sure, we should check
$\overline{\mathcal{F}}\subset
\tilde{\mathcal{F}}$.
\end{remark}

F tilde is the translation invariant
\begin{exercise}
\label{vertex-algebras-exercise-Fbar-subset-Ftilde}
If $[T,a(z)]=\partial_za(z)$,
$[T,b(z)]=\partial_zb(z)$,
then
$$
[T,a(z)_{(n)}b(z)]=\partial_z(a(z)_{(n)}b(z))
.
$$
\end{exercise}

\begin{theorem}
\label{vertex-algebras-theorem}
\begin{enumerate}
\item For any $a(z) \in \tilde{\mathcal{F}}$,
\label{vertex-algebras-item-verte-algeb-secon-defin-1}
$$
\begin{aligned}
a(z)|0\rangle&=v+z\cdot
Tv+\frac{z^2}{2}T^2v+\ldots\\
&=e^{2T}\qquad (v\in V)
\end{aligned}
$$
Thus we define
$$
\begin{aligned}
s: \tilde{\mathcal{F}} &\longrightarrow V \\
a(z) &\longmapsto a(z)|0\rangle|_{z=0}
\end{aligned}
$$

\item
\label{vertex-algebras-item-verte-algeb-secon-defin-2}
$$
s(z(z)_{(n)}b(z)=a_{(n)}s(b(z)).
$$
\item Let $\mathcal{G} \subset
\tilde{\mathcal{F}}$ be
\label{vertex-algebras-item-verte-algeb-secon-defin-3}
such that $s(\mathcal{G})=V$, and suppose
$a(z)$ is local with all $b(z) \in
\mathcal{G}$.
``If you commute with a large enough bunch of
guys,
you are close enough to being zero.''
If $s(a(z))=0$, then $a(z)=0$.
\end{enumerate}
\end{theorem}

\begin{proof}
\begin{enumerate}
\item We need to prove $a(z)|0\rangle$ does
not have negative
powers of $z$. So suppose it does, i.e. there
is $n\geq 0$ such that
$a_{(n)}|0\rangle z^{-n-1}\neq 0$.

Idea: use translation invariance to pair $n$
and $-n$, and
then the fact that $a(z)$ is a quantum field
to get some vanishing.

So first we translate and find that:
$$
Tv(z)=Ta((z)|0\rangle=(Ta(z)-a(z)T)|0\rangle=
\partial_z
a(z)|0\rangle
$$
so that
\begin{equation}
\label{vertex-algebras-equation-star}
v_n=\frac{1}{n}Tv_{n-1},\qquad (n \neq 0).
\end{equation}
So, $v_n \neq 0$ for some $n<0 \implies
v_{n-1}$ also not zero.
But we also know by $a(z)$ being a quantum
field
that there exists $N$ such that
$a_{(n)}|0\rangle=0$ for all $n \geq N$.
So, we'd have a contradiction.
So $v(z)=\sum_{n \geq 0}v_nz^n$.

Now Eq. \ref{vertex-algebras-equation-star} also implies
$$
v_1=Tv_0,v_0=\frac{1}{2}Tv_1,\ldots,
v_n=\frac{1}{n!}T^nv_0.
$$
So $a(z)|0\rangle=e^{zT}v_0$ $(v_0 \in V)$.

\item Consider
$$
a(w)_{(n)}b(w)|0\rangle=
\text{Res}_z(a(z)b(w)i_{z,w}(z-w)^n-
\underbrace{b(w)a(z)i_{w,z}(z-w)^n}_{=0}
$$
where that vanishing is because the $i_{w,z}$
expands the argument
in positive powers, and $b(w)a(z)$ also has
only positive powers
since $a(z)|0\rangle$ also has only positive
powers
(and the residue picks the coefficient of the
power  $-1$).
So we obtain
$$
=\text{Res}_za(z)b(w)i_{z,w}(z-w)^n.
$$
But
$$
\begin{aligned}
s(a(w)_{(n)}b(w))&=\left(
\text{Res}_zs(z)b(w)i_{z,w}(z-w)
^n|0\rangle\right)|_{w=0}\\
&=\text{Res}_za(z)b(w)z^n|0\rangle|_{z=0}\\
&=\text{Res}_z z^n a(z)(s(b(w)))\\
&=\text{Res}_z z^n
\sum_{k}a_{(k)}z^{-k-1}s(b)\\
&=a_{(n)}s(b).
\end{aligned}
$$

\item By locality, there exists $N$ such that
$$
\begin{aligned}
(z-w)^Na(z)b(w)
|0\rangle&=(z-w)^Nb(w)a(z)|0\rangle\\
&=(z-w)^Nb(w)e^{2T}s(a)=0
\end{aligned}
$$
$$
\begin{aligned}
(z-w)^Na(z)b(w)|0\rangle |_{w=0}&=0\\
\implies z^Na(z)s(b)&=0\\
\implies a(z)v&=0\\
\implies a(z)&=0
\end{aligned}
$$
\end{enumerate}
\end{proof}

\medskip\noindent
{\bf Corollary.}
\begin{enumerate}
\item $S|_{\mathcal{F}'}$ is injective.
\item $S|_{\overline{\mathcal{F}}}$ is
surjective.
\item $S|_{\mathcal{F}}$ is isomorphism
$\overline{\mathcal{F}}\to V$.

So we define the inverse (linear) map $Y:V
\to \overline{\mathcal{F}}$,
i.e. for all $a \in V$, $Y(a,z)$ denotes
$s^{-1}(a) \in \overline{\mathcal{F}}$.

\item
$(V,|0\rangle,T,\overline{\mathcal{F}})$
is a vertex algebra,
and
$$
Y(a_{(n)}b,z)=Y(a,z)_{(n)}Y(b,z)\qquad
\forall a,b \in V.
$$
\end{enumerate}

So we have
$$
\xymatrix{
\mathcal{F}\ar[r]&\overline{\mathcal{F}}
\ar@{->>}[d]\ar[r]&
\mathcal{F}'\ar@{^{(}->}[d]\ar[r]&
\tilde{\mathcal{F}}\ar[d]^s\\
&  V\ar[r]^{\text{id}}& V\ar[r]^{\text{id}}
&V
}
$$
\begin{proof}
\begin{enumerate}
\item Now we apply part
\ref{vertex-algebras-item-verte-algeb-secon-defin-3}
 in the theorem to $a(z) \in \mathcal{F}'$,
$\mathcal{G}=\overline{\mathcal{F}}$.
If $s(a(z))=0$, then  $a(z)=0$. So
$s:\mathcal{F}' \to V$ is injective
(as in the diagram).

\item By iterating the property in the
theorem that
$a_{(n)}s(b(z))=s(a(z)_{(n)}b(z))$ so we can
take out ``one by one''
in the following product to obtain
$$
a^1_{(n_1)}\ldots a_{(n_s)}^s |0\rangle
=s(a^1(z)_{(n)}\ldots a^s(z)_{(n_s)}I_V
$$
This proves the double headed arrow in the
diagram.

\item Now define $Y(-,z):V \to
\overline{\mathcal{F}}$.
$$
s(Y(a,z)_{(n)}Y(b,z))=a_{(n)}s(Y(b,z))=a_{(n)
}b.
$$
\item Then all the axioms in our definition
of vertex algebra are
satisfied!
\end{enumerate}
\end{proof}

\begin{remark}
\label{vertex-algebras-remark-a-relation}
Let $(V,|0\rangle,T,\mathcal{F})$ be a vertex
algebra.
$\overline{\mathcal{F}}$ and $Y$ as above.
$\forall a,b \in V$,
\begin{equation}
\label{vertex-algebras-equation-star2}
[Y(a,z),Y(b,w)]=\sum_{j=0}^{N-1}c_j(w)
\partial_w^{(j)}\delta(z,w).
\end{equation}
$c_j(w)=Y(a,w)_{(j)}Y(b,w)=Y(a_{(j)}b,w)$.
\end{remark}

Which is like something we had written
before.
But now we know more. Let us extract
the $z^{-m-1},w^{-n-1}$ coefficient of Eq.
\ref{vertex-algebras-equation-star2}.
We obtain
$$
LHS=\sum_{m,n}[a_{(m)}z^{-m-1},b_{(n)}w^{-n-
1}]=[a_{(m)},b_{(n)}]
$$
$$
RHS=\sum_{j}\left(\sum_{k}(a_{(j)}b)_{(k)}
w^{-k-1}\right)
\left(\sum_{s}\binom{s}{j}z^{-s-1}w^{s-j}
\right)
$$
The coefficient of $z^{-m-1}w^{-n-1}$ of RHS
is, $s=m$,
$$
\sum_{j}\binom{m}{j}w^{m-j}(a_{(j)}b)_{(k)}
w^{-k-1}
=\sum_{j}\binom{m}{j}(a_{(j)}b)_{(m+n-j)}
$$
So we put this as a proposition:
\begin{proposition}[Commutator formula]
\label{vertex-algebras-proposition-commutator-formula}
In a vertex algebra $(V,|0\rangle,T,Y)$, we
have
$$
[a_{(m)},b_{(n)}]=\sum_{j \geq
0}\binom{m}{j}(a_{(j)}b)_{(m+n-j)}
\qquad \forall a,b \in V,\quad m,n\in
\mathbb{Z}
$$
\end{proposition}

\begin{exercise}
\label{vertex-algebras-exercise-its-what-it-is}
Prove that in a vertex algebra
\begin{equation}
\label{vertex-algebras-equation-star3}
[Y(a,z)Y(b,w)i_{z,w}-Y(b,w)Y(a,z)i_{w,z}](z-
w)^n
=\sum_{j \geq
0}Y(a_{(n+j)}b,w)\partial_w^{(j)}\delta(z,w).
\end{equation}
Hint. Prove that he left hand side is local.
\end{exercise}

By extracting coefficients of Eq.
\ref{vertex-algebras-equation-star3}, we obtain
for all $a,b,c \in V$ and $m,n,k \in
\mathbb{Z}$:
\begin{equation}
\label{vertex-algebras-equation-Borcherds-identity}
\begin{aligned}
&\sum_{j \geq
0}\binom{m}{j}(a_{(n+j)}b)_{(m+k-j)}c\\
&=\sum_{j \geq 0}(-1)^j \binom{n}{j}
\left(a_{(m+n-j)}b_{(k+j)}c-(-1)^nb_{(n+k-j)}
a_{(m+j)}c\right)
\end{aligned}
\end{equation}

This is called {\it Borcherd's identity}, and
is the key
ingredient in our second definition of vertex
algebra
(see \cite[Proposition 4.8(b)]{VAFB}):

\begin{definition}
\label{vertex-algebras-definition-vertex-algebra2}
A {\it vertex algebra} is a vector space $V$,
a nonzero vector $|0\rangle\in V$,
and a set of bilinear products
$V \times V \to V$, $a,b \mapsto a_{(n)}b$,
$n \in \mathbb{Z}$
such that
\begin{enumerate}
\item $\forall a,b \in V$ $\exists N$ such
that $a_{(n)}b=0$
$\forall n \geq N$.
\item $\forall  a \in V$,
$|0\rangle_{(n)}a=\delta_{n,-1}a$.
\item $\forall a \in V,$ $a_{(-1)}|0\rangle
=a$ and $a_{(\geq 0)}|0\rangle=0$.
\item Equation
\ref{vertex-algebras-equation-Borcherds-identity} holds
$\forall a,b,c \in V$, $m,n,k \in
\mathbb{Z}$.
\end{enumerate}
\end{definition}

\section{Calculating vertex algebras}
\label{vertex-algebras-section-calculating-vertex-algebras}


We begin with some notation.

\begin{definition}
\label{vertex-algebras-definition-lambda-bracket}
Let $V$ be a vertex algebra and
$a,b \in V$.
We package the elements
$a_{(0)}b,a_{(1)}b,\ldots$
into a series
$$
[a_{\lambda}b]=\sum_{j \geq 0}a_{(j)}b
\frac{\lambda^j}{j!}
$$
called the {\it $\lambda$-bracket} or {\it
operator product expansion (OPE)}
of $a$ and $b$.
\end{definition}

Recall that elements in $V$ correspond to
fields $Y(a,z)$ and $Y(b,z)$.
They are local, and the coefficients of their
bracket are
$c_j(w)=Y(a_{(j)}b,w)$. So we put that
information in the $\lambda$-bracket.

We also denote by $\,:\!ab\!:\,$ the vector
correponding to the normally ordered
product $\,:\!Y(a,z)Y(b,z)\!:\,$. Again, by
last time's theorem,
$\,:\!ab\!:\,=a_{(-1)}b$.

\begin{example}
\label{vertex-algebras-example-heise-verte-algeb-defin}
$V=H=\mathbb{C}[x_1,x_2,\ldots]$,
$\mathcal{F}=\{h(z)\}$ where
$h(z)=\sum_{n \in \mathbb{Z}}z^{-n-1}$ (so in
this example we are denoting
$h_n \equiv h_{(n)}$), where, as before,
$$
h_n=\begin{cases}
n\frac{\partial }{\partial x_n}\qquad &n>0 \\
x_n\qquad &n<0
\end{cases}
$$
Recall we completed $\mathcal{F}$ to
$\overline{\mathcal{F}}$
and discovered
$$
\xymatrix@R=0.2 em{
V\ar[r]^{\simeq}&\overline{\mathcal{F}}\\
a(z)|0\rangle |_{z=0}&a(z)\ar@{|->}[l]^s\\
a \ar@{|->}[r]_Y&Y(a,z)
}
$$
For $I_V \in \overline{\mathcal{F}}$,
where
 \begin{align*}
I_V&=\sum I_{V(n)}z^{-n-1}\\
I_{V(n)}&=\begin{cases}
I_V\qquad &\text{if }n=-1 \\
0\qquad &\text{if }n \neq -1
\end{cases}
\end{align*}
we get $I_V |0\rangle|_{z=0}=|0\rangle$. So,
$$
Y(|0\rangle,z)=I_V
$$
This shows that the vacuum $|0\rangle$
corresponds to the identity.

Next $h(z)$ in our example. We have
$$
\begin{aligned}
h(z)|0\rangle&=\sum_{ n \in \mathbb{Z}}h_n
1z^{-n-1}\\
&=\sum_{ k \geq 1}x^kz^{k-1}+\sum_{k \geq
1}k\frac{\partial }{\partial x_k}1\\
\implies h(z)|0\rangle|_{z=0}&=x_1\\
\implies S(h(z))&=x_1\\
\implies Y(x,z)&=h(z).
\end{aligned}
$$
So, $x_1$ corresponds to the sort of
``generating field'' $h(z)$.

What about $x_2$? Well,
$h(z)=h(z)_{(-1)}I_V$.
There was a lemma (maybe Remark
\ref{vertex-algebras-remark-nth-product}?) that
$$
a(z)_{(-2)}I_V=\partial_za(z).
$$
So $\partial_zh(z) \in
\overline{\mathcal{F}}$,
$$
\begin{aligned}
a(\partial_zh(z))&=\partial \sum_{k \geq
1}x_kz^{k-1}\Big|_{z=0}\\
&=\sum_{k\geq 1}(k-1)x_kz^{k-2}\Big|_{z=0}\\
&=x_2.
\end{aligned}
$$
So
$$
Y(x_{n+1},z)=\partial_z^{(n)}h(z).
$$
Continuing:
$$
Y(x_3,z)=\frac{1}{2}\partial_z^2h(z),
$$
and in general
$$
Y(x_{n+1},z)=\partial_z^{(n)}h(z).
$$
The next question is: what about $x_1^2$?
The trick here is that $x_1$ is identified
with $h$, so
$$
\begin{aligned}
Y(x_1^2,z)&=Y(x_1\cdot x_1,z)\\
&=Y(h_{(-1)}x_1,z)\\
&=Y(x_1,z)_{(-1)}Y(x,z)\\
&=\,:\!h(z)h(z)\!:\,.
\end{aligned}
$$

Recall that $L(z)=\frac{1}{2}
\,:\!h(z)h(z)\!:\,$. We have claimed
(Exercise \ref{vertex-algebras-exercise-hard})
that $[L_m,L_n]=(m-n)L_{m+n}$ when $m,n,m+n
\neq 0$
where
$$
\begin{aligned}
L(z)&=\sum_{m \in \mathbb{Z}}L_mz^{-m-2}\\
&=\sum_{n \in \mathbb{Z}}L_{(n)}z^{-n-1}.
\end{aligned}
$$

\medskip\noindent
What about the $\lambda$-bracket? Let's
denote
$h=x_1$, (so $Y(h,z)=h(z)$,
$L=\frac{1}{2}x_1^2$
and $Y(L,z)=L(z)$).

Well, long ago we computed that
$$
\begin{aligned}
[h(z),h(w)]&=\partial_w \delta(z,w)I_V\\
&=\sum_{j \geq
0}Y(h_{(j)}h,w)\partial_w^{(j)}\delta(z,w).
\end{aligned}
$$
Then
$$
\begin{aligned}
Y(h_{(0)}h,w)&=0\\
Y(h_{(1)}h,w)&=I_V\\
Y(h_{\geq 2)}h,w)&=0\\
\implies h_{(1)}h&=|0\rangle, \text{ and }\\
h_{(j)}h=0,\qquad j=0\text{ and }j\geq 2.
\end{aligned}
$$
So we see that
$$
[h_\lambda h]=\lambda |0\rangle.
$$

We often omit  ``$|0\rangle$'' from these
computations,
as if it were ``1''.

So
$$
[h_\lambda h]=\lambda.
$$
\medskip\noindent
$[h_\lambda |0\rangle]=0=[|0\rangle_\lambda
h]$
because $[h(z),I_V]=0$.
In fact $[x_\lambda
|0\rangle]=[|0\rangle_\lambda x]=0$ for all
$x \in V$.

\medskip\noindent
What about $h_{(0)}h$? Well, $h_{(0)}=0$ by
definition of $h(z)$.
So $h_{(0)}h=0$. Next,
$h_{(j)}=j\frac{\partial }{\partial x_j}$ for
$j>0$, so
$$
h_{(1)}h= \frac{\partial }{\partial
x_1}(x_1)=1=|0\rangle.
$$
So $[h_\lambda h]=\lambda$ again.

What I really want is $[L_\lambda L]=$? Let's
use
Borcherd's formula
\ref{vertex-algebras-equation-Borcherds-identity}.

First let's compute $(h_{(-1)}h)_{(0)}L$,
i.e. we
are putting $n=-1$, $m=0$ and $k=0$.

$$
[\text{missing computations}]
$$
\begin{exercise}
\label{vertex-algebras-exercise-conta}
Check $(h_{(-1)}h)_{(1)}h=2h$ and
$(h_{(-1)}h)_{\geq 2}h=0$.
Thus $[L_\lambda h]=Th+\lambda h$.
\end{exercise}
If we wanted, we could continue to calculate
$[L_\lambda L]$ by putting
$c=\frac{1}{2}x_1^2$ above
instead of $c=x_1$.

We find
$$
[L_\lambda L]=\underbrace{TL}_{L_{(0)}L}
+\underbrace{2\lambda L}_{L_{(1)}L}
+\underbrace{\frac{\lambda^3}{12}|0\rangle}
_{L_{(3)}L}.
$$
This means
\begin{equation}
\label{vertex-algebras-equation-looks-important}
[L(z),L(w)]=(\partial_wL(w))\delta(z,w)
+2L(w)\partial_w\delta(z,w)
+\frac{1}{12}\partial_w^3\delta(z,w)I_V.
\end{equation}
So we can extract coefficients:
$$
LHS=\sum_{m,n}[L_m,L_n]w^{-m-2}z^{-n-2}
$$
Coefficient of $w^{-m-2}z^{-n-2}$ on RHS is:
$$
\begin{aligned}
&\sum_{k}(-k-2)L_kw^{-k-3}\sum z^{-a-1}w^a\\
&+2\sum_{k}L_kw^{-k-2}\sum_{a}az^{-a-1}w^{a-
1}
+\frac{1}{12}\sum_{a}a(a-1)(a-2)z^{-a-1}w^{a-
3}.
\end{aligned}
$$
So actually we found
$$
[L_m,L_n]=(m-n)L_{m+n}+\frac{m^3-m}{12}
\delta_{m,-n}I_V.
$$
The moral is: one tries to define
$L=\frac{1}{2}h^2$ naively,
expecting $[L_m,L_n]=(m-n)L_{m+n}$.
Need to ``normally order/normalise'' to make
$L$
well-defined, now $L=\frac{1}{2}\,:\!h
h\!:\,$.
The ``cost'' is that the expected relation
$[L_m,L_n]=(m-n)L_{m+n}$ gets altered by an
``anomaly'' $\frac{m^3-m}{12}I_V$.
\end{example}

Right, so we showed how to compute this using
Bochner's
formula, but there's actually another way
using $\lambda$-bracket:

\begin{theorem}
\label{vertex-algebras-theorem-lambda-brackets}
Let $V$ be a vertex algebra, $a,b,c \in V$.
Then
\begin{align}
\label{vertex-algebras-equation-lambda-brackets}
[Ta_\lambda b]&=-\lambda[a_\lambda b]\\
[a_\lambda Tb]&=(T+\lambda)[a_\lambda b]\\
T(\,:\!ab\!:\,)&=\,:\!(Ta)b\!:\,+\,:\!a(Tb)
\!:\,\\
[b_\lambda a]&=-[a_{-\lambda -T}b]\\
[a_\lambda \,:\!bc\!:\,]&=\,:\![a_\lambda
b]c\!:\,+\,:\!b[a_\lambda c]\!:\,
+\int_0^\lambda [[a_\lambda b]_\mu c]d\mu\\
[a_\lambda[b_\mu c]]-[b_\mu [a_\lambda
c]]&=[[a_\lambda b]_{\lambda+\mu}c].
\end{align}
\end{theorem}

\begin{proof}
For the first one use $\partial_z
\delta(z,w)=-\partial_z \delta(z,w)$.

For the last one notice that
$$
\begin{aligned}
-[a_{-\lambda -
T}b]&=\sum_{j}\frac{1}{j!}(-\lambda-T)
^j(a_{(j)}b)
\end{aligned}
$$
if $[a_\lambda b]
=\sum_{j}\frac{\lambda^j}{j!}b$.
\end{proof}

\begin{remark}
\label{vertex-algebras-remark-how-to-use}
By noting that $[a_\lambda
b]=\text{Res}_ze^{\lambda z}Y(a,z)b$,
these identities can be proved pretty
efficiently.
\end{remark}

That is,

\begin{exercise}
\label{vertex-algebras-exercise-use-the-5th}
Use the fifth equation to prove what we
proved with Bochner's formula.
\end{exercise}

And another one:

\begin{exercise}
\label{vertex-algebras-exercise}
Let $V=H=\mathbb{C}[x_1,x_2,\ldots]$ again.
Define $B=\frac{1}{2}\,:\!h h\!:\,+\beta Th$
with $\beta \in \mathbb{C}$.
Confirm that
$$
[B_\lambda B]=TB+2\lambda
B=\frac{1}{12}(1-12\beta^2)|0\rangle.
$$
\end{exercise}

\begin{definition}
\label{vertex-algebras-definition-Virasoro-Lie-algebra}
The {\it Virasoro Lie algebra} is
$$
\text{Vir}=\bigoplus_{n \in
\mathbb{Z}}\mathbb{C}L_n \oplus \mathbb{C}C
$$
with Lie bracket
$$
[L_m,L_n]=(m-n)L_{m+n}+\frac{m^3-m}{12}
\delta_{m,-n}C,\qquad
[C,\text{Vir}]=0.
$$
\end{definition}
By the discussion above (all we've said so
far!),
there is a representation $\rho$ of
$\text{Vir}$ in $H$
given by
$$
\rho(\underbrace{L_n}_{\substack{\text{as an
abstract} \\
\text{object in Vir}}})
=\underbrace{L_n}_{\substack{\text{as a
complicated} \\ \text{operator}}},
\qquad \rho(C)=I_H.
$$
Also $H$ carries a representation
$\rho_\beta$ of $\text{Vir}$,
$$
\rho_\beta(L_n)=B_n,\qquad
\rho_\beta(C)=(1-12\beta^2)I_V
$$
If $(M,\rho)$ is a representation of
$\text{Vir}$, in which
 $\rho(C)=c\cdot I_M$ for some scalar $c$, we
 say $M$ has a
{\it central charge} $c$.

Looks like we have a second example of a
vertex algebra.
Consider $V=H$, $|0\rangle=1$, $T=T$ from
before, but this time put
$F=\{B(z)\}$. Then all the axioms of the
first definition of vertex
algebra are satisfied, except the second:
$B_{(n_1)}B_{(n_2)}\ldots B_{(n_s)}|0\rangle$
might not span all of $V=H$.
So take $V=\text{span}\{\text{these
monomials}\}\subset H$.
One can check that $T(V) \subset V$, and
$B(z)$ is a quantum field on $V$.
So $(V,1,T,\{B(z)\})$ is a vertex algebra,
called the
{\it Virasoro vertex algebra}.

\section{The charged free fermions: a vertex
superalgebra}
\label{vertex-algebras-section-charges-free-fermions}

Also known as:
\begin{itemize}
\item The Clifford (vertex) algebra.
\item The Dirac sea.
\item The bc system.
\end{itemize}

First we consider a vector
superspace with basis $\varphi$ and
$\varphi^*$
$$
\mathfrak{a}=\mathbb{C}\varphi +\mathbb{C}
\varphi^*
$$
(where $\varphi$ and $\varphi^*$ are odd,
so this is a O,2-dimensional superspace).

(A vector superspace is a
$\mathbb{Z}/2$-graded vector space,
i.e. a vector space split in two pieces,
$V=V_0\oplus V_1$, which we call even and
odd.)

We consider a bilinear form
$$
\begin{aligned}
\left\langle{}\cdot,\cdot\right\rangle{}: \mathfrak{a}\times
\mathfrak{a}
&\longrightarrow  \mathbb{C}\\
\left\langle{}\varphi^*,\varphi\right\rangle{} &=1\\
\left\langle{}\varphi,\varphi^*\right\rangle{}&=1\text{ (why?
see below)}
\end{aligned}
$$
Now consider
$$
\tilde{\mathfrak{a}}=\underbrace{t^{1/2}
\mathfrak{a}[t,t^{-1}]}_{\text{odd}}
\oplus
\underbrace{\mathbb{C}K}_{\text{even}}.
$$
with Lie bracket
$$
\begin{aligned}
[at^m,b^n]&=\delta_{m,-n}\left\langle{}a,b\right\rangle{}K,\\
[K,\tilde{\mathfrak{a}}]&=0.
\end{aligned}
$$
(Notice the small difference from the
Heisenberg case: there is no $m$ before the
$\delta$!)

To see why
$\left\langle{}\varphi,\varphi^*\right\rangle{}=1$,
notice that in a Lie superalgebra we always
want
$[x,y]=(-1)(-1)^{p(x)p(y)}[y,x]$
where $p$ denotes the parity.
Since both  $at^m$ and $bt^n$ are odd,
we have $[at^m,bt^n]=[bt^n,at^m]$.
Then apply the definition of bracket.

Since $a,b$ odd here,
we say $\left\langle{}a,b\right\rangle{}=\left\langle{}b,a\right\rangle{}$
means
$\left\langle{}\cdot,\cdot\right\rangle{}$ is skew-super
symmetric.

\begin{definition}
\label{vertex-algebras-definition-supersymmetric}
A bilinear form
$$
(\cdot,\cdot):U \times U \to \mathbb{C}
$$
is
{\it supersymmetric} if
$(b,a)=(-1)^{p(a)p(b)}(a,b)$,
and {\it skew-supersymmetric} if
 $(b,a)=-(-1)^{p(a)p(b)}(a,b)$.
\end{definition}

\begin{exercise}
\label{vertex-algebras-exercise-Lie-superalgebra}
Invent the definition of Lie superalgebra,
and confirm that if $U=U_0 \oplus U_1$
is a vector superspace,
$\text{End}(U)$ is a Lie superalgebra
with $[X,Y]:=XY-(-1)^{p(X)p(Y)}YX$.
\end{exercise}

\medskip\noindent
Let's build a Fock-type representation of
$\tilde{\mathfrak{a}}$.

Construction 1.
$$
\tilde{\mathfrak{a}_+}=
\bigoplus_{n>0}t^n\mathfrak{a}
\oplus\mathbb{C}K,
\qquad (n \in \frac{1}{2}+\mathbb{Z},\text{
recall})
$$
$\tilde{\mathfrak{a}_+}\subset
\tilde{\mathfrak{a}}$ is
a superalgebra.

Consider $U(\tilde{\mathfrak{a}})$
and $U(\tilde{\mathfrak{a}})\subset
U(\tilde{\mathfrak{a}})$.

Recall that the PWB theorem says
that $U(\mathfrak{g})$ ``looks like''
polynomials on $\mathfrak{g}$.

When $\mathfrak{g}$ is pure even,
this says:
take $\{a_1,a_1,\ldots\}$ a basis of
$\mathfrak{g}$,
then
$$
\{a_{i_1}^{m_1},a_{i_2}^{m_2}\ldots
a_{i_s}^{m_s}:
i_1\leq i_2\leq \ldots\leq i_s,m_j \geq 1\}.
$$
For $\mathfrak{g}=\mathfrak{g}_0 \oplus
\mathfrak{g}_1$
super Lie algebra,
the statement is similar,
only now, we take each $a^i$ homogeneous
($\in \mathfrak{g}_0$ or $\mathfrak{g}_1$):
$$
\{a_{i_1}^{m_1}a_{i_2}^{m_2}\ldots
a_{i_s}^{m_s},
i_1\leq i_2 \leq \ldots \leq i_s, m_j\geq 1,
\text{ if }a_{i_j}\text{ odd, }m_j=1\}.
$$
The point is that if we put an odd guy
twice it becomes one.

Why? Because in $U(\mathfrak{g})$, $X \in
\mathfrak{g}_1$ should
satisfy $X X + X X=[X,X]=0$.
So  $X^2=0$. So we only allow
odd basis vectors to appear 0 or 1 times
each.

\medskip\noindent
Dfine
$F=U(\tilde{\mathfrak{a}})
\otimes_{U(\tilde{\mathfrak{a}}_+)}
\mathbb{C}|0\rangle$,
where $at^n\cdot|0\rangle$ for all
$n>\frac{1}{2}$, and
$K|0\rangle=|0\rangle$.

[Picture of $F$]

We can introduce a bi-grading of $F$:
$$
F=\bigoplus_{\substack{c \in \mathbb{Z} \\
\Delta \in \frac{1}{2}\mathbb{Z}_+}}
F_\Delta^c
$$
where, for a monomial
$$
\varphi^*_{-n}\varphi^*_{-n_2}\ldots
\varphi_{-n_s}^*
\varphi_{-m},\varphi_{-m_2}\ldots
\varphi_{-m_t}|0\rangle
$$
we define its {\it energy} by
$\Delta=\sum_in_i+\sum_jm_j$
and its {\it charge} by $c=s-t$.

The point is that $F$ is an
infinite-dimensional
vector space, but we get a reasonable mental
picture of what it looks like.

What is $\dim F_\Delta^c$? Let's write…

\medskip\noindent
So we have a vector superspace $F$,
and an action on $F$ of
$\tilde{\mathfrak{a}}$.
Let's build some quantum fields!

There are going to be two: $\varphi(z)$
and $\varphi^*(z)$.
$$
\varphi(z)=\sum_{n \in
\frac{1}{2}+\mathbb{Z}}\varphi_nz^{-n-
\frac{1}{2}}
$$
(Notice the exponent on $z$ is an integer!)
$$
\varphi^* (z)=\sum_{n \in
\frac{1}{2}+\mathbb{Z}}
\varphi^*_n z^{-n-\frac{1}{2}}.
$$
Now let's compute
$[\varphi^*(z),\varphi(w)]$.

$$
\begin{aligned}
[\varphi^*(z),\varphi(w)]&=\sum_{m,n \in
\frac{1}{2}+\mathbb{Z}}
[\varphi^*_m,\varphi_n]z^{-m-\frac{1}{2}}w^{-
n-\frac{1}{2}}\\
&=\sum_{n \in
\frac{1}{2}+\mathbb{Z}}Kz^{n-\frac{1}{2}}w^{-
n-\frac{1}{2}}\\
&=\sum_{\substack{r \in \mathbb{Z} \\
r=n-\frac{1}{2}}}z^rw^{-r-1}\\
&=\delta(z,w).
\end{aligned}
$$
So
$F,|0\rangle,\mathcal{F}=\{\varphi(z),
\varphi^*(z)\}$,
$T:F\to F$, is a vertex superalgebra.
Here $T:F \to F$ has to satisfy
$T|0\rangle=0$ and
$[T,\varphi(z)]=\partial_z\varphi(z)$,
$$
\begin{aligned}
\sum [T,\varphi_n]z^{-n-\frac{1}{2}}&=
\sum (-n-\frac{1}{2})\varphi_n
z^{-n-\frac{3}{2}}\\
&=\sum
(-k+\frac{1}{2})\varphi_{k-1}z^{-k-\frac{1}
{2}}
\end{aligned}
$$
after putting $n=k-1$. So
$$
[T,\varphi_n]=-(n-\frac{1}{2})\varphi_{n-1}.
$$
Inductively, this determines
how $T$ acts on $F$.

\medskip\noindent
In $\lambda$-bracket notation, this just says
$$
[\varphi^* _\lambda\varphi]=|0\rangle.
$$
Skew-symmetry says $[b_\lambda
a]=-(-1)^{p(a)p(b)}[a_{-\lambda -T}b]$.

For example,
$$
\begin{aligned}
[h_\lambda h]&=\lambda |0\rangle,\\
[h_\lambda
h]&=-[h_{-\lambda-T}h]=-(-\lambda-T)
|0\rangle\\
&=T\lambda |0\rangle.
\end{aligned}
$$
In this case $[\varphi_\lambda \varphi^*
]=+|0\rangle$ too.

Notice $\sum \varphi_n z^{-n-\frac{1}{2}}=
\sum \varphi_{(m)}z^{-m-1}$.

So
$$
\begin{aligned}
\varphi_{(m)}&=\varphi_{m-\frac{1}{2}}\text{
and}\\
\varphi^*_{(m)}&=\varphi_{m-\frac{1}{2}}.
\end{aligned}
$$
So
$$
\begin{aligned}
\,:\!\varphi^*\varphi\!:\,&=\varphi_{(-1)}
^*\varphi_{(-1)}|0\rangle\\
&=\varphi_{-\frac{1}{2}}^*\varphi_{-\frac{1}
{2}}|0\rangle.
\end{aligned}
$$
I.e.
$$
\varphi=\varphi_{(-1)}|0\rangle=\varphi_{-
\frac{1}{2}}|0\rangle
$$
Now we can put the former picture but with
the
associated operators using the
second definition of vertex algebra:
[Picture]

Let's compute some brackets: [lots of
computations]

So,
$$
[\alpha_\lambda \alpha]=\lambda |0\rangle.
$$
Just like $[h_\lambda h]=\lambda |0\rangle$.
This relation is the basis of the
``boson-fermion correspondence'', i.e.
consider the subspace $U \subset F$ defined
as
 $$
U=\text{span}\{\alpha_{n_1}\alpha_{n_2},
\ldots,\alpha_{n_s}|0\rangle
:n_{i} \in \mathbb{Z}\}.
$$

We claim that since $[T,\alpha(z)]=\partial_z
\alpha(z)$,
$[T,\alpha_n]=-n\alpha_{n-1}$.

So $T(U) \subset U$. Put
$\mathcal{F}_1=\{\alpha(z)\}$.
Then $(U,|0\rangle,T|_{U},\{\alpha(z)\})$ is
a vertex algebra
(an honest one, not super).

\begin{proposition}
\label{vertex-algebras-proposition-isomorphic-to-Heisenberg}
$U$ is isomorphic to the Heisenberg
vertex algebra. (In particular,
$U \simeq \mathbb{C}[x_1,x_2,\ldots]$
as a vector space.)
\end{proposition}

\begin{proof}
Next time.
\end{proof}

Since $[\alpha_0,\varphi_n]=-\varphi_n$,
$[\alpha_0,\varphi_m^*]=\varphi_m^*$,
and $\alpha_0 |0\rangle=0$, in fact the
charge of a monomial is exactly the
eigenvalue of $\alpha_0$ acting on it.

$$
\begin{aligned}
\alpha_0\cdot\varphi^*_{-m}\varphi^*_{-m_2}
\varphi_{-n}|0\rangle
&=[\alpha_0,\varphi^*_{-m},\varphi^*_{-m_2},
\varphi_{-n}]|0\rangle\\
&=\underbrace{(+1+1-1)}_{\text{charge}}
\varphi_{-m_1}^*\varphi_{-m_2}\varphi_{-n}
|0\rangle.
\end{aligned}
$$
From $[\alpha_\lambda \alpha]=\lambda
|0\rangle$,
we get
$$
[\alpha_m,\alpha_n]=m\delta_{m,-n}I_F\text{
(as for $h$)}
$$
So $[\alpha_0,\alpha_n]=0$ for all $n \in
\mathbb{Z}$.
Thus  $\alpha_0 |_{ U}=0$ and thus $U \subset
F^{(0)}$.

\begin{proposition}
\label{vertex-algebras-proposition-U-is-F0}
$U=F^{(0)}$.
\end{proposition}

\begin{proof}
Next time.
\end{proof}

\medskip\noindent
Let's take two copies of $F$,
i.e. $F_2=F \otimes F$
with fields
$\varphi^1(z),\varphi^2(z),\varphi^{1*}(z),
\varphi^{2*}(z)$
$$
[\varphi^{i*}_{\lambda}\varphi^j]=\delta_{ij}
.
$$
Basis of $F_2$ is
$$
\varphi_{-n_1}^1\ldots
\varphi_{-n_2}^1\varphi_{-m_1}^2\ldots
\varphi_{-m_t}^2\varphi^{*-1}_{-p_1}\ldots
\varphi_{-p_u}^{*-1}\varphi_{-q_1}^{*2}
\ldots\varphi_{-q_v}^{*2}|0\rangle.
$$
Let
$\alpha^{ij}=\,:\!\varphi^{j*}\varphi^i\!:\,
$.
$$
[\alpha^{ij}_\lambda \alpha^{k\ell}]=?
$$
(certainly $[\alpha^{ii}_\lambda
\alpha^{ii}]=\lambda |0\rangle$.)
In general
$$
[\varphi^k_\lambda
\alpha^{ij}]=\text{computation}=\delta_{ki}
\varphi^j.
$$
Similarly
$$
[\varphi^{k*}_\lambda
\alpha^{ij}=-\delta_{kj}\varphi^{i*}.
$$
$$
\begin{aligned}
[\alpha^{ij}_\lambda
\varphi^k]&=-\delta_{ik}\varphi^j\\
[\alpha^{ij}_\lambda
\varphi^{k*}]&=+\delta_{j k}\varphi^{i*}.
\end{aligned}
$$

We would like to consider
$$
F_n=F^{\otimes n}=F \otimes \ldots \otimes F
$$
where on the $i$-th factor we have
$\varphi_i,\varphi^*_i$.

Then we have relations among the generators:
$$
[\varphi^i_\lambda \varphi_j^* ]=[\varphi^*
_j\lambda\varphi_i]
=\delta_{ij}|0\rangle
$$
Today let's define
$$
\alpha_{ij}=\,:\!\varphi_i \varphi^*_j\!:\,
$$
We notice that
\begin{itemize}
\item the $\varphi_i$ behave like $e_i$,
\item the $\alpha_{ij}$ behave like $E_{ij}$,
\item the $\varphi_i^*$ behave like $e_i^*$.
\end{itemize}

$$
\text{[Computations of some
$\lambda$-brackets]}
$$

Recall, the $E_{ij}$ span a Lie algebra:
$\mathfrak{gl}_n$.
It has a representation on
$\mathbb{C}^n=\left\langle{}e_1,\ldots,e_n\right\rangle{}$
by $E_{ij}e_k=\delta_{ik}e_i$,
and on $(\mathbb{C}^n)^*$ we have
$$
[\text{computations}]
$$
$$
E_{ij}e_k^* =-\delta_{ki}e_j^*.
$$
Let $A \in \mathfrak{gl}_n$. Let's write
$$
\alpha^A=\sum_{i,j}a_{ij}\alpha_{ij}\in F_n
$$
We'd like to compute $[\alpha^A_\lambda
\alpha^B]$.

Some computations done in lecture are the
proof of

\begin{theorem}
\label{vertex-algebras-theorem-alpha-is-like-matrices}
For $\alpha_{ij}=\,:\!\varphi_i
\varphi_j^*\!:\, \in F_n$,
and $\alpha^A$ as above,
$$
[\alpha^A_\lambda\alpha^B]=\alpha^{[A,B]}
+\lambda\text{Tr}(AB)|0\rangle.
$$
\end{theorem}

Which says that $\alpha^A$ behaves like
matrices,
but with that correction term.

\medskip\noindent
Nex, let $\mathfrak{g}\subset
\mathfrak{gl}_n$
be a Lie subalgebra.
For $A,B \in \mathfrak{g}$, $[A,B] \in
\mathfrak{g}$
also, and so the set
$$
\mathcal{F}=\{\alpha^A:A \in \mathfrak{g}
\cup \{|0\rangle\}\}
$$
is ``closed under $\lambda$-brackets''.

More precisely, for
$$
F_n \supset
V=\text{span}\{\alpha^{A_1}_{(n_1)}
\ldots\alpha^{A_s}_{(n_s)}:
n_i \in \mathbb{Z}, A_i \in \mathfrak{g}\},
$$
with $|0\rangle=|0\rangle$, $T=T$ and
$\mathcal{F}=\{\alpha^A(z):A \in
\mathfrak{g}\}\cup \{I\}$
is a vertex algebra.

\section{Universal affine vertex algebra}
\label{vertex-algebras-section-unive-affin-verte-algeb}

Recall the construction of the affine vertex
algebra,
Definition
\ref{vertex-algebras-definition-affine-Lie-algebra}.

Let $(M,\rho)$ be a representation of
$\hat{\mathfrak{g}}$.

For $a \in \mathfrak{g}$, define
\begin{equation}
\label{vertex-algebras-equation-quant-field-unive-verte}
a^M(z)=\sum_{n \in
\mathbb{Z}}\rho(at^n)z^{-n-1}
\in\text{End}(M)[\![z^{\pm 1}]\!].
\end{equation}
\begin{definition}
\label{vertex-algebras-definition-smooth-g-module}
A {\it smooth $\hat{\mathfrak{g}}$-module}
is a $\hat{\mathfrak{g}}$-module $(M,\rho)$
such that for each $m \in M$ there is $N \in
\mathbb{Z}$
such that $\rho(at^n)m=0$ for all
$a \in \mathfrak{g}$ and $n \geq N$.
\end{definition}

Notice that on a smooth module $M$,
the fields $a^M(z)$ are quantum fields.

We may calculate:
$$
\begin{aligned}
[a^m(z),b^m(w)]&=\sum_{m,n}[\rho(at^m),
\rho(bt^n)]z^{-m-1}w^{-n-1}\\
&=\sum_{m,n}z^{-m-1}w^{-n-1}
([a,b]t^{m+n}+m\delta_{m,-n}(a,b)\rho(K))\\
&=\sum_{m,n}([a,b]t^{m+n}w^{-(m+n)-1})z^{-m-
1}w^m
+\sum_{m,n}z^{-m-1}w^{-n-1}\delta_{m,-n}(a,b)
\rho(K)\\
&=[a,b]^m(w)\delta(z,w)
+(a,b)\partial_w\delta(z,w)\rho(K).
\end{aligned}
$$

\noindent
Now suppose the $K$ acts by some constant,
i.e.
$\rho(K)=k I_M$.
Then our quantum fields $a^M(z)$ are
mutually local, with exponent 2, i.e.
$(z-w)^2[a^M(z),b^M(w)]=0$.

\medskip\noindent
As a special case, we may take
$M$ to be a sort of ``Fock space'':
$$
\hat{\mathfrak{g}}_+=\mathfrak{g}[t]\oplus
\mathbb{C}K\subset
\hat{\mathfrak{g}}.
$$
Let $\hat{\mathfrak{g}}_+$
act on $\mathbb{C} |0\rangle$ by
$$
\begin{aligned}
at^m |0\rangle&=0 \qquad \forall m \geq 0,\\
K\cdot |0\rangle&=k |0\rangle\qquad
(\text{$k \in \mathbb{C}$ called {\it
level}}).
\end{aligned}
$$
Consider
$$
V^k(\mathfrak{g})=U(\hat{\mathfrak{g}})
\otimes_{U(\hat{\mathfrak{g}}_+)}
\mathbb{C}|0\rangle.
$$

\begin{remark}
\label{vertex-algebras-remark-heise-speci-case-smoot-const}
$H=\mathbb{C}[x_1,x_2,\ldots]$ is an example
of this construction with
$\mathfrak{g}$ one-dimensional and $(h,h)=1$.
\end{remark}

\begin{exercise}
\label{vertex-algebras-exercise-Vk-is-smooth}
$V^k(\mathfrak{g})$ is a smooth
$\hat{\mathfrak{g}}$-module,
in which $K \mapsto  k \text{Id}$.
\end{exercise}

\begin{proof}
To prove  $V^k(\mathfrak{g})$ is smooth
we need to show that for every
formal power series of endomorphisms of
$V^k(\mathfrak{g})$ of the form Equation
\ref{vertex-algebras-equation-quant-field-unive-verte},
the coefficient operators vanish at every $v$
for sufficiently
large $n$. But to define these
power series we need to find a copy
of  $\hat{\mathfrak{g}}$ inside
$V^k(\mathfrak{g})$.
\end{proof}

The set $\mathcal{F}=\left\{a(z)=\sum_{n \in
\mathbb{Z}}
(at)^nz^{-n-1}:a \in \mathfrak{g}\right\}$
are mutually local quantum fields.

Define $T:V^k(\mathfrak{g}) \to
V^k(\mathfrak{g})$
by the relation $T |0\rangle=0$
and $[T,at^m]=-mat^{m-1}$.
$V^k(\mathfrak{g})$ is a vertex algebra
called the {\it universal affine vertex
algebra}
of level $k$ associated with $\mathfrak{g}$.

[Picture of $V^k(\mathfrak{g})$ for
$\mathfrak{g}=\mathfrak{sl}_2$.]

We have a bigrading, vertical grading
$\Delta$ (energy) is
$$
\Delta(a^1_{(-n_1)}a^2_{(-n_2)}\ldots
a^s_{(-n_s)}|0\rangle
=n_1+n_2+\ldots+n_s.
$$

Notice that for $a \in \mathfrak{g}$, we have
$$
\begin{aligned}
\mathcal{F}\ni a(z)&\overset{s}{\mapsto }
a(z)|0\rangle |_{z=0}\\
&=\sum_{m}at^m |0\rangle
z^{-m-1}\Big|_{z=0}\\
&=(at^{-1})|0\rangle+(at^{-2})z|0\rangle+
\ldots\\
&=at^{-1} |0\rangle.
\end{aligned}
$$
Thus we can think of
$\{a t^{-1} |0\rangle:a \in
\mathfrak{g}\}\subset V^k(\mathfrak{g})$
as a ``copy'' of $\mathfrak{g}$ inside
$V^k(\mathfrak{g})$.
Let's {\bf abuse notation} and write $a$ for
$at^{-1}|0\rangle$.

Then $at^{-2}|0\rangle=Ta$, also
$(at^{-1})(bt^{-1})|0\rangle=\,:\!ab\!:\,$,
etc.
$e=et^{-1}|0\rangle=e_{(-1)}|0\rangle$,
$\Delta(e)=1$.

In fact $\Delta(a)=1$ for all $a \in
\mathfrak{g}$,
$\Delta(Ta)=2$,
$\Delta(\,:\!ab\!:\,)=2$, etc.

In this example, i.e.
$\mathfrak{g}=\mathfrak{sl}_2$,
$\mathfrak{sl}_2$ is itself a
$\mathbb{Z}$-graded Lie algebra,
\begin{align*}
\begin{aligned}
w(E)&=2\\
W(H)&=0\\
w(F)&=-2
\end{aligned}
\qquad
\begin{aligned}
[H,E]&=2E\\
[H,H]&=0\\
[H,F]&=-2F.
\end{aligned}
\end{align*}

(That is, $\mathfrak{g}=\bigoplus_{n \in
\mathbb{Z}}\mathfrak{g}_n$,
$[\mathfrak{g}_m,\mathfrak{g}_n]=\mathfrak{g}
_{m+n}$.)
$w=$ eigenvalue of $\text{ad}(H)$.
This induces a $\mathbb{Z}$-grading on
$V^k(\mathfrak{sl}_2)$,
compatible, i.e.
$$
w(a^1_{(-n_1)}\ldots a^s_{(-n_s)}|0\rangle=
2\#(E)-2\#(F).
$$

\medskip\noindent
As for the character,
$$
\begin{aligned}
\chi(q,u)&=\sum_{\Delta,w}\dim
V^k(\mathfrak{g})_\Delta^w q^\Delta u^w\\
&=1+q(u^2+1+u^{-2})
+q^2(u^4+2u^2+2+2u^{-2}+u^{-4}+\ldots
\end{aligned}
$$
Then by the generating function argument we
have discussed
before,
$$
\chi(q,u)=\prod_{n=1}^\infty \frac{1}
{(1-u^2q^n)(1-q^n)(1-u^{-2}q^n)}.
$$
Can discard $u$ to get vertical grading
$$
\chi(q)=\prod_{n=1}^\infty
\frac{1}{(1-q^n)^3}.
$$

\medskip\noindent
[Missing: we changed a little the definition
of $V^k(M)$;
now it also depends on a bilinear form $B$.]

Here's another example:

\begin{example}
\label{vertex-algebras-example-fermi-unive-verte-algeb}
$\mathfrak{g}=\mathfrak{gl}_n$, with
$B(X,Y)=\text{Tr}(XY)$,
and $k=1$.
\end{example}

\begin{remark}
\label{vertex-algebras-remark-resca-killi-form-coxet-numbe}
In general
$V^1(\mathfrak{g},kB)=V^k(\mathfrak{g},B)$.
If $\mathfrak{g}$ is finite-dimensional
simple, we
typically take $B(a,b)=\frac{1}{2
h^\vee}\kappa(a,b)$
where $\kappa$ is the Killing form. (This $B$
has the property that $B(\theta,\theta)=2$
for {\bf long} roots $\theta \in \Delta
\subset \mathfrak{h}^*$.)
(In fact, in a simple f.d. there's
all invariant forms are proportional.)
Then we write this vertex algebra
as $V^k(\mathfrak{g})$ without specifying
$B$.
\end{remark}

So perhaps that's why last time we didn't put
$B$.

\medskip\noindent
Back to the Fermion algebra $F_n$, and now
denoting $J$ instead of $\alpha$,
we built fields
$\{J^A:A \in \mathfrak{gl}_n\}$ with
$\lambda$-bracket
$[J^A_\lambda
J^B]=J^{[A,B]}+\lambda\text{Tr}(AB)
|0\rangle$.
This suggests a relation with
$V^1(\mathfrak{gl}_n,B)$,
where $B(X,Y)=\text{Tr}(X,Y)$.

Consider $\mathcal{F}=\{J^A(z):A \in
\mathfrak{gl}_n\}\subset \mathbb{F}_{F_n}
\cong F_n$.
Let $V=\text{span}\{J^{A_1}_{(n_1)}\ldots
J^{A_s}_{(n_s)}|0\rangle:
A_i \in \mathfrak{gl}_n,n_i \in
\mathbb{Z}\}\subset F_n$.

Then $(V,\mathcal{F},T|_{V},|0\rangle)$ is a
vertex algebra.
$V \subset F_n$. Notice that $V \neq F^n$
since $V$ consists only of even fields,
indeed, $V=\left\langle{}\,:\!\varphi_i
\varphi_j^*\!:\,\right\rangle{}$.

Here $[T,\varphi_{(n)}]=-n\varphi_{(n-1)}$.
$$
T(\varphi^{i_1}_{(n_1)}\ldots\varphi^{i_s}
_{(n_s)}|0\rangle)
=-\sum_{j=1}^s n_j\varphi_{(n_1)}^{i_1}
\ldots vo_{(n_j-1)}^{i_j}\ldots
\varphi_{(n_s)}^{i_s}|0\rangle.
$$
Does $V=F^{(0)}_n$ then?

Let's examine $n=1$ first.
$\mathfrak{g}=\mathbb{C}1$,
$B(1,1)=1$, $k=1$. We just remarked
that $V^1(\mathfrak{g},B)=H$ is
Heisenberg. And
$F_1=F=\left\langle{}\varphi,\varphi^*\right\rangle{}$
which we have drawn previously.
We also computed
the character as an infinite product.

\medskip\noindent
We have $F^{(0)} \ni J = \,:\!\varphi
\varphi^*\!:\,$
with $\lambda$-bracket relation $[J_\lambda
J]=\lambda |0\rangle$.
Recall $J(z)=\sum_{n \in
\mathbb{Z}}J_{(n)}z^{-n-1}=
\sum_{n \in \mathbb{Z}}J_n z^{-n-1}$.
The $\lambda$-bracket relation implies
$$
\begin{aligned}
[J(z),J(w)]&=\partial_w \delta(z,w)I\\
\implies [J_m,J_n]&=m \delta_{m,-n}I.
\end{aligned}
$$
This implies that $F^{(0)}$ is a
representation
of the oscillator Lie algebra
$\hat{\mathfrak{a}}=\bigoplus_{m \in
\mathbb{Z}}\mathbb{C} h_n \oplus
\mathbb{C}K$,
in which $h_m \mapsto  J_m$, $k \mapsto  I$
and $h_m |0\rangle$
for all $m \geq 0$.

\begin{proposition}
\label{vertex-algebras-proposition-f0-isomo-polyn-repre}
As representations of $\mathfrak{a}$,
$F^{(0)}\simeq H=\mathbb{C}[x_1,x_2,\ldots]$.
\end{proposition}

\begin{proof}
Using $H \simeq
U(\hat{\mathfrak{a}})
\otimes_{U(\hat{\mathfrak{a}_+})}\mathbb{C}
1$,
by its universal property, there exists a
morphism of
 $\hat{\mathfrak{a}}$-representations
$f: H \to F^{(0)}$, such that
$f(1)=|0\rangle$.

By Exercise \ref{vertex-algebras-exercise-H-is-irreducible}
we know $H$ is irreducible, so
$f:H \to F^{(0)}$ is injective.
Is it an isomorphism?
Consider  $\tilde{F}=F^{(0)}/f(H)$.
We want to prove $\tilde{F}=\{0\}$.
The trick is to consider
$L=\frac{1}{2}\,:\!J J\!:\,$.

We have seen that
$$
F^{(m)}=\{v \in F: J_0v=mv\}
$$
(since $[J_\lambda \varphi]=-\varphi$ and
$[J_\lambda \varphi^*]=+\varphi^*$,
$[J_0,\varphi_n]=-\varphi_n$).

Write $L(z)=\sum_{n \in
\mathbb{Z}}L_nz^{-n-2}$,
then
$L_0=\frac{1}{2}J_0^2+\sum_{m>0}J_{-m}J_m$.
This implies (for example) that
$L_0|0\rangle=0$.

One may compute
$$
\begin{aligned}
[L_\lambda
\varphi]&=T\varphi+\frac{1}{2}
\lambda\varphi\\
[L_\lambda\varphi^* ]&=T\varphi^*
+\frac{1}{2}\lambda \varphi^*.
\end{aligned}
$$

\noindent
If you expand these in terms of coefficients,
you find
$[L_0,\varphi_n]=-n\varphi_n$,
$[L_0,\varphi_n^*]=-n\varphi_n^*$.

From these relations we conclude that
$\Delta(\text{monomial})=L_0$-eigenvalue
(monomial), i.e.,
the (charge, energy)-grading is the
eigenspace grading by $(J_0,L_0)$.
Suppose $\tilde{F}\neq 0$.
Then $\exists \overline{v} \in \tilde{F}$ for
which
$J_m \overline{v}=0$ for all $m>0$.
Indeed, by Lemma (?) the
$\mathbb{Z}_+$-grading
of $F^{(0)}$ is inherited by $\tilde{F}$.
Since $|0\rangle \in f(H)$,
$\tilde{F}=\bigoplus_{\Delta \geq
N}\tilde{F}_{\Delta}$
for some $N>0$. Let $\overline{v} \in
\tilde{F}_N$,
$\overline{v}\neq 0$.
For $m>0$, $J_m \overline{v} \in
\tilde{F}_{N-m}=0$.

Key point: $J_0 \overline{v}$, because
$\tilde{F}$ is quotient of $F^{(0)}$.
Hence $L_0\overline{v}=0$.
But $L_0=\Delta$, and we know the only
element of $F^{(0)}$
with $\Delta=0$ is $|0\rangle$.
This contradiction implies $\tilde{F}=0$,
hence $F^{(0)}=f(H) \simeq H$.
\end{proof}

In the process we saw that
$\Delta$ coincides with $L_0$,
where $L=\frac{1}{2}\,:\!J J\!:\,$,
and $J=f(h)$. So consider
$L=\frac{1}{2}\,:\!h h_i\!:\,$.

We have $[L_\lambda h]=Th+\lambda h$
so $[L_0,h_n]=-nh_n$ in particular,
and
$$
\begin{aligned}
L_0(h_{-n_1},h_{-n_2},\ldots,h_{-n_s}
|0\rangle&=
L_0(x_{n_1}x_{n_2}\ldots x_{n_s}|0\rangle)
&=\left(\sum_{j}n_j x_{n_1}\ldots
x_{n_s}|0\rangle\right)
\end{aligned}
$$

\noindent
Using this we obtain
 $$
\chi_{F^{(0)}}(q)=\sum_{\Delta}\dim(F^{(0)}
_{\Delta})q^{\Delta}
=\prod_{m=1}^{\infty}\frac{1}{1-q^m}.
$$

In $F^{(m)}$, let us denote

$$
|m\rangle=\begin{cases}
\varphi_{-1/2}\varphi_{-3/2}\varphi_{-5/2}
\ldots
\varphi_{-\frac{2(-m)-1}{2}}|\rangle\qquad
&\text{if }m<0 \\
\varphi^*_{-1/2}\varphi^*_{3/2}\ldots
\varphi^*_{-\frac{2m-1}{2}}\qquad &\text{if
}m>0
\end{cases}
$$

We observe that
$\Delta(|m\rangle)=\frac{m^2}{2}$
and $F^{(m)}_{m^2/2}=\mathbb{C} |m\rangle$.

\begin{proposition}
\label{vertex-algebras-proposition-isomo-repre-all-m}
As an $\hat{\mathfrak{a}}$-module,
$F^{(m)}\simeq H^m$, and is irreducible.
\end{proposition}

\begin{proof}
Same as above.
\end{proof}

This gives a formula
$$
\chi_{F^{(m)}}(q)=\sum_{\Delta}\dim
F^{(m)}_\Delta q^\Delta
=q^{m/2}\prod_{k=1}^\infty \frac{1}{1-q^k}.
$$
As a corollary, we obtain again
Jacobi's triple product formula. (Again!)

$$
\prod_{m=1}^\infty(1+yq^{m-1/2})(1+y^{-1}
q^{m-1/2})
=\prod_{k=1}^\infty\frac{1}{1-q^k}
\left(\sum_{n \in
\mathbb{Z}}y^nq^{n^2/2}\right).
$$

\section{Another presentation the charged
free Fermions}
\label{vertex-algebras-section-anoth-prese-fermi}

The idea is to consider $X = \bigoplus_{i \in
\mathbb{Z}} \mathbb{C} e_i$,
a countably infinite dimensional space,
and define
\begin{equation}
\label{vertex-algebras-equation-semi-inifinite}
\Lambda^{\infty/2}=
\left\{\substack{\text{span of symbols of the
form} \\
e_{i_0}\wedge e_{i_1}\wedge e_{i_1} \wedge
e_{i_3}\wedge\ldots\\
\text{where }\exists N\text{ s.t. }\forall n
\geq  N,
i_{n+1}=i_n+1}\right\}\Big/
\substack{\text{relation that} \\ e_i\wedge
e_j=- e_j \wedge e_i\\
\text{``wherever it occurs''}}.
\end{equation}
That is, the indices can dance around as they
like
but at some point they will become
consecutive.
Clearly a basis of $\Lambda^{\infty/2}$ is
given
by those ``semi-infinite monomials''
for which
$i_0 < i_1< i_2<\ldots$.

For any semi infinite monomial
$\underline{e}=e_{i_0}\wedge e_{i_1}\wedge
\ldots$
there exists a number $m$ such that
$i_j=-m+j$ for all  $j \gg 0$.
This number is called the {\it charge} of
$\underline{e}$.

We have a decomposition in vector spaces
$$
\Lambda^{\infty/2}=\bigoplus_{m \in
\mathbb{Z}}\Lambda^{\infty/2,(m)}
$$
called {\it charge decomposition}.

Let's introduce some operators in
$\Lambda^{\infty/2}$:
$$
\begin{aligned}
\varphi_{(n)}: \Lambda^{\infty/2}
&\longrightarrow \Lambda^{\infty/2} \\
\varphi_{(n)}(\underline{e}) &= e_n \wedge
\underline{e}
\end{aligned}
$$
i.e. we just put $e_n$ at the beginning of
the
monomial.

Next
$$
\begin{aligned}
\varphi_{(n)}^*: \Lambda^{\infty/2}
&\longrightarrow \Lambda^{\infty/2} \\
\varphi^*_{(n)} & =\sum_{k \geq 0}(-1)^k
\delta_{n,i_k}
e_{i_0}\wedge e_{i_1}\wedge \ldots \wedge
\underbrace{\widehat{e_{i_k}}}_{\text{remove
this}}\wedge \ldots
\end{aligned}
$$
so in analogy, we remove this term.
The $(-1)^k$ accounts for moving around the
unwanted
term to the beginning of the monomial.

\begin{exercise}
\label{vertex-algebras-exercise-something-like}
Something like
$$
\varphi_{(m)}\varphi^*_{(n)}
+\varphi_{(n)}^*\varphi_{(m)}=
\delta_{m,+n}I_{\Lambda^{\infty/2}}
$$
\end{exercise}

Let's introduce some quantum fields now:

$$
\varphi(w)=\sum_{ n \in
\mathbb{Z}}\varphi_{(n)}w^{-n-1}
$$
which vanishes for very large $n$ because
we eventually get to the ``consecutive''
region,
and
$$
\varphi^*(w)=\sum_{ n \in
\mathbb{Z}}\varphi^*_{(n)}w^{n}
$$
which vanishes for very negative $n$ since
the
star operators are defined to give zero
if the $e_n$ is not found in the monomial.

Notice that the convention
$\varphi^*(w)=\sum \varphi^*_{(n)}w^n$,
instead of $\sum \varphi^*_{(n)}w^{-n}$ is so
that
$\varphi^*(w)$ is a quantum field
(indeed, by our convention on how we write
quantum fields,
see Definition
\ref{vertex-algebras-definition-quantum-field}).
The convention $\varphi^*(w)=\sum
\varphi^*_{(n)}w^n$
instead of, say, $\sum \varphi_{(n)}^*
w^{n+1}$
is so that the equation in
Exercise \ref{vertex-algebras-exercise-something-like}
turns into a nice relation
$$
[\varphi(z),\varphi^*(z)]=\delta(z,w)I,
$$
(i.e. $[\varphi_\lambda \varphi^*
]=|0\rangle$, just like in the
charged fermions $F^{\text{ch}}$!).

\begin{exercise}
\label{vertex-algebras-exercise-relations-for-these-fields}
These fields will match with our previous
$\varphi,\varphi^*$ as
$$
\begin{aligned}
\sum \varphi_{(n)}z^{-n-1}&=\sum \varphi_n
z^{-n-1/2}
: \varphi_n=\varphi_{n-1/2}\\
\sum \varphi_{(n)}^*z^n&=\sum \varphi_n^*
z^{-n-1/2}
: \varphi_n^*=\varphi_i
\end{aligned}
$$
\end{exercise}

\medskip\noindent
The space $\Lambda^{\infty/2}$ has a
super-structure
with parity $\Lambda^{\infty/2,(m)}$ = the
parity of $m$.
Setting
$|0\rangle=e_0 \wedge e_1 \wedge e_2 \ldots$,
$\mathcal{F}=\{\varphi(z),\varphi^*(z)\}$,
$T$ what it needs to be,
we get a vertex superalgebra.
In fact,
$$
F^{\text{ch}}=U(\hat{\mathfrak{a}})
\otimes_{U(\hat{\mathfrak{a}_+})}
\mathbb{C}|0\rangle
\xrightarrow{\simeq}\Lambda^{\infty/2}
$$
is an isomorphism of vertex algebras.
Indeed, if $n_j,m_k \geq 1$,
$$
\begin{aligned}
&\varphi_{(n-1)}\varphi_{(-n_2)}
\ldots\varphi_{(-n_s)}
\varphi^*_{(m_1)}\ldots
\varphi^*_{(m_t)}|0\rangle
&=\pm e_{-n_1}\wedge e\ldots \wedge e_{-n_s}
\wedge \underbrace{(e_0 \wedge e_1 \wedge
\ldots)
}_{\substack{\text{but with} \\ e_{m_1}\ldots
e_{m_t}\\
\text{missing}}}
\end{aligned}
$$

\noindent
But these two pictures are bases of
$F^{\text{ch}}$ and $\Lambda^{\infty/2}$,
so we have our linear isomorphism
$F^{\text{ch}} \to \Lambda^{\infty/2}$.
Just need to check carefully $\varphi$ and
$\varphi^*$
are compatible with it.

\medskip\noindent
Pauli exclusion principle. In an ensemble
state of fermions,
two cannot occupy the same state,
(so they are modelled by odd variables
like ``$e_n$'' such that $e_n \wedge e_n=0$.)

Electron waves should obey
``quantized'' wave-equation,
which is something like
$$
\begin{aligned}
\partial_t^2 \psi&=(- \overline{h}^2
\nabla+m)\psi\\
\partial_t
\psi&=\sqrt{-\overline{h}\nabla+m}\psi.
\end{aligned}
$$

\noindent
A solution to this is, instead of
$$
\sqrt{\partial_{x x}+m},
$$
try to define
$$
\sqrt{(\partial_{x x}+m)I_4}
$$
where $I_4$ is the identity $4 \times 4 $
matrix.
Then the ``square root'' of $-\overline{h}^2
\nabla +m$
becomes a matrix valued differential
operator,
denoted sometimes as $\not \partial +
\gamma$,
where $\gamma$ is an explicit $4 \times 4$
matrix.
In $\partial_t \psi=\ldots$, $\psi$ has
become
a vector-valued function on space,
with a splitting in spin up and spin down
parts
of the electron wavefunction, and also an
``unphysical'' part. The latter
can be interpreted as the positrons.

Problem: the Dirac equation has positive and
negative
energy solutions.

\medskip\noindent
Let's examine the neutral part
of $\Lambda^{\infty/2(0)}$.
Consider the element
$$
\underline{e}=e_{-4}\wedge e_{-2}\wedge
e_{-1}
\wedge e_0 \wedge e_1 \wedge e_4 \wedge e_5
\wedge\underbrace{ e_7 \wedge e_8 \wedge
e_9\wedge\ldots
}_{\text{consecutive part}}
$$
The successive differences are
$$
2\quad 1 \quad  1\quad 1 \quad 3 \quad 1
\quad 2 \quad 1
\quad  1 \quad  \ldots
$$
Since the sequence will stabilize at 1,
we might as well subtract 1 and obtain the
finite
sequence
$$
(1\quad 0\quad 0\quad 0\quad 2\quad 0\quad 1)
$$
In fact, you could reconstruct
$\underline{e}$
from this list, and knowing that
$\text{charge}(\underline{e})=0$.

Form the partial sums
$$
1000201 \to 4333311
$$
these numbers are non-increasing,
so define uniquely a partition of some
integer
(in this case 18),
and if we compare with
$$
\underline{e}=\pm
\varphi_{-4}\varphi_{-2}\varphi_{-1}
\varphi^*_2\varphi_3^*\varphi_6^*|0\rangle
$$
which has energy
$$
\Delta=(4-1/2)+(2-1/2)+(1-1/2)
+(2+1/2)+(3+1/2)+(6+1/2)=18.
$$

\noindent
We denote by $\underline{e}_\lambda$
the semi-infinite monomial of charge 0 given
by
this procedure.

\medskip\noindent
So we have a basis
$\{\underline{e}_\lambda|\lambda\text{
integer partitions}\}$
of $F^{(0)}\simeq H =
\mathbb{C}[x_1,x_2,x_3,\ldots]$,
which also has the basis
$\{\underline{x}_\lambda=x_1^{\lambda_1}
\ldots
x_5^{\lambda_5}|
\lambda \text{ integer partitions}\}$
where
$\lambda=(1^{\lambda_1},2^{\lambda_2},\ldots,
s^{\lambda_s})$,
that is, $1$ appears $\lambda_1$ times and so
on.
We do {\bf not} have
$\underline{e}_\lambda=\underline{x}
_\lambda$!

\section{Schur polynomials}
\label{vertex-algebras-section-schur-polynomials}

The Schur polynomials are given by
the exponential generating function of $\sum
z^kx_k$.
More precisely,

\begin{definition}
\label{vertex-algebras-definition-schur-polynomial}
We set  $S_k(\underline{x})$ by
$$
\sum_{k \geq 0}z^k S_k(\underline{x})
=\text{exp}\left(\sum_{k \geq 1}z^k
x_k\right).
$$
\end{definition}

For a partition
$\lambda=(1^{\lambda_1}2^{\lambda_2}\ldots
k^{\lambda_k})$,
define
$$
S_\lambda(\underline{x})
=\det
\begin{pmatrix}
S_{\lambda_1}& S_{\lambda_1+1}&\cdots &
S_{\lambda_1+k-1}\\
S_{\lambda_2-1}& S_{\lambda_2}&\cdots & \\
\vdots &  &  &  \vdots \\
S_{\lambda_k-k+1}&  &  &  S_{\lambda_k}
\end{pmatrix}
$$

Recall we defined $X=\bigoplus_{i \in
\mathbb{Z}}\mathbb{C}e_i$
and
$\Lambda^{\infty/2}=\left\langle{}e_{i_0}
\wedge\ldots\right\rangle{}$.
Let $g:X \to X$ be an invertible
endomorphism.
We would like to
define an endomorphism
$$
\begin{aligned}
R(g): \Lambda^{\infty/2} &\longrightarrow
\Lambda^{\infty/2} \\
R(g)(e_{i_0}\wedge e_{i_1}\wedge\ldots) &=
(ge_{i_0})\wedge(ge_{i_1})\wedge \ldots
\end{aligned}
$$
For some $g$, $\varepsilon g$,
$g(e_n)=e_{-n}$,
$R(g)$ does not make sense
(or does not send  $\Lambda^{\infty/2}\to
\Lambda^{\infty/2}$).

As a note, if $X$ were finite-dimensional,
say $X=\left\langle{}e_1,\ldots,e_n\right\rangle{}$,
then $R(g):\Lambda^{k+1}X \to \Lambda^{k+1}X$
has matrix entry
$$
(e_{i_0}\wedge e_{i_1}\wedge\ldots\wedge
e_{i_k})
\to(e_{j_0}\wedge e_{j_1} \wedge \ldots
\wedge e_{j_k}),
$$
(with strictly increasing indices on both
sides,
up to a sign maybe)
given by the determinant of the
$(k+1)\times(k+1)$
matrix given by selecting the rows
$i_0,i_1,\ldots,i_k$, and the columns
$j_0,j_1,\ldots,j_k$ of the matrix of $g$.

\medskip\noindent
For our $X$ and $\Lambda^{\infty/2}$,
$R(g)$ makes sense if $g$ is such that
$$
g(e_i)-e_i \in \text{span} \{ e_{>i}\}\qquad
\forall i.
$$

\begin{proposition}
\label{vertex-algebras-proposition-schur-polynomials}
$\underline{e}_\lambda=
S_\lambda(\underline{x})$
where
\end{proposition}

\begin{proof}
Write $\underline{e}_\lambda=P(x)$.
We wish to show $P(x)=S_\lambda(x)$.
Introduce new variables $y_1,y_2,\ldots$ and
consider
$$
F(y)=\text{exp}
\left(\sum_{k>0}y_k\frac{\partial }{\partial
x_k}\right)P(x)\Big|_{x=0}
$$
First we observe that
$$
F(y)=P(x_1+y_1,x_2+y_2,\ldots)|_{x=0}=P(x+y)
|_{x=0}=P(y)
$$
since in general exponentiaing a differential
operator
gives shifting by the coefficient, i.e.
$e^{a \frac{\partial }{\partial
t}}f(t)=f(t+a)$.

Considering
$$
P(x) \in \mathbb{C}[x_1,x_2,\ldots]=H\simeq
F^{(0)},
$$
observe that for $k>0$,
$$
\begin{aligned}
\frac{\partial }{\partial x_k}&=h_k
=(\,:\!\varphi \varphi^*\!:\,)_k
=\sum_{i \in
\mathbb{Z}}\varphi_{k+1}\varphi^*_i
\end{aligned}
$$

\noindent
i.e. it's a sum of ``remove $e_i$'' and
``insert $e_{k+i}$''.

Denote
$$
\begin{aligned}
\Lambda_k: X &\longrightarrow X \\
\Lambda_k(e_n) &=e_{n+k}\qquad  \forall n \in
\mathbb{Z}
\end{aligned}
$$

\noindent
$R(\Lambda_k)=\frac{\partial }{\partial
x_k}$,
since
$\Lambda^{\infty/2(0)}\xrightarrow{\simeq}H$.

Therefore
$$
F(y)=R \text{exp}\left(\sum_{k>0}y_k
\Lambda_k\right)
\underline{e}_\lambda\Big|_{\text{coef. in
}|0\rangle}.
$$
Notice that $R \text{exp}(\Lambda_k)$ is of
the form
$g(e_i)-e_i$.

Finally, notice that $\Lambda_k=\Lambda_1^k$,
so
$$
\begin{aligned}
&R \text{exp}\left(\sum_{k>0}y_k
\Lambda_k\right)
\underline{e}_\lambda\Big|_{\text{coef. in
}|0\rangle}\\
&=R \text{exp}\left(\sum y_k
\Lambda_1^k\right)\underline{e}_\lambda\\
&=R(\left(S_k
\Lambda_k\right)\underline{e}_\lambda.
\end{aligned}
$$
\end{proof}


\section{The Tate extension and the Japanese
cocycle}
\label{vertex-algebras-section-tate-exten-japan-cocyc}

\noindent
Recall $X=\bigoplus_{n \in
\mathbb{Z}}\mathbb{C}e_n$
and $\Lambda^{\infty/2}$.
(You can be floppy and think that
$\Lambda^{\infty/2}$
is the exterior power of $X$,
though that's not completely right.)
Define
$$
\mathfrak{gl}_\infty
=\left\{(a_{ij}):\substack{i,j \in
\mathbb{Z},a_{ij}\in \mathbb{C} \\
\text{all but finitely many of the $a_{ij}$
vanish}}\right\},
$$
add and multiply as usual.

$\text{GL}_\infty=\{I+(a_{ij}):a_{ij}
\in\mathfrak{gl}_\infty\}$,
with $I=(I_{ij})$ and $I_{ij}=\delta_{ij}$.
A typical element of $\mathfrak{gl}_\infty$
is $E_{ij}$ with
$E_{ij}(e_k)=\delta_{jk}e_i$.

Last time we used shift operators
$\Lambda_k:e_n\mapsto e_{n+k}$ for all $n$.
These are {\bf not} in
$\mathfrak{gl}_\infty$,
but {\bf are} in
$$
\widetilde{\mathfrak{gl}_\infty}
=\left\{(a_{ij}):\substack{\exists  N s.t.
a_{ij}=0 \\ \text{ whenever }
|i-j|>N}\right\}.
$$
Then $\Lambda_k=\sum E_{i+k,k}\in
\widetilde{\mathfrak{gl}_\infty}$.

\begin{remark}
\label{vertex-algebras-remark-gltilde-is-Lie-algebra}
$\widetilde{\mathfrak{gl}_\infty}$ is an
associative
Lie algebra (hence a Lie algebra with
commutator).
This is because the condition of finiteness
in the definition
excludes the possibility of infinite sums.
\end{remark}


We have a represenation $r$ of
$\mathfrak{gl}_\infty$
on $\Lambda^{\infty/2}$ by
$$
r:E_{ij}\mapsto \varphi_i\varphi_j^*
$$
But this doesn't extend to
$\widetilde{\mathfrak{gl}_\infty}$.
For instance, $r(\Lambda_0)=\sum_{j \in
\mathbb{Z}}\varphi_{-j}\varphi^*_j$
diverges.

A ``solution'' is to use normal order,
so $r(\Lambda_0)=\sum_{j \in
\mathbb{Z}}\,:\!\varphi_{-j}\varphi^*_j\!:\,
$,
is now well-defined, but $r$ is no longer a
representation.

\medskip\noindent
Today: more conceptual point of view.
Consider the vector space
$\mathbb{C}(\!(t)\!)=X$.
(Which is {\it uncountably}
infinite-dimensional.)
Give $X$ a linear topology
with a base of open neighbourhoods
of $0$ being $t^N\mathbb{C}[\![t]\!]$
for $N \in \mathbb{Z}$.
The idea is that
$t,t^2,t^3,\ldots$ ``tends to $0$''
in this topology.

\begin{definition}
\label{vertex-algebras-definition-tate}
A {\it Tate vector space} is a linearly
topologised vector space $X$ and a set
$\mathcal{L}$
of linear subspaces $L \subset X$
(called {\it lattices}) such that
\begin{enumerate}
\item (Separated.) Any neighbourhood of $0$
contains
a lattice.

\item (Exhaustive.) Every $x \in X$
is contained in some lattice.

\item (Commensurable.) For all $L_1,L_2 \in
\mathcal{L}$,
$L_1 \cap L_2$ has finite codimension
in $L_1$ (and in $L_2$). (This says
that any two lattices cannot be
``infinitely far apart''.)

\item (Complete.) For all $L_1,L_2 \in
\mathcal{L}$,
all vector subspaces
$L_1 \cap L_2 \subset S \subset L_1+L_2$
then $S \in \mathcal{L}$.

\item (Another completeness.) By the
universal property of limit,
we know there exists a map
$X \to \lim_{\substack{L \in \mathcal{L} \\ L
\to 0}} X/L$.
We ask this is an isomorphism.
\end{enumerate}
\end{definition}

\begin{example}
\label{vertex-algebras-example-tate-vector-space}
\begin{enumerate}
\item
Laurent series is an example of a Tate vector
space.
Let $X=\mathbb{C}(\!(t)\!)$ with
 $$
\mathcal{L}=\left\{L\subset X:\exists N\text{
s.t. }
t^N \mathbb{C}[\![t]\!] \subset L \subset
t^{-N}\mathbb{C}[\![t]\!]\right\}.
$$
We can also say that $\mathcal{L}$ is the
unique structure such that $\mathcal{L} \ni
\mathbb{C}[\![t]\!]$.

\item If $\mathcal{L}=\{0\}$,
then $\mathcal{L}=\{L \subset X |
\dim(L)<\infty\}$.

\item If $\mathcal{L}\ni X$, then
$\mathcal{L}=\{L \subset
X|\dim(X/L)<\infty\}$.
\end{enumerate}
\end{example}

\noindent
Since we introduced a topology on $X$,
we can consider the continuous endomorphisms
$\text{End}_{\text{cont}}(X)$.
Let's denote throughout this section
 $$
\text{End}(X)=\{f \in
\text{End}_{\text{cont}}(X)
:\exists U,V \in \mathcal{L}, f(U)\subset
V\}.
$$
Relative to the ``basis''
$\{t^n|n \in \mathbb{Z}\}$,
the matrix of $f \in \text{End}(X)$
looks like
$$
\begin{pmatrix}
*&0\\
*&*
\end{pmatrix}
\begin{pmatrix}
0\\
*
\end{pmatrix}
$$
since the vectors in $V\subset
t^m\mathbb{C}[\![t]\!]$
have zeroes before a certain index $n$.

\medskip\noindent
\begin{align*}
\text{End}_c(X)&=\{f \in \text{End}(X)|
\exists  V \in \mathcal{L}
\text{ s.t. }f(X) \subset V\}
=\left\{\begin{pmatrix}
0&0\\
*&*
\end{pmatrix}\right\}\\
\text{End}_d(X)&=\{f \in
\text{End}(X)|\exists U \in \mathcal{L}
\text{ s.t. }f(U)=0\}
=\left\{\begin{pmatrix}
*&0\\
*&0
\end{pmatrix}\right\}\\
\text{End}_f(X)&=\text{End}_c(X) \cap
\text{End}_d(X)
=\left\{\begin{pmatrix}
0&0\\
*&0
\end{pmatrix}\right\}.
\end{align*}

\begin{remark}
\label{vertex-algebras-remark-trace-in-End}
For $f \in \text{End}_f(X)$,
the trace $\text{Tr}(f) \in \mathbb{C}$
is well-defined because
the intersection of the diagonal
with the lower-left part of any
matrix in $\text{End}_f(X)$ is finite.
For $\text{End}_c$ and $\text{End}_d$
trace is not well-defined.
\end{remark}

\noindent
We introduce the map

$$
\begin{aligned}
p:\text{End}_c(X) \oplus \text{End}_d(X)
&\longrightarrow \text{End}(X)\to  0\\
(f,g) &\longmapsto f+g
\end{aligned}
$$

\noindent
which is clearly surjective since we can put
$$
\begin{pmatrix}
A&0\\
B&C
\end{pmatrix}=\begin{pmatrix}
A&0\\
B&0
\end{pmatrix}+
\begin{pmatrix}
0&0\\
0&D
\end{pmatrix}.
$$
We can complete this to a short exact
sequence
(of vector spaces)
by putting
$$
\begin{aligned}
i:\text{End}_f(X)  &\longrightarrow
\text{End}_c \oplus \text{End}_d \\
h &\longmapsto (h,-h).
\end{aligned}
$$

\noindent
Considering the trace
$\text{Tr}:\text{End}_f(X) \to \mathbb{C}$,
we can form the pushout $L$ in the
category of
{\bf vector spaces}
$$
\xymatrix{
0\ar[r]&\text{End}_f\ar[r]\ar[d]_{\text{Tr}}
\ar@{}[dr]|-{\ulcorner}&
\text{End}_c(X) \oplus
\text{End}_d(X)\ar[d]\ar[r]&
\text{End}(X)\ar[r]&0\\
&\mathbb{C}\ar[r]& L
}
$$
Concretely,
$$
L=\frac{(\text{End}_c(X) \oplus
\text{End}_d(X))\oplus \mathbb{C}}
{\left\langle{}(h,-h)=(0,0,\text{Tr}(h))|
\forall  h \in \text{End}_f(X)\right\rangle{}}.
$$
Notice that the map $i$
is {\bf not} a Lie algebra map:
for this it would have to be a Lie algebra
map
in both entries, which is true in the first
entry
but not on the second.
Meanwhile, $p$ {\it is} a morphism of Lie
algebras.
We could correct $i$ to map $h\mapsto  (h,h)$
and then change $p$ to be $\beta-\gamma$,
but then $p$ would stop being a Lie algebra
morphism.
So this is not a short exact sequence of Lie
algebras;
indeed it's not sensible to think of
short exact sequences of Lie algebras
since Lie algebras do not form an Abelian
category.

But it is a short exact sequence of vector
spaces
and we can form its exact sequence pushout by
the following
exercise. (Also see Stacks Project tag
\href{https://stacks.math.columbia.edu/tag/%
010I}{010I}.)

\begin{exercise}
\label{vertex-algebras-exercise-pushout}
Let
$$
\xymatrix{
0\ar[r]&A\ar[r]&B\ar[r]&C\ar[r]&0
}
$$
be a short exact sequence of vector spaces
and $t:A \to K$ a linear map of vector
spaces.
Let $X$ be the pushout
$$
\xymatrix{
A\ar[r]^i\ar[d]_t\ar@{}[dr]|-{\ulcorner}&
B\ar[d]\\
K\ar[r]&X.
}
$$
Explicitly,
$$
X=\frac{K \oplus B}{\left\langle{}(t(a),0)-(0,i(a))|a
\in A\right\rangle{}}.
$$
Show that there exists a short exact sequence
 $$
\xymatrix{
0\ar[r]&K\ar[r]&X\ar[r]&C\ar[r]&0
}
$$
such that
$$
\xymatrix{
0\ar[r]&A\ar[r]\ar[d]_{t}\ar@{}[dr]|-
{\ulcorner}&
B\ar[d]\ar[r]&
C\ar[r]\ar[d]^{\text{id}}&0\\
0\ar[r]&K\ar[r]& X\ar[r]& C\ar[r]&0
}
$$
commutes.
\end{exercise}

\begin{proof}
To define a map $X \to C$ we first notice
that all we have to do is define the map on
$K \oplus B$ at points not coming from $A$,
i.e. whose entries are not of the form
$(i(a),t(a))$ for $a \in A$.
Indeed, those are cancelled by the pushout
definition
as a quotient. Further, we already have a map
$\varphi:B \to C$
defined on elements not of the form $i(a)$
for $a \in A$
by exactness of the given short exact
sequence.
Then we must define $(k,b) \mapsto
\varphi(b)$
for the diagram to commute.

The resulting exact sequence is exact once
we know that $K \to B$ in the pushout
is given by $k \mapsto (k,0) \text{mod
}\sim$.
\end{proof}

Then we obtain
$$
\xymatrix{
0\ar[r]&\text{End}_f\ar[r]\ar[d]_{\text{Tr}}
\ar@{}[dr]|-{\ulcorner}&
\text{End}_c(X) \oplus
\text{End}_d(X)\ar[d]\ar[r]&
\text{End}(X)\ar[r]\ar[d]^{=}&0\\
0 \ar[r]&\mathbb{C}\ar[r]& L\ar[r]&
\text{End}(X)\ar[r]&0
}
$$
\medskip\noindent
Notice that the exact sequence
$$
\xymatrix{
0\ar[r]&\text{End}_f\ar[r]^i
&\text{End}_d \oplus \text{End}_f\ar[r]^p
&\text{End}\ar[r]&0
}
$$
is a bit more than
just a short exact sequence of vector spaces.
Notice that $\text{End}_d,\text{End}_c
\subset \text{End}$
are ideals.
Indeed,
if $f \in \text{End}_c$ and $g \in
\text{End}$,
then $f \circ g \in \text{End}_c$ obviously,
$f(X) \subset V$ so $g \circ f (X) \subset
g(V)$,
continuity of $g$
(plus axioms of Tate vector space)
implies that there exists $\widetilde{V} \in
\mathcal{L}$
such that $g(V) \subset \widetilde{V}$.

\begin{exercise}
\label{vertex-algebras-exercise-spell-this-out}
Spell this out.
\end{exercise}

So $g \circ f \in \text{End}_c$.

\begin{exercise}
\label{vertex-algebras-exercise-ideal-too}
$\text{End}_d \subset \text{End}$ is an ideal
too.
\end{exercise}

So $\text{End}_c, \text{End}_d$ are
$\text{End}$-modules
(as associative algebras and as Lie
algebras),
and
$$
\xymatrix{
0\ar[r]&\text{End}_f\ar[r]&\text{End}_d
\oplus \text{End}_c\ar[r]&
\text{End}\ar[r]&0
}
$$
is a short exact sequence of End-modules.

\medskip\noindent
Let $\mathfrak{g}$ be a Lie algebra, $E$ a
$\mathfrak{g}$-module
with $E \xrightarrow{p}\mathfrak{g}$
surjective
$\mathfrak{g}$-module morphism.
Consider
$$
\xymatrix{
0\ar[r]&\mathfrak{a}\ar[r]^i&E\ar[r]^p&
\mathfrak{g}\ar[r]&0
}
$$
i.e. $\mathfrak{a}=\Ker(p)$.

Then one obtains a symmetric bilinear form
defined by
$$
(a,b)=p(a)b+p(b)a.
$$
\begin{exercise}
\label{vertex-algebras-exercise-bilinear-form-extension}
This form is (1) symmetric and
(2) $\mathfrak{g}$-invariant.
\end{exercise}

\begin{exercise}
\label{vertex-algebras-exercise-central-extension}
Check that the data above,
satisfying $(\cdot,\cdot)=0$
is the same as a central extension of
$\mathfrak{g}$
by $\mathfrak{a}$,
i.e. Lie bracket $[\cdot,\cdot]:E \times E
\to E$
compatible with $[\cdot,\cdot]$ on
$\mathfrak{g}$
such that $\mathfrak{a}\subset E$
is central.
{\bf Hint.} Set $[a,b]^E=p(a)b$.
\end{exercise}

\noindent
As a particular case of this,
if we have
\begin{itemize}
\item $\mathfrak{g}$ a Lie algebra,
\item $\mathfrak{a}_1, \mathfrak{a}_2 \subset
\mathfrak{g}$
two ideals such that
$\mathfrak{g}=\mathfrak{a}_1 +
\mathfrak{a}_2$,
\item setting $\mathfrak{a}_0=\mathfrak{a}_1
\cap \mathfrak{a}_2$,
a linear map $\mathfrak{a}_0
\xrightarrow{T}K$
such that $T([x_1,x_2])=0$ for all
$x_1 \in \mathfrak{a}_1$ and $x_2 \in
\mathfrak{a}_2$.
\end{itemize}

\noindent
Then we can form
$$
\xymatrix{
0\ar[r]&\mathfrak{a}_0\ar[d]_T\ar[r]^i
&\mathfrak{a}_1 \oplus
\mathfrak{a}_2\ar[d]\ar[r]^p
&\mathfrak{g}\ar[r]\ar[d]^{\text{id}}&0\\
0\ar[r]& K\ar[r]&
X\ar[r]&\mathfrak{g}\ar[r]&0
}
$$
\noindent
with pushout at level of vector spaces
and $i:\alpha \mapsto (\alpha,-\alpha)$
and $p:(\beta,\gamma)\mapsto \beta+\gamma$.

\begin{exercise}
\label{vertex-algebras-exercise-central-extension-duplicate-2}
Then $X$ will become a Lie algebra
with $K$ central and
$\mathfrak{a}_1 \oplus \mathfrak{a}_2 \to X$
a map of Lie algebras.
\end{exercise}

\noindent
The main point of this construction
is to use $T([x_1,x_2])$
to confirm that
$$
T(p(a)b+p(b)a)=0\qquad \forall a,b \in
\mathfrak{a}_1 \oplus \mathfrak{a}_2.
$$
Indeed, write $a=(a_1,a_2)$ and
$b=(b_1,b_2)$.
Then
$$
\begin{aligned}
&T(([a_1,a_2,b_1],[a_1+a_2,b_2])+([b_1+b_2,
b_1],[b_1+b_2,a_2])\\
&=T([a_1,b_1]+[a_2,b_1]+[b_1,a_1]+[b_2,a_1],
[a_1,b_2]+[a_2,b_2]+[b_1,a_2]+[b_2,a_1])\\
&=T(i(\underbrace{[a_2,b_1]}
_{\in[\mathfrak{a}_1,\mathfrak{a}_2]}+
\underbrace{[b_2,a_1]}_{\in[\mathfrak{a}_2,
\mathfrak{a}_1]}))\\
&=0\qquad \text{by hypothesis.}
\end{aligned}
$$

\medskip\noindent
We apply this, in particular, to
$$
\mathfrak{g}=\text{End}(X)
\qquad \mathfrak{a}_1=\text{End}_d(X)
\qquad \mathfrak{a}_2=\text{End}_c(X)
\qquad \mathfrak{a}_0=\text{End}_f(X)
$$
$$
X \text{ our Tate vector space}\qquad
T \text{ trace}.
$$

\noindent
For the construction above to work
we still have to check that
$$
\text{Tr}([A_d,A_c])=0
$$
whenever $\text{Ad} \in \text{End}_d$ and
$A_c \in \text{End}_c$.
You would think this is obvious,
but it isn't.
(I.e., $\text{Tr}(AB-BA)=0$ for $A,B \in
\text{End}(\text{finite-dimensional
vector space})$.)

Let us sketch the proof.
Let $f \in \text{End}_c$ and $f(X) \subset V$
($V \in \mathcal{L}$),
and $g \in \text{End}_d$ and $g(U)=0$
($U \in \mathcal{L}$).
Then
$$
X \xrightarrow{f}\underbrace{f(X)}_{\subset
V}\xrightarrow{g}
\underbrace{gf(X)}_{\subset g(V)}.
$$
And
$$
X \xrightarrow{g}g(X) \to
\underbrace{fg(X)}_{\subset V}.
$$
So $\text{Im}[f,g] \subset V+g(V)$
So $[f,g](U \cap f^{-1}(U))=0$.
$$
\underbrace{f^{-1}g^{-1}(0)}_{\subset
f^{-1}(U)}
\xrightarrow{f}\underbrace{g^{-1}(0)}
_{\subset U}\xrightarrow{g}0
$$
and
$$
\underbrace{g^{-1}f^{-1}(U)}_{\subset U}
\xrightarrow{g}f^{-1}(0)\xrightarrow{f}0.
$$

\noindent
Now the idea is to argue that
$$
\text{Tr}[f,g]=\text{Tr}_Q[f,g],
$$
where
$$
Q=(V+g(V))/(U\cap f^{-1}(U))
$$
is a finite-dimensional vector space.
But on finite dimensional vector spaces
the trace of commutator vanishes.

\medskip\noindent
You might imagine that $\text{Tr}[f,g]=0$
whenever $[f,g]\in \text{End}_f(X)$,
but this is false.

\begin{example}
\label{vertex-algebras-example-trace-commu-not-vanis}
Let $A \subset \text{End} (X)$ be a
commutative subalgebra.
(For instance, if $X= \mathbb{C}((t))$,
then $f(t) \in \mathbb{C}((t))$
we have $\mu_f \in \text{End}(X)$
and $\mu_f(g)=fg$.)
Then I claim, if $f,g \in A$,
$f=f_c+f_d$,  $g=g_c +g_d$,
that $[f_c,g_c] \in \text{End}_f$.
Obviously  $[f_c,g_c] \in \text{End}_c$,
but also
$$
\begin{aligned}
[f_c,g_c]&=[f-f_d,g-g_d]\\
&=\underbrace{[f,g]}_{=0}
\underbrace{-[f,g_d]-[f_d,g]+[f_d,g_d]}_
{\substack{\in\text{End}_d \text{ because
}\\\text{End}_d \overset{\text{ideal}}
{\subset} \text{End}}}.
\end{aligned}
$$

\noindent
I claim that $\text{Tr}[f_c,g_c]$ might not
vanish.

In fact let $X=\mathbb{C}((t))$,
$f=\mu_{t^{-2}}$, $g=\mu_{t^2}$
(or with any $N \geq 0$ in place of 2).

To fix $f_c,g_c$, lwt's make the following
choice:
 $$
\begin{aligned}
\pi: X &\longrightarrow X \\
\pi(t^n) &= \delta_{n\geq 0}t^n.
\end{aligned}
$$

\noindent
Then set $f_c=\pi \circ f$, etc.
Let's compute $[f_c,g_c]$.
[Picture]. We obtain $\text{Tr}[f_c,g_c]=2$.
\end{example}

\noindent
The next proposition is
``Tate's definition of the residue''.
Tate was trying to generalize the Residue
Theorem,
i.e. that $\sum_{p \in C}\text{Res}_p
\omega=0$
for points on a curve $C$ and a differential
form $\omega$.

\begin{proposition}
\label{vertex-algebras-proposition-trace}
For any choice of splittings,
$f=f_c+f_d$, etc, and for all $f,g \in
\mathbb{C}((t))$,
$$
\text{Tr}[f_c,g_c]=\text{Res}_tf\cdot dg.
$$
\end{proposition}

\section{Representing the endomorphisms
algebra on the charged fermions}
\label{vertex-algebras-section-end-rep-fermions}

We want to represent $\text{End}(X)$ on $
\Lambda^{\infty/2}$.
We start presenting the naive idea.
Let $x \in X$ and $\varphi \in X^*$.
We already have
$$
\rho(x)=x \wedge(-) \in
\text{End}(\Lambda^{\infty/2})
$$
and
$$
\rho(\varphi)=\sum_{j=0}^\infty(-1)^j
\varphi(x_{i_j})x_{i_0}\wedge
x_{i_1}\wedge\ldots\wedge \widehat{x_{i_j}}
\wedge \ldots.
$$

\noindent
We already saw that
$\rho(\hat{x}_i)\rho(\hat{\varphi}_j)+
\rho(\hat{\varphi}_j)\rho(\hat{x}_i)
=\delta_{ij}I$
(here I'm identifying
$\hat{x}_i \equiv t^i \in \mathbb{C}((t))=X$,
 $\varphi_i \in X^*$ is
$\hat{\varphi}_i (\sum c_jt^j)=c_j$).
See Exercise \ref{vertex-algebras-exercise-something-like}.

If $f \in \text{End}_d$, we can think of $f$
as
$$
f=\sum_{i=0}^a x_i \otimes \varphi_i,
$$
where $\varphi_i\to 0$ as $i \to \infty$.

Here $x_i \to 0$ means
$\forall  N \geq 0$ $\exists n_0$ such that
$x_n \in t^N \mathbb{C}[[t]]$ for all $n\geq
n_0$.
And $\varphi_i \to 0$ means $\forall  N \geq
0$ $\exists n_0$
such that
$\varphi_n(t^{-N}\mathbb{C}[[t]])=0$ for all
$n\geq n_0$.

One can confirm that for $f \in
\text{End}_d$,
$$
\rho_d(f):=\sum_i\rho(x_i)\rho(\varphi_i)
$$
acts a finite sum,
when applied to any fixed vector of
$\Lambda^{\infty/2}$.
So
$$
\rho_d:\text{End}_d \to
\text{End}(\Lambda^{\infty/2}X)
$$
is well defined.

For  $\text{End}_c \ni f$, we can write
$$
f=\sum_i x_i \otimes \varphi_i,
$$
where $x_i \to 0$.
Now $\rho_d(f)$ makes no sense, but
$$
\rho_c(f):=\sum_{i=0}^\infty\rho(\varphi_i)
\rho(x_i)
$$
does.

\begin{exercise}
\label{vertex-algebras-exercise-representation-of-end}
Confirm that $\rho_d: \text{End}_d \to
\text{End}(\Lambda^{\infty/2})$
and $\rho_c:\text{End}_c \to
\text{End}(\Lambda^{\infty/2}X)$
are morphisms of Lie algebras.
\end{exercise}

\noindent
(Neither of these morphisms work at the
level of associative algebras.)

\medskip\noindent
We notice that,
for $f \in \text{End}_f(X)$,
$$
\rho_d(f)-\rho_c(f)=\text{Tr}(f).
$$
This means
$$
\text{End}_d \oplus \text{End}_c
\xrightarrow{(\rho_d,\rho_c)}
\text{End}(\Lambda^{\infty/2}X)
$$
descends to a morphism
$$
\xymatrix{
\text{End}_d \oplus
\text{End}_c\ar[r]^{(\rho_d,\rho_c)}\ar[d]
&\text{End}(\Lambda^{\infty/2}X)\\
\mathfrak{gl}^\flat (X).\ar[ur]_\rho
}
$$

\begin{theorem}
\label{vertex-algebras-theorem-tate}
Let $X$ be a Tate vector space.
There exists a natural representation
of $\mathfrak{gl}^\flat(X)$ on
$\Lambda^{\infty/2}(X)$.
\end{theorem}

\noindent
Next time we'll apply this to the case
$X$ itself is already a Lie algebra!

\section{The 26-dimensionality of the
universe}
\label{vertex-algebras-section-26-dimensionality-universe}

Last time: $X$ a Tate vector space (for us
$X=\mathbb{C}((t))$).
We say that
$\text{End}(X)=\text{End}_c(X)+\text{End}
_d(X)$
comes with a canonical central extension
(as a Lie algebra
$$
\xymatrix{
0\ar[r]&\mathbb{C}\ar[r]&
\mathfrak{gl}(X)^\flat\ar[r]&
\text{End}(X)\ar[r]&0
}
$$
with
$$
\mathfrak{gl}(X)^\flat=
\frac{\text{End}_c \oplus \text{End}_d \oplus
\mathbb{C}}{
\left\langle{}(h,-h,0)=(0,0,\text{Tr}(h))|h \in
\text{End}_f(X)\right\rangle{}},
$$
and there is a canonical representation of
$fl(X)^\flat$ on the semi-infinite wedge
space
$\Lambda^{\infty/2}(X)=\left\langle{}i_{i_0}\wedge
e_{i_1}\wedge\ldots\right\rangle{}$
defined by
\begin{equation}
\label{vertex-algebras-equation-rep}
\begin{aligned}
\rho_d(\underbrace{\sum_ix_i\varphi_i}
_{\in\text{End}_d?})&=
\sum_i\rho(x_i)\rho(\varphi_i)\\
\rho_c(\underbrace{\sum_ix_i\varphi_i}_{\in
\text{End}_c})&=
-\sum_i\rho(\varphi_i)\rho(x_i),
\end{aligned}
\end{equation}
where
$$
\rho(x)=x\wedge(-),\qquad
\rho(\varphi)=
\sum_{j=0}^\infty(-1)^j \varphi(e_{i_j})
e_{i_0}\wedge e_{i_1}\wedge \ldots \wedge
\widehat{e_{i_j}}\wedge\ldots
$$


\noindent
Now suppose $X$ itself were a Lie algebra.

\begin{example}
\label{vertex-algebras-example-with-X-lie-algebra}
$X=\mathbb{C}((t))$ as above, but we identify
$t^m \rightsquigarrow L_m$ in Vir @ $c=0$.
So  $[t^m,t^n]=(m-n)t^{m+n}$.
To avoid confusion, let's write $t^m$ as
$L_m$ in fact.
In such situation, we have a linear map
$\text{ad}:X \to \text{End}(X)$.

We can pull back $\mathfrak{gl}(X)^\flat$
$$
\xymatrix{
0\ar[r]&\mathbb{C}\ar[r]
&\mathfrak{gl}(X)^\flat\ar[r]^\pi
&\text{End}(X)\ar[r]&0\\
&& \hat{X}\ar[r]\ar[u]& X\ar[u]^{\text{ad}}.
}
$$
Explicitly,
$$
\hat{X}=
\{(x,A) \in X \oplus \mathfrak{gl}(X)^\flat
|\text{ad}(X)=\pi(A) \text{ in
}\text{End}(X)\}.
$$
\end{example}

\begin{exercise}
\label{vertex-algebras-exercise-pushout2}
Similarly to Exercise \ref{vertex-algebras-exercise-pushout},
show $\hat{X}$ fits into a sequence
$$
\xymatrix{
0\ar[r]&\mathbb{C}\ar[r]&\hat{X}\ar[r]&
X\ar[r]&0.
}
$$
\end{exercise}
So we get a (canonical!) central extension of
$X$:
$$
\xymatrix{
0\ar[r]&\mathbb{C}\ar[d]^{\text{id}}\ar[r]
&\mathfrak{gl}(X)^\flat\ar[r]^\pi
&\text{End}(X)\ar[r]&0\\
0\ar[r]&\mathbb{C}\ar[r]&\hat{X}\ar[u]\ar[r]&
X\ar[u]^{\text{ad}}\ar[r]&0
}
$$
For $X=\text{Der}(\mathbb{C}[t^{\pm 1}]=$
centerless Vir,
we will find the canonical central extension
is
Vir  @ $c=-26$. (i.e. charge is  $-26$.)

\noindent
What might be amazing here is that any
infinite-dimensional
Lie algebra has a central extension --- just
by being infinite-dimensional.

\medskip\noindent
We'll exploit
$$
\Lambda^{\infty/2}(X)
=\left\langle{}L_{m_0}\wedge L_{m_1}\wedge
L_{m_2}\wedge\ldots\right\rangle{}.
$$
We have quantum fields
$$
\begin{aligned}
\varphi(w)&=\sum_n\varphi_nw^{-n-1},\qquad
\varphi=L_n \wedge(-),\\
\varphi^* (w)&=\sum_n \varphi_n^* w^n,\qquad
\varphi_n^* (\underline{L})=
\sum(-1)^j\varphi(L_{i_j})L_{i_0}\wedge
L_{i_1}\wedge\ldots
\end{aligned}
$$

\noindent
Now for
$$
L_m \in X,\qquad  L_m :L_n \mapsto
(m-n)L_{m+n},
$$
$$
\text{ad}(L_m)=\underbrace{\sum_{n \in
\mathbb{Z}}(m-n) L_{m+n}L_n^*}
_{\in \text{End}_c + \text{End}_d}.
$$
To represent $\text{ad}(L_m)$ on
$\Lambda^{\infty/2}$
we have to split $\text{ad}(L_m)$ into pieces
in $\text{End}_d$ and $\text{End}_c$ and
send each part to their corresponding pieces
according to Equation \ref{vertex-algebras-equation-rep}.

\medskip\noindent
Let's work at the level of fields.

$$
\begin{aligned}
L(w)&=\sum_m \rho(\text{ad}(L_m)) w^{-m-2}\\
&=\sum_{m,n}(m-n)\rho(L_{m+n}L_n^*)w^{-m-2}\\
&=\sum_{m,n}(m-n)\,:\!\varphi_{m+n}
\varphi^*_n\!:\,w^{-m-2}\\
&=\sum_{m,n}(m-n)\,:\!(\varphi_{m+n}w^{-(m+n)
-?})
(\varphi_n^*w^{+n+?})\!:\,
\end{aligned}
$$
where
$$
\rho(L_{m+n}L_n^* =\begin{cases}
\rho(L_{m+n})\rho(L_n^*)\qquad &\text{when
$n\ll 0$} \\
-\rho(L^*_n)\rho(L_{m+n})\qquad &\text{when
}n\gg 0.
\end{cases}
$$
Notice that
$$
\begin{aligned}
\partial_w\varphi(w)&=\partial_w \sum
\varphi_{m+n}w^{-m+nj}\\
&=-\sum(m+n)\varphi_{m+n}w^{-(m+n)-1}\\
&=\sum n \varphi_n^* w^{n-1}
\end{aligned}
$$

\noindent
since $m-n=(m+n)-2n$,
so it seems we should consider
$$
-\,:\!(\partial \varphi)\varphi^*\!:\,
-2\,:\!\varphi(\partial \varphi^*)\!:\,k
$$

\medskip\noindent
So, we have a vertex superalgebra
$V=\Lambda^{\infty/2}$, with quantum fields
$\varphi$ and $\varphi^*$, OPE relation
$[\varphi_\lambda \varphi^*]=1$.

We are defining a field
$$
L=-\,:\!(\partial \varphi)\varphi^*:\,
-2\,:\!\varphi(\partial \varphi^*)\!:\,,
$$
(which is a quantum field, $L(w)=\sum_m
L_mw^{-m-2}$),
adn we want to know/check
$$
[L_m,L_n]=(m-n)L_{m+n}+\delta_{m,-n}
\frac{m^3-m}{12}CI_{\Lambda^{\infty/2}}.
$$
So, what's $c$? Let's do it.

\begin{itemize}
\item (Method 1.) Brute force.
\item (Method 2.) Use mathematica (Thielman's
package).
\item (Method 3.) Recall $\alpha=\,:\!\varphi
\varphi^*\!:\,$,
satisfies $[\alpha_\lambda \alpha]=\lambda$
and
$L^0=\frac{1}{2}\,:\!\alpha \alpha\!:\,$
is Virasoro @ $c=1$.
\end{itemize}
$$
[L^0_\lambda L^0]-TL^0+2\lambda
L^0+\frac{\lambda^3}{12}.
$$
$$
\begin{aligned}
L^0&=\frac{1}{2}\,:\!\alpha \alpha\!:\,\\
&=\frac{1}{2}
\,:\!(\,:\!\varphi
\varphi^*\!:\,)(\,:\!\varphi
\varphi^*\!:\,)\!:\,\\
&=\ldots\\
&=\frac{1}{2}(\,:\!(T\varphi)\varphi^*\!:\,+
\,:\!\varphi(T\varphi^*)\!:\,.
\end{aligned}
$$

How to see this?
Recall Borcherds identity (Equation
\ref{vertex-algebras-equation-Borcherds-identity}):
$$
\begin{aligned}
&\sum_{j \geq
0}\binom{m}{j}(a_{(n+j)}b)_{(m+k-j)}c\\
&=\sum_{j \geq 0}(-1)^j \binom{n}{j}
(a_{(a+m-j)}b_{(k+j)}c-(-1)^nb_{(n+k-j)}
a_{(m+j)}c.
\end{aligned}
$$
So put
$$
\begin{matrix}
a=\varphi & m=0\\
b=\varphi^*  & n=-1\\
c=\,:\!\varphi \varphi ^*\!:\,& k=-1.
\end{matrix}
$$

Then
$$
LHS=(a_{(-1)}b)_{(-1)}c=\,:\!(\,:
\!\varphi\varphi^*\!:\,)
(\,:\!\varphi\varphi^*\!:\,)\!:\,
$$
and
$$
RHS=\sum_{j \geq
0}(\varphi_{(-i-j)}\varphi^*_{(-1+j)}
\varphi_{(-1)}\varphi^*
+\varphi^*_{(-2-j)}\varphi_{(j)}\varphi_{(-1)
}\varphi^*).
$$
$$
\text{[Picture]}
$$
$$
RHS=\varphi_{(-1)}\varphi^*_{(-1)}\alpha+
\varphi_{(-2)}\varphi^*_{(0)}\alpha
+\varphi^*_{(-2)}\varphi_{(0)}\alpha.
$$

Note:
$$
\varphi_{(0)}\alpha=-\alpha_{(0)}\varphi=+
\varphi.
$$
In general
$$
b_{(0)}a=-\sum_{j\geq
0}\frac{1}{j!}T^j(a_{(j)}b)
=-a_{(0)}b+T(\text{stuff}).
$$

\begin{remark}
\label{vertex-algebras-remark-quotient}
Recall (ref?) that if $V$ is any vertex
algebra,
$V/TV$,
$[\overline{a},\overline{b}]=a_{(0)}b$,
is a Lie algebra.
\end{remark}

\noindent
So,
$$
RHS=\varphi_{(-1)}\varphi^*_{(-1)}\alpha
-\varphi_{(-2)}\varphi^*+\varphi^*_{(-2)}
\varphi.
$$
Now, since in general
$$
b_{(n)}a=-(-1)^{p(a)p(b)}\sum_{j \geq
0}\frac{(-1)^j}{j!}
T^j(a_{(n+j)}b,
$$
then
$$
\varphi^*_{(-2)}\varphi=\sum_{j\geq
0}\frac{(-1)^j}{j!}
T^j(\varphi_{(-2+j)}\varphi^*.
$$
{\bf But this led to mistaken calculations.}
I was trying to do the following: we know
that
$\alpha=\,:\!\varphi \varphi^*\!:\,$
satisfies
$[\alpha_\lambda \alpha]=\lambda$,
and $L^0=\frac{1}{2}\,:\!\alpha \alpha\!:\,$
is
$\text{Vir}[L^0_\lambda L^0]=TL^0+2\lambda
L^0+\frac{\lambda^3}{12}$
($c=1$).
Also, if $B=L^0+kT_\alpha$ then
$$
[B_\lambda B]=TB+2\lambda
B+\frac{\lambda^3}{12}c_k,\qquad
c_k=1-12k^2.
$$
Writing $B$ in terms of $vp,\varphi^*$,
one gets \underline{something} like
$$
B=b\,:\!(T\varphi)\varphi^*\!:\,
+(1-b)\,:\!\varphi(T\varphi^*)\!:\,.
$$
The correct answer is: if
$L=-\,:\!(T\varphi)\varphi^*\!:\,-2\,:
\!\varphi(T\varphi^*)\!:\,$,
we may verify that
$$
[L_\lambda L]=TL+2\lambda L-
\frac{26}{12}\lambda^3.
$$

\medskip\noindent
In summary, today we did the following.
For $\varphi, \varphi^*$ odd,
$[\varphi_\lambda \varphi^*]=1$,
we define the operator
$$
L=-\,:\!(T\varphi)\varphi^*\!:\,=2\,:
\!\varphi(T\varphi^*)\!:\,.
$$
Then compute and find out that
 $$
[L_\lambda L]=TL+2\lambda
L+\frac{c}{12}\lambda^3,
$$
where $c=-26$.

That is, after computing the central
extension
of the Virasoro algebra, we turn to the
vertex algebra language to find that the
central
charge is $-26$.

Even more explicitly: by the Tate vector
space construction
we have the central extension
$$
\xymatrix{
0\ar[r]&\mathbb{C}C\ar[r]&\mathfrak{gl}(X)
^\flat\ar[r]&\mathfrak{gl}(X)\ar[r]&0
}
$$
where $W=\mathbb{C}((t))$ is the Witt
algebra,
with bracket $[L_m,L_n]=(m-n)L_{m+n}$.

And then we do the pullback in the following
way:
$$
\xymatrix{
0\ar[r]&\mathbb{C}K\ar[d]^{\text{id}}\ar[r]
&\mathfrak{gl}(W)^\flat\ar[r]^\pi
&\text{End}(X)\ar[r]&0\\
0\ar[r]&\mathbb{C}K\ar[r]&W \oplus \mathbb{C}
K\ar[u]\ar[r]
& W\ar[u]^{\text{ad}}\ar[r]&0
}
$$
Then we get a representation in which $K$
goes to the identity,
namely
$$
\mathfrak{gl}(X)^\flat \text{ acts on }
\Lambda^{\infty/2},
\qquad K \mapsto \text{Id}.
$$
So, the vertex algebra language allows us to
compute
the central charge of the new algebra we
obtained,
$W \oplus \mathbb{C}K$,
and realise it's $-26$. So the bracket in the
new algebra is
 $$
[L_m,L_n]=(m-n)L_{m+n}+\frac{m^3-m}{12}(-26)
K.
$$

\section{Lattices and the neighbor
construction}
\label{vertex-algebras-section-lattices}

\begin{definition}
\label{vertex-algebras-definition-lattice}
A {\it lattice} (for us) is a discrete
subgroup $L \subset \mathbb{R}^n$,
$L \simeq \mathbb{Z}^n$, and such that
$(\alpha,\beta) \in \mathbb{Z}$
for all $\alpha, \beta \in L$,
where $(\cdot,\cdot):\mathbb{R}^n \times
\mathbb{R}^n \to \mathbb{R}$
is the standard bilinear form
$((u_1,\ldots,u_n),(v_1,\ldots,v_n))=\sum
u_iv_i$.
\end{definition}

\noindent
In particular, $L$ contains a basis of the
ambient space $\mathbb{R}^n$.

\begin{example}
\label{vertex-algebras-example-lattices}
\begin{enumerate}
\item $\mathbb{Z}^n \subset \mathbb{R}^n$ is
a lattice.
\item $A_2$, the root lattice of
$\mathfrak{sl}_3$, is a lattice.
Recall this looks like a triangular tiling
with
$(\alpha_1,\alpha_1)=2$,
$(\alpha_2,\alpha_2)=2$
and $(\alpha_1,\alpha_2)=-1$, which says
that the angle between $\alpha_1$ and
$\alpha_2$ is $120$ degrees.
\end{enumerate}
\end{example}

\begin{definition}
\label{vertex-algebras-definition-even-lattice}
A lattice is called {\it even} if
$(\alpha,\alpha)\in 2 \mathbb{Z}$
for all $\alpha \in L$.
\end{definition}

\noindent
(That is, the product of every element with
{\it itself} is even,
not that $(\alpha,\beta)\in 2 \mathbb{Z}$ for
all $\alpha,\beta \in L$.)

So $A_2$ is even. In fact, whenever the basis
elements
are even, say
$(\alpha,\alpha),(\beta,\beta)\in
2\mathbb{Z}$,
then the lattice is even since
$(\alpha \pm \beta,\alpha \pm
\beta)=(\alpha,\alpha)
\pm 2(\alpha,\beta)+(\beta,\beta) \in
2\mathbb{Z}$ too.

\begin{definition}
\label{vertex-algebras-definition-dual-lattice}
The {\it dual} of $L$ is
$$
L^\vee=\{ x \in \mathbb{R}^n|(x,\alpha) \in
\mathbb{Z}\forall \alpha \in L\}.
$$
\end{definition}

\noindent
Clearly $L \subset L^\vee$.
But $L^\vee$ might be strictly larger.

\begin{exercise}
\label{vertex-algebras-exercise-discriminant}
$L^\vee/L$ is a finite group.
\end{exercise}

\noindent
Notice that for  $\gamma,\delta \in L^\vee$
it might happen that $(\gamma,\delta) \not\in
\mathbb{Z}$.

For $L=A_2$ we find that
$$
L^\vee = A_2 \cup (\omega+A_2) \cup  (2\omega
+ A_2).
$$
So $L^\vee/L \simeq \mathbb{Z}/3$ as groups.

For $L = \mathbb{Z}^n$, we have
$L^\vee=\mathbb{Z}^n=L$, that is,
$\mathbb{Z}^n$ is self dual.
On the other hand $\mathbb{Z}^n$ is not even.

Are there any other even self-dual lattices
(other than $\{0\}$)?

The simplest nontrivial example is
$$
E_8=\{(x_1,x_2,\ldots,x_8) \in \mathbb{Z}^8
\cup (\frac{1}{2}+\mathbb{Z})^8|
\sum_{i=1}^8 \equiv 0 \text{mod }2\}.
$$

\begin{lemma}
\label{vertex-algebras-lemma-e8}
$E_8$ is even and self-dual.
\end{lemma}

\begin{theorem}
\label{vertex-algebras-theorem-self-dual-lattices}
\begin{enumerate}
\item If $L \subset \mathbb{R}^n$ is even and
self dual,
then $8|n$

\item If, also, $n=8$, then $L \simeq E_8$.
\end{enumerate}
\end{theorem}

\begin{remark}
\label{vertex-algebras-remark-discriminant-form}
Let $L$ be an even lattice.
On the group $D=L^\vee/L$ we define
$$
\begin{aligned}
q: D &\longrightarrow \mathbb{Q}/\mathbb{Z}
\\
q(a) &=(\alpha,\alpha)/2 \text{ mod
}\mathbb{Z}.
\end{aligned}
$$

\noindent
Then $q$ is a well-defined quadratic form.

Indeed, for $\beta \in L$,
$$
q(\alpha+\beta)=\frac{(\alpha+\beta,\alpha+
\beta)}{2}
=\underbrace{\frac{(\alpha,\alpha)}{2}}_{=
q(\alpha)}
+\underbrace{\frac{2(\alpha,\beta)}{2}}_{\in
\mathbb{Z}}
+\underbrace{\frac{(\beta,\beta)}{2}}_{\in
\mathbb{Z}}.
$$
We call $(D,q)$ the {\it discriminant form}
of $L$.
\end{remark}

\medskip\noindent
Now we explain the neighbour
construction/orbifold.
Let $L$ be a self-dual even lattice.
Let $\phi:L \to \mathbb{Z}/2$ a homomorphism
(i.e. if $\{e_1,\ldots,e_n\}$ is a basis of
$L$,
set $\phi(e_i)=\varepsilon_i \in \{0,1\}$
and $\phi(\sum m_i e_i)=\sum m_i
\varepsilon_i \text{ mod }2$.

Assume $\phi$ is nontrivial and surjective,
so that its kernel
$L_0 = \Ker \phi \subset L$ has index 2.
That is, $L/L_0 \simeq \mathbb{Z}/2$.

Now we have $L_0 \underbrace{\subset
}_{\text{index 2}}L
= L^\vee
\underbrace{\subset}_{\substack{\text{index
2} \\ \text{exercise}}}
L_0^\vee$.

\begin{remark}
\label{vertex-algebras-remark-discriminant-iso}
$D=L_0^\vee/L_0 \simeq
\mathbb{Z}/2\times\mathbb{Z}/2$.
\end{remark}

\noindent
Why? If $\omega \in L_0^\vee$, then
$2\omega$… Exercise.

What about $(D,q)$? I claim there are two
possibilities
(up to $\simeq$):
\begin{equation}
\label{vertex-algebras-equation-cases}
\begin{pmatrix}
0&0\\
0&1/2
\end{pmatrix}\qquad
\begin{pmatrix}
0&1/4\\
0&3/4
\end{pmatrix}
\end{equation}

\noindent
Exercise.

Suppose (by choice of $\phi$) we are in the
first case.
Then $L=L_0 \cup  (\alpha+L_0)$,
$q(\beta)=0$, i.e.
$(\beta,\beta)\equiv 0 \text{ mod }2$,
$$
\begin{pmatrix}
L_0&\beta+L_0\\
\alpha+L_0&\gamma+L_0
\end{pmatrix}.
$$
Define $L^{\text{orb}(\phi)}=L_0 \cup
(\beta+L_0)$.

\begin{lemma}
\label{vertex-algebras-lemma-even-lattice}
$L^{\text{orb}(\phi)}$ is an even lattice.
\end{lemma}

\begin{proof}
Exercise. Idea: $(L_0,L_0) \subset
\mathbb{Z}$ and
$(L_0,\beta+L_0)\subset \mathbb{Z}$ by
definition.
\end{proof}

\noindent
For $\beta,\beta' \in \beta+L_0$,
$$
\begin{aligned}
(\beta',\beta)&=(\beta+\alpha,\beta),\qquad
\alpha \in L_0\\
&=\underbrace{(\beta,\beta)}_{\in \mathbb{Z}}
+\underbrace{(\alpha,\beta)}_{\in \mathbb{Z}}
\in \mathbb{Z}
\end{aligned}
$$

\noindent
Also,
$(L^{\text{orb}(\phi)})^\vee=L^{\text{orb}
(\phi)}$.
Indeed
$$
L_0 \underbrace{\subset}_{\text{index
2}}L^{\text{orb}(\phi)}
\underbrace{\subset}_{\text{by above}}
(L^{\text{orb}(\phi)})^\vee
\underbrace{\subset}_{\text{index 2}}
L_0^\vee.
$$

\begin{theorem}[Niemer]
\label{vertex-algebras-theorem-niemer}
If we consider the set $\Gamma_n$ of all
even self-dual lattices of rank $n$ ($8|n$)
as a graph, with edge whenever there
exists an operation $L \to
L^{\text{orb}(\phi)}$,
then $\Gamma_n$ is connected for all $n$.
\end{theorem}

\begin{remark}
\label{vertex-algebras-remark-reverse-orbifold}
There is a reverse orbifold construction.
\end{remark}

\noindent
Let $L$ be even self dual, $L=L_0 \cup  L_1$,
and
$$
\begin{pmatrix}
L_0&L_+\\
L_1&L_-
\end{pmatrix}
$$
$L_0^\vee=L_0 \cup  L_1 \cup  L_+ \cup  L_-$,
say $L_+$ is $q(\alpha)=0$ for all $alp \in
L_+$
and $q(\beta)=1/2$ for all $\beta \in L_-$
since we are fixed in the first case of
Equation \ref{vertex-algebras-equation-cases}.

Define
\begin{align*}
\psi: L^{\text{orb}(\phi)} &\longrightarrow
\mathbb{Z}/2 \\
 \psi(\alpha)&=
\begin{cases}
0\qquad \text{for }\alpha \in L_0& \\
1\qquad &\text{for }\alpha \in L_+.
\end{cases}
\end{align*}

\noindent
By definition,
 $\Ker(\psi)=(L^{\text{orb}(\phi)})_0=L_0
 \subset L^{\text{orb}(\phi)}$.
So $(L^{\text{orb}(\phi)})_0^\vee=L_0^\vee$
and the discriminant form
$(L^{\text{orb}(\phi)})_0^\vee/L^{\text{orb}
(\phi)}_0$
brecovers the same picture.
$$
(L^{\text{orb}(\phi)})^{\text{orb}(\phi)}=L.
$$
$$
\begin{matrix}
L
\underbrace{0}_{L^{\text{orb}(\phi)}_0=L_0} &
0 & M\\
0&1/2
\end{matrix}
$$
\begin{definition}
\label{vertex-algebras-definition-root-system-lattice}
Let $L \subset \mathbb{R}^n$ be an even
lattice.
The {\it root system} of $L$ is
$$
\Delta(L)=\{\alpha \in L|(\alpha,\alpha)=2\}.
$$
\end{definition}

\begin{proposition}
\label{vertex-algebras-proposition-root-syste-latti-lie}
The root system of an even lattice is a root
systems.
\end{proposition}

\begin{example}
\label{vertex-algebras-example-e8-root-system}
The root system of $E_8$ (the lattice)
is $E_8$ (the root system).
\end{example}

\noindent
If $\Delta(L_1) \not \simeq \Delta(L_2)$
then $L_1\not\simeq L_2$.

In fact, in $\mathbb{R}^{16}$ there exists
two even self-dual lattices,
with root systems $E_8 \oplus E_8$ and
$D_{16}$.

For rank 24, there exist (coincidentally) 24
even self-dual lattices
called  {\it Neimeier lattices}.
They are constructed as a graph going from
one another
by applying the orbifold construction on
different homomorphisms $\phi$.
Eventually the graph looks like this:
$$
\xymatrix{
\bullet \ar@{-}[dr]\\
& A_1^{24}\ar@{-}[r]& \Lambda\\
\bullet \ar@{-}[ur]
}
$$
where $\Lambda$ is the {\it Leech lattice},
which has {\it empty} root system,
that is, $\Lambda$ contains $0$, and no
vectors of norm $2$,
and 196560 (?) vectors of squared norm 4.

\medskip\noindent
There is something called the Siegel mass
formula.
Consider the orthogonal group of a lattice,
namely
$$
\begin{aligned}
\text{Aut}(L)&=\{g \in
\text{GL}_n(\mathbb{R})
|(g \alpha,g \beta)=(\alpha,\beta)\forall
\alpha,\beta \in L\text{ and }
g(L) \subset L\}\\
&=\{g \in O_n(\mathbb{R})|g(L) \subset L\},
\end{aligned}
$$
which is a finite group.

Then
$$
\sum_{L \in \Gamma_n}\frac{1}{\#
\text{Aut}(L)}
=\frac{|B_{n/2}|}{n}
\prod_{1 \leq j \leq  n/2}\frac{|B_{2j}|}{4j}
$$
where $B_k$ is a Bernoulli number.
(See Wikipedia page for Neimeier lattice.)



\section{Vertex algebra modules}
\label{vertex-algebras-section-vertex-algebra-modules}

\noindent
Let $V$ be a vertex algebra, so we have,
for every $v \in V$ a quantum field
$Y(a,z)$ on $V$, also $|0\rangle \in V$
and $T:V \to V$,
such that various properties hold,
in particular
$$
Y(a_{(n)}b,z)=Y(a,z)_{(n)}Y(b,z),\qquad
Y(|0\rangle,z)=\text{Id},
\qquad [T,Y(a,z)]=\partial_zY(a,z).
$$

\begin{definition}
\label{vertex-algebras-definition-v-module}
A {\it V-module} is a vector space $M$ and,
for all $a \in V$, a quantum field $Y^M(a,z)$
on $M$
(i.e. $Y^M(a,z)=\sum_{n \in
\mathbb{Z}}a^M_{(n)}z^{-n-1}$
such that $\forall x \in M$, $\exists N$
such that $a_{(n)}x=0$ for all $n \geq N$)
such that
\begin{enumerate}
\item
$Y^M(a_{(n)}b,z)=Y^M(a,z)_{(n)}Y^M(b,z)$.
\label{vertex-algebras-item-vmod1}

\item $Y^M(|0\rangle,z)=I_M$.
\label{vertex-algebras-item-vmod2}

\item
\label{vertex-algebras-item-vmod3}
Ambiguity! We might either omit (3), or…
demand $M$ come with an endomorphism $T_M:M
\to M$ such that
$[T_M,Y^M(a,z)]=\partial_z Y^M(a,z)$.
\end{enumerate}
\end{definition}

But we will not use \ref{vertex-algebras-item-vmod3}, i.e.
for us a V-module just satisfies
\ref{vertex-algebras-item-vmod1} and \ref{vertex-algebras-item-vmod2}.

\medskip\noindent
Let $\mathfrak{g}$ be a finite-dimensional
Lie algebra,
$(\cdot,\cdot):\mathfrak{g} \times
\mathfrak{g} \to \mathbb{C}$
an invariant symmetric bilinear form, and
$\hat{\mathfrak{g}}$ the
corresponding affine Lie algebra, i.e.
$\hat{\mathfrak{g}}=\mathfrak{g}[t^{\pm 1}]
\oplus \mathbb{C}K$,
$[at^m,bt^n]=[a,b]t^{m+n}+m\delta_{m,-n}(a,b)
K$,
$K$ central.
Recall the affine vertex algebra
$V^k(\mathfrak{g})=U(\hat{\mathfrak{g}})
\otimes_{U(\hat{\mathfrak{g}}_+)}
\mathbb{C}|0\rangle$,
$K |0\rangle=k |0\rangle$,
$\mathfrak{g}[t]|0\rangle=0$,
which is a
$\hat{\mathfrak{g}}$-representation with
quantum fields
for all $a \in \mathfrak{g}$,
$$
Y(a_{-1}|0\rangle,z)=a(z)=\sum_{n \in
\mathbb{Z}}at^m \cdot z^{-m-1},
$$
where we denote $a_m=at^m$.

Recall that the Lie bracket definition for
$\hat{\mathfrak{g}}$ implies that
$$
[a(z),b(w)]=[a,b](w)\delta(z,w)+k(a,b)
\partial_w\delta(z,w)I,
$$
i.e. in $\lambda$.$[a_\lambda
b]=[a,b]+\lambda(a,b)k|0\rangle$.

Recall also that a
$\hat{\mathfrak{g}}$-module $M$ is called
smooth if $\forall x \in M$ and $a \in
\mathfrak{g}$ there
is $N$ such that $at^n x=0$ for all $n \geq
N$.

It $M$ is a smooth
$\hat{\mathfrak{g}}$-module,
then
$$
Y^M(a,z)=\sum_{n \in
\mathbb{Z}}(at^n)|_{M}z^{-n-1}
$$
is a quantum field.

If $M$ is a smooth
$\hat{\mathfrak{g}}$-module of level $k$
(i.e. $K|_{M}=k I_M$), then
$$
[Y^M(a,z),Y^M(b,w)]
=Y^M([a,b],w)\delta(z,w)+(a,b)k\partial_w
\delta(z,w)I_M.
$$
The following proposition says that there is
an equivalence of
categories between smooth level $k$
$\hat{\mathfrak{g}}$-modules
and $V^k(\mathfrak{g})$-modules:

\begin{proposition}
\label{vertex-algebras-proposition-equivalence-vmods}
If $M$ is a smooth level $k$
$\hat{\mathfrak{g}}$-module
then there exists a structure of
$V^k(\mathfrak{g})$-module on $M$
such that $Y^M(a_{-1}|0\rangle,z)=Y^M(a,z)$
for all $a \in \mathfrak{g}$.
\end{proposition}

\noindent
As a special case,
$\mathfrak{g}=\mathfrak{h}$ a
finite-dimensional
vector space,
$(\cdot,\cdot):\mathfrak{h}\times\mathfrak{h}
\to\mathbb{C}$
some nondegenerate symmetric bilinear form,
$k=1$.
Then we get
$V^k(\mathfrak{g})=H(\mathfrak{h})$
Heisenberg.
Let $\mu \in \mathfrak{h}$, then we have the
Fock module
$H^\mu =
U(\hat{\mathfrak{h}})
\otimes_{U(\hat{\mathfrak{h}}_+)}\mathbb{C}
|\mu\rangle$,
$$
ht^m|\mu\rangle=\begin{cases}
0\qquad &\text{if }m\geq 1 \\
(\mu,h)|\mu\rangle\qquad &\text{if }m=0.
\end{cases}
$$
As a vector space,
$H^\mu\cong\mathbb{C}[h^{i_1}_{(n_1)},h^{i_2}
_{(n_2)}|
\substack{i=1,\ldots,\dim \mathfrak{h} \\
n=1,2,\ldots}]|\mu\rangle$.

$H^\mu$, $\mu=0$ is the Heisenberg vertex
algebra itself,
the $H^\mu$ for various $\mu$ are all modules
over $H=H^0$.

We saw an example in
$F^{\text{ch}}=\bigoplus_{c \in
\mathbb{Z}}F^c$.
[Picture of upside-down parabola.
The center column is
 $F^{(0)}\cong H \cong H(\mathbb{C})\cong
 \mathbb{C}[h_{-1},h_{-2},\ldots]
\cong H^0$.
In fact we can show that the $H^i$ are
V-modules
over $H^0$.]

\begin{remark}
\label{vertex-algebras-remark-super}
The fact that $F= \bigoplus_{c \in
\mathbb{Z}}F^c$ is
a vertex superalgebra implies that $F^c$ with
$a \in F^0$ acting on $F^c$
by restriction of $Y(a,z)$, is a
$F^0$-module.
\end{remark}

\noindent
Can we go the other way and build a vertex
(super)algebra from an infinite sum of
$H(\mathfrak{h})$-modules? Yes.

\begin{remark}[Technical]
\label{vertex-algebras-remark-technical}
In $H(\mathbb{C})$, we defined
$L=\frac{1}{2}\,:\!hh\!:\,$,
which is a Virasoro vector.
$$
[L_\lambda L]=TL+2\lambda L
+\frac{\lambda^3}{12}|0\rangle.
$$
($L(z)=\sum_{n \in \mathbb{Z}}L_nz^{-n-2}$.)
We observed that $L_{-1}=\frac{1}{2}\sum_{k
\in \mathbb{Z}}h_{-1-k}h_k$
coincides with $T$.
\end{remark}

\begin{exercise}
\label{vertex-algebras-exercise-check-it}
Check it.
\end{exercise}

\noindent
More generally, if $\{a^1,\ldots,a^n\}$ is a
basis of $\mathfrak{h}$,
and $\{b^1,\ldots,b^n\}$ the dual basis
(i.e., $(a^i,b^j)=\delta_{i,j}$).
Define
$$
L=\frac{1}{2}\sum_{i=1}^{\dim
\mathfrak{h}}\,:\!a^ib^i\!:\,.
$$

\begin{exercise}
\label{vertex-algebras-exercise-virasoro-vector}
$[L_\lambda L]=TL+2\lambda
L+\frac{\lambda^3}{12}cI$,
where $c=\dim(\mathfrak{h})$. We call this
a {\it Virasoro or conformal vector} of
central charge.
$\dim (\mathfrak{h})$.
\end{exercise}

\begin{exercise}
\label{vertex-algebras-exercise-once-again}
$L_{-1}=T$ once again.
\end{exercise}

\noindent
(Consider $\mathfrak{h} \otimes
\mathfrak{h}$, the element
$\sum_i a_i \otimes b_i$ is well-defined
(i.e. depends on $\mathfrak{h}$ and
$(\cdot,\cdot)$
but not the choice of basis $\{a^i\}$.)

So if we chose a different basis $\{a^i\}$
we'd get the same
$L \in H(\mathfrak{h})$.

\begin{definition}
\label{vertex-algebras-definition-conformal-vertex-algebra}
If $V$ contains a Virasoro vector $L$,
such that $T=L_{-1}$,
we call $V$ (more precisely, the pair
$(V,L)$)
a {\it conformal vertex algebra} of ventral
charge $c$.
\end{definition}

\noindent
Now we can reinterpret Exercise
\ref{vertex-algebras-exercise-virasoro-vector}
as saying that the Heisenberg (?) algebra is
conformal.

\begin{remark}
\label{vertex-algebras-remark-third-axiom-vmod}
If $(V,L)$ is a conformal vertex algebra,
and $M$ a $V$-module, then we may define
$T_M:=L_{-1}^M$, i.e.
(where $Y^M(L,Zz)=\sum L^m_n z^{-n-2}$)
then we will have
$$
[T_M,Y^M(a,z)]=\partial_zY^M(a,z).
$$
(So for conformal vertex algebras, axiom
\ref{vertex-algebras-item-vmod3} is superfluous).
\end{remark}

\noindent
Observe that in the Fock space $H^\mu$, we
have
$$
\begin{aligned}
L_{-1}|\mu\rangle&=\frac{1}{2}\sum_i
\left(\sum_{k \in
\mathbb{Z}}\,:\!a^i_{-1-k}b^i_k\!:\,\right)
\left(\sum_{k \in
\mathbb{Z}}a^i_{-1-k}b^i_k\right),
\qquad \text{since they commute}\\
&=\sum_i(a_{-1}b_0^i+a_0^ib_{-1})
|\mu\rangle\\
&=\frac{1}{2}\sum_i(a^i_{-1}|\mu\rangle(b^i,
\mu)
+b^i_{-1}|\mu\rangle (a^i,\mu)\\
&=
\frac{1}{2}(\mu_{-1}|\mu\rangle+\mu_{-1}
|\mu\rangle)
\qquad \substack{\text{lin. alg. fact} \\
\sum_i(\mu,b^i)a^i=\mu}\\
&=\mu_{-1}|\mu\rangle\\
&=\,:\!\underbrace{\mu}
_{\substack{\text{element
of} \\ \text{Heis. alg.}}}
\underbrace{\mu}_{\substack{\text{generator
of} \\ \text{Fock module}}}\!:\,.
\end{aligned}
$$

\noindent
It's also fun to do:
$$
\begin{aligned}
L_0|\mu\rangle&=\frac{1}{2}
\sum_i\left(\sum_k\,:\!a^i_{-k}
b^i_k\!:\,\right)|\mu\rangle\\
&=\frac{1}{2}\sum_i a_0^ib_0^i |\mu\rangle\\
&=\frac{1}{2}\sum_i(b^i,\mu)(a^i,\mu)
|\mu\rangle
\qquad \substack{\text{lin. alg. fact} \\
\sum_i(x,a^i)(y,b^i)=(x,y)}\\
&=\underbrace{\frac{1}{2}(\mu,\mu)}_{\text{a
number}}|\mu\rangle.
\end{aligned}
$$

\medskip\noindent
Let us try to build a vertex (super)algebra
$V_S=\bigoplus_{\mu \in S}H^\mu$, where $S$
is some
subset $S \subset \mathfrak{h}$ (assume $0
\in S$)
such that $T=L_{-1}$, where
$L=\frac{1}{2}\,:\!a^ib^i\!:\, \in H^0$.
In particular
\begin{equation}
\label{vertex-algebras-equation-t-mu}
T|\mu\rangle=\,:\!\mu|\mu\rangle\!:\,\qquad
\text{for all }\mu \in S.
\end{equation}

\noindent
Let $\Gamma^\mu(z)$ denote $Y(|\mu\rangle,z)$
now.
Equation \ref{vertex-algebras-equation-t-mu} implies
$$
\partial_z\Gamma^\mu(z)\underbrace{=}
_{\substack{\text{holds
for}
 \\\text{vertex
algebras} }}[T,\Gamma^\mu(z)]
\underbrace{=}_{\text{now}}
\,:\!\mu(z)\Gamma^\mu(z)\!:\,.
$$
So
\begin{equation}
\label{vertex-algebras-equation-t-mu2}
\partial_z\Gamma(z)=\,:\!\mu(z)\Gamma(z)\!:\,
\end{equation}

\noindent
(Note that this is formally similar to
$\partial_zf(z)=\mu(z)f(z)$, whose solution
is
$f(z)=ce^{\int \mu(z)dz}$. So informally we
should have $\Gamma(z)=c\,:\!e^{\int
\mu(z)dz}\!:\,$.)

So we kind of have
$$
\begin{aligned}
\mu(z)&=\sum_{n \in \mathbb{Z}}\mu_nz^{-n-1},
\int\mu(z)dz&=\mu_0\cdot
\underbrace{\text{log}z}_{\substack{\text{not
very} \\
\text{well defined} }}+\sum_{n \neq
0}\frac{z^{-n}}{-n}\mu_n,\\
\text{exp}(\mu_0\text{log}(z))=z^{\mu_0}.
\end{aligned}
$$

Maybe this means
\begin{equation}
\label{vertex-algebras-equation-maybe-this-means}
\Gamma(z)=c
\text{exp}\left(\sum_{j<0}\frac{z^{-j}}{-j}
\mu
j\right)
\text{exp}\left(\sum_{j>0}\frac{z^{-j}}{-j}
\mu_j\right)z^{\mu_0}.
\end{equation}
This crazy argument turns out to get us to
the right definition.

\medskip\noindent
Now we describe the precise argument.
Firstly $h_0|\mu\rangle=(h,\mu)|\mu\rangle$
for all $h \in \mathfrak{h}$,
$h_{\geq 1}|\mu\rangle=0$.
So
 $$
\underbrace{[h_\lambda |\mu\rangle]}
_{\sum_jh_j|\mu\rangle\cdot\frac{\lambda^j}
{j!}}=(h,\mu)|\mu\rangle.
$$
This implies that
$[h_m,\Gamma^\mu(w)]=(h,\mu)w^m\Gamma^\mu(w)
$.
In particular
$[h_0,\Gamma^\mu(w)]=(h,\mu)\Gamma^\mu(w)$.
It follows that
$$
\Gamma^\mu(w):H^\beta \to H^{\beta+\mu}((w)).
$$
(where $H^\beta$ is where elements of $H$ act
by pairing with $\beta$, i.e.
$\forall x \in H^\beta$, $h_0x=(h,\beta)x$
for all $h \in \mathfrak{h}$.)

Recall that all Fock spaces are spaces of
polynomials,
they are only different by the constant by
which $h_0$ acts.
So they are all isomorphic as Lie algebras
over the oscillator Lie algebra
$\mathfrak{a}'=\bigoplus_{n \neq
0}\mathbb{C}h_n \oplus \mathbb{C}K$.
This means we can introduce
$$
u:H^\beta \to H^{\beta+\mu}
$$
such that
\begin{enumerate}
\item $u(|\beta\rangle)=|\beta+\mu\rangle$,
\item $u(h_nx)=h_nu(x)$ for all  $n \neq 0$
(extend linearly).
\end{enumerate}

\noindent
Consider on $H^\beta$
\begin{equation}
\label{vertex-algebras-equation-quantum-field-vmod}
A=u^{-1}\text{exp}\left(-\sum_{j<0}\frac{z^{-
j}}{-j}\mu_j\right)\Gamma^\mu(z)
\text{exp}\left(-\sum_{j>0}\frac{z^{-j}}{-j}
\mu_j\right)
\end{equation}
$H^\beta \to H^\beta((z))$.
(Which already looks similar to Equation
\ref{vertex-algebras-equation-maybe-this-means}!)

\begin{exercise}
\label{vertex-algebras-exercise-a-and-hn}
$[h_n,A]=0$ for all $n \neq 0$.
\end{exercise}

\noindent
So
$A|_{H^{\beta}}=p_{\mu,\beta}(z)I_{H^\beta}$.

\begin{proposition}
\label{vertex-algebras-proposition-exponentials}
$$
\Gamma^\mu(z)|_{H^\beta}=C_{\mu,\beta}
\text{exp}\left(\sum_{j<0}\frac{z^{-j}}{-j}
\mu_j\right)
\text{exp}\left(\sum_{j>0}\frac{z^{-j}}{-j}
\mu_j\right)z^{(\mu,\beta)}u.
$$
\end{proposition}

\begin{proof}
Consider $\Gamma^\mu(z)|\beta\rangle=
\sum_{n \in
\mathbb{Z}}\Gamma^\mu_{(n)}|\mu\rangle
z^{-n-1}$.
Take two elements $a,b$ with conformal
weights given by
 $L_0a=\Delta(a)a$ and $L_0b=\Delta(b)b$.
Then $L_0(a_{(n)}b)=\Delta(a)+\Delta(b)-n-1$
(exercise: deduce this from the commutator
formula).

We computed earlier that
$$
\Delta(|\beta\rangle)=\frac{(\beta,\beta)}{2}
,\qquad
\Delta(|\beta+\mu\rangle)=\frac{(\beta+\mu,
\beta+\mu)}{2},
$$
and
$$
\Delta(|\mu\rangle)=\frac{(\mu,\mu)}{2}
$$
This implies
$$
\Gamma^\mu_{(-(\beta,\mu)-1)}|\beta\rangle
=P_{\mu,\beta}(z)|\beta+\mu\rangle.
$$
It follows that
$$
P_{\mu,\beta}(z)=c z^{(\mu,\beta)}.
$$
\end{proof}

\medskip\noindent
Here's a summary of this section so far (new
lecture).
$\mathfrak{h}$ complex vector space with
nondegenerate symmetric bilinear form
$(\cdot,\cdot)$.
$V=\bigoplus_{\mu \in S}H^\mu$ vertex
(super)algebra?
extending $H^0$-module structure $H^0
\hookrightarrow H^\mu$.
We saw (putting
$\Gamma^\mu(z)=T(|\mu\rangle,z))$
$$
\Gamma^\mu(z):H^\beta \to
H^{\beta+\mu}[[z^{\pm 1}]],
$$
so $S$ should be closed under addition
(we hadn't said this so far!)
and
\begin{equation}
\label{vertex-algebras-equation-quantum-field-vmod2}
\Gamma^\mu(z)|_{H^\beta}=
C_{\mu,\beta}
\text{exp}\left(\sum_{n<0}\frac{z^{-n}}{-n}
\mu_n\right)
\text{exp}\left(\sum_{n>0}\frac{z^{-n}}{-n}
\mu_n\right)
z^{(\mu,\beta)}u
\end{equation}
(see Equation
\ref{vertex-algebras-equation-quantum-field-vmod}).
Here $u$ is the unique linear map $u:H^\beta
\to H^{\beta+\mu}$
such that
$u(|\beta\rangle)=|\beta+\mu\rangle$,
and $u$ commutes with the action of $h_n$ for
all
$n \in \mathbb{Z}\setminus \{0\}$.
(Intuitively, it takes a monomial and takes
it to
``the same'' monomial.)

Let's assume the constant $C_{\mu,\beta}$ are
all nonzero.
Also want $z^{(\mu,\beta)}$ to be an integer
power of $z$,
so $(\alpha,\beta) \in \mathbb{Z}$ for all
$\alpha,\beta \in S$,
i.e. $S \subset \mathfrak{h}$ should be an
\underline{integral lattice}.
Call $S$ as $Q$ from now on.

In particular
$$
\begin{aligned}
H^\beta&=\mathbb{C}[t^{-1}\mathfrak{h}[t^{-1}
]]\otimes|\beta\rangle,\\
H^{\mu+\beta}&=\mathbb{C}[t^{-1}\mathfrak{h}
[t^{-1}]]\otimes
|\mu+\beta\rangle,
\end{aligned}
$$

\noindent
and
$$
u(\underbrace{
h^{i_1}_{-n_1}h^{i_2}_{-n_2}\ldots
h^{i_s}_{-n_s}\otimes |\beta\rangle}
_{\substack{h_0\text{ acts by} \\ \text{the
scalar }(h,\beta) \in \mathbb{C}}}
=\underbrace{
h^{i_1}_{-n_1}h^{i_2}_{-n_2}\ldots
h^{i_s}_{-n_s}\otimes | \mu+\beta\rangle}
_{\substack{h_0\text{ acts by} \\ \text{the
scalar }(h,\mu+\beta)
\in \mathbb{C}}}.
$$

\noindent
(We can picture this as a lattice with a
vector space
sitting at every point, and a map from every
vector space
with tag $\beta$ to the vector space with tag
$\beta+\mu$.
The vector spaces are infinite-dimensional
and
are drawn as cones with, e.g.,
$|\beta\rangle$ and $|\beta+\mu\rangle$
in the vertex of the cone.)

\medskip\noindent
Notice that in the vector space
$H^\beta$, $h_0=(h,\beta)$ constant, we can
write
$z^{(\mu,\beta)}I_{H^\beta}$ as $z^{\mu_0}$,
therefore, let's write a symbol $e^\mu$ for
the map(s) $u$,
more precisely
$$
e^\mu:V \to V,
$$
whose restriction to $H^\beta$ is
$C_{\mu,\beta}$ times
the unique $u:H^\beta \to H^{\beta+\mu}$ as
above.

Let's also start writing
$|\beta\rangle$ as $e^\beta$.
So now elements of $V$ are linear
combinations
of vectors of the form
$h^{i_1}_{-n_1}\ldots h^{i_s}_{-n_s} \otimes
e^\beta$,
for $\beta \in Q$.
$$
\begin{aligned}
e^\mu(h^{i_1}_{-n_1}\ldots h^{i_s}_{-n_s}
\otimes e^\beta)
&=C_{\mu,\beta} h^{i_1}_{-n_1}\ldots
h^{i_s}_{-n_s}\otimes e^{\beta+\mu}\\
z^{\mu_0}(h^{i_1}_{-n_1}\ldots
h^{i_s}_{-n_s}\otimes e^\beta)
&=z^{(\mu,\beta)}h^{i_1}_{-n_1}\ldots
h^{i_s}_{-n_1} \otimes e^\beta.
\end{aligned}
$$

\noindent
Now
\begin{equation}
\label{vertex-algebras-equation-quantum-field-vmod3}
\Gamma^\mu(z)=e^\mu
\text{exp}\left(\sum_{n<0}\frac{z^{-n}}{-n}
\mu_n\right)
\text{exp}\left(\sum_{n>0}\frac{z^{-n}}{-n}
\mu_n\right)
z^{\mu_0}.
\end{equation}

\noindent
Note this expression works in all of $V$,
and not in a fixed $H^\beta$, as Equation
\ref{vertex-algebras-equation-quantum-field-vmod2}
does!

\begin{exercise}
\label{vertex-algebras-exercise-quantum-field-vmod}
Each $\Gamma^\mu(z)$ is a quantum field.
\end{exercise}

\noindent
Are the $\Gamma^\mu(z)$ mutually local?
We claim that
\begin{equation}
\label{vertex-algebras-equation-mutually-local-vmode}
\Gamma^\mu(z)\Gamma^\beta(w)=
i_{z,w}(z-w)^{(\mu,\beta)}\,:\!\Gamma^\mu(z)
\Gamma^\beta(w)\!:\,
\end{equation}
We can sketch a proof of this:
$$
\begin{aligned}
LHS&=e^\mu e^\beta z^? w^?
\text{exp}\left(\sum_{n<0}\frac{z^{-n}}{-n}
\mu_n\right)
\underbrace{\text{exp}\left(\sum_{n>0}
\frac{z^{-n}}{-n}\mu_n\right)
\text{exp}\left(\sum_{n<0}\frac{w^{-n}}{-n}
\beta_n\right)
}_{\substack{\text{to pass to RHS we have to
swap these terms}\\
\text{(at least morally)}}}
\text{exp}\left(\sum_{n>0}\frac{w-n}{-n}
\beta_n\right)
\end{aligned}
$$

\noindent
Consider $h \in \mathfrak{h}$, $(h,h)=1$ for
simplicity,
$\mu=\beta=h$. Now we claim
$$
\begin{aligned}
\left[ \sum_{m>0}\frac{z^{-m}}{-m}h_m,
\sum_{n<0}\frac{w^{-n}}{-n}h_n \right] &=
-\sum_{n>0}\frac{1}{n}\left(\frac{w}{z}
\right)^n(\mu,\beta)\\
&=\text{log}(1-\frac{w}{z})(\mu,\beta)
\end{aligned}
$$

\noindent
Where the logarithm here is defined by
$\text{log}(1-x)=-x+\frac{1}{2}x^2+\frac{1}
{3}x^3+\ldots$.
$(\mu,\beta)\text{log}(1-\frac{w}{z})$ goes
into the exponential
and becomes $i_{z,w}(z-w)^{(\mu,\beta)}$.
(Check the details of this argument!)

\medskip\noindent
Consider this argument fixing conformal
weights:
$$
\Delta(a_{(n)}b)=\Delta(a)+\Delta(b)-n-1.
$$
If $a=|\mu\rangle$, $b=|\beta\rangle$, then
$$
\begin{aligned}
\Delta(|\beta\rangle)=(\beta,\beta)/2,&\qquad
\Delta(|\mu\rangle)=(\mu,\mu)/2\\
\Delta(|\mu+\beta\rangle)=(\mu+\beta,\mu+
\beta)/2.
\end{aligned}
$$

\noindent
Knowing that
$|\mu\rangle_{(n)}|\beta\rangle\in
H^{\mu+\beta}$
and
$$
\begin{aligned}
\Delta(|\mu\rangle_{(n)}|\beta\rangle)&=
\frac{(\mu,\mu)}{2}+\frac{(\beta,\beta)}{2}-
n-1\\
&=\frac{(\mu+\beta,\mu+\beta)}{2},
\end{aligned}
$$
IF
$$
(\mu,\beta)=-n-1 \iff n=-(\mu,\beta)-1.
$$
Therefore,
$$
\Gamma^\mu_{(-(\mu,\beta)-1)}|\beta\rangle=
C\cdot|\mu+\beta\rangle.
$$
$$
\,:\!(\partial_z^{(\mu,\beta)}\Gamma^\mu(z)
\Gamma^\beta(z)\!:\,
=CT^{\mu+\beta}(z)
$$
if $(\mu,\beta)\geq 0$.

One finds
$$
\begin{aligned}
\,:\!\Gamma^\mu(z)\Gamma^\beta(w)\!:\,&=e^\mu
e^\beta(\text{junk}).\\
\,:\!\Gamma^\beta(w)\Gamma^\mu(z)\!:\,&=
e^\beta
e^\mu(\text{same junk}).
\end{aligned}
$$

\noindent
$$
\Gamma^\mu(z)\Gamma^\beta(w)-(-1)^{p(\mu)
p(\beta)}\Gamma^\beta(w)\Gamma^\mu(z)
$$
is local if and only if
$$
e^\mu
e^\beta=(-1)^{p(\mu)p(\beta)}(-1)^{(\mu,
\beta)}e^\beta
e^\mu.
$$
(Here the $(-1)^{(\mu,\beta)}$ appeared
because we changed the
order of $z$ and $w$ in Equation
\ref{vertex-algebras-equation-mutually-localvmode}.)
That is,
\begin{equation}
\label{vertex-algebras-equation-quantum-relation}
C_{\mu,\beta}=(-1)^{p(\mu)p(\beta)}(-1)
^{(\mu,\beta)}C_{\beta,\mu}.
\end{equation}

In case $Q=\mathbb{Z}$, $(1,1)=1$,
we have
$$
C_{1,1}=-(-1)^{p(1)p(1)}C_{1,1},
$$
so $p(1)=1$ is forced on us! That is,
$F^{\text{ch}}=\bigoplus_{m \in
\mathbb{Z}}H^m$ can only be
a vertex \underline{superalgebra},
which we already knew (i.e. $F^\text{ch}$ is
a vertex superalgebra,
but now we know it couldn't have been
otherwise).

In fact, for all $\mu$, we get
$$
C_{\mu,\mu}=(-1)^{p(\mu)^2+(\mu,\mu)}C_{\mu,
\mu},
$$
so we \underline{must} assign a parity
$(\mu,\mu)\text{ mod }2$
to $H^\mu$.

In particular, if we take $Q$ an even
lattice,
$V$ will be an honest vertex algebra.

\begin{exercise}
\label{vertex-algebras-exercise-cocycle}
Deduce (from Borcherds identity, or
otherwise)
that $C:Q \times Q\to \mathbb{C}^\times$
must (in order for $V$ to be a vertex
superalgebra)
\begin{equation}
\label{vertex-algebras-equation-cocycle}
C(\alpha,\beta+\gamma)C(\alpha,\beta)
C(\alpha,\beta)C(\alpha+\beta,\gamma)
\end{equation}
for all $\alpha,\beta,\gamma \in Q$. (This is
a cocycle condition.)

This apparently necessary condition on $C$
was not very evident
from the locality of
$\Gamma^\mu,\Gamma^\beta,\ldots$ Discuss.
\end{exercise}

\noindent
There's a few other things we need apart from
Equation \ref{vertex-algebras-equation-cocycle}.
Namely, that $C(0,\alpha)=C(\alpha,0)=1$
(normalization condition)
and Equation \ref{vertex-algebras-equation-quantum-relation}.

Equation \ref{vertex-algebras-equation-cocycle} says $C$ is a
{\it group cocycle}.
More precisely, consider the short exact
sequence of
(not necessarily abelian) groups
$$
\xymatrix{
1\ar[r]&\mathbb{C}^\times\ar[r]&\hat{Q}\ar[r]
&Q\ar[r]&1,
}
$$
a central extension, i.e.
$\hat{Q}=\mathbb{C}^\times \times Q$
as a set, $\mathbb{C}^\times$ central, and
the product given by
$(a,\alpha)(b,\beta)=(c(\alpha,\beta)ab,
\alpha+\beta)$
which, to be associative,
$C:Q \times Q \to \mathbb{C}^\times$ must
satisfy
Equation \ref{vertex-algebras-equation-cocycle}, that is,
$C:Q \times Q \to \mathbb{C}^\times$
is a 2-cocycle
(with coefficients in the trivial $Q$-module
$\mathbb{C}^\times$).
(Also, $\underbrace{(1,0)}_{e \in
\hat{Q}}\cdot(1,\alpha)=(1,\alpha)=
(1,\alpha)(1,0)\iff
C(0,\alpha)=C(\alpha,0)=1$.)

\begin{remark}
\label{vertex-algebras-remark-isomorphism}
If we set up a linear isomorphism
$\eta:V \to V$, of the form
$\eta|_{H^\beta}=\eta(\beta) \in
\mathbb{C}^\times$ for all $\beta$,
then $\eta:Q \to \mathbb{C}^\times$ arbitrary
function,
conjugating $C$ by $\eta$, i.e.
$$
\tilde{C}(\alpha,\beta)=
\eta(\alpha,\beta)C(\alpha,\beta)\eta^{-1}
(\alpha)\eta^{-1}(\beta)
$$
i.e.,
$$
\tilde{C}(\alpha,\beta)=C(\alpha,\beta)
\frac{\eta(\alpha+\beta)}
{\eta(\alpha)\eta(\beta)}.
$$
\end{remark}

This corresponds to modifying $C$ by the
coboundary $d\eta$.
It's the same vertex superalgebra, though.

So really vertex superalgebra structures on
$V$
are classified by
$$
H^{2}(Q,\mathbb{C}^\times)=\frac{Z^2(Q,
\mathbb{C}^\times)}
{B^2(Q,\mathbb{C}^\times)}.
$$

\medskip\noindent
By the universal coefficient theorem,
if $G$ is any group and $M$ any trivial
$G$-module,
$$
\xymatrix{
0\ar[r]&\text{Ext}^1_\mathbb{Z}(H_{n-1}(G,
\mathbb{Z})\ar[r]
&H^{n}(G,M)\ar[r]
&\Hom_{\mathbb{Z}}(H_n(G,\mathbb{Z}),M)\ar[r]
&0
}
$$
In this case, $G=Q\simeq \mathbb{Z}^m$,
and $H_k(Q,\mathbb{Z}) \simeq \Lambda^kQ$.
So
$$
\xymatrix{
0\ar[r]&
\text{Ext}^1_{\mathbb{Z}}(Q,\mathbb{C}
^\times)\ar[r]&
H^{2}(Q,\mathbb{C}^\times)\ar[r]&
\Hom_{\mathbb{Z}}(\Lambda^2Q,\mathbb{C}
^\times)\ar[r]&0
}
$$
Now since $\mathbb{C}^\times$ is a divisible
group (algebra fact),
hence injective as a $\mathbb{Z}$-module,
so
$\text{Ext}^1_{\mathbb{Z}}(Q,\mathbb{C}
^\times)=0$.
(See [Weibel, 6.1.5] for more on universal
coefficients and group homology.)
So we have an isomorphism
$$
H^{2}(Q,\mathbb{C}^\times) \to
\Hom_{\mathbb{Z}}(\Lambda^2Q,\mathbb{C}
^\times).
$$

Concretely, the map is $C \mapsto  B$, where
$$
B(\alpha,\beta)=\frac{C(\alpha,\beta)}
{C(\beta,\alpha)}
$$
Recall that our $C$ was required to satisfy
$$
C(\alpha,\beta)=(-1)^{\frac{(\alpha,\alpha)}
{2}\frac{(\beta,\beta)}{2}
+(\alpha,\beta)}C(\beta,\alpha)
$$
hence the associated $B$ is
$$
B(\alpha,\beta)=(-1)^{\frac{\alpha,\alpha}{2}
+(\alpha,\beta)}
$$
In conclusion, on $V_Q=\bigoplus_{\beta \in
Q}H^\beta$,
there exists a unique conformal vertex
superalgebra structure
which extends that on $H^0$.

\medskip\noindent
$Q=\mathbb{Z}$ gave $V_Q=F^{\text{ch}}$.
$Q=\mathbb{Z}\sqrt{2}$ gives $V_Q\simeq
L_1(\mathfrak{sl}_2)$,
simple quotient of $V^1(\mathfrak{sl}_2)$,
affine $\mathfrak{sl}_2$.

Consider $\mathfrak{g}=B_2$, which was two
different lengths of weight vectors,
so as a lattice it's just $Q=\mathbb{Z}^2$.
So $V_Q=V_\mathbb{Z} \otimes
V_\mathbb{Z}=F^{\text{ch}}_2$,
the charged fermions in two variables.
Compare this with affine $B_2$, $L_1(B_2)$.
Here $B_2$ sits (horizontally)
as the charge 1 piece. So they are different!

Finally, it turns out that the vertex algebra
associated
to the Leech lattice $V_{\text{Leech}}$ is
related
to the moonshine module $V^{\text{natural}}$
by
an orbifold step.
In a paper by [Jethro-?] it was shown that in
rank 24 it's possible to pass from any vertex
algebra to $V_{\text{Leech}}$ by a single
orbifold step.


\section{Representations of vertex algebras}
\label{vertex-algebras-section-repre-verte-algeb}

\noindent
In this section we aim at describing the
irreducible representations of lattice vertex
algebras $V_Q$.
By restriction, $M$ is a $H$-module,
where $H=H^0(\mathfrak{h})$ is the Heisenberg
vertex algebra
(and also a Fock representation of the
Heisenberg Lie algebra,
which has a vertex algebra
structure, unlike other Fock representations)
associated with the commutative Lie algebra
$\mathfrak{h}=\mathbb{C}\otimes_\mathbb{Z}Q$
with the induced bilinear form
$(\cdot,\cdot):\mathfrak{h}\times\mathfrak{h}
\to\mathbb{C}$.

A few lectures ago, I stated
(ref?)

\begin{proposition}
\label{vertex-algebras-proposition-past}
$H(\mathfrak{h})$-modules are the same as
smooth $\hat{\mathfrak{h}}$-modules.
\end{proposition}

\noindent
This is a special case of a construction that
works
for any $\mathfrak{h}$ abelian Lie algebra.
Also recall

\begin{proposition}
\label{vertex-algebras-proposition-past2}
$V^k(\mathfrak{g})$-modules = $\{\text{smooth
level $k$
$\hat{\mathfrak{g}}$-modules}\}$.
\end{proposition}

\noindent
Let $\mathfrak{h}=\mathbb{C}$ now for
simplicity
(general case is similar) and let
$\mathfrak{a}=\bigoplus_{n \neq
0}\mathbb{C}h_n \oplus \mathbb{C} K \subset
\hat{\mathfrak{h}}$. (We have excluded 0!)

\begin{proposition}
\label{vertex-algebras-proposition-repre-heise}
Consider $\mathfrak{a}$-modules with
$\mathbb{Z}$-grading
$M=\bigoplus_{k \in \mathbb{Z}}M_k$ with
degrees bounded from below,
that is such that $h_n(M) \subset M_{k-n}$
and $m_k$ for $k$ sufficiently negative.
We call this a {\it positive energy module}.
(Will this work for any smooth??)
$\mathfrak{a}$-module is isomorphic to a
direct
sum of copies of the Fock space
$F=\mathbb{C}[x_1,x_2,x_3,\ldots]$
with
$$
h_n=
\begin{cases}
n\frac{\partial }{\partial x_n}\qquad &n>0 \\
x_{-n}\qquad &n>0.
\end{cases}
\qquad \qquad K=1.
$$
\end{proposition}

\begin{proof}
Let $M$ be an $\mathfrak{a}$-module.
Define in general the {\it space of vacuum
vectors} by
$$
\Omega_M=\{x \in M:h_nx=0\;\mathfrak{a}n\geq
1\}.
$$
Let $\Omega=\Omega_M$ and define
$$
\begin{aligned}
\varphi: F \otimes \Omega &\longrightarrow M
\\
p(\underline{x})\otimes v &\longmapsto
p(\underline{h})v.
\end{aligned}
$$

\noindent
We claim that $\varphi$ is injective. Indeed,
if $v \in \Ker(\varphi)\subset F \otimes
\Omega$, then
there exists $v' \in 1 \otimes \Omega \cap
\Ker$.
To see why think of $v$ as a finite sum of
polynomials
times elements in the module, take the
highest degree of the polynomials,
apply as many operators as you need to get
down to the vacuum,
this gives a sum of vacuum vectors.
(We had used this sort of argument to show
irreducibility of modules)

[Rest of proof made in class.]
\end{proof}



\medskip\noindent
Now let $M$ be a $V_Q$-module which is
positive
energy as an $\mathfrak{a}$-module.
So $M \cong F \otimes \Omega$, as above.
We have the action of quantum fields
$$
h^M(z)=\sum_{n \in \mathbb{Z}}h_{(n)}^M
z^{-n-1}
$$
on $M$, for every $h \in \mathfrak{h}$.
We also have
$$
\Gamma^M_\alpha(z)=\sum_n\Gamma^m_{\alpha(n)}
z^{-n-1}
=Y^M(e^\alpha,z)\qquad \forall \alpha\in Q.
$$

We have
$$
[h^M_{(m)},\Gamma^\alpha_M(z)=
(h,\alpha)z^m\Gamma^\alpha_M(z),
\qquad \forall m \in \mathbb{Z},h \in
\mathfrak{h}, \alpha \in Q.
$$
(We saw this relation in $V_a$ itself a few
lessons ago.)

From this we deduce
$$
\Gamma^\alpha_M(z)=
\text{exp}\left(\sum_{j<0}\frac{z^{-j}}{-j}
\alpha_j\right)
\text{exp}\left(\sum_{j>0}\frac{z^{-j}}{-j}
\alpha_j\right)
A_\alpha(z)
$$
(exactly as previously),
where the coefficients of $A_\alpha(z)$
commute with all
$h_n$ {\bf for $n \neq 0$}.
In other words, it commutes with
$\mathfrak{a}$ as defined above.

Since $T=L_{-1}$ in $V_Q$,
we have (as before) that
$$
\partial_z\Gamma^\alpha_M(z)=\,:\!\alpha(z)
\Gamma^\alpha_M(z)\!:\,
$$
which implies (as before) that
$$
A_\alpha(z)=Ce_M^\alpha z^{\alpha_0}.
$$

So
$$
\Gamma^\alpha_M(z)=e^\alpha
\text{exp}(\sum\ldots)
\text{exp}(\sum\ldots)
z^{\alpha_0},
$$
where $e_M^\alpha:\Omega_M \to \Omega_M$.

In particular $\alpha_0$ must act on $\Omega$
by integer eigenvalues
for all $\alpha \in Q$,
$\alpha_0|_{\Omega_\mu}=(\alpha,\mu)$.
So
 $$
\Omega=\bigoplus_{\mu \in Q^\vee} \Omega_\mu.
$$
Furthermore (as before),
$e_M^\alpha:\Omega_\mu \to
\Omega_{\mu+\alpha}$.
From Borcherds identity, we can deduce
\begin{equation}
\label{vertex-algebras-equation-cocycle-m}
e^\alpha_M
e^\beta_M=c(\alpha,\beta)e_M^{\alpha+\beta}
\qquad
\forall \alpha,\beta \in Q.
\end{equation}
(Compare with Equation
\ref{vertex-algebras-equation-cocycle}.)

\noindent
In particular, if $x \in \Omega_\mu$ nonzero,
then $V_Q\cdot x \subset M$ is of the form
$$
\bigoplus_{\alpha \in
Q}H^{\mu+\alpha}e^\alpha x.
$$
If $M$ is irreducible, then $M=
\bigoplus_{\alpha \in
Q}H^{\mu+\alpha}e^\alpha x$.

As before, we can picture this as a lattice
(say, triangular)
with a Fock module sitting in every point.
$V_Q$ is just
the direct sum of these modules. We can also
picture
the dual lattice $Q^\vee$, labelled $Q+\mu$.

(I think the upshot is that the weight
modules
are indexed by the quotient $Q^\vee/Q$.)

\medskip\noindent
Finally, for each $\mu \in Q^\vee/Q$,
construct
(or confirm it exists) $V_{\mu+Q}$.
Recall the cocycle $Q \times Q \to
\mathbb{C}^\times$,
$C(\alpha,\beta) \in \mathbb{C}^\times$
had to satisfy Equation
\ref{vertex-algebras-equation-quantum-relation}, namely:
 $$
C(\alpha+\beta,\gamma)C(\alpha,\beta)
=C(\alpha,\beta+\gamma)C(\beta,\gamma),
$$
etc, $\Lambda^2Q \to \mathbb{C}^\times$,
$$
\frac{C(\alpha,\beta)}{C(\beta,\alpha)}
=(-1)^{(\alpha,\alpha)(\beta,\beta)+(\alpha,
\beta)}.
$$

For an explicit construction of
$C(\alpha,\beta)$ see
\cite[Section 5]{KVAB}.

\medskip\noindent
The idea is to consider
$$
V_{Q^\vee}=\bigoplus_{\mu \in
Q^\vee}H^\mu\cdot e^\mu.
$$
Define an action
 $$
V_Q \times V_{Q^\vee} \to V_{Q^\vee}[[z^{\pm
1}]]
$$
by a similar formula,
(extend $C:Q \times Q \to \mathbb{C}^\times$
to $C:Q\times Q^\vee \to \mathbb{C}^\times$)
get a $V_Q$-module $V_{Q\vee}$.
Then
$$
V_{Q^\vee}=\bigoplus_{\mu \in
Q^\vee/Q}V_{\mu+Q}
$$
where the $V_{\mu+Q}$ are irreducible
modules.




\medskip\noindent
\noindent
Let $\mathfrak{g}$ be a Lie algebra and
$(\cdot,\cdot)$ a symmetric invariant
bilinear form.
We have the affine Lie algebra
$\hat{\mathfrak{g}}=\mathfrak{g}[t^{\pm
1}\oplus \mathbb{C}K$.
We defined a smooth
$\hat{\mathfrak{g}}$-module
and level $k$ (which means that $K \mapsto  k
\text{Id}$
$\hat{\mathfrak{g}}$-modules,
and the vertex algebra $V^k(\mathfrak{g})$.
Again recall that $V^k(\mathfrak{g})$-modules
are in correspondence with smooth level $k$
$\hat{\mathfrak{g}}$-modules.

Consider $U$ an associative algebra over
$\mathbb{C}$ with
$U_0 \supset U_1 \supset U_2 \supset\ldots$ a
decreasing
filtration by subspaces $U_p$, such that
\begin{enumerate}
\item
\label{vertex-algebras-item-separating}
(Separating.) $\bigcap_{p}U_p=0$.
[The idea here is that the $U_p$ may be
thought of
as a topological basis, and this condition
means
the algebra, as a topological algebra,
separates points:
there's an open set containing zero but not
any given point.]

\item
\label{vertex-algebras-item-l}
For any $u \in U$ and $p\geq 0$, there exists
$\ell$
wuch that  $U_\ell \cdot u \subset U U_p$.
(think: for any $p$, no matter how big,
there exists $\ell$, maybe much bigger, such
that  …)
[The idea here is that we want to turn the
$U_p$ into ideals,
just by multiplying by $U$ --- indeed, $U
U_p$ is an ideal.]

\item (Exhaustive.) $U_0= U$? Maybe not
necessary.
\label{vertex-algebras-item-exhaustive}
\end{enumerate}

\noindent
We say a left $U$-module $M$ is {\it smooth}
if
for all $x \in M$ there exists $p$ such that
$U_p\cdot x=0$.

\medskip\noindent
Given $U$ and $\{U_p\}$ as above, we can form
the {\it completion} $U^c$
as a categorical limit of the diagram

$$
\xymatrix{
U/U U_0&U/U U/1 \ar[l]& U/ U U_2 \ar[l]&
\cdots \ar[l].
}
$$
Can think of $U^c$ as
$$
\left\{\sum_is_i:\sum_is_i\text{ mod }U U_p
\text{is finite for each $p$}\right\}.
$$
Formally,
$$
U^c=\prod_{p \in \mathbb{Z}_{\geq 0}}U/U U_p
=\left\{(u_0,u_1,u_2,\ldots):\substack{u_i
\in U/U U_i \\
\text{and }\overline{u_{i+1}}=u_i}\right\}.
$$
One can then set $s_i=u_i-u_{i-1}$, then
$$
(u_0,u_1,u_2,\ldots)=(s_0,s_0+s_1,s_0+s_1+
s_2,\ldots)
\text{``=''}\sum_is_i
$$
and we recover the first description.

Clearly the action of $U$ on a smooth
$U$-module $M$
extends to a well-defined action of $U^c$.
Now let $a(w)=\sum_{n \in
\mathbb{Z}}a_{(n)}w^{-n-1}
\in U^c[[w^{\pm}]]$,
we call $a(w)$ a {\it continuous
distribution} if
for every $p \geq 0$ there exists $N \geq 0$
such that
$a_{(n)}\in U U_p$ for all $n \geq N$.
We denote the vector space of continuous
distrinutions by
$U^c[[w]]^{cont}$.

Then it's clear that
\begin{lemma}
\label{vertex-algebras-lemma-conti-distr-quant}
If $M$ is a smooth $U$-module
(denote action of $a_{(n)}$ on $M$ by
$a^M_{(n)}\in \text{End}(M)$,
and $a(w)$ a continuous distribution, then
$$
a^M(w)=\sum_n a_{(n)}^Mw^{-n-1}
$$
is a quantum field on $M$.
\end{lemma}

\begin{proposition}
\label{vertex-algebras-proposition-duplicate-2}
If $a(w),b(w)\in U^c[[w^{\pm 1}]]^{cont}$,
the $n$-th product
 $$
a(w)_{(n)}b(w)
=\text{Res}_z(a(z)b(w)i_{z,w}-b(w)a(z)i_{w,z}
)(z-w)^n
$$
is well defined and lies in $U^c[[w^{\pm
1}]]^{cont}$.
\end{proposition}

\begin{proof}
Analogous to Exercise
\ref{vertex-algebras-exercise-most-important}.
\end{proof}

\begin{definition}
\label{vertex-algebras-definition-local}
Let $a(w),b(w) \in U^c[[w^{\pm 1}]]^{cont}$.
We say the $a(w),b(w)$ is
 \begin{enumerate}
\item {\it weakly local} if there exists  $N$
such that
$a(w)_{(n)}b(w)=0$ for all $n \geq N$.
\item {\it local} if there exists $N$ such
that
$(z-w)^N[a(z),b(w)]=0$.
\end{enumerate}
\end{definition}

\noindent
Obviously local implies weakly local.

Let $\mathcal{F}\subset
U^c[[w^{\pm}]]^{cont}$ be a subset such that
all pairs of elements are local.
Let $\overline{\mathcal{F}}\subset
U^c[[w^{\pm 1}]]^{cont}$ be the
closure of $\mathcal{F}$.
 $$
\begin{aligned}
\mathcal{F}^0&=\{1_{U^c}\},\qquad
\mathcal{F}^k=\left\{a(w)_{(n)}b(w):
\substack{a
\in \mathcal{F}, \\
b \in \mathcal{F}^{k-1}, \\
n \in \mathbb{Z}}\right\}.
\end{aligned}
$$

\noindent
The analogue of Dong's lemma \ref{vertex-algebras-lemma-dong}
guarantees all
$a,b \in \overline{\mathcal{F}}$ are local.

Then we obtain a vertex algebra:

\begin{theorem}
\label{vertex-algebras-theorem-vertex-algebra-continuous}
The vector space $\overline{\mathcal{F}}$,
equipped with
$|0\rangle=1_{U^c}$,
$Y(a(w),x)=\sum_{n \in
\mathbb{Z}}a(w)_{(n)}x^{-n-1}$,
$Ta(w)=\partial_wa(w)=a(w)_{(-2)}1_{U^c}$ is
a vertex algebra.
\end{theorem}

\begin{proof}
Exercise.
\end{proof}

\medskip\noindent
Now let $M$ be a smooth $U$-module
(thought of as a $U^c$-module.)

Let's assign
$$
Y^M(a(w),x)=a^M(x)=\sum_{n \in
\mathbb{Z}}a^M_{(n)}x^{-m-1}.
$$

\begin{proposition}
\label{vertex-algebras-proposition-smooth-f-module}
$(M,Y^M)$ is an
$\overline{\mathcal{F}}$-module.
\end{proposition}

\begin{proof}
$$
Y(a(w)_{(n)}b(w),x)=Y(a(w),x)_{(n)}Y(b(w),x),
$$
because $\overline{\mathcal{F}}$ is a vertex
algebra.

So

$$
\begin{aligned}
Y^M(a(w)_{(n)}b(w),x)&=a^M(x)_{(n)}b^M(x)
\qquad
\text{by definition}\\
&=Y^M(a(w),x)_{(n)}Y^M(b(w),x).
\end{aligned}
$$
\end{proof}

\noindent
[It looks as if this construction was a
slight generalization
of our earlier work for vertex algebras over
vector spaces
to vertex algebras over associative algebras;
now we have to blend in this topological
algebra notions -
completion, etc.]

\medskip\noindent
Now  let's look at the case of
$\hat{\mathfrak{g}}$-modules and
$V^k(\mathfrak{g})$.
We have
$\hat{\mathfrak{g}}=\mathfrak{g}[t^{\pm
1}] \oplus \mathbb{C}K$.
Let
$$
U=U(\hat{\mathfrak{g}})/(K-k_\text{id}).
$$
Then $U$-modules are the same as level $k$
$\hat{\mathfrak{g}}$-modules.
Now consider
$U_p=\text{Im}(\mathfrak{g}t^{\geq p})\subset
U$.
I claim $U$ with $U_p$ so defined satisfies
the conditions
(\ref{vertex-algebras-item-separating}) and (\ref{vertex-algebras-item-l}).
Indeed, (\ref{vertex-algebras-item-separating}) is given by
PBW theorem
since every monomial has a highest degree
term.
As for (\ref{vertex-algebras-item-l}), e.g., for $u=at^{-5}$,
$p=10$, if I take
$\ell=15$, then $u_\ell \ni bt^{15}$,
$[bt^{15},at^{-5}]=[b,a]t^{10}$.

Now, in $U^c[[w^{\pm 1}]]^{cont}$ we have,
for each
$a \in \mathfrak{g}$, the continuous
distribution
$$
a(w)=\sum_{n \in \mathbb{Z}}at^n w^{-n-1}.
$$
By the commutation relation
$[at^m,bt^n]=[a,b]t^{m+n}+m\delta_{m,-n}(a,b)
k
1$
we have $(z-w)^2[a(z),b(w)]=0$
so these are mutually local,
($[a(z),b(w)]=[a,b](w)\delta(z,w)+(a,b)k 1
\partial_w \delta(z,w)$.)

So we get a vertex algebra
$\overline{\mathcal{F}}$
from $\mathcal{F}=\{a(w):a \in
\mathfrak{g}\}$
and any smooth  $\hat{\mathfrak{g}}$-module
of level $k$ gives
an $\overline{\mathcal{F}}$-module.

\medskip\noindent
In particular
$V^k(\mathfrak{g})
=U(\hat{\mathfrak{g}})
\otimes_{U(\hat{\mathfrak{g}}_+)}\mathbb{C}K$
(which happens to be a vertex algebra),
being a smooth level $k$
$\hat{\mathfrak{g}}$-module is an
$\overline{\mathcal{F}}$-module.

In $M=V^k(\mathfrak{g})$, for $a \in
\mathfrak{g}$,
$a(w)$ (from $\overline{\mathcal{F}}$) acts
via
$Y^M(a(w),x)=a(x)=T(a_{(-1)}|0\rangle,x)$.

$$
\xymatrix{
a(w)\ar@{|->}[r]&Y(a(w),z)\\
\hat{\mathcal{F}}\ar[r]\ar[rd]& \{Y(a,z):a
\in V^k(\mathfrak{g})\}
\ar[d]
& a(z)\ar@{|->}[d]\\
&V^k(\mathfrak{g})&a(z)|0\rangle|_{z=0}
}
$$
which is the state-field correspondence,

$$
\xymatrix{
a(w)_{(n)}b(w)\ar@{|->}[r]&Y(a(w),z)_{(n)}
Y(b(w),z)\\
\hat{\mathcal{F}}\ar[r]\ar[rd]_\phi&
\{Y(a,z):a \in V^k(\mathfrak{g})\}
\ar[d]
& a(z)_{(n)}b(w)\ar@{|->}[d]\\
&V^k(\mathfrak{g})&a_{(n)}b
}
$$
So $\phi$ is a homomorphism of vertex
algebras
$$
\begin{aligned}
\phi(a_{(n)}b)&\mapsto \phi(a)_{(n)}\phi(b)\\
1 & \mapsto |0\rangle.
\end{aligned}
$$

\noindent
Clearly $\phi$ is surjective. Want to show
$\Ker(\phi)=0$.

\medskip\noindent
We can characterize $V^k(\mathfrak{g})$ as
the universal
highest weight level $k$
$\hat{\mathfrak{g}}$-module,
i.e. if $M$ is a $\hat{\mathfrak{g}}$-module,
with $x \in M$ such that
\begin{enumerate}
\item $Kx = kx$,
\item $at^{\geq 0}x=0$
\end{enumerate}

\noindent
then there exists a unique
$\hat{\mathfrak{g}}$-module map
$$
\begin{aligned}
\psi: V^k(\mathfrak{g}) &\longrightarrow M \\
|0\rangle &\longmapsto x.
\end{aligned}
$$

\begin{lemma}
\label{vertex-algebras-lemma-f-ghat-module}
I claim
\begin{enumerate}
\item $\overline{\mathcal{F}}$ is a
$\hat{\mathfrak{g}}$-module,
in which $x=1_{U^c}$ satisfies
\item $Kx=k$
\item $\mathfrak{g}t^{\geq 0}x=0$.
\end{enumerate}
\end{lemma}

\begin{proof}
\begin{enumerate}
\item $\overline{\mathcal{F}}$ is a
$\hat{\mathfrak{g}}$-module.
For $c(w) \in \overline{\mathcal{F}}$,
define $\rho(at^m)c(w)=a(w)_{(m)}c(w)$.
Then
$$
\begin{aligned}
[\rho(at^m),\rho(bt^n)]&=[a(w)_{(m)},b(w)
_{(n)}\\
&=\sum_{j \geq
0}\binom{m}{j}(a(w)_{(j)}b(w))_{(m+n-j)}\\
&=(*).
\end{aligned}
$$

\noindent
But for $a,b \in \mathfrak{g}$, we have
$[a(z),b(w)]=\ldots$ see above.
So
$$
\begin{aligned}
a(w)_{(0)}b(w)&=[a,b](w)\\
a(w)_{(1)}b(w)&=k(a,b)1_{U^c}\\
a(w)_{\geq 2}b(w)&=0.
\end{aligned}
$$

\noindent
Then
$$
\begin{aligned}
(*)&=[a,b](w)_{(m+n)}+mk(a,b)1_{U^c}(w)_{(m+
n-1)}\\
&=\rho([a,b]t^{m+n}+m\delta_{m,-n}(a,b)K).
\end{aligned}
$$

\noindent
So $\rho$ is a representation of
$\hat{\mathfrak{g}}$ on
$\overline{\mathcal{F}}$.

Obviously $\rho$ is level $k$
and obviously $x=1_{U^c}(w) \in
\overline{\mathcal{F}}$
is highest weight.
$a(w)_{(j)}1(w)=0$ for all $j\geq 0$.

So we have maps
$$
\overline{\mathcal{F}}
\underbrace{\xrightarrow{\phi}}
_{\substack{\text{VA homo, hence} \\
\hat{\mathfrak{g}}\text{-mod map}}}
V^k(\mathfrak{g})
\underbrace{\xrightarrow{\psi}}_{
\hat{\mathfrak{g}}\text{-mod map}}
\overline{\mathcal{F}}.
$$
Can check $\psi \circ \phi=\text{id}$.
So $\overline{\mathcal{F}}\simeq
V^k(\mathfrak{g})$
as vertex algebras.
\end{enumerate}
\end{proof}


\section{Characters of vertex algebras and
their modules}
\label{vertex-algebras-section-characters-of-vertex-algebras}

\noindent
{\bf Disclaimer!} There may be several
computational mistakes in this section.

\medskip\noindent
Prelude. Let $q=e^{2\pi i\tau}$, $\tau \in
\mathbb{C}$,
$\text{Im}(\tau)>0$ (so  $|q|<1$).
Consider
$$
\begin{aligned}
\Delta(\tau)&=q\prod_{n=1}^\infty(1-q^n)^{24}
\\
&=\sum_{n=1}^\infty a_n q^n\\
&=q-24q^2+252q^3+\ldots
\end{aligned}
$$

\noindent
Then $\sum_{d|n}d^{11}\equiv a_n\text{ mod
}691$.

\medskip\noindent
Let $V$ be a vertex algebra,
and $L \in V$ a conformal vector.
Let $L(Z)=\sum L_nz^{-n-2}$.
Suppose furthermore that $L_0$ acts on $V$
by integer eigenvalues with finite
dimensional eigenspaces.
Write $V=\bigoplus_{n \in \mathbb{Z}}V_n$,
where $V_n=\{v \in V:L_0v=nv\}$
and write $\chi_V = \sum_{n \in
\mathbb{Z}}\dim(V_n)q^n$.

\begin{example}
\label{vertex-algebras-example-heisenberg-character}
Recall that $V=H=\mathbb{C}[x_1,x_2,\ldots]$
with conformal vector
$L=\frac{1}{2}\,:\!hh\!:\,$.
We saw that $L_01=0$, $[L_0,h_n]=nh_n$.So
$$
L_0(h_nv)=h_n(L_0v)-nh_nv.
$$
This means: take a vector with eigenvalue
$\Delta$,
then apply again, you'll get an eigenvector
with eigenvalue
$\Delta-n$. In symbols: $v \in V_\Delta$,
then
$L_0v \in V_{\Delta-n}$.

So $\Delta(x_1^{m_1}x_2^{m_2}\ldots
x_k^{m_k}1)=\sum i m_i$
and
$$
\begin{aligned}
\chi_H&=\prod_{m=1}^\infty(1-q^m)^{-1}\\
&=\sum_{N=0}^\infty p(N)q^N,
\end{aligned}
$$

\noindent
where $p(N)$ is the number of integer
partitions of $N$.
\end{example}

\begin{example}
\label{vertex-algebras-example-rank-n-heisenberg}
If $\mathfrak{h}$ is an $n$-dimensional
vector space
$(\cdot,\cdot)$ a nondegenerate symmetric
bilinear form,
$\mathfrak{h}$ abelian Lie algebra,
$\hat{\mathfrak{h}}$
affine Lie algebra. Then $H=H(\mathfrak{h})$
is the {\it rank $n$ Heisenberg algebra}.
Given a basis $\{a^1,\ldots,a^n\}$ of
$\mathfrak{h}$
and a dual basis $\{b^1,\ldots,b^n\}$
i.e. $(a^i,b^i)=\delta_{ij}$,
(so we need $(\cdot,\cdot)$ for this),
we compute Virasoro vector to be
$L=\frac{1}{2}\sum_{i=1}^n\,:\!a^ib^i\!:\,$.
Then
$$
[L_\lambda L]=TL+2\lambda
L+\frac{\lambda^3}{12}c,
$$
where the central charge $c=n$ is the rank.
Then
$$
\chi_{H(\mathfrak{h})}=\prod_{m=1}^\infty
\frac{1}{(1-q^m)^n}.
$$
\end{example}

\begin{example}
\label{vertex-algebras-example-latti-verte-algeb-chara}
Let $Q \subset \mathbb{R}^n$ be an even
integral lattice.
Consider
$\mathfrak{h}=\mathbb{R}^n
\otimes_{\mathbb{R}} \mathbb{C}=\mathbb{C}^n$
with $(\cdot,\cdot)$ inherited.
We have $V=V_Q=\bigoplus_{\alpha \in
Q}H^\alpha(\mathfrak{h})$.

As part of our construction of $V_Q$ we
computed that
$$
L_0|\alpha\rangle=\frac{(\alpha,\alpha)}{2}
|\alpha\rangle.
$$
Hence $|\alpha\rangle\in
V_{\frac{(\alpha,\alpha)}{2}}$.
(Also $[L_0,h_n]=-nh_n$
for all $h \in \mathfrak{h}$ and all $n \in
\mathbb{Z}$.)

Also $h_{-n_1}^{i_1}h_{-n_2}^{i_2}\ldots
h_{-n_s}^{i_s}|\alpha\rangle$
lies in $V_{\Delta}$,
$\Delta=\frac{(\alpha,\alpha)}{2}
+\sum_{j=1}^sn_j$. (Just as above: we apply
the operator again
and obtain another eigenvector, this time
with eigenvalue given by that sum.)
Therefore
$$
\chi_{V_Q}=\sum_{\alpha \in
Q}q^{\frac{(\alpha,\alpha)}{2}}
\prod_{m=1}^\infty \frac{1}{(1-q^m)^n}.
$$
This makes sense, since (for now) $Q$ is
assumed to be positive-definite.

In fact
$$
\chi_{V_Q}=\prod_{m=1}^\infty\frac{1}{(1-q^m)
^n}\cdot
\underbrace{\sum_{\alpha \in
Q}q^{\frac{(\alpha,\alpha)}{2}}}_{
\substack{\text{theta} \\ \text{function}}}.
$$
The series on RHS is called the {\it theta
function}
of the lattice $Q$, denoted $\theta_Q(\tau)$.
(Think of it as a function that knows how
many vectors of each
norm the lattice has.)
More precisely, for $q=e^{2\pi i \tau}$,
$$
\theta_Q(\tau)=\sum_{\alpha \in
Q}q^{\frac{(\alpha,\alpha)}{2}}
=\sum_{\alpha \in Q}e^{\pi i
(\alpha,\alpha)\tau}.
$$

As detour, notice that if $Q$ is integral,
but not even,
then some $\frac{(\alpha,\alpha)}{2} \not \in
\mathbb{Z}$,
but are in $\frac{1}{2}\mathbb{Z}$.
As an example of this, $Q=\mathbb{Z} \subset
\mathbb{R}$ is
an odd lattice,
$$
\theta_Q(\tau)=
\sum_{n \in \mathbb{Z}}q^{\frac{n^2}{2}}
=1+2q^{\frac{1}{2}}+2q^2+2q^{\frac{a}{2}}+
\ldots.
$$
\end{example}

\begin{remark}
\label{vertex-algebras-remark-theta-function-converges}
Note that for any positive definite $Q$ (even
or odd)
the series $\theta_Q(\tau)$ converges for all
$\text{Im}(\tau)>0$.
(Easy application of some test or other.)

If $Q$ is even, then only integer powers of
$q$
appear, and we have
$$
\theta_Q(\tau+1)=\theta_Q(\tau)\qquad \forall
\tau\in \mathbb{H}
$$
where $\mathbb{H}=\{\tau \in
\mathbb{C}:\text{Im}(\tau)>0\}$.

Can then also think of $\theta_Q$ as a
holomorphic
function (of $q$) on the disk
$D=\{q:|q|<1\}$.

If $Q$ is odd, then
$\theta_Q(\tau+2)=\theta_Q(\tau)$.
\end{remark}

\medskip\noindent
Many lectures ago
(see Section \ref{vertex-algebras-section-theta-functions}),
I proved for $Q=\mathbb{Z}$
that $\theta_Q\left(-\frac{1}{\tau}\right)
=\sqrt{\frac{\tau}{i}}\theta_Q(\tau)$
using the Poisson summation formula.

What about for other $Q$?

\begin{proposition}
\label{vertex-algebras-proposition-theta-funct-self-dual}
If $Q$ is self-dual of rank $n$ (and even
(?)),
then
$$
\theta_Q\left(-\frac{1}{\tau}\right)=
(\tau/i)^n\theta_Q(\tau).
$$
\end{proposition}

\begin{proof}
This is also an application of Poisson
summation.
\end{proof}

\noindent
More generally,

\begin{proposition}
\label{vertex-algebras-proposition-duplicate-3}
If $Q$ is integral with dual $Q^\vee$,
then
$$
\theta_Q\left(-\frac{1}{\tau}\right)
=(\tau/i)^n \sum_{\mu \in
Q^\vee/Q}s_{0,\mu}\cdot \theta_{\mu+Q}(\tau)
$$
where the $s_{0,\mu} \in \mathbb{C}^\times$
are nonzero complex numbers
and
$\tau_{\mu+Q}(\tau)=\sum_{\beta\in\mu+Q}
q^{(\beta,\beta),2}$.
\end{proposition}

\noindent
As opposed to the usual case, this function
has very
fractional powers of $q$. However, it is
well-defined as
a function on $\mathbb{H}$.

\begin{proposition}
\label{vertex-algebras-proposition-duplicate-4}
If $Q$ is integral, then for all $\lambda \in
Q^\vee/Q$,
$\theta_{\lambda+Q}\left(-\frac{1}{\tau}
\right)
=(\tau/i)^n \sum_{\mu \in
Q^\vee/Q}s_{\lambda\mu}\cdot\theta_{\mu+Q}
(\tau)$.
\end{proposition}

\noindent
Again proved by Poisson summation.
In the process an explicit formula for
$\{S_{\lambda\mu}\}$ emerges.
I can't remember, but it's something like $$
S_{\lambda\mu}=\frac{1}{|Q^\vee/Q|^{1/2}}
e^{2\pi
i(\lambda,\mu)}
\qquad (?)
$$
Thinking of $Q^\vee/Q$ as a finite list
$\{1,2,\ldots,N\}$, $N=|Q^\vee/Q|$, then
$\{S_{\lambda\mu}\}_{\lambda,\mu=1,\ldots,n}$
becomes a matrix,
the {\it $S$-matrix} of $Q$.

Since
$-1\frac{1}{\left(-\frac{1}{\tau}\right)}=
\tau$,
we (naively) should have $S^2=I$ identity
matrix.

There is a slight problem:
$\theta_{\mu+Q}(\tau)$
might not all be distinct.

\begin{example}
\label{vertex-algebras-example-a2-theta-function}
Consider $Q=A_2$, the triangular lattice.
We know that $Q^\vee/Q\simeq\mathbb{Z}/3$.
Actually, $Q^\vee/Q=\{Q,\mu+Q,2\mu+Q\}$.
In fact, two of the theta-functions
associated to these
three cosets are actually the same:
$$
\begin{aligned}
\theta_Q&=1+6q+\ldots\\
Q_{\mu+2}&=3q^{\frac{2}{3}}+\ldots\\
Q_{2\mu+Q}&=3q^{-\frac{2}{3}}+\ldots.
\end{aligned}
$$

\noindent
We can refine the $\theta_{\mu+Q}(\tau)$
to make them linearly independent.

Let $z \in \mathfrak{h}$, and define
$$
\theta_{\mu+Q}(\tau,z)=\sum_{\beta \in
\mu+Q}q^{\frac{(\beta,\beta)}{2}}
e^{2\pi i(z,\beta)}.
$$
Now
$$
\theta_{\mu+Q}:\mathbb{H}\times\mathfrak{h}
\to\mathbb{C}
$$
(convergent and holomorphic).
\end{example}

\begin{proposition}
\label{vertex-algebras-proposition-duplicate-5}
For all $\lambda \in Q^\vee/Q$,
$$
\theta_{\lambda+Q}\left(-\frac{1}{\tau},
\frac{z}{\tau}\right)
=(\tau/i)^ne^{\frac{\pi i(z,z)}{\tau}}
\sum_{\mu \in
Q^\vee/Q}S_{\lambda,\mu}\theta_\mu(\tau,z)
$$
\end{proposition}

\begin{proof}
Also by Poisson summation, same
$s_{\lambda,\mu}$.
But now $\{\theta_{\mu+Q}(\tau,z)|\mu \in
Q^\vee/Q\}$
are linearly independent. So $S$ makes sense
(as a matrix).
So it no longer holds that $S^2=I$.
We apply $S$-transformations as follows
$$
\begin{aligned}
(\tau,z)&\overset{s}{\mapsto
}\left(-\frac{1}{\tau},\frac{z}{\tau}\right)
\mapsto
\left(-\frac{1}{-\frac{1}{\tau}},\frac{z/
\tau}{-1/\tau}\right)
=(\tau,-z)
\overset{s}{\mapsto
}\left(-\frac{1}{\tau},-\frac{z}{\tau}\right)
\overset{s}{\mapsto }(\tau,z).
\end{aligned}
$$

\noindent
So $S^4=I$.
\end{proof}

\noindent
Recall that
$$
\chi_{V_Q}(\tau)=\frac{1}{\prod_{m=1}
^\infty(1-q^m)^n}\theta_Q(\tau).
$$

Recall the definition of the Dedekind $\eta$
function
(Theorem
\ref{vertex-algebras-theorem-dedekind-eta-function}).
Why do we need to put the factor $q^{1/24}$?
The reason is that without it, it is not a
modular form,
whereas with it, it is.

\begin{lemma}
\label{vertex-algebras-lemma-duplicate-2}
$$
\eta\left(-\frac{1}{\tau}\right)
=\sqrt{(\tau/i)}\eta(\tau)
$$
(that's why we normalised by $q^{1/24}$).
\end{lemma}

\noindent
So
$$
\prod_{m=1}^\infty (1-q^m)^n
=q^{-n/24}\eta(\tau)^n
$$
$$
\chi_{V_Q}(\tau)=\frac{\theta_Q(\tau)}
{\eta(\tau)^n}q^{n/24}.
$$

Let's correct the definition of
$\chi_{V_Q}(\tau)$.
Note that $n=\text{rank}(Q)$ is the central
charge $c$ of $V_Q$.

\begin{definition}
\label{vertex-algebras-definition-chara-confo-verte-algeb}
The {\it character} of a conformal vertex
algebra $V$ is
$$
\chi_V(\tau)=q^{-\frac{c}{24}}\sum_{n \in
\mathbb{Z}}\dim(V_n)q^n.
$$
If $M$ is a $V$-module,
we also define
$\chi_M(\tau)=q^{-\frac{c}{24}}\sum_{n\in!!}
\dim(M_n)q^n$,
where $!!$ are the eigenvalues on $M$ of
$L_0$ (maybe not in $\mathbb{Z}$).

We write
$$
M_\Delta=\{m \in M:L_0m=\Delta m\}.
$$
A more uniform notation is
$$
\begin{aligned}
\chi_V(\tau)&=\text{Tr}_Vq^{L_0-c/24},\\
\chi_M(\tau)&=\text{Tr}_Mq^{L_0-c/24}.
\end{aligned}
$$
\end{definition}

\begin{lemma}
\label{vertex-algebras-lemma-character-v-modules}
For $\mu \in Q^\vee/Q$, we have the
$V_Q$-module $V_\mu$.
Then
$$
\text{Tr}_{V_\mu}q^{L_0-c/24}=
\frac{\theta_\mu(\tau)}{\eta(\tau)^n}
$$
and
$$
\boxed{\chi_M(\tau,z)=\text{Tr}_Me^{2\pi
iz_0}q^{L_0-c/24}.}
$$
Indeed in
$V_\lambda=\bigoplus_{\beta\in\lambda+Q}
H^\beta(\mathfrak{h})$,
$z \in \mathfrak{h}$. In
$H^\beta(\mathfrak{h})$,
$z_0|_{H^\beta}=(\beta,z)I_{H^\beta}$.
So $e^{2\pi i(z,\beta)}=e^{2\pi i z_0}$.
$$
\chi_{V_\lambda}\left(0\frac{1}{\tau},
\frac{z}{\tau}\right)
=e^{\frac{\pi i(z,z)}{\tau}}
\sum_{\mu \in
Q^\vee/Q}S_{\lambda,\mu}\chi_\mu(\tau,z).
$$
\end{lemma}




\medskip\noindent
(New lecture.) Let's do a summary of what we
have done
in this section so far.
For $Q \subset \mathbb{R}^n$ an integral
lattice,
we defined a function
$$
\theta_Q:\underbrace{\mathbb{H}}_{\{\text{Im}
>0\}}x
\mathfrak{h} \to \mathbb{C}
$$
by
$$
\theta_Q(\tau,z)=\sum_{\alpha \in
Q}q^{\frac{(\alpha,\alpha)}{2}}
e^{2\pi i(\alpha,z)},
$$
where $q=e^{2\pi i \tau}$, $z \in
\mathfrak{h}=\mathbb{C}^n$.
We explained that this series converges since
$|q|<1$ and it
is raised to a quadratic term.

More generally, if
$[\mu] \in Q^\vee/Q$, then
$$
\theta_\mu(\tau,z)=\sum_{\alpha \in
\mu+Q}q^{\frac{(\alpha,\alpha)}{2}}
e^{2\pi i (\alpha,z)}
$$
similarly.

\medskip\noindent
If $Q$ is self-dual,
$$
\theta_Q\left(-\frac{1}{\tau},\frac{z}{\tau}
\right)
=(\tau/i)^{n/2}e^{\frac{\pi i (z,z)}{\tau}}
\theta_Q(\tau,z)
$$
by the Poisson summation formula.
More generally, for $Q$ non self-dual,
$$
\theta_\lambda\left(-\frac{1}{\tau},\frac{z}
{\tau}\right)
=(\tau/i)^{n/2}e^{\frac{\pi(z,z)}{\tau}}
\sum_{\mu \in
Q^\vee/Q}S_{\lambda\mu}\theta_\mu(\tau,z),
$$
where $S=(S_{\lambda\mu})$ is some matrix,
$S^4=I$.

\begin{remark}
\label{vertex-algebras-remark-t-matrix}
There is also the $T$-matrix, corresponding
to
$(\tau,z)\mapsto (\tau+1,z)$.
\begin{align*}
S&\xymatrix{\ar@{<->}[r]&}
\begin{pmatrix}
0&-1\\
1&0
\end{pmatrix}\in \text{SL}_2\mathbb{Z}\\
T & \xymatrix{\ar@{<->}[r]&}
\begin{pmatrix}
1&1\\
0&1
\end{pmatrix}\in \text{SL}_2(\mathbb{Z}).
\end{align*}

\noindent
In fact, these generate
$\text{SL}_2(\mathbb{Z})$.
\end{remark}

$$
\begin{aligned}
T&:\underbrace{e^{2\pi i \tau}}_{q}\to
e^{2\pi i(\tau+1)}
=e^{2\pi i \tau}=q\\
T&:\underbrace{e^{2\pi i\tau a}}_{q^a}\to
e^{2\pi i a \tau+2\pi i a}=q^ae^{2\pi i a}.
\end{aligned}
$$

\medskip\noindent
In $\theta_{\mu+Q}(\tau,z)$,
we have powers
$q^{\frac{(\alpha,\alpha)}{2}}$ such that
$\alpha \in \mu+Q$.
Note that if $\mu,\mu' \in \mu+Q$, $\mu'-\mu
\in Q$,
$$
\frac{(\mu+\alpha,\mu+\alpha)}{2}
=\frac{(\mu,\mu)}{2}
+\underbrace{\frac{2(\mu,\alpha)}{2}}_{\in
\mathbb{Z}}
+\underbrace{(\alpha,\alpha)}_{\in
\mathbb{Z}}
$$
since $Q$ is even.

So $T$ acts on $\theta_\mu(\tau,z)$
by a constant. Indeed, by the constant
$e^{2\pi i (\mu,\mu)/2} \in
\mathbb{C}^\times$.

So we have a $T$-matric, which is diagonal,
with these numbers
$$
\{e^{2\pi i (\mu,\mu)/2}|\mu \in Q^\vee/Q\}
$$
on the diagonal.

In $\text{SL}_2(\mathbb{Z})$, $S^4=I$,
$(ST)^3=I$.
These relations hold for our matrices $S$ and
$T$.

\begin{proposition}
\label{vertex-algebras-proposition-latti-sl2-repre}
An even integral lattice $Q$
gives a representation of
$\text{SL}_2(\mathbb{Z})$
of dimension $|Q^\vee/Q|$.
\end{proposition}

\noindent
Also earlier we noted that
$$
\chi_{\mu+Q}(\tau,z)
=\frac{\theta_{\mu+Q}(\tau,z)}{\eta(\tau)^n}
$$
is the graded dimension of the $V_Q$-module
$V_{\mu+Q}$.
More precisely,
$V=V_Q$ acts on $M=V_{\mu+Q}$.
$z \in \mathfrak{h}\subset H\subset V_Q$,
so $z_0$ acts on $V_{\mu+Q}$,
also $L_0$ acts on $V_{\mu+Q}.$
And
$$
\chi_{\mu+Q}(\tau,z)=\text{Tr}_{V_{\mu+Q}}
e^{2\pi
i z_0}
q^{L_0-c/24}
$$
where the central charge is
$c=n=\text{rank}(Q)$.

\medskip\noindent
Recall (compare with Definition
\ref{vertex-algebras-definition-modular-form}) that

\begin{definition}
\label{vertex-algebras-definition-modular-form2}
A {\it modular form} of weight $k$ is a
holomorphic function
$f(\tau)$ of $\tau \in \mathbb{H}$ such that
\begin{enumerate}
\item
$$
f\left(\frac{a\tau+b}{c\tau+d}\right)
=(c\tau+d)^kf(\tau)
$$
for all
$\begin{pmatrix}
a&b\\
c&d
\end{pmatrix}\in \text{SL}_2(\mathbb{Z})$.

\item The Fourier series expansion of $f$
(which exists since it must be periodic):
 $$
f(\tau)=\sum_{n=0}^\infty a_nq^n
$$
has no $a_{<0}$ terms.
\end{enumerate}
\end{definition}

\medskip\noindent
Now we explain a construction:
let $\Lambda=\{m+n\tau|m,n \in
\mathbb{Z}\}\subset\mathbb{C}$
a ``lattice spanned by 1 and $\tau$''.
(It's not a lattice in the sense of Section
\ref{vertex-algebras-section-lattices}.)
Consider
$$
\sum_{\lambda \in \Lambda\setminus 0}
\frac{1}{\lambda^m}=G_m(\Lambda).
$$
\begin{enumerate}
\item $G_m(\Lambda)$ can be thought of as
$G_m(\tau)$.
\item The infinite series converges if $m
\geq 3$.
Inside the disk $D_{2^N}$, there are
on the order of $(2^N)^2$ points.
These points have
$$
\left|\frac{1}{\lambda^m}\right|<\frac{1}
{(2^{N-1})^m}.
$$
$$
\sum\left|\frac{1}{\lambda^m}\right|
=\sum_N\sum_{N/2<|\lambda|<N}
\frac{1}{\lambda^m}
<2^m \sum_N\frac{2^{2N}}{2^{m N}}
$$
converges for $m \geq 3$.
\end{enumerate}

Think of $G_m(\tau)$ ($m\geq 3$).
Consider
$\begin{pmatrix}
a&b\\
c&d
\end{pmatrix}$,
$a,b,c,d$ integers,
and  $a\tau+b \in \Lambda$, $c\tau+d \in
\Lambda$.
E.g. $\begin{pmatrix}
2&5\\
1&3
\end{pmatrix}$.

We claim that if $\det \begin{pmatrix}
a&b\\
c&d
\end{pmatrix}=1$,
then
$$
\{m(a\tau+b)+n(c\tau+d)|m,n\in
\mathbb{Z}\}=\Lambda.
$$
So, denoting $\sum'=\sum_{(m,n) \in
\mathbb{Z}^2\setminus(0,0)}$,
$$
\begin{aligned}
G_k(\tau)&=\sum'\frac{1}{(m\tau+n)^k}\\
&=\sum'\frac{1}{(m(a\tau+b)+n(c\tau+d))^k}\\
&=(c\tau+d)^{-k}\sum'\frac{1}{\left(
m\frac{a\tau+b}{c\tau+d}+n\right)^k}\\
&=(c\tau+d)^{-k}G_k\left(\frac{a\tau+b}
{c\tau+d}\right).
\end{aligned}
$$

\noindent
So
$G_k\left(\frac{a\tau+b}{c\tau+d}\right)
G_k(\tau)$
for all $\begin{pmatrix}
a&b\\
c&d
\end{pmatrix}\in \text{SL}_2(\mathbb{Z})$.
This is called an {\it Eiseistein series} of
weight $k$.

\begin{remark}
\label{vertex-algebras-remark-firs-eisestein-series}
If $k$ odd, then $\frac{1}{\lambda^m}$ and
$\frac{1}{(-\lambda)^m}$
cancel, and $G_k(\tau)=0$, nontrivial.
So $G_4(\tau)$ is the first Eisenstein
series.
\end{remark}

\medskip\noindent
Now let's talk about the Fourier series of
$G_{2k}(\tau)$.
Start with
\begin{equation}
\label{vertex-algebras-equation-euler}
(1-x^2)(1-\frac{x^2}{4})(1-\frac{x^2}{9}
\ldots=\frac{\sin(\pi
x)}{x},
\end{equation}

\noindent
which may be proved by applying Liouville's
theorem
to the quotient LHS/RHS.

In fact, this was Euler's method for proving
the identity
$\sum_{n=1}^\infty
\frac{1}{n^2}=\frac{\pi^2}{6}$.

\medskip\noindent
Do log derivative of Equation
\ref{vertex-algebras-equation-euler}.
$$
\pi \frac{\sin(\pi x)}{\cos(\pi x)}
=\frac{d}{dx}\text{log}\left(\prod_{n=1}
^\infty
(1-\frac{x^2}{n^2})\right)=\ldots
$$

\begin{theorem}
\label{vertex-algebras-theorem-g2k}
$$
\begin{aligned}
G_{2k}(\tau)&=(-1)^k\frac{\pi^{2k}}{(2k)!}
B_{2k}\cdot
E_{2k}(\tau),\\
E_{2k}(\tau)&=1-\frac{4k}{B_{2k}}\sum_{m=1}
^\infty
\frac{m^{2k-1}q^m}{1-q^m}.
\end{aligned}
$$
\end{theorem}

\noindent
Here $B_{2k}$ are Bernoulli numbers,
$$
\sum_{n=0}^\infty B_n\frac{t^n}{n!}
=\frac{t}{e^t-1}.
$$
The first of them are
 $$
\begin{matrix}
2k:& 0 & 2 & 4 & 6 & 8 & 10 & 12 & 14\\
B_{2k}:& 1 & 1/6 & 1/30 & 1/42 & -1/30 & 5/66
& -691/2730 & 7/6
\end{matrix}
$$
So
$$
E_4(\tau)=1-240 \sum_{m=1}^\infty \frac{m^3
q^m}{1-q^m},
$$
$$
G_4(\tau)=\frac{\pi^4}{45}E_4(\tau).
$$

\begin{definition}
\label{vertex-algebras-definition-space-holom-modul-forms}
Let $\mathcal{M}^{\text{hol}}_{2k}$ be the
vector space
of holomorphic modular forms of weight $k$.
\end{definition}

\begin{proposition}
\label{vertex-algebras-proposition-dimension-mhol2k}
$\dim \mathcal{M}^{\text{hol}}_{2k}$ is
 $$
\begin{array}{c|cccccccccc}
2k:& 0 & 2 & 4 & 6 & 8 & 10 & 12 & 14 & 16 &
18\\
\dim\mathcal{M}^{\text{hol}}_{2k}:
& 0 & 0 & 1 & 1 & 1 & 1 & 2 & 1 & 2 & 2
\end{array}
$$
and
$\dim\mathcal{M}^{\text{hol}}_{2k+12}=
\dim\mathcal{M}^{\text{hol}}_{2k}+1$.
\end{proposition}

\begin{proof}[Proof sketch]
Let $f(\tau) \in
\mathcal{M}^{\text{hol}_{2k}}$
and consider
$$
\xi(\tau)=\frac{d}{d\tau}\text{log}f(\tau)=
\frac{f'(\tau)}{f(\tau)}.
$$
Since
$$
\begin{aligned}
f(-1/\tau)&=\tau^{2k}f(\tau),\\
\xi(-1/\tau)=\frac{2k}{\tau}+\xi(\tau),
\end{aligned}
$$

\noindent
apply Cauchy-formula-type argument
$\int_{\partial \mathcal{D}}\xi(\tau)d\tau
=$ numer of zeroes of $f -$ number of poles
of  $f$ in $\mathcal{D}$.

Then integrate along the boundary of the
region above the unit circle,
between the vertical lines $x=-1/2,1/2$ and
bounded above by
some horizontal segment at height $N$.
In conclusion, the vertical integrals cancel
each other
and $\frac{k}{12}\geq N(i\infty)$,
where $N(i\infty)=$ order of the zero of $f$
at $i\infty$.
\end{proof}

\noindent
One has, if $f \in
\mathcal{M}_{2k}^{\text{hol}}$,
$g \in \mathcal{M}_{2\ell}^{\text{hol}}$,
$fg \in
\mathcal{M}^{\text{hol}}_{2(k+\ell)}$.
So $E^3_4 \in \mathcal{M}^{\text{hol}}_{12}$
and
$E_{12}\in \mathcal{M}_{12}^{\text{hol}}$.

Then do
$$
\begin{aligned}
E_{12}-E^3_4&=\frac{65520}{691}(q-24q+252q-
\ldots)\\
&=\frac{65520}{691}q\prod_{n=1}^\infty(1-q^N)
^{24}.
\end{aligned}
$$

\begin{remark}
\label{vertex-algebras-remark-dedekind}
$$
\begin{aligned}
\Delta(\tau)&=q\prod_{n=1}^\infty(1-q^n)^{24}
\\
&=\eta(\tau)^{24}.
\end{aligned}
$$

\noindent
In fact,
$\bigoplus_{k=0}^\infty\mathcal{M}_{2k}
^{\text{hol}}
\simeq\mathbb{C}[E_4,E_{12}]$.
\end{remark}

\medskip\noindent
Next time. Let $Q$ be unimodular of rank 24.
So
$$
P(\tau,z)=\frac{\theta_Q(\tau,z)}{\eta(\tau)
^{24}}
$$
satisfies
$$
P\left(-\frac{1}{\tau},\frac{z}{\tau}\right)
=e^{\frac{\pi i (z,z)}{\tau}}P(\tau,z).
$$

\section{A non-modular form}
\label{vertex-algebras-section-a-non-modular-form}

\noindent
$$
\begin{aligned}
G_2(\tau)&=\frac{\pi^2}{3}E_2(\tau)\\
E_2(\tau)&=1-24\sum_{n=0}^\infty\frac{n
q^n}{1-q^n}.
\end{aligned}
$$

\noindent
The series $\sum_{\lambda \in
\Lambda\setminus 0}\frac{1}{\lambda^2}$
does not converge absolutely, but converges
conditionally.

We could define, i.e.,
$$
G_2(\tau)=\sum_{ n \neq 0}\frac{1}{n^2}
+\sum_{m \neq 0}\sum_{n \in \mathbb{Z}}
\frac{1}{(m\tau+n)^2}.
$$
$$
[\text{Picture of dot arrangement in plane}]
$$
A modular transformation on $\Lambda$
sends the rows here into some diagonal lines
$$
[\text{Picture of diagonal dor arrangement in
plane}]
$$
and so the argument that proves modularity of
$G_{2k}(\tau)$, $k\geq 2$, in this case shows
$$
G_2\left(\frac{a\tau+b}{c\tau+d}\right)
=(c\tau+d)^2[\underbrace{\leavevmode}_{
\substack{\sum_{\lambda \in \Lambda \setminus
0}\frac{1}{\lambda^2}, \\
 \text{but summed in }\\
\text{different order}}}].
$$
\begin{proposition}
\label{vertex-algebras-proposition-e2}
$$
E_2(-1/\tau)=\tau^2E_2(\tau)+\frac{12\tau}
{2\pi
i}.
$$
\end{proposition}

\begin{remark}
\label{vertex-algebras-remark-bizzare-fact}
This relation can be used to prove that
$\eta(\tau)=q^{1/24}\prod_{n=1}^\infty(1-q^n)
$
is modular of weight $1/2$.
($\eta(-1/\tau)=(\tau/i)^{1/2}\eta(\tau)$.)
Idea: consider
$\frac{\eta'(\tau)}{\eta(\tau)}$
and relate it to $E_2(\tau)$.
\end{remark}

\noindent
$$
\bigoplus_{k=0}^\infty\mathcal{M}^{\text{hol}
}_{2k}
\cong\mathbb{C}[G_4,G_6].
$$

\medskip\noindent
Recall we were working with an even self-dual
lattice $Q$
of rank (say) 24. [This case is actually
simpler thank
ranks 8 or 16 since we don't get fractional
powers of $q$.]
We had
$$
\theta_Q(\tau,z)=\sum_{\alpha \in
Q}q^{\frac{(\alpha,\alpha)}{2}}
e^{2\pi i(\alpha,z)}.
$$
$$
\theta_Q\left(-\frac{1}{\tau},\frac{z}{\tau}
\right)
=e^{\frac{\pi i
(z,z)}{\tau}}(\tau/i)^{n/2}\theta_Q(\tau,z)
$$
Let's consider
$$
P_Q(\tau,z)=e^{\frac{-2\pi
i(z,z)}{24}E_2(\tau)}
\theta_Q(\tau,z).
$$

We can check that the non-modularity of
$E_2(\tau)$
``cancels'' the $e^{\frac{\pi i
(z,z)}{\tau}}$ factos,
and we get
$$
P_Q\left(-\frac{1}{\tau},\frac{z}{\tau}
\right)=(\tau/i)^{n/2}P_Q(\tau,z).
$$

\begin{proof}
$$
\begin{aligned}
e^{-2\pi i
\frac{\left(\frac{z}{\tau},\frac{z}{\tau}
\right)}{24}
E_2\left(-\frac{1}{\tau}\right)}
&=e^{2\pi i \frac{(z,z)}{24\tau^2}
\left(\tau^2E_2(\tau)+\frac{12\tau}{2\pi
i}\right)}\\
&=e^{-2\pi i
\frac{(z,z)}{24}E_2(\tau)-\frac{(z,z)}{2\tau}
},
\end{aligned}
$$

\noindent
which is close enough.
\end{proof}

\medskip\noindent
Let's Taylor series expand $P_Q$ (in
$\mathbb{Z}$).
$$
\begin{aligned}
\theta_Q(\tau,z)
&=\sum_{\alpha \in
Q}q^{\frac{(\alpha,\alpha)}{2}}
\left(1+2\pi i(\alpha,z)
+\frac{(2\pi i)^2}{2}(\alpha,z)^2
+\frac{(2\pi
i)^2}{6}(\alpha,z)^3+\ldots\right)\\
&=\sum_{\alpha \in
Q}q^{\frac{(\alpha,\alpha)}{2}}
+2\pi i\underbrace{\sum_{\alpha\in
Q}q^{\frac{(\alpha,\alpha)}{2}}(\alpha,z)
}_{\substack{=0\text{ since} \\
\alpha,-\alpha\text{ cancel }\forall \alpha
\in
Q}}
+\frac{(2\pi i)^2}{2}\sum_{\alpha \in
Q}q^{\frac{(\alpha,\alpha)}{2}}
(\alpha,z)^2+\ldots
\end{aligned}
$$

\noindent
Meanwhile,
$$
e^{-2\pi i \frac{(z,z)}{24}E_2(\tau)}
=1-\frac{2\pi i}{24}(z,z)E_2(\tau)
+\frac{(2\pi
i)^2}{24^2}(z,z)^2E_2(\tau)^2+\ldots
$$
Multiply to get
$$
P_Q(\tau,z)=\theta_Q(\tau)+
\left(\theta_Q(\tau)\frac{-2\pi
i}{24}(z,z)E_2(\tau)
+\frac{(2\pi i)^2}{2}
\sum_{\alpha \in
Q}q^{\frac{(\alpha,\alpha)}{2}}(\alpha,z)
^2\right)+\ldots
$$
Expanding further the $z^2$ term:
$$
\begin{aligned}
-\frac{2\pi i}{24}\underbrace{(1-24q+\ldots)
}_{E_2(\tau)}
\sum_{\alpha \in
Q}q^{\frac{(\alpha,\alpha)}{2}}(z,z)
+\frac{(2\pi i)^2}{2}\sum
q^{\frac{(\alpha,\alpha)}{2}}(\alpha,z)^2
\end{aligned}
$$

\noindent
write the $q$-series expansion:
gives $(2\pi i)^2$ times
\begin{equation}
\label{vertex-algebras-equation-e14}\begin{aligned}
&-\frac{(z,z)}{24}\left(1-24q
+\sum_{(\alpha,\alpha)=2}q+O(q^2)\right)\\
&+\frac{1}{2}\left(
\sum_{(\alpha,\alpha)=2}q(\alpha,z)^2+O(q^2)
\right)\\
&=-\frac{(z,z)}{24}
+q\left(\frac{1}{2}\sum_{(\alpha,\alpha)=2}
(\alpha,z)^2
-\frac{(z,z)}{24}(N-24)\right)
+O(q^2)
\end{aligned}
\end{equation}

\noindent
where $N=\#\{\alpha|(\alpha,\alpha)=2\}$.

\medskip\noindent
What good is this?
Since
$P_Q\left(-\frac{1}{\tau},\frac{z}{\tau}
\right)
=\tau^{12}P_Q(\tau,z)$ (for $n=24$).
We have ($S$ and $T$-invariance, which
generate
the whole modular group)
$$
\begin{aligned}
P_Q\left(\frac{a\tau+b}{c\tau+d},\frac{z}
{c\tau+d}\right)
&=(c\tau+d)^{12}P_Q(\tau,z).\\
P_Q(\tau,z)&=
\sum_{i=0}^\infty P_Q^i(\tau)z^i,
\end{aligned}
$$

\noindent
one finds
$$
P_Q^i\left(\frac{a\tau+b}{c\tau+d}\right)
=(c\tau+d)^{12+i}P_Q^i(\tau).
$$
In particular the quadratic (in $z$)
term lies in $\mathcal{M}^{\text{hol}}_{14}$,
which is 1-dimensional, spanned by
$E_{14}(\tau)=1-24q+\ldots$.

For us, this means
that $\text{coef}(q)=-24\text{coef}(1)$ in
the last equality of
Equation \ref{vertex-algebras-equation-e14}.
I.e.,
\begin{equation}
\label{vertex-algebras-equation-coefficients}
\frac{1}{2}\sum_{(\alpha,\alpha)=2}(\alpha,z)
^2
-\frac{(z,z)}{24}(N-24)=(z,z).
\end{equation}
$$
\frac{1}{2}\sum_{(\alpha,\alpha)=2}(\alpha,z)
^2=(z,z)\frac{N}{24}.
$$
Recall that
$$
\Delta=\{\alpha\in Q|(\alpha,\alpha)=2\}
$$
is a simply laced root system.
$\Delta$ is either $A_n$, $D_n$, $E_6$, $E_7$
or $E_8$
(or an orthogonal sum of such components).
A simple root system $\Delta$ has
an associated invariant called the dual
Coxeter number
$h^\vee$. By definition, if $r=r_\alpha$ is
the reflection in $\alpha \in \Delta$,
then $(r_\alpha x,r_\alpha y)=(x,y)$ for all
$x,y \in \mathbb{R}^n$ ambient space.
In particular, $(\cdot,\cdot)$ is invariant
for the action of the Weyl group
$W=\left\langle{}r_\alpha|\alpha \in \Delta\right\rangle{}$,
or rather, $W \subset O(\mathbb{R}^n)$.
The following lemma claims that this form is
unique.

\begin{lemma}
\label{vertex-algebras-lemma-uniqueness-of-bilinear-form}
If
$\left\langle{}\cdot,\cdot\right\rangle{}:\mathbb{R}
^n\times\mathbb{R}^n\to\mathbb{R}$
invariant under $W$, then
$\left\langle{}\cdot,\cdot\right\rangle{}=\alpha(\cdot,\cdot)
$
for some scalar $\alpha$.
\end{lemma}

\noindent
OK, so consider
$$
\left\langle{}x,y\right\rangle{}=\frac{1}{2}
\sum_{\alpha \in \Delta}(\alpha,x)(\alpha,y).
$$
This is $W$-invaraint, and so is some
multiple of $(\cdot,\cdot)$.
The definition of $h^\vee$ is that
$$
\left\langle{}x,y\right\rangle{}=h^\vee(x,y).
$$
In particular,
$$
\frac{1}{2}\sum_{\alpha \in
\Delta}(\alpha,z)^2=h^\vee(z,z),
$$
which appears in Equation
\ref{vertex-algebras-equation-coefficients}.

\begin{example}
\label{vertex-algebras-example-coxeter-numbers}
For $A_n$, $h^\vee=n+1$, for $E_7$, 18,
for $D_n$, $2n-2$.
\end{example}

\noindent
So, suppose
$$
Q \supset \Delta=\{\alpha|(\alpha,\alpha)=2\}
=\Delta_1\cup \Delta_2\cup \ldots\cup
\Delta_k
$$
where $\Delta_i$ is a simple root system.

Let $\beta\in\Delta_i$, our formula implies
$$
\begin{aligned}
h_i^\vee(\beta,\beta)&=(\beta,\beta)\frac{N}
{24}\\
\implies h_i^\vee&=\frac{N}{24}.
\end{aligned}
$$

\noindent
that is,
$$
N=\sum_{i=1}^k|\Delta_i|,
$$
and
$$
h_i^\vee=\frac{N}{24}
$$
for each component $\Delta_i$.

\begin{example}
\label{vertex-algebras-example-simple-root-systems}
If $\Delta=\Delta_i$ is simple and type
$A_n$,
say $A_{24}$. We have $h^\vee=25$,
$$
\begin{aligned}
\dim
\mathfrak{sl}_{25}&=25^2-1=(25+1)(25-1)\\
N&=25^2-1-24=(24)(26)-24=25(24).
\end{aligned}
$$
\end{example}

\begin{remark}
\label{vertex-algebras-remark-other-solutions}
\begin{itemize}
\item  The empty set is vacuously a solution.
\item $E_8^3$: $|\Delta_{E_8}|=240$,
$h^\vee=30$,
$h^\vee=\frac{3\times 240}{24}=30$, is a
solution.
\item $A_2^{24}$, $h^\vee=2$ for all
components,
$\frac{|\Delta|}{24}=\frac{24(2)}{24}=2$.
\end{itemize}
\end{remark}

\noindent
Consider $\mathfrak{g}$ a semisimple Lie
algebra with root system $\Delta$,
$\mathfrak{g}=\mathfrak{g}
_1\oplus\mathfrak{g}_2
\oplus \ldots
\oplus \mathfrak{g}_k$,
$\Delta(\mathfrak{g}_i)=\Delta_i$.

Note that if $\Delta\neq \emptyset$, then
$\text{rank}(\Delta)=24$.
So we can rewrite $\#\Delta$ as $\dim
\mathfrak{g}-24$.
so
\begin{equation}
\label{vertex-algebras-equation-hi}
h_i^\vee=\frac{\dim(\mathfrak{g})}{24}-1.
\end{equation}

\noindent
There exist at most ``partitions of 24 into 3
colours
A, D and E'' solutions to Equation
\ref{vertex-algebras-equation-hi}.

Observe: $\ell:$ rank, $d:$ dimension,
$h^\vee:$ $h^\vee$.
For ADE, $h^\vee=O(\ell)$, $d=O(\ell^2)$.
Use this to bound $\dim(\mathfrak{g})$.

Turns out there are 23 solutions (24 if we
include the
empty root system)
$A_{24},E_8^3,A_1^{24},A_2^{12}$, etc.

\begin{theorem}
\label{vertex-algebras-theorem-solutions}
Each nonempty solution $\Delta$ is the root
system
of exactly one self-dual even lattice.
\end{theorem}

\begin{proof}
Construct $Q_\Delta$ by orbifolding from
$E_8^3$.
Compute $|\text{Aut}(Q_\Delta)|$,
 $$
RHS(\text{mass formula})-\sum_{\Delta \neq
\emptyset}
\frac{1}{|\text{Aut}(Q_\Delta)|}=\frac{1}
{10^{18}}.
$$
\end{proof}

\medskip\noindent
The subalgebra $\mathfrak{g}_i\subset
\mathfrak{g}$
induces a homomorphism of vertex algebras
$$
L_1(\mathfrak{g}_i) \to V_Q.
$$
Note that $V_Q$ has one irreducible module,
namely $V_Q$.
Let $V=V_0 \oplus V_1 \oplus V_2 \oplus
\ldots$
be a vertex algebra with this parity,
$$
\begin{aligned}
V_0&=\mathbb{C}|0\rangle\\
V_1&=\mathfrak{g}\text{ some Lie algebra}
\end{aligned}
$$

\noindent
It turns out that
$\mathfrak{g}=\mathfrak{g}_1 \oplus
\mathfrak{g}_2
\oplus \ldots \oplus \mathfrak{g}_k$
induces a set of embeddings
$$
L_{k_i}(\mathfrak{g}_i) \to V,\qquad  k_i \in
\mathbb{Z}_{\geq 1}.
$$
One can prove that
$$
\frac{h_i^\vee}{k_i}
=\frac{\dim\mathfrak{g}}{24}-1
$$
now. If you allow $k_i\geq 2$,
there are no longer 23 solutions
but 221 solutions.
71 vertex algebra. This is called
``Schellekens list''.




% --- End vertex-algebras.tex ---

% --- Begin springer.tex ---
\chapter{Springer theory}
\label{chapter-springer}




\section{Introduction}
\label{springer-section-intro}

\noindent
Sections \ref{springer-section-flag-variety}
to \ref{springer-section-affine-weyl-group}
(except the $\mathfrak{sl}_n$ exercises
section)
are AI-transcribed from a pdf written
by Jethro.

Summary of Section \ref{springer-section-flag-variety}:
\begin{enumerate}
\item (Setup.) $\mathfrak{g}$ simple Lie
algebra, $G$ connected algebraic group
whose Lie algebra is $\mathfrak{g}$.
The mental image should be $G=\text{SL}_n$
and $\mathfrak{g}=\mathfrak{sl}_n$,
and maybe even $n=2$.
In fact, for $n=2$ the Borels
are those fixing a line, so
$\mathcal{B}=\mathbb{P}^1$,
and for general $n$ a Borel consists
of the endomorphisms of $\mathbb{C}^n$
which map each $F_i$ into itself.

\item ($\mathcal{B}$.) 
This is the set of
Borel subalgebras. It is a smooth
variety with an action of $G$.
In fact, the action is transitive,
so $\mathcal{B} = G/\operatorname{Stab}$,
and it turns out that
$\operatorname{Stab} = B$ where
$B$ is the Lie group whose Lie algebra
is $\mathfrak{b}$.

\item (Nilpotent cone $\mathcal{N}$.)
The set of nilpotent elements of $\mathfrak{g}$
is a subvariety $\mathcal{N} \subset
\mathfrak{g}$. The adjoint action of $G$
preserves $\mathcal{N}$ and splits it
into a finite disjoint union of
nilpotent orbits we denote by $\mathbb{O}$.

\item (Springer fibers.) They are the fibers
of the projection 
$\tilde{\mathcal{B}} \to \mathfrak{g}$ in
$$
\xymatrix{
&  \widetilde{\mathcal{B}}
=\tilde{\mathfrak{g}}\ar[dl]^{\mu}\ar[dr]\\
\mathfrak{g}&&\mathcal{B}.
}
$$
We write the springer fibers as
$\mathcal{B}_x$ for $x \in \mathfrak{g}$.

\item (Kinds of Springer fibers.)
\begin{enumerate}
\item (Regular semisimple.) In this case
the Springer fiber is finite of size
$|W|$ and has an action
of the Weyl group $W$ of $G$.

\item (Regular but not semisimple.)
These elements are nilpotent
and lie in a single nilpotent
orbit called the {\it principal orbit}.
The Springer fibers are a single point.

\item (Non-regular nilpotent elements.)
The Springer fiber can be higher dimensional,
like the fiber above zero.
In general the Weyl group does 
not act on these fibers.
\end{enumerate}

\item (Nilpotent Springer resolution.)
$\mu^{-1}(\mathcal{N}):=\widetilde{\mathcal{N}}$,
contained in $\widetilde{\mathcal{B}}$,
is a resolution of $\mathcal{N}$.
Further, $\widetilde{\mathcal{N}}
=T^*\mathcal{B}$.

\item (Steinberg variety.)
$\widetilde{\mathcal{N}}\times_{\mathcal{N}}
\widetilde{\mathcal{N}}$
is the Steinberg variety,
and here we consider correspondences.
$Z = \coprod_{N}
T^*_Y(\mathcal{B} \times \mathcal{B})$,
a union of conormal bundles.

\item (Not-Springer correspondence.)
Using a slightly confusing construction
involving the definition of the Weyl group,
we obtain a map $\Phi$ from $W$ to
the set of nilpotent orbits, which is not
the Springer correspondence since the latter
identifies irreducible representations of $W$
with nilpotent orbits.

\item (Springer fibers and nilpotent orbits.)
The Springer fiber $\mathcal{B}_x$
carries an action of the stabilizer group
$G(x)$ and in fact $\mathbb{O}=G/G(x)$,
and $\tilde{\mathbb{O}} \simeq 
G \times_{G(x)}\mathcal{B}_x$.
\end{enumerate}

\medskip\noindent
Here's the roadmap to understand the
toy case of $\text{SL}_2$:

\begin{enumerate}
\item Set
$G=\operatorname{SL}_2$ and
$\mathfrak g=\mathfrak{sl}_2$.

\item Identify the Weyl group:
$W=S_2$.

\item Identify the flag variety:
$\mathcal B=\mathbb P^1$.

\item Interpret points of $\mathcal B$
as lines in $\mathbb C^2$.

\item Interpret Springer fibres
as invariant lines.

\item Work out the regular semisimple
Springer fibre.

\item See the $S_2$-action on that fibre.

\item Work out the regular nilpotent
Springer fibre.

\item Work out the zero Springer fibre.

\item Compare the three fibres.

\item Understand why
$\widetilde{\mathcal N}=T^*\mathbb P^1$.

\item Relate this example to the general
Springer picture.
\end{enumerate}

\medskip\noindent
Summary of Section \ref{springer-section-weyl-action}:
\begin{enumerate}
\item After recalling Borel-Moore homology,
``to get representations of the Weyl group
from the Steinberg variety,
the hing to do is look at
$H^{BM}_{\operatorname{top}}(Z)
=H^{BM}_{4|\Delta^+|}(Z)$''.

\item Recall that $H^{BM}_{\operatorname{top}}$
of a variety is generated by the fundamental
classes of the irreducible components,
$[V_1],…,[V_k]$. ``It is proven
somewhere that the $|W|$ distinct
conormal bundles that make up the
Steinberg variety $Z$ are exactly
its irreducible components, so
$H^{BM}_{\operatorname{top}}(Z)
\cong \bigoplus_{w \in W}\mathbb{C}[w]$''.
\end{enumerate}


\section{Lecture February 2026}
\label{springer-section-lecture-feb}

\noindent
All varieties and algebras defined over the
complex numbers.

\medskip\noindent
Simple Lie algebra $\mathfrak{g}$. Nilpotent
cone $\mathcal{N} \subset
\mathfrak{g}$,
$\mathcal{N}=\{ x \in \mathfrak{g}:
\text{ad}(x)\text{ nilpotent}\}$
[we've already seen this object].

$G$ (connected complex) Lie group with Lie
algebra $\mathfrak{g}$. There is the adjoint
action of $G$ on $\mathfrak{g}$. (If $G
\subset \text{Mat}_{n\times n}(\mathbb{C})$,
and $\mathfrak{g}= \text{Lie} G\subset
\text{Mat}_{n \times n}(\mathbb{C})$,
$G$ acts on $\mathfrak{g}$ by conjugation,
$\text{ad}(g)\cdot X:=g X g^{-1}$.)

$$
\mathfrak{g}=\coprod (\text{infinitely many
adjoint orbits}),
$$
but
\begin{equation}
\label{springer-s-label-001}
\mathcal{N}=\coprod(\text{finitely many
adjoint orbits}).
\end{equation}

\noindent
Equation
\ref{springer-s-label-001}
is proved via a sequence of results:
\begin{enumerate}
\item Jacobian-Morozov theorem
(each nilpotent $e \in \mathfrak{g}$ is
contained in an $\mathfrak{sl}_2$-triple
$\{e,h,f\}$),

\item ``Conjugacy theorems of (a) Kostant
and (b) Malcer''.

\item Classification of the possible
orbits of $h$ in the $\mathfrak{sl}_2$-triple
$\xymatrix{\ar@{~>}[r]&}$ wieghed Dynkin
diagrams.
\end{enumerate}

\medskip\noindent
Springer correspondence:
$$
\text{Irreps of $W$}
\leftrightarrow
\text{Nilpotent orbits}.
$$
\begin{example}
\label{springer-example-sln-nilpotent}
$\mathfrak{g}=\mathfrak{sl}_n$, then nilp $e
\in \mathfrak{sl}_n$. Put $e$ in Jordan
canonical form
$$
\left[\substack{\text{Matrix decomposed in
Jordan blocks,}\\
\text{each block with all zero but 1's over
the diagonal}\\
\text{- because it's nilpotent it looks like
that}}\right]
$$
That is, a partition of $n$. $[n]$ correspond
to orbit called principal nilpotent orbit (or
``regular nilpotent orbit''). In particular
$[1^n]$ corresponds to $e=0$.

So \# of nilpotent orbits = $p(n) $ =
\# integer partitions of $n$.

For $\mathfrak{sl}_n$, Weyl group is $W
\simeq S_n$. A conjugacy class in $S_n$ is
given by ``cycle shape'' determined by
lengths of cycles. [Example.]

If $G$ finite group, we know that
\# conjugacy classes = \# irreps.

For $S_n$, this also = $p(n)$. [In general
it's not clear how to assign conjugacy
classes to irreps; $S_n$ is an exception]
\end{example}

\noindent
For $\mathfrak{g}=E_8$, $\# W \simeq 700
\times 10^6$,
\# conjugacy classes in  $W=112$.
\# nilpotent orbits in $\mathcal{N}(E_8)=70$.

\medskip\noindent
Key objects. Start with $\mathfrak{g}$. A
Borel subalgebra $\mathfrak{b}\subset
\mathfrak{g}$
is a maximal (wrt inclusion) solvable
subalgebra. Denote by $\mathcal{B}$ the set
of all Borels in $\mathfrak{g}$.

\noindent
Facts. By the adjoint action of $G$, all
Borels are conjugate. ($G$ acts transitively
in $\mathcal{B}$.)

This $\implies \mathcal{B}
\overset{\text{sets}}{\simeq}G/
Stab(\mathfrak{b}_0)
\subset G$
the Borel subgroup corresp.

For $\mathfrak{sl}_n$, can take
$$
\mathfrak{b}+0=\left\{
\begin{pmatrix} \ast & \ldots{} & \\
& \ast & \ast \\
& & \ddots \\
0 & & \ldots{} & \ast\end{pmatrix} \right\}
$$
(upper triangular matrices).

And $B_0=$ same but in $\text{SL}_n$
invertible, det=1, upper triangular

\noindent
Aside. If $\mathfrak{g}$ has rank $\ell$,
root system  $\Delta$,
$\mathfrak{g}=\mathfrak{n}_-\oplus
\mathfrak{h}\oplus \mathfrak{n}_+$
and $\mathfrak{b}=\mathfrak{h} \oplus
\mathfrak{n}_+$
is an example of a Borel algebra. Dimension
is $m=\ell+\frac{1}{2}|\Delta|$.

\medskip\noindent
Back to main point. Since all Borels are
conjugate, they all have the same dimension.
$\mathcal{B}\subset
\text{Gr}(\mathfrak{g},m)$,
the Grassmannian. $\mathcal{B}\subset
\text{Gr}(\mathfrak{g},m)$ is an algebraic
subvariety [which may not be smooth]. Since
$\mathcal{B}$ has a transitive $G$-action, a
local homeo with $\mathbb{C}^m$ near any
point $\implies$ local homeos everywhere, so
$\mathcal{B}$ is smooth. This $\mathcal{B}$ is
called the flag variety $\mathcal{B} \simeq
G/B$.

Over $\text{Gr}(V,m)$ in general one has the
{\it tautological bundle}, i.e. over the
$m$-dimensional vector space $U \subset V$,
over the point $[U] \in \text{Gr}(V,m)$ one
`puts' $U$ itself.

\medskip\noindent
We do the same for $\mathcal{B}$. Define
$\tilde{\mathfrak{g}}=
\{(\mathfrak{b},x):\mathfrak{b} \in
\mathcal{B}, x \in \mathfrak{b}\}$.
Then $(\mathfrak{b},x) \mapsto  \mathfrak{b}$
defines $\tilde{\mathfrak{g}}\to \mathcal{B}$
and $\tilde{\mathfrak{g}}$ is a sort of
tautological vector bundle of $B$.

Can also map
$$
\begin{aligned}\tilde{\mathfrak{g}}
&\longrightarrow \mathfrak{g} \\
(\mathfrak{b},x) &\longmapsto x
\end{aligned}
$$
$$
\xymatrix{
&\tilde{\mathfrak{g}}\ar[dl]_{\mu}\ar[dr]
^{\pi}\\
\mathfrak{g}&&\mathcal{B}.
}
$$

\medskip\noindent
\{It looks like the projection $\mu$
is more interesting than $\pi$: fix an
element of the Lie algebra and look at all
the Borels that contain this element; indeed,
0 is contained in every Borel and thus we get
the ``zero section'' on top of 0.\}

\begin{definition}
\label{springer-definition-springer-fibre}
Let $x \in \mathfrak{g}$, the subset
$\mu^{-1}(x) \subset
\tilde{\mathfrak{g}}$
is called a {\it Springer fibre}.
\end{definition}

\begin{example}
\label{springer-example-flag-variety}
$\mu^{-1}(0)=\mathcal{B}$ [since all borels
contain zero, this is kind of a zero section
of the bundle, i.e. a copy of $\mathcal{B}$],
the flag variety.
\end{example}
Other extreme: Let $x \in \mathfrak{g}$ be a
regular semisimple element. Then
$$
|\mu^{-1}(x)|=|W|,
$$
finite set in bijection w/ Weyl group.

$$
\left[
\substack{\text{Picture of plane containing
the nilpotent cone $\mathcal{N}$} \\
\text{with $\mathcal{B}$ sitting over
origin,}\\
\text{a fiber of 4 points on top a point
labelled ss reg}\\
\text{and other varieties on top of other
points in the base plane.}}\right]
$$
$h \in \mathfrak{g}$ regular semisimple, then
$h$ sits in a unique Cartan $\mathfrak{h}$,
then $\mathfrak{h}$ sits inside $|W|$
disjoint Borels.

\begin{remark}
\label{springer-remark-pincipal-orbit}
For $x \in \mathbb{O}_{\text{prin}}\subset
\mathcal{N}$
(principal orbit), $\mu^{-1}(x)=(\text{one
point})$. But for $x$ in smaller orbits,
$\mu^{-1}(x)$ becomes a higher dimensional
variety. In fact
\begin{equation}
\label{springer-equation-important}
\dim\mu^{-1}(x)=\frac{1}{2}(\dim
\mathcal{N}-\dim \mathbb{O})
\end{equation}
for $x \in \mathbb{O}$.
\end{remark}

\medskip\noindent
Consider $\mu|_{\mu^{-1}(\mathcal{N})}$ i.e.
$$
\tilde{\mathcal{N}}=\{(\mathfrak{b},x):
\mathfrak{b}
\in \mathcal{B},
x \in \mathfrak{b}\text{ nilpotent}\}.
$$
[Now we consider a diagram as before but
restricted (and abuse notation denoting the
maps with the same letters:]
$$
\xymatrix{
&
\tilde{\mathcal{N}}\ar[dr]^{\pi}\ar[dl]_{\mu}
\\
\mathcal{N} & & \mathcal{B}.
}
$$
If $\mathfrak{g}=\mathfrak{n}_-+\mathfrak{h}+
\mathfrak{n}_-$,
$x \in \mathfrak{b}$ nilpotent $\iff x \in
\mathfrak{n}
=[\mathfrak{b},\mathfrak{b}]$,
$\mathfrak{b}=\mathfrak{h}+\mathfrak{n}_-$,
$[\mathfrak{b},\mathfrak{b}]=\mathfrak{n}_-$.

Notice that
$\tilde{\mathcal{N}}\xrightarrow{\pi}
\mathcal{B}$
is also a vector bundle [fibres are vector
spaces of the same dimension]. Now fibre over
$\mathfrak{b}$ is $\mathfrak{n}\subset
\mathfrak{b}$.

Claim.
$$
\xymatrix{
\tilde{\mathcal{N}}\ar[d]\\\mathcal{B}
}
$$
is isomorphic to the cotangent bundle.
$\tilde{\mathcal{N}}=T^*_{\mathcal{B}}$
(bundle of 1-forms).
$$
\dim \mathcal{B}
=\dim G - \dim B
=\dim \mathfrak{g}- \dim \mathfrak{b}.
$$
[It's not hard to see this bundle is indeed
the cotangent bundle, for instance we check
that dimension must be
$$
\begin{aligned}
\dim \mathfrak{g}+\dim \mathfrak{b}
&=\dim \mathfrak{n}_-\\
&=\dim \mathfrak{n}_+\\
&=\text{rk}(\tilde{\mathcal{N}})
\end{aligned}
$$
since $\mathfrak{g}
=\mathfrak{n}_-+\underbrace{\mathfrak{h}+
\mathfrak{n}_+}_{\mathfrak{b}}$.]

[Kaledon's theory - when we pushforward a
constant sheaf on a resolution of varieties
where the upstairs variety has symplectic
structure. See perverse sheaves, Equation
\ref{springer-equation-important}.]

\medskip\noindent
Let $h$ be regular semisimple, $h \in
\mathfrak{h}$ (we know that $\mathfrak{h}$ is
unique). How many $\mathfrak{b}$ Borels
contain $\mathfrak{h}$? Recall $N_G(T)/T
\simeq W$ [$T$ torus].

\medskip\noindent
Let $\mathfrak{b}_1,\mathfrak{b}_2$ be two
Borels, $h \in
\mathfrak{b}_1,\mathfrak{b}_2$. There exists
$g$ such that
$g\mathfrak{b}g^{-1}=\mathfrak{b}$.

\begin{remark}
\label{springer-remark-flag}
$G$ acts on
$\mathcal{B}\xrightarrow{\simeq}G/B$
$$
\xymatrix{
\mathfrak{b}\ar@{|->}[d]\ar[r]&
g_0B\ar@{|->}[d]\\
g\mathfrak{b}g^{-1}\ar[r]&g g_0B.
}
$$
Consider the restricted action of $B$ on
$G/B$ and its orbits. $B\setminus G/B$
``double cosets''.
\end{remark}

\noindent
Fact. There exist finitely many double
cosets, in fact $|W|$. Indeed
$$
\xymatrix@R=0.1em{
W\ar@{-}[r]
&NG(T)/T\ar@{-}[r]
&  B\setminus G / B \ar@{-}[r]
& \mathcal{B}\times \mathcal{B}/\Delta G\\
&g \ar@{|->}[r]& g\\
& & BgB \ar@{|->}[r]
&(gB,eB)\text{ or something}
}
$$
Recall that $T \subset B \subset G$.
[Concerning the rightmost group: the action
if $g(\mathfrak{b}_1,\mathfrak{b}_2)
=(g\mathfrak{b}_1g_1,g\mathfrak{b}_2g^{-1})
$.]

[$W$ can be simply thought of as the group
generated by the reflections determined by
the roots; $N_G(T)/T$ is similar, recall that
$T$ is the\ldots{}]

\medskip\noindent
[We shall see one last variety, the Steinberg
variety.]

Recall
$$
\xymatrix{
\tilde{\mathcal{N}}=
\{(\mathfrak{b},x):x \in \mathcal{N},
\mathfrak{b} \ni x\}
\ar[d]^{\mu}\\
\mathcal{N}.
}
$$
Define [a fibered product]
$$
\xymatrix{
Z=\tilde{\mathcal{N}}\times_{\mathcal{N}}
\tilde{\mathcal{N}}
=- \{(\mathfrak{b}_1,\mathfrak{b}_2,x):
x \in \mathcal{N},
\mathfrak{b}_1,\mathfrak{b}_2 \ni x\}
\ar[d]^{\mu^{(2)}}\\
\mathcal{N}.
}
$$
If $\mathbb{O}\subset \mathcal{N}$ is a
nilpotent orbit, $x \in \mathbb{
\mathcal{O}}$,
$\dim \mu^{-1}(x)=\frac{1}{2}(\dim
\mathcal{N} - \dim \mathbb{ O})$.

$Z$ is called the {\it Steinberg variety}.

$$
\xymatrix{
\mu^{-1}(x)\ar[d]
^{\substack{\text{fibration}
\\ \mu}}\\
\mathbb{O}
}
$$

\medskip\noindent
Let $Y_W \subset \mathcal{B} \times
\mathcal{B}$,
the $\Delta G$-orbit correspond to $w \in W$.
Define $Z_W = T^*_{Y_W}(\mathcal{B} \times
\mathcal{B})$
the  ``conormal bundle'' of $Y_n \subset
\mathcal{B} \times \mathcal{B}$
by, $C_p=\{ \varphi \in T^*_pM:
\varphi(T_p N)=0\}$

[We have two ways of partitioning $Z$, one is
by in which orbit an elements lies, i.e.
using
$$
\xymatrix{
Z\ar[d]^{\mu^{(2)}}\\
\mathbb{O},
}
$$
(observe that dimension = $\dim
\mathbb{O}+2\cdot \frac{1}{2}
(\dim \mathcal{N}-\dim \mathbb{ O}=\dim
\mathcal{N}$.


and another using parametrized by the Weyl
group as in Equation
\ref{springer-equation-partition-w}

\medskip\noindent
Claim.
$$
Z \cong \coprod_{w \in W}Z_w
$$


\section{Summary of lecture 1}
\label{springer-section-summary-1}

\begin{itemize}
\item Lie group $G$ with lie algebra
$\mathfrak{g}$. The set of Borel subalgebras
of $\mathfrak{g}$ is a smooth algebraic
variety $\mathcal{B}$.

\item Define $\tilde{\mathfrak{g}}=
\{(\mathfrak{b},x):\mathfrak{b} \in
\mathcal{B}, x \in \mathfrak{b}\}$.
Then $(\mathfrak{b},x) \mapsto  \mathfrak{b}$
defines $\tilde{\mathfrak{g}}\to \mathcal{B}$
and $\tilde{\mathfrak{g}}$ is a sort of
tautological vector bundle of $B$.

Can also map
$$
\begin{aligned}\tilde{\mathfrak{g}}
&\longrightarrow \mathfrak{g} \\
(\mathfrak{b},x) &\longmapsto x
\end{aligned}
$$
$$
\xymatrix{
&\tilde{\mathfrak{g}}\ar[dl]_{\mu}\ar[dr]
^{\pi}\\
\mathfrak{g}&&\mathcal{B}.
}
$$

\medskip\noindent
\{It looks like the projection $\mu$
is more interesting than $\pi$: fix an
element of the Lie algebra and look at all
the Borels that contain this element; indeed,
0 is contained in every Borel and thus we get
the ``zero section'' on top of 0.\}

\item Restrict those maps to
$$
\tilde{\mathcal{N}}=\{(\mathfrak{b},x):
\mathfrak{b}
\in \mathcal{B},
x \in \mathfrak{b}\text{ nilpotent}\}.
$$

\item $\tilde{\mathcal{N}}$ is the cotangent
bundle of $\mathcal{B}$!
\end{itemize}

\section{Meeting June 12}
\label{springer-section-june-12}

\noindent
\begin{itemize}
\item $\widetilde{\mathcal{B}} \to \mathcal{B}$
is the tautological bundle,
mapping $\mathfrak{b} \mapsto  \mathfrak{b}$
because $\mathfrak{b}$ is a vector space
and a point of $\mathcal{B}$.

\item There's also $\widetilde{\mathcal{N}}
\hookrightarrow \widetilde{\mathcal{B}}$,
which turns out to be the resolution
of the nilpotent cone
$\mathcal{N} \subset\mathfrak{g}$,
which is a singular variety.
How does it work? Ah, because
there's the other projection
$\widetilde{\mathcal{B}} \to \mathfrak{g}$
because the fibers of $\mathcal{B}$
are the Borels, each of which is a subset
of $\mathfrak{g}$, so that map makes sense.

\item (Conormal story.)
Maybe $\widetilde{\mathcal{N}} \simeq
T^*\mathcal{B}$, and maybe this is
the conormal bundle
(the kernel of the cotangent bundle of
the ambient variety)
of… $\mathcal{N}$ I think, i.e. in the
photo we see $T^*  \mathcal{B}
\xrightarrow{\pi}\mathcal{N}$.

\item (Perverse sheaf.) There is a constant
sheaf which pushed downstairs becomes a
perverse sheaf.
$\pi_*\mathbb{Q}_{T^*\mathcal{B}}$.
\end{itemize}
\section{The flag variety and its friends
and cousins}
\label{springer-section-flag-variety}

\noindent
Let $\mathfrak g$ be a simple Lie algebra and
let $G$ be a connected algebraic
group whose Lie algebra is $\mathfrak g$. For
example if
$\mathfrak g = \mathfrak{sl}_n$ we can take
$G = \operatorname{SL}_n$ which
happens to be simply connected. We might also
take $G$ to be the `adjoint
group', which in this case is the quotient of
$\operatorname{SL}_n$ by its
centre (isomorphic to $\mathbb Z/n$) which
consists of $n$th roots of unity
multiples of the identity matrix.

\medskip\noindent
Let $\mathcal B$ be the set of all Borel
subalgebras (Boreis) in
$\mathfrak g$. All Boreis have the same
dimension
$m = |\Delta^+| = (\dim(\mathfrak g) -
\operatorname{Rank}(\mathfrak g))/2$,
so $\mathcal B$ is contained in the
Grassmannian variety
$\operatorname{Gr}(\mathfrak g, m)$ of
$m$-planes in the vector space
$\mathfrak g$ (which is a projective
variety). The condition on an $m$-plane
of being a Borel is a collection of
polynomial equations, so $\mathcal B$ is a
closed subvariety. The group $G$ acts on
$\mathcal B$, which quickly leads to
a proof that $\mathcal B$ is a smooth
variety. It is a non trivial fact that
$G$ acts transitively on $\mathcal B$, so
$\mathcal B$ is connected,
furthermore $\mathcal B \cong
G/\operatorname{Stab}_G(\mathfrak b)$ where
$\mathfrak b$ is any fixed Borel and
$\operatorname{Stab}_G(\mathfrak b)$
denotes the stabiliser in $G$ of $\mathfrak
b$.

\medskip\noindent
From the definition of a Borel, the
stabiliser group of a Borel $\mathfrak b$
coincides with the Lie subgroup $B \subset G$
with Lie algebra $\mathfrak b$.
So $\mathcal B \cong G/B$.

\medskip\noindent
We denote by $\widetilde{\mathfrak g}$ the
total space of the tautological
bundle of $\mathcal B$, i.e., the set of
pairs $(x, \mathfrak b)$ where
$x \in \mathfrak b$. Why do we denote it
$\widetilde{\mathfrak g}$? Well,
because we want to think of
$\widetilde{\mathfrak g}$ as `lying over'
$\mathfrak g$ via the projection
$\mu : \widetilde{\mathfrak g} \to \mathfrak
g$ defined by
$(x, \mathfrak b) \mapsto x$. This map has
some interesting geometry. For an
element $x \in \mathfrak g$ we denote by
$\mathcal B_x$ its preimage in
$\widetilde{\mathfrak g}$. These preimages
are called Springer fibres. Each
Springer fibre $\mathcal B_x$ is a subset of
$\mathcal B$ cut out by the
closed condition $x \in \mathfrak b$, thus
the Springer fibres are compact.
Note that they need not be smooth since they
do not inherit $G$-actions from
$\mathcal B$.

\medskip\noindent
We recall the definition of an element $x \in
\mathfrak g$ being `regular
semisimple'. The characteristic polynomial
$$
\det(tI - \operatorname{ad}(x)) = t^N +
c_{N-1}(x)t^{N-1} + \cdots + c_1(x)t + c_0(x)
$$
is divisible by $t^\ell$ where $\ell$ is the
rank of $\mathfrak g$ (this is
essentially the definition of the rank), and
then $c_\ell(x)$ is some nonzero
polynomial on $x \in \mathfrak g$ which is
zero in some but not all points.
The points where this polynomial is nonzero
are the `regular semisimple'
elements. Remark: In Kac's notes, such
elements are simply called `regular',
it is a fact that they are semisimple. In 
\cite{ChrissGinzburg}
(see Definition 3.1.3 there) the
word regular is used in a different way, so
be careful. In \cite{ChrissGinzburg} an element $x$
is said to be regular if the centraliser
subalgebra $Z_{\mathfrak g}(x)$ has
dimension equal to $\ell$. Rather weirdly it
turns out that the set of regular
elements is the union of the set of the set
of regular semisimple elements,
and the principal nilpotent orbit.

\medskip\noindent
If $x$ is regular semisimple then the
Springer fibre $\mathcal B_x$ consists
of a finite set of points, with cardinality
equal to $|W|$ (the cardinality of
the Weyl group). This is [1, Proposition
3.1.36]. The geometric description of
$W$ is used in a fundamental way: namely $W =
N_G(T)/T$. Furthermore each of
these finite Springer fibres carries a simply
transitive action of $W$.

\medskip\noindent
In the case of $\mathfrak g =
\mathfrak{sl}_n$ there is quite a nice way to
understand this result. Recall that $\mathcal
B$ is naturally isomorphic in
this case to the `flag variety' consisting of
flags
$$
0 = F_0 \subset F_1 \subset F_2 \subset
\ldots \subset F_n = \mathbb C^n.
$$
(Such a flag $F$ gives a Borel $\mathfrak b$
consisting of those endomorphisms
of $\mathbb C^n$ that map each $F_i$ into
itself). Now let
$x \in \mathfrak{sl}_n$ be regular
semisimple: Its action on $\mathbb C^n$ is
by a collection of pairwise distinct
eigenvalues, so it specifies an
(unordered) collection of lines in $\mathbb
C^n$. If we choose an order for
these eigenspaces, and then set $F_i$ to be
the sum of the first $i$ of them,
then we get a flag whose corresponding Borel
contains $x$ as a diagonal matrix
(with its eigenvalues arranged in the chosen
order). It turns out that the
$n!$ choices of an order give the $n! =
|S_n|$ points of $\mathcal B_x$.

\medskip\noindent
The next set of points we consider in
$\mathfrak g$ are those which are
regular but not semisimple. All of these
elements are nilpotent and lie in a
single adjoint orbit called the principal (or
regular) nilpotent orbit. The
Springer fibre $\mathcal B_x$ over a
principal nilpotent element $x$ consists
of a single point. (FIND THE STATEMENT in
\cite{ChrissGinzburg})

\medskip\noindent
The set of all nilpotent elements of
$\mathfrak g$ is a subvariety
$\mathcal N \subset \mathfrak g$. The adjoint
action of $G$ on $\mathfrak g$
preserves $\mathcal N$ and it splits into a
disjoint union of nilpotent
orbits, it turns out there are finitely many
nilpotent orbits. The principal
nilpotent elements constitute one orbit,
which is dense in $\mathcal N$, i.e.,
makes up most of $\mathcal N$. The remaining
elements of $\mathcal N$ split
into a disjoint union of a bunch of adjoint
orbits of smaller dimension. The
dimension of $\mathcal N$ (and of the
principal orbit) is $2|\Delta^+|$ by the
way.

\medskip\noindent
Although $\mathcal B_x$ is smaller for $x$
principal nilpotent than for $x$
regular semisimple, the pattern does not
continue for $x$ in smaller nilpotent
orbits. In fact for such $x$ the Springer
fibre $\mathcal B_x$ can become a
higher dimensional variety. In fact we can
see that it must by passing to the
extreme case of $x = 0$, then $\mathcal B_0$
consists of all Boreis, and so is
a copy of $\mathcal B$ itself. For $x \in
\mathcal N$ the structure and
dimension of $\mathcal B_x$ depends on the
nilpotent orbit in which $x$ lies.

\medskip\noindent
The Weyl group $W$ does not act in a natural
way on the non regular Springer
fibres $\mathcal B_x$, but there is a way to
use the action on the regular
semisimple Springer fibres, and a sort of
limit procedure, to induce a
$W$-action on the Borel-Moore homology
$H^{BM}_*(\mathcal B_x)$.

\medskip\noindent
The union of all the nilpotent Springer
fibres is
$\mu^{-1}(\mathcal N) \subset
\widetilde{\mathfrak g}$, we denote this set
$\widetilde{\mathcal N}$. Recall that we have
a projection
$\nu : \widetilde{\mathfrak g} \to \mathcal
B$ with
$\nu^{-1}([\mathfrak b]) = \mathfrak b$. It
is not hard to show that $\nu$
restricts to a projection
$\widetilde{\mathcal N} \to \mathcal B$ under
which
the preimage of $\mathfrak b$ is $\mathfrak n
= [\mathfrak b, \mathfrak b]$.
So in fact $\widetilde{\mathcal N}$ is a
vector bundle over $\mathcal B$.
Furthermore $\widetilde{\mathcal N} \cong
T^*\mathcal B$.

\medskip\noindent
We consider another variety, called the
Steinberg variety and denoted $Z$. It
is
$$
Z = \widetilde{\mathcal N} \times_{\mathcal
N} \widetilde{\mathcal N} = \{(x, \mathfrak
b_1, \mathfrak b_2) \mid x \in \mathcal
N\text{ and }\mathfrak b_1, \mathfrak b_2 \in
\mathcal B_x\}.
$$
Obviously it projects to $\mathcal N$ with
fibre over $x$ being
$\mathcal B_x \times \mathcal B_x$. Another
way to view $Z$ is as a subset of
$\widetilde{\mathcal N} \times
\widetilde{\mathcal N} \cong T^*(\mathcal B
\times \mathcal B)$.

\medskip\noindent
In fact $G$ acts on $\mathcal B$ and hence
(via the diagonal action) on
$\mathcal B \times \mathcal B$. It turns out
that the diagonal $G$-orbits are
in natural bijection with $B\backslash G/B$,
which in turn is in bijection
with $W$:
$$
N_G(T)/T \longleftrightarrow B\backslash B
\longleftrightarrow (\mathcal B \times
\mathcal B)/\Delta G.
$$
If we denote by $Y_w \subset \mathcal B
\times \mathcal B$ the $G$-orbit
corresponding to $w \in W$, and denote by
$Z_w$ the conormal bundle
$T^*_{Y_w}(\mathcal B \times \mathcal B)$ of
$Y_w$ (or rather its total space)
then it turns out $Z$ equals the union of the
$Z_w$ [1, Proposition 3.3.4].
Notice that the dimension of the fibres of a
normal or conormal bundle equals
the codimension of the subvariety in
question. Thus the dimension of the total
space $Z_w$ equals that of $\mathcal B \times
\mathcal B$, i.e.,
$2|\Delta^+|$, same as the dimension of
$\mathcal N$.

\medskip\noindent
Let's return to thinking about $Z$ `over'
$\mathcal N$. To fix notation, here
is the diagram
$$
\xymatrix{
& Z \ar[dl] \ar[dr]
\ar[dd]^{\mu} & \\
\widetilde{\mathcal N} \ar[dr]_{\pi} & &
\widetilde{\mathcal N} \ar[dl] \\
& \mathcal N &
}
$$
For each nilpotent orbit $\mathbb{O} \subset \mathcal
N$ we consider
$\widetilde{\mathbb{O}} = \pi^{-1}(
\mathbb{O})$ (this is a
fibration over $\mathbb{O}$ with all fibres
isomorphic to $\mathcal B_x$ for some
representative $x \in \mathbb{O}$), and
$Z_{\mathbb{O}} = \mu^{-1}(
\mathbb{O})$ (which is similarly a
fibration with fibres isomorphic to
$\mathcal B_x \times \mathcal B_x$).

\medskip\noindent
Each subvariety $Z_w \subset Z$ `lies over' a
uniquely defined nilpotent orbit
$\mathbb{O} \subset \mathcal N$. This is stated in [1,
Remark 3.3.22], but I didn't
understand very well the reason for it. There
is a $G$-action on $Z_w$
compatible with that on $Y_w \subset \mathcal
B \times \mathcal B$. By
definition $Y_w$ is a single $G$-orbit, but I
don't know if $Z_w$ is a single
orbit or not. If it is, then it would be
clear that $\mu(Z_w)$ could not be a
union of more than one $G$-orbit $\mathbb{O}$ in
$\mathcal N$.... Anyway this means
there is a well defined function
$$
\Phi : W \to \{\text{nilpotent orbits}\}.
$$
This function should be compatible with
adjoint action of $G$ (which induced
the adjoint action of $W$ on the left hand
side), so elements of $W$ in the
same conjugacy class are sent to the same
nilpotent orbit. We might wonder how
far from a bijection is this function. For
$\mathfrak g = E_8$ there are 70
nilpotent orbits, but in the corresponding
Weyl group there are 112 conjugacy
classes. For $\mathfrak g = \mathfrak{sl}_n$
we know that both the set of
nilpotent orbits and the set of conjugacy
classes in $W \cong S_n$ are
parametrised in natural ways by partitions of
the integer $n$. As I understand
it, the function above is a bijection
compatible with these parametrisations.

\medskip\noindent
The function $\Phi$ above should not be
confused with the `Springer
correspondence' which is a function from
irreducible representations of $W$ to
nilpotent orbits. Of course the number of
irreducible representations of $W$
is the same as that of conjugacy classes. For
a general group one does not
expect any natural bijection between
conjugacy classes and irreducible
representations, but $S_n$ is a famous
exception.

\medskip\noindent
Indeed the Springer fibre $\mathcal B_x$
carries an action of the stabiliser
subgroup $G(x) = \{g \in G \mid gxg^{-1} =
x\}$, and we have $\mathbb{O} \cong G/G(x)$
and $\widetilde{\mathbb{O}} \cong G \times_{G(x)}
\mathcal B_x$. Next $Z_{\mathbb{O}}$ is presented
as
$$
Z_{\mathbb{O}} \cong G \times_{G(x)} (\mathcal B_x
\times \mathcal B_x).
$$
The Springer fibres are known to be connected
(this is not easy to prove, it
seems) but are not irreducible. So we talk

\section{Case study: Borel subalgebras of 
sl2}
\label{springer-section-borel-subalgebras-sl2}

\noindent
Upshot: Borel subalgebras of $\mathfrak{sl}_2$
are line-stabilizing subalgebras.
That is, subalgebras $\mathfrak{b}_L$ such that
$XL \subset L$ for some line $L$
and all $X \in \mathfrak{b}_L$.
For example,
$$
\mathbb C e_1
\longleftrightarrow
\left\{
\begin{pmatrix}
a&b\\
0&-a
\end{pmatrix}
\right\},
\qquad
\mathbb C e_2
\longleftrightarrow
\left\{
\begin{pmatrix}
a&0\\
c&-a
\end{pmatrix}
\right\}.
$$

That's why we have:
$$
\mathcal{B}(\mathfrak{sl}_2)
\cong
\mathbb{P}^1.
$$
Now, for a regular semisimple element
like
$$
x = \begin{pmatrix}1&0\\ 0&-1\end{pmatrix},
$$
we may check that the $x$-invariant
lines are exactly $\mathbb{C}e_1$
and $\mathbb{C}e_2$. It's a finite set
of lines, that is, of Borel subalgebras.

Similarly, for a regular nilpotent element
like
$$
e = \begin{pmatrix}0&1\\ 0&0\end{pmatrix},
$$
we can check that there is only
one $e$-invariant line: $\mathbb{C}e_1$.
There's only one nilpotent element that
is not regular in this case: the zero matrix.
Of course, this fixes every line, so
the Springer fibre is all of $\mathcal{B}$.

\medskip\noindent
Let $V=\mathbb C^2$ and
$\mathfrak g=\mathfrak{sl}(V)=\mathfrak{sl}_2$.
We want a concrete model for the flag variety
of $SL_2$, which is the set
of Borel subalgebras.

A point of $\mathbb P^1=\mathbb P(V)$ is a
line $L\subset V$. To such a line we associate
the subalgebra
$$
\mathfrak b_L
=
\{X\in\mathfrak{sl}(V):X(L)\subset L\}.
$$
In words, $\mathfrak b_L$ consists of the
traceless endomorphisms of $V$ preserving the
line $L$.

For example, take $L=\mathbb C e_1$. Write
$$
X=
\begin{pmatrix}
a&b\\
c&-a
\end{pmatrix}
\in\mathfrak{sl}_2.
$$
Then
$$
Xe_1=(a,c).
$$
The condition $X(\mathbb C e_1)\subset
\mathbb C e_1$ is equivalent to $c=0$.
Therefore
$$
\mathfrak b_{\mathbb C e_1}
=
\left\{
\begin{pmatrix}
a&b\\
0&-a
\end{pmatrix}
:a,b\in\mathbb C
\right\},
$$
the upper triangular Borel subalgebra.

Similarly, for $L=\mathbb C e_2$ we have
$$
Xe_2=(b,-a).
$$
The condition $X(\mathbb C e_2)\subset
\mathbb C e_2$ is equivalent to $b=0$.
Therefore
$$
\mathfrak b_{\mathbb C e_2}
=
\left\{
\begin{pmatrix}
a&0\\
c&-a
\end{pmatrix}
:a,c\in\mathbb C
\right\},
$$
the lower triangular Borel subalgebra.

Thus lines in $V$ give Borel subalgebras of
$\mathfrak{sl}_2$. Conversely, every Borel
subalgebra of $\mathfrak{sl}_2$ is obtained
this way. Hence 
if $\mathcal{B}$ is the set
of Borel subalgebras of $\mathfrak{sl}_2$,
we have
$$
\mathcal B
\cong
\mathbb P^1.
$$

Under this identification, a point of the flag
variety may be viewed concretely as a line
$L\subset\mathbb C^2$, and the corresponding
Borel subalgebra is the stabilizer Lie algebra
of that line.

\section{Component groups}
\label{springer-section-component-groups}

\noindent
The [?, Proposition 3.3.27]

\medskip\noindent
Coll McG on the two (main) component
groups....

\medskip\noindent
Remark:
\cite{ChrissGinzburg} on page 176 discuss the component
group and why its action can be
dispensed with in the type A case. See also
p. 170.

\section{Weyl group actions on homology of
Springer fibres}
\label{springer-section-weyl-action}

\noindent
For regular semisimple $x \in \mathfrak g$
there is a $W$-action on the
Springer fibre $\mathcal B_x$.

\medskip\noindent
Let $Y$ be a set, then we consider subsets $X
\subset Y \times Y$, and a
`convolution operation':
$$
X \circ X' = \{(y_1, y_3) \mid \text{for some
}y_2 \in Y\text{ both }(y_1, y_2) \in X\text{
and }(y_2, y_3) \in X'\}.
$$
This can be done more generally, but let's
try to keep things as simple as
possible. Notice that this operation
generalises composition of functions
$f : Y \to Y$. Indeed any such function
defines a subset
$X_f = \{(y, f(y)) \mid y \in Y\}$, and then
$X_f \circ X_g = X_{g \circ f}$.

\medskip\noindent
A slightly more general construction is to
start with a surjection $Y \to N$
and then consider composition of subsets of
$Z = Y \times_N Y$. In Springer
theory, this machinery is applied to the
example in which $Y \to N$ is
$\widetilde{\mathfrak g} \to \mathfrak g$ or
$\widetilde{\mathcal N} \to \mathcal N$.

\medskip\noindent
Suppose, in the setting above, that the same
element $(y_1, y_3)$ of the
composition were obtained via two different
$y_2$. Then one would like to keep
track of this multiplicity somehow. We might
define $C_{fin}[Y]$ to be the
vector space of finite linear combinations of
points of $Y$, and then extend
the pointwise convolution
$$
(y_1, y_2) \circ (y'_2, y_3) = \delta_{y_2,
y'_2}(y_1, y_3)
$$
linearly to $C_{fin}[Y]$ in the only natural
way. Identifying $(y_1, y_2)$
with an elementary matrix with 1 at
coordinates $(y_1, y_2)$, this convolution
is just matrix multiplication.

\medskip\noindent
For the applications to Springer theory it is
not enough to only work with
finite subsets of $Z$. Instead we work with
some variant of homology. In fact
we work with Borel-Moore homology, but the
reason for this is not really clear
at the outset.

\medskip\noindent
To do: work out what happens for usual
homology, or if it works....

\medskip\noindent
A zero homology class is nothing but a finite
collection of points, and one
can interpret the convolution procedure of
0-cycles as
$$
\alpha \circ \alpha' = (\pi_{13})_*
(\pi^*_{12}(\alpha) \cap
\pi^*_{23}(\alpha')).
$$
This looks like we have converted the
convolution into operations in homology,
but in fact we have not, not for ordinary
homology at least! The pullback
$\pi^*_{12}(\alpha)$ takes the form of a
finite set of points (namely
$\alpha$) times $Z$, and this is not a
homology class since $Z$ is not
compact!

\medskip\noindent
Recall, roughly, that the singular homology
of a `reasonable' topological
space $X$ is computed as follows: $C_i(X)$
denotes the abelian group of finite
formal linear combinations $\sum_k
a_k\sigma_k$ of continuous maps
$\sigma_k : \Delta_i \to X$. There are
differentials $C_i(X) \to C_{i-1}(X)$
and homology $H_i(X)$ is given by ker /im.

\medskip\noindent
The space of Borel-Moore chains is larger:
$C^{BM}_i(X)$ consists of possibly
infinite formal linear combinations $\sum_k
a_k\sigma_k$ of continuous maps
$\sigma_k : \Delta_i \to X$, subject to the
restriction that for any fixed
compact set $K \subset X$, the set of $k$ for
which $\sigma_k$ meets $K$ and
$a_k \neq 0$, is finite. Clearly if $X$ is
compact then
$C^{BM}_i(X) = C_i(X)$.

\medskip\noindent
Singular cohomology (with coefficients in
$R$) is given by ker /im for the
cochain complex $C^i(X) =
\operatorname{Hom}(C_i(X), R)$. (I think it's
enough
for us to take all our homologies and
cohomologies with coefficients in
$R = \mathbb C$.) The last theory is
cohomology with compact support:
$C^i_c(X)$ consists of those cochains $\alpha
\in C^i(X)$ for which there
exists some compact $K \subset X$ such that
$\alpha(\sigma) = 0$ for all
chains $\sigma$ supported in $X\backslash K$.

\medskip\noindent
Poincare duality: For a smooth manifold (not
necessarily compact) of real
dimension $d$ there are natural isomorphisms
$$
H^{BM}_{d-i}(X) \cong H^i(X)
\quad\text{and}\quad H_{d-i}(X) \cong
H^i_c(X).
$$
A potential confusion: the first of these
isomorphisms suggests that working
with BM homology is essentially equivalent to
working with ordinary
cohomology, which is surely less exotic and
thus preferable. Maybe this is
true for smooth varieties, but the important
thing to keep in mind is that our
example the Steinberg variety is mildly
non-smooth, and this makes a big
difference. The non-smoothness is quite mild
indeed: it is in the same sense
that the complex algebraic variety $xy = 0$,
which is a union of two
intersecting complex lines, is non-smooth.
But still, it means we can't rely
on Poincare duality, at least not in the
traditional sense.

\medskip\noindent
To get representations of the Weyl group from
the Steinberg variety, the thing
to do is to look at the top degree
$$
H^{BM}_{top}(Z) = H^{BM}_{4|\Delta^+|}(Z).
$$
(The top degree is equal to the real
dimension of $Z$, which is twice the
complex dimension of $Z$. The dimension of
$Z$ equals that of $\mathcal N$.
The complex dimension of $\mathcal N$ is
$2|\Delta^+|$.)

\medskip\noindent
An interesting feature of BM homology is that
it can detect irreducible
components of a complex variety. The complex
variety $X$ might well be
connected, but if it is equidimensional with
several irreducible components of
complex dimension $n$, then the top degree BM
homology is $H^{BM}_{2n}(X)$ and
this one is of dimension equal to the number
of irreducible components.

\medskip\noindent
This shocked me somewhat at first. It means,
in particular, that BM homology
is not homotopy invariant. Now, no one ever
said that it was..., but it's the
sort of thing I lazily assumed.

\medskip\noindent
There is a sort of excision result for BM
homology, which implies in
particular that the BM homology in top degree
of $X$ equals that of $X$ minus
the intersections of the irreducible
components. Since the intersections have
real codimension 2 (so here is where the
complex variety structure enters),
the remainder is smooth with each irreducible
component contributing now a
connected component. It is proven somewhere
that the $|W|$ distinct conormal
bundles that make up the Steinberg variety
$Z$ are exactly its irreducible
components. So $H^{BM}_{top}(Z) \cong
\bigoplus_{w \in W} \mathbb C[w]$.

\medskip\noindent
It will turn out that the convolution algebra
structure on $H^{BM}_*(Z)$ has
$H^{BM}_{top}(Z)$ as a subalgebra. Sure
enough, this algebra will turn out to
be isomorphic to $\mathbb C[W]$. The homology
$H^{BM}_*(\mathcal B_x)$ of any
particular Springer fibre should now become a
$\mathbb C[W]$-module in a
natural way via this sort of construction.

\medskip\noindent
TO BE COMPLETED....

\section{Slodowy's monodromy construction}
\label{springer-section-slodowy-monodromy}

\noindent
Let $x \in \mathcal N$ be a nilpotent element
and $\mathbb{O} = G \cdot x$ its orbit.
Recall the construction by Slodowy of a
`transverse slice' $S$ to $\mathbb{O}$ at $x$,
i.e., $T_x\mathfrak g = T_x\mathbb{O} \oplus T_xS$.
Explicitly $S = x + \mathfrak g^y$,
where we used an $\mathfrak{sl}_2$-triple. We
recall the Springer resolution
$\mu : \widetilde{\mathfrak g} \to \mathfrak
g$ and we consider its
restriction $\mu : \widetilde S \to S$.
Slodowy proves $\widetilde S$ is
smooth [?, 3.4]. The proof uses
transversality of $S$ with $G$-orbits in a
way
I do not understand, but seems to be
standard.

\medskip\noindent
We take a small closed ball $D \subset S$,
centered on $x$. Set
$K = \mu^{-1}(D) \subset \widetilde S$.
Because $\mu$ is proper we have that
$K$ is compact. Remember that $\mu$ is a
$|W|$-to-1 finite map over the set of
regular semisimple elements of $S$ (which is
most elements). But there is in
general a nontrivial intersection of $S$ with
$\mathcal N$ (the exception is
when $\mathbb{O}$ is the principal orbit, so that
transversality of $S$ with $\mathbb{O}$ means
transversality with $\mathcal N$ in fact).

\medskip\noindent
We also take a small open ball $H \subset
\mathfrak h$ around 0, consider the
preimage $\theta^{-1}(H)$ in $\widetilde S$
under the map
$\theta : \widetilde{\mathfrak g} \to
\mathfrak h$, and intersect it with $K$.
Notice that $\theta^{-1}(0) =
\widetilde{\mathcal N}$. Recall\footnote{I
haven't yet written out this part.}
$$
\theta : \widetilde{\mathfrak g} \to
\mathfrak h
$$
is defined either with $\mathfrak h$ the
`abstract Cartan subalgebra', or else
as follows: we fix a choice of Borel
$\mathfrak b_0$ and a Cartan in it
$\mathfrak h_0$, our point in $\widetilde S$
is a pair $(s, \mathfrak b)$
where $s \in S$ and $\mathfrak b$ is some
Borel containing $s$. We choose an
element $g \in G$ that conjugates $\mathfrak
b$ to $\mathfrak b_0$, and define
$\theta(s)$ to be the projection of $g(s) \in
\mathfrak b_0$ to
$\mathfrak h_0$.

\medskip\noindent
My mental image: $K$ is a potato-like blob,
consisting of a smooth surface
together with its interior. There is a map
$\theta : K \to \mathfrak h$ whose
zero level set is the intersection of
$\widetilde{\mathcal N}$ with $K$.
Inside this level set sits the compact
manifold $\mathcal B_x$. The preimage
of $H$ I view as taking a slight fattening of
$K \cap \theta^{-1}(0)$. The
result of Slodowy (which is not visually
obvious) is that all these fibres
$\theta^{-1}(h) \cap K$ are homotopic
manifolds with boundary.

\medskip\noindent
We intersect to get $\theta^{-1}(H) \cap K
\to H$. There are two claims: (1)
the boundary $\partial K$ is a manifold, and
(2) the map above is a
differentiable fibration with isomorphic
fibres.

\medskip\noindent
The interior of $K$ is an open set in
$\widetilde S$, so I guess it's
obviously a manifold. Slodowy shows $\partial
K$ is a manifold using a
standard application of `transversality': the
$\mathbb C^\times$-action allows
one to show that $\partial K$ is transversal
to $\mu$. The claim about the
fibration follows from Ehresmann's lemma
which says that a proper surjective
submersion is a locally trivial fibration.
Notice that Slodowy is applying a
version of Ehresmann's lemma for `manifold
with boundary'.

\medskip\noindent
The stated properties seem to be easily
verified, except for the `zeroth
property' which would be that the space up
stairs is a manifold. We know $K$
is a manifold with boundary, and we are
intersecting with an open, but the
result is supposed to have compact fibres. So
what I need to understand is why
the fibres are compact manifolds, not open
manifolds nor manifolds with
boundary. Apparently transversality with
$\mathcal N$ is relevant. Maybe it's
also like... a version of Ehresmann's lemma
where submersion property is
imposed in the interior and boundary of a
compact set, INSIDE a larger
manifold, surjecting to another.

\medskip\noindent
Idea (maybe misguided): Using differential
topology to show a bunch of fibres
are homotopic, and then getting a $W$-action
on cohomology of the fibre from a
geometric action on a larger space / related
space / how exactly? Anyway, is
this picture, perhaps, one of the motivations
for people wanting to bring
homotopical methods into representation
theory? If so, where to look for this
POV?


\section{The Springer sheaf}
\label{springer-section-springer-sheaf}

\noindent
We consider the pushforward $\mu_*(\mathbb
Q[\dim \widetilde{\mathcal N}])$ of
the shift of the constant sheaf $\mathbb Q$,
along
$\mu : \widetilde{\mathcal N} \to \mathcal
N$. We recall the fibres
$\mathcal B_x$ are all projective varieties,
so $\mu_* = \mu_!$ which is
useful at some points. The inequality
$$
\dim(G \cdot x) + 2 \dim(\mathcal B_x) = \dim
\widetilde{\mathcal N}
$$
quite easily proves that $\mu_*(\mathbb
Q[\dim \widetilde{\mathcal N}])$ is a
perverse sheaf. These are ubiquitous, so I
imagine this is very important, but
so far I do not know what this fact is used
for.

\medskip\noindent
The dimension condition above generalises to
the following notion: A proper
morphism $f : X \to Y$ is called semismall if
$Y$ can be given a
stratification $\{Y_i\}$ such that $\dim Y_i
+ 2 \dim f^{-1}(y) \leq \dim X$
for all $y \in Y_i$. Then [2, Corollary 2.2]
semismallness of a proper map
$f : X \to Y$ is equivalent to perversity of
$f_*(\mathbb Q[\dim X])$.

\medskip\noindent
It is easy to think of examples of proper
maps which are not semismall (I
suppose, though correct me if I'm wrong, that
the map $Z \to \mathcal N$ from
the Steinberg variety is not semismall).
Kaledin's theorem, though, says that
a resolution $f : X \to Y$ by symplectic $X$
is semismall. So the fact that
$\widetilde{\mathcal N} \cong T^*\mathcal B$
makes semismallness of the
Springer resolution immediate.

\section{The affine Weyl group}
\label{springer-section-affine-weyl-group}

\noindent
There is an algebra homomorphism
$$
\mathbb C[\widehat W] \cong K_G(Z)
$$
between the group algebra of the affine Weyl
group $\widehat W$, and the
$G$-equivariant $K$-theory of the Steinberg
variety. ``The $K$-theory'' and
``equivariant $K$-theory'' of a variety $X$
are invariants distilled from the
categories $\operatorname{Coh}(X)$ and
$\operatorname{Coh}_G(X)$ respectively.
Distilled in the sense that they are the
Grothendieck groups thereof: i.e.,
one mods out by the relation that an
extension of two objects is set equal to
their direct sum.

\medskip\noindent
I think I had my first inkling of why the
affine Weyl group would appear in
all this. This group or something related to
it, arises as a homotopy group of
a set of regular elements in $\mathfrak g$ or
$\mathfrak h$ or somesuch thing.
The reason I ask is that Ginzburg relates the
affine Weyl group to equivariant
$K$-theory of the Steinberg, while Slodowy's
approach interprets Springer's
$W$-reps as monogromy action of $W$
interpreted as a fundamental group. In
fact in Slodowy's approach the fundamental
groups of $\mathfrak h'$ and
$\mathfrak h'/W$ appear, but the objects
under consideration are trivial
upstairs, and so the reps, a priori of
$\pi_1(\mathfrak h')$, actually factor
to $W$. I could imagine some other object
which does not trivialise upstairs,
giving reps of the braid group. But then
naturally I wonder if Ginzburg's
construction can be thought of in these
terms... but the idea is a non-starter
unless there's actually some relationship
between the affine Weyl group and
the braid group.

\medskip\noindent
[See Orlik, Terao, `Arrangements of
hyperplanes'.]

\medskip\noindent
Consider the complement $M \subset \mathbb
C^\ell$ of the hyperplanes
$z_i = z_j$. Better to picture this space as
collections of $\ell$ distinct
labeled points in $\mathbb C$. It is
essentially clear that $\pi_1(M)$ is the
pure braid group $PB_\ell$. The braid group
$B_\ell$ is the group of all
braids on $\ell$ labeled strands, it
obviously surjects to $S_\ell$, and
$PB_\ell$ is the kernel. This is the group of
braids that start and end at the
same configuration. The complement of
hyperplanes in
$\mathfrak h \subset \mathfrak{sl}_n$ is the
intersection of $M$ with the
locus $\sum z_i = 0$. This corresponds to
sets of $\ell$ distinct points whose
centre of gravity is fixed at 0. Clearly the
homotopy group is unchanged by
this modification.

\medskip\noindent
We can consider the quotient $M/S_\ell$,
whose homotopy group should be
$B_\ell$ I think, the full braid group.

\medskip\noindent
QUESTION: Is there a close relation between
the full braid group $B_\ell$ and
the affine Weyl group of $S_\ell$?

\medskip\noindent
Remark: there is a Wikipedia article on the
``double affine braid group'',
which is said (in type A) to be the
fundamental group of the space of $n$
distinct points in the 2-torus.

\medskip\noindent
I found some notes on the internet about the
Coxeter-style presentation, and
Bernstein presentation, of the affine Weyl
group. I think by a student of
Geordie Williamson. I seem not to have saved
the note, so try to find it
again. See whatsapp auto-convo.

\section{June 25, 2026}
\label{springer-section-june25-2026}

\noindent
Springer fibres 
(and Weyl group action on their
(co)homology).

\medskip\noindent
Let $G$ be a simple Lie group
and $\mathfrak{g}=\text{Lie} G$.
Root system $\Delta \subset \mathfrak{h}^*$,
$\mathfrak{g}=\mathfrak{n}_- \oplus
\mathfrak{h} \oplus \mathfrak{n}_+$,
where $\mathfrak{h}$ is Cartan
and $\mathfrak{h}\oplus \mathfrak{n}_+$
is Borel (a Borel is a maximal solvable
subalgebra).
$\mathcal{B}=$ set of all borels in
$\mathfrak{g}$.

$G$ acts on $\mathcal{B}$ (easy)
transitivel (not so easy but OK).
$\operatorname{Stab}_G(\mathfrak{b})=B$,
a Borel subgroup.
So $\mathcal{B}\simeq G/B$.
Also $\mathcal{B} \subset \text{Gr}
(\dim \mathfrak{g}, \dim \mathfrak{b})$,
so it's a compact (projective)
algebraic variety.

$\tilde{\mathfrak{g}}=
\{(\mathfrak{b},x):\mathfrak{b} \in 
\mathcal{B}, x \in \mathfrak{b}\}$ is
the tautological bundle.
$$
\xymatrix{
&  \tilde{\mathfrak{g}}
\ar[dl]^{\mu}\ar[dr]\\
\mathfrak{g}&  &  \mathcal{B}
}
$$
Preimages $\mu^{-1}(x) \subset \tilde{\mathfrak{g}}$
are the Springer fibres.

{\bf Fact.} If $x \in \mathfrak{g}$
is regular semisimple,
$\# (\mu^{-1}(x))=\# W$ where
$W$ is the Weyl group.

\begin{example}
\label{springer-example-sln-borels}
$\mathfrak{g} = \mathfrak{sl}_n$,
then $\mathcal{B}$ is isomorphic to the
{\it actual} flag variety
$G/B \simeq \mathcal{F}
=\{0=V_0 \subset V_1\subset
… \subset V_n = \mathbb{C}^n:
\dim V_j = j\}$.
\end{example}

\noindent
Let $x \in \mathfrak{sl}_n$ be
regular semisimple, so $x \in
\text{End}(\mathbb{C}^n)$ has distinct
eigenvalues. Each eigenvalue
has an eigenspace $\mathbb{C}v_{\lambda_i}$.
Let $F_i = \bigoplus_{r=1}^i
\mathbb{C}v_{\lambda_r}$.
This gives a flag $F \in \mathcal{F}$
and $x$ acts on it by diagonal
matrix
$$
\begin{pmatrix}
\lambda_1\\
& \lambda_2\\
& & \ddots\\
& & & \lambda_n
\end{pmatrix}
$$
$$
\mathcal{F} \leftrightarrow \mathcal{B}
$$
$$
F \leftrightarrow \mathfrak{b},
$$
and $x \in \mathfrak{b}$ (exercise).

$\mu^{-1}(x) = $ all $F \in \mathcal{F}$
obtained by ordering $\lambda,…,\lambda_n$
in different ways.

\medskip\noindent
$\mathcal{N}=\{x \in \mathfrak{g}:
\text{ad}_x \text{ nilpotent} \}
\supset \mathcal{O}_{\operatorname{pr}}$,
a principal $G$-orbit.

$\dim \mathcal{N} =
\dim\mathcal{O}_{\operatorname{pr}}
= \dim \mathfrak{g} -
\dim \mathfrak{g}^e$,
$e \in \mathbb{O}_{\operatorname{pr}}$
$= | \Delta|=
2 |\Delta_+|$.

$\mathcal{N}= \bigcup$ nilpotent orbits.
For $x \in \mathbb{O}_{\operatorname{pr}}$,
$\mu^{-1}(x) = *$ one point.
For $x=0$, $\mu^{-1}(x) = \mathcal{B}$.
In general, $\mathcal{B}_x=
\mu^{-1}(x)$ are some (possibly
singular) projective algebraic varieties.
(Recall that for the $\mathfrak{sl}_n$
case there are no other Springer
fibres than the regular semisimple
and 0, but in general there are more,
as we see in the next example.)

\begin{example}
\label{springer-example-e8-springer-fibres}
$\mathfrak{g} = E_8$,
$\mathcal{N}\setminus
\mathbb{O}_{\operatorname{pr}}
= \bigcup$ nilpotent orbits,
and the largest one is
the so-called subregular
orbit $\mathbb{O}_{\operatorname{sr}}$.
$x \in \mathbb{O}_{\operatorname{sr}},
\mu^{-1}(x)$ is a union of $\mathbb{P}^1$
intersecting just as the $E_8$ 
Dynkin diagram.
\end{example}

\noindent
Write
$\widetilde{\mathcal{N}}=\mu^{-1}(\mathcal{N})
\subset \tilde{\mathfrak{g}}$.
$$
\xymatrix{
\tilde{\mathfrak{g}}\ar[d]
&  \tilde{\mathcal{N}}
\simeq T^*\mathcal{B}
\ar@{^{(}->}[l]\ar[d]\\
\mathcal{B}\ar@{=}[r]&  \mathcal{B}.
}
$$

\begin{remark}
\label{springer-remark-symplectic}
$\mathfrak{g}$ comes with a Poisson structure,
$\mathcal{N}$ is a Poisson variety,
$T^*M$ is always a symplectic manifold.
(Symplectic leaves = coadjoint orbits.)
$\mu :T^*\mathcal{B} \to N$
symplectic singularity.
\end{remark}

\noindent
Steinberg variety.
$Z = \{(x,\mathfrak{b}_1,\mathfrak{b}_2:
x \in \mathfrak{g}, \mathfrak{b}_1,
\mathfrak{b}_2 \in \mathcal{B},
x \in \mathfrak{b}_1,
x \in \mathfrak{b}_2\}$.
$$
\xymatrix{
Z\ar[d]^\pi\\
\mathcal{N}
}
$$
It's a fibred product: $Z = 
\widetilde{\mathcal{N}}\times_\mathcal{N}
\widetilde{\mathcal{N}}$.

Over any point, we get the
product of the Springer fiber with itself,
$\pi^{-1}(x) = \mathcal{B}_x
\times \mathcal{B}_x$.

\noindent
Fact. $2 \dim \mathcal{B}+
\dim(G\cdot x) = \dim \mathcal{B}
= \dim \mathcal{N}=2 |\Delta_+|$
$\forall x \in \mathcal{N}$.

This implies, for orbits
$\mathbb{O} \subset \mathcal{N}$,
$Z \supset\pi^{-1}(\mathbb{O})$
has dimension $\dim \mathbb{O}
+ 2 \dim \mathcal{B}_x=
2 |\Delta_+|$.

An ``explanation'' of this coincidence:
think about $Z \subset \widetilde{
\mathcal{N}}\times \widetilde{\mathcal{N}}
=T^* \mathcal{B} \times T^*\mathcal{B}
=T^*(\mathcal{B}\times \mathcal{B})$.

\begin{definition}
\label{springer-definition-conormal-bundle}
If $N \subset M$ submanifold,
the {\it conormal bundle} is
$T^*_NM \subset T^*M$ is 
$\{(n,\xi) \in T^*M:
n \in N, \xi \in T^*_n M,
\xi(T_n N) = 0\}$.
\end{definition}

\noindent
$T^*_NM$ is a vector bundle over $N$,
$\dim T^*_NM= \dim M$ always.

\begin{example}
\label{springer-example-conormal-bundle-r}
$M= \mathbb{R}$, $T^*M \simeq \mathbb{R}$.
if $N= \mathbb{R}\setminus\{0\}$,
$T^*_NM$ is the zero section.
If $N = \{0\}$, then $T^*_NM =
T^*_0\mathbb{R}$,
they're both 1-dimensional.
\end{example}

\noindent
Consider $\mathbb{R}^{\times}$,
which acts on $\mathbb{R}$ by multiplication.
$\mathbb{R} = U \coprod \{*\}$.
$\bigcup_{\mathcal{O}}T^*_{\mathcal{O}}
\mathbb{R}=$ union of two axes.

\medskip\noindent
Obviously $G$ acts on $\mathcal{B}\times
\mathcal{B}$ (diagonal action),
and it turns out that
$\mathcal{B}\times\mathcal{B}=\coprod_{w \in
W}Y_w$, $W$-many $G$-orbits.

\begin{proposition}
\label{springer-proposition-union-of-conormal-bdls}
$\widetilde{\mathcal{N}}\times
\widetilde{\mathcal{N}}\supset Z
= \coprod_{w}T^*_{Y_w}
(\mathcal{B}\times\mathcal{B})$.
\end{proposition}

\noindent
The fact that they all have $\dim
= 2 |\Delta_+|$ kind of explains the
remark about $\dim \mathcal{B}$.

\begin{exercise}
\label{springer-exercise-lands-in-single-orbit}
Each $T^*_{Y_w}(\mathcal{B} \times \mathcal{B})
\xrightarrow{\pi}$ lands in 
a single orbit $\mathbb{O}$.
\end{exercise}

\noindent
So we have a function
$$
\Phi:W \to \{\operatorname{nilpotent
orbits}\}.
$$
\begin{example}
\label{springer-example-e8-nilpotent-orbits}
For $E_8$, $W$ has 112 conjugacy classes,
which maps to 70 nilpotent orbits.
\end{example}

\medskip\noindent
Next. $W$ acts on ``(co)homology''
of the $\mathcal{B}_x$, i.e.
each $x \in \mathbb{O} \subset \mathcal{N}$
gives us some representations of $W$.

\begin{remark}
\label{springer-remark-w-not-act-springer-fibre}
$W$ does not act on $\mathcal{B}_x$.
\end{remark}

\noindent
Two approaches
(There are two standard ways to construct the
Springer $W$-action on the homology of
Springer fibres: one using the Steinberg
variety and convolution, and another using
Slodowy slices and monodromy):
\cite{ChrissGinzburg} explain
that
$$
H^{BM}_{\operatorname{top}}(Z)
\simeq \mathbb{C}[W]
$$
and $H^{BM}_{\operatorname{top}}
(Z)$ acts on $H^{BM}_{\operatorname{top}}
(\mathcal{B}_x)$
gives the $W$-action.

That is,
top BM homology of the Steinberg variety Z
has a basis indexed by Weyl group elements,
and its convolution product matches the
multiplication in the group algebra.

\noindent
Finite set $X$, $n$ elements.
$Z = X \times X$,
$(x,y) *(z,w)
= \delta_{y,z}(x,w)$
extend linearly to
$$
\mathbb{C}[Z]
= \text{Mat}_{n}.
$$
Recover multiplication of matrices.
How to do $X$ infinite?
$$
\xymatrix{
&  X\times X\times X
\ar[dl]^{\pi_{12}}\ar[dr]^{\pi_{23}}
\ar[dd]^{\pi_{13}}\\
X\times X& & X\times X\\
& X \times X
}
$$

\begin{definition}[Informal]
\label{springer-definition-informal}
$\alpha \in H_*(Z), \beta \in H_*(Z)$,
$\alpha * \beta
= \pi_{13,*}(\pi_{12}^* (\alpha)
\cap \pi_{23}^*(\beta)$.
\end{definition}

\noindent
The problem here is that if $X$ is not
compact, $\pi_{12}$ not proper,
and $\pi_{12}^*$ not well-defined.

Borel-Moore homology works
with infinite (but ``locally finite'')
sets of simplices.
So $\pi^*_{12}$ makes sense in $H_*^{BM}$.

In our case,
$X = \widetilde{\mathcal{N}}$,
and we work over $\mathcal{N}$.
So really $Z \neq X \times X$,
rather $X = $ Springer,
$Z = $ Steinberg,
$Z = X \times_{\mathcal{N}}X$.

For $H_*^{BM}$, $\pi_*$ is not
in general well-defined.
But in our case, being in
$\pi_{12}^* (\alpha)
\cap \pi_{23}^*(\beta)$.

\medskip\noindent
Poincaré duality says that for any
orientable $d$-manifold $M$,
$$
H^i(M) \simeq H_{d-i}^{BM}(M)
$$
$$
H^i_c(M) \simeq H_{d-i}(M).
$$
So one might imagine we could just use 
usual cohomology. But $Z$ is singular, not a
manifold, so we don't have Poincaré
duality!

Also $H^{BM}_{\operatorname{top}}
(Z) = $
General fact: if $X$ is equidimensional
possibly singular complex variety
then $H^{BM}_{\operatorname{top}}(X)
=H_0(X \setminus $intersection of
irreducible components$)$.

\medskip\noindent
Second approach (Slodowy).
For $x \in \mathbb{O} \subset \mathcal{N}$,
take a transverse slice in 
$\mathfrak{g}$, $S \subset \mathfrak{g}$,
$x \in S$, i.e.
$T_x \mathfrak{g} \simeq T_x
\mathbb{O} \oplus T_x S$.

\begin{example}
\label{springer-example}
$S = x + \mathfrak{g}^y$,
$(x,h,y)$ $\mathfrak{sl}_2$-triple.
\end{example}

\noindent
Slodowy considers a disc $D$ in $S$,
centered on $x$.
$K = \mu^{-1}(D)\subset \mathcal{B}_x$.
$$
\xymatrix{
&  \tilde{g}
\ar[dl]^{\mu}\\
\mathfrak{g}&  &  \tilde{S}\ar@{^{(}->}[ul]
\ar[dl]^{\mu}\\
&  S \ar@{^{(}->}[ul]
}
$$
$$
\xymatrix{
&\tilde{\mathfrak{g}}\ar[dl]
\ar[dr]\\
\mathfrak{g}\ar[dr]_{CST}&  &  \mathfrak{h}
\ar[dl]^{\pi}\\
& \mathfrak{h}/W.
}
$$
where CST is Chevalley-Shepperd-Todd.

This is related to GIT quotients:
$$
\mathbb{C}[\mathfrak{g}]^G
\simeq\mathbb{C}[\mathfrak{h}]^W,
$$
$$
\mathfrak{g} \to \mathfrak{g}//G
\simeq\mathfrak{h}//W.
$$
$U \subset \mathfrak{h}\setminus H$,
where $H$ is a union of hyperplanes.
$\pi:U \to \pi(U)$ is a nice $W:1$-map.

Consider $\sigma^{-1}(U) \cap K$.
So some differential topology
(Ehresmann, submersion $\to$ locally
trivial fibration, bla, bla),
get action of $\pi_1(U)$ on fibres $F$
which are homotopic to $\mathcal{B}_x$.

$\pi_1(U) \simeq$ affine Weyl group
$\hat{W}$, but the action of $\hat{W}$
factors to an action of $W$,
where $W$ is the finite Weyl group.

The idea is that we have
an action of the fundamental
group on the cohomology of the fibers
of this differential-topology-obtained
fibration.

\section{Yun's description of
$T^*\mathcal B$}
\label{springer-section-yuns-description}

Yun's notes \cite[\S 1.2.1]{YunSpringer}
state the Springer resolution in a very
concrete way:
the cotangent bundle $T^*\mathcal B$
classifies pairs
$$
(e,B)
$$
where $e \in \mathcal N$ and $B$ is a Borel
subgroup such that
$$
e \in \mathfrak n_B,
$$
where $\mathfrak n_B$ is the nilpotent
radical of $\operatorname{Lie}B$.
Thus, in the notation of these notes,
$$
T^*\mathcal B
\simeq
\widetilde{\mathcal N}
=
\{(e,\mathfrak b):
\mathfrak b\in\mathcal B,\,
e\in\mathfrak n_{\mathfrak b}\}.
$$

One way to remember the construction is this.
At $\mathfrak b\in\mathcal B$,
$$
T_{\mathfrak b}\mathcal B
\simeq
\mathfrak g/\mathfrak b.
$$
Using the Killing form to identify
$\mathfrak g\simeq\mathfrak g^*$, we get
$$
T^*_{\mathfrak b}\mathcal B
\simeq
(\mathfrak g/\mathfrak b)^*
\simeq
\mathfrak b^\perp.
$$
For a Borel subalgebra,
$$
\mathfrak b^\perp=\mathfrak n_{\mathfrak b}.
$$
This is the promised identification of the
cotangent fibre with the nilradical.

The Springer resolution is then the forgetful
map
$$
\mu:T^*\mathcal B\simeq\widetilde{\mathcal N}
\longrightarrow \mathcal N,
\qquad
(e,\mathfrak b)\longmapsto e.
$$
So $\mu$ forgets the Borel, not the cotangent
vector.

\section{The zero section and the map to
$\mathfrak{g}$}

The point that was confusing me is therefore
the following.

Start with
$$
\widetilde{\mathfrak{g}}
=
\{(x,\mathfrak{b}):
\mathfrak{b}\in \mathcal{B},
x\in\mathfrak{b}\}.
$$
The map to $\mathfrak{g}$ is
$$
\mu(x,\mathfrak{b})=x.
$$
Thus $\mu$ forgets the Borel
$\mathfrak{b}$, not the element $x$.

If we fix one Borel $\mathfrak{b}$,
then the whole fibre of
$\widetilde{\mathfrak{g}}\to\mathcal{B}$
over $\mathfrak{b}$ maps to the whole
subalgebra $\mathfrak{b}$:
$$
\{(x,\mathfrak{b}):x\in\mathfrak{b}\}
\xrightarrow{\mu}
\mathfrak{b}.
$$
This is where all the nonzero elements
of $\mathfrak{b}$ appear.

Now restrict to the nilpotent cone.
Then
$$
\widetilde{\mathcal{N}}
=
\{(x,\mathfrak{b}):
\mathfrak{b}\in\mathcal{B},
x\in\mathfrak{b}\cap\mathcal{N}\}.
$$
For a Borel, the nilpotent elements in
its Lie algebra are precisely its
nilradical:
$$
\mathfrak{b}\cap\mathcal{N}
=
\mathfrak{n}_{\mathfrak{b}}.
$$
So
$$
\widetilde{\mathcal{N}}
=
\{(x,\mathfrak{b}):
x\in\mathfrak{n}_{\mathfrak{b}}\}.
$$

This space is identified with
$T^*\mathcal{B}$. Under this
identification, the cotangent vector is
the $x$-coordinate. More precisely, the
cotangent fibre over $\mathfrak{b}$ is
identified with the nilradical:
$$
T^*_{\mathfrak{b}}\mathcal{B}
\simeq
\mathfrak{n}_{\mathfrak{b}}.
$$
Therefore a point of $T^*\mathcal{B}$
over $\mathfrak{b}$ should be thought of
as $(x,\mathfrak{b})$, with
$x\in\mathfrak{n}_{\mathfrak{b}}$.

The zero section is not obtained by
fixing $\mathfrak{b}$ and allowing all
$x\in\mathfrak{n}_{\mathfrak{b}}$.
That would be the whole cotangent fibre.
The zero section is obtained by setting
the cotangent vector equal to zero:
$$
\mathcal{B}
\subset
T^*\mathcal{B},
\qquad
\mathfrak{b}\longmapsto (0,\mathfrak{b}).
$$
Applying $\mu$ gives
$$
\mu(0,\mathfrak{b})=0.
$$
So the zero section maps to the single
point $0\in\mathfrak{g}$.

On the other hand, the whole cotangent
fibre over $\mathfrak{b}$ maps to the
nilradical:
$$
T^*_{\mathfrak{b}}\mathcal{B}
\simeq
\mathfrak{n}_{\mathfrak{b}}
\xrightarrow{\mu}
\mathfrak{n}_{\mathfrak{b}}.
$$
For the standard Borel in
$\mathfrak{sl}_3$, this fibre contains
both $E_{12}$ and $E_{12}+E_{23}$.
They are nonzero cotangent vectors over
the same Borel. They are not points of
the zero section.

Thus the slogan is: the zero section maps
to zero,
$$
\mu(\mathcal{B})=0.
$$
Here $\mathcal{B}$ means the zero section
inside $T^*\mathcal{B}$. On the other
hand, the whole fibre over
$\mathfrak{b}$ maps to the nilradical:
$$
T^*_{\mathfrak{b}}\mathcal{B}
\xrightarrow{\mu}
\mathfrak{n}_{\mathfrak{b}}.
$$
The confusion comes from mixing up the
zero section with the whole fibre over a
fixed Borel.


% --- End springer.tex ---

% --- Begin seminars-others.tex ---
\chapter{Other Seminars}
\label{chapter-seminars-others}















\section{Neutrinos}
\label{seminars-others-section-neutrinos}

\noindent
Hiroshi Nunokawas, PUC-Rio.
Friday Seminar, Seminar Name.
June 27, 2025.

\medskip\noindent
{\bf Abstract.} Hiroshi will come and tell us
everything we (not) wanted to know
about these mysterious particles, and are not
going to be afraid to ask. In
particular, about the neutrino oscillation,
and the great matrices.

\bigskip\noindent

Protons and neutrons have very similar mass
of $m_p\approx 940$ MeV, while
electrons have mass of $m_e \approx 0.5$ MeV.
MeV is $10^{6}$ electronvolts,
where one eV is approximately $1.6\times
10^{-19}$ J. This is standard in high
energy physics, they use electronvolts
instead of Joules. Recall that
2J=1N$\times$1m.

Most of the things we see are protons since
they are so much larger than
electrons. But protons nor neutrons are
elementary particles.

Here's the standard model:

\medskip

\begin{tabular}{c c c c}
Quarks & $\begin{bmatrix} u\\d
\end{bmatrix}_L $ & $\begin{bmatrix} c\\s
\end{bmatrix}_L$ & $\begin{bmatrix} t\\b
\end{bmatrix}_L$\\
Leptons & $\begin{bmatrix} \nu_e\\e^-
\end{bmatrix}_L $ & $\begin{bmatrix}
\nu_\mu\\ \mu^-
\end{bmatrix}_L$ & $\begin{bmatrix}
\nu_\tau\\ \tau^- \end{bmatrix}_L$\\
Generation & 1st & 2nd & 3rd\\
Bosons & $g$, $\gamma$ , $\omega^\pm$, $z$ \\
Higgs Bosson & $H$
\end{tabular}

\medskip

It is very particular that nature repeats
itself three times. The $L$ in those
matrix actually means left-handed, and
accounts for chirality. Only left-handed
fermions have weak interaction. Right-handed
have electromagnetic interaction,
gravitational interaction, but not weak
interaction.

And then there's neutrinos. They have
negative helicity (chirality). Being
left-handed, mathematically, means to have
helicity $-1$. I think this means
that the spin is left-handed. But chirality
and helicity are not the same:
helicity is observer-dependent, and chirality
is not. Almost all neutrinos we
can see (\% 99.99999…) have negative
helicity, but not all of them.

Consider the following:
$$
n+\nu_e\leftrightarrow p+e^-
$$
But it's not completely correct: we'd better
put $d$ instead of $n$, and $u$
instead of $p$: the $d$ and $u$ quarks,
instead of the neutrons and protons.

Now consider the following reaction: a
neutron decays into a proton, an electron
and an antineutrino:
$$
n\to p+e^+\overline{\nu}_e
$$
Protons is very stable, that's why we are
here. But neutron decays in only 15
minutes.

By experimental data, we can conclude that
neutrinos' mass is consistent with
zero. But if they have mass, it should be
much smaller than the electron's $m_v
\leq 0.5$ eV. And the electron is already the
lightest fermion!

If the mass of the neutrino was zero, i.e.
$m_\nu=0$, then $v_\nu=c$ in vacuum,
which would imply that
$$
\xymatrix{
\nu_e\ar[r]^{L}&\nu_e\ar[r]&\nu_e\\
0:00&0:00&0:00
}
$$
meaning: time doesn't pass! And this means
the state of the particle cannot
change.

\section{Classification of admissible
W-algebras of low growth}
\label{seminars-others-section-w-algebras-of-low-growth}

\noindent
Igor Alarcon Blatt. June 10, 2026.
Algebra Seminar, IMPA.

\medskip\noindent
{\bf Abstract.} 
Admissible W-algebras form an important class
of vertex algebras, closely related to affine
Kac-Moody algebras and conformal field
theory. A natural invariant of such algebras
is (the asymptotic behaviour of) their
normalised characters. We classify all
admissible W-algebras whose normalised
characters have asymptotic growth less than
1. The proof relies on the computation of the
relevant data together with a theorem of Kac
and Wakimoto, leading to a characterisation
in terms of Virasoro minimal models.

\medskip\noindent
Goal: classify a certain class of W-algebras.
W-algebras are a kind of vertex algebra.

[A lot of content missing.
Some ideas are: minimal models
of Virasoro algebras,
assymptotic data]

\medskip\noindent
Recall that partitions of $n$
correspond with nilpotent orbits.

\begin{theorem}[Arakawa, 2015]
\label{seminars-others-theorem-arakawa}
For $k$ an admissible level with
denominator $u$, there is a unique
nilpotent orbit $\mathbb{O}_n$
suh that $W_k(\mathfrak{g},f)$
is rational for $f \in \overline{
\mathbb{O}_n}$. The orbits $\mathbb{O}_n$
correspond to the partitions listed below.
\end{theorem}

\medskip\noindent
Upshot: use a lemma 
(by Kac-Wakimoto, perhaps in
``Quantum reduction
for affine superalgebras'') which asserts
that if the assymptotic growth
satisfies $>1$ then there is a direct
sum expression.



\section{Vertex Algebras and the Monster}
\label{seminars-others-section-vertex-algebras-monster}

\noindent
David Melo. September 18, 2026.
Students' Colloquium. Room 228.
Friday, 15:20--16:20.

\medskip\noindent
{\bf Abstract.}
The classification of finite simple groups is
one of the most remarkable achievements in
algebra. Its proof spans about 500 articles,
involves about 100 mathematicians, and fills
more than 10,000 pages. The work surrounding
the classification also uncovered exceptional
structures connecting several areas of
mathematics.

A particularly striking example is the
relationship among the Monster group, sphere
packings, modular forms, and conformal field
theory through vertex algebras. The Leech
lattice gives the optimal sphere packing in
dimension $24$. A conformal field theory is
constructed from this lattice. Borcherds's
proof of monstrous moonshine explains
remarkable coincidences between the
Monster and Ramanujan's $j$-function. This
talk will introduce the exceptional objects
and outline their relationship.

\medskip\noindent
{\bf Introduction.}
The Monster is the largest sporadic finite
simple group. Its order is
$$
\begin{aligned}
|\mathbb{M}|
&=2^{46}3^{20}5^9 7^6 11^2 13^3 17\\
&\phantom{={}}\cdot19\cdot23\cdot29\cdot31
\cdot41\cdot47\cdot59\cdot71.
\end{aligned}
$$
Some of the coefficients in the Fourier
expansion of Ramanujan's $j$-function have
unexpected connections with dimensions of
representations of the Monster. Monstrous
moonshine explains these numerical
coincidences.

The Monster is the automorphism group of the
moonshine vertex algebra.

\medskip\noindent
{\bf Applications of the modular group.}
The modular group has applications in the
following areas.

\begin{enumerate}
\item
It parametrizes complex structures on a
torus.

\item
It occurs in Diophantine equations, including
the proof of Fermat's last theorem.

\item
It is used in cryptography.

\item
It helps compute elliptic integrals.

\item
It governs Fourier expansions of modular
functions and modular forms.
\end{enumerate}

\medskip\noindent
{\bf Modular invariants.}
The modular group is
$$
\operatorname{SL}_2(\mathbb{Z}).
$$
Its action on the upper half-plane comes from
translation and inversion.
$$
\tau\longmapsto\tau+1,
\qquad
\tau\longmapsto-\frac{1}{\tau}.
$$

A modular function is a meromorphic invariant
function for this action. A modular form of
weight $k$ transforms by
$$
f\left(\frac{a\tau+b}{c\tau+d}\right)
=(c\tau+d)^k f(\tau).
$$
The matrix with entries $a,b,c,d$ belongs
to the modular group.

Examples include Eisenstein series, theta
functions, and the Dedekind eta function.
Their coefficients often count solutions of
arithmetic equations. They were a key
ingredient in the proof of Fermat's last
theorem.

The $j$-function is the simplest nonconstant
modular function. Its Fourier expansion is
$$
j(\tau)=q^{-1}+744+196884q
+21493760q^2+\cdots.
$$

\medskip\noindent
{\bf Classification of simple groups.}
Let $G$ be a finite simple group. Then $G$ is
one of the following.

\begin{enumerate}
\item
A cyclic group of prime order.

\item
An alternating group.

\item
A group of Lie type.

\item
One of the $26$ sporadic groups.
\end{enumerate}

The periodic table of finite simple groups
displays the sporadic groups and their
relationships.

$$
\xymatrix{
& \mathbb{M}\ar[dl]\ar[d]\ar[dr] & \\
\mathbb{B} &
\mathrm{Fi}_{24}' &
\mathrm{Co}_1 \\
& \mathrm{Fi}_{23}\ar[ur] &
\mathrm{Co}_2\ar[u]
}
$$

\medskip\noindent
{\bf Lattices.}
An even lattice is a discrete subgroup
$L\subset\mathbb{R}^n$ such that
$$
\left\langle{} x,x \right\rangle{}
\in2\mathbb{Z}
$$
for every $x\in L$. Its dual lattice consists
of vectors in $\mathbb{R}^n$ whose pairings
with every element of $L$ are integral.

The roots of an even lattice are the vectors
of norm $2$. They have the smallest possible
positive norm.

\medskip\noindent
{\bf The Leech lattice.}
The Leech lattice comes from coding theory.
It gives the optimal sphere packing in
dimension $24$, as was proved in 2016. It is
the unique even unimodular lattice of rank
$24$ without roots.

The identity
$$
1^2+2^2+\cdots+24^2=70^2
$$
is another striking occurrence of the number
$24$.

Show a slide with a matrix whose rows form a
basis of the Leech lattice.

\medskip\noindent
{\bf Conway's group.}
The Conway group is
$$
\mathrm{Co}_1=
\operatorname{Aut}(\Lambda)/\{1,-1\}.
$$
It is a sporadic simple group. The Monster is
not obtained from $\mathrm{Co}_1$ by
adjoining one element. Rather,
$\mathrm{Co}_1$ occurs as a quotient of a
subgroup of the Monster.

\medskip\noindent
{\bf Group representations.}
{\it Definition.} A representation of $G$ on
$V$ is a homomorphism
$$
\rho:G\longrightarrow\operatorname{GL}(V).
$$
A morphism of representations is a linear map
$f:V\longrightarrow W$ such that
$$
f\rho(g)=\sigma(g)f
$$
for every $g\in G$.

Examples include the following.

\begin{enumerate}
\item
The group $\operatorname{Aut}(\Lambda)$ acts
on $\mathbb{R}^{24}$. The quotient
$\mathrm{Co}_1$ acts projectively.

\item
The Monster has an irreducible representation
of dimension $196883$.

\item
The symmetric group $S_n$ acts on
$\mathbb{R}^n$ by permuting coordinates.

\item
The circle has one-dimensional
representations indexed by $n\in\mathbb{Z}$.
$$
e^{i\theta}\longmapsto e^{in\theta}.
$$

\item
The groups $A_4$ and $A_5$ act on
$\mathbb{R}^3$ as rotation groups.
\end{enumerate}

\medskip\noindent
{\bf Putting a ring on it.}
For a group $G$, the group algebra
$\mathbb{k}[G]$ is the free vector space on
$G$, with multiplication extending the group
operation. Representations of $G$ are
equivalent to $\mathbb{k}[G]$-modules.

\medskip\noindent
{\bf Reducibility.}
{\it Definition.} A representation is
reducible if it has a nonzero proper
subrepresentation. In a semisimple category,
it is then a direct sum of two nonzero
representations. Otherwise, it is
irreducible.

\begin{theorem}
\label{seminars-others-theorem-finite-group-representations}
Let $G$ be finite and let $\mathbb{k}$ have
characteristic zero. There are finitely many
irreducible representations $V_j$, and every
finite-dimensional representation is a direct
sum of such representations.
\end{theorem}

\begin{theorem}[Peter--Weyl]
\label{seminars-others-theorem-peter-weyl}
Let $G$ be a compact topological group. Then
$$
L^2(G)=\bigoplus_{j\in\Sigma}
V_j\otimes V_j^*
$$
as an orthogonal decomposition of Hilbert
spaces.
\end{theorem}

For $G=S^1$, this gives
$$
L^2(S^1)=\bigoplus_{n\in\mathbb{Z}}
\mathbb{C}e^{in\theta}.
$$
The number of irreducible complex
representations of a finite group equals its
number of conjugacy classes.

\medskip\noindent
{\bf The Virasoro Lie algebra.}
{\it Definition.} A Lie algebra is a vector
space with a bilinear bracket satisfying
antisymmetry and the Jacobi identity.

On $\mathbb{C}^{\times}$, put
$$
L_n=-z^{n+1}\frac{d}{dz}.
$$
These vector fields satisfy
$$
[L_m,L_n]=(m-n)L_{m+n}.
$$
Their unique nontrivial central extension has
relations
$$
[L_m,L_n]=(m-n)L_{m+n}
+\frac{m^3-m}{12}\delta_{m+n,0}C.
$$
This is the Virasoro algebra. It is the
quantum analogue of local conformal
transformations and is central in string
theory. The central extension accounts for
normal-ordering anomalies.

\medskip\noindent
{\bf Positive-energy representations.}
{\it Definition.} A positive-energy
representation of the Virasoro algebra is a
vector space $V$. The element $C$ acts by a
scalar $c$, $L_0$ acts diagonally with
spectrum bounded below, and every eigenspace
of $L_0$ has
finite dimension.

Its renormalized character is
$$
\chi_V(q)=\sum_{h\in\mathbb{R}}
\dim(V^{L_0=h})q^{h-c/24}.
$$

\medskip\noindent
{\bf Vertex algebras.}
{\it Definition.} A conformal vertex algebra
consists of a vertex algebra $V$, a vacuum
vector $|0\rangle$, and a Virasoro vector $L$
with
$$
T=L_{-1}.
$$
The vector $L$ is the energy-momentum tensor.

\medskip\noindent
{\bf Borcherds's identity.}
For $a,b,c\in V$ and $m,n,k\in\mathbb{Z}$,
the identity is
$$
\begin{aligned}
&\sum_{j\geq0}\binom{m}{j}
(a_{(n+j)}b)_{(m+k-j)}c\\
&=\sum_{j\geq0}(-1)^j\binom{n}{j}
\left(a_{(m+n-j)}b_{(k+j)}c\right.\\
&\qquad\left.{}-(-1)^n
b_{(n+k-j)}a_{(m+j)}c\right).
\end{aligned}
$$

\medskip\noindent
{\bf Examples and applications.}
For a simple Lie algebra $\mathfrak{g}$ and a
complex number $k$, the universal affine
vertex algebra is $V^k(\mathfrak{g})$. Its
weight-one subspace is naturally
$\mathfrak{g}$. If $f\in\mathfrak{g}$ is
nilpotent, quantum Hamiltonian reduction
gives the vertex algebra
$W^k(\mathfrak{g},f)$. An even lattice $L$
gives a lattice vertex algebra
$V_L$.

Vertex algebras have applications to knot
theory and braid relations, integrable
partial differential equations in
symplectic geometry,
$D$-modules on Riemann surfaces,
combinatorial product-sum identities, and
theoretical physics.

\medskip\noindent
{\bf Bosonic strings.}
Let
$$
\mathcal{H}=\mathbb{C}[x_{-1},x_{-2},\ldots]
$$
and let $|0\rangle=1$. The Heisenberg modes
act by
$$
\begin{aligned}
(x_{-m})_{(n)}P&=x_{-m}P\qquad(n=-m),\\
(x_{-m})_{(n)}P&=n
\frac{\partial P}{\partial x_{-n}}
\qquad(n>0).
\end{aligned}
$$
For the standard conformal vector, the
central charge is $1$ and
$$
L_0x_{-k}=kx_{-k}.
$$
The character is
$$
\chi_{\mathcal{H}}(q)=
\prod_{m\geq1}\frac{1}{1-q^m}.
$$

\medskip\noindent
{\bf Lattice vertex algebras.}
For an even lattice $L$, the lattice vertex
algebra $V_L$ combines the Heisenberg algebra
with a copy of the group algebra of $L$.
Show a picture of the lattice, with a
Heisenberg cone at each lattice point.

\medskip\noindent
{\bf Moonshine.}
Starting with the Leech lattice vertex
algebra $V_\Lambda$, orbifolding by the
involution
induced by $-1$ gives the moonshine module
$$
V^{\natural}=V_\Lambda^+
\oplus V_\Lambda^{T,+}.
$$
The second summand comes from a twisted
module.
The automorphism group of $V^{\natural}$ is
the Monster.

\medskip\noindent
{\bf Extra ingredients.}
The no-ghost theorem states that critical
bosonic string theory in dimension $26$ has
no negative-norm physical states. Borcherds's
proof constructs a Lie algebra from a
semi-infinite cohomology group for
$$
V^{\natural}\otimes V_{II_{1,1}}.
$$
This construction is closely related to
current research on semi-infinite cohomology.



% --- End seminars-others.tex ---

\part{Smooth Geometry and Topology}
\label{part-smooth-geometry-and-topology}

% --- Begin smooth-manifolds.tex ---
\chapter{Smooth Manifolds}
\label{chapter-smooth-manifolds}




\section{Topological manifolds}
\label{smooth-manifolds-section-topological-manifolds}

\begin{definition}
\label{smooth-manifolds-definition-topological-manifold}
A {\it topological manifold} is a topological
space which is locally homeomorphic to an
open ball in $\mathbb{R}^n$.
\end{definition}

\medskip\noindent
Manifolds can be smooth (of a given
``differentiability class''),
real analytic, or topological (continuous).

\begin{definition}
\label{smooth-manifolds-definition-group-action-topological-manifold}
An action of a group on a manifold is
silently assumed to be continuous. Let $G$
be a group acting on a set $M$. The
\emph{stabilizer} of $x \in M$ is the
subgroup of all elements in $G$ that fix
$x$. An action is \emph{free} if the
stabilizer of every point is trivial. The
\emph{quotient space} $M/G$ is the space of
orbits, equipped with the following
topology: an open set $U \subset M/G$ is
open if its preimage in $M$ is open.
\end{definition}

\begin{exercise}
\label{smooth-manifolds-exercise-finite-free-quotient-manifold}
Let $G$ be a finite group acting freely on a
Hausdorff manifold $M$. Show that the
quotient space $M/G$ is a topological
manifold.
\end{exercise}

\begin{proof}
%
\end{proof}

\begin{exercise}
\label{smooth-manifolds-exercise-finite-nonfree-quotient-not-manifold}
Construct an example of a finite group $G$
acting non-freely on a topological manifold
$M$ such that $M/G$ is not a topological
manifold.
\end{exercise}

\begin{proof}
Consider the manifold $\mathbb{R}^2$ and a
group $G$ generated by a rotation about the
origin. This makes the origin the only fixed
point of $G$. Any other point has a
neighbourhood in which it is the only
element in its orbit; such a neighbourhood
is mapped diffeomorphically onto the
quotient, giving a coordinate chart.

However, any neighbourhood of the origin
contains at least two $G$-related elements.
Suppose there is a coordinate chart of
$\bar{0}$ in the quotient. The preimage of
this open set is an open set in
$\mathbb{R}^2$ containing $0$. The
composition of the quotient map and the
coordinate chart produce a non-injective map
from $\mathbb{R}^2$ to $\mathbb{R}^2$. If
such a map was smooth at $0$, by the inverse
function theorem there would be a
neighbourhood of $0$ in which the map would
be invertible, a contradiction.
\end{proof}

\begin{exercise}
\label{smooth-manifolds-exercise-quotient-by-minus-one}
Consider the quotient of $\mathbb{R}^2$ by
the action of $\{\pm 1\}$ that maps $x$ to
$-x$. Is the quotient space a topological
manifold?
\end{exercise}

\begin{proof}
Yes. A homeomorphism between
$\mathbb{R}^2/\{\pm 1\}$ and $\mathbb{R}^2$
is $re^{i\theta} \mapsto re^{i2\theta}$,
where the angle $\theta$ in the domain is
in the interval $[0,\pi)$ and
$r \in [0,\infty)$. This map is clearly
bijective, and its continuous inverse is
$re^{i\varphi} \mapsto re^{i\varphi/2}$ for
$\varphi \in [0,2\pi)$.
\end{proof}

\begin{exercise}
\label{smooth-manifolds-exercise-homeomorphism-group-transitive}
Let $M$ be path connected, Hausdorff
topological manifold and $G$ a group of all
its homeomorphisms. Prove that $G$ acts
transitively.
\end{exercise}

\begin{exercise}
\label{smooth-manifolds-exercise-closed-subgroup-manifold}
Prove that any closed subgroup
$G \subset \mathsf{GL}(n)$ of a matrix
group is homeomorphic to a manifold, or find
a counterexample.
\end{exercise}

\medskip\noindent
In the above definition of a manifold, it is
not required to be Hausdorff. Nevertheless,
in most cases, manifolds are tacitly assumed
to be Hausdorff.

\begin{exercise}
\label{smooth-manifolds-exercise-nonhausdorff-manifold}
Construct an example of a non-Hausdorff
manifold.
\end{exercise}

\begin{exercise}
\label{smooth-manifolds-exercise-r2-quotient-z2-manifold}
Show that $\mathbb{R}^2/\mathbb{Z}^2$ is a
manifold.
\end{exercise}

\begin{proof}
Let $\bar{z}$ be a point in
$\mathbb{R}^2/\mathbb{Z}^2$. Its preimage is
the lattice $\{z+(a,b):a,b \in \mathbb{Z}\}$.
A ball of radius $\frac{1}{2}$ centered at
any representative of $\bar{z}$ contains
only one representative of any other class,
so that the restriction of the projection is
bijective (and continuous by definition of
quotient topology). The inverse map is also
continuous by definition of quotient
topology: an open set in the ball on
$\mathbb{R}^2$ is mapped to an open set in
the quotient because its preimage is open.
\end{proof}

\begin{exercise}
\label{smooth-manifolds-exercise-irrational-action-not-manifold}
Let $\alpha$ be an irrational number. The
group $\mathbb{Z}^2$ acts on $\mathbb{R}$ by
the formula $t \mapsto t + m + n\alpha$.
Show that this action is free, but the
quotient $\mathbb{R}/\mathbb{Z}^2$ is not a
manifold.
\end{exercise}

\begin{proof}
This action is obviously free since any
nonzero pair of integers moves any number
$t$. The idea is to show that the orbit of
every point is dense. This would mean that
any open set in the quotient is homeomorphic
to $\mathbb{R}$, preventing it from having
local neighbourhoods homeomorphic to balls,
i.e. being a topological manifold.
\end{proof}

\section{Differentiable maps}
\label{smooth-manifolds-section-differentiable-maps}

\begin{definition}
\label{smooth-manifolds-definition-affine-map}
Let $U, V \subset \mathbb{R}^n$ be open
subsets. An affine map is a sum of linear
map $\alpha$ and a constant map. Its linear
part is $\alpha$.
\end{definition}

\begin{definition}
\label{smooth-manifolds-definition-differentiable-map}
Let $U \subset \mathbb{R}^m$, $V \subset
\mathbb{R}^n$ be open subsets. A map
$f:U \to V$ is called differentiable if it
can be approximated by an affine one at any
point: that is, for any $x \in U$, there
exists an affine map $\varphi_x:\mathbb{R}^m
\to \mathbb{R}^n$ such that
$$
\lim_{x_1 \to x}
\frac{|f(x_1)-\varphi_x(x_1)|}{|x-x_1|}=0
$$
\end{definition}

\begin{definition}
\label{smooth-manifolds-definition-differential}
Differential, or derivative, of a
differentiable map $f:U \to V$ is the linear
part of $\varphi_x$.
\end{definition}

\begin{definition}
\label{smooth-manifolds-definition-diffeomorphism}
Diffeomorphism is a differentiable map $f$
which is invertible, and such that $f^{-1}$
is also differentiable. A map
$f:U \to V$ is a local diffeomorphism if each
point $x \in U$ has an open neighbourhood
$U_1 \ni x$ such that
$f:U_1 \to f(U_1)$ is a diffeomorphism.
\end{definition}

\medskip\noindent
Chain rule says that a composition of two
differentiable functions is differentiable,
and its differential is composition of their
differentials.

\medskip\noindent
Chain rule implies that differential of a
diffeomorphism is invertible. Converse is
also true (in a sense):

\section{Inverse function theorem}
\label{smooth-manifolds-section-inverse-function-theorem}

\begin{theorem}
\label{smooth-manifolds-theorem-inverse-function}
Let $U, V \subset \mathbb{R}^n$ be open
subsets, and $f:U \to V$ a differentiable
map. Suppose that the differential of $f$ is
everywhere invertible. Then $f$ is locally a
diffeomorphism.
\end{theorem}

\begin{proof}
Step 1: Let $x \in U$. Without restricting
generality, we may assume that $x=0$,
$U=B_r(0)$ is an open ball of radius $r$,
and in $U$ one has
$$
\frac{|f(x_1)-\varphi_x(x_1)|}{|x-x_1|}<1/2.
$$
Replacing $f$ with $-f \circ (D_0f)^{-1}$,
where $D_0f$ is differential of $f$ in $0$,
we may assume also that $D_0f=-\mathrm{Id}$.

Step 2: In these assumptions,
$$
|f(x)+x|<1/2|x|,
$$
hence $\psi_s(x) := f(x)+x-s$ is a
contraction. This map maps $B_{r/2}(0)$ to
itself when $s < r/4$. By Banach fixed point
theorem, $\psi_s(x)=x$ has a unique fixed
point $x_s$, which is obtained as a solution
of the equation $f(x)+x-s=x$, or,
equivalently, $f(x)=s$. Denote the map
$s \mapsto x_s$ by $g$.

Step 3: By construction, $fg=\mathrm{Id}$. Applying
the chain rule again, we find that $g$ is
also differentiable.
\end{proof}

\medskip\noindent
Usually, diffeomorphisms are assumed smooth
(infinitely differentiable). A smooth
version of this result is left as an
exercise.

\section{Critical points and critical values}
\label{smooth-manifolds-section-critical-points-and-critical-values}

\begin{definition}
\label{smooth-manifolds-definition-critical-point-value}
Let $U \subset \mathbb{R}^m$, $V \subset
\mathbb{R}^n$ be open subsets, and $f:U \to
V$ a smooth function. A point $x \in U$ is a
critical point of $f$ if the differential
$D_xf:\mathbb{R}^m \to \mathbb{R}^n$ is not
surjective. Critical value is an image of a
critical point. Regular value is a point of
$V$ which is not a critical value.
\end{definition}

\begin{theorem}
\label{smooth-manifolds-theorem-sard-lemma}
(Sard's lemma)
The set of critical values of $f$ is of
measure $0$ in $V$.
\end{theorem}

\medskip\noindent
We leave this theorem without a proof. We
won't use it.

\begin{definition}
\label{smooth-manifolds-definition-smooth-submanifold}
A subset $M \subset \mathbb{R}^n$ is an
$m$-dimensional smooth submanifold if for
each $x \in M$ there exists a neighbourhood
$U \subset \mathbb{R}^n$ containing $x$ and a
diffeomorphism from $U$ to an open ball $B
\subset \mathbb{R}^n$ which maps $U \cap M$
to an intersection $B \cap \mathbb{R}^m$ of
$B$ and an $m$-dimensional linear subspace.
\end{definition}

\medskip\noindent
Clearly, a smooth submanifold is a
(topological) manifold.

\begin{theorem}
\label{smooth-manifolds-theorem-regular-fiber-submanifold}
Let $U \subset \mathbb{R}^m$, $V \subset
\mathbb{R}^n$ be open subsets,
$f:U \to V$ a smooth function, and $y \in V$
a regular value of $f$. Then $f^{-1}(y)$ is
a smooth submanifold of $U$.
\end{theorem}

\begin{proof}
Let $x \in U$ be a point in $f^{-1}(y)$.
It suffices to prove that $x$ has a
neighbourhood diffeomorphic to an open ball
$B$, such that $f^{-1}(y)$ corresponds to a
linear subspace in $B$. Without restricting
generality, we may assume that $y=0$ and
$x=0$.

The differential $D_0f:\mathbb{R}^n \to
\mathbb{R}^m$ is surjective. Let
$L:=\ker D_0f$, and let
$A:\mathbb{R}^n \to L$ be any map which acts
on $L$ as identity. Then
$D_0f \oplus A:\mathbb{R}^n \to
\mathbb{R}^m \oplus L$ is an isomorphism of
vector spaces. Therefore, $\Psi: f \oplus A$
mapping $x_1$ to $f(x_1) \oplus A(x_1)$ is a
diffeomorphism in a neighbourhood of $x$.
However, $f^{-1}(0)=\Psi^{-1}(0 \oplus L)$.
We have constructed a diffeomorphism of a
neighbourhood of $x$ with an open ball
mapping $f^{-1}(0)$ to $0 \oplus L$.
\end{proof}

\begin{lemma}[Corollary]
\label{smooth-manifolds-lemma-regular-maps-solutions}
Let $f_1, \dots, f_m$ be smooth functions on
$U \subset \mathbb{R}^n$ such that the
differentials $df_i$ are linearly
independent everywhere. Then the set of
solutions of equations
$f_1(z)=f_2(z)=\dots=f_m(z)=0$ is a smooth
$(n-m)$-dimensional submanifold in $U$.
\end{lemma}

\begin{definition}
\label{smooth-manifolds-definition-smooth-hypersurface}
A {\it smooth hypersurface} is a closed
codimension $1$ submanifold.
\end{definition}

\begin{exercise}
\label{smooth-manifolds-exercise-smooth-hypersurface-regular-value}
Prove that a smooth hypersurface in $U$ is
always obtained as a solution of an equation
$f(z)=0$, where $0$ is a regular value of a
function $f:U \to \mathbb{R}$.
\end{exercise}

\section{Smooth manifolds}
\label{smooth-manifolds-section-smooth-manifolds}

\begin{definition}
\label{smooth-manifolds-definition-cover-refinement}
An open cover of a topological space $X$ is a
family of open sets $\{U_i\}$ such that
$\bigcup_i U_i=X$. A cover $\{V_i\}$ is a
refinement of a cover $\{U_i\}$ if every
$V_i$ is contained in some $U_i$.
\end{definition}

\medskip\noindent
Any two covers $\{U_i\}$, $\{V_i\}$ of a
topological space admit a common refinement.

\begin{definition}
\label{smooth-manifolds-definition-atlas}
Let $M$ be a topological manifold. A cover
$\{U_i\}$ of $M$ is an atlas if for every
$U_i$ we have a map $\varphi_i:U_i \to
\mathbb{R}^n$ giving a homeomorphism of
$U_i$ with an open subset in $\mathbb{R}^n$.
The transition maps
$$
\varphi_{ij}:\varphi_i(U_i \cap U_j) \to
\varphi_j(U_i \cap U_j)
$$
are induced by the above homeomorphisms.
\end{definition}

\begin{definition}
\label{smooth-manifolds-definition-differentiability-class}
A function $\mathbb{R} \to \mathbb{R}$ is of
differentiability class $C^i$ if it is
$i$ times differentiable, and its $i$-th
derivative is continuous. A map
$\mathbb{R}^n \to \mathbb{R}^m$ is of
differentiability class $C^i$ if all its
coordinate components are. A smooth
function/map is a function/map of class
$C^\infty$.
\end{definition}

\begin{definition}
\label{smooth-manifolds-definition-smooth-atlas}
An atlas is smooth if all transition maps
are smooth (of class $C^\infty$, i.e.,
infinitely differentiable), smooth of class
$C^i$ if all transition functions are of
differentiability class $C^i$, and real
analytic if all transition maps admit a
Taylor expansion at each point.
\end{definition}

\begin{definition}
\label{smooth-manifolds-definition-smooth-structure}
A smooth structure on a manifold (of class
$C^\infty$ or $C^i$) is an atlas of class
$C^\infty$ or $C^i$ considered up to the
above equivalence. A smooth manifold is a
topological manifold equipped with a smooth
structure.
\end{definition}

\medskip\noindent
Terrible, isn't it?

\begin{definition}
\label{smooth-manifolds-definition-smooth-function-manifold}
A smooth function on a manifold $M$ is a
function $f$ whose restriction to a chart
$(U_i,\varphi_i)$ gives a smooth function
$f \circ \varphi_i^{-1}:\varphi_i(U_i)\to
\mathbb{R}$ for each open subset
$\varphi_i(U_i)\subset \mathbb{R}^n$.
\end{definition}

\medskip\noindent
There are several ways to define a smooth
manifold. The above way is most standard.
It is not the most convenient one but you
should know it. Two other ways, via sheaves
of functions and via Whitney's theorem, are
presented later.

\begin{definition}
\label{smooth-manifolds-definition-smooth-embedded-submanifold}
Let $M$ be a smooth manifold. A subset
$N \subset M$ is called a smoothly embedded
submanifold of dimension $m$ if for every
point $x \in N$ there is a neighbourhood
$U \subset M$ diffeomorphic to an open ball
$B \subset \mathbb{R}^n$, such that this
diffeomorphism maps $U \cap N$ onto a linear
subspace of $B$ of dimension $m$.
\end{definition}

\begin{exercise}
\label{smooth-manifolds-exercise-extend-sheaf-on-submanifold}
Let $(M,\mathcal{F})$ be a smooth manifold
and let $N \subset M$ be a smoothly embedded
submanifold. Consider the space
$\mathcal{F}'(U)$ of smooth functions on
$U \subset N$ that are extendable to
functions on $M$ defined on some
neighbourhood of $U$.
\begin{enumerate}
\item Show that $\mathcal{F}'$ is a sheaf.
\item Show that this sheaf defines a smooth
structure on $N$.
\item Show that the natural embedding
$(N,\mathcal{F}') \to (M,\mathcal{F})$ is a
morphism of manifolds.
\end{enumerate}
\end{exercise}

\begin{proof}
\begin{enumerate}
\item To see that $\mathcal{F}'$ is a
presheaf, fix an open set $U \subset N$ and
a function $f \in \mathcal{F}'(U)$. This
means that $f$ can be extended to a
function $\tilde{f}$ on $M$ defined on some
neighbourhood of $U$. Then the restriction
of $f$ to any open subset $U_1 \subset U$
can be extended to the same function
$\tilde{f}$ on $M$ defined on the same
neighbourhood of $U$, which is also a
neighbourhood of $U_1$. This says that
$f|_{U_1}\in \mathcal{F}'(U_1)$.

To check that $\mathcal{F}'$ is a sheaf,
consider a cover $\{U_i\}$ of $U$ and choose
$f_i \in \mathcal{F}'(U_i)$ for all $i$
satisfying
$$
f_i|_{U_i \cap U_j}=f_j|_{U_i \cap U_j},
\qquad \forall i,j.
$$
This means that every $f_i$ can be extended
to a function $\tilde{f}_i$ on $M$ defined
on some neighbourhood $\tilde{U}_i \subset
M$ of $U_i$. Consider
$\tilde{U}=\bigcup_i \tilde{U}_i$; we must
construct a smooth function on all of
$\tilde{U}$ from the $\tilde{f}_i$.

Suppose there is a partition of unity
$\{\nu_i\}$ subordinate to the cover
$\{\tilde{U}_i\}$. Then each
$\tilde{f}_i \nu_i$ is a smooth function
defined on $\tilde{U}$, and so is the
function $F=\sum_i \tilde{f}_i \nu_i$.

To conclude, we must show that the
restriction of $F$ to any $U_j$ coincides
with $f_j$. Let $x \in U_j$ for some $j$.
Then
$$
F(x)=\sum_i \tilde{f}_i(x)\tilde{\nu}_i(x)
=\sum_i f_i(x)\tilde{\nu}_i(x)
=f_j(x)\sum_i \tilde{\nu}_i(x)=f_j(x)
$$
since the original functions $f_i$ coincide
in the intersections.

\item Suppose that $N$ is a smoothly
embedded submanifold of dimension $m$.
Since $N$ is a smoothly embedded
submanifold, at every point of $N$ there is
a neighbourhood $U$ of $M$ homeomorphic to
a ball $B$ in $\mathbb{R}^n$ such that
$U \cap N$ is mapped to a linear subspace of
$B$. Since $M$ is a smooth manifold we may
suppose, restricting to a smaller open set
if necessary, that the same $(U,\mathcal{F})$
is isomorphic to $(\mathbb{R}^n,\mathcal{F}'')$.

Let $V:=\varphi(U \cap N)$ and
$f'' \in \mathcal{F}''(V)$. Then $f''$ may
be smoothly extended to a function on
$\varphi(U)\cong \mathbb{R}^m$: define an
extension $\tilde{f}''(x,y)=f''(x)$; then
the partial derivatives with respect to the
new variables vanish. Then $\tilde{f}''$
corresponds to a smooth function on $U$ by
the isomorphism $(U,\mathcal{F})\cong
(\mathbb{R}^m,\mathcal{F}'')$. This shows
that the function $f''$ corresponds to a
function on $U \cap N$ that may be extended
to a neighbourhood of $M$, meaning that it
is an element of $\mathcal{F}'(U \cap N)$.

Conversely, a function
$f' \in \mathcal{F}'(U \cap N)$ may be
smoothly extended to a function on some open
set of $M$ by definition. Intersecting such
a set with $U$ and restricting smooth
functions, we may suppose it is isomorphic
to $(\mathbb{R}^n,\mathcal{F}'')$. Then
$\varphi$ maps the extension of $f'$ to a
smooth function on $\mathbb{R}^n$, whose
restriction to $V$ is an element of
$\mathcal{F}''(V)$.

\item To check that the natural embedding
$(N,\mathcal{F}') \xrightarrow{\Psi}
(M,\mathcal{F})$ is a morphism of manifolds,
we must check that it is a continuous map
satisfying $f \circ \Psi \in
\mathcal{F}'(\Psi^{-1}(U))$ for any open set
$U \subset M$ and $f \in \mathcal{F}(U)$.

Continuity is immediate since $N$ is
equipped with the subspace topology. The
second condition is also immediate by
definition of $\mathcal{F}'$.
\end{enumerate}
\end{proof}

\begin{exercise}
\label{smooth-manifolds-exercise-same-image-embedded-submanifolds}
Let $N_1,N_2$ be two manifolds and let
$\varphi_i:N_i \to M$ be smooth embeddings.
Suppose that the image of $N_1$ coincides
with that of $N_2$. Show that $N_1$ and
$N_2$ are isomorphic.
\end{exercise}

\begin{proof}
The smooth structures of $N_i$ are given by
the previous exercise: they are the sheaves
$\mathcal{F}'_i$, so that
$(N_i,\mathcal{F}'_i)$ are locally
isomorphic to $(\mathbb{R}^{n_i},\mathcal{F}'')$
for some $n_i$. Notice that $n_1=n_2$
because the image of these embeddings in $M$
coincides. That means there is a
neighbourhood $U \subset M$ diffeomorphic to
a ball $B$ in $\mathbb{R}^m$ such that
$N_i \cap U$ is mapped to the same linear
subspace of $B$ of dimension
$n_1=n_2=:n$.

Local homeomorphisms between $N_1$ and $N_2$
may be obtained by composing the local
homeomorphisms with $\mathbb{R}^n$ given by
each smooth structure:

$$
\xymatrix{
V_1 \subset N_1 \ar@{-->}[r] \ar[d]
& N_2 \supset V_2 \ar[d] \\
\mathbb{R}^n \ar[r]
& \mathbb{R}^n
}
$$

The fact that these local homeomorphisms
map smooth functions to smooth functions
follows from the fact that they are
compositions of diffeomorphisms and that the
dimensions coincide.
\end{proof}

\medskip\noindent
By the above problem, in order to define a
smooth structure on $N$, it suffices to
embed $N$ into $\mathbb{R}^n$. As it will be
clear later, every manifold is embeddable
into $\mathbb{R}^n$ (assuming it admits a
partition of unity). Therefore, in place of
a smooth manifold, we can use manifolds
that are smoothly embedded into
$\mathbb{R}^n$.

\section{Covering maps}
\label{smooth-manifolds-section-covering-maps}

\begin{definition}
\label{smooth-manifolds-definition-covering-map}
Let $\phi:\tilde{M}\to M$ be a continuous
map of manifolds (or CW complexes). We say
that $\phi$ is a covering if $\phi$ is
locally a homeomorphism, and for any
$x \in M$ there exists a neighbourhood
$U \ni x$ such that $\phi^{-1}(U)$ is a
disconnected union of several manifolds
$U_i$ such that the restriction
$\phi|_{U_i}$ is a homeomorphism from
$U_i$ to $U$.
\end{definition}

\begin{definition}
\label{smooth-manifolds-definition-properly-discontinuous}
Let $\Gamma$ be a discrete group
continuously acting on a locally compact
topological space $M$. This action is called
properly discontinuous if the space of
orbits of $\Gamma$ is Hausdorff.
\end{definition}

\begin{lemma}
\label{smooth-manifolds-lemma-local-diffeomorphism-finite-fibers-covering}
If a local diffeomorphism has the same
finite number of preimages of any point,
then it is a covering.
\end{lemma}

\begin{theorem}
\label{smooth-manifolds-theorem-properly-discontinuous-covering}
Let $\Gamma$ be a discrete group acting on
$M$ freely and properly discontinuously.
Then $M \to M/\Gamma$ is a covering.
\end{theorem}

\medskip\noindent
This result is left as an exercise.

\medskip\noindent
\begin{example}
\label{smooth-manifolds-example-circle-covering}
The map $x \mapsto nx$ on a circle $S^1$ is
a covering.
\end{example}

\begin{example}
\label{smooth-manifolds-example-torus-covering}
For any non-degenerate integer matrix
$A \in \operatorname{End}(\mathbb{Z}^n)$,
the corresponding map of a torus $T^n$ is a
covering.
\end{example}

\begin{lemma}
\label{smooth-manifolds-lemma-covering-constant-fiber-size}
Let $\phi:\tilde{M}\to M$ be a covering,
with $M$ connected. Then the number of
preimages $|\phi^{-1}(m)|$ is constant in
$M$.
\end{lemma}

\begin{proof}
Since $\phi^{-1}(U)$ is a disconnected union
of several copies of $U$, this number is a
locally constant function of $m$.
\end{proof}

\begin{definition}
\label{smooth-manifolds-definition-degree-covering}
Let $\phi:\tilde{M}\to M$ be a covering,
with $M$ connected. Assume that $\tilde{M}$
and $M$ are oriented and that $\phi$
preserves the orientation. The number
$|\phi^{-1}(m)|$ is called the degree of
$\phi$.
\end{definition}

\begin{lemma}
\label{smooth-manifolds-lemma-compact-covering-finite-degree}
Any covering $\phi:\tilde{M}\to M$ with
$\tilde{M}$ compact has finite degree.
\end{lemma}

\begin{proof}
Take $U$ in such a way that
$\phi^{-1}(U)$ is a disconnected union of
several copies of $U$, and let $x \in U$.
Then $\phi^{-1}(x)$ is discrete, and since
$\tilde{M}$ is compact, any discrete subset
of $\tilde{M}$ is finite.
\end{proof}

\section{Orientable manifolds}
\label{smooth-manifolds-section-orientable-manifolds}

\begin{exercise}
\label{smooth-manifolds-exercise-orientable-surfaces}
Show that a regular surface $S \subset
\mathbb{R}^3$
is orientable if and only if there is a
differentiable
map $N:S \to \mathbb{R}^3$ such that
$N(p)\perp T_pS$
and $|N(p)|=1$ for all $p \in S$.
\end{exercise}

\begin{proof}
($\implies$) Take the normalization of the
cross product of an orientable basis.
It is nowhere vanishing because the
determinant,
which is positive by orientability,
is the size of the cross product.

($\impliedby$) Take the determinant of the
matrix $(v,w,N)$
where $v,w$ is any basis of $S$.
This is nowhere zero differentiable 2-form on
$S$.
\end{proof}

\section{Immersions and submersions}
\label{smooth-manifolds-section-immersions-and-submersions}

\noindent
For any smooth manifold submersion
$\pi:\tilde{M}\to M$ there is a bundle
isomorphism I love:
\begin{equation}
\label{smooth-manifolds-equation-subme-bundl-decom}
\pi^*TM \oplus \Ker \pi\cong
T\tilde{M}
\end{equation}
 where $\Ker\pi$ is the {\it kernel bundle},
 whose
fiber is the kernel of $\pi_*$, and is also
called the {\it vertical bundle}.
Proving that $\Ker\pi$ is in fact a bundle
and confirming the isomorphism
\ref{smooth-manifolds-equation-subme-bundl-decom} is a simple
exercise.

We call $\pi$ a {\it riemannian submersion}
only when $\pi^*TM$ is isometric to
$TM$, which makes sense only when $\tilde{M}$
and $M$ are Riemannian manifolds.

Following notation of
\cite{Milnor-Characteristic-Classes}, for any
subbundle
$\xi \subset\eta$ over a Riemannian manifold,
the {\it orthogonal complement
bundle} $\xi^\perp$ can be defined fiberwise
as the orthogonal complement of the
fiber and shown to be a bundle as in the case
of the kernel bundle. Then Milnor
defines the {\it normal bundle} when $\xi$ is
the tangent bundle of a
submanifold. Finally this construction is
noted to work just as well for any
isometric immersion $f:M\to \tilde{M}$, so
that
\begin{equation}
\label{smooth-manifolds-equation-immer-bundl-decom}
f^*T\tilde{M}\cong TM\oplus T^\perp M
\end{equation}
Further, in problem 3-B we are asked to
define for $\xi \subset \eta$ the {\it
quotient bundle} and show it's a bundle. And
if $\eta$ has a Riemannian metric,
then $\xi^\perp\cong\eta/\xi$. So that's why
we sometimes say that we can define
the normal bundle using a Riemannian metric
but we don't need to.

An interesting thought is: given a subbundle,
is there a manifold whose tangent
bundle is that subbundle?

\section{Regular level-set theorem}
\label{smooth-manifolds-section-regular-level-set-thm}

\noindent
This is from \cite[p. 105]{tus}.

\begin{theorem}[Regular level-set]
\label{smooth-manifolds-theorem-regular-level-set}
Is $\pi:\widetilde{M}\to M$
is a smooth submersion then every point
in $\pi^{-1}(p)$ has a local chart of
the form
$$
(x^1,…,x^{n-k},c_1,…,c_k)
$$
where $x^i$ are coordinate functions of
$\widetilde{M}$ and $c_i$ are constants.
\end{theorem}

\begin{proof}
We have the following situation on charts:
$$
\xymatrix{
\widetilde{M}\ar[r]\ar[d]^{\pi}& (x^1,…,x^n)
\ar[d]^{\tilde{\pi}}\\
M \ar[r]& (\pi^1(x),…,\pi^k(x)).
}
$$
Now introduce a new chart:
$$
(x^1,…,x^{n-k},\pi^1(x),…,\pi^k(x)),
$$
which is indeed a local diffeomorphism
since its differential is
$$
\begin{pmatrix}
\text{Id}& 0
d\pi & d\pi
\end{pmatrix},
$$
since $d\pi$ has rank $k$ we can swap columns
so that the lower-right piece of $d\pi$
has nonzero determinant.

\end{proof}

\section{Embeddings}
\label{smooth-manifolds-section-embeddings}

\begin{exercise}
\label{smooth-manifolds-exercise-p2-in-r4}
Seja \(F:\mathbb{R}^3 \to \mathbb{R}^4\)
dada por
\[F(x,y,z)=(x^2-y^2,xy,xz,yz),\qquad
(x,y,z)=p \in \mathbb{R}^3.\]
Seja \(S^2 \subset \mathbb{R}^3\) a esfera
unitária com centro na origem \(0 \in
\mathbb{R}^3\). Oberve que a restrição \(
\varphi:= F|_{S^2}\) é tal que
\(\varphi(p)=\varphi(-p)\), e considere a
aplicação \(\tilde{\varphi}:\mathbb{R}P^2
\to
\mathbb{R}^4\) dada por
\[\tilde{\varphi}([p])=\varphi(p),\qquad
[p]\text{=clase de equivalência de
\(p=\{p,-p\}\)}\]

Prove que
\begin{enumerate}
\item \(\tilde{\varphi}\) é uma imersão.
\item \(\tilde{\varphi}\) é biunívoca;
junto com (a) e a compacidade de
\(\mathbb{R}P^2\), isto implica
que \(\tilde{\varphi}\) é um mergulho.
\end{enumerate}
\end{exercise}

\begin{proof}
\begin{enumerate}
\item Considere a carta \(\{z=1\}\). A
representação coordenada de
\(\tilde{\varphi}\) vira
\[(x,y) \longmapsto (x^2-y^2,xy,x,y)\]
cuja derivada como mapa \(\mathbb{R}^2 \to
\mathbb{R}^4\) é
\[\begin{pmatrix} 2x& -2y\\y&x\\1&0\\0&1
\end{pmatrix}\]
que é injetiva. Agora pegue a carta
\(\{x=1\}\). Então a representão coordenada
de \(\tilde{\varphi}\) vira
\[(y,z) \longmapsto (1-y^2,y,z,yz)\]
e tem derivada
\[\begin{pmatrix} -2y&0\\1&0\\0&1\\z&y
\end{pmatrix}\]
que também é injetiva. Seguramente algo
análogo acontece na carta \(\{y=1\}\).

\item \(\tilde{\varphi}\) é injetiva. Pegue
dois pontos \(p_1:=[x_1:y_1:z_1]\) e
\(p_2:=[x_2:y_2:z_2]\) e suponha que
\(\tilde{\varphi}(p_1)=\tilde{\varphi}(p_2)\)
.
I.e.,
\[x_1^2-y_1^2=x_2^2-y_2^2,\qquad
x_1y_1=x_2y_2,\qquad x_1z_1=x_2z_2,\qquad
y_1z_1=y_2z_2\]
Suponha primeiro que \(z_1 \neq 0\). Segue
que
\[x_1=\frac{z_2}{z_1}x_2,\qquad
y_1=\frac{z_2}{z_1}y_2\]
logo
\[x_2^2-y_2^2=x_1^2-y_1^2=
\left(\frac{z_2}{z_1}\right)^2(x_2^2-y_2^2)
\implies
z_2=z_1\implies x_1=x_2,\qquad y_1=y_2\]

\noindent
Em fim, uma imersão injetiva com domínio
compacto é um mergulho porque é fechada:
pegue um fechado no domínio, vira compacto,
imagem é compacta, que é fechado. Pronto.
\end{enumerate}
\end{proof}

\begin{exercise}
\label{smooth-manifolds-exercise-extension-vector-fields}
Dada uma subvariedade \(M \subseteq
\tilde{M}\) uma subvariedade mergulhada e
\(X \in \mathfrak{X}(M)\). Mostre que existe
um aberto \(U \subset \tilde{M}\) contendo
\(M\) e um campo
\(\tilde{X} \in \mathfrak{X}(U)\) tal que
\(\tilde{X}|_{M}=X\).
Caso \(M\) seja subconjunto fechado de
\(\tilde{M}\), prove que \(U\) pode ser
tomado igual a \(\tilde{M}\). Se \(M\) não
é subconjunto fechado de \(\tilde{M}\),
pode não existir extensão de \(X\)
definida em todo \(\tilde{M}\).
\end{exercise}

\begin{proof}
Take submanifold coordinates for $M$,
i.e. coordinates where $M\cap U
=(x_1,...,x_n,0,...,0)$.
The local expression of $X$ is
$\sum a^i(p)\partial_i$.
Then we just literally extend this to a
vector
field in all of $U$ by putting
$\tilde{X}=\sum
a^i(\pi_M(\tilde{p}))\partial_i$.
\end{proof}

\section{Lie groups}
\label{smooth-manifolds-section-lie-groups}

\begin{definition}
\label{smooth-manifolds-definition-lie-group}
A {\it Lie group} is a smooth manifold $G$
with a group structure such that the map
$$
G\times G\to G,\qquad
(g,h)\mapsto gh^{-1}
$$
is smooth.
Equivalently, multiplication and inversion
are smooth.
A {\it Lie group homomorphism} is a smooth
group homomorphism.
\end{definition}

\medskip\noindent
The basic examples are the additive groups
$\mathbb{R}^n$ and $\mathbb{C}^n$, the
multiplicative groups $\mathbb{R}^*$ and
$\mathbb{C}^*$, the circle group
$S^1\subset\mathbb{C}$, and the tori
$$
T^n=S^1\times\cdots\times S^1.
$$
The fundamental matrix examples are
$$
\operatorname{GL}(n,\mathbb{R}),\qquad
\operatorname{GL}(n,\mathbb{C}),
$$
which are open subsets of the corresponding
matrix vector spaces. Other standard examples
are
$$
\operatorname{SL}(n,\mathbb{R}),\quad
\operatorname{O}(n),\quad
\operatorname{SO}(n),\quad
\operatorname{U}(n),\quad
\operatorname{SU}(n),\quad
\operatorname{Sp}(2n,\mathbb{R}).
$$
Here
$$
\operatorname{Sp}(2n,\mathbb{R})=
\{A\in \operatorname{GL}(2n,\mathbb{R}):
A^TJA=J\},
$$
where
$$
J=\begin{pmatrix}0&-I_n\\ I_n&0\end{pmatrix}.
$$

\begin{theorem}
\label{smooth-manifolds-theorem-closed-subgroup}
Every closed subgroup of a Lie group is an
embedded Lie subgroup.
In particular, every closed subgroup of
$\operatorname{GL}(n,\mathbb{R})$ or
$\operatorname{GL}(n,\mathbb{C})$ is a Lie
group.
\end{theorem}

\medskip\noindent
This theorem is the clean reason why the
usual matrix groups above are Lie groups.
For instance $\operatorname{O}(n)$ is the
closed subgroup cut out by $A^TA=I$, while
$\operatorname{SL}(n,\mathbb{R})$ is cut out
by $\det A=1$.

\begin{definition}
\label{smooth-manifolds-definition-lie-algebra}
Let $G$ be a Lie group. Its {\it Lie algebra}
is
$$
\mathfrak g=T_eG.
$$
If $u\in\mathfrak g$, there is a unique
left-invariant vector field $U$ on $G$ with
$U_e=u$. The bracket on $\mathfrak g$ is
then defined by
$$
[u,v]=[U,V]_e.
$$
\end{definition}

\begin{exercise}
\label{smooth-manifolds-exercise-rough-vector-field}
Let $G$ be a Lie group.
Let $X:G \to TG$ be a section
of the projection $TG \to G$,
not necessarily smooth.
Show that if $X$ is left invariant
(i.e. $dL_g(X)=X \circ L_g$
for all $g \in G$, then $X$ is smooth.
\end{exercise}

\begin{proof}
We have to show that
$X(f)$ is a smooth function on $M$
for any $f \in C^\infty(M)$.
Indeed,
$$
\begin{aligned}
X(f)g&=X_g(f)\\
&=X_e \circ L_g(f)\\
&=dL_g(X_e)(f)\\
&=X_e(f \circ L_g),
\end{aligned}
$$

\noindent
which is smooth since $X_e$
is not varying while
$f \circ L_g$ is smooth.
\end{proof}


\medskip\noindent
For a matrix Lie group $G\subset
\operatorname{GL}(n,\mathbb{F})$, its Lie
algebra can be computed as
$$
\mathfrak g=
\{A\in M_n(\mathbb{F}):\exp(tA)\in G
\text{ for all }t\in\mathbb{R}\}.
$$
This gives
$$
\mathfrak{gl}(n,\mathbb{F})=M_n(\mathbb{F}),
\qquad
\mathfrak{sl}(n,\mathbb{F})=
\{A:\operatorname{tr}A=0\},
$$
$$
\mathfrak{o}(n)=\{A:A^T+A=0\},
\qquad
\mathfrak{u}(n)=\{A:A^*+A=0\},
$$
$$
\mathfrak{su}(n)=
\{A:A^*+A=0,\operatorname{tr}A=0\},
$$
and
$$
\mathfrak{sp}(2n,\mathbb{R})=
\{A:A^TJ+JA=0\}.
$$
For matrix groups, the Lie bracket is just
the commutator
$$
[A,B]=AB-BA.
$$

\begin{definition}
\label{smooth-manifolds-definition-lie-exponential}
Let $G$ be a Lie group and let
$u\in\mathfrak g$. Let $U$ be the
left-invariant vector field determined by
$u$. The integral curve of $U$ through $e$
is defined for all $t\in\mathbb{R}$ and is a
one-parameter subgroup. We define
$$
\exp(u)=\gamma_u(1),
$$
or equivalently
$$
\exp(tu)=\gamma_u(t).
$$
\end{definition}

\begin{proposition}
\label{smooth-manifolds-proposition-exponential-properties}
The Lie group exponential map
$\exp:\mathfrak g\to G$ satisfies:
\begin{enumerate}
\item $\exp(0)=e$.
\item For each $u\in\mathfrak g$, the curve
$t\mapsto \exp(tu)$ is a one-parameter
subgroup of $G$.
\item $d(\exp)_0:\mathfrak g\to\mathfrak g$
is the identity. Hence $\exp$ is a local
diffeomorphism from a neighborhood of
$0\in\mathfrak g$ to a neighborhood of
$e\in G$.
\item If $F:G\to H$ is a Lie group
homomorphism, then
$$
F(\exp_G u)=\exp_H(dF_e(u)).
$$
\item If $[u,v]=0$, then
$$
\exp(u+v)=\exp(u)\exp(v).
$$
\end{enumerate}
\end{proposition}

\medskip\noindent
For a matrix Lie group the exponential is
the usual matrix exponential:
$$
\exp(A)=e^A=
\sum_{k=0}^\infty \frac{A^k}{k!}.
$$
In general $\exp$ need not be globally
surjective and it does not usually satisfy
$\exp(u+v)=\exp(u)\exp(v)$ unless $u$ and
$v$ commute.

\begin{definition}
\label{smooth-manifolds-definition-lie-group-action}
A smooth left action of $G$ on a manifold
$M$ is a smooth map
$$
G\times M\to M,\qquad (g,p)\mapsto g\cdot p,
$$
satisfying $e\cdot p=p$ and
$(gh)\cdot p=g\cdot(h\cdot p)$.
For $p\in M$, the {\it orbit map} is
$$
\Phi_p:G\to M,
\qquad
\Phi_p(g)=g\cdot p.
$$
The orbit is $G\cdot p=\Phi_p(G)$.
The stabilizer is
$$
G_p=\{g\in G:g\cdot p=p\}.
$$
\end{definition}

\begin{definition}
\label{smooth-manifolds-definition-fundamental-vector-field}
For $u\in\mathfrak g$, the associated
{\it fundamental vector field} on $M$ is
$$
u_M(p)=
\left.\frac{d}{dt}\right|_{t=0}
\exp(tu)\cdot p.
$$
\end{definition}

\begin{theorem}
\label{smooth-manifolds-theorem-orbit-theorem-lie-actions}
Let a Lie group $G$ act smoothly on a
manifold $M$. Then every orbit $G\cdot p$ is
an immersed submanifold of $M$ and the orbit
map
$$
\Phi_p:G\to G\cdot p
$$
is a submersion. Moreover $G_p$ is a closed
Lie subgroup of $G$, and the natural map
$$
G/G_p\to G\cdot p,
\qquad
[g]\mapsto g\cdot p
$$
is a diffeomorphism, where the orbit carries
its immersed-submanifold smooth structure.
\end{theorem}

\medskip\noindent
In particular, differentiating the orbit map
at the identity gives
$$
T_p(G\cdot p)=\operatorname{Im}(d\Phi_p)_e
=\{u_M(p):u\in\mathfrak g\}.
$$
The kernel of $(d\Phi_p)_e$ is the stabilizer
algebra
$$
\mathfrak g_p=
\{u\in\mathfrak g:u_M(p)=0\},
$$
so infinitesimally one has
$$
T_p(G\cdot p)\cong \mathfrak g/\mathfrak g_p.
$$
This is the precise sense in which an orbit
is locally parametrized by $\mathfrak g$, but
with stabilizer directions killed.

\medskip\noindent
For the coadjoint action on $\mathfrak g^*$,
the convention
$$
\langle \operatorname{Ad}^*_g\xi,v\rangle
=
\langle \xi,\operatorname{Ad}_{g^{-1}}v\rangle
$$
gives
$$u_{\mathfrak g^*}(\xi)=
\left.\frac{d}{dt}\right|_{t=0}
\operatorname{Ad}^*_{\exp(tu)}\xi.
$$
Hence, for the coadjoint orbit
$\mathcal O_\xi$, the orbit theorem gives
$$
T_\xi\mathcal O_\xi=
\{u_{\mathfrak g^*}(\xi):u\in\mathfrak g\}.
$$
The sign convention above implies
$$
\langle u_{\mathfrak g^*}(\xi),v\rangle
=
-\langle \xi,[u,v]\rangle.
$$
This is the KKS symplectic form
on the coadjoint orbits,
see Symplectic Geometry Section
\ref{symplectic-geometry-section-kks}.

\medskip\noindent
Let's introduce the example of
the action of $\text{SO}_3$ in
$\mathbb{R}^3$, which will appear
again in symplectic geometry.
The main upshot is that
``infinitesimal rotations act by
cross-product''. That is, an infinitesimal
rotation, which is a rather abstract
gadget that we imagine simply as a tiny
rotation, puts a vector $v \in \mathbb{R}^3$
to some other vector infinitesimally
near-by. We can only keep track of
the speed of this tiny movement,
which is the infinitesimal
generator of an element
$u \in\mathfrak{so}_3 \simeq \mathbb{R}^3$,
applied to the vector $v$.
The key observation is that
$u_{\mathfrak{so}_3}(v)
=u \times v $.

At the risk of being redundant:
If $v\in\mathbb R^3$, then $A(v)$ is the
instantaneous velocity of the point $v$ under
this infinitesimal rotation. But the
instantaneous velocity of rotation with
angular velocity vector $u$ is
$$
u\times v.
$$
Thus the matrix $A$ acts on vectors as
$$
A(v)=u\times v.
$$
In other words, under the identification
$u\leftrightarrow \hat u$, the skew-symmetric
matrix $\hat u$ is the operator
$$
v\mapsto u\times v.
$$
Now the question is: why?
At least intuitively, a rotation
puts $v$ to another vector contained
in some sphere determined by the axis $u$.
Thus the velocity of the rotation
will be perpendicular to the sphere,
and thus the cross-product fits in
perfectly. Formally, this reduces
to the identification 
$\mathfrak{so}_3 \simeq \mathbb{R}^3$
we use
in the next exercise.

\medskip\noindent
This is a special feature of dimension $3$.
Algebraically, skew-symmetric matrices are the
same as bivectors:
$$
\mathfrak{so}(3)\cong \Lambda^2\mathbb R^3.
$$
But in dimension $3$, the Hodge star
identifies bivectors with vectors:
$$
\Lambda^2\mathbb R^3\cong\mathbb R^3.
$$
Thus
$$
\mathfrak{so}(3)
\cong
\Lambda^2\mathbb R^3
\cong
\mathbb R^3.
$$
The cross product is precisely this
dimension-three identification written as an
operation.


\begin{exercise}
\label{smooth-manifolds-exercise-so3-adjoint-action}
Consider the linear isomorphism
$\mathbb{R}^3\to\mathfrak{so}(3)$, given by
$$
v=(x,y,z)\mapsto \hat v:=
\begin{pmatrix}
0&-z&y\\
z&0&-x\\
-y&x&0
\end{pmatrix}
$$

\begin{enumerate}
\item Describe the Lie bracket on
$\mathbb{R}^3$ induced by the commutator in
$\mathfrak{so}(3)$, and the inner product in
$\mathfrak{so}(3)$ that corresponds to the
canonical inner product in $\mathbb{R}^3$.

\item Describe the
$\operatorname{SO}(3)$-action on
$\mathbb{R}^3$ corresponding to the adjoint
action, its orbits, as well as its
infinitesimal generators. Find (without any
calculation!) a description of the coadjoint
action on $\mathbb{R}^3$ identified with
$(\mathbb{R}^3)^*$ through the canonical
inner product.
\end{enumerate}
\end{exercise}

\begin{proof}
\begin{enumerate}
\item Save doing the computations,
the Lie bracket is the cross product
and the innter product is proportional
to the trace (not the determinant!).

\item First we need to find
$\text{Ad}_gu$ for
$g \in \text{SO}_3$
and $u \in \mathfrak{so}_3$.
The first thing I do is track the definition.
I get:
$$
\text{Ad}_gu
=\frac{d}{dt}\Big|_{t=0}
g^{-1}\text{exp}(tu)\cdot g
=g^{-1}\frac{d}{dt}\Big|_{t=0}\text{exp}(tu)\cdot
g=g^{-1}ug.
$$
But that's not good enough.
While I would dare to go to coordinates
and try to figure out some pattern…
it might also be sort of clear
that there's a way of using the
identification in the first item.
The answer is AI-crafted:
To see that the adjoint action corresponds to
the usual action of $\operatorname{SO}(3)$ on
$\mathbb{R}^3$, let $R\in\operatorname{SO}(3)$.
We claim that
$$
R\hat vR^{-1}=\widehat{Rv}.
$$
Indeed, both sides are linear maps
$\mathbb{R}^3\to\mathbb{R}^3$, so it is
enough to apply them to an arbitrary vector
$w$ (this was basically my 
argument above differentiating exp). We get
$$
(R\hat vR^{-1})(w)
=
R(\hat v(R^{-1}w))
=
R(v\times R^{-1}w).
$$
Since $R\in\operatorname{SO}(3)$ preserves
the cross product,
$$
R(v\times R^{-1}w)
=
(Rv)\times(RR^{-1}w)
=
(Rv)\times w.
$$
But this is exactly
$$
\widehat{Rv}(w).
$$
Therefore
$$
R\hat vR^{-1}=\widehat{Rv}.
$$
The orbit of $u \in \mathfrak{so}_3$
under the adjoint action is 
$\mathcal{O}_u=\{\text{Ad}_gu:g \in
\text{SO}(3)\}$. By the discussion
so far we know these are spheres
in $\mathbb{R}^3$.

Then I can tell that the infinitesimal
generators $u_{\mathfrak{so}_3}$
are tangent vectors to those spheres,
but that's no answer to the question
``what are the infinitesimal generators?''.
The real answer comes from the observation
that precedes this example; namely 
that ``infinitesimal rotations
act by cross product'':
Under the identification
$\mathbb R^3\cong\mathfrak{so}(3)$,
a matrix $\hat u\in\mathfrak{so}(3)$ acts on
an abstract vector $v\in\mathbb R^3$ by
cross product with the corresponding vector
$u$:
$$
\hat u(v)=u\times v.
$$
This confirmed trivially by
computing the cross product using the
identification given above.
Now by the definition of infinitesimal
generator we have:

\end{enumerate}
\end{proof}

\section{Principal bundles}
\label{smooth-manifolds-section-principal-bundles}

Upshot:
a principal bundle is a fiber bundle such
that, if you choose a point in a fiber,
then the map $g \mapsto p\cdot g$
identifies that fiber with $G$.
It is surjective by transitivity and
injective by freeness.
The chosen point plays the role of the
identity.

\begin{definition}
\label{smooth-manifolds-definition-principal-bundle}
Let $G$ be a Lie group.
A {\it principal $G$-bundle}
\[
\pi:P\to M
\]
is a smooth fiber bundle with a smooth
right action of $G$ on $P$ such that:
\begin{enumerate}
\item the action preserves fibers,
\item it is free and transitive on each
fiber, and
\item locally $P|_U \cong U\times G$
as $G$-spaces,
with $G$ acting on itself by right
multiplication.
\end{enumerate}
\end{definition}

Equivalently,
the transition functions of a local
trivialization take values in $G$.

In particular,
the frame bundle $\operatorname{Fr}(M)$
from
hodge-birational-atoms-2026.tex
is a principal
\(\operatorname{GL}(n,\mathbf R)\)-bundle.

% --- End smooth-manifolds.tex ---

% --- Begin differential-topology.tex ---
\chapter{Differential Topology}
\label{chapter-differential-topology}





\section{Transversality theorems}
\label{differential-topology-section-transversality-theorems}

\noindent
A simple case of the following definition is
when $f$ is the inclusion and $X,W$ are both
submanifolds of $Y$.

\begin{definition}
\label{differential-topology-definition-transversal-map}
Let $X,Y$ be smooth manifolds and let
$f \in C^\infty(X,Y)$.
Let $W \subset Y$ be a submanifold.
We say that $f$ is {\it transverse} to $W$,
and write $f \pitchfork W$, if for every
$x \in f^{-1}(W)$ we have
$$
d_xf(T_xX)+T_{f(x)}W=T_{f(x)}Y.
$$
\end{definition}

\noindent
This imposes restrictions on the dimensions
of $X$ and $W$:
$$
\dim X+\dim W \geq \dim Y.
$$
This may also be thought of as a
generalization of the regular value
condition.  If $W$ is a single point $y$,
then $f \pitchfork \{y\}$ says exactly that
$$
d_xf(T_xX)=T_yY,
\qquad x \in f^{-1}(y),
$$
that is, $df$ is surjective at every point of
the fiber.

\begin{theorem}
\label{differential-topology-theorem-preimage-transverse}
Let $f:X \to Y$ be smooth and let
$W \subset Y$ be a submanifold.  If
$f \pitchfork W$, then $f^{-1}(W)$ is a
submanifold of $X$ and
$$
\operatorname{codim}_X f^{-1}(W)
=
\operatorname{codim}_Y W.
$$
\end{theorem}

\begin{proof}
This is the regular value theorem in local
coordinates.  Near a point of $W$ choose
coordinates in which $W$ is cut out by the
last $c$ coordinates, where
$c=\operatorname{codim}W$.  The map to these
last $c$ coordinates has surjective
differential along $f^{-1}(W)$ exactly because
of transversality.  Hence its zero set is a
submanifold of codimension $c$.
\end{proof}

\section{Course plan and review}
\label{differential-topology-section-course-plan-review}

\begin{reference}
These notes follow Vinicius Ramos' course on
differential topology.  Standard references
are \cite{milnordt}, \cite{muk}, and
\cite{hirsch}.
\end{reference}

\noindent
The course plan was the following.

\begin{enumerate}
\item Review of manifolds.

\item Transversality: Sard's theorem,
Whitney topologies, and approximation.

\item Intersection theory and index.

\item Morse theory.

\item Additional possible topics:
$h$-cobordism, low-dimensional topology, and
Poincar\'e in dimensions $n\geq 5$.
\end{enumerate}

\noindent
We use the standard convention that a smooth
manifold is Hausdorff, second countable, and
locally Euclidean, equipped with a maximal
smooth atlas.  The dimension is locally
constant, hence constant on connected
components.  There are topological atlases
which do not contain a smooth subatlas.

\begin{theorem}
\label{differential-topology-theorem-smoothing-manifolds}
Every $C^k$-manifold, for
$k=1,\ldots,\infty$, is $C^k$-diffeomorphic
to a $C^\infty$-manifold.
\end{theorem}

\begin{proof}
We record the statement as one of the basic
smoothing theorems used in the course.  The
proof is not included here.
\end{proof}

\begin{theorem}
\label{differential-topology-theorem-compatible-smoothings}
Let $1\leq \ell \leq k \leq \infty$.
If two $C^k$-manifolds are
$C^\ell$-diffeomorphic, then they are
$C^k$-diffeomorphic.
\end{theorem}

\begin{proof}
We omit the proof.  The point is that lower
regularity compatibility is enough to compare
the induced smooth structures.
\end{proof}

\begin{theorem}
\label{differential-topology-theorem-partition-of-unity}
Let $M$ be a smooth manifold and let
$\{U_i\}$ be an open cover.  Then there is a
smooth partition of unity subordinate to this
cover.
\end{theorem}

\begin{proof}
This uses paracompactness of smooth
manifolds.  Every topological manifold is
paracompact, and smooth bump functions allow
one to refine locally finite covers by smooth
functions.
\end{proof}

\section{Tangent and cotangent bundles}
\label{differential-topology-section-tangent-cotangent-bundles}

\noindent
Let $M$ be a smooth manifold.  The tangent
space $T_pM$ can be defined either as
velocities of curves through $p$ or as
derivations at $p$.  In a chart
$(U,\varphi)$ with
$\varphi=(x^1,\ldots,x^n)$, a basis is
$$
\left\{
\frac{\partial}{\partial x^1}\Big|_p,
\ldots,
\frac{\partial}{\partial x^n}\Big|_p
\right\}.
$$
The cotangent space is
$$
T_p^*M=(T_pM)^*
=\operatorname{Hom}(T_pM,\mathbb R),
$$
with dual basis
$$
\left\{dx^1|_p,\ldots,dx^n|_p\right\},
\qquad
 dx^i|_p
\left(
\frac{\partial}{\partial x^j}\Big|_p
\right)=\delta^i_j.
$$
Changing charts changes these bases; there is
no canonical basis of $T_pM$ or $T_p^*M$.

\begin{exercise}
\label{differential-topology-exercise-change-of-basis}
Let $(U,\varphi)$ and $(V,\psi)$ be charts
near $p$, with coordinates
$x^1,\ldots,x^n$ and $y^1,\ldots,y^n$.
Show that
$$
\frac{\partial}{\partial x^j}
=
\sum_{i=1}^n
\frac{\partial y^i}{\partial x^j}
\frac{\partial}{\partial y^i}.
$$
\end{exercise}

\noindent
The tangent bundle is
$$
TM=\bigsqcup_{p\in M}T_pM.
$$
For every chart $(U,\varphi)$ there is a
bijection
$$
\begin{aligned}
U\times \mathbb R^n
&\longrightarrow \pi^{-1}(U)\\
(p,(v_1,\ldots,v_n))
&\longmapsto
\sum_{i=1}^n v_i
\frac{\partial}{\partial x^i}\Big|_p.
\end{aligned}
$$
Using these bijections one topologizes $TM$
and gives it a smooth structure of dimension
$2n$.  The same construction gives the
cotangent bundle
$$
T^*M=\bigsqcup_{p\in M}T_p^*M.
$$
If $M$ is only $C^k$, then $TM$ is naturally
$C^{k-1}$.

\begin{definition}
\label{differential-topology-definition-parallelizable}
A smooth $n$-manifold $M$ is
{\it parallelizable} if
$$
TM\cong M\times \mathbb R^n.
$$
\end{definition}

\begin{example}
\label{differential-topology-example-parallelizable-manifolds}
The manifolds $\mathbb R^n$ and $S^1$ are
parallelizable.  Every orientable
3-manifold is also parallelizable, although
that result is much deeper.
\end{example}

\section{Immersions, submersions, and embeddings}
\label{differential-topology-section-immersions-submersions-embeddings}

\begin{definition}
\label{differential-topology-definition-differential-map}
Let $f:M\to N$ be smooth.  The
{\it derivative} or {\it differential} of
$f$ at $p$ is the linear map
$$
Df_p:T_pM\to T_{f(p)}N.
$$
In the curve model of tangent vectors,
$Df_p[\gamma]=[f\circ\gamma]$.  In the
derivation model,
$$
(Df_pv)(g)=v(g\circ f),
\qquad g\in C^\infty(N).
$$
\end{definition}

\noindent
In local coordinates $x^j$ on $M$ and $y^i$
on $N$, this is represented by the Jacobian
matrix:
$$
Df_p
\left(\frac{\partial}{\partial x^j}\Big|_p
\right)
=
\sum_i
\frac{\partial f_i}{\partial x^j}(p)
\frac{\partial}{\partial y^i}\Big|_{f(p)}.
$$

\begin{definition}
\label{differential-topology-definition-immersion-submersion-embedding}
Let $f:M\to N$ be smooth.
\begin{enumerate}
\item $f$ is an {\it immersion} at $p$ if
$Df_p$ is injective.

\item $f$ is a {\it submersion} at $p$ if
$Df_p$ is surjective.

\item $f$ is an {\it embedding} if it is an
injective immersion and the inverse
$f(M)\to M$ is continuous, where $f(M)$ has
the subspace topology from $N$.
\end{enumerate}
\end{definition}

\noindent
We often write
$$
M\rightlooparrow N,
\qquad
M\hookrightarrow N,
\qquad
M\twoheadrightarrow N
$$
for immersions, embeddings, and submersions.
The continuity condition in the definition of
embedding is necessary: an injective immersed
curve may accumulate on itself in the target,
so that the inverse from its image is not
continuous.

\begin{remark}
\label{differential-topology-remark-embedding-submanifold}
If $f:M\to N$ is an embedding, then $f(M)$
inherits a unique smooth structure such that
$f:M\to f(M)$ is a diffeomorphism.  Thus
embeddings are precisely the maps whose
images are honest submanifolds.
\end{remark}

\begin{exercise}
\label{differential-topology-exercise-same-embedded-image}
Let $N_1,N_2$ be manifolds and let
$\varphi_i:N_i\to M$ be smooth embeddings.
Suppose their images coincide.  Show that
$N_1$ and $N_2$ are diffeomorphic.
\end{exercise}

\section{Regular values}
\label{differential-topology-section-regular-values}

\begin{definition}
\label{differential-topology-definition-regular-value}
Let $f:M\to N$ be smooth.  A point $y\in N$
is a {\it regular value} of $f$ if $f$ is a
submersion at every point of $f^{-1}(y)$.
Equivalently, $Df_x$ is surjective for all
$x\in f^{-1}(y)$.
\end{definition}

\begin{theorem}[Regular value theorem]
\label{differential-topology-theorem-regular-value}
If $y$ is a regular value of $f:M\to N$, then
$f^{-1}(y)$ is a submanifold of $M$ of
dimension
$$
\dim M-\dim N,
$$
provided the fiber is nonempty.
\end{theorem}

\begin{proof}
This is the implicit function theorem.  Take
charts around $x\in f^{-1}(y)$ and around
$y$, both centered at the relevant points.
In coordinates we get a map
$\Phi:\mathbb R^m\to\mathbb R^n$ with
surjective derivative at $0$.  After a local
change of coordinates, $\Phi$ becomes the
projection
$$
\mathbb R^n\times \mathbb R^{m-n}
\longrightarrow \mathbb R^n.
$$
Hence $\Phi^{-1}(0)$ is locally
$\{0\}\times \mathbb R^{m-n}$.
\end{proof}

\section{Vector bundles and sections}
\label{differential-topology-section-vector-bundles-sections}

\begin{definition}
\label{differential-topology-definition-vector-bundle}
Let $E$ and $M$ be manifolds and let
$\pi:E\to M$ be a surjective submersion.
We say that $\pi$ is a {\it vector bundle} of
rank $r$ if every fiber
$E_p=\pi^{-1}(p)$ is a vector space and for
every $p\in M$ there is a neighborhood $U$
and a diffeomorphism
$$
\varphi:\pi^{-1}(U)\to U\times\mathbb R^r
$$
which is linear on each fiber and commutes
with the projections to $U$.
\end{definition}

\begin{example}
\label{differential-topology-example-basic-vector-bundles}
The tangent bundle, cotangent bundle, and the
trivial bundle $M\times\mathbb R^r\to M$ are
vector bundles.
\end{example}

\begin{definition}
\label{differential-topology-definition-section-vector-bundle}
A {\it section} of a vector bundle
$\pi:E\to M$ is a smooth map
$s:M\to E$ such that $\pi\circ s=\text{id}_M$.
\end{definition}

\noindent
A vector field is a section of $TM$.  A
1-form is a section of $T^*M$.  If a section
is compactly supported, its support has
compact closure.  Partitions of unity are the
basic tool for producing global sections from
local data.

\section{Sard theorem}
\label{differential-topology-section-sard-theorem}

\begin{definition}
\label{differential-topology-definition-critical-point}
Let $f:X^m\to Y^n$ be smooth.  A point
$x\in X$ is a {\it critical point} if
$Df_x$ does not have maximal rank.  A point
$y\in Y$ is a {\it critical value} if
$y=f(x)$ for some critical point $x$.
\end{definition}

\begin{theorem}[Sard]
\label{differential-topology-theorem-sard}
Let $f:X\to Y$ be smooth.  The set of
critical values of $f$ has measure zero in
$Y$.  In particular, the regular values of
$f$ are dense.
\end{theorem}

\begin{proof}
We only record the structure of the proof.
In Euclidean coordinates one considers the
sets where all partial derivatives up to a
certain order vanish.  Taylor's theorem shows
that the image of each such set can be
covered by boxes whose total volume is
arbitrarily small.  The manifold case follows
by covering the domain with countably many
coordinate charts.
\end{proof}

\begin{remark}
\label{differential-topology-remark-sard-finite-regularity}
Sard's theorem also holds for $C^\ell$ maps
provided
$$
\ell>\max(m-n,0).
$$
This regularity threshold is important in the
full statement.
\end{remark}

\section{Jets}
\label{differential-topology-section-jets}

\begin{definition}
\label{differential-topology-definition-jet-equivalence}
Let $f,g:X\to Y$ be smooth maps.  We say
that $f$ and $g$ have the same $k$-jet at
$x\in X$ if, in local coordinates around $x$
and around $f(x)=g(x)$, all partial
derivatives up to order $k$ agree at $x$.
The equivalence class is denoted $j_x^kf$.
\end{definition}

\begin{definition}
\label{differential-topology-definition-jet-space}
The {\it $k$-jet space} $J^k(X,Y)$ is the set
of all $k$-jets $j_x^kf$ of smooth maps
$X\to Y$.
There are natural source and target maps
$$
\alpha:J^k(X,Y)\to X,
\qquad
\beta:J^k(X,Y)\to Y.
$$
\end{definition}

\noindent
The space $J^k(X,Y)$ is a smooth manifold.
For $k=0$ one has $J^0(X,Y)=X\times Y$.
For $k=1$, a point of $J^1(X,Y)$ consists of
$x\in X$, $y\in Y$, and a linear map
$T_xX\to T_yY$.
Thus
$$
J^1(X,Y)
\cong
\operatorname{Hom}(TX,TY)
$$
in the appropriate fibered sense.

\begin{definition}
\label{differential-topology-definition-jet-extension}
For a smooth map $f:X\to Y$, its
{\it $k$-jet extension} is the map
$$
j^kf:X\to J^k(X,Y),
\qquad
x\longmapsto j_x^kf.
$$
\end{definition}

\section{Whitney topologies}
\label{differential-topology-section-whitney-topologies}

\noindent
The space $C^\infty(X,Y)$ admits several
natural topologies.  The weak Whitney
topology controls a map and finitely many
of its derivatives on compact subsets.  The
strong Whitney topology allows the control to
vary over locally finite families of compact
sets.  These topologies are designed so that
statements such as ``transverse maps are
dense'' become precise.

\begin{definition}
\label{differential-topology-definition-delta-ball}
Let $\delta:X\to (0,\infty)$ and let
$f\in C^\infty(X,Y)$, after choosing a
metric on a jet bundle.  A $\delta$-ball
around $f$ consists of those maps $g$ whose
relevant jets stay within distance
$\delta(x)$ from the jets of $f$ at every
$x$.
\end{definition}

\begin{definition}
\label{differential-topology-definition-baire-space}
A topological space is a {\it Baire space} if
a countable intersection of open dense sets
is dense.
\end{definition}

\noindent
The strong Whitney topology is a Baire
topology on $C^\infty(X,Y)$ in the usual
situations.  This is why genericity
statements can be proved by intersecting
countably many open dense conditions.

\section{Parametric transversality}
\label{differential-topology-section-parametric-transversality}

\begin{theorem}[Parametric transversality]
\label{differential-topology-theorem-parametric-transversality}
Let $F:X\times B\to Y$ be smooth and let
$W\subset Y$ be a submanifold.  If
$F\pitchfork W$, then for almost every
$b\in B$ the map
$$
F_b:X\to Y,
\qquad
F_b(x)=F(x,b),
$$
is transverse to $W$.
\end{theorem}

\begin{proof}
Let $Q=F^{-1}(W)$.  Since $F\pitchfork W$,
Theorem \ref{differential-topology-theorem-preimage-transverse}
shows that $Q$ is a submanifold of
$X\times B$.  Consider the projection
$\pi:Q\to B$.  By Sard's theorem, almost
every $b\in B$ is a regular value of $\pi$.
For such $b$, the condition that $b$ be a
regular value of $\pi$ is equivalent to
$F_b\pitchfork W$.
\end{proof}

\begin{theorem}[Thom transversality]
\label{differential-topology-theorem-thom-transversality}
Let $X,Y$ be smooth manifolds and let
$W\subset J^k(X,Y)$ be a submanifold.  Then
$$
\{f\in C^\infty(X,Y):j^kf\pitchfork W\}
$$
is residual in $C^\infty(X,Y)$, and in the
strong Whitney topology it is dense.  If
$W$ is closed, this set is also open.
\end{theorem}

\begin{proof}
The proof reduces locally to the parametric
transversality theorem.  One perturbs a map
by adding polynomial terms in coordinate
balls.  The finite-dimensional space of
polynomial perturbations supplies the
parameter space $B$, and the evaluation map
from $X\times B$ to the jet space is made
transverse to $W$.  A partition of unity and
a Baire category argument globalize the
construction.
\end{proof}

\section{Multijets and Whitney theorems}
\label{differential-topology-section-multijets-whitney-theorems}

\noindent
Multijet spaces keep track of jets of a map
at several distinct points at once.  They are
used to impose conditions on pairs or tuples
of points, for example to avoid double points
or to control self-intersections.

\begin{theorem}[Whitney immersion theorem]
\label{differential-topology-theorem-whitney-immersion}
Every smooth $n$-manifold admits an immersion
into $\mathbb R^{2n-1}$.
\end{theorem}

\begin{proof}
Start with an embedding into some
$\mathbb R^N$.  Project generically to a
lower-dimensional Euclidean space.  The bad
projections are those for which the
differential drops rank somewhere; this is a
transversality condition on the 1-jet.  Thom
transversality shows that a generic
projection avoids the bad locus as long as
the target has dimension at least $2n-1$.
\end{proof}

\begin{theorem}[Whitney embedding theorem]
\label{differential-topology-theorem-whitney-embedding}
Every smooth $n$-manifold admits an embedding
into $\mathbb R^{2n}$.
\end{theorem}

\begin{proof}
Again start from an embedding into
$\mathbb R^N$ and project generically.  The
immersion condition is controlled by the
1-jet.  Injectivity is controlled by the
2-point multijet: one must avoid the diagonal
in $\mathbb R^{2n}\times\mathbb R^{2n}$ for
pairs of distinct points.  The relevant bad
sets have sufficiently large codimension for
generic projections to avoid them.
\end{proof}

\section{Manifolds with boundary}
\label{differential-topology-section-manifolds-with-boundary}

\begin{definition}
\label{differential-topology-definition-half-space}
Let
$$
\mathbb H^n=\{(x_1,\ldots,x_n)\in
\mathbb R^n:x_n\geq 0\}.
$$
A {\it smooth manifold with boundary} is a
space locally modeled on open subsets of
$\mathbb H^n$ with smooth transition maps.
\end{definition}

\noindent
The boundary $\partial M$ consists of those
points which in a half-space chart map to
$x_n=0$.  The interior consists of the points
which map to $x_n>0$.  This is independent
of the chart.

\begin{theorem}[Sard with boundary]
\label{differential-topology-theorem-sard-boundary}
Let $f:X\to Y$ be smooth, where $X$ is a
manifold with boundary.  If $a\in Y$ is a
regular value of both $f$ and
$f|_{\partial X}$, then $f^{-1}(a)$ is a
manifold with boundary and
$$
\partial(f^{-1}(a))=f^{-1}(a)\cap \partial X.
$$
\end{theorem}

\begin{proof}
Apply the regular value theorem in boundary
charts.  In interior charts this is the usual
argument.  In boundary charts the local model
is the intersection of a regular level set
with a half-space.
\end{proof}

\begin{theorem}[Thom transversality with boundary]
\label{differential-topology-theorem-thom-boundary}
Let $X$ be a manifold with boundary and let
$W\subset Y$ be a submanifold.  Generically,
a smooth map $f:X\to Y$ is transverse to
$W$ and $f|_{\partial X}$ is also transverse
to $W$.
\end{theorem}

\begin{proof}
Apply the ordinary Thom transversality
theorem simultaneously to $X$ and to
$\partial X$.  The boundary condition is an
extra transversality condition on the
restriction.
\end{proof}

\section{Orientation}
\label{differential-topology-section-orientation}

\begin{definition}
\label{differential-topology-definition-orientation-vector-space}
An {\it orientation} of a finite-dimensional
real vector space is a choice of one of the
two connected components of the set of
ordered bases.  Equivalently, two ordered
bases define the same orientation if the
change-of-basis determinant is positive.
\end{definition}

\begin{definition}
\label{differential-topology-definition-orientation-manifold}
An {\it orientation} of a smooth manifold
$M$ is a smoothly varying choice of
orientation on each tangent space $T_pM$.
\end{definition}

\noindent
If $M$ is an oriented manifold with boundary,
then $\partial M$ receives the boundary
orientation by the ``outward normal first''
rule: an ordered basis $(v_1,\ldots,v_{n-1})$
of $T_p\partial M$ is positive if
$(\nu,v_1,\ldots,v_{n-1})$ is a positive
basis of $T_pM$, where $\nu$ is an outward
normal vector.

\section{Intersection theory}
\label{differential-topology-section-intersection-theory}

\begin{definition}
\label{differential-topology-definition-intersection-number}
Let $X^n$ and $Y^m$ be oriented manifolds,
let $W^k\subset Y$ be an oriented
submanifold, and let $f:X\to Y$ be smooth and
transverse to $W$.  If
$$
\dim X+\dim W=\dim Y,
$$
then $f^{-1}(W)$ is a finite set when $X$ is
compact.  The {\it intersection number}
$I(f,W)$ is the signed count of this set.
\end{definition}

\noindent
The sign at a point $x\in f^{-1}(W)$ is
determined by comparing the orientation of
$T_{f(x)}Y$ with the orientation obtained
from
$$
df_x(T_xX)\oplus T_{f(x)}W.
$$

\begin{proposition}
\label{differential-topology-proposition-intersection-homotopy}
If $f_0,f_1:X\to Y$ are homotopic through
maps transverse to $W$ at the ends, then
$$
I(f_0,W)=I(f_1,W).
$$
\end{proposition}

\begin{proof}
Let $F:X\times[0,1]\to Y$ be a homotopy
between $f_0$ and $f_1$, perturbed relative
to the boundary so that $F\pitchfork W$.
Then $F^{-1}(W)$ is an oriented compact
1-manifold whose boundary is
$f_1^{-1}(W)-f_0^{-1}(W)$.  The signed number
of boundary points of a compact oriented
1-manifold is zero.
\end{proof}

\section{Degree}
\label{differential-topology-section-degree}

\begin{definition}
\label{differential-topology-definition-degree-map}
Let $X$ and $Y$ be compact oriented
$n$-manifolds, with $Y$ connected.  For a
smooth map $f:X\to Y$ and a regular value
$y\in Y$, the {\it degree} of $f$ is
$$
\deg(f)=\sum_{x\in f^{-1}(y)}\operatorname{sign}
(Df_x).
$$
\end{definition}

\begin{proposition}
\label{differential-topology-proposition-degree-well-defined}
The degree is independent of the chosen
regular value and invariant under homotopy.
\end{proposition}

\begin{proof}
This is the intersection number of $f$ with
the point $y\subset Y$.  Independence of $y$
and homotopy invariance follow from the
homotopy invariance of intersection number.
\end{proof}

\begin{example}
\label{differential-topology-example-degree-circle}
For maps $S^1\to S^1$, the degree is the
winding number.  If $f:S^1\to S^1$ is given
by $f(z)=z^d$, then $\deg(f)=d$.
\end{example}

\section{Fundamental theorem of algebra}
\label{differential-topology-section-fundamental-theorem-algebra-degree}

\begin{proposition}
\label{differential-topology-proposition-fta-degree}
Every nonconstant complex polynomial has a
complex root.
\end{proposition}

\begin{proof}
Let $p(z)$ be a polynomial of degree $d>0$.
For $R$ large, the map
$$
S_R^1\to S^1,
\qquad
z\longmapsto \frac{p(z)}{|p(z)|}
$$
has the same degree as
$z\mapsto z^d/|z|^d$, namely $d$.  If $p$ had
no root in the disk, this boundary map would
extend to the disk and hence would have
degree zero.  This is impossible.
\end{proof}

\section{Winding number and Jordan--Brouwer}
\label{differential-topology-section-winding-jordan-brouwer}

\noindent
For a smooth closed curve
$\gamma:S^1\to \mathbb R^2\setminus\{p\}$,
the winding number about $p$ is the degree of
$$
S^1\to S^1,
\qquad
z\longmapsto
\frac{\gamma(z)-p}{\|\gamma(z)-p\|}.
$$
It is constant on connected components of the
complement of the curve.

\begin{theorem}[Jordan--Brouwer separation]
\label{differential-topology-theorem-jordan-brouwer}
If $M\subset S^n$ is a compact connected
embedded hypersurface, then
$S^n\setminus M$ has exactly two connected
components, and $M$ is the boundary of each.
\end{theorem}

\begin{proof}
We only record the idea from the course.  One
uses degree or mod-two intersection number
with paths crossing the hypersurface.  The
intersection parity changes when crossing
$M$, and this produces two complementary
regions.  A local chart shows that there are
at least two local sides, and the global
intersection argument shows there are no more
than two.
\end{proof}

\section{Borsuk--Ulam theorem}
\label{differential-topology-section-borsuk-ulam}

\begin{theorem}[Borsuk--Ulam]
\label{differential-topology-theorem-borsuk-ulam}
Every continuous map $f:S^n\to\mathbb R^n$
identifies a pair of antipodal points: there
exists $x\in S^n$ such that $f(x)=f(-x)$.
\end{theorem}

\begin{proof}
The smooth case follows from degree theory by
considering the odd map
$g(x)=f(x)-f(-x)$.  If $g$ never vanished,
then $g/\|g\|:S^n\to S^{n-1}$ would give an
odd map to a lower-dimensional sphere, which
contradicts the corresponding degree parity
statement.  The continuous case follows by
approximation.
\end{proof}

\section{Poincare--Hopf theorem}
\label{differential-topology-section-Poincare-Hopf-theorem}

\begin{definition}
\label{differential-topology-definition-isolated-zero-vector-field}
Let $V\in\mathfrak X(X)$ be a vector field.
A point $x\in X$ with $V(x)=0$ is an
{\it isolated zero} if there is a
neighborhood $U$ of $x$ such that
$V(y)\neq0$ for all $y\in U\setminus\{x\}$.
\end{definition}

\begin{definition}
\label{differential-topology-definition-index-vector-field}
Let $x$ be an isolated zero of $V$ on an
$n$-manifold.  Choose coordinates identifying
a small ball around $x$ with $B_r(0)$.
The {\it index} of $V$ at $x$ is
$$
\operatorname{ind}(V,x)
=
\deg\left(
\frac{\widetilde V}{\|\widetilde V\|}
:\partial B_r(0)\to S^{n-1}
\right).
$$
\end{definition}

\begin{theorem}[Poincare--Hopf]
\label{differential-topology-theorem-Poincare-Hopf}
Let $X$ be a compact manifold without
boundary and let $V\in\mathfrak X(X)$ have
isolated zeroes.  Then
$$
\sum_{x\in Z(V)}\operatorname{ind}(V,x)
=\chi(X).
$$
\end{theorem}

\begin{proof}
We give the course-level summary.  Choose a
triangulation or a Morse function and build a
vector field whose zeroes correspond to cells
or critical points.  The local index at each
zero gives the sign with which the cell
contributes.  The sum of these signed local
contributions is the Euler characteristic.
A homotopy argument shows that the sum is
independent of the chosen vector field.
\end{proof}

\section{Hopf theorem}
\label{differential-topology-section-Hopf-theorem}

\begin{theorem}[Hopf]
\label{differential-topology-theorem-Hopf}
Let $X$ be a compact, oriented, connected
$n$-manifold without boundary.  Let
$f,g:X\to S^n$ be smooth maps.  Then
$$
f\sim g
\quad\Longleftrightarrow\quad
\deg(f)=\deg(g).
$$
\end{theorem}

\begin{proof}
One direction is homotopy invariance of the
degree.  Conversely, choose a regular value
$p\in S^n$ for both maps.  The preimages
$f^{-1}(p)$ and $g^{-1}(p)$ are finite signed
sets.  Equality of degrees says that these
signed sets have the same algebraic count.
Using paths and local modifications, one
cancels pairs of opposite signs and moves the
remaining points so that the two maps have
matching inverse images of $p$.  Then the
maps are homotopic on the complement of small
balls and the local degree data fills in the
balls.
\end{proof}

\section{Pontryagin--Thom construction}
\label{differential-topology-section-pontryagin-thom}

\noindent
The Pontryagin--Thom idea is to convert a
submanifold with framing data into a map to a
sphere.  If $W^k\subset M^n$ has a framed
normal bundle, choose a tubular neighborhood
$\nu(W)\cong W\times\mathbb R^{n-k}$.  Collapse
the complement of this neighborhood to the
base point of $S^{n-k}$, and use the framing
to map each normal disk radially to the
sphere.  This produces a map
$$
M\to S^{n-k}.
$$
Cobordism of framed submanifolds corresponds
to homotopy of the resulting maps.

\section{Morse functions}
\label{differential-topology-section-morse-functions}

\begin{definition}
\label{differential-topology-definition-morse-function}
Let $M$ be a smooth manifold.  A smooth
function $f:M\to\mathbb R$ is a
{\it Morse function} if every critical point
is nondegenerate, i.e. the Hessian at that
point is a nondegenerate quadratic form.
\end{definition}

\begin{theorem}[Morse lemma]
\label{differential-topology-theorem-morse-lemma}
Let $p$ be a nondegenerate critical point of
$f:M\to\mathbb R$.  Then there are local
coordinates around $p$ in which
$$
f=f(p)-x_1^2-\cdots-x_\lambda^2
+x_{\lambda+1}^2+\cdots+x_n^2.
$$
The number $\lambda$ is the {\it Morse index}
of $p$.
\end{theorem}

\begin{proof}
The proof is a smooth diagonalization of the
quadratic part of $f$ near the critical
point.  After subtracting the constant
$f(p)$, the first-order terms vanish.  The
nondegenerate Hessian can be diagonalized,
and the remaining higher-order terms are
absorbed by a smooth change of coordinates.
\end{proof}

\begin{theorem}
\label{differential-topology-theorem-morse-functions-generic}
Morse functions are generic in
$C^\infty(M,\mathbb R)$.
\end{theorem}

\begin{proof}
A function fails to be Morse when its 1-jet
meets the zero section in a nontransverse way,
or equivalently when the corresponding
2-jet lies in the degenerate Hessian locus.
This is a jet transversality condition, so
Thom transversality gives the result.
\end{proof}

\section{Handle decomposition}
\label{differential-topology-section-handle-decomposition}

\noindent
Morse theory describes the topology of a
manifold by studying sublevel sets
$$
M^a=\{x\in M:f(x)\leq a\}.
$$
If there are no critical values in
$[a,b]$, then $M^a$ and $M^b$ are
diffeomorphic.  If one crosses a single
critical point of index $\lambda$, then the
sublevel set changes by attaching a
$\lambda$-handle
$$
D^\lambda\times D^{n-\lambda}.
$$

\begin{theorem}
\label{differential-topology-theorem-handle-attachment}
Let $f:M\to\mathbb R$ be Morse and suppose
$c$ is a critical value with one critical
point of index $\lambda$.  For small
$\varepsilon>0$, the sublevel set
$M^{c+\varepsilon}$ is obtained from
$M^{c-\varepsilon}$ by attaching a
$\lambda$-handle.
\end{theorem}

\begin{proof}
This is the local normal form from the Morse
lemma together with the gradient flow.  The
negative directions give the core
$D^\lambda$, and the positive directions give
the cocore $D^{n-\lambda}$.
\end{proof}

\section{Cancellation and four-manifolds}
\label{differential-topology-section-cancellation-four-manifolds}

\noindent
Handle decompositions are not unique.  A
central operation is cancellation: under the
right incidence condition, a
$\lambda$-handle and a $(\lambda+1)$-handle
can cancel each other.  Geometrically, this
corresponds to modifying the Morse function
so that a pair of critical points disappears.

\begin{theorem}[Cancellation principle]
\label{differential-topology-theorem-cancellation-principle}
If the attaching sphere of a
$(\lambda+1)$-handle intersects the belt
sphere of a $\lambda$-handle transversely in
one point, then the two handles cancel.
\end{theorem}

\begin{proof}
We only record the geometric idea.  The
single transverse intersection means that the
gradient trajectory between the two critical
points is unique and has the correct
orientation.  One can then change the Morse
function in a neighborhood of this trajectory
and remove the two critical points.
\end{proof}

\noindent
The course ended by pointing toward the
classification of 4-manifolds.  In dimension
4, handle decompositions are especially rich:
0-, 1-, 2-, 3-, and 4-handles interact with
intersection forms, framings, and Kirby
calculus.  This is the bridge from Morse
theory and differential topology to the
classification problems for smooth
4-manifolds.


% --- End differential-topology.tex ---

% --- Begin algebraic-topology.tex ---
\chapter{Algebraic Topology}
\label{chapter-algebraic-topology}






\medskip
\noindent
{\bf Organization of these notes.}
These notes combine material from several courses and sets of
lecture notes.  The first block develops homotopy theory of spaces
(coverings, homotopy groups, fibrations, cofibrations, and model
structures).  The second block develops the standard structural
results (Blakers--Massey, Freudenthal, Whitehead and Postnikov
towers, homotopy fibers and cofibers, and representability).
The final block introduces spectral sequences, especially the
Serre spectral sequence and its applications.  Some subsections
are still lecture-note fragments; these are left visibly marked
rather than silently completed.


\section{Covering spaces}
\label{algebraic-topology-section-covering-spaces}

\begin{proposition}[Unique lifting property]
\label{algebraic-topology-proposition-unique-lifting-property}
\cite{hat}, Proposition 1.34. Given a
covering space $p:\tilde{X} \to X$ and a
map $f:Y \to X$, if two lifts
$\tilde{f}_1,\tilde{f}_2:Y \to \tilde{X}$ of
$f$
agree at one point of $Y$ and $Y$ is
connected, then $\tilde{f}_1$ and
$\tilde{f}_2$ agree on all of $Y$.
\end{proposition}

\begin{proof}
We aim to show that the set of points where
$\tilde{f}_1$ and $\tilde{f}_2$
agree is both open and closed in $Y$. Let $y
\in Y$ and consider an evenly
covered neighbourhood $U$ of $f(y)$, so that
$p^{-1}(U)$ is a disjoint union of
sets, each of which is homeomorphic to $U$
via $p$. Let $\tilde{U}_1$ and
$\tilde{U}_2$ be the sheets containing
$\tilde{f}_1(y)$ and $\tilde{f}_2(y)$,
respectively. By continuity of $\tilde{f}_1$
and $\tilde{f}_2$, the intersection
$\tilde{f}^{-1}_1(\tilde{U}_1) \cap
\tilde{f}_2^{-1}(\tilde{U}_2):=N$ is
 an open neighbourhood of $y$.

Since  $\tilde{U}_1$ and $\tilde{U}_2$ are
mapped homeomorphically to the set
$U \subset X$, we can only have that either
$\tilde{U}_1$ coincides
with $\tilde{U}_2$, or they are disjoint. If
they are disjoint, then the
neighbourhood $N$ shows that the set of
points where $\tilde{f}_1$ and
$\tilde{f}_2$ agree is closed; that's because
in this case they don't agree at a
point, so that we are on the complement of
the set where they agree, and we have
shown it's an open set (every point has a
neighbourhood contained in the set).

If, on the other hand, $\tilde{U}_1$ and
$\tilde{U}_2$ are the same, which
happens when $\tilde{f}_1(y)=\tilde{f}_2(y)$,
then $\tilde{f}_1$ and
$\tilde{f}_2$ agree throughout $N$, so that
we have shown that the set where
$\tilde{f}_1$ and $\tilde{f}_2$ coincide is
open.
\end{proof}

\begin{definition}[Properly discontinuous
action]
\label{algebraic-topology-definition-prope-disco-actio}
See \cite{hat} p. 72. The action of a group
$G$ on a space $Y$ is called
{\it properly discontinuous} or {\it covering
space action} if each $y \in Y$
has a neighbourhood $U$ such that $g_1(U)
\cap g_2(U)\neq \emptyset$ for any
$g_1,g_2 \in G$ implies $g_1=g_2$.
\end{definition}

\begin{proposition}
\label{algebraic-topology-proposition-deck-trans-act-prope}
The action of deck transformations on a
covering space is a properly
discontinuous action.
\end{proposition}

\section{Definitions in homotopy theory}
\label{algebraic-topology-section-defin-homot-theor}

This is taken from \cite{TomDieck}.

\begin{itemize}
\item Let $X$ and $Y$ be topological spaces
and
 $f,g:X\to Y$ continuous maps.
An {\it homotopy} from $f$ to $g$ is a
continuous map
$$
H:X\times[0,1]\to Y,\qquad(x,t)\mapsto
H(x,t)=H_t(x)
$$
\noindent
such that $f(x)=H(x,0)$ and $g(x)=H(x,1)$ for
all $x\in X$.
We denote this situation by $f\simeq g$.
The homotopy relation $\simeq$ is an
equivalence relation
on the set of continuous maps $X\to Y$.
A homotopy of maps $H_t:X\to Y$ is called
{\it relative to $A\subset X$} if $H_t|_A$ is
constant.


\item Topological spaces and homotopy classes
of maps form a quotient category
of $\text{Top}$, the {\it homotopy category
$\text{hTop}$}, where composition
of homotopy classes is induced by composition
of representing maps. If $f:X\to
Y$ represents an isomorphism in
$\text{hTop}$, then $f$ is called a
{\it homotopy equivalence} or
{\it $\text{h}$-equivalence}. In explicit
terms this means $f:X\to
Y$ is a homotopy equivalence if there exists
$g:Y\to X$, a
{\it homotopy inverse of $f$}, such that $gf$
and $fg$ are both
homotopic to the identity. Spaces $X$ and $Y$
are called
{\it homotopy equivalent} or of the same {\it
homotopy
type} if there exists a homotopy equivalence
$X\to Y$. A space is
{\it contractible} if it is homotopy
equivalent to a point. A map
$f:X\to Y$ is {\it null homotopic} if it is
homotopic to a constant
map.

\end{itemize}

\section{n-th homotopy
group of a pointed space}
\label{algebraic-topology-section-n-th-homot-group-point-spac}

\begin{definition}
\label{algebraic-topology-definition-homotopy-group}
Let $(X,x_0)$ be a pointed topological space
and $s_0\in S^n$. The elements of
the {\it $n$-th homotopy group} are homotopy
classes of maps
$(S^n,s_0)\to (X,x_0)$. Equivalently, they
are homotopy classes of maps
$(I^n,\partial I^n)\to (X,x_0)$. (Homotopies
are required to preserve the base
points, $s_0\mapsto x_0$ or $\partial
I^n\mapsto x_0$.)

Also,
$$
\pi_n(X,*)=[(I^n,\partial
I^n),(X,\{*\})]\cong[I^n/\partial I^n,X]^0
$$
where $[X,Y]$ denotes the set of homotopy
classes $[f]$ of maps $[f]:X\to Y$.
\end{definition}

\section{n-th relative
homotopy group of
a pointed pair of spaces}
\label{algebraic-topology-section-relative-homotopy-group}

\noindent
Let $A$ be a subspace of $X$ and $x_0\in A$.
The elements of the {\it relative homotopy
group $\pi_n(X,A,x_0)$}
are homotopy classes of maps $(I^n,\partial
I^n,J^{n-1})\to (X,A,x_0)$
where $J^{n-1}$ is the union of all but one
face of $I^n$.
That is,
$$
\pi_{n+1}(X,A,*)=[(I^{n+1},\partial
I^{n+1},J^n),(X,A,x_0)].
$$

The elements of such a group are homotopy
classes of based maps
$D^n\to X$ which carry the boundary $S^{n-1}$
into $A$.
Two maps $f,g$ are called {\it homotopic
relative to $A$}
if they are homotopic by a
basepoint-preserving homotopy
 $F:D_n\times[0,1]\to X$ such that, for each
 $p$ in $S^{n-1}$ and $t$ in $[0,1]$,
the element $F(p,t)$ is in $A$. Ordinary
homotopy groups are
recovered for the case in which $A=\{x_0\}$.

\begin{remark}
\label{algebraic-topology-remark-relative-homotopy}
This construction is motivated by studying the kernel of the
map induced by inclusion
\[
i_*:\pi_n(A,x_0)\longrightarrow\pi_n(X,x_0).
\]
\emph{[The rest of this lecture remark was not recorded.]}
\end{remark}

\section{Commutativity of higher homotopy
groups}
\label{algebraic-topology-section-commu-highe-homot-group}

\begin{proposition}
\label{algebraic-topology-proposition-commu-highe-homot-group}
$\pi_n(X,x_0)$ is an abelian group for $n>1$.
\end{proposition}

\section{Action of fundamental groupoid
\texorpdfstring{$\Pi_1(X)$}{Π₁(X)} of
\texorpdfstring{$X$}{X} on
\texorpdfstring{$\pi_{n}(X)$}{π₁(X)}}
\label{algebraic-topology-section-actio-funda-group}

\begin{itemize}
\item A {\it groupoid} is a small category in
which every morphism is an isomorphism.
\item The {\it fundamental groupoid}
$\Pi_{1}(X)$
is the category whose objects
are the points of  $X$ and
morphisms are homotopy classes of paths
joining two points.
Composition is defined doubling the speed.
\end{itemize}

The fundamental groupoid acts on homotopy
group functorially.

\section{Homotopy groups of a product of
spaces
are products of homotopy groups of factors}
\label{algebraic-topology-section-homot-group-produ-space}

\section{Higher homotopy groups of a covering
space
are the same as those of a base space}
\label{algebraic-topology-section-highe-homot-group-cover}

\section{Long exact sequence of homotopy
groups associated to a pair}
\label{algebraic-topology-section-long-exact-seque-homot-grou}

For any pair $(X,A,x)$ we have a long exact
sequence
$$
\xymatrix@C=1.5em{
\pi_n(A,x_0)\ar[r]^{i_*}
& \pi_{n}(X,x_0)\ar[r]^{j_*}
& \pi_{n-1}(X,A,x_0)\ar[r]^{\partial}
& \pi_{n-1}(A,x_0)\ar[r]
& \cdots\ar[r]
& \pi_0(X,x_0)
}
$$
where $i$ and $j$ are the inclusions
$(A,x_0)\hookrightarrow(X,x_0)$ and
$(X,x_0,x_0)\hookrightarrow(X,A,x_0)$. The
map $\partial$ comes from restricting maps
$(I^n,\partial I^n,J^{n-1})\to (X,A,x_0)$ to
$I^{n-1}$, or by restricting maps
$(D^n,S^{n-1},s_0)\to (X,A,x_0)$. The map,
called the {\it boundary map}, is a
homomorphism when $n>1$.

\section{n-connected
space}
\label{algebraic-topology-section-n-connected-space}

A space $X$ with basepoint $x_0$ is called
{\it $n$-connected} if $\pi_i(X,x_0)=0$ for
$i\leq n$. Thus 0-connected means
path-connected and 1 connected means
simply-connected.

\section{n-connected
pair}
\label{algebraic-topology-section-n-connected-pair}

A pair $(X,A)$ is {\it $n$-connected} if
$\pi(X,A,x_0)=0$ for $i\leq n$.

\section{n-equivalence of spaces}
\label{algebraic-topology-section-n-equivalence-of-spaces}

A map $f:X\to Y$ is an {\it  $n$-equivalence}
if the induced map on homotopy groups
$f_{*}:\pi_{i}(X)\to \pi_{i}(Y)$ is an
isomorphism for $i<n$ and a surjection for
$i=n$.

\begin{remark}
\label{algebraic-topology-remark-what-i-picke-up-lectu-if}
\begin{itemize}
\item $(X,A)$ is $n$-connected if and only if
$A\hookrightarrow X$ is an $n$-equivalence.
\item $f:X\to Y$ is an $n$-equivalence if and
only if when we replace $f$ by a cofibration
$X\to M_{f}$, $(M_{f},X)$ is $n$-connected.
\end{itemize}
\end{remark}

\section{Compactly generated spaces}
\label{algebraic-topology-section-compactly-generated-spaces}

(What I picked up from the lecture.) If $X$
is compactly generated, then $X$ is
weakly Hausdorff if the diagonal subset
$\Delta_X\subset X\times X$ is
$k$-closed.

Now some other notes from the lectures:

In $\text{CGWH}$, $\Hom(X,Y)$ is a space with
the compact-open topology.
This is a compactly generated space,
$\mathbf{k}(\Hom(X,Y))$.

(Also see
\href{https://en.wikipedia.org/wiki/%
Currying#Function_spaces}{wiki
on currying})
$$
\begin{aligned}
\text{Map}(X,Y):=\text{ the space of maps
}X\to Y.\\
\text{Map}(X\times
Y,Z)\cong\text{Map}(X,\text{Map}(Y,Z))\\
\Hom(X\times Y,Z)\cong
\Hom(X,\text{Map}(Y,Z))
\end{aligned}
$$
In the last line, product is product in
$\text{CGWH}$, not in $\text{Top}$.
The functor $-\times Y$ is left adjoint to
$\text{Map}(Y,-)$.

We don't care so much about $\text{Top}$. We
care much more about $\text{CGWH}$,
the full subcategory of $\mathit{Top}$ on
 {\it compactly generated weakly Hausdorff}
 spaces.

\begin{definition}[Is this right?]
\label{algebraic-topology-definition-is-this-right}
$X$ is {\it compactly generated} if, for any
subset $C\subset X$, and for all continuous
maps $f:K\to X$ from compact Hausdorff
spaces, $$\text{if } f^{-1}(C) \text{is
closed in }K\text{, then } C\text{ is
closed}.$$
\end{definition}
{\bf Claim.}
If $X$ is compactly generated, then $X$ is
weakly Hausdorff if the diagonal subset
$\Delta_X\subset X\times X$ is $k$-closed.


{\bf Explanation from the book of May}
The ordinary category of spaces
allows pathology that obstructs a clean
development of the foundations.
The homotopy and homology groups of spaces
are supported on compact subspaces, and
it turns out that if one assumes a separation
property
that is a little weaker
than the Hausdorff property,
then one can refine the point-set topology of
spaces to eliminate
such pathology without changing these
invariants.

One major source of point-set level pathology
can be passage to quotient spaces.
Use of compactly generated topologies
alleviates this.

\begin{proposition}
\label{algebraic-topology-proposition-if-x-compa-gener-x-y}
If $X$ is compactly generated and $\pi:X\to
Y$ is a quotient map,
then $Y$ is compactly generated if and only
if $(\pi\times \pi)^{-1}(\Delta Y)$
 is closed in $X\times X$.
\end{proposition}

\noindent
The interpretation is that a quotient space
of a
compactly generated space by a “closed
equivalence relation”
is compactly generated.

\subsubsection{Weakly Hausdorff space}
\label{algebraic-topology-subsection-weakly-housdorff-space}

\subsection{Cartesian closedness of the
category $CGWH$}
\label{algebraic-topology-subsection-carte-close-categ-cgwh}
See "Compactly generated spaces" by Charles
Rezk.

\section{Wedge}
\label{algebraic-topology-section-wedge}

\section{Smash product of spaces, its
associativity in
\texorpdfstring{$\text{CGWH}$}{CGWH}}
\label{algebraic-topology-section-smash-produ-space-its-assoc}

The appropriate analogue of the Cartesian
product in the category of based spaces is
the {\it smash product} $X\wedge Y$ defined
by
$$X\wedge Y=X\times Y/X\vee Y.$$
Here $X\vee Y$ is viewed as the subspace of
$X\times Y$ consisting of those pairs $(x,y)$
such that either $x$ is the basepoint of $X$
or $y$ is the basepoint of $Y$.

\begin{example}[Basic example]
\label{algebraic-topology-example-basic-example}
	The smash product of two circles is a
	sphere, as may be seen by collapsing the
	generating circles on a revolution torus
	(see
	\href{https://en.wikipedia.org/wiki/%
	Smash_product#/media/File:
	Visualization_of_the_smash_product_of_%
	two_circles.gif}{Wikipedia}).
\end{example}

\section{Reduced and unreduced suspension}
\label{algebraic-topology-section-reduc-unred-suspe}
We also have the {\it suspension of pointed
spaces}, which is like usual suspension but
also collapsing the distinguished point,
which has become an interval:
$$\Sigma X=(I\times X)/(t,x)\sim (0,y)\sim
(1,y)\;\forall y\in X.$$

\section{Path and loop spaces of a pointed
space}
\label{algebraic-topology-section-path-loop-space-point-space}
[This is defined in Hatcher close to the
homotopy fiber construction--- it's related
to Blakers-Massey, Freudenthal theorem.]

\begin{definition}[Path fibration]
\label{algebraic-topology-definition-path-fibration}
An important case of the homotopy fiber
construction
is when the given function $f$ is the
inclusion of the basepoint $b_0$ into $B$.
Then $E_{f}$ is the space $PB$ of paths in
$B$ starting at $b_{0}$,
and $p:PB\to B$ sends each path to its
endpoint.
The fiber $p^{-1}(b_{0})$ is the loopspace
$\Omega B$
consisting of all loops in  $B$ based at
$b_{0}$.
Since $PB$ is contractible by progressively
truncating paths,
the long exact sequence of homotopy groups
for the path
fibration yields (another proof that)
$$
\pi_{n}(X,x_{0})\cong \pi_{n-1}(\Omega
X,x_{0})\qquad \forall n.
$$
\end{definition}

\section{For pointed spaces: smash with any
\texorpdfstring{$Y$}{Y}
is adjoint to
\texorpdfstring{$\text{Map}(Y,-)$}{Map(Y,-)}}
\label{algebraic-topology-section-point-space-smash-any-y}

[From
\href{https://en.wikipedia.org/wiki/%
Smash_product#}{Wikipedia}]
Adjoint functors make the analogy between the
tensor product
and the smash product more
precise. In the category of $R$-modules over
a commutative ring $R$, the tensor
functor $(-\otimes_{A} A)$ is left adjoint to
the internal Hom functor
$\Hom(A,-)$, so that
$$
\Hom(X\otimes A,Y)\cong\Hom(X,\Hom(A,Y)).
$$

In the category of pointed spaces, the smash
product plays the role the tensor product in
this formula: if $A$, $X$ are compact
Hausdorff then we have an adjunction
$$\text{Maps}_{*}(X\wedge A,Y)\cong
\text{Maps}_{*}(X,\text{Maps}_{*}(A,Y))$$
where $\text{Maps}_{*}$ denotes continuous
maps that send basepoint to basepoint, and
$\text{Maps}_{*}(A,Y)$ carries the
compact-open topology.

\section{Adjointness of Sigma and Omega}
\label{algebraic-topology-section-adjointness-of-and}
In particular, taking $A$ to be the unit
circle $S^{1}$, we see that the reduced
suspension functor $\Sigma$ is left adjoint
to the loop space functor $\Omega$:
$$\text{Maps}_{*}(\Sigma X,Y)\cong
\text{Maps}_{*}(X,\Omega Y).$$

[From lectures:] Then we have
$$\Hom_{\text{CGWH}_*}(\Sigma X,\Sigma
X)\cong\Hom_{\text{CGWH}_*}(X,\Omega\Sigma
X)$$
where $\Sigma X=S\wedge X$ and $\Omega\Sigma
X=\text{Map}(S^1,\Sigma X)$. That is,
$S^1\wedge-$ is adjoint to
$\text{Map}(S^1,-)$.

\section{Hurewicz homomorphism}
\label{algebraic-topology-section-Hurewicz-homomorphism}

\begin{definition}[Hurewicz homomorphism]
\label{algebraic-topology-definition-hurewicz-homomorphism}
$$
\begin{aligned}
h: \pi_{n}(X) &\longrightarrow H_{n}(X) \\
[f] &\longmapsto f_{*}(\alpha)
.\end{aligned}
$$
\end{definition}
where $\alpha$ is a generator of $
H_{n}(S^{n})$.

\section{Absolute and relative versions of
Hurewicz theorem}
\label{algebraic-topology-section-absol-relat-hurew-theor}

\begin{theorem}[Hurewicz]
\label{algebraic-topology-theorem-Hurewicz-absolute}
If $X$ is $n$-connected for some $n \geq 2$,
them $\tilde{H}_i(X)=0$ for $i<n$
and $\pi_{n}(X)\cong H_n(X)$.
\end{theorem}

\begin{lemma}
\label{algebraic-topology-lemma-hurew-homom-gives-homot-equiv}
If $X$ and $Y$ are simply connected CW-spaces
and  $f:X \to Y$ induces
isomorphisms $H_n(X) \cong H_n(Y)$ for all
$n$, then it is an homotopy
equivalence.
\end{lemma}

\begin{theorem}[Hurewicz]
\label{algebraic-topology-theorem-Hurewicz-relative}
If $(X,A)$ is $(n-1)$-connected, i.e.
$\pi_{n}(X,A)=0$ for $i \leq n-1$, for
some $n\geq 2$ and $A$ is simply connected,
then $H_i(X,A)=0$ for $i\leq n-1$
and $\pi_{n}(X,A)\cong H_n(X,A)$.
\end{theorem}
Perhaps $X$ should be $(n-1)$-connected for
the relative version too---
they are stated together in Hatcher.

\section{Weak factorization systems}
\label{algebraic-topology-section-weak-factorization-systems}

Here we follow \cite{Riehl}.

\begin{definition}
\label{algebraic-topology-definition-weak-factorization-system}
A {\it weak factorization system
$(\mathcal{L},\mathcal{R})$} on a category
$\mathcal{M}$ is comprised of two clases of
morphisms $\mathcal{L}$
and $\mathcal{R}$ so that

\begin{enumerate}
\item Every morphism in $\mathcal{M}$ may be
factored as a morphism in $\mathcal{L}$
followed by a morphism in $\mathcal{R}$:
$$
\xymatrix{
\bullet\ar[rr]^{f}\ar[rd]_{\mathcal{L}
\ni\ell}
&& \bullet\\
& \bullet\ar[ru]_{r\in\mathcal{R}}
}
$$

\item The maps in $\mathcal{L}$ have the {\it
left lifting property} with respect to each
map in $\mathcal{R}$ and equivalently the
maps in $\mathcal{R}$ have the {\it right
lifting property} with respect to each map in
$\mathcal{L}$, that is, any commutative
square
$$
\xymatrix{
\bullet\ar[d]_{\mathcal{L}\ni\ell}\ar[r]
& \bullet\ar[d]^{r\in\mathcal{R}}\\
\bullet\ar@{-->}[ur]\ar[r]
& \bullet
}
$$
admits a diagonal filler as indicated making
both triangles commute.
When this lift is unique,
we say the factorization system is {\it
orthogonal}.

\item The classes $\mathcal{L}$ and
$\mathcal{R}$ are each closed under
retracts in the arrow category:
given a commutative diagram
$$
\xymatrix{
\bullet\ar[r]\ar[d]_{t}\ar@/^1pc/@{-}[rr]
\ar@/^2pc/@{-}[rr]
& \bullet\ar[r]\ar[d]^{s}
& \bullet\ar[d]^{t}\\
\bullet\ar[r]\ar@/_1pc/@{-}[rr]\ar@/_2pc/@{-}
[rr]
& \bullet\ar[r]
& \bullet
}
$$

\noindent
if $s$ is in that class then so is its
retract $t$.
\end{enumerate}
\end{definition}

\begin{exercise}[3.1.8, Riehl]
\label{algebraic-topology-exercise-3-1-8-rihel}
Verify that the class of morphisms
$\mathcal{L}$ characterized
by the left lifting property against a fixed
class of morphisms $\mathcal{R}$ is closed
under coproducts,
closed under retracts, and contains the
isomorphisms.
\end{exercise}

\section{Model category structure on a
category}
\label{algebraic-topology-section-model-category-structure}

\begin{definition}[Lecture]
\label{algebraic-topology-definition-lecture}
A {\it model structure} on a category
$\mathcal{A}$
is a choice of subcategories
$\mathcal{W},\mathcal{C},\mathcal{F}$
called {\it weak-equivalences}, {\it
cofibrations} and {\it fibrations}
with the following properties:

\begin{enumerate}
\item[0.] All (finite) small limits an
colimits.

\item \textbf{(2 of 3)} Given
$\bullet\overset{f}{\to}\bullet\overset{g}
{\to}\bullet$,
if either 2 out of 3 among $f,g,f\circ g$
are in $\mathcal{W}$ then all of them are.

\item
$(\mathcal{C}\cap\mathcal{W},\mathcal{F})$
and
$(\mathcal{C},\mathcal{F}\cap\mathcal{W})$
are both weak factorization systems.

$(\mathcal{B},\mathcal{D})$ is a weak
factorization system. That is,

\begin{enumerate}
\item Any morphism in $\mathcal{A}$ can be
factored
as a morphism in $\mathcal{B}$
followed by a morphism in $\mathcal{D}$.

\item Lifts:

$$
\xymatrix{
\bullet\ar[d]_{f}\ar[r]&\bullet\ar[d]^{g}\\
\bullet\ar@{-->}[ur]^{\exists}\ar[r]&\bullet
}
$$

\item[(c')] Notice that the axiom of retracts
is not necessary. $r'\in\mathcal{R}\iff$ it
satisfies (b) for all $\ell\in\mathcal{L}$.
\end{enumerate}
\end{enumerate}
\end{definition}

\begin{definition}
\label{algebraic-topology-definition-item}
\begin{itemize}
\item $X$ is {\it fibrant} if
$X\to\text{pt}$.
\item $X$ is {\it cofibrant} if
$$
\xymatrix{
X\ar@{->>}[r]&X
}
$$
\item $X$ is {\it bifibrant} if
$$
\xymatrix{
0\ar@{^{(}->}[r]&X\ar@{->>}[r]&\text{pt}
}
$$
\end{itemize}
\end{definition}

\medskip\noindent
Two examples of model structures
are Str\o m and Quillen model structures.

\begin{example}[Two interesting model
category structures on $\text{CGWH}$]
\label{algebraic-topology-example-two-inter-model-categ-struc}
\begin{enumerate}
\item {\it Hurewicz-Str\o m model structure}.
\begin{itemize}
\item Cofibrations = Huerwicz cofibrations.
\item Fibrations = Hurewicz fibrations.
\item Weak equivalences = homotopy
equivalences.
\end{itemize}
\item {\it Quillen model structure}. Defined
on $\text{Top}$.
\begin{itemize}
\item Cofibrations = "retracts of relative
cell complexes".
\item Fibrations = Serre fibrations.
\item Weak equivalences = weak homotopy
equivalences.
\end{itemize}
\end{enumerate}
\end{example}

\subsection{Hurewicz fibrations and
cofibrations}
\label{algebraic-topology-subsection-hurew-fibra-cofib}

\section{Fibrations}
\label{algebraic-topology-section-fibrations}

\begin{definition}[Hatcher]
\label{algebraic-topology-definition-hatcher}
A map $p:E\to B$ is said to have the {\it
homotopy lifting property} with respect to a
space $X$ if, given a homotopy $g_t:X\to B$
and a map $\tilde{g}_0:X\to E$ lifting $g_0$,
so $p\tilde{g}_0=g_0$, then there exists a
homotopy $\tilde{g}_t:X\to E$ lifting $g_t$.

The {\it lift extension property for a pair
$(Z,A)$} asserts that every map $X\to B$ has
a lift $Z\to E$ extending a given lift
defined on the subspace $A\subset Z$. The
case $(Z,A)=(X\times I,X\times\{0\})$ is the
homotopy lifting property.

A {\it fibration} is a map $p:E\to B$ having
the homotopy property with respect to all
spaces $X$.
\end{definition}

\begin{definition}[\href{https://ncatlab.org%
/nlab/show/lift}{nLab}]
\label{algebraic-topology-definition-https-ncatl-org-nlab}
	A morphism $i$ has the {\it left lifting
	property with respect to a morphism $p$} and
	$p$ has the {\it right lifting property with
	respect to $i$} if for each morphisms $f$
	and $g$, if the outer square in the
	following diagram commutes, there exists
	$\phi$ (I think not necessarily unique)
	completing the diagram:
	$$
	\xymatrix@R=2.5em{
		A\ar[r]^{f}\ar[d]_{i}&X\ar[d]^{p}\\
		B\ar[r]_{g}\ar@{-->}[ur]^{\phi}&Y
	}
	$$
\end{definition}

\begin{definition}[\href{https://ncatlab.org%
/nlab/show/homotopy+lifting+property}{nLab}]
Let $C$ be a category with products and with
interval object $I$. A morphism $E\to B$ has
the {\it homotopy lifting property} if it has
the right lifting property with respect to
all morphisms of the form $(\text{id},0):Y\to
Y\times I$.
\label{algebraic-topology-definition-homot-lifti-prope-nlab}

	This means that for all commuting squares
	$$
	\xymatrix{
		Y\ar[r]^{f}\ar[d]&E\ar[d]^{p}\\
		Y\times I\ar[r]_{F}\ar@{-->}[ur]^{\sigma}&B
	}
	$$
	there exists a morphism $\sigma:Y\times I\to
	E$ such that both triangles in the former
	diagram commute.

	A {\it fibration} (also called {\it Hurewicz
	fibration}) is a mapping $p:E\to B$
	satisfying the homotopy lifting property for
	all spaces $X$.
\end{definition}

\subsubsection{Cofibrations}
\label{algebraic-topology-subsection-cofibrations}

\begin{definition}[Hatcher]
\label{algebraic-topology-definition-hatcher-2}
	If in the definition of a fibration as a map
	satisfying the homotopy lifting property we
	reverse the direction of the arrows, we
	obtain the dual notion of a {\it
	cofibration}. This is a map $i:A\to B$
	(which is $E\to B$ backwards) satisfying the
	following property: (we used to have an
	homotopy from some space $X$ into the base,
	now we have:) given an homotopy $g_{t}:A\to
	X$ and (a lift)  $\tilde{g}:B\to X$ such
	that $\tilde{g}_{0}i=g_{0}$, there exists a
	homotopy $\tilde{g}_{0}:B\to X$ such that
	$\tilde{g}_{t}i=g_{t}$.

	In the special case that $i$ is the
	inclusion of a subspace, this is the
	homotopy extension property, and the next
	proposition says that this is indeed the
	general case. So a cofibration is an
	inclusion satisfying the homotopy extension
	property.
\end{definition}

\begin{proposition}[4H.1]
\label{algebraic-topology-proposition-4h-1}
	If $i:A\to B$ is a cofibration, then $i$ is
	injective, and in fact a homeomorphism onto
	its image.
\end{proposition}

\begin{remark}[From Hatcher p. 461]
\label{algebraic-topology-remark-from-hatcher-p-461}
The duality between subobjects and quotient
objects is clear for abelian groups,
where subobjects are kernels and quotient
objects cokernels. The strict
topological analog of a kernel is a fiber of
a fibration. Dually, the
topological analog of a cokernel is the {\it
cofiber} $B/A$ of a
cofibration $A\hookrightarrow B$. If we make
an arbitrary map  $f:A\to B$ into a
cofibration $A\hookrightarrow M_{f}$, the
cofiber is the mapping cone
$C_{f}=M_{f}/(A\times {0}$.
\end{remark}

\begin{definition}[\href{https://%
en.wikipedia.org/wiki/
Homotopy_extension_property}{wiki}]
\label{algebraic-topology-definition-https-en-wikip-org-wiki}
	Let $X$ be a topological space and let
	$A\subset X$. We say that the pair $(X,A)$
	has the {\it homotopy extension property}
	if, given a homotopy $f_\bullet:A\to Y^I$
	and a map $\tilde{f}_0:X\to Y$ such that
		$$\tilde{f}_0\circ\iota=f_0$$\
		(so $\tilde{f}$ is the lift of $f_0:A\to
		Y$) then there exists an {\it extension} of
		$f_\bullet$ to a homotopy
		$\tilde{f}_\bullet:X\to Y^I$ such that
		$\tilde{f}_\bullet\circ\iota=f_\bullet$.

		That is,
		$$
		\xymatrix{
			A\ar[r]^{f_\bullet}\ar[d]_{\iota}
			&Y^I\ar[d]^{\pi_0}\\
			X\ar@{-->}[ur]^{\tilde{f}_\bullet}
			\ar[r]_{\tilde{f}_0}&Y
		}
		$$
		So there's some
		\href{https://en.wikipedia.org/wiki/%
		Currying#Function_spaces}{currying}
		to make usual homotopies $f_\bullet:A\times
		I\to Y$ look like $f_\bullet:A\to Y^I$. As
		said in our lectures: ``a homotopy $X\times
		I\to Y$ is the same as a map $X\to
		\text{Map}(I,Y)$".

	A map $\iota:A\to X$ is a {\it cofibration}
	if it has the homotopy extension property
	with respect to all topological spaces  $Y$.
\end{definition}

\subsubsection{Mapping cylinder}
\label{algebraic-topology-subsection-mapping-cylinder}

\begin{definition}[\href{https://ncatlab.org%
/nlab/show/mapping+cylinder}{nLab}]
\label{algebraic-topology-definition-https-ncatl-org-nlab-2}
Given a continuous map $f:X\to Y$ of
topological spaces, one can define its {\it
mapping cylinder} as a pushout
$$
\xymatrix{
		X\ar[r]^{f}\ar[d]&Y\ar[d]\\
		X\times
		I\ar[r]_{(\sigma_0)_*(f)}&\text{Cyl}(f)
	}
$$
	in $\text{Top}$, where $I=[0,1]$ and
	$\sigma:X\to X\times I$ is given by
	$x\mapsto(x,0)$.

	Set theoretically, the mapping cyllinder is
	usually represented as que quotient space
	$$\left(X\times I\coprod Y\right)\Big/\sim$$
	where $\sim$ is the smallest equivalence
	relation identifying $(x,0)\sim f(x)$ for
	all $x\in X$.
\end{definition}

\begin{definition}[\href{https://%
en.wikipedia.org/wiki/Mapping_cylinder}{wiki}
]
\label{algebraic-topology-definition-https-en-wikip-org-wik-2}
The {\it mapping cylinder} of a function $f$
between topological spaces $X$ and $Y$ is the
quotient
$$M_f=(([0,1]\times X)\amalg Y)\big/\sim$$
where $\amalg$ denotes disjoint union, and
$\sim$ is the equivalence relation generated
by
$$(0,x)\sim f(x)\text{ for each }x\in X.$$
{\it That is, the mapping cylinder $M_f$ is
obtained by gluing one end of $X\times[0,1]$
to $Y$ via the map $f$.} Notice that the
``top" of the cylinder $\{1\}\times X$ is
homeomorphic to $X$, while the ``bottom" is
the space $f(X)\subset Y$.

Looks the mapping cylinder is just deforming
$X$ to $Y$ putting $X$ inside $Y$ via $f$…
\end{definition}

\begin{definition}[From homework]
\label{algebraic-topology-definition-from-homework}
	Let $f:X\to Y$ be a map. Let
	$M_f=X\times[0,1]\cup_fY$ be the {\it
	mapping cylinder of $f$}, i.e. the pushout
	of
	$X\overset{\cong}{\to}X\times\{0\}
	\hookrightarrow
	X\times[0,1]$ and of $f:X\times Y$.
\end{definition}

\subsection{Serre fibrations}
\label{algebraic-topology-subsection-serre-fibrations}

\begin{definition}[Hatcher, p. 376]
\label{algebraic-topology-definition-hatcher-p-376}
		The map $p:E\to B$ is said to have the {\it
		homotopy lifting property for a pair}
		$(X,A)$ if each homotopy $f_t:X\to B$ lifts
		to a homotopy $\tilde{g}_t:X\to E$ starting
		with a given lift $\tilde{g}_0$ and
		extending a given lift $\tilde{g}_t:A\to
		E$. In other words, the homotopy lifting
		property for $(X,A)$ is the lift extension
		property for $(X\times I,X\times\{0\}\cup
		A\times I)$.

	The important observation here is that
	homotopy lifting property for disks is
	equivalent to the homotopy lifting property
	for all CW pairs $(X,A)$. A map $p:E\to B$
	satisfying the homotopy lifting property for
	disks is sometimes called a {\it Serre
	fibration}.
\end{definition}

\subsection{Fiber bundles}
\label{algebraic-topology-subsection-fiber-bundles}

Here I include some more notes of the
discussion following Serre fibrations on
Hatcher.

	A {\it fiber bundle} structure on a space
	$E$, with fiber $F$, consists of a
	projection map $p:E\to B$ such that each
	point $B$ has a neighbourhood $U$ for which
	there is a homeomorphism $h:p^{-1}(U)\to
	U\times F$ making the following diagram
	commute
	$$
	\xymatrix{
		p^{-1}(U)\ar[rr]^{h}\ar[dr]_{p}
		&&U\times F\ar[dl]\\
		&U
	}
	$$
\begin{example}
\label{algebraic-topology-example-proje-space-yield-inter}
Projective spaces yield interesting fiber
bundles. In the real case we have the
familiar covering spaces $S^n\to\mathbb{R}
P^n$ with fiber $S^0$. Over the complex
numbers the analog of this is a fiber bundle
$S^1\to S^{2n+1}\to\mathbb{C} P^n$. Here
$S^{2n+1}$ is the unit sphere in
$\mathbb{C}^{n+1}$ and $\mathbb{C} P^n$ is
viewed as the quotient space of $S^{2n+1}$
under the equivalence relation
$(z_0,\ldots,z_n)\sim\lambda(z_0,\ldots,z_n)$
for $\lambda\in S^1$. The projection
$p:S^{2n+1}\to\mathbb{C} P^n$ sends
$(z_0,\ldots,z_n)$ to its equivalence class
$[z_0,\ldots,z_n]$.

To see that the local triviality condition
for fibre bundles is satisfied, …

	The constructino of the bundle $S^1\to
	S^{2n+1}\to\mathbb{C} P^n$ also works when
	$n=\infty$, so there is a fiber bundle
	$S^1\to S^\infty\to \mathbb{C} P^\infty$.

	The case $n=1$ is particularly interesting
	since $\mathbb{C} P^1=S^2$ and bundle
	becomes $S^1\to S^3\to S^2$ with fiber,
	total space, and base all speres. This is
	known as the {\it Hopf bundle}. The
	projection $S^3\to S^2$ can be taken to be
	$(z_0,z_1)\mapsto
	z_0/z_1\in\mathbb{C}\cup\{\infty\}=S^2$.
	(That is, seeing $S^2$ as the one-point
	compactification of $\mathbb{C}$.)

	In polar coordinates we may see $S^3$ as the
	union of several tori. Stereorgraphic
	projection yields the following figure:
	% Figure omitted: hopf-bundle.
	The limiting cases $T_0$ and $T_\infty$
	correspond to the unit circle in the
	$xy$-plane and the $z$-axis under the
	stereographic projection.	Each torus
	$T_\rho$ is aunion of circle fibers. These
	fiber circles have slope 1 on the torus,
	winding around once longitudinally and once
	meridionally. As $\rho$ goes to $0$ or
	$\infty$ the fiber circles approach the
	circles $T_0$ and $T_\infty$, which are also
	fibers. The figure below shows four tori
	decomposed into fibers:
	% Figure omitted: hopf-bundle2.

	Replacing the field $\mathbb{C}$ by the
	quaternions $\mathbb{H}$, the same
	constructions yield fiber bundles $S^3\to
	S^{4n+3}\to\mathbb{H} P^{n}$ over
	quaternionic projective spaces $\mathbb{H}
	P^n$. Here the fiber $S^3$ is the unit
	quaternions, and $S^{4n+3}$ is the unit
	sphere in $\mathbb{H}^{n+1}$. Taking $n=1$
	gives a second Hopf bundle $S^3\to S^7\to
	S^4=\mathbb{H} P^1$.

	Another Hopf bundle $S^7\to S^{15}\to S^8$…
\end{example}

\subsection{Homotopy equivalences and weak
homotopy equivalences}
\label{algebraic-topology-subsection-homot-equiv-weak-homot}

\begin{definition}[Weak homotopy equivalence]
\label{algebraic-topology-definition-weak-homotopy-equivalence}
	Any map $f:X\to Y$ that induces isomorphisms
	on all homotopy groups ( $\forall n\geq 0$
	and all choices of basepoint in $X$).
\end{definition}

\subsection{Any Hurewicz fibration is Serre
fibration}
\label{algebraic-topology-subsection-any-hurew-fibra-serre}
\subsection{Any relative CW-complex is a
cofibration in Quillen model structure}
\label{algebraic-topology-subsection-any-relat-cw-compl-cofib}
\subsection{Any map can be replaced by…}
\label{algebraic-topology-subsection-any-map-can-be-replaced-by}
\subsubsection{an homotopy equivalence
followed by Hurewicz fibration}
\label{algebraic-topology-subsection-homot-equiv-follo-hurew}
\subsubsection{a Hurewicz cofibration
followed by homotopy equivalence}
\label{algebraic-topology-subsection-hurew-cofib-follo-homot}
\subsubsection{a Quillen cofibration followed
by weak equivalence}
\label{algebraic-topology-subsection-quill-cofib-follo-weak}
This one follows from CW-approximation
theorem

\section{Homotopy theory of spaces (B)}
\label{algebraic-topology-section-homotopy-theory-of-spaces-b}

\subsection{Blakers-Massey theorem}
\label{algebraic-topology-subsection-blakers-massey-theorem}

As in Hatcher and in terms of homotopy
pullbacks and pushouts. (Be careful if you
refer to REsk's paper for the latter version,
since in his paper the $n$ in definition of
$n$-connectiviy is shifted by $1$.
\begin{theorem}[Blakers-Massey. Hatcher 4.23,
Tom Dieck 6.4.1]
\label{algebraic-topology-theorem-blake-masse-hatch-4-23-tom}
	Let $X$ be a CW complex decomposed as the
	union of subcomplexes $A$ and $B$ with
	nonempty connected intersection $C=A\cap B$
	(that is also a CW complex?). If $(A,C)$ is
	$m$-connected and $(B,C)$ is $n$-connected,
	$m,n\geq 0$, then the map
	$$\pi_{i}(A,C)\to \pi_{i}(X,B)$$
is an isomorphism for $i<m+n$ and a
surjection for $i=m+n$.
\end{theorem}

\begin{proposition}[Quotient theorem.
Hatcher, 4.28and Tom Dieck 6.10.2]
\label{algebraic-topology-proposition-quoti-theor-hatch-4}
	If a CW pair $(X,A)$ is $r$-connected and
	$A$ is $r,s \geq 0$, then the map
	$\pi_{i}(X,A)\to \pi_{n}(X/A)$ induced by
	the quotient map $X\to X/A$ is an
	isomorphism for $i\leq r+s$ and a surjection
	for $i=r+s+1$.
\end{proposition}
\begin{proof}
	Using cones.
\end{proof}

\subsection{Freudenthal suspension theorem}
\label{algebraic-topology-subsection-freud-suspe-theor}

\begin{lemma}[Freudenthal suspension theorem,
Hatcher 4.24]
\label{algebraic-topology-lemma-freud-suspe-theor-hatch-4}
	The suspension map $\pi_{i}(S^{n})\to
	\pi_{i+1}(S^{n+1})$ is an isomorphism for
	$i<2n-1$ and a surjection for $i=2n-1$. More
	generally, this holds for the suspension
	$\pi_{i}(X)\to \pi_{i+1}(SX)$ whenever $X$
	is an $(n-1)$-connected CW complex.
\end{lemma}
\begin{theorem}[Freudenthal suspension
theorem, lecture notes]
\label{algebraic-topology-theorem-freud-suspe-theor-lectu}
		Let $X$ be $\ell$-connected space. Then
		$X\to\Omega\Sigma X$ is a
		$(2\ell+1)$-equivalence, that is,
	$$\pi_i(X)\to \pi_{i+1}(\Sigma X),$$
	is a bijection for $i<2\ell+1$ and a
	surjection for $i=2\ell+1$..
\end{theorem}
\begin{proof}[Proof 2 from lectures]
	Consider
	$$
	\xymatrix{
		X\ar[r]\ar[d]&CX\ar[d]\\
		CX\ar[r]&\Sigma X
	}
	$$
	Then we use relative homotopy long exact
	sequence with $(X,CX)$ to get
	$\pi_i(CX,X)\cong\pi_{i-i}(X)$, which is
	zero for $0\leq i\leq \ell+1$. Then use
	relative homotopy exact sequence for the
	pair $(\Sigma X,CX)$. then we get that
	$\pi_i(\Sigma X,CX)=\pi_i(\Sigma X)$. And
	then if you use it for $(\Sigma X, X)$ and

	But it also turns out that $\pi_i(\Sigma
	X)=\pi_{i-1}(\Omega\Sigma X)$ because
	$$
\begin{aligned}
		\pi_n(Z)&=\Hom_{\text{hTop},*}(S^n,Z)=
		\Hom(S^1\wedge
		S^{n-1},Z)=\text{Hom}(S^{n-1},\Omega
		Z)=\pi_{n-1}(\Omega,Z)
	\end{aligned}
$$.
	And then since $CX\hookrightarrow \Sigma X$
	we get an arrow $\pi_i(CX,X)\to\pi_i(\Sigma
	X,CX)$ which is isomorphism for $0\leq i
	\leq 2\ell +1$ and surjective for
	$i=2\ell+2$.

	So apply Blakers-Massey an ell equalities to
	get maps fro
	$\pi_{i-1}(X)\to\pi_{i-1}(\Omega\Sigma X)$
	for $i$ as desired.
\end{proof}
\begin{lemma}
\label{algebraic-topology-lemma-if-x-conne-then-x-1-conne}
	If $X$ is $\ell$-connected, then $\Sigma X$
	is $(\ell+1)$-connected for $\ell\geq0$.
\end{lemma}
$$\begin{matrix}
	\text{space}&&S^0&\Sigma
	S^0=S^1&\Sigma^2S^0=S^2&\Sigma^3S^0=S^3&
	\cdots&\Sigma^nS^0=S^n\\
	\text{conectedness}&&-1&0&1&2&\cdots&(n-1)
\end{matrix}$$

\subsubsection{\texorpdfstring{$\pi_{i}(S^{n}
)=0$}{πᵢ(Sⁿ)}
for \texorpdfstring{$i<n$}{i<n}}
\label{algebraic-topology-subsection-s-for-i-n}

\begin{lemma}
\label{algebraic-topology-lemma-s-n-is-n-1-connected}
	$S^n$ is $(n-1)$-connected.
\end{lemma}
Back to Hopf fibration:
$$S^1\hookrightarrow S^3\to S^2$$
we get
$$0=\pi_2(S^3)\to\pi_2(S^2)\overset{\cong}
{\to}\pi_1(S^1)\to\pi_1(S^3)=0,$$
so
$$\mathbb{Z}=\pi_2(S^2).$$
Now consider a map $S^n\to S^n$. We get a map
$CS^n\to CS^n$ (in general, for $f:X\to Y$ we
have $(t,x)\mapsto(t,f(x))$ in the cones). We
also have $CS^n\to CS^n/S^n=S^{n+1}$.

Now if we take $\text{id}:S^n\to S^n$ we
shall get $\text{id}:S^{n+1}\to S^{n+1}$.
Think about this like
$\pi_1(S^1)\to\pi_2(S^2)$. Now from
Freudenthal we get
$\pi_{i-1}(X)\to\pi_i(\Sigma X)$ is
surjective because $i=0$. From Hopf fibration
we have $\pi_2(S^2)=\mathbb{Z}$. So we have a
surjective map $\mathbb{Z}\to\mathbb{Z}$. So
it's an isomorphism and we conclude that
$\text{id}_{S^2}$ is a generator of
$\pi_2(S^2)$.
\begin{lemma}
\label{algebraic-topology-lemma-since-s-n-n-1-conne-we-have}
	Since $S^n$ is $(n-1)$-connected, we have
	$$\pi_i(S^n)\to\pi_{i+1}(S^{n+1})$$
	is isomorphism for $i\leq 2(n-1)=2n-1$ and
	epimorphism form $i=2n-1$. (We just shift
	the indices of the Freudenthal theorem by
	one.)
\end{lemma}
\begin{lemma}
\label{algebraic-topology-lemma-n-s-n-z-id-s-n-as-gener}
	$\pi_n(S^n)=\mathbb{Z}$ with
	$\text{id}_{S^n}$ as generator.
\end{lemma}

\subsubsection{$\pi_{n}(S^{n})=\mathbb{Z}$}
\label{algebraic-topology-subsection-n-s-n-z}
\begin{lemma}
\label{algebraic-topology-lemma-n-k-s-n-n-k-1-s-n-1-isomo}
	$\pi_{n+k}(S^n)\to\pi_{n+k+1}(S^{n+1})$ is
	isomorphism for $k\leq n-1$ and epimorphism
	for $k=n-1$.
\end{lemma}
So for example
$$\pi_4(S^3)=\pi_5(S^4)=\pi_6(S^5).$$
And in fact they are $\mathbb{Z}/2$. This is
what are called the {\it $k$th stable
homotopy groups of a sphere}. And more in
general, we take any space and apply
$\Omega\Sigma$ enough times, and the homotopy
will start to stabilize.

Or for example from
$$S^1\hookrightarrow S^3\to S^2$$
we get
$$0=\pi_i(S^1)\to
\pi_i(S^3)\overset{\cong}{\to}\pi_i(S^2)
\to\pi_{i-1}(S^2)=0$$
So $\pi_3(S^2)\cong\mathbb{Z}$ in case you
were wondering.
{\bf Claim.}[Serre]
	$\pi_{4n-1}(S^{2n})\cong\mathbb{Z}
	\oplus$finite
	abelian. And $\pi_i(S^k)$ is finite abelian
	in all other cases.

\subsection{Whitehead theorem that weak
equivalence between CW-complexes is a
homotopy equivalence}
\label{algebraic-topology-subsection-white-theor-weak-equiv}

\begin{theorem}[Whitehead]
\label{algebraic-topology-theorem-whitehead}
	Weak homotopy equivalence for CW complexes
	is homotopy equivalence.

	An inclusion inducing isomorphism on all
	homotopy groups is a deformation retract.
\end{theorem}


and related theorem that weak equivalence
between bifibrant objects is a homotopy
equivalence.
\subsection{Cellular approximation theorem}
\label{algebraic-topology-subsection-celul-appro-theor}

\begin{definition}[Cellular map]
\label{algebraic-topology-definition-cellular-map}
	A map between CW complexes $f:X\to Y$ is
	called {\it cellular} if $f(X^{n})\subset
	Y^{n}$ for all $n$.
\end{definition}
\begin{theorem}
\label{algebraic-topology-theorem-every-map-betwe-cw-compl}
	Every map between CW complex is homotopic to
	a cellular map.
\end{theorem}
\begin{theorem}[CW approximation]
\label{algebraic-topology-theorem-cw-approximation}
	Every space has a CW approximation, ie. it
	is weakly equivalent to a CW-complex.
	If $X$ is path-connected, the CW-complex can
	be chosen to have a single 0-cell.

	If $(X,A)$ is an $n$-connected CW-pair (so
	$A$ is a sub-CW-complex of $X$, made up of
	cells from $X$) there exists a CW-pair
	$(Z,A)\simeq (X,A)\text{rel}A$ with cells in
	$Z\setminus A$ of dimension greater than
	$n$.
\end{theorem}


\section{Postnikov towers}
\label{algebraic-topology-section-postnikov-towers}
Let $X$ be a CW complex. There exist spaces
$$
\xymatrix{
&\vdots\ar[d]\\
&Y_2\ar[d]\\
&Y_1\ar[d]\\
X\ar[uur]\ar[ur]\ar[r]&Y_0
}
$$
that are like the original space
(I think this means that the arrows in the
above diagram induce isomorphisms in
homotopy groups for smaller indices)
but with trivial homotopy groups for $k\geq
n$.
They are obtained by attaching to the
previous one
(or to $X$ itself?)
cells of dimension $n+2$
(because this kills the $n+1$-th homotopy
group
so that only the $n$-th homotopy group
remains nontrivial).

The inclusions
$X\hookrightarrow Y_{n}\hookrightarrow
Y_{n-1}\hookrightarrow
\ldots\hookrightarrow Y_{1}=K(\pi_{1}(X),1)$
 are converted to fibrations,
and in fact the fiber of $Y_{k}\to Y_{k-1}$
is a $K(\pi_{k}(X),k)$-space.

\begin{remark}
\label{algebraic-topology-remark-postn-tower-some-kind-facto}
It looks like Postnikov towers posses an
interesting property
of being a factorization system for some
category.
See
\url{https://ncatlab.org/nlab/show/Postnikov%
+system}.
Recall that this means that the maps in that
category
will be factored through maps in the
subcategories of right
and left morphisms provided by the Postnikov
tower.
\end{remark}


\section{Whitehead towers}
\label{algebraic-topology-section-whitehead-towers}

Now following Laurentium Maxim,
"Lecture notes on homotopy theory and
applications".

\begin{definition}[\href{https://%
people.math.wisc.edu/~lmaxim/754notes.pdf}
{754notes.pdf}]
\label{algebraic-topology-definition-https-peopl-math-wisc}
A {\it Whitehead tower} of $X$ is a sequence
of fibrations
$$
\xymatrix{
\cdots\ar[r]&X_n\ar[r]&X_{n-1}\ar[r]&\cdots
\ar[r]&X_0=X
}
$$
such that
\begin{itemize}
\item $X_n$ is $n$-connected.
\item $\pi_{q}(X_{n})=\pi_{q}(X)$ for $q\geq
n+1$.
\item The fiber of $X_n\to  X_{n-1}$ is a
$K(\pi_{n}(X),n-1)$-space.
\end{itemize}
\end{definition}

\begin{lemma}
\label{algebraic-topology-lemma-white-tower-exist-cw-compl}
Whitehead towers exist for CW complexes.
\end{lemma}
\begin{proof}
	Consider
	$$\text{hofib}\hookrightarrow X\to
	\text{"$X$ without $\pi_{k}$, $k>n$}"$$
	We immediately get that the second arrow
	induces an isomorphism on $\pi_{i}$ for
	$i\leq n$, while the first arrow induces
	isomorphisms for $i>n$. Define $
	X_{n}:=\text{hofib}$.

	Now notice that $X_n$ and $X_{n+1}$ have
	isomorphic homotopy groups except perhaps
	the $(n+1)$-th, because for $X_{n}$ it is
	already the same as that of $X$, but for
	$X_{n+1}$ it is still zero. So, from the
	fibration
	$$\text{hofib} \hookrightarrow X_{n+1}\to
	X_{n}$$
	we get an isomorphism
	$$
	\xymatrix@C=.9em@R=.9em{
		\cdots\ar[r]
		&\pi_{n+2}(X_{n})\ar[r]\ar[d]_{\cong}
		&\pi_{n+1}(\text{hofib} )\ar[d]^{=}\ar[r]
		&\pi_{n+1}(X_{n+1})\ar[r]\ar[d]^{=}
		&\pi_{n+1}(X_{n})\ar[r]\ar[d]_{\cong}
		&\pi_{n}(\text{hofib} )\ar[r]
		&\pi_{n}(X_{n+1})\ar[r]\ar[d]^{=}
		&\cdots\\
		&\pi_{n+2}(X)&0&0
		&\pi_{n+1}(X)\ar[ur]_{\cong}
		&&0&
	}
	$$
	showing that the homotopy fiber of this
	fibration is in fact $K(\pi_{n+1}(X),n)$.

\end{proof}

\subsection{How to kill higher homotopy
groups without changing lower homotopy groups
and lower homology}
\label{algebraic-topology-subsection-how-kill-highe-homot}

Looks like Blakers-Massey is needed for this.
The following are lecture notes corresponding
to post-Blakers-Massey and pre-Postnikow
towers part our the course.

Glue a disk to a space and what happens to
homotopy groups?
$$
\xymatrix{
	S^{n-1}\ar[r]^{(n-1)\text{-equiv}}\ar[d]
	_{0\text{-equiv}}
	&D^n\ar[d]\\
	X\ar[r]&X\cup D^n\ar@{}[ul]|{\ulcorner}
}
$$
Assume $X$ is connected. We get a map from
the vertical arrows
$$
\xymatrix{
	\pi_i(D^n,S^{n-1})\ar[r]&\pi_i(X\cup D^n,X)
}
$$
which is $(n-1)$-equivalence by
Blakers-Massey. So, by attaching $\sqcup D^n$
we can kill $\pi_{n-1}(X)$, that is,
$X\cup(\sqcup D^n)$ has trivial $\pi_{n-1}$.

Now notice that

% since $S^{n-1}\to D^n$ is an
% $(n-1)$-equivalence, which is the case
% since
$$
\xymatrix{
	0=\pi_i(D^n)\ar[r]
	&\pi_i(D^n,S^{n-1})\ar[r]^{\cong}
	&\pi_{i-1}(S^{n-1})\ar[r]
	&\pi_{i-1}(D^n)=0
}
$$
that is, $\pi_i(D^n,S^{n-1})=0$ for $i\leq
n-1$. This implies that $\pi_i(X\cup
D^n,X)=0$ for $i\leq n-1$.

Also by homotopy long exact sequence,
$$\pi_{n-1}(X)\to \pi_i(X\cup D^n)\text{ is
sujrective}$$
$$\pi_i(X)\to \pi_{i}(X\cup D^n)\text{ is
isomorphism for }i\leq n-2.$$
So what we have thus far is
$$
\xymatrix{
	\pi_n(X\cup D^n)\ar[r]
	&\pi_{n-1}(X)\ar[r]
	&\pi_{n-1}(X\cup D^n)\ar[r]
	&0=\pi_{n-1}(X\cup D^n)
}
$$
Notice that $\pi_n(X\cup D^n,X)$ is not
ingeneral cyclic (counterexample $S^1\cup
D^2$ taking unieversal cover which is real
line with spheres on integers, homotopy
equivalent to $\bigvee_\mathbb{Z} S^2$ which
is not finitely generated).

\begin{quote}
	So basically attaching a disk via $f$ will
	kill $[f]$ inside $\pi_n(X)$ this is called
	{\it killing} an element of the homotopy
	group.
\end{quote}

\begin{proposition}
\label{algebraic-topology-proposition-any-cw-compl-x-x-x-equi}
	For any CW-complex $X$, $X^\ell\to X$ is an
	$\ell$-equivalence.
\end{proposition}

\begin{remark}
\label{algebraic-topology-remark-we-have-used-x-long-exact}
	We have used that for $A\hookrightarrow X$
	from long exact sequence of relative
	homotopy groups we get $\pi_n(X,A)=0$.
\end{remark}

\begin{lemma}
\label{algebraic-topology-lemma-let-i-g-d-i-x-g-d-i-x-then}
	Let $i\leq \ell$ and $g: D^i\to X$,
	$g:\partial D^i\to X^\ell$. Then there is a
	homotopy $\text{rel }\partial D^i$ to a map
	with $\text{img}\subset X^\ell$.
\end{lemma}

\begin{theorem}[Cellular approximation
theorem]
\label{algebraic-topology-theorem-cellu-appro-theor}
	Let $X$ and $Y$ be CW-complexes and
	$\xi:Y\to X$ be any map. Then $\xi$ is
	homotopic to a {\it cellular map}, that is,
	a map $\psi:Y\to X$, such that for all
	$\ell$, $\psi Y^0\subset X^\ell$.
\end{theorem}

We also have
\begin{proposition}
\label{algebraic-topology-proposition-let-n-2-then}
Let $n\geq 2$. Then
$$
\pi_n\left(
\bigvee_{k\in I}S^n\right)
=\bigoplus_{k\in I}\pi_n(S^n)
=\bigoplus_{k\in I}\mathbb{Z}=Z^{\oplus I}
$$
\end{proposition}

\begin{proposition}
\label{algebraic-topology-proposition-finite-i}
First notice that for finite $I$,
$$X^n=X^{n+1}=\vee_{k\in I}S^n$$
by looking carefully. Then
$$\pi_n(X,X^{n+1})=0=\pi_{n+1}(X,X^{n+1})$$
so
$$
\bigoplus_{k\in I}\mathbb{Z}
=\prod_{k\in I}\pi_n(S^n)
=\pi_n(X)=\pi_n(X^{n+1})
=\pi_n(X^n)=\pi_n\left(
\vee_{k\in I}S^n\right)
$$
and for the infinite case it also works,
using finite compactness of the CW complex.
\end{proposition}

\section{Homotopy theory of spaces (C)}
\label{algebraic-topology-section-homotopy-theory-of-spaces-c}

\subsection{Homotopy pushouts and pullbacks}
\label{algebraic-topology-subsection-homot-pusho-pullb}
The idea is simply to construct the homotopy
version of pushouts and pullbacks. So let's
remember that
\begin{definition}
\label{algebraic-topology-definition-item-2}
	\begin{itemize}
		\item A {\it pullback} of the morphisms $f$
		and $g$ consists of an object $P$ and two
		morphisms $p_1:P\to X$ and $p_2:P\to Y$
		satisfying the following universal
		property:
		$$
		\xymatrix{
			Q\ar@/^1pc/[rrd]^{q_2}\ar@/_1pc/[ddr]
			_{q_1}
			\ar@{-->}[dr]^{\phi}\\
			&P\ar[r]^{p_2}\ar@{}[rd]|{\lrcorner}
			\ar[d]_{p_1}&Y\ar[d]^{g}\\
			&X\ar[r]_{f}&Z
		}
		$$
		\item A {\it pushout} of the morphisms $f$
		and $g$ consists of an object $P$ and two
		morphisms $i_1:P\to X$ and $i_2:P\to Y$
		satisfying the following universal
		property:
		$$
		\xymatrix{
			Z\ar[r]^{g}\ar[d]_{f}
			&Y\ar[d]^{i_2}\ar@/^1pc/[ddr]^{j_2}\\
			X\ar[r]_{i_1}\ar@/_1pc/[drr]_{j_1}
			&P\ar@{}[ul]|{\ulcorner}\ar@{-->}[dr]
			^{\phi}\\
			&&Q
		}
		$$
	\end{itemize}
\end{definition}
\begin{remark}
\label{algebraic-topology-remark-other-names-pusho-cofib-prod}
Other names for the pushout are {\it
cofibered product of $X$ and $Y$} (especially
in algebraic categories when $i_1$ and $i_2$
are monomorphisms), or {\it free product of
$X$ and $Y$} with $Z$ {\it amalgamated sum},
or more simply an {\it amalgamation} or {\it
amalgam of $X$ and $Y$.}
\end{remark}
So basically
\begin{definition}[Homotopy pullback and
homotopy pushout]
\label{algebraic-topology-definition-homot-pullb-homot-pusho}
	We just ask that the diagrams above commute
	up to homotopy?
\end{definition}
\begin{remark}[Concrete construction of
homotopy pushout]
\label{algebraic-topology-remark-concr-const-homot-pusho}
	We also need to ask in our definition that
	there is some sort of homotopy-naturality.
	For the case of the homotopy pushout it
	should look something like this:
		$$
		\xymatrix{
			\bullet\ar[rr]\ar[dd]\ar[dr]
			&&\bullet\ar[d]^{\simeq}\\
			&A\ar[r]\ar[d]&B\ar[d]\ar@/^1pc/[ddr]\\
			\bullet\ar[r]^{\simeq}
			&C\ar[r]\ar@/_1pc/[drr]
			&\bullet\ar@{}[ul]|{\ulcorner}\ar[dr]
			^{\simeq}\\
			&&&\bullet
		}
		$$
	so everything is commuting up to homotopy.
	What are the formal definitions? Probably
	rather involved. But a concrete definition
	of the {\it homotopy pushout} according to
	the diagram above is
	$$B\sqcup C\sqcup A\times
	I\Big/\substack{A\times 0\sim B \\ A\times
	I\sim C}$$
\end{remark}
\begin{remark}[From lecture notes]
\label{algebraic-topology-remark-from-lecture-notes}
		The pullback of
	$$
	\xymatrix{
		&Z^I\times_ZY\ar[d]\\
		X\ar[r]&Z
	}
	$$
	is the homotopy pullback of
	$$
	\xymatrix{
		&Y\ar[d]\\
		X\ar[r]&Z
	}
	$$
\end{remark}
\begin{definition}[From lecture notes. Is
this right?]
\label{algebraic-topology-definition-lectu-notes-right}
	The {\it homotopy pullback} of a diagram
		$$
		\xymatrix{
			&Y\ar[d]\\
			X\ar[r]&Z
		}
		$$
		is
		$$
		\xymatrix{
			X\times_{\text{ev}_0}Z^I\times_{\text{ev}
			_1}Y
			\ar[d]\ar[r]&Y\ar[d]\\
			X\ar[r]&Z
		}
		$$
		Intuitively, for any $x\in X$ and $y\in Y$
		this object has the space of paths
		connecting $x$ and $y$.
\end{definition}

\subsection{Homotopy fibers and cofibers}
\label{algebraic-topology-subsection-homot-fiber-cofib}

Before we go, what is a fiber? Perhaps it is
the same as a kernel. And a cofiber may be
the same as a cokernel. Recall:
\begin{definition}[Kernel and cokernel]
\label{algebraic-topology-definition-kernel-and-cokernel}
	\begin{itemize}
		\item The {\it kernel} of a morphism is
		that part of its domain which is sent to
		zero. Formally, in a category with an
		initial object 0 and pullbacks, the {\it
		kernel $\ker f$} of a morphism $f:A\to B$
		is the pullback $\ker(f)\to A$ along $f$ of
		the unique morphism $0\to B$

		More explicitly, this characterizes the
		object $\ker(f)$ as \textit{the} object
		(unique up to isomorphism) that satisfies
		the following universal property:
		\begin{quote}
			for every object $C$ and every morphism
			$h:C\to A$ such that $f\circ h=0$ is the
			zero morphism, there is a unique morphism
			$\phi:C\to\ker(f)$ such that
			$h=p\circ\phi$.
		\end{quote}
		$$
		\xymatrix{
			C\ar@{-->}[dr]^{\phi}\ar@/_1pc/[ddr]_{h}\\
			&\ker(f)\ar[d]\ar[r]&0\ar[d]\\
			&A\ar[r]_{f}&B\ar@{}[ul]|{\ulcorner}
		}
		$$
		\item In a category with a terminal object
		1, the {\it cokernel} of a morphism $f:A\to
		B$ is the pushout (arrows $h$ and $\phi$
		apply if terminal object is zero)
		$$
		\text{coker}(f):=1\sqcup_AB
		\qquad\qquad
		\xymatrix{
			A\ar[r]^{f}\ar[d]\ar@{}[dr]|{\lrcorner}
			&B\ar[d]\ar@/^1pc/[rdd]^{h}\\
			1\ar[r]&\text{coker}(f)\ar@{-->}[dr]
			^{\phi}\\
			&&C
		}
		$$
		In the case when the terminal object is in
		fact zero object, one can, more explicitly,
		characterize the object $\text{coker}(f)$
		with the following universal property:
		\begin{quote}
			for every object $C$ and every morphism
			$h:B\to C$ such that $h\circ f=0$ is the
			zero morphism, there is a unique morphism
			$\phi:\text{coker}(f)\to C$ such that
			$h=\phi\circ i$.
		\end{quote}
\end{itemize}
\end{definition}
The {\it homotopy fiber} and the  {\it
homotopy cofiber} are probably just the
homotopic versions of the kernel and
cokernel.

\begin{definition}[From lecture notes, looks
incomplete…]
\label{algebraic-topology-definition-lectu-notes-looks-incom}
	The {\it homotopy fiber} if $f:Y\to Z$ is
	the pullback of
		$$
		\xymatrix{
			&Y\ar[d]^{f}\\
			pt\ar[r]&Z
		}
		$$
		$F\subset Z^I\times_ZY\to Z$, where $F$ is
		the space of paths starting at $x$ and
		ending at the same point $f(y)$.
\end{definition}

\begin{definition}[Homotopy fiber, a concrete
construction in Hatcher]
\label{algebraic-topology-definition-homot-fiber-concr-const}
	Given a map $f:A\to B$, let $E_f$ be the
	space of pairs $(a,\gamma)$ where $a\in A$
	and $\gamma:I\to B$ is a path in $B$ with
	$\gamma(0)=f(a)$. So: take a point in $A$,
	and map it to some point in $B$, say $f(a)$.
	And then take any path in $B$ starting at $f
	(a)$.
\begin{proposition}
\label{algebraic-topology-proposition-map-p-e-f-b-given-p}
	The map $p:E_f\to B$ given by
	$p(a,\gamma)=\gamma(1)$ is a fibration.
\end{proposition}
We can regard $A$ as a subspace of $E_{f}$
consisting of pairs $(a,\gamma)$ with
$\gamma$ the constant path at $f(a)$, and
$E_f$ deformation retracts onto this subspace
by restricting all the paths $\gamma$ to
shorter and shorter initial segments. The map
$p:E_f\to B$ restricts to $f$ on the subspace
$A$, so we have factored the map $f:A\to B$
as the composition $A\hookrightarrow E_f\to
B$ of a homotopy equivalence and a fibration.

We can also think of this construction as
extending $f$ to a fibration $E_f\to B$ by
enlarging its domain to a homotopy equivalent
space.

The fiber $F_f$ of $E_f\to B$ is called the
{\it homotopy fiber} of $f$. It consists of
the pairs $(a,\gamma)$ with $a\in A$ and
$\gamma$ a path in $B$ from $f(a)$ to a
basepoint $b_0\in B$.
\end{definition}

\subsection{Homotopy fiber of homotopy fiber}
\label{algebraic-topology-subsection-homot-fiber-homot-fiber}

It is interesting to see what happens when
the process of forming homotopy fibers is
iterated. Given a fibration $p:E\to B$ with
fiber $F=p^{-1} (b_0)$, we know that the
inclusion of $F$ into the homotopy fiber
$F_{p}$ is a homotopy equivalence (because
all fibers in homotopy fibrations are
homotopy equivalent?) Well, beacause
\begin{proposition}[4.66]
\label{algebraic-topology-proposition-4-66}
	If $p:E\to B$ is a fibration, then the
	inclusion $E\hookrightarrow E_p$ is a fiber
	homotopy equivalence (a fiber-preserving map
	of the total spaces that has an "fiber
	homotopy inverse" i.e. composition is
	homotopic to the identity via
	fiber-preserving maps). In particular, the
	homotopy fibers are homotopy equivalent to
	the actual fibers.
\end{proposition}

Recall that the homotopy fiber $F_p$ consists
of pairs $(e,\gamma)$ with $e\in E$ and
$\gamma$ a path in $B$ from $p(e)$ to $b_0$.
The inclusion $F\hookrightarrow E$ extends to
a map $i:F_{p}\to E$, $i(e,\gamma)=e$ and
\textit{this map is obviously a fibration}.
In fact it is the pullback via $p$ of the
path fibration $PB\to B$… what could this
mean? Maybe this:
$$
\xymatrix{
	F_{p}\ar[r]^{i}\ar[d]\ar@{}[dr]|{\lrcorner}
	&E\ar[d]^{p}\\
	PB\ar[r]&B
}
$$
Perhaps this justifies why, up to
equivalence, the fiber of this map $i$ must
be $\Omega B$. In the end we get the
following sequence:
$$
\xymatrix@C=1.5em{
	\cdots\ar[r]&\Omega^{2}B\ar[r]&\Omega
	F\ar[r]
	&\Omega E\ar[r]&\Omega
	B\ar[r]&F\ar[r]&E\ar[r]&B
}
$$
where any two consecutive maps form a
fibration, up to homotopy equivalence, and
all the maps to the left $\Omega B$ are
obtained by applying the functor $\Omega$ to
the later maps. The long exact sequence of
homotopy groups for any fibration in the
sequence coinincides with the long exact
sequence for $F\hookrightarrow E\to B$.

\subsection{Homotopy cofiber of homotopy
cofiber}
\label{algebraic-topology-subsection-homot-cofib-homot-cofib}

The dual construction to the sequence above
is
$$
\xymatrix@C=1.5em{
	X\ar[r]&Y\ar[r]&M_{f}\ar[r]&\Sigma X\ar[r]
	&\Sigma Y\ar[r]&\Sigma
	M_{f}\ar[r]&\Sigma^{2}X\ar[r]
	&\cdots
}
$$

\subsection{The two Dold-Puppe sequences,
their (co)-exactness}
\label{algebraic-topology-subsection-two-dold-puppe-seque}

\subsection{Fiber of a (Serre) Hurewicz
fibration is a (weak) homotopy fiber}
\label{algebraic-topology-subsection-fiber-serre-hurew-fibra}

\subsection{Fiber bundles are Serre
fibrations}
\label{algebraic-topology-subsection-fiber-bundl-serre-fibra}
We've been talking a lot about Hurewicz
fibrations. Let's talk about Serre
fibrations. Notice that H. fibration
$\implies$ S. fibration. What is the most
natural example of a Serre fibration?

\begin{proposition}[Hatcher, 4.48]
\label{algebraic-topology-proposition-hatcher-4-48}
	Let $E$ be a fiber bundle with fiber $F$.
	Then $f$ is a Serre fibration.
\end{proposition}
\begin{proof}
	What sdoes it mean to be a Serre fibration?
	It means that
	$$
	\xymatrix{
		I^n\ar[r]\ar[d]&E\ar[d]\\
		I^{n+1}=I^n\times I\ar[r]\ar@{-->}[ur]&B
	}
	$$
	So if $\mathcal{U}$ is a covering of $B$
	such that $f^{-1}U\cong U\times F$. By
	Lebesgue lemma, there is a $\delta>0$ such
	that for all $x\in I^{n+1}$, the ball
	$B(x,\delta)$ lies in some $f^{-1}U$ for
	some $U$.

	Then we subdivide $I^{n+1}$ in smaller cubes
	of the same size with diameter $<\delta$.
	So, each the image of each cube lies in some
	$U\in\mathcal{U}$.

	Then
	$$
	\xymatrix{
		I^n\ar[d]\ar[r]&F\times U\ar@{->>}[d]\\
		I^{n+1}\ar[r]\ar@{-->}[ur]&U
	}
	$$
	has a lift for every little square because
	$$
	\xymatrix{
		X\ar[r]\ar[d]&U\ar[d]\\
		X\times I\ar[r]\ar@{-->}[ur]&pt
	}
	$$
	is always a fibration (think about this) and
	because pullbacks of fibrations are
	fibrations:
	$$
	\xymatrix{
		U\times
		F\ar@{->>}[r]\ar@{->>}[d]&U\ar@{->>}[d]&\\
		F\ar@{->>}[r]&pt
	}
	$$.
	Then we may just add up the squares because
	$$
	\xymatrix{
		D^n\ar[d]\\
		D^n\times I
	}
	$$
	and we're done.
\end{proof}

\subsection{Long exact sequence of homotopy
groups associated to Serre fibration}
\label{algebraic-topology-subsection-long-exact-seque-homot}
\begin{theorem}[Hatcher 4.41]
\label{algebraic-topology-theorem-hatcher-4-41}
	Suppose $p:E\to B$ is a Serre fibration.
	Choose basepoints $b_{0}\in B$ and $x_{0}\in
	F=p^{-1}(b_{0})$. Then the map
	$p_{*}:\pi_{n}(E,F,x_{0})\to
	\pi_{n}(B,b_{0})$ is an isomorphism for all
	$n\geq 1$.

	Hence, if $B$ is path connected, there is a
	long exact sequence
$$
\xymatrix{
	\cdots\ar[r]&\pi_n(F)\ar[r]&\pi_n(E)\ar[r]
	&\pi_n(B)\ar[r]&\pi_{n-1}(F)\ar[r]
	&\pi_{n-1}(E)\ar[r]&\cdots\\
	\cdots\ar[r]&\pi_n(F)\ar[r]\ar[u]^{=}
	&\pi_n(E)\ar[r]\ar[u]^{=}
	&\pi_n(E,F)\ar[r]\ar[u]^{\cong}\ar[u]_{p_{*}
	}
	&\pi_{n-1}(F)\ar[r]\ar[u]^{=}
	&\pi_{n-1}(E)\ar[r]\ar[u]^{=}&\cdots
}
$$
\end{theorem}

\subsection{Fiber bundles over manifolds are
Hurewicz fibrations}
\label{algebraic-topology-subsection-fiber-bundl-over-manif}

\section{Homotopy theory of spaces (D)}
\label{algebraic-topology-section-homotopy-theory-of-spaces-d}

\subsection{Eilenberg-MacLane spaces}
\label{algebraic-topology-subsection-eilenberg-maclane-spaces}
A space $X$ having jut one nontrivial
homotopy group $\pi_{n}(X)\cong G$ is called
an {\it Eilenberg-MacLane space} $K(G,n)$. We
can build a CW complex $K(G,n)$ by the
familiar construction of attaching cells to a
wedge of spheres, acounting for generators
and relations, and then attaching further
cells to kill higher homotopy groups.

\begin{proposition}[4.30]
\label{algebraic-topology-proposition-4-30}
	The homotopy type of a CW complex $K(G,n)$
	is uniquely determined by $G$ and $n$.
\end{proposition}
\begin{proof}[Idea of proof]
	Attaching cells and Whithead theorem.
\end{proof}

\subsection{Representable functors and Yoneda
lemma}
\label{algebraic-topology-subsection-repre-funct-yoned-lemma}

\begin{definition}[Representable functor]
\label{algebraic-topology-definition-representable-functor}
	Let $\mathcal{C}$ be a category. A functor
	$F:\mathcal{C}^{\text{op}}\to \text{Sets}$
	is called {\it representable} if there
	exists an object $B=B_{F}$ in $\mathcal{C}$
	with the property that there is a
	\textit{natural} isomorphism of functors
	$$\varphi:\mathcal{C}(-,B_{F})\to F$$
	where $\mathcal{C}(-,B_{F})$ is the set of
	arrows from $-$ to $B_{F}$.

	One ussually expresses the naturality
	condition for a map $f:X\to Y$ with the
	following diagram:

	$$
	\xymatrix{
		\mathcal{C}(X,B)\ar[r]^{\varphi_{X}}\ar[d]
		_{f^{*}}
		&F(X)\ar[d]^{f^{*}}\\
		\mathcal{C}(Y,B)\ar[r]_{\varphi_{Y}}&F(Y)
	}
	$$
\end{definition}

\begin{remark}
\label{algebraic-topology-remark-apply-b-ident-map-b-b-we}
	By applying $\varphi_{B}$ to the identity
	map  $B\to B$ we obtain a special element
	$\gamma=\varphi_{B}(\text{id})\in F(B)$,
	which is \textit{generic} in the sense that
	any element $\xi \in F(X)$ can be obtained
	as $\xi=f^{*} (\gamma)$, for a suitable
	$f:X\to B$. Indeed, one can take
	$f=\varphi^{-1}_{X}(\xi)$ and apply
	naturality to check that
	$\xi=f^{*}(\gamma)$.
\end{remark}


\begin{lemma}[Yoneda,
\href{https://en.wikipedia.org/wiki/%
Yoneda_lemma}{wiki}]
\label{algebraic-topology-lemma-yoned-https-en-wikip-org-wiki}
	Let $F$ be a functor from a locally small
	category $\mathcal{C}$ to $\text{Set}$. Then
	for each object $X$ of $\mathcal{C}$, the
	natural transformations
	$\text{Nat}(\Hom(X,-),F)$ are in one-to-one
	correspondence with the elements of $F(X)$,
	that is
	$$\text{Nat}(\Hom(X,-), F)\cong F(X)$$
Moreover, this isomorphism is natural in $A$
and $F$ when both sides are regarded as
functors from
$\mathcal{C}\times\text{Set}^{\mathcal{C}}$
to $\text{Set}$. ($\text{Set}^{\mathcal{C}}$
denotes de category of functors from
$\mathcal{C}$ to $\text{Set}$.)

	There is a contravariant version of Yoneda
	lemma asserting that if $F$ is a
	contravariant functor from $\mathcal{C}$ to
	$\text{Set}$,
	$$\text{Nat}(\Hom(-,X), F)\cong F(X).$$
\end{lemma}
\begin{lemma}
\label{algebraic-topology-lemma-nat-x-y-hom-x-y}
	$\text{Nat}
	(\Hom(-,X),\Hom(-,Y))\cong\text{Hom}(X,Y)$.
\end{lemma}
\begin{remark}
\label{algebraic-topology-remark-corre-x-x-fully-faith-corre}
	The correspondence $X\mapsto\Hom(-,X)$ is
	fully faithful, that is, the correspondence
	$\Hom(X,X')\to\text{Nat}(\text{Hom}(-,X),
	\text{Hom}(-,X'))$
	is injective and bijective.
\end{remark}

\subsection{Brown's representability theorem}
\label{algebraic-topology-subsection-brown-s-repre-theor}

\subsubsection{In lecture13.pdf}
\label{algebraic-topology-subsection-in-lecture13-pdf}

See
\href{https://www.math.ru.nl/~gutierrez/%
files/Lecture13.pdf}{Lecture
13: Representable functors and the Brown
Representability Theorem.pdf}

Let us summarize the discussion so fat.
Suppose that
$$F:\text{Ho}(\text{Top}_*)^{\text{op}}\to
\text{Sets}_{*}$$
is a functor having the following two
properties:
\begin{enumerate}
	\item[(i)] $F\left( \bigvee_{i\in I}X_{i}
	\right) \to \prod_{i\in I}F(X_{i})$ is an
	isomorphism for any family of pointed spaces
	$\{X_{i}\} $.
\item[(ii)] "Pushout where one of the arrows
is a cofibration goes to weak pullback".
$F(B\cup_{A}C)\to F(B)\times_{F(A)} F(C)$ is
a surjection for any two cofibrations $A\to
B$ and $A\to C$.
\end{enumerate}

Then we also have
\begin{enumerate}
	\item[(iii)] $F(*)=*$
	\item[(iv)] If we have a pushout of pointed
	spaces
		$$
		\xymatrix{
			A\ar[r]\ar[d]&X\ar[d]\\
			B\ar[r]&Y
		}
		$$
		in which $i$ is a cofibration, then
		$F(Y)\to F(B)\times_{F(A)}F(X)$ is a
		sujrection.
	\item[(v)] If
	$\xymatrix@C=1.5em{A\ar@<.5ex>[r]\ar@<-
	.5ex>[r]&B\ar[r]&C}$
	is a coequalizer of pointed spaces in which
	the two maps form a cofibration $A\vee A\to
	B$, then $F(C)\to F(B)$ maps surjectively to
	the equalizer of
	$\xymatrix@C=1.5em{F(B)\ar@<.5ex>[r]\ar@<-
	.5ex>[r]&F(A)}$.
	\item[(vi)] If $X_{0}\to X_{1}\to X_{2}\to
	\ldots$ is a sequence of cofibrations, then
	$F(X)\to \lim_{\leftarrow}$ is a sujection.
\end{enumerate}

\begin{theorem}[Brown representability
theorem]
\label{algebraic-topology-theorem-brown-repre-theor}
	Let $F$ be a contravariant functor from the
	homotopy category of parallel connected
	CW-complexes to pointed sets. If $F$
	satisfies conditions (i) and (ii) above (for
	any pointed connected CW-complexes $X_{i}$,
	$A$, $B$, $C$), then $F$ is representable.
\end{theorem}

\begin{remark}[So what is a classifying
space?]
\label{algebraic-topology-remark-so-what-is-a-classifying-space}
	The theorem says that there is a space
	$B=B_{F}$ (itself a pointed CW-complex) for
	which there is a natural isomorphism
	$$\varphi :[X,B_{F}]_{*}\overset{\cong
	}{\longrightarrow}F(X)$$
	for any pointed CW-complex $X$. This space
	$B_{F}$ is called a {\it classifying space}
	for $F$. Recall also that when such
	$\varphi$ exists, it is completely
	determined by a \textit{generic} element
	$\gamma\in F(B_{F})$.

	The classifying space together with the
	genereic element is unique up to homotopy.
\end{remark}

Now let's sketch how the proof of Brown
representability theorem goes.

\begin{definition}
\label{algebraic-topology-definition-b-spher-class-space-f}
	$(B,\gamma)$ is a {\it spherical classifying
	space} for $F$ if $\gamma\in F(B)$ is an
	element inducing an isomorphism
	$$\gamma_{*}:[S^{n},B]_{*}\to
	F(S^{n}),\qquad \gamma_{*}(f)=f^{*}
	(\gamma)$$
	for any sphere $S^{n}$, $n>0$.
\end{definition}

\begin{remark}
\label{algebraic-topology-remark-if-b-1-1-b-2-2-two-such-sphe}
	If $(B_{1},\gamma_{1})$ and
	$(B_{2},\gamma_{2})$ are two such spherical
	classifying spaces in the category
	$\mathcal{C}$, and $f:B_{1}\to B_{2}$ is a
	map with $f^{*} (\gamma_{2})=\gamma_{1}$,
	then $f$ induces isomorphisms
	$$\pi_{n}(B_{1})\to \pi_{n}(B_{2})$$
	between the homotopy groups for each $n>0$.

	Since $B_{1}$ and $B_{2}$ are pointed
	connected CW-complexes, this means that $f$
	is a weak homotopy equivalence, hence a
	homotopy equivalence by Whitehead theorem.
\end{remark}

The proof of Brown representability theorem
is based in the following two lemmas:

\begin{lemma}
\label{algebraic-topology-lemma-let-x-be-point-cw-compl-let}
	Let $X$ be a pointed CW-complex and let $\xi
	\in F(X)$. Then there exists a spherical
	classifying space $(B,\gamma)$ for $F$ with
	a cofibration  $f:X\to B$ with $f^{*}
	(\gamma)=\xi$.
\end{lemma}

And in fact,

\begin{lemma}
\label{algebraic-topology-lemma-any-spher-class-space-b-f}
	Any spherical classifying space $(B,\gamma)$
	for $F$ is a classifying space.

	(Thus, $\gamma_{*}:[X,B]_{*}\to F(X)$ is an
	isomorphism for any pointed connected
	CW-complex $X$, not just for spheres.
\end{lemma}

\subsubsection{In Hatcher}
\label{algebraic-topology-subsection-in-hatcher}

\begin{theorem}[4.57]
\label{algebraic-topology-theorem-4-57}
	There are natural bijections
	$T:\left\langle{}X,K(G,n) \right\rangle{} \to H^{n}(X;G)$
	for all CW complexes $X$ and all $n>0$, with
	$G$ any abelian group. Such a $T$ has the
	form $T([f])=f^{*}$ for a certain
	distinguised class $\alpha\in
	H^{n}(K(G,n);G)$.
\end{theorem}

\begin{definition}[Fundamental class]
\label{algebraic-topology-definition-fundamental-class}
	A class $\alpha\in H^{n}(K(G,n);G)$ with the
	property stated in the theorem is called a
	{\it fundamental class}.
\end{definition}


\begin{theorem}[4E.1]
\label{algebraic-topology-theorem-4e-1}
	Every reduced cohomology theory on the
	category of basepointed CW complexes and
	basepoint-preserving maps has the form
	$h^{n} (X)=\left\langle{}X,K_{n} \right\rangle{} $ for some
	$\Omega$-spectrum $\{K_{n}\} $.
\end{theorem}


\subsection{Representability of generalized
cohomology theories}
\label{algebraic-topology-subsection-repre-gener-cohom-theor}

\subsection{$H^{n}(-,G)$ is represented by
$K(G,n)$ together with a chosen element in
$H^{n}(K(G,n),G)$}
\label{algebraic-topology-subsection-h-n-g-repre-k-g-n-toget}

Reduced cohomology is representable since

\begin{enumerate}
	\item[(i)]
	$\widetilde{H}^{n}(\bigvee_{i}X_{i})=
	\prod_{i}
	\widetilde{H}_{n}(X_{i})$.
	\item[(ii)] Also holds.
\end{enumerate}

What is the classifying space? Well, we have
$$[S^{p},B]=H^{n}(S^{p})=\begin{cases}
	0\qquad &p\neq n\\
	\mathbb{Z}\qquad &p=n
\end{cases}$$
which means that $B\cong K(\mathbb{Z},n)$.
More explicityly, we see that the homotopy
groups of the classifying space must coincide
with the reduced cohomology groups, which are
zero in all but the $n$-th dimension.

\subsection{Eilenberg-MacLane spaces can be
construced from Moore spaces by killing
higher homotopy groups}
\label{algebraic-topology-subsection-eilen-macla-space-can}
\subsection{Definition of $\Omega$-spectrum}
\label{algebraic-topology-subsection-definition-of-spectrum}


\subsection{Hurewicz homomorphism}
\label{algebraic-topology-subsection-hurewicz-homomorphism}
\begin{definition}[Hurewicz homomorphism]
\label{algebraic-topology-definition-hurewicz-homomorphism-2}
	$$
\begin{aligned}
		h: \pi_{n}(X) &\longrightarrow H_{n}(X) \\
		[f] &\longmapsto f_{*}(\alpha)
	.\end{aligned}
$$
\end{definition}
where $\alpha$ is a generator of $
H_{n}(S^{n})$.

\subsection{Absolute and relative versions of
Hurewicz theorem}
\label{algebraic-topology-subsection-absol-relat-versi-hurew}
\begin{theorem}[Hurewicz]
\label{algebraic-topology-theorem-hurewicz}
	If $X$ is $n$-connected for some $n\geq 2$,
	then $\tilde{H}_{i}(X)=0$ for $i<n$ and
	$\pi_{n}(X)\cong H_{n}(X)$.
\end{theorem}
\begin{lemma}
\label{algebraic-topology-lemma-if-x-y-simpl-conne-cw-space}
	If $X$ and $Y$ are simply connected
	CW-spaces and $f:X\to Y$ induces
	isomorphisms $H_{n}(X)\cong H_{n}(Y)$ for
	all $n$, then it is an homotopy equivalence.
\end{lemma}
\begin{theorem}[Relative Hurewicz]
\label{algebraic-topology-theorem-relative-hurewicz}
	If $(X,A)$ is $(n-1)$-connected, so
	$\pi_{n}(X,A)=0$ for $i\leq n-1$, for some
	$n\geq 2$ and $A$ is simply connected, then
	$H_{i}(X,A)=0$ for $i\leq n-1$ and
	$\pi_{n}(X,A)\cong H_{n}(X,A)$.

	\textit{Note. Perhaps $X$ should be
	$(n-1)$-connected for the relative version
	too---they are stated together in Hatcher.}
\end{theorem}
\begin{theorem}[Relative Hurewicz for general
$A$]
\label{algebraic-topology-theorem-relat-hurew-gener}
	If $(X,A)$ is an $(n-1)$-connected pair of
	path-connected spaces with $n\geq 2$ and
	$A\neq \emptyset $, then
	$h':\pi_{n}'(X,A,x_{0})\to H_{n}(X,A)$ is an
	isomorphism and $H_{i}(X,A)=0$ for $i<n$.
\end{theorem}



\section{Exact couple. Spectral sequence of
an exact couple.}
\label{algebraic-topology-section-exact-coupl-spect-seque}

These notes are based on S. Burkin lecture
notes and Hatcher.

\begin{enumerate}
	\item Suppose a space $X$ is expressed a
	union of a sequence of subspaces
		$$\ldots\subset X_p\subset X_{p+1}\subset
		\ldots X.$$
		Such a sequence is called a {\it filtration
		of $X$}. For example, $X$ could be a CW
		complex and $X_p$ its $p$-skeleton.

	\item Given such a filtration, for every $p$
	we have a pair $(X_p,X_{p-1})$ that yields a
	long exact sequence. All these long exact
	sequences can be arranged in a staricase
	fashion:
$$
\xymatrix@C=.7em@R=.8em{
	H_{p+q}(X_p)\ar@<.5ex>[r]\ar[r]\ar[d]_{i}
	&H_{p+q}(X_p/X_{p-1})\ar@<.5ex>[r]\ar[r]
	\ar[d]
	&H_{p+q-1}(X_{p-1})\ar[r]\ar[d]_{i}
	&H_{p+q-1}(X_{p-1}/X_{p-2})\ar[d]\ar[r]
	&\cdots\\
	H_{p+q}(X_{p+1})\ar[r]\ar[d]_{i}
	&H_{p+q}(X_{p+1}/X_{p})\ar[r]\ar[d]
	&H_{p+q-1}(X_{p})\ar@<.5ex>[r]\ar[r]_{j}
	\ar[d]_{i}
	&H_{p+q-1}(X_{p}/X_{p-1})\ar@<.5ex>[r]\ar[d]
	\ar[r]_{k}
	&\cdots\\
	\cdots&\cdots&\cdots&\cdots
}
$$
	\item If you put all the absolute groups
	$H_n(X_p)$ in a big $A=\bigoplus_{n,p\in
	\mathbb{N}}H_n(X_p)$ and all the relative
	$H_n(X_p,X_{p-1})$ in a big
	$E=\bigoplus_{n,p\in
	\mathbb{N}}H_n(X_p,X_{p-1})$, you can write
	all those long exact sequences in a compact
	triangle diagram called an {\it exact
	couple} triangle diagram
	$$
	\xymatrix{
		A\ar[rr]^{i}&&A\ar[dl]^{j}\\
		&E\ar[ul]^{k}
	}
$$
here the maps $i,j$ and $k$ are just the maps
from the long exact sequences in the other
big diagram (which makes the exact couple
actually \textit{exact}).

	\item There is a nice way to define new
	groups $A'$ and $E'$ and maps between them
	that yield a {\it derived couple}, which is
	another exact couple.	Namely,
		\begin{itemize}
			\item $A^1_{n,p}:=H_n(X_p)$
			\item $E^1_{n,p}:=H_n(X_p,X_{p-1})$.
			\item $d:=jk$
			\item $E':=\ker d/\text{img}d$, the
			homology of $E$ respect to $d$.
			\item $A':=i(A)\subset A$.
			\item $i':=i|_{A'}$.
			\item $j'(ia):=[ja]\in E'$.
			\item $k'[e]:=ke$.
		\end{itemize}
	anyway, they are well defined and yield the
	exact
	$$
	\xymatrix@C=1.5em@R=1em{
		A'\ar[rr]^{i'}&&A'\ar[dl]^{j'}\\
		&E'\ar[ul]^{k'}
	}
	$$

	\item The sequence $E,E',\ldots$ is called a
	{\it spectral sequence}.
\end{enumerate}

\section{First quadrant. Convergence}
\label{algebraic-topology-section-first-quadrant-convergence}

[Explanation that when the spectral sequence
is defined only in the first quadrant it
eventually stabilizes and $E_{\infty}$-page
is well-defined.]

\section{Homological and cohomological
spectral sequences}
\label{algebraic-topology-section-homol-cohom-spect-seque}

\noindent
\emph{[The introductory discussion for this section is missing from
the notes.]}

\begin{remark}
\label{algebraic-topology-remark-diffe-spect-seque-homol-go}
	The differentials in the spectral sequence
	for homology go up-left and for cohomology
	they go right-down.
\end{remark}

\section{Serre spectral sequence}
\label{algebraic-topology-section-serre-spectral-sequence}

\section{Homological version}
\label{algebraic-topology-section-homological-version}

\begin{theorem}[Serre spectral sequence for
homology]
\label{algebraic-topology-theorem-serre-spect-seque-homol}
	Let $F\to X\to B$ be a fibration with $B$
	path-connected. If $\pi_1(B)$ acts trivially
	on $H_*(F;G)$, then there is a spectral
	sequence $\{E^r_{p,q},d_r\}$ with
	\begin{itemize}
		\item
			$$d^r_{p,q}:E^r_{p,q}\to
			E^r_{p-r,q+r-1}\qquad \text{and} \qquad
			E^{r+1}_{p,q}=\ker d^r_{p,q}/\text{img
			}d^r_{?}$$
		\item The groups
			$$F_pH_n:=\text{img}(H_n(X_p)\to H_n(X))$$
		where the map $H_{n}(X_p)\to H_{n}(X)$ is
		just the induced map by inclusion
		$X_p\hookrightarrow X$, form a filtration
			% \item Stable terms $E^\infty_{p,n-p}$
			% are isomorphic to the successive
			% quotients $F_n^p/F^{p-1}_n$ in a
			% filtration
		$$0\subset F_0H_n\subset \ldots\subset
		F_nH_n=H_n(X;G)$$
		of $H_n(X;G)$ such that
		$$E^\infty_{p,n-p}\cong
		F_pH_n/F_{p-1}H_n.$$
		Another way to write this is
		$$E_{p,q}^\infty=F_pH_{p+q}\Big/F_{p-1}
		H_{p+q}.$$
	\item $$E^2_{p,q}\cong H_p(B;H_q(F;G))$$
\end{itemize}
\end{theorem}

\begin{remark}
\label{algebraic-topology-remark-case-serre-fibra-p-n-map}
	In the case of a Serre fibration, for $p\gg
	n$ the map $H_{n}(X_p)\to H_{n}(X)$ is an
	isomorphism. The group
	$\bigoplus_{p}F^p_n/F^{p-1}_n$ is called the
	{\it associated graded group of the filtered
	group $H_n(X)$}.
\end{remark}

\begin{remark}[Quote from Hatcher]
\label{algebraic-topology-remark-quote-from-hatcher}
	If the coefficient group $G$ is a field then
	$H_{n}(X;G)$ is the direct sum
	$\bigoplus_{p} E^\infty_{p,n-p}$.
\end{remark}

\begin{remark}
\label{algebraic-topology-remark-can-we-recov-h-n-e-gener}
	Can we recover $H^{n}(E)$ in general? No.
\end{remark}

\section{Cohomological version}
\label{algebraic-topology-section-cohomological-version}

\begin{theorem}[Serre spectral sequence for
cohomology]
\label{algebraic-topology-theorem-serre-spect-seque-cohom}
	Let $F\to X\to B$ be a Serre fibration with
	$\pi_0(B)=0$ such that the action of
	$\pi_1(B)$ on $H^*(F;G)$ is trivial.

	Then there is a spectral sequence
	$\{E_r^{p,q},d_r\}$ such that
\begin{itemize}
	\item $$d_r:E_r^{p,q}\to
	E_r^{p+r,q-r+1}\qquad\text{and} \qquad
	E_{r+1}^{p,q}=\ker d_r/\text{img }d_r.$$
	\item The groups
		$$F^pH^n:=\text{img} (F^pH^n\to
		F^nH^n=H^n(X))$$
	form a filtration
	$$H^{n}(X)=F^0H^n\supset F^1H^n\supset
	\ldots\supset F^nH^n\supset \emptyset $$

	of $H^{n}(X,G)$ such that
	$$E_{\infty}^{p,q}\cong
	F^pH^{p+q}/F^{p+1}H^{p+q}.$$
\item $$E^{p,q}_2\cong H^p(B,H^{q}(F,G)).$$
\end{itemize}
\end{theorem}

\section{Product in cohomological Serre
spectral sequence}
\label{algebraic-topology-section-produ-cohom-serre-spect}
Hatcher, Spectral sequences p. 543.

The Serre spectral sequence for cohomology
becomes much more powerful when cup products
are brought into the picture. For this we
need to consider cohomology with coefficients
in a ring $R$ rather than just a group $G$.
What we will show is that the spectral
sequence can be provided with bilinear
products
$$E^{p,q}_{r}\times E^{s,t}_{r}\to
E^{p+s,q+t}_{r}$$
for $1\leq r\leq \infty$ satisfying the
following properties
\begin{itemize}
	\item Each differential $d_{r}$ is a
	derivation, satisfying
		$$d(xy)=(dx)y+(-1)^{p+q} x(dy)\qquad
		\forall x\in E^{p,q}_{r}$$
	This implies that the product on page $r$
	induces a product in page $r+1$, and the
	product in page $\infty$ shall be the one
	induced by previous products.
	\item The product on the second page,
	$E^{p,q}_{2}\times E^{s,t}_{2}\to
	E^{p+s,q+t}_{2}$ is $(-1)^{qs}$ times the
	standard (standard? because as I remember
	cup product did not move the coeffients,
	right?) cup product
		$$H^{p}(B;H^{q}(F;R))\times
		H^{s}(B;H^{t}(F;R))\to
		H^{p+s}(B;H^{q+t}(F;R))$$
		sending a pair of cocycles $(\phi,\psi)$ to
		$\phi\smile \psi$ where coefficients are
		multiplied via the cup product
		$H^{q}(F;R)\times H^{t}(F;E)\to
		H^{q+t}(F;R)$.
	\item The cup product in $H^{*}(X;R)$
	restricts to maps $F^{m}_{p}\times
	F^{n}_{s}\to F^{m+n}_{p+s}$. These induce
	quotient maps $F^{ m}/F^{m}_{p+1}\times
	F^{n}_s/F^{n}_{p+s}\to
	F^{m+n}_{p+s}/F^{m+n}_{p+s+1}$ that coincide
	with the product $E^{p,m-p}_{\infty}\times
	E^{s,n-s}_{\infty}\to
	E^{p+s,m+n-p-s}_{\infty}$.
\end{itemize}

\section{Wang spectral sequence}
\label{algebraic-topology-section-wang-spectral-sequence}

Let's use lecture notes and
\href{https://en.wikipedia.org/wiki/%
Spectral_sequence#Wang_sequence}{wikipedia}
to understand Wang spectral sequence.

Consider a fibration over a sphere
$F\hookrightarrow E\to B$. According to our
theorem for Serre spectral sequences, we know
that the $n$-th page looks as follows:
\begin{align*}
\begin{array}{c|cccccccc}
	n&H_n(F)&&&&&&&H_{n}(F)\\
	 &\vdots&&&&&&&\vdots\\
	2&H_2(F)&&&&&&&H_2(F)\\
	1&H_1(F)&&&&&&&H_1(F)\\
	0&H_0(F)&&&&&&&H_0(F)\\
	\hline\\
	 &0&&&\ldots&&&&n
\end{array}
\end{align*}

We can see there can be nontrivial
differentials only on the $n$th page, so that
$E^{n+1}=E^n$. Moreover, the differentials
that don't vanish is are of the form
$$
\xymatrix{
E^n_{n,q}\ar[r]^{d^n_{n,q}}&E^n_{0,q+n-1}
}
$$

For $d_{n,q-n}^n$ we obtain
$$
\xymatrix{
	0\ar[r]&\ker
	d^n_{n,q-n}\ar@{^{(}->}[r]^{\text{inclusion}
	}
	&E^n_{n,q-n}\ar[r]^{d^n_{n,q-n}}
	&E^n_{0,q-1}\ar@{->>}[r]^{\text{quotient}}
	&\text{coker}d^n_{n,q-n}\ar[r]&0
}
$$
But all that is just
$$
\xymatrix@R=1em{
	0\ar[r]&\ker
	d^n_{n,q-n}\ar@{^{(}->}[r]^{\text{inclusion}
	}
	\ar[d]^{=}
	&E^n_{n,q-n}\ar[r]^{d^n_{n,q-n}}\ar[d]^{=}
	&E^n_{0,q-1}\ar@{->>}[r]^{\text{quotient}}
	\ar[d]^{=}
	&\text{coker}d^n_{q-n}\ar[r]\ar[d]^{=}&0\\
	0\ar[r]&\dfrac{\ker
	d^n_{n,q-n}}{\underbrace{\text{img
	}d^{n-1}_{n,q-n}}_{=0}}
	\ar[r]\ar[d]^{=}
	&E^2_{n,q-n}\ar[r]\ar[d]^{=}
	&E^2_{0,q-1}\ar[r]\ar[d]^{=}
	&\dfrac{E^n_{n,q-n}}{\text{img }d^n_{n,q-n}}
	\ar[d]^{=}\ar[r]&0\\
	0\ar[r]&E^\infty_{n,q-n}\ar[r]&H_{q-n}(F)
	\ar[r]
	&H_{q-1}(F)\ar[r]&E^{\infty}_{0,q-1}\ar[r]&0
}
$$
This is the first "half" of the Wang
sequence. For the other half recall that by
the theorem of Serre spectral sequence for
homology we know that
$$E^{\infty}_{p,n-p}=F_pH_n/F_{p-1}H_n$$
for a filtration on the $n$-th homology group
of the total space
$$\emptyset=F_{-1}\subset F_0 \subset
F_1\subset \ldots\subset F_n=H_{n}(E).$$
And recall that we may also write
$$E^\infty_{p,q}=F_pH_{p+q}/F_{p-1}H_{p+q}.$$
Then we can see that
$$
\xymatrix@R=1em{
	0\ar[r]&E^\infty_{0,q}\ar[d]^{=}\ar[r]^{i^*}
	&H_{q}(E)\ar[r]&E^\infty_{n,q-n}\ar[r]\ar[d]
	^{=}
	&0\\
	&\dfrac{F_0H_{q}(F)}{\underbrace{F_{-1}H_{q}
	(F)}_{=0}}
	&&\dfrac{F_nH_{q}(E)}{F_{n-1}H_{q}(E)}
}
$$
But
$$E^\infty_{0,q}=\dfrac{F_0H_{q}(F)}
{\underbrace{F_{-1}H_{q}(F)}_{=0}}\quad
\overset{\text{why?}}{=}H_{q}(F)$$

This is the other "half" of the Wang
sequence. It only remains to put both halves
together and conclude that
$$
\xymatrix@C=1.5em{
	\cdots\ar[r]&H_{q}(F)\ar[r]^{i^*}&H_{q}(E)
	\ar[r]
	&E^{\infty}_{n,q-n}\ar[r]^{d^n_{n,q-n}}
	&H_{q-n}(F)\ar[r]&H_{q-1}(F)\ar[r]
	&E^\infty_{0,q-1}\ar[r]&\cdots
}
$$
But we may remove the $E^\infty$ terms to get
$$
\xymatrix@C=1.5em{
	\cdots\ar[r]&H_{q}(F)\ar[r]^{i^*}&H_{q}(E)
	\ar[r]
	&H_{q-n}(F)\ar[r]^{d^n_{n,q-n}}&H_{q-1}(F)
	\ar[r]
	&H_{q-1}(E)\ar[r]&H_{q-n-1}(F)\ar[r]&\cdots
}
$$

\section{Extra: spectral sequences for chain
complexes and the Fr\"olicher spectral
sequence}
\label{algebraic-topology-section-extra-spect-seque-chain}
Let's follow Griffiths and Harris to
construct a nice generalization of this
theory.

A {\it chain complex} $(K^\bullet,d)$ is a
sequence of abelian groups (if working in the
category $\text{ab}$)
$$
\xymatrix@C=1.5em{
	K^{0}\ar[r]^{d}&K^{1}\ar[r]^{d}&K^2\ar[r]
	^{d}&\cdots
}
$$
such that the differential $d$ satisfies
$d\circ d=0$. We have {\it cohomology groups}
$$H^{p}(K^*)=\frac{\ker d:K^p\to
K^{p+1}}{\text{img }d:K^{p-1}\to K^p}$$
and the {\it cohomology ring}
$$H^{*}(K^*)=\bigoplus_{p\geq 0}
H^{p}(K^*).$$

A {\it subcomplex} $(J^*,d)$ is given by
subgroups $J^*\subset K^*$
 with $dJ^*\subset J^*$.
The {\it quotient complex} $(L^*,d)$ is also
well-defined by $L^*=K^*/J^*$ and the obvious
differential. this construction yields a {\it
short exact sequence of complexes}
$$
\xymatrix{
	0\ar[r]&J^*\ar[r]&K^*\ar[r]&L^*\ar[r]&0
}
$$
and by the fundamental theorem of homological
algebra, this yields
$$
\xymatrix{
	\cdots\ar[r]&H^{p}(J^*)\ar[r]&H^{p}(K^*)
	\ar[r]
	&H^{p}(K^*)\ar[r]&H^{p+1}(J^*)\ar[r]&\cdots
}
$$

Generalizing the notion of subcomplex, a {\it
filtered complex} $(F^pK^*,d)$ is a
decreasing sequence of subcomplexes
$$K^*=F^0K^*\supset
F^1K^*\supset\cdots\supset F^nK^*\supset
F^{n+1}K^*=\{0\}$$
and the spectral sequence will generalize the
long exact sequence in cohomology.

The {\it associated graded complex} to a
filtered complex $ (F^pK^*)$ is the complex
$$\text{Gr}K^*=\bigoplus_{p>0}
F^pK^*/F^{p+1}K^* $$
with the obvious differential. The filtration
$F^pK^*$ on $K^*$ also induces a filtration
$F^pH^*(K^*)$ on cohomology by (I think)
$$F^pH^q(K^*)=\frac{\ker d:F^pK^q\to
F^pK^{q+1}}{\text{img }d:F^pK^{q-1}\to
F^pK^q}$$
This means that the $F^pH^q(K^*)$ are
subgroups of $H^{q}(K^*)$. This filtration in
turn induces an {\it associated graded in
cohomology}
$$\text{Gr}H^*(K^*)=\bigoplus_{p,q}
F^pH^q(K^*)\big/F^{p+1}H^q(K^*)
$$

\begin{definition}
\label{algebraic-topology-definition-spect-seque-seque-e-r}
	A {\it spectral sequence} is a sequence
	$\{E_r,d_r\}$ with $r\geq 0$ of bigraded
	groups
	$$E_r=\bigoplus_{p,q\geq 0} E^{p,q}_r $$
together with differentials
$$d_r:E^{p,q}_r\to E^{p+r,q-r+1}_r,\qquad
d^2_r=0$$
such that
$$H^{*}(E_r)=E_{r+1}$$
I just wonder what is the cohomology of a
bigraded group, right? there are several
chain complexes and cohomologies, but ok…
\end{definition}

\begin{proposition}[existence and description
of the chain complex spectral sequence]
\label{algebraic-topology-proposition-exist-descr-chain-compl}
	Let $K^*$ be a filtered complex. then there
	exists a spectral sequence $\{E_r\}$ with
	$$
\begin{aligned}
		E^{p,q}_0&=F^pK^{p+q}/F^{p+1}K^{p+q},\\
		E^{p,q}_1&=H^{p+q}\left( F^pK^*/F^{p+1}K^*
		\right)\\
		E^{p,q}_\infty&=F^pH^{p+q}(K^*)\big/F^{p+1}
		H^{p+q}(K^*)
	\end{aligned}
$$
and we just say that the spectral sequence
{\it abuts} to $H^{*}(K^*)$.
\end{proposition}

\paragraph{Explanation.} We have defined four
objects:
\begin{enumerate}
	\item A filtered complex.
	\item A graded complex associated to the
	filtered complex whose entries are direct
	sums of quotients of the filtered complex.
	\item A filtration on the cohomology of the
	complex given by the cohomologies of every
	complex of the filtered complex.
	\item A graded complex associated to the
	filtration on the cohomology whose entries
	are direct sums running over the cohomology
	groups and the filtration entries.
\end{enumerate}
And the proposition says is that there is a
spectral sequence such that
\begin{enumerate}
	\item The $(p,q)$-entry of the zero page is
	the (first) graded at $p+q$.
	\item The $(p,q)$-entry of the first page is
	the $p+q$ cohomology of the graded.
	\item The $(p,q)$-entry of the $\infty$-page
	is the graded of the cohomology
 (ie. the second graded).
\end{enumerate}

Now consider the following situation. A {\it
double complex} is a bigraded group
$$K^{*,*}=\bigoplus_{p,q\geq 0}  K^{p,q}$$
equipped with two differentials
$$d:K^{p,q}\to K^{p+1,q}$$
$$\delta:K^{p,q}\to K^{p,q+1}$$
satisfying $d^2=\delta^2=0$ and $d\delta
+\delta d=0$.

There is an {\it associated single complex}
$(K^*,d)$ defined by
$$K^n=\bigoplus_{p+q=n} K^{p,q} $$
with differential
$$D=d+\delta.$$
there are two filtrations on the associated
single complex:
$$'F^pK^n=\bigoplus_{\substack{p'+q=n\\p'\geq
p}}K^{p',q},$$
	$$''F^qK^n=\bigoplus_{\substack{p+q''=
	n\\q''\geq
	q}}K^{p,q''}.$$
each of these filtrations yields a spectral
sequence according to our proposition, thus
both abutting to the cohomology of the total
space $H^{*}(K^*)$.

Now we turn to Voisin:

\begin{proposition}[Voisin, 8.25]
\label{algebraic-topology-proposition-voisin-8-25}
	Let $K^\cdot$ be the simple complex
	associated to a double complex
	$(K^{p,q},D_1,D_2)$. Then if
	$F^pK^n=\bigoplus_{r\geq o,r+s=n} K^{r,s}$
	is the filtration introduced in section
	8.3.1, the spectral sequence associated to
	$F$ has first terms given by
	\begin{itemize}
		\item $E^{p,q}_0=K^{p,q}$, $d_0=(-1)^pD_2$.
		\item $E^{p,q}_1=H^{q}(K^{p,\cdot})$,
	\end{itemize}
	and the differential $d_1:H^{q}(K^{p,\cdot
	})\to H^{q}(K^{p+1})$ is induced by the
	morphism of complexes $D_1:K^p\to K^{p+1}$.

\end{proposition}

Now to the case that concerns us. Let's go
back to Griffiths and Harris. Let $M$ be a
complex manifold and define double complex
$$K^{p,q}=A^{p,q}(M)$$
with differentials $d$ and $\bar\partial$. I
wonder why the first one is $ d$… it appears
more intuitive to choose $\partial$.

The associated single complex is just the de
Rham complex since
$$A^{n}=\bigoplus_{p+q=n} A^{p,q}.$$
The de Rham differential $d$ satisfies
$d=\partial+\bar\partial$, so I definitely
would like to take $\partial$ instead of $d$
as I said above… I guess I'm not certain
whether $\partial$ induces a chain complex, or
if
$\partial\bar\partial+\bar\partial\partial$.
Then the corresponding filtrations yield the
{\it Fr\"olicher spectral sequences}, both of
which abut to $H^{*}_{\text{dR}}(m)$.

Now if $M$ is compact K\"ahler then every
class $[a]\in 'E^{p,q}_1\cong
H^{p,q}_{\bar\partial}(m)$ has a harmonic
representative for the
$\bar\partial$-laplacian (following
Huybrechts,
$\Delta_{\bar\partial}=
\bar\partial^*\bar\partial+
\bar\partial\bar\partial^*$).
Shouldn't it be the \textit{second}
filtration, $''E^{p,q}$, ie. the one
associated to $\bar\partial$?

Now recall one the K\"ahler identities (prop.
3.1.12 in Huybrechts):
$$2\Delta_{\bar\partial}=\Delta_d:=d^*d+dd^*,
$$
which implies that $da=0$. I would like to
check why this holds, ie. why
$\bar\partial$-harmonic $\implies$
$d$-harmonic $\implies$ closed. This means
that the differentials of the spectral
sequence associated to the first
differential, $d$, are all zero, so
\textit{the spectral sequence degenerates in
the first page}, that is
$$'E_1\cong 'E_2\cong \ldots\cong
'E_\infty.$$
in this case the corresponding filtration is
the {\it Hodge filtration}:
$$
\begin{aligned}
	F^pH^{n}_{\text{dR}}(M)&\cong
	H^{n,0}(m)\oplus\ldots\oplus H^{p,n-p}(M)\\
				       &=\bigoplus_{\substack{p'+q=n \\
				       p'\geq p}}  H^{p',q}(M)
\end{aligned}
$$
\begin{remark}
\label{algebraic-topology-remark-if-m-not-k-ahler-spect-seque}
If $M$ is not K\"ahler the spectral sequence
need not degenerate at the first page. the
{\it Iwasawa manifold} given by a quotient of
a Lie group of complex matrices is an example
where $'E_1\neq 'E_2$, but no example is
known there ${'E_2\neq 'E_\infty}$
\end{remark}

\begin{remark}[Altan]
\label{algebraic-topology-remark-altan}
	It appears that we have only argued that the
	differential on the first page of the
	Fr\"olicher spectral sequence is zero. Other
	spectral seqeunces like the Gysin spectral
	sequence also have vanishing differentials
	on the first pages, but this does not imply
	that the first and the infinity page
	coincide. How can we explain that this is
	actually the case for the Fr\"olicher
	spectral sequence?
\end{remark}

\paragraph{Answer}It looks to me like the
proof reduces to Dolbeault´s theorem. Not
only on the first page can we use that any
form has a $\bar\partial$-harmonic
representative so that the differential
$d_1=\bar\partial$ (or $\partial$?) vanishes,
but also on any other page since all the
$E_{r}$ are just subquotients of some
Dolbeault cohomology group---in the end we
just gave take a representative (of the
representative of the representative…) that
is $\bar\partial$-harmonic. And this all just
made me think:

\begin{remark}[Dani]
\label{algebraic-topology-remark-dani}
	Then what was all that
	$\bar\partial$-harmonic implies $d$-closed
	discussion in Griffiths and Harris for?
\end{remark}

\paragraph{Atlan's answer to Altan's
question} We have shown that
$$E_1=\bigoplus_{p,q} H^{q}(X,\Omega^q)$$
and we also know that
$$E_\infty=\bigoplus_{p,q} E_\infty^{p,q} $$
So we get an isorphism between these pages by
the fact that the filtration on the de Rham
complex is the Hodge filtration. Voisin
justifies this via
\begin{proposition}[Voisin, 7.5]
\label{algebraic-topology-proposition-voisin-7-5}
Let $F^p\mathcal{A}^k(X)$ be the set of
complex differential forms which are sums of
forms of type $(r,k-r)$ with $r\geq p$ at
every point. Then
$$F^pH^{k}(X,\mathbb{C})=\dfrac{\ker
d:F^pA^k(X)\to F^pA^{k+1}(X) }{\text{img
}d:F^pA^{k-1}(X)\to F^pA^k(X)}$$
\end{proposition}

\begin{remark}
\label{algebraic-topology-remark-why-all-meani-voisi-claim}
	Why is all of this meaningful? Voisin claims
	that this result is very nearly
	equivalent to the Hodge theorem, which says
	that
\begin{equation*}\label{algebraic-topology-eq:hodge-decom}
	 H^{k}(X,\mathbb{C})=\bigoplus_{p+q=k}
	 H^{p,q}(X).
\end{equation*}
While this precise equation cannot be derived
from the Fr\"olicher spectral sequence, what
we \textit{did} obtain is
$$\bigoplus_{p,q}F^pH^{p+q}(X,\mathbb{C})/
F^{p+1}H^{p+q}(X,\mathbb{C})=E^{p,q}_{\infty}
=E^{p,q}_{1}=H^q(X,\Omega^p).$$

\paragraph{Question.}
Why are Hodge theorem and the conclusion from
the degeneration of the Fr\"olicher spectral
sequence nearly equivalent?
\end{remark}





% --- End algebraic-topology.tex ---

% --- Begin differential-geometry.tex ---
\chapter{Differential Geometry}
\label{chapter-differential-geometry}




\section{Riemannian metrics}
\label{differential-geometry-section-Riemannian-metrics}

\begin{definition}
\label{differential-geometry-definition-product-metric}
Let $M_1$ and $M_2$ be Riemannian manifolds.
The {\it product metric} on
$M_1\times M_2$ is $\pi_1^*g_1+\pi_2^*g_2$.
\end{definition}

\begin{definition}
\label{differential-geometry-definition-flat-torus}
$S^1\times\ldots\times S^1$ with the product
metric is called the {\it flat
torus}.
\end{definition}

\medskip\noindent
Let $G$ be a Lie group equipped with a
bi-invariant metric.
Then for any $U,V,X \in \mathfrak{g}$,
\begin{equation}
\label{differential-geometry-equation-lie-group-bi-invar-metri}
\left\langle{}[U,X],V\right\rangle{}=-\left\langle{}U,[V,X]\right\rangle{}.
\end{equation}

\noindent
This follows from
$$
\begin{aligned}
\left\langle{}U,V\right\rangle{}&=
\left\langle{}dR_{\text{exp}(tX)}\circ
dL_{\text{exp}(tX)^{-1}}U,
dR_{\text{exp}(tX)}\circ
dL_{\text{exp}(tX)^{-1}}V\right\rangle{}\\
&=\left\langle{}dR_{\text{exp}(tX)}U,
dR_{\text{exp}(tX)}V\right\rangle{}
\end{aligned}
$$

\noindent
since both $U$ and $V$ are left invariant.
The key observation is that the map
$R_{\text{exp}(tX)}$
is exactly the flow of $X$.
Thus, differentiating and evaluating at $t=0$
we obtain
$$
0=\left\langle{}[U,X],V\right\rangle{}+\left\langle{}U,[V,X]\right\rangle{},
$$
where the $V$ and $U$ appear since evaluating
$R_{\text{exp}(tX)}$ at $t=0$ is the
identity.

A similar construction shows that the
derivative of the map
$\text{Ad}:G \to \text{End}(\mathfrak{g})$
is the map $\text{ad}:\mathfrak{g} \to
\text{End}(\mathfrak{g})$,
$X \mapsto [X,\cdot]$.
(Indeed, to compute this differential we use
the curve definition,
$\gamma(t)=\text{exp}(tX)$…)

\medskip\noindent

\begin{exercise}[Warped product]
\label{differential-geometry-exercise-wraped-product}
See \cite{pet}. Para $(M,g_M)$ e $(N,g_N)$ e
$f:M \to \mathbb{R}_+$, definimos o
 {\it warped product} como sendo o produto
 cartesiano $M\times N$ com a métrica
 $g_M+f^2g_N$.

Mostre que se os fatores são completos, o
warped product é completo.
\end{exercise}

\begin{definition}
\label{differential-geometry-definition-gradient}
Let  $f:M\to\mathbb{R}$. Define the {\it
gradient} of $f$ as the (unique) vector
field $\nabla f$ such that
\begin{equation}
\label{differential-geometry-equation-gradient}
\left\langle{}\nabla f,X\right\rangle{}=Xf, \qquad \forall
X\in\mathfrak{X}(M).
\end{equation}
\end{definition}

Notice that for a frame $E_i$ with dual frame
$E^i$ we have
$$
\begin{aligned}
\left\langle{}\nabla f,X\right\rangle{}&=g_{ij}E^i\odot
E^j(\nabla f,X)
=g_{ij}(\nabla f)^iX^j\\
Xf&=df(X)=(df)_iE^i(X)=(df)_iX^i\\
\implies& g^{ij}(df)_iX^j=(\nabla f)^iX^i
\end{aligned}
$$
so substituting $X$ with a basic vector $E_i$
and $(df)_i$ with $\partial_if$ we
obtain
\begin{equation}
\label{differential-geometry-equation-gradient-in-coordinates}
(\nabla f)^i=g^{ij}\partial_if \iff \nabla
f=g^{ij}\partial_ifE_j
\end{equation}

\begin{definition}
\label{differential-geometry-definition-divergence}
For any $X\in \mathfrak{X}(M)$,
$$
\text{div}(X):=\text{tr}(v\mapsto \nabla_vX)
$$
\end{definition}

\begin{lemma}
\label{differential-geometry-lemma-Laplacian-of-product}
\cite{doc}, Chapter III, Exercise 9(b).
\begin{equation}
\label{differential-geometry-equation-Laplacian-of-product}
\Delta(fg)=f \Delta g + g \Delta
f+2\left\langle{}\nabla f,\nabla g\right\rangle{}
\end{equation}
\end{lemma}

\begin{proof}
Just compute.
\end{proof}

\begin{lemma}
\label{differential-geometry-lemma-divergence-and-volume}
\begin{equation}
\label{differential-geometry-equation-divergence-and-volume}
di_X\text{Vol}=\text{div}X\text{Vol},\qquad
\text{for all }X\in \mathfrak{X}(M).
\end{equation}
\end{lemma}

\begin{proof}
Pick a geodesic frame $E_i$ (see Lemma
\ref{differential-geometry-lemma-geodesic-coordinates}) and
its dual coframe $\varepsilon^i$, i.e.
satisfying
$\varepsilon^i(E_j)=\delta_{ij}$. Then
$\operatorname{Vol}=\varepsilon^1
\wedge \ldots \wedge \varepsilon^n$. Then for
any $X=X^iE_i \in
\mathfrak{X}(U)$,

$$
\begin{aligned}
i_X\operatorname{Vol}
&=
\varepsilon^1\wedge\ldots\wedge
\varepsilon^n(X^iE_i,\cdot,\ldots,\cdot)
=X^i\varepsilon^1\wedge\ldots\wedge
\varepsilon^n(E_i,\cdot,\ldots,\cdot)
\end{aligned}
$$
How to compute that? Recall that for
top-forms we have
$$
\varepsilon^1\wedge\ldots\wedge
\varepsilon^n(Z_1,\ldots,Z_n)
=\det(\varepsilon^i(Z_j))
$$
so for example if $n=3$
\begin{align*}\varepsilon^1\wedge
\varepsilon^2\wedge\varepsilon^3(E_1,Z_2,Z_3)
&=\begin{vmatrix}
\varepsilon^1(E_1)&\varepsilon^1(Z_2)
&\varepsilon^1(Z_3)\\
\varepsilon^2(E_1)&\varepsilon^2(Z_2)&
\varepsilon^2(Z_3)\\
\varepsilon^3(E_1)&\varepsilon^3(Z_2)&
\varepsilon^3(Z_3)
\end{vmatrix}
=\begin{vmatrix}
1&\varepsilon^1(Z_2) &\varepsilon^1(Z_3)\\
0&\varepsilon^2(Z_2)&\varepsilon^2(Z_3)\\
0&\varepsilon^3(Z_2)&\varepsilon^3(Z_3)
\end{vmatrix}\\
&=\varepsilon^2(Z_2)\varepsilon^3(Z_3)-
\varepsilon^2(Z_3)\varepsilon^3(Z_2)
=\varepsilon^2\wedge\varepsilon^3(Z_2,Z_3),\\
\varepsilon^1\wedge\varepsilon^2\wedge
\varepsilon^3(E_2,Z_2,Z_3)
&=\begin{vmatrix}
0&\varepsilon^1(Z_2) &\varepsilon^1(Z_3)\\
1&\varepsilon^2(Z_2)&\varepsilon^2(Z_3)\\
0&\varepsilon^3(Z_2)&\varepsilon^3(Z_3)
\end{vmatrix}=-
\varepsilon^1\wedge\varepsilon^3(Z_2,Z_3)
\end{align*}
and so on. When we sum over all $i$, we get
$$
X^i\varepsilon^1\wedge\ldots\wedge
\varepsilon^n(E_i,\cdot,\ldots,\cdot)
=\sum_{i}(-1)^{i+1}
X^i\varepsilon^1\wedge\ldots\wedge\widehat{
\varepsilon^i}
\wedge\ldots\wedge\varepsilon^n.
$$
Now take exterior derivative of that, we get
$$
\begin{aligned}
d i_X\operatorname{Vol}
&=\sum_i(-1)^{i+1}(dX^i)
\wedge\varepsilon^1\wedge\ldots
\wedge\widehat{\varepsilon^i}
\wedge\ldots\wedge\varepsilon^n\\
&+\sum_{i}(-1)^{i+1}X^i\wedge
d(\varepsilon^1\wedge\ldots\wedge\widehat{
\varepsilon^i}\wedge\ldots
\wedge\varepsilon^n)\end{aligned}
$$
And then the first term actually is
$$
\begin{aligned}
\sum_i(-1)^{i+1}
E_jX^i\varepsilon^j\wedge\varepsilon^1
\wedge\ldots
\wedge\widehat{\varepsilon^i}
\wedge\ldots\wedge\varepsilon^n
&=E_iX^i\operatorname{Vol}
\end{aligned}
$$
while the second term vanishes because look,
$$
\begin{aligned}d \varepsilon^i(E_j,E_k)
&=E_j\varepsilon^i(E_k)-E_k\varepsilon^i(E_j)
-\varepsilon^i([E_j,E_k])\\
&=-\varepsilon^i(\nabla_{E_j}E_k-\nabla_{E_k}
E_j)\qquad
\text{torsion!}
\end{aligned}
$$
which vanishes at $p$ because we said that
this geodesic frame would have
 vanishing Christoffel symbols at $p$. So we
 conclude:
\[di_X
\operatorname{Vol}=E_iX^i\operatorname{Vol}\]

Now you just have to think what is
divergence:
$$
\begin{aligned}
\operatorname{div}X&=\sum_{i}
\left\langle{}\nabla_{E_i}X^jE_j,E_i\right\rangle{}
=\sum_{i}\left\langle{}E_iX^jE_j,E_i\right\rangle{}+
X^j\left\langle{}\nabla_{E_i}E_j,E_i\right\rangle{}\\
&=\sum_i
\left\langle{}E_iX^jE_j,E_i\right\rangle{}=
\sum_iE_iX^j\left\langle{}E_j,E_i\right\rangle{}=E_iX^i
\end{aligned}
$$
again using that we are a geodesic frame with
vanishing covariant
derivative at $p$.
\end{proof}

















\section{Connections}
\label{differential-geometry-section-connections}

\begin{theorem}[Miracle]
\label{differential-geometry-theorem-miracle}
There exists a unique metric torsion-free
connection on any Riemannian manifold.
\end{theorem}

\begin{proof}
$$
\begin{aligned}
2\left\langle{}\nabla_XY,Z\right\rangle{}
&=X\left\langle{}Y,Z\right\rangle{}+Y\left\langle{}X,Z\right\rangle{}-
Z\left\langle{}X,Y\right\rangle{}\\
&-\left\langle{}X,[Y,Z]\right\rangle{}+\left\langle{}Y,[Z,X]\right\rangle{}+
\left\langle{}Z,[X,Y]\right\rangle{}.
\end{aligned}
$$
\end{proof}

\begin{example}
\label{differential-geometry-example-euclidean-connection}
The usual connection in $\mathbb{R}^n$ is
given by
$$
\nabla_XY=X(Y^i)\partial_i
$$
\end{example}

\begin{proposition}[Everything I know about
connections]
\label{differential-geometry-proposition-exist-pullb-conne}
If $f:M \to N$ is a map of manifolds
and $E$ is a vector bundle on $N$
with a connection $\nabla$,
there exists a unique connection $\nabla^f$
on $f^*E$ such that
$$
\nabla^f_XY\circ f=\nabla_{f_*X}Y
$$
for $X \in TM$ and $Y \in \Gamma(E)$.
\end{proposition}

\begin{lemma}
\label{differential-geometry-lemma-compo-pullb-conne}
\begin{reference}
\cite{aulas}, p.5
\end{reference}
If $g:M_1\to M_2$ and $f:M_2\to M_3$ are
smooth maps of
manifolds and $E$ is a vector bundle over
$M_3$, then
$$
g^* (f^* (E))=(f\circ g)^* (E)\qquad
\text{and}\qquad
(\nabla^f)^g=\nabla^{f\circ g}
$$
\end{lemma}

\begin{proposition}
\label{differential-geometry-proposition-existence-of-levi-civita}
On any Riemann manifold
there exists a unique metric and torsion-free
connection called the {\it Levi-Civita
connection}.
\end{proposition}

\begin{proposition}[Symmetry and
compatibility]
\label{differential-geometry-proposition-symme-compa}
Let $f:M \to N$ be a map of smooth manifolds
and suppose $M$ is Riemannian.
Then the pullback of the LC connection is
symmetric, i.e.
\begin{equation}
\label{differential-geometry-equation-pullb-conne-symme}
\nabla^f_X f_*Y-\nabla^f_Yf_*X =
f_*[X,Y],\qquad
\forall X,Y \in \mathfrak{X}(N)
\end{equation}
and compatible with the pullback Riemannian
metric.
\end{proposition}

The following three lemmas are fundamental
properties of isometric
immersions and will be used in Section
\ref{differential-geometry-section-isometric-immersions}. (Cf.
Lista 3, Exercise 2.)

\begin{lemma}
\label{differential-geometry-lemma-tange-compo-pullb-conne-ii}
If $f:M \to \tilde{M}$ is an isometric
immersion, then
$f_*(\nabla_XY)=(\nabla^f_Xf_*Y)^\top$.
\end{lemma}

\begin{proof}
We have to check that $(\nabla^f_Xf_*Y)^\top$
is a torsion-free and compatible
connection on $TM$.
\end{proof}

\begin{lemma}
\label{differential-geometry-lemma-isometric-immersion-along-curve}
Let $f:M \to \tilde{M}$ be an isometric
immersion, $c$ is a smooth curve in $M$
and $Y\in\mathfrak{X}_c$. Then
$$
f_*\nabla_{\frac{d}{dt}}^\gamma Y
=\left(\nabla^{f\circ\gamma}_{\frac{d}{dt}}
f_*Y\right)^\top
$$
\end{lemma}

\begin{proof}
This is just a mixture of Lemmas
\ref{differential-geometry-lemma-compo-pullb-conne}
 and \ref{differential-geometry-lemma-tange-compo-pullb-conne-ii}.
\end{proof}

The following result will be reminded when we
define totally geodesic
submanifolds, cf. Definition
\ref{differential-geometry-definition-total-geode-subma},
addressing the case when also the normal
component vanishes:

\begin{lemma}
\label{differential-geometry-lemma-geodesics-of-submanifold}
Let $f:M \to \tilde{M}$ be an isometric
immersion and $\gamma:I\to M$ a smooth
curve parametrized by arclength. Show that
$\gamma$ is a geodesic of M
if and only if its acceleration on
$\tilde{M}$ is perpendicular to $M$, that is,
$$
\tilde{\nabla}^{f\circ\gamma}_{\frac{d}
{dt}}(f\circ\gamma)'\perp
f_*(T_{\gamma(t)}M)
$$
\end{lemma}

\begin{proof}
Decompose the vector field
$\tilde{\nabla}^{f\circ\gamma}_{\frac{d}
{dt}}(f\circ\gamma)'$
in tangent and
normal part. Then apply Lemma
\ref{differential-geometry-lemma-tange-compo-pullb-conne-ii}.
\end{proof}

\begin{exercise}[Parallel transport
determines connection]
\label{differential-geometry-exercise-paral-trans-deter-conne}
\begin{reference}
\cite[Chapter II, Exercise 2]{doc}
\end{reference}
Let $X,Y$ be vector fields in a Riemannian
manifold $M$.
Let $p \in M$ and $c:I \to M$ an integral
curve of $X$ passing through $p$.
Show that the Riemannian connection of $M$ is
given by
$$
(\nabla_XY)(p)
=\frac{d}{dt}\Big|_{t=t_0}
(P^{-1}_{c;t_0;t}(Y(c(t)))
$$
where $P_{c;t_0;t}:T_{c(t_0)}M \to T_{c(t)}M$
is the parallel transport along $c$ from
$t_0$ to $t$.
\end{exercise}

\begin{proof}
Let $e_i$ be a basis of $T_{c(t_0)}M$.
Since parallel transport is a bijection
(its inverse is parallel transport
along the same curve but in reversed
direction;
+ uniqueness of ODE?),
we have that $P_t(e_i)$ is a basis of
$T_{c(t)}M$.
Thus we can write $P_tY$ in terms of this
basis, i.e.
$P_tY = Y^i_tP_te_i$.

Then
$$
\begin{aligned}
\nabla_XY&=\nabla_{c'}Y=\nabla_{c_*\frac{d}
{dt}}Y
=\nabla_{\frac{d}{dt}}^{c}Y \circ c\\
&=\nabla^c_{\frac{d}{dt}}Y^i_tP_te_i
=\frac{d}{dt}Y^i_tP_te_i+
Y^i_t\underbrace{\nabla_{\frac{d}{dt}}P_te_i}
_{=0}
=\frac{d}{dt}P_tY.
\end{aligned}
$$
\end{proof}

\begin{exercise}
\label{differential-geometry-exercise-unique-connection-torsion}
Show that if $\nabla$ is any connection
on $M$, then there exists a unique
torsion-free connection on $M$
which has the same geodesics as $\nabla$.
\end{exercise}

\begin{proof}
Suppose that $\nabla'$  and $\nabla''$
are two such connections. Let
$A(X,Y)=\nabla'_XY-\nabla''_XY$.
Then $A$ is a symmetric tensor.
By tensoriality, the condition of
having the same geodesics gives $A(X,X)=0$.
Since $A$ is symmetric, we may polarize
and obtain that $A \equiv 0$.

For existence let
$\nabla'=\nabla-\frac{1}{2}T$
where $T$ is the torsion of $\nabla$.
\end{proof}

\begin{exercise}
\label{differential-geometry-exercise-two-torsion-free-connections}
$\nabla$ and $\nabla'$ are two torsion
free connections on $M$ with the same
geodesics up to reparametrization
if and only if there exists a 1-form
$\omega$ such that
$\nabla_XY-\nabla'_XY
=\omega(X)Y-\omega(Y)X$.
\end{exercise}

\begin{proof}

\end{proof}

\begin{exercise}
\label{differential-geometry-exercise-constant-curve-connection}
Let $M$ be a Riemannian manifold and $p$ a
point in $M$. Consider the constant curve
$f:I \to M$ given by $f(t)=p$. Let $V$ be a
differentiable vector field along $f$, that
is, a differentiable map from $I$ to $T_pM$.
Show that $\frac{DV}{dt}=\frac{dV}{dt}$, that
is, the covariant derivative coincides with
the usual derivative.
\end{exercise}

\begin{proof}
$$
\begin{aligned}
\frac{DV}{dt}&=\nabla^f_{\frac{d}{dt}}V
=\nabla^f_{\frac{d}{dt}}V^i(\partial_i\circ
f)
=\underbrace{\frac{d}{dt}V^i(t)
\partial_i\circ
f}_{\text{usual derivative}}
+V^i(t)\nabla_{\frac{d}{dt}}^f\partial_i
\circ f.
\end{aligned}
$$

\noindent
Note that the last term vanishes since
$$
\nabla^f_{\frac{d}{dt}}\partial_i \circ f
= \nabla_{f_*\frac{d}{dt}}\partial_i
=\nabla_0\frac{d}{dt}\partial_i=0
$$
since $f$ is constant.
\end{proof}

\begin{exercise}
\label{differential-geometry-exercise-geodesics-tangent-bundle}
It's possible to introduce a Riemannian
metric on the tangent bundle
of a Riemannian manifold. Let $(p,v) \in TM$
and $V,W$ tangent vectors
of $TM$ at $(p,v)$. Let
$$
\alpha:t \to (p(t),v(t)),\qquad  \beta:s \to
(q(s),w(s))
$$
with $p(0)=q(0)=p$, $v(0)=w(0)=v$
and $V=\alpha'(0)$, $W=\beta'(0)$.
Define an inner product in $TM$ by
\begin{equation}
\label{differential-geometry-equation-tangent-bundle-metric}
\left\langle{}V,W\right\rangle{}_{(p,v)}
=\left\langle{}\pi_*(V),\pi_*(W)\right\rangle{}_p
+\left\langle{}\nabla^p_{\frac{d}{dt}}v(0),
\nabla_{\frac{d}{ds}^q}w(0)\right\rangle{},
\end{equation}
where $\pi:TM \to M, \pi(p,v)=p$.

\begin{enumerate}
\item Show that this inner product is well
defined
and induces a Riemannian metric on $TM$.

\item A vector $(p,v) \in TM$ is orthogonal
(in the metric above) to the fiber
$\pi^{-1}(p) \cong TM$
is called a {\it horizontal vector}. A curve
$$
t \to (p(t),v(t))
$$
in $TM$ is called {\it horizontal} if its
tangent vector is
horizontal for all $t$.
Show that the curve $t \to (p(t),v(t))$
is horizontal if and only if the vector field
$v(t)$
is parallel along $p(t)$ in $M$.

\item Show that the geodesic field is a
horizontal vector field.

\item Show that the trajectories of the
geodesic field
are geodesics of $TM$ in the above metric.
\end{enumerate}
\end{exercise}

\begin{proof}
\begin{enumerate}
\item We just have to check that the right
term in
Equation \ref{differential-geometry-equation-tangent-bundle-metric}
does not depend on the choice of $\alpha$ and
$\beta$.
In principle the covariant derivatives could
be different
for different choices of curves, but this is
not the
case by the hypothesis that $\alpha'(0)=V$
and $\beta'(0)=W$.
Indeed, the covariant derivative
$\nabla_{\frac{d}{dt}}^pv(0)$
determines the last $n$ coordinates of
$\alpha'(0)$.

\item We compute $\left\langle{}\alpha',V\right\rangle{}$ for
a curve $\alpha$ on $TM$
and a vertical vector field $V$, i.e. $V$ is
contained in $TM$.
In Equation
\ref{differential-geometry-equation-tangent-bundle-metric},
the first term will always vanish by
verticality,
while the second depends on
$\nabla_{\frac{d}{dt}}^pv$.
This vanishes if and only if $v$ is parallel
along $p$.

\item This is immediate from the previous
item
since the velocity vector fields of geodesics
are parallel by definition.

\item Let $\alpha=(p,v)$ be a trajectory of
the geodesic field,
that is $v=p'$ and $v'=0$.
To show that $\alpha$ is a geodesic of $TM$
we must show that
$\nabla_{\alpha'}^{TM}\alpha'=0$.
Note that
$$
\begin{aligned}
2\left\langle{}\nabla_{\alpha'}^{TM}\alpha',
\alpha'\right\rangle{}
&=\frac{d}{dt}\left\langle{}\alpha',\alpha'\right\rangle{}\\
&=\frac{d}{dt}\left(\left\langle{}\pi_*\alpha',
\pi_*\alpha'\right\rangle{}
+\left\langle{}\nabla_{\frac{d}{dt}}^pv,
\nabla_{\frac{d}{dt}}^pv\right\rangle{}\right)\\
&=\frac{d}{dt}\left\langle{}p',p'\right\rangle{}=0
\end{aligned}
$$
since $p$ is a geodesic of $M$.
But $\alpha'$ cannot be zero because $p$ is a
geodesic.
\end{enumerate}
\end{proof}

\begin{exercise}
\label{differential-geometry-exercise-lie-group-bi-invar-metri}
Let $G$ be a Lie group and $\mathfrak{g}$ its
Lie algebra.
Let $X \in \mathfrak{g}$. The trajectories of
$X$
determine a map
$\varphi:(-\varepsilon,\varepsilon) \to G$
with $\varphi(0)=e$ and
$\varphi'(t)=X(\varphi(t))$.

\begin{enumerate}
\item Show that $\varphi(t)$ is defined for
all $t \in \mathbb{R}$
and that $\varphi(t+s)=\varphi(t)\varphi(s)$.
$\varphi$ is called a {\it 1-parameter
subgroup} of $G$.
\end{enumerate}
\end{exercise}

\begin{proof}
\begin{enumerate}
\item Recall that $X \in \mathfrak{g}$
determines a vector field on $G$
by left translation, that is,
$X(g):=L_{g,*}X(e)$ where $X:=X(e)$.
To show that $\varphi$ is defined for all $t$
we may show that there is $\varepsilon>0$
such that for all $t$,
$\varphi(t+\varepsilon)$ is defined.
Suppose that $\varphi$ is a local solution of
$X$ (at $e$).
Pick any $t$ for which $\varphi$ is defined
and consider the translation
$L_{\varphi^{-1}(t)}$, which we use to
pushforward $X(\varphi(t))$
back at $e$.
But $X$ was constructed using left
translations,
so that the vector we obtain is again $X(e)$.
Integrating again we have a local solution up
to the same $t$.
Pushforward now using $L_{\varphi(t)}$.

To show that
$\varphi(t+s)=\varphi(t)\varphi(s)$
fix $t$. Then $\varphi(t+s)$ is an integral
curve of $X$
passing through $\varphi(t)$ at $s=0$.
But so is $\varphi(t)\varphi(s)$ since
$\varphi(0)=e$.
By uniqueness of solutions of ODEs we are
done.
\end{enumerate}
\end{proof}


\section{Riemannian submersions and
immersions}
\label{differential-geometry-section-riema-subme-immer}

For any smooth manifold submersion
$\pi:\tilde{M}\to M$ there is a bundle
isomorphism I love:
\begin{equation}
\label{differential-geometry-equation-subme-bundl-decom}
\pi^*TM \oplus \Ker \pi\cong
T\tilde{M} \end{equation}
 where $\Ker\pi$ is the {\it kernel bundle},
 whose
fiber is the kernel of $\pi_*$, and is also
called the {\it vertical bundle}.
Proving that $\Ker\pi$ is in fact a bundle
and confirming the isomorphism
\ref{differential-geometry-equation-subme-bundl-decom} is a simple
exercise.

We call $\pi$ a {\it riemannian submersion}
only when $\pi^*TM$ is isometric to
$TM$, which makes sense only when $\tilde{M}$
and $M$ are Riemannian manifolds.

Following notation of
\cite{Milnor-Characteristic-Classes}, for any
subbundle
$\xi \subset\eta$ over a Riemannian manifold,
the {\it orthogonal complement
bundle} $\xi^\perp$ can be defined fiberwise
as the orthogonal complement of the
fiber and shown to be a bundle as in the case
of the kernel bundle. Then Milnor
defines the {\it normal bundle} when $\xi$ is
the tangent bundle of a
submanifold. Finally this construction is
noted to work just as well for any
isometric immersion $f:M\to \tilde{M}$, so
that
\begin{equation}
\label{differential-geometry-equation-immer-bundl-decom}
f^*T\tilde{M}\cong TM\oplus T^\perp M
\end{equation}
Further, in problem 3-B we are asked to
define for $\xi \subset \eta$ the {\it
quotient bundle} and show it's a bundle. And
if $\eta$ has a Riemannian metric,
then $\xi^\perp\cong\eta/\xi$. So that's why
we sometimes say that we can define
the normal bundle using a Riemannian metric
but we don't need to.

An interesting thought is: given a subbundle,
is there a manifold whose tangent
bundle is that subbundle?

\section{Geodesics}
\label{differential-geometry-section-geodesics}

\begin{example}[Geodesics of sphere]
\label{differential-geometry-example-geodesics-of-Sn}
The geodesics of $S^n$ are
\begin{equation}
\label{differential-geometry-equation-geodesics-of-Sn}
\gamma(t)=\cos t p+\sin t v
\end{equation}
for $p \in S^n$ and $v\in T^1_pS^n$.

Derivando como uma simples curva em
$\mathbb{R}^{n+1}$, vemos que
$\gamma''=-\gamma$, o que significa que
$(\gamma'')^{\top}=0$, i.e. $\gamma$ é
uma geodésica de $\mathbb{S}^n$.

Sendo essa uma geodésica partindo de um ponto
arbitrário numa direção
arbitrária, concluimos por unicidade das
geodésicas e \textit{rescaling lemma}
que todas as geodésicas de $\mathbb{S}^n$ são
como $\gamma$.
\end{example}

\begin{example}[Isometry group of sphere]
\label{differential-geometry-example-isometry-group-of-Sn}
The isometry group of $S^n$ is
$$
\text{O}(n):=\{A\in\text{GL}(n+1):
AA^{\mathbf{T}}=\text{Id}\}
$$
It's immediate that $\mathcal{O}(n+1)\subset
\text{Isom}(S^n)$ since orthogonal
transformations preserve euclidean product:
$$
\begin{aligned}
A A^{\mathbf{T}}&=\text{Id}\iff \sum_k
A_{ik}A_{jk}=\delta_{ij}
\iff Ae_i\cdot Ae_j=\delta_{ij}\\
\iff Av\cdot Aw&= A\left( v^ie_i\right) \cdot
A\left( w^je_j\right)
=v^iw^j (Ae_i)\cdot (Ae_j)
=v^i\cdot w^j\delta_{ij}=v\cdot w
\end{aligned}
$$
For the other contention let $A:S^n \to S^n$
be an isometry.

Vamos mostrar que $A$ é a restrição de uma
função $\tilde{A} \in
\mathsf{O}(n+1)$. Defina
$$
\begin{aligned}
\tilde{A}: \mathbb{R}^{n+1} &\longrightarrow
\mathbb{R}^{n+1} \\
(r,\theta) &\longmapsto rA(1,\theta)\\
0&\longmapsto 0
\end{aligned}
$$
Se mostramos que $\tilde{A}$ é uma isometria
linear, é claro que ela é um
elemento de $\mathsf{O}(n+1)$ pela conta
anterior. De fato, basta mostrar que
$\tilde{A}$ é uma isometria, pois toda
isometria de espaços de Banach que fixa
a origem é linear (\cite{braga} Teo. 7.11).

Para ver que $\tilde{A}$ é uma isometria de
$\mathbb{R}^{n+1}$,
afirmo que a distância de $p$ a $q$ está
totalmente determinada
pelas normas  $\|p\|$ e $\|q\|$, e pela
distancia esférica entre
$\frac{p}{\|p\|}$ e $\frac{q}{\|q\|}$. Note
que essa afirmação é na verdade
um problema de geometria plana, pois todas
essas quantidades podem ser descritas
dentro do único plano que contém $0$, $p$ e
$q$.

Acabou que essa afirmação é simplesmente a
lei dos cosenos, já que a distância
esférica entre $\frac{p}{\|p\|}$ e
$\frac{q}{\|q\|}$ é exatamente o angulo
entre $p$ e $q$ (poque essa distância é um
segmento de círculo máximo!):
\[\text{lei dos cosenos:} \qquad
d(p,q)^2=\|p\|^2+\|q\|^2-2\|p\|\|q\| \cos
\angle(p,q) \]

Em fim, $\tilde{A}$ é uma isometria porque
$d_{\mathbb{R}^{n+1}}(p,q)=d_{\mathbb{R}^{n+
1}}(\tilde{A}p,\tilde{A}q)$
pelo
argumento anterior.
\end{example}

\begin{proposition}[Totally Normal
Neighbourhoods]
\label{differential-geometry-proposition-total-norma-neigh}
For all $p \in M$ there is a neighbourhood $W
\ni p$ and a number $\delta>0$
 such that for every $q \in W$, $B_\delta(q)$
 is normal and
 $W \subset B_\delta(q)$.
\end{proposition}

\medskip\noindent
Now we discuss distance, geodesics and
exponential.
First notice that if $p$ and $q$ are two
points sufficiently near each other,
$d(p,q)=|\text{exp}^{-1}_p(q)|$.
Indeed, if $q=\gamma_w(1)$ for some geodesic
with
velocity $w \in T_pM$ at $p$ (this is what
we mean by ``sufficiently near''),
we have

$$
\begin{aligned}
d(p,q)=\int_0^1|\gamma'|=\int_0^1|w|=|w|.
\end{aligned}
$$

\begin{exercise}
\label{differential-geometry-exercise-diffe-dista-funct}
Let $d_p:=d(p,\cdot)$. Compute $|\nabla d_p|$
and the integral
curves of $\nabla d_p$ inside
$B_\varepsilon(p)\setminus\{p\}$.
\end{exercise}

\begin{proof}
Suppose that $\gamma$ is a geodesic joining
$p$ and $q$
with $|\gamma'|=1$.
Let $\sigma$ be a curve contained in the
geodesic sphere $S_{d(p,q)}(p)$.
Then
$$
\left\langle{}\nabla d_p,\sigma'\right\rangle{}
=\sigma'(d_p\circ \sigma)
=\sigma'(\underbrace{d(p,q)}_{\text{cte}})=0.
$$
Thus $\nabla d_p=\lambda\gamma'$.
Now
$$
\left\langle{}\nabla d_p,\gamma'\right\rangle{}
=\lambda |\gamma'|=\lambda
$$
for some $\lambda \in \mathbb{R}$.
But
$$
\gamma'd_p=\frac{d}{dt}d_p \circ \gamma'
=\frac{d}{dt}t=1,
$$
so $\lambda=1$, and then $|\nabla d_p|=1$.
It is clear that the integral curves of
$\nabla d_p$
are unit speed geodesics.
\end{proof}

\begin{lemma}
\label{differential-geometry-lemma-geodesic-coordinates}
At every point $p$ in $M$ there is an
orthonormal frame
$\{E_i\} \subset \mathfrak{X}(U)$ such that
\begin{equation}
\label{differential-geometry-equation-norma-coord-vanis-chris}
\Big(\nabla_{E_i}E_j\Big)_p=0
\end{equation}
(This means the Christoffel symbols of the
connection vanish at the point $p$.)
\end{lemma}

\begin{proof}
This frame can be obtained by taking geodesic
coordinates at the point, an
orthonormal base $\{e_i\}$ of $T_pM$, and
taking parallel transport of the
vectors $e_i$ along radial geodesics
emanating from $p$. This immediately
ensures that $E_i$ is orthonormal since
parallel transport preserves angles.

To check that Christoffel symbols vanish at
$p$ we do as follows.
Take a random vector $v \in T_pM$ and its
geodesic
$\gamma_v(t)=\operatorname{exp}_p(tv)$. Then
$$
\begin{aligned}
0&=\nabla_{\frac{d}{dt}}^{\gamma}\gamma'
=\nabla^{\gamma}_{\frac{d}{dt}}v^i(E_i\circ
\gamma)
\end{aligned}
$$
where the $v=(v^1,\ldots,v^n)$. Indeed: this
is very silly but, since the
coordinate chart of geodesic coordinates is
$\operatorname{exp}_p^{-1}$,
the coordinate representation of $\gamma$ in
this chart is as simple as
$$
\hat{\gamma}(t)=(\underbrace{\varphi}
_{\text{chart}}
\circ \gamma)(t)
=\operatorname{exp}_p^{-1}
\operatorname{exp}_p(tv)=tv.
$$
And the composition $E_i \circ \gamma$ just
means that we take our local
frame \textit{along $\gamma$}. Continue:
$$
\begin{aligned}
&=v^i\nabla^{\gamma}_{\frac{d}{dt}}
E_i\circ\gamma
=v^i\nabla_{\gamma_{v,*} \frac{d}{dt}}E_i\\
&=v^i\nabla_{v^jE_j}E_i=v^iv^j\nabla_{E_j}
E_i\\&
=v^iv^j\Gamma_{ji}^kE_k
\end{aligned}
$$
along $\gamma$. Now choose $v=e_1$. You get
$\Gamma_{11}^k=0$ for all $k$ along
$\gamma_{e_1}$. Now choose $v=e_2$, then
$\Gamma_{22}^k=0$ along $\gamma_{e_2}$, so at
least at $p$ they both vanish. And now choose
$v=e_1+e_2$. You get
$$
\begin{aligned}
0=(v^1)^2 \underbrace{\Gamma_{1 1}^k}_{=0}
+v^1v^2 \Gamma_{1 2}^k+v^2v^1\Gamma_{2 1}^k
+(v^2)^2\underbrace{\Gamma_{2 2}^k}_{=0}
\end{aligned}
$$
So $\Gamma_{12}^k=0$ since Levi-Civita is
torsion-free, i.e. symmetric.
And so on. So the all Christoffel symbols
vanish at the same time at $p$.
\end{proof}

\begin{lemma}[Gauss' lemma]
\label{differential-geometry-lemma-Gauss}
Let $p\in M$ and let $v\in T_pM$ be such
that $\exp_p$ is defined at $v$. Then, after
identifying $T_v(T_pM)\cong T_pM$, we have
$$
\left\langle{}d_v\exp_p(v),d_v\exp_p(w)
\right\rangle{}
=
\left\langle{}v,w\right\rangle{}
$$
for every $w\in T_pM$.
\end{lemma}

\begin{proof}[Sketch of proof]
First suppose $w\perp v$. Let
$\alpha(s)$ be a curve in the Euclidean
sphere of radius $|v|$ in $T_pM$ such that
$\alpha(0)=v$ and $\alpha'(0)=w$. Consider
$$
f(s,t)=\exp_p(t\alpha(s)).
$$
For each fixed $s$, the curve
$t\mapsto f(s,t)$ is a geodesic, and its
speed is $|\alpha(s)|=|v|$, independent of
$s$.

We want to show that the radial direction is
orthogonal to the spherical direction after
applying $d\exp_p$. For this, compute
$$
\frac{d}{dt}\left\langle{}f_t,f_s
\right\rangle{}
=
\left\langle{}f_{tt},f_s\right\rangle{}
+\left\langle{}f_t,f_{st}\right\rangle{}.
$$
The first term vanishes because
$t\mapsto f(s,t)$ is a geodesic. For the
second term, use symmetry of the pullback
connection and the fact that
$
\left\langle{}f_t,f_t\right\rangle{}
$
is
constant in $s$:
$$
\left\langle{}f_t,f_{st}\right\rangle{}
=
\frac{1}{2}\frac{d}{ds}
\left\langle{}f_t,f_t\right\rangle{}
=0.
$$
Thus $\left\langle{}f_t,f_s\right\rangle{}$
is constant in $t$. At $t=0$ it is zero,
because $f(s,0)=p$ is constant. Therefore
$$
\left\langle{}d_v\exp_p(v),d_v\exp_p(w)
\right\rangle{}
=0
$$
whenever $w\perp v$.

If $w$ is radial, say $w=av$, then
$d_v\exp_p(w)=a\gamma'(1)$ for
$\gamma(t)=\exp_p(tv)$, so
$$
\left\langle{}d_v\exp_p(v),d_v\exp_p(w)
\right\rangle{}
=
a|v|^2
=
\left\langle{}v,w\right\rangle{}.
$$
The general case follows by decomposing
$w=w^\parallel+w^\perp$.
\end{proof}

\begin{proposition}
\label{differential-geometry-proposition-naturality-of-exponential}
If $f:M\to \tilde{M}$ is an isometry,
$$
\xymatrix{
B_\varepsilon(0)\subset T_pM \ar[r]^{d_p f}
\ar[d]_{\text{exp}_p} &
B_\varepsilon(0)\subset T_{f(p)}\tilde{M}
\ar[d]^{\text{exp}_{f(p)}} \\
B_\varepsilon(p)\subset M \ar[r]_{f} &
B_\varepsilon(f(p))\subset \tilde{M}
}
$$
\end{proposition}

\begin{proof}
Isometry maps geodesics to geodesics by Lemma
\ref{differential-geometry-lemma-isome-maps-geode-geode},
then uniqueness of solutions
of ODEs.
\end{proof}

\begin{definition}
\label{differential-geometry-definition-killing}
A {\it Killing vector field} is a vector
field
whose flow diffeomorphisms are isometries.
\end{definition}

\begin{exercise}
\label{differential-geometry-exercise-killing}
$X$ is a Killing vector field if and only if
$$
\left\langle{}\nabla_YX,Z\right\rangle{} +
\left\langle{}\nabla_ZX,Y\right\rangle{}=0.
$$
\end{exercise}

\begin{exercise}
\label{differential-geometry-exercise-killi-field-zero-order}
Let $X$ be a Killing field in a connected
Riemannian manifold $M$.
Suppose there is a point $q \in M$ such that
$X(q)=0$ and $\nabla_YX(q)=0$ for any $Y(q)
\in T_qM$.
Show that $X=0$.
\end{exercise}


\begin{exercise}
\label{differential-geometry-exercise-maximum-principle}
\cite{doc}, Chapter III, Exercise 12. If $f$
is {\it subharmonic}, i.e.
$\Delta f\geq 0$, on $M$
compact connected $\implies$ $f$ constant.
\end{exercise}

\begin{proof}
To prove the theorem for subharmonic
functions first we show that in fact they
are harmonic via Eq.
\ref{differential-geometry-equation-divergence-and-volume}
on $X:= \nabla f$ and Stokes:
$$
\int_M \Delta f \operatorname{Vol}=\int_M
\operatorname{div}X \operatorname{Vol}
=\int_M d(i_X
\operatorname{Vol})=\int_{\partial M}i_X
\operatorname{Vol}=0
$$
meaning that the non-negative function
$\Delta f$ is in fact 0, i.e. $f$ is
harmonic. Now we do it again for
$X:=\nabla(f^2/2)$:
$$
\int_M \Delta(f^2/2)\operatorname{Vol}
=\int_M d(i_X \operatorname{Vol})
=\int_{\partial M}i_X \operatorname{Vol}
=0
$$
And then apply Eq.
\ref{differential-geometry-equation-Laplacian-of-product}:
$$
0=\int_M \Delta(f^2/2)\operatorname{Vol}
=\int_M f \Delta f \operatorname{Vol}
+\int_M \left\langle{}\nabla f,\nabla
f\right\rangle{}\operatorname{Vol}
$$
First one vanishes because $f$ is harmonic,
so second one is zero which says
 $f$ is constant!
\end{proof}


\section{Curvature}
\label{differential-geometry-section-curvature}

The {\it curvature tensor} is
\begin{equation}
\label{differential-geometry-equation-curvature-tensor}
R(X,Y)Z=\nabla_X\nabla_YZ-\nabla_Y\nabla_XZ-
\nabla_{[X,Y]}Z
\end{equation}

\begin{proposition}
\label{differential-geometry-proposition-symme-curva-tenso}
The curvature tensor is a tensor and
\begin{enumerate}
\item It is antisymmetric in the first two
and last two entries,
\item
$\left\langle{}R(X,Y)Z,W\right\rangle{}=\left\langle{}R(Z,W)X,
Y\right\rangle{}$,
and
\item Bianchi identity.
\end{enumerate}
\end{proposition}

Given a dimension-2 vector subspace
$\sigma\subset T_pM$, the {\it sectional
curvature} of $M$ in $\sigma$ is
\begin{equation}
\label{differential-geometry-equation-sectional-curvature}
K(\sigma)=\frac{\left\langle{}R(u,v)v,u\right\rangle{}}
{\left\langle{}u,u\right\rangle{}\left\langle{}v,v\right\rangle{}-\left\langle{}u,
v\right\rangle{}^2}
\end{equation}
for any base $\{u,v\}$ of $\sigma$.

\begin{lemma}
\label{differential-geometry-lemma-secti-curva-well-defin}
Sectional curvature does not depend on the
choice of base.
\end{lemma}

\begin{proposition}
\label{differential-geometry-proposition-curva-like-tenso-deter}
Let $R$ and $R'$ be two type $(0,4)$ tensors
on a vector space $\mathbb{V}$. If
their associated sectional curvatures $K$ and
$K'$ are equal, then so are $R$
and $R'$.
\end{proposition}

\begin{proof}
Polarize and make lots of computations, cf.
\cite{doc}.
\end{proof}

If a manifold has constant sectional
curvature $c$,
\begin{equation}
\label{differential-geometry-equation-constant-curvature}
\left\langle{}R(X,Y)Z,W\right\rangle{}=
c(\left\langle{}X,Z\right\rangle{}\left\langle{}Y,W\right\rangle{}-\left\langle{}X,
W\right\rangle{}\left\langle{}Z,Y\right\rangle{})
\end{equation}
since the right-hand-side of Eq.
\ref{differential-geometry-equation-constant-curvature} is a
so-called curvature tensor, i.e. a tensor
satisfying the symmetries of $R$, and
applying Proposition
\ref{differential-geometry-proposition-curva-like-tenso-deter}.

The {\it Ricci tensor} is
\begin{equation}
\label{differential-geometry-equation-Ricci-tensor}
\begin{aligned}
\text{Ric}: TM\times TM &\longrightarrow
\mathbb{R} \\
\text{Ric}(X,Y)&=\frac{1}{n-1}\text{tr}
(Z\mapsto
R(Z,X)Y)
\end{aligned}
\end{equation}
It is symmetric.

The {\it Ricci curvature} is
\begin{equation}
\label{differential-geometry-equation-Ricci-curvature}
\begin{aligned}
\text{Ric}: T^1M &\longrightarrow \mathbb{R}
\\
\text{Ric}(X)=\text{Ric}(X,X)&=\frac{1}{n-1}
\sum_{i=1}^{n-1}K(X,E_i)
\end{aligned}
\end{equation}
where $\{X,E_1,\ldots,E_{n-1}\}$ is an
orthonormal basis.

The {\it scalar curvature} is
\begin{equation}
\label{differential-geometry-equation-scalar-curvature}
\begin{aligned}
\text{scal}: M &\longrightarrow \mathbb{R} \\
\text{scal}(p)&=\frac{1}{n}\text{tr}
\text{Ric}
=\frac{1}{n}\sum_i\text{Ric}(E_i,E_i)
\end{aligned}
\end{equation}

\begin{lemma}
\label{differential-geometry-lemma-scala-curva-mean-secti}
$$
\text{scal}(p)=\frac{1}{n}\sum_{i}\text{Ric}
(e_i,e_i)
=\frac{1}{n(n-1)}\sum_{i,j}K(e_i,e_j)
=\frac{1}{\text{Vol}(S^{n-1})}
\int_{\sigma\subset
T_pM}K(\sigma)
$$
\end{lemma}

\medskip\noindent
The following result is very important.

\begin{lemma}[Symmetry lemma for curvature]
\label{differential-geometry-lemma-symmetry-lemma-curvature}
Let $f:U \subset \mathbb{R}^2 \to M$
be a map into a Riemannian manifold and $V
\in \mathfrak{X}^f$.
Then
$$
\nabla_{\partial_s}\nabla_{\partial_t}V
-\nabla_{\partial_t}\nabla_{\partial_s}V
=R(f_s,f_t)V
$$
where $f_s:=f_*\partial_s$,
$f_t:=f_*\partial_t$.
More generally, the pullback of $R$ is the
Riemannian curvature tensor of the pullback
vector bundle
for any affine vector bundle $(E,\nabla)$ on
$M$.
\end{lemma}


\begin{exercise}
\label{differential-geometry-lemma-curvatures-of-product-metric}
Lista 5, Exercise 1. Let $(M_1,g_1)$ and
$(M_2,g_2)$ be Riemannian manifolds and
let $M_1\times M_2$ be equipped with the
product metric $g:=g_1\oplus g_2$. Show
that the curvature tensor, the Ricci
curvature and the scalar curvature of $g$
are given by
\begin{enumerate}
\item $R=\pi_1^*R_1+\pi_2^*R^2$,
\item
$\text{Ric}=\pi_1^*\text{Ric}_1+
\pi_2^*\text{Ric}_2$,
\item
$\text{scal}=\pi_1^*\text{scal}_1+
\pi_2^*\text{scal}_2$
\end{enumerate}
where $\pi_i:M_1\times M_2\to M_i$ are the
projections.
\end{exercise}

\begin{example}[Curvature of complex
projective space]
\label{differential-geometry-example-curvature-of-CPn}
Defina uma métrica Riemanniana  em
$\mathbb{C}^{n+1}\setminus\{0\}$ do
seguinte modo. Se $Z \in
\mathbb{C}^{n+1}\setminus\{0\}$ e $V,W \in
T_Z
(\mathbb{C}^{n+1}\setminus\{0\})$,

$$
\left\langle{}V,W\right\rangle{}_Z=\frac{\text{Re}(V,W)}{(Z,
Z)}
$$
onde
$$
(Z,W)=z_0\overline{w}_0+\ldots+
z_n\overline{w}_n
$$

é o produto hermitiano em $\mathbb{C}^{n+1}$.
Observe que a métrica
$\left\langle{}\cdot,\cdot\right\rangle{}$ restrita a
$S^{2n+1}\subset
\mathbb{C}^{n+1}\setminus\{0\}$ coincide com
a métrica induzida por
$\mathbb{R}^{2n+2}$.
\begin{enumerate}

\item Mostre que, para todo $0\leq \theta
\leq 2\pi$, $e^{i\theta}:S^{2n+1}\to
S^{2n+1}$ é uma isometria, e que, portanto, é
possível definir uma métrica
Riemanniana em $\mathbb{P}^n(\mathbb{C})$ de
modo que a submersão $f$ seja
Riemanniana.

\item Mostre que, nesta métrica, a curvatura
seccional de
$\mathbb{P}^n(\mathbb{C})$ é dada por
$$
K(\sigma)=1+3\cos^2\varphi
$$
onde $\sigma$ é gerado pelo par ortonormal
$X$, $Y$, $\cos
\varphi=\left\langle{}\overline{X},i\overline{Y}
\right\rangle{}$,
e $\overline{X}$,
$\overline{Y}$ são os levantamentos
horizontais de $X$ e $Y$, respetivamente. Em
particular, $1 \leq  K(\sigma) \leq  4$.
\end{enumerate}
\end{example}

\section{Jacobi Fields}
\label{differential-geometry-section-jacobi-fields}

\begin{proposition}[Everyday Jacobi field]
\label{differential-geometry-proposition-everyday-jacobi-field}
\cite{doc} Chapter V, Proposition 2.4. Let
$\gamma:[0,a]\to M$ be a normalized
geodesic and let $J$ be a Jacobi field
through $\gamma$ with $J(0)=0$.
Let $p:=\gamma(0)$ $w:=J'(0)$ and
$v:=\gamma'(0)$.
Consider the variation
$$
f(s,t):=\text{exp}_{p}(t(v+sw))
$$
Then the associated Jacobi field is $J$. So
in particular the field
$$
J(t)=d_{tv}\text{exp}_p(tw)
$$
is a Jacobi field.

\end{proposition}

\begin{remark}
\label{differential-geometry-remark-everyday-jacobi}
The variation can also be taken
with respect to
any a curve $\sigma$ on $T_vT_{p}M$ passing
through $v$
at $s=0$ with velocity $w$.
\end{remark}

\begin{example}[Jacobi Fields in constant
curvature]
\label{differential-geometry-example-jacob-field-const-curva}
If $M$ has constant curvature $K$ and
$\gamma$ is a normalized geodesic, a
Jacobi field is the solution of the equation
\begin{equation}
\label{differential-geometry-equation-jacob-field-const-curva}
J''+KJ=0
\end{equation}
This follows from Eq.
\ref{differential-geometry-equation-constant-curvature} since it
implies that
for any $X\in\mathfrak{X}_\gamma$, which we
may take orthogonal to $\gamma'$,
$$
\left\langle{}X,R_{\gamma'}J\right\rangle{}
=K(\left\langle{}X,J\right\rangle{}\left\langle{}\gamma',
\gamma'\right\rangle{}
-\left\langle{}X,\gamma'\right\rangle{}\left\langle{}J,
\gamma'\right\rangle{})
$$
so that $R_{\gamma'}J=KJ$. The solutions of
Eq.
\ref{differential-geometry-equation-jacob-field-const-curva} can be
computed taking a
parallel vector field $W$ along $\gamma$
orthogonal to $\gamma'$ and of unit
length. Then
\begin{equation}
\label{differential-geometry-equation-jacob-field-const-curva-2}
J(t)=\begin{cases}
\sin(\sqrt{K}t)W(t)\qquad & K>0\\
tW(t)\qquad &K=0\\
\text{sinh}(\sqrt{-K}t)W(t)\qquad &K<0
\end{cases}
\end{equation}
provided that $J(0)=0$, $J'(t)\perp
\gamma'(t)$ and $J'(0)=W(0)$.
\end{example}

\begin{proposition}
\label{differential-geometry-proposition-jacob-field-taylo-expan}
Differentiating a Jacobi field enough times,
we arrive to
\begin{equation}
\label{differential-geometry-equation-jacob-field-taylo-expan}
\|J(t)\|^2=\|w\|^2t^2-\frac{1}{3}\left\langle{}R(w,v)
v,w\right\rangle{}t^4+O(t^4).
\end{equation}
\end{proposition}

Sometimes it's useful to work with

\begin{definition}
\label{differential-geometry-definition-Jacobi-tensor}
The {\it Jacobi tensor} is the map
$$
\begin{aligned}
	\overline{J}_t:  &\longrightarrow
	T_{\gamma(t)}M \\
	 \overline{J}_t(w)&=d_{tv}\text{exp}_p(tw)
\end{aligned}
$$
\end{definition}

\begin{exercise}
\label{differential-geometry-exercise-Jacobi-tensor-derivative}
Show that $\bar{J}'(0)=\text{Id}$.
\end{exercise}

\begin{proof}
Just notice that since
$\bar{J}_t=d_{tv}\text{exp}_p(tw)$ is linear
on
$tw$ (because it's a differential) we have
$$
\begin{aligned}
\overline{J}'(0)&=\frac{d}{dt}\Big|_{t=0}
(d_{tv}\text{exp}_p(tw))\\
&=\frac{d}{dt}\Big|_{t=0}t\cdot
(d_{tv}\text{exp}_p(w))\\
&=d_{tv}\text{exp}_p(w)|_{t=0}+
t\cdot(\text{something})|_{t=0}\\
&=\text{Id}_{T_pM}
\end{aligned}
$$
\end{proof}

\begin{remark}
\label{differential-geometry-remark-dimen-space-jacob-field}
Lembre o seguinte fato geral mostrado em
aula:
\begin{equation}
\label{differential-geometry-equation-dimen-space-ortho-jacob}
\dim\{J \in \mathfrak{X}^J_\gamma:J(0)=0, J
\perp \gamma\}=n-1
\end{equation}
Isso segue das seguintes observações:
\begin{enumerate}
\item $\dim\mathfrak{X}^J_\gamma=2n$ porque
são soluções de $n$ equações
diferenciais ordinárias de segunda ordem,
i.e. cada campo de Jacobi está
determinado pelas condições iniciais $J(0)$ e
$J'(0)$.
\item $\dim \{J \in
\mathfrak{X}^J_\gamma:J(0)=0\}=n$.
\item $\dim \{J \in
\mathfrak{X}^J_\gamma:J\perp \gamma'\}=2n-2$.
Para
confirmar isso note que pelas simetrias de
$R$ e a equação de Jacobi tem-se
que $\left\langle{}J,\gamma'\right\rangle{}''=0$, pelo que
$\left\langle{}J,\gamma'\right\rangle{}=a+bt$
para dois números reais $a,b$. Segue que
qualquer $J \in
\mathfrak{X}^J_\gamma$ se escreve como
$J(t)=a\gamma'(t)+bt\gamma'(t)+\hat{J}(t)$
para algum $\hat{J}$
perpendicular a  $\gamma'$. (Supondo que
$|\gamma'|=1$.) Então se $J \perp
\gamma$, temos que  $a=b=0$, então tiramos
dois números do $2n$ que
tínhamos.
\item Segue
\ref{differential-geometry-equation-dimen-space-ortho-jacob}.
\end{enumerate}

\end{remark}


\section{Conjugate points}
\label{differential-geometry-section-conjugate-points}

\begin{definition}
\label{differential-geometry-definition-conjugate-points}
A point $p \in M$ is {\it conjugate} to $q
\in M$ along $\gamma$ if $\gamma$ is
a geodesic joining $p$ and $q$ and there is a
Jacobi field
$J\in\mathfrak{X}_\gamma^J$ such that
$J(0)=0$ and  $J(\ell)=0$.
\end{definition}

\begin{proposition}[Critical points of
exponential]
\label{differential-geometry-proposition-critical-points-exp}
\cite{doc} Chapter V, Proposition 3.5. A
vector $v \in T_pM$ is a critical
point of $\text{exp}_p$ if and only if the
corresponding
$\gamma_v(1)$ is a conjugate point of $p$
along $\gamma_v$.
Moreover, in such points, the dimension of
the kernel of
$\text{exp}_p$ at this vector is the
multiplicity
of the conjugate point $\gamma_v(1)$
(i.e. the dimension of the space of Jacobi
fields vanishing
at the endpoints).
\end{proposition}

\begin{proof}
($\implies$) Suppose that $v$ is a critical
point of $\text{exp}_p$ and
let $w \in \ker d_v \text{exp}_p$ then a
Jacobi field along $\gamma$
vanishing at the endpoints is given by
$$
J(t):=d_{tv}\text{exp}_ptw
$$
($\impliedby$) If you have a critical point
along $\gamma$, then the
differential of $\text{exp}_p$ has kernel.

(Mulitplicity.) The proof every element of
the kernel gives a Jacobi
 field vanishing at the points. Linearly
 dependent elements will give
 linearly dependent Jacobi fields.
\end{proof}

\begin{exercise}
\label{differential-geometry-exercise-jacobi-nonpositive-k}
\begin{reference}
\cite[Chapter V, Section 3, Exercise 3]{doc}
\end{reference}
Let $M$ be a Riemannian manifold
of nonpositive sectional curvature.
Show $M$ has no conjugate points.
\end{exercise}

\begin{proof}
Differentiate twice $\langle J,J\rangle$
to find $\langle J',J\rangle$ is an increasing function.
But $\langle J',J\rangle(0)=0=
\langle J',J\rangle(\ell)$,
so $\langle J',J\rangle
=\langle J,J\rangle$' is zero,
and also vanishes at $t=0$.
\end{proof}


\section{Isometric immersions}
\label{differential-geometry-section-isometric-immersions}

Let $f:M \to \tilde{M}$ be an isometric
immersion of Riemannian manifolds. When
computing the covariant derivative along two
vector fields
$X,Y \in\mathfrak{X}(M)$ we get a vetor field
along $X$ that may be split into
its normal and tangent part thanks to
Equation
\ref{differential-geometry-equation-immer-bundl-decom}. By Lemma
\ref{differential-geometry-lemma-tange-compo-pullb-conne-ii} we see
that the
 tangent part is
in fact the Levi-Civita connection of $M$
since it is a symmetric and torsion
free connection on the tangent bundle. Since
the Levi-Civita connection of
$\tilde{M}$ is torsion free, the normal part
is a
symmetric tensor we call {\it second
fundamental form}:
\begin{equation}
\label{differential-geometry-equation-second-fundamental-form}
\nabla^{f}_X f_*Y=f_*\nabla_XY+\alpha(X,Y)
\end{equation}

\medskip\noindent

Now let's differentiate a normal section
$\xi\in T^\perp M$.
We call the tangent part {\it shape
operator}. By Lemma
\ref{differential-geometry-lemma-normal-connection} below, the
normal part is a connection
we call the {\it normal connection}.
\begin{equation}
\label{differential-geometry-equation-shape-operator}
\nabla^f_X\xi=-A_\xi X+\nabla^\perp_X\xi
\end{equation}
\begin{lemma}
\label{differential-geometry-lemma-normal-connection}
$\nabla^\perp$ is a connection on $T^\perp
M$.
\end{lemma}
Notice the adjointness relation
\begin{equation}
\label{differential-geometry-equation-secon-funda-form-shape}
\left\langle{}\xi,\alpha(X,Y)\right\rangle{}=-\left\langle{}A_\xi
X,Y\right\rangle{}
\end{equation}
which follows from computing
$0=X\left\langle{}Y,\xi\right\rangle{}$.

The {\it Gauss equation}
$$
(R_{\nabla^f}(X,Y)Z)^\top=R_\nabla(X,Y)Z-
A_{\alpha(Y,Z)}X+A_{\alpha(X,Z)}Y
$$
may be derived by substituting Eq.
\ref{differential-geometry-equation-second-fundamental-form}
into Equation
\ref{differential-geometry-equation-curvature-tensor}. But most
likely we will use the
following form:
\begin{equation}
\label{differential-geometry-equation-Gauss}
\left\langle{}\tilde{R}(X,Y)Z,W\right\rangle{}=\left\langle{}R(X,Y)Z,
W\right\rangle{}
-\left\langle{}\alpha(X,Z),\alpha(Y,W)\right\rangle{}+
\left\langle{}\alpha(X,W),\alpha(Y,Z)\right\rangle{}
\end{equation}
Substituting for $X=W$ and $Y=Z$ we obtain
the equivalent (by Proposition
\ref{differential-geometry-proposition-curva-like-tenso-deter})
equation
\begin{equation}
\label{differential-geometry-equation-Gauss-sectional-curvature}
K(X,Y)=\tilde{K}(X,Y)+\left\langle{}\alpha(X,Y),
\alpha(Y,Y)\right\rangle{}-\|\alpha(X,Y)\|^2
\end{equation}

Notice that we are seeing the determinant of
the second fundamental form, which
is the product of its eigenvalues, which may
be defined as the {\it principal
curvatures}. In particular, this shows that
the Gaussian curvature of a surface
in $\mathbb{R}^3$, defined as the determinant
of the second fundamental form,
coincides with sectional curvature when
substituting $\tilde{K}\equiv0$ for
 the flat metric on euclidean space. (!!!)

In the case of a hypersurface, the normal
vector $\alpha(X,Y)$ must be
proportional to the unit normal $N$, and by
Eq.
\ref{differential-geometry-equation-secon-funda-form-shape} the
proportionally
factor is exactly $\left\langle{}A_N X,Y\right\rangle{}$.
Using this and evaluating
 at an orthonormal pair $X,Y$ we obtain
\begin{equation}
\label{differential-geometry-equation-gauss-secti-curva-shape}
K(X,Y)=\tilde{K}(X,Y)+\left\langle{}AX,
X\right\rangle{}\left\langle{}AY,Y\right\rangle{}-\left\langle{}AX,Y\right\rangle{}^2
\end{equation}
In this case we are looking at the
determinant of the matrix
\begin{equation}
\label{differential-geometry-equation-matrix-of-shape-operator}
\begin{pmatrix}
\left\langle{}AX,X\right\rangle{}&\left\langle{}AX,Y\right\rangle{}\\
\left\langle{}AX,Y\right\rangle{}&\left\langle{}AY,Y\right\rangle{}
\end{pmatrix}
\end{equation}
which is the matrix representation of the
shape operator $A$ to the plane
spanned by $X,Y$. (That's just a particular
case of the simple fact that the
representation of the matrix $A_i^j$ for
$Ae_i=A_i^je_j$ is given by
$\left\langle{}Ae_i,e_j\right\rangle{}=A^k_i\left\langle{}e_k,
e_j\right\rangle{}=A^j_i$.)

\begin{example}[Sectional curvature of
sphere]
\label{differential-geometry-example-sectional-curvature-of-sphere}
Using Eq.
\ref{differential-geometry-equation-gauss-secti-curva-shape} we can
prove that the sectional
curvature of $S^n$ is 1 once we make sure
that the shape operator of the sphere
is the identity. And this follows from the
fact that the connection on
$\mathbb{R}^n$ is just directional
derivative, see Example
\ref{differential-geometry-example-euclidean-connection}.

More explicitly, the normal component of
$\tilde{\nabla}_X N$ will vanish (see
discussion in Section
\ref{differential-geometry-section-Riccati-equation}), so that
$\tilde{\nabla}_XN=-A_N X$. But the
connection in Euclidean space is just
differentiating the vector field, that is,
$\tilde{\nabla}_X N=N X$. And
finally, who's $N$? It's $-\text{id}$ if you
chose the inner normal. So
 $\tilde{\nabla}_X N=-X$.
\end{example}

\begin{exercise}
\label{differential-geometry-exercise-Gauss-sectional}
Use Eq.
\ref{differential-geometry-equation-gauss-secti-curva-shape}
to show that no
hypersurface $M \hookrightarrow
\mathbb{R}^{n+1}$ can have constant negative
curvature.
\end{exercise}

\begin{proof}
After substituting $\tilde{K}(X,Y)=0$, we see
that Eq.
\ref{differential-geometry-equation-gauss-secti-curva-shape} is
determinant of the second fundamental form.
 But… how to argue that this determinant
 can't be always negative?
\end{proof}

\begin{example}[Sectional curvature of
hyperbolic space]
\label{differential-geometry-example-hyperbolic-space}
\cite{doc}, Chapter VIII, Exercise 3. The
unit normal of hyperbolic space in the
hyperboloid model is also the identity. This
allows for a similar computation,
but Eq.
\ref{differential-geometry-equation-gauss-secti-curva-shape}
doesn't hold
exactly because the
metric is not Riemannian. In fact, the normal
vector has norm $-1$, giving that
hyperbolic space has constant curvature $-1$.
\end{example}

\begin{definition}
\label{differential-geometry-definition-total-geode-subma}
A submanifold $M\subset\tilde{M}$ is {\it
totally geodesic} if its second
fundamental form vanishes.
\end{definition}

\begin{lemma}
\label{differential-geometry-lemma-totally-geodesic-submanifold}
A submanifold is totally geodesic if and only
if every geodesic of $M$ is a
geodesic of $\tilde{M}$.
\end{lemma}

\begin{proof}
The direct implication is clear: if $\alpha
\equiv 0$ and we take any curve
contained in $M$, it will have zero
acceleration in $M$ and $\tilde{M}$
simultaneously. For the converse we check
that all the eigenvalues
of $\alpha$ are zero. Consider any tangent
vector $v$ to $M$ and let $\gamma_v$
be the geodesic in $M$ with initial velocity
$v$. By hypothesis this geodesic
will also be a geodesic of $\tilde{M}$, so
that $\alpha(v,v)=0$. Diagonalizing
$\alpha$ we see it must vanish identically.
\end{proof}

Lemma \ref{differential-geometry-lemma-geode-subma-chara}
says what happens in
general: a geodesic of the submanifold
doesn't need to be a geodesic of the
ambient manifold, but it will always have
acceleration normal to the
submanifold.

\begin{exercise}
\label{differential-geometry-exercise-fixed-point-set-isome-tota}
\cite{pet}, Proposition 5.6.5. Let $f:M\to M$
be an isometry. Show that the
connected component of the fixed points of
$f$ is a totally geodesic
submanifold.
\end{exercise}

\begin{proof}
The trick is to show that the eigenvectors of
$df$ are in bijection with the
fixed point set of $f$. For this consider an
eigenvector, a geodesic realising
the vector, map it with $f$ and notice it
must coincide with the original
geodesic because $f$ is an isometry and the
point is a fixed point. This shows
that the fixed point set is a submanifold
since it is parametrized as a linear
subspace under geodesic coordinate charts. It
also shows that the fixes point
set is totally geodesic.
\end{proof}

\begin{definition}
\label{differential-geometry-definition-minimal-immersion}
An immersion $f:M\to\tilde{M}$ is called {\it
minimal} if the trace of its shape
operator vanishes at all points of $M$ for
any choice of normal vector field.
\end{definition}

\begin{definition}
\label{differential-geometry-definition-minim-free-bound-immer}
A minimal surface $f:\Sigma\to
\mathbb{B}^3\subset \mathbb{R}^3$
is of {\it free boundary} if it is orthogonal
to the boundary of $S^2$.
\end{definition}

\begin{exercise}[Mean curvature and Hessian]
\label{differential-geometry-exercise-mean-curvature}
Lista 4, Exercício 8. Let $(M,g)$ be a
Riemannian manifold. Suppose that there
is a function $f \in C^\infty(M)$ such that
$0 \in \mathbb{R}$ is a regular
value of $f$ and let $\Sigma:=f^{-1}(0)$.
\begin{enumerate}
\item Let $N=\frac{\nabla f}{|\nabla f| \in
\mathfrak{X}^\perp(\Sigma)}$. Mostre
que
$$
\left\langle{}\alpha_{\Sigma}(X,Y),N\right\rangle{}=-
\frac{\text{Hess}(f)(X,Y)}{|\nabla
f|},
$$
para todos $X,Y \in \mathfrak{X}(\Sigma)$.

\item Show that the mean curvature of
$\Sigma$ is given by
$$
H_{\Sigma}=-\frac{1}{n}\mathsf{div}
\left(\frac{\nabla
f}{|\nabla f|}\right)
$$
\end{enumerate}
{\it Note.} The notation in item (1)
identifies fields using the differential of
the inclusion map, as is customary.
\end{exercise}

\begin{proof}
\begin{enumerate}
\item One one hand,
 $$
\left\langle{}\alpha(X,Y),\nabla
f\right\rangle{}=-\left\langle{}Y,A_{\nabla f}X\right\rangle{}
$$
by Eq. \ref{differential-geometry-equation-secon-funda-form-shape}.
On the other hand,
$$
\text{Hess}(f)(X,Y)=\left\langle{}\tilde{\nabla}_X
\nabla f,Y\right\rangle{}=
\left\langle{}-A_{\nabla f}X+\nabla_X^\perp \nabla
f,Y\right\rangle{},
$$
then multiply by $\frac{1}{|\nabla f|}$.

\item To compute the divergence as the trace
of the gradient we need an
orthonormal frame of $M$. So let $E_i$ be the
completion of
$\frac{\nabla f}{|\nabla f|}$ to an
orthonormal frame of $M$ at a point of
$\Sigma$. Then:
\begin{equation}
\label{differential-geometry-equation-diver-norma-gradi}
\begin{aligned}
\text{div}\left(\frac{\nabla f}{|\nabla
f|}\right)
&=\sum_{i=1}^{n-1}
\left\langle{}\tilde{\nabla}_{E_i}\frac{\nabla
f}{|\nabla f|},E_i\right\rangle{}
+\underbrace{\left\langle{}\tilde{\nabla}
_{\frac{\nabla
f}{|\nabla f|}}
\frac{\nabla f}{|\nabla f|},\frac{\nabla
f}{|\nabla f|}\right\rangle{}}
_{=\frac{1}{2}\frac{\nabla f}{|\nabla f|}
\left\langle{}\frac{\nabla f}{|\nabla
f|},\frac{\nabla f}{|\nabla f|}\right\rangle{}=0}
\\
&=\sum_{i=1}^{n-1} \left\langle{}-A_{\frac{\nabla
f}{|\nabla f|}}E_i,E_i\right\rangle{}\\
&=-\text{tr}\left(X \mapsto A_{\frac{\nabla
f}{|\nabla f|}}X\right)\\
&=-nH_{\Sigma}(p)
\end{aligned}
\end{equation}
\end{enumerate}
\end{proof}

Notice that since the Laplacian of a function
$f$ is defined as divergence of
gradient, that is, as the trace of $Z\mapsto
\nabla_Z\nabla f$,
we have for the ambient manifold Laplacian
\begin{equation}
\label{differential-geometry-equation-immersion-Laplacian}
\Delta_{\tilde{g}}
f=\text{div}_{\tilde{g}}(\nabla f)
=\text{tr}(Z \mapsto \tilde{\nabla}_Z\nabla
f)
=\sum_{i=1}^{n-1}\left\langle{}\tilde{\nabla}_{E_i}
\nabla
f,E_i\right\rangle{}
+\left\langle{}\tilde{\nabla}_{\nabla f}\nabla
f,\nabla f\right\rangle{}
\end{equation}
where we must add an extra term since a frame
for the ambient manifold
also needs the normal vector. Then we have
$$
\Delta_{\tilde{g}}f-\text{Hess}_{\tilde{g}}
(\nabla
f,\nabla f)
=\sum_{i=1}^{n-1}\left\langle{}\tilde{\nabla}_{E_i}
\nabla
f,E_i\right\rangle{}
$$
If we multiply both sides by
$\frac{1}{|\nabla f|}$, on the right hand
side we
arrive at the mean curvature $H$ by Eq.
\ref{differential-geometry-equation-diver-norma-gradi} (after
applying Leibniz rule),
 and on the left hand side…
\begin{equation}
\label{differential-geometry-equation-mean-curva-level-sets}
H=\frac{1}{|\nabla
f|}\left(\Delta_{\tilde{g}}
f-\text{Hess}_{\tilde{g}}f
\left(\frac{\nabla f}{|\nabla
f|},\frac{\nabla f}{|\nabla f|}\right)\right)
\end{equation}

{\bf Caveat of Eq.
\ref{differential-geometry-equation-mean-curva-level-sets}.} This
is
Lucas Ambrozio's formula from his talk on
Spheres with Minimal Equators (see
Seminars, Section
\ref{differential-geometry-seminars-semin-secti-spher-minim}).
Now it's almost perfect except for
putting the $\frac{1}{|\nabla f|}$ inside the
Hessian…

\begin{example}[Equatorial spheres]
\label{differential-geometry-example-equatorial-spheres}
Show that equatorial spheres are minimal
hypersurfaces of $S^n$. (See
Seminars, Definition
\ref{differential-geometry-seminars-definition-equator}.)
\end{example}

\begin{proof}
To use Eq.
\ref{differential-geometry-equation-mean-curva-level-sets} we just
need to compute
the submanifold Laplacian $\Delta_gf$ where
$f(x):=\left\langle{}x,v\right\rangle{}$ for a fixed
$v$ which is just the normal vector defining
the equator, i.e. the equator is
just  $\Sigma:=f^{-1}(0)$, the normal plane
to $v$ intersected with $S^n$. First
we need to compute the gradient of $f$, which
by Eq.
\ref{differential-geometry-equation-gradient-in-coordinates} is
just
$$
\nabla
f=g^{ij}\partial_ifE_j=\sum_{i,j}g^{ij}v^iE_j
$$
So that's a problem because in practice we
don't know the coefficients of the
metric, nor do we ever want to compute them.
What this really says is: you don't
know what's the gradient. The next step would
be to compute the Laplacian by
differentiating the gradient
according to Eq.
\ref{differential-geometry-equation-immersion-Laplacian}:
$$
\begin{aligned}
\Delta_gf&=\sum_{i=1}^{n-1}\left\langle{}\nabla_{E_i}
\nabla
f,E_i\right\rangle{}
\end{aligned}
$$
Which just takes me back to Eq.
\ref{differential-geometry-equation-diver-norma-gradi}. So, need
to compute second fundamental form.

The only easy way I see to do this is use the
fact that the geodesics of the
sphere are well-known by Example
\ref{differential-geometry-example-geodesics-of-Sn}. We just need
to
note that the equator is in fact a sphere and
that any curve in the equator is a
geodesic of the equator if and only if it is
a geodesic of the whole $S^n$ (cf.
Lemma
\ref{differential-geometry-lemma-totally-geodesic-submanifold}).
But this is obvious by
definition of equator: it's the unitary
vectors in a hyperplane passing through
the origin, and it is equipped with the
metric induced from $S^n$, which in turn
is induced from $\mathbb{R}^{n+1}$, so that
geodesics of the equator are of the
form Eq. \ref{differential-geometry-equation-geodesics-of-Sn} by
the same reasons as why the geodesics
of $S^n$ of that form.

The other way to do this, following
\cite{doc}, Chapter VI, Exercises 8 (see also
Exercise \ref{differential-geometry-exercise-Clifford-torus}) is
computing the shape operator of the
submanifold with respect to
$\mathbb{R}^{n+1}$ first, which is easy
because
Euclidean space is flat, and then computing
the shape operator with respect to
$S^n$. To do this we just need to compute
inner products of the form
$\left\langle{}A_{N_k}E_i,E_j\right\rangle{}$ since these are
exactly the entries of $A_{N_k}$
in the basis $E_i$ (cf. Eq.
\ref{differential-geometry-equation-matrix-of-shape-operator}). But
I will
do this in the next example (maybe some day
do it for $S^n$ too).
\end{proof}

\begin{exercise}[Clifford torus]
\label{differential-geometry-exercise-Clifford-torus}
\cite{doc}, Chapter VI, Exercises 2 and 8.
Show that the map
$\mathbf{x}:\mathbb{R}^2\to\mathbb{R}^4$
given by
$$
\mathbf{x}(\theta,\varphi)
=\frac{1}{\sqrt{2}}(\cos\theta,\sin\theta,
\cos\varphi,\sin\varphi),
\qquad(\theta,\varphi)\in\mathbb{R}^2
$$
is an immersion of $\mathbb{R}^2$ in the unit
sphere $S^3\subset\mathbb{R}^4$,
whose image is a torus $T^2$ with zero
sectional curvature in the induced
metric. Finally show that it is a minimal
submanifold of $S^3$. This is called
the {\it Clifford torus}.
\end{exercise}

\begin{proof}
This map is an immersion since the
derivatives of sine and cosine cannot vanish
simultaneously. Its image is a torus since
the map can be seen as a map
from $S^1\times S^1$ to $\mathbb{R}^4$.

It is natural to think that this manifold is
the flat torus from Example
\ref{differential-geometry-example-flat-torus}. And it is, because
it is diffeomorphic to a product.
We can use the immersion $\mathbf{x}$ on a
small domain as a local
parametrization of $T^2$.  When
differentiating the parametrization we notice
that the tangent vectors are decomposed as a
pair of vectors tangent to one
circle or to the other. Then the metric
induced from $\mathbb{R}^4$ and then by
$S^3$ is exactly the product metric from
Definition
\ref{differential-geometry-definition-product-metric}. By Lemma
\ref{differential-geometry-lemma-curvatures-of-product-metric}, the
curvature tensor on $S^1\times
S^1$ vanishes identically since the curvature
tensor of any 1-manifold vanishes
by antisymmetry in the first or two last
entries. This means that this torus is
flat. This says that its first fundamental
form has determinant $-1$ when seen
as a submanifold of $S^3$.

To see it is a minimal submanifold of $S^3$
it suffices to show it is totally
geodesic Lemma
\ref{differential-geometry-lemma-totally-geodesic-submanifold}. In
other words, we may
equivalently show that $\alpha\equiv0$ or
that the geodesics of this torus are
geodesics of $\mathbb{R}^n$. However, this
time it is not so immediate what are
the geodesics of $T^2$.

The other way to do this, following
\cite{doc}, Chapter VI, Exercises 8 (see also
Exercise \ref{differential-geometry-exercise-Clifford-torus}) is
computing the shape operator of the
submanifold with respect to $\mathbb{R}^4$
first, which is easy because
Euclidean space is flat, and then computing
the shape operator with respect to
$S^3$. To do this we just need to compute
inner products of the form
$\left\langle{}A_{N_k}E_i,E_j\right\rangle{}$ since these are
exactly the entries of $A_{N_k}$
in the basis $E_i$ (cf. Eq.
\ref{differential-geometry-equation-matrix-of-shape-operator}).

Strategy: we need the following ingredients:
a point $p$ in the submanifold, a
basis $E_i$ of the tangent space of the
submanifold and a basis $N_k$ of the
normal space. For every basic normal vector
$N_k$ we compute its shape operator
as a matrix. Then we pass to $S^n$. We should
only consider the normal vectors
$N_k$ that are tangent to $S^n$---this way
they are normal to the subsubmanifold
$\Sigma$ as a submanifold of $S^n$. Then the
trace of the shape operator with
respect to those normals should be the mean
curvature.

Differentiating $\mathbf{x}$ as a local
parametrization we see that
$$
e_1=\partial_\theta=(-\sin\theta,\cos\theta,
0,0),\qquad
e_2=\partial_\varphi=(0,0,-\sin\varphi,
\cos\varphi)
$$
are an orthonormal basis of $T_pT^2$, while
$$
n_1=\frac{1}{\sqrt{2}}(\cos\theta,\sin\theta,
\cos\varphi,\sin\varphi),
\qquad
n_2=\frac{1}{\sqrt{2}}(-\cos\theta,-
\sin\theta,\cos\varphi,\sin\varphi)
$$
are an orthonormal basis of $T_p^\perp T^2$.
Only $n_2$ is tangent to $S^3$ at
$p$ since $n_1$ is in fact the position
vector $p$, which is normal to $S^3$.

Now observe that by Eqs.
\ref{differential-geometry-equation-secon-funda-form-shape} and
\ref{differential-geometry-equation-second-fundamental-form},
$$
\left\langle{}A_{n_2}e_i,e_j\right\rangle{}=\left\langle{}n_k,
\alpha(e_i,e_j)\right\rangle{}
=\left\langle{}n_k,\nabla^{\mathbb{R}^{n+1}}_{e_i}e_j
-\nabla^{T^2}_{e_i}e_j\right\rangle{}
$$
Now
$$
\nabla^{\mathbb{R}^{n+1}}_{\partial_\theta}
\partial_\theta=
-\cos\theta\partial_\theta-
\sin\theta\partial_\theta
$$
so that
$$
\left\langle{}n_2,-\cos\theta\partial_\theta-
\sin\theta\partial_\theta\right\rangle{}
=\frac{1}{\sqrt{2}}
$$
An analogous computation for
$\partial_\varphi$ shows that
$$
A_{n_2}=\frac{1}{\sqrt{2}}\begin{pmatrix}
1&0\\
0&-1
\end{pmatrix}
$$
\end{proof}

Here are some comments after reading
\cite{Jost}, Section 5.4. Let $f:M\to N$ be
an isometric immersion. This allows to use
the vectors
$$
e_\alpha:=f_*\left(\frac{\partial }{\partial
x^\alpha}\right)=
\frac{\partial f^i}{\partial
x^\alpha}\frac{\partial }{\partial f^i},
$$
(where $f^i$ are local coordinates of $N$
near $f(x)$ using that $f$ is an
isometry) as a basis of
$T_{f(x)}N$. Then, by Eq.
\ref{differential-geometry-equation-second-fundamental-form},
checking that
$f(M)$ is then a minimal submanifold, i.e.
that the trace of its second
fundamental form vanishes, is to say that
\begin{equation}
\label{differential-geometry-equation-isometric-immersion-basis}
(\nabla^{N}_{e_\alpha}e_\alpha)^\perp=0
\end{equation}
where we are implicitly summing over
$\alpha$.

Choosing normal coordinates $\partial_\alpha$
at $M$, using Eq.
\ref{differential-geometry-equation-norma-coord-vanis-chris} to
obtain
that the tangent part to the submanifold
vanishes at the point, we get
\begin{equation}
\label{differential-geometry-equation-minim-subma-norma-coord}
\begin{aligned}
(\nabla^N_{e_\alpha}e_\alpha)^\perp&=
\nabla^N_{e_\alpha}e_\alpha\\
&=\nabla^N
_{\frac{\partial f^i}{\partial
x^\alpha}\frac{\partial }{\partial f^i}}
\frac{\partial f^j}{\partial
x^\alpha}\frac{\partial }{\partial f^j}\\
&=\frac{\partial^2 f^j}{(\partial
x^\alpha)^2}\frac{\partial }{\partial f^j}
+\frac{\partial f^i}{\partial
x^\alpha}\frac{\partial f^j}{\partial
x^\alpha}
\Gamma^j_{ik}\frac{\partial }{\partial
f^j}\quad \text{for }j=1,\ldots,n.
\end{aligned}
\end{equation}
And it is claimed that in arbitrary
coordinates this is just
\begin{equation}
\label{differential-geometry-equation-harmonic-map}
\Delta_Mf^j+\gamma^{\alpha
\beta}(x)\Gamma^j_{ik}(f(x))
\frac{\partial f^i}{\partial
x^\alpha}\frac{\partial f^k}{\partial
x^\beta}=0
\quad \text{for }j=1,\ldots,n.
\end{equation}
(It would be a nice exercise to perform the
passage from normal to arbitrary
coordinates.)
\begin{definition}
\label{differential-geometry-definition-harmonic-map}
A map of manifolds satisfying Eq.
\ref{differential-geometry-equation-harmonic-map} is
called {\it harmonic}.
\end{definition}
Notice that this is not asking that the
coordinate
functions of $f$ must be harmonic. Also
notice that this says that images of
harmonic isometric immersions are minimal by
definition.

\begin{lemma}
\label{differential-geometry-lemma-fibers-of-harmonic-submersion}
Let $f:M\to N$ be a smooth (isometric,
right?) submersion of
Riemannian manifolds. $f$ is
harmonic if and only if its fibers are
minimal submanifolds of $M$.
\end{lemma}

\begin{proof}
Here's my idea: do the same trick as in Eqs.
\ref{differential-geometry-equation-isometric-immersion-basis} and
\ref{differential-geometry-equation-minim-subma-norma-coord} but
now using
$f^{-1}$ instead of $f$, which we may do
since $f$ is a submersion. Since it is
also an isometry, we get the same result,
again with $f^{-1}$ instead of $f$.
Then the claim becomes: $f$ is harmonic if
and only $f$ its local inverse is
harmonic. I think this might follow from
taking the Laplacian of the composition
$f f^{-1}=\text{Id}$, which must vanish
because the identity is harmonic. A
Leibniz rule should appear (modulo some other
metric term? see Eq.
\ref{differential-geometry-equation-Laplacian-of-product}).
\end{proof}

\medskip\noindent

Now we prove that the images of minimal
isometric immersions are critical points
of the volume functional.

\begin{exercise}
\label{differential-geometry-exercise-minim-isome-immer-criti}
Consider an isometric immersion $f_0:(M,g)
\to(\tilde{M},\tilde{g})$ and a
vector field $\xi \in f^*T\tilde{M}$.

Define a variation by:
$$
\begin{aligned}
f: (-\varepsilon,\varepsilon)\times M
&\longrightarrow \tilde{M} \\
f(s,p) &=\text{exp}_{p}(s \xi)
\end{aligned}
$$

Then $f_s:=f(s,\cdot):M \longrightarrow
\tilde{M}$ is an immersion (isometric at
$s=0$ by hypothesis). Now we define the
volume functional, which simply
computes the volume of $M_s$:
$$
S(s):= \int_{M}\operatorname{Vol}_{f_s
^*\tilde{g}}
=\int_M f_s^*\operatorname{Vol}_{\tilde{g}}.
$$
Compute $S'(0)$ and show that critical points
of $S$ are immersed submanifolds with
vanishing mean curvature for every normal
vector.
\end{exercise}

\begin{proof}
How to express volume in any coordinate
system?
\[\sqrt{\det (f^*
_sg)_{ij}}dx^1\wedge\ldots\wedge dx^n\]
Let's differentiate:
$$
\begin{aligned}
\frac{d}{ds}S(s)&=\frac{d}{ds}
\int_M\operatorname{Vol}_{f^*_s\tilde{g}}
=\frac{d}{ds}\int_M \sqrt{\det (f^*
_sg)_{ij}}dx^1\wedge\ldots\wedge dx^n\\
&=\int_M \frac{d}{ds}\sqrt{\det (f^*
_sg)_{ij}}dx^1\wedge\ldots\wedge dx^n
\end{aligned}
$$
We must differentiate the square root of the
determinant of a matrix. By
Exercise
\ref{differential-geometry-exercise-derivative-of-determinant},
$$
\frac{d}{dt}\det A(t)=
\det A(t) \cdot
\operatorname{tr}\left(A(t)^{-1}\cdot
\frac{d}{dt}A(t)\right)
$$
and using that we see that
$$
\frac{d}{dt}\sqrt{\det A(t)}
=\frac{1}{2}\sqrt{\det A(t)}
\cdot\operatorname{tr}\left(A(t)^{-1}\cdot
\frac{d}{dt}A(t)\right)
$$
So we put $A(s)=(f^*_s\tilde{g})_{ij}$. The
good news is that square root of the
determinant part is exactly the local
coordinate function of
$\operatorname{Vol}_{f_0^*\operatorname{Vol}
_{\tilde{g}}}=\operatorname{Vol}_g$.
Then we only have to integrate that other
function, which hopefully is related
to the mean curvature.

Let's compute $$
\begin{aligned}
\frac{d}{ds}\Big|_{s=0}(f_s^*\tilde{g})_{ij}
&=\frac{d}{ds}\Big|_{s=0}f_s^*\tilde{g}
\left(\partial_i,\partial_j\right)
=\frac{d}{ds}\Big|_{s=0}\tilde{g}
(f^s_*\partial_i,f^s_*\partial_j)\\
&=\tilde{g}\left(\nabla^f_{\partial_s}
f^s_*\partial_i,f^s_*\partial_j\right)
+\tilde{g}
\left(f^s_*\partial_i,\nabla^f_{\partial_s}
f_*^s\partial_j\right)\Big|_{s=0}\\
&=\tilde{g}\left(\nabla^f_{\partial_s}
f^0_*\partial_i,f_*^0\partial_j\right)
+\tilde{g}\left(f^0_*\partial_i,
\nabla^f_{\partial_s}f_*^0\partial_j\right)\\
&=\tilde{g}\left(\nabla^f_{\partial_i}
f^0_*\partial_s,f_*^0\partial_j\right)
+\tilde{g}\left(f^0_*\partial_i,
\nabla^f_{\partial_j}f_*^0\partial_s\right)
\qquad  \text{symmetry lemma} \\
&=\tilde{g}(\nabla^f_{\partial_i}\xi,
f_*^0\partial_j)
+\tilde{g}(f^0_*\partial_i,
\nabla^f_{\partial_j}\xi)\end{aligned}
$$
Where
$\nabla^f$ is the pullback connection. The
point is that we arrived at the
variational field $\xi$. If we assume (for
now) that $\xi$ is normal to $M$ then
upon differentiation we get: $$
\nabla^f_{\partial_i}\xi=-A_{\xi}\partial_i+N
$$
where the normal component $N$ vanishes since
it is orthogonal to the basic
tangent vectors of $M$. We conclude that
$$
\begin{aligned}
\frac{d}{ds}\Big|_{s=0}(f_s^*\tilde{g})_{ij}
&=\tilde{g}\left(-A_{\xi}\partial_i,
\partial_j\right)
+\tilde{g}\left(\partial_i,-A_{\xi}
\partial_j\right)\\
&=\tilde{g}(\xi,\alpha(\partial_i,\partial_j)
)
+\tilde{g}(\alpha(\partial_j,\partial_i),\xi)
\\
&=2\tilde{g}(\xi,\alpha(\partial_i,
\partial_j))\qquad\qquad
\text{$\alpha$ is symmetric}\\
&=-2\tilde{g}(A_{\xi}\partial_i,\partial_j)
\end{aligned}
$$

Now we have to multiply by the inverse
matrix, which is just the inverse of the
metric in $M$ because we are at $s=0$. Then
we take trace and obtain the mean
curvature, defined as the trace of the shape
operator:
$$
\begin{aligned}
\operatorname{tr}\Big(g^{ij}g(A
\partial_j,\partial_k)\Big)
=g^{ij}g(A \partial_j,\partial_i)
=g^{ij}g(A^\ell_j\partial_\ell,\partial_i)
=g^{ij}A^\ell_jg_{\ell i}
=\delta^{i}_jA^j_i
=A^j_j
\end{aligned}
$$
We conclude that
$$
\frac{d}{ds}\Big|_{s=0}S(s)=-\int_M H
\operatorname{Vol}_{g}
$$
\end{proof}

\medskip\noindent

Here are two other important equations in the
theory of isometric immersions.
The {\it Codazzi equation} is
\begin{equation}
\label{differential-geometry-equation-Codazzi}
(R_{\nabla^f}(X,Y)Z)^\perp
=(\nabla^\perp_X\alpha)(Y,Z)-(\nabla^\perp_Y,
\alpha)(X,Z)
\end{equation}

Codazzi equation may help us prove… something
cool about bundles I don't quite
remember.

And the {\it Ricci equation} is
\begin{equation}
\label{differential-geometry-equation-Ricci}
\left\langle{}R_{\nabla^f}(X,Y)\xi,\eta\right\rangle{}
=\left\langle{}R_{\nabla^\perp}(X,Y)\xi,\eta\right\rangle{}
-\left\langle{}[A_\xi,A_\eta]X,Y\right\rangle{}
\end{equation}

Ricci equation may help prove that the normal
bundle of a high codimension
submanifold is trivial.

\begin{theorem}[Fundamental theorem of
submanifolds]
\label{differential-geometry-theorem-funda-theor-subma}
See \cite{pet}, Exercise 3.4.21. Let $M$ be a
simply connected Riemannian
manifold and
$(E^p,\widetilde{\left\langle{}\cdot,\cdot\right\rangle{}},
\tilde{\nabla})$
be a
bundle over $M$ with $\tilde{\nabla}$
compatible with
$\widetilde{\left\langle{}\cdot,\cdot\right\rangle{}}$.
Let $\tilde{\alpha}:\mathfrak{X}(M)\times
\mathfrak{X}(M)\to E$ be a symmetric
tensor. Suppose there is a constant $c$ such
that
$(E,\tilde{\nabla},\tilde{\alpha})$ satisfies
Gauss equation for constant
curvature
\ref{differential-geometry-equation-Gauss-constant-curvature},
Codazzi equation
\ref{differential-geometry-equation-Codazzi} and Ricci equation
\ref{differential-geometry-equation-Ricci}.
Then there exists a unique isometric
immersion $f:M \to \mathbb{Q}_c^{n+p}$ and a
bundle isomorphism $\varphi:E \to
T^\perp_pM$ such that
$\varphi\circ\tilde{\alpha}=\alpha_f$.
\end{theorem}

\medskip\noindent
Now I put some exercises.

\begin{exercise}[Singularities of a Killing
vector field]
\label{differential-geometry-exercise-singu-killi-vecto-field}
\begin{reference}
\cite[Exercise 12, Chapter VI]{doc}
\end{reference}
Let $X$ be a Killing vector field on a
Riemannian manifold $M$.
Let $N=\{p \in M: X(p)=0\}$. Show that
\begin{enumerate}
\item If $p\in N$, $V \subset M$ is a normal
neighbourhood of $p$,
and $q \in N \cap V$, then the radial
geodesic segment
joining $p$ to $q$ is contained in $N$.
Conclude that $\gamma \cap V\subset N$.

\item If $p \in N$, there exists a
neighbourhood $V \subset M$
of $p$ such that $V \cap N$ is a submanifold
of $M$.

\item The codimension of any connected
component of $N$ in $M$ is even.
Use that a sphere admitting a
nowhere-vanishing vector field
has odd dimension.
\end{enumerate}
\end{exercise}

\begin{proof}
\begin{enumerate}
\item Consider the flow $\varphi_t$ of $X$
for small $t$.
Since it is an isometry, it must map
geodesics to geodesics.
Since both $p$ and $q$ are in $N$, the two
geodesics
$\gamma$ and $\varphi_t(\gamma)$ intersect at
$p$ and at $q$.
But the exponential is bijective in $V$, so
this cannot happen
unless $X$ vanishes throughout the segment of
$\gamma$ joining $p$ and $q$.
This means that $X$ vanishes throughout
$\gamma$ since
we can take any point in $\gamma$ not between
$p$ and $q$,
apply the flow, which would give another
geodesic but
whose starting segment coincides with
$\gamma$,
and by uniqueness of ODEs such a geodesic
must be $\gamma$ itself.

\item Por indução usando o inciso anterior.

\item Let $p \in N$ and $q$ in a normal
neighbourhood of $p$ in $M$.
Since $N$ is a submanifold we have the
decomposition
$T_qM\simeq T_q N\oplus T_q^\perp N$.
Consider the submanifold
$\tilde{N}:=\text{exp}_q T_q^\perp N$,
and a geodesic sphere contained in this
submanifold.
$X$ cannot vanish along this sphere.
Since $X$ is Killing, it is orthogonal to
such sphere
by $\frac{d}{dt}\left\langle{}X,\gamma'\right\rangle{}=0$
and $\left\langle{}X(q),\gamma'(q)\right\rangle{}=0$.
Thus we have constructed a nowhere-vanishing
vector field
in a hypersphere of $\tilde{N}$,
which shows that $\dim \tilde{N}=\text{codim
}N$ is even.
\end{enumerate}
\end{proof}

\section{Harmonic maps}
\label{differential-geometry-section-harmonic-maps}

{\bf Goal.} Understand harmonic maps $T^2 \to
S^3$. Then gauge theory and
algebraic geometry (spectral data).

Let $M$ be a compact Riemann surface, $G$ be
a compact Lie group with a
bi-invariant metric and $f:M\to G$ a smooth
map.

\begin{example}[To keep in mind]
\label{differential-geometry-example-torus}
$M=T^2$ and $G=\text{SU}(2)\cong S^3$.
\end{example}

\begin{definition}
\label{differential-geometry-definition-connection}
A {\it connection} is a map
$$
\nabla^A:\Omega^{0}(M;G)\to \Omega^{1}(M;E)
$$
where
$\Omega^{p}(M;E)=\Gamma(M,\Lambda^{p}
T^*M\otimes
E)$.
\end{definition}

Consider the pullback bundle of
$(TG,\nabla^{LC})$ over $M$, which we denote
by
$(E,\nabla^A)$.

Locally, we can express
$$
\nabla^A_{\text{loc}}=d+A
$$
where $A$ is the {\it connection matrix}, a
mtrix of 1-forms.

Extend  $\nabla^A$ to an operator
$$
d_A:\Omega^{p}(M;E)\to \Omega^{p+1}(M;E)
$$
by Leibniz rule
$$
d_A(\alpha\wedge\beta=d\alpha\wedge\beta+(-1)
^{|\alpha|}\alpha\wedge
d_A\beta
$$
And we have a {\it curvature} given by
$$
F_A=d^2_A\in\Omega^{2}(M,;\text{End}(E))
$$

\begin{lemma}
\label{differential-geometry-lemma-derivative-of-differential}
$d_A(df)=0$
\end{lemma}

\begin{proof}
If $\alpha\in\Omega^{1}(M;E)$ and
$X,Y\in\mathfrak{X}(M)$,
$$
d_A(\alpha)(X,Y)=\nabla_X^A(\alpha(Y))-
\nabla_Y^A(\alpha(X))-\alpha([X,Y])
$$
So for $\alpha=df$… (use that LC is
torsion-free!)
\end{proof}

\begin{definition}
\label{differential-geometry-definition-Hodge-star}
The {\it Hodge star operator} is
$$
d^*_A=-*d_A*
$$
where
$$
\begin{aligned}
*: \Omega^{p}(M;E) &\longrightarrow
\Omega^{2-p}(M;E) \\
\alpha\otimes s &\longmapsto *\alpha\otimes s
\end{aligned}
$$
since our manifold is dimension 2.
\end{definition}

\begin{definition}
\label{differential-geometry-definition-energy-functional}
$$
E(f)=\frac{1}{2}\int_M|df|^2d\text{Vol}
$$
\end{definition}

\begin{definition}
\label{differential-geometry-definition-harmonic-map-duplicate-2}
A {\it harmonic map} of smooth manifolds
$f:M\to N$ is a critical point of the
energy functional.
\end{definition}

\begin{lemma}
\label{differential-geometry-lemma-criti-point-energ-funct}
The critical points of $E$ satisfy
$$
\tau(f):=\text{tr}_g(\nabla df)=0
$$
where $\nabla df$ is the {\it induced
connection} on $T^*M\otimes E$, i.e.
$$
\nabla(df)(X,Y)=\nabla_X^A(df(Y))-
df(\nabla_X^gY)
$$
\end{lemma}

\begin{lemma}
\label{differential-geometry-lemma-harmonic}
$f$ is harmonic if and only if $d_A^*(df)=0$.
\end{lemma}

\begin{proof}
Choose coordinates such that $\nabla^g_XY$
vanishes at the point. Then
$$
\tau(f)=\nabla^A_{e_1}(df(e_1))+
\nabla^A_{e_1}(df(e_2))
$$
Let's compute $d_A^*(df)=-*d_A*(df)$. First
we find that $*df=e^2s_1-e^1s_2$ for
a dual frame $e^i$ and now no longer writing
the tensor product. Then we get
that
$$
\begin{aligned}
d_A(*df)&=de^2\otimes s_1-e^2\nabla^A(s_1)-
\end{aligned}
$$

\end{proof}

\begin{remark}
\label{differential-geometry-remark-harmonic-functions}
In the case of a function
$f:\Omega\subset\mathbb{C}\to \mathbb{R}$ we
obtain
the usual Laplacian (Exercise!).
\end{remark}

\section{Riccati equation}
\label{differential-geometry-section-Riccati-equation}

Take any point and any tangent vector at that
point. Take the geodesic and take
the geodesic sphere of raduis say $r$ with
center in the point. It's
a hypersurface. Normal vectors are all
multiples of the unit normal vector. But
who's that, the unit normal? Ah, it's the
speed of the geodesic.

So? The normal derivative of any vector field
tangent to the sphere with respect
to the unit normal $\gamma'(r):=\nu$ is…
$$
0=X\left\langle{}\nu,\nu\right\rangle{}=2\left\langle{}\nabla^\perp_X
\nu,\nu\right\rangle{}
$$
that is, normal derivative of $\nu$ with
respect to $X$ is perpendicular to
$\nu$. But normal derivative is normal so
it's actually a multiple of $\nu$,
which makes $\nu$ have norm zero or else
$\nabla^\perp_X \nu=0$.

Now consider everyday Jacobi variation
\ref{differential-geometry-proposition-everyday-jacobi-field}
 $f(s,t):=\text{exp}_p(t(v+sw))$ with Jacobi
 field
$J(t)=d_{tv}\text{exp}_p(tw)$. Differentiate
with using the pullback connection
along $f$:
$$
\begin{aligned}
\nabla_{\frac{d}{dt}}^\gamma
J(t)&=\nabla_{\partial_t}^ff_s=
\nabla_{\partial_t}^f
\partial_sf\\
&=\nabla_{\partial_s}^f \partial_t
f=\nabla_{J}^f \gamma' = -A_{\gamma'}f
\end{aligned}
$$
using that by our discussion above the normal
derivative component vanishes. We
conclude that
\begin{equation}
\label{differential-geometry-equation-jacob-field-deriv-shape}
J'=AJ
\end{equation}
Finally we differentiate that to obtain
$$
J''=A'J+AJ'=A'J+A^2J
$$
that is,
\begin{equation}
\label{differential-geometry-equation-Riccati}
A'+A^2+R_{\gamma'}=0
\end{equation}
since we can let $J(r)$ be any vector
whatsoever and make this construction
work.
\begin{exercise}
\label{differential-geometry-exercise-limit-of-shape-operator}
Show that $\lim_{s\to0}sA(s)=\text{Id}$.
\end{exercise}

\begin{proof}[Solution]
By Rafael. First we present a naive
computation. We will use that
$A(t)(J(t))=J'(t)$,  $J(0)=0$ and $J'(0)\neq
0$. Let $\varepsilon>0$ and let
$s>0$ be such that
$$
\frac{\|J(s)-J(0)-sJ'(0)\|}{s}
<\frac{\varepsilon}{2}
$$
and
$$
\|J'(s)-J'(0)\|<\frac{\varepsilon}{2}.
$$
Then
$$
\begin{aligned}
\left\|A(s)(J(s))-\frac{J(s)}{s}\right\|&
=\left\|J'(s)-\left(\frac{J(s)-J(0)}{s}
\right)\right\|\\
&\leq
\|J'(s)-J'(0)\|+\frac{\|J(s)-J(0)-sJ'(0)\|}
{s}\\
&\leq\varepsilon/2+\varepsilon/2=\varepsilon
\end{aligned}
$$
This shows that $A=\frac{1}{s}J$ for small
$s$, but we have differentiated as if
we were in $\mathbb{R}^n$. A similar
computation but picking $s$ such that
$$
\frac{|\left\langle{}J(s),e_i(s)\right\rangle{}-\left\langle{}J(0),
e_i(0)\right\rangle{}
-s\left\langle{}J'(0),e_i(0)\right\rangle{}|}{s}<\varepsilon/
2
$$
and
$$
|\left\langle{}J'(0),e_i(0)\right\rangle{}-\left\langle{}J'(s),e_i(s)
\right\rangle{}|<\varepsilon/2
$$
for a parallel frame $e_i$ along $\gamma$
allows to conclude as desired.
\end{proof}

\section{Hopf-Rinow Theorem}
\label{differential-geometry-section-Hopf-Rinow-theorem}

\begin{definition}
\label{differential-geometry-definition-non-extendable}
A Riemannian manifold $M$ is {\it extendable}
if there is another Riemannian manifold
$\tilde{M}$
such that $M$ is isometric to a proper open
subset of $\tilde{M}$.
\end{definition}

\begin{theorem}[Hopf-Rinow]
\label{differential-geometry-theorem-Hopf-Rinow}
The following are equivalent for any $p \in
M$:
\begin{enumerate}
\item The domain of the exponential map
$\text{exp}_p$ is $T_pM$.
\item Closed and bounded subsets are compact.
\item $(M,d)$ is a complete metric space.
\item $M$ is geodeiscally complete (i.e.
every geodesic is maximally defined in
all of $\mathbb{R}$).
\item Existence of exhaustions.
\end{enumerate}
Moreover, each of the former conditions imply
that there exists a minimizing
geodesic joining any two points.
\end{theorem}

\begin{exercise}
\label{differential-geometry-exercise-divergent-curves}
\begin{reference}
\cite[Exercise 3.5, Chapter VII]{doc}
\end{reference}
A {\it divergent curve} in a Riemannian
manifold $M$
is a smooth map $\gamma:[0,\infty) \to M$
such that for any compact $K \subset M$
there exists $t_0\in(0,\infty)$
such that $\gamma(t) \not\in K$ for $t>t_0$.
The length of $\gamma$ is defined by
$$
\ell(\gamma)=\lim_{t \to \infty}
\int_0^t|\gamma'(t)|dt.
$$

Show that $M$ is complete if and only if
the length of any divergent curve is
infinite.
\end{exercise}

\begin{proof}
Suppose $M$ is complete and $\gamma$ is a
divergent curve
with finite length. Choose a ball $B$ of
radius $\ell(\gamma)$
containing $\gamma$. The closure of $B$ is
compact,
but $\gamma$ is contained in it.

For the converse …
\end{proof}

\begin{exercise}
\label{differential-geometry-exercise-ray}
A geodesic $\gamma:[0,\infty) \to M$
in a Riemannian manifold $M$ is a {\it ray
starting at $\gamma(0)$}
if it is minimizing between $\gamma(0)$ and
$\gamma(s)$
for all $s \in (0,\infty)$.
Let $M$ be complete, non-compact and $p \in
M$.
Construct a ray starting at $p$.
\end{exercise}

\begin{proof}
Since $m$ is not compact and complete,
we can pick a sequence $p_n \in m\setminus
\bar{b}_n(p)$
and minimizing geodesics $\gamma_n$ from $p$
to $p_n$.
we would like to consider a curve by taking
limit
and show it's minimizing.
however pointwise convergence would not
ensure smoothness of the limit curve.
instead, we consider the normalized initial
velocity vectors
of the curves $\gamma_n$ and take the limit
inside
the compact unit sphere in the tangent space
and apply the exomponential map.
This defines a geodesic $\gamma$ which
is minimizing for any $s \in (0,\infty)$.
Indeed, supposing that our geodesics
are normalized, we have
$$
\begin{aligned}
s &= \ell(\gamma(s))\\
&< s + d(\gamma_n(s),\gamma(s))\\
&< \ell(\gamma_n(s)) +
d(\gamma_n(s),\gamma(s)),
\end{aligned}
$$

\noindent
To conclude consider a curve $\alpha$
connecting $p$ and $\gamma(s)$,
we extend it to a curve $\alpha_n$ connecting
$p$ and $\gamma_n(s)$.
Since $\gamma_n$ is minimizing,
$$
\ell(\gamma_n(s))\leq \ell(\alpha_n)
=\ell(\alpha)+d(\gamma_n(s),\gamma(s)).
$$
The conclusion follows since we showed that
the
length of $\gamma$ up to $\gamma(s)$
is smaller than the length of $\alpha$ up to
$\gamma(s)$
plus a term that converges to zero as $n\to
\infty$.
\end{proof}

\begin{exercise}
\label{differential-geometry-exercise-contraction-and-completeness}
Let $M$ and $\overline{M}$ be Riemannian
manifolds
and $f:M \to \overline{M}$ a diffeomorphism.
Suppose $\overline{M}$ is complete and there
is a constant $c>0$
such that
$$
|v|\geq c|d_pf(v)|
$$
for all $p \in M$ and $v \in T_pM$.
Show that $M$ is complete.
\end{exercise}

\begin{proof}
Let $p_n$ be a Cauchy sequence in $M$.
The given inequality is used to prove
that $f(p_n)$ is also Cauchy in
$\overline{M}$,
and thus converges to a point $f(p)$.
Convergence of the preimage follows from
continuity of $f$.
\end{proof}

\begin{exercise}
\label{differential-geometry-exercise-local-isometry-hopf-rinow}
\begin{reference}
\cite[Exercise 8, Chapter VIII]{doc}
\end{reference}
Let $M$ be a complete Riemannian manifold,
$\overline{M}$ connected Riemannian manifold
and $f:M \to \overline{M}$ a local isometry.
Suppose that any two points of $\overline{M}$
can be connected be a unique geodesic
of $\overline{M}$.
Show that $f$ is injective and surjective
(and thus a global isometry).
\end{exercise}

\begin{proof}
For injectivety consider the fiber of
any point of $\overline{M}$. Since $M$
is complete, we can connect any two distinct
points in the fiber by a geodesic of $M$.
Since $f$ is a local isometry, it maps
this geodesic to a closed geodesic of
$\overline{M}$, which is impossible:
two nearby points along a closed geodesic
have to geodesics joining them, the short
way and the long way along the geodesic.

For surjectivity suppose there is a point
outside the image of $f$. Connect such a
point with a point inside the image of $f$
with the geodesic of $\overline{M}$
that connects them. At the point in the image
of $f$,  $f^{-1}$ takes this geodesic segment
to a geodesic segment of $M$,
which can be extended
throughout $\mathbb{R}$ since $M$ is
complete. But then the whole geodesic in
$\overline{M}$ would be contained in the
image of the geodesic in $M$ since both
are determined by their initial velocities
and point.
\end{proof}

\begin{exercise}
\label{differential-geometry-exercise-bounded-flow-speed}
\begin{reference}
\cite[Exercise 11, Chapter VII]{doc}
\end{reference}
Let $M$ be a complete Riemannian manifold
and let $X$ be a smooth vector field on $M$.
Suppose that $|X(p)|<c$ for all $p \in M$.
Show that the integral curves of $X$
are defined for all $t$.
\end{exercise}

\begin{proof}
The basic idea is:
why if the manifold is complete
every geodesic,
characterized by having constant
velocity,
is maximally defined in $\mathbb{R}$?
The same trick should work in this case.

Suppose that an integral curve $\varphi$
is defined maximally at $(a,b)$.
Choose a sequence converging to $b$,
which is a Cauchy sequence in $\mathbb{R}$.
This gives a Cauchy sequence in $M$
along $\varphi$ by continuity.
By completeness of $M$, $\varphi(t_n) \to p$
for some $p \in M$.
If $X(p) \neq 0$ we can integrate locally
to find a flow $\psi$ at $p$,
which must coincide with $\varphi$ in
the intersection of both curves,
thus extending $\varphi$.
If $X(p)=0$…
\end{proof}

\begin{exercise}
\label{differential-geometry-exercise-homogeneous-manifold}
A Riemannian manifold $M$ is
{\it homogeneous} if for any $p,q \in M$
there is an isometry of $M$ that takes
$p$ to $q$. Show that any homogeneous
manifold is complete.
\end{exercise}

\begin{proof}
Let $\gamma$ be any geodesic. Consider
a geodesic ball $B_\varepsilon(\gamma(0))$
of radius $\varepsilon$
about $\gamma(0)$. Let $t \neq 0$
such that $\gamma(t)$ is well defined.
Take $\gamma(0)$ to $\gamma(t)$
with an isometry. Such an isometry
also takes $B_\varepsilon(\gamma(0))$
to $B_\varepsilon(\gamma(t))$.
This shows we can extend $\gamma(t)$
to $\gamma(t+\varepsilon/2)$
for any $t$ in the domain of definition
of $\gamma$, i.e. $\gamma$ is a complete
geodesic.
\end{proof}

\section{Hadamard's theorem}
\label{differential-geometry-section-Hadamard}

\noindent
Proposition
\ref{differential-geometry-proposition-exp-nonpositive-curvature-expanding}
is the main tool
to prove Hadamard's theorem. It
is basically an application
of Rauch's Theorem \ref{differential-geometry-theorem-rauch-clean}.
The idea is to disguise the image
of $w$ under the differential of
the exponential as a Jacobi field,
and compare it with the constant Euclidean
field. Since the curvature is nonpositive,
Rauch's theorem flips the inequality
and gives the result. But to apply
Rauch we need that the angle between
$w$ and $v$ is the same as the angle
of our comparison vectors, which are
perpendicular. So, we need to split
$w$ into tangential and perpendicular
parts with respect to $v$, and apply
Gauss lemma, which, recalling its
precise statement, allows us to
forget the tangential part ---
exponential will preserve the
lengths of the tangential part,
and preserve orthogonality
of the orthogonal part.

Having 
Proposition
\ref{differential-geometry-proposition-exp-nonpositive-curvature-expanding}
at hand, we use 
Proposition 
\ref{differential-geometry-proposition-local-diffeomorphism-covering}
to obtain that the exponential map
is a covering space. But since $M$
is simply connected, the identity map
is a covering space, so that the
identity factors through the exponential,
providing an inverse smooth map
with $\mathbb{R}^n\cong T_pM$.

Of course,
to apply 
Proposition
\ref{differential-geometry-proposition-local-diffeomorphism-covering}
that we have to use
Proposition
\ref{differential-geometry-proposition-critical-points-exp}
and
Exercise \ref{differential-geometry-exercise-jacobi-nonpositive-k}
to show that the exponential is a local
diffeomorphism.

\begin{proposition}[Differential of the
exponential in nonpositive curvature]
\label{differential-geometry-proposition-exp-nonpositive-curvature-expanding}
Let $M$ be a complete Riemannian manifold
with $K\leq 0$. Fix $p\in M$. Then
$\exp_p:T_pM\to M$ is defined on all of
$T_pM$, and for every $v\in T_pM$ and every
$w\in T_v(T_pM)\cong T_pM$,
$$
|d_v\exp_p(w)|\geq |w|.
$$
\end{proposition}

\begin{proof}
Completeness implies, by Hopf-Rinow, that
$\exp_p$ is defined on all of $T_pM$.

Fix $v\in T_pM$ and decompose
$w=w^\parallel+w^\perp$ into the part
parallel to $v$ and the part perpendicular to
$v$. By Gauss' lemma,
$d_v\exp_p(w^\parallel)$ is orthogonal to
$d_v\exp_p(w^\perp)$, and the radial part
preserves length:
$$
|d_v\exp_p(w^\parallel)|=|w^\parallel|.
$$
Hence
$$
|d_v\exp_p(w)|^2=
|w^\parallel|^2+
|d_v\exp_p(w^\perp)|^2.
$$
So it is enough to prove
$$
|d_v\exp_p(w^\perp)|\geq |w^\perp|.
$$

Let $\gamma(t)=\exp_p(tv)$ and consider the
Jacobi field
$$
J(t)=d_{tv}\exp_p(tw^\perp).
$$
Then $J(0)=0$, $J'(0)=w^\perp$, and
$J(t)\perp\gamma'(t)$. Compare $J$ with the
Euclidean Jacobi field
$$
\tilde{J}(t)=tw^\perp.
$$
Since $K\leq 0=K_{\mathbb{R}^n}$, Rauch
comparison gives
$$
|J(1)|\geq |\tilde{J}(1)|
=|w^\perp|.
$$
Thus
$$
|d_v\exp_p(w)|^2
=
|w^\parallel|^2+
|d_v\exp_p(w^\perp)|^2
\geq
|w^\parallel|^2+|w^\perp|^2
=
|w|^2.
$$
\end{proof}

\begin{proposition}
\label{differential-geometry-proposition-local-diffeomorphism-covering}
\begin{reference}
\cite[Chapter VII, Lemma 3.3]{doc}
\end{reference}
Let $f:\tilde M\to M$ be a local
diffeomorphism of Riemannian manifolds.
Suppose that $\tilde M$ is complete and that
$$
|d_pf(v)|\geq |v|
$$
for every $p\in\tilde M$ and every
$v\in T_p\tilde M$. Then $f$ lifts curves.
In particular, $f$ is a covering map.
\end{proposition}

\begin{proof}
Let $\gamma:[0,a]\to M$ be a smooth curve
and let $\tilde p\in \tilde M$ be such that
$f(\tilde p)=\gamma(0)$. Since $f$ is a
local diffeomorphism, there is a unique local
lift $\tilde\gamma$ of $\gamma$ starting at
$\tilde p$. Let $[0,b)$ be the maximal
interval on which this lift is defined.

For $t<b$ we have
$$
\gamma'(t)=d_{\tilde\gamma(t)}f
(\tilde\gamma'(t)).
$$
By hypothesis,
$$
|\tilde\gamma'(t)|\leq |\gamma'(t)|.
$$
Thus the length of $\tilde\gamma|_{[s,t]}$
is bounded by the length of
$\gamma|_{[s,t]}$. In particular, if
$t_i\to b$, then $\tilde\gamma(t_i)$ is a
Cauchy sequence. Since $\tilde M$ is complete,
it converges to some point $\tilde q$.

Because $f$ is a local diffeomorphism near
$\tilde q$, the lift can be continued past
$b$. This contradicts maximality unless
$b=a$. Therefore every curve admits a lift.

Finally, a local diffeomorphism with curve
lifting is a covering map.
\end{proof}

\begin{theorem}[Hadamard]
\label{differential-geometry-theorem-Hadamard}
If $M$ is a complete simply connected
manifold with $K\leq 0$, then the
exponential map is a diffeomorphism.
\end{theorem}

\begin{proof}

\end{proof}

\section{Manifolds of constant sectional
curvature}
\label{differential-geometry-section-manif-const-secti-curva}

The following theorem and proof where
explained by Dimitri Korshunov (IMPA) at
Geometric Structures Seminar on July 3, 2025.

{\bf Abstract.} Crystallographic group is a
discrete co-compact subgroup of the
group of isometries of a Euclidean space. A
theorem of Bieberbach says that it
contains a subgroup of translations of finite
index. Or, in other words, any
compact flat Riemannian manifold is a finite
quotient of a flat torus. We will
prove this theorem following a streamlined
geometric argument of Vinberg.

\begin{theorem}[Bieberbach]
\label{differential-geometry-theorem-Bieberbach}
In particular, any flat manifold is a
quotient of a torus.
\end{theorem}

\begin{proof}[Sketch of proof]
First we show that it suffices to find only
one pure translation, that is, an
inductive argument shows that if we have a
pure translation, then we can find
others. Perhaps this is the hardest part.

Then we obtain some bounds in operator norm
for the pure translations, this part
looks straightforward.  One of these bounds
is that $\|[A,B]-\text{Id}\|\leq
2\|\text{Id}-A\|\|\text{Id}-B\|$. The other
of the bounds is roughly as follows.
For $g=Ax+a$ and $h=Bx+b$, we have that
$[g,h]=[A,B]x+c$ for some $c$ whose norm
we shall bound by $r$ later.

We use the first bound to show that if two
elements in the group are close to
the identity, then their commutator will be
even closer. This shows that there
is a finite set of small generators
$g_1,\ldots,g_k$. Each of them is of the
form $g_i=A_ix+a_i$.

To arrive at a contradiction suppose that
there are no pure translations to
obtain a contradiction.

Let $r(g)$ be the norm of the translational
part of a transformation $g$. Let
$r$ be the maximum of the translational parts
of the $g_i$.
Let $S:=\{g:r(g)<r\}$. It's a finite set. Let
$g_0$ be minimal in $S$. After
some discussion we show that the linear part
of $g_0$ commutes with the linear
part of $g_i$ for all $i$. In fact, this
implies that $g_0$ itself is in the
center of the group.

Finally we obtain a contradiction with
compactness… simple argument.
\end{proof}

\section{First Variation Formula}
\label{differential-geometry-section-first-variation}

\noindent
The notion of variation is beautiful because
it lets us look at a curve by placing it
inside a smooth family of nearby curves. The
important object is not only the family
itself, but the vector field along the
original curve which records, at each time
$t$, the infinitesimal direction in which the
curve is being moved. This is the
variational field.

\begin{definition}
\label{differential-geometry-definition-variation}
A {\it variation} of a curve
$\gamma:[0,a]\to M$ is a smooth map
$f:(-\varepsilon,\varepsilon)\times[0,a]\to
M$ such that $f(0,t)=\gamma(t)$. We write
$\gamma_s(t):=f(s,t)$ for the curves in the
variation, and
$f_s=f_*\partial_s$ and $f_t=f_*\partial_t$.
The {\it variational field} of $f$ along
$\gamma$ is
$$
V(t)=f_s(0,t).
$$
\end{definition}

\noindent
The energy of a curve is similar
to the length, but instead of integrating
the norm of the speed,
we integrate the norm squared.
For a given variation, the energy
function computes the energy of each of
the curves in the family.

\begin{definition}
\label{differential-geometry-definition-energy}
The {\it energy function} of a variation
$f:(-\varepsilon,\varepsilon)\times[0,a]\to
M$ is
$$
E(s)=\int_0^a\left\langle{}f_t(s,t),f_t(s,t)\right\rangle{}\,dt.
$$
\end{definition}

\noindent
Recall Schwarz inequality from functional
analysis:
$$
\left(\int_0^afg\right)^2\leq
\int_0^af^2\int_0^ag^2
$$
Setting $f\equiv1$ and $g=f_t$ we obtain
\begin{equation}
\label{differential-geometry-equation-energy-bounds-length}
\ell(\gamma)\leq aE(\gamma)
\end{equation}

\medskip\noindent
As a preliminary remark, observe that,
in contrast with the second variation
formula, the first variation formula
does not need evaluation at $s=0$,
and, importantly, we do not need to assume
that $\gamma$ is a geodesic.

\begin{theorem}[First Variation Formula]
\label{differential-geometry-theorem-first-variation-formula}
If  $f:(-\varepsilon,\varepsilon)\times[0,a]
\to M$ is a variation of
the geodesic $\gamma:[0,a]\to M$, then
\begin{equation}
\label{differential-geometry-equation-first-variation-formula}
\frac{1}{2}E'(0)=-\int_0^a\left\langle{}V,
\gamma''\right\rangle{}+\left\langle{}V,\gamma'\right\rangle{}|_{0}^a
+\sum_{i=1}^{k-1}\left\langle{}V,\gamma'(t_{i}^-)-
\gamma'(t_i^+)\right\rangle{}
\end{equation}
\end{theorem}

\begin{proof}
If $\gamma$ is smooth,
$$
E'(s)=
\left(\int_0^a\left\langle{}f_t,t_s\right\rangle{}\right)_s
=2\int_0^a\left\langle{}f_{ts},f_t\right\rangle{}.
$$
By the Symmetry Lemma,
$f_{ts}=f_{st}$.
Now since the connection is metric,
we can make the $\left\langle{}f_{st},f_t\right\rangle{}$
term appear by differentiating
$\left\langle{}f_s,f_t\right\rangle{}$
with respect to $t$:
$$
\left\langle{}f_s,f_t\right\rangle{}_t=\left\langle{}f_{st},
f_t\right\rangle{}+\left\langle{}f_s,f_{tt}\right\rangle{}.
$$
Substituting we obtain
$$
\frac{1}{2}E'(s)
=\left\langle{}f_s,f_t\right\rangle{}|_0^a
-\int_0^a\left\langle{}f_s,f_{tt}\right\rangle{}.
$$

If $\gamma$ is not smooth, the integral of
$\left\langle{}f_s,f_t\right\rangle{}_t$ must be
done part by part, obtaining the sum in Eq.
\ref{differential-geometry-equation-first-variation-formula}.
\end{proof}

\begin{exercise}
\label{differential-geometry-exercise-first-var-submfd-dist}
\begin{reference}
\cite[XIX]{doc}
\end{reference}
Let $M$ be a complete Riemannian manifold,
and let $N$ be a closed submanifold of $M$.
Let $q \in M\setminus N$. Show that
there exists a point $p_0 \in N$
such that the distance $d(q,N)$ is
realised at $p_0$, and that
any minimizing geodesic joining $p_0$ to $q$
is orthogonal to $N$ at $p_0$.
\end{exercise}

\begin{proof}
The existence of $p_0$ is a consequence of
completeness of $M$ and closedness of $N$.
By definition of $d(q,N)$ as an infimum, for
each $n \in \mathbb{Z}_+$ we can choose a
point $p_n \in N$ such that
$$
d(q,p_n)<d(q,N)+\frac{1}{n}.
$$
Thus the sequence $(p_n)$ is contained in a
closed metric ball centered at $q$. Since $M$
is complete, Hopf-Rinow implies that this
closed ball is compact. Hence, after passing
to a subsequence, $p_n$ converges to some
$p_0 \in M$. Since $N$ is closed and
$p_n \in N$, we have $p_0 \in N$. Finally,
by continuity of the distance function,
$$
d(q,p_0)=\lim_n d(q,p_n)=d(q,N).
$$

Let $\gamma:[0,a] \to M$ be a minimizing
geodesic from $q$ to $p_0$, parametrized by
arc length. Take $v\in T_{p_0}N$ and choose a
curve $c:(-\varepsilon,\varepsilon)\to N$
with $c(0)=p_0$ and $c'(0)=v$. Extend $v$ to
a vector field $V$ along $\gamma$ such that
$V(0)=0$ and $V(a)=v$, and choose a variation
$f$ of $\gamma$ with variational field $V$,
$f(s,0)=q$, and $f(s,a)=c(s)$.

Since $p_0$ realizes the distance from $q$ to
$N$, every curve $f(s,\cdot)$ joins $q$ to a
point of $N$, so
$$
\ell(f(s,\cdot))\geq d(q,N)=a.
$$
By Schwarz inequality,
$E(s)\geq \ell(f(s,\cdot))^2/a\geq a=E(0)$.
Thus $E'(0)=0$. The first variation formula
gives
$$
0=\left\langle{}V,\gamma'\right\rangle{}|_0^a
=\left\langle{}v,\gamma'(a)\right\rangle{}.
$$
Since $v\in T_{p_0}N$ was arbitrary, $\gamma$
is orthogonal to $N$ at $p_0$.
\end{proof}

\begin{exercise}
\label{differential-geometry-exercise-endpoint-variation}
Let $\gamma:[0,a]\to M$ be a smooth curve and
let $c:(-\varepsilon,\varepsilon)\to M$ be a
smooth curve with $c(0)=\gamma(a)$. Show that
there exists a variation
$f:(-\delta,\delta)\times[0,a]\to M$ of
$\gamma$ such that
$$
f(s,0)=\gamma(0),\qquad f(s,a)=c(s).
$$
Moreover, show that its variational field
$V(t)=f_s(0,t)$ may be chosen so that
$V(0)=0$ and $V(a)=c'(0)$.
\end{exercise}

\section{Second Variation Formula}
\label{differential-geometry-section-second-variation}

\begin{theorem}[Second Variation Formula]
\label{differential-geometry-theorem-second-variation-formula}
Let $f:(-\varepsilon,\varepsilon)\times[0,a]
\to M$ be a variation of a smooth
geodesic $\gamma$ with perhaps piecewise
smooth variation field $V$. Then
\begin{equation}
\label{differential-geometry-equation-second-variation-formula}
E''(0)=\int_0^a |V'|^2-\int_0^a \left\langle{}R_{\gamma'}V,V\right\rangle{}+
\left\langle{}\gamma'(t),\nabla_{\partial_s}V(t)
\right\rangle{}|_{0}^a+
\sum_{i}\left\langle{}V(t_i),V'(t^-_i)-V'(t^+_i)
\right\rangle{}
\end{equation}
\end{theorem}

\begin{proof}
Differentiating the First Variation
Formula we obtain
$$
\begin{aligned}
\frac{1}{2}E''(s)&
=\left(
\left\langle{}f_s,f_t\right\rangle{}|_{0}^a
-\int_0^a\left\langle{}f_s,f_{tt}\right\rangle{}\right)_s\\
&=(\left\langle{}f_s,f_t\right\rangle{}|_{0}^a)_s
-\int_0^a\left\langle{}f_{ss},f_{tt}\right\rangle{}
-\int_0^a\left\langle{}f_s,f_{tts}\right\rangle{}
\end{aligned}
$$
Now substitute
using the Symmetry Lemma for Curvature
\ref{differential-geometry-lemma-symmetry-lemma-curvature}:
$$
R(\gamma',V)\gamma'
=f_{t
ts}-f_{tst}
$$
to get
$$
\frac{1}{2}E''(s)=
(\left\langle{}f_s,f_t\right\rangle{}|_{0}^a)_s
-\int_0^a \left\langle{}f_{ss},f_{tt}\right\rangle{}
-\int_0^a\left\langle{}f_s,R(s,t)f_t\right\rangle{}
-\int_0^a\left\langle{}f_s,f_{t s t}\right\rangle{}
$$

Now use that the connection is metric
to make the term $\left\langle{}f_s,f_{tst}\right\rangle{}$ 
appear by differentiating 
$\left\langle{}f_s,f_{t s}\right\rangle{}$
with respect to $t$:
$$
\left\langle{}f_s,f_{ts}\right\rangle{}_t
=\left\langle{}f_{st},f_{ts}\right\rangle{}
+ \left\langle{}f_s,f_{t s t}\right\rangle{}.
$$
We obtain
$$
\frac{1}{2}E''(s)=
(\left\langle{}f_s,f_t\right\rangle{}|_{0}^a)_s
-\int_0^a \left\langle{}f_{ss},f_{tt}\right\rangle{}
-\int_0^a\left\langle{}f_s,R(s,t)f_t\right\rangle{}
+\int_0^a\left\langle{}f_{s t},f_{t s}\right\rangle{}
-\left\langle{}f_s,f_{t s}\right\rangle{}|_0^a.
$$
Now note that $\left\langle{}f_s,f_{t s}\right\rangle{}|_0^a$
cancels out with the very first term,
which is
$$
(\left\langle{}f_s,f_t\right\rangle{}|_0^a)_s
=\left\langle{}f_{s s},f_t\right\rangle{}|_0^a
+\left\langle{}f_s,f_{t s}\right\rangle{}|_0^a,
$$
giving
$$
\frac{1}{2}E''(s)=
\left\langle{}f_{ss},f_t\right\rangle{}|_{0}^a
-\int_0^a \left\langle{}f_{ss},f_{tt}\right\rangle{}
-\int_0^a\left\langle{}f_s,R(s,t)f_t\right\rangle{}
+\int_0^a\left\langle{}f_{s t},f_{t s}\right\rangle{}.
$$
Since $\gamma$ is a geodesic the
$\left\langle{}f_{ss},f_{tt}\right\rangle{}$ vanishes
when we evaluate at $s=0$,
and thus we obtain the desired formula
in case $\gamma$
is smooth.

As before, for piecewise smooth $\gamma$
we obtain a sum at the integration
by parts (i.e. using metric connection)
step.
\end{proof}

\begin{remark}
\label{differential-geometry-remark-second-variation-non-geodesic}
If $\gamma$ is not a geodesic we may
use again that the connection
is metric to deal with the
$\left\langle{}f_{ss},f_{tt}\right\rangle{}$ term:
$$
\left\langle{}f_{ss},f_t\right\rangle{}_t
=\left\langle{}f_{ss t},f_{t}\right\rangle{}
+\left\langle{}f_{ss},f_{tt}\right\rangle{}
$$
gives
$$
\begin{aligned}
\frac{1}{2}E''(s)&=
\left\langle{}f_{ss},f_t\right\rangle{}|_{0}^a
-\left\langle{}f_{ss},f_t\right\rangle{}|_0^a
+\int_0^a \left\langle{}f_{ss t},f_t\right\rangle{}\\
&-\int_0^a\left\langle{}f_s,R(s,t)f_t\right\rangle{}
+\int_0^a\left\langle{}f_{s t},f_{t s}\right\rangle{}.
\end{aligned}
$$
So, a second-variation-like formula 
for non geodesic curves is
$$
\frac{1}{2}E''(s)=
\int_0^a \left\langle{}f_{ss t},f_t\right\rangle{}
-\int_0^a\left\langle{}f_s,R(s,t)f_t\right\rangle{}
+\int_0^a\left\langle{}f_{s t},f_{t s}\right\rangle{}.
$$
\end{remark}

\begin{definition}
\label{differential-geometry-definition-index-form}
The {\it index form} is
\begin{equation}
\label{differential-geometry-equation-index-form}
I_a(V,W):=\int_0^a
\left\langle{}V',W'\right\rangle{}-\int_0^a\left\langle{}R_{\gamma'}V,
W\right\rangle{}
\end{equation}
for any $V,W \in \mathfrak{X}_\gamma$.
\end{definition}

\begin{remark}
\label{differential-geometry-remark-index-form-alternative-formula}
Using that the connection is metric,
$$
\begin{aligned}
\left\langle{}V,W'\right\rangle{}'&=\left\langle{}V',W'\right\rangle{}+
\left\langle{}V,W^{\prime\prime}\right\rangle{}\\
\implies
\left\langle{}V',W'\right\rangle{}&=\left\langle{}V,W'\right\rangle{}'-
\left\langle{}V,W^{\prime\prime}\right\rangle{}
\end{aligned}
$$
and substituting in the definition we obtain
that
\begin{equation}
\label{differential-geometry-equation-index-form-alter-formu}
\begin{aligned}
I_a(V,W)&=\int_0^a\left\langle{}V,W'\right\rangle{}'-\int_0^a
\left\langle{}V,W^{\prime\prime}\right\rangle{}-\int_0^a
\left\langle{}R_{\gamma'}V,W\right\rangle{}\\
&=\left\langle{}V,W'\right\rangle{}|_{0}^a-\int_0^a
\left\langle{}V,W^{\prime\prime}\right\rangle{}-\int_0^a\left\langle{}R_{\gamma'}
W,V\right\rangle{}\\
&=\left\langle{}V,W'\right\rangle{}|_{0}^a-\int_0^a
\left\langle{}V,W^{\prime\prime}+R_{\gamma'}W\right\rangle{}
\end{aligned}
\end{equation}
\end{remark}

\begin{exercise}
\label{differential-geometry-exercise-two-submanifolds}
\cite{doc} Chapter IX, Exercise 5. Sejam
$N_1$ e $N_2$ duas subvariedades
 fechadas e disjuntas de uma variedade
 Riemanniana compacta.
\begin{enumerate}
\item Mostre que a distância entre $N_1$ e
$N_2$ é realizada por uma geodésica
	$\gamma$ perpendicular a ambas $N_1$ e
	$N_2$.
\item Mostre que, para qualquer variação
ortogonal $h(t,s)$ de $\gamma$, com
$h(0,s) \in N_1$ e $h(\ell,s) \in N_2$,
tem-se para a fórmula da segunda
variação a seguinte expressão
$$
\frac{1}{2}E''(0)=I_{\ell}(V,V)+
\left\langle{}V(\ell),A_{\gamma'(\ell)}^{(2)}V(\ell)
\right\rangle{}
-\left\langle{}V(0),A^{(1)}_{\gamma'(0)}(V(0))\right\rangle{}
$$
\end{enumerate}
\end{exercise}

\begin{proof}
\begin{enumerate}
\item Since both submanifolds are compact,
there exists a minimizing geodesic
$\gamma$ joining them. Then we can apply the
first variation formula. Consider a
variation that fixes $\gamma$ everywhere but
in a small neighbourhood of one of
the contact points. We realise by Eq.
\ref{differential-geometry-equation-first-variation-formula}
that it must be orthogonal. A similar
procedure proves the other endpoint also
arrives orthogonally.

\item By the previous item, we
know that $\gamma'$ is orthogonal to the
curves $h(0,s)$ and $h(\ell,s)$. Then
we can differentiate using the formula for
normal sections at the endpoints of
$\gamma$:
$$
\tilde{\nabla}_{V}\gamma'=A^{(i)}_{\gamma'}V+
\nabla^{\perp,i}_V\gamma'
$$
When computing
$V\left\langle{}\gamma',\gamma'\right\rangle{}$ we obtain
$\left\langle{}\nabla^{\perp,i}_V\gamma',
\gamma'\right\rangle{}=0$
since the shape operator is
tangent to the submanifolds. If $M_1$ and
$M_2$ are hypersurfaces we conclude
that the normal vector
$\nabla^{\perp,i}_V\gamma'$ vanishes since it
must be a
multiple of $\gamma'$. Otherwise I'm not sure
why this term would vanish.

Now, other than the Index form, in the Second
Variation Formula
\ref{differential-geometry-equation-second-variation-formula} we
have
$\left\langle{}\nabla_{\partial_s}V,\gamma'\right\rangle{}$.
Notice that
$$
\left\langle{}f_s,f_t\right\rangle{}_s=\left\langle{}f_{ss},
f_t\right\rangle{}+\left\langle{}f_s,f_{st}\right\rangle{}
$$
The left-hand-side will vanish because the
variation is orthogonal, giving the
equality we want (modulo a sign) since
$f_{st}$ is precisely
$\tilde{\nabla}_V\gamma'$ by the Symmetry
Lemma.
\end{enumerate}
\end{proof}


\section{Bonnet-Myers Theorem}
\label{differential-geometry-section-Bonnet-Myers-theorem}

\begin{theorem}[Bonnet-Myers]
\label{differential-geometry-theorem-Bonnet-Myers}
If $M$ is a complete manifold with
$\text{Ric} \geq \frac{1}{r^2}$, then
$\text{diam}M\leq \pi r$. In particular, $M$
is compact.
\end{theorem}

\begin{proof}[Sketch of proof]
Let $p,q \in M$ be points with $d(p,q)>\pi
r$. We will construct a Jacobi
field vanishing at the endpoints, so that
these points are conjugate and thus
any geodesic $\gamma$ joining them cannot be
minimizing.

Let $w_i$ be vectors in $T_pM$ which along
with $\gamma'(0)$ form an orthonormal
basis. Define $W_i$ to be the parallel
transport of $w_i$ along $\gamma$ and
$J_i:=\sin(\pi t/\ell)W_i$.

Computing the second variation formula, which
reduces to the index of $J_i$,
which we may write as in Eq.
\ref{differential-geometry-equation-index-form-alter-formu} to
obtain
$$
\begin{aligned}
\frac{1}{2}E''(0)&=-\int_0^\ell\left\langle{}J,J''-
R_{\gamma'}J\right\rangle{}
\approx\int_0^\ell \sin(\pi
t)(1+K(W_i,\gamma'(t))
\end{aligned}
$$
Summing over all indices $i$ we arrive at the
Ricci curvature, which must be
greater or equal than $1/r^2$ by hypothesis.
In comparing with the length of the
curve $\ell>\pi r$, we find that the
integrand must be
negative, which forces one of the indices to
be negative.
\end{proof}

\begin{lemma}
\label{differential-geometry-lemma-bonne-myers-impli-compa-unive}
\cite{doc}, Chapter IX, Corollary 3.2. Let
$M$ be a complete Riemannian manifold
with $\text{Ric}\geq \delta>0$. Then the
universal covering of $M$ is compact
and $M$ has finite fundamental group.
\end{lemma}

\begin{proof}

The universal cover with the pullback metric
satisfies the same curvature
hypothesis as $M$, making it compact by
Bonnet-Myers Theorem
\ref{differential-geometry-theorem-Bonnet-Myers}. To confirm that
the deck transformation group is
finite notice that if we had an infinite
orbit, we could find a limit point
whose projection to the base would not admit
a regular neighbourhood.
\end{proof}


\begin{exercise}[Exercício 3]
\label{differential-geometry-exercise-posit-scala-but-no-posit}
Encontre um exemplo de variedade suave que
admite alguma métrica Riemanniana com
curvatura escalar positiva, mas não admite
uma métrica Riemanniana com Ricci
positivo.
\end{exercise}

\begin{proof}
Considere $S^2\times S^1$. A métrica produto
tem curvatura escalar positiva,
pois é a soma das curvaturas escalares em
cada fator, onde a curvatura escalar
de $S^2$ é constate $1$ quanto que a
curvatura escalar de $S^1$ é constante
zero. Suponha que $S^2\times S^1$ admite uma
métrica com $\text{Ric}>0$. Como
$S^2\times S^1$ é compacta, o fibrado
unitário dela é compacto. Isso significa
que existe uma constante $\delta$ tal que
$\text{Ric}\geq \delta>0$. O
recobrimento universal de $S^2 \times S^1$
com a métrica de
recobrimento satisfaz as mesmas hipóteses de
curvatura, de modo que podemos
aplicar o teorema de Bonnet-Myers para
concluir que é compacto e que o grupo de
transformações de coberta, ou seja, o grupo
fundamental de $S^2\times S^1$ é
finito.
\end{proof}

\section{Weinstein's theorem}
\label{differential-geometry-section-Weinstein-theorem}

\begin{theorem}[Weinstein]
\label{differential-geometry-theorem-Weinstein}
Let $M^n$ be oriented, compact, $K>0$ and $f
\in \text{Iso}(M)$. Suppose that
if $n$ is even, $f$ preserves orientation,
and if $n$ is odd, $f$ reverts
orientation. Then $f$ has a fixed point.
\end{theorem}

\begin{proof}
       A fórmula da segunda variação é a
       fonte de inspiração, olhe:
$$
E''(0)=\int |\dot V|^2-\int\left\langle{}R_ {\dot
\gamma}V,V\right\rangle{}+\left\langle{}\dot\gamma,
\nabla_{\partial_s}f_s\right\rangle{}|_{0}^a
$$
e pense: o que preciso para tirar algo bacano
dessa formulinha?

De alguma maneira nos damos conta de que isso
tem a ver com os pontos fixos de
uma isometria $f$. Suponha que $f$ não tem
pontos fixos e busquemos uma
contradição. Defina
$$
\begin{aligned}
g: M &\longrightarrow \mathbb{R} \\
g(q) &=d(q,f(q))
\end{aligned}
$$

então como $M$ é compacta e essa função é
contínua, ela tem um ponto mínimo que
chamamos de $p$. Temos que:
$$
0<g(p)=\operatorname{min}g\leq g(q)\qquad
\forall q.
$$

Vamos usar a segunda fórmula da variação para
achar um ponto que avaliado em $g$
fica menor do que $g(p)$. Seja $\gamma$ a
geodésica minimizante entre $p$ e
$f(p)$ parametrizada por comprimento de arco
(que existe porque $M$ compacta
implica completa).

Agora pegue um vetor $v \in T_pM$ unitário e
ortogonal a $\dot \gamma$ (para que
o termo com $R$ na fórmula da seg. var. fique
$K$), transporte ele paralelamente
ao longo de $\gamma$ para obter $V$, e pegue
a variação por geodésicas
\[h(s,t)=\operatorname{exp}_{\gamma(t)}(sV(t)
)=\gamma_{V(t)}(s).\]

Então a 2nda fórmula fica beleza: como o
campo é paralelo, o termo $|\dot V|^2$
morre, o termo da curvatura fica $K$ e o
termo com $\nabla_{\partial_s}f_s$
morre também (porque as curvas verticais são
geodésicas).

Então como estamos em $K>0$, acaba que
$E''(0)<0$, então existe uma vizinhança
de $0$  onde as curvas $f(s,t)$ tem menor
energia que $f(0,t)$ (para toda $s$
nessa  vizinhança). Agora lembre que a
$\ell^2 \leq dE$, e que no caso de
geodésicas minimizantes se da a igualdade.
Então fica que
$$
\frac{1}{d}\ell^2(c^s)\leq
E(s)<E(0)=\frac{1}{d}\ell(\gamma)^2=\frac{1}
{d}d(p,f(p))^2
=\frac{1}{d}(\operatorname{min}g)^2
$$
lado esquerdo. Defina $\beta(s)$ como a curva
``vertical" no tempo 0, ie.
$\beta(s)=f(s,0)$. Queremos mostrar que
$$
\frac{1}{d}d(\beta(s),f(\beta(s))^2\leq
\frac{1}{d}\ell^2(c^s)
$$
Já que por definição, do lado esquerdo temos
$\frac{1}{d}g(\beta(s))^2$, dando
uma contradição com a equação anterior. Para
confirmar essa desigualdade e
concluir a prova só devemos mostrar que a
curva $c^s$ liga $\beta(s)$ e
$f(\beta(s))$, pois a distância entre esses
dois pontos é menor do que o
comprimento de qualquer curva que liga esses
pontos.

Ou seja, queremos que
$$
c^s(0)=\beta(s),\qquad
c^s(d)=f(\beta(s))\qquad \forall s
$$
Ou seja, queremos que
$$
f \circ \beta=\gamma_{V(d)}
$$
Note que $\beta(s)=\gamma_{V(0)}(s)$. Então
derive em $s=0$: \[f \circ
\beta=f\circ\gamma_{V(0)}=\gamma_{V(d)}\iff
f_{*,p}V(0)=V(d)=PV(0)\] onde $P$ é
o transporte paralelo ao longo de $\gamma$ de
$p=\gamma(0)$ a $\gamma(d)$.  Isso
implica que queremos que
\[P^{-1}f_{*,p}V(0)=V(0)\]

Com um desenho muito convincente (ver
Manfredo p. 225, esse agumentinho tá
tranquilo) fica que como $f$ é uma isometria
\[f_{*,p}\gamma'(0)=\gamma'(d)\]
Decorre daí que
$P^{-1}f_{*,p}\gamma'(0)=P^{-1}\gamma'(d)=
\gamma'(0)$.
Então
podemos fixar nossa atenção em
$$
A:=P^{-1}f_{*,p}|_{\gamma'(p)^\perp}\in
\mathsf{O}(n-1)
$$
Nosso objetivo é analisar o que acontece
quando $A$ tem um ponto fixo, $V(0)$,
pois essa é uma condição que queremos que
seja verdade para que $f\circ \beta =
\gamma_{V(d)}$.

Agora segue o argumento dos eigenvalores que
diz assim. Defina
$A:=P^{-1}f_{*,p}$. Suponha que $A$ tem um
ponto fixo. Os autovalores vem dados
assim:
$$
\operatorname{Spec}A=\{\underbrace{\lambda_1,
\overline{\lambda_1},\lambda_2,
\overline{\lambda_2}}_{\text{par}
},-1,\ldots,-1,1,\ldots,1\}
$$
Ou seja, os eigenvalores que não são 1 ou
$-1$ não afetam o determinante. Daí
concluímos que
\begin{itemize}
\item Se $n-1$ é ímpar, $f$ preserva
orientação. Por que? Deixa os complexos do
lado que não importam. Então pensa: se só
tenho $-1$s, então tenho uma
quantidade ímpar de $-1$s porque a dimensão é
ímpar. Respira. Nesse caso $f$
inverte orientação. Mas: $f$ deve ter pelo
menos um $1$! Então a quantidade de
$-1$s  é par. Então não: $f$ preserva
orientação.
\item se $n-1$ é par, $f$ reverte orientação.
\end{itemize}
E essas são as condições que precisamos para
que $A$ tenha um ponto fixo e tudo
funcione. Então bota isso aí na hipótese do
teorema.
\end{proof}

\section{Synge's theorem}
\label{differential-geometry-section-Synge-theorem}

\begin{theorem}[Synge]
\label{differential-geometry-theorem-Synge}
Let $M^n$ be compact and $K>0$. Then
\begin{enumerate}
\item If $n$ is even,
$$
\pi_1(M)=\begin{cases}
	1\qquad & \text{if $M$ is orientable}\\
	\mathbb{Z}_2\qquad & \text{if $M$ is not
	orientable}
\end{cases}
$$
\item If $n$ is odd, then $M^n$ is
orientable.
\end{enumerate}
\end{theorem}

\begin{proof}
First suppose $n$ is even and $M$ is
orientable. Since $M$ is compact we must
have that $K\geq \delta>0$. Now consider the
universal cover $\tilde{M}$. Then
it also satisfies the curvature bound (with
the pullback metric),
making it compact as well by Bonnet-Myers
\ref{differential-geometry-theorem-Bonnet-Myers}.
Now pick a deck transformation, which must be
orientation-preserving if we choose the
orientation on $\tilde{M}$ (simply
connected is always orientable) making the
projection orientation-preserving.
Then $f$ must have a fixed point by
Weinstein's theorem \ref{differential-geometry-theorem-Weinstein},
which applies since we have shown that
$\tilde{M}$ is compact.

However, deck transformations cannot have
fixed points unless they are the
identity: this is a consequence of the Unique
Lifting Property
 \ref{differential-geometry-proposition-unique-lifting-property},
 which is a
slightly more general statement than the
uniqueness of curve lifts that we
used for Klinenberg's lemma exercise
\ref{differential-geometry-exercise-l7-1}.
More exactly: given a covering space and two
maps from any space into the base,
these coincide if they agree at one point
only.

If $M$ is not orientable we can show that the
two-sheeted orientable cover of
$M$ is its universal covering by the same
arguments as above. Then the group of
deck transformations is the group of two
elements $\mathbb{Z}_2$ since the
covering is two-sheeted.

Now suppose $n$ is odd and consider the
two-sheeted orientable cover of $M$.
Then by Weinstein's theorem we can say that
no deck transformation can be
orientation reversing. But since $M$ is
non-orientable, the single deck
transformation generating the group of deck
transformations, which is
$\mathbb{Z}_2$, is not orientation
preserving: if it was, then the quotient of
the two-sheeted orientable covering by the
action of the deck transformations
would induce an orientation on $M$.
\end{proof}

\begin{exercise}
\label{differential-geometry-exercise-posit-ricci-but-no-posit}
Encontre um exemplo de variedade suave que
admite uma métrica Riemanniana com
Ricci positivo, mas não admite uma métrica
Riemanniana com curvatura seccional
positiva.
\end{exercise}

\begin{proof}
Considere $S^3\times S^1$. A métrica produto
tem $\text{Ric}>0$, já que ele é a
soma dos tensores de Ricci em cada fator de
acordo ao Exercício
\ref{differential-geometry-exercise-curvatures-in-product} (a
curvatura de Ricci de $S^3$ é 1,
enquanto que a curvatura escalar de $S^1$ é
zero). Agora suponha que existe uma
 métrica com $K>0$. Então pelo teorema de
 Synge \ref{differential-geometry-theorem-Synge}, como
$S^3\times S^1$ é compacta, de dimensão par e
orientável, o grupo fundamental
dela deveria ser trivial, mas não é o caso.
\end{proof}

\section{Index lemma}
\label{differential-geometry-section-index-lemma}

\begin{lemma}[Index]
\label{differential-geometry-lemma-index}
\cite{doc} Chapter X, Section 2, Lemma 2.2.
Let $\gamma:[0,a]\to M$ be a
geodesic without conjugate points (this means
$\gamma(t)$ is not conjugate to
$\gamma(0)$ on $(0,a]$). Let
$J\in\mathfrak{X}_\gamma^J$ with
$\left\langle{}J,\gamma'\right\rangle{}=0$ and let $V$ be a
piecewise differentiable vector
field along $\gamma$ with
$\left\langle{}V,\gamma'\right\rangle{}=0$. Suppose
$J(0)=0=V(0)$ and
that there is  $t_0\in(0,a]$ such that
$J(t_0)=V(t_0)$. Then
$$
I_{t_0}(J,J)\leq I_{t_0}(V,V)
$$
and equality holds if and only if $V=J$ on
$[0,t_0]$.
\end{lemma}

\section{Rauch's comparison theorem}
\label{differential-geometry-section-Rauch}

We motivate Rauch's comparison theorem with a
result from ordinary
differential equations.

\begin{theorem}[Sturm]
\label{differential-geometry-theorem-Sturm}
Suppose $f,\tilde{f},K,\tilde{K}:[0,a]\to
\mathbb{R}$ are smooth functions
satisfying the following conditions.
$\tilde{f} > 0$ in $(0,a]$,
$f(0)=\tilde{f}(0)=0$,
$f'(0)=\tilde{f}'(0)$ and
$$
f''+Kf=0,\qquad
\tilde{f}''+\tilde{K}\tilde{f}=0.
$$
If $K\leq \tilde{K}$, then $f/\tilde{f}$ is
non decreasing. Moreover, if there
is any point $t \in [0,a]$ where
$f(t)=\tilde{f}(t)$, then $f \equiv
\tilde{f}$
and $K \equiv \tilde{K}$.
\end{theorem}

\begin{proof}
Since we want to show $f/\tilde{f}$ is non
decreasing we are interested in
computing the sign of the derivative of this
quotient. We thus calculate
$f'\tilde{f}-f \tilde{f}'$. We only have
information about second derivatives,
so we differentiate to obtain
$$
(f'\tilde{f}-f
\tilde{f}')'=(\tilde{K}-K)f\tilde{f}
$$
We may integrate to find that
\begin{equation}
\label{differential-geometry-equation-Sturm}
f'\tilde{f}-f \tilde{f}'=\int_0^t
(f'\tilde{f}-f \tilde{f}')'
=\int_0^t (\tilde{K}-K)f\tilde{f}
\end{equation}
If we show that $f\geq 0$ we get our result.
So suppose that $t_0$ is the
smallest number $t_0 \in (0,a]$ where $f$ is
zero. Then $f'(t_0)<0$ since before
that $f$ is positive, which in turn is
because $f'(0)=\tilde{f}'(0)$ and
$\tilde{f}$ is strictly positive. But this
contradicts Equation
\ref{differential-geometry-equation-Sturm} since the integrand was
positive up to $t_0$.

Now suppose that $f(t_0)=\tilde{f}(t_0)$ for
some $t_0 \in (0,a]$. Then Equation
\ref{differential-geometry-equation-Sturm} vanishes at $t_0$, and
we conclude that $K \equiv
\tilde{K}$ in $[0,t_0]$. Somehow this implies
that $f$ also equals $\tilde{f}$…

Here's ChatGPT in action: define
$h(t):=f(t)-\tilde{f}(t)$. Notice that on
$(0,t_0]$ it satisfies $h''+Kh=0$ with
$h(0)=h'(0)=0$. But so does constant zero
function, hence  $h\equiv0$. To extend to all
of $[0,a]$ use uniqueness of
ODEs again.
\end{proof}

\begin{theorem}[Rauch]
\label{differential-geometry-theorem-rauch-comparison}
Let $\gamma:[0,a]\to M$ and
$\tilde{\gamma}:[0,a]\to \tilde M^{n+p}$
be geodesics, and let
$J\in\mathfrak{X}^J_\gamma$ and
$\tilde J\in\mathfrak{X}^J_{\tilde\gamma}$
be Jacobi fields with comparable initial
conditions, that is,
$$
|\gamma'|=|\tilde\gamma'|,
\qquad
J(0)=0,\qquad \tilde J(0)=0,
$$
$$
\left\langle{}J'(0),\gamma'(0)\right\rangle{}
=
\left\langle{}\tilde J'(0),
\tilde\gamma'(0)\right\rangle{},
$$
and
$$
|J'(0)|=|\tilde J'(0)|.
$$
Assume that $\tilde\gamma$ has no conjugate
points on $(0,a]$ and that, on $(0,a]$,
$$
K(\gamma',J)
\leq
\tilde K(\tilde\gamma',\cdot).
$$
Then
$$
\frac{|J|}{|\tilde J|}
$$
is nondecreasing. In particular,
$$
|J|\geq |\tilde J|.
$$
Moreover, if $|\tilde J(r)|=|J(r)|$ for
some $r\in(0,a]$, then
$$
K(\gamma',J)=\tilde K(\tilde\gamma',
\tilde J)
$$
on $(0,r]$.
\end{theorem}

\begin{remark}
\label{differential-geometry-remark-rauch-comparison-notation}
The notation $\tilde K(\tilde\gamma',\cdot)$
means the sectional curvature of every plane
containing $\tilde\gamma'$. Thus the
hypothesis says that the sectional curvature
of the plane generated by $\gamma'$ and $J$
is bounded above by the corresponding
sectional curvatures in $\tilde M$.
\end{remark}

\begin{remark}
\label{differential-geometry-remark-jacobi-orthogonality-initial}
For any Jacobi field $J$ along a geodesic
$\gamma$, the function
$$
t\mapsto \left\langle{}J(t),\gamma'(t)
\right\rangle{}
$$
has zero second derivative. Indeed,
$\gamma''=0$ and
$$
J''=-R(J,\gamma')\gamma',
$$
so the curvature symmetries give
$$
\left\langle{}J'',\gamma'\right\rangle{}=0.
$$
Thus
$$
\left\langle{}J,\gamma'\right\rangle{}=a+bt.
$$
If $J(0)=0$, then $a=0$, and if
$J'(0)\perp\gamma'(0)$, then $b=0$. Hence
$J\perp\gamma'$ for all $t$.
\end{remark}

\begin{remark}
\label{differential-geometry-remark-rauch-real-jacobi-equation}
The real Jacobi equation is
$$
J''+R(J,\gamma')\gamma'=0.
$$
The notation
$$
J''+KJ=0
$$
is only literally correct in dimension two,
in constant curvature, or when the curvature
operator preserves the line spanned by $J$.
In the clean Rauch theorem, $K_M(\gamma',J)$
means the sectional curvature of the plane
spanned by $\gamma'$ and $J$.
\end{remark}

\begin{theorem}[Rauch, Sturm-related version]
\label{differential-geometry-theorem-rauch}
Let $\gamma$ be a geodesic of $M$ and
$\tilde{\gamma}$ a geodesic of
$\tilde{M}$. Suppose that $J \in
\mathfrak{X}^J_\gamma$ and $\tilde{J}\in
\mathfrak{X}^J_{\tilde{\gamma}}$ satisfy the
following conditions. $J(0)=0$,
$\tilde{J}(0)=0$, $|J'(0)|=|\tilde{J}'(0)|$.
This follows from matching the initial
tangential components:
$$
\left\langle{}J'(0),\gamma'(0)\right\rangle{}
=
\left\langle{}\tilde{J}'(0),
\tilde{\gamma}'(0)\right\rangle{}.
$$
Also,
$$
J''+K J=0,\qquad
\tilde{J}''+\tilde{K}\tilde{J}=0
$$
If $\tilde{\gamma}$ has no conjugate points
and $\tilde{K}\geq K$, then
$|J|/|\tilde{J}|$ is nondecreasing. Moreover,
if $|J|$ and $|\tilde{J}|$
coincide in one point, then $|J|\equiv
|\tilde{J}|$ and $K \equiv \tilde{K}$.
\end{theorem}

\begin{proof}
Let $f:=|J|^2$ and
$\tilde{f}:=|\tilde{J}|^2$. I think the main
difference with
Sturm is that the condition on Jacobi fields
is not a condition on the
derivatives of the functions $f$ and
$\tilde{f}$ we just defined. Note that
$$
f'=2\left\langle{}J',J\right\rangle{},\qquad
J''=2(\left\langle{}J',J'\right\rangle{}+\left\langle{}J'',J\right\rangle{}).
$$
While we may substitute $-KJ$ on the second
bracket, we still have the first
one. Perhaps this is the reason to define
$U(t):=J(t)/|J(r)|$ and
$\tilde{U}(t):=\tilde{J}(t)/|\tilde{J}(r)|$.
But wait, what's that $r$? There's
a trick in this proof.

Notice that $U$ and $\tilde{U}$ are Jacobi
fields and $|U(r)|=|\tilde{U}(r)|$.
\end{proof}

\begin{proposition}
\label{differential-geometry-proposition-contraction}
Sejam $p \in M$, $\tilde{p} \in \tilde{M}$ e
$i:T_pM \to T_{\tilde{p}}\tilde{M}$
uma isometria. Então $f:=\text{exp}^{-1}_p
\circ i \circ
\text{exp}_{\tilde{p}}$ é uma contração
métrica, i.e. $|f_*w|\leq |w|$
para todo $w \in T_pM$.

Mas ainda, se $\text{exp}_p^{-1}$ está
definida numa bola totalmente
convexa, $f$ é uma contração métrica.
\end{proposition}

\begin{proof}
Write an abitrary $w\in T_p M$ as a Jacobi
variational field of the variation
$f(s,t)=\text{exp}_p(t(v+sw)$. The resulting
field on $\tilde{M}$ is a
Jacobi field with respect to $iw$ using that
$i$ is an isometry. The conditions
of Rauch theorem are satisfied and the
desired inequality is obtained.

For the metric result we just integrate along
a curve realising the distance and
compute.
\end{proof}

\section{Moore's theorem}
\label{differential-geometry-section-Moore}

\begin{theorem}[Moore]
\label{differential-geometry-theorem-Moore}
Let $M^n \subset \tilde{M}^{n+p}$ be a
compact submanifold of a Hadamard
manifold $\tilde{M}$. If $K\leq \tilde{K}$,
then $p\geq n$.
\end{theorem}

\section{Focal points}
\label{differential-geometry-section-focal-points}

\begin{definition}
\label{differential-geometry-definition-focal-point}
Let $M$ be a submanifold of $\tilde{M}$ and
$p$ a point not in $M$. We say $p$
is a {\it focal point} of $M$ if there is a
geodesic $\gamma$ orthogonal to $M$
and a Jacobi field
$J\in\mathfrak{X}^J_\gamma$ such that
$$
J'(0)\in T_{\gamma(0)}M,\qquad
J'(0)+A_{\gamma'(0)}J(0)\in
T_{\gamma(0)}^\perp M
$$
\end{definition}

\begin{proposition}
\label{differential-geometry-proposition-focal-point-singu-norma}
The focal points of $M$ are the singularities
of $\text{exp}|_{T^\perp M}.$
\end{proposition}

\section{Morse index theorem}
\label{differential-geometry-section-morse-index}

\begin{definition}[Index of index form]
\label{differential-geometry-definition-index-index-form}
$$
i(I_a):=\text{sup}\{\dim L \subset
\mathfrak{X}_\gamma:I_a |_{L \times L}<0\}
$$
\end{definition}

For the proof of Morse Index Theorem
\ref{differential-geometry-theorem-Morse-index} we will use the
following objects
\begin{equation}
\label{differential-geometry-equation-index-theorem-spaces}
\begin{aligned}
\mathcal{V}_a&=\{V \in
\mathfrak{X}_\gamma:\text{d.p.p., }V(0)=0,
V(a)=0\}\\
\mathcal{V}_t^+&:=\{V \in
\mathcal{V}_t:V(t_i)=0\}\\
\mathcal{V}_t^-& :=\{V \in
\mathcal{V}_t:V|_{[t_j,t_{j+1}]}\in
\mathfrak{X}_{\gamma|_{[t_j,t_{j+1}]}}^J\}
\end{aligned}
\end{equation}

\begin{proposition}
\label{differential-geometry-proposition-kernel-of-index-form}
\cite{doc}, Chapter X, Proposition 2.3. The
kernel of the index form i
$$
\Ker I_t:=\{V \in
\mathcal{V}_t:I_t(V,W)=0\forall
W\in\mathcal{V}_t\}=
V_t \cap \mathfrak{X}_\gamma^J
$$
that is, it is composed of smooth Jacobi
fields along $\gamma$.
\end{proposition}

\begin{proof}[Sketch of proof]
The point here is to show that if $V \in
\Ker$ then it is (1) a Jacobi
 field in each interval where $V$ is smooth,
 and (2) smooth in the whole
 interval.

We use the fact that $V$ is in the kernel to
show that (1) the norm of
$V-R_{\gamma'}V$ is zero, and (2) the left
and right derivatives of $V$
coincide. The first part involves multiplying
by a positive function $f$ that
vanishes at the singular points; this will
make the ``singular part" of the
index form, i.e. the sum, vanish, so that we
get only with the integral part,
which is the norm of the Jacobi equation,
which will vanish since we are in the
kernel. The second part I'm not sure---why do
we need to show that $V$ is
differentiable, or $C^2$ (since it's the
solution of an ODE).
\end{proof}

\begin{lemma}
\label{differential-geometry-lemma-index-form-degen-iff-conju}
\cite{doc}, Chapter X, Corollary 2.4. The
index form is degenerate if and only
if the initial and final points of $\gamma$
are conjugate. In this case, the
nullity, i.e. dimension of $\Ker I_a$
coincides with the nullity of $\gamma(a)$
as a conjugate point of $\gamma(p)$, that is,
the dimension of the space of
Jacobi fields vanishing at the endpoints.
\end{lemma}

\begin{proof}
Just by Proposition
\ref{differential-geometry-proposition-kernel-of-index-form} and by
the fact that
$\mathcal{V}_a$ is defined as the smooth
vector fields vanishing in the
endpoints.
\end{proof}

\begin{proposition}
\label{differential-geometry-proposition-V-is-direct-sum}
\cite{doc}, Chapter XI, Proposition 2.5.
$\mathcal{V}_a=\mathcal{V}^+ \oplus
\mathcal{V}^-$.
\end{proposition}

\begin{proof}
Assume that the partition
$0<t_1<\ldots<t_k=a$ of $[0,a]$ is such that
every
segment $\gamma|_{[t_{i-1},t_i]}$ is
contained in a totally geodesic
neighbourhood, so that no such segment can
have conjugate points.

Let $V \in \mathcal{V}_a$ and pick a field $J
\in \mathcal{V}^-$ defined as the
piecewise Jacobi field with pairwise boundary
conditions $J(t_j)=V(t_j)$. The
assumtion that $\gamma|_{[t_{i-1},t_i]}$ has
no conjugate points ensures the
uniqueness of such $J$. Then
$V-J\in\mathcal{V}^+$ and we obtain the
result.
\end{proof}

\begin{theorem}[Morse Index]
\label{differential-geometry-theorem-morse-index}
O índice $i(I_t)$ é finito e igual ao número
de pontos conjugados em $[0,t)$ ao
longo de $\gamma$.
\end{theorem}

\begin{proof}
\begin{enumerate}
\item Define
$$
\begin{aligned}
\mathcal{V}_a&=\{V \in
\mathfrak{X}_\gamma:\text{d.p.p. }V(0)=0,
V(a)=0\}\\
\mathcal{V}_t^+&:=\{V \in
\mathcal{V}_t:V(t_i)=0\}\\
\mathcal{V}_t^-& :=\{V \in
\mathcal{V}_t:V|_{[t_j,t_{j+1}]}\in
\mathfrak{X}_{\gamma|_{[t_j,t_{j+1}]}}^J\}
\end{aligned}
$$
\item Note que
$$
\mathcal{V}_t=\mathcal{V}_t^+ \oplus
\mathcal{V}_t^-
$$
\item Note que
$$
\mathcal{V}_t^-=\bigoplus_{j=1}^{i-1}
T_{\gamma(t_j)}M
$$
porque para qualquer escolha de vetores
 $v_1 \in T_{\gamma(t_1)}M,\ldots,v_{i-1}\in
 T_{\gamma(t_{i-1})}M$ existe
 um único campo $J \in \mathcal{V}^-$ tal que
 $J(t_j)=v_j$, por ser a solução
 da EDO com condições de contorno.
\item Usamos Proposition
\ref{differential-geometry-proposition-kernel-of-index-form}.
\item …
\end{enumerate}
\end{proof}



\begin{theorem}[Jacobi's theorem]
\label{differential-geometry-theorem-Jacobi-theorem}
Se $\gamma$ é minimizante, então não pode ter
pontos conjugados. More generally,
if $\gamma:I \to M$ is a geodesic and $
\gamma(a)$ is conjugate to $\gamma(0)$
through $\gamma$ then for every $\delta>0$,
$I_{a+\delta}\not \geq 0$.
\end{theorem}

\begin{proof}
Se $\gamma$ é minimizante, a segunda fórmula
da variação é positiva para todo
campo em $\mathcal{V}=\{V\in
\mathfrak{X}_\gamma:V \text{ d.p.p. },
V(a)=0,V(b)=0\}$. Ou seja, a forma do índice
é uma forma bilinear positiva, e
portanto tem signatura zero, e pelo Teorema
do Índice de Morse
\ref{differential-geometry-theorem-Morse-index} o número de pontos
conjugados é zero.
\end{proof}

It is possible to prove this theorem without
using Morse Index Theorem
\ref{differential-geometry-theorem-Morse-index} by constructing a
variation where the index form is
negative, see Exercise \ref{differential-geometry-exercise-l7-4}.

\section{Cut locus}
\label{differential-geometry-section-cut-locus}

\begin{definition}
\label{differential-geometry-definition-cut-point}
O {\it cut point} de  $p$ ao longo de
$\gamma$ é $\gamma(\rho(\gamma))$
onde
$$
\rho(\gamma)=\text{sup}\{t:\gamma|_{[0,t]}
\text{
é minimizante}\}
$$

\end{definition}

\begin{proposition}
\label{differential-geometry-proposition-cut-point-chara}
\cite{doc} Chapter XIII, Proposition 2.2. Let
$\gamma$ be a minimizing geodesic joining $p$
and $q$. Then $q$ is the cut
point of $p$ if and only if either of the
following conditions hold:
\begin{enumerate}
\item $q$ is the first conjugate point of $p$
along $\gamma$.
\item There is another, distinct, minimizing
geodesic $\sigma$ joining $p$ and
 $q$.
\end{enumerate}
\end{proposition}

\begin{proof}[Sketch of proof]
($\implies$) Since $q:=\gamma(r)$ is the cut
point of $p$, we know that for every point
$\gamma(r+\varepsilon)$ along $\gamma$ after
$q$ there must be some other geodesic
$\gamma_\varepsilon$ which realizes the
distance between $p$ and
$\gamma(r+\varepsilon)$

Taking the limit as $\varepsilon \to 0$ we
obtain another geodesic
(because we are taking limit inside $S^n
\subset T_pM$, so we get convergence;
and because ``things will behave well under
the limit")
which puts us in the second condition of the
theorem,
unless the limit geodesic is the same as
$\gamma$ setwise.

Now there's a critical observation:
the two geodesics $\gamma$ and
$\gamma_\varepsilon$ meet at the point
 $\gamma(r+\varepsilon)$, so you can write
 $\gamma_\varepsilon(r+\delta(\varepsilon))=
 \gamma(r+\varepsilon)$
 for some function $\delta(\varepsilon)$
 (that could be negative, doesn't
 matter).
 But since $\gamma_\varepsilon$ is the
 minimizing
 geodesic from $\gamma(0)$ to this point,
 then
 $|\delta(\varepsilon)|<\varepsilon$.

Right, the point is that defining
$$
u_\varepsilon:=(r+\varepsilon)v,\qquad
u'_\varepsilon:=r+\delta(\varepsilon))
v_\varepsilon
$$
we see that $\text{exp}_p$ is not injective,
since these two vectors
hit that point.

This being true for every $\varepsilon>0$,
there is no neighbourhood of $vr$
where $\text{exp}_p$ is injective---much less
a diffeomorphism.
So it cannot be a local diffeomorphism, and
we know that critical points of
$\text{exp}_p$ are conjugate points of $p$.
(Remember: if
$\text{exp}_p$ has a critical point, its
differential has kernel,
and the map $d_{tv}\text{exp}_pw$, for $w$ in
the kernel, gives
a Jacobi field. See Proposition
\ref{differential-geometry-proposition-critical-points-exp})
\end{proof}

\begin{exercise}
\label{differential-geometry-exercise-diametros}
Calcule o diâmetro de $S^2$, $\mathbb{T}^2$,
$\mathbb{R}P^{2}$.
\end{exercise}

\begin{proof}[Solution]
O diâmetro de $S^2$ pode ser calculado via o
teorema de Bonnet Myers
\ref{differential-geometry-theorem-bonnet-myers}: nenhuma geodésica
é minimizante depois de atingir
comprimento $\pi r$, e temos uma geodésica
que atinge esse comprimento: qualquer
uma!

O diâmetro de $\mathbb{T}^2$ é $1$. Isso é
por simples geometria
euclidiana: é o diâmetro do cubo! Por
definição, a métrica de  $\mathbb{T}^2$ é
a induzida pela projeção quociente.

O diâmetro de $\mathbb{R}P^{2}$ é $\pi/2$.
Qualquer geodésica que liga dois
pontos a distância $\pi/2$ é minimizante,
pois estamos na métrica esférica e
podemos pegar cartas esféricas onde dois
pontos a essa distância estão contidos.
Por outro lado, se tivéssemos dois pontos a
distância maior, a geodésica
esférica que percorre o ponto antípoda ao
inicial, chega no ponto final mais
rapidamente que inicial; portanto geodésicas
de comprimento maior que $\pi/2$
não são minimizantes.
\end{proof}

\begin{lemma}
\label{differential-geometry-lemma-cut-point-reflexive}
If $q$ is the cut point of $p$ through
$\gamma$, then $p$ is the cut point of
$p$ through $\gamma$.
\end{lemma}

\begin{lemma}
\label{differential-geometry-lemma-not-cut-locus-minimizing}

\end{lemma}

\begin{proposition}
The function
$$
\begin{aligned}
	\rho: T^1_p &\longrightarrow [0,+\infty] \\
	v &\longmapsto \text{cut point along
	}\gamma_v
\end{aligned}
$$
is continuous.
\end{proposition}

\begin{lemma}
\label{differential-geometry-lemma-cut-locus-closed}
Because it's the image of the unit sphere in
$T_pM$, which is closed.
\end{lemma}

\begin{lemma}
\label{differential-geometry-lemma-cut-locus-topology}

\end{lemma}

\begin{lemma}
\label{differential-geometry-lemma-outsi-cut-locus-exist-minim}
\cite{doc} Chapter XIII, Corollary 2.8. If $q
\in M\setminus C_m(p)$, there exists a
 unique minimizing geodesic joining
$p$ and $q$.
\end{lemma}

\begin{proof}
By Hopf-Rinow \ref{differential-geometry-theorem-Hopf-Rinow}, we
know that
there exists a minimizing geodesic joining
any two points. If there was more
than one such geodesic, by
Proposition \ref{differential-geometry-proposition-cut-point-chara}
we get that $q$ is the cut point of  $q$
along either of these geodesics,
a contradiction.
\end{proof}

The former Lemma
\ref{differential-geometry-lemma-outsi-cut-locus-exist-minim}
shows that the exponential map $\text{exp}_p$
 is injective outside the cut locus of $p$.
This motivates the following definition.

\begin{definition}
\label{differential-geometry-definition-injectivity-radius}
The {\it injectivity radius} of a manifold
$M$ is
$$
i(M)=\text{inf}\{d(p,C_m(p)):p \in M\}
$$
\end{definition}

\section{Bishop-Gromov Theorem}
\label{differential-geometry-section-bishop-gromov}

\begin{theorem}[Bishop-Gromov]
\label{differential-geometry-theorem-Bishop-Gromov}
Let $p \in M$ and $0\leq t \leq
i(p)=d(p,C_m(p)$. Denote $B_t(p)$ the ball of
raduis $t$ centered in $p$ and $B_{t,k}$ the
ball of radius $t$ in a space of
constant curvature $k$.

If $\text{Ric}\geq k$, then
$\text{Vol}(B_t(p))/\text{Vol}(B_{t,k})$ is a
non
increasing function.

Moreover, if there are numbers  $0<s<r\leq
i(p)$ such that
$$
\frac{\text{Vol}(B_s(p))}{\text{Vol}(B_{s,k}}
=
\frac{\text{Vol}(B_r(p))}{\text{Vol}(B_{r,k})
}
$$
then $B_r(p)=B_{r,k}$ isometrically.
\end{theorem}

\begin{proof}
Use the exponential map a parametrization of
the geodesic sphere with radius
$r$ and centre in $p\in M$. (This
parametrization works at least locally, but
why
have we said in lecture that ``it's a global
chart''?) Then
$$
\text{Vol}(S^n_r(p))=\int_{S^n_r(p)}
\text{Vol}_{S^n_r(p)}=
\int_{S^n_r(0)}\text{exp}_p^*\text{Vol}
_{S^n_r}=
\int_{S^n_r(0)}|\det d_{rv}
\text{exp}_p|\text{Vol}_{S^n_r(0)}
$$
This determinant can be given by a basis of
tangent vectors to $S_r(0)$, each of
which yields a Jacobi field via Proposition
\ref{differential-geometry-proposition-everyday-jacobi-field}.
These Jacobi fields are the columns
of the matrix $\mathbb{J}$, which is the
Jacobian matrix of the exponential we
are interested in computing.

Now we want to differentiate this with
respect to $r$ because we aim to apply
Sturm's theorem \ref{differential-geometry-theorem-Sturm}. By
Exercise
\ref{differential-geometry-exercise-derivative-of-determinant}, we
know that the derivative of the
determinant of a one-parameter family of
invertible matrices $\mathbb{J}(r)$ is
given by $\det
\mathbb{J}(r)\text{tr}(\mathbb{J}^{-1}(r)
\mathbb{J}'(r))$.

By the discussion in Section
\ref{differential-geometry-section-Riccati-equation}, we know that
$AJ=J'$ where $A$ is the shape operator with
respect to the unit normal
$\gamma'(r)$ and $J$ is a Jacobi field
defined by
\ref{differential-geometry-proposition-everyday-jacobi-field}.
This gives $A=J'J^{-1}$.

Thediagonal entries of the matrix
$\mathbb{J}$ are vectors of this kind, so
that
 its trace, which depends only on the
 diagonal by Exercise
\ref{differential-geometry-exercise-trace-indep-coord}, is
precisely the trace of
the shape operator $A$, which we have defined
as $H/(n-1)$ where
$H$ is called the mean curvature.

This data translates to the following
equation
\begin{equation}
\label{differential-geometry-equation-Bishop-Gromov-proof1}
\det\mathbb{J}'=\det\mathbb{J}H
\end{equation}
Now recall that any self-adjoint operator can
be expressed uniquely as a
multiple of the identity plus a symmetric
traceless matrix (ref?). This allows
to write
\begin{equation}
\label{differential-geometry-equation-Bishop-Gromov-proof2}
A=H\text{Id}+A_0
\end{equation}
with $A_0$ symmetric and traceless.
Differentiating Eq.
\ref{differential-geometry-equation-Bishop-Gromov-proof2} and using
that the shape operator satisfies
Riccati equation \ref{differential-geometry-equation-Riccati} we
shall obtain that
$$
H'+H^2+\mathcal{R}=0
$$
where
$\mathcal{R}:=\text{Ric}(\gamma')+|A_0|^2$.
Then
\begin{equation}
\label{differential-geometry-equation-Sturm-type1}
\det\mathbb{J}''+\mathcal{R}\det\mathbb{J}=0,
\det\mathbb{J}(0)=,\det\mathbb{J}'(0)=1
\end{equation}
\bigskip
Notice that in the case of constant curvature
$k$, Eq.
\ref{differential-geometry-equation-Bishop-Gromov-proof1} becomes
\begin{equation}
\label{differential-geometry-equation-Bishop-Gromov-proof3}
\det\overline{\mathbb{J}}'=
\det\overline{\mathbb{J}}k
\end{equation}
which in turn yields
\begin{equation}
\label{differential-geometry-equation-Sturm-type2}
\det\overline{\mathbb{J}}''+
k\det\overline{\mathbb{J}}=0
\end{equation}
Now we want to compare the volume of the
sphere in $M$ with the volume of a
sphere of the same raduis in the space of
constant curvature
$\mathbb{Q}_k^{n+1}$:
$$
\begin{aligned}
\frac{\text{Vol}(S_r^{n}(p))}{\text{Vol}
(S^{n}_{r,k})}&=
\frac{\int_{S^n(0)}\det
\mathbb{J}}{\int_{S^n(0)}
\det\overline{\mathbb{J}}}
=\frac{1}{\text{Vol}(S^n(0))}
\int_{S^n(0)}\frac{\det\mathbb{J}}
{\det\overline{\mathbb{J}}}
\end{aligned}
$$
By Sturm Theorem \ref{differential-geometry-theorem-Sturm}, which
we may apply since both
$\det\mathbb{J}$ and
$\det\overline{\mathbb{J}}$ satisfy Eqs.
 \ref{differential-geometry-equation-Sturm-type1} and
 \ref{differential-geometry-equation-Sturm-type2}, and moreover
they both vanish at $r=0$ and their
derivatives are $1$ at $r=0$ by Exercise
\ref{differential-geometry-exercise-Jacobi-tensor-derivative}, and
of course, since we are supposing
that $\text{Ric}M\geq k$, we conclude that
the integrand is a non {\bf
decreasing} function.
\begin{remark}[Pregunta]
\label{differential-geometry-remark-pregunta2}
Según Sturm, el cociente de las funciones en
cuestión, en este caso,
$\frac{\det\mathbb{J}}
{\det\overline{\mathbb{J}}}$
debería ser una función no
decreciente, pero el resultado que buscamos
es que sea no {\bf creciente}.
\end{remark}
\bigskip
Now we turn to computing the ratios of the
volumes of balls of radius $r$.
$$
\begin{aligned}
\frac{\text{Vol}(B_r(p))}{\text{Vol}(B_{r,k})
}&=
\frac{\int_{S^n(0)}\int_0^r\det\mathbb{J}(t)
dt\text{Vol}S^n(0)}
{\text{Vol}(S^n(0))\int_0^r
\det\overline{\mathbb{J}}(t)dt}\\
&=\frac{1}{\text{Vol}(S^n(0))}
\frac{\int_0^r}{}
\end{aligned}
$$
\begin{remark}[Preguntas]
\label{differential-geometry-remark-pregunta3}
I have two questions:
\begin{enumerate}
\item We have said in lecture that we may
apply Fubini's theorem by Gauss'
lemma. Why? By Gauss' Lemma the normal
vectors are orthogonal to the spheres,
the constructions related to Riccati equation
are valid (cf. Section
\ref{differential-geometry-section-Riccati-equation}). By why does
this allow to use Fubini?
\item We can pull out the volume in the
denominator because the volumes of
spheres in the space of constant curvature
somehow does not depend on the radius
 $r$. But how exactly? I think  $\det
 \overline{\mathbb{J}}(r)$ does depend
on $r$ (in the space of constant curvature),
see Eq.
\ref{differential-geometry-equation-Bishop-Gromov-proof3}.
\end{enumerate}
\end{remark}
Then the equation continues to
$$
\begin{aligned}
&=\frac{1}{\text{Vol}(S^n(0))}
\int_{S^n(0)}\frac{\int_0^r \det
\mathbb{J}(t)dt \text{Vol}S^n(0)}
{\int_0^r\det\overline{J}dt}\\
&=\frac{1}{\text{Vol}(S^n(0))}
\int_{S^n(0)}\frac{\int_0^r \frac{\det
\mathbb{J}(t)}
{\det\overline{\mathbb{J}}(t)}
\det\overline{\mathbb{J}}(t)dt
\text{Vol}S^n(0)}
{\int_0^r\det\overline{\mathbb{J}}dt}\\
&=\frac{1}{\text{Vol}(S^n(0))}
\int_{S^n(0)}\frac{\int_0^r \frac{\det
\mathbb{J}(t)}
{\det\overline{\mathbb{J}}(t)}d\mu}{\mu[0,r]}
\text{Vol}S^n(0)
\end{aligned}
$$
where we have introduced a measure
$\mu:=\det \overline{\mathbb{J}}dt$.

By our discussion before, the integrand is a
non decreasing function, almost as
required.

Finally, notice that by the rigidity part on
Sturm's Theorem
\ref{differential-geometry-theorem-Sturm}, the condition
$$
\frac{\text{Vol}(B_s(p))}{\text{Vol}(B_{s,k}}
=
\frac{\text{Vol}(B_r(p))}{\text{Vol}(B_{r,k})
}
$$
implies that
$\det\mathbb{J}=\det\overline{\mathbb{J}}$
and $k=\mathcal{R}$.
This makes $A_0=0$ from Eq.
\ref{differential-geometry-equation-Bishop-Gromov-proof2} so that
$A$ is
proportional to the identity. By some mixture
of Eqs.
\ref{differential-geometry-equation-jacob-field-const-curva}
and
\ref{differential-geometry-equation-jacob-field-deriv-shape}, this
means that the
Jacobi fields along $\gamma$ are of the form
Eq.
\ref{differential-geometry-equation-jacob-field-const-curva-2}. In
turn, this makes the map
of Proposition \ref{differential-geometry-proposition-contraction}
an isometry. {\bf Who is $i$ then?}
\end{proof}

\section{Cheng's Theorem}
\label{differential-geometry-section-Cheng-theorem}

Cheng's theorem says what happens in case of
equality in the conditions of
Bonnet-Myers Theorem
\ref{differential-geometry-theorem-Bonnet-Myers}.
\begin{theorem}[Cheng]
\label{differential-geometry-theorem-Cheng}
Let $M$ be a complete Riemannian manifold
with $\text{Ric}\geq 1$. If
$\text{diam}M=\pi$ then $M=S^n$
isometrically.
\end{theorem}

\begin{proof}
By Bonnet-Myers Theorem
\ref{differential-geometry-theorem-Bonnet-Myers}, $M$ is a compact
manifold
with diameter $\pi$. This means that $M$ is
the closure of any ball of radius
$\pi$. Then, by Bishop-Gromov Theorem
\ref{differential-geometry-theorem-Bishop-Gromov}, it's enough
to show the rigidity condition holds for
radius $\pi$. That is, we must show
that there is $s<\pi$ such that
$$
\frac{\text{Vol}(B_s(p))}{\text{Vol}(B_{s,1}}
=
\frac{\text{Vol}(B_\pi(p)}{\text{Vol}B_{\pi,
1}}
$$
Consider two points $p_1$ and $p_2$ at
distance $\pi$. Then the balls of radius
$\pi/2$ with center at these points cannot
intersect. This says that
$$
\text{Vol}M\geq B_{\pi/2}(p_1)+B_{\pi/2}(p_2)
$$
But since the volume of $M$ is exactly the
volume of $B_\pi(p)$ for any $p$, we
see that
$$
\frac{\text{Vol}(B_\pi(p))}{2}\leq
B_{\pi/2}(p_i),\qquad i=1,2
$$
\end{proof}

\section{Toponogov's theorem}
\label{differential-geometry-section-Toponogov}

\begin{definition}
\label{differential-geometry-definition-geodesic-triangle}
\cite{doc}, p. 285. A {\it geodesic triangle}
is the set formed by three
minimizing geodesics $\gamma_i$ so that
$\gamma_{i+1}(0)=\gamma_i(\ell(\gamma_i))$
cyclically. (Just to note that for
some authors geodesic triangles are
minimizing by definition. Looks like for
\cite{meyers} it's not the case.)
\end{definition}

First notice that the following result can be
proved using Rauch's theorem
\ref{differential-geometry-theorem-rauch-comparison}:

\begin{theorem}[Toponogov, local version]
\label{differential-geometry-theorem-Toponogov-local}
Let $o,p_1,p_2 \in B$ where $B$ is a totally
convex ball of $M$. Let $\gamma_1$
and $\gamma_2$ be the normalized geodesics
joining $o$ to $p_1$ and to $p_2$.

Suppose $\tilde{o}, \tilde{p}_1,\tilde{p}_2
\in \tilde{B}$ is a triple on
another manifold $\tilde{M}$ such that the
distances from $\tilde{o}$ to
$\tilde{p}_1$ and from $\tilde{o}$ to
$\tilde{p}_2$ coincide with those in $M$.

Then for any $t \in [0,\ell(\gamma_1)]$ and
$s \in [0,\ell(\sigma)]$,
$$
\tilde{d}(\tilde{\gamma}_1(t),\tilde{\gamma}
_2(s))\leq
d(\gamma_1(t),\gamma_2(s))
$$
\end{theorem}

\begin{proof}
This is an application of Proposition
\ref{differential-geometry-proposition-contraction} for $i$ as
given by the condition of mapping
$\gamma'_1(0)$ to $\gamma_2'(0)$ and the
angle
between them to the corresponding vectors and
angle on $\tilde{M}$.
\end{proof}

The global version of that Theorem
\ref{differential-geometry-theorem-Toponogov-local} is the
following. Notice it coincides with our
intuition from comparing plane triangles
with spherical ones; spherical are fatter.

\begin{theorem}[Toponogov, hinge version]
\label{differential-geometry-theorem-Toponogov-hinge}
$M$ complete, $K \geq k$. $\gamma_1,\gamma_2$
normalized geodesics with
$\gamma_1(0)=\gamma_2(0)$. Suppose $\gamma_1$
is minimizing.

If $k>0$ suppose additionally that
$\ell(\gamma_2) \leq \pi/\sqrt{k}$ (because
that's where conjugate points appear).

Suppose that
$\tilde{\gamma}_1,\tilde{\gamma}_2$ is a
corresponding hinge in
$\mathbb{Q}_k^2$, that is,
$\ell(\gamma_i)=\ell(\tilde{\gamma}_i)$ and
$\angle(\gamma_1'(0),\gamma_2'(0)=
\angle\tilde{\gamma}_1'(0),\tilde{\gamma}
_2'(0)$.
Then
\begin{equation}
\label{differential-geometry-equation-Toponogov-hinge1}
d(\gamma_1(\ell_1),\gamma_2(\ell_2)) \leq
\tilde{d}(\tilde{\gamma}_1(\ell_1),
\tilde{\gamma}_2(\ell_2))
\end{equation}
\end{theorem}

This is not intuitive for me: the distance
between the remaining side should be
fatter in the manifold that is more curved.

The following version is equivalent to the
hinge version and is intuitive for
me:

\begin{theorem}[Toponogov]
\label{differential-geometry-theorem-Toponogov-hinge2}
Let $M$ be complete with $K\geq k$. If
$\gamma_i$ is a minimizing geodesic
triangle in $M$, then there is a unique
minimizing geodesic triangle
$\tilde{\gamma}_i$ in $\mathbb{Q}_k^2$ with
$\ell(\tilde{\gamma}_i)=\ell(\gamma_i)$ for
$i=0,1,2$, and such that
\begin{equation}
\label{differential-geometry-equation-Toponogov-hinge2}
\tilde{d}(\tilde{\gamma}_1(t),\tilde{\gamma}
(s))\leq
d(\gamma_1(t),\gamma_2(s)
\end{equation}
for all $t\in[0,\ell(\gamma_1)]$ and
$s\in[0,\ell(\gamma_2)]$.
\end{theorem}

Another equivalent version:
\begin{theorem}[Toponogov]
\label{differential-geometry-theorem-Toponogov-hinge3}
Let $M$ be complete with $K\geq k$. If
$\gamma_i$ is a minimizing geodesic
triangle in $M$, then there is a unique
minimizing geodesic triangle
$\tilde{\gamma}_i$ in $\mathbb{Q}_k^2$ with
$\ell(\tilde{\gamma}_i)=\ell(\gamma_i)$ for
$i=0,1,2$ and such that
\begin{equation}
\label{differential-geometry-equation-Toponogov-hinge3}
\tilde{d}(\tilde{o},\tilde{\gamma}(t))\leq
d(o,\gamma_0(t))
\qquad \forall t\in[0,\ell(\gamma_0)]
\end{equation}
\end{theorem}
All these versions are equivalent and we will
only prove the following one:
\begin{theorem}[Toponogov, metric version]
\label{differential-geometry-theorem-Toponogov-metric}
Let $M$ be a complete manifold and $K \geq
k$. Let $o,p_1,p_2 \in M$, distinct
points. Let $\gamma_1$ and $\gamma_2$ be
normalized geodesics joining
$o$ to $p_1$ and to $p_2$. Suppose $\gamma_1$
is minimizing (in the prove we
show that $\gamma_2$ is minimizing too).
Let $\gamma_0$ be a nonconstant geodesic
joining $p_1$ to $p_2$ and suppose
$$
\ell(\gamma_0) \leq \ell(\gamma_1)
+\ell(\gamma_2)
$$
(which is a weaker condition than asking that
$\gamma_0$ is minimizing as well).

If $k>0$ suppose additionally that
$\ell(\gamma_0) \leq \pi/\sqrt{k}$.

Then there exists a geodesic triangle
$\{\tilde{\gamma}_i:i=1,2,0\}$ in
$\mathbb{Q}_k^2$ such that
$\ell(\tilde{\gamma}_i) = \ell(\gamma_i)$ and
\begin{equation}
\label{differential-geometry-equation-Toponogov-metric}
\tilde{d}(\tilde{o},\tilde{\gamma}_0(t))
\leq  d(o,\gamma_0(t))
\qquad \forall t \in[0,\ell(\gamma_0)]
\end{equation}
\end{theorem}

\begin{proof}[Outline of proof]
This is the one we proved in class. Define
distance functions
$\rho:=d(\sigma,\cdot)$ and
$\tilde{\rho}:=\tilde{d}(\tilde{\sigma},
\cdot)$.
\end{proof}

\begin{lemma}
\label{differential-geometry-lemma-equality-Toponogov}
If equality holds in Eq.
\ref{differential-geometry-equation-Toponogov-hinge1}, then
$\gamma_1$ and
$\gamma_2$ are the boundary of a totally
geodesic surface $\Delta \subset M$
with $K_\Delta \equiv k$.
\end{lemma}

\begin{proof}
Using the map from Cartan's lemma. Use
Rauch's theorem \ref{differential-geometry-theorem-rauch-comparison} to
compare $M$ with $\mathbb{Q}_k$, and again to
compare $\Delta$ with $M$. As you
might expect, use Gauss' equation
\ref{differential-geometry-equation-Gauss} to show that the second
fundamental form of $\Delta$ vanishes. The
essence of the proof is that
$\alpha(\gamma',\gamma')=0$ for the geodesics
emanating from $p$, which is more
or less straightforward, and then showing
that also $\alpha(\sigma',\sigma')=0$
for geodesics emanating from another vertex
of the triangle, which requires to
``look at things from another point of
view''.
\end{proof}

\begin{theorem}[Toponogov, angle version]
\label{differential-geometry-theorem-Toponogov-angle}
\cite{Cheeger-Ebin} Theorem 2.2.(A). Let $M$
be a complete manifold with
$K_M\geq H$.

Let $(\gamma_1,\gamma_2,\gamma_3)$ determine
a geodesic triangle in $M$. Suppose
$\gamma_1,$ $\gamma_3$ are minimal and if
$H>0$, suppose $\ell(\gamma_2)\leq
\frac{\pi}{\sqrt{H}}$. Then $M^H$, the simply
connected 2-dimensional space of
constant curvature $H$, there exists a
geodesic triangle
$(\tilde{\gamma}_1,\tilde{\gamma}_2,
\tilde{\gamma}_3)$
such that
$\ell(\gamma_i)=\ell(\tilde{\gamma}_i)$ and
$\tilde{\alpha}_1 \leq\alpha_1$,
$\tilde{\alpha}_3 \leq \tilde{\alpha}_3$.

Except in case $H>0$ and
$\ell(\gamma_i)=\frac{\pi}{\sqrt{H}}$ for
some $i$, the
triangle in $M^H$ is uniquely determined.
\end{theorem}

\section{Gromov's theorem}
\label{differential-geometry-section-Gromov-theorem}

\begin{theorem}[Gromov]
\label{differential-geometry-theorem-Gromov}
$M^n$ complete, $K \geq 0$, then $\pi$ admits
a set of $3^n$ generators. In
particular, it is finitely generated.
\end{theorem}

\begin{proof}
Consider the universal cover $\tilde{M}$ of
$M$ with the pullback metric. Fix a
point $x \in \tilde{M}$ and consider the set
$$
\{f \in \Gamma:d(x,f(x))<r\},\qquad r>0
$$
By \ref{differential-geometry-proposition-deck-trans-act-prope}
every point must have a disjoint
neighbourhood …
\end{proof}

\section{Cartan's Theorem}
\label{differential-geometry-section-Cartan-theorem}

\begin{definition}
\label{differential-geometry-definition-free-homotopy-set}
The {\it free homotopy set} is the set of
homotopy classes of loops based on any
point of $M$.
\end{definition}

\begin{definition}
\label{differential-geometry-definition-geodesic-loop}
A {\it geodesic loop} is a geodesic whose
starting and ending points are the
same. It may not be differentiable at such
point.
\end{definition}

\begin{definition}
\label{differential-geometry-definition-closed-geodesic}
A {\it closed geodesic} a smooth geodesic
loop.
\end{definition}

\begin{theorem}[Cartan]
\label{differential-geometry-theorem-Cartan}
Let $M^n$ be compact, and $\omega \in
\hat{\pi}_1M$ nontrivial. Then there is a
closed geodesic $\gamma$ such that
$[\gamma]=\omega$.
\end{theorem}

\begin{proof}
Consider a sequence of loops whose lengths
converge to the number
$$
\ell:=\text{inf}\{\ell(c):c \in \omega\}>0
$$
Consider a sequence $c_n$ of curves in
$\omega$ such that their lengths converge
to $\ell$. Then apply Arzelà-Ascoli Theorem
\ref{real-analysis-theorem-Arzela-Ascoli}
 to ensure they converge to some curve $c$
 and show this is a geodesic
in $\omega$.

First suppose that each of the $c_n$ is a
piecewise geodesic parametrized by
arclength (why can we do this?). We may then
pick the supremum $L$ of all the
lengths since we are assuming that their
lengths decrease. Then
$$
d(c_n(t),c_n(s)) \leq \ell
\left(c_n|_{[t,s]}\right) \leq  N |t-s|
$$
This says that $\{c_n\}$ is equicontinuous
(the same $N$ works for all $n$) as a
family of maps from
\end{proof}

\section{Preissman's Theorem}
\label{differential-geometry-section-Preissman-theorem}

\begin{definition}
\label{differential-geometry-definition-translation}
An isometry $f \in \text{Iso}(M)$ is called a
{\it translation} if there exists
a geodesic of $M$ such that
$f(\gamma)=\gamma$. That is, $f \circ \gamma$
is a
reparametrization of $\gamma$.
\end{definition}

\begin{theorem}[Preissman]
\label{differential-geometry-theorem-Preissman}
Let $M^n$ be compact with $K<0$. If $1 \neq
H \subset \pi_1(M)$ is abelian,
then $H=\mathbb{Z}$.
\end{theorem}

\begin{theorem}[Preissman]
\label{differential-geometry-theorem-m-compa-k-negat-then-pi1}
If $M$ is compact and $K<0$, then
$\pi_{1}(M)$ is not abelian.
\end{theorem}

\section{Byers' Theorem}
\label{differential-geometry-section-Byers-theorem}

\begin{lemma}
\label{differential-geometry-lemma-deck-fixes-geode-then-not}
Let $M^n$ be a complete manifold with $K<0$.
Suppose that every element of
$\text{Deck}(\pi)$, where $\pi$ is the
universal covering, fixes the same
 geodesic $\tilde{\gamma}$. Then $M$ is not
 compact.
\end{lemma}

\begin{theorem}[Byers]
\label{differential-geometry-theorem-Byers}
Let $M$ be compact, $K<0$ and $H$ soluble
subgroup of $\pi_1(M)$, $H\neq \{e\}$.
Then $H$ is cyclic infinite. Moreover,
$\pi_1(M)$ does not have any cyclic
subgroup of finite index.
\end{theorem}

\section{Cheeger-Gromoll Siplitting Theorem}
\label{differential-geometry-section-chegg-gromo-split-theor}

\begin{theorem}[The Strong Maximum Principle]
\label{differential-geometry-theorem-strong-maximum-principle}
\cite{pet} Theorem 7.1.7. If $f:(M,g) \to
\mathbb{R}$ is continuous and
subharmonic, then $f$ is constant in a
neighbourhood of every local maximum. In
particular, if $f$ has a global maximum, then
$f$ is constant.
\end{theorem}

\begin{proof}
If $\Delta f>0$ then we can find a support
function. Notice that Prof. Florit
shows the theorem for functions satisfying
$\Delta f \leq 0$.
\end{proof}

The following lemma follows from the lemma
above and gives a bound on the
Laplacian of the distance function.

\begin{lemma}[Calabi]
\label{differential-geometry-lemma-Calabi}
If $\text{Ric} \geq 0$, for
$\rho:=d(p,\cdot)$, it holds that
$\Delta\rho \leq (n-1)/\rho$ on $M\setminus
C_m(p)\cup\{p\}$.
\end{lemma}


\begin{theorem}[Cheeger-Gromoll]
\label{differential-geometry-theorem-splitting}
Let $M$ be complete with  $\text{Ric}\geq0$.
If $M$ has a line, then $M$ is
isometric to $N\times\mathbb{R}$ for some
Riemannian manifold $N$.
\end{theorem}

\section{Exercises: basic constructions}
\label{differential-geometry-section-exercises-constructions}

\begin{exercise}
\label{differential-geometry-exercise-trace-indep-coord}
Show that the trace is independent of
coordinates.
\end{exercise}

\begin{proof}[Solution]
Let $A$ be an endomorphism of a vector space.
By definition
$$
\text{tr}A=A_i^i
$$
Now if $P$ is any invertible matrix,
$$
\text{tr}PAP^{-1}=P^j_iA_j^k\overline{P}_k^i=
P^j_i\overline{P}_k^iA_j^k
=\delta_k^jA_j^k=A_j^j
$$
where $\overline{P}$ denotes the indices of
the inverse of $P$.
\end{proof}

The following two lemmas are preparatory
results for Exercise
\ref{differential-geometry-exercise-derivative-of-determinant}.

\begin{lemma}
\label{differential-geometry-lemma-for-derivative-of-determinant1}
Let
$\gamma_1,...,\gamma_m:(a,b)\to\mathbb{R}^n$
be smooth curves and define
$f:(a,b)\to\mathbb{R}$ by
$f(t)=\det(\gamma_1(t),\ldots,\gamma_n(t))$.
Then
$$
f'(t_0)
=\sum_{k=1}^{m}\det(\gamma_1(t_0),\ldots,
\gamma'_k(t_0),\ldots,\gamma_m(t_0)).
$$
\end{lemma}

\begin{proof}
This follows from multilinearity of the
determinant using the definition of
derivative of a map $g:\mathbb{R}^n \to
\mathbb{R}^m$ as given by
$$
g(x+h)-g(x)=L(x)h+\varepsilon(x;h)
$$
where $\varepsilon$ is $o(h)$, meaning that
$\lim_{h \to 0}
\frac{\varepsilon(x;h)}{h}=0$.

More exactly, since $\gamma_k$ are smooth we
may write
$$
\gamma(t_0+h)=\gamma_k(t_0)+h\gamma'_k(t_0)+
\varepsilon_k(h)
$$
We substitute into the determinant…
\end{proof}

\begin{lemma}
\label{differential-geometry-lemma-for-derivative-of-determinant2}
Let $M \in \text{Mat}_{n\times
n}(\mathbb{R})$ be any matrix. Let
$f:\mathbb{T}\to \mathbb{R}$ be defined by
$f(t):=\det(\text{Id}+tM)$. Then
$f'(0)=\text{tr}(M)$.
\end{lemma}

\begin{proof}
This follows easily from Lemma
\ref{differential-geometry-lemma-for-derivative-of-determinant1}.
\end{proof}

\begin{exercise}
\label{differential-geometry-exercise-derivative-of-determinant}
Let $A(t)$ be a one parameter family of
matrices. Compute $\det A(t)'$.
\end{exercise}

\begin{proof}[Naive solution]
The determinant of a matrix is given by the
product of its eigenvalues
$\lambda_i$. Taking logarithms we see that
$$
\text{log}\det
A(t)=\text{log}\prod_{i}\lambda_i=\sum_{i}
\text{log}\lambda_i
$$
and differentiating we obtain
$$
\frac{\det A(t)'}{\det A(t)}=\sum_i
\frac{\lambda_i'}{\lambda_i}
=\text{tr} A'(t)A(t)^{-1}.
$$
with the caveat that we must confirm whether
the functions $\lambda_i$ are
differentiable. (Still pending…)
\end{proof}

\begin{proof}[Solution]
Consider the determinant as a function
$\det:\text{GL}(n,\mathbb{R})\to\mathbb{R}$.
Is total derivative at
$A \in \text{GL}(n,\mathbb{R})$ is a linear
map
$\text{GL}(n,\mathbb{R})\to \mathbb{R}$ which
when evaluated at a matrix
$H\in\text{GL}(n\mathbb{R})$ is defined by
$$
\det(A)'H=\lim_{t\to0}\frac{\det(A+tH)-
\det(A)}{t}
$$
Now define $f(t):=\det(A+tH)$ so that the
latter equation reads $f'(0)$. Then
observe that
$$
f(t)=\det(A+tH)=\det(A(\text{Id}+tA^{-1}H)=
\det(A)\det(\text{Id}+tA^{-1}H).
$$
By Lemma
\ref{differential-geometry-lemma-for-derivative-of-determinant2} we
conclude that
$f(t)=\det(A)\text{tr}(A^{-1}H)$. Letting
$H=A(t)'$ we obtain the result since
by the chain rule
$\det(A(t))'=\det(A(t))'A(t)'$.
\end{proof}

\begin{exercise}
\label{differential-geometry-exercise-levi-civit-conne-produ}
Let $P$ and $Q$ be Riemannian manidols. Sow
that the Levi-Civita connection of
the product metric is given by
$\nabla_{Y_1+Y_2}(X_1+X_2)=\nabla^P_{Y_1}X_1+
\nabla^Q_{Y_2}X_2$.
\end{exercise}

\begin{exercise}
\label{differential-geometry-exercise-total-geode-subma-produ}
Let $P$ and $Q$ be Riemannian manifolds. Show
that for any $p \in P$ and $q \in
Q$, $P\times \{q\}$ and $\{q\}\times P$ are
totally geodesic
submanifolds of $P \times Q$.
\end{exercise}

\begin{proof}[Solution]
Let $v \in T_{(p,q)}(P\times\{q_0\})$ and
consider the geodesic
 $\text{exp}_p^{P\times Q}(tv)$. Then
its acceleration vanishes since the
Levi-Civita connection of the product metric
is given by Exercise
\ref{differential-geometry-exercise-levi-civit-conne-produ}, and
the
component along $Q$ vanishes.
\end{proof}

\section{Exercises: global differential
geometry}
\label{differential-geometry-section-exerc-globa-diffe-geome}

\begin{exercise}
\label{differential-geometry-exercise-Klingenberg-lemma}
\cite{doc}, Chapter X, Exercise 1. Seja $M$
uma variedade Riemanniana
completa com curvatura seccional $K \leq
K_0$ onde $K_0$ é uma constante
positiva. Sejam $p,q \in M$ e seja $\gamma_0$
e $\gamma_1$ duas geodésicas
distinas unindo $p$ a $q$ com
$\ell(\gamma_0) \leq  \ell(\gamma_1)$.
Admita que $\gamma_0$ é homotópica a
$\gamma_1$, isto é, existe uma família
contínua de curvas $\alpha_t$, $t \in [0,1]$
tal que $\alpha_0 = \gamma_0$
e $\alpha_1=\gamma_1$. Prove que existe $t_0
\in [0,1]$ tal que
$$
\ell(\gamma_0)+\ell(\alpha_{t_0}) \geq
\frac{2\pi}{\sqrt{K_0} }
$$
\end{exercise}

\begin{exercise}
\label{differential-geometry-exercise-Hadamard-via-Kingenberg}
Use Klingenberg lemma from Exercise
\ref{differential-geometry-exercise-Klingenberg-lemma} to prove
Hadamard's theorem.
\end{exercise}

\begin{proof}
Suponha que $M$ é uma variedade completa,
simplesmente conexa e com curvatura
positiva. Por completitude, sabemos que para
qualquer $p \in M$ o domínio de
$\text{exp}_p$ é todo $T_pM$. Como $K\leq 0$,
sabemos que $M$
não pode ter pontos conjugados, de modo que
$\text{exp}_p$ é um
difeomorfismo local. Também por completitude
sabemos que
$\text{exp}_p$ é sobrejetiva: qualquer ponto
$q$ está ligado a $p$
mediante uma geodésica, e essa geodésica
coincide com uma geodésica partindo de
$p$ dada como a imagem de uma reta em $T_pM$
sob $\text{exp}_p$.

Portanto, o desafio é provar que
$\text{exp}_p$ também é injetiva.
Suponha que não é o caso, i.e. considere dois
pontos $v_1$ e $v_2$ em
$T_pM$ tais que
$\text{exp}_pv_1=\text{exp}_pv_2:=q$. As
imagens das retas geradas por $v_1$ e $v_2$
sob $\text{exp}_p$ são
duas geodésicas $\gamma_1$ e $\gamma_2$ em
$M$ ligando $p$ e $q$.

Como $M$ é simplesmente conexa, existe uma
homotopia entre $\gamma_1$ e
$\gamma_2$. Podemos aplicar o lema de
Klingenberg para obter uma curva
$\alpha_{K_0}$ tal que $$
\ell(\gamma_1)+\ell(\alpha_{K_0}) \geq
2\pi/\sqrt{K_0} $$ para qualquer $K_0>0$.
Isso mostra que o comprimento das
curvas na homotopia não está limitado, o que
não é possível já que a homotopia é
uma função contínua definida em um compacto,
pelo qual a sua imagem debe ser
compacto e portanto limitado. {\bf Errado!
Por exemplo, um círculo percorrido
miles de vezes tem comprimento
arbitrariamente grande e está contido num
compacto.}
\end{proof}

\begin{exercise}
\label{differential-geometry-exercise-geode-sem-ponto-conju-e}
$\gamma$ geodésica sem pontos conjugados em
$[0,a]$. Então $\gamma$ é localmente
minimizante.
\end{exercise}

\begin{proof}
Temos dois casos:
\begin{enumerate}
\item Se $E''(0)<0$ usamos o teorema do
índice. Isto é, não é possível que
 $E''(0)<0$ porque isso implica que o índice
 da forma do índice é diferente
 de zero, ou seja, existem pontos conjugados.
\item Se $E''(0)=0$,
\begin{enumerate}
\item[(0)] Como $\gamma$ não tem pontos
conjugados, então $I$ não tem
nulidade em todo  $\mathcal{V}$.
\item $I_{\mathcal{V}^-}$ não tem índice nem
nulidade. Segue do ponto anterior +
algum passo na prova.
\item $I_{\mathcal{V}^-}$ não tem índice.
\item 2+3 dão que $I_{\mathcal{V}^-}>0$.
\item Logo $I>0$.
\end{enumerate}
\end{enumerate}
\end{proof}

\begin{exercise}
\label{differential-geometry-exercise-inter-minim-hyper}
Let $(M^n,g)$ be a complete Riemannian
manifold with $\text{Ric}_M>0$.
Suppose $M_1$, $M_2$ are minimal (compact?)
complete hypersurfaces. Show that
$M_1 \cap M_2 \neq \emptyset$.
\end{exercise}

\begin{proof}
Let $\gamma$ be a minimizing geodesic joining
$p \in M_1$ to $q\in M_2$ and
minimizing the distance between $M_1$ and
$M_2$. Let $e_1,\ldots,e_{n-1}\in
T_pM$ be such that along with $\gamma'(0)$
form an orthonormal basis of $T_pM$.
Let $E_i$ be the parallel transport of $e_i$
along $\gamma$. Notice that
$E_i(\ell)$ is tangent to $M_2$ since it is
orthogonal to $\gamma'(\ell)$ and
$M_2$ is a hypersurface, but it is not
immediate that the variational curve will
be contained in $M_2$ if we define the
variation using the exponential map. This
is essential for the second variational
formula to vanish since $\gamma$ is
minimizing among curves joining $M_1$ and
$M_2$.

We need to produce the variation not from the
vector field but from the
variational curves. Take rectifying
coordinates for $M_1$ at $p$. Define the
curve $c^i$ as the $i$-th coordinate curve on
the whole neighbourhood of $p$
along $\gamma$.

Now take the velocity of $c^i$ for some
$t>0$. Parallel transport this vector
along $\gamma$ until it reaches a rectifying
neighbourhood of $q$. Along the
path of $\gamma$ we can use the exponential
to define a variation. Once in the
rectifying neighbourhood of $q$, we must
define by hand the variational curves
to ensure that the curve at $q$ is contained
in $M_2$.

Suppose that $\gamma(t_0)$ is in the
rectifying neighbourhood of $q$. Consider a
smooth function $\varphi:[t_0,\ell]\to
\mathbb{R}$ that is zero at $t_0$ and $1$
at $\ell$. We can smoothly deform the
exponential curves to a curve contained in
$M_2$ at $t=\ell$ by defining
$$
c^s_i(t):=(1-\varphi)\text{exp}_{\gamma(t)}
sE_i(t)+\varphi
sE_i(t)
$$
Notice that the variational vector is
orthogonal to $M_2$ at $q$, so that we may
apply Exercise
\ref{differential-geometry-exercise-two-submanifolds}. Thus we can
express the second
variation formula in such a way that when
summing over all indices we obtain the
mean curvatures of each submanifold, which
vanish, thus obtaining a
contradiction with the hypothesis that
$\text{Ric}>0$.
\end{proof}

\begin{exercise}
\label{differential-geometry-exercise-even-dimen-posit-k-compa}
Prove que $M^{2n},K_M>0$ compacta, então ela
é simplesmente conexa.
\end{exercise}

\begin{exercise}
\label{differential-geometry-exercise-odd-dimen-posit-k-orien}
$M^{2n+1}$, $K_M>0$ compacta $\implies$
orientável.
\end{exercise}

\begin{proof}
I want to show that $M^{2n+1}$ is
diffeomorphic to its double orientable cover
 $\tilde{M}$. Since $M$ is odd-dimensional,
 so is $\tilde{M}$.
\end{proof}

O exercício anterior usa Exercício 6 da lista
6. Ou, equivalentemente,

\begin{exercise}
\label{differential-geometry-exercise-inver-orien-short}
Se $M^{2n+1}$ tem $K_M >0$ e $\gamma:S^1 \to
M$ inverta orientação, então é
possível encurtar $\gamma$ na sua classe de
homotopia.
\end{exercise}

\begin{exercise}
\label{differential-geometry-exercise-produ-compa-does-not-admit}
Let $N,M$ be compact. Then $N \times M$ does
not admit a metric with $K<0$.
\end{exercise}

\begin{proof}
If both $N$ and $M$ are simply connected then
they cannot be compact by
Hadamard's theorem. Assume first that both
are not simply connected. Then each
has a nontrivial homotopy class, and such a
pair generates an abelian subgroup
of $\pi_1(M \times N)$ that is not isomorphic
to $\mathbb{Z}$, a contradiction
with Preissman Theorem
\ref{differential-geometry-theorem-Preissman}.

Now suppose that $N$ is simply connected and
$M$ is not. The universal cover of
$M\times N$, if it admits a metric with
negative sectional curvature,
must be a Hadamard manifold, so that
$\tilde{M} \times N \cong \mathbb{R}^n$.
Then the induced map on cohomology is
injective map ({\bf why?})
$\pi^* :H^{k}(M)\to H^{k}(\tilde{M} \times
N)=H^{k}(\mathbb{R}^n)=0$. This means
by Poincaré duality that $0=H^{n}(M)\cong
H_0(M)\neq 0$, a contradiction.
\end{proof}

\begin{exercise}
\label{differential-geometry-exercise-SnxR-no-positive-Ric}
Show that $S^n \times \mathbb{R}$ does not
admite a complete metric with
$\text{Ric}>0$.
\end{exercise}

\begin{proof}
If it did, then by Preissman theorem
\ref{differential-geometry-theorem-Preissman} then $S^n\times
\mathbb{R}\cong N \times \mathbb{R}$, but the
upshot is that that the latter
isomorphism is an isometry, and the
$\mathbb{R}$ factor is flat, so we can't
have strictly positive Ricci.
\end{proof}

\begin{exercise}
\label{differential-geometry-exercise-compl-witho-compa-set-isom}
Let $(M^n,g)$ be complete and such that there
exists a compact set
 $K\subset M$ such that $M\setminus K \cong
 \mathbb{R}^n \setminus \tilde{K}$
for some compact set $\tilde{K} \subset
\mathbb{R}^n$. If $\text{Ric}M \geq 0$,
then $M^n \cong \mathbb{R}^n$.
\end{exercise}

\begin{proof}
Since $M\setminus K$ is isometric to
$\mathbb{R}^n\setminus \tilde{K}$, there
must be a line (oops! this is not obvious).
But suppose for now that there {\it
is} a line on $M$. Then by the Splitting
theorem \ref{differential-geometry-theorem-splitting} there
is an isomorphism $M \cong N \times
\mathbb{R}$. ({\it My original idea:}
 If we manage to show that $N$ also has a
 line we would have that
 $M \cong N_2 \times \mathbb{R}^2$.  if we
 can further do this process and find
 that all $N_i$ have lines,
then we end up with $M \cong N_k \times
\mathbb{R}^k$ for some manifold $N_k$
without lines.)

Let's fix some notation. Let
$f:M\setminus K \to \mathbb{R}^n
\setminus\tilde{K}$ be the given isometry and
$\varphi:N^{n-1}\times\mathbb{R} \to M$ the
isometry given by the Splitting
theorem. Consider the restriction $f \circ
\varphi |_{ N^{n-1}\times \{ a\}}$
for some $a \in \mathbb{R}$. Observe que
$N^{n-1}\times \{ a\}$ é totalmente
geodésica (cf. Exercise
\ref{differential-geometry-exercise-total-geode-subma-produ})
, e portanto também a sua imagem sob  $f
\circ \varphi$. Toda subvariedade
totalmente geodésica de $\mathbb{R}^n$ é um
plano, e acabamos.

However, notice that it is not immediate from
the
isomorphism $f$ that $M$ has a line. Fix a
line $L$ in $\mathbb{R}^n$. Then its
inverse image under $f$ is minimizing in $M
\setminus K$, but there could be a
curve that minimizes distance between some
two points of $f^{-1}(L)$ that is not
contained in $f^{-1}(L)$.

Suppose that $f^{-1}(L)$ is not minimizing in
$M$. Then there are two points $j$
and $-j$ in $f^{-1}(L)$ and a minimizing
geodesic $\sigma$ joining them and
intersecting $K$. Then the path that joins
any point before $j$ and any point
after $j$ on $\sigma$, before $\sigma$
intersects $K$, is a minimizing segment
contained in $M\setminus K$ and should thus
be contained in $f^{-1}(M)$. This
shows that $f^{-1}(K)$ accumulates near $K$,
which is not possible.
\end{proof}

\begin{exercise}
\label{differential-geometry-exercise-S1xS1xR-admite-curvatura-1}
Mostre que $S^1 \times S^1 \times \mathbb{R}$
{(\bf não?)} admite uma métrica de
 curvatura $K\equiv-1$. {\it Dica.} Teorema
 de Preissman
\ref{differential-geometry-theorem-Preissman}.
\end{exercise}

\begin{proof}
Porque $\mathbb{Z} \times \mathbb{Z}$ é
abeliano mas não é isomorfo a
$\mathbb{Z}$.
\end{proof}

\begin{exercise}
\label{differential-geometry-exercise-rpnxr-admit-curva-posit}
$\mathbb{R}P^{n}\times \mathbb{R}P^{m}$
admite uma métrica com
curvatura positiva?
\end{exercise}

\begin{proof}
Se $m,n$ são os dois ímpares, são
orientáveis, então o produto é orientável e
de
dimensão par, e pelo Teorema de Synge
\ref{differential-geometry-theorem-Synge} deve ser
simplesmente conexa, absurdo.

Se as duas tem dimensão par, considere o
recobrimento duplo orientável, que deve
ser compacto (pode usar Bonnet-Myers
\ref{differential-geometry-theorem-Bonnet-Myers}), com curvatura
positiva e dimensão par, e portanto é
simplesmente conexo. Ou seja, trata-se do
recobrimento universal. Porém, ele é um
recobrimento duplo e portanto a fibra
tem dois elementos, e desse jeito o
$\text{Deck}\cong
\pi_{1}(\mathbb{R}P^{n}\times\mathbb{R}P^{m})
$
deve ser $\mathbb{Z}_2$,
mas não é.

E se uma é de dimensão par e a outra ímpar… é
igualzinho que no caso anterior.
\end{proof}

\begin{exercise}
\label{differential-geometry-exercise-Sn-Tn-CPn-existem-curvaturas}
Considere $S^n,\mathbb{T}^n$ e
$\mathbb{C}P^{n}$. Existem métricas com
$K>0$, $K
\leq 0$, $K<0$, $K \geq 0$, $\text{Ric}>0$,
$\text{Ric} \geq 0$? Faça uma
tabela.
\end{exercise}

\begin{exercise}
\label{differential-geometry-exercise-Hadamard-manifold-sphere}
Let $M^n$ be a Hadamard manifold and
$S^{n-1}_r \subset M$. Mostre que
 a curvatura média satisfaz $H_S\geq
 \frac{1}{r}$. {\it Dica.} A curvatura média
da esfera em $\mathbb{R}^n$ é $\frac{1}{r}$.
\end{exercise}

\begin{proof}
Considere uma geodésica que liga $p$, o
centro da esfera, e o bordo dela.
Considere um referencial de campos de Jacobi
$J_i$. Lembre que cada um dele
satisfaz $A J=J'$. Então
$$
H=\sum_{i}\left\langle{}A
\frac{J_i}{|J_i|},J_i\right\rangle{}=
$$
\end{proof}

\begin{exercise}
\label{differential-geometry-exercise-product-of-compact-and-Rn}
$M^n$ compact, $\text{Ric} \geq 0$, then
$\tilde{M}=N \times \mathbb{R}^n$ for
some compact manifold $N$, where  $\tilde{M}$
is
the universal cover of $M$.
\end{exercise}

\begin{proof}
If $\tilde{M}$ is compact we are done.
So suppose it is not.
Then there exists a
ray. We want to construct
a line to use Splitting Theorem
\ref{differential-geometry-theorem-splitting}.
Consider the fundamental domain of the cover
$\tilde{M}$,
which exists since
$M$ is compact: we may take a finite covering
of normal neighbourhoods of $M$,
so that the preimage under $\pi$ is a compact
set containing at least one
representative of each class.

Let $p_i$ be a sequence along the ray
$\gamma$
so that $p_i$ goes to infinity.
For each of them there exists a
deck transformation $F_i$ that puts $p_i$
inside
$K$. Consider the ray $F_i \circ \gamma$.
\end{proof}

\begin{exercise}
\label{differential-geometry-exercise-rays-semispaces}
Lista 8, Exercício 12. Let $M$ be a complete
Riemannian manifold not compact
with $K \geq 0$. Let $\sigma$ be a ray of
$M$. Define for $t>0$ the {\it
semispace}
$$
H_t^{\sigma}:=M\setminus
\bigcup_{s>0}\Big(\sigma(t+s),s).
$$
Show that $H_t^{\sigma}$ is totally convex,
that is, is a geodesic segment
 $\alpha:[0,1]\to M$ has its endpoints in
 $H_t^{\sigma}$, i.e.
$\alpha(0),\alpha(1) \in H_t^{\sigma}$, then
that segment is completely
contained in $H_t^{\sigma}:\alpha([0,1])
\subset H_t^{\sigma}$. {\it Hint.}
Suppose that it's false and use Toponogov.
\end{exercise}

\begin{proof}
Suppose that there is a geodesic
$\gamma:[0,1]\to M$ with $\gamma(0),\gamma(1)
\in M$, but it intersects $H_t^{\sigma}$.
Consider the triangle joining $p$,
$\sigma(t+s)$ and some other point.

Choose the other point so that you obtain a
contradiction using the structure of
$H_t^{\sigma}$ of embedded balls---if a point
is in any ball, it must be in all
the next ones.
\end{proof}

\section{Lista 3}
\label{differential-geometry-section-lista-3}

\begin{exercise}[Curvas minimizantes]
\label{differential-geometry-exercise-curvas-minimizantes}
\begin{enumerate}
\item Seja $\gamma$ uma curva suave por
partes parametrizada por comprimento de
 arco, conectando $p$ a $q$, pontos de uma
 variedade Riemanniana $(M,g)$.
 Mostre que se $d(p,q)=\ell(\gamma)$, então
 $\gamma$ é uma geodésica.
 Dizemos que $\gamma$ realiza a distância
 entre $p$ e $q$.
\item Suponha que $\gamma,\sigma:[0,2]\to M$
são geodésicas distintas
 e satisfazem: $\gamma(0)=\sigma(0):=p$,
 $\gamma(1)=\sigma(1):=q$,
$\gamma$ e $\sigma$ realizam a distância
entre $p$ e $q$.
 Mostre que $\gamma$ não realiza a distância
 entre $p$ e $\gamma(1+s)$
para nenhum $s>0$.
\end{enumerate}
\end{exercise}

\begin{proof}
\begin{enumerate}
\item Primeiro suponha que $\gamma$ é suave.
Podemos ver que trata-se de uma
geodésica usando seu campo de velocidades
como campo variacional na primeira
fórmula da variação:
$$
E'(0)=\int_0^a |\gamma'|=0 \implies
\gamma'\equiv 0
$$
Uma geodésica quebrada não pode ser
minimizante porque em uma vizinhança
totalmente normal (cf.
\ref{differential-geometry-proposition-total-norma-neigh})
do ponto singular podemos achar uma geodésica
radial (suave)
 que acurta a distância entre qualquer ponto
 antes do ponto singular
 e qualquer ponto depois do ponto singular.
\item Se as duas geodésicas se intersectam
não tangencialmente, podemos
construir um caminho mais curto entre $p$ e
$\gamma(1+s)$ ao longo de $\sigma$,
a menos de suavizar a quina perto de $q$. Se
as curvas se intersectam
 tangencialmente devem ser a mesma.
\end{enumerate}
\end{proof}

\section{Lista 5}
\label{differential-geometry-section-lista-5}

\begin{exercise}
\label{differential-geometry-exercise-l5-5}
\cite{doc} Capítulo XII, Exercício 6. Uma
geodésica $\gamma:[0,\infty)\to M$ em
uma variedade Riemannana $M$ é um {\it raio
partindo de $\gamma(0)$} se ela é
minimizante entre $\gamma(0)$ e $\gamma(s)$
para todo  $s \in (0,\infty)$.
Admita que $M$ é completa, não compacta, e
seja $p \in M$. Mostre que $M$ contém
um raio partindo de $p$.
\end{exercise}

\begin{proof}
Como $M$ é compacta e completa, ela não pode
ser limitada. Então
existe uma sequência de pontos $p_i$  tal que
$\lim_{i \to \infty}
d(p_i,p)=\infty$. Para cada ponto podemos
pegar uma geodésica minimizante
$\gamma_n$ ligando $p$ e $p_n$. Podemos
associar a cada geodésica um vetor
unitário $v_i \in S^n \subset T_pM$ tal que
$\gamma_i(t)=\text{exp}_p(tv_i)$.

Como $S^n \subset T_pM$ é compacta existe um
vetor $v$ limite da sequência
$v_i$. A geodésica
$\gamma(t):=\text{exp}_p(tv)$ é um raio, pois
por
continuidade de $\text{exp}_p$ e da distância
Riemanniana,
$$
d(\gamma(0),\gamma(t))=d\left(p,\text{exp}
_p(tv)\right)
=d\left(p,\text{exp}_p\left(t\lim_{i \to
\infty} v_i\right)\right)
=\lim_{i \to \infty} d(p,\gamma_i (t))
$$
Pegando $i$ suficientemente grande, teremos
que $\gamma_i$ é minimizante entre
$p$ e $\gamma_i(t)$. Isso significa que
$d(p,\gamma_i(t))$ está dada como o
comprimento de $\gamma_i$ até esse ponto, e
pegando o limite concluímos que o
comprimento de $\gamma$ é a distância entre
$\gamma(0)$ e $\gamma(t)$.
\end{proof}

\begin{definition}
\label{differential-geometry-definition-line-riemannian-manifold}
Let $(M,g)$ be a Riemannian manifold. A
geodesic $\gamma:\mathbb{R} \to M$ is
called a {\it line} if
$d(\gamma(t),\gamma(s))=|t-s|$ for all $t,s
\in
\mathbb{R}$.
\end{definition}

\begin{exercise}
\label{differential-geometry-exercise-l5-6}
Mostre que toda métrica completa em $S^n
\times \mathbb{R}$ admite uma linha.
\end{exercise}

\begin{proof}
Considere o conjunto $S^n\times
\mathbb{R}\setminus(S^n \times\{0\})$.
Trata-se de um aberto disconexo, e portanto
não fechado, nem compacto nem
limitado.  Dentro desse conjunto podemos
pegar duas sequências de pontos $p_i$ e
$q_i$, onde a segunda coordenada de $p_i$ tem
signo positivo, e a segunda
coordenada de $q_i$ tem  signo negativo para
toda $i$. Considere a geodésica
$\gamma_i$ que liga $p_i$ com $q_i$, que deve
passar pela esfera $S^n\times
\{0\}$. Então obtemos uma sequência $v_i$ de
vetores no fibrado tangente
unitário de $S^n\times\{0\}$ dadas como as
velocidades das geodésicas
$\gamma_i$. Como tal fibrado é compacto,
achamos um limite $v$ dessa sequência
que por um argumento análogo ao exercício
anterior, converge a um vetor cuja
geodésica associada é uma linha. ($\gamma$ é
minimizante, e como está
parametrizada por comprimento de arco obtemos
a condição dada na Definição
\ref{differential-geometry-definition-line-riemannian-manifold}.
\end{proof}

\section{Lista 6}
\label{differential-geometry-section-lista-6}

\begin{exercise}
\label{differential-geometry-exercise-posit-scala-but-no-posit-duplicate-2}
Encontre um exemplo de variedade suave que
admite alguma métrica Riemanniana com
curvatura escalar positiva, mas não admite
uma métrica Riemanniana com Ricci
positivo.
\end{exercise}

\begin{proof}
Se o problema for para $\text{Ric}\geq
\delta>0$, poderíamos usar o teorema de
Bonnet-Myers \ref{differential-geometry-theorem-Bonnet-Myers} como
segue. Um produto de uma esfera com
uma reta tem curvatura escalar positiva, mas
não admite uma métrica com
$\text{Ric}\geq\delta>0$ porque não é
compacto.

{\bf Ideia errada:} Considere $S^1\times
S^1$. É claro que na métrica produto
temos curvatura escalar positiva (erro! As
variedades de dimensão 1 tem tensor
de curvatura nulo pela antisimetria nas
primeiras (ou últimas) duas entradas!)
pelo Exercício
\ref{differential-geometry-exercise-curvatures-in-product}. Como é
uma superfície, a
curvatura de Ricci coincide com a curvatura
seccional. (Ricci é o promédio das
curvaturas seccionais.) A curvatura seccional
é a curvatura Gaussiana. Se
$S^1\times S^1$ tivesse uma métrica com
curvatura Gaussiana positiva, pelo
teorema de Gauss-Bonnet teria característica
de Euler positiva, mas a sua
característica de Euler é zero. No nosso
curso, aplicamos o teorema de Synge
\ref{differential-geometry-theorem-Synge} para uma variedade
orientável de dimensão par com curvatura
escalar positiva, que deve ter grupo
fundamental trivial, uma contradição com o
fato de que $\pi_{1}(S^1\times
S^1)=\mathbb{Z}\times\mathbb{Z}$.

Nem tudo está perdido: considere
$\mathbb{R}P^{2}$. Na métrica induzida pela
projeção quociente desde $S^2$ temos
curvatura escalar constante 1, e é uma
superfície, de modo que a curvatura seccional
coincide com a curvatura de Ricci,
aplicamos o teorema de Synge
\ref{differential-geometry-theorem-Synge}, mas agora
$\mathbb{R}P^{2}$
não é orientável, então seu grupo deveria ser
$\mathbb{Z}_2$… pois é.

Depois reparei que esse argumento está
totalmente errado: no caso das
superfícies nem só a curvatura escalar
coincide com a Ricci… também com a
seccional!

 $S^3\times S^1$. Então sim: pelo teorema de
Synge \ref{differential-geometry-theorem-Synge}, uma variedade
orientável de dimensão par com
curvatura {\it escalar} positiva deveria ter
grupo fundamental trivial, que não
é o caso. Mas já não é verdade que a
curvatura de Ricci coincide com a curvatura
 seccional; poderíamos ter um plano com
 curvatura seccional negativa e ainda ter
Ricci positivo. Temos que descartar a
possibilidade de ter um plano com
curvatura seccional não positiva.
\end{proof}

\begin{exercise}
\label{differential-geometry-exercise-posit-ricci-but-no-posit-duplicate-2}
Encontre um exemplo de variedade suave que
admite uma métrica Riemanniana com
Ricci positivo, mas não admite uma métrica
Riemanniana com curvatura seccional
positiva.
\end{exercise}

\begin{proof}
Considere $S^3\times S^1$. A métrica produto
tem $\text{Ric}>0$, já que ele é a
soma dos tensores de Ricci em cada fator de
acordo ao Exercício
\ref{differential-geometry-exercise-curvatures-in-product} (a
curvatura de Ricci de $S^3$ é 1,
enquanto que a curvatura escalar de $S^1$ é
zero). Agora suponha que existe uma
 métrica com $K>0$. Então pelo teorema de
 Synge \ref{differential-geometry-theorem-Synge}, como
$S^3\times S^1$ é compacta, de dimensão par e
orientável, o grupo fundamental
dela deveria ser trivial, mas não é o caso.
\end{proof}

\begin{exercise}
\label{differential-geometry-exercise-two-total-geode-subma}
Let $(M,g)$ be a complete, connected
Riemannian manifold with positive
curvature. Let $M_1$ and $M_2$ be two totally
geodesic submanifolds such that
$\dim M_1+\dim M_2\geq n$. Show that $M_1\cap
M_2\neq \emptyset$.
\end{exercise}

\begin{proof}[Solução]
Suponha que a interseção não é vazia. Então
existe uma geodésica minimizante não
trivial que minimiza a distância entre $A$ e
$B$. Sabemos que essa geodésica
intersecta tanto $A$ quanto $B$
ortogonalmente. Queremos ver que $E''(0)<0$,
ou
seja que $\gamma$ não pode ser minimizante,
uma contradição.

Suponha que conseguimos construir uma
variação de $\gamma$ tal que todos os
pontos iniciais $f(s,0)$ estão em $A$, e o
mesmo acontece com os pontos finais,
i.e. $f(s,\ell) \in B$. Então as derivadas
das curvas $f(s,0)$ e $f(s,\ell)$ em
$s=0$ são perpendiculares a  $\gamma'(0)$ e
$\gamma'(\ell)$. Ou seja, temos que
$\left\langle{}\gamma'(0),\nabla_{\partial_s}
V\right\rangle{}|_{0}^\ell=0$.

Suponha ainda que essa variação tem campo
variacional paralelo. Então acabou
porque a segunda fórmula da variação diz que
\[E''(0)=-\int_0^\ell
\left\langle{}R_{\gamma'}V,V\right\rangle{}<0\]

Vamos construir essa variação. O passo
inicial é fácil: pegamos qualquer vetor
$v \in T_a A$ e transportamos paralelamente
ao longo de $\gamma$ para obter o
campo paralelo $V \in \mathfrak{X}_\gamma$.
Como $A$ é totalmente geodésica, a
variação
$f(s,t):=\text{exp}_{\gamma(0)}(sV(0))$ fica
dentro de $A$.

Para concluir basta ver que $ V(\ell) \in
T_bB$, já que $B$ também é totalmente
geodésica. Isso segue da última hipótese.
Podemos realizar esse processo para
cada vetor básico de  $T_a A$, obtendo $\dim
A$ vetores linearmente
independentes em $T_bB$ (já que o transporte
paralelo é uma isometria). Se
nenhum deles ficasse em $T_bB$, poderíamos
construir um espaço de dimensão $\dim
A$ estritamente contido em $T_b M\setminus
T_b B$, absurdo pois isso implicaria
que $\dim A< \text{codim} B$, porém, $\dim A+
\dim B \geq n$.

\textbf{Edit.} Depois de reler a prova do
exercício seguinte, parece que não era
necessário usar que $\gamma$ intersecta
ortogonalmente $A$ e $B$, pois a
variação que usei foi por geodésicas, então
$\nabla_{\partial_s}V=0$.
\end{proof}

\begin{exercise}
\label{differential-geometry-exercise-free-homotopy-variation}
Seja $M^{2n}$ uma variedade Riemanniana de
dimensão par, completa, orientável e
com curvatura seccional $K>0$. Seja $\gamma$
uma geodésica fechada em $M$ de
comprimento $\ell(\gamma)$. Mostre que
existem curvas livremente homotópicas a
$\gamma$ em $M$, arbitrariamente próximas de
$\gamma$, que possuem comprimento
menor que $\ell(\gamma)$.
\end{exercise}

\begin{proof}
Considere qualquer vector unitário ortogonal
a $\dot \gamma$, defina $V\in
\mathfrak{X}_{\gamma}''$ como sendo o
transporte paralelo de $v$ ao longo de
$\gamma$. Note que $\dot V=0$. Defina a
variação
$f(s,t)=\text{exp}_{\gamma(t)}sV(t)$. Note
que
$\nabla_{\partial_s}f_s=0$.

A segunda fórmula da variação nos diz que
\[E''(0)=-\int_0^a \left\langle{}R_{\dot
\gamma}V,V\right\rangle{}=-\int_0^a K<0\] já que
$K>0$. Como $\gamma$ é uma geodésica,
temos que para $s$ pequeno
\[\frac{1}{a}\ell^2(c^s)\leq E(s) < E(0)=
\frac{1}{a}\ell^2(\gamma)\] Note que essa
variação pode não preservar o
fechamento das curvas. Para resolver isso
olhemos ao seguinte desenho:

Fica claro que se $V(0)=V(a)$ as curvas $c^s$
são fechadas para toda $s$. Então
a condição que precisamos é que
\[P^\gamma_{0,a}V(0)=V(a)\] onde $P$ é o
transporte paralelo. Dito de outra forma,
basta ver que $P$ tem um ponto
fixo---além de $\gamma'(0)=\gamma'(a)$, que
já é um ponto fixo. Esse argumento é
uma imitação da prova do teorema de
Weinstein: como $P$ preserva orientação, o
determinante dele deve ser positivo, ou seja,
o produto dos seus autovalores
deve ser positivo. Restringindo-nos ao
subespaço $\gamma'(0)^\perp$, que tem
dimensão ímpar, concluímos que para que o
produto dos autovalores de $P$ seja
positivo, deve ter pelo menos um que seja
$1$.
\end{proof}

\section{Lista 7}
\label{differential-geometry-section-lista-7}

\begin{exercise}
\label{differential-geometry-exercise-minim-impli-no-conju-point}
Seja $\gamma:[0,a]\to M$ uma geodésica em uma
variedade Riemanniana $M$.

 Prove que se $\gamma$ é minimizante, então
 $\gamma$ não possui pontos
 conjugados em  $(0,a)$. Encontre um exemplo
 de geodésica $\gamma:[0,a] \to M$
sem pontos conjugados que não é minimizante.
\end{exercise}

\section{Lista 8}
\label{differential-geometry-section-lista-8}

\begin{exercise}
\label{differential-geometry-exercise-l8-1}
Prop. 2.12 do capítulo XIII, \cite{doc}. Seja
$p \in M$.
Suponha exista um ponto $q \in C_m(p)$ que
realiza a distância de $p$ a
$C_m(p)$. Então:
\begin{enumerate}
\item ou existe uma geodésica
minimizante $\gamma$ de $p$ a $q$ ao longo da
qual $q$ é
conjugado a $p$,
\item ou existem exatamente duas geodésicas
minimizantes $\gamma$ e $\sigma$ de $p$ a
$q$; além disto,
$\gamma'(\ell)=-\sigma'(\ell)$,
$\ell=d(p,q)$.
\end{enumerate}
\end{exercise}

\begin{proof}
No final da secção 1 do Capítulo XIII está
especificado que as variedades que se
consideram no capítulo são completas.
Portanto podemos supor que existe
uma geodésica minimizante $\gamma$ ligando
$p$ e $q$.

Então podemos aplicar a Proposição
\ref{differential-geometry-proposition-cut-point-chara}:
 um ponto $q$ é o cut point de $p$ ao longo
 de uma
geodésica minimizante $\gamma$ se e somente
se
alguma das seguintes condições é verdadeira:
(a) $q$ é o primeiro ponto conjugado a $p$ ao
longo de $\gamma$, ou
(b) existem duas geodésicas minimizantes
ligando $p$ e $q$.

Se (a) é verdadeira, terminamos. Se (b) é
verdadeira, temos que existe outra
geodésica minimizante $\sigma$ ligando $p$ e
$q$.

{\bf Intento pessoal (não funcionou):}
considere a variação por geodésicas
$$
f(s,t):=\text{exp}_{\gamma(t)}
(s\text{exp}_{\gamma(t)}^{-1}\sigma(t))
$$
cujo campo de Jacobi
$$
J(t)=d_{t\text{exp}_{\gamma(t)}^{-1}\sigma(t)
}
\text{exp}_p(\text{exp}_{\gamma(t)}^{-1}
\sigma(t))
$$
se anula em $t=0,1$. Pensei que esse campo de
Jacobi estava bem definido pelo
Corolário 2.8, Cap. XIII (Lema
\ref{differential-geometry-lemma-outsi-cut-locus-exist-minim}), que
assegura
 que $\text{exp}_p$ é injetiva fora do cut
 locus de $p$. Porém, precisaria
 que a exponencial ao longo de $\gamma$,
 $\text{exp}_{\gamma(t)}$, for injetiva fora
 do cut locus de $p$.
Essa condição seria garantida se
$d(p,q)=i(M)$, ou seja, se a exponencial
 $\text{exp}_{p'}$ for injetiva na bola de
 raio $d(p,q)$ para qualquer
$p' \in M$. Em conclusão: minha variação não
está bem definida.
(Se estivesse, com
isso terminaria o exercício, pois obtemos que
o único caso em que $p$ não é
conjugado a $q$ é se
$\gamma'(\ell)=-\sigma'(\ell)$, i.e. se o
campo de Jacobi
associado à variação é nulo.)

\bigskip

A variação certa é produzida no \cite{doc}
supondo que
$\gamma'(\ell)\neq -\sigma'(\ell)$; isso
implica que existe um vetor
 $V \in T_q M$ tal que
$\left\langle{}\gamma'(\ell),V\right\rangle{}< 0$ e
$\left\langle{}\sigma'(\ell),V\right\rangle{}< 0$.

Uma maneria simples de conferir a existência
desse vetor $V$ é olhando para o
plano gerado por $\gamma'(\ell)$ e
$\sigma'(\ell)$ (usando que são linearmente
independentes). O conjunto de vetores $V$
nesse plano satisfazendo
$\left\langle{}\gamma'(\ell),V\right\rangle{}<0$ é um
semiespaço, e o mesmo acontece com
 os vetores satisfazendo
 $\left\langle{}\sigma'(\ell),V\right\rangle{}<0$. Esses dois
 semiespaços
devem ter interseção não vazia precisamente
porque
$\gamma'(\ell)\neq -\sigma'(\ell)$.

Agora usamos o fato de que $p$ e $q$ não são
conjugados para
achar uma vizinhança do vetor $\ell
\gamma'(0)$ onde $\text{exp}_p$ é um
difeomorfismo e assim levantar alguma
curva $r$ que realize o vetor $V$ ao espaço
tangente,
 digamos $v:(-\varepsilon,\varepsilon) \to
 T_q M$.

A variação então está dada por
$$
f(s,t):=\text{exp}_p\left(\frac{t}{\ell}v(s)
\right)
$$
Podemos aplicar a primeira fórmula da
variação, na qual: o termo com a
integral é zero porque $\gamma$ é uma
geodésica; o termo da sumatoria é zero
porque trata-se de uma variação suave; e o
primeiro termo dos extremos é zero
porque a variação fixa o ponto de partida.
Obtemos que
$$
E'(0)=\left\langle{}V,\gamma'(0)\right\rangle{}<0
$$
Note que Manfredo escreve a equação anterior
usando o funcional de distância.
Isso também é válido: de fato a primeira
fórmula da variação pode ser
deduzida de maneira análoga (feito em sala
para variações próprias)
 para o funcional de distância no caso de
 geodésicas parametrizadas por
comprimento de arco (cf. Teo. 6.3
\cite{ler}).

O fato da derivada do funcional de
comprimento ser negativa nos diz que para
valores perto de $s=0$ podemos achar curvas
com menor comprimento. Mais
precisamente, existe um $s>0$ tal que
$\ell(f(s,\cdot))<\ell(\gamma)$.

É claro que o mesmo procedimento genera uma
variação $\tilde{f}$ de $\sigma$, e
podemos supor que a mesma $s>0$ faz
$\ell(\tilde{f}(s,\cdot))<\ell(\sigma)$.

Concluímos seguindo o argumento do Professor
Manfredo: se
$\ell(\gamma_s)=\ell(\sigma_s)$, então temos
duas geodésicas com o mesmo
comprimento ligando $p$ e
$\gamma_s(\ell)=\sigma_s(\ell)$; isso
significa
que existe um ponto $\tilde{t} \in (0,\ell]$
tal que $\gamma_s(\tilde{t})$ é
 o cut point de $p$. Porém, isso contradiz o
 fato de que $q$ realiza a
distância entre $p$ e o seu cut locus.

Finalmente, se
$\ell(\gamma_s)<\ell(\sigma_s)$, segue que o
cut point de $p$
 ao longo de $\sigma_s$ está a distância
 menor do que $q$, que de novo contradiz
a nossa hipótese.
\end{proof}

\begin{exercise}
\label{differential-geometry-exercise-l8-2}
Proposição 2.13, Cap. XIII \cite{doc}. Se a
curvatura seccional $K$ de uma
variedade Riemanniana completa $M$ satisfaz
$$
0< K_{\text{min}}\leq K\leq K_{\text{max}},
$$
Então
\begin{enumerate}
\item $i(M) \geq \pi/\sqrt{K_{\text{max}}}$,
ou
\item existe uma geodésica fechada $\gamma$
em $M$, cujo comprimento é menor do
que o de qualquer outra geodésica fechada em
$M$, tal que
$$
i(M)=\frac{1}{2}\ell(\gamma)
$$
\end{enumerate}
\end{exercise}

\begin{proof}
Suponha que (1) não é verdadeiro.

Note que pelo Teorema de Bonnet-Myers
\ref{differential-geometry-theorem-Bonnet-Myers}, $M$ é compacta e
portanto $C_m(p)$ é compacto para
todo $p$. Isso nos permite achar dois pontos
$p$ e $q$ cuja distância é o raio
de injetividade  $i(M)$.

Então podemos usar o exercício anterior para
$p$ e $q$.  Primeiro vamos ver o
que acontece no caso da segunda possibilidade
daquele exercício, i.e. que
existam exatamente duas geodésicas ligando
$p$ e $q$, digamos $\gamma$ é
$\sigma$, tais que
$\gamma'(\ell)=-\sigma'(\ell)$ onde
$\ell=d(p,q)$.  Considere
$\gamma * \overline{\sigma}$, onde $*$ é a
concatenação de curvas e
$\overline{\sigma}$ é a curva $\sigma$
percorrida em sentido contrário. É claro
que $\gamma * \overline{\sigma}$ é uma
geodésica fechada de comprimento $2
d(p,q)$. Note que por definição
$d(p,q)=\ell=i(M)$, como queríamos.

Essa geodésica fechada tem a distância mínima
entre as geodésicas fechadas, já
que se tivéssemos alguma outra com distância
menor, podemos pegar um ponto
qualquer nela e o seu cut point ficaria a
distância exatamente a metade do
comprimento do laço (pois existem duas
geodésicas que chegam nele: cada
metade do laço partindo em direções opostas
do ponto inicial escolhido),
contradizendo o fato de que $i(M)=d(p,q)$.

Portanto, basta descartar o primeiro caso do
exercício anterior. Para chegar a
 uma contradição, suponha que $p $ é
 conjugado a $q$, i.e. que existe um campo
 de Jacobi $J$ que se anula em $p$ e em $q$.
Podemos usar o teorema de Rauch
\ref{differential-geometry-theorem-rauch-comparison} para comparar esse campo
com $\tilde{J}$,
um campo em $S^n_{K_{\text{max}}}$, a esfera
de curvatura constante
 $K_{\text{max}}$. Como $K \leq
 K_{\text{max}}$,
concluímos que $\tilde{J}$ deve se anular em
$\ell=d(p,q)=i(M)$. Absurdo, pois
estamos supondo que (1) não é verdadeiro,
i.e. que
 $i(M) < \pi/\sqrt{K_{\text{max}}}$.
\end{proof}

\begin{exercise}
\label{differential-geometry-exercise-l8-3}
\cite{doc}, Capítulo XIII, Proposição 3.4. Se
a curvatura seccional $K$ de uma
variedade Riemanniana $M^n$, compacta,
orientável e de dimensão par, satisfaz
$0<K \leq 1$, então $i(M) \geq \pi$.
\end{exercise}

\begin{proof}
Para obter uma contradição, suponha que
existe um ponto $p \in M$ tal que
 $d(p, C_m(p))<\pi$. Como $M$ é compacta,
 sabemos que $C_m(p)$ é compacto e
portanto existe uma geodésica $\gamma$
ligando $p$ com $q \in C_m(p)$ e
realizando a distância $d(p,C_m(p))$.

Pelo teorema de Rauch, sabemos que qualquer
geodésica não pode ter pontos
conjugados antes de alcançar comprimento
$\pi$. Explicitamente, se $J$ é um
campo de Jacobi ao longo de alguma geodésica
$\gamma$ tal que $J(0)=0$ e
$J(\ell(\gamma))=0$, comparando com um campo
$\tilde{J}$ em $S^n$ tal que
 $\tilde{J}(0)=0$, $|\tilde{J}'(0)|=|J'(0)|$
 e
$\left\langle{}J,\gamma'\right\rangle{}=\left\langle{}\tilde{J},
\tilde{\gamma}\right\rangle{}$,
concluímos que
 $|\tilde{J}|\leq |J|$ ao longo de $\gamma$.
 Como as geodésicas de $S^n$ não
 tem pontos conjugados antes de atingir
 comprimento $\pi$, concluímos que
$J$ não pode se anular antes desse ponto.

Isso mostra que, como estamos supondo que
$d(p,C_m(p))<\pi$, devemos estar no
caso (2) do Exercício \ref{differential-geometry-exercise-l8-1},
i.e. existem exatamente duas
geodésicas $\gamma$ e $\sigma$ ligando $p$ e
$q$, e
$\gamma'(\ell)=-\sigma'(\ell)$. Repetindo o
procedimento do Exercício anterior
 \ref{differential-geometry-exercise-l8-3}, podemos usar essas duas
 geodésicas para achar uma
geodésica fechada de comprimento $2i(M)$,
cujo comprimento é menor do que o
comprimento de qualquer outra geodésica
fechada.

\begin{remark}
\label{differential-geometry-remark-curva-posit-impli-maior-que}
Na prova do teorema de Synge
\ref{differential-geometry-theorem-Synge}, o Professor
Manfredo afirma que ``Como $M$ é compacta e
tem curvatura positiva, $K\geq
\delta>0$"; que não me parece imediato, pois
a curvatura seccional não é uma
função definida em $M$. Porém, não preciso
mostrar isso para garantir a
existência do loop geodésico (via o exercício
anterior); apenas é necessário
que $M$ seja compacta para garantir a
existência do ponto $q$.
\end{remark}

\medskip

Para concluir seguimos a prova de \cite{doc}.
Vamos a usar o exercício 6 da
lista 6, (cf.
\ref{differential-geometry-exercise-free-homotopy-variation}), onde
mostrei que, nestas condições, i.e.
$M$ de dimensão par, completa, orientável e
com curvatura seccional positiva,
existem curvas livremente homotópicas a
qualquer geodésica fechada $\gamma$, que
possuem comprimento menor que $\ell(\gamma)$.
Lembre que essas curvas foram
construídas mediante a função exponencial
aplicada a um campo paralelo ao longo
de $\gamma$, de forma que são curvas suaves;
vou usar esse fato depois.

A ideia é chegar numa contradição mostrando
que existe uma terceira geodésica
ligando $p$ e $q$. Essa geodésica se obtém
como o limite de uma sequência de
geodésicas associada à variação dada pelo
Exercício 6 da Lista 6.

Seja $c_s$ uma das curvas da variação com
comprimento estritamente
menor do que $\ell(\gamma)$. Pegue o ponto
inicial $c_s(0):=\tilde{p}_s$ e o
ponto $\tilde{q}_s$ em $c_s$ tal que
$d(\tilde{p}_s,\tilde{q}_s)$ é máxima ao
longo de $c_s$. Existe uma geodésica
$\tilde{\gamma}_s$ que liga
$\tilde{p}_s$ e $\tilde{q}_s$.

A seguir mostrarei que tomando limite quando
$s \to 0$, obteremos uma terceira
geodésica minimizante $\tilde{\gamma}_0$ que
liga $p$ e $q$, que não é possível
de novo pelo Exercício \ref{differential-geometry-exercise-l8-1}.

De fato, podemos dar explicitamente
$\tilde{\gamma}_0$ como sendo $exp_q(tw)$
onde $w$ é o vetor limite das velocidades
iniciais de cada $\tilde{\gamma}_s$,
supondo que elas são de tamanho 1 para
assegurar a existência do limite por
compacidade do fibrado tangente unitário.
Note que a geodésica obtida minimiza a
distância entre $p$ e $q$: a geodésica liga
$p$ com  $q$ porque, por um lado, o
limite dos pontos iniciais é $p$ por
definição, e, por outro lado, porque os
pontos finais são os pontos que maximizam a
distância dentro de cada geodésica.
A geodésica limite é minimizante por
continuidade da função distância.

Como cada curva $c_s$ é diferenciável, temos
um vetor tangente a ela em cada
ponto $\tilde{q}_s$. Como $\tilde{\gamma}$ é
minimizante, quando aplicamos a
fórmula da primeira variação,vemos que
$\tilde{\gamma}'_s$ é ortogonal a
$c_s'(\tilde{q}_s)$. Como o a métrica é
contínua, concluímos que
$\tilde{\gamma}_0'(q) \perp \gamma'(q)$.
\end{proof}

\begin{exercise}
\label{differential-geometry-exercise-l8-11}
Seja $M$ uma variedade Riemanniana completa,
de curvatura seccional não
negativa. Sejam $\gamma,\sigma:[0,\infty) \to
M$ geodésicas tais que
$\gamma(0)=\sigma(0)$. Se $\gamma$ é um raio
e
$\angle(\gamma'(0),\sigma'(0))<\frac{1}{2}
\pi$,
então $\lim_{t \to \infty}
\rho(\sigma(0),\sigma(t))=\infty$ (i.e.
$\sigma$ vai para o infinito).
\end{exercise}

\begin{proof}[Meu intento de prova]
Seja $t \in \mathbb{R}$. Considere uma
geodésica minimizante $\tau_t$ ligando
$\gamma(t)$ com $\sigma(t)$. Defina o angulo
$\beta_t:=\angle(-\gamma'(t),\tau'_t(0))$ e o
número $s_t$ como sendo o tempo em que
$\tau_t$ chega em $\sigma$, i.e.
$\tau_t(s_t)=\sigma(t)$. Ver Figura

Como $0 \leq  K$, podemos comparar a
``hinge'',
i.e. a informação do angulo e vetores
incidentes no ponto $\gamma(t)$,
$(-\gamma'(t),\tau_t'(0),\beta)$ com uma
hinge no espaço
euclidiano. Obtemos que
$$
\rho_{\mathbb{R}^n}(\tilde{\gamma}(0),
\tilde{\tau}_t(s_t))\leq
\rho(\gamma(0),\tau_t(s_t))=\rho(\sigma(0),
\sigma(t))
$$
Podemos calcular a distância
$\rho_{\mathbb{R}^n}(\tilde{\gamma}(0),
\tilde{\tau}_t(s_t))$
usando a lei de
cosenos euclidiana desde que conheçamos
$\beta_t$ e $\ell(\tau_t)$:
$$
\rho_{\mathbb{R}^n}(\tilde{\gamma}(0),
\tilde{\tau}_t(s_t))^2=
t^2+\ell(\tau_t)^2-2t\ell(\tau_t)\cos \beta_t
$$
Para calcular $\ell(\tau_t)$ podemos usar o
Teorema de Toponogov de novo, dessa
vez comparando a hinge $(\gamma'(0),
\sigma'(0),\alpha)$ com uma correspondente
no espaço euclidiano:
Obtemos que
$$
\rho_{\mathbb{R}^n}(\tilde{\gamma}(t),
\tilde{\sigma}(t))\leq
\ell(\tau_t)
$$
Como $\alpha$ é constante conforme $t$ muda,
podemos fixar $\tilde{\gamma}$ e
$\tilde{\sigma}$ para fazer a nossa
comparação. Conforme  $t$ avança, nos
movemos a velocidade constante ao longo de
$\tilde{\gamma}$ en direção ao
infinito, e embora não sabemos se a distância
$\rho(\sigma(0),\sigma(t))$ cresse
 ou diminue, a distância do nosso interesse
$\rho_{\mathbb{R}^n}(\tilde{\gamma}(t),
\tilde{\sigma}(t))$
vai sempre crescendo, e concluimos que
diverge ao infinito.

Portanto o problema acaba se conseguimos
achar uma cota inferior positiva para
$\beta_t$. Parece que isso é equivalente a
provar que a distância entre
$\gamma$ e $\sigma$ está acotada
inferiormente por uma constante positiva…

A lei de cosenos do lado Euclideano nos diz
que
$$
\rho_{\mathbb{R}^n}(\tilde{\gamma}(t),
\tilde{\sigma}(t))^2=
t^2+|\tilde{\sigma}(t)|^2-2t
|\tilde{\sigma}(t)|\cos \alpha
\leq t^2+|\tilde{\sigma}(t)|^2-2t
|\tilde{\sigma}(t)|
$$
Supondo que
$\tilde{\gamma}(0)=\tilde{\sigma}(0)=0 \in
\mathbb{R}^n$.
Parece que para medir essa distância
deveríamos conhecer
$|\tilde{\sigma}(t)|=\rho(\sigma(0),\sigma(t)
)$...
 o que parece levar o problema de volta ao
problema inicial.
\end{proof}

\begin{proof}[Prova de \cite{Cheeger-Ebin}]
Por desigualdade triangular, para quaisquer
$t,s \in \mathbb{R}$,
$$
\rho(\sigma(t),\sigma(0)) \geq
\rho(\gamma(s),\sigma(t))-\rho(\gamma(s),
\gamma(0)).
$$
Como o lado esquerdo não depende de $s$,
vemos que basta tomar o limite quando
$s\to \infty$ e mostrar que esse limite está
inferiormente limitado por um
número que depende linearmente de $t$, i.e.
$t \cos\alpha$.

Considere uma geodésica minimizante
$\tau_{s,t}$ ligando $\gamma(s)$ e
$\sigma(t)$. Podemos usar o teorema de
Toponogov \ref{differential-geometry-theorem-Toponogov-angle}
na versão de ângulos (que está enunciada em
\cite{Cheeger-Ebin}) já que duas das
geodésicas no triângulo
$\gamma(0),\gamma(s),\sigma(t)$ são
minimizantes por
hipótese. Como a cota inferior da curvatura é
0 não temos problemas com o
comprimento de $\sigma$.

Obtemos que existe um triângulo Euclidiano
com lados do mesmo comprimento tal
que os ângulos correspondentes são menores o
iguais que os ângulos com que
começamos.

Denote $A:=\rho(\gamma(0),\gamma(s))=s$,
$B:=\ell(\sigma(t))=t$ e
 $C:=\rho(\gamma(s),\sigma(t))=\ell(\tau_{s,
 t})$.
Nosso objetivo é limitar inferiormente $C-A$
conforme  $s \to \infty$.
 Aplicamos a lei de cosenos no lado
 Euclidiano para obter
$$
\begin{aligned}
C^2&=A^2+B^2-2AB\cos\tilde{\alpha}\\
\implies C^2-A^2&=B^2-2AB\cos\tilde{\alpha}\\
\implies
C-A&=\frac{B^2-2AB\cos\tilde{\alpha}}{C+A}
\end{aligned}
$$
Onde só $A=s$ e $C$ dependem de $s$. O
comportamento assintótico desse quociente
quando $s \to \infty$ é claro
geometricamente: como $A$ e $C$ são dois
lados de
um triângulo, eles crescem
proporcionadamente. Formalmente, como $C \leq
A+B$,
obtemos $C+A \leq 2A+B$. Daí
$$
\begin{aligned}
\frac{1}{C+A}&\geq \frac{1}{2A+B}\\
\implies \frac{B^2-2AB\cos
\tilde{\alpha}}{C+A}&\geq
 \frac{B^2-2AB \cos \tilde{\alpha}}{2A+B}
=\frac{A}{A}\frac{B^2/A-2B\cos
\tilde{\alpha}}{2+B/A}
\end{aligned}
$$
que (a menos de analizar o que acontece se o
numerador é negativo…)
 converge a $B\cos \tilde{\alpha}$ conforme
 $s \to \infty$.
 Concluímos que $C-A$ deve
convergir a
$-t\cos\tilde{\alpha}\leq-t\cos\alpha$.
Tirando o signo obtemos a
desigualdade desejada.
\end{proof}

\clearpage

% --- End differential-geometry.tex ---

% --- Begin symplectic-geometry.tex ---
\chapter{Symplectic geometry}
\label{chapter-symplectic-geometry}





\noindent
See
for the original notes of Henrique Bursztyn's
course.

\section{Symplectic vector spaces}
\label{symplectic-geometry-section-symplectic-vector-spaces}

\begin{proposition}
\label{symplectic-geometry-proposition-symplectic-basis}
Every symplectic vector space $(V,\omega)$
admits a {\it symplectic basis},
i.e. a basis $e_1,…,e_n,f_1,…,f_n$
such that
$$
\omega(e_i,e_j)=0,\qquad
\omega(e_i,f_j)=1,\qquad \omega(f_i,f_j)=0.
$$
\end{proposition}

\begin{proof}
Like Gram-Schmidt, start with any vector
$e_1$,
then there exists another vector $f_1$
such that $\omega(e_1,f_1)=1$.
The space generated by $e_1,f_1$
is symplectic, so we get a normal
decomposition…
\end{proof}

\begin{definition}
\label{symplectic-geometry-definition-symplectic-orthogonal}
Let $(V,\omega)$ be a symplectic vector
space and let $W \subseteq V$ be a
subspace. The {\it symplectic orthogonal}
of $W$ is
$$
W^\omega=\{v \in V \mid \omega(v,w)=0
\;\forall\, w \in W\}.
$$
The restriction $\omega|_W$ has kernel
$$
\ker(\omega|_W)=W \cap W^\omega.
$$
\end{definition}

\begin{lemma}
\label{symplectic-geometry-lemma-symplectic-orthogonal-dimension}
If $(V,\omega)$ is symplectic, then
$$
\dim W + \dim W^\omega = \dim V.
$$
\end{lemma}

\begin{proof}
The map
$$
\omega^\flat:V \longrightarrow V^*,\qquad
v \longmapsto \omega(v,\cdot)
$$
is an isomorphism. Under this map,
$W^\omega$ corresponds to $\operatorname{Ann}(W)$,
so
$$
\dim W^\omega=\dim \operatorname{Ann}(W)
=\dim V-\dim W.
$$
\end{proof}

\medskip\noindent
The linear algebra behind the dimension count
is just the following. If $W \subseteq V$,
then $\operatorname{Ann}(W)$ is the kernel of
the restriction map
$$
V^* \longrightarrow W^*,
\qquad
\varphi \longmapsto \varphi|_W.
$$
This map is surjective, since every linear
functional on $W$ extends to one on $V$ by
extending a basis of $W$ to a basis of $V$.
Hence rank-nullity gives
$$
\dim \operatorname{Ann}(W)
=
\dim V-\dim W.
$$
So, in the symplectic situation, once
$\omega^\flat:V\to V^*$ identifies
$W^\omega$ with $\operatorname{Ann}(W)$,
the formula above is immediate.

\begin{definition}
\label{symplectic-geometry-definition-symplectic-subspaces}
Let $W \subseteq V$ be a subspace of a
symplectic vector space $(V,\omega)$.
\begin{itemize}
\item $W$ is {\it isotropic} if
$$
\omega|_W=0,$$
equivalently if
$$
W\subseteq W^\omega.
$$

\item $W$ is {\it coisotropic} if
$$
W^\omega\subseteq W.
$$
Equivalently, the quotient
$W/W^\omega$ carries an induced symplectic
form.

\item $W$ is {\it Lagrangian} if $W$ is
maximal isotropic. Equivalently, $W$ is
isotropic and
$$
\dim W=\frac{1}{2}\dim V.
$$
Equivalently,
$$
W=W^\omega.
$$

\item $W$ is {\it symplectic} if
$\omega|_W$ is nondegenerate,
i.e. $(W,\omega|_W)$
is a symplectic vector space. Equivalently,
$$
W\cap W^\omega=\{0\}.
$$
\end{itemize}
\end{definition}

\medskip\noindent
To prove the equivalences
in the definitions above, the
it might be convenient to use 
Proposition 
\ref{symplectic-geometry-proposition-symplectic-basis}.
Write
$$
V=\langle e_1,f_1,\ldots,e_n,f_n\rangle
$$
with
$$
\omega(e_i,f_i)=1
$$
and all other pairings between basis vectors
zero, except those forced by skew-symmetry.
Each pair $\langle e_i,f_i\rangle$ should be
thought of as a symplectic plane.

\medskip\noindent
For example,
$$
W=\langle e_1\rangle\subset \mathbb R^4
$$
is isotropic. But
$$
W^\omega=\langle e_1,e_2,f_2\rangle,
$$
is not isotropic, since
$$
\omega(e_2,f_2)=1.
$$
Thus $W^\omega$ does not behave like an
ordinary Euclidean orthogonal complement.
It kills the symplectic partners of the
vectors in $W$, but it keeps the unused
symplectic planes.

\begin{exercise}
\label{symplectic-geometry-exercise-equivalent-definitions-symplectic-subspaces}
Prove that the equivalent formulations in
Definition
\ref{symplectic-geometry-definition-symplectic-subspaces}
are indeed equivalent. In particular, prove:
\begin{enumerate}
\item $\omega|_W=0$ if and only if
$W\subseteq W^\omega$.

\item If $W^\omega\subseteq W$, then
$\omega|_W$ descends to a symplectic form on
$W/W^\omega$.

\item A subspace is maximal isotropic if and
only if it is half-dimensional isotropic.

\item A subspace is Lagrangian if and only if
$W=W^\omega$.

\item The restriction $\omega|_W$ is
nondegenerate if and only if
$W\cap W^\omega=\{0\}$.
\end{enumerate}
\end{exercise}

\begin{remark}
If $W$ is symplectic, then
$$
V=W \oplus W^\omega.
$$
If $W$ is isotropic, then
$$
\dim W \leq \frac{1}{2}\dim V.
$$
If $W$ is coisotropic, then
$$
\dim W \geq \frac{1}{2}\dim V.
$$
\end{remark}

\begin{exercise}
\label{symplectic-geometry-exercise-symplectic-orthogonal}
Show that
$$
(W^\omega)^\omega = W
$$
and
$$
(W_1 \cap W_2)^\omega = W_1^\omega + W_2^\omega.
$$
\end{exercise}

\begin{example}
\leavevmode
\begin{itemize}
\item One-dimensional subspaces are
isotropic.
\item Codimension-one subspaces are
coisotropic.
\item If
$$
\{e_1,\ldots,e_n,f_1,\ldots,f_n\}
$$
is a symplectic basis of $V$, then
\(\operatorname{span}\{e_1,\ldots,e_k\}\)
is isotropic, and it is Lagrangian when
$k=n$.
\item If $(V_1,\omega_1)$ and
$(V_2,\omega_2)$ are symplectic vector
spaces and $T:V_1 \to V_2$ is linear, then
$\operatorname{graph}(T) \subseteq
V_1 \times \bar V_2$ is Lagrangian if and
only if $T$ is a symplectomorphism.
\item If $W$ is a vector space, then
$W \oplus W^*$ with
$$
\omega((w,\alpha),(w',\alpha'))
=\alpha'(w)-\alpha(w')
$$
is symplectic, and $W$ and $W^*$ are
Lagrangian subspaces.
\end{itemize}
\end{example}

\begin{exercise}
\label{symplectic-geometry-exercise-omega-power-n}
Let $V$ be a vector space and $\omega$ a
2-form on $V$.
Show that $\omega$ is nondegenerate if and
only if
$\omega^n$ is not zero.
\end{exercise}

\begin{proof}
\begin{reference}
\href{https://math.stackexchange.com/
questions/417365/
how-to-show-omegan-is-a-volume-form-in-a-
symplectic-manifold-m-omega}
{Stack Exchange}
\end{reference}
If $\omega$ is nondegenerate,
by Proposition
\ref{symplectic-geometry-proposition-symplectic-basis}
there is a symplectic basis
$e_1,…,e_n,f_1,…,f_n$
such that $\omega=\sum e^i \wedge f^i$.
Then
$$
\omega^n=n! e^1\wedge…\wedge e^n \wedge
f^1\wedge… \wedge f^n.
$$
Hence evaluating in
$e_1,…,e_n,f_1,…,f_n$
we see $\omega^n\neq 0$.

For the converse we suppose that $\omega$
is degenerate and prove that $\omega^n=0$.
It follows from the following property of
interior product:
$$
i_v\omega^n=n i_v\omega \wedge
\omega^{n-1}\qquad \forall  v \in V.
$$
Supposing that $\omega$ is degenerate we
have $v \in V$ such that $i_v\omega=0$,
which gives $i_v \omega^n=0$.
Completing $v$ to a basis $v_2,…,v_n$
we see that
$$
\omega^n(v,v_2,…,v_n)
=i_v\omega^n(v_2,…,v_n)=0.
$$
\end{proof}

\section{Hamiltonian vector fields, Poisson
brackets}
\label{symplectic-geometry-section-hamil-vecto-field-poiss}

\noindent
We know that a symplectic form on a manifold
$M$ gives us a map
$$
\begin{aligned}
\omega^\sharp: TM &\longrightarrow T^*M \\
X &\longmapsto i_X\omega
\end{aligned}
$$

\noindent
which says that for $f \in C^\infty(M)$,
$$
i_{X_f}\omega=df.
$$
\begin{definition}
\label{symplectic-geometry-definition-poisson-bracket}
$$
\begin{aligned}
\{\cdot,\cdot\}: C^\infty(M) \times
C^\infty(M) &\longrightarrow C^\infty(M) \\
\{f,g\}&=\omega(X_f,X_g)=df(X_g)=\mathcal{L}
_{X_f}g=-\mathcal{L}_{X_g}f
\end{aligned}
$$
is the {\it Poisson bracket}.
\end{definition}

\section{Symplectic structure of the
cotangent bundle}
\label{symplectic-geometry-section-sympl-struc-cotan-bundl}

\noindent
Upshot. A tangent vector to $T^*M$
records two infinitesimal motions: the base
point moves in $M$, and the covector moves in
the cotangent fiber. The tautological 1-form
ignores this second motion. It takes the
covector which is the point of $T^*M$ and
evaluates it on the motion of the base point
using the differential of the projection
map of the bundle.

\medskip\noindent
Let $\pi:T^*M \to M$ be the cotangent bundle.
The {\it tautological 1-form} $\lambda \in
\Lambda^1(T^*M)$
is defined by
$$
\lambda_\xi(v)=\xi(d\pi_\xi(v)),\qquad \xi
\in T^*M,\; v \in T_\xi T^*M.
$$
The canonical symplectic form on $T^*M$ is
$$
\omega_{\text{can}}=-d\lambda.
$$
If $(q^i)$ are coordinates on $M$
and $(q^i,p_i)$ are the induced coordinates
on $T^*M$, then
$$
\lambda=\sum_i p_i dq^i,\qquad
\omega_{\text{can}}=\sum_i dq^i \wedge dp_i.
$$

\begin{exercise}
\label{symplectic-geometry-exercise-cotan-bundl-canon-form}
Let $\alpha \in \Lambda^1(M)$ and view it as
a section
$\alpha:M \to T^*M$.
\begin{enumerate}
\item Show that $\alpha^*\lambda=\alpha$.
\item Conclude that
$\alpha^*\omega_{\text{can}}=-d\alpha$.
\item Show that the image $\alpha(M) \subset
T^*M$ is Lagrangian
if and only if $\alpha$ is closed.
\item In particular, show that the zero
section is Lagrangian.
\end{enumerate}
\end{exercise}

\begin{proof}
\begin{enumerate}
\item Pick $v \in TM$.
$$
\alpha^*\lambda(v)=\lambda(\alpha_*(v))
=\xi(\pi_*(\alpha_*(v))
$$
where $\alpha_*(v) \in T_\xi(M)$.
Now since $\alpha$ is a section,
$\pi \circ \alpha=\text{id}$,
and thinking about it, $\xi=\alpha$!
\end{enumerate}
\end{proof}

\begin{definition}
\label{symplectic-geometry-definition-cotangent-lift}
Let $\varphi:M \to N$ be a diffeomorphism.
The {\it cotangent lift} of $\varphi$ is
the map
$$
\widehat{\varphi}:T^*M \longrightarrow T^*N
$$
defined as follows.
If $\xi \in T_p^*M$, then
$\widehat{\varphi}(\xi)$ is the covector at
$\varphi(p)$ given by
$$
\widehat{\varphi}(\xi)(w)
=\xi(d\varphi_p^{-1}(w)),
\qquad w \in T_{\varphi(p)}N.
$$
\end{definition}

\noindent
In other words, the cotangent lift moves
the base point by $\varphi$ and transports
the covector by the inverse differential.

\begin{proposition}
\label{symplectic-geometry-proposition-cotangent-lift-symplectic}
The cotangent lift preserves the tautological
1-form:
$$
\widehat{\varphi}^*\lambda_N=\lambda_M.
$$
In particular,
$$
\widehat{\varphi}^*\omega_{\text{can},N}
=\omega_{\text{can},M},
$$
so cotangent lifts are symplectomorphisms.
\end{proposition}

\begin{proof}
Let $\xi \in T_p^*M$ and
$v \in T_\xi T^*M$.
The cotangent lift covers $\varphi$, so the
diagram
$$
\xymatrix{
T^*M \ar[r]^{\widehat{\varphi}}
\ar[d]_{\pi_M}
& T^*N \ar[d]^{\pi_N} \\
M \ar[r]_{\varphi} & N
}
$$
commutes.
By definition of the tautological 1-form,
$$
(\widehat{\varphi}^*\lambda_N)_\xi(v)
=\widehat{\varphi}(\xi)
\left(
d\pi_N(d\widehat{\varphi}_\xi(v))
\right).
$$
By commutativity of the diagram,
$$
d\pi_N(d\widehat{\varphi}_\xi(v))
=d\varphi_p(d\pi_M(v)).
$$
Applying the definition of
$\widehat{\varphi}(\xi)$,
we get
$$
(\widehat{\varphi}^*\lambda_N)_\xi(v)
=\xi(d\pi_M(v))=(\lambda_M)_\xi(v).
$$
The statement for $\omega_{\text{can}}$
follows by applying $-d$.
\end{proof}

\noindent
This is why cotangent lifts show up so often:
every diffeomorphism, or every action, on $M$
gives a canonical symplectic
diffeomorphism, or symplectic action,
on $T^*M$.
If $X$ is a vector field on $M$ and
$X_{T^*M}$ is the infinitesimal vector
field of the lifted flow on $T^*M$, then
$$
i_{X_{T^*M}}\omega_{\text{can}}
=d\left(\lambda(X_{T^*M})\right),
$$
so the cotangent lift is Hamiltonian with
Hamiltonian function
$$
H_X(\xi)=\xi(X_{\pi(\xi)}).
$$

\begin{exercise}
\label{symplectic-geometry-exercise-cotangent-lift-char}
Let $F:T^*M \to T^*M$ be a
diffeomorphism such that
$$
F^*\lambda=\lambda.
$$
Show that $F$ is the cotangent lift of a
diffeomorphism $\varphi:M \to M$.

\noindent
Hint: first show that $F$ preserves the
vertical distribution
$\ker d\pi \subset T(T^*M)$.
Then show that $F$ sends fibers of
$T^*M$ to fibers.
\end{exercise}

\section{Moser's trick}
\label{symplectic-geometry-section-moser-trick}

\noindent
Moser's trick is an application of a
technical formula concerning so-called
time-dependent vector fields, which
are 1-parameter families of vector fields.

\begin{lemma}
\label{symplectic-geometry-lemma-lie-deriv-time-depen-vecto}
\begin{reference}
\cite[Proposition 22.15]{les}
\end{reference}
Let $M$ be a smooth family and $X_t$ a
one-parameter family of vector fields
on $M$, i.e. a {\it time-dependent vector
field}.
Let $A_t$ be a family of tensors indexed by
the same parameter $t$.
Then
$$
\frac{d}{dt}\Big|_{t=t_0}\varphi^* _tA_t
=\varphi^*_t\left(\mathcal{L}_{X_t}A_t+
\frac{d}{dt}\Big|_{t=t_0}A_t\right).
$$
\end{lemma}

\begin{theorem}[Moser's trick]
\label{symplectic-geometry-theorem-moser-trick}
Let $\omega_t$ be cohomologous
symplectic forms on compact $M$,
$0\leq t\leq 1$.
There exists a diffeomorphism
$\varphi$ such that $\varphi^*\omega_1
=\omega_0$.
\end{theorem}

\begin{proof}
Since $\omega_t$ are cohomologous,
$$
0=\frac{d}{dt}[\omega_t]
=[\dot\omega_t]
$$
which says that $\dot\omega_t$ is exact,
say $\dot\omega_t=d\eta_t$.
Now define vector fields $X_t$
such that
$$
i_{X_t}\omega_t=-\eta_t.
$$
We use compactness of $M$
to make sure the flows of $X_t$
exist in $[0,1]$
The formula for 
derivative of the time-dependent flows reads
$$
\frac{d}{dt}\varphi_t^*\omega_t
=\varphi_t^* (L_{X_t}\omega_t+\dot \omega_t)
=-d\eta_t+d\eta_t
$$
since $L_{X_t}\omega_t = di_{X_t}\omega_t=
-d\eta_t$. At $t=0$ we get $\omega_0$,
and since the derivative is zero,
at $t=1$ we obtain 
$\varphi_1^*\omega_1=\omega_0$.
\end{proof}

\begin{theorem}[Moser's trick for volume
forms]
\label{symplectic-geometry-theorem-moser-trick-volume}
Let $M$ be a smooth manifold and $V_1, V_0$
two cohomologous volume forms.
Then there exists a diffeomorphism $\varphi:M
\to M$
such that $\varphi^*V_1=V_0$.
\end{theorem}

\begin{proof}
Suppose that $V_1-V_0=d\eta$ for some $\eta
\in \Lambda^{n-1}(M)$.
To construct a time-dependent vector field
consider the forms
$$
V_t:=tV_1+(1-t)V_0,\qquad t \in [0,1].
$$
Then we may associate a unique vector field
$X_t$ to $-\eta$
via the isomorphism $i_\bullet
V_t:\mathfrak{X}(M) \to \Lambda^{n-1}(M)$.
This is a time-dependent vector field,
so that by Lemma
\ref{symplectic-geometry-lemma-lie-deriv-time-depen-vecto}
we see that the pullback forms
$\varphi^*_tV_t$ are constant:

$$
\begin{aligned}
\frac{d}{dt}\Big|_{t=t_0}\varphi^*_tV_t
&=\varphi^*_t
\left(\mathcal{L}_{X_t}V_t
+\frac{d}{dt}\Big|_{t=t_0}V_t\right)\\
&=\varphi^*_t
\left(di_{X_t}V_t+i_{X_t}dV_t
+\frac{d}{dt}\Big|_{t=t_0}(tV_1+(1-t)V_0)
\right)\\
&=\varphi^*_t
\left(d(-\eta)
+V_1-V_0\right)\\
&=\varphi^*_t(-(V_1-V_0)+V_1-V_0)\\
&=0.
\end{aligned}
$$

\noindent
The results follows by taking $t=1$.
\end{proof}

\section{Darboux theorem}
\label{symplectic-geometry-section-darboux-theorem}

\begin{theorem}[Darboux]
\label{symplectic-geometry-theorem-darboux}
Let $(M,\omega)$ be a symplectic manifold and
$p \in M$.
There exist local coordinates
$x^1,…,x^n,y^1,…,y^n$ at $p$ such that
$$
\omega=\sum_{i=1}^n dx^i \wedge dy^i.
$$
\end{theorem}

\begin{proof}
Take any chart $(\varphi,U)$ at $p$.
It's enough to show that there is a
symplectomorphism
$\psi:(U,\varphi^*\omega)\to(U,
\omega_{\text{can}})$.

Since both $\varphi^*\omega$ and
$\omega_{\text{can}}$
are closed, so is their difference,
so that by Poincaré lemma, restricting the
neighbourhood if necessary,
we may assume that
$d\eta=\varphi^*\omega-\omega_{\text{can}}$
for some 1-form $\eta$.
The result follows by applying the exact same
procedure as in
Moser's trick as follows.
Let $\omega_t:=t
\varphi^*\omega+(1-t)\omega_{\text{can}}$.
Then $\omega_t$ is symplectic for all $t$,
which is not immediate:
in general the middle points of such 
a linear homotopy need not be nondegenerate.
In this case it works because we are working
locally: nondegeneracy is an open condition
in the sense that if the determinant
of a form is nonzero at a point,
it will be nonzero at near-by points.
Thus shrinking the domain if necessary
we can safely assume that $\omega_t$ is
nondegenerate.
Assign to every $t$ the vector field
$X_t:=\omega_t^\sharp(-\eta)$.
Then apply Lemma
\ref{symplectic-geometry-lemma-lie-deriv-time-depen-vecto}.
\end{proof}

\section{Weinstein's Lagrangian submanifold
theorem}
\label{symplectic-geometry-section-weins-lagra-subma}

\noindent
It was Weinstein who said that everything is
a Lagrangian submanifold.
The following theorem says that Lagrangian
submanifolds are
defined intrinsically in the sense that they
can all
be thought of as embedded in their own
cotangent bundle:
if $L \subset M$ is a Lagrangian submanifold,
there is are neighbourhoods $L \subseteq U_0
\subset M$
and $L \subset U_1 \subset T^*L$
and a symplectomorphism $\varphi:U_0 \to
U_1$.

Recall what it means to be a Lagrangian
vector subspace: $W \subset V$ is 
Lagrangian when $W = W^\omega$.
The key fact we need from
Lagrangian subspaces is this
(Exercise 
\ref{symplectic-geometry-exercise-lagrangian-complements}):
at the cost of choosing a metric
and making the construction non-canonical,
the symplectic orthogonal is in fact
a direct sum complement of $W$.
But more is true: the resulting
complement turns out to be Lagrangian too,
and it's identified with the dual space
of the original Langrangian subspace,
and, further, the splitting of the
original vector space has a canonical
symplectic form induced by this
decomposition. 

Showing those facts is the first step
towards Weinstein Lagrangian
submanifold theorem.
Then we need to show the corresponding
facts apply to submanifolds in a
fiberwise version of the exercise
(Lemma 
\ref{symplectic-geometry-lemma-lagrangian-complements-fiber}).
After that we use a generalized
version of Darboux theorem which essentially
is the statement of Weinstein's theorem.
The upshot is that the Lagrangian lemma
gives us the condition we need
to apply Generalized Darboux theorem.

\begin{exercise}[Lagrangian complements]
\label{symplectic-geometry-exercise-lagrangian-complements}
Let $(V,\omega)$ be a symplectic vector
space, let $W\subset V$ be a subspace, and
let $J$ be an $\omega$-compatible complex
structure. Let
$$
g(u,v)=\omega(u,Jv)
$$
be the associated inner product. Then
$$
J(W^\omega)=W^{\perp_g}.
$$
Consequently, every coisotropic subspace
has an isotropic complement. In particular,
every Lagrangian subspace $L\subset V$ has
a Lagrangian complement $L'$:
$$
V=L\oplus L'.
$$
Moreover, if $L'$ is any Lagrangian
complement of $L$, then there is a natural
identification
$$
L'\simeq L^*
$$
which induces a symplectomorphism
$$
V\simeq L\oplus L^*,
$$
where $L\oplus L^*$ is equipped with the
canonical symplectic form
$$
\omega_{\operatorname{can}}
((l,\alpha),(l',\alpha'))
=
\alpha(l')-\alpha'(l).
$$
\end{exercise}

\begin{proof}
The first part of the exercise
is straightforward from the definitions.
For the coisotropic subspace part
we need to use that $(W^\omega)^\omega=W$,
which follows from the dimension formula
of the symplectic orthogonal
(recall that while the symplectic
orthogonal does not give
direct sum decompositions,
at least the sum of the dimensions
add up to the dimension of the ambient space,
a key tool in these matters of symplectic
linear algebra; 
Lemma
\ref{symplectic-geometry-lemma-symplectic-orthogonal-dimension}).

\medskip\noindent
I'll leave this AI-part for now:
Now let $L'$ be any Lagrangian complement
of a Lagrangian subspace $L$, so
$$
V=L\oplus L'.
$$
Define
$$
\kappa:L'\longrightarrow L^*
$$
by
$$
\kappa(x)=\omega(x,-)|_L.
$$
This map is injective. Indeed, if
$\kappa(x)=0$, then
$$
\omega(x,l)=0
\qquad
\text{for all }l\in L.
$$
Hence $x\in L^\omega=L$. But also
$x\in L'$, and $L\cap L'=\{0\}$. Thus
$x=0$. Since
$$
\dim L'=\dim L^*,
$$
the map $\kappa$ is an isomorphism.

\medskip\noindent
Using $\kappa$, define
$$
F:V=L\oplus L'\longrightarrow L\oplus L^*
$$
by
$$
F(l+x)=(l,\kappa(x)).
$$
We now check that $F$ is symplectic. Since
both $L$ and $L'$ are Lagrangian,
$$
\omega(l+x,l'+x')
=
\omega(l,x')+\omega(x,l').
$$
On the other hand,
\begin{align*}
\omega_{\operatorname{can}}
(F(l+x),F(l'+x'))
&=
\omega_{\operatorname{can}}
((l,\kappa(x)),(l',\kappa(x')))\\
&=
\kappa(x)(l')-\kappa(x')(l)\\
&=
\omega(x,l')-\omega(x',l)\\
&=
\omega(x,l')+\omega(l,x').
\end{align*}
Thus
$$
F^*\omega_{\operatorname{can}}=\omega.
$$
So
$$
V\simeq L\oplus L^*
$$
as symplectic vector spaces.
\end{proof}


\begin{lemma}[Lagrangian complements
fiberwise]
\label{symplectic-geometry-lemma-lagrangian-complements-fiber}
Let $L\subset (M,\omega)$ be a
Lagrangian submanifold. Then there is an
isomorphism of symplectic vector bundles
over $L$
$$
TM|_L\simeq TL\oplus T^*L,
$$
where $TL\oplus T^*L$ is equipped with the
canonical symplectic form
$$
\Omega((X,\alpha),(Y,\beta))
=
\beta(X)-\alpha(Y).
$$
Moreover, this isomorphism restricts to
the identity on $TL$.
\end{lemma}

\begin{proof}
Since $L$ is Lagrangian, $TL$ is a
Lagrangian subbundle of the symplectic
vector bundle $TM|_L$. Choose a Lagrangian
complement $E$ to $TL$, so that
$$
TM|_L=TL\oplus E.
$$
Thus we have
$$
\xymatrix{
& TM|_L \ar[d]\\
& L
}
\qquad\text{with}\qquad
TM|_L=TL\oplus E.
$$

\medskip\noindent
The symplectic form identifies $E$ with
$T^*L$. Indeed, define
$$
b:E\longrightarrow T^*L
$$
by
$$
b(e)=\omega(e,-)|_{TL}.
$$
This map is injective: if $b(e)=0$, then
$e$ pairs trivially with $TL$. Since also
$E$ is Lagrangian, $e$ pairs trivially with
$E$. Hence $e$ pairs trivially with all of
$$
TM|_L=TL\oplus E,
$$
so $e=0$. Since $E$ and $T^*L$ have the
same rank, $b$ is an isomorphism.

\medskip\noindent
Using $b$, we identify
$$
TM|_L=TL\oplus E
\simeq
TL\oplus T^*L.
$$
Under this identification, the symplectic
form becomes
$$
\Omega((X,\alpha),(Y,\beta))
=
\beta(X)-\alpha(Y),
$$
up to the harmless sign convention in the
definition of $b$. Choosing the sign as
above gives the stated formula. The
isomorphism is the identity on $TL$.
\end{proof}

\noindent
The roadmap to the following generalization
of Darboux theorem is as follows.
First we show that if two symplectic
forms on the same manifold agree
along a submanifold $Q$,
then there are symplectomorphic
neighbourhoods of $Q$. That is,
there is a diffeomorphism
of a neighbourhood of $Q$ to another
neighbourhood of $Q$ taking one symplectic
form to another. Of course, this uses
Moser's trick, which needs a non-local
version of Poincaré lemma:

\begin{lemma}[Relative Poincaré]
\label{symplectic-geometry-lemma-relative-poincare}
Let $Q$ be an embedded submanifold of $M$
and $i:Q \hookrightarrow M$ the inclusion.
Let $\alpha \in \Omega^k(M)$ bet a closed
form such that $i^*\alpha=0$.
Then there is a neighbourhood
$U \supset Q$ and $\eta \in \Omega^{k-1}(U)$
such that $\alpha=d \eta$.
\end{lemma}

But the generalization of Darboux theorem
can be stated slightly more generally than 
that: we can allow the two symplectic 
forms to live in two different manifolds.
Then, instead of asking that the two
forms agree along the submanifold $Q$,
we simply require that there be
a bundle isomorphism of the pullback
bundles. In fact, we construct
an diffeomorphism of neighbourhoods
of $Q$ whose derivative is the given
bundle isomorphism.

Once again: the upshot about 
Lemma 
\ref{symplectic-geometry-lemma-lagrangian-complements-fiber}
is that it will provide us with the
bundle isomorphism we need to apply
Relative Darboux and obtain Weinstein's
Theorem.

\begin{theorem}[Relative Darboux theorem]
\label{symplectic-geometry-theorem-relative-darboux}
Let $(M_0,\omega_0)$ and $(M_1,\omega_1)$
be symplectic manifolds, and let
$$
i_0:Q\hookrightarrow M_0,
\qquad
i_1:Q\hookrightarrow M_1
$$
be embeddings. Suppose that there is an
isomorphism of symplectic vector bundles
$$
\Phi:TM_0|_Q\longrightarrow TM_1|_Q
$$
such that
$$
\Phi|_{TQ}=\operatorname{id}_{TQ}.
$$
In other words,
$$
\xymatrix{
TM_0|_Q \ar[rr]^-{\Phi} \ar[dr]
&& TM_1|_Q \ar[dl]\\
& Q &
}
$$
commutes and $\Phi$ is the identity on
$TQ$. Then there are neighborhoods $U_0$
of $Q$ in $M_0$ and $U_1$ of $Q$ in
$M_1$, and a symplectomorphism
$$
\varphi:(U_0,\omega_0)\longrightarrow
(U_1,\omega_1)
$$
such that
$$
\varphi|_Q=\operatorname{id}_Q.
$$
\end{theorem}

\begin{proof}
By the tubular neighborhood theorem, after
shrinking neighborhoods, the bundle map
$\Phi$ is the derivative along $Q$ of a
diffeomorphism
$$
\psi:U_0\longrightarrow U_1
$$
such that
$$
\psi|_Q=\operatorname{id}_Q,
\qquad
d\psi|_Q=\Phi.
$$
Thus
$$
\xymatrix{
U_0 \ar[rr]^-{\psi}
&& U_1\\
& Q \ar[ul] \ar[ur] &
}
$$
commutes.

\medskip\noindent
Since $\Phi$ is symplectic, the forms
$$
\omega_0
\qquad\text{and}\qquad
\psi^*\omega_1
$$
agree along $Q$. Put
$$
\eta=\psi^*\omega_1-\omega_0.
$$
Then $\eta$ is closed and vanishes along
$Q$. By the relative Poincare lemma, after
shrinking $U_0$, there is a $1$-form
$\sigma$ such that
$$
\eta=d\sigma,
\qquad
\sigma|_Q=0.
$$

\medskip\noindent
Consider
$$
\omega_t=\omega_0+t\eta.
$$
After shrinking $U_0$, all $\omega_t$ are
symplectic. Define $X_t$ by
$$
\iota_{X_t}\omega_t=-\sigma.
$$
Since $\sigma|_Q=0$, we have
$$
X_t|_Q=0.
$$
Thus the flow $\rho_t$ of $X_t$ fixes $Q$.

\medskip\noindent
Moser's computation gives
\begin{align*}
\frac{d}{dt}\rho_t^*\omega_t
&=
\rho_t^*
\left(
\mathcal L_{X_t}\omega_t+\dot\omega_t
\right)\\
&=
\rho_t^*
\left(
d(\iota_{X_t}\omega_t)+\eta
\right)\\
&=
\rho_t^*(-d\sigma+d\sigma)\\
&=0.
\end{align*}
Hence
$$
\rho_1^*(\psi^*\omega_1)=\omega_0.
$$
Therefore
$$
\varphi=\psi\circ \rho_1
$$
satisfies
$$
\varphi^*\omega_1=\omega_0.
$$
Since both $\psi$ and $\rho_1$ fix $Q$,
the map $\varphi$ fixes $Q$.
\end{proof}

\begin{theorem}[Weinstein's Lagrangian
neighborhood theorem]
\label{symplectic-geometry-theorem-weinstein-lagrangian}
Let $L\subset (M,\omega)$ be an embedded
Lagrangian submanifold. Then there are a
neighborhood $U_0$ of the zero section in
$T^*L$, a neighborhood $U_1$ of $L$ in
$M$, and a symplectomorphism
$$
\varphi:(U_0,\omega_{\operatorname{can}})
\longrightarrow (U_1,\omega)
$$
such that
$$
\varphi|_L=\operatorname{id}_L.
$$
Equivalently,
$$
\xymatrix{
T^*L \ar@{-->}[rr]^-{\varphi}
&& M\\
& L \ar[ul]^{0_L} \ar[ur]_{i} &
}
$$
near $L$.
\end{theorem}

\begin{proof}
Apply Theorem
\ref{symplectic-geometry-theorem-relative-darboux} to the two
embeddings
$$
\xymatrix{
T^*L && M\\
& L \ar[ul]^{0_L} \ar[ur]_{i} &
}
$$
Thus we take
$$
M_0=T^*L,
\qquad
M_1=M,
\qquad
Q=L.
$$

\medskip\noindent
Along the zero section of $T^*L$ there is
the canonical identification
$$
T(T^*L)|_L\simeq TL\oplus T^*L.
$$
The canonical symplectic form on this
bundle is
$$
\Omega((X,\alpha),(Y,\beta))
=
\beta(X)-\alpha(Y).
$$

\medskip\noindent
By Lemma
\ref{symplectic-geometry-lemma-lagrangian-complements-fiber},
there is also an isomorphism of symplectic
vector bundles
$$
TM|_L\simeq TL\oplus T^*L
$$
which is the identity on $TL$. Therefore
we get an isomorphism
$$
\Phi:T(T^*L)|_L\longrightarrow TM|_L
$$
of symplectic vector bundles, covering the
identity of $L$, and satisfying
$$
\Phi|_{TL}=\operatorname{id}_{TL}.
$$
In diagram form,
$$
\xymatrix{
T^*L && M\\
& L \ar[ul]^{0_L} \ar[ur]_{i} &
}
$$

\medskip\noindent
Now Theorem
\ref{symplectic-geometry-theorem-relative-darboux} gives
neighborhoods
$$
U_0\subset T^*L,
\qquad
U_1\subset M,
$$
and a symplectomorphism
$$
\varphi:(U_0,\omega_{\operatorname{can}})
\longrightarrow (U_1,\omega)
$$
such that
$$
\varphi|_L=\operatorname{id}_L.
$$
This is Weinstein's theorem.
\end{proof}

\medskip\noindent
Here's an AI-generated application
of this theorem concerning deformations
of a Lagrangian submanifold 
by Lagrangian submanifolds.

Let $L_t\subset M$ be a smooth family of
Lagrangian submanifolds with $L_0=L$.
Choose parametrizations
$$
\iota_t:L\to M
$$
with image $L_t$, and set
$$
V=
\frac{d}{dt}\bigg|_{t=0}\iota_t .
$$
Thus $V$ is a vector field along $L$.

\medskip\noindent
The tangential part of $V$ only
reparametrizes $L$. Hence the true
infinitesimal deformation is the class of
$V$ in
$$
TM|_L/TL.
$$
Since $L$ is Lagrangian, the symplectic
form gives a canonical identification
$$
TM|_L/TL\simeq T^*L.
$$
Explicitly, the class of $V$ is sent to
the $1$-form
$$
\alpha_V\in\Omega^1(L)
$$
defined by
$$
\alpha_V(X)=\omega(V,X),
\qquad X\in TL.
$$
This is well-defined because
$\omega|_L=0$.

\medskip\noindent
Now use the Lagrangian condition. Since
each $L_t$ is Lagrangian, we have
$$
\iota_t^*\omega=0
$$
for every $t$. Differentiating at $t=0$
gives
$$
0=
\frac{d}{dt}\bigg|_{t=0}\iota_t^*\omega
=
\iota_0^*(\mathcal L_V\omega).
$$
By Cartan's formula,
$$
\mathcal L_V\omega
=
d(\iota_V\omega)+\iota_Vd\omega.
$$
Since $d\omega=0$, this becomes
$$
0=d(\iota_0^*\iota_V\omega).
$$
But
$$
\iota_0^*\iota_V\omega=\alpha_V.
$$
Therefore
$$
d\alpha_V=0.
$$
Thus every infinitesimal Lagrangian
deformation gives a closed $1$-form.

\medskip\noindent
So far Weinstein's theorem has not been
used. The previous argument only gives a
necessary condition. Weinstein's theorem
gives the converse. It identifies a
neighborhood of $L$ in $M$ with a
neighborhood of the zero section in
$T^*L$, by a symplectomorphism
$$
\varphi:U_0\to U_1
$$
where
$$
U_0\subset T^*L,
\qquad
U_1\subset M.
$$
Therefore a sufficiently small $1$-form
$\alpha$ on $L$ determines a nearby
submanifold of $M$, namely
$$
\varphi(\Gamma_\alpha).
$$
Here ``sufficiently small'' means that
$$
\Gamma_\alpha\subset U_0.
$$

\medskip\noindent
In the cotangent model, the standard
calculation says
$$
\Gamma_\alpha\text{ is Lagrangian}
\quad\Longleftrightarrow\quad
d\alpha=0.
$$
Since $\varphi$ is a symplectomorphism,
this implies that every sufficiently small
closed $1$-form gives an actual nearby
Lagrangian submanifold of $M$. In
particular, if $\alpha$ is closed, then
$t\alpha$ gives such a deformation for
$|t|$ small.

\medskip\noindent
Finally, Hamiltonian-trivial deformations
correspond to exact $1$-forms. Indeed, if
the family is generated by a Hamiltonian
vector field $X_H$, then, up to sign
convention,
$$
\iota_{X_H}\omega=dH.
$$
Restricting to $L$ gives
$$
\alpha_{X_H}=d(H|_L).
$$
Thus Hamiltonian-trivial infinitesimal
deformations give exact $1$-forms.

\medskip\noindent
Consequently, modulo Hamiltonian-trivial
deformations, the infinitesimal
deformations of $L$ are measured by
$$
\frac{Z^1(L)}{B^1(L)}
=
H^1(L,\mathbb R).
$$
The role of Weinstein's theorem is to turn
the infinitesimal calculation into a local
classification: nearby Lagrangians are
locally the same thing as small closed
$1$-forms on $L$.

\medskip\noindent
Let us finish the section
discussing why Weinstein's theorem is a
useful local model in Misha's work.
(AI-written passage.)

\medskip\noindent
The theorem says that if
$L\subset (M,\omega)$ is Lagrangian, then a
neighborhood of $L$ in $M$ is symplectomorphic
to a neighborhood of the zero section in
$T^*L$. In other words, near $L$ there is no
extra ambient symplectic geometry:
the local model is forced by $L$ itself.
See \emph{Seminars}, Theorem
\ref{symplectic-geometry-theorem-weins-lagra-neigh}.

\medskip\noindent
This is exactly what is used in the
Lagrangian blow-down construction. If
$S\subset M$ is a Lagrangian sphere, then
Weinstein identifies a neighborhood of $S$
with a standard cotangent neighborhood.
One can then remove this standard piece and
glue in the corresponding orbifold piece.
This gives the Lagrangian blow-down; see
\emph{Seminars}, Definition
\ref{symplectic-geometry-definition-lagrangian-blow-down}.
The inverse operation is the Lagrangian
blow-up.

\medskip\noindent
Thus Weinstein is the local normal form
which makes the surgery possible. The global
arguments then use Moser's trick and Darboux
coordinates to show that the construction is
independent of choices; see
\emph{Seminars}, Theorems
\ref{symplectic-geometry-theorem-lagrangian-blow-up-well-defined}
and \ref{symplectic-geometry-theorem-same-for-blow-down}.
This is also why the Lagrangian
Teichmüller space appears naturally: it keeps
track of symplectic forms for which a fixed
submanifold remains Lagrangian.

\medskip\noindent
There is a parallel global story for
holomorphic Lagrangian fibrations. If
$(M,\Omega)$ is holomorphically symplectic
and
$$
\pi:M\to B
$$
is a holomorphic Lagrangian fibration, then
the fibers are half-dimensional and
$\Omega$ restricts to zero on them; see
\emph{Complex Geometry}, Definition
\ref{symplectic-geometry-definition-holomorphic-lagrangian-fib}.
By holomorphic Liouville, the smooth compact
fibers are complex tori. Hence a Lagrangian
fibration is a family of complex tori.

\medskip\noindent
In Misha's Shafarevich/isogeny picture, this
torus structure is the key. Degenerate
twistor deformations preserve the same
Lagrangian fibration; see
\emph{Complex Geometry}, Definition
\ref{symplectic-geometry-definition-degenerate-twistor-def}.
Different members of the family can then be
viewed as torsors over the same relative
torus fibration. After a finite base change,
these torsors acquire sections and become
trivial; see \emph{Complex Geometry},
Theorem \ref{symplectic-geometry-theorem-degenerate-twistor-cover}.
The graph of the resulting isomorphism
produces the desired symplectic
correspondence.

\medskip\noindent
So Weinstein's theorem explains the local
reason why Lagrangian submanifolds are
controlled by cotangent geometry:
$$
TM|_L/TL\simeq T^*L.
$$
Misha's work uses this local idea in global
ways: for Lagrangian blow-up and blow-down,
it gives the standard local model for
surgery; for Lagrangian fibrations, the
fibers become tori, and the global problem
turns into one about torsors, finite covers,
and correspondences.

\section{Coadjoint orbits and KKS theorem}
\label{symplectic-geometry-section-kks}

\noindent
There is an obvious way
of defining the coadjoint action
of a Lie group $G$ on the dual of its Lie
algebra, $\mathfrak{g}^*$. This gives
the coadjoint orbits:
$$
\mathcal{O}_\eta=
\{\text{Ad}^*_g(\eta):g \in G\}
$$
for $\eta \in \mathfrak{g}^*$.

As we show next,
the tangent space at any
point $\xi $of a coadjoint orbit is given by
infinitesimal vector fields, that is, 
$$
T_{\xi}\mathcal{O}=\{u_{\mathfrak{g}^*}(\xi):
u \in \mathfrak{g}\}.
$$
The reason for this expression
is essentially that the orbit map
$G \to \mathcal{O}$
is in general a submersion, 
and writing the definition
of the differential of this map
basically gives the definition
of the infinitesimal vector fields
$u_{\mathfrak{g}^*}$.

More precisely
By the general orbit theorem for smooth Lie
group actions
(Smooth Manifolds Theorem
\ref{symplectic-geometry-smooth-manifolds-theorem-orbit-theorem-lie-actions})
, each orbit is an immersed
submanifold and the orbit map
$$
\Phi_\xi:G\to \mathcal O_\xi,
\qquad
\Phi_\xi(g)=\operatorname{Ad}_g^*\xi
$$
is a submersion. Hence
$$
T_\xi\mathcal O_\xi
=
\operatorname{Im}(d\Phi_\xi)_e.
$$
But for $u\in\mathfrak g$,
$$
(d\Phi_\xi)_e(u)
=
\left.\frac{d}{dt}\right|_{t=0}
\operatorname{Ad}_{\exp(tu)}^*\xi
=
u_{\mathfrak g^*}(\xi).
$$

\noindent
\begin{definition}
\label{symplectic-geometry-definition-kks}
We define {\it Kostant--Kirillov--Souriau
(KKS) form} pointwise
on $T_\xi \mathcal{O}_\eta$ by
$$
\omega_\xi(u_{\mathfrak{g}^*}(\xi),
v_{\mathfrak{g}^*}(\xi)):=\xi([u,v])
=-u_{\mathfrak{g}^*}(\xi)v
=v_{\mathfrak{g}^*}(\xi)(u).
$$
\end{definition}

\begin{theorem}[KKS]
\label{symplectic-geometry-theorem-kks}
$\xi \mapsto  \omega_\xi$ is 
a well-defined 
(does not depend on the
choice of $u,v$) symplectic on
$\mathcal{O}$.
\end{theorem}

\begin{proof}
Non degeneracy is easy.
Then show that $\omega$ is $G$-invariant.
Closedness follows from Jacobi identity.
\end{proof}

\begin{exercise}
\label{symplectic-geometry-exercise-kks-so3-sphere}
Consider $G=\operatorname{SO}(3)$.
Identify $\mathfrak{so}(3)$ with
$\mathbb{R}^3$ by sending $u\in\mathbb{R}^3$
to the matrix $\hat u$ such that
$\hat u(v)=u\times v$.
Using the Euclidean inner product, identify
$\mathfrak{so}(3)^*$ with $\mathbb{R}^3$.

Show that the coadjoint orbits are the
spheres
$$
S_r^2=\{\xi\in\mathbb{R}^3:|\xi|=r\}.
$$
Then show that the KKS form on $S_r^2$ is
$$
\omega_\xi(u_{\mathfrak g^*}(\xi),
v_{\mathfrak g^*}(\xi))
=
\xi\cdot(u\times v),
$$
where
$$
u_{\mathfrak g^*}(\xi)=u\times \xi.
$$
Finally, prove that
$$
\omega=\frac{1}{r}\operatorname{Area}_{S_r^2}.
$$
In particular,
$$
\int_{S_r^2}\omega=4\pi r.
$$
\end{exercise}

\section{Hamiltonian actions, moment map
and symplectic reduction}
\label{symplectic-geometry-section-reduction}

\noindent
To warm-up we define real symplectic and
Hamiltonian actions.

\begin{definition}
\label{symplectic-geometry-definition-one-parameter-actions}
\begin{itemize}
\item An $\mathbb{R}$-action is {\it
symplectic}
if $\psi_t^*\omega=\omega
\iff L_X \omega=0$,
that is, $X$ is a symplectic field.
\item An $\mathbb{R}$-action is {\it
hamiltonian}
if $X$ is Hamiltonian,
that is, if there exists
a function $H \in C^\infty(M)$ such that
$X=X_H$.
\end{itemize}
\end{definition}

\noindent
To define moment map suppose
there exists a map $\hat{\mu}$
such that
$$
\xymatrix@C=3em@R=.7em{
&C^\infty(M)\ar[dd]&H\ar@{|->}[dd]\\
\mathfrak{g}\ar[dr]\ar[ur]^{\hat{\mu}}\\
u \ar@{|->}[dr] & \mathfrak{X}(M) & X_H\\
& u_M=X_u \ar@{=}[ur]
}
$$
Note that if  $\hat{\mu}$ exists, it's not
uniquely determined, and that we may suppose
that it is linear.

It's equivalent to define $\hat{\mu}$
as a map
$\hat{\mu}:\mathfrak{g}\to C^\infty(M)$
or as a map
$$
\begin{aligned}
\mu: M &\longrightarrow \mathfrak{g}^* \\
\hat{\mu}(u)(x) &= \left\langle{}\mu(x),u\right\rangle{}.
\end{aligned}
$$

\begin{definition}
\label{symplectic-geometry-definition-moment-map}
\begin{itemize}
\item A $G$-action is symplectic if
$\psi_g^*\omega=\omega$ for all $g \in G$,
that is, if
$$
\xymatrix{
\psi:G\ar[r]\ar[dr]&\text{Dif}(M)\\
&\text{Symp}(M,\omega)\ar@{^{(}->}[u].
}
$$

\item A symplectic action of $G$ on
$(M,\omega)$ via $\psi$
is {\it weakly Hamiltonian}
if there exists a map
$\mu:M \to \mathfrak{g}^*$
(or equivalently,
$\hat{\mu}:\mathfrak{g} \to C^\infty(M)$
linear) such that
$$
i_{u_M}\omega=d\left\langle{}\mu(p),u\right\rangle{},\qquad
(\text{resp. } u_M=X_{\hat{\mu}(u)}).
$$

\item A weakly Hamiltonian action is
{\it Hamiltonian} if it is also
$\mu$-equivariant for the action of $G$ on
$M$ and the coadjoint action on
$\mathfrak{g}^*$, that is,
$$
\mu \circ \psi_g=(\text{Ad}^*)_g(\mu).
$$
In this case we call $\mu$ a {\it moment
map}.
\end{itemize}
\end{definition}

\noindent
This first exercise shows how
the function $\hat{\mu} = H$
looks like for the simplest
action on a symplectic manifold:

\begin{exercise}
\label{symplectic-geometry-exercise-linear-moment-map-rotation}
Let $V=\mathbb{R}^2$ with symplectic form
$\omega=dx\wedge dy$.
Consider the action of $S^1$ on $V$
given by rotations,
$$
e^{it}(x,y)=
(x\cos t-y\sin t,x\sin t+y\cos t).
$$
Identifying $\mathfrak{s}^1=\mathbb{R}$,
let $u=1$.
\begin{enumerate}
\item Compute the infinitesimal vector field
$$
u_V(x,y)=
\frac{d}{dt}\Big|_{t=0} e^{it}(x,y).
$$
\item Compute the $1$-form $i_{u_V}\omega$.
\item Find a function $H:\mathbb{R}^2
\to \mathbb{R}$ such that
$$
dH=i_{u_V}\omega.
$$
\item Check that
$$
H(v)=\frac{1}{2}\omega(uv,v).
$$
Here we identify $T_vV$ with $V$, and $uv$
denotes the infinitesimal action of $u$ on
$v$.
\end{enumerate}
\end{exercise}

\begin{proof}
We have
$$
\begin{aligned}
u_V(x,y)
&=
\frac{d}{dt}\Big|_{t=0}
(x\cos t-y\sin t,
x\sin t+y\cos t)\\
&=(-y,x).
\end{aligned}
$$
Let $w=(a,b)$. Since
$$
\omega((p,q),(a,b))=pb-qa,
$$
we get
$$
\begin{aligned}
(i_{u_V}\omega)(w)
&=\omega((-y,x),(a,b))\\
&=-yb-xa.
\end{aligned}
$$
Thus
$$
i_{u_V}\omega=-x\,dx-y\,dy.
$$
Therefore
$$
H(x,y)=-\frac{1}{2}(x^2+y^2)
$$
satisfies
$$
dH=i_{u_V}\omega.
$$
Finally, we identify $T_vV$ with $V$. Since
$$
uv=(-y,x),\qquad v=(x,y),
$$
we get
$$
\begin{aligned}
\frac{1}{2}\omega(uv,v)
&=
\frac{1}{2}\omega((-y,x),(x,y))\\
&=
-\frac{1}{2}(x^2+y^2).
\end{aligned}
$$
Hence
$$
H(v)=\frac{1}{2}\omega(uv,v).
$$
\end{proof}

\noindent
The following exercise shows
how the moment map, i.e.
$\hat{\mu}=H$ looks for any
linear symplectic action.

\begin{exercise}
\label{symplectic-geometry-exercise-linear-moment-map}
Let $(V,\omega)$ be a symplectic vector space
and let $G$ act linearly on $V$ preserving
$\omega$.
Write
$$
u_M(v)=X_u(v)=uv=\frac{d}{dt}\Big|_{t=0}
\exp(tu)v
$$
for the infinitesimal action of $u$ on $v$.
We identify $T_vV$ with $V$ and,
for $u \in \mathfrak{g}$, define
$$
H_u(v)=\frac{1}{2}\omega(uv,v).
$$
Show that
$$
dH_u=i_{u_V}\omega,
$$
where $u_V(v)=uv$ is the infinitesimal
vector field associated to $u$.
\end{exercise}
\begin{proof}
Let $w \in V$. Since $V$ is a vector
space, we identify $T_vV$ with $V$.
Thus $w$ is a tangent vector at $v$,
and $dH_u(v)(w)$ is the directional
derivative of $H_u$ at $v$ in the
direction $w$. Therefore
$$
\begin{aligned}
dH_u(v)(w)
&=
\frac{d}{dt}\Big|_{t=0}
H_u(v+tw)\\
&=
\frac{1}{2}
\frac{d}{dt}\Big|_{t=0}
\omega(u(v+tw),v+tw)\\
&=
\frac{1}{2}\omega(uw,v)
+
\frac{1}{2}\omega(uv,w).
\end{aligned}
$$
Here we used that $\omega$ is bilinear:
differentiating a bilinear expression
means differentiating in the first slot
and in the second slot.

Since the action preserves $\omega$,
we have
$$
\omega(e^{tu}a,e^{tu}b)=\omega(a,b)
$$
for all $a,b \in V$.
Differentiating at $t=0$ gives
$$
\omega(ua,b)+\omega(a,ub)=0.
$$
Taking $a=w$ and $b=v$, we obtain
$$
\omega(uw,v)+\omega(w,uv)=0.
$$
Since $\omega$ is skew-symmetric, this
gives
$$
\omega(uw,v)=\omega(uv,w).
$$
Therefore
$$
dH_u(v)(w)=\omega(uv,w).
$$
But $u_V(v)=uv$, so
$$
dH_u(v)(w)
=
(i_{u_V}\omega)_v(w).
$$
Thus
$$
dH_u=i_{u_V}\omega.
$$
\end{proof}

\begin{exercise}
\label{symplectic-geometry-exercise-tangent-moment-level-set}
Let $(M,\omega)$ be a symplectic manifold
and let $G$ act on $M$ with moment map
$\mu:M\to\mathfrak{g}^*$.
Assume that $0$ is a regular value of
$\mu$.
Show that, for every $p \in \mu^{-1}(0)$,
we have
$$
T_p\mu^{-1}(0)
=
(T_p(Gp))^\omega.
$$
Here $(T_p(Gp))^\omega$ denotes the
symplectic orthogonal of the tangent
space to the orbit.
\end{exercise}

\begin{proof}
For the inclusion $\subseteq$
we need to show that $\omega(v,w)=0$
for any $v \in T_p\mu^{-1}(0)$
and $w \in T_p(G\cdot p)$.
A $w$ looks like: take a 1-parameter
subgroup of $G$, which of course
we want to see as something of the kind
$\text{exp}(tu)\cdot p$ and differentiate
at $t=0$. And of course this is the
infinitesimal generator $u_M = X_u$
for $u \in \mathfrak{g}$.
On the other hand, a $v$ looks like:
take a curve $\gamma \subset \mu^{-1}(0)$
and differentiate at $t=0$.

By the moment map property we see that
$$
\omega(u_M,v)=d\left\langle{}\mu,u\right\rangle{}_p(v).
$$
How to compute that differential?
I remember that the differential of a
function can be computed by
using a curve that realises the vector,
i.e. $df(v) = \frac{d}{dt}\Big|_{t=0}
f\circ \gamma(t)$.
That is,
$$
d\left\langle{}\mu(p),u\right\rangle{}v
=\frac{d}{dt}\Big|_{t=0}
\left\langle{}\mu(\gamma(t)),u\right\rangle{}=0.
$$

For the $\supseteq$ inclusion
suppose that $v$ is such that
$\omega(u_M,v)=0$ for all
$u \in \mathfrak{g}$.
To show that $v \in T_p\mu^{-1}(0)$
we need to show that
$v \in \Ker d\mu$ by the
level set theorem.
This is what the moment map equation says:
$0=i_{u_M}\omega(v)
=d\left\langle{}\mu,u\right\rangle{}_p(v)$,
and the latter differential is in fact
$d(\text{ev}_u \circ \mu)$.
Since $\operatorname{ev}_u$ is linear, its
differential is again $\operatorname{ev}_u$.
Hence
$$
d\left\langle{}\mu,u\right\rangle{}_p(v)
=
\left\langle{}d\mu_p(v),u\right\rangle{}.
$$
Since this vanishes for every
$u \in \mathfrak{g}$, the covector
$d\mu_p(v)\in\mathfrak{g}^*$ vanishes
identically. Thus $d\mu_p(v)=0$.
\end{proof}

\noindent
The following definition
and exercise are the
key to symplectic reduction:

\begin{definition}
\label{symplectic-geometry-definition-basic-form}
Let $M$ be a manifold and $\omega \in
\Omega^k(M)$.
Suppose that $\pi:M \to B$ is a surjective
submersion with connected fibers.
We say that $\omega$ is {\it basic} (with
respect
to $\pi$) is there exists a form
\(\overline{\omega}\in \Omega^k(B)\) such
that
\(\pi^*\overline{\omega}\).
\end{definition}

\begin{exercise}
\label{symplectic-geometry-exercise-basic-form}
\begin{enumerate}
\item Show that $\omega$ is basic if and only
if
$i_X \omega=0$ and $L_X\omega=0$
for all vector fields $X$ tangent
to the fibers of $\pi$.
In particular, if $\omega$ is closed,
show that it is basic if
\(\Ker(\pi_*) \subseteq \Ker(\omega)\)
(pointwise in $M$).

\item
\label{symplectic-geometry-item-reduction-form}
Suppose that $\omega$ is a closed
2-form on $M$ and
\(\Ker(\pi_*)=\Ker(\omega)\).
Show that $\omega=\pi^*\overline{\omega}$
and \(\overline{\omega}\in\Omega^2(B)\)
is symplectic.

\item (Application to reduction.)
Let $(M,\omega)$ be a symplectic
manifold and 
$\iota:N\hookrightarrow M$
a submanifold such that
$D=TN \cap TN^\omega \subset TN$
has constant rank
(e.g. $N$ could be coisotropic).
We saw in class that $D$ is an
integrable distribution
(by Frobenius);
suppose that the leafspace
$B:=N/\sim$ is smooth so that the
natural projection
$\pi:N \to B$ is a submersion.
Show that $B$ inherits a unique symplectic
form $\omega_{\text{red}}$
with the property that
$\pi^*\omega_{\text{red}}=\iota^*\omega$.
\end{enumerate}
\end{exercise}

\begin{proof}
\begin{enumerate}
\item The idea is to define a form 
downstairs by the equation
$$
\overline{\omega}(\bar{X}_1,\dots,\bar{X}_k)
:=\omega(X_1,\dots,X_k)
$$
where every $X_i$ is a lift of $\bar{X}_i$.
To construct such lifts the first thing
is to note that since $\pi$ is a surjective
submersion we immediately get a 
rough vector field on the total space.
But it's not immediate that it is smooth.
To construct a smooth lift we can use
local submersion charts for $\pi$,
and just lift the vector field in
the base by simply putting the same
local coordinate functions in the first
slots upstairs. This is clearly a smooth
vector field on the local chart for
the upstairs manifold.
Then we use a partition of unity to extend
it.

To check that $\bar{\omega}$ is well
defined suppose that $X$ and $X'$
are lifts of $\bar{X}$. Then their
difference must vertical. Since $i_Z\omega=0$
for vertical $Z$, we see that
$\omega(X-X',\cdot)=0$ which says
that $\bar{\omega}$ has the same value
for either lift.

Observe that we have defined a map
$$
\bar{\omega} : 
\mathfrak{X}(B) \times \mathfrak{X}(B)
\longrightarrow C^\infty(M),
$$
i.e. with the wrong codomain.
We need to prove that the function
$\bar{\omega}(\bar{X},\bar{Y})$
is constant along the fibers,
for which we use
the condition that $L_Z\omega =0$
for vertical $Z$.
Given any such $Z$, $L_Z\omega=0$
says that $\varphi^*_t\omega=\omega$
for all $t$, where $\varphi$ is
the flow of $Z$.
In other words,
$\omega(\varphi_{*,t}X,\varphi_{*,t}Y)
=\omega(X,Y)$
for any $X,Y \in \mathfrak{X}(M)$.
This is all we need:
given any pair of points $p,p'$ 
in some fiber $\pi^{-1}(b)$,
there exists a vector field $Z \in\ker\pi_*$
whose flow maps $p$ to $p'$.
More precisely,
since the fibers are connected, 
any two points
in a fiber can be joined by 
a path inside the
fiber. Covering this path by finitely many
submersion charts, we see that the two points
are connected by a finite composition 
of flows of vertical vector fields.

We have constructed a smooth map
$\bar{\omega}:\mathfrak{X}(B)\times
\mathfrak{X}(B) \to C^\infty(B)$.
It is also $C^\infty(B)$-linear
(i.e. tensorial) since
any smooth function on $B$ can be
extended to a smooth function on $M$
using submersion charts, and then
we use tensoriality of $\omega$.

It follows that $\bar{\omega}$ is
a smooth differential form on $B$,
and $\pi^* \bar{\omega}=\omega$
by construction:
taking submersion charts,
the action of $\pi_*$ is the identity
on the nonvanishing coordinates.

\item If $\omega$ is closed,
then it's clear that $\bar{\omega}$
will be closed too since exterior
differential commutes with pullback.
It's also immediate that
$\bar{\omega}$ is nondegenerate
because by passing to the quotient
and assuming $\Ker\pi_*=\Ker\omega$
we have literally thrown away
all the vectors in the kernel.

\item Pulling back $\omega$ from
$M$ to $N$ we obtain a closed form
$\iota^*\omega$ which need not be
nondegenerate.
However, the fibers of the leaf-space
projection $\pi$ are by definition
the vectors in $\Ker \omega$.
Thus we are in the situation
of Item (\ref{symplectic-geometry-item-reduction-form}).

For uniqueness suppose
that $\omega_{\text{red}}' 
\in \Omega^2(B)$
also satisfies 
$\pi^*\omega_{\text{red}}'=\iota^*\omega$.
Then
$\pi^*(\omega_\text{red}-\omega_\text{red}')
=0$,
but since $\pi$ is a surjective
submersion, the pullback map $\pi^*$
is injective and thus
$\omega_\text{red}=\omega_\text{red}'$.
\end{enumerate}
\end{proof}

\begin{exercise}
\label{symplectic-geometry-exercise-reduced-symplectic-form}
Let $(M,\omega)$ be a symplectic manifold
with a Hamiltonian action of $G$ and
moment map $\mu:M\to\mathfrak{g}^*$.
Assume that $0$ is a regular value of
$\mu$ and that $G$ acts freely and
properly on $\mu^{-1}(0)$.

Let
$$
Z=\mu^{-1}(0)
$$
and let
$$
\pi:Z\longrightarrow Z/G
$$
be the quotient map.
Show that there exists a unique $2$-form
$\omega_{\mathrm{red}}$ on $Z/G$ such that
$$
\pi^*\omega_{\mathrm{red}}
=
\iota^*\omega,
$$
where $\iota:Z\hookrightarrow M$ is the
inclusion.
\end{exercise}

\begin{proof}
By Exercise
\ref{symplectic-geometry-exercise-basic-form},
we know that all we have to do is check that
$\Ker \pi_* \subset \Ker \iota^*\omega$,
that is, that every vertical vector field
$Z$ gives the zero 1-form when contracted
with $\omega$, i.e. $i_Z\omega=0$.

We know that every vertical vector field
is locally given as $u_Z$ for some
$u \in \mathfrak{g}$.
By the moment map equation,
we shall obtain that $i_{u_Z}\omega=0$
if we show that $\left\langle{}\mu(\cdot),u\right\rangle{}$
is a constant function on $Z$,
which is clear since $Z= \mu^{-1}(0)$.
\end{proof}

\noindent
Note that the original intention of the
exercise above was to use its
preceding exercise (but of course
my solution is also correct).

\begin{exercise}
\label{symplectic-geometry-exercise-reduced-form-is-symplectic}
In the setting of Exercise
\ref{symplectic-geometry-exercise-reduced-symplectic-form},
show that the reduced form
$\omega_{\mathrm{red}}$ is symplectic.
That is, show that
$$
d\omega_{\mathrm{red}}=0
$$
and that $\omega_{\mathrm{red}}$ is
nondegenerate.
\end{exercise}

\begin{proof}
Closedness is immediate
by commutation of pullback and
exterior differential since $\omega$
is closed.

For nondegeneracy,
by Exercise \ref{symplectic-geometry-exercise-basic-form}
we just need to check the inclusion
$\Ker \iota^*\omega \subset \Ker \pi_*$.
In words, if a vector $v$ is in the
kernel of $\omega$, it has to be vertical.
By Exercise
\ref{symplectic-geometry-exercise-tangent-moment-level-set}
we know that
$$
T_pZ = T_p(G\cdot p)^\omega
$$
for any $p \in Z$.

Taking symplectic orthogonals
we get
$$
(T_pZ)^\omega=T_p(G\cdot p).
$$
Now just observe that the vertical
vectors are, by definition 
since $\pi$ is the projection of
a group quotient, the vectors
tangent to the orbits.
\end{proof}

\noindent
We have just proved

\begin{theorem}[Marsden-Weinstein, Megre]
\label{symplectic-geometry-theorem-marsden-weinstein-megre}
Let $G$ act on $(M,\omega)$ in a Hamiltonian
way
with moment map $\mu:M \to \mathfrak{g}^*$.
Suppose $0 \in \mathfrak{g}^*$ is a regular
value.
By the equivariance of $\mu$, we have an
action of $G$ on $\mu^{-1}(0)$
given as follows. For $x \in \mu^{-1}(0)$,
$\mu(\psi_g(x))=\text{Ad}_g^*(\mu(x))=0$.
Suppose that this action is regular (this
will be defined later),
that $\mu^{-1}(0)/G$ is smooth, and that the
projection is a submersion.
Then there exists a unique symplectic form
$\omega_{\text{red}}$ in $\mu^{-1}(0)/G$
such that
$$
\iota^*\omega=\pi^*\omega_{\text{red}}.
$$
\end{theorem}

\noindent
That is, the pullback of the symplectic form
to the
0 level-set of $\mu$ passes to the quotient:
$$
\xymatrix{
Z=\mu^{-1}(0)\ar[d]^{\pi}\ar@{^{(}->}[r]
^{\iota}&
(M,\omega)\\
(\mu^{-1}(0)/G,\omega_{\text{red}})
}
$$

\noindent
A slightly more general setting is the
following.
Hamiltonian action of $G$ on $(M,\omega)$,
and moment map
$\mu:M \to \mathfrak{g}^*$.
$\xi \in \mathfrak{g}^*$ regular, $G_\xi
\subseteq G$
stabilizer of $\xi$.
Then $G_\xi$ acts on the level set
$\mu^{-1}(\xi)$, so we obtain:
$$
\xymatrix{
\mu^{-1}(\xi)\ar[d]^{\pi}\ar@{^{(}->}[r]
^-{\iota}&
(M,\omega)\\
(\mu^{-1}(\xi)/G_\xi,\omega_{\text{red}
}).
}
$$



\medskip\noindent
Now we pass to applications.
The next exercise shows the
construction of $\mathbb{CP}^{n-1}$ from
$\mathbb{C}^n$.

\begin{exercise}
\label{symplectic-geometry-exercise-projective-space-as-reduction}
Let $\mathbb{C}^n$ have its standard
symplectic form
$$
\omega_0=\frac{i}{2}\sum_{j=1}^n
dz_j\wedge d\bar z_j.
$$
Consider the diagonal $S^1$-action
$$
e^{it}\cdot(z_1,\dots,z_n)
=
(e^{it}z_1,\dots,e^{it}z_n).
$$
\begin{enumerate}
\item Show that this action is Hamiltonian
with moment map
$$
\mu(z)=
-\frac{1}{2}\sum_{j=1}^n |z_j|^2.
$$

\item For $r>0$, identify
$$
\mu^{-1}\left(-\frac{r^2}{2}\right)
$$
with the sphere $S^{2n-1}_r\subset
\mathbb{C}^n$.

\item Show that
$$
\mu^{-1}\left(-\frac{r^2}{2}\right)/S^1
\simeq \mathbb{P}^{n-1}_{\mathbb{C}}.
$$

\item Conclude that $\mathbb{P}^{n-1}$
inherits a symplectic form from reduction.
This is the Fubini--Study form, up to scale.
\end{enumerate}
\end{exercise}


\noindent
The next exercise is the 
Hamiltonian-mechanics example.

\begin{exercise}
\label{symplectic-geometry-exercise-so3-action-cotangent}
Consider the group $\text{SO}(3)$ acting
on $T^*\mathbb{R}^{3}$ by the cotangent
lift of the usual action of
$\text{SO}(3)$ on $\mathbb{R}^{3}$.
\begin{enumerate}
\item For $u\in\mathfrak{so}(3)$, compute
the corresponding infinitesimal generator
$u_{T^*\mathbb{R}^3}\in\mathfrak{X}
(T^*\mathbb{R}^{3})$.
\item Identify $\mathfrak{so}(3)$ with
$\mathbb{R}^3$ as in Lista 5. Show that,
with this identification, we have
$u_{T^*\mathbb{R}^3}(q,p)
=(u\times q,u\times p)$.
\item Identifying $\mathfrak{so}(3)^*
\cong(\mathbb{R}^{3})^*
\cong\mathbb{R}^{3}$ using the usual
inner product, show that the moment map
for the action of $\text{SO}(3)$ on
$T^*\mathbb{R}^{3}$ is $\mu(q,p)=q\times
p$. Conclude (by Noether's theorem) that
if $V\in \mathcal{C}^\infty(\mathbb{R}^3)$
is $\text{SO}(3)$-invariant, then the
flow of the Hamiltonian $H(q,p)=\frac{p^2}
{2m}+V(q)$ preserves ``angular momentum''
$q \times p$.
\end{enumerate}
\end{exercise}

\begin{proof}
\begin{enumerate}
\item To survive without getting too
involved in the definitions, I observe
the following. $\text{SO}(3)$ acts
by rotations on $\mathbb{R}^3$.
This means that the tangent vector
to any 1-parameter orbit is a tangent
vector to a sphere containing the
1-parameter orbit.
In other words, we have an expression
for the 1-parameter orbit
$p \cos t + u_{\mathbb{R}^3} \sin t$
where $u_{\mathbb{R}^3}$
is the $\mathbb{R}^3$ component
of $u_{T^*\mathbb{R}^3}$ at
$p \in \mathbb{R}^3$.
This just says that $u_{\mathbb{R}^3}$
is orthogonal to $p$.
Since the induced map on the cotangent
space is essentially given by the
pushforward of the action, which is linear,
then it's exactly the same ---
thus the $T^*\mathbb{R}^3$
component of $u_{T^*\mathbb{R}^3}$
must also be orthogonal to
the covector where it is based.

Now let's get too involved in the
definitions.
By definition, for
$(p,q) \in T^*\mathbb{R}^3$
and $u \in \mathfrak{so}_3$,
$$
u_{T^*\mathbb{R}^3}
=\frac{d}{dt}\Big|_{t=0}
e^{tu}(p,q).
$$
To find the integral curve $\varphi(t)
=e^{tu}$ recall how it is constructed:
the vector $u \in T_e G$
gives the left-invariant vector field
$X(g)=d_eL_g(u)$, which
has an integral curve at the identity
$\varphi_u(t)$ defined by the properties
$\varphi_u'(0)=u$ and $\varphi_u(0)=e$.
We denote $\varphi_u(t)=e^{tu}$.
Thus differentiating the integral
curve by definition gives
the left-invariant vector field $X$
along $\varphi_u$,
which is the differential of a linear
map and thus it acts by itself,
that is $X(\varphi(t))
=d_eL_{\varphi_u(t)}(u)
=\varphi_u(t)u$.

Differentiating and evaluating at $t=0$
we see that $u_{T^*\mathbb{R}^3}
=(up,uq)$
since, in fact, the action on the
cotangent part is also given by
simple matrix multiplication.

More explicitly, let $A\in \text{SO}(3)$.
The cotangent lift is encoded by
$$
\xymatrix{
T^*\mathbb{R}^3 \ar[r]^{\widehat A}
\ar[d]_{\pi}
& T^*\mathbb{R}^3 \ar[d]^{\pi}\\
\mathbb{R}^3 \ar[r]_{A}
& \mathbb{R}^3
}
$$
and sends $(q,\xi)$ to $(Aq,A\cdot \xi)$,
where
$$
(A\cdot \xi)(v)=\xi(A^{-1}v).
$$
Use the Euclidean inner product to
identify covectors with vectors. Thus
write
$$
\xi(v)=\langle p,v\rangle.
$$
Here $p$ is the vector representing
$\xi$ by the Riesz representation
theorem.
By the definition of adjoint from
Definition \ref{linear-algebra-definition-adjoint},
we have
$$
\langle A^{-1}v,p\rangle
=\langle v,(A^{-1})^\dag p\rangle.
$$
Therefore
$$
\begin{aligned}
(A\cdot \xi)(v)
&=\xi(A^{-1}v)\\
&=\langle p,A^{-1}v\rangle\\
&=\langle A^{-1}v,p\rangle\\
&=\langle v,(A^{-1})^\dag p\rangle\\
&=\langle (A^{-1})^\dag p,v\rangle.
\end{aligned}
$$
Thus the covector $A\cdot \xi$ is
represented by $(A^{-1})^\dag p$.
Since $A\in\text{SO}(3)$, we have
$A^\dag=A^{-1}$, and hence
$(A^{-1})^\dag=A$. Therefore, after this
identification,
$$
A\cdot(q,p)=(Aq,Ap).
$$
It follows that, for
$u\in\mathfrak{so}(3)$,
$$
u_{T^*\mathbb{R}^3}(q,p)
=
\frac{d}{dt}\Big|_{t=0}
(e^{tu}q,e^{tu}p)
=(uq,up).
$$
\item We now identify $u$ with a vector
and compute $Uq$. Write
$$
U=
\begin{pmatrix}
0&-c&b\\
c&0&-a\\
-b&a&0
\end{pmatrix}
\longleftrightarrow (a,b,c).
$$
If $q=(q_1,q_2,q_3)$, then
\begin{align*}
Uq
&=
\begin{pmatrix}
0&-c&b\\
c&0&-a\\
-b&a&0
\end{pmatrix}
\begin{pmatrix}
q_1\\ q_2\\ q_3
\end{pmatrix}\\
&=
\begin{pmatrix}
-cq_2+bq_3\\
cq_1-aq_3\\
-bq_1+aq_2
\end{pmatrix}.
\end{align*}
These are the coordinates of
$$
(a,b,c)\times(q_1,q_2,q_3).
$$
Thus $Uq=u\times q$. The same computation
for $p$ gives $Up=u\times p$. Therefore
$$
u_{T^*\mathbb{R}^3}(q,p)
=(u\times q,u\times p).
$$

\item There is of course the tour-de-force
path in which we do everything in
coordinates and just confirm that
the moment map and equivariance equations
hold. But let's not do that.

I see two other ways: (1) try to use somehow
that the symplectic form in $T^*\mathbb{R}^3$
is the differential of the tautological
form, and (2) try to use that everything
is linear.

Let's do (2). Recall that $\mu$
takes values in dual Lie algebra.
Thus the vector $q\times p$ we are given
is to be understood as a traceless
anitymmetric matrix. The map
$T^*\mathbb{R}^3 \ni (q,p) 
\overset{\hat{\mu}}{\mapsto } 
(q\times p)u$
is most likely linear.
So the exterior differential is…
probably itself!
\end{enumerate}
\end{proof}

\section{Noether's principle}
\label{symplectic-geometry-section-Noether-principle}

\begin{theorem}[Noether's principle]
\label{symplectic-geometry-theorem-Noether-principle}
Let $(M,\omega,\mu)$ be a Hamiltonian
$G$-space with Hamiltonian
$H\in C^\infty(M)$. $H$ is $G$-invariant if
and only if $\mu$ is preserved by
the Hamiltonian flow.
\end{theorem}

\begin{proof}
one line
\end{proof}


\section{Toric varieties}
\label{symplectic-geometry-section-toric-varieties}

\noindent
For any compact connected symplectic manifold
$(M^{2n},\omega)$ with an action
of a torus $\mathbb{T}^n$ and moment map
$\mu:M \to \mathbb{R}^n
=(\mathfrak{t}^n)^*$, the image of $\mu$
is a polytope. By [Guilleimin-Sternberg] and
[Atiyah], this polytope is
convex. Further, if the dimension of the
torus is exactly half
the dimension of the manifold, by [Delzant],
the moment polytope
determines the manifold, i.e. there is a
one-to-one
correspondence between
Delzant polytopes and symplectic toric
varieties.
Moduli space of polygons is an example of
this,
see \cite{moduli-polygons}.

In fact, this situation of having a torus
action of half the dimension
is what we may call a ``totally integrable
system''.

\section{(Completely) integrable systems}
\label{symplectic-geometry-section-completely-integrable-systems}

\noindent
Let $(M,\omega)$ be a symplectic manifold.
Recall that

\begin{itemize}
\item
$\{f,g\}=\omega(X_g,X_f)=L_{X_f}g=-L_{X_g}f$.
\item $X_{\{f,g\}}=[X_f,X_g]$
\end{itemize}

\noindent
If $\{f,g\}=0$ then the Hamiltonian flows of
$f$ and $g$ commute.

\begin{definition}
\label{symplectic-geometry-definition-compl-integ-syste}
Let $(M^{2n},\omega)$ and $H \in
C^\infty(M)$.
$(M,\omega,H)$ is a {\it completely
integrable system}
(in Liouville's sense)
if there exist functions $f_1,\ldots,f_n \in
C^\infty(M)$ such that
\begin{enumerate}
\item $H=f_1,\ldots,f_n$ are linearly
independent
($df_1,\ldots,df_n$ are linearly independent
in every point).

\item $\{ f_i, f_j\}=0$ for all $i,j$.
\end{enumerate}
\end{definition}

\noindent
Note that $F=(f_1,\ldots,f_n):M \to
\mathbb{R}^n$ is a submersion.
For all $c \in \mathbb{R}^n$, $F^{-1}(c)$
are Lagrangian submanifolds.
$F^{-1}(c)$ is invariant by the flow of
$X_{f_i}$ for $i=1,\ldots,n$.

Explanation: $F^{-1}(c)$ is a level-set
submanifold of $M$,
whose tangent space is generated by the
$df_i$.
These are all in the kernel of $dF$ by the
integrability condition,
namely $dF=(df_1,\ldots,df_n)$,
$df_i(X_{f_j})=\{f_i,f_j\}=0$.

The idea is that finding these ``conserved
quantities''
we manage to describe the Hamiltonian
dynamics completely.
Showing that a system is integrable is a
victory.
The moduli space of polygons is an example,
and I think that so is
$(\mathbb{C}^2)^{[n]}$.

There is another method of using
``bihamiltonian actions''
(two Hamiltonian actions yielding the same
dynamics)
/Lax pairs (defining the system using
eigenvalues of matrices).
Examples of this are Toda, Caolengo-Moser.

\medskip\noindent
The following theorem says that the connected
component of a symplectic manifold with an
integrable system


\begin{theorem}[Arnold-Liouville]
\label{symplectic-geometry-theorem-arnold-liouville}
Let $(M^{2n},\omega)$, $H=f_1,\ldots,f_n$ be
completely intgrable system
and $F:M \to \mathbb{R}^n$,
$F=(f_1,\ldots,f_n)$.
Let $M_c$ be the connected component of
$F^{-1}(c)$.
Then
\begin{enumerate}
\item
\label{symplectic-geometry-item-toric1}
If $M_c$ is compact, then $M_c \simeq
\mathbb{T}^n=\mathbb{R}^n/2\pi\mathbb{Z}^n=
\{(\theta_1,\ldots,\theta_n)
\text{
mod }2\pi\}$.

\item
\label{symplectic-geometry-item-toric2}
In angular coordinates, the Hamiltonian flow
is linear, that is,
$\frac{d}{dt}^{-1}(t)=\alpha \in
\mathbb{R}^n$
implies that $\theta(t)=\theta_0+\alpha t$.
\end{enumerate}
\end{theorem}

The proof of item \ref{symplectic-geometry-item-toric1} is given
by
the following proposition
putting $N=M$ and
$X_1=X_{f_1},\ldots,X_n=X_{f_n}$.

\begin{proposition}
\label{symplectic-geometry-proposition-torus}
Let $N$ be complete, connected, of dimension
$n$.
Let $X_1,\ldots,X_n$ be linearly independent
vector fields on $N$
such that $[X_i,X_j]=0$ for all $j$.
Then $N \simeq \mathbb{T}^n$.
\end{proposition}

\begin{proof}
First we present the steps:
\begin{enumerate}
\item Flows of $X_1,\ldots,X_n$ define an
action of $\mathbb{R}^n$ on $N$.

\item This action is transitive. Then $N
\simeq \mathbb{R}^n/\Gamma$
where $\Gamma$ is the stabilizer.

\item Since $\dim N=n$, then  $\Gamma$ is a
lattice of $\mathbb{R}^n$.

\item Since $N$ is compact, then $\Gamma$ is
of maximum rank, that is,
$\Gamma=2\pi \mathbb{Z}^n \subset
\mathbb{R}^n$ .
\end{enumerate}


Let $\varphi^i_t$ be the flow of $X_i$ (which
are complete since $N$ is
compact).
Since $[X_i,X_j]=0$, then $\varphi_t^i
\varphi_s^j=\varphi_s^j \varphi^i_t$.
Thus we have an action of $\mathbb{R}^n$ on
$N$
given by moving along the flows of the vector
fields, that is
$$
\underbrace{(t_1,\ldots,t_n)}_{\mathbf{t}}
\cdot
x
=\underbrace{\varphi_{t_1}^1 \circ \ldots
\circ \varphi_{t_n}^n(x)}_{
\varphi_{\mathbf{t}}(x)}
$$
To prove transitivity
we can show that the orbit maps of this
action are surjective.
Indeed, the image of the orbit map is open
since it is a local
diffeomorphism (it's derivative is given by
the vector fields,
which are linearly independent!),
and it is closed by a taking a point in the
boundary and
using the flow at this point (and
commutativity of the flows).

To check that the stabilizer $\Gamma$ of this
action is trivial
first note that all stabilizers are equal
(not only conjugate) since $\mathbb{R}^n$ is
abelian.
Further, since the orbit map is a local
diffeomorphism,
$\Gamma$ is discrete. Thus, it must be a
lattice of rank $k \leq n$.

To show that in fact $k=n$ consider a linear
map  $T$
that maps the generators of $\Gamma$ to the
first $k$ canonical
vectors of $\mathbb{R}^n$. By the following
diagram, the arrow on the bottom
is well defined as in fact a bijective local
diffeomorphism,
that is, a diffeomorphism.
$$
\xymatrix{
\mathbb{R}^k \times\mathbb{R}^{n-k} \simeq
\mathbb{R}^n\ar[r]^{T}\ar[d]_{p}
\ar[dr]^{\psi_x \circ T}&
\mathbb{R}^n\ar[d]^{\psi_x=
\substack{\text{orbit}
\\ \text{map}}}\\
\mathbb{T}^k \times
\mathbb{R}^{n-k}\ar[r]_{\overline{T}}&
N
}
$$
Thus the factor $\mathbb{R}^{n-k}$ must be a
point since $N$ is
compact by hypothesis.
\end{proof}

\noindent
Item \ref{symplectic-geometry-item-toric2} from Theorem
\ref{symplectic-geometry-theorem-arnold-liouville}
is obtained by considering the flows
on the latter diagram and is left as an
exercise.

\noindent
In the case that $N$ is not compact we just
conclude that
it $M_c$ is a product of a torus with an
Euclidean space.

Perhaps the simplest example of an integrable
system
is given by taking an open ball about the
origin $B \subseteq \mathbb{R}^n$
and $B \times
\mathbb{T}^n=\{(I_1,\ldots,I_n,\phi_1,\ldots,
\phi_n)\}$
with the symplectic form $\omega=\sum_{i=1}^n
dI_1 \wedge d\phi_i$
with $\{I_1,I_j\}=0$.

The following theorem says that in fact the
{\bf symplectic form} of every
level set of the kind $M_c$, which we just
proved
that is a torus, may is described locally as
such a product.

\begin{theorem}
\label{symplectic-geometry-theorem-local}
There is a neighbourhood $U$ of $M_c$ and
a symplectomorphism
$$
\xymatrix{
\sum d \mathcal{I} \wedge d\phi_i&\ar@{~>}[l]
B \times
\mathbb{T}^n\ar[r]^{r}_{\simeq}\ar[d]_{p}&U
\subseteq M\ar[d]^{F}
\ar@{~>}[r]
& \omega\\
&B\ar[r]_{\simeq}&F(U)
}
$$
\end{theorem}

\noindent
The coordinates
$(I_1,\ldots,I_n,\phi_1,\ldots,\phi_n)$
are called {\it action-angle} coordinates.

The case when $h(I)$ does not depend on the
angle coordinates $\phi_i$
we can solve the Hamiltonian equations
easily:
$$
\begin{cases}
\dot I_i=0 \\
\dot \phi_i=-\frac{\partial h}{\partial I}(I)
\end{cases}
$$
so that $\phi_i(t)=\phi_0+\alpha t$ for some
constants $\alpha$.
This is why integrable systems are actually
easy to solve!

\section{Vertex algebras}
\label{symplectic-geometry-section-vertex-algebras}

Mônadas, álgebras sobre mônadas. Exemplo:
ultrafiltros (Manes, 1976). Campo
quântico, como serie de Fourier. Exemplo:
Heisenberg algebra.

Problema: o producto de campos quânticos
não
é bem definido.


% --- End symplectic-geometry.tex ---

% --- Begin geometric-prequantization.tex ---
\chapter{Geometric prequantization}
\label{chapter-geometric-prequantization}





\section{Overview}
\label{geometric-prequantization-section-overview}

\noindent
These notes are about the first step in
geometric quantization.  The starting point
is
a symplectic manifold $(M,\omega)$,
regarded as
the phase space of a classical mechanical
system.  The goal is to attach to it a
Hilbert
space and a rule sending classical
observables
$f \in C^\infty(M)$ to operators on that
Hilbert space.

\medskip\noindent
The construction treated here is only
prequantization.  It gives the correct formal
commutator rule, but the resulting Hilbert
space is usually too large.  The later step,
not treated here, is to impose a
polarization.

\medskip\noindent
We use the following convention.  The
Hamiltonian vector field of $f$ is defined by
$$
\iota_{X_f}\omega=-df,
$$
and the Poisson bracket is
$$
\{f,g\}=-\omega(X_f,X_g).
$$
With this convention $X_f(g)=-\{f,g\}$.

\section{The Dirac rule}
\label{geometric-prequantization-section-dirac-rule}

\noindent
In classical mechanics an observable is a
smooth function on the phase space.  The
Hamiltonian flow of $H$ evolves another
observable $f$ by
$$
\frac{d}{dt}f=\{f,H\}.
$$
Thus the Poisson bracket is the classical
operation governing the time evolution of
observables.

\medskip\noindent
In quantum mechanics observables are
self-adjoint operators on a Hilbert space
$\mathcal H$.  If $H_q$ is the quantum
Hamiltonian, the Schrodinger equation is
$$
i\hbar\frac{d\psi}{dt}=H_q\psi.
$$
For a time-independent observable $A$, the
expected value in the state $\psi$ is
$$
\left\langle{} A \right\rangle{}_\psi
=\left\langle{} A\psi,\psi \right\rangle{}.
$$
The commutator controls its evolution:
$$
\frac{d}{dt}
\left\langle{} A \right\rangle{}_\psi
=
\frac{1}{i\hbar}
\left\langle{} [A,H_q]\psi,\psi
\right\rangle{}.
$$
This is the basic reason the commutator is
the quantum analogue of the Poisson bracket.

\begin{definition}
\label{geometric-prequantization-definition-dirac-prequantization}
Let $(M,\omega)$ be a symplectic manifold.
A {\it prequantization} of $M$ consists of a
Hilbert space $\mathcal H$ and a linear map
$$
Q:C^\infty(M,\mathbb R)\longrightarrow
\operatorname{End}(\mathcal H)
$$
such that
\begin{enumerate}
\item $Q(1)=\operatorname{Id}$, and
\item
$$
Q(\{f,g\})=\frac{1}{i\hbar}[Q(f),Q(g)]
$$
for all $f,g\in C^\infty(M,\mathbb R)$.
\end{enumerate}
\end{definition}

\noindent
A full quantization should also include a
minimality condition saying that a complete
set of commuting classical observables goes
to a complete set of commuting quantum
observables.  This is where
prequantization is
not yet enough.

\section{Hermitian line bundles}
\label{geometric-prequantization-section-hermitian-line-bundles}

\begin{definition}
\label{geometric-prequantization-definition-hermitian-line-bundle}
A {\it Hermitian line bundle} over $M$ is a
complex line bundle $L\to M$ together with a
smooth Hermitian product $h$ on the fibers.
Thus $h(s,t)$ is a smooth complex-valued
function whenever $s,t$ are smooth sections.
\end{definition}

\noindent
We write
$$
\Omega^k(M,L)=
\Gamma(M,\Lambda^kT^*M\otimes L)
$$
for the space of $L$-valued differential
forms.

\begin{definition}
\label{geometric-prequantization-definition-connection-line-bundle}
A {\it connection} on $L$ is a linear map
$$
\nabla:\Gamma(M,L)\longrightarrow
\Omega^1(M,L)
$$
such that
$$
\nabla(fs)=df\otimes s+f\nabla s
$$
for all $f\in C^\infty(M)$ and all sections
$s$ of $L$.
For a vector field $X$ we write
$$
\nabla_Xs=\iota_X\nabla s.
$$
\end{definition}

\begin{definition}
\label{geometric-prequantization-definition-unitary-connection}
A connection $\nabla$ on a Hermitian line
bundle $(L,h)$ is {\it unitary} if
$$
dh(s,t)=h(\nabla s,t)+h(s,\nabla t)
$$
for all smooth sections $s,t$.
\end{definition}

\noindent
A connection extends to $L$-valued forms by
$$
\nabla(\alpha\otimes s)
=d\alpha\otimes s-\alpha\wedge \nabla s.
$$
Then $\nabla^2$ is $C^\infty(M)$-linear, so
there is a complex 2-form $\Omega_\nabla$
such that
$$
\nabla^2s=\Omega_\nabla s.
$$
This form is the curvature of $\nabla$.

\begin{lemma}
\label{geometric-prequantization-lemma-local-curvature-line-bundle}
Let $e$ be a nowhere vanishing local section
of $L$ and suppose
$$
\nabla e=\theta\otimes e.
$$
Then, on the same open set,
$$
\Omega_\nabla=d\theta.
$$
If $e$ is unitary and $\nabla$ is unitary,
then $\theta$ and $\Omega_\nabla$ are purely
imaginary forms.
\end{lemma}

\begin{proof}
For vector fields $X,Y$ we have
$$
\begin{aligned}
\nabla_X\nabla_Ye
&=\nabla_X(\theta(Y)e)\\
&=X(\theta(Y))e+	heta(Y)\theta(X)e.
\end{aligned}
$$
Subtracting the same expression with $X$ and
$Y$ interchanged, and subtracting
$\nabla_{[X,Y]}e$, gives
$$
\Omega_\nabla(X,Y)e
=d\theta(X,Y)e.
$$
Hence $\Omega_\nabla=d\theta$.

If $e$ is unitary, then $h(e,e)=1$.  Since
$\nabla$ is unitary,
$$
0=dh(e,e)=h(\nabla e,e)+h(e,\nabla e)
=\theta+\overline\theta.
$$
Thus $\theta$ is purely imaginary, and so is
$d\theta$.
\end{proof}

\begin{definition}
\label{geometric-prequantization-definition-first-chern-class}
The {\it first Chern class} of a Hermitian
line bundle with unitary connection is
$$
c_1(L)=\left[\frac{1}{2\pi i}
\Omega_\nabla\right]\in H^2(M,\mathbb R).
$$
\end{definition}

\noindent
Chern--Weil theory says this class is
independent of the chosen unitary connection.
Moreover it is integral, i.e. it lies in the
image of $H^2(M,\mathbb Z)$ in real
cohomology.

\begin{theorem}[Weil integrality theorem]
\label{geometric-prequantization-theorem-weil-integrality}
\begin{reference}
See \cite[Lecture 12]{wang} and
\cite[Section 8.3]{wood}.
\end{reference}
Let $\beta$ be a closed real 2-form on $M$
whose de Rham cohomology class is integral.
Then there exists a Hermitian line bundle
$L\to M$ with unitary connection $\nabla$
such that
$$
\Omega_\nabla=2\pi i\beta.
$$
Equivalently, $c_1(L)=[\beta]$.
\end{theorem}

\begin{proof}
We omit the proof.  It is the standard
Cech--de Rham construction of a line bundle
from an integral cocycle, followed by the
construction of a compatible unitary
connection.
\end{proof}

\section{Prequantum line bundles}
\label{geometric-prequantization-section-prequantum-line-bundles}

\begin{definition}
\label{geometric-prequantization-definition-prequantizable}
A symplectic manifold $(M,\omega)$ is
{\it prequantizable} if
$$
\left[\frac{\omega}{2\pi\hbar}\right]
\in H^2(M,\mathbb Z).
$$
A {\it prequantum line bundle} is a
Hermitian line bundle $(L,h)$ with a unitary
connection $\nabla$ whose curvature is
$$
\Omega_\nabla=\frac{\omega}{i\hbar}.
$$
\end{definition}

\noindent
The sign in the integrality condition is
irrelevant: if a class is integral, so is its
negative.  By Theorem
\ref{geometric-prequantization-theorem-weil-integrality}, the
integrality condition is exactly what is
needed for the existence of a prequantum line
bundle.

\medskip\noindent
When $M$ is compact, the prequantum Hilbert
space is the $L^2$-completion of the smooth
sections of $L$ with respect to
$$
\left\langle{} s,t\right\rangle{}
=
\frac{1}{(2\pi\hbar)^n}
\int_M h(s,t)\frac{\omega^n}{n!}.
$$
The measure $\omega^n/n!$ is the Liouville
measure.

\begin{definition}
\label{geometric-prequantization-definition-prequantum-operator}
Let $f\in C^\infty(M,\mathbb R)$.
The {\it prequantum operator} associated to
$f$ is
$$
Q(f)=-i\hbar\nabla_{X_f}+m_f,
$$
where $m_f$ denotes multiplication by $f$.
\end{definition}

\section{The Kostant--Souriau formula}
\label{geometric-prequantization-section-kostant-souriau-formula}

\begin{lemma}
\label{geometric-prequantization-lemma-hamiltonian-volume-preserving}
For every $f\in C^\infty(M)$,
$$
\mathcal L_{X_f}\omega=0
$$
and hence
$$
\mathcal L_{X_f}\omega^n=0.
$$
\end{lemma}

\begin{proof}
By Cartan's formula and the definition of
$X_f$,
$$
\mathcal L_{X_f}\omega
=d\iota_{X_f}\omega+
\iota_{X_f}d\omega
=-d^2f+0=0.
$$
The statement for $\omega^n$ follows by the
Leibniz rule for the Lie derivative.
\end{proof}

\begin{lemma}
\label{geometric-prequantization-lemma-integral-poisson-bracket-zero}
Assume that $M$ is compact without boundary.
Then, for all $f,g\in C^\infty(M)$,
$$
\int_M \{f,g\}\frac{\omega^n}{n!}=0.
$$
\end{lemma}

\begin{proof}
Since $X_f(g)=-\{f,g\}$ and
$\mathcal L_{X_f}\omega^n=0$, we have
$$
\{f,g\}\omega^n
=-\mathcal L_{X_f}(g\omega^n).
$$
By Cartan's formula,
$$
\mathcal L_{X_f}(g\omega^n)
=d\iota_{X_f}(g\omega^n)+
\iota_{X_f}d(g\omega^n).
$$
The second term vanishes for degree reasons,
because $g\omega^n$ is a top-degree form.
Thus $\{f,g\}\omega^n$ is exact, and the
integral vanishes by Stokes' theorem.
\end{proof}

\begin{proposition}
\label{geometric-prequantization-proposition-prequantum-self-adjoint}
Assume that $M$ is compact without boundary.
For real $f$, the operator $Q(f)$ is formally
self-adjoint on smooth sections.
\end{proposition}

\begin{proof}
The multiplication operator $m_f$ is formally
self-adjoint because $f$ is real.  It remains
to check $-i\hbar\nabla_{X_f}$.
By compatibility of $\nabla$ with $h$,
$$
X_fh(s,t)=h(\nabla_{X_f}s,t)
+h(s,\nabla_{X_f}t).
$$
Integrating and using Lemma
\ref{geometric-prequantization-lemma-integral-poisson-bracket-zero} for
real and imaginary parts of $h(s,t)$, we get
$$
\int_M X_fh(s,t)\frac{\omega^n}{n!}=0.
$$
Hence
$$
\int_M h(\nabla_{X_f}s,t)\frac{\omega^n}{n!}
=-
\int_M h(s,\nabla_{X_f}t)\frac{\omega^n}{n!}.
$$
Thus $\nabla_{X_f}$ is formally skew-adjoint,
and $-i\hbar\nabla_{X_f}$ is formally
self-adjoint.
\end{proof}

\begin{theorem}[Kostant--Souriau]
\label{geometric-prequantization-theorem-kostant-souriau}
\begin{reference}
See \cite[Lecture 12]{wang} and
\cite[Chapter 23]{hallq}.
\end{reference}
Let $(M,\omega)$ be prequantizable, and let
$(L,h,\nabla)$ be a prequantum line bundle.
Then
$$
Q(\{f,g\})=\frac{1}{i\hbar}[Q(f),Q(g)]
$$
for all $f,g\in C^\infty(M,\mathbb R)$.
Thus $f\mapsto Q(f)$ is a prequantization in
the sense of Definition
\ref{geometric-prequantization-definition-dirac-prequantization}.
\end{theorem}

\begin{proof}
Let $h=\{f,g\}$.  First,
$$
{}[X_f,X_g]=-X_h.
$$
Indeed,
$$
\iota_{[X_f,X_g]}\omega
=\mathcal L_{X_f}\iota_{X_g}\omega
=-\mathcal L_{X_f}dg
=-d(X_fg)=dh,
$$
because $X_f(g)=-h$.  Since
$\iota_{X_h}\omega=-dh$, this gives
$[X_f,X_g]=-X_h$.

Now compute the commutator on smooth
sections.  We use
$$
{}[\nabla_X,m_g]=m_{X(g)}
$$
and the curvature identity
$$
{}[\nabla_X,\nabla_Y]
=\nabla_{[X,Y]}+\Omega_\nabla(X,Y).
$$
Therefore
$$
\begin{aligned}
{}[Q(f),Q(g)]
&=(-i\hbar)^2[\nabla_{X_f},\nabla_{X_g}]\\
&\quad -i\hbar m_{X_f(g)}
+i\hbar m_{X_g(f)}.
\end{aligned}
$$
Since $X_f(g)=-h$ and $X_g(f)=h$, the last
two terms give $2i\hbar m_h$.  For the
first term we use the prequantum curvature:
$$
\begin{aligned}
{}[\nabla_{X_f},\nabla_{X_g}]
&=\nabla_{-X_h}
+\frac{1}{i\hbar}\omega(X_f,X_g)\\
&=-\nabla_{X_h}-\frac{1}{i\hbar}h.
\end{aligned}
$$
Thus
$$
\begin{aligned}
{}[Q(f),Q(g)]
&=(-i\hbar)^2
\left(-\nabla_{X_h}
-\frac{1}{i\hbar}m_h\right)
+2i\hbar m_h\\
&=\hbar^2\nabla_{X_h}+i\hbar m_h\\
&=i\hbar(-i\hbar\nabla_{X_h}+m_h)\\
&=i\hbar Q(h).
\end{aligned}
$$
Dividing by $i\hbar$ proves the formula.
Linearity is immediate from the definition of
$Q$, and $Q(1)=\operatorname{Id}$ because
$X_1=0$.
\end{proof}

\section{What remains}
\label{geometric-prequantization-section-what-remains}

\noindent
The prequantum Hilbert space is usually too
large.  Geometric quantization reduces it by
choosing a polarization and keeping only
sections that are covariantly constant along
that polarization.  For example, on Kahler
manifolds this leads to holomorphic sections
of the prequantum line bundle.

\medskip\noindent
There is also a structural obstruction to
asking for a perfect quantization of all
classical observables: the Groenewold--van
Hove theorem says that the full Dirac rule
cannot hold for all polynomial observables on
standard phase space.  This is one reason why
prequantization is only the first step.


% --- End geometric-prequantization.tex ---

% --- Begin seminars-differential-geometry.tex ---
\chapter{Seminars in Differential Geometry}
\label{chapter-seminars-differential-geometry}















\section{A primer on symplectic groupoids}
\label{seminars-differential-geometry-section-prime-sympl-group}

\noindent
Camilo Angulo, Jilin University.
Geometric Structures Seminar, IMPA.
February 13, 2025.

\medskip\noindent
{\bf Abstract.}
In the late 17th century, Simeon Denis
Poisson discovered an operation that
helped encoding and producing conserved
quantities. This operation is what we
now know as a Lie bracket, an infinitesimal
symmetry, but what is its global
counterpart? Symplectic groupoids are one
possible answer to this question. In
this talk, we will introduce all the basic
concepts to define symplectic
groupoids, and their role in Poisson
geometry. We will discuss key examples, and
applications. The talk will be accessible to
those familiar with differential
geometry, but no prior knowledge of groupoids
will be assumed.


\medskip\noindent
{\bf Part 1. Poisson geometry.}

Hamiltonian formalism. Recall that being a
conserved quantity \(f \in
C^\infty(X)\) is the same thing as
\(\{H,f\}=0\).

\begin{itemize} \item We have seen that it is
always possible to take quotient
of a symplectic manifold with a group action
to obtain a  Poisson
manifold.  \item Then we have found a way to
produce a symplectic foliation
from a 2-vector \(\pi \in
\mathfrak{X}^2(M):=\Lambda^{2}(TM)\).  \item

\begin{remark}
\[\left\{ \text{Lie algebra on
\(\mathfrak{g}\)}
\right\} \xrightarrow{1-1}\left\{
\text{Linear Poisson bracket on
\(\mathfrak{g}^*\)}  \right\}  \]
\end{remark}

\item We saw very nice examples
of foliation that have to do with Lie
algebras. So \(\mathfrak{b}^*_3\) which
gives the ``open book foliation", and
\(\mathfrak{e}^*\) that gives a foliation
by cylinders.
\end{itemize}

\medskip\noindent
{\bf Part 2. Symplectic realizations.}

Consider
\[(\Sigma,\omega) \xrightarrow{\mu}(M,\pi)\]


\iffalse\begin{tikzcd}
	T_p\Sigma
	\arrow[r,"d_p\mu"]\arrow[d,"\omega^\flat"]&
	T_xM\\
	T_p^*\Sigma &
	T^*_pM\arrow[l,"(d\mu)^*",swap]\arrow[u,"\pi
	^\sharp"]
\end{tikzcd}\]\fi
So that
\[\pi ^\sharp=d_p\mu \circ \omega^{-1}
\circ(d\mu)^*\]



\begin{lemma}
\(\dim (\Sigma) \geq  2 \dim (M) -
\operatorname{rk}(\pi_x)\) for all \(x \in
M\).
\end{lemma}

\begin{proof}
Done in seminar.
\end{proof}

\begin{example}
\((\mathbb{R}^2,0)\). So the map
$$
\begin{aligned}
	(\mathbb{R}^4,dx \wedge du + dy \wedge dv)
	&\longrightarrow  \mathbb{R}^2\\
	(x,y,u,v) &\longmapsto (x,y)
\end{aligned}
$$

\end{example}

\begin{exercise}
Find the symplectic realization \(\omega\) in
\((\mathbb{R}^4, \omega )
\xrightarrow{\mu}(\mathbb{R}^2,(x^2+y^2)
\partial_x
\wedge \partial_y\)

$$
\begin{aligned}
	(\mathbb{R}^4, \omega) &\longrightarrow
	(\mathbb{R}^2,(x^2+y^2) \partial_x \wedge
	\partial _y) \\
	(x,y,z,w) &\longmapsto (x,y)
\end{aligned}
$$

Also find the symplectic realization of
\(\operatorname{ a f f }^*\) with
\(\{x,y\}=x\).
\end{exercise}

\medskip\noindent
{\bf Part 3: Groupoids}

\medskip\noindent
{\bf Motivation}
\begin{enumerate}
\item Fundamental grupoid: objects are points
in the manifold and arrows are paths.
\item The action of \(S^1\) on \(S^2\) by
rotation does not give a nice quotient
because there are two singular points.
Consider the groupoid  \(S^1 \times S^2\) of
orbits. These are the arrows. The points are
just the points of \(S^2\).
\item Consider a foliation (like Möbius
foliation of circles; where there is a
singular circle, the soul). You can do the
same thing as in fundamental groupid
leafwise. Arrows then are equivalence classes
of paths that live inside leaves. This is
called monodromy of a foliation. Again
objects are points.
\item You can take a quotient of monodromy
using a connection given by the foliation.
This allows to identify certain paths between
the leaves. This is called {\it holonomy (of
a foliatiation)}. (So you can make this
notion match the usual holonomy given by
riemmanian connection.)

{\bf Upshot.}
	So the point is taking some sort of function
	space on these groupoids you can gather the
	information given by the non-smooth quotient
	(like in the case of the sphere rotating).
	So this groupoid motivation says how to get
	some structure that resembles a non smooth
	quotient.
\item Last motivation: the grid of squares
has a tone of symmetries. If you restrict to
just a few squares you loose so many
symmetries. But there's a grupoid hidden in
there that tells you what you intuition knows
about this finite grid of squares.
\end{enumerate}

\begin{definition}
A {\it groupoid} is a category where all
morphisms are invertible.
\end{definition}

So there is a kind of product among the
objects, given by composition but:
 not every two pair of objects can be
 multiplied!---only those whose
source and target match. So that's the lance
about grupoids.

Just so you make sure you understand: the
groupoid \(G\) is the morphisms of the
category. The objects are points (of a
manifold).

\begin{definition}
{\it Lie grupoid} is when the following
diagram is inside category of smooth
manifolds and \(s,t\) (source and target
maps) are submersions:
$$
\xymatrix{
G^{(2)}\ar[r]^m& G\ar[r]^ s_t & M\ar[r]^u&G
}
$$
\end{definition}

\begin{proposition}[Properties of Lie
groupoids]
\begin{itemize}
\item \(m\) is also a submersion.
\item \(i\) (inversion) is a diffeomorphism.
\item  \(u\) (unit=identity) is an embbeding.
\end{itemize}
\end{proposition}

\begin{definition}
Consider \(x \in M\) and the inverse image of
source map: \(s^{-1}(x)=\{\text{arrows that
start at \(x\)} \}\). Now if you act with
\(t\) on this set you get {\it the orbit} of
\(x\): \(\{\text{\(y \in M\) such that there
is an arrow from \(x\) to \(y\)} \}\).

And there also {\it an isotropy} \(G_x=\{g
\in G: \text{\(g\) goes from \(x\) to \(x\)}
\}\)
\end{definition}

\begin{example}
\begin{enumerate}
\item \(G=M\),  \(M=M\).
 \item Lie groups.
	\item Lie group bundles.
\item \(G=M \times M\), \(M =M\).
 \item Fundamental groupoid. Isotropy group
 is fundamental group! And orbit
 is…\clearpage universal cover!
\item Subgroupoids.
\item  Foliations.
\item If you have a normal group action of
\(G\) on \(M\) you construct a groupoid
action with groupoid \(G \times M\) and
objects \(M\), with product given on the
group part of the product. Orbits are orbits.
Isotropy group is isotropy group.
 \item Principal bundles.
\end{enumerate}
\end{example}

{\bf Back to Poisson.}

There's also a notion of Lie algebroid. Which
is strange. But the point is that to every
Poisson manifold there is a Lie algebroid.

So the question is whether there is a Lie
groupoid associated to that Lie algebroid.
Not always.

\medskip\noindent
{\bf Big question}[Fernandez and ?]
When a symplectic manifold is integrable?


(Remember that integrating means go from
algebra(oid) to group(oid).

And the point is that:

\medskip\noindent
When you {\it can} go back, you get a
\textit{ symplectic groupoid}.

\begin{remark}
Look for Kontsevich's notes on Weinstein!
\end{remark}

\begin{remark}
History: Weinstein did this intending to do
quantization (geometric?) on Poisson
manifolds. (That involves a \(C^*\)algebra
coming from the symplectic groupoid.)
\end{remark}

\begin{definition}
A {\it symplectic groupoid} is a groupoid
\(G, M\) together with  \(\omega \in
\Omega^2(M)\) such that \(\omega\) is
symplectic and multiplicative, meaning that
\(\partial  \omega =0\), that is, \(\iff
m^*\omega=\operatorname{pr}_1^*\omega+
\operatorname{pr}^*_2\omega\in
\Omega^{2}(G^{(2)})\iff\) take two vectors \(
X_k,Y_k \in TG\), and \(\omega(X_0 \star
Y_0,X_1\star Y_1) =  \omega(X_0,Y_0) + \omega
(X_1,Y_1)\)
\end{definition}

\begin{theorem}
If \((G, \omega)\) is a symplectic groupoid,
then
\begin{enumerate}
\item \(\exists !\) poisson structure on
\(M\)
\item for which \(t: G \to M\) is a
symplectic realization,
\item  Leaves are connected components of
orbits,
\item \(\mathsf{Lie}(G)\cong T^*M\) via \(X
\mapsto -u ^*(i_X \omega\).
\end{enumerate}
\end{theorem}

\begin{remark}
Look for Alejandro Cabrera, Kontsevich. There
are two things one is de Rham and the other…
from the future: simplicial?
\end{remark}

\medskip\noindent
{\bf Upshot.} The obstruction to knowing when
symplectic groupoid exists is ``variation of
symplectic form \(\omega=(1+t^2)
\omega_{S^2}\)". So how does the symplectic
group vary from leaf to leaf. So there are
two situations in which the thing doesn't
work.

\section{Spheres with minimal equators}
\label{seminars-differential-geometry-section-spheres-with-minimal-equators}

\noindent
Lucas Ambrozio, IMPA.
Differential Geometry Seminar, IMPA.
June 24, 2025.

\medskip\noindent
{\bf Abstract.} We will discuss the
connection between Riemannian metrics on the
sphere with respect to which all equators are
minimal hypersurfaces, and
algebraic curvature tensors with positive
sectional curvatures.

\bigskip\noindent

\begin{definition}
\label{seminars-differential-geometry-definition-equator}
An {\it $(n-k)$-equator} orthogonal to $\Pi$
is
$$
\Sigma_\Pi:=\{p\in\mathbb{S}^n:\left\langle{}p,
x\right\rangle{}=0\forall
x\in\Pi\}
$$
for $\Pi$ a $k$-dimensional linear subspace
of $\mathbb{R}^{n+1}$.
\end{definition}

\begin{remark}
\label{seminars-differential-geometry-remark-equat-total-geode-usual-metr}
Equators are totally geodesic hypersurfaces
with the usual sphere metric, which
implies they are minimal hypersurfaces.
\end{remark}

{\bf Problem.} Characterize the set
$\mathcal{M}_k(U)$ of metrics $g$ on an
open set $U\subset\mathbb{S}^n$ such that all
$k$-equators $\Sigma_\Pi$ with
 $\Sigma\cap U\neq\emptyset$ yield are
 minimal hypersurfaces $\Sigma\cap U$
on $(U,g)$.

\begin{remark}
\label{seminars-differential-geometry-remark-why-not-Rn}
This problem can be thought of as a problem
of finding metrics on $\mathbb{R}^n$
such that $k$-planes are minimal. To see why
project the $k$-equators to
$T_p\mathbb{S}^n$ and pullback those metrics
to the sphere.
\end{remark}

\medskip\noindent

Let $g\in\mathcal{M}_k(U)$ for
$U\subset\mathbb{S}^n$ open and $n\geq 2$.

\begin{theorem}[Beltrami, Schäfli]
\label{seminars-differential-geometry-theorem-Beltrami-Schafli}
If $k=1$ then $g$ has constant sectional
curvature.
\end{theorem}

\begin{theorem}[Hongan]
\label{seminars-differential-geometry-theorem-Hongan}
If $1<k<n-1$ then $g$ has constant sectional
curvature.
\end{theorem}

Then Hongan also managed to produce a
classification of these metrics for
$k=n-1$.

\begin{remark}
\label{seminars-differential-geometry-remark-linea-inver-prese-equat-mg}
If $T\in \text{GL}(n+1,\mathbb{R})$, then
$$
\begin{aligned}
\varphi: \mathbb{S}^n &\longrightarrow
\mathbb{S}^n \\
x &\longmapsto \frac{Tx}{|Tx|}
\end{aligned}
$$
is a diffeomorphism that maps $k$-equators
into $k$-equators. Thus if
$g\in\mathcal{M}_k(\mathbb{S}^n)$ then so is
$\varphi(T)^*g$.
\end{remark}

\begin{theorem}
\label{seminars-differential-geometry-theorem-equivariant-bijection}
There exists a $\text{GL}(n+1,\mathbb{R})$
equivariant bijection
$$
\mathcal{M}_{n-1}(\mathbb{S}^n)
\leftrightarrow
\text{Curv}_+(\mathbb{R}^{n+1})
$$
where the set on the right-hand-side is the
set of algebraic curvature tensors
(also called curvature-like, i.e. with the
same symmetries as the Riemannian
curvature tensor) on $\mathbb{R}^{n+1}$ with
positive sectional curvature.

The group action is given as follows for
$T\in\text{GL}(n+1,\mathbb{R})$:
$$
(R\cdot
T)(x,y,z,w)=\frac{1}{|\det(T)|^{\frac{1}{n+1}
}}R(Tx,Ty,Tz,Tw)
$$
\end{theorem}

The point is that
$\text{Curv}_+(\mathbb{R}^{n+1})$ is an open
cone on a linear
space. Here are two simple corollaries:

\begin{lemma}
\label{seminars-differential-geometry-lemma-corollaries}
\begin{enumerate}
\item $\mathcal{M}_{n+1}(\mathbb{S}^n)$ is in
bijection with an open positive
cone of an
$\frac{n(n+2)(n+1)^2}{12}$-dimensional real
vector space.
\item Every metric on
$\mathcal{M}_{n-1}(\mathbb{S}^n)$ is
invariant by the
antipodal map.
\end{enumerate}
\end{lemma}

{\bf Algorithm.} From any $R\in
\text{Curv}_p(\mathbb{R}^{n+1})$ we obtain a
symmetric positive definite
(positive-definitiness comes from the
positiveness
of the curvature of $R$) 2-tensor  $k_R$
satisfying
$$
(k_R)_p(v,v)=R(pv,pv)>0
$$
Also, $k_R$ has the  {\it Killing property},
i.e. that $\overline{\nabla}k(X,X,X)=0$ for
all
$X\in\mathfrak{X}(\mathbb{S}^n)$.

Then we define a positive function on
 $\mathbb{S}^n$ by
\begin{equation}
\label{seminars-differential-geometry-equation-def-DR}
D_R:=\left(\frac{d\text{Vol}_{k_R}}{dV_g}
\right)^{\frac{4}{n-1}}
\end{equation}
and finally a Riemannian metric on
$\mathbb{S}^n$ in
$\mathcal{M}_{n-1}\mathbb{S}^n$ by
$$
g_R=\frac{1}{D_R}k_R
$$
And to go back, for $g \in
\mathcal{M}_{n-1}(\mathbb{S}^n)$ define a
positive
function on $\mathbb{S}^n$
$$
F_g:=\left(\frac{dV_g}{dV_{\overline{g}}}
\right)^{\frac{4}{n-1}}
$$
Then let $k_g:=\frac{1}{F_g}g>0$, which is a
positive definite Killing 2-tensor,
from which we may define
$R_g\in\text{Curv}_+(\mathbb{R}^{n+1})$ with
$R_g(pv,pv)=(k_g)_p(v,v)$ for all $p,v\in
T\mathbb{S}^n$.

More corollaries:
\begin{lemma}
\label{seminars-differential-geometry-lemma-corollaries2}
\begin{enumerate}
\item $g\in\mathcal{M}_{n-1}(\mathbb{S}^n)$
is analytic because it is a Killing
tensor on $\mathbb{S}^n$, which are
well-known.
\item If $g$ is left-invatiant on
$\mathbb{S}^3$, seen as unit quaternions,
then
$g\in\mathcal{M}_2(\mathbb{S}^3$. Moreover,
for $a\geq b\geq c>0$,
$$
aL_i\odot L_i+bL_j\odot L_j+cL_k\odot L_k=k
$$
is Killing, $k>0$,  $D_k$ constant and thus
$g=\frac{1}{\text{const.}}k
\in\mathcal{M}_2(\mathbb{S}^3)$.
\item $R$ curvature tensor of
$(\mathbb{C}P^{2},g_{FS})$. We may not
remember
what's the curvature tensor, but we know the
sectional curvature is $1\leq
\text{sec}(R)\leq 4$,
$$
(k_R)_p(v,w)=\overline{g}(v,w)+3\overline{g}
(Jp,v)\overline{g}(Jp,w)
$$
and $D_R=4^{\frac{4}{3-1}}=4$, so that by
\ref{seminars-differential-geometry-equation-def-DR} we obtain
$g_R=\frac{1}{4}k_R$, which is a Berger
metric on $\mathbb{S}^3$ with scalar
curvature 0.
\end{enumerate}
\end{lemma}

\bigskip
Now define
$$
\Sigma_V=\{p\in\mathbb{S}^n:\left\langle{}p,v\right\rangle{}=
0\}=V^{-1}(0)
$$
where $V(x):=\left\langle{}x,v\right\rangle{}$ for all
$x\in\mathbb{S}^n$. Then the normal
vector field is $\nabla V/|\nabla V|_g$, and
the second fundamental form is
given by
$$
A=\frac{1}{|\nabla V|_g}\text{Hess}_gV
$$
and its mean curvature by
\begin{equation}
\label{seminars-differential-geometry-equation-mean-curva-level-sets-V}
H=\frac{1}{|\nabla
V|_g}\left(\Delta_gV-\text{Hess}
_gV\left(\frac{\nabla
V}{|\nabla V|},\frac{\nabla V}{|\nabla
V|}\right)\right)
\end{equation}
For every $v \in \mathbb{S}^n$ and
$p\in\Sigma_V$, we see that $H_{\Sigma_V}=0$
iff
$$
|\nabla
V|^2_g(p)\Delta_gV(p)-\text{Hess}_gV(\nabla
V(p),\nabla V(p))=0
$$
And for $\overline{g}$,
$$
\text{Hess}_{\overline{g}}V+V\overline{g}=
0\implies
\text{Hess}_{\overline{g}}V(X,X)=0
$$
for all $X\in T_p\mathbb{S}^n$ and
$p\in\Sigma_v$. Then
$$
J_g(X,Y,Z)=g(\nabla_XY-\overline{\nabla}_XY,
Z)
$$
$$
J_g(X,Y,\nabla
V)=\text{Hess}_{\overline{g}}-\text{Hess}
$$
{\bf Problems.}
\begin{enumerate}
\item Similar story for
$\mathbb{C}P^{n},\mathbb{H}P^n$?
\item Complete metrics on $\mathbb{R}^n$ with
minimal hyperplanes.
\item Find geometric invariants of metrics on
$\mathcal{M}_{n-1}(\mathbb{S}^n)$
(may be useful to study $(M^n,g)$, $n\geq 4$,
$\text{sec}>0$.
\end{enumerate}

\section{On wrapped Floer homology barcode
entropy and hyperbolic sets}
\label{seminars-differential-geometry-section-wrapp-floer-homol-barco}

\noindent
Rafael Fernandes, UC Santa Cruz.
Differential Geometry Seminar, IMPA.
September 4, 2025.

\medskip\noindent
{\bf Abstract.} In this talk, we will discuss
the interplay between the
wrapped Floer homology barcode and
topological entropy. The concept of barcode
entropy was introduced by Çineli, Ginzburg,
and Gürel and has been shown to be
related to the topological entropy of the
underlying dynamical system in various
settings. Specifically, we will explore how,
in the presence of a topologically
transitive, locally maximal hyperbolic set
for the Reeb flow on the boundary of
a Liouville domain, barcode entropy is
bounded below by the topological entropy
restricted to the hyperbolic set.

\medskip\noindent

Let $M^n$ be a manifold.
$\omega \in \Omega^2(M)$ is {\it symplectic}
if $d \omega=0$ and it is
nondegenerate.

\begin{example}
\label{seminars-differential-geometry-example-Rn-is-symplectic}
$\mathbb{R}^n$ is symplectic with canonical
Darboux form.
\end{example}

Recall the definition of Hamiltonian vector
field
associated to a function $H \in C^\infty(M)$.

\begin{definition}
\label{seminars-differential-geometry-definition-nonde-diffe}
A diffeomorphism $\varphi:M \to M$ is called
{\it non-degenerate} if
$\Phi(\varphi)\cap\Delta \subset M\times M$
(pitchfork, i.e. transversal intersection!).
\end{definition}

Let $M^{2n}$ be a closed symplectic manifold.
Arnold's conjecture says
\begin{enumerate}
\item If $\varphi=\varphi_H$ (Hamiltonian
flow) is nondegenerate, then
$$
\# \text{Fix}(\varphi_H) \geq
\sum_{i=0}^{2n}\dim H_i(M,k)=\dim H_*(M,k)
$$
\item If $\varphi=\varphi_H$ is degenerate,
then
$$
\# \text{Fix}(\varphi)\geq \text{Cl}(M)+1
$$
where $\text{Cl}(M)$ is the maximum number of
homology classes
we can add before getting to zero.
\end{enumerate}

Why do we care? Because
$$
\#
\text{Fix}(\varphi_H)
\leftrightarrow\{\text{1-periodic
orbits of $X_H$}\}
$$

Idea by Floer. Construct an invariant that
would
say something about periodic orbits.

\medskip\noindent
{\bf Question.} Can Floer theory capture
other ``dynamical information''?
(Other than the periodic orbits.)

\medskip\noindent
A {\it persistence module} is a pair
$(V,\Pi)$, where $V=\{V_t\}_{t \in
\mathbb{R}}$ is a family of
$\mathbb{F}$-vector spaces
and $\Pi=\{\Pi_{st}\}_{s \leq  t}$ is a
family of maps such that
\begin{enumerate}
\item $\Pi_{ss}=, \Pi_{ts}\circ
\Pi_{rt}=\Pi_{rs}$.
\item $\exists  s \subset \mathbb{R}$ such
that
$\Pi_{s t}$ is an isomorphism for $s,t$ in
the same connected component of
$R\setminus S$.
\item $\Pi_{s t}$ have finite rank.
\item $\exists s_0$ $V_s=\{0\}$, $s \leq
s_0$.
\item $V_t= \lim_{s \to t} V_s$ (lower
limit!!)
\end{enumerate}

\begin{theorem}
\label{seminars-differential-geometry-theorem-persi-modul-sum-integ}
Any persistence module is a sum of integral
persistence modules,
$$
(V,\Pi) \cong \bigoplus_{I \in B(V)}F(I).
$$
\end{theorem}

\begin{example}
\label{seminars-differential-geometry-example-hart-and-sphere}
Heart and sphere. There is a noise in the
persistence module
of the heart due to an unnecessary critical
point.
\end{example}

$(M^{2n},\omega)$ a {\it Liouville domain} is
a compact
symplectic manifold and $X \in
\mathfrak{X}(M)$
with $X \cap\partial M$ (pitchfork, i.e.
transversal intersection!)
pointing outwards
and preserved by the symplectic form, i.e.
$\mathcal{L}_X \omega=\omega$
($\omega=d \alpha$).

When we restrict $\omega$ to the boundary,
we obtain a contact form and get some
interesting dynamics.

A Lagrangian $(L,\partial L) \subset
(M,\partial M)$
is {\it asympotically conical} if
\begin{enumerate}
\item $\partial L \subset \partial M$ is
Legendrian.
\item $L \cap [1-\varepsilon,1] \times
\partial M=
[1-\varepsilon,1] \times \partial L$.
\end{enumerate}

\begin{remark}
\label{seminars-differential-geometry-remark-Hamiltonian}
Take a Hamiltonian $H: \hat{M}\to \mathbb{R}$
such that
$$
\begin{cases}
H(r,x)=h(r)\qquad &r=1 \\
H(r,x)=rT-B\qquad &
\end{cases}
$$
then $X_H=h'(r)R_\alpha$.

For $L_0,L,A\subset$ Lagrangians, $H$ linear
at infinite,
then $A_{H}^{L_0 \to L_1}$,
$$
\begin{aligned}
A_{H}^{L_0 \to L_1}: P_{L_0\to L_1}
&\longrightarrow \mathbb{R} \\
\gamma &\longmapsto \int_0^1 \gamma^* \alpha-
\int_0^1H(x(t))dt
\end{aligned}
$$
where $P_{L_0\to L} =\{\gamma:[0,1] \to
\hat{M}:\gamma(0) \in L_0,
\gamma(1) \in L_1\}
$ is the set of chords.
\end{remark}

\begin{remark}
\label{seminars-differential-geometry-remark-1-chords-are-crticial-points}
$\text{crit}(A_{H}^{L_0 \to
L_1})=\{\text{1-chords of $X_H$ from
 $L_0$ to $L_1$}\}$.
\end{remark}

Putting a metric on $P_{L_0 \to L_1}$ we can
consider $\varphi:\mathbb{R} \times
[0,1] \to \hat{M}$, solutions of some PDE
which is some kind
of generalization of a gradient, $- \nabla
A_{H}^{L_0 \to L_1}$.
These solutions can be put in a moduli space
$$
\tilde{\mathcal{M}}(x_-,x_+,H,J)=
\{\varphi\text{
solutions s.t. …}\}
$$
Then we define a boundary operator
$\partial$.
\begin{theorem}
\label{seminars-differential-geometry-theorem-homology}
$\partial^2=0$
\end{theorem}

So that we have a homology, called {\it
wrapped Floer homology}
$HW^t(H,L_0,L_1,J)$

\begin{remark}
\label{seminars-differential-geometry-remark-embedding}
We have $H \leq  K \rightsquigarrow
HW^t(H,L_0,L_1,J) \to
HW^t(K,L_0,L_1,K)$.
\end{remark}

\begin{definition}
\label{seminars-differential-geometry-definition-direct-limit}
For $t \geq 0$
$$
HW^t(M,L_0,L_1)=\underline{\lim }_H
HW^t(H,L_0,L_1,J)
$$
\end{definition}

(Where we have taken direct limit.)

Taking direct limit of the homology,
we make sure the homology theory is
independent of the
choice of objects (I think, complex structure
and Hamiltonian)
we used to construct it.


\begin{proposition}
\label{seminars-differential-geometry-proposition-persistence-module}
$t \to HW^t(M,L_0,L_1)$ is a peristence
module $B(M,L_0,L_1).$
\end{proposition}

Finally we can define {\it barcode entropy}.
Fix $\varepsilon>0$,
$t \geq 0$,
$$
b_\varepsilon(M,L_0,L_1,t)=
\# \{\text{ of bars in }B(M,L_0,L_1)
\text{with length $\geq \varepsilon$
and start before $t$}\}
$$
Then
$$
\bar{h}^{HW}(M_0,L_0,L_1)=
\lim_{\varepsilon \to 0} \text{lim sup}_{t
\to \infty}
\frac{\text{log}^+(b_\varepsilon(M,L_0,L_1,t)
}{t}
$$

\medskip\noindent
Consider a contact manifold
$(\Sigma,\lambda,L_0,L_1)$
and $A$ the Lagrangians. (Example raising
question of filling.)

\begin{theorem}[M '24]
\label{seminars-differential-geometry-theorem-entro-indep-filli}
$\bar{h}^{HW}$ is independent of the filling.
\end{theorem}

\begin{theorem}[M '24]
\label{seminars-differential-geometry-theorem-bound}
$\bar{h}^{HW}(M,L_0,L_1) \leq
h_{\text{top}}(\alpha)$.
\end{theorem}

\begin{theorem}[M '25]
\label{seminars-differential-geometry-theorem-restriction}
Consider $(M,L_0,L_1)$.
Let $K$ be a compact topologically transitive
hyperbolic set
for the Reeb flow $\alpha$. Assume $W_\delta
^s (q) \subset \partial L_0$,
$W_\delta^s(p) \subset \partial L_1$. Then
$$
\bar{h}^{HW}(M,L_0,L_1)\geq
h_{\text{top}}(\alpha |_{K})>0.
$$
\end{theorem}
Which says it captures dynamics beyond
unconditional phenomena.
In lower dimensions these tend to coincide,
but in higher dimension we don't
know. This is related to
 ``the sup over hyperbolic sets […]''

Here's a conjecture:
$$
\text{sup}_{L_0,L_1}\bar{h}^{HW}(M,L_0,L_1)=
h_{\text{top}}(\alpha).
$$
\medskip\noindent
{\bf Extra comments.}
One of the aims is to describe topological
entropy $h_\text{top}$
using Floer theory. Theorems by
Çineli-Ginzburg-Gürel show bounds
of topological entropy and barcode entropy
(one of which is for
hyperbolic sets).

There is a notion of  {\it admissible
Hamiltonian and Reeb vector fields}
which is related to some asymptotical
behaviour  ``linear at infinity''.
I understand that admissible vector fields
give
the interesting chords for the Floer homology
construction.

\section{Revisiting cotangent bundles}
\label{seminars-differential-geometry-section-revisiting-cotangent-bundles}

\noindent
Mieugel Cueca, KU Leuven.
Symplectic Geometry Joint Seminar, IMPA.
September 5, 2025.

\medskip\noindent
{\bf Abstract.} Cotangent bundles provide key
examples of symplectic
manifolds. On the other hand, one can think
of Lie groupoids as generalizations
of manifolds. In this context, Alan Weinstein
constructed their cotangent
bundles and proved that they are so-called
symplectic groupoids. In this talk, I
will recall this construction and explain
what happens when one replaces a Lie
groupoid with a Lie 2 (or n)-groupoid. If
time permits, I will exhibit some of
their main applications. This is joint work
with Stefano Ronchi.

\medskip\noindent

Recall the basic properties of the cotangent
bundle $T^*M$ for a symplectic manifold:
\begin{enumerate}
\item It's a vector bundle.
\item $\left\langle{}\cdot,\cdot\right\rangle{}:TM \otimes
T^*M \to \mathbb{R}_M$
is the dual pairing.
\item $\omega_{\text{can}}\in \Omega^2(T^*M)$
is symplectic.
\item $\mathcal{L}_\varepsilon
\omega_{\text{can}}=\omega_{\text{can}}$,
$\varepsilon$ Euler vector field.
\end{enumerate}

\medskip\noindent
{\bf Main goal.} Reproduce the above for Lie
$n$-groupoids.

For $n=0$ we get the above situation. For
$n=1$ [Duzud-Weinstein], [Prezun],
for $n\geq 2$ ? and we care about $n=2$.

\begin{definition}
\label{seminars-differential-geometry-definition-groupoid}
A {\it Lie $n$-groupoid}
$\mathcal{G}:\Delta^{\text{op}}\to
\mathsf{Man}$
such that
$$
P_{e,j}:\mathcal{G}_{\ell}\to \Lambda_j^\ell
\mathcal{G}
$$
are surjective submersions $\forall  \ell,j$
are differomorphisms for $\ell > n$.
\end{definition}

\begin{remark}
\label{seminars-differential-geometry-remark-generators}
This sort of manifolds-valued presheaf
category is generated by
$$
\begin{aligned}
d_i^\ell:\mathcal{G}_\ell \to
\mathcal{G}_{\ell-1}&\text{face maps}
0 \leq  i,j \leq  \ell\\
s_j^\ell:\mathcal{G}_\ell \to
\mathcal{G}_{\ell+1}&  \qquad
\text{degeneracies}
\end{aligned}
$$
\end{remark}

\medskip\noindent
The tangent is a functor, it satisfies
$$
T_\bullet(\mathcal{G})=T_k(\mathcal{G})=
T\mathcal{G}_k
$$
(It looks like $T$ preserves diagrams.)

\medskip\noindent
{\bf Dold-Kan.} The category $\mathsf{SVect}$
of simplicial vector spaces
has objects
$$
\xymatrix{
\mathbb{V}_\bullet&
\mathbb{V}_n\ar[r]&\cdots
\ar[r]&\mathbb{V}_2\ar[r]^{\text{3 arrows}}&
\mathbb{V}_1\ar[r]^{\text{2
arrows}}&\mathbb{V}_0
}
$$
where the $\mathbb{V}_i$ are vector spaces.

There is a functor
$$
\mathsf{SVect}\overset{N}{\to}\{\text{chain
complexes}\geq 0\}
$$
$$
\mathbb{V}_{\bullet}\to N\mathbb{V}=
\left(N_\ell \mathbb{V}\Ker
P_{\ell,\ell},\partial=d_\ell\right)
$$

\begin{theorem}[Dold-Kan]
\label{seminars-differential-geometry-theorem-Dold-Kan}
That's an equivalence of categories. [Confirm
this!]
\end{theorem}

Those categories are monodial:
$$
(\mathsf{SVect},\otimes),\qquad
(\mathbb{V}_\bullet \otimes \mathbb{W})_\ell
=\mathbb{V}_\ell \otimes \mathbb{W}_\ell
$$
$$
(\mathsf{ch}_{\geq 0},\otimes),\qquad (V
\otimes W)_i=
\bigoplus_{\ell +k=1}V_\ell \otimes W_k
$$
And $N$ is Lax monoidal with Lax structure
given by the
Eilenberg-Zilber map, though we won't explain
the details of this.

There are duals given by internal Hom:
$$
\begin{aligned}
\mathbb{V}^{n*}&=\underline{\Hom}(\mathbb{V},
B^n\mathbb{R})
\end{aligned}
$$
Where the internal Hom is given by
$$
\underline{\Hom}(\mathbb{V},B^n\mathbb{R})
_\ell
=\Hom_{\mathsf{SVect}}(\mathbb{V}
\otimes\Delta_n[\ell],B^n\mathbb{R})
$$
for an object
$\Delta[\ell]=\mathbb{R}[\Delta[\ell]]$.


\medskip\noindent
{\bf Properties.}
\begin{enumerate}
\item $\mathbb{V}^{n*}$ isa simplicial vector
bundle.
\item $N(V^{n*})$ and $N(\mathbb{V})^*[n]$ is
a quasi isomorphism.
\item
$\left\langle{}\cdot,\cdot\right\rangle{}:\mathbb{V}\otimes
\mathbb{V}^{n*}\to
B^n\mathbb{R}$ is non-degenerate on homology.
\item $\mathbb{V}\hookrightarrow
(\mathbb{V}^{n*})^{n*}$
Mont. a eq.
\end{enumerate}

\medskip\noindent
{\bf The vector bundle case.}
$$
\text{Maps}(\Delta[i],\mathcal{G})_k=
\Hom_{\mathsf{SSet}}(\Delta[i]\times\Delta[k]
,\mathcal{G})
$$
\begin{proposition}
\label{seminars-differential-geometry-proposition-on-an-n-Lie-groupoid}
Let $\mathcal{G}$ be a Lie $n$-groupoid.
\begin{enumerate}
\item $\text{Maps}(\Delta[i],\mathcal{G})$
Lie $n$-groupoids ME $\mathcal{G}$.
\item
$\text{Maps}(\Delta[i],\mathcal{G})_0=
\mathcal{G}_i$.
\item $\text{ev}:\Delta[i]\times
\text{Maps}(\Delta[i],\mathcal{G})\to
\mathcal{G}$.
\end{enumerate}
\end{proposition}

[Staircase looking diagram.]

\begin{definition}
\label{seminars-differential-geometry-definition-what-is-this}
$\mathcal{G}_\bullet$.
$$
T_i^{n*}\mathcal{G}=\Hom_{\mathsf{SVect}}(1^*
\Pi_{\Delta[i]}(T\mathcal{G}),
B^n\mathbb{R}_{\mathcal{G}_i})
$$
$$
(d_j,F)_{K|d_j\mathcal{G}}(x^a)=(F_k)
|_{\mathcal{G}}(x^{\delta_ja}).
$$
\end{definition}

\begin{proposition}
\label{seminars-differential-geometry-proposition-properties-of-that}
$\mathcal{G}$ Lie $n$-groupoid, then
$T^{n*}\mathcal{G}$ satisfy
\begin{enumerate}
\item is a vector bundle $n$-groupoid
\item dual to $T\mathcal{G}$
$$
\left\langle{}\cdot,\cdot\right\rangle{}:T\mathcal{G} \otimes
T^{n*}\mathcal{G}
\to B^n\mathbb{R}_{\mathcal{G}}
$$
non-degenerate on homology.
\item $n$-shifted symplectic
$$
T^{n*}_n\mathcal{G}\overset{p}{\to}
T^*\mathcal{G}_n
$$
and $p^*\omega_{\text{can}}$.
\end{enumerate}
\end{proposition}

[More computations I missed]

\section{Holomorphic extensions of s-proper
Lie groupoids}
\label{seminars-differential-geometry-section-holom-exten-s-prope-lie}

\noindent
Rui L. Fernandes, Instituto Superior Técnico
(Lisboa)/University of Illinois
at Urbana-Champaign.
Symplectic Geometry Seminar, IMPA.
September 17, 2025.

\medskip\noindent
{\bf Abstract.} Every smooth manifold admits
a compatible analytic
structure, and a classical result of
Whitney–Bruhat states that any analytic
manifold has a holomorphic extension. Lie
groups also admit compatible analytic
structures, and another classical result, due
to C. Chevalley, shows that any
compact Lie group has a holomorphic extension
to a complex Lie group. D.
Martínez Torres has shown that any proper Lie
groupoid admits a compatible
analytic structure. I will discuss an
extension of the classical results of
Whitney–Bruhat and Chevalley, establishing
that any s-proper Lie groupoid has a
holomorphic extension. This talk is based on
recent joint work with Ning Jiang
(arXiv:2508.18036).

\medskip\noindent

\begin{theorem}[Whitney-Beuhat]
\label{seminars-differential-geometry-theorem-Whitney-Beuhat}
If $M$ is analytic there exists a complex
manifold  $M_\mathbb{C}$
together with an analytic map $i: M \to
M_\mathbb{C}$
which is totally real ($T_HM_{\mathbb{C}}=TM
\oplus J(TM)$)
and
\begin{enumerate}
\item For every complex manifold $X$ and
every
analytic map $\phi:N \to X$ there exists
$\phi^*:U \to X$
holomorphic contained open $M \subset U
\subset M_\mathbb{C}$.
\item If $\psi:V \to X$ holomorphic on $M
\subset V \subset M_\mathbb{C}$ then
$\psi=\phi^*$ where $\phi=\psi \circ i$ on a
possibly
smaller open contained in $U \cap V$.
\end{enumerate}
\end{theorem}

\begin{theorem}[Chevalley]
\label{seminars-differential-geometry-theorem-Chevalley}
If $G$ is a Lie group, there exists a complex
Lie group $G_{\mathbb{C}}$ and a morphism
$i:G \to G_\mathbb{C}$
satisfying the following universal property.
For every complex Lie group $H$ and morphism
$\phi:G \to H$ there exists a unique
holomorphic map $\phi^*:G_\mathbb{C} \to H$
such that
$\phi=\phi^* \circ i$.
\end{theorem}

Here's the construction of $G_\mathbb{C}$:

$$
\xymatrix{
N\subset \tilde{G}\ar[r]\ar[d]&G^*\ar[d]\\
\tilde{G}/N \simeq
G\ar[r]_{i}&G_\mathbb{C}=G^* /\overline{i^*N}
}
$$
where $\tilde{G}$ is the universal cover
group of $G$,
 $G^*$ is the group that integrates to
the complexification of the Lie algebra of
$G$,
i.e.
$\text{Lie}(G^*)=\mathfrak{g}_\mathbb{C}$
and $\overline{i^*(N)}$ is the smallest
closed
normal complex Lie subgroup of $G^*$
containing $i^*(N)$.

\begin{example}
\label{seminars-differential-geometry-example-complexification-groupoid}
$$
\xymatrix{
\widehat{\text{SL}_2(\mathbb{R})}\ar[d]
\ar[dr]\\
\text{SL}_2(\mathbb{R})\ar@{^{(}->}[r]&
\text{SL}_2(\mathbb{C})
}
$$
which is a simple case,
but consider instead
$\tilde{\text{SL}_2}(\mathbb{R})\times
\mathbb{R}$
and the normal subgroup
$N=\left\langle{}a\right\rangle{}\times \left\langle{}\lambda\right\rangle{}$
for irrational $\lambda$.
Then we obtain
$$
\xymatrix{
\widehat{\text{SL}_2(\mathbb{R})}
\times\mathbb{R}\ar[r]\ar[d]
&\text{SL}_2(\mathbb{C})\times\mathbb{C}
\ar[d]\\
\widehat{\text{SL}_2(\mathbb{R})}
\times\mathbb{R}/N\simeq
G\ar[r]&
G_\mathbb{C}
}
$$
and $G_\mathbb{C}$ is 3 complex dimensions!
It is not
$\text{SL}_2(\mathbb{C})\times\mathbb{C}$.
\end{example}

\section{Lagrangian blow-up and blow-down for
4-dimensional symplectic manifolds
}
\label{seminars-differential-geometry-section-blow-up-down}

\noindent
Misha Verbitsky, IMPA.
Geoemtric Structures Seminar, IMPA.
October 23, 2025.

\medskip\noindent
{\bf Abstract.} The usual (complex geometric)
blow-up has a symplectic
version, called the symplectic cut. I will
introduce a variant of this
construction, which is valid only in
symplectic category, called ``Lagrangian
blow-up'', and its inverse, called Lagrangian
blow-down. Lagrangian blow-down
takes a Lagrangian sphere in a 4-dimensional
symplectic manifold and contracts
it to a symplectic orbifold with a double
point. Given a symplectic orbifold
with a double point, Lagrangian blow-up
produces a symplectic manifold, and this
construction is inverse to the Lagrangian
blow-down. I will explain why these
constructions are functorial, that is,
defined on the corresponding symplectic
Teichmuller spaces. This is used to prove an
orbifold counterexample to a famous
conjecture, sometimes attributed to
Donaldson, who asked whether any symplectic
form on a K3 surface is compatible with a
Kahler structure. I will prove that
this is false for orbifolds: some K3
orbifolds admit symplectic forms not
compatible with any Kahler structure. The
same argument is used to produce a
countable family of Lagrangian spheres in a
K3 surface which are not Lagrangian
isotopic to special Lagrangian spheres. This
is a joint work with Michael Entov.



\medskip\noindent


Consider the total space of the cotangent
bundle of $\mathbb{C}P^1$.
It's $\mathcal{O}(-2)$.
Let $\gamma \in H^{0}(\mathcal{O}(2))$.
$\gamma$ defines a fiberwise linear function
on $T^*\mathbb{C}P$.

The blow-up of $\mathbb{C}^2/\pm 1$, which is
the orbifold $\mathbb{C}^2$ with
 a doble point, is the cotangent bundle $T^*
 \mathbb{C}P^1$.
That is,
$T^*\mathbb{C}P^1 \to \mathbb{C}^2/\pm 1$.
This is called {\it crepant resolution}.
Consider the $dx \wedge dy$, which is
holomorphically symplectic
on $\mathbb{C}^2$ and descends to the
quotient $\mathbb{C}^2/\pm 1$
as $\Omega_0$.

\begin{proposition}
\label{seminars-differential-geometry-proposition}
$\pi:T^*\mathbb{C}P^1 \to \mathbb{C}^2/\pm
1$.
Then $\pi^*\Omega_0$ is holomorphically
symplectic.
\end{proposition}

\begin{proof}
Actually, it gives the standard symplectic
form on $T^*\mathbb{C}^2$.
\end{proof}

\medskip\noindent
Now we shall define the Langrangian blow-up.
It's a procedure that starts with a
symplectic manifold
with double points and outputs a symplectic
manifold
without double points.

\begin{definition}[Lagrangian blow up]
\label{seminars-differential-geometry-definition-lagrangian-blow-up}
Let $(B,\omega)$ be the standard ball in
$\mathbb{R}^n$,
$dz_1 \wedge dz_2 +dz_3 \wedge dz_4$.
Consider the symplectic orbifold
$(B,\omega)/\pm 1$.
Consider $M$ be a symplectic orbifold
with double point singularities.
Darboux coordinates in a neighbourhood of $0$
with $(B,\omega)/\pm 1$ resolve
singularities.
\end{definition}

This definition has the caveat that there ar
e lots of choices involved.
We shall later that these choices do not
alter the result,
namely Theorem
\ref{seminars-differential-geometry-theorem-lagrangian-blow-up-well-defined}
.

To define the Lagrangian blow-down we first
recall

\begin{theorem}[Weinstein's Lagrangian
neighbourhood]
\label{seminars-differential-geometry-theorem-weins-lagra-neigh}
Let $X \subset M$ be a compact Lagrangian
submanifold in $(M,\omega)$.
Then there exists a neighbourhood $U$ of $X
\subset M$
which is symplectomorphic to a neighbourhood
of $X$
in $X \subset T^*M$.
\end{theorem}

\noindent
Let $S \subset M$ be a Lagrangian sphere
(diffeomorphic to sphere
and Lagrangian).

\begin{definition}
\label{seminars-differential-geometry-definition-lagrangian-blow-down}
The {\it Lagrangian blow-down} is obtained
from $M$
by removing $W \subset M$ and gluing $W_0$ in
its place,
where $W_0$ is such that

$$
\xymatrix{
W
\ar@{^{(}->}[r]\ar[d]&T^*\mathbb{C}P^1\ar[d]
\\
W_0\ar@{^{(}->}[r]&\mathbb{C}^2/\pm 1.
}
$$
\end{definition}

\medskip\noindent
Now we recall Teichmüller spaces.

Let $\Gamma(\Lambda^2M)$ be the space of all
2-forms on a manifold $M$.
The space of symplectic forms is a subspace
of the closed 2-forms on $M$,
which in turn is a subspace of
$\Gamma(\Lambda^2M)$.
Equip this space with the $C^\infty$
topology.
The {\it Teichmüller space} is
$\text{Symp}/\text{Diff}_0$.

Two symplectic structures are called {\it
isotopic} if they lie
in the same orbit of $\text{Diff}_0$.
The {\it period map}
$\text{Per}:\text{Teich}_s \to
H^{2}(M,\mathbb{R})$
maps a symplectic structure to its cohomology
class.

\begin{theorem}[Moser, 1965]
\label{seminars-differential-geometry-theorem-period-map-is-local-diffeo}
The period map is a local diffeomorphism.
\end{theorem}

\noindent
This implies that $\text{Teich}_s$ is smooth.

\begin{theorem}[Moser's trick]
\label{seminars-differential-geometry-theorem-mosers-trick}
Let $\omega_t$, $t \in S$ be a smooth family
of symplectic
structures, parametrized by a connected
manifold $S$.
Assume that the cohomology class
$[\omega_t]$ is constant.
Then all $\omega_t$ are diffeomorphic.
\end{theorem}

\noindent
This means that the fibers of $\text{Symp}\to
H^{2}(M)$
are the orbits (of the action of
$\text{Diff}_0$, I suppose).

\medskip\noindent
Consider a Lagrangian submanifold $L \subset
(M,\omega)$.
Let us assume for simplicity (though the
results
can also be obtained without this) that
$b_1(L)=0$,
i.e. first cohomology group is zero.

Let $\text{Symp}_L$ be the set of symplectic
forms vanishing on $L$,
and $\text{Diff}_0^L$ the connected component
of diffeomorphisms
preserving $L$.
Define
$\text{Teich}_L=\text{Symp}_L/\text{Diff}
_0^L$.

\begin{theorem}
\label{seminars-differential-geometry-theorem}
Let $V \subset H^{2}(M,\mathbb{R})$ be the
space of all
cohomology classes on $M$ which vanish on
$L$,
and $\text{Per}_L:
\text{Symp}_L/\text{Diff}^L_0 \to V$
take $(M,\omega)$ to the cohomology class of
$\omega$.
Then Per is a local diffeomorphism.
\end{theorem}

\begin{proof}
We used a version of Moser's trick: for
$\omega_t$ a family
of symplectic forms with $\omega_t|_{L}=0$,
with $[\omega_t]$ constant, then $\omega_t$
are $\text{Diff}_0^L$-isotropic.
\end{proof}

\noindent
As a corollary we obtain that this
Teichmüller space
classifies Lagrangian submanifolds:

\begin{lemma}
\label{seminars-differential-geometry-lemma-techm-class-lagra-subma}
Assuma that $L$ and $L'$ are Lagrangian
submanifolds in $(M,\omega)$,
and $\varphi \in \text{Diff}_0(M)$ a smooth
isotopy such that
$\varphi(L)=L'$. Then $(\omega,L)$ and
$(\varphi^*\omega,L)$
represent the same point in $\text{Teich}_L$
if and only if
$L$ and $L'$ are Lagrangian isotopic in
$(M,\omega)$.
\end{lemma}

\noindent
Again: points of this Teichmüller space are
isotopy classes
of Lagrangian submanifolds.

\medskip\noindent
\begin{theorem}
\label{seminars-differential-geometry-theorem-blow-down-diffeomorphism}
Let $L$ be a Lagrangian 2-sphere in a K3
surface  $M$
and $\hat{M}$ the Lagrangian blow-down of
$M$.
The blow-up/blow-down
is a diffeomorphism $\text{Teich}_L(M) \to
\text{Teich}(\hat{M})$.
\end{theorem}

\medskip\noindent
Now we prove the theorem which asserts that
the Lagrangian
blow-up is independent of choices.

\begin{theorem}
\label{seminars-differential-geometry-theorem-lagrangian-blow-up-well-defined}
Let $(\hat{M},\hat{\omega})$ be a orbifold.
Then its blow-up
$(M,L)$ defines a point in $\text{Teich}_L$
independent of the choices.
\end{theorem}

\begin{proof}
By taking local Darboux coordinates, and then
Moser's trick.
\end{proof}

\noindent
Actually, the same happens for blow-downs

\begin{theorem}
\label{seminars-differential-geometry-theorem-same-for-blow-down}
Let $(M,L,\omega)$ be a 4-manifold with a
Lagrangian sphere $L$, and
$(\hat{M},\hat{\omega)}$ be its Lagrangian
blow-down.

Then the corresponding point in
$\text{Teich}_{\hat{M}}$ is
independent from the choices made.
\end{theorem}

\begin{proof}
Similar.
\end{proof}

\medskip\noindent
A conjecture by Donaldson ways that all
symplectic structures on a K3
surface are compatible with a Kähler
structure. This conjecture
(although it's commonly believed to be true)
is still open.
When trying to prove that this would hold for
orbifolds,
Misha and collaborators realised it's
actually false.

Note: Seiberg-Witten invariants may be used
to tell whether a symplectic
form is compatible with a complex structure.

\begin{theorem}
\label{seminars-differential-geometry-theorem-conje-false-orbif}
There exists an orbifold symplectic K3
surface $M$ with a single
double point not admitting an orbifold Kähler
structures.
\end{theorem}

\begin{proof}
\noindent
Step 1. Consider
$$
\text{Teich}_{\text{Kähler-symplectic}}
(\hat{M})
=\{\eta \in H^{2}(\hat{M}):\eta^2>0\}.
$$
As follows from [Amerik-V., 2005], the
Teichmüller space of symplectic structures
compatible
with a hyperkähler structure is Hausdorff and
connected
(the same argument works for Kähler K3
orbifolds).

Essentially, if there wasn't any orbifold
as required, the Teichmüller space
classifying
Lagrangian submanifolds would be connected,
which
is impossible by [Seidel 2000].
\end{proof}

\noindent
In fact, Seidel's result is essentially
mirror symmetry.

\medskip\noindent
Now we discuss special Lagrangian
submanifolds.

\begin{definition}
\label{seminars-differential-geometry-definition-phase}
Let $(M,I,\omega,\Omega)$ be a Calabi-Yau
manifold, where
$\Omega \in \Lambda^{n,0}_I(M)$ is
nondegenerate and $\omega$ is Kähler.
Let $L$ be a Lagrangian submanifold.
Define the {\it phase} as
$$
\begin{aligned}
\text{phase}: L &\longrightarrow \Lambda^nL
\otimes \mathbb{C} \\
x &\longmapsto \Omega|_{T_xL}.
\end{aligned}
$$
\end{definition}

\noindent
A better definition is:
$$
\begin{aligned}
\text{phase}: L &\longrightarrow
S^1=\text{U}(1) \\
x &\longmapsto \frac{\Omega|_{T_xL}}{|\Omega
T_xL|}.
\end{aligned}
$$

\begin{definition}
\label{seminars-differential-geometry-definition-special-lagrangian}
A special Lagrangian submanifold is the phase
is constant.
\end{definition}

\noindent
Special Lagrangian submanifolds have
deformation space of
dimension 1, it's unobstructed. They are
calibrated, so minimal.

\begin{theorem}
\label{seminars-differential-geometry-theorem-lagrangians-of-k3}
$L \subset$ K3 is Lagrangian is isotopic to a
special Lagrangian
if and only if its blow-down is Kähler type.
\end{theorem}

\medskip\noindent
One of the basic results about special
Lagrangian
submanifolds is

\begin{lemma}[Hitchin]
\label{seminars-differential-geometry-lemma-hitchin}
If $(M,\Omega)$ is holomorphically symplectic
and $X \subset M$
is Lagrangian with respect to
$\text{Re}\Omega$ and $\text{Im}\Omega$
then $X$ is complex.
\end{lemma}

\noindent
As corollaries:

\begin{lemma}
\label{seminars-differential-geometry-lemma-corollary-hitchin}
Let $(M,I,J,K,g)$ be a hyperkähler manifold
of real dimension $4n$,
considered as a Kähler manifold
$(M,I,\omega_I)$ and $\phi:\Omega^n \in
\Lambda^{2n,0}(M,I)$ the corresponding
holomorphic volume form.
Consider a Lagrangian submaniold $S \subset
(M,\omega_I)$
such that $a \omega_J + b\omega_K|_S=0$ for
some nonzero real numbers
$a,b \in \mathbb{R}$ such that $a^2+b^2=1$.
Then $S$ is special Lagrangian, and,
moreover,
it is holomorphic with respect to the complex
structure
$-a K+b J$.
\end{lemma}

\begin{lemma}
\label{seminars-differential-geometry-lemma-corollary2-hitchin}
Let $\omega_I,\omega_J$ and $\omega_K$ on
$M$, with $\dim M=4$.
Let $S \subset (M,\omega_I)$ be Lagrangian.
Then $S$ is special Lagrangian if and only if
$S$ is complex analytic
on $a J + b K$ for two nonzero numbers  $a,b
\in \mathbb{R}$ such that
$a^2+b^2=1$.
\end{lemma}

\noindent
Now an auxiliary claim:

\begin{lemma}
\label{seminars-differential-geometry-lemma-k3-type}
$\Omega$ is holomorphic symplectic
\end{lemma}

\begin{theorem}
\label{seminars-differential-geometry-theorem-duplicate-2}
Let $S\subset M$ be a Lagrangian sphere in a
K3 surface.
Then  $S$ is isotropic to a special
Lagrangian if and only if
$M$ is of Kähler type.
\end{theorem}

\section{Systolic inequality for tori}
\label{seminars-differential-geometry-section-systolic-inequality}

\noindent
Dimitri Korshunov, IMPA.
Minimal Surfaces end of course presentation.
October 24, 2025.

\medskip\noindent
This is due to L. Guth.

\begin{definition}
\label{seminars-differential-geometry-definition-systole}
Given a Riemannian manifold,
a {\it systole} $\text{Sys}(M,g)$ is the
infimum of the lengths
of noncontractible paths.
\end{definition}

\begin{theorem}[Gromov]
\label{seminars-differential-geometry-theorem-gromov}
For any metric on a torus $T^n$,
$$
\text{Sys}(T^n,g) \leq C_n
\text{Vol}(T^n,g)^{\frac{1}{n}},
$$
where $C_n$ is a constant depending only on
the dimension
of the torus $n$.
\end{theorem}

\noindent
This also holds for $\mathbb{R}P^n$. We will
use induction,
for which we need the following stronger
statement:
given a number bounded by half the systole,
there exists a ball
which is {\it big}:

\begin{theorem}[Guth]
\label{seminars-differential-geometry-theorem-guth}
For $(T^n,g)$, given
$R<\frac{1}{2}\text{Sys}(T^n,g)$,
then there exists a point $p$ on $T^n$
such that the volume of ball of radius $r$
with center in $p$
satisfies
$$
\text{Vol}B(p,R) \geq C_nR^n.
$$
\end{theorem}

\noindent
It was conjectured by Gromov that $C_n=1$.

\medskip\noindent
Before stating the theorem that we will
prove, we introduce the
following definition that will allow us to
measure cohomological cycles:

\begin{definition}
\label{seminars-differential-geometry-definition-measu-cohom-cycle}
Let $\alpha \in H^{1}(M_g,\mathbb{Z}_2)$.
$$
L(\alpha)
=\text{inf}\{|\lambda|:\lambda\text{
1-cycles, }\alpha(\lambda)\neq 0\}.
$$
\end{definition}

\noindent
The theorem that we will prove, which implies
Theorem \ref{seminars-differential-geometry-theorem-guth}
is the following

\begin{theorem}
\label{seminars-differential-geometry-theorem-bound-classes}
There exists a constant $\varepsilon_n>0$
such that given a manifold $(M,g)$, and
$\alpha_1,\ldots,\alpha_n \in
H^1(M,\mathbb{Z}_2)$,
\begin{enumerate}
\item $L(\alpha_i)>2R$
\label{seminars-differential-geometry-item-positive}
\item
\label{seminars-differential-geometry-item-cup} $\alpha_1 \cup  \alpha_2
\cup \ldots \cup \alpha_n \neq 0 \in
H^n(M,\mathbb{Z})$.
\end{enumerate}

\noindent
Then there exists $p \in M$ such that
$$
\text{Vol}B(p,R) \geq \varepsilon_nR^n.
$$
\end{theorem}

\noindent
In particular $T^n$ and $\mathbb{R}P^n$
satisfy (\ref{seminars-differential-geometry-item-cup}).

Before we proceed to the proof we need the
following definition.
Essentially we ask that the surface minimizes
volume in its homology
class up to some constant $\delta$.

\begin{definition}
\label{seminars-differential-geometry-definition-delta-minimal-surface}
A  {\it $\delta$-minimal surface} is an
embedded hypersurface $Z$
such that for any $Z'$ such that $[Z'-Z] = 0$
in  $H_{n-1}(M,\mathbb{Z}_2)$,
then  $\text{Vol}Z < \text{Vol}Z'+\delta$.
\end{definition}

\noindent
Moreover, we need the following

\begin{lemma}[Stability]
\label{seminars-differential-geometry-lemma-stability}
Let $(M,g)$ de closed Riemannian manifold and
$\alpha \in H^1(M,\mathbb{Z}_2)$.
Suppose $Z$ is Poincaré dual to $\alpha$ and
that $Z$ is $\delta$-minimal.
If $R<\frac{1}{2}L(\alpha$) then the
following
holds for any $p \in M$:
$$
|Z \cap \text{Vol}B(p,\frac{R}{2}|\leq
\frac{2}{R}|\text{Vol}B(p,R)|+\delta.
$$
\end{lemma}

\begin{proof}[Proof of the stability lemma]
\medskip\noindent
Step 1. Using commas formula, namely
$$
\text{Vol}(B_R)=\int_0^R S_tdt,
$$
we see that there exists $t$ such that
$\frac{R}{2}<t<R$ and
$$
|S(p,t)|_R< \frac{2}{R}|\text{Vol}B(p,R)|.
$$

\medskip\noindent
Step 2. Consider
$[Z \cap B(p,t)] \in
H^{n-1}(B(p,t),S(p,t),\mathbb{Z})$.
This can be interpreted geometrically since
we are using relative homology.

\medskip\noindent
Step 3. We claim that $[Z \cap B(p,t)]=0$.
By contradiction,
if it weren't trivial, then there would exist
a cycle $H_1(B,\mathbb{Z})$
such that $Z' \cap \gamma \neq 0$ since
Poincaré pairing is nondegenerate ---
namely any cocycle that is not zero pairs
nontrivially with some cycle.
Then we will just split the cocycle in very
small parts,
and we will reach a contradiction with the
fact that
$R<\frac{1}{2}L(\alpha)$.

We will state this as another lemma.
Essentially it says that any
cycle in homology can be decomposed as a sum
of cycles
with small lengths.

\begin{lemma}[Curve factoring, Gromov]
\label{seminars-differential-geometry-lemma-curve-factoring-gromov}
Let $(M,g)$ be a Riemannian manifold and
$\gamma$ a cycle in $B(p,t)$.
Then $[\gamma]=\sum \sigma_i$ is such that
 $L(\sigma_i)<2t+a$ for all $i$.
\end{lemma}

\begin{proof}
We spit a cycle (a curve) into some quantity
of intervals
such that each has length smaller than $a$,
then add paths joining the center point of
length $t$,
so we count $t$ twice: to reach the point and
to get away from it.
\end{proof}

\noindent
We apply this to $\gamma$ and reach a
contradiction
with the hypothesis that
$R<\frac{1}{2}L(\alpha)$ making $a$ smaller
and
smaller.

\medskip\noindent
Step 4. Finally we use the minimality of $Z$
and the property
derived from Commas lemma to arrive at a
contradiction.
\end{proof}

\noindent
Now we prove Theorem
\ref{seminars-differential-geometry-theorem-bound-classes}:

\begin{proof}
By induction. The case $n=1$ is trivial.
Now take the last one of our classes,
$\alpha_n$.
Let $Z$ be Poincaré dual of $\alpha_n$ and
$\delta$-minimal.
Notice that $\alpha_1\cup \ldots \cup
\alpha_{n-1}\neq 0$
since if it were, then $\alpha_1\cup
\ldots\cup \alpha_n$ would also vanish.

By induction hypothesis we get a point $p \in
Z$
such that $\text{Vol} B_Z(p,\frac{R}{2}) >
R^{n-1}$.
This means that , $|Z \cap
B(p,\frac{R}{2})|>\varepsilon_n(R/2)^{n-1}$.

By the Stability Lemma \ref{seminars-differential-geometry-lemma-stability},
$$
\begin{aligned}
\frac{2}{R}|B(p,R)|&>\varepsilon_n\frac{R^{n-
1}}{2^{n-1}}-\delta\\
\implies
|B(p,B)|&>\frac{\varepsilon_n}{2^n}R^n-
\frac{\delta
R}{2}
> \frac{\varepsilon_n}{2^n}R^n.
\end{aligned}
$$
\end{proof}

\section{From higher Lie groupoids to higher
Lie algebroids}
\label{seminars-differential-geometry-section-highe-group-highe-algeb}

\noindent
Matias del Hoyo, UFRJ.
Symplectic Geometry Seminar, IMPA.
October 29, 2025.

\medskip\noindent
{\bf Abstract.}
This is a report of an ongoing collaboration
with A. Cabrera, where we develop a
higher Lie theory. Starting with simple,
concrete examples, we present the
differentiation functor in a geometric,
explicit way, using simplicial methods
and differential graded algebras. We describe
how our construction extends the
classical Lie functors, relate our project
with previous works, and show a van
Est theorem linking the cohomology at the
global and infinitesimal levels.

\medskip\noindent

Upshot: Integration-differentiation relation
between Lie groups and Lie algebras
in higher versions. Groupoids are the
counterparts of Poisson manifolds, and in
the next level there are 2-groupoids and
Courant algebroids. How to integrate a
2-groupoid to a Courant algebroid? This is an
open question which motivated this
work. See arXiv
\href{http://arxiv.org/abs/2309.14105}%
{2309.14105}.

\medskip\noindent
Let $M$ be a smooth manifold.

\begin{remark}
\label{seminars-differential-geometry-remark-smoot-funct-ring-fully-faith}
$M \mapsto C^\infty M$ is fully faithful.
\end{remark}

\noindent
Which allows us to think of manifolds as
algebraic objects.
Notice that
$\Omega^kM=\{k\text{-forms}\}
=\Lambda^k_{C^\infty M}\text{Der}(C^\infty
M,C^\infty M)$,
which is (algebraically) locally spanned
by the exact 1-forms.

Recall that the exterior differential is
the only operator $d:\Omega^kM \to
\Omega^{k+1}M$
such that $d(\alpha \wedge \beta)=d\alpha
\wedge \beta +
(-1)^{|\alpha|}\alpha \wedge d\beta$
and that $df(X)=Xf$.

This makes $(\Omega M,\wedge,d)$ a
differential graded algebra.
Recall that $\text{dom}(X,Y)=X\omega(Y)-Y
\omega(X)-\omega([X,Y])$.
We also have de Rham cohomology.

Consider a submanifold $S \subset M$. Recall
that $\mathcal{I}_S/\mathcal{I}^2_S$ is
essentially
the sections of the normal bundle.
In particular, for $M \subset M \times M$
we can see that the embedding vanishes on the
diagonal,
so there is no linear part in the Taylor
expansion,
and the second term is in fact the de Rham
differential
$df=\pi_1^*f-\pi_2^*f \text{ mod
}\mathcal{I}^2_{\Delta(M)}$.

In fact, this construction can be extended to
a simplicial manifold
(which I think is just a presheaf valued on
category of manifolds)
using projections and diagonals,
$$
\xymatrix{
\cdots & M\times M \times M
\ar @< 3pt> [r]
\ar @< 0pt> [r]
\ar @<-3pt> [r] & M \times M \ar @< 2pt> [r]
\ar @<-2pt> [r]& M.
}
$$
To this we can apply the $C^\infty$ functor
and obtain
a cosimplicial ring (not a dga?).
The idea here
is that de Rham differential is an
alternating sum of pullbacks;
as other differentials will also be.

\medskip\noindent
Now observe that for a finite-dimensional
vector space $V$,
$$
\left\{\substack{\text{Lie bracket} \\
V \times V
\xrightarrow{[\cdot,\cdot]}V}\right\}
\leftrightarrow
\left\{\substack{\text{differential on} \\
\Lambda V^*}\right\},
$$
making $\Lambda V^*$ into a differential
graded algebra.

Recall that given a Lie algebra
$\mathfrak{g}$,
we have $(\Lambda^\bullet \mathfrak{g}^*
,d)=\text{CE}(\mathfrak{g})$,
the Chevalley Eilenberg algebra,
which gives Lie algebra cohomology.

\medskip\noindent
[Missing part on Simplicial approach,
exponential map,
Simplicial manifold nerve]

\medskip\noindent
Lie algebroid: $\Delta \to M$ vector bundle
with $[\cdot,\cdot]$ on $\Gamma A$,
$A \xrightarrow{\rho}TM$ Leibniz.

\begin{remark}
\label{seminars-differential-geometry-remark-lie-algebroids-differentials}
$$
\left\{\substack{\text{Lie algebroids on} \\
E \to M}\right\}
\leftrightarrow
\left\{\substack{\text{differentials} \\
\text{on }\Lambda\Gamma E^*}\right\}.
$$
\end{remark}

\noindent
This gives a differential graded algebra
$$
\xymatrix{
\Delta\ar[r]&M\ar[r]&\text{CE}(A)
}
$$
and
$$
\xymatrix{
C^\infty M\ar[r]^{\tilde{\rho}}&\Gamma
E^*\ar[r]^{\widetilde{[\cdot,\cdot]}}
&\Gamma E^* \wedge \Gamma E^*\ar[r]&\cdots.
}
$$

From Lie groupoids to Lie algebroids:
Lie functor/differentiation.

$$
\Omega(G)=\Gamma \Lambda^\bullet
\underbrace{\supset}
_{\text{sub dga}}
\Gamma \Lambda^\bullet T^s
\xrightarrow{-/G}\text{CE}(A_G).
$$
There's the following ecosystem around Lie
groupoids:
$$
\xymatrix{
&\text{Lie groupoids}
\ar[dl]\ar@{->>}[d]\ar[dr]^{\text{nerve}}\\
\text{Lie algebroids}&
\substack{\text{Differentiable stacks} \\
\supset \text{orbifolds}}
&\text{simplicial manifolds.}
}
$$

\medskip\noindent
Lie groupoids are models for stacks,
which are geometric objects solving
classification
problems with objects identified up to
isomorphism.
Relaxing the identification to a weaker
notion
of equivalence we must introduce higher
versions of stacks.

Moving higher,
$$
\xymatrix{
\text{Lie groupoid}\ar[r]\ar[d]_{\text{Lie}}
&\text{Simplicial manifolds}\ar[d]^?\\
\text{Lie
algebroids}\ar[r]&\text{dg-manifolds}
}
$$
dg algebra is a  dga $(A=\bigoplus_{k=1}^n
A_k \to M,\wedge,d)$
starting at $C^\infty M$, and continuing
until you get to
the sections, namely
$$
\xymatrix{
C^\infty M \ar[r]_d & A_2 \ar[r]& \cdots &
\cong \Gamma A.
}
$$
[dH-Cabrera] provide the solution to this
problem,
namely what to put in the right vertical
arrow:

\begin{theorem}[dH-Cabrera]
\label{seminars-differential-geometry-theorem-dh-cabrera}
$$
\frac{(C^\infty_N(G),\cup,\sum(-1)A^*_i)}{J+
\delta
J}
=\text{CE}(A_G)
$$
where $A_G$ is a higher Lie algebroid
and $J$ is defined by: $f$ is in $J$ if it
has higher vanishing.
\end{theorem}


Groupoids are the counterparts of Poisson
manifolds,
and in the next level there are 2-groupoids
and Courant algebroids.
How to integrate a 2-groupoid to a Courant
algebroids?
This is an open question which motivated this
work.

\section{Hyperbolic surfaces and Mirzakhani's
curve counting}
\label{seminars-differential-geometry-section-hyper-surfa-mizak-curve}

\noindent
Viveka Erlandson, University of Bristol, UK.
8th Brazilian School of Dynamical Systems,
IMPA.
November 4, 2025.

\medskip\noindent
{\bf Abstract.}
It’s a classical problem to study the growth
of the number of periodic orbits of
bounded length $L$ under a geodesic flow on a
manifold. When it comes to
(closed) hyperbolic surfaces Huber proved
that the number is asymptotic to $e^L
L$ and this results has since been
generalized in many directions, maybe most
famously by Margulis to negatively curved
manifolds as well as more general
flows. However, a fundamentally different
result is due to Mirzakhani who
counted the subset of those closed geodesics
on hyperbolic surfaces which have
no self-intersections (or more generally,
those inside a fixed mapping class
group orbit) and showed that in this setting
the growth is asymptotic to a
polynomial in L of degree only depending on
the topology of the surface. In this
mini-course we will give the necessary
background on hyperbolic surfaces to
understand these results, give an idea of the
proof of Mirzakhani’s theorem and
introduce some tools used such as measured
laminations and train tracks. Time
allowing we present some generalizations (for
example, to non-simple closed
geodesics or to various metrics) and
applications to the distribution of closed
geodesics.

\medskip\noindent
Our goal is to count certain closed geodesics
on hyperbolic surfaces.
By ``certain'' we main, for example, simple
or fixed topological type;
we will define these properly. Also, we may
treat more general
types of curves than only hyperbolic, e.g.
combinatorial, algebraic.
As ``closed geodesics'', we mean free
homotopy classes of curves.

As historical note,
results on coarse polynomial growth for
simple curves
were are due to Birman-Serres, Rees in the
80's.
In 2001, McShane-Runn studied asymptotic
growth
for simple curves on punctured torus.
Asymptoric growth on any hyperbolic
surface was proved by Mirzakhani for simple
curves in 2008
and then in 2016 for any topological type of
curve.

Here's an outline of this minicourse:
\begin{itemize}
\item Introduction, motivation.
\item Measured laminations, train tracks,
Thurston measure.
\item  (Geodesic) currents.
\item Curve counting theorem.
\end{itemize}

\medskip\noindent
Throughout $S$ will be a topological
(connected, orientable
of finite type, i.e. finite genus). Let
$g<\infty$ be the genus
and $r< \infty$ the number of punctures.
We assume that $\chi(S)=2-2g-r<0$,
which is equivalent to saying that $S$
admits hyperbolic metrics.
We will not consider the pair of pants,
since there's only three simple closed
geodesics on it,
making it uninteresting for our means.

Let $X$ be a hyperbolic surface homeomorphic
to $S$.
Consider the universal cover
$\tilde{X}=\mathbb{H}^2$ of $X$,
and
$\pi_1(X)=\Gamma\underbrace{<}
_{\substack{\text{discrete}
\\ \text{free}}}
\text{PSL}(2,\mathbb{R})=\text{Isom}^+
(\mathbb{H}^2)$.
Recall that a $4g$-gon yields a genus $g$
surface $S_g$ by
adequate identification of its sides.

For us, a {\it curve} $\gamma$ is a free
homotopy class
of a closed immersion $\mathbb{S}^1
\hookrightarrow S$,
which are closed $X$-geodesics.
In turn, a {\it multi-curve}
is a formal sum $\sum_{i=1}^na_i\gamma_i$
for $\gamma_i$ distinct curves and $a_i>0$,
$a_i \in \mathbb{Q}$.
The {\it length} of a curve $\gamma$ is
$\ell_X(\gamma)$ is the length of
$X$-geodesic representing $\gamma$,
and the length of a multi-curve is the sum of
the lengths
of the curves in the formal sum.

We won't use this: we can think of $X$ as
an element of Teichmüller space
$\mathcal{T}(S)\cong \mathbb{R}^{6g-6+2r}$
parametrizing complex structures on $S$.

\medskip\noindent
Notice that given $X \in\mathcal{T}(S)$,
the number of curves with length bounded by a
given number $L$,
$\#\{\gamma|\ell_X(\gamma)\leq L\}$,
is finite. This can be checked going to the
universal cover.
However, as $L \to \infty$, the number also
tends to infinity.

Huber, Hekhal (59' or later) showed that if
$X$ is hyperbolic
on $S_{g,r}$, then
$\#\{\gamma|\ell_X(\gamma)\leq L\}\sim
\frac{e^L}{L}$
asymptotically, that is,
$$
\lim_{L\to \infty}\frac{\#
\{\gamma|\ell_X(\gamma)\leq L\}}{e^L/L}=1.
$$

\medskip\noindent
What is we restrict to a fixed topological
type?

Recall that the {\it mapping class group} of
$S$ is
$\text{Map}(S):=\text{Homeo}^+(S)/
\text{homotopy}$.
Then $\gamma$ and $\gamma_0$ are of the same
(topological)
type if $\gamma \in \text{Map}(S)\gamma_0$.

We define the {\it (geometric) intersection}
of $\alpha,\beta$
distinct curves by
$$
i(\alpha,\beta)=\text{min}\{\#(\alpha' \cap
\beta'|
\substack{\alpha'\text{ homotopic to }\alpha
\\
\beta'\text{ homotopic to }\beta}\}.
$$
(In fact, this minimum is always realized by
geodesics,
but it's more a topological definition.)
The {\it self-intersection} of a curve is
$$
i(\alpha,\alpha)=\frac{1}{2}\text{min}
\{\#(\alpha' \cap
\alpha''|\alpha',\alpha''\text{
homotopic to }\alpha\}.
$$

We say $\gamma$ is {\it simple} if
$i(\gamma,\gamma)=0$.

It is known that $i(\cdot,\cdot)$ is
$\text{Map}(S)$-invariant,
and that $\{\gamma:i(\gamma,\gamma)=k\}$ is a
finite union
of topological types.

\begin{theorem}[Mirzakhani, '08, '16]
\label{seminars-differential-geometry-theorem-mirzakhani}
Let $X$ be a hyperbolic surface, $X \cong
S_{g,r}$
and $\gamma_0$ any (multi-)curve.
Then
$$
\#\{\gamma\text{ type
}\gamma_0|\ell_X(\gamma)\leq L\}
\sim CL^{6g-6+2r},
$$
where $C>0$,
$C=\frac{c(\gamma_0)}{b_\gamma}B(X)$.
\end{theorem}

\noindent
As a corollary, we obtain
$$
\begin{aligned}
\#\{\gamma\text{ simple}|\ell_X(\gamma)\leq
L\}&
\sim\text{const.}L^{6g-6+2r}\\
\#\{\gamma|i(\gamma,\gamma)\leq
k,\ell_X(\gamma)\leq L\}
&\sim (\text{const. depending on
$k$})L^{6g-6+2r}.
\end{aligned}
$$

\medskip\noindent
As an example, consider the torus with simple
(multi-)curves.
If we put a puncture it becomes a hyperbolic
curve.
Forgetting the puncture, we have an Euclidean
flat torus,
of which we know the geodesics: they are
straight lines
with rational slope in the plane (before
quotienting).
We see that the amount of geodesics of length
bounded by $L$
is $\frac{1}{2}\mathbb{Z}^2 \cap B((0,0),L)$,
which is in
fact $\sim \pi L^2$.
A result by Mc-Shane-Rivn, 2001 shows that
 $$
\{\text{(primitive) simple curves}\}\sim
\frac{\pi}{6}L^2.
$$
Where primitive (most likely) excludes
counting
a curve twice as a multicurve.


\medskip\noindent
Let us introduce a measure on $\mathbb{R}^2$.
Let $L>0$. Then
$$
m^L=\frac{1}{L^2}\sum_{p \in
\mathbb{Z}^2}\delta_{\frac{1}{L}p}
\xrightarrow{L \to
\infty}\text{Leb}_{\mathbb{R}^2}.
$$
For $p=(a,b)$,
$\frac{1}{L}p=(\frac{1}{L}a,\frac{1}{L}b)$.
$$
\begin{aligned}
U_{\mathbb{E}}&=\{x \in \mathbb{R}^2:|x|\leq
1\}\\
m^L(U_{\mathbb{E})}&=\frac{1}{L^2}\#\{p \in
\mathbb{Z}^2
\left|\frac{1}{L}p\right|\leq 1\}\\
&=\frac{1}{L^2}\#\{p \in \mathbb{Z}^2:|p|\leq
L\}\\
&=\frac{1}{L^2}\#\mathbb{Z}^2 \cap
B_{\mathbb{R}^2}((0,0),L)\\
&\to\text{Leb}_{\mathbb{R}^2}(U_{\mathbb{E})}
=\pi.
\end{aligned}
$$

\medskip\noindent
Let's go back to study the geodesics on a
torus.
First we recall a fact, which is that
on any hyperbolic surface $S$ there exists a
compact $K \subset S$ such that all
simple closed geodesics stay in $K$.

Now consider a geodesic given by a curve of
irrational slope,
which is dense in the torus. However,
it is not dense on a hyperbolic torus
$S_{1,1}$.
We may approximate this geodesic by $q_n \in
\mathbb{Q}$ with
$q_n \to \alpha$. Each $q_n \rightsquigarrow
\gamma_n$
is a closed geodesic on $S_{1,1}$.
The lamination in this case is the Hausdorff
limit
of the $\gamma_n$.

\begin{definition}
\label{seminars-differential-geometry-definition-lamination}
A {\it lamination} $L$ on a (hyperbolic)
surface $S$
is a compact subset of $S$ consisting of
simple, pairwise
disjoint complete geodesics.

A {\it measured lamination} $\lambda$ is
$\lambda=(L,\lambda)$ is a lamination $L$
together with a {\it transverse measure}
$\lambda$:
a measure is assigned to each arc $I$
transverse to $L$.

We write $ML(S)$ for the space of all
measured laminations on $S$.
\end{definition}

\noindent
While the definition of lamination depends on
a choice
of hyperbolic metric on the surface, the
different $ML(S)$
that we obtain with different metrics are
actually homeomorphic,
and we thus do not distinguish between them.

As an example, every simple multi-curve is a
measured lamination.

Morally we think of measured laminations as
limits of simple multicurves:
simple multi-curves are dense in $ML(S)$.

\medskip\noindent
Now we will discuss train tracks, which will
give us a combinatorial model for measured
laminations.

\begin{definition}
\label{seminars-differential-geometry-definition-train-track}
A {\it train track} $\tau$ on $S$ is an
embedded (trivalent)
graph such that at every vertex the three
half edges
meet with the same tangent line and such that
the complementary region (cc. $S \setminus
\tau$)
are not disk
(with a certain caveat: it can be a
contractible
region as long as the vertices cannot
preserve
it's smooth incidence structure
along the homotopy),
a cyllinder (ring) nor a punctured disk.

A {\it set of weights} on $\tau$ is an
assignment
of a non-negative real number $w_e$
to each edge $e$. They satisfy the {\it
switch equations}
if at every vertex we have $w_a=w_b+w_c$
(where $b$ and $c$ ``stem'' from $a$).

We denote $w(\tau)$ the solutions to the
switch equations
and $w_{\mathbb{Z}}(\tau)$ the integral
solutions.

We say a measured lamination $\lambda \in
ML(S)$
is {\it carried} by $\tau$ ($\lambda < \tau$)
if its support can be homotopied into $\tau$.
We denote
$$
\begin{aligned}
ML(\tau)&=\{\lambda \in ML(S):\lambda <
\tau\}\\
ML_{\mathbb{Z}}(\tau)&=\{\lambda \in
ML_{\mathbb{Z}}(S):\lambda <\tau\}\\
ML_{\mathbb{Z}}(S)&=\{\lambda \in
ML(S):\substack{\text{
simple multi-curve} \\ \text{with integral
coefficients}}\}.
\end{aligned}
$$
\end{definition}

\medskip\noindent
Now we will explain the fact that there are
one-to-one correspondences
$$
\begin{aligned}
ML(\tau)&\xymatrix{\ar@{<->}[r]&}w(\tau)\\
ML_{\mathbb{Z}}(\tau)&\xymatrix{\ar@{<->}[r]&
}w_{\mathbb{Z}}(\tau).
\end{aligned}
$$

\noindent
In a simple case where we pick a multicurve
$\lambda \in ML_{\mathbb{Z}}(\tau)$, the
associated measure
is just counting the intersection number
of the curve with a perpendicular segment
near
the train track edge.

The following argument appears to be a
converse construction.
Given a solution to the switch equations,
i.e. an element in $\omega(\tau)$, we can
produce
a multicurve $\lambda \in
ML_{\mathbb{Z}}(\tau)$.
(The method is essentially putting as many
strands
as the weights say near the vertex;
this yields formal sums of curves.)
This also works for formal sums with rational
coefficients.
And for non-negative real coefficients, we
take a thickening of $\tau$.

\begin{theorem}[Thurston]
\label{seminars-differential-geometry-theorem-thurston-train-tracks}
The weights on $\tau$ uniquely determine a
transverse measure on $L$.
\end{theorem}

\medskip\noindent
Facts:
\begin{enumerate}
\item Every $\lambda \in MS(S)$
is carried by some train track $\tau$ and
determines
a unique point in $\omega(\tau)$.
\item There are finitely many train tracks
$\tau_i$ such that
every $\lambda \in ML(S)$ is carried by at
least $\tau_i$.
From this it follows that we can write
$ML(S)=ML(\tau_1)\cup  ML(\tau_2)\cup
\ldots\cup ML(\tau_n)$
and $ML(\tau)\approx \omega(\tau)\approx
\mathbb{R}^{E-V}$
(under very mild conditions on the train
tracks).

Here is an exercise: suppose $S$ is closed.
Show that $E-V \leq  6g-6$, and
$E-V=6g-6$ if $\tau$ is maximal (which means
that all complementary
regions are triangles).

\item $ML(\tau_i)\cap ML(\tau_j)$ for $i \neq
j$ has dimension
$<6g-6$.
\end{enumerate}

Given those facts,
$$
\begin{aligned}
ML(S)&=ML(\tau_1)\dot\cup \ldots \dot\cup
ML(\tau_n)\\
&=\omega(\tau_1)\dot\cup \ldots\dot\cup
\omega(\tau_n).
\end{aligned}
$$

\noindent
In particular, $ML(S)\approx
\mathcal{E}^{6g+6L+2r}$
($\omega(\tau_1)\approx\mathbb{R}^{6g+6L-2r}
$).
We may now introduce a measure on the space
of laminations:

\begin{definition}
\label{seminars-differential-geometry-definition-thurston-measure}
The {\it Thurston measure} on $ML(S)$ is
$$
\lim_{L\to\infty}\frac{1}{L^{6g-6+2r}}
\sum_{\lambda
\in ML_{\mathbb{Z}}(S)}
\delta_{\frac{1}{L}\lambda}
$$
\end{definition}

\noindent
One checks that this is a Lebesgue measure
(in the limit) as follows:
$$
\begin{aligned}
ML_{\mathbb{Z}}(S)&=\bigcup_{i=1}
^nML_{\mathbb{Z}}(\tau_i)\\
&\frac{1}{L^{6g-6+2r}}\sum_{\lambda \in
ML_{\mathbb{Z}}\tau_1}
\delta_{\frac{1}{L}\lambda}\\
&\frac{1}{L^{6g-6+2r}}\sum_{(\omega)\in
\omega_{\mathbb{Z}}(\tau_i)}
\delta_{\frac{1}{L}\omega}\\
&\frac{1}{L^{6g-6+2r}}\sum_{p \in
\mathbb{Z}^{6g-6+2r}}
\delta_{\frac{1}{L}p}
\end{aligned}
$$

\noindent
on $\mathbb{R}^{6g-6+2r}$.

Notice that there is an action of the mapping
class group
on the space of laminations: a multicurve is
taken to another multicurve.

The following theorem was used for proving
the speaker's results:

\begin{theorem}[Masur '85]
\label{seminars-differential-geometry-theorem-masur-85}
$m_{Thu}$ is $\text{Map}(S)$-invariant
and is ergodic with respect to
$\text{Map}(S)$.
\end{theorem}

\medskip\noindent
Now we discuss (geodesic) currents.
Consider de space $\mathcal{G}(\tilde{S})$
of unoriented geodesics on
$\tilde{S}=\mathbb{H}^2$
as $(S^1\times S^1\setminus\Delta)/(a,b)\sim
(b,a)$.
This  is well defined (even topologically,
i.e.
independently of the choice of metric)
since two geodesics on Poincaré plane are
determined by their endpoints.

\begin{definition}
\label{seminars-differential-geometry-definition-geodesic-current}
A {\it (geodesic) current} is a
$\pi_1(S)$-invariant
Radon measure on $\mathcal{G}(S)$.
\end{definition}

\begin{example}
\label{seminars-differential-geometry-example-multicurves-are-currents}
Any (multi-)curve is a current.
\end{example}

\noindent
We shall see that not only curves, but also
measured laminations
and in fact the Teichmüller space are
contained in currents:
$$
\begin{aligned}
\{\text{curves}\}&\subset C(S)\\
ML(S)&\subset C(S)\\
\mathcal{T}(S)&\subset C(S)
\end{aligned}
$$

Let $X$ be a hyperbolic metric on $S$.
Let $U=\{\lambda \in
ML(S):\ell_X(\lambda)\leq 1\}$.
If we let $m^L(U_X)=\frac{1}{L^{6g-6}}
\#\{\gamma \in
ML_{\mathbb{Z}}(S)|\ell_X(\gamma)\leq  L\}$,
then $m^L(U_X) \to m_{\text{Thu}}(U_X):=B(X)$
 as $L \to \infty$.

\medskip\noindent
In the following, we think os simple
curves as measured laminations and currents.
Recall that a measured lamination is
a union of disjoint curves.
$$
\xymatrix@C=4em{
ML_{\mathbb{Z}}(S)\ar@{^{(}->}[r]&\overline{
\{\text{simple
multicurves}\}}\ar@{^{(}->}[d]\ar@{^{(}->}[r]
^-{\substack{\text{equality for} \\
\text{closure only!}}}
&ML(S)\ar@{^{(}->}[d]\\
&\overline{\{\text{multi-curves}\}}\ar@{^{(}->}[r]^-{
\substack{\text{equality for} \\
\text{closure only!}}}
&CS=\{\text{geodesic currents}\}.
}
$$

\medskip\noindent
Facts:
\begin{enumerate}
\item (Banahon.) The geometric intersection
number of curves
extends to
$$
i:C(S) \times C(S) \to \mathbb{R}_{\geq 0}
$$
continuously, bi-linear,
$\text{Map}(S)$-invariant.

\item $ML(S)=\{\mu \in
C(S):i(\mu,\mu)=0\}\rightsquigarrow
m_{\text{Thu}}$ measure on $C(S)$
(supported on $ML(S)$).

\item $C_{\text{fin}}(S)=\{\mu \in C(S):
(\mu,\lambda)>0 \forall \lambda \in C(S)\},$
the {\it filling currents}, is an open dense
subset of $C(S)$.

\item (Banahan.) For all $X$ hyperbolic
metric on $S$,
there exists a unique $\lambda_X \in C(S)$
such that
$i(\lambda_X,\gamma)=\ell_X(\gamma)$ for all
curves $\lambda$
(Liouville current) $\rightsquigarrow $
$\mathcal{T}(S) \subset C(S)$.
This way we have a notion of hyperbolic
length of currents $\mu \in C(S)$:
$\ell_X(\mu)=i(\lambda_X,\mu)$.
\end{enumerate}

\begin{definition}
\label{seminars-differential-geometry-definition-nice-length-function}
A length function on curves on $S$
is {\it nice} if it extends to a function
$F:C(S) \to \mathbb{R}_{\geq 0}$ wich is
continuous, homogeneous
($F(c\cdot\mu)=cF(\mu)\;\forall c>0$,
positive ($F(\mu)>0\;\forall \mu\neq 0$.
\end{definition}

\begin{example}
\label{seminars-differential-geometry-example-nice-length-functions}
Following are examples of nice length
functions:
\begin{itemize}
\item Hyperbolic length.
\item (Otal.) Negatively curved Riemannian
metric.
\item (E-Parlier-Souto.) (Singular)
Riemannian metric.
\item (Duchen-Leinhein-Rafi.) Euclidean cone
surface.
\item (Banachan.) $\sigma \in C_{fill}(S)$,
$\ell_{\sigma}=(\sigma,M)$.
\item Word length w.r.t. finite generating
set $\pi_1(S)$.
\end{itemize}
\end{example}

\medskip\noindent
Let us go back to our original problem
of counting curves of a specific type.
Fix $\gamma_0$ a (multi-)curve (not
necessarily simple).
$\gamma$ is of type $\gamma_0$ if
$\gamma \in \text{Map}(S)\cdot \gamma_0$.

We want asymptotics of
$\#\{\gamma\text{ type
}\gamma_0:\ell_X(\gamma)\leq L\}$.
Fix $\gamma_0$ and $L>0$. Consider
$m^L_{\gamma_0}=\frac{1}{L^{6g-6}}
\sum_{\gamma\text{ type
}\gamma_0}\delta_{\frac{1}{L}\gamma}$
measure on $C(\gamma)$.

Note that $m^L_{\gamma_0}(U_X)=
\frac{1}{L^{6g-6}}\#
\{\gamma\text{ type
}\gamma_0:\ell_X(\gamma)\leq L\}$.

We are ready to state the main theorem:

\begin{theorem}[Mirzakhani '08, E-Souto '22]
\label{seminars-differential-geometry-theorem-mirzakhani-e-souto}
For any $\gamma_0$ on $S=S_g$
$$
m^L_{\gamma_0}\to\frac{c(\gamma_0)}{b_g}
m_{\text{Thu}},\qquad
L\to \infty.
$$
\end{theorem}

\noindent
Mirzakhani proved it for simple curves,
while we now presented a general version
using currents.
We write the constants as
$$
\begin{aligned}
c(\gamma_0)&=m_{\text{Thu}}(\{\lambda \in
ML(S)/\text{Stab}(\gamma_0)
|i(\lambda,\gamma_0)\leq 1\},\\
b_g&=\sum_{\alpha_0 \in
ML_\mathbb{Z}/\text{Map}(S)}c(\alpha_0).
\end{aligned}
$$

\noindent
Mirzakhani used another expression.
This one has the virtue of being purely
combinatorial.

Mirzakhani showed the following example.
Consider the genus 2 surface $S_2$.
Then
$$
\lim_{L \to \infty}
\frac{\#\{\text{non-separating
}\gamma|\ell_X(\gamma)\leq L\}}{
\#\{\text{separating
}\gamma|\ell_X(\gamma)\leq L\}}
=48.
$$
This may be interpreted by saying that when
taking a closed simple geodesic,
the probability of it being separating is
$\frac{1}{48}$.

\section{The unstable flow of a partially 
hyperbolic diffeomorphism with an 
application to measure rigidity}
\label{seminars-differential-geometry-section-unstable-flow-hyp-dyn}

\noindent
Bruno Santiago (UFF?)

Qui 21 mai 2026, 15:30

\medskip\noindent
{\bf Abstract.}
Let $f$ be a partially hyperbolic
diffeomorphism preserving a splitting
$
TM = E^s \oplus E^c \oplus E^u,
$
where the unstable bundle $E^u$ is
one-dimensional.
When $f$ is homogeneous, like the time-one
map of the geodesic flow on a surface of
constant negative curvature or a linear
automorphism of the torus, the unstable
foliation admits a natural parametrization
that makes invariant measures with
absolutely continuous conditionals
(the so-called $u$-Gibbs measures)
correspond exactly to invariant measures
of this flow.
How can this picture be generalized when
$f$ is not homogeneous?
In this talk, I will explain a construction
of a suitable suspension flow of $f$ for
which the unstable foliation becomes the
orbit foliation of a new flow, called the
unstable flow.
Ergodic $u$-Gibbs measures are in
one-to-one correspondence with measures
simultaneously invariant under both flows
and ergodic for the suspension.
As an application, I will deduce the
following measure rigidity statement:
if $E^c$ is expanding and the bundle
$E^s \oplus E^u$ is integrable, then any
$u$-Gibbs measure whose entropy exceeds its
unstable entropy is necessarily an SRB
measure.
This is joint work with Sébastien
Alvarez, Sylvain Crovisier, Martin
Leguil, and Davi Obata.

\begin{remark}[AI-generated overview]
\label{seminars-differential-geometry-remark-ai-overview-seminar-unstable-flow}
The talk was about translating a geometric
object into a dynamical one. We start with a
partially hyperbolic diffeomorphism
$$
f:M\to M.
$$
This means that there is a $Df$-invariant
splitting
$$
TM=E^s\oplus E^c\oplus E^u.
$$
Here $Df$-invariant means that
$$
Df_x(E_x^s)=E_{f(x)}^s,\qquad
Df_x(E_x^c)=E_{f(x)}^c,\qquad
Df_x(E_x^u)=E_{f(x)}^u.
$$
Thus $f$ moves points, while $Df$ moves the
corresponding tangent directions in a way
compatible with the splitting.

\medskip\noindent
The key simplifying hypothesis is that
$E^u$ is one-dimensional. Hence the unstable
distribution integrates to curves
$$
W^u(x),
$$
called unstable leaves. The first geometric
idea is that these leaves should not merely
be thought of as a foliation, but as the
orbits of a flow. Thus one wants to construct
an unstable flow $\varphi_t^u$ whose orbit
through $x$ is exactly $W^u(x)$.

\medskip\noindent
The nontrivial point is the choice of
parametrization
[because the passage from the orbits
to a vector field is not integration,
but differentiation, which in principle
is easy --- once you have
the right parametrization!]. 
A one-dimensional leaf can
be parametrized in many different ways, and
the dynamics of $f$ selects a preferred one.
If $x$ and $y$ lie on the same unstable leaf,
the function
$$
\rho_x^u(y)
$$
compares the expansion of $Df^n$ in the
unstable direction at $x$ and at $y$. Thus
$\rho_x^u$ is a dynamical density along
$W^u(x)$.

\medskip\noindent
This density defines a coordinate on the
unstable leaf:
$$
H_x^u(y)
=
\int_x^y
\rho_x^u(z)\,d\operatorname{Leb}_x^u(z).
$$
The inverse map
$$
\Phi_x^u=(H_x^u)^{-1}
$$
then gives the natural parametrization of the
unstable leaf. This is the geometric heart of
the construction.

\medskip\noindent
The measure-theoretic part enters because in
dynamics one studies not only where an orbit
goes, but also its long-time statistical
distribution. For a point $x$, the empirical
measures
$$
\mu_{n,x}
=
\frac{1}{n}
\sum_{j=0}^{n-1}
\delta_{f^j(x)}
$$
record the proportion of time the orbit of
$x$ spends in different regions of $M$.
Equivalently, for a measurable subset
$A\subset M$, one has
$$
\mu_{n,x}(A)
=
\frac{1}{n}
\#\{0\leq j<n:f^j(x)\in A\}.
$$
Thus $\mu_{n,x}(A)$ is the proportion of
time, up to time $n$, that the orbit of $x$
spends inside $A$. This is where time enters:
we are not measuring the size of $A$ itself,
but rather how frequently the forward orbit
$$
x,\ f(x),\ f^2(x),\ldots
$$
visits $A$.

\medskip\noindent
A $u$-Gibbs measure is an invariant measure
which is smooth in the unstable direction.
More precisely, if
$$
\{\mu_x^u\}
$$
is the disintegration of $\mu$ along unstable
leaves, then
$$
\mu_x^u\ll \operatorname{Leb}_x^u.
$$
So the measure may be globally complicated,
but leaf-by-leaf along $W^u$ it has a density
with respect to ordinary leafwise length.

\medskip\noindent
The main theorem of the talk says that such
$u$-Gibbs measures can be characterized after
passing to a suspension space. The original
map $f$ is discrete-time, while the suspension
turns it into a continuous-time flow
$\psi_t$. The theorem identifies $u$-Gibbs
measures with lifted measures which are
invariant both under the suspension flow
$\psi_t$ and under the unstable flow
$\varphi_t^u$.

\medskip\noindent
Finally, the application concerns SRB
measures. An SRB measure is stronger than a
$u$-Gibbs measure: it is meant to describe
the statistics seen by a positive-Lebesgue
set of initial conditions. In the theorem
from the end of the talk, if the center
direction is expanding and a $u$-Gibbs measure
has more entropy than what is accounted for
only by the unstable direction, namely
$$
h_\mu(f)>
\int_M \log \lambda_x^u\,d\mu(x),
$$
then the measure must be SRB.

\medskip\noindent
In the language of the Heap Project, this is
not like constructing a scheme from equations;
it is more like extracting an invariant
object from asymptotic data. Just as a Hilbert
polynomial records asymptotic information
about graded pieces, an invariant measure
records asymptotic information about the
distribution of an orbit. The slogan of the
talk is therefore:
$$
\text{unstable expansion}
\quad\leadsto\quad
\text{dynamical density}
\quad\leadsto\quad
\text{unstable flow}
\quad\leadsto\quad
\text{measure rigidity}.
$$
\end{remark}

\noindent
Let $S:M \to M$ be a $C^r(M)$ map.
Let $V$ be a vector field
inducing a foliation
whose leaves form a $C^r$ foliation
on $M$.
The idea is
that this dimension-1 
foliation is dilated by the
diffeomorphism.

\begin{example}
\label{seminars-differential-geometry-example-sl3-spectrum}
\begin{itemize}
\item $A \in \text{SL}_3(\mathbb{Z})$,
with spectrum
$\text{spec}(A)
= \{0<\lambda_A^s<1<
\lambda^c_A
<\lambda^u_A\}$.
$W_A^u
= $ folitation with orbits
of irrational flux $\varphi_t$
of $f_A:\mathbb{T}^3 \to \mathbb{T}^3$.
\end{itemize}
\end{example}

\begin{definition}[AI]
\label{seminars-differential-geometry-definition-u-gibbs-measure}
Let $f:M \to M$ be a partially hyperbolic
diffeomorphism, and let
$\mathcal F^u$ be its unstable foliation.
A measure
$$
\mu \in \mathcal P_f^{\operatorname{reg}}(u)
$$
is called a $u$-Gibbs measure if its
conditional measures along unstable leaves
are absolutely continuous with respect to
Lebesgue measure on the leaves.

More precisely, if
$$
\{\mu_x^u\}_{x \in M}
$$
denotes the disintegration of $\mu$ along
the unstable foliation, then
$$
\mu_x^u \ll \operatorname{Leb}_x^u
$$
for $\mu$-almost every $x \in M$.
\end{definition}

\begin{remark}
\label{seminars-differential-geometry-remark-stable-directions-dyn}
$$
TM = E^s \oplus E^c \oplus E^u.
$$
\end{remark}

\begin{remark}
\label{seminars-differential-geometry-remark-leb-ae-ai}
For Lebesgue almost everywhere,
$x \in M$, $n_k \to \infty$,
$$
\mu_{n_k,x}
:=
1 \sum_{n_k=0}^{n_k-1}
\delta_{f_k} \to \mu
$$
which implies that $\mu \in 
\text{Leb}^u$.
\end{remark}

\noindent
I'm very lost, here's AI-corrected
above remark:

\begin{remark}
\label{seminars-differential-geometry-remark-lebesgue-ae}
The phrase ``Lebesgue almost every $x$''
means almost every point with respect to
the Riemannian volume measure on $M$.

Given $x \in M$, its empirical measures are
$$
\mu_{n,x}
:=
\frac{1}{n}
\sum_{j=0}^{n-1}
\delta_{f^j(x)}.
$$
A measure $\mu$ is called physical if its
basin
$$
B(\mu)
:=
\left\{
x \in M :
\mu_{n,x} \to \mu
\right\}
$$
has positive Lebesgue measure, where the
convergence is weak-$*$ convergence of
probability measures.
\end{remark}

\begin{lemma}[Ledrappier density]
\label{seminars-differential-geometry-lemma-ledrappier}
$$
\frac{d\mu^u_x}{
d\text{Leb}_x^u}
=c(X)\rho_x^u
$$
where $\rho_x^u:w^u(x)\to \mathbb{R}$,
$$
\rho_x^u(y)
=
\lim_{n \to \infty} 
\frac{\|D\bar{f}^n(y)V^u\|}{
D\bar{f}^n(x) V^u}
$$
is the dynamic density function.
\end{lemma}
\begin{lemma}[Ledrappier density, AI]
\label{seminars-differential-geometry-lemma-ledrappier-density}
Let $\mu$ be a $u$-Gibbs measure. Then the
conditional measure $\mu_x^u$ is absolutely
continuous with respect to $\operatorname{Leb}_x^u$,
and its density is of the form
$$
\frac{d\mu_x^u}{d\operatorname{Leb}_x^u}(y)
=
c(x)\rho_x^u(y),
$$
where $c(x)>0$ is a normalizing constant and
$$
\rho_x^u:W^u(x)\to \mathbb R
$$
is given by
$$
\rho_x^u(y)
=
\lim_{n\to\infty}
\frac{
\|D\bar f^{-n}(y)|_{E_y^u}\|
}{
\|D\bar f^{-n}(x)|_{E_x^u}\|
}.
$$
Equivalently, in the one-dimensional unstable
case,
$$
\rho_x^u(y)
=
\lim_{n\to\infty}
\frac{
\|D\bar f^n(x)v_x^u\|
}{
\|D\bar f^n(y)v_y^u\|
},
$$
where $v_x^u\in E_x^u$ and $v_y^u\in E_y^u$
are unit vectors.
\end{lemma}

\medskip\noindent
{\bf Properties of EDD.}
\begin{itemize}
\item 
$$
\rho_y^(u)
=\frac{j}{\rho^u_x(y)}
$$
\item 
$$
z,y \in W^u(x)
\implies 
\rho^u_x(z)
=\rho_x^u(y)\rho_y^u(z)
$$
\item 
$$
\rho^u_{f(x)}(f(y))
=\frac{\lambda_x^u}{\lambda_y^u}
\rho_x^u(y).
$$
\end{itemize}

\medskip\noindent
{\bf Affine structure.} 
$$
H_x^u:W^u(x) \to \mathbb{R}
$$
$$
H_x^u(y)
=\int_x^y \rho_x^u(z)
d\text{Leb}_x^u(z)
$$
$$
\Phi_x^\mu
:=(H_x^u)^{-1}
:\mathbb{R} \to W^u(x).
$$

Change of coordinates:
$$
H_y^u \circ\Phi_x^u(t)
=\rho_y^u(x)H
+H_y^u(x).
$$
\begin{definition}[Unstable flux]
\label{seminars-differential-geometry-definition-unstable-flux}
$$
\begin{aligned}
\tilde{\varphi}_t:
M \times \mathbb{R} &\longrightarrow 
M \times \mathbb{R} \\
(x,\alpha) &\longmapsto 
(\Phi^u_x(e^{-x}t),
\alpha
+\text{log}\rho_x^u\Phi_t^u(e^{-\alpha}t))
\end{aligned}
$$
\end{definition}

\noindent
Let's check why this, in fact,
defines a flux:

\begin{lemma}
\label{seminars-differential-geometry-lemma-defines-a-flux}
$$
\tilde{\varphi}_t \circ \tilde{\varphi}_s
(x,\alpha)=\tilde{\varphi}_{t+s}(x,\alpha)
$$
\end{lemma}

\begin{proof}
Provided
\end{proof}

\begin{theorem}
\label{seminars-differential-geometry-theorem-gibbs-measure}
$\mu \in \text{Gibbs}_{\text{reg}}^u(f)$,
then there exists a unique
$$
\hat{\mu} 
\in \mathcal{P}_{\psi}^{\text{reg}}
(\hat{\mu}) \cap
\mathcal{P}_\varphi(\hat{\mu})
$$
\end{theorem}

$$
D(f)
=\{(x,\alpha):0\leq \alpha<\lambda_x^u\}
\subset\hat{\mu}
$$
$$
\hat{\mu}
=\frac{1}{
\int_M \text{log} \lambda_x^u
d\mu(x)}\mu\otimes\text{Leb}_{\mathbb{R}}
(\mathbb{T}^{-x}(A)).
$$

\medskip\noindent
{\bf Applications of Theorem
\ref{seminars-differential-geometry-theorem-gibbs-measure}}

\noindent
Back to the totally partially hyperbolic
context, suppose there is a decomposition
$$
TM^3 = E^s \oplus E^c \oplus E^u
$$
such that
$$
0 \leq \|Df(x)|_{E^s}\|
<I<\|Df(x)|_{E^c}\|
<\|Df(x)|_{E^u}\|.
$$
Then $f$ is Anosov partially hyperbolic.
Thus, there exist foliations
$W^s,W^c,W^u$ such that
$E^s,E^c$ and $W^u$ and
$W^{cs}, W^{du}$
$$
F^s \oplus E^c \oplus E^c \oplus E^u.
$$

\begin{definition}
\label{seminars-differential-geometry-definition}
$\mu \in \mathcal{P}_f(u)$
is SRB if all the conditionals
$$
\{\mu_x^{c\mu}\}_{x \in M}
$$
along $W^{c\mu}(\alpha)$
satisfy
$$
\mu_x^{c\mu}\ll \text{Leb}_x^{c\mu}.
$$
\end{definition}

\begin{theorem}[Sinoi-Rulh-Bouso]
\label{seminars-differential-geometry-theorem-srb}
There exists a unique SRB measure
$\mu$.
\end{theorem}

\noindent
Conjecture (ALOS, Katz): 
if $E^s \oplus E^u$
are not integrable then any
$$
\mu \in \text{Gibbs}_{\text{reg}}^{\mu}(f)
$$
is SRB.

\begin{theorem}[B]
\label{seminars-differential-geometry-theorem-gibbs-b}
If $E^s \oplus E^u$ is integrable
and
$$
\mu \in \text{Gibbs}_{\text{reg}}^u(u)
$$
satisfies
$$
h_\mu(f)
> \int \text{log}\lambda_x^u d\mu(x),
$$
then $\mu$ is SRB.
\end{theorem}

\medskip\noindent
{\bf Some ideas.}

 \noindent
$x \in M \mapsto \text{log}\lambda_x^u$
is $C^{1+\alpha}$.

Descição das $W^*$
$$
\widetilde{W}^s(x,\alpha)
=\{(x,\alpha+\theta_x^{s}(y)
:y \in W^s(x)\}
$$
where
$\theta(\alpha)=\text{log}\lambda_x^u$
and
$\theta_x^s(y)
=\sum_{\ell=0}^{\infty}
\theta(f^\ell y)
-\theta(f^\ell x)$.

\section{Weil--Petersson homeomorphisms 
and maximal surfaces in anti-de Sitter space}
\label{seminars-differential-geometry-section-weil-petersson}

\noindent
Graham Smith

May 21, 2026.

\medskip\noindent
{\bf Abstract.}
A classe de Weil--Petersson é uma classe
notável de homeomorfismos do círculo que se
encontra na fronteira de vários campos da
matemática contemporânea, incluindo a teoria
de Teichm\"uller, a teoria de
Schramm--Loewner, a geometria hiperbólica e
assim por diante.
Neste trabalho, apresentamos uma nova
perspectiva, estudando essa classe do ponto
de vista da geometria anti-de Sitter.
Mostramos que um homeomorfismo do círculo é
Weil--Petersson se e somente se seu gráfico
limitar uma superfície máxima completa,
semelhante ao espaço, em
$\mathrm{AdS}^{2,1}$ de área finita
renormalizada.

Este é um trabalho conjunto com F. Diaf,
A. Moriani, R. Smai e E. Trebeschi.
2604.17804.

\medskip\noindent
We are interested in Weil-Petersson
Teichmüller space $\mathcal{T}_0$
This is a very interesting object
that intersects several branches of 
mathematics:
\begin{itemize}
\item Renormalized area of minimal surfaces
in $\mathbb{H}^3$.
\item Stochastic Loewner evolution.
\item Geometric computer vision.
\item Berezin Quantization and Bott-Virasoro
group.
\item Planar Coulomb gasses.
\end{itemize}

\medskip\noindent
{\bf Loewner transport.}
Let  $\Omega \subset \hat{\mathbb{C}}$
be a Jordan domain (interior of Jordan
curve).
Let $p,q \in \Omega$.
For example, $\Omega=\mathbb{H}^2$,
$p=0$ and $q=\infty$.
Let $\gamma:[0,1]\hookrightarrow
\overline{\mathbb{H}^+}$
with $\gamma(0)=0$ and $\gamma(t) \in
\mathbb{H}^+$ for all $t>0$.

For all $t$ define
$$
\Omega_t:\mathbb{H}^+\setminus
\gamma([0,t])
$$
$$
g_t:\mathbb{H}^t \to \Omega_t
$$
be the unique uniform map
$$
g_t(z)
=z + \frac{2\varphi(t)}{z}
+ o \left(\frac{1}{z}\right)
$$
where $\varphi$ is real valued and
increasing. Parametrize $\gamma$
by $\varphi$:
$$
g_t(z)=z+ \frac{2t}{z}
+o \left(\frac{1}{z}\right).
$$
Caratheory:
$$
W(t)=g^{-1}_t(\gamma(t))
$$
$$
W:\underbrace{[0,L]}_{\text{domain of }
\gamma} \to \mathbb{R}.
$$

The {\it Lowener transform} (1923)
is the map
$$
\gamma \mapsto W
$$
and $W$ is the {\it driving function},
invertible.

\begin{remark}
\label{seminars-differential-geometry-remark-brownian-loewner}
If $W$ is Brownian we obtain a 
stochastic loewner evolution.
Yilin Wane.
\end{remark}

\medskip\noindent
{\it Loewner energy:}
$$
I(\gamma)
:= \frac{1}{2}
\int_0^L W'(t)^2 dt
$$
\begin{remark}
\label{seminars-differential-geometry-remark-loewner-energy-independent}
Jordan curves have a well-defined
Loewner enery
[some remark I didn't understand concerning
$\Gamma$ Jordan curve, $p,q \in \Gamma$…]
\end{remark}

\noindent
In this talk we shall take this
as a definition:

\begin{theorem}[Wang]
\label{seminars-differential-geometry-theorem-wang}
If $\Gamma \subseteq \hat{ \mathbb{C}}$
is a Jordan curve then
$I(\Gamma) < \infty$ if and only if
$\Gamma$ is in $\mathcal{T}_0$,
the {\it Weil-Petersson Teichmüller space}.
\end{theorem}

\medskip\noindent
{\bf Universal Teichmüller space.}
Let
$B_1^{\infty}$ be the open unit ball in
$L^\infty(\mathbb{H}^+)$.
Let $\mu \in B_1^\infty$,
$$
\mu_0(z)=
\begin{cases}
\mu(z)\qquad &\text{ in }\mathbb{H}^+ \\
0\qquad &\text{ in }\mathbb{H}^-.
\end{cases}
$$

The upshot of the following theorem
is that these functions behave
like holomorphic functions:

\begin{theorem}[Measurable Riemann Mapping
Theorem]
\label{seminars-differential-geometry-theorem-measurable-rmt}
There exists a unique $f_\mu:\hat{\mathbb{C}}
\to \hat{\mathbb{C}}$ such that
\begin{enumerate}
\item $f_\mu(0)$, $f_\mu(1)=1$,
 $f_\mu(\infty)=\infty$.
\item 
$$
\frac{f_{\mu,\bar{z}}}{
f_{\mu z}}=\mu_0.
$$
\end{enumerate}
\end{theorem}

\noindent
Recall:

\begin{definition}
\label{seminars-differential-geometry-definition-quasiconformal-map}
$f:\Omega(\subseteq\mathbb{C})
\to \mathbb{C}$
is {\it quasiconformal} if
\begin{enumerate}
\item $f$ is a local homeomorphism,
\item $f$ has locally $L^2$ 
dist. derivatives.
\item There exists an $L^\infty$
function $\mu$ such that
$\|\mu\|_{L^\infty}<1$ such that
$$
f_{\bar{z}}=f_z\cdot \mu
$$
and $\mu$ is the {\it Beltrami differential} 
of $f$.
\end{enumerate}
\end{definition}

\noindent
$\mu_0$ is the ``extension by $0$''.
There's also:
$$
\mu_s(z)
=\begin{cases}
\mu(z)\qquad &z \in \mathbb{H}^+ \\
\overline{\mu(\bar{z})}\qquad &
z \in \mathbb{H}^-
\end{cases}
$$
and by the Measurable Riemann Mapping 
Theorem, there exists a unique
$f^u:\hat{\mathbb{C}} \to \hat{\mathbb{C}}$
such that $\sim$
$$
f_{\bar{z}}^\mu
=\mu_s f^\mu.
$$

[Picture which I think
says $f_\mu$ maps the upper half plane
to a Jordan region]

\begin{definition}
\label{seminars-differential-geometry-definition-quasicircle}
A Jordan curve $\Gamma$ is a
{\it quasicircle} whenever
$$
\Gamma=f_\mu(\hat{\mathbb{R}})
$$
\end{definition}

\noindent
Symmetry:
[Picture showing $f^\mu$
putting the upper half plane
to the whole plane.]

$f^\mu$ restricts to a homeomorphism
of $\hat{\mathbb{R}}$,
$$
\varphi:\hat{\mathbb{R}}
\to \hat{\mathbb{R}}
$$
is quasisymmetic whenever
$$
\varphi=f^\mu|_{\hat{\mathbb{R}}}
$$
for some $\mu$.

$$
\xymatrix{
& B_1\subseteq L^\infty
(\mathbb{H}^+)\ar[dl]\ar[dr]\\
\text{Quasicircles}
\ar[rr]_{\text{Homeo}}
&&
\text{Quasisymmetries}
}
$$
where the labelled homeomorpihsm
is the {\it conformal Welding}.
If $\mu_1, \mu_2 \in B_1$, then
$$
\text{Im}(f_{\mu_1})
= \text{Im}(f_{\mu_2})
\iff
f^{\mu_1}|_{\hat{\mathbb{R}}}
=f^{\mu_2}|_{\hat{\mathbb{R}}}
$$

Bers was very interested in this triangle,
and introduced what he called
``universal Teichmüller space''.

$$
\begin{aligned}
\mathcal{T}&=\text{QCirc}/
\text{PSL}(2,\mathbb{C})\\
&=\text{QSymm}/\text{PSL}(2,\mathbb{R})
\times \text{PSL}(2,\mathbb{R})\\
&=B^\infty_1/\sim.
\end{aligned}
$$

Bers showed that
\begin{itemize}
\item $\mathcal{T}$ is a Banach
manifold.
\item Every $TM$ space embeds smoothly
in $\mathcal{T}$.
\end{itemize}

\noindent
Every Teichmüller space has a 
Weil-Petersson metric. Does $\mathcal{T}$
carry a metric which is induced
by a Weil-Petersson metric? No.

Takhtajan-Teo ($\sim$2004):
$\mathcal{T}$ carries a Hilbert
manifold structure $\mathcal{T}_1$,
it has $\mathbb{R}$ connected components,
contains every Teichmüller space
and its metric (induced by the
Hilbert manifold structure)
induces a Weil-Petersson
metric.

\begin{definition}
\label{seminars-differential-geometry-definition-wp-teichmuller}
$\mathcal{T}_0 \subseteq T_1$
is the connected component of $I_0$.
\end{definition}

\medskip\noindent
{\bf Properties.}

\begin{itemize}
\item $\mathcal{T}_0$ has a
unique homogeneous Kähler metric.
\item Kähler potential is $I$.
\item Among the dozens
of characterizations of this space,
we have:
$\Gamma$ is $\mathcal{T}$
if and only if $\Gamma$
is chord arch and arc-length
parametrization is $H^{3/2}$,
there exists $c$ such that
$$
d_{\text{ext}}(p,q)
\leq  d_{\text{int}}(p,q)
\leq cd_{\text{ext}}(p,q).
$$
\end{itemize}

\noindent
Here's another characterization 
of this space:

\begin{theorem}[Bishop]
\label{seminars-differential-geometry-theorem-bishop-teichmmuller}
$\Phi \in \mathcal{T}_0$ if and only if
$\Gamma$ bounds a minimal surface
$\Sigma$ in $\mathbb{H}^3$
of finite topology and finite
renormalized area, which
is the integral of the determinant
of the second fundamental form
(and can be infinite).
\end{theorem}

\medskip\noindent
{\bf Reminder on anti de Sitter space.}

$$
\text{AdS}^{2,1}
=
\{(x_1,…,x_4)
: x_1^2+x_2^2-x_3^2-x_4^2=-1\}
/\{\pm  1\}.
$$
Also, define for any $M$
$$
\text{Det}(M):=
-\frac{1}{2}
\text{Tr}(J_0M J_0M^t)
$$
for
$$
J_0=
\begin{pmatrix}0&-1\\ 
1&0\end{pmatrix}.
$$
Then
$$
\begin{aligned}
\text{AdS}^{2,1}&=
\{M : \text{Det}M=0, M\neq 0\}
/\{\pm I_0\}
\end{aligned}
$$

\noindent
and
$$
\begin{aligned}
\delta_{\infty}\text{AdS}^{2,1}&=
\{M: \text{Det}(M)=0,M \neq 0\}\\
&=\{M:\text{rk}M=1\}/\mathbb{R}^*\\
&=\underbrace{\mathbb{R}P^1}_{\text{image}} 
\times \underbrace{\mathbb{R}P^1}
_{\text{kernel}}.
\end{aligned}
$$

\noindent
A homeomorphism 
$\varphi:\mathbb{R}P^1 \to \mathbb{R}P^1$
gives a graph $\Lambda_\varphi
\subset \partial_{\infty}\text{AdS}
^{2,1}$ which is spacelike.

Anti de Sitter space is a lot nicer
than hyperbolic space:

\begin{theorem}[Labourie-Toul-Wolf]
\label{seminars-differential-geometry-theorem-lab-tou-wol}
Every spacelike curve $\Lambda$
in $\partial_{\infty}\text{AdS}^{2,1}$
bounds a unique complete maximal disk.
\end{theorem}

\begin{theorem}[Graham-et al.]
\label{seminars-differential-geometry-theorem-graham-teich}
$\varphi$ is WP if and only if
$D_\varphi$ has finite renormalized area.
\end{theorem}

\section{On multiplicities in length spectra of semi-arithmetic surfaces}
\label{seminars-differential-geometry-section-multiplicities-length-spectra-semi-arithmetic-surfaces}

\noindent
Misha Belolipetsky.

Ter 26 mai 2026, 15:30.

SALA 236.

\medskip\noindent 
{\bf Abstract.} Mostramos
que determinadas superfícies semi-aritméticas
de Riemann têm um crescimento exponencial das
multiplicidades médias em seu espectro de
comprimento. Anteriormente, o crescimento
exponencial das multiplicidades médias era
conhecido apenas para superfícies
aritméticas, embora os resultados
experimentais que indicavam a possibilidade
de exemplos não aritméticos já estivessem
disponíveis. Explicarei alguns detalhes
interessantes de nossa prova e indicarei
algumas questões em aberto relacionadas. A
palestra é baseada em um trabalho conjunto
recente com Gregory Cosac, Cayo Dória e
Gisele Teixeira Paula.

\medskip\noindent
Zoll metrics $\leftrightarrow$
systole has large mutiplicity.

\medskip\noindent
Our case $\leftrightarrow$
``large geodesics have
very large multiplicity''.

\medskip\noindent
Plan of the talk:
\begin{enumerate}
\item Introduction to semiarithmetic
surfaces.
\item 
\end{enumerate}


\medskip\noindent
Let $S=\mathbb{H}^2/\Gamma$
be a Riemann surface, i.e.
the hyperbolic plane modulo a 
Fuschian group, which is a finite discrete
subgroup of $\text{PSL}(2,\mathbb{R})$
of finite coarea.
We introduce the number $\mathcal{N}(\ell)$,
the number of closed geodesics on $S$
of length $\leq \ell$.

\begin{theorem}[Huber, Selberg]
\label{seminars-differential-geometry-theorem-huber-selberg}
$\mathcal{N}(\ell)
\sim \frac{\text{exp}\ell}{\ell}$.
\end{theorem}

\noindent
This number is well-defined: i.e.
there are only finitely many closed
geodesics in a hyperbolic surface.

\begin{remark}
\label{seminars-differential-geometry-remark-nlog}
$\mathcal{N}(\text{log} x)
\sim \frac{x}{\text{log} x}$.
\end{remark}

\noindent
The length set is
$$
\{0<\ell_1<\ell_2<…\}
$$
Also we denote
$$
\mathcal{N}'(\ell)
=\#\{n : \ell_n <\ell\}
$$
and $g(\ell)$ is the multiplicity of
$\ell_n$.

\begin{definition}
\label{seminars-differential-geometry-definition-mean-multiplicity}
The {\it mean multiplicity} 
$\langle g(\ell)\rangle$ is a continuous
monotone function
$$
\mathcal{N}(\ell)
=\sum_{\ell_n<\ell}g(\ell_n)
\sim
\sum_{\ell_n \leq  \ell}
\langle f(\ell)\rangle
$$
\end{definition}

$$
\langle g(\ell)\rangle
\sim
\frac{\mathcal{N}(\ell)}{\mathcal{N}'
(\ell)}.
$$
\medskip\noindent
In a hyperbolic surface
$S = \mathbb{H}^2/\Gamma$,
closed geodesics correspond to
hyperbolic elements $\gamma \in \Gamma$,
$|\text{tr}\gamma|>2$.
Similarly, lengths correspond to traces.
We have
$$
\ell(\gamma)
=2 \text{arcosh}
\left(\frac{|\text{tr}\gamma|}{2}\right).
$$
$$
\mathcal{L}(\Gamma)
=\{|\text{tr}(\gamma)|:\gamma \in \Gamma,
\gamma\text{ hyp.}\}
$$
\begin{example}
\label{seminars-differential-geometry-example-modular-group}
$\Gamma = \text{PSL}_2(\mathbb{Z})$,
$\mathcal{L}(\Gamma) \subset \mathbb{N}$
(in particular, discrete in $\mathbb{R}$).
\end{example}

$$
\#(\mathcal{L}(\Gamma)
\cap[2,T])<T.
$$
Each length $\ell$
produces a unique trace
$\leq  2 \text{cosh}\left(\frac{\ell}{2}
\right)$.
$$
\mathcal{N}'(\ell)
\leq 2 \text{cosh}(\ell/2)\leq 2e^{\ell/2},
$$
$$
\langle g(\ell)\rangle
\sim\frac{\mathcal{N}(\ell)}{
\mathcal{N}'(\ell)}
\geq \frac{\text{exp}(\ell/2)}{2\ell}
$$

\begin{theorem}[Luo-Sernak (1994)]
\label{seminars-differential-geometry-theorem-luo-sarnak}
Instead of discretness
of $\mathcal{L}(\Gamma)$ in $\mathbb{R}$,
it is enough
$$
\# (\mathcal{L}(\Gamma)
\cap [N-1,N]
\leq \text{const.}
$$
for any $N \in \mathbb{N}$.
\end{theorem}

\noindent
This is called the
{\it bounded clustering property (BC)}.

\begin{example}
\label{seminars-differential-geometry-example-triangle24oo}
$\Gamma=\Delta(2,4,\infty)$.
It is well-known this group is generated
by
$$
\begin{pmatrix}0&1\\ -1&0\end{pmatrix},
\qquad 
\begin{pmatrix}1&\sqrt{2}\\ 0&1\end{pmatrix}.
$$
$\mathcal{L}(\Gamma) \supset 2 \mathbb{N}
\cup \sqrt{2}\mathbb{N}$.
Non discrete in $\mathbb{R}$,
but satisfies B. C.
\end{example}

\begin{remark}
\label{seminars-differential-geometry-remark-bc-holds}
BC holds for ``some''
arithmetic groups.
\end{remark}

\begin{itemize}
\item BC $\implies $ exponential group
of $\langle g(\ell)\rangle$.
\item All arithmetic Fuchsian groups have 
BC.
\item Conjecture: BC iff group is $\Gamma$
is arithmetic.
\item Non-compact case: Gerinska-Leuzinger
(2008).
\item Open in general.
\end{itemize}

\medskip\noindent
Let's try to define arithmetic 
and semi-arithmatic groups.

Let $\Gamma \subset \text{SL}_2(\mathbb{R})$
be a Fuschian group.
Consider the {\it invariant trace field}.
It's the field generated by the traces
of the elements in the group:
$$
K:=\mathbb{Q}\bigl(\{\text{tr}(\gamma^2):
\gamma \in \Gamma\}\bigr).
$$

Recall that a totally real field
is such that all complex embeddings are real.
More precisely:

\begin{definition}
\label{seminars-differential-geometry-definition-totally-real-field}
A number field $K$ is {\it totally real} if 
every field embedding
$K \hookrightarrow \mathbb{C}$ has image 
contained in $\mathbb{R}$.
\end{definition}

\begin{definition}
\label{seminars-differential-geometry-definition-semiarithmetic}
Let $\Gamma$ be a Fuchsian group
of finite covolume. $\Gamma$ is called
{\it semiarithmetic} is $K$ is totally
real and
$\text{tr}(\gamma) \in \mathcal{O}_K$
--- the integers of $K$.
\end{definition}

$$
\Gamma^{(2)}
\subset \text{SL}_2(\mathcal{O}_K)
\underbrace{\subset}_{\text{lattice}}
\text{SL}_2(K \otimes_\mathbb{Q}\mathbb{R})
=\text{SL}_2(\mathbb{R})^2
\times 
\underbrace{\text{SO}(3)^{[K:\mathbb{Q}]-r}}_{
\text{compact}}
$$
That is $\text{SL}_2(\mathbb{R})^r$
is a noncompact semisimple piece,
and $\text{SO}(3)^{[K:\mathbb{Q}]-r}$
is compact. Then we take projection,
which ammounts to ``modding out the
noncompact part''.

  $$
  \xymatrix{
  \mathrm{SL}_2(K \otimes_{\mathbb{Q}}
  \mathbb{R}) \ar[r]^-{\pi_{\infty_1}}
  \ar[dr]_{\pi_{\infty_2}} &
  \mathrm{SL}_2(\mathbb{R})^r
\times\text{SO}(3)^{[K:\mathbb{Q}]-r} \\
  & \mathrm{SL}_2(\mathbb{R})
  }
  $$

Now for $r = 1$ corresponds with 
$\Gamma$ being arithmetic.
For  $r>1$, is semi-arithmetic,
$r$ is the arithmetic dimension.
More precisely,

\begin{definition}
\label{seminars-differential-geometry-definition-arithmetic-group}
Let $\Gamma$ be a Fuchsian group of finite
covolume, and let $K$ be its invariant trace
field. We say that $\Gamma$ is
{\it arithmetic} if $\Gamma^{(2)}$ is
commensurable with $\mathrm{SL}_2(\mathcal{O}_K)$.
\end{definition}

What is happening here is that one starts
with a group of matrices over the number
field $K$ and then views it inside the much
larger real Lie group
$\mathrm{SL}_2(K \otimes_\mathbb{Q}\mathbb{R})$.
After extending scalars, this ambient group
splits into one or more noncompact
$\mathrm{SL}_2(\mathbb{R})$ factors and a
product of compact factors. The projection
diagram is the way of choosing the real
factor that gives the actual action on the
hyperbolic plane, while the other factors are
auxiliary arithmetic data. In the case
$r=1$, there is only one noncompact factor,
and that is exactly the arithmetic case;
when $r>1$, the same construction gives the
semi-arithmetic situation.

\medskip\noindent
Now,
$$
\xymatrix{
\{|\text{tr}\gamma|:\gamma \in \Gamma^{(2)}\}
\ar[r]^{\sigma_i}\ar[dr]&  \text{unbounded}\\
&C[0,2]
}
$$
satisfies BC by the pigeon hole.

\begin{example}[Golden triangle]
\label{seminars-differential-geometry-example-golden-triangle}
$$
\Gamma = \Delta(2,5,\infty)
=\langle \begin{pmatrix}0&-1\\
-1&0\end{pmatrix},
\begin{pmatrix}1&\gamma\\ 0&1\end{pmatrix}
\rangle,
$$
where $\gamma=\frac{1+\sqrt{5}}{2}$.
Then
$$
\Gamma\underbrace{\subset}_{\text{thin}}
\text{SL}_2(\mathbb{Z}[\sqrt{5}])
\underbrace{\subset}_{\text{arithm.}}
\text{SL}_2(\mathbb{R})^2.
$$
Here $K=\mathbb{Q}(\sqrt{5})$,
$\sigma_1=\text{id}:K \hookrightarrow 
\mathbb{R}$,
$\sigma_2:\sqrt{5} \to -\sqrt{5}$.
$\sigma_1(\Gamma) = \Gamma
\subset\text{SL}_2(\mathbb{R})$ - Fuschian
subgroup.
$\sigma_2(\Gamma)\subset
\text{SL}_2(\mathbb{R})$ - not discrete.
\end{example}

\noindent
And that's it! Since this was particularly
difficult to grasp, let us
include an AI-summary:

\medskip\noindent
{\bf On multiplicities in length spectra of
semiarithmetic surfaces.}
This was a talk by Misha Belolipetsky on
joint work with Gregory Cosac, Cayo Dória
and Gisele Teixeira Paula. The main result
is that certain semiarithmetic Riemann
surfaces have exponential growth of the
mean multiplicities in their length spectra.
Previously this phenomenon was known for
arithmetic surfaces; the point of the talk
was that it also occurs in some genuinely
semiarithmetic examples.

\medskip\noindent
The geometric object is a hyperbolic surface
$$
S=\mathbb{H}^2/\Gamma,
$$
where $\Gamma$ is a Fuchsian group, i.e. a
discrete subgroup of
$\operatorname{PSL}_2(\mathbb{R})$ of finite
coarea. We write $\mathcal{N}(\ell)$ for the
number of closed geodesics on $S$ of length
at most $\ell$. The Huber--Selberg theorem
says
$$
\mathcal{N}(\ell)\sim \frac{e^\ell}{\ell}.
$$
Thus the total number of closed geodesics
grows exponentially. The question of the talk
is subtler: if we list only the distinct
lengths
$$
0<\ell_1<\ell_2<\cdots,
$$
and write
$$
\mathcal{N}'(\ell)=\#\{n:\ell_n<\ell\},
$$
then how large is the average multiplicity
of a length? If $g(\ell_n)$ denotes the
number of closed geodesics with length
$\ell_n$, then morally
$$
\langle g(\ell)\rangle
\sim
\frac{\mathcal{N}(\ell)}{\mathcal{N}'(\ell)}.
$$
So the problem becomes: there are very many
closed geodesics, but how many different
lengths do they produce?

\medskip\noindent
The bridge from geometry to arithmetic is
given by traces. Closed geodesics correspond
to hyperbolic conjugacy classes in
$\Gamma$. If $\gamma\in\Gamma$ is
hyperbolic, then $|\operatorname{tr}\gamma|>2$
and
$$
\ell(\gamma)
=
2\operatorname{arccosh}
\left(\frac{|\operatorname{tr}\gamma|}{2}\right).
$$
Thus length is controlled by trace. In
particular, many different conjugacy classes
can have the same trace, hence the same
length. For example, in the modular group
$\operatorname{PSL}_2(\mathbb{Z})$, traces
are integers. Therefore the number of
possible trace values up to size
$2\cosh(\ell/2)$ is at most on the order of
$e^{\ell/2}$. Combining this with
Huber--Selberg gives
$$
\langle g(\ell)\rangle
\geq
\frac{e^\ell/\ell}{e^{\ell/2}}
=
\frac{e^{\ell/2}}{\ell}.
$$
This is the basic mechanism: if arithmetic
forces the trace set to be sparse, then many
geodesics must share the same length.

\medskip\noindent
Luo and Sarnak observed that one does not
need the trace set to be literally discrete
in $\mathbb{R}$. It is enough to have the
bounded clustering property:
$$
\#(\mathcal{L}(\Gamma)\cap[N-1,N])
\leq C
$$
for all integers $N$. Here
$\mathcal{L}(\Gamma)$ is the set of absolute
trace values of hyperbolic elements of
$\Gamma$. This condition says that traces
may cluster, but only in a uniformly bounded
way. Bounded clustering is enough to force
exponential growth of the average
multiplicity.

\medskip\noindent
This is where the commutative algebra enters.
Given a Fuchsian group $\Gamma$, its invariant
trace field is
$$
K=
\mathbb{Q}
\bigl(\{\operatorname{tr}(\gamma^2):
\gamma\in\Gamma\}\bigr).
$$
The semiarithmetic condition says that $K$ is
totally real and that the traces are
``integers'' in $K$, meaning
$$
\operatorname{tr}(\gamma)\in\mathcal{O}_K.
$$
Here $\mathcal{O}_K$ is the ring of integers
of the number field $K$. This is exactly the
notion of integrality from commutative
algebra: an element is integral over a ring
if it satisfies a monic polynomial with
coefficients in that ring. Thus
$$
\mathcal{O}_K
=
\{x\in K:
x\text{ satisfies a monic polynomial in }
\mathbb{Z}[t]\}.
$$
Equivalently, $\mathcal{O}_K$ is the
integral closure of $\mathbb{Z}$ inside
$K$. This was the important conceptual point:
these are not ordinary integers, but they
are the elements of $K$ that behave like
integers relative to $\mathbb{Z}$.

\medskip\noindent
The reason this helps is that algebraic
integers in a fixed number field form a
lattice-like object. Since $K$ is totally
real, it has several real embeddings
$$
\sigma_i:K\hookrightarrow\mathbb{R}.
$$
So one trace has several real shadows
$\sigma_i(\operatorname{tr}\gamma)$. One of
these embeddings gives the actual Fuchsian
group acting on $\mathbb{H}^2$, while the
other embeddings provide arithmetic
constraints. In the arithmetic case there is
only one noncompact real factor. In the
semiarithmetic case there may be more than
one noncompact factor, and this number is
called the arithmetic dimension. The proof
uses the fact that the traces lie in
$\mathcal{O}_K$ and that their conjugates are
controlled. By a pigeonhole/lattice-counting
argument this gives bounded clustering of
trace values.

\medskip\noindent
A guiding example is the golden triangle
group
$$
\Gamma=\Delta(2,5,\infty).
$$
Its trace field is
$$
K=\mathbb{Q}(\sqrt{5}),
$$
with two real embeddings:
$$
\sigma_1(\sqrt{5})=\sqrt{5},
\qquad
\sigma_2(\sqrt{5})=-\sqrt{5}.
$$
The first embedding gives the actual
Fuchsian group. The second embedding gives
another real realization, which is not
discrete. This is typical of the
semiarithmetic situation: the group still has
arithmetic trace data, but it is not
arithmetic in the classical sense.

\medskip\noindent
The main story of the talk can therefore be
summarized as follows. Hyperbolic geometry
turns closed geodesics into conjugacy classes.
The length of a closed geodesic is determined
by a trace. If the traces live inside the
integer ring $\mathcal{O}_K$ of a number
field, commutative algebra gives arithmetic
control over how those traces can be
distributed. This control implies bounded
clustering, and bounded clustering forces
large average multiplicities in the length
spectrum. The new theorem says that this
mechanism works not only for arithmetic
surfaces, but also for certain
semiarithmetic surfaces.

\section{Widths, index, intersection
and isospectrality}
\label{seminars-differential-geometry-section-widths-index-intersection}

\noindent
Jared Marx-Kuo, June 16, 2026,
IMPA, Differential Geometry Seminar.

\medskip\noindent
{\bf Abstract}
In this talk, I will discuss a series of
works on Gromov's $p$-widths,
$\{\omega_p\}$, on surfaces.
For ambient dimensions larger than $2$,
$\omega_p$ morally realizes the area of
an embedded minimal hypersurface of
index $p$. Historically, this was used
to prove the existence of infinitely
many minimal hypersurfaces in closed
Riemannian manifolds.
In ambient dimension $2$, the width
$\omega_p$ realizes the length of a
finite union of possibly immersed
closed geodesics:
$$
\omega_p
=
\sum_i m_i \ell(\gamma_i).
$$
Heuristically, one expects
$$
p
=
\sum_i \operatorname{Ind}(\gamma_i)
+
\#\{\text{vertices}\},
$$
where the vertices are the intersection
points of the union of geodesics.
Jointly with Lorenzo Sarnataro and
Douglas Stryker, we prove upper bounds
for the index and for the number of
vertices, making progress towards this
heuristic. Along the way, we prove a
generic regularity statement for
immersed geodesics.
If time allows, we will also discuss
the isospectral problem for the
$p$-widths, and explain how surfaces
provide a convenient setting to
investigate it.



\begin{definition}
\label{seminars-differential-geometry-definition-minimal-hypersurface}
$Y^n$ is {\it minimal} if it's critical
for area, i.e. $\{Y_t\}$ with $Y_0=Y$,
$\frac{d}{dt}\Big|_{t=0}\operatorname{Area}
Y_t=0$, which is related to mean curvature,
$\leftrightarrow H_Y \equiv0$.
\end{definition}

\noindent
Applications:

\begin{itemize}
\item Inverse problems: how can we
relate the areas of the fibers of a
foliation with the ambient Riemannian
metric?

\item Physics: energy $\leftrightarrow$ area.
\end{itemize}

\noindent
The 1-dimensional version is geodesics.
We're interested in finding the curve
of minimal length in its homology class.

\medskip\noindent
{\bf Sweepouts.}
 $(M^{n+1},g)$ closed Riemannian manifold.
$$
\Phi:X^p \to Z_n(M)
$$
where we think of $Z_n(M)$ as the 
``space of hypersurfaces''.

\begin{example}
\label{seminars-differential-geometry-example-1-sweepout}
$\Phi:[-1,1] \to Z_1(S^2)$ given by
$\Phi(t) = \{x = t\} \cap\{x^2+y^2+z^2=1\}$.
This is a foliation of $S^2$ by parallel
circles. For every point in the space,
there is a point in the domain
the interval, whose image is a fiber passing
through the given point.
\end{example}

\noindent
The full definition generalizes this
to $p$ points: if $\forall \{q_1,…,q_p\}$
there is $x \in X$ such that $\Phi(x)
\supseteq \{q_1,…,q_p\}$.

\begin{example}
\label{seminars-differential-geometry-example-rp2-sweepout}
$\Phi:\mathbb{R}P^2 \to Z_1(S^2)$,
$\Phi[a:b:c] = \{ax+by+cz=0\}\cap S^2$.
is a 2-sweepout.
\end{example}

\noindent
Now consider
$$
P_p = \{\Phi:X^p \to Z_n(M):
\Phi \operatorname{sweepout}\}
$$
\begin{definition}
\label{seminars-differential-geometry-definition-gromov-p-width}
Let $p$ be a positive integer.
The {\it Gromov $p$-width} is
$$
\omega_p(M,g) = \text{inf}_{\Phi \in P_p}
\text{sup}_{x \in \operatorname{Dom}(\Phi)}
\operatorname{Area}(\Phi(x)).
$$
\end{definition}

\noindent
Which Gromov explained is some sort
of ``non-linear spectrum of area
functional''.
It is positive and, $\omega_p \leq W_{p+1}$.

\begin{theorem}
\label{seminars-differential-geometry-theorem-weyl-sweepouts}
$$
\lim_{p \to \infty}\lambda_p
p^{\frac{-2}{n+1}}
= \frac{\alpha_{n+1}}{\text{Vol}(M,g)}.
$$
\end{theorem}

\begin{theorem}[GGLMN]
\label{seminars-differential-geometry-theorem-gglm}
$$
\lim_{p \to \infty}\omega_p
p^{\frac{-1}{n+1}}
=B_{n+1}\text{Vol}(M,g)^{\frac{n}{n+1}.}
$$
$\omega_p$ give existence of minimal surfaces.

\end{theorem}

\noindent
The idea is to see how this number
grows as $p$ varies.
The upshot is that this connects
with minimal surfaces!

\begin{theorem}[M-Neves-Z]
\label{seminars-differential-geometry-theorem-mnz}
$(M^{n+1},g)$, $3 \leq  n+1 \leq  ?$,
$g$ ``generic'', then
$$
\omega_p(M,g)
=\operatorname{Area}(\Sigma_p),
$$
$\Sigma_p$ minimal embedded with multiplicity
1.
\end{theorem}

\begin{lemma}[Corollary, Yau's conjecture]
\label{seminars-differential-geometry-lemma-yau-conjecture-minimal-surfaces}
Every closed $(M^{n+1},g)$,
$3 \leq  n+1 \leq ?$
has infinitely minimal surfaces.
\end{lemma}

\noindent
Other results:
\begin{itemize}
\item Existence of many CMCs (D,M-Z).
\item Existence of many PMCs (G-MK).
\end{itemize}

\noindent
We look for objects realising the $p$-widths.

\noindent
Questions:
\begin{enumerate}
\item What about $1=2$?
\item Can we compute $\omega_p$ for nice
$(M,g)$?
\item When does $\{\omega_p(M,g)\}_{p=1}^\infty$
determine $(M,g)$?
Can we hear the shape of a drum?
\end{enumerate}

\begin{theorem}[C-M]
\label{seminars-differential-geometry-theorem-cm1}
$(M^2,g)$
closed smooth. $\forall  p$ there exists
$\{\gamma_p\}_{p=1}^{N_p}$
closed immersed geodesics
$\{m_p\}_{p=1}^{N_p}\subset \mathbb{Z}^+$
such that
$$
\omega_p(M^2,g)
=\sum_{i=1}^{N_p}m_{i,p}\ell(\gamma_{i,p}.
$$
\end{theorem}

\begin{theorem}[C-M]
\label{seminars-differential-geometry-theorem-cm2}
$$
\omega_p(S^2,g_{\operatorname{round}})
=2\pi \operatorname{floor}(\sqrt{p}).
$$
\end{theorem}

\noindent
Contrast:
\begin{theorem}[Ambrosio-M-N]
\label{seminars-differential-geometry-theorem-amn}
$(\mathbb{R}P^2,g_{\operatorname{round}})$
is $\omega_p$-spectrally rigid, that is,
if $(M^{n+1},g)$ is such that
$\omega_p(M^{n+1},g)=\omega_p(\mathbb{R}P^2,g_{\operatorname{round}})$
for all $p$ then
$(M^{n+1},g)$ is isometric to
$(\mathbb{R}P^2,g_{\operatorname{round}})$.
\end{theorem}

\medskip\noindent
{\bf Zoll metrics.}

 \begin{definition}
\label{seminars-differential-geometry-definition-zoll-metric}
A {\it Zoll metric} is such
that all geodesics are simple 
(they do not autointersect) and of the
same length.
\end{definition}

\noindent
There is a Zoll metric on $S^2$ which is
not the round metric, found by Zoll
solving a very particular ODE.

\begin{theorem}[G-F]
\label{seminars-differential-geometry-theorem-gf-moduli}
For a fixed positive number $\ell_0$
there is an infinite-dimensional moduli
space of Zoll metrics on $S^2$, $Z$
with
$$
T_{g_{\operatorname{round}}}Z
\cong
\{f:S^2 \to \mathbb{R}:
f(-x) = -f(x)\}.
$$
\end{theorem}

\medskip\noindent
{\bf Index and intersection on surfaces.}
Recall that for $3 \leq  n+1 \leq  ?$,
$$
\omega_p = \operatorname{Area}(\Sigma_p)
$$
and
$$
p = \operatorname{Ind}(\Sigma_p).
$$

Note that the latter condition about
the index cannot be true for $n=2$.

The idea is that points of intersection
amog $\bigcup_{i=1}^{N_p}\gamma_{i,p}$
constribute to ``generalized index'' $=p$.

Conjecture: if $m_{i,p}=1$,
$$
p=\sum_{i=1}^{N_p}\operatorname{Ind}(\gamma_{i,p}+
\alpha
$$
where $\alpha$ is the number of vertices
(i.e. intersections)
of $\bigcup_{i=1}^{N_p}\gamma_{i,p}$.

\begin{theorem}[MK-S-S]
\label{seminars-differential-geometry-theorem-mkss}
If $\omega_p =
\sum_{i=1}^{N_p}m_{i,p}(\gamma_{i,p})$
with $\gamma_{i,p}$ immersed geodesics,
then
$$
\sum_{i=1}^{N_p}\operatorname{Ind}(\gamma_{i,p})
\leq p,
$$
$$
\alpha \leq  p,
$$
where $\alpha$ is as above.
\end{theorem}

\begin{theorem}
\label{seminars-differential-geometry-theorem-geodesics-intersect}
For generic $g$ on $M^2$,
any finite union of geodesics has
intersections
whcih look like 2 lines at a point.
\end{theorem}

\section{Singular Riemannian Foliations,
Variational Problems and Principles of
Symmetric Criticalities}
\label{seminars-differential-geometry-section-singular-riemannian-foliations}

\medskip\noindent
{\bf Conference information.}

\href{%
https://sites.google.com/view/%
60thanniversary-dmat-puc-rio/Home
}{International Conference Celebrating
the 60th Anniversary of the Department
of Mathematics at PUC-Rio}

August 12--14, 2026,
RDC Auditorium, PUC-Rio,
Rio de Janeiro.

\medskip\noindent
{\bf Speaker.}

Marcos Alexandrino, USP.

\medskip\noindent
{\bf Date and time.}

August 13, 2026, 13:30--14:30.

\medskip\noindent
{\bf Title.}

Singular Riemannian Foliations,
Variational Problems and Principles of
Symmetric Criticalities.

\medskip\noindent
{\bf Abstract.}

A singular foliation on a complete
Riemannian manifold $M$ is called a
singular Riemannian foliation if its
leaves are locally equidistant. An example
is the partition of $M$ into the orbits of
an isometric action. The talk investigates
problems in the calculus of variations on
compact Riemannian manifolds equipped with
singular Riemannian foliations having
special properties. The examples under
consideration include isoparametric
foliations, singular Riemannian foliations
on Euclidean fiber bundles, and partitions
into the orbits of Lie groups acting by
isometries. This is joint work with
Leonardo F. Cavenaghi, Diego Corro, and
Marcelo K. Inagaki. The talk is intended
for a general audience interested in
differential geometry and analysis.


\medskip\noindent
{\bf Main goals.}

\begin{enumerate}
\item[(I)] Present the new class of
singular foliations adapted to
variational problems (AVP). This class
should include homogeneous foliations,
as well as nonhomogeneous examples such
as isoparametric foliations and SRF on
Euclidean fiber bundles.

\item[(II)] Define the new concept of
$\mathcal F$-symmetric criticality in
order to include interesting operators
$J:W^{1,p}(M)\to\mathbb R$, such as
$$
\begin{aligned}
J_\lambda(f)
&=
\mathcal M(\|f\|^p)
\int_M
\mathcal L(
|\nabla f|^2,f,x)\,\mathrm{vol}
\\
&\quad
-\frac{\lambda}{a}
\left(
\int_M F(f,x)\,\mathrm{vol}
\right)^{r+1}.
\end{aligned}
$$
We also aim to establish
$\mathcal F$-principles of symmetric
criticality.

\item[(III)] Apply the
$\mathcal F$-principles to find
infinitely many weak solutions
$\{u_i\}_{i=1}^{\infty}$ of the
variational problem
$$
dJ_\lambda(u_i)=0.
$$
\end{enumerate}

\bigskip
\section{Rafael Fernandes, Differential Geometry
Seminar}
\label{seminars-differential-geometry-section-fernandes-today}
\medskip\noindent
{\bf Speaker.}

Rafael Fernandes, UC Santa Cruz.

\medskip\noindent
{\bf Date.}

September 11, 2026.

\medskip\noindent
{\bf Seminar.}

Differential Geometry Seminar, IMPA.

\medskip\noindent {\bf Title.} On the growth
rate of Reeb orbits on star-shaped
hypersurfaces


\medskip\noindent
{\bf Abstract.}

In this talk, I will discuss the growth rate
of Reeb orbits with respect to their periods
on certain contact manifolds. We focus on
fiberwise star-shaped hypersurfaces in the
cotangent bundle of a closed manifold. This
class of dynamical systems includes, for
example, geodesic flows. I will explain how
the string topology of the base manifold
implies that the number of Reeb orbits with
period at most grows at least like the prime
numbers, that is, on the order of . The
topological condition is the existence of a
certain non-nilpotent homology class of the
free loop space (with respect to the
Chas–Sullivan product). This is joint work
with João Pering.

\medskip\noindent
The following theorem links topology and
dynamics:

\begin{theorem}[Bangart-Hingston]
\label{seminars-differential-geometry-theorem-bangart-hingston}
Let $(M^n,g)$, $n \geq 2$, closed.
If $\pi_1(M)$ is infinite and abelian,
then
$$
\lim_{t \to \infty} N(t)
\frac{\operatorname{log}(t)}{t}>0
\iff
N(t) \geq C \frac{t}{\operatorname{log} t}.
$$
\end{theorem}

\noindent
{\bf Idea.} Lustenik-Shcnrelmann
theory + construction of
some higher homotopy groups of free loop
spaces.

$\phi_s:S^\bullet M \to S^\bullet M$, bijective
correspondence between geodesics of length
up to $t$ and closed orbits of $\phi_s$
with period up to $t$.

\medskip\noindent
Note. $S^\bullet M \subset T^\bullet M$.

\medskip\noindent
Question: does the growth hold for more
general higher surfaces of contact type of
$T^\bullet M$?

\begin{definition}
\label{seminars-differential-geometry-definition-fiberwise-star-shaped}
$D \subset T^*M$ is {\it fiberwise
star-shaped} if $D \cap T^*_+M$ is
starshaped and $\partial D = \Sigma$ is
transverse to the radial direction.
\end{definition}

\noindent
$$
\Lambda M
=
\coprod_{\eta \in H_1(M,\mathbb{Z})}
\Lambda_\eta M,
$$
$$
H_*(\Lambda M) =
\bigoplus_{\eta \in H_1(M,\mathbb{Z})}
H_*(\Lambda_\eta M).
$$

\begin{definition}
\label{seminars-differential-geometry-definition-h-inf-non-nilpotent}
$\beta \in H_*(\Lambda_\eta M)$
is called {\it $H_1$-infinity non-nilpotent} 
if $\beta^m \neq  0$ for all $m \geq  1$
and $\eta \neq  \in H_1(M,\mathbb{Z})/Tor$,
i.e. $\eta $ is not torsion.
\end{definition}

\begin{theorem}[F-Pering]
\label{seminars-differential-geometry-theorem-f-pering}
Let $M^n$ be a closed manifold,
$n>2$. If there exists $\beta$
$H_1$-infinity non-nilpotent, then
($\eta=\eta(\beta)$)
$$
\operatorname{lim inf}_{t \to \infty}
N^\eta_\Sigma(t)
\frac{\operatorname{log}t}{t}>0
$$
for all $\Sigma$ fiberwise star-shaped.
\end{theorem}

\begin{example}
\label{seminars-differential-geometry-example-f-pering}
\begin{enumerate}
\item $S^1 \times M^n$, $n>1$.
\item $M \times P$, where $M,P$ satisfy the
theorem.
\item $H$-spaces within $\pi_1$ infinite.
\item Principal $S^1$-bundles
$S^1 \to M \to B$ within
$\pi_0(B) = 0$.
\item $M$ that admit a homologically non
trivial $S^1$-action.
\end{enumerate}
\end{example}

\medskip\noindent
{\bf Symplectic homology and string
topology.} $W^{\partial M}$ Liouville domain
with $c_1(W) =0$.

Very roughly, symplectic homology is…

\begin{theorem}[A-S]
\label{seminars-differential-geometry-theorem-as}
$$
SH_*(T^*M) \simeq
H_*(\Lambda M)
$$
is an algebra isomorphism.
\end{theorem}

\medskip\noindent
{\bf Local symplectic homology.}
Let $X$ be an isolated 1-periodic orbit
of a hamiltonian $H$ in the completion of
$W$, $\hat{W}$.



\section{Ambient geometry via min-max
widths of embedded circles}
\label{seminars-differential-geometry-section-ambient-circle-widths}

\noindent
Ivan Miranda de Almeida, IMPA.
Differential Geometry Seminar, IMPA.
September 22, 2026, 15:30, room 236.

\medskip\noindent
This is joint work with Lucas Ambrozio
and Rafael Montezuma.

\medskip\noindent
{\bf Basic pipeline of the project.}

The guiding idea is to extract global
information about the ambient metric from
an extremely simple configuration: two
points on an embedded circle.

Start with an embedded circle $\Gamma$ and
look at unordered pairs $\{p,q\}$ on it.
The ambient metric assigns to each pair the
number $d_g(p,q)$. Instead of studying one
pair, consider sweepouts by such pairs. A
sweepout is a continuous movie
$\{p_t,q_t\}$ of pairs: initially the two
points coincide, and as $t$ varies the arc
of $\Gamma$ between them grows continuously
from a point until it covers the whole
circle. For each such movie take the largest
ambient distance $d_g(p_t,q_t)$ attained,
and then minimize this maximum over all
sweepouts. This gives the number
$\mathcal{S}_g(\Gamma)$.

The min-max procedure forces this number to
be realized by a geometrically special pair,
namely a critical pair for the ambient
distance function. One can then compare
$\mathcal{S}_g(\Gamma)$ with invariants of
the whole ambient sphere, such as its
Birkhoff width.

Thus the rough pipeline is: begin with two
points; organize pairs into sweepouts; take
a min-max of their ambient distances; obtain
a critical geometric configuration; compare
the resulting width with global invariants
of the ambient manifold; and use sufficiently
strong equality or extremality statements to
seek rigidity of the ambient metric.

In particular, one possible kind of final
result in this program is a rigidity theorem:
under suitable extremal hypotheses, the
width data may force a metric on a sphere to
be geometrically special, potentially even
the round metric. This last sentence is a
motivation for the direction of the project,
not a statement of the theorem being proved
in this talk.

\medskip\noindent
{\bf First motivation: pairs of points.}

Let $\Gamma$ be an embedded circle.
Consider the space $P$ of subsets
$\{p,q\}\subset\Gamma$ with at most two
points. Equivalently, these are ordered
pairs modulo interchange of the two
entries. The space $P$ is a M\"obius band.
Its boundary consists of the singletons
$\{p\}$, that is, the pairs with $p=q$.

If all boundary points are identified to
one point, the quotient is
$
P^*\cong\mathbb{R}P^2.
$
This is the reason both the M\"obius band
and $\mathbb{R}P^2$ appear in the setup.

If $\Gamma\subset(M,g)$, define
$
D(\{p,q\})=d_g(p,q).
$
Thus the min-max problem is for the
ambient distance between pairs of points
of the embedded circle.

\medskip\noindent
{\bf Critical points without smoothness.}

The distance function need not be smooth
at the cut locus. Grove--Shiohama
introduced a substitute for the usual
notion of critical point.

\begin{definition}
Fix $p\in M$. A point $q\neq p$ is a
{\it critical point of the distance
function from $p$} if, for every nonzero
$v\in T_qM$, there is a minimizing
geodesic $\gamma$ from $q$ to $p$ such
that
$
\angle(v,\dot\gamma(0))\leq\frac{\pi}{2}.
$
\end{definition}

\noindent
The point is that this notion gives a
deformation theory for distance functions
even where the distance is not smooth:
regions containing no critical points can
be deformed through level sets, in the
spirit of Morse theory.

\begin{theorem}[Grove--Shiohama, 1977]
\label{seminars-differential-geometry-theorem-grove-shiohama-diameter}
Let $(M,g)$ be complete and connected.
If
$
\sec_g\geq 1
$
and
$
\operatorname{diam}(M,g)>\frac{\pi}{2},
$
then $M$ is homeomorphic to a sphere.
\end{theorem}

\medskip\noindent
{\bf Sweepouts by pairs of points.}

\begin{definition}
A {\it sweepout of $\Gamma$ by pairs of
points} is a continuous family
$
\{p_t,q_t\}\subset\Gamma,
\qquad
t\in[0,1],
$
with $p_0=q_0$ and $p_1=q_1$, together
with a continuous choice of arcs
$C_t\subset\Gamma$ bounded by
$\{p_t,q_t\}$ such that
$
C_0=\emptyset,
\qquad
C_1=\Gamma.
$
\end{definition}

\noindent
Equivalently, a sweepout is a
homotopically nontrivial loop in
$P^*\cong\mathbb{R}P^2$ based at the
collapsed boundary point.

\begin{definition}
The {\it min-max width} of the embedded
circle $\Gamma\subset(M,g)$ is
$
\mathcal{S}_g(\Gamma)
=
\inf_{\{p_t,q_t\}}
\max_{t\in[0,1]}
d_g(p_t,q_t),
$
where the infimum is over sweepouts.
\end{definition}

\begin{theorem}[Ambrozio--Montezuma--Santos]
\label{seminars-differential-geometry-theorem-ams-width-critical-pair}
For every smoothly embedded circle
$\Gamma$ in a complete Riemannian
manifold,
$
\mathcal{S}_g(\Gamma)>0.
$
Moreover, there is a critical pair
$\{p,q\}\subset\Gamma$ for $D$ such that
$
d_g(p,q)=\mathcal{S}_g(\Gamma).
$
\end{theorem}

\noindent
Here ``critical'' is the nonsmooth
distance-functional analogue of the
usual variational notion. Their second
main theorem gives a finer description
when $\Gamma$ bounds a totally convex
disc and there is no stable
free-boundary geodesic: nontrivial
critical pairs lie at distance at least
$\mathcal{S}_g(\Gamma)$, and a pair
realizing the width satisfies strong
restrictions on its minimizing
geodesics.

\medskip\noindent
{\bf Motivation for the present work.}

Fix the ambient Riemannian manifold
$(M,g)$. How large can the width of an
embedded circle be? Clearly,
$
\sup_{\Gamma\subset M}
\mathcal{S}_g(\Gamma)
\leq
\operatorname{diam}(M,g),
$
where the supremum is over embedded
circles.

The ellipsoid example suggests how one
might approach the diameter: take an
elongated ellipsoid, choose a large
embedded circle passing through
diametrically opposite points, and use a
foliation to organize the relevant
pairs.

\medskip\noindent
{\bf The Birkhoff invariant.}

For a Riemannian two-sphere, a Birkhoff
sweepout is a continuous one-parameter
family of closed curves that sweeps the
sphere from one point to another. One
can obtain the standard model by taking
latitudes on a sphere and composing with
a degree-one map.

The Birkhoff min-max invariant is
$
w_B(S^2,g)
=
\inf_{\{\gamma_t\}}
\max_{t\in[0,1]}
L_g(\gamma_t),
$
where the infimum is over Birkhoff
sweepouts.

The present problem is therefore to
compare two min-max constructions:
the width of an embedded circle, built
from distances between pairs of points,
and the Birkhoff invariant, built from
lengths of closed curves sweeping the
ambient sphere.

\medskip\noindent
{\bf An ambient upper bound.}

\begin{theorem}[Ambrozio--Miranda--Montezuma]
\label{seminars-differential-geometry-theorem-amm-ambient-upper-bound}
For every Riemannian sphere $(S^2,g)$
and every embedded circle
$\Gamma\subset S^2$,
$
\mathcal{S}_g(\Gamma)
\leq
\frac{1}{2}w_B(S^2,g).
$
\end{theorem}

\noindent
For the round sphere of radius $R$ and a
great circle $\Gamma$,
$
\mathcal{S}_g(\Gamma)=\pi R,
\qquad
w_B(S^2,g)=2\pi R.
$
Thus equality holds in this example.
More generally, equality holds for convex
spheres with antipodal symmetry.

\medskip\noindent
{\bf Sweepout induced by a Morse function.}

\begin{theorem}[Induced sweepout]
\label{seminars-differential-geometry-theorem-induced-sweepout-morse}
Let
$
f:S^1\longrightarrow\mathbb{R}
$
be a Morse function such that each critical
level contains a unique critical point.
Then there exists a sweepout
$\{p_t,q_t\}$ of $S^1$ such that
$
f(p_t)=f(q_t)
$
for every $t\in[0,1]$.
\end{theorem}

\medskip\noindent
{\bf AI summary of the paper.}


The paper develops a min-max invariant of
embedded circles and then asks what this
apparently one-dimensional invariant can say
about the ambient Riemannian manifold.

\medskip\noindent
{\bf Step 1. Define the invariant.}
Start with an embedded circle
$\Gamma\subset(M,g)$ and the ambient
distance function on unordered pairs of
points of $\Gamma$. Consider continuous
one-parameter families
$
t\longmapsto\{p_t,q_t\},
$
called sweepouts, in which the arc between
the two points starts as a point and grows
until it covers the whole circle. Define
$\mathcal{S}_g(\Gamma)$ by taking, for each
sweepout, the largest distance
$d_g(p_t,q_t)$ attained and then minimizing
this maximum over all sweepouts.

\medskip\noindent
{\bf Step 2. Show that the min-max number
comes from an actual special pair.}
For every sweepout, look at the moment when
$d_g(p_t,q_t)$ is largest. Then minimize this
largest value over all sweepouts. A priori,
$\mathcal{S}_g(\Gamma)$ is just a number
produced by this optimization procedure.
The important theorem says that there is an
actual pair $\{p,q\}\subset\Gamma$ with
$
d_g(p,q)=\mathcal{S}_g(\Gamma)
$
which is critical for the distance function.
Thus the min-max construction produces a
geometrically special pair of points, not
just a number. The ordinary derivative of
the distance function cannot always be used,
because distance is not smooth at the cut
locus. This is why the paper uses the
Grove--Shiohama notion of critical point for
distance functions.

\medskip\noindent
{\bf Step 3. Produce explicit sweepouts in
order to prove upper bounds.}
Since $\mathcal{S}_g(\Gamma)$ is an infimum
over all sweepouts, to prove an upper bound
for it we only need to exhibit one sweepout
for which all the distances
$d_g(p_t,q_t)$ are controlled. Thus we need
a systematic way to construct useful
sweepouts.

This is where Morse functions enter. Suppose
$
f:S^1\longrightarrow\mathbb{R}
$
is a Morse function with one critical point
on each critical level. Look at pairs of
points on the same level of $f$, so that
$
f(p_t)=f(q_t).
$
As the level changes, the two points move
around the circle. The Morse assumptions
allow these pairs to be chosen continuously
through the critical levels, in such a way
that the arc between them starts as a point
and eventually covers the whole circle.
Therefore these pairs form a sweepout.

This becomes useful in the ambient sphere
when $f$ comes from a family of level curves,
for example from a foliation. Restricting
$f$ to the embedded circle $\Gamma$ produces
a sweepout of $\Gamma$. If the geometry of
the level curves controls the distance
between the two intersection points, then it
controls $\mathcal{S}_g(\Gamma)$ from above.

\medskip\noindent
{\bf Step 4. Use the constructed sweepouts to
compare with classical geometric quantities.}
Once we know how to manufacture a sweepout
from a useful family of level curves, we can
estimate the distances between its two
points. This converts known geometric
constructions into upper bounds for
$\mathcal{S}_g(\Gamma)$.

For a plane curve, this compares the new
min-max width with the classical width of
the curve. On a Riemannian two-sphere, start
with a Birkhoff sweepout of the whole sphere
by closed curves. From it one constructs a
sweepout by pairs on $\Gamma$. Controlling
the distance between the two points by the
length of the corresponding Birkhoff curve
gives
$$
\mathcal{S}_g(\Gamma)
\leq
\frac{1}{2}w_B(S^2,g).
$$
Thus an invariant defined using pairs of
points on one embedded circle is controlled
by a min-max invariant of the whole ambient
sphere. The round sphere attains equality,
so the estimate is sharp.

\medskip\noindent
{\bf Step 5. Ask what equality tells us
about the geometry.}
An inequality becomes especially interesting
when we understand when equality can occur.
For the round sphere of radius $R$ and a
great circle,
$$
\mathcal{S}_g(\Gamma)=\pi R,
\qquad
w_B(S^2,g)=2\pi R,
$$
so equality holds in the previous estimate.
The project therefore asks whether equality,
or other extremal width behavior, forces the
ambient metric to have special geometry.
This is the rigidity direction of the work.
The philosophy is that very simple data from
distances between pairs of points may be
strong enough, in extremal situations, to
characterize much more of the ambient
metric.

\medskip\noindent
{\bf Step 6. Stop fixing one circle and ask
what all circle widths tell us about the
whole manifold.}
So far $\Gamma$ was fixed. Now consider a
complete noncompact manifold $(M,g)$ and ask
whether there is one constant $C$ such that
$$
\mathcal{S}_g(\Gamma)\leq C
$$
for every embedded circle $\Gamma\subset M$.
This is a global condition on the ambient
manifold, expressed entirely using the
widths of its embedded circles.

The theorem says that this happens precisely
under strong large-scale restrictions. In
particular, $M$ has at most two ends, and
there is a map from $M$ to a finite graph
whose fibers all have uniformly bounded
diameter. Very roughly, when viewed from far
away, $M$ can therefore be compressed onto a
one-dimensional graph without having to
collapse sets of arbitrarily large diameter.

This is the connection with Urysohn width:
the widths of all embedded circles detect
whether the ambient manifold is, in a
large-scale quantitative sense, essentially
one-dimensional.

\medskip\noindent
Putting the paper together, start with two
points on a circle. Let the pair move through
a sweepout. Apply min-max to their ambient
distance. Obtain a special critical pair.
Construct useful sweepouts using Morse
functions and level-set families. Use them to
compare the new width with classical global
invariants. Study equality to obtain rigidity
information. Finally, vary the embedded
circle and use all of these widths together
to detect the large-scale geometry of the
ambient manifold.



\section{Bifurcations of the Clifford torus as
Willmore surfaces in Berger spheres}
\label{seminars-differential-geometry-section-caio-rodrigues-willmore-berger}

\noindent
Caio B. Rodrigues.
Differential Geometry Seminar, IMPA.
October 6, 2026, 15:30, room 224.

\medskip\noindent
{\bf Abstract.}
The Clifford torus in a Berger sphere with
parameter $\tau$ is a critical point of the
Willmore functional for every $\tau>0$,
yielding a smooth path of Willmore surfaces.
By estimating the Morse index along this
path, we apply bifurcation theory to produce
new symmetric Willmore tori emerging from
the Clifford torus.

\medskip\noindent
{\bf Outline.}
\begin{enumerate}
\item The Willmore functional and Berger
spheres.
\item Hopf tori and elastica curves in the
sphere.
\item First and second variations; stability
of the Clifford torus.
\item Heuristics: when to expect a
bifurcation.
\item Bifurcation of the Clifford torus.
\item Main ideas of the proof.
\end{enumerate}

% --- End seminars-differential-geometry.tex ---

% --- Begin xxii-egd-2026.tex ---
\chapter{XXII Escola de Geometria Diferencial}
\label{chapter-xxii-egd-2026}


Diferencial}

\medskip\noindent
{\bf Conference information.}

\href{https://egd.ufpi.edu.br/home}{%
XXII Escola de Geometria Diferencial}

July 13--17, 2026,
Gran Hotel Arrey,
Teresina, Piau\'i, Brazil.


\section{Geometric topology of
3-manifolds}
\label{xxii-egd-2026-section-geometric-topology-three-manifolds}

\medskip\noindent
{\bf Advanced minicourse.}

\'Alvaro Ramos,
Universidade Federal do Rio Grande do Sul.

\medskip\noindent
{\bf Idea of the course.}
\begin{enumerate}
\item Lecture 1. Topology and
geometryization of surfaces, 3-manifolds
and Poincaré conjecture (Ch 1), and decompositions
of 3-manifolds (Ch 2).
\item Lecture 2. Seiferet fibered spaces (Ch
2).
\item Lecture 3 (most important).
The geometrization conjecture and the work
of Thurston (Ch 4).
\item Lectura 4. The work of Perelman:
proving geometrization via the Ricci flow
(Ch 6).
\end{enumerate}

\medskip\noindent
What are all (orientable, closed) possible
topologies for 3-manifolds? This is a
{\it very} difficult question. Even
with geometrization!

Poincaré was interested in this question
around 1900. He introduced the fundamental
group and homology groups in this research.
A question by Poincaré: since we can
distinguish topologies using algebra, can
this algebra be sufficient to classify
topologies? Here enteres Poincaré
homology sphere.

\medskip\noindent
{\bf Poincaré conjecture.}
If $M$ is a closed, orientable
simply connected 3-manifold, then $M$
is homeomorphic to $S^3$.

\medskip\noindent
{\bf Formal start of the course.}

\begin{definition}[Connected sum]
\label{xxii-egd-2026-definition-connected-sum}
Let $M_1$ and $M_2$ be $3$-manifolds
and $B_1 \subset M_1$, $B_2 \subset M_2$
be (topological) balls. Choose
$\varphi:S^1=\partial B_1 \to S^2 =
\partial B_2$ an orientation-reversing
diffeomorphism. Then
$$
M_1 \# M_2=….
$$
\end{definition}

\begin{proposition}[Properties of connected
sum]
\label{xxii-egd-2026-proposition-prop-connected-sum}
\begin{enumerate}
\item The orientation of the connected sum
is compatible with the orientation of both
$M_1$ and $M_2$.
\item (Commutativity.)
$M_1 \# M_2 = M_2 \# M_1$.
\item (Associativity.)
\item $S^3$ is the identity element.
\item There are no inverses.
\item (Something about the fundamental
group.)
\end{enumerate}
\end{proposition}

\begin{theorem}[Kneser-Milnor decomposition,
1929]
\label{xxii-egd-2026-theorem-km-decomp}
Let $M$ be a closed oriented 3-manifold.
Then there are closed oriented prime
$3$-manifolds such that $M$ is decomposed as
a connected sum.
\end{theorem}

\begin{theorem}[Kneser 1929]
\label{xxii-egd-2026-theorem-kneser29}
Furthermore, the nontrivial factors in this
decomposition are unique up to reoedering ad
orientation preserving diffeomorphisms.
\end{theorem}

\medskip\noindent
{\bf Consequences.} To understand the
topology of 3-manifolds it sufficies
to understand prime 3-manifolds.

But how do we identify prime 3-manifolds?

\begin{definition}[Essential sphere]
\label{xxii-egd-2026-definition-essential-sphere}
A sphere $S \subset M$ is called
{\it  essential} if it is not the boundary
of a 3-ball in $M$.
\end{definition}

\begin{example}
\label{xxii-egd-2026-example-not-prime-essential}
If $M = M_1 \# M_2$ is not prime
then there exists an essential sphere in
$M$.
\end{example}

\begin{theorem}
\label{xxii-egd-2026-theorem-essential-sphere-connected-sum}
If a 3-manifold contains a separating
essential sphere if and only if $M=M_1 \#
M_2$ decomposes as a connected sum.
\end{theorem}

\begin{definition}
\label{xxii-egd-2026-definition-irreducible-mfd}
An orientable 3-manifold is said to be
irreducible if $M$ does not admit any
essential sphere (irreducible implies
prime).
\end{definition}

\begin{example}
\label{xxii-egd-2026-example-essential-prime}
If $M = S^2 \times S^1$, then
$S = S^2 \times \{1\}$ is an essential
sphere,
 but $M$ is prime.
\end{example}

\begin{theorem}
\label{xxii-egd-2026-theorem-unique-prime-not-irreduc}
$S^2 \times S^1$ is the unique orientable
prime 3-manifold that is not irreducible.
\end{theorem}

\begin{theorem}[Decomposition by spheres]
\label{xxii-egd-2026-theorem}
Every closed, orientable 3-manifold can be
uniquely
expressed as a connected sum of irreducible
manifolds $M_i \neq S^3$ and factors
diffeomorphic to $S^2 \times S^1$
\end{theorem}

\noindent
Thus: to classify 3-manifolds we need to
focus on irreducible manifolds.

\begin{theorem}[Kneser (1929), Milnor (1962)]
\label{xxii-egd-2026-theorem-kneser29-milnor62}
Found a canonical decomposition of any
closed 3-manifold along essential spheres.
\end{theorem}

\medskip\noindent
The 70s: keep on decomposing!
Cutting along tori.

\begin{definition}[Incompressible torus]
\label{xxii-egd-2026-definition-incompressible-torus}
A torus $T$ in a 3-manifold is {\it
incompressible} if any embedded disk that
intersects your atorus along the boundary,
$D \cap T = \partial D$, satisfies that
$\partial D$ is trivial in $\pi_1(T)$.
\end{definition}

Equivalently (Loop theorem,
Papakyriakopoulos 1957) $T$ is
incompressible if and only if it is
$\pi_1$-injective.

\medskip\noindent
The following definition captures the
building blocks of 3-manifolds;
it's very important!

\begin{definition}[Atoroidal]
\label{xxii-egd-2026-definition-atoroidal}
A compact 3-manifold $P$ with either empty
or toroidal boundary is called
{\it atoroidal} if any (immersed)
incompressible torus is homotopoic to a
boundary component.
\end{definition}

\medskip\noindent
{\bf Consequence:} a closed 3 manifold is
atoroidal if $\pi_1(M)$… oops!

\begin{theorem}[JSJ decomposition v1]
\label{xxii-egd-2026-theorem-jsj-decompositionv1}
Let $M$ be an irreducible, closed and
orientable 3-manifold. Then tere exists a
(possibly empty) collection
$\mathcal{T}=\{T_1,T_2,…,T_n\}$ of
disjointly embedded, incompressible tori
such that each component $P$ of $M$ cut
along $\mathcal{T}$ satisfies wither:
\begin{itemize}
\item $P$ is atoroidal, or
\item  $P$ is Seifert fibered (foliated by
circles, defined below).
\end{itemize}
\noindent
Furthermore, any collection of such tori
with the minimal numer of elements is
unique, up to isotopy.
\end{theorem}

\medskip\noindent
{\bf Next steps.} 
Examples of compact 3-manifolds with
toroidal boundary (baby steps).
Defining and classifiying Seifert fibered
spaces (relatively easy) and unserstanding
atoroidal manifolds (hardest part).

\medskip\noindent
{\bf Examples.}
Let $M = S^3$ and let $K$ be a knot
in $S^3$ (an embedded closed curve).
Take $N$ a tubular neighbourhood around $K$.
So $P = S^3 \setminus N$ is a compact
orientable 3-manifold with toroidal boundary
(a not complement).

But we could have also considered any link
$L$ (a collection of disjoint knots),
or even let $M$ be another 3-manifold.

\begin{example}[Complement of the unknot]
\label{xxii-egd-2026-example-complement-unknot}
Let $S^1 \subset S^3$ be the unknot.
Then $S^3 \setminus S^1 =
\mathbb{R}^3\setminus\mathbb{R}=…=
S^1\times \mathbb{D}$,
the solid tours.
\end{example}

\begin{example}[Complement of Hopf link]
\label{xxii-egd-2026-example-complement-hopf-link}
$S^3\setminus L = T^2 \times I$.
\end{example}

\begin{example}[Guing examples along
their boundaries]
\label{xxii-egd-2026-example-gluing-examples}

\end{example}

\medskip\noindent
{\bf Seifert fibered spaces.}

\begin{remark}
\label{xxii-egd-2026-remark-seifert}
\begin{itemize}
\item A Seifert fibered space will be a
compact 3-manifold $M$ possibly with
boundary which admits a foliation by
circles.
\item Studied by Herbert Seifert in the
1930s, with the intention of getting closer
to the topological classification of closed
3-manifolds.
\end{itemize}
\end{remark}

\noindent
{\bf First step: fibrations of the solid
torus.} For a pair of coprime  $p,q$,
consider a $p,q$-curve in $\partial
\mathbb{D}\times S^1$ and think of it as a
curve in $\mathbb{D}\times S^1$.

\begin{definition}
\label{xxii-egd-2026-definition-seibert-fibered}
A {\it Seifert fibered space} is a compact
(possibly with boundary) 3-manifold
that admits a decomposition into disjoint
simple closed curves such that each of this
fibers has a tubular neighbourhood made of
fibers which is isomorphic to $T(p,q)$ for
some pair $(p,q)$.
\end{definition}

\begin{example}[Hopf fibration]
\label{xxii-egd-2026-example-hopf-fibration}
\end{example}

\begin{example}[Half turn flat manifold]
\label{xxii-egd-2026-example-half-turn-flat-mfd}

\end{example}

\begin{theorem}[Seifert, 1933]
\label{xxii-egd-2026-theorem-seifert33}
Let $\pi:M \to X$ be a Seifert fibration
with $M$ compact (possibly with boundary)
and orientable. Then $\pi$ (and, in
particular, $M$) can be obtained from the
construction below:
\begin{enumerate}
\item Let $S$ be a surface.
\item LEt $D_1,…,D_n$ be disjoint disks
in $S$ and $S'=S\setminus(D_1 \cup  …
\cup D_n)$.
\item Let $N$ be the orientable $S^1$ bundle
over $S'$. (If $S'$ is orienbtable,
$N=S' \times S^1$, otherwise 
$N=S' \tilde{\times}S^1$.)
$N$ is orientable, compact and each
component of $\partial N$ is a torus.
\item Fill the toroidal boundaries using
Dehn fillings.
\end{enumerate}
\end{theorem}

\medskip\noindent
{\bf Dehn fillings.} (Missing.)

\medskip\noindent
{\bf Summary so far.}
Closed 3-manifolds $\to$ Prime factors
$\to$ $S^2 \times S^1$ or Irreducible
$\to$ JSJ composition
$\to$ either Seifert fibered pieces
(we saw there is a way of classifying these
spaces, very roughly: we saw how to
construct them), or atoroidal manifolds.

Again: we need to understand the atoroidal
manifolds.

\medskip\noindent
{\bf Recall.}

\begin{theorem}[JSJ decomposition, v2]
\label{xxii-egd-2026-theorem-jsj-decompositionv2}
Let $M$ be a compact orientable irreducibly
3 manifold, then there is a collection of
tori susch that each component $N$ of $M$
cut along $\mathcal{T}=$collection of tori
is either Seifert fibred or atoroidal.
Furthermore…
\end{theorem}

\begin{example}
\label{xxii-egd-2026-example-jsj-decompositionv2}
\begin{itemize}
\item If $M$ is already Seifert fibered or
atoroidal, $\mathcal{T} = \emptyset$.
\item  (missing)
\item (The torus bundle over $S^1$.)
The action of
$\operatorname{SL}(2,\mathbb{Z})$
by multiplication descends to the torus
$T^2 = \mathbb{R}^2/\mathbb{Z}^2$.
\end{itemize}
\end{example}

\begin{definition}
\label{xxii-egd-2026-definition-torus-bundle}
The {\it torus bundle over $S^1$} 
with monodromy $A$ is
$$
T^2_Z = 
\frac{T^2 \times [0,1]}{(x,1) \sim
(A(x),0)}.
$$
We say it's
\begin{itemize}
\item periodic if…
\item irreducible if…
\item Anosov if…
\end{itemize}
\end{definition}

\noindent
{\bf Description.} (Intuitive description of
the three kinds of torus bundles over
$S^1$.)



\medskip\noindent
{\bf Thurston.} Inspired by R. Riley (who
had found a complete hyperbolic metric of
finite volume in the complement of the
figure-eight knot in 1975), Thurston was
motivated by this question to ``bring
geometry back to topology''. He started to
study knot complements,
and ``started to see hyperbolic geometry
everywhere --- except in Seifert fibered
spaces''.

\begin{theorem}[Thurston]
\label{xxii-egd-2026-theorem-thurston}
Any knot in $S^3$ is either:
\begin{itemize}
\item A torus knot (a $(p,q)$-curve)
\item A satellite knot
\item Or a hyperbolic knot (i.e. a knot
whose complement admits a hyperbolic metric
of finite volume).
\end{itemize}
\end{theorem}

\noindent
The complement of torus knot is always
Seifert fibred. The complement of a
satellite knot always has an essential
torus. Otherwise, you can find a complete
hyperbolic metric in there.
Maybe: atoroidal + (…) $\implies $
hyperbolic?

Let's bring in the hard stuff:

\begin{definition}[Model geometry]
\label{xxii-egd-2026-definition-model-geometry}
A {\it model geometry} is a pair 
$\mathbb{X} = (X,G)$ where $X$ is a simply
connected manifold and $G$ Lie group acting
on $X$ by diffeomorphisms such that
\begin{enumerate}
\item The action is transitive,
\item The stabilizers are compact
\item $G$ is maximal wrt 1 and 2
\end{enumerate}

\noindent
Two model geometries are equivalent if there
is a diffeomorphisms of the spaces
compatible with the actions.
\end{definition}

\noindent
Geometric interpretation: $G$ is the
isometry group of $X$.

\begin{definition}[Geometric structure]
\label{xxii-egd-2026-definition-geometric-structure}
A manifold $M$ is {\it modelled} on a
geometry $\mathbb{X}=(X,G)$ is $M$ is
diffeomorphic to $X/\Pi$ where $\Pi$ is a
discrete subgroup of $G$ acting on $X $
freely. We say $\mathbb{X} $ is a {\it model
geometry} for $M$.
\end{definition}

\noindent
Equivalently, (covering condition).

\begin{definition}[missing]
\label{xxii-egd-2026-definition-missing}
?
\end{definition}

\noindent
The geometrization of surfaces can be read
in terms of model structures. Geometric
structure means you are a quotient, model
geometry is the simply connected case ($S^2,
\mathbb{H}^2$ or $\mathbb{R}^2$).

A nontrivial connected sum $M_1 \# M_2$ does
not admit a geometric structure…

The following theorem says which are the
model geometries of dimension 3
{\it that matter} --- that admit compact
quotients.

\begin{theorem}[Model geometries]
\label{xxii-egd-2026-theorem-model-geometries}
There are only $8$ model geometries of
dimension 3 that admit compact quotients:
$$
\mathbb{R}^3, S^3, S^2 \times \mathbb{R},
\mathbb{H}^2 \times \mathbb{R},
\operatorname{Nil}_3,
\widetilde{\operatorname{SL}}(2,\mathbb{R}),
\operatorname{Sol}_3
$$
\end{theorem}
\begin{theorem}[Seifert fibered]
\label{xxii-egd-2026-theorem-seifert-fibered}
Let $M$ be closed 3 manifold, Seifert
fibered.. Then $M$ is geometric and 
$\mathbb{X}$ is given by…
\end{theorem}

\noindent
Now asume $M$ {\it is} geometric:

\begin{theorem}[Classification, closed case]
\label{xxii-egd-2026-theorem-classification-closed}
Let $M$ be a closed geometric 3-manifold.
Then the underlying geometry is unique.
Moreover
\begin{itemize}
\item it's Seifert fibered if and only if…
\item it's a torus bundle of Anosov type if
and only it…
\item Otherwise $M$ is atoroidal and
$\mathbb{X}= \mathbb{H}^2$. 
\end{itemize}
\end{theorem}

\begin{theorem}[Classification, compact case]
\label{xxii-egd-2026-theorem-classification-compact}
Let $M$ be compact geometric 3-mfd with
nonempty toroidal boundary. Then the model
geometry $\mathbb{X}$ satisfies
\begin{itemize}
\item $\mathbb{X}=\mathbb{H}^3$ iff $M$ is
atoroidal.
\item $\mathbb{X}=\mathbb{H}^2 \times \mathbb{R}$
or
$\widetilde{\operatorname{SL}}(2,\mathbb{R})$
iff $M$ is Seifert fibered (with $\chi
<0$). The geometry is not unique!
\end{itemize}
\end{theorem}

\medskip\noindent
{\bf Wait a minute…}
How can it be that $\operatorname{Sol}_3$
are nwither Seifert fibered or atoroidal?

\medskip\noindent
{\bf The (canonical) torus decomposition.}

\begin{definition}[Torus decomposition]
\label{xxii-egd-2026-definition-torus-decomposition}
A minimcal collection of tori
$\mathcal{T}$ decomposing $M$ into Seifert
fibered, atoroidal pieces or (up to double
covers) torus bundles of Anosov type
is called the {\it torus decomposition} of
$M$.
\end{definition} 

\noindent
Why does it exist?… 

\medskip\noindent
{\bf The last obstruction.}
Let $\mathbb{K}$ be the Klein bottle.
There is a compact, orientable 3-manifold
denoted by $\mathbb{K} \tilde{\times} I$
which is an orientable $I$-bundle over
$\mathbb{K}$. Moreover,
$\partial(\mathbb{K}\tilde{\times}I)=T$
is a single torus.
Also $\mathbb{K}\tilde{\times}I$ is Seifert
fibered. So $\mathbb{K}\tilde{x}I$
has trivial torus decomposition. However, it
is not geometric.

When this appears in the decomposition…
we need to deal with it.

\medskip\noindent
{\bf Dealing with
$\mathbb{K}\tilde{\times}I$.}

…

\begin{definition}[Geometric decomposition]
\label{xxii-egd-2026-definition-geometric-decomposition}
A minimal collection of surfaces of zero
Euler characteristic $\mathcal{T}_g=
\{T_1,…,T_n,K_1,…,K_m\}$
decomposing $M$ into Seifert fibered,
atoroiidal pieces or (up to covers)…
\end{definition}


\medskip\noindent
{\bf Geometrization conjecture.}
Every component in the geometric
decomposition of $M$ is geometric.

(Picture of proof strategy.)

At this point, to prove geometrization, one
needs: elliptization conjecture and
hyperbolization conjecture.

Actually, the elliptization conjecture
is stronger than the elliptization
conjecture.

\medskip\noindent
{\bf Thurston's main result.}

\begin{definition}
\label{xxii-egd-2026-definition-haken}
A compact 3 manifold is {\it Haken} if it is
irreducible and contains an embedded,
two0sided, incompressible surface of genus
$g \geq 1$.
\end{definition}

\noindent
The following is Thurston's main result:

\begin{theorem}[Thurston]
\label{xxii-egd-2026-theorem-thurston-hyperbolization}
Hyperbolization conjecture holds for Haken
manifolds. In particular, geometrization
holds for Haken manifolds.
\end{theorem}

\noindent
This gave him the Fields medal. Since
Thurston's geometrization programme
includes Poincaré conjecture, after this,
many people stopped trying to disprove
Poincaré conjecture --- at the time it has
like half and half who believed it was true
and false!

\medskip\noindent
{\bf Warm up for last session: the
curve-shortening flow.}
Start with a simple closed curve
$\gamma \subset \mathbb{R}^2$.
Let $\vec{\kappa}_t(s)$ be the curvature vector.
This translates to a PDE whose solution
is a deformation of $\gamma_0$.
By Grayson 1987, evntually, $\gamma_{t_0}$
becomes convex.
This looks more and more like a circle.
``In finite time, the flow generates a
singularity, converging to a {\it round
point}.''
This is the 1-dimensional analogue of the
mean curvature flow. We are ``distributing''
the curvature --- evenly distributing the
curvature.
Something similar happens with the heat
flow. Here it's distibution of temperature:
$$
u_t = \Delta u.
$$
Hot points get colder and cold points
warmer. The heat flow evens out the
distribution of termperature to a solution
of constant temperature.

Richard Hamilton proposed to prove
geometrization using analysis --- exactly
what happened with temperature in the heat
flow but with curvature.

\medskip\noindent
{\bf Hamilton.} The goal is: start with a
manifold with metric $g_0$. Find a family of
metrics $g_t$ that even out curvature
by some parabolic flow (like the heat flow).
Hopefully we converge to a locally
homogeneous metric, this would make $M$
geometric, thus proving geometrization.

The Ricci flow is:
$$
\frac{d}{dt}g_t = -2
\operatorname{Ric}_{g_t}.
$$

\begin{remark}
\label{xxii-egd-2026-remark-ricci-metric}
We can regard $\operatorname{Ric}(X,Y)$
as a function depending of the metric, i.e.
a functional on the space of metrics.
\end{remark}

\medskip\noindent
{\bf Ricci flow.}
A {\it solution to the Ricci flow} is
a solution of the Ricci flow equation.
We also call solutions {\it Ricci flows}.

The fundamental remark is that this is a PDE
on the space of Riemannian metrics of $M$.
That is, we are changing metrics, NOT
topology.

\begin{example}[Ricci flows]
\label{xxii-egd-2026-example-ricci-flows}
\begin{itemize}
\item $M= S^3$, $g_0$ the round metric of
constant curvature 1. Since the Ricci flow
preserves isometries, the solution can only
be a rescaling $g_t=r(t)g_0$.
Since $\operatorname{Ric}$ is invariant
under scalings, $\operatorname{Ric}_{g_t}
=2 g_0$.

Actually, the metric collapses to a point.
This is called the {\it round point}.

\item  $M=S^2 \times S^1$, $g_0+d\theta^2$.
In this case the metric collapses to
$0+d\theta^2$, the metric of the circle! We
lost dimension.

\item A connected sum $M=M_1 \# M_2$
connected by a long thin neck ($S^2 \times
[0,1]$. The neck has a lot of Ricci
curvature. This region shrinks faster, and
eventually could collapse to a line, or to a
point. Or one of the $M_i$ could also
collapse.
\end{itemize}
\end{example}

\medskip\noindent
{\bf Dealing with singularities.}

\begin{theorem}[Hamilton, 1982]
\label{xxii-egd-2026-theorem-hamilton-1982}
Let $M$ be any closed 3-manifold and $g_0$
a Riemannian metric. There exists a unique
solution $g_t$ of the Ricci flow with
$g(0)=g_0$, defined over a maxixmal interval
$[0,T)$ for $T \in (0,\infty]$.
Furthermore, if $T< \infty$, then
the curvature must blow up,
in the sense that
$\lim \operatorname{sup}_{t \to T}K(t) 
= \infty$ where we should not worry about
what $K$ is right now.
\end{theorem}

\begin{theorem}[Hamilton, 1982]
\label{xxii-egd-2026-theorem-hamilton82-2}
Let $M$ be closed 3-manifold with a metric
$g_0$ with positive Ricci curvature.
Then $T< \infty$.
Moreover, after volume normalization,
the metric converges to a metric
of constant positive sectional
curvature.
Consequently, $M$ is diffeomorphic to a
quotient of $S^3$ by a finite subgroup
of $\operatorname{Iso}(S^3$ acting freely
(for simply connected, this is Poincaré
conjecture).
\end{theorem}

\medskip\noindent
{\bf Classifying singularities.}
This is Hamilton's second paper.
A round point singularity can be
renormalized to obtain a quotient of $S^3$.
What are other possible singularities, what
to do with them?

\begin{definition}[Hamilton, 1995]
\label{xxii-egd-2026-definition-hamilton95}
\begin{itemize}
\item Type I. If $T< \infty$ and
$\operatorname{sup}(T-t)K(t)<\infty$.
\item Type II(a). $T<\infty$ but
$\operatorname{sup}(T-t)K(t)=\infty$.
\item Type II(b). $T<\infty$ but
$\operatorname{sup}tK(t)=\infty$.
\item Type III. $T=\infty$ and
$\operatorname{sup} t_k(t)<\infty$.
\end{itemize}
\noindent
The first two are called {\it singular},
and the last two {\it eternal}.
\end{definition}

\noindent
Hamilton's {\it non-collapsing assumption} 
(1995): let $\rho$ denote the injectivity
radius of the manifold, the infimum
of the injectivity radius function.
The assumption is that there is a constant
$c>0$ such that
$$
\sqrt{K(t)}\rho(t) \geq c.
$$
In other words, the injectivity radius
cannot … This avoids dimension collapse.

\medskip\noindent
{\bf Singular flows.}
\begin{theorem}[Hamilton 1995]
\label{xxii-egd-2026-theorem-hamilton95}
For a singular flow under the non-collapsing
assumtion, a sequence of dilatons around a
blow-up point converges to quotients
of one of the following:
\begin{itemize}
\item  For Type I, $S^3$ or $S^2 \times \mathbb{R}$
\item For Type II…
\end{itemize}
\end{theorem}

\noindent
The meaning: if we find $S^3/\Gamma$,
then $M$ is geometric.
If we find $(S^2 \times \mathbb{R})\Gamma$,
wither $M$ is geometric, or this is just
a neck region of $M$ (we will need to 
do a surgery in this case).
Finally, if we get the cigar 
$(\Sigma^2 \times \mathbb{R})/\Gamma$,
nobody knew what to do. Hamilton
conjectured this last case would never
appear.

\medskip\noindent
{\bf Hamilton's surgery process
(with lies).}
Start with a closed $3$-manifold.
Let Ricci flow until we get a cylinder
singularity. There is a region $U \subset M$
diffeomorphic to $S^2 \times [0,1]$
that generated the singularity. Chop it off.
Then we have two components with holes.
Glue back balls, called {\it standard caps}.
Finding the standard caps is hard analysis
--- finding a metric that is smooth and
survives the flow. There are several
parameters.
In the end,
this decomposes $M$ as a connected sum.

The lies are: the surgery happens
not at the original metric but
$\varepsilon$-close to time $T$.
We don't remove the whole region, just
one sphere. The resulting manifold
could still be connected, so we could
be removing a loop instead of decomposing
into two connected sums.

\medskip\noindent
{\bf Eternal flows.}
Assume a solution $g_t$ exists for all
$t \geq 0$ (Types IIb  and III).
There is no singularity, but 
does $g_t$ converge?

\begin{theorem}[Hamilton, 1999]
\label{xxii-egd-2026-theorem-hamilton99}
Let $g_t$ be Type III solution of RG in $M$.
Then $M $ splits into a thick part and a
thin part, glued along incompressible tori,
ther
\begin{itemize}
\item On the thick part the rescalings
converge to complete hyperbolic metrics of
finite volume.
\item On the thin part, the manifold
collapses and is topologically a graph
manifold.
\end{itemize}
\end{theorem}

\noindent
A {\it graph manifold} is either
\begin{enumerate}
\item $\operatorname{Sol}_3$ manifold.
\item A compact manifold with possibly empty
toroidal boundary whose all JSJ components
are Seifert fibered.
\item Or a finite connected sum of the above
two.
\end{enumerate}

\noindent
So: graph manifolds are geometric (up to
connected sums; we decompose them into
geometric pieces).

In conclusion, if
the Ricci flow on $M$ is Type III,
then $M$ is geometric.

\medskip\noindent
{\bf The status of Hamilton's progra.}
Consider a Ricci flow in a closed,
orientable 3-manifold $M$:
\begin{enumerate}
\item Type 1. Blow up $S^3/\Gamma$
or $S^2\times S^1/\Gamma$: $M$ is geometryc,
yay!
\item Type II(a). The cigar $\Sigma^2 \times
\mathbb{R}$: don't know.
\item Type ii(b): don't know
\item Type III. Geometrization holds, yay!
\end{enumerate}

\noindent
Thus: we need to make Ricci flow work after
surgery.
And prove Hamilton's non-collapsing
assumtion.
And prove that Type II flows never appears,
or change the program.

\medskip\noindent
{\bf Enter Perelman.}
Posted three papaers in arXiv between 2002
and 2003 and finished Hamilton's program,
proving Poincaré's conjecture.

\medskip\noindent
{\bf 2002 paper: entropy formula for Ricci
flow and its geometric applications.}

\begin{itemize}
\item Introduces the concepts of
$\mathcal{W}$-entropy and reduced volume,
which are ``monotonic'' along the glow.
\item With this monotonicity, he proves the
$\kappa$-non-collapsing theorem, which
implies that Hamilton's non-collapsing
assumption actually holds.
\item Then, he proves the cigar soliton
can never occus as blow-up limit (by volume
growth).
\end{itemize}

\noindent
At the end of the paper he sketches a proof
of geometrization saying only a finite
number of surgergies is necessary 
(on every finite time interval).
That is, the (need of)
surgeries do not accumulate.
This requires essentially no new ideas.
``It is likely that by passing to the limit in
this construction one would get a
canonically defined Ricci flow therough
singularities, but at the moment I don't
have a proof of that.''

\medskip\noindent
{\bf 2003a paper: Ricci flow with a surgery
on three-manifolds.}
``This is a technical paper.''
In the paper, he redefined to Hamilton's
Ricci flow: to correct an unjustified
statement by Hamilton that Perelman
``was unable to fix''.

\begin{itemize}
\item The most important result is the
Canonical Neighbourhood theorem.
Local control of singular points,
changing perspective from global
to local.
\item This allows to prove that singular
times never accumulate.
\end{itemize}

\noindent
Then he implements the Ricci flows with
surgeries. Start with $M$ closed and
consider a solution of Ricci flow with
surgeries.
\begin{enumerate}
\item (First case: finite time extinction.)
The flow becomes empty after a finite number
of surgeries, this decomposes $M$ in
geometric pieces. Yay!
\item (Second case: eternal flow.)
Otherwise, the flow is eternal. Perelman
obtains a thick-thin decomposition as
Hamilton in 1999. He proves the thick part
is hyperbolic and he explains (in general
lines) why the thin part is again a graph
manifold (the proof, however, is deferred to
another paper, since ``it has nothing to do
with Ricci flow''. This is all which is
missing to prove geometrization.
\end{enumerate}

\medskip\noindent
{\bf 2003b paper: Finite extinction
for the solutions to the Ricci flow on
certain 3-manifolds.}

\begin{theorem}[Perelman's shortcut]
\label{xxii-egd-2026-theorem-perelmans-shortcut}
It $\pi_1(M)$ is finite, then the RF with
surgeries of $(M,g_0)$ becomes extinct in
finite time.
\end{theorem}

\begin{lemma}[Corollary]
\label{xxii-egd-2026-lemma-elliptization}
Elliptization conjecture and Poincaré
conjecture are true.
\end{lemma}

\begin{proof}
Assume $\pi_1(M)$ is finite.
Since the flow becomes extinct in finite
time, there are only finite number of
surgeries (2003a). So eventually we
decompose $M$…
\end{proof}

\medskip\noindent
{\bf What about hyperbolization?}
This is called {\it Perelman's gap}.
He never proved the claim about graph
manifolds---it had nothing to do with
Ricci flow!
Did he know how to do it but just
chose not to?

This was completed by joint effor of several
goups (Kleiner-Lott, Morgan-Tian, etc…)






\section{AI Meets Geometry Workshop}
\label{xxii-egd-2026-section-ai-meets-geometry}

\medskip\noindent
Edward Hirst, Henrique S\'a Earp and
Tom\'as S. R. Silva,
Universidade Estadual de Campinas.

\medskip\noindent
{\bf Plan.}
\begin{enumerate}
\item Lecture 1. AI overview: understand
what is AI (theoretical framework).
\item All the rest: workshop.
\end{enumerate}

\noindent
We assume zero knowledge in computer
science or AI at all.

\medskip\noindent
{\bf Introduction.} 
Questions about geometry, curvature,
are conditions on PDEs. Other questions,
like
a connection on a vector bundle, or tensors
which carry other geometric information.
This also ammounts to posing conditions on PDEs.

Consider the difference between proving
that any Kähler class contains a Ricci-flat
metric, contrasted with: what is the value
of the metric at a given point?
Sometimes knowing that a solution exists may
not be enough.

This course is a motivation to apply these
techniques to your own problems, improve
them, adapt them. Get in touch with the
speakers after the workshop is over.

The techniques we are going to show are
particularly amenable for our problems:
finding a special embedding that has
prescribed curvature, etc --- the 
Brazilian tradition.
The idea is to scale it up.

\medskip
\medskip\noindent
{\bf Lecure 1: AI Overview.
(Dr. Edward Hirst.)}

\medskip\noindent
{\bf Plan.}
\begin{enumerate}
\item What is AI.
\item AI subfields.
\item Summary.
\end{enumerate}

\medskip\noindent
{\bf What is AI.}
AI $\implies $ Methods of computational
statistics. (Computer science + statistics.)

AI methods can:

\begin{enumerate}
\item Approximating functions.
\item Recognizing patterns.
\item Search solution spaces.
\end{enumerate}

\noindent
Developed through python libraries such as
PyTorch, TensorFlow and JAX.

AI is optimal for mathematics, since data is
{\it cheap} to ample at scale, and noiseless
(as opposed to, say, biology).

Most modern hype surrounds LLM Agents,
large AI models trained on all human
knowledge to use tools themselves
(instead of attempting at producing
knowledge,
they can now use tools themselves, which
makes them more reliable --- agents instead
of just LLM).

\medskip\noindent
{\bf Machine Learning subfields.}
\begin{enumerate}
\item Supervised learning: fit a map
$\{\operatorname{inputs}\} \to
\{\operatorname{outputs}\}$
from labelled examples. Tasks include
classification and regression.

\item Unsupervised learning: analyse
unlabelled data to find structure (typical
goals are clustering, anomaly detection and
dimensional reduction.

\item Reinforcement learning: learn a
sequence of actions that maximizes
long-term (reward in a state space, this is
useful for search control) Reinforcement
learning: learn a sequence of actions that
maximizes long-term (reward in a state
space, this is useful for sear cha control).
\end{enumerate}

\medskip\noindent
Supervised learning.
\begin{enumerate}
\item (Choose data set.) Start with labelled pairing
$(x_i, y_i)$, for example, the coefficients
of an equation and a target invariant.

\item (Choose an architecture.) Choose an
ansatz of a function to approximate the map,
$f_{\theta}$ with many trainable parameters
$\theta$.

\item Then define a ``loss function''
(we use this to measure how well is the
problem being solved by this AI method).
$\mathcal{L}(f_{\theta},y)$ measuring prediction error.

\item Use an ``optimizer'' (the algorithmic
processes on updating the parameters,
based in stochastic gradient descent).

\item Very important: check generalisation
on held-out data to asses performance, then
analyse and interpret.
\end{enumerate}

\medskip\noindent
Standard loss functions.
\begin{itemize}
\item Loss choice depends on the task,
the architecture, and desired solution.
\item The optimiser follows gradients,
schematically $\theta$ <--- $\theta-
\eta \nabla_{\theta}\mathcal{L} $.
\end{itemize}

\medskip\noindent
PINN losses.
For PDE problems, the loos can encode the
equation
itself rather than labelled outputs.
\begin{itemize}
\item Differential losses are computed with
automatic differentiation.
\item Boundary or initial conditions become
extra penalty terms.
\item Extra constraints can repel known
trivial solutions or enforce
symmetries.
\end{itemize}

\medskip\noindent
{\bf Supervised Learning: Architectures.}

\medskip\noindent
Neural networks.
\begin{itemize}
\item A neural network is a highly-parametrized
non-linear map built from affine maps and
activation functions, layering linearity and
non-linearity builds complexity and
expressibility.

\item Universal approximation theorem say that
sufficiently wide or deep networks can
approximate continuous functions on compact
sets;
this is an exstence result (not a training
guarantee).
\end{itemize}

\medskip\noindent
Transformers.
This was the next hype in the AI wave.
They are the basis of LLM. GPT is a very
large thing of this kind of model.
\begin{itemize}
\item Transformers, the basis for LLM,
combine neural subarchitectures.

$$
\operatorname{Attention}(Q,K,V)
:=\operatorname{softmax}
\left(\frac{Q K^T}{\sqrt{d_k}}\right)V.
$$
(it has three inputs, that softmax thing,
the matrix $V$ again on the RHS…).
Notably successful on sequential data such
as text, code, and tool instructions.

\item A general empirical rule of AI is that
arquitectures on tasks… [missing].
\end{itemize}

\medskip\noindent
Supervised models for geometry.
\begin{itemize}
\item Many geometric objects can be cheaply
sampled to provide large datasets:manifold
points, moduli space points, tensors,
immersion map coefficients…
\item Supervised models can approximate at
varying complexity to build intuition.
\item Trained models provide fast predictors
for scanning large parameters spaces,
ranking promising regions before exact
analysis; and for extrapolations beyond
known domains.
\end{itemize}

\medskip\noindent
Interpretability methods.
As geometers, we look for some mathematical
insight from this computational process.
\begin{itemize}
\item Interpretability is analysis of
trained models to extract mathematical
insight.
\item Common methods include: feature
importance, saliency, symbolic regression.
\item Symmetry-aware architectures can
test for, and uncover, invariances or
equivalences present in the data,solution or
task.
\item Trained models can predict verifiable
counter examples.
\item Recently PDE methods can now prove
nearby existence and error bounds on AI
approximation.
\end{itemize}

\medskip\noindent
{\bf Unsupervised learning.}

\medskip\noindent
Unsupervsed analysis.
\begin{itemize}
\item UNsupervised learning has no paired
output to train against. Their methods are
effectively computationally-intensive
statistical analysis.
\item Data is of the form $x_i$ so we could
take any input or output dataset from
supervised learning, looking to uncover
mathematical structure.
\end{itemize} 

\medskip\noindent
Unsupervised Technique Classes.
\begin{itemize}
\item Historically: two dominant classes of
techniques:
\begin{itemize}
\item Clustering:
\item Dimensional-reduction
\end{itemize}
\item Unsupervised methods are often used
throughout the pipeline on a AI-project:
analysing the data before training,
debugging training processes, analysing
trained models.
\end{itemize}

\medskip\noindent
Unsupervised Learning: Clustering.

\medskip\noindent
Standard methods:
\begin{itemize}
\item The canonical clustering method is
$K$-means: choose $K$ random centres, assign
each datum to its nearest centre, update the
centres, reassign, iterate to convergence.
\item  this can developed into Kernel,
Hierarchical, Spectral, or Density methods…
useful for irregular cluster shapes, or when
outliers matter.
\end{itemize}

\noindent
Topological data analysis (TDA).
(maybe unsupervised learning)
\begin{itemize}
\item Most common method is persistent
homology.
Treat datapoints as embedded 0-simplices,
grow balls of increasing radii
simulataneously where intersections occur,
update the growing simplicial complex,
track the homology of the complex for
persistent features.
\end{itemize}

\medskip\noindent
{\bf Unsupervised learning: dimensional
reduction}

\medskip\noindent
Standard methods
\begin{itemize}
\item The canonical dimensional reduction
method is PCA: …. linear projection onto
directions of largest variance: simple and
interceptable.
\item NNs can be used via
Autoencoders:…output data equals input, NN
learns identity map through a
lower-dimensional bottleneck, giving the
latent representation.
\item Other non-linear methods include t-SNE
and UMAP.
…mainly visualization tools, that minimise
some difference measure with a
lowe-dimensional representation of the data
distribution.
\end{itemize}

\noindent
Unsupervised methods for geometry.
\begin{itemize}
\item Identify redundant degrees of freedom:
coordinates, submanifolds, invariants.
\item Reveal effective coordinates on
high-dimensional families of solutions, and
symmetry-reduced descriptions.
\item Visualise geometric objects in
lower-dimensions.
\end{itemize}


\medskip\noindent
{\bf Reinforcement Learning: training}
(the most we'll hear about.)

\medskip\noindent
RL setup
\begin{itemize}
\item For RL problems you have a {\it state
space} $\{s_t\}$ of potential solutions, and
an {\it action space} $\{ a_t\}$ of ways to
perturb these solutions.
\item The goal is to train a {\it policy}
which optimally selects the next action
$a_t$ given the current state $s_t$.
\item Need a method of evaluating states,
which dictates the {\it reward} $r_t$
allocated per step in training.
\item The training process looks to maximise
long-term cumulative reward.
\end{itemize}

\medskip\noindent
Discrete vs continuous.
\begin{itemize}
\item The state and action spaces may be
discrete, continuous or a mix.
\begin{itemize}
\item Discrete spaces arise in
combinatorics, symbolic manipulation, and
graph search.
\item Continuous spaces arise in PDE
parameter tuning, shape optimisation and
manifold search.
\end{itemize}
\item RL useful when a good solution can be
found from iterated perturbations from
generatic starting points.
\end{itemize}

\medskip\noindent
{\bf Reinforcement learning: Methods}

\medskip\noindent
Standard methods
\begin{itemize}
\item The canonical RL method is {\it
Deep-Q}: the state space may be contiuous or
discrete, the cation space is discrete.
Value-based, such that NN predicts success
of actions given a state, the polyicy is
then a function of this (often argmax).
\item These can develop into Polycy-based
methods: … both state and action spaces may
be continuous. The policy function is
replace directly with a NN, which is
trained over trial rules.
\item These can develop further to
Model-based methods. The transition to the
next state is now probabilistically
dependent on the function, which may be
approximated with a trained NN. The reward
function may also be approximated by another
NN. Both used for simulation in training.
\end{itemize}

\noindent
RL for geometry.
\begin{itemize}
\item RL methods do not require gradients,
so useful where gradients are infeasible or
expensive --- especially for non-continuous
or singular functions.
\item Compatibility with discrete spaces
means… [missing]
\end{itemize}

\medskip\noindent
{\bf Reinforcement learning. Genetis
algorithms.}

\medskip\noindent
Genetic algorithms
\begin{itemize}
\item Predecerssors to RL, genetic algorithm
mantain a {\it population} of potential
solutions, which are {\it mutated} by small
perturbations (alike RL actions) to produce
the next generation. These are evaluated and
{\it fittest} survive.
\item {\it Crossover} can be introduced
where some interpolation function is sed to
make offspring for the next generation from
pairs in the previous.
\item These are gradient-free search
methods, so they remain useful when the
objective is noisy, discontinuous,
combinatorial, or hard to differentiate.
\end{itemize}

\medskip\noindent
Symbolic regression.
\begin{itemize}
\item This is the application of genetic
algorithm to mathematical expressions
(instead of genetics… we use some math).
\item The population list of potential
expressions, that can be mutated by a list
of actions involving a basis of functions,
e.g. $\{+,\times,\operatorname{exp}\}$,
etc…
\item Expressions are evaluated on the data
to give fitness score.
\item Can weight against complexity of
expressions; directly interpretable outputs.
\end{itemize}

\medskip\noindent
Summary.

\begin{itemize}
\item AI methods are computational and
statistical, they are useful for
approximation, pattern finding, search;
there are {\it many} methods suited to
different problems. (HERE!!! I think 
I need to ask Edward what is the method I
should use for my project!)
\item AI roughly splits into supervised,
unsupervised (first function on losses) and
reinforcement methods (Analyses data and
interprets trained models)
(perturbatively improves solutions allowing
discrete regimes)
\item Geometry is super optimal for this.
\item AI methods are informative or
uncovering relationships, building
intuition, finding counter-examples and more
recently providing rigorous proofs.
\end{itemize}


\medskip\noindent
\medskip\noindent
{\bf Lecture 2: supervised learning ---
representing and fitting geometric geometric
structures.}

\medskip\noindent
{\bf Three topics:} 
\begin{enumerate}
\item Euler's formula for polytopes
\item Ricci-flat CY metrics
\item (some holonomy UNICAMP-related topic.)
\end{enumerate}

\medskip\noindent
{\bf Neural network recap.}
A neural network is a function
$f_{\theta}:\mathbb{R}^n \to \mathbb{R}^m$.
\begin{itemize}
\item (Layers.) $v_k = \sigma(W_k v_{k-1}
+b_k)$, parameters $\theta=\{W_k,b_k\}$.
We are looking for the parameters that best
described the system.
\item  $\sigma$ a nonlinearity (ReLU, GELU).
\item $f_{\theta}=…$
\item  Train by descent…
\item Universal approximation…
\end{itemize}

\medskip\noindent
This distinction organises the whole
lecture:
\begin{itemize}
\item (A) Supervised learning: learn from
labelled examples. Given pairs $(x,y)$,
minimise $\|f_{\theta}(x)-y\|^2$.
Both of todays labs: Euler's formula
and CY metrics.

\item (B) Self-supervised… (missing)
\end{itemize}

\medskip\noindent
{\bf MNIST.} A collection of hand-written
numbers to be classified. It's like the
``hello-world''.
As we said yesterday, geometry gives us
better toy data: exact, cheap, unlimited,
noiseless.

\medskip\noindent
{\bf Euler formula.}
\begin{enumerate}
\item Sample many random polytopes, record
$(f_1,f_2,f_3)= (V,E,F)$.
\item Do not tell the network the formula.
\item Train by MSE the map
$$
f_\theta:(V,E) \mapsto F.
$$
\item Ask: what did we learn?
\end{enumerate}

\noindent
A deliberately trivial target… (missing)

\medskip\noindent
{\bf What it learns --- and why we care.}
\begin{itemize}
\item It recovers $F=E-V+1:$
weights $\approx (-1,+1)$, bias
$\approx 2$. The network discovered Euler's
formula.
\item Noiseless data $\implies $
we can check exactly that it got the maths.
\item Template established: encode as
vectors $\to$ pick a loss $\to$ train
$\to$ verify.
\end{itemize}

\noindent
{\bf The real version.}
(The original project) Reflexive polytopes
$\leftrightarrow$ CY3, Sasakian metrics..
(incomplete)

\medskip\noindent
{\bf Scaling up: does it find Euler-Poincaré
formula?} (Missing.)

\medskip\noindent
{\bf Workshop --- my own notes here.} 
The Jupyter notebook does the following:
\begin{enumerate}
\item Create de data of 205 polytopes.
represented by their
  face numbers $(V,E,F)$.
\item The learning task:  approximate
$F$ given the data of $V,E$ for lots
of polytopes.
\item Training.
\end{enumerate}

\noindent
The stochastic process here is
gradient descent. At some point
we deduce the formula.

\medskip\noindent
{\bf CY example.}
We want to compute the Ricci-flat metric
of a CY3. Recall that this metric 
is a solution of Monge-Ampère
equation: Ricci flat $\iff$
$\omega^3 = c \Omega \wedge
\overline{\Omega}$.

We'll be working with Fermat's quintic.

Since we are doing supervised learning
we need to kinds of variables
(``dogs and cats'').

Donaldson in 2005 in {\it Some numerical
results in complex differential geometry}
he approximates the metric. Very 
roghly, his algorithm is:

\begin{enumerate}
\item Pick a clever ansatz for the Kahler
potential. This potential is obtained from
sections of an ample line bundle on the
threefold.

\item He procduces a way to find sections of
the line bundle having higher and higher
degree, at infinity this sequence converges
to a potential of a Ricci-flat metric.
\end{enumerate}

\noindent
This procedure is super expensive.
For degree 12 it takes about 50 hours
of super-computing time. For degree-20
it takes 35 years --- it's not feasible.

\medskip\noindent
Out approach:
\begin{itemize}
\item Input: encode rach point
$p=(z_0,…,z_4)$.
\item Target= the determinant of the metric
(roughly)
\item Train: a very small neural network.
\item Learn from a SMALL sample.
\end{itemize}

\noindent
Here we use the mean squared function.
Then we need the evaluation method, which
is this $W$.

Own/AI notes:
\begin{enumerate}
\item Pick points in the quintic.
\item At each point compute the determinant
of the metric up to order 3 (cheap) and
at order 4 (expensive, but ok).
\item Ask the thing to estimate.
How good is the estimate? We use mean
squared error to measure how good/bad is the
estimate.
\item Improve the rule somehow.
\item Apply the estimate to some points that
were hidden during the training (but for
which we also computed the determinants).
\item Finally, ``check geometrically''
whether the prdicted metric
is close to the Ricci-flat ``using
$\sigma$''.
\end{enumerate}

\noindent
How good is the learned metric:
we must measure its Ricci-flatness. 

See Ashmore, He Ovrut: Machine Learning
Calabi-Yau metrics.

\begin{remark}
\label{xxii-egd-2026-remark-not-standard}
This is not the state-of-the-art way to compute
Ricci-flat metrics. The state-of-the-art
is without labels. It's
``physics-informed neural network''.
Find corrections to FS ambient metric based
on some set of ansatz and a tailored loss
function. Design an ad-hoc loss function
that accounts for the Monge-Ampère failure,
etc. This is what the team used
to compute the $G_2$-holonomy metrics.
\end{remark}




\medskip\noindent
{\bf $G_2$ metric computation.}
Recall what a $G_2$-manifold is… (missing).

Pipeline
\begin{enumerate}
\item cymetric $\implies $ Ricci-flat metric
on the quintic base.
\item lift it to the circle bundle
\item Generate a lot of data point and train
aa network.
\end{enumerate}

\noindent
This is supervised approach.
``First neural appoximation of a
$G_2$-structure on a CY''.

\medskip\noindent
{\bf What's next?} Flows, construct 
$G_2$-holonomy metrics.

See Heyes, Hirst Sa Earp {\it Neural and
numerical methods for $G_2$-structures on
contact Calabi-Yau 7-manifolds}.

\medskip\noindent
{\bf Summary.}
\begin{enumerate}
\item Encode the structure as
tensors/vectors.
\item Choose a loss: (A) supervised (today's
examples), (B) self-supervised.
\item Train with autodiff on the chosen
signal. (The glorified regression.)
\item Verify with an independent check
(Euler's coefficients, the
$\sigma$-measure).
\end{enumerate}



\medskip\noindent
{\bf Lecture 3: Neural PDEs.} Dr. Edward
Hirst. Today we'll do unsupervised
(=self-supervised); tomorrow
reinforcement.

The idea is: now there's no Dolandson's
algorithm to check the current data.

\medskip\noindent
{\bf Overview}.
\begin{enumerate}
\item Physics-informed neural netwoks (PINs)
\begin{itemize}
\item Losses
\item Arquitectures.
\end{itemize}
\item Minimal Willmore surfaces.
\begin{itemize}
\item The Willmore conjecture.
\item Neural Willmore flow.
\end{itemize}
\item Hands-on Python

\item Summary.
\end{enumerate}

\medskip\noindent
{\bf Neural networks review.}
It's a function with a lot of parameters,
and a function $\mathcal{L}$ which
somehow evaluates. And then there's 
some correction, improvement.
 
\medskip\noindent
{\bf PINN models}

Edward's Geometric PINN arquitectures:
there's two: one is built locally from local
patches, but it carries some error that
mathematicians don't like. The second one
solves this relying on a global embedding.

\medskip\noindent
{\bf Willmore conjecture.}
That the integral of the mean curvature
of a surface embedded in $\mathbb{R}^3$ (or
another space form?) is some multiple of
$\pi$. This was proved by Neves et al. in
2012.a

\medskip\noindent
{\bf Neural Willmore Flow.}
See Arxiv:2604.04321.
Problem setup:…

The objective is to optimize the Willmore
energy of a torus embedding in
$\mathbb{R}^3$.
The notebook shows how, as we increase
the epoch (=time of training), the embedding
approaches more and more to the optimal
embedding. This is shown is the sequence of
tori, which pass from a thin torus which is
not optimal to a fatter one which is.

``You have solved a geometric PDE in 10
minutes''. That is, this is numerical
evidence that the Willmore conjecture is
true --- which was proved in 2012 by
Neves-Marques in Annals. But the higher genus
situation is an open problem…

This uses gradients --- it's a method for
solving geometric PDEs.


\medskip\noindent
{\bf Addendum by Sá Earp.}
Transactions of the AMS --- a journal
that would accept AI-oriented submisions.
Prior submission by the tem: 
Physics letter B -- impact factor 5.
They were first ever to implement this in 
differential geometry -- in April.
Also see Daniel Platt's work on
PDEs and ML.


\medskip\noindent
{\bf Last lecture:
Reinforcement learning for geometry.}

\medskip\noindent
{\bf Plan} 
\begin{enumerate}
\item Part I. RL for mathematics.
The toy model for RL: finding geodesics
on a surface.
\item Part II. Line aragements: Freeness,
Saito Tera's conjecture.
\end{enumerate}

\medskip\noindent
{\bf General philosophy.}
\begin{enumerate}
\item (Verify.) Lean, Coq, Isabelle.
Dour-Colour, Kepler's conjecture, Liquid
Tensor. Answers: valid/invalid.
\item (Solve.) Produce a proof.
AplhaGeometry/AlphaProof IMO solver, 2024.
FirsProof project. Benchmarks of unpublished
research lemmas.
\item (Discover - today!) Suggest patterns,
examples, conjectures. Davies et al. Nature
2001. AlphaTensor, Nature 2022. No statement
to prove.
\end{enumerate}

\noindent
Today we'll explore how can a computer
explore, not merely verify.

\medskip\noindent
{\bf From labels to rewards.}

\begin{itemize}
\item (self-)supervised.
\item Unsupervised.
\item Reinforcement. No fixed data, an agent
acts. A reward signal. Search, games,
construction.
\end{itemize}

\noindent
In mathematics we usually have a clear
{\it target property} to satisfy ---
which is exactly a reward
$\implies $ many mathematical {\it search}
problems are RL problems.

The idea is just like playing chess:
what's the best next move?

\medskip\noindent
{\bf How does RL work?}
An agent interacts with an environment
\begin{itemize}
\item observe a state $s_t$,
\item pick an action $a_t
\sim\pi_\theta(a|s_t)$
\item receive a reward $r_t$,
\item move to $s_{t+1}$
\end{itemize}

\noindent
Learn $\pi_\theta$ maximizing the expected
return
$$
J(\theta)= \mathbb{E}_{\pi_\theta}
\left[ \sum_t \gamma^t r_t
\right]
$$
We need a {\it policy}: a statistical model
that predicts the possibility of obtaining a
good reward.

\medskip\noindent
{\bf Why RL fits mathematical discovery.}
Many searches share one shape:
\begin{itemize}
\item The object is built step by step:
steps of path, lines, generators, moves.
\item there is a target property (shortest
path, freenes, bound to beat)
\item can measure progress, even without
proof.
\end{itemize}

That is we need: actions + terminal
condition + reward.

\medskip\noindent
{\bf A guiding example: finding geodesics.}
Build a path, one step at a time.
On $(M,g)$ fix $p,q \in M$, we want to find
a geodesic joining $p$ and $q$. Discretize
the point choosing a sequence of points.

Do a Markov decision process:
\begin{itemize}
\item State: nodes placed so far.
\item Action: place the next node on $M$
\item Reward: minus the path's discrete
energy
$$
r=-\sum_{t} \ell(x_i, x_{i+1})^2
$$
on the {\it finished path}.
\end{itemize}

\noindent
The return is
$$
-E[\gamma]=-\frac{1}{2}\int_0^1
\|\gamma'\|^2_g dt
$$
maximizing it…

\medskip\noindent
{\bf The simplest solver: the cross-entropy
method.}
A small policy $\pi_\theta$ proposes paths
(the node coordinates). Then loop
\begin{enumerate}
\item Sample a batch of paths from
$\pi_\theta$.
\item Score each by its length $L$.
\item Keep the {\it elite} --- the shortest
path.
\item  ?
\end{enumerate}

\medskip\noindent
{\bf Rediscovery, then discovery.}

\begin{itemize}
\item (Sanity -- a known answer.)

\item (Discovery -- no closed form.)
\end{itemize}

\medskip\noindent
{\bf The template generalizes.} See papers:
AlphaTensor (2002), Learning to unkot
(2021), Ribbon knots (2025).

\medskip\noindent
{\bf Discovery vs verification.}
The machine will propose a solution to the
problem, but we need to verify it's precise.

\begin{itemize}
\item Discover --- cheap, fallible.
Optimization, RL, heuristics.
\item Verification --- exact, final.
Symbolic/exact computation and classical
proof {\it certify} or refute.
\end{itemize}


\medskip\noindent
{\bf Line arrangements.}
Start with a collection of lines,
derive a lattice.
Define the module of logarithmic
derivations of the arrangement.
Derivations that are invariant with the
arrangement. Think of derivations
as tangent vectors to the arrangement.
A derivation  $\theta=(p,q,r)$ is given
by three polynomials. It is logarithmic
if it is tangent to every line of the
arrangement. That 's why it tells us about
the singularities, because at singularities
these derivations are not well-defined.

The {\it Euler derivation} is always
in this module, it has degree 1.
This gives a decomposition of the derivation
module into the Euler part and the
syzygy part.

When this module is free, we say the
arrangement is free. More precisely,
it is a direct sum of twisted modules
with degrees $(1,d_1,d_2)$ where the
1 comes from Euler's derivation.

\begin{theorem}[Saito, 80]
\label{xxii-egd-2026-theorem-saito}
Line arrangement is free iff
$$
\det(…)
$$

\end{theorem}

\begin{theorem}[Terao 1982]
\label{xxii-egd-2026-theorem-terao}
If $\mathcal{A}$ is free with exponents
$(1,d_1,d_2)$ it characteristic
polynomial factos.
\end{theorem}

\noindent
There is a cheap test to check whether
this is satisfied, so we can discard
some candidates easily.

\noindent
{\bf Terao's conjecture.} 
Freeness is determined by the combinatorics
of the arrangement.
Open even in $\mathbb{P}^2$. Verified for
14 lines (Barakat et al. 2021).

We want to do automatic searches, the
problem explodes combinatorially.

Yes/No is not good!! It's better to have
levels of how good/bad.

{\bf Upshot.} Translate Saito's criterion
to some notion of being parallel. This
translates the yes/no to: how close to
parallel? That's the $\mathfrak{S}$ in the
paper.

Question: How do you choose the line!?!?

Bootstrap for $n \geq  14$…

\medskip\noindent
{\bf Open problems.}



\section{An Allen-Cahn approach to the
Plateau problem}
\label{xxii-egd-2026-section-allen-cahn-plateau}

\medskip\noindent
Marco Guaraco, Imperial College London.



\medskip\noindent {\bf Abstract.} Modelling
non-separating surfaces as level sets of
solutions to an elliptic equation, requires
thinking of these as sections of a line
bundle rather than as usual functions. This
perspective introduces an arena full of
interesting new problems for elliptic
functionals. In this talk I would like to
present a few of these problems, show you how
they arise naturally from geometric settings,
and discuss my recent work with Marco Badran
and Manuel Del Pino on Plateau’s problem.

\medskip\noindent
{\bf Upshot.}
Instead of solving the minimal surface
equation, we solve the Allen-Cahn
equation with some parameter $\varepsilon$,
and let $\varepsilon \to 0$.
We obtain minimal surfaces
from the level sets of the Allen-Cahn
equation.
The author is more recently working
on the ``boundary model''.


\medskip\noindent
{\bf AI Main idea.}
The Allen--Cahn approach does not
identify the minimal surface with a
solution of the PDE itself. For every
small $\varepsilon>0$, one solves
$$
\varepsilon^2\Delta u_\varepsilon
-W'(u_\varepsilon)=0
$$
and considers its transition interface
$$
\Sigma_\varepsilon,
{u_\varepsilon=0}.
$$
As $\varepsilon\to0$, these level sets
converge, in an appropriate sense, to a
minimal surface $\Sigma$ spanning the
prescribed boundary $\gamma$. Thus the
solution of the Plateau problem is
obtained as the limit of the zero level
sets of Allen--Cahn solutions.


\medskip\noindent
The Allen-Cahn is a PDE. We
shall study the level sets of PDEs.

\medskip\noindent
De Giorgi observed in 1970 an analogy
between semilinear elliptic equations 
(their level sets) and
minimal surfaces.

Recall that we say that a surface is minimal,
which means that the mean curvature $H$
is constant zero. Recall that $H$
is the average of the principal curvatures.
They are locally minimizers of the
area functional. In general they
can also be saddle points of the area
functional, like the equators on the sphere.

\medskip\noindent
{\bf Semilinear elliptic equations
(Allen-Cahn).}
Assume  $f(u)=W'(u)$ where 
$W(u)=\frac{(1-u^2)^4}{4}$.
Scaling factor: $0<\varepsilon \ll 1$.
More precisely,
$$
\varepsilon^2\Delta u - W'(u)=0.
$$
Solutions are critical point of [some
energy functionsl].

General principle: level sets $u=c$ give
minimal surfaces as $\varepsilon \to 0$.
That's why think that
level sets of semilinear
elliptic equations are in correspondence
with minimal surfaces.

\medskip\noindent
{\bf Minimisers.} Modica-Mortola (1977).
Level sets of bounded-energy minimisers
converge locally to an area-minimising
surface
as $\varepsilon \to 0$.

Allen-Cahn minimisers:
$u_\varepsilon$ minimises $E_\varepsilon$
(energy), give perimeter minimisers.

\medskip\noindent
{\bf Pacard-Ritoré (2003).} 
A way to go the other way around:
from an input surface $\Sigma \subset M$,
nondegenerate, separating and minimal,
we obtain an Allen-Cahn output, i.e.
a solution of an Allen-Cahn equation.

We think of the equation as a nonlinear
operator which gives the laplacian minus 
the $W$, etc. To solve this we linearize
the operator: take its derivative and 
apply some known method.
What's the invertibility theory
of the linearized equation?
In one dimension we have the homogeneous
Allen-Cahn ODE:
$$
q''+W'(q)=0.
$$
The only obstacle of the invertibility
theory is the parts that are tangential to
$q$.

We need to introduce another operator
to understand the invertibility theorym.
This is the Jacobi operator.
Assuming the surface is nondegerate
means that the Jacobi operator is invertible.

Then we are going to use some Newton method
to find a solution of the equation.
Approximate solution:
$$
\omega_h(x,y)
=
q(y-h(x))
$$

 \medskip\noindent
{\bf Newton method.} We want to find
a small perturbation $\varphi$ so that
$\omega+\varphi$ is a solution.
Let 
$$
S(u):= \varepsilon^2 \Delta u-W'(u).
$$
Then you can decompose the equation
that you want to solve into
a linear part and another.

\medskip\noindent
{\bf The other direction (again) 
for closed manifolds.}

\begin{theorem}[G, Finite index convergence]
\label{xxii-egd-2026-theorem-finite-index-convergence}
If $E_\varepsilon(u_\varepsilon)\leq  C$
and $\operatorname{ind}(u_\varepsilon) \leq
I$ then $\{u_\varepsilon=0\} \to \Sigma$,
a smooth embedded minimal hyupersurface for
$\dim M \leq  7$. This uses GMT regularity.
\end{theorem}

\begin{theorem}[Gaspar-G., min-max and Weyl
law]
\label{xxii-egd-2026-theorem-minmax-weyl}
[Missing]
\end{theorem}

\begin{theorem}[Chodosh-Mantouldis]
\label{xxii-egd-2026-theorem-cm}
[Missing]
\end{theorem}


\medskip\noindent
{\bf Level set regularity.}
In general, zero level sets of solutions
of the Allen-Cahn equation are not smooth
(regular). Luckily, not all solutions are
created equal!

Wang-Wei and Chodosh-Mantouldis: if you assume very mild properrties
of the solutions then we can expect them to
be regular. They do the other way aroung
of what we did above of splitting
the equation into linear and remaining part.

\medskip\noindent
{\bf Fixed boundary problems -
Plateau problem.}
 $\Sigma$ minimal, $\partial \Sigma =
\gamma$,
$\gamma \subset \mathbb{R}$.

Separating sets: a real-valued
function defined on $S^1$ must intersect
the circle an even amount of times.
This can be modelled with a cylinder,
and it's visually clear that the intersection
number must be even.
But if we put a Moebius band instead,
we can get a single intersection point!
This requires us to think of sections
of line bundles instead of functions.
This is the idea of Frölich-Struew (1990).

Let $X = \mathbb{R}^3/\gamma$. Algebraic
topology fives a unique real line bundle $L
\to X$ with monodromy $-1$ around $\gamma$.

If $W$ is even, the Allen-Cahn equation
…[is well defined?]

\begin{theorem}[FS]
\label{xxii-egd-2026-theorem-fs}
sections minimising the Allen-Cahn
energy converge, as $\varepsilon \to 0$,
to a minimal surface minimising in compact
subsets of $\mathbb{R}^3\setminus \gamma$
(uses Modica-Mortola convergence).
\end{theorem}

\begin{theorem}[G-Lynch, 2024]
\label{xxii-egd-2026-theorem-g-lynch}
[Missing]
\end{theorem}

\noindent
Then the speaker would be working
on a more recent work on the
``boundary model''.

\begin{theorem}[Bardan-Del Pino-G]
\label{xxii-egd-2026-theorem-bdpg}
There exists a unique odd solution
$\bar{U}$ of 
$$
\Delta \bar{U}
-4 |z|^2 W'(\bar{U})=0,\qquad 
\{\bar{U}=0\}=
\{y=0\},\qquad 
\partial_y \bar{U}>0.
$$
\end{theorem}

\begin{theorem}[Bardan-Del Pino-G]
\label{xxii-egd-2026-theorem-bdpg2}
Let $f$ be a bounded Hölder
section whose tangential component
$f^\top(x):= \int_{\mathbb{R}}f(x,z)
q'(z)dz$ decays in an admissible way.
Then there is a unique bounded solution
$\varphi$… [missing].
\end{theorem}

\section{Non-K\"ahler Hodge--Lefschetz
theory and the Bianchi identity}
\label{xxii-egd-2026-section-non-kahler-hodge-lefschetz}

\medskip\noindent
Arpan Saha,
Universidade Estadual de Campinas.

\medskip\noindent
{\bf Abstract.}

Being K\"ahler imposes severe constraints
on the cohomology of compact complex
manifolds, such as the Hard Lefschetz
property. It is natural to ask how far
these constraints extend beyond K\"ahler
geometry. We report on ongoing joint work
with Mario Garc\'ia Fern\'andez and Ra\'ul
Gonz\'alez Molina which abstracts the
notion of a variation of Hodge--Lefschetz
structure. Evidence suggests that such a
structure exists on distinguished
subspaces of moduli spaces of
Bismut--Ricci-flat metrics which are
pluriclosed up to source terms. These
subspaces may be regarded as replacements
for the K\"ahler cone, with affine
structure modelled on a subspace of the
$(1,1)$ Aeppli cohomology of the compact
complex manifold.

\medskip\noindent
What are non-Kähler CY manifolds?
CY means that trivial canonical bundle,
or, equivalently, a holomorphic volume form.

Compact CY arise as background solutions in
string theory.
Smoothly moving on the moduli space (the
Kähelr cone) may correspond to birational
transformations og the CY3 called flops.

Furushima-Oguiso 93: Flops can destroy
the Kähler property.

Question: what is the replacement for the
Kähler cone in the non-Kähler setting?

\medskip\noindent
Plan
\begin{enumerate}
\item Abstrac Lefschetz structures.
\item Beyond Kähler manifolds.
\item Examples via the Bianchi
identity.
\end{enumerate}

\medskip\noindent
{\bf Abstract Lefschetz structures.} 
Let $(X,J,\omega)$ be a compact Kähler
m$d$-manifold. 
We are intereseted in real classes.
$$
E:= \bigoplus_{q=0}^d
H^{q,q}(X,\mathbb{R}).
$$

The Kähler package includes
three properties,
a {\it pseudo-polarized Lefschetz
structure}:
\begin{enumerate}
\item The Hard Lefschetz property
(there's an $\mathfrak{sl}_2$ representation
with highest weight $d$ and unique highest
weight vector up to scaling.
\item Poincaré duality: there is a
nondegenerate symmetric pairing
$\mathcal{P} \in E^* \otimes E ^*$.
\item Hodge-Riemannan bilinear relations
(this is if we drop the term ``pseudo''
in ``pseudo-polarized Lefschetz structure'')
$\mathcal{P}(* \cdot,\cdot)$ is positive
definite.
\end{enumerate}

\medskip\noindent
Now we consider the Kähler cone -
this leads to variations of the Lefschetz
structure:
$$
E:= \bigoplus H^{q,q}(X,\mathbb{R})\times
\mathcal{K} \to \mathcal{K}
$$
where $\mathcal{K}$ is the Kähler cone.
This comes equipped with a conection $D$
satisfying some properties.

By Yau's theorem, for every Kähler class
there's a unique Ricci-flat representative.
Such representative induces pseud-polarised
Lefschetz structures of weight $d$ on the
fibres of $E_{[\omega]}$.

\begin{definition}
\label{xxii-egd-2026-definition}
A {\it variation of pseudo-polarised
Lefschetz structure} of weight
$d$ over $\mathcal{K}$ is a bundle
$E \to \mathcal{K}$ os pseudo-polarised
Lefschetz structures of weight $d$equipped
with a connection $D$ satisfying 1, 2 and 3.
(These conditions are related to the
connection $D$ associated of the bundle
\end{definition}

In the case of threefolds, the $(1,1)$
classes may be encodedd 

In the case $X$ is Kähler CYd, we find that
\begin{itemize}
\item The Kähler cone $\mathcal{K} \simeq$
the moduli space of Kähler Ricci-flat
$\omega$.
\item The part of the VPLS generated
by $(1,1)$ classes may be encoded
in the affine structure $D^{\mathcal{K}}$…
\end{itemize}

\medskip\noindent
The C condition is…

The V condition is…

The D condition (also called hypothesis
E for Euler, since this function replaces
the Euler vector field) is…

The L (for Lefschetz) condition is…

Hypothesis $L$ encodes the hard
Lefschetz property.

\begin{proposition}[Fernandez-Molina-Saha,
in progress]
\label{xxii-egd-2026-proposition-fmas}
Given C, V, E and L holding for some
system $\mathfrak{S}$ of PDEs on $X$,
with moduli space of solutions $\mathcal{K}$,
then
$$
E=?
$$
has the structure of a VPLS…
\end{proposition}

\medskip\noindent
{\bf Bismut-Bianchi system (examples).}
[Missing]

\section{Recent developments on
$G_2$-monopoles}
\label{xxii-egd-2026-section-g2-monopoles}

\medskip\noindent
Daniel Gomes Fadel,
Universidade de S\~ao Paulo.
Short communication.

\medskip\noindent
{\bf Abstract.}

$G_2$-monopoles solve a gauge-theoretic
PDE on a noncompact $G_2$-manifold. They
arise by dimensional reduction of the
eight-dimensional
$\operatorname{Spin}(7)$-instanton
equation and are special critical points
of an intermediate energy functional
related to the Yang--Mills--Higgs energy.
Donaldson and Segal suggested in 2009
that an enumerative invariant of
noncompact $G_2$-manifolds might be
obtained by counting $G_2$-monopoles and
should be related to conjectural
invariants counting rigid coassociative
cycles. Oliveira supported this program
in 2014 by constructing families of
$G_2$-monopoles on two Bryant--Salamon
manifolds. These families are
parametrized by the positive mass of the
monopoles. As the mass tends to infinity,
the monopoles concentrate along the
unique compact coassociative submanifold.
We discuss recent joint work with
Oliveira on limits of sequences of
$G_2$-monopoles with arbitrarily large
mass on asymptotically conical
$G_2$-manifolds, and its relation with
recent work of Parise--Pigati--Stern and
Li.

\medskip\noindent
{\bf AI Notes.} These notes are based
on a single photo and my discussion with
ChatGPT.


\medskip\noindent
\textbf{Main idea.}
The talk studies $G_2$-monopoles on noncompact
$G_2$-manifolds, especially the Bryant--Salamon
examples, and investigates the limit as the
\emph{mass} of the monopole tends to infinity.

\medskip\noindent
\textbf{Monopoles.}
A $G_2$-monopole consists of a principal
$G$-bundle $P \to X$, a connection $A$, and
a Higgs field $\Phi \in \Gamma(\operatorname{ad}P)$
satisfying
\[
*(F_A \wedge \psi)=d_A\Phi,
\qquad \psi=*\,\varphi.
\]
The Higgs field has a zero set, and for
finite-mass monopoles its norm $|\Phi|$
approaches a constant at infinity, called
the \emph{mass}.

\medskip\noindent
\textbf{Main result.}
As the mass becomes large, the monopoles
concentrate along an \emph{associative integral
cycle}. The zero set of the Higgs field detects the
support of this cycle, while away from a singular
set the gauge theory effectively abelianizes.

For the Bryant--Salamon families, the singular
behavior is especially simple, and the limiting
cycle is the compact associative zero section.


\medskip\noindent
\textbf{Geometric interpretation.}
Associative submanifolds are calibrated by the
parallel $G_2$ $3$-form. Hence they are
volume-minimizing and therefore minimal.
The talk establishes a bridge between
non-abelian gauge theory and calibrated
geometry:
\[
\text{large-mass } G_2\text{-monopoles}
\quad\Longrightarrow\quad
\text{associative calibrated cycles}.
\]

\medskip\noindent
\textbf{Take-home message.}
The asymptotic behavior of monopoles encodes
distinguished minimal submanifolds of
$G_2$-manifolds, providing a gauge-theoretic
construction of associative cycles.

\medskip\noindent
{\bf Copyied message from the photo.} 
``Summary. Large-mass $G_2$-monopoles
$\implies $ charged coassociative
integral cycles, Higgs zeros
detect the charged support,
abelianization holds on
$X\setminus (\mathcal{S} \cup  
\mathcal{O}_3)$.''

\section{First integrals of the geodesic
flow on nilpotent Lie groups}
\label{xxii-egd-2026-section-first-integrals-geodesic-flow}

\medskip\noindent
Viviana del Barco,
Universidade Estadual de Campinas.

\medskip\noindent
{\bf Abstract.}

A symmetric Killing tensor on a
Riemannian manifold is a symmetric tensor
whose covariant derivative has vanishing
complete symmetrization. Such tensors
generalize Killing fields and define first
integrals of the equations of motion which
are conserved along geodesics.

Parallel symmetric tensors, symmetric
products of Killing fields, and their
linear combinations give decomposable
symmetric Killing tensors. It is generally
difficult to determine whether a manifold
admits indecomposable symmetric Killing
tensors.

We first introduce the geometry of
Riemannian two-step nilpotent Lie groups
modelled on Euclidean space and develop
intuition for their geodesic flow. Using
the algebraic structure of these spaces,
we present recent results on the existence
and characterization of left-invariant
symmetric Killing tensors on Lie groups
with invariant metrics, focusing on
two-step nilpotent and almost abelian
groups. We give explicit examples of
indecomposable tensors and discuss their
relation to the underlying Lie algebra.

\medskip\noindent
Plan
\begin{enumerate}
\item Nilpotent Lie algebras.
Define a nilpotent Lie algebra
and a nilpotent Lie group (because
we want to discuss geodesics).
\item Riemannian Nilpotent Lie groups and
their geodesics.
\item Geodesic flow on Riemannian Lie groups.
\item First integrals of the geodesic
flow.
\end{enumerate}

\medskip\noindent
{\bf Nilpotent Lie algebras.}
What is a two-step nilpotent Lie algebra?
For us, it is a pair $(\mathbb{R}^n,
[\cdot,\cdot])$ where
$$
[\cdot,\cdot]:\mathbb{R}^n\times \mathbb{R}^n
\to \mathbb{R}^n
$$
which is bilinear, skew symmetric and
$$
[x,[y,z]]=0 \qquad \forall x,y,z \in
\mathbb{R}^n.
$$
Examples:
\begin{itemize}
\item The trivial example is when the bracket
is zero.

\item For $n=3$ put
$$
[x,y]=(0,0,x_1y_2-x_2y_1).
$$
This is called the {\it Heisenberg Lie
algebra}. For the standard basis,
$$
[e_1,e_2]=e_3,[e_1,e_2]=0=[e_2,e_3].
$$
\item For $n=6$, put
$$
[x,y]=[0,0,…]
$$

\item Etc.
\end{itemize}

\noindent
Now we introduce an inner product 
$\left<\cdot,\cdot\right>$ on
$\mathfrak{n}=(\mathbb{R}^n,[\cdot,\cdot])$.
Combining $[\cdot,\cdot]$ and
$\left<\cdot,\cdot\right>$
we define a linear transformation
$$
\begin{aligned}
\nabla: \mathbb{R}^n &\longrightarrow 
\mathfrak{so}(\mathbb{R}^n,[\cdot,\cdot]) \\
x &\longmapsto \nabla_x
\end{aligned}
$$
which satisfies Koszul's identity.

Examples:
\begin{itemize}
\item If the bracket is zero,
then the RHS of Koszul vanishes, and so
does $\nabla$.

\item In Heisenberg: matrices…
\end{itemize}

\medskip\noindent
Now we introduce a group structure on
$\mathbb{R}^n$:
$$
\begin{aligned}
\cdot: \mathbb{R}^n\times \mathbb{R}^n &\longrightarrow  
\mathbb{R}^n\\
(p,q) &\longmapsto p\cdot q:=
p+q+ \frac{1}{2}[p,q].
\end{aligned}
$$
Examples:
\begin{itemize}
\item Null bracket: usual additive
$\mathbb{R}^n$.

\item Heisenberg:
$$
(p,q)=(p_1+q_1,p_2+q_2,
p_3+p_3+\frac{1}{2}(p_1q_2-p_2q_1)
$$
\end{itemize}

\noindent
Notice that the product is differentiable
as a map from usual $\mathbb{R}^n$.
Consider the left-invariant vector field 
$X^\vee$
generated by $q \in \mathbb{R}^n$.
In fact, we can express $X^\vee$
in terms of the usual Lie bracket.
This is how we can modify euclidean geometry
to another one, e.g. using the Heisenberg
bracket. Now we introduce a Riemannian
metric that makes the left-translations
isometries.

Once we have a covariant derivative
we can introduce geodesics by the usual
equation. We can also consider the geodesic
flow. We can consider first integrals
of the geodesics flow, that is
$$
\begin{aligned}
F: \mathbb{R}^n\times\mathbb{R}^n
&\longrightarrow \mathbb{R} \\
(p,v) &\longmapsto F(p,v).
\end{aligned}
$$

\noindent
The question is: how do we construct
first integrals?

Examples:
\begin{itemize}
\item A constant function.
\item  $F(p,v)=\langle v,v\rangle_p$.
\item Vector fields satisfying the Killing
equation
$$
\langle\nabla_v X,v\rangle=0
\qquad \forall v.
$$
\end{itemize}

\medskip\noindent
{\bf Quadratic first integrals.}
We are interested in constructing functions
that are constant along geodesics.

\begin{definition}
\label{xxii-egd-2026-definition-killing-tensor}
A symmetric $2$-tensor is a {\it Killing
tensor} if $F(p,v):=\langle S_p,v\rangle_p$
is a first integral of the geodesic flow.
\end{definition}

\noindent
Problem A: characterizations of symmetric
Killing 2-tensors (left-invariant) on Lie
groups.

\medskip\noindent
Problem B: whan does a Riemannian manifold
admit indecomposable symmetric Killing
tensors?

\medskip\noindent
Construction:
\begin{itemize}
\item If $\xi_1,\xi_2$ are Killing vector
fields on $(N,\langle\cdot,\cdot\rangle)$
we can set
$$
\langle S_pv,w\rangle:=
\langle(\xi_1\cdot\xi_2,v,w\rangle
=…
$$
This is the symmetryc product of Killing
vector fields, the symmetric 2-tensor $S$ is
a Killing tensor.

\item Any linear combination of two such
fields is another one.

\item The Riemannian metric is a symmetric
tensor.
\end{itemize}

\medskip\noindent
All the construction above
can be used also for any
nilportent lie group acting simply
transitively by isometries on $(M,g)$.

Previous results:
\begin{itemize}
\item Takeuchi, Thompson (80):
spaces of constant sectional curvature, the
algebra of symmetric Killing tensors is
generated by Killing fields.
\item Heil, Moroianu, Semmelmann ('00):
metric on torus without Killing fields.
\item There are no general results:
it's interesting to study the case of
symmetric spaces (see Bolsinov, Matveev,
Miranda, Tabachnikov).
\item Viviana et al: 2-step nilpotent or almost abelian
nilpotent Lie… [missing]
\end{itemize}

\medskip\noindent
{\bf PbA.}
Let  $(N,\langle\cdot,\cdot\rangle)$
be a 2-step Lie group with a left-invaraint
metric,
simply connected and $\mathfrak{n}$ 
its Lie algebra, then
invariant symmetric 2-tensors $S$
are in correspondence with $S \in
\operatorname{End}(\mathfrak{n})$
such that $\langle Sx,y\rangle=
\langle x,Sy\rangle$.

Split the Lie algebra in two pieces,
``the image of the bracket map''
(not precise, but think about it that way
for today)
and the orthogonal complement.
That is
$$
\mathfrak{n}=\mathfrak{v}\oplus^\perp 
\mathfrak{z}
$$
A result says that $S$ is killing
iff and only if [two endomorphisms]
are both Killing separately.
This allows for the construction of several
first integrals.

Now we go to problem B.
Let $(N,\langle\cdot, \cdot \rangle)$
be a 2-step Lie group with left-invariant
metric, simply connected and
$\mathfrak{n}$ its Lie algebra. Let $S \in
\operatorname{Sym}^2(\mathfrak{n})$
with
$$
S = \begin{pmatrix}S_{\mathfrak{n}}&
S^d\\ S^d&S_{\mathfrak{z}}\end{pmatrix}.
$$
We showed that $S^d$ and $S_{\mathfrak{z}}$
are always decomposable.
Thus the main problem is $S_{\mathfrak{v}}$.
\begin{theorem}
\label{xxii-egd-2026-theorem-decomposable-sv}
$S_{\mathfrak{v}}$ is decomposable
if and only if there is an endomorphism
(we don't include it here)
$T(S_{\mathfrak{v}})$ is built from
$S_{\mathfrak{v}}$ and $[\cdot,\cdot]$
extends to a skew-symmetric derivation of
$\mathfrak{n}$.
\end{theorem}

\noindent
They have used strongly the description
of the isometry group given by
Killing vector fields. This happens
because we are on nilpotent Lie groups,
where Killing fields are well-known.

\begin{lemma}[Corollary]
\label{xxii-egd-2026-lemma-left-invariant-killing2}
Every left-invariant symmetric Killing
2-tensor left-invariant on a 2-step
nilpotent Lie group
of dimension $\leq 7$ is decomposable.
\end{lemma}

\noindent
This is not true in dimension $\geq 8$.

\medskip\noindent
{\bf Final remarks.} 
\begin{itemize}
\item Nilpotent groups are easier
because we have a precise description of the
isometry group.

\item Almost abelian Lie groups are somwhoe
simpler
thatn nilpotent Lie groups (though we
don't have a precise description of the
isometry group).

\item Ongoing work: symmetric Killing tensors
on symmetric spaces.
\end{itemize}

\section{Extracting subcomplexes on
Carnot groups}
\label{xxii-egd-2026-section-subcomplexes-carnot-groups}

\medskip\noindent
Francesca Tripaldi,
University of Leeds.

\medskip\noindent
{\bf Abstract.}
Differential complexes encode geometric
structures by means of differential
operators and cohomological invariants.
In sub-Riemannian geometry, intrinsic
subcomplexes such as the Rumin complex
often replace the de Rham complex because
they better reflect the underlying graded
structure. This talk considers how to
extract differential subcomplexes adapted
to Carnot groups. It introduces spectral
complexes constructed from spectral
sequences and discusses some of their
applications.

\medskip\noindent
{\bf Upshot.} We care about
sub-Riemannian geometry because there
are situations in which not every direction
is allowed. Like a car parallel-parking, 
because cars cannot move sideways.
Carnot groups are local models of
sub-Riemannian manifolds.
Tripaldi constructs the differential
complexes used to study these objects.

\medskip\noindent
{\bf Plan}
\begin{enumerate}
\item SubRiemannian manifolds
\begin{enumerate}
\item Riemannian vs subRiemannian
\item Carnot groups
\item Differential geometry on subRiemannian
manifolds
\end{enumerate}
\item Diff --- oops!
\end{enumerate}

\noindent
Smooth manifold $M$ of dimension $n$.
Horizontal distibution $H$ of rank
$m<n$, i.e. $H_p \subset T_pM$.
The horizontal distibution is
bracket-generating.

\begin{example}
\label{xxii-egd-2026-example-heisenberg-subbundle}
$X_1=\partial_x - \frac{y}{2}\partial_z$,
$X_2=…$, $X_3=…$ $\implies [X_1,X_2]=X_3$.
\end{example}

\noindent
{\bf Upshot on subRiemannian manifolds:}
 you cannot move along the non
horizonatal direction.

\begin{definition}[Admissible curve]
\label{xxii-egd-2026-definition-admissible-curve}

\end{definition}

\noindent
We are interested in the subRiemmanian
version of $\mathbb{R}^n$.

\begin{definition}
\label{xxii-egd-2026-definition-graded-lie-algebra}
A Lie algebra $\mathfrak{g}$ is graded
if it can be written…

It is called {\it stratified} when $\mathfrak{g}$
admits a grading in which the first layer is
bracket generating:
$$
\mathfrak{g}=V_1 \oplus … \oplus V_s
, \qquad 
[V_1,V_i]=V_{i+1},
\qquad V_s \neq 0, V_k =0, k>s.
$$
\end{definition}

\begin{definition}[Carnot group,
IA-generated]
\label{xxii-egd-2026-definition-carnot-group}
A Carnot group is a connected, simply
connected nilpotent Lie group $G$ whose
Lie algebra admits a stratification
\[
\mathfrak g=V_1\oplus\cdots\oplus V_s,
\qquad [V_1,V_i]=V_{i+1}.
\]
The horizontal distribution is obtained
by translating $V_1$ by left multiplication.
\end{definition}

\noindent
We want:
\begin{itemize}
\item Hodge laplacian,
\item Pullback [of some kind]
\item Pushforward (? of some kind?)
\end{itemize}

\noindent
The easiest example of a Carnot group
is $\mathbb{H}^1$, the Heisenberg group.
The role of stratification is that I can
come up with homogeneous dilations --- 
I put a weight on the vectors of the Lie
algebra. So: stratification $\to$ good
dilation that works. This is called
the {\it homogeneous weight} of the vectors,
and by duality, also of the covectors.

So: I have a covector, I assign a weight to
it. The idea is: two covectors with
different weight, they shall not be
dependent on each other.

Now we extend the idea of weight saying:
a zero form (function) has weight zero.

Usual notation for weighted forms:
 $$
\Omega^{h,w}(\mathbb{G})
=\{\alpha \in \Omega^\bullet(\mathbb{G}):
\operatorname{deg}(\alpha)=h,
w(\alpha)=w\}.
$$
This object is bigraded.

We can decompose the exterior differential
$d$ using the weights. Differentiating
a form gives an $h+1$
form with the same weight.

We want to test for hyper(hypo?)-ellipticity.
Principal symbol is not invertible when
testing for (hypo)ellipticity (Rockland
condition fails.

First problem: [missing].

\begin{definition}[Intrinsic (Pansu)
derivative]
\label{xxii-egd-2026-definition-pansu-derivative}
Given two Carnot groups and a map between
them…
\end{definition}

\noindent
Second problem: the exterior derivative
does not commute with the $P$-pullback.

\medskip\noindent
Third problem.
Let's try to integrate on these things.
It won't work! Not all the forms work.
A general form would give
$$
\int_{\operatorname{graph}\varphi}\alpha=0.
$$

\medskip\noindent
{\bf The Rumin complex on Heisenberg groups.}
Using the Lefschetz map (like in Kähler
geometry) we can construct the complex
of {\it intrinsic forms}.

It works good on Heisenberg! The Hodge-Laplacian
is homogeneous, satisfies the R** condition,
hypo(hyper?) elliptic, (self adjoint?).

\begin{theorem}[Rumin ;00]
\label{xxii-egd-2026-theorem-rumin00}

\end{theorem}

\noindent
Also there's a Stokes theorem, and other
results.

Let's try to extend this to arbitrary Carnot
groups.

\begin{definition}[Rumin complex]
\label{xxii-egd-2026-definition-rumin-complex}

\end{definition}

\noindent
Is it good? No.

After Heisenberg, there is the Engel group,
$\mathbb{G}=(\mathbb{R}^4,*)$.

\medskip\noindent
{\bf Using spectral sequences.}
Can we glue differentials?

\section{Stable extremal domains for the
first Dirichlet eigenvalue}
\label{xxii-egd-2026-section-stable-extremal-domains}

\medskip\noindent
Ivaldo Nunes,
Universidade Federal do Maranh\~ao.

\medskip\noindent
{\bf Abstract.}
This talk studies stable extremal domains
for the first Dirichlet eigenvalue of the
Laplacian. Such domains are classified in
the $2$-sphere and in higher-dimensional
spheres under the assumption that either
their boundary is minimal or their first
eigenvalue equals the dimension of the
sphere. The talk also presents topological
bounds for stable domains in general
compact Riemannian surfaces, assuming
either nonnegative total Gaussian
curvature or sufficiently small volume.
This is joint work with
Marcos P. Cavalcante.

\medskip\noindent
{\bf Baby explanation.}
For a smooth domain $\Omega$, consider
the Laplacian acting on functions
$u:\Omega\to\mathbb{R}$ satisfying the
Dirichlet boundary condition
$u|_{\partial\Omega}=0$. The first
Dirichlet eigenvalue $\lambda_1(\Omega)$
is the smallest number $\lambda>0$ for
which there exists a nonzero function $u$
such that
\[
-\Delta u=\lambda u
\qquad\text{in }\Omega.
\]
Thus, changing the shape of $\Omega$
changes the number $\lambda_1(\Omega)$.

A deformation of $\Omega$ is obtained by
moving its boundary infinitesimally along
a vector field, while preserving its
volume. The domain is called
{\it extremal} if the first-order change
of $\lambda_1(\Omega)$ vanishes for every
such deformation. This is equivalent to
saying that the first eigenfunction
reaches zero at the boundary with the
same normal slope everywhere:
\[
\partial_\nu u
=
\text{constant on }\partial\Omega.
\]
The problem is therefore to classify
domains satisfying this strong symmetry
condition.
\medskip\noindent
{\bf Plan.}
\begin{enumerate}
\item Definition of extremal domains.
\item Some previous results and examples.
\item Stability of extremal domains.
\item Classification of stable extremal
domains on $S^2$ [their result].
\item Partial results on $S^n$,
for $n \geq 3$:
\begin{enumerate}
\item The case $\partial \Omega$ minimal.
\item The case $\lambda_1 (\Omega)=n$.
\end{enumerate}
\item Index estimate for extremal domains
in $S^n$ with $\lambda_1(\Omega)=n.$
\item (Last item missing…)
\end{enumerate}

\medskip\noindent
{\bf Introduction.}
Let $(M^n,g)$ be a complete Riemannian
manifold.
For each domain $\Omega \subset M$ with
smooth boundary $\partial \Omega$, let
$\lambda_1(\Omega)$ be the first eigenvalue
of the Laplacian operator on $\Omega$ with
(zero) Dirichlet boundary condition.

Given a number $0<m<
\operatorname{Vol}(M,g)$,
we can consider $\lambda_1$ as a
function… [Missing, AI explanation:
as a function on 
$\{\operatorname{domains}\Omega\subset M
\operatorname{with fixed volume}\}$].

\medskip\noindent
{\bf Faber-Krahn inequality.}
Let $\Omega\subset \mathbb{R}^n$…[Missing]

\medskip\noindent
{\bf What are the critical points
of $\lambda_1$?}
A deformation given by a vector field
is {\it extremal} if it is
a critical point of $\lambda_1$,
i.e. when the first derivative vanishes.

\medskip\noindent
{\bf Some results }
\begin{itemize}
\item (Serrin's theorem.) 
If the overdetermined … then $\Omega$ is a 
geodesic ball.
\item (Kumerasan et al.) 
Serrin's theorem is also true in the
open hemisphere $S^n_+$ and the hyperbolic
space.
\item (Pacard-Sicbaldi, 09.)
\item (Delay-Sicbaldi, '15.)
\item (Sicbaldi, '13.)
\item (Sicbaldi, '10, Sicbaldi-Schenk '12.)
\item (Sicbaldi '12, Dai-Morbabito-Sicbaldi.)
\item (Espinar-Mazet, '17.)
\end{itemize}

\begin{example}[Shklover, 2000]
\label{xxii-egd-2026-example-shklover00}
\end{example}

\medskip\noindent
{\bf Stability.}
\begin{definition}[Stability]
\label{xxii-egd-2026-definition-duplicate-2}
We say an extramal domain
 $\Omega \subset M$ is {\it stable} 
if the second derivative is nonegative for
any volume preserving local deformation
$\Omega_t$ of $\Omega$ in $M$.
\end{definition}

\noindent
The basic tool to compute this is via the
second variation formula. Consider the
extremal domain. Let $\varphi$ be the first
positive eigenfunction of the Laplacian
operator on $\Omega$.

\begin{theorem}[Shinmakura, Cavalcante, N.]
\label{xxii-egd-2026-theorem-scn}
For any volume preserving local deformation
$\Omega_t = f_t (\Omega)$ of $\Omega$ in $M$
we have
$$
\operatorname{(formula)}
$$
…
\end{theorem}

\medskip\noindent
{\bf Consequence of second variation
formula.}
An extremal domain  $\Omega$ is stable
if and only if the quadratic form $Q$ is
nonnegative. ($Q$ 
looks like a complicated integral.)

\medskip\noindent
{\bf Characterisation of extremal domains.}
Now we start presentind the speaker's
results. The first one is the classification
of these domains on the 2-sphere:

\begin{theorem}[Cavalcante-N, 2024]
\label{xxii-egd-2026-theorem-cn24}
Let $\Omega$ be an extremal domain in $S^2$
with the canonical metric. If $\Omega$ is
stable, then $\Omega$ is a geodesic disk.
\end{theorem}

\noindent
The next result is a partial classification:

\begin{theorem}[CN-24]
\label{xxii-egd-2026-theorem-cn24-2}
Let $\Omega\subset S^n$ be extremal, $n \geq
3$. If $\Omega$ is stable and $\partial
\Omega$ is compact minimal hypersurface,
then $\Omega$ is a hemisphere.
\end{theorem}

\noindent
The next result is a partial result in the
special case when the eigenvalue is the
dimension of the sphere.

\begin{theorem}[CN]
\label{xxii-egd-2026-theorem-cn-3}
Let $\Omega \subset S^n$ be extremal with
$\lambda_1 = n$. If $\Omega$ is stable,
then $\Omega$ is a hemisphere.
\end{theorem}

\noindent
As a corollary, all the examples of expremal
domains in $S^3$ by Hines, Kolesar and
Mcgrath (2025) are unstable.

\begin{theorem}[CN]
\label{xxii-egd-2026-theorem-cn-4}
[Missing]
\end{theorem}

\medskip\noindent
{\bf Topological control for stable extremal
domains in $(M^2,g)$.}
Let $(M^2,g)$ and $\Omega \subset M$ be an
orientable compact Riemannian surface and a
stable extremal domain, respectively.

\begin{theorem}[Cavalcante-N, 24]
\label{xxii-egd-2026-theorem-cn24-5}
If (condition on Gauss curvature
$\int K$) then (options for genus and …?
\end{theorem}

\noindent
The inspiration for these results
come from Ross and Vergasta.

\medskip\noindent
{\bf Proofs of theorems.}
[Missing…]




\section{Geometry of critical constant
mean curvature surfaces in hyperbolic
$3$-spaces}
\label{xxii-egd-2026-section-critical-cmc-hyperbolic}

\medskip\noindent
Laurent Hauswirth,
Universit\'e Gustave Eiffel.

\medskip\noindent
{\bf Abstract.}
This talk surveys old and new results on
surfaces whose mean curvature takes a
critical value. The examples include
minimal surfaces in Euclidean $3$-space,
constant mean curvature one surfaces in
hyperbolic $3$-space, minimal surfaces in
the Heisenberg group, and constant mean
curvature one-half surfaces in
hyperbolic $2$-space times the real line.

\medskip\noindent
{\bf The critical mean curvature.}

Ther eis a critical value $H_0$ for
the mean curvature, separating two regimes
into $M^3(c)$:
\begin{itemize}
\item (Supercritical case ($|H|>H^0)$.)
comact CMC surfaces exist (e.g. spheres
of constant mean curvature).
\item Missing…
\end{itemize}

\medskip\noindent
{\bf Lawson correspondence.}
There is a lcoal isometric correspondence
between minimal surfaces $M \subset
\mathbb{R}^3$ and constant mean curvature
$H=1 $ surfaces $M \subset \mathbb{H}^3$.

Reminder of Gauss-Codazzi equation.

The {\it cousin} of $M$ is $M^*:=
(M, ds^2, S_+)$ where $S_+=
S+\operatorname{Id}$,
which satisfy Gauss-Codazzi (roughly).

(AI paragraph.)
The Lawson correspondence associates to a
minimal surface in $\mathbb{R}^3$ a
constant mean curvature one surface in
$\mathbb{H}^3$, called its cousin. Thus,
by classifying properly embedded $H=1$
surfaces in hyperbolic space---with the
horosphere as the basic model
(i.e. plane corresponds to horosphere)
---one can
deduce corresponding rigidity and
classification results for minimal
surfaces in Euclidean space.

\medskip\noindent
{\bf Gauss map.} The gauss map $g$ satisfies
[an equation].

For the cousin, this is map gives no
geometric information.

\medskip\noindent
{\bf Classification of properly
embedded $H=1$ surfaces.}
\begin{theorem}[Collin,-, Rosenberg, 01]
\label{xxii-egd-2026-theorem-cr01}
\begin{itemize}
\item The horosphere is the only simply
connected…
\end{itemize}
\end{theorem}

\begin{proposition}[Collin,-,R 01]
\label{xxii-egd-2026-proposition-cr01}
Let $M$ be a properly embedded $H=1$ surface
with $\partial_\infty M=\{p\}$, then $M$
is a horocylinder.
\end{proposition}

\begin{proof}
(In French.)
\end{proof}

\begin{proposition}
\label{xxii-egd-2026-proposition-cr01-2}
Let $E$ be a poperly embedded $H=1$
annular end with $\partial_\infty M \neq 
S_{\infty}$, then $E$ is a regular end,
i.e. $G$ extend meromorphically to the
puncture.
\end{proposition}

\noindent
Regular end. Introduced by Schoen.
Existence of a limit tangent plane
at the infinity of $E$.


\medskip\noindent
{\bf Several slides missing.}

 \medskip\noindent
{\bf Fundamental theorem.}

\begin{theorem}[Daniel]
\label{xxii-egd-2026-theorem-daniel}
Necessary and sufficient conditions
to get an isometric immersion with
$(ds^2,S,T,\nu)$ with $\|T\|^2+\nu^2=1$
for all…
\end{theorem}

\begin{theorem}[Daniel, 2007]
\label{xxii-egd-2026-theorem-daniel7}
Isometric correspondence between simply
connected minimal surfaces of $Nil$ and
simply connected $H=1/2$ surfaces of
$\mathbb{H}^2\times\mathbb{R}$.
\end{theorem}

\noindent
Fernandez-Mira: finda  geometric Gauss map 
$G:\Sigma \to \mathbb{H}^2$.

Examples: the rotational bowl 
and the hyperbolic bowl family.

\medskip\noindent
{\bf Conjecture.}
A simply connected properly embedded
$H=1/2$ surface is a horocylinder or an
entire graph.

\begin{theorem}[-,Rosenberg,Spruck]
\label{xxii-egd-2026-theorem-duplicate-2}
Any complete $H=1/2$ multigraph
$\nu>0$ is an entire graph on $\mathbb{H}^2$.
\end{theorem}

\begin{theorem}[Manzano-Perez-Rodriguez]
\label{xxii-egd-2026-theorem-mpr}
A stable immersed, parabolic $H=1$ surface
is an entire graph or a horocylinder.
\end{theorem}

\begin{theorem}[-,Alvaro Ramos]
\label{xxii-egd-2026-theorem-ar}
If $\partial_\infty M = \{p\}$ (or is a set
of discontinuous points) then a properly
embedded $H=1$ surface is an entire graph or
a horocylinder.
\end{theorem}

\medskip\noindent
{\bf Geometry of graphs.}

\begin{theorem}[Menezes,Rodriguez,-]
\label{xxii-egd-2026-theorem-mr}
Let $M \subset \mathbb{H}^2 \times
\mathbb{R}$ be an entire graph with $H=
\frac{1}{2}$. If $M \subset \mathbb{H}^2
\times[0,+\infty)$, then its asymptotic
boundary $\partial_\infty M \subset
(\partial_{\infty} \mathbb{H}^2 )\times \mathbb{R}$
is either empty, line, two lines… (incomplete)
\end{theorem}

\begin{theorem}[Menezes, Rodriguez, -]
\label{xxii-egd-2026-theorem-mr2}
(Missing.)
\end{theorem}

\begin{proposition}[Fernandez-Mira]
\label{xxii-egd-2026-proposition-fm}
If
$$
X^* =(F^* ,h^* ):M \to \mathbb{H}^2 \times
\mathbb{R}
$$
is a conformal immersion with CMC
$H=\frac{1}{2}$ and 
$$
X = (F,h):M \to Nil
$$
the isometric cousin then there exists
$\varepsilon  = \pm 1$ such that
$$
Y = (F,\varepsilon h ^* ):M \to \mathbb{L}^3
$$
… (missing).
\end{proposition}

\medskip\noindent
{\bf The Gauss map.}
In $\mathbb{L}^3$, the Gauss map of a CMC
$H=1/2$ space like surface take its values
in $\mathbb{H}^2$ and is harmonic.
The holomorphic quadratic differential is
given by (formula). Any harmonic Gauss map
induces $Y=(F,h) \Sigma \to \mathbb{L}^3$
with induced metric $ds^2= \tau |dz|^2$.

\medskip\noindent
{\bf Choi-Tribergs theorem-82}

\medskip\noindent
{\bf Several slides missing…}

\medskip\noindent
{\bf Asymptotic consequences.}
Along a non-compact horizontal level curve
$\gamma^*=M \cap \{x_3=s\}$ with convex
sister curve $\gamma \subset M \subset
\mathbb{L}^3$ (so $k_g(\gamma) \to 0$.

\noindent
{\bf Proposition 4.1.} If $|\omega| \geq
\varepsilon>0$ along $\gamma^*$ then
$\gamma^*$ has exactly one point on
the asymptotic boundary.

\noindent
{\bf Proposition 4.4.} If $\omega \to 0$
uniformly along $\gamma^*$ (embedded curve)
then $\gamma^*$ again has exactly one point
on the asymptotic boundary… (missing ending)
.

\section{From Gaussian to
$Q$-curvature: a conformal viewpoint}
\label{xxii-egd-2026-section-gaussian-q-curvature}

\medskip\noindent
Rayssa Caju,
Universidad de Chile.

\medskip\noindent
{\bf Abstract.}
Gaussian curvature is central to the local
and global geometry of surfaces. Its
transformation under conformal changes
connects geometry with complex analysis
and partial differential equations. This
talk asks how that theory extends to higher
dimensions and discusses curvature
quantities arising in conformal geometry,
culminating in $Q$-curvature. Originally
introduced as a fourth-order analogue of
Gaussian curvature in dimension four,
$Q$-curvature also has higher-dimensional
versions. The talk presents results for
the associated equation in compact and
noncompact settings, as well as open
problems and directions for further work.


\medskip\noindent
{\bf Introduction.} This talk explains
a 5-year programme in conformal geometry.

\begin{definition}
\label{xxii-egd-2026-definition-conformal-metrics}
Two metrics on a manifold are conformal
if they are related by a positve smooth
function $f$:
$\tilde{g}=fg$.
\end{definition}

\noindent
This means the metrics have the same angles
but not the same lengths.
Conformal geometry studies quantities
or operators that are invariant under
conformal changes.

\medskip\noindent
{\bf 2-dimensions.}
Let $(M^2,g)$ be compact and $K$ its
Gaussian curvature.

\noindent
Question: given a surface, can we modify the
metric conformally 
so that we obtain something with a
``nice shape'', i.e. homogeneous and with
constant $K$.

\noindent
Question 2: why conformally? Because
conformal changes are easiest in geometry.

If $\tilde{g}=e^{2u}g$, then
$$
\tilde{K}=e^{-2u}(K-\Lambda_gu).
$$

\begin{theorem}[Uniformization]
\label{xxii-egd-2026-theorem-uniformization-conformal}
There is $\tilde{g} \in [g]$
with constant Gaussian curvature.
\end{theorem}

\noindent
Solving this problem is equivalent to
solving a PDE:
$$
-\Lambda_g u + K
=
\tilde{K} e^{2u}
$$
on $M$.

Integrating the equation Laplace dies,
and then we use Gauss-Bonnet. Thus the sign
of $\tilde{K}$ is directly related to the
topology.

Thus,
$g \to \int K$ is a conformal invariant.

We know that the uniformization theorem
leads to the three plane geometries.
This is a complete classification of simply
connected surfaces in dimension 2,
answering the question of Poincaré.

\medskip\noindent
{\bf Compactification.}
Cohn-Vossen-Huber: a complete surface
$(M^2,g)$ with Gauss curvature in $L^1$ has
a conformal compactification $\bar{M} = 
M \cup  \{ q_1,…,q_l\}$ as a compact Riemann
surface and
$$
2\pi\chi(M)
=
\int_M K dA + \sum_{k=1}^l\nu_k
$$
where the $\nu_k$ are ``constants of
compactification''.


\medskip\noindent
{\bf 3-dimensions (and higher).}
Perhaps every person that works in conformal
geometry is used to working with the Yamabe
problem. Our main focus is to understand
{\it another} curvature measure.

Now consider a compact manifold $M^n,g$ with
$n>2$, and let $R_g$ be the scalar
curvature.

Under a conformal change
$\tilde{g} = u^{\frac{4}{n-2}}g$, then
$$
R_{\tilde{g}}=
u^{-\frac{n+2}{n-2}}\left(-\frac{4(n-1)}{n-2}\Delta_g
u + R_g u\right).
$$

{\bf Yamabe problem.} Finding a conformal
metric
$$
\tilde{g} = u^{\frac{4}{n_02}}g
$$
with constant scalar curvature.

Note that the exponent in the PDE above
appears in the Sobolev immersion.

\begin{itemize}
\item The problem is variational.
\item The sign of $R_{\tilde{g}}$ agrees
with the sign of the {\it Yamabe invariant} 
$$
Y(M,g)=
\operatorname{inf}_{\tilde{g} \in [g]}
\frac{\int_M
R_{\tilde{g}}dV_{\tilde{g}}}{(\operatorname{vol}\tilde{g})^{\frac{n-2}{n}}}
$$
\item Compacness conjecture: the set of
solutions is compact up to dimension 2.
\begin{itemize}
\item Khuri, Marques and Shoen 09, Brendle
08, and Brendle-Marques 09: figured out the
set is compact up to dimension 24 (for some
reason) and is not compact for dimension
$\geq  25$.
\end{itemize}
\end{itemize}

\medskip\noindent
{\bf 4 dimensions: what is the
$Q$-curvature.}
\begin{itemize}
\item $(M^4,g)$ = compact 4-dimensional
manifold.
\item Paneitz and Branson: fourth order
operator --- $P_g$ and a ``new'' geometric
quantity called the {\it Q-curvature}.
\item $Q$-curvature is defined by
$$
Q_g = -\frac{1}{12}
\left(\Delta_g R_g - R_g^2+3
|\operatorname{Ric}_g|^2\right)
$$
where $R_g$ is the scalar curvature and
$\operatorname{Ric}_g$ the Ricci curvature
tensor (of $g$).

\item In dimension 4, as an application of
the Chern-Gauss-Bonnet theorem we verify:
$$
\int_M\left(Q_g + \frac{1}{8}|W_g|^2_g
\right)d V_g
= 4 \pi^2\chi(M).
$$

\item Analogy between $Q_g$ and Gaussian
curvature of a surface!

\item $|W_g|^2 dV_g$ is a {\it pointwise
conformal} invariant. (The Weyl tensor
is pointwise conformal invariant.)
\item $g \to \int_M Q_g dv_g$
is a {\it conformal invariant}.
\end{itemize}

\begin{example}
\label{xxii-egd-2026-example-q-curvature}
\begin{itemize}
\item $S^4$: $Q_{g_{S^4}}=3$.
\item $\mathbb{H}^4$:
$Q_{g_{\mathbb{H}^4}}=3$.
\item Einstein metric: $Q_g > 0$. Then, by
Chern-Gauss-Bonnet, the manifold has
positive Euler characteristic.
Thus the full classification depending on
this quantity (the $Q$-curvature) is not
reasonable.
\item Any Einstein 4-manifold must have
nonnegative Euler characteristic.
\item  The same holds for manifolds with
positive $Q$-curvature.
\end{itemize}
\end{example}

\medskip\noindent
{\bf Paneitz operator.}
This operator is actually why the notion of
$Q$-curvature appeared. It's basically the
Weyl Laplacian with lower order terms:
$$
P_gu= \Delta_g^2 + \operatorname{div}_g
\left(\frac{2}{3}R_g g - 2
\operatorname{Ric}_g\right)du.
$$

\begin{itemize}
\item $\{\operatorname{constant
functions}\}\subset \Ker P_g$.
\item (Missing)
\end{itemize}

\noindent
Question: can we generalize the
uniformization theorem for higher
dimensions? Try to find metrics with
constant $Q$-curvature in 4 dimensions.

\begin{theorem}[Gursky 99, Chang-Yang 95]
\label{xxii-egd-2026-theorem-ggy}
If the $P$ is positive and the kernel
has only constant functions, and with a
hypothesis and another reasonable condition,
then $Q_g$ is constant.
\end{theorem}

\medskip\noindent
{\bf Applications in dimension 4.}
\begin{itemize}
\item On a conformally flat 4-manifold,
we have $\int_M Q= 4\pi^2\chi (M)$.
\item Chang, Qing, Yang, 00:
compactification of 4-dimensional locally
conformally flat manifolds when $Q \in L^1$.
\item See the survey ``What is
$Q$-curvature?''.
\end{itemize}

\medskip\noindent
{\bf In search of conformal invariants.}
Between the 80s and 00 they were trying to
generalize the Paneitz operator and
$Q$-curvature to every dimension. See
Feffer, Graham, Brenson (80s-00s).

If $\tilde{g} =
u^{\frac{4}{n-4}}g$…(missing)

They managed to find a generalization:

$$
P_g = \Delta_g^2 + \operatorname{div}
\left(\operatorname{something}\right)
+ \frac{n-4}{2}Q_g
$$
and the $Q$-curvature is given by
$$
Q_g=\operatorname{ugly}.
$$

\medskip\noindent
{\bf $Q$-curvature problem: the compact
case.}

\noindent
{\bf Q-curvature problem.} Given a manifold
of sidimension $\geq 5$, can I find a
conformal metric with constaant
$Q$-curvature?

Again we reduce the problem to a PDE (of
fourth order):
$$
P_gu=\frac{n(n-4)}{\operatorname{missing}}…
$$
Being variational, it means that we can
associate a function to the PDE and we try
to minimize this function.
Consider the {\it energy functional}
associated to $Q$:
$$
\mathcal{E}_g(u) =
\mathcal{E}\left(u^{\frac{4}{n-4}}g\right)
=\frac{2}{n-4}\frac{\langle
P_gu,u\rangle_{L^2}}{\|u\|^2L^{\frac{2n}{n-4}P_g}}
$$

\begin{itemize}
\item Gursky, Hang and Lin 16: introduced
another conformal invariant in the
fourth-order context:
$$
Y_4(M,g)
=
\operatorname{inf} \{\mathcal{E}_g(u)
: u \in H^2(M)\setminus \{0\}\}.
$$
\end{itemize}

\noindent
As in the Yamabe problem, there are many
solution to the $Q$-curvature equation!

\begin{itemize}
\item If $(M,g)$ is conformal to the sphere,
one has lots of solutions.
Classicaly, in the negative case, we have
only one solution.

\item In the $Q$-curvature case, we get that
even in the negative case, we have lots of
solution --- we expect more solutions to
appear (at expense of having more
constrainsts):
Bettiol, Piccione, Sirre 21: even
nonisometric conformal metrics with the same
constant negative curvature can exit.

\item How many solutions can we
find? 

\item Wei and Zhao 09: conjecture false for
$n \geq  25$.
\item Li 19, Xiong 15 compacntess in some
spacial casses.
\end{itemize}

\begin{remark}[Curious note]
\label{xxii-egd-2026-remark-curious}
The sixth-order case changes the threshold
dimension: compacntess holds in $7 \leq n
\leq  26$ blow-up examples (missing ending).
\end{remark}

\noindent
We have the counter example for dimension
higher than 24. Let's try to classify.

If I have a fixed metric closed to the
metric, I can find many metric with
arbitrarily large energy, i.e. the conformal
class has lots of high energy and high
volume metrics. That is, we can find
solutions with these properties. This shows
that we have the non-compactness property
for high-dimensional case.

The way to solve these PDEs is via bubbles.
We ``glue bubbles'' and obtain solutions
that are counter examples to the compactness
conjecture.

The problem we have been working the most on
is the noncompact setting. Consider a
manifold. We want a conformal metric with
constant Q-curvature. Remove a set: look for
a conformal Q-constant manifold
$M\setminus\Lambda$.

Relevant quantities: sign of $R_g$, sign
$Q_g$ and the dimension of $\Lambda$
(the ``singular set'').

The firs contribution in this problem was
due Hyder and Sire 20, considering $\Lambda$
has some sort of dimension. But we worked on
the case where $\Lambda$ is a set of points.

Near each of the singularities, I want to
understand the local behaviour of the
solution. They figured out the solutions
would be assymptotic to what we called
Delaunay solutions, i.e. this PDE has been
solved, solutions are known.

OK, if we understand the asympotics, we can
produce solutions.
Back in 2022, we considered a manifold with
one singularity only. Let's try to glue a ``
model'' (picture of double-bubble looking
object). We deform the manifold with a ring
function, producing a neck that we use to
glue the double bubble.

In the case of the sphere, Andrade, Wey, Ye
24 did exactly this.

Now that we understand that we have
solutions for the singular problem, we can
ask: how many solutions are there. We only
can understand this because there is a
theory for the sphere.

In a subsequent work published last year,
the speaker and collaborator understand all
the moduli space of solutions. It's a real
analytic variety with formal dimension $\#
\Lambda \geq 3$.

Can we understand this variety a little
more? The speaker and collaborators tried an
end-to-end construction. (Picture of
multiple-bubble object + gluing.) We
conclude that the moduli space is
non-contractile, i.e. topologically
nontrivial.

All the inspiration for these results is due
to the asymptotics. 

\medskip\noindent
{\bf Open problems}.

\begin{itemize}
\item Getting a full picture of the
noncompactness phenomena. (With Rodrigo, and
dos Santos.)

\item Compact + noncompact case. Put the
compacntess conjecture together with the
noncompacntess case! the importance of the
Weyl tensor is essential --- it allows us to
do the gluing. Question: can we construct
something that is non Delaunay for
dimensions higher than 25?

\item Negative case: $Q_g <0$, and $R_g<0$.
\item Mixed cases: $Q_c>0$ and $R_g<0$.
\item Is is possible to compute
$Q$-curvature invariants? (Bray and Neves do
this in $\mathbb{R}P^3$ and so on.)
Can we do it in the same way here?
But the theory of Yamabe invariant for
$Q$-curvature is only set for dimensions
greater than 6. This is not well-known yet,
work in progress.
\end{itemize}

\noindent
Special thanks to a collaborator of the
speaker (a collaborator in {\it all} her
works!). Speaker started in PDEs in Paraiba,
and then did a sandwich with Fernando… and
switched to the dark side of the force.

They're putting a group together for
geometry (PDEs and analysis) in Chile Is the same
threshold for the compactness conjecture for
the Y. See
pilot 
conference: geometric analysis in the Andes.




\section{Nearly $G_2$-manifolds and
$G_2$-Laplacian co-flows}
\label{xxii-egd-2026-section-nearly-g2-laplacian-coflows}

\medskip\noindent
Jakob Stein,
Universidade Estadual de Campinas.

\medskip\noindent
{\bf Short communication.}

\medskip\noindent
{\bf Abstract.}
Nearly $G_2$-structures on compact
$7$-manifolds define positive Einstein
metrics and are critical points, up to
scale, of geometric flows on co-closed
$G_2$-structures. This talk studies their
stability under suitably normalized
$G_2$-Laplacian co-flows. The normalization
accounts for rescaling, as normalized Ricci
flow does for Einstein metrics. Despite
this normalization, many nearly
$G_2$-structures are unstable. In
particular, those arising naturally from
$3$-Sasakian geometry are unstable, and
the standard nearly $G_2$-structure on the
round $7$-sphere is an unstable critical
point of high index. This is joint work
with Jason Lotay.


\medskip\noindent
{\bf Geometric flows.}
\begin{itemize}
\item Want to find special geometric
structures on manifolds.

\item Stable critical points of a geometric
flow (parabolic) PDE.

\item Analytically like the heat equation
$$
\frac{\partial }{\partial t}u=-\Delta u
$$
and we want to find the stable points of the
PDE. So which states converge to the stable
point.
\end{itemize}

\begin{example}[Ricci flow]
\label{xxii-egd-2026-example-ricci-flow}
Start with a metric $g_0$,
and a one-parameter family of metrics $g_t$,
 evolve the metrics with
the Ricci function:
$$
\frac{\partial }{\partial
t}g=-2\operatorname{Ric}g.
$$

\begin{itemize}
\item Feature: conformal scaling, we stay in
the same conformal class.
\item Eg. Einstein metric:
$\operatorname{Ric}(g_0)=\kappa g_0$,
$\kappa>0$:
$$
g_t= (1-2\kappa t) g_0.
$$
\item Finite-time singularity as $t \to
\frac{1}{2 \kappa}$.
\end{itemize}
\end{example}

\begin{example}[Ricci flow normalized]
\label{xxii-egd-2026-example-ricci-flow-normalised}
$$
g_t= (1-2\kappa t) g_0 + 2\kappa g.
$$
\end{example}

Higher dimensions?

\medskip\noindent
{\bf $G_2$-Laplacian co-flow} (Karigiannis,
McKay, Tsui.) A positive closed family of
4-forms on $M^7$ (determins a metric $g_t$)
solving
$$
\frac{\partial }{\partial t}\psi 
= \Delta_g \psi.
$$

Normalised version (this has to do with
conformality):
$$
\frac{\partial }{\partial t}\psi 
= \Delta_g \psi+ -\kappa^2 \psi.
$$

\begin{itemize}
\item Critical points: $n G_2$-structure (n
for nearly).
\end{itemize}

\medskip\noindent
{\bf Nearly $G_2$-structure.}
We have introduced:
$$
d(* \psi)=\kappa \psi.
$$
\begin{itemize}
\item Assicated metric $g$ is positive
Einstein.
\item Metric cone $dr^2 + r^2g$ has holonomy
$\subseteq \operatorname{Spin}(7)$.
\item Admit Killing spinors.
\end{itemize}

\begin{example}[G2 structures]
\label{xxii-egd-2026-example-g2-structures}
\begin{itemize}
\item $S^7 \subset \mathbb{R}^8$.
\item 3-Sasaki metrics.

Proper n$G_2$-structures: unique Killing
spinor (up to orientation) associate to $g$.

\item ``squashed '' Einstein metrics
(Jensen)
\item Aloff Wallach spaces
$\operatorname{SU}(2)/\operatorname{U}(1)_{p,q}$

\item Berger space
$\operatorname{SO}(2)/\operatorname{SO}(3)$.

Classified in homogeneous case by Friedrich,
Kath, Moroianu, Semmelman 1997.
\end{itemize}
\end{example}

\medskip\noindent
{\bf Modified Laplacian Co-Flow.}
\begin{itemize}
\item Problem: LCF is non-parabolic, even up
to diffeomorphisms (deTurk trick.
\item No kown short-time
existence/uniquecness via parabolic methods
\item no Shi-type esitamtes i.e. a opriori
bound on 2nd derivatives fives bound on
higher derivatives.
\end{itemize}

\noindent
{\bf Modified Laplacian co-flow (Grigorian).} 
A positive closed family of 4-forms $\psi_t$
on $M^7$ (determins a metric $g_t$) solving
$$
\frac{\partial }{\partial t}\psi
=
\Delta_g \psi + \frac{1}{2}d
((5 \kappa \gamma - 7\tau_0) * \psi)
+ \frac{5}{2}(1-\gamma)\kappa^2\psi
$$
here the $\gamma$ and the $+\frac{5}{2}…$
are the modification proposed by the speaker
and collaborator. $\tau_0$ is ``scalar
torsion'' of $\psi$, $\gamma$ constant,
$1\gg  \gamma - 2>0$.

\begin{itemize}
\item Only equivarian modification of mLCF
such that critical point n$G_2$, $d(*
\psi)=\kappa \psi$ is stable under scaling
(Bedulli-Vezzoni).
\item Existence/uniqueness, estimates
(Grigorian, Chen).
\end{itemize}

\medskip\noindent
{\bf Main results (Lota-S.)}
As I said, we want these things to be
stable. We came up with a spectral condition
such that some given nearly $G_2$ manifold
is stable.

\begin{itemize}
\item  Let $\Omega_{27}^4$ be component of
$\Omega^4(M)$ poinwise isometric to
$\operatorname{Sym}^2_0$.
\end{itemize}

\begin{theorem}[Lotay-S]
\label{xxii-egd-2026-theorem-lotays}
nmLCF is linearly stable at the
n$G_2$-structure only if the spectrum of the
operator $d *$ is not contained in 
$(-\kappa , -\frac{5}{2} (\gamma -
1)\kappa)$.
\end{theorem}


\begin{itemize}
\item Unstable conditions for flow.
\item Compare instability range (missing..)
\end{itemize}

\begin{itemize}
\item Index of critical point: dim of
unstable subspace
\end{itemize}

\begin{theorem}[Lotay-S]
\label{xxii-egd-2026-theorem-lotays2}
The standard n$G_2$ structure on $S^7$ 
for nmLCF, with index at least 7047.
\end{theorem}


\medskip\noindent
{\bf Proof.}
\begin{itemize}
\item $S^7$ admits orientation-reversing
involution, anit-commutes with $d *:
\Omega^4 \to \Omega^4 \implies $ can compute
spectral data on $\Omega^4$ using
$O(7)$ rep. theory of Laplacian (Bär 19).
\item Find lower bound on multiplicities
using explicit decomposition of $d *
$spectra on exact $\Omega^4$… (missing).
\end{itemize}

\medskip\noindent
{\bf 3-Sasaki Ansatz.}
\begin{itemize}
\item Sasaki structure: is $M^7$ admits
Sasaki structure, exists locally free action
of $\operatorname{SU}(2)$ on $M$ with
$N=M/\operatorname{SU}(2)$ an Riemmanian
orbifold with submersion metric $g_N$.
\item (missing)
$$
g=a_0^2g_N + a_1^2
\eta_1^2+a_2^2\eta_2^2+a_3^2\eta_3^2.
$$
(missing other stuff)
\end{itemize}

\begin{theorem}[Lotay-S]
\label{xxii-egd-2026-theorem-lotays3}
In the family $\psi$ with $a_1=a_2$,
given scale $\kappa >0$
\begin{itemize}
\item For any initial condition, nLCF
exists and converges to $\psi_{\pm
}^{\kappa}$.
\item Given $\gamma> 2$, $\psi_{\pm }^\kappa$
are unstable under nmLCG and have index one.
\end{itemize}
\end{theorem}

\medskip\noindent
{\bf Summary.}
We've made the natural modification of a
very primising family of flows which was
though to be a replacement of Ricci flow in
this scenario. Produces interesting class of
metric with special holonomy. Doing the
obvious thing, it fails quite drastically
for a ver natural class of nearly $G_2$
critical points that we now. We do that by
some spectral estimates. Most of the
examples are unstable. We conjecture the
rest of the examples are also unstable. We
have some preliminary work that suggests
that. There appears to be other unstable
points for this critical ansatz - ODE.
More precisely:

\medskip\noindent
{\bf Open questions.}
\begin{itemize}
\item Conjecture all n$G_2$ are unstable for
nmLCF.
\item Appears to be other (unstable)
critical points of nmLCF in the 3-Sasaki
Ansatz.
\item Dimensional reduction (hypersymplectic
flow).
\end{itemize}

\section{Min-max widths associated with
distance functions}
\label{xxii-egd-2026-section-min-max-distance-widths}

\medskip\noindent
Rafael Montezuma,
Universidade Federal do Cear\'a.

\medskip\noindent
{\bf Abstract.}
Min-max theory for the area functional can
be viewed as Morse theory on a space of
hypersurfaces, with widths appearing as
distinguished critical values. This talk
introduces a min-max width associated with
the distance between pairs of points on a
smoothly embedded circle in a complete
Riemannian manifold. The construction uses
a Morse--Lusternik--Schnirelmann theory for
the distance function and extends the
classical width of plane curves. The talk
also studies Riemannian analogues of plane
curves of constant width, comparison
inequalities with other geometric
invariants, and rigidity in the equality
cases.

\medskip\noindent
{\bf Sweepouts as movies of pairs of points
(AI).}

A sweepout can be thought of as a movie of
two points moving continuously around the
circle. Initially, the two points coincide.
They then move apart, so that the arc between
them grows from a single point, through
larger and larger arcs, until it becomes the
whole circle.

At each time $t$, the endpoints determine a
pair ${p_t,q_t}\subset\gamma$. The entire
continuous family of such pairs is called a
sweepout. For each sweepout, one considers
the largest distance attained between its
two points. The min--max width
$\mathcal{S}_g(\gamma)$ is the smallest such
maximum among all possible sweepouts.


\medskip\noindent
{\bf Introduction.} We started from the
idea, a few years ago, that the numbers that
come up in minmax theory for area functional
can be viewed as nonlinear spectrum of the
family.

\medskip\noindent
{\bf Two classical geometric quantities.}

\medskip\noindent
{\bf Birkoff invaiant.}

\begin{theorem}[Birkhof, 1917]
\label{xxii-egd-2026-theorem-birkof17}
Every Riemannian sphere contains a closed
geodesic.
\end{theorem}

\noindent
He used what we are going to call
the Birkhoff sweepout.
He found a nice critical value of the
length functional: for every sweepout (a
function from sphere to manifold), for each
one we  maximize the length and then
minimize among sweepouts. This is a critical
value of the length.
This is (realized) by a curve in the
manifold.

\medskip\noindent
{\bf Width of plane curves.}
Even older than Birkhoff's. It's the
following quantity: given a plane curve, for
each vector consider the shortest distance
between two parallel distance which contains
the curve.

\noindent
Plane curves of constant width. These are
the curves such that $w_v(\gamma) =
w(\gamma)$ for all $v \in S^1$. They are
interesting: the length of these curves is
determined by this quantity!

\medskip\noindent
{\bf A new invariant related to the last
two.}
It's a generalization of the width of planar
curves, and it's related to Birkhoff's
invariant in a fun way.

This invariant came up in a work 3 years ago
with Lucas Ambrozio and R. Santos (now at
King's College).
\begin{itemize}
\item width of embedded circles in
Riemannian manifolds.
\item Embedded circles of constant width.
\end{itemize}

\noindent
The width is generalized as a min-max
invariant.
Which is a critical point of the distance
function defined on the space of pairs of
points of the curve $\gamma \subset M^n$:
$$
d(\{p,q\})…
$$

The min-max width of $\gamma \subset M$ is
the number: the maximum distance of all
sweepouts and the minimize, just as in the
Birkhoff invariant.

$$
\mathcal{S}_g(\gamma)
=\operatorname{min} \operatorname{max}
$$


The sweepouts are the pairs in the space of
pairs of points of $\gamma$ if:
\begin{itemize}
\item the two pairs are the same at the
beginning and at the end,
\item maps parametrized by $t$ are
continuous
\item consider the arcs joining the two
points. The first one is a point and the
last one the whole curve.
\end{itemize}

\noindent
``So those are the sweepouts.''

\begin{theorem}[Ambrozio, Miranda, Montezuma]
\label{xxii-egd-2026-theorem-amm}
Every closed simple planar curve $\gamma
\subset \mathbb{R}^2$ satisfies
$$
\mathcal{S}_g(\gamma) \leq  w (\gamma).
$$
\end{theorem}

\begin{itemize}
\item Equality holds for convex curves.
\end{itemize}

\begin{theorem}[Ambrozio, Miranda, Montezuma]
\label{xxii-egd-2026-theorem-amm2}
Let $(S^2,g)$ be a Riemannian sphere. For every
embedded circle $\gamma$ in $S^2$,
$$
\mathcal{S}_g (\gamma) \leq 
\frac{1}{2}w_B(S^2,g)
$$
\end{theorem}


\begin{itemize}
\item This inequality is sharp: the standard
sphere attains it.
\item If $\varphi:\gamma \to \gamma$
is continuous without fixed points,
then every sweepout by pairs of points has a
point of the form $\{p , \varphi(p)\}$
so we can estaimate: $\mathcal{S}_g \geq
\operatorname{min}_{p \in \gamma} d_g(p ,
\varphi(p))$.
The theorem can be an effective tool to
estimate the Birkhoff invariant.
\end{itemize}

\medskip\noindent
{\bf Constructing sweepouts of a circle out
of a given family of curves.}
The first case says that the with of a plane
curve is bounded from below by the minmax
with.

Estimate the number $\mathcal{S}(\gamma)$ from above by
the vertical direction (figure).

\medskip\noindent
{\bf Sweepout by pairs of points of the
circle induced by a Morse function.}

The construction can be extended to:
\begin{theorem}[Miranda, Ambrozio, M]
\label{xxii-egd-2026-theorem-amm3}
Let $s:S^1 \to \mathbb{R}$ be a Morse
function such that all critical level sets
contain a single critical point.
Then there exists a sweepout of the circle
such that both points belong to the level
set, i.e. a sweepout $\{p_t,q_t\}$ such
that…
\end{theorem}

\begin{proof}
For this proof sketch assume that
$\{\Sigma_s\}$ is a foliation…
\end{proof}


\medskip\noindent
{\bf Non-compact surfaces with embedded
circles of uniformly bounded widths
(applications).}
Once we understand that this number is tied
with these other variational quantities.
We are curious about these other things:
what are the examples of noncompact spaces
such that the width of embedded circles is
not uniformly bounded.

\begin{theorem}[Ambrozio, Mirand, M.]
\label{xxii-egd-2026-theorem-amm4}
Let $(M^n,g)$ be a complete non-compact
Riemannian manifold. The following are
equivalent:
\begin{enumerate}
\item $\operatorname{sup}
\mathcal{S}(\gamma)$ for $\gamma$ embedded
circle is finite.
\item number of ends of $M$ is $\leq  $2 and
there exists a finite graph $G$
(1-dimensional CW-complex) and continuous
$f:M \to G$ such that the supremum of the
diameter $f^{-1}(g)$ for $g \in G$ is
finite.
\end{enumerate}
\end{theorem}

\noindent
This geometric restriction here amounts to a
macroscopic dimension in Urysohn's theory.

Let's discuss three ingredients that we use
in the proof of this theorem --
characterization of noncompact spaces for
which all curves have uniformly bounded
width.

\medskip\noindent
{\bf Ruling out three ends.}
Consider the map $\varphi:\gamma_n \to
\gamma_n$ sdefined on the long three-legged
starfish such that
\begin{itemize}
\item $\varphi$ maps one red component to
the other.
\item and similarly for the other colors
(figure…).
\item (missing…)
\end{itemize}

\noindent
So this is how you rule out the space having
three ends.

A similar picture shows that you can also
bound distance to rays in those spaces:

\medskip\noindent
{\bf Bounding distance to rays.}
Let $\sigma$ be a ray, and $\gamma$ be an
embedded circle whose sides are close to
minimizing geodesics (one them --the base--
contained in $\sigma$).

\begin{enumerate}
\item If $d$ and the base are long (figure)
then $\mathcal{S}(\gamma)$ is large.

The level sets of the Busemann function of
the ray have uniformly bounded diameters.
\end{enumerate}

\medskip\noindent
{\bf Constructing sweeputs (asuming 2).}
``This proposition is vague:''

\begin{proposition}
\label{xxii-egd-2026-proposition-2-equiv}
2 is equivalent to the existence of a
continuous function $f:M \to \mathbb{R}$
such that $\operatorname{diam}f^{-1}(t) \leq  C$ for all
$t$, i.e. the diameters of levels are
uniformly bounded.
\end{proposition}

\noindent
If $\gamma $is nice in such way so that the
composition of $f$ restricted to $\gamma$
has a single critical point in each critical
level set, our methods give a sweepout
of $\gamma$ such that $f(p_t)=f(q_t)$.
Then the hypothesis on $f$ yields
$$
\mathcal{S}(\gamma)\leq 
\operatorname{max} \{d(p_t,q_t)\}\leq C.
$$
In general, an approximation argument
on $(f \circ \gamma)$ is applied and we
combine
\begin{itemize}
\item a sweepout such that $|f(p_t)-f(q_t)|$
is small, and
\item a compacness argument which yields
(inequality involving diameter).
\end{itemize}


\medskip\noindent
{\bf Noncompact-related section.}
The systole is at least length of
homotopical nontrivial curves. We prove that
if $\sigma$ is a homotopically nontrivial
curve of length
$L_g(\sigma)=\operatorname{sys}(M,g)$ then
its width is half of the systole,
$$
\mathcal{S}_g(\sigma) =
\frac{1}{2}\operatorname{sys}(M,g).
$$
This follows from the fact that for any pair
points of $\sigma$ at leas on of the arcs
they determine is minimizing. Otherwise we
obtain a contradiction with $\sigma$ being
contractible.

\begin{lemma}[Corollary, AMM]
\label{xxii-egd-2026-lemma-amm7}
Let $(M^n,g)$, $n \geq 2$ be a Riemannian
manifold with $\pi_1(M) \neq 0$ suppose
there exists a continuous function $f :M \to
\mathbb{R}$ such that the supremum of
lengths of level sets is bounded unofrmly by
$C < \infty$. If there exists a nontrivial
curve whose length equals the systole, then
the systole is bounded by twice the
constant:
$$
\operatorname{sys}(M,g) \leq  2 C.
$$
\end{lemma}

\noindent
Indeed
$$
C \geq  \mathcal{S}_g(\sigma)
=
\frac{1}{2}L_g(\sigma)=\frac{1}{2}\operatorname{sys}(M,g)
$$
for any geodesic (something) the systole.

\medskip\noindent
{\bf Rigidity results for Zoll spheres.}
The inequality $\mathcal{S}(\gamma) \leq
\frac{1}{2}w_B(S^2,g)$ is not rigid.
E.g. metric in sphere with anitpodal
isometry, then for some closed geodesic,
there minmax number equals half the Birkhoff
invariant.
There are many examples where some embedded
circle does this.

But if we restrict ourselves to the class of
Zoll spheres (i.e. all geodesics are closed
and have the same lengths; a very large
class of Riemannian spheres).
In this settings, if there is a closed
geodesics such that the minmax = half
Birkhoff, then sphere must be round.

\begin{theorem}[AMM]
\label{xxii-egd-2026-theorem-amm8}
Precisely that above.
\end{theorem}

\begin{theorem}[AMM]
\label{xxii-egd-2026-theorem-amm9}
Zoll, intrically rotationally symmetric
and
$$
\operatorname{supp} \mathcal{S}(\gamma)
$$
for $\gamma$ embedded equals half
$w_B(S^2,g)$ then $g$ is round.
\end{theorem}

\noindent
For the proof of these rigidity results we
use the following lemma, which says that in
a Zoll sphere if we have equality of sup
embedded and Birk then every closed geodesic
contains a pair of points for which the
distance equals half the Birkhoff invariant.
Moreover, in this case,
$\operatorname{diam}(S^2,g)=\frac{1}{2}w_B(S^2,g)$,
i.e. the diameter equals half the Birkhoff
invariant.

\medskip\noindent
{\bf Another minmax invariant.}
Minmax distance width of a sphere (w/
I.Ribeiro). $\mathbb{C}P^2$ case.
A minmax number that is nontrivial critical
value of the distnace functional in the
space of…(missing).
This number coincides with the classical
width.

\begin{theorem}[AMM]
\label{xxii-egd-2026-theorem-amm10}
Let $(S^2,g)$ be a Zoll sphere.
Then
$$
w_B(S^2,g) \geq  2W_d(S^2,g)
$$
and equality holds if and only if
$(S^2,g)$ has constant Gaussian curvature.
\end{theorem}

\section{Finite-index CMC hypersurfaces
in low dimensions}
\label{xxii-egd-2026-section-finite-index-cmc}

\medskip\noindent
Ivan Miranda,
Instituto de Matem\'atica Pura e Aplicada.

\medskip\noindent
{\bf Short communication.}

\medskip\noindent
{\bf Abstract.}
Let $X$ be a six-dimensional Riemannian
manifold with nonnegative sectional
curvature which splits as the product of a
closed manifold and a Euclidean factor.
Every complete, finite-index, nonminimal
CMC hypersurface immersed in $X$ is
compact. This answers a question of
do Carmo for this class of ambient spaces
and extends results in lower dimensions.
As a consequence, the talk completes the
classification of two-sided weakly stable
CMC hypersurfaces in the round sphere
$S^6$ and real projective space
$\mathbb{R}P^6$.

\medskip\noindent
{\bf CMC hyoersurfaces} They are codimension
one immersed submanifolds…

\medskip\noindent
{\bf Weak stability}. The notion of
stability we are interested is closely
related to the classical isoperimetric
problem.

We look by variations of hypersifaces.
\begin{enumerate}
\item The first variation is CMC.
\item The second variation (formula with
integrals, has Ric and $|A|^2$.

Weakly stable CMC hypersurface, see Barbosa
and do Carmo.

One can define the notion of volume
preserving variation ofor non-compact and
immersed two-sided hypersurfaces. We search
for hyupersurfaces that minimize are up to
second order. Algebraic volume between
$\Sigma_0 $ and $\Sigma_t$.

Do Carmo-Barbosa noticed that if we take a
piece of a cylinder of height $h$ is
stable,a nd eventually becomes unstable.
They compute threshold of stability.
Physical exeriments with cylinders confirms
it.

\medskip\noindent
{\bf Strong stability.}
A CMC hypersurface is a critical point for
the … 

\begin{itemize}
\item One can presicribe the mean curvature
$H \in [0,+\infty)$ searching for critical
points $\Sigma^{n-1}\subset X^n$ with area
$A(\Sigma)$ and enclosing a volume
$V(\Sigma)$ using the functional
$F_H(\Sigma)=A(\Sigma) - H V(\Sigma)$.
\item  The strong stability inequality.
(another integral inequality with Ric and
$|A|^2$.)
\end{itemize}
\end{enumerate}

\medskip\noindent
{\bf Applications.}
\begin{itemize}
\item The isoperimetric problem (Ritoré-Ros,
Viana).
\item Foliations by CMCs (Ilias-Nelli-Soret).
\item Maximum principles at infinity
(Ros-Rosenberg).
\item Min-max existence for CMC
hypersurfaces (Mazurowski-Zhou).
\item Compacness arfuments
(Rosenberg-Souam-Toubiana).
\end{itemize}

\medskip\noindent
{\bf do Carmo's question.}
Question: Are complete noncompact weakly
stable CMC hypersurfaces immersed in
Euclidean $(n+1)$-dimensional space
necessarily minimal?

Several works show this holds in low
dimension. (Silviera for dimension 3,
Elbert-N-L for 4,5. Recenlty Cheng-Hong-Li
for 6.) The general question is open.

\begin{itemize}
\item Volume growth conditions (do
Carmo-Alencar, do Camo-Zhou,
Ilias-Nelli-Soret).
\item In $\mathbb{R}^7$, a non-minimal
example would have only one end (Fu-Li).
\item C. Pontale, B. Nelh partially anwsers
this question.
\end{itemize}

\medskip\noindent
{\bf The weak index.}

\begin{theorem}[M]
\label{xxii-egd-2026-theorem-m}
6-dimensional related.
\end{theorem}

\noindent
This is related to:

\begin{theorem}[Classification of CMC]
\label{xxii-egd-2026-theorem-classification-cmc}
A complete, two-sided weakly stable CMC
imppersed in $\mathbb{R}^{n+1}$
is either a round sphere or an affine
hyperplane.

Immersed in $S^{n+1}$ is a geodesic sphere
if $n+1 \in \{ 3,4,5,6\}$.

In $\mathbb{R}P^{n+1}$, geodesic sphere,
Clifford hypersurface or… in dimensions…
\end{theorem}

\medskip\noindent
{\bf Chodosh' question.} If $\Sigma$ is
strongly stable CMC hypersurfaced immersed
in hyperbolic space with mean curvature
$|H|\geq  n$ must $\Sigma$ be a horosphere?

Positive answer in $\mathbb{H}^3$, partial
results in dimensions 4,5.

\begin{theorem}[M]
\label{xxii-egd-2026-theorem-m2}
Partial results in dimension 6. Let
$(X^6,g)$ be Riemannian with sectional
curvature bounded from below by $-1$.
Suppose tha ction of the isometry group on
$X$ is cocompact. Then every complete,
finite index CMC hypersurface immersed in
$X$…
\end{theorem}

\noindent
The proof relies on a reduction lemma, which
states that for a cocompact action… In fact,
this can be used for both theorems since
both use cocompact actions. We want to prove
compactness. We do it by contradiction,
conclude that $M$ would have finite volume.

Highlight: surve of B. Nelli on do Carmo's
question.
\section{Solitons for mean curvature flow
in $H^2 \times \mathbb{R}$}
\label{xxii-egd-2026-section-mcf-solitons-h2r}

\medskip\noindent
Jo\~ao Paulo dos Santos,
Universidade de Bras\'ilia.

\medskip\noindent
{\bf Abstract.}
Solitons for mean curvature flow are
solutions which evolve through ambient
isometries. This talk studies solitons in
$H^2 \times \mathbb{R}$, with emphasis on
vertical translators and rotators. It
reviews rotational, parabolic and
hyperbolic symmetric translators, then uses
helicoidal symmetry to construct
rotator-translators and relate them to
vertical translators. It concludes with
the asymptotic geometry of vertical
translators and a description of the
vertical and horizontal components of
their boundary at infinity.

\medskip\noindent
{\bf Basic idea (AI).}

A translator is a special solution of mean
curvature flow which does not change shape.
The whole surface moves rigidly in one fixed
direction. (There's also rotators.)

João studies such surfaces in
$\mathbb{H}^2\times \mathbb{R}$. He imposes
symmetries, such as rotational or helicoidal
symmetry, in order to construct explicit
examples. He then studies their global shape
and how they behave near infinity.


\medskip\noindent
{\bf Mean curvature flow.}
An immersion $f:M^n \to \tilde{M}^{n+1}$,
and then $F:X \times I \to \tilde{M}$
determined by a differential equation on the
mean curvature.

\begin{itemize}
\item A soliton solution for the MCF is
given by $F(p,t)=\Gamma_t(f(p))$ where
$\Gamma_t$ is a one-parameter subgroup of
the group of isometries of $\tilde{M}$.

\item Result by Hungerbuhler-Smoczyk 00.
(This gives a good way to produce nice
solutions.)
\end{itemize}

\begin{example}[Euclidean space]
\label{xxii-egd-2026-example-euclidean-space-soliton}
\begin{itemize}
\item Evolution by translations
\item  In this case $\xi = e_3$ and $\Sigma$
satisfying $H= \langle\xi,\eta\rangle$ is
called a translator.
\item Trivial examples: vertical planes.
\end{itemize}
\end{example}

\noindent
We wanted to leave euclidean space --- find
other solutions. There's a lack of examples
in hyperbolic space.

\begin{itemize}
\item Considered the upper-half plane model.
\item Hyperbolic translations along the
$x_3$-axis defined by…
\item …
\end{itemize}

\noindent
Produced at least two examples in hyperbolic
space.

\begin{itemize}
\item Symmetric vertical tranlatos are built
from foliations by circles, horocycles, and
…
\end{itemize}

\medskip\noindent
{\bf Rotators for the MCF in $\mathbb{H}^2
\times \mathbb{R}$.}

\begin{itemize}
\item Goal: rotators with helicoidal
symmetry.
\end{itemize}

\begin{theorem}[de Lima-Ramos-dS 25]
\label{xxii-egd-2026-theorem-dlrds25}
Let $\Sigma=X(\mathbb{R}^2)$ be a helicoidal
surface of pitch $h>0$ in $\mathbb{H}^2
\times \mathbb{R}$. Then $\Sigma$ is a
rotator for the MCF if and only if…
\end{theorem}

\medskip\noindent
{\bf Recap of translators.}
Translators with rotational symmetry: the
ball goes to infinity in height and
horizontally, same for the catenoid. …

\medskip\noindent
{\bf Asymptotic boundary of vertical
translators in $\mathbb{H}^2 \times
\mathbb{R}$.}
\begin{itemize}
\item Goal: describe the vertical and
horizontal asymptotic.
\item Our approach is inspired by ideas and
techniques developed to study the asymptotic
boundary of minimal and CMC surfaces:
Nelli-Sa Earp-Santos-Toubiana 08, Sa
Earp-Toubiana 17, Kloeckner-Mazzeo 17.
\item Key ingredient: tangency principle and
symmetric vertical translators and barriers.
\end{itemize}

\noindent
Idea: approach infinity, there are two
ways:…


\begin{theorem}[Pipoli-dS-Tinglia 25]
\label{xxii-egd-2026-theorem-pdst25}
Let $M$ be a properly immersed translator in
$\mathbb{H}^2 \times \mathbb{R}$ with
(possibly empty) compact boundary….
\end{theorem}


\medskip\noindent
{\bf Work in progress.}

\begin{itemize}
\item Goal: understand the asymptotic value
of the angle function $\nu =
\langle\eta,\partial_h\rangle$ at the corner
asymptotic boundary.
\item Motivation: for symmetric examples,
one has
$$
\nu \to \pm \frac{1}{\sqrt{2}}
$$
by Lima-Pipoli 25.
\end{itemize}

\section{Tensile minimal surfaces}
\label{xxii-egd-2026-section-tensile-minimal-surfaces}

\medskip\noindent
Magdalena Rodr\'iguez,
Universidad de Granada.

\medskip\noindent
{\bf Abstract.}
In architecture, minimal surfaces belong
to the structural family of tensile
surfaces. These are lightweight,
form-active structural membranes held in
tension by cables, masts, or air pressure,
creating durable and efficient shapes.
Their study advanced significantly in the
twentieth century, mainly through the work
of Frei Otto and his team. This talk
introduces a method for producing new
minimal surfaces bounded by asymptotic
lines of constant curvature. They provide
highly efficient tensile structures in
architecture and also solve certain
symmetric thread problems.

\medskip\noindent
{\bf Notes taken at the talk.}

\medskip\noindent
{\bf Introduction.}
In architecture, minimal surfaces
belong to the family of {\it tensile
structures}: tension, but no compression
nor bending. Originally used by Frei Otto
and Institute for Lightweight Structures
(Univ. Stuttgart). Masts transfer the
force to the ground. See ``Otto eye''
new Stuttgart train station. Books of
pictures: IL1 … IL18 (for Institute
Lightweight).

Tensile catenoid: two collections of points
which are joined by cables, and soap film.
See speaker's joint work with Romane
Boutillier and Laurent Hauswirth.

Tensile $k$-noid. See Luquesio Jorge and
Bill Meeks (Top. 1983).

Last year in Granada
there was a meeting with architects and
engineers: GASM 2025 at IMAG. See videos of
talks. An introduction for how to use
geometry in architecture. Suppose we have a
surface. In order to have better structure
--better support-- we usually discretize the
surface. There are different ways of doing
that: by considering lines in the grid,
curvature lines. This is expensive to
construct. Also recommended is: considering
geodesics (which minimize 
geodesic curvature); flat lamelas, easy to
construct, cheap; lamela tangent to the
surface. Also:
consider asymptotic lines, where normal
curvature is zero; asymptotic line is
orthogonal to the surface.
See surfaces by
Eike Schling-Gridshell-TU Munich, with
flat steal lamelas.

See speaker's joint work with Boutillier and
Hauswirth (orthogonal beams along asympotic
lines and quadrilateral surfaces).

\medskip\noindent
{\bf Construction.}
Once we have the motivation for
constructing these surfaces, let's do it.

\medskip\noindent
{\bf Conjugate minimal surfaces.}
Let $M \subset \mathbb{R}^3$ be minimal,
$X=(x_1,x_2,x_3)$ which are harmonic,
which means they are the real or imaginary
part of a holomorphic map $h$.
$dh$ is called {\it height differential}.

Reminder on Weierstrass data: the Gauss map
+ height differential. We can recover the
minimal surface from this data.

The {\it conjugate surface} is obtained
by considering the Weiestrass data
and rotating the height differential.
(AI: the classical example is helicoid and
catenoid, they are conjugate.)

There's an interesting relation between
certain lines and curves on a minimal
surface and its conjugate:

\begin{lemma}
\label{xxii-egd-2026-lemma-straight-line}
$\Gamma \subset M$ (resp $M^*$) vertical
straight line $\iff$ $\Gamma ^* \subset M^*$ 
(resp $M$) horizontal
curve of symmetry.
\end{lemma}

\noindent
Asymptotic lines go to curvature lines…

(AI: the helicoid contains the vertical
$z$-axis; on its conjugate surface, the
catenoid, this becomes the horizontal waist
circle.)

At the boundary, straight lines in $\partial
M \leadsto$ symmetry curves in $\partial
M^*$.

\begin{theorem}[Joachimstahl]
\label{xxii-egd-2026-theorem-j}
$S_1, S_2$…
\end{theorem}

\begin{theorem}[Dierkes, Hildebrandt,
Küster, Wohlrab]
\label{xxii-egd-2026-theorem-dhkw}
$\Gamma \subset \mathbb{R}^3\setminus
\mathbb{B}$ piecewise smooth Jordan curve,
$\partial \Gamma \in S^2 \implies \exists M
\subset \mathbb{R}^3 \setminus \mathbb{B}$
minimal disk such that $\partial M = \Gamma
\cup  \ell$, $\ell \subset S^2$ and $M \perp
S^2$ along $\ell$. Moreover, $M$ is minimum
of the energy (and area), so $\ell$ is
smooth.
\end{theorem}

\noindent
We had to work a little more to show that
this is a good solution.

\medskip\noindent
{\bf Construction of tensile catenoids.}
Long explanation with figures and
computations.

\section{Dynamical stability of
translating solitons in hyperbolic space}
\label{xxii-egd-2026-section-dynamical-stability-solitons}

\medskip\noindent
Ronaldo Freire de Lima,
Universidade Federal do Rio Grande do
Norte.

\medskip\noindent
{\bf Short communication.}

\medskip\noindent
{\bf Abstract.}
This talk discusses the dynamical
stability of radial graphical solutions
to mean curvature flow in hyperbolic
space. In joint work with \'Alvaro Ramos,
the speaker establishes the dynamical
stability of horospheres, viewed as
translating solitons, among radial
graphical solutions to mean curvature
flow in $\mathbb{H}^{n+1}$. This is the
hyperbolic counterpart of a theorem of
Clutterbuck, Schn\"urer and Schulze, who
proved the dynamical stability of
translating paraboloids among graphical
solutions to mean curvature flow in
$\mathbb{R}^{n+1}$.

\medskip\noindent
{\bf Mean curvature flow.}
It's a one parameter family of immersions 
$F_t:M \to \bar{M}$
satisfying
$$
\frac{\partial F_t}{\partial t}=H_t
$$
where $H_t$ is the mean curvature of $F_t$.

\medskip\noindent
{\bf Three celebrities.}
Grim reaper, bowl soliton and translating
catenoid (parametrized by the radius of the
neck, must see picture).

\begin{definition}[Graphical solution
of the mean curvature flow]
\label{xxii-egd-2026-definition-graphical}

\end{definition}

\begin{lemma}
\label{xxii-egd-2026-lemma-graphical}
A solution is graphic iff solves
the PDE…
\end{lemma}

\medskip\noindent
{\bf Dynamical stability of the bowl
soliton.}

\begin{theorem}[Clitterbuck-Schnürer-Schulze
2007]
\label{xxii-egd-2026-theorem-css07}

\end{theorem}

\medskip\noindent
{\bf Graphical solutions in hyperbolic
space.}

\begin{definition}[Graphical solution in
hyperbolic space]
\label{xxii-egd-2026-definition-graphical-hyperbolic}

\end{definition}

\begin{theorem}[Unterberger, 2003]
\label{xxii-egd-2026-theorem-unter}
\begin{itemize}
\item A graphical solution to MCF in
hyperbolic space is graphical if and only if
solves the PDE…
\item For any locally Lipschitz function…
\end{itemize}
\end{theorem}


\medskip\noindent
{\bf White's avoidance principle in space
forms}

\begin{theorem}[White 2024]
\label{xxii-egd-2026-theorem-white24}
Consider a pair of smooth surfaces, possibly
noncompact, properly immersed moving by mean
curvature flow…
(AI completion:
If their initial distance in Euclidean
space is positive, then they remain disjoint at all
subsequent times.)
\end{theorem}

\begin{theorem}[Ramos et al.]
\label{xxii-egd-2026-theorem-ramos}
Horosphere-related.
\end{theorem}

\section{Alexandrov theorem for
crystalline bubbles}
\label{xxii-egd-2026-section-alexandrov-crystalline}

\medskip\noindent
Mario Santilli,
Universit\`a dell'Aquila.

\medskip\noindent
{\bf Abstract.}
This talk discusses recent progress on
bubbling phenomena for almost critical
points of anisotropic surface energies.
Consider uniformly convex norms
$\phi_h$ on $\mathbb{R}^{n+1}$ which
converge to an arbitrary norm $\phi$.
The talk studies the rigidity of
$L^1$-accumulation points of sequences
of finite-perimeter sets
$E_h \subset \mathbb{R}^{n+1}$ which are
volume-constrained almost critical points
of the anisotropic surface energy
associated with $\phi_h$. Such limits are
finite unions of disjoint, possibly
mutually tangent, $\phi$-Wulff shapes.

\medskip\noindent
{\bf Upshot.}
For the isoperimetric function critical
points are the same as minimal points.

\medskip\noindent
{\bf Introduction.}

\begin{theorem}[Alexandrov, 1958]
\label{xxii-egd-2026-theorem-alexandrov58}
LEt $\Omega \subset \mathbb{R}^{n+1}$
bounded connected with $C^2$ boundary. If the
scalar mean curvature is constant on
$\partial \Omega$, then $\Omega$ is a ball.
\end{theorem}

\noindent
Bubbles are round because they are solutions
of the classical isoperimetric problem
--- they are the least-surface-energy case.

A set of finite perimeter 
$E \subset \mathbb{R}^{n+1}$ has $C^1$
boundary, we can apply divergence theorem.
$E$ has {\it finite perimeter} if 
$$
\left\{\int_E \operatorname{div} X\right\}
$$

There is a formulation (perhaps too
abstract)
of the divergence theorem in terms of Riesz
representation theorem.

(An important result from De Giorgi.)

Upshot: the Wulff shapes (related to the
dual norm) are the solution of the
isoperimetric problem.

\medskip\noindent
{\bf Several slides missing.}

\begin{theorem}[Perpendicularity theorem, -,
et al.]
\label{xxii-egd-2026-theorem-perpendicularity}

\end{theorem}

\begin{theorem}[Alexandrov 1958, He-L-Ma-Ge
2009]
\label{xxii-egd-2026-theorem-alexandrov-anisotropic}
Anisoropic version of Alexandrov theorem?
\end{theorem}

\begin{theorem}[Delgadino-Maggi 2017]
\label{xxii-egd-2026-theorem-dm17}
Critical points of the isoperimetric
functional are possible tangential balls.
\end{theorem}

\begin{theorem}[De Rosa-Kolasinski-S,2019]
\label{xxii-egd-2026-theorem-drks}
\end{theorem}

\noindent
``Morgan's result was the starting point
of everything'':

\begin{theorem}[Morgan, 2005]
\label{xxii-egd-2026-theorem-morgan}

\end{theorem}

\medskip\noindent
Francesco Maggi: take sequence of uniformly
convex norms approaching any given norm
$\varphi$.
Now take a sequence of sets $E_h$
having anisotropic mean curvature converging
to a constant $\lambda$, that is
$$
\int_{\partial
\Omega_h}|H^{\varphi_h}-\lambda|\to 0
$$
What are the accumulation points wrt
convergence in $L^1$?


\medskip\noindent
{\bf Koiso-Palmer-Delaunay bubbling.}
…

\begin{theorem}[S, 2026]
\label{xxii-egd-2026-theorem-s26}
A general bubbling for any norm(?)
\end{theorem}

\medskip\noindent
{\bf Work in progress.}
Could this theorem be used to identify the
limit of a volume preserving crystalline
mean curvature flow in $\mathbb{R}^3$
starting from non convex data?




\section{Scalar curvature stability of
hyperbolic 3-manifolds}
\label{xxii-egd-2026-section-scalar-curvature-stability}

\medskip\noindent
Vanderson Lima,
Universidade de Bras\'ilia.

\medskip\noindent
{\bf Abstract.}
One consequence of Perelman's work on
Ricci flow is a comparison and rigidity
result for closed hyperbolic
$3$-manifolds. Let $(M,h)$ be such a
manifold, and let $g$ be a Riemannian
metric on $M$ whose scalar curvature is
at least $-6$. Then
$$
\operatorname{vol}(M,g)
\geq \operatorname{vol}(M,h).
$$
Moreover, equality implies that the two
metrics are isometric. This answered the
three-dimensional case of a conjecture of
Schoen. The talk discusses quantitative
versions of the rigidity statement. Its
main tool is the theory of Ricci flows
whose initial conditions are singular
spaces. This is joint work with Ben Lowe.

\medskip\noindent
{\bf Scalar curvature and volume.}
Reminder on definitions.

\medskip\noindent
{\bf Schoen's (volume) conjecture 
(1989).}
Suppose we have a closed hyperbolic manifold
of dimension $\geq  3$. If $g$ is a
Riemannian metric on $M$ such that
$\operatorname{Sc}_g \geq -n(n-1)$
then
$$
\operatorname{vol}_g(M) \geq  
\operatorname{vol}_h(M).
$$
This is still open, but there are some
partial results:
\begin{itemize}
\item (Besson-Coutois-Gallot, 1991.)
\item (Besson-Courtois-Gallot, 1995.) 
The conjecture holds if
$\operatorname{Ric}_g \geq -(n-1)$.
Moreover, $\operatorname{vol}_g(M)$
if and only if…
\item (Perelman/Anderson, 2003.)
The conjecture is true if $n=3$.
Moreover, $\operatorname{vol}_g(M)
= \operatorname{vol}(M)$ if and only if,
$(M^3,g)$ is isometric to $(M^3,h)$.
\end{itemize}

\medskip\noindent
{\bf Ricci flow.}
(Hamilton, 1982)
$(M,\hat{g})$ closed,
$$
\frac{\partial g_t}{\partial t}
=
-2 \operatorname{Ric}(g_t),
$$
$g_0=\hat{g}$.

\medskip\noindent
{\bf Proof of Schoen's conjecture
in dimension 3.}

\medskip\noindent
{\bf The stability problem.}
Question. Let $(M^3,h)$ be a closed
hyperbolic manifold. Let $g$ be a Riemannian
metric on $M^3$ such that
$\operatorname{Sc}\geq -6$ and …

\medskip\noindent
{\bf Flexibility on scalar curvature lower
bounds.}
\begin{itemize}
\item Bubbles (Gromov-Lawson; Schoen-Yau).
\item Splines (Gromov-Lawson;
Schoen-Yaul; Ilmanen). ``Tentacles''.
\item Drawstrings (Lee-Naber-Neumayer;
Kazaras-Xu). Take a closed geodesic
and somehow try to collapse the geometry
near it, and far from it nothing is
happening. This is a phenomenon of
codimension at least 2.
\end{itemize}

\medskip\noindent
{\bf The (counter)examples of
Kazaras-Song-Xu.}
\begin{itemize}
\item In 2007, Agol-Storm-Thurston
conjectured that for a closed Riemannian
3-manifold such that $\operatorname{Sc}_g
\geq  -6$, one should have
$E_{\operatorname{vol}}(M,g) \leq 2$.
\end{itemize}

\noindent
Here
\begin{definition}[AI]
The volume entropy of a closed Riemannian
manifold $(M,g)$ is
\[
E_{\operatorname{vol}}(M,g)
=
\lim_{R\to\infty}
\frac{1}{R}
\log
\operatorname{vol}_{\widetilde g}
\bigl(
B_{\widetilde g}(\widetilde p,R)
\bigr),
\]
where $(\widetilde M,\widetilde g)$ is the
universal cover of $(M,g)$ and
$\widetilde p\in\widetilde M$.
\end{definition}

\begin{theorem}[Kazaras-Song-Xu, 2023]
\label{xxii-egd-2026-theorem-ksx}
Family of counter examples.
\end{theorem}

\noindent
A byproduct of this construction is tht
there exist metrics $g$ with 
$\operatorname{Sc}_g \geq  -6$ and volumes
arbitrary close to $\operatorname{vol}_h$
but such that …

\medskip\noindent
{\bf Hausdorff distance reminder.}

\medskip\noindent
{\bf Back to stability.}
We are looking for limits in the GH
distance. Spaces can collapse and lower
dimension.

\medskip\noindent
{\bf Gromov compactness theorem.}
Here's a simple version of the actual
Gromov compactness theorem:

\begin{theorem}[Gromov (pre)compactness]
\label{xxii-egd-2026-theorem-gct}
Given a sequence of closed Riemannian
manifolds $(M^n_i,g_i)$ such that for all
$i$,
$$
\operatorname{Ric}_{g_i} \geq -C
$$
and
$$
\operatorname{vol}_{g_i}(B_{g_i}(x,1))
\geq v
$$
for all $x \in M_i$,
we may pass to a subsequence to ibtain GH
convergence to $(X,d)$ to some compact limit
metric space $(X,d)$ of Hausdorff dimension
$n$ whose singular set $\mathcal{S}$ is of
dimension $\leq  n-2$.
\end{theorem}

\noindent
This is a local result of regularity,
the space is a manifold outside of a
``small'' set of singularities.

\medskip\noindent
{\bf Ricci flow and singular spaces I}

\begin{theorem}[Simon 2009; Simon-Topping
2016; McLeod-Topping 2018]
\label{xxii-egd-2026-theorem-sstmt}
$M^3$ closed and a set of metrics $g_i$
with 
$\operatorname{Ric}_{g_i} \geq  C$
and
$\operatorname{vol}_{g_i}(B_{g_i}(x,1))>v$
for all $x \in M$. Then there is convergent
to 
$(X,d)$ which is homeomorphic to $M$.
\end{theorem}


\medskip\noindent
{\bf Main results.}

\begin{theorem}[L-Lowe, 2026]
\label{xxii-egd-2026-theorem-ll}
Let $C>0$ and $(M^3,h)$ be closed
hyperbolic. Let $\{g_i\}$ be such that
$\operatorname{Ric}_{g_i}\geq -C$ and
$\operatorname{Sc}_{g_i} \geq -6$,
and $\operatorname{vol}_{g_i} \to 
\operatorname{vol}_h$.
Then there is stability wrt GH toppology,
that is
$d_{GH}((M,g_i),(M,h) \to 0$ as $i \to \infty$.
\end{theorem}

\medskip\noindent
{\bf The $d_p$ distance.}
\begin{itemize}
\item For any ($M^n,g)$, we have the
following characterization of its distance
$$
d_g(x,y)=
\operatorname{sup}\{|f(x)-f(y)|:
f \in W^{1,\infty}(M),
\|\nabla_gf\| \leq 1\}.
$$
\item For $p>n$, we set
$$
d_g(x,y)=
\operatorname{sup}\{|f(x)-f(y)|:
f \in W^{1,\infty}(M),
\int_M|\nabla_gf|^pd\operatorname{vol}_g \leq 1\}.
$$
defines a distance on $M$ and for all $x,y
\in M$,
$$
d_{p,g}(x,y) \to d_g(x,y), p\to \infty.
$$

\item For $\mathbb{R}^n$,
$$
d_{p, euc}(x,y) = C_{n,p}
a|x-y|^{1-n/p}
$$
for constants $C_{n,p} \to 1$.
\item $\mathbb{H}^n$ has finite diameter!
\end{itemize}

\medskip\noindent
{\bf Ricci flow and singular spaces II.}

\begin{theorem}[Lee-Naber-Neumayer, 2020]
\label{xxii-egd-2026-theorem-lnn20}
Fir $n \in \mathbb{N}$, $p \geq  n+1$
and $\tau_0 >0$.
Then there is $\nu_0(n,p,\tau_0)>0$ such
that: is $M^n$ is a closed smooth manifold
and $g_i$ is a sequence of Riemannian
metrics such that
$$
\operatorname{Sc}_g \geq 
-n(n-1),\qquad 
\nu(g,\tau_0) \geq 
-\nu_0
$$
then $(M^n,d_{p,g},d\operatorname{vol}_{g_i}$
converge in the measured GH sense to
rectifiable Riemannian space
$(X,g,d\operatorname{vol}_g)$
which is homeomorphic to $M$.
\end{theorem}

$$
\xymatrix{
(M^n,g) \ar[d]^{GH \operatorname{conv}.}
\ar[r]^{\operatorname{Ric} flow}
&\{(M^n,g_{i,t}\}\ar[d]^{\operatorname{Ham}
\operatorname{conv}.}\\
?&  \ar[l]\{(M^n,g_t)\}
}
$$
\medskip\noindent
{\bf Perelman's entropy.}
\begin{itemize}
\item Perelman's entropy…
\item Perelman's $\nu$-functional…
\item Non-collapsing…
\end{itemize}

\medskip\noindent
{\bf Main results II.}
\begin{theorem}[L-Lowe, 2026]
\label{xxii-egd-2026-theorem-ll26-2}
[Seemingly same hypothesis as the last
theorem by speaker]. Then
$$
d_p((M,g_i),(M,h)) \to …
$$
\end{theorem}

\noindent
``We drop the Ricci condition but
pay a price for that''.






\section{Personal highlights}
\label{xxii-egd-2026-section-personal-highlights}

\begin{itemize}
\item I understood the basic idea of
``variational spectral geometry''.
Consider the Laplace operator on smooth
functions of a manifold. It has a preferred
eigenvalue $\lambda_1(M)$. As we consider
variations of $M$, we can think of 
$\lambda_1$ as a function on the 
space of variations of $M$.
We are interested in the critical points
of this function --- which are surfaces
of course.

\item A soliton is a surface $\Sigma$ such that
when we apply the (mean curvature) flow
to it, any other surface in such 
one-parameter family of surfaces $\Sigma_t$
is related to $\Sigma$ via a diffeomorphism
(AI: isometries or homotheties)
of the ambient manifold. Translators
are solitons related by translations,
and rotators by rotators.

\item Bubble-like shapes appear because
they are ``stable'' points of the mean
curvature flow, which takes a shape either
to a sphere or to a point. So these
sausage-like shapes appear.

\end{itemize}
% --- End xxii-egd-2026.tex ---

% --- Begin seiberg-witten.tex ---
\chapter{Seiberg-Witten invariants}
\label{chapter-seiberg-witten}


\section{Introduction}
\label{seiberg-witten-section-introduction}

\noindent
Course by Misha Verbitsky at IMPA,
summer 2026, which I attended.
This short overview follows the eight
lecture slide sets, ending March 2, 2026.
\href{http://verbit.ru/IMPA/SW-2026/}
{Course materials}.
References below give lecture and slide
numbers in these files.

\medskip\noindent
{\bf Big picture.}
Construct invariants of smooth
four-manifolds by studying solutions of
geometric differential equations modulo
symmetry. The course first builds the
objects needed to write the equations,
then studies their moduli spaces, and
finally translates them into holomorphic
geometry on K\"ahler surfaces.

\section{Clifford algebras and spinors}
\label{seiberg-witten-section-clifford}

\noindent
Lecture 1, slides 2--14.
A Clifford algebra encodes a quadratic
form in an associative multiplication:
$$
uv+vu=-2g(u,v)1.
$$
Its modules provide the spinor spaces.
The lecture develops the grading,
filtration, tensor products and
periodicity of Clifford algebras, then
introduces the spin group.
The purpose is to find the algebraic
objects on which the later differential
operators will act.

\section{Spin structures}
\label{seiberg-witten-section-spin}

\noindent
Lecture 2, slides 3--11;
Lecture 3, slides 3--6.
The spin group covers the special
orthogonal group. A spin structure lifts
the oriented orthonormal frame bundle
through this covering; the associated
spin representation gives a spinor bundle.
The obstruction is the second
Stiefel--Whitney class $w_2$.
This moves the algebra of Lecture 1 from
a single vector space to a manifold.

\section{Spin-c structures}
\label{seiberg-witten-section-spin-c}

\noindent
Lecture 3, slides 7--9.
Allowing an additional circle factor gives
$$
\operatorname{Spin}^c(n)
=\frac{\operatorname{Spin}(n)\times U(1)}
{\{(1,1),(-1,-1)\}}.
$$
The lifting problem now asks whether
$w_2$ has an integral lift, rather than
whether it vanishes. Complex vector
bundles supply natural examples.
This enlargement lets the construction
work beyond spin manifolds.

\section{Twistor spaces in dimension four}
\label{seiberg-witten-section-twistors}

\noindent
Lecture 4, slides 5--19.
The Hodge star splits two-forms into
self-dual and anti-self-dual parts.
The twistor space is the sphere bundle
of the self-dual part: its points encode
compatible orthogonal complex structures
on tangent spaces.
Misha describes spin-c structures using
line bundles of degree one on twistor
fibres. The Gysin sequence and Euler
class then enter his proof that every
closed oriented smooth four-manifold
admits a spin-c structure.
See also Witten, Section 2, for this
existence statement and its use.

\section{Connections and the Dirac operator}
\label{seiberg-witten-section-dirac}

\noindent
Lecture 5, slides 2--21;
Lecture 6, slides 7--9.
After constructing connections on
associated bundles, compose covariant
differentiation with Clifford
multiplication to obtain the Dirac
operator. The Weitzenb\"ock formula
expresses its square using the connection
Laplacian and curvature terms.
Integrating this identity gives vanishing
results for harmonic spinors; its
spin-c version supplies estimates for
the later equations.

\section{The Seiberg-Witten equations}
\label{seiberg-witten-section-equations}

\noindent
Lecture 6, slides 10--11;
Lecture 7, slides 4--8.
Fix a spin-c structure. The unknowns are
a connection $A$ on its determinant line
bundle and a positive spinor $\psi$:
$$
D_A\psi=0,\qquad F_A^+=q(\psi).
$$
Here $F_A^+$ is self-dual curvature and
$q$ is the quadratic map obtained from
Clifford multiplication.
Thus the spinor is harmonic for a
connection whose curvature it also
controls. See Witten, equation (2.1).

\section{From solutions to invariants}
\label{seiberg-witten-section-invariants}

\noindent
Lecture 7, slides 8--12.
Identify solutions under gauge
transformations. Curvature estimates
control the solutions; the resulting
moduli space carries a universal line
bundle whose Chern class is integrated
to produce an invariant. In dimension
zero this becomes a signed count.

\medskip\noindent
For precision, the standard construction
uses generic perturbations and an
orientation choice. Metric independence
holds for $b_2^+>1$; when $b_2^+=1$,
one must keep track of chambers.
Here $b_2^+$ is the dimension of the
self-dual harmonic two-forms.
See Witten, Section 2; this qualification
is suppressed in the slides.

\section{K\"ahler surfaces and applications}
\label{seiberg-witten-section-kahler}

\noindent
Lecture 8, slides 2--4 and 9--22.
On a K\"ahler surface, spinors become
bundle-valued differential forms.
The equations lead to holomorphic line
bundles and sections; a Kazdan--Warner
equation determines the remaining metric
data. This turns the analytic moduli
problem into a holomorphic one.
The lecture applies this description to
minimal surfaces of general type and
also discusses blow-ups, connected sums,
and Taubes' relation with symplectic
geometry and pseudoholomorphic curves.
See also Witten, Section 4.

For the different role of connections in
families of cohomology groups, see
\ref{variations-of-hodge-structures-%
section-gauss-manin} and
\ref{variations-of-hodge-structures-%
section-quantum-periods-background}.
The four-manifold invariants considered
here are distinct from the period geometry
of Seiberg--Witten curves.

\section{Exercises and references}
\label{seiberg-witten-section-references}

\noindent
The three exercise sheets cover Clifford
algebras, principal bundles, and the
Serre spectral sequence, respectively.
They supply the algebra and bundle
cohomology used in the lectures.

\medskip\noindent
Primary source: M. Verbitsky,
{\it Seiberg-Witten invariants},
IMPA 2026, lectures 1--8 and
exercise sheets 1--3, linked above.
This is a summary, not a verbatim
transcription; inconsistent numerical
normalizations in the slides are not
reproduced.

\medskip\noindent
Published reference: E. Witten,
``Monopoles and four-manifolds'',
{\it Mathematical Research Letters}
{\bf 1} (1994), 769--796,
Sections 2--4.
\href{https://arxiv.org/abs/hep-th/9411102}
{Author's version}.

% --- End seiberg-witten.tex ---

\part{Complex Geometry}
\label{part-complex-geometry}

% --- Begin complex-geometry.tex ---
\chapter{Complex Geometry}
\label{chapter-complex-geometry}





\section{Complex structures on vector spaces}
\label{complex-geometry-section-complex-structure}

\noindent
Let's discuss complex structures at the level
of linear algebra.

\begin{definition}
\label{complex-geometry-definition-compl-struc-vecto-space}
An {\it almost complex structure} on a vector
space $V$
is an endomorphism $I \in \text{End}(V)$ such
that $I^2 = -\text{Id}_V$.
\end{definition}

\begin{lemma}
\label{complex-geometry-lemma-almos-compl-struc-then-compl}
A vector space endowed with an almost complex
structure
admits a natural structure of a complex
vector space.
\end{lemma}

\begin{proof}
$(a+ib)v:=av + b I(v)$.
\end{proof}

\noindent
Conversely, given a complex vector space,
multiplication by $i$ produces an almost
complex structure.
Thus the two notions become the same.

\medskip\noindent
Now we introduce the complexification of a
real vector space.
Let $V_{\mathbb{C}}:=V
\otimes_{\mathbb{R}}\mathbb{C}$.
This means that elements of $V_{\mathbb{C}}$
are finite sums of elements of the form $v
\otimes (a+iv)$
for $a,b \in \mathbb{R}$ and $i = \sqrt{-1}$.

Notice that although $(V,I)$ is also a
complex vector space,
it is a different one from $V_\mathbb{C}$.
The complexification $V_\mathbb{C}$
is much more interesting since it gives us
the $(p,q)$-decomosition
when we take the exterior algebra.

Further, $V_\mathbb{C}=V \oplus i V$.
We define {\it conjugation} of vectors in
$V_{\mathbb{C}}$ by
$\bar{v}=\overline{v \otimes \lambda} := v
\otimes \bar{\lambda}$

This means that $\overline{V}$ is just $V$
but
with the scalar product $\lambda
v:=\bar{\lambda}v$.
for any $v \in V_{\mathbb{C}}:=V
\otimes_{\mathbb{R}}\mathbb{C}$.
This makes the equation
$\overline{\Lambda^{1,0}(V)}=\Lambda^{0,1}(V)
$,
which appears in the next proposition,
obvious.
(See \cite[p. 25]{huc}.)

\begin{lemma}
\label{complex-geometry-lemma-i-eigenspace-decomposition}
We have a decomposition
$$
V_\mathbb{C}=V^{1,0}\oplus V^{0,1}
$$
given by the eigenspaces of $I$.
$\overline{\Lambda^{1,0}}=\Lambda^{0,1}$
holds by definition of conjugation.
\end{lemma}


Now let's discuss the $p,q$ decomposition.
The upshot is the following:
$$
T^2(V \oplus  W)=
T^2(V)\otimes T^0(W)
\oplus T^1(V)\otimes T^1(W)
\oplus T^0(V)\otimes T^2(W).
$$
Explicitly, the decomposition of a general
element on the LHS is
$$
\begin{aligned}
(v_1,w_1)\otimes(v_2,w_2)
&=(v_1,0)\otimes(v_2,w_2)+(0,w_1)\otimes(v_2,
w_2)\\
&=(v_1,0)\otimes(v_2,0)+(v_1,0)\otimes(0,w_2)
\\
&+(0,w_1)\otimes(v_2,0)+(0,w_1)\otimes(0,w_2)
\end{aligned}
$$
which is valid because the tensor product of
a vector space
$T^2X$ is just the pairs $x \otimes y$
subject
to relations that $(x_1+x_2)\otimes y=x_1
\otimes y+x_2\otimes y$
and addition for $y$ and scalar product.
So we just identify elements
$$
\begin{aligned}
(v_1,0)\otimes(v_2,0)
&\leftrightarrow
v_1 \otimes v_2 \in T^2(V) \otimes
T^0(W):=T^{2,0}(V \oplus W)\\
(v_1,0)\otimes(0,w_2)
&\leftrightarrow
v_1 \otimes w_2 \in T^1(V) \otimes T^1(W):=
T^{1,1}(V\oplus W)\\
(0,w_1)\otimes(v_2,0)
&\leftrightarrow
w_1 \otimes v_2 \in T^1(W) \otimes  T^1(V)
\cong T^1(V) \otimes  T^1(W):=T^{1,1}(V\oplus
W)\\
(0,w_1)\otimes(0,w_2)
&\leftrightarrow
w_1\otimes w_2 \in T^0(V) \otimes
T^2(W):=T^{0,2}(V \oplus  W).
\end{aligned}
$$

\noindent
We have proved that any element in $T^2(V
\oplus W)$
can be written as an element in the vector
space
$T^2(V)\otimes T^0(W)
+ T^1(V)\otimes T^1(W)
+ T^0(V)\otimes T^2(W)$.
Since the intersection of the factors is
disjoint,
the sum can be taken direct.

\medskip\noindent
Notice that we have two copies of
$T^{1,1}(V\oplus W)$.
By the binomial theorem and isomorphisms
of the kind $A \otimes B \cong B \otimes A$,
we readily see that
 $$
T^k(V \oplus W)
\cong
\bigoplus_{p+q=k}\binom{k}{p}T^p(V)\otimes
T^q(W).
$$
Of course the binomial coefficient matters
little
when the underlying field is $\mathbb{C}$.

The next step is to take the quotient on the
LHS
by the relation $v \otimes v=0$,
i.e., the ideal generated by $v \otimes v$.
When we compute $(v_1,w_1)\otimes(v_1,w_1)$
as above,
this relation becomes $v_1 \otimes v_1=0$
and $w_1\otimes w_2=0$ on the first and last
components.
If we want the middle component to vanish as
well
we must ask that the isomorphism
$T^1(W) \otimes  T^1(V)\cong T^1(V) \otimes
T^1(W)$
exchanges sign. Then we obtain that
$$
\Lambda^2(V \oplus W)
=\Lambda^{2,0}(V\oplus W)
\oplus\Lambda^{1,1}(V\oplus W)
\oplus\Lambda^{0,2}(V \oplus W)
$$
by defining $\Lambda^{p,q}(V \oplus
W):=\Lambda^p(V)\otimes \Lambda^q(W)$.

To make sure that this generalizes to all
$p,q$
write, e.g. $p+q=3$,
$$
\begin{aligned}
(v_1,w_1)\otimes(v_1,w_1)\otimes(v_1,w_1)
&=(v_1,0)\otimes(v_1,0)\otimes(v_1,0)\\
&+(v_1,0)\otimes(v_1,0)\otimes(0,w_1)\\
&+…
\end{aligned}
$$

\noindent
so that imposing the relation $v \otimes v$
on $T^3(V \oplus W)$
implies imposing it on every factor.
[Further arguments needed here… but looks
like just boring computations…]

\medskip\noindent
The main result of Hodge theory is that
when we take $T^{1,0}(M)\oplus T^{0,1}(M)$,
the eigenspaces of the complex structure,
and pass to complex singular cohomology of a
Kähler manifold we get
\begin{equation}
\label{complex-geometry-equation-hodge-theorem}
H^{k}(M,\mathbb{C})
=\bigoplus_{p+q=k}H^{p,q}(M)
\end{equation}

\noindent
where $H^{p,q}(M)$ is the set of differential
forms
represented by a closed form in
$\Lambda^{p,q}(M)$.

\medskip\noindent
It also happens that
$\overline{\Lambda^{p,q}(V)}=\Lambda^{q,p}(V)
$,
which is due to the fact that
$\overline{\Lambda^{1,0}}=\Lambda^{0,1}$.


\begin{exercise}
\label{complex-geometry-exercise-1-1-forms-are-i-invariant}
Show that $(1,1)$-forms are invariant under
the action of
the complex structure $I$.
And conversely, a 2-form is of type $(1,1)$
only if it is $I$-invariant.
\end{exercise}

\begin{proof}
Let $\alpha\wedge\beta \in \Lambda^{1,1}(M)$.
Recall that
$\Lambda^{1,1}(M)=\Lambda^{1,0}(M)
\otimes\Lambda^{0,1}(M)$,
i.e. $\alpha \in \Lambda^{1,0}(M)$ and
$\beta \in \Lambda^{0,1}(M)$.
For any vectors $v,w$ tangent to some point
$p \in M$,
$$
\begin{aligned}
I(\alpha\wedge\beta)(v,w)
&=(\alpha\wedge\beta)(Iv,Iw)\\
&=\alpha(Iv)\beta(Iw)-\beta(Iv)\alpha(Iw)\\
&=-\sqrt{-1}\sqrt{-1}(\alpha(v)\beta(w)-
\beta(v)\alpha(w))\\
&=\alpha(v)\beta(w)-\beta(v)\alpha(w)\\
&=(\alpha\wedge\beta)(v,w).
\end{aligned}
$$

\noindent
The converse follows easily by the direct sum
splitting of $\Lambda^2(M)$ and noting that
a $(2,0)$ or $(0,2)$-form $\gamma$
satisfies $I\cdot\gamma(v,w)=-\gamma(v,w)$
in an analogous computation.
\end{proof}

\begin{remark}
\label{complex-geometry-remark-1-0-vectors}
By simply applying the complex structure,
one can check that a standard $1,0$ vector
is
$$
X-i JX
$$
where $X \in T_p M$. Similarly, a
$0,1$-vector is
$$
X+i J X.
$$
 
\end{remark}

\noindent
In particular this implies that $\omega$,
the fundamental form of the Riemannian metric
$g$
and the complex structure $I$ is a
$(1,1)$-form.
So this shows that Kähler classes are in
$H^{1,1}(M)$.






\noindent
Now we shall continue discussing complex
linear algebra up
to Lefschetz hard theorem
\ref{complex-geometry-theorem-lefschetz-decomposition}
and the Hodge-Riemann relations
\ref{complex-geometry-theorem-hodge-riemann-relation}.


\begin{definition}
\label{complex-geometry-definition-lefschetz-operator}
Let $(V,g,I)$ be an almost complex vector
space
equipped with an hermitian metric $g$.
The {\it Lefschetz operator} is
$$
\begin{aligned}
L: \Lambda^\bullet(V_{\mathbb{C}}^\vee)
&\longrightarrow
\Lambda^\bullet(V_{\mathbb{C}}^\vee) \\
\alpha &\longmapsto \omega\wedge\alpha.
\end{aligned}
$$
\end{definition}

\noindent
Recall that a metric on a vector space
induces a metric on the dual vector space
by defining the dual basis of an orthonormal
basis
to be orthonormal.
And for higher exterior powers all we do is
declare that wedge products of elements of
the dual orthonormal basis
(in ascending order)
form an orthonormal basis; more precisely
if $e_1,…,e_n$ is orthonormal for $V$ then
the elements of the form
$e_{i_1}\wedge…\wedge e_{i_k}$ with
$i_1<…<i_k$
are an orthonormal basis for
$\Lambda^k(V^\vee)$.


\begin{exercise}
\label{complex-geometry-exercise-hodge-star-existence}
Let $V$ be an oriented Riemannian vector
space with
orientation form $\text{Vol}$.
There is a unique operator
$*:\Lambda^\bullet(V^\vee) \to
\Lambda^\bullet(V^\vee)$
called the {\it Hodge star} such that
$$
\alpha\wedge*\beta=g(\alpha,\beta)\text{Vol}
$$
for any $\alpha,\beta \in \Lambda^k(V^\vee)$.
\end{exercise}

\begin{proof}
First prove uniqueness,
then prove that mapping a basic form
$e^{i_1}\wedge… \wedge e^{i_k}$
to $e^{j_1}\wedge…\wedge e^{j_{n-k}}$
where $(i_1,…,i_k,j_1,…,j_{n-k})$ is a
permutation of
$(1,…,n)$ will work (up to sign maybe).
\end{proof}

\noindent
Intuition on the Hodge star operator is that
it maps a form to its ``complement'' with
respect to the volume form.
For instance, $*dx=dy \wedge dz$ in
$\mathbb{R}^3$.

\begin{lemma}
\label{complex-geometry-lemma-hodge-star-isometry}
Hodge star is an isometry.
\end{lemma}

\begin{proof}
Consider two basic covectors $e^I$ and $e^J$.
Then $*e^I$ and $*e^J$ are determined by $I$
and $J$.
By definition of the metric on
$\Lambda^\bullet(M)$,
$\left\langle{}e^I,e^J\right\rangle{}$ is zero if $I\neq J$
and 1 otherwise.
Similarly, $\left\langle{}*e^I,*e^J\right\rangle{}$ is zero
if $I\neq J$,
while, if $I=J$, then it will clearly be 1.
\end{proof}

\begin{lemma}
\label{complex-geometry-lemma-hodge-star-squared}
\begin{reference}
\cite[Lemma 5.5]{voi}
\end{reference}
$$
*^2=(-1)^{k(n-k)}\text{Id}.
$$
\end{lemma}

\begin{proof}
We prove that
$$
\left\langle{}\alpha-*^2\alpha,\beta\right\rangle{}=0\qquad
\forall \beta \in \Lambda^k(M).
$$
$$
\begin{aligned}
\left\langle{}\alpha,\beta\right\rangle{}\text{Vol}
&=\left\langle{}*\alpha,*\beta\right\rangle{}\text{Vol}\\
&=*\beta\wedge *^2\alpha\\
&=(-1)^{\text{deg}(*\alpha)
\text{deg}(\beta)}*^2\alpha\wedge *\beta\\
&=\left\langle{}(-1)^{(n-k)k}*^2\alpha,
*\beta\right\rangle{}\text{Vol}.
\end{aligned}
$$
\end{proof}

\noindent
In the setting of linear algebra alone,
there is no confusion on which inner product
we use to define adjoint operators.

\begin{definition}
\label{complex-geometry-definition-dual-lefschetz}
The adjoint operator of $L$ with respect to
the metric
is called {\it dual Lefschetz operator}, i.e.
it is
the unique operator $\Lambda$ such that
$$
\left\langle{}L \alpha,\beta\right\rangle{}=\left\langle{}\alpha,
\Lambda \beta\right\rangle{}.
$$
\end{definition}

\begin{exercise}
\label{complex-geometry-exercise-dual-lefsc-hodge-twist}
$\Lambda=*^{-1}\Lambda *$.
\end{exercise}

\begin{proof}
$$
\begin{aligned}
\left\langle{}L \alpha,\beta\right\rangle{}
&=L \alpha\wedge*\beta\\
&=\omega\wedge\alpha\wedge*\beta\\
&=\alpha\wedge\omega\wedge*\beta\\
&=\alpha\wedge L * \beta\\
&=\alpha\wedge * *^{-1}L * \beta\\
&=\left\langle{}\alpha,*^{-1}L *
\beta\right\rangle{}\text{Vol}.
\end{aligned}
$$
\end{proof}

\noindent
Next we define $H$, another operator on
$\Lambda^\bullet$
called the {\it counting operator},
which is just the identity multipliied by the
codegree of a form:
$$
H|_{\Lambda^kV}:=(k-n)\text{id}.
$$
\begin{proposition}
\label{complex-geometry-proposition-sl2-relat-exter-algeb}
Let $(V,\left\langle{}\cdot,\cdot\right\rangle{})$ be an
euclidean vector space
endowed with an compatible almost complex
structure $I$.
Then
$$
[H,L]=2L,\qquad [H,\Lambda]=-2\Lambda,\qquad
[L,\Lambda]=H,
$$
that is, there is a representation of
$\mathfrak{sl}_2$
on $\Lambda^\bullet(V)$.
\end{proposition}

\begin{definition}
\label{complex-geometry-definition-primitive-form}
Let $(V,\left\langle{}\cdot,\cdot\right\rangle{})$ be an
euclidean
vector space with a compatible almost complex
structure $I$.
{\it Primitive forms} are the elements of
$\Ker \Lambda$.
The set of primitive forms of degree $k$
is denoted $P^k \subset \Lambda^k(V^\vee)$.
\end{definition}

\begin{theorem}[Lefschetz decomposition]
\label{complex-geometry-theorem-lefschetz-decomposition}
$$
\Lambda^k(V^\vee)
=
\bigoplus_i L^i(P^{k-2i})
$$
holds and is called the {\it Lefschetz
decomposition}.
Moreover, the decomposition is orthogonal
with respect to
$\left\langle{}\cdot,\cdot\right\rangle{}$.
\end{theorem}

\noindent
The upshot about this theorem
is that it is proved using the fact
that $\Lambda^\bullet(V^\vee)$ is an
$\mathfrak{sl}_2$-representation.

\begin{definition}
\label{complex-geometry-definition-hodge-riemann-pairing}
\begin{reference}
\cite[Definition 1.2.35]{huc}
\end{reference}
Let $(V,\left\langle{}\cdot,\cdot\right\rangle{})$
be an euclidean vector space with a
compatible almost complex structure.
The {\it Hodge-Riemann pairing} is

$$
\begin{aligned}
Q: \Lambda^k(V^\vee)\times \Lambda^k(^\vee)
&\longrightarrow
\Lambda^{2n}(V^\vee)\cong \mathbb{R} \\
Q(\alpha,\beta)&
=(-1)^{\frac{k(k-1)}{2}}
\alpha\wedge\beta\wedge\omega^{n-k}.
\end{aligned}
$$
\end{definition}

\begin{theorem}[Hodge-Riemann bilinear
relation]
\label{complex-geometry-theorem-hodge-riemann-relation}
$Q$ is positive definite on $P^{p,q}$, more
precisely,
$$
\sqrt{-1}^{p-q}Q(\alpha,\bar{\alpha})
=(n-(p+q))!\left\langle{}\alpha,
\alpha\right\rangle{}_{\mathbb{C}}>0,
$$
for $\alpha \in P^{p,q}\setminus\{0\}$ and
$p+q\leq n$.
\end{theorem}

\medskip\noindent

\section{$\bar\partial$-Poincaré lemma,
Dolbeault and Hartogs theorems}
\label{complex-geometry-section-dbar-poincare}
\noindent
We start by defining $\partial$ and
$\bar\partial$.
Take a real-valued smooth function
$f\in C^\infty(V,\mathbb{R})$.
The de Rham differential is given in local
coordinates $x,y$ by
$$
df=\partial_xf\,dx+\partial_yf\,dy.
$$
On the complexified exterior algebra define
$$
\partial_z
=
\frac{1}{2}(\partial_x-i\partial_y),
\qquad
\partial_{\bar z}
=
\frac{1}{2}(\partial_x+i\partial_y),
$$
and
$$
dz=dx+i\,dy,
\qquad
d\bar z=dx-i\,dy.
$$
Then
$$
\begin{aligned}
\partial_zf\,dz
+\partial_{\bar z}f\,d\bar z
&=
\frac{1}{2}
(\partial_xf-i\partial_yf)(dx+i\,dy)\\
&\quad+
\frac{1}{2}
(\partial_xf+i\partial_yf)(dx-i\,dy)\\
&=
\partial_xf\,dx+\partial_yf\,dy\\
&=df.
\end{aligned}
$$

\begin{definition}
\label{complex-geometry-definition-dbar}
Let $X$ be an almost complex manifold.
If
$$
d:
\mathcal A^k_{X,\mathbb C}
\longrightarrow
\mathcal A^{k+1}_{X,\mathbb C}
$$
is the $\mathbb C$-linear extension of the
exterior differential, define
$$
\partial
=
\Pi^{p+1,q}\circ d
$$
on $\mathcal A^{p,q}_X$, and define
$$
\bar\partial
=
\Pi^{p,q+1}\circ d
$$
on $\mathcal A^{p,q}_X$, where
$\Pi^{p,q}$ is the projection onto type
$(p,q)$.
\end{definition}

For a real or complex-valued function $f$,
$$
\partial f=\partial_zf\,dz,
\qquad
\bar\partial f
=
\partial_{\bar z}f\,d\bar z.
$$
If $f=u+iv$, the condition
$\bar\partial f=0$ is the Cauchy--Riemann
equation. Thus a complex-valued function is
holomorphic if and only if
$\bar\partial f=0$.

We shall need the basic relations between
these operators. First, here is the original
coordinate computation.

\begin{proof}
Since real partial derivatives commute,
$$
\begin{aligned}
\partial_z \partial_{\bar{z}}
&=\left(\frac{1}{2}(\partial_x - i
\partial_y)\right)
\left(\frac{1}{2}(\partial_x+i
\partial_y)\right)\\
&=\frac{1}{4}(\partial_x^2
+i \partial_x\partial_y-i\partial_y\partial_x
-i^2 \partial_y^2)\\
&=\frac{1}{4}(\partial_x^2+\partial_y^2)\\
&=\frac{1}{4}(\partial_x^2
-i \partial_x\partial_y+i\partial_y\partial_x
-i^2 \partial_y^2)\\
&=\left(\frac{1}{2}(\partial_x + i
\partial_y)\right)
\left(\frac{1}{2}(\partial_x-i
\partial_y)\right)\\
&=\partial_{\bar{z}}\partial_z.
\end{aligned}
$$

\noindent
so that
$$
\begin{aligned}
\partial \bar\partial f
&=\partial(\partial_{\bar{z}}f d\bar{z})\\
&=\partial_z \partial_{\bar{z}}f dz \wedge
d\bar{z}.
\end{aligned}
$$

\noindent
And similarly,

$$
\begin{aligned}
\bar\partial \partial f
&=\partial_{\bar{z}}(\partial_z f dz)\\
&=\partial_{\bar{z}}\partial_zf d\bar{z}
\wedge dz\\
&=-\partial_z\partial_{\bar{z}}fdz\wedge
d\bar{z}
\end{aligned}
$$

\noindent
and thus
$$
\partial\bar\partial=-\bar\partial \partial
$$
on functions.
\end{proof}

\begin{proposition}
\label{complex-geometry-proposition-leibniz-for-del}
For a $(p,q)$-form $\alpha$,
$$
\partial^2=0,
\qquad
\bar\partial^2=0,
\qquad
\partial\bar\partial
=
-\bar\partial\partial.
$$
Moreover,
$$
\partial(\alpha\wedge\beta)
=
\partial\alpha\wedge\beta
+
(-1)^{p+q}
\alpha\wedge\partial\beta,
$$
and
$$
\bar\partial(\alpha\wedge\beta)
=
\bar\partial\alpha\wedge\beta
+
(-1)^{p+q}
\alpha\wedge\bar\partial\beta.
$$
\end{proposition}

\medskip\noindent

\noindent
We first explain why the
$\bar\partial$-Poincaré lemma is the local
statement needed for Dolbeault's theorem.
Fix $p$ on a complex manifold $M$.
For each open set $U\subset M$, let
$\Omega^p_M(U)$ be the holomorphic
$p$-forms on $U$, and let
$\mathcal A^{p,q}_M(U)$ be the smooth
$(p,q)$-forms on $U$.
A differential form is a section of a
bundle of exterior powers of the cotangent
bundle. Assigning these sections to every
open $U$, with restriction maps, gives the
corresponding sheaf of sections.

Taking global sections means evaluating the
sheaf on $M$:
$$
\Gamma(M,\mathcal A^{p,q}_M)
=
\mathcal A^{p,q}_M(M).
$$
Dolbeault's theorem identifies sheaf
cohomology with the cohomology of the global
smooth-form complex:
$$
H^q(M,\Omega^p_M)
\cong
H^{p,q}_{\bar\partial}(M).
$$

\noindent
Recall that the stalk $\mathcal F_x$
consists of germs $[s]_x$ of sections near
$x$. Two sections have the same germ if they
agree on some smaller neighbourhood of $x$.
In particular,
$$
[s]_x=0
$$
if and only if $s$ vanishes on some
neighbourhood of $x$.

A complex of sheaves is exact precisely when
it is exact on every stalk. Thus a section
in a kernel only needs to lift locally near
each point. Different neighbourhoods may
have different lifts.

For a sheaf map
$\varphi:\mathcal F\to\mathcal G$, the image
presheaf may fail the gluing axiom. Its
sheafification is the image sheaf. A section
$s\in\mathcal G(U)$ belongs to this image
sheaf precisely when some open cover
$U=\bigcup_i U_i$ admits local lifts
$t_i\in\mathcal F(U_i)$ with
$$
\varphi(t_i)=s|_{U_i}.
$$
The images agree on overlaps, but the lifts
need not glue. Thus a surjective sheaf map
guarantees local lifts. The global-sections
functor asks for a lift on all of $M$, which
may not exist. This is why global sections
can lose right exactness.

\begin{theorem}
[$\bar\partial$-Poincaré lemma]
\label{complex-geometry-theorem-dbar-poincare}
Let $\eta$ be a smooth $(p,q)$-form with
$q>0$ on a neighbourhood of a point $x$ in
a complex manifold. If
$\bar\partial\eta=0$, then there is a smaller
neighbourhood $V$ of $x$ and a smooth
$(p,q-1)$-form $\beta$ on $V$ such that
$$
\eta|_V=\bar\partial\beta.
$$
\end{theorem}

\noindent
The theorem gives exactness on stalks.
Indeed, take
$$
[\alpha]_x
\in
\ker\left(
\bar\partial:
\mathcal A^{p,q}_{M,x}
\to
\mathcal A^{p,q+1}_{M,x}
\right).
$$
Choose a representative $\alpha$ on a
neighbourhood of $x$. The equality
$[\bar\partial\alpha]_x=0$ lets us shrink
until $\bar\partial\alpha=0$. The
$\bar\partial$-Poincaré lemma then gives a
smaller neighbourhood and a form $\beta$
such that
$$
\alpha=\bar\partial\beta.
$$
Therefore
$$
[\alpha]_x=\bar\partial[\beta]_x.
$$
The reverse inclusion follows from
$\bar\partial^2=0$.

In degree zero,
$\bar\partial\alpha=0$ means that the
coefficients of the $(p,0)$-form are
holomorphic. Hence
$$
0\longrightarrow\Omega^p_M
\longrightarrow\mathcal A^{p,0}_M
\xrightarrow{\bar\partial}
\mathcal A^{p,1}_M
\xrightarrow{\bar\partial}\cdots
$$
is exact. This is the Dolbeault resolution
of $\Omega^p_M$.

Applying global sections gives
$$
\mathcal A^{p,0}_M(M)
\xrightarrow{\bar\partial}
\mathcal A^{p,1}_M(M)
\xrightarrow{\bar\partial}
\mathcal A^{p,2}_M(M)
\longrightarrow\cdots.
$$
Its cohomology is, by definition,
$H^{p,q}_{\bar\partial}(M)$.
Local primitives need not agree on overlaps,
so the cohomology of global sections can be
nonzero even though the sheaf sequence is
exact.

We now compare the two resolutions from
the lecture. An injective sheaf
$\mathcal I$ has the extension property:
every map into $\mathcal I$ extends across
an inclusion. Every sheaf $\mathcal F$
admits an injective resolution
$
0\longrightarrow\mathcal F
\longrightarrow\mathcal I^0
\longrightarrow\mathcal I^1
\longrightarrow\cdots.
$
Sheaf cohomology is defined by taking global
sections first and then cohomology:
$
H^q(M,\mathcal F)
=
H^q\left(
\Gamma(M,\mathcal I^\bullet)
\right).
$
For Dolbeault's theorem we take
$\mathcal F=\Omega^p_M$.

On the other hand, the
$\bar\partial$-Poincaré lemma gives the
Dolbeault resolution
$
0\longrightarrow\Omega^p_M
\longrightarrow\mathcal A^{p,0}_M
\xrightarrow{\bar\partial}
\mathcal A^{p,1}_M
\xrightarrow{\bar\partial}\cdots.
$
The following is the fineness input.

\begin{proposition}
\label{complex-geometry-proposition-smooth-forms-fine}
\begin{reference}
\cite[Proposition 4.36]{voi}
\end{reference}
The sheaves $\mathcal A^{p,q}_M$ of smooth
differential forms are fine, hence acyclic.
\end{proposition}

\begin{proof}
They are modules over smooth functions, and
partitions of unity act by multiplication.
\end{proof}

Thus we have the two resolutions
$
0\longrightarrow\Omega^p_M
\longrightarrow\mathcal A^{p,0}_M
\xrightarrow{\bar\partial}
\mathcal A^{p,1}_M
\xrightarrow{\bar\partial}\cdots
$
and
$
0\longrightarrow\Omega^p_M
\longrightarrow\mathcal I^0
\longrightarrow\mathcal I^1
\longrightarrow\cdots.
$
The cohomology of global sections of the
first is, by definition,
$H^{p,q}_{\bar\partial}(M)$.
The cohomology of global sections of the
second is, by definition,
$H^q(M,\Omega^p_M)$.

The general comparison theorem for acyclic
resolutions says that both compute the same
cohomology. Therefore
$
H^{p,q}_{\bar\partial}(M)
\cong
H^q(M,\Omega^p_M).
$
This is Dolbeault's theorem.

The full general statement is Voisin,
Proposition 4.32. More importantly, our
version, including the proof, is Homological
Algebra, Proposition~\ref{%
homological-algebra-proposition-432-acyclic}.

In degree one the mechanism is especially
transparent. Let
$$
\mathcal Z^1
=
\ker\left(
\bar\partial:
\mathcal A^{p,1}_M
\to
\mathcal A^{p,2}_M
\right).
$$
The $\bar\partial$-Poincaré lemma gives a
short exact sequence
$$
0\longrightarrow\Omega^p_M
\longrightarrow\mathcal A^{p,0}_M
\xrightarrow{\bar\partial}
\mathcal Z^1
\longrightarrow0.
$$
The long exact sequence contains
$$
\Gamma(M,\mathcal A^{p,0}_M)
\xrightarrow{\bar\partial}
\Gamma(M,\mathcal Z^1)
\longrightarrow
H^1(M,\Omega^p_M)
\longrightarrow
H^1(M,\mathcal A^{p,0}_M).
$$
The last group vanishes by fineness, so
$$
H^1(M,\Omega^p_M)
\cong
\frac{\Gamma(M,\mathcal Z^1)}
{\bar\partial\Gamma(M,\mathcal A^{p,0}_M)}
=
H^{p,1}_{\bar\partial}(M).
$$
Repeating this with the successive kernel
sheaves gives the comparison in every
degree.

\medskip\noindent
We now prove the
$\bar\partial$-Poincaré lemma first in one
complex variable.

\begin{lemma}[One-dimensional
$\bar\partial$-Poincaré lemma]
\label{complex-geometry-lemma-dbar-poincare-one-variable}
Let $\eta$ be a smooth $(0,1)$-form on a
neighbourhood of a closed disk
$\overline D\subset\mathbb C$. Then there is
a smooth function $u$ on $D$ such that
$$
\bar\partial u=\eta.
$$
\end{lemma}

\begin{proof}
Write $z=x+iy$ and
$$
\eta=g\,d\bar z.
$$
For a smooth function $u$,
$$
\bar\partial u
=
\frac12(u_x+i u_y)d\bar z.
$$
Every $(0,1)$-form in complex dimension one
is automatically $\bar\partial$-closed.

The proof has three steps. First we expand
on a torus in Fourier modes. Second, every
nonconstant closed Fourier mode is
$\bar\partial$-exact. Third, we embed the
disk in a large elliptic curve, extend the
form with a bump function, solve there, and
restrict back to the disk.

We recall the Fourier input. In
$L^2(0,2\pi]$, the functions
$$
\frac{e^{int}}{\sqrt{2\pi}},
\qquad n\in\mathbb Z,
$$
form an orthonormal system. To prove that
they form a Hilbert basis, put
$$
A
=
\operatorname{span}
\left\{
e^{int}:n\in\mathbb Z
\right\}
\subset C(S^1,\mathbb C).
$$
The algebra $A$ contains constants, is
closed under conjugation, and separates
points. Stone--Weierstrass therefore gives
$$
\overline A^{\|\cdot\|_\infty}
=
C(S^1,\mathbb C).
$$
On $(0,2\pi]$ we have
$$
\|f-g\|_2
\leq
\sqrt{2\pi}\|f-g\|_\infty.
$$
Thus uniform approximation implies
$L^2$ approximation. Continuous functions
are dense in $L^2(S^1)$ by Real Analysis,
Proposition~\ref{%
real-analysis-proposition-cc-dense-lp},
so
$$
\overline A^{L^2}=L^2(S^1).
$$
Hence every $L^2$ function has a Fourier
expansion
$$
f(t)=\sum_{n\in\mathbb Z}a_ne^{int},
$$
where
$$
a_n
=
\frac1{2\pi}
\int_0^{2\pi}
f(t)e^{-int}\,dt.
$$

Now take the elliptic curve
$$
E
=
\mathbb C/
(2\pi\mathbb Z+2\pi i\mathbb Z).
$$
Its Fourier modes are
$$
e_{n,m}(x,y)=e^{inx+imy},
\qquad
(n,m)\in\mathbb Z^2.
$$
Their span is again a
conjugation-closed algebra containing
constants and separating points.
Stone--Weierstrass and density of continuous
functions in $L^2(E)$ show that these modes
form a Hilbert basis.

Every smooth function on $E$ has an
expansion
$$
g
=
\sum_{n,m\in\mathbb Z}
a_{n,m}e_{n,m}.
$$
The coefficients are
$$
a_{n,m}
=
\frac1{(2\pi)^2}
\int_E
g(x,y)e^{-inx-imy}\,dx\,dy.
$$
A differential form is expanded by expanding
its coefficients. Thus
$$
\eta
=
\sum_{n,m}
a_{n,m}e_{n,m}\,d\bar z.
$$

The Dolbeault operator preserves Fourier
frequency. Indeed,
$$
\bar\partial e_{n,m}
=
\frac{in-m}{2}
e_{n,m}\,d\bar z.
$$
Therefore every nonzero mode is exact. More
precisely, if $(n,m)\neq(0,0)$, then
$$
a_{n,m}e_{n,m}\,d\bar z
=
\bar\partial
\left(
\frac{2a_{n,m}}{in-m}e_{n,m}
\right).
$$
Separate the constant coefficient
$$
c=a_{0,0}.
$$
Define
$$
f
=
\sum_{(n,m)\neq(0,0)}
\frac{2a_{n,m}}{in-m}e_{n,m}.
$$
Then
$$
\bar\partial f
=
(g-c)d\bar z,
$$
and consequently
$$
\eta=c\,d\bar z+\bar\partial f.
$$

For completeness, smoothness of the Fourier
series justifies termwise differentiation.
Repeated integration by parts gives rapid
decay of the coefficients. In the more
elementary estimate used in the lecture,
bounded fourth derivatives give
$$
|a_{n,m}|
\leq
\frac{C_x}{|n|^4}
$$
for $n\neq0$, and
$$
|a_{n,m}|
\leq
\frac{C_y}{|m|^4}
$$
for $m\neq0$. Put
$$
k=\max(|n|,|m|).
$$
With $C=\max(C_x,C_y)$,
$$
|a_{n,m}|
\leq
\frac{C}{k^4}.
$$
The set of lattice points with
$\max(|n|,|m|)=k$ is the boundary of a
square and has $8k$ points. Hence the
Fourier series is absolutely and uniformly
convergent. After one derivative each term
gains at most one factor $k$, so the
corresponding shell is bounded by
$$
\frac{8C}{k^2}.
$$
By the Weierstrass M-test
(Real Analysis,
Theorem~\ref{real-analysis-theorem-m-test}),
the derivative series converges uniformly.
Then the differentiation of series theorem
(Real Analysis,
Theorem~\ref{%
real-analysis-theorem-differentiation-series})
justifies differentiating the series term by
term.

It remains to pass from the torus back to
the disk. Choose $L$ large enough that
$\overline D$ lies in the interior of the
square
$$
[-L/2,L/2]\times[-L/2,L/2].
$$
The quotient by
$L\mathbb Z+iL\mathbb Z$ gives an elliptic
curve in which the disk embeds without
identifying its points.

Choose a smooth bump function $\chi$,
supported in the given neighbourhood of
$\overline D$, and equal to one near
$\overline D$. Extend
$$
\widetilde\eta=\chi\eta
$$
by zero and periodically to the elliptic
curve. The form $\widetilde\eta$ agrees with
$\eta$ on $D$ and is automatically
$\bar\partial$-closed in complex dimension
one. The torus calculation gives
$$
\widetilde\eta
=
c\,d\bar z+\bar\partial f.
$$
On the disk, $\bar z$ is a well-defined
function and
$$
d\bar z=\bar\partial\bar z.
$$
Therefore
$$
\eta
=
\bar\partial
\left(
c\bar z+f|_D
\right).
$$
This is the required primitive.
\end{proof}

\begin{remark}
\label{complex-geometry-remark-dbar-poincare-boundary}
The closed-disk hypothesis means that the
form is smooth on a neighbourhood of the
closed disk, while the primitive is produced
on the interior. The local
$\bar\partial$-Poincaré lemma follows by
choosing a smaller disk around the point.
The stronger statement on an arbitrary open
disk, without extension across the boundary,
is a separate result.
\end{remark}

\medskip\noindent
We now repeat the construction in several
variables. The important point is that the
one-variable solution operator can be
applied in one coordinate while treating
the remaining coordinates as parameters.

Let $T_i$ denote the one-variable primitive
operator in the coordinate $z_i$. Thus
$$
\partial_{\bar z_i}(T_i g)=g.
$$
The same bump function is used for every
value of the other parameters. Therefore,
for $j\neq i$,
$$
\partial_{\bar z_j}(T_i g)
=
T_i(\partial_{\bar z_j}g).
$$
In particular, solving in one variable
preserves holomorphic dependence on the
other variables.

We first see the mechanism in two variables.
Let
$$
\alpha
=
a\,d\bar z_1+b\,d\bar z_2
$$
be a smooth $(0,1)$-form with
$\bar\partial\alpha=0$. We want a function
$u$ such that
$$
\partial_{\bar z_1}u=a,
\qquad
\partial_{\bar z_2}u=b.
$$
Construct $u=u_1+u_2$. First set
$$
u_1=T_1a.
$$
Then
$$
\partial_{\bar z_1}u_1=a.
$$
Subtract its full Dolbeault derivative:
$$
\alpha-\bar\partial u_1
=
r\,d\bar z_2,
$$
where
$$
r
=
b-\partial_{\bar z_2}u_1.
$$
To keep the first equation solved when we
add the second correction, we need
$$
\partial_{\bar z_1}u_2=0.
$$
We therefore check that the remainder is
holomorphic in $z_1$. Using commutativity of
mixed derivatives,
$$
\partial_{\bar z_1}r
=
\partial_{\bar z_1}b
-
\partial_{\bar z_2}a.
$$
On the other hand,
$\bar\partial\alpha=0$ is equivalent to
$$
\partial_{\bar z_1}b
=
\partial_{\bar z_2}a.
$$
Thus
$$
\partial_{\bar z_1}r=0.
$$
Now set
$$
u_2=T_2r.
$$
Then
$$
\partial_{\bar z_2}u_2=r
$$
and, because $T_2$ preserves holomorphic
dependence on $z_1$,
$$
\partial_{\bar z_1}u_2=0.
$$
Hence
$$
\alpha=\bar\partial(u_1+u_2).
$$

The same correction procedure works in
$n$ variables. Let
$$
\alpha
=
\sum_{j=1}^n a_j\,d\bar z_j.
$$
Finding $u$ with
$\bar\partial u=\alpha$ means solving the
$n$ equations
$$
\partial_{\bar z_j}u=a_j.
$$
Suppose the first $k-1$ equations have
already been solved. Write the remainder as
$$
R_{k-1}
=
\alpha
-
\bar\partial
(u_1+\cdots+u_{k-1})
=
\sum_{j=k}^n
r_j\,d\bar z_j.
$$
At the $k$th step we want a correction
$u_k$ satisfying
$$
\partial_{\bar z_k}u_k=r_k
$$
and, for every $i<k$,
$$
\partial_{\bar z_i}u_k=0.
$$
The latter conditions mean exactly that the
correction keeps all previous equations
solved.

The remainder is still closed:
$$
\bar\partial R_{k-1}=0.
$$
For a general $(0,1)$-form
$$
R
=
\sum_{\ell=1}^n
c_\ell\,d\bar z_\ell,
$$
closedness gives the compatibility equations
$$
\partial_{\bar z_i}c_j
=
\partial_{\bar z_j}c_i.
$$
In $R_{k-1}$ the coefficient of
$d\bar z_i$ is zero for every $i<k$.
Thus, for $i<k$ and $j\geq k$,
$$
\partial_{\bar z_i}r_j=0.
$$
Every remaining coefficient is holomorphic
in every variable whose equation has already
been solved.

Set
$$
u_k=T_kr_k.
$$
Then
$$
\partial_{\bar z_k}u_k=r_k
$$
and, for $i<k$,
$$
\partial_{\bar z_i}u_k=0.
$$
Define
$$
R_k
=
R_{k-1}
-
\bar\partial u_k.
$$
There is now no term involving
$d\bar z_1,\ldots,d\bar z_k$. The remainder
is still closed, and the same compatibility
argument prepares the next step. After $n$
corrections,
$$
R_n=0.
$$
Therefore
$$
\alpha
=
\bar\partial
(u_1+\cdots+u_n).
$$

For a general $(0,q)$-form the same
induction works coefficientwise. Suppose
that after the first $k-1$ corrections the
closed remainder $R_{k-1}$ has no factor
$d\bar z_1,\ldots,d\bar z_{k-1}$ and its
coefficients are holomorphic in
$z_1,\ldots,z_{k-1}$. Separate the terms
containing $d\bar z_k$:
$$
R_{k-1}
=
d\bar z_k\wedge B+C.
$$
Write
$$
B
=
\sum_J b_J\,d\bar z_J.
$$
Define the correction by applying $T_k$ to
each coefficient:
$$
\beta_k
=
\sum_J
(T_kb_J)\,d\bar z_J.
$$
Differentiating gives
$$
\begin{aligned}
\bar\partial\beta_k
&=
\sum_J\sum_{i=1}^n
\partial_{\bar z_i}(T_kb_J)
\,d\bar z_i\wedge d\bar z_J\\
&=
d\bar z_k\wedge B
+
\sum_J\sum_{i\neq k}
\partial_{\bar z_i}(T_kb_J)
\,d\bar z_i\wedge d\bar z_J.
\end{aligned}
$$
The first term removes precisely the
$d\bar z_k$ part. For $i<k$, the induction
hypothesis gives
$$
\partial_{\bar z_i}(T_kb_J)=0.
$$
Thus no earlier factor is reintroduced.
Define
$$
R_k
=
R_{k-1}
-
\bar\partial\beta_k.
$$
Then $R_k$ has no factor
$d\bar z_1,\ldots,d\bar z_k$ and remains
closed. Closedness implies that every
remaining coefficient $r_J$ satisfies
$$
\partial_{\bar z_i}r_J=0
$$
for $i\leq k$. Hence $R_k$ satisfies the
same induction hypothesis with $k$ in place
of $k-1$.

After the last correction, no
$d\bar z_j$ factor is available. Since
$q>0$, every term of a $(0,q)$-form needs at
least one such factor. Therefore
$$
R_n=0.
$$
Adding the corrections gives a
$(0,q-1)$-form $\beta$ satisfying
$$
\alpha=\bar\partial\beta.
$$

Finally, for a general $(p,q)$-form, write
$$
\eta
=
\sum_I dz_I\wedge\alpha_I,
$$
where each $\alpha_I$ is a $(0,q)$-form.
Since $\bar\partial dz_I=0$, the equality
$\bar\partial\eta=0$ implies
$$
\bar\partial\alpha_I=0
$$
for every $I$. Solve
$$
\bar\partial\gamma_I=\alpha_I
$$
for each $I$ and put
$$
\beta
=
(-1)^p
\sum_I
dz_I\wedge\gamma_I.
$$
Then
$$
\bar\partial\beta=\eta.
$$
This proves the
$\bar\partial$-Poincaré lemma in all
bidegrees.

\medskip\noindent
We finish with Hartogs' extension theorem.

\begin{theorem}[Hartogs extension theorem]
\label{complex-geometry-theorem-hartogs-extension}
Let $n\geq2$ and let
$K\subset\mathbb C^n$ be compact, with
$\mathbb C^n\setminus K$ connected. Every
holomorphic function on
$\mathbb C^n\setminus K$ extends uniquely to
a holomorphic function on $\mathbb C^n$.
\end{theorem}

\begin{proof}
The idea is to repair a smooth extension.
Choose a smooth function
$\widetilde f$ on all of $\mathbb C^n$
which agrees with $f$ outside a compact set.
Its error is
$$
\alpha=\bar\partial\widetilde f.
$$
This is a smooth, compactly supported,
$\bar\partial$-closed $(0,1)$-form.

Using the standard affine embedding
$\mathbb C^n\subset\mathbb P^n$, regard
$\alpha$ as a smooth $(0,1)$-form on
$\mathbb P^n$ by extending it by zero across
the hyperplane at infinity. We need a global
smooth function $\varphi$ satisfying
$$
\bar\partial\varphi=\alpha.
$$

There are two ways to justify this step.

The first uses Hodge theory. Projective space
is compact Kähler and
$$
H^1_{\mathrm{dR}}
(\mathbb P^n,\mathbb C)=0.
$$
Hodge decomposition gives
$$
H^1_{\mathrm{dR}}
(\mathbb P^n,\mathbb C)
\cong
H^{1,0}_{\bar\partial}(\mathbb P^n)
\oplus
H^{0,1}_{\bar\partial}(\mathbb P^n).
$$
Therefore
$$
H^{0,1}_{\bar\partial}(\mathbb P^n)=0.
$$
Since $\alpha$ is
$\bar\partial$-closed, it is
$\bar\partial$-exact.

The second proof makes the role of the
$\bar\partial$-Poincaré lemma explicit.
Cover $\mathbb P^n$ by small polydisks
$U_i$. On every $U_i$ the
$\bar\partial$-Poincaré lemma gives a local
function $\varphi_i$ such that
$$
\bar\partial\varphi_i=\alpha|_{U_i}.
$$
On overlaps,
$$
\bar\partial(\varphi_i-\varphi_j)=0,
$$
so
$$
g_{ij}=\varphi_i-\varphi_j
$$
is holomorphic. The functions $g_{ij}$ form
a Čech one-cocycle with values in
$\mathcal O_{\mathbb P^n}$. Using the
standard vanishing
$$
H^1(\mathbb P^n,\mathcal O_{\mathbb P^n})
=
0,
$$
this cocycle is a coboundary. Thus there are
holomorphic functions $h_i$ with
$$
g_{ij}=h_i-h_j.
$$
The corrected local primitives
$\varphi_i-h_i$ agree on overlaps and glue
to a global smooth function $\varphi$ with
$$
\bar\partial\varphi=\alpha.
$$

Because $\mathbb P^n$ is compact,
$\varphi$ is bounded on $\mathbb C^n$.
Outside the compact support of $\alpha$,
$$
\bar\partial\varphi=0,
$$
so $\varphi$ is holomorphic there.

The geometric point where $n\geq2$ enters is
that through a point outside a sufficiently
large ball one can choose a complex affine
line in an extra direction which stays
outside the ball. Restrict $\varphi$ to such
a line. It is a bounded entire function of
one complex variable, so Liouville's theorem
makes it constant.

The union of affine complex lines which miss
the compact set contains the complement of
another compact set. Hence $\varphi$ is
constant outside a compact set. Subtracting
this constant, we may assume that $\varphi$
has compact support.

Now set
$$
F=\widetilde f-\varphi.
$$
Then
$$
\bar\partial F=0,
$$
so $F$ is holomorphic on all of
$\mathbb C^n$. Outside a compact set,
$\widetilde f=f$ and $\varphi=0$, so
$F=f$ there. Since
$\mathbb C^n\setminus K$ is connected, the
identity theorem gives equality on the whole
original domain. Uniqueness follows from the
same identity theorem.
\end{proof}

\medskip\noindent
\section{Complex manifolds}
\label{complex-geometry-section-complex-manifolds}

\noindent
Finally we arrive at complex manifolds.

\begin{definition}
\label{complex-geometry-definition-complex-manifold}
A {\it complex manifold} $M$ is a smooth
manifold admitting an open cover
$\{U_\alpha\}$ and coordinate maps
$\varphi_\alpha:U_\alpha\to\mathbb{C}^n$ such
that $\varphi_\alpha\circ\varphi_\beta^{-1}$
is holomorphic on
$\varphi_\beta(U_\alpha\cap
U_\beta)\subset\mathbb{C}^n$ for all
$\alpha,\beta$.
\end{definition}

\begin{example}
\label{complex-geometry-example-projective-space}
The standard coordinate charts of projective
space
$\mathbb{P}^n$ are
$$
U_i:=\{(z_0:…:z_n): z_i \neq 0\}
$$
along with the maps
$$
\begin{aligned}
\varphi_i: U_i &\longrightarrow \mathbb{C}^n
\\
(z_0:…:z_n) &\longmapsto
\left(\frac{z_0}{z_i},…,\frac{z_{i-1}}{z_i},
\frac{z_{i+1}}{z_i},…,\frac{z_n}{z_i}\right).
\end{aligned}
$$

\noindent
Transition functions are given by … which are
clearly holomorphic.
\end{example}

\noindent
As a warning note that we can take the charts
to be
$U_i=\{z_i=1\}$,
but when considering $U_i \cap U_j$
we must be careful since this is {\it not}
$\{z_i=1\}\cap \{z_j=1\}$.
For example in $\mathbb{P}^1$, this would say
that the intersection
$U_0 \cap U_1$
is the single point $(1:1)$, but what about
$(1:2)$?

\begin{definition}
\label{complex-geometry-definition-almost-complex-manifold}
An {\it almost complex manifold} $M$ is a
smooth manifold together with
an endomorphism of the real tangent bundle
$I:TM \to TM$
such that $I^2=-\text{Id}$.
\end{definition}

\begin{proposition}
\label{complex-geometry-proposition-complex-structure}
A complex manifold has a natural almost
complex structure.
\end{proposition}

\begin{proof}
Define bundle maps $I_i$ on local
trivializations $U_i$
by rotating vectors $90$ degrees in
counter-clockwise
direction at every copy of $\mathbb{C}$ in
the tangent space $\mathbb{C}^n$.
The upshot is that since $X$ is a complex
manifold
we can interpret the $I_i$ as multiplication
by $\sqrt{-1}$
in the complex vector space structure of the
tangent space.

Then the changes of coordinates satisfy the
condition
needed to glue to a global endomorphism,
namely
$I_j=g_{ij}I_ig_{ij}^{-1}$, trivially because
the $g_{ij}$ are holomorphic and thus
$\mathbb{C}$-linear.

\noindent
{\bf Addendum.} It looks like we can just
directly
define the $I_i$ by multiplication by
$\sqrt{-1}$.
\end{proof}


\noindent
But what about the converse?
What's the condition for an almost complex
manifold to be complex?

\begin{definition}
\label{complex-geometry-definition-integ-almos-compl-struc}
An almost complex structure on a smooth
manifold $X$
is called {\it integrable} if it is induced
from a
complex structure on $X$ as in Proposition
\ref{complex-geometry-proposition-complex-structure}.
\end{definition}

\begin{theorem}[Newlander-Nirenberg]
\label{complex-geometry-theorem-newlander-nirenberg}
An almost complex structure $I$
on a smooth manifold is integrable if and
only if
$[T^{0,1},T^{0,1}]\subset T^{0,1}$.
\end{theorem}

Other equivalent conditions…

The {\it tensor} is given by
$$
N_I[X,Y]=[X,Y]+I[IX,Y]+I[X,IY]-[IX,IY]
$$




\medskip\noindent
Now let's discuss Dolbeault cohomology.


\begin{proposition}
\label{complex-geometry-proposition-dolbeault-proposition}
\begin{reference}
\cite[Proposition 2.6.11]{huc}
\end{reference}
The space $H^0(X,\Omega^p_X)$ of holomorphic
$p$-forms on a complex manifold is the
subspace
$\{\alpha \in
\mathcal{A}^{p,0}(X):\bar\partial=0\}$.
\end{proposition}

\begin{proof}
The problem is local.
We use that $f dz_{i_1}\wedge…\wedge
dz_{i+p}$ is holomorphic
if and only if $f$ is holomorphic, which is
equivalent to
$$
\bar\partial(fdz_{i_1}\wedge…\wedge dz_{i_p}
=\sum \frac{\partial f}{\partial
\overline{z}_i}d\overline{z}_i
\wedge dz_{i_1}\wedge…\wedge dz_{i_p}=0.
$$
\end{proof}

\noindent
For a more general statement we must define
the {\it $(p,q)$-Dolbeault cohomology} by
 $$
H^{p,q}(X)_{\bar\partial}
:=\frac{\Ker
(\bar\partial:\mathcal{A}^{p,q}(X) \to
\mathcal{A}^{p,q+1}(X)}
{\text{Im}(\bar\partial:\mathcal{A}^{p,q-1}
(X)
\to \mathcal{A}^{p,q}(X))}.
$$
Then

\begin{theorem}[Dolbeault]
\label{complex-geometry-theorem-dolbeault}
\begin{reference}
\cite[Corollary 2.6.21]{huc}, \cite[Theorem
4.2]{voi}
\end{reference}
The Dolbeault cohomology of $X$ computes the
cohomology
of the sheaf $\Omega^p_X$, i.e.
$H^{p,q}_{\bar\partial}(X)=H^q(X,\Omega^p_X)
$.
\end{theorem}



\section{Kähler manifolds, Kähler identities}
\label{complex-geometry-section-kahler}


\begin{definition}
\label{complex-geometry-definition-hermitian-structure}
A {\it hermitian structure} on a complex
manifold $M$
is a Riemannian metric $g$ that is also an
isometry for
the complex structure $I$, that is
$g(I\cdot,I\cdot)=g(\cdot,\cdot)$.
\end{definition}

\begin{definition}
\label{complex-geometry-definition-fundamental-form}
If $(M,I,g)$ is a hermitian manifold,
$\omega(\cdot,\cdot):=g(I\cdot,\cdot)$
is called the {\it fundamental form}.
$(M,I,g)$ is called {\it Kähler} if its
fundamental form
is closed.
\end{definition}

\noindent
As in with vector spaces endowed with almost
complex structures,
we have the operators $L,\Lambda,*$ defined
on the
cotangent spaces of a complex manifold.

\begin{remark}
\label{complex-geometry-remark-lefschetz-operator-trace}
If $\alpha$ is a 2-form, we have
$$
\Lambda \alpha
=
\operatorname{tr}_{\omega}\alpha
=
g(\alpha,\omega).
$$
More generally, there is an expression
for $\Lambda$ in terms of contractions:
$$
\Lambda \alpha = \frac{1}{2}
\sum_i J e_i \bullet (e_i \bullet \alpha)
$$
where $\bullet$ denotes contraction 
(in lack of a better symbol at this time).
\end{remark}

\begin{proposition}
\label{complex-geometry-proposition-contractions-commutators}
The following relations hold, where $\bullet$
denotes contraction:
\begin{enumerate}
\item $[X \bullet,\Lambda]=0$.
\item $[X \bullet,L]=JX^\flat \wedge$
\end{enumerate}
\end{proposition}

\begin{proof}
\begin{enumerate}
\item Exercise.
\item Uses that $\bullet$ is a derivation.
\end{enumerate}
\end{proof}

\noindent
As a corollary we have that Lefschetz
decomposition
passes to exterior algebra
of a hermitian manifold:

\begin{lemma}
\label{complex-geometry-lemma-lefsc-decom-manif}
Let $(X,g)$ be an hermitian manifold.
Then there exists a direct sum decomposition
of vector bundles
$$
\Lambda^k(X)=\bigoplus_{i\geq
0}L^i(P^{k-2i}X),
$$
there $P^{k-2i}X$ is the kernel of $\Lambda$.
\end{lemma}

\medskip\noindent
Apart from the linear-algebraic operators,
on manifolds we have other operators.

\begin{definition}
\label{complex-geometry-definition-adjoi-de-rham-diffe}
The {\it adjoint de Rham differential} is
$$
d^*:=(-1)^{n(k+1)+1}*\circ d\circ *
:\mathcal{A}^k(M) \to \mathcal{A}^{k-1}(M)
$$
where, as always, $\mathcal{A}^k(M)$ denotes
the
sections of $\Lambda^k(M):=\Lambda^k(T^*M)$.
\end{definition}

\begin{definition}
\label{complex-geometry-definition-laplace}
The {\it Laplace operator} is
$$
\Delta=d d^*+d^*d.
$$
\end{definition}

\noindent
If $M$ is a complex manifold,
then
$$
d^*=-*\circ d \circ *.
$$

Similarly we define
$$
\partial^*:=-* \circ \partial \circ *,
\qquad
\bar\partial^*:=-* \circ \bar\partial \circ
*,
$$
which yield the Laplacians
$$
\begin{aligned}
\Delta_\partial&=\partial \partial^*
+\partial^*\partial,\\
\Delta_{\bar\partial}&=\bar\partial
\bar\partial^*+\bar\partial^*\bar\partial.
\end{aligned}
$$
 

\medskip\noindent
Here's the baby version
of the Kähler identities as
explained by Beatricce, maybe based in
A. Moroianu's notes. First,
a definition of $d^c$ which {\bf applies
only for Kähler manifolds}:
$$
d^c \alpha
=
\sum_i J e^i \wedge \nabla_{e_i}\alpha,
$$
which uses a definition of Levi-Civita
on forms using permutations.

\begin{lemma}[Baby Kähler identities]
\label{complex-geometry-lemma-baby-kahler-identities}
\begin{enumerate}
\item $[L,d^*]=d^c$.
\item $[L,d]=0$.
\end{enumerate}

\noindent
Which, applying adjoints
(adjoint of contraction is wedge!), 
easily yield
\begin{enumerate}
\item $[\Lambda,d]=-(d^c)^*$.
\item $[\Lambda,d^*]=0$.
\end{enumerate}
\end{lemma}

\begin{proof}
For the second one, use that
$$
\nabla_X (\Lambda\alpha)=L \nabla_X\alpha.
$$

\end{proof}


The most important Kähler identity is:

\begin{theorem}[Kähler identity]
\label{complex-geometry-theorem-kaher-identity}
$$
\Lambda=2 \Lambda_{\bar\partial}
=2\Lambda_{\partial}.
$$
\end{theorem}



\medskip\noindent
The general definition of $d^c$ is
of course not the one above,
but the following.

\begin{definition}
\label{complex-geometry-definition-dc}
On a complex manifold $(M,I)$,
$$
d^c:=I^{-1}d I=-I d I.
$$
\end{definition}

Notice (prove) that
$$
d^c=\frac{i}{2}(\bar\partial-\partial)
$$
so that
$$
dd^c=i(\partial \bar\partial).
$$
We may also define
$$
{d^c}^*:=-* \circ d^c \circ *.
$$

\noindent
Let us put the book's version
of the Kähler identities for completeness:

\begin{proposition}[Kähler identities]
\label{complex-geometry-proposition-kahler-identities}
Let $X$ be a complex manifold endowed with a
Kähler metric $g$. Then
\begin{enumerate}
\item $[\bar\partial,L]=[\partial,L]=0$,
$[\bar\partial^*,\Lambda]=[\partial^*,
\Lambda]=0$.
\item  Etc.
\item
$\Delta_\partial=\Delta_{\bar\partial}=
\frac{1}{2}\Delta$.
\item $\Delta$ commutes with
$*,\partial,\bar\partial,\partial^*,
\bar\partial^*,L$
and $\Lambda$.
\end{enumerate}
\end{proposition}



\noindent
Now let's introduce Fubini-Study form.
The following discussion of psh functions is
because
to pass from a potential, i.e. a function,
to a $(1,1)$-form we apply $dd^c$ (or
$\partial\bar\partial$).
Now if the potential is psh,
then the Hessian will be positive definite,
which is just the same as saying that
$dd^cf$ is positive definite $(1,1)$-form.

\begin{definition}
\label{complex-geometry-definition-positive-11form}
Let $V$ be a real vector space equipped with
a complex structure $I$.
A $(1,1)$-form $\omega$ is {\it positive}
if $\omega(x,Ix)\geq 0$ for any nonzero $x
\in V$.
(Notice the $I$ is in the second entry!)
It is called {\it Hermitian} or {\it
semipositive}
if $\omega(x,Ix)>0$ for all nonzero $x \in
V$.
\end{definition}

\noindent
See \cite[Lemma 3.1.7]{huc}
to find Huybrecht's definition of a closed
{\it positive definite}
$(1,1)$-form, namely if
$\omega=\frac{i}{2}\sum h_{ij}dz_i \wedge
d\bar{z}_j$
such that $(h_{ij}(x))$ is a positive
definite hermitian matrix for any $x \in X$.

The point is that this condition
is the same as saying that the associated
Riemannian metric is in fact positive
definite.
Indeed, $\omega(\cdot,I\cdot)=g(\cdot,\cdot)$
which is positive definite.

Finally observe that numerical positivity of
the
Kähler class,
i.e. positivity with respect to the
intersection form,
follows exactly from this fact: we shall have
that
$\int_Y \omega^p>0$ for any analytic cycle
$Y$.

\begin{definition}
\label{complex-geometry-definition-psh}
A function $f$ on a complex manifold is
{\it plurisubharmonic} if $dd^cf$ is
a positive $(1,1)$-form.
\end{definition}

\noindent
Again: psh is equivalent to the Hermitian
matrix
$$
\left(\frac{\partial^2 f}{\partial
z_i\partial \bar{z}_j}\right)
$$
being positive semidefinite.
ChatGPT: it is also equivalent to: its
restriction to
any complex line being subharmonic ($\Delta
f\geq 0$).

\begin{exercise}
\label{complex-geometry-exercise-logl}
Let $z_1,…,z_n$ be the standard coordinates
on $\mathbb{C}^n$,
and $l(z_1,…,z_n):=\sum|z_i|^2$.
Prove that
$dd^c\text{log}l=\frac{dd^cl}{l}-\frac{dl
\wedge d^cl}{l^2}$
on $\Omega=\mathbb{C}^n\setminus 0$.
Prove that the Hermitian form $dd^c
\text{log} l$
is degenerate and can be written as
$(-\sqrt{-1})\sum_{i=1}^{n-1}
\alpha_i\xi_i\wedge \bar{\xi}_i$,
where $\xi_i$ is a local frame in
$\Lambda^{1,0}(\Omega)$
and $\alpha_i$ are positive functions.
\end{exercise}

\begin{proof}
For the first equation we just use the fact
(why?) that both $d^c$ and $d$ satisfy the
chain rule.
For the second equation I computed that
\begin{equation}
\label{complex-geometry-equation-ddcl1}
dd^cl=\sqrt{-1}\sum dz_i\wedge d\bar{z}_i.
\end{equation}
But I also need to compute
$d^cl=\frac{\sqrt{-1}}{2}(\bar\partial-
\partial)l$
and $dl=(\partial+\bar\partial)l$.
I find that $dl \wedge d^cl=\sqrt{-1}\partial
l \wedge \bar\partial l$.
But I also found that
$\partial l=\sum \bar{z}_idz_i$ and
$\bar\partial l=\sum z_i d\bar{z}_i$.
So I get
\begin{equation}
\label{complex-geometry-equation-ddcl2}
dl\wedge d^cl=\sqrt{-1}
\left(\sum \bar{z}_i dz_i\right)
\wedge
\left(\sum z_i d \bar{z}_i\right)
=\sqrt{-1}\sum\bar{z}_iz_jdz_i\wedge
d\bar{z}_j
\end{equation}

\noindent
Now we join Equations \ref{complex-geometry-equation-ddcl1}
and \ref{complex-geometry-equation-ddcl2}
to find the coefficient functions with
respect to $dz_i \wedge d\bar{z}_i$.
From Equation \ref{complex-geometry-equation-ddcl1} we get
coefficient functions $\delta_{ij}$,
and from Equation \ref{complex-geometry-equation-ddcl2} we get
$\bar{z}_iz_j$ so
$$
dd^c\text{log}l=\sqrt{-1}\sum
\left(\frac{\delta_{ij}}{l}
-\frac{\bar{z}_iz_j}{l^2}\right)
dz_i \wedge d\bar{z}_j.
$$
I'm not sure how to prove that  the
coefficient functions
are positive… Or that the form is degenerate…
But fortunately this can be checked at
\cite[p. 118]{huc}.
\end{proof}

\noindent
The point of this exercise is to show that
the form
$dd^c\ell$ is almost an Hermitian form ---
but it has a kernel.
I think the point of the next exercises in
this
\href{http://verbit.ru/IMPA/Surfaces-2025/%
assign-surfaces-2025-04.pdf}{
handout}
is to show that the kernel is actually
the ``radial directions''.
And the point of that is to take quotient
by $\mathbb{C}^*$ action, which is literally
passing to the projective space.

A bit more formally:

\begin{definition}
\label{complex-geometry-definition-basic-form}
Let $M$ be a manifold and $\omega \in
\Omega^k(M)$.
Suppose that $\pi:M \to B$ is a surjective
submersion with connected fibers.
We say that $\omega$ is {\it basic} (with
respect
to $\pi$) is there exists a form
$\overline{\omega}\in \Omega^k(B)$ such that
$\pi^*  \overline{\omega}$.
\end{definition}

\begin{exercise}
\label{complex-geometry-exercise-basic-forms-characterization}
\begin{enumerate}
\item Show that $\omega$ is basic if and only
if
$i_X \omega=0$ and $L_X\omega=0$
for all vector fields $X$ tangent
to the fibers of $\pi$.
In particular, if $\omega$ is closed,
show that it is basic if
$\Ker (\pi_*) \subseteq \Ker(\omega)$
(pointwise in $M$).

\item Suppose that $\omega$ is a closed
2-form on $M$ and $\Ker(\pi_*)=\Ker(\omega)$.
Show that $\omega=\pi^*\overline{\omega}$
and $\overline{\omega}\in\Omega^2(B)$
is symplectic.

\item (Application to reduction.)
\end{enumerate}
\end{exercise}

\begin{proof}
\begin{enumerate}
\item Since $\pi$ is surjective, for any
vector $\bar{X}$ tangent to $B$ there is
a vector $X$ tangent to $M$ such that
$\pi_*X=\bar{X}$.
Define a form
$\overline{\omega}(\bar{X}^1,…,\bar{X}^k):=
\omega(X^1,…,X^k)$.

To show this does not depend on the choice of
lift of $\bar{X}$
just notice that the tangent space of $M$ can
be written
$T_pM \simeq T_pF_b \oplus \pi^*(T_bB)$ where
$\pi(p)=b$.
Since $\omega$ kills vectors tangent to the
fibers,
$\overline{\omega}$ is well-defined.

Notice that there can be different choices of
$p$ such that
$\pi(p)=b$. We must show that the definition
is also
independent of this choice.
Let $X \in T_pM$ and $Y \in T_qM$ be such
that
$\pi(p)=\pi(q)=b \in B$ and $\pi_*X=\pi_*Y$.
We wish to show that
$\omega(…,X,…)=\omega(…,Y,…)$.

This a little more involved so I'll just
sketch it for now.
If there is a vector field $V$ tangent to the
fibers
whose flow $\varphi$ joins $p$ and $q$, we
are done,
since we pull back $\omega$ along this flow,
and by the condition of the Lie derivative,
i.e. $L_V\omega$ for $V$
tangent to the fiber, we see that
$\varphi^*\omega=\omega$.
(Indeed,
$L_V\omega=\frac{d}{dt}\varphi^*\omega$,
so the derivative is zero and at $t=0$ we
have $\omega$.)

Note that the choice of $Y$ within $T_qM$
is superfluous since we already know that
$\omega$
is constant within every tangent space at a
poi

…
\end{enumerate}
\end{proof}




\begin{exercise}
\label{complex-geometry-exercise-kahler-form-conformal-class}
Let $X$ be a connected complex manifold of
dimension $n>1$
and let $g$ be a Kähler metric.
Show that $g$ is the only Kähler metric in
its conformal class,
i.e. if $g'=e^fg$ is Kähler then $f$ is
constant.
\end{exercise}

\begin{exercise}
\label{complex-geometry-exercise-lefsc-injec-1-forms}
Let $X$ be a connected complex manifold and
$h$ be a Kähler metric.
Let $\omega$ be the associated Kähler form.
Show that if $\dim X \geq  2$, the wedge
product with $\omega$
is injective on 1-forms.

Show that if $\dim X \geq  2$ and $\phi$ is a
differentiable
function with values in $\mathbb{R}^+$ such
that $\phi h$ is also
a Kähler metric, then $\phi$ is constant.
\end{exercise}

\begin{proof}
Take local symplectic coordinates, i.e. such
that
$\omega = \sum_idz^i \wedge d \bar{z}^i$.
Let $\alpha$ be any 1-form and suppose it is
written
in those coordinates as
$\alpha=\sum_ia_idz^i +
\sum_i\bar{a}_id\bar{z}^i$.
Then
$$
\begin{aligned}
\alpha\wedge\omega&=
\left(\sum_ia_idz^i\right)
\wedge
\left(\sum_i dz^i \wedge d\bar{z}^i\right)
+
\left(\sum_i \bar{a}_i dz^i\right)
\wedge
\left(\sum_i dz^i \wedge d\bar{z}^i\right)\\
&=
\sum_{i,j}a_i dz^i \wedge dz^j \wedge
d\bar{z}^j
+
\sum_{i,j}\bar{a}_i d\bar{z}^i \wedge dz^j
\wedge d\bar{z}^j.
\end{aligned}
$$

\noindent
We see that if $\alpha \wedge \omega=0$
then both its $1,2$ and $2,1$ parts are zero,
which foces $a_i$ and $\bar{a}_i$ to vanish,
i.e. $\alpha=0$.
\end{proof}
\begin{exercise}
\label{complex-geometry-exercise-holom-vecto-bundl-kahle}
\begin{reference}
\cite[Exercise 2, Chapter 3]{voi}
\end{reference}
Let $E$ be a holomorphic vector bundle of
rank $r$
over a complex manifold $X$.
Show that if $L$ is a holomorphic line bundle
on $X$,
then $\mathbb{P}(E^* \otimes L^*)=
\mathbb{P}(E^*)$
but the line bundle
$\mathcal{O}_{\mathbb{P}(E^*\otimes L^*)}(1)$
is isomorphic to
$\mathcal{O}_{\mathbb{P}(E^*)}(1) \otimes
\pi^*L$,
where $\pi:\mathbb{P}(E^*)\to X$ is the
structural morphism.
\end{exercise}

\begin{exercise}
\label{complex-geometry-exercise-kahler-iff-i-parallel}
$(M,I,g)$ is Kähler if and only if
$\nabla^{LC}I=0$.
\end{exercise}





\section{Hodge theory}
\label{complex-geometry-section-hodge-theory}

\noindent
On a compact hermitian manifold $(X,g)$ we
may define
$$
(\alpha,\beta):=\int_Xg_{\mathbb{C}}(\alpha,
\beta)*1,
$$
which is an hermitian product on
$\mathcal{A}_{\mathbb{C}}^*(X)$.

\begin{lemma}
\label{complex-geometry-lemma-adjoints}
\begin{reference}
\cite[Lemma 3.2.3]{huc}
\end{reference}
Let $X$ be a compact hermitian manifold.
Then $\partial^*$ and $\bar\partial^*$
are adjoints to $\partial$ and $\bar\partial$
with respect to
$(\cdot,\cdot)$.
\end{lemma}

\begin{proof}
We have
$$
\begin{aligned}
(\partial\alpha,\beta)
&=\int_X\left\langle{}\partial\alpha,
\beta\right\rangle{}_\mathbb{C}\text{Vol}
=\int_X \partial\alpha\wedge*\beta,\\
(\alpha,\partial^*\beta)
&=\int_X\left\langle{}\alpha,-
*\partial*\beta\right\rangle{}_{\mathbb{C}}\text{Vol}
=\int_X\alpha\wedge\partial*\beta.
\end{aligned}
$$

\noindent
Now we also have
$\partial(\alpha\wedge\beta)
=\partial\alpha\wedge\beta+
\alpha\wedge\partial\beta$,
and the integral of the first vanishes by
Stokes theorem.
\end{proof}

\begin{lemma}
\label{complex-geometry-lemma-harmonic-iff-closed}
Let $X$ be a compact hermitian manifold.
A form is $\bar\partial$-harmonic if and only
if
it is $\bar\partial$-closed and
$\bar\partial^*$ closed.
The same holds for $\partial$,
and in fact, on any compact oriented
Riemannian manifold $(M,g)$,
a similar result holds for the usual
$d$-Laplacian.
\end{lemma}

\begin{proof}
$$
\begin{aligned}
(\Delta \alpha,\alpha)&=((d
d^*+d^*d)\alpha,\alpha)\\
&=(d d^*\alpha,\alpha)+(d^*d\alpha,\alpha)\\
&=(d^*\alpha,d^*\alpha)+(d\alpha,d\alpha)\\
&=\|d^*\alpha\|^2+\|d\alpha\|^2.
\end{aligned}
$$
\end{proof}

\noindent
The following is the first of the symmetries
of the Hodge diamond:

\begin{lemma}
\label{complex-geometry-lemma-hodge-commutes-with-laplacian}
\begin{reference}
\cite[Lemma A.0.14]{huc}
\end{reference}
For any Riemannian manifold $(M,g)$,
not necessarily compact, $*\Delta=\Delta*$
and $*:\mathcal{H}^k(M,g) \cong
\mathcal{H}^{m-k}(M,g)$.
\end{lemma}

\noindent
In the hermitian case, we have for a compact
hermitian connected
manifold $(M,g)$, that the pairing
$$
\begin{aligned}
\mathcal{H}^{p,q}_{\bar\partial}
\times\mathcal{H}^{n-p,n-q}_{\bar\partial}
 &\longrightarrow \mathbb{C} \\
(\alpha,\beta) &\longmapsto
\int_X\alpha\wedge\beta
\end{aligned}
$$

\noindent
is non-degenerate and thus
$$
\mathcal{H}^{p,q}_{\bar\partial}(X,g)
\cong
\mathcal{H}^{n-p,n-q}_{\bar\partial}(X,g),
$$
which is called {\it Serre duality},
the second of the symmetries of the Hodge
diamond.
We shall later generalize this Serre duality
to other vector bundles.

There is a third symmetry of the Hodge
diamond,
given by the Lefschetz operator.
Since $[L,\Delta]=0$, one of the Kähler
identities,
we see that $L$ maps harmonic forms to
harmonic forms.
Further, since $L^{n-k}:\Lambda^k(V^\vee) \to
\Lambda^{2n-k}(V^\vee)$
is bijective we have that the induced map
$$
L^{n-k}:\mathcal{H}^{p,k-p}(X,g)
\to \mathcal{H}^{n+p-k,n-p}(X,g)
$$
is injective. It is also bijective since…
thus it is an isomorphism.

\medskip\noindent
Now we are almost ready for Hodge theorem.

\begin{proposition}
\label{complex-geometry-proposition-decom-ortho}
\begin{reference}
\cite[Proposition 3.2.2]{huc}
\end{reference}
Let $(X,g)$ be a compact hermitian manifold.
The degree, bidegree and Lefschetz
decomspositions
are orthogonal with respect to
$(\cdot,\cdot)$.
\end{proposition}

\noindent
The bidegree decomposition passes to
$\partial$ and $\bar\partial$-harmonic forms
on any compact hermitian manifold,
but only for Kähler manifolds does it pass to
$d$-harmonic forms.

\begin{proposition}
\label{complex-geometry-proposition-decom-harmo-forms}
Let $(X,g)$ be a hermitian manifold, not
necessarily compact.
\begin{enumerate}
\item $\mathcal{H}^{k}_{\bar\partial}(X,g)
=\bigoplus_{p+q=k}\mathcal{H}^{p,q}
_{\bar\partial}(X,g)$
and
$\mathcal{H}^{k}_\partial(X,g)
=\bigoplus_{p+q=k}\mathcal{H}^{p,q}
_\partial(X,g)$.

\item If $(X,g)$ is Kähler then both
decompositions coincide with
$\mathcal{H}^{k}(X,g)
=\bigoplus_{p+q=k}\mathcal{H}^{p,q}(X,g)$.
In particular,
$\mathcal{H}^k(X,g)_{\mathbb{C}}
=\mathcal{H}^k_{\bar\partial}(X,g)
=\mathcal{H}^k_{\partial}(X,g)$.
\end{enumerate}
\end{proposition}

\noindent
We finally arrive at the Hodge decomposition.
The first result concerns only compact
Riemannian manifolds,
and ``requires some hard, but by now
standard, analysis'':

\begin{theorem}[Hodge decomposition]
\label{complex-geometry-theorem-hodge-decomposition}
Let $(M,g)$ be a compact Riemannian manifold.
Then
$$
\mathcal{A}^k(M) =
d(\mathcal{A}^{k-1}(M))
\oplus
\mathcal{H}^k(M,g)
\oplus
d^* (\mathcal{A}^{k+1}(M))
$$
which is orthogonal with respect to
$(\cdot,\cdot)$.
Moreover, the space $\mathcal{H}^k(M)$ is
finite-dimensional.
\end{theorem}

\noindent
And now we also have
\begin{theorem}[Hodge decomposition]
\label{complex-geometry-theorem-hodge-decomposition-hermitian}
Let $(X,g)$ be a compact hermitian manifold.
Then
$$
\mathcal{A}^{p,q}(X)=
\partial\mathcal{A}^{p-1,q}(X)
\oplus
\mathcal{H}^{p,q}_\partial(X,g)
\oplus
\partial^*\mathcal{A}^{p+1,q}(X)
$$
and
$$
\mathcal{A}^{p,q}(X)=
\bar\partial\mathcal{A}^{p-1,q}(X)
\oplus
\mathcal{H}^{p,q}_{\bar\partial}(X,g)
\oplus
\bar\partial^*\mathcal{A}^{p+1,q}(X).
$$
If $(X,g)$ is Kähler, then
$\mathcal{H}^{p,q}_\partial(X,g)=\mathcal{H}
^{p,q}_{\bar\partial}(X,g)$.
\end{theorem}

\noindent
Upshot: if $X$ is Kähler we can pass from
usual
harmonic $(p,q)$-forms to
$\bar\partial$-harmonic $(p,q)$-forms,
which takes us to $H^{p,q}(X)$.
To pass from $H^k(X)$ to $\mathcal{H}^k(X)$
and from
$\mathcal{H}^{p,q}_{\partial}(X)$ to
$H^{p,q}(X)$
we have the following two lemmas
which in fact are corollaries from the
corresponding decompositions
(Theorems \ref{complex-geometry-theorem-hodge-decomposition}
and
\ref{complex-geometry-theorem-hodge-decomposition-hermitian},
respectively).

\begin{lemma}
\label{complex-geometry-lemma-de-rham-to-harmonic}
\begin{reference}
\cite[Corollary A.0.17]{huc}
\end{reference}
Let $(M,g)$ be a compact oriented Riemannian
manifold.
The natural map $\mathcal{H}^k(M) \to
H^k(M,\mathbb{R})$
that associates to any harmonic form its
cohomology class is
an isomorphism.
\end{lemma}

\begin{proof}[Misha's]
Step 1 (harmonic analysis, no proof).
$\Omega^k(M)=\bigoplus_{\lambda\operatorname{eigenvalue}\Delta}
\Omega_\lambda^k(M)$.

\noindent
Step 2. $[\Delta,d]=0 \implies $ $d$
preserved the eigenspaces (see Misha's
notes, or hodge-theory.tex).

This gives a subcomplex
$$
0 \to \Omega_0^0(M)
\to \Omega^1_\lambda(M)
\to …
$$

\noindent
Step 3. The cohomology of the subcomplex
above is trivial, $H^*_\lambda=0$. That is,
every closed form is exact since for a
$\lambda$-eigenform $\eta$ we have,
$\Delta \eta = \lambda \eta$ and then
$\eta = \frac{1}{\lambda}d d^*\eta$.
$\Omega_0^k = \Ker \Delta$.

\noindent
Step 4. $H^*(M) = \bigoplus_\lambda
H_\lambda ^*$.

\end{proof}

\begin{lemma}
\label{complex-geometry-lemma-harmo-pq-dolbe-pq}
Let $(X,g)$ be a compact hermitian manifold.
Then the canonical projection
$\mathcal{H}_{\bar\partial}^{p,q}(X,g)\to
H^{p,q}(X)$
is an isomorphism.
\end{lemma}

\medskip\noindent
Before we pass to Hodge theorem we state one
final key result

\begin{lemma}[ddbar Lemma]
\label{complex-geometry-lemma-ddbar}
Let $X$ be a compact Kähler manifold.
If $\alpha$ is a $d$-closed $(p,q)$-form,
it is equivalent that $\alpha$ is
\begin{enumerate}
\item $d$-exact,
\item $\partial$-exact,
\item $\bar\partial$-exact,
\item $\partial \bar\partial$-exact.
\end{enumerate}
\end{lemma}

\noindent
There is another version in terms of
$dd^c:=
I^{-1}dI=\frac{1}{2}(\partial-\bar\partial)$.

\begin{lemma}[ddc Lemma]
\label{complex-geometry-lemma-ddc}
\begin{reference}
\cite[Lemma 3.A.22]{huc}
\end{reference}
Let $X$ be a compact Kähler manifold.
Any $d$-closed and $d^c$-exact form form is
$dd^c$-exact.
\end{lemma}

\begin{proof}
\begin{enumerate}
\item $\alpha$ $d^c$-exact and $d$-closed 
implies beta is $\partial \bar\partial$-closed, 
where $\alpha=d^c\beta$.
\item $\alpha- i \bar\partial \beta$
is $\partial$-exact.
\item $\alpha-i \bar\partial\beta$
is $d$-closed.
\item By $\partial\bar\partial$-lemma,
$\alpha - i \bar\partial$ is
$\partial$-exact.
\item Thus $\alpha$ is $\bar\partial$-exact.
\item Thus alpha is 
$\partial\bar\partial$-exact,
i.e. $dd^c$-exact.
\end{enumerate}
\end{proof}

\begin{remark}
\label{complex-geometry-remark-ddbar-importance}
$d\bar\partial$ lemma makes
the different cohomologies 
(Aeppli, Bott-Chern, Dolbeault)
be isomorphic. Though, to Beatricce it's
not clear that it gives a Hodge decomposition
in a nonKähler manifold.
\end{remark}

\noindent
That's all the ingredients for a
(very rough) idea on why the Hodge theorem
holds.
Upshot:
$$
H^k(X,\mathbb{C})
= \mathcal{H}^k(X)
= \bigoplus_{p+q=k}\mathcal{H}^{p,q}(X)
=
\bigoplus_{p+q=k}\mathcal{H}^{p,q}
_{\bar\partial}(X)
=\bigoplus_{p+q=k}H^{p,q}(X).
$$
\begin{enumerate}
\item Start with singular complex cohomology,
\item pass to harmonic forms
via Lemma \ref{complex-geometry-lemma-de-rham-to-harmonic}
(this is where heavy real analysis enters;
what Huybrechts calls Hodge decomposition),
\item decompose harmonic forms in
$(p,q)$-types
via Proposition
\ref{complex-geometry-proposition-decom-harmo-forms}
(using Kählerness),
\item use that $\mathcal{H}^{p,q}\cong
\mathcal{H}^{p,q}_{\bar\partial}$
by the same Proposition
\ref{complex-geometry-proposition-decom-harmo-forms}
(using Kählerness),
\item go back to Dolbeault cohomology
via Lemma \ref{complex-geometry-lemma-harmo-pq-dolbe-pq}
(again Hodge decomposition, now for
$\bar\partial$).
\end{enumerate}

\begin{theorem}[Hodge]
\label{complex-geometry-theorem-hodge-theorem}
Let $(X,g)$ be a compact Kähler manifold.
Then there exists a decomposition
$$
H^k(X,\mathbb{C})
= \bigoplus_{p+q=k}H^{p,q}(X)
$$
which does not depend on the Kähler
structure.
Moreover, with respect to complex conjugation
on
$H^\bullet(X,\mathbb{C})=H^\bullet(X,
\mathbb{R})\otimes\mathbb{C}$
we have
$\overline{H^{p,q}(X)}=H^{q,p}(X)$,
and, by Serre duality on harmonic forms,
$H^{p,q}(X) \cong H^{n-p,n-q}(X)^\vee$.
\end{theorem}

\begin{exercise}
\label{complex-geometry-exercise-pic-of-pn-is-z}
\begin{reference}
\cite[Exercise 3.2.11]{huc}
\end{reference}
Show that $\text{Pic}(\mathbb{P}^n) \simeq
\mathbb{Z}$.
\end{exercise}

\noindent
The next step is to check
that most results we've developed so far for
complex vector spaces and complex manifolds
descend to cohomology in Kähler manifolds.

\medskip\noindent
{\bf Beatricce's explanation.}
Let's put this alternative route as well.
From the most important Kähler identity,
explained a while ago, we get:

\begin{lemma}[Corollary of the most
important Kähler identity]
\label{complex-geometry-lemma-hodge-kahler-important}
Any cohomology class can be represented
as the sum of Dolbeault classes,
giving the decomposition
$$
H^i
=
\bigoplus_{p+q=i}H^{p,q}_{\bar\partial},
\qquad
\overline{H^{p,q}}=H^{q,p}.
$$
\end{lemma}

\begin{proof}
We use harmonic representatives
from Hodge theorem:
$$
0=\Delta \alpha =
\sum \Delta \alpha_{p,q}
=
2 \sum \Delta_{\bar\partial}\alpha_{p,q},
$$
and $\Delta \bar{\alpha}=\overline{\Delta
\alpha}$.
\end{proof}

\noindent
It also follows that
\begin{enumerate}
\item Odd Betti numbers are even.
\item Even Betti numbers are $\geq 1$.
\end{enumerate}



\medskip\noindent
Let's put here the part on the cones until
I find a better place.

Recall that $\text{NS}(X)$ is the quotient
group
$\text{Pic}(X)/H^1(\mathcal{O}_X)\subset
H^2(X,\mathbb{Z})
\subset H^2(X,\mathbb{R})$
coming from the exponential exact sequence.
This has an intersection form.
In the case of a K3 we readily see that
$h^1=0$ so $\text{NS}=\text{Pic}$.
In fact, the intersection form coincides with
the usual cup product on 2-cocycles of
singular cohomology.

For a Kähler manifold we have
\begin{itemize}
\item The {\it Kähler cone}. Classes
in $H^{1,1}(X,\mathbb{R})$ that are
represented by a
Kähler form.

\item The {\it positive cone}. Classes
$\text{NS}(X)_\mathbb{R}
:=\text{NS}(X) \otimes_\mathbb{Z}\mathbb{R}$
with $\alpha^2>0$.

\item The {\it ample cone}. Classes
$\text{NS}(X)_\mathbb{R}$
with $\alpha^2>0$ and $\alpha.C>0$ for all
algebraic curves 0.
(See
\ref{algebraic-geometry-theorem-%
Nakai-Moishezon-criterion}.)
This is true, but hard to prove, see
\cite[p. 150]{huk}.
Another description is that the ample cone
is the set of classes in
$\text{NS}(X)_\mathbb{R}$
that are finite sums $\sum a_iL_i$
for $L_i \in \text{NS}(X)$ ample
and $a_i \in \mathbb{R}_{>0}$.

\item The {\it nef cone}. Classes (in
$\text{NS}(X)?)$
with $\alpha.C=\int_C\alpha\geq 0$ for any
algebraic curve.
There is also a notion of nef for analytic
cycles
which basically says that the form is
approximated
by Kähler forms.
In \cite[Theorem 4.7(ii)]{Demailly-Paun} we
have that
a $(1,1)$-class is nef (in the analytic
sense) if and only if
$\int_Y \alpha\wedge\omega^{p-1}\geq 0$ for
every
irreducible analytic set $Y \subset X$,
$p=\dim Y$
and $\bar{\omega}$ Kähler class.

\end{itemize}

\noindent
The upshot about Demailly-Pa\u un theorem
is that Kähler classes are numerically
positive
(recall that ``numerically'' means w.r.t.
intersection form)
on {\it analytic} cycles. This distinguishes
complex analytic geometry from algebraic
geometry.
This is why Demailly-Pa\u un say that
their theorem is a generalization of
Nakai-Moishezon:
because the latter is for {\it algebraic}
cycles,
indeed, found in Hartshorne's book.

The proof goes roughly as follows.
It's clear that $\mathcal{K} \subset
\mathcal{P}$
($\mathcal{P}$ is the analytically
numerically positive cone)
because, since the Kähler form
comes from the Riemannian metric
and the curves are complex,
which means that they are invariant under
the complex structure,
we have integrals
$\int_Y\omega$ for curves $Y \subset X$
but $\omega(v,Jv)>0$ for a basis $v,Jv$
of $T_pY$.
That is to say, there are no complex
Lagrangian curves.

But that's elementary.
The point of the theorem is proving that
a form in $\overline{\mathcal{K}}\cap
\mathcal{P}$,
i.e. in the closure of the Kähler cone
and which is also numerically analytically
positive,
is in fact Kähler.

So, the theorem says that the boundary of the
Kähler cone does
not intersect the numerically analytically
positive cone.

\begin{theorem}[Demailly-Pa\u un]
\label{complex-geometry-theorem-demailly-paun}
Let $X$ be a compact Kähler manifold.
The Kähler cone is a connected component
of the cone of 1,1-classes numerically
positive on
analytic cycles.
\end{theorem}

\begin{exercise}
\label{complex-geometry-exercise-demailly-paun}
Let $M$ be a K3 surface such that
$\text{rk}\text{NS}(M)<20$.
Using Demailly-Paun theorem,
prove that there exists a non-zero vector
$\eta \in H^{1,1}(M)$ which satisfies
$\eta^2=0$
and belongs to the boundary of the Kähler
cone.
\end{exercise}

\begin{proof}
Let $\alpha \in
\text{NS}(X)^{\perp}\setminus\{0\}$,
which is nonempty since $\dim
\text{NS}(X)<20=h^{1,1}$.
By the Hodge index theorem
\ref{complex-geometry-theorem-hodge-index}
the intersection form on
$H^{1,1}(X,\mathbb{R})$
has index $(1,19)$.
If we assume that $X$ is projective
then we know that the restriction of
the Fubini-Study form $\omega$ is integral.
Since $\omega$ is positive (by being a Kähler
form),
we see by the index of the intersection form
that $\alpha$ must be negative.

Define $\eta_t:=t\alpha+(1-t)\omega$ for
$t\in[0,1]$.
Then $\eta_t^2=t^2\alpha^2+(1-t)^2\omega^2$.
Since at $t=0$ we have a positive number
and at  $t=1$ a negative number,
we know that there exists  $t_0$
such that $\eta_{t_0}^2=0$
by continuity of the intersection form.

Now we need to show that $\eta_{t_0}$ is
in the boundary of the Kähler cone.
For this it's enough to show that for every
$\varepsilon>0$
$\eta_{t_0+\varepsilon}$ is in the Kähler,
which
by Demailly-P\u aun
means that we have to show
that it is positive on analytic cycles.

By our choice of $\omega$ and $\alpha$
we see that
$$
\int_Y \eta_{t_0+\varepsilon}
=\int_Y
(t_0+\varepsilon)\alpha+(1-(t_0+\varepsilon))
\omega
$$
and the integral of $\alpha$ vanishes
since $\alpha$ is in the orthogonal
complement
of $\text{NS}(X)$, which contains
all analytic curves by $X$ being projective
(Chow theorem).
\end{proof}






\section{Complex tori and abelian varieties}
\label{complex-geometry-section-complex-tori-abelian}

\begin{definition}
\label{complex-geometry-definition-complex-torus}
A {\it complex torus} of dimension $g$ is a
quotient
$
T=\mathbb{C}^g/\Lambda,
$
where $\Lambda\subset\mathbb{C}^g$ is a
lattice of real rank $2g$.
\end{definition}

\noindent
The addition on $\mathbb{C}^g$ descends to
$T$, so a complex torus is in particular a
complex Lie group. The lattice condition
also implies that $T$ is compact.

\begin{example}
\label{complex-geometry-example-elliptic-curve-torus}
For $\tau\in\mathbb{C}$ with positive
imaginary part,
$
E_\tau=
\mathbb{C}/(\mathbb{Z}+\tau\mathbb{Z})
$
is a one-dimensional complex torus.
It is an elliptic curve.

This is the same kind of compact curve used
in the proof of the one-dimensional
$\bar\partial$-Poincaré lemma,
Lemma
\ref{complex-geometry-lemma-dbar-poincare-one-variable}:
there we embedded a disk into a sufficiently
large elliptic curve and solved the global
problem by Fourier series.
\end{example}

\begin{definition}
\label{complex-geometry-definition-abelian-variety}
An {\it abelian variety} over
$\mathbb{C}$ is a projective complex torus.
An abelian variety of complex dimension $2$
is called an {\it abelian surface}.
\end{definition}

\begin{remark}
\label{complex-geometry-remark-complex-torus-not-abelian}
Not every complex torus is projective.
Thus a quotient
$
\mathbb{C}^g/\Lambda
$
is automatically a compact complex torus,
but it is an abelian variety only when it is
projective.
\end{remark}

\begin{example}
\label{complex-geometry-example-product-elliptic-abelian}
If $E_1,\ldots,E_g$ are elliptic curves,
then
$
E_1\times\cdots\times E_g
$
is projective and hence is an abelian
variety. In particular,
$
E\times E
$
is an abelian surface.
\end{example}

\noindent
Abelian varieties are the
higher-dimensional analogues of elliptic
curves that will appear most often below.
For background on complex tori and their
relation with line bundles and projectivity,
see \cite{huc}.

\section{Lefschetz theorems}
\label{complex-geometry-section-lefschetz-theorems}

\begin{theorem}[Lefschetz theorem on
(1,1)-classes]
\label{complex-geometry-theorem-11-classes}
Let $X$ be a compact Kähler manifold.
An integral class
$\alpha\in H^2(X,\mathbb{Z})$
is the first Chern class of a
holomorphic line bundle if and only if
its image $\alpha_{\mathbb{C}}$ in
$H^2(X,\mathbb{C})$ has type $(1,1)$.
Thus the map
$$
c_1:\operatorname{Pic}(X)
\longrightarrow H^{1,1}(X,\mathbb{Z})
$$
is surjective, where
$$
H^{1,1}(X,\mathbb{Z})
=\{\alpha\in H^2(X,\mathbb{Z}):
\alpha_{\mathbb{C}}\in H^{1,1}(X)\}.
$$
This definition includes torsion classes,
whose images in complex cohomology vanish.
\end{theorem}

\begin{proof}
Consider the exponential sequence of
sheaves
$$
0\longrightarrow\mathbb{Z}
\longrightarrow\mathcal{O}_X
\xrightarrow{\exp(2\pi i\,\cdot)}
\mathcal{O}_X^*
\longrightarrow 1.
$$
Here $\mathcal{O}_X^*$ is the sheaf of
nowhere vanishing holomorphic functions.
The kernel of the exponential map is
the locally constant integer functions.
Surjectivity is local: a nowhere
vanishing holomorphic function admits
a holomorphic logarithm near each point.

The long exact cohomology sequence
contains
$$
H^1(X,\mathcal{O}_X^*)
\xrightarrow{c_1}H^2(X,\mathbb{Z})
\xrightarrow{j}H^2(X,\mathcal{O}_X).
$$
The first group is
$\operatorname{Pic}(X)$: its cocycles
are transition functions of holomorphic
line bundles. The connecting map is
the integral first Chern class.
Concretely, lift transition functions
$g_{ij}$ locally to holomorphic
functions $f_{ij}$ with
$g_{ij}=\exp(2\pi i f_{ij})$.
Their coboundary
$f_{ij}+f_{jk}-f_{ik}$ is integer valued
on triple overlaps and represents $c_1$.
Exactness therefore gives
$$
\operatorname{im}c_1=\ker j.
$$

We now identify $j$. Dolbeault gives
$$
H^2(X,\mathcal{O}_X)
\cong H^{0,2}_{\bar\partial}(X).
$$
The inclusion $\mathbb{Z}\to\mathcal{O}_X$
factors through the constant sheaf
$\mathbb{C}$. On the de Rham and
Dolbeault resolutions, the map induced
by $\mathbb{C}\to\mathcal{O}_X$ is
projection to the $(0,k)$ component.
Indeed, for a complex $k$-form $\eta$,
$$
(d\eta)^{0,k+1}
=\bar\partial(\eta^{0,k}).
$$
Thus this projection commutes with the
differentials and induces the required
map on cohomology.

Compactness and the Kähler hypothesis
give the Hodge decomposition
$$
H^2(X,\mathbb{C})
=H^{2,0}(X)\oplus H^{1,1}(X)
\oplus H^{0,2}(X).
$$
The harmonic representative decomposes
into harmonic forms of these types.
Consequently $j(\alpha)$ is precisely
the $(0,2)$ component
$\alpha_{\mathbb{C}}^{0,2}$.

Since $\alpha$ is integral, its complex
image comes from real cohomology.
Complex conjugation therefore gives
$$
\alpha_{\mathbb{C}}^{2,0}
=\overline{\alpha_{\mathbb{C}}^{0,2}}.
$$
Hence $j(\alpha)=0$ if and only if
both the $(0,2)$ and $(2,0)$ components
vanish, which is equivalent to
$\alpha_{\mathbb{C}}\in H^{1,1}(X)$.
Together with exactness, this proves
the theorem.
\end{proof}

\noindent
This implies that
$\text{NS}(X):=\text{Pic}(X)/\ker c_1$
is isomorphic to $H^{1,1}(X,\mathbb{Z})$.
The ``realification'' of the integer classes,
$H^{1,1}(X,\mathbb{Z}) \otimes \mathbb{R}
=\text{NS}(X) \otimes \mathbb{R}$
is a vector subspace
$H^{1,1}(X) \cap H^2(X,\mathbb{R})$.

Following \cite{huc}, we now have:

\begin{definition}
\label{complex-geometry-definition-picard-number}
Let $X$ be a compact complex manifold.
Then the rank of the image of
$\text{Pic}(X) \to H^2(X,\mathbb{R})$
is called the {\it Picard number} of
$X$ and is denoted by $\rho(X)$.
\end{definition}

\noindent
This means that if $X$ is Kähler,
by Lefschetz' theorem on $(1,1)$-classes,
we must have
$\rho(X) = \text{rk} H^{1,1}(X,\mathbb{Z})
= \text{rk} \text{NS}(X)$.

\begin{definition}
\label{complex-geometry-definition-jacobian}
Let $X$ be a compact complex manifold.
Its {\it Jacobian} is the kernel of the map
$\text{Pic}(X) \to H^2(X,\mathbb{Z})$
and is denoted $\text{Pic}^0(X)$.
\end{definition}

\noindent
Using the exponential sequence
the Jacobian can be (for compact $X$)
described
as the quotient
$$
\text{Pic}^0(X)=
\frac{H^1(X,\mathcal{O}_X)}{
H^1(X,\mathbb{Z})}.
$$
Tracing back the exponential sequence
(use a whole page to write the long
exact sequence and go back step by step),
this isomorphism boils down to
$H^0(\mathcal{O}_X) \to H^0(\mathcal{O}_X^*)$
being surjective.
For compact $X$ this is true:
non zero global holomorphic
functions are non zero constants,
each of which was a well-defined
logarithm, which is the inverse map.
For example, the logarithm of $-1$
is $\text{log}(-1)=\text{log} e^{i\pi}
=i\pi$.

\medskip\noindent
For a compact Riemann surface, there is a
concrete map into the Jacobian that will be
useful later.

\begin{definition}
\label{complex-geometry-definition-abel-jacobi}
Let $C$ be a compact Riemann surface and
choose a point $p_0\in C$. The
{\it Abel--Jacobi map} is
$$
\begin{aligned}
a:C&\longrightarrow\text{Pic}^0(C),\\
p&\longmapsto
\mathcal{O}_C(p-p_0).
\end{aligned}
$$
More generally, for $d\geq 1$ there is a
map
$$
\begin{aligned}
a_d:\text{Sym}^d(C)
&\longrightarrow\text{Pic}^d(C),\\
p_1+\cdots+p_d
&\longmapsto
\mathcal{O}_C(p_1+\cdots+p_d).
\end{aligned}
$$
\end{definition}

\noindent
The symmetric power appears because a divisor
does not remember an ordering of its points.
Thus a point of $\text{Sym}^d(C)$ is an
effective divisor of degree $d$. The
Abel--Jacobi map remembers only the
associated line bundle.

\begin{remark}
\label{complex-geometry-remark-abel-jacobi-fibres}
Fix a line bundle $L$ of degree $d$.
The fibre of $a_d$ over $L$ is the complete
linear system
$
|L|=\mathbb{P}(H^0(C,L)).
$
Indeed, its points are precisely the
effective divisors $D$ such that
$
\mathcal{O}_C(D)\cong L.
$
See \cite{Donaldson} for the classical
Riemann-surface viewpoint.
\end{remark}

\begin{example}
\label{complex-geometry-example-abel-jacobi-elliptic}
Let $E$ be an elliptic curve, with origin
$0\in E$. Under the standard identification
$
\text{Pic}^0(E)\cong E,
$
the degree-$3$ Abel--Jacobi map can be
translated from $\text{Pic}^3(E)$ to
$\text{Pic}^0(E)$ and written as
$$
\begin{aligned}
\text{Sym}^3(E)&\longrightarrow E,\\
[p,q,r]&\longmapsto p+q+r.
\end{aligned}
$$
Thus the map used in the crystallographic
description of $\mathbb{CP}^2$ is an
Abel--Jacobi map.
\end{example}

\begin{lemma}
\label{complex-geometry-lemma-jacobian-is-torus}
If $X$ is a compact Kähler manifold,
then $\text{Pic}^0(X)$ is in a natural
way a complex torus of dimension $b_1(X)$.
\end{lemma}

\begin{proof}
We have to check that $H^1(X,\mathbb{Z})$
is a lattice in $H^1(X,\mathcal{O}_X)$.
i.e. a subgroup isomorphic
to $\mathbb{Z}^r$
and such that 
$H^1(X,\mathbb{Z}) 
\otimes_\mathbb{Z}
\mathbb{R}
\cong H^1(X,\mathcal{O}_X)$.
By Dolbeault
$H^1(\mathcal{O}_X)
=H^{0,1}(X)$.
Then prove that $H^{0,1}(X)
\cong H^1(X,\mathbb{R})
\cong H^1(X,\mathbb{Z}) \otimes \mathbb{R}$.
\end{proof}

\noindent
Doing the same for other preimages
(not just the vector zero, which 
gives the kernel)
of the map
$H^1(\mathcal{O}_X^*) \to H^2(\mathbb{Z})$
actually shows that $\text{Pic}(X)$
is a torus fibration
(for $X$ with torsion-free $H^1(\mathbb{Z})$,
i.e. $H^1(\mathbb{Z})$ is a free-group,
i.e. a lattice)
!
$$
\xymatrix{
\text{Pic}^\alpha(X)\ar[r]\ar[d]
&\text{Pic}(X)\ar[d]\\
\{\alpha\}\ar[r]
&H^2(X,\mathbb{Z})
}
$$

\begin{definition}
\label{complex-geometry-definition-albanese-torus}
Albanese torus…
\end{definition}

\begin{theorem}[Hodge index]
\label{complex-geometry-theorem-hodge-index}
The index of the intersection pairing on
$H^{1,1}(X,\mathbb{R})$
is $(1,\rho(X)-1)$.
\end{theorem}

\begin{theorem}[Hard Lefschetz Theorem]
\label{complex-geometry-theorem-hard-lefschetz}
\begin{reference}
\cite[p. 122]{gri}
\end{reference}
The map
$$
L^k:H^{n-k}(M) \to H^{n+k}(M)
$$
is an isomorphism.
If we define the {\it primitive} cohomology
as
$$
P^{n-k}(M) = Ker L^{k+1}
$$
we have
$$
H^m(M)=\bigoplus_k L^kP^{m-2k}(M),
$$
called the Lefschetz decomposition.
\end{theorem}

\section{Vector bundles, connections,
curvature and Chern classes}
\label{complex-geometry-section-vector-bundles}

\begin{lemma}[Vector bundle construction from
fibers]
\label{complex-geometry-lemma-vector-bundle-from-fibers}
\begin{reference}
\cite[Lemma 3.4]{lec}
\end{reference}
Suppose $M$ is a complex manifold,
and for each $p \in M$ we are given a
$k$-dimensional vector space $E_p$.
Let $E=\coprod_{p \in M}E_p$ and let $\pi:E
\to M$ be the obvious projection.
Suppose further that we are given
\begin{enumerate}
\item an indexed open cover $U_i$ of $M$
\item for each $i$, a bijection $
\psi_i:\pi^{-1}(U_i) \to U_i
\times\mathbb{C}^k$,
\item for each $i,j$ with $U_i\cap U_j \neq
\emptyset$
a holomorphic map $g_{ij}:U_i \cap U_J \to
\text{GL}(k,\mathbb{C})$
such that $\psi_i\circ
\psi_j^{-1}(p,v)=(p,g_{ij}(p)v$.
\end{enumerate}
\noindent
Then $\pi:E \to M$ has a unique structure as
a holomorphic vector bundle
with $(U_i,\psi_i)$ as holomorphic local
trivializations,
and with $g_{ij}$ as transition functions.
\end{lemma}

\begin{proposition}[Vector bundle
construction from cocycle]
\label{complex-geometry-proposition-vecto-bundl-cocyc}
Let $M$ be a complex manifold and let $U_i$
be an open cover of $M$.
Suppose that for each $i,j$ such that $U_i
\cap U_j \neq \emptyset$
we are given a holomorphic map $g_{ij}:U_i
\cap U_j \to \text{GL}(k,\mathbb{C})$
such that
$$
g_{ij}(p)g_{jk}(p)g_{ki}(p)=\text{Id}
$$
for any $p \in U_i\cap U_j\cap U_k$.

Then there is a rank-$k$ holomorphic vector
bundle $E \to M$
that admits a holomorphic trivialization over
each set $U_i$
with transition functions given by $g_{ij}$.
\end{proposition}

\begin{proposition}[Holomorphic Vector Bundle
Isomorphism Criterion]
\label{complex-geometry-proposition-isomo-vecto-bundl}
\begin{reference}
\cite[Proposition 3.7]{lec}
\end{reference}
Two vector bundles $E$ and $F$
with transition functions $g_{ij}$ and
$h_{ij}$
are isomorphic if and only if
$g_{ij}=A^{-1}_ih_{ij}A_j$
for some $A_i:U_i \to
\text{GL}_n(\mathbb{C})$.
\end{proposition}

\begin{proof}
The idea is to define a map in local
trivializations
and confirm that this glues to a global
bundle map.
This means: invent
$\varphi_i:E(U_i):=\Gamma_E(U_i) \to F(U_i)$,
i.e. a local map of sections
(yes, a map of sections, since a bundle map
is constructed as map of the total spaces,
which are points and vectors,
but we need to preserve the projection,
so all we care is about the vectors).
Then check that on $U_i\cap U_j:= U_{ij}$
we have $\varphi_i=\varphi_j$.
This literally creates a bundle map.

Try to do it for line bundles,
so $A_i$ are just functions.
Take a local frame $e_i$ for $E$
and a local frame $f_j$ for $F$.
Change coordinates and apply the map,
apply the map and change coordinates
- this should give both sides of the cocycle
condition
of the proposition.
\end{proof}

\begin{exercise}[Tensor product bundle]
\label{complex-geometry-exercise-tensor-product-bundle}
\begin{reference}
\cite[Exercise 2.2.1]{huc}
\end{reference}
Let $E$ and $F$ be vector bundles determined
by the cocycles
$g_{ij}: U_i \cap U_j \to
\text{GL}(r,\mathbb{C})$
and $\tilde{g}_{ij}:U_i \cap U_j \to
\text{GL}(\tilde{r},\mathbb{C})$.
Verify the cocycle description for
\begin{itemize}
\item The direct sum $E \oplus F$,
\item the tensor product $E \otimes F$ and
\item the dual bundle $E^\vee$.
\end{itemize}
\end{exercise}

\begin{proof}
Let's do the tensor product.
I think that, slightly more generally than in
the
formulation of the question, if $U_i$
is a trivializing cover for $E$ and
$\tilde{U}_j$ is a trivializing cover for
$F$,
we should consider the cover $U_i\cap
\tilde{U}_j$.
Let us denote such a cover simply by $U_i$.

Define a cocycle by
$$
\begin{aligned}
h_{ij}: U_i\cap U_j &\longrightarrow
\text{GL}(r\tilde{r},\mathbb{C}) \\
p &\longmapsto g_{ij}(p)\otimes
\tilde{g}_{ij}(p).
\end{aligned}
$$

\noindent
The key observation is that while formally
the image of $h_{ij}$
lies in $\text{GL}(\mathbb{C}^r \otimes
\mathbb{C}^{\tilde{r}})$,
we can regard such a space as
$\text{GL}(r\tilde{r},\mathbb{C})$
by defining, for $A \in
\text{GL}(r,\mathbb{C})$
and $A' \in \text{GL}(\tilde{r},\mathbb{C})$,
the map $A(v^ie_i) \otimes
\tilde{A}(\tilde{v}^j\tilde{e}_j)
=v^i\tilde{v}^j Ae_i \otimes
\tilde{A}\tilde{e}_j$
where $e_i$ and $\tilde{e}_j$ are bases
of $\mathbb{C}^r$ and
$\mathbb{C}^{\tilde{r}}$.
This is indeed a linear transformation of the
$r\tilde{r}$-dimensional $\mathbb{C}$-vector
space
generated by the basis $e_i \otimes
\tilde{e}_j$.

Thus the cocycle condition follows simply by
the cocycles of $E$ and $F$, and
commutativity of the
coordinate functions.
\end{proof}

\begin{example}[Determinant bundle]
\label{complex-geometry-example-determinant-bundle}
Let $\pi:E \to M$ be a vector bundle of rank
$r$.
The {\it determinant bundle} of $E$ is given
by the following data.
A space $\tilde{E}:=\coprod_{p \in
M}\Lambda^rE_p$
with the obvious projection map
and topology induced from such map.
Since we have a cover $U_i$ of $M$
along with local trivializations
$U_i\overset{\varphi_i}{\simeq}
U_i\times\mathbb{C}^r$
we may just postcompose $\varphi_i$ with the
functor $\Lambda^r$
on the second entry, giving trivializations
$$
\xymatrix{
\pi(U_i)\ar[r]^-{\varphi_i}_-{\simeq}
& U_i \times
\mathbb{C}^r\ar[r]^-{\text{Id}\times
\Lambda^r}
& U_i \times \Lambda^r\mathbb{C}^r=U_i \times
\mathbb{C}
}
$$
We can also apply $\Lambda^r$ to the
transition functions
$g_{ij}$ on the intersections
$U_{ij}:=U_i\cap U_j$
$$
\xymatrix{
U_{ij}\ar[r]^-{\varphi_i}&U_{ij} \times
\mathbb{C}^r
\ar[d]^{g_{ij}\in\text{GL}_r(\mathbb{C})}\\
U_{ij}\ar[r]^-{\varphi_j}&U_{ij}
\times\mathbb{C}^r
}
$$
giving maps in $\text{GL}_1(\mathbb{C})$.
Recall that $\Lambda^r$ is functorial on
linear maps
$g_{ij}\in \text{GL}_n(\mathbb{C})$ because
it is given by the determinant of $g$,
and the determinant is multiplicative.
We thus obtain the cocycle condition on the
maps $g_{ij}$
after applying $\Lambda^r$.
\end{example}

\begin{definition}
\label{complex-geometry-definition-connection-vector-bundle}
Let $E$ be a vector bundle.
A {\it connection} on $E$ is an
$\mathbb{R}$-linear map
$\nabla:H^0(E)\to H^0(E) \otimes \Omega^1(M)$
satisfying the identity
$\nabla(fs)=dfs + f\nabla s$.
\end{definition}

\noindent
For a complex vector bundle, we use a
complex linear connection and the
Leibniz rule for smooth complex valued
functions.
Choose a smooth local frame
$e_1,\ldots,e_r$.
The {\it connection matrix}
$\theta=(\theta_{ij})$ in this frame is
the matrix of complex valued one-forms
determined by
$$
\nabla e_j=\sum_i\theta_{ij}\otimes e_i.
$$
For $s=\sum_j f_j e_j$, the Leibniz rule
then gives
$$
\nabla s
=\sum_i\left(df_i+\sum_j\theta_{ij}f_j\right)
\otimes e_i.
$$
Thus, in this frame, $\nabla=d+\theta$.

\noindent
A local section $s:U\to E$ of a
holomorphic vector bundle is
{\it holomorphic} if it is a holomorphic
map. Equivalently, in a holomorphic
local frame $e_1,\ldots,e_r$, its
coefficients in $s=\sum_j f_j e_j$ are
holomorphic functions.
The bundle Dolbeault operator is given
in such a frame by
$
\bar\partial s
=\sum_j\bar\partial f_j\otimes e_j.
$
This is independent of the holomorphic
frame because the transition functions
are holomorphic.

\begin{lemma}
\label{complex-geometry-lemma-connection-holomorphic-sections}
Let $E\to M$ be a holomorphic vector
bundle and let $\nabla$ be a complex
linear connection on its underlying
smooth bundle. Put
$\nabla^{0,1}=\pi^{0,1}\nabla$,
where $\pi^{0,1}$ projects one-forms
onto their $(0,1)$ part.
Then $\nabla^{0,1}=\bar\partial$ if and
only if $\nabla^{0,1}s=0$ for every
local holomorphic section $s$.
\end{lemma}

\begin{proof}
If $\nabla^{0,1}=\bar\partial$ and
$s=\sum_j f_j e_j$ is holomorphic,
then each $f_j$ is holomorphic, so
$\bar\partial f_j=0$.
Thus $\nabla^{0,1}s=\bar\partial s=0$.

Conversely, assume that $\nabla^{0,1}$
annihilates every local holomorphic
section. Choose a holomorphic local
frame $e_1,\ldots,e_r$.
Each $e_j$ is a holomorphic section,
so $\nabla^{0,1}e_j=0$.
Every smooth section on this open set
has the form $s=\sum_j f_j e_j$ with
smooth coefficients $f_j$.
Taking the $(0,1)$ part of the Leibniz
rule gives
$$
\begin{aligned}
\nabla^{0,1}s
&=\sum_j\bigl(
\bar\partial f_j\otimes e_j
+f_j\nabla^{0,1}e_j\bigr)\\
&=\sum_j\bar\partial f_j\otimes e_j\\
&=\bar\partial s.
\end{aligned}
$$
Such frames cover $M$, proving the
equality of the two operators.
\end{proof}

\noindent
The word \emph{local} matters: global
holomorphic sections need not span
the fibres, and may all be zero.

\begin{remark}
\label{complex-geometry-remark-connection-matrix-type}
If $\nabla^{0,1}=\bar\partial$ and
$e_1,\ldots,e_r$ is a holomorphic local
frame, then $\theta$ is a matrix of
$(1,0)$-forms. Indeed, the lemma gives
$\nabla^{0,1}e_j=0$, hence
$$
0=\sum_i\theta_{ij}^{0,1}\otimes e_i.
$$
Since the $e_i$ form a basis in each
fibre, every $\theta_{ij}^{0,1}$ is zero.
Thus $\theta_{ij}=\theta_{ij}^{1,0}$.
This assertion requires a holomorphic
frame; it need not hold in an arbitrary
smooth frame.
\end{remark}

\noindent
Perhaps this notion defined on a seminar
talk by Beatrice Brienza is the same:
let $(M,I,g)$ be a hermitian manifold. A
linear connection $\nabla$ on $TM$ is
{\it hermitian} if
$\nabla g=0$ and $\nabla I =0$.

\begin{definition}
\label{complex-geometry-definition-hermitian-connection}
\begin{reference}
\cite[Definition 4.2.9]{huc}
\end{reference}
Let $(E,h)$ be a hermitian vector bundle.
A connection $\nabla$ on $E$ is called
{\it hermitian} if
$d(h(s,s'))=h(\nabla s,s')+h(s,\nabla s')$.
\end{definition}

\begin{definition}
\label{complex-geometry-definition-holomorphic-connection}
A connection $\nabla$ on a vector bundle $E$
is said to be {\it compatible with the
complex structure $I$} if
$\pi^{0,1}\nabla s=\bar\partial s$
where $\pi^{0,1}:\mathcal{A}^1(M)
\to \mathcal{A}^{0,1}(M)$ is the projection.
\end{definition}

\begin{proposition}
\label{complex-geometry-proposition-chern-connection}
On a holomorphic vector bundle
equipped with a hermitian metric there exists
a unique hermitian connection
compatible with the complex structure
called the {\it Chern connection}.
\end{proposition}

\begin{definition}
\label{complex-geometry-definition-curvature-connection}
Let $\nabla$ be a connection on a smooth
complex vector bundle $E\to M$.
Extend it to $E$-valued differential
forms by the rule
$$
\nabla(\alpha\otimes s)
=d\alpha\otimes s
+(-1)^k\alpha\wedge\nabla s,
$$
where $\alpha$ is a $k$-form.
The {\it curvature} of $\nabla$ is
$$
R^\nabla=\nabla^2.
$$
On sections this is linear over smooth
functions, so it is an
$\operatorname{End}(E)$-valued two-form.
In a local frame with connection
matrix $\theta$, the
{\it curvature form} is the matrix
$\Theta=(\Theta_{ij})$ of two-forms
defined by
$$
R^\nabla e_j=\sum_i\Theta_{ij}\otimes e_i.
$$
Applying the extended Leibniz rule gives
$$
\begin{aligned}
\Theta&=d\theta+\theta\wedge\theta,\\
(\theta\wedge\theta)_{ij}
&=\sum_k\theta_{ik}\wedge\theta_{kj}.
\end{aligned}
$$
\end{definition}

\begin{lemma}
\label{complex-geometry-lemma-chern-curvature-type}
For the Chern connection of a
holomorphic hermitian vector bundle,
the entries of $\Theta$ have type
$(1,1)$.
In a holomorphic local frame,
$$
\Theta=\bar\partial\theta
=\bar\partial(h^{-1}\partial h),
$$
where $h$ denotes the matrix
$(h(e_i,e_j))$ of the hermitian metric
in that frame, with the metric
conjugate linear in its first argument.
Compatibility with the holomorphic
structure alone does not suffice.
\end{lemma}

\begin{proof}
Choose a holomorphic local frame.
We already know that $\theta$ has type
$(1,0)$. Put $h_{ij}=h(e_i,e_j)$.
Metric compatibility says
$$
dh_{ij}
=h(\nabla e_i,e_j)+h(e_i,\nabla e_j).
$$
Substitute
$\nabla e_j=\sum_k\theta_{kj}\otimes e_k$.
Conjugate linearity in the first
argument gives
$$
dh_{ij}
=\sum_k\overline{\theta_{ki}}h_{kj}
+\sum_k h_{ik}\theta_{kj}.
$$
In matrix notation this is
$$
dh=\overline{\theta}^{\,t}h+h\theta.
$$
Since $\theta$ has type $(1,0)$,
$\overline{\theta}^{\,t}h$ has type
$(0,1)$, while $h\theta$ has type
$(1,0)$. Thus the $(1,0)$ part is
$$
\partial h=h\theta.
$$
The positive definite matrix $h$ is
invertible. Multiplying on the left
by $h^{-1}$ therefore gives
$$
\theta=h^{-1}\partial h.
$$
Differentiating $h^{-1}h=I$ gives
$$
\partial(h^{-1})
=-h^{-1}(\partial h)h^{-1}.
$$
Since $\partial^2h=0$, it follows that
$$
\begin{aligned}
\partial\theta
&=-h^{-1}(\partial h)h^{-1}
  \wedge\partial h\\
&=-\theta\wedge\theta.
\end{aligned}
$$
Consequently,
$$
\Theta
=\partial\theta+\bar\partial\theta
+\theta\wedge\theta
=\bar\partial\theta
=\bar\partial(h^{-1}\partial h),
$$
which has type $(1,1)$.
A change of frame conjugates $\Theta$
by a matrix of functions, so its type
does not depend on the frame.
\end{proof}

\begin{lemma}
\label{complex-geometry-lemma-curvature-trace-closed}
The trace $\operatorname{tr}\Theta$ of
the curvature form is a global closed
two-form. In particular, for a line
bundle, $\Theta$ itself is a global
closed two-form.
\end{lemma}

\begin{proof}
Under a change of frame by $g$, the
curvature matrix becomes
$g^{-1}\Theta g$, so its trace is
unchanged.
Moreover,
$$
\operatorname{tr}(\theta\wedge\theta)
=\sum_{i,j}\theta_{ij}\wedge\theta_{ji}
=0,
$$
because terms with reversed indices
cancel and diagonal terms vanish.
Thus, in each local frame,
$$
\operatorname{tr}\Theta
=d(\operatorname{tr}\theta),
$$
and $d^2=0$ proves closedness.
For a line bundle, the matrix has one
entry, so $\operatorname{tr}\Theta=\Theta$.
In this case $\theta\wedge\theta=0$,
and more explicitly,
$$
\Theta=d\theta,\qquad d\Theta=0.
$$
The local forms $\theta$ need not glue
to a global one-form, so this does not
say that $\Theta$ is globally exact.
\end{proof}

\noindent
In higher rank, $\Theta$ itself need
not be closed entry by entry.
Also, type $(1,1)$ alone does not imply
closedness.

\begin{definition}
\label{complex-geometry-definition-first-chern-form}
Let $(E,h)$ be a holomorphic hermitian
vector bundle, and let $\Theta$ be the
curvature matrix of its Chern connection.
Its {\it first Chern form} is
$$
c_1(E,h)=\frac{i}{2\pi}
\operatorname{tr}\Theta.
$$
Its {\it real first Chern class} is
the de Rham cohomology class
$$
c_1(E)_{\mathbb{R}}
=\left[\frac{i}{2\pi}
\operatorname{tr}\Theta\right]
\in H^2_{\mathrm{dR}}(M,\mathbb{R}).
$$
For a line bundle $L$, this becomes
$$
c_1(L)_{\mathbb{R}}
=\left[\frac{i}{2\pi}\Theta\right].
$$
This is the real image of the integral
first Chern class; differential forms
do not detect its torsion part.
\end{definition}

\begin{lemma}[Integrality for line bundles]
\label{complex-geometry-lemma-first-chern-integral}
For a holomorphic hermitian line bundle
$L$, the class
$$
\left[\frac{i}{2\pi}\Theta\right]
$$
is the image of an integral class
$c_1(L)\in H^2(M,\mathbb{Z})$.
This integral class is the connecting
class of $L$ in the smooth exponential
sequence.
\end{lemma}

\begin{proof}
Let $\mathcal{C}^{\infty}_{\mathbb{C}}$
be the sheaf of smooth complex valued
functions, and let
$\mathcal{C}^{\infty,*}_{\mathbb{C}}$
be its subsheaf of nowhere vanishing
functions. The smooth exponential
sequence is
$$
0\longrightarrow\mathbb{Z}
\longrightarrow
\mathcal{C}^{\infty}_{\mathbb{C}}
\xrightarrow{\exp(2\pi i\,\cdot)}
\mathcal{C}^{\infty,*}_{\mathbb{C}}
\longrightarrow 1.
$$
It is exact because smooth logarithms
exist locally. Partitions of unity
imply that the higher cohomology of
$\mathcal{C}^{\infty}_{\mathbb{C}}$
vanishes. Hence the connecting map is
an isomorphism
$$
H^1(M,\mathcal{C}^{\infty,*}_{\mathbb{C}})
\xrightarrow{\delta}
H^2(M,\mathbb{Z}).
$$
Smooth line bundles are classified by
the group on the left, using their
transition functions. Define
$c_1(L)=\delta([L])$.

We must check that its real image is
the curvature class with our sign
convention. Choose a good open cover
$U_i$, so all nonempty finite
intersections are contractible, and
smooth unitary frames $e_i$ for $L$.
Write $e_j=g_{ij}e_i$ on overlaps.
The functions $g_{ij}$ take values in
the unit circle, so choose real smooth
functions $f_{ij}$ satisfying
$$
g_{ij}=\exp(2\pi i f_{ij}).
$$
On triple overlaps,
$$
n_{ijk}=f_{ij}+f_{jk}-f_{ik}
$$
is integer valued and locally constant.
By construction of the connecting map,
this integer cocycle represents $c_1(L)$.

Write $\nabla e_i=\theta_i\otimes e_i$.
The Leibniz rule gives
$$
\theta_j-\theta_i
=g_{ij}^{-1}dg_{ij}
=2\pi i\,df_{ij}.
$$
Set $A_i=\frac{i}{2\pi}\theta_i$.
These are real one-forms because the
frames are unitary and the connection
is hermitian. Thus
$$
A_j-A_i=-df_{ij},\qquad
dA_i=F,\qquad
F=\frac{i}{2\pi}\Theta.
$$

These equations identify the de Rham
class of $F$ with the real image of
the cocycle $n$. To see the sign
explicitly, use the Čech--de Rham
complex. Denote its Čech differential
by $\check\delta$ and its total
differential on Čech degree $p$ by
$D=\check\delta+(-1)^p d$.
The equations above say
$$
DA=-df+F,\qquad Df=n-df.
$$
Consequently
$$
F-n=D(A-f).
$$
Thus $F$ and $n$ represent the same
class under the Čech--de Rham
comparison, proving integrality.

Finally, inclusion of holomorphic
functions into smooth functions gives
a map from the holomorphic exponential
sequence to the smooth one.
Naturality of the connecting maps
identifies this integral class with
the one used in the proof of
Theorem \ref{complex-geometry-theorem-11-classes}.
\end{proof}

\begin{remark}
\label{complex-geometry-remark-weil-and-lefschetz}
There are three related statements here.

Lemma \ref{complex-geometry-lemma-first-chern-integral}
starts with a holomorphic hermitian line
bundle and identifies its normalized
curvature class with the real image of
$c_1(L)=\delta([L])$.
The same proof works for any smooth
hermitian line bundle with unitary
connection. The exponential sequence
makes $\delta([L])$ integral by
construction; the additional
Cech--de Rham comparison identifies it
with the curvature class.

The Weil integrality theorem gives the
converse with prescribed curvature.
If $\beta$ is a closed real two-form
on a smooth manifold $M$ and $[\beta]$
lies in the image of $H^2(M,\mathbb{Z})$,
there is a smooth hermitian line bundle
with unitary connection such that,
in the curvature convention used here,
$
\frac{i}{2\pi}\Theta_\nabla=\beta.
$
This realizes the form itself, not just
its cohomology class.
Indeed, the smooth exponential sequence
classifies smooth complex line bundles
by $H^2(M,\mathbb{Z})$.
Choose a bundle realizing an integral
lift of $[\beta]$, a hermitian metric,
and a unitary connection $\nabla_0$.
The preceding comparison gives a real
one-form $\eta$ satisfying
$
\beta-\frac{i}{2\pi}\Theta_{\nabla_0}
=d\eta.
$
Then
$
\nabla=\nabla_0-2\pi i\eta
$
is unitary and has the required
curvature. No complex structure or
Kähler hypothesis is needed here.

The Lefschetz $(1,1)$ theorem
(Theorem \ref{complex-geometry-theorem-11-classes})
instead concerns holomorphic line
bundles on a compact Kähler manifold
$X$. It says that
$
c_1:\operatorname{Pic}(X)
\twoheadrightarrow H^{1,1}(X,\mathbb{Z}).
$
Here the target consists of integral
classes whose complex images have
type $(1,1)$, including torsion classes.
The holomorphic exponential sequence
gives this surjectivity; the map need
not be injective. Thus every integral
class has a smooth line bundle, and
the $(1,1)$ condition characterizes
which classes have a holomorphic one.
Weil additionally prescribes a unitary
connection realizing a given integral
closed form.
\end{remark}

\noindent
These forms are real because a hermitian
connection has skew-hermitian curvature
on real tangent vectors, so
$\operatorname{tr}\Theta$ is purely
imaginary. Closedness was proved above.
The class is independent of the metric:
for two connections, their difference
$A=\nabla'-\nabla$ is a global
endomorphism-valued one-form, and locally
$A=\theta'-\theta$. Consequently,
$$
\operatorname{tr}\Theta'
-\operatorname{tr}\Theta
=d(\operatorname{tr}A),
$$
so their first Chern forms differ by
an exact form.

\begin{remark}[Chern polynomial]
\label{complex-geometry-remark-chern-polynomial}
The trace formula is the first coefficient
of a determinant. For a rank-$r$ bundle
with curvature matrix $\Theta$, Chern--Weil
theory gives the \emph{Chern polynomial}
in real cohomology:
$$
\begin{aligned}
c_t(E)_{\mathbb{R}}
&=\left[\det\left(
I+t\frac{i}{2\pi}\Theta\right)\right]\\
&=1+\sum_{k=1}^r t^k c_k(E)_{\mathbb{R}}.
\end{aligned}
$$
Here $t$ is a formal variable: this is a
polynomial of degree at most $r$, not an
infinite series. Products of matrix entries
are wedge products, so the coefficient of
$t^k$ is a closed $2k$-form. Its cohomology
class is independent of the connection and
is the real image of the integral class
$$
c_k(E)\in H^{2k}(M,\mathbb{Z}).
$$
Evaluating at one gives the \emph{total Chern
class}. For rank two, explicitly,
$$
c_2(E)_{\mathbb{R}}
=\left[\det\left(
\frac{i}{2\pi}\Theta\right)\right].
$$
See Huybrechts, \emph{Complex Geometry:
An Introduction}, Section 4.4,
``Chern Classes'' \cite[Section 4.4]{huc}.
\end{remark}

\section{Hodge theory with values in a vector
bundle and Serre duality}
\label{complex-geometry-section-bundle-valued-hodge-serre}

\begin{reference}
See \cite[Section 4.1]{huc} for Hodge theory
with values in a holomorphic vector bundle
and \cite[Proposition 4.1.15]{huc} for
Serre duality. For the Kähler identities
with values in a bundle, see
\cite[Lemma 5.2.3]{huc}; the curvature
correction is discussed in the same section.
\end{reference}

\noindent
Let $X$ be a compact complex manifold
equipped with a hermitian metric and let
$(E,h)$ be a holomorphic hermitian vector
bundle. We write
$
\mathcal A^{p,q}(X,E)
$
for smooth $(p,q)$-forms with values in $E$.
The Chern connection decomposes as
$
\nabla=\partial_E+\bar\partial_E.
$
The operator $\bar\partial_E$ is determined
by the holomorphic structure of $E$ and
satisfies
$
\bar\partial_E^2=0.
$
Hence one defines
$$
H^{p,q}_{\bar\partial}(X,E)
=
\frac{\ker(\bar\partial_E:
\mathcal A^{p,q}(X,E)\to
\mathcal A^{p,q+1}(X,E))}
{\operatorname{im}(\bar\partial_E:
\mathcal A^{p,q-1}(X,E)\to
\mathcal A^{p,q}(X,E))}.
$$

\begin{theorem}[Hodge theorem for
bundle-valued Dolbeault cohomology]
\label{complex-geometry-theorem-bundle-valued-hodge}
Let $\Delta_{\bar\partial_E}$ be the
$\bar\partial_E$-Laplacian. Then
$$
H^{p,q}_{\bar\partial}(X,E)
\cong
\mathcal H^{p,q}_{\bar\partial_E}(X)
:=
\ker\Delta_{\bar\partial_E}.
$$
In other words, every Dolbeault class with
values in $E$ has a unique harmonic
representative.
\end{theorem}

\noindent
This part only uses elliptic Hodge theory;
$X$ need not be Kähler. If $X$ is Kähler,
the ordinary Kähler identities have
bundle-valued analogues.

\begin{theorem}[Nakano identities]
\label{complex-geometry-theorem-nakano-identities}
Let $X$ be Kähler and let
$\Lambda=(\omega\wedge-)^*$. Then
$$
[\Lambda,\bar\partial_E]
=-i\,\partial_E^*,
\qquad
[\Lambda,\partial_E]
=i\,\bar\partial_E^*.
$$
Consequently the two bundle-valued
Laplacians satisfy the
Bochner--Kodaira--Nakano identity
$$
\Delta_{\bar\partial_E}
=
\Delta_{\partial_E}
+
[i\Theta(E),\Lambda],
$$
where $\Theta(E)$ is the curvature of the
Chern connection.
\end{theorem}

\noindent
This is the essential difference from the
scalar case: the two Laplacians need not
coincide. They do coincide when the curvature
term vanishes, for example for the trivial
flat bundle. Also, for a general holomorphic
bundle,
$
\nabla^2=\Theta(E)\neq0,
$
so $\nabla$ does not define a de Rham
complex. The natural generalization above is
therefore Dolbeault cohomology with values in
$E$; de Rham cohomology with coefficients in
$E$ requires a flat connection.

\medskip\noindent
Now we apply the bundle-valued Hodge theorem
to Serre duality. The hermitian metrics on
$X$ and $E$ define a conjugate-linear
bundle-valued Hodge star
$$
\bar *_{E}:
\mathcal A^{p,q}(X,E)
\longrightarrow
\mathcal A^{n-p,n-q}(X,E^\vee).
$$
Together with harmonic representatives this
gives the following perfect pairing.

\begin{theorem}[Serre duality]
\label{complex-geometry-proposition-serre-duality}
Let $X$ be a compact complex manifold of
dimension $n$ and let $E$ be a holomorphic
vector bundle on $X$. Then the pairing
$$
H^{p,q}_{\bar\partial}(X,E)
\times
H^{n-p,n-q}_{\bar\partial}(X,E^\vee)
\longrightarrow
\mathbb C,
\qquad
([\alpha],[\beta])
\longmapsto
\int_X \alpha\wedge\beta,
$$
where the $E\otimes E^\vee$ factors are
evaluated, is nondegenerate. Equivalently,
$$
H^{p,q}_{\bar\partial}(X,E)
\cong
H^{n-p,n-q}_{\bar\partial}(X,E^\vee)^\vee.
$$
\end{theorem}

\begin{proof}
By Theorem
\ref{complex-geometry-theorem-bundle-valued-hodge}, choose
the harmonic representative $\alpha$ of a
nonzero class. Then
$\bar *_{E}\alpha$ is an
$E^\vee$-valued harmonic form of type
$(n-p,n-q)$, and
$$
\int_X\alpha\wedge\bar *_{E}\alpha
=
\|\alpha\|_{L^2}^2>0.
$$
Hence the pairing is nondegenerate.
\end{proof}

\noindent
Using Dolbeault's theorem,
$$
H^{p,q}_{\bar\partial}(X,E)
\cong
H^q(X,\Omega_X^p\otimes E),
$$
Serre duality becomes
$$
H^q(X,\Omega_X^p\otimes E)
\cong
H^{n-q}(X,\Omega_X^{n-p}\otimes
E^\vee)^\vee.
$$
In particular, taking $p=0$ gives the usual
form
$$
H^q(X,E)
\cong
H^{n-q}(X,K_X\otimes E^\vee)^\vee.
$$

\section{Line bundles on $\mathbb{P}^n$}
\label{complex-geometry-section-line-bundles-on-pn}

\begin{definition}
\label{complex-geometry-definition-line-bundles-in-pn}
$$
\mathcal{O}(1):=\mathcal{O}(-1)^\vee,\qquad
\mathcal{O}(d):=\bigotimes_d \mathcal{O}(1).
$$
\end{definition}

\noindent
To show that the sections of $\mathcal{O}(d)$
are degree-$d$ homogeneous polynomials we
first
show that the fibers are degree-$d$
homogeneous functions
(not necessarily polynomials)
and then pass to the global situation.

\begin{lemma}
\label{complex-geometry-lemma-fiber-od-homog-degre-d-funct}
The fiber $\mathcal{O}(d)_p$ is in bijection
with
degree-$d$ homogeneous functions on $p$.
\end{lemma}

\begin{proof}
So far we can only say that the fibers of
$\mathcal{O}(-1)\otimes \mathcal{O}(1)$ are
the fibes of $\mathcal{O}$.
This means that for a fixed $f \in
\mathcal{O}(1)_p$
we have a linear functional $z \mapsto
f\otimes z$
on $\mathcal{O}(-1)_p$. Any linear functional
on $\mathcal{O}(-1)_p$
can be realised this way.

To generalize this for degree $d$ we do as
follows.
First note that
$$
\Hom(\mathcal{O}(-1)_p,\mathbb{C})
\otimes\Hom(\mathcal{O}(-1)_p,\mathbb{C})
\cong
\Hom(\mathcal{O}(-1)_p\otimes\mathcal{O}(-1)
_p,\mathbb{C}).
$$
(This is not true for general modules,
but in this case it is using a basis.)

By the universal property of tensor product
(see Commutative Algebra Theorem
\cite{theorem-unive-prope-tenso-produ})
we know that
$\Hom(\mathcal{O}(-1)_p \otimes … \otimes
\mathcal{O}(-1)_p,\mathbb{C})$
is the set of multilinear maps
$\text{Lin}(\mathcal{O}(-1)_p
\times…\times\mathcal{O}(-1)_p,\mathbb{C})$.
Any such map $L$ gives a degree-$d$ function
(a polynomial) $P$ by
$P(x):=L(x,…,x)$.
In fact, since the fiber $\mathcal{O}(-1)_p$
is one-dimensional,
we can reconstruct $L$ from $P$ by defining
the multilinear map on a basis (which is
composed by a single vector $x$)
by $L(x,…,x):=P(x)$. This independs on the
choice
of basis since if $x'=\lambda x$ and we
define $L'(x',…,x'):=P(x')$,
We see that $\lambda^d L'(x,…,x)=L'(x',…,x')=
=P(x')=P(\lambda x)=\lambda^d P(x)=\lambda^d
L(x,…,x)$.
\end{proof}

\noindent
To pass from the local situation in Lemma
\ref{complex-geometry-lemma-fibers-of-od}
to a section we do as follows.

\begin{exercise}
\label{complex-geometry-exercise-polynomials-are-sections}
\begin{reference}
\cite[Exercise 2.2.8]{huc}, and
\cite[Lemma 3.30]{lec} for the answer.
\end{reference}
Show that any non-trivial homogeneous
degree-$d$ polynomial $0 \neq f \in
\mathbb{C}[z_0,…,z_n]_d$ can be considered as
a non-trivial section of $\mathcal{O}(d)$ on
$\mathbb{P}^n$.
\end{exercise}

\begin{proof}
Let $P$ be a homogeneous degree-$d$
polynomial. We define a section fiberwise and
then show it is holomorphic. For $p \in
\mathbb{P}^n$ we restrict $P$ to the
1-dimensional vector space that $p$ is in
$\mathbb{C}^{n+1}$. To show this is
holomorphic we can evaluate our section on a
local frame for $\mathcal{O}(d)$, which gives
a holomorphic function since $P$ is a
polynoimal (see \cite[Proposition 3.35]{lec}.
\end{proof}

\begin{theorem}
\label{complex-geometry-theorem-secti-od-degre-d-homog-poly}
Sections of $\mathcal{O}(d)$ are in
correspondence with homogeneous degree-$d$
polynomials.
\end{theorem}

\begin{proof}
It only remains to prove the direct
implication.
Let $s \in H^0(\mathcal{O}(d))$.
Then $s(p) \in \mathcal{O}(d)_p$ is a
homogeneous degree-$d$ function $f_p$
by Lemma
\ref{complex-geometry-lemma-fiber-od-homog-degre-d-funct}
defined on the vector space $p$.

Now define a holomorphic homogeneous
degree-$d$ function $f$ on
$\mathbb{C}^{n+1}\setminus\{0\}$ as follows.
For any $z \in
\mathbb{C}^{n+1}\setminus\{0\}$ we can take
$p$
to be the line passing through $z$
and assign the number $f_p(z)$
(since $z \in p$).

Finally we show that $f$ must be a polynomial
by considerations on its Taylor series
after extending to $\mathbb{C}^{n+1}$ via
Hartog's theorem
(see \cite[Theorem 3.36]{lec}).
\end{proof}


\begin{exercise}
\label{complex-geometry-exercise-line-bundle-trivial-iff}
\begin{reference}
\cite[Exercise 2.2.5]{huc}
\end{reference}
Let $L$ be a holomorphic line bundle on a
compact complex manifold $X$.
Show that $L$ is trivial if and only if $L$
and its dual $L^\vee$
admit non-trivial global sections.
\end{exercise}

\begin{proof}
Direct implication is easy: trivial bundle
has 1-dimensional space of sections.
For the converse choose nonvtrivial sections
$s$ of $L$
and $s^\vee$ of $L^\vee$.
Then $s \otimes s^\vee$ which is a section
of the trivial bundle.
If such a section has a zero, then it must be
zero throughout $X$.
\end{proof}

\section{Line bundles and divisors}
\label{complex-geometry-section-line-bundles}

\noindent
A simple way of associating a line bundle to
an irreducible
hypersurface $Y\subset X$ is by noting that
by Weierstrass Division Theorem
\ref{complex-analysis-theorem-%
weierstrass-division-theorem}
the ideal sheaf of $Y$ is locally generated,
i.e. there is an open cover $U_i$ of $X$
such that $\mathcal{I}(U_i)=(f_i)$.
This means that at intersections we have
$f_j \in (f_i)$ and $f_i\in (f_j)$,
which says that $f_i$ and $f_j$ are
associates,
i.e. there is a unit $u_{ij} \in
\mathcal{O}^*_X(U)$
(the units in the ring of holomorphic
functions
are nowhere vanishing functions)
such that $f_i=u_{ij}f_j$.
Thus $u_{ij}:U_i\cap U_j \to
\text{GL}(1,\mathbb{C})$
define a cocycle in $H^1(X,\mathcal{O}^*_X) =
\text{Pic}(X)$.

\begin{definition}
\label{complex-geometry-definition-divisor}
A {\it Weil divisor} is a formal sum
$$
D=\sum a_i Y_i
$$
where $a_i\in \mathbb{Z}$ and $Y_i$ are
irreducible hypersurfaces.
\end{definition}

\noindent
An irreducible hypersurface $Y$ clearly
defines a Weil divisor,
and so does a reducible reduced hypersurface
with irreducible
components $Y_i$ putting $a_i=1$.
To allow for coefficients $>1$ we may admit
nonreduced hypersurfaces,
so that we have larger multiplicites.

This can also be presented as follows.
Given a holomorphic function $g$, we define
{\it order of vanishing of $g$ at $P \in Y$}
to be the largest number $\text{ord}_{P,Y}g$
such that
$$
g=f_i^{\text{ord}_{P,Y}g}h,\qquad h \in
\mathcal{O}^*_{X,P},
$$
where $f_i$ is an irreducible function
vanishing along $Y$ near $P$,
i.e. a generator of the local ring of $Y$.
Further, by coherence/preservation relative
primes in nearby stalks
(see \cite[p. 696]{gri})
the order of $g$ is constant for any $P \in
Y$.

Thus for any holomorphic function $g$
we can define the divisor
$$
\text{div}(g):=\sum_{Y}\text{ord}_Y(g)Y
$$
where the sum ranges in all irreducible
hypersurfaces of $X$.

Similarly, for a meromorphic function
$g=g_1/g_2$
we define the order of vanishing of $g$ at
$Y$ to be
$$
\text{ord}_Y(g)=\text{ord}_Y(g_1)-\text{ord}
_Y(g_2).
$$
The associated divisor to a meromorphic
function
is defined identically as that of a
holomorphic function,
and we observe that there might appear
negative coefficients.
Divisors obtained in this way from
meromorphic functions
are called {\it principal divisors}.

This shows how to assign a divisor from any
meromorphic function,
i.e. it's a map
$$
H^0(X,\mathcal{M}^*) \to \text{Div}(X).
$$
Note that if $g$ is in fact a nowhere
vanishing meromorphic or
holomorphic function,
the corresponding divisor is trivial,
i.e. it is in the kernel of our map.
In fact,
$$
H^0(X,\mathcal{M}^*/\mathcal{O}^*) \simeq
\text{Div}(X).
$$
Indeed, given a divisor we can assign a
section of
$\mathcal{M}^*/\mathcal{O}_X^*$ as follows.
If the divisor is an irreducible hypersurface
locally given by $V(f_i)$
we assign the meromorphic function given
locally as $f_i$.
If the divisor has several irreducible
connected components
we consider the product of the irreducible
generators
at every point.

\medskip\noindent
It remains to connect $\text{Div}(X)$ and
$\text{Pic}(X)$.
The idea is that a section $f$ of
$\mathcal{M}^*/\mathcal{O}_X^*$
defines a cocycle in $H^1(\mathcal{O}_X^*)$.
Indeed, $f$ is given, by definition of
quotient sheaf (?),
by the data of an open cover $U_i$ and
meromorphic functions
$$
f_i=\frac{g_i}{h_i} \in \mathcal{M}^*
$$
which coincide at the intersections $U_i \cap
U_j$,
that is,
$$
f_i|_{U_i\cap U_j} = f_j|_{U_i\cap U_j}.
$$
But since we are in the quotient ring,
equality means that
$$
f_i f_j^{-1} \in \mathcal{O}_X^*.
$$
This defines a cocycle, and thus a map
$$
\text{Div}(X) \to \text{Pic}(X).
$$
This map can also be though of as induced
from
$$
\xymatrix{
0\ar[r]
&\mathcal{O}^*_X\ar[r]
&\mathcal{M}^*\ar[r]
&\mathcal{M}^*/\mathcal{O}_X^*\ar[r]
&0
}
$$
via the cohomology long exact sequence giving
$$
\xymatrix{
\cdots\ar[r]
& H^0(\mathcal{M}^*_X)\ar[r]
&H^0(\mathcal{M}^*_X/\mathcal{O}^*_X)\ar[r]
&H^1(\mathcal{O}_X^*)\ar[r]
&\cdots
}
$$
Indeed, the kernel of the map $\text{Div}(X)
\to \text{Pic}(X)$
is given by the principal divisors,
i.e. those coming from meromorphic sections.

In general there is no way to ensure this map
is an isomorphism.

\medskip\noindent
Now we discuss how to obtain submanifolds
from sections.
We will see that $\text{Div}(X)$
is isomorphic to the group of line bundles
equipped with a meromorphic section.


Bear in mind the following observation for a
rank-$r$
vector bundle $E \to X$.
Suppose $(\varphi_i,U_i)$ are local
trivializations of $E$,
i.e. $\varphi_i:\pi^{-1}(U_i)
\xrightarrow{\simeq}U_i \times \mathbb{C}^k$.
Take a local section $s :U_i \to
\pi^{-1}(U_i)$.
Take a local basis $e_i$ of $\mathbb{C}^k$.
Then
$$
s_i(x)=f_i^k(x)\varphi_i^{-1}(x,e_k)
$$
for some functions $f_i^k:U_i \to
\mathbb{C}$.

At an intersection $U_i \cap U_j$,
we have
$$
s_i(x)=f_i^k(x)\varphi_i^{-1}(x,e_k)
=
s_j(x)=f_j^k(x)\varphi_j(x,e_k).
$$
Applying $\varphi_i$ we see that
$$
f_i^k e_k = f_j^k g_{ij}(e_k).
$$

In the case of a line bundle this just says
that sections
are given by local functions $f_i:U_i \to
\mathbb{C}$ satisfying
$$
f_i=g_{ij}f_j
$$
where $g_{ij}$ is the cocycle defining the
line bundle.

Let $L$ be a line bundle.
A {\it meromorphic section} of $L$ is a
section
of $\mathcal{O}(L)
\otimes_\mathcal{O}\mathcal{M}$,
which is given by a collection of meromorphic
functions
$s_i \in \mathcal{M}(U_i)$
such that
$$
s_i=g_{ij}s_j.
$$
That is, sections of the line bundle viewed
as a module over $\mathcal{O}$
with coefficients that can be meromorphic
functions.

The orders of the coefficient meromorphic
functions
is well defined along hypersurfaces of $X$,
and thus
we can define the
{\it divisor $\text{div}(s)$ associated to
the section $s$ of $L$}
by
$$
\text{div}(s)
:= \sum_Y\text{ord}_Y(s)Y.
$$

Conversely, given any divisor $D$
we can assign a line bundle $L$ as before,
namely if the divisor is given by the
meromorphic functions $f_i$
we put $g_{ij}=f_i/f_j$.
But now we observe
the existence of a distinguished section $s$
such that $\text{div}(s)=L$.
Indeed, we note that
$$
f_i=f_if_j^{-1}f_j=g_{ij}f_j,
$$
which is a section of $L$,
and in fact (check…) $\text{div}(s)=L$.

We conclude that $L$ is the line bundle
associated to a divisor
if and only if it has a meromorphic section
not identically zero.
It is the line bundle associated to an
effective divisor
if and only if it has a holomorphic section
not identically zero
--- whose zero set is the divisor.

\medskip\noindent
Upshot. We have more than the submanifold
associated to a section:
we have the divisor associated to a
meromorphic section,
which gains to becoming less geometrically
intuitive
the virtue of being an element of a group.

\medskip\noindent
Now we turn to line systems.
Take a divisor
$$
D=\sum a_iV_i.
$$
The divisor
$$
D+(f)
$$
is linearly equivalent to $D$ --- their
difference is principal.
If
$$
D+(f)\geq 0,
$$
then $f$ must be holomorphic away from $V_i$,
and on $V_i$ it can have poles with the
condition
$$
\text{ord}_{V_i}(f) \geq  -a_i.
$$
Denote $\mathcal{L}(D)$ the space of such
functions.

Note that if $s_0$ is a section such that
$(s_0)=D$, then multiplication by $s_0$ gives
an identification
$$
\mathcal{L}(D) \xrightarrow{\otimes s_0}
H^0(X,L).
$$

\begin{definition}
\label{complex-geometry-definition-line-system}
We denote by
$$
|D|\subset \text{Div}(X)
$$
the set of all effective divisors linearly
equivalent to $D$.
\end{definition}

\noindent
$\mathcal{L}(D)=H^0(X,L)$ is almost the same
thing as $|D|$.
At least in the compact case,
two functions in $f,f' \in \mathcal{L}(D)$
give the same divisor when $f=\lambda f'$
for some constant function $\lambda \in
\mathbb{C}$.

Thus
$$
|D|=\mathbb{P}H^0(X,L).
$$
This is called a {\it complete linear
system}.
More generally, a {\it linear system} is a
linear
subspace of $\mathbb{P}H^0(X,L)$.
The {\it base locus} of $|D|$
is the set of points where all the global
sections
generating $H^0(X,L)$ vanish.
We will see in Section
\ref{complex-geometry-section-kodaira-embedding}
the relationship between line systems and
embeddings to projective space.

\begin{theorem}[Bertini]
\label{complex-geometry-theorem-bertini}
The generic divisor of a linear system
is smooth away from the base locus.
\end{theorem}

\medskip\noindent
Now let's introduce the first adjunction
formula,
which says that when a divisor is a smooth
hypersurface,
it's associated line bundle is its
holomorphic normal bundle.

\begin{theorem}[Adjunction formula I]
\label{complex-geometry-theorem-adjunction-formula-i}
\begin{reference}
\cite[p.146]{GH}
\end{reference}
If $X$ is a smooth hypersurface of a smooth
complex manifold $Y$, then
\begin{equation}
\label{complex-geometry-equation-adjunction-i}
\mathcal{O}_Y(X)|_{X} \simeq
\mathcal{N}_{X/Y},
\end{equation}
where $\mathcal{N}_{X/Y}:=T_Y'/T_X'$.
\end{theorem}

\begin{proof}
If $X$ is locally defined
by the vanishing of $f_i$ on an open cover
$U_i$ of $Y$,
then $\mathcal{O}_Y(X)$ is given by the
cocycle $f_{ij}:=f_if_j$.

To prove Equation \ref{complex-geometry-equation-adjunction-i}
it's enough to show that
$\mathcal{O}_Y(X)|_{X}\otimes
\mathcal{N}^\vee_{X/Y}$
is trivial, which can be done by constructing
a nonvanishing section.
Such a section is given by
$$
df_i=d(f_{ij}f_j)=df_{ij}\cdot
f_j+f_{ij}df_j,
$$
since the first term vanishes along $X$
(because $f_j$ does)
and the remaining terms define precisely a
section
of ${T'}^\vee_Y/{T'}^\vee_X$,
the conormal bundle ---
indeed, sections of such a bundle
are holomorphic differencials vanishing along
$X$.
\end{proof}

\noindent
Now let's find the canonical bundle of
$\mathbb{P}^n$.
By the above discussion, we need the inverse
map of
$$
\text{Div}(\mathbb{P}^n) \to \{L \in
\text{Pic}(X):
\text{$L$ has a nonzero meromorphic
section}\},
$$
that is, we need to find a nonzero
meromorphic section $s$
of the canonical bundle $\Omega^{n,0}=K$.
Then the divisor will be given by
$$
\sum_Y \text{ord}_Ys [Y].
$$
A meromorphic section is given by local
sections $s_i$
so that
$$
s_i=g_{ij}s_j
$$
where $g_{ij}$ is the cocycle of $K.$

Consider in the canonical chart $z_0\neq 0$
of $\mathbb{P}^n$
the local meromorphic section of $K$
$$
s_0=
\frac{dw_1}{w_1}\wedge\frac{dw_2}{w_2}\wedge…
\wedge\frac{dw_n}{w_n}
$$
where $w_i=z_i/z_0$ are the coordinate
functions in $z_0\neq 0$.

In any chart $z_j\neq 0$ we obtain the local
expression
$$
(-1)^j\frac{dw_0'}{w_0'}
\wedge…
\wedge \widehat{\frac{dw_j'}{w_j'}}
\wedge…
\wedge \frac{dw_n'}{w_n'}
$$
according to the canonical change of
coordinates, which says
$$
w_i=\frac{w_i'}{w'0}, i\neq j,\qquad
w_j=\frac{1}{w_0'}.
$$

The point is that we constructed a
meromorphic section which has
a single pole along each hyperplane $z_i=0$,
which means
$$
K=-(n+1)H.
$$
\medskip\noindent
The next adjunction formula is the key to
computing
the canonical bundle of submanifolds.
It is a consequence of this linear algebra
fact:

\begin{lemma}
\label{complex-geometry-lemma-detrminant-of-ses}
If
$$
\xymatrix{
0\ar[r]
&E'\ar[r]
&E\ar[r]
&E''\ar[r]
&0
}
$$
is a short exact sequence of vector bundles
(of possibly different ranks),
then
$$
\det(E)\cong \det(E')\otimes \det(E'').
$$
\end{lemma}

\begin{proof}
First notice that each of the three
determinant bundles gives fiberwise the
top-exterior
power of the corresponding vector spaces,
thus varying the rank of the exterior power
taken.
Then define a map
$$
\begin{aligned}
x_1 \wedge…\wedge x_{\text{rk}E'}
&\otimes x_{\text{rk}E'+1}\wedge…\wedge
x_{\text{rk}E'+\text{rk}E''}\\
&\qquad \qquad \mapsto x_1 \wedge…\wedge
x_{\text{rk}E'}
\wedge x'_{\text{rk}E'+1}\wedge…\wedge
x'_{\text{rk}E'+\text{rk}E''}
\end{aligned}
$$
where $x'_{\text{rk}E'+i}$ maps to
$x_{\text{rk}E'+i}$
under the surjection given by the exact
sequence.

It's clear that such a map is injective:
if any monomial in the image is zero, then
one of the
factors must be zero, and then the element
we began with also vanishes.
Being two vector spaces of dimension 1,
the map must also be surjective.
Check this map is well-defined!
This gives an isomorphism in every fiber.

In fact the isomorphism above acts like 1.
To show that this glues to a bundle
isomorphism
we have to check some cocycle condition
as in Proposition
\ref{complex-geometry-proposition-isomo-vecto-bundl}.
In this case the transformation functions
$A_i$
are just 1 by definition of the morphism.
The cocycle condition follows from the fact
that
any short exact sequence splits,
so on $U_{ij}$ we have the isomorphism
$E(U_{ij}) \simeq E'(U_{ij})\oplus
E''(U_{ij})$.
Recall that the cocycle of the direct sum of
bundles is given by
$\begin{pmatrix}h'_{ij}&0\\
0&h''_{ij}\end{pmatrix}$
so that when taking determinants
we obtain that
$g_{ij}=h'_{ij}h''_{ij}:=h_{ij}$
(since taking determinant bundle is taking
determinant of
the transition functions).
Thus the transition functions of the two
bundles in question
actually coincide,
and the cocycle condition in Proposition
\ref{complex-geometry-proposition-isomo-vecto-bundl}
is the trivial equality $g_{ij}=h_{ij}$.
\end{proof}

\noindent
Here's a cleaner version of this proof:

\begin{proof}
Choose an open cover $\{U_i\}$ on which the
sequence splits, so that
\[
E|_{U_i}\simeq E'|_{U_i}\oplus E''|_{U_i}.
\]
Pick local frames $(e'_{i,1},\dots,e'_{i,r})$
for $E'|_{U_i}$ and
$(e''_{i,1},\dots,e''_{i,s})$ for
$E''|_{U_i}$, and use the splitting to obtain
a frame $(e'_{i,\bullet},e''_{i,\bullet})$ of
$E|_{U_i}$.

On overlaps $U_{ij}$ the transition matrix of
$E$ in these adapted frames has
block diagonal form
\[
g^E_{ij}=
\begin{pmatrix}
g'_{ij} & 0\\
0 & g''_{ij}
\end{pmatrix},
\]
hence
\[
\det(g^E_{ij})=\det(g'_{ij})\,\det(g''_{ij}).
\]
Thus the transition functions of $\det(E)$
coincide with those of
$\det(E')\otimes\det(E'')$, giving an
isomorphism of line bundles. In these
frames it is given by
\[
(e'_{i,1}\wedge\cdots\wedge e'_{i,r})\otimes
(e''_{i,1}\wedge\cdots\wedge e''_{i,s})
\longmapsto
e'_{i,1}\wedge\cdots\wedge e'_{i,r}\wedge
e''_{i,1}\wedge\cdots\wedge e''_{i,s}.
\]

If a different local splitting is chosen, the
adapted frames differ by a
matrix of the form
\[
\begin{pmatrix}
I & B\\
0 & I
\end{pmatrix},
\]
whose determinant is $1$. Therefore the
induced map on determinant lines is
unchanged, showing that the above isomorphism
is independent of the splitting
and hence canonical.
\end{proof}

\begin{theorem}[Adjunction formula II]
\label{complex-geometry-theorem-adjunction-ii}
Let $X$ be a complex submanifold of a complex
manifold $Y$.
Then
\begin{equation}
\label{complex-geometry-equation-adjunction-ii}
(K_Y\otimes \mathcal{O}_Y(X))|_{X}\cong K_X.
\end{equation}
\end{theorem}

\begin{proof}
From the holomorphic tangent bundle exact
sequence we get
$$
\xymatrix{
0\ar[r]
&T'_X\ar[r]
&T'_Y|_{X}\ar[r]
&N_{X/Y}:=T'_Y/T'_X\ar[r]
&0
}
$$
Taking determinant of this exact sequence we
obtain the result.
\end{proof}


\medskip\noindent
Now we turn to Chern classes of line bundles.

Slogan: the first Chern class of a line
bundle
is the Poincaré dual of the divisor's
fundamental class
and the cohomology class of the curvature
form
of any connection of $L$.
It is an integral $(1,1)$-curvature form.

The {\it curvature form} is the differential
2-form
that represents the curvature operator with
respect to a given connection
on the line bundle.
This can also be defined for affine vector
bundles of any rank.

\begin{proposition}
\label{complex-geometry-proposition-first-chern-class}
\begin{reference}
\cite[p. 141]{gri}
\end{reference}
\begin{enumerate}
\item For any line bundle $L$ with curvature
form $\Theta$,
$$
c_1(L)=\left[\frac{\sqrt{-1}}{2\pi}
\Theta\right]
\in H^2_{dR}(X).
$$

\item If $L=[D]$ for some $D \in
\text{Div}(M)$,
$c_1(L)$ is the Poincaré dual of the
fundamental class of $D$.
\end{enumerate}
\end{proposition}

\medskip\noindent
Finally let's just put some more exercises
here.

\begin{exercise}[Curves on surfaces and the
Hodge index theorem]
\label{complex-geometry-exercise-curve-surfa-hodge-index}
\begin{reference}
\cite[Exercise 2, Chapter 11]{voi}
\end{reference}
Let $S$ be a smooth projective surface and
let $\phi:S \to \mathbb{P}^n$
be a holomorphic map satisfying the following
properties:
\begin{itemize}
\item There exists a curve $C \subset S$
such that $\phi(C)$ is a finite set of
points.
(Here $C$ is a not necessarily connected
finite union
of irreducible hypersurfaces $C_i$ which are
not necessarily smooth.)

\item There exists an open set of $S$ where
$\phi$ is an embedding.
\end{itemize}
\begin{enumerate}
\item Show that the class
$h=\phi^*c_1(\mathcal{O}_{\mathbb{P}^n}(1))$
satisfies $\left\langle{}h,h\right\rangle{} >0$ where
$\left\langle{}\cdot,\cdot\right\rangle{}$ is the
intersection form on $H^2(S,\mathbb{Z})$.

\item Show that $\left\langle{}[C_i],h\right\rangle{}=0$ for
any $i$.
Deduce from the Hodge index theorem that the
intersection form
$\left\langle{}\cdot,\cdot\right\rangle{}$ is negative
definite on the subgroup
of $H^2(S,\mathbb{Z})$ generated by the
$[C_i]$'s.

\item Show that $\left\langle{}C_i,C_j\right\rangle{}\geq 0$
for $i\neq j$.
Deduce then from the above that the $[C_i]$'s
are independent in $H^2(S,\mathbb{Z})$.
\end{enumerate}
\end{exercise}

\begin{proof}
\begin{enumerate}
\item Here's an outline of how to prove this
result.
First notice that $h$ is the class of a line
bundle
since $\varphi^*c_1=c_1\varphi^*$.
Thus $\left\langle{}h,h\right\rangle{}$ is the intersection
number
of two divisors in the class $h$.

To compute their intersection number
we use that $\varphi$ is an embedding on an
open set of $S$,
say $U \simeq \varphi(U) \subset
\mathbb{P}^n$.
To compute the intersection number we take
two generic hyperplanes $H_1$ and $H_2$
in $\mathbb{P}^n$ and intersect them with
$\phi(U)$.
The dimension of the resulting submanifold is
zero.
Thus the intersection is a finite number of
points.

\item Since $C_i$ is mapped to a set of
points in $\mathbb{P}^n$,
we compute the intersection number as above,
which gives zero
for generic choice of $H_1,H_2$. Hodge index
theorem tells us that
the signature of $\left\langle{}\cdot,\cdot\right\rangle{}$
is $(1,\rho-1)$ on
$H^2(S,\mathbb{Z})$. Since $h$ is positive
and $C_i$ is orthogonal to $h$,
it's necessary that $C_i$ are negative
vectors.

\item This is by definition/construction of
the intersection form on a surface:
for two irreducible divisors of the surface,
the intersection number is the number of
intersection points
of the divisors counted with multiplicity.
Such a number cannot be negative.
The fact that the $C_i$'s are independent
follows from a linear-algebra
argument.
\end{enumerate}
\end{proof}

\section{Kodaira Vanishing theorem}
\label{complex-geometry-section-kodaira-vanishing}

\noindent
In \cite[p. 150]{gri}
we see that the curvature form of the
tautological bundle
is given by $\partial\bar\partial \text{log}
|z|^2$
where $z$ is a nonzero section.
Not giving further details for now,
this implies that the Chern class of the
hyperplane bundle
is the Fubini-Study form in $\mathbb{P}^n$.
We already discussed that this form is
positive
in the sense of Definition
\ref{complex-geometry-definition-positive-11form}.

\begin{definition}
\label{complex-geometry-definition-positive-line-bundle}
\begin{reference}
\cite[p. 148]{gri}
\end{reference}
Let $M$ be a compact Kähler manifold.
A line bundle $L \to M$ is {\it positive}
if there exists a metric on $L$ with
curvature form $\Theta$
such that $(\sqrt{-1}/2\pi)\Theta$ is a
strictly positive $(1,1)$-form
(Definition
\ref{complex-geometry-definition-positive-11form}); here
strict positivity means
$
\frac{i}{2\pi}\Theta(v,Iv)>0
$
for every nonzero real tangent vector $v$.
\end{definition}

\begin{theorem}[Kodaira-Nakano Vanishing
Theorem]
\label{complex-geometry-theorem-kodaira-nakano-vanishing}
\begin{reference}
\cite[p. 154]{gri};
\cite[Proposition 5.2.2]{huc}
\end{reference}
[Let $M$ be a compact Kähler manifold of
complex dimension $n$.]
If $L \to M$ is a positive line bundle,
then
$$
H^q(M,\Omega^p(L)) = 0,\qquad
p+q>n=\dim_\mathbb{C} M.
$$
\end{theorem}

\begin{proof}
See \cite[Proposition 5.2.2]{huc} and
\cite[p. 154]{gri}.
The idea is to show that an $L$-valued
harmonic $(p,q)$-form must vanish when
$p+q>n$.

Choose a Hermitian metric on $L$ whose
Chern curvature $\Theta$ is strictly
positive after multiplication by $i$.
Set
$$
\omega=i\Theta.
$$
This is a real, closed, strictly positive
$(1,1)$-form, hence a Kähler form.
We use its Kähler metric on $M$.
This choice is allowed because the
cohomology group does not depend on the
metric used to find harmonic
representatives.

To avoid confusing the line bundle $L$
with a Lefschetz operator, write
$$
L_\omega(\alpha)=\omega\wedge\alpha,
\qquad
\Lambda=L_\omega^*.
$$
The curvature operator $i\Theta$ acts
on forms by wedge multiplication, so it
is exactly $L_\omega$. The Lefschetz
commutator from Proposition
\ref{complex-geometry-proposition-sl2-relat-exter-algeb}
therefore gives, on $(p,q)$-forms,
$$
[i\Theta,\Lambda]
=[L_\omega,\Lambda]
=(p+q-n)\operatorname{id}.
$$
Thus the Nakano identity
(Theorem \ref{complex-geometry-theorem-nakano-identities})
becomes
$$
\Delta_{\bar\partial_L}
=
\Delta_{\partial_L}
+(p+q-n)\operatorname{id}.
$$

By Dolbeault's theorem and the
bundle-valued Hodge theorem
(Theorem \ref{complex-geometry-theorem-bundle-valued-hodge}),
every class in $H^q(M,\Omega^p(L))$
has a harmonic representative $\alpha$.
In particular,
$\Delta_{\bar\partial_L}\alpha=0$.
Taking the global $L^2$ inner product
with $\alpha$ yields
$$
\begin{aligned}
0
&=\left\langle{}
\Delta_{\bar\partial_L}\alpha,\alpha
\right\rangle{}_{L^2}\\
&=\|\partial_L\alpha\|_{L^2}^2
+\|\partial_L^*\alpha\|_{L^2}^2\\
&\quad +(p+q-n)\|\alpha\|_{L^2}^2.
\end{aligned}
$$
The first two terms are nonnegative.
If $p+q>n$, the coefficient in the last
term is strictly positive, so
$\alpha=0$. Hence every cohomology
class is zero.
\end{proof}

\noindent
Here's the same result as it appears in
\cite[sea], where it is not proved.

\begin{theorem}[Kodaira Vanishing]
\label{complex-geometry-theorem-Kodaira-vanishing-sea}
\begin{reference}
\cite[21.5.8]{sea}
\end{reference}
Suppose $k$ is a field of characteristic 0,
and $X$ is a
smooth projective $k$-variety. Then for any
ample invertible sheaf $L$,
$H^{i}(X,K_X \otimes L)=0$ for $i>0$.
\end{theorem}

\noindent
For completeness we put another vanishing
theorem
concerning higher-rank vector bundles
(probably due to Serre though Griffith-Harris
don't say so).
The idea is that twisting an ample line
bundle with itself
enough times and tensoring with any vector
bundle $E$
will kill positive cohomology.

\begin{theorem}
\label{complex-geometry-theorem-vanis-highe-rank-verct-bund}
\begin{reference}
\cite[p. 159]{gri}
\end{reference}
Let $M$ be a compact, complex manifold and $L
\to M$
a positive line bundle.
Then for any holomorphic vector bundle $E$,
there exists $\mu$ such that
$$
H^q(M,L^\mu \otimes E)=0\qquad \text{for
}q>0, \mu\geq \mu_0.
$$

\end{theorem}

\section{Lefschetz hyperplane sections
theorem}
\label{complex-geometry-section-lefschetz-hyperplane}

Lefschetz hyperplane theorem holds for
any $n$-dimensional compact complex manifold
$M$
and $V \subset M$ a smooth hypersurface with
$L=[V]$ positive.
An example of this is when $M$ is projective
and $V$ is the pullback of the hyperplane
bundle
(which, as explained in Section
\ref{complex-geometry-section-kodaira-vanishing},
is positive since its curvature form is
Fubini-Study),
that is, $V=M \cap H$ for a generic
hyperplane $H \subset \mathbb{P}^N$.

\begin{theorem}[Lefschetz hyperplane
sections]
\label{complex-geometry-theorem-lefschetz-hyperplane}
The map
$$
H^q(M,\mathbb{Q}) \to H^q(V,\mathbb{Q})
$$
induced by the inclusion $V
\xymatrix{\ar@{^{(}->}[r]&} M$
is an isomorphism for $q \leq  n-2$
and injective for $q=n-1$.
\end{theorem}

\begin{proof}
Looks accessible! ``Purely topological''.
\end{proof}

\begin{exercise}
\label{complex-geometry-exercise-hypersurface-as-hyperplane-section}
Let \(Y\subset \mathbb P^n\) be a hypersurface
defined by a homogeneous polynomial
\(F(z_0,\dots,z_n)\) of degree \(d\). Show that
\(Y\) is naturally identified with a hyperplane
section of the \(d\)-uple Veronese embedding
\[
\nu_d:\mathbb P^n\longrightarrow \mathbb P^N.
\]
\end{exercise}

\begin{proof}
Recall that the \(d\)-uple Veronese embedding
is defined by all monomials of degree \(d\):
\[
\nu_d([z_0:\cdots:z_n])
=
[\cdots : z^I : \cdots],
\]
where
\[
z^I=z_0^{i_0}\cdots z_n^{i_n},
\qquad
|I|=d.
\]
Since \(F\) is homogeneous of degree \(d\), it
can be written as a linear combination of these
monomials:
\[
F=\sum_{|I|=d} a_I z^I.
\]
Let \((X_I)\) be the homogeneous coordinates on
the target projective space \(\mathbb P^N\),
corresponding to the monomials \(z^I\). Define
a hyperplane \(H\subset \mathbb P^N\) by
\[
H=\left\{\sum_{|I|=d} a_I X_I=0\right\}.
\]
Then, for a point \([z]\in\mathbb P^n\), we have
\[
\nu_d([z])\in H
\]
if and only if
\[
\sum_{|I|=d} a_I z^I=0.
\]
But the latter condition is exactly
\[
F(z)=0.
\]
Therefore
\[
Y
=
\{F=0\}
=
\nu_d^{-1}(H).
\]
Since \(\nu_d\) identifies \(\mathbb P^n\) with
its image, this gives
\[
Y
\simeq
\nu_d(\mathbb P^n)\cap H.
\]
Thus \(Y\) is a hyperplane section of the
Veronese image of \(\mathbb P^n\).
\end{proof}

\begin{exercise}
\label{complex-geometry-exercise-lefschetz-hypersurface}
Let $Y\subset \mathbb P^n$ be a smooth
hypersurface. Use the Lefschetz hyperplane
theorem to show that the restriction map
$
H^k(\mathbb P^n,\mathbb Z)
\longrightarrow
H^k(Y,\mathbb Z)
$
is an isomorphism for $k<n-1$, and is
injective for $k=n-1$. Conclude that all
cohomology of $Y$ below the middle degree is
inherited from $\mathbb P^n$, and that new
cohomology can only appear in degree $n-1$.
\end{exercise}

\begin{proof}
Suppose $Y$ is defined by a homogeneous
polynomial
$
F(z_0,\dots,z_n)
$
of degree $d$. By the previous exercise,
using the $d$-uple Veronese embedding
$
\nu_d:\mathbb P^n\longrightarrow \mathbb P^N,
$
the hypersurface $Y$ is identified with a
hyperplane section of the Veronese image:
$
Y\simeq \nu_d(\mathbb P^n)\cap H
$
for some hyperplane $H\subset \mathbb P^N$.

Since $\nu_d$ is an embedding, the image
$\nu_d(\mathbb P^n)$ is isomorphic to
$\mathbb P^n$. Therefore we may apply the
Lefschetz hyperplane theorem to the smooth
projective variety
$
X=\nu_d(\mathbb P^n)
$
and its smooth hyperplane section
$
X\cap H\simeq Y.
$

The Lefschetz hyperplane theorem gives that
the restriction map
$
H^k(X,\mathbb Z)\longrightarrow
H^k(X\cap H,\mathbb Z)
$
is an isomorphism for
$
k<\dim_\mathbb C X-1
$
and is injective for
$
k=\dim_\mathbb C X-1.
$
Since $\dim_\mathbb C X=n$, we obtain that
$
H^k(\mathbb P^n,\mathbb Z)
\longrightarrow
H^k(Y,\mathbb Z)
$
is an isomorphism for $k<n-1$, and injective
for $k=n-1$.

Thus, below the middle degree, the cohomology
of $Y$ is the same as the cohomology inherited
from $\mathbb P^n$. Since $Y$ has complex
dimension $n-1$, its middle degree is $n-1$.
Therefore any new cohomology can only appear
in degree $n-1$.
\end{proof}

\section{Complex varieties}
\label{complex-geometry-section-complex-varieties}

For a definition of sheaf
see Algebraic Geometry Definition
\ref{algebraic-geometry-definition-sheaf}.

\begin{lemma}
\label{complex-geometry-lemma-O-X-is-a-sheaf}
Let $X$ be a complex manifold of complex
dimension $n$.

The functor
$$
\begin{aligned}
\mathcal{O}_X: \mathit{Open}_X^{\text{op}}
&\longrightarrow \mathit{Set}
\\
U &\longmapsto \mathcal{O}_X(U):=\{f:U\to
\mathbb{C}^n:
f \text{ is holomorphic}\}\\
i:V \hookrightarrow V &\longmapsto
\mathcal{O}_X(i):\mathcal{O}_X(U) \to
\mathcal{O}_X(V)
\end{aligned}
$$
where $\mathcal{O}_X(i)$ is given by
restriction,
is a sheaf.
\end{lemma}

\begin{proof}
$\mathcal{O}_X$ is a presheaf,
i.e. a functor,
by definition: restriction
maps the identity to the identity
and preserves compositions.

To check the condition of being a sheaf let
$U_i$ be an open cover of some open set $U$
of $X$
and $f_i:U_i \to X$ sections of
$\mathcal{O}_X(U_i)$,
that is, holomorphic functions, which
coincide pairwise in intersections.
Then it is obvious that the function $f:U \to
\mathbb{C}$,
$x \mapsto f_i(x)$ for any $U_i \ni x$
is well defined. It is also holomorphic
since holomorphicity is a condition defined
on
open sets.
\end{proof}

\medskip\noindent
A meromorphic function is a function that
can be written locally as a quotient of two
holomorphic functions:

\begin{definition}
\label{complex-geometry-definition-meromorphic-functions}
\begin{reference}
\cite[Definition 2.1.8]{huc}
\end{reference}
A {\it meromorphic function} on a complex
manifold $X$
is a map
$$
f:X \to \coprod_{x \in
X}\text{Frac}(\mathcal{O}_{X,x})
$$
which associates to any $x\in X$
an element $f_x \in
\text{Frac}(\mathcal{O}_{X,x})$
such that for any $x_0 \in X$
there exists $U\ni x$ open and $f,g
\in\mathcal{O}(U)$
such that $f_x=g/h$ for all $x \in U$.
The {\it sheaf of meromorphic functions} is
$\mathcal{K}_X$.
\end{definition}


\medskip\noindent
Now we pass to the theory of analytic
subvarieties.

\begin{definition}
\label{complex-geometry-definition-analytic-subvariety}
An {\it analytic set} or {\it subvariety} $Y$
is a subset of a complex manifold $X$
given locally as the zeros of a finite
collection of holomorphic functions.
\end{definition}

\begin{lemma}
\label{complex-geometry-lemma-structure-sheaf-of-subvariety}
Let $Y \subset X$ be an analytic subvariety
of a complex manifold $X$
equipped with the subspace topology.
Then
$$
\begin{aligned}
\mathcal{O}_Y: \mathit{Open}_Y^{\text{op}}
&\longrightarrow \mathit{Set} \\
U\cap Y &\longmapsto \{f|_{U \cap Y}:f \in
\mathcal{O}_X(U)\}
\end{aligned}
$$

\noindent
is a sheaf.
\end{lemma}

\begin{proof}
This is a presheaf for the same reasons as
the sheaf $\mathcal{O}_X$ is a presheaf.
The fact that it is a sheaf follows
from the fact that $\mathcal{O}_X$ is a
sheaf and by definition of subspace topology;
namely for any family of functions
defined over an open cover $\{U_i \cap Y\}$
of some open set  $U\cap Y$ we can
reconstruct a function
on $U$ that restricted to $U \cap Y$ is as
desired.
\end{proof}

Recall that any map of sheaves of groups
has well-defined kernel, see Lemma
\ref{algebraic-geometry-lemma-sheaves-%
valued-on-groups-have-kernels}.

\begin{definition}
\label{complex-geometry-definition-ideal-sheaf}
Let $Y \subset X$ be an analytic subvariety
of a complex manifold $X$.
The {\it ideal sheaf} of $Y$, denoted
$\mathcal{I}_Y$,
is the kernel
of the restriction map $\mathcal{O}_X \to
\mathcal{O}_Y$.
\end{definition}

\noindent
This just says that $\mathcal{I}_X(U)$ is the
set of functions $f \in \mathcal{O}_X$ that
vanish under the restriction map, i.e.
that vanish at points of $Y$.

Since the restriction map is surjective,
we have the following short exact sequence:
$$
\xymatrix{
0\ar[r]&\mathcal{I}_Y\ar[r]&\mathcal{O}
_X\ar[r]&\mathcal{O}_Y\ar[r]&0
}
$$

\medskip\noindent
Notice that for any open set $U\subset X$,
$\mathcal{I}_Y(U)$ is an ideal of
$\mathcal{O}_X(U)$. (Indeed,
if a function vanishes at $Y$ then its
product
with any other function will also vanish.)

\begin{definition}
\label{complex-geometry-definition-irreducible-variety}
An analytic variety $V\subset U\subset
\mathbb{C}^n$ is {\it irreducible} if
$V$ cannot be written as the union of two
distinct analytic varieties
$V_1,V_2\subset U$, both distinct to $V$.
\end{definition}

\medskip\noindent
We will show that for the case of an
irreducible hypersurface
the ideal sheaf is locally generated by a
single element.
This will allow us to construct a line bundle
from the defining functions
of the hypersurface, which is in general not
possible
for an arbitrary smooth hypersurface,
--- although it suffices to ask for
orientability in the
smooth case.

We need two ingredients:
(1) the fact that the stalks $\mathcal{O}_p$
of the
sheaf of holomorphic functions are UFDs,
and (2) the following lemma.

\begin{lemma}
\label{complex-geometry-lemma-ideal-irred-hyper-heigh-1}
The ideal
(correction: this should be the stalk)
of an irreducible analytic hypersurface is a
height-1 prime ideal.
\end{lemma}

\begin{proof}
We already know that irreducible varieties
correspond with prime ideals.

To show the ideal $I$ defining the variety
$Y$ has height 1
suppose that $0 \subset \mathfrak{p}
\subseteq I$ is prime.
Then $V(\mathfrak{p})$ is an analytic variety
that contains $Y$.
At regular points of both $V(\mathfrak{p})$
and $Y$,
both are smooth manifolds, but since $Y$ is
of codimension 1 and
$V(\mathfrak{p})$ does not equal all of  $Y$,
we conclude that they coincide.
Thus, in a neighbourhood
of a regular point we have
$\mathfrak{p}=I(V(\mathfrak{p}))=I$ by the
Nullstellensatz.
\end{proof}

\begin{lemma}
\label{complex-geometry-lemma-stalk-ideal-sheaf-princ-hyper}
The stalks of the ideal sheaf of
an irreducible hypersurface
are principal rings.
\end{lemma}

\begin{proof}
Since $\mathcal{O}_p$
is a UFD, we pick any element
$f \in \mathcal{I}_p$, and factor it as a
product of irreducibles.
Since $\mathcal{I}$ is prime, some of the
irreducible factors,
say $f_i$, is in $\mathcal{I}_p$.
The ideal generated by $f_i$ is prime since
$\mathcal{O}_p$
js a UFD (see
\href{https://stacks.math.columbia.edu/tag/%
034T}{034T}).
But by Lemma
\ref{complex-geometry-lemma-ideal-irred-hyper-heigh-1},
$\mathcal{I}_p$ is a height-1 prime ideal, so
the prime ideal $(f_i)$
must equal all of $\mathcal{I}_p$.
\end{proof}

\begin{lemma}
\label{complex-geometry-lemma-ideal-sheaf-irred-hyper-line}
The ideal sheaf $\mathcal{I}$
of the irreducible hypersurface $Y \subset X$
corresponds to a unique line bundle.
\end{lemma}

\begin{proof}
Let $p \in U$ and consider for any stalk
$\mathcal{I}_p$
its generator $[f_p]$,
which has a representative $(f_p,U_p)$.
These form an open cover of $X$,
coincide up to a unit in the intersections as
follows.
These units shall give us a cocycle.

For any $U_i\cap U_j \neq 0$
we have
[This is slightly imprecise.
We may extend $f_j|_{U_i\cap U_j}$ to $U_i$
and $f_i|_{U_i\cap U_j}$ to $U_j$,
and then restrict back to $U_i\cap U_j$.
This works because we can extend holomorphic
functions by their local power series
expansion!],
$f_j\in (f_i)$, so there exist functions
$g_{ij}$ such that
$$
f_i=g_{ij}f_j.
$$
These satisfy the cocycle condition since
$$
f_if_jf_k=g_{ij}f_jg_{jk}f_kg_{ki}f_i.
$$
By Proposition
\ref{complex-geometry-proposition-vecto-bundl-cocyc} we have a
line
bundle.
\end{proof}

\noindent
Finally we have concluded that the ideal
sheaf
of a hypersurface is a line bundle. Now we
can take its dual to define:

\begin{definition}
\label{complex-geometry-definition-sheaf-assoc-effec-carti}
Let $D:=S$ be an effective Cartier divisor,
that is,
an analytic hypersurface of a complex
manifold.
The line bundle $\mathcal{O}(D)$ is the dual
line bundle
of the ideal sheaf of $D$.
\end{definition}

\section{Chow's theorem}
\label{complex-geometry-section-chow-theorem}

\begin{theorem}[Chow]
\label{complex-geometry-theorem-chow}
Any analytic subvariety of $\mathbb{C}P^{n}$
is algebraic.
\end{theorem}

\begin{proof}
First we prove it for hypersurfaces.
A hypersurface in $\mathbb{P}^n$ defines a
line bundle
along with a section vanishing on the
hypersurface.
Such a line bundle must be $\mathcal{O}(d)$
for some $d$
by Exercise \ref{complex-geometry-exercise-pic-of-pn-is-z}.
Sections of $\mathcal{O}(d)$ are homogeneous
degree-$d$ polynomials
by Theorem
\ref{complex-geometry-theorem-secti-od-degre-d-homog-poly}.

For higher codimension the challenge is find
adequate planes
in $\mathbb{P}^n$, so that we can project
$V$, $\dim V=k$
to $\mathbb{P}^{k+1}$ and apply the
hypersurface version.
Pick any point $p\in \mathbb{P}^n$ not in
$V$.
We can choose a $(n-k-1)$-plane not
intersecting $V$ and containing $p$.
Inside such $\mathbb{P}^{n-k-1}$ we can pick
a $\mathbb{P}^{n-k-2}$
not touching $p$.
Now pick a complementary $\mathbb{P}^{k+1}$
in $\mathbb{P}^n$,
i.e. $\mathbb{P}^n=\mathbb{P}^{k+1}\oplus
\mathbb{P}^{n-k-2}$.

Now project $V$ to $\mathbb{P}^{k+1}$ and
apply Chow for hypersurfaces.
We obtain a polynomial $F$ defining the
projection of $V$.
Extend $F$ to $\tilde{F}$ on all
$\mathbb{P}^n$ by the rule
$\tilde{F}(x_0,…,x_n)=F(x_0,…,x_{k+1})$.
Note that since $p$ is not in $V$,
$\tilde{F}$ does not vanish in $p$.
Thus $V$ is the intersection of all such
polynomials $F$,
one for each $p$, which ensures that no
points other than those
of $V$ are in the common zeroes.
\end{proof}

\section{Kodaira Embedding theorem}
\label{complex-geometry-section-kodaira-embedding}

\noindent
We introduced line systems in Section
\ref{complex-geometry-section-line-bundles}.
Let $L$ be a holomorphic bundle on a complex
manifold $X$.
A point $x \in X$ is
a {\it base point} of $L$ if $s(x)=0$ for all
$s\in H^{0}(X,L)$
and the {\it base locus} $\text{Bs}(L)$
is the set of all base points of $L$.

Let $s_0,\ldots,s_N\in H^{0}(X,L)$ are a
basis
for some subspace $E \subset H^0(X,L)$.
Define a map
$$
\begin{aligned}
\varphi_L: X\setminus\text{Bs}(L)
&\longrightarrow \mathbb{P}^N \\
x &\longmapsto [s_0(x):\ldots:s_N(x)]
\end{aligned}
$$

\noindent
which is well-defined since we excluded
the base locus,
and holomorphic since we chose holomorphic
sections of $L$ for the coordinates.
Further, the pullback of the hyperplane
bundle of $\mathbb{P}^N$ is $L$,
i.e.
$$
\varphi^*_L \mathcal{O}_{\mathbb{P}^N}(1)
\simeq L|_{X\setminus\text{Bs}(L)}
$$
since $s_i=0$ is the equation of a hyperplane
in $\mathbb{P}^N$,
so, it is a section of $H$.
Its pullback under $\varphi$ gives the
section $s_i$
whose associated divisor $(s_i)$
corresponds to the line bundle $L$ ---
the line bundle associated to a divisor
associated to a section of the line bundle
is the line bundle itself.

\medskip\noindent
The question is: when is $\varphi$ an
embedding?

\begin{definition}
\label{complex-geometry-definition-ample}
\begin{reference}
\cite[Definition 2.3.28]{huc}
\end{reference}
A line bundle $L$ on a complex manifold is
called {\it ample} if for some $k>0$
and some linear system in $H^{0}(X,L^k)$ the
associated map $\varphi$ is an
embedding.
\end{definition}

\begin{theorem}[Kodaira Embedding]
\label{complex-geometry-theorem-Kodaira-embedding}
\begin{reference}
\cite[Theorem 10.12]{lec}
\end{reference}
Suppose $M$ is a compact complex manifold.
A holomorphic line bundle $L\to M$ is
ample if and only if it is positive.
Thus $M$ is projective if and only if it
admits a positive holomorphic line bundle.
\end{theorem}

\begin{proof}
We must show that $\varphi$ separates points
and has
surjective differential everywhere…
\end{proof}

\noindent
Compare this theorem with Nakai-Moishezon
Criterion
\ref{algebraic-geometry-theorem-%
Nakai-Moishezon-criterion}.

\begin{exercise}
\label{complex-geometry-exercise-l-ample-impli-lotim-globa}
Let $L$ be an ample bundle on a K3 surface
$M$.
Prove that $L^{\otimes 2}$ is
globally generated
(that is, for each $x \in M$ there exists a
section $h \in
H^{0}(L^{\otimes 2})$ which does not vanish
in $x$).
\end{exercise}

\begin{proof}
\begin{reference}
\href{https://math.stackexchange.com/%
questions/5081481/
ample-line-bundle-on-a-k3-surface-implies-l-
otimes-2-is-base-point-free/
5081492?noredirect=1#comment10928708_5081492}
{StackExchange}
\end{reference}
As an outline for exposition:
\begin{enumerate}
\item Show that $L$ has sections
(by ampleness+Kodaira Vanishing+Riemann-Roch
for surfaces),
pick a section $C$ and note that if
$L^{\otimes2}$ had base points they would be
in $C$.

\item  Using Kodaira Vanishing again we show
that every section of the line bundle
restricted to $C$
comes from a section of the line bundle on
$M$.
Thus it's enough to find a nonvanishing
section on the restricted bundle.

\item Then we aim to apply Riemann-Roch on
$C$.
If we can show that $h^0((L^{\otimes 2})^\vee
\otimes \omega_C)$
and $h^0((L^{\otimes 2}(-p))^\vee \otimes
\omega_C)$
both vanish, we apply Riemann-Roch and see
that
$h^0(L^{\otimes 2})> h^0(L^{\otimes 2}(-p)$,
which proves the result since that says
that not every section is a section vanishing
at $p \in C$.

\item To compute the vanishings of the $h^0$
in the previous step we show they have
negative degrees
(genus formula + degree of canonical bundle
of a curve).
\end{enumerate}

First we need to show that $L$ has sections.
To see this I will show that since $L$ is
ample, $L^2>0$. Indeed, if $mL$ is very ample
we apply by Riemann-Roch, Kodaira Vanishing
and the fact that $\chi(\mathcal{O}_X)=2$ on
a K3 to obtain $h^0(mL)=2+\frac{1}{2}m^2L^2$
which must be positive since $mL$ has
sections by being very ample. Notice $L^2$
cannot be zero as $h^0(mL)=2$ would imply we
have an embedding of a K3 surface into
$\mathbb{P}^1$.

Applying again Riemann-Roch and Kodaira
vanishing we see that $h^0(L)>0$. Let $C$ be
a section of $L$.

Notice that if $L^{\otimes 2}$ had any base
points, they would have to be on $C$. Indeed,
since $C$ is the vanishing set of a section
$s\in H^{0}(L)$, we must have $s \otimes s
\in H^{0}(L^{\otimes 2})$ vanishing in any
base point of $L^{\otimes 2}$.

Then we consider the ideal exact sequence for
$C$ and tensor by $L^{\otimes 2}$ (which
preserves left-exactness because line bundles
are locally free of rank 1, that is, we are
tensoring by the underlying ring) to obtain
$$ \xymatrix{ 0\ar[r]&L^{\otimes
2}(-C)\ar[r]&L^{\otimes 2}\ar[r]& L^{\otimes
2}\otimes\mathcal{O}_C\ar[r]&0 } $$ Notice
that the term on the left is actually $L$
since the ideal sheaf
$\mathcal{O}_M(-C)=\mathcal{I}_C$ is dual to
$L$ because $C$ is defined as the vanishing
set of a section of $L$. Passing to
cohomology we obtain $$ \xymatrix{
H^{0}(L^{\otimes 2})\ar[r]&H^{0}(L^{\otimes
2}\otimes \mathcal{O}_C)\ar[r]& H^{1}(L) } $$
where the latter term vanishes by Kodaira
Vanishing because $L$ is ample. This says
that every section of $L^{\otimes 2}|_C$ is
the restriction to $C$ of a section of
$L^{\otimes 2}$. In turn this shows that it's
enough to show that the bundle $L^{\otimes
2}$ has no base points along $C$.

Now pick a point $p\in C$. Saying
that $p$ is not a base point of $L^{\otimes
2}|_{C}$ is the same as saying that not every
section of $L^{\otimes 2}|_{C}$ vanishes at
$p$, that is, that $h^0(L^{\otimes
2}|_{C}(-p))<h^0(L^{\otimes 2}|_{C})$,
which concludes the exercise.

The key step to prove
the last inequality is showing that
$h^0(\omega_C\otimes(L^{\otimes2}|_{C})^\vee)
=0=h^0(\omega_C\otimes
(L^{\otimes2}|_{C}(-p))^\vee)$. For this we
need to see these line bundles as Weil
divisors, so that their degree is the sum of
their multiplicities. First note that by the
genus formula for (possibly singular) curves
on smooth surfaces and by the fact that $M$
is smooth, we see that $2p_a(C)-2=(L.L)$ so
that
$$
4p_a(C)-4=(2L.L)=\text{deg}_C(L^{\otimes 2})
$$
Notice that since $L^2>0$ we exclude the
cases that $p_a(C)=0,1$.

By Riemann-Roch on the curve
$C$, we know that
$\text{deg}\omega_C=2p_a-2$. Then
$$
\text{deg}_C
(\omega_C\otimes (L^{\otimes2})^\vee)
=\text{deg}_C(\omega_C)-\text{deg}
_C(L^{\otimes
2})
=2p_a-2-(4p_a-4)<0
$$
and a similar computation works for
$L^{\otimes 2}(-p)$ as it has the same
degree of $L^{\otimes 2}$ minus 1. This
implies that neither of these bundles
can have sections since any section would
provide a linearly equivalent
effective divisor
--- degree is preserved under linear
equivalence,
and divisors given by sections are effective.
Therefore
$h^0(\omega_C\otimes
(L^{\otimes2}|_{C})^\vee)=0
=h^0(\omega_C\otimes (L^{\otimes
2}|_{C}(-p))^\vee)$ as claimed.

Finally we apply Riemann-Roch to $L^{\otimes
2}$ and $L^{\otimes 2}(-p)$ to
obtain that $h^0(L^{\otimes
2})>h^0(L^{\otimes 2}(-p))$.
\end{proof}

\section{Riemann-Roch formulas}
\label{complex-geometry-section-riemann-roch}

\noindent
We want to know the dimension of
$H^0(X,\mathcal{O}(D))$,
where $X$ is a curve and $D$ is a divisor on
$X$.
To refresh out memory recall that if $D$ is a
smooth hypersurface
(in this case this just means that it's a set
of points
without multiplicities)
then that $\mathcal{O}(D)$ is given by the
cocycle
$g_{ij}=f_j/f_i$ where $D=\text{div}(f_i)$ on
$U_i$.
The $g_{ij}$ give a cocycle, which in turn is
a line bundle on $X$.
If $e_i$ is a local basis then we have
$e_i=g_{ij}e_j$
by definition of transition function.
A section is given locally by $h_ie_i$ and it
has to glue,
that is, $h_je_j=h_ie_i=e_if_j/f_ie_j$
so $h_j/f_j=h_i/f_i$.
That's literally the data of a meromorphic
function,
namely some local meromorphic functions which
coincide in intersections.

What's the vanishing of this meromorphic
function along $D$?
Now suppose $D$ is an arbitrary Weil divisor,
so the $f_i$ have some multiplicities.
If $h_i$ has no zeroes, then the multiplicity
of the new
meromorphic function remains unaltered, but
$h_i$ can also vanish
along $D$ and then the vanishing becomes as
positive as we want it.
But it cannot have poles, the least positive
it can be is
the prescribed negativity of $D$.

This is why looking for $h^0(D)$ is just as
much as looking for
meromorphic functions $h$ on $X$ such that
$$
\text{div}h + D \geq 0\qquad \text{(is
effective).}
$$



\begin{theorem}[Riemann-Roch for curves]
\label{complex-geometry-theorem-Riemann-Roch-for-curves}
Let $D$ be a divisor on a compact Riemann
surface $X$, that is, $D$ is a
collection of points with multiplicities on
$X$.
\begin{equation}
\label{complex-geometry-equation-riema-roch-riema-surfa}
h^0(D)-h^0(K-D)=d-g+1
\end{equation}
\end{theorem}

\begin{exercise}
\label{complex-geometry-exercise-no-holom-vecto-field-g}
Prove that if $X$ is a complex smooth curve
with $g\geq 2$ then $h^0(T_X)=0$.
\end{exercise}

\begin{proof}
First you notice that the dimension of the
holomorphic 1-forms is the genus.
This is just because the genus $p_a$ is
defined as $h^1(\mathcal{O}_X)$, which
is just $h^0(K_X)$ by Serre duality. Then you
say well if you have a holomorphic
vector field, pair it with the nonzero
holomorphic 1-form. This gives a nonzero
function, but it must be constant because $X$
is compact. Then it is actually
nowhere vanishing. This says the vector field
cannot vanish anywhere, which
means the tangent bundle of the curve is
trivial. Apparently elliptic curves,
i.e. $g=1$ are the
\end{proof}

\begin{theorem}[Riemann-Roch for line bundles
on surfaces]
\label{complex-geometry-theorem-riema-roch-line-bundl-surfa}
Let $L$ be a line bundle on a complex surface
$X$, and $K_X=\Omega^2X$ the
canonical bundle of $X$. Then
\begin{equation}
\label{complex-geometry-equation-riema-roch-line-bundl-surf}
\chi(L)=\chi(\mathcal{O}_X)+\frac{1}{2}(L,L-
K_X)
\end{equation}
\end{theorem}

\begin{lemma}[Genus formula]
\label{complex-geometry-lemma-genus-formula}
Let $C$ be any curve on a surface $S$. Then
\begin{equation}
\label{complex-geometry-equation-genus-formula}
2p_a(C)-2=(C.C+K_S).
\end{equation}
\end{lemma}

\begin{proof}
By Remark
\ref{complex-geometry-remark-euler-characteristic-additive}
below we can take Euler characteristic
of the ideal sheaf exact sequence of $C$
to obtain
$$
\chi(\mathcal{O}(-C))+\chi(\mathcal{O}_C)=
\chi(\mathcal{O}_S).
$$
Since $p_a(C)=h^0(K_S)
=1-\chi(\mathcal{O}_C)$, we
obtain that
$p_a(C)=1-\chi(\mathcal{O}_S)-
\chi(\mathcal{O}(-C))$.
To compute $\chi(\mathcal{O}(-C))$ we use
Riemann-Roch formula for curves on surfaces
to get
$$
\chi(\mathcal{O}(-C))=\chi(\mathcal{O}_S)
+\frac{(\mathcal{O}(-C).
\mathcal{O}(-C)-K_S)}{2}
$$
Taking duals on both entries of the
intersection form and substituting in our
previous expression for $p_a(C)$ we obtain
the result.
\end{proof}

\begin{remark}
\label{complex-geometry-remark-euler-characteristic-additive}
To prove that Euler characteristic is
additive, i.e. that for an exact sequence
$$
\xymatrix{
0\ar[r]
&E\ar[r]
&F\ar[r]
&G\ar[r]
&0
}
$$
we must have
$$
\chi(F)=\chi(E)+\chi(G),
$$
we take the long exact sequence in cohomology
and use the general fact that for
a finite exact sequence of vector spaces
$$
\xymatrix{
0\ar[r]
&V_0\ar[r]
&V_1\ar[r]
&\cdots\ar[r]
&V_n\ar[r]
&0
}
$$
we must have
$$
\sum_i(-1)^i\dim V_i,
$$
which basically follows by dimension theorem.
\end{remark}

\begin{exercise}
\label{complex-geometry-exercise-singular-curve-on-k3}
Let $C$ be a singular, irreducible complex
curve on a K3 surface $X$. Prove that
$(C.C)\geq 0$.
\end{exercise}

\begin{proof}
This is just an application of Eq.
\ref{complex-geometry-equation-genus-formula}. Since
$K_X=\mathcal{O}_X$ it suffices to show that
the arithmetic genus
$p_a(C)$ is not zero. This follows from the
fact that $C$ is singular, since its
arithmetic genus, defined as 
$p_a(C):=1-\chi(\mathcal{O}_C)$,
satisfies
$$
p_a(C)=g(C) + h^0(\delta)
$$
where $g(C)$ is the geometric genus of $C$,
defined as
the genus of the normalization $\tilde{C}$,
i.e. $g(C):=g(\tilde{C})$,
while $\delta:=\nu_*\mathcal{O}_{\tilde{C}}
/\mathcal{O}_C$
where $\nu:\tilde{C} \to C$ is the
normalization map.

So we have to prove that equation
and also prove that $\delta$
always has sections…
\end{proof}

\section{Multi-valued functions and branched covers}
\label{complex-geometry-section-multi-valued-functions-and-branched-covers}

\noindent
The objective of this section is to
show that every algebraic variety
over $\mathbb{C}$ is smooth
away from a divisor.
That is, $X_{\text{sing}}$
is a Zariski-closed set, an algebraic
subvariety.
In particular, this gives a stratification
$X\supset X^1_{\text{sing}}
\supset X^2_{\text{sing}}…$.

\medskip\noindent
An elliptic function is either a square root
of a polynomial, or its antiderivative. This
is not a well-defined global function. But
locally it does make sense. This is what
Riemann was working on when he defined
manifolds.

\begin{definition}
\label{complex-geometry-definition-multi-valued-function}
{\it 
The graph of 
a multi-valued function} on $M$ is a closed,
irreducible complex subvariety
$Z \subset M \times \mathbb{C}$ such that
the projection $Z \to M$ is locally a
diffeomorphism outside of a closed, nowhere
dense subset in $Z$.
\end{definition}

\noindent
This is supposed to be reminiscent
of Noether's Normalization lemma!!!

\begin{example}
The logarithm is a multi-valued function on
$\mathbb{C}$. Indeed, let $Z$ be the graph of
the exponent $y=e^x$ in $\mathbb{C}^2$. The
projection to $x$ expresses all branches of
logarithm $x=\log y$ as functions of $y$.

The point is that the logarithm
is a local diffeomorphism, it can be
locally inverted, but not globally.
\end{example}

\begin{example}
$y \mapsto \sqrt{y}$ is a multi-valued
function. Indeed, the graph of $y=x^2$
projected to $x$ gives both branches of
$\sqrt{y}$.
\end{example}

\noindent
The following theorem introduces
finite coverings. Outside of the so-called
branching locus they are a covering.
Unfortunately I can't follow the proof
(perhaps it has mistakes).
I think it's just application of inverse
function theorem to a hypersurface
of $\mathbb{C}^n$.

\begin{theorem}
\label{complex-geometry-theorem-monic-polynomial-multivalued}
Consider a subvariety $Z \subset
\mathbb{C}^{n+1}$ given by a monic
polynomial equation $P(t)=0$, with
$P(t)\in \mathcal{O}_{\mathbb{C}^n}[t]$.
Assume that $P(t)$ is irreducible. Then $Z$
is a graph of a multi-valued function.
Moreover, the projection of $Z$ to
$\mathbb{C}^n$ (to the first $n$
coordinates) is a diffeomorphism at
$(z,t)$ if and only if $P'(z)\neq 0$.
\end{theorem}

\begin{proof}
Step 1: We shall represent points of
$\mathbb{C}^{n+1}$ by pairs $(z,t)$, with
$z=(z_1,\dots,z_n)$. Let $\pi:Z\to \mathbb{C}^n$
be the standard projection along $t$. By
the inverse function theorem, $Z \subset
\mathbb{C}^{n+1}$ is a smooth submanifold in
a neighbourhood of any point $(z,t)\in Z$
with $dP(z,t)\neq 0$.

Step 2: This implies that $Z$ is smooth
outside of the set of all $(z,t)\in
\mathbb{C}^{n+1}$ such that
$dP(z,t)(\xi,\tau)=0$ for all
$\xi \in \langle d/dz_1,\dots,d/dz_n\rangle$
and $\tau \in \langle d/dt\rangle$.
Let
$$
P(z,t)=t^k+\sum_{i=0}^{k-1} t^i a_i(z).
$$
Then
$$
dP(z,t)=kt^{k-1}dt
+ \sum_{i=0}^{k-1} t^i da_i(z)
+ \sum_{i=1}^{k-1} i t^{i-1} a_i(z) dt.
$$
For $|t|\gg 0$, the leading term
$kt^{k-1}dt+t^{k-1}da_{k-1}(z)$ dominates
the rest, and it is non-zero because its
$dt$ component is non-zero.

Step 3: Let $z \in Z$ be a smooth point,
and $(\xi,\tau)\in T_zZ$. Then
$\pi:W \to \mathbb{C}^n$ is invertible
whenever $W$ does not contain a vector
$(0,\tau)$, equivalently, when
$dP(z,t)(0,\tau)\neq 0$. This is equivalent
to
$$
\frac{dP(z,t)}{dt}\neq 0.
$$

Step 4: Let
$P'(z,t):=\frac{dP(z,t)}{dt}$.
To prove that $Z$ defines a multi-valued
function, it remains to show that
$P'(z,t)$ is not identically zero on $Z$.
Since $P(z,t)$ is irreducible, and
$\mathcal{O}_{\mathbb{C}^{n+1}}$ is
factorial, the ring
$\mathcal{O}_{\mathbb{C}^{n+1}}/(P)$ has no
zero divisors. Then the Hilbert
Nullstellensatz implies that any function
$f \in \mathcal{O}_{\mathbb{C}^{n+1}}$
which vanishes on $Z$ is divisible by
$P(z,t)$. Therefore $P'(z,t)$ does not
vanish on $Z$, because it is a polynomial
of smaller degree.
\end{proof}

\medskip\noindent
Now we make a detour
to discuss
symmetric polynomials and discriminant.

\begin{definition}
\label{complex-geometry-definition-symmetric-polynomial}
A symmetric polynomial
$P(z_1,\dots,z_n)\in \mathbb{C}[z_1,\dots,z_n]$
is a polynomial which is invariant under the
symmetric group $\Sigma_n$ acting on
$\mathbb{C}[z_1,\dots,z_n]$ in the usual
way.
\end{definition}

\begin{definition}
\label{complex-geometry-definition-elementary-symmetric-polynomial}
Consider the polynomial
$$
P(z_1,\dots,z_n,t):=\prod_{i=1}^n(t+z_i)
=\sum e_it^i,
$$
with $e_i \in \mathbb{C}[z_1,\dots,z_n]$.
Then $e_i$ are called elementary symmetric
polynomials on $z_1,\dots,z_n$.
\end{definition}

\begin{theorem}
\label{complex-geometry-theorem-symmetric-polynomial-generated}
Every symmetric polynomial on
$z_1,\dots,z_n$ can be polynomially
expressed through the elementary symmetric
polynomials.
\end{theorem}

\medskip\noindent
Proof left as an exercise.

\begin{definition}
\label{complex-geometry-definition-discriminant}
Consider the symmetric polynomial
$$
\prod_{i\neq j}(z_i-z_j).
$$
Discriminant of the polynomial
$$
P(z_1,\dots,z_n,t):=\prod_{i=1}^n(t-z_i)
$$
is $\prod_{i\neq j}(z_i-z_j)$ considered as
a polynomial of its coefficients.
\end{definition}

\begin{example}
The discriminant of the quadratic polynomial
$t^2+bt+c$ is $b^2-4c$.
\end{example}

\begin{example}
The discriminant of the cubic polynomial
$t^3+bt^2+ct+d$ is
$$
b^2c^2-4c^3-4b^3d-27d^2+18bcd.
$$
\end{example}

\begin{lemma}
\label{complex-geometry-lemma-polynomial-no-multiple-roots}
A polynomial has no multiple roots if and
only if its discriminant is non-zero.
\end{lemma}

\begin{lemma}[Corollary]
\label{complex-geometry-lemma-derivative-on-roots-and-discriminant}
Let $P(t)\in k[t]$ be a polynomial over an
algebraically closed field, and $D$ its
discriminant. Then the derivative $P'(t)$
does not vanish on all roots of $t$ if and
only if $D\neq 0$.
\end{lemma}

\begin{proof}
Let $\alpha$ be a root of $P$. Then
$P(t)=P_\alpha(t)(t-\alpha)$, and
$P'(t)=P'_\alpha(t)(t-\alpha)+P_\alpha(t)$.
Therefore $P'(\alpha)=0$ if and only if
$P_\alpha(t)=0$. This is equivalent to
$P(t)$ being divisible by $(t-\alpha)^2$.
\end{proof}

\medskip\noindent
We already knew much about this map,
but the following theorem guarantees
$\pi|_Z$ is a local diffeomorphism
away from the branch locus.
Recall: $P$ is a degree-$d$ polynomial,
so every point has $d$ preimages,
each of which is a root.

\begin{theorem}
\label{complex-geometry-theorem-discriminant-ramified-covering}
Consider a subvariety $Z \subset
\mathbb{C}^{n+1}$ given by an irreducible
monic polynomial equation $P(t)=0$ of
degree $k$, with
$P(t)\in \mathbb{C}[z_1,\dots,z_n][t]$,
and let $\pi:Z \to \mathbb{C}^n$ be the
projection to the coordinates
$z_1,\dots,z_n$. Assume that $P(t)$ is
irreducible. Denote by $D(z)$ the
discriminant of $P(t)$, considered as a
polynomial function on $(z_1,\dots,z_n)$,
and let $U \subset \mathbb{C}^n$ be the set
of all $z \in \mathbb{C}^n$ such that
$D(z)\neq 0$. Then the intersection
$Z \cap \pi^{-1}(U)$ is smooth and the
projection
$\pi:Z\cap \pi^{-1}(U)\to U$ is a $d$-sheeted
covering.
\end{theorem}

\begin{proof}
Step 1: Since $P(t)$ is irreducible, it has
no common divisors with $P'(t)$. Since the
polynomial ring is factorial,
irreducibility implies that the ideal
$(P(t))$ is prime, hence $P'(t)$ does not
identically vanish on $Z$. Therefore $U$ is
open and nowhere dense. By Corollary 1, for
any $z \in U$, the polynomials $P(z,t)$ and
$P'(z,t)$ have no common roots. Therefore
$dP(z,t)/dt\neq 0$ on
$Z \cap \pi^{-1}(U)$, and the set
$Z=\{(z,t)\mid P(z,t)=0\}$ is smooth outside
of zeros of $D(z)$.

Step 2: Let $z \in Z\cap \pi^{-1}(U)$ and
$(\xi,\tau)\in T_{(z,t)}Z$. Then
$\pi:T_{(z,t)}Z\to \mathbb{C}^n$ is
invertible whenever $T_{(z,t)}$ does not
contain a vector $(0,\tau)$, equivalently,
when $dP(z,t)(0,\tau)\neq 0$. This is
equivalent to
$dP(z,t)/dt\neq 0$.

Step 3: The map
$\pi:Z\cap\pi^{-1}(U)\to U$ is locally a
diffeomorphism, and each point has precisely
$d$ preimages. Therefore, it is a covering.
\end{proof}

\noindent
Note that the theorem only applies
to hypersurfaces --- and not every
variety can be realised as a hypersurface.

\begin{lemma}[Corollary]
\label{complex-geometry-lemma-birational-finite-map-to-covering}
Let $X$ be an algebraic variety. Then there
exists a birational, finite map
$X \to Z$, a divisor $D \subset Z$, and a
divisor $D_1 \subset \mathbb{C}^n$, such
that $Z \setminus D$ is a $d$-sheeted
covering of $\mathbb{C}^n \setminus D_1$.
\end{lemma}

\begin{lemma}[Corollary]
\label{complex-geometry-lemma-algebraic-variety-smooth-point}
Every algebraic variety $X$ over $\mathbb{C}$
has a smooth point. Moreover, non-smooth
points of $X$ are contained in a proper
algebraic subvariety of $X$.
\end{lemma}

\begin{proof}
Indeed, every birational map is an
isomorphism outside of a divisor.
\end{proof}


\section{Symplectic isogenies and
Shafarevich}
\label{complex-geometry-section-symplectic-isogenies-shaf}

These are notes from Misha Verbitsky's talk
{\it Symplectic isogenies for hyperkahler
manifolds and Shafarevich conjecture},
given at IMPA on May 7, 2026.

\medskip
Abstract.
A correspondence between hyperkahler
manifolds is a symplectic isogeny if it is
generically finite over both factors and its
push-pull operation sends the holomorphic
symplectic form to a nonzero holomorphic
symplectic form.
The conjecture, attributed to Shafarevich by
Buskin, says that hyperkahler manifolds in
the same deformation family with isomorphic
rational period lattices should be related
by such an isogeny.
The talk explains degenerate twistor
families and shows that projective
hyperkahler manifolds in the same such
family are symplectically isogenous.

\medskip
Absolute upshot.
Same degenerate twistor family is a much
more structured condition than being merely
deformation-equivalent.
Misha proves the Shafarevich-type isogeny
statement first in this structured setting:
the common Lagrangian fibration gives torus
fibers and torsors, which can be trivialized
after finite base change.
The resulting graph produces the desired
symplectic correspondence.

\medskip
Upshot.
Find the simplest Hodge-theoretic piece
forced by the holomorphic symplectic form:
the transcendental lattice $V^{tr}$.
Then prove that this piece behaves well:
it is usually irreducible, and in the
degenerate case has only one controlled
defect.
See Theorem
\ref{complex-geometry-theorem-transcendental-k3-type}.

\medskip
{\bf Motivation: ADE contractions.}

Let $X$ be a smooth K3 surface and let
$C \subset X$ be a smooth rational curve.
Since $K_X=\mathcal{O}_X$, adjunction gives
$$
2g(C)-2=C^2+K_X\cdot C=C^2.
$$
Thus a smooth rational curve on a K3 surface
has self-intersection $C^2=-2$.
Conversely, a smooth curve with $C^2=-2$
has genus $0$.
Such a curve is called a $(-2)$-curve.

\medskip
A $(-2)$-curve is not obtained by blowing up
another smooth K3 surface.
Rather, it can be contracted:
$$
X \longrightarrow Y.
$$
The target $Y$ is usually singular.
For a single $(-2)$-curve, the image point is
a Du Val singularity of type $A_1$.
Locally this is the rational double point
$$
xy=z^2,
$$
or equivalently $\mathbb{C}^2/\{\pm 1\}$.

\medskip
More generally, suppose
$C_1,\ldots,C_r \subset X$ are $(-2)$-curves
whose intersection matrix is negative
definite.
Then the whole configuration can be
contracted.
The dual graph has one vertex for each
$C_i$ and an edge when two curves meet
transversely once.
For rational double points this graph is one
of the ADE Dynkin diagrams:
$$
A_n,\qquad D_n,\qquad E_6,E_7,E_8.
$$

\medskip
Conversely, the minimal resolution of a
Du Val singularity replaces the singular
point by a configuration of $(-2)$-curves.
Their intersection graph is the corresponding
Dynkin diagram.
Thus the slogan is
$$
\text{Du Val singularity}
\longleftrightarrow
\text{ADE configuration of }
(-2)\text{-curves}.
$$
The lattice generated by these curves is the
corresponding root lattice.

\medskip
In the motivating picture, one contracts a
Dynkin configuration
$$
f:X\longrightarrow Y
$$
and then resolves the singular surface $Y$:
$$
\widetilde{Y}\longrightarrow Y.
$$
If $X$ is not already the chosen resolution,
the common singular surface $Y$ gives a
relation between $X$ and $\widetilde{Y}$.
One studies this relation by the
correspondence
$$
X \xleftarrow{\ p_1\ }
X\times_Y\widetilde{Y}
\xrightarrow{\ p_2\ }\widetilde{Y}.
$$
This is the prototype for replacing a
birational map by a cycle in a product.

\medskip
This is the end of the motivating example.
The real subject of the talk is to replace
maps between hyperkahler manifolds by
correspondences which transport holomorphic
symplectic forms.

\medskip
{\bf Correspondences and symplectic
isogenies.}

\begin{definition}
\label{complex-geometry-definition-correspondence-verbitsky}
Let $M$ be a complex variety.
A {\it cycle} in $M$ is a subvariety
together with a positive integer
multiplicity on each irreducible component.
Let $A,B$ be irreducible compact complex
varieties of dimension $n$.
A {\it correspondence} from $A$ to $B$ is an
$n$-dimensional cycle
$$
Z\subset A\times B
$$
such that every irreducible component of
$Z$ projects surjectively to both $A$ and
$B$.
\end{definition}

\noindent
Thus a correspondence is a multivalued map.
When using cohomology, we often think of the
cycle $Z$ as its Poincare dual differential
form. The caveat is that pushforward of a
differential form is generally a current.
For holomorphic forms in this setting, the
pushforward is holomorphic away from the
exceptional locus and extends by Hartogs.

\begin{definition}
\label{complex-geometry-definition-symplectic-isogeny}
Let $A,B$ be holomorphically symplectic
manifolds and let
$Z\subset A\times B$ be a correspondence.
Let $\pi_1,\pi_2$ be the projections from
$Z$ to $A$ and $B$.
The correspondence $Z$ is a
{\it symplectic isogeny} if, for a
holomorphic symplectic form $\Omega_A$ on
$A$,
$$
(\pi_2)_*(\pi_1^*\Omega_A)\neq 0.
$$
\end{definition}

\begin{lemma}
\label{complex-geometry-lemma-sympl-isog-transcendental}
Let $M_1,M_2$ be holomorphically symplectic
manifolds of K3 type, and let
$Z\subset M_1\times M_2$ be a symplectic
isogeny.
Then the push-pull map
$$
\Phi_Z(\eta):=(\pi_2)_*(\pi_1^*\eta)
$$
induces an isomorphism of rational
transcendental Hodge lattices
$$
\Phi_Z:T(M_1)\longrightarrow T(M_2).
$$
\end{lemma}

\noindent
Here $T(M)$ means the transcendental Hodge
lattice $V^{tr}\subset H^2(M,\mathbb{Q})$.
This is not literally the complex torus
attached to a weight-one Hodge structure.
The torus story is the model case: isogenies
of tori are the same as isomorphisms of
rational weight-one Hodge structures.
For K3-type Hodge structures, the analogous
object is $T(M)$.

\medskip
Recall also the dictionary between complex
tori and Hodge structures; see
{\it mumford-tate-groups-in-hodge-theory.tex}
for the Hodge-theoretic background.
If $T$ is a compact complex torus, then
$H_1(T,\mathbb{Z})$ carries an integral
Hodge structure of weight $1$.
Conversely, such a Hodge structure
$$
V_{\mathbb{C}}=V^{1,0}\oplus V^{0,1}
$$
with a lattice $V_{\mathbb{Z}}\subset V$,
gives the complex torus
$$
T=V^{1,0}/V_{\mathbb{Z}}.
$$
After tensoring the lattice with
$\mathbb{Q}$, this is the category of
complex tori up to isogeny.

\begin{definition}
\label{complex-geometry-definition-isogeny-tori}
Let $T_1,T_2$ be compact complex tori of the
same dimension.
An {\it isogeny} from $T_1$ to $T_2$ is a
subtorus
$$
T\subset T_1\times T_2
$$
such that the two projections
$T\to T_1$ and $T\to T_2$ are coverings.
We say that $T_1$ and $T_2$ are isogenous
if such a subtorus exists.
\end{definition}

\medskip
Briefly, the Mumford-Tate group records the
rational algebraic symmetries forced by a
Hodge structure.
See Section
\ref{complex-geometry-mumford-tate-section-mumford-tate-group}
and Definition
\ref{mumford-tate-definition-mumford-tate-%
group}.
If $V$ is a rational Hodge structure and
$$
\rho:U(1)\longrightarrow
GL(V_{\mathbb{R}})
$$
is the Hodge circle action, then $MT(V)$ is
the smallest algebraic subgroup of $GL(V)$
defined over $\mathbb{Q}$ whose real points
contain $\rho(U(1))$.

\begin{proposition}
\label{complex-geometry-proposition-mt-generated-galois}
Let $V$ be a rational Hodge structure.
For each
$\sigma\in\operatorname{Aut}_{\mathbb{Q}}
(\mathbb{C})$,
let $U(1)^\sigma$ denote the Galois conjugate
of the Hodge circle in
$GL(V_{\mathbb{C}})$.
Then $MT(V)_{\mathbb{C}}$ is the Zariski
closure of the subgroup generated by all
$U(1)^\sigma$.
\end{proposition}

\begin{theorem}
\label{complex-geometry-theorem-transcendental-hodge-lattice}
Let $V_{\mathbb{Q}}$ be a rational Hodge
structure of weight $d$.
Let $V^{tr}\subset V_{\mathbb{Q}}$ be the
smallest rational subspace such that
$$
V^{d,0}\subset V^{tr}_{\mathbb{C}}.
$$
Equivalently, $V^{tr}_{\mathbb{C}}$ is
generated by all Galois conjugates of
$V^{d,0}$.
Then $V^{tr}$ is a Hodge substructure.
\end{theorem}

\begin{definition}
\label{complex-geometry-definition-polarization-hodge}
A {\it polarization} on a rational Hodge
structure $V$ of weight $k$ is a
$U(1)$-invariant nondegenerate bilinear form
$$
h\in V_{\mathbb{Q}}^*\otimes
V_{\mathbb{Q}}^*
$$
which is symmetric or skew-symmetric
depending on the parity of $k$, and which
satisfies the Hodge-Riemann positivity
condition.
\end{definition}

\noindent
This is Misha's formulation.
In
Definition
\ref{complex-geometry-mumford-tate-definition-polarization},
Andrey packages the same positivity using a
morphism to $\mathbb{Q}(-k)$ and the Weil
operator.

\medskip
The basic source of such polarizations is the
Hodge-Riemann pairing.
For a compact Kahler manifold $X$, a Kahler
class $\omega$ gives, on primitive
$H^k(X)$, the pairing
$$
(\alpha,\beta)\longmapsto
\int_X\alpha\wedge\beta\wedge\omega^{n-k}.
$$
After the usual sign and Weil-operator
normalization, the Hodge-Riemann bilinear
relations say this pairing is positive.
Thus the abstract polarization in Hodge
theory is the linear-algebra shadow of the
intersection pairing from complex geometry.

\begin{definition}
\label{complex-geometry-definition-simple-semisimple-object}
Let $\mathcal{A}$ be an abelian category.
An object $V$ is {\it simple} if its only
subobjects are $0$ and $V$.
It is {\it semisimple} if it is a finite
direct sum of simple objects.
\end{definition}

\noindent
The upshot is that polarizations prevent
extensions from hiding inside a Hodge
structure.
If $W\subset V$ is a Hodge substructure, the
orthogonal complement $W^\perp$ is again a
Hodge substructure, so $V$ splits as
$W\oplus W^\perp$.

\begin{proposition}
\label{complex-geometry-proposition-polarizable-hs-semisimple}
Every polarizable rational Hodge structure
is semisimple.
\end{proposition}

\noindent
In particular, the rational Hodge structures
$H^k(X,\mathbb{Q})$ of a compact Kahler
manifold are semisimple.
Indeed, the Lefschetz decomposition reduces
the statement to primitive cohomology, where
the Hodge-Riemann bilinear relations give
the required polarization.

\begin{definition}
\label{complex-geometry-definition-k3-type-hodge-structure}
A rational Hodge structure of weight $2$
is of {\it K3 type} if
$$
V_{\mathbb{C}}=
V^{2,0}\oplus V^{1,1}\oplus V^{0,2}
$$
and $\dim V^{2,0}=1$.
\end{definition}

\begin{example}
\label{complex-geometry-example-bbf-k3-type}
Let $M$ be an irreducible holomorphic
symplectic manifold.
Then $H^2(M,\mathbb{Q})$ is a Hodge
structure of K3 type.
It carries the Beauville-Bogomolov-Fujiki
quadratic form
$$
q:H^2(M,\mathbb{Q})\times H^2(M,\mathbb{Q})
\longrightarrow \mathbb{Q}.
$$
Fujiki's formula says that, if
$\dim_{\mathbb{C}}M=2n$, then
$$
\int_M\alpha^{2n}=c\,q(\alpha,\alpha)^n
$$
for a positive rational constant $c$.
\end{example}

\noindent
Thus, for hyperkahler manifolds, the BBF
form plays the role of the natural quadratic
form on the K3-type Hodge structure
$H^2(M)$.
It is not literally the ordinary
Hodge-Riemann pairing on all cohomology, but
it is the replacement on $H^2$ which controls
the period lattice and the Torelli theory.

\begin{theorem}
\label{complex-geometry-theorem-transcendental-k3-type}
Let $(V,q)$ be a pseudo-polarized rational
Hodge structure of K3 type.
Let $V^{tr}\subset V$ be the transcendental
Hodge sublattice, i.e. the smallest rational
Hodge substructure containing $V^{2,0}$.
Then exactly one of the following happens.
\begin{enumerate}
\item If $V^{1,1}_{\mathbb{Q}}$ contains a
vector $x$ with $q(x,x)>0$, then
$q|_{V^{tr}}$ is a polarization and
$V^{tr}$ is irreducible.
\item If $q|_{V^{1,1}_{\mathbb{Q}}}$ is
negative definite, then $V^{tr}$ is
irreducible and
$q|_{V^{tr}}$ is a pseudo-polarization.
\item If $q|_{V^{1,1}_{\mathbb{Q}}}$ is
negative semidefinite with one-dimensional
kernel, then
$$
V^{1,1}_{\mathbb{Q}}\cap
(V^{1,1}_{\mathbb{Q}})^\perp
=\langle\eta\rangle,
\qquad q(\eta,\eta)=0.
$$
In this case $\langle\eta\rangle\subset
V^{tr}$ is the only nonzero proper Hodge
substructure of $V^{tr}$, and
$V^{tr}/\langle\eta\rangle$ is irreducible.
\end{enumerate}
\end{theorem}

\noindent
The third case is the strange one.
The transcendental lattice is no longer
irreducible, but it fails in the smallest
possible way: it has a unique isotropic
rational line, and the quotient by this line
is irreducible.

\medskip
Upshot: for K3-type Hodge structures the
transcendental part is almost simple.
Usually it is irreducible.
In the degenerate case it has exactly one
visible defect, the isotropic line
$\langle\eta\rangle$.
Thus a correspondence which is nonzero on
$V^{2,0}$ is forced to control essentially
all of $V^{tr}$.
This is why the condition
$(\pi_2)_*\pi_1^*\Omega\neq 0$ is so strong.

\begin{lemma}
\label{complex-geometry-lemma-sympl-isog-controls-tm}
The point of the symplectic isogeny
condition is that the correspondence is
nonzero on the line $H^{2,0}$.
Since $T(M)$ is the smallest rational Hodge
substructure generated by this line, the
correspondence is forced to identify the
transcendental lattices:
$$
T(M_1)\cong T(M_2).
$$
\end{lemma}

\begin{theorem}[Witt]
\label{complex-geometry-theorem-witt-quadratic-spaces}
Let $(V,q)$ be a nondegenerate rational
quadratic space.
Let $W_1,W_2\subset V$ be subspaces such that
$q|_{W_i}$ is nondegenerate.
If $W_1$ and $W_2$ are isomorphic as
quadratic spaces, then
$W_1^\perp$ and $W_2^\perp$ are isomorphic
as quadratic spaces.
\end{theorem}

\begin{lemma}
\label{complex-geometry-lemma-isogeny-gives-h2-hodge-isom}
Let $M_1,M_2$ be projective hyperkahler
manifolds in the same deformation class.
If $M_1$ and $M_2$ are symplectically
isogenous, then
$$
H^2(M_1,\mathbb{Q})\cong
H^2(M_2,\mathbb{Q})
$$
as rational Hodge structures.
\end{lemma}

\noindent
Indeed, the isogeny identifies the
transcendental Hodge lattices $T(M_i)$.
The algebraic parts are the orthogonal
complements with respect to the BBF form.
Witt's theorem then gives the corresponding
isomorphism on the complements, hence on all
of $H^2$.

\begin{theorem}[Torelli for K3 surfaces]
\label{complex-geometry-theorem-torelli-k3}
Let $S_1,S_2$ be K3 surfaces.
Then $S_1\cong S_2$ if and only if their
integral Hodge structures on $H^2$ are
isomorphic, with the Kahler cone condition.
\end{theorem}

\begin{theorem}[Buskin]
\label{complex-geometry-theorem-buskin-k3-isogeny}
Let $S_1,S_2$ be projective K3 surfaces.
If
$$
H^2(S_1,\mathbb{Q})\cong
H^2(S_2,\mathbb{Q})
$$
as rational Hodge structures, then $S_1$ and
$S_2$ are symplectically isogenous.
\end{theorem}

\noindent
This is essentially a Hodge-conjecture
statement for $S_1\times S_2$.
An isomorphism of rational Hodge structures
on $H^2$ gives a rational Hodge class in
$$
H^2(S_1,\mathbb{Q})^*
\otimes H^2(S_2,\mathbb{Q})
\subset H^4(S_1\times S_2,\mathbb{Q}).
$$
To produce a symplectic isogeny, one must
realize this Hodge class by an algebraic
cycle, i.e. by a correspondence.
That is exactly the kind of assertion the
Hodge conjecture predicts.

The Shafarevich conjecture says:
if $M_1,M_2$ are hyperkahler manifolds in
the same deformation class and their rational
Hodge structures on $H^2$ are isomorphic,
then $M_1$ and $M_2$ are symplectically
isogenous.

\medskip
Very briefly: Torelli says integral periods
recover a K3 surface.
Buskin says rational periods for projective
K3 surfaces recover a symplectic isogeny.
Shafarevich asks for the analogous rational
period statement for hyperkahler manifolds.
The earlier lemma gives only the reverse
direction: an isogeny forces an isomorphism
of rational $H^2$ Hodge structures.

\begin{definition}
\label{complex-geometry-definition-holomorphic-lagrangian-fib}
Let $(M,\Omega)$ be a holomorphically
symplectic manifold of complex dimension
$2n$.
A holomorphic map
$$
\pi:M\longrightarrow B
$$
is a {\it holomorphic Lagrangian fibration}
if its fibers have dimension $n$ and
$\Omega$ restricts to zero on the smooth
part of every fiber.
\end{definition}

\noindent
By the holomorphic Liouville theorem, the
smooth compact fibers are complex tori.
Thus a holomorphic Lagrangian fibration is a
family of complex tori over its smooth base
locus.

\medskip
{\bf Degenerate twistor deformations.}

Now let $\pi:M\to X$ be a holomorphic
Lagrangian fibration, with $M$ irreducible
holomorphic symplectic and holomorphic
symplectic form $\Omega$.
Usually $X=\mathbb{P}^n$, but this is not
essential here.
Let $\omega_{FS}$ be a Kahler form on $X$.
For $\lambda\in\mathbb{C}$, put
$$
\Omega_\lambda=
\Omega+\lambda\pi^*\omega_{FS}.
$$

\begin{lemma}
\label{complex-geometry-claim-degenerate-twistor-form}
The closed complex form $\Omega_\lambda$
defines a complex structure on the smooth
manifold $M$ for which $\Omega_\lambda$ is a
holomorphic symplectic form.
\end{lemma}

\noindent
This is the basic degenerate twistor
deformation.
We deform the complex structure by adding a
form pulled back from the base of the
Lagrangian fibration.

\begin{definition}
\label{complex-geometry-definition-degenerate-twistor-def}
The family of holomorphic symplectic
manifolds obtained from the forms
$$
\Omega_\lambda=
\Omega+\lambda\pi^*\omega_{FS},
\qquad \lambda\in\mathbb{C},
$$
is called the {\it degenerate twistor
deformation} associated to the Lagrangian
fibration $\pi:M\to X$.
\end{definition}

\noindent
Basic properties:
\begin{enumerate}
\item If the added form on $X$ is exact, the
deformation is trivial by Moser's theorem.
\item When $M$ is compact and of Kahler type,
these deformations are limits of ordinary
twistor deformations.
\item Each member still carries a
holomorphic Lagrangian fibration over the
same base $X$.
\item The fibers and the base are unchanged;
only the complex structure on the total
space is varied.
\item Algebro-geometrically, this is a
Tate-Shafarevich deformation: it is described
by a class in
$$
H^1(X,\operatorname{Aut}_\pi(M)),
$$
where $\operatorname{Aut}_\pi(M)$ is the
sheaf of fiberwise automorphisms.
\end{enumerate}

\begin{definition}
\label{complex-geometry-definition-torsor}
Let $G\to X$ be a group object over $X$.
A {\it torsor} over $G$ is a space
$P\to X$ with a fiberwise action of $G$ such
that, locally on $X$, $P$ is isomorphic to
$G$ with its translation action.
Equivalently, $G$ acts freely and
transitively on the fibers of $P\to X$.
\end{definition}

\noindent
In the Lagrangian fibration setting, the
smooth fibers are tori.
Different fibrations in the same degenerate
twistor family can be viewed as torsors over
the same relative torus fibration.

\medskip
Upshot.
The culminating goal is to show that any two
projective hyperkahler manifolds in the same
degenerate twistor family are symplectically
isogenous.
To construct this isogeny, first make the two
torus fibrations literally the same after a
finite cover of the base.
Then the graph of the resulting isomorphism,
after compactification, gives the desired
correspondence.

\begin{theorem}
\label{complex-geometry-theorem-degenerate-twistor-cover}
Let $M_1,M_2\to X$ be projective
hyperkahler manifolds in the same degenerate
twistor family.
Then there is an open subset $U\subset X$
and a finite cover $\widetilde{U}\to U$ such
that
$$
M_1\times_X\widetilde{U}
\cong
M_2\times_X\widetilde{U}
$$
over $\widetilde{U}$.
\end{theorem}

\noindent
This is the theorem where torsors enter.
After finite base change, the torsors acquire
sections, hence become trivial.
The graph of the resulting isomorphism over
an open set is then compactified to produce
the correspondence, giving the symplectic
isogeny.

\section{Hypercomplex manifolds}
\label{complex-geometry-section-hypercomplex-manifolds}

\begin{exercise}
\label{complex-geometry-exercise-SU2-invariant-forms}
Let $M$ be a compact hypercomplex manifold of
real dimension 4, equipped with a
quaternionic Hermitian structure, and $V$ the
space of closed
$\text{SU}(2)$-invariant 2-forms. Prove that
$V$ is finite-dimensional.
\end{exercise}

\begin{proof}[Proof by Arpan.]
Consider the Hodge star operator with respect
to the Riemannian metric on $M$.
\begin{enumerate}
\item A closed anti-self-dual form is
harmonic.
\item The space of harmonic forms is
finite-dimensional.
\item An $\text{SU}(2)$-invariant closed form
is
anti self-dual.
\begin{enumerate}
\item
$\Lambda^+=\text{span}(\omega_I,\omega_J,
\omega_K)$.
$M$ is a manifold of real dimension 4,
which means that $\dim\Lambda^2(M)=6$ and
thus
$\dim\Lambda^+(M)=3$
(here we can check using a basis of
$\mathbb{R}^6$
that $*$ permutes the basic covectors with no
fixed points,
so that it has zero trace;
in turn this means that $\dim \Lambda^+ +
\dim\Lambda^-=0$).
Further, each of $\omega_I$, $\omega_J$ and
$\omega_K$
squared and multiplied by $\frac{1}{2}$ give
the Riemannian
volume form, and therefore the three are in
$\Lambda^+(M)$.
Clearly they are linearly independent and
thus span $\Lambda^+(M)$.

\item $\omega_I,\omega_J$ and $\omega_K$ are
not fixed by the
$\text{SU}(2)$ action
Indeed, $J\omega_I=-\omega_I$ and similarly
for other cases. Therefore there are no
$\text{SU}(2)$-fixed
points in $\Lambda^+$.

\item
$\text{SU}(2)\Lambda^+(M)=\Lambda^+(M)$.
This is clear since $J\omega_I = -\omega_I$
and so on.
Since $*$ is an isomorphism of $\Lambda^+(M)$
it follows that
$\text{SU}(2)\Lambda^-(M)=\Lambda^-(M)$.

\medskip\noindent
We conclude that if $\alpha$ is
$\text{SU}(2)$-invariant (i.e. a fixed point
of
the action of $\text{SU}(2)$ on $\Lambda^2$)
its positive part will vanish, so that
$\alpha$ is anti-self-dual. More precisely,
if
$\alpha = \alpha^+ +\alpha^-$ is
$\text{SU}(2)$-invariant and $A \in
\text{SU}(2)$,
$$
\alpha^+ +\alpha^- = \alpha =
A\alpha=A\alpha^+ + A\alpha^-,
$$
so $A\alpha^+=\alpha^+$, but $\text{SU}(2)$
has no fixed points in $\Lambda^+$, so
$\alpha^+=0$.
Therefore $\alpha \in \Lambda^-$ as required.
\end{enumerate}
\end{enumerate}
\end{proof}

\begin{exercise}
\label{complex-geometry-exercise-close-funda-forms-almos}
Let $(M,K,J,K,g)$ be an almost hypercomplex
Hermitian structure
on a K3 surface, and $\omega_I, \omega_J,
\omega_K$
its fundamental forms.

Suppse that $\omega_I$ is closed.
Prove that $\omega_J,\omega_K$ are closed, or
find a counterexample.
\end{exercise}

\noindent
The idea of the following solution
(provided entirely by AI)
is that we can ``rotate''
(in the $J,K$-plane within $\text{End}(TM)$)
two of the complex structures while
preserving the
quaternionic relations.
Any orthonormal triple of such endomorphisms
should be expected to satisfy the
quaternionic relations (algebra),
but the closedness condition on the
associated fundamental
forms will be modified (differential).

\begin{proof}
Let $(M,I,J,K,g)$ be a hyperkähler structure
on a K3 surface.
Define
$$
\begin{aligned}
\tilde{J}&=\cos \theta J+ \sin \theta K\\
\tilde{K}&=-\sin \theta J + \cos \theta K
\end{aligned}
$$

\noindent
Then for any $v \in TM$,
$$
\begin{aligned}
\tilde{J}^2v&=\cos J(\cos \theta Jv+ \sin
\theta Kv)
+\sin \theta K (\cos J v + \sin Kv)\\
&=\cos^2J^2v+\cos \theta \sin \theta J K v
+ \sin \theta \cos \theta K J v + \sin^2
\theta K^2 v\\
&=-\text{Id} v
\end{aligned}
$$

\noindent
And likewise $\tilde{K}^2=-\text{Id}$.
Also,
$$
\begin{aligned}
I \tilde{J} \tilde{K}v&=
I(\cos\theta J(-\sin \theta Jv+\cos \theta
Kv)
+\sin \theta K (- \sin Jv+ \cos \theta Kv))\\
&=I (-\cos \theta \sin \theta J^2 v
+ \cos^2 \theta J K v
- \sin^2 \theta K J v
+\sin \theta \cos \theta K^2 v)\\
&=-\cos \theta \sin \theta I J^2 v
+ \cos^2 \theta I J K
- \sin^2 \theta I K J v
+ \sin \theta \cos \theta I K^2 v\\
&=-\cos^2v- \sin^2v=-\text{Id}v.
\end{aligned}
$$

\noindent
Thus $(I,\tilde{J},\tilde{K})$ is an almost
hypercomplex
structure on $M$.
However,
$$
\begin{aligned}
d
\omega_{\tilde{J}}&=d(g(\cdot,\tilde{J}\cdot)
)\\
&=d(g(\cdot,\cos \theta J\cdot + \sin \theta
K\cdot))\\
&=d(\cos \theta \omega_J + \sin \theta
\omega_K)\\
&=d \cos \theta \wedge \omega_J + d \sin
\theta \wedge \omega_K\\
&=-\sin \theta
d\theta\wedge\omega_J+\cos\theta
d\theta\wedge\omega_K\\
&=d \theta \wedge \omega_{\tilde{K}},
\end{aligned}
$$

\noindent
which is not always zero if $\theta$ is not
constant.
\end{proof}

\medskip\noindent
Now we introduce an operator
similar to $d^c$ for hypercomplex manifolds.
Let $d = \partial + \bar\partial$ w.r.t.
$I$. We define
$$
\partial_J:= (-1)^k J \bar\partial J.
$$
\begin{proposition}
\label{complex-geometry-proposition-twisted-j}
\begin{enumerate}
\item $\partial_J: \Omega^{p,q}_I(M)
\to \Omega^{p+1,q}_I(M)$.

\item $\partial_J^2=0$.

\item $\partial \partial_J
+ \partial_J \partial=0$.
\end{enumerate}
\end{proposition}

\medskip\noindent
We also have a suitable cohomology
theory for hypercomplex manifolds,
introduced by [LGV]:
$$
H_{\operatorname{qBC}}^{\bullet,\bullet}
=
\frac{\Ker \partial \cap \Ker \partial_J}{
\operatorname{Im} \partial \partial_J}.
$$

\medskip\noindent
Now consider a hermitian metric $g$
on a hypercomplex manifold 
$(M,I,J,K)$, that is, the three
complex structures are isometries of $g$.
Denote by $\omega$ the fundamental
form w.r.t. $I$. Then
$$
\omega \in
\Omega^{1,1}_I
\cap \Omega^{2,0+0,2}_J
\cap \Omega^{2,0+0,2}_J
:= A^{1,1},
$$
where $\Omega^{2,0+0,2}$ denotes 
$\Omega^{2,0}\oplus \Omega^{0,2}$.
To check this we simply observe that
$\omega \in \Omega^{2,0+0,2}_J \iff
J\omega=-\omega$ and compute.
$A^{1,1}$ are called the
{\it Salamon 2-forms}.

If $\alpha \in A^{1,1}$ and $\alpha$
is positive, i.e.
$\alpha(Ix,x)>0$ for all $x$,
then $g(x,y):=\alpha(Ix,x)$
is a hypercomplex hermitian metric.

\medskip\noindent
Now let
$$
\Omega=\frac{\omega_J+i \omega_J}{2}
\in \Omega^{2,0}_I(M).
$$
\begin{proposition}
\label{complex-geometry-proposition-properties-omega}
\begin{enumerate}
\item $J\overline{\Omega}=\Omega$. (Easy.)
\item (Nondegeneracy.)
$\Omega(JX,X)>0$ for all $x\neq 0$.
\end{enumerate}
\end{proposition}

\noindent
Thus, the set of hypercomplex metrics
is in one-to-one correspondence with
the space of $(2,0)$-forms satisfying
the two properties above.

\begin{definition}
\label{complex-geometry-definition-quaternionic-positive}
$\alpha \in \Omega^{2,0}_I(M)$
is {\it quaternionic positive} 
if $J \bar{\alpha}=\alpha$
and $\alpha(JX,X)>0$
for all nonzero $X$.
\end{definition}

\section{Hyperkähler manifolds}

\begin{definition}
A hyperkähler manifold is a Riemannian 
manifold
$(M,g)$ endowed with three complex structures
$I,J,K$ satisfying
$$
I^2=J^2=K^2=IJK=-1,
$$
such that $(M,g,L)$ is Kähler for each
$L \in \{I,J,K\}$.
\end{definition}

\begin{remark}
Equivalently, a hyperkähler manifold is a
Riemannian manifold whose holonomy satisfies
$$
\operatorname{Hol}(g)\subseteq \operatorname{Sp}(n).
$$
\end{remark}

\begin{remark}
If $(M,g,I,J,K)$ is hyperkähler, then
$$
\Omega_I=\omega_J+i\omega_K,
$$
where
$$
\omega_L(X,Y)=g(LX,Y),
$$
is a holomorphic symplectic form on $(M,I)$.
\end{remark}

\begin{definition}
A {\it holomorphic symplectic manifold} is a 
complex
manifold $(X,I)$ endowed with a closed
nondegenerate holomorphic $2$-form
$$
\Omega \in H^0(X,\Omega_X^2).
$$
\end{definition}

\noindent
Equivalently, a holomorphically symplectic
form is a closed nondegenerate
$(2,0)$-form. Indeed, if $\Omega$ has type
$(2,0)$ and $d\Omega=0$, then
$\bar\partial\Omega=0$, hence $\Omega$ is
holomorphic.

Recall that a complex $2$-form 
$\Omega$ has type $(2,0)$
if and only if
$$
\Omega(Iv,w)=i\,\Omega(v,w)
$$
for all tangent vectors $v,w$.

\medskip\noindent
The implication
$$
\text{hyperkähler}
\Longrightarrow
\text{holomorphically symplectic}
$$
is immediate from the construction
$\Omega_I=\omega_J+i\omega_K$.

Conversely, if $X$ is compact Kähler and
holomorphically symplectic, then $K_X$ is
trivial, hence $c_1(X)=0$. By Yau's theorem,
every Kähler class contains a unique Ricci-flat
Kähler metric. For this metric, the holomorphic
symplectic form is parallel, so
$$
\operatorname{Hol}(g)\subseteq \operatorname{Sp}(n).
$$
Therefore $X$ admits a hyperkähler metric.

\begin{remark}
If an almost complex manifold admits a closed
nondegenerate $(2,0)$-form $\Omega$, then the
almost complex structure is integrable.

Indeed, since $\Omega$ has type $(2,0)$,
$$
T^{0,1}=\ker(\Omega).
$$
The kernel of a closed form is involutive, so
$T^{0,1}$ is involutive. By Newlander--Nirenberg,
the almost complex structure is integrable.
\end{remark}

\begin{remark}
If $\nabla I=0$ for the Levi-Civita connection
of $g$, then $(M,g,I)$ is Kähler. In particular,
on a hyperkähler manifold the forms
$\omega_I,\omega_J,\omega_K$ are parallel and
therefore closed.
\end{remark}

\begin{exercise}
\label{complex-geometry-exercise-symplectic-forms-on-k3}
Let $\omega_1,\omega_2$ be symplectic
forms on a K3 surface,
satisfying $\omega_1 \wedge \omega_2=0$.
Prove that $\omega_1^2 = c\omega_2^2$
for some $c \in \mathbb{R}$,
or find a counterexample.
\end{exercise}

\begin{proof}
We define a complex-valued symplectic form
$$
\Omega=\omega_1 + i \omega_2.
$$
(Indeed, this form is nondegenerate
by complex dimension, 
and also clearly closed.)
The idea is: if $\Omega$ were holomorphic,
then $0=\Omega^2 = \omega_1^2-\omega_2^2$
where 
the middle term $i\omega_1 \wedge \omega_2$
vanishes using the exercise's hypothesis.
Also, $\Omega^2$ would be zero by
holomorphicity and dimsneion: 
$\Omega^0\in H^{4,0}\cong H^0(X,\Omega^4)=0$.

But the problem is that $\Omega$ is not
automatically holomorphic.
The idea would be to show that
$$
\Omega(Iv,w)=i\Omega(v,w),
$$
but this is not immediate because 
$\omega_1$ and $\omega_2$
are just real symplectic.
\end{proof}

\noindent
Though still unsolved,
at least the exercise motivates
the following definition.
But don't get lost, there's several
subtleties here: all we know is: if $\Omega$
was holomorphic and complex dimension
was 2, then $\Omega^2 = 0$, which is 
the first condition in the definition.
And in defining $\Omega$ as in the exercise,
we have nondegeneracy of the real and
imaginary parts, which matches
the second part of the definition.
BUT, these conditions do not make $\Omega$
holomorphic in general.
I think the exercise needs an application
of Torelli theorem, as suggested in the
slides.

\begin{definition}
\label{complex-geometry-definition-c-symplectic}
A {\it C-symplectic} form on a complex
manifold $M$ of complex dimension $n$
is a complex-valued 2-form such that
$\Omega^{n+1} =0$
and $\Omega^n \wedge \overline{\Omega}^n$
is a nondegenerate volume form.
\end{definition}


\medskip\noindent
{\bf Hyperkähler with torsion.}
Recall the following definition from the
previous section:
$$
\Omega:=\frac{\omega_J+i \omega_J}{2}
\in \Omega^{2,0}_I(M).
$$
Also see
Proposition
\ref{complex-geometry-proposition-properties-omega}.

\begin{definition}
\label{complex-geometry-definition-hkt}
Let $(M,g,I,J,K)$ be a hypercomplex
hermitian manifold. We call it
{\it HKT} if any of the following
properties holds:
\begin{enumerate}
\item ([HP].) The Bismut connections
are the same:
$$
\nabla_I^B=\nabla_J^B=\nabla_K^B=\nabla^B.
$$

\item ([GP].) $\partial \Omega=0$.

\item ([BS].)
$$
d\omega \in
\bigcap_{L \in S}
\Omega_L^{2,1}+\Omega_L^{1,2}
=
\bigcap_{L \in \{I,J,K\}}
\Omega^{2,1}_L
+
\Omega^{1,2}_L
$$
where $S=\{a I+b J+ c K:a^2+b^2+c^2=1\}$
is the sphere of complex structures.
\end{enumerate}
\end{definition}

\noindent
Thus $\Omega$ defines a
$H^{2,0}_{\operatorname{qBC}}$ class.

\medskip\noindent
{\bf HKT Einstein.}
\begin{definition}
\label{complex-geometry-definition-phi-hkt-einstein}
Let $(M,g,I,J,K)$ be hyperhermitian.
Define
$$
\phi(\omega)
=
\frac{\rho-J\rho}{2}
$$
where $\rho$ is the Chern-Ricci
form of $(I,g)$.
\end{definition}

\noindent
This $\phi$ is nothing but the
projection to the
$\Omega^{2,0+0,2}$ component w.r.t. $J$.

\begin{lemma}
\label{complex-geometry-lemma-phi-is-a11}
$\phi(\omega) \in A^{1,1}$.
\end{lemma}

\begin{definition}
\label{complex-geometry-definition-hkt-einstein}
A HKT structure is {\it HKT Einstein} 
if 
$$
\phi(\omega)=\lambda \omega,\qquad \lambda
\in \mathbb{R}.
$$
\end{definition}

\noindent
Note that the Chern scalar curvature
of $(L,g)$ satisfies $S=2n \lambda$,
and, by a computation,
$\lambda \geq 0$.

\medskip\noindent
Now observe that for any hypercomplex
manifold there is a $(1,0)$-form for $I$
$\alpha$ such that
$$
\partial \overline{\Omega}^n
=
\alpha \wedge \overline{\Omega}^n.
$$
In the HKT case,
$\theta=\alpha+\overline{\alpha}$.

We define the following hypercomplex version
of the curvature form:
$$
[-\partial_J\alpha]
:=c_1^{\operatorname{qBC}}.
$$

\begin{definition}
\label{complex-geometry-definition-hkt-ricci-flow}
Let $(I,J,K,g,M)$ be a HKT manifold.
A 1-parameter family of HKT metrics
$\omega(t)$ is a solution of the
{\it HKT Ricci flow} if $\omega(0)=\omega_0$
and
$$
\partial_t \omega
=
-\phi(\omega).
$$
\end{definition}

\noindent
This flow was already introduced
by [BGV] as a flow of hyperhermitian
matrics to study the quaternionic
Monge-Ampère equation.

A computation shows
$-\phi(\omega)$
can be written as $S+Q$
where $S$ is the Chern Ricci form
and $Q$ is some form.






\section{Holonomy}
\label{complex-geometry-section-holonomy}

\noindent
Let $M$ be a smooth manifold
and $\nabla$ a linear connection.
Any loop $\gamma$ gives a map
$P_\gamma:T_pM \to T_pM$.

Define
$$
\text{Hol}_p(\nabla)
=\{P_\gamma:\text{loops starting at }p\}
\subseteq \text{GL}(n,\mathbb{R}),
$$
which indeed is a subgroup of
$\text{GL}(n,\mathbb{R})$
since parallel transport is invertible
by changing the direction of the
parametrization.
If $M$ is connected, $\text{Hol}_p(\nabla)$
is independent of $p$.

\begin{example}
\label{complex-geometry-example-holonomy}
\begin{enumerate}
\item For $\nabla^{\text{LC}}$
we have $\text{Hol}(\nabla^{\text{LC}})
\subset O(n,\mathbb{R})$.

\end{enumerate}
\end{example}

\begin{theorem}[Berger, 1955]
\label{complex-geometry-theorem-berger}
If $(M,g)$ a complete Riemannian
simply connected
manifold,
not locally a product
and not a symmetric space,
then $\text{Hol}(\nabla^{\text{LC}})$
can only be one of the following:
\begin{enumerate}
\item $\text{SO}(n)$ oriented Riemannian
manifold.
\item $\text{SU}(n/2)$ Calabi-Yau.
\item $Sp(n/4)$ Hyperkähler.
\item $Sp(n/4)Sp(1)$ quaternionic Kähler.
\item $U(n/2)$ Kähler.
\item $G_2$ exceptional geometry -+.
\item $\text{Spin}(7)$ exceptional
geometry 8.
\end{enumerate}
\end{theorem}

\section{Bismut connection}
\label{complex-geometry-section-bistmuth}

\noindent
Let $(M,g)$ be a Riemannian manifold.
Let $\nabla^g$ be LC connection and
$\nabla$ any other metric connection.
Define $A:= \nabla-\nabla^g$,
which is a tensor.
Let $A(X,Y,Z)=g(A(X,Y),Z)$.

\begin{lemma}
\label{complex-geometry-lemma-torsion-connections}
$A(X,Y,Z)+A(X,Z,Y)=0$.
\end{lemma}

\begin{proof}
Exercise.
\end{proof}

\noindent
This shows that there is a bijective
correspondence between metric connections
and tensors $A:\mathfrak{X}(M) \times
\mathfrak{X}(M) \to \mathfrak{X}(M)$
satisfying $A(X,Y,Z)+A(X,Z,Y)=0$.

We can also reconstruction the torsion
from $A$, and viceversa:

\begin{lemma}
\label{complex-geometry-lemma-torsion-connections2}
$T(X,Y)=A(X,Y)-A(Y,X)$.
\end{lemma}

\begin{proof}
Exercise.
\end{proof}

\begin{definition}
\label{complex-geometry-definition-skew-torsion}
A metric connection $\nabla$ has
{\it skew torsion} when
$A \in \Omega^3(M)$.
\end{definition}

\begin{theorem}[Strominger 1986, Bismut 1989]
\label{complex-geometry-theorem-strominger-bismut}
Let $(M,J,g)$ be a Hermitian manifold.
There exists a unique metric connection with
skew torsion such that $\nabla J=0$.
\end{theorem}

\begin{proof}
Let $\nabla=\nabla^g + \frac{1}{2}H$,
where $H \in \Omega^3(M)$.
The condition $\nabla J=0$ shall
determine the rest of the proof.
Since $g$ is metric, we know that
$\nabla \omega=0 \iff \nabla J=0$.
We use a permutation expression for
$d \omega$, that the tensor $N_J$ vanishes
to conclude that
$H(X,Y,Z) = -d \omega(JX,JY,JZ)$.
\end{proof}

\begin{lemma}[Corollary]
\label{complex-geometry-lemma-kahler-bismut}
$(J,g)$ is Kähler if and only if
$\nabla^g H =0 $.
\end{lemma}

\begin{proof}
$(J,g)$ is Kähler if and only if $H=0$
if and only if $\nabla^g J=0$,
by definition of the Bismut connection (and
the reverse by uniqueness).
\end{proof}




These notes document Beatrice Brienza's
work based on a few talks by her.

\begin{theorem}[Gauduchon]
\label{complex-geometry-theorem-gauduchon-connections}
On $TM$ there exists an affine line
of Hermitian connections known as
{\it canonical} or
{\it Gauduchon connections}
$$
g(\nabla_X^tY,Z)
=g(\nabla_X^{\text{LC}}Y,Z)
+\frac{t-1}{4}
d^c\omega(X,Y,Z)
+\frac{t+1}{4}
d^c(X,JY,JZ).
$$
\end{theorem}

\noindent
Note that:
\begin{enumerate}
\item $\nabla^{\text{LC}}$ is Hermitian
if and only if $d \omega=0$.
\item $d \omega \neq 0 \implies
\forall t \in \mathbb{R}$,
$\nabla^t$ has torsion. This leads to
``Hermitian geometry with torsion''.
\item $t=1$ gives the Chern connection,
i.e. the unique Hermitian connection
such that $\nabla^{0,1}=\bar\partial$
(see Definition
\ref{complex-geometry-definition-chern-connection}).
\item $t=-1$ gives the Bismut connection,
defined next, satisfying
$$
g(\nabla_X^BY,Z)
=g(\nabla^{\text{LC}}_XY,Z)
-\frac{1}{2}d^c\omega(X,Y,Z),
$$
where $\omega$ is the {\it Bismut torsion}.
\end{enumerate}

\noindent
Upshot: Chern connection is adapted
for the complex structure,
Bistmut connection is adapted
for skew torsion,
both equal Levi-Civita connection in
Kähler case.

\begin{definition}[Stve, Bis, Fl, 2002]
\label{complex-geometry-definition-bismut-connection}
The {\it Bismut connection} is the
unique Hermitian connection with
totally skew symmetric torsion, i.e.
$g(T^B(X,Y),Z) \in \Omega^3(M)$,
where of course
$T^B(X,Y)=\nabla^B_XY-\nabla^B_YX-[X,Y]$
is the usual torsion.
\end{definition}

After lowering an index of the torsion,
we obtain a tensor $T$:
$$
T \in \Lambda^2(T^*M) \otimes T^*M
\simeq \underbrace{T^*M}_{\text{vectorial
type}} \oplus\underbrace{\Lambda^3T^*M}_{
\text{skew type}}
\oplus \underbrace{\tau}_{\text{
terra incognita}}
$$
and we are interested in the case
where $T$ is in $\Lambda^3T^*M$.

\medskip\noindent
It turns out that
$$
\rho^B=0 \iff
\text{Hol}(\nabla^B) \subseteq
\text{SU}(n)
$$
where $\rho^B$ is the
Bismut Ricci curvature,
$$
\rho^B(X,Y)
=\frac{1}{2}\sum_i^{2n}
R^B(X,Y,e_i,J e_i)
\in \Omega^2(M),
$$
in fact $\rho^B \in H^2$ is
the Chern class of $M$.

\begin{definition}
\label{complex-geometry-definition-bismut-hermite-einstein}
A {\it Bismut Hermite Einstein} manifold
is when
$$
dH=d d^c\omega=0,\qquad
\rho^B=0.
$$
\end{definition}

\noindent
And $\rho^B=0 \iff \text{Hol}(\nabla^B)
\subseteq \text{SU}(n)$.

\medskip\noindent
In dimension 4, a BHE
can only be a Kähler Ricci flat manifold
or $S^1 \times \text{SU}(2)$. On
dimension 6, ABSL.

\begin{definition}
\label{complex-geometry-definition-almost-contact}
An {\it almost contact manifold}
is a triple $(\varphi,\eta,\xi)$
where $\varphi \in \text{End}(TM)$,
$\eta \in \Omega^1(M)$,
$\xi \in \mathfrak{X}(M)$,
$\varphi^2 = -\text{Id} + \eta \otimes \xi$,
$\eta(\xi)=1$.
\end{definition}

\noindent
This means that
$$
TM = \mathcal{F}_{\xi}\oplus Ker\xi.
$$

We then introduce a metric
compatible with $\varphi$.




\section{Instantons}
\label{complex-geometry-section-instantons}

\noindent
Let $X$ be a compact 4-dimensional compact
oriented Riemannian manifold, $G$ a compact
Lie group (e.g. $U(n), \text{SO}(n),Sp(n)$).

\begin{definition}
\label{complex-geometry-definition-g-instanton}
A {\it $G$-instanton} on $X$
is a $G$-vector bundle $E$
with a connection $\nabla$ such that its
curvature $F_{\nabla}$ is self-dual, i.e.
$*F_{\nabla}=F_{\nabla}$
where $*$ is Hodge operator.
\end{definition}

ADHM studied the case when $X=S^4$,
and $X=\mathbb{C}P^2$ was also studied in a
paper by N. Buchdahl.

\section{Twistor space}
\label{complex-geometry-section-twistor-space}

\noindent

Let $X$ be a compact oriented 4-dimensional
Riemannian manifold. Recall that since the
Hodge star operator $*:\Lambda^2\to\Lambda^2$
has eigenvalues $\pm 1$, we have a
decomposition
$$
\Lambda^2=\Lambda^+ \oplus \Lambda^- \cong
\text{SO}(T_X)
$$
Consider the spherization of $\Lambda^+$:
$$
Z:=S\Lambda^+=\{S \in \Lambda^+:\text{unit
vectors}\}
$$
$(m,s) \in Z$, $m \in X$, $s_m \in
\Lambda^+_m$,

$$
T_{(m,s)}Z=
T_mX \oplus_{LC}T_sS \Lambda^+_m
\cong T\mathbb{C}P^1
$$
which has a complex structure
$\mathcal{I}_{\mathbb{C}P^1}$
(metric orientation).

\begin{theorem}
\label{complex-geometry-theorem-twistor-space}
$(Z,\mathcal{I}_Z=(s,\mathcal{I}_{\mathbb{C}
P^1}))$
is a complex manifold of complex dimension 3.
\end{theorem}

\noindent
We also have a projection $p:Z \to X $
whose fibers are $\mathbb{C}P^1$,
called {\it real lines}.

\begin{example}
\label{complex-geometry-example-twistor-space}
\begin{itemize}
\item $X=S^4$ then $Z=\mathbb{C}P^3$
(projective).
\item $X=\mathbb{C}P^2$ then
$Z=\mathbb{F}_3(\mathbb{C})$ (projective).
\end{itemize}
\end{example}

Let $Z$ be a 3-dimension complex manifold.
Let $L \xymatrix{\ar@{^{(}->}[r]&} Z$
be a rational curve with a normal bundle
$N_L \cong \mathcal{O}(1) \oplus
\mathcal{O}(1)$.

\begin{remark}
\label{complex-geometry-remark-small-deformations}
If $Z \supset W$ where $W$ is a complex
submanifold,
the small deformations of $W$ are encoded by
$H^0(N_W)$ and $H^1(N_W)$.
\end{remark}

Let $\mathcal{M}$ be the space of
deformations of $L$.
Since $H^1(N_L)=0$ then
$TM \simeq H^0(N_L) = H^0(\mathcal{O}(1)
\oplus \mathcal{O}(1))$,
so $\dim_\mathbb{C}TM=4$.

Let $\sigma:Z \to Z$ be an antiholomorphic
involution with no fixed points but $L$ ($L$
is fixed). There is an induced
antiholomorphic involution $\tilde{\sigma}$
on $\mathcal{M}$.

\begin{remark}
\label{complex-geometry-remark-antiholomorphic-involution}
\begin{itemize}
\item We can recover $X$ as a connected
component
of the fixed point set of $\tilde{\sigma}$ in
$\mathcal{M}$.
\item We can recover the
conformal structure.
\end{itemize}
\end{remark}

\medskip\noindent
Now we discuss Ward-Penrose transform.
Let $E$ be a $G$-instanton
(i.e. a vector bundle with a connection such
that
$*F_\nabla=F_\nabla$) on $X$.
Let $\phi:V \to \overline{V}^*$ be an
Hermitian form.

There is a correspondence (the {\it Ward
correspondence})
which allows to construct a holomorphic
vector bundle on $Z$
such that $\tilde{E}$ restricted to the fiber
of $p:Z \to X$
(the fibers of $p$ are $\mathbb{C}P^1$)
trivially.

Conversely, starting from a hermitian
holomorphic vector bundle $\tilde{E}$
on $Z$ which restricts trivially to the
fibers of $p:Z \to X$.

\medskip\noindent
{\bf Problem (ADHM).} Describe all
$G$-instanton bundles on $S^4$
with $r \geq 2$, $c_2=k$ (fixed second Chern
class).

This makes us want to study the vector
bundles on $\mathbb{C}P^3=Z$.

See Horrseks for the original construction of
monads.
Upshot: ``All bundles can be represented by
the cohomology of some monad''.

\begin{definition}[Monad]
\label{complex-geometry-definition-monad}
Let $U,V,W$ be complex vector spaces over
$\mathbb{C}P^1$.
Consider the sequence $h \circ a=0$, where
$a,b$are bundle homomorphisms:
$$
\xymatrix{
0\ar[r]&U(-1):=U \otimes \mathcal{O}(1)
\ar[r]^a
& V \otimes \mathcal{O}
\ar[r]^b
& W(1):=W \otimes \mathcal{O}(1)\ar[r]&0
}
$$
Let $E:=\frac{\Ker b}{\text{Im}a}$.
\end{definition}

\medskip\noindent
Let's study the case of $\mathbb{C}P^3$.
If $Z_i$ are homogeneous coordinates on
$\mathbb{C}P^3$,
$$
a=\sum_i A_iZ_i,\qquad b=\sum_jB_jZ_j
$$
the {\it monad condition} is
$$
B_jA_i+A_jB_j=0
$$
[which gives the conditions for a hyperkähler
metric on the moduli space].

\begin{theorem}[Barth, Barth-Hulek]
\label{complex-geometry-theorem-bbh}
Let $E$ be a holomorphic vector bundle of
rank $r \geq 2$
on $\mathbb{C}P^3$ with a $c_2=k$ and
$E|_L=$trivially.
Then $E$ is the cohomology of a monad.
Also, it satisfies some vanishing conditions:
$$
H^0(E)=H^0(E^*)=H^1(E(-2))=H^1(E^*(-2))=0.
$$
\end{theorem}






\section{Stability}
\label{complex-geometry-section-stability}

\noindent
This section outlines some concepts
from a paper by Verbitsky from approx. 2002.

A {\it Gauduchon metric} is such that
$\partial\bar\partial \omega^{n-1}=0$. The
degree of a holomorphic vector bundle is
given by $$ \text{deg}_g(\mathcal{F}) =
\int_M c_1(\det \mathcal{F}) \wedge
\omega^{n-1}, $$ which is well defined when
$g$ is Gauduchon.

\begin{definition}
\label{complex-geometry-definition-slope-and-semistability}
The {\it slope} of a vector bundle is
$$
\mu_g(\mathcal{F}) = \frac{\text{deg}_g
\mathcal{F}}{\text{rk}\mathcal{F}}.
$$
A vector bundle $\mathcal{F}$ is {\it stable}
if for any coherent subsheaf $\mathcal{E}
\subset \mathbb{F}$
$$
\mu_g(\mathcal{E}) < \mu_g(\mathcal{F}).
$$
It is called {\it semistable} if we put $\le
$
instead.
\end{definition}

\begin{example}
\label{complex-geometry-example-line-bundles-are-stable}
Line bundles are stable.
\end{example}

\begin{theorem}[Hitchin-Kobayashi
correspondence]
\label{complex-geometry-theorem-hitchin-kobayashi}
Let $(M,g)$ be a compact complex manifold
with $g$ Gauduchon.
Then an hermitian vector bundle $(F,h)$
is Hermitian-Einstein,
i.e. it satisfies
$$
\sqrt{-1}\Lambda\Theta=\gamma\text{id}
$$
if and only if
it is polystable, i.e.
it is the sum of stable bundles
with the same slope.
\end{theorem}

\noindent
The aim of Misha's paper is to
classify stable bundles in positive elliptic
fibrations using Hitchin-Kobayashi
correspondence. Positive elliptic fibrations
are $\pi:M \to X$ with an action of a torus
on the fibers, and such that
$\pi^*\omega_X=d\eta$ where $\omega_X$ is a
Kähler form on $X$. Such a metric is called
preferred if it is invariant on the fibres
and is a Riemannian submersion.

The main theorem is the following:

\begin{theorem}
\label{complex-geometry-theorem-stability-misha}
For a positive elliptic fibration $\pi:M \to
X$,
if $(B,h)$ is a stable bundle on $M$,
$$
B \simeq L \otimes \pi^*B_0
$$
where $L$ is a line bundle and
$B_0$ is stable bundle on $X$.
\end{theorem}

\noindent
The upshot is that a locally conformal
hyperkähler manifold
has an associated elliptic fibration…


\section{Deformation theory (introduction)}
\label{complex-geometry-section-deformations}

\noindent
Here I discuss deformations in the category
of complex spaces.

Let's just do a summary of \cite[Chapter
9]{Voisin}.

Let $\mathcal{X}$ be a complex manifold, $B$
a complex manifold
and $\phi:\mathcal{X} \to B$ a holomorphic
map.
We call $\phi$ a {\it family} if it is a
proper holomorphic submersion.

A theorem by Ehresmann says that any
``deformation'' of
smooth real manifolds is trivial:

\begin{theorem}[Ehresmann]
\label{complex-geometry-theorem-ehresmann}
\begin{reference}
\cite[Theorem 9.3]{voi}
\end{reference}
If $\phi:\mathcal{X} \to B$ is a
proper submersion of differentiable manifolds
(no complex structure), with $B$ contractible
and $0 \in B$
some distinguished point,
then the total space $\mathcal{X}$ is trivial
as a
manifold over $B$, i.e.
there exists a diffeomorphism
 $$
T:\mathcal{X}\xrightarrow{\cong}X_0 \times B
$$
of manifolds over $B$, i.e. commuting with
projection.
\end{theorem}

\begin{proof}
The main tool is the theorem of tubular
neighbourhoods,
which says that there is a diffeomorphism of
a nieghbourhood $W$ of $X_0$ in $\mathcal{X}$
and a neighbourhood of $X_0$ in its normal
bundle.
(We have a normal bundle because $X_0$ is a
submanifold of $\mathcal{X}$.)

In particular this gives a differentiable
retraction
$T_0:W \to X_0$, which allows us to consider
the map
$$
(T_0,\phi):W \to X_0 \times B.
$$
Since $T_0$ is a retraction
(it's the identity when restricted to $X_0$)
and $\phi$ is a submersion,
the map $(T_0,\phi)$ has nonsingular
differential along $X_0$.
Then we can invert it locally along $X_0$,
and since $X_0$ is compact we obtain a whole
neighbourhood $W'\subset W \subset
\mathcal{X}$
where the differential is nonzero and thus
$(T_0,\phi)|_{W'}$ is an embedding
(using the inverse function theorem).

Finally we restrict the image of
$(T_0,\phi)|_{W'}$
to a neighbourhood $U\subset B$ of $0$,
using properness of $\phi$ to ensure
that $\phi^{-1}(U) \subset W'$ --- though I'm
not exactly sure how.
We have constructed a diffeomorphism
$\phi^{-1}(U) \to X_0 \times U$
commuting with the projection.

This proof was only local; the global version
is in \cite{voi} though it's not used in
the rest of the text.
\end{proof}

In the complex situation, if
$\phi:\mathcal{X} \to B$ is a family
of complex manifolds and $0 \in B$, then
``up to replacing $B$ by a neighbourhood of
0'',
there exists a smooth trivialization
$T=(T_0,\phi):\mathcal{X} \to X_0 \times B$
such that the fibres of $T_0$ are complex
submanifolds of $\mathcal{X}$.

Now let $\phi:\mathcal{X} \to B$ be a family
of complex manifolds.
Recall the pullback bundle $\phi^*B$ and the
fact
that the differential $\phi_*$ can be thought
of as a gadget
that maps tangent vectors to $\mathcal{X}$
not to vectors tangent to $B$ but
to vectors in the pullback bundle $\phi^*B$.
Restricting everything to the fibre
$X=\phi^{-1}(0)$ we get
\begin{equation}
\label{complex-geometry-equation-short-exact-seque-defor}
\xymatrix{
0\ar[r]&T_X\ar[r]&T_{\mathcal{X}|X}\ar[r]&
\phi^*T_{B|X}\ar[r]&0.
}
\end{equation}
There's what seems to be a really nice result
saying that
$H^1(\mathcal{F})$ parametrizes isomorphisms
classes
of extensions of $\mathcal{F}$ by the trivial
bundle,
but for now let me just define it roughly as
follows.
The category of $\mathcal{O}_X$ is most
surely abelian
with enough injectives. This means that
objects
have associated complexes, and we can take
cohomologies of these complexes.
This says, roughly, that short exact
sequences of vector bundles have associated
cohomology
long exact sequences. Then the
Kodaira-Spencer map
is the map
$$
\rho:T_{B,0}=H^0(X,\phi^* T_{B|X})\to
H^1(X,T_X)
$$
induced by the long exact sequence associated
to
\ref{complex-geometry-equation-short-exact-seque-defor}.
Here the equality $T_{B,0}=H^1(X,\phi^*
T_{B|X})$
is just because the pullback of the tangent
bundle of $B$
to any fiber is trivial.

\medskip\noindent
Now I will explain first order deformations.
Consider a family $\mathcal{X} \to B$.
We introduce a complex analytic subscheme
$B_\varepsilon$ of $B$
giving an ideal sheaf $\mathcal{M}^2_0$
where $\mathcal{M}_0$ is the maximal ideal
of $\mathcal{O}_{B,0}$ consisting of the
holomorphic
functions vanishing at 0.
This means that a function in
$\mathcal{M}^2_0$
is (a finite sum of terms) of the form $f=gh$
for $g,h \in \mathcal{M}_0$,
which means that $df=d(gh)=hdg+gdh=0$ at $0$
since both $g$ and $h$ vanish at $0$.
That is, this sheaf contains functions that
vanish
at $0$ and whose first derivative vanishes at
$0$.

This means that $B_\varepsilon$ is a point:
it's the vanishing locus of the set of
functions
that vanish at a single point.
However, it is endowed with an unusual
structure sheaf.
Let $v$ be a tangent vector to $B$ at $0$ and
$f$ a function in the strange sheaf
$\mathcal{M}_0/\mathcal{M}_0^2$.
Then by the computation above, since $v$ is a
differential
operator, $v$ annihilates $f$.
This means that this sheaf detects not only
the point $0$
but also any tangent vector at $0$.
Let's denote
$B_\varepsilon=\Spec
\mathbb{C}[\varepsilon]/(\varepsilon^2)$
just because its structure sheaf is
$\mathbb{C}[\varepsilon]/(\varepsilon^2)$.

In fact, this is so throughout deformation
theory:
$\Spec k[\varepsilon]/(\varepsilon^2)$ is
topologically a point!
But it's a {\it fat point}: the support of
its
structure sheaf is a point, but as a scheme,
it's not a point.
And the same happens with the family
$\mathcal{X}$ over
$\Spec k[\varepsilon]/(\varepsilon^2)$,
topologically it's just the fibre $X$!
But its structure sheaf is different:
``it contains the directions in $\mathcal{X}$
which
are normal to $X$''.

\medskip\noindent
Perhaps the most important part of this
construction
comes when considering a fat point
$B_\varepsilon \subset B$.
This gives a first-order deformation
$\phi^{-1}(B_\varepsilon)) \subset
\mathcal{X}$
$$
\xymatrix{
X\ar[r]\ar[d]
&\phi^{-1}(B_\varepsilon) \subset
\mathcal{X}\ar[d]\\
\ast\ar[r]
&B_\varepsilon \subset B
}
$$
Using the constructions above, one shows that
to any such deformation we can associate a
cocycle
in $H^1(X,T_X)$.
Thus, first order deformations are ``exactly
parametrised''
by the group $H^1(X,T_X)$.
This is the main idea behind Kodaira-Spencer
map,
which is well-explained in
\href{https://en.wikipedia.org/wiki/Kodaira
% E2%80%93Spencer_map}{Wikipedia}.

I will present a similar construction
to Kodaira-Spencer map
that works for embedded deformations
instead of just deformations.
(That's with the intention of understanding
the
forthcoming exercise Exercise
\ref{complex-geometry-exercise-deformations-curve-in-k3}.)

Since we are in complex space category
I follow \cite[p. 224]{voi} for the main
setup,
but this is also a mixture
of \cite[p. 147, 150]{Sernesi}.

The result is in fact quite reasonable
geometrically:
the dimension of directions in which we can
deform a
submanifold should be non other than
the dimension of the normal bundle.

\section{Deformation theory}
\label{complex-geometry-section-deformation-theory}

\begin{definition}
\label{complex-geometry-definition-deformation}
Let $X$ be a complex manifold.
A {\it deformation} of $X$ is a holomorphic
submersion
$\pi:\mathcal{X} \to B$ of complex manifolds
such that $\pi^{-1}(0)\cong X$ for some $0
\in B$.

If $\mathcal{X}$ and $\mathcal{Y}$ are two
deformations
of $X$ over the same base,
a {\it morphism of deformations} is a map
$f:\mathcal{X} \to \mathcal{Y}$ such that the
diagram
$$
\xymatrix{
& X\ar[dl]\ar[dr]\\
\mathcal{X}\ar[dr]\ar[rr]^f&&\mathcal{Y}
\ar[dl]\\
&B
}
$$
commutes. The {\it space of deformations} of
$X$
is the set of deformations modulo
the equivalence relation of two deformations
being isomorphic. A priori we don't know
if this set has any structure.

Let $Y$ be another complex manifold.
A {\it deformation of $X$ in $Y$}
is a deformation of $X$ such that
$\mathcal{X}\subset Y\times B$.

The {\it fat point} is a ringed space
$(B_\varepsilon,\mathbb{C}[\varepsilon]/
(\varepsilon^2))$
whose underlying topological space
$B_\varepsilon$ is a point
and structure sheaf is the ring of dual
numbers
$\mathbb{C}[\varepsilon]/(\varepsilon^2)$.

A {\it first order deformation} of $X$ is a
deformation such that
the base $B$ is a fat point.
\end{definition}

The following proposition is a
mixture of ideas from \cite[p. 224]{voi}
and \cite[Examples
3.2.4(ii)]{Sernesi-deformations}.

\begin{proposition}
\label{complex-geometry-proposition-kodai-spenc-embed-defor}
Let $X$ be a smooth hypersurface in a complex
manifold $Y$.
For every deformation of $X$ in $Y$
with base of dimension 1 we can construct a
section of the bundle
$N_{X/Y}$. Different deformations give
different sections up to rescaling.
\end{proposition}

\begin{proof}
Let
$$
\xymatrix{
X\ar[r]\ar[d]
&\mathcal{X}\subset Y\times B\ar[d]^\pi\\
0\ar[r]
&B
}
$$
be a deformation of $X$ in $Y$
and suppose that $\dim B=1$.
This means that $\mathcal{X}$ is a complex
submanifold of
the product $Y \times B$
and $\pi:\mathcal{X} \to B_\varepsilon$ is
a holomorphic submersion.

Since $X$ is a hypersurface,
we can find an open cover $V_i$
of $M$ such that $V_i \cap X:= U_i$
is given locally by the vanishing of
a function $f_i \in \mathcal{O}_Y(V_i)$.
In terms of coordinates we find that
that $V_i \cap X$ is given by
$(z_1,…,z_{n-1},0)$.

Since $\pi$ is a submersion,
by the holomorphic inverse function theorem
we obtain for every $V_i \subset \mathcal{X}$
(such that $V_i \cap X \neq  \emptyset$)
a biholomorphism from some $B_i \subset B$
containing $0$ to $V_i$.
This gives coordinates $(0,…,b)$ to $V_i$.
We conclude that $V_i \cong U_i\times B_i$
along $X$.

Now we ``restrict'' to a first order
deformation.
Consider the fat point subscheme
$B_\varepsilon$ of $B$
and let
$X_\varepsilon=\pi^{-1}(B_\varepsilon)$,
which is a ringed space whose topological
space
is the same as $X$, but its
structure sheaf is induced from
$B_\varepsilon$.

(I think that to prove this part I have
to look carefully at the induced sheaf
$\pi^*(\mathbb{C}[\varepsilon]/
(\varepsilon^2)$.)
Since we are on a local trivialization,
the structure sheaf of the product $U_i
\times B_i$
is the tensor product of the structure
sheaves
(reference?),
so that the structure sheaf of
$X_\varepsilon$ at
$U_i\times B_\varepsilon$ is given by
$$
\mathcal{O}_{X_\varepsilon}(U_i)=
 \mathcal{O}_Y(U_i) \otimes
 \mathbb{C}[\varepsilon]/(\varepsilon^2).
$$
Functions here look like finite sums of
functions of the form $f+ \varepsilon g$ for
$f,g \in \Gamma(U_i,\mathcal{O}_{U_i})$.

The next key step is to argue that
$X_\varepsilon$ is locally cut by
functions of the kind $f_i + \varepsilon g_i$
for some $g_i \in \mathcal{O}_M(U_i)$.
Here the $f_i$ are the same as those which
define $X$.
I think a proof is that from the ideal sheaf
exact sequence of $X$ in $Y$
$$
\xymatrix{
0\ar[r]
&\mathcal{I}_X\ar[r]
&\mathcal{O}_Y\ar[r]
&\mathcal{O}_X\ar[r]
&0
}
$$
we obtain after tensoring by
$\mathbb{C}[\varepsilon]/(\varepsilon^2)$
$$
\xymatrix{
0\ar[r]
&\mathcal{I}_X \otimes
\mathbb{C}[\varepsilon]/(\varepsilon^2)\ar[r]
&\mathcal{O}_Y \otimes
\mathbb{C}[\varepsilon]/(\varepsilon^2)\ar[r]
&\mathcal{O}_X \otimes
\mathbb{C}[\varepsilon]/(\varepsilon^2)\ar[r]
&0.
}
$$
(But is the functor $-\otimes
\mathbb{C}[\varepsilon]/(\varepsilon^2)$
exact? Why? So far we didn't impose a
flatness condition on $\pi$.)
The sheaf in the right
coincides with the structure sheaf of
$X_\varepsilon$ when
evaluated on a trivial open set of the form
$U_i \otimes B_\varepsilon$.
This shows that the ideal sheaf of
$X_\varepsilon$
as a ringed space is
$\mathcal{I}_X
\otimes\mathbb{C}[\varepsilon]/
(\varepsilon^3)$.
Since $\mathcal{I}_X$ is locally principal
generated by $f_i$,
we conclude that indeed $X_\varepsilon$
must be cut locally by functions of the form
$f_i + \varepsilon g_i$
for some $g_i \in \mathcal{O}_Y(U_i)$

Recall that the line bundle associated to $X$
is $\mathcal{O}_Y(X)$, given by the cocycle
$f_i/f_j:=f_{ij}$.
A version of adjunction formula
(more precisely, \cite[Adjunction formula I,
p. 146]{GH})
establishes that this line bundle,
when restricted to $C$, is in fact
the normal bundle of $X$ in $Y$, defined as
que quotient
bundle $T_Y'/T_X':=N_{X/Y}$.
In symbols, $\mathcal{O}_Y(X)|_{X}\cong
N_{X/Y}$.

Similarly, the line bundle associated to
$\mathcal{X}$
is given by the cocycle
$$
F_{ij}=f_{ij}+\varepsilon g_{ij}
$$
for $g_{ij}=\frac{g_i-f_{ij}g_j}{f_j}$.
Indeed, note that
$\frac{1}{f_j+\varepsilon
g_j}=\frac{1}{f_j}-\varepsilon
\frac{g_j}{f^2_j}$
and compute $F_{ij}=(f_i+\varepsilon
g_j)/(f_j+\varepsilon g_j)$.

Then
$$
f_i+\varepsilon g_i=(f_{ij}+\varepsilon
g_{ij})(f_j+\varepsilon g_j)
$$
gives
$$
g_i=f_{ij}g_j+g_{ij}f_j
$$
and since $f_j$ vanishes along $X$,
we are left with $g_i=f_{ij}g_j$
which is how a section of the sheaf $N_{X/Y}$
is obtained.
\end{proof}

\begin{exercise}
\label{complex-geometry-exercise-deformations-curve-in-k3}
Let $C$ be a smooth genus $g$ curve which
can be embedded in a K3 surface $M$, and
$\mathcal{X}$
``the family of all deformations'' of $C$ in
$M$.
Prove that $\dim \mathcal{X} \leq g$.
\end{exercise}

\begin{proof}
By Adjunction formula II
(Theorem \ref{complex-geometry-theorem-adjunction-ii})
and $M$ being a K3 surface
we get that $h^0(C,N_{C,M})=g$.

If the deformation space $X$ exists and is a
complex manifold,
the map constructed in
Proposition
\ref{complex-geometry-proposition-kodai-spenc-embed-defor}
would be a holomorphic (why holomorphic!?)
injective map $X \to H^0(C,N_{C/Y})$.
This obliges $\dim X \leq  h^0(C,N_{C/Y})=g$.
Indeed, if not, the differential must have a
kernel,
and integrating along a direction in the
kernel
we obtain the germ of a complex curve along
which
the map is constant, which is impossible
since it is injective.
\end{proof}







% --- End complex-geometry.tex ---

% --- Begin hodge-theory.tex ---
\chapter{Hodge Theory}
\label{chapter-hodge-theory}


\author{Misha Verbitsky}
\date{IMPA, 2026}



These notes follow the slides for Misha
Verbitsky's 2026 course at IMPA. The
course materials are available on the
\href{%
http://verbit.ru/IMPA/Kahler-2026/%
}{course page}.

\section{Lecture 1: complex manifolds}
\label{hodge-theory-section-lecture-one-complex}

\noindent
This lecture was given on August 10,
2026. Erich K\"ahler described his work
as a golden thread leading into the
palpitating heart of mathematics.

\medskip\noindent
{\bf Complex structures.}

\begin{definition}
\label{hodge-theory-definition-complex-operator}
Let $V$ be a real vector space. A
{\it complex structure operator} on $V$
is an automorphism
$$
I:V\longrightarrow V
$$
satisfying
$$
I^2=-\operatorname{Id}_V.
$$
\end{definition}

The action of $I$ extends
multiplicatively to every tensor space
built from $V$ and $V^*$. Thus, for
example,
$$
\begin{aligned}
I(&v_1\otimes\cdots\otimes v_k
\otimes w_1\otimes\cdots\otimes w_l)
\\
&=I(v_1)\otimes\cdots\otimes I(v_k)
\otimes I(w_1)\otimes\cdots\otimes
I(w_l).
\end{aligned}
$$

The following elementary facts will be
used repeatedly.

\begin{enumerate}
\item The eigenvalues of $I$ are
$\sqrt{-1}$ and $-\sqrt{-1}$.

\item There is an $I$-invariant positive
definite scalar product on $V$. Indeed,
starting with $g_0$, one may take
$$
g(x,y)=g_0(x,y)+g_0(Ix,Iy).
$$

\item The operator $I$ is orthogonal for
this scalar product, since
$$
g(Ix,Iy)=g(x,y).
$$

\item The operator $I$ is diagonalizable
over $\mathbb{C}$.

\item The two eigenvalues occur with the
same multiplicity.
\end{enumerate}

In a suitable complex basis of
$V\otimes_{\mathbb{R}}\mathbb{C}$, the
matrix of $I$ is
$$
\begin{bmatrix}
\sqrt{-1}\operatorname{Id}&0\\
0&-\sqrt{-1}\operatorname{Id}
\end{bmatrix}.
$$
In a suitable real basis it is block
diagonal, with every block equal to
$$
\begin{bmatrix}
0&-1\\
1&0
\end{bmatrix}.
$$

\begin{definition}
\label{hodge-theory-definition-vector-hodge}
Let $(V,I)$ be a real vector space with a
complex structure. Its
{\it Hodge decomposition} is
$$
V\otimes_{\mathbb{R}}\mathbb{C}
=V^{1,0}\oplus V^{0,1},
$$
where $V^{1,0}$ and $V^{0,1}$ are the
$\sqrt{-1}$ and $-\sqrt{-1}$
eigenspaces of $I$, respectively.
\end{definition}

The dual space has the analogous Hodge
decomposition. Moreover, $V^{1,0}$
determines $I$: the operator acts by
$\sqrt{-1}$ on $V^{1,0}$ and by
$-\sqrt{-1}$ on its conjugate. We thus
obtain a bijection between complex
structures on $V$ and subspaces
$$
W\subset V\otimes_{\mathbb{R}}\mathbb{C}
$$
such that
$$
\dim_{\mathbb{C}}W
=\frac{1}{2}\dim_{\mathbb{R}}V,
\qquad
W\cap\overline{W}=0.
$$

\begin{definition}
\label{hodge-theory-definition-hermitian-metric-vector}
An $I$-invariant positive definite scalar
product $g$ on $(V,I)$ is a
{\it Hermitian metric}. The triple
$(V,I,g)$ is a {\it Hermitian space}.
\end{definition}

If $g$ is Hermitian, the bilinear form
$$
\omega(x,y)=g(x,Iy)
$$
is skew-symmetric. Indeed,
$$
\begin{aligned}
\omega(x,y)
&=g(x,Iy)\\
&=g(Ix,I^2y)\\
&=-g(Ix,y)\\
&=-\omega(y,x).
\end{aligned}
$$

\begin{definition}
\label{hodge-theory-definition-hermitian-form-vector}
The skew-symmetric form $\omega$ above is
the {\it Hermitian form} of $(V,I,g)$.
\end{definition}

Any two members of the triple
$(I,g,\omega)$ determine the third.

\medskip\noindent
{\bf Almost complex manifolds.}

\begin{definition}
\label{hodge-theory-definition-almost-complex}
An {\it almost complex structure} on a
manifold $M$ is an endomorphism
$$
I:TM\longrightarrow TM
$$
such that
$$
I^2=-\operatorname{Id}_{TM}.
$$
The pair $(M,I)$ is an
{\it almost complex manifold}.
\end{definition}

On $\mathbb{C}^n$, write
$$
z_j=x_j+\sqrt{-1}y_j.
$$
The standard almost complex structure is
given by
$$
I\left(\frac{\partial}{\partial x_j}\right)
=\frac{\partial}{\partial y_j},
\qquad
I\left(\frac{\partial}{\partial y_j}\right)
=-\frac{\partial}{\partial x_j}.
$$

In real dimension two, an almost complex
structure is equivalent to an oriented
conformal structure.

\begin{exercise}
\label{hodge-theory-exercise-almost-complex-conformal}
Prove this equivalence.
\end{exercise}

\begin{definition}
\label{hodge-theory-definition-holomorphic-function-ac}
A function
$$
f:M\longrightarrow\mathbb{C}
$$
on an almost complex manifold is
{\it holomorphic} if
$$
df\in\Lambda^{1,0}(M).
$$
\end{definition}

Some almost complex manifolds have no
nonconstant holomorphic functions even
locally. A standard example is $S^6$ with
its $G_2$-invariant almost complex
structure.

\begin{theorem}
\label{hodge-theory-theorem-holomorphic-equivalences}
Let $U\subset\mathbb{C}^n$ be open and let
$$
f:U\longrightarrow\mathbb{C}
$$
be differentiable. The following
conditions are equivalent.

\begin{enumerate}
\item The function $f$ is holomorphic.

\item At every point, the differential
$$
df:T_xU\longrightarrow\mathbb{C}
$$
is complex linear.

\item The restriction of $f$ to every
complex affine line is holomorphic as a
function of one complex variable.

\item Near every point, $f$ is represented
by a convergent complex Taylor series.
\end{enumerate}
\end{theorem}

\begin{proof}
The first two conditions are equivalent by
definition. A $(1,0)$-form restricts to a
$(1,0)$-form on every complex line. The
converse follows from the elementary fact
that a complex-valued linear form on
$(V,I)$ has type $(1,0)$ if and only if its
restriction to every $I$-invariant real
two-plane has type $(1,0)$. This proves the
equivalence with the third condition. A
convergent complex Taylor series is
holomorphic, while the converse follows
from the Cauchy formula.
\end{proof}

\begin{exercise}
\label{hodge-theory-exercise-line-restrictions}
Prove the linear-algebraic fact used in the
proof of Theorem
\ref{hodge-theory-theorem-holomorphic-equivalences}.
\end{exercise}

{\bf Sheaves and ringed spaces.}

\begin{definition}
\label{hodge-theory-definition-presheaf-functions}
A {\it presheaf of functions} on a
topological space $M$ assigns a subring
$$
\mathcal{F}(U)\subset C(U)
$$
to every open set $U\subset M$, in such a
way that restriction from $U$ to any open
subset $U_1$ maps $\mathcal{F}(U)$ into
$\mathcal{F}(U_1)$.
\end{definition}

\begin{definition}
\label{hodge-theory-definition-sheaf-functions}
A presheaf $\mathcal{F}$ is a
{\it sheaf of functions} if it also has the
following local-to-global property. Given
an open cover $U=\bigcup_iU_i$, a function
on $U$ belongs to $\mathcal{F}(U)$ whenever
its restriction to every $U_i$ belongs to
$\mathcal{F}(U_i)$.
\end{definition}

Continuous, smooth, and real analytic
functions form sheaves. Constant functions
and bounded functions form presheaves that
are not sheaves.

\begin{exercise}
\label{hodge-theory-exercise-presheaves-not-sheaves}
Explain why constant functions and bounded
functions fail the sheaf condition.
\end{exercise}

\begin{definition}
\label{hodge-theory-definition-ringed-space-course}
A {\it ringed space} $(M,\mathcal{F})$ is a
topological space with a sheaf of functions.
A {\it morphism of ringed spaces}
$$
\Psi:(M,\mathcal{F})
\longrightarrow(N,\mathcal{F}')
$$
is a continuous map such that, for every
open $U\subset N$ and every
$f\in\mathcal{F}'(U)$,
$$
\Psi^*f=f\circ\Psi
\in\mathcal{F}(\Psi^{-1}(U)).
$$
An {\it isomorphism of ringed spaces} is a
homeomorphism for which both directions are
morphisms.
\end{definition}

If $M$ is a $C^i$ manifold, the sheaf
$C^i$ makes $(M,C^i)$ a ringed space. Every
$C^i$ map induces a morphism between the
corresponding ringed spaces.

\begin{proposition}
\label{hodge-theory-proposition-smooth-ringed-maps}
Morphisms between the ringed spaces
$(M,C^i)$ and $(N,C^i)$ are precisely the
$C^i$ maps from $M$ to $N$.
\end{proposition}

\begin{proof}
One implication follows because a
composition of $C^i$ functions is $C^i$.
Conversely, work in charts. A map between
open subsets of $\mathbb{R}^n$ is $C^i$ if
the pullback of every $C^i$ function is
$C^i$, since this condition applies in
particular to the coordinate functions.
\end{proof}

\begin{definition}
\label{hodge-theory-definition-smooth-ringed-manifold}
A {\it smooth manifold of class $C^i$} is a
topological manifold with a sheaf of
functions which is locally isomorphic, as a
ringed space, to an open subset of
$(\mathbb{R}^n,C^i)$.
\end{definition}

\begin{definition}
\label{hodge-theory-definition-chart-ringed}
A {\it chart}, or
{\it coordinate system}, on an open subset
of such a manifold is an isomorphism with
an open subset of $(\mathbb{R}^n,C^i)$.
\end{definition}

\begin{definition}
\label{hodge-theory-definition-diffeomorphism-ringed}
A {\it diffeomorphism} is a homeomorphism
which induces an isomorphism of the
corresponding ringed spaces.
\end{definition}

These definitions agree with the usual
definitions in terms of atlases.

\begin{definition}
\label{hodge-theory-definition-holomorphic-function-cn}
A {\it holomorphic function} on
$\mathbb{C}^n$ is a function whose
differential is complex linear, or
equivalently has type $(1,0)$.
\end{definition}

Holomorphic functions form a sheaf.

\begin{definition}
\label{hodge-theory-definition-complex-manifold-ringed}
A {\it complex manifold} is a ringed space
locally isomorphic to an open ball in
$\mathbb{C}^n$ with its sheaf of
holomorphic functions.
\end{definition}

Equivalently, it is a manifold covered by
complex coordinate balls whose transition
functions are holomorphic.

\begin{definition}
\label{hodge-theory-definition-holomorphic-ac-map}
A map
$$
\Psi:(M,I)\longrightarrow(N,J)
$$
of almost complex manifolds is
{\it holomorphic} if
$$
\Psi^*\Lambda^{1,0}(N)
\subset\Lambda^{1,0}(M).
$$
Equivalently, $d\Psi$ is complex linear.
\end{definition}

For the standard almost complex structures,
this is equivalent to holomorphicity of all
coordinate components of $\Psi$.

\begin{definition}
\label{hodge-theory-definition-integrable-ac}
An almost complex structure $I$ on $M$ is
{\it integrable} if every point has a
neighborhood diffeomorphic to an open set
in $\mathbb{C}^n$ under a diffeomorphism
which identifies $I$ with the standard
almost complex structure.
\end{definition}

An integrable almost complex structure
defines a sheaf $\mathcal{O}_M$ of
holomorphic functions and makes
$(M,\mathcal{O}_M)$ a complex manifold.
Conversely, a complex manifold determines
an integrable almost complex structure. In
fact, its sheaf of holomorphic functions
determines the subbundle
$$
d\mathcal{O}_M
=\Lambda^{1,0}(M)
\subset\Lambda^1(M,\mathbb{C}),
$$
and this subbundle determines $I$.

\section{Assignment 1: complex structures}
\label{hodge-theory-section-assignment-one-complex}

\noindent
This assignment was issued on August 10,
2026. Written solutions were requested for
the following Monday, with class discussion
on the Wednesday after that.

\begin{exercise}
\label{hodge-theory-exercise-assignment-one-gl}
Prove that the group
$\operatorname{GL}(n,\mathbb{C})$ is
connected.
\end{exercise}

\begin{definition}
\label{hodge-theory-definition-assignment-complex}
A {\it complex structure} on a real vector
space $V$ is an operator
$I\in\operatorname{End}(V)$ satisfying
$I^2=-\operatorname{Id}$.
\end{definition}

\begin{exercise}
\label{hodge-theory-exercise-orthogonal-complex-four}
Let $V$ be a four-dimensional real vector
space with a Euclidean metric, and let $S$
be the space of orthogonal complex
structures on $V$. Prove that $S$ is a
disconnected union of two two-dimensional
spheres.
\end{exercise}

\begin{exercise}
\label{hodge-theory-exercise-all-complex-components}
Let $V$ be an even-dimensional real vector
space, and let $S$ be the space of complex
structures on $V$. Prove that $S$ has two
connected components.
\end{exercise}

\begin{exercise}
\label{hodge-theory-exercise-orthogonal-components}
Let $V$ be an even-dimensional real vector
space with a Euclidean metric $g$, and let
$S_g$ be the space of orthogonal complex
structures on $V$. Prove that $S_g$ has two
connected components.
\end{exercise}

\begin{exercise}
\label{hodge-theory-exercise-two-forms-complex}
Let $V$ be a four-dimensional real vector
space, and let $U$ be the space of
skew-symmetric two-forms $\Omega$ on
$V\otimes_{\mathbb{R}}\mathbb{C}$ such that
$$
\Omega\wedge\Omega=0,
\qquad
\Omega\wedge\overline{\Omega}\ne0.
$$
Let $\mathbb{P}U$ be the projectivization
of $U$, namely the quotient of
$U\setminus\{0\}$ by the standard
$\mathbb{C}^*$-action. Construct a
$\operatorname{GL}(V)$-invariant bijection
between $\mathbb{P}U$ and the space of
complex structures on $V$.
\end{exercise}

\begin{proof}
Set $V_{\mathbb C}
=V\otimes_{\mathbb R}\mathbb C$.
We first recall that a nonzero two-form
$\Omega\in\Lambda^2V_{\mathbb C}^*$
is decomposable if and only if
$\Omega\wedge\Omega=0$.
Indeed, there exists a basis
$e^1,\ldots,e^4$ of $V_{\mathbb C}^*$
such that
\[
\Omega
=
\lambda_1e^1\wedge e^2
+
\lambda_2e^3\wedge e^4.
\]
Therefore
\[
\Omega\wedge\Omega
=
2\lambda_1\lambda_2
e^1\wedge e^2\wedge e^3\wedge e^4.
\]
Thus, if $\Omega\wedge\Omega=0$,
one of the $\lambda_i$ vanishes,
and hence $\Omega$ is decomposable.

Now let $\Omega\in U$. Write
$\Omega=\alpha\wedge\beta$ and define
\[
L_\Omega
=
\ker\Omega
=
\ker\alpha\cap\ker\beta.
\]
This is a two-dimensional complex
subspace of $V_{\mathbb C}$.
The condition
$\Omega\wedge\overline{\Omega}\ne0$
says that
$\alpha,\beta,\overline{\alpha},
\overline{\beta}$ are linearly
independent. Consequently,
\[
L_\Omega\cap\overline{L_\Omega}=0,
\qquad
V_{\mathbb C}
=
L_\Omega\oplus\overline{L_\Omega}.
\]

Define a complex-linear endomorphism of
$V_{\mathbb C}$ by
\[
J_\Omega|_{L_\Omega}
=
-i\,\operatorname{id},
\qquad
J_\Omega|_{\overline{L_\Omega}}
=
i\,\operatorname{id}.
\]
It commutes with complex conjugation,
so it is the complexification of a real
endomorphism $J_\Omega:V\to V$.
Moreover,
\[
J_\Omega^2=-\operatorname{id},
\]
so $J_\Omega$ is a complex structure
on $V$.

Since $L_{c\Omega}=L_\Omega$ for every
$c\in\mathbb C^*$, the complex structure
$J_\Omega$ depends only on the projective
class $[\Omega]$. Denote by
$\mathcal J(V)$ the set of complex
structures on $V$. We obtain a map
\[
\Phi:\mathbb PU\longrightarrow\mathcal J(V),
\qquad
[\Omega]\longmapsto J_\Omega.
\]

Conversely, let $J$ be a complex
structure on $V$. Its complexification
gives a decomposition
\[
V_{\mathbb C}
=
V_J^{1,0}\oplus V_J^{0,1},
\]
where $V_J^{1,0}$ and $V_J^{0,1}$ are
the $i$- and $-i$-eigenspaces,
respectively. Choose a nonzero form
\[
\Omega_J
\in
\Lambda^{2,0}V_{\mathbb C}^*.
\]
Equivalently, $\Omega_J$ is a nonzero
two-form with
$\ker\Omega_J=V_J^{0,1}$.
Since $V_J^{1,0}$ has complex dimension
two,
\[
\Omega_J\wedge\Omega_J=0.
\]
Furthermore, the decomposition of
$V_{\mathbb C}$ implies
\[
\Omega_J\wedge\overline{\Omega_J}\ne0.
\]
Thus $\Omega_J\in U$. Since
$\Lambda^{2,0}V_{\mathbb C}^*$ is
one-dimensional, the class $[\Omega_J]$
is independent of the choice of
$\Omega_J$.

These two constructions are inverse.
Finally, if $g\in\operatorname{GL}(V)$,
then
\[
\ker(g\cdot\Omega)=g(\ker\Omega),
\]
and hence
\[
J_{g\cdot\Omega}
=
gJ_\Omega g^{-1}.
\]
Therefore $\Phi$ is a
$\operatorname{GL}(V)$-equivariant
bijection.
\end{proof}


\begin{definition}
\label{hodge-theory-definition-conformal-structure}
Two Riemannian metrics $g_1$ and $g_2$ on a
smooth manifold are {\it conformal}, or
{\it conformally equivalent}, if there is a
smooth positive function $\lambda$ such that
$$
g_1=\lambda g_2.
$$
A {\it conformal structure} is a metric up
to conformal equivalence.
\end{definition}

\begin{exercise}
\label{hodge-theory-exercise-surface-conformal}
Let
$$
f:M\longrightarrow N
$$
be an orientation-preserving diffeomorphism
between almost Hermitian manifolds of real
dimension two. Prove that $f$ is
holomorphic if and only if it preserves the
conformal structures.
\end{exercise}

\begin{exercise}
\label{hodge-theory-exercise-no-almost-complex}
Construct an oriented even-dimensional
manifold which admits no almost complex
structure.
\end{exercise}
\section{Lecture 2: real analytic
manifolds}
\label{hodge-theory-section-lecture-two-real-analytic}

\noindent
This lecture was given on August 14, 2026.

\begin{definition}
\label{hodge-theory-definition-real-analytic-function}
A {\it real analytic function} on an open
set $U\subset\mathbb{R}^n$ is a function
which has a convergent Taylor expansion near
every point. Thus, locally,
$$
\begin{aligned}
f(z_1+t_1,&\ldots,z_n+t_n)\\
&=\sum_{i_1,\ldots,i_n}
a_{i_1,\ldots,i_n}
t_1^{i_1}\cdots t_n^{i_n}.
\end{aligned}
$$
The real parameters $t_j$ are sufficiently
small for the expansion under consideration.
\end{definition}

Real analytic functions form a sheaf.

\begin{definition}
\label{hodge-theory-definition-real-analytic-manifold}
A {\it real analytic manifold} is a ringed
space locally isomorphic to an open ball in
$\mathbb{R}^n$ equipped with its sheaf of
real analytic functions.
\end{definition}

Grauert's description of real analytic
manifolds in terms of complexifications is
developed below.

\medskip\noindent
{\bf Involutions and fixed points.}

\begin{definition}
\label{hodge-theory-definition-involution}
An {\it involution} on $M$ is a map
$$
\iota:M\longrightarrow M
$$
satisfying
$$
\iota^2=\operatorname{Id}_M.
$$
\end{definition}

\begin{exercise}
\label{hodge-theory-exercise-linear-involution}
Prove that every linear involution on a real
vector space is diagonalizable with
eigenvalues $1$ and $-1$.
\end{exercise}

\begin{theorem}
\label{hodge-theory-theorem-fixed-set-smooth}
Let $M$ be a smooth manifold and let $\iota$
be a smooth involution. The fixed point set
of $\iota$ is a smooth submanifold.
\end{theorem}

\begin{proof}
Fix a point $m$ in the fixed point set. By
the inverse function theorem, $k$ functions
whose differentials are linearly independent
at $m$ define coordinates near $m$.

Suppose that $d\iota_m$ has $k$
eigenvalues equal to $1$ and $n-k$ equal to
$-1$. Choose coordinates $x_1,\ldots,x_n$
such that the first $k$ covectors at $m$ are
invariant and the remaining covectors are
anti-invariant. Set
$$
y_j=x_j+\iota^*x_j
$$
for $j\leq k$, and set
$$
y_j=x_j-\iota^*x_j
$$
for $j>k$. Up to harmless constant factors,
the $dy_j$ remain linearly independent at
$m$. They therefore give invariant local
coordinates, and the fixed point set is
locally cut out by
$$
y_{k+1}=\cdots=y_n=0.
$$
\end{proof}

\medskip\noindent
{\bf Real structures.}

\begin{definition}
\label{hodge-theory-definition-real-structure-vector}
A {\it real structure} on a complex vector
space $V$ is an $\mathbb{R}$-linear
involution $\iota$ such that
$$
\iota(\lambda x)
=\overline{\lambda}\iota(x).
$$
\end{definition}

\begin{proposition}
\label{hodge-theory-proposition-fixed-real-form}
Let $V^\iota$ be the fixed subspace of a
complex vector space with a real structure.
Then
$$
\dim_{\mathbb{R}}V^\iota
=\dim_{\mathbb{C}}V,
\qquad
V=V^\iota\otimes_{\mathbb{R}}\mathbb{C}.
$$
\end{proposition}

\begin{proof}
Suppose that $x_1,\ldots,x_n$ are linearly
independent in $V^\iota$ over $\mathbb{R}$.
Applying $\iota$ to a complex linear
relation among them and averaging shows that
the real parts of all coefficients vanish.
Applying the same argument after multiplying
the relation by $\sqrt{-1}$ shows that their
imaginary parts vanish. Hence the natural map
$$
V^\iota\otimes_{\mathbb{R}}\mathbb{C}
\longrightarrow V
$$
is injective.

It is also surjective. For $v\in V$, both
$$
\frac{v+\iota(v)}{2}
$$
and
$$
\frac{v-\iota(v)}{2\sqrt{-1}}
$$
belong to $V^\iota$, and together they
recover $v$ by complex linear combination.
\end{proof}

Let $\mathcal{C}_\iota$ be the category of
complex vector spaces with real structures,
where morphisms are complex linear maps
commuting with the involutions. Let
$\mathcal{C}_{\mathbb{R}}$ be the category
of real vector spaces.

\begin{theorem}
\label{hodge-theory-theorem-real-structure-category}
The functor
$$
\begin{aligned}
\mathcal{C}_\iota
&\longrightarrow\mathcal{C}_{\mathbb{R}},\\
(V,\iota)&\longmapsto V^\iota
\end{aligned}
$$
is an equivalence of categories.
\end{theorem}

\begin{proof}
A quasi-inverse sends a real vector space
$W$ to
$$
W\otimes_{\mathbb{R}}\mathbb{C}
$$
with the involution
$$
w\otimes\lambda
\longmapsto
w\otimes\overline{\lambda}.
$$
The two composites are naturally isomorphic
to the identity by Proposition
\ref{hodge-theory-proposition-fixed-real-form}.
\end{proof}

\begin{definition}
\label{hodge-theory-definition-antiholomorphic-map}
A smooth map
$$
\Psi:(M,I)\longrightarrow(N,J)
$$
is {\it antiholomorphic} if
$$
d\Psi\circ I=-J\circ d\Psi.
$$
A complex-valued function is
{\it antiholomorphic} if its conjugate is
holomorphic.
\end{definition}

\begin{exercise}
\label{hodge-theory-exercise-antiholomorphic-function}
Prove that an antiholomorphic function on
$M$ defines an antiholomorphic map from $M$
to $\mathbb{C}$.
\end{exercise}

\begin{exercise}
\label{hodge-theory-exercise-antiholomorphic-pullback}
Prove that $\Psi$ is antiholomorphic if and
only if
$$
\Psi^*\Lambda^{0,1}(N)
\subset\Lambda^{1,0}(M).
$$
\end{exercise}

\begin{exercise}
\label{hodge-theory-exercise-antiholomorphic-functions}
Let $\iota$ be a smooth self-map of a
complex manifold. Prove that $\iota$ is
antiholomorphic if and only if it pulls
every locally defined holomorphic function
back to an antiholomorphic one.
\end{exercise}

\begin{definition}
\label{hodge-theory-definition-real-structure-manifold}
A {\it real structure} on a complex
manifold $M$ is an antiholomorphic
involution
$$
\iota:M\longrightarrow M.
$$
\end{definition}

\begin{proposition}
\label{hodge-theory-proposition-real-local-coordinates}
Let $M$ be a complex manifold with a real
structure $\iota$. For every
$x\in M^\iota$, there is an
$\iota$-invariant coordinate neighborhood
with holomorphic coordinates
$z_1,\ldots,z_n$ such that
$$
\iota^*z_j=\overline{z_j}.
$$
\end{proposition}

\begin{proof}
Any basis of $\Lambda_x^{1,0}(M)$ occurs as
the differentials at $x$ of suitable
holomorphic coordinates: start from an
arbitrary chart and apply a complex linear
change of coordinates.

The differential $d\iota_x$ is a real
structure. A real basis of the cotangent
space to $M^\iota$ therefore determines a
complex basis of $T_x^*M$. Choose
holomorphic coordinates $u_1,\ldots,u_n$
with these prescribed differentials. The
functions
$$
z_j=u_j+\iota^*\overline{u_j}
$$
are holomorphic and satisfy the required
identity. After shrinking the neighborhood,
they form a coordinate system.
\end{proof}

\begin{definition}
\label{hodge-theory-definition-compatible-real-atlas}
An atlas is {\it compatible with the real
structure} if every chart which meets the
fixed point set has the form given by
Proposition
\ref{hodge-theory-proposition-real-local-coordinates}.
\end{definition}

\begin{proposition}
\label{hodge-theory-proposition-fixed-real-analytic}
Let $M^\iota$ be the fixed point set of a
real structure on a complex manifold $M$.
A complex atlas compatible with the real
structure restricts to a real analytic
atlas on $M^\iota$.
\end{proposition}

\begin{proof}
The transition maps between compatible
charts commute with conjugation. They
therefore preserve the real locus and
restrict there to real-valued real analytic
maps.
\end{proof}

\begin{proposition}
\label{hodge-theory-proposition-all-real-fixed-loci}
Every real analytic manifold arises as the
fixed point set of a real structure on a
complex manifold.
\end{proposition}

\begin{proof}
Take a locally finite real analytic atlas by
relatively compact balls $U_i$ with
transition maps $\Psi_{ij}$. For sufficiently
small $\varepsilon>0$, their convergent
Taylor series extend the transition maps to
complex neighborhoods
$$
\widetilde{U}_i
=U_i\times B_\varepsilon
\subset\mathbb{C}^n,
$$
where the second factor lies in the
imaginary directions. These complexified
transition maps glue the neighborhoods to a
complex manifold $M_{\mathbb{C}}$. Complex
conjugation is compatible with the gluing
and gives a real structure whose fixed point
set is the original manifold.
\end{proof}

\begin{definition}
\label{hodge-theory-definition-complexification}
Let $M_{\mathbb{R}}$ be a real analytic
manifold. A {\it complexification} of
$M_{\mathbb{R}}$ is a complex manifold
$M_{\mathbb{C}}$ with a real structure whose
fixed point set is $M_{\mathbb{R}}$.
\end{definition}

There are several complexifications of the
circle. They include
$$
\mathbb{C}/\mathbb{Z}^2
$$
with $z\mapsto\overline z$, the space
$\mathbb{C}^*$ with
$$
z\longmapsto\overline z^{-1},
$$
and $\mathbb{C}P^1$ with
$$
[x:y]\longmapsto
[\overline{x}:\overline{y}].
$$
Their fixed loci are respectively
$\mathbb{R}/\mathbb{Z}$, the unit circle,
and $\mathbb{R}P^1$.

\begin{definition}
\label{hodge-theory-definition-real-analytic-tensor}
A tensor on a real analytic manifold is
{\it real analytic} if, in real analytic
coordinates, it is a sum of coordinate
tensor monomials with real analytic
coefficients.
\end{definition}

\begin{definition}
\label{hodge-theory-definition-same-germ}
Let $K\subset M$ and $K_1\subset M_1$ be
identified closed subsets of complex
manifolds. The manifolds $M$ and $M_1$
{\it have the same germ along $K$} if there
are neighborhoods $U$ and $U_1$ of the
identified set and a biholomorphism
$$
U\longrightarrow U_1
$$
which is the identity on $K$.
\end{definition}

\begin{definition}
\label{hodge-theory-definition-germ-manifold}
The {\it germ of $M$ along $K$} is the
equivalence class of neighborhoods of $K$
under the relation in Definition
\ref{hodge-theory-definition-same-germ}.
\end{definition}

Let $\mathcal{C}$ be the category of complex
manifolds with real structures and
holomorphic maps commuting with those
structures.

\begin{theorem}[Grauert]
\label{hodge-theory-theorem-grauert-real-analytic}
The category of real analytic manifolds is
equivalent to the category of germs of
objects $(M,\iota)$ of $\mathcal{C}$ along
their fixed point sets $M^\iota$.
\end{theorem}

\begin{exercise}
\label{hodge-theory-exercise-grauert-equivalence}
Prove Theorem
\ref{hodge-theory-theorem-grauert-real-analytic}.
\end{exercise}

The theorem is associated with Hans Grauert,
who lived from 1930 to 2011.

\begin{definition}
\label{hodge-theory-definition-antiholomorphic-tensor}
An {\it antiholomorphic tensor} on $(M,I)$
is a holomorphic tensor on $(M,-I)$.
\end{definition}

A function is holomorphic on $(M,I)$ if and
only if it is antiholomorphic on $(M,-I)$.
Complex conjugation identifies the sheaves
of holomorphic and antiholomorphic
functions. If $\iota$ is a real structure,
then $\iota^*$ sends holomorphic functions
and tensors to antiholomorphic ones.

\begin{theorem}
\label{hodge-theory-theorem-extend-analytic-tensors}
Let $M_{\mathbb{C}}$ be a complexification
of $M_{\mathbb{R}}$. Every real analytic
tensor $\psi_{\mathbb{R}}$ extends to a
holomorphic tensor $\psi_{\mathbb{C}}$ on a
neighborhood of $M_{\mathbb{R}}$ such that
$$
\iota^*\psi_{\mathbb{C}}
=\overline{\psi_{\mathbb{C}}}.
$$
Conversely, every holomorphic tensor near
$M_{\mathbb{R}}$ satisfying this identity
restricts to a real analytic real tensor on
$M_{\mathbb{R}}$.
\end{theorem}

\begin{proof}
A holomorphic tensor restricts to a real
analytic tensor. The displayed identity
makes the restriction real. Conversely,
write a real analytic tensor in local
coordinates and replace the convergent real
Taylor series of its coefficients by their
complexifications. After shrinking the
neighborhoods, the resulting holomorphic
tensors agree on overlaps and glue.
\end{proof}

\medskip\noindent
\section{Lecture 2 bis: Frobenius theorem}
\label{hodge-theory-section-lecture-two-bis-frobenius}

\noindent
This lecture was given on August 17, 2026.

\begin{definition}
\label{hodge-theory-definition-distribution}
A {\it distribution}, or
{\it Pfaffian system}, on a manifold is a
subbundle
$$
B\subset TM.
$$
\end{definition}

Let
$$
\Pi:TM\longrightarrow TM/B
$$
be the quotient map. If $X$ and $Y$ are
sections of $B$, then
$$
[fX,Y]=f[X,Y]-Y(f)X.
$$
It follows that $\Pi([X,Y])$ is
$C^\infty(M)$-linear in both $X$ and $Y$.

\begin{definition}
\label{hodge-theory-definition-frobenius-form}
The resulting bundle map
$$
\Lambda^2B\longrightarrow TM/B
$$
is the {\it Frobenius bracket}, or
{\it Frobenius form}, of $B$.
\end{definition}

\begin{definition}
\label{hodge-theory-definition-involutive-distribution}
A distribution is {\it holonomic},
{\it involutive}, or {\it integrable} if
its Frobenius form vanishes.
\end{definition}

\medskip\noindent
{\bf Submersions and foliations.}

\begin{definition}
\label{hodge-theory-definition-submersion-immersion}
A smooth map
$$
\pi:M\longrightarrow M'
$$
is a {\it submersion} if $d\pi$ is
surjective at every point, and an
{\it immersion} if $d\pi$ is injective at
every point.
\end{definition}

\begin{proposition}
\label{hodge-theory-proposition-submersion-local-form}
Let $\pi:M\longrightarrow M'$ be a
submersion. Every point of $M$ has a
neighborhood $U$ isomorphic to a product
$V\times W$ in such a way that
$$
\pi|_U:V\times W\longrightarrow W
$$
is the second projection.
\end{proposition}

\begin{proof}
This is the local form supplied by the
inverse function theorem.
\end{proof}

\begin{definition}
\label{hodge-theory-definition-vertical-tangent}
The {\it vertical tangent bundle} of a
submersion $\pi:M\longrightarrow M'$ is
$$
T_\pi M=\ker d\pi.
$$
\end{definition}

\begin{theorem}[Ehresmann]
\label{hodge-theory-theorem-ehresmann-fibration}
A smooth proper submersion is a locally
trivial smooth fibration. In particular,
this holds for a smooth submersion between
compact manifolds.
\end{theorem}

\begin{theorem}[Frobenius]
\label{hodge-theory-theorem-frobenius}
Let $B\subset TM$ be a subbundle. It is
involutive if and only if every point has a
neighborhood $U$ and a smooth submersion
$$
\pi:U\longrightarrow V
$$
such that
$$
B|_U=T_\pi U.
$$
\end{theorem}

One implication is immediate from the local
normal form for a submersion. The fibers of
$\pi$ are the local integral submanifolds of
$B$.

\begin{definition}
\label{hodge-theory-definition-leaf-foliation}
A {\it leaf} of $B$ is a maximal connected
immersed submanifold tangent to $B$ at every
point. If $B$ is integrable, its collection
of leaves is the {\it foliation} associated
with $B$.
\end{definition}

Leaves need not be closed. To prove the
Frobenius theorem, it is enough to find an
integral submanifold through every point;
locally, the required submersion is then the
projection to the leaf space.

The problem was known as Pfaff's problem
before Cartan used the name Frobenius theorem
in 1922. Jacobi obtained an early weaker
result in 1837, Deahna supplied an important
technical argument in 1844, Clebsch treated
the generic case in 1866, and Frobenius gave
the general reformulation and solution in
1877. A historical account appears in
Thomas Hawkins's book on Frobenius in
context. The mathematicians highlighted in
the slides were Carl Gustav Jacobi
(1804--1851), Alfred Clebsch (1833--1872),
and Ferdinand Georg Frobenius (1849--1917).

\medskip\noindent
{\bf Ehresmann connections.}

\begin{definition}
\label{hodge-theory-definition-ehresmann-connection}
Let $\pi:M\longrightarrow Z$ be a smooth
submersion. An {\it Ehresmann connection}
is a subbundle $T_{\rm hor}M\subset TM$
such that
$$
TM=T_{\rm hor}M\oplus T_\pi M.
$$
\end{definition}

Given a path
$$
\gamma:[0,a]\longrightarrow Z,
$$
parallel transport for an Ehresmann
connection is a smoothly varying family of
diffeomorphisms
$$
V_t:\pi^{-1}(\gamma(0))
\longrightarrow\pi^{-1}(\gamma(t))
$$
whose trajectories are tangent to
$T_{\rm hor}M$.

\begin{proposition}
\label{hodge-theory-proposition-ehresmann-transport}
If $\pi$ is proper, parallel transport
exists along every compact path.
\end{proposition}

\begin{proof}
The horizontal lift is the solution of an
ordinary differential equation. Compactness
of the inverse image of the path prevents
the solution from escaping in finite time.
\end{proof}

This also gives the standard proof of
Theorem \ref{hodge-theory-theorem-ehresmann-fibration}:
parallel transport along short radial paths
in a coordinate ball trivializes the
submersion over that ball.

\medskip\noindent
{\bf Proof of the Frobenius theorem.}

Suppose first that a Lie group $G$ acts on
$M$ and its infinitesimal vector fields span
a subbundle $B\subset TM$. The $G$-orbits
are integral submanifolds, so $B$ is
integrable. Also, if $X$ and $Y$ commute,
their local flows commute. This can be seen
in coordinates for which $X$ is a coordinate
vector field.

\begin{exercise}
\label{hodge-theory-exercise-commuting-flows}
Prove that commuting vector fields have
commuting local flows.
\end{exercise}

Thus it is enough to construct locally a
basis of commuting vector fields for an
involutive distribution $B$. Fix a point
and choose a local submersion
$$
\sigma:M\longrightarrow M_1
$$
whose fibers are transverse to $B$. More
explicitly, in local coordinates take a
linear map to
$\mathbb{R}^{\operatorname{rank}B}$ whose
differential restricts to an isomorphism on
$B$ at the chosen point. After shrinking,
$$
d\sigma|_{B_x}:B_x
\longrightarrow T_{\sigma(x)}M_1
$$
is an isomorphism at every nearby point.

A vector field $X$ is projectable under
$\sigma$ if $d\sigma(X)$ is constant along
the fibers of $\sigma$. Projectable vector
fields are closed under the Lie bracket,
and
$$
d\sigma([X,Y])
=[d\sigma(X),d\sigma(Y)].
$$
Equivalently, if $\underline{X}$ is the
projected vector field, then
$$
\mathcal{L}_X(\sigma^*f)
=\sigma^*\mathcal{L}_{\underline{X}}f
$$
for every $f\in C^\infty(M_1)$.

Let $\zeta_1,\ldots,\zeta_k$ be coordinate
vector fields on $M_1$. There are unique
fields $\xi_1,\ldots,\xi_k$ in $B$ with
$$
d\sigma(\xi_i)=\zeta_i.
$$
They form a local basis of $B$. Since $B$
is involutive,
$$
[\xi_i,\xi_j]\in B.
$$
On the other hand,
$$
d\sigma([\xi_i,\xi_j])
=[\zeta_i,\zeta_j]=0.
$$
The restriction of $d\sigma$ to $B$ is
injective, so all these brackets vanish.
The commuting flows define a local action
whose orbits are the integral manifolds of
$B$, proving the theorem.

\section{Assignment 2: Frobenius form}
\label{hodge-theory-section-assignment-two-frobenius}

\noindent
This assignment was issued on August 17,
2026, with the same submission and
discussion rules as Assignment 1.

\begin{exercise}
\label{hodge-theory-exercise-dense-torus-flow}
Find a vector field on the two-dimensional
torus $T^2$ such that every orbit of its
flow is dense.
\end{exercise}

\begin{exercise}
\label{hodge-theory-exercise-bracket-growth}
Construct a four-manifold $M$ and a rank
two distribution $B\subset TM$ such that
$[B,B]$ has rank three and $[[B,B],B]$ has
rank four.
\end{exercise}

\begin{definition}
\label{hodge-theory-definition-contact-structure}
A {\it contact structure} on a
$(2n+1)$-dimensional manifold $M$ is a rank
$2n$ distribution $B\subset TM$ for which
$TM/B$ is a trivial line bundle and the
Frobenius form
$$
\Lambda^2B\longrightarrow TM/B
$$
is nondegenerate.
\end{definition}

\begin{exercise}
\label{hodge-theory-exercise-rescale-contact-form}
Let $\theta$ be a one-form on a
$(2n+1)$-dimensional manifold such that
$$
\theta\wedge(d\theta)^n
$$
is nondegenerate. Let $f$ be a nowhere
vanishing smooth function and set
$\theta'=f\theta$. Prove that
$$
\theta'\wedge(d\theta')^n
$$
is nondegenerate.
\end{exercise}

\begin{exercise}
\label{hodge-theory-exercise-kernel-contact}
Let $\theta$ be a one-form on a
$(2n+1)$-dimensional manifold such that
$\theta\wedge(d\theta)^n$ is nondegenerate.
Prove that
$$
\ker\theta\subset TM
$$
is a contact distribution.
\end{exercise}

\begin{exercise}
\label{hodge-theory-exercise-contact-defining-form}
Let $M$ be an odd-dimensional manifold and
let $B\subset TM$ be a contact distribution.
Prove that there is a one-form $\theta$ such
that
$$
B=\ker\theta
$$
and $\theta\wedge(d\theta)^n$ is
nondegenerate.
\end{exercise}

\begin{exercise}
\label{hodge-theory-exercise-contact-sphere}
Construct a contact structure on
$S^{2n+1}$ for every positive integer $n$.
\end{exercise}

\section{Lecture 3: the
Newlander--Nirenberg theorem}
\label{hodge-theory-section-lecture-three-newlander}

\noindent
This lecture was given on August 17, 2026.
It begins with the real analytic manifolds,
real structures, compatible coordinates,
and complexifications constructed in
Lecture 2.

\begin{lemma}
\label{hodge-theory-lemma-holomorphic-real-locus}
Let $X\subset\mathbb{C}^n$ be an open ball
invariant under conjugation, and let
$$
X_{\mathbb{R}}=X\cap\mathbb{R}^n.
$$
A holomorphic tensor on $X$ which vanishes
on $X_{\mathbb{R}}$ vanishes identically.
\end{lemma}

\begin{proof}
Every holomorphic coefficient which vanishes
on the real locus has all derivatives zero
there. Its Taylor series therefore vanishes,
and analytic continuation gives the result.
\end{proof}

\begin{definition}
\label{hodge-theory-definition-real-analytic-ac}
An almost complex structure is
{\it real analytic} if its defining
endomorphism of $TM$ is a real analytic
tensor.
\end{definition}

\begin{lemma}[Corollary]
\label{hodge-theory-lemma-complexified-ac-square}
Let $(M,I)$ be a real analytic almost
complex manifold. If $I_{\mathbb{C}}$ is
the holomorphic extension of $I$ to a
complexification, then
$$
I_{\mathbb{C}}^2
=-\operatorname{Id}.
$$
\end{lemma}

\begin{proof}
The holomorphic tensor
$$
I_{\mathbb{C}}^2+
\operatorname{Id}
$$
vanishes on the real locus, so Lemma
\ref{hodge-theory-lemma-holomorphic-real-locus} applies.
\end{proof}

\medskip\noindent
{\bf The underlying real analytic
manifold.}

A holomorphic map between open subsets of
$\mathbb{C}^n$ is real analytic when viewed
as a map between open subsets of
$\mathbb{R}^{2n}$. Indeed, its real and
imaginary parts have convergent Taylor
series in the real variables.

\begin{definition}
\label{hodge-theory-definition-underlying-real-analytic}
The {\it underlying real analytic manifold}
$M_{\mathbb{R}}$ of a complex manifold $M$
is the same manifold with its holomorphic
transition functions viewed as real analytic
maps.
\end{definition}

The real analytic functions on
$M_{\mathbb{R}}$ are the convergent power
series generated locally by holomorphic and
antiholomorphic functions. In holomorphic
coordinates, this is the same as taking
convergent power series in the real and
imaginary parts of the coordinates.

\begin{definition}
\label{hodge-theory-definition-conjugate-manifold}
The {\it complex conjugate manifold}
$\overline{M}$ of $(M,I)$ is the same real
manifold equipped with $-I$. Its holomorphic
functions are the antiholomorphic functions
on $M$.
\end{definition}

\begin{proposition}
\label{hodge-theory-proposition-product-complexification}
Let $M$ be a complex manifold. A
complexification of its underlying real
analytic manifold is
$$
M_{\mathbb{C}}=M\times\overline{M},
$$
with the real structure
$$
\tau(x,y)=(y,x).
$$
\end{proposition}

\begin{proof}
The fixed point set is the diagonal and is
naturally identified with
$M_{\mathbb{R}}$. Holomorphic functions from
the first factor and antiholomorphic
functions from the second restrict to the
corresponding functions on the diagonal.
Their convergent power series generate its
sheaf of real analytic functions.
\end{proof}

\medskip\noindent
{\bf Formal integrability.}

Recall that a complex manifold and an
integrable almost complex manifold are
equivalent notions, as explained after
Definition \ref{hodge-theory-definition-integrable-ac}.

\begin{definition}
\label{hodge-theory-definition-formally-integrable}
Let
$$
TM\otimes_{\mathbb{R}}\mathbb{C}
=T^{1,0}M\oplus T^{0,1}M
$$
be the Hodge decomposition of the
complexified tangent bundle. The almost
complex structure is
{\it formally integrable} if
$$
[T^{1,0}M,T^{1,0}M]
\subset T^{1,0}M.
$$
\end{definition}

\begin{definition}
\label{hodge-theory-definition-nijenhuis-tensor}
The Frobenius form of $T^{1,0}M$, viewed as
a tensor measuring the component of
$[T^{1,0}M,T^{1,0}M]$ in $T^{0,1}M$, is the
{\it Nijenhuis tensor}.
\end{definition}

\begin{proposition}
\label{hodge-theory-proposition-integrable-formal}
Every integrable almost complex structure is
formally integrable.
\end{proposition}

\begin{proof}
In holomorphic coordinates,
$T^{1,0}M$ is generated by the commuting
fields
$$
\frac{\partial}{\partial z_1},\ldots,
\frac{\partial}{\partial z_n}.
$$
Its Frobenius form therefore vanishes.
\end{proof}

\begin{theorem}[Newlander--Nirenberg]
\label{hodge-theory-theorem-newlander-nirenberg}
An almost complex structure is integrable
if and only if it is formally integrable.
\end{theorem}

The name Nijenhuis tensor originates in a
1955 paper of Jan Schouten and Kentaro Yano
on the geometric meaning of its vanishing.
The slides also highlighted Albert Nijenhuis.
The real analytic version of the theorem was
proved in 1951 by Beno Eckmann and Alfred
Fr\"olicher. The latter surname is often
incorrectly written with an additional
letter.

The smooth theorem appeared in the sole
paper of August Newlander, written jointly
with Louis Nirenberg and published in 1957.
In a later interview Nirenberg recalled that
Andr\'e Weil and Shiing-Shen Chern drew his
attention to the problem. The task was to
recognize Cauchy--Riemann operators in
arbitrary coordinates. The first approach
removed error terms one at a time and worked
in complex dimension two, but failed in
higher dimensions. The eventual proof used
a fully nonlinear change of variables to
remove all error terms simultaneously.

\medskip\noindent
{\bf The real analytic case.}

Let $(M,I)$ be real analytic and formally
integrable. Extend $I$ holomorphically to a
complexification $M_{\mathbb{C}}$. After
shrinking, Lemma
\ref{hodge-theory-lemma-complexified-ac-square} gives a
holomorphic decomposition
$$
TM_{\mathbb{C}}
=T^{1,0}M_{\mathbb{C}}
\oplus T^{0,1}M_{\mathbb{C}}.
$$
Formal integrability on the real locus and
Lemma \ref{hodge-theory-lemma-holomorphic-real-locus}
show that both distributions are involutive
near that locus. The Frobenius theorem gives
their $(1,0)$ and $(0,1)$ foliations.

When $I$ is already integrable, one may use
the complexification
$$
M_{\mathbb{C}}=M\times\overline{M}.
$$
If $\pi$ and $\overline{\pi}$ are the two
projections, the fibers of
$\overline{\pi}$ are the $(1,0)$ leaves and
the fibers of $\pi$ are the $(0,1)$ leaves.
Along the diagonal, this splitting is the
usual Hodge splitting. It follows in
particular that an integrable almost complex
structure is real analytic with respect to
its underlying real analytic atlas.

Holomorphic functions on
$M_{\mathbb{C}}$ which are constant along
the $(0,1)$ leaves restrict to holomorphic
functions on $(M,I)$, because their
differentials vanish in every $(0,1)$
direction.

\begin{theorem}
\label{hodge-theory-theorem-real-analytic-newlander}
Every real analytic, formally integrable
almost complex structure is integrable.
\end{theorem}

\begin{proof}
The two foliations above are transverse.
Consequently, the complexification is
locally a product of their leaf spaces,
say
$$
M_{\mathbb{C}}=M'\times M''.
$$
Functions on $M'$ lift to functions on the
product which are constant along the
$(0,1)$ leaves. Choose
$$
n=\frac{1}{2}\dim_{\mathbb{R}}M
$$
such functions $f_1,\ldots,f_n$ whose
differentials are linearly independent at a
given point. Their restrictions to $M$ are
holomorphic. The inverse function theorem
shows that they form local complex
coordinates, and their transition functions
are holomorphic by construction.
\end{proof}

\section{Assignment 3: almost complex
structures and integrability}
\label{hodge-theory-section-assignment-three-integrable}

\noindent
This assignment was issued on August 24,
2026, with the same submission and
discussion rules as the earlier assignments.

\begin{definition}
\label{hodge-theory-definition-hodge-component-operator}
Let $(M,I)$ be an almost complex manifold
and let
$$
A:\Lambda^*M\longrightarrow\Lambda^*M
$$
be linear. The {\it Hodge components} of
$A$ are the operators $A^{p,q}$ such that
$$
A=\sum_{p,q}A^{p,q}
$$
and
$$
A^{p,q}(\Lambda^{i,j}(M))
\subset\Lambda^{i+p,j+q}(M).
$$
\end{definition}

\begin{exercise}
\label{hodge-theory-exercise-four-hodge-components}
Prove that the de Rham differential on an
almost complex manifold has at most four
nonzero Hodge components:
$$
d=d^{2,-1}+d^{1,0}+d^{0,1}+d^{-1,2}.
$$
\end{exercise}

\begin{exercise}
\label{hodge-theory-exercise-cartan-formulas}
Let $\theta$ be a one-form on a smooth
manifold $M$, and let $X,Y$ be vector
fields.

\begin{enumerate}
\item Let
$$
\iota_X:\Lambda^i(M)
\longrightarrow\Lambda^{i-1}(M)
$$
be contraction with $X$. Prove Cartan's
formula
$$
\mathcal{L}_X=\iota_Xd+d\iota_X.
$$

\item Prove that
$$
\begin{aligned}
X(\theta(Y))
&=\theta([X,Y])\\
&\quad+Y(\theta(X))\\
&\quad+d\theta(X,Y).
\end{aligned}
$$

\item Deduce the equivalent identity
$$
\begin{aligned}
Y(\theta(X))
&=-\theta([X,Y])\\
&\quad+X(\theta(Y))\\
&\quad-d\theta(X,Y).
\end{aligned}
$$

\item Prove the formula
$$
\begin{aligned}
d\theta(X,Y)
&=X(\theta(Y))-Y(\theta(X))\\
&\quad-\theta([X,Y]).
\end{aligned}
$$
\end{enumerate}
\end{exercise}

For the first part, use the general
properties of derivations. Deduce the later
identities successively from Cartan's
formula.

\begin{exercise}
\label{hodge-theory-exercise-integrability-components}
Let $(M,I)$ be an almost complex manifold
and write
$$
d=d^{2,-1}+d^{1,0}+d^{0,1}+d^{-1,2}.
$$

\begin{enumerate}
\item Prove that $d^{2,-1}$ and
$d^{-1,2}$ are $C^\infty(M)$-linear.

\item Suppose that
$$
d^{2,-1}=d^{-1,2}=0.
$$
Prove that $I$ is formally integrable.

\item Suppose that
$$
(d^{0,1})^2=0.
$$
Prove that $I$ is formally integrable.
\end{enumerate}
\end{exercise}

Use the formula in Exercise
\ref{hodge-theory-exercise-cartan-formulas} for the
second part and $d^2=0$ for the third.

\begin{exercise}
\label{hodge-theory-exercise-nijenhuis-hodge-component}
Let $(M,I)$ be an almost complex manifold.

\begin{enumerate}
\item Prove that
$$
d^{-1,2}:\Lambda^{1,0}(M)
\longrightarrow\Lambda^{0,2}(M)
$$
is dual to the conjugate Nijenhuis tensor
$$
\overline{N}:\Lambda^2(T^{0,1}M)
\longrightarrow T^{1,0}M.
$$

\item Prove that
$$
d^{-1,2}(d\varphi)=0
$$
for every holomorphic function $\varphi$.
\end{enumerate}
\end{exercise}

\begin{exercise}
\label{hodge-theory-exercise-surjective-nijenhuis}
Let $(M,I)$ be an almost complex manifold.
Suppose that the Nijenhuis tensor
$$
N:\Lambda^2(T^{1,0}M)
\longrightarrow T^{0,1}M
$$
is surjective. Prove that $(M,I)$ has no
nonconstant local holomorphic functions.
\end{exercise}

Use Exercise
\ref{hodge-theory-exercise-nijenhuis-hodge-component}.

\begin{definition}
\label{hodge-theory-definition-left-invariant-complex}
Let $G$ be a real Lie group with Lie algebra
$\mathfrak{g}=T_eG$. The {\it left action}
of $G$ on itself is
$$
L_g(x)=gx.
$$
An almost complex structure on $G$ is
{\it left invariant} if it is preserved by
$L_g$ for every $g\in G$.
\end{definition}

\begin{exercise}
\label{hodge-theory-exercise-left-invariant-integrable}
Let $G$ be a real Lie group, let
$\mathfrak{g}=T_eG$, and set
$$
\mathfrak{g}_{\mathbb{C}}
=\mathfrak{g}\otimes_{\mathbb{R}}
\mathbb{C}.
$$

\begin{enumerate}
\item Prove that every complex structure
operator on $\mathfrak{g}$ extends uniquely
to a left-invariant almost complex structure
on $G$.

\item Let $I$ be a complex structure on
$\mathfrak{g}$ and write
$$
\mathfrak{g}_{\mathbb{C}}
=\mathfrak{g}^{1,0}
\oplus\mathfrak{g}^{0,1}.
$$
Prove that the induced left-invariant
structure is formally integrable if and
only if
$$
\mathfrak{g}^{1,0}
\subset\mathfrak{g}_{\mathbb{C}}
$$
is a Lie subalgebra.
\end{enumerate}
\end{exercise}

\begin{exercise}[Optional]
\label{hodge-theory-exercise-left-invariant-no-functions}
Construct a left-invariant almost complex
structure $I$ on a Lie group $G$ such that
$(G,I)$ has no nonconstant local
holomorphic functions.
\end{exercise}

Use Exercise
\ref{hodge-theory-exercise-surjective-nijenhuis}.

\begin{exercise}
\label{hodge-theory-exercise-holomorphic-generators}
Let $(M,I)$ be an almost complex manifold.
Suppose that $\Lambda^{1,0}(M)$ is generated
by differentials of holomorphic functions.
Without using Theorem
\ref{hodge-theory-theorem-newlander-nirenberg}, prove
that $I$ is integrable.
\end{exercise}

\section{Lecture 4: K\"ahler manifolds}
\label{hodge-theory-section-lecture-four-kahler}

\noindent
Note: in this lecture we proved
the existence on Fubini-Study metric
on $\mathbb{C}P^n$ (and apparently also
in any symmetric space)
using $U(n+1)$-invariance.
Namely, a theorem states that there
is a bijection between $G$-invariant
metrics on $M$ and
$\operatorname{St}_x(G)$-invariant
metrics on $T_xM$.

For symmetric spaces, there are no
odd invariant tensors,
but $d \omega$ would be such.
And thus $\omega$ is closed.

We also proved that a complex
submanifold of a Kähler manifold is also
Kähler.

\medskip\noindent
{\bf }This lecture was given on August 21,
2026.

\medskip\noindent
{\bf Homogeneous spaces.}

\begin{definition}
\label{hodge-theory-definition-lie-group-action}
A {\it Lie group} is a smooth manifold with
a group structure for which multiplication
and inversion are smooth. A
{\it smooth action} of a Lie group $G$ on a
manifold $M$ is a smooth map
$$
G\times M\longrightarrow M
$$
which is a group action.
\end{definition}

\begin{definition}
\label{hodge-theory-definition-homogeneous-space}
A $G$-manifold $M$ is a
{\it homogeneous space} if the action is
transitive. For $x\in M$, the
{\it stabilizer}, or
{\it isotropy subgroup}, is
$$
\operatorname{Stab}_x(G)
=\{g\in G\mid g(x)=x\}.
$$
\end{definition}

\begin{proposition}
\label{hodge-theory-proposition-homogeneous-quotient}
If $G$ acts transitively on $M$, then
$$
M\cong G/H,
$$
where $H=\operatorname{Stab}_x(G)$ for any
chosen point $x\in M$.
\end{proposition}

\begin{proof}
The surjective orbit map
$$
G\longrightarrow M,
\qquad
g\longmapsto g(x),
$$
has fibers equal to the left cosets of $H$.
\end{proof}

If $g(x)=y$, the stabilizers of $x$ and $y$
are conjugate.

\begin{definition}
\label{hodge-theory-definition-isotropy-representation}
The {\it isotropy representation} at
$x\in G/H$ is the natural action of
$\operatorname{Stab}_x(G)$ on $T_xM$.
\end{definition}

\begin{definition}
\label{hodge-theory-definition-invariant-tensor}
A tensor on a homogeneous space $G/H$ is
{\it invariant} if it is preserved by the
diffeomorphisms induced by all elements of
$G$.
\end{definition}

An isotropy-invariant tensor $\Phi_x$ on
$T_xM$ can be transported to any
$y=g(x)$ by the action of $g$. This is
independent of the choice of $g$, because
two choices differ by an element of the
stabilizer.

\begin{theorem}
\label{hodge-theory-theorem-invariant-isotropy-tensors}
For a homogeneous space $M=G/H$, invariant
tensors on $M$ are in bijection with
isotropy-invariant tensors on $T_xM$.
\end{theorem}

\medskip\noindent
{\bf Hermitian and K\"ahler manifolds.}

\begin{definition}
\label{hodge-theory-definition-hermitian-metric-manifold}
A Riemannian metric $g$ on an almost complex
manifold is {\it Hermitian} if
$$
g(IX,IY)=g(X,Y).
$$
\end{definition}

Every Riemannian metric $g$ produces a
Hermitian metric
$$
h(X,Y)=g(X,Y)+g(IX,IY).
$$

\begin{definition}
\label{hodge-theory-definition-hermitian-form-manifold}
The {\it Hermitian form} of an almost
Hermitian manifold $(M,I,g)$ is
$$
\omega(X,Y)=g(X,IY).
$$
\end{definition}

The Hermitian form is skew-symmetric and
has Hodge type $(1,1)$. As in the vector
space case, any two members of
$(I,g,\omega)$ determine the third.

\begin{definition}
\label{hodge-theory-definition-kahler-manifold}
A complex Hermitian manifold $(M,I,g)$ is
{\it K\"ahler} if its Hermitian form is
closed. In this case $\omega$ is the
{\it K\"ahler form}, and
$$
[\omega]\in H^2(M,\mathbb{R})
$$
is its {\it K\"ahler class}.
\end{definition}

The slides recalled Erich K\"ahler
(1906--2000) and Andr\'e Weil
(1906--1998), including a childhood
photograph of Andr\'e with Simone Weil.

\medskip\noindent
{\bf Invariant metrics.}

\begin{theorem}
\label{hodge-theory-theorem-transitive-sphere-form}
Let $G$ act linearly on $V$ and transitively
on the unit sphere of a positive definite
quadratic form $g$. Every $G$-invariant
symmetric bilinear form on $V$ is a scalar
multiple of $g$.
\end{theorem}

\begin{proof}
The form $g$ is invariant because $G$
preserves its unit sphere. If $g'$ is
another invariant quadratic form, then
$$
x\longmapsto\frac{g'(x,x)}{g(x,x)}
$$
is constant on the unit sphere by
transitivity and invariant under nonzero
real scaling. It is therefore constant on
$V\setminus\{0\}$.
\end{proof}

\begin{exercise}
\label{hodge-theory-exercise-transitive-irreducible}
Let $V$ be a representation of $G$. Suppose
that $G$ acts transitively on a sphere in
$V$. Prove that $V$ is irreducible.
\end{exercise}

\begin{exercise}
\label{hodge-theory-exercise-real-schur-form}
Prove the following real form of Schur's
lemma. If $V$ is an irreducible real
representation of $G$ and $g$ is an
invariant positive definite symmetric
bilinear form, then every invariant
symmetric bilinear form is proportional to
$g$.
\end{exercise}

\medskip\noindent
{\bf The Fubini--Study metric.}

The unitary group $U(n+1)$ acts transitively
on $\mathbb{C}P^n$, and the stabilizer of a
point is $U(n)\times U(1)$.

\begin{theorem}
\label{hodge-theory-theorem-fubini-study}
There is a $U(n+1)$-invariant Riemannian
metric on $\mathbb{C}P^n$. It is unique up
to a positive scalar and is K\"ahler.
\end{theorem}

\begin{proof}
At a point, the tangent space is the standard
representation of $U(n)$. Its standard
positive form is isotropy invariant and
therefore extends to an invariant metric by
Theorem
\ref{hodge-theory-theorem-invariant-isotropy-tensors}.
The isotropy group acts transitively on the
unit sphere, so uniqueness follows from
Theorem \ref{hodge-theory-theorem-transitive-sphere-form}.

Let $\omega$ be the Hermitian form of this
metric. The three-form $d\omega$ is
isotropy invariant. Since the isotropy group
contains $-\operatorname{Id}$, every
invariant tensor of odd covariant degree
vanishes. Thus $d\omega=0$.
\end{proof}

\begin{definition}
\label{hodge-theory-definition-fubini-study}
The invariant metric in Theorem
\ref{hodge-theory-theorem-fubini-study} is the
{\it Fubini--Study metric}; its Hermitian
form is the {\it Fubini--Study form}.
\end{definition}

\medskip\noindent
{\bf Projective manifolds.}

\begin{definition}
\label{hodge-theory-definition-complex-submanifold}
A smooth submanifold $X\subset M$ of a
complex manifold is a
{\it complex submanifold} if
$$
I(TX)\subset TX.
$$
The inclusion is then a
{\it complex embedding}.
\end{definition}

\begin{definition}
\label{hodge-theory-definition-projective-manifold}
A complex manifold is {\it projective} if
it admits a complex embedding into some
$\mathbb{C}P^n$.
\end{definition}

A complex submanifold of a K\"ahler manifold
is K\"ahler: restrict both the Hermitian
metric and its closed Hermitian form. Hence
every projective manifold is K\"ahler.

\begin{definition}
\label{hodge-theory-definition-projective-algebraic}
A subvariety of $\mathbb{C}P^n$ is
{\it complex algebraic} if it is the common
zero locus of homogeneous polynomials.
\end{definition}

\begin{theorem}[Chow]
\label{hodge-theory-theorem-chow-projective}
Every complex submanifold of
$\mathbb{C}P^n$ is complex algebraic.
\end{theorem}

\begin{definition}
\label{hodge-theory-definition-integral-kahler-class}
A K\"ahler manifold has an
{\it integral K\"ahler class} if its K\"ahler
class lies in the image of
$$
H^2(M,\mathbb{Z})
\longrightarrow H^2(M,\mathbb{R}).
$$
\end{definition}

Since
$$
H^2(\mathbb{C}P^n,\mathbb{R})
\cong\mathbb{R},
$$
the Fubini--Study form can be scaled to
have integral cohomology class. Its
restriction shows that every projective
manifold has a K\"ahler metric with integral
K\"ahler class.

\begin{theorem}[Kodaira embedding]
\label{hodge-theory-theorem-kodaira-embedding}
Every compact K\"ahler manifold with an
integral K\"ahler class is projective.
\end{theorem}

The slides conclude with the hierarchy of
classes of almost complex manifolds. Every
projective manifold is both K\"ahler and
Moishezon. Every K\"ahler manifold is both
complex and symplectic. Every Moishezon
manifold is complex. Finally, complex and
symplectic manifolds both carry almost
complex structures. The slides refer to
Moishezon manifolds as algebraic spaces.

\section{Lecture 5: the Levi--Civita
connection}
\label{hodge-theory-section-lecture-five-levi-civita}

\noindent
Maybe the highlight of this lecture
(to which I arrived late)
was the proof of the existence
and uniqueness of the LC connection
via a non-conventional proof:
existence of metric connections is easy,
and each is determined by its torsion
via an isomorphism proved
combinatorially, which is the heart
of the proof. Taking preimage of zero
under such isomorphism we obtain
a metric torsion-free connection.

\medskip\noindent
This lecture was given on August 24, 2026.

\medskip\noindent
{\bf Connections and parallel transport.}

\begin{definition}
\label{hodge-theory-definition-vector-bundle-connection}
A {\it connection} on a vector bundle $B$
over $M$ is an operator
$$
\nabla:\Gamma(B)
\longrightarrow
\Gamma(B\otimes\Lambda^1M)
$$
satisfying
$$
\nabla(fb)=b\otimes df+f\nabla b.
$$
For a vector field $X$, contraction with
$X$ gives the covariant derivative
$\nabla_Xb$.
\end{definition}

A connection on $B$ induces a connection on
$B^*$, characterized by
$$
\begin{aligned}
d\left\langle{}b,\beta\right\rangle{}
&=\left\langle{}\nabla b,
\beta\right\rangle{}\\
&\quad+
\left\langle{}b,
\nabla\beta\right\rangle{}.
\end{aligned}
$$
The same symbol $\nabla$ is conventionally
used for both. The Leibniz rule
$$
\begin{aligned}
\nabla(b_1\otimes b_2)
&=(\nabla b_1)\otimes b_2\\
&\quad+b_1\otimes\nabla b_2
\end{aligned}
$$
then defines connections on all tensor
bundles built from $B$ and $B^*$.

\begin{theorem}
\label{hodge-theory-theorem-parallel-section-line}
Let $B$ be a vector bundle with connection
over the real line. For each $x$ and each
$b_x\in B_x$, there is a unique parallel
section $b$ satisfying
$$
\nabla b=0,
\qquad
b|_x=b_x.
$$
\end{theorem}

\begin{proof}
In a trivialization this is the existence
and uniqueness theorem for the ordinary
differential equation
$$
\frac{db}{dt}+A(t)b=0.
$$
\end{proof}

\begin{definition}
\label{hodge-theory-definition-parallel-transport}
Let
$$
\gamma:[0,1]\longrightarrow M
$$
be a smooth path from $x$ to $y$. The
{\it parallel transport} of $b_x\in B_x$
along $\gamma$ is the vector $b_y\in B_y$
obtained by solving
$$
\nabla_{\dot\gamma}b=0,
\qquad
b(0)=b_x,
$$
and taking $b_y=b(1)$.
\end{definition}

\begin{definition}
\label{hodge-theory-definition-holonomy-group}
The {\it holonomy group} of a connection at
$x$ is the group of automorphisms of $B_x$
obtained by parallel transport around all
loops based at $x$.
\end{definition}

\medskip\noindent
{\bf Lie algebras and tensors.}

If $V$ is a representation of a Lie algebra
$\mathfrak{g}$, then $V^*$ becomes a
representation by
$$
\left\langle{}a(x),\lambda
\right\rangle{}
=-
\left\langle{}x,a(\lambda)
\right\rangle{}.
$$
Tensor products become representations by
the Leibniz rule
$$
a(x\otimes y)
=a(x)\otimes y+x\otimes a(y).
$$
Thus $\mathfrak{g}$ acts on all tensor
powers of $V$ and $V^*$.

\begin{definition}
\label{hodge-theory-definition-lie-preserves-tensor}
The Lie algebra $\mathfrak{g}$
{\it preserves a tensor} $\Phi$ if
$$
a(\Phi)=0
$$
for every $a\in\mathfrak{g}$.
\end{definition}

\begin{definition}
\label{hodge-theory-definition-orthogonal-lie-algebra}
Let $h$ be a nondegenerate symmetric
bilinear form on $V$. The
{\it orthogonal Lie algebra}
$\mathfrak{so}(V,h)$ consists of the
endomorphisms $a$ satisfying
$$
h(a(x),y)=-h(x,a(y)).
$$
\end{definition}

In an $h$-orthonormal basis these are the
skew-symmetric matrices.

\begin{proposition}
\label{hodge-theory-proposition-two-forms-so}
Use $h$ to identify $V$ with $V^*$ and let
$$
\tau:V^*\otimes V^*
\longrightarrow V^*\otimes V
$$
be the resulting isomorphism. Then
$$
\tau(\Lambda^2V^*)
=\mathfrak{so}(V,h).
$$
\end{proposition}

\begin{proof}
The bilinear form corresponding to
$a\in\operatorname{End}(V)$ is
$$
(x,y)\longmapsto h(a(x),y).
$$
It is skew-symmetric precisely when $a$
satisfies the defining identity for
$\mathfrak{so}(V,h)$.
\end{proof}

\begin{definition}
\label{hodge-theory-definition-parallel-tensor}
A tensor $\Psi$ built from a vector bundle
$B$ and its dual is {\it parallel}, or
{\it preserved by $\nabla$}, if
$$
\nabla\Psi=0.
$$
\end{definition}

This condition holds if and only if every
parallel transport preserves $\Psi$.

\begin{lemma}[Corollary]
\label{hodge-theory-lemma-parallel-holonomy}
A tensor is parallel if and only if it is
invariant under the holonomy group.
\end{lemma}

\begin{definition}
\label{hodge-theory-definition-orthogonal-connection}
Let $h$ be a positive definite form on $B$.
A connection is {\it orthogonal} if
$$
\nabla h=0.
$$
\end{definition}

\begin{definition}
\label{hodge-theory-definition-unitary-connection}
Let $(B,I,h)$ be a Hermitian vector bundle.
A connection is {\it unitary}, or
{\it Hermitian}, if
$$
\nabla I=0,
\qquad
\nabla h=0.
$$
\end{definition}

\medskip\noindent
{\bf Torsors and affine spaces.}

\begin{definition}
\label{hodge-theory-definition-torsor}
A {\it torsor} over a group $G$ is a set
$X$ with a free and transitive $G$-action.
\end{definition}

\begin{definition}
\label{hodge-theory-definition-torsor-morphism}
A {\it morphism of torsors} from
$(X,G,\rho)$ to $(X',G',\rho')$ is a pair
of maps
$$
\Psi_X:X\longrightarrow X',
\qquad
\Psi_G:G\longrightarrow G',
$$
where $\Psi_G$ is a homomorphism and
$$
\Psi_X(\rho(g,x))
=\rho'(\Psi_G(g),\Psi_X(x)).
$$
\end{definition}

These objects and morphisms form the
category of torsors.

\begin{definition}
\label{hodge-theory-definition-affine-space}
An {\it affine space} is a torsor over a
vector space $V$. The vector space $V$ is
its {\it linearization}.
\end{definition}

If $\nabla$ and $\nabla_1$ are connections
on $B$, their difference is an
$\operatorname{End}(B)$-valued one-form.
Conversely, adding any
$$
A\in\Lambda^1M\otimes
\operatorname{End}(B)
$$
to a connection produces another
connection. Hence the space of connections
on $B$ is an affine space over
$$
\Lambda^1M\otimes\operatorname{End}(B).
$$

\begin{proposition}
\label{hodge-theory-proposition-affine-orthogonal}
Let $B$ carry a scalar product. The space of
orthogonal connections on $B$ is an affine
space over
$$
\Lambda^1M\otimes\mathfrak{so}(B).
$$
\end{proposition}

\begin{proof}
The difference $A$ of two connections acts
on a bilinear form $h$ by
$$
\begin{aligned}
(A(h))(x,y)
&=-h(A(x),y)\\
&\quad-h(x,A(y)).
\end{aligned}
$$
Thus two orthogonal connections differ by a
one-form with values in
$\mathfrak{so}(B)$, and every such change
preserves orthogonality.
\end{proof}

More generally, let $B$ carry a tensor
$\Phi$ and let
$$
\mathfrak{g}\subset\operatorname{End}(B)
$$
be the Lie algebra preserving $\Phi$. The
space of connections preserving $\Phi$ is
an affine space over
$$
\Lambda^1M\otimes\mathfrak{g}.
$$

\medskip\noindent
{\bf The de Rham differential.}

\begin{definition}
\label{hodge-theory-definition-de-rham-algebra-course}
The {\it de Rham algebra} of $M$ is the
graded algebra
$$
\Lambda^*M
=\Gamma(\Lambda^*T^*M)
$$
with the wedge product. Its degree $i$
elements are the
{\it differential $i$-forms}.
\end{definition}

In degree zero one has
$$
\Lambda^0M=C^\infty(M).
$$

\begin{theorem}
\label{hodge-theory-theorem-de-rham-characterization}
There is a unique sequence of operators
$$
C^\infty(M)\xrightarrow{d}\Lambda^1M
\xrightarrow{d}\Lambda^2M
\xrightarrow{d}\cdots
$$
such that $d$ is the ordinary differential
on functions, $d^2=0$, and, for a form
$\eta$ of degree $p$,
$$
d(\eta\wedge\xi)
=d\eta\wedge\xi
+(-1)^p\eta\wedge d\xi.
$$
\end{theorem}

\begin{definition}
\label{hodge-theory-definition-de-rham-differential}
The operator $d$ in Theorem
\ref{hodge-theory-theorem-de-rham-characterization} is
the {\it de Rham differential}.
\end{definition}

\begin{proposition}[Cartan formula]
\label{hodge-theory-proposition-cartan-one-form}
For a one-form $\eta$ and vector fields
$X,Y$, one has
$$
\begin{aligned}
d\eta(X,Y)
&=\mathcal{L}_X(\eta(Y))\\
&\quad-\mathcal{L}_Y(\eta(X))\\
&\quad-\eta([X,Y]).
\end{aligned}
$$
\end{proposition}

\begin{proof}
The right hand side is alternating and
satisfies the Leibniz rule in $\eta$. It
vanishes when $\eta=df$, because
$$
\mathcal{L}_{[X,Y]}f
=\mathcal{L}_X\mathcal{L}_Yf
-\mathcal{L}_Y\mathcal{L}_Xf.
$$
These properties identify it with the de
Rham differential on one-forms.
\end{proof}

\medskip\noindent
{\bf Torsion.}

\begin{definition}
\label{hodge-theory-definition-torsion-vector}
Let $\nabla$ be a connection on $TM$. Its
{\it torsion} is the tensor
$$
T_\nabla(X,Y)
=\nabla_XY-\nabla_YX-[X,Y].
$$
\end{definition}

The induced connection on $T^*M$ gives a
bundle map
$$
\operatorname{Alt}\circ\nabla:
\Lambda^1M\longrightarrow\Lambda^2M.
$$
Cartan's formula gives
$$
\begin{aligned}
(d\eta-&\operatorname{Alt}
\circ\nabla\eta)(X,Y)\\
&=\eta(T_\nabla(X,Y)).
\end{aligned}
$$
Thus the operator
$$
d-\operatorname{Alt}\circ\nabla
$$
on one-forms is dual to the usual
vector-valued torsion tensor. In particular,
it is $C^\infty(M)$-linear.

\begin{definition}
\label{hodge-theory-definition-levi-civita}
Let $(M,g)$ be a Riemannian manifold. A
connection on $TM$ is
{\it Levi--Civita} if it is orthogonal and
torsion-free.
\end{definition}

\begin{theorem}
\label{hodge-theory-theorem-fundamental-riemannian}
Every Riemannian manifold has a unique
Levi--Civita connection.
\end{theorem}

The slides recalled Gregorio
Ricci-Curbastro (1853--1925) and his former
student Tullio Levi-Civita (1873--1941).
Their joint work pioneered tensor calculus.
Ricci-Curbastro signed that publication with
the shortened surname Ricci, a convention
which has caused some later confusion.

\medskip\noindent
{\bf Linearized torsion.}

The space of all connections is affine, and
torsion is an affine map from this space to
$$
\Lambda^2M\otimes TM.
$$
If a connection is changed by
$$
A\in\Lambda^1M\otimes
\operatorname{End}(TM),
$$
its torsion changes by the
antisymmetrization
$$
\begin{aligned}
\operatorname{Alt}_{12}(A)(X,Y)
&=A(X)Y-A(Y)X.
\end{aligned}
$$

\begin{definition}
\label{hodge-theory-definition-linearized-torsion}
The {\it linearized torsion} is the map
$$
\begin{aligned}
T_{\rm lin}:\Lambda^1M\otimes
\operatorname{End}(TM)
&\longrightarrow
\Lambda^2M\otimes TM,\\
A&\longmapsto
\operatorname{Alt}_{12}(A).
\end{aligned}
$$
\end{definition}

\begin{proposition}
\label{hodge-theory-proposition-orthogonal-exists}
Every vector bundle with a scalar product
admits an orthogonal connection.
\end{proposition}

\begin{proof}
Choose a trivializing cover with local
orthonormal frames, and on each member choose
a connection preserving its frame. A
partition of unity gives an affine
combination of these local connections. The
result is a global orthogonal connection.
\end{proof}

\begin{proof}[Proof of Theorem
\ref{hodge-theory-theorem-fundamental-riemannian}]
Choose an orthogonal connection $\nabla_0$.
By Proposition
\ref{hodge-theory-proposition-affine-orthogonal}, the
space of orthogonal connections is affine
over
$$
\Lambda^1M\otimes\mathfrak{so}(TM).
$$
The metric identifies this bundle with
$$
\Lambda^1M\otimes\Lambda^2M.
$$
On this affine subspace, the linearized
torsion is
$$
\begin{aligned}
T_{\rm lin}:\Lambda^1M\otimes\Lambda^2M
&\longrightarrow
\Lambda^2M\otimes\Lambda^1M,\\
A&\longmapsto
\operatorname{Alt}_{12}(A).
\end{aligned}
$$
It is an isomorphism. The two bundles have
the same rank, so it is enough to prove
injectivity. In covariant form, an element
of the kernel is symmetric in its first two
arguments and skew-symmetric in its last
two. The proof is a permutation argument.
Let $\sigma$ be the cyclic permutation
$$
\sigma(x,y,z)=(y,z,x).
$$
The two symmetries give
$$
\sigma(\eta)=-\eta.
$$
Since $\sigma^3=\operatorname{Id}$, applying
the permutation three times gives
$$
\eta=-\eta,
$$
so $\eta=0$.

It follows that an orthogonal connection is
uniquely determined by its torsion. In
particular,
$$
\nabla
=\nabla_0
-T_{\rm lin}^{-1}(T_{\nabla_0})
$$
is the unique orthogonal connection with
zero torsion.
\end{proof}

\section{Lecture 9: the Bismut
connection}
\label{hodge-theory-section-lecture-nine-bismut}

\noindent
This lecture was given on April 9, 2025.

\medskip\noindent
{\bf Intrinsic torsion.}

Let $\Phi$ be a tensor on a manifold and
let $\nabla$ preserve $\Phi$. Denote by
$$
\mathfrak{a}(M)
\subset\operatorname{End}(TM)
$$
the bundle of Lie algebras consisting of
the endomorphisms $A$ satisfying
$$
A(\Phi)=0.
$$
A second connection $\nabla_1$ preserves
$\Phi$ if and only if
$$
\nabla-\nabla_1
\in\Lambda^1M\otimes\mathfrak{a}(M).
$$
Thus the connections preserving $\Phi$
form an affine space over
$$
\Lambda^1M\otimes\mathfrak{a}(M).
$$

\begin{definition}
\label{hodge-theory-definition-intrinsic-torsion-space}
The {\it space of intrinsic torsion} for
$\mathfrak{a}(M)$-valued connections is
the quotient bundle
$$
\mathcal{T}_{\mathfrak{a}}
=\frac{\Lambda^2M\otimes TM}
{\operatorname{Alt}_{12}
(\Lambda^1M\otimes\mathfrak{a}(M))}.
$$
\end{definition}

\begin{definition}
\label{hodge-theory-definition-intrinsic-torsion}
Assume that $\mathfrak{a}(M)$ is a vector
bundle. The {\it intrinsic torsion} of a
connection preserving $\Phi$ is the image
of its torsion in
$\mathcal{T}_{\mathfrak{a}}$.
\end{definition}

\begin{theorem}
\label{hodge-theory-theorem-intrinsic-torsion}
The intrinsic torsion is independent of
the choice of connection preserving
$\Phi$. There is a torsion-free connection
preserving $\Phi$ if and only if the
intrinsic torsion vanishes.
\end{theorem}

\begin{proof}
If two such connections differ by
$$
A\in\Lambda^1M\otimes\mathfrak{a}(M),
$$
their torsions differ by
$\operatorname{Alt}_{12}(A)$. This proves
the first assertion. The second follows by
changing a connection by a preimage of its
torsion under the linearized torsion map.
\end{proof}

\medskip\noindent
{\bf Decomposition of third tensor powers.}

For $1\leq i<j\leq n$, let
$$
\operatorname{Alt}_{ij},
\operatorname{Sym}_{ij}:
V^{\otimes n}\longrightarrow V^{\otimes n}
$$
be antisymmetrization and symmetrization in
the $i$-th and $j$-th factors.

\begin{proposition}
\label{hodge-theory-proposition-third-tensor}
For every vector space $V$, one has
$$
\begin{aligned}
V^{\otimes3}
={}&\operatorname{Sym}^3V
\oplus\Lambda^3V\\
&\oplus\operatorname{Alt}_{12}
(V\otimes\operatorname{Sym}^2V)\\
&\oplus\operatorname{Sym}_{12}
(V\otimes\Lambda^2V).
\end{aligned}
$$
\end{proposition}

\begin{proof}
Every tensor $x$ satisfies
$$
x=\operatorname{Sym}_{ij}(x)
+\operatorname{Alt}_{ij}(x).
$$
First apply this to the last two factors to
obtain
$$
V^{\otimes3}
=V\otimes\operatorname{Sym}^2V
\oplus V\otimes\Lambda^2V.
$$
Then apply it to the first two factors.
The space
$$
\operatorname{Sym}_{12}
(V\otimes\operatorname{Sym}^2V)
$$
is symmetric in all three factors, while
$$
\operatorname{Alt}_{12}
(V\otimes\Lambda^2V)
$$
is antisymmetric in all three factors.
These are $\operatorname{Sym}^3V$ and
$\Lambda^3V$, respectively.
\end{proof}

\begin{lemma}[Corollary]
\label{hodge-theory-lemma-symplectic-torsion-quotient}
There is a natural isomorphism
$$
\frac{\Lambda^2V\otimes V}
{\operatorname{Alt}_{12}
(V\otimes\operatorname{Sym}^2V)}
\cong\Lambda^3V.
$$
\end{lemma}

\begin{proof}
Antisymmetrizing the first two factors in
Proposition
\ref{hodge-theory-proposition-third-tensor}
gives
$$
\begin{aligned}
\Lambda^2V\otimes V
={}&\Lambda^3V\\
&\oplus\operatorname{Alt}_{12}
(V\otimes\operatorname{Sym}^2V).
\end{aligned}
$$
\end{proof}

\medskip\noindent
{\bf Symplectic connections.}

For a connection on $T^*M$, define
$$
\tau_\nabla(\eta)
=\operatorname{Alt}(\nabla\eta)-d\eta.
$$
On one-forms this is the tensor dual to the
torsion. It extends to differential forms
as a graded derivation. Thus, if $\lambda$
has degree $p$, then
$$
\begin{aligned}
\tau_\nabla(\lambda\wedge\mu)
={}&\tau_\nabla(\lambda)\wedge\mu\\
&+(-1)^p\lambda\wedge
\tau_\nabla(\mu).
\end{aligned}
$$

\begin{definition}
\label{hodge-theory-definition-almost-symplectic}
An {\it almost symplectic structure} is a
nondegenerate two-form.
\end{definition}

\begin{definition}
\label{hodge-theory-definition-symplectic-connection}
Let $(M,\omega)$ be almost symplectic. A
{\it symplectic connection} is a connection
on $TM$ satisfying
$$
\nabla\omega=0.
$$
\end{definition}

Such connections exist: local connections
preserving symplectic frames can be joined
by an affine partition-of-unity argument.

\begin{lemma}
\label{hodge-theory-lemma-symplectic-torsion-form}
Let $\nabla$ be a symplectic connection and
use $\omega$ to identify its torsion with
$$
T_\omega
\in\Lambda^2M\otimes\Lambda^1M.
$$
Then
$$
\operatorname{Alt}(T_\omega)
=-\frac{1}{2}d\omega.
$$
\end{lemma}

\begin{proof}
Since $\nabla\omega=0$, one has
$$
\tau_\nabla(\omega)=-d\omega.
$$
Apply the graded derivation identity to the
two covariant factors of $\omega$. The two
contractions with the torsion have the same
antisymmetrization. Each is
$\operatorname{Alt}(T_\omega)$, and hence
$$
\tau_\nabla(\omega)
=2\operatorname{Alt}(T_\omega).
$$
\end{proof}

\begin{theorem}
\label{hodge-theory-theorem-almost-symplectic-torsion}
Let $(M,\omega)$ be almost symplectic. Use
$\omega$ to identify the torsion of a
symplectic connection with
$$
T_\omega
\in\Lambda^2M\otimes\Lambda^1M.
$$
Then
$$
\operatorname{Alt}_{123}(T_\omega)
=-\frac{1}{2}d\omega.
$$
Conversely, every tensor
$$
u\in\Lambda^2M\otimes\Lambda^1M
$$
satisfying
$$
\operatorname{Alt}_{123}(u)
=-\frac{1}{2}d\omega
$$
is the torsion of a symplectic connection.
\end{theorem}

\begin{proof}
The symplectic Lie algebra is identified by
$\omega$ with symmetric two-tensors:
$$
\mathfrak{sp}(TM)
\cong\operatorname{Sym}^2(\Lambda^1M).
$$
Under this identification, linearized
torsion becomes
$$
\begin{aligned}
\operatorname{Alt}_{12}:
\Lambda^1M\otimes
\operatorname{Sym}^2(\Lambda^1M)
\longrightarrow{}&
\Lambda^2M\otimes\Lambda^1M.
\end{aligned}
$$
Its kernel is
$\operatorname{Sym}^3(\Lambda^1M)$.
Lemma
\ref{hodge-theory-lemma-symplectic-torsion-quotient}
therefore gives the exact sequence
$$
\begin{aligned}
0\longrightarrow{}&
\operatorname{Sym}^3(\Lambda^1M)\\
\longrightarrow{}&
\Lambda^1M\otimes
\operatorname{Sym}^2(\Lambda^1M)\\
\xrightarrow{\operatorname{Alt}_{12}}{}&
\Lambda^2M\otimes\Lambda^1M\\
\xrightarrow{\operatorname{Alt}_{123}}{}&
\Lambda^3M\longrightarrow0.
\end{aligned}
$$
Thus $\Lambda^3M$ is the space of
intrinsic torsion of symplectic
connections. Lemma
\ref{hodge-theory-lemma-symplectic-torsion-form}
identifies the intrinsic torsion with
$-d\omega/2$, proving both assertions.
\end{proof}

\medskip\noindent
{\bf Torsion of unitary connections.}

\begin{proposition}
\label{hodge-theory-proposition-unitary-torsion-types}
Let $(M,I,\omega)$ be a Hermitian complex
manifold and let $\nabla$ preserve $I$ and
$\omega$. Use the Riemannian metric to
regard the torsion as an element of
$$
\Lambda^2M\otimes\Lambda^1M.
$$
Then it belongs to
$$
\begin{aligned}
\mathcal{W}={}&
\Lambda^{2,0}(M)\otimes\Lambda^{0,1}(M)\\
&\oplus\Lambda^{0,2}(M)
\otimes\Lambda^{1,0}(M)\\
&\oplus\Lambda^{1,1}(M)
\otimes\Lambda^1M.
\end{aligned}
$$
\end{proposition}

\begin{proof}
Integrability of $I$ implies
$$
[T^{1,0}M,T^{1,0}M]\subset T^{1,0}M.
$$
Since $\nabla I=0$, one also has
$$
\nabla_X(T^{1,0}M)\subset T^{1,0}M
$$
for every vector field $X$. It follows that
the torsion lies in
$$
\begin{aligned}
&\Lambda^{2,0}(M)\otimes T^{1,0}M\\
&\quad\oplus\Lambda^{0,2}(M)
\otimes T^{0,1}M\\
&\quad\oplus\Lambda^{1,1}(M)\otimes TM.
\end{aligned}
$$
The Riemannian metric has type $(1,1)$, so
it identifies $T^{1,0}M$ with
$\Lambda^{0,1}(M)$ and $T^{0,1}M$ with
$\Lambda^{1,0}(M)$.
\end{proof}

The slides recalled Jean-Michel Bismut,
who was born in 1948.

\begin{theorem}[Bismut]
\label{hodge-theory-theorem-bismut-connection}
Let $(M,I,\omega)$ be a Hermitian complex
manifold. There is a unique connection
$\nabla$ preserving $I$ and $\omega$ whose
torsion is a three-form. Its torsion is
$$
T_\nabla=-\frac{1}{2}I(d\omega).
$$
\end{theorem}

\begin{definition}
\label{hodge-theory-definition-bismut-connection}
The connection in Theorem
\ref{hodge-theory-theorem-bismut-connection} is the
{\it Bismut connection}.
\end{definition}

\begin{proof}
There are two identifications
$$
TM\longrightarrow\Lambda^1M,
$$
one given by $g$ and the other by $\omega$.
Write $T_g$ and $T_\omega$ for the two
covariant forms of the same torsion. Define
$I_3$ by
$$
I_3(x\otimes y\otimes z)
=x\otimes y\otimes I(z).
$$
Then
$$
I_3(T_g)=T_\omega.
$$
By Theorem
\ref{hodge-theory-theorem-almost-symplectic-torsion},
$$
\operatorname{Alt}(I_3(T_g))
=-\frac{1}{2}d\omega.
$$
The left hand side is independent of the
unitary connection. Its linearization
therefore vanishes:
$$
\operatorname{Alt}\circ I_3\circ
T_{\rm lin}=0
$$
on
$$
\Lambda^1M\otimes\mathfrak{u}(TM).
$$

By Proposition
\ref{hodge-theory-proposition-unitary-torsion-types},
unitary torsion belongs to $\mathcal{W}$.
The linearized torsion map
$$
T_{\rm lin}:
\Lambda^1M\otimes\mathfrak{u}(TM)
\longrightarrow\mathcal{W}
$$
is injective by the same permutation
argument used for the Levi--Civita
connection. Dimension counting gives the
exact sequence
$$
\begin{aligned}
0\longrightarrow{}&
\Lambda^1M\otimes\mathfrak{u}(TM)\\
\xrightarrow{T_{\rm lin}}{}&
\mathcal{W}\\
\xrightarrow{\operatorname{Alt}\circ I_3}{}&
\Lambda^{2,1}(M)\oplus\Lambda^{1,2}(M)\\
\longrightarrow{}&0.
\end{aligned}
$$
Surjectivity of the second map follows by
antisymmetrizing tensors in
$I_3(\mathcal{W})$.

Let
$$
\mathcal{U}
=\Lambda^{2,1}(M)\oplus\Lambda^{1,2}(M)
$$
be the three-forms contained in
$\mathcal{W}$. On three-forms one has
$$
\operatorname{Alt}(I_3(\eta))
=\mathcal{W}_I(\eta),
$$
where the Weil operator is given by
$$
\begin{aligned}
\mathcal{W}_I(\eta)(x,y,z)
={}&\eta(Ix,y,z)\\
&+\eta(x,Iy,z)\\
&+\eta(x,y,Iz).
\end{aligned}
$$
Its restriction from $\mathcal{U}$ to
itself is bijective. Consequently there is
a unique $\sigma\in\mathcal{U}$ satisfying
$$
\operatorname{Alt}(I_3(\sigma))
=-\frac{1}{2}d\omega.
$$

Choose a unitary connection $\nabla_0$ and
let $T_0\in\mathcal{W}$ be its torsion.
Exactness gives a unique
$$
A\in\Lambda^1M\otimes\mathfrak{u}(TM)
$$
such that
$$
T_{\rm lin}(A)+T_0=\sigma.
$$
The connection $\nabla_0+A$ is unitary and
has totally antisymmetric torsion. The Weil
operator calculation on forms of types
$(2,1)$ and $(1,2)$ gives
$$
\sigma=-\frac{1}{2}I(d\omega).
$$
\end{proof}

\begin{lemma}[Corollary]
\label{hodge-theory-lemma-kahler-levi-civita-unitary}
On a K"ahler manifold, the Levi--Civita
connection satisfies
$$
\nabla I=0,
\qquad
\nabla\omega=0.
$$
\end{lemma}

\begin{proof}
If $d\omega=0$, the Bismut connection has
zero torsion. It is orthogonal, so uniqueness
of the Levi--Civita connection identifies
the two connections.
\end{proof}

There is also a direct intrinsic-torsion
proof. For a unitary connection $\nabla$ on
a K"ahler manifold, Theorem
\ref{hodge-theory-theorem-almost-symplectic-torsion}
gives
$$
\operatorname{Alt}(I_3(T_\nabla))=0.
$$
The exact sequence in the proof of Theorem
\ref{hodge-theory-theorem-bismut-connection} implies
that $T_\nabla$ is in the image of
$T_{\rm lin}$. Hence
$$
\nabla-T_{\rm lin}^{-1}(T_\nabla)
$$
is a torsion-free unitary connection and
therefore is the Levi--Civita connection.

\section{Lecture 9-bis: harmonic forms and
cohomology}
\label{hodge-theory-section-lecture-nine-bis-harmonic}

\noindent
In this lecture I learnt that
we can define harmonic forms on any
chain of vector spaces equipped with
an inner product. Define $d^*$
as always, satisfying 
$\langle d \alpha, \beta\rangle
=
\langle \alpha, d^*\beta\rangle$,
and the Laplacian, as always by
$d d^* + d^* d$.
The key observation is that
$$
H^k = \frac{\ker d}{\operatorname{Im} d}
\simeq
\ker d \cap \operatorname{Im} (d)^\perp
\simeq
\ker d \cap \ker d^*
\simeq
\ker \Delta,
$$
which says that evert cohomology
class is represented by a harmonic
vector.

\medskip\noindent
{\bf Laplacians on a complex.}

\begin{definition}
\label{hodge-theory-definition-complex-vector-spaces}
A {\it complex of vector spaces} is a
sequence of vector spaces and linear maps
$$
\cdots\xrightarrow{d}C^{i-1}
\xrightarrow{d}C^i
\xrightarrow{d}C^{i+1}
\xrightarrow{d}\cdots
$$
satisfying $d^2=0$. It is {\it finite} if
only finitely many of its vector spaces are
nonzero.
\end{definition}

\begin{definition}
\label{hodge-theory-definition-cohomology-complex}
The {\it cohomology} of $(C^*,d)$ in degree
$i$ is
$$
H^i(C^*)
=\frac{\ker(d|_{C^i})}{d(C^{i-1})}.
$$
\end{definition}

\begin{definition}
\label{hodge-theory-definition-laplacian-complex}
Suppose that every $C^i$ has a positive
definite inner product. If $d^*$ denotes
the adjoint of $d$, the {\it Laplacian} is
$$
\Delta_d=dd^*+d^*d.
$$
\end{definition}

The Laplacian commutes with $d$, and hence
is an endomorphism of the complex.

\begin{definition}
\label{hodge-theory-definition-harmonic-vector}
A vector in $C^i$ is {\it harmonic} if it
belongs to $\ker\Delta_d$.
\end{definition}

For every $x\in C^i$, one has
$$
\begin{aligned}
\left\langle{}x,\Delta_dx
\right\rangle{}
={}&\left\langle{}dx,dx
\right\rangle{}\\
&+\left\langle{}d^*x,d^*x
\right\rangle{}.
\end{aligned}
$$
Therefore
$$
\ker\Delta_d=\ker d\cap\ker d^*.
$$

The adjoint identities give
$$
\ker d=(\operatorname{im}d^*)^\perp,
\qquad
\ker d^*=(\operatorname{im}d)^\perp.
$$
When the relevant images are closed, they
also give
$$
\operatorname{im}d
=(\ker d^*)^\perp,
\qquad
\operatorname{im}d^*
=(\ker d)^\perp.
$$

\begin{lemma}[Corollary]
\label{hodge-theory-lemma-harmonic-representatives}
Suppose that the vector spaces $C^i$ are
finite-dimensional. More generally, suppose
that the images of the differentials are
closed in the relevant Hilbert spaces. Then
every cohomology class has a unique harmonic
representative, and
$$
H^*(C^*)\cong\ker\Delta_d.
$$
\end{lemma}

\begin{proof}
Orthogonal decomposition gives
$$
\begin{aligned}
\ker d
&=\operatorname{im}d
&\quad\oplus
(\ker d\cap(\operatorname{im}d)^\perp)\\
&=\operatorname{im}d
\oplus(\ker d\cap\ker d^*)\\
&=\operatorname{im}d
\oplus\ker\Delta_d.
\end{aligned}
$$
The quotient of $\ker d$ by
$\operatorname{im}d$ is therefore naturally
isomorphic to $\ker\Delta_d$.
\end{proof}

\medskip\noindent
{\bf Laplacians on differential forms.}

A metric $g$ on a vector space induces a
metric on every tensor power by
$$
\begin{aligned}
&g(x_1\otimes\cdots\otimes x_k,
y_1\otimes\cdots\otimes y_k)\\
&\qquad
=\prod_{j=1}^kg(x_j,y_j).
\end{aligned}
$$
Consequently a Riemannian metric on $M$
induces the $L^2$ inner product
$$
\left\langle{}\alpha,\beta
\right\rangle{}_{L^2}
=\int_M g(\alpha,\beta)
\operatorname{Vol}_M
$$
on differential forms.

\begin{definition}
\label{hodge-theory-definition-laplacian-forms}
The {\it Laplacian on differential forms}
on a Riemannian manifold is
$$
\Delta=dd^*+d^*d.
$$
\end{definition}

\begin{theorem}[Hodge theorem]
\label{hodge-theory-theorem-hodge-spectral}
Let $M$ be a compact Riemannian manifold.
The Hilbert space $L^2(\Lambda^*M)$ has an
orthonormal basis of eigenforms of
$\Delta$. Every eigenspace is
finite-dimensional, and the spectrum is
discrete. The inverse of $\Delta$ on
$\operatorname{im}\Delta$ is continuous in
the $L^2$ topology.
\end{theorem}

\begin{theorem}[Elliptic regularity]
\label{hodge-theory-theorem-laplacian-regularity}
Every $L^2$ eigenform of the Laplacian on a
compact Riemannian manifold is smooth.
\end{theorem}

The analogous spectral and regularity
statements hold for the Dolbeault Laplacian
on a compact Hermitian manifold.

The slides recalled Fritz Noether
(1884--1941), the brother of Emmy Noether.

\medskip\noindent
{\bf De Rham cohomology.}

\begin{definition}
\label{hodge-theory-definition-de-rham-cohomology}
The {\it de Rham cohomology} of $M$ is
$$
H^i(M)
=\frac{\ker(d|_{\Lambda^iM})}
{d(\Lambda^{i-1}M)}.
$$
\end{definition}

\begin{definition}
\label{hodge-theory-definition-harmonic-form}
A differential form $\alpha$ is
{\it harmonic} if
$$
\Delta\alpha=0.
$$
The space of harmonic $i$-forms is denoted
by $\mathcal{H}^i(M)$.
\end{definition}

A harmonic form is both $d$-closed and
$d^*$-closed. The natural map
$$
\mathcal{H}^i(M)\longrightarrow H^i(M)
$$
is injective. Indeed, if
$\alpha=d\beta$ is harmonic, then
$$
\begin{aligned}
\left\langle{}\alpha,\alpha
\right\rangle{}
&=\left\langle{}\alpha,d\beta
\right\rangle{}\\
&=\left\langle{}d^*\alpha,\beta
\right\rangle{}\\
&=0.
\end{aligned}
$$

\begin{theorem}[Hodge decomposition]
\label{hodge-theory-theorem-hodge-de-rham}
For a compact Riemannian manifold, the
natural map
$$
\mathcal{H}^i(M)\longrightarrow H^i(M)
$$
is an isomorphism.
\end{theorem}

\begin{proof}
The identities $d^2=0$ and $(d^*)^2=0$
imply
$$
[d,\Delta]=0.
$$
Thus $d$ preserves every eigenspace of the
Laplacian. Let
$$
\Lambda^*(M)
=\mathop{\widehat{\bigoplus}}_\lambda
\Lambda^*_{\lambda}(M)
$$
be the eigenspace decomposition. Each
$\Lambda^*_{\lambda}(M)$ is a complex under
$d$.

Suppose that $\lambda\neq0$ and that
$\eta\in\Lambda^*_{\lambda}(M)$ is closed.
Then
$$
\begin{aligned}
\lambda\eta
&=dd^*\eta+d^*d\eta\\
&=dd^*\eta.
\end{aligned}
$$
It follows that
$$
\eta=d(\lambda^{-1}d^*\eta).
$$
Hence the nonzero eigenspaces have no
cohomology. The zero eigenspace is
$\mathcal{H}^*(M)$, which proves the
claim.
\end{proof}

Poincar\'e duality follows by applying the
Hodge star to harmonic forms.

\medskip\noindent
{\bf Dolbeault cohomology.}

\begin{definition}
\label{hodge-theory-definition-dolbeault-cohomology}
The {\it Dolbeault cohomology} of a complex
manifold $M$ is
$$
H^{p,q}_{\bar\partial}(M)
=\frac{\ker(\bar\partial|
_{\Lambda^{p,q}(M)})}
{\bar\partial\Lambda^{p,q-1}(M)}.
$$
\end{definition}

On a Hermitian manifold, define
$$
\Delta_{\bar\partial}
=\bar\partial\bar\partial^*
+\bar\partial^*\bar\partial
$$
and
$$
\mathcal{H}^{p,q}_{\bar\partial}(M)
=\ker(\Delta_{\bar\partial}|
_{\Lambda^{p,q}(M)}).
$$

\begin{theorem}
\label{hodge-theory-theorem-hodge-dolbeault}
Let $M$ be a compact Hermitian complex
manifold. The natural map
$$
\mathcal{H}^{p,q}_{\bar\partial}(M)
\longrightarrow
H^{p,q}_{\bar\partial}(M)
$$
is an isomorphism.
\end{theorem}

\begin{proof}
The identities
$$
\bar\partial^2=0,
\qquad
(\bar\partial^*)^2=0
$$
give
$$
[\bar\partial,
\Delta_{\bar\partial}]=0.
$$
Therefore $\bar\partial$ preserves the
eigenspaces of $\Delta_{\bar\partial}$.

Let $\lambda\neq0$, and let $\eta$ be a
$\bar\partial$-closed eigenform with
eigenvalue $\lambda$. Then
$$
\begin{aligned}
\lambda\eta
&=\bar\partial\bar\partial^*\eta
+\bar\partial^*\bar\partial\eta\\
&=\bar\partial\bar\partial^*\eta.
\end{aligned}
$$
Consequently
$$
\eta
=\bar\partial
(\lambda^{-1}\bar\partial^*\eta).
$$
The nonzero eigenspaces contribute no
Dolbeault cohomology. The zero eigenspace is
$\mathcal{H}^{p,q}_{\bar\partial}(M)$.
\end{proof}

% --- End hodge-theory.tex ---

% --- Begin variations-of-hodge-structures.tex ---
\chapter{Variations of Hodge Structures}
\label{chapter-variations-of-hodge-structures}


\section{Introduction}
\label{variations-of-hodge-structures-section-introduction}

These notes follow Misha Verbitsky's
2024 course at IMPA.
\href{http://verbit.ru/IMPA/VHS-2024/}
{Course materials} include ten lecture
slide sets, three assignments and an exam.
References below use lecture and slide
numbers. The presentation is a concise
account of the course, supplemented by
period integrals and Picard--Fuchs theory.

A family of varieties gives a local system
of cohomology groups. Its flat connection
records topological transport, while its
Hodge filtration varies holomorphically.
Period maps describe that variation;
monodromy records transport around loops.
The course applies this structure to
Noether--Lefschetz theory.

\section{Hodge structures and polarization}
\label{variations-of-hodge-structures-section-hodge-structures}

Lecture 1, slides 2--5, 7--11, 17--20.
A real Hodge structure of weight $k$ is a
finite-dimensional real vector space with
$$
V_{\mathbb C}=\bigoplus_{p+q=k}V^{p,q},
\qquad
\overline{V^{p,q}}=V^{q,p}.
$$
A rational or integral structure specifies
compatible rational coefficients or a
lattice. The circle acts on $V^{p,q}$
with weight $p-q$; tensor products add
bidegrees.

For a compact K\"ahler manifold, harmonic
representatives give the Hodge decomposition
of cohomology. A polarization is a
nondegenerate rational bilinear pairing,
with parity determined by the weight,
compatible with the decomposition and
satisfying Hodge--Riemann positivity.
For a smooth projective variety, an ample
class and the Lefschetz decomposition
supply polarizations on primitive
cohomology. Orthogonal complements then
show that polarizable rational Hodge
structures form a semisimple category.

Weight-one structures connect the theory
to complex tori. Polarizable integral
structures of type $(1,0)+(0,1)$ correspond
to abelian varieties; passing to rational
coefficients corresponds to working up to
isogeny. Using $H^1$ makes this
correspondence contravariant.

\section{Period domains and period maps}
\label{variations-of-hodge-structures-section-period-maps}

Lecture 2, slides 3--5, 8--18, 25--26.
Use the decreasing filtration convention
$$
F^p=\bigoplus_{a\geq p}V^{a,k-a}.
$$
Fixing Hodge numbers gives a flag variety.
The conjugacy and positivity conditions,
and the equations imposed by a fixed
polarization, define the period domain.
A marking identifies fibre cohomology with
one fixed vector space. The period map
then sends a variety to its Hodge
filtration in this space.

Torelli questions ask how much of the
variety this Hodge data determines.
The lecture discusses tori, curves, K3
surfaces and hyperk\"ahler manifolds.
Markings, polarizations and the relevant
quotients must be specified in each case.
For polarized Hodge structures of K3 type,
the line $V^{2,0}$ determines the filtration.
The period domain consists of lines
satisfying
$$
Q(\omega,\omega)=0,
\qquad Q(\omega,\overline\omega)>0.
$$
When $Q$ has signature $(2,n)$, this is
also the space of oriented positive
real two-planes.

\section{Local systems and Gauss--Manin}
\label{variations-of-hodge-structures-section-gauss-manin}

Lecture 3, slides 2--10.
A local system is a locally constant
sheaf of finite-dimensional vector spaces.
It determines a vector bundle with a flat
connection; conversely, parallel sections
of a flat connection form a local system.
On a connected base, parallel transport
around loops gives its monodromy
representation.

Let $\pi:X\longrightarrow S$ be a smooth
proper submersion. Ehresmann's theorem
identifies nearby fibres differentiably.
Their cohomology groups form the local
system $R^k\pi_*\mathbb C$. Its associated
flat bundle carries the Gauss--Manin
connection.

One construction differentiates a closed
relative form along a lift of a vector
field on the base, then takes its fibrewise
cohomology class. Changing the lift adds
a vertical Lie derivative, which is exact
on closed forms by Cartan's formula.
Thus the induced connection is canonical,
although the chosen lift is not.
Flatness means that transport depends
only on the homotopy class of a path;
it does not require trivial monodromy.

\section{Variations and transversality}
\label{variations-of-hodge-structures-section-variations}

Lecture 3, slides 11--19.
A variation of Hodge structures consists
of a local system with real, rational or
integral structure, and a holomorphic
Hodge filtration on its associated bundle.
Each fibre is a Hodge structure, and the
flat connection satisfies Griffiths
transversality:
$$
\nabla F^p\subset F^{p-1}\otimes\Omega_S^1.
$$
A polarization is parallel and polarizes
each fibre. Smooth projective families
provide the basic examples, using relative
ample classes and primitive cohomology.

The filtration varies holomorphically but
need not be parallel. Transversality says
that differentiation lowers its index by
at most one. On a simply connected base,
flat transport identifies all fibres and
the filtration defines a holomorphic
period map. Otherwise this description
holds on the universal cover, with
monodromy equivariance.

For a complementary treatment, see
\ref{mumford-tate-groups-in-%
hodge-theory-section-variations-of-hs}
in {\it Mumford--Tate groups in Hodge
 theory}.

\section{Fixed parts and semisimplicity}
\label{variations-of-hodge-structures-section-fixed-part}

Lecture 4, slides 4--17;
Lecture 5, slides 6--15.
For a polarized variation over a compact
connected complex base, the fixed part
theorem says that the Hodge components of
a global parallel section are parallel.
Thus monodromy-invariant classes inherit
a Hodge structure. In particular, trivial
monodromy forces the variation to be
constant under these hypotheses.

The proof uses the Chern connection on
the associated graded Hodge bundle,
the components of Gauss--Manin that change
Hodge type, and curvature estimates.
Plurisubharmonic norms become constant on
a compact base. Applied to a Hom variation,
the theorem gives rigidity: a flat map
preserving Hodge structures at one point
preserves them everywhere.

Deligne's semisimplicity theorem, in the
integral polarized setting of Lecture 5,
gives a decomposition over $\mathbb C$:
$$
V=\bigoplus_i L_i\otimes_{\mathbb C}W_i.
$$
Here $L_i$ have irreducible monodromy and
carry complex variations, while $W_i$ are
constant complex Hodge structures.
The decomposition respects Hodge types;
it need not respect the integral lattice.
Lecture 5 also states the extension to
smooth quasiprojective bases, using the
corresponding fixed part theorem.

\section{Hodge groups and K3 type}
\label{variations-of-hodge-structures-section-hodge-groups}

Lecture 6, slides 2--8, 12--15.
The smallest rational algebraic subgroup
containing the Hodge circle is the Hodge
group. The slides call this the
Mumford--Tate group; the convention using
the full Deligne torus also records weight.
Keeping this distinction explicit avoids
confusing the two group conventions.

Tensor invariants recover the Hodge group:
it fixes the rational tensors invariant
under the circle. Hodge substructures
correspond to invariant subspaces.
See \ref{mumford-tate-groups-in-%
hodge-theory-section-mumford-tate-group}
for the group conventions and
\ref{mumford-tate-groups-in-%
hodge-theory-section-tensor-construction}
for tensor constructions.

For a simple polarizable rational Hodge
structure of K3 type, its Hodge
endomorphism algebra is a number field.
Zarhin's theorem makes this field either
totally real or CM; the corresponding
Hodge group is described by orthogonal
or unitary symmetries. The course applies
this to the transcendental Hodge structure,
the smallest rational Hodge substructure
containing the holomorphic top forms.

\section{Noether--Lefschetz loci}
\label{variations-of-hodge-structures-section-noether-lefschetz}

Lecture 3, slide 20;
Lecture 7, slides 6--12;
Lecture 10, slides 4--6, 19.
Fix a rational cohomology class by flat
transport on a simply connected open set.
The condition that it remain of type
$(p,p)$ defines a Hodge locus. In degree
two, these loci detect additional divisor
classes and jumps of the Picard rank.

The Noether--Lefschetz theorem says that a
very general smooth surface of degree
$d\geq4$ in $\mathbb P^3$ has Picard rank
one. Here very general means outside a
countable union of proper algebraic loci.
The hyperplane class supplies that one
class. Irreducibility of monodromy on
primitive cohomology excludes a persistent
proper subsystem of additional divisor
classes. The alternative that all
primitive cohomology has type $(1,1)$ is
excluded by the existence of holomorphic
two-forms for $d\geq4$.

\section{Lefschetz pencils and vanishing
cycles}
\label{variations-of-hodge-structures-section-vanishing-cycles}

Lecture 8, slides 2--11;
Lecture 9, slides 2--18;
Lecture 10, slides 11--19.
A pencil is a one-parameter family of
divisors from a two-dimensional space of
sections. A Lefschetz pencil has smooth
base locus and only ordinary double
points in singular members. Blowing up
the base locus resolves the rational map
to $\mathbb P^1$.

Near a nondegenerate critical point,
the local model is
$$
z_1^2+\cdots+z_{n+1}^2=t.
$$
An $n$-sphere in a smooth fibre contracts
as $t$ approaches zero. Its class is a
vanishing cycle; sweeping it along a path
to zero gives a Lefschetz thimble.
Picard--Lefschetz theory expresses local
monodromy as a rank-one modification
controlled by intersection with this
cycle. The sign depends on dimension
and intersection conventions.

The vanishing cycles generate vanishing
cohomology. For the universal family of
smooth hyperplane sections, monodromy on
this part is irreducible. Applied to the
Veronese embedding, this supplies the
irreducibility used in the proof of the
Noether--Lefschetz theorem.

Picard--Lefschetz theory describes transport
around singular fibres. Picard--Fuchs
equations describe the differential
variation of periods over regular values.
Both concern the same Gauss--Manin local
system from complementary viewpoints.

\section{Periods, Picard--Fuchs equations
and quantum corrections}
\label{variations-of-hodge-structures-section-quantum-periods-background}

\noindent
{\bf Periods and differentiation.}
Over regular values of $f$, follow a
locally constant homology class
$\gamma_t$ in the fibres $C_t=f^{-1}(t)$.
For a relative cohomology class $\omega$,
$$
\Pi(t)=\int_{\gamma_t}\omega_t,
\qquad
\Pi'(t)=\int_{\gamma_t}\nabla_{\partial_t}
\omega.
$$
Gauss--Manin is the connection that
makes this differentiation possible.
Reference: Movasati, Proposition 13.1,
equation (13.6), cited below.
For the local system and Hodge filtration,
see Sections \ref{variations-of-hodge-structures-section-gauss-manin}
and \ref{variations-of-hodge-structures-section-variations}.

\medskip\noindent
{\bf Picard--Fuchs equations.}
In a rank-$r$ cohomology bundle, the
classes $\omega,\nabla\omega,\ldots,
\nabla^r\omega$ are dependent over the
field of meromorphic functions.
Integrating such a relation gives
$$
\sum_{j=0}^{m}b_j(t)\Pi^{(j)}(t)=0,
\qquad m\leq r.
$$
Reference: Movasati, Section 13.5,
equations (13.10)--(13.12).
For compact genus-$g$ curves, $r=2g$.
For affine curves, punctures can add
homology classes, so genus alone does
not give the full rank.
An elliptic curve supplies the familiar
rank-two compact case. The curve is a
fibre of the family; it need not itself
be a moduli space. Periods of a nonzero
holomorphic one-form along a marked
basis of cycles determine its complex
structure through their ratio. This is
the weight-one period map, rather than
an identification of the fibre with the
space parametrizing the family.

For an action integral of a one-form
$\lambda$ on a surface, write locally
$$
d\lambda=df\wedge\eta.
$$
Then differentiation of
$A(t)=\int_{\gamma_t}\lambda$ gives
$$
A'(t)=\int_{\gamma_t}\eta.
$$
This is the Gelfand--Leray derivative.
The rank bound applies directly to the
period of $\eta$. An equation for $A'$
can give an equation for $A$ with one
additional derivative; identifying the
action with a cohomological period
requires the appropriate setup.

\medskip\noindent
{\bf Quantization.}
The relation $[q,p]=h$ replaces commuting
coordinates by noncommuting generators.
The quantum-period construction discussed
in Garay's seminar introduces formal
corrections to the classical area.
For the passage from classical observables
to operators,
see the Dirac rule and Kostant--Souriau
formula in
\ref{geometric-prequantization-%
section-dirac-rule} and
\ref{geometric-prequantization-%
section-kostant-souriau-formula}
in {\it Geometric prequantization}.
For quantum tori and the distinction
between deformation quantization and
quantum cohomology, see
\ref{homological-mirror-symmetry-%
section-where-this-sits} and
\ref{homological-mirror-symmetry-%
section-tori-and-mumford}
in {\it Homological mirror symmetry}.
For the reported quantum-period
construction and its unresolved details,
see \ref{seminars-algebraic-geometry-%
section-garay-quantum-periods}.
These quantum corrections are additional
structure, not part of the definition of
a classical variation of Hodge structures.

\medskip\noindent
{\bf Connections and gauge theory.}
Seiberg--Witten theory studies equations
for connections and spinors modulo gauge
transformations. See
{\it Seiberg--Witten invariants}, Sections
\ref{seiberg-witten-section-dirac} and
\ref{seiberg-witten-section-invariants}.
There the connections are part of the
unknown geometric data. In period theory,
the Gauss--Manin connection describes
variation of fibre cohomology.
The cited sections concern invariants of
four-manifolds; they do not treat
Seiberg--Witten curves and their periods.

\medskip\noindent
{\bf Published reference.}
H. Movasati, {\it A Course in Hodge
Theory: With Emphasis on Multiple
Integrals}, International Press, 2021.
Definition 10.16 (Gelfand--Leray form);
Proposition 13.1 (differentiating periods);
Section 13.5 (Picard--Fuchs equations).
\href{https://w3.impa.br/~hossein/%
myarticles/hodgetheory.pdf}
{Author's text}, consulted
October 14, 2026.

% --- End variations-of-hodge-structures.tex ---

% --- Begin k3.tex ---
\chapter{K3 Surfaces}
\label{chapter-k3}




\section{K3 surfaces}
\label{k3-section-k3-surfaces}

\begin{reference}
\cite[Definition 1.1]{huc-k3}
\end{reference}

\begin{definition}[Huybrechts]
\label{k3-definition-k3-huybrechts}
\textit{A K3 surface over $k$} is a complete
non-singular variety $X$ of dimension two
such that
$$
\Omega^2_{X/k}\simeq\mathcal{O}_X.
$$
and
$$
H^1(X,\mathcal{O}_X)=0.
$$
\end{definition}

\noindent
By definition, a variety over a field $k$ is
complete if the morphism
$X\to\operatorname{Spec}(k)$ is proper.
The variety $X$ over $k$ is non-singular if
the cotangent sheaf $\Omega_{X/k}$ is locally
free of rank $\dim(X)$. This is equivalent
to
$X_{\overline{k}}=X\times_k\overline{k}$
being regular; see, for example, [375,
Proposition 6.2.2].

\begin{definition}[Misha]
\label{k3-definition-k3-misha}
\textit{A K3 surface} is a compact complex
surface $M$
with
$b_1(M)=0$ and
$c_1(M,\mathbb{Z})=0$.
\end{definition}

\begin{remark}
This implies that the canonical
bundle is trivial.
\end{remark}

\begin{proposition}
The Hodge diamond of a K3 surface
is
\[
\begin{matrix}
&&1&&\\
&0&&0&\\
1&&20&&1\\
&0&&0&\\
&&1&&
\end{matrix}
\]
so in particular
$H^{2,0}(M)$ is one-dimensional.
\end{proposition}

\begin{remark}
The middle cohomology is the main
cohomological invariant of a K3
surface.
The intersection form lives on
$H^2(M,\mathbb{Z})$.
\end{remark}

\medskip\noindent
In this exercise we shall find
an example of a non-algebraic K3 surface.
Recall that the algebraic dimension $a(X)$
of a variety $X$ is the transcendence
degree of its field of meromorphic
functions over $\mathbb{C}$.

\begin{exercise}
\label{k3-exercise-k3-algebraic-dimension}
Prove that there exists a K3 surface
of algebraic dimension 1.
\end{exercise}

\begin{proof}
The main idea is that a surface
is algebraic if and only if $a(X) = 2$.
We need to construct a surjection
$f:X \to \mathbb{P}^1$
for some nonprojective $X$.

Pulling back
$\mathcal{M}(\mathbb{P}^1)=\mathbb{C}(t)$
to $X$ we obtain an injection
$f^*\mathbb{C}(t)\hookrightarrow\mathcal{M}(X)$,
and so $a(X) \geq 1$.

Showing that $X$ is nonprojective
thus gives $a(X) = 1$.

My idea so far is to use a Kummer surface
$K$. For starters, since $K$ is defined
as the normalization of an abelian
surface $A$ quotiented by an involution
$\iota$, we can deal only with $A/\iota$
since normalization is a dominant map.

Then we argue that quotient by $\iota$
also makes no difference in
$\mathcal{M}(A)$.

But I'm not sure how to create the map
$K \to \mathbb{P}^1$, or
$A \to \mathbb{P}^1$.

Also not sure how to prove $K$ is not
projective.

\medskip\noindent
{\bf Day 2.} Today I did the following:
let $E = \mathbb{C}/\mathbb{Z}^2$.
Regular functions on $E$ are double
periodic functions, i.e. such that
$f(z) = f(z+n e_1+ m e_2)
=f(z+n+im)$ for all $n,m \in \mathbb{Z}$.
They are all constant by Liouville.
But there can be doubly periodic
{\it meromorphic} functions: having poles
makes them be unbounded.

In fact, there exist nontrivial meromorphic
functions such as the Weierstrass
something function.
Being nontrivial, it has positive
transcendence degree, for no nonconstant
function can be constructed from 
complex numbers.

Now consider $E \times E \to E$.
I aim to show that $E \times E$ is not
projective. By Huybrechts HK paper,
it's enough to show that the intersection
form is negative semidefinite, i.e. that
$L^2 \leq 0$ for all $L \in
\operatorname{Pic}_{E \times E}$,
which is a criterion that works for
hyperkähler manifolds such as the K3
surface $E \times E$.

I still need to use (and prove)
the criterion that
projective iff $a(X)=2$.
But $a(X)$ cannot be greater than 2
if $X$ is a surface?
So: not projective, not 2; surface, not
$\geq 2$. Elliptic: $a(X) \geq 1$.

\medskip\noindent
{\bf Day 3.} Maybe I should just
try to prove that $a(E\times E)<2$
somehow…

\medskip\noindent
{\bf Day 4.} The idea is: $E\times E$
is not projective because
Pic of a product is direct sum of Pics,
and Pic of a torus $E$ is trivial because
$E$ is a curve. 
Wait: I think
$$
H^0(X,\mathcal{O}_X) = 0
\to
H^0(X,\mathcal{O}_X^*) = 0
\to
H^1(X,\mathbb{Z}) = \mathbb{Z} \oplus \mathbb{Z}
\to
H^1(X,\mathcal{O}_X)
\to
H^1(X,\mathcal{O}_X^*)=\operatorname{Pic}(X)
\to
H^2(X,\mathbb{Z}) = \mathbb{Z}
\to
H^2(X,\mathcal{O}_X)=0
$$


Thus $E \times E$
satisfies the ampleness criterion
of Huybrechts, and it is not projective.

Then I need to prove that this implies
$a(E \times E) \neq 2$.
\end{proof}

\section{Curves on K3 surfaces}
\label{k3-section-curves-on-k3}

\noindent
My progress on this
exercise still has three black boxes:
\begin{enumerate}
\item Construction of normalization
from gluing affine pieces.
\item Finite scheme morphism preserves Euler
characteristic.
 
I know that the normalization map 
$\nu:\widetilde C\to C$
is finite by Noether normalization
(and gluing, which is my current work).
The roadmap is:
Finite morphism $\implies$
affine morphism $\implies$
preserves cohomology, i.e.
for any sheaf $\mathcal{F}$ on
$\widetilde C$ we have
$$
H^i(\widetilde C, \mathcal{F})
\cong H^i(C,\nu_*\mathcal{F}).
$$
That uses Leray Spectral Sequence,
apparently.
But it obviously gives
$\chi(\mathcal{O}_{\widetilde{C}})
=\chi (\nu_*\mathcal{O}_{\widetilde{C}})$
which is what we really need.
\item If the local ring is regular
i.e. $\mathfrak{m}/\mathfrak{m}^2
=\dim \mathcal{O}_p$
then the variety is smooth
in the sense that it has 
holomorphic coordinates near $p$.
If I presented this on Friday I would have to
prove that.
\end{enumerate}

\begin{exercise}
\label{k3-exercise-singular-curve-on-k3}
Let $C$ be a singular, irreducible complex
curve on a K3 surface $X$. Prove that
$(C.C)\geq 0$.
\end{exercise}

\begin{proof}
By the Genus Formula
(Lemma
\ref{complex-geometry-lemma-genus-formula})
we have
$$
2p_a(C)-2 = C \cdot (C + K_X).
$$
Since $X$ is a K3 surface, the canonical
divisor is linearly equivalent to zero.
Hence
$$
C^2 = 2p_a(C)-2.
$$
Thus it suffices to prove that
$p_a(C) \geq 1$.

\medskip\noindent
Let
$$
\nu : \widetilde C \longrightarrow C
$$
be the normalization of $C$. We recall its
construction. Choose an affine open covering
$U_i = \Spec(A_i)$ of $C$. Since $C$ is
irreducible, all the rings $A_i$ have the
same function field $K(C)$. Let
$\overline A_i$ be the integral closure of
$A_i$ in $K(C)$. Then the affine schemes
$$
\Spec(\overline A_i)
$$
glue to a normal curve $\widetilde C$,
and the inclusions
$$
A_i \subset \overline A_i
$$
glue to a finite birational morphism
$\nu : \widetilde C \to C$
(finiteness follows
from Noether normalization lemma,
i.e. Commutative Algebra
Lemma \ref{commutative-algebra-lemma-noether-normalization},
and birationality is basically by
construction of the normalization
as the spectrum of the integral
closure, i.e. the fraction fields
are the same and the induced
map on the original ring is the identity;
check the definition of birational
morphism, Commutative Algebra Definition
\ref{commutative-algebra-definition-birational-morphism}.).

We have an exact sequence of sheaves
$$
0 \to \mathcal O_C
  \to \nu_*\mathcal O_{\widetilde C}
  \to \delta \to 0,
$$
where
$$
\delta =
\nu_*\mathcal O_{\widetilde C}/\mathcal O_C.
$$
Taking Euler characteristics gives
$$
\chi(\mathcal O_C)
=
\chi(\nu_*\mathcal O_{\widetilde C})
-
\chi(\delta).
$$
Since $\nu$ is finite (this is
a black box for now)
$$
\chi(\nu_*\mathcal O_{\widetilde C})
=
\chi(\mathcal O_{\widetilde C}).
$$

Now observe that, at smooth
points $p \in C$, the stalks
$\mathcal{O}_{C,p}$ are
UFDs by Complex Analysis
Lemma \ref{complex-analysis-lemma-stalk-is-ufd}
--- due to Weierstrass Preparation
Theorem, in turn due to the Gauss lemma.
By Commutative Algebra Lemma
\ref{commutative-algebra-proposition-ufd-integrally-closed},
UFDs are integrally closed,
which means that the stalks
of $\tilde{C}$ and $C$ are isomorphic.
Thus, the stalks of $\delta$
can only be non trivial at non smooth
points of $C$. There can only
be finitely many non smooth points
in a projective variety,
and therefore $\delta$
is a skyscraper sheaf with finite support.
Thus
$$
\chi(\delta)=h^0(\delta).
$$
Therefore
$$
\begin{aligned}
p_a(C)
&= 1-\chi(\mathcal O_C) \\
&= 1-\chi(\mathcal O_{\widetilde C})
   + h^0(\delta) \\
&= g(\widetilde C)+h^0(\delta),
\end{aligned}
$$
where in the last equality
we we can replace
$1-\chi(\mathcal{O}_{\widetilde{C}})$
with the geometric genus
$g(\widetilde{C}):=
h^0(\omega_{\widetilde{C}})$
by applying Serre duality,
thereby using that $\widetilde{C}$
is smooth since it is normal.
To prove that normal implies
smooth we use
Commutative Algebra
Proposition
\ref{commutative-algebra-proposition-integrally-closed-dvr}
to get from integrally closed
that the local ring is regular.
Then we need Hartshorne
to get differential-geometric smoothness
from regular local ring.

Now $g(\widetilde C)\geq 0$. Since $C$ is
singular, the normalization is not an
isomorphism, so $\delta\neq 0$. Since
$\delta$ is a nonzero skyscraper sheaf, it
has a nonzero global section. Hence
$$
h^0(\delta)\geq 1.
$$
Thus
$$
p_a(C)=g(\widetilde C)+h^0(\delta)\geq 1.
$$
Consequently
$$
C^2=2p_a(C)-2\geq 0.
$$
\end{proof}



\section{Period domain}
\label{k3-section-period-domain}

\begin{definition}
\label{k3-definition-k3-lattice}
The {\it K3 lattice} is the even
unimodular lattice
$$
\Lambda_{K3}
=
U^{\oplus 3}\oplus E_8(-1)^{\oplus 2}
$$
of signature $(3,19)$, where $U$ is the
hyperbolic plane with intersection matrix
$$
\begin{pmatrix}
0&1\\
1&0
\end{pmatrix},
$$
and $E_8(-1)$ is the negative-definite
$E_8$ lattice.
For every K3 surface $X$, the lattice
$H^2(X,\mathbb{Z})$ equipped with its
intersection pairing is isomorphic to
$\Lambda_{K3}$.
\end{definition}

\begin{definition}
\label{k3-definition-period-k3-surface}
Let $X$ be a K3 surface.
The {\it period} of $X$ is the
weight-two integral Hodge structure
$$
H^2(X,\mathbb{C})
=
H^{2,0}(X)\oplus H^{1,1}(X)
\oplus H^{0,2}(X)
$$
on $H^2(X,\mathbb{Z})$, together with
the intersection pairing
$$
(\alpha,\beta)=\int_X\alpha\wedge\beta.
$$
Equivalently, after choosing a {\it
marking}, i.e. a lattice isomorphism
$\phi:H^2(X,\mathbb{Z})\to\Lambda_{K3}$,
the period is the line
$$
\phi_{\mathbb{C}}(H^{2,0}(X))
\subset\Lambda_{K3}\otimes\mathbb{C}.
$$
\end{definition}

\begin{definition}
\label{k3-definition-period-domain-k3-surfaces}
Let $\Lambda_{K3}$ be the K3 lattice.
The {\it period domain} of marked K3
surfaces is
$$
\Omega_{\Lambda_{K3}}
=
\left\{
[\sigma]\in
\mathbb{P}(\Lambda_{K3}\otimes\mathbb{C})
\ \middle|\
(\sigma,\sigma)=0,\quad
(\sigma,\overline{\sigma})>0
\right\}.
$$
Thus $\Omega_{\Lambda_{K3}}$ parametrizes
the weight-two Hodge structures on
$\Lambda_{K3}$ which are compatible with
its intersection pairing and have
one-dimensional $(2,0)$-part.
\end{definition}

\begin{definition}[Misha]
Let $V=H^2(M,\mathbb{R})$ with the
intersection form.
The \textit{period domain} is
\[
\mathbb{P}\mathrm{er}=
\{[v]\in \mathbb{P}H^2(M,\mathbb{C})
\mid (v,v)=0,\ (v,\bar{v})>0\}.
\]
\end{definition}

\begin{remark}
This is the open subset of the
quadric cut out by $(v,v)=0$.
The positivity condition chooses
the correct connected component.
\end{remark}

\begin{proposition}[Misha]
The period point of a K3 surface
is the line $H^{2,0}(M)$.
\end{proposition}

\begin{remark}
So the period map sends a complex
structure to its Hodge line in
$H^2(M,\mathbb{C})$.
The old notes identify this
space with a Grassmannian of
positive oriented planes.
\end{remark}

\section{Local Torelli}
\label{k3-section-local-torelli}

\begin{definition}[Misha]
\label{k3-definition-period-map-misha}
Let $\mathrm{Comp}$ be the space of complex
structures $I$ on $M$ for which $(M,I)$ is
a K3 surface, and let $\mathrm{Diff}_0(M)$
be the connected component of the identity
in the diffeomorphism group of $M$.
Set
$$
\mathrm{Teich}
=
\mathrm{Comp}/\mathrm{Diff}_0(M).
$$
The \textit{period map} is
$$
\begin{aligned}
\mathrm{Per}:\mathrm{Teich}
&\longrightarrow \mathbb{P}\mathrm{er},\\
I&\longmapsto H^{2,0}(M,I).
\end{aligned}
$$
\end{definition}

\begin{proposition}
The period map for K3 surfaces is
a local diffeomorphism.
\end{proposition}

\begin{remark}
This is the deformation-theoretic
statement that connects complex
structure with the Hodge subspace
$H^{2,0}(M)$.
\end{remark}

\begin{proposition}
On a K3 surface, a holomorphic
symplectic form determines the
complex structure locally.
\end{proposition}

\section{Lefschetz}
\label{k3-section-lefschetz}

\begin{remark}
Morse theory enters the proof of
Lefschetz hyperplane section via
a strictly plurisubharmonic
function on the affine complement.
\end{remark}

\begin{proposition}
If $Z\subset \mathbb{C}P^n$ is
smooth and $H$ is a hyperplane,
then the topology of
$Z\setminus H$ is controlled by
the critical points of an
appropriate Morse function.
\end{proposition}

\begin{remark}
The outline is:
\begin{enumerate}
\item Work on the affine
complement $\mathbb{C}P^n\setminus H$.
Take the standard potential
$f_1=\sum |z_i|^2$ on
$\mathbb{C}^n$.
\item Its complex Hessian
$dd^cf_1$ is positive.
So $f_1$ is strictly
plurisubharmonic.
\item Perturb $f_1$ slightly.
Strict plurisubharmonicity is
open in the $C^2$ topology and
Morse functions are dense.
Hence choose a function that is
both Morse and strictly
plurisubharmonic on
$Z\setminus (Z\cap H)$.
\item Use the gradient flow of
that function.
Each critical point gives a
stable manifold.
The stable manifolds are the
pieces attached to the lower
level set.
\item The stable manifolds have
real dimension equal to the
Morse index, and for a
strictly plurisubharmonic
function the index is at most
$\dim_{\mathbb{C}} Z$.
\item Therefore the open
complement is built from the
hyperplane section by attaching
handles of controlled index.
This gives the desired
connectivity and cohomological
vanishing range.
\end{enumerate}
\end{remark}

\begin{remark}
This is the point of the old
notes: one identifies
the hyperplane section as a
deformation retract of the
complement of the stable
manifolds.
That turns the Lefschetz
statement into a Morse
computation.
\end{remark}

\begin{thebibliography}{9}
\bibitem{huyk3}
D. Huybrechts,
\textit{Lectures on K3 Surfaces}.
\end{thebibliography}

% --- End k3.tex ---

% --- Begin mumford-tate-groups-in-hodge-theory.tex ---
\chapter{Mumford-Tate groups in Hodge theory}
\label{chapter-mumford-tate-groups-in-hodge-theory}



\noindent
Minicourse by Andrey Soldatenkov, UNICAMP,
Brazil,
IMPA Summer School, 2026.




\noindent
Upshot. The Mumford-Tate group is
the stabilizer of the Hodge classes,
and is constant on the complement of the
Hodge locus
(see Definition
\ref{mumford-tate-groups-in-hodge-theory-definition-generic-mt-group}).

\section{Plan}
\label{mumford-tate-groups-in-hodge-theory-section-plan}

\begin{enumerate}
\item Motivation: cohomology of algebraic
varieties.
\item Definition. Hodge structures,
Mumford-Tate group.
\item Characterizations of the MT groups and
relations with representation
theory.
\item Variations of Hodge structures and
moduli spaces.
\item Dichotomy: abelian vs non-abelian HS.
\item The Kuga-Satake construction.
\end{enumerate}

\section{Introduction}
\label{mumford-tate-groups-in-hodge-theory-section-intro}

$X\subset \mathbb{C}P^N$ smooth complex
subvariety,
$\dim_\mathbb{C}=n$.
First recall we have singular cohomology,
$H^k(X,\mathbb{C})$,
which is isomorphic to the cohomology
of the constant sheaf
$\underline{\mathbb{C}}_X$.
This cohomology is nonzero for $0\leq k \leq
2n$.

Recall. $U \subset X$ open,
$\Gamma(U,\mathbb{C})=\{f:U \to \mathbb{C}:
\text{$f$ is locally constant}\}
=\prod_{\pi_0(U)}\mathbb{C}$.

\begin{example}
\label{mumford-tate-groups-in-hodge-theory-example-singular-cohomology}
\begin{enumerate}
\item
$X=\mathbb{C}P^n$,
$$
H^k(\mathbb{C}P^n,\mathbb{C})=
\begin{cases}
  \mathbb{C}\qquad &k=2m, 0 \leq  m \leq  n
  \\
  0\qquad &\text{otherwise.}
\end{cases}
$$
A way to prove this is using the CW
decomposition of $\mathbb{C}P^n$.

\item $X \subset \mathbb{C}P^2$ hypersurface
of degree $d$;
$X$ a Riemann surface of genus
$g=\frac{(d-1)(d-2)}{2}$.
Then
$H^0(X,\mathbb{C})=\mathbb{C}$,
$H^1(X,\mathbb{C})=\mathbb{C}^{2g}$,
$H^2(X,\mathbb{C})=\mathbb{C}$.
\end{enumerate}
\end{example}

\noindent
We also have the following additional data
(a Hodge structure) on $H^k(X,\mathbb{C})$:
\begin{itemize}
\item A lattice
$H^k(X,\mathbb{Z})/\text{torsion}\subset
H^k(X,\mathbb{C})$,
\item A $(p,q)$-decomposition,
$H^k(X,\mathbb{C})=\bigoplus_{p+q=k}H^{p,q}
(X)$.
\end{itemize}

Why is it useful?
\begin{itemize}
\item It gives restrictions on possible Betti
numbers of algebraic varieties.
(So it may tell us that certain complex
variety cannot be algebraic,
for example.)

\item It $f:X \to Y$ is a morphism of
algebraic varieties,
then $f^*:H^k(Y,\mathbb{C}) \to
H^k(X,\mathbb{C})$
preserves the Hodge structure.
\end{itemize}

\noindent
As an example of the latter statement,

\begin{example}
\label{mumford-tate-groups-in-hodge-theory-example-preserving-hodge-structure}
Let $X \subset \mathbb{C}P^4$ be a
``general'' hypersurface
of degree 5. Then there exists no abelian
variety
(i.e. a projective variety that is
biholomorphic
to $\mathbb{C}P^N/\Lambda$ where $\Lambda$ is
a lattice;
so, a complex torus that is also a projective
variety)
$A$ that admits a dominant birational map
onto $X$
$f:A\xymatrix{\ar@{-->}[r]&}X$.
\end{example}

\section{The p,q decomposition}
\label{mumford-tate-groups-in-hodge-theory-section-p-q-decomposition}

\noindent
We use the de Rham complex.
Let $\Omega_X^k$ be the sheaf of holomorphic
$k$-forms.
The de Rham complex is
$$
\Omega_{dR}^\bullet
=(0 \to \mathcal{O}_X \to
\Omega_X^1\xrightarrow{d}\cdots
\to\Omega_X^n \to 0)
$$

$\Omega_{dR}^\bullet$ is a resolution of
$\mathbb{C}_X$.
Therefore $H^k(X,\mathbb{C})\cong
H^k(X,\Omega_{dR}^\bullet$.

Let's define a subcomplex:
$$
F^p\Omega_{dR}^\bullet=
(0 \to \cdots \to \Omega_X^p \to O_X^{p+1}
\to \cdots \to \Omega_X^n \to 0)
$$
It's a subcomplex $F^p\Omega_{dr}^\bullet
\subset \Omega_{dR}^\bullet$
since the sheaves coincide when
$F^p\Omega_{dR}$ are nonzero
and they inject otherwise (because it's the
zero sheaf).

This gives
$H^k(X,F^p\Omega_{dR}^\bullet)
\xrightarrow{(\ast)}
H^k(X,\Omega_{dR}^\bullet)=H^k(X,\mathbb{C})
$.

\begin{definition}
\label{mumford-tate-groups-in-hodge-theory-definition-hodge-filtration}
The {\it Hodge filtration} is
$F^pH^k(X)=\text{Im}(\ast)$.
\end{definition}

\noindent
Hodge theory tells us the map $(\ast)$ is in
fact injective.

Let $\Lambda^{p,q}_X$ be the sheaf of
$C^\infty$-forms
on $X$ of type $(p,q)$.
They are given locally by
$\sum_{\substack{|I|=p \\ |J|=q}}
\alpha_{I J}dz_I \wedge d\overline{z}_J$,
for $\alpha_{I J}\in C^\infty_X$.

Then we have an acyclic resolution of
$\Omega_X^p$:
$$
0 \to \Omega_X^p \hookrightarrow
\Lambda^{p,0}_X
\xrightarrow{\bar{\partial}}
\Lambda^{p,1}_X
\xrightarrow{\overline{\partial}}
\cdots
\to
\Lambda^{p,n}_X\to 0.
$$

$\Lambda^{\bullet,\bullet}_X$ is a fine
resolution of $\Omega_X^\bullet$.

Then $H^k(X,\mathbb{C})\cong
H^k(X,\text{Tot}\Lambda^{\bullet,\bullet}_X$,
which can be computed by a spectral sequence.
The first page of such spectral sequence is
given by
$E^{p,q}_1=H^q(X,\Omega_X^p)$.
This converges to $H^k(X,\mathbb{C})$.
This is the Hodge-to-de Rham spectral
sequence.

\medskip\noindent
Since our manifolds are projective they admit
a Kähler metric
$\omega$ induced by the inclusion $X
\xrightarrow{i}\mathbb{C}P^N$,
that is, $\omega=i^* (\text{Fubini-Study
metric on }\mathbb{C}P^N)$.

Then
$\Lambda_X^k=\bigoplus_{p+q=k}\Lambda_X^{p,q}
$,
$\Lambda^\bullet_X$ becomes an elliptic
complex.

We have
$$
… \to \Lambda_X^{k-1}\xrightarrow{d}
\Lambda_X^k\to …,
\qquad
…\to\Lambda_X^k
\xrightarrow{d^*}\Lambda_X^{k-1}\to…
$$
where $d^*$ is the adjoint of $d$ w.r.t.
$\omega$.
Then $\Lambda=d d^*+d^*d$ is an elliptic
operator.
$\mathcal{H}^k=\Ker(\Lambda|_{\Lambda^k_X})
=\Ker(d|_{\Lambda^k_X})\cap
\Ker(d^*|_{\Lambda^k_X})$
are the harmonic forms.

Consider a natural map from the harmonic
forms of type $(p,q)$
to the cohomology:
$$
\mathcal{H}^{p,q}=\mathcal{H}
^k\cap\Lambda^{p,q}_X
\xrightarrow{(\ast \ast)} H^q(X,\Omega_X^p).
$$
Fact: since $d\omega=0$ ($X$ is Kähler),
$(\ast\ast)$ is an isomorphism.

This means that for Kähler manifolds
$$
\dim H^k(X,\mathbb{C})
\leq
\sum_{p+q=k}H^q(X,\Omega_X^p)
\underbrace{\leq}_{(\ast \ast)}
\dim H^k(X,\mathbb{C})
$$
This implies that the Hodge-to-de Rham
spectral sequence
degenerates at $E_1$. Therefore,
$$
H^k(X,\mathbb{C})=\mathcal{H}^k=
\bigoplus_{p+q=k}\mathcal{H}^{p,q}\cong
H^q(X,\Omega_X^p).
$$

The Hodge filtration is
$$
F^pH^k(X,\mathbb{C})=\bigoplus_{\substack{p'+
q'=k
\\ p' \geq p}}
H^{p',q'}(X).
$$

\section{Symmetries of the p,q decomposition}
\label{mumford-tate-groups-in-hodge-theory-section-symme-p-q-decom}

\begin{enumerate}
\item Since
$\overline{\Lambda^{p,q}_X}=\Lambda^{q,p}_X$,
we have
$\mathcal{H}^{q,p}=\overline{\mathcal{H}^{p,
q}}$.
This means that if $k \equiv 1 (\text{mod
}2)$,
$H^k(X,\mathbb{C})=\mathcal{H}^{k,0} \oplus …
\oplus
\mathcal{H}^{\frac{k+1}{2},\frac{k-1}{2}}
\oplus \mathcal{H}^{0,k}\oplus … \oplus
\mathcal{H}^{\frac{k-1}{2},\frac{k+1}{2}}$.
This means that the $k$-th Betti number is
even,
$b_k(X)\equiv 0 \text{mod }2$.

\item (Poincare duality.) We have a perfect
pairing
$$
\begin{aligned}
H^k(X,\mathbb{C}) \otimes
H^{2n-k}(X,\mathbb{C})
&\longrightarrow \mathbb{C} \\
[\alpha]\otimes[\beta] &\longmapsto
\int_X\alpha\wedge\beta
\end{aligned}
$$

\noindent
Notice that if
$\alpha \in \mathcal{H}^{p,q}$
then $\beta \in \mathcal{H}^{n-p,n-q}$.
This induces a perfect pairing
$$
\mathcal{H}^{p,q}\otimes\mathcal{H}^{n-p,n-q}
\to \mathbb{C}.
$$

\item (Polarization and the Lefschetz
operator.)
The polarization is the Kähler class of the
Kähler form.
By $X$ being projective we have that the
Kähler class is integral.
Moreover, it is the Poincar'e dual of the
hyperplane section class.
That is, let $h \in
H^2(\mathbb{C}P^n,\mathbb{Z})$
be the class of a hyperplane,
then $i^*h=[\omega] \in H^2(X,\mathbb{Z})$.
We have $\omega \in \mathcal{H}^{1,1}$
and $[\omega]\in H^2(X,\mathbb{Z}) \cap
H^{1,1}(X)$.
The Lefschetz operator is
$$
\begin{aligned}
L_\omega: H^{p,q}(X) &\longrightarrow
H^{p+1,q+1}(X) \\
[\alpha] &\longmapsto
[\alpha\wedge\omega]=[\alpha]\cup [\omega].
\end{aligned}
$$

\noindent
Lefschetz theorem says
\begin{enumerate}
\item
$L_\omega^k:H^{n-k}(X,\mathbb{Q})
\xrightarrow{\sim}
H^{n+k}(X,\mathbb{Q})$
is an isomorphism for $0\leq k\leq n$

\item The dual of $L_\omega$ is
$$
\begin{aligned}
\Lambda_\omega: H^{p,q}(X) &\longrightarrow
H^{p-1,q-1}(X) \\
[\alpha] &\longmapsto [i_\alpha \omega].
\end{aligned}
$$

$$
[L_\omega,\Lambda_\omega]=\Theta \in
\text{End}(H^\bullet(X,\mathbb{Q}))
$$
$$
\Theta|_{H^k(X,\mathbb{Q})}
=(k-n)\text{Id}
$$
Then $L_\omega, \Lambda_\omega$ and $\Theta$
span a subalgebra of
$\text{End}(H^\bullet(X,\mathbb{Q}))$
isomorphic to $\mathfrak{sl}_2$.
This allows us to use what we know about the
representation theory of $\mathfrak{sl}_2$.
Let
$H^k_{\text{prim}}(X,\mathbb{Q})=
(\Lambda|_{H^k(X,\mathbb{Q})})$.
Then $H^m(X,\mathbb{Q})=\bigoplus_{i \geq
0}L_\omega^i
H^{m-2i}_{\text{prim}}(X,\mathbb{Q})$
for $0 < m \leq n$.
(I think this corresponds to the usual
weight space decomposition.)

[Picture of Hodge diamond.
Reflection by vertical axis is complex
conjugation,
180-degree rotation is Poincar'e duality,
$p+q=$constant is a horizontal line,
reflection along horizontal axis is Lefschetz
theorem.
Warning! This depends on conventions of how
we draw the diamond.]
\end{enumerate}
\end{enumerate}

\section{The Hodge-Riemann relations}
\label{mumford-tate-groups-in-hodge-theory-section-hodge-riemann}

\noindent
For all $[\alpha]\in
H^k_{\text{prim}}(X,\mathbb{C})\cap
H^{p,q}(X)$
we have
$$
i^{p-q}(-1)^{\frac{k(k-1)}{2}}
\int_X \alpha \wedge
\overline{\alpha}\wedge\omega^{n-k}>0.
$$
(Here $i$ is the imaginary unit.)

Define a pairing on
$H^k_{\text{prim}}(X,\mathbb{C})$:
$$
\psi([\alpha],[\beta])
=(2\pi i)^{-k}
(-1)^{\frac{k(k-1)}{2}}
\int_X \alpha\wedge\beta\wedge\omega^{n-k}
$$
which is $(-1)^k$-symmetric.

\begin{definition}
\label{mumford-tate-groups-in-hodge-theory-definition-weil-operator}
The {\it Weil operator} $C \in
\text{End}(H^k(X,\mathbb{C}))$ is given by
$C|_{H^{p,q}}(X)=(i)^{p-q}\text{Id}$.
\end{definition}

\noindent
Let $Q([\alpha],[\beta])=(2\pi i)^k
\psi(C[\alpha],[\beta])$.

Then $Q$ is symmetric (exercise) and positive
on $H^k(X,\mathbb{R})$.
Positive is just the Hodge-Riemann relation.

\section{Hodge structures and Mumford-Tate
groups}
\label{mumford-tate-groups-in-hodge-theory-section-hodge-struc-mumfo-tate-grou}
\label{mumford-tate-groups-in-hodge-theory-section-mt-hodge-structures}

\begin{definition}
\label{mumford-tate-groups-in-hodge-theory-definition-rational-hodge-structure}
A {\it rational Hodge structure}
($\mathbb{Q}$-HS)
of weight $k\in\mathbb{Z}$ is a
finite-dimensional
$\mathbb{Q}$-vector space $V$
and a decomposition $V \otimes_{\mathbb{Q}}
\mathbb{C}
=\bigoplus_{p+q=k}V^{p,q}$
such that
$V^{q,p}=\overline{V^{p,q}}\forall p,q$.

A $\mathbb{Q}$-HS (without fixed weight)
is a $\mathbb{Q}$-vector space $V$
with a decomposition
$V=V_{k_1}\oplus … \oplus V_{k_n}$
where $V_{k_i}$ is a $\mathbb{Q}$-HS
of weight $k_i$.

Analogously, define $\mathbb{Z}$-HS,
$\mathbb{R}$-HS, etc
(i.e. take $V$ to be a finitely generated
$\mathbb{Z}$-module, $\mathbb{R}$-vector
space, etc.)
\end{definition}

\begin{example}
\label{mumford-tate-groups-in-hodge-theory-example-hodge-structures}
\begin{enumerate}
\item $X \subset \mathbb{C}P^N$ smooth
subvariety,
then $H^k(X,\mathbb{Z})$ is a $\mathbb{Z}$-HS
of weight $k$.

\item The {\it Tate HS} is
$\mathbb{Z}(1):=2\pi i\mathbb{Z}
=\Ker(\mathbb{C}\xrightarrow{\text{exp}}
\mathbb{C}^*$,
since $\mathbb{C}=
\mathbb{Z}(1)\otimes_{\mathbb{Z}}\mathbb{C}=
\mathbb{Z}(1)^{-1,1}$,
so this is a $\mathbb{Q}$-HS of weight $-2$.

Note that
$H^2(\mathbb{C}P^1,\mathbb{Z})=\mathbb{Z}$
sits inside
$H^2(\mathbb{C}P^1,\mathbb{C}=\mathbb{Z}
\otimes\mathbb{C}
\cong \mathbb{C}
=H^{1,1}(\mathbb{C}P^1)$.

Analogously,
$\mathbb{Q}(1)=2\pi i\mathbb{Q} \subset
\mathbb{C}$
and $\mathbb{Q}(1)\otimes
\mathbb{C}=\mathbb{Q}(1)^{-1,-1}$
is of weight $-2$.
\end{enumerate}
\end{example}

\section{The Deligne torus}
\label{mumford-tate-groups-in-hodge-theory-section-deligne-torus}

\begin{definition}
\label{mumford-tate-groups-in-hodge-theory-definition-deligne-torus}
$\mathbb{S}$ is the algebraic group such that
$\mathbb{S}(\mathbb{R})$ is the
$\mathbb{C}^*=\mathbb{C}\setminus\{0\}$.
\end{definition}

\noindent
The group of real points
$\mathbb{S}(\mathbb{R})$ is a real Lie group.

Note that
$$
\begin{aligned}
\mathbb{C}^*&=
\{(x,y) \in \mathbb{R}^2:x^2+y^2\neq 0\}\\
&=\{(x,y,y)\in \mathbb{R}^2:(x^2+y^2)t=1\},
\end{aligned}
$$

\noindent
which allows to see $\mathbb{C}^\ast$
as the vanishing locus of the polynomial
$(x^2+y^2)t-1$, so to see it as an algebraic
variety.
Now
$$
\begin{aligned}
\mathbb{S}(\mathbb{C})&=
\{(x,y,t)\in \mathbb{C}^3:(x^2+y^2)t=1\}\\
&=\{(x,y,t)\in
\mathbb{C}^3:\underbrace{(x+iy)}_{z}
\underbrace{(x-iy)}_{w}t=1\}\\
&=\{(z,w,t)\in\mathbb{C}^3:zwt=1\}\\
&=
\{(z,w)\in\mathbb{C}^2:z\neq 0, w\neq 0\}
\cong \mathbb{C}^*\times\mathbb{C}^*.
\end{aligned}
$$

\medskip\noindent
Let $V$ be a $\mathbb{Q}$-HS of weight $k$.
Define a representation over $\mathbb{R}$
$\rho:\mathbb{S}(\mathbb{R}) \to
\text{GL}(V\otimes\mathbb{R})$
as follows:
$\forall z \in \mathbb{C}^*\forall v \in
V^{p,q}$
$\rho(z)\cdot v=z^p \overline{z}^qv$
if $v \in V \otimes \mathbb{R}$,
then $v = \sum_{ p+q=k}v^{p,q},$
$v^{q,p}=\overline{v^{p,q}}$,
and then
$\overline{\rho\cdot v}
=\sum_{p+q=k}z^p \overline{z}^q
=\sum_{p+q=k}\overline{z}^pz^qv^{q,p}
=\rho(z)\cdot v$,
so it is in fact a representation.
But why? Why is this equality what we need
to make sure it is a representation?

Observe that $z \in
\mathbb{R}^*\subset\mathbb{S}(\mathbb{R})=
\mathbb{C}^*$,
the eigenvalue is
$\rho(z)\cdot v= z^{p+q}\cdot v=z^k\cdot v$.
This motivates the following:

Conversely,
given a $\mathbb{Q}$-vector space $V$
and a representation
$\rho:\mathbb{S}(\mathbb{R}) \to
\text{GL}(V\otimes_{\mathbb{Q}}\mathbb{R})$
of $\mathbb{R}$-groups, such that
$\forall r \in \mathbb{R}^*$
$\rho (r)=r^k\cdot \text{Id}$,
we have
$V \otimes \mathbb{C}=\bigoplus_{p,q}V^{p,q}$
where $v \in V^{p,q}$,
$\rho(z,w)=z^pw^q\cdot v$.
Since $\rho$ is a representation of
$\mathbb{R}$-groups
we have $V^{q,p}=\overline{V^{p,q}}$,
so by the computation above we have
$V^{p,q}=0$ when $p+q \neq k$,
then $V$ becomes a $\mathbb{Q}$ =HS
of weight $k$.

In conclusion, a $\mathbb{Q}$-HS is the same
thing as a
$\mathbb{Q}$-vector space $V$ and a
representation
of $\mathbb{R}$-groups
$\rho:\mathbb{S}(\mathbb{R}) \to
\text{GL}(V\otimes_{\mathbb{Q}}\mathbb{R})$
such that $\rho |_{\mathbb{R}^*}$
is defined over $\mathbb{Q}$.

(We may also replace $\mathbb{Q}$ with
$\mathbb{Z}$, etc. in this construction.)

\medskip\noindent
Back to the Tate group
$\mathbb{Q}(1)=2\pi i \mathbb{Q} \subset
\mathbb{C}$,
define for all $z \in \mathbb{C}^*$ and
$v \in \mathbb{Q}(1) \otimes \mathbb{C}$
by $\rho(z)\cdot v=|z|^{-2}\cdot v$.

We have operations on HS:
$V_1 \oplus V_2, V_1 \otimes V_2,
\Hom(V_1,V_2)$.
The $\mathbb{Q}$-HS form an abelian category.
Let $\mathbb{Q}(m)=\mathbb{Q}(1)^{\otimes m}$
when $m \geq 0$,
$\mathbb{Q}(-1)=\mathbb{Q}(1)^*$
and $\mathbb{Q}(-m)=\mathbb{Q}(-1)^{\otimes
m}$
for $m \geq 0$.

If $V$ is a $\mathbb{Q}$-HS, then
$V(m)=V \otimes \mathbb{Q}(m)$
is the {\it Tate twist}.

\begin{example}
\label{mumford-tate-groups-in-hodge-theory-example-chern-class}
\begin{enumerate}
\item $X \subset \mathbb{C}P^n$ subvariety.
Notice that while $\mathbb{C}$ is the
algebraic closure of $\mathbb{R}$,
such a closure can be obtained by
choosing the imaginary unit $i$ or $-i$,
so this is not canonical.

However, the first Chern class is canonical:
$$
\xymatrix{
0\ar[r]
&\mathbb{Z}(1)\ar[r]
&\mathbb{C}\ar[r]^{\text{exp}}
&\mathbb{C}^*\ar[r]
&0
}
$$
gives the connecting homomorphism
$H^1(X,\mathbb{C}^*
)\xrightarrow{c_1}H^2(X,\mathbb{Z}(1))$
which is in fact a HS of weight 0
(because $H^2$ has weight 2 and
$\mathbb{Z}(1)$ has
weight $-2$).

Analogously, the $p$-th chern class is
$c_p(A\text{ coh. sheaf}) \in
H^{2p}(X,\mathbb{Q}(p))$.

\item $cl:CH_\mathbb{Q}^p(X) \to
H^{2p}(X,\mathbb{Q})$
the cycle class map, where
$$
CH^p_\mathbb{Q}(X)=\left\{\sum_{\substack{Z_i
\subset X \\
\text{subvar.}}}\alpha_i[Z_i]:\alpha_i \in
\mathbb{Q}\right\}
\Big/\text{rational equiv.}
$$
\end{enumerate}
\end{example}

Consider
$$
\xymatrix{
Z\ar@{^(->}[r]^i& X\\
\tilde{Z} \ar[ur]_j\ar[u]
}
$$
where
$Z$ is a subvariety of codimension $p$ and
$\tilde{Z}$ is a resolution of singularities.
Then we have the pushforward of the
fundamental class,
$j_*[\tilde{Z}] \in H_{2n-2p}(X,\mathbb{Q})$
where $[\tilde{Z}]\in
H_{2n-2p}(\tilde{Z},\mathbb{Q})
\cong H^{2n-2p}(\tilde{Z}, \mathbb{Q})^*$.
Then the Poincare dual of
$j_*[\tilde{Z}]:=d[Z]
\in H^{2p}(X,\mathbb{Q}) \cap H^{p,p}(X)$.
So $H_{2n-2p}(\tilde{Z},\mathbb{Q})$
is a HS of weight $2(p-n)$
and  $[\tilde{Z}] \in H_{p-n,p-n}$.

\begin{definition}
\label{mumford-tate-groups-in-hodge-theory-definition-space-of-hodge-classes}
The {\it space of Hodge classes}
is $H^{2p}_{\text{Hdg}}(X)
=H^{2p}(X,\mathbb{Q})
\cap H^{p,p}(X)$
\end{definition}



\noindent
Then we have
$cl:CH^p_{\mathbb{Q}}(X)
\to H^{2p}_{\text{Hdg}}(X)
\cong
\Hom_{\mathbb{Q}-HS}(\mathbb{Q}(-p),H^{2p}(X,
\mathbb{Q})$.
The Hodge conjecture is that cl
is surjective onto the space of Hodge
classes.

\section{Polarizations}
\label{mumford-tate-groups-in-hodge-theory-section-polarizations}

\noindent
Let $V$ be a $\mathbb{Q}$-HS of weight $k$
and $\rho:\mathbb{S}(\mathbb{R})\to
\text{GL}(V\otimes_{\mathbb{Q}}\mathbb{R})$
the corresponding representation of the
Deligne torus.
The Weil operator in terms of $\rho$ is
$C=\rho(i)$.

\begin{definition}
\label{mumford-tate-groups-in-hodge-theory-definition-polarization}
A {\it polarization} on $V$ is a morphism
of $\mathbb{Q}$-HS is a $(-1)^k$-symmetric
morphism
$\psi:V \otimes V \to \mathbb{Q}(-k)$
such that the bilinear form
$Q:V_{\mathbb{R}}\otimes V_{\mathbb{R}}\to
\mathbb{R}$
given by $Q(x,y)=(2\pi i)^k\psi(Cx,y)$
is symmetric and positive definite.
\end{definition}

\begin{example}
\label{mumford-tate-groups-in-hodge-theory-example-polarization}
$V=H^k_{\text{prim}}(X,\mathbb{Q})$
with $\psi$ the Hodge-Riemann pairing.
\end{example}

\noindent
Observe: if $V$ is polarizable
(i.e. it admits a polarization;
not every HS admits a polarization e.g. HS
on nonprojetive varieties)
then $V$ is semisimple
$V=V_1 \oplus  … \oplus  V_m$,
$V_i$ is simple.

Assume $W \subset V$ is a sub-HS,
then $W^{\perp_\psi} \subset V$
is a sub-HS.
Then
$0 \underbrace{=}_{\substack{Q \\
\text{positive}\\\text{def.}}}
W \cap W^{\perp_\psi} \subset V$,
Then $V=W \oplus  W^{\perp_\psi}$.

\section{The Mumford-Tate group}
\label{mumford-tate-groups-in-hodge-theory-section-mumford-tate-group}

\noindent
Let $V$ be a $\mathbb{Q}$-HS
and $\rho:\mathbb{S}(\mathbb{R}) \to
\text{GL}(V \otimes \mathbb{R})$
a representation of $\mathbb{R}$-groups.

Consider $G \subset \text{GL}(V)$ such that
$\text{Im}(\rho) \subset G(\mathbb{R})$.

(The idea is that the Mumford-Tate group
will recover the rational structure of $V$,
because the representation $\rho$ knows
nothing
about this rational structure, which must
exist
since $V$ is a rational vector space.)

\begin{definition}
\label{mumford-tate-groups-in-hodge-theory-definition-mumford-tate-group}
The {\it Mumford-Tate group} is
$$
MT(V)=\bigcap_{\substack{G \subset
\text{GL}(V) \\
\text{subgroup s.t.}\\
\text{Im}(\rho) \subset G(\mathbb{R})}}G
=
\text{smallest $Q$-subgroup of $\text{GL}(V)$
containing $\text{Im}\rho$}.
$$
\end{definition}

\begin{remark}
\label{mumford-tate-groups-in-hodge-theory-remark-Hodge-group}
We can also consider $U(1) \subset
\mathbb{S}(\mathbb{R})$
and $\rho':=\rho|_{U(1)}: U(1) \to
\text{GL}(V \otimes \mathbb{R})$.
Then the {\it Hodge group} is
$\text{Hdg}(V)=$
smallest $Q$-subgroup of $\text{GL}(V)$ such
that
$\text{Im}(\rho')\subset
\text{Hdg}(V)(\mathbb{R})$.

$\text{Hdg}(V)$
is always smaller than $MT(V)$.
\end{remark}

\begin{example}
\label{mumford-tate-groups-in-hodge-theory-example-mumford-tate-group}
\begin{enumerate}
\item $MT(\mathbb{Q}(1))=\mathbb{Q}^*$
$\mathbb{Q}(1)=2\pi i \mathbb{Q} \subset
\mathbb{C}$,
$\mathbb{Q}(1) \otimes \mathbb{R}=2\pi
i\mathbb{R}\subset \mathbb{C}$,
$\text{GL}(\mathbb{Q}(1)\otimes
\mathbb{R})=\mathbb{R}^*$,
since $z \in \mathbb{S}(\mathbb{R})$ acts as
$|z|^{-2}\cdot \text{Id}$.

\item $\mathbb{Q}(0)=\mathbb{Q} \subset
\mathbb{C}$,
$MT(\mathbb{Q}(0)) = \{1\}$.
In general, if $V$ is of weight $k$,
then $\mathbb{Q}^* =$ center of
$\text{GL}(V)\subset MT(V)$
and $MT(V)$ is generated by  $\mathbb{Q}^*$
and $Hdg(V)$.
\end{enumerate}
\end{example}


\section{Tensor construction}
\label{mumford-tate-groups-in-hodge-theory-section-tensor-construction}

\noindent
Let $V$ be a $\mathbb{Q}$-HS with $MT(V)
\subset \text{GL}(V)$
and $\rho:\mathbb{S}(\mathbb{R}) \to
MT(V)(\mathbb{R})$.
Then
$$
T^\bullet(V)=\bigoplus_{e,f \geq 0}
V^{\otimes e}\otimes (V^* )^{\otimes f}
$$
is a $MT(V)$-representation.

\begin{proposition}
\label{mumford-tate-groups-in-hodge-theory-proposition-sub-HS-iff}
A finite-dimensional subspace $W \subset
T^\bullet(V)$
is a sub-HS if and only if $W$ is
a $MT(V)$-subspace.
\end{proposition}

\begin{proof}
$(\impliedby)$. If $W$ is a
$MT(V)$-subrepresentation,
then from $\rho$ we can compose with the
representation that $MT(V)$ is
to obtain
$\rho':\mathbb{S}(\mathbb{R}) \to \text{GL}(W
\otimes \mathbb{R})$.
Note that $\rho|_{\mathbb{R}^*}$ is defined
over $Q$,
We get that $W$ is a sub-HS.

($\implies$). Assume that $W \subset
T^\bullet(V)$
is a sub-HS, $G - \text{Stab}(W) \subset
\text{GL}(V)$
is a $\mathbb{Q}$-subgroup.
Since $W$ is a sub-HS,
the action of $\mathbb{S}(\mathbb{R})$ on
$T^\bullet(V)$
preserves $W$,
then $\text{Im}(\rho) \subset G(\mathbb{R})$
and thus $MT(V)\subset G$,
which implies that $W$ is a
$MT(V)$-subrepresentation.
\end{proof}


As a corollary,

\begin{lemma}
\label{mumford-tate-groups-in-hodge-theory-lemma-mumford-tate-invariant-iff}
$x \in T^\bullet(V)$ is $MT(V)$-invariant
if and only if $X$ is a $(0,0)$ Hodge
element.
\end{lemma}

\begin{proposition}
\label{mumford-tate-groups-in-hodge-theory-proposition}
Assume that $V$ is a $\mathbb{Q}$-HS of
weight $k$
and $\psi:V \otimes V \to \mathbb{Q}(-k)$ is
a polarization.
Then
$$
Hdg(V) \subset
\begin{cases}
\text{SO}(V,\psi)\qquad &\text{ if }k\equiv
0(\text{mod }2) \\
Sp(V,\psi)\qquad &\text{ if }k \equiv
1(\text{mod }2)
\end{cases}
$$
\end{proposition}

\begin{proof}
$\psi: V \otimes V \to \mathbb{Q}(-k)$,
where the $\mathbb{Q}(-k)$ is a trivial
$Hdg(V)$-module,
so that $\psi$ is a morphism of
$Hdg(V)$-modules.
The action of $Hdg(V)$ preserves the
symmetric form
$\psi$ if $k \equiv 0 (\text{mod }2)$
and the antisymmetric form $\psi$
if $k \equiv 1 (\text{mod }2)$.
\end{proof}

\begin{example}
\label{mumford-tate-groups-in-hodge-theory-example-elliptic-curve}
Let $V = H^1(E,\mathbb{Q})$ where $E$ is an
elliptic curve
and $\psi$ the Hodge-Riemann pairing.
Then $V\otimes \mathbb{C}=V^{1,0}\oplus
V^{0,1}$
and $\dim V^{1,0}=1$.
$Hdg(V) \subset Sp(V,\psi) \cong
\text{SL}_2(\mathbb{Q})$.

Observe thati, in general, if $V$ is
polarizable,
then $V=V_1 \oplus  … \oplus  V_m$, $V_i$ are
irreducible $MT(V)$-representations,
so $MT(V)$ is reductive.

Back to our elliptic curve example,
$Hdg(V) \subset \text{SL}_2(\mathbb{Q})$ is
reductive.

There are two possibilities:
\begin{enumerate}
\item $Hdg(V)=\text{SL}^2(\mathbb{Q})$,
which implies that $MT(V) = \text{GL}_2(V)$.
This happens when $E$ is generic in the
moduli space.

\item $Hdg(V)$ is properly contained in
$\text{SL}_2(\mathbb{Q})$.
Then $Hdg(V)$ is a 1-dimensional torus.
Then $\text{End}_{\mathbb{Q}-HS}(V) \neq
\mathbb{Q}$,
which implies that $E$ has complex
multiplication.
\end{enumerate}
\end{example}

\begin{definition}
\label{mumford-tate-groups-in-hodge-theory-definition-cm-type}
A HS $V$ is of  {\it CM-type} if $MT(V)$ is
abelian
(such HS defines a special point in the
moduli space of HS.)
\end{definition}


\section{Summary so far}
\label{mumford-tate-groups-in-hodge-theory-section-summary}

\noindent
Recall:
\begin{itemize}
\item A $\mathbb{Q}$-HS is a finite
dimensional $\mathbb{Q}$
vector space $V$ with a representation of
$\mathbb{R}$-group
$\rho:\underbrace{\mathbb{S}(\mathbb{R})}
_{\cong\mathbb{C}^*}
\to \text{GL}(V \otimes \mathbb{R})$
such that $\rho|_{\mathbb{R}^*}$ is defined
over $\mathbb{Q}$.

This gives a Hodge structure
$V\otimes \mathbb{C}=\bigoplus_{p,q}V^{p,q}$,
$V^{p,q}=\overline{V^{q,p}}$.

$V$ is of weight $k$ if
$\forall  r \in \mathbb{R}^k$,
$\rho(r)=r^k\cdot \text{Id}$.
$\forall z \in \mathbb{S}(\mathbb{R})$,
$\rho(z)|_{V^{p,q}}=z^p\overline{v}
^q\text{Id}$

\item $MT(V)=\bigcap_{\substack{G \in
\text{GL}(V) \\
\text{s.t.}\text{Im}(\rho) \subset
G(\mathbb{R})}}G$
is the Mumford-Tate group.

\item $Hdg(V)=\bigcap_{\substack{
G \in \text{GL}(V)\\
\text{s.t.}\text{Im}(\rho|_{\text{U}(1)})
\subset
G(\mathbb{R})}}G$
is the Hodge group.

\item $T^\bullet V=\bigoplus_{e,f \geq 0}
V^{\otimes e}\otimes(V^* )^{\otimes f}$.
A sub-HS of $T^\bullet V$ is the same thing
as a
sub-$MT(V)$-representation.
$T^\bullet V^{MT(V)}=$ $MT(V)$-invariants.
$T^\bullet V \cap(T^*V)^{0,0}=$ the space of
Hodge $(0,0)$-classes.

\item (Polarization.)
Assume that $V$ has weight $k$
$\psi:V \otimes V \to \mathbb{Q}(-k)$ a
$\mathbb{Q}$-HS morphism
such that the following form is positive
definite:
$$
\begin{aligned}
Q: V_{\mathbb{R}}\otimes V_{\mathbb{R}}
&\longrightarrow \mathbb{R} \\
Q(x,y) &=(2\pi i)^k \psi(Cx,y)
\end{aligned}
$$

\noindent
where $C=\rho(i) \in Hdg(V)(\mathbb{R})$,
then $MT(V),Hdg(V)$ are reductive.

\end{itemize}

\noindent
``We are interested in polarizable Hodge
structures
because that is the case of cohomologies of
algebraic varieties.''

\section{Some facts about reductive groups}
\label{mumford-tate-groups-in-hodge-theory-section-reductive-groups}

\noindent
Let $V$ be a finite dimensional vector space
over $k$, $\text{char}(k)=0$
and $G \subset \text{GL}(V)$ an irreducible
algebraic $k$-group.

\begin{remark}
\label{mumford-tate-groups-in-hodge-theory-remark-connected}
$MT(V)$ and $Hdg(V)$ are connected.
\end{remark}

\noindent
Since by definition $G$ is a subgroup of
$\text{GL}(V)$,
we can view $V$ as a faithful representation
(because the inclusion has no kernel).
We say $G$ is {\it reductive} if $V$ is
semisimple,
i.e.  $V$ is a direct sum of irreducible
representations
$V=V_1 \oplus  … \oplus  V_n$.
Equivalently, Every $G$-representation is
semisimple.
Equivalently, the center  $Z^0(G)$ is an
algebraic torus
and $G/Z(G)$ is semisimple.
Equivalently, $G(\mathbb{C})$ admits a
compact real form.

$\sigma\in \text{Aut}(G(\mathbb{C}))$
is a {\it Cartan involution}
if $\sigma^2=\text{id}$ and $G^\sigma=\{g\in
G(\mathbb{C}):\sigma(\bar{g})=g \}$
is compact.

\begin{example}
\label{mumford-tate-groups-in-hodge-theory-example-polarizable-hs}
Let $V$ be a polarizable $\mathbb{Q}$-HS of
weight $k$.
$C^2=\rho(i^2)=\rho(-1)\in
Z(\underbrace{Hdg(V)}_{G})$,
$G=Hdg(V)$.

Note that  $\sigma=Ad_C:G \to G, g \mapsto
CgC^{-1}$,
$\text{id}=Ad_C^2, g \mapsto  C^2g C^{-2}=g$.

We claim that $\sigma$ is a Cartan involution
on $Hdg(V)=G$.
Indeed,  $Q$ defines a Hermitian product on
$V \otimes \mathbb{C}$.
$g \in G^\sigma$,
$$
\begin{aligned}
Q(gx,\overline{gy})&=
Q(gx,\sigma(g)\bar{y})\\
&=(2\pi i)^k
\psi(\underbrace{Cg}_{\in
Hdg}x,\underbrace{Cg}_{\in
Hdg}C^{-1}\bar{y})\\
&=(2\pi i)^k\psi(x,C^{-1}\bar{y})\\
&=(2\pi i)^k\psi(Cx,\bar{y})\\
&=Q(x,\bar{y}).
\end{aligned}
$$

\noindent
Thus, $G^\sigma \subset \underbrace{U(V
\otimes \mathbb{C}}_{\text{compact}}$
as claimed.
\end{example}

\begin{remark}
\label{mumford-tate-groups-in-hodge-theory-remark-polarizable}
A $\mathbb{Q}$-HS is polarizable if and only
if
$G=Ad_C$ is a Cartan involution.
\end{remark}

\section{Reductive groups and their
representations}
\label{mumford-tate-groups-in-hodge-theory-section-reduc-group-their-repre}

\noindent
Let $V$ be a finite-dimensional vector space
over $k$, $\text{char}k=0$.
$T^\bullet V=\bigoplus_{e,f \geq 0}V^{\otimes
e}\otimes (V^* )^{\otimes f}$.

\begin{theorem}
\label{mumford-tate-groups-in-hodge-theory-theorem-reduc-group-repre}
Let $G\subset \text{GL}(V)$ be a reductive
group.
Then
\begin{enumerate}
\item Any finite-dimensional representation
of $G$
is a subgroup of $(T^\bullet V)^{\oplus  N}$
for some  $N>0$.

\item Assume that $H \subset G$
is an algebraic subgroup.
There exists a $G$-representation $W$
and a 1-dimensional subspace $L \subset W$
such that $H=Stab_W(L)$.

\item If $H$ is reductive,
then there exists a $G$-representation $W$
and $x \in W$ such that $H=Stab_W(x)$.
\end{enumerate}
\end{theorem}

\noindent
We need some preparations for the proof:

\begin{enumerate}
\item $G\subset \text{GL}(V)$
Let $k[G]$ be the algebra of regular
functions on  $G$.

$$
\xymatrix{
\text{GL}(V)\ar[rr]\ar[ddrr]&&\text{End}(V)
\oplus  \text{End}(V^* )\\
g \ar@{|->}[rd]\\
&(g,(g^t)^{-1})&
\{(A,B) : A \cdot
B^t=\text{Id}\}\ar@{^{(}->}[uu].
}
$$
So we have
$$
\text{Sym}(V \otimes V^* )^{\otimes 2}
\to\!\!\!\!\!\to K[\text{GL}(V)]
\to\!\!\!\!\!\to K[G].
$$

\item $K[G]$ is a Hopf algebra. Indeed, we
have
\begin{align*}
\begin{aligned}
G\times G &  \xrightarrow{m}G\\
\{1\}&\xymatrix{\ar@{^{(}->}[r]^e&} G\\
G &  \xrightarrow{inv}G
\end{aligned}\qquad
\begin{aligned}
K[G] &\xrightarrow{\Delta}K[G]\otimes K[G]\\
K[G]&  \xrightarrow{\varepsilon}K\\
K[G]&\xrightarrow{coinv}K[G].
\end{aligned}
\end{align*}

\noindent
and
$$
\xymatrix{
G \times G \times G\ar[r]^{m \times
\text{id}}\ar[d]^{\text{id} \times m}
&G \times G\ar[d]^{m}\\
G \times G\ar[r]^{m}
&G
}
\qquad
\xymatrix{
K[C]\ar[r]^{\Delta}\ar[d]^{\Delta}
&K[G]^{\otimes 2}\ar[d]^{\text{id} \otimes
\Delta}\\
K[G]^{\otimes 2}\ar[r]^{\Delta \otimes
\text{id}}
&K[G]
}
$$

Then $V$ is a $G$-representation, i.e. a {\it
comodule}.
We have
$$
\begin{aligned}
G \times V^*&\to V^*\\
\text{Sym}(V) & \to \text{Sym}(V) \otimes
K[G]\\
V & \to V \otimes K[G]\\
\mu:V & \to V \otimes K[G]
\end{aligned}
$$

and the following diagram commutes
\begin{equation}
\label{mumford-tate-groups-in-hodge-theory-equation-diagram}
\xymatrix{
V\ar[r]^{\mu}\ar[d]_{\mu}
&V \otimes K[G]\ar[d]^{\mu \otimes
\text{id}}\\
V \otimes K[G]\ar[r]_{\text{id} \otimes
\Delta}
&V \otimes K[G] \otimes K[G].
}
\end{equation}

\noindent
Equation \ref{mumford-tate-groups-in-hodge-theory-equation-diagram} defines a
$K[G]$-comodule
structure on $K[G]$.

\item Any $G$-representation is a union
of finite-dimensional subrepresentations.
(In particular this applies to $K[G]$.)

Indeed, let $\{a_i\}$ be a basis of $K[G]$.
Fix $x \in V$ and denote
$$
\mu(x)=\sum_{i, \text{finite}} v_i \otimes
e_i
\qquad \text{ and } \qquad
\Delta e_i=\sum_{j,k}a_{ijk}e_j \otimes
e_k,\qquad
a_{ijk}\in k.
$$
Then
$$
\begin{aligned}
\sum_i\mu(v_i)\otimes e_i&=
\sum_{i,j,k}a_{i j k}v_i \otimes e_j \otimes
e_k\\
&=\sum_i\left(\sum_{j,k}
a_{i j k}v_k \otimes v_k \otimes e_j \right)
\otimes e_i.
\end{aligned}
$$

\noindent
$\implies \mu(v_i)=\sum_{j,k}a_{k j i}v_k
\otimes e_j$.
$W=\text{span}(x,v_i)\implies \mu:W \to W
\otimes K[G]$,
$\dim W < \infty$.
\end{enumerate}

Now we can prove Theorem
\ref{mumford-tate-groups-in-hodge-theory-theorem-reduc-group-repre}.

\begin{proof}
\begin{enumerate}
\item Let $W$ be a finite-dimensional
$G$-representation.
It is enough to show that $W$
is a subrepresentation of $K[G]^{\oplus N}$
for some $N>0$.
\begin{equation}
\label{mumford-tate-groups-in-hodge-theory-equation-diagram2}
\xymatrix{
W\ar[r]^{\mu}\ar[d]_{\mu}
&W \otimes K[G]\ar[d]^{\mu \otimes
\text{id}}\\
W \otimes K[G]\ar[r]_{\text{id} \otimes
\Delta}
& W \otimes K[G] \otimes K[G]
}
\end{equation}

\noindent
Equation \ref{mumford-tate-groups-in-hodge-theory-equation-diagram2}
means that $\mu$ defines a morphism of
$K[G]$-comodules
$$
W \xymatrix{\ar@{^{(}->}[r]&}\underbrace{W
\otimes K[G]}_{
K[G]^{\oplus \dim W}}.
$$
\item Let $H \subset G$ be an algebraic
$k$-subgroup.
[We want to show] There exists a
$G$-representation $W$
and $L \subset W$ with $\dim L=1$ such that
$H=Stab_W(L)$.

$k[H]=k[G]/I$ for an ideal  $I$.
For $f \in K[G],h,g \in G$ we have
$gf(h)=f(g^{-1}h)$.

Suppose that  $H=V(I)$ for an ideal  $I$.
Note that $g \in H$ if and only if $g\cdot
I=I$
since $g\cdot H=V(g I)$
since $0=f(x)=f(g^{-1}gx)=gf(gx)=0$.
That is, $H=Stab_{K[G]}(I)$.

Let $f_1,…,f_m$
be generators of $I$.
Then there exists $\tilde{W} \subset k[G]$,
a finite-dimensional $G$-subgroup
with $f_1,…,f_m \in W$,
$H=Stab_{\tilde{W}}(\tilde{W} \cap I)$.
Consider $W = \Lambda^d\tilde{W}$
and $L=\Lambda^d(W \cap I)$, which satisfy
the conditions.


\item If $H$ is reductive,  $H=Stab_W(L)$
$W=L \oplus  W'$,
$W \otimes W^* =(\underbrace{L \otimes
L^*}_{\text{trivial}}
\oplus (…)$.
Let $0\neq x \in L \otimes L^*$.
Then $H=Stab_{W \otimes W^*}Stab(x)$.

\end{enumerate}
\end{proof}

\noindent
As a corollary,

\begin{lemma}
\label{mumford-tate-groups-in-hodge-theory-lemma-mt-reduc-polar-hs}
If $V$ is a polarizable $\mathbb{Q}$-HS,
then [the Mumford-Tate group is reductive]
$MT(V)=$ Stabilizer of all $(0,0)$-Hodge
classes
in $T^\bullet V$.
\end{lemma}

\begin{proof}
Let $G=MT(V)$ and
$$
G'=\bigcap_{\substack{x \in
(0,0)\text{-Hodge}\\\text{classes}}}Stab(x)
\subset \text{GL}(V).
$$
$G \subset G'$.

By part 3 of Theorem \ref{mumford-tate-groups-in-hodge-theory-theorem},
$G=Stab_W(x)$ where $W$ is a representation
of $\text{GL}(V)$.
By part 1 of Theorem \ref{mumford-tate-groups-in-hodge-theory-theorem},
$W \subset (T^\bullet V)^{\oplus N}$.
$x$ is fixed by $G=MT(V)$,
then  $x$ is a $(0,0)$-Hodge class.
Thus $G' \subset G$.
\end{proof}

\noindent

[Approximate comment]
We have shown that the representations of
$MT(V)$
live in that tensor algebra.
We would like to find a universal object
such that any Hodge structure is realised as
a
representation of such group.
This is analogous to Galois theory,
where we take an inverse limit over Glois
groups
of field extensions,
which provides a universal object and
an equivalence of categories.

\section{Short summary}
\label{mumford-tate-groups-in-hodge-theory-section-summary2}

$V$ a $\mathbb{Q}$-Hs, polarizable.
$\rho:\mathbb{S}(\mathbb{R})\to \text{GL}(V
\otimes \mathbb{R})$,
$MT(V)\subset\text{GL}(V)$, the smallest
subgroup over $\mathbb{Q}$ such that
$MT(V)(\mathbb{R})$ contains
$\text{Im}(\rho)$.
Polarizable $\implies$ $MT(V)$ is reductive.
$MT(V)=$ stabilizer of all $(0,0)$ Hodge
classes
in $T^\bullet V=\bigoplus_{e,f\geq
0}V^{\otimes e}\otimes(V^* )^{\otimes f}$.

\section{Absolute Mumford-Tate group}
\label{mumford-tate-groups-in-hodge-theory-section-absolute}

Let $V' \subset V$ be a sub-HS.
$MT(V) \subset
Stab_V(V')=\left\{\begin{pmatrix}\ast&\ast\\
0&\ast\end{pmatrix} \right\} $.

We also have a map
$MT(V) \to\!\!\!\!\!\to MT(V')$.

\medskip\noindent
The absolute MT group should be something
like the inverse limit
$\lim_{\xymatrix{&\ar[l]}}MT(V) $.

\medskip\noindent
Let $(\mathbb{Q}-HS)$ be the category
of $\mathbb{Q}$-HS.
\begin{enumerate}
\item $\mathbb{Q}$-linear abelian category.

\item Has $\otimes$, which is commutative
with unit $1=\mathbb{Q}(0)$,
which a Hodge structure such that
$\text{End}(\mathbb{Q}(0))=\mathbb{Q}$.
We also have duality, $\forall V\exists V^*$
and $\mathbb{Q}(0)=q \to V \otimes V^*$
which is the ``diagonal embedding'',
and $V \otimes V^* \to 1=\mathbb{Q}(0)$,
which is the trace.

\item There exists a functor
$\omega:(\mathbb{Q}-HS)\to
(\mathbb{Q}-v.s.)$,
forgetful functor.
$\omega$ is called {\it fibre functor}.
A $\otimes$-functor that is exact and
faithful
(injective on Hom's).
\end{enumerate}

\noindent
The pair $(\mathbb{Q}$-HS, $\omega)$ is
called
a {\it neutral Tannakian category}.

\begin{theorem}
\label{mumford-tate-groups-in-hodge-theory-theorem-tannakian}
A neutral Tannakian category
$(\mathcal{C},\omega)$
(where $\mathcal{C}$ is $K$-linear,
$\text{char}K=0$)
is equivalent to the category of finite
dimensional representations of an affine
group scheme over $K$
$Spc K[G]$, $K[G]$ Hopf algebra.
$G$ is a pro-alg. group called the {\it
Tannakian fundamental group}.
\end{theorem}

\noindent
Apply this to $(\mathbb{Q}-HS),\omega$,
$\implies $ $G=MT_{\mathbb{Q}-HS}$ is the
{\it absolute MT group}.

Let  $(\mathbb{Q}-HS)^{pol}$ be the full
$\otimes$-subcategory
of polarizable HS.
$MT_{(\mathbb{Q}-HS)^{\text{pol}}}$ is
pro-reductive.

\section{How to construct the Tannakian
fundamental group}
\label{mumford-tate-groups-in-hodge-theory-section-tannakian-fundamental-group}

\noindent
We have $\omega:\mathcal{C}\to (K-v.s.)$.
Let $\text{Aut}(\omega)^{\otimes}$ be the
automorphisms
ot $\omega$ as a $\otimes$-functor.
That is, if $g \in
\text{Aut}(\omega)^{\otimes}$,
we have $\forall  X \in \mathcal{C}$, $g_x\in
\text{GL}(\omega(X))$.

We also have:

\begin{itemize}
\item
Compatibility with morphisms,
i.e. for every $f:X \to Y$ we have
$$
\xymatrix{
  \omega(X)\ar[r]^{\omega(f)}\ar[d]_{g_X}
&\omega(Y)\ar[d]^{g_Y}\\
\omega(X)\ar[r]_{\omega(f)}
&\omega(Y)
}
$$
commutes, i.e. $g_Y \circ \omega(f)=\omega(f)
\circ g_X$.

\item Compatibility with $\otimes$:
$$
\begin{aligned}
g_{X \otimes Y}&=g_X \otimes g_Y\\
g_1&=1 \in K^* = \text{GL}(K),
  g_{X^*}=(g^*_X)^{-1}.
\end{aligned}
$$
\end{itemize}

\noindent
$\text{Aut}(\omega)^\otimes$ is the group of
$K$-points of
the Tannakian fundamental group.

\medskip\noindent
[Our category $(\mathbb{Q}-HS)^{\text{pol}}$
is equivalent to the
category of representations of this group.
But the problem is that this group is huge.
So we shall restrict to a smaller
subcategory.]

Let $V \in (\mathbb{Q}-HS)^{\text{pol}}$
and $\left\langle{}V\right\rangle{}$ the full
$\otimes$-subcategory generated by $V$.
Objects are sub-HS of $T^\bullet V$ and their
direct sums.

What is the Tannakian fundamental group of
$\left\langle{}V\right\rangle{}$?

Consider $\omega:\left\langle{}V\right\rangle{}\to
(\mathbb{Q}-v.s.)$.
For $g \in \text{Aut}(\omega)^{\otimes}$
we have that $g_V \in \text{GL}(V)$
is uniquely determined by $g$.

Let $x \in T^\bullet V \cap(T^\bullet
V)^{0,0}$.
We have an embedding
$$
\begin{aligned}
1=\mathbb{Q}(0)&\xymatrix{\ar@{^{(}->}[r]&}
T^\bullet
V\\
1 & \mapsto x.
\end{aligned}
$$

\noindent
We have $g_{\mathbb{Q}(0)}=1 \implies g_V \in
\text{Stab}_{T^\bullet V}(x)$.
Thus  $g_V \in MT(V)$.
We also have the converse, so we see that
the Tannakian fundamental group of
$\left\langle{}V\right\rangle{}$
is $MT(V)$.
This could be taken as an equivalent
definition of $MT(V)$


\section{Hodge structures of abelian
varieties}
\label{mumford-tate-groups-in-hodge-theory-section-hodge-struc-abeli-varie}

\noindent
Let $(\mathbb{Q}-HS)^{\text{ab}}$
be the full $\otimes$-subcategory of
$(\mathbb{Q}-HS)^{\text{pol}}$
genertated by $H^1(A,\mathbb{Q})$
where $A$ is an abelian varieties.

Recall that if $A$ is a projective variety it
is biholomorphic to $\mathbb{C}^n/\Lambda$
and $H^1(A,\mathbb{C})=H^{1,0}(A) \oplus
H^{0,1}(A)$.

\medskip\noindent
[Suppose that you have some Hodge structure,
say, on the cohomology of your favourite
variety.
Does it belong to this category or not?
It is not clear --- there might be some
complicated Hodge structures in there]

Let $V=H^1(A,\mathbb{Q})$, $MT(V) \subset
\text{GL}(V)$.
Let
$\mathfrak{mt}(V)=\text{Lie}(MT(V))\subset
\mathfrak{gl}(V)=V \otimes V^*$,
a sub-HS (of $\mathfrak{gl}(V)$, which is
always a Hodge structure).
That is, the Lie algebra of $MT(V)$ always
carries a Hodge structure.
This HS may only have types
$(-1,1),(0,0),(1,-1)$
because $((1,0)+(0,1)) \otimes
((-1,0),(0,-1))$.

\begin{proposition}
\label{mumford-tate-groups-in-hodge-theory-proposition-lie-algebra-of-mtv}
Let $V \in (\mathbb{Q}-HS)^{\text{ab}}$.
Then the HS on $\mathfrak{mt}(V)$ may have
only Hodge types
$(-1,1),(0,0),(1,-1)$.
\end{proposition}

\begin{proof}
We may assume that $V \subset T^\bullet W$
where $W=H^1(A,\mathbb{Q})$ by the tensor
construction.
Then we must have $MT(W)
\to\!\!\!\!\!\to MT(V)$,
which on the level of Lie algebras becomes
$\mathfrak{mt}(W)
\to\!\!\!\!\!\to\mathfrak{mt}(V)$,
which is a surjection of HS.
\end{proof}

\noindent
The next example shows how we may use this
result.

\begin{example}
\label{mumford-tate-groups-in-hodge-theory-example-hodge-type-prop}
\begin{enumerate}
\item Let $X=$ a K3 surface, $X \subset
\mathbb{P}^3$, a smooth quartic.
Then $K_X=\mathcal{O}_X$,
$h^{2,0}=h^{0,2}=1$,
$H^2(X,\mathbb{C})=H^{2,0}\oplus
H^{1,1}\oplus H^{0,2}$.
The HS is $V=H^2(X,\mathbb{Q})$
and the intersection product is a map
$V \otimes V \to \mathbb{Q}(-2)$.
Recall that the Hodge group is the image of
$U(n)$ under the representation
from the Deligne torus.
In this case we have
$MT(V) \supset Hdg(V)\subset
\text{SO}(V,\psi)$.
Also
$$
\mathfrak{mt}(V)=\underbrace{\mathfrak{hdg}
(V)}_{\text{Lie}(Hdg(V))}
\oplus
\underbrace{\mathbb{Q}(0)}_{\text{center
of }\mathfrak{gl}(V)}.
$$
And also $(1,-1),(0,0),(-1,1)$.
Now since $H^{2,0}$ is 1-dimensional we can
compute
$$
\Lambda^2(V \otimes \mathbb{C})
=\underbrace{(H^{2,0}\otimes H^{1,1})}_{3,1}
\oplus \underbrace{(\Lambda^2H^{1,1}\oplus
H^{2,0}\otimes H^{0,2})}_{2,2}
\oplus\underbrace{(H^{1,1}\otimes
H^{0,2})}_{1,3}
$$
\item $X \subset \mathbb{C}P^3$ hypersurface
of degree $\geq 5$.
Then $h^{2,0}>1$, $V=H^2(X,\mathbb{Q})$,
$\mathfrak{hdg}(V) \subset
\mathfrak{so}(V,\psi)$.
Hodge types:
$(2,2),(1,-1),(0,0),(-1,1),(-2,-2)$.
\end{enumerate}
\end{example}

\section{Variations of Hodge structures}
\label{mumford-tate-groups-in-hodge-theory-section-variations-of-hs}

\noindent
A {\it family} is $\pi:\mathcal{X} \to B$
with $\mathcal{X},B$ complex
manifolds, $\pi$ is a submersion, $B$ is
connected
and $\mathcal{X} \xymatrix{\ar@{^{(}->}[r]&}
\mathbb{P}^N \times B$.

Then $R^k\pi_*\mathbb{Z}/\text{tors.}$ is a
local system over $B$.
In fact,
$(R^k\pi_*\mathbb{Z})_t=H^k(\mathcal{X}_t,
\mathbb{Z})$
where $\mathcal{X}_t=\pi^{-1}(t)$.
That is, $R^k\pi_*\mathbb{Z}$ is a sheaf with
the information
of the cohomologies of the fibres.

\medskip\noindent
Consider
$$
\xymatrix{
0\ar[r]
&\pi^*\Omega^1_B\ar[r]
&\Omega^1_\mathcal{X}\ar[r]
&\Omega_{\mathcal{X}/B}^1\ar[r]
&0
}
$$
Assume $\dim B=1$
(This is used to prove the property of the
Gauss-Manin connection, though it can be
also be done by restricting to curves.)
Then
$$
\xymatrix{
0\ar[r]
&\pi^*\Omega^1_B \otimes
\Omega_{\mathcal{X}}^{k-1}\ar[r]
&\Omega_{\mathcal{X}}^k\ar[r]
&\Omega_{\mathcal{X}/B}^k\ar[r]
&0
}
$$
\begin{equation}
\label{mumford-tate-groups-in-hodge-theory-equation-diagram3}
\xymatrix{
0\ar[r]
&\pi^*\Omega_B^1 \otimes
\Omega_{\mathcal{X}/B}^\bullet[-1]\ar[r]
&\Omega_{\mathcal{X}}^\bullet\ar[r]
&\Omega^\bullet_{\mathcal{X}/B}\simeq
\pi^{-1}\mathcal{O}_B\ar[r]
&0
}
\end{equation}

We can construct the vector bundle
$$
\mathcal{V}_k=(R^k\pi_*\mathbb{C})
\otimes\mathcal{O}_B
=R^k\pi_*(\pi^{-1}\mathcal{O}_\Big)
\simeq
R^k\pi_*(\Omega^\bullet_{\mathcal{X}/B}).
$$

On $\mathcal{V}_k$ we have a flat connection
from Equation \ref{mumford-tate-groups-in-hodge-theory-equation-diagram3}.
We get
$$
\xymatrix{
  R^k\pi_*(\Omega^\bullet_{\mathcal{X}/B}
  \ar[r]
&R^{k+1}\pi_*(\pi^*\Omega^1_B \otimes
\Omega^\bullet_{\mathcal{X}/B}[-1]\ar[r]
&\mathcal{V}_k\ar[r]^{\nabla}
&\Omega_B^1 \otimes
R^k\pi_*\Omega^\bullet_{\mathcal{X}/B}
=\Omega^1_B \otimes \mathcal{V}^k.
}
$$
The map $\nabla$ is called the {\it
Gauss-Manin} connection.

\medskip\noindent
Define
$F^pV_k=R^k\pi_*(\Omega_{\mathcal{X}/B}^{\geq
p})$.
Then $F^p\mathcal{V}_k \subset \mathcal{V}_k$
is a holomorphic subbundle.

The {\it Hodge filtration} on
$\mathcal{V}_k$.
\begin{enumerate}
\item $\nabla:\mathcal{V}_k \to \Omega_B^1
\otimes \mathcal{V}_k$
a flat connection.

\item The Hodge filtration
$$
0 \subset … \subset F^p\mathcal{V}_k
\subset F^{p-1}\mathcal{V}_k
\subset…\subset V_k.
$$
\item (Griffiths transversality.)
When we restrict this construction to $F^p$
the relation we obtain is {\it Griffiths
transversality}.
$\nabla(F^k\mathcal{V}_k) \subset
\Omega^1_B \otimes F^{p-1}\mathcal{V}_k$.
\end{enumerate}

\begin{definition}
\label{mumford-tate-groups-in-hodge-theory-definition-polarization-on-vk}
A {\it polarization} on $\mathcal{V}_k$
is defined by a section of
$R^2\pi_*\mathbb{Z}(1)$.
\end{definition}

\noindent
By the exponential sequence
$$
\xymatrix{
0\ar[r]
&\mathbb{Z}(1)\ar[r]
&\mathbb{C}\ar[r]^{\text{exp}}
&\mathbb{C}^*\ar[r]
&1
}
$$
we obtain maps
$$
R^1\pi_*\mathbb{C}^*\to
R^2\pi_*\mathbb{Z}(1),
$$
and
$$
\mathcal{O}_{\mathbb{P}^N}(1) \in
H^1(\mathcal{X},\mathbb{C}^*)
\to H^0(B,R^1\pi_*\mathbb{C}^*)
\to H^0(B,R^2\pi_*\mathbb{Z}(1)).
$$

From there we get
$$
\psi:V_B \otimes V_B \to \mathbb{Z}_B(-k),
$$
where $V_B=R^k\pi_*\mathbb{Z}/\text{tors.}$,
a morphism of local systems that defines
polarizations on the fibres of $\pi$.

\begin{definition}
\label{mumford-tate-groups-in-hodge-theory-definition-z-variation}
A {\it $\mathbb{Z}$-variation of
Hodge-Structure over $B$}
is a $\mathbb{Z}$-local system $V_B$ with a
fibre  $\mathbb{Z}^r$ for some $r\geq 0$
and a finite filtration
of holomorphic subbundles on $\mathcal{V}=V_B
\otimes \mathcal{O}_B$
$$
0\subset…\subset F^p\mathcal{V}
\subset F^{p-1}\mathcal{V}
\subset…\subset \mathcal{V}
$$
[which satisfies Griffiths transversality]
such that
\begin{enumerate}
\item for all $t \in B$ $(V_{B,t},F^\bullet
\mathcal{V}_t)$
is a $\mathbb{Z}$-HS.
\item $\nabla(F^p\mathcal{V})\subset
\Omega^1_B \otimes F^{p-1}\mathcal{V}$
where $\nabla$ is the flat connection on
$\mathcal{V}$.
\end{enumerate}

\end{definition}

\section{The Hodge locus}
\label{mumford-tate-groups-in-hodge-theory-section-hodge-locus}

\noindent
[Next we shall try to study the MT group
on families. Since MT essentially is the
satbilizer of the Hodge
classes, we look for those. There are fibres
that may
have more Hodge classes than others. This
leads to the following concept.]

Let $B$ be a complex manifold and $V_B$ a
$\mathbb{Z}$-local system
with a VHS. $F^p\mathcal{V}$ of weight $2k$.
We have
$$
\xymatrix{
\text{Tot}(\mathcal{V})\ar[r]^p
& B\\
\text{Tot}(V_B)\ar@{^{(}->}[u]\ar[ur]
_{\text{covering}}.
}
$$
Note that $\text{Tot}(V_B)=\amalg W_i$
where $W_i$ are the connected components of
$\text{Tot}(V_B)$.

In this setting, a Hodge class is a point $x$
in one of $W_i$
such that $x \in F^k\mathcal{V} \cap
\overline{F^k\mathcal{V}}$
$\iff$
$x \in F^k\mathcal{V} \cap W_i$.

Let $Z_i=\text{Tot}(F^k\mathcal{V})\cap W_i$.
Note that $Z_i$ is a closed subvariety (maybe
singular)
of $\text{Tot}(\mathcal{V})$.

If the projections of these $Z_i$, $p(Z_i)$,
is not equal to $B$,
then $p(Z_i)$ is in the Hodge locus.
That is, the {\it Hodge locus} is
$$
\begin{aligned}
B^{\text{Hdg}}(V_B)&=\bigcup_{i:p(Z_i)\neq
B}p(Z_i)\\
B^{\text{Hdg}}=\bigcup_{\substack{e,f \\
\text{countable}\\\text{union}}}
B^{\text{Hdg}}(\underbrace{T^{e,f}V_B}_{
V_B^{\otimes e}\otimes(V^*_B)^{\otimes f}}).
\end{aligned}
$$

\noindent
[Our goal is to consider the complement of
the Hodge locus and study there the MT
group.]

\section{Third summary}
\label{mumford-tate-groups-in-hodge-theory-section-summary-duplicate-2}

\noindent
Recall: $\mathbb{Z}$-VHS of weight $k$
\begin{itemize}
\item $B$ complex manifold,

\item $V_B$ a local system over $B$ with
fibre $\mathbb{Z}^2$,
$\mathcal{V}=V_B \otimes \mathcal{O}_B$,

\item The Hodge filtration $F^p\mathcal{V}
\subset \mathcal{V}$
(filtration by holomorphic subbundles)
$$
0=F^{k+1}\mathcal{V} \subset F^k\mathcal{V}
\subset…\subset F^0\mathcal{V}=\mathcal{V}.
$$

\item $\nabla:\mathcal{V} \to \Omega^1_B
\otimes \mathcal{V}$
the flat connection induced by $V_B$.
Griffiths transversality condition
$\nabla(F^p\mathcal{V})
\subset\Omega_B^1\otimes
F^{p-1}\mathcal{V}$.

\item Polarization given by a morphism of
local systems
$\psi:V_B \otimes V_B \to \mathbb{Z}_B$
such that for every $t \in B$, $\psi_t$
defines a polarization
on $(V_{B,t},F^p\mathcal{V}_t)$.
$\psi$ is $(-1)^k$-symmetric.
We have a representation
$\rho:\mathbb{S}(\mathbb{R}) \to
\text{GL}(V_{B,t}\otimes \mathbb{R})$,
$C=\rho(1)$,
$Q(x,y)=(2\pi i)^k \psi(Cx,y)$,
$Q>0$ on  $V_{B,t} \otimes \mathbb{R}$.

\item Hodge locus.
Assume weight = $2k$.
$x \in V_{B,t}$ is a Hodge class if and only
if
$x \in F^k\mathcal{V} \cap V_{B,t}$.

 $$
\xymatrix{
\text{Tot}(V_B)\ar@{^{(}->}[r]\ar@{=}[d]
&\text{Tot}(\mathcal{V})\ar[d]^{p}\\
\amalg W_i\ar[r]
&B
}
$$
where $W_i$ are the connected components and
$p|_{W_i}\to B$ is a covering.

Let $Z_i=W_i\cap \text{Tot}(F^k\mathcal{V})$,
which is a closed complex subvariety of
$\text{Tot}(\mathcal{V})$.
There are two possibilities:
\begin{enumerate}
\item Either $Z_i=W_i$ and $p(Z_i)=B$,
\item or $\dim Z_i=\dim W_i= \dim B \implies
p(Z_i)$
is nowhere dense. [This one corresponds to
the Hodge locus.]
\end{enumerate}

\noindent
$$
B^{Hdg}(V_B)
=\bigcup_{p(Z_i)\neq B}p(Z_i)
$$
Then the Hodge locus is
$$
B^{Hdg}=\bigcup_{e,f \geq
0}B^{Hdg}(T^{e,f}V_B
$$
where $T^{e,f}V_B=V_B^{\otimes e}\otimes
(V_B^* )^{\otimes f}$.

Note that $B \neq  B^{Hdg}$
because $B^{Hdg}$ is a countable union of
nowhere dense subsets.

\end{itemize}

\section{Hodge locus (continued)}
\label{mumford-tate-groups-in-hodge-theory-section-hodge-locus-cont}

\begin{proposition}
\label{mumford-tate-groups-in-hodge-theory-proposition-proper}
$\forall  i$, $p|_{Z_i}:Z_i \to B$
is a proper map.
Thus $p(Z_i)\subset B$ is a closed
subvariety.
\end{proposition}

\begin{proof}
Let $x \in Z_i$.
Define
$$
A=(2\pi i) \psi(x,x)
\underbrace{=}_{\substack{\text{$x$ is of} \\
\text{type $(k,k)$}}}
Q(x,x)\geq 0
$$
$A$ does not depend on the choice of $x \in
Z_i$
because $\psi$ is flat, i.e. $\nabla \psi=0$.

Let $S_A=\{z \in
\text{Tot}(\mathcal{V}^{k,k})
:Q(z,\bar{z})=A\}$.
Then $p|_{S_A}:S_A \to B$
is a sphere bundle $\implies $ proper map.
Since $Z_i$ is a closed subset of $S_A$,
we conclude that $p|_{Z_i}$ is also a proper
map.
\end{proof}

\begin{remark}
\label{mumford-tate-groups-in-hodge-theory-remark-quasi-projective}
When $B$ is quasi-projective
$p(Z_i)$ is algebraic
(Cattani-Deligne-Kaplan).
\end{remark}

\noindent
If $W_i=Z_i \implies W_i$ spans
a sub-local system on $V_B$,
that is in  $\mathcal{V}^{k,k}$.

\begin{definition}
\label{mumford-tate-groups-in-hodge-theory-definition-hodge-tate-local-system}
$V_B^{HT}$ is the biggest sub-local system in
$V_B$
that is Hodge-Tate, i.e. in
$\mathcal{V}^{k,k}$.
\end{definition}

\noindent
Let $t_0 \in B\setminus B^{Hdg}$.
Denote $MT_{t_0}$ the MT group of the
corresponding Hodge structure,
i.e.
$MT_{t_0}=MT(V_{B,t},F^\bullet\mathcal{V}
_{t_0})$.
We do parallel transport along the loops in
$B$ based in $t_0$,
and obtain an action of the fundamental
group:
the monodromy representation is
$\mu:\pi_1(B,t_0) \to \text{GL}(V_{B,t_0})$.

\begin{proposition}
\label{mumford-tate-groups-in-hodge-theory-proposition-monodromy}
$MT_{t_0}$ contains a finite index subgroup
of $\text{Im}(\mu)$.
\end{proposition}

\begin{proof}
Recall that $MT_{t_0}$ is the stabilizer of
all
$(0,0)$ Hodge classes in
all tensor powers, i.e. in $(T^\bullet
V_B)_{t_0}$.

There exists Hodge classes
$x_j \in (T^{e_j,f_j}V_B)_{t_0}\cap
T^{e_j,f_j}\mathcal{V})^{0,0}_{t_0}$
where $j=1,…,m$
such that $MT_{t_0}=\bigcap_{j=1}^m
Stab(x_j)$.

Note that $t_0 \not\in B^{Hdg}\implies
x_j\in(T^{e_j,f_j}V_B)^{HT}_{t_0}$,
which is a sublocal system.
Now consider the monodromy action:
$\pi_1(B,t_0)$ acts on
$(T^{e_j,f_j}V_B)^{HT}_{t_0}$.
By Proposition \ref{mumford-tate-groups-in-hodge-theory-proposition-proper},
all $x_j$'s have finite $\pi_1(B,t_0)$
orbits.
Thus there exists a finite index subgroup
$\pi_1(B,t_0)$ that stabilizes all $x_j$'s.
Thus, the image of this subgroup under the
monodromy representation lies in $MT_{t_0}$
(because it stabilizes all $x_j$'s).
\end{proof}

\noindent
\begin{definition}
\label{mumford-tate-groups-in-hodge-theory-definition-generic-mt-group}
For any $t \in B$, define
$$
G_t=\bigcap_{x \in (T^\bullet
V_B)_t^{HT}}Stab(x).
$$
Then $G_t$ define a local system of algebraic
groups
with fibre $G$ called the {\it generic MT
group}.
\end{definition}

For any $t \in B\setminus B^{Hdg}$,
$G_t=MT_t$
and for any $t \in B^{Hdg}$, $MT_t \subset
G_t$

\begin{example}
\label{mumford-tate-groups-in-hodge-theory-example-hypersurface}
Let $B$ be the space of smooth hypersurfaces
$X \subset \mathbb{C}P^n$
of degree $d$.
Then $B$ is a Zariski open subset of
$\mathbb{C}P^{\binom{n+d}{n}-1}$,
i.e. it is a quasi-projective variety.

Let $\pi:\mathcal{X} \to B$ be the universal
family.

Assume that  $n=2r$.
$V_B=R^{2r-1}\pi_\ast\mathbb{Z}$
is a $\mathbb{Z}$-VHS$/B$.

$\psi:V_B \otimes V_B \to \mathbb{Z}_B$
is skew-symmetric,
$\mu:\pi_1(B,t_0) \to Sp(V_{B,t_0},\psi)$.

\begin{theorem}[Kazhdan-Margulis]
\label{mumford-tate-groups-in-hodge-theory-theorem-kazhdan-margulis}
$\text{Im}(\mu)$ is Zariski dense in
$Sp(V_{B,t_0}\otimes \mathbb{Q},\psi)$.
\end{theorem}

And a corollary is that
\begin{lemma}
\label{mumford-tate-groups-in-hodge-theory-lemma-generic-hodge-is-sp}
The generic Hodge group of the family
$\mathcal{X}$
is $Sp(V_{B,t_0} \otimes \mathbb{Q},\psi)$.
\end{lemma}

\medskip\noindent
Consider $n=4$, $d=5$, i.e.
$X \subset \mathbb{C}P^5$ a quintic 3-fold.
This is a Calabi-Yau, i.e. $K_X =
\mathcal{O}_X$.
Thus $h^{3,0}=1,h^{2,1}=101,h^{1,2}=101$,
$h^{0,3}=1$.

Assume  $X=\mathcal{X}_t$ for $t \in B
\setminus B^{Hdg}$,
$Hdg(X)=Sp(H^3,\psi)$,
$\text{Lie}(Hdg(X))\cong S^2H^3(3)$
where $S^2$ is the symmetric square, and it
is twisted by 3.

In the Hodge structure $S^2H^3(3)$
we have Hodge classes of types
$(3,-3),(2,-2),(1,-1),(0,0),…,(-3,3)$.
Thus $H^3(X,\mathbb{Q}) \notin
(\mathbb{Q}-HS)^{ab}$.
\end{example}

\section{Period domains and period maps}
\label{mumford-tate-groups-in-hodge-theory-section-periods}

\noindent
[A variation of Hodge structures is just a
bundle.
The classifying space of these bundles is
called period.]

\medskip\noindent
Let $V_B$ be a local system over $B$.
Pass to the universal covering
$u:\tilde{B}\to B$.
Consider $u^* V_B=V \otimes
\mathbb{Z}_{\tilde{B}}$.
Let $V=\mathbb{Z}^2$,
$\mathcal{V}=V \otimes
\mathcal{O}_{\tilde{B}}$.
$F^p\mathcal{V} \subset \mathcal{V}$,
$t_0 \in \tilde{B}\setminus B^{Hdg}$,
$Hdg(V_{B,u(t_0)})=G$, a reductive group.
We also have $\rho_0:U(1) \to G(\mathbb{R})$.

\begin{definition}
\label{mumford-tate-groups-in-hodge-theory-definition-connected-component}
Let $\mathcal{D}$ be the connected component
of
$\Hom(U(1),G(\mathbb{R}))$ containing
$\rho_0$.
\end{definition}

\begin{lemma}
\label{mumford-tate-groups-in-hodge-theory-lemma-d-orbit}
$\mathcal{D} \cong
G(\mathbb{R})^0/Stab(\rho_0)$,
where $G(\mathbb{R})$ acts on
$\Hom(U(1),G(\mathbb{R}))$ by conjugation on
the image,
i.e. for $g \in G(\mathbb{R})$ we put
$(g\cdot \rho_0)(z)=g\rho_0(z)g^{-1}$.
\end{lemma}

\begin{proof}
$\text{Im}(\rho_0)$ is a compact abelian
subgroup
of $G(\mathbb{R})$, thus $\text{Im}(\rho_0)$
is contained in a
maximal compact torus in $G(\mathbb{R})$.

Fact: $G$-reductive implies that all maximal
compact
tori are conjugate by the action of
$G(\mathbb{R})^0$.

$\rho:U(1) \to G(\mathbb{R})$ such that
$[\rho] \in \mathcal{D}$,
then there exists $g \in G(\mathbb{R})$ such
that
$g\cdot\rho \in T_0$,
$[g,\rho]\in \mathcal{D}$,
$\Hom(U(1),T_0)$ is a discrete topological
space,
thus $g\rho=\rho_0 \iff \rho=g^{-1}\rho_0$.
\end{proof}

\medskip\noindent
Let $K=Stab(\rho_0)\subset G(\mathbb{R})^0$.

$C=\rho_0(i)$, $\sigma=Ad_C$
is a Cartan involution on $G(\mathbb{R})$.
$G^\sigma = \{g \in
G(\mathbb{C}):\bar{g}=(g)\}$
is compact.
We also have that $\sigma|_K=\text{id}
\implies K\subset G^\sigma
\implies K$ is compact.

Consider the Lie algebra of $G$, i.e.
$\mathfrak{g}=\text{Lie}(G)$.
Consider the Hodge decomposition induced by
$\rho_0$, namely
$\mathfrak{g}_\mathbb{C}=\bigoplus_{p \in
\mathbb{Z}}\mathfrak{g}^{p,-p}$.
We see that the action is trivial on the
component $(0,0)$:
$$
\begin{aligned}
z \in U(1)&\implies \bar{z}=z^{-1},\\
\rho_0(z)|_{\mathfrak{g}^{p,-p}}&=z^{2p}
\text{Id}\\
&\implies \mathfrak{g}^{0,0}\cap
\mathfrak{g}_\mathbb{R}=\text{Lie} K.
\end{aligned}
$$

\noindent
Note that $F^p\mathfrak{g}_\mathbb{C}
=\bigoplus_{r \geq p}\mathfrak{g}^{r,-2}$,
$F^0\mathfrak{g}_\mathbb{C}=$ Lie subalgebra
that preserves the Hodge
filtration on $V_{B,t_0}$.
$P \subset G(\mathbb{C})$ be such that
$\text{Lie}(P)=F^0\mathfrak{g}_\mathbb{C}$.

\begin{definition}
\label{mumford-tate-groups-in-hodge-theory-definition-dhat}
$\hat{\mathcal{D}}=G(\mathbb{C})/P$ where $P$
is the stabilizer
of the Hodge filtration on $V_{B,t_0}\otimes
\mathbb{C}$.
\end{definition}

\noindent
$\hat{\mathcal{D}}$ is a flag variety.

[We want to see how to embed $\mathcal{D}$ in
$\hat{\mathcal{D}}$.]
Consider $\text{Lie}(P) \cap
\text{Lie}(G(\mathbb{R}))
=F^0\mathfrak{g}_\mathbb{C} \cap
\mathfrak{g}_\mathbb{R}
=\text{Lie}(K)\implies $
Let $x_0=[1] \in \hat{\mathcal{D}}$
be the $G(\mathbb{R})^\bullet$-orbit of $x_0$
is
$\mathcal{D}=G(\mathbb{R})^\bullet/K
\implies \mathcal{D} \subset
\hat{\mathcal{D}}$.

Note that
$$
\begin{aligned}
\dim_\mathbb{R}(\mathcal{D})
&=\dim_\mathbb{R} \mathfrak{g}_\mathbb{R}
-\dim_\mathbb{R}(\mathfrak{g}^{0,0}\cap
\mathfrak{g}_\mathbb{R})\\
&=\dim_\mathbb{C}\mathfrak{g}_\mathbb{C}-
\dim\mathbb{C}\mathfrak{g}^{0,0}\\
&=2\cdot
\dim_\mathbb{C}\left(\bigoplus_{r<0}
\mathfrak{g}^{r,-r}\right)\\
&=2\cdot \dim_\mathbb{C}\hat{\mathcal{D}}\\
&=\dim_\mathbb{R}\hat{\mathcal{D}}.
\end{aligned}
$$

\noindent
Thus we have an open embedding $\mathcal{D}
\subset \hat{\mathcal{D}}$.

\medskip\noindent
Over $\tilde{B}$ we have a flat bundle
$V\otimes \mathcal{O}_{\tilde{B}}$
and a filtration $F^p \mathcal{V} \subset
\mathcal{V}$
by holomorphic subbundles.
Thus we get a period map
$p:\tilde{B}\to \mathcal{D} \subset
\hat{\mathcal{D}}$.

Griffiths transversality is [means] the
following condition
$$
dp:T\tilde{B}\to p^*\mathcal{H}
$$
where $\mathcal{H}$ is given as follows.
$F^{-1}\mathfrak{g}_\mathbb{C}/
F^0\mathfrak{g}_\mathbb{C}$
defines a homogeneous
subbundle of $T\hat{\mathcal{D}}$ which we
call
the {\it horizontal subbundle} $\mathcal{H}$.

\medskip\noindent
If
\begin{equation}
\label{mumford-tate-groups-in-hodge-theory-equation-hermitian}
\mathfrak{g}_\mathbb{C}=\mathfrak{g}^{-1,1}
\oplus
\mathfrak{g}^{0,0}
\oplus \mathfrak{g}^{1,-1}
\end{equation}

\noindent
then
$F^{-1}\mathfrak{g}_\mathbb{C}=\mathfrak{g}
_\mathbb{C}$
which implies that
$\mathcal{H}=T\hat{\mathcal{D}}$.
Thus Griffiths transversality holds
automatically
for any holomorphic map $\tilde{B} \to
\mathcal{D}$.

\begin{proposition}
\label{mumford-tate-groups-in-hodge-theory-proposition-hermitian}
If Equation \ref{mumford-tate-groups-in-hodge-theory-equation-hermitian} holds,
then $\mathcal{D}$ is a Hermitian symmetric
space.
\end{proposition}

\begin{proof}
The $Ad_C$ [Weil operator] acts on
$\mathfrak{g}^{1,-1}\oplus
\mathfrak{g}^{-1,1}$ as $-\text{Id}$.
Thus $Ad_C$ gives a holomorphic isometry
fixing $x_0\in \mathcal{D}$
and acting as $-\text{Id}$ on
$T_{x_0}\mathcal{D}$.
\end{proof}


\section{The Kuga-Satake construction}
\label{mumford-tate-groups-in-hodge-theory-section-kuga-satake}

\noindent
Let $X$ be a compact hyperkähler manifold
(aka IHS manifold).
Recall that this means that $X$ is a
compact smooth manifold, simply connected
equipped with a Riemannian metric $g$ such
that
$\text{Hol}(\nabla^g)=\text{Sp}(n)$
where $\nabla^g$ is the Levi-Civita
connection
and $\text{Sp}(n)$ is the compact form of
$\text{Sp}(2n,\mathbb{C})$,
i.e.
$\text{Aut}(\mathbb{H}^n,\left\langle{}\cdot,
\cdot\right\rangle{})$,
the quaternions equipped with the canonical
quaternionic Hermitian product
$\left\langle{}v,w\right\rangle{}:=\sum_{i=1}
^nv_i\overline{w_i}$.

The holonomy principle implies
that there exist complex structures $I,J$ and
$K$
such that $I J=J I=K$.

Let $X_I$ be the complex manifold $X$
equipped with $I$.
Again by the holonomy princple, there is a
holomorphically symplectic form
$\sigma_I \in H^0(X,\Omega_{X_I}^2)$.
In fact, this is unique up to scaling and
$H^0(X,\Omega_{X_I}^2)\cong
\mathbb{C}\sigma_I$.
Also, $\sigma_I^n$ is a holomorphic volume
form,
so that $K_{X_I}\cong\mathcal{O}_{X_I}$,
so that these are Calabi-Yau manifolds.
$\dim_{\mathbb{C}}X_I=2n$.

\begin{example}[IHS manifolds]
\label{mumford-tate-groups-in-hodge-theory-example-hk}
\begin{enumerate}
\item $S$ complex K3 surfaces. E.g. quartics
in $\mathbb{C}P^3$
[by CY theorem there is a hyperkähler metric;
it's difficult to construct such a metric
explicitly].

\item Let $S$ be a complex K3 surface.
Then $\text{Hilb}^n(S)$, the Hilbert scheme
of length-$n$ 0-dimensional subschemes of
$S$.

\item Let $A =\mathbb{C}^2/\Lambda$.
Consider the Albanese map
$\text{Hilb}^{n+1}(A)\to A$.
Then the fiber is called a {\it Generalized
Kummer variety}
$K^nA:=Alb^{-1}(0)$.

\item O'Grady: exceptional examples of
dimension 6, 10.
\end{enumerate}
\end{example}

\section{Cohomology of an IHS manifold}
\label{mumford-tate-groups-in-hodge-theory-section-cohomology-ihs}

\noindent
$V:=H^2(X,\mathbb{Q})$ is a K3-type Hodge
structure.
Since
$H^0(X,\Omega_{X_T}^2)=\mathbb{C}\sigma$
where $\sigma$ is the holomorphic symplectic
form,
we must have $h^{2,0}=h^{0,2}=1$.

\begin{theorem}[Beauville-Bogomolov-Fujiki]
\label{mumford-tate-groups-in-hodge-theory-theorem-bbf}
There exists a constant $c_X \in \mathbb{Q}$
and a quadratic form
$q \in S^2V^*$ of signature $(3,b_2(X)-3)$
such that $\forall  a \in V$
$$
\int_X a^{2n}=c_X q(a)^n \qquad \text{(Fujiki
relation)}
$$
where $\int_X$ is the pairing with the
fundamental class
defined by $I$.
\end{theorem}

[We have the integral of a polynomial on $V$,
and it happens to be the $n$-th power of a
quadratic form $q$.
In fact, the integral does not depend on the
choice of complex
structure.]


We call $q$ the  {\it BBF form}.
The HS on $H^2$ is given by
$$
\rho:U(1) \to \text{SO}(V\otimes
\mathbb{R},q).
$$

If $\omega$ is a Kähler form on $X$,
then $q([\omega])>0$.
If the class of  $\omega$ is rational,
i.e. $[\omega] \in H^2(X,\mathbb{Q})$,
then
$H^2(X,\mathbb{Q})_{\text{prim}}=[\omega]
^\perp$
is polarized by $q$.

Since $\omega$ is Kähler, we have the
Lefschetz triple
$$
\begin{aligned}
L_\omega&=[\omega] \cup  -,\qquad
\Lambda_\omega=\text{dual of }L_\omega,
\qquad [L_\omega,\Lambda_\omega]=\Theta,
\end{aligned}
$$

\noindent
[which give us a representation of
$\mathfrak{sl}_2$].

Consider $L_\omega, \Lambda_\omega$ for all
Kähler forms
on all deformations of  $(X,I)$.
The smallest Lie subalgebra of
$\text{End}(H^\bullet(X,\mathbb{R}))$
that contains all $L_\omega,\Lambda_\omega$?

\section{Mukai extension}
\label{mumford-tate-groups-in-hodge-theory-section-mukai-extension}

\noindent
Let $(V,q)$, $\tilde{V}=V \oplus U$,
where $U$ is the hyperbolic plane
$q_n=\begin{pmatrix}0&1\\ 1&0\end{pmatrix}$.
Let $\tilde{q}=q \oplus q_n$.
We have a grading
$$
\tilde{V}=\underbrace{\mathbb{Q}e_0}_{0}
\oplus
\underbrace{V}_{2}
\oplus\underbrace{\mathbb{Q}e_4}_{4}.
$$

Then
$$
\begin{aligned}
\mathfrak{so}(\tilde{V},\tilde{q})
&\simeq \Lambda^2\tilde{V}\\
&=\underbrace{e_0 \wedge V}_{\substack{-2 \\
\Lambda_\omega \in }}
\oplus\underbrace{(\underbrace{\Lambda^2V}
_{\mathfrak{so}(V,q)}
\oplus e_0 \wedge e_n)}_{0}
\oplus \underbrace{e_n\wedge V}_{\substack{2
\\ L_\omega \in}}
\end{aligned}
$$

\noindent
is in fact the minimal Lie subalgebra
containing
the Lefschetz operators:

\begin{theorem}[Verbitsky, Looijenga-Lunts]
\label{mumford-tate-groups-in-hodge-theory-theorem-vll}
The smallest subalegbra of
$\text{End}(H^\bullet(X,\mathbb{R}))$
containing $L_\omega, \Lambda_\omega$
is isomorphic to
$\mathfrak{so}(\tilde{V},\tilde{q})$.
\end{theorem}

\noindent
This implies that we have an action
of $\mathfrak{so}(V,q)$ on
$H^k(X,\mathbb{Q})$ for all $k$.
Thus, $H^k(X,\mathbb{Q})$ is a representation
of
$\text{Spin}(X,q)$ for $k \equiv 0
(\text{mod})$ of $\text{SO}(V,q)$.

The HS of $H^k(X,\mathbb{Q})$ are induced by
$\rho:U(1) \to \text{Spin}(V\otimes
\mathbb{R},q)$.

[The idea is to understand the cohomology
as a representation of
$\mathfrak{so}(\tilde{V},\tilde{q})$.]

Let $(\mathbb{Q}-HS)^{\text{tori}}$ be the
full
$\otimes$-subcategory of $(\mathbb{Q}-HS)$
generated
by $H^1(T,\mathbb{Q})$ for tori
$T=\mathbb{C}/\Lambda$.
We have
$$
(\mathbb{Q}-HS)^{\text{tori}}\supset
(\mathbb{Q}-HS)^{ab}
\subset (\mathbb{Q}-HS)^{pol}.
$$

\begin{theorem}[Kurnosov-Soldatenkov-
Verbitsky]
\label{mumford-tate-groups-in-hodge-theory-theorem-ksv}
Let $X$ be an ISH manifold.
Then for all $k$,
$H^k(X,\mathbb{Q}) \in (Q-HS)^{\text{tori}}$.
If $X$ is projective then
$$
H^k(X,\mathbb{Q}) \in (\mathbb{Q}-HS)^{ab}.
$$
\end{theorem}

\section{Idea of the proof}
\label{mumford-tate-groups-in-hodge-theory-section-proof}

\noindent
[We need to understand cohomology groups
as representations of $Spin(V,q)$.]

\begin{enumerate}
\item (Clifford algebras.)
Let $(V,q)$ be a quadratic vector space over
$\mathbb{Q}$.
The Clifford algebra $Cl(V,q)=\bigoplus_{i
\geq 0}V^{\otimes i}/J$
where $J$ is the ideal generated by
$v \otimes v-q(v)\cdot 1$ for all $v \in V$.

In $Cl(V,q)$, $v^2=q(v)\cdot 1$,
$uv-vu=2q(u,v)\cdot 1$
Also,
$$
\alpha(v_1 \otimes … \otimes v_k)
=(-1)^k v_1 \otimes … \otimes v_k,\qquad
\beta(v_1 \otimes … \otimes v_k)
=v_k \otimes v_{k-1} \otimes … \otimes v_1
$$
descend to $Cl(V,q)$.

Foro all  $a \in cl(V,q)$ write
$\bar{a}=\alpha\beta(a)$,
$Cl(V,q)=Cl^0(V,q) \oplus Cl^1(V,q)$
is a $\mathbb{Z}/2$-grading.

Denote $Cl(V,q)=\mathcal{C}$ and
$\mathcal{C}^\times$
the invertible elements of $\mathcal{C}$.
For all $x \in V$ such that $q(x) \neq 0$,
$x \in \mathcal{C}^\times$ since
$x^2=q(x)\cdot 1 \implies
x^{-1}=\frac{x}{q(x)}$.

We have a natural embedding $V
\xymatrix{\ar@{^{(}->}[r]&} \mathcal{C}$.
Let $\{ x \in
\mathcal{C}^\times:\alpha(x)Vx^{-1}\subset
V\}:=G$.

If $x \in V$, $q(x) \neq 0$, $y \in V$,
$$
\begin{aligned}
\alpha(x) y x^{-1}&=-xy\frac{x}{q(x)}\\
&=(yx-2q(x,y)\frac{x}{q(x)}\\
&=y-2\frac{q(x,y)}{q(x)}x,
\end{aligned}
$$

\noindent
i.e. a reflection.

$x \in G$, $y \in V$,
$$
\begin{aligned}
q(\alpha(x)yx^{-1})&=\underbrace{\alpha(x)
yx^{-1}}_{V}\alpha(x)yx^{-1}\\
&=-\alpha(\alpha(x)yx^{-1})\alpha(x)yx^{-1}\\
&=xy\alpha(x)^{-1}\alpha(x)yx^{-1}\\
&=q(y).
\end{aligned}
$$

\noindent
We get a representation $\mu:G \to O(V,q)$.
$\mu$ is surjective (because $\text{Im}(\mu)$
contains reflections
[and there's a theorem that $O(V,q)$ is
generated by reflections]).

$\Ker(\mu)=\mathbb{Q}^*\cdot 1$.
[The kernel is quite big so let's make it a
bit smaller.]
Define a norm $N:G \to \mathbb{Q}^*$ by
$N(a)=a\cdot \bar{a}$.
Define $\Ker(N):=Pin(V,q)$.
Then $Pin(V,q) \cap
\mathcal{C}^0=\text{Spin}(V,q)$.
Then $\mu:\text{Spin}(V,q)\to
\text{SO}(V,q)$.
Thus, $\Ker \mu=\{\pm 1\}$,
[the double cover we were looking for].

\item (Key observation.)
The morphism $\rho:U(1) \to
\text{SO}(V\otimes \mathbb{R},q)$
defining the HS on $H^2(X,\mathbb{Q})$
lifts to $\text{Spin}(V\otimes
\mathbb{R},q)$.

To see why, write
$$
V \otimes \mathbb{R}=
\underbrace{(V^{2,0}\oplus
V^{0,2})}_{V^+_{\mathbb{R}}}
\oplus V^{1,1}_\mathbb{R}.
$$
Then $\dim V^+_\mathbb{R}=2$.

We have a double cover
$$
\xymatrix{
\rho:U(1)\ar[rd]_{2:1}\ar[r]&
\text{SO}(V_{\mathbb{R}}^+)\ar@{=}[d]\\
& U(1).
}
$$

We have
$$
Cl(V \otimes \mathbb{R},q)
= \underbrace{ Cl(V_\mathbb{R}^+)}_{4-\dim}
\otimes Cl(V^{1,1}_\mathbb{R}).
$$
Then $V_\mathbb{R}^+=\left\langle{}e,f\right\rangle{}$ --- an
orthonormal basis,
$Cl^0(V_\mathbb{R}^+)=\left\langle{}1,ef\right\rangle{}\simeq
\mathbb{R}^2$.

$$
\xymatrix{
\text{Spin}(V_\mathbb{R}^+)\cong U(1)
\ar[d]^{2:1}\\
\text{SO}(V_\mathbb{R}^+).
}
$$

\item (Definition of the Hodge structure.)
[What's the use of lifting?]
We have a lift $\tilde{\rho}:U(1) \to
\text{Spin}(V\otimes \mathbb{R},q)$.
Note that $Cl(V_\mathbb{R}^+)$ is two copies
of the standard
representation of $U(1)$.

$$
Cl(V \otimes \mathbb{R},q)
\simeq
\bigoplus \text{standard reps of }U(1) \simeq
\text{Spin}(V_\mathbb{R}^+).
$$
Denote $H:=Cl(V,q)$,
$\implies H \otimes \mathbb{C}=H^{1,0} \oplus
H^{0,1}$.
Then $H \in (\mathbb{Q}-HS)^{\text{tori}}$.

\item (Constructing the embedding.)
We want to construct $V
\xymatrix{\ar@{^{(}->}[r]&} H^*$.
[We should use some bilinear form on the
Clifford algebra.]
In fact, there exists a non-degenerate
pairing
$\tau:\mathcal{C} \times \mathcal{C} \to
\mathbb{Q}$,
$\tau(x,y)=\text{Tr}(x\bar{y})$,
where we consider $x\bar{y}$ as an
endomorphism of $\mathcal{C}$
acting by left multiplication and
Tr is the trace of this endomorphism.
In fact, $\tau(x,y)=\tau(y,x)$,
$\tau(ax,y)=\tau(x,\bar{a}y)$ and
$\tau$ is Pin-invariant, i.e.
$\forall g \in \operatorname{Pin}(V,q)$,
$\tau(ax,ay)=\tau(x,y)$.

Analogously, on $H=\mathcal{C}^{\oplus m}$
we have $\tau:H \otimes H \to \mathbb{Q}$.

To construct the embedding, for any $x \in V$
let $\omega_x(u,v)=\tau(xu,v)$.
Then $\omega_x(u,v)=-\omega_x(v,u)$,
$\omega_x \in \Lambda^2H^*$.
This defines $\varphi:V \to \Lambda^2H^*$,
$x \mapsto \omega_x$, the embedding we
wanted.

$\varphi$ is a morphism of
$\text{Spin}(V,q)$-modules
[representations].
Thus, $\varphi$ is a morphism of Hodge
structures.

[This finishes the work for K3 surfaces.
For general HK manifolds we need to work a
bit more,
how to embed all cohomology groups?
In fact, this is where we will use the $m$
in $H=\mathcal{C}^{\oplus m}$.]

\item (Extra work needed for general HK
manifolds.)
Let $\dim H=4N$ and $\text{Im}(\varphi)=W$.
For $p \in S^{2N}H^*$,
$p(\omega)=\omega^{2N}\in
\Lambda^{2N}(H^*)\simeq\mathbb{Q}$.

[$p$ is a polynomial, when we restrict it:]
$p|_W \in (S^{2N}W^* )^{\text{Spin}(V,q)}$.
The classical invariant theory implies [that
this polynomial
is a power of the quadratic form, i.e.]
$p = A \cdot q^N, A \in \mathbb{Q}$.
So $p(\omega)=\omega^{2N}=A q(\omega)^N$,
a ``Fujiki relation''.

$$
\begin{aligned}
\mathfrak{so}(\tilde{V},\tilde{q})&=
\underbrace{V}_{\substack{2 \\ \Lambda_x\in}}
\oplus
\underbrace{\mathfrak{so}(\tilde{V},\tilde{q}
)}_{0}
\oplus \underbrace{V}_{\substack{2 \\
L_x\in}}
\end{aligned}
$$

\noindent
Let $\eta \in \Lambda^\bullet H^*$, $x \in
V$,
\begin{equation}
\label{mumford-tate-groups-in-hodge-theory-equation-rep}
L_x \eta=\omega_x \wedge \eta,\qquad
\Lambda_x\eta=i_\eta \omega_x
\end{equation}
[Where $i$ denotes interior multiplication].
$$
\begin{aligned}
\varphi: V &\xymatrix{\ar@{^{(}->}[r]&}
\Lambda^2H^*\\
x &\longmapsto \omega_x.
\end{aligned}
$$

\begin{proposition}
\label{mumford-tate-groups-in-hodge-theory-proposition-rep}
Equation \ref{mumford-tate-groups-in-hodge-theory-equation-rep} define a $
\mathfrak{so}(\tilde{V},\tilde{q})$-
representation
on $\Lambda^\bullet H^*$.
The corresponding
$\text{Spin}(\tilde{V},\tilde{q})$-
representation
is faithful.
\end{proposition}

In conclusion,
for all $\mathfrak{so}(\tilde{V},\tilde{q})$
representation $U$
exists $m>0$ such that for
$H=\mathcal{C}^{\oplus m}$,
$U \subset \Lambda^\bullet H^*$
as a $\mathfrak{so}(\tilde{V},\tilde{q})$
subrepresentation.
$U$ is contained in a direct sum of
 $(\Lambda^\bullet \mathcal{C}^* )^{\otimes
 i}
\simeq \Lambda^\bullet({\mathcal{C}^*
}^{\oplus i})$.

Take $U = H^2(X,\mathbb{Q})$,
qed an embedding of
$\text{Spin}(\tilde{V},\tilde{q})$
representations,
$H^\bullet(X,\mathbb{Q})
\xymatrix{\ar@{^{(}->}[r]&} \Lambda^\bullet
H^*$
is an embedding of HS.
\end{enumerate}

[Thus the HS on the cohomology of $X$ are
abelian.]








\section{References}
\label{mumford-tate-groups-in-hodge-theory-section-references}

\begin{enumerate}
\item Deligne, Milne, Ogus, Shih,
``Hodge cycles, motives and Shimura
varieties''.
We have seen mostly Chapter 1 and a bit of
Chapter 2.

\item Deligne:
\begin{itemize}
\item
``La conj. de Weil pour les surfaces K3''

\item ``Théorie de Hodge II''

\item ``Variétes de Shimura: …''
\end{itemize}

\item Kingler: ``Hodge theory, between
algebraicity and transcendence
+ reference therein''.
\end{enumerate}

% --- End mumford-tate-groups-in-hodge-theory.tex ---

% --- Begin seminars-complex-geometry.tex ---
\chapter{Seminars in Complex Geometry}
\label{chapter-seminars-complex-geometry}















\section{Everything you always wanted to know
about polygons but were too
afraid to ask}
\label{seminars-complex-geometry-section-every-you-awlwa-wante-know}

\noindent
Alessia Mandini, UFF.
GAAG, IMPA.
September 22 and 23, 2025.

\medskip\noindent
{\bf Abstract.} Moduli spaces of polygons
form a family of Kähler
manifolds that can be constructed as a Kähler
reduction of coadjoint orbits.
These spaces have deep connections to various
areas of mathematics, including
symplectic and algebraic geometry as well as
representation theory. In this
talk, I will define these spaces and explore
their connections. I will then
discuss wall-crossing phenomena in these
spaces and demonstrate how it can be
used to determine their cohomology rings.
Finally, I will introduce the
hyperkähler analogue of these spaces, known
as hyperpolygon spaces, and describe
some of their generalizations.

\medskip\noindent



Plan of the talk
\begin{enumerate}
\item Polygons in $\mathbb{R}^3$ and
relations with
other moduli spaces.
\item Polygons in other spaces.
\end{enumerate}

Let $\alpha=(\alpha_1,\ldots,\alpha_n)\in
\mathbb{R}^n_{>0}$.
Consider $S^2$ with its usual symplectic,
Kähler form.
Also consider the product of several spheres.
There's a Hamiltonian action of
$\text{SO}(3)$ by rotations.
The moduli space of polygons is
the symplectic reduction obtained with this
Hamilatonian action,
$$
M(\alpha)=\prod
S^2_{\alpha_i}/\!/\text{SO}(3).
$$

Why is it called the space of polygons?
We think that
$$
[v_1,\ldots,v_n] \in M(\alpha) \iff
\sum_{i=1}^nv_i=0
$$
are polygons.

When smooth, $M(\alpha)$ is a
$(n-3)$-complex-dimensional
Kähler manifold. Consider
$$
\varepsilon_I(\alpha)=\sum_{i \in I}\alpha_i
-\sum_{j \in I^c}\alpha_j
$$
for some index set $I \subseteq \{1,\ldots,
n\}$.
$M(\alpha)$ is {\it smooth} if
$\varepsilon_I(\alpha)
\neq 0$ for all $I\subseteq \{1,\ldots,n\}$.
If so, we say $\alpha$ is {\it generic}.

\begin{theorem}[Kapovich-Millson]
\label{seminars-complex-geometry-theorem-Kapovich-Millson}
$M(\alpha)$ is a complex analytic space
with (eventually) isolated singularity
(homogeneous quadratic cones).
\end{theorem}

\begin{example}
\label{seminars-complex-geometry-examples-polygon-moduli-spaces}
\begin{enumerate}
\item $(n=3.)$ Then $M(\alpha)$ is either
empty or a point.

\item $(n=4.)$ $M(\alpha)$ is either empty or
a sphere.
\end{enumerate}
\end{example}

\begin{remark}[Hausmann-Knutson]
\label{seminars-complex-geometry-remark-Hausmann-Knutson}
Let $M_n$ be the space of all $n$-gons
modulo rigid motions. Then $M_n$
can be equipped with a Poisson structure
for which $M(\alpha)$ are the symplectic
leaves.

\medskip\noindent
{\bf Polygons as quiver varieties.}
Consider the star-shaped quiver,
which is a distinguished point with some
points around it,
and arrows from every point to the
distinguished one.
Let the distinguished point be
$V_0=\mathbb{C}^2$ and
the rest $\mathbb{C}$. Then a representation
of this quiver
is
$$
\text{Rep}Q=\bigoplus_{i=1}^n
\Hom(\mathbb{C},\mathbb{C}^2)=\mathbb{C}^{2n}
.
$$
Now put
$$
K=(U(2) \times U(1)^n)/\Delta
$$
where $\Delta$ is the diagonal $S^1$.
This gives a Hamiltonian action on
$\mathbb{C}^{2n}$
as follows. For
$$
(A,\lambda_1,\ldots,\lambda_n)\cdot
(q_1,\ldots,q_n)
$$
we map
$q_i \mapsto  A^{-1}q_i \lambda_i$.
That is
$$
\begin{aligned}
\mu: \mathbb{C}^{2n} &\longrightarrow K^* \\
(q_1,\ldots,q_n) &\longmapsto
\left(\sum_{n=1}^n (q_iq_i^*
),\ldots,\frac{1}{2}|q_i|^2\ldots\right)
\end{aligned}
$$

\begin{theorem}[Hausmann-Knutson]
\label{seminars-complex-geometry-theorem-Hausmann-Knutson}
$$
\mathbb{C}^{2n}/\!/_{(0,\alpha)}K=M(\alpha)
$$
\end{theorem}

\begin{proof}
Uses
$$
\begin{aligned}
\mathbb{R}^3  &\xrightarrow{\simeq}
\mathfrak{su}(2)^* \\
v_i &\longmapsto (q_iq_i^*)_0
\end{aligned}
$$
\end{proof}
\end{remark}

$$
\xymatrix{
&\mathbb{C}^{2n}\ar[dl]_{/\!/_\alpha U(1)^n}
\ar[dr]^{/\!/U(2)}\\
\prod S^2 \ar[dr]_{/\!/_\alpha \text{SO}(3)}
 &  &
 \text{Gr}(2,\mathbb{C}^n)\ar[dl]^{/\!/
 _\alpha
 U(1)^n}\\
&M(\alpha).
}
$$
So we have (I think former symplectic
reduction)
$$
\mu_{U(1)^n}:\text{Gr}(2,\mathbb{C}^n)\to
\mathbb{R}^n
$$
\medskip\noindent
{\bf Walls and wall crossing.}
The walls are
$$
W_I=\{\alpha \in
\mathbb{R}^n_{>0}:\varepsilon_I(\alpha)=0\}
$$
for $I \subseteq \{1,\ldots,n\}$.

To understand wall crossing suppose we have
a wall $W_I$, with $a^c$ in the wall,
$\alpha^+$ on one side and $\alpha^-$ on the
other.

Not that
$$
\varepsilon_I(\alpha^+)>0 \iff
\sum_{k \in I}\alpha_i^+>
\sum_{ j \in I^c}\alpha_j^+
$$
Define
$$
\begin{aligned}
M_{I^c}(\alpha^+)
&=\{(v_1,\ldots,v_n) \in M(\alpha^+):
v_i=\lambda v_j \forall i,j \in I^c,
\lambda>0\}\\
&=M(\tilde{\alpha}),\qquad \tilde{\alpha}=
\left(\alpha_{i_1},\ldots,\alpha_{i_k},
\sum_{j
\in I^c}\alpha_j\right)
\end{aligned}
$$
Then $\varepsilon_I(\alpha^+)<0$
implies $M_I(\alpha^-) \subseteq
M(\alpha^-)$.

$$
\xymatrix{
&\tilde{M} \subseteq E=M_I \times
M_{I^c}\ar[dd]^{\text{blow up}}\\
M(\alpha^+)\ar[dr]&  &  M(\alpha^-)\ar[dl]\\
&M(\alpha^c)
}
$$
\begin{remark}
\label{seminars-complex-geometry-remark-modul-space-polyg-propo}
\begin{itemize}
\item $M(\alpha)\cong M(\sigma(\alpha))$
for $\sigma$ a permutation on the order of
sets.
\item $M(\alpha)$ is conformally
symplectomorphic
to $M(\lambda \alpha)$ for all $\lambda>0$.
\end{itemize}
\end{remark}

\begin{definition}
\label{seminars-complex-geometry-definition-short-and-long}
Let $\alpha \in \mathbb{R}^n_{>0}$.
$I \subseteq \{1,\ldots,n\}$ is {\it short}
if $\varepsilon_I(\alpha)<0$ and {\it long}
if  $\varepsilon_I(\alpha)>0$.
\end{definition}

\begin{example}
\label{seminars-complex-geometry-example-wall-crossing}
We did an example of wall crossing.
The relevant manifolds $M_{I^c}(\alpha)$ and
$M_I(\alpha)$
were projective spaces.
\end{example}

\medskip\noindent
Now recall our first quotient
$$
\xymatrix{
\mu_{U(1)}^{-1}(\alpha)\subseteq
\text{Gr}(2,\mathbb{C}^n)
\ar[d]^{/\!/U(1)^n}\\
M(\alpha)
}
$$
Let $c_i$ be the first Chern
classes associated to the
$n$ $S^1$-bundles above
(by taking reduction in stages).

\begin{theorem}[Haussmann-Knutson,M.]
\label{seminars-complex-geometry-theorem-Haussmann-Knutson-M}
The $c_i$ generate the cohomology of
$M(\alpha)$.
\end{theorem}

\begin{theorem}[Guillemin-Stenberg]
\label{seminars-complex-geometry-theorem-Guillemin-Stenberg}
In this setup, for $\alpha$ generic,
$$
H(M)=\mathbb{C}[c_1,\ldots,c_n]/\text{Ann}
(\text{Vol}(M(\alpha)))
$$
i.e., $Q(c_1,\ldots,c_n) \in
\text{Ann}(\text{Vol}M(\alpha))$
$\iff$ $Q\left(\frac{\partial }{\partial
\alpha_i},\ldots,
\frac{\partial }{\partial \alpha_1}\right)
\text{Vol}(M(\alpha))=0$
\end{theorem}

\begin{theorem}[Takakura, The Koi]
\label{seminars-complex-geometry-theorem-Takakura-The-Koi}
$$
\text{Vol}(M(\alpha))
=-\frac{(2\pi)^{n-3}}{(n-3)!}
\sum_{I\text{
long}}(-1)^{n-|I|}\varepsilon_I(\alpha)^{n-3}
.
$$
\end{theorem}

\begin{example}
\label{seminars-complex-geometry-example-Delta1}
According to Example
\ref{seminars-complex-geometry-example-wall-crossing} (which I
did not copy) we find that
$\alpha_1 \in \Delta_1$ gives
$\text{Vol}(M(\alpha_1))=2\pi^3(1-2\alpha_3)
^2$
and
$$
H(M(\alpha_1))=\mathbb{C}[c_3]/(c_3^3)
$$
\end{example}

\medskip\noindent
{\bf Polygon game.}
Let $G$ be a Lie group and $\mathfrak{g}$ its
Lie algebra.
Then the co-adjoint orbits
$\mathfrak{g}^*  \supseteq
\mathcal{O}_{\xi_i}$
are symplectic manifolds with the KKS form.
Then
$$
\begin{aligned}
G\text{ acts on }\Pi \mathcal{O}_{\xi_i}
&\longrightarrow \mathfrak{g}^* \\
(A_i,\ldots\alpha_n) &\longmapsto \sum A_i
\end{aligned}
$$
Then the quotient
$$
M(\xi)=\Pi \mathcal{O}_{\xi_i}/\!/_0 G
$$
generalizes the previous construction,
which we get with $G=\text{SU}(2)$.

It could be interesting to investigate
which of the next constructions can be
generalized
to other Lie groups.

\begin{theorem}[Sotillo-Floentins-Gadihl]
\label{seminars-complex-geometry-theorem-Sotillo}
Wall-crossing for $\text{SU}(m)$.
\end{theorem}

\medskip\noindent
{\bf Bending action.}
We consider a polygon of $n$ sides
and put some diagonals that don't intersect.
We introduce some
notion of ``bending'' that allows
to define a Hamiltonian function.
In turn, this defines a torus action on
$M(\alpha)$
and a moment map.

For $n=2$ we obtain the moment map
$$
\begin{aligned}
\mu: M(\alpha) &\longrightarrow \mathbb{R}^2
\\
p &\longmapsto (\ell_1(p),\ell(p))
\end{aligned}
$$
and we can see the moment polytope in
$\mathbb{R}^2$.

\begin{remark}
\label{seminars-complex-geometry-remark-nonva-nonin-diago}
Any system of $(n-3)$ non-vanishing
and non-intersecting diagonals determine
a torus action of $T^{n-3}$ on $M(\alpha)$.
\end{remark}

\medskip\noindent
{\bf Relations of polygon spaces to other
moduli spaces.}
Representations of the fundamental group
of the punctured sphere in $SU(2)$, i.e.
$$
\text{Rep}(\pi_1(S^2\setminus
\{p_1,\ldots,p_n\},\text{SU}(2))/\text{SU}(2)
\cong
M(\alpha)
$$
This can be put very explicitly:
$$
\begin{aligned}
&\text{Rep}(\pi_1(S^2\setminus
\{p_1,\ldots,p_n\},\text{SU}(2))\\
&=\{(g_1,\ldots,g_n)\in
\text{SU}(2)^n:g_1\cdot\ldots\cdot
g_n=\text{Id},
t_rg_i=2\cos\pi \alpha_i\}
\end{aligned}
$$
See [Agapito-Godinho]. That's all we
will say about this example.

\medskip\noindent
Now consider the moduli space of parabolic
bundles.
Consider a holomorphic bundle of rank 2 over
$\mathbb{C}P^{2}$.
Choose some points $D=\{x_1,\ldots,x_n\}$ in
$\mathbb{C}P^{2}$
and a flag
$E_{x_i}=E^{x_i,1}\supseteq
E^{x_i,2}\supseteq \{0\}$
for every $x_i \in D$. Each of the
$E_{x_i,j}$ is isomorphic to $\mathbb{C}$
(1-dimensional).
A {\it quasi-parabolic bundle}
is an holomorphic bundle $E$
with such a choice of flag.
A {\it parabolic bundle} is a quasi-parabolic
bundle with a choice of parabolic weights
$0< \beta_1(x_i)< \beta_2(x_i)<1$.

There is also a notion of stability:
$$
\begin{aligned}
\text{pdeg}E&=\text{deg}E+ \sum_{i=1}^n
(\beta_1(x_i)+\beta_2(x_i))\\
\mu(E)&=\frac{\text{pdeg}(E)}{\text{rank}(E)}
\end{aligned}
$$
We say $E$ is {\it (semi)stable} if
$µ(E)> \mu(L)$ ($\mu(E) \geq \mu(L)$)
for any $L \subseteq E$ parabolic
subbundle.

Then we obtain that
$\mathcal{M}_{\pi,d}(\beta)$
is the moduli space of (semi)-stable
parabolic bundles on $\mathbb{P}^1$
of rank $r$ and degree $d$.
In particular $\mathcal{M}_{2,0}(\beta)$
is the moduli space of (semi)-stable
parabolic
bundles on $\mathbb{P}^1$ of rank 2 and
degree 0
holomorphically trivial.

\begin{theorem}[Jeffrey, Godinho, M.]
\label{seminars-complex-geometry-theorem-Jeffrey-Godinho}
For generic $ \alpha$,
$M(\alpha)$ is diffeomorphic to
$\mathcal{M}_{2,0}(\beta)$
whenever
$\alpha_i=\beta_2(x_i)-\beta_1(x_i)$.
\end{theorem}

\begin{proof}[Idea of proof]
The correspondence can be made quite explicit
as a map
$$
\begin{aligned}
M(\alpha) &\longrightarrow
\mathcal{M}_{2,0}(\beta) \\
[q_1,\ldots,q_n] &\longmapsto
\substack{E=\mathbb{C}P^{1}\times\mathbb{C}^2
 \\ E_{x_i}\supset E_{x_i,1}\supset
 X_{x_i,2}\supset \{0\}}\end{aligned}
$$

\end{proof}

\medskip\noindent
Now consider the quiver variety we described
above:
$$
\text{Rep}Q=\bigoplus_{i=1}^n
\Hom(\mathbb{C},\mathbb{C}^2)=\mathbb{C}^{2n}
.
$$
Now put
$$
K=(U(2) \times U(1)^n)/\Delta
$$
where $\Delta$ is the diagonal $S^1$.

But this time put
$$
\text{Rep}\tilde{Q}=\bigoplus_{i=1}^n
\Hom(\mathbb{C},\mathbb{C}^2)
\oplus \bigoplus_{i=1}^n
\Hom(\mathbb{C}^2,\mathbb{C})
$$
We also have an action of $K$, and we get
$$
K\text{ acts on }\text{Rep}\tilde{Q} = T^*
\mathbb{C}^{2n},
$$
a so-called {\it hyperHamiltonian action}.
The quotient
$$
X(\alpha)=/\!/\!/_{\substack{(0,\alpha) \\
(0,0)}}K=
\frac{\mu_{\mathbb{R}}1(0,\alpha) \wedge
\mu^{-1}_{\mathbb{C}}(0,0)}{K}
$$
called {\it hyperpolygon space}.

Here
$$
\begin{aligned}
\mu_\mathbb{R}:T^*\mathbb{C}^{2n}&\to
\mathcal{K}^* \\
\mu_{\mathbb{C}}:T^*\mathbb{C}^{2n}&  \to
\mathcal{K}_{\mathbb{C}}^* \\
(q_1,\ldots,q_n,p_1,\ldots,p_n)&  \mapsto
(\sum_{i=1}^n(p_i,q_i)_0,\ldots)
\end{aligned}
$$
It turns out that $X(\alpha)$ is smooth
if and only if $\varepsilon_1(\alpha)=0$ for
all
$I \subseteq \{ 1,\ldots,n\}$.
When smooth,  $X(\alpha)$ is a hyperkähler
manifold
(non-compact, unfortunately),
$M(\alpha)=\{[0_{1i},0] \in X(\alpha)\}$.

\begin{theorem}[Boalch]
\label{seminars-complex-geometry-theorem-Boalch}
$X(\alpha)$ is the moduli space of polygons
for $\text{GL}(2,\mathbb{C})$.
\end{theorem}

\medskip\noindent
{\bf Parabolic Higgs bundles.}
Let $E$ be a parabolic bundle as before, i.e.
$E \in \mathcal{M}_{0,2}(\beta)$.
A {\it Higgs field on $E$} is
$$
\phi \in H^{0}(\mathbb{P}^1,\text{SPEnd}(E)
\otimes K_{\mathbb{P}^1}(D)))
$$
where a {\it strongly parabolic endomorphism}
is  $f:E \to E$
such that $f(E_{x_i})\subseteq E_{x_{i+1}}$
for all $i$.
(Fix details!)

A {\it parabolic Higgs bundle} is (probably
the pair $(E,\phi)$).
A parabolic Higgs bundle $(E,\phi)$ is {\it
(semi)stable}
if  $\mu(E) > \mu(L)$ ($\mu(E) \geq \mu(L)$)
for all
$L$ parabolic Higgs subbundle.

Then $\mathcal{N}_{r,a}^{0,1}(\beta)$ is the
moduli
space of parabolic Higgs bundles over
$\mathbb{P}^1$
with rank $r$ and degree $d$ (with fixed
determinant,
and traceless; two notions that we will not
define here).


\begin{theorem}[Goldinho, M.; Biswar,
Florentino, Godinho, M.]
\label{seminars-complex-geometry-theorem-GBSFG}
Let $\alpha$ be generic, let $\beta$
be such that
$\alpha_i=\beta_2(x_i)-\beta_1(x_i)$.
Then the hyperpolygon space $X(\alpha)$
is symplectomorphic to the moduli space
of parabolic HIggs bundles over
$\mathbb{P}^1$, $\beta$-stable,
holomorphically trivial (with fixed
determinant and trace-free).


\begin{proof}[Idea of proof]
We can define a correspondence
$$
\begin{aligned}
X(\alpha) &\longrightarrow \mathcal{H}(\beta)
\\
[p,q] &\longmapsto
\substack{E=\mathbb{C}P^{1}\times
\mathbb{C}^2
 \\ E_{x_i}\simeq E_{x_i,1}\supset
 E_{x_i,2\supset \{0\}}}
\end{aligned}
$$
No we can use the moment map condition
that the sum of the residues is zero,
where $\text{Res}_{x_i}\phi=(p_i,q_i)$
where $\phi$ is the unique function
that satisfies the last equality by
the residue theorem.
\end{proof}
\end{theorem}

\medskip\noindent
{\bf Wall-crossing for moduli spaces of
parabolic bundles.}
See ``Translation''from Thaddeus.
What happens is that there is an $S^1$-action
on $X(\alpha)$ given by
$$
\lambda\cdot[p,q]=[\lambda p,q].
$$
(This can be paralleled with the $S^1$-action
on $\mathcal{H}(\beta)$,
given by
$\lambda\cdot[E,\phi]=[E,\lambda\phi]$.)

This action has a moment map
$$
\mu_{S^1}([p,q])=\frac{1}{2}\sum_{i}|p_i|^2
$$
Fixed points are [Konno]
$$
X_S=\{[p,q] \in X(\alpha):
S,S^c \text{ are straight, }p_j=0 \forall j
\in S\}
$$
for all $|S|\geq 2$ short. Here $S\subseteq
\{1,\ldots,n\}$
is {\it straight} if $q_i=\lambda_j q_j$
for $\lambda_j>0$ $\forall  i,j \in S$
(which in the case of polygons means the
edges are aligned).

Let $U_S$ be the flow down from $X_S$
Then, crossing a wall  $W_I$ as defined
above,
i.e. $W_I=\{\alpha \in
\mathbb{R}^n_{>0}:\varepsilon_S(\alpha)=0\}$,
``replaces'' $U_S$ by $U_{S^c}$.

Let $\alpha$ and $\tilde{\alpha}$ be generic,
then $X(\alpha)$ is diffeomorphic to
$X(\tilde{\alpha})$.
The wall crossing for $M(\alpha)$
involves
$$
\begin{aligned}
M_S(\alpha^+)&=U_S \cap M(\alpha^+)\\
M_{S^c}(\alpha^-)&=U_{S^c}\cap M(\alpha)
\end{aligned}
$$
This is thought of as a Mukai transform.

\medskip\noindent
{\bf Generalization.}
See [Florentino, Godinho, Sotillo] for wall
crossing.
See also [Fisher] PhD thesis, [Fisher,
Rayan].
[Hausel et al], [Rayan-Schopnick].

\section{Vertex algebras and special holonomy
on quadratic Lie algebras}
\label{seminars-complex-geometry-section-verte-algeb-speci-holon}

\noindent
Mario García Fernandez, ICMAT.
GAAG, IMPA.
September 24, 2025.

\medskip\noindent
{\bf Abstract.}
The chiral de Rham complex (CDR) is a sheaf
of vertex algebras on any smooth
manifold, introduced by Malikov, Schechtman
and Vaintrob, which provides a
formal quantization of the non-linear sigma
model in mathematical physics.
Motivated by the algebra of chiral symmetries
in two-dimensional superconformal
field theories, vertex algebra embeddings on
the CDR have been studied for
special holonomy Riemannian manifolds, thanks
mainly to the work of Heluani,
Zabzine, and collaborators, with interesting
applications to the elliptic genus.
In these lectures, we will discuss extensions
of some of these results to the
case of special holonomy manifolds with
skew-torsion. The presence of torsion
typically allows for continuous symmetries in
the geometry, with an enhanced
interplay with Lie theory and algebra, as
well as the application of techniques
from generalized geometry.


\medskip\noindent
Based on joint work with Luis Álvarez Cónsul,
Andoni De Arriba de la Hera. arXiv:2012.01851
(IMRN '24).

\medskip\noindent
{\bf Motivation.}
Let $(M^n,g)$ be a Riemannian spin manifold
with
parallel spinor $\nabla^g \varphi=0$.
Then $\text{hol}(g) \subset G_\varphi
\subseteq \text{SO}(n)$,
and $\text{Ric}(g)=0$.

There is a construction by Markov-S-V that
puts a sheaf
of vertex algebras $\mathcal{V} \to M^n$.
How to construct special embeddings of
trivial vertex algebras
in the cohomology of $\mathcal{V}$, i.e.
$\mathcal{V}_\varphi \hookrightarrow
H^*(\mathcal{V})$
by Heluani et al.

\medskip\noindent
Applications. Construction of topological
invariants.
Elliptic genus, [Borisov-L].
How to understand mirror symmetry using
vertex algebras [Borisov].
Holography [Witten].

\medskip\noindent
{\bf Geometry in algebra.}
We shall do geometry in quadratic Lie
algebras.
Recall that a {\it quadratic Lie algebra} is
 $(\mathfrak{g},[\cdot,\cdot],
 (\cdot,\cdot))$ where
$(\mathfrak{g},[\cdot,\cdot])$ is a Lie
algebra
(in this course we use $\mathbb{R}$ as base
field)
and $(\cdot,\cdot)$ is a symmetric bilinear
invariant form.

\begin{example}
\label{seminars-complex-geometry-example-quadratic-Lie-algebra}
Let $R$ be a Lie algebra with
$\left\langle{}\cdot,\cdot\right\rangle{}_R: R \otimes R \to
\mathbb{R}$.
Take $\mathfrak{g}=R \oplus R^*$,
and pick $H \in \Lambda^{3}R^*$.
Put $H(a,b,c)=\left\langle{}[a,b],c\right\rangle{}$, and
define
$$
[v+\alpha,w+\varphi]=[v,w]-\varphi([v,-])+
\alpha([v,-])+H(v,w,-)
$$
for $v,w \in R$ and $\varphi,\alpha \in R
^*$. This turns out to be a bracket.

If you like geometry you can pick $K$
compact, $\text{Lie}K=R$,
this is ``Generalized geometry on $T K \oplus
T^* K$''
à la Hitchin.
\end{example}

\medskip\noindent
QLA:
\begin{itemize}
\item Courant algebroids /$\{*\}$.
\item  Symplectic supermanifolds.
\end{itemize}

\medskip\noindent
\begin{definition}
\label{seminars-complex-geometry-definition-generalizaed}
\begin{enumerate}
\item A {\it (generalized) metric} on
$(\mathfrak{g},(\cdot,\cdot))$ is
$G \in \text{End}(\mathfrak{g})$ with
$G^2=\text{Id}$, $(G,G)=(\cdot,\cdot)$,

This gives $\mathfrak{g}=V_+ \oplus V_-$
where $G$ acts as identity on the
first term and as $-\text{Id}$ on the second
one.
$(G_-,-)$ is a non-degenerate pairing.
$(-,-)_{V_I}$ non-degenerate tensor.

\item A divergence $\varphi \in \mathfrak{g}
^*$.

\item A {\it connection} is
$$
D: \mathfrak{g} \to \mathfrak{g}^* \otimes
\mathfrak{g}
$$
such that $(D_ab,c)+(b,D_ac)=0$.
So that $D \in \mathfrak{g}^*  \otimes
\Lambda^{2}\mathfrak{g}$
(where we identify $\mathfrak{g}$ with its
dual using the pairing).

$D$ is  {\it compatible} with $G$ if
$[D_a,G]=0$. This condition says that $D$
splits into four operators:
$$
\xymatrix{
D^+_+&  &  D_+^-\\
&  D\ar[ul]\ar[ur]\ar[dl]\ar[dr]\\
D^+_-& & D_-^-
}
$$
\item  Let $a \in \mathfrak{g}$, $a=a_++a_-$.
A {\it generalized connection} satisfies
$D_a^+=[a_-,b_+]_+ +[a_+,b_-]_-$. (This will
happen to things
living in $\mathcal{D}^0$, see below.)
\end{enumerate}
\end{definition}

\begin{definition}
\label{seminars-complex-geometry-definition}
Given $D$ define
\begin{enumerate}
\item $\varphi_D(a)=-\text{tr}Da$.
\item $\text{Torsion}(AX):T_D \in
\Lambda^{2}\mathfrak{g}^*$,
$$
T_D(a,b,c)=(D_ab-D_ba-[a,b],c)+(D_ca,b).
$$
(Just copying the definition of torsion.)
\end{enumerate}
\end{definition}

\begin{lemma}
\label{seminars-complex-geometry-lemma}
Define
$\mathcal{D}^0(G,\varphi)=\{D:G\text{-
compatible
}, \varphi_D=\varphi\}$.
Then this is non-empty, but is not a point.
Furthermore, $\forall  D \in
\mathcal{D}^0(G,\varphi)$,
$$
D_{a_-}b_+=[a_-,b_+]_+,\qquad
D_{a_+}b_-=[a_+,b_-]_-.
$$
\end{lemma}

(Checking non-emptyness is just writing out
the equations.)

\medskip\noindent
Given $G$ consider $Cl (V_+)$,
 $a_+\cdot a_+=(a_+,a_+)$ for $a_+ \in V_+$.

Fix an irreducible representation for
$Cl(V_+)$: $S_+$.
Given  $D \in \mathcal{D}^0(G,\varphi)$:
\begin{enumerate}
\item $D_-^{\delta_+} \in V_-^* \otimes
\text{End}((S_+)\supset
V_-^*\otimes \Lambda^{2}V_+ \ni D_-^+$.
\item An operator like Dirac operator:
$$
\begin{aligned}
\underline{D}^+: S_+ &\longrightarrow S_+ \\
\zeta &\longmapsto \sum_{j=1}^{\dim
V}e_j^+\cdot D_{e_j^+}\zeta.
\end{aligned}
$$
\end{enumerate}

\begin{lemma}
\label{seminars-complex-geometry-lemma-independendant}
$D^{S_+}$ and $\underline{D}^+$ are
independnt of
$D \in \mathcal{D}^0(G,\varphi)$.
\end{lemma}

\begin{definition}
\label{seminars-complex-geometry-definition-Killing-spinor-equations}
$(G,\varphi,\zeta)$, $\varphi \in S_+$
satisfies
{\it Killing spinor equations} if
$$
KSE\qquad
\underbrace{D_-^{S_+}\zeta=0,}
_{\text{Gravitino
Eq.}}\qquad
\underbrace{\underline{D}^+\varphi=0
}_{\text{Dilatino Eq.}}$$
\end{definition}

Expectation.
\begin{enumerate}
\item
$\mathcal{M}=\{(G,\varphi,\zeta):KSE\}/
\mathfrak{g}$
special metric.
\item Given a solution $(G,\varphi,\zeta)$
then
$$
V_{(g,\varphi,\zeta)}\hookrightarrow
V^k(\mathfrak{g}).
$$
(Where I think $V^k(\mathfrak{g})$ is the
Kac-Moody affinization.)
\end{enumerate}

\begin{remark}
\label{seminars-complex-geometry-remark-Lie-algebra}
In the case $T K \oplus T^*K: \mathcal{M}$
2-stack. See [Bursztyn].
\end{remark}

\begin{lemma}
\label{seminars-complex-geometry-lemma-Ricci-tensors}
Associated to $(G,\varphi)$ there are
well-defined ``Ricci tensors''
$\text{Ric}^+ \in V_- \otimes V_+$,
$\text{Ric}^- \in V_+ \otimes V_-$.
$$
\text{Ric}^+(a_-,b_+)=\text{Tr}(C_+ \to
R_D(c_+,a_-)b_+)
$$
$$
R_D(a,b)=[D_a,D_b]_c-D_{[a,b]}c,\qquad D \in
\mathcal{D}^0(G,\varphi).
$$
\begin{remark}
\label{seminars-complex-geometry-remark-in-geometric-setup}
In geometric setup, the vanishing of the
Ricci corresponds to
the motion equations of some physical
supersymmetry theory.
\end{remark}

\begin{proposition}
\label{seminars-complex-geometry-proposition-KSE}
If $(G,\varphi,\zeta)$ is solution of KSE,
then
 $\text{Ric}^+_{G,\phi}=0$.
If $[\varphi,G]=0 \implies
\text{Ric}^-_{G,\varphi}=0$.
\end{proposition}
\end{lemma}

\begin{proof}
$\text{Ric}^+(a-,-)\cdot
\varphi=[\underline{D}^+,D_{a_-}^{S_+}\zeta
-D_{D_+a_-}^{S_+}\zeta=0$.
$[\varphi,G]=0 \implies
\text{Ric}^+(a_-,b_+)=\text{Ric}^-(b_+,a_-)$.
\end{proof}

\medskip\noindent
{\bf Generalized Ricci flow.}
Finding solutions to these equations.
Evaluating by
$$
G_t^{-1}\partial_tG_t=-2(\text{Ric}^+-
\text{Ric}^-)
$$
$\Hom(V_+,V_-) \oplus \Hom(V_-,V_+)$.
$G G+G - G=0$,
$G=\text{Id}_{V_+}-\text{Id}_{V_-}$.

\begin{exercise}
\label{seminars-complex-geometry-exercise}
\begin{enumerate}
\item Prove STE for GRF.
\item Assuming there exists a solution of
KSE, prove long
time existence and convergence. (See theorem
by Streets, Jordan, GF;
and a book by Streets, GF.)
\end{enumerate}
\end{exercise}

\begin{remark}
\label{seminars-complex-geometry-remark-sigma-model}
For physicists the generalized Ricci flow
(GRF) is the GRF of a $2d$
$\sigma$-model with target a compact Lie
group
$K$,  $\mathfrak{g}=R \oplus R^*$
\end{remark}

Let's explain something in this situation.
consider a compact Lie group $K$ and
$\mathfrak{g}=R \oplus R^*$.
A generalized metric: $\mathfrak{g}=V_+
\oplus V_-$,
$(-,-)|_{V_\pm}$ non-degenerate,
$(-,-)|_{V_+}>0$.
$V_+=\mathcal{C}^b\{X+g(X)\}$,
$X \in R$, $g \in S^2(R^*)$, $b \in
\Lambda^{2}R^*$,
$H=H_0+\overline{\partial}$.

\medskip\noindent
{\bf Case $\dim V_+=2n$.}
Complex pure spinor $\varphi$ on $V_+$ is
equivalent $V_+ \otimes \mathbb{C}=\ell
\oplus \overline{\ell}$,
where $\ell=\{0 \in V_+ \oplus \mathbb{C}:0
\cdot\zeta=0\}$.
$(\ell,\ell)=0$,
$(\overline{\ell},\overline{\ell})=0$.

This gives a complex structure
$J_\zeta:V_+ \to V_+$, $J_\zeta=i
\text{Id}_\ell-i
\text{Id}_{\overline{\ell}}$.

\begin{lemma}
\label{seminars-complex-geometry-lemma-solutions-KSE}
$(G,\varphi,\zeta)$ with $\zeta$ pure
satisfies KSE if and only if
\begin{enumerate}
\item $[\ell,\ell]\subset \ell$.
\item Take bars
$\{\varepsilon_j,\overline{\varepsilon}_j\}
^n_{j=0}$,
$\varepsilon_j \in \ell$,
$\overline{\varepsilon}_j\in \overline{\ell}$
$(\varepsilon_j,\overline{\varepsilon}_j) =
s_{j,k}$.
$\sum_{j=1}^n[\varepsilon_j,
\overline{\varepsilon}_j=-J_\zeta
\varphi_+$.
\end{enumerate}
\end{lemma}
(This sum is a moment map condition.)

Consider $\mathcal{L}=\{\ell \subset
\mathfrak{g} \otimes \mathbb{C}:
\dim\ell=n,\ell\cap\overline{0},(-,-)|_{\ell
\oplus \overline{\ell}}
\text{non-degenerate}\}$.
This is a complex manifold.
$T_\ell \mathcal{L}
\cong \Hom(\ell,\overline{\ell}\oplus V_-
\oplus \mathbb{C})$,
$V_- \otimes \mathbb{C}=(\ell \oplus
\overline{\ell})^\perp$
Pick a vector $\dot \ell = \dot J + \dot G
\in
\Hom(\ell,\overline{\ell}\oplus V_- \oplus
\mathbb{C})$
That is, $\dot J \in \Hom(\ell,
\overline{\ell})$,
$\dot J :V_+ \to V_+$,
$\dot J J+ J \dot J=0$. (Deformations of the
complex structure.)
$G \in \Hom(V_+,V_-) \oplus \Hom(V_-,V_+)$.

\begin{proposition}[Romero]
\label{seminars-complex-geometry-proposition-Romero}
$\mathcal{L}$ has a pseudo-Kähler structure
preserved by the $\mathfrak{g}$-action and
there exists
a moment map
$$
\begin{aligned}
\mu: \mathcal{L} &\longrightarrow
\mathfrak{g}^* \\
\ell &\longmapsto
\frac{i}{2}\sum_{j=1}^n[\varepsilon_j,
\overline{\varepsilon}_j],-),
\end{aligned}
$$
which is the quantity we mentioned in Lemma
above.
\end{proposition}

\begin{proof}
Use the natural complex structure on the
space of complex structures
(found by Fujiki),
$$
\text{tr}_{V_+}(J \dot G_2,\dot
G_1)-\text{tr}_{V_+}(J \dot J_{\varepsilon}
\dot J_1).
$$
\end{proof}

Problem: prove that
$\mathcal{M}=\{(G,\varphi,\zeta):\zeta\text{
puse}\}/\mathfrak{g}$
is a pseudo-Kähler manifold.

\begin{enumerate}
\item Pseudo-Kähler.
\item Shifter symplectic stuff.
\end{enumerate}

\begin{example}
\label{seminars-complex-geometry-example-SU2xSU1}
Take $K=\text{SU}(2)\times\text{U}(1)\cong
S^3\times S^1$
and $\mathfrak{g}=R \oplus R^*$.
Take generators $v_1,v_2,v_3$ of
$\text{SU}(2)$ and $v_4$ of $\text{U}(1)$.
We have $[v_2,v_3]=-v_1$, $[v_3,v_1]=-v_2$,
$[v_1,v_2]=-v_3$.
Put $H_\ell=\ell v^{123}$. $\ell \in
\mathbb{R}$. $x,a \in \mathbb{R}_{>0}$.
$$
g_{x,a}=\frac{a}{x}\left(\sum_{i=1}
^3\omega^{\otimes
2}
+x^2 (v^4)^{\otimes 2}\right)
$$
$$
V_+=\{ x+g_{xa}\}\subset \mathfrak{g}
$$
$$
I_X v_4=xv_1,\qquad U_X v_2=v_3, \qquad
\varphi=-xv_4
$$
\end{example}

\begin{exercise}
\label{seminars-complex-geometry-exercise-solution-of-KSE}
If $\ell=\frac{a}{x} \implies $ solution of
KSE.
\end{exercise}

Furthermore $[\varphi_+,\ell] \subset \ell$
(holomorphic divergence).

The $a$ parameter is naturally complexified
by  $b=yv^{23}$.
$y+ia=z$.
$$
V=\text{log}
\left(\frac{\omega^2_{x,n}}{v^{1234}}\right)=
\text{log}a.
$$
KSV hyperbolic metric, metric on
$\mathbb{H}$.

\medskip\noindent
{\bf $\dim V_+=7$}.
A real spinor on $V_+$ is equivalent to $\phi
\in (\Lambda^{3}V_+^* )_{>0}$.
The action of $\text{GL}(\mathbb{R}^7)$ on
$\Lambda^{3}(V_+^*)$.
The space of spinors $S_+=\mathbb{R}^8
=\mathbb{R}^7\oplus\mathbb{R}
\left\langle{}\varphi\right\rangle{}$.
$\phi(x,y,z)=\left\langle{}x\cdot y\cdot z\cdot
\zeta,\zeta\right\rangle{}$.

\begin{remark}
\label{seminars-complex-geometry-remark-7}
$v^1,\ldots,v^7 \cong \mathbb{R}^6 \oplus
\mathbb{R}$.
$\phi=(v^{12}+v^{34}+v^{56})\wedge v^7+
\text{Re}((v^1+i
v^2)\wedge(v^3+iv^4)\wedge(v^5+iv^6))$.
\end{remark}

\begin{example}
\label{seminars-complex-geometry-example-SU2-SU2-R}
Take $R=\text{SU}(2)\oplus \text{SU}(2)\oplus
\mathbb{R}$.
$\mathfrak{g}=R \oplus R^*$.
$e,s \in \mathbb{R}$, $H=sv^{123}+\ell
v^{456}$.
$\phi=\omega \wedge \zeta+\Omega^+$,
$\zeta=\sqrt{\varepsilon/\ell}v^7$,
$\omega=\sqrt{s\ell}(v^{14}+v^{25}-v^{36})$,
$\Omega^+=\sqrt{s^3}v^{123}+\ell
\sqrt{s}v^{156}-\ell \sqrt{s}v^{345}$.
\end{example}

\section{Non-Kähler Hodge Lefschetz theory
and the Bianchi identity}
\label{seminars-complex-geometry-section-non-kahle-hodge-lefsc-theor}

\noindent
Arpan Saha, UNICAMP.
Geometric Structures Seminar, IMPA.
September 25, 2025.

\medskip\noindent
{\bf Abstract.}

Being Kähler imposes severe constraints on
the cohomology of compact complex
manifolds such as the Hard Lefschetz
property, and the question of how far this
generalises beyond the class of Kähler
manifolds has been of great interest for
a while. In this talk, I shall report on
ongoing joint work with Mario García
Fernández and Raúl González Molina that
abstracts out the definition of a
variation of Hodge--Lefschetz structure and
provides evidence that, under
certain natural assumptions, such a structure
exists more generally on
distinguished subspaces within moduli spaces
of Bismut--Ricci-flat metrics that
are pluriclosed up to source terms. In
particular, these distinguished subspaces
may be regarded as replacements for the
Kähler cone, with affine structure
modelled on a subspace of the (1,1) Aeppli
cohomology of the compact complex
manifold.

\medskip\noindent
{\bf Broad motivation.}
As we move on the parameter space associated
to a CY manifold we encounter walls.
The Kähler cone is contained in this
parameter space, and we may cross
its walls. When we cross a wall
we obtain a birational transformation
(I think this means that the moduli
on either side are birational).
But the Kähler condition is not a
birational invariant.

\medskip\noindent
{\bf Hodge-Lefschetz theory.}
For $(X,J,\omega)$ Kähler,
the cohomology $H^{\bullet}(X)$ comes with:
\begin{itemize}
\item the Lefschetz operator $L$,
given by $[\omega]\wedge-$.

\item Hodge star operator.
\item $\Lambda=* L *$ satisfying
the $\mathfrak{sl}_2$ relations:
$$
[L,\Lambda]=H,\qquad [H,L]=2L,\qquad
[H,\Lambda]=-2\Lambda.
$$
These relations are equivalent to
the Hard Lefschetz property that
$$
L^q:H^{d-q} \to H^{d+q}.
$$
\item The Poincaré pairing
(which is a different pairing from Lefschetz)
of the form
$$
H^{d-q}\cong (H^{d+q})
$$
given by integration.

\item Hodge-Riemann bilinear relations
which is a positive definite form given by
$\int_X\cdot \wedge *\cdot$.

\end{itemize}

\noindent
Moving about on the Kähler cone gives
a variation of Hodge-Lefschetz structure.

\begin{definition}
\label{seminars-complex-geometry-definition-varia-hodge-lefsc-struc}
A {\it variation of Hodge Lefschetz structure
(VHLS)}
of weight $d$ over a manifold $K$ consists of
\begin{itemize}
\item a graded real vector bundle
$E=\bigoplus_{q=0}^d E^q \to K$
with $\text{rk}_\mathbb{R}E^0=1$.

\item An isomorphism $E^0 \otimes TK
\xrightarrow{\simeq}E^1$
(this is reminiscent of the usual definition
of variations of Hodge structures, and
we will not use it in this talk).

\item Endomorphism fields
$L, \Lambda,H \in
\mathcal{A}^0(K,\text{End}(E))$
satisfying the $\mathfrak{sl}_2$ relations
that
$$
[L,\Lambda]=H,\qquad [H,L]=2L,\qquad
[H,\Lambda]=-2\Lambda.
$$
and that
$$
H|_{E_q}=(2q-d)|_{\text{id}_{E^q}}.
$$
\item An involution $* \in
\mathcal{A}^0(K,\text{End}(E))$
such that $* L *=\Lambda$.

\item A nondegenerate symmetric pairing
$P \in \mathcal{A}^0(\mathcal{K},E^* \otimes
E^*)$
w.r.t. $L$ and $*$ are self-adjoint.

\item (like the Gauss-Manin connection)
a flat connection $D:
\mathcal{A}^0(K,E)\to\mathcal{A}^1(K,E)$
such that $DH=0=DP$.
\end{itemize}
The VHLS is said to be {\it positive} if
$P(*-,-)$
is positive definite.
\end{definition}

\noindent
{\bf Calabi's dream beyond Kähler manifolds.}
Let $(X,\Omega)$ be a compact Calabi-Yau
manifold,
where $\Omega$ is the volume form.
In general such a manifold is not Kähler.

\begin{example}[Non-Kähler manifolds]
\label{seminars-complex-geometry-example-non-kahler}
These manifolds are not Kähler:
Heisenberg algebra quotiented by some
lattice,
torus bundle, Hopf surface.
\end{example}

\noindent
By [Yau], the Kähler cone is the moduli
space of Kähler Ricci-flat metrics.

We look for PDEs on  $X$ such that
\begin{itemize}
\item Assumption C (Calabi). Moduli
space of alutions is an open set $K$
of an affine space modelled on
a subspace $\mathbb{V}$ of
$(1,1)$-Aeppli cohomology group,
which is given by
$$
H^{1,1}_A(X)=\frac{\Ker
\partial\overline{\partial}|_{\Omega^{1,1}}}
{\text{Im}
\partial+\text{Im}\overline{\partial}}
$$

\item Assumtion D (dilaton).
A volume function $v:K\to \mathbb{R}_{>0}$
such that $\underline{d}\text{log} V$
is nowhere zero (where $\underline{d}$
denotes
Aeppli differential), and the bilinear form
$g_K:=-Dd \text{log}V$
is nondegenerate (basically the Hessian
metric).
(Ideally this would be positive definite.)

\item Assumption V (vector field).
$Z=g_K^{-1}d \text{log}V$
is such that $DZ$ is invertible
(when thought of as an endomorphism)
$g_K(Z,Z)=d \in \mathbb{Z}_{>0}$,
 $$
(D_{(DZ)^{-1}})^{d+1}V=0,
$$
\end{itemize}

\begin{remark}
\label{seminars-complex-geometry-remark-pairing}
$$
(D_{(DZ)^{-1}})^dV=\int_X \wedge \wedge
\wedge.
$$
\end{remark}

\begin{theorem}[García Fernández-González
Molina-AS]
\label{seminars-complex-geometry-theorem-GGA}
If assumptions C, D and V all hold with  $d
\leq 3$,
then $K$ admits a VHLS of weight $d$.
\end{theorem}

\begin{proof}[Idea of proof for $d=3$]
$$
E=\underbrace{\mathbb{R}_K}_{E^0} \oplus
\underbrace{TK}_{E^1}
\oplus \underbrace{T^*K}_{E^2} \oplus
\underbrace{\mathbb{R}_K^*}_{E^3}.
$$
Define $L$ in each of the terms of the direct
sum as
$$
L1=Z,\qquad Lv=D^2_vV,\qquad
L_\alpha=i_Z\alpha
$$
and $*$ as
$$
*1=V 1^\vee,\qquad  *v=VD^2_0
\text{log}V=Vg_K(v,-).
$$
\end{proof}

\medskip\noindent
Recall the Hull-Strominger
system, which is given by two conditions
on the connection and the metric
$$
d\left(\|\Omega\|_\omega
\frac{\omega^{n-1}}{(n-1)!}\right)=0,\qquad
F_\theta \wedge
\frac{\omega^{n-1}}{(n-1)!}=0.
$$
This determines $d
d^c\omega=\alpha\left\langle{}F_\theta \wedge
F_\theta\right\rangle{}$.

\section{Formality of Dolbeault cohomology}
\label{seminars-complex-geometry-section-forma-dolbe-cohom}

\noindent
Misha Verbitsky, IMPA.
Geometric Structures Seminar, IMPA.
October 29, 2025.

\medskip\noindent
{\bf Abstract.}
A quasi-isomorphism of differential graded
algebras (DGA) is a multiplicative
map inducing an isomorphism on cohomology. A
DGA is called  formal if it can be
connected by a chain of  quasi-isomorphisms
to its cohomology algebra. For the
de Rham algebra of a compact manifold,
formality is an important  topological
property which has many implications (such as
vanishing of Massey products). For
the Dolbeault DGA of a compact complex
manifold, the situation is not well
understood, unless it is Kahler. The
Dolbeault DGA of a torus, symmetric spaces
and hyperkahler manifolds are known to be
formal.  We prove that the Dolbeault
DGA of a complex nilmanifold is formal only
if it is a torus, and the Dolbeault
algebra of (0,p)-forms is formal if and only
if the complex structure is
abelian. This is a joint work with Tommaso
Sferruzza.

\medskip\noindent
Upshot: the Dolbeault dga of a nilmanifold
$M$ is never formal.

\medskip\noindent
If a space is formal, Massey products vanish.
Unfortunately, Massey products are not
well-defined.
Let $a,b,c \in \Lambda^*(M)$ be closed forms
such that their cohomology classes satisfy
$[a][b]=[b][c]=0$.
Let  $\alpha,\gamma \in \Lambda^*(M)$
be such that $d(\alpha)=a \wedge b$ and
$d(\gamma)b \wedge c$.
Then the {\it Massey product} of $a,b,c$
is the class  $[\alpha \wedge
c-(-1)^{\tilde{a}}a \wedge \gamma]$
which is only defined up to
 $\left\langle{}\text{img}L_a + \text{img}
 L_c\right\rangle{}$
where $L_a$ and $L_c$ are multiplication by
$a$ and $c$ operators.

The {\it Heisenberg nilmanifold} is the
quotient of
the Heisenberg group $G$ of upper triangular
matrices
with $1$ on the diagonal over
$G_{\mathbb{Z}}$
of the same group with integer entries.
The Heisenberg nilmanifold fibers over a
torus
with fiber a circle:
$$
\xymatrix{
0\ar[r]&S^1\ar[r]&G/G_{\mathbb{Z}}\ar[r]&
\mathbb{R}^2/\mathbb{Z}^2\ar[r]&0
}
$$
which in fact comes from the central
extension
$$
\xymatrix{
0\ar[r]&\mathbb{R}\ar[r]&G\ar[r]&\mathbb{R}
^2\ar[r]&0.
}
$$


\begin{lemma}
\label{seminars-complex-geometry-lemma-heise-nilma-masse-produ}
Massey products on  $G/G_{\mathbb{Z}}$ are
non-zero.
\end{lemma}

\begin{proof}
Step 1. There is an action of $G$ in the
cohomology of $G$.
It's not hard to see that the cohomology
classes on $G/G_{\mathbb{Z}}$
can be represented by right $G$-invariant
forms,
and the cohomology of $G/G_{\mathbb{Z}}$ is
equal to the cohomology of the complex
of right-$G$-invariant forms on $G$.

\medskip\noindent
Step 2. In fact, this complex is the same as
the
Chevalley-Eilenberg complex of the Lie
algebra $\mathfrak{g}$ of $G$,
and the cohomology is the Lie algebra
cohomology.

\medskip\noindent
Step 3. From a basis $a,b,t$ of
$\mathfrak{g}$ with the only nontrivial
commutator $[a,b]=t$ and $\alpha,\beta,\tau$
dual basis in $\mathfrak{g}^*$
with the only nontrivial differential
$d\tau=\alpha\wedge\tau$,
we obtain a basis $\alpha \wedge \beta,
\alpha \wedge \tau, \beta \wedge \tau$
in $\Lambda^2(\mathfrak{g}^*)$ with
$d|_{\Lambda^2\mathfrak{g}^*}=0$.
This gives $\text{rk}
H^1(G/G_{\mathbb{Z}})=2$
and $\text{rk}H^2(G/G_{\mathbb{Z}})=2$.

\medskip\noindent
Step 4. Compute the Massey product of
$\alpha,\beta,\tau$:
it does not vanish.
\end{proof}

\medskip\noindent
Digression: this is Sullivan's motivation for
working on dgas.
Consider the class of spaces whose homotopy
and homology
groups are vectors spaces over $\mathbb{Q}$
(and it's better to also assume that $\pi_1$
is nilpotent).
There is a equivalence of categories between
old category
and homotopy category of rational homotopy
spaces.
And in fact it's also equivalent to the
(homotopy?)
category of dga algebras.

\medskip\noindent
Now we move to differential graded algebras
(dga).
A morphism of dgas (which I think only means
that commutes with
the differentials) is called {\it
quasi-isomorphism} if it
induces isomorphisms on all $H^i(A^\bullet)$.

A dga is {\it formal} if it is connected to
$H^\bullet(A^\bullet)$
by a chain of quasi-isomorphisms.

Recall the $dd^c$-lemma. It states that for
the twisted
differential $d^c=Jd J^{-1}$, if a form
$\alpha$ is $d$-closed and $d^c$-exact,
(or $\partial$-exact and
$\overline{\partial}$-exact),
then $\alpha\in \text{img}dd^c$.

\begin{lemma}
\label{seminars-complex-geometry-lemma-quasi-isomorphisms}
The map $(\Lambda^\bullet(M)_{d^c},d)\to
(\Lambda^\bullet(M),d)$
is a quasi-isomorphism.
\end{lemma}

\begin{proof}
Uses Hodge decomposition of a form to prove
induced maps
on cohomology are injective. Then prove
surjectivity.
\end{proof}

\begin{lemma}
\label{seminars-complex-geometry-lemma-kahler-formal}
A compact Kähler manifold is formal (i.e. the
usual cohomology algebra
is formal).
\end{lemma}

\begin{theorem}
\label{seminars-complex-geometry-theorem-kahler-dolbeault-formal}
$M$ Kähler then
$(\Lambda^\bullet(M),\overline{\partial})$ is
formal.
\end{theorem}

\medskip\noindent
A manifold $M$ is {\it nilmanifold} if it has
an action of aa nilpotent Lie
group. All nilmanifolds are obtained as
qutient spaces $M=G/\Gamma$
for a cocompact lattice in a nilpotent Lie
group.

I gather there are two things:
one is an {\it integrable complex structure
on a Lie algebra}, which
is a subalgebra
$\mathfrak{g}^{1,0}\subset \mathfrak{g}
\otimes_\mathbb{R} \mathbb{C}$
such that $\mathfrak{g}^{1,0} \oplus
\overline{g^{1,0}}
=\mathfrak{g} \otimes \mathbb{C}$.
Integrability is
$[\mathfrak{g}^{1,0},\mathfrak{g}^{1,0}]
\subset \mathfrak{g}^{1,0}$.

And then there's a {\it complex nilmanifold},
which is $M=G/\Gamma$
with a complex structure $I$ such that $G$
has
a right-invariant complex structure and $G
\to M$ is holomorphic.
In this case we say $I$ is {\it invariant}.
Not all complex structures on a nilmanifold
are invariant:
a compact torus admits a non-Kähler complex
structure that is not invariant.
A Kodaira surface, a locally trivial
holomorphic fibration over $T$ with fiber
$T'$ with non-trivial Chern class, for $T$
and $T'$ elliptic curves,
is a nilmanifold.

Recall the Chevalley-Eilenberg cohomology
on $\Lambda^\bullet\mathfrak{g}^*$ where the
Lie bracket becomes
a differential by Cartan formula.
More generally,

\begin{proposition}
\label{seminars-complex-geometry-proposition-cheva-eilen-diffe}
Let $w \in \Hom(\Lambda^2V,V)$.
Consider the dual map $d_w:V^* \to
\Lambda^2V^*$.
Extend this map to $d_w:\Lambda^kV^*  \to
\Lambda^{k+1}V^*$
using Leinbiz rule. Then $d^2_w=0$ if and
only if $w$
defines the Lie algebra structure on $V$.
\end{proposition}

\noindent
We may identify the Chevalley-Eilenberg
complex of the Lie algebra  $\mathfrak{g}$
with the
complex of left-invariant differential forms
in its
Lie group $G$. This defines a natural
homomorphism of dgas
$(\Lambda^\bullet\mathfrak{g}^*,d) \to
(\Lambda^\bullet(G/\Gamma),d)$
for any discrete group $\Gamma \subset G$.

\begin{theorem}[Nomizu]
\label{seminars-complex-geometry-theorem-nomizu}
Let $M=G/\Gamma$ be a nilmanifold.
Then the map above is a quasi-isomorphism.
\end{theorem}

As a corollary, we see that nilmanifolds,
with exception of tori, are never formal.

\begin{proposition}
\label{seminars-complex-geometry-proposition}
Let $(\mathfrak{g},I$) be a complex structure
on a Lie algebra. Then we may consider the
Dolbeault
cohomology of
$\partial_w:\Lambda^\bullet(\mathfrak{g}^*)
\to\Lambda^2(\mathfrak{g}^*)$.
Then $[\cdot,\cdot]_{\partial_w}$ is defined
as follows: let
$x,y \in \mathfrak{g}^{1,0}$, $x',y' \in
\mathfrak{g}^{0,1}$.
Then
$$
\begin{aligned}
[x,y]_{\partial_w}&=[x,y]\\
[x,x']_{\partial_w}&=[x,x']^{1,0}\\
[x',y']_{\partial_w}&=0.
\end{aligned}
$$
\end{proposition}

\noindent
This implies that $(CE,\overline{\partial})$
is formal if and only if
$[\cdot,\cdot]=0$.

\medskip\noindent
We intend to apply the construction made
before for the usual
Lie algebra cohomology but for Dolbeault
cohomology.
Unfortunately there is no analogue of
Nomizu's theorem Theorem
\ref{seminars-complex-geometry-theorem-nomizu}.

\noindent
[Missing slides]

\medskip\noindent
We say a dga has {\it Poincaré duality} if
the product
map $H^i(A^* ) \times H^{n-i}(A^*)\to
H^n(A^*)=\mathbb{C}$
is a non-degenerate pairing in cohomology.
A morphism of dgas $(B^*,d)\to(A^*,d)$, where
$(B^*,d)$ admits
Poincaré duality, that induces injective maps
on cohomology
is called a {\it domination}.

\begin{theorem}[Milivojevic,Stelzig,Zoller
arXiv
\href{http://arxiv.org/abs/2306.12364}%
{2306.12364}]
\label{seminars-complex-geometry-theorem-msz}
Let $\psi:(B^*,d)\to (A^*,d)$ be a domination
of dga.
If $(B^*,d)$ is not formal, then $(A^*,d)$ is
not formal.
\end{theorem}

\begin{theorem}[Barberis-Dotti-V., 2007]
\label{seminars-complex-geometry-theorem-bdv7}
A complex nilmanifold has trivial canonical
bundle,
trivialized by an invariant holomorphic top
form.
\end{theorem}

As corollaries we get that the natural dga
morphism
$(\Lambda^\bullet\mathfrak{g}^*,
\overline{\partial}_w)\to
(\Lambda^\bullet
G/\Gamma,\overline{\partial})$ is a
domination,
and that the Dolbeault dga of a nilmanifold
$M$ is never formal.

\section{Dolbeault cohomology of
nilmanifolds}
\label{seminars-complex-geometry-section-dolbe-cohom-nilma}

\noindent
Misha Verbitsky, IMPA.
Geometric Structures Seminar, IMPA.
November 6, 2025.

\medskip\noindent
{\bf Abstract.}
A nilmanifold is a (left) quotient of a
nilpotent Lie group by a cocompact
lattice. If this Lie group is equipped with a
left-invariant complex structure,
the quotient is called a complex nilmanifold.
Examples of complex nilmanifolds
include Kodaira-Thurston surface and many
other non-Kahler complex manifolds.
The differential graded algebra of invariant
differential forms on a nilmanifold
has the same cohomology as a nilmanifold, as
shown by Nomizu. I am going to
explain a Dolbeault cohomology version of
this theorem. A Lie algebra of a
nilmanifold admits a natural rational
structure, induced by the cocompact
lattice. If the complex structure operator,
acting on the Lie algebra, has
rational coefficients, Console and Fino
proved that the Dolbeault cohomology of
a nilmanifold is equal to the Dolbeault
cohomology of its invariant forms. I
will describe an attempt to prove this result
for a general nilpotent Lie
algebra, based on number-theoretic arguments.

\medskip\noindent
Let $G$ be a nilpotent Lie group.
[For me this would mean that the Lie algebra
of $G$ is
nilpotent as in Lie Algebras Definition
\ref{lie-algebras-definition-nilpo-%
solva-lie-algeb},
but according to Misha this would be the case
for $G$
connected only; so there's must some other
notion
of nilpotent Lie groups.]
A {\it nilmanifold} $M$ is a manifold with a
transitive action
of a nilpotent Lie group. This means that
$M=G/\Gamma$.

\begin{theorem}[]
\label{seminars-complex-geometry-theorem}

\end{theorem}

\noindent
Let $G$ be a real Lie group
and $I \in \text{End}(TG)$ with
$I^2=-\text{Id}$.
Suppose $I$ is right-invariant.
Then $I$ is integrable iff
$[T^{1,0}G,T^{1,0}G]\subset T^{1,0}G$
iff
$\mathfrak{g}^{1,0}\subset\mathfrak{g}
\otimes\mathbb{C}$
where $\mathfrak{g}^{1,0}$ is the
$\sqrt{-1}$-eigenspace of $I$.

A {\it complex nilmanifold} is
$(G,I)/\Gamma$.

\begin{example}
\label{seminars-complex-geometry-example-iwasawa-manifold}
An {\it Iwasawa manifold} is a quotient of
the set
of upper triangular matrices with 1 on the
diagonal
(this group is $U_3(\mathbb{C})$) by a
lattice $\Gamma$.
For example, we can take
$\Gamma=U_3(\mathbb{Z}[\sqrt{-1}])$.
\end{example}

\noindent
The complex structure $I$ is bi-invariant iff
it is paralelizable.

\medskip\noindent
Let $E$ be an elliptic curve with an ample
line bundle $L$.
We put an action on $\text{Tot}L$ given by
multiplication by
$\lambda \in \mathbb{C}$ for $|\lambda|>0$.

\begin{theorem}
\label{seminars-complex-geometry-theorem-de-rham-cheva-eilen}
De Rham cohomology of a nilmanifold coincides
with the
CE cohomology of $\mathfrak{g}$.
\end{theorem}

\begin{theorem}[Console-Fino]
\label{seminars-complex-geometry-theorem-console-fino}
For complex nilmanifolds,
$H^{\bullet,\bullet}(\mathfrak{g}^\bullet,
\overline{\partial}$
and
$H^{\bullet,\bullet}_{\overline{\partial}}(M)
$
are isomorphic.
\end{theorem}

\noindent
There's several results addressing a converse
situation
(when the algebras are isomorphic) and other
situations
involving fibrations.

\medskip\noindent
Now we make a sudden change to deformation
theory.

\begin{lemma}
\label{seminars-complex-geometry-lemma-duplicate-2}
Complex spaces on $V=\mathbb{R}^{2n}$ are
uniquely determined
by $W \in V \otimes_{\mathbb{R}}\mathbb{C}$
such that $W \oplus \overline{ W}=V \otimes
\mathbb{C}$
which is an element in $\text{Gr}(n, V
\otimes \mathbb{C})$.
\end{lemma}

\begin{definition}
\label{seminars-complex-geometry-definition-duplicate-2}
Let $\mathfrak{g}$ be a Lie algebra and
$\mathfrak{g}_\mathbb{C}$
its complexification.
The set of complex structures on
$\mathfrak{g}$ is $\text{Comp}(\mathfrak{g})$
which is $\{ W \in
\text{Gr}_0(\mathfrak{g}_\mathbb{C}:
W \subset W \subset
\mathfrak{g}_{\mathbb{C}}\text{ is a
subalgebra}\}$
since being a subalgebra means
$[\mathfrak{g}^{1,0},\mathfrak{g}^{1,0}]
\subset \mathfrak{g}^{1,0}$.
\end{definition}

We can define a variety out of this:
$$
\overline{\text{Comp}(\mathfrak{g}
_{\mathbb{C}}}
=\{W \in
\text{Gr}(n,\mathfrak{g}_{\mathbb{C}}:
\text{$W$ subalgebra}\}
$$
is a complex subvareity of
$\text{Gr}(n,\mathfrak{g}_{\mathbb{C}})$.


\begin{definition}
\label{seminars-complex-geometry-definition-field-of-defition}
Let $X \subset \mathbb{C}P^n$ be a projective
subvariety.
Consider the smallest field $k \supset
\mathbb{Q}$ such that
any ideal $I$ of $X$ [ideal in the coordinate
algebra of $X$?]
can be generated by polynomials with
coefficients
in $k[k^n]$ is called a {\it field of
definition}.
\end{definition}

\noindent
Apparently there's some discussion about this
notion
not being easily defined.

\begin{lemma}
\label{seminars-complex-geometry-lemma-duplicate-3}
$\text{Gal}_{\mathbb{Q}}(k)$ acts on
$\mathbb{P}_k^n$.
Let $X \subset \mathbb{P}^n$ be a projective
subvariety
and $\Gamma \subset
\text{Gal}_{\mathbb{Q}}(k)$
the stabilizer of $X$.
Then the field of definition of $X=\{\lambda
\in k:\Gamma(\lambda)=\lambda\}$.
\end{lemma}

\noindent
We have a simpler definition of field of
definition for a point
(which is apparently the former definition
when the variety is a point)

\begin{definition}
\label{seminars-complex-geometry-definition-field-of-definition-point}
For a point $(z_1,\ldots,z_n) \in
\mathbb{C}^n$, its
{\it field of definition} is the field
generated by $z_1,\ldots,z_n$.
\end{definition}

\noindent
So, for instance, a rational point has
transcendence degree 0.

\begin{remark}
\label{seminars-complex-geometry-remark}
Let $X$ be an $n$-dimensional algebraic
variety
defined over $\mathbb{Q}$ or any finite
extension of $\mathbb{Q}$,
and let $\pi:X\to \mathbb{C}^n$ be finite
dominant over $\mathbb{Q}$.
\end{remark}


\noindent
We have a fun theorem:

\begin{theorem}
\label{seminars-complex-geometry-theorem-fun-theorem}
Let $X$ be an irreducible complex algebraic
variety defined over $\mathbb{Q}$.
Then the group
$\text{Aut}_\mathbb{Q}(\mathbb{C})$ acts
transitively
on points of maximal transcendence degree in
$X$.
\end{theorem}

\begin{definition}
\label{seminars-complex-geometry-definition-algebraic-point}
We say that a point $x \in X$ is {\it
algebraic} if
its transcendence degree over $k_X$ is zero.
\end{definition}

\noindent
It's not very clear whether the variety $X$
is projective or affine.

\begin{lemma}
\label{seminars-complex-geometry-lemma-algebraic-points-are-dense}
The set of algebraic points in
$X\subset\mathbb{C}P^n$ is dense
in the analytic topology.
\end{lemma}

\begin{theorem}
\label{seminars-complex-geometry-theorem-aut-orbit-of-point}
Let $X$ be a complex (affine) algebraic
variety defined over $\mathbb{Q}$.
Consider an
$\text{Aut}_{\mathbb{Q}}\mathbb{C}$-orbit of
a point
$p \in X$. Then the closure of its orbit is a
minimal subvariety $Z \subset X$ defined
over $\mathbb{Q}$ and containing $p$.
\end{theorem}

\noindent
Now we introduce a Galois action on Dolbeault
cohomology.

\begin{theorem}
\label{seminars-complex-geometry-theorem-duplicate-2}
$\mathfrak{g}$ nilpotent Lie algebra, $M$
corresponding nilmanifold,
$\sigma \in
\text{Aut}_{\mathbb{Q}}(\mathbb{C})$,
and let $J=\sigma(I) \in \text{Comp}(G)$.
Then $\dim H^{p,q}_{\overline{\partial}}(M,I)
=\dim H^{p,q}_{\overline{\partial}}(M,J)$ for
all $p,q$.
\end{theorem}

From this it follows that

\begin{lemma}
\label{seminars-complex-geometry-lemma-console-fino-complex-structure}
For a nilmanifold with complex structure $I$,
if the field of definition of $I$ is
$\mathbb{Q}$,
then the Console-Fino map
$H_{\overline{\partial}}^{p,q}(\mathfrak{g})
\to
H_{\overline{\partial}}^{p,q}(M)$ is an
isomorphism.
\end{lemma}

\noindent
The idea is that if we know the result for
algebraic
points then we'll have it for transcendental
points as well.
Namely,

\begin{lemma}
\label{seminars-complex-geometry-lemma-if-conso-fino-algeb-then-hold}
If Console-Fino is an isomorphism for all
algebraic complex structures,
then it holds for all transcendental complex
structures too.
\end{lemma}

\begin{proof}
Let $I \in \text{Comp}(\mathfrak{g})$ be
trascendental.
Consider the variety
$\overline{\text{Aut}_{\mathbb{Q}}\mathbb{C}}
$,
which is algebraic and thus contains
algebraic points.
Then Console-Fino is an isomorphism on such
points.
Then it is so in a neighbourhood of this
point.
Then it is so in the orbit of this
neighbourhood.
\end{proof}

\noindent
Consider $(\mathfrak{g},I)$ with $I$ with
number field $k$.
Consider
$\tilde{\mathfrak{g}}=\bigoplus_{\sigma \in
\text{Gal}(k:\mathbb{Q})}
(\mathfrak{g},\sigma(I))$…
this constructions leads to a Lie algebra
$(\tilde{\mathfrak{g}},\tilde{I})$
equipped with a complex structure defined
over $\mathbb{Q}$.
This satisfies Console-Fino. We can lift
$(\tilde{M},\tilde{I})$
since it remains Dolbeault harmonic --- which
is the main tool used by
Colsole-Fino in their theorem.
This is how to proof Console-Fino for complex
structures over a closed field.

\section{Holomorphic symmetric differentials,
semiample bundles,
fundamental groups}
\label{seminars-complex-geometry-section-holom-symme-diffe-semia}

\noindent
Ernesto Mistretta, Università di Padova.
Algebra Seminar, IMPA.
November 12, 2025.

\medskip\noindent
{\bf Abstract.}
We will present some recent results and some
open conjectures relating
holomorphic symmetric differentials and the
geometry or topology of a compact
complex manifold, we will illustrate the
difference in the projective, or
Kaehler, or compact complex cases. Then we
will apply some elementary
constructions on semiample vector bundles to
partially answer some of these
questions.



\medskip\noindent



\noindent
As a motivation,
$$
\xymatrix{
&\text{differential forms}\ar[ld]\ar[rd]\\
\text{topology}&&\text{geometry}
}
$$
$X$ compact complex manifold, $E$
(holomorphic) vector bundle.
$L$ line bundle, the  {\it Kodaira-Itaka
dimension} is
$\text{KI}(X,L)=k$ if $h^0(X,L^{\otimes
m})=\mathcal{O}(m^k)$.

{\it Holomorphic $p$-differentials} are
elements of
$H^0(X,\text{Sym}^m\Omega^p_X)$ for some
$m>0$.

A conjecture by Mumford: suppose that $X$ is
a compact
Kähler manifold, then $X$ is rationally
connected if and only if
for every  $m>0$ and for every $p>0$,
$H^0(X,\text{Sym}^m\Omega^p_X)=0$.

Some relevant results are the following.

\begin{theorem}[Campana-Demailly-Peterwell,
2015]
\label{seminars-complex-geometry-theorem-cdp}
Suppose $X$ is a smooth and projective
variety over $\mathbb{C}$.
$X$ is rationally connected if and only if
for every
ample line bundle there exists a positive
constant $C_A>0$,
such that $\forall p>0$, $\forall K>0$ and
$\forall m>C_A K$,
$$
H^0(X,\text{Sym}^m\Omega^p_X \otimes A^K)=0.
$$
\end{theorem}

\noindent
Recall that if $X$ is a Kähler manifold and
rationally connected,
then $X$ is projective and $\pi_1(X)=0$.

\begin{theorem}[Brunarbe-Campana, 2016]
\label{seminars-complex-geometry-theorem-bc16}
Let $X$ be compact Kähler.
Suppose that $\forall m>0$, $\forall p>0$,
$H^0(X,\text{Sym}^m\Omega^p_X)=0$.
Then $X$ is projective and $\pi_1(X)=0$.
\end{theorem}

\noindent
What if we restrict to $p=1$?
$H^0(X,\text{Sym}^m\Omega^1_X)$?

A conjecture by Esnault says: let $X$ be a
compact Kähler manifold.
Suppose that $\forall m>0$,
$H^0(X,\text{Sym}^m\Omega^1_X)=0$.
Then $|\pi_1(X)|<\infty$.

\begin{theorem}[Burnabbe-Klinger-Totsaro]
\label{seminars-complex-geometry-theorem-bkt}
Suppose $X$ is a compact Kähler manifold.
Suppose that there exists a
finite-dimensional
representation of $\rho:\pi_1(X)\to
\text{GL}_r(k)$
over some field $k$ such that
$|\pi(\pi_1(X))|=\infty$.
Then there exists $m>0$ such that
 $H^0(X,\text{Sym}^m\Omega_X^1) \neq 0$.
\end{theorem}

\noindent
Notice the relationship between Theorem
\ref{seminars-complex-geometry-theorem-bkt}
and the conjecture by Esnault.

\medskip\noindent
Conjecturally, we have:
\begin{enumerate}
\item No holomorphic symmetric
$p$-differentials for all $p$
if and only if $X$ is rationally connected.
\item No holomorphic symmetric
1-differentials, them $|\pi_1(X)|<\infty$.
\end{enumerate}

\noindent
Question. Suppose there are some (many)
nonzero holomorphic symmetric
1-differentials, does this imply that
$|\pi_1|=\infty$?
What if $X$ is compact Kähler? And if $X$ is
projective?

Let's first answer the easiest of these
questions.
Let $X$ be compact Kähler with
$H^0(X,\Omega^1_X)\neq 0$.
By Hodge theorem,
$b_1(X)=\dim_\mathbb{Q}H^1(X,\mathbb{Q})
=2\dim_\mathbb{C}H^0(X,\Omega_X^1)$.
Also,
$H_1(X,\mathbb{Z})=\text{Ab}(\pi_1(X))$.
In particular, if $b_1 \neq 0$,
then
$H_1(X,\mathbb{Z})=\mathbb{Z}^{b_1}\oplus$
torsion.

\begin{theorem}[Brotbeck-Drondeau,2018]
\label{seminars-complex-geometry-theorem-brotbeck-drondeau18}
$X \subseteq \mathbb{P}^N$ very gen. complete
intersection
of high degree such that $\dim X\leq N/2$.
Then $\Omega_X^1$ is ample.
(In particular, $\text{Sym}^m\Omega_X^1$
is globally generated for $\gg m$.)
Now if $\dim X\geq 2$, then
$\pi_1(X)=\pi_1(\mathbb{P}^N)=0$.
\end{theorem}

\noindent
In this case $\omega_X=\det(\Omega_X^1)$ is
an ample line bundle.
So the Kodaira dimension of $X$ is
$\text{kod}(X)=\text{kod}(X,\omega_X)=\dim(X)
$.

\medskip\noindent
Now we explain joint work with Franccesco
Exposito from Padova.

First, a conjecture (by E-M).
Suppose $X$ is a compact complex manifold
with $\text{kod}(X)<\dim(X)$. If $\Omega_X^1$
is (weakly) semiample, then
$|\pi_1(X)|=\infty$.

The evidence is:
\begin{enumerate}
\item If $\dim X=1,2$, holds because of
classification.
\item If $X$ is projective, then the
conjecture follows from
a result of Hörng, which states that
$\Omega_X^1$ is semiample,
then there exists $X'\to X$ finite étale
cover
such that $X'\sim_{\text{bitat}}Y\times A$
where $A$
is an abelian variety and  $\dim
Y=\text{kod}(X)$.
\item If $X$ is Kähler, the conjecture (E-M)
follows
from a conjecture of Wu-Zheng from 2002,
which says that
if  $\Omega^1_X$ is (NEF) semiample (which
implies NEF), then
there exists $X' \to X$ finite étale that
admits a smooth fibration
$X' \to Y$ in complex tori with $\dim
Y=\text{kod}X$.
In this case the fundamental groups of the
tori are injected
in the fundamental group of $X$, i.e.,
$\pi_1(F) \hookrightarrow \pi_1(X)$.
(It appears that this is proved via the Serre
homotopy
long exact sequence, though the Kähler
hypothesis may needed, and
it's not clear how…)
This makes $|\pi_1(X)|=\infty$.

\item If $X$ is a compact complex manifold
and $\text{kod}(X)=0$,
and $\Omega_X^1$ is (weakly) semiample, then
$|\pi_1(X)|=\infty$.
\end{enumerate}

\noindent

\begin{definition}
\label{seminars-complex-geometry-definition-duplicate-3}
Let $E$ be a vector bundle on $X$, a compact
complex manifold.
Recall the construction
$$
\xymatrix{
\mathbb{P}(E)\ar[d]^{\pi}\ar@{-}[r]&\pi^*
E\ar@{->>}[r]
&\mathcal{O}_{\mathbb{P}(E)}(1)\\
X\ar@{-}[r]&E
}
$$
$$
\pi_*(\mathcal{O}_{\mathbb{P}(E)}(m))=
\text{Sym}^m(E).
$$
\begin{enumerate}
\item $E$ is {\it strongly semiample} if
there exists $m>0$
such that $\text{Sym}^mE$ is globally
generated.
\item $E$ is {\it weakly semiample} if
$\mathcal{O}_{\mathbb{P}(E)}(1)$ is semiample
line bundle
(i.e. $\mathcal{O}_{\mathbb{P}(E)}(m)$ is
globally generated
for some $m\gg 0$).
\end{enumerate}
\end{definition}

\noindent
Notice that these two definitions are not
equivalent.
Strongly implies weakly, but not conversely.
To see that strongly implies weakly, consider
the pullback
$$
\xymatrix{
\pi^*\text{Sym}^nE=\text{Sym}^n\pi^*E\ar[r]
\ar[d]&\mathbb{P}(E)\ar[d]^{\pi}\\
\text{Sym}^mE\ar[r]&X
}
$$
then
 $$
\text{Sym}^m\pi^* E
\to\!\!\!\!\!\to\mathcal{O}
_{\mathbb{P}(E)}(m).
$$
To see that weakly does not imply strongly,
suppose $X$ is a smooth projective
surface with $\text{kod}(X)=0$.
Consider $E=\Omega_X^1$ semiample. We have
that $\Omega_X^1$
is strongly semiample if and only if $X$ is
an abelian surface.
Also, if $X$ is a hyperelliptic surface, then
$\Omega_X^1$
is weakly semiample.
This means that $X$ hyperelleptic surface is
the counterexample.

\begin{theorem}[E-M]
\label{seminars-complex-geometry-theorem-em}
Suppose that $X$ is compact complex and that
$\text{kod}(X)=0$. Then
\begin{enumerate}
\item if $\Omega_X^1$ is strongly semiample
then
$X$ is a parallelalizable complex manifold.
\item If $\Omega_X^1$ is weakly semiample,
then $X$ is a finite étale Galois quotient
$X\cong P/G$
of a parallelizable manifold $P$.
\end{enumerate}
\end{theorem}

\begin{definition}
\label{seminars-complex-geometry-definition-paralellizable}
$X$ is {\it parallelizable} is
$\Omega_X^1\cong\mathcal{O}_X$.
\end{definition}

\noindent
A $X$ is a compact complex manifold is
parallelizable if and only if
$X \cong H/\Gamma$ for $H$ a complex Lie
group and $\Gamma\subset H$
a discrete subgroup.
Further, it is Kähler if and only if $H$ is a
complex torus.

\section{Hyperkahler manifolds with 
finite group of birational automorphisms}
\label{seminars-complex-geometry-section-hk-misha-26}

\noindent
Misha Verbitsky, 20 May, 2026. IMPA.

\medskip\noindent {\bf Abstract.} Seja $M$ uma
coletânea hyperkahler, $b_2(M)\ge 6$, e $D$ seu
espaço de deformação, $I\in D$ uma estrutura
complexa e $H_I$ a rede Hodge de $(M,I)$, ou seja,
o conjunto de todas as classes de cohomologia
inteiras do tipo $(1,1)$. Usando o teorema de
Vinberg sobre redes geradas por reflexões,
provamos que há apenas um número finito de
classes de isomorfismo de redes de Hodge $H_I$
de grau $\ge 3$, de modo que $(M,I)$ seja projetivo
e tenha um grupo finito de automorfismos
birracionais. Usamos essa observação para
mostrar que qualquer família não trivial de
deformações projetivas de $M$ tem um conjunto
denso de estruturas complexas com um grupo
infinito de automorfismos birracionais. Este
é um trabalho conjunto com Ekaterina Amerik e
Andrey Soldatenkov.  


\medskip\noindent
{\bf Upshot.}
The main result we shall prove is the
following.
Let $M$ be a K3 surface and let $\mathcal{H}$
be its polarized moduli space.
There exist finitely many points
$z_1,\dots,z_n \in \mathcal{H}$
such that for any
$m \in \mathcal{H}\setminus \bigcup_{i=1}^n \{z_i\}$
with $\rho(M)\leq 3$, one has
$\lvert \operatorname{Aut}(M)\rvert=\infty$.

In plain language: for most members of the
family, once the Picard rank is small enough,
the surface has infinitely many
automorphisms in ``almost every case'';
only finitely many special
parameter values have infinite automorphism
group.

My favorite idea from the talk is the
following, which I think is called
Noether-Lefschetz result:

\begin{quotation}
  ``Torelli tells us that there are roughly as many
K3 surfaces as Hodge decompositions. Fix a
polarization $p \in H^2(M,\mathbb{Z})$ with
$p^2>0$. There is a group acting on
$\mathbb{P}\operatorname{er}$, namely
$\operatorname{SO}(H^2(M,\mathbb{Z}),p)$, and the
quotient is quasiprojective. This is precisely the
moduli space of polarized K3 surfaces.

\noindent
Now suppose you have a sublattice
$\Lambda \subset H^2(M,\mathbb{Z})$. The
Noether-Lefschetz component is
\[
H_{\Lambda}
\;=\;
\{\,M \in \mathbb{P}\operatorname{er}/\operatorname{SO}
\mid \Lambda \subset H^{1,1}(M)\,\}.
\]
The classes in such quotient correspond
with the K3 that have finite automorphism
group!"
\end{quotation}


\medskip\noindent
The talk starts reviewing several notions
and results like:

\begin{theorem}
A K3 surface $M$ is projective if and only if
there exists a class $\lambda \in H^{1,1}(M,\mathbb{Z})$
such that $\int_M \lambda \wedge \lambda > 0$.
\end{theorem}

\begin{remark}
From the Riemann-Roch formula it follows that
any $(-2)$-class $\lambda$ is either effective or
its negative $-\lambda$ is effective, that is,
equal to the fundamental class of a curve.
\end{remark}

\begin{remark}
This implies that a class in the orthogonal
complement of a $(-2)$-class cannot be Kähler.
\end{remark}

\noindent
The follwing theorem uses the Demailly-Paun
characterization of Kählerness:

\begin{theorem}
Let $M$ be a K3 surface, and $S \subset
H^{1,1}(M,\mathbb{Z})$ the set of all
$(-2)$-classes. Consider the decomposition of
$\operatorname{Pos}(M) \setminus
\bigcup_{s\in S} s^\perp$ into connected
components. Then the Kähler cone of $M$ is one
of the connected components.
\end{theorem}

\begin{theorem}
Let $M$ be a K3 surface. The group
$\operatorname{Aut}(M)$ of holomorphic
isometries of $M$ is naturally identified with
the subgroup of $O(H^2(M,\mathbb{Z}))$ which
preserves the Kähler cone.
\end{theorem}

\begin{definition}
A group is called virtually free abelian if it
contains a free abelian subgroup of finite
index.
\end{definition}

\begin{proposition}[Oguiso]
Let $M$ be a K3 surface, and $H^{1,1}(M,\mathbb{Z})$
its Neron-Severi lattice. Then the following
alternatives occur.
\begin{enumerate}
\item[(a)] If the intersection form on
$H^{1,1}(M,\mathbb{Z})$ is negative definite, then
the kernel of $\Psi: \operatorname{Aut}(M)\to
\operatorname{Aut}(H^{2,0}(M))=\mathbb{C}^*$ is finite,
and the image is virtually free abelian.
\item[(b)] If the intersection form on
$H^{1,1}(M,\mathbb{Z})$ is degenerate, the image and
the kernel of the map $\Psi$ are virtually free
abelian.
\item[(c)] If $M$ is projective, the image of $\Psi$
is finite.
\end{enumerate}
\end{proposition}

\begin{lemma}
Let $M$ be a projective K3 surface with Picard
rank $1$. Then $\operatorname{Aut}(M)$ is finite.
\end{lemma}

\begin{proof}
Indeed, for a projective K3, the map
$\operatorname{Aut}(M)\to O(H^{1,1}(M,\mathbb{Z}))$
has finite kernel (Proposition 1).
\end{proof}

\begin{definition}
Let $(\Lambda,q)$ be a quadratic lattice, that is,
$\mathbb{Z}^n$ equipped with an integer-valued
quadratic form. We say that $(\Lambda,q)$
represents $n\in \mathbb{Z}$ if there exists
$x\in \Lambda$ such that $q(x,x)=n$.
\end{definition}

\begin{lemma}
Let $M$ be a projective K3 surface with Picard
rank $2$. Then $\operatorname{Aut}(M)$ is finite if
its Neron-Severi lattice represents zero, and
contains an infinite cyclic subgroup of finite
index otherwise.
\end{lemma}

\begin{proof}
Since $M$ is projective, $\operatorname{Aut}(M)$ is
finite if and only if its group $\operatorname{Aut}_{\Omega}(M)\subset
\operatorname{SO}(H^{1,1}(M,\mathbb{Z}))$ of symplectic
automorphisms is finite (Proposition 1). The group
of isometries of a quadratic lattice $(\Lambda,q)$
of signature $(1,1)$ is finite if $(\Lambda,q)$
represents zero. It contains an infinite cyclic
subgroup of finite index otherwise, by Pell's
classification of solutions of an equation
$x^2-ay^2=1$.
\end{proof}

\noindent
It looks like a ``divisorial
reflection'' is just a reflection
w.r.t. a  $(-2)$-curve.

\begin{definition}
Let $v \in H^{1,1}(M,\mathbb{Z})$ be a $(-2)$-class.
Consider the map
\[
x \mapsto x-\frac{2(x,v)}{(v,v)}v,
\]
where $(\cdot,\cdot)$ is the intersection form.
This map is called the divisorial reflection.
\end{definition}

\begin{lemma}
The group generated by divisorial reflections acts
transitively on the set of connected components of
\[
\operatorname{Pos}(M)\setminus \bigcup_{s\in S}s^\perp,
\]
where $S$ is the set of all $(-2)$-classes.
\end{lemma}

\noindent
We denote by $\operatorname{Ref}(M)$ the subgroup
of $\operatorname{SO}(H^2(M,\mathbb{Z}))$
generated by divisorial reflections.

\noindent
Clearly, it is a normal subgroup in the group of
all Hodge isometries of $H^2(M,\mathbb{Z})$.

\begin{lemma}
The group $\operatorname{Aut}_{\Omega}(M)$ of
symplectic automorphisms of a K3 surface is
determined, up to an isomorphism, by its
Neron-Severi quadratic lattice.
\end{lemma}

\begin{proof}
It is identified with the subgroup
$\operatorname{St}(\operatorname{Kah}(M))$ of
$\operatorname{SO}(H^{1,1}(M,\mathbb{Z}))$
preserving the Kähler cone. However,
$\operatorname{Ref}(M)$ acts on
$\operatorname{SO}(H^2(M,\mathbb{Z}))$ by Hodge
isometries, transitively on all connected
components of
$\operatorname{Pos}(M)\setminus
\bigcup_{s\in S}s^\perp$.
Therefore, the action of $\operatorname{Ref}(M)$ is
transitive on all groups of form $\operatorname{St}(A)$,
where $A$ is a connected component in
$\operatorname{Pos}(M)\setminus
\bigcup_{s\in S}s^\perp$, and $\operatorname{St}(A)$
its stabilizer.
\end{proof}

\begin{remark}
The group $O(m,n)$, $m,n > 0$, has $4$ connected
components. We denote the connected component of $1$
by $SO^+(m,n)$. We call a vector $v$ positive if its
square is positive.
\end{remark}

\begin{definition}
Let $V$ be a vector space with quadratic form $q$
of signature $(1,n)$, $\operatorname{Pos}(V)=\{x\in
V \mid q(x,x)>0\}$ its positive cone, and
$\mathbb{P}V$ projectivization of $\operatorname{Pos}(V)$.
Denote by $g$ any $SO(V)$-invariant Riemannian
structure on $\mathbb{P}V$. Then $(\mathbb{P}V,g)$ is
called hyperbolic space, and the group $SO^+(V)$ the
group of oriented hyperbolic isometries.
\end{definition}

\begin{example}
For any K3 surface, $\mathbb{P}\operatorname{Pos}(M)$
is a hyperbolic space, and $\operatorname{Aut}(M)$
acts on this space by isometries.
\end{example}

\begin{theorem}[Sterk, 1985]
Let $M$ be a projective K3 surface. Then
$\operatorname{Aut}(M)$ acts with finitely many orbits
on the set of the walls of the Kähler cone.
\end{theorem}

\noindent
The proof may be found in a book by Kneser.

\medskip\noindent
When $M$ is a Calabi-Yau manifold, Sterk's
theorem is known as the Morrison-Kawamata
conjecture, formulated by Morrison in 1993,
and proven by Kawamata for elliptic Calabi-Yau
manifolds in 1997. It was proven by B. Totaro
for log-Calabi-Yau surfaces (2010), and by
Amerik-V. for hyperkähler manifolds (2015).

\begin{remark}
This is an AI-generated explanatory remark.
The ample cone is usually described by ample line
bundles, so its natural generators have integral
first Chern classes. Rational classes appear here
because one often passes to the rational span of
the cone when intersecting with the Kähler cone:
on a projective K3 surface, ample classes are
exactly the rational Kähler classes.
\end{remark}

\begin{lemma}
Let $\operatorname{Amp}(M):=\operatorname{Kah}(M)\cap H^{1,1}(M,\mathbb{Q})$
denote the ample cone of a projective K3
surface, and $\operatorname{Amp}_{\mathbb{R}}(M):=
\operatorname{Kah}(M)\cap H^{1,1}(M,\mathbb{Q})\otimes_{\mathbb{Q}}\mathbb{R}$
denote the corresponding cone in
$H^{1,1}(M,\mathbb{Q})\otimes_{\mathbb{Q}}\mathbb{R}$.
Consider its projectivization
$\mathbb{P}\operatorname{Amp}_{\mathbb{R}}(M)$ as a
subset in the hyperbolic space $\mathbb{P}\operatorname{Pos}(M)$.
Then $\mathbb{P}\operatorname{Amp}_{\mathbb{R}}(M)$
has finite volume if and only if
$\operatorname{Aut}(M)$ is finite.
\end{lemma}

\begin{proof}
By construction,
$\mathbb{P}\operatorname{Amp}_{\mathbb{R}}(M)$ is a
hyperbolic polyhedron (with some vertices on
absolute). If $\operatorname{Aut}(M)$ is finite, this
polyhedron has finite number of faces, hence its
volume is finite. Conversely, if $\operatorname{Aut}(M)$
is infinite, $\mathbb{P}\operatorname{Amp}_{\mathbb{R}}(M)$
contains infinite number of copies of its fundamental
domain, hence its volume is infinite.
\end{proof}

\begin{theorem}[Vinberg]
Let $M$ be a projective K3 surface of
Picard rank $\geq 3$. Then $\operatorname{Aut}(M)$
is finite if and only if the reflection group
$\operatorname{Ref}(M)$ has finite index in
$\operatorname{SO}(H^{1,1}(M,\mathbb{Z}))$.
\end{theorem}

\begin{proof}
Step 1: If $\operatorname{Aut}(M)$ is finite, the
ample cone has finite volume, hence the action of
$\operatorname{Ref}(M)$ on the hyperbolic space
$\mathbb{P}\operatorname{Amp}_{\mathbb{R}}(M)$ has
finite volume. This implies that $\operatorname{Ref}(M)$
is a lattice subgroup in
$\operatorname{SO}(H^{1,1}(M,\mathbb{R}))$, hence has
finite index in $\operatorname{SO}(H^{1,1}(M,\mathbb{Z}))$.

Step 2: Conversely, assume that $\operatorname{Ref}(M)$
has finite index in $\operatorname{SO}(H^{1,1}(M,\mathbb{Z}))$.
Then its Coxeter polyhedron is finite (Vinberg).
However, $\mathbb{P}\operatorname{Amp}_{\mathbb{R}}(M)$
is its Coxeter polyhedron.
\end{proof}

\begin{definition}
\label{seminars-complex-geometry-definition-reflective}
We call a quadratic lattice with
$\operatorname{Ref}(\Lambda)$ of finite index in
$\operatorname{SO}(\Lambda)$ reflective.
\end{definition}

\noindent
Torelli tells us that there are roughly as many
K3 surfaces as Hodge decompositions. Fix a
polarization $p \in H^2(M,\mathbb{Z})$ with
$p^2>0$. There is a group acting on
$\mathbb{P}\operatorname{Per}$, namely
$\operatorname{SO}(H^2(M,\mathbb{Z}),p)$, and the
quotient is quasiprojective. This is precisely the
moduli space of polarized K3 surfaces.

\noindent
Now suppose you have a sublattice
$\Lambda \subset H^2(M,\mathbb{Z})$. The
Noether-Lefschetz component is
\[
H_{\Lambda}
\;=\;
\{\,M \in \mathbb{P}\operatorname{Per}/\operatorname{SO}
\mid \Lambda \subset H^{1,1}(M)\,\}.
\]
The classes in such quotient correspond
with the K3 that have finite automorphism
group!

\begin{lemma}
\label{seminars-complex-geometry-lemma-noether-lefschetz-components}
Let $H$ be the moduli space of polarized K3
surfaces. Then there is a finite set of
quasiprojective subvarieties
$Z_1,\dots,Z_n \subset H$ such that for any K3
surface $M$ outside of $\bigcup_i Z_i$ with Picard
rank $\geq 3$, the group $\operatorname{Aut}(M)$ is
infinite.
\end{lemma}

\begin{proof}
Step 1: Fix a polarization on $M$, that is, an
integer vector $p \in H^2(M,\mathbb{Z})$ such that
its square is positive. By Torelli theorem,
$H=\mathbb{P}\operatorname{Per}/\operatorname{SO}(H^2(M,\mathbb{Z}),p)$,
where $\operatorname{Per}$ is the space of all Hodge
structures on $H^2(M)$ such that $p$ has type
(1,1), and $\operatorname{SO}(H^2(M,\mathbb{Z}),p)$
the group of isometries preserving $p$. For any
sublattice $\Lambda \subset H^2(M,\mathbb{Z})$
containing $p$, the corresponding
Noether-Lefschetz component $H_{\Lambda}\subset H$
is the image of all $I \in \operatorname{Per}$ such
that $H^{1,1}(M,I)\subset \Lambda$. By Baily-Borel
theorem, $H_{\Lambda}$ is a quasiprojective
subvariety in $H$.

Step 2: Let $Z_i$ be the set of Noether-Lefschetz
components associated with reflective lattices
$\Lambda \subset H^2(M,\mathbb{Z})$. Clearly, any
$M$ which has Picard rank $\geq 3$ and is not in
$Z_i$ has infinite automorphism group.
\end{proof}

\begin{theorem}[K. Oguiso, 2003]
\label{seminars-complex-geometry-theorem-oguiso-2003}
Let $M_t$ be a non-trivial quasiprojective family
of K3 surfaces over a curve $C$, with general
member of Picard rank $\geq 3$. Then
$|\operatorname{Aut}(M_t)|=\infty$ for a dense set
of $t \in C$.
\end{theorem}

\begin{proof}
Remove from $C$ all proper subsets which belong to
the Noether-Lefschetz components $Z_i$. If
$C \not\subset Z_j$, we are done. Otherwise, there
is a dense subset of points in $C$ where the Picard
rank jumps, and these points are not in $Z_i$.
\end{proof}

\begin{remark}
\label{seminars-complex-geometry-remark-oguiso-lower-picard}
This theorem is also valid for families with
general member of Picard rank $\geq 1$, but the
proof for lower Picard values is more elementary.
\end{remark}

\begin{remark}
\label{seminars-complex-geometry-remark-oguiso-vinberg}
Oguiso's proof does not use Vinberg's theorem, he
uses explicit lattice-theoretic arguments.
\end{remark}

\begin{theorem}[Denisi, Onorati, Rizzo, Viktorova]
\label{seminars-complex-geometry-theorem-denisi-onorati-rizzo-viktorova}
Let $M_t$ be a non-trivial quasiprojective family
of hyperkähler manifolds of known type over a
curve $C$. Then $|\operatorname{Bir}(M_t)|=\infty$
for a dense set of $t \in C$, where
$\operatorname{Bir}(M_t)$ denotes the group of
birational automorphisms.
\end{theorem}

\begin{remark}
\label{seminars-complex-geometry-remark-dorv-lattices}
Their proof relies on known properties of the
cohomology lattices of hyperkähler manifolds.
\end{remark}

\begin{theorem}[Amerik, Soldatenkov, V.]
\label{seminars-complex-geometry-theorem-amerik-soldatenkov-v}
This theorem holds for any family of hyperkähler
manifolds.
\end{theorem}

\begin{remark}
\label{seminars-complex-geometry-remark-asv-vinberg}
Our proof is based on Vinberg's theorem, in fact,
it is pretty much identical to the proof for K3
surfaces that I gave.
\end{remark}

\section{Multiplicity of the fibers in
Lagrangian fibrations}
\label{seminars-complex-geometry-section-multiplicity-fibers-lagrangian-fibrations}

\noindent
Misha Verbitsky.
Ter 02 jun 2026, 15:30.

\medskip\noindent
Seja \(M\) um coletor hyperkahler de
holonomia m\'axima e
\(\pi : M \to \mathbf{P}^n\) uma fibração
lagrangiana. Uma fibra \(F\) de \(\pi\) \'
e m\'ultipla de multiplicidade \(p\) se um
pequeno disco transversal a \(F\) encontrar
uma fibra geral de \(\pi\) em \(p > 1\)
pontos. Provarei que as fibras m\'ultiplas
n\~ao ocorrem em codimens\~ao 1 se e
somente se o pullback de uma se\c{c}\~ao de
hiperplano for primitivo (indivis\'ivel) na
cohomologia de \(M\). Em seguida,
explicarei como usar essa observa\c{c}\~ao
para provar que n\~ao h\'a fibras m\'ultiplas
em codimens\~ao 1. Os argumentos se
baseiam em uma vers\~ao forte do
desaparecimento de Kodaira-Nakano, v\'alido
para feixes de linhas com curvatura
semipositiva, e na estabilidade do feixe
tangente. Este \'{e} um trabalho conjunto
com Ljudmila Kamenova.

\medskip\noindent
It's a mystery why any elliptic fibration
over a K3 surface has no multiple fibers.
We will explain how to generalize this result
to higher dimensions. By some
theorem, elliptic fibrations shall
be Lagrangian fibrations.

\medskip\noindent

\begin{theorem}[Calabi--Yau]
\label{seminars-complex-geometry-theorem-calabi-yau}
A compact, K\"ahler, holomorphically
symplectic manifold admits a unique
hyperk\"ahler metric in any K\"ahler class.
\end{theorem}

\begin{definition}
\label{seminars-complex-geometry-definition-hyperkahler-manifold}
For the rest of this talk, a hyperk\"ahler
manifold is a compact, K\"ahler,
holomorphically symplectic manifold.
\end{definition}

\begin{definition}
\label{seminars-complex-geometry-definition-maximal-holonomy}
A hyperk\"ahler manifold \(M\) is called
maximal holonomy, or IHS, if
\(\pi_1(M)=0\) and \(H^{2,0}(M)=\mathbf{C}\).
\end{definition}

\medskip\noindent
\textbf{Bogomolov's decomposition.}
Any hyperk\"ahler manifold admits a finite
covering which is a product of a torus and
several hyperk\"ahler manifolds of maximal
holonomy.

\medskip\noindent
Further on, all hyperk\"ahler manifolds are
assumed to be of maximal holonomy.

\begin{theorem}[Matsushita]
\label{seminars-complex-geometry-theorem-matsushita-lagrangian}
Let \(M\) be a hyperk\"ahler manifold of
maximal holonomy, and let
\(\pi : M \to X\) be a surjective
holomorphic map, with \(0 < \dim X < \dim M\).
Then \(\pi\) is a Lagrangian fibration
(that is, it has holomorphic Lagrangian
fibers).
\end{theorem}

\begin{theorem}[Hwang]
\label{seminars-complex-geometry-theorem-hwang-smooth-base}
In these assumptions, \(X\) is biholomorphic
to \(\mathbf{C}P^n\) when it is smooth.
\end{theorem}

\medskip\noindent
The conjecture is that \(X\) is biholomorphic
to \(\mathbf{C}P^n\) when it is normal.

\begin{theorem}[Matsushita]
\label{seminars-complex-geometry-theorem-matsushita-normal-base}
Let \(M\) be a hyperk\"ahler manifold of
maximal holonomy, and let
\(\pi : M \to X\) be a Lagrangian fibration,
with \(X\) normal. Then
\(H^*(X,\mathbf{Q}) \cong H^*(\mathbf{C}P^n,
\mathbf{Q})\).
\end{theorem}

\begin{remark}
\label{seminars-complex-geometry-remark-general-fibers-abelian}
General fibers of \(\pi\) are Abelian
varieties (projective complex tori), by
Arnold-Liouville. Conversely, as shown by
Hwang-Weiss, any Lagrangian complex torus in
\(M\) is a fiber of a Lagrangian fibration.
\end{remark}

\begin{remark}
\label{seminars-complex-geometry-remark-fibration}
A fibration \(\pi : M \to X\) is a proper
surjective holomorphic map of complex
manifolds. We shall always tacitly assume
that the fibration is proper and
equidimensional, that is, the irreducible
components of all fibers have dimension
\(\dim M - \dim X\). A special fiber of
\(\pi\) is a fiber containing a critical
point.
\end{remark}

\begin{definition}
\label{seminars-complex-geometry-definition-multiplicity-special-fiber}
Let \(Z\) be an irreducible component of a
special fiber of a fibration
\(\pi : M \to X\), and \(z=\pi(Z)\). The
multiplicity of \(Z\) is the rank of
\(\pi^*(\mathcal{O}_X/\mathfrak{m}_z)\) in a
general point of \(Z\), where
\(\mathfrak{m}_z \subset \mathcal{O}_X\) is
the maximal ideal of \(z\).
\end{definition}

\begin{remark}
\label{seminars-complex-geometry-remark-multiplicity-one-dimensional}
Suppose that \(\dim X=1\). Locally in \(M\),
around a general point of \(m \in Z\), there
is a coordinate system \(x_1,\ldots,x_n\),
such that \(\pi(x_1,\ldots,x_n)=x_1^d\).
In this case, \(d\) is multiplicity of
\(Z\).
\end{remark}

\begin{remark}
\label{seminars-complex-geometry-remark-equivalent-definition}
Another definition of multiplicity is more
geometric, but it is equivalent to the one
given above; the equivalence is left as an
exercise.
\end{remark}

\begin{definition}
\label{seminars-complex-geometry-definition-multiplicity-transversal-disk}
Let \(Z\) be an irreducible component of a
special fiber of a fibration
\(\pi : M \to X\), and \(z=\pi(Z)\).
Consider a small disc \(D \subset M\) of
dimension \(\dim X\) transversal to \(Z\)
and intersecting it in \(m\). The
multiplicity of \(Z\) in \(m \in M\) is the
number of intersection points between \(D\)
and a general fiber of \(\pi\) which is
sufficiently close to \(Z\).
\end{definition}

\begin{theorem}[Campana-Kamenova-Verbitsky]
\label{seminars-complex-geometry-theorem-campana-kamenova-verbitsky}
Let \(\pi : M \to X\) be an Abelian
fibration (that is, a fibration with general
fiber a complex torus), and assume that
\(M\) is K\"ahler. Let \(\mu_i\) be the
multiplicity of irreducible components of
\(\pi^{-1}(z)\) for some \(z \in X\), and
\(\gcd(\mu_i)\) their greatest common
divisor. Then \(\gcd(\mu_i)=\min \mu_i\).
\end{theorem}

\begin{definition}
\label{seminars-complex-geometry-definition-discriminant}
Let \(\pi : M \to X\) be a surjective
holomorphic map of complex manifolds, and
\(D \subset X\) its set of critical values,
which is known as the discriminant, or the
discriminant divisor (it has codimension
1). We say that \(\pi\) has no multiple
fibers in codimension 1 if for a general
point \(x \in D\), the fiber \(\pi^{-1}(x)\)
has a component with multiplicity 1.
\end{definition}

\begin{theorem}
\label{seminars-complex-geometry-theorem-elliptic-k3-no-multiple-fibers}
Let \(\pi : M \to X\) be an elliptic
fibration on a K3 surface. Then \(\pi\) has
no multiple fibers.
\end{theorem}

\medskip\noindent
The main result today is a generalization of
this theorem.

\begin{theorem}
\label{seminars-complex-geometry-theorem-main-no-multiple-fibers}
Let \(\pi : M \to X\) be a Lagrangian
fibration on a hyperk\"ahler manifold. Then
\(\pi\) has no multiple fibers in
codimension 1.
\end{theorem}

\begin{lemma}
\label{seminars-complex-geometry-lemma-h2-torsion-free}
Let \(M\) be a hyperk\"ahler manifold of
maximal holonomy. Then \(H^2(M)\) is
torsion-free.
\end{lemma}

\begin{proof}
The universal coefficients formula gives the
exact sequence:
$$
0 \to \operatorname{Ext}_{\mathbb{Z}}^1
\bigl(H_1(M;\mathbb{Z}),\mathbb{Z}\bigr)
\to H^2(M;\mathbb{Z}) \to
\operatorname{Hom}_{\mathbb{Z}}
\bigl(H_2(M;\mathbb{Z}),\mathbb{Z}\bigr)
\to 0.
$$
Since \(H_1(M,\mathbb{Z})=0\) for a maximal
holonomy hyperk\"ahler manifold, this gives
an isomorphism
\[
H^2(M,\mathbb{Z}) =
\operatorname{Hom}_{\mathbb{Z}}
\bigl(H_2(M;\mathbb{Z}),\mathbb{Z}\bigr),
\]
hence the torsion vanishes.
\end{proof}

\begin{lemma}
\label{seminars-complex-geometry-lemma-primitive-classes-lagrangian}
Let \(\pi : M \to \mathbf{C}P^n\) be a
Lagrangian fibration on a hyperk\"ahler
manifold of maximal holonomy, and
let \(H \subset \mathbf{C}P^n\) be a
hyperplane section. Then the following
assertions are equivalent.
\begin{itemize}
\item The homology class of
\(\pi^{-1}(H)\) is primitive.
\item The map \(\pi\) has no multiple
fibers in codimension 1.
\end{itemize}
\end{lemma}

\begin{proof}
Step 1: The implication (ii) \(\Rightarrow\)
(i) is proven later today using the ETMDPS
vanishing theorem. The converse implication:
let \(D_1\) be an irreducible component of
the discriminant \(D\); assume it has
multiplicity \(\mu\). We need to show that
the preimage of the hyperplane section \(H\)
is not primitive.

Step 2: Arguing ad absurdum, assume that
\(D_1\) is homologous to \(k[H]\), but
\([\pi^{-1}(H)]\) is primitive in
\(H^2(M,\mathbb{Z})\). Since
\([D_1] = k[H]\), \(\pi^{-1}(D_1)\) is the
zero set of a section
\(s \in H^0(\mathbf{C}P^n,\mathcal{O}(k))\).
The pullback of this section has multiplicity
\(\mu\), since \(\pi^*(\mathcal{O}(1))\) is
primitive, \(k\) is divisible by \(\mu\).
The universal coefficients formula implies
that \(H^2(M,\mathbb{Z})\) is torsion-free.
Therefore, \(\pi^*s = s_0^\mu\), where
\(s_0 \in H^0(M,\pi^*\mathcal{O}(k/\mu))\).
Since \(\pi^*\mathcal{O}(k/\mu)\) is
trivial on fibers of \(\pi\), the section
\(s_0\) is a pullback of a section
\(s_1 \in H^0(\mathbf{C}P^n,\mathcal{O}(k/\mu))\).
This is impossible, because \(s\) and hence
\(s_1\) vanishes on \(D_1\), which has
degree \(k\).
\end{proof}

\begin{theorem}[Fujiki]
\label{seminars-complex-geometry-theorem-fujiki}
Let \(\eta \in H^2(M)\), and \(\dim M = 2n\),
where \(M\) is hyperk\"ahler. Then
\[
\int_M \eta^{2n} = cq(\eta,\eta)^n,
\]
for some primitive integer quadratic form
\(q\) on \(H^2(M,\mathbb{Z})\), and
\(c>0\) a rational number.
\end{theorem}

\begin{definition}
\label{seminars-complex-geometry-definition-bbf-form}
This form is called the Bogomolov-Beauville-
Fujiki form. It is defined by the Fujiki's
relation uniquely, up to a sign. The sign is
determined from the following formula
\[
\lambda q(\eta,\eta)=
\int_X \eta \wedge \eta \wedge \Omega^{n-1}
\wedge \overline{\Omega}^{n-1}
- \frac{n-1}{2n}
\left(
\int_X \eta \wedge \Omega^{n-1}
\wedge \overline{\Omega}^{n}
\right)
\left(
\int_X \eta \wedge \Omega^{n}
\wedge \overline{\Omega}^{n-1}
\right)
\]
where \(\Omega\) is the holomorphic symplectic
form, and \(\lambda>0\).
\end{definition}

\begin{remark}
\label{seminars-complex-geometry-remark-bbf-signature}
\(q\) has signature \((3,b_2-3)\). It is
negative definite on primitive forms, and
positive definite on \(\langle \Omega,
\overline{\Omega}, \omega \rangle\), where
\(\omega\) is a K\"ahler form.
\end{remark}

\begin{definition}
\label{seminars-complex-geometry-definition-holomorphic-euler}
Let \(B\) be a holomorphic vector bundle (or
a coherent sheaf). The holomorphic Euler
characteristic is
\[
\chi(B) := \sum_i (-1)^i H^i(M,B).
\]
\end{definition}

\begin{theorem}[Riemann-Roch-Hirzebruch]
\label{seminars-complex-geometry-theorem-rriemann-roch-hirzebruch}
Let \(M\) be a compact complex manifold, and
\(B\) a holomorphic vector bundle. Then
\(\chi(B)\) can be expressed through the
the Chern classes of \(TM\) and \(B\),
\[
\chi(B)=\int_M td(TM)\wedge ch(B)
\]
where \(td\) is the the Todd polynomial on
Chern classes of \(TM\),
\[
td(M)=1+\frac{1}{2}c_1+\frac{1}{12}(c_1^2+c_2)
+\frac{1}{24}c_1c_2+\frac{1}{720}
(-c_1^4+4c_1^2c_2+c_1c_3+3c_2^2-c_4)+\ldots
\]
and \(ch(B)\) its Chern character,
\[
ch(B)=1+c_1+\frac{1}{2}(c_1^2-2c_2)
+\frac{1}{6}(c_1^3-3c_1c_2+3c_3)+\ldots
\]
\end{theorem}

\noindent
The idea of the following theorem
is to look at $\text{SO}(H^2)$.
The function $\chi(L)$ is an isometry
of BBF form, which makes it a polynomial
on $q(v,v)$. Also, isometries of $H^2$
are correspond with automorphisms of $M$.

\begin{theorem}[Huybrechts]
\label{seminars-complex-geometry-theorem-huybrechts}
Let \(M\) be a hyperk\"ahler manifold,
\(\dim_{\mathbf{C}} M = 2n\), and \(L\) a
holomorphic line bundle. Then
\[
\chi(L)=\sum a_i q(c_1(L))^i,
\]
where the coefficients \(a_i\) are constants
depending on the topology of \(M\).
\end{theorem}

\begin{proof}
Step 1: Let \(A^*\) be the subalgebra in
cohomology generated by \(H^2(M)\). Then
\[
A^{2i} \cong \operatorname{Sym}^i(H^2(M))
\]
up to the middle degree, and
\[
A^{n+i} \cong \operatorname{Sym}^{n-i}
\bigl(H^2(M)\bigr);
\]
there is an \(O(H^2(M))\)-action on
cohomology, and the multiplication is
\(O(H^2(M))\)-invariant (V., 1995).

Step 2: All Chern classes of \(TM\) are
\(O(H^2(M))\)-invariant, but there is only
one (up to a constant multiplier)
\(O(H^2(M))\)-invariant functional on
\(\operatorname{Sym}^{2i}(H^2(M))\). On the
class \(\eta^{2i} \in H^{4i}(M)\) this
functional takes value \(q(\eta,\eta)^i\).
Therefore, all \(L\)-dependent coefficients
in the Hirzebruch-Riemann-Roch formula for
\(\chi(L)\) are expressed through
\(q(c_1(L))\).
\end{proof}

\begin{lemma}[Corollary]
\label{seminars-complex-geometry-corollary-huybrechts-chi}
Let \(L\) be a line bundle on a hyperk\"ahler
manifold \(M\), \(\dim_{\mathbf{C}} M = 2n\).
Assume that \(q(c_1(L))=0\). Then
\(\chi(L)=n+1\).
\end{lemma}

\begin{proof}
Indeed, \(\chi(L)=\chi(\mathcal{O}_M)=n+1\),
with the second equality implied by
Bochner's vanishing theorem.
\end{proof}

\begin{definition}
\label{seminars-complex-geometry-definition-semipositive}
A real \((1,1)\)-form \(\eta\) on a complex
manifold \(M\) is called semipositive if
\(\eta(x,Ix)\geq 0\) for all real tangent
vectors \(x\).
\end{definition}

\medskip\noindent
The following theorem was rediscovered
several times during 1990-ies (Enoki,
Takegoshi, Morougane). Its most general form
(which we do not use) is due to Demailly,
Peternell and Schneider.

It is similar to Kodaira-Nakano
vanishing theorem: instead of having that
$H^i(L \otimes K)$ will vanish for ample $L$,
we have an estimate for how big it can be:

\begin{theorem}[ETMDPS]
\label{seminars-complex-geometry-theorem-etmdps-vanishing}
Let \((M,I,\omega)\) be a compact K\"ahler
manifold, \(\dim_{\mathbf{C}} M = n\), \(K\)
its canonical bundle, and \(L\) a
holomorphic line bundle on \(M\) equipped
with a Hermitian metric \(h\). Assume that
the curvature \(\Theta\) of \(L\) is a
positive form on \(M\). Then the wedge
multiplication operator \(\eta \mapsto
\omega^i \wedge \eta\) induces a surjective
map
\[
H^0(\Omega^{n-i}_M \otimes L)
\xrightarrow{\ \omega^i\ }
H^i(K \otimes L).
\]
Here \(\omega\) is considered as an element
in \(H^1(\Omega^1_M)\), and multiplication by
\(\omega\) maps
\[
H^k(\Omega^{n-l}_M \otimes L)
to H^{k+1}(\Omega^{n-l+1}_M \otimes L).
\]
\end{theorem}

\begin{lemma}
\label{seminars-complex-geometry-lemma-global-sections-vanish}
Let \(B\) be a vector bundle equipped with a
filtration \(0 = B_0 \subset B_1 \subset
\ldots \subset B_k = B\). Assume that
\(H^0(B_i/B_{i-1}) = 0\) for all \(i\).
Then \(H^0(B)=0\).
\end{lemma}

\begin{theorem}
\label{seminars-complex-geometry-theorem-line-bundle-trivial-on-fibers}
Let \(M\) be a hyperk\"ahler manifold
admitting a Lagrangian fibration
\(\pi : M \to X\), and \(H\) a line bundle on
\(X\). Let \(L\) be a line bundle such that
\(L^{\otimes k} = \pi^*H\). Then \(L\) is
trivial on all smooth fibers of \(\pi\).
\end{theorem}

\begin{proof}
Step 1: Let \(F\) be a smooth fiber of
\(\pi\), which is an abelian variety by
Arnol'd-Liouville. Then \(T^*M|_F\) is an
extension of a trivial bundle \(TF\) with
another trivial bundle \(NF = T^*F\). For
any non-trivial line bundle
\(L \in \operatorname{Pic}_0(F)\), we have
\(H^0(L\otimes TF)=0\) and
\(H^0(L\otimes NF)=0\), which implies that
\(H^0(L\otimes T^*M|_F)=0\). Similarly one
\(H^0(L\otimes \wedge^k T^*M|_F)=0\)
(Lemma 2).
Step 2: Unless \(L\) is trivial on \(F\), we
have \(H^0(L\otimes \wedge^k T^*M|_F)=0\),
which implies that \(H^0(L\otimes \wedge^*M)
=0\). By Enoki-Mourugane-Takegoshi-Demailly-
Peternell-Schneider theorem, this implies
that \(H^i(L)=0\), hence \(\chi(L)=0\),
contradicting the formula \(\chi(L)=n+1\).
\end{proof}

Consider a fibration \(M \to \mathbf{C}P^n\)
with fiber \(T^n\). This gives an exact
sequence
\[
0 \to TF \to TM|_F \to \pi^*T\mathbf{C}P^n|_F
\to 0.
\]
The point is that when we have a
decomposition into trivial bundles, it is
very easy to control the sections.
The usage appears to be that a torsion
bundle tensored with a trivial bundle cannot
have sections.

\begin{proposition}
\label{seminars-complex-geometry-proposition-fiberwise-monodromy-line-bundle}
Let \(M\) be a hyperk\"ahler manifold
admitting a Lagrangian fibration
\(\pi : M \to X\), and \(H\) a line bundle on
\(X\). Let \(L\) be a line bundle such
that \(L^{\otimes k} = \pi^*H\). Then \(L\)
admits a connection \(\nabla\) which is flat
on each restriction \(L|_F\) to the fiber of
\(\pi\).
\end{proposition}

\begin{proof}
Choose a constant metric \(h^k\) on
\(L^{\otimes k}|_F = \mathcal{O}_F\) and let
\(h\) be its \(k\)-th root, which is a metric
on \(L|_F\). Since \(h^k\) is constant on
each fiber, its curvature vanishes, and the
Chern connection \(\nabla\) associated with
\(h\) is also flat.
\end{proof}

\begin{definition}
\label{seminars-complex-geometry-definition-fiberwise-monodromy}
Fiberwise monodromy of \(L\) is its
monodromy on the fibers of \(\pi\).
\end{definition}

\begin{remark}
\label{seminars-complex-geometry-remark-pullback-monodromy}
Clearly, \(L = \pi^*L_0\) if and only if the
fiberwise monodromy of \(L\) on each fiber
is trivial.
\end{remark}

\begin{remark}
\label{seminars-complex-geometry-remark-flat-connection}
Since \(\mathcal{O}(1)\) is flat on
\(\mathbf{C}^n = \mathbf{C}P^n \setminus H\),
we also obtain a flat connection on
\(L|_{\pi^{-1}(\mathbf{C}P^n\setminus H)}\).
The main theorem would follow if we prove
that the monodromy of this flat connection is
trivial on \(\pi^{-1}(z)\), where
\(z \in D\) is a general point of the
discriminant.
\end{remark}

\section{G4 fluxes and Hodge theory}
\label{seminars-complex-geometry-section-g4-fluxes-and-hodge-theory}

\noindent
Daniel Lopez, UFF.
GADEPS Focused Meeting June 30, 2026, IMPA.

\medskip\noindent
A flux $G$ is a primitive Hodge class,
i.e. an element of $H^{2,2}(X)
\cap H^4(X,\mathbb{Z})$.

\medskip\noindent
Conjecture: if $G$ is general,
 $|G|\gg \chi(X)/12$.

\medskip\noindent
Extra note:
$h^{3,1}$ is the dimension of the moduli
of fourfolds.

\noindent
The second method (the first
one is Griffiths basis)
uses Poincaré duality to pass
from (rational) Hodge classes
to (rational) homology classes and in
this language formulate hodge conjecture.
Namely, that any rational Hodge homology
class is algebraic.

\begin{example}
\label{seminars-complex-geometry-example-fermat-sextic}
The {\it Fermat sextic fourfold} 
$X \subset \mathbb{P}^5$ is the hyupersurface
given by
$$
F(x) = X_0^6+….+x_5^6.
$$
It has a lot of symmetries;
its automorphism group
is given as a semidirect product
and has an action
of symmetric group on automorphisms.

The Griffiths basis is
$$
\omega_\beta=
\operatorname{Res}
\left(\frac{x^\beta \Omega}{F^{k+1}}\right)
\in H^{4-k,k}
$$
where $\beta$ is the multiindex
$x^\beta= x_0^{\beta_0}…x_5^{\beta_5}$.
This information allows us to compute
the Hodge numbers of $X$.

Shioda proved in 1979 that $X$ satisfies
the Hodge conjecture over $\mathbb{Q}$.
Braun, Fortin, Lopez, Villaflor, in 2024,
provided a basis for $H^{2,2}(X) \cap
H^4(X,\mathbb{Q})$ given by algebraic cycles
divided into 3 different types.
\end{example}

\begin{example}
\label{seminars-complex-geometry-example-weighted-36-fermat}
Consider the {\it weighted 36 Fermat} 
given by
$$
F=x_0^{36}+x_1^{36}+x_2^{36}+x_3^4
+x_4^4+x_5^3=0
$$
in $\mathbb{P}_{1,1,1,9,12,12}$.

It's not known whether the Hodge conjecture
holds for this variety.
\end{example}



% --- End seminars-complex-geometry.tex ---

% --- Begin reading-the-masters.tex ---
\chapter{Reading the Masters}
\label{chapter-reading-the-masters}




\noindent
Personal notes on important mathematical
articles I read, one paragraph at a time.

\medskip

\section{Saint-Donat: first impressions}
\label{reading-the-masters-section-saint-donat}

\noindent
B. Saint-Donat, \textit{Projective Models
of K-3 Surfaces} \cite{saint-donat}.
Notes from the introduction, 9 September
2026. Here the K3 surfaces are projective
over $\mathbb{C}$.

\medskip\noindent
\textit{My first impressions, with minimal
corrections.}
I gather that the cool thing about
Saint-Donat's paper is that a line bundle
$L$ with $L^2>0$ provides a morphism to
projective space if $|L|$ is nonempty and
has no fixed components. I first understood
this as a projective embedding, but that
requires more. If I remember correctly,
Nakai--Moishezon says that $L$ is ample if
$L^2>0$ and $L\cdot C>0$ for every
irreducible curve $C$. This condition on
curves still matters on a K3; ampleness is
also different from very ampleness.

Saint-Donat also says he derives some other
theorem involving no s\'e qu\'e
non-hyperelliptic curves. I do not yet
understand that part.

And there was a third thing\ldots

\medskip\noindent
\textbf{AI clarification.}
The second point is an analogue of the
Enriques--Petri theorem about equations of
non-hyperelliptic canonical curves. For
non-hyperelliptic K3 projective models,
Saint-Donat proves
generation of the homogeneous ideal by
quadrics outside explicitly described
exceptions; see
\cite[Theorem 7.2]{saint-donat}.
The third point is that, if $L$ is ample,
then $L^{\otimes n}$ is basepoint-free for
$n\geq 2$ and very ample for $n\geq 3$;
see \cite[Theorem 8.3]{saint-donat}.

\medskip\noindent
\textbf{My question.}
So every line bundle on every variety
defines a birational map to projective
space?

\medskip\noindent
\textbf{AI answer.}
No. Let $X$ be an integral projective
variety over $\mathbb{C}$ and $L$ a line
bundle. If $H^0(X,L)=0$, the construction
using global sections gives no map.
Otherwise, a basis $s_0,\ldots,s_N$ of
$H^0(X,L)$ gives a rational map
$$
\begin{aligned}
\varphi_{|L|}:X
&\dashrightarrow\mathbb{P}^N,\\
x&\longmapsto[s_0(x):\cdots:s_N(x)].
\end{aligned}
$$
This formula is defined wherever the
sections do not all vanish. If $L$ is
basepoint-free, it defines a morphism on
all of $X$. If $L$ is very ample, the
morphism is a closed embedding.

Being \textit{birational onto its image}
means that the map restricts to an
isomorphism between dense open subsets
of $X$ and its image. This is an extra
property, not part of the meaning of
``rational map.''

For example, the trivial line bundle
$\mathcal{O}_{\mathbb{P}^1}$ has only
constant global sections. Its complete
linear system maps $\mathbb{P}^1$ to a
point, so the map is a morphism but is
not birational onto its image.

\section{Apéry: Fano threefolds and links}
\label{reading-the-masters-section-apery-fano-links}

\noindent
Roger Apéry, \textit{Étude de certaines
hypersurfaces}, Chapter III of his 1947
thesis. In modern language, Apéry studies
Picard-rank-one Fano threefolds and
birational maps obtained by blowing up a
curve on each side. The striking point is
that the numerical relations he derives
are essentially the numerical theory of
type-II Sarkisov links: his ``curves of the
second kind'' are the flopping curves.
Compare Section~\ref{reading-the-masters-section-mmp-and-sarkisov},
especially Theorem~\ref{reading-the-masters-theorem-sarkisov}.
There the HP states the modern Sarkisov
program: birational maps of Mori fibre
spaces are decomposed into Sarkisov links.
Apéry is studying a concrete rank-one Fano
case of this picture, where one blows up a
curve, passes through the flopping curves,
and contracts a curve on the other side.
His tables also contain the genus-$12$ Fano
threefold $V_{22}$ and locate the step in
Fano's argument that missed it.

\medskip\noindent
\textbf{Sergey's reconstruction.}
Sergey Galkin's September 2026 working
draft makes Chapter III readable in modern
terms: it gives the French text beside an
English translation, supplies a dictionary
from Apéry's terminology to current Fano
geometry, rewrites the birational
constructions as blow-up--flop--blow-down
Sarkisov links, and checks Apéry's fifty
table rows against the modern
classification. All fifty rows are
accounted for by the modern
classification, although Apéry's tables
are incomplete; Sergey also identifies the
few misprints and one false non-existence
statement in the chapter.

\section{Engel--de Gaay Fortman--Schreieder:
Mumford degenerations}
\label{reading-the-masters-section-mumford-survey}

\noindent
Philip Engel, Olivier de Gaay Fortman and
Stefan Schreieder,
\emph{Combinatorics and Hodge Theory of
Degenerations of Abelian Varieties:
A Survey of the Mumford Construction}.
These notes refer to the supplied version
dated January 28, 2026: Introduction,
pp.~2--3, and Example 3.31.

\medskip\noindent
{\bf Main idea.}
We have a family of abelian varieties
depending on a parameter, and we want
to describe what happens when the
parameter reaches a value where the
varieties degenerate. In the elliptic
example, the curves are quotients of a
cylinder, but the rule for identifying
points stops making sense at the missing
parameter value. Mumford's idea is to
{\bf extend the covering family of
cylinders, extend its integer action,
then take the quotient.}
Toric geometry supplies this enlargement
from a periodic tiling. Taking the quotient
then gives a family including a singular
limiting fibre: in the simplest example,
a projective line with two points
identified. The construction works for
abelian varieties in arbitrary dimension,
not only elliptic curves; see the
Introduction, pp.~2--3, Construction 3.3,
and Example 3.31. The calculations below
explain the initial quotient; understanding
the toric enlargement is the next step.

\medskip\noindent
{\bf Outline}
\begin{itemize}
\item First we take the "two step quotient"
$\mathbb{C}/\mathbb{Z}+\tau\mathbb{Z}$.
Step 1:
$\mathbb{C}/\mathbb{Z}\approx\mathbb{C}^*$,
the "cylinders" appear. Second step: use
exponential coordinate and the action
$w\mapsto q^nw$ to obtain
$\mathbb{C}/\mathbb{Z}+\tau\mathbb{Z}
\approx\mathbb{C}^{\ast}/q^{\mathbb{Z}}$.
\item we introduce the $xy=q$ notation by
considering the family before taking the
(in additive notation) $z\mapsto z+\tau n$
part of the quotient.
\item Via the $xy=q$ trick, we pass from
$$
\mathbb{C}^{\ast}\times\Delta^{\ast}
\longrightarrow\Delta^{\ast},
\qquad (w,q)\longmapsto q,
$$
to
$$
\begin{aligned}
\{(x,y)\in\mathbb{C}^2:|xy|<1\}
&\longrightarrow\Delta,\\
(x,y)&\longmapsto xy.
\end{aligned}
$$
This allows us to see something
above 0. But to obtain the family of
elliptic curves, we still have to take
that quotient. But we can't. So somehow
we introduce "one chart for every
integer n"
\end{itemize}

\medskip\noindent
Reading checkpoint, October 9, 2026:
the quotient and cylinder explanation
records the part worked through so far.
The toric extension remains to be understood.

\medskip\noindent
{\bf The elliptic curve case.}
Fix $\tau\in\mathbb{C}$ with
$\operatorname{Im}\tau>0$. The lattice
$\mathbb{Z}+\tau\mathbb{Z}$ consists of
all $m+n\tau$, with $m,n\in\mathbb{Z}$,
and
$$
E_\tau=\mathbb{C}/
(\mathbb{Z}+\tau\mathbb{Z}).
$$
Taking this quotient means identifying
$z$ with every $z+m+n\tau$.
First impose only the integer translations.
The exponential is surjective onto
$\mathbb{C}^*=\mathbb{C}\setminus\{0\}$,
and two points have the same exponential
exactly when their difference is an integer.
It therefore induces a biholomorphism
$$
\mathbb{C}/\mathbb{Z}
\longrightarrow\mathbb{C}^*,
\qquad [z]\longmapsto e^{2\pi iz}.
$$
Here biholomorphic means that the map and
its inverse are holomorphic.

\medskip\noindent
Writing $z=x+iy$, integer translations
identify $x$ with $x+1$ and leave $y$ free.
Thus, as topological spaces,
$$
\mathbb{C}/\mathbb{Z}
\cong(\mathbb{R}/\mathbb{Z})\times\mathbb{R}
\cong S^1\times\mathbb{R}.
$$
This is an infinite cylinder. In
$$
e^{2\pi i(x+iy)}
=e^{-2\pi y}e^{2\pi ix},
$$
$x$ modulo integers gives the angle and
$y$ determines the positive radius.

\medskip\noindent
Now impose the remaining translations
$z\mapsto z+n\tau$. Set
$$
q=e^{2\pi i\tau},\qquad 0<|q|<1.
$$
To translate this action, recall that
$w=e^{2\pi iz}$. Thinking of $z$ as a
logarithm of $w$ is right, with the
normalization
$$
z=\frac{\log w}{2\pi i}.
$$
This formula uses a local choice of
logarithm: there is no single-valued
holomorphic logarithm on all of
$\mathbb{C}^*$. Two choices differ by
$2\pi i k$, with $k\in\mathbb{Z}$, so
the resulting values of $z$ differ by
$k$. Thus $w$ determines $z$ modulo
$\mathbb{Z}$, exactly as in the first
quotient $\mathbb{C}/\mathbb{Z}$.

Now compute the multiplicative coordinate
of the translated point $z+n\tau$:
$$
\begin{aligned}
e^{2\pi i(z+n\tau)}
&=e^{2\pi iz}e^{2\pi in\tau}\\
&=w(e^{2\pi i\tau})^n\\
&=wq^n.
\end{aligned}
$$
Thus translation by $n\tau$ in the
additive coordinate $z$ becomes
multiplication by $q^n$ in the
multiplicative coordinate $w$.
This is a direct calculation using
the exponential, unpacking the quotient
$\mathbb{C}^*/q^{\mathbb{Z}}$ appearing
in Example 3.31 of the survey.
Hence
$$
E_\tau\cong\mathbb{C}^*/q^{\mathbb{Z}},
\qquad
q^{\mathbb{Z}}=\{q^n:n\in\mathbb{Z}\}.
$$
This notation identifies all the points
$q^nw$ with one another.
A family of lattices means that the
generators vary: here $1$ stays fixed
while $\tau$ varies. As
$\operatorname{Im}\tau\to\infty$, $q\to0$.
The multiplicative quotient cannot simply
be evaluated at $q=0$: multiplication
by zero is not an automorphism of
$\mathbb{C}^*$.

\medskip\noindent
{\bf From the family of cylinders to $xy=q$.}
For one fixed nonzero $q$, we have a
cylinder $\mathbb{C}^*$. Before taking
the integer quotient, the whole family is
$$
\mathbb{C}^*\times\Delta^*\longrightarrow\Delta^*,
\qquad (w,q)\longmapsto q,
\qquad \Delta^*=\{q\in\mathbb{C}:0<|q|<1\}.
$$
Here $w$ specifies a point on the cylinder
and $q$ specifies which fibre it belongs to.
We keep the multiplicative coordinate $w$
used above.

Set
$$
x=w,\qquad y=\frac{q}{w}.
$$
Then $xy=q$. Conversely, if $xy=q\ne0$,
we recover $w=x$ and $y=q/x$.
Thus, over nonzero $q$, this is exactly
our original family in different coordinates.
The parameter is still $q=e^{2\pi i\tau}$;
the equation $xy=q$ is not a new identity
deduced from the exponential.

The advantage is that we can now allow
$q=0$: take
$$
\{(x,y)\in\mathbb{C}^2:|xy|<1\}
\longrightarrow\Delta,\qquad (x,y)\longmapsto xy,
\qquad \Delta=\{q\in\mathbb{C}:|q|<1\}.
$$
Its fibre over $q\ne0$ is the cylinder
parametrized by $x\mapsto(x,q/x)$.
Its fibre over $0$ is
$$
\{xy=0\}=\{x=0\}\cup\{y=0\},
$$
two affine lines meeting at the origin,
not yet two projective lines.
This extends the family of cylinders.
Next we extend the integer action.

\medskip\noindent
{\it Relation to the paper.}
The coordinate explanation above is an
AI-assisted unpacking, understood in our
discussion on October 12, 2026.
In Example 3.31, p.~35, the survey instead
obtains charts $\mathbb{C}^2_{(x_n,y_n)}$
from the toric fan, states their gluing,
and gives the parameter map $u=x_ny_n$
(its $u$ is our $q$). It does not explicitly
present the introductory change of
coordinates $x=w$, $y=q/w$ in this form.

\medskip\noindent
{\bf Extending the action over zero.}
Start with the action $w\mapsto q^nw$,
where $n\in\mathbb{Z}$ and $q\ne0$
stays fixed. Recall our coordinates
$$
x=w,\qquad y=\frac{q}{w},\qquad xy=q.
$$
After applying the action, both
coordinates change:
$$
\begin{aligned}
x'&=q^nw=q^nx,\\
y'&=\frac{q}{q^nw}=\frac{y}{q^n}.
\end{aligned}
$$
Since $q=xy$, we can also write
$$
\begin{aligned}
x'&=(xy)^nx=x^{n+1}y^n,\\
y'&=\frac{y}{(xy)^n}=x^{-n}y^{1-n}.
\end{aligned}
$$
Thus $x^{n+1}y^n$ is only the new
$x$-coordinate. The action on the pair is
$$
(x,y)\longmapsto
\bigl(x^{n+1}y^n,x^{-n}y^{1-n}\bigr).
$$
The product remains $x'y'=xy=q$, so
the action stays in the same fibre.
In particular, the generator $n=1$ is
$$
(x,y)\longmapsto(qx,y/q)
=\bigl(x^2y,1/x\bigr).
$$

On the original cylinders, $x=w\ne0$,
so $1/x$ causes no problem. But we
enlarged the family to
$$
\{(x,y)\in\mathbb{C}^2:|xy|<1\}
\longrightarrow\Delta,\qquad
(x,y)\longmapsto xy.
$$
This includes new points with $x=0$,
all over $q=0$. These were not points
of the original cylinders; they need
not have an old coordinate
$w\in\mathbb{C}^*$. The formula $x=w$
identifies the original family with
the part over $q\ne0$.

At the new points with $x=0$, the
generator's formula $1/x$ is undefined.
Thus the action works on the original
family, but has not yet been extended
to all the added points.

This is a direct coordinate calculation
for the action defining the quotient in
Example 3.31 of the survey. The extension
using glued charts in that example,
p.~35, remains to be understood.

\medskip\noindent
{\bf Next steps}
Next time I'll understand the
"one chart for every integer n" part.
I also need to understand when the tiling
comes in.

\section{Engel--Mauri: hyperk\"ahler SYZ}
\label{reading-the-masters-section-engel-mauri-syz}

\noindent
Philip Engel and Mirko Mauri,
\emph{Hyperk\"ahler SYZ conjecture},
\href{https://arxiv.org/html/2610.12277v1}
{arXiv:2610.12277v1}, October 8, 2026.
The discussion below summarizes the
preprint; its proof has not been
independently checked here.

\medskip\noindent
{\bf Main question.}
When does a compact irreducible
hyperk\"ahler manifold admit a
Lagrangian fibration? In complex
dimension $2n$, this is a map
$$
f:X\longrightarrow B,
\qquad \dim_{\mathbb{C}}B=n,
$$
whose general fibres are complex tori
of dimension $n$, on which the
holomorphic symplectic form vanishes.
Singular fibres are allowed.
The authors' Theorem 1.1 asserts that
every nef line bundle is semiample.
For a nontrivial bundle with
Beauville--Bogomolov--Fujiki square
$q(c_1(L))=0$, this gives such a
fibration.

\medskip\noindent
{\bf What does the line bundle do?}
Recall the question in the Saint-Donat
section: how do sections of a line
bundle give a map?
Semiample means that, for some $m>0$,
sections of $L^{\otimes m}$ have no
common zero. A basis then gives
$$
x\longmapsto[s_0(x):\cdots:s_N(x)].
$$
The absence of common zeros makes this
defined everywhere.
Different points may have the same
image; their collection is a fibre.
Thus having enough sections to define
a map does not mean having enough to
embed $X$ into projective space.

Nef is a positivity condition: in the
projective case it means
$L\cdot C\geq0$ for every curve $C$.
For a compact K\"ahler manifold, use
$c_1(L)$ in the closure of the
K\"ahler cone.
Isotropic means $q(c_1(L))=0$, although
$c_1(L)\neq0$. The form $q$ is the
higher-dimensional counterpart of the
intersection form on a K3 surface.
These conditions concern the class of
$L$; semiampleness concerns actual
sections. Passing from the former to
the latter is the problem.

\medskip\noindent
{\bf The K3 picture to keep in mind.}
On a projective complex K3 surface,
a nontrivial nef line bundle with
$L^2=0$ is basepoint-free and is a
positive multiple of an elliptic fibre.
See \cite[Chapter 2, Proposition 3.10]
{huc-k3}; see also Corollary 3.12 there
for semiampleness of every nef bundle.
Geometrically, we obtain a surface
fibred by genus-one curves.
This is a useful companion to the
positive-square situation in the
Saint-Donat section: square zero
points toward a fibration.
For a fourfold, the corresponding
picture has a surface as base and
two-dimensional complex tori as fibres.

\medskip\noindent
{\bf Where degeneration and rotation enter.}
The authors use a maximal degeneration
to obtain special Lagrangian tori
(Theorem 3.2).
These are initially real submanifolds.
Hyperk\"ahler rotation changes the
complex structure while keeping the
metric. Under the cohomological
condition in Proposition 2.5, it makes
the tori holomorphic Lagrangian.
This is not automatic in higher
dimension.
Deformation arguments then yield
the required fibration
(Section 4).

\medskip\noindent
{\bf Connection with my current reading.}
For the Mumford and Kummer questions,
the useful theme is that degeneration
can reveal tori which are difficult
to see on the original manifold.
Mumford's construction gives concrete
degenerating families; here degeneration
is part of an existence argument.
This suggests a question to keep in
mind in the crystal project: can the
geometry of the degeneration identify
a candidate isotropic class or torus
fibre?

The theorem's hypotheses already
require a smooth hyperk\"ahler
manifold and a nef line bundle.
Consequently, it does not by itself
show that the Stanley--Reisner scheme
of the nine-vertex triangulation of
$\mathbb{CP}^2$ is smoothable, or that
a smoothing is a generalized Kummer
fourfold. Those remain separate
questions in our project.


\section{Manin: theta functions,
quantum tori and Heisenberg groups}
\label{reading-the-masters-section-manin-quantum-tori}

\noindent
Yuri I. Manin,
\emph{Theta Functions, Quantum Tori
and Heisenberg Groups},
\href{https://arxiv.org/abs/math/0011197}
{arXiv:math/0011197v1}.

\medskip\noindent
basic idea:

classical setting: instead of just an
algebraic group, work with quotient, this
is theory by Mumford.

Manin: introduce a poisson structure and
see how multiplication in the ring of
functions looks like now.

catch: the deformation might be encoded in
a lattice parameter q, which, in the case
of the Engel's survey is the whole point:
q can't be zero so we extend "by hand" the
deformation to 0.


\section{Huybrechts: compact hyperk\"ahler manifolds}
\label{reading-the-masters-section-huybrechts-hk-basic-results}

\noindent
Daniel Huybrechts, \emph{Compact Hyperk\"ahler Manifolds:
Basic Results}, \href{https://arxiv.org/abs/alg-geom/9705025}
{arXiv:alg-geom/9705025} (1997).
Reading: Introduction and Sections~1, 5--7.

\subsection{Historical and bibliographical background}

\noindent
A selective historical reading list. The first group
contains works cited in Huybrechts' 1997 bibliography;
the second adds other foundational and later papers.
Bracketed numbers in the first group are Huybrechts'
original reference numbers.

\begin{itemize}
\item \textbf{Beauville [2] (1983).}
\emph{Vari\'et\'es k\"ahleriennes dont la premi\`ere
classe de Chern est nulle},
\emph{J. Differential Geom.} 18, 755--782.
The foundational examples (Hilbert schemes of points
on K3 surfaces, generalized Kummer varieties)
and the quadratic form on $H^2$.
See Huybrechts Sections~1.9 and 2.1--2.3.

\item \textbf{Bogomolov [6] (1974), [7] (1978).}
\emph{On the decomposition of K\"ahler manifolds
with trivial canonical class},
\emph{Math. USSR-Sb.} 22, 580--583;
\emph{Hamiltonian K\"ahler manifolds},
\emph{Sov. Math. Dokl.} 19, 1462--1465.
Decomposition and deformations in the background
of the theory (Huybrechts Sections~1.1, 1.12).

\item \textbf{Fujiki [21] (1987).}
\emph{On the de Rham Cohomology Group of a Compact
K\"ahler Symplectic Manifold},
\emph{Adv. Stud. Pure Math.} 10, 105--165.
Structure of cohomology and the relation
$\int_X\alpha^{2n}=c_Xq_X(\alpha)^n$.
See Huybrechts Sections~1.6--1.11.

\item \textbf{Verbitsky [51] (1996).}
\emph{Cohomology of compact hyperk\"ahler
manifolds and its applications},
\emph{Geom. Funct. Anal.} 6(4), 602--611.
Cohomological structure in higher dimensions.
See Huybrechts Sections~1.6, 1.9, 8.

\item \textbf{Mukai [39] (1984).}
\emph{Symplectic structure of the moduli space of
sheaves on an abelian or K3 surface},
\emph{Invent. Math.} 77, 101--116.
Holomorphic symplectic moduli spaces of sheaves,
an important source of examples
(Huybrechts Section~2.10).

\item \textbf{Burns--Rapoport [12] (1975).}
\emph{On the Torelli Problem for K\"ahlerian K3 Surfaces},
\emph{Ann. Sci. \'Ec. Norm. Sup\'er.} 8, 235--274.
A key K3 Torelli precedent for the arguments
in Huybrechts Section~4.

\item \textbf{Tian [48] (1987) and Todorov [49] (1989).}
Unobstructedness of deformations with trivial
canonical bundle; invoked in Huybrechts Section~1.12.
See his bibliography for the complete titles.

\item \textbf{Huybrechts [28] (1997).}
\emph{Birational symplectic manifolds and their
deformations}, cited as forthcoming in
\emph{J. Differential Geom.} 45.
Precursor to the birational results in Section~4.
\end{itemize}


\medskip\noindent
{\bf Additional foundational works not cited in the 1997 survey
(or published subsequently).}
\begin{itemize}
\item E. Calabi, \emph{On K\"ahler manifolds with vanishing canonical class},
in \emph{Algebraic Geometry and Topology} (1955), 78--89.
The Ricci-flat K\"ahler metric problem.
\item S.-T. Yau, \emph{On the Ricci curvature of a compact
K\"ahler manifold and the complex Monge--Amp\`ere equation, I},
\emph{Comm. Pure Appl. Math.} 31 (1978), 339--411,
\href{https://doi.org/10.1002/cpa.3160310304}{DOI}.
Solves the Calabi conjecture.
\item K. G. O'Grady, \emph{Desingularized moduli spaces of
sheaves on a K3}, \emph{J. Reine Angew. Math.} 512
(1999), 49--117,
\href{https://arxiv.org/abs/alg-geom/9708009}{preprint};
and \emph{A new six-dimensional irreducible symplectic variety},
\emph{J. Algebraic Geom.} 12 (2003), 435--505,
\href{https://arxiv.org/abs/math/0010187}{preprint}.
New deformation types beyond Beauville's two series.
\item J.-P. Demailly and M. P\u{a}un,
\emph{Numerical characterization of the K\"ahler cone
of a compact K\"ahler manifold},
\emph{Ann. of Math.} 159 (2004), 1247--1274,
\href{https://doi.org/10.4007/annals.2004.159.1247}{DOI}.
Numerical description of K\"ahler cones.
\item M. Verbitsky, \emph{Mapping class group and a global
Torelli theorem for hyperk\"ahler manifolds},
\emph{Duke Math. J.} 162 (2013), 2929--2986.
Birational Torelli via Teichm\"uller spaces and periods.
\item E. Amerik and M. Verbitsky,
\emph{Rational curves on hyperk\"ahler manifolds},
\emph{Int. Math. Res. Not.} (2015), 13009--13045,
\href{https://arxiv.org/abs/1401.0479}{preprint}.
MBM classes and walls of the K\"ahler cone.
\item E. Markman, \emph{A survey of Torelli and monodromy
results for holomorphic-symplectic varieties},
\href{https://arxiv.org/abs/1101.4606}{arXiv:1101.4606}
(2011). A useful later survey (not an original theorem)
organizing monodromy and birational cones.
\end{itemize}

\section{Fukaya: Lagrangian foliations and noncommutative mirror symmetry}

\noindent
\textbf{Paper.} Kenji Fukaya, \emph{Floer homology of Lagrangian foliation and noncommutative mirror symmetry I} (Kyoto University preprint 98-08, 1998). The preprint studies the relationship between Lagrangian foliations, Floer homology, and noncommutative mirror symmetry. A Lagrangian foliation decomposes a symplectic manifold into Lagrangian leaves. This is a provisional description based on the bibliographic record, not a verified reading of the preprint.

\medskip\noindent
\textbf{Why we care.} Sergey pointed to Fukaya after describing two ways to present the noncommutative torus: the irrational rotation algebra associated with the action of $\mathbb Z\theta$ on $\mathbb R/\mathbb Z$, and the algebra of the corresponding irrational foliation of the two-torus. He said that Rieffel's Morita equivalence relates these descriptions (and, in the rotation picture, the parameters $\theta$ and $1/\theta$). He also mentioned complex structures on these algebras, holomorphic bundles, and derived categories related to an elliptic curve $E_\tau$. Separately, following Parshin, he described the Tate curve $E_q=\mathbb C^*/q^{\mathbb Z}$ as a moduli space of certain representations of a quantum torus with relation $XY=qYX$. This suggests two bridges worth understanding: from foliations and noncommutative tori to elliptic curves via mirror symmetry (the reason to read Fukaya), and from quantum-torus representations to the Tate curve (Parshin's picture). Sergey then asked whether restricting quantum-torus algebras to cones, and gluing over a fan, could make sense of $q\to0$ in a manner reminiscent of Mumford. We have not established that Fukaya's preprint proves any part of that proposed degeneration; it is a suggested conceptual reference.

\medskip\noindent
\textbf{Our noncommutative dream (speculative).} In \texttt{crystal} we study degenerations of quotients by crystallographic groups through Tate--Mumford constructions. A possible noncommutative analogue starts with a quantum torus, with generators $X,Y$ satisfying $XY=qYX$, restricts to cone subalgebras so that a limit as $q\to0$ may exist, and studies moduli of suitable representations. Whether this yields a Mumford-like degeneration or relates to our Kummer fourfold is an open question, not an established result.

\medskip


% --- End reading-the-masters.tex ---

\part{Algebraic Geometry}
\label{part-algebraic-geometry}

% --- Begin ringed-spaces.tex ---
\chapter{Ringed Spaces}
\label{chapter-ringed-spaces}





\section{Sheaves}
\label{ringed-spaces-section-presheaves-stalks-sheaves}
\label{ringed-spaces-section-sheaves}

\medskip\noindent
{\bf Source.}
The material in this section is transcribed verbatim from
Chapter 5 of Jethro van Ekeren's \emph{Homological Algebra.
IMPA 2026-2} course notes
\cite[Chapter 5]{vanEkerenHomologicalAlgebra2026}.
Only the \LaTeX\ formatting and HP labels have been added.

\subsection{Sheaves on topological spaces}
\label{ringed-spaces-section-presheaves}
\label{ringed-spaces-section-stalks}

In this section we review sheaves over a topological space. The basic idea is that a sheaf
consists of objects of a category attached to the open sets in a topological space, in a way
that should be “consistent” in some sense. Actually this describes the notion of presheaf,
the notion of sheaf being a little more than this.

\begin{definition}
\label{ringed-spaces-definition-presheaf}
Let $X$ be a topological space and $\mathcal C$ a category (which for now might as
well be $\mathcal C = \operatorname{Set}$). A presheaf $\mathcal F$ on $X$ consists of an assignment
$U\mapsto\mathcal F(U)$ of an object $\mathcal F(U)\in\mathcal C$ to each open $U\subset X$, and a
morphism $\operatorname{Res}_{U,V}:\mathcal F(U)\to\mathcal F(V)$ to each inclusion of open sets
$V\hookrightarrow U$, such that
\begin{itemize}
\item $\operatorname{Res}_{U,U}=\operatorname{Id}$,
\item whenever $W\subset V\subset U$ we have
$\operatorname{Res}_{V,W}\circ\operatorname{Res}_{U,V}=\operatorname{Res}_{U,W}$.
\end{itemize}
The elements of $\mathcal F(U)$ are called the sections of $\mathcal F$ over $U$, though this terminology only makes
sense if $\mathcal A$ is a concrete category, i.e., if objects of $\mathcal A$ are “sets with structure”. The collection
of all presheaves on $X$ (with values in $\mathcal C$) themselves form a category, with a morphism
$\varphi:\mathcal F\to\mathcal G$ consisting of a collection
$\varphi_U:\mathcal F(U)\to\mathcal G(U)$ for each open $U$, compatible with all restrictions, i.e.,
\[
\xymatrix{
\mathcal F(U)\ar[r]^{\varphi_U}\ar[d]_{\operatorname{Res}}&
\mathcal G(U)\ar[d]^{\operatorname{Res}}\\
\mathcal F(V)\ar[r]_{\varphi_V}&
\mathcal G(V).
}
\]
\end{definition}

\begin{remark}
In Chapter 2 we have associated with the topological space $X$ a certain
category $\operatorname{Open}(X)$. One notices that the definition of presheaf can be summarised very
succinctly as follows: a presheaf on $X$ with values in $\mathcal C$ is a functor
\[
\operatorname{Open}(X)^{\mathrm{op}}\to\mathcal C.
\]
Under this correspondence a morphism of presheaves is, as expected, a natural transformation of such functors.
\end{remark}

The most primitive example, and one of the most important, is the following: Let $A\in\mathcal C$
be a fixed object. We assign $\mathcal F(U)=A$ for every open $U\subset X$, and set
$\operatorname{Res}_{U,V}=\operatorname{Id}$ for all $V\subset U$. A presheaf of this form is known as a constant presheaf on $X$.

Another class of examples come from considering functions. Let $X$ be a topological space
and $R$ a commutative ring, then
\[
\mathcal F(U)=\text{the set of functions }f:U\to R
\]
defines a presheaf on $X$, where the restriction maps are literally restriction of functions. If
$R$ carries a topology, such as $\mathbb R$ with the metric topology or $\mathbb Z$ with the discrete topology,
then
\[
\mathcal F(U)=\text{the set of continuous functions }f:U\to R
\]
also defines a presheaf on $X$, which we denote $\mathcal C_X$. If $X$ is a smooth manifold, and $R=\mathbb R$,
then we have the notion of smooth function, and
\[
\mathcal F(U)=\text{the set of smooth functions }f:U\to\mathbb R
\]
defines a presheaf on $X$, which we denote $\mathcal C^\infty_X$. A similar definition can be made for $X$ a
complex manifold, $R=\mathbb C$, and sets of holomorphic functions.

To build intuition, I suggest imagining sheaves for the case $\mathcal C=\operatorname{Ab}$ or even $\mathcal C=\operatorname{Vec}$.
Most of the examples we study later will be with values in an abelian category. But as we
see already in the case of $\mathcal C^\infty_X$, for which the most natural target category is the category
commutative rings, there are plenty of natural and relevant examples where we must allow
targets which are not abelian categories.

\begin{definition}
\label{ringed-spaces-definition-stalk}
Let $\mathcal F$ be a presheaf on $X$. For any point $p\in X$ we define the stalk $\mathcal F_p$
to be the colimit (also known in this case as “direct limit”)
\[
\operatorname{colim}_{p\in U}\mathcal F(U)
\]
over all $U$ containing $p$, along the morphisms $\operatorname{Res}$. Elements of the stalk are called germs.
\end{definition}

Let us understand this definition more concretely: germs of $\mathcal F$ at $p$ consist of pairs $(U,s)$
where $U$ is an open containing $p$ and $s\in\mathcal F(U)$. Two such pairs represent the same germ if
there is an open subset of the intersection on which the restrictions of the sections coincide.
If our presheaf is taken over an abelian category, then addition of germs $(V_1,s_1)$ and $(V_2,s_2)$
is carried out by restriction to $V_1\cap V_2$ and addition of the sections $s_1$ and $s_2$ in
$\mathcal F(V_1\cap V_2)$.
Near the beginning of the course I set the formulation of the definition of a colimit as an
exercise. Let's recall the general definition.

\begin{definition}
Let $J$ be a diagram in a category $\mathcal C$, that is, a collection of objects
$A_i\in\mathcal C$ and a collection of morphisms $\alpha_{ij}:A_i\to A_j$. We consider the
undercategory of $J$, which consists of: a choice of an object $X\in\mathcal C$ and a collection
of compatible morphisms $f_i:A_i\to X$ for each $A_i$ (compatible means
$f_j\circ\alpha_{ij}=f_i$ whenever $\alpha_{ij}:A_i\to A_j$ is part of $J$). A colimit of $J$
is an initial object in the undercategory of $J$, that is, an object $(X,\{f_i\})$ such that
for any other object $(X',\{f'_i\})$, there exists a unique morphism $X\to X'$ making
a bunch of obvious diagrams commute.
\end{definition}

In order to be sure that the stalks of a presheaf $\mathcal F$ exist, we ought to require that $\mathcal F$ take
values in a category $\mathcal C$ in which colimits exist, or at least in which filtered colimits exist. A
filtered colimit is a colimit over a filtered diagram, and a filtered diagram is one in which (1)
for any $X$ and $Y$ in the diagram, there exists $Z$ in the diagram with morphisms $X\to Z$
and $Y\to Z$ and (2) if $a,b:X\to Y$ are two morphisms in the diagram, there exists $Z$ in
the diagram and $c:Y\to Z$ such that $c\circ b=c\circ a$. Certainly the diagrams relevant for the
definition of stalks are filtered. And certainly for categories like $\operatorname{Set}$, $\operatorname{Ab}$,
$\operatorname{Vec}$, the necessary colimits exist.

\begin{definition}
\label{ringed-spaces-definition-sheaf}
A sheaf is a presheaf $\mathcal F$ such that the following condition holds: Let
$S\subset X$ be an open, and let $\{V_i\}_{i\in I}$ be an open cover of $U$, i.e.,
$V_i\subset U$ for each $i\in I$ and $\bigcup_{i\in I}V_i=U$, then
\begin{enumerate}
\item if $s,t\in\mathcal F(U)$ satisfy
$\operatorname{Res}_{U,V_i}(s)=\operatorname{Res}_{U,V_i}(t)$ for each $i\in I$, then $s=t$,
\item if $s_i\in\mathcal F(V_i)$ is given for each $i\in I$, and if
$\operatorname{Res}_{V_i,V_i\cap V_j}(s_i)=
\operatorname{Res}_{V_j,V_i\cap V_j}(s_j)$ for each $i,j\in I$, then there exists
$s\in\mathcal F(U)$ such that $\operatorname{Res}_{U,V_i}(s)=s_i$ for each $i\in I$.
\end{enumerate}
\end{definition}

Intuitively the definition says that sections over pieces of a cover of $U$ that “match” on
intersections, come from a unique section over $U$. Or, even more tersely, “local sections can
be glued”.

\begin{example}
Let $X$ be a topological space. The presheaf of all functions on $X$ is a sheaf.
Indeed a function is uniquely specified by its values on points, and so a function defined
compatibly on members of a cover $\{V_i\}$ of $U$ automatically defines (uniquely) a function on
$U$. The presheaf $\mathcal C_X$ of continuous functions on $X$ is similarly a sheaf, because the unique
gluing of local continuous functions is again continuous. Similarly also the presheaf
$\mathcal C^\infty_Z$ of smooth functions is a sheaf.
\end{example}

Every sheaf is a presheaf by definition, but not every presheaf satisfies the “gluing axiom”
that characterises sheaves. The constant presheaf furnishes an example: Indeed consider the
constant presheaf $A^\circ_X$ on some space $X$ and consider a pair of disjoint opens
$U_1,U_2\subset X$. Assuming $A$ has more than one element, we can consider the sections
$a_1\in A^\circ_X(U_1)$ and $a_2\in A^\circ_X(U_2)$, and if $a_2\neq a_1$ then these sections
will not be restrictions of a common section in $A^\circ_X(U_1\cup U_2)$.

\medskip\noindent
{\it Figure 5.1: Some sections of the constant presheaf might not glue on $U_1\cup U_2$}


It is easy to see, by the way, that the stalks of the constant presheaf $A^\circ_X$ are all isomorphic
to $A$.
There is a natural way to convert any presheaf into a sheaf. This process is called sheafification
and it consists intuitively of “adding glued together local sections” which don't already
exist in the presheaf. Formally, we consider the disjoint union
\[
\operatorname{\acute Et}(\mathcal F)=\coprod_{p\in X}\mathcal F_p.
\]
For open $U\subset X$ we can consider functions $f:U\to\operatorname{\acute Et}(\mathcal F)$ such that
$f(p)\in\mathcal F_p$ for all $p\in U$. In other words sections over $U$ of the natural projection
$\pi:\operatorname{\acute Et}(\mathcal F)\to X$. We can also consider the subset of such functions
with the following property: For all $p\in U$ there exists a neighbourhood $V\subset U$ of $p$
and a section $s\in\mathcal F(V)$, such that $f(q)=(V,s)\in\mathcal F_q$ for all $q\in V$.
Call the latter set of functions $\mathcal F^+(U)$. Then the claim is:

\begin{proposition}
The assignment $U\mapsto\mathcal F^+(U)$ defines a sheaf on $X$.
\end{proposition}

\begin{remark}
What on earth does “$\operatorname{\acute Et}$” mean? It stands for étale, which is French for
something. In fact the set $\operatorname{\acute Et}(\mathcal F)$ can be equipped with a certain topology, and then it is
called the étale space of the presheaf. The topology is the one generated by the following
base of open sets: For $U\subset X$ open and $s\in\mathcal F(U)$ we define
\[
[U,s]=\{(p,s(p))\mid p\in U\}\subset\operatorname{\acute Et}(\mathcal F)
\]
(see [6, Chap. II, Exercise 1.13]). This topology is quite difficult to visualise, and if you
are not interested, don't bother trying. The point of this construction is that the functions
$f:U\to\operatorname{\acute Et}(\mathcal F)$ that we singled out above to make up $\mathcal F^+(U)$ -- these are precisely the
continuous sections of $\pi:\operatorname{\acute Et}(\mathcal F)\to X$ over $U$.
\end{remark}

Let us see what sheafification does to the constant presheaf. Intuitively a section of the
sheafification over $U$ consists of assigning an element of $A$ to each connected component
of $U$. Indeed the stalks of $A^\circ_X$ are copies of $A$, and with a little work we can see that the
sheafification $A_X$ assigns to open $U\subset X$ the set of continuous functions $U\to A$ where $A$ is
endowed with the discrete topology. In other words, locally constant $A$-valued functions.
I would like to discuss vector bundles, but at first rather superficially.

\begin{definition}
A smooth manifold of dimension $n$ is a topological space\footnote{technically a second coutable Hausdorff topological space} $X$ equipped
with an open cover $\{U_i\}_{i\in I}$ and continuous maps $\psi_i:U_i\to\mathbb R^n$ which are homeomorphisms to
their images, and such that
\[
\psi_i\circ\psi_j^{-1}:\psi_j(U_i\cap U_j)\to\psi_i(U_i\cap U_j)
\]
is smooth for all $i,j\in I$.
\end{definition}

\begin{definition}
Let $X$ be a topological space. A (real) vector bundle of rank $k$ over $X$ is
a topological space $E$ and a surjective map $\pi:E\to X$, together with an open cover
$\{U_i\}_{i\in I}$ of $X$ and, for each $i\in I$, a homeomorphism
$\varphi_i:\pi^{-1}(U_i)\to U_i\times\mathbb R^k$, such that
\begin{enumerate}
\item for each $i\in I$, the diagram
\[
\xymatrix{
\pi^{-1}(U_i)\ar[rr]^{\varphi_i}\ar[dr]_{\pi}&&U_i\times\mathbb R^k\ar[dl]^{p_1}\\
&U_i
}
\]
commutes (in this way the fibres
$\pi^{-1}(p)=\varphi_i^{-1}(p\times\mathbb R^k)$ become real vector spaces via
transport of structure),
\item for each $i,j,\in I$ the composition
\[
\varphi_i\circ\varphi_j^{-1}:
\varphi_j(\pi^{-1}(U_i\cap U_j))\to
\varphi_i(\pi^{-1}(U_i\cap U_j))
\]
is continuous and its restriction to each fibre is $\mathbb R$-linear.
\end{enumerate}
\end{definition}

\begin{definition}
Let $X$ be a smooth manifold of dimension $n$. A smooth vector bundle of
rank $k$ over $X$ is the same data as in Definition 5.1.6 with the additional requirement that
each $\varphi_i\circ\varphi_j^{-1}$ is smooth.
\end{definition}

Let $\pi:E\to X$ be a vector bundle. We may define a presheaf $\mathcal E$ by
\[
\mathcal E(U)=\{s:U\to\pi^{-1}(U)\mid s\text{ continuous and }\pi\circ s=\operatorname{Id}_U\}.
\]
As for the “sheaves of functions” examples, this presheaf is already a sheaf. Indeed if
$E=X\times\mathbb R$ then $\mathcal E(U)$ is just the set of continuous functions on $U$ and construction recovers
$\mathcal C_X$.
I recall these constructions mostly to emphasise one basic issue: Over a point $p\in X$ the
vector bundle $E$ has a fibre $E_p=\pi^{-1}(p)\cong\mathbb R^k$, and the étale space
$\operatorname{\acute Et}(\mathcal E)$ has a stalk $\mathcal E_p$, and
these two objects are very different from one another. The stalk is essentially the set of all
germs of $\mathbb R^k$-valued functions at $p$ which is vastly bigger than $\mathbb R^k$.
I skipped these definitions in class.
The étale space is actually more similar in spirit to the notion of covering space than to the
notion of vector bundle.

\begin{definition}
Let $X$ be a topological space. A covering space over $X$ is a topological
space $\widetilde X$ together with a surjection $\pi:\widetilde X\to X$ such that for every $p\in X$ there is a
neighbourhood $U\subset X$ of $p$ such that $\pi^{-1}(U)$ is a disjoint union of sets homeomorphic to $U$.
\end{definition}

\begin{definition}
Let $X$ be a topological space. An étale space over $X$ is a topological space
$E$ and a continuous surjection $\pi:E\to X$, with the property that for every $e\in E$ there is a
neighbourhood $\widetilde U\subset E$ of $e$ such that the restriction
$\pi|_{\widetilde U}:\widetilde U\to\pi(\widetilde U)$ is a homeomorphism.
\end{definition}

We notice that every covering space is an étale space, but the reverse might not hold true.
Let us mention one more fun example of a sheaf, the “skyscraper sheaf”. This example is
somehow analogous to the Dirac delta function.

\begin{definition}
\label{ringed-spaces-definition-skyscraper}
Let $\mathcal A$ be an abelian category and $A\in\mathcal A$, and let $p\in X$. The skyscraper
sheaf at $p$, with stalk $A$, is defined as follows:
\[
\mathcal F(U)=
\begin{cases}
A,&p\in U,\\
0,&p\notin U,
\end{cases}
\]
with restriction map $\operatorname{Res}_{U,V}=\operatorname{Id}_A$ if $p\in V\subset U$ and $0$ otherwise.
\end{definition}

\subsection{Morphisms of sheaves}

A morphism of sheaves is just a morphism of presheaves. We will begin by figuring out
what monomorphisms and epimorphisms of sheaves of sets look like, and then we pass to
sheaves with values in an abelian category and discuss kernels and cokernels. The first thing
to notice is that any morphism $\phi:\mathcal F\to\mathcal G$ of presheaves induces a morphism
$\phi_p:\mathcal F_p\to\mathcal G_p$ of stalks for each point $p$.

\begin{definition}
Let $\mathcal F,\mathcal G$ be sheaves of sets and let $\phi:\mathcal F\to\mathcal G$ be a morphism. Then:
\begin{enumerate}
\item we say that $\phi$ is injective if $\phi_p:\mathcal F_p\to\mathcal G_p$ is injective for each point $p\in X$,
\item we say that $\phi$ is surjective if $\phi_p:\mathcal F_p\to\mathcal G_p$ is surjective for each point $p\in X$.
\end{enumerate}
\end{definition}

A subtle and important point should be noted. These definitions make sense for presheaves $\mathcal F$
and $\mathcal G$, but they are not very useful notions. Indeed consider the following fact: two distinct
morphisms $\phi_1,\phi_2:\mathcal F\to\mathcal G$ might induce the same morphism of stalks
$\mathcal F_p\to\mathcal G_p$ for every $p\in X$. To see how this can occur it helps to remember that a presheaf can fail to be a sheaf
in two ways, (1) that local sections fail to glue, or (2) that local sections glue but fail to do
so uniquely.
It is straightforward to translate these notions back into a language which makes no reference
to stalks:

\begin{proposition}
Let $\mathcal F,\mathcal G$ be sheaves of sets and let $\phi:\mathcal F\to\mathcal G$ be a morphism. Then:
\begin{enumerate}
\item the morphism $\phi$ is injective if and only if
$\phi_U:\mathcal F(U)\to\mathcal G(U)$ is injective for every open $U\subset X$,
\item the morphism $\phi$ is surjective if and only if for every $p\in X$ and every neighbourhood
$U$ of $p$ and $t\in\mathcal G(U)$, there exists a neighbourhood $V\subset U$ of $p$ and
$s\in\mathcal F(V)$ such that $\phi_V(s)=\operatorname{Res}_{U,V}(t)$.
\end{enumerate}
\end{proposition}

Let us denote by $s_p$ the germ of a section $s$ at $p$.

\begin{proof}
For the first part. In one direction, assume $\phi$ is injective. If $s_1,s_2\in\mathcal F(U)$ satisfy
$\phi(s_1)=\phi(s_2)$ then $\phi_p(s^1_q)=\phi_p(s^2_q)$ for all $q\in U$. Injectivity of $\phi$ now means that
$s^1_q=s^2_q$ at all $q\in U$, and hence $s_1=s_2$ by the sheaf axiom, thus $\phi_U$ is injective. In the other direction,
if there were some $p\in X$ at which $\phi_p$ were not injective, then let $s^1_p\neq s^2_p$ be two distinct
germs in $\mathcal F_p$ such that $\phi_p(s^1_p)=\phi_p(s^2_p)$. Choosing representatives $(U^1,s_1)$ and $(U^2,s_2)$ of the
germs, and passing to a subset $V$ of their intersection as necessary, we obtain two sections
in $\mathcal F(V)$ with the same image in $\mathcal G(V)$.
For the second part. The stated condition is essentially the spelling out of the surjectivity
of the colimit, so the equivalence is immediate.
\end{proof}

For the reason alluded to above, the corresponding statement for presheaves is not true.

\begin{remark}
The proofs of both parts of Proposition 5.2.1 proceed by manipulating
elements of sets, and one might ask how such arguments generalise to other categories. There
exists a notion of concrete category, which informally is a category whose objects are sets with
structure, and which formally is a category $\mathcal C$ equipped with a faithful functor $U:\mathcal C\to\operatorname{Set}$
(think of the forgetful functor $U:\operatorname{Ab}\to\operatorname{Set}$). It is not quite true that Proposition 5.2.1
holds for sheaves valued in $\mathcal C$ for any concrete category $\mathcal C$, but it is true with one additional
condition. If $\mathcal C$ is a concrete category and if $U$ has a left adjoint $L:\operatorname{Set}\to\mathcal C$, and if $L(*)\in\mathcal C$
is a compact object, then the proposition holds. A compact object $C\in\mathcal C$ means an object
for which for any filtered diagram $J$ the canonical morphism
\[
\operatorname{colim}_{i\in I}\operatorname{Hom}_{\mathcal C}(C,X_i)
\to
\operatorname{Hom}_{\mathcal C}\left(C,\operatorname{colim}_{i\in I}X_i\right)
\]
is an isomorphism.
\end{remark}

Finally

\begin{proposition}
Let $\mathcal F,\mathcal G$ be sheaves of sets and let $\phi:\mathcal F\to\mathcal G$ be a morphism. Then:
\begin{enumerate}
\item the morphism $\phi$ is a monomorphism if and only if it is injective,
\item the morphism $\phi$ is an epimorphism if and only if it is surjective.
\end{enumerate}
\end{proposition}

\begin{proof}
Just some hints for the proof. For surjectivity consider a skyscraper sheaf. For
injectivity consider the following construction: let $U\subset X$ be an open subset, and consider
the “indicator sheaf” $*_U$ defined by $*_U(V)$ is a point $*$ if $V\subset U$ and $\varnothing$ otherwise.
\end{proof}

\begin{remark}
Proposition 5.2.1 relates the stalk-wise definitions of injectivity and surjectivity
to some equivalent characterisations. The previous remark observes that the proposition
holds for a broad class of concrete categories, which includes $\operatorname{Ab}$, $\operatorname{Grp}$, $\operatorname{CRing}$. But in
any of these categories for which epimorphisms do not coincide with surjections (say $\operatorname{CRing}$)
these results are of limited relevance.
\end{remark}

Now let $\mathcal A$ be an abelian category and let $\operatorname{Sh}_{\mathcal A}(X)$ be the category of sheaves on $X$ with
values in $\mathcal A$.

\begin{proposition}
The category $\operatorname{Sh}_{\mathcal A}(X)$ is an abelian category. The kernel and cokernel
of a morphism are given by the stalk-wise kernel and cokernel.
\end{proposition}

For $\mathcal F,\mathcal G\in\operatorname{Sh}_{\mathcal A}(X)$ we have the set
$\operatorname{Hom}(\mathcal F,\mathcal G)$. As we know this consists of a bunch of
compatible morphisms of objects of $\mathcal A$. Since all these morphisms live in abelian groups,
$\operatorname{Hom}(\mathcal F,\mathcal G)$ acquires an abelian group structure. There is another notion of Hom for sheaves
though, sometimes denoted $\mathcal Hom(\mathcal F,\mathcal G)$.

Firstly, let $\mathcal F$ be a sheaf on $X$ and $U\subset X$ an open subset. We have
$\mathcal F(U)$ which is just some object of $\mathcal A$, but we have another entity $\mathcal F|_U$ a sheaf on $U$ defined by
\[
\mathcal F|_U(V)=\mathcal F(V)\text{ for each open }V\subset U.
\]
\label{ringed-spaces-definition-restriction-sheaf}
The sheaf Hom $\mathcal Hom(\mathcal F,\mathcal G)$ is defined analogously.
$\mathcal Hom(\mathcal F,\mathcal G)$ is a sheaf of abelian groups on $X$, given by
\[
\mathcal Hom(\mathcal F,\mathcal G)(U)=\operatorname{Hom}(\mathcal F|_U,\mathcal G|_U)
\text{ for each open }V\subset U.
\]
Notice that trying to define the sheaf by
\[
\mathcal Hom(\mathcal F,\mathcal G)(U)=\operatorname{Hom}(\mathcal F(U),\mathcal G(U))
\]
would not work as an element of $\operatorname{Hom}(\mathcal F(U),\mathcal G(U))$ does not contain the information needed
to define a restriction in $\operatorname{Hom}(\mathcal F(V),\mathcal G(V))$ for some open subset $V\subset U$.

\begin{definition}
\label{ringed-spaces-definition-local-system}
A local system on $X$ is a sheaf $\mathcal F\in\operatorname{Sh}_{\mathcal C}(X)$ with the following property:
there is a fixed object $A\in\mathcal C$ and, for every $p\in X$ there is a neighbourhood $U$ of $p$ and an
isomorphism $\mathcal F|_U\cong A_U$.
\end{definition}

The constant sheaf is an obvious example of such a thing, but there are other examples too.
Let $X=\mathbb C^\times=\mathbb C\setminus\{0\}$ and let
\[
\mathcal F(U)=\{ax^{2/3}\mid a\in\mathbb C\}.
\]
The function $f(x)=x^{2/3}$ is not globally defined without a branch cut (in which case it
becomes discontinuous across the cut). But in any sufficiently small ball we can define the
function unambiguously, more precisely in
\[
B_\varepsilon(a)=\{x\in\mathbb C\mid |x-a|<\varepsilon\}
\]
for $a\in\mathbb C^\times$ and $0<\varepsilon<|a|$.
We obtain a sheaf which is isomorphic to the constant sheaf $\mathbb C_U$ when restricted to any of
these balls, but which has no nonzero sections over $X$ and so is definitely not isomorphic to
$\mathbb C_X$.
A comment on specifying a sheaf on a base. See [15, Section 6.30] for a careful discussion.
We recall that a topological space consists of a set $X$ and a set $\tau$ of subsets $U\subset X$ called
the open subsets, such that $\tau$ satisfies certain axioms. By definition a base of $\tau$ is a subset
$\mathcal B\subset\tau$ such that every $U\in\tau$ is equal to the union of a (possibly infinite) collection of sets
in $\mathcal B$. The open balls form a base of the usual metric topology of $\mathbb R^n$ for example.

\medskip\noindent
{\it Figure 5.2: For each $i$ we have $\mathcal F|_{U_i}=\mathbb C_{U_i}$, while
$\mathcal F(\bigcup_i U_i)=\{0\}$}

To specify a sheaf it is sufficient to specify it on a base of the topology. If one wants to be
formal, it is possible to define the notion of “sheaf on $\mathcal B$” (where $\mathcal B$ is a base of a topology),
and then prove a proposition:

\begin{proposition}
Let $\mathcal F$ be a sheaf on a topological space $(X,\tau)$ with base $\mathcal B$, then the
restriction of $\mathcal F$ to the open sets in $\mathcal B$ defines a sheaf on $\mathcal B$, this operation defines a functor
from the category of sheaves on $\tau$ to the category of sheaves on $\mathcal B$, and this functor is an
equivalence of categories.
\end{proposition}

\begin{example}
Let $X$ be a topological space, and consider the sheaf $\mathcal F$ of $\mathbb R$-valued functions
on $X$. There is a sub presheaf $\mathcal F^{\mathrm{bdd}}$ consisting of bounded functions, i.e.
$\mathcal F^{\mathrm{bdd}}(U)$ consists of functions $f:U\to\mathbb R$ such that there exists $C>0$
such that $|f(p)|<C$ for all $p\in U$.
This is not a sheaf however, intuitively because “boundedness is not a local property”. Let
it be an exercise to formulate the notion of “locally bounded function” and confirm that the
sheafification of $\mathcal F^{\mathrm{bdd}}$ consists of locally bounded functions.
\end{example}

\medskip\noindent
Resolution of constant sheaf by singular cohomology.

\begin{definition}
\label{ringed-spaces-definition-singular-cochain-sheaf}
Recall $C_p(X)$ the free $\mathbb Z$-module generated by $p$-chains $\sigma:\Delta^p\to X$, and
the corresponding dual groups $C^p(X)$. It is clear that $U\mapsto C^p(U)$ defines a presheaf on
$X$. In fact it is not a sheaf (unless $p=0$), so we define the sheaf of singular $p$-cochains
$\mathcal C^p_{\mathrm{sing}}$ to be the sheafification of this presheaf. Since the groups $C^p(U)$ come with compatible
differentials $d:C^p(U)\to C^{p+1}(U)$ for all $U$, the $\mathcal C^p_{\mathrm{sing}}$ assemble into a complex of sheaves
\[
0\to\mathcal C^0\to\mathcal C^1\to\mathcal C^2\to\cdots.
\]
\end{definition}

What is the kernel of the morphism $d:\mathcal C^0\to\mathcal C^1$. We note that $\mathcal C^0$ is nothing other than
the sheaf of all ($\mathbb Z$-valued) functions on $X$. An element of $\mathcal C^1(U)$ is a gadget which assigns
values to paths $\gamma:[0,1]\to U$. The differential sends a function $f\in\mathcal C^0(U)$ to the element of
$\mathcal C^1(U)$ that maps $\gamma$ to the difference $f(\gamma(1))-f(\gamma(0))$. In order for $df=0$ to hold, then, it
must be the case that $f$ is constant on all path connected components of $U$. What we have
proved is that $\ker(d)=\mathbb Z_X$.

Next we would like to see that $\mathcal C^0\to\mathcal C^1\to\mathcal C^2$ is exact, i.e.,
$\ker(d)=\operatorname{im}(d)$ at this point. I don't want to get bogged down in topological pathologies, so I will assume $X$ is a smooth
manifold. Thus every point $p$ has a neighbourhood homeomorphic to a ball $B_\varepsilon(0)\subset\mathbb R^n$. Let
$\gamma$ be a $1$-simplex inside such $U$ with endpoint $\gamma(0)=p_0$ and $\gamma(1)=p_1$. For each $q\in U$ we
consider $\sigma_q$ the $1$-simplex moving radially from $p$ to $q$. If
$\alpha\in\ker(d)(U)\subset\mathcal C^1(U)$ then
\[
\alpha([\gamma])=\alpha([\sigma_{p_1}])-\alpha([\sigma_{p_0}]).
\]
We define $f\in\mathcal C^0(U)$ by $f(q)=\alpha([\sigma_q])$, and easily confirm that $df=\alpha$. So exactness is
established.

A similar argument can be used to prove exactness in all degrees $p>0$. But in fact we can
deduce this rather directly from what we already know, and in fact the argument presented
above, while instructive, was not really necessary. Indeed since $B_\varepsilon(0)$ is contractible, we know
the complex $C_\bullet(U)$, and therefore also $C^\bullet(U)$, is acyclic in degree $p>0$. The corresponding
statement for sheaves now follows because for every $p\in X$ and every neighbourhood of $p$
there is a sub neighbourhood homeomorphic to $B_\varepsilon(0)$, so the maps $d$ on stalks have to factor
through the maps $d$ on such balls.
\label{ringed-spaces-proposition-singular-cochains-resolve-constant-sheaf}

Let us pause to reflect on the peculiar nature of the sheaves $\mathcal C^p_{\mathrm{sing}}$ for $p>0$. Consider a
$p$-chain $\sigma_0$ in $C_p(U)$ where $U\subset X$ is some open. Suppose the image of $\sigma$ is not a single
point. We can consider two cochains $\alpha_1,\alpha_2\in C^p(U)$ which assign different values to $\sigma_0$ and
the same value to every other $p$-chain. Upon passing to $C^p(V)$ for any $V$ which does not
wholly contain the image of $\sigma_0$, the restrictions
$\operatorname{Res}_{U,V}\alpha_1$ and $\operatorname{Res}_{U,V}\alpha_2$ become equal. In
particular their images in every stalk are equal. This means that in the sheafification $\mathcal C^p$,
the cochains $\alpha_1$ and $\alpha_2$ actually become equal. Weird, right? Indeed the value of a cochain
on any particular “macroscopic” chain is not a local property, and so is obliterated by the
sheafification process. Only microscopic information remains.

Here is a fun example, which will better communicate what I want to say about the nature
of the $\mathcal C^p_{\mathrm{sing}}$. Consider $X=\mathbb R$ and set $A_p=0$ unless $p$ is of the form
$1/n$ for $n\in\mathbb Z_{\geq1}$, in which case $A_p=\mathbb C$. Let
$E=\coprod_{p\in\mathbb R}A_p$ with the obvious projection $\pi:E\to\mathbb R$, and consider
the sheaf of arbitrary sections
\[
s:\mathbb R\to E.
\]
Then the stalk at $0$ is the vector space of “tails”, namely the vector space of sequences
\[
(a_1,a_2,a_3,a_4,\ldots)
\]
modulo the equivalence relation $\sim$ that says
\[
(a_1,a_2,a_3,a_4,\ldots)\sim(b_1,b_2,b_3,b_4,\ldots)
\]
if there exists $N\geq1$ such that $a_n=b_n$ for all $n<N$. Every sequence is equivalent to one
that begins with an arbitrarily long string of zeroes. So it feels somehow impossible to write
down a nonzero element of this abelian group, and yet it is definitely not $\{0\}$.

\subsection{Image functors of sheaves}
\label{ringed-spaces-section-image-functors}
\label{ringed-spaces-section-pushforward}
\label{ringed-spaces-section-pullback}

Let $X,Y$ be topological spaces, and $f:X\to Y$ a continuous function. Let $\mathcal F$ be a sheaf
on $X$. It is possible to push $\mathcal F$ forward to $Y$ as follows: Let $V\subset Y$ be an open set, then
$f^{-1}(V)\subset X$ is open, we define a presheaf $f_*\mathcal F$ on $Y$ by
\[
f_*\mathcal F(V):=\mathcal F(f^{-1}(V)).
\]
\label{ringed-spaces-definition-pushforward}
One verifies that this is actually a sheaf.
The pushforward is always left exact, but not in general right exact. That is, let
\[
0\to\mathcal F_1\to\mathcal F_2\to\mathcal F_3\to0
\]
be a short exact sequence of sheaves on $X$, then we claim that
\[
0\to f_*\mathcal F_1\to f_*\mathcal F_2\to f_*\mathcal F_3.
\]
We have seen that injectivity/surjectivity can be tested on stalks, and germs are equivalence
classes of local sections. The thing about surjectivity is that no matter how small a neighbourhood
$V$ of $y\in Y$ we take, the preimage $f^{-1}(V)$ might be “big”, and so surjectivity
of $f_*\mathcal F_2(V)\to f_*\mathcal F_3(V)$ might fail for all neighbourhoods of $y$. The case of injectivity is
different: An injection of sheaves really is injective on every open, and therefore injectivity
of $\mathcal F_1\to\mathcal F_2$ implies injectivity of $f_*\mathcal F_1\to f_*\mathcal F_2$.
Now we describe pullback $f^{-1}\mathcal G$ of a sheaf $\mathcal G$ on $Y$. Since the image $f(U)\subset Y$ of open
$U\subset X$ is not open in general, how to turn a presheaf on $Y$ into a presheaf on $X$ is not so
clear. The definition, in fact, is the following: Let $\mathcal G$ be a presheaf on $Y$, and let $U\subset X$ be
open. Consider the set of all open $V\subset Y$ containing $f(U)$. We form the colimit over these,
defining
\[
f^{-1}\mathcal G(U)=\operatorname{colim}_{V\supset f(U)}\mathcal G(V).
\tag{5.3.1}
\]
If $\mathcal G$ is a sheaf, it happens that $f^{-1}\mathcal G$ might not be a sheaf. So for sheaves $\mathcal G$ on $Y$ we define
$f^{-1}\mathcal G$ to be the sheafification of the presheaf pullback described above.
\label{ringed-spaces-definition-pullback}
The definition of pullback of a sheaf is actually easier if we use the stalk-wise definition.
The first point to note is that the stalk itself is a special case of the pullback. Consider the
continuous map $i:*\to Y$ defined by $i(*)=p$. Then $i^{-1}\mathcal G$ by its definition is a colimit of
$\mathcal G(V)$ over neighbourhoods $V$ of $p$, which is exactly what the stalk $\mathcal G_p$ is.
If we consider the composition $*\to X\to Y$ now, then we easily see that
\[
(f^{-1}\mathcal G)_p=\mathcal G_{f(p)}.
\]
We note that the pullback functor $f^{-1}$ is exact. Since exactness can be tested at the level
of stalks, this is transparent from the observation above about stalks. See [15, Lemma
17.3.3(c)].
It turns out $(f^{-1},f_*)$ forms a left-right adjoint pair between categories of sheaves. Recall
this means there are natural isomorphisms
\[
\operatorname{Hom}(f^{-1}\mathcal G,\mathcal F)\cong\operatorname{Hom}(\mathcal G,f_*\mathcal F).
\]
If we unpack what the two sides mean, the isomorphism becomes clear. This is a good
exercise.
Here is a little more explanation. A morphism $\mathcal G\to f_*\mathcal F$ consists of a compatible system of
morphisms (indexed by open $V\subset Y$):
\[
\mathcal G(V)\to f_*\mathcal F(V)=\mathcal F(f^{-1}(V)).
\]
Now we wish to describe morphisms $f^{-1}\mathcal G\to\mathcal F$. Since we have defined
$f^{-1}(\mathcal G)$ as a sheafification, we should comment on this before starting.

\begin{remark}
The sheafification functor $(-)^+$, and the functor $(-)^\circ$ sending a sheaf to its
presheaf of sections (or rather, the functor of forgetting that the presheaf is a sheaf) is itself
an example of an adjoint pair. Namely for any presheaf $\mathcal G$ and sheaf $\mathcal F$ on $X$, we have
\[
\operatorname{Hom}(\mathcal G^+,\mathcal F)\cong\operatorname{Hom}(\mathcal G,\mathcal F^\circ).
\]
\end{remark}

This means that to specify a morphism $f^{-1}\mathcal G\to\mathcal F$ of sheaves is really the same thing as to
specify a morphism $f^{-1}\mathcal G\to\mathcal F$ where this time $f^{-1}\mathcal G$ stands for the presheaf (5.3.1).
Such a morphism consists of a compatible system of morphisms of vector spaces (indexed by open
$U\subset X$):
\[
\operatorname{colim}_{V\supset f(U)}\mathcal G(V)\to\mathcal F(U).
\]
Let us restrict to a fixed $U$ for a moment. By definition of colimit, a morphism as above is
the same thing as a collection of compatible morphisms
\[
\mathcal G(V)\to\mathcal F(U)
\]
running over all open $V\supset f(U)$. Switching perspective and viewing $V$ as fixed, this is the
data of a compatible system of morphisms running over all $U$ such that $f(U)$ is contained
in $V$. But $f(U)\subset V$ is the same thing as saying $U\subset f^{-1}(V)$. Such a system is completely
and uniquely determined by the morphism $\mathcal G(V)\to\mathcal F(f^{-1}(V))$. So we have proved the
adjointness property. Moreover we have seen that at its core the adjointness of $(f^{-1},f_*)$
hinges on the change of perspective between all open $V$ containing $f(U)$, and all open $U$
contained in $f^{-1}(V)$.
We reserve the symbol $\Gamma(X,-)$ for $\pi_*(-)$ the pushforward along the surjection $\pi:X\to *$
to a point. In other words $\Gamma(X,\mathcal F)$ is $\mathcal F(X)$. As for pushforwards in general, the functor
\[
\Gamma(X,-):\operatorname{Sh}(X)\to\operatorname{Vec}
\]
is left-exact. If we could show that $\operatorname{Sh}(X)$ has enough injectives (which it does), then we
could compute derived functors $R^n\Gamma(X,-)$. We'll do this in the next chapter.


\section{Glueing ringed spaces}
\label{ringed-spaces-section-glueing-ringed-spaces}

\noindent
Let's start with Hartshorne's exercise,
then I'll paste the Stacks Project version.

An interesting special case is when the
family $X_i$ is arbitrary, but the $U_{ij}$
and $\varphi_{ij}$ are all empty.
Then the scheme $X$ is called the
{\it disjoint union} of the $X_i$,
and is denoted $\coprod X_i$.

\begin{exercise}
\label{ringed-spaces-exercise-gluing-schemes}
\begin{reference}
\cite[Exercise II.2.12]{hart}.
\end{reference}
Let $\{X_i\}$ be a family of schemes
(possibly infinite).
For each $i \neq j$,
suppose given an open subset $U_{ij}
\subset X_i$,
and let it have the induced scheme
structure. Suppose also given for each
$i \neq j$ an isomorphism of schemes
$\varphi_{ij}:U_{ij}\to U_{ji}$
such that (1) for each $i,j$,
$\varphi_{ji}=\varphi_{ij}^{-1}$,
and (2) for each $i,j,k$,
$\varphi_{ij}(U_{ij}\cap U_{ik}
=U_{ji}\cap U_{jk}$ and
$\varphi_{ik}= \varphi_{jk}\circ \varphi_{ij}$
on $U_{ij}\cap U_{ik}$.
Then show that there is a scheme $X$,
together with morphisms
$\psi_i:X_i \to X$ for each $i$,
such that (1) $\psi_i$ is an isomorphism
of $X_i$ onto an open subsceme of $X$,
(2) the $\psi_i(X_i)$ cover $X$,
(3) $\psi_i(U_{ij}) = \psi_i(X_i)
\cap\psi_j(X_j)$ and
(4) $\psi_i = \psi_j \circ \varphi_{ij}$
on $U_{ij}$.
\end{exercise}

\begin{proof}
Step 1. Define the underlying topological 
space:
$$
X = (\coprod_{i} X_i)/\sim
$$
where $x_i \in X_i$ and $x_j \in X_j$
are identified if $x_i \in U_{ij}$,
$x_j \in U_{ji}$ and $\varphi_{ij}(x_i)=x_j$.
Equip this space with the quotient topology.
This defines maps of topological spaces
(i.e. continuous maps)
$$
\psi_i: X_i \to X,\qquad 
x_i \mapsto [x_i]
$$
which are just the inclusions followed
by the quotient projections.
To prove $\psi_i$ is an embedding
we must show it's a homeomorphism, which
follows from it being bijective and open.
For bijective we only need to prove
injectivity. It might look immediate since
$X_i \hookrightarrow \coprod_j X_j$ is an
embedding, and projecting would only glue
points of $X_i$ with points of other $X_j$.
The nuance is that there could be
``zigzags'': $x_i$ from $X_i$ glues to
$x_j \in X_j$, which glues to $x_k \in X_k$
and that one back to $x_i' \in X_i$,
breaking injectivity.

To fix this let's be more precise:
$x_i$ being glued to $x_j$ means that
$\varphi_{ij}(x_i) = x_j$.
Similarly $\varphi_{j k}(x_j) = x_k$
and $\varphi_{ki}(x_k)=x_i'$.
But the cocycle condition says
$\varphi_{ik}=\varphi_{jk}\varphi_{ij}$,
which means that both $x_i'$ and $x_i$
map to $x_k$. But $\varphi_{ik}$ is a
homeomorphism, so it is injective, so
$x_i'=x_i$.

To conclude injectivity from this argument
suppose that $x_i, x_i' \in X_i$
are such that $\psi_i(x_i) = \psi_i(x_i')$.
This means that there is a zigzag as above,
involving possibly more than three maps,
since no two points of $X_i$ are
glued directly in the quotient.
Then we would just apply the cocycle
condition inductively and conclude as above.

Having proved injectivity, it still
remains to show that $\psi_i$ is an open map.
So we pick an open set $V$ of
$X_i$ and check that its image is open.
Since $X$ has the quotient topology,
this means that $\psi_i(V)$ is open
if its preimage under the projection map
$q$ is open in the coproduct $\coprod X_i$.
The nuance is that such preimage
may be spread over other $X_j$, 
$j\neq i$. However, since $\varphi_{ij}$
are homeomorphisms, the union of all
the ``spread-over'' parts remains a union
of (arbitrarily many) open sets, which is
open.

Step 2. Define the structure sheaf of $X$.
Let $U$ be an open set of $X$. Define
a presheaf that maps $U$ to 
$$
\mathcal{O}_X(U)
:=
(\bigotimes \mathcal{O}_{X_i}(U \cap X_i))
/\sim
$$
(since coproduct in the category
of rings is direct sum
by Exercise
\ref{commutative-algebra-exercise-coproduct-in-rings})
where two sections $s_i \in
\mathcal{O}_{X_i}(U \cap \psi_i(X_i))$
and $s_j \in \mathcal{O}_{X_j}(U \cap
\psi_j(X_j))$ 
are identified… it's not clear yet
why we need the identification.

To show this is a sheaf suppose we have
an inclusion $U \hookrightarrow V$.
Then we have a restriction map
$\rho_{U,V}:\mathcal{O}_X(V)
\to \mathcal{O}_X(U)$
given by the restrictions of
every section obtained after projecting
to each $\mathcal{O}_{X_i}(U 
\cap \psi_i(X_i))$.
This preserves compositions
and maps the identity
morphism to the identity morphism
after projecting to each summand,
and thus also on $X$.
More explicitly: for the identity
$\text{id}:U \to U$ we obtain a map
of sections $\rho_{U,U}$ putting
$(s_i) \mapsto (\rho^{X_i}_{U \cap X_i,
U\cap X_i}s_i)$
but the latter maps are the identity
because each $\mathcal{O}_{X_i}$ is indeed
a presheaf. Composition and the 
glueing condition for sheaves are
confirmed similarly.

Step 3. The conditions (1) to (4)
from the exercise are confirmed by
construction.
\end{proof}

\noindent
This section comes from Stacks Project
tag 
\href{https://stacks.math.columbia.edu/tag/01JA}{01JA}. See also

Let $I$ be a set.
For each $i \in I$ let $(X_i, \mathcal{O}_i)$ be
a locally ringed space. (Actually the construction that
follows works equally well for ringed spaces.)
For each pair $i, j \in I$ let $U_{ij} \subset X_i$
be an open subspace.
For each pair $i, j \in I$, let
$$
\varphi_{ij} : U_{ij} \to U_{ji}
$$
be an isomorphism of locally ringed spaces.
For convenience we assume that $U_{ii} = X_i$.
For each triple $i, j, k \in I$ assume that
\begin{enumerate}
\item we have
$\varphi_{ij}^{-1}(U_{ji} \cap U_{jk}) =  U_{ij} \cap U_{ik}$, and
\item the diagram
$$
\xymatrix{
U_{ij} \cap U_{ik} \ar[rr]_{\varphi_{ik}} \ar[rd]_{\varphi_{ij}} & &
U_{ki} \cap U_{kj} \\
& U_{ji} \cap U_{jk} \ar[ru]_{\varphi_{jk}}
}
$$
is commutative.
\end{enumerate}
(In particular, it follows that $\varphi_{ii} = \text{id}_{X_i}$.)
Let us call a collection
$(I, (X_i)_{i\in I}, (U_{ij})_{i, j\in I}, (\varphi_{ij})_{i, j\in I})$
satisfying the conditions above a glueing data.

\begin{lemma}
\label{ringed-spaces-lemma-glue}
\begin{slogan}
If you have two locally ringed spaces, and a subspace of the first one
is isomorphic to a subspace of the other, then you can glue them together
into one big locally ringed space.
\end{slogan}
Given any glueing data of locally ringed spaces there
exists a locally ringed space $X$ and open subspaces
$U_i \subset X$ together with isomorphisms
$\varphi_i : X_i \to U_i$ of locally ringed spaces such that
\begin{enumerate}
\item $X=\bigcup_{i\in I} U_i$,
\item $\varphi_i(U_{ij}) = U_i \cap U_j$, and
\item $\varphi_{ij} =
\varphi_j^{-1}|_{U_i \cap U_j} \circ \varphi_i|_{U_{ij}}$.
\end{enumerate}
The locally ringed space $X$ is characterized by the following
mapping properties: Given a locally ringed space $Y$ we have
\begin{eqnarray*}
\Mor(X, Y) & = & \{ (f_i)_{i\in I} \mid
f_i : X_i \to Y, \ f_j \circ \varphi_{ij} = f_i|_{U_{ij}}\} \\
f & \mapsto & (f|_{U_i} \circ \varphi_i)_{i \in I} \\
\Mor(Y, X) & = &
\left\{
\begin{matrix}
\text{open covering }Y = \bigcup\nolimits_{i \in I} V_i\text{ and }
(g_i : V_i \to X_i)_{i \in I}
\text{ such that}\\
g_i^{-1}(U_{ij}) = V_i \cap V_j
\text{ and }
g_j|_{V_i \cap V_j} = \varphi_{ij} \circ g_i|_{V_i \cap V_j}
\end{matrix}
\right\} \\
g & \mapsto &
V_i = g^{-1}(U_i), \ g_i = \varphi_i^{-1} \circ g|_{V_i}
\end{eqnarray*}
\end{lemma}

\begin{proof}
We construct $X$ in stages.
As a set we take
$$
X = (\coprod X_i) / \sim.
$$
Here given $x \in X_i$ and $x' \in X_j$ we say
$x \sim x'$ if and only if $x \in U_{ij}$, $x' \in U_{ji}$
and $\varphi_{ij}(x) = x'$. This is an equivalence relation
since if $x \in X_i$, $x' \in X_j$, $x'' \in X_k$, and $x \sim x'$ and
$x' \sim x''$, then $x' \in U_{ji} \cap U_{jk}$, hence by condition (1) of
a glueing data also $x \in U_{ij} \cap U_{ik}$ and
$x'' \in U_{ki} \cap U_{kj}$ and by condition (2)
we see that $\varphi_{ik}(x) = x''$. (Reflexivity and symmetry
follows from our assumptions that $U_{ii} = X_i$ and
$\varphi_{ii} = \text{id}_{X_i}$.)
Denote $\varphi_i : X_i \to X$
the natural maps. Denote $U_i = \varphi_i(X_i) \subset X$.
Note that $\varphi_i : X_i \to U_i$ is a bijection.

\medskip\noindent
The topology on $X$ is defined by the rule that
$U \subset X$ is open if and only if $\varphi_i^{-1}(U)$
is open for all $i$. We leave it to the reader to verify
that this does indeed define a topology.
Note that in particular $U_i$ is open since $\varphi_j^{-1}(U_i)
= U_{ji}$ which is open in $X_j$ for all $j$.
Moreover, for any open set $W \subset X_i$ the image
$\varphi_i(W) \subset U_i$ is open because
$\varphi_j^{-1}(\varphi_i(W)) = \varphi_{ji}^{-1}(W \cap U_{ij})$.
Therefore $\varphi_i : X_i \to U_i$ is a homeomorphism.

\medskip\noindent
To obtain a locally ringed space we have to construct the
sheaf of rings $\mathcal{O}_X$. We do this by glueing the
sheaves of rings $\mathcal{O}_{U_i} := \varphi_{i, *} \mathcal{O}_i$.
Namely, in the commutative diagram
$$
\xymatrix{
U_{ij} \ar[rr]_{\varphi_{ij}} \ar[rd]_{\varphi_i|_{U_{ij}}}
& &
U_{ji} \ar[ld]^{\varphi_j|_{U_{ji}}} \\
& U_i \cap U_j &
}
$$
the arrow on top is an isomorphism of ringed spaces,
and hence we get unique isomorphisms of sheaves of rings
$$
\mathcal{O}_{U_i}|_{U_i \cap U_j}
\longrightarrow
\mathcal{O}_{U_j}|_{U_i \cap U_j}.
$$
These satisfy a cocycle condition as in Sheaves,
Section \ref{ringed-spaces-sheaves-section-glueing-sheaves}.
By the results of that section we obtain a sheaf of rings
$\mathcal{O}_X$ on $X$ such that $\mathcal{O}_X|_{U_i}$
is isomorphic to $\mathcal{O}_{U_i}$ compatibly with
the glueing maps displayed above.
In particular $(X, \mathcal{O}_X)$ is a locally ringed
space since the stalks of $\mathcal{O}_X$ are equal
to the stalks of $\mathcal{O}_i$ at corresponding
points.

\medskip\noindent
The proof of the mapping properties is omitted.
\end{proof}

\begin{lemma}
\label{ringed-spaces-lemma-glue-schemes}
\begin{slogan}
Schemes can be glued to give new schemes.
\end{slogan}
In Lemma \ref{ringed-spaces-lemma-glue} above, assume that all
$X_i$ are schemes. Then the resulting locally ringed
space $X$ is a scheme.
\end{lemma}

\begin{proof}
This is clear since each of the $U_i$ is a scheme
and hence every $x \in X$ has an affine neighbourhood.
\end{proof}

\noindent
It is customary to think of $X_i$ as an open subspace of
$X$ via the isomorphisms $\varphi_i$. We will do this in
the next two examples.



\nocite{*}

% --- End ringed-spaces.tex ---

% --- Begin schemes.tex ---
\chapter{Schemes}
\label{chapter-schemes}





\section{Schemes}
\label{schemes-section-schemes}

For a definition of presheaf,
see Categories, Definition
\ref{categories-definition-presheaf}.

\begin{definition}
\label{schemes-definition-sheaf}
Let $X$ be a topological space.
\begin{enumerate}
\item A {\it sheaf $\mathcal{F}$ of sets on
$X$} is a presheaf
of sets which satisfies the following
additional property: Given
any open covering $U = \bigcup_{i \in I} U_i$
and any collection
of sections $s_i \in \mathcal{F}(U_i)$, $i
\in I$ such that
$\forall i, j\in I$
$$
s_i|_{U_i \cap U_j} = s_j|_{U_i \cap U_j}
$$
there exists a unique section $s \in
\mathcal{F}(U)$ such that
$s_i = s|_{U_i}$ for all $i \in I$.
\item A {\it morphism of sheaves of sets} is
simply a
morphism of presheaves of sets.
\item The category of sheaves of sets on $X$
is denoted
$\Sh(X)$.
\end{enumerate}
\end{definition}


\medskip\noindent
Let $X$ be a topological space. Let $x \in X$
be a point.
Let $\mathcal{F}$ be a presheaf of sets on
$X$.
The {\it stalk of $\mathcal{F}$ at $x$} is
the set
$$
\mathcal{F}_x
=
\colim_{x\in U} \mathcal{F}(U)
$$
where the colimit is over the set of open
neighbourhoods
$U$ of $x$ in $X$. The set of open
neighbourhoods is
partially ordered by (reverse) inclusion:
We say $U \geq U' \Leftrightarrow U \subset
U'$.
The transition maps in the system are
given by the restriction maps of
$\mathcal{F}$.
See Categories, Section
\ref{categories-section-posets-limits}
for notation and terminology regarding
(co)limits over systems.
Note that the colimit is a directed colimit.
Thus it is easy to describe $\mathcal{F}_x$.
Namely,
$$
\mathcal{F}_x
=
\{
(U, s)
\mid
x\in U, s\in \mathcal{F}(U)
\}/\sim
$$
with equivalence relation given by $(U, s)
\sim (U', s')$ if and only if
there exists an open $U'' \subset U \cap U'$
with $x \in U''$ and
$s|_{U''} = s'|_{U''}$. Given a pair $(U, s)$
we sometimes denote
$s_x$ the element of $\mathcal{F}_x$
corresponding to the equivalence
class of $(U, x)$. We sometimes use the
phrase
``image of $s$ in $\mathcal{F}_x$'' to denote
$s_x$.
For example, given two pairs $(U, s)$ and
$(U', s')$ we sometimes
say ``$s$ is equal to $s'$ in
$\mathcal{F}_x$'' to indicate
that $s_x = s'_x$. Other authors use the
terminology
``germ of $s$ at $x$''.

\section{Closed immersions}
\label{schemes-section-immersions}

\noindent
Closed immersions of ringes spaces are in
\href{https://stacks.math.columbia.edu/tag/%
01C1}{01C1}.

\begin{definition}
\label{schemes-definition-closed-immersion}
A map of ringed spaces
$i:(Z,\mathcal{O}_Z) \to (X,\mathcal{O}_X)$
is a {\it closed immersion}
if $i$ is injective,
its image is closed (I think these two are
``closed immersion of topological spaces'')
and the induced map $\mathcal{O}_X \to
i_*\mathcal{O}_Z$
is surjective; denote its kernel by
$\mathcal{I}$.
To make sure that if $(X,\mathcal{O}_X)$ is a
scheme
then so is $(Z,\mathcal{O}_Z)$ with add the
extra condition
that the $\mathcal{O}_X$-module $\mathcal{I}$
is
locally generated by sections.
\end{definition}

\medskip\noindent
\section{Affine schemes}
\label{schemes-section-affine-schemes}

\noindent
The main objective of this section is
to define the structure sheaf of the
spectrum of a ring, which
was already introduced in 
Commutative Algebra 
\ref{commutative-algebra-section-spectrum}.
The resulting pair
$(\Spec R, \mathcal{O}_{\Spec R}$
will be an affine scheme.

\section{Reduced schemes}
\label{schemes-section-reduced}

\noindent
So far a reduced scheme,
for me,
is just a scheme where there are
no ``repetitions'':
$(x^2)$ is not reduced because
it gives the same information as $(x)$.
Algebraically this means that the
coordinate ring is reduced,
which is defined by asking
that there are no nilpotent elements,
i.e. elements such that some power vanishes.
(There's no definition of reduced ring in
Stacks Project
since it's taken as part of the basic
algebraic knowledge).

The definition is on stalks but the first
lemma
shows it's equivalent to define it on any
open set.
This is in virtue of some ``injectivity''
lemma
from the ring at $U$ to the product of all
stalks parametrized in $U$;
essentially that if a section vanishes when
projected to all
the stalks then it's zero.

\begin{definition}
\label{schemes-definition-reduced}
Let $X$ be a scheme. We say $X$ is {\it
reduced} if every local ring
$\mathcal{O}_{X, x}$ is reduced.
\end{definition}

\begin{proof}
Assume that $X$ is reduced.
Let $f \in \mathcal{O}_X(U)$ be a section
such that $f^n = 0$.
Then the image of $f$ in $\mathcal{O}_{U, u}$
is zero for
all $u \in U$. Hence $f$ is zero, see
Sheaves, Lemma
\ref{schemes-sheaves-lemma-sheaf-subset-stalks}.
Conversely, assume that $\mathcal{O}_X(U)$ is
reduced
for all opens $U$. Pick any nonzero element
$f \in \mathcal{O}_{X, x}$.
Any representative $(U, f \in
\mathcal{O}(U))$  of $f$ is nonzero and
hence not nilpotent. Hence $f$ is not
nilpotent in $\mathcal{O}_{X, x}$.
\end{proof}

\section{Normalization}
\label{schemes-section-normalization}

\noindent
It might be interest to check
the construction of
normalization in Stacks Project:
it is not done via gluing.

\begin{definition}
\label{schemes-definition-normal-scheme}
\begin{reference}
\cite[Exercise II.3.8]{hart}
\end{reference}
A scheme is {\it normal} 
if all of its local rings are integrally
closed.
\end{definition}

\begin{exercise}[Existence of normalization]
\label{schemes-exercise-normalization-schemes}
\begin{reference}
\cite[Exercise II.3.8]{hart}
\end{reference}
Let $X$ be an integral scheme.
For each open affine subset
$U = \Spec A$ of $X$,
let $\tilde{A}$ be the integral closure
of $A$ in its quotient field,
and let $\tilde{U} = \Spec \tilde{A}$.
Show that one can glue the schemes $\tilde{U}$
to obtain a normal integral scheme
$\tilde{X}$, called
the {\it normalization} of $X$.
\end{exercise}

\begin{proof}
The first challenge is to define
the $U_{ij}$. If $U_i = \Spec(\widetilde{A_i})$,
then the most natural is to define
$U_{ij}$ as the spectrum of the 
integral closure of the polynomial ring
modulo the
ideal $I_i + I_j$, which is the ideal
that corresponds
to the intersection of the original
spectrums. This looks like it needs
further reading.
\end{proof}

\medskip\noindent
Here's AI's solution (probably).

Let
$$
\nu : \widetilde C \longrightarrow C
$$
be the normalization of $C$. We recall its
construction. Choose an affine open covering
$U_i=\Spec(A_i)$ of $C$. Since $C$ is
irreducible, all $A_i$ have the same
function field $K(C)$. Let $\overline A_i$
be the integral closure of $A_i$ in $K(C)$.

If $U_{ij}=U_i\cap U_j$, then, after
refining the cover if necessary, we may
assume
$$
U_{ij}=D(f)\subset U_i
$$
for some $f\in A_i$. Its coordinate ring is
$(A_i)_f$. The integral closure of
$(A_i)_f$ in $K(C)$ is
$$
(\overline A_i)_f,
$$
because integral closure commutes with
localization. Thus the normalization over
$U_{ij}$ is
$$
\Spec((\overline A_i)_f).
$$
Doing the same computation from $U_j$ gives
the same subring of $K(C)$, namely the
integral closure of $\mathcal O_C(U_{ij})$
inside $K(C)$. Hence the affine schemes
$\Spec(\overline A_i)$ have canonical
identifications over the intersections
$U_i\cap U_j$. These identifications satisfy
the cocycle condition, since they are all
induced by equality inside $K(C)$.

Therefore the affine schemes
$\Spec(\overline A_i)$ glue to a curve
$\widetilde C$, and the inclusions
$$
A_i\subset \overline A_i
$$
glue to a finite birational morphism
$$
\nu:\widetilde C\to C.
$$
The curve $\widetilde C$ is normal by
construction.


\section{Proj}
\label{schemes-section-proj}

\noindent
(Stacks Project
\href{https://stacks.math.columbia.edu/tag/%
01M3}{01M3},
Proj of a graded ring.)
This section may be taken as the preamble to
the question: what is projective space
anyway?
Yes, projective space is Proj of a graded
ring $S$.
And yes, projective space is a scheme.
So in particular it's locally affine.

(Stacks Project
\href{https://stacks.math.columbia.edu/tag/%
01MX}{01MX},
Functoriality of Proj.)
This is reminiscent of the olden days.
The point is that Proj is covariant,
i.e. for rings  $S' \subset S$ we have
$\text{Proj}S \subset \text{Proj}S$
(under the right conditions,
which is the point of this section),
meaning that, indeed,
{\bf Projs are projective schemes}.






\medskip\noindent
Noe let's discuss invertible sheaves on Proj
(twists!).

\noindent
As for the construction of the twisted
sheaves:
for a projective space $\text{Proj} S$,
we just compute the twisted $S$ module
$S(m)$ by the formula $S(m)_d=S_{m+d}$
and turn that into a sheaf
$\widetilde{S(m)}$.

Here's the construction of the twisted
sheaves $\mathcal{O}_X(n)$.
The point is that there is a good way
to pass from an $S$-module $M$ to an
$\mathcal{O}_X$-module, called
$\widetilde{M}$.
So basically you just define the twists
by $M(d)_n=M_{n+d}$ and apply that
construction.

Pick an element of degree $d$.
Then think of the module generated by this
element.
The least degree piece of such a module
is the degree $d$ piece of the original
module,
that is we have $M(-d)$.










\section{Restricting sheaves to subschemes}
\label{schemes-section-restricting-sheaves}

\noindent
Upshot about restricting sheaves to
subschemes:
the sheaves on $Z$ correspond essentially
(in the categorical sense)
to sheaves on $X$ with support on $Z$.

More explicitly, from
\href{https://stacks.math.columbia.edu/tag/%
01QY}{01QY}:
the pushforward of the inclusion (a closed
immersion of schemes)
$i:Z \to X$, with kernel on induced map
$\mathcal{O}_Z \to \mathcal{O}_X$
denoted by $I$, is an exact, fully faithful
functor
$$
i_*:\QCoh(\mathcal{O}_Z) \to
\QCoh(\mathcal{O}_X)
$$
with essential image the quasi-coherent
$\mathcal{O}_X$ modules
$\mathcal{G}$ such that
$\mathcal{I}\mathcal{G}=0$.

It makes sense to think of this operation as
$\mathcal{F}|_Z=i^*\mathcal{F}=\mathcal{F}
\otimes\mathcal{O}_Z$
for $\mathcal{F}\in\QCoh(X)$.
Indeed: tensoring by the ideal sheaf would
kill the information on $Z$,
and tensoring with the structure sheaf does
the opposite
--- it's the sheaves on $X$ that are no fun
outside $Z$,
so basically sheaves on $Z$.

\medskip\noindent
Here's a more detailed discussion:

In understanding the cohomology of sheaves
when
we restrict them to submanifolds (see Complex
Geometry
Exercise
\ref{complex-geometry-exercise-%
deformations-curve-in-k3},
I found Stacks Project
\href{https://stacks.math.columbia.edu/tag/%
01AW}{01AW}
section on closed immersions and abelian
sheaves.
First we have a few propositions stating that
sheaves on the subscheme $Z \subset X$ can be
pushed
to sheaves on $X$.

And then there's the converse:
what I would like to call the restriction
functor $\mathcal{H}$
but rather is called the {\it sub-module of
sections with support in $Z$}.
In Remark
\href{https://stacks.math.columbia.edu/tag/%
01AY}{01AY}
(and then in Remark
\href{https://stacks.math.columbia.edu/tag/%
0G6N}{0G6N}
for ringed spaces; these seem to be analogue
constructions…)
we see that for an abelian sheaf
$\mathcal{F}$ on $X$
we can define an abelian sheaf
$\mathcal{H}_Z(\mathcal{F})$
which although is a sheaf on $X$
we can just think it's a sheaf on $Z$ -
that's beacause
$$
\mathcal{H}_Z(\mathcal{F})(U)
=\{s \in \mathcal{F}(U))| \text{the support
of $s$ is contained in }Z\cap U\}.
$$
We can just think that it's the functions
that vanish outside the subvariety.
Which doesn't really make sense:
a function vanishing outside the submanifold
will most likely
vanish also in the submanifold! Just by
continuity,
the submanifold is in the closure of the open
set!
I guess what I mean to say is that I don't
mind if
the functions do not vanish outside $Z$…
I just want to consider them as functions on
$Z$…

\medskip\noindent
Here's another attempt in understanding how
to restrict sheaves to submanifolds.
First and foremost, lacking a reference this,
I understand the symbol $\mathcal{F}|_Z$ for
a sheaf
on $X$ and a subscheme  $Z \subset X$ to mean
the
pullback of $\mathcal{F}$ by the inclusion.
Then look at
\href{https://stacks.math.columbia.edu/tag/%
01QY}{01QY}
(more basically, as above,
\href{https://stacks.math.columbia.edu/tag/%
01AW}{01AW}…
in fact there appear to be several version of
this statement)
to recall that sheaves of the subscheme $Z$
should correspond
to sheaves on $X$ with support on $Z$.
{\bf So essentially $\mathcal{F}|_Z$ should
be a version of $\mathcal{F}$
but with support on $Z$.}
The key point here is that this is non other
than
$\mathcal{F} \otimes \mathcal{O}_Z$.
Why? Because when multiply with the ideal
sheaf of $Z$ we obtain
$(\mathcal{F} \otimes
\mathcal{O}_Z)\mathcal{I}=0$
since sections of $\mathcal{I}$
vanish in
$\mathcal{O}_Z=\mathcal{O}_X/\mathcal{I}$.
That is, the support of $\mathcal{F} \otimes
\mathcal{O}_Z$ indeed
contained in $Z$: if the sections become zero
when
multiplied by sections of $\mathcal{I}$
(sections of $\mathcal{I}$ are only nonzero
outside $Z$)
it means that all the action was inside $Z$.
that outside $Z$

I think the upshot here is that structure
sheaves
act as some sort of ``contour'' for the
geometric object:
they vanish {\it outside} the geometric
object,
they tell us what is the shape of the object
by describing its surroundings, rather than
its interior.
This is exactly how algebraic/analytic sets
are defined,
and this why, when we want to restrict a
sheaf
to a subscheme, we tensor with the structure
sheaf:
tensor by the ideal sheaf would kill the
information on $Z$,
and tensoring with the structure sheaf does
the opposite.













\section{Flat modules}
\label{schemes-section-flat-modules}

\noindent
The essential technical property
for for defining flatness
is the preservation of exact sequences.
Right-exactness is true in general;
it follows from currying in category
of commutative rings.
But the functor $-\otimes_R N$ does not
preserve injectivity of maps---
that's the point of flatness.

\begin{lemma}[Internal Hom for R-modules]
\label{schemes-lemma-hom-from-tensor-product}
For any three $R$-modules $M, N, P$,
$$
\Hom_R(M \otimes_R N, P) \cong \Hom_R(M,
\Hom_R(N, P))
$$
\end{lemma}

\begin{proof}
An $R$-linear map $\hat{f}\in \Hom_R(M
\otimes_R N, P)$ corresponds to an
$R$-bilinear map $f : M \times N \to P$. For
each $x\in M$ the mapping $y\mapsto f(x, y)$
is $R$-linear by the universal
property. Thus $f$ corresponds to a
map $\phi_f : M \to \Hom_R(N, P)$. This map
is $R$-linear since
$$
\phi_f(ax + y)(z) =
f(ax + y, z) = af(x, z)+f(y, z) =
(a\phi_f(x)+\phi_f(y))(z),
$$
for all $a \in R$, $x \in M$, $y \in M$ and
$z \in N$. Conversely, any
$f \in \Hom_R(M, \Hom_R(N, P))$ defines an
$R$-bilinear
map $M \times N \to P$, namely $(x, y)\mapsto
f(x)(y)$.
So this is a natural one-to-one
correspondence between the
two modules
$\Hom_R(M \otimes_R N, P)$ and $\Hom_R(M,
\Hom_R(N, P))$.
\end{proof}

\begin{lemma}[Tensor product is right exact]
\label{schemes-lemma-tensor-product-exact}
Let
$$
\begin{aligned}
M_1\xrightarrow{f} M_2\xrightarrow{g} M_3 \to
0
\end{aligned}
$$
be an exact sequence of $R$-modules and
homomorphisms, and let $N$ be any
$R$-module. Then the sequence
\begin{equation}
\label{schemes-equation-2ndex}
M_1\otimes N\xrightarrow{f \otimes 1}
M_2\otimes N \xrightarrow{g \otimes 1}
M_3\otimes N \to 0
\end{equation}
is exact. In other words, the functor $-
\otimes_R N$ is
{\it right exact}, in the sense that
tensoring
each term in the original right exact
sequence preserves the exactness.
\end{lemma}

\begin{proof}
For every $R$-module $P$
we apply the functor $\Hom(-, \Hom(N, P))$ to
the first exact
sequence. We obtain
$$
0 \to
\Hom(M_3, \Hom(N, P)) \to
\Hom(M_2, \Hom(N, P)) \to
\Hom(M_1, \Hom(N, P))
$$
which is exact by Lemma \ref{schemes-lemma-hom-exact}
(1).
By Lemma \ref{schemes-lemma-hom-from-tensor-product}
this becomes the sequence
$$
0 \to \Hom(M_3 \otimes N, P) \to
\Hom(M_2 \otimes N, P) \to \Hom(M_1 \otimes
N, P)
$$
which is therefore also exact. Then using
Lemma \ref{schemes-lemma-hom-exact} (1) again, we
arrive at the desired exact sequence.
\end{proof}

\begin{remark}
\label{schemes-remark-tensor-product-not-exact}
However, tensor product does NOT preserve
exact sequences in general.
In other words, if $M_1 \to M_2 \to M_3$ is
exact, then it is not necessarily true that
$M_1 \otimes N \to M_2 \otimes N \to M_3
\otimes N$
is exact for arbitrary $R$-module $N$.
\end{remark}

\begin{example}
\label{schemes-example-tensor-product-not-exact}
Consider the injective map $2 : \mathbf{Z}\to
\mathbf{Z}$
viewed as a map of $\mathbf{Z}$-modules.
Let $N = \mathbf{Z}/2$. Then the induced map
$\mathbf{Z} \otimes \mathbf{Z}/2 \to
\mathbf{Z} \otimes \mathbf{Z}/2$
is NOT injective. This is because for
$x \otimes y\in \mathbf{Z} \otimes
\mathbf{Z}/2$,
$$
(2 \otimes 1)(x \otimes y) = 2x \otimes y = x
\otimes 2y = x \otimes 0 = 0
$$
Therefore the induced map is the zero map
while $\mathbf{Z} \otimes N\neq 0$.
\end{example}

\begin{definition}
\label{schemes-definition-flat-module}
For $R$-modules $N$, if the
functor $-\otimes_R N$ is exact, i.e.
tensoring
with $N$ preserves all exact
sequences, then $N$ is said to be {\it flat}
$R$-module.
We will discuss this later in Section
\ref{schemes-section-flat}.
\end{definition}

\medskip\noindent
Epiphany:
in the category of commutative rings
pushout is tensor product.
So think of $\Spec $ as a functor
from $\mathit{CRing}^{\text{op}}$ to $\Sch$,
then pushout goes to pullback,
and what's an example of a pullback?
Fibre!
So, the coordinate ring of a fiber
is essentially given by the residue field
at the point that parametrizes it!
(tensored with the coordinate ring
of the deformation space).

There is a lot of information on Stacks
Project about flatness.
It looks like the heart of the concept is
captured in the commutative-algebraic notion
of preserving
exact sequences:

\begin{definition}
\label{schemes-definition-flat}
Let $R$ be a ring.
\begin{enumerate}
\item An $R$-module $M$ is called {\it flat}
if whenever
$N_1 \to N_2 \to N_3$ is an exact sequence of
$R$-modules
the sequence $M \otimes_R N_1 \to M \otimes_R
N_2 \to M \otimes_R N_3$
is exact as well.
\item An $R$-module $M$ is called {\it
faithfully flat} if the
complex of $R$-modules
$N_1 \to N_2 \to N_3$ is exact if and only if
the sequence $M \otimes_R N_1 \to M \otimes_R
N_2 \to M \otimes_R N_3$
is exact.
\item A ring map $R \to S$ is called {\it
flat} if
$S$ is flat as an $R$-module.
\item A ring map $R \to S$ is called {\it
faithfully flat} if
$S$ is faithfully flat as an $R$-module.
\end{enumerate}
\end{definition}

\medskip\noindent
Recall that a module $M$ over a ring $R$ is
{\it flat} if the functor
$-\otimes_R M : \text{Mod}_R \to
\text{Mod}_R$ is exact. A ring map
$R \to A$ is said to be {\it flat} if $A$ is
flat as an $R$-module.



% --- End schemes.tex ---

% --- Begin algebraic-geometry.tex ---
\chapter{Algebraic Geometry}
\label{chapter-algebraic-geometry}





\section{Algebraic sets}
\label{algebraic-geometry-section-algebraic-sets}

\noindent
The {\it affine space} $\mathbb{A}^n$ is the
Cartesian
product of the field $n$ times.
An {\it algebraic set} is the zero set of an
ideal
in $k[x_1,…,x_n]$.

$V(I)$ and lots of properties…

\begin{lemma}
\label{algebraic-geometry-lemma-irreducible-variety-prime-ideal}
An algebraic set is irreducible if and only
if
its defining ideal is prime.
\end{lemma}

\begin{proof}
First notice that in
the case that $Y$ is an irreducible
subvariety,
this ideal is prime: if $fg \in I_Y(U)$,
then $V(f) \cup  V(g) \supset Y\cap U$
so that $(V(f) \cup  V(g))\cap Y = Y\cap U$,
and then it must be true that $V(f)\cap
Y\subseteq V(g)\cap Y$
or that $V(f)\cap Y \supseteq V(g)\cap Y \cap
U$
since $Y$ is irreducible, and thus, say,
$V(f)\supseteq Y \cap U$,
that is, $I(V(f))\subseteq Y \cap U$
so that $f \in I(Y \cap U)$.
\end{proof}

\begin{theorem}[Weak Nullstellensatz]
\label{algebraic-geometry-theorem-weak-nullstellensatz}
(Misha's formulation.) Let $A \subset
\mathbb{C}^n$
be an affine variety and $\mathcal{O}_A$ the
ring of polynomial
functions on $A$.
Then every maximal ideal in $A$
is an ideal of a point $a \in A$.
\end{theorem}

\noindent
The idea behind the radical idea is that
we can take roots of elements:

\begin{definition}
\label{algebraic-geometry-definition-radical-ideal}
$I \subset R$ is {\it radical}
if for every $x \in R$ we have $x^n \in I
\implies x \in I$.
\end{definition}

\begin{definition}
\label{algebraic-geometry-definition-nilpotent}
An element $a \in R$ is called {\it
nilpotent}
if $a^n = 0$.
\end{definition}

\begin{exercise}
\label{algebraic-geometry-exercise-no-nilpotents-iff-radical}
$A/I$ has no nilpotents if and only if $I$ is
a radical ideal.
\end{exercise}

\noindent
We may understand the Nullstellensatz
as saying that algebraic subsets of an affine
variety
are the radical ideals of its coordinate
algebra.

\begin{theorem}[Nullstellensatz]
\label{algebraic-geometry-theorem-nullstellensatz}
$$
I(V(I))=\sqrt{I}
$$
where $I$ is an ideal of $k[x_1,…,x_n]$.
\end{theorem}

\begin{proof}
Prove the Weak Nullstellensatz first
and then use Rabinowitsch trick
(which uses localization).
\end{proof}

\noindent
(Prove this:)
$k[x_1,…,x_n]/I$ has no nilpotents if and
only if
$I = \sqrt{I}$.

\medskip\noindent
Now we discuss how to decompose a variety in
irreducible components.

\begin{definition}
\label{algebraic-geometry-definition-irreducible-variety}
An algebraic set is {\it irreducible}
if it cannot be decomposed as a union of two
properly
contained algebraic sets.
\end{definition}

\begin{lemma}
\label{algebraic-geometry-lemma-irred-iff-coord-ring-reduc}
An affine variety is irreducibly if and only
if
its coordinate ring has no zerodivisors.
\end{lemma}

\begin{proof}
A zero divisor gives an expression $fg=0$.
\end{proof}

\begin{exercise}
\label{algebraic-geometry-exercise-surje-map-irred-varie}
If $\varphi:A \to B$ is a surjective morphism
of varieties
and $B$ is irreducible, so is $A$.
\end{exercise}


To prove the decomposition in
irreducible ideals we must discuss
Noetherianity.

\begin{definition}
\label{algebraic-geometry-definition-noetherian-ring}
A ring $R$ is {\it noetherian}
if any ascending chain of ideal stabilizes.
\end{definition}

\begin{theorem}[Hilbert's basis theorem]
\label{algebraic-geometry-theorem-hilberts-basis}
If $R$ is a Noetherian ring, so is $R[x]$.
Equivalently, every ideal in $R[x]$ is
finitely generated.
\end{theorem}

\begin{lemma}
\label{algebraic-geometry-lemma-noeth-iff-ideal-finit-gener}
A ring $R$ is Noetherian
if and only if
all its ideals are finitely generated.
\end{lemma}

\begin{proof}
Try!
\end{proof}

\begin{theorem}[Hilbert's basis theorem,
Misha's version]
\label{algebraic-geometry-theorem-hilberts-basis-misha}
Any finitely generated ring is Noetherian.
\end{theorem}

\noindent
The idea is to use this theorem
to show that any algebraic set may be
decomposed
as a finite union of irreducible varieties.
The idea of the proof is to reduce the
original algebraic set
to a union of two proper subvarieties, and
repeat the process until we arrive at
irreducible varieties.

\begin{definition}
\label{algebraic-geometry-definition-noeth-topol-space}
A topological space is {\it Noetherian}
if any descendending sequence of closed
subsets stabilizes.
\end{definition}

\begin{definition}
\label{algebraic-geometry-definition-zariski-topology}
{\it Zariski topology} is the topology
whose closed sets are algebraic sets.
\end{definition}

\begin{exercise}[Misha]
\label{algebraic-geometry-exercise-varie-zaris-compa}
Any variety is Zariski compact.
\end{exercise}

\noindent
To prove that any variety has a decomposition
as a finite union of irreducible components,
we first show pointwise that every point
is contained in a finite union of irreducible
components.
With existence at hand we just prove
in a similar way the existence of the
decomposition
for all the variety.

\medskip\noindent
I don't know where to put this yet:

\begin{definition}
\label{algebraic-geometry-definition-degree}
\begin{reference}
\href{https://en.wikipedia.org/wiki/%
Degree_of_an_algebraic_variety}{Wikipedia}.
\end{reference}
The {\it degree} of an affine or projective
variety $X$ of dimension $n$ is the number of
points
of intersection of $X$ and $n$ hyperplanes in
general position.
\end{definition}

\begin{exercise}
\label{algebraic-geometry-exercise-hypersurfaces-of-pn}
\begin{reference}
\cite[Exercise I.2.8]{hart}
\end{reference}
A projective variety $Y \subseteq
\mathbb{P}^n$
has dimension $n-1$ if and only if
it is the zero set of a single irreducible
homogeneous polynomial
$f$ of positive degree.
\end{exercise}

\begin{proof}
$Y$ has dimension $n-1$ if and only if there
is a chain
of irreducible subsets $Y \supset Y_2\supset…
\supset Y_{n-1}$
if and only if there is a chain of
prime ideals $I(Y) \subset I_2
\subset…\subset I_{n-1}$.
But $S:=k[x_0,…,x_n]$ has Krull dimension
$n+1$,
so that $I(Y)$ has height 1
(not considering the irrelevant ideal) and
thus principal
since $S$ is a UFD.
\end{proof}

\begin{proof}
($\implies$) Suppose $Y$ ir irreducible
and $\dim Y=n-1$. Pick any nonzero
homogeneous $g \in I(Y)$.
Factor $g$ into homogeneous irreducible
facts and choose one such facto $f$.
Then $f$ is homogeneous, irreducible and
vanishes on $Y$,
so $Y \subseteq V(f)$.
Both $Y$ and $V(f)$ are irreducible
projective
varieties and have the same dimension $n-1$,
hence $Y=V(I(f))$.

($\impliedby$) Suppose $Y=V(f)$ for some
nonconstant homogeneous polynomial $f$.
Then it has dimension $n-1$.
\end{proof}

\section{Regular functions, rational
functions and morphisms}
\label{algebraic-geometry-section-regular-functions}

\noindent
We start with the intuitive definition, yet
different from Harshorne's,
which works well for affine varieties.
A {\it regular} or {\it polynomial map}
is $\varphi:V \subset \mathbb{A}^n \to W
\subset \mathbb{A}^m$
if there exist $P_1,…,P_m \in k[x_1,…,x_n]$
such that
$\varphi(a_1,…,a_n)=(P_1(a_1,…,a_n),…,
P_m(a_1,…,a_n))$
and $\varphi(V) \subset W$.
If $\varphi$ is polynomial then there is an
induced
map in coordinate rings
$\varphi^*:k[W] \to k[V]$ given by pullback.
This map is well defined since
$\varphi^*(I(W))\subset I(V)$.

A polynomial map $\varphi:V \to W$ is an {\it
isomorphism}
if it has a left and right inverse that is
also a polynomial function.
$V \cong W \iff k[V] \cong k[W]$ for affine
varieties.

\medskip\noindent
Here is Hartshorne's definition:

\begin{definition}
\label{algebraic-geometry-definition-regular-function}
Let $Y$ be a quasi-affine variety
(i.e. an open set of an affine variety).
A function $Y \to k$ is called {\it regular}
at $P \in Y$
if there is an open set $U \ni P$ and
$g,h \in k[x_1,…,x_n]$ such that
$f=g/h$ on $U$. $f$ is {\it regular} if it is
regular
at all points of $Y$.
\end{definition}

\noindent
And the definition for projective varieties:

\begin{definition}
\label{algebraic-geometry-definition-regul-funct-proje}
Let $Y \subset \mathbb{P}^N$ be a
(quasi)projective variety.
A function $f:Y \to k$ is regular at $P \in
Y$
if there is an open set $U \ni P$ such that
$f=g/h$ for two homogeneous polynomials of
the same degree,
with $h$ nonzero in $U$.
\end{definition}

\noindent
Note that although homogeneous polynomials
are not functions on $\mathbb{P}^n$
(because we can rescale points, but the
functions will give different outputs),
the quotient of two homogeneous polynomials
is (because the rescaling factor
cancels out).

\begin{definition}
\label{algebraic-geometry-definition-morphism}
A {\it morphism} of varieties $\varphi:V \to
W$
is a function that pulls back polynomial
functions
$f:W \to k$ to polynomial functions
$\varphi^*f:V \to k$.
\end{definition}

\begin{exercise}
\label{algebraic-geometry-exercise-morphisms-def}
$\varphi$ is a morphism of affine varieties
if and only if
it is a polynomial function.
\end{exercise}

\begin{proof}
If $\varphi$ is polynomial then it is a
morphism:
if $f:W \to k$ is polynomial,
then
$\varphi^*f(a_1,…,a_n)=f(P_1(a_1,…,a_n),…,
P_m(a_1,…,a_n))$
is also polynomial.
Conversely, if $\varphi$ is a morphism
we see that its coordinate functions are
polynomials
by computing $y_i\varphi^*$ where $y_i:W \to
k$
is $i$-th coordinate function (which is
clearly a polynomial map).
\end{proof}

\medskip\noindent
Now we introduce the notion of local ring
(stalk).
Here's Olivier's (intuitive) definition:
Let $V$ be an affine variety (:= irreducible
algebraic set)
and $p \in V$.
$$
\mathcal{O}_p(V)
:= \{f \in k(V):\exists g,h \in k[V]
\text{ s.t. }f=g/h \text{ and } h(p) \neq
0\}\subset k(V).
$$

Compare with Hartshorne's:

\begin{definition}
\label{algebraic-geometry-definition-local-ring}
Let $Y$ be a variety (affine or projective)
and $P \in Y$.
The {\it local ring} of $P$ on $Y$ is
the ring of germs of regular functions on $Y$
near $P$.
\end{definition}

\begin{exercise}
\label{algebraic-geometry-exercise-local-ring-is-local}
The local ring is local with maximal ideal
given by functions vanishing at $P$.
\end{exercise}


\begin{exercise}
\label{algebraic-geometry-exercise-stalk-affine-varieties}
Show that if $V$ is an affine variety,
$\mathcal{O}_p(V) =
k[V]_{\mathfrak{m}_p}$
where $\mathfrak{m}_p$ is the set of regular
functions
that vanish in $p$.
\end{exercise}

\begin{proof}
Consider the obvious map: take a quotient of
classes
$[g]/[h]$ in the localization of the
coordinate ring of $V$
and map it to the rational function $g/h$.
This is well defined because if $g-g' \in
I(V)$,
then they obviously coincide in any open set
of $V$
because they coincide in all of $V$, hence
forming the same germ.
The key fact is that $\mathcal{O}_p(V)$ is
the germs of rational
functions of $V$.
It is injective because if a germ $g/h$ is
zero, then there's
an open of $V$ where $g$ vanishes, implying
that $g=0$
since it is a polynomial (or holomorphic)
map.
It is obviously surjective.
\end{proof}

\begin{definition}
\label{algebraic-geometry-definition-degree-zero-localization}
Let $S=\bigoplus_{d\geq 0} S_d$ be a
graded ring and let $f \in S$ be
homogeneous.
The localization $S_f$ is again graded by
\[
\deg(a/f^m)=\deg(a)-m\deg(f)
\]
for homogeneous $a \in S$.
We write $(S_f)_0$ for its degree-zero
part.
Thus $(S_f)_0$ consists of the fractions
$a/f^m$ such that
\[
\deg(a)=m\deg(f).
\]
More generally, if $\mathfrak{p}$ is a
homogeneous prime of $S$, we write
$S_{(\mathfrak{p})}$ for the degree-zero
part of the localization of $S$ at the
homogeneous elements outside
$\mathfrak{p}$.
\end{definition}

\begin{proposition}
\label{algebraic-geometry-proposition-standard-open-proj}
Let $Y=\operatorname{Proj} S$ and let
$f \in S$ be homogeneous.
Then the open set
\[
D_+(f)=\{\mathfrak{p}\in Y:f\notin
\mathfrak{p}\}
\]
is affine, with coordinate algebra
\[
\mathcal{O}_Y(D_+(f))=(S_f)_0.
\]
In other words, restricting to the chart
where $f$ does not vanish corresponds to
localizing at $f$ and taking degree zero.
\end{proposition}

\begin{theorem}
\label{algebraic-geometry-theorem-functions-projective-case}
\begin{reference}
\cite[Theorem 3.4]{har}
\end{reference}
Let $Y \subseteq \mathbb{P}^n$ be a
projective variety
with homogeneous coordinate ring $S(Y)$.
Then:
\begin{enumerate}
\item $\mathcal{O}(Y)=k$,
\item
$\mathcal{O}_P=S(Y)_{(\mathfrak{m}_P)}$,
where $\mathfrak{m}_P$ is the ideal generated
by the homogeneous
$f \in S(Y)$ not vanishing in $P$,
\item $K(Y)\cong S(Y)_{((0))}$.
\end{enumerate}
\end{theorem}

\noindent
Now let's discuss the equivalence of
categories between
affine varieties and finitely generated
reduced algebras.

\begin{proposition}
\label{algebraic-geometry-proposition-coordinate-algebra}
The coordinate algebra of an affine variety
$V$,
i.e. $k[x_1,…,x_n]/I(V)$ is isomorphic to
the ring of regular functions on $V$, i.e.
$\mathcal{O}(V)$.
\end{proposition}
The following proposition is also at
\href{https://stacks.math.columbia.edu/tag/%
01I1}{01I1}

\begin{proposition}
\label{algebraic-geometry-proposition-irreducible-prime}
Let $V,W \subset \mathbb{A}^n$ are
irreducible algebraic set.
There is a bijection
$$
\Hom_{k-alg}(k[W],k[V]))
\xrightarrow{bij}\underbrace{\Mor(V,W)}_{
\text{polynomial maps}}.
$$
\end{proposition}

\begin{proof}
We map
$$
\begin{aligned}
\tilde{\varphi}:
\underbrace{k[W]}_{=k[y_1,…,y_m]/I(W)}
&\longrightarrow
\underbrace{k[V]}_{=k[x_1,…,x_n]/I(V)} \\
 y_1&\longmapsto f_1\\
&\vdots \\
y_m & \longmapsto f_m
\end{aligned}
$$

\noindent
to
$$
\begin{aligned}
\varphi: V &\longrightarrow W \\
(x_1,…,x_n) &\longmapsto (f_1(x_1,…,x_n), …,
f_m(x_1,…,x_n)).
\end{aligned}
$$

\noindent
We check that this is well-defined, i.e.
that it doesn't depend on the choice of
representatives,
namely the $f_i$. But choosing other $f_i$
which differ
by an element of the ideal would not make a
difference
when evaluating at $V$ because the ideal is
precisely $I(V)$.

The inverse map we already described above.
\end{proof}

\noindent
This shows that we have an equivalence of
categories:

\begin{lemma}
\label{algebraic-geometry-lemma-algeb-varie-categ-equiv}
There is an equivalence of categories between
finitely generated reduced
$k$-algebras
and irreducible algebraic sets.
\end{lemma}


\medskip\noindent
Now let's discuss rational functions.
Again let's start with a more reasonable
explanation.
Since $k[V]$ for an irreducible affine
algebraic set is a domain,
there exists a fraction field, which is the
set of {\it rational functions}.
Functions here look like $f/g$ for $f,g \in
k[V]$,
which are not defined in $V(g)$.
Since this can lead to different ways
of defining the same function
--- the correct definition shall be that
we take the greatest set where the function
can be defined.

\begin{definition}
\label{algebraic-geometry-definition-rational-ring}
The {\it ring of rational
functions} is the set of regular functions
defined
on open sets with the identifications that
$(U,f)\sim(V,g)$
if they coincide in an intersection. (Like
the stalk,
but there's no point.)
\end{definition}

Compare this to Voisin's way of describing
a meromorphic function:
it's a function that can be written locally
as a quotient of two holomorphic functions.



\section{Dimension: Krull dimension,
transcendence degree}
\label{algebraic-geometry-section-dimension}

\begin{definition}
\label{algebraic-geometry-definition-irreducible-topological}
A nonempty subset $Y$ of a topological space
$X$
is {\it irreducible} if it cannot be
expressed
as a union of two proper subsets,
each of which is closed in $Y$.
\end{definition}

\begin{definition}
\label{algebraic-geometry-definition-topol-space-dimen}
If $X$ is a topological space,
the {\it dimension} of $X$ is the supremum of
all integers $n$
such that there exists a chain
$Z_0 \subset Z_1\subset … \subset Z_n$
of distinct irreducible closed subsets of
$X$.
The {\it dimension of an algebraic variety}
is its dimension as a topological space.
\end{definition}

\begin{definition}
\label{algebraic-geometry-definition-krull-dimension}
In a ring $A$, the {\it height} of a prime
ideal $\mathfrak{p}$
is the supremum of all integers $n$
such that there exists a chain
$\mathfrak{p}_0
\subset\mathfrak{p}
_1\subset…\subset\mathfrak{p}_n=\mathfrak{p}$
of distinct prime ideals.
The {\it Krull dimension} of $A$ is the
supremum of the heights
of all prime ideals.
\end{definition}

\begin{proposition}
\label{algebraic-geometry-proposition-dimension}
If $Y$ is an affine algebraic set,
the dimension of $Y$ is equal to the
dimension
of its coordinate ring.
\end{proposition}

\begin{theorem}
\label{algebraic-geometry-theorem-dimen-trans-degre}
Let $k$ be a field and let $B$ be an integral
domain
which is a finitely generated $k$-algebra.
Then:
\begin{enumerate}
\item the dimension of $B$ is equal to the
transcendence degree
of the quotient field $K(B)$ of $B$ over $k$;
\item For any prime ideal $\mathfrak{p}$ in
$B$, we have
$$
\text{height }\mathfrak{p}+ \dim B /
\mathfrak{p}=\dim B.
$$

\end{enumerate}
\end{theorem}

\begin{exercise}
\label{algebraic-geometry-exercise-affine-twisted-cubic}
Let $Y\subseteq \mathbb{A}^3$ be the set
$Y= \{(t,t^2,t^3:t \in k\}$.
Show that $Y$ is an affine variety of
dimension 1.
Find generators for the ideal $I(Y)$.
Show that $A(Y)$ is isomorphic to a
polynomial ring in one variable over $k$.
We say that $Y$ is given by the {\it
parametric representation}
$x=t$, $y=t^2$, $z=t^3$.
\end{exercise}

\begin{proof}
It's immediate that $x^2-y$ and $x^3-z$ must
be in
the ideal defining $Y$. Let
$J=(x^2-y,x^3-z)$.
Consider the ring map
$$
\begin{aligned}
\varphi:k[x,y,z] &\longrightarrow k[t] \\
x &\longmapsto t\\
y &\longmapsto t^2\\
z &\longmapsto t^3
\end{aligned}
$$

\noindent
Since $\varphi$ is given literally by the
parametrization of $Y$,
it follows from definition that the ideal of
$Y$
is the kernel of $\varphi$.

It's clear the in the kernel of $\varphi$
lie the polynomials $x^2-y$ and $x^3-z$.
But how to make sure that $J$ is the complete
kernel?

The procedure by AI is the following.
First notice that there is an isomorphism
$k[x,y,z]/J \cong k[x]$
given by $f(x,y,z)+J \mapsto f(x,x^2,x^3)$.
This shows that $k[x,y,z]/J$ is a domain, and
thus $J$ is prime.
Further, the projection $k[x,y,z] \to
k[x,y,z]/J$,
which obviously has kernel $J$,
is exactly the map I introduced before from
$k[x,y,z]$ to $k[t]$ when we see $k[x,y,z]/J$
as $k[t]$.
Since the kernel of such map is the ideal of
$Y$ by definition,
we are done.
\end{proof}

\section{Smooth points}
\label{algebraic-geometry-section-smooth-points}

\noindent
The notion of nonsingularity for affine
varieties is what we usually call
``the Jacobian criterion'',
i.e. it's given by the vanishing
of partial derivatives.
We shall see how to translate
that to commutative algebra,
which allows for a definition applicable
to non-affine varieties.

\begin{definition}
\label{algebraic-geometry-definition-nonsingular-affine}
Let $Y \subseteq \mathbb{A}^n$ be ann affine
variety,
and let $f_1,…,f_t \in R=k[x_1,…,x_n]$ be a
set of generators
for the ideal of $Y$.
$Y$ is {\it nonsingular at a point $P \in Y$}
if the rank of the matrix
$\left(\frac{\partial f_i}{\partial
x_j}(P)\right)$
is $n-r$, where $r$ is the dimension of $Y$.
$Y$ is {\it nonsingular} if it is nonsingular
at every point.
\end{definition}

\noindent
It can be checked that this definition is
independent of the choice of generators.

\begin{definition}
\label{algebraic-geometry-definition-residue-field}
Let $R$ be a ring and $\mathfrak{m}$ a
maximal ideal of $R$.
The {\it residue field} is  $R/\mathfrak{m}$.
\end{definition}

\begin{definition}
\label{algebraic-geometry-definition-regular-local-ring}
Let $R$ be a noetherian local ring with
maximal ideal $\mathfrak{m}$
and residue field $k:=R/\mathfrak{m}$.
$R$ is a {\it regular local ring} if
$\dim_k \mathfrak{m}/\mathfrak{m}^2=\dim R$.
\end{definition}

\noindent
Upshot. $\mathfrak{m}$ is the ideal of
functions vanishing at $P$.
This means that the constant term of $h \in
\mathfrak{m}$ is zero.
Differentiate $h$.
Now the constant term is what originally was
the linear part.
If $h'$ has no constant term, that is, if $h$
vanishes at $P$
up to order 1,
then $h$ can be seen as a product $h=fg$ for
$f,g \in \mathfrak{m}$.
Indeed, the smallest nonzero term of $h$ is
of degree 2,
so we can make it look like a product of two
things of degree 1.
Conversely, products $h=fg$ vanish up to
order 1 since
the first-order coefficient is given by the
derivative
(Taylor expansion)
and $dh=df\cdot g+fdg =0$ at $P$.

Then we can interpret
$\mathfrak{m}/\mathfrak{m}^2$ as follows.
Among all the functions which vanish at $P$,
i.e. inside $\mathfrak{m}$,
we identify any two functions if their
difference
is in $\mathfrak{m}^2$,
that is, if their difference vanishes up to
order 1 at $P$,
that is, if their linear parts coincide at
$P$.

Thus, $\mathfrak{m}/\mathfrak{m}^2$ is the
space
of linear parts of polynomials vanishing at
$P$.

\begin{theorem}[Jacobian criterion]
\label{algebraic-geometry-theorem-nonsi-jacob-iff-regul-local}
Let $Y \subseteq \mathbb{A}^n$ be an affine
variety.
Let $P \in Y$ be a point.
Then $Y$ is nonsingular at $P$
if and only if
the local ring $\mathcal{O}_{P,Y}$ is a
regular local ring.
\end{theorem}

\begin{proof}
Still hard to grasp completely,
but an important insight is that the map
$$
\begin{aligned}
\nabla: k[x_1,…,x_n]  &\longrightarrow k^n \\
 f &\longmapsto \nabla f=\left(\frac{\partial
 f}{\partial x_1},
…,\frac{\partial f}{\partial x_n}\right)
\end{aligned}
$$

\noindent
has kernel $\mathfrak{m}^2$ and is surjective
when restricted to $\mathfrak{m}$
(because $\nabla(x_i-p_i)=(0,…,0,1,0,…0)$, so
we construct a base of $k^n$)
so we have the isomorphism
$$
\mathfrak{m}/\mathfrak{m}^2=k^n,
$$
which we understand as the space of linear
directions at $P$.
\end{proof}

\begin{definition}
\label{algebraic-geometry-definition-nonsingular}
Let $Y$ be any variety,
$Y$ is {\it nonsingular} at a point $P \in Y$
if the local ring $\mathcal{O}_{P,Y}$ is a
regular local ring,
that is,
if $\dim_k\mathfrak{m}/\mathfrak{m}^2=\dim
\mathcal{O}_{P,Y}$.
\end{definition}

\noindent
Actually the Jacobian criterion
also works for projective hypersurfaces
(by simply looking at the what the
differential of a homogeneous polynomial is):

\begin{lemma}[Jacobian test for projective hypersurfaces]
\label{algebraic-geometry-lemma-projective-jacobian}
Let
$$
  X=V(F)\subset \mathbb P^n
$$
be a projective hypersurface, where
$F\in k[X_0,\ldots,X_n]$ is homogeneous.
Let
$$
  p=[a_0:\cdots:a_n]\in X.
$$
Then $p$ is singular on $X$ iff
$$
  \frac{\partial F}{\partial X_0}(p)
  =
  \cdots
  =
  \frac{\partial F}{\partial X_n}(p)
  =
  0.
$$
Equivalently, $p$ is smooth iff at least
one homogeneous partial derivative of
$F$ does not vanish at $p$.
\end{lemma}

\begin{proof}
Choose an affine chart containing $p$,
say $X_0\neq 0$. Put
$$
  x_i=\frac{X_i}{X_0}
  \qquad
  i=1,\ldots,n.
$$
On this chart the hypersurface is defined
by the dehomogenized equation
$$
  f(x_1,\ldots,x_n)
  =
  F(1,x_1,\ldots,x_n).
$$
By the affine Jacobian criterion, $p$ is
singular iff
$$
  \frac{\partial f}{\partial x_i}(p)=0
  \qquad
  i=1,\ldots,n.
$$
But these partials are exactly the
corresponding homogeneous partials
$$
  \frac{\partial F}{\partial X_i}(p),
  \qquad
  i=1,\ldots,n.
$$

It remains only to include the missing
partial with respect to $X_0$. Since $F$
is homogeneous, Euler's formula gives
$$
  \sum_{i=0}^n
  X_i\frac{\partial F}{\partial X_i}
  =
  dF,
$$
where $d=\deg F$. At points of $X$ we
have $F=0$. Hence
$$
  \sum_{i=0}^n
  X_i\frac{\partial F}{\partial X_i}
  =
  0.
$$
If $X_0\neq 0$ and all the other partials
vanish, then the $X_0$-partial vanishes
too. Thus the affine chart criterion is
equivalent to the vanishing of all
homogeneous partial derivatives.
\end{proof}


\section{Smooth points of curves}
\label{algebraic-geometry-section-smooth-points-curves}

\noindent
Singularity theory for plane curves
is particularly important for exam questions.
It can also be developed slightly differently
from the general construction so we put it
in this separate section.

\medskip\noindent
Upshot: the multiplicity of a point
is how many tangent lines there are at the
point. A point is smooth if and only if
it has multiplicity 1.

\medskip\noindent
The multiplicity of a plane curve $F$ at the
origin
is the degree of the lowest nonzero
homogeneous term $F_m$
of $F=F_m+F_{m+1}+…+F_n$.
The tangent lines at the origin are the
linear factors of $F_m$, which decomposes in
linear components
by the following homogenization trick:

\begin{exercise}
\label{algebraic-geometry-exercise-factor-in-linear-terms}
Let $f_m$ be a homogeneous polynomial of
degree $m$.
Show that $f_m$ is a product of $m$ linear
factors
$f_m=\prod (b_ix+c_iy)$.
\end{exercise}

\begin{proof}
Upshot: fix $y=1$, factor as product of
linear coefficients
using fundamental theorem of algebra, and
then rehomogenize.
The last is because a 1-variable polynomial
$F$ is homogeneous of degree $m$ if and only
if
$F(\lambda x)=\lambda^mF(x)$. So if you have
$F(x)=f(x,1)$
you rehomogenize by putting
$y^mF(x/y)=y^mf(x/y,1)$.
Which is strange because $y^mF(x/y)$ is in
$k[x,y,1/y]$,
but that doesn't matter because when you plug
it in you get
 $$
y^mf(x/y,1)=y^mc\prod(\frac{x}{y}-\alpha_i)=
c\prod(x-\alpha_iy),
$$
(where $c$ is a constant coming from making
the product having no
coefficients on $x$ after using fundamental
theorem of algebra)
which is already a product of linear
polynomials.
We conclude by showing that in fact
$y^mf(x/y,1)=f(x,y)$,
but that's just literally
$$
y^mf(x/y,1)=y^m\sum_{i}a_i \frac{x^i}{y^i}
=\sum_ia_i x^iy^{m-i}.
$$
\end{proof}

\noindent
Thus tangent lines can have multiplicity.
A {\it node} is an ordinary double point,
which means that there are two tangent lines
and they are not the same.
A {\it cusp} is a double point with a double
tangent line.

\begin{exercise}[Intersection multiplicity]
\label{algebraic-geometry-exercise-intersection-multiplicity}
If $Y,Z \subseteq \mathbb{A}^2$ are two
distinct curves,
given by equations $f$ and $g$, and if $P \in
Y \cap Z$,
we define the {\it intersection multiplicity}
$(Y \cdot Z)_P$
of $Y$ and $Z$ at $P$ to be the length of the
$\mathcal{O}_P$-module $\mathcal{O}_P/(f,g)$.
\begin{itemize}
\item Show that $(Y\cdot Z)_P$ is finite,
and $(Y\cdot Z)_P \geq  \mu_P(Y) \mu_P(Z)$.
\end{itemize}
\end{exercise}


\medskip\noindent
The upshot about Bézout theorem is that the
``size'' of the intersection
of two curves is the product of their
degrees.
In nice cases this is just the number of
intersection points.
In even nicer cases one of the curves is a
line,
in which case the product of degrees is
just the degree of the curve,
so that we compute the degree of a curve by
{\it the number of intersection points with a
hyperplane}.

But the notion of ``size'' of the
intersection
has to be made precise, which is done via the
notion of multiplicity.
The idea is that intersecting with a line
$\ell=ax +by+c$ where, say,  $b \neq 0$,
can be done by substituting $y$
for $-\alpha x-b$.
This makes the ``resultant''
$f_\ell(x):=f(x,\alpha x+\beta)$
be a polynomial in one variable,
which can be factored as a product of linear
terms, each of the roots of which
are the intersection points,
and are now equipped with a multiplicity.

\begin{definition}
\label{algebraic-geometry-definition-multiplicity-line}
The {\it multiplicity} or {\it intersection
index}
of $f\in k[x,y]\setminus k$ and
the line $\ell=ax+by+c$ (where $a \neq 0$ or
$b\neq 0$) is
$$
(\ell,f)_p =
\begin{cases}
0\qquad &\text{if }p \not\in V(f,\ell) \\
\infty \qquad &\text{if }p \in V(\ell) \text{
and }V(\ell) \subset V(f)\\
m_i \qquad & \text{if }p=(x_i,\alpha
x_i+\beta)
\end{cases}
$$
where $f_\ell(x):=f(x,\alpha x+\beta)$
factors as $c\prod(x-x_i)^{m_i}$.

If $V(\ell)\not \subset V(f)$, the integer
$$
\text{deg}(f)-\sum_{i=1}^rm_i
$$
is the {\it multiplicity of $\ell,f$ at
infinity}.
\end{definition}

\begin{exercise}
\label{algebraic-geometry-exercise-asymptotic-direction}
If $f=\sum_{i=0}^df_i \in k[x,y]$ for $f_i$
homogeneous of degree $i$,
an {\it asymptotic direction} of $f$ is
$ax+by$
which divides $f_d$.

Show that  $m_\infty \neq 0$ if and only if
$y-\alpha x$
is an asymptotic direction of $f$.
\end{exercise}

\begin{proof}
This uses a theorem using the resultant…
\end{proof}

\begin{theorem}[Bézout]
\label{algebraic-geometry-theorem-bezout}
Let $k$ be an algebraically closed field.
Let  $C,P \subset \mathbb{P}_k^2$ be two
curves without
irreducible components in common
and suppose that $C \cap D= \{p_1,…,p_r\}$
with intersection multiplicities  $m_1,…,m_r$
in $p_1,…,p_r$.
Then
$$
\text{deg}f\text{deg}g=\sum_{i=1}^r m_i.
$$
\end{theorem}


\begin{proposition}
\label{algebraic-geometry-proposition-multiplicity-point}
Let $f \in k[x,y]\setminus k$ and $p \in
V(f)$.
There exists an integer $m_p(f)\geq 1$ called
the {\it multiplicity of $p$ in $f$}
such that for every line $\ell$ passing
through  $p$
 $$
(\ell,f)_p \geq m
$$
and $(\ell,f)_p>m$ for at most $m$ lines and
at least one line.
\end{proposition}

\begin{definition}
\label{algebraic-geometry-definition-tangent-line}
Let $p=(x_0,y_0) \in V(f)$.
Write $f(x+x_0,y+y_0)=f_m(x,y)+\text{lower
degree terms} $,
i.e. translate $f$ so that the origin is in
its zero set
and write as homogeneous polynomials.
Then put
$f_m(x,y)=\prod_{i=1}^r(a_ix+b_iy)^{m_i}$.
The lines $\ell_i=a_ix+b_iy$ are called
{\it tangent lines} to $f$ and each has a
multiplicity $m_i$.
\end{definition}

\begin{definition}
\label{algebraic-geometry-definition-smoot-ordin-point-curve}
A point $p \in V(f)$ is {\it smooth} if
$m_p(f)=1$
and  {\it ordinary} if all its tangent lines
have multiplicity 1.
\end{definition}

\noindent
Thus, ordinary points are singular,
but these are mild singularities.
Essentially they are points where $m$
straight lines intersect transversally.

Definition of intersection index of two plane
curves
without common components.
Definition of intersection index of two plane
curves
with common components.
Properties of intersection index (7 of them).
These 7 properties characterize uniquely the
intersection index.

The multiplicity of a root $\alpha$ of a
power series
$f\in k[[x]]$,
i.e. the greatest integer such that
$(x-\alpha)^r|f$,
equals $\dim_k k[[x]]/(x^r)=r$.
(\href{https://www.youtube.com/live/%
EspNBiSU8F4?si=UEcdxrysAS47ZA9H&t=5230}{
Here}'s why.)
This leads to the alternative definition
of intersection index as
$\dim_k
k[[x,y]]/(f(x-x_0,y-y_0),g(x-x_0,y-y_0))$.
This satisfies the 7 properties of the
intersection index.


\noindent
\begin{proposition}[Jacobian criterion]
\label{algebraic-geometry-proposition-jacobian-criterion}
The singular locus of $V(f)$ is
$V(f,\partial_x,\partial_y)$.
\end{proposition}



\begin{theorem}[Plücker formulas]
\label{algebraic-geometry-theorem-plucker-formulas} Let $C
\subset \mathbb{P}^2$ be an irreducible curve
of degree $d$ whose only singularities are
$\delta$ nodes and $\chi$ ordinary cusps
(cusps such that the tangent line intersects
with multiplicity 3 at the singular point).
Let \begin{itemize} \item $i$ be the number
of inflection lines (lines $\ell$ such that
$(\ell,C)_p \geq 3$, and, in this case, do
not intersect another point with multiplicity
$\geq 2$), and

\item $d^*$ the number or lines tangent to
$C$ passing through a point $p \not\in C$
such that $p$ is not in any tangent lines at
the singular points, tangent inflection lines
or bitangent lines (lines that are tangent to
two different points). \end{itemize}

\noindent Then $$
\begin{aligned}
d(d-1)&=d^*+2\delta+3\chi\\
3d(d-2)&=i+6\delta+8\chi. \end{aligned}
$$
\end{theorem}


\section{Zariski tangent space}
\label{algebraic-geometry-section-zariski-tangent-space}

\noindent
Upshot: the Zariski tangent space is the set
of vectors killed simultaneously by the
differentials of all the generators of the
ideal.

\medskip\noindent
Let $X$ be an affine variety over $k$,
and let $p\in X$ be a $k$-point.
Let $\mathcal O_{X,p}$ be the local
ring at $p$, and let $\mathfrak m_p$
be its maximal ideal.

\begin{definition}
\label{algebraic-geometry-definition-zariski-cotangent}
The Zariski cotangent space of $X$
at $p$ is
$$
  T_p^*X:=\mathfrak m_p/\mathfrak m_p^2.
$$
\end{definition}

\begin{definition}
\label{algebraic-geometry-definition-zariski-tangent}
The Zariski tangent space of $X$
at $p$ is
$$
  T_pX
  :=
  \operatorname{Hom}_k
  (\mathfrak m_p/\mathfrak m_p^2,k).
$$
\end{definition}

The slogan is:
$T_pX$ is the space of first-order
directions at $p$ which preserve the
equations of $X$.

\begin{proposition}
\label{algebraic-geometry-proposition-jacobian-tangent}
Let
$$
  X=V(f_1,\ldots,f_r)
  \subset \mathbb A_k^n,
$$
and let $p=(p_1,\ldots,p_n)\in X$.
Then, under the identification
$T_p\mathbb A^n\simeq k^n$,
$$
  T_pX=\ker J(p),
$$
where
$$
  J(p)
  =
  \left(
    \frac{\partial f_j}{\partial x_i}(p)
  \right)_{j,i}.
$$
Equivalently, $v=(v_1,\ldots,v_n)$
belongs to $T_pX$ iff
$$
  \sum_{i=1}^n
  \frac{\partial f_j}{\partial x_i}(p)
  v_i
  =
  0
$$
for every $j=1,\ldots,r$.
\end{proposition}

\begin{proof}
A vector $v=(v_1,\ldots,v_n)\in k^n$
defines a derivation at $p$
$$
  D_v:k[x_1,\ldots,x_n]\to k
$$
by
$$
  D_v(g)
  =
  \sum_{i=1}^n
  \frac{\partial g}{\partial x_i}(p)
  v_i.
$$
This is the algebraic version of
evaluating $dg_p$ on $v$.

For $v$ to be tangent to $X$, the
equations defining $X$ must remain true
to first order. Hence $D_v$ must vanish
on the ideal $(f_1,\ldots,f_r)$.

It is enough to impose this on the
generators. Thus
$$
  D_v(f_j)=0
$$
for every $j$. But
$$
  D_v(f_j)
  =
  \sum_{i=1}^n
  \frac{\partial f_j}{\partial x_i}(p)
  v_i.
$$
Therefore $v\in T_pX$ exactly when
it lies in the kernel of the Jacobian
matrix evaluated at $p$.
\end{proof}

\begin{remark}
\label{algebraic-geometry-remark-residue-field}
The definition
$$
  T_pX
  =
  \operatorname{Hom}_k
  (\mathfrak m_p/\mathfrak m_p^2,k)
$$
assumes that $p$ is a $k$-point, i.e.
that its residue field is
$$
  \kappa(p)=k.
$$
In general, the correct definition is
$$
  T_pX
  =
  \operatorname{Hom}_{\kappa(p)}
  (\mathfrak m_p/\mathfrak m_p^2,\kappa(p)).
$$

Thus the only change is:
replace $k$ by the residue field
$\kappa(p)$.

Over an algebraically closed field $k$,
closed points are exactly $k$-points.
So, for classical varieties 
(zero-sets of ideals, affine or projective)
over $\mathbb C$, one may safely write
$$
  T_pX
  =
  \operatorname{Hom}_{\mathbb C}
  (\mathfrak m_p/\mathfrak m_p^2,\mathbb C).
$$

Scheme-theoretically there are also
non-closed points, such as generic
points. At such points, the residue
field need not be $k$.
\end{remark}



\section{Smoothness exercises}
\label{algebraic-geometry-section-smoothness-exercises}

\noindent
This is the result of a study session.
For now I'll leave it separate from the notes
in the previous section --- task: merge
the two sections.

\medskip\noindent
Let
$$
X=V(f_1,\dots,f_r)\subset \mathbb{A}^n.
$$
The Zariski tangent space at $p\in X$ is
$$
T_pX=\ker J(p),
$$
where
$$
J(p)=
\left(\frac{\partial f_i}{\partial x_j}(p)\right).
$$
Thus
$$
\dim T_pX=n-\operatorname{rank}J(p).
$$
The point $p$ is smooth if and only if
$$
\dim T_pX=\dim_pX.
$$
Equivalently,
$$
\operatorname{rank}J(p)=n-\dim_pX,
$$
the local codimension of $X$ in the ambient affine
space.

\medskip\noindent
For a hypersurface
$$
X=\{F=0\}\subset \mathbb{A}^n,
$$
a point $p\in X$ is smooth if and only if
$$
dF_p\neq 0.
$$
Equivalently, at least one partial derivative of $F$
is nonzero at $p$. A point is singular if
$$
F(p)=0,
\qquad
F_{x_1}(p)=\cdots=F_{x_n}(p)=0.
$$

\medskip\noindent
For a plane curve singularity at the origin, write
$$
f=f_m+f_{m+1}+\cdots+f_d,
$$
where each $f_i$ is homogeneous of degree $i$ and
$f_m\neq 0$. Then
$$
\operatorname{mult}_0(f)=m,
$$
and the tangent cone is
$$
TC_0V(f)=V(f_m).
$$
The multiplicity is the degree of the tangent cone,
counted with multiplicity. The number of distinct
tangent directions may be smaller.

\begin{exercise}
\label{algebraic-geometry-exercise-smoothness-node}
Classify the singularity of
$$
C=\{y^2=x^2+x^3\}\subset \mathbb{A}^2
$$
at the origin using multiplicity and tangent cone.
\end{exercise}

\begin{proof}
Write
$$
f=y^2-x^2-x^3.
$$
The lowest homogeneous part is
$$
f_2=y^2-x^2=(y-x)(y+x).
$$
Thus the multiplicity is $2$ and the tangent cone is
$$
TC_0C=\{y=x\}\cup \{y=-x\}.
$$
There are two distinct tangent directions, so the
origin is a node.
\end{proof}

\begin{exercise}
\label{algebraic-geometry-exercise-zariski-tangent-node}
For
$$
C=\{y^2=x^2+x^3\},
$$
compute the Zariski tangent space at the origin and
compare it with the tangent cone.
\end{exercise}

\begin{proof}
The local ring is
$$
\mathcal{O}_{C,0}
=
\left(k[x,y]/(y^2-x^2-x^3)\right)_{(x,y)}.
$$
Let $\mathfrak m=(x,y)$. Since the defining equation
has no linear part, it imposes no linear relation in
$\mathfrak m/\mathfrak m^2$. Therefore
$$
\mathfrak m/\mathfrak m^2
=\langle \bar x,\bar y\rangle_k,
$$
so
$$
\dim_k\mathfrak m/\mathfrak m^2=2.
$$
Equivalently, the Jacobian vanishes at the origin, so
$$
T_0C=\mathbb{A}^2.
$$
But the tangent cone is
$$
TC_0C=\{y^2-x^2=0\},
$$
a union of two lines inside $\mathbb{A}^2$.
The tangent space sees only linear equations, while the
tangent cone sees the first nonzero homogeneous part.
\end{proof}

\begin{exercise}
\label{algebraic-geometry-exercise-cusp-local-ring-not-regular}
Let
$$
C=\{y^2=x^3\}\subset \mathbb{A}^2.
$$
Show that $\mathcal{O}_{C,0}$ is not regular.
Compare this with the normalization
$$
\mathbb{A}^1\to C,
\qquad
 t\mapsto (t^2,t^3).
$$
\end{exercise}

\begin{proof}
We have
$$
\mathcal{O}_{C,0}
=
\left(k[x,y]/(y^2-x^3)\right)_{(x,y)}.
$$
The maximal ideal is
$$
\mathfrak m=(x,y).
$$
The equation $y^2-x^3$ has no linear part, so it gives
no relation in $\mathfrak m/\mathfrak m^2$. Therefore
$$
\mathfrak m/\mathfrak m^2
=\langle \bar x,\bar y\rangle_k,
\qquad
\dim_k\mathfrak m/\mathfrak m^2=2.
$$
But $C$ is a curve, so
$$
\dim \mathcal{O}_{C,0}=1.
$$
Hence
$$
\dim_k\mathfrak m/\mathfrak m^2
\neq
\dim \mathcal{O}_{C,0},
$$
and the local ring is not regular.

The parametrization induces
$$
k[x,y]/(y^2-x^3)\hookrightarrow k[t],
\qquad
x\mapsto t^2,
\quad
 y\mapsto t^3.
$$
Thus the coordinate ring is $k[t^2,t^3]\subset k[t]$.
The element $t$ is integral over $k[t^2,t^3]$, since
$$
t^2-x=0,
$$
but $t\notin k[t^2,t^3]$. Thus the cusp is not normal,
and its normalization is $\mathbb{A}^1$.
\end{proof}

\begin{exercise}
\label{algebraic-geometry-exercise-node-xy-normalization}
Let
$$
C=\{xy=0\}\subset \mathbb{A}^2.
$$
Compute $\dim_k\mathfrak m/\mathfrak m^2$ at the
origin and describe the normalization geometrically.
\end{exercise}

\begin{proof}
The local ring at the origin is
$$
\left(k[x,y]/(xy)\right)_{(x,y)}.
$$
The maximal ideal is $\mathfrak m=(x,y)$. Since $xy$
has no linear part, it gives no relation in
$\mathfrak m/\mathfrak m^2$. Hence
$$
\mathfrak m/\mathfrak m^2
=\langle \bar x,\bar y\rangle_k,
\qquad
\dim_k\mathfrak m/\mathfrak m^2=2.
$$
The curve has dimension $1$, so the origin is not
regular.

Geometrically, $C$ is the union of the two coordinate
axes. Its normalization separates the two branches:
$$
\widetilde C=\mathbb{A}^1\sqcup \mathbb{A}^1.
$$
On coordinate rings, the normalization map is
$$
k[x,y]/(xy)\longrightarrow k[x]\times k[y],
$$
where
$$
x\longmapsto (x,0),
\qquad
 y\longmapsto (0,y).
$$
The origin becomes two points, one on each component.
\end{proof}

\begin{exercise}
\label{algebraic-geometry-exercise-projective-jacobian-criterion-plane-curve}
Let
$$
Y=V(F)\subset \mathbb{P}^2
$$
be a plane curve, with $F$ homogeneous of degree $d$.
Show that
$$
Y\text{ is singular at }[p]
\iff
F_X(p)=F_Y(p)=F_Z(p)=0.
$$
\end{exercise}

\begin{proof}
The condition is projectively well-defined because
$F_X,F_Y,F_Z$ are homogeneous of degree $d-1$.
Thus their vanishing does not depend on the choice of
homogeneous representative.

Work in the chart $Z\neq 0$ and write
$$
f(x,y)=F(x,y,1).
$$
Then $Y$ is singular at $[a:b:1]$ if and only if
$$
f_x(a,b)=f_y(a,b)=0.
$$
But
$$
f_x=F_X(x,y,1),
\qquad
f_y=F_Y(x,y,1).
$$
Thus singularity in the affine chart gives
$$
F_X(a,b,1)=F_Y(a,b,1)=0.
$$
Since $F$ is homogeneous, Euler's formula gives
$$
XF_X+YF_Y+ZF_Z=dF.
$$
At a point of $Y$ we have $F=0$. Since also
$F_X=F_Y=0$ and $Z=1$, we get
$$
F_Z=0.
$$
The converse is immediate from the affine Jacobian
criterion in the chart.
\end{proof}

\begin{exercise}
\label{algebraic-geometry-exercise-cone-over-smooth-plane-curve-singularities}
Let
$$
Y=\{F(X,Y,Z)=0\}\subset \mathbb{P}^2
$$
be a smooth plane curve of degree $d>1$. Let
$$
C(Y)=\{F(x,y,z)=0\}\subset \mathbb{A}^3
$$
be the affine cone. Show that
$$
\operatorname{Sing}C(Y)=\{0\}.
$$
\end{exercise}

\begin{proof}
The singular locus of the cone is
$$
\{F=F_x=F_y=F_z=0\}\subset \mathbb{A}^3.
$$
If $q=(a,b,c)\neq 0$ is singular, then $[q]$ is a
point of $Y$. In an affine chart, the vanishing of all
partials of $F$ implies the vanishing of the chart
partials. Thus $[q]$ would be a singular point of
$Y$, contradiction.

At the origin, $F(0)=0$ because $F$ is homogeneous of
positive degree. Since $d>1$, each partial derivative
has degree $d-1\geq 1$, so all partials vanish at the
origin. Thus the origin is singular.
\end{proof}



\section{Blow-up}
\label{algebraic-geometry-section-blow-up}

\noindent
The blow-up of the origin in $\mathbb{A}^2$ is
$$
\operatorname{Bl}_0\mathbb{A}^2
=
\{((x,y),[u:v])\in \mathbb{A}^2\times
\mathbb{P}^1:xv=yu\}.
$$
The projection
$$
\pi:\operatorname{Bl}_0\mathbb{A}^2\to
\mathbb{A}^2
$$
is the first projection. The exceptional divisor is
$$
E=\pi^{-1}(0)\cong \mathbb{P}^1.
$$

\medskip\noindent
In the chart $u\neq 0$, write $[u:v]=[1:v]$.
Then
$$
y=xv.
$$
In the chart $v\neq 0$, write $[u:v]=[u:1]$.
Then
$$
x=uy.
$$
Thus the blow-up replaces the origin by the set of
possible tangent directions.

\medskip\noindent
Let $C=\{f(x,y)=0\}\subset \mathbb{A}^2$ be a
plane curve passing through the origin. The total
transform is
$$
\pi^{-1}(C).
$$
The strict transform is
$$
\widetilde C=
\overline{\pi^{-1}(C\setminus\{0\})}.
$$
In the chart $y=xv$, if
$$
f(x,xv)=x^m g(x,v),
$$
where $x$ does not divide $g$, then
$$
\widetilde C=\{g=0\}.
$$
Indeed, away from the exceptional divisor $x=0$, the
factor $x^m$ is invertible, so the equation is just
$g=0$. Taking the Zariski closure keeps the equation
$g=0$ and allows $x=0$.

\medskip\noindent
For a plane curve the divisor formula is
$$
\pi^*C=\widetilde C+mE,
$$
where $m$ is the multiplicity of $C$ at the origin.
Equivalently, $m$ is the multiplicity with which
$E$ appears in the total transform.

\begin{exercise}
\label{algebraic-geometry-exercise-blow-up-parabola}
Blow up the origin and compute the strict transform
of
$$
C=\{y=x^2\}\subset \mathbb{A}^2.
$$
\end{exercise}

\begin{proof}
Use the chart $y=xv$. Then
$$
y-x^2=xv-x^2=x(v-x).
$$
Thus
$$
\pi^*C=E+\widetilde C,
$$
where
$$
E=\{x=0\},
\qquad
\widetilde C=\{v=x\}.
$$
The strict transform meets $E$ at
$$
x=0,
\qquad
v=0.
$$
So it meets $E$ in one point. The strict transform is
smooth, since it is the line $v=x$ in the chart.
The original curve was already smooth; the blow-up
only records its tangent direction.
\end{proof}

\begin{exercise}
\label{algebraic-geometry-exercise-cuspidal-curve-blow-up}
Let $Y$ be the cuspidal curve $y^2=x^3$ in
$\mathbb{A}^2$.
Blow up the point $O=(0,0)$, let $E$ be the
exceptional curve,
and let $\widetilde Y$ be the strict transform
of $Y$.
Show that $E$ meets $\widetilde Y$ in one point,
and that $\widetilde Y\cong \mathbb{A}^1$.
\end{exercise}

\begin{proof}
Use the chart $y=xv$. Then
$$
y^2-x^3=x^2v^2-x^3=x^2(v^2-x).
$$
Therefore
$$
\pi^*Y=2E+\widetilde Y,
\qquad
\widetilde Y=\{v^2=x\}.
$$
The exceptional divisor is $E=\{x=0\}$, so
$$
E\cap \widetilde Y=\{x=0,v^2=0\}.
$$
This is one point, with multiplicity $2$.
The strict transform is smooth because
$$
F=v^2-x
$$
has
$$
dF=2v\,dv-dx,
$$
which is nowhere zero on $F=0$.
Moreover, the equation $x=v^2$ parametrizes
$\widetilde Y$ by $v$, so
$$
\widetilde Y\cong \mathbb{A}^1.
$$
\end{proof}

\begin{exercise}
\label{algebraic-geometry-exercise-blow-up-tacnode}
Blow up the origin in
$$
C=\{y^2=x^4\}\subset \mathbb{A}^2.
$$
Show that the strict transform has a node.
\end{exercise}

\begin{proof}
Again use $y=xv$. Then
$$
y^2-x^4=x^2(v^2-x^2).
$$
Thus
$$
\pi^*C=2E+\widetilde C,
\qquad
\widetilde C=\{v^2=x^2\}.
$$
The strict transform meets $E=\{x=0\}$ at
$$
v^2=0,
$$
so at one point. But
$$
v^2-x^2=(v-x)(v+x).
$$
Thus the strict transform has two distinct tangent
lines at this point. Hence it has a node.
\end{proof}

\begin{exercise}
\label{algebraic-geometry-exercise-blow-up-node}
Blow up the origin in
$$
C=\{y^2=x^2+x^3\}\subset \mathbb{A}^2.
$$
Show that the strict transform meets the exceptional
curve in two distinct smooth points.
\end{exercise}

\begin{proof}
In the chart $y=xv$,
$$
y^2-x^2-x^3=x^2(v^2-1-x).
$$
Thus
$$
\pi^*C=2E+\widetilde C,
\qquad
\widetilde C=\{v^2-1-x=0\}.
$$
Intersecting with $E=\{x=0\}$ gives
$$
v^2-1=0.
$$
Thus the strict transform meets $E$ at the two points
$$
v=1,
\qquad
v=-1.
$$
The strict transform is smooth because
$$
F=v^2-1-x
$$
satisfies
$$
dF=2v\,dv-dx,
$$
which is nonzero at both points.
\end{proof}

\begin{exercise}
\label{algebraic-geometry-exercise-blow-up-y3-x5}
Blow up the origin in
$$
C=\{y^3=x^5\}\subset \mathbb{A}^2.
$$
Show that the exceptional divisor appears with
multiplicity $3$, and that the strict transform has a
cusp-type double point.
\end{exercise}

\begin{proof}
Use $y=xv$. Then
$$
y^3-x^5=x^3(v^3-x^2).
$$
Thus
$$
\pi^*C=3E+\widetilde C,
\qquad
\widetilde C=\{v^3=x^2\}.
$$
The strict transform meets $E$ at $(x,v)=(0,0)$.
At this point the lowest homogeneous part of
$$
F=v^3-x^2
$$
is
$$
F_2=-x^2.
$$
So the strict transform has multiplicity $2$ and one
repeated tangent line. This is cusp-type, not a node.
\end{proof}

\medskip\noindent
The blow-up of the origin in $\mathbb{A}^3$ is
$$
\operatorname{Bl}_0\mathbb{A}^3
=
\{((x,y,z),[u:v:w])\in \mathbb{A}^3\times
\mathbb{P}^2:\operatorname{rank}
\begin{pmatrix}
x&y&z\\
u&v&w
\end{pmatrix}
\leq 1\}.
$$
Equivalently, the $2\times 2$ minors vanish:
$$
xv-yu=0,
\qquad
xw-zu=0,
\qquad
yw-zv=0.
$$
The exceptional divisor is
$$
E\cong \mathbb{P}^2.
$$

\begin{exercise}
\label{algebraic-geometry-exercise-blow-up-ordinary-double-surface}
Blow up the origin in the surface
$$
X=\{z^2=xy\}\subset \mathbb{A}^3.
$$
Compute the strict transform in the chart
$y=xu$, $z=xv$.
\end{exercise}

\begin{proof}
Substitution gives
$$
z^2-xy=x^2v^2-x^2u=x^2(v^2-u).
$$
Therefore
$$
\pi^*X=2E+\widetilde X,
\qquad
\widetilde X=\{v^2=u\}.
$$
This chart of the strict transform is smooth, since
$$
F=v^2-u
$$
has
$$
dF=2v\,dv-du,
$$
which is nowhere zero.
\end{proof}

\begin{exercise}
\label{algebraic-geometry-exercise-blow-up-quadratic-cone}
Blow up the origin in
$$
X=\{x^2+y^2+z^2=0\}\subset \mathbb{A}^3.
$$
Compute the strict transform in the chart
$y=xv$, $z=xw$.
\end{exercise}

\begin{proof}
We get
$$
x^2+y^2+z^2=x^2(1+v^2+w^2).
$$
Thus
$$
\pi^*X=2E+\widetilde X,
\qquad
\widetilde X=\{1+v^2+w^2=0\}.
$$
The intersection with $E=\{x=0\}$ is
$$
\{x=0,1+v^2+w^2=0\}.
$$
This is part of the projectivized tangent cone inside
$E\cong \mathbb{P}^2$.
The strict transform is smooth in this chart: the only
possible critical point of $1+v^2+w^2$ would be
$v=w=0$, but this point is not on the strict transform.
\end{proof}

\begin{exercise}
\label{algebraic-geometry-exercise-cone-over-smooth-curve-blow-up}
Let
$$
Y=\{F(X,Y,Z)=0\}\subset \mathbb{P}^2
$$
be a smooth plane curve of degree $d>1$.
Let
$$
C(Y)=\{F(x,y,z)=0\}\subset \mathbb{A}^3
$$
be the affine cone over $Y$.
Show that $C(Y)$ is singular only at the vertex, and
that the blow-up of the vertex resolves this
singularity. Show that the exceptional divisor of the
strict transform is naturally isomorphic to $Y$.
\end{exercise}

\begin{proof}
The singular locus of the cone is cut out by
$$
F=F_x=F_y=F_z=0.
$$
If $q=(a,b,c)\neq 0$ is a singular point of
$C(Y)$, then $[q]=[a:b:c]$ is a point of $Y$.
On an affine chart, say $Z\neq 0$, the equation is
$f(x,y)=F(x,y,1)$, and
$$
f_x=F_X(x,y,1),
\qquad
f_y=F_Y(x,y,1).
$$
Thus the vanishing of the cone partial derivatives
would make $[q]$ a singular point of $Y$, impossible.
Therefore the only possible singular point is the
origin.

Since $F$ is homogeneous of degree $d>1$, all partial
 derivatives are homogeneous of degree $d-1\geq 1$.
Hence they vanish at the origin. Thus the vertex is
singular.

Now use the description
$$
\operatorname{Bl}_0\mathbb{A}^3
=
\{(p,[\ell])\in \mathbb{A}^3\times \mathbb{P}^2:
 p\in \ell\}.
$$
If $p=t\ell$, then by homogeneity
$$
F(p)=t^dF(\ell).
$$
The total transform has the exceptional factor $t^d$.
Removing it gives the strict transform
$$
\widetilde{C(Y)}
=
\{(p,[\ell]):p\in \ell, F(\ell)=0\}.
$$
Over the exceptional divisor $p=0$ this becomes
$$
\{0\}\times Y.
$$
So the exceptional divisor of the strict transform is
naturally isomorphic to $Y$.

Moreover, the strict transform is the total space of
$$
\mathcal{O}_{\mathbb{P}^2}(-1)|_Y.
$$
Indeed, over $[\ell]\in Y$, the fiber is the line
$\ell$. Since $Y$ is smooth and a vector bundle over a
smooth variety is smooth, the strict transform is
smooth.
\end{proof}

\medskip\noindent
Now consider the blow-up of a point in the projective
plane. Let
$$
\pi:X=\operatorname{Bl}_p\mathbb{P}^2\to
\mathbb{P}^2.
$$
For example, take $p=[1:0:0]$. Then
$$
X=
\{([X:Y:Z],[U:V])\in \mathbb{P}^2\times
\mathbb{P}^1:YV=ZU\}.
$$
The exceptional divisor is
$$
E=\pi^{-1}(p)\cong \mathbb{P}^1.
$$
It parametrizes tangent directions at $p$.

\medskip\noindent
Let $H$ be the pullback class of a line in
$\mathbb{P}^2$:
$$
H=\pi^*[\text{line}].
$$
If $L$ is a line not passing through $p$, then its
strict transform has class $H$ and is disjoint from
$E$.
Therefore
$$
H^2=1,
\qquad
H\cdot E=0.
$$
If $L$ is a line passing through $p$, then locally it
has multiplicity $1$ at $p$, so
$$
\pi^*L=\widetilde L+E.
$$
Since $\pi^*L$ has class $H$, we get
$$
\widetilde L=H-E.
$$
Take two distinct lines $L,L'$ through $p$. Downstairs
they meet only at $p$. Upstairs their strict transforms
meet the exceptional curve at different points,
because they have different tangent directions.
Thus
$$
\widetilde L\cdot \widetilde L'=0.
$$
Since both have class $H-E$,
$$
0=(H-E)^2=H^2-2H\cdot E+E^2.
$$
Using $H^2=1$ and $H\cdot E=0$, we obtain
$$
E^2=-1.
$$
Thus
$$
H^2=1,
\qquad
H\cdot E=0,
\qquad
E^2=-1.
$$

\medskip\noindent
A final sanity check is the blow-up of a smooth
divisor. Blow up $\mathbb{A}^2$ along
$$
D=\{x=0\}.
$$
The ideal is principal, $I=(x)$. Blowing up a
Cartier divisor is trivial. Indeed, the Rees algebra is
$$
\bigoplus_{n\geq 0}I^nt^n
=
k[x,y]\oplus (x)t\oplus (x^2)t^2\oplus\cdots
=
k[x,y][xt].
$$
There is only one normal direction, so the
projectivized normal direction is $\mathbb{P}^0$ over
each point. Hence no new exceptional geometry appears.

\section{Rational maps}
\label{algebraic-geometry-section-rational-maps}

\noindent
The category of varieties and dominant
rational maps is equivalent to the category
of finitely generated field extensions of $k$
with the arrows reversed.

\section{Normalization}
\label{algebraic-geometry-section-normalization}

\noindent
``All bounded meromorphic functions are
holomorphic''. For exam curves:
find a parametrization of the rational curve.

For the abstract definitions see
Commutative Algebra Section
\ref{commutative-algebra-section-normalization}.
Here we explain how to
parametrize rational curves
and the compute normalizations of some
varieties.

A curve $C$ is rational if its
normalization is $\mathbb P^1$.
A parametrization of $C$ is then the
normalization map
$$
  \nu:\mathbb P^1\to C.
$$

In practice, for plane curves, one often
finds this map using a pencil of lines.

Suppose
$$
  C=V(f)\subset \mathbb A^2
$$
has a singular point at $p=(0,0)$.
Consider the lines through $p$:
$$
  y=tx.
$$
Substitute this into the equation:
$$
  f(x,tx)=0.
$$

Since $p$ is singular, the factor $x=0$
appears with multiplicity at least $2$.
(See Definition
\ref{algebraic-geometry-definition-smoot-ordin-point-curve},
which would match the regular local ring
definition of smoothnes via
the Jacobian criterion,
Proposition
\ref{algebraic-geometry-proposition-jacobian-criterion}.)
Remove this factor. The remaining
equation gives the other intersection
points of the line with $C$.

If this remaining equation can be solved
rationally for $x$ in terms of $t$, then
one gets
$$
  x=x(t),
  \qquad
  y=tx(t).
$$
This gives a parametrization of $C$.

For a singular plane cubic, this method
usually works immediately: a line through
the singular point meets the cubic with
multiplicity $2$ at the singular point,
so there is only one remaining point.

Example: for
$$
  C: y^2=x^2+x^3,
$$
put $y=tx$. Then
$$
  t^2x^2=x^2+x^3.
$$
Thus
$$
  x^2(t^2-1-x)=0.
$$
Ignoring the factor $x^2$, which is the
singular point, gives
$$
  x=t^2-1.
$$
Hence
$$
  y=t(t^2-1).
$$
So a parametrization is
$$
  t\mapsto (t^2-1,t(t^2-1)).
$$

Warning: guessing monomials
$x=t^a$, $y=t^b$ is not a general
method. It works for some cusps, but not
for nodes.

\medskip\noindent
Now I include some
Hartshorne/half-theoretical
exercises I did before.

Recall from commutative algebra
that a (affine or projective) 
variety is called {\it normal at $P$}
is $\mathcal{O}_P$ is integrally closed.
It is called {\it normal} if every point
is normal.

\begin{exercise}
\label{algebraic-geometry-exercise-normal-varieties}
If $Y$ is an affine variety,
$\mathcal{O}_P$ is integrally closed
for all $P \in Y$ if and only if
$\mathcal{O}_Y$ is integrally closed.
\end{exercise}

\begin{proof}
For the converse
let $f/g \in \text{Frac}\mathcal{O}_P$
be integral.
Then there exists $P \in \mathcal{O}_P[t]$
monic killing $f/g$.
We want to show that $f/g \in \mathcal{O}_P$
for which it's enough to put $f/g$ in
$\mathcal{O}_Y$ using the closure-integrality
of $Y$.
Suppose
$$
P=\sum \frac{f_i}{g_i}t^i.
$$
Since this polynomial is not in
$\mathcal{O}_Y[t]$, we cannot use it directly
to prove integrality of $f/g$;
in fact $f/g$ might not be integral at all.
So we get rid of the fractions
considering the least common multiple
of the $g_i$.
But if we just multiply by the lcm
we get a polynomial that is not monic.
We must consider
$$
P\cdot lcm^{\text{deg}P}
\left(\frac{f}{g}\right)
=\sum lcm^{\text{deg}P}\frac{f_i}{g_i}
\left(\frac{f}{g}\right)^i
=\frac{f_i}{g_i}
\left(lcm \frac{f}{g}\right)^{\text{deg}P}
+\frac{f_i}{g_i}lcm
\left(lcm\frac{f}{g}\right)^{\text{deg}P-1}
+…
+\frac{f_0}{g_0}lcm^{\text{deg}P},
$$
which is monic over $\mathcal{O}_Y$
and kills $lcm \frac{f}{g}$,
which implies that $\frac{f}{g}
\in \mathcal{O}_Y$
and thus $g$ is never zero
(or simply a unit of the coordinate algebra),
and thus $f/g$ is in $\mathcal{O}_P$.

The direct implication is actually easier.
Let $f/g \in \text{Frac}\mathcal{O}_Y$
be integral, i.e. there is a polynomial
$P \in \mathcal{O}_Y[t]$ killing $f/g$.
Pick any point $p \in Y$.
Then $P$ can be taken to be a polynomial
in $\mathcal{O}_p[t]$, which in fact
kills the fraction
$\overline{f/g} \in \text{Frac}\mathcal{O}_p$
in the local ring of $p$.
Since $\mathcal{O}_p$ is integrally closed,
$f/g \in \mathcal{O}_P$.
Thus, thinking of $f/g$ as a meromorphic
function, we see that $g(p) \neq  0$ for all $p$,
meaning the fraction $f/g$ is in fact
regular. Alternatively, thinking of $f/g$
as an element of the localization by the
maximal ideal $\mathfrak{m}_P$,
saying that $f/g \in \mathcal{O}_P$
means that $g \not\in \mathfrak{m}_P$
for all $P$, which is to say that $g$
is in no maximal ideal of $k[Y]$,
which means that $g$ is a unit.
\end{proof}

\begin{exercise}
\label{algebraic-geometry-exercise-normal-varieties-examples}
\begin{enumerate}
\item Show that every conic in $\mathbb{P}^2$
is normal.
\item Show that the quadric surfaces
$Q_1,Q_2$ in $\mathbb{P}^3$ given by
$Q_1:xy-zw$ and $Q_2:xy=z^2$ are normal.
\item Show that the cuspidal cubic
$y^2=x^3$ in $\mathbb{A}^2$
is not normal.
\end{enumerate}
\end{exercise}

\begin{proof}
\begin{enumerate}
\item We need to show that every conic
in $\mathbb{P}^2$ is rational.
\item
\item The element $x^2/y$ has a minimal
polynomial in the local ring at the origin,
namely $P(t)=t-x$, but it clearly
does not lie in the local ring since $y$
vanishes at $(0,0)$.
\end{enumerate}
\end{proof}

\begin{exercise}
\label{algebraic-geometry-exercise-universal-prop-normalization}
Let $Y$ be an affine variety.
Show that there is a normal affine
variety $\tilde{Y}$, and a morphism
$\pi:\tilde{Y}\to Y$,
with the property that whenever $Z$
is a (affine) normal variety and 
$\varphi:Z \to Y$ is a dominant morphism,
then there is a unique morphism
$\theta:Z \to \tilde{Y}$ such that
$\varphi=\pi \circ \theta$.
\end{exercise}

\begin{proof}
Let $\tilde{Y}=\Spec \tilde{A}$
where $\tilde{A}$ is the integral closure
of $A=\mathcal{O}_Y$.
Then by Exercise
\ref{algebraic-geometry-exercise-normal-varieties}
it is normal.
Now if the dominant map $\varphi:Z \to Y$
induces an injective algebra morphism
$\varphi^\#:A \to B$
and $B$ is integrally closed,
we can put $\theta^\#:\tilde{A} \to B$
by the inclusion.
Since any element $x \in \tilde{A}$
is the root of a unique monic polynomial
$P \in A[t] \subset B[t]$,
and since $\theta^\#(x)$ must be
the root of a monic polynomial in $B[t]$
because $B$ is integrally closed,
we conclude that $\theta^\#(x)$
is also a root of $P$.
… but what if the roots are different?

A correct solution is as follows.
The inclusion $A \xymatrix{\ar@{^{(}->}[r]&}
B$ induces an inclusion on fraction fields
$\text{Frac}A \xymatrix{\ar@{^{(}->}[r]&}
\text{Frac} B$.
The key observation is that by definition
an element in the integral closure 
$\tilde{A}$ of $A$ is an element in
the fraction field $\text{Frac} A$.
This says that any integral element
can already be seen as an element in
$\text{Frac}B$.
But since $B$ is integrally closed and the
minimal polynomial of any such element
has coefficients in $A \subset B$,
we are done: an integral element of $A$
is mapped to an element of $B$ making
the following diagram commute
$$
\xymatrix{
A\ar[dr]_{\pi^\sharp}
\ar@{^{(}->}[r]^{\varphi^\sharp}
&  B\\
&  \tilde{A}\ar@{-->}[u]^{\theta^\sharp}
}
$$
For uniqueness, suppose that
$\theta':\tilde{A} \to B$ is another
map making the above diagram commute.
Then for any $x=a/b$ integral over $A$ 
we have
$\theta'(b)\theta'(x)=\theta'(a)$.
But this coincides on $A$
with $\theta^\sharp$
we introduced before, giving
$\theta^\sharp(b)(\theta'(x)-\theta^\sharp)
=0$. Since there are no zerodivisors on $B$,
we obtain $\theta\equiv\theta^\sharp$.
\end{proof}

\begin{exercise}
\label{algebraic-geometry-exercise-projectively-normal-vars}
A projective variety $Y \subseteq
\mathbb{P}^n$ is {\it projectively normal} 
(with respect to the given embedding)
if its homogeneous coordinate ring
$S(Y)$ is integrally closed.
\begin{enumerate}
\item If $Y$ is projectively normal
then $Y$ is normal.
\item There are normal varieties 
(i.e. stalks are integrally closed)
in projective space which are not 
projectively normal.
For example, let $Y$ be the twisted
quartic curve in $\mathbb{P}^3$
given parametrically by
$(x,y,z,w)=(t^4,t^3u,tu^3,u^4)$.
\item Show that the twisted quartic curve
$Y$ above is isomorphic to $\mathbb{P}^1$,
which is projectively normal. Thus projective
normality depends on the embedding.
\end{enumerate}
\end{exercise}

\begin{proof}
\begin{enumerate}
\item Take an affine chart $D_+(f)$.
By Proposition
\ref{algebraic-geometry-proposition-standard-open-proj}
the coordinate algebra of $Y$
at this open set is given by
$(S(Y)_f)_0$, the degree-0 elements
in the localization $S(Y)_f$.
By Commutative Algebra Proposition
\ref{algebraic-geometry-prop-localization-integrally-closed},
localization of an integrally closed
ring remains integrally closed,
thus $S(Y)_f$ is integrally closed.
Let $A=(S(Y)_f)_0$.
If $x \in \text{Frac}(A)$ is integral
over $A$, then it is integral over
$S(Y)_f$.
Thus $x \in S(Y)_f$.
But $x$ is a quotient of degree-0
elements, hence has degree 0.
Therefore $x \in (S(Y)_f)_0=A$.

\item By the next item, the twisted quartic
is isomorphic to $\mathbb{P}^1$, hence is
normal.
It remains to see that it is not
projectively normal in this embedding.
Its homogeneous coordinate ring is
\[
R=k[t^4,t^3u,tu^3,u^4]\subset k[t,u].
\]
The element $t^2u^2$ lies in
$\text{Frac}(R)$, since
\[
t^2u^2=\frac{(t^3u)(u^4)}{tu^3}.
\]
Moreover it is integral over $R$, since
\[
(t^2u^2)^2=(t^3u)(tu^3)\in R.
\]
But $t^2u^2\notin R$: the exponent
$(2,2)$ is not in the semigroup generated
by
\[
(4,0),(3,1),(1,3),(0,4).
\]
Thus $R$ is not integrally closed.

\item Let
\[
\varphi:\mathbb{P}^1\to Y
\]
be the map
\[
[t:u]\longmapsto [t^4:t^3u:tu^3:u^4].
\]
We construct its inverse.
On the open subset of $Y$ where
$(x,y)\neq (0,0)$, define
\[
[x:y:z:w]\longmapsto [x:y].
\]
On the open subset where
$(z,w)\neq (0,0)$, define
\[
[x:y:z:w]\longmapsto [z:w].
\]
These opens cover $Y$.
On the image of $\varphi$ we have
\[
[x:y]=[t^4:t^3u]=[t:u]
\]
and
\[
[z:w]=[tu^3:u^4]=[t:u].
\]
Thus the two local formulae agree on
overlaps and define a morphism
$Y\to \mathbb{P}^1$ inverse to
$\varphi$.
Hence $Y\cong \mathbb{P}^1$.
Finally, the standard embedding of
$\mathbb{P}^1$ has homogeneous
coordinate ring $k[t,u]$,
which is integrally closed.
Thus $\mathbb{P}^1$ is projectively
normal.
\end{enumerate}
\end{proof}



\section{Divisors, invertible sheaves}
\label{algebraic-geometry-section-divisors-line-bundles}

\begin{definition}
\label{algebraic-geometry-definition-weil-divisor}
\begin{reference}
\href{https://stacks.math.columbia.edu/tag/%
0BE2}{0BE2},
\cite[p. 130]{hart}
\end{reference}
A {\it Weil divisor} on the variety $X$ is a
formal sum of irreducible
subvarieties over the integers.

More formally, a {\it prime divisor} is an
integral closed
subscheme $Z \subset X$ of codimension 1.
A {\it Weil divisor} is a formal sum $D=\sum
n_Z Z$ where
the sum is over prime divisors of $X$
and the collection
$\{Z|n_Z \neq  0\}$ is locally finite.
If all coefficients are non-negative we say
the divisor is {\it effective}.
\end{definition}

\noindent
A reasonable notion of order of vanishing
should then be introduced. In the case of
Riemann surfaces
this can be done using the order of
vanishings
of the denominator and numerator functions in
$f=g/h$ where $f \in K(X)^*$.
A more general construction is a bit more
involved.

\begin{definition}
\label{algebraic-geometry-definition-principal-weil-divisor}
\begin{reference}
\href{https://stacks.math.columbia.edu/tag/%
0be0}{0be0}
\end{reference}
Let $X$ be a Noetherian integral scheme.
Let $f \in R(X)^*$.
The {\it principal Weil divisor associated to
$f$}
is the Weil divisor
$$
\text{div}(f)=\text{div}_X(f)
=\sum \text{ord}_Z(f)[Z]
$$
where the sum is over prime divisors and
$\text{ord}_Z(f)$ is as in Definition ?.
This makes sense by Lemma ? [a lemma about
finiteness of the sum].
\end{definition}

\begin{definition}
\label{algebraic-geometry-definition-weil-divisor-class-group}
\begin{reference}
\href{https://stacks.math.columbia.edu/tag/%
0BE4}{0BE4}
\end{reference}
Let $X$ be a locally Noetherian scheme.
The {\it Weil divisor class group} of $X$,
denoted
$\text{Cl}(X)$, is the quotient
of the group of Weil divisors by the subgroup
of principal
Weil divisors.
\end{definition}

\noindent
The next section,
\href{https://stacks.math.columbia.edu/tag/%
02SE}{02SE},
shows there is an injective homomorphism
$\text{Pic}(X) \to \text{Cl}(X)$
which is also surjective if $X$ is normal
and the local rings of $X$ are UFDs
(which happens for smooth varieties).

\medskip\noindent
Earlier, in
\href{https://stacks.math.columbia.edu/tag/%
01WQ}{01WQ},
we have

\begin{definition}
\label{algebraic-geometry-definition-cartier-divisor}
Let $S$ be a scheme.
\begin{itemize}
\item A {\it localy principal closed
subscheme} of $S$
is a closed subscheme whose sheaf of ideals
is locally generated by a single element.

\item An {\it effective Cartier divisor} on
$S$
is a closed subscheme $D \subset S$ whose
ideal sheaf
$\mathcal{I}_D \subset \mathcal{O}_S$ is an
invertible
$\mathcal{O}_S$-module.
\end{itemize}
\end{definition}

\noindent
Followed by Lemma
\href{https://stacks.math.columbia.edu/tag/%
01WS}{01WS},
which essentially says that a Cartier
divisor is cut locally by a single function.
After this we would show the well-known
construction of sheaf associated to an
effective Cartier divisor
and eventually construct all group operations
and subtleties
to arrive at Lemma
\href{https://stacks.math.columbia.edu/tag/%
01X0}{01X0},
which says that effective Cartier divisors on
a scheme
are the same as invertible sheaves with fixed
regular global section.

\medskip\noindent
Here's two basic facts I'm not sure how to
prove:
\begin{enumerate}
\item
$$
\mathcal{O}_{\mathbb{P}^n}(1)\cong
\mathcal{O}_{\mathbb{P}^n}(H)
$$
where $\mathcal{O}_{\mathbb{P}^n}(1)$ is
defined
as dual of the tautological bundle,
and $H$ is any hyperplane (zeroes of a
homogeneous degree-1 polynomial
in $n+1$ variables) seen as an effective
Cartier divisor,
so that $\mathcal{O}_{\mathbb{P}^n}(H)$ is
its associated line bundle.

\item
$$
H^0(\mathbb{P}^n,\mathcal{O}_{\mathbb{P}^n}
(1))
\cong
k[x_0,…,x_n]_1,
$$
that is, the space of sections of
$\mathcal{O}_{\mathbb{P}^n}(1)$
is isomorphic to the space of homogeneous
polynomials
of degree 1 in $n+1$ variables
(each of which is a hyperplane in
$\mathbb{P}^n$).
I think this might be
\href{https://www.youtube.com/live/%
bXhgeRZMu8M?si=5BPdnECUREsCUwpV&t=4601}{
here}.
\end{enumerate}

\begin{exercise}
\label{algebraic-geometry-exercise-degree-d-surface}
Show that a hypersurface in $\mathbb{P}^3$
given by the zeroes of a homogeneous
polynomial
of degree $d$ has degree $d$,
i.e., if $S \subset \mathbb{P}^3$, $S=V(f)$
with
$f$ homogeneous of degree $d$,
then
 $$
\mathcal{O}_S(1).\mathcal{O}_S(1)=d.
$$
\end{exercise}

\begin{proof}
The exercise is kind of by definition no?
The intersection form counts the size of the
intersection when
the two classes are just two hypersurfaces
intersecting transversally (as is the case of
two
hyperplane sections in general position,
both of which have $\mathcal{O}_S(1)$ as
associated line bundle). The intersection
has $d$ points because its a line in
$\mathbb{P}^3$
intersecting a degree-$d$ polynomial.
And that's it.
\end{proof}

\begin{exercise}
\label{algebraic-geometry-exercise-intersection-lines}
Let $S \subset \mathbb{P}^3$ be a surface
of degree $d$ and $L \subset S$ a line.
Compute $L^2$ in $S$. Check that when $d=3$,
$L^2=-1$!
\end{exercise}

\begin{proof}
We apply the genus formula:
$$
(L.L+K_S)=2p_a-2
$$
By $L$ being a line we know that $p_a=0$
and that the intersection number of $L$
with any divisor of degree $d'$ will be $d'$.
Indeed, $L.d'H=d'(L.H)=d'$
because $L.H$ is the degree of the line
(by definition of degree, namely
the intersection number of $L$ with as many
hyperplanes as its dimension).
Thus
$$
\begin{aligned}
-2=(L.L+K_S)&=L^2+(L.K_S)=L^2+(L.\mathcal{O}
_S(d-4))
=L^2+d-4\\
\iff L^2=-d+2.
\end{aligned}
$$

\noindent
where the canonical bundle of $S$
is computed using adjunction formula
$K_S=(K_{\mathbb{P}^3}+N)|_S
=(\mathcal{O}_{\mathbb{P}^3}(-4)+\mathcal{O}
_{\mathbb{P}^3}(d))$.
\end{proof}

\begin{exercise}[Bézout]
\label{algebraic-geometry-exercise-bezout}
Show that if $C_d$ and $C_{d'}$ are two plane
curves
(i.e. contained in $\mathbb{P}^2$)
of degrees $d$ and $d'$ then
$C_d.C_{d'}=dd'$.
\end{exercise}

\noindent
One fun application of Bézout's theorem
is the following method for computing
degrees of line bundles in curves:

\begin{lemma}
\label{algebraic-geometry-lemma-line-bundle-degree}
Let $C\subset \mathbb P^n$ be a
projective curve. Then
$$
  \deg \mathcal O_C(m)=m\deg C.
$$
In particular, if $C\subset \mathbb P^2$
is a plane curve of degree $d$, then
$$
  \deg \mathcal O_C(m)=md.
$$
\end{lemma}

\begin{proof}
By definition, $\mathcal O_C(1)$ is the
restriction of the hyperplane bundle
$\mathcal O_{\mathbb P^n}(1)$ to $C$.

A section of $\mathcal O_C(1)$ is obtained
by restricting a hyperplane equation to
$C$. For a general hyperplane $H$, the
zero divisor of this section is
$$
  H\cap C.
$$
By definition, this divisor has degree
$\deg C$. Hence
$$
  \deg \mathcal O_C(1)=\deg C.
$$

Now
$$
  \mathcal O_C(m)
  =
  \mathcal O_C(1)^{\otimes m}.
$$
Since degree is additive under tensor
products of line bundles on curves, we get
$$
  \deg \mathcal O_C(m)
  =
  m\deg \mathcal O_C(1)
  =
  m\deg C.
$$
\end{proof}

\medskip\noindent
The Euler sequence on $\mathbb P^n$ is
$$
  0
  \to
  \mathcal O_{\mathbb P^n}
  \to
  \mathcal O_{\mathbb P^n}(1)^{\oplus n+1}
  \to
  T_{\mathbb P^n}
  \to
  0.
$$

The dual Euler sequence is
$$
  0
  \to
  \Omega^1_{\mathbb P^n}
  \to
  \mathcal O_{\mathbb P^n}(-1)^{\oplus n+1}
  \to
  \mathcal O_{\mathbb P^n}
  \to
  0.
$$

\begin{example}
\label{algebraic-geometry-example-canonical-projective-space}
We compute the canonical bundle of
$\mathbb P^n$.

From the dual Euler sequence
$$
  0
  \to
  \Omega^1_{\mathbb P^n}
  \to
  \mathcal O_{\mathbb P^n}(-1)^{\oplus n+1}
  \to
  \mathcal O_{\mathbb P^n}
  \to
  0,
$$
we take determinants. Since
$$
  \omega_{\mathbb P^n}
  =
  \det \Omega^1_{\mathbb P^n},
$$
we get
$$
  \omega_{\mathbb P^n}
  =
  \det
  \left(
    \mathcal O_{\mathbb P^n}(-1)^{\oplus n+1}
  \right)
  \otimes
  \det(\mathcal O_{\mathbb P^n})^{-1}.
$$
Hence
$$
  \omega_{\mathbb P^n}
  =
  \mathcal O_{\mathbb P^n}(-n-1).
$$

In particular,
$$
  \omega_{\mathbb P^2}
  =
  \mathcal O_{\mathbb P^2}(-3).
$$

Therefore, if
$C\subset \mathbb P^2$ is a plane curve
of degree $d$, adjunction gives
$$
  \omega_C
  =
  \left(
    \omega_{\mathbb P^2}
    \otimes
    \mathcal O_{\mathbb P^2}(d)
  \right)\big|_C
  =
  \mathcal O_C(d-3).
$$
\end{example}


\begin{proof}
Just notice that the line bundles associated
to
$C_d$ and $C_{d'}$ are $dH$ and $d'H$,
then use linearity of the bilinear of form.
To check why the divisor associated to $C_d$
is indeed $dH$ note that the intersection
number
of $C_d$ and a line $H$ is $d$ because $C_d$
is given by the roots of a polynomial of
degree $d$,
thus $C_d.H=d$, and since the Picard group of
$\mathbb{P}^2=\mathbb{Z}H$
we must have $C_d \sim dH$.
\end{proof}
\section{Adjunction formulas}
\label{algebraic-geometry-section-adjunction-formulas}

There are several statements called
adjunction formula
in different texts. All of them concern
``subvarieties'',
that is, closed embedded subschemes.

\begin{exercise}[Genus formula for a curve on
a surface]
\label{algebraic-geometry-exercise-genus-formu-curve-surfa}
Let $C \to X$ be a closed embedded subscheme
of dimension $1$ (as a topological space,
i.e. pure dimension)
inside a smooth surface $X$.
Then
$2p_a-2=(\mathcal{O}_X(C),\mathcal{O}_X(C))$.
\end{exercise}

\begin{proof}
Consider the ideal sheaf exact sequence
$$
\xymatrix{
0\ar[r]&\mathcal{O}_X(-C)\ar[r]&\mathcal{O}
_X\ar[r]&\mathcal{O}_C\ar[r]&0
}
$$
This sequence splits since there is an
obvious inverse morphism
to the inclusion $\mathcal{O}_X(-C)\to
\mathcal{O}_X$, namely
mapping a function $f$ to
Then $\mathcal{O}_X\cong \mathcal{O}_X(-C)
\oplus \mathcal{O}_C$.
 $\chi(\mathcal{O}_X)=\chi$
\end{proof}


\section{Linear systems}
\label{algebraic-geometry-section-linear-systems}

\noindent
A basic fact is that
$$
H^0(\mathcal{O}(d)) \cong
\{\text{homogeneous polynomials of degree
$d$}\}.
$$
And more generally,
$$
|D|=\mathbb{P}H^0(\mathcal{O}(D)).
$$

\begin{exercise}
\label{algebraic-geometry-exercise-hodge-numbers-blow-up-9pts}
Let $C_1,C_2$ be two cubic curves in
$\mathbb{P}^2$.
Compute the Hodge numbers of the blowup of
the 9 points
of intersection.
\end{exercise}

\begin{proof}
The solution is:
$$
\begin{aligned}\xymatrix@C=.5em@R=.5em{
&  &   1\\
&  0&  &  0\\
0&&  10&  &  0\\
&  0&  &  0\\
&  &  1
}\end{aligned}
$$
Which follows from:
\begin{itemize}
\item $h^{0,0}=1$ because $X$ is connected.
\item $h^{1,0}=0$ because $X$ is simply
connected.
\item $h^{2,0}=0$ because the canonical
bundle of $X$ has no sections.
\item $h^{1,1}=10$ because in the exponential
exact sequence
we get $\text{Pic} \cong H^2(X,\mathbb{Z})$
(by vanishing irregularity and genus)
and we know
$\text{Pic}=\mathbb{Z}H \oplus
\mathbb{Z}E_1\oplus…\oplus \mathbb{Z}E_9$.
\end{itemize}
\end{proof}

\begin{exercise}
\label{algebraic-geometry-exercise-line-bundle-degree}
Let $C\subset \mathbb P^n$ be a
projective curve. Explain what
$\mathcal O_C(1)$ means and compute its
degree.
\end{exercise}

\begin{proof}
By definition,
$$
  \mathcal O_C(1)
  =
  \mathcal O_{\mathbb P^n}(1)|_C.
$$
Thus $\mathcal O_C(1)$ is the hyperplane
bundle restricted to $C$.

A section of $\mathcal O_C(1)$ is obtained
by restricting a hyperplane equation to
$C$. Its zero divisor is
$$
  C\cap H,
$$
where $H\subset \mathbb P^n$ is a
hyperplane.

Therefore
$$
  \deg \mathcal O_C(1)=\deg C.
$$
More generally,
$$
  \deg \mathcal O_C(m)=m\deg C.
$$
This is Bezout in line-bundle language.
\end{proof}

\begin{exercise}
\label{algebraic-geometry-exercise-plane-section-dimension}
Let $C\subset \mathbb P^2$ be a plane
curve not contained in any line. Compute
$h^0(C,\mathcal O_C(1))$ and describe
$|\mathcal O_C(1)|$.
\end{exercise}

\begin{proof}
The global sections of
$\mathcal O_{\mathbb P^2}(1)$ are the
linear forms
$$
  aX+bY+cZ.
$$
Thus
$$
  H^0(\mathbb P^2,\mathcal O(1))
  =
  \langle X,Y,Z\rangle.
$$

Restricting to $C$, we get sections
$$
  X|_C,\quad Y|_C,\quad Z|_C
$$
of $\mathcal O_C(1)$.

They are linearly independent. Indeed,
if
$$
  aX+bY+cZ
$$
vanished on all of $C$, then $C$ would be
contained in the line
$$
  aX+bY+cZ=0,
$$
contradicting the assumption.

Hence
$$
  h^0(C,\mathcal O_C(1))=3.
$$
Therefore
$$
  |\mathcal O_C(1)|
  =
  \mathbb P(H^0(C,\mathcal O_C(1)))
  \simeq
  \mathbb P^2.
$$
\end{proof}

\begin{exercise}
\label{algebraic-geometry-exercise-sections-map}
Let $L$ be a line bundle on a variety $X$,
and let $s_0,\ldots,s_N$ be global
sections with no common zero. Describe
the morphism induced by these sections.
\end{exercise}

\begin{proof}
The sections define a morphism
$$
  \varphi:X\to \mathbb P^N
$$
by
$$
  p\mapsto [s_0(p):\cdots:s_N(p)].
$$
The condition that the sections have no
common zero guarantees that this point of
projective space is defined for every
$p\in X$.

In particular, for
$C\subset \mathbb P^2$ and
$L=\mathcal O_C(1)$, the sections
$$
  X|_C,\quad Y|_C,\quad Z|_C
$$
give the original embedding
$$
  C\hookrightarrow \mathbb P^2.
$$
\end{proof}

\begin{exercise}
\label{algebraic-geometry-exercise-conic-pullback}
Let
$$
  C=V(Y^2-XZ)\subset \mathbb P^2.
$$
Use the parametrization
$$
  \nu:\mathbb P^1\to C,
  \qquad
  [U:V]\mapsto [U^2:UV:V^2].
$$
Compute $\nu^*\mathcal O_C(1)$ and the
pulled-back coordinate sections.
\end{exercise}

\begin{proof}
The coordinate sections of
$\mathcal O_C(1)$ are
$$
  X|_C,\quad Y|_C,\quad Z|_C.
$$
Pulling back by $\nu$ gives
$$
  \nu^*X=U^2,
  \qquad
  \nu^*Y=UV,
  \qquad
  \nu^*Z=V^2.
$$
These are homogeneous polynomials of
degree $2$ on $\mathbb P^1$. Therefore
$$
  \nu^*\mathcal O_C(1)
  \cong
  \mathcal O_{\mathbb P^1}(2).
$$

Moreover,
$$
  H^0(C,\mathcal O_C(1))
$$
pulls back to the subspace
$$
  \langle U^2,UV,V^2\rangle
  \subset
  H^0(\mathbb P^1,\mathcal O(2)).
$$
This is all of
$H^0(\mathbb P^1,\mathcal O(2))$.
Hence
$$
  h^0(C,\mathcal O_C(1))=3.
$$
\end{proof}

\begin{exercise}
\label{algebraic-geometry-exercise-twisted-cubic}
Let $C\subset \mathbb P^3$ be the
twisted cubic, parametrized by
$$
  \nu:\mathbb P^1\to C,
  \qquad
  [U:V]\mapsto
  [U^3:U^2V:UV^2:V^3].
$$
Compute $\nu^*\mathcal O_C(1)$,
$h^0(C,\mathcal O_C(1))$, and
$|\mathcal O_C(1)|$.
\end{exercise}

\begin{proof}
The coordinate sections pull back to
$$
  U^3,\quad U^2V,\quad UV^2,\quad V^3.
$$
These are homogeneous polynomials of
degree $3$ on $\mathbb P^1$. Hence
$$
  \nu^*\mathcal O_C(1)
  \cong
  \mathcal O_{\mathbb P^1}(3).
$$

The four sections
$$
  U^3,\quad U^2V,\quad UV^2,\quad V^3
$$
form a basis of
$$
  H^0(\mathbb P^1,\mathcal O(3)).
$$
Therefore
$$
  h^0(C,\mathcal O_C(1))=4.
$$
Thus
$$
  |\mathcal O_C(1)|
  \simeq
  \mathbb P^3.
$$
The induced morphism is the original
twisted cubic embedding
$$
  C\hookrightarrow \mathbb P^3.
$$
\end{proof}

\begin{exercise}
\label{algebraic-geometry-exercise-nodal-cubic}
Let $C$ be the nodal cubic
$$
  C: Y^2Z=X^2Z+X^3.
$$
Its normalization is
$$
  \nu:\mathbb P^1\to C,
$$
given by
$$
  [U:V]\mapsto
  [V(U^2-V^2):U(U^2-V^2):V^3].
$$
Compute $\nu^*\mathcal O_C(1)$ and
explain why
$$
  h^0(C,\mathcal O_C(1))=3.
$$
\end{exercise}

\begin{proof}
The coordinate sections pull back to
$$
  V(U^2-V^2),
  \qquad
  U(U^2-V^2),
  \qquad
  V^3.
$$
These are homogeneous polynomials of
degree $3$ on $\mathbb P^1$. Hence
$$
  \nu^*\mathcal O_C(1)
  \cong
  \mathcal O_{\mathbb P^1}(3).
$$

However,
$$
  h^0(C,\mathcal O_C(1))
$$
is not equal to
$$
  h^0(\mathbb P^1,\mathcal O(3))=4.
$$
Only the sections on the normalization
which descend to $C$ are allowed.

The node has two preimages:
$$
  [1:1],
  \qquad
  [-1:1].
$$
A section descending to $C$ must take
compatible values at these two points.

The pulled-back coordinate sections span
the $3$-dimensional subspace
$$
  \langle
  V(U^2-V^2),
  U(U^2-V^2),
  V^3
  \rangle
  \subset
  H^0(\mathbb P^1,\mathcal O(3)).
$$
Therefore
$$
  h^0(C,\mathcal O_C(1))=3.
$$
\end{proof}

\begin{exercise}
\label{algebraic-geometry-exercise-cuspidal-cubic}
Let $C$ be the cuspidal cubic
$$
  C: Y^2Z=X^3.
$$
Its normalization is
$$
  \nu:\mathbb P^1\to C,
  \qquad
  [U:V]\mapsto [U^2V:U^3:V^3].
$$
Compute $\nu^*\mathcal O_C(1)$ and
$h^0(C,\mathcal O_C(1))$.
\end{exercise}

\begin{proof}
The coordinate sections pull back to
$$
  U^2V,
  \qquad
  U^3,
  \qquad
  V^3.
$$
These are homogeneous polynomials of
degree $3$. Hence
$$
  \nu^*\mathcal O_C(1)
  \cong
  \mathcal O_{\mathbb P^1}(3).
$$

The sections descending from $C$ are
spanned by
$$
  U^2V,\quad U^3,\quad V^3.
$$
This is a $3$-dimensional subspace of
$$
  H^0(\mathbb P^1,\mathcal O(3)).
$$
Therefore
$$
  h^0(C,\mathcal O_C(1))=3.
$$
\end{proof}

\begin{exercise}
\label{algebraic-geometry-exercise-plane-quartic-canonical}
Let $C\subset \mathbb P^2$ be a smooth
plane quartic. Compute $K_C$,
$\deg K_C$, $h^0(K_C)$, and describe the
canonical map.
\end{exercise}

\begin{proof}
By adjunction for plane curves,
$$
  K_C=\mathcal O_C(d-3).
$$
Here $d=4$, so
$$
  K_C=\mathcal O_C(1).
$$
Since $C$ has degree $4$,
$$
  \deg K_C
  =
  \deg \mathcal O_C(1)
  =
  4.
$$
Also a smooth plane quartic has genus
$$
  g=\frac{(4-1)(4-2)}2=3.
$$
Thus
$$
  h^0(K_C)=g=3.
$$
The canonical linear system is
$$
  |K_C|=|\mathcal O_C(1)|.
$$
Therefore the canonical map is the given
embedding
$$
  C\hookrightarrow \mathbb P^2.
$$
\end{proof}

\section{Ample line bundles}
\label{algebraic-geometry-section-ample-line-bundles}


\begin{exercise}
\label{algebraic-geometry-exercise-ample-bundle-on-K3}
Let $L$ be an ample bundle on a K3 surface
$M$. Prove that
$\mathcal{L}^{\otimes 2}$ is globally
generated (that is, for each $x\in M$
there exsits a section $h \in
H^{0}(L^{\otimes 2})$ which does not vanish
in
$x$).
\end{exercise}

\begin{proof}
This just asks that the $n$ in Definition
\ref{algebraic-geometry-definition-ample} is $2$ for all
$x\in X$. Because, again, $x\in X_s$ means
that $s(x)\neq 0$ because if it was,
then we could somehow write $s$ as a product
of a vanishing function on
$\mathfrak{m}_x$ and a local frame of
$\Gamma(X,\mathcal{L})$. But I guess for
the exercise do this: a line bundle is {\it
ample} if there is $n$ such that the
canonical embedding (cf Lemma
\ref{algebraic-geometry-lemma-map-into-proj}) is an embedding,
i.e.
that $\mathcal{L}^{\otimes n}$ is {\it very
ample}. (Interestingly, the notion
very ampleness is defined in morphisms.tex.)
\end{proof}


\section{Riemann-Roch theorem for curves}
\label{algebraic-geometry-section-riemann-roch}

\noindent
Let $D=\sum n_p p$ be a divisor on a curve.
How many rational functions $z \in K^*$
satisfy
$\text{div}(z) \geq -D$? That is, rational
functions
which can only have poles of order $\leq n_p$
at the points $p$ with $n_p>0$,
and can only have zeros of order $\geq
|n_p|$
at the points $p$ with $n_p<0$.

For example, fix a point $p$ in the curve.
Then $\mathcal{O}_C(p)$ is the space
of rational functions with a simple pole at
$p$.
$\mathcal{O}_C(2p)$ is the space of rational
functions
with a pole of order at most 2 at $p$, and so
on.
Note that $\mathcal{O}_C(p) \subset
\mathcal{O}_X (2p)
\subset \mathcal{O}(3p)…$.
That is, the more poles we allow, this space
of functions will be larger.

In classical (Riemann's) notation, this space
was denoted
$L(D)$, and its dimension $\ell(D)$.
Modern notation reads $H^0(\mathcal{O}_C(D))$
instead of $L(D)$ and
$h^0(\mathcal{O}_C(D))$
instead of $\ell(D)$.

Riemann-Roch formula for curves is
\begin{equation}
\label{algebraic-geometry-equation-riemann-roch-curves}
h^0(D)-h^0(K-D)=\text{deg}(D)-g+1,
\end{equation}

\noindent
where $K$ is a divisor associated to the
curve canonically.
This is done by choosing a meromorphic
section $s$
of the bundle $\Omega_C^1$
and taking $K:=\text{div}(s)$.
Then one shows that this is independent of
the choice of $s$.

Note that by Serre duality
$h^1(D)=h^0(K-D)$.

\medskip\noindent
Here's the hands-on criterion for
computing the genus of a plane curve:

\begin{proposition}
\label{algebraic-geometry-proposition-plane-genus}
Let $C\subset \mathbb P^2$ be a plane
curve of degree $d$. Then its arithmetic
genus is
$$
  p_a(C)=\frac{(d-1)(d-2)}{2}.
$$
\end{proposition}

\begin{proof}
By definition,
$$
  p_a(C)=1-\chi(\mathcal O_C).
$$
Since $C$ is cut out by one homogeneous
polynomial of degree $d$, we have the
exact sequence
$$
  0
  \to
  \mathcal O_{\mathbb P^2}(-d)
  \to
  \mathcal O_{\mathbb P^2}
  \to
  \mathcal O_C
  \to
  0.
$$
Taking Euler characteristics gives
$$
  \chi(\mathcal O_C)
  =
  \chi(\mathcal O_{\mathbb P^2})
  -
  \chi(\mathcal O_{\mathbb P^2}(-d)).
$$
We use the standard formula
$$
  \chi(\mathcal O_{\mathbb P^2}(m))
  =
  \binom{m+2}{2}.
$$
Thus
$$
  \chi(\mathcal O_{\mathbb P^2})=1
$$
and
$$
  \chi(\mathcal O_{\mathbb P^2}(-d))
  =
  \binom{2-d}{2}
  =
  \frac{(d-1)(d-2)}{2}.
$$
Therefore
$$
  \chi(\mathcal O_C)
  =
  1-\frac{(d-1)(d-2)}{2}.
$$
Hence
$$
  p_a(C)
  =
  1-\chi(\mathcal O_C)
  =
  \frac{(d-1)(d-2)}{2}.
$$
\end{proof}

\begin{exercise}
\label{algebraic-geometry-exercise-riemann-roch-canonical}
Apply Riemann-Roch to $D=K_C$.
\end{exercise}

\begin{proof}
For $D=K_C$, Riemann-Roch gives
$$
  h^0(K_C)-h^0(0)
  =
  \deg K_C-g+1.
$$
We know
$$
  h^0(K_C)=g
  \qquad
  \text{and}
  \qquad
  h^0(0)=1.
$$
Hence
$$
  g-1=\deg K_C-g+1.
$$
Therefore
$$
  \deg K_C=2g-2.
$$
\end{proof}

\begin{exercise}
\label{algebraic-geometry-exercise-high-degree-divisor}
Let $D$ be a divisor on $C$ with
$$
  \deg D>2g-2.
$$
Compute $h^0(D)$.
\end{exercise}

\begin{proof}
Since
$$
  \deg K_C=2g-2,
$$
we have
$$
  \deg(K_C-D)<0.
$$
A line bundle of negative degree has no
nonzero global sections, so
$$
  h^0(K_C-D)=0.
$$
Thus Riemann-Roch gives
$$
  h^0(D)=\deg D-g+1.
$$
\end{proof}

\begin{exercise}
\label{algebraic-geometry-exercise-p1-sections}
Compute $h^0(\mathbb P^1,\mathcal O(d))$
for $d\geq 0$.
\end{exercise}

\begin{proof}
On $\mathbb P^1$, we have $g=0$.
Also
$$
  \deg \mathcal O(d)=d.
$$
Since $d\geq 0>2g-2=-2$, Riemann-Roch
gives
$$
  h^0(\mathcal O(d))
  =
  d-0+1
  =
  d+1.
$$
Therefore
$$
  h^0(\mathbb P^1,\mathcal O(d))=d+1.
$$
\end{proof}

\begin{exercise}
\label{algebraic-geometry-exercise-elliptic-point}
Let $E$ be an elliptic curve, and let
$p\in E$. Compute
$$
  h^0(E,\mathcal O_E(p)).
$$
\end{exercise}

\begin{proof}
An elliptic curve has genus $1$, and
$$
  K_E\sim 0
$$
(because the standard top-degree
holomorphic form $dz^1 \wedge…\wedge dz^n$
descends to $\mathbb{C}^n/\mathbb{Z}^n$).

Apply Riemann-Roch to $D=p$:
$$
  h^0(p)-h^0(K_E-p)
  =
  \deg p-1+1.
$$
Since $K_E\sim 0$, this becomes
$$
  h^0(p)-h^0(-p)=1.
$$
But
$$
  \deg(-p)<0,
$$
so
$$
  h^0(-p)=0.
$$
Therefore
$$
  h^0(E,\mathcal O_E(p))=1.
$$
\end{proof}

\begin{exercise}
\label{algebraic-geometry-exercise-plane-quartic-genus}
Let $C\subset \mathbb P^2$ be a smooth
plane quartic. Compute $g(C)$,
$K_C$, $\deg K_C$, and $h^0(K_C)$.
\end{exercise}

\begin{proof}
For a smooth plane curve of degree $d$,
$$
  g(C)=\frac{(d-1)(d-2)}2.
$$
Here $d=4$, so
$$
  g(C)=\frac{3\cdot 2}{2}=3.
$$

By adjunction,
$$
  K_C=\mathcal O_C(d-3).
$$
Thus
$$
  K_C=\mathcal O_C(1).
$$

Since $C$ has degree $4$,
$$
  \deg \mathcal O_C(1)=4.
$$
Therefore
$$
  \deg K_C=4.
$$

Finally,
$$
  h^0(K_C)=g(C)=3.
$$
\end{proof}

\begin{exercise}
\label{algebraic-geometry-exercise-plane-quintic}
Let $C\subset \mathbb P^2$ be a smooth
plane quintic. Compute:
$$
  g(C),\quad
  K_C,\quad
  \deg K_C,\quad
  h^0(K_C),
  \quad
  h^0(\mathcal O_C(1)),
$$
$$
  \deg \mathcal O_C(2),
  \quad
  h^0(\mathcal O_C(2)).
$$
Also describe the morphism induced by
$|\mathcal O_C(1)|$.
\end{exercise}

\begin{proof}
Since $d=5$,
$$
  g(C)
  =
  \frac{(5-1)(5-2)}2
  =
  6.
$$

By adjunction,
$$
  K_C=\mathcal O_C(5-3)
  =
  \mathcal O_C(2).
$$

Since
$$
  \deg \mathcal O_C(1)=5,
$$
we get
$$
  \deg K_C
  =
  \deg \mathcal O_C(2)
  =
  10.
$$

Also
$$
  h^0(K_C)=g(C)=6.
$$

The sections of $\mathcal O_C(1)$ are
the restrictions of $X,Y,Z$, so
$$
  h^0(\mathcal O_C(1))=3.
$$

Finally,
$$
  h^0(\mathcal O_C(2))
  =
  h^0(K_C)
  =
  6.
$$

The linear system $|\mathcal O_C(1)|$
induces the original embedding
$$
  C\hookrightarrow \mathbb P^2.
$$
\end{proof}

\section{Surfaces}
\label{algebraic-geometry-section-surfaces}

Main example:

$$
  S=\operatorname{Bl}_p\mathbb P^2.
$$

Let

$$
  H=\pi^*(\text{line}),
  \qquad
  E=\text{exceptional divisor}.
$$

Then

$$
  \operatorname{Pic}(S)
  =
  \mathbb ZH\oplus \mathbb ZE.
$$

The intersection table is

$$
  H^2=1,
  \qquad
  H\cdot E=0,
  \qquad
  E^2=-1,
$$
where
$H$ is the total transform of a line,
$E\simeq \mathbb P^1$ is the exceptional
  curve, and
$E^2=-1$. To check that $E^2 = -1$,
let $L$ be a line through $p$, and pull
it back to the blow-up. We obtain:
total transform = strict transform $+ E$,
that is: $H = \tilde{L} + E$,
which gives $1 = H^2 = L^2 + L\cdot E + E^2
=1+1+E^2$, since the strict
transform does intersect $E$.

The canonical divisor is computed by
$$
K_S=\pi^*K_{\mathbb P^2}+E,
$$
which may be seen by pulling back
a top form on $\mathbb{P}^2$,
and when writing its coordinates
at the blow up, like in the routine
exercises, an extra term appears,
i.e. something of the kind $x dx \wedge dy$
(see Exercise
\ref{algebraic-geometry-exercise-canonical-blow-up}).

Since
$$
K_{\mathbb P^2}=-3H,
$$
we get
$$
K_S=-3H+E.
$$
More generally, if we blow up points
$p_1,\ldots,p_r$, then
$$
 K_S=-3H+E_1+\cdots+E_r.
$$
If $C\subset \mathbb P^2$ is a curve of
degree $d$ and multiplicity $m$ at $p$,
then its strict transform has class

$$
\widetilde C=dH-mE.
$$

For several blown-up points,

$$
\widetilde C
=
dH-m_1E_1-\cdots-m_rE_r.
$$
Slogan:

blowing up adds an exceptional divisor
$E$ with $E^2=-1$, and strict transforms
are computed by subtracting the
multiplicity times $E$.

\begin{exercise}
\label{algebraic-geometry-exercise-canonical-blow-up}
Explain locally why, for the blow-up of
a point on a surface,
$$
  K_S=\pi^*K_X+E.
$$
\end{exercise}

\begin{proof}
Use local coordinates
$$
  \pi(x,u)=(x,y)=(x,xu).
$$
A local generator of the canonical bundle
downstairs is
$$
  dx\wedge dy.
$$
Pull it back:
$$
  \pi^*(dx\wedge dy)
  =
  dx\wedge d(xu).
$$
Since
$$
  d(xu)=u\,dx+x\,du,
$$
we get
$$
  \pi^*(dx\wedge dy)
  =
  dx\wedge (u\,dx+x\,du)
  =
  x\,dx\wedge du.
$$
Thus the pulled-back volume form has one
extra zero along
$$
  E=\{x=0\}.
$$
Zeros of top forms define the canonical
divisor. Hence
$$
  K_S=\pi^*K_X+E.
$$
\end{proof}

\begin{exercise}
\label{algebraic-geometry-exercise-blowup-strict-transform}
Let $C\subset \mathbb P^2$ be a curve
of degree $d$ with multiplicity $m$ at
$p$. What is the class of the strict
transform of $C$ in
$\operatorname{Bl}_p\mathbb P^2$?
\end{exercise}

\begin{proof}
The total transform has class $dH$.
Since $C$ has multiplicity $m$ at $p$,
the exceptional divisor appears with
multiplicity $m$ in the total transform.
Thus
$$
  \pi^*C=\widetilde C+mE.
$$
Therefore
$$
  \widetilde C=dH-mE.
$$
\end{proof}

\begin{exercise}
\label{algebraic-geometry-exercise-blowup-line}
On $S=\operatorname{Bl}_p\mathbb P^2$,
compute the self-intersection of the
strict transform of a line through $p$.
\end{exercise}

\begin{proof}
A line through $p$ has degree $1$ and
multiplicity $1$ at $p$. Hence its
strict transform has class
$$
  H-E.
$$
Therefore
$$
  (H-E)^2
  =
  H^2-2H\cdot E+E^2.
$$
Using
$$
  H^2=1,
  \qquad
  H\cdot E=0,
  \qquad
  E^2=-1,
$$
we get
$$
  (H-E)^2=1-1=0.
$$
\end{proof}

\begin{exercise}
\label{algebraic-geometry-exercise-blowup-conic}
On $S=\operatorname{Bl}_p\mathbb P^2$,
compute the self-intersection of the
strict transform of a conic through $p$
with multiplicity $2$.
\end{exercise}

\begin{proof}
The conic has degree $2$ and multiplicity
$2$ at $p$. Thus its strict transform has
class
$$
  2H-2E.
$$
Hence
$$
  (2H-2E)^2
  =
  4H^2-8H\cdot E+4E^2.
$$
Using the intersection table, this gives
$$
  (2H-2E)^2=4-4=0.
$$
\end{proof}

\begin{exercise}
\label{algebraic-geometry-exercise-two-point-blowup}
Let
$$
  S=\operatorname{Bl}_{p,q}\mathbb P^2,
$$
where $p$ and $q$ are distinct points.
Let $H=\pi^*(\text{line})$, and let
$E_p,E_q$ be the exceptional divisors.
Compute the intersection table and
$K_S$.
\end{exercise}

\begin{proof}
The Picard group is generated by
$$
  H,\quad E_p,\quad E_q.
$$
The intersection table is
$$
  H^2=1,
  \qquad
  E_p^2=E_q^2=-1,
$$
and
$$
  H\cdot E_p=H\cdot E_q=0,
  \qquad
  E_p\cdot E_q=0.
$$
The exceptional divisors are disjoint,
so $E_p\cdot E_q=0$.

The canonical divisor is
$$
  K_S=\pi^*K_{\mathbb P^2}+E_p+E_q.
$$
Since $K_{\mathbb P^2}=-3H$, we get
$$
  K_S=-3H+E_p+E_q.
$$
\end{proof}

\begin{exercise}
\label{algebraic-geometry-exercise-two-point-curves}
On
$$
  S=\operatorname{Bl}_{p,q}\mathbb P^2,
$$
compute the classes and self-intersections
of the following curves:
a line through $p$ but not $q$;
the line through both $p$ and $q$;
a conic through $p$ and $q$.
\end{exercise}

\begin{proof}
A line through $p$ but not $q$ has class
$$
  H-E_p.
$$
Its self-intersection is
$$
  (H-E_p)^2=1-1=0.
$$

The line through both $p$ and $q$ has
class
$$
  H-E_p-E_q.
$$
Its self-intersection is
$$
  (H-E_p-E_q)^2=1-1-1=-1.
$$

A conic through $p$ and $q$ has class
$$
  2H-E_p-E_q.
$$
Its self-intersection is
$$
  (2H-E_p-E_q)^2=4-1-1=2.
$$
\end{proof}

\begin{exercise}
\label{algebraic-geometry-exercise-minus-one-curves}
Find the basic $(-1)$-curves on
$$
  \operatorname{Bl}_{p,q}\mathbb P^2.
$$
\end{exercise}

\begin{proof}
The exceptional curves themselves have
self-intersection $-1$:
$$
  E_p^2=E_q^2=-1.
$$
Also, the strict transform of the line
through both points has class
$$
  H-E_p-E_q,
$$
and
$$
  (H-E_p-E_q)^2=-1.
$$
Thus the basic $(-1)$-curves are
$$
  E_p,\qquad E_q,\qquad H-E_p-E_q.
$$
\end{proof}


\begin{exercise}
\label{algebraic-geometry-exercise-quartic-surface}
Let $X\subset \mathbb P^3$ be a smooth
quartic surface. Compute $\dim X$,
$\deg X$, $K_X$, and $H^2$, where
$H=\mathcal O_X(1)$. Classify $X$.
\end{exercise}

\begin{proof}
Since $X$ is a hypersurface in
$\mathbb P^3$, we have
$$
  \dim X=2.
$$
Its degree is
$$
  \deg X=4.
$$
By adjunction,
$$
  K_X
  =
  \mathcal O_X(4-3-1)
  =
  \mathcal O_X.
$$
Also
$$
  H^2=\deg X=4.
$$
Thus $X$ is a smooth surface with trivial
canonical bundle. Moreover, smooth
complete intersections in projective
space satisfy $H^1(X,\mathcal O_X)=0$.
Hence $X$ is a K3 surface.
\end{proof}

\begin{exercise}
\label{algebraic-geometry-exercise-two-three-surface}
Let
$$
  X=(2,3)\subset \mathbb P^4
$$
be a smooth complete intersection.
Compute $\dim X$, $\deg X$, $K_X$,
$H^2$, $K_X^2$, and
$h^0(X,\mathcal O_X(1))$. Classify $X$.
\end{exercise}

\begin{proof}
The codimension is $2$, hence
$$
  \dim X=4-2=2.
$$
The degree of a complete intersection is
the product of the degrees:
$$
  \deg X=2\cdot 3=6.
$$
The normal bundle is
$$
  N_{X/\mathbb P^4}
  =
  \mathcal O_X(2)\oplus \mathcal O_X(3).
$$
Adjunction gives
$$
  K_X
  =
  \mathcal O_X(2+3-4-1)
  =
  \mathcal O_X.
$$
Thus
$$
  K_X^2=0.
$$
Also
$$
  H^2=\deg X=6.
$$

The global sections of $\mathcal O_X(1)$
are the linear forms on $\mathbb P^4$
restricted to $X$. No nonzero linear form
vanishes on all of $X$, because the ideal
of $X$ is generated in degrees $2$ and
$3$. Therefore
$$
  h^0(X,\mathcal O_X(1))
  =
  h^0(\mathbb P^4,\mathcal O(1))
  =
  5.
$$
The linear system $|\mathcal O_X(1)|$
gives the original embedding
$$
  X\hookrightarrow \mathbb P^4.
$$

Since $K_X\simeq \mathcal O_X$ and
$H^1(X,\mathcal O_X)=0$, this is a
complete-intersection K3 surface.
\end{proof}

\begin{exercise}
\label{algebraic-geometry-exercise-two-two-two-surface}
Let
$$
  X=(2,2,2)\subset \mathbb P^5
$$
be a smooth complete intersection.
Compute $\dim X$, $\deg X$, $K_X$,
$H^2$, $K_X^2$, and
$h^0(X,\mathcal O_X(1))$. Classify $X$.
\end{exercise}

\begin{proof}
The codimension is $3$, hence
$$
  \dim X=5-3=2.
$$
The degree is
$$
  \deg X=2\cdot 2\cdot 2=8.
$$
By adjunction,
$$
  K_X
  =
  \mathcal O_X(2+2+2-5-1)
  =
  \mathcal O_X.
$$
Hence
$$
  K_X^2=0.
$$
Also
$$
  H^2=\deg X=8.
$$

Again, no nonzero linear form vanishes
on all of $X$, since the ideal is
generated in degree $2$. Therefore
$$
  h^0(X,\mathcal O_X(1))
  =
  h^0(\mathbb P^5,\mathcal O(1))
  =
  6.
$$
The linear system $|\mathcal O_X(1)|$
gives the original embedding
$$
  X\hookrightarrow \mathbb P^5.
$$

Thus $X$ is a complete-intersection
K3 surface. (For some reason it turns out
that complete intersections have
$h^1(\mathcal{O}_X)=0$).
\end{proof}

\begin{exercise}
\label{algebraic-geometry-exercise-surface-intersection}
Let $X\subset \mathbb P^n$ be a
projective surface, and let
$H=\mathcal O_X(1)$. Show that
$$
  H^2=\deg X.
$$
Deduce that
$$
  \mathcal O_X(a)\cdot\mathcal O_X(b)
  =
  ab\deg X.
$$
\end{exercise}

\begin{proof}
The class $H$ is the hyperplane class
restricted to $X$. Thus $H^2$ is the
intersection of $X$ with two general
hyperplanes. By definition, this number
is the degree of $X$. Hence
$$
  H^2=\deg X.
$$
Since
$$
  \mathcal O_X(a)=aH,
  \qquad
  \mathcal O_X(b)=bH,
$$
we get
$$
  \mathcal O_X(a)\cdot\mathcal O_X(b)
  =
  (aH)\cdot(bH)
  =
  abH^2
  =
  ab\deg X.
$$
\end{proof}

\section{Schemes}
\label{algebraic-geometry-section-shemes}

For a definition of presheaf,
see Categories, Definition
\ref{categories-definition-presheaf}.

\begin{definition}
\label{algebraic-geometry-definition-sheaf}
Let $X$ be a topological space.
\begin{enumerate}
\item A {\it sheaf $\mathcal{F}$ of sets on
$X$} is a presheaf
of sets which satisfies the following
additional property: Given
any open covering $U = \bigcup_{i \in I} U_i$
and any collection
of sections $s_i \in \mathcal{F}(U_i)$, $i
\in I$ such that
$\forall i, j\in I$
$$
s_i|_{U_i \cap U_j} = s_j|_{U_i \cap U_j}
$$
there exists a unique section $s \in
\mathcal{F}(U)$ such that
$s_i = s|_{U_i}$ for all $i \in I$.
\item A {\it morphism of sheaves of sets} is
simply a
morphism of presheaves of sets.
\item The category of sheaves of sets on $X$
is denoted
$\Sh(X)$.
\end{enumerate}
\end{definition}


\medskip\noindent
Let $X$ be a topological space. Let $x \in X$
be a point.
Let $\mathcal{F}$ be a presheaf of sets on
$X$.
The {\it stalk of $\mathcal{F}$ at $x$} is
the set
$$
\mathcal{F}_x
=
\colim_{x\in U} \mathcal{F}(U)
$$
where the colimit is over the set of open
neighbourhoods
$U$ of $x$ in $X$. The set of open
neighbourhoods is
partially ordered by (reverse) inclusion:
We say $U \geq U' \Leftrightarrow U \subset
U'$.
The transition maps in the system are
given by the restriction maps of
$\mathcal{F}$.
See Categories, Section
\ref{categories-section-posets-limits}
for notation and terminology regarding
(co)limits over systems.
Note that the colimit is a directed colimit.
Thus it is easy to describe $\mathcal{F}_x$.
Namely,
$$
\mathcal{F}_x
=
\{
(U, s)
\mid
x\in U, s\in \mathcal{F}(U)
\}/\sim
$$
with equivalence relation given by $(U, s)
\sim (U', s')$ if and only if
there exists an open $U'' \subset U \cap U'$
with $x \in U''$ and
$s|_{U''} = s'|_{U''}$. Given a pair $(U, s)$
we sometimes denote
$s_x$ the element of $\mathcal{F}_x$
corresponding to the equivalence
class of $(U, x)$. We sometimes use the
phrase
``image of $s$ in $\mathcal{F}_x$'' to denote
$s_x$.
For example, given two pairs $(U, s)$ and
$(U', s')$ we sometimes
say ``$s$ is equal to $s'$ in
$\mathcal{F}_x$'' to indicate
that $s_x = s'_x$. Other authors use the
terminology
``germ of $s$ at $x$''.

\medskip\noindent
Now let's discuss closed immersions.

\noindent
Closed immersions of ringes spaces are in
\href{https://stacks.math.columbia.edu/tag/%
01C1}{01C1}.

\begin{definition}
\label{algebraic-geometry-definition-closed-immersion}
A map of ringed spaces
$i:(Z,\mathcal{O}_Z) \to (X,\mathcal{O}_X)$
is a {\it closed immersion}
if $i$ is injective,
its image is closed (I think these two are
``closed immersion of topological spaces'')
and the induced map $\mathcal{O}_X \to
i_*\mathcal{O}_Z$
is surjective; denote its kernel by
$\mathcal{I}$.
To make sure that if $(X,\mathcal{O}_X)$ is a
scheme
then so is $(Z,\mathcal{O}_Z)$ with add the
extra condition
that the $\mathcal{O}_X$-module $\mathcal{I}$
is
locally generated by sections.
\end{definition}

\medskip\noindent
Now let's discuss restricting sheaves to
subschemes.

Upshot about restricting sheaves to
subschemes:
the sheaves on $Z$ correspond essentially
(in the categorical sense)
to sheaves on $X$ with support on $Z$.

More explicitly, from
\href{https://stacks.math.columbia.edu/tag/%
01QY}{01QY}:
the pushforward of the inclusion (a closed
immersion of schemes)
$i:Z \to X$, with kernel on induced map
$\mathcal{O}_Z \to \mathcal{O}_X$
denoted by $I$, is an exact, fully faithful
functor
$$
i_*:\QCoh(\mathcal{O}_Z) \to
\QCoh(\mathcal{O}_X)
$$
with essential image the quasi-coherent
$\mathcal{O}_X$ modules
$\mathcal{G}$ such that
$\mathcal{I}\mathcal{G}=0$.

It makes sense to think of this operation as
$\mathcal{F}|_Z=i^*\mathcal{F}=\mathcal{F}
\otimes\mathcal{O}_Z$
for $\mathcal{F}\in\QCoh(X)$.
Indeed: tensoring by the ideal sheaf would
kill the information on $Z$,
and tensoring with the structure sheaf does
the opposite
--- it's the sheaves on $X$ that are no fun
outside $Z$,
so basically sheaves on $Z$.

\medskip\noindent
Here's a more detailed discussion:

In understanding the cohomology of sheaves
when
we restrict them to submanifolds (see Complex
Geometry
Exercise
\ref{complex-geometry-exercise-%
deformations-curve-in-k3},
I found Stacks Project
\href{https://stacks.math.columbia.edu/tag/%
01AW}{01AW}
section on closed immersions and abelian
sheaves.
First we have a few propositions stating that
sheaves on the subscheme $Z \subset X$ can be
pushed
to sheaves on $X$.

And then there's the converse:
what I would like to call the restriction
functor $\mathcal{H}$
but rather is called the {\it sub-module of
sections with support in $Z$}.
In Remark
\href{https://stacks.math.columbia.edu/tag/%
01AY}{01AY}
(and then in Remark
\href{https://stacks.math.columbia.edu/tag/%
0G6N}{0G6N}
for ringed spaces; these seem to be analogue
constructions…)
we see that for an abelian sheaf
$\mathcal{F}$ on $X$
we can define an abelian sheaf
$\mathcal{H}_Z(\mathcal{F})$
which although is a sheaf on $X$
we can just think it's a sheaf on $Z$ -
that's beacause
$$
\mathcal{H}_Z(\mathcal{F})(U)
=\{s \in \mathcal{F}(U))| \text{the support
of $s$ is contained in }Z\cap U\}.
$$
We can just think that it's the functions
that vanish outside the subvariety.
Which doesn't really make sense:
a function vanishing outside the submanifold
will most likely
vanish also in the submanifold! Just by
continuity,
the submanifold is in the closure of the open
set!
I guess what I mean to say is that I don't
mind if
the functions do not vanish outside $Z$…
I just want to consider them as functions on
$Z$…

\medskip\noindent
Here's another attempt in understanding how
to restrict sheaves to submanifolds.
First and foremost, lacking a reference this,
I understand the symbol $\mathcal{F}|_Z$ for
a sheaf
on $X$ and a subscheme  $Z \subset X$ to mean
the
pullback of $\mathcal{F}$ by the inclusion.
Then look at
\href{https://stacks.math.columbia.edu/tag/%
01QY}{01QY}
(more basically, as above,
\href{https://stacks.math.columbia.edu/tag/%
01AW}{01AW}…
in fact there appear to be several version of
this statement)
to recall that sheaves of the subscheme $Z$
should correspond
to sheaves on $X$ with support on $Z$.
{\bf So essentially $\mathcal{F}|_Z$ should
be a version of $\mathcal{F}$
but with support on $Z$.}
The key point here is that this is non other
than
$\mathcal{F} \otimes \mathcal{O}_Z$.
Why? Because when multiply with the ideal
sheaf of $Z$ we obtain
$(\mathcal{F} \otimes
\mathcal{O}_Z)\mathcal{I}=0$
since sections of $\mathcal{I}$
vanish in
$\mathcal{O}_Z=\mathcal{O}_X/\mathcal{I}$.
That is, the support of $\mathcal{F} \otimes
\mathcal{O}_Z$ indeed
contained in $Z$: if the sections become zero
when
multiplied by sections of $\mathcal{I}$
(sections of $\mathcal{I}$ are only nonzero
outside $Z$)
it means that all the action was inside $Z$.
that outside $Z$

I think the upshot here is that structure
sheaves
act as some sort of ``contour'' for the
geometric object:
they vanish {\it outside} the geometric
object,
they tell us what is the shape of the object
by describing its surroundings, rather than
its interior.
This is exactly how algebraic/analytic sets
are defined,
and this why, when we want to restrict a
sheaf
to a subscheme, we tensor with the structure
sheaf:
tensor by the ideal sheaf would kill the
information on $Z$,
and tensoring with the structure sheaf does
the opposite.














\medskip\noindent
Now let's dicuss Proj of a graded ring.

\noindent
(Stacks Project
\href{https://stacks.math.columbia.edu/tag/%
01M3}{01M3},
Proj of a graded ring.)
This section may be taken as the preamble to
the question: what is projective space
anyway?
Yes, projective space is Proj of a graded
ring $S$.
And yes, projective space is a scheme.
So in particular it's locally affine.

(Stacks Project
\href{https://stacks.math.columbia.edu/tag/%
01MX}{01MX},
Functoriality of Proj.)
This is reminiscent of the olden days.
The point is that Proj is covariant,
i.e. for rings  $S' \subset S$ we have
$\text{Proj}S \subset \text{Proj}S$
(under the right conditions,
which is the point of this section),
meaning that, indeed,
{\bf Projs are projective schemes}.






\medskip\noindent
Noe let's discuss invertible sheaves on Proj
(twists!).

\noindent
As for the construction of the twisted
sheaves:
for a projective space $\text{Proj} S$,
we just compute the twisted $S$ module
$S(m)$ by the formula $S(m)_d=S_{m+d}$
and turn that into a sheaf
$\widetilde{S(m)}$.

Here's the construction of the twisted
sheaves $\mathcal{O}_X(n)$.
The point is that there is a good way
to pass from an $S$-module $M$ to an
$\mathcal{O}_X$-module, called
$\widetilde{M}$.
So basically you just define the twists
by $M(d)_n=M_{n+d}$ and apply that
construction.

Pick an element of degree $d$.
Then think of the module generated by this
element.
The least degree piece of such a module
is the degree $d$ piece of the original
module,
that is we have $M(-d)$.











\medskip\noindent
Lost note on logarithmic derivations:

\noindent
Just to record a notion of logarithmic
derivation
from Vladimiro Benedetti, Danielle Faenzi and
Simone.

Let $F$ be a homogneous polynomial of degree
$d$.
The {\it logarithmic derivations} of $F$ are
$$
\text{Der}_U(-\text{log}F)
=\{\partial \in \text{Der}_U
:\partial F \in \left\langle{}F\right\rangle{}_U\}
$$
where $U=\mathbb{C}[x_1,…,x_e]$.










\medskip\noindent
Lost notes on gonality.

\noindent
From a talk by Juliana Coelho at Bandoleiros
2025:

A curve $C$ is {\it $k$-gonal} if there is a
map $C \to \mathbb{P}^1$ of degree $k$.

A curve is {\it stable} if it is nodal and
$\text{Aut}(C)$ is finite.
Equivalently, the rational components of $C$
meet the
rest of the curve in at least three points.

We may study the moduli of $k$-stable curves,
or the space of $k$-gonal structures.

\section{Coherent and reflexive sheaves}
\label{algebraic-geometry-section-coherent-reflexive-sheaves}

Unless stated otherwise, $X$ is a smooth
integral variety over $\mathbb{C}$ and
the sheaves are coherent.

\medskip\noindent
\textbf{Coherent, locally free, bundle.}
A coherent sheaf is locally described by
finitely many generators and relations.
On an affine open $\operatorname{Spec}R$
it corresponds to a finite $R$-module.
It is \emph{locally free of rank $r$} if
each point has a neighbourhood $U$ with
$$
\mathcal{F}|_U\cong\mathcal{O}_U^{\oplus r}.
$$
These are exactly the sheaves of sections
of algebraic vector bundles of rank $r$.
A line bundle has rank one; its sheaf is
also called \emph{invertible}.
The rank of a coherent sheaf is its
dimension over the function field at the
generic point. Rank $r$ alone does not
imply local freeness.

\medskip\noindent
\textbf{Torsion is about multiplication.}
Over a domain $R$, an element $m$ of an
$R$-module $M$ is torsion if $fm=0$ for
some nonzero $f\in R$.
The module is \emph{torsion-free} if
only $m=0$ can satisfy this condition.
It is a \emph{torsion module} if every
element is torsion. A sheaf has these
properties when its stalks do.

For example, take
$$
R=\mathbb{C}[x],\qquad M=R/(x).
$$
Then $x\overline{1}=0$ although
$x\ne0$ in $R$ and $\overline{1}\ne0$
in $M$. All elements of $M$ are scalar
multiples of $\overline{1}$; this is not
the only nonzero element.
Moreover $M\cong\mathbb{C}$ as a ring,
so there are no nonzero nilpotents.
Torsion concerns the action of $R$ on
$M$; nilpotence concerns powers inside
a ring. They are different notions.

For a proper closed subscheme $Z$ of $X$,
$\mathcal{O}_Z$, viewed on $X$, is torsion.
In contrast $\mathcal{I}_Z$ is
torsion-free, since it is a subsheaf of
$\mathcal{O}_X$. The sequence is
$$
0\longrightarrow\mathcal{I}_Z
\longrightarrow\mathcal{O}_X
\longrightarrow\mathcal{O}_Z
\longrightarrow0.
$$
The ideal sheaf of a curve in
$\mathbb{P}^3$ need not be locally free.

\medskip\noindent
\textbf{Dual, evaluation, reflexive.}
For modules, put
$$
M^*=\operatorname{Hom}_R(M,R),\qquad
M^{**}=\operatorname{Hom}_R(M^*,R).
$$
An element $m$ defines a functional on
$M^*$ by evaluating functionals at $m$:
$$
\begin{aligned}
\operatorname{ev}:M&\longrightarrow M^{**},\\
\operatorname{ev}(m)(\varphi)&=\varphi(m).
\end{aligned}
$$
This construction requires no basis.
It is natural: for $u:M\to N$,
$$
u^{**}\operatorname{ev}_M
=\operatorname{ev}_N u.
$$
Thus the familiar finite-dimensional
isomorphism $V\cong V^{**}$ comes with
a distinguished map. Reflexivity asks
whether this particular map remains an
isomorphism for modules and sheaves.

\begin{definition}
\label{algebraic-geometry-definition-reflexive}
Let $X$ be an integral locally Noetherian
scheme. Let $\mathcal{F}$
be a coherent $\mathcal{O}_X$-module.
The {\it reflexive hull}
of $\mathcal{F}$ is the
$\mathcal{O}_X$-module
$$
\mathcal{F}^{**}
=\mathcal{H}om_{\mathcal{O}_X}(
\mathcal{H}om_{\mathcal{O}_X}(\mathcal{F},
\mathcal{O}_X), \mathcal{O}_X)
$$
We say $\mathcal{F}$ is {\it reflexive}
if the natural map
$j : \mathcal{F} \longrightarrow
\mathcal{F}^{**}$
is an isomorphism.
\end{definition}

\begin{lemma}
\label{algebraic-geometry-lemma-reflexive-torsion-free}
Let $X$ be an integral locally Noetherian
scheme. Let $\mathcal{F}$
be a coherent $\mathcal{O}_X$-module.
\begin{enumerate}
\item If $\mathcal{F}$ is reflexive, then
$\mathcal{F}$
is torsion free.
\item The map $j : \mathcal{F}
\longrightarrow \mathcal{F}^{**}$
is injective if and only if $\mathcal{F}$
is torsion free.
\end{enumerate}
\end{lemma}

The evaluation map above is the natural
map in the definition. Torsion-freeness
requires injectivity; reflexivity also
requires surjectivity.

Here is the proof that reflexive implies
torsion-free. If $fm=0$ with $f\ne0$,
then for every $\varphi\in M^*$,
$$
f\varphi(m)=\varphi(fm)=0.
$$
Since $R$ is a domain, $\varphi(m)=0$.
Thus $\operatorname{ev}(m)=0$, so the
injectivity of evaluation gives $m=0$.
Locally free sheaves are reflexive,
because finite free modules equal their
double duals by evaluation.

The converses fail in general.
For a line $L\subset\mathbb{P}^3$,
$$
\mathcal{I}_L^{**}\cong\mathcal{O},
\qquad \mathcal{I}_L\ne\mathcal{O}.
$$
Thus $\mathcal{I}_L$ is torsion-free
but not reflexive.
On a smooth surface, reflexive sheaves
are locally free. On a smooth threefold,
they are locally free outside a subset
of codimension at least three, so on
$\mathbb{P}^3$ only finitely many points
can fail local freeness.
See the Stacks Project,
\href{%
https://stacks.math.columbia.edu/tag/0AV0%
}{Tag 0AV0} and
\href{%
https://stacks.math.columbia.edu/tag/0AVT%
}{Tag 0AVT}.

\medskip\noindent
\textbf{Twists and sections.}
On projective space, the parentheses
denote tensoring with a line bundle:
$$
\begin{aligned}
\mathcal{F}(a)&=
\mathcal{F}\otimes\mathcal{O}(a),\\
\mathcal{F}(a)(b)&=\mathcal{F}(a+b),\\
\mathcal{O}(a)^*&=\mathcal{O}(-a).
\end{aligned}
$$
Locally free sheaves are flat: on each
trivializing open set, tensoring makes
copies of the original exact sequence.
See Commutative Algebra, Lemma
\ref{commutative-algebra-%
lemma-free-tensor-copies}.
Thus tensoring with $\mathcal{O}(a)$ is
exact and does not change rank. Moreover
$$
\operatorname{Hom}(\mathcal{O}(-a),E)
\cong H^0(E(a)).
$$
This is why a section of $E(a)$ becomes
an arrow $\mathcal{O}(-a)\to E$.
The image is \emph{saturated} if the
quotient is torsion-free. Injectivity
alone does not guarantee this.

For example, multiplication by the
equation of a hyperplane $H$ gives an
injection $\mathcal{O}(-1)\to\mathcal{O}$.
It is not saturated: the quotient
$\mathcal{O}_H$ is torsion on the
ambient projective space.

\begin{lemma}
\label{algebraic-geometry-lemma-line-to-torsion-free}
On an integral variety, a map from a
line bundle to a torsion-free coherent
sheaf is either zero or injective.
\end{lemma}

\begin{proof}
A nonzero map is nonzero at the generic
point, since its target has no torsion.
There it is a nonzero map out of a
one-dimensional vector space, hence
injective. Its kernel would therefore
be torsion, but a line bundle has none.

Locally, the argument is just this:
a map $R\to M$ sends $1$ to some $m$.
If a nonzero $f$ is in its kernel, then
$fm=0$. Torsion-freeness forces $m=0$,
so the map is zero.
\end{proof}

\section{Chern classes of coherent sheaves}
\label{algebraic-geometry-section-chern-coherent-sheaves}

Over $\mathbb{C}$, these are the same
characteristic classes encountered in
differential geometry:
$$
c_i(E)\in H^{2i}(X,\mathbb{Z}).
$$
Algebraically they also have cycle-class
versions. On $\mathbb{P}^3$, put
$H=c_1(\mathcal{O}(1))$ and write
$$
c_1(E)=cH,\qquad c_2(E)=eH^2.
$$
We often identify a Chern class with
its integer coefficient in this basis.
For rank two, twisting gives
$$
\begin{aligned}
c_1(E(a))&=(c+2a)H,\\
c_2(E(a))&=(e+ac+a^2)H^2.
\end{aligned}
$$
One can remember this by adding $aH$
to each of the two formal Chern roots.
If a section of $E(a)$ has zero scheme
$C$ of pure codimension two, then
$$
[C]=c_2(E(a)),\qquad
\deg C=e+ac+a^2.
$$
Thus the second Chern class computes
the degree of the zero curve.
For the section construction, see the
Lefschetz workshop notes, Section
\ref{iv-jornadas-lefschetz-2026-%
section-jardim-problems}.

Chern classes of a coherent sheaf on
smooth projective space are defined
using a finite resolution by bundles.
A rank-two vector bundle has $c_3=0$.
A rank-two reflexive sheaf need not:
its higher Chern classes can retain
information about its singularities.

More explicitly, a locally free resolution
$V_\bullet\to F$ gives
$$
c(F)=\prod_i c(V_i)^{(-1)^i}.
$$
For an exact sequence
$0\to A\to B\to C\to0$, this yields
$$
c(B)=c(A)c(C).
$$
On $\mathbb{P}^3$ the calculation takes
place in $\mathbb{Z}[H]/(H^4)$.
This makes resolutions a practical way
to compare possible middle sheaves;
see the line--conic calculation in the
Lefschetz workshop notes, Section
\ref{iv-jornadas-lefschetz-2026-%
section-jardim-problems}.

\section{Sheaf Ext and duality}
\label{algebraic-geometry-section-sheaf-ext-duality}

Global $\operatorname{Ext}^1(B,A)$
classifies extensions with fixed end
terms. Zero is the split extension;
a nonzero class need not have a locally
free middle sheaf. See Homological
Algebra, Section
\ref{homological-algebra-section-extensions}
and Lemma
\ref{homological-algebra-%
lemma-extension-class-retraction}.

For sheaves, a locally free resolution
computes a sheaf Hom complex, whose
hypercohomology gives global Ext.
Locally free sheaves need not be
projective objects for global Hom.
If $E$ is locally free, then
$$
\operatorname{Ext}^i(E,A)
\cong H^i(X,E^*\otimes A).
$$
For example, on $\mathbb{P}^3$,
$$
\operatorname{Ext}^3
(\mathcal{O},\mathcal{O}(-4))
\cong\mathbb{C},
$$
although all the positive sheaf Exts
of $\mathcal{O}$ vanish.

\medskip\noindent
\textbf{Local freeness and reflexivity.}
The sheaf $\mathcal{E}xt^i(F,G)$ is the
sheafification of the presheaf
$$
U\longmapsto
\operatorname{Ext}^i_{\mathcal{O}_U}
(F|_U,G|_U).
$$
On an affine open $U=\operatorname{Spec}R$,
for coherent sheaves corresponding to
modules $M,N$, it is the sheaf associated
to $\operatorname{Ext}^i_R(M,N)$.
In particular, its stalk is
$$
\mathcal{E}xt^i(F,G)_x
\cong\operatorname{Ext}^i_{\mathcal{O}_{X,x}}
(F_x,G_x).
$$
Thus it records a local obstruction,
rather than one global extension space.

\begin{proposition}
\label{algebraic-geometry-proposition-threefold-sheaf-ext}
Let $X$ be a smooth projective threefold
and $F$ a coherent torsion-free sheaf.
Put
$$
\mathcal{E}^i=
\mathcal{E}xt^i(F,\mathcal{O}_X).
$$
Then $F$ is locally free if and only if
$$
\mathcal{E}^1=\mathcal{E}^2=0.
$$
It is reflexive if and only if
$\mathcal{E}^2=0$ and
$\operatorname{Supp}\mathcal{E}^1$
is a finite set of points, possibly
empty. For a reflexive $F$, this becomes:
$F$ is locally free if and only if
$\mathcal{E}^1=0$.
\end{proposition}

More generally, if a torsion-free $F$
locally has a resolution of length one,
$$
0\longrightarrow V_1\longrightarrow V_0
\longrightarrow F\longrightarrow0
$$
with $V_i$ locally free, the condition
$\mathcal{E}^2=0$ is automatic. In this
setting the two tests are particularly
simple: vanishing of $\mathcal{E}^1$
means locally free, and finite support
of $\mathcal{E}^1$ means reflexive.
This applies to extensions of ideal
sheaves of locally complete intersection
curves by line bundles.

Why does this work? At a point of a
smooth threefold the local ring is
regular. The Auslander--Buchsbaum
formula relates the length of a minimal
free resolution to depth:
$$
\operatorname{pd}F_x+\operatorname{depth}F_x
=\dim\mathcal{O}_{X,x}.
$$
The last nonzero term of a minimal
resolution gives a nonzero local Ext
in that degree. Vanishing of the
positive Ext groups therefore detects
a free stalk. Reflexivity requires
depth at least two wherever the local
ring has dimension at least two.
On a threefold this allows length-one
resolutions at closed points, but
requires freeness in codimension two.
That is exactly the finite-support
condition above, together with
$\mathcal{E}^2=0$.

The hypotheses matter. If $p$ is a
point of $\mathbb{P}^3$, then
$$
\begin{aligned}
\mathcal{E}xt^1
(\mathcal{I}_p,\mathcal{O})&=0,\\
\mathcal{E}xt^2
(\mathcal{I}_p,\mathcal{O})&\cong
\mathcal{O}_p.
\end{aligned}
$$
Yet $\mathcal{I}_p$ is neither locally
free nor reflexive. Thus vanishing of
sheaf Ext in degree one alone does not
test an arbitrary torsion-free sheaf.

The general support criterion is in
Huybrechts--Lehn,
\emph{The Geometry of Moduli Spaces
of Sheaves}, Propositions 1.1.6 and
1.1.10. It says, for a torsion-free
sheaf on a smooth variety, that
reflexivity is equivalent to
$$
\operatorname{codim}_X
\operatorname{Supp}
\mathcal{E}xt^i(F,\mathcal{O}_X)
\geq i+2\qquad(i>0).
$$
On a threefold, a support of codimension
at least four must be empty.

\medskip\noindent
\textbf{Line-bundle cohomology.}
\label{algebraic-geometry-paragraph-projective-line-cohomology}
For every integer $m$,
$$
H^i(\mathbb{P}^n,\mathcal{O}(m))=0
\qquad(0<i<n).
$$
This is not a vanishing theorem for
arbitrary coherent sheaves, nor for top
cohomology. For example,
$$
H^n(\mathbb{P}^n,\mathcal{O}(-n-1))
\cong\mathbb{C}.
$$
On $\mathbb{P}^3$ the intermediate
vanishing of line bundles is what
identifies the $H^1_*$ of a bundle and
its quotient by a line sheaf.

\medskip\noindent
\textbf{Computing extension spaces.}
\label{algebraic-geometry-paragraph-curve-extension-duality}
For coherent $F,G$ on a smooth projective
$n$-fold, Serre duality gives
$$
\operatorname{Ext}^i(F,G)
\cong\operatorname{Ext}^{n-i}
(G,F\otimes\omega_X)^\vee.
$$
Here $\vee$ denotes the vector-space
dual. In particular, on $\mathbb{P}^3$,
$$
\begin{aligned}
&\operatorname{Ext}^1
(\mathcal{I}_C(c+a),\mathcal{O}(-a))\\
&\quad\cong\operatorname{Ext}^2
(\mathcal{O}(-a),\mathcal{I}_C(c+a-4))^\vee
\\
&\quad\cong H^2
(\mathbb{P}^3,\mathcal{I}_C(c+2a-4))^\vee.
\end{aligned}
$$
The $-4$ comes from
$\omega_{\mathbb{P}^3}=\mathcal{O}(-4)$;
the second $a$ comes from dualizing
$\mathcal{O}(-a)$ in the first argument.
Finally the ideal sequence and the
vanishing just recalled give
$$
H^2(\mathbb{P}^3,\mathcal{I}_C(m))
\cong H^1(C,\mathcal{O}_C(m)).
$$
Thus extension spaces can be computed
from cohomology on the curve.

\section{Reduced schemes}
\label{algebraic-geometry-section-reduced}

\noindent
So far a reduced scheme,
for me,
is just a scheme where there are
no ``repetitions'':
$(x^2)$ is not reduced because
it gives the same information as $(x)$.
Algebraically this means that the
coordinate ring is reduced,
which is defined by asking
that there are no nilpotent elements,
i.e. elements such that some power vanishes.
(There's no definition of reduced ring in
Stacks Project
since it's taken as part of the basic
algebraic knowledge).

The definition is on stalks but the first
lemma
shows it's equivalent to define it on any
open set.
This is in virtue of some ``injectivity''
lemma
from the ring at $U$ to the product of all
stalks parametrized in $U$;
essentially that if a section vanishes when
projected to all
the stalks then it's zero.

\begin{definition}
\label{algebraic-geometry-definition-reduced}
Let $X$ be a scheme. We say $X$ is {\it
reduced} if every local ring
$\mathcal{O}_{X, x}$ is reduced.
\end{definition}


\begin{proof}
Assume that $X$ is reduced.
Let $f \in \mathcal{O}_X(U)$ be a section
such that $f^n = 0$.
Then the image of $f$ in $\mathcal{O}_{U, u}$
is zero for
all $u \in U$. Hence $f$ is zero, see
Sheaves, Lemma
\ref{algebraic-geometry-sheaves-lemma-sheaf-subset-stalks}.
Conversely, assume that $\mathcal{O}_X(U)$ is
reduced
for all opens $U$. Pick any nonzero element
$f \in \mathcal{O}_{X, x}$.
Any representative $(U, f \in
\mathcal{O}(U))$  of $f$ is nonzero and
hence not nilpotent. Hence $f$ is not
nilpotent in $\mathcal{O}_{X, x}$.
\end{proof}

\section{Fibres and specialization}
\label{algebraic-geometry-section-fibres-specialization}

\noindent
For flatness of modules, see
Commutative Algebra, Section
\ref{commutative-algebra-%
section-flat-modules}.
Here we study the fibres of a family
and how a point on its generic fibre
can specialize to its special fibre.

\medskip\noindent
Epiphany:
in the category of commutative rings
pushout is tensor product.
So think of $\Spec $ as a functor
from $\mathit{CRing}^{\text{op}}$ to $\Sch$,
then pushout goes to pullback,
and what's an example of a pullback?
Fibre!
So, the coordinate ring of a fiber
is essentially given by the residue field
at the point that parametrizes it!
(tensored with the coordinate ring
of the deformation space).
See Commutative Algebra, Proposition
\ref{commutative-algebra-%
proposition-preimage-of-point}
and Remark
\ref{commutative-algebra-%
remark-preimage-as-fibered-product}.
There we take the reduced fibre of
affine varieties; for scheme-theoretic
fibres we keep the tensor product,
without removing its nilpotents.

\begin{remark}
\label{algebraic-geometry-remark-units-formal-series}
A formal power series
$$
f=a_0+a_1s+a_2s^2+\cdots
\in\mathbb{C}[[s]]
$$
is a unit if and only if $a_0\neq0$.
It need not be constant: its inverse
may itself be an infinite formal series.

If $fg=1$, comparison of constant
terms forces $a_0g(0)=1$, so $a_0\neq0$.
Conversely, write $f=a_0(1+w)$ with
$w\in s\mathbb{C}[[s]]$. Then
$$
f^{-1}=\frac{1}{a_0}
\sum_{n=0}^{\infty}(-w)^n.
$$
Each coefficient is a finite sum, since
$w^n$ is divisible by $s^n$.
This is the geometric-series argument;
no analytic convergence is needed.

See Atiyah--Macdonald,
{\it Introduction to Commutative Algebra},
Chapter 1, Exercise 5(i).
Over a general commutative coefficient
ring, the constant term must be a unit,
rather than merely nonzero.
\end{remark}

\begin{exercise}[The two fibres of
a formal family]
\label{algebraic-geometry-example-two-point-base}
Let
$$
R=\mathbb{C}[[t]],\qquad
K=\mathbb{C}((t)).
$$
\begin{enumerate}
\item Show that $\operatorname{Spec}R$
has exactly two points:
$\eta=(0)$ and $s=(t)$.
\item Show that $\eta$ is generic
and $s$ is closed.
\item Compute their residue fields.
\item For an affine family
$\operatorname{Spec}A\to
\operatorname{Spec}R$,
describe its generic and special
fibres using these residue fields.
\end{enumerate}
The base has two points; a family over
it has a fibre over each.
\end{exercise}

\begin{proof}
\begin{enumerate}
\item By definition, points of
$\operatorname{Spec}R$ are prime ideals
of $R$, not elements of $R$.
We determine all of them.

First, $(0)$ is prime because $R$ is
a domain: the product of two nonzero
power series has a nonzero leading term.

Now let $\mathfrak{p}$ be a nonzero
prime ideal and choose
$0\neq f\in\mathfrak{p}$.
Write
$$
f=a_nt^n+a_{n+1}t^{n+1}+\cdots=t^n u,
\qquad a_n\neq 0.
$$
The series $u$ has nonzero constant
term $a_n$, so it is a unit by Remark
\ref{algebraic-geometry-remark-units-formal-series}.
This is the valuation picture from
Commutative Algebra, Example
\ref{commutative-algebra-%
example-laurent-series-valuation}.

Since a proper ideal cannot contain
a unit, $n>0$ and $u\notin\mathfrak{p}$.
As $t^n u=f\in\mathfrak{p}$ and
$\mathfrak{p}$ is prime, we obtain
$t^n\in\mathfrak{p}$.
Applying primality again gives
$t\in\mathfrak{p}$, so
$(t)\subset\mathfrak{p}$.
But evaluation at zero gives
$$
R/(t)\cong\mathbb{C}.
$$
Hence $(t)$ is maximal, and necessarily
$\mathfrak{p}=(t)$.
It is itself prime because its quotient
is a field. Therefore there are exactly
two prime ideals, hence two points:
$$
\eta=(0),\qquad s=(t).
$$
\item Recall that in the Zariski topology,
$$
\overline{\{\mathfrak{p}\}}
=V(\mathfrak{p})
=\{\mathfrak{q}\in\operatorname{Spec}R:
\mathfrak{p}\subseteq\mathfrak{q}\}.
$$
Indeed, a closed set $V(I)$ contains
$\mathfrak{p}$ precisely when
$I\subseteq\mathfrak{p}$; it then
contains every prime containing
$\mathfrak{p}$. Thus $V(\mathfrak{p})$
is the smallest closed set containing
that point.

Every prime ideal contains $(0)$, so
$$
\overline{\{\eta\}}=\operatorname{Spec}R.
$$
This is why $\eta$ is called generic.
Since $(t)$ is maximal, the only prime
containing it is itself, so
$\overline{\{s\}}=\{s\}$:
the point $s$ is closed.
\item Their residue fields are
$$
\kappa(\eta)=K,\qquad
\kappa(s)=\mathbb{C}.
$$
To check the first equality, recall that
the residue field is the local ring
modulo its maximal ideal:
$$
\kappa(\eta)=R_{(0)}/(0)R_{(0)}.
$$
Localizing at $(0)$ inverts every
element outside $(0)$, hence every
nonzero element of $R$. Therefore
$$
R_{(0)}=\operatorname{Frac}(R)=K.
$$
Its maximal ideal is $(0)$, so
$$
\kappa(\eta)=K/(0)=K.
$$
We quotient the local ring $R_{(0)}$,
not $R$ itself.

\item For an affine family
$$
\mathcal{X}=\operatorname{Spec}A
\longrightarrow\operatorname{Spec}R,
$$
the two fibres are
$$
\mathcal{X}_{\eta}
=\operatorname{Spec}(A\otimes_R K),
\qquad
\mathcal{X}_s
=\operatorname{Spec}(A/tA).
$$
The first makes $t$ invertible; the
second sets $t=0$. The generic fibre
does not mean setting $t$ equal to
a particular nonzero complex number.
For the definition of a fibre, see
Stacks Project, Definition 26.18.4,
\href{https://stacks.math.columbia.edu/tag/%
01K0}{Tag 01K0}.
\end{enumerate}
\end{proof}

\begin{exercise}[A generic-fibre point
specializes by taking its closure]
\label{algebraic-geometry-exercise-follow-point-to-fibre}
With $R$ and $K$ as above, consider
$$
\mathcal{X}
=\operatorname{Spec}(R[x,y]/(xy-t)).
$$
This exercise illustrates specialization
by a direct calculation.
\begin{enumerate}
\item Write down its generic and special
fibres. Check that $P=(t,1)$ is a
$K$-point of the generic fibre.
\item Inside the total space impose
$x=t$ and $y=1$, obtaining a closed
subscheme $Z$. Show that
$$
R[x,y]/(xy-t,x-t,y-1)\cong R.
$$
\item Thus $Z$ has two points. Identify
its generic point with $P$. Find the
coordinates of its closed point by
reducing modulo $t$.
\item Check that this closed point
satisfies the special-fibre equation.
\item Show that $Z$ is the closure of
$P$ in the total space: use that the
generic point of $\operatorname{Spec}R$
is dense and that $Z$ is closed.
\end{enumerate}
The goal is to see how a point that is
closed in the generic fibre can acquire
a specialization in the total space.
No singularity assumption is needed
for this exercise.
\end{exercise}

\begin{proof}
\begin{enumerate}
\item
Put
$$
A=R[x,y]/(xy-t).
$$
We compute fibres by tensoring over
$R$ with the residue field of the
base point, as in Exercise
\ref{algebraic-geometry-example-two-point-base}
and Stacks Project, Definition 26.18.4,
\href{https://stacks.math.columbia.edu/tag/%
01K0}{Tag 01K0}.

For the special point, the residue
field is $\mathbb{C}=R/(t)$. Hence
$$
A\otimes_R\mathbb{C}
\cong A/tA
\cong\mathbb{C}[x,y]/(xy).
$$
Tensoring with $\mathbb{C}$ here does
not leave $A$ unchanged: the tensor
product is over $R$, not over
$\mathbb{C}$. In the tensor product,
$$
t\otimes1=1\otimes(t\cdot1)=0,
$$
because $t$ acts as zero on $R/(t)$.
Thus the relation $xy=t$ becomes
$xy=0$. The special fibre is
$$
\mathcal{X}_s
=\operatorname{Spec}
(\mathbb{C}[x,y]/(xy)).
$$

For the generic point, the residue
field is $K=\mathbb{C}((t))$. Thus
$$
B:=A\otimes_R K
\cong K[x,y]/(xy-t),
$$
and the generic fibre is
$\mathcal{X}_{\eta}=\operatorname{Spec}B$.
Here $t$ becomes invertible, not equal
to $1$. We can nevertheless set
$z=y/t$, obtaining
$$
B\cong K[x,z]/(xz-1).
$$
This is a change of coordinates over
$K$, not the substitution $t=1$.

The notation $P=(t,1)$ specifies
coordinates in the affine plane over
$K$. It does not denote the ideal
$(t,1)$, which contains $1$ and is
therefore the whole ring.
Since $t\cdot1=t$, these coordinates
satisfy the generic-fibre equation.
Evaluation gives a surjective map
$$
B\longrightarrow K,\qquad
f\longmapsto f(t,1).
$$
Its kernel is
$$
\mathfrak{m}_P=(x-t,y-1).
$$
Indeed, quotienting by these relations
replaces $x$ by $t$ and $y$ by $1$,
leaving exactly $K$. Consequently
$$
B/\mathfrak{m}_P\cong K.
$$

To connect this with the definition of
the residue field, localizing at $P$
inverts every $g\notin\mathfrak{m}_P$,
that is, every $g$ with $g(t,1)\neq0$.
Evaluation therefore extends to
$$
B_{\mathfrak{m}_P}\longrightarrow K,
\qquad
\frac{f}{g}\longmapsto
\frac{f(t,1)}{g(t,1)}.
$$
This map is surjective, since constants
already give all of $K$.
Its value is zero exactly when
$f(t,1)=0$, or $f\in\mathfrak{m}_P$.
Its kernel is therefore
$\mathfrak{m}_P B_{\mathfrak{m}_P}$.
Thus
$$
\begin{aligned}
\kappa(P)
&=B_{\mathfrak{m}_P}/
  \mathfrak{m}_P B_{\mathfrak{m}_P}\\
&\cong B/\mathfrak{m}_P
\cong K.
\end{aligned}
$$
Localizing can change the ring, but
does not change this quotient:
the denominators already become
nonzero, hence invertible, elements
of the field $K$.
\end{enumerate}
\end{proof}



\begin{exercise}[A unit cannot be
cancelled by higher-order terms]
\label{algebraic-geometry-exercise-dvr-orders-contradiction}
Work in
$$
V=\mathbb{C}[[s]],\qquad
\mathfrak{m}=(s).
$$
For a nonzero series, let $\nu$ be its
smallest exponent, as in
Commutative Algebra, Example
\ref{commutative-algebra-%
example-laurent-series-valuation}.
\begin{enumerate}
\item For
$$
x=s^3(2+s),\qquad y=5+s^2,
$$
find their reductions modulo $(s)$,
their valuations, and which is a unit.
\item Set $t=s^2$. Express $x^2$ as
a power of $t$ times a unit of $V$.
Compare $\nu(x^2)$ with $\nu(t^2)$.
\item Suppose $u,r\in V$, with $u$
a unit. Show that
$$
0=t^2u+t^3r
$$
is impossible.
Hint: divide by $t^2$, then reduce
modulo $(s)$.
\end{enumerate}
The last item illustrates how orders
of vanishing can give a contradiction:
a higher-order term cannot cancel
a lower-order term with unit coefficient.
\end{exercise}

\begin{proof}
\begin{enumerate}
\item Reduction modulo $(s)$ keeps
only the constant term, so
$$
\overline{x}=0,\qquad \overline{y}=5.
$$
The valuation is the smallest exponent
with nonzero coefficient. Since
$$
x=2s^3+s^4,\qquad y=5s^0+s^2,
$$
we obtain
$$
\nu(x)=3,\qquad \nu(y)=0.
$$
The constant term counts as exponent
zero.

By Remark
\ref{algebraic-geometry-remark-units-formal-series},
$y$ is a unit in $V$ because its
constant term is $5\neq0$.
In contrast, $x$ is not a unit because
its constant term is zero.
The reduction $\overline{y}=5$ is also
a unit in the residue field, whereas
$\overline{x}=0$ is not.
\item Since $t=s^2$, we have
$$
x^2=s^6(2+s)^2=t^3(4+4s+t).
$$
The factor $4+4s+t=4+4s+s^2$ has
nonzero constant term $4$, so it is
a unit in $V$ by Remark
\ref{algebraic-geometry-remark-units-formal-series}.
Moreover,
$$
\nu(t^2)=\nu(s^4)=4,\qquad
\nu(x^2)=6.
$$
Thus $\nu(x^2)>\nu(t^2)$.
\item Suppose
$$
0=t^2u+t^3r=t^2(u+tr).
$$
Since $V$ is a domain and $t^2\neq0$,
cancelling $t^2$ gives $u=-tr$.
As $t=s^2$, reduction modulo $(s)$
then gives
$$
\overline{u}=0.
$$
But $u$ is a unit, so reducing
$uu^{-1}=1$ gives
$$
\overline{u}\,\overline{u^{-1}}=1
$$
in $V/(s)=\mathbb{C}$.
Thus $\overline{u}$ cannot be zero.
This is the contradiction.
\end{enumerate}
\end{proof}


\begin{exercise}[Singularity forces
an impossible reduction]
\label{algebraic-geometry-exercise-singular-point-dvr-bridge}
This model exercise combines
specialization with the DVR
contradiction. It is not a calculation
from the Grunbaum draft.

Let $R=\mathbb{C}[[t]]$ and fix any
polynomial $Q\in\mathbb{C}[X,Y,Z]$.
Consider the family
$$
\mathcal{Y}=\operatorname{Spec}
(R[X,Y,Z]/(F)),
$$
where
$$
F=XY-t^2Z+t^3Q(X,Y,Z).
$$
Let $V=\mathbb{C}[[s]]$,
$L=\mathbb{C}((s))$, and $t=s^e$
for an integer $e\geq1$.

Suppose a point on the generic fibre
after extending scalars to $L$ has
coordinates
$$
P=(a(s),b(s),c(s)),\qquad a,b,c\in V,
$$
with
$$
\overline{a}=\overline{b}=0,\qquad
\overline{c}\neq0.
$$
Here $X,Y,Z$ are variables in the
equations; $a,b,c$ are the series
obtained by evaluating them at $P$.
Assume this extension over $V$ has
already been supplied. The goal is
to show that $P$ cannot be singular.

\begin{enumerate}
\item Explain why these coordinates
give a map
$\operatorname{Spec}V\to\mathcal{Y}$
over $R$. Identify the image $P_0$
of the closed point by reducing the
coordinates. Check that $P_0$ lies
on the special fibre $XY=0$.
Which of $a,b,c$ are units in $V$?

\item Suppose $P$ is singular.
The hypersurface Jacobian criterion
requires $F_X(P)=F_Y(P)=F_Z(P)=0$.
We will only need the first two.
Differentiate in $X,Y$, holding $t$
fixed, and substitute $a,b,c$.
Explain why the resulting equalities
are equalities in $V$, even though
singularity is tested over $L$.

\item Use those two equalities to show
that $a,b\in t^3V$. Write
$$
a=t^3\alpha,\qquad b=t^3\beta,
\qquad \alpha,\beta\in V.
$$
This is stronger than merely knowing
that $a,b$ reduce to zero.

\item Substitute into $F(P)=0$,
cancel $t^2$, and reduce modulo $(s)$.
Show that the resulting equality
contradicts the fact that $c$ is
a unit. You should not need to know
the coefficients of $Q$.

\item Identify the role of each
assumption: having a specialization,
knowing its nonzero coordinate,
lying on the family, and being
singular. Which step supplied the
extra divisibility that made the
contradiction possible?
\end{enumerate}
\end{exercise}

\section{Hilbert polynomial}
\label{algebraic-geometry-section-Hilbert-polynomial}

\noindent
The following lemma will be made obsolete by
the more general
Lemma
\ref{algebraic-geometry-lemma-numerical-polynomial-from-euler}.

\begin{lemma}
\label{algebraic-geometry-lemma-hilbert-polynomial}
Let $k$ be a field. Let $n \geq 0$. Let
$\mathcal{F}$ be a coherent sheaf
on $\mathbf{P}^n_k$. The function
$$
d \longmapsto \chi(\mathbf{P}^n_k,
\mathcal{F}(d))
$$
is a polynomial.
\end{lemma}

\begin{proof}
We prove this by induction on $n$. If $n =
0$, then
$\mathbf{P}^n_k = \Spec(k)$ and
$\mathcal{F}(d) = \mathcal{F}$.
Hence in this case the function is constant,
i.e., a polynomial
of degree $0$. Assume $n > 0$. By
Lemma
\ref{algebraic-geometry-lemma-euler-chara-exten-base-field}
we may assume $k$ is infinite. Apply
Lemma \ref{algebraic-geometry-lemma-exact-sequence-induction}.
Applying Lemma
\ref{algebraic-geometry-lemma-euler-characteristic-additive}
to the twisted sequences
$0 \to \mathcal{F}(d - 1) \to \mathcal{F}(d)
\to i_*\mathcal{G}(d) \to 0$
we obtain
$$
\chi(\mathbf{P}^n_k, \mathcal{F}(d)) -
\chi(\mathbf{P}^n_k, \mathcal{F}(d - 1)) =
\chi(H, \mathcal{G}(d))
$$
See Remark
\ref{algebraic-geometry-remark-exact-seque-induc-cohom}.
Since $H \cong \mathbf{P}^{n - 1}_k$
by induction the right hand side is a
polynomial.
The lemma is finished by noting that any
function
$f : \mathbf{Z} \to \mathbf{Z}$ with the
property that the map
$d \mapsto f(d) - f(d - 1)$ is a polynomial,
is itself a polynomial.
We omit the proof of this fact (hint: compare
with
Algebra, Lemma
\ref{algebra-lemma-numerical-polynomial}).
\end{proof}

\begin{definition}
\label{algebraic-geometry-definition-hilbert-polynomial}
Let $k$ be a field. Let $n \geq 0$. Let
$\mathcal{F}$ be a coherent sheaf
on $\mathbf{P}^n_k$. The function
$d \mapsto \chi(\mathbf{P}^n_k,
\mathcal{F}(d))$ is called the
{\it Hilbert polynomial} of $\mathcal{F}$.
\end{definition}

\noindent
The Hilbert polynomial has coefficients in
$\mathbf{Q}$ and not
in general in $\mathbf{Z}$. For example the
Hilbert polynomial
of $\mathcal{O}_{\mathbf{P}^n_k}$ is
$$
d \longmapsto {d + n \choose n} =
\frac{d^n}{n!} + \ldots
$$
This follows from the following lemma and the
fact that
$$
H^0(\mathbf{P}^n_k,
\mathcal{O}_{\mathbf{P}^n_k}(d)) = k[T_0,
\ldots, T_n]_d
$$
(degree $d$ part) whose dimension over $k$ is
${d + n \choose n}$.

\begin{lemma}
\label{algebraic-geometry-lemma-hilbert-polynomial-H0}
Let $k$ be a field. Let $n \geq 0$. Let
$\mathcal{F}$ be a coherent sheaf
on $\mathbf{P}^n_k$ with Hilbert polynomial
$P \in \mathbf{Q}[t]$.
Then
$$
P(d) = \dim_k H^0(\mathbf{P}^n_k,
\mathcal{F}(d))
$$
for all $d \gg 0$.
\end{lemma}

\begin{proof}
This follows from the vanishing of cohomology
of high enough twists
of $\mathcal{F}$. See
Cohomology of Schemes,
Lemma
\ref{algebraic-geometry-coherent-lemma-coherent-projective}.
\end{proof}


\medskip\noindent
For completeness I include earlier notes
from \cite{har} on the matter.

The fact that $M$ is finitely generated is
what makes the following two
definitions make sense.

\begin{definition}
\label{algebraic-geometry-definition-Hilbert-function}
The {\it Hilbert function} of a finitely
generated graded $S=k[x_0,\ldots,x_r]$
-module $M$ is
$$
H_M(d)=\dim_kM_d
$$
\end{definition}

\begin{definition}
\label{algebraic-geometry-definition-sysygy}
Define $F_0$ to be the free $S$-module on the
generators of $M$. Elements in the
 kernel $M_1$ of the inclusion are called
 {\it sysygies}. By Hilbert's basis
theorem, $M_1$ is also finitely generated, so
we may choose a set of generators
and repeat this process.
\end{definition}

\begin{theorem}[Hilbert Syzygy Theorem]
\label{algebraic-geometry-theorem-Hilbert-syzygy}
\begin{reference}
\cite[Theorem 1.1]{sys}
\end{reference}
Any finitely generated $S$-module $M$ has a
finite graded free resolution
$$
\xymatrix{
0\ar[r]&F_m\ar[r]^{\varphi_m}&\ar[r]&F_{m-q}
\ar[r]&\cdots\ar[r]&
F_1\ar[r]^{\varphi_1}&F_0
}
$$
Moreover, we may take $m\leq r+1$, the number
of variables in $S$.
\end{theorem}

\begin{lemma}
\label{algebraic-geometry-lemma-Hilbert-function}
Suppose that $S=k[x_0,\ldots,x_r]$ is a
polynomial ring. If the graded
$S$-module $M$ has finite free resolution
$$
\xymatrix{
0\ar[r]&F_m\ar[r]^{\varphi_m}&F_{m-1}
\ar\cdots[r]&F_1\ar[r]^{\varphi_1}&
\ar[r]&F_0
}
$$
with each $F_i$ a finitely generated free
module,
$F_i=\bigoplus_{j}S(-a_{i,j})$, then
\begin{equation}
\label{algebraic-geometry-equation-Hilbert-function}
H_M(d)=\sum_{i=0}(-1)^i\sum_{j}\binom{r+d-
a_{i,j}}{r}
\end{equation}
\end{lemma}

\begin{lemma}
\label{algebraic-geometry-lemma-hilbe-funct-becom-polyn}
There is a polynomial $P_M(d)$ called the
{\it Hilbert polynomial} such that, if
$M$ has free resolution as above, then
$P_M(d)=H_M(d)$ for
$d\geq\text{max}_{i,j}\{a_{i,j}-r\}$.
\end{lemma}

\begin{proof}
When $d$ satisfies this bound then the
binomial coefficients in Eq.
\ref{algebraic-geometry-equation-Hilbert-function} are
polynomials of degree $r$ in $d$.
\end{proof}

\begin{theorem}[Hilbert-Serre]
\label{algebraic-geometry-theorem-Hilbert-Serre}
\begin{reference}
\cite[I, Theorem 7.5]{hart}
\end{reference}
Let $M$ be a finitely generated graded
$S=k[x_0,\ldots,x_n]$. Then there exists
a unique polynomial $p_M$ such that
$p_M(\ell)=\dim S_\ell$ for large enough
$\ell$.
\end{theorem}

\begin{definition}
\label{algebraic-geometry-definition-Hilbert-polynomial}
\begin{reference}
\cite[I, p. 52]{har}
\end{reference}
The polynomial $P_M$ of Hilbert-Serre Theorem
\cite{Hilbert-Serre} is the {\it
Hilbert polynomial} of the finitely generated
$k[x_0,\ldots,x_n]$-module $M$.
\end{definition}

\begin{definition}
\label{algebraic-geometry-definition-degre-proje-varie}
\begin{reference}
\cite[p. 52]{hart}
\end{reference}
If $Y\subset \mathbb{P}^n$ is an algebraic
set of dimension $r$, we define the
{\it degree of $Y$} to be $r!$ times the
leading coefficient of the Hilbert
polynomial of the homogeneous coordinate ring
$S(Y)$.
\end{definition}

\begin{exercise}
\label{algebraic-geometry-exercise-very-ample-bundl-self-inte}
\begin{reference}
\cite[V, Exercise 1.2]{har}
\end{reference}
Let $H$ be a very ample divisor on the
surface $X$, corresponding to a
projective embedding
$X\subseteq\mathbb{P}^N$. If we write the
Hilbert
polynomial of $X$ as
$P(z)=\frac{1}{2}az^2+bz+c$, show that
$a=H^2$,
$b=\frac{1}{2}H^2+1-\pi$, where $\pi$ is the
genus of a nonsingular curve
representing $H$, and $c=1+p_a$.
\end{exercise}

\section{Nakai-Moishezon Criterion}
\label{algebraic-geometry-section-Nakai-Moishezon-criterion}

\begin{theorem}[Nakai-Moishezon Criterion]
\label{algebraic-geometry-theorem-Nakai-Moishezon-criterion}
\begin{reference}
\cite[V, Theorem 1.10]{hart}
\end{reference}
A divisor $D$ on the surface $X$ is ample if
and only if $D^2>0$ and $D.C>0$ for
all irreducible curves $C$ in $X$.
\end{theorem}

\begin{proof}
The direct implication is easy: since $D$ is
ample,  $mD$ is very ample for some
$m$, so that $m^2D^2$ is the
self-intersection number of $mD$. By exercise
\ref{algebraic-geometry-exercise-very-ample-bundl-self-inte},
 $D^2$ is the leading coefficient of the
 Hilbert polynomial of $X$ as a
subscheme of $\mathbb{P}^n$. This means that
$D^2$ is twice the leading
coefficient of the Hilbert polynomial of a
projective variety for large enough
$m$, so that it must be a positive number
(it's the dimension of one of the
graded components of the coordinate ring of
the surface).
\end{proof}

\section{Deformation theory}
\label{algebraic-geometry-section-deformation-theory}

\noindent
\cite{Sernesi-Overview} has a great way of
motivating deformation
theory via the moduli problem.
Here let's just put the following important
meaning for
``deformation theory'' (as opposed to
studying flat families!):

\begin{quote}
``{\it Deformation theory} is the study of
infinitesimal deformaitons
as a tool to understand the local structure
of the moduli space.
The goal is to be able to describe the
restriction
of the universal family to a small
neighborhood of $m \in \mathcal{M}$,
or, more precisely, its restriction to the
germ of $M$ at $m$.

What is interesting here is that we can study
order and infinitesimal
deformations even though the functor $F$ is
not representable or simply we don’t
yet know it is. This is the most frequent
case. Such an investigation will
reveal the infinitesimal properties at $[X]$
of a yet unknown global structure
on $M$ which will be hopefully understood at
a subsequent stage of the
investigation.  In order words it turns out
to be possible and convenient to
separate the {\it global moduli problem} from
the {\it local moduli problem},
and deformation theory studies the latter,
with the purpose of constructing a
family of deformations of a given object
parametrized by the spectrum of a local
ring, and having properties as close as
possible to a universal property.''
\end{quote}




\noindent
I start by reading Stacks Project.

\medskip\noindent
The first notion is {\it thickening of ringed
spaces},
which I ultimately think of as a closed
subscheme
$(X,\mathcal{O}_X) \to (X',\mathcal{O}_{X'})$
with a nilpotent ideal sheaf,
which in an imprecise way means
that $\mathcal{I}_X^n=0$ for some $n$,
and in a precise way its given on sections.

\medskip\noindent
The following definition is from
[lucas-defos], which in turn comes from
\cite{Sernesi-deformations}

\begin{definition}
\label{algebraic-geometry-definition-deformation}
Let $X$ be an algebraic $\mathbb{C}$-scheme.
\begin{enumerate}
\item A {\it deformation} of $X$ is a
Cartesian diagram $\xi$
$$
\xymatrix{
X\ar[r]\ar[d]&\mathcal{X}\ar[d]^{\pi}\\
\Spec \mathbb{C}\ar[r]^s&S
}
$$
where $\pi$ is a flat surjective morphism of
algebraic $\mathbb{C}$-schemes and
$S$ is connected. (Recall that flatness
accounts for ``continuity''.)
\item A {\it local deformation} of $X$ is a
deformation $\xi$ where
$S=\Spec A$ for $A$ a noetherian local
$\mathbb{C}$ algebra with residue
field $\mathbb{C}$.
\item An {\it infinitesimal deformation of
$X$} is a local deformation with $A$
an artinian local $\mathbb{C}$-algebra with
residue field $\mathbb{C}$. $X$ is
called {\it rigid} if all infinitesimal
deformations are trivial.
\item An {\it inifinitesimal deformation of
order $n$} is an infinitesimal
deformation when $S=\Spec
(\mathbb{C}[t]/(t^{n+1})$.
\end{enumerate}
\end{definition}

{\bf Upshot.} An interpretation of the
so-called {\it dual numbers}
$k[t]/(t^2)$ (see
\cite{Hartshorne-deformation}) as the tangent
space of
something is thinking of Taylor polynomials:
after quotienting by $t^2$ we loose
the tails of the polynomials and are left
with the first derivative information
only - this justifies the use of the ``first
order deformations'' term.

There is the following interpretation in
[continued-fractions] p. 39: the space
of first order deformation classes of $X$ is
$D(\mathbb{C}[t]/(t^2)$. This is
said to ````represent'' the tangent space
$\mathbb{T}^1_X$ of the hypothetical
deformation space of $X$''. (I put double
quotations because the word
``represent'' is quoted in the original
text.) Further, if $X$ is nonsingular
and compact, then
$\mathbb{T}^1_X=H^{1}(X,T_X)$.

Which basically I interpret as: the dimension
of the deformation space of a
smooth compact variety is $H^{1}(X,T_X)$.

\medskip\noindent
After a few months of the previous text in
this section
I think I have understood the construction of
Kodaira-Spencer map for embedded
deformations,
which is \cite[Proposition
3.2.1(ii)]{Sernesi-deformations}.

But I will actually briefly outline a
particular case
of this, which is easier: \cite[Examples
3.2.4(ii)]{Sernesi-deformations}.
This the case when $X \subset Y$ is a Cartier
divisor.
Then $X$ is cut out locally by a set of
functions $f_i$
defined on an affine cover $U_i$ of $Y$.
Then we have a line bundle defined by the
transition functions $f_{ij}=f_i/f_j$.
This is the line bundle $\mathcal{O}_Y(X)$,
which
by the adjunction formula is also the normal
bundle $N_{X/Y}$
of $X$ in $Y$.

Now consider a first order deformation
$$
\xymatrix{
X\ar[r]\ar[d]
&\mathcal{X}\subset Y\times
B_\varepsilon\ar[d]^\pi\\
0\ar[r]
&B_\varepsilon
}
$$

The great thing about first order
deformations
is that the structure sheaf of $\mathcal{X}$
on an affine open $U_i$ of $Y$
looks like the tensor product
$\mathcal{O}_{U_i}\otimes
\mathbb{C}[\varepsilon]/(\varepsilon^2)$
(\href{https://stacks.math.columbia.edu/tag/%
01I4}{01I4}).
By definition of tensor product this just
says
that sections look like $f+\varepsilon g$
for $f,g \in \Gamma(\mathcal{O}_{U_i})$.

If $\mathcal{X}$ is a Cartier divisor on $Y
\times B_\varepsilon$
(it is; but I'm missing a formal proof;
the intuition is that an ``infinitesimal
thickening'' $\mathcal{X}$
of $X$ has the same underlying topological
space
but different scheme structure;
\href{https://stacks.math.columbia.edu/tag/%
01JL}{01JL}
\href{https://stacks.math.columbia.edu/tag/%
01I4}{01I4}),
it is given as the vanishing locus
of some local sections $f_i+\varepsilon g_i$.
(Perhaps the $f_i$ coincide with the
distinguished
section of $X$.)
The line bundle associated to $\mathcal{X}$
is given by the cocycle
$$
F_{ij}=f_{ij}+\varepsilon g_{ij}
$$
for $g_{ij}=\frac{g_i-f_{ij}g_j}{f_j}$.
Indeed, note that
$\frac{1}{f_j+\varepsilon
g_j}=\frac{1}{f_j}-\varepsilon
\frac{g_j}{f^2_j}$
and compute $F_{ij}=(f_i+\varepsilon
g_j)/(f_j+\varepsilon g_j)$.

Then
$$
f_i+\varepsilon g_i=(f_{ij}+\varepsilon
g_{ij})(f_j+\varepsilon g_j)
$$
gives
$$
g_i=f_{ij}g_j+g_{ij}f_j
$$
and since $f_j$ vanishes along $X$,
we are left with $g_i=f_{ij}g_j$
which is how a section of the sheaf $N_{X/Y}$
is constructed by definition of sheaf.


\begin{example}
\label{algebraic-geometry-example-modul-space-nonsi-riema}
\begin{reference}
\cite[Example 3.1]{continued-fractions}
\end{reference}
Fix $g \geq 2$. The dimension of the
deformation space of a nonsingular
projective curve $X$ is $3g-3$. This is ``the
dimension of the moduli space of
curves of genus $g$''.

We can compute this number by Riemann-Roch
formula on the bundle $-K_X$.
Indeed, since $X$ is a curve,
$\Omega_X^1=K_X$ and thus $-K_X=T_X$. We get
$$
\begin{aligned}
h^0(-K)-h^0(K-(-K))&=\text{deg}(-K)-g+1\\
&=-2g+2-g+1\qquad \substack{\text{degree
additive and}\\ \text{deg}K=2g-2}\\
&=-3g+3
\end{aligned}
$$
Now by Serre duality
$h^0(K-(-K))=h^1(-K)=h^1(T_X)$, and it turns
out that a
Riemann surface of genus $g \geq 2$ has no
holomorphic vector fields, so that
$h^0(-K)=h^0(T_X)=0$.
\end{example}

\begin{exercise}
\label{algebraic-geometry-exercise-deformations-of-curves-in-K3}
Let $C$ be a smooth genus $g$ curve which can
be embedded in a K3 surface  $M$,
and $X$ the family of all deformations of $C$
in $M$.
\begin{enumerate}
\item Prove that $\dim X \leq  g$.
\item Let $\mathcal{X}_g$ be the space of all
curves of genus $g$ (smooth?)
which can be possibly embedded to a K3
surface. Prove that each irreducible
component $Z$ of $\mathcal{X}_g$ satisfies
$\dim_\mathbb{C} \leq  g+19$. Deduce
that there exists a compact complex curve
which cannot be embedded in a K3
surface.
\end{enumerate}
\end{exercise}

\begin{exercise}
\label{algebraic-geometry-exercise}
Let $C$ be a smooth curve embedded in a K3
surface $X$. Show that the dimension
of the deformation space of $C$ is $\leq g$.
\end{exercise}

\begin{proof}
% \begin{enumerate}
% \item
The deformation space of a variety is the
space of isomorphism classes of
deformations as explained above. It turns out
that there is a way to associate
1-cocyles of the tangent sheaf to
deformations, so that in fact the deformation
space $\text{Def}_1$ is isomorphic to
$H^{1}(X,\mathcal{T}_X)$ for any variety
$X$.

For our curve $C$ we thus know that the
dimension of the space of deformations
(deformations not necessarily contained in
$X$) is $h^1(\mathcal{T}_C)$.  The
family of deformations of $C$ that are
contained in $X$ is the Hilbert scheme of
curves with fixed Hilbert polynomial $P(t)$
after quotienting by $\mathbb{C}s$,
where $s$ is the section whose vanishing
locus is  $C$.
This says that the number we are looking for
is $h^{0}(X,\mathcal{O}(C))-1$.

Before showing that in fact
$h^0(X,\mathcal{O}(C))=1+g$,
let's explain why this Hilbert scheme
construction actually works.
Intuitively, any curve lying in the same
component
as $C$ would be deformation-equivalent to $C$
--- at least we know that it's Hilbert
polynomial is the same.
But it's not clear what's the deformation
involving these two
(this question should be addressed!).
Further, this doesn't explain Vladimiro's
suggesion of ``taking quotient by
$\mathbb{C}s$''.
To me it looks like I would just compute the
dimension
of the tangent space of $\text{Hilb}^P(M)$ at
$C$.

Now I will show that
$h^0(X,\mathcal{O}(C))=1+g$ (see
\cite[Lemma 1.2.1, Remark 1.2.2]{huk}).

Consider the ideal sheaf exact sequence
twisted by
$\mathcal{O}(C)$:
$$
\xymatrix{
0\ar[r]&\mathcal{O}_X\ar[r]&\mathcal{O}(C)
\ar[r]
&\mathcal{O}_X(C) \otimes
\mathcal{O}_C=\mathcal{O}(C)|_{C}\ar[r]&0
}
$$
By $X$ being a K3 we know that
$H^{1}(\mathcal{O}_X)=0$, so that we have the
short exact sequence in cohomology
$$
\xymatrix{
0\ar[r]&H^{0}(\mathcal{O}_X)\ar[r]&H^{0}
(\mathcal{O}(C))\ar[r]
&H^{0}(\mathcal{O}(C)|_{C})\ar[r]&0
}
$$
so that
$h^0(\mathcal{O}(C))=1+h^1(\mathcal{O}(C)
|_{C})$.
A version of adjunction formula says
$\omega_C \cong (K_X
\otimes\mathcal{O}(C)|_{C}$, and using that
$K_X=\mathcal{O}_X$ we obtain
$h^0(\mathcal{O}(C)|_{C})=h^0(\omega_C)=g$.

To

\medskip\noindent
For the record I put other thoughts I went
through in solving this exercise.

Recall from \cite[p. 146]{gri} that the
normal bundle
$\mathcal{N}$ of a hypersurface of a smooth
variety $X$ satisfies
$\mathcal{N}^\vee\cong\mathcal{O}_X(-C)|_{C}
$.
Taking duals we get that
$\mathcal{N}\cong\mathcal{O}_X(C)|_{C}$.

By adjunction formula
$2g-2=(\mathcal{O}(C),\mathcal{O}(C))$.
Applying Riemann-Roch
to $\mathcal{O}(C)$ (which is by definition
the dual of the ideal sheaf of $C$,
which is a line bundle on $X$), we obtain
that
$\chi(\mathcal{O}(C))=2+\frac{1}{2}
(\mathcal{O}(C),\mathcal{O}(C))$.
Then
$\chi(\mathcal{O}(C))=g+1$.

Now recall that
$\chi(\mathcal{O}(C))=h^0(\mathcal{O}(C)
-h^1(\mathcal{O}(C))+h^2(\mathcal{O}(C))$. By
Serre duality and $X$ being a K3
surface we see that $h^2(\mathcal{O}(C))\cong
h^{0}(\mathcal{O}(-C))$, which is
the ideal sheaf of $C$. Any section of such a
sheaf would vanish along $C$, and
since $X$ is compact we conclude there cannot
be any such section.

Now we show that also $h^1(\mathcal{O}(C))=0$
to conclude that
$h^0(\mathcal{O}(C))=g+1$.

We thus conclude that $h^0(\mathcal{O}(C))$




Since $C$ is smooth we can use the normal
exact sequence
$$
\xymatrix{
0\ar[r]&\mathcal{T}_C\ar[r]&\mathcal{T}
_X|_{C}\ar[r]&\mathcal{N}\ar[r]&0
}
$$
Taking Euler characteristic we see that
$\chi(\mathcal{T}_C)+\chi(\mathcal{N})=
\chi(\mathcal{T}_X|_{C})$.


This means that we should be done once we
compute $\chi(\mathcal{T}_X|_{C})$.
For this we can use Riemman-Roch formula for
coherent sheaves on a curve, which
tells us that
$$
\text{deg}(\mathcal{T}_X|_{C})=
\chi(\mathcal{T}_X|_{C})
-\text{rk}(\mathcal{T}_X|_{C})
\cdot\chi(\mathcal{O}_C)
$$
But of course we know that since $C$ is a
curve, by Serre duality we get that
$\chi(\mathcal{O}_C)=1-g$. (Indeed:
$h^1(\mathcal{O}_C)=h^0(\Omega^1_C):=p_a(C)
$.)
And now the question is what is
the degree. Apparently this is just the first
Chern class $c_1$ of the
restricted tangent bundle. So what is it? And
what is $h^0(\mathcal{T}_C)$ if
the genus is 0 or 1, and that's it.

I
know that $h^0(\mathcal{T}_X)=0$ and
$h^1(\mathcal{T}_X)=20$ by $X$ being a K3
surface, but I'm not sure what happens when
we restrict to $C$.

If $g \geq 2$ we know that
$h^0(\mathcal{T}_C)=0$, so that
$\boxed{\chi(\mathcal{T}_C)=h^1(\mathcal{T}
_C)}$.
\bigskip
To compute the
latter Euler characteristic, which is given
by definition by
$\chi(\mathcal{T}_X|_{C})=h^0(\mathcal{T}
_X|_{C})-h^1(\mathcal{T}_X|_{C})$,
we
first note that $h^0(\mathcal{T}_X|_{C})=0$,
because this is the dimension of
the global holomorphic vector fields on $X$
restricted to $C$, which is constant along
$C$ since $X$ is
smooth. And then this number is actually zero
by Hodge numbers of a K3 surface.

and taking cohomology long exact sequence we
obtain
$$
\xymatrix{
\cdots\ar[r]&H^{0}(T_X)\ar[r]&H^{0}
(\mathcal{N})\ar[r]&H^{1}(T_C)\ar[r]&\\
H^{1}(T_X)\ar[r]&H^{1}(\mathcal{N})\ar[r]&
\cdots
}
$$
\medskip\noindent
Or is it $h^1(\mathcal{N})$? See
\cite{HarrMorr}. If this was the case, then
we
can use adjunction formula as above to get
that
$2g-2=(\mathcal{N},\mathcal{N})=\text{deg}_C
\mathcal{N}$. Then we may find
$h^0(\mathcal{N})$ via Riemann-Roch:
$$
h^0(\mathcal{N})-h^0(K_C-\mathcal{N})=
\text{deg}\mathcal{N}-g+1
$$
Note that
$h^0(N_C-\mathcal{N})=h^1(-K_C+N+K_C)=
h^1(\mathcal{N})$
via Serre
duality, and by Riemann-Roch on a surface as
above we see that
$$
g+1=\chi(\mathcal{N})=h^0(\mathcal{N})-
h^1(\mathcal{N})
\implies
h^1(\mathcal{N})=-g-1+h^0(\mathcal{N})
$$
so that
$$
\begin{aligned}
h^0(\mathcal{N})-(-g-1+h^0(\mathcal{N}))&=
\text{deg}\mathcal{N}-g+1\\
\implies \text{oops! I lost
$h^0(\mathcal{N})$ in this operation…}
\end{aligned}
$$
Maybe if I just use normal exact sequence and
realise that
$$
h^0(\mathcal{N})=h^0(\mathcal{T}_X|_{C})-
h^0(\mathcal{T}_C)
$$
I know that $h^0(\mathcal{T}_C)=0$ for $g>1$,
so the question is how to compute
the restricted holomorphic vector fields.
% \end{enumerate}
\end{proof}

\medskip\noindent
Look for the ``Perfect obstruction theory''.
This is used by Bruzzo for deformations of
(super)stacks,
work with E. Paiva and ---.

\medskip\noindent
Until I have a better candidate of where to
put this…

\begin{definition}
\label{algebraic-geometry-definition-extension}
\begin{reference}
\cite[p.7]{Sernesi-deformations}
\end{reference}
Let $A \to R$ be a ring homomorphism
(i.e. $R$ is an $A$-algebra).
An {\it $A$-extension of $R$ by $I$} is a
short exact sequence
$$
\xymatrix{
0\ar[r]
&I\ar[r]
&E\ar[r]^\varphi
&R\ar[r]
&0.
}
$$
where $E$ is an $R'$-algebra
and $\varphi$ is a morphism of $A$-algebras
whose kernel is $I$.
We also require that $I^2=(0)$;
this condition implies that $I$ has the
sturcture of $R$-module.
\end{definition}

\begin{example}
\label{algebraic-geometry-example-dual-numbers-extension}
$$
\xymatrix{
0\ar[r]
&(\varepsilon)\ar[r]
&A[\varepsilon]\ar[r]
&A\ar[r]
&0
}
$$
where $A[\varepsilon]$ is the algebra of dual
numbers over $A$.
\end{example}

\section{Cotangent modules}
\label{algebraic-geometry-section-cotangent-modules}

\noindent
My aim in this section is to cover the
minimal
theory needed to apply results in \cite{jan2}
to my deformation problem. The main theorem
I expect to use is the following:

\begin{theorem}
\label{algebraic-geometry-theorem-deformations-threefolds}
\begin{reference}
\cite[Theorem 5.7]{jan2}
\end{reference}
If $\mathcal{K}$ is a 3-dimensional manifold
then
$$
\dim T^1_{\mathbb{P}(K)}
=11d_3+5e_3+3e_4+e_{\geq 5}+c_{\geq
6}+5f_1^{(3)}+2f_1^{(4)}+h^2(\mathcal{K}).
$$
\end{theorem}

\noindent
I need to do some interpretation first.
From \cite[p. 1]{jan3}:

\begin{quotation}
  ``For each $\mathbb{C}$-algebra $A$,
there is a cohomology theory providing
modules $T^i_A$.
However, only $T^1_A$ and $T^2_A$ are
relevant for the deformation
theory of $\Spec A$. The module $T^1_A$
collects the infinitesimal
deformations of $A$ and $T^2_A$ contains the
obstructions
for lifting these deformations to decent
parameter spaces.
Eventually, $\text{Der}_{\mathbb{C}}(A,A)$
will be called $T_A^0$."
\end{quotation}

\cite[p. 3]{jan3}:
\begin{quotation}
``In general, for a finitely generated
$\mathbb{C}$-algebra $A$,
the modules $T_A^i$ allow the following ad
hoc definitions:
Let $P=\mathbb{C}[x_0,…,x_n]$ mapping onto
$A$ so that $A \simeq P/I$
for an ideal $I$. The $T_A^1$ is the cokernel
of the natural map
$\text{Der}_{\mathbb{C}}(P,P) \to
\Hom_P(I,A)$.
Moreovr, […] and we obtain $T^2_A$ as the
cokernel
of the induced map
$\Hom_P(P^m,A)\to\Hom_A(R/R_0,A)$.''
\end{quotation}

\begin{remark}[dani]
\label{algebraic-geometry-remark-dani1}
Let's think more about that definition of
$T^1_A$
as the cokernel of the ``natural map''
$\text{Der}_\mathbb{C}(P,P) \to \Hom_P(I,A)$.

Pick a derivation $\chi \in
\text{Der}_{\mathbb{C}}(P,P)$.
We want a map $I \to A=P/I$.
Pick an element $a \in I$ and map it to
$\chi(a)$.
And that's it. It might be zero or not.
If the $\chi$ we picked maps any $a \in I$ to
some element in $I$,
then this $\chi$ is in the kernel of the
natural map.
So $T^1_A$ is the image of the natural map
modulo the kernel.
\end{remark}

\noindent
Unfortunately this doesn't say much about the
following crucial fact:
the $T^i_A$ modules are
$\mathbb{Z}^{n+1}$-graded.
But here goes an explanation.
We use the notation that $e_p$ is a ``face
not in the complex'',
which of course corresponds to a generating
monomial $x^p$
of the SR ideal $I$.
It has a degree, of course (given by the
``support'' of $p$).
Now here's a rather sloppy step but:
since the cotangent modules are built upon
$I$,
who is given by $e_p$,
we just notice that the degree of $e_p$
along with a (very natural, too) notion of
degree among relations,
induce a grading in all the cotangent modules
$T_A^i$.
Thus, the grading is morally just the degree
of the monomials,
which is in fact the support of $p$, a vector
$\mathbf{c} \in \mathbb{Z}^{n+1}$.

Next in line is the decomposition
$\mathbf{c}=\mathbf{a}-\mathbf{b}$
into positive minus negative part. But how
exactly…?








\section{Fano varieties}
\label{algebraic-geometry-section-Fano-varieties}

\begin{definition}
\label{algebraic-geometry-definition-Fano-variety}
A {\it Fano variety} is a projective variety
with $-K_X$ ample.
\end{definition}

\noindent
The idea of the Fano index is that the
larger the index, the simpler is the variety.
In a talk by Jõao, it is defined as the
larger integer
dividing $-K_X$ in $\text{Pic}(X)$.
\begin{definition}
\label{algebraic-geometry-definition-Fano-index}
$$
r(X):=\text{min}\{r:\frac{c_1(X)}{r}\in
H^{2}(X,\mathbb{Z})\}
$$
\end{definition}

\begin{exercise}
\label{algebraic-geometry-exercise-fano-vanis-highe-cohom}
By Kodaira vanishing theorem
\ref{algebraic-geometry-theorem-Kodaira-vanishing},
you can show that the cohomology $H^{i}(X,L)$
for
a Fano variety $X$ vanishes. You just have to
put $L=\mathcal{O}(k)$ with $k\geq
-r$, where $r$ is the Fano index.
\end{exercise}

\begin{exercise}
\label{algebraic-geometry-exercise-Pic-H2-Fano}
Show that  $\text{Pic}(X)\cong
H^{2}(X,\mathbb{Z})$ holds for Fano
varieties.
\end{exercise}

\begin{remark}
\label{algebraic-geometry-remark-deriv-categ-fano-3-folds}
\begin{reference}
Marcos Jardim, CIMPA 2025 Florianópolis,
Lecture 2.
\end{reference}
If $H^3(X,\mathbb{Z})=0$ of a Fano 3-fold,
then its derived category is
generated by 4 elements.
\end{remark}

\medskip\noindent
Araujo-Druel em 2017 classify
Fano foliations relating the Fano index (of
the foliation)
with its rank.
João found conditions
(integrability of the leaves)
to describe del Pezzo foliations
(a kind of Fano foliation)
with Picard rank 1.
Araujo (2019) classifies the leaves of a
algebraically integrable
del Pezzo foliation (with no restriction on
the (log
canonica) singularities!).


\section{Quivers}
\label{algebraic-geometry-section-quivers}

\begin{definition}
\label{algebraic-geometry-definition-quiver}
A {\it quiver} is a set of vertices $Q_0$, a
set of arrows $Q_1$ equipped with
the maps of source $s$ and target $t$ that to
each arrow they assign the point
that is source or target of the arrow.
\end{definition}

\begin{definition}
\label{algebraic-geometry-definition-representation-of-quiver}
A {\it representation} of a quiver is a set
of finite dimensional vector spaces
equipped with maps between them realising a
given quiver (incomplete…).
\end{definition}

There is a notion of projective
representation, which I missed to write. But
it
is analogous to the injective representation:

\begin{definition}
\label{algebraic-geometry-definition-injec-repre-quive}
Given a quiver $Q$, the {\it injective
representation} of $Q_0$ is given by, for
$i \in Q_0$,
$$
I(i)_j=\begin{cases}
k\qquad &i=j \\
k^{d'}\qquad &j \neq i
\end{cases}
$$
where $d'$ is the number of paths from $j$ to
$i$.
\end{definition}

\section{Stacks}
\label{algebraic-geometry-section-stacks}

My first definition of stack can be extracted
from

\begin{definition}
\label{algebraic-geometry-definition-superstack}
A {\it superstack} is a stack over
the étale site $\text{SSch}$ of superschemes,
i.e. it is a category fibered in groupoids
over the category of superschemes,
the latter equipped with the
étale topology,
satisfying the descent condition.
\end{definition}

Here are some other definitions:

\begin{definition}
\label{algebraic-geometry-definition-algebraic-stacks}
Let  $\mathfrak{X}$ be a stack over
$\text{Sch}_{\text{ét}}$.
An {\it algebraic space} is
such that there exists morphism
$\mathcal{U} \to \mathfrak{X}$
where $\mathcal{U}$ is a scheme, that is
schematic, étale and injective (check this
one).

$\mathfrak{X} \to y$ is {\it representable}
if
there exists a scheme $\mathcal{U}$ and a map
$\mathcal{U} \to y$ such that the
fibered product
$$
\xymatrix{
\mathcal{U} \times_y
\mathfrak{X}\ar[r]\ar[d]\ar@{}[dr]|-
{\lrcorner}&\mathfrak{X}\ar[d]\\
\mathcal{U}\ar[r]&y
}
$$
is an algebraic space.

Finally, a  stack is {\it algebraic} (resp.
{\it Deligne-Mumford})
is there exists a
representable surjective morphism
$\mathcal{U} \to \mathfrak{X}$
that is smooth (resp. étale).

A {\it stable map} over a projective
variety $X$ is an element of the first
Chow group $\beta \in A_1$, where
 $(C,g)$ is an algebraic curve and
$f:C \to X$ with $[f(C)]=\beta$.
\end{definition}

\noindent
The curves that are points under this map
(contractible) are {\bf stable}.

\section{Moduli}
\label{algebraic-geometry-section-moduli}

\noindent
Pseudo-admissible covers may be the way
Mumford orifinally developed stability
notions.

\noindent
This is a two-session minicourse by Ugo
Bruzzo
in LEGALzinho 2025.

Let $X$ be a smooth complex projective
variety,
$X \overset{j}{\hookrightarrow}\mathbb{P}^n$.
$\mathcal{O}_{\mathbb{P}^n}(1)$ is the class
of hyperplanes of $\mathbb{P}^n$.
$\mathcal{L}=j^*\mathcal{O}_{\mathbb{P}^n}(1)
$.
$c_1(\mathcal{L})\in H^2(X,\mathbb{Z})$ is
a {\it polarization}.

Let $\mathcal{E}$ be a coherent
$\mathcal{O}_X$-module.
For example, $\mathcal{E}$ may be a locally
free $\mathcal{O}_X$-module.
Recall that $H^k(X,\mathcal{E})$ are finite
dimensional vector spaces
and they vanish for $k>\dim X=n$.
Recall that the {\it Euler characterstic} of
$\mathcal{E}$
is given by
$\chi(\mathcal{E})=\sum_{k=0}^n(-1)
^k\dim_{\mathbb{C}}H^k(X,\mathcal{E})$.

The {\it reduced Hilbert polynomial} of
$\mathcal{E}$ is
$$
p_{\mathcal{E}}(k)=\frac{1}{\text{rk}
\mathcal{E}}
\chi(\mathcal{E}\otimes_{\mathcal{O}_X}
\mathcal{L}^k).
$$

\begin{definition}
\label{algebraic-geometry-definition-stable-module}
Let $\mathcal{E}$ be a coherent
$\mathcal{O}_X$-module
with no torsion and positive rank.
$\mathcal{E}$ is {\it stable} if for all $f
\subset \mathcal{E}$
($0 \text{rk} f<\text{rk}\mathcal{E}$),
$$
p_f(k)<p_{\mathcal{E}}(k)\qquad \text{for }
k\gg 0
$$
and {\it semi-stable} if we put $\leq$
instead of $<$.

\begin{lemma}
\label{algebraic-geometry-lemma-stable}
If $\mathcal{E}$ is stable then $\mathcal{E}$
is irreducible
($\mathcal{E} \neq  \mathcal{E}_1 \oplus
\mathcal{E}_2$)
and $\text{Aut}(\mathcal{E})\simeq
\mathbb{C}^*$.
\end{lemma}

\end{definition}

The {\it functor of moduli} is
$$
\begin{aligned}
\mathcal{M}: \text{Sch}_\mathbb{C}
&\longrightarrow \mathsf{Set} \\
S
&\longmapsto
\left\{
\substack{\text{stable family}  \\ \text{of
sheaves with} \\
\text{Hilbert polynomial} f\\
\text{parametrized by $S$}}
\right\}\Big/\sim.
\end{aligned}
$$

\medskip\noindent
Riemann-Roch.
Let $C$ be a smooth, projective, connected
curve
of genus $g$.
Let $\mathcal{E}$ be a coherent sheaf over
$H^2(C,\mathbb{Z})$
with generator of $H^2(X,\mathbb{Z})$
$w=[pt] \in H^2(C,\mathbb{Z})$.
$c_1(\mathcal{E})=\text{deg}\mathcal{E}\cdot
w$.

\begin{theorem}
\label{algebraic-geometry-theorem-euler-characteristic-moduli}
$$
\begin{aligned}
\chi(\mathcal{E})
&=\int_C(1+(1-g)w)(\text{rk}\mathcal{E}+
\text{deg}\mathcal{E}w)\\
&=\int_C
\text{deg}\mathcal{E}+w+(1-g)\text{rk}
\mathcal{E}\cdot
w\\
&=\text{deg}(\mathcal{E})+1-g
\text{rk}\mathcal{E}.
\end{aligned}
$$
\end{theorem}

\begin{example}
\label{algebraic-geometry-example-hilbert-polynomial-sheaf}
\end{example}

\begin{theorem}
\label{algebraic-geometry-theorem-moduli-suave}
Se $(r,d)$, $r>0$ are coprime,
the functor of moduli of stable bundles
over $C$ of rank $r$ and degree $d$
is representable, and the moduli space
$\mathcal{M}_C(\mathcal{E},d)$ is smooth.
\end{theorem}

\noindent
In the non coprime case this does not
necessarily hold.

\begin{definition}
\label{algebraic-geometry-definition}
$\mathcal{M}$ is a {\it grosseiro} moduli
space
if the funtor is not representable.
\end{definition}

\noindent
This means that although there is not
a universal family,
there is still a correspondence
of the closed points of $\mathcal{M}$ and the
equivalence classes of semistable sheaves
of rank $r$ and degree $d$.

\medskip\noindent
Now let's fix $\dim X=2$.
Let $\mathcal{E}$ be a stable
$\mathcal{O}_X$-module.
Let $r=\text{rk}\mathcal{E}$,
and the {\it discriminant}
$$
\Delta(\mathcal{E})=c_2(\mathcal{E})-
\frac{\mathcal{E}-1}{2r}c_1(\mathcal{E})^2
\in H^4(X,\mathbb{Q})\cong\mathbb{Q}.
$$
Here enter $c_1(\mathcal{E}) \in
H^2(X,\mathbb{Z})$
and $c_2(\mathcal{E}) \in H^4(X,\mathbb{Z})$.

\begin{theorem}[Bogomolov]
\label{algebraic-geometry-theorem-bogomolov}
$\Delta(\mathcal{E}) \geq 0$.
\end{theorem}

\noindent
It could be interesting to study the proof of
this theorem.

\medskip\noindent
Miyaoka (1987). Consider a curve $C$ and a
bundle $E$.
We can for any fibre, which is a vector space
$V$,
consider its projectivization
$V\setminus\{0\}/\mathbb{C}^*=\mathbb{P}(V)$.
Then we can bundlize this construction ad
obtain
$\mathbb{P}E \to C$.
Consider $\mathcal{O}_{\mathbb{P}E}(1)$.
We define
$$
\lambda=rc_1(\mathcal{O}_{\mathbb{P}E}(1))
-\pi^*c_1(E) \in H^2(\mathbb{P}E,\mathbb{Z}).
$$

\begin{theorem}
\label{algebraic-geometry-theorem-semistable-nef}
$E$ is semistable if $\lambda$ is
semicorrente efectiva.
\end{theorem}

\noindent
\begin{definition}
\label{algebraic-geometry-definition-nef}
Consider again a projective smooth connected
curve $\Gamma$ and let $f:\Gamma \to
\mathbb{P}E$.
We say $\lambda$ is {\it nef} if
$$
\int_{\Gamma} f^*\lambda\geq 0.
$$
\end{definition}


\noindent
The following theorem was proved by Nakayama
(who is a different author than that of
Nakayama's lemma)
in 1999 but was unkown to Bruzo and Daniel
HR.
\begin{theorem}[Nakayama '99, Bruzzo-HR '04]
\label{algebraic-geometry-theorem-bruzzo-daniel-nakayama}
The following are equivalent:
\begin{enumerate}
\item $\lambda$ is nef.
\item For all $f:\Gamma \to X$,
$f^*E$ is semi-stable.
\item $E$ is semistable with respect to
the polarization, and  $\Delta(E)=0$.
\end{enumerate}
\end{theorem}

The polarization says how the surface is
embedded in projective space,
it is $H=c_1(\mathcal{L})$. The stability
notion is
$$
\mu(E)=\frac{c_1(E)\cdot
H^{n-1}}{\text{rk}E}.
$$
\medskip\noindent
They have proved that for Higgs bundles 1
$\iff$ 2, $3 \implies 1,2$.
And also for $1,2 \implies 3$ but only for
$r=2$ using the Higgs Grassmannian,
but the problem is very complicated for
general rank.

\medskip\noindent
Second lecture. I think this might be a more
general
definition of the Hilbert polynomial:
$$
\sum_{k=0}^n(-1)^k \dim
H^k(X,\mathcal{E})=\chi(\mathcal{E})
=\int_X\text{ch}(\mathcal{E})\text{td}X
$$
And we also have:
$$
\chi(\mathcal{E}\otimes \mathcal{L}^*)
=\int_X (1+(1-g)w)(r+kr
\underbrace{c_1(\mathcal{L})}_{\text{deg}
\mathcal{L}\cdot
w}
+\text{deg}\mathcal{E}\cdot w.
$$
And then
$$
p_{\mathcal{E}}=k \text{deg} \mathcal{L}
+1-g+\mu(\mathcal{E}).
$$



\medskip\noindent
Now we discuss Higgs bundles.
Let $X$ be ac onnected complex projective
variety.

\begin{definition}
\label{algebraic-geometry-definition-higgs-bundle}
A {\it Higgs bundle} on $X$
is a pair $\mathfrak{C}=(E,\phi)$ with
\begin{itemize}
\item $E$ a vector bundle on $X$
\item $\phi:$ a morphism (connection-like…)
\end{itemize}
\end{definition}

For the following see Nitsuze, FGA explained.
There following scheme is also representable
(I think this is the ``universal quotient'')

We have also consider the Grassmanization of
a vector bundle.
We obtained a short exact sequence
and tried to compute the tensor product with
the holomorphic differentials.
The idea was to pullback the grassman bundle
given a map $f:Y \to X$.
The result is not satisfactory, the pulled
back
thing may be strange. This is called the
{\it Higgs Grassmannian}.


Here's some notion also studied by Campana,
Demailly:
numerically flat bnef bundle $E$.
Then we have that
\begin{lemma}
\label{algebraic-geometry-lemma-proof-hugo-bruzzo}
If $c_1=0$ then numerically flat bundle iff
semistable.
\end{lemma}

\medskip\noindent
This is the beginning of a talk by Alan Muniz
at Bandoleros 2025 titled
``Rank 2-bundles on $\mathbb{P}^3$ with odd
determinant''.
(Work in progress with A. Fontes.)

Problem: Classify rank-2 bundles on
$\mathbb{P}^3$.

We consider bundles $E$ with chern classes
$c_1=-1$, $c_2=2n$ and $c_3=0$.
Let $\mathcal{B}(-1,2n)$ be the moduli space
of $\mu$-stable rank 2-bundles.

\begin{itemize}
\item $\mathcal{B}(-1,2)$ is irreducible of
dimension 11
[Hartshorne-Sols '82, Manolache '81].

\item $\mathcal{B}(-1,4)$ has two components
[Barnica-Manolache '85].

\item $c_2\geq 6$, Tikhomirov (2014)
Fontes-Jardim (2023, 2024, 2025).
$\mathcal{B}(-1,6)$ has $\geq 4$ components.
\end{itemize}

\medskip\noindent
Now let's discuss Serre correspondence.

How can we describe these spaces?
Let $E$ be a vector bundle of rank 2
on $\mathbb{P}^3$ and $\sigma \in H^0(E(r))$.
Then $C=(\sigma_0)\subset \mathbb{P}^3$ is
a curve * equipped
with an ``extension class''
$\zeta \in H^0(\omega_C(4-2r-c_1))$.
In fact, this correspondence
is 1-1, and we call it Serre correspondence.
We have that $\text{deg}C=c_2(E(r))$ and
$p_a(C)=1=\frac{c_2(E(r))(c_1(E(r))-4)}{2}$
To pass from the curve to the bundle
and viceversa we use
$$
\xymatrix{
0\ar[r]
&\mathcal{O}(-r)\ar[r]
&E\ar[r]
&\mathcal{I}_C(c_1+r)\ar[r]
&0
}
$$


\medskip\noindent
Upshot about Deligne-Mumford stacks:
the sit below a stack via a map that is
étale.


















\medskip\noindent
Now let's discuss framed moduli.

\noindent
Framed instantons on the blow-up of
$\mathbb{P}^3$ at a point.

Let $X$ be a projective variety and fix a
polarization $H$.
Let $\mathcal{E}$ be a sheaf on $X$.
A {\it framed pair} is a
$(\mathcal{F},\alpha)$ where
$\mathcal{F}$ is a coherent sheaf on $X$
and $\alpha$ is a {\it framing datum}, i.e.
a map $\alpha:\mathcal{F} \to \mathcal{E}$.

Set
$$
\mathcal{E}(\alpha)
=\begin{cases}
	1\qquad &\alpha\neq 0 \\
	0\qquad &\text{otherwise} .
\end{cases}
$$
Define the Hilbert polynomial by
$$
P(k)=\chi(\mathcal{F} \otimes
\mathcal{O}(kH)).
$$
For fied $\delta \in \mathbb{Q}[t]$ rational
positive,
define
$$
P_{(\mathcal{F},\alpha)}=P_{\mathcal{F}}(k)-
\varepsilon(\alpha)\delta.
$$

We can also define induced framings from
subsheaves $\mathcal{F}' \subset
\mathcal{F}$, by asking that the inclusion
commutes
with the framing datum $\alpha$.
Then we put
$$
\varepsilon(\alpha')
=\begin{cases}
	1\qquad &\mathcal{F}' \subseteq Ker \alpha
	\\
	0\qquad &\text{otherwise} .
\end{cases}
$$
We say that $\mathcal{F}$ is {\it
$\delta$-(semi)stable}
if for any subpair $(\mathcal{F}',\alpha')$
with
$\text{rk}\mathcal{F}'=r'<r=\text{rk}
\mathcal{F}$
one has
$$
\frac{1}{r'}P_{(\mathcal{F}',\alpha')}(k)
\leq  \frac{1}{r}P_{(\mathcal{F},\alpha)}(k).
$$
If $\text{deg}(\delta)\leq \dim X$,
there are moduli functors
$\mathcal{M}^{st}(X,\mathcal{E},\delta)
\subseteq
\mathcal{M}^{\text{ss}}(X,\mathcal{E},\delta)
$.

\begin{remark}
\label{algebraic-geometry-remark-degree-greater-than-dimension}
If $\text{deg} \delta>\dim X$
then $\alpha$ is injective or zero, and
$\mathcal{F}$ is semistable.
This gives quot schemes.
\end{remark}

\noindent
Huybrechts-L. In $\dim X\leq 2$
there exists a quasi-projective fine moduli
space
for $\delta \in \mathbb{Q}[t]$ and $p \in
\mathbb{Q}[k]$.

\medskip\noindent
Then there is a notion of good framing by
Bruzz-M,
depending on $D$, an effective divisor on
$X$.
They proved that having good framing is
enough
to obtain a fine quasi-projective moduli
space.

\begin{definition}
\label{algebraic-geometry-definition-instanton}
A coherent sheaf $\mathcal{E}$ on
$X=\tilde{\mathbb{P}^3}$
is an {\it instanton} if it is
$\mu_{\mathcal{L}}$-semistable,
$c_1(\mathcal{E})=0$ and
$H^0(\mathcal{E})=0=H^0(Ec(-3,2))$,
$H^1(\mathcal{E}(-2,1))=0=H^2(\mathcal{E}(-2,
1))$
and
$H^3(\mathcal{E}(0,-1))=H^2(\mathcal{E}(-1,1)
)$.
\end{definition}

\noindent
Some examples are locally free sheaves of
rank 2 (via Harrshorne-Serre
correspondence),
torsion-free rank 2 constructed via
elementary
transformations.

\medskip\noindent
General result: for rank 2 instantons on Fano
threefolds
that have a plane $E$,
$\mathcal{E}|_{E}$ is $\mu$-semistable
if and only if $h^1(\mathcal{E}(-2,1) \otimes
\mathcal{O}(-E))=0$.




\section{Stability}
\label{algebraic-geometry-section-stability}

\noindent
Stability helps us select the
building blocks of an abelian
category, every object in the
category can be reconstructed
from those.
Wall and chamber structure: inside a
compact set $K \subset
\text{Stab}(\mathcal{T})$,
only finitely many walls appear.
If $C$ is a connected component of
the complement of the walls in $K$,
then the stable objects are constant
on $C$.
Think of an ant walking in
$\text{Stab}(\mathcal{T})$ with a
backpack of objects and a detector:
when it crosses a wall and some object
stops being stable, the red light fires.
Stable objects change only when you
cross a wall.

\begin{center}
\setlength{\unitlength}{1mm}
\begin{picture}(98,78)(0,0)
\scriptsize

% axes
\linethickness{0.35pt}
\put(6,7){\vector(1,0){88}}
\put(88,7){\vector(0,1){68}}
\put(92,10){$\beta$}
\put(91,68){$\alpha$}

% Gieseker/PT chamber boundary
\linethickness{0.45pt}
\qbezier(18,7)(18,13)(24,15)
\qbezier(24,15)(31,16)(38,19)
\qbezier(38,19)(49,23)(60,22)

\qbezier(60,22)(57,27)(57,35)
\qbezier(57,35)(50,39)(47,45)
\qbezier(47,45)(41,53)(52,56)
\qbezier(52,56)(58,57)(68,58)
\qbezier(68,58)(75,62)(77,74)

% diagonal wall, blue in the photo
\linethickness{0.9pt}
\qbezier(22,73)(46,40)(70,7)

% semicircular wall, red in the photo
\linethickness{0.75pt}
\qbezier(45,7)(45,23)(60,22)
\qbezier(60,22)(75,23)(75,7)

% marked point and labels
\linethickness{0.4pt}
\put(59.35,21.35){\circle*{1.3}}
\put(61.5,23.5){$(\beta^\nu,\alpha^\nu)$}

\put(20,36){Gieseker chamber}
\put(47,59){PT chamber}
\put(39,42){$\theta^-_\nu$}

\end{picture}
\end{center}

Let $\mathcal{A}$ be an abelian
category with Grothendieck group
$K_0(\mathcal{A})$.

\begin{definition}
\label{algebraic-geometry-definition-stability-function}
An additive homomorphism
$Z: K_0(\mathcal{A}) \to \mathbb{C}$
is a {\it stability function} if for all
$0 \neq E \in \mathcal{A}$,
$$
\text{Im}Z(E) \geq 0\qquad
\text{and}\qquad
\text{Re}Z(E) <0 \text{ if }
\text{Im}Z(E)=0.
$$
The {\it slope} of $E \in \mathcal{A}$
with respect to $Z$ is
$$
M(E)=
\begin{cases}
-\frac{\text{Re}Z(E)}{
\text{Im}Z(E)}\qquad &
\text{if }\text{Im}Z(E) \neq 0 \\
\infty\qquad &
\text{if }\text{Im}Z(E)=0.
\end{cases}
$$
A nonzero object $E \in \mathcal{A}$
is called {\it stable} if for all
proper non trivial subobjects
$F \subset E$, we have
$$
M(F) < M(E)
$$
and {\it semistable} if we put
$\leq $ instead.
\end{definition}

\begin{definition}
\label{algebraic-geometry-definition-stability-condition}
A pair $(\mathcal{A},Z)$ is called
a {\it stability condition} if
\begin{enumerate}
\item $Z$ is a stability function, and
\item every $0 \neq E \in \mathcal{A}$
has a Harder-Narasimhan filtration.
\end{enumerate}
\end{definition}

\begin{lemma}
\label{algebraic-geometry-lemma-harder-narasimhan-existence}
Let $Z:K_0(\mathcal{A})\to \mathbb{C}$
be a stability function.
Assume that
\begin{enumerate}
\item $\mathcal{A}$ is Noetherian, and
\item the image of $\text{Im}Z$
is discrete in $\mathbb{R}$.
\end{enumerate}
Then every $0 \neq E \in \mathcal{A}$
has a Harder-Narasimhan filtration
with respect to $Z$.
\end{lemma}

\noindent
Let $\mathcal{D}$ be a triangulated
category. Fix a finite rank lattice
$\Lambda$, a surjective homomorphism
$v:K_0(\mathcal{D}) \to \Lambda$
and a norm $\|\cdot\|$ in $\Lambda$.

\begin{definition}
\label{algebraic-geometry-definition-bridgeland-stability}
A {\it Bridgeland stability condition}
$\mathcal{D}$ is a pair $(\mathcal{A},Z)$
where $\mathcal{A}$ is the heart of a
bounded $t$-structure on $\mathcal{D}$,
such that
\begin{enumerate}
\item $(\mathcal{A},Z)$ is a stability
condition,
\item the following support property
holds:
$$
\text{inf}\left\{
\frac{|Z(E)|}{
\|v(E)\|}
:0 \neq  E \in \mathcal{A}
\text{semistable}\right\}
>0.
$$
\end{enumerate}
\end{definition}

\begin{example}
\label{algebraic-geometry-example-stability-curve}
Let $X$ be a smooth projective
curve and $\mathcal{D}=D^b(X)$.
Take $\mathcal{A}=\Cohstack(X)$ as the
heart of the standard
$t$-structure.
Let
$$
Z(E)=-\deg(E)+i\,\mathrm{rk}(E).
$$
Then $\sigma$-stability coincides
with slope stability.
\end{example}

\noindent
Let $Stab^{Br}_\Lambda(\mathcal{D})$
be the set of Bridgeland stability
conditions and let
$Stab_\Lambda(\mathcal{D})$ be the set
of stability conditions on $\mathcal{D}$
with respect to $(\Lambda,v)$.

\begin{theorem}[Bridgeland]
\label{algebraic-geometry-theorem-bridgeland}
The natural map
$$
Stab_\Lambda(\mathcal{D})
\to \Hom(\Lambda,\mathbb{C})
$$
is a local homeomorphism.
$Stab_\Lambda(\mathcal{D})$ is a complex
manifold and
its dimension is the rank of $\Lambda$.
In particular,
$$
Stab_\Lambda^{Br}(\mathcal{D})
\subset
Stab(\mathcal{D}).
$$
\end{theorem}

\noindent
Let $(\mathcal{A},Z)$ be a stability
condition and $M$ be the slope with
respect to $Z$.
For each $\beta \in \mathbb{R}$
we define the following pair of
subcategories on $\mathcal{A}$:
$$
\begin{aligned}
\mathcal{T}^\beta(A,Z)&=\mathcal{T}^\beta
=\{\mathcal{E} \in \mathcal{A}
: M(\mathcal{G})>\beta
\text{ whenever }
\mathcal{E} \to\!\!\!\!\!\to
\mathcal{G}\}\\
\mathcal{F}^\beta(\mathcal{A},Z)
&=\mathcal{F}^\beta
=\{\mathcal{E} \in \mathcal{A}
:M(\mathcal{F})\leq \beta
\text{ for every }
0 \neq  \mathcal{F} \subset \mathcal{E}\}.
\end{aligned}
$$

\begin{definition}
\label{algebraic-geometry-definition-tilt}
The {\it tilt}
$\mathcal{A}^*
=\left\langle{}\mathcal{F}[1],\mathcal{T}\right\rangle{}
\subset\mathcal{D}$
of $\mathcal{A}$ with respect to
$(\mathcal{T}^\beta,\mathcal{F}^\beta)$
is the smallest extension closed full
subcategory of $\mathcal{D}$ containing
$\mathcal{T}^\beta$ and
$\mathcal{F}^\beta[1]$.
\end{definition}

\noindent
$\mathcal{A}^*$ is the heart of a bounded
$t$-structure on $\mathcal{D}$.

\begin{definition}
\label{algebraic-geometry-definition-gieseker-stability}
Let $\mathcal{E}$ be a torsion free
sheaf.
We say that $\mathcal{E}$ is
{\it Gieseker semistable} if
for every subsheaf $F \subset \mathcal{E}$
with $0 < \mathrm{rk}(F) <
\mathrm{rk}(\mathcal{E})$,
$$
\frac{P_F(m)}{\mathrm{rk}(F)}
\leq
\frac{P_{\mathcal{E}}(m)}{
\mathrm{rk}(\mathcal{E})}
\qquad \text{for } m \gg 0.
$$
It is {\it Gieseker stable}
if we replace $\leq$ by $<$.
Here $P_{\mathcal{E}}(m)$ is the
Hilbert polynomial of $\mathcal{E}$.
\end{definition}

\begin{definition}
\label{algebraic-geometry-definition-2-gieseker-stability}
Let $\mathcal{E}$ be a torsion free
sheaf.
We say that $\mathcal{E}$ is
{\it 2-Gieseker semistable} if
for every subsheaf $F \subset \mathcal{E}$
with $0 < \mathrm{rk}(F) <
\mathrm{rk}(\mathcal{E})$,
$$
\frac{P_F(m)-\chi(F)}{\mathrm{rk}(F)}
\leq
\frac{P_{\mathcal{E}}(m)-\chi(\mathcal{E})}
{\mathrm{rk}(\mathcal{E})}
\qquad \text{for } m \gg 0.
$$
It is {\it 2-Gieseker stable}
if we replace $\leq$ by $<$.
\end{definition}

\noindent
The Gieseker chamber is the region
where this is the relevant stability.
The PT chamber is the neighboring region
where the PT condition is the relevant
one.

\begin{definition}[PT stability]
\label{algebraic-geometry-definition-pt-stability}
\end{definition}

\begin{theorem}
\label{algebraic-geometry-theorem-pt-stable-triple}
Assume $\gcd(\mathrm{ch}_0,
\mathrm{ch}_1,\mathrm{ch}_2)=1$.
Then every PT stable object comes from
a stable triple $(B,P,\eta)$:
\begin{enumerate}
\item $B$ is a 2-Gieseker
      semistable sheaf with
      $\mathrm{hd}(B)\geq 1$,
\item $P$ is a $0$-dimensional sheaf,
\item $\eta \in
      \Hom_{D^b(X)}(P,B[2])=
      \Ext^2(P,B)$ is nontrivial.
\end{enumerate}
Moreover,
$$
A_\eta := \operatorname{cone}(\eta)[-1].
$$
\end{theorem}

\begin{remark}
\label{algebraic-geometry-remark-wall-crossing-triples}
Stable triples generalize rank $1$
PT stable points.
\end{remark}

\begin{theorem}[Jardim et al.]
\label{algebraic-geometry-theorem-single-wall-crossing}
There is a single wall crossing
that takes the DT moduli space to
the PT moduli space.
Equivalently, the Gieseker and PT
chambers are separated by one wall.
\end{theorem}

\begin{proof}
\begin{enumerate}
\item Understand the stable
objects at the boundary,
labelled $\Theta_v^-$
in the picture.
\item Study the asymptotic stable
objects, that is,
what you find if you go to
infinity.
\item Identify what the PT
chamber is.
This may also be an asymptotic
analysis.
\item Show that there are no other
walls between the $\Theta$ wall
and infinity on the Gieseker side,
and no other walls in the
direction of the PT side.
\end{enumerate}
\end{proof}

\medskip\noindent
History.
\begin{itemize}
\item Pandharipande and Thomas
(2009) suggested that DT/PT
correspondence could be a wall
crossing phenomenon in
Bridgeland stability space.
\item Bayer (2009) established
that the rank $1$ case is a
wall crossing for polynomial
stability conditions.
\item Toda.
\item \dots
\end{itemize}

\medskip\noindent
Case studies.

\noindent
Let $X$ be a smooth threefold with
Picard rank $1$.
\begin{itemize}
\item Collapsing wall.
For
$v=(r,0,0,-n)$,
$$
\mathcal{G}(v)=
\mathrm{Quot}^n(
\mathcal{O}_X^{\oplus r})/\!/
\mathrm{GL}(r),
$$
while $\mathcal{T}(v)=\emptyset$.
\item Fake wall.
For $v=(v_0,v_1,v_2,v_{3,\max})$,
$\mathcal{G}(v)=\mathcal{T}(v)$.
This is a wall that changes nothing.
\end{itemize}

\noindent
Now set $X=\mathbb{P}^3$.
\begin{itemize}
\item For
$v=(2,-1,-\tfrac{1}{2},-\tfrac{1}{6})$,
$\mathcal{G}(v)$ is birational to a
$\mathbb{P}^1$-fibration over
$\mathbb{P}^3\times\mathbb{P}^3$,
while $\mathcal{T}(v)\cong\mathbb{P}^3$.
The wall crossing induces a
retraction from the $6$-dimensional
space $\mathcal{G}(v)$ to
$\mathbb{P}^3\times\mathbb{P}^3$,
and
$\mathcal{T}(v)\simeq \mathbb{P}^3$.
\end{itemize}

\medskip\noindent
There is a map
$$
\xymatrix{
\mathcal{G}(-v)\ar[dr]^\gamma
&&
\mathcal{T}(v)\ar[dl]^\tau
\\
&
\mathcal{M}_{\bar\beta,\bar\alpha,s}(v)
&
}
$$
Pavel and Tajakka show that
the PT moduli stack and the
Bridgeland moduli stack at the wall
are projective schemes.
So the map above is a map of
schemes.
\noindent
Formula of the map missing.

\medskip\noindent
Some questions:
\begin{itemize}
\item (Existence.) Given a projective
variety $X$, are there stability
conditions on $D^\text{b}(X)$?
Positive answer for surfaces,
Fano 3folds with Picard rank 1,
abelian 3folds, quintic 3folds,
products of curves, $\mathbb{P}^n$,
\dots
\item (Moduli spaces.) Is
$M_\sigma(v)$ a projective scheme?
Cannot use usual git techniques to
study. Yes for K3, $\mathbb{P}^2$,
and $\mathbb{P}^1\times \mathbb{P}^1$.
The parameter space is a stack
(ref. missing, involves Toda).
\item What do stable objects
look like?
\end{itemize}

\medskip\noindent
\medskip\noindent
Following are very sketchy notes picked
up along the way; don't hessitate to skip

\begin{enumerate}
\item Stable objects in an abelian category
are the ``building blocks":
we can reconstruct the whole category from
them.
\item  An abelian subcategory (hart)
$\subset$ a
triangulated
\item Stability defined via stability f
unction on $\mathcal{A}$.
\item (Example.) $\mathcal{A}=\text{Coh}X$ is
heart
of $D^b(X)$ with stability function.
\item Bridgeland Stabl:= the stability
conditions are a complex manifold
of complex codimension $\text{rk}\Lambda$:
$$
\mathcal{Z}:\text{Stab}(\mathcal{T})
\longrightarrow
\text{Hom}(\Lambda,\mathbb{C})
$$
\item There's a chamber structure;
moduli space changes across chambers.
\item Picture: blue + black are walls.
Q. What are $\beta$ and $\alpha$?
\item thm: bridgeland stable
= gieseker stable ?
\item Q. slope stability
= bridgeland stability?
A. Not always.
\item DT/PT correspondance:
only one wall between PT and G chambers
\medskip
\item Polynomial stability function.
This is an asymptotic version of BS.
\item There are some $\rho$'s.
Arrangements of $\rho_i$ are polynomial
stability conditions on a threefold.
\item Pata-Thomas introduced stability
for rank 1 objects.
Bayer compares the---wall. Q.
Same for Bridge S---only one wall?
\medskip
\item Def. A {\it stable triple}: when
$\text{gcd}(\text{ch}_0,\text{ch}_2,\text{ch}
_3)=1$,
every PT stable object
comes from three conditions (missing).
\item What happens when
you cross the blue wall?
Both $\mathcal{G}$ and $\mathcal{T}$ are
projective.
\medskip
\item For $X$ smooth threefold with
$\text{rk}\text{Pic}=1$, $\mathcal{G}(v)$,
$v=(r,0,0,-n)$, $\mathcal{G}(v)$ is a known
sheaf object and
$\mathcal{T}=\emptyset$.
\item $X$ sm 3 rkpic1, Fake wall;
$\mathcal{G}=\mathcal{T}$.
\item (Extra.) red circle is a wall for a
weaker form of stability.
\end{enumerate}

\begin{definition}[Talk at impa]
\label{algebraic-geometry-definition-slope-stability}
A rank-2 sheaf $\mathcal{F}$ is
{\it semistable (stable)} if
$H^{0}(\mathcal{F}(t))=0$ for $-t \geq(>)
\frac{c_1F}{2}$
\end{definition}

\noindent
Compare with

\begin{definition}[moduli-curves.tex]
\label{algebraic-geometry-definition-semistable}
Let $f : X \to S$ be a family of curves.
We say $f$ is a {\it semistable
family of curves} if
\begin{enumerate}
\item $X \to S$ is a prestable family of
curves,
and
\item $X_s$ has genus $\geq 1$ and
does not have a rational tail for all $s \in
S$.
\end{enumerate}
\end{definition}

\section{Instantons}
\label{algebraic-geometry-section-instantons}


\begin{itemize}
\item The twistor diagram.
$$
\xymatrix{
&  \ell_\infty \ar@{.>}[dr]\ar[dl]\\
\mathbb{C}P^{2} \ar@{^{(}->}[rr]& &
\mathbb{C}P^{3}\ar[dd]
_{\substack{\text{twistor}
\\ \text{map}}}\\ \\
\mathbb{C}^2\ar@{^{(}->}[u]
\mathbb{R}^4 \ar[u]\ar[rr]& &S^4
}
$$

\item Instantons are solutions to some
Yang-Mills solution.
\item Expository paper
of Donaldson arXiv:2205.08639 \item ADHM
construction, 1978. The first appearance of
Algebraic Geometry in Mathematical Physics.
See Hitchin-Kobayashi correspondence.
\item Donaldson: ``unashamedly
computational''.
\item Expository paper by Simon.
\item Take bundle $(E,\nabla)$ with
an anti-self dual
(ADS) connection on $S^4$ and pullback to
$\mathbb{C}P^{3}$ via the twistor map

$$
\tau[x:y:z:w]=[x+jy:z+jw] \qquad \text{note:
$\tau^{-1}(p)=\mathbb{C}P^{1}$}
$$
Then:
\begin{itemize}
\item Restriction to fibres are trivial.
\item Invariant under anti holomorphic
involution (check this formula!)
$[x:y:z:w] \mapsto [-y:x:-w:z]$.
\item $\mathcal{E}$ is also an instanton
bundle.
\item Penrose Transform:
 $H^1(\mathcal{E}(-2)) \cong \Ker \Delta=0$
 where
$\Delta$ is a Laplacian.
\item Definition of instanton sheaf
on $\mathbb{P}^3$ via $c_1(E)=0$ and some
vanishing of cohomologies.
\item Passage from
[differential equations? algebraic geometry?]
to linear algebra: via {\it monads}.
The point is that instanton sheaf is
equivalent to  ``$E$ being the cohomology
of a linear monad''; theorem by Horps
in the 60's, and is the main tool used
by ADHM. Indeed, ADHM equations come from
the cohomology sequences of the so-called
monads.
\item Mathamatical inst bund:=
locally free instanton sheaves.
\end{itemize}
\end{itemize}

\noindent
{\bf Properties.}
\begin{itemize}
\item The only instanton of rank 1 on
$\mathbb{P}^3$ is
the structure sheaf.
\item non-trivial rank 2 locally free
instanton sheaf is  $\mu_0$ stable
\item double dual is locally free and
also instanton
\item non-trivial rank 2 instanton sheaf
is Fieseker stable
therefore it makes sense to define
moduli space of instanton sheaves as
an open subset of
$\mathcal{M}(c)=\mathcal{G}(k,0,2,0)$.
\end{itemize}

Then studied the irreducibility (Tikhomirov)
and smoothness (Jardim-Verbitsky, 2014. Uses
``3rd hyperkähler quotient'') of
$\mathcal{I}(c)$, the moduli space of rank 2
locally free instanton sheaves of charge $c.$
But nobody likes this results; want new
proofs.

In contrast, $\mathcal{M}(c)$ of rank
2-instanton sheaves of charge $c$ is not
irreducible in general! $\mathcal{M}(1)$ and
$\mathcal{M}(2)$ are irreducible,
$\mathcal{M}(3)$ has exactly 2 irreducible
components of dimension 21; $\mathcal{M}(4)$
has 4 irredicuble components: the locally
free is irreducible, and the other 3 that
intersect the closure of the locally free,
$\overline{\mathcal{I}(c)}$. 3 components of
dimension 29 and one of dimension 32

\begin{remark}
\label{algebraic-geometry-remark-gieseker}
In general the $\mu$ moduli space is
not projective, but the Gieseker is.
\end{remark}

\noindent
{\bf Is $\mathcal{M}(c)$ connected?} Use
$\mathbb{C}^*$ action. True for $c \leq 4$;
every component intersects
$\overline{\mathcal{I}(c)}$ in this range!

\begin{definition}[Elementary transformation]
\label{algebraic-geometry-definition-elementary-transformation}
$F$ of rank 2 locally free instanton,
$Q$ of rank 0 instanton with 1 dimensional
sheaf $h^p Q(-2)$ for $p = 0,1$.
So in this conditions if we have an
epimorphism
 $$
F \overset{\varphi,\text{ surj}}{\to}\} Q
$$
we get that $\Ker \varphi$ is an instanton.
\end{definition}

So  we might be interested in
$$
E \hookrightarrow E^{*
*}\xrightarrow{\text{surj.}} Q_E
$$
Let's have a look at
$\mathcal{M}(3)=\overline{\mathcal{I}(3)}\cup
\overline{C(0,1,3}$. Consider a generic line
bundle of degree 0 on a planar cubic (which
is encoded in that we have intersection of
something of cimension 1 and something of
simension 3), $L inn \text{Pic}^0(C)$ so we
have an epimorphism
$$
\xymatrix{
0\ar[r]&E\ar[r]&\mathcal{O}_{\mathbb{P}^3}
\oplus
\mathcal{O}_{\mathbb{P}^3}\ar[r]&(i_*
L)\ar[r]&0
}
$$
yielding a bundle $E$ as previously outlined.

So we are studying the components via
pushforwards of sheaves on complete
intersection curves inside $\mathbb{P}^3$

``when we do alementary transformation of
rational (not rationl?) we get
something on the boundary of locally free.''

\medskip\noindent
{\bf Perverse instanton sheaves.} Like an
instnaton sheaf in monad description but with
some restrictions on the cohomologies. This
leads to definition of {\it $0$-rimensional
instanton}, a perverse instanton such that
$\mathcal{H}^0=0$.

\medskip\noindent
{\bf Instantons and quivers.}

\medskip\noindent
{\bf Framed instantons.} Fix a line $j:L\to
\mathbb{P}^3$ and $E$ a perverse
sheaf; a {\it framing} at $L$ is an
isomorphis [missing].

Apply GIT to the ADHM data to construct a
moduli space
$$
\mathcal{P}(r,c)=\mathcal{V}(r,c)^{\text{st}}
/\!/\text{GL}(V_c)
$$
and it will follow that $\mathbb{P}(?,?)$ is
connected.

\medskip\noindent

Considerations on quaternionic spaces lead to
generalization of what has been discussed so
far to higher dimensions. I.e. an {\it
instanton sheaf} on $\mathbb{P}^n$ is… [other
cohomological conditions] sheaf

Why are instantons interesting? They are the
simplest; may provide examples for Bridgeland
stability.

In Kuznetsov (2012) and Faenzi (2014)
introduced {\it rank 2 instanton bundles on
Fano 3-folds}. An {\it instanton bundle} on
$X$ is a $\mu$-stable … and some Chern class
is called the {\it charge}. An {\it instanton
sheaf} (introduced by Marcos-Gaia) is ….

 ``Since we are imposing $\mu$-stability on
the defintion we can consider the moduli
$\mathcal{I}_X(c) \subset
\mathcal{G}_X(2,-r_X,c,0)$''.

There's also monad representations as an
ingredient.

\medskip\noindent
Here are two questions that invite us to join
the instanton fever:

\medskip\noindent
{\bf Task 1.} Construct rank 2-instanton
sheaves that do not deform into locally free
ones, and obtain the new irreducible
components of $\mathcal{G}(2,-r_X,c,0)$.

\medskip\noindent
{\bf Task 2.} Nonlocally instanton sheaves
that can be deformed into non-locally free
ones: the {\it instanton boundary}
$\overline{\mathcal{I}(c)}/\mathcal{I}(X)
$[formula
right?]


\medskip\noindent
{\bf Recipe to construct your own instanton.}
\begin{enumerate}
\item (Make a bunch of instantons.) Find an
appropriate curve to
use Serre correspondence to find some
rank 2 instanton sheaf:
$$
\xymatrix{
0\ar[r]&\mathcal{O}_{\mathbb{P}^3}(-1)\ar[r]&
\underbrace{E}_{\exists
}\ar[r]&\mathcal{I}_{\sqcup
\text{lines}}\ar[r]&0
}
$$
where $E$ is a locally free instanton
of charge = the number of lines $-1$.

\item (Is your family of instantons generic?)
You look at the $\text{Ext}$s.
``Therefore, the family of instantons only
defines a locally closed subset
within a generically smooth irreducible
component of $\mathcal{G}$''.

\item ``Find suitable rank 0 instanton
sheaves
 to perform an elementary
transformations on the examples obtained in
Step 1.'' But they are non-locally
free.

\item Now I have my instantons, I know they
are locally free: but how to prove
that the elementary transformations deform
to locally free ones? Looks like
the challenge is to prove that the
deformation is locally free.
\end{enumerate}

In the papers by the group there are several
particular cases when the
deformations are locally free. But they don't
have a general result that would
work for 3-folds.

\medskip\noindent
{\bf What you need to call a thing an
instanton.}
\begin{itemize}
\item Minimal cohomology possible;
try to kill as much cohomology as you can.
\item Fixing $c_1$ (which may determine other
Chern classes).
\item Some stability condition
like $\mu$-stability or quiver stability.
Here is an example that is not
$\mu$-semistable:
$T\mathbb{P}^3(-1)\oplus\Omega_{\mathbb{P}^3}
(1)$
\item Whenever possible,
look for a monadic representation.
(The monadic
representation comes from ADHM ---
the beginnings of this theory.
And it's still here!)
\end{itemize}

\section{Jardim-Muniz}
\label{algebraic-geometry-section-jardim-muniz}

\noindent
Reflexive sheaves
are vector bundles except for a small
locus.

For the definition and basic properties,
see Definition \ref{algebraic-geometry-definition-reflexive}
and Lemma \ref{algebraic-geometry-lemma-reflexive-torsion-free}
in Section
\ref{algebraic-geometry-section-coherent-reflexive-sheaves}.

\begin{remark}[Talk at IMPA, 11 June]
\label{algebraic-geometry-remark-reflexive-talk}
Torsion could also be defined so that the
sheaf can inject onto its dual. In this talk
we discussed the moduli space of
reflexive/torsion-free sheaves, which turned
out to be parametrized by $c_1, c_2$ and
$c_3$. This was denoted by $R(c_1,c_2,c_3)$.
Actually I think it may have been Manolache
that proved the existence of this moduli
space.

Alan Muniz

Nesta palestra discutiremos a classificação
de feixes reflexivos de posto dois e seus
espaços de módulos. Apresentaremos algumas
ferramentas básicas usadas na construção e
determinação de tais feixes. Aplicaremos
estas técnicas para o caso de feixes com
segunda classe de Chern igual a quatro,
obtido recentemente em colaboração com Marcos
Jardim.
\end{remark}




\medskip\noindent
Now let's discuss
distributions on manifolds.

Here's the abstract from a talk by Marcos
Jardim at Geometric Structures:

``I will revise the work done over the past
10 years with various collaborators on
distributions and foliations on 3-folds,
especially on the projective space, with a
focus on properties of the tangent sheaf and
singular scheme."

Here are two key ideas: if the distribution
is codimension 1 we can write:
$$
\xymatrix{
0\ar[r]&F\ar[r]&TX\ar[r]^{\omega}&I_Z \otimes
L\ar[r]&0
}
$$
where $L$ is a line bundle and
$\omega \in H^{0}(\Omega_X \otimes L)$, and
$Z=\{p:\omega(p)=0\}$.

When codimension is 2 then $\mathcal{D}$ is
given by
a holomorphic vector field
$\nu$: $T_p=\left\langle{}\nu(p)\right\rangle{}$.

It can be encoded as an exact sequence
$$
\xymatrix{
	0\ar[r]&L\ar[r]^{\nu}&TX\ar[r]&N\ar[r]&0
}
$$
where $L$ is a line bundle
and $\nu \in H^{0}(TX \otimes L^\vee)$;
  $Z=\{p | \nu(p) = 0\}$.

\begin{remark}
\label{algebraic-geometry-remark-stauration}
Saturation means that $Z \subset X$ is a
union
of curves and points.
\end{remark}

And again, distributions are parametrized by
Chern classes.

Two interesting open questions:
\begin{enumerate}
\item {\bf Conjeture.} if $\mathcal{D}$ is a
codimension 1 foliation of degree
$d$ on $\mathbb{P}^3$, then $c_2(F)\leq
d^2-d+1$ and
bound is attained
by rational foliations of type $(1,d+1)$.
(True for $d \leq 2$.)
\item {\bf Conjecture (with Pepe Seade).}
$\mathcal{D}$ is a codimension 1
foliation on a smooth projective 3-fold,
then $\text{Sing}\mathcal{D}$ is connected.
\end{enumerate}

\begin{theorem}[Jardim-Muniz]
\label{algebraic-geometry-theorem-jardim-muniz}
Conditions on Chern classes used to
understand moduli space $R(c_1,c_2,c_3)$.
$c_2=4$ gives (?). For $c_3\leq 6$, possible
``spectrum" exists…
\end{theorem}









\section{Campana}
\label{algebraic-geometry-section-campana}

A \textit{special} variety should be thought
of
as the farthest notion possible from being of
general type.

\begin{definition}[Campana]
\label{algebraic-geometry-definition-special-variety}
$X$ complex projective.
$X$ is \textit{special} if $\forall i \leq  p
\leq \dim X$,
for all line bundle $L \subset \Omega^P_X$,
$\kappa(L)<P$ (Kodaira dimension).
\end{definition}

\noindent
A conjecture by Campana says that
$X$ is special if and only if there exists
$f:\mathbb{C}\to X$
Zariski dense.
Compare this with Kobayashi hyperbolicity!
There are no lines on …?

The main tool for proving this kind of
``Bloch-Ochai" statements (i.e.
something of the kind: if $h^0(\Omega_X)>
\dim X$
then there are no nonconstant
$f:\mathbb{C}\to X$
Zariski dense)
is by passing through the Albanese variety.

In a talk [fill details]
we study a theorem of this kind but with the
bound
given by $q^+(X)$ defined as the supremum of
the $h^0(\hat{X},\pi^*\Omega_{X}$
for $\hat{X}\to X$ some (étale) cover.
Can we prove a BO stement with the condition
$q^+(X)>\dim X$?

Notion of special comes in: indeed
 $$
q^+(X)>\dim X \implies X \text{ not special
}\implies
\not\exists f:\mathbb{C}\to X \text{Zariski
dense}.
$$

The following theorem is not at all easy to
prove:
\begin{theorem}[Campana]
\label{algebraic-geometry-theorem-campana}
$\hat{X}\xrightarrow{\text{étale} }X$, then
$\hat{X}$ is special iff $X$ is special.
\end{theorem}

\noindent
My question: why does the abelianess of
the Albanese variety is important in this
construction?





% --- End algebraic-geometry.tex ---

% --- Begin stanley-reisner.tex ---
\chapter{Stanley-Reisner schemes}
\label{chapter-stanley-reisner}





\section{Stanley-Reisner schemes}
\label{stanley-reisner-section-sr}

\noindent
Here I follow \cite{luk1} for the basic definitions; these
are consistent with other sources.

Let $[n]:=\{0,\ldots,n\}$ be the set of all
positive integers from 0 to $n$, and let
$\Delta_n$ denote the power set of $[n]$. A
subset $\mathcal{K} \subseteq \Delta_n$ is
called a {\it simplicial complex} if for
every $f \in \mathcal{K}$ all the subsets of
$f$ are also in $\mathcal{K}$. Elements of
$\mathcal{K}$ are called {\it faces}.

Let $S=k[x_0,\ldots,x_n]$ be the polynomial ring in $n+1$ variables
over an algebraically closed field $k$.
If $a=\{i_1,\ldots,i_k\}\in \Delta_n$ we write $x_a \in S$ 
for the square-free monomial $x_{i_1}\cdot\ldots\cdot x_{i_k}$.
The {\it Stanley-Reisner} ideal
of a simplicial complex $\mathcal{K}$ is
$$
I_{\mathcal{K}}=\langle x_p : p \in \Delta_n\setminus \mathcal{K} \rangle
\subseteq S.
$$
The {\it Stanley-Reisner ring} of $\mathcal{K}$ is
$A_{\mathcal{K}}=S/I_{\mathcal{K}}$
and the {\it Stanley-Reisner scheme} is $\text{Proj} A_{\mathcal{K}}$.

\section{Stanley-Reisner scheme of Grünbaum sphere and a determinantal CY3}
\label{stanley-reisner-section-grunbaum-sphere}

\noindent
A Gröbner basis for the Stanley-Reisner ideal
of Grünbaum-Shreedharan's strange complex
$\mathcal{M}$ is 
$$
\begin{aligned}
x_6x_7x_8,
x_4x_6x_8, x_3x_7x_8, x_3x_5x_7, x_3x_4x_8,
x_2x_7x_8, \\ 
x_2x_5x_7, x_2x_5x_6,
x_2x_4x_7, x_2x_4x_6, x_1x_4x_6,\\ 
x_1x_4x_5,
x_1x_3x_8, x_1x_3x_6, x_1x_3x_5, x_1x_2x_5
\end{aligned}
$$
as computed in gfan and also GAP, and also as
shown in \cite{luk1}.

\section{GN variety}
\label{stanley-reisner-section-gn}

\noindent
There is a Calabi-Yau threefold in
$\mathbb{P}^7$ of degree 20 that could be the
smoothing of $\text{SR}(\mathcal{M})$. (As
far as I know there could exist others,
though they aren't very well-known.) It was
first studied by Gulliksen-Neg\aa ard
\cite{Gulliksen-Negard}. Among the papers I
read often for this project, it is mentioned
in \cite[Table 1 and last
paragraph]{CGKK-Aritmethically-Gorenstein-CY-in-P7}
and  \cite{cascade}. Additionally, this
variety has been studied in mirror-symmetry
context, see \cite{Jockers}.

In trying to immitate the computations I did
for the twisted cubic, I would like to find
the equations of the GN variety. In
\cite{cascade} we see its description as the
``$3\times 3$ minors of $4\times 4$ matrix
with linear forms''.

Back in May 19, 2025 I computed what might be ideal of the GN variety.
That is, as the $3 \times 3$ minors of the matrix
\begin{equation}
\label{stanley-reisner-equation-matrix}
\begin{bmatrix}
x_1 & x_2 & x_3 & x_4 \\
x_5 & x_6 & x_7 & x_8 \\
x_2 & x_3 & x_4 & x_5 \\
x_6 & x_7 & x_8 & x_1
\end{bmatrix}
\end{equation}
Those minors are
\[
\begin{aligned}
&x_3^2 x_5 - x_2 x_4 x_5 - x_2 x_3 x_6 + x_1 x_4 x_6 + x_2^2 x_7 - x_1 x_3 x_7, \\
&-x_3 x_6^2 + x_3 x_5 x_7 + x_2 x_6 x_7 - x_1 x_7^2 - x_2 x_5 x_8 + x_1 x_6 x_8, \\
&-x_3^2 x_6 + x_2 x_4 x_6 + x_2 x_3 x_7 - x_1 x_4 x_7 - x_2^2 x_8 + x_1 x_3 x_8, \\
&x_4 x_6^2 - x_4 x_5 x_7 - x_3 x_6 x_7 + x_2 x_7^2 + x_3 x_5 x_8 - x_2 x_6 x_8, \\
&x_3 x_4 x_5 - x_2 x_5^2 - x_2 x_4 x_6 + x_1 x_5 x_6 + x_2^2 x_8 - x_1 x_3 x_8, \\
&-x_1 x_2 x_5 + x_1^2 x_6 - x_4 x_6^2 + x_4 x_5 x_7 + x_2 x_6 x_8 - x_1 x_7 x_8, \\
&-x_1 x_2^2 + x_1^2 x_3 - x_3 x_4 x_6 + x_2 x_5 x_6 + x_2 x_4 x_7 - x_1 x_5 x_7, \\
&x_1 x_3 x_5 - x_1 x_2 x_6 + x_5 x_6^2 - x_5^2 x_7 - x_3 x_6 x_8 + x_2 x_7 x_8, \\
&x_4^2 x_5 - x_3 x_5^2 - x_2 x_4 x_7 + x_1 x_5 x_7 + x_2 x_3 x_8 - x_1 x_4 x_8, \\
&-x_1 x_3 x_5 + x_1^2 x_7 - x_4 x_6 x_7 + x_4 x_5 x_8 + x_3 x_6 x_8 - x_1 x_8^2, \\
&-x_1 x_2 x_3 + x_1^2 x_4 - x_4^2 x_6 + x_3 x_5 x_6 + x_2 x_4 x_8 - x_1 x_5 x_8, \\
&x_1 x_4 x_5 - x_1 x_2 x_7 + x_5 x_6 x_7 - x_5^2 x_8 - x_4 x_6 x_8 + x_2 x_8^2, \\
&x_4^2 x_6 - x_3 x_5 x_6 - x_3 x_4 x_7 + x_2 x_5 x_7 + x_3^2 x_8 - x_2 x_4 x_8, \\
&-x_1 x_3 x_6 + x_1 x_2 x_7 - x_4 x_7^2 + x_4 x_6 x_8 + x_3 x_7 x_8 - x_2 x_8^2, \\
&-x_1 x_3^2 + x_1 x_2 x_4 - x_4^2 x_7 + x_3 x_5 x_7 + x_3 x_4 x_8 - x_2 x_5 x_8, \\
&x_1 x_4 x_6 - x_1 x_3 x_7 + x_5 x_7^2 - x_5 x_6 x_8 - x_4 x_7 x_8 + x_3 x_8^2
\end{aligned}
\]
Now in theory I would just do all I did for the twisted cubic,
but now for GN, find a nice weight vector in for which the 
initial ideal is the ideal of $\text{SR}(\mathcal{M})$, and win.

But several problems arise:
\begin{enumerate}
\item First I have to double check that ideal for GN. I'm not
sure if those $3 \times 3$ minors of the $4 \times 4$ matrix I used are
good. {\bf Edit:} I think it's right. At least
by skimming through my references on GN, which are:
\cite{Jockers, Jockers2, Bertin, Gross-moduli-abelian}, I don't find
a matrix exactly but the claim is that a ``generic'' matrix
of $4 \times 4$ with linear forms in the entries should do the job.
(So that ChatGPT matrix above has indeed linear forms as entries.)
But most importantly \cite{Bertin} put a betti tally like
the ones one computes in Macaulay2 and indeed it's the same
as the one I get.

\item Most importantly, and this is what happened that other time
I tried to do this (when by the way I was being led by
ChatGPT to do exactly what I'm doing now, only now I
know a little more what's going on), is that the Gröbner fan
of GN is not computable. It's too much data and there's no output.
So it's not clear how to find the weight vector.
\end{enumerate}

One thing to try: obtain more freedom
on the $4\times 4$ matrix, see if you can change
that and still find the same good candidate.

\medskip\noindent
A breakthrough discovery:
the minimal resolution of GN and $\text{SR}(\mathcal{M})$ is not the same.

After this observation (and now knowing if
minimal resolutions of fibres of flat maps are isomorphic),
I'm not sure how to proceed. Possible ways are:

\begin{itemize}
\item Keep trying to prove/find out somehow 
(maybe asking: Eduardo Esteves thinks it's not true,
and recommended that I asked someone at URFJ)
whether of not the minimal resolution claim is true.
I may also look for other invariants on the fibers of a flat map,
see Stacks Project 
\href{https://stacks.math.columbia.edu/tag/02NM}{02NM}.

\item Keep trying to find the smoothing,
but I exhausted my computational attempts to find a perfect pairing
between the sets of Gröbner bases,
as well as my attempts to get around the size
of the Gröbner fan and just compute it via gfan.
\end{itemize}

\medskip\noindent
As a note, in \cite{luk2}, where the table
where GN was found lies, it is claimed that
there is  ``evidence'' that this table is
complete. The table counts arithmetically
Gorenstein CY3s. Adding up, either I find a
smoothing that is not the table, or prove
that it doesn't exist.


\section{6-vertex triangulation of RP2}
\label{stanley-reisner-section-6-vertex-triangulation-of-rp2}

\noindent
Here I read \cite{KuhnelBanchoff1983}.

\medskip\noindent
The icosahedron quotiented by antipodal
map gives a triangulation of 6 vertices
of $\mathbb{R}P^2$.
This can be embedded in $\mathbb{R}^5$ 
inside the 5-simplex $\Delta^5$,
which has 6 vertices.
This embedding has the properties
of symmetry (icosahedral$/-\text{Id}$),
duality (three vertices of $\Delta^5$ determine
a triangle of $\mathbb{R}P^2_6$ if and only if
the remaining do not),
tightness (too complicated) and
secant tangency (a chord joining
two points of $\mathbb{R}P^2$ is contained
in a boundary 4-simplex of $\Delta^5$.

Then consider Veronesse embedding
$[x_0:x_1:x_2]\mapsto [x_0^2:x_1^2:x_2^2:x_0x_1:x_0x_2:x_1x_2]$.
Like $\mathbb{R}P^2_6$ is contained in the boundary of $\Delta^5$,
this embedding is contained in $S^4$.
It also satisfies properties of symmetry (of $\mathbb{R}^3$),
duality, tightness and secant-tangency.

Each of the third and fourth properties
characterizes $\mathbb{R}P^2_v$ among
the smooth embeddings of $\mathbb{R}P^2$
in $\mathbb{R}^5$ not lying in an affine hyperplane.
``Pohl and Kuipar have in fact shown that
any {\it topological} embedding of the real
projective plane into $\mathbb{R}^5$ not lying in 
a hyperplane which has the tightness property
alone must either be $\mathbb{R}P^2_v$ or $\mathbb{R}P^2_v$.''

With respect to the complex projective plane,
the situation has been less satisfactory.

\section{9-vertex triangulation of CP2}
\label{stanley-reisner-section-9-vertex-triangulation-of-cp2}



\medskip\noindent
{\bf September 25.}
Summary of Calla's talk.
The objective is to construct a deformation 
of hyperkähler generalized Kummer fourfolds
starting from a special fiber.
The special fiber is the smoothing
of a Stanley-Reisner variety.

Proving this result would 
``imply that a general hyperkähler generalized kummer
fourfold with polarization of degree 2 and fixed divisibility
is contained in 36 quartics''.

\noindent
{\bf Questions.}
\begin{itemize}
\item How did they know that this funny simplicial
complex gives the special fiber of the degeneration
upon smoothing??
\item How does that result about being contained
in 36 quadric follows?
\item How come the special fiber of a degeneration
of fourfolds is a surface?
\end{itemize}

\medskip\noindent
Now let's go over the details of Calla's construction.
In \cite{KuhnelBanchoff1983} it was shown that there is a 9-vertex
triangulation $\Sigma$ of $\mathbb{C}P^2$.
This triangulation has 36 4-simplices.
Notice that the definition of SR scheme for Calla
might be different from ours: the {\it Stanley-Reisner variety}
of $\Sigma$ is the union of 36 copies of $\mathbb{P}^4$,
one for each 4-simplex of $\Sigma$.
Its ideal is generated by 36 quartics.

Part 1. We denote by $\Sigma$ the
9 vertex triangulation of  $\mathbb{P}^2$,
and $\text{SR}_{36}$ its corresponding
Stanley-Reisner scheme.

\begin{enumerate}

\item Blow up 9 coordinate points in $\mathbb{P}^8$ 
$(1:0:\ldots:0),\ldots,(0:\ldots:0:1)$.
(Recall that by the construction of the
Stanley-Reisner scheme associated to
a simplicial complex, in this case with 9
vertices, the vertices of the SR scheme
are precisely the coordinate points of $\mathbb{P}^8$.)
This introduces 9 copies of $\mathbb{P}^7$.
Proper transform of  $\text{SR}_{36}$
intersects $\mathbb{P}^7$ in a union
of  20 $\mathbb{P}^3$s, denoted $T_{20}$.
This does not depend on which copy of $\mathbb{P}^7$ 
we look at.

As a note, recall that the proper
transform of a subvariety $Y \subset X$ 
is defined as the {\it closure} of
$\pi^{-1}(Y\setminus Z)$ where
the blow up occurs at $Z$. In this case
$Z$ is the 9 vertices and $Y=\text{SR}_{36}$.
Calla claims that
that the closure of the proper
transform of $\text{SR}_{36}$ 
intersects each of the exceptional divisors
(one $\mathbb{P}^7$ for every vertex)
in a union of 20 $\mathbb{P}^3$ 
regardless of which exceptional divisor we
look at. We call that intersection $T_{20}$,
and it is inside $\mathbb{P}^7$.

\item Now we turn to modify $T_{20}$ 
to satisfy Friedman conditions.
Blow up 8 coordinate points in $\mathbb{P}^7$.
The proper transform of $T_{20}$ intersects
$\mathbb{P}^6$ in a union of 10 $\mathbb{P}^2$ s,
which we call $S_{10}$.
This {\bf does} depend on which copy of
$\mathbb{P}^7$ and $\mathbb{P}^6$ 
we are looking at!
$S_{10}$ is a Stanley-Reisner variety
and in fact is a triangulation of a sphere.
(A polyhedron!?)

\item Now we proceed to smooth $S_{10}$.
It has seven vertices, which we blow up,
introducing seven copies of $\mathbb{P}^5$.
Each of the $\mathbb{P}^2$ s becomes
a del Pezzo 6.
Looking at the picture we
see that there are vertices with
different valency. For each vertex
where $d$ $\mathbb{P}^2$s meet,
we glue in a del Pezzo surface of 
degree $d$. 
We get $\tilde{S}_{10}$, which is a 
simple normal crossing surface,
with trivial canonical bundle
and satisfies $d$-stability 
Friedman conditions.

Then there exists a Kulikov model
of type III degenerations of K3s
(of degree $10$ in $\mathbb{P}^7$)
with special fiber $\tilde{S}_{10}$.
(This is local smoothing over a
small disc, does not give a projective
family of K3s.)
\end{enumerate}

{\bf Key move.}
$S_{10}$ is a triangulation of a sphere.
Blow up the singularities and look
how does the proper transform of $S_{10}$ 
intersect the exceptional divisors.
To produce a smoothing,
we glue in del Pezzo surfaces
at the exceptional divisors.

Now we shall do the same for $T_{20}$.
(Which is still not the blow up
of $\text{SR}_{36}$, but the intersection
of the proper transform of $\text{SR}_{36}$ 
with any of the 9 exceptional divisors.)
We will glue in our projecive degeneration of K3s.
Notice that Friedman only gave us existence
of local model, so we must find equations.

For now let's just say that there's a bunch of procedures and
the resulting $\tilde{T}_{20}$ is a normal crossing and 
has trivial canonical bundle.
Then we can run smoothing procedure and find equations of family
of CY3s in $\mathbb{P}^7$ by noting that $T_{20} \subset 10 \mathbb{P}^4$s.
We should find generators of the ideal of
the variety $10\mathbb{P}^4$
among minors of $3\times 5$ matrix, similarly to surface case.

\medskip\noindent
{\bf October 19.}
It looks as if Calla was defining the SR variety of the triangulation
in a different way to my definition in Section \ref{stanley-reisner-section-sr}.
Today I checked that the ideal generated by 36 quartics
(one for every facet in the triangulation)
is not the same as the SR ideal.
How to tell which definition is ``correct''?
Correct for what? To answer I must get better understanding
of the deformation process from the Kummer perspective,
but it's not clear where this idea that the SR scheme
corresponding to the 9-vertex triangulation
would give the special fiber of a degeneration of
Kummer fourfolds!

I also noticed today that in \cite{KuhnelBanchoff1983}
there is a mention of $\mathcal{M}$!
I think $\mathcal{M}$ might be the vertex figure of $\mathbb{C}P^2_9$!


\section{Calla's talk}
\label{stanley-reisner-section-calla-talk}

\noindent
In this section I write literally Calla's talk.

\medskip\noindent
Aim: construct a type III
degeneration of varieties of generalised Kummer type 4-folds
with polarisation degree 2,
by starting from a given special fibre and constructing a smoothing.

Corollary. Proving this result will imply that a 
general HK variety of this type with polarisation
degree 2 (\& fixed divisibility) is contained
in 36 quartics.
(In particular, it is not contained in a cubic.)

\medskip\noindent
Reminder of what we mean by type III degen:

Figure with three strings of text:
\begin{itemize}
\item "Moduli space of gen. Kum. type
4-folds w/ fixed polarization" enclosed in a circle
in the center,
\item "Boundary of BB compactification
is made up of 1 dim components 
intersecting in dim 0" with an arrow
pointing at to lines intersecting the circle
of the last item, and
\item "Type III degen has special fibre corresponding
to this point", which has an arrow pointing
at the point of intersection of the last item.
\end{itemize}

\medskip\noindent
Method: 
\begin{itemize}
\item We give a candidate $\text{SR}_{36}$ for
the special fibre of this degen.
\item Since $\text{SR}_{36}$ is badly singular,
we make some modifications,
in order to be able to smoothe it.

\end{itemize}
\section{Moduli space of generalized Kummer type fourfolds}
\label{stanley-reisner-section-moduli-kummer}

\noindent
According to Calla's talk,
the whole point of constructing the smoothing
of $\text{SR}(\mathcal{M})$ is that it
should help us to smooth $\mathbb{C}P^2_9$,
which in turn should lie in the boundary
of the moduli space of generalized Kummer type fourfolds
with polarization degree 2.

Here's a quote from Calla's talk:
``Aim: construct a type III
degeneration of varieties of generalized
Kummer type 4-folds with polarization degree 2,
by starting from a given special fibre
and constructing a smoothing.''

\medskip\noindent
Let's address the following question.
Even if there exists a smoothing of $\text{SR}(\mathbb{C}P^2_9)$,
why would it be a generalized Kummer variety?

First we go over \cite[Theorem 0.2.1]{Fausk}:

\begin{theorem}
\label{stanley-reisner-theorem-fausk-021}
A smoothing, if it exists, of the Stanley-Reisner scheme
of a triangulation of the 3-sphere yields a Calabi-Yau 3-fold.
\end{theorem}

\begin{proof}
Two key steps: show using \cite[Theorem 6.1]{Bayer-Eisenbud}
that the canonical bundle of SR(sphere) is trivial.
Then show using Hartshorne that triviality of
canonical bundle stays the same under deformations.
\end{proof}

\noindent
Sure, that says that if I smooth $\text{SR}(M)$ 
then I have a CY 3fold. That's only the vertex-figure
of $\mathbb{C}P^2_9$.
Still, this looks far from saying that a smoothing
{\it of $\mathbb{C}P^2_9$} is a generalized Kummer 4fold.


\section{Dual complex, Hironaka theorem and Kummer type manifolds}
\label{stanley-reisner-section-dual-complex-hironaka-kummer-type-manifolds}

\noindent
These are notes from 
\href{http://verbit.ru/MATH/TALKS/Dual-complex-ddc_IMPA-2026.pdf}{
a seminar by Misha Verbitsky}
on his work with Nikon K.

\begin{definition}
\label{stanley-reisner-definition-snc}
Let $X \subset Y$ be a divisor in a smooth manifold.
We say that $X$ has normal crossings if
each of its components is smooth, and in a sufficiently
small neighbourhood of every point of $X$ the manifold
 $Y$ has coordinate system such that $X$ is identified with a union of
coordinate hyperplanes.
It has {\it simple normal crossings} if it has normal crossings and 
for each collection of irreducible components of the divisor,
their intersection is connected (possibly empty).
\end{definition}

\begin{theorem}[Hironaka resolution]
\label{stanley-reisner-theorem-hironaka}
Let $X$ be a complex variety.
Then there is a sequence of blow-ups with
smooth centers (=variety we are blowing up)
which results in a smooth resolution
$\pi:\tilde{X} \to X$ and the exceptional set of $\pi$ is SNC.
\end{theorem}

\begin{theorem}[Hironaka resolution of a pair]
\label{stanley-reisner-theorem-hironaka-pair}
Let $X$ be a complex variety and $Y \subset X$ a subvariety.
Then there is a sequence of blow-ups with smooth centers
which results in a smooth resultion
and the exceptional set of $\pi$ is a SNC divisor
and the preimage of $Y$ is also a SNC divisor.
\end{theorem}

\begin{definition}[Stepanov, 2004]
\label{stanley-reisner-definition-dual-complex}
Let $X \subset Y$ be a SNC divisor.
The {\it dual complex} of $X$ is a simplicial complex
whose vertices are the irreducible components $X_1,…X_n$,
and the $k$-simplices correspond to
intersections of $k+1$ of the irreducible components.
A {\it boundary} of a $k$-simplex is a $k-1$-simplex
obtained as the intersection of the divisors except one.
\end{definition}

\begin{theorem}[Stepanov]
Let $x \in X$ be an isolated singularity
and $\tilde{X}$ its smooth resolution such that
the preimage of $x$ is a SNC divisor.
Then the homotopy type of its dual complex is independent from
the choice of a resolution.
\end{theorem}

\noindent
Stepanov's theorem was generalized by
Arapura-Bakhtary-Wlodarczyk.

\begin{theorem}[ABW]
\label{stanley-reisner-theorem-abw}
Let $Y \subset X$ be a subvariety,
and $\tilde{X} \xrightarrow{\pi}X$ the resolution of the pair
$Y \subset X$.
Let $D$ be the exceptional set of $\pi$ (which is a SNC)
and $\tilde{Y}$ the set of all components
$D' \subset D$ such that $\pi(D')\subset Y$.
Then the homotopy type of the dual complex of $\tilde{Y}$
is independent from the choice of resoution.
\end{theorem}

\begin{definition}
\label{stanley-reisner-definition-kummer-type}
We say that $Y \subset X$ is of {\it Kummer type} 
if $b_1(C(X,Y))=0$ where $C(X,Y)$ is the dual complex
of this pair, and $b_1$ the rational Betti number.
\end{definition}

\begin{theorem}
\label{stanley-reisner-theorem-kummer-type}
Let $X \subset M$ be a complex subvariety of Kummer type,
and $G$ a finite group holomorphically acting on $M$
and preserving $X$.
Then $X/G \subset M/G$ is also Kummer type.
\end{theorem}

\medskip\noindent
Now we discuss $dd^c$-lemma and the dual complex.

\begin{theorem}
\label{stanley-reisner-theorem-kummer-type-manifolds}
Let $X \subset M$ be compact complex ($M$ is not necessarily compact).
Assume that $X$ is Fujiki class C
(i.e. bimeromorphic to Kähler) (normal) of Kummer type,
that is, $b_1(S)=0$ where $S$ is the dual complex
of the pair $X \subset M$.
If $\eta$ is an exact $(1,1)$-form,
then there exists a function $f$ defined
in a neighbourhood of $X$ such that $\eta=dd^cf$.
\end{theorem}

\begin{remark}
\label{stanley-reisner-remark-normality}
Normality of $X$ may not be used.
Recall that normal= when all meromorphic functions are bounded.
\end{remark}

\begin{proof}
\noindent
Step 1. Consider the resolution $\tilde{X} \subset \tilde{M}$
With some extra blow ups we may assume that all the 
components of $\tilde{X}$ are Kähler.
Then we apply $dd^c$-lemma on any of the components.

\medskip\noindent
Step 2. $dd^cf=0$ means that it is a sum of holomorphic
and antiholomorphic,
which means that $\ker dd^c$ are
constant functions. There's also maximum princple
which implies that $f$ restricted to $\pi^{-1}(x)$ for any
$x$ in $X$ is constant.

\medskip\noindent
Step 3. Choose a function $f_i$ which satisfies
$dd^c f_i=\eta|_{D_i}$ on each of these components.
Then we prove that this satisfies the cocycle condition,
and we thus from $\eta$ we obtain a canonical element 
$\mu_\eta \in H^1(S,\mathbb{C})$.
[I think this cocycle condition is used
to glue the $f_i$ to a global function].

\medskip\noindent
Step 4. Extending $f$ to a neighbourhood…
\end{proof}

\noindent
Then there is a statement about the $dd^c$ lemma not
holding if $X$ is not of Kummer type: there is a form
that is not $dd^c$-exact:

\begin{theorem}
\label{stanley-reisner-theorem-ddc-counterexample}
Let $X \subset M$be a compact, positive-dimensional
subvariety of a manifold $M$, which is not of Kummer type.
Assume that there exists a non-contractible path in the
dual complex of the resolution of the pair, $S$,
which visits a component $D_0 \subset \tilde{X}$
only once, and the projection of $D_0$ to $M$
is positive-dimensional.
Assume $\dim M=2$.
Then there exists an exact $(1,1)$-form $\eta$
on a neighbourhood of $X$ such that $\eta \neq  dd^cf$ for any 
function $f$ on $X$.
\end{theorem}

\medskip\noindent
For the proof of ABW theorem we need the following results,
which were proved in 2002 in D. Abramovish, K. Karu, K. Matsuky,
J. Wlodarczyk, Torication and factorization of birational maps.

\begin{theorem}[Weak factorization theorem]
\label{stanley-reisner-theorem-weak-factorization-theorem}
Let $\varphi: X \dashrightarrow X'$ be a bimeromorphic map
of compact complex manifolds and $U \subset X$, $U' \subset X'$
the complements to the exceptional sets.
Then $\varphi$ can be decomposed 
into a sequence of blow-ups and blow-downs with
smooth centers.
\end{theorem}

\begin{remark}
\label{stanley-reisner-remark-meromorphic}
So what does it mean to be meromorphic? It means that there is
$$
\xymatrix{
&  Z \subset X \times X'\ar[dl]\ar[dr]\\
X & & X'
}
$$
where projections are proper, dominant and 1-to-1. [Maybe something
else I couldn't write down.
\end{remark}

\noindent
We need a version with boundary:

\begin{theorem}
\label{stanley-reisner-theorem-weak-factorization-for-pairs}
Let $\varphi:X \dashrightarrow X'$ be a bimeromorphic map
of compact complex manifolds, and $U \subset X$, $U' \subset X'$
be complements to the exceptional sets.
Consider SNC divisors $D \subset X$, $D' \subset X'$
and let $D_1,…, D_n$ and $D_1',…,D_n'$
be their irreducible components.
Assume that $\varphi$ takes $X\setminus D_i$
to $X' \setminus D_i'$
and this bimeromorphic map is proper
(that is, the closure of its grap in 
$X \setminus D_i \times X'\setminus D'$
projects properly to the components $X \setminus D_i$
and $X'\setminus D_i$).
Then $\varphi$ can be decomposed into a sequence
of blow-ups and blow-downs with smooth
centers, and the exceptional sets of blow-ups
have SNC with the proper preimages of $\cup  D_i$.
\end{theorem}

\noindent
Now we can prove Theorem \ref{stanley-reisner-theorem-abw},
which says that the homotopy type of the dual complex
is independent from the resolution.

\begin{proof}[Proof of Theorem \ref{stanley-reisner-theorem-abw}]
\noindent
Step 1. Take to resolutions of the pair.
Then they are bimeromorphic.
If we can show that the homotopy type of the dual
complex stays the same in each of the steps of the 
weak factorization theorem, then we are done.

\medskip\noindent
Step 2. We analyse the different possible cases of center in the 
first step. The first case is easy: the dual complex remains the same.
The second case I didn't understand.
The third case boils down to gluing a simplex to the 
original dual complex.
\end{proof}

\noindent
Then there is a $G$-equivariant version of this theorem
and of the weak factorization theorem.

\medskip\noindent
Finally we sketch Theorem \ref{stanley-reisner-theorem-kummer-type}.
Using $G$-invariant Hironaka, we
resolve, take quotient, take product, resolve, etc…
This would give a sequence of dual complexes…

\medskip\noindent
Take Hilbert scheme of torus, projected to symmetric
product of elliptic curve…
$$
\xymatrix{
(T^n)^{[n]}\ar[d]\\ \text{Sym}E.
}
$$


\section{Dual complex of a degeneration}
\label{stanley-reisner-section-dual-complex}

\noindent 
Here I review \cite{KLSV} to
understand the relation between
$\mathbb{C}P_9^2$ and the moduli of
generalized Kummer varieties.

``It turns out that the right class of
singularities that still allow the definition
of a meaningful dual complex is dlt. In other
words, the correct context for defining and
intrincis dual complex associated to a
degeneration is that of minimal dlt
degeneration. The minimal dlt model
$\mathcal{X}/\Delta$ is not unique, but
changing the model has no effect on
$|\Sigma|$.''

``Relevant for us is the fact that associated
to a $K$-trivial [I guees, trivial canonical
bundle] degenertion $\mathcal{X}/\Delta$
there is an intrinsic (depending only on
$\mathcal{X}^* /\Delta^*$) topological space
that we call (following [MN15]) {\it the
essential skeleton} $\text{Sk}(X)$ associated
to the degeneration. For a minimal dlf
degeneration of $K$-trivial varieties, the
essential skeleton $\text{Sk}(\mathcal{X})$
can be identified with the topological
realization $|\Sigma|$ of the dual complex
(cf Nicaise-Xu).''

``The purpose of this section is to make some
remarks on the structure of the essential
skeleton $\text{Sk}(\mathcal{X})$ for a
degeneration of hyper-Kähler manifolds''.

``[…] we note that the
$\text{Sk}(\mathcal{X})$ for a MUM
degeneration of Calaby-Yau $n$-folds is
predicted to be the sphere $S^n$. This is
true in dimension 2 by Kulikov's theorem, and
in dimensions 3 (unconditional) and 4
(assuming additionally that degeneration is
normal crossings by Kollár-Xu.''

``6.2.6. To conclude, mirror symmetry (via
SYZ conjecture and Kontsevich-Soibelman)
predicts that the essential skeleton
$\text{Sk}(\mathcal{X})$ for a MUM
degeneration is $S^n$ and respectively
$\mathbb{C}P^n$ for Calabi-Yau's and
respectively hyper-Kähler's. The following
result is a weaker version of this statement,
saying that it holds in a cohomological
sense.''

\begin{theorem}
\label{stanley-reisner-theorem-skeleton-cohomology}
\begin{reference}
\cite[Theorem 6.14]{KLSV}
\end{reference}
Let $\mathcal{X}/\Delta$ be a minimal dlt
degeneration of $K$-trivial varieties.
Assume that the general fiber $X_t$ 
is a simply connected $K$-trivial variety.

Assume that the degeneration is maximal unipotent.
Then
\begin{enumerate}
\item $H^* (\text{Sk}(\mathcal{X}),\mathbb{Q})
\cong \prod_iH^* (S^{k_i},\mathbb{Q})
\times_jH^* (\mathbb{C}P^{l_j},\mathbb{Q}))$,
where $k_i$ represent the dimensions
of the Calabi-Yau factos and $2l_j$ 
the dimensions of the hyper-Kähler factos
in the Beauville-Bogomolov decomposition of the general fiber $X_t$.

\item Conversely, the cohomology algebra
of the skeleton $\text{Sk}(\mathcal{X})$ 
determines the type of the Beauville-Bogomolov
decomposition of $X_t$.
\end{enumerate}
\end{theorem}



\section{Degenerations of Hilbert schemes (notes on Calla's thesis)}
\label{stanley-reisner-section-degenerations-hilbert}

\noindent
Notes on Calla's thesis on
degenerations of Hilbert schemes after
\href{https://www.youtube.com/watch?v=Lf4NYTvnO9I&t=969s}
{this video}.

The aim is to construct ``good'' degenerations
of Hilbert schemes of points over semistable
degenerations of surfaces. This was motivated by
the study of the Hilbert schemes of points
on ``maximally degenerate K3 surfaces'',
i.e. K3 surfaces that are on the boundary of
the compactified moduli space of K3s.

Basic setup: a flat projective family
of surfaces $X \to C \cong \mathbb{A}^1$ over
an algebraically closed field of characteristic 0.
Consider Étale local model given by
$\Spec k[x,y,z,t]/(xyz-t)$.
The general fibers are smooth, and
the special fibre $X_0$ over $0 \in C$ 
has simple normal crossing singularity.
We picture this as the three axes $x,y,z$ in $X_0$ 
intersecting at the origin.

Let $X^0= X\setminus X_0$ and $C^0=C\setminus \{0\}$.
Consider the Hilbert scheme of each fiber.
We have removed the limit over 0.
The question is: how can we compactify the Hilbert scheme
$\text{Hilb}^n(X^0/C^0)$ after we removed that fiber?
How to add a ``good limit''?

Starting from a ``big'' deformation space,
namely a fiber product $X \times_{\mathbb{A}^1}\mathbb{A}^n$,
Calla produces $X[n]$ after a sequence of blow-ups,
and later takes a quotient using a group action.
Thus produces a stack $\mathfrak{X}$
that contains the modifications
on the singular fiber
as well as the original fibers.

\begin{theorem}[Calla]
\label{stanley-reisner-theorem-calla}
The stack of Li-Wu stable length in 0-dimensional
subschemes in $\mathfrak{X}$ is Deligne-Mumford
and proper over $C$.
\end{theorem}

\noindent
It turns out this is a specific example
where things work out better than expected,
essentially by the way $X[n]$ was constructed.
The general case is not so ``nice''
(see [Maulik-Ranganathan]).

\medskip\noindent
{\bf Note.} At some point the main focus
of Calla's work turned to the Hilbert schemes.
I wonder if the generalised Kummer fourfolds studied in the
$\mathbb{C}P^2_9$ topic are actually examples
of these Hilbert schemes, or if it is a genuine
change of subject.


% --- End stanley-reisner.tex ---

% --- Begin surfaces.tex ---
\chapter{A tour through algebraic surfaces}
\label{chapter-surfaces}



\noindent
Minicourse by
Aline Zanardini
(EPFL, Switzerland),
IMPA Summer School, 2026.

\medskip\noindent

\bigskip\noindent


\section{Plan}
\label{surfaces-section-plan}

\begin{enumerate}
\item Basics.
\item Birational maps and classification.
\item Elliptic surfaces.
\item Halphen surfaces.
\item K3 surfaces and friends.
\item Research questions.
\end{enumerate}

\section{References}
\label{surfaces-section-references}

\begin{enumerate}
\item A. Beauville. Complex algebraic
surfaces.
\item R. Miranda. Overview of algebraic
surfaces.
\item M. Reid. Chapters on algebraic
surfaces.
\end{enumerate}

\section{Introduction}
\label{surfaces-section-intro}

\noindent
Everything will be over the complex numbers.
By ``surface'' we mean ``algebraic surface'',
but we drop the term ``algebraic''.
Informally, a surface is a smooth projective
algebraic variety of dimension 2.
We may think of surfaces also as
compact connect complex manifolds of
dimension 2.

$\mathcal{M}(S)$, meromorphic functions on
$S$, has transcendence degree
2 over $\mathbb{C}$.
That is, for $p,q \in S$
there exists $f \in \mathcal{M}(S)$
such that $f(p)\neq f(q)$
and for all $p \in S$ there exists $f,g \in
\mathcal{M}(S)$
such that $(f,g)$ give local coordinates at
$p$.

\begin{remark}
\label{surfaces-remark-trans-impli-surfa-def}
The remark on transcendence just made implies
that $S$
is smooth algebraic of dimension 2.
Indeed, compactness implies that there exists
meromorphic functions  $\varphi_1,…,\varphi_n
\in \mathcal{M}(S)$
such that $S
\xymatrix{\ar@{^{(}->}[r]^{alg}&}\mathbb{P}
^n$,
$p \mapsto (1:\varphi_1(p):…:\varphi_n(p))$
and Chow theorem
imply that $S = V(f_1,…,f_k)$ with
$f_i \in \mathbb{C}[x_0,x_1,…,x_n]$.
\end{remark}

\begin{example}
\label{surfaces-example-surfaces}
\begin{enumerate}
\item $S = V(f) \subset \mathbb{P}^3$
with $f \in \mathbb{C}[x_0,x_1,x_2,x_3]$
homogeneous of degree $d$.
For example, $f=x_0^d+x_1^d+x_2^d+x_3^d$.
Take $d=1$, then we can assume that $f=x_i$,
and then $S\simeq\mathbb{P}^2$.
If $d=2$ we can assume that
$f=x_0x_1-x_2x_3$ and
one can verify that $S \simeq \mathbb{P}^1
\times \mathbb{P}^1$.
If $d=3$, $S \simeq Bl_6\mathbb{P}^2$.
These are all rational surfaces.

If $d=4$, $S$ is a K3 surface.

\item (Complete intersections of type
$(d_1,…,d_{n-1})$ in $\mathbb{P}^n$.)
Let $f_1,…,f_{n-2} \in
\mathbb{C}[x_0,x_1,…,x_n]$ homogeneous
of degrees $d_1,…,d_{n-2}$.
Then $S=V(f_1,…,f_{n-2})$.

E.g., intersections of type $(2,3)$ in
$\mathbb{P}^4$
or of type $(2,2,2)$ in $\mathbb{P}^5$, which
are K3 surfaces.

\item (Ruled surfaces.)
Take any curve $C$ and consider  $C \times
\mathbb{P}^1$
(or anything birational to it).
[This is interesting because we can consider
the projection to $C$,
and we get a fibration.]
\end{enumerate}
\end{example}

\section{Curves on surfaces (divisors)}
\label{surfaces-section-curves-on-surfaces}

\noindent
We will use: $\text{Div}(S)/\sim \simeq
\text{Pic}(X)=H^1(S,\mathcal{O}_S^*)$
where $\sim$ is linear equivalence.

$D \subset S$ is a divisor if and only if $D$
is one of the following:
\begin{enumerate}
\item (Cartier divisor.)
[I can think of $D$ given locally as the zero
locus of some function.]
$D:f=0$ for  $0 \neq  f \in \mathcal{M}(S)$ +
gluing condition.
This is equivalent to saying that
$D|_f=\text{div}(f)$.

\item (Weil divisors.)
$D=\sum n_iC_i$ for $n_i \in \mathbb{Z}$,
$C_i \subset S$ irreducible curves,
where there are only finitely many nonzero
coefficients $n_i$.

\item (Line bundle.)
$D$ effective (i.e. $n_i\geq 0$), we
associate
$\mathcal{O}_S(D)$, a holomorphic line bundle
on $S$ along
with a nonzero section [whose zero locus
defines $D$.
Recall that $U \mapsto  \mathcal{O}_S(D)(U)$
is ``given''
by $g(z,w)$ such that  $fg$ is holomorphic
where $D|_U=\text{div}(f)$.
And we extend by linearity.
\end{enumerate}

\noindent
[And we say that two divisors are linearly
equivalent if their difference is $div(f)$
for some $f$.]

\begin{example}
\label{surfaces-example-divisors}
\begin{enumerate}
\item $\text{Pic}(\mathbb{P}^2)=\mathbb{Z}$
generated by $\mathcal{O}_{\mathbb{P}^2}(1)$,
the dual of the tautological bundle, or,
equivalently,
the class of a hyperplane $H$,
whose sections are polynomials of degree 1
(in the appropriate number of variables).

\item $\text{Pic}(\mathbb{P}^n)=\mathbb{Z}$
is generated by
$\mathcal{O}_{\mathbb{P}^n}(1)$.
We have
$\mathcal{O}_{\mathbb{P}^n}(d):=\mathcal{O}
_{\mathbb{P}^n}^{\otimes
d}$.
For $n=2$, the sections of $\mathcal{O}(d)$
are
plane curves of degree 2.
\end{enumerate}
\end{example}

\noindent
Suppose that $\pi:S
\xymatrix{\ar@{^{(}->}[r]&} \mathbb{P}^n$.
The {\it hyperplane class} is  $H|_S:=
\pi^*\mathcal{O}_{\mathbb{P}^n}(1)$.
We denote
$\mathcal{O}_S(1)=\mathcal{O}_S(H|_S)$.
This is always a non-trivial class.

We also have the {\it canonical class},
which is the class of any divisor $K$
such $\mathcal{O}_S(K)=\Omega^2_S$,
the sheaf of holomorphic 2-forms,
informally its elements look like
$f(z,w)dx\wedge dw$.
We denote this as $K$,  $K_S$ or $\omega_S$.

\begin{example}[Canonical class of projective
plane]
\label{surfaces-example-canonical-class}
Take coordinates $Z=X_1/X_0$ and $W=X_2/X_0$
on $U_0: X_0 \neq 0$.
Likewise put $U=X_0/X_1$, $V=X_2/X_1$ at
$U_1$.
Then $dZ\wedge dW=-U^{-3}dU\wedge dV$.
Thus, the canonical class of $\mathbb{P}^2$
is linearly equivalent to $-3H$.
In general, $K_{\mathbb{P}^n}\sim -(n+1)H$.

In practice, we can use this to compute $K_S$
for any  $S \xymatrix{\ar@{^{(}->}[r]&}
\mathbb{P}^n$
via $\omega_S=\omega_{\mathbb{P}^n}\otimes
\Lambda^{n-2}N_S/\mathbb{P}^n$.
\end{example}


\section{Numerical invariants}
\label{surfaces-section-numerical-invariants}

\begin{itemize}
\item (Hodge decomposition.)
$$
H^k(S,\mathbb{C})=\bigoplus_{p+q=k}
\underbrace{H^q(S,\Omega^p_S)}_{
H^{p,q}(S)}.
$$
Setting $h^{p,q}=\dim H^{p,q}$ + symmetries,
we have the Hodge diamond
$$
\begin{aligned}\xymatrix@C=.5em@R=.5em{
&  &   h^{0,0}\\
&  h^{0,1}&  &  h^{1,0}\\
h^{0,2}&&  h^{1,1}&  &  h^{2,0}\\
&  h^{1,2}&  &  h^{2,1}\\
&  &  h^{2,2}
}\end{aligned}
\qquad =\qquad
\begin{aligned}\xymatrix@C=.5em@R=.5em{&  &
1\\
&  q&  &  q\\
p_g&&  h^{1,1}&  &  p_g\\
&  q&  &  q\\
&  &  1
}\end{aligned}
$$
where we call $q$ the {\it irregularity}
$p_g$ the {\it geometric genus}
and the {\it Euler characteristic}
is the alternating sum of Betti numbers.
Notice that $h^{1,1}\geq 1$.
\end{itemize}

\section{The Néron-Severi group and Lefschetz
(1,1) theorem}
\label{surfaces-section-neron-severi}

\noindent
There is always a map
$$
\begin{aligned}
c_1: \text{Pic}(S) &\longrightarrow
H^2(S,\mathbb{Z})
\overset{\substack{\text{Poincaré} \\
\text{dual}}}{\cong}
H_2(S,\mathbb{Z}) \\
\mathcal{L} &\longmapsto c_1(\mathcal{L})
\end{aligned}
$$

\noindent
We define
$\text{NS}(X):=\text{Im}c_1/\text{torsion}$.
It is a finitely generated abelian group
isomorphic to $\mathbb{Z}^{\rho}$
where $\rho$ is the Picard rank.

% \begin{theorem}[Lefschetz (1,1)-classes]
% \label{surfaces-theorem-lefschetz-11}

\begin{theorem}
\label{surfaces-theorem-form}
$$
\text{NS}(S)=H^2(S,\mathbb{Z})\cap H^{1,1}(S)
$$
$\text{NS}(S)$, $\text{Div}(S)/\sim$ and
$\text{Pic}(S)$
have more structure!
There exists a unique bilinear form
$$
\begin{aligned}
\cdot: \text{Div}(S)\times \text{Div}(S)
&\longrightarrow \mathbb{Z} \\
(C,D) &\longmapsto C\cdot D
\end{aligned}
$$

\noindent
such that
\begin{enumerate}
\item If $C $ and $D$ are smooth curves which
intersect
transversally, then $C.D=\# C \cap D$.
\item If $C \sim C'$, then $C\cdot D = C'
\sim D$.
\end{enumerate}
\end{theorem}

\begin{definition}
\label{surfaces-definition-form}
$C.D=\mathcal{O}_S(C).\mathcal{O}_S(D)
=\left\langle{}c_1(\mathcal{O}_S(D) \cup
c_1(\mathcal{O}_S(D)),
[S]_{\text{fund}}\right\rangle{}$,
and we say that $C.C=\text{deg}
\mathcal{N}_{C/S}$.
\end{definition}

\noindent
[An important invariant of a surface
is the self-intersection number of the
canonical class.]
Particular case: $K^2=c_1^2$
where $c_1^2$ is the first Chern number.

\begin{example}
\label{surfaces-example-intersection}
\begin{enumerate}
\item $K_{\mathbb{P}^2}^2=(-3 H).(-3 H)=9$.
\item (Degree of $S$.)
Let $S \xymatrix{\ar@{^{(}->}[r]&}
\mathbb{P}^n$.
Recall our notation above for
$\mathcal{O}_S(1).\mathcal{O}_S(1)$.
If $n=3$,
the degree $d$ of $S$ is the degree of
$\mathcal{N}_{S/\mathbb{P}^3}$.
Thus the adjunction becomes
$\omega_S=\mathcal{O}_S(d-4$

\item $C$ smooth curve, then
$f:S \to\!\!\!\!\!\to C$,
$F$ fiber, then  $F^2=0$.

\item $g=\tilde{S}\overset{d}{\to}D$,
then $D_1,D_2 \in \text{Div}(S)\implies
g^*(D_1)\cdot g^*(D_2)=d(D_1.D_2)$.
\end{enumerate}
\end{example}

\begin{definition}
\label{surfaces-definition-nef-and-ample}
A divisor $D \subset S$ is {\it nef} if and
only if
for each $C \subset S$ irreducible,
$D\cdot C \geq 0$.
We say that $D$ is ample if and only if
$D^2>0$ and for each irreducible $C \subset
S$
we have $D.C \geq 0$.
\end{definition}

\section{Big theorems}
\label{surfaces-section-big-theorems}

\begin{theorem}
\label{surfaces-theorem}
Let $S$ be a surface.
\begin{enumerate}
\item (Noether's formula.)
$$
12\chi = K^2+e=c_1^2+c_2
$$
where $e$ is the Euler number.

\item (Riemann-Roch.) Given $D \in
\text{Div}(S)$,
$$
\chi(\mathcal{O}_S(D))=\frac{D.(D-K)}{2}+
\chi(\mathcal{O}_S).
$$

\item (Genus formula.) Given $C \subset S$
smooth,
$$
2g-2=C.(C+K).
$$
\end{enumerate}
\end{theorem}

\begin{example}
\label{surfaces-example-big-theorems}
\begin{enumerate}
\item For $\mathbb{P}^2$, $K^2=9$ and $e=3$.
Thus $12\chi = 12 \implies \chi =1$.
The irregularity is
$q=H^0(\mathbb{P}^2,\Omega^1)=0$.
The arithmetic genus is $p_g=\dim
H^0(\mathbb{P}^2,\Omega^2)=0$.
Thus, the Hodge diamond is
$$
\xymatrix@C=.5em@R=.5em{
& & 1\\
& 0 && 0\\
0 && 1 && 0\\
& 0 && 0\\
& & 1
}
$$
\item Any rational surface has a Hodge
diamond
$$
\xymatrix@C=.5em@R=.5em{
& & 1\\
& 0 && 0\\
0 && \rho && 0\\
& 0 && 0\\
& & 1
}
$$
The Noether formula gives
$12=K_S^2+\rho+2 \implies K_S^2=10-\rho$.

\item $C \subset \mathbb{P}^2$ smooth of
degree $d$,
then $2g-2=d^2-3d$ and
$g=\frac{d^2-3d+2}{2}=\frac{(d-1)(d-2)}{2}
=\binom{d-1}{2}$.
In particular,
if $d=1$ or  $2$,
then  $g=g(C)=0$
if $d=3$ then  $g(C)=1$.
\end{enumerate}
\end{example}

\section{Birational maps and classification
theorem}
\label{surfaces-section-birat-maps-class-theor}

\noindent
Question: given two surfaces $S$ and $S'$,
when are $S$ and $S'$ birational/biregular or
isomorphic?

\begin{itemize}
\item Rules surfaces
(all birational to $C\times \mathbb{P}^1$ for
some $C$.
\item Classify minimal (models of) surfaces
by looking at the positive of $K$
(and related invariants
$q$, $p_g$ and $\kappa$).
\end{itemize}

\begin{center}\begin{tabular}{c|c|c|c|c}
$K$ & $\kappa$ & $p_g$ &  $q$  & Structure\\
\hline\hline
$K^2>0$&2&&&General type\\ \hline
$K^2=0$&1&&&Elliptic*\\\hline
$K=0$&0&
$\substack{1\\0\\1\\0 \\ }$&
$\substack{2 \\1\\0\\0 }$
&$\substack{
\text{abelian} \\
\text{hyperelliptic}\\
\text{K3}\\
\text{Enriques}}$\\\hline
&$-\infty$&$\substack{0 \\ 0}$&
$\substack{\geq 1 \\ 0}$&
$\substack{\text{ruled surfaces} \\
\text{rational}}$
\end{tabular}\end{center}

\medskip\noindent
The following proposition is also a
definition:

\begin{proposition}
\label{surfaces-proposition-birational}
Let $S$ and $S'$ be two surfaces, then the
following
statements about $S$ and $S'$ are equivalent:
\begin{enumerate}
\item $S$ and $S'$ are birational
\item $\mathcal{M}(S)\simeq \mathcal{M}(S')$
as algebras (over $\mathbb{C}$).
\item There exists $U \subset S$, $V \subset
S'$
opens such that $U \simeq V$.
\end{enumerate}
\end{proposition}

\section{The blow up of the projective plane
at a point}
\label{surfaces-section-blow-up}

\noindent
Consider $p=(0:0:1)$ and
$$
\begin{aligned}
\varphi_p: \mathbb{P}^2\setminus \{p\}
&\longrightarrow \mathbb{P}^1 \\
(x_0:x_1:x_2) &\longmapsto (x_0:x_1)\\
q &  \longmapsto L_{pq}
\end{aligned}
$$

\noindent
Consider also
$$
\Gamma:=\left\{((x_0:x_1:x_2),(y_0:y_1)) \in
\mathbb{P}^2\times \mathbb{P}^1
: \det\begin{pmatrix}x_0&x_1\\
y_0&y_1\end{pmatrix}=0\right\}.
$$

By construction,
$$
\pi:\Gamma \to \mathbb{P}^2
$$
is an isomorphism over $U=\mathbb{P}^2
\setminus \{p\}$
and $\pi^{-1}(p)=\{p\}\times \mathbb{P}^1
\simeq\mathbb{P}^1:=E$.

\begin{definition}
\label{surfaces-definition-blow-up}
$\pi:\Gamma \to \mathbb{P}^2$ as above is
called
the {\it blowup} of $\mathbb{P}^2$ at $p$
(also $\Gamma$). The curve $E$ is called
the {\it exceptional divisor}.
\end{definition}

\noindent
Note that we have the following local
description
(chart $x_2=1$):
$\{(x,\ell) \in \mathbb{C}^2\times
\mathbb{P}^1: x \in \ell\}
\to \mathbb{C}^2$.
That is, the exceptional divisor may be
identified with
the zero section of the tautological bundle
$\mathcal{O}_{\mathbb{P}^1}(-1)$,
which means that the self intersection of $E$
is 1,
i.e. $E^2=1$.

\begin{remark}
\label{surfaces-remark-blowup}
\begin{enumerate}
\item The surface $Bl_p\mathbb{P}^2:=\Gamma$
is the Hirzebruch surface $\mathbb{F}_1$ (or
$\Sigma_1$).

\item Choosing local coordinates, we can talk
about
$Bl_p S$ for $p \in S$.

\item In fact, $\pi:Bl_pS \to S$
is characterized by the following universal
property:
if $f:\tilde{S}\to S$ is birational such that
$f^{-1}$
is not defined at $p \in S$, then $f$
factorizes
as $f:\tilde{S}
\underbrace{\xrightarrow{g}}_{\text{bir}}
\hat{S}:=Bl_pS \xrightarrow{\varepsilon} S$
where $\varepsilon^{-1}(S\setminus \{p\})
\simeq S \setminus \{p\}$
and $\varepsilon^{-1}(p) \simeq
\mathbb{P}^1$.
\end{enumerate}
\end{remark}

\begin{proposition}
\label{surfaces-proposition-blowup}
Consider $\pi:\tilde{S}=Bl_pS \to S$
as above, with exceptional divisor $E$,
\begin{enumerate}
\item $E^2=-1$,

\item $E\cdot \pi^*D=0$ for all $D \in
\text{Div}(S)$,

\item If $C \subset S$ has mutliplicity $m$
at $p$, then
$\tilde{C}:= \overline{\pi^{-1}(C) \setminus
E}\subset\tilde{S}$
satisfies
$\tilde{C}\cdot E=m$ and
$\tilde{C}^2=C^2-m^2$

\item $K_{\tilde{S}}=\pi^* K_S+E\implies
K_{\tilde{S}}^2=K_S^2-1$

\item $q,p_g$ and  $\chi$ are the same for
$\tilde{S}$ as for $S$

\item $e$,  $h^{1,1}$ and $b_2$ all go up by
1

\item
$\text{Pic}(\tilde{S})=\pi^*(\text{Pic}(S))
\oplus
\mathbb{Z}E$.
\end{enumerate}
\end{proposition}

\section{Castelnuovo's criterion}
\label{surfaces-section-castelnuovo-criterion}

\noindent
[It turns out that if you find a surface with
a $-1$ curve,
then that surface is the blow up of another!]

\begin{theorem}
\label{surfaces-theorem-castelnuovo}
If on a surface $\tilde{S}$ one finds a curve
$E \simeq \mathbb{P}^1$
with $E^2=-1$, then $\tilde{S}=Bl_pS$
for some $p \in S$.
\end{theorem}

\begin{definition}
\label{surfaces-definition-minimal-surface}
A surface $S$ is {\it minimal}
if it has no $(-1)$-curves
(i.e. a curve $\simeq \mathbb{P}^1$ with
self-intersection -1).

(If and only if
$f:S \xymatrix{\ar@{-->}[r]&}\tilde{S}$ and
$S$ minimal,
then $f$ is an isomorphism.)
\end{definition}


\section{Mori's point of view}
\label{surfaces-section-moris-point-of-view}

\noindent
If $S$ is a surface such that $K_S$ is not
nef
[so, by definition of nef there must exist
an irreducible curve that intersects
the canonical divisor negatively, and this
curve
in fact is a $(-1)$ curve:]
then there exists a $(-1)$-curve on $S$.




\section{Kodaira dimension}
\label{surfaces-section-kodaira-dimension}

\noindent
Let $S$ be a surface.

\begin{definition}
\label{surfaces-definition-kodaira-dimension}
For each $n \geq 0$,
we define the {\it pluri-genera}
$$
P_n:= \dim H^0(S,\mathcal{O}_S(n K))
$$
and the {\it Kodaira dimension}
$$
\kappa:=
\begin{cases}
-\infty&\qquad \text{if }P_n=0\forall n \\
\text{smallest $k$ such that }
P_n=O(n^k)&\qquad \text{otherwise}
\end{cases}
$$
\end{definition}

\begin{remark} \label{surfaces-remark-kodaira} (By
Daniel.) In one of Misha's course we
have that the function $\dim H^0(n K)$ is
known to be polynomial for all
projective varieties (Birkar, Cascini, Hacon,
McKernan). Conjecturally it is
always polynomial. Thus, the definition of
Kodaira dimension is nothing
more than the degree of this polynomial.

Further, in \cite[Definition 2.2.26]{huc} we
see
that this should be equivalent to defining it
as the
transcendence degree over $\mathbb{C}$
of the fraction field of the ring
$\bigoplus_{m \geq 0}H^0(X,K_X^{\oplus m})$,
which is endowed with a natural product
which, in general for
vector bundles $E$ and $F$, maps
$H^0(X,E) \otimes H^0(X,F) \to H^0(X,E
\otimes F)$.
\end{remark}

\section{Results of Mori's theory}
\label{surfaces-section-results-of-mori}

\begin{theorem}
\label{surfaces-theorem-one}
Given a surface $S$,
there exists a chain
$$
S=S^0 \xrightarrow{\sigma_1} S^1 \to \cdots
\xrightarrow{\sigma_N}S_N=S'
$$
where each $\sigma_i$ is the contraction of a
$(-1)$ curve $E_i$
and $S'$ satisfies
\begin{enumerate}
\item $K_{S'}$ is nef or
\item $S'$ is a $\mathbb{P}^1$-bundle
over a curve or
\item $S' \simeq \mathbb{P}^2$.
\end{enumerate}
\end{theorem}

\begin{theorem}
\label{surfaces-theorem-two}
If a surface $S$ is such that
$K_S$ is not nef then there exists $\varphi:S
\to X$,
with $\dim X=0,1$ or $2$, contracting at
least
one curve to a point, such that
$-K_S\cdot C>0$ for every curve $C$ contained
in a fiber of $\varphi$.
\end{theorem}

Note:
$$
\dim X=
\begin{cases}
0\xymatrix{\ar@{~>}[r]&}&-K \text{ is
ample}\\
1\xymatrix{\ar@{~>}[r]&}&S \text{ is a conic
bundle}\\
2\xymatrix{\ar@{~>}[r]&}& S=Bl_{p_1,…,p_n}.
\end{cases}
$$


\section{Examples and the cubic surface
revisited}
\label{surfaces-section-examples-and-cubic}

\begin{enumerate}
\item Slogan: ``blowups resolve
singularities''.
Consider the plane cubic
$$
C=\{(x_0:x_1:x_2) \in \mathbb{P}^2:
\underbrace{x_1^2x_2-x_0^2(x_0+x_2)=}_{:=f_C}
0\},
$$
the point $p=(0:0:1)$
and $\pi:Bl_p\mathbb{P}^2 \to \mathbb{P}^2$,
then
$$
\pi^{-1}(C)=
\begin{cases}
f_C=0\qquad & \\
x_0y_0=y_0x_1\qquad &
\end{cases}
$$
[This curve has a simple node.
In the blowup, the curve will intersect
the exceptional divisor in two points ---
we have ``separated'' the singularity into
two points]

\item $Bl_1\mathbb{P}^1 \times \mathbb{P}^1
\simeq Bl_1 \mathbb{F}_1
\simeq Bl_2\mathbb{P}^2$.
[Consider two curves $L_1$ and $L_2$ on
$\mathbb{P}^1 \times \mathbb{P}^1$.
Their strict transforms along with the
exceptional divisor
become $(-1)$ curves.
Then contract $\tilde{L}_1$, this gives
$\mathbb{F}_1$,
the Hirzebruch surface, which is just the
blowup
at a point of $\mathbb{P}^2$.
$\mu(\tilde{L}_2)=-1$, $\mu(E\setminus
\tilde{L}_1)=0$.]

\item Consider two cubics $C_1$ and $C_2$.
They intersect in 9 points.
Do
$$
\begin{aligned}
\mathbb{P}^2 &\xymatrix{\ar@{-->}[r]&}
\mathbb{P}^1 \\
p &\longmapsto (C_1(p):C_2(p))
\end{aligned}
$$

\noindent
Let $S=Bl_{p_1,…,p_9}\mathbb{P}^2$.
Then
\begin{enumerate}
\item $K_S^2=0$ [Because $K^2$ of
$\mathbb{P}^2$ is 9,
and for every blowup we substract 1!]
\item $\text{Pic}(S)$  ``$=$''
$\pi^*\text{Pic}(\mathbb{P}^2)
\oplus\mathbb{Z}E_1 \oplus … \oplus
\mathbb{Z}E_q$.
\item $-K_S=-\pi^* (K_{\mathbb{P}^2})
-E_1-…-E_q$ is nef.
(In fact, this is the class of a fiber of $S
\to \mathbb{P}^1$.)
\end{enumerate}

\item (Continuation of previous item.)
$-K_S$ nef + genus formula  $\implies $
$C^2 \geq -2$ for every $C\subset S$
irreducible and rational.
This is an elliptic surface!

[We have taken a 1-dimensional linear system.
But why not take a higher-dimensional linear
system?]

\item Fix $p_1,…,p_6$ in  $\mathbb{P}^2$ in
general position
(not three in a line, not six in a conic).
Consider $f_1,f_2,f_3,f_4 \in
\mathbb{C}[x_0,x_1,x_2]$
homogeneous of degree $3$ vanishing at $p_i$
and consider
$$
\begin{aligned}
\mathbb{P}^2
&\xymatrix{\ar@{-->}[r]&}\mathbb{P}^3 \\
p &\longmapsto (f_1(p):f_2(p):f_3(p):f_4(p)).
\end{aligned}
$$

\noindent
The resolved morphism $Bl_6\mathbb{P}^2 \to
\mathbb{P}^3$
is an embedding with image of degree $3$
[a cubic surface on $\mathbb{P}^3$].
\end{enumerate}



\section{Elliptic surfaces}
\label{surfaces-section-elliptic-surfaces}

\begin{definition}
\label{surfaces-definition-elliptic-surface}
An {\it elliptic surface}
is a surface $S$ together with $f:S\to C$
holomorphic
(with connected fibers)
from $S$ to a smooth curve $C$
such that the generic fiber
is a smooth curve of genus 1.

Moreover, we assume that there
do not exist $(-1)$-curves on fibers of $f$.
\end{definition}

Warning: I don't assume the existence
of a section
[a distinguished point on the fibers],
and I'm not saying that
all fibers are smooth.

\medskip\noindent
Singular fibers have been classified by
Kodaira:
$I_0,I_1,(m)I_n,II, IV,IV^*=\tilde{E}_6,
III^*=\tilde{E}_7, II^*=\tilde{E}_8$ and
$I_n$.

\begin{example}
\label{surfaces-example-pencils-of-cubics}
$$
\left\{\substack{\text{pencils} \\
\text{of cubics}\\
\text{w/ smooth}\\
\text{member}}\right\}
\leftrightarrow
\left\{\substack{\text{RES} \\ \text{w/
section}}\right\}.
$$
$$
\mathcal{P} \mapsto S_{\mathcal{P}}.
$$
 Question: given $\mathcal{P}$ and
 $S_{\mathcal{P}}$
how many singular fibers does
$S_{\mathcal{P}}$ have?
\end{example}

\begin{lemma}
\label{surfaces-lemma-discriminant-locus}
The discriminant locus in $\mathbb{P}^9$
[this $\mathbb{P}^9$ parametrizes cubics in
$\mathbb{P}^2$ ]
has degree 12.
\end{lemma}

\noindent
Note that the Hodge diamond $S_{\mathcal{P}}$
is
$$
\xymatrix@C=.5em@R=.5em{
& & 1\\
& 0 && 0\\
0 && 10 && 0\\
& 0 &&0\\
& & 1
}
$$
(Notice that adding up the center column
gives 12.)

\begin{example}
\label{surfaces-example-pencil-of-conics}
[We take a conic $C$
and a triple line $3L$
(it has to be triple so that it is a pencil
of conics)
which is tangent to $C$ at an inflection
point.
Then we blow up the point of intersection 9
times,
each time picking a point in the intersection
of the strict transform of the original cubic
and the triple line.]

Blowing up 9 times give
a surface $S_{\mathcal{P}}$ with
a $II^*$ fiber.
\end{example}

\medskip\noindent
Question: How is Kodaira's classification
obtained?

\begin{enumerate}
\item If $F=\sum n_iC_i$ is singular fiber,
then the intersection form restricts
to
$\text{span}(\{C_i\})$.
This restriction is negative semi-definite
with a one-dimensional kernel spanned by $F$.

\item If $F$ is irreducible and $p_a(F)=1$,
then $F=$smooth,nodal or cuspidal.
 If $F$ is reducible, then each $C_i$ is
a $(-2)$-curve.
\end{enumerate}

\noindent
This two items combined imply that the
dual graph [vertices are curves,

two vertices joined if they intersect]
is a Dynkin diagram (ADE type)
[by a result of graph theory].




\section{The Weistrass model}
\label{surfaces-section-weierstrass-model}

\noindent
Let $S \to C$ be an elliptic surface
with a section.
A choice of a section defines
a divisor of degree 1 on the generic fiber
$S_\eta$,
i.e. a point $p$.

Looking at $H^0(S_\eta,np)$
for every $n$, we get:
$$
\begin{aligned}
H^0(p)&=\left\langle{}1\right\rangle{},\qquad
H^0(2p)=\left\langle{}1,x\right\rangle{},
\qquad H^0(3p)=\left\langle{}1,x,y\right\rangle{},\\
…,&\qquad
H^0(6p)=\left\langle{}1,x,y,x^2,xy,x^3,y^2\right\rangle{},
\end{aligned}
$$

\noindent
so there must exist a relation.
Up to further tricks,
we get to $y^2=x^3+ax+b$.

Making the construction global
gives the Weierstrass model.

\begin{itemize}
\item Any $S \to C$ is birational to a
Weierstrass fibration $W \to C$
(where $W$ can be singular, while $C$ is
smooth)
where all fibers are irreducible and have
$p_a=1$.
\end{itemize}

\noindent
Geometrically, $W$ is obtained from $S$
by contracting all fiber components that
do not meet the chosen section

[Once we show that we can do this,
it will suffice to study Weierstrass
fibrations.]

[From now on $W$ is a Weierstrass fibration.]
Given $\pi_i:W \to C$ (with a section
$\sigma$)
consider the fundamental line bundle
$\mathcal{L}=(\pi_*N_{\sigma/K})^\vee=
(R^1\pi_*\mathcal{O}_W)^\vee$.
Then for all $n \geq 2$
we have a splitting
$\pi_*\mathcal{O}_{W}(n\sigma)\simeq
\mathcal{O}_C \oplus \mathcal{L}^{-2}
\oplus … \oplus \mathcal{L}^{-n}$.
In particular, we have a section
$$
\pi^*\pi_*\mathcal{O}_W(3\sigma)
\to\!\!\!\!\!\to\mathcal{O}
_W(3\sigma)
$$
gives a map $f:W \to
\mathbb{P}(\pi_*\mathcal{O}_W(3\sigma))$
that exhibits $W$ as a relative cubic in a
$\mathbb{P}^2$-bundle
over $C$.

A global equation is given by
$$
Y^2Z=X^3+AXZ^2+BZ^3
$$
where $(Z,X,Y,A,B)$ are interpreted
as sections of
$(\mathcal{O}_C,\mathcal{L}^2,\mathcal{L}^3,
\mathcal{L}^4,\mathcal{L}^6)$.
$\sigma$ is given by $Z=X=0$.

\begin{remark}
\label{surfaces-remark-singular-fibers}
The singular fibers of a minimal resolution
of $W$
are determined by the order of vanishing of
$A,B$ and
$\Delta=4A^3+27B^2$.
\end{remark}

\medskip\noindent
Consequences:
\begin{enumerate}
\item $\omega_W=\pi^*(\omega_C \otimes
\mathcal{L})$
implies that $K_W^2=0$.
$e=R\cdot \text{deg} \mathcal{L}=$
number of singular fibers.

\item $\kappa(W) \leq  1$.

\item $\text{deg} \mathcal{L}=\chi(W)$,
$e=R\cdot \text{deg} \mathcal{L}=$
number of singular fibers.

\item If $C=\mathbb{P}^1$,
then
$\mathcal{L}=\mathcal{O}_{\mathbb{P}^1}(d)$
for some $d$.

$$
\begin{cases}
d=1\leftrightarrow&\text{(general)
RES w/section} \\
d=2\leftrightarrow&\text{ K3s,
}K_W=0\\
3\geq 3\leftrightarrow&
\kappa=1\text{ and }K_W=\alpha F
\text{for some }\alpha>0.
\end{cases}
$$
\end{enumerate}

\medskip\noindent
Invariants. $\pi:W \to C$.
Set $g=g(C)$.

\begin{center}\begin{tabular}{c|c c}
& $q$ & $p_g$\\
\hline
$W$ not a product &  $g$ & $ g+
\text{deg}\mathcal{L}-1$ \\
product& $g+1$ &  $g+\text{deg}\mathcal{L}$.
\end{tabular}\end{center}

\noindent
We can also compute the pluri-genera $P_n$
and deduce:
\begin{itemize}
\item $g=0$: done.

\item $g=1$:
 $$
\begin{cases}
\mathcal{L}=\mathcal{O}_C\qquad &E\times E \\
\mathcal{L}\text{ is torsion of order 2,3,4
or 6 hyp}\\
\text{deg}(\mathcal{L})\geq 1\qquad &\kappa=1
\end{cases}
$$

\item $g\geq 2$: $\kappa=1$.
\end{itemize}

\begin{proposition}
\label{surfaces-proposition-minimal-surface}
Let $S$ be a minimal surface with $\kappa=1$.
Then
\begin{enumerate}
\item $K^2_S=0$ and
\item there exists $f:S
\to\!\!\!\!\!\to C$ elliptic
fibration.
\end{enumerate}
\end{proposition}

\begin{proof}
\begin{itemize}
\item $S$ contains a curve
$E= \sum n_iC_i$
($n_i \in \{ 0,1\}, \text{supp}(E)$ is
connected)
with $E \overset{num}{\not \sim}0$
and $E^2=K_S E=0$.

In fact $K_S \overset{num}{\sim} rE$ for some
$r \in \mathbb{Q}$.

\item Some multiple of $E$ moves in an
elliptic pencil.
Some multiple of $E$ is the movable part of
$\tilde{r}K_S$
for some $\tilde{r} \in \mathbb{Q}$.
\end{itemize}
\end{proof}

\section{Examples}
\label{surfaces-section-examples}

\begin{itemize}
\item (The Fermat quartic.)
Consider $S \subset \mathbb{P}^3$
given by $x_0^4+x_1^4+x_2^4+x_3^4=0$
and for each $(\lambda:\mu) \in
\mathbb{P}^1$,
let
$$
C_{(\lambda:\mu)}
=\begin{cases}
\lambda(x_0^2+\zeta^2x_1^2-\mu(x_2^2-
\zeta^2x_3^2=0\\
\mu(x_0^2-\zeta^2x_1)+\lambda(x_2^2+
\zeta^2x_3^2)=0
\end{cases}
$$
where $\zeta \in \mathcal{M}_8$
(an eight root of unity), primitive.
These curves live in $S$ and
$$
\begin{aligned}
\pi: S &\longrightarrow \mathbb{P}^1 \\
(x_0:x_1:x_2:x_3) &\longmapsto
(x_0^2-\zeta^2x_1^2:-x_2^2-\zeta x_3^2)
\end{aligned}
$$

\noindent
is an elliptic surface.

\item Consider a smooth plane quartic $Q$
and choose a point $p \not\in Q$.
$$
\xymatrix{
S \ar[d]^{Bl_{q_1,q_2}Y \text{ where
$\pi^{-1}(p)=\{q_1,q_2\}$}}\\
Y \simeq Bl_7\mathbb{P}^2\ar[d]_\pi^{2:1}\\
\mathbb{P}^2 \supset Q.
}
$$
where $\pi^{-1}(\ell)$ has $g=1$.

[This produces a rational elliptic surface
with a section]
\end{itemize}

\noindent
Notice that there exist elliptic surfaces
with multiple fibers
(hence without a section).

\begin{example}
\label{surfaces-example-multiple-fibers}
Start with three lines in the plane
intersecting pairwise, and a conic
intersecting tangentially
once every line [Fano plane].
We blow up these three points and obtain
a type $2I_3$ fiber.

In general, we take pencils of sextics
with nodes at 9 points. These are in
correspondence
with Halpen surfaces.
\end{example}





\section{Halphen surfaces}
\label{surfaces-section-halphen-surfaces}

\noindent
[These are a class of rational surfaces.]

\begin{proposition}
\label{surfaces-proposition-res}
Let $\pi:S \to C$ be an elliptic surface with
$S$ rational.
Then
\begin{enumerate}
\item $C=\mathbb{P}^1$, $b_2(S)=10$.

\item $\pi$ has at most one multiple fibre
$mF$
(necessarily type  $m I_n$
for some $n$ and some $m$.

\item either $\pi$ has no multiple fiber
and $-K_S \sim F$ (any fiber),
or $\pi$ has one muliple fiber
of multiplicity $m$ and $-m K_S \sim S_p$

\item if $\pi$ has no multiple fibers, there
exists a section.
\end{enumerate}
\end{proposition}

\begin{proof}[Proof of 2. and 3.]
Canonical bundle formula:
\begin{equation}
\label{surfaces-equation-canonical-bundle-formula}
\omega_S=\pi^*
(\omega_{\mathbb{P}^1}\otimes\mathcal{O}
_{\mathbb{P}^1}(1))
\otimes \mathcal{O}_S\left(
\sum_p(m(p)-1)S_p)\right).
\end{equation}
Then, given any smooth fiber $F$,
for any $n \in \mathbb{N}$,
we have $n K_S \sim n\left(\sum_p\frac{
m(p)-1}{m(p)}-1\right)^\alpha F$.

Note that $\alpha<0$ because no multiple
of $K_S$ can be effective.
Thus, $m(p)=1$ except for at most $1$ at $p$.

Why is
\ref{surfaces-equation-canonical-bundle-formula}
true?
In general, for any relatively minimal
elliptic surface $\pi:S \to C$
with multiple fibers
$m_1F_1,…, m_k F_k$:
$$
\omega_S
=\pi^* (\omega_C \otimes
\underbrace{R^1\pi_*\mathcal{O}_S}
_{\text{deg}\chi})^\vee)
\otimes \mathcal{O}_S\left(\sum
m(p)-1S_p\right).
$$
Why? Because $\pi_*\omega_S$ is locally free,
thus $\omega_S=\pi^*(\omega_C \otimes
\mathcal{L})\otimes\mathcal{O}_S(D)$
where $D$ is the zero divisor
of $\pi^*\pi_*\omega_S\to \omega_S$.
Argue that $D=\sum n_iF_i$
with $n_i< m_i$ using adjunction:
$$
\mathcal{O}_{F_i}\simeq \omega_{F_i}
=\omega_S
\otimes\underbrace{\mathcal{O}_{F_i}(F_i)}
_{\text{torsion}}
=\mathcal{O}_{F_i}(F_i)^{\otimes (n_i+1)}
$$
thus $n_i+1$ has to divide $m_i$.
\end{proof}

\begin{remark}
\label{surfaces-remark-rational-elliptic-surface}
Let $\pi:S \to \mathbb{P}^1$ be a rational
elliptic surface (RES)
and note that $K_S^2=0$.
So, $S$ is not minimal and there exists a
$(-1)$-curve.
By Proposition \ref{surfaces-proposition-res},
there are two possibilities:
\begin{enumerate}
\item There are no multiple fibers,
$M.S_n=1$.
\label{surfaces-item-no-multiple-fibers}

\item There exists a multiple fiber $mF$
\label{surfaces-item-halphen}
and $MS_n=m$,
\end{enumerate}
\end{remark}

\begin{definition}
\label{surfaces-definition-halphen-surface}
A RES as in Item \ref{surfaces-item-halphen}
is called a {\it Halpen surface} of index
$m$.
Alternatively, a rational surface  $S$ is
a {\it Halphen surface} if there exists $m$
such that $|-m K_S|$ is one dimensionsional
and has no fixed component and no base
points.
The smallest such $m$ is the {\it index}.
(We include $m=1 \iff$ there exists a
section.)
\end{definition}

\begin{proposition}
\label{surfaces-proposition-halphen}
If $\pi:S \to \mathbb{P}^1$
is a Halphen surface
of index $m$, there exists $f:S \to
\mathbb{P}^2$
birational such tat $\pi \circ f^{-1}$ is a
{\it Halpen pencil}
of index $m$
(i.e. a pencil of plane curves
with degree $3M$ with 9 base points,
each of multiplicity $m$).
Conversely, the resolution
of any such pencil gives a Halpen surface.
\end{proposition}

\begin{example}
\label{surfaces-example-halphen}
$Q:y^2+xz=0$,
$L:y=0$,
$C:(y^2+xz)(\alpha y+z)+\beta yx^2$,
$\alpha,\beta \neq 0$.

The pencil $\lambda(4Q+L)+\mu(3C)=0$
is a Halphen pencil of index 3.

[The intersection of $C$, $Q$ and $L$
consists of two points,
$p_1=(0:0:1)$ and $p_2=(1:0:0)$.]
Blow up at $p_1$ seven times
and at $p_2$ two times.
[We obtain one the exceptional Dynkin diagram
fibers.]
\end{example}

\begin{example}
\label{surfaces-example-halpen2}
Consider any smooth cubic $C$.
Take $p_1,…,p_9 \in C$ such that
$p_1 \oplus … \oplus p_9$ is $m$-torsion.
Then there exists a curve $D$ of degree $3m$
passing
though $p_1,…,p_9$ with multiplicity $m$,
thus $\lambda(D)+\mu(mC)=0$ is a Halpen
pencil.
\end{example}

\section{Singular fibers}
\label{surfaces-section-singular-fibers}

Recall that
$$
\left\{\substack{\text{Halphen pencils} \\
\text{$\lambda C_{3m}+\mu(mC)=0$}}\right\}
\leftrightarrow
\left\{\substack{
\text{Halphen surfaces} \\
\pi:S \xrightarrow{|-m
K_S|}\mathbb{P}^1}\right\}
$$
$$
C_{3m}\leftrightarrow F\in |-m K_S|.
$$

\begin{proposition}[Z]
\label{surfaces-proposition-singular-fibers}
If $F$ is of type $II^*$, $III^*$ or  $IV^*$,
then $C_{3m}$ is non-reduced.
\end{proposition}

\begin{proof}
$R\subset S$ rational smooth $\implies $ $R^2
\geq  -2$,
$\implies $ exceptional curves in $S$ can
only appear
in disjoint chains.
Now, if $F$ is as in the statement,
then at least one curve in [picture of fiber]
is not exceptional.
Here are the fibers:
\end{proof}

\begin{proposition}[Z]
\label{surfaces-proposition-halpen}
Let $S \to \mathbb{P}^1$ be a Halphen surface
of index 2
and assume that $\pi$ contains a fiber $F$ of
type $II^*$.
Then there exists $\varphi:S \to
\mathbb{P}^2$ such that
the sextic $\varphi(F)$ is one of the
following:
[Picture of conic of multiplicity 3,
line of multiplicity 3 intersecting cubic,
two lines of multiplicity 3,
quadric and line of multiplicity 4,
a line and a line with multiplicity 5].
\end{proposition}

\begin{example}
\label{surfaces-example-halphen3}
$C:y^2z=x(x-2)(x-\alpha z)$, $\alpha \neq
0,1$.
$L:z=0$,
$\tilde{L}:x=0$,
$\mathcal{P}:\lambda(L+5\tilde{L})+\mu(2C)=
0$.

[A point of intersection is $p_1=(0:1:0)$,
where the line is tangent at the inflection
point of $C$.
Another point of intersection is
$p_2=(0:0:1)$.]
Blow up $p_1$ 4 times and $p_2$ 5 times.
\end{example}

\section{Constructions}
\label{surfaces-section-construction}
\begin{enumerate}
\item
$$
\xymatrix{
X\ar[d]^{m:1}\\
S \ar[d]^{|-m K_S|\ni F \text{ smooth}}\\
\mathbb{P}^1
}
$$
For  $m=2$,  $X$ is a K3. For $m\geq  3$,
$\kappa(X)=1$ and contains a $(-m)$ $P^1$
section.


\item Type $II$ degenerations of K3 surfaces
and degenerations of Enriques surfaces
(see Alexeev and Engel).

\item It $\pi:S \to \mathbb{P}^1$ is a
general Halphen surface
of index 2 and we blow up we get a Coble
surface.

\item  The Halpen transform
$$
\xymatrix{
S \ar[d]^{|-m K_S|}\\ \mathbb{P}^1
}
$$
$|F+M|$ curves of $g=1$ through $q$.
Contract $M$ and blow up $q$,
$\xymatrix{\ar@{~>}[r]&}$ we obtain a RES
with section.

\item (Halpne.)
Given $p_1,…,p_8 \in \mathbb{P}^2$,
the locus of the 9th base point of a Halpen
pencil of index $m$
is a curve of degree $3\psi(m)$
having ordinary multiple points at
$p_1,…,p_8$
of multiplicity $\psi(m)$,
where
 $$
\psi(m)=m^2\left(1-\frac{1}{q_2^2}\right)
\left(1-\frac{1}{q_2^2}\right)…
$$
where $m=q_1^{r_1}q_2^{r_2}…$
prime factorization. $m=2 \implies
\psi(m)=3$.

\end{enumerate}


\section{K3 surfaces}
\label{surfaces-section-k3}

\begin{definition}
\label{surfaces-definition-k3}
A surface $S$ is a {\it K3} if
$\Omega^2=\mathcal{O}_S$ and
$h^1(S,\mathcal{O}_S)=0$.
\end{definition}

\noindent
So, $\omega_S=\mathcal{O}_S$.
Also $\Omega^2=\mathcal{O}_S$ means that
there exists
a global nowhere zero holomorphic 2-form
$\sigma$.

\begin{example}
\label{surfaces-example-k3}
\begin{enumerate}
\item $S \subset \mathbb{P}^3$ smooth
quartic.
$$
\xymatrix{
0\ar[r]
&\mathcal{O}_{\mathbb{P}^3}(-4)\ar[r]
&\mathcal{O}_{\mathbb{P}^3}\ar[r]
&\mathcal{O}_S\ar[r]
&0
}
$$
The long exact sequence in cohomology
gives $h^1(S,\mathcal{O}_S)=0$.
Adjunction gives
$\omega_S=\omega_{\mathbb{P}^3}
\otimes\mathcal{O}_{\mathbb{P}^3}(4)|_S
=\mathcal{O}_S$.

\item $\pi:S \xrightarrow{2:1}\mathbb{P}^2
\supset B$
smooth sextic.
$\pi_*\mathcal{O}_S
=\mathcal{O}_{\mathbb{P}^2}\oplus\mathcal{O}
_{\mathbb{P}^2}(-3)$
gives $h^1(S,\mathcal{O}_S)=0$.
(If
$\mathcal{L}=\mathcal{O}_{\mathbb{P}^2}(-3)$,
$B$ is a section of $\mathcal{L}^{-2}$.)
Since $\omega_S=\pi^*
(\omega_{\mathbb{P}^2}\otimes\mathcal{L}^{-1})$
this implies that $\omega_S \simeq
\mathcal{O}_S$.
\begin{remark}
\label{surfaces-remark-ade-singularities-on-k3}
We can allow ADE singularities on $B$ and
take a minimal resolution
$$
\xymatrix{
\tilde{S}\ar[d]_-{2:1}&S\text{ K3
surface}\ar[d]^{2:1}
\ar[l]_{\substack{ \text{min}\\\text{res}
}}\\
\mathbb{P}^2\ar@{-->}[d]_{\substack{B \\
\text{sextic}\\
\text{w/ 9 nodes}}}& \text{ Halphen index
2}\ar[d]^{|-2K_Y|\ni F}\\
\mathbb{P}^1\ar[r]&\mathbb{P}^1
}
$$
\end{remark}

\item $S$ a K3 surface with a fixed point
free involution $i$
then $S/i$ is an {\it Enriques surface} [a
friend of K3],
$p_g=q=0,2K=0$.

\item (Kummer surfaces.) $A$ abelian surface
(e.g. $E \times E^1$)
$a \mapsto  -a$ has 16 fixed points.
$Kum(A):=\widetilde{A/(-1)}$ is a K3 surface.
\end{enumerate}
\end{example}

\noindent
The Hodge diamond of a K3:
$$
\xymatrix@C=.5em@R=.5em{
& & 1\\
& 0 && 0\\
1 && 20 && 1\\
& 0 && 0\\
& & 1
}
$$
[The $H^2$ lattice is]
$$
H^2(S,\mathbb{Z})=E_8(-1)^{\oplus 2} \oplus
U^{\oplus 3},
$$
an even unimodular lattice of signature
$(\underbrace{3}_{>0},\underbrace{19}_{<0})$.

\section{Curves on a K3 surface}
\label{surfaces-section-curves-on-k3}

Let $S$ be a K3 surface and $C \subset S$ a
curve.
By adjunction formula, $C^2=2p_a(C)-2$, so
$C^2 \geq -2$.

In particular, if $C$ is smooth and $C^2=-2$,
$C=\mathbb{P}^1$.

Assume that $C$ is smooth and irreducible,
and set $g=g(C)$,
$\mathcal{L}=\mathcal{O}_S(C)$.

Then $\mathcal{L}^2=2g-2$,
$h^0(S,\mathcal{L})=g+1$.
Moreover, if $g \geq 1$,
then $|C|$ is base point free and the induced
morphism
$S \to \mathbb{P}^9$ restricts to the
canonical map on $C$.

\begin{definition}
\label{surfaces-definition-polarized-k3}
A {\it polarized K3 surface} of degree $2d$
is a pair
$(S,\mathcal{L})$ with $S$ a K3 and
$\mathcal{L}$ an
ample (ind in Pic) line bundle of degree
$\mathcal{L}^2=2d$
\end{definition}

\begin{proposition}
\label{surfaces-proposition-polarized-k3}
For any $g \geq 3$ there exists a polarized
K3
of degree $2g-2$ in $\mathbb{P}^9$.
\end{proposition}

\begin{example}
\label{surfaces-example-polarized-k3}
\begin{itemize}
\item $2d=4$. Smooth quartic in
$\mathbb{P}^3$
or a double cover of $\mathbb{P}^1 \times
\mathbb{P}^1$ branched
aloong a $(4,4)$ curve.

\item $2d=2$.  $S
\xrightarrow{2:1}\mathbb{P}^2 \supset$
sextic.
\end{itemize}
\end{example}

\section{Elliptic K3s}
\label{surfaces-section-elliptic-k3}

\begin{enumerate}
\item
$$
\xymatrix{
&Kum(E \times E^1)\ar[dl]\ar[dr]\\
\mathbb{P}^1\simeq E/(-1)& & \mathbb{P}^1
\simeq E^1/(-1)
}
$$

\item
$$
\xymatrix{
S \ar[r]^{\substack{\text{min} \\
\text{res}}}
&\tilde{S} \ar[r]^{2:1}
&\mathbb{P}^1 \supset \substack{\text{sextic}
\\ \text{w/ node}}.
}
$$
The elliptic pencil comes from the pencil of
lines in $\mathbb{P}^1$
through the node.

\item
$$
\xymatrix{
S \text{ elliptic}\ar[d]^{2:1}\\
\underbrace{F_1 \cup
F_2}_{\text{smooth}}\quad
X \text{ a RES w/ section}\ar[d]\\
\mathbb{P}^1
}
$$
\end{enumerate}

\begin{proposition}
\label{surfaces-proposition-elliptic-k3}
Let $S$ be a K3 surface.
\begin{enumerate}
\item $S$ is elliptic $\iff \exists  L \in
\text{Pic}(S)$ with $L^2=0$.

\item $S$ is elliptic with a section $\iff U
\xymatrix{\ar@{^{(}->}[r]^{
\text{primitive}}&} \text{NS}(S)$ [where $U$
is the hyperbolic lattice].

\item (If the Picard rank is big enough, the
surface will be elliptic.)
If $\rho(S) \geq 5$, then $S$ is elliptic
$\rho(S) \geq  12 \implies \exists $ a
section.
\end{enumerate}
\end{proposition}

To visualize the first point of the
proposition,
consdier the following example.

\begin{example}
\label{surfaces-example-smooth-quartic-elliptic}
A smooth quartic in $\mathbb{P}^3$ is
elliptic iff it contains a line.
\end{example}

\begin{remark}
\label{surfaces-remark-elliptic-k3}
\begin{itemize}
\item A K3 may admit more than one elliptic
fibration.

\item Since the base is always $\mathbb{P}^1$
[as you remember from the examples],
there will always be a singular fiber.
If $\underbrace{S}_{K3} \to C$ with $g(C) >
0$
then  $h^0(\Omega_C)>0$
but $H^0(\Omega_C)
\xymatrix{\ar@{^{(}->}[r]&} H^0(\Omega_S)$.

\item Any
$\underbrace{S}_{K3}
\xymatrix{\ar@{->>}[r]^{\text{flat}}&}
\mathbb{P}^1$
is an elliptic fibration.
\end{itemize}
\end{remark}

\section{Weierstrass model for K3}
\label{surfaces-section-weierstrass-model-k3}

\noindent
For an elliptic K3 $S \to \mathbb{P}^1$ with
section,
$R^1\pi_*\mathcal{O}_S =
\mathcal{O}_{\mathbb{P}^1}(-2)$.
So, a Weierstrass model is given by
$$
y^1=x^3 + A(t)x + B(t),\qquad t \in
\mathbb{P}^1
$$
with $\text{deg}A=8,\text{deg}B=12$.

\medskip\noindent
The volume form can be written locally as
$$
\frac{dx\wedge dt}{2y}.
$$
\section{Automorphisms}
\label{surfaces-section-automorphisms}

\noindent
[Automorphisms will act in differential
forms, cohomology, etc.
A fundamental question is if an automorphism
will
preserve the volume form or not.]

$$
\xymatrix{
&  \text{symplectic }\iff \sigma \text{ is
preserved}\\
\text{2 types}\ar@{-}[ur]\ar@{-}[dr]\\
&  \text{non- symplectic otherwise}
}
$$
(Purely non-symplectic if it acts as
multiplication by a
primitive $n-$-th root of 1.)

\noindent
We restrict to finite order.

\begin{proposition}
\label{surfaces-proposition-restrictions-on-orders}
Let $S$ be a K3 surface
with volume form $\sigma$.
\begin{enumerate}
\item If $f \in \text{Aut}(S)$ acts as
$f^*\sigma=\underbrace{\zeta_s}
_{\text{primitive}}\sigma$,
then $\varphi(n) \leq  20$ and all $n \neq
20$ satisfying
this appear.

If $n=2$, $\text{Fix}(f)$ contains only
curves
(no isolated points) or $\emptyset$.

If $n \geq 3$, $\text{Fix}(f)
=C_g \coprod R_1 \coprod … \coprod R_k
\coprod \{p_1,…,p_m\}$,
$g \geq 0$. (Nikulin, Sarti+Artebani,
Hoffmann Brandhorst.)

\item If $f^*\sigma=\sigma$,
then $\text{Fix}(f)=\{p_1,…,p_k\}$ with $1
\leq k \leq 8$
and $\text{ord}(f) \leq 8$.
\end{enumerate}
\end{proposition}

\begin{center}
\begin{tabular}{c|c c c c c c c}
n & 2&3&4&5&6&7&8\\
\hline
$\text{Fix}(f)$&8&6&4&4&2&3&2\\
$\rho(S)$ &9&13&15&17&17&19&19
\end{tabular}
\end{center}

\noindent
(See Salgado, Garbagnati.
Comparin-Prieto-Troncoso-Montero (on
involutions).)

\begin{example}
\label{surfaces-example-automorphisms}
Consider the elliptic K3 given by
[the Weierstrass model …]
$$
y^2=x^3+t^2x+t^{10}
$$
[we have written in this in a simpler way,
check degrees]
and the automorphism
$$
f:(x,y,t)\mapsto
(\zeta_7x,\zeta_7^5y,-\zeta_7t)
$$
(order 14).

\noindent
Note:
$$
\begin{aligned}
A(t,s)&=t^2s^6\\
B(t,s)&=t^{10}s^2\\
\Delta(t,s)&=4A^3+27B^2
=\text{const. }s^4t^6(s^{14}-t^{14}).
\end{aligned}
$$

\begin{itemize}
\item At $t=0$,  $a=2,b=10,\delta=6$.
$I_0^*$.

\item At $s=0$, $a=6, b=2, \delta=4$. $IV$.

\item At each root of $s^{14}-t^{14}$,
$a=b=0,\delta=1$, $I_1$.
\end{itemize}

\noindent
[Pictures of the fibers and how the
automorphisms act on them - fixed components
etc.]
\end{example}


\section{Summary}
\label{surfaces-section-summary}

\noindent
Slogan. We can parametrize the smooth pencils
of cubics as long
as we remove the RESs with bad singular
fibers.

$$
\begin{aligned}
\text{PGL}_3\text{ acts on }
\left\{
\substack{
\text{smooth} \\
\text{pencils of}\\
\text{plane cubics}}\right\}
&\leftrightarrow
\left\{\substack{
\text{RES} \\ \text{
w/ a section}}\right\}\\
\mathcal{P}
&  \mapsto
S_{\mathcal{P}}.
\end{aligned}
$$

\noindent
Question: can we parametrize orbits?
Does the geometry of the RESs
play any role?

\noindent
Answer: R. Miranda gives an answer
in the '80s using GIT.

$$
\begin{aligned}
\left\{
\substack{\text{RES } \\ \text{w/
section}}\right\}&\leftrightarrow
\left\{\substack{\text{Weierstrass eq.} \\
y^2=x^3+A_4(t)+B_6(t)}\right\}
\end{aligned}
$$

\noindent
Question. Can we parametrize such Weierstrass
equations?

\noindent
Answer. R. Miranda also gives an answer via
GIT.

\section{Quick intro to GIT}
\label{surfaces-section-git}

Consider a projective embedding
$X
\xymatrix{\ar@{^{(}->}[r]^{\mathcal{L}}&}
\mathbb{P}^N$
using a line bundle $\mathcal{L}$,
where $X$ and $\mathcal{L}$ equipped with an
action of a group $G$.
[The idea is to use this embedding and line
bundle
to construct the algebraic quotient.]

\begin{definition}
\label{surfaces-definition-git}
$$
X//G:=
\text{Proj}
\left(
\bigoplus_{d \geq 0}
H^0(X,\mathcal{L}^{\otimes d})^G\right).
$$
\end{definition}

\noindent
We have a map
$$
X \xrightarrow{\pi}X//G.
$$
[Notice that $\pi$ is not well-defined
at some places (this
is what we meant before with
the RESs with bad singular fibers
that we must remove),
namely where all the invariant sections
vanish. These points have a name:]

\begin{itemize}
\item $x \in X$ is {\it unstable} if and only
if
$\pi(x)$ is not defined.
\item  $X \supset X^{ss} \supset X^s$.
\end{itemize}

$$
\begin{aligned}
X^{ss}&\text{``='' points in $X$ where $\pi$
is defined}\\
X^s &  \text{``=''}\substack{
\text{semistable points w/ closed} \\
\text{orbit and finite stabilizer}}.
\end{aligned}
$$

\noindent
[The first step in constructing the GIT
quotient
is understanding what are the unstable and
semistable points.
This is what the naive questions above are
supposed to motivate:
we intend to find the unstable and semistable
points
of the corresponding GIT problem
in terms of the fibers of the surfaces.]

\medskip\noindent
Problem. $X=$ Grassmanian of lines
in
$H^0(\mathbb{P}^2,\mathcal{O}_{\mathbb{P}^2}
(3))$.
$G=\text{SL}(3)$ acts on $X$,
$\mathcal{L} \leftrightarrow $
Plücker embedding.

\medskip\noindent
Miranda proved:
\begin{enumerate}
\item A pencil $\mathcal{P}$ of plane cubics
is stable
if and only if $\mathcal{P}$ is smooth
and $S_{\mathcal{P}}$ contains only
reduced fibers.

\item Moreover, if $\mathcal{P}$ is smooth,
then $\mathcal{P}$ is semistable if and only
if
$S_{\mathcal{P}}$ has no fibers of type
$II^*, III^*$ or $IV^*$.
\end{enumerate}



\section{Natural generalizations}
\label{surfaces-section-natural-generalizations}

\begin{enumerate}
\item What about Halphen pencils/surfaces?

\item \leavevmode
[Recall the construction above on $X$ =
Grasmannian of lines
in
$H^0(\mathbb{P}^2,\mathcal{O}_{\mathbb{P}^2}
(3))$:
why
only in $\mathbb{P}^2$? Why only degree $3$?]
What about higher dimensional linear systems
in higher dimension?

\item What about a generalization of the
Weierstrass models?

\item [(1')] What  about higher genus
fibrations?
\end{enumerate}

\noindent
Now we give answers to those questions:

\begin{enumerate}
\item For $m=2$ [recall there is an $m$
associated to Halphen surfaces],
part of my PhD thesis.
For any $m$ (joint w/ Masafumi Hattori),
$$
\mathcal{P}
\mapsto  S_{\mathcal{P}}
\mapsto \mathbb{P}^1.
$$
\begin{theorem}
\label{surfaces-theorem-duplicate-2}
If $lct(S_{\mathcal{P}},F) > \frac{1}{2m}$
(resp. $\geq$)
for any fiber $F$,
then $\mathcal{P}$ is
stable (resp. semistable).
Type
$I_n$, II, III, IV, II$^*$
\end{theorem}

\begin{tabular}{c|c c c c c c c c c c}
type & $I^*_n$&$I_n$
&II&III&IV&$II^*$&$III^*$&$IV^*$&$m I_n$\\
\hline
lct&1/2&1&5/6&3/4&2/3&1/6&1/4&1/3&$1/m$
\end{tabular}

\noindent
For $t \in [0,1]$,
we compute
$$
\begin{aligned}
K_{\tilde{X}}+t\tilde{D}-\mu^*
(K_{\mathbb{C}^2}+tD
&=(1-2t)E_1+(2-3t)E_2+(4-6t)E_3
\end{aligned}
$$

\item Consider
$$
\underbrace{\text{Gr}(K,H^0(\mathbb{P}^n,
\mathcal{O}_{\mathbb{P}}(d)}_{
X_{K,d,n}}
\mathbb{x}\text{SL}(n+1)
$$
(HZ) $\mathcal{L} \in X_{K,d,n}$
is GIT unstable/not stable if and only if
there exists a choice of generators for
$\mathcal{L}$, $H_1,…,H_{k+1}$
such that
$H_1+…+H_{k+1}$ is GIT unstable/not stable.
[The action, which is compatible
also in symmetric power, exterior algebra,
etc.,
at the level of degree ¿? hypersurfaces will
help
determine the stable points.]
This implies (by Hacking, Kim-Lee) that the
pair
$$
\left(
\mathbb{P}^n,
\frac{n+1}{d(k+1)}(H^1+…+H_k)\right)
$$
is not lc (not klt).

Work in progress (HPZ).
The converse where the converse of the last
implication
holds if $n=2, d=3$
and for any $K$.

\medskip\noindent
Now we consider
$$
\left\{
\substack{\text{nets of cubics} \\
\text{through $p_1,…,p_7$.} }\right\}
$$
$$
\xymatrix{
Bl_{p_1,…,p_7}\mathbb{P}^2
\ar[r]^{\substack{\text{res} \\\text{sing.}
}}
\ar@{-->}[rd]_{2:1}
& Y \ar[d]^{2:1}\\
&\mathbb{P}^2 \supset
\underbrace{B}_{\text{quartic}}
}
$$

$M$ is GIT (semi)stable if and only if so is
$B$.
\end{enumerate}

\noindent
Now we introduce a new question that is not
in the above list:

\begin{enumerate}
\item [(4)] Let
$$
\xymatrix{
X\text{ RES w/ section}
\ar[d]^{\text{Blow up }f_1}\\
Y \ar[d]_f^{2:1}\\
\mathbb{P}^2 \supset B+p.
}
$$
Question. Given $X$,
classify all pairs $(B,p)$.

This morning (CGMVZ).
We did this for sufaces with $6$ double
fibers.
([For completeness:]
6II, 6$I_2$, 2$I_2+4 II$,
$2II+4I_2$,
$3I_2+3II$.)

\end{enumerate}

\noindent
Now let's describe the higher genus scenario:

\begin{enumerate}
\item [(1')]
Two constructions:
\begin{enumerate}
\item
$$
\xymatrix{
&\tilde{S}\ar[d]^{\text{Blow up
}f^{-1}(p)}\ar[ddl]\\
&S \text{ a K3}\ar[d]\\
\mathbb{P}^1 & \mathbb{P}^2 \cup  B \text{
sextic }+p
}
$$

\item (With Livia Campo.)
$T_{a,b}$ toric 3-fold with weight matrix
[we didn't introduce this gadget yet]
$$
\begin{pmatrix}
u&v & x & y & z\\
1& 1 & 1 & 0 & -a & -3a+b\\
0 & 0 & 1 & 1 & 3 \end{pmatrix}
$$
This gives
$$
\xymatrix{
T_{a,b}\ar[r]&\mathbb{P}^1_{(u,v)}\\
X\ar@{^{(}->}[u]\ar[ur]_{\substack{\text{of
bidigree} \\
\binom{-6a+2b}{a}}}
}
$$
$$
\leftrightarrow
$$
$$
\xymatrix{
X_{a,b}\ar[d]^{2:1}\\
\mathbb{F}_a
}
$$
When $a=0$ we can then parametrize
certain divisors in $\mathbb{P}^1 \times
\mathbb{P}^1$.

The case $(a,b)=(0,1)$
is the Horikawa model of the rational surface
$Bl_{13}\mathbb{P}^2$ where
the 13 points are the base locus of a pencil
of quartics
with a common node.
\end{enumerate}
\end{enumerate}

\noindent
[Finally we address the problem of
generalizing the
Weierstrass model:]

\begin{enumerate}
\item [(3)]
For Halpen surfaces,
at least 2 approaches.
$$
\xymatrix{
S \ar[d]^{|-m K_S|}\\
\mathbb{P}^1.
}
$$
$\xymatrix{\ar@{~>}[r]&}$
$$
\xymatrix{
S \ar[d]^{\substack{\text{gen. of} \\ \text{
deg. }m}}\\
\mathbb{P}^1 \times \mathbb{P}^1.
}
$$
We have an equation $y^m=…$.

The second approach is (j/ with S.
Rollenske).
[Diagram of blowing up.]
For $m=3$, we obtain $S$ as a relative cubic
inside a $\mathbb{P}^2$
bundle over $\mathbb{P}^1$.
\end{enumerate}


% --- End surfaces.tex ---

% --- Begin cremona-transformations.tex ---
\chapter{Factoring Cremona transformations}
\label{chapter-cremona-transformations}


\noindent
Factoring Cremona transformations.
Minicourse by Carolina Araujo and
Sokratis Zikas.
OPEGA 2026, UFPE.

\medskip\noindent
Notes at


\medskip\noindent
\textbf{Abstract.}
O grupo de Cremona em dimensão n é o grupo
das transformações birracionais do espaço
projetivo de dimensão n. O celebrado teorema
de Noether-Castelnuovo (1871-1901) afirma que
o grupo de Cremona em dimensão 2 é gerado
pelos automorfismos lineares e por uma única
transformação quadrática. Em dimensões
superiores, não há uma descrição
simples do
grupo de Cremona em termos de geradores, e a
situação é bem mais complicada. Por outro
lado, técnicas de geometria birracional, em
particular o MMP (Minimal Model Program),
fornecem uma maneira de fatorar
transformações de Cremona como
composições de
elos elementares. Essa teoria, conhecida como
"Programa de Sarkisov", tem se mostrado
extremamente útil no estudo do grupo de
Cremona em dimensão superior. Neste
minicurso, faremos uma introdução ao
MMP e ao
Programa de Sarkisov.


\section{Cremona group}
\label{cremona-transformations-section-cremona-group}

\noindent
Let $\text{Cr}_n(k):=\text{Aut}_k
k(x_1,…,x_n)$.

\begin{definition}
\label{cremona-transformations-definition-rational-map}
$f:X \dashrightarrow Y$ is {\it rational}
if it is defined in an open dense set $U_X
\subseteq X$.
It is {\it birational} if it admits an
inverse;
equivalently, if there is $U_Y \subseteq Y$
open dense
such that $f:U_X \to U_Y$ is an isomorphism.
\end{definition}

\noindent
Fix $n =2$ and consider the map
$(x_1,x_2) \mapsto
\left(\frac{1}{x_1},\frac{1}{x_2}\right)$.

\begin{example}
\label{cremona-transformations-example-stand-quadr-trans}
$$
\begin{aligned}
\sigma: \mathbb{P}^2 &\dashrightarrow
\mathbb{P}^2\\
(x_0,x_1,x_2) &\longmapsto
(x_1x_2,x_0x_2,x_0x_1).
\end{aligned}
$$

\noindent
Then $\sigma=\sigma^{-1}$, with
$$
\begin{aligned}
\{x_i=0\}\mapsto p_i,\qquad \qquad
\{x_0=0\}\mapsto (1:0:0).
\end{aligned}
$$

\noindent
$\sigma$: standard quadratic transformation.
\end{example}

\medskip\noindent
\textbf{Main question.}
Question (Enrique, 1895):
Is  $\text{Cr}_n(k)$ simple?
[i.e. does it have any nontrivial normal
subgroups?]
Do there exists any homomorphisms
$\text{Cr}_n(k) \to H?$

$\text{PGL}_{n-1}(k) \cong
\text{Aut}(\mathbb{P}^n)
\subseteq \text{Cr}_n(k)$.

\medskip\noindent
\textbf{Strategy.}
How to find homomorphism to a group $H$,
$\text{Cr}_n(k) \to H$.

\begin{enumerate}
\item Find a set of generators $G$ by
$\text{Cr}_n(k)$.

\item Get a set $R$ of relators.
$\text{Cr}_n(k)=\left\langle{}G|R\right\rangle{}$.

\item Map the generators to arbitrary
elements of $H$
and check that the relators are mapped to
$1_H$.
\end{enumerate}

\begin{theorem}[Noether-Castelnuovo, 1872]
\label{cremona-transformations-theorem-noether-castelnuovo}
$\text{Cr}_2(\mathbb{C})$ is generated by two
automorphisms of $\mathbb{P}^2$ and $\sigma$.
\end{theorem}

\begin{theorem}[Gizatulin, 1983]
\label{cremona-transformations-theorem-gizatulin}
Description of the relators with respect to
two generators
$\{\text{PGL}_3(k),\sigma\}$.
\end{theorem}

\begin{remark}
\label{cremona-transformations-remark-dim-2}
Contat-Lamy, 2010. $\text{Cr}_2(k)$ is not
simple for $k=\bar{k}$.
Lonjoi, 2017: any field.
[This is the proof that uses an action on an
infinite-dimensional hyperbolic space.
The same technique does not work for higher
dimensions.]
\end{remark}

\begin{theorem}[Hudson, 1927 \& Pan, 1999]
\label{cremona-transformations-theorem-hudson-pan}
Any set of [nonlinear] generators for
$\text{Cr}_n(k)$, $n\geq 3$,
is uncountable.
\end{theorem}

\noindent
… and we don't know any relators!


\section{MMP and Sarkisov theory}
\label{cremona-transformations-section-mmp-and-sarkisov}

\noindent
\textbf{Main idea.}
MMP is an algorithm.

$$
\xymatrix{
X:\substack{\text{smooth} \\
\text{projective}}
\ar[r]
&  \text{MMP}
\ar[r]
&  X_{\text{min}}: \substack{\text{mildly
singular} \\ \text{projective}}
}
$$
so that
\begin{itemize}
\item  $X \sim_{bir} X_{\text{min}}$
\item $X_{\text{min}}$ is the ``simpler''
than $X$
\begin{itemize}
\item $K_{X_{\text{min}}}$ is ``more
positive'',
\item $\rho(X_{\text{min}}) \leq  \rho(X)$.
\end{itemize}
\item The process is realized in
``elementary''steps.
\end{itemize}

\noindent
The outputs of MMP
are of two types:
\begin{enumerate}
\item minimal models, or
\item Mori fiber spaces (Mfs)
\end{enumerate}

\medskip\noindent
\textbf{Sarkisov program.}
An algorithm for decomposing birational maps
among Mfs into ``simpler'' maps.

[The idea of these maps is to copy the action
of
the standard quadratic transformation
$\sigma$
from Example \ref{cremona-transformations-example-stand-quadr-trans}:
we blow up three points (vertices of a
triangle, the $p_i$),
then we get a ``hexagonal'' arangement,
and contract three lines (the other three)
to get back at a triangle).]

$$
\xymatrix{
\text{triangle}\ar@{~>}[r]
&\text{hexagon}\ar@{~>}[r]
&\text{triangle}
}
$$
$$
\xymatrix{
\underbrace{\mathbb{P}^2}_{\text{triangle}}
\ar@{~>}[r]
&\underbrace{\mathbb{F}_1}_{\text{Hirzebruch}
}
\ar@{~>}[r]
&\ast
\ar@{~>}[r]
&\mathbb{P}^1 \times \mathbb{P}^1
\ar@{~>}[r]
&\ast
\ar@{~>}[r]
&\ast
\ar@{~>}[r]
&\underbrace{\mathbb{P}^2}_{\text{triangle}}
}
$$

\begin{theorem}[Corb, 1995 and Hacon-McKeman,
2011]
\label{cremona-transformations-theorem-corcb-hacon-mckeman}
Any birational map among Mfs can be factored
as
a composition of Sarkisov links
if and only if
the Sarkisov links are generators of the
groupoid
$$
\underbrace{\text{BirMori}
}_{\supseteq \text{Cr}_n(k)}
=\left\{f:X\dashrightarrow Y
\substack{f\text{ birational} \\ X,Y\text{
Mfs,}\\
X,Y \sim_{\text{bir}}\mathbb{P}^n}\right\}
$$
\end{theorem}

\begin{theorem}[Blanc-Lamy-Zimmerman]
\label{cremona-transformations-theorem-blanc-lamy-zimmerman}
Description of the relators among the
Sarkisov links.
\end{theorem}

\begin{theorem}[BLZ, 2021]
\label{cremona-transformations-theorem-blz-2021}
$\text{Cr}_n(k)$ is not simple for $n \geq
3$,
$\varphi_i$ only appearing in relations of
the form
$\varphi_i \circ X \circ \varphi_i^{-1} \circ
X^{-1}=\text{id}$,
$\varphi_i^2=\text{id}$.
$$
\begin{aligned}
\text{BirMor}(\mathbb{P}^n) &\longrightarrow
\mathbb{Z}/2\mathbb{Z} \\
\varphi &\longmapsto 1\\
\text{links}\neq \varphi&\longmapsto 0
\end{aligned}
$$

\end{theorem}

\section{Review of surfaces}
\label{cremona-transformations-section-surfaces}

\begin{definition}
\label{cremona-transformations-definition-surface}
A {\it surface} is a smooth projective
variety
of dimension 2.
\end{definition}

\begin{example}
\label{cremona-transformations-example-surfaces}
\begin{enumerate}
\item $S=\mathbb{P}^2$

\item
 $$
\xymatrix{
S=\mathbb{P}^1\times \mathbb{P}^1
\ar@{^{(}->}[r]
&z(xy-zw)\subset \mathbb{P}^3_{(x:y:z:w)}
}
$$

\item Blowup de $\mathbb{P}^2$ em $p$.
[Picture showing an arbitrary line
$\mathbb{P}^1$ contained
in the projective plane $\mathbb{P}^2$.
Given $P$ not in the line,
we can project any point $Q$ not in the line
to the line.
This gives a rational map
$\pi_P:\mathbb{P}^2\setminus
\{P\}\dashrightarrow \mathbb{P}^1$.
Consider the graph $\Gamma$ of $\pi_P$,
it is contained in $\mathbb{P}^2\setminus\{
P\} \times \mathbb{P}^1$.
Its closure, $\overline{\Gamma}$
is contained in $\mathbb{P}^2\times
\mathbb{P}^1$.
We have projections $p:\mathbb{P}^2\times
\mathbb{P}^1 \to \mathbb{P}^2$
and $q:\mathbb{P}^2\times P^1\to
\mathbb{P}^1$.
Note that $E=p^{-1}(P) \cong \mathbb{P}^1$
and $S\setminus E
\xrightarrow{\cong}\mathbb{P}^1\setminus\{P\}
$.
More generally,
$$
\tilde{S}=Bl_PS \xrightarrow{\pi}S
$$
$$
E=\pi^{-1}(P)\cong \mathbb{P}^1 \to P
$$
for $P \in S$.
\end{enumerate}
\end{example}

\begin{definition}
\label{cremona-transformations-definition-picard}
$S$ surface,
$\text{Pic}(S) = \text{Div}(S)/\sim \cong
\{\text{invertible sheaves}\}/\sim$.

$$
\text{Div}(S)
=\left\{\sum_{\text{finite}}n_iD_i:
n_i \in \mathbb{Z},
D_i \subset S\text{ irreducible
curve}\right\}.
$$
$$
\text{Div}(S) \iff D-D'=\text{div}(f)=(f)_0 -
(f_\infty),\qquad
f \in K(S)\setminus \{0\}.
$$
\end{definition}



\section{The role of the blowup in surface
birational geometry}
\label{cremona-transformations-section-role-of-blowup}

\noindent
Let $S$ and $S'$ be two (projective
nonsingular) surfaces.
\begin{itemize}
\item Any birational morphism $f:S \to S'$
is a composition of blowups.

\item Any birational map $f:S \dashrightarrow
S'$
factors as
$$
\xymatrix{
&  W \ar[dl]^{p}\ar[dr]_q\\
S\ar@{-->}[rr]&& S'
}
$$
where $W$ is a surface and $p,q$ are
compostions
of blowups.
\end{itemize}

\noindent
[We wish to find the simplest variety in the
same birational class.]

\begin{definition}
\label{cremona-transformations-definition-minimal-surface}
A surface $S$ is called {\it minimal} if
[if whenever we have a birational map from
$S$ to another surface then
that morphism is an isomorphism]
if any birational morphism $S \to S'$ is an
isomorphism.
\end{definition}

\section{Intersection theory for surfaces}
\label{cremona-transformations-section-intersection-surfaces}

\begin{theorem}
\label{cremona-transformations-theorem-intersection-form}
Let $S$ be a surface.
There exists a unique symmetric bilinear form
$\text{Pic}(S) \times \text{Pic}(S)\to
\mathbb{Z}$
such that
if $C$ and  $C'$ are two curves intersecting
transversally,
then $[C]\cdot[C']=\#(C \cap C')$.
\end{theorem}

\begin{example}
\label{cremona-transformations-example-intersection}
\begin{enumerate}
\item $S=\mathbb{P}^1$,
$\text{Pic}(S)=\mathbb{Z}[H]$ where
$H$ is a line (hyperplane).
$H^2=1$.

\item $S=\mathbb{P}^1 \times \mathbb{P}^1$,
$\text{Pic}(S)=\mathbb{Z}[H_1] \oplus
\mathbb{Z}[H_2]$
there $H_i$ is the class determined by either
projection.
We have
$$
\begin{cases}
H_1^2=0&=H^2_2 \\
H_1\cdot H_2 &= 1.
\end{cases}
$$
In this case $C \subset S$, $C=Z(F)\subset
\mathbb{P}^1_{(x:y)}
\times \mathbb{P}^1_{(z:w)}$,
$F \in k[x,y,z,w]$
bi homogeneous of degree $(d_1,d_2)$.

\item $S=Bl_P\mathbb{P}^2
\xrightarrow{\pi}\mathbb{P}^2 $,
$\text{Pic}(S)=\mathbb{Z}[\pi^*H]
\oplus\mathbb{Z}[E]$.
$$
\begin{cases}
\pi^*H \cdot \pi^*H&=H\cdot H=1\\
\pi^*H\cdot E&=0\\
E^2&=?
\end{cases}
$$
$$
\pi^* \ell=\tilde{\ell}+E,\qquad  \ell \text{
line }\subset \mathbb{P}^2
$$
$$
\begin{aligned}
E^2&=E\cdot(\pi^* \ell - \tilde{\ell})
=\underbrace{E\cdot \pi^*\ell}_{0}
-\underbrace{E\cdot \tilde{\ell}}_{1}=-1.
\end{aligned}
$$
\end{enumerate}
\end{example}

\noindent
In general [for blowups],
let
$$
\begin{aligned}
\pi:\tilde{S}=Bl_PS&\longrightarrow S \\
\mathbb{P}^1 \cong E &\longmapsto P
\end{aligned}
$$

$$
\text{Pic}(\tilde{S})=\pi^*\text{Pic}(S)
\oplus \mathbb{Z}E
$$
$$
\begin{cases}
\pi^*D\cdot\pi^*D'&=D\cdot D \\
\pi^*D\cdot E&=0\\
E^2&=-1.
\end{cases}
$$
\begin{definition}
\label{cremona-transformations-definition-minus1-curve}
A {\it $(-1)$ curve} in a surface $S$ is a
curve $C \subset S$
such the $C \cong \mathbb{P}^1$ and $C^2=-1.$
\end{definition}

\begin{theorem}[Castelnuovo's contractibility
theorem]
\label{cremona-transformations-theorem-castelnuovo}
Let $S$ be a surface and $C\subset S$ a
$(-1)$-curve.
Then there exists a surface $S' \ni P$ such
that
$S \cong Bl_PS'$ and via this isomorphism
$C$ is an exceptional curve.
\end{theorem}

\noindent
Notation. We say that $S\cong Bl_PS' \to S'$
is a contraction of $C$.

As a corollary,
\begin{lemma}
\label{cremona-transformations-lemma-corollary-castelnuovo}
A surface $S$ is minimal if $S$ does not
contain any $(-1)$ curves.
\end{lemma}

\section{MMP for surfaces}
\label{cremona-transformations-section-mmp-for-surfaces}

Let $S$ be a (nonsingular projective)
surface.

$$
\xymatrix{
&\text{Does $S$ contain any $(-1)$ curves?}
\ar[dl]\ar[dr]\\
\substack{\text{Não} \\ \text{Stop!}\\
\text{$S$ is a minimal}\\\text{surface}}
&  &
\substack{\text{Sim} \\ \text{Pick a $(-1)$
curve}\\
\text{Apply
Castelnuovo}}\ar@/^2.5pc/[ul]^{\substack{
\text{substitute} \\ S=S'\\\text{where
}S=Bl_PS'\\
\text{and repeat}}}
}
$$

\begin{definition}
\label{cremona-transformations-definition-numerical-equivalence}
Let $S$ be a surface and $D,D' \in
\text{Div}(S)$.
$D$ and $D'$ are {\it numerically equivalent}
if for any curve $C \subset S$
we have $D\cdot C= D'\cdot C$.
$\text{Num}(S)=\text{Div}/
\equiv=\text{Pic}(S)/\equiv$.
\end{definition}

\begin{theorem}
\label{cremona-transformations-theorem-num}
$\text{Num}(S)$ is a finite-rank free abelian
group.
Such a rank is called the {\it Picard number}
$\rho(S):=\text{rk}(\text{Num}(S))$.
\end{theorem}

\begin{remark}
\label{cremona-transformations-remark-pic-number}
[The rank of the blowup is one larger.]
$\tilde{S}=Bl_PS$ Then
$$
\begin{aligned}
\text{Pic}(\tilde{S})&=\pi^*\text{Pic}(S)
\oplus \mathbb{Z}\cdot E\\
\text{Num}(S)&=\pi^*\text{Num}(S) \oplus
\mathbb{Z}[E].
\end{aligned}
$$
\end{remark}

\noindent
[The following examples shows how we may
arrive
at different surface…]

\begin{example}
\label{cremona-transformations-example-arrive-at-different-surfaces}
Let $S= Bl_{P,Q}\mathbb{P}^2$ and denote
$E_P$ and $E_Q$ the exceptional divisors of
$P$ and $Q$.
Let $\ell$ be the line containing $P$ and
$Q$.
Consider a lift  $\tilde{\ell} \cong
\mathbb{P}^1$.
Then $\pi^* \ell = \tilde{\ell}+E_P+E_Q$
and $(\tilde{\ell})^2=-1$.
[One may show that in fact these are the only
$(-1)$ curves.]
Applying Castelnuovo we obtain a contraction
of $\tilde{\ell}$,
$S \to \mathbb{P}^1\times \mathbb{P}^1$,
where $\tilde{\ell}$ is mapped to a point
$R \in \mathbb{P}^1 \times\mathbb{P}^1$.
[It is AG jargon that the blowup of two
points is $\mathbb{P}^1\times
\mathbb{P}^1$
with …]
Since $\mathbb{P}^1 \times \mathbb{P}^1 \cong
Q \subset \mathbb{P}^3$,
[a doubly ruled surface?]
[We also have a map $Q \xrightarrow{\pi_R}
\mathbb{P}^2$.
It looks like the two lines passing through
$R$
(since it is a doubly ruled surface)
are mapped to two points under  $\pi_R$.]
\end{example}

$$
\xymatrix{
\substack{S \\
\text{rational}\\\text{surface}\\
\text{i.e., }S \sim_{bir}\mathbb{P}^2}
\ar@{~>}[r]^{\text{MMP}}
&\mathbb{P}^2,\mathbb{P}^1\times
\mathbb{P}^1, \mathcal{F}_n,
\qquad  n\geq 2.
}
$$
\begin{definition}
\label{cremona-transformations-definition-scroll}
A {\it scroll} is a surface $S
\xrightarrow{\pi}\mathbb{P}^1$
such that all fibers are isomorphic to
$\mathbb{P}^1$.
$S= \mathbb{P}_{\mathbb{P}^1}(\mathcal{E})$
where $\mathcal{E}$ is a vector bundle of
rank 2.

$$
\mathcal{E}=\mathcal{O}_{\mathbb{P}^1}\oplus
\mathcal{O}_{\mathbb{P}^1}(n)
$$
$$
\mathbb{P}(\mathcal{E}) \cong
\mathbb{P}(\mathcal{E} \otimes
\mathcal{O}_{\mathbb{P}^1}(n)).
$$
\end{definition}
$$
\mathbb{F}_n=\mathbb{P}_{\mathbb{P}^1}
(\mathcal{O}
\oplus \mathcal{O}(n)).
$$
\begin{exercise}
\label{cremona-transformations-exercise-inter-theor-hirze}
$\text{Pic}(\mathbb{F}_n)=\mathbb{Z}\cdot[f]
\oplus \mathbb{Z}\cdot[E]$
$$
\begin{cases}
f^2&=0\\
f\cdot E &= 1\\
E^2&=-n.
\end{cases}
$$

\end{exercise}

\section{Summary}
\label{cremona-transformations-section-summary}

$$
\xymatrix{
\substack{S\text{ rational} \\
\text{surface}}
\ar@{~>}[r]^{\text{MMP}}
&  \substack{S=\mathbb{P}^2\\\text{or} \\
S=\mathbb{F}_m, \quad
\substack{m=0 \\ \text{or }m \geq 2}}.
}
$$

\begin{exercise}
\label{cremona-transformations-exercise-rational}
Given $n_0>1$,
construct $S$ such that there exists MMP
$\leftrightarrow$
$\mathbb{F}_m$ for all $m=0$ or $2 \leq  m
\leq  n_0$.
\end{exercise}

\medskip\noindent
Consider $\mathbb{F}_m$ and $\sigma^2=-m$.
If $m\geq 1 \exists !$ curve $C \subset
\mathbb{F}_m$
such that $C^2<0$  $(C=\sigma)$.
If $m=0$, $\mathbb{F}_0=\mathbb{P}^1 \times
\mathbb{P}^1$.


$$
\xymatrix{
&\substack{\text{variety containing }
\\\pi^*
f,E,\widetilde{f_P}\cong\mathbb{P}^1\\
\tilde{\sigma}}
\ar[dl]^{Bl_P}_\pi
\ar[dr]^{\substack{\text{contraction} \\
\text{of $\widetilde{f_P}$}}}
\ar[d]\\
\substack{\mathbb{F}_m \\
\text{containing}\\
f,f_P,P}
\ar[d]\ar@{-->}[rr]^{\text{Sarkisov link}}
&\mathbb{P}
&\substack{\mathbb{F}_n \\
\text{containing}\\
\tilde{E},\tilde{\tilde{\sigma}},Q\\
n=?}
\ar[d]\\
\mathbb{P}^1&& \mathbb{P}^1.
}
$$
$$
\begin{aligned}
(\tilde{f}_P)^2&=\tilde{f}_P\cdot(\pi^*f-E)=-
1\\
\pi^*f_P …
\end{aligned}
$$

\noindent
[Computations show that $n=m+1$.]

\section{Positivity of the canonical class }
\label{cremona-transformations-section-positivity-canonical}

\noindent
Now we want to reinterpret the MMP
in terms of ``positivity of the canonical
class''.

\begin{definition}
\label{cremona-transformations-definition-canonical-class}
Let $S$ be a surface and $\Omega^1_S$
the Kähler differentials on $S$.
The canonic sheaf is
$\omega_S=\Lambda^2\Omega_S^1$
[because the dimension of $S$ is 2].
Thus, $[\omega_S]=K_S \in
\text{Pic}(S)=\text{Div}/\sim$
[is well-defined].
\end{definition}

\noindent
More concretely,
$\omega$ is a rational 2-form in $S$,
$\text{div}(\omega)=(\omega)_0-(\omega)
_\infty$.

\begin{example}
\label{cremona-transformations-example-canonical-class-of-p2}
\begin{enumerate}
\item $S=\mathbb{P}^2$,
$\text{Pic}(S)=\mathbb{Z}\cdot[H]$.
[We look for the zeroes and poles of the
form.]
[Choose an affine open]
$U=\mathbb{A}^2_{\left(
\underbrace{\frac{x}{z}}_{u},
\underbrace{\frac{y}{z}_{t}}\right) }\subset
\mathbb{P}^2$
where $\omega=du \wedge dt$.

Let $V=\mathbb{A}^2_{\left(
\frac{y}{x},\frac{z}{x}\right)}$.
[We look for the coefficient function of]
$\omega= ? dv \wedge ds$.

In $U \cap V$, [computations…]
Thus, $\omega=-\frac{1}{s^3}dv \wedge ds$.
[Thus, this form has a pole of degree 3 at
infinity.]

We conclude that $K_{\mathbb{P}^2}=-3H$.

\item $S=\mathbb{P}^1 \times \mathbb{P}^1$,
$\text{Pic}(S)=\mathbb{Z}\cdot H_1 \oplus
\mathbb{Z}\cdot H_2$,
[Exercise:] $K_S=-2H_1 - 2H_2$.
\begin{lemma}[Adjunction formula]
\label{cremona-transformations-lemma-adjunction}
Let $X$ be a smooth projective variety
and $Y \subset X$ an irreducible
hypersurface.
Then $K_Y=(K_X+Y)|_Y$.
\end{lemma}

\noindent
[Then]
$$
S=\mathbb{P}^1 \times \mathbb{P}^1 \cong S_2
\subset \mathbb{P}^3
$$
$$
K_S = (-4H+2H)|_S=-2H|_S
$$

\item $S= \mathbb{F}_m$,
$\text{Pic}(\mathbb{F}_m)=\mathbb{Z}\cdot f
\oplus \mathbb{Z}\cdot\sigma$… [More
computations].
\end{enumerate}
\end{example}

\medskip\noindent
\begin{definition}
\label{cremona-transformations-definition-positive}
A divisor $D \in \text{Div}(S)$
is {\it nef} if [it has nonegative
intersection number with
any curve in $S$] $D\cdot C \geq 0$ for any
curve $C \subset S$.
\end{definition}

\begin{example}
\label{cremona-transformations-example-positive-divisors}
\begin{enumerate}
\item (A very ample divisor.) Consider
$$
\xymatrix{
S \ar@{^{(}->}[r]^f
&\mathbb{P}^n
}
$$
where
$\text{Pic}(\mathbb{P}^n)=\mathbb{Z}\cdot[H]
$,
then $f^*H$ is very ample since
$f^*H\cdot C = f_*(C)\cdot H >0$.

\item (Ample divisor.) Recall that $D$ is
ample if $mD$
is very ample for some $m>0$.
Let $f:S \to \mathbb{P}^m$ be a morphism,
$D= f^* H$, $C \subset S$.
Then
$$
D\cdot C
=H \cdot f_\ast(C)
=
\begin{cases}
>0\qquad &\text{if }f_*(C)\text{ is a curve}
\\
=0\qquad &\text{if }f(C) = pt.
\end{cases}
$$
[Qualquer $D$ que seja o pullback da classe
do hiperplano
sob alguma $f$ é dito de semi-amplo.]
$D$ is called {\it semi-ample}.
\end{enumerate}
\end{example}

\section{Kleimann amplness theorem}
\label{cremona-transformations-section-kleimann}

\noindent
``nef= limit of ample''.

\medskip\noindent
Recall that we defined
$$
\text{Num}(S) = \text{Div}(S)/\equiv =
\mathbb{Z}^{\rho(S)}.
$$
Now let
$$
N'(S):= \text{Num}(S)
\otimes_{\mathbb{Z}}\mathbb{R}=\mathbb{R}
^{\rho(S)}.
$$
[Now we consider the following cones:]

$\overline{NE(S)}$,
the closed cone generated by all classes
$[C]$ of curves $C \subset S$,
called the {\it Mori cone}.
$\overline{NE(S)}\subset N'(S)$.

Let
$D \in \text{Div}(S)$.
$D$ is nef if and only if
$D\cdot\alpha \geq 0\forall \alpha \in
\overline{NE(S)}$.

$\text{Nef}(S)$ is the closed convex cone
generated by
the nef divisors = $(\overline{NE_1(S)})^\vee
\subset N^1(S)$.

\begin{theorem}[Kleimann's amplitude theorem]
\label{cremona-transformations-theorem-kleimann}
Let $D \in \text{Div}(S)$.
$D$ is ample if and only if $D\cdot \alpha>0$
for all
$\alpha \in \overline{NE_1(S)}\setminus\{0\}$
if and only if $D \in
\text{Int}(\text{Nef}(S))$.
\end{theorem}

\begin{definition}
\label{cremona-transformations-definition-face-extremal-rays}
Let $C \subset \mathbb{R}^n$ a closed convex
cone.
A {\it face} $F \subset N$ is a subcone such
that
$u,v \in N$ and $u+v \in F$ imply
that $u,v \in F$.

An {\it extremal ray} of $N$ is a face of
[maximal?] dimension.
\end{definition}

\begin{example}
\label{cremona-transformations-example-cone}
\begin{enumerate}
\item
$S= \mathbb{F}_m$,
$N_1(\mathbb{R}^2)$.
$K_S \cdot f=-2$,
$K_S \cdot \sigma=n-2$.

\item (The extremal ray of the exceptional
divisor
is an extremal ray.)

$$
\xymatrix{
\tilde{S}=Bl_PS\ar[r]^{\pi}
&S\\
\mathbb{P}^1\ar[r]\ar@{^{(}->}[u]
&pt\ar@{^{(}->}[u].
}
$$
Then $\mathbb{R}_{\geq 0}E$ is an extremal
ray of $\overline{NE_1(\tilde{S})}$.

$K_S\cdot E=-1$.
[ $(-1)$-curves always have $-1$ intersection
number with the canonical
divisor.]
\end{enumerate}
\end{example}

\begin{exercise}
\label{cremona-transformations-exercise-adjunction-again}
$C \subset S$, $C$ a $(-1)$-curve if and only
if
$C^2 <0$ and $C\cdot K_S<0$
if and only if $C^2=-1$ and  $K_S\cdot C=-1$.
\end{exercise}

\begin{remark}
\label{cremona-transformations-remark-minimal-surface}
If $K_S$ is nef, then  $S$ is a minimal
surface.
$\mathbb{P}^2$ is minimal but
$K_{\mathbb{P}^2}=-3H$.
\end{remark}

\begin{definition}
\label{cremona-transformations-definition-minimal-model}
$S$ is  {\it minimal model} if $K_S$ is nef.
\end{definition}

\begin{theorem}[Mori's cone theorem]
\label{cremona-transformations-theorem-mori-cone}
Let $S$ be a surface.
$$
\begin{aligned}
\overline{NE}_1(S)&=\overline{NE}_1(S)^{K_S
\geq 0}
+\sum_i\mathbb{R}_{\geq 0}[C_i]\\
&=\overline{NE}_1(S)^{K_S+A \geq 0}
+ \sum_{\text{finite}}\mathbb{R}_{\geq
0}[C_i], \qquad C_i\cong
\mathbb{P}^1.
\end{aligned}
$$
\end{theorem}

$$
\left[
\substack{\text{
Figure of $N^1(S)$ split into two regions} \\
\text{by $K_S$. One region is positive}\\
\text{and the other negative. In the
boundary}\\
\text{we see $K_S + \varepsilon A$, ample.}}
\right]
$$

\begin{theorem}[Classification of the $C_i$
's]
\label{cremona-transformations-theorem-classification}
\begin{itemize}
\item $C_i^2<0 \implies C_i$ is a $(-1)$
curve.
\item $C_i^2 = 0 \implies $ $S \cong
\mathbb{F}_m$ and $C_i=f$.
\item  $C_i^2 >0 \implies S \cong
\mathbb{P}^2$ and $C_i=\lambda H$.
\end{itemize}
\end{theorem}



\section{Summary }
\label{cremona-transformations-section-summary-duplicate-2}

\noindent
[Today: interperetation of the classical MMP
in modern language.]

\medskip\noindent
Recall:
\begin{itemize}
\item $D \in \text{Div}(S)$ is {\it nef} if
$D\cdot C \geq 0$ for all $C \subset S$.

\item [Take the blow up of a surface at a
point:]
$\tilde{S}=Bl_PS \xrightarrow{\pi}S$ with
exceptional divisor $E$,
then $K_S\cdot E=-1$ and $E^2=-1$.

\item $S= \mathbb{F}_m$, $K_S\cdot f=-2$.
$f^2=0$.

\item\leavevmode [The Mori cone is the convex
cone generated
by the curves in $S$:]
$\overline{NE}_1(S)=
\overline{\left\langle{}\text{curves
in $S$}\right\rangle{}}
\subset \underbrace{N^1(S)}_{=\mathbb{R}^e}
=\text{Num}(S) \otimes_\mathbb{Z}
\mathbb{R}$.

\item Theorem. Let $S$ be a surface. Then
\begin{enumerate}
\item (Cone theorem.) [
Picture of the hyperplane $K_S$ limiting two
regions,
a positive and a negative one. The negative
side is locally
polyhedral - while on the positive
side it may be round. We also consider a
small
perturbation $K_S+\varepsilon A$.]

\item (Contraction theorem.)
Any face on the negative side of the Mori
cone admits a contraction.
[See definitions below.]

\item (Classification of contractions
of $K_S$-negative extrmal rays.)
\begin{itemize}
\item If $C^2<0$, then  $\text{cont}_R=Bl_P$.
\item If $C^2=0$, then  $S= \mathbb{F}_m$
and $\text{cont}_R:\mathbb{F}_m \to
\mathbb{P}^1$.
\item If $C^2 \geq 0$,
then $S \cong \mathbb{P}^1$.
\end{itemize}
\end{enumerate}

\item \leavevmode [A face is a hyperplane
that encosta no cone.]
A {\it face} $F$ of a cone is such that  $v,u
\in \overline{ NE}(S)$
and $v+u \in F$ then $u,v \in F$.

\begin{definition}
\label{cremona-transformations-definition-contraction-of-face}
Let $F \subset \overline{NE}_1(S)$ be a face.
A {\it contraction} of $F$ is a morphism
$\text{cont}_F:S \to X$
with connected fibers,
$X$ normal such that
given a curve $C \subset S$,
$$
\text{cont}_F(C) = pt \iff
[C] \in F.
$$
\end{definition}

\begin{remark}
\label{cremona-transformations-remark-contractions}
$\text{cont}_F$ may or not exist, but if it
exists it is unique.
\end{remark}
\end{itemize}

\section{Modern version of MMP}
\label{cremona-transformations-section-modern-mmp}

\noindent
[Instead of asking whether there are any
$(-1)$ curves,
now we ask if the canonical divisor is
positive:]

$$
\xymatrix{
&\text{Is $K_S$ nef?}
\ar[dl]\ar[dr]\\
\substack{\text{Yes} \\ \text{Stop!}\\
\text{$S$ is a minimal}\\\text{model}}
&  &
\substack{\text{No} \\
\text{Pick a $K_S$ negative}\\
\text{extremal ray}\\
\text{and take}\\
f=\text{cont}_R:S \to X.}
\ar[dr]\ar[dl]\\
&
\substack{\dim X=2 \\
\text{Then} \\
S=Bl_P\underbrace{S'}_{X}}
\ar[uu]^{\substack{S:=S' \\
\rho(S')\\
\rho(S)=1}}.
&&
\substack{\dim X <2 \\
\text{Stop!}\\
\text{In this case}\\
f:S \to X\\
\text{is a Mori fiber space}\\
\text{second item in}\\
\text{classification part of thm}}.
}
$$


\begin{example}
\label{cremona-transformations-example-hirzebruch}
$S=\mathbb{F}_m$,
we computed that $K_S=-2\sigma -(2+m)f$.
$K_S\cdot f=-2$,
$K_S\cdot \sigma = 2m-2-m=m-2$.
\end{example}

\noindent
[Notice the diagram is more complicated now:
there are more options.
Why is this formulation convenient?]

Now fix a divisor $D \in \text{Pic}(S)$.
We may consider now the $D$-MMP:
$$
\xymatrix{
&\text{Is $D$ nef?}
\ar[dl]\ar[dr]\\
\substack{\text{Yes} \\ \text{Stop!}\\
\text{$S$ is a $D$-minimal}\\\text{model}}
&  &
\substack{\text{No} \\
\text{Pick a $D$ negative}\\
\underbrace{\text{extremal ray}}_{\text{may
not exist!}}\\
\text{and take}\\
f=\text{cont}_R:S \to X.}
\ar[dr]^{\substack{\text{type} \\\text{A}}}
\ar[dl]_{\substack{\text{type} \\
\text{B}}}\\
&
\substack{\dim X=2 \\
\text{Then} \\
f:S \to S'}
\ar[uu]^{\substack{
X\text{ may be singular}\\
D:=f_* D}}
&&
\substack{\dim X <2 \\
\text{Stop!}\\
\text{In this case}\\
f:S \to X\\
\text{is a $D$-Mori fiber space}\\
\text{second item in}\\
\text{classification part of thm}}.
}
$$

Some remarks:
\begin{itemize}
\item The final result of MMP (or $D$-MMP)
may not be unique,
but its type (A or B)
depends only on the birational class of $S$
(or  $D=K_S$).

\item When MMP ends in B
(i.e.  $S$ is a [rational] ruled surface)
the result is in general not unique.
\end{itemize}

\begin{example}
\label{cremona-transformations-example-mmp-rational-surface}
When $S$ is a rational surface,
the possible results are wither
$\mathbb{P}^2 \to pt$
or
$\mathbb{F}_m \to \mathbb{P}^1$.
\end{example}


\medskip\noindent
When does a $D$-MMP exist?

\begin{itemize}
\item $S$ toric.
($D$-MMP exists for any  $D$.)

\item $S$ Mori dream space (MDS).
($D$-MMP exists for any $D$.)
[This is the reason for the name.]
If $\text{Pic}(S)=\mathbb{Z}^{\oplus m}$,
the {\it Cox ring} is
 $\text{Cox}(S)=\bigoplus_{D \in
 \text{Pic}(S)}H^0(S,D)$.
$S$ is a {\it Mori dream space} if
$\text{Cox}(S)$ is finitely generated.

\item Let $S$ be any surface,
$D=K+\varepsilon A$,
$A$ ample.
[With such a $D$ we always guarantee that the
program exists.]

\begin{remark}
\label{cremona-transformations-remark-ample}
If $D \in \overline{NE}(S)$,
the $D$-MMP always ends in a $D$-minimal
model.
\end{remark}
\end{itemize}

\noindent
[Recall, the classical version consists of
a sequence of blowdowns (or blowups)
starting from $S$ and finishing in $S'$.]
$$
\xymatrix{
S \ar[r]&\cdots
\ar[r]
&S' \ar[d]_g\\
& &\underbrace{X}_{\substack{\text{$D$-ample}
\\ \text{model}}}
}
$$
$D'=g^*AX$, then
$D'$ is nef (semi-ample).

\medskip\noindent
[Extra: A K3 is a MDS is it has a trivial
automorphism group
(for $D$-MMP).]

\section{Geography of ample models}
\label{cremona-transformations-section-geography}

\begin{example}
\label{cremona-transformations-example-toric-surface}
[The following is a toric surface;
we worked on this two lectures ago.]
$S= Bl_{P,Q}\mathbb{P}^2$
[The cone has dimension 3, but we can make a
transverse section.
Imagine that the rest of the cone is behind
the screen:]
$$
\xymatrix{
  &           &  \tilde{\ell}
\ar@{-}[dl]\ar@{-}[dr]\\
  &  \tilde{\ell}_Q\ar@{-}[ddl]
\ar@{-}[dr]\ar@{-}[rr]&  &  \tilde{\ell}_{P}
\ar@{-}[ddr]\ar@{-}[dl]\\
  &                          & H\\
E_P\ar@{-}[urr]\ar@{-}[uur]
\ar@{-}[rrrr]& & & &
E_Q\ar@{-}[ull]\ar@{-}[uul]
}
$$

[$H$ should be in the center of the of the
quadrilateral
with vertices $E_P$,  $E_Q$  $\tilde{\ell}_P$
and $\tilde{\ell}_Q$.
$Nef(S)$ is the triangle with vertices
 $\tilde{\ell}_Q, \tilde{\ell}_P$ and $H$.]
\end{example}

$$
\xymatrix@C=8em{
S\ar[r]^-{f}_{\text{blowups}}
&\underbrace{S'}_{\substack{
D\text{-minimal}\\\text{model}}}
\ar[r]^-{\varphi |mD|}
&\underbrace{X}_{\substack{
\text{$D$-ample} \\ \text{model}}}=X(S,D).
}
$$
\begin{definition}
\label{cremona-transformations-definition-mori-equivalence}
Let $D, D' \in \overline{NE}(S)$.
$D \sim_{\text{Mori}}D'$ if
the morphisms
$$
S \to X(S,D)\text{ and }
S \to X(S,D')
$$
are the same
(up to $X(S,D) \cong X(S,D')$).

\end{definition}

[Inside $Nef(S)$ the model is itself.
In the lower edge of the triangle,
the model is $\mathbb{P}^2 \to pt$.
In the edge from $\tilde{\ell}_p$ to
$\tilde{\ell}$,
$\mathbb{P}^1 \times \mathbb{P}^1
\xrightarrow{p_1}\mathbb{P}^1$.
In the edge $\tilde{\ell}_p$ to $E_Q$,
$\mathbb{F}_1 \to \mathbb{P}^1$.]

[This figure shows us all the Mfs that may be
obtained from $S$.
The results are codified in the faces of the
cone decomposition.
The strategy of the proof of Sarkisov theorem
is to move around in the boundary of the
cone;
that is, passing from one Mfs to another.]

\begin{theorem}[Sarkisov program]
\label{cremona-transformations-theorem-sarkisov}
Any birational map of Mfs is a composition of
{\it Sarkisov links}, which may be either of
the following:
\begin{enumerate}
\item
$$
\xymatrix{
&\mathbb{F}_1
\ar[dl]\ar[dr]&\\
\mathbb{P}^2\ar[d] & & \mathbb{P}^1\\
\operatorname{pt} & &
}
$$

\item
$$
\xymatrix{
&  \tilde{S}\ar[dl]_{Bl_P}\ar[dr]\\
\mathbb{F}_m \ar[d]\ar@{-->}[rr]&&
\mathbb{F}_{m \pm 1}\ar[d]\\
\mathbb{P}^1 & & \mathbb{P}^1
}
$$
\item
$$
\xymatrix{
&\mathbb{F}_1\ar[dl]\ar[dr]\\
\mathbb{P}^1 & & \mathbb{P}^2\ar[d]\\
& & \operatorname{pt}
}
$$

\item  $\mathbb{F}_0 = \mathbb{P}^1 \times
\mathbb{P}^1$,
$$
\xymatrix{
\mathbb{F}_0\ar[d]_{p_1} \ar@{=}[r]&
\mathbb{F}_0\ar[d]^{p_2}\\
\mathbb{P}^1&\mathbb{P}^1
}
$$
\end{enumerate}
\end{theorem}

\noindent
[In fact, the Sarkisov links appear when
we pass from one chamber to another in the
Mori cone.]




\section{Proof of Sarkisov theorem}
\label{cremona-transformations-section-proof-sarkisov}

\begin{itemize}
\item
Let $D \in N^1(S)_d$ be of the form $K_S +
A$, $A$ ample.

\item A birational map $f:S \to T$ is called
a
 {\it $D$-minimal model} if it is the result
 of a $D$-MMP.

\item A map $g:S \to B$ is called an
{\it ample model of $D$} if it factors from a
$D$-minimal model
$f:S \to T$ followed by the contraction of
$f_*D$.

\item (Geography of ample models.)
We pick two ample divisors $A_1,…,A_k, K_S
\not\in \overline{NE}(S)$,
 $$
\mathcal{C}=
\{D=A_0K_S + \sum_{i=1}na_iA_i:
a_0,…,a_k \geq 0\}\cap \overline{NE}(S).
$$
$D_1,D_2$ are Mori equivalent if they have
the same ample model.

\item A {\it Mori chamber} is an equivalence
class
$(\mathcal{A}_i)_{i \in I}$ with ample model
$f:S \to T_i$. $\mathcal{A}_I$ is a {\it big
chamber}
if $f_i$ is birational.
\end{itemize}

\begin{proposition}
\label{cremona-transformations-proposition-mori-chambers}
\begin{enumerate}
\item The chamber decomposition
[of the Mori cone $\mathcal{C}$ in Mori
chambers]
is finite.
(Cone theorem + $K_S \not\in \overline{
\text{NE}}(S)$.)

\item The chambers are convex subcones.
$(\mathcal{A}_i=
\left\langle{}f_i^*\text{Nef},f_i\text{-exceptional
divisors}\right\rangle{}$.

\item $\mathcal{A}_i$ big chamber if and only
if
\label{cremona-transformations-item-3}
$\dim \mathcal{A}_i = \dim \mathcal{C}$.

\item If $\mathcal{A}_i$ is big,
$\overline{\mathcal{A}_i}=
\{D \in \mathcal{C}:f_i:S \to T_i\text{ is a
minimal $D$-model}\}$.
(ample $\to$ nef).

\item If
$\overline{\mathcal{A}_i}\cap\mathcal{A}_j
\neq \emptyset$
there exists $\rho_{ij}:T_i \to T_j$ making
the diagram
$$
\xymatrix{
A \ar[r]^{f_i}\ar[dr]_{f_j}&  T_i\\
& T_j
}
$$
commute and $\rho(T_i)-\rho(T_j)
=\dim(\mathcal{A}_i)-
\dim(\overline{\mathcal{A}_i}
-\dim(\overline{\mathcal{A}_i} \cap
\mathcal{A}_j)$.
\end{enumerate}
\end{proposition}


\begin{proof}[Proof of Sarkisov program]
[Start with a birational map of Mfs]
$$
\xymatrix{
& W\ar[dl]_{p_1}\ar[dr]^{p_2}\\
S_1\ar@{-->}[rr]^{\varphi}\ar[d]_{\eta_1}
&&S_2\ar[d]^{\eta_2}\\
B_1
&&B_2
}
$$
$\mathcal{C}_w$, there exist $D_1,D_2 \in
\mathcal{C}$
such that $\eta_i \circ p_i$ is an ample
model of $D_i$.
By (\ref{cremona-transformations-item-3}), $D_i \in \partial
\overline{\text{NE}}(W)$.

[The idea is to show that there exists a way
that joins
those two points which contained completely
in the
boundary.
After showing it exists,
we show that this path is in fact a Sarkisov
link decomposition.]

The visible boundary is
$$
\partial^+ \mathcal{C}
=\left\{D \in \partial \mathcal{C}:
\substack{ D= \partial \mathcal{C} \cap
[K_W,D^0 ]\\
\text{for some }D^0 \in
\text{Int}(\mathcal{C}) }\right\}
\subseteq \partial\mathcal{C}.
$$
[Geometrically, we know that the canonical
divisor
is not in the cone -- because it is not
ample.]

\begin{lemma}
\label{cremona-transformations-lemma-boundary}
\begin{itemize}
\item The visible boundary is contained in
$\partial \overline{
\text{NE}}(W)$.

\item $\dim \partial^+\mathcal{C}=\dim
\mathcal{C}-1$.
\item $\partial^+\mathcal{C}$ is connected.
\end{itemize}
\end{lemma}

\begin{proof}
$$
\begin{aligned}
D \in \partial^+ \mathcal{C}_1 D&=
t K_W+(1-t) A_0,\qquad  t \in [0,1)\\
A_0 \in \text{Int}\mathcal{C}
&  \implies A_0 = \sum a_i A_i,\qquad a_i>0,
\forall i=1,…,k\\
D\not\in \partial \mathcal{D} &\implies
D \in \partial \overline{\text{NE}}(W)
\end{aligned}
$$
where $\mathcal{D}:=\{a_0K_W+\sum a_iA_i:a_i
\geq 0\}$
which appears in the definition of
$\mathcal{C}$.
\end{proof}

\noindent
Let $V_2 \subseteq N^1(W)$
be a 2-dimensional affine space.
Take a slice of the cone:
$\mathcal{C} \cap V_2=\mathcal{P}$.
[This space has to be general in the sense
that]
$\text{codim}_{V_2}(\mathcal{A}_i \cap V_2)=
\text{codim}_{\mathcal{C}}\mathcal{A}_i$.
[We obtain a polygon with a chamber
decomposition.
Let's see what happens at a point $\Theta$
where two
(or more) chambers meet.
Zooming in, we see a some lines
$\mathcal{O}_0,…,\mathcal{O}_k$
meeting at $\Theta$, which determine chambers
$\overline{\mathcal{A}}_1,
\overline{\mathcal{A}}_2,…,
\overline{\mathcal{A}}_k$.]

[The $\mathcal{O}_i$ are the intersections of
the sheaves.
The $\mathcal{O}_i$ are not chambers
themselves:]

\begin{proposition}
\label{cremona-transformations-proposition-o-are-not-chambers}
 \begin{enumerate}
\item $0<i<k$, $\mathcal{O}_i \subseteq
\mathcal{A}_i$ or
$\mathcal{O}_i \subseteq \mathbb{A}_{i+1}$.

\item $D_0 \in \text{ rel int}\mathcal{O}_0$,
the ample model of $D_0$ is a Mfs.

\item $k<3$.

\item
Let $\mathcal{A}$ be the chamber such that
$\Theta \in \mathcal{A}$.
If there exists $i$ such that
$\dim(\overline{\mathcal{A}}_k \cap
\mathcal{A})=0$
then there exists a Sarkisov link between the
Mfs of Proposition
\ref{cremona-transformations-proposition-mori-chambers}.
\end{enumerate}
\end{proposition}

\begin{proof}
$$
\begin{aligned}
\text{Type I}
&\leftrightarrow k=2, \mathcal{O}_q
\subseteq \mathcal{A}_2\\
\text{Type II}&\leftrightarrow k=3\\
\text{Type III}&\leftrightarrow
k=2,\mathcal{O}_q \subseteq \mathcal{A}_1\\
\text{Type IV}&\leftrightarrow k=1.
\end{aligned}
$$
\end{proof}
\end{proof}






\section{About simplicity of the Cremona
group}
\label{cremona-transformations-section-simplicity-Cremona-group}

\medskip\noindent
\textbf{Título.}
On simplicity of the Cremona groups.

\medskip\noindent
\textbf{Palestrante.}
Sokratis Zikas (IMPA).

\medskip\noindent
\textbf{Data.}
27 Fevereiro 2026.

\medskip\noindent
\textbf{Resumo.}
The Cremona group or rank $n$ is the group of
birational transformations of the projective
space of dimension $n$.
One of the central questions regarding these
groups, dating all the way back to the 19th
century and settled only recently, is whether
they are simple.
In this talk I will review the history of the
problem, as well as give a brief overview of
the theory of elementary relations of
Blanc-Lamy-Zimmermann, that led to the proof
of the non-simplicity of higher-rank Cremona
groups.
We will also see how this theory was later
used to prove even more exotic properties of
these groups.

Recall that the {\it Cremona group} is
defined by
$$
\text{Cr}_n(k)=
\{f:\mathbb{P}^n \dashrightarrow
\mathbb{P}^n:
f\text{ is birational}\}.
$$
Question by Enriques (1865):
is $\text{Cr}_n(k)$ simple?

We know that $\text{Cr}_n(k)$ is not simple
in the following cases:
\begin{enumerate}
\item $n=2$, $k$ any field.
Contat-Lamy '13, Lonjou '16.

\item $n=2$, $k$ perfect field
such that there exists a Galois orbit of size
8.
Lamy-Zimmermann '20.

\item $n \geq 3, k=\mathbb{C}$.
Blanc-Lamy-Zimmermann '21.
\end{enumerate}

[The second and third item they use a
different method
to define the automorphisms.]

\begin{theorem}
\label{cremona-transformations-theorem-cr3}
There is an isomorphism
$\text{Cr}_3(\mathbb{C}) \simeq
G \ast (\ast_J \mathbb{Z}/2\mathbb{Z})$
where $J$ is uncountable.
\end{theorem}

With corollaries:

\begin{lemma}
\label{cremona-transformations-lemma-corollary}
$\text{Aut}(\text{Cr}_3(\mathbb{C}))$
is not generated by inner automorphisms
nor automorphisms of the field $\mathbb{C}$.
\end{lemma}

Déserti '06:
$\text{Aut}(\text{Cr}_2(\mathbb{C}))$
is generated by such automorphisms.

\medskip\noindent
[We shall prove Theorem \ref{cremona-transformations-theorem-cr3}.]

Define a groupoid by
$$
\underbrace{\text{BirMori}(\mathbb{P}^n)}
_{\text{Cr}_n(\mathbb{C})}
=
\left\{f:X \dashrightarrow Y
:\substack{X,Y \text{ Mfs}, X,Y
\sim_{\text{bir}}\mathbb{P}^n \\
f\text{ birational}}\right\}.
$$
\begin{definition}
\label{cremona-transformations-definition-rank-r-fibration}
Let $r \geq 1$.
A morphism  $\eta:X \to B$ is a {\it rank $r$
fibration}
if
\begin{enumerate}
\item $\dim X > \dim B, \rho(X)-\rho(B) = r$.
\item $X$ is a MDS over $B$
\item $-K_X$ is $\eta$-big.
\item The result of any
$D$-MMP over $B$ is
$\mathbb{Q}$-factorial is terminal.
\item $\exists \Delta_B \geq 0$
such that $(B,\Delta_B)$ is kld.
\end{enumerate}
\end{definition}

In the first case there exists a 2-ray game.
In case (2) the 2-ray game ends.
In the last 3 cases the 2-ray game
stays in the category of MMP.

\begin{itemize}
\item Rank 1 fibration = Mfs.
\item Rank 2 fibration = Sarkisov link.
\end{itemize}

\begin{theorem}[BLZ]
\label{cremona-transformations-theorem-blz2}
The relations among Sarkisov links are
generated by
by the dominated relations of rank 3
fibrations.
\end{theorem}

$$
\text{BirMori}=
\left\langle{}
\substack{\text{Sarkisov} \\ \text{links}}
:\substack{\text{relations given by} \\
\text{rank-3 fibrations}}\right\rangle{}.
$$

Let $(g,d) \in
\{(2,8),(6,9),(10,10),(11,14)\}$,
$C \in U_{g,d}\subseteq H_{g,d}^{\leq}$,
$J=\bigcup U_{g,d}$,
$X=Bl_C \mathbb{P}^3 \dashrightarrow
\mathbb{P}^3$.

\begin{proposition}
\label{cremona-transformations-proposition-rank-2}
Let $X \to \mathbb{P}^3 \to
\Spec(\mathbb{C})$
be a rank-2 fibration.
$$
\xymatrix{
X\ar@{-->}[rr]\ar[d]
&&X\ar[d]\\
\ar[dr]\mathbb{P}^3\ar@{-->}[rr]^{\chi_C}
&&\mathbb{P}^3 \ar[dl]\\
&\Spec(\mathbb{C}).
}
$$
with $\chi_C^2=\text{id}$.
\end{proposition}

$$
\xymatrix{
X\ar[dd]\ar[dr]&&& X\ar[dl]\ar[dd]\\
&Z \ar[r]^{t \mapsto -t}& Z\\
\mathbb{P}^3 \ar@{-->}[rrr]&&&\mathbb{P}^3.
}
$$
\begin{proposition}
\label{cremona-transformations-proposition-no-rank-3}
There are no rank 3 fibrations
dominated $X \to \Spec(\mathbb{C})$.
\end{proposition}








\clearpage
\section{Automorphisms of quartic surfaces
and Cremona transformations}

\medskip\noindent
\textbf{Speaker.}
Carolina Araujo.

\medskip\noindent
\textbf{Place.}
IMPA.

\medskip\noindent
\textbf{Date.}
24 September, 2024.

\medskip\noindent
\textbf{Motivation.}

$X$ a smooth hypersurface of degree $d $,
$X\subset \mathbb{P}^{n+1}$. We want to
understand the group $\operatorname{Aut}(X)$.
These are invertible polynomial maps from $X$
to itself. $X$ is defined by a single
polynomial equation of degree $d$.

\begin{theorem}[Matsumura-Mousky,
1964]\leavevmode
	Except in two special cases, all
	automorphisms of $X$ come from automorphisms
	of the ambient space:
	
	If $(n,d)\neq (1,3),(2,4)$, there is a
	surjective map
	$$
	\operatorname{Aut}(\mathbb{P}^{n+1},X)
	\overset{\pi}{\longrightarrow}
	\operatorname{Aut}(X)
	$$
	where $\operatorname{Aut}(X,Y)$ means
	automorphisms of $X$ that stabilize  $Y$.
\end{theorem}

\medskip\noindent
\textbf{Exceptional cases.}

Let's look at the exceptionals cases.

\begin{enumerate}
	\item $(n,d)=(1,3)$. In this case
	$C=X_3\subset \mathbb{P}^2$ is an elliptic
	curve and we have
		$\operatorname{Aut}(C)$ is a
		semidirect product of $C$ and
		$\mathbb{Z}_m$, with $m=2,4,6$.
and
 $$\operatorname{Aut}(\mathbb{P}^2,C)\text{
 is finite.} $$

 Elements in $C$ are translations of the
 torus. The {\it translation of
 $x$ with respect to $p$ and $(x:y:z)$} is
 done by intersecting the curve with the line
 that joins $p$ and $(x:y:z)$ and reflecting.
 So we have a map
  $$
  t_p(x:y:z)=
  (F_1(x,y,z),F_2(x,y,z),F_3(x,y,z)
  $$
  We have created a {\it Cremona
  transformation} (=biholomorphic birational
  map?).
 
  \begin{definition}[Cremona group]
	$$\operatorname{Bir}(\mathbb{P}^n)=\{\varphi
	:\mathbb{P}^n
	\overset{\operatorname{bir}}
	{\longrightarrow}
	\mathbb{P}^n:\text{bimeromorphic
	map} \}$$
\end{definition}

And we then have the surjective map
$$
\operatorname{Bir}(\mathbb{P}^2,C)
\overset{\pi}{\longrightarrow}
\operatorname{Aut}(C)
$$


	\item $(n,d)=(2,4)$. Here
	$S=X_4=\mathbb{P}^3$ is a smooth quartic
	surface.

		\paragraph{Problem} (Gizatullin) Which
		automorphisms of $S$ are induced (not
		necesarily by automorphisms of
		$\mathbb{P}^3$) but at least by Cremona
		transformations of $\mathbb{P}^3$? ie. are
		restrictions of
		$\varphi\in
		\operatorname{Bir}(\mathbb{P}^3,S)$

\begin{remark}
	Related to K3 surface structure,
	$$
	\operatorname{Bir}(\mathbb{P}^3,S)
	\overset{\pi}{\longrightarrow}
	\operatorname{Bir}(S)\cong
	\operatorname{Aut}(S)$$
\end{remark}

$S$ is a K3 surface.
$H^{2}(S,\mathbb{Z})\cong \mathbb{Z}^{22}$
and the Picard group acts on this lattice.

\begin{definition}[Picard rank of $S$]
	$$
	\rho(S)=
	\operatorname{rk}(\operatorname{Pic}(S))
	\in\{1,\ldots,20\}
	$$
\end{definition}

\begin{itemize}
\item If $S$ is very general, then
$\rho(S)=1$ and
$\operatorname{Aut}(S)=\{1\}$.

\item We are interested in $\rho(S)\geq 2$.
\end{itemize}
\end{enumerate}

\begin{example}[Oguiso, 2012]\leavevmode 
	\begin{enumerate}
		\item $\rho(S)=2$. Any cremona
		transformation that stabilizes the quadric
		is the identity:
			$$\operatorname{Aut}(S)=\mathbb{Z}\qquad
			\operatorname{Bir}(\mathbb{P}^3,S)=\{1\}$$
	
			
		\item $\rho(S)=3$,
			$$
			\operatorname{Aut}(S)=
			\mathbb{Z}_2*\mathbb{Z}_2*\mathbb{Z}_2
			$$
			All automorphisms are induced by Cremona
			transformations, ie. there is a surjective
			map
			$$
			\operatorname{Bir}(\mathbb{P}^3,S)
			\overset{\pi}{\longrightarrow}
			\operatorname{Aut}(S)
			$$

		\item (Paiva, Quedo 2023) Constructed
		surfaces with $\rho(S)=2$,
		$\operatorname{Aut}(S)=\mathbb{Z}_2$ and
		$\operatorname{Bir}(\mathbb{P}^3,S)=\{1\}$.
	\end{enumerate}
\end{example}

\begin{theorem}[A-Paiva-(Socrates)
Zika]\leavevmode
	Solution of Giztallin S problem for
	$\rho(S)=2$.
\end{theorem}

\begin{remark}
	Not exactly, but "there is a moduli space
	for K3 surfaces of dimension $20-\rho(S)$".
	The generic case is $\rho(S)=1$.
\end{remark}

\medskip\noindent
\textbf{Last application.}

\begin{definition}
	A {\it K3 surface} is a smooth
	projective surface that is simply connected
	and has a nowhere vanishing symplectic form
	$\omega\in H^{0}(S,\Omega^2_S)$
\end{definition}

\medskip\noindent
\textbf{Lattices of $S$.}

A {\it lattice} is a
finitely-generated abelian group with a
nondegenerate pairin.

\begin{enumerate}
	\item $H^{2}(X,\mathbb{Z})\cong \mathbb{Z}
	\overset{\pi}{\hookleftarrow}
	\operatorname{Pic}(S)
	$. And if we tensor this with $\mathbb{C}$
	we get $H^{2}(C,\mathbb{C})$, which admits a
	Hodge decomposition.

		Let's study automorphisms of a K3 surface.
		Let $g\in\operatorname{Aut}(S)$. It yields
		an element $g^*$ that acts on cohomology,
		ie $g^*\in\mathcal{O}(H^{2}(X,\mathbb{Z}))$
		preserving the Hodge decomposition. This is
		called {\it Hodge isometry}.

\begin{theorem}[Global Torelli
theorem]\leavevmode
	\begin{itemize}
	\item The correspondence $g\mapsto g^*$ is
	injective
	
	\item If
	$\varphi\in\mathcal{O}(H^{2}(X,\mathbb{Z}))$
	is a Hodge isometry preserving the ample
	classes, then $\varphi=g^*$ for some
	$g\in\operatorname{Aut}(S)$.
	\end{itemize}
\end{theorem}
\end{enumerate}

\begin{example}
	$S\subset \mathbb{P}^3$ smooth quartic
	surface with $\rho(S)=2$. Using the
	equivalence of line bundles module
	isomorphism and curves modulo intersection,
	$$\operatorname{Pic}(S)= \left\langle{}H,C\right\rangle{} $$
	where $H$ is a hyperplane section of
	$\mathbb{P}^3$. Then
	$$Q=\begin{pmatrix} H^2& H\cdot C\\H\cdot C&
	C^2 \end{pmatrix} =\begin{pmatrix} 4&b\\b&2c
	\end{pmatrix} $$
	using that $2g(c)-1$, so the number in the
	lower-right on RHS is even.
\end{example}

\begin{remark}
	Hodge Index Them $Q$ is rank (1,1).
\end{remark}

\begin{definition}[Discriminant]
	$$r=\operatorname{disc}(S)=-\det Q$$
\end{definition}

\begin{proposition}
	$S$ general K3 surface (not needed that it
	is a quartic) with $\rho(S)=2$. Then
	$$\operatorname{Aut}(S)=\begin{cases}
		\{1\}\qquad &\text{(finite)}  \\
		\mathbb{Z}_2\qquad &\text{(finite,
		Dani-Ana)} \\
		\mathbb{Z}\qquad&\text{(infinite, Oguiso
		1)}  \\
		\mathbb{Z}_2*\mathbb{Z}_2\qquad
		&\text{(infinite, Dani-Ana)}
	\end{cases}$$
	where the first two are characterized by
	containing a en elliptic curve or a
	rational? curve, ir. $\exists
	D\in\operatorname{Pic}(S)$ such that
	$D^2=0,-2$. On the other hand, the last two
	cases are distinguised by $\not\exists
	D\in\operatorname{Pic}(S)$ such that
	$D^2=0,-2$.

	$$\exists
	\sigma\in\operatorname{Aut}(S)\text{ of
	order 2}\iff \exists  \text{ ample
	}A\in\operatorname{Pic}(S) \text{ such that
	}A^2=2.$$
\end{proposition}

\begin{remark}
	In low Picard numbers there are no
	symplectic involutions…?
\end{remark}

So to understand the surface we want to
understand those bundles and that is all in
the discriminant (not in the lattice
itself!).

\begin{remark}
	$\exists D\in\operatorname{Pic}(S)$ s.t.
	$D^2=k \iff x^2-ry^2=4k$  has integer
	solutions.
\end{remark}

Given the quadratic form we can find the
automorphisms.

\medskip\noindent
\textbf{Main theorem.}

\begin{theorem}[A-Paiva-Zika]\leavevmode
	$S\subset \mathbb{P}^3$ general smooth
	quartic surface with $\rho(S)=2$ and
	$\operatorname{disc}(S)=r$.
\begin{enumerate}
	\item (Negative answer to Gizatulla's
	problem) If $r>57$ or  $r=52$ then
		 $$
		 \operatorname{Bir}(\mathbb{P}^3,S)=\{1\}
		 $$
		 So we cannot realize any automorphism as a
		 Cremona transformation.

	\item If $r\leq 57$ and $r \neq 52$, then we
	get the full automorphism group, ie. a
	surjective map
		$$
		\operatorname{Bir}(\mathbb{P}^3,S)
		\overset{\pi}{\longrightarrow}
		\operatorname{Aut}(S)
		$$
\end{enumerate}
\end{theorem}

\medskip\noindent
\textbf{Birational geometry.}

Take the case of
$$
\xymatrix{
\operatorname{Bir}(\mathbb{P}^2)=
\left\langle{}\operatorname{Aut}(\mathbb{P}^2),
q\right\rangle{}\\
\operatorname{Aut}(\mathbb{P}^2)
\ar[u]_{\hookrightarrow}
}
$$
and the map
$$
\begin{aligned}
	q: \mathbb{P}^2
	&\overset{\operatorname{bir}}
	{\longrightarrow}
	\mathbb{P}^2
	\\
	(x:y:z) &\longmapsto
	\left(\frac{1}{x}:\frac{1}{y}:
	\frac{1}{z}\right)=(yz:xz:xy)
\end{aligned}
$$
which is well-known (Noether-Castelnuovo). So
that is a decomposition of automorphisms and
quadratics. Now in greater dimension, $n\geq
3$ we have

\begin{theorem}[Sakisov Program]\leavevmode
	(The theorem is much more general) Any
	$\varphi
	\in\operatorname{Pic}(\mathbb{P}^n)$ can be
	factorized with  {\it Sarkisov
	links}  $\varphi_i$:
	$$\xymatrix{
		\mathbb{P}^n
		\ar[r]^-{\varphi_1}
		\ar@/^2pc/[rrrr]^-{\varphi}
		&
		X_1\ar[r]^-{\varphi_2}\ar[d]
		&
		X_2\ar[r]\ar[d]
		&
		\cdots
		\ar[r]
		&
		X_k=\mathbb{P}^n\ar[d]
		\\
		&T_1&T_2&&
	}$$
\end{theorem}

Now look at $n=3$,
$\operatorname{Bir}(\mathbb{P}^3,S)\subset
\operatorname{Bir}(\mathbb{P}^3)$. This is a
Calabi-Yau pair:

\begin{definition}%
[Calabi-Yau pair]\leavevmode 
	A pair $(X,D)$
	\begin{itemize}
	\item Terminal projective variety.
	\item $K_X+D\sim0$ that is, a meromorphic
	top form that does not vanish on
	hypersurface and has simple pole on $D$.
	Then $(X,D)$ is called {\it log
	canonical}.
	\end{itemize}

Now take two Calabi-Yau pairs $(X,D_X)$ and
$(Y,D_Y)$.

 $$
\begin{aligned}
	\operatorname{div}(\omega_{D_X})&=-D_X\\
	\operatorname{div}(\omega_{D_Y})&=-D_Y
\end{aligned}
$$
We say that a birrational map $f:X\to Y$ is
{\it volume preserving} if
$f_*\omega_{D_X}=\omega_{D_Y}$.

\end{definition}


\begin{theorem}[Volume-Preserving]\leavevmode
	Everything like in the Sarkisov theorem but
	now maps are volume-preserving.
\end{theorem}

In our case, $(\mathbb{P}^3,S)$, we can
classify the v.p. Sakisov links from
$(\mathbb{P}^3,S)$. It starts by blowing up a
curve $C\subset S$. But this curve has genus
and degree very restricted, it's something
like
$$
(g(C),\operatorname{deg}(C))\in
\{(0,1),(0,2),\ldots,(11,10),(14,11)\}
$$

So for the main theorem, it was checked that
if $r>57$ there are no curves from the list.
And in the second item of the main theorem,
there exist these curves, for instance curve
$(14,11)$ for rank 56, then produce a link
that starts by blowing it up and magically
gives you the Cremona transformation that
restricts with automorphism with which you
started.

\begin{remark}
	So perhaps we expect the answer to G.
	problem to be almost never.
\end{remark}

\begin{remark}
	How does that blowing-up work?
\end{remark}

$$
\xymatrix{
X\ar[r]^-{\text{flops}}
\ar[d]_-{\operatorname{Bl}_C}
&
X\ar[d]^-{\text{contract?}}
\\
C\subset S\subset
\mathbb{P}^3\ar@{-->}[r]&
\mathbb{P}^3\supset
S'\supset C'
}
$$
So for example in case $(2,8)$ you obtain
something of the same type.

\section{Birational geometry of Calabi-Yau
pairs}

\medskip\noindent
\textbf{Speaker.}
Carolina Araujo.

\medskip\noindent
\textbf{Place.}
IMPA.

\medskip\noindent
\textbf{Date.}
January 31, 2025.

\vspace{2em}
\medskip\noindent
\textbf{Pairs $(X,D)$.}
Itaca's program (1970's): $U$ complex
algebraic variety. We want to find invariants
like Kodaira dimension. Compactify $U
\rightsquigarrow X$, and consider
$X\setminus U:=D$ (the boundary). The
$\Omega^q_X(\operatorname{log}D)$

\begin{theorem}[Itaca, 1977]\leavevmode
Kodaira dimension of the pair $(X,D)$ does
not depend on the choice of compactification
(and it exists ;)
\end{theorem}

\begin{definition}[Lu-Zhang, 2017]\leavevmode
$(X,D)$ is
{\it Brody-hyperbolic}
{\it (Mori-hyperbolic)} if there
is no nonconstant holomorphic (analytic)
morphism  $f:\mathbb{C} \to X\setminus D$
and the same holds for the open strata of
$D$.
\end{definition}

\begin{remark}
$(X,D)$ is Brody-hyperbolic then $K_X +
D$ is ample.
\end{remark}

\begin{theorem}[Sraldi, 2019]\leavevmode
$(X,D)$ Mori-hyperbolic then $K_X+D$ is
nef.
\end{theorem}

\begin{remark}\leavevmode
How to make sense of ``the birational
geometry of $(X,D)$"
\end{remark}

\medskip\noindent
\textbf{Calabi-Yau pairs.}

\begin{definition}%
[Calabi-Yau pair]\leavevmode
We relax a condition: we won't ask that $X$
is smooth, but that it is a terminal
projective variety. And also that $K_X+D
\sim 0$ (I think this ``would be equivalent
to" Kodaira dimension 0). Also ask that
$(X,D)$ is log canonica.
\end{definition}

\begin{remark}\leavevmode
The condition $K_X+D \sim 0$ says that
there is unique up to scaling volume form
$\omega_D \in \Omega^m_{\mathbb{C}(X)}$
that does not vanish and has a simple pole
along $D$, i.e.
$\operatorname{div}(\omega_D)=-D$.
\end{remark}

Now we study the birational geometry of CY
pairs:

\begin{definition}\leavevmode
Let $(X,D_X),(Y,D_Y)$ CY pairs and $f: X
\dashrightarrow Y$ a birational map. We get
a pullback
$f_*:\Omega^n_{\mathbb{C}(X)}
\overset{\cong}{\to}
\Omega^n_{\mathbb{C}(Y)}$. Then $f:(X,D_X)
\dashrightarrow (Y,D_Y)$ is
{\it volume preserving} if
$f_*(\omega_{D_X})=\lambda \omega_{D_Y}$.
\end{definition}

\begin{remark}[Valuative
characterization]\leavevmode

\end{remark}

\begin{definition}\leavevmode
$(X,D)$ CY pair,
$$\operatorname{Bir}(X,D):=\{f \in
\operatorname{Bir}(X) : f:(X,D)
\dashrightarrow (X,D) \text{ is volume
preserving} \}$$
\end{definition}

\medskip\noindent
\textbf{The Sarkisov-Program.}

$\operatorname{Bir}(\mathbb{P}^n)$ Cremona
group.

\begin{theorem}[Noether-Castelnuovo,
1870-1901]\leavevmode
The Cremona group of $\mathbb{P}^2$ has a
nice set of generators: maps of degree 1 and
$q$:
$$\operatorname{Bir}(\mathbb{P}^2) =
\left\langle{}\operatorname{Aut}(\mathbb{P}^2),
q\right\rangle{}$$
where $q$ is the {\it standard
quadratic transformation} given by $(x:y:z)
\overset{q}{\mapsto }(yz:xz:xy)$
\end{theorem}

You might think that because it has such a
nice group generators it'd be easy to study
this group. It's not the case.

Now in dimension 3:

\begin{theorem}%
[Hilda Hudson (1927)]\leavevmode
$\operatorname{Bir}(\mathbb{P}^n)$ does not
admit a set of generators of bounded degree.
\end{theorem}

\begin{theorem}[Sarkisov Program, Corti 1995;
Hazon-McKunam 2013 for $n \geq
4$]\leavevmode
	$$\xymatrix{
		\mathbb{P}^n
		\ar[r]^-{\varphi_1}
		\ar[d]
		\ar@/^2pc/[rrrr]^-{\varphi}
		&
		X_1\ar[r]^-{\varphi_2}\ar[d]
		&
		X_2\ar[r]\ar[d]
		&
		\cdots
		\ar[r]
		&
		X_k=\mathbb{P}^n\ar[d]
		\\
	\{\operatorname{pt}\}
	&T_1&T_2&&\{\operatorname{pt}\}
	}$$
	The intermediate varieties $X_i/T_i$ are
	called {\it Mori-fiber spaces},
	and $\varphi_i$ are
	{\it Sarkisov links}. So the
theorem is that a map (what map?) can be
factorized by Sarkisov links.
\end{theorem}

\begin{theorem}%
[Corti-Kalog(?) 2016]\leavevmode
Any volume-preserving birrational map between
Mori-fiber CY pairs is a composition of
volume-preserving Sarkisov links. Now every
intermidiate variety admits a divisor making
	it a CY pair:
	$$\xymatrix{
		(X,D_X)\ar[r]^-{\varphi_1}\ar[d]
		\ar@/^2pc/[rrrr]^-{\varphi}
		&
		(X_1,D_1)\ar[r]^-{\varphi_2}\ar[d]
		&
		(X_2,D_2)\ar[r]\ar[d]
		&
		\cdots
		\ar[r]
		&
		(Y,D_Y)\ar[d]
		\\
	T_X	&T_1&T_2&&T_Y
	}$$
\end{theorem}

\begin{remark}\leavevmode
\begin{itemize}
\item The Sarkisov links
$$
\xymatrix{
 \mathbb{P}^n\ar@{-->}[r]\ar[d]
 &
 X_1\ar[d]
 \\
 \operatorname{pt}&T_1
}
$$
are not classified.

\item Depending on $D$, the
volume-preserving Sarkisov links
$$
\xymatrix{
 \mathbb{P}^n\ar@{-->}[r]\ar[d]
 &
 X_1\ar[d]
 \\
 \operatorname{pt}&T_1
}
$$
can be classified.
\end{itemize}
\end{remark}

\begin{theorem}[A-Coti-Massareti]\leavevmode
$D \subset \mathbb{P}^n$ a \textit{general}
hypersurface of degree $n+1$. If you want
to study the birational group of the pair
$(\mathbb{P}^n,D)$,
$\operatorname{Bir}(\mathbb{P}^n,D)$, then
this group is not interesting. Because
$\operatorname{Bir}(\mathbb{P}^n,D)=
\operatorname{Aut}(\mathbb{P}^n,D)$.

And then also consider $D$ smooth instead
of general, and with Picard rank
$\rho(D)=1$.
\end{theorem}

\begin{theorem}[A-C-M]\leavevmode
$D \subset \mathbb{P}^3$ general (A1
singularity and $\rho(S)=1$) quartic with a
singularity.
$\operatorname{Bir}(\mathbb{P}^3,D)$ is a
semidirect product of $\mathbb{G}$ and
$\mathbb{Z}/2\mathbb{Z}$.
Here $\mathbb{G}$ is ``form of
$\mathbb{G}m$ over $\mathbb{C}(x,y)$".
\end{theorem}

\section{Comparison with the hyperk\"ahler
story}
\label{cremona-transformations-section-comparison-hyperkahler}

\medskip\noindent
Both stories say that birational
geometry is controlled by cones and by
wall crossing.
But the cones are different.

\medskip\noindent
In the Mori and Sarkisov story, the
basic cone is the Mori cone
$$
\overline{NE}_1(X),
$$
generated by curve classes.
Its dual is the nef cone.
The canonical class $K_X$ is the
driving functional.
The $K_X$-negative extremal rays are
contractible, and their contractions
produce divisorial contractions,
flips, or Mori fiber spaces.

\medskip\noindent
For Sarkisov, one studies a chamber
decomposition of a cone of divisors.
Two divisors lie in the same chamber
when they define the same ample model.
Moving from one chamber to an
adjacent chamber changes the
birational model.
The rank two wall crossings are the
Sarkisov links.
Thus
$$
\text{Mori chambers}
\Rightarrow
\text{Sarkisov links}
\Rightarrow
\text{factorizations}.
$$

\medskip\noindent
In the Calabi--Yau pair version,
the same Sarkisov philosophy remains,
but the boundary divisor and volume
form are part of the data.
The maps are volume-preserving
birational maps of Calabi--Yau pairs,
and the elementary pieces are
volume-preserving Sarkisov links.

\medskip\noindent
In the hyperk\"ahler story, the
canonical class is trivial.
So $K$ cannot drive the process.
Instead, the cone is the positive cone
inside $H^{1,1}(M)$, defined using
the Beauville--Bogomolov--Fujiki form
$q$:
$$
\mathcal C_M
=
\{x\in H^{1,1}(M):q(x)>0\}^{\circ}.
$$
Inside it sits the K\"ahler cone
$\mathcal K_M$.
Amerik--Verbitsky describe a locally
finite wall-and-chamber decomposition
of $\mathcal C_M$, with walls
orthogonal to MBM classes.

\medskip\noindent
If
$$
\varphi:M\dashrightarrow M'
$$
is bimeromorphic, then
$$
\varphi^*\mathcal K_{M'}
$$
is another chamber in the same
positive cone.
Crossing a wall gives a small
symplectic contraction and a
wall-crossing flop.
Amerik--Soldatenkov--Verbitsky show
that these wall crossings give a
finite decomposition into Hu--Yau
transformations.
In codimension two these are modeled
by Mukai flops.

\medskip\noindent
The dictionary is:
$$
\begin{array}{c|c}
\text{Mori/Sarkisov} &
\text{hyperk\"ahler} \\
\hline
\text{Mori cone, nef cone} &
\text{positive cone, K\"ahler cone} \\
\text{extremal rays, faces} &
\text{MBM walls} \\
\text{Mori chambers} &
\text{K\"ahler-type chambers} \\
\text{MMP contractions} &
\text{symplectic small contractions} \\
\text{Sarkisov links} &
\text{wall-crossing flops} \\
\text{Mori fiber spaces} &
\text{hyperk\"ahler models} \\
K_X\text{-negative geometry} &
K_M=0\text{ symplectic geometry} \\
\text{volume-preserving CY links} &
\text{Hu--Yau transformations}
\end{array}
$$

\medskip\noindent
Thus Carolina's side is the
MMP/Sarkisov chamber geometry of
Mori fiber spaces and Calabi--Yau
pairs.
Misha's side is the K\"ahler cone
wall geometry of hyperk\"ahler
manifolds.
The common pattern is that a
complicated birational map becomes a
path through chambers of a cone, and
the elementary steps are the wall
crossings.


% --- End cremona-transformations.tex ---

% --- Begin birational-maps-commutative-algebra.tex ---
\chapter{Birational maps from the commutative algebra viewpoint}
\label{chapter-birational-maps-commutative-algebra}


algebra viewpoint}

\noindent
Minicourse by Hamid Hassanzadeh,
OPEGA 2026, UFPE.

\medskip\noindent
O estudo da equivalência de variedades
algébricas a partir da biracionalidade é um
dos temas centrais da Geometria Algébrica.
Nessa direção, a classificação de variedades
pode ser vista como parte da chamada
geometria birracional. De modo análogo,
podemos concentrar-nos na estrutura dos mapas
racionais que definem biracionalidades entre
variedades — este será o ponto de partida do
minicurso. Apresentarei aplicações de métodos
de Álgebra Comutativa, como o uso de syzygies
no estudo dos ideais de base de mapas
racionais, e discutirei critérios algébricos
que fornecem ferramentas computacionais para
verificar a biracionalidade de tais mapas.
Introduziremos também alguns métodos no
Macaulay2 para testar biracionalidade. Além
disso, veremos como diferentes teorias de
multiplicidades — como mixed multiplicity e
j-multiplicity — desempenham um papel
importante na abordagem algébrica dos mapas
racionais.


\section{Plan}
\label{birational-maps-commutative-algebra-section-plan}

\begin{itemize}
\item A. Simis, A. Doria, - 2012 (Adv. Math)
Criterio de Biracionalidade
\item Simis, - 2012: Plane Cremona.
\item Simis, -, 2017. Bounds for degrees of
Birational maps with CM graph.
\item Simis, Chardin 2021:
degree of rational map versus syzygy.
\item Mostafazadeh, 2024: loci de
$\text{Bir}(X)$.
\item \texttt{RationalMaps.m2}
\end{itemize}

\section{Introdução}
\label{birational-maps-commutative-algebra-section-intro}

\noindent
Let $R$ be a ring.
We are interested in $Mod_R$ and $R$-algebras
categories.
The first one is abelian, while the second
one isn't.
There are other diferences.

We are interested in discussing flat
$R$-algebras
and flat $R$-modules.

\begin{theorem}[Grothendieck Generic freeness
theorem]
\label{birational-maps-commutative-algebra-theorem-grothendieck}
Let $R$ be noetherian domain, $\varphi:R \to
S$,
$S$ a finitely generated $R$-algebra.
Then there exists $f \in R$
such that $\varphi:R_f \to S_{\varphi(f)}$,
$S_{\varphi(f)}$ is a free $R_f$-module.
$S_{\varphi(f)}$ is a flat $R_f$-module.
\end{theorem}

\begin{theorem}[Dimension of fibers theorem
(Matsumura 15.1)]
\label{birational-maps-commutative-algebra-theorem-dimension-of-fibers}
$\varphi:R \to S$, $R$ Noetherian,
$S$ Noetherian, $Q \in \Spec(S)$,
$p=\varphi^{-1}(Q) \in \Spec(R)$,
\begin{enumerate}
\item $ht Q \leq  ht p + \dim
\frac{S_Q}{pS_Q}$
[Shafarevich section 3].
\item If $\varphi$ is flat then
$\text{ht}Q=\text{ht}p+\dim
\frac{S_Q}{pS_Q}$.
\end{enumerate}

\noindent
If $Q \not\ni \varphi(f) \implies $
$S_Q$ is a free $R_p$-module (by
Grothendieck's theorem).
$\dim S_Q= \dim R_p+\dim \frac{S_Q}{p S_Q}$.
\end{theorem}

\noindent
Topologia de Zariski.
$$
\Spec(S)=\{Q:Q \text{ primo}\}
$$
$$
\text{fechado }=V(I)= \{ Q: Q \supseteq I\}.
$$

$S$ is a domain, $\Spec(S)$ is an irreducible
space,
$\implies $ open nonempty sets are dense.


\medskip\noindent
$k$ field. An algebraic set is  $V(I) = Z(I)$
for an ideal
$I \subset k[x_1,…,x_n]$. $I$ is prime if and
only if
$X:=V(I)$ is irreducible.
Recall that a polynomial map is given
by polynomial functions in the entries.
Any such map yields in the category of
algebras a map
$$
\begin{aligned}
k[Y]=A[Y]=\mathcal{O}_Y=\frac{k[y_1,…,y_m]}
{J}
&\xrightarrow{\varphi^*}
\frac{k[x_1,…,x_n]}{I}\\
y_i&\mapsto \overline{f_i(x)}
\end{aligned}
$$

\noindent
where $Y:=V(J)$.
$\varphi^*(g)=0$ $\forall  g \in J$,
$g \in J \iff g(f_1,…,f_m)\in I$
where $f_i$ are the coordinate polynomial
functions of $\varphi$.

$\varphi$ is an isomorphism if it has a right
and left inverse morphism.
$\varphi:X \to X$ is an automorphism
if it is an isomorphism.

\medskip\noindent
Question: determine the automorphism group
$\text{Aut}(X)$.

If $X=\mathbb{A}^1$, a morphism
$$
\begin{aligned}
\varphi: k[x] &\longrightarrow k[x] \\
x &\longmapsto f(x)
\end{aligned}
$$

\noindent
has an inverse if and only if there exists
$f(x)$ such that $f(g(x))=x$.
[This forces $f$ to have degree 1 since
composition $f \circ g$ must have degree
$\text{deg} f \cdot \text{deg} g = \text{deg}
x=1$.]
So $f(x)=ax+b$, $a\neq 0$.
Thus, we can parametrize the automorphism by
two numbers
$a,b$ with $a \neq 0$. This is a quasi-affine
variety.

For $X=\mathbb{A}^2$
the situation is a little more involved.
de Jonequeres, Triangular group.
Hierzch Jung 1942, chre.
Wout Van der Kulk, 1953, chr arbitrary.
$$
\text{Aut}(\mathbb{A}^2)
=\left\langle{}\text{linear, triangular}\right\rangle{}
$$

For $\text{Aut}(\mathbb{A}^3)$,
the question if whether
$\text{Aut}(\mathbb{A}^3)=
\left\langle{}\text{linear, triangular}\right\rangle{}$.
This is false in general

\begin{definition}
\label{birational-maps-commutative-algebra-definition-tame}
An auotomorphism $\sigma \in
\text{Aut}(\mathbb{A}^3)$
is called {\it tame} if $\sigma \in
\left\langle{}\text{linear, triangular}\right\rangle{}$
and {\it wild} otherwise.
\end{definition}

\noindent
Nagata's conjecture is that
a $\sigma:k[x,y,z] \to k[x,y,z]$,
$(x,y,z) \mapsto
(x+(x^2-yz)z,y+2(x^2yz)x+(x^2-yz)^2z,z)$
is wild.

Shiguero Kuroda 2014. Introduce a monomial
order
$\leq $ in $k[x,y,z]$.
[The following is not a criterion;
this just says "if this happens, then the
morphism is wild",
but not an iff.]
Let
$$
\begin{aligned}
\sigma: k[x,y,z] &\longrightarrow k[x,y,z] \\
\leavevmode &\longmapsto (f_1,f_2,f_3)
\end{aligned}
$$

\noindent
$\sigma$ is wild if
\begin{enumerate}
\item $\text{int}(f_1),\text{int}(f_2),
\text{int}(f_3)$
are linearly dependent over $\mathbb{Z}$,
pairwise linear over $\mathbb{Z}$.
\item $\text{int}(f_{i_1}) \neq  p
\text{int}(f_{i_2})+q(\text{int}(f_{i_3})$
for $p,q \in \mathbb{Z}_{\geq 0}$.
\end{enumerate}

\medskip\noindent
Jacobian conjecture:
a polynomial map $\varphi:\mathbb{C}^n \to
\mathbb{C}^n$
[with coordinate functions]
$(f_1(\mathbf{x}),…,f_n(\mathbf{x}))$
is an automorphism if and only if
$\det
\text{Jacob}(\varphi)=\text{constant}\neq
0$.

The forward implication is clear:
if  $\varphi \psi=\text{id}$
then
$\text{Jac}(\varphi)(\psi)\text{Jac}(\psi)=
\text{Jac}(\text{id})=1$.

[Bass-Conell-Wright, 1982]
Let $\varphi:\mathbb{C}^n \to \mathbb{C}^n$
a polynomial map with
$\det \text{Jac}\varphi=$ constant. Then
$\varphi$ is invertible (its inverse is a
polynomial map)
[note that by the holomorphic inverse
function theorem
we already know that there exists a local
inverse
that is a series, but here we see it's
actually polynomial]
if and only if $\varphi$ is injective
if and only if
$\varphi$ is proper
($\varphi^{-1}(\text{compact})=\text{compact}
$).

\begin{theorem}
\label{birational-maps-commutative-algebra-theorem-jacobian-conjecture}
If the Jacobian conjecture is valid for all
$n$ and
for any polynomial of degree $\leq 3$,
then it's valid.
\end{theorem}

\begin{theorem}
\label{birational-maps-commutative-algebra-theorem-jacobian-conjecture2}
If the Jacobian conjecture is valid over
$\mathbb{C}$
then it is valid over any domain.
\end{theorem}


\section{Rational maps of projective
varieties}
\label{birational-maps-commutative-algebra-section-ratio-maps-proje-varie}

\noindent
$X=V(I) \subseteq \mathbb{P}^n$,
$I \subseteq k[x_0,…,x_n]$ homogeneous.

\begin{definition}
\label{birational-maps-commutative-algebra-definition-presentation-of-map}
A rational map is $\varphi:X \dashrightarrow
Y$
for $X\subseteq \mathbb{P}^n$, $Y\subseteq
\mathbb{P}^m$
given by $(f_0:…:f_m)$
is {\it presented} by
$f_i \in R=k[\mathbf{x}]/I_X$,
$\text{deg}f_i=\text{deg}f_j$,
$f_i$ homogeneous.
\end{definition}

\noindent
A map can have another presentation:
$$
\begin{aligned}
\varphi:X &\dashrightarrow Y \\
\leavevmode &(g_0:…:g_m)
\end{aligned}
$$

\noindent
For all $p \in \text{Dom}\varphi$,
 $$
(f_0(p):…:f_m(p))=(g_0(p):…:g_m(p)).
$$

If
$$
I_2
\begin{pmatrix}f_0(\mathbf{x})&f_m(\mathbf{x}
)\\
g_0(\mathbf{x})&g_m(\mathbf{x})\end{pmatrix}
\in I_X
$$
[then the presentations are equivalent].

\begin{exercise}
\label{birational-maps-commutative-algebra-exercise-standard-quadratic}
$X=V(y^3-x^2z) \subseteq \mathbb{P}^2$,
$\sigma:X \dashrightarrow X$, $(yz:xz:xy)$.
[Show that] $(xz:y^2:x^2)$
[is an equivalent presentation, that is,]
$$
I_2
\begin{pmatrix}
yz & xz & xy\\
xz & y^2 & x^2 \end{pmatrix}
\in (y^3-x^2z).
$$

\end{exercise}

\noindent
What are the possible presentations of a
given rational map?
$$
f_0g_j=f_jg_0\qquad \text{ in }I_X.
$$
If $f_0 \neq 0 \implies
g_j=\left(\frac{g_0}{f_0}\right) f_j$.

\medskip\noindent
[We see that there is a 1-1 correspondence
between]
$$
\text{Presentations of } \varphi
\leftrightarrow
\Hom_R(I,R).
$$
From
$$
\xymatrix{
0\ar[r]
&I\ar[r]
&R\ar[r]
&R/I\ar[r]
&0
}
$$
we obtain
$$
\xymatrix{
0\ar[r]
&\underbrace{\Hom(R/I,R)}_{=0}\ar[r]
&\Hom(R,R)\ar[r]
&\Hom(I,R)\ar[r]
&\text{Ext}^1(R/I,R)\ar[r]
&\underbrace{\text{Ext}^1(R,R)
}_{=0}}
$$
And $\Hom(R/I,R)=\ast \iff \text{grade}(I)
\geq 1$.

If $\text{Ext}^1(R/I,R)=0 \implies R
\simeq\Hom(I,R)$.
Note that $\text{Ext}^1(R/I,R) \iff
\text{grade}(I,R)\geq 2$.

[This would help us understand
$\text{Bir}(X)_d$.
$$
\text{Domin}\varphi=\bigcup
\text{Domin}(\text{presentation})=X.
$$


\section{Image and graph of a rational map}
\label{birational-maps-commutative-algebra-section-image-graph}

[Given a rational map]
$\varphi:X \dashrightarrow Y$, $(f_0:…:f_m)$,
the ideal of definition of the image of
$\varphi$
is
$$
\mathcal{K}
=\{h(y) \in k[y_0,…,y_m]
:h(f_0(p),…,f_m(p)), \forall p \in
\text{Dom}\varphi\}.
$$
[The coordinate ring of the image has a
special name:]
$$
\frac{k[y_0:…:y_m]}{\mathcal{K}}
\simeq k[f_0,…,f_m]
\subseteq\frac{k[x_0,…,x_n]}{I_X}
$$
is called {\it special fiber $\mathcal{I}_I$
of the ideal $I\subseteq R$}.

The {\it analytic spread} is
$$
\dim(\text{Im}\varphi)
=\dim \mathcal{I}_I-1
=\ell_R(I)-1.
$$

Para um anel em geral
$f_0,…,f_m \in R$,
$\mathcal{I}_I=k[f_0t,…,f_mt] \subseteq
R[t]$.
$$
\text{ht}(I) \leq  \ell(I) \leq  \dim R.
$$

One can also prove via Dedekind-Mertens
lemma,
which says that
$c(f)^nc(fg)=c(g)c(f)^{n+1}$,
where $c$ is the content of a polynomial
as in Gauss lemma,
that
$$
\ell(I_j) \leq  \ell(I) + \ell(J)-1.
$$

\medskip\noindent
The {\it graph} of $\varphi$ is
$$
\Gamma_\varphi
= \overline{\{(p_0:…:p_n)\times
(f_0(p):…:f_m(p))
\subseteq X \times \mathbb{P}^m\}:
p \in \text{Dom}(\varphi)\}}
$$
$$
\{
h(\mathbf{x},\mathbf{y})
\in k[x_0,…,x_n,y_0,…,y_m]
:h(\mathbf{p},\mathbf{f}(\mathbf{x}))=0
\forall p \in \text{Dom}(\varphi)\}.
$$

The  {\it ideal of definition of the graph}
is given by homogeneous polynomial equations
of a presentation of $\varphi$.

The {\it Rees algebra} is
$$
\frac{R[y_0,…,y_m]}{J}
\simeq R[f_0t,…,f_\mapsto ] \subseteq R[t]
$$
where $J = $ polynomial equations of
$f_0,…,f_m$.


\section{Finite maps}
\label{birational-maps-commutative-algebra-section-finite-maps}

\noindent
[Another concept before we get to birational
maps.]

\begin{definition}
\label{birational-maps-commutative-algebra-definition-finite-map}
A rational map $\varphi:X \dashrightarrow
\mathbb{P}^m$
presented by $(f_0:…:f_m)$
is called a {\it finite map} if and only if
$\ell(I) = \dim R$.
$I$ maximal analytic spread.
\end{definition}

\noindent
[Keep in mind that dimension of image is
analytic spread minus 1,
and dimension of domain is $\dim R-1$.]

\medskip\noindent
[The following is another version
of Hilbert's Nullstellensatz in commutative
algebra:]

\begin{theorem}[Matsumura 5.6]
\label{birational-maps-commutative-algebra-theorem-matsumura-5-6}
$R$ Noetherian domain containing a field $k$,
$$
\dim (R) = \text{trdeg}_k Q(R).
$$
\end{theorem}

$$
[\text{Equations involving fraction fields of
coordinate rings}]
$$
We obtain that
$\varphi:X \dashrightarrow Y$
is finite if and only if
$\dim(\varphi^{-1}(q))=0$ for generic $q \in
Y$.
$\text{deg}(\varphi)=\# \{p :
\varphi(p)=q\}$.

\medskip\noindent
How can we compute (algebraically)
$\text{deg}(\varphi)$?

$$
\begin{aligned}
\text{deg}Z&=\sum_{\substack{z\text{ closed}
\\ \text{
in $Z$}}}
\lambda(\mathcal{O}_{Z,z}\text{deg}(Z)\\
&=
\sum_{ p \not \supseteq I}
\lambda \left(\left(
\frac{R}{I_q}\right)_p\right)
e\left(\frac{R}{p}\right)
\end{aligned}
$$

\noindent
since
$$
\{p:(f_0(p):…:f_m(p))
=(q_0:…:q)m): p \in \text{Domin}(\varphi)\},
$$
$$
I_q=I_2
\begin{pmatrix}f_0(\mathbf{x})&f_m(\mathbf{x}
)\\
q_0&q_m\end{pmatrix},\qquad  V(I_q).
$$
$$
I_q=Q_1\cap…\cap Q_t\cap\underbrace{…\cap
Q_s}_{
\sqrt{Q_i}\supseteq I}.
$$

\begin{exercise}
\label{birational-maps-commutative-algebra-exercise-saturation}
$(I_q:I^\infty)=Q_1\cap…\cap Q_t$.
\end{exercise}

$$
\text{deg}(\varphi)
=e\left(\frac{R}{I_q:I^\infty}\right).
$$


\medskip\noindent
Integral closure (1,82).

$$
\xymatrix{
&  V \ar[dr]^{\pi_L}\\
X\ar@{-->}[ur]^{(g_0,…,g_N)}\ar@{-->}[rr]
^{\varphi=(f_0,…,f_m)}
& & Y
}
$$
where $\pi_L$ is a linear map.
$$
\underbrace{(f_0,…,f_m)}_{J}
\subseteq \underbrace{(g_0,…,g_N)}_{I}.
$$


The center of the projection $\pi_L$
does not intersect $V$ if and only if $J$
is reduction of $I$, that is,
$I \subseteq \overline{J}$.

\medskip\noindent
Another way of stating this result is:
given $J \subseteq I$,
$$
J \subseteq \mathcal{F}_I
=k[g]
=\frac{k[y_0,…,y_m]}{k},
\ell_1,…,\ell_m.
$$
$J \subseteq I$ reduction if and only if
$(\ell_0,…,\ell_m)$ is $m$-primary.



\section{Summary}
\label{birational-maps-commutative-algebra-section-summary}

\noindent
Let $x \subset \mathbb{P}^n$ and $y \subset
\mathbb{P}^m$.
$\varphi:x \dashrightarrow y$, $(f_0:…:f_m)$,
$r=k[x]=k[x_0,…,x_n]/I_x$,
$I=(f_0,…,f_n) \subseteq R$.

We defined $\varphi$ to be {\it finite} if
$\ell(I) = \dim r$.
$[q(x),q(y)]<\infty$.
$\#\{p : p \in x: \varphi(p) = q,q\text{
general in }y\}$.

$$
i_q=i_2
\begin{pmatrix} f_0(\mathbf{x}&  \cdots &
f_n(\mathbf{x})\\
q_0& \cdots & q_m \end{pmatrix}
\subseteq r
$$
$$
e\left(\frac{R}{I_q:I^\infty}\right)
=\text{deg}(\varphi)
$$
[is the topological degree of $\varphi$.]

Eisenbud-Ulrich 2007: raw ideal.
$I_q:I^\infty=$ morphism fiber ideal.
Correspondence fiber ideal
$$
\xymatrix{
&\underbrace{\Gamma}_{\substack{\text{
graph} \\ \text{of }X}} \subseteq X
\times\mathbb{P}^m
\ar[dr]_{\pi_2}\ar[dl]^{\substack{\pi_1 \\
\text{(birational)}}}\\
X \ar@{-->}[rr]& & Y \ni q
}
$$

$$
\begin{aligned}
\pi_2^* :\frac{k[Y]}{I_Y}&\to
R_I=\frac{y_0,…,y_m]}{J}
y_i&\longmapsto \mathbf{y_i}=f_it \in
R[\mathbf{f}t].
\end{aligned}
$$

$$
C_q=(I_qtR_I :_{R_I}I^{\infty}t) \cap R,
\qquad  R \subseteq R_I.
$$
\begin{exercise}
\label{birational-maps-commutative-algebra-exercise-raw-ideal}
$C_q=(I_qtR_I:R_i^{\infty}t)
=\bigcup_{n=1}^\infty (I_q I^{n-1}:_R I^n)
\subseteq(I_q:I^\infty)$.
\end{exercise}

$$
v(q)=\{p:(I_q)_p:\text{ is not a reduction of
$I_p$}\}.
$$
$$
I_q(q:\psi)=(\underbrace{I_q:I}_{\text{raw
ideal}}
\subseteq C_q \subseteq (I_q:I^\infty)
$$
Question: when to $(I_q:I)$ and $C_q$ have
the
same saturation?

If $I$ has a [crossed: linear generalized]
raw
 and $I$ is locally of linear type
$\implies
(I_q:I)^{\text{sat}}=C_q^{\text{sat}}$.




\medskip\noindent
[Simis Chardin, 2021.]
$Q(X)=R_i{\underline{0}}=\left\{\frac{a}{b}:
b\neq
0\right\}$.
$R_{(\underline{0})}=\left\{\frac{a}{b}:b\neq
0,b\text{
homogeneous}\right\}$.

 $$
k[x]_x=k[x,x^{-1}], R_{((0))}=k,
$$
$$
k(X)=R_{((0))}.
$$

\begin{lemma}
\label{birational-maps-commutative-algebra-lemma-degree}
Let $A = \bigoplus_{i=0}^\infty A_i$
be a graded domain over a field $A_0=F$
with $\dim A=1$.
Then $e(A)=[K(A):F]$.
\end{lemma}

\begin{proof}
There exists $d$ such that $d_F A_i=e$
for all $i \geq d$,
$A_d=F\left\langle{}a_1,…,a_e\right\rangle{}$.
Show that $K(A)= F
\left\langle{}\frac{a_1}{a_1},…,\frac{a_e}{a_1}
\right\rangle{}$
(exercise).
\end{proof}

\noindent
[Then]
$$
\begin{aligned}
e(S/J \otimes_B Q(Y))
&=[K(S/J \otimes Q(Y)):Q(Y)]\\
&=[K(S/J)(K_0):K(Y)(K_0)]\\
&=[K(\Gamma):K(Y)]\\
&=\text{deg} \varphi.
\end{aligned}
$$


[The ring $S/J \otimes_B Q(Y)$ is always
Cohen-Macaulay of rank 1.]

\begin{theorem}
\label{birational-maps-commutative-algebra-theorem-minimal}
Suppose that $J$ is prime minimal over an
ideal $I$
generated by elements of degrees $d_1 \geq
d_2 \geq …$.
$$
\text{deg}\varphi\leq  d_1…d_{t}e(R)
$$
where $t=\dim X$.

\end{theorem}

\noindent
We have the following corollary:
\begin{lemma}
\label{birational-maps-commutative-algebra-lemma-minimal}
$\text{deg}(\varphi)\leq  d_1…
d_{d_\text{max}}e(R)$.
\end{lemma}

\section{Eisenbud-Goto conjecture}
\label{birational-maps-commutative-algebra-section-eisenbud-goto-conjecture}

\noindent
1982.
Let $X \subseteq \mathbb{P}^n$ irreducible
and ``non-degenerated''.

$$
\begin{aligned}
reg(X) &\leq  \text{deg}(X) -
\text{codim}(X)\\
reg(X)&=
\text{min}\{r:H^1(X,\mathcal{O}_X(r-i))=
0\forall
i >0\}\\
&=\text{min} \{r:
H^2_{\underline{m}}(R)_{r-i}=0 \forall
i>0\}.
\end{aligned}
$$

[Regularity controls properties of the ideal
such as the degree of the generators in a
Gröbner base.]

$$
reg(R/I_X)
=\text{max}
\{t_i+i:t_i=\text{max}\{d_{ji}\}\}
$$
$$
reg(R/f)=\text{deg}f-(n+1)+n=\text{deg}f-1.
$$

[Notice the term ``non-degenerate''. It means
that
the ideal $I_X$ in
$R=\frac{k[x_0,…,x_n]}{I_X}$
does not contain any linear terms.]

$$
\begin{aligned}
e\left(\frac{S}{J}\otimes Q(Y)\right) &\geq
n+1-
\dim_{Q(Y)}(\tilde{J}_{\underbrace{1}
_{\text{deg}
X}}
+reg(\tilde{J})-2.\\
\text{deg}\varphi\geq
n+q-\dim_{Q(Y)}(\tilde{J}_1)
=\text{ Jacobian dual rank}.\\
reg(\tilde{J}) \geq
\text{deg}(\text{gens})\geq 2\\
\text{deg}\varphi\geq n+1-r
\text{jdrank}(\varphi).
\end{aligned}
$$

\noindent
[The conjecture is false in general,
but in some cases holds, such as in the
Cohen-Macaulay case.]


\section{Summary}
\label{birational-maps-commutative-algebra-section-summary-duplicate-2}

Recall our notations:
\begin{itemize}
\item $\varphi:X \dashrightarrow Y$,
$(f_1,…,f_n)$.

\item $I=(f_0,…,f_n)$.

\item $R=k[x_0,…,x_n]/I_X$.

\item $\mathcal{R}_I(R)=R[y_0,…,y_n]/J=S/J$.

\item $B=k[y_0,…,y_m]$.

\item $k[Y]=k[y_0,…,y_m]/J\cap B=B/q$/

\item $S/J \otimes_B Q(Y)$, $Q(Y)=Q(B/q)$. CM
dim 1, domain.

 \item $S/J \otimes )B Q(Y)
=\frac{Q(Y)[x_0,…,x_n]}{\underbrace{\tilde{J}
}_{=\text{Im}\theta}}$,
$$
\xymatrix{
0\ar[r]
&J\ar[r]
&S\ar[r]
&S/J\ar[r]
&0
}
$$
$$
\xymatrix{
J \otimes Q(Y)\ar[r]^{\theta}
&S \otimes Q(Y)\ar[r]
&S/J \otimes_B Q(Y)\ar[r]
&0
}
$$
$S=k[\mathbf{x},\mathbf{y}]/I_X$.

\item
$e\left(\frac{Q(Y)[\mathbf{x}]}{\tilde{J}}
\right)=\text{deg}\varphi$.

\item
$\text{djrank}=\dim_{Q(Y)}(\tilde{J})_1$

\item $\text{deg}(\varphi)\geq
(n-\text{jdrank}(\varphi))
+\text{reg}(\tilde{J})-1)$.
\end{itemize}



\begin{theorem}
\label{birational-maps-commutative-algebra-theorem-jacob-rank-birat}
$n=\text{jdrank}(\varphi)$ if and only if
$\varphi$ is birational.
\end{theorem}

\begin{proof}
$\text{deg}(\varphi)=e(Q(Y)[x_0,…,x_n]/
\tilde{J})$.

($\implies$) If $\text{jdrank}(\varphi)=n
\implies
\tilde{J}_1=(\ell_1,…,\ell_n)
\implies e(A)=e(Q(Y)[x]/\tilde{J})$
where $A:=S/J\otimes_BQ(Y)$.
$A$ domain,  $\dim=1 \implies  \tilde{J}'=0
\implies e(A)=1=\text{deg}(\varphi)$.

($\impliedby$) $e(A)=1 \implies $ACM
$\implies $ quasi-unmixed
$\implies$ $\hat{A}$ is equidimensional
and $A$ is analytically unmixed
$\implies $ $A$ is a regular ring.
$A=\frac{Q(Y)[\mathbf{x}]}{\tilde{J}}$ (See
[Bruns-Herzog §2]).
[Thus $\tilde{J}$ is linear,]
$\tilde{J}=(\ell_1,…,\ell_n)$
since $\text{codim}\tilde{J}=n$.
\end{proof}

\begin{proposition}
\label{birational-maps-commutative-algebra-proposition-regular-local-ring}
Let $(A,\mathfrak{m})$ a regular local ring
and $I\subseteq \mathfrak{m}$.
If  $A/I$ is regular then $I$ is generated by
a part
of a system of parameters.
\end{proposition}

\section{How to compute the Jacobian dual
rank}
\label{birational-maps-commutative-algebra-section-computing-jacobian-dual-rank}

$$
R=\frac{k[x]}{I_X}\supseteq I=(f_0,…,f_m),
$$
$$
\mathcal{I}_I(R)=\frac{R[\mathbf{y}]}{J}.
$$
$$
\left\langle{}J_{(1,\ast)}\right\rangle{}
=\left\langle{}\underbrace{g_1}_{\in R[\mathbf{y}]},
…,\underbrace{g_t}_{R[\mathbf{y}]}\right\rangle{}
$$
lifts $g_1(\overline{Q}_i)$,
$Q_i \in k[\mathbf{x},\mathbf{y}]_{i,\ast)}$.
$$
k[\mathbf{y}]\ni
\begin{pmatrix}
\frac{\partial Q_1}{\partial x_0}& \cdots
&\frac{\partial Q_t}{\partial x_0}\\
\vdots &  &  \vdots \\
\frac{\partial Q_1}{\partial x_n}&  \cdots
&  \frac{\partial Q_t}{\partial x_n}
 \end{pmatrix}
=M
$$
Jacobian dual matrix.
$$
-:k[\mathbf{y}]\to \frac{k[t]}{I_Y}/
$$
$$
\psi=
\begin{pmatrix}
\overline{\frac{\partial Q_1}{\partial x_0}}
&  \cdots &
\overline{\frac{\partial Q_1}{\partial
x_0}}\\
\vdots &  &  \vdots \\
\overline{\frac{\partial Q_1}{\partial x_n}}
&  \cdots &
\overline{\frac{\partial Q_t}{\partial x_n}}
 \end{pmatrix}
\in M_{(n+1) \times t}(k[Y]).
$$
$$
\tilde{\psi}
=\begin{pmatrix}
\frac{\partial Q_1}{\partial x_0}    \\
&  \ddots & \\
&  &  \frac{\partial Q_t}{\partial x_n}
\end{pmatrix}
\Big|_{(f_0(\mathbf{x}),…,f_m(\mathbf{x})}
\in M_{(n+1)t}(R).
$$

\begin{proposition}
\label{birational-maps-commutative-algebra-proposition-rank-and-birationality}
$\text{rk}\psi=\text{rk}\tilde{\psi}=
\text{jrank}(\varphi)
=\dim_{Q(Y)}\tilde{J}_1$,
$\text{rk}\psi=n$ if and only if
$\varphi$ is birational.
\end{proposition}

\begin{proof}
[Compute minors, add a row that is missing…
this way we obtain a candidate for the
inverse map…]
\end{proof}

\section{Extra}
\label{birational-maps-commutative-algebra-section-extra}

\noindent
Consider
$$
\text{Bir}(X)_d
= \{ \varphi:X \dashrightarrow
X,(f_0,…,f_n):\text{deg} f_i=d\}.
$$
$$
f_0=d_0x^d+a_1x_0^{d-1}x_1
+a_2x_0^{d-1}x_1x_2…
$$
Let $H_d \subseteq
\mathbb{A}^{(n+1)\times\binom{n+d}{d}}$
be the points whose corresponding map is
birational.
We have a map $H_d \to \text{Bir}(X)_d$.
[We have proved that] $H_d$ is a
quasi-projective variety
and $\text{Bir}(X)_d$ is a constructible set.

[If the Cremona map has degree $d$,
what's the degree of the inverse map?]
$\varphi^{-1}\leq d^{n-1}$.
A proof of this is due to Gabber.

\section{Extra 2}
\label{birational-maps-commutative-algebra-section-extra2}

$$
D=\underbrace{\sum a_i P_i}_{E}
- \sum\underbrace{b_j}_{\in \mathbb{Z}_{\geq
0}}
\underbrace{q_j}_{\text{proj}},
\qquad \text{proj}(R)=X.
$$
$$
\mathcal{O}_X(D)=\tilde{M}
$$
$$
\begin{aligned}
\mathcal{O}_X(D)&=\tilde{M}\\
\text{prod}E&=\prod p_i^{[a_i]}=(f^{a_i}:…:
f_k^{a_i})\\
\text{prod}F&=\prod q_j^{[b_j]}\\
\text{prod}E^{\ast \ast}&=
\Hom(\Hom(\text{prod}F,R',R))\\
\text{dual}&=\text{prod}E^{\ast \ast}\otimes
\Hom(\text{prod}(D^{\ast \ast},R)\\
M&=\Hom(\text{dual},R).
\end{aligned}
$$

\noindent
D. OO(D). divisorToModule(D)

\medskip\noindent
Divisor (Karl Schewede,…) [is the name of
this Macaulay2 package].
$M$ rk 1, torsion-free. embedAsIdeal(M).

$$
\xymatrix{
\bigoplus R(-e_j)
\ar[r]^{
\begin{pmatrix} a_1\\
\vdots \\
a_n \end{pmatrix} }_\psi
&  R^n(-d_i)
\ar[r]
& M\ar[d]^{m_i \mapsto  f_i}\ar[r]& 0\\
& & R
}
$$
$(\underline{f}) \in \Ker \psi^t$.

$I=$.
$$
\psi^t:R^n(td_i) \to R(fe_j)
$$
$I=(\underline{f}) \in R^n(td_i)_a$.

mapToProjectiveSpace(D),
Map defined by base of $I_a$.









% --- End birational-maps-commutative-algebra.tex ---

% --- Begin hodge-birational-atoms-2026.tex ---
\chapter{Hodge theory, birational geometry and atoms 2026}
\label{chapter-hodge-birational-atoms-2026}


and atoms 2026}

\medskip\noindent
{\bf Conference information.}

\href{%
https://sites.google.com/unicamp.br/%
campinas-miami-math-2026/%
}{Campinas--Miami Mathematical Meeting 2026}

May 11--13, IMECC/Unicamp,
Campinas-SP.

\href{%
https://impa.br/evento/
conference-on-hodge-theory-
birational-geometry-and-atoms/
}{Conference on Hodge theory,
birational geometry and atoms}

May 14--15, IMPA,
Rio de Janeiro.

\medskip\noindent
{\bf Organizing Committee.}

Leonardo Cavenaghi, Lino Grama,
Ludmil Katzarkov, Jaqueline Mesquita,
Ricardo Miranda, Misha Verbitsky.

\section{Collapsing Flat
SU(2)-Bundles to
Spherical 3-Manifolds}
\label{hodge-birational-atoms-2026-section-spherical-3-manifolds}

\medskip\noindent
{\bf Speaker.}

Eder M. Correa
(Campinas, Brazil)

\medskip\noindent
{\bf Title.}

Collapsing Flat
SU(2)-Bundles to
Spherical 3-Manifolds.

\medskip\noindent
{\bf Abstract.}

In this talk, we present a geometric
mechanism for the emergence of
spherical $3$-manifolds from the
superspace of Riemannian metrics
associated with flat
SU(2)-bundles over closed orientable
hyperbolic surfaces. We show that any
spherical $3$-manifold can be realized
as a boundary point in the
Gromov-Hausdorff closure of this
superspace through a collapsing
process. Furthermore, we discuss how
the convergence toward the spherical
limit is controlled by the order of
the fundamental group of the target
manifold and the systolic geometry of
the hyperbolic base. Finally, we
illustrate these results by showing
how the Bolza surface provides the
optimal error bound for the
convergence toward the Poincar\'e
homology sphere.

\medskip\noindent
{\bf Main theme.}

Spherical $3$-manifolds appear as
Gromov--Hausdorff boundary points of
superspaces of suitable
$\operatorname{SU}(2)$-bundles over
closed hyperbolic surfaces.

\medskip\noindent
{\bf Spherical manifolds.}

\begin{definition}
\label{hodge-birational-atoms-2026-definition-riemannian}
A Riemannian manifold $M^n$ is
{\it spherical} if it is complete and
locally isometric to the round sphere
$S^n$.
Equivalently, it has constant
sectional curvature $+1$ and
universal cover $S^n$.
\end{definition}

If $M$ is connected, then
$$
        M \cong S^n/\Gamma
$$
where $\Gamma<\operatorname{Isom}(S^n)$
is finite and acts freely.

In dimension $3$ we use
$$
        S^3\cong \operatorname{SU}(2),
$$
via the unit quaternions.
Thus one gets spherical quotients
$$
        \operatorname{SU}(2)/\Gamma
$$
when $\Gamma<\operatorname{SU}(2)$ is
finite and acts by left multiplication.

The finite subgroups of
$\operatorname{SU}(2)$ are the binary
cyclic, binary dihedral, binary
tetrahedral, binary octahedral and
binary icosahedral groups.
This is the ADE list through the
McKay correspondence.

See the constant curvature discussion in
Section
\ref{%
differential-geometry-%
section-curvature%
},
especially Equation
\ref{%
differential-geometry-%
equation-constant-curvature%
}.
For the ADE side, see Section
\ref{representation-theory-section-sl2c}.

\medskip\noindent
{\bf Warning.}

The general spherical space-form
statement is $S^3/\Gamma$, with
$\Gamma<\operatorname{SO}(4)$ finite and
free.
The ADE statement is the clean
single-action case
$\Gamma<\operatorname{SU}(2)$.

\medskip\noindent
{\bf Main theorem.}

\begin{theorem}
\label{hodge-birational-atoms-2026-theorem-spherical-gh}
Let
$$
        S=\operatorname{SU}(2)/\Gamma
$$
be a spherical $3$-manifold, with
$\Gamma<\operatorname{SU}(2)$ finite.
Then there exists a closed orientable
hyperbolic surface $\Sigma$ and a flat
principal $\operatorname{SU}(2)$-bundle
$$
        \operatorname{SU}(2)\to P\to\Sigma
$$
such that $S$ occurs as a boundary point
of the Gromov--Hausdorff closure of
$$
        {\mathcal S}(P)
        =
        \operatorname{Met}(P)/
        \operatorname{Diff}(P).
$$
Equivalently, there are smooth metrics
$g_i$ on $P$ such that
$$
        (P,d_{g_i})\to (S,d_{g_0})
$$
in Gromov--Hausdorff distance.
\end{theorem}

Here $g_0$ is the spherical quotient
metric.
I am recording the theorem in the form
understood from the talk.
The exact statement should be checked
against the speaker's paper.

The point is that a fixed smooth
manifold $P$ can carry metrics which
collapse, in the Gromov--Hausdorff
sense, to a lower dimensional spherical
space form.

\medskip\noindent
{\bf Holonomy and monodromy.}

Let $\pi:P\to M$ be a principal
$G$-bundle with connection $\theta$.
Fix $m\in M$ and $u\in P_m$.
If $\gamma$ is a loop based at $m$,
let $\widetilde\gamma$ be its
horizontal lift starting at $u$.
Then $\widetilde\gamma(1)$ lies in
$P_m$, so there is a unique
$g_\gamma\in G$ with
$$
        \widetilde\gamma(1)
        =
        u\cdot g_\gamma .
$$

The holonomy group is
$$
        \operatorname{Hol}_u(\theta)
        =
        \{g_\gamma\}.
$$
The restricted holonomy group
$\operatorname{Hol}_u^0(\theta)$ is
obtained from null-homotopic loops.
For a connected base it is the
identity component of the holonomy
group.

The monodromy map is
$$
        m_\theta:\pi_1(M,m)\to
        \operatorname{Hol}_u(\theta)/
        \operatorname{Hol}_u^0(\theta),
        \qquad
        [\gamma]\mapsto [g_\gamma].
$$
So the quotient by restricted holonomy
records the discrete part of holonomy
seen by the fundamental group.

\medskip\noindent
{\bf Ambrose--Singer theorem.}

Let $\Omega_\theta$ be the curvature
form of $\theta$.
The Ambrose--Singer theorem says that
$$
        \mathfrak{hol}_u(\theta)
        =
        \operatorname{Lie}
        \operatorname{Hol}_u^0(\theta)
$$
is generated by curvature.
More exactly, take a point $v\in P$
reachable from $u$ by a horizontal
curve, and take horizontal vectors
$X,Y\in T_vP$.
Then the curvature value
$$
        \Omega_\theta(v)(X,Y)
        \in\mathfrak g
$$
transported back to $u$ contributes to
$\mathfrak{hol}_u(\theta)$.
The span of all such values is the
holonomy Lie algebra.

For a vector bundle connection $\nabla$,
this is the formula
$$
        \mathfrak{hol}_x(\nabla)
        =
        \operatorname{span}\{
        P_\gamma^{-1}R_y(X,Y)P_\gamma
        \}.
$$
Here $\gamma:x\to y$, $P_\gamma$ is
parallel transport, and
$X,Y\in T_yM$.

Thus curvature gives the infinitesimal
part of holonomy, while monodromy
records the residual discrete part.

\medskip\noindent
{\bf Riemann--Hilbert correspondence.}

For the smooth discussion here, the
relevant case is the flat case:
$$
        \Omega_\theta=0.
$$
Then parallel transport depends only on
the homotopy class of the path.
Choosing $m\in M$ and $u\in P_m$ gives
the holonomy representation
$$
        \rho_\theta:\pi_1(M,m)\to G,
        \qquad
        [\gamma]\mapsto g_\gamma .
$$

The Riemann--Hilbert correspondence
identifies flat principal $G$-bundles
with representations of the
fundamental group:
$$
        \left\{
        \hbox{flat principal }G
        \hbox{-bundles on }M
        \right\}/\cong
        \longleftrightarrow
        \operatorname{Hom}(\pi_1(M,m),G)/G .
$$
The group $G$ acts on representations
by conjugation.

Thus two flat bundles with connection
$(P,\theta)$ and $(Q,\omega)$ are
equivalent if there exists a principal
$G$-bundle isomorphism
$$
        \varphi:P\to Q
$$
such that
$$
        \theta=\varphi^*\omega .
$$
On the representation side this says
that the holonomy representations
$\rho$ and $\rho'$ are conjugate:
$$
        \rho'=g\rho g^{-1}.
$$
If a point in the fiber over $m$ is
fixed, the conjugacy ambiguity
disappears.

Conversely, a representation
$$
        \rho:\pi_1(M,m)\to G
$$
defines a flat principal bundle by
$$
        \widetilde M\times_\rho G .
$$
This is the flat case of the
Riemann--Hilbert dictionary.

\medskip\noindent
{\bf Collapse sequence.}

The paper uses a Kaluza--Klein type
family of metrics
$$
        g_{P,n}
        =
        {1\over n}\pi^*(h_\Sigma)
        +
        (\theta,\theta)_{
        \operatorname{SU}(2)}
$$
on the total space $P$.
As $n\to\infty$, the horizontal
directions shrink and one expects the
Gromov--Hausdorff limit to be the
spherical quotient determined by the
holonomy.

\medskip\noindent
{\bf Gromov--Hausdorff distance.}

Let $(X,d_X)$ and $(Y,d_Y)$ be compact
metric spaces.
A correspondence is a subset
$$
        R\subset X\times Y
$$
whose two coordinate projections are
surjective.

A correspondence is like a generalized
graph.
The graph of a map $f:X\to Y$ is a
correspondence only when $f$ is
surjective.
In general, $R$ behaves like a
multivalued graph.

The distortion of $R$ is
$$
        \operatorname{dis}(R)
        =
        \sup\{
        |d_X(x,x')-d_Y(y,y')|:
        (x,y),(x',y')\in R
        \}.
$$

The Gromov--Hausdorff distance is
$$
        d_{GH}(X,Y)
        =
        {1\over 2}
        \inf_R\operatorname{dis}(R),
$$
where $R$ runs through all
correspondences between $X$ and $Y$.

A small-distortion correspondence is an
almost distance-preserving matching.

\medskip\noindent
{\bf Isometry classes.}

The Gromov--Hausdorff distance is
invariant under isometry.
If
$$
        X\cong X',
        \qquad
        Y\cong Y',
$$
then
$$
        d_{GH}(X,Y)=d_{GH}(X',Y').
$$
Hence $d_{GH}$ descends to isometry
classes:
$$
        d_{GH}([X],[Y])=d_{GH}(X,Y).
$$

Moreover, for compact metric spaces,
$$
        d_{GH}([X],[Y])=0
        \quad\Longleftrightarrow\quad
        X\cong Y.
$$
So compact metric spaces modulo
isometry form a genuine metric space.

\medskip\noindent
{\bf Superspace.}

Let $P$ be a smooth manifold and
denote by $\operatorname{Met}(P)$ the
space of smooth Riemannian metrics on
$P$.
The diffeomorphism group acts by
pullback.
The superspace of $P$ is
$$
        {\mathcal S}(P)
        =
        \operatorname{Met}(P)/
        \operatorname{Diff}(P).
$$
Thus two metrics define the same point
when they differ by a diffeomorphism.

Assume now that $P$ is compact.
Each metric $g$ defines the compact
metric space $(P,d_g)$, so there is a
map
$$
        {\mathcal S}(P)\to{\mathcal{GH}},
        \qquad
        [g]\mapsto[(P,d_g)].
$$
It is well-defined because if
$g=\phi^*g'$, then
$$
        \phi:(P,d_g)\to(P,d_{g'})
$$
is an isometry.

Conversely, by Myers--Steenrod, an
isometry between smooth Riemannian
manifolds is smooth.
Thus superspace embeds in the
Gromov--Hausdorff moduli space of
compact metric spaces.

\medskip\noindent
{\bf Riemannian submersion lemma.}

A technical lemma used in the proof
gives a Riemannian submersion
$$
        f:(P,g)\to
        (\operatorname{SU}(2)/\Gamma,g_0).
$$
Here $g_0$ is the quotient of the
round metric on
$\operatorname{SU}(2)\cong S^3$.

Recall that
$f:(P,g)\to(N,h)$ is a Riemannian
submersion if $f$ is a submersion and
$$
        df_p:
        (\operatorname{Ker}df_p)^\perp
        \to T_{f(p)}N
$$
is an isometry for every $p\in P$.
Thus
$$
        T_pP
        =
        \operatorname{Ker}df_p
        \oplus
        (\operatorname{Ker}df_p)^\perp
$$
is the vertical-horizontal splitting.

For horizontal vectors $X,Y$ one has
$$
        g(X,Y)=g_0(df(X),df(Y)).
$$

This lemma is the bridge between the
principal-bundle construction and the
Gromov--Hausdorff collapse:
the ADE quotient is not only a
holonomy target, but also the metric
quotient sitting below $(P,g)$.

\medskip\noindent
{\bf Examples.}

The Poincare homology sphere is
$$
        \Sigma(2,3,5)\cong S^3/I^*
$$
where
$I^*<\operatorname{SU}(2)$ is the
binary icosahedral group of order
$120$.
It double covers the usual icosahedral
rotation group
$A_5<\operatorname{SO}(3)$.
Thus the quotient is by the binary
icosahedral group, not by the ordinary
icosahedral group.

It is a homology sphere:
$$
        H_*(\Sigma(2,3,5),\mathbb{Z})
        \cong H_*(S^3,\mathbb{Z}).
$$
but it is not simply connected.
Its fundamental group is $I^*$.

The Bolza surface is the smooth
completion of
$$
        y^2=x^5-x.
$$
It is a two-fold cover of
$\mathbb{P}^1$ branched over
$$
        0,\ \infty,\ 1,\ -1,\ i,\ -i.
$$
By Riemann--Hurwitz, a double cover of
$\mathbb{P}^1$ branched over six points has
genus
$$
        g={6-2\over 2}=2.
$$

The Bolza surface is the unique
{\it global} maximum of the systole
function on the moduli space
${\mathcal M}_2$, after giving every
Riemann surface its uniformizing
hyperbolic metric.
In other words, among closed
hyperbolic surfaces of genus $2$,
it maximizes the length of the
shortest closed geodesic.

\medskip\noindent
{\bf Points to check.}

\begin{enumerate}
\item Is the theorem stated for all
spherical 3-manifolds, or only for the
single-action quotients
$\operatorname{SU}(2)/\Gamma$?
\item In the construction, which
representation
$\rho:\pi_1(\Sigma)\to\Gamma$ is chosen?
\item What is the exact collapse
sequence $g_i$ on $P$?
\item Where is the Riemannian submersion
lemma used in the proof?
\end{enumerate}

\section{Updates on flows of
geometric structures}
\label{hodge-birational-atoms-2026-section-updates-on-flows-of-geometric-structures}

\medskip\noindent
{\bf Speaker.}

Henrique Sa Earp
(Campinas, Brazil)

\medskip\noindent
{\bf Title.}

Updates on flows of geometric
structures.

\medskip\noindent
{\bf Abstract.}

Aiming at a public with interests
among Riemannian and complex
geometry,
Lie groups and parabolic PDEs,
I will recall some results towards a
general analytical theory for flows
of $H$-structures,
obtained with Fadel,
Loubeau and Moreno.
Then, as a concrete example,
I will present some recent progress
on $SU(2)$-flows on 4-manifolds,
initiating a classification
programme of ``parabolic'' flows,
based on a representation-theoretic
method originally due to Bryant
in the context of $G_2$ geometry,
obtained with Fowdar.

\medskip\noindent
{\bf Baby summary.}

The talk asks which geometric
structures can be moved by a
well-posed PDE.
First one sets up the general
$H$-structure framework.
Then one studies $SU(2)$ examples and
uses representation theory to list
the allowed flow terms.

\medskip\noindent
{\bf H-structures.}

\begin{definition}
\label{hodge-birational-atoms-2026-definition-let-m}
Let $M^n$ be a smooth manifold and
let $H\subset \operatorname{GL}(n,
\mathbb{R})$.
An {\it $H$-structure} on $M$ is a
reduction of the frame bundle
$\operatorname{Fr}(M)$ from
$\operatorname{GL}(n,\mathbb{R})$
to $H$.
Equivalently, it is a section of
$\operatorname{Fr}(M)/H$.
If $\xi_0$ is a reference tensor,
then $H=\operatorname{Stab}(\xi_0)$.
\end{definition}

$$
\xymatrix{
P:=\operatorname{Fr}(M)
\ar[dd]
\ar[dr]^H
&
\\
& N:=P/H
\ar[dl]_\pi
\\
M\ar@/^-1pc/@{-->}[ur]_{\sigma}
&
}
$$
and there is an associated
tensor bundle picture
$$
        \mathcal{F}\subset
        \bigoplus \mathcal{T}^{p,q}
        \to M,
        \qquad
        \pi^*\mathcal{F}\to N.
$$
Here $TP$ splits into vertical and
horizontal parts.
Write
$$
\xymatrix{
TP \ar[dd] \ar[dr]
&
\\
&
TN:=\nu\oplus\mathcal H\ar[dl]_\pi\\
M\ar@/^-1pc/@{-->}[ur]_{\sigma}
}
$$
to record the tangent-bundle picture.

In prose: an $H$-structure is a way to
equip $M$ with extra geometric
data that is modeled on a fixed tensor
or algebraic structure with symmetry
group $H$.
The section
$\sigma:M\to P/H$ picks, at each point,
an adapted frame modulo $H$.
Equivalently, one can encode the same
information by a tensor field
$\xi$ on $M$ whose stabilizer is
exactly $H$; the bundle $P/H$ is the
space of all such local choices.

\begin{example}
\label{hodge-birational-atoms-2026-example-harmonic-flow2}
For an $H$-structure $\sigma$ with
torsion tensor $T$, one may consider
the energy
$$
        E(\sigma)=\int_M |T|^2 .
$$
The harmonic flow is the negative
gradient flow of this energy.
Its purpose is to deform $\sigma$
toward a more special or torsion-free
structure.

Then it is natural to consider the
Ricci flow of $H$-structures.
\end{example}

\begin{definition}
\label{hodge-birational-atoms-2026-definition-harmonic-h}
Let $H\subset \operatorname{SO}(n)$.
A {\it harmonic $H$-structure} is a
critical point of
$$
        E(\xi)
        =
        {1\over 2}
        \int_M |T|^2 .
$$
Here $\xi$ is the corresponding
$H$-structure and $T$ is its torsion
tensor.
\end{definition}

\begin{theorem}[$\varepsilon$-regularity]
\label{hodge-birational-atoms-2026-theorem-epsilon-regularity2}
Let $(M,g)$ be closed and let
$\xi(t)$ be a harmonic $H$-flow on
$M\times[0,\tau)$.
Fix $0<r_M<\operatorname{inj}(M,g)$ and
$E_0\in(0,\infty)$.
If the local energy of $\xi(t)$ on a
geodesic ball $B_{r_M}(y)$ is small
enough, then torsion stays controlled
on a smaller ball for a short time.
\end{theorem}

\begin{theorem}[Energy gap]
\label{hodge-birational-atoms-2026-theorem-energy-gap2}
Let $(M,g)$ be a closed Riemannian
$n$-manifold admitting a compatible
$H$-structure.
Then there is a constant
$\varepsilon(M,g,H)>0$ such that if a
compatible harmonic $H$-structure
$\xi$ satisfies
$$
        E(\xi)
        =
        \|\nabla \xi\|_{L^2(M)}^2
        <
        \varepsilon(M,g,H),
$$
then $\xi$ is torsion-free:
$$
        \nabla \xi = 0.
$$
\end{theorem}

The upshot is that the flow is being
used to drive the structure into the
small-energy regime where the
$\varepsilon$-regularity theorem
forces torsion to disappear.
So the long-time picture is meant to
end at a torsion-free limit.

One can also reverse the procedure.
Let $(M,g)$ be such that there is no
torsion-free structure in the homotopy
class.
Then the question becomes whether the
flow still has a useful limit, or
whether it must break down.

I will recall results toward a general
theory for flows of H-structures,
obtained with Fadel, Loubeau and
Moreno.
Then I will present recent results on
SU(2)-flows on 4-manifolds, starting a
classification programme for
parabolic flows.
The method is representation-theoretic
and goes back to Bryant in G2 geometry;
the recent work is joint with Fowdar.

\medskip\noindent
The key concept to construct these
flows is the Fadel--Loubeau--Moreno--
SE 24 theorem, and also the Karigiannis
08 theorem: the ``diamond action''.
In this language the evolution equation
is
$$
        \partial_t \eta = A \diamond \eta
$$
with
$$
        A\equiv A(t)=S(t)+C(t),
        \qquad
        S(t)\in S^2,
        \qquad
        C(t)\in \Lambda^2.
$$

The freedom in choosing $A$ is the
freedom to choose a differential
invariant of the $H$-structure.
Bryant's 2003 point is that, for
$H\subset \operatorname{O}(n)$, the
space $V_k(\mathfrak h)$ of $k$th order
invariants is characterized by
$$
        V_k(\mathfrak h)\oplus
        (\Lambda^1\otimes S^{k+1})
        =
        (S^2\oplus\mathfrak h^\perp)
        \otimes S^k,
$$
where
$S^k=S^k(\Lambda^1)$ and
$\mathfrak h^\perp=\Lambda^2/\mathfrak h$.
The first space
$V_1(\mathfrak h)=\Lambda^1\otimes
\mathfrak h^\perp$ is the intrinsic
torsion.
The next piece $V_2(\mathfrak h)$ is
spanned by $\nabla T$ and
$\operatorname{Rm}$,
with overlap coming from the
Bianchi-type identity
$$
        \nabla T\diamond \eta
        =
        \nabla^2\eta+\hbox{l.o.t.}
$$
so the lower-order terms are not
independent.

In the simplest case $H=\operatorname{O}(n)$,
one has
$$
        V_1(\mathfrak{so}(n))=0.
$$
Also,
$$
        V_2(\mathfrak{so}(n))
        =
        \mathbf R
        \oplus
        S^2_0(\mathbf R^n)
        \oplus W,
$$
and
$$
        \operatorname{End}(\mathbf R^n)_
        {\mathfrak{so}(n)}
        =
        \mathbf R
        \oplus
        S^2_0(\mathbf R^n).
$$

\medskip\noindent
{\bf SU(2) flows.}

Any $SU(2)$-flow can be written as
$$
        \partial_t \omega = A\diamond\omega .
$$
The point of the theory is that the
possible choices of $A$ form a finite
dimensional space of differential
invariants.
So one can classify the flows.
For $SU(2)$ this leads to a list of
quadratic torsion functionals.
There are 15 of them, built from
terms of the form
$$
        \int_M g(a_i,a_j)\,\mathrm{vol}
        \qquad\hbox{and}\qquad
        \int_M g(a_i,J_j a_k)\,
        \mathrm{vol}.
$$
Among them are the Dirichlet energy
$$
        \int_M |T|^2\,\mathrm{vol}
$$
and the complex-geometric energy
$$
        \int_M |N_J|^2\,\mathrm{vol}.
$$
So the classification problem is
really a representation-theoretic
bookkeeping problem: write down the
allowed invariant tensors, then choose
a linear combination that gives a
parabolic flow.

\medskip\noindent
In other words, the flow is not
arbitrary.
It is built from the finite list of
quadratic expressions that the
representation theory allows.
The hard part is to choose the
coefficients so that the evolution
equation is still geometric and
parabolic.
That is why the talk keeps returning
to invariant tensors, torsion, and
curvature.

\medskip\noindent
{\bf Favorite flow.}

Consider the two-parameter rich flow
$$
        \partial_t\omega=
        \Bigl(
        -2\operatorname{Ric}(g)
        -\lambda_1\mathcal{L}_{
        \sum_{i=1}^2 J_i a_i}g
        +\lambda_2\sum_{i=1}^3
        \operatorname{div}(a_i)
        \Bigr)
        \diamond\omega .
$$
It has short-time existence and
uniqueness when
$$
        -4\sqrt 5<
        \lambda_2+\lambda_1/2<0 .
$$
This is the favorite flow.

\medskip\noindent
{\bf Questions and outlook.}

Can one add terms involving $W^+$
to the flow and still keep short-time
existence?
Are there good choices of such
curvature terms?
The Lie-derivative term can be more
generally taken as
$$
        \sum_{i,j}\lambda_{ij}
        \mathcal{L}_{J_i a_j}g.
$$
What lower-order terms should one
choose?
Are there flows preserving torsion
classes, for instance Laplacian-type
flows?
What about long-time existence,
singularity formation, and solitons?
More broadly, can the community
classify the good flows for other
groups $H$?

\medskip\noindent
{\bf Audience questions.}

Is one tacitly assuming the flow is
quasi-linear?
Is that crucial?

The answer is that one wants to
control low-order derivatives of the
torsion.
Once those are under control, the
rest follows.
So the point is to control the
derivatives of the norm of the
torsion.

Would it make sense to study flows
that are not quasi-linear?
For example, one could take the
Riemann curvature tensor, contract it
in many ways, and ask what happens.

Why would one do that?
The lower-order terms depend on which
PDE one chooses.
The hope is to use this as technology
for solving geometric PDEs.
That is the perspective here.

\medskip\noindent
{\bf Baby summary.}

The talk is about moving geometric
structures by PDEs.
The structure has torsion, and the
flow tries to reduce it.
Representation theory tells you which
terms you are allowed to use.
In the best cases, the flow is
well-posed and drives the structure
toward something more rigid.
The $SU(2)$ case is especially rich,
because there are many explicit
quadratic terms to choose from.

\section{Atom applications and
MMP}
\label{hodge-birational-atoms-2026-section-atom-applications-and-mmp}

\medskip\noindent
{\bf Speaker.}

Ludmil Katzarkov

\medskip\noindent
{\bf Title.}

Atom applications and MMP.

\medskip\noindent
{\bf Abstract.}

In this talk we introduce the
theory of Atoms.
Applications and connection with
MMP will be discussed.

\medskip\noindent
{\bf Birational invariants.}

\begin{enumerate}
\item LG-modules.
\item Orlov spectra.
\item Differential equations to the
monads.
\item Road map to the twistor
examples.
\end{enumerate}

\medskip\noindent
{\bf Orlov spectrum.}

Take a generator $g$ of
$D^b(X)$ and generate by taking
sums, summands, triangles and
shifts.
Let $t(g)$ be the minimum number of
triangles needed to get to any
object in $D^b(X)$.

The Orlov spectrum is the list of
generation times of
$D^b(X)$.
In the shorthand from the talk,
$$
D^b(X)=
\langle
t_1(g),\dots,t_n(g)
\rangle.
$$

\begin{enumerate}
\item $D^b(\mathbb{P}^1)$ has Orlov
spectrum $(1,2)$.
\item $D^b(E)$ has Orlov spectrum
$(1,2,2,4)$.
\end{enumerate}

Conjecture: this number is the
dimension of the variety.
This was confirmed in dimension $1$.

There are gaps in the spectrum.
A conjecture by Ballard, Favero and
Katzarkov says that if the gap is
greater than $\dim X-2$, then
$X$ is not rational.

\medskip\noindent
{\bf Differential equations.}

Now we take a shortcut and use
differential equations.
Keep in mind the Landau--Ginzburg
model, but something else will
happen.
Take a Fano manifold $X$ and
consider it as an algebraic
variety, but also as a symplectic
manifold.

\medskip\noindent
{\bf Main objective.}

If $X$ is a variety over
$\mathbb{C}$,
then birational geometry asks when
$$
\mathbb{C}(X)
\cong
\mathbb{C}(x_1,\dots,x_n),
\qquad
n=\dim X.
$$
This is the rationality problem in
its simplest form.

\medskip\noindent
{\bf Mirror symmetry.}

For smooth $X$, homological mirror
symmetry says
$$
D^b(X)
\leftrightarrow
\operatorname{FS}(Y,W),
$$
where $Y$ is an open symplectic
manifold and $W$ is a function.
The category on the right is built
from Lagrangian submanifolds, and
morphisms are given by Floer
homology.

\begin{example}
For $X=\mathbb{P}^2$ one has
$$
D^b(\mathbb{P}^2)=
\langle
\mathcal{O},
\mathcal{O}(1),
\mathcal{O}(2)
\rangle.
$$
The mirror is
$$
Y=(\mathbb{C}^*)^2,
\qquad
W=x+y+\frac{1}{xy}.
$$

The observation made by them is that
the singularity theory on the
symplectic side recovers, my guess,
the birational data on the algebraic
side.

\end{example}

\medskip\noindent
{\bf Blow-up.}

Now turn to birational geometry.
What happens when you blow up a
point of $\mathbb{P}^2$?

$$
D^b(\operatorname{Bl}_p \mathbb{P}^2)=
\langle
\mathcal{O},
\mathcal{O}(1),
\mathcal{O}(2),
E
\rangle,
$$

and on the symplectic side we get
$I_8$
(singularity).

\medskip\noindent
{\bf Cubic threefold.}

Next example:
$X_3 \subset \mathbb{P}^4$.
The classical approach is by studying
$$
H^{2,1}(X_3)/
H_3(X_3,\mathbb{Z}),
$$
that is, the intermediate Jacobian.
This gives a moduli space called the
Theta divisor
$$
\Theta.
$$

\begin{example}
\label{hodge-birational-atoms-2026-example-hms-point}
From the HMS point of view,
we have a symplectic manifold
and a potential
$$
W=
\frac{(x+y+z)^3}{xyz}+z.
$$
Some singular K3 surfaces appear.
There is also a complicated thing
at infinity.

Here one writes
$$
D^b(X_3)=
\langle
\mathcal{A},
E_1,
E_2
\rangle.
$$
\end{example}

\begin{theorem}[(-, Prizholkovsky, 2005)]
\label{hodge-birational-atoms-2026-theorem-let-x}
Let $X$ be an equidimensional Fano
with $\Pic X = \mathbb{Z}$,
and suppose $X$ is not
$\mathbb{P}^3$.
Consider the monodromy action on
near-by fibers of the symplectic
side.
The monodromy is a $3\times 3$
matrix.
If $\mu$ is strictly unipotent,
then $X$ is rational.
\end{theorem}

\medskip\noindent
{\bf Differential equations.}

Now we shall take a shortcut:
we use differential equations.
Keep in mind the
Landau--Ginzburg model,
but something else will happen.
Take a Fano manifold $X$ and
consider it as an algebraic
variety, but also as a symplectic
manifold.

We have de Rham cohomology, and also
a noncommutative Hodge package
$$
H_{nc}=
\bigoplus_{\lambda_i}
H_{\lambda_i}.
$$
The formula for $E_{\lambda_i}$ in
the notes is too unclear to
reconstruct exactly, so I do not
guess it here.

\begin{theorem}[Kontsevich,
\label{hodge-birational-atoms-2026-theorem-pantev-yu}
Pantev, Yu]
The object
$$
\bigoplus_{\lambda_i}
H_{\lambda_i}
$$
is a birational invariant.
\end{theorem}

The quantum differential equation is
$$
\left(
\frac{\partial}{\partial u}
-\frac{K}{u^2}
+\frac{G}{u}
\right)
\psi(u)=0.
$$
where
$$
\psi=u^{\delta}g(1,\dots).
$$

\begin{theorem}
\label{hodge-birational-atoms-2026-theorem-let-x2}
Let $X \subset \mathbb{P}^N$
be a Fano manifold of dimension
$\geq 3$.
Then
$$
\delta_{\max}
=
\dim X -1
\cdots
$$
If MMP works, and we can
...
then you might be able to conclude
that this $7$-dimensional cubic is
Picard as well.
\end{theorem}

\medskip\noindent
{\bf G-equivariant Hodge atoms.}

\medskip\noindent
{\bf Part 1.}

Review of coarse Hodge atoms.

\medskip\noindent
Let $X/\mathbb{C}$ be a smooth
projective variety of dimension at
least $2$.
For instance, one can think of a
compact K\"ahler manifold.
Let $G<\operatorname{Aut}(X)$ be a
finite subgroup.
We call $X$ a $G$-variety if $G$
acts generically freely on $X$.
The pair $(X,G)$ is non-$G$-
linearizable if there is no linear
$G$-action on $\mathbb{P}^n$ and no
$G$-equivariant birational map
between $X$ and $\mathbb{P}^n$.

\medskip\noindent
The equivariant version of the
coarse Hodge atoms keeps track of
the $G$-action on the Hodge
package.
The natural coarse version is
obtained by decomposing the atomic
package into $G$-isotypical
pieces.

\medskip\noindent
{\bf Application.}

These invariants give obstructions
to $G$-linearizability.
In particular, they produce many
examples of non-$G$-linearizable
$G$-actions on projective
varieties.

\section{Recent progresses in the
theory of atoms}
\label{hodge-birational-atoms-2026-section-recent-progresses-in-the-theory-of-atoms}

\medskip\noindent
{\bf Speaker.}

Leonardo Cavenaghi

\medskip\noindent
{\bf Title.}

Recent progresses in the theory of
atoms.

\medskip\noindent
{\bf Abstract.}

In this talk we report on new
results related to the theory of
atoms obtained in collaboration
with L. Katzarkov and M.
Kontsevich. They will appear in the
text titled "9.5 lectures in the
theory of atoms", in the final
stage of preparation.

\medskip\noindent
{\bf Upshot.}

The point is to package the
Hodge-theoretic and quantum
information into atoms.
These atoms carry both the
representation-theoretic and
birational content.

\medskip\noindent
{\bf Relation with
Mumford--Tate groups.}

This fits the notes in
\ref{mumford-tate-%
section-mt-hodge-structures}.
There the Mumford--Tate group
controls Hodge symmetries.
The same Hodge package is also
developed in the complex-geometry
notes.
Here the Hodge atoms are a more
refined bookkeeping device:
they split the Hodge package into
pieces on which
$
\operatorname{Hod}(\mathbb{k})
$
and its equivariant variants act.

\medskip\noindent
{\bf G-equivariant atoms and
quantum cohomology.}

\medskip\noindent
{\bf Theorem.}

For any multi-degree hypersurface
$$
X:=X(1,1,1,1)\subset(\mathbb{P}^1)^4,
$$
for instance the blow-up of
$\mathbb{P}^1\times\mathbb{P}^1
\times\mathbb{P}^1\times\mathbb{P}^1$,
the invariant $\mathbb{Z}/2$-action
by swapping the first and second,
and the third and fourth factors,
is not $\mathbb{Z}/2$-linearizable.
The same holds for
$$
X\times\mathbb{P}^1,
$$
where $\mathbb{Z}/2$ acts trivially
on the $\mathbb{P}^1$ factor.

\medskip\noindent
{\bf Coarse atoms.}

For this we introduce a coarse
version of the theory of
$G$-equivariant atoms.

\medskip\noindent
{\bf Notation.}

Take your projective complex variety  $X$
and its rational
cohomology
$$
V:=H^*(X,\mathbb{Q}),
$$
seen as a supervector space.
\begin{itemize}
\item Coordinates:
$$
\mathbf{t}=\{t_a\}
$$
formal dual super coordinates.
\item Novikov ring:
to track effective curve classes
$\beta \in \overline{\operatorname{NE}}
(X,\mathbb{Z})$.
The small Novikov ring is
$$
\mathcal{N}_X
=\left\{\sum c_\beta q ^\beta:\beta
\in \overline{\text{NS}}(X,\mathbb{Z})
\right\}
$$
\noindent
And then
$$
\text{Nov}_X
=\mathcal{N}_X [[t_a,a \text{is even}]]
\hat{\otimes}\Lambda(ta,a\text{ is odd})
$$

We ``enrich'' the ring with those
power series.
\item Coefficient ring:
formal power series.
\item Take a homogeneous basis of $V$:
$$
\{\Delta_0,\Delta_1,…
,\Delta_\rho,\Delta_{\rho+1},…
,\Delta_r\}
$$
Think that
$t_a \leftrightarrow \Delta_a$
are (super) coordinates.
\end{itemize}
\noindent
Question: what is a generating function
to the GW invarirants of $X$?
A $\phi_X \in \text{Nov}_X$
with properties! (WDVV equations)

The Gromov--Witten potential allows
us to define a new product in this
ring that recovers the classical
Poincar\'e cup product.
\begin{itemize}
\item Classical part:
intersection theory.
\item Quantum correction:
curve counting, enumerative
geometry.
\end{itemize}

\medskip\noindent
{\bf Big quantum cohomology.}

\begin{itemize}
\item Space:
$$
V_{\mathrm{big}}=
V\otimes_{\mathbb{Q}}\mathcal{N}_X.
$$
\item Product: third-order
differential operator.
\item Deformation: full deformation
in $\mathbf{t}$.
\end{itemize}

\medskip\noindent
{\bf Small quantum cohomology.}

\begin{itemize}
\item Space:
$$
V\otimes\mathcal{N}_X.
$$
\item Product:
[missing]
\item Deformation:
[missing]
\end{itemize}

\begin{example}
For $X=\mathbb{P}^2$, the cohomology
has three generators:
$$
\{t_0,t_1,t_2\}=
\{\Delta_0,\Delta_1,\Delta_2\}
=\{1,H,P\},
$$
Now since
$H_2(\mathbb{P}^2,\mathbb{Z})
=\mathbb{Z}[\ell]$,
we think $q^\beta \leftrightarrow
q^d$.
The Gromov--Witten potential is
$$
 \phi_X(\mathbf{t},q)
=\underbrace{
\frac{1}{2}t_0^2t_1
+\frac{1}{2}t_0t_1^2}
_{\text{classical part}}
+\underbrace{\sum_{d=1}^{+ \infty}
N_d (t_2^{3d-1})/(3d-1)!
q^3 e^{dt_1}}_{\text{quantum part}}
$$
where $N_d$ is the number of rational curves
of degree $d$ passing through $3d-1$
points in general position in $\mathbb{P}^1$.
We can write the term $q^3e^{dt_1}$ as
$$
q^de^{dt_1}=(qe^{t_1})^d.
$$
As a first remark note that,
 a priori, $\phi_{\mathbb{P}^2}
\in \text{Nov}_X$,
but because we have that element 
$q^de^{dt_1}$, …,
we can actually see that
$$
\phi_{\mathbb{P}^2} \in
\mathbb{Q}[[qe^{t_1},t_2]].
$$

A second remark says how to compute the
quantum part:
$$
N_d= \left\langle{}\underbrace{p,…,p}_{3d-1}\right\rangle{}
_{0,3d-1,\beta=[Z?]}
$$

\end{example}
Kontsevich noted (a while ago) that
WDW $\iff$ recursion formula for computing
$N_d$.

We have chosen 
$$
\Delta_1,…,\Delta_{\rho}
$$
separately. Assume that this is a basis
of $\text{NS}(X) \subset H^2(X,\mathbb{Q})$.
Goal: quantum multiplication, very
geometrically, builds products on the tangent
bundle of some space.

For $\mathbb{P}^1$, the GW potential
is an entire function.

More generally,
GW potential has a positive radius of
convergence [check this…].

\medskip\noindent
{\bf Questions.}

\begin{enumerate}
\item Analytic germ.
\item Tangent space product.
\item Logarithmic tangent bundle.
\end{enumerate}


\medskip\noindent
Now consider
$$
q^\beta
=\prod_{i=1}^\ell q_i^{\int \Delta_i} \beta
$$
and
$$
\mathcal{N}_X [[ (t_a)_{a \not \in
\{1,…,\ell\}}]]
\otimes_{\mathbb{Q}}
\Lambda(t_a; a \text{ is odd})
=\text{Nov}_X^{\text{red}}.
$$
The idea is to take a space such that
its ring of functions is precisely
$\text{Nov}_X^{\text{red}}$,
we call it
$\text{Spf}(\text{Nov}_X^{\text{red}}$,
and do
$$
\Delta_a \leftrightarrow
\frac{\partial }{\partial t_a}
\text{ \& }
\Delta_a \leftrightarrow
q_s \frac{\partial }{\partial q_i}
$$
$$
a \not \in \{ 1,…,\ell\},
q_i \frac{\partial }{\partial q_i}
q^\beta)=
\left(\int_\beta \Delta_i\right)
q^{\beta}.
$$
Passing from the formal picture to the
rigid picture, one can show how to actually
have a positive radius of convergence
for $\phi_X \leftrightarrow$
replace with
$\text{Spf}\text{Nov}_X^{\text{red}}$.

Now we are ready to define the quantum
multiplication.
It's given by derivatives of third order:
$$
\Delta_a \star \Delta_b 
:= \sum_{c,d}
\frac{\partial^3}{\partial t_a \partial
t_b \partial t_c}\eta^{dd^c}\Delta_c.
$$

Big quantum product relations.

\medskip\noindent
{\bf General logarithmic tangent
bundle.}

Idea: to have a germ of a manifold
such that the tangent space is the
cohomology vector space and the
quantum product is an associative
product in this vector space.

For $X/\mathbb{C}$ smooth
projective, let
$$
V:=H^*(X,\mathbb{Q}).
$$
Let $\Phi_X$ be the
Gromov--Witten potential.
Then the third derivatives
$\partial^3\Phi_X$ yield a unital,
commutative, associative product on
the tangent bundle of a germ
$$
M\subset H^*(X,\mathbb{C}).
$$
This is the Frobenius-manifold
picture behind quantum cohomology.

\medskip\noindent
The problem is that convergence is
most of the time hypothetical.
The strategy is to introduce an
auxiliary non-Archimedean field
$\mathbb{k}=\overline{\mathbb{Q}}((T^Q))$
(which is a Puiseaux seris ring) 
of characteristic $0$
to ensure convergence.
The idea is to replace the germ
$M$ by a manifold $B_X$ where we do
have convergence for the
Gromov--Witten invariants.
More precisely, for $b \in B_X$,
$$
T_bB_X\simeq H^*(X,\mathbb{k},\star_b).
$$
That is: there exists a concrete
supermanifold
with analytic topology $B_X$, for which
we claim: $T_bB_X$ has the structure
of a commutative associative algebra (WDVV).

\begin{definition}
\label{hodge-birational-atoms-2026-definition-manifold-b}
The manifold $B_X$ is called an
{\it $A$-model manifold}.
\end{definition}

Now we put an Euler vector field
on $B_X$.
If $\star$ is a tensor field on $B_X$
which is actually a Higgs field,
then $\operatorname{Eu}$ scales the
product.
That is,
$$
\mathcal{L}_{\operatorname{Eu}}(\star)=\star.
$$
Equivalently, we have an
automorphism
$$
k_b:=\operatorname{End}_{\operatorname{Eu}}
(\star_b)
\in \operatorname{End}
\bigl(H^*(X,\mathbb{k})\bigr).
$$
That is the starting point of atom
theory, from our perspective.
There are other starting points.

\medskip\noindent
{\bf Disjoint unions.}

Thom--Sebastiani principle.
If a variety is a disjoint union
$$
X\sqcup Y,
$$
its Gromov--Witten potential
uncouples as a sum of the factors:
$$
\Phi_{X\sqcup Y}
=
\Phi_X \oplus \Phi_Y.
$$
This is important because we need
to study blow-ups, the main
operations in birational geometry.

\begin{theorem}[Iritani, 2003]
\label{hodge-birational-atoms-2026-theorem-let-x3}
Let $\widetilde X$ be the blow-up
of $X$ along a smooth subvariety
$Z\subset X$ of codimension
$r\geq 2$.
There exists a canonical analytic
identification between open
domains
$$
\psi:
B_{\widetilde X}\supset U
\xrightarrow{\sim}
V\subset
B_X\times B_Z\times\cdots\times B_Z.
$$
Moreover, one can choose the open
domains so that the Euler vector
field and the product are
preserved:
$$
\psi^*(*)=*,
\qquad
\psi^*(\operatorname{Eu})
=\operatorname{Eu}.
$$
\end{theorem}

\medskip\noindent
{\bf Hodge action and fixed locus.}

$\operatorname{Hod}(\mathbb{k})$
acts on $B_X$ preserving the quantum
product $*$ and the Euler field
$\operatorname{Eu}$.

On coordinates:
\begin{itemize}
\item $U_1$ (N\"eron--Severi):
trivial action, corresponding to
algebraic classes.
\item $U_2, U_{\mathrm{odd}}$:
scaling by Hodge weights.
\end{itemize}

\begin{definition}
\label{hodge-birational-atoms-2026-definition-hodge-fixed}
The {\it Hodge fixed locus}
$$
B_X^{\operatorname{Hod}}\subset B_X
$$
is the subvariety where
$\operatorname{Hod}(\mathbb{k})$
acts trivially.
Geometrically, it is the locus
where all non-$(p,p)$ coordinates
vanish.
\end{definition}

$$
\mathcal{HS}\simeq\operatorname{Rep}
\operatorname{Hod}
$$
for $\mathbb{Q}$-polarizable pure
Hodge structures.

We can decompose the tangent space
to $B_X$ using generalized
eigenspaces:
$$
T_bB_X=
\bigoplus_{\lambda\in\operatorname{Spec}
(k_b)}V_\lambda.
$$
The catch is that the group is
pro-reductive.
This means we can write, in fact,
$$
T_bB_X=
\bigoplus_{\lambda\in\operatorname{Spec}
(k_b)}
\bigoplus_{\gamma}
V_\lambda\cap V_\gamma.
$$
The summands
$$
\bigoplus_{\gamma}
V_\lambda\cap V_\gamma
$$
are called {\it coarse Hodge atoms}.

Each atom $\alpha$ corresponds to a
$\overline{\mathbb{Q}}$-linear
representation
$$
E^\alpha
$$
of
$\operatorname{Hod}_{\overline{\mathbb{Q}}}$,
independent of the point $b$.

\medskip\noindent
{\bf Numerical invariants.}

The Hochschild grading $(p-q)$ on
$E^\alpha\otimes_{\overline{\mathbb{Q}}}
\mathbb{C}$ defines:
\begin{enumerate}
\item Rank of Hodge classes:
$$
\rho_\alpha
=
\dim
\left(V_\lambda\cap H_{p,q}\right).
$$
\item Hodge polynomial.
It is characterized by the
coefficient:
$$
P_\lambda(t)=
\dim\left(
V_\lambda\cap H^{p-q=n}(X)
\right).
$$
\end{enumerate}

\begin{theorem}[KKPY
\label{hodge-birational-atoms-2026-theorem-non-rationality}
Non-rationality criterion]
Let $X$ be an $n$-dimensional
smooth complex projective variety.
If there is a coarse Hodge atom in
the Hodge atomic composition of $X$
that does not appear in the Hodge
atomic composition of any variety
of dimension $\leq n-2$, then $X$
is not rational.
\end{theorem}

\medskip\noindent
{\bf Presence of $G$-actions.}

The action on the variety defines
an action on the cohomology.
Now one does not look at
$\operatorname{Hod}(\mathbb{k})$
alone, but at
$$
\operatorname{Hod}(\mathbb{k})
\times G,
$$
which acts on $B_X$.
Restricting to
$$
B_X^{G\times\operatorname{Hod}},
$$
a coarse $G$-atom $\alpha$ is the
isomorphism class of a
$G\times\operatorname{Hod}(\mathbb{k})$-
representation on
$\mathcal{B}_{\lambda,b}$.

\medskip\noindent
{\bf $G$-invariant Hodge classes.}

$$
\rho_\alpha^G:=
\dim_{\mathbb{C}}
\left(
E^\alpha\otimes_{\overline{\mathbb{Q}}}
\mathbb{C}
\right)^{
G\times\operatorname{Hod}
}.
$$

\medskip\noindent
{\bf Equivariant constraint.}

Since the algebraic unit $1$ is
always $G\times\operatorname{Hod}$
invariant, we have
$$
\rho_\alpha^G\geq 1.
$$

\medskip\noindent
{\bf Obstructions to
$G$-linearizable actions.}

\medskip\noindent
{\bf Atomic constraints in low
dimensions ($\leq 1$).}

Atoms originating from points
$(G/H)$ or curves
$(G\times_H C)$ satisfy:
\begin{itemize}
\item Rank bound:
$$
\rho_\alpha^G\in\{1,2\}.
$$
\item Rigidity:
if $\rho_\alpha^G$ is constant,
there is only one invariant Hodge
class.
\end{itemize}

\begin{example}
The hypersurface
$$
X(1,1,1,1)\subset(\mathbb{P}^1)^4.
$$
Rationality:
$X$ is rational, the blow-up of
$\mathbb{P}^1\times\mathbb{P}^1
\times\mathbb{P}^1$
along an elliptic curve $E$.
The $G$-action is
$G=\mathbb{Z}/2$, swapping factors
1 $\leftrightarrow$ 2 and
3 $\leftrightarrow$ 4.

The atomic decomposition of $X$
yields three primary coarse
$G$-atoms:
$$
\begin{array}{c|c|c|c}
\lambda & \rho & \rho^G & P_\alpha(t)\\
16 & 1 & 1 & 1\\
4 & 4 & 2 & 4\\
0 & 5 & 3 & t+3+t^{-1}
\end{array}
$$
Note that
$$
\rho_{\mathrm{total}}^G=1+2+3=6,
$$
tracking the $G$-invariant Hodge
classes.

The atom of the elliptic curve
must contribute.
It contributes with:
\begin{itemize}
\item points $(1,1)$ or $(2,1)$.
\item curve $E$:
$(2,2)$.
\end{itemize}

\end{example}

Target atom $\lambda=0$:
$$
(\rho,\rho^G)=(5,3).
$$

If the action was linearizable,
its atoms must be constructible
from atoms of $\mathbb{P}^3$ and
$G$-invariant subvarieties of
dimension $\leq 1$.

The contradiction:
...

\begin{theorem}[Rationality criterion]
\label{hodge-birational-atoms-2026-theorem-x-smooth}
If $X$ is a smooth complex
projective variety of dimension
$n\geq 2$, and the maximal generic
spectral dimension on the Hodge
locus $B_X^{\operatorname{Hod}}$
satisfies
$$
\mu(X)>n-2,
$$
then $X$ is not rational.
\end{theorem}

\medskip\noindent
{\bf Quantum differential equation.}

Over the non-Archimedean field,
recall that we have a manifold and
then we put the quantum
differential equation.

There is a bundle $(\mathcal{H},\nabla)$
with a flat connection $\nabla$
and we have the equation
$$
\nabla_{u \frac{\partial }{\partial u}}
=\frac{\partial }{\partial u}
+ \frac{1}{u} \text{Eu}\star +
\text{Gr}
$$
where $\text{Gr}$ is the grading operation.
This is the {\it quantum differential
equation}.

We solve it via by an inf
$$
\psi(u) = u^{-5}e^{\lambda/y}
(\psi+…)
$$
\begin{definition}
\label{hodge-birational-atoms-2026-definition-dim-spec}
$$
\dim_{\Spec}(\lambda)
=-2\text{min}\left\{
s \in \mathbb{Q}:
\substack{\exists \psi(u)
\text{ solving} \\ \text{QDE
associated to $\lambda$ w exp. $s$}}\right\}
$$

\begin{theorem}[inf]
\label{hodge-birational-atoms-2026-theorem-inf2}
At a very general point
$$
\dim_{\Spec}(0)
$$
obstructs rationality.
If $X$ is a smooth cubic threefold
then $\lambda \in \{0 , \pm 1\}$.
Solution to QDE on the direction of $0$
$s_1=-1/6$, $s_2=-5/6$.
$\dim_0 = 5/3>3-2=1$ then $X$ is irrational.
\end{theorem}

\end{definition}

After a big analysis, we get for
the cubic, indeed,
$$
\frac{5}{3}.
$$

\section{Finite simple subgroups of the
real Cremona group of rank three}
\label{hodge-birational-atoms-2026-section-finite-simple-subgroups-real-cremona-group-rank-three}

\medskip\noindent
{\bf Speaker.}

Antoine Pinardin

\medskip\noindent
{\bf Title.}

Finite simple subgroups of the real
Cremona group of rank three.

\medskip\noindent
{\bf Abstract.}

Very little is known about the
classification of finite subgroups of
Cremona in dimension three. It is
natural to start with the case of
simple groups, and this step was
achieved by Prokhorov in 2009 over the
field of complex numbers. In the work
I will present, we show that the only
non-cyclic finite simple subgroups of
the real Cremona group of rank three
are A5 and A6. This is a joint
project with I. Cheltsov and
Y. Prokhorov.

\medskip\noindent
{\bf Upshot.} To prove the theorem
of classification of non-cyclic finite simple
subgroups of
$\text{Cr}_3(\mathbb{R})$,
i.e. Theorem \ref{hodge-birational-atoms-2026-theoremcheltsov-prokhorov},
which says they are $A_5$ and $A_6$,
we introduce group actions on the varieties.
Note that there is no mention of such
actions in the notes of Opega II.

\medskip\noindent
{\bf AI-generated proof strategy outline.}

The bridge theorem identifies a
finite subgroup of the Cremona group
with a rational $G$-variety, modulo
equivariant birational change of
model.
From there, equivariant MMP reduces
the problem to a $G$-Mori fiber
space, and in the cases discussed
here to a $G$-Fano model.
The big diagram records this passage
from the original action to the
fiber/base decomposition.
Once that reduction is in place, the
remaining possibilities are checked
against the known Fano-type examples
and the Sarkisov-link factorization
picture.

\medskip\noindent
{\bf Notes.}

Take a field of characteristic $0$.
Define
$$
\operatorname{Cr}_n(k)
=
\operatorname{Bir}(\mathbb{P}^n)
=
\operatorname{Aut}_k
(k(x_1,\ldots,x_n)).
$$
This is a huge group.
A question of Serre is
whether any finite
subgroup injects in 
$\text{Cr}_3(\mathbb{C})$.

For $n=1$, $k=\mathbb{C}$, we have
$$
\operatorname{Cr}_1(\mathbb{C})
=
\operatorname{PGL}_2(\mathbb{C}),
$$
whose finite subgroups were
classified by Klein:
$$
Z_n,\ D_n,\ A_4,\ S_4,\ A_5.
$$

For $n=2$, $k=\mathbb{C}$:
\begin{itemize}
\item Bertini, 1877.
\item Bangalo--Beauville, 2000:
subgroups isomorphic to
$\mathbb{Z}_2$
up to conjugation.
\item
Belzance, 2006:
finite abelian subgroups.
\item Dolgachev--Iskovskikh, 2009:
``all'' (there are conjugation
subtleties) finite subgroups of
$\operatorname{Cr}_2(\mathbb{C})$.
\end{itemize}

\[
\begin{array}{c|ccc}
& A_5 & A_6 &
\operatorname{PSL}_2(\mathbb{F}_7)
\\
\hline
\text{Conj.\ classes}
& 3 & 1 & 2 \\
\text{Models}
&
\mathbb{P}^2,\ dP_5,\ \mathbb{P}^1\times
\mathbb{P}^1
&
\mathbb{P}^2
&
\substack{\mathbb{P}^2,\\ 
\ dP_2
\overset{2:1}{\to}
\{x^3y+y^3z+z^3x=0\}
\subset \mathbb{P}^2\\
\text{i.e. a double cover of}\\
\text{Klein's quartic}}
\end{array}
\]

For $n=3$, $k=\mathbb{C}$, this is
widely open.

\begin{itemize}
\item Progress for abelian groups,
2025.
\item (Prokhorov, 2009.)
Noncyclic finite simple subgroups:

\begin{itemize}
\item $A_5$: infinitely many.
\item $A_6$: finitely many? (not known!)
\item $\operatorname{PSL}_2(7)$:
finitely many? (not known!)
\item $A_7$: one,
$\mathbb{P}^3$.
\item $\operatorname{SL}_2(8)$:
one,
Prokhorov's threefold.
\item $\operatorname{PSpin}_4(3)$:
two,
$\mathbb{P}^3$ and
Birkhodt's quartic.
\end{itemize}
\end{itemize}

\medskip\noindent
Further notes: 
\begin{itemize}
\item For Prokhorov's threefold:
$$
\operatorname{LGr}(4,9)\cap
\mathbb{P}^8
\subset
\mathbb{P}^{15}.
$$

\item For Birkhodt's quartic:
$$
\{x_0^4-x_0\sum_{i=1}^3 x_i^3
+\cdots\}
\subset
\mathbb{P}^4.
$$
\item For $A_5$:
book:
``Cremona groups and the
icosahedron'',
Cheltsov and Yomdin, 2016.
\end{itemize}


\medskip\noindent
{\bf Purpose of today.}

Take $k=\mathbb{R}$.

For $n=1$,
there is no non-cyclic finite group
contained in
$\operatorname{Cr}_1(\mathbb{R})$.

For $n=2$,
$A_5$ acts on
$\mathbb{P}^2$,
$\mathbb{P}^1\times\mathbb{P}^1$,
and $dP_5$.
This is due to Yaskovsky.

For $n=3$:

\begin{theorem}
\label{hodge-birational-atoms-2026-theorem-cheltsov-prokhorov}
[Cheltsov, Prokhorov,
Pinardin, 2025]
The non-cyclic finite simple
subgroups of
$\operatorname{Cr}_3(\mathbb{R})$
are $A_5$ and $A_6$.
\end{theorem}
\medskip\noindent
{\bf Method of proof.}

\begin{definition}
\label{hodge-birational-atoms-2026-definition-g-variety}
An {\it $G$-variety} is a pair
$(X,\rho)$ with $X$ projective and
$\rho:G\to\operatorname{Aut}(X)$.
\end{definition}

\begin{definition}
\label{hodge-birational-atoms-2026-definition-g-bir}
An {\it $G$-birational map} is
$\varphi:(X,\rho)\dashrightarrow
(X',\rho')$
such that
$\varphi\rho(g)\varphi^{-1}=\rho'(g)$
for all $g\in G$.
If such a $\varphi$ exists,
then $X$ and $X'$ are
{\it $G$-birationally equivalent}.
\end{definition}

\begin{theorem}
\label{hodge-birational-atoms-2026-theorem-finite-subgroups}
Finite subgroups of
$\operatorname{Cr}_n(k)$,
up to conjugacy,
correspond to rational $G$-varieties
for finite $G$,
up to equivariant birational maps.
\end{theorem}

\begin{definition}
\label{hodge-birational-atoms-2026-definition-g-mori-fiber-space}
We may assume $X$ is a
{\it $G$-Mori fiber space}.
That is, $X$ is a terminal $G$-variety,
all invariant divisors are 
$\mathbb{Q}$-Cartier, and:
\begin{itemize}
\item There exists
$\pi:X\to B$ surjective and $G$-equivariant
with $\dim B<\dim X$.
\item $K_X$ is $\pi$-ample.
\item
$\operatorname{rk}\operatorname{Pic}^G(X)$
\[
=1+
\operatorname{rk}\operatorname{Pic}^{G_B}(B).
\]
\item The fibers are connected.
\end{itemize}
\end{definition}

In particular, if $B=\operatorname{pt}$,
then $X$ is a minimal $G$-Fano
variety.

\begin{remark}
\label{hodge-birational-atoms-2026-remark-1-k}
Since
$$
\xymatrix{
1\ar[r]&  K\ar[r]\ar[d]&  G \ar[r]& 
G_B\ar[d]& 1\\
&  \text{fibers of }\pi& & B
}
$$
we may assume that $X$ is Fano.

For $\mathcal{L}_0$ and simple
groups, we may assume $X$ is Fano.
\end{remark}


\begin{remark}
\label{hodge-birational-atoms-2026-remark-cremona-real-containd-complx}
$\text{Cr}_3(\mathbb{R}) \subset 
\text{Cr}_3(\mathbb{C})$.
\end{remark}

\medskip\noindent
$A_5$ acts on:
\begin{itemize}
\item $\mathbb{P}^3$.
\item The quadric
$$
x_0^2+x_1^2+x_2^2+x_3^2=x_4^2
\subset \mathbb{P}^4.
$$
\item $\mathbb{P}^1\times Q$,
where
$$
Q=
\{x_0^2+x_1^2+x_2^2=x_3^2\}.
$$
\item The Segre cubic.
\end{itemize}

\medskip\noindent
$A_6$ acts on:
\begin{itemize}
\item The Segre cubic.
\item
$$
\begin{cases}
\sum x_i^3=0,\\
\sum x_i=0
\end{cases}
\subset \mathbb{P}^5.
$$
\end{itemize}

The group $S$ acts on $X$ by
permuting the coordinates.

$$
\xymatrix{
&& \widetilde{X}
\ar[dl]_{5\ \text{pts}}
\ar[dr]^{\text{contract $10$ lines}}\\
& \mathbb{P}^3 && X
}
$$

\medskip\noindent
The first part of this project
was obtaining the following result:

\begin{proposition}
\label{hodge-birational-atoms-2026-proposition-x-real}
If $X$ is a real rational
threefold and
$$
G\in
\{
\operatorname{PSL}_2(7),
\operatorname{SL}_2(8),
A_7,
\operatorname{PSpin}_4(3)
\},
$$
then there is no faithful action of
$G$ on $X$.
\end{proposition}

\noindent
In the following diagram ``BAD''
means how difficult it was to obtain
the result.

$$
\xymatrix@C=1.2em@R=2.5em{
&& \leavevmode
\ar[dl]_{OK}
\ar[dr]^{BAD}
&&
\\
& \substack{BIG\text{, i.e. one of}\\
\{\text{SL}_28, A_7,PSp_43\}}
&&
G=\operatorname{PSL}_2(7)
\ar[dl]
\ar[dr]^{BAD\ BAD}
&
\\
&& X\ \text{Gor.}
&&
X\ \text{non-Gor.}
\ar[dl]
\ar[dr]^{BAD\ BAD\ BAD}
&
\\
&&& |-K_X|\neq\emptyset
&&
|-K_X|=\emptyset
&
}
$$

\begin{remark}[Terpenan-Simmerman]
\label{hodge-birational-atoms-2026-remark-terpenan-simmerman}
Let
$$
X_{\mathbb{C}}=
\{x_0y_0+x_1y_1+x_1y_2=0\}
\subset
\mathbb{P}^2_{\mathbb{C}}
\times
\mathbb{P}^2_{\mathbb{C}}.
$$
$X_\mathbb{C}$ has a real form $X_\mathbb{R}$
such that $\text{PSL}_2(7) \subset
\text{Aut}(X_\mathbb{R})$.

Moreover, $X_\mathbb{C}$ is rational.

This means that $
X_\mathbb{C}$ has a real form 
$X_\mathbb{R}$ with an
action of a group ---
$\operatorname{PSL}_2(7)$.
The variety $X_\mathbb{C}$ is rational, but
$X_\mathbb{R}$ is not.
\end{remark}

\medskip\noindent
In this extra time we shall explain
more details about the
uppermost leftmost item in the big
diagram above. This was 
deduced from the work of Prokhov
(and something else).

Let $I$ be BIG, i.e.
$\in \{\text{SL}_2(8),A_7,PSp_4(3)\}$.
Now, $A_7$ acts on $X_\mathbb{C} =
\mathbb{P}^3$,
$\text{SL}_2(8)$ acts on Prokhov's
threefold, $LG(4,9)\cap \mathbb{P}^8
\subset \mathbb{P}^{15}$,
and finally $PSp_4(3)$ acts on
$\mathbb{P}^3$ and Birkhodt's threefold
$$
\left\{x_0^4-x_0 \sum_1^3 x_i^3
+3 \prod_0^3 x_i\right\}.
$$
Over $\mathbb{R}$, we want to list Fano
threefolds $X$ such that $G \subset
\text{Aut}(X)$.

\section{On metric connections with
parallel and closed skew torsion}
\label{hodge-birational-atoms-2026-section-on-metric-connections-with-parallel-and-closed-skew-torsion}

\medskip\noindent
{\bf Speaker.}

Paul Schwahn

\medskip\noindent
{\bf Title.}

On metric connections with parallel
and closed skew torsion.

\medskip\noindent
{\bf Abstract.}

This talk is divided in three parts.
First, I review structural properties
and examples of geometries with
parallel skew torsion.
Second, I take a closer look at those
geometries where the torsion is in
addition closed; this equips every
tangent space with the structure of a
metric Lie algebra. These manifolds
always locally split as a product into
well understood factors, allowing a
complete local classification.
Third, I investigate various
$G$-structures defining a geometry of
this type, with a focus on almost
Hermitian structures, extending recent
results of Barbaro-Pediconi-Tardini.
This is joint work with Andrei
Moroianu.

\medskip\noindent
{\bf Table.}

\[
\begin{array}{|l|l|l|l|l|l|}
\hline
\parbox[t]{7em}{\raggedright
$(M,g)$}
&
\nabla^\tau
&
\parbox[t]{4em}{\raggedright
$\mathfrak h$}
&
\underline{\operatorname{Hol}
(\nabla^\tau)}
&
\parbox[t]{5em}{\raggedright
$
\underline{\operatorname{Stab}}
$
\\
$(\tau)$}
&
d\tau=0?
\\
\hline
\parbox[t]{7em}{\raggedright
$G/H$ isotropy-irred.,\\
non-symmetric}
&
\text{canonical}
&
\parbox[t]{4em}{\raggedright
$\mathfrak h$}
&
\mathfrak h
&
\parbox[t]{5em}{\raggedright
$\mathfrak j$ or\\
$\mathfrak g_2$\\
(Berger)}
&
\text{No}
\\
\hline
\parbox[t]{7em}{\raggedright
$(H\times H)/H$, $\mathfrak h$,\\
$H^\mathbb{C}/H$\\
(bi-invariant metric)}
&
\text{Cartan-Schouten}
&
\nabla^t,\ t=0
&
\parbox[t]{4em}{\raggedright
$\mathfrak h$}
&
\parbox[t]{5em}{\raggedright
0 on $(H\times H)/H$,\\
$t=\pm1$ or $\mathfrak h$}
&
\text{Yes}
\\
\hline
\end{array}
\]

\medskip\noindent
{\bf Special cases.}

\[
\begin{array}{|l|l|l|l|l|l|}
\hline
\text{Geometry}
&
\nabla^\tau
&
\mathfrak h
&
\underline{\operatorname{Stab}}(\tau)
&
d\tau=0?
\\
\hline
\shortstack[l]{6-dimensional\\
nearly K\"ahler}
&
\text{characteristic}
&
\underline{\mathfrak{su}(3)}
&
\parbox[t]{13em}{\raggedright
$\underline{\mathfrak{so}(3)}$\\
or $\underline{\mathfrak{su}(3)}$\\
$(S^3\times S^3)$\\
$\underline{\mathfrak{u}(1)\oplus
\mathfrak{u}(1)}$\\
$(\operatorname{SU}(3)/T^2)$\\
$\underline{\mathfrak{su}(2)\oplus
\mathfrak{u}(1)}$\\
(twistor space)}
&
\underline{\mathfrak{su}(3)}
&
\text{No}
\\
\hline
\shortstack[l]{7-dimensional\\
nearly parallel\\
$G_2$}
&
\text{connection}
&
\underline{\mathfrak g_2}
&
\parbox[t]{13em}{\raggedright
$\mathfrak{so}(3)_{\mathrm{inv}}$\\
(Berger space)\\
$\mathfrak{so}(4)$\\
(3-or(5)-Sos)\\
or $\underline{\mathfrak g_2}$}
&
\underline{\mathfrak g_2}
&
\text{No}
\\
\hline
M^3
&
\nabla^g+t\,\operatorname{vol}_g
&
\underline{\mathfrak{so}(3)}
&
\parbox[t]{12em}{\raggedright
0\ (S$^3$, irred.)\\
$\underline{\mathfrak u}(1)$\ 
(-Sasakian)\\
or $\underline{\mathfrak{so}(3)}$}
&
\underline{\mathfrak{so}(3)}
&
\text{Yes}
\\
\hline
\end{array}
\]


\noindent
Curiosity: the 
$\mathfrak g_2$ (Berger) appearing
in the first table
is the very interesting
space
$\operatorname{SO}(5)/
\operatorname{SO}(3)_{\mathrm{irr}}$.

\medskip

\begin{definition}
\label{hodge-birational-atoms-2026-definition-geometry-parallel}
A {\it geometry with parallel skew
torsion} (PST) is a Riemannian
manifold
$(M,g)$ with a metric connection
$\nabla^\tau$
whose torsion $T^\tau$ is skew,
that is,
$T^\tau\in\Omega^3(M)$,
and
\[
\nabla^\tau T^\tau=0.
\]
\end{definition}

\medskip\noindent
{\bf Notation.}

We write
\[
\nabla^\tau=\nabla^g+\tau,
\qquad
\tau\in\Omega^3(M),
\qquad
T^\tau=2\tau.
\]
For a form $\alpha$,
\[
\nabla^\tau_X\alpha
=
\nabla^g_X\alpha
+(\tau_X)\cdot\alpha.
\]

\begin{example}
\label{hodge-birational-atoms-2026-example-manifolds-killing}
\begin{enumerate}
\item Manifolds with Killing spinors:
nearly K\"ahler $NK^6$,
nearly parallel $G_2$,
$(\alpha)$-Sasaki,
and $3$-$(\alpha,\delta)$-Sasaki.
\item Naturally reductive homogeneous
spaces.
Let $M=G/H$ and assume that
\[
\mathfrak g=\mathfrak h\oplus\mathfrak m.
\]
This gives a canonical connection
$\nabla$ with
\[
T(X,Y)=-[X,Y]_{\mathfrak m},
\qquad
R(X,Y)Z=
-[[X,Y]_{\mathfrak h},Z].
\]
Any $G$-invariant tensor is parallel.
In particular,
\[
\nabla T=0,
\qquad
\nabla R=0.
\]
If $g$ is a $G$-invariant metric with
\[
g([X,Y]_{\mathfrak m},Z)
+g([X,Z]_{\mathfrak m},Y)=0,
\]
then
\[
(M,g,\tau=\tfrac12T)
\]
is a PST.
\item Cartan-Schouten connections.
Fix a compact Lie group $G$
with a bi-invariant metric $g$.
Set
\[
\sigma(X,Y,Z)=g([X,Y],Z).
\]
Then $\sigma\in\Omega^3(G)$.
It is both left- and right-invariant.
Now define
\[
\nabla^t=\nabla^g+\tfrac t2\,\sigma.
\]
For $t=\pm1$, the connection
$\nabla^t$ is flat.
Moreover, $\nabla^{-1}$
parallelizes the left-invariant vector
fields, and $\nabla^{+1}$
parallelizes the right-invariant
vector fields.
\end{enumerate}
\end{example}

\medskip\noindent
{\bf Properties of PST.}

Let $(M,g,\tau)$ be a PST. Then:

\begin{itemize}
\item
$\operatorname{Hol}(\nabla^\tau)
\subseteq
\operatorname{Stab}_{O(n)}(\tau)$.
\item
$R^\tau$ is pair-symmetric, so
\[
R^\tau\in \operatorname{Sym}^2
\Lambda^2T^*M.
\]
Hence, if
\[
\mathfrak{hol}(\nabla^\tau)
\subseteq
\mathfrak h,
\]
then $\nabla^\tau$ is an
$\mathfrak h$-instanton.
Then
\[
\Lambda^2
\cong
\mathfrak h\oplus\mathfrak h^\perp,
\qquad
R^\tau(\mathfrak h^\perp)=0.
\]
That is, $\nabla^\tau$ is an instanton.
\item
Cleiton--Moriano--Semmelman, 2021.
Fix
\[
\mathfrak{hol}(\nabla^\tau)
\subseteq
\mathfrak h
\subseteq
\underline{\operatorname{Stab}}(\tau)
\subseteq
\underline{\mathfrak{so}}(n).
\]
Then $(M,g,\tau)$ admits a foliation
with totally geodesic leaves.
The local leaf space is a PST
\[
(N,g_N,\sigma)
\cong
\prod_i (N_i,g_i,\sigma_i),
\]
where
$\underline{\operatorname{Stab}}(\sigma_i)$
acts irreducibly on $T N_i$.
Moreover,
\[
\tau=\tau_{\mathrm{leaf}}
+\pi^*\sigma-A,
\]
where $A$ is the O'Neill tensor.
\end{itemize}

\medskip\noindent
{\bf Closed torsion.}

What if we impose
that the torsion is closed,
i.e.\ $d\tau=0$?

\begin{lemma}
\label{hodge-birational-atoms-2026-lemma-let-m}
Let $(M,g,\tau)$ be a PST.
Then the following are equivalent.
\begin{itemize}
\item $d\tau=0$.
\item $\nabla^g\tau=0$.
\item $R^\tau$ satisfies the first
Bianchi identity.
\item For all $X\in TM$,
$(\tau_X)\diamond\tau=0$.
\item For all $p\in M$,
$(T_pM,g_p,\tau_p)$ is a
metric Lie algebra.
\end{itemize}
\end{lemma}

\begin{definition}
\label{hodge-birational-atoms-2026-definition-pst-geometries}
PST geometries with closed torsion are
called {\it PSCT geometries}.
\end{definition}

Barbaro--Pediconi--Tardini, 2024,
study Hermitian PSCT geometry.

\begin{theorem}
\label{hodge-birational-atoms-2026-theorem-any-psct}
Any PSCT geometry locally splits.
More precisely,
\[
(M,g,\tau)
\cong
\prod_i (M_i,g_i,\tau_i)
\times
(\mathbb{R}^k,g_{\mathrm{Euc}},0)
\times
\prod_j (N_j,h_j,\eta_j),
\]
where each $(M_i,g_i,\tau_i)$ is
irreducible under
$\underline{\operatorname{Stab}}(\tau_i)$
and $\tau_i\neq 0$,
and each $(N_j,h_j,\eta_j)$ is
irreducible and non-flat.
\end{theorem}

\medskip\noindent
{\bf Parallelized geometries.}

Which $G$-structures are parallelized
by a PSCT connection?
A natural list is the Berger list:
\[
U(m),\ SU(m),\ \operatorname{Spin}(m),
\ \operatorname{Spin}(m)\times
\operatorname{Spin}(1),
\ G_2,\ \operatorname{Spin}(7).
\]

\medskip\noindent
{\bf Hermitian torsion classes.}

For a $U(m)$-structure,
the intrinsic torsion lies in
\[
TM\otimes \Lambda^{(2,0)+(0,2)}.
\]
The Gray--Hervella classes are:
\[
W_1=
\Lambda^{3,0}\oplus\Lambda^{0,3},
\qquad
W_3=
\Lambda^{2,1}_0\oplus\Lambda^{1,2}_0,
\qquad
W_4=TM\wedge\omega.
\]
The class $W_2$ is the remaining
Gray--Hervella summand.
In this notation:
\begin{itemize}
\item $W_3\oplus W_4$ iff integrable.
\item $W_1\oplus W_3\oplus W_4$
iff skew torsion.
\item $W_1\oplus W_2\oplus W_3$
iff semi-K\"ahler.
\item $W_1$ iff nearly K\"ahler.
\end{itemize}

\medskip\noindent
{\bf General PSCT manifolds.}

Now we explain how a general PSCT
manifold looks.

\begin{theorem}
\label{hodge-birational-atoms-2026-theorem-let-m}
Let $(M,g,\tau,J)$ be a Hermitian PSCT
manifold.
Then locally
\[
(M,g,\tau)\cong
\underbrace{(H,g_H,-\tfrac12\sigma)}
_{\text{Lie group with }\nabla^{-1}}
\times
\prod_i
\underbrace{(S_i,g_i,\alpha_i
\operatorname{vol}_{g_i})}
_{\text{3-dimensional}}
\times
\prod_j
\underbrace{(N_j,g_j,0)}
_{\text{K\"ahler}}.
\]
Each $S_i$ is an $\alpha_i$-Sasaki
manifold, with Reeb vector field
\[
\xi_i\in\mathfrak X(S_i).
\]
This splitting is not preserved by $J$.
Instead, $J$ preserves
\[
\xi_i^\perp\subset TS_i,
\qquad
\mathfrak h\oplus
\operatorname{span}
\{\xi_1,\dots,\xi_r\},
\qquad
TN.
\]
\end{theorem}

\begin{theorem}[Barbaro--Pediconi--Tardini,
\label{hodge-birational-atoms-2026-theorem-2024-m}
2024]
If $(M,g,\tau,J)$ is Hermitian PSCT,
then there exists a Cartan subalgebra
\[
\underline{\mathfrak t}\subseteq
\mathfrak h
\]
such that $J$ preserves
\[
\underline{\mathfrak t}\oplus
\operatorname{span}
\{\xi_1,\dots,\xi_r\},
\]
and it projects to an integrable
complex structure on $H/T$.
\end{theorem}

Now assume $\dim M>4$.

\begin{theorem}
\label{hodge-birational-atoms-2026-theorem-m-g2}
If $(M,g,\tau,J)$ is a PSCT structure
in $W_1\oplus W_4$ or in $W_3$, then
it is K\"ahler.
\end{theorem}

\begin{theorem}
\label{hodge-birational-atoms-2026-theorem-m-g2-duplicate-2}
If $(M,g,\tau,J)$ is a PSCT structure
in $W_1\oplus W_3$, then
\[
(M,g,\tau)\cong
(H,g_H,-\tfrac12\sigma)
\times
\prod_i (N_i,h_i,0),
\]
and
\[
\underline{\mathfrak h}
\neq
\mathfrak{su}(2)\oplus\mathbb R.
\]
\end{theorem}

These manifolds are classified.
The $J$ are not classified.
One can have different almost
complex structures, but at least
we know what the manifolds are.

\section{Flag varieties as GKM
spaces}
\label{hodge-birational-atoms-2026-section-flag-varieties-as-gkm-spaces}

\medskip\noindent
{\bf Speaker.}

Giovane Galindo

\medskip\noindent
{\bf Title.}

Flag varieties as GKM spaces

\medskip\noindent
{\bf Abstract.}

In this talk, we present flag varieties
as GKM spaces and use their Lie-theoretic
structure to describe the associated
combinatorial data explicitly.
This perspective enables computations in
equivariant cohomology, including the
equivariant Gromov--Witten invariants.
This will facilitate simpler atom
computations and enable the solution of
problems in $G$-equivariant birational
geometry.
This is ongoing joint work with
Ludmil Katzarkov.

\medskip\noindent
For GKM varieties, we can compute GW
invariants in the equivariant
cohomology ring via a torus action and
the graph associated to the GKM
variety.

\medskip\noindent
{\bf Big quantum cohomology.}

\begin{definition}
\label{hodge-birational-atoms-2026-definition-smooth-projective}
A smooth projective variety $X$ is
given.
A {\it genus zero stable map} is a map
\[
f:(C;x_1,\dots,x_n)\to X
\]
where:
\begin{enumerate}
\item $C$ is a projective connected
nodal curve of genus $0$;
\item $x_1,\dots,x_n\in C$ are smooth
marked points;
\item $\operatorname{Aut}(f)$
is finite.
\end{enumerate}
\end{definition}

The moduli space of stable maps is
\[
\overline{\mathcal M}_{0,n}(X,\beta).
\]
It is a Deligne--Mumford stack.

\begin{definition}
\label{hodge-birational-atoms-2026-definition-evaluation-map}
The {\it evaluation map} is
\[
\operatorname{ev}_i:
\overline{\mathcal M}_{0,n}(X,\beta)
\to X
\]
given by
\[
\operatorname{ev}_i(f)=f(x_i).
\]
\end{definition}

\begin{definition}
\label{hodge-birational-atoms-2026-definition-gromov-witten}
The {\it Gromov--Witten invariant} is
\[
\left\langle{}
\phi_1,\dots,\phi_n
\right\rangle{}^X_{0,n,\beta}
=
\int_{
[\overline{\mathcal M}_{0,n}(X,\beta)]^{
\mathrm{vir}}
}
\prod_{i=1}^n \operatorname{ev}_i^*(\phi_i).
\]
\noindent where the $\phi_i$ are
cohomology classes on $X$.
\end{definition}

\begin{definition}
\label{hodge-birational-atoms-2026-definition-graded-basis}
For a graded basis
$(\alpha_k)$ of $H^*(X)$, put
\[
\tau=\sum_k T_k\alpha_k.
\]
The {\it big quantum product} is
\[
\phi_i\star_\tau\phi_j
=
\sum_{k,n,\beta}
\left\langle{}
\phi_i,\phi_j,\phi_k,\tau^n
\right\rangle{}^X_{0,3+n,\beta}
\frac{\phi^k}{n!}Q^\beta.
\]
\end{definition}

These are formal power series.

\medskip\noindent
{\bf GKM spaces.}

Flag varieties carry a torus action with
enough fixed-point and one-dimensional
orbit data to make GKM localization
effective.

\medskip\noindent
{\bf Flag varieties.}

\begin{definition}
\label{hodge-birational-atoms-2026-definition-let-g}
Let $G$ be a connected complex reductive
group and $B\subset G$ a Borel subgroup.
The {\it flag variety} $G/B$
parametrizes
complete flags in the standard
representation.
\end{definition}

The $T$-action on $G/B$ is induced by a
maximal torus $T\subset B$.
Its $T$-fixed points are indexed by the
Weyl group $W$.

\begin{remark}
\label{hodge-birational-atoms-2026-remark-flag-variety}
For a flag variety, the $T$-fixed points
are isolated and the one-dimensional
$T$-orbits are controlled by the roots of
$G$.
This is the GKM situation.
\end{remark}

The point of flag varieties is that we
can study them using the generalized
eigenspace decomposition.
In particular, we can study their
GW invariants this way.

\medskip\noindent
{\bf GKM data.}

\begin{definition}
\label{hodge-birational-atoms-2026-definition-t-space}
A $T$-space $X$ is a {\it GKM space}
if the fixed-point set $X^T$ is finite
and the one-dimensional orbits are
isolated in the sense required by GKM
theory.
\end{definition}

For $X=G/B$, the combinatorial model is
the moment graph:
the vertices are $W$ and the edges are
given by reflections.
Each edge is labeled by a root.

\begin{remark}
\label{hodge-birational-atoms-2026-remark-moment-graph}
The moment graph encodes the Bruhat
order, the root system, and the
localization data needed for equivariant
cohomology computations.
\end{remark}

\medskip\noindent
{\bf Equivariant cohomology.}

\begin{definition}
\label{hodge-birational-atoms-2026-definition-t-equivariant}
The {\it $T$-equivariant cohomology}
of $X$ is
\[
H_T^*(X):=
H^*(ET\times_T X).
\]
\end{definition}

For GKM spaces, localization identifies
$H_T^*(X)$ with tuples of polynomial
functions on the fixed points satisfying
edge-compatibility relations.
For flag varieties, this gives a
combinatorial description of
Schubert classes.

\begin{proposition}
\label{hodge-birational-atoms-2026-proposition-schubert-basis}
The Schubert basis of $H_T^*(G/B)$ is
compatible with the Weyl group
parametrization of fixed points, and the
restriction maps to $X^T$ are determined
by the GKM graph.
\end{proposition}

\medskip\noindent
{\bf Quantum part.}

Let $X$ be a smooth projective variety.
We now switch to the small quantum
cohomology.

\begin{definition}
\label{hodge-birational-atoms-2026-definition-let-t}
Let $T_i,T_j\in H^*(X)$.
The {\it small quantum product} is
\[
T_i\star_{\mathrm{small}}T_j
=
\sum_{a,b}
\sum_{\beta\in \overline{\mathrm{NE}}
(X,\mathbb{Z})}
\langle T_i,T_j,T_a\rangle_{3,\beta}
q^\beta g^{ab}T_b.
\]
\end{definition}

The ring
\[
(QH^*(X),\star_{\mathrm{small}})
\]
is called the small quantum cohomology.

The point is that these are no longer
series but finite sums, so we can do
actual computations.

\medskip\noindent
{\bf Flag varieties.}

\begin{definition}
\label{hodge-birational-atoms-2026-definition-flag-variety}
A {\it flag variety} is a complex
homogeneous space
\[
G^{\mathbb{C}}/P,
\]
where $G^{\mathbb{C}}$ is a connected,
simply connected semisimple complex Lie
group and $P\subset G^{\mathbb{C}}$ is a
parabolic subgroup, that is, a Lie
subgroup containing a Borel subgroup.
\end{definition}

\begin{definition}
\label{hodge-birational-atoms-2026-definition-let-g2}
Let $G/B$ be a flag variety.
The {\it Weyl group} $W$ is the finite
reflection group generated by the simple
reflections $S_\alpha$ associated to
the simple roots $\alpha\in \Pi$.
It carries the length function
\[
\ell:W\to \mathbb{Z}_{\geq 0},
\]
where for $w\in W$ the value $\ell(w)$
is the length of any reduced expression
of $w$ as a product of simple
reflections.
\end{definition}

\begin{example}
\label{hodge-birational-atoms-2026-example-let-3}
Let
\[
\mathrm{Fl}_3=SL_3/B
\]
be the space of full flags in
\(\mathbb{C}^3\):
\[
\{0\subset V_1\subset V_2\subset
\mathbb{C}^3\}.
\]
It can also be described as the
projectivized tangent bundle of the
projective plane:
\[
\mathrm{Fl}_3\cong \mathbb{P}(T_{\mathbb{
P}^2})
\cong \mathbb{P}^1\times \mathbb{P}^2.
\]
The point is to use the root space
decomposition to encode the geometry.
The flag variety \(\mathrm{Fl}_3\)
embeds into
\[
\mathbb{P}^2\times (\mathbb{P}^2)^*
\]
as a smooth incidence hypersurface of
bidegree \((1,1)\).
It is cut out by the incidence relation
\[
x_0y_0+x_1y_1+x_2y_2=0.
\]
The set of positive roots is
\[
R^+=\{\alpha,\beta,\alpha+\beta\}.
\]
\[
W=\langle s_\alpha,s_\beta\rangle
\cong S_3.
\]
\end{example}

\medskip\noindent
{\bf Schubert basis.}

Given an element $w\in W$, define the
Schubert cell of $w$ in the flag variety
$G/B$ by
\[
C_w=\omega_0 B\omega_0^{-1}\cdot
wB/B.
\]
The Schubert variety is
\[
X_w=\overline{C_w}.
\]
The Schubert class
\[
\sigma_w=[X_w]\in H^{2\ell(w)}(X).
\]
The set
\[
\{\sigma_w\}_{w\in W}
\]
is a homogeneous basis for $H^*(X)$.
In particular,
\[
H^*(X)=\sum_{w\in W}\mathbb{Q}\sigma_w.
\]

\begin{example}
\label{hodge-birational-atoms-2026-example-2-have}
For \(\mathrm{Fl}_2\), we have the
following basis for our cohomology:
\[
\sigma_{\mathrm{id}},
\sigma_{s_\alpha},
\sigma_{s_\beta},
\sigma_{s_\beta s_\alpha}
\in H^2,
\]
\[
\sigma_{s_\alpha s_\beta}
\in H^3,
\qquad
\sigma_{s_\alpha s_\beta s_\alpha}
\in H^6.
\]
\end{example}

A nonzero linear functional
\[
\alpha\in\mathfrak h^*
\]
is called a root of $\mathfrak g$
with respect to $\mathfrak h$ if the
corresponding root space
\[
\mathfrak g_\alpha:=
\{X\in\mathfrak g:
[H,X]=\alpha(H)X
\text{ for all }H\in\mathfrak h\}
\]
is nonzero.

The Lie algebra $\mathfrak g$
decomposes as a direct sum of weight
spaces for the adjoint action of
\(\mathfrak h\):
\[
\mathfrak g
=
\mathfrak h
\oplus
\bigoplus_{\alpha\in R}
\mathfrak g_\alpha.
\]

\begin{remark}
\label{hodge-birational-atoms-2026-remark-root-data}
The root data of $G^{\mathbb{C}}/P$
is what makes the GKM description
effective.
\end{remark}

\begin{definition}
\label{hodge-birational-atoms-2026-definition-let-t2}
Let
\[
T\cong (\mathbb{C}^*)^r
\]
be an algebraic torus.
A smooth projective variety $X$ with
an effective $\mathbb{T}$-action is
called {\it GKM}, after Goresky,
Kottwitz, and
MacPherson, if the following hold:
\begin{enumerate}
\item $X^{\mathbb{T}}$ is the fixed
point set, and it is finite.
\item The one-dimensional orbits are
isolated.
\end{enumerate}
\end{definition}

\begin{definition}
\label{hodge-birational-atoms-2026-definition-let-x}
Let $X$ be a GKM variety.
The {\it GKM graph} of $X$ is the graph
\[
\Gamma=(V(\Gamma),E(\Gamma))
\]
where:
\begin{enumerate}
\item the vertices are the
$\mathbb{T}$-fixed points;
\item the edges are the closures of the
one-dimensional orbits.
\end{enumerate}
\end{definition}

\begin{remark}
\label{hodge-birational-atoms-2026-remark-flag-variety2}
For a flag variety $X=G/B$ there is a
natural action of a torus $T\subset G$.
The fixed points are the points
$wB/B$ for $w\in W$.

Given $w\in W$ and a root
$\alpha\in R$, the subgroup
$G_\alpha\cong SL_2$ defined by
$\alpha$ gives an invariant
$\mathbb{P}^1$:
\[
wG_\alpha B/B.
\]
It connects $wB/B$ and
$ws_\alpha B/B$.
\end{remark}

\begin{remark}
\label{hodge-birational-atoms-2026-remark-gkm-graph}
The GKM graph of our running example
has six vertices.
\end{remark}

\begin{definition}
\label{hodge-birational-atoms-2026-definition-x-gkm}
If $X$ is a GKM space with a
$T$-action, there is an induced action
on
\[
\overline{\mathcal M}_{g,n}(X,\beta).
\]
Given $t\in T$, define
\[
t(C,p_1,\dots,p_n,f)
=
(C,p_1,\dots,p_n,t\circ f).
\]

We denote by
$\overline{\mathcal M}_{g,n}(X,\beta)^T$
the {\it fixed point set} of this
action.
\end{definition}

\begin{definition}
\label{hodge-birational-atoms-2026-definition-n-0}
For $n\ge 0$, an {\it $n$-marked tree}
$\gamma$ consists of:
\begin{enumerate}
\item a finite nonempty set
$V(\gamma)$, the vertices;
\item a subset
$E(\gamma)\subset \binom{V(\gamma)}{2}$
such that the graph
$(V(\gamma),E(\gamma))$ is a tree;
\item a map
$\psi:\{1,\dots,n\}\to V(\gamma)$;
\item a map assigning an effective
curve class to each vertex;
\item a map assigning a positive
degree to each edge.
\end{enumerate}
\end{definition}

\begin{lemma}
\label{hodge-birational-atoms-2026-lemma-marked-trees}
The marked trees that occur are in
bijection with the fixed points of
$\overline{\mathcal M}_{0,n}(X,\beta)$
under the $\mathbb{T}$-action.
\end{lemma}

\begin{theorem}
\label{hodge-birational-atoms-2026-theorem-let-x5}
Let $X$ be a smooth projective GKM
space, and fix a splitting
\[
s:H^*(X,\mathbb{Q})
\to
H^*_{\mathbb{T}}(X,\mathbb{Q}).
\]
Then the equivariant Gromov--Witten
invariants are given by the
localization formula, as a sum over
the fixed marked trees.
\end{theorem}

\begin{remark}
\label{hodge-birational-atoms-2026-remark-fl-3}
For $Fl_3$, the equivariant Gromov--
Witten invariants
\[
\langle \gamma_i,\gamma_j\rangle_{2,\beta}
\]
are encoded by a matrix.
The point is that the localization
theorem lets us compute these
invariants.

Once GW invariants are computable, the
quantum cohomology is computable as
well.
For $Fl_3$, this gives the matrix for
small quantum multiplication.
Hence we get six Hodge atoms of
dimension $1$ for $Fl_3$.
\end{remark}

\medskip\noindent
{\bf Detour.}

The following map:
\[
\mathrm{Fl}_3\to \mathrm{Fl}_3,
\qquad
(V_1,V_2)\mapsto
(V_2^\perp,V_1^\perp)
\]
defines a $\mathbb{Z}_2$-action on
\(\mathrm{Fl}_3\).
It permutes the generators of the
Picard group.
This gives the following Hodge diamond:

\[
\begin{matrix}
1
\end{matrix}
\]

To compute the $G$-Hodge atoms for this
variety, we now look for the invariant
Hodge class that gives us the maximum
amount of eigenvalues.
So now it is invariant classes, in
contrast with the usual atoms.

The $\mathbb{Z}_2$ atoms for
\(\mathrm{Fl}_3\) are:

\[
\begin{tabular}{c|c}
$\rho$ & $\rho_{\mathbb{Z}_2}$ \\
\hline
$\lambda_1$ & $2,1$ \\
$\lambda_2$ & $2,2$ \\
$\lambda_3$ & $1,1$ \\
$\lambda_3$ & $1,1$
\end{tabular}
\]

The idea is to compute the matrix
multiplication in quantum cohomology
using the ambient space.

\begin{theorem}
\label{hodge-birational-atoms-2026-theorem-let-x5-duplicate-2}
Let $X$ be a smooth projective GKM
space, and let
$E\to X$ be a $\mathbb{T}$-equivariant
vector bundle generated by global
sections.
Let $Y\subset X$ be a generic zero
locus.
Then the Gromov--Witten invariants of
$Y$ are computed by a localization
formula.
This is the same kind of computation as
before, but now for complete
intersections.

In particular, one can compute the
matrix for quantum multiplication.
\end{theorem}

\begin{example}
\label{hodge-birational-atoms-2026-example-construct-del}
We construct a del Pezzo surface $S$
directly from the canonical geometry of
the flag variety $\mathrm{Fl}_3$.
We can compute the matrix
multiplication.
\end{example}

\medskip\noindent
{\bf Takeaway.}

The Lie-theoretic structure of a flag
variety turns the geometry into
combinatorics, and GKM methods make
equivariant calculations explicit.

\section{Geometric aspects of
conditional measures}
\label{hodge-birational-atoms-2026-label-hodge-birational}

\medskip\noindent
{\bf Speaker.}

Renata Possobon

\medskip\noindent
{\bf Title.}

Geometric aspects of
conditional measures

\medskip\noindent
{\bf Abstract.}

From a logistical problem
formulated in the 18th century
emerged what we now know as
Optimal Transport Theory, whose
development has been particularly
intense in recent decades.
By providing a natural geometric
framework for comparing measures,
this theory establishes deep
connections with several areas,
including Geometry, Probability,
Analysis, Dynamical Systems,
Economics, and Data Science.
In this talk, we will discuss the
concept of disintegration of
measures from the perspective of
optimal transport problems, aiming
to investigate the structure and
geometric arrangement of the
associated conditional measures.
In particular, we will analyze
rigidity conditions for
disintegration maps, as well as
propose an approach to characterize
foliations induced by
disintegration.
This talk is based on joint work
with Florentin M\"unch and Christian
Rodrigues.

\medskip\noindent
{\bf Optimal transport theory.}

Let $(X,d)$ be a metric space.
We study the space
$\mathcal{P}(X)$
of Borel probability measures on $X$.

Optimal transport began with a
logistical problem posed by
Caspard Monge in 1781.

The main idea is to transport mass
from one point to another at
minimal cost.
Think of soil in a pile in the garden
being moved to a hole made somewhere
else.

The soil is represented by a
distribution
$\mu\in\mathcal{P}(X)$,
and the hole, or destination,
by another distribution
$\nu\in\mathcal{P}(X)$.

Let $X$ and $Y$ be Radon spaces,
and let
$\mu\in\mathcal{P}(X)$,
$\nu\in\mathcal{P}(Y)$.
Fix a Borel cost function
$c:X\times Y\to[0,\infty]$.
We minimize
\[
\gamma\mapsto
\int_{X\times Y}
c(x,y)\,d\gamma(x,y)
\]
among all transport plans with
marginals $\mu$ and $\nu$.

\section{Rational Cubic Fourfolds and
Their Relation to K3 Surfaces}
\label{hodge-birational-atoms-2026-label-hodge-birational2}

\medskip\noindent
{\bf Speaker.}

J\'er\'emy Gu\'er\'e

\medskip\noindent
{\bf Title.}

Rational Cubic Fourfolds and Their
Relation to K3 Surfaces

\medskip\noindent
{\bf Abstract.}

I will briefly review the construction
of atoms, specifically addressing
certain technical difficulties
regarding the behavior of Hodge
structures under Iritani's blow-up
formula. Subsequently, I will
introduce a new atomic invariant and
provide a proof for the following
theorem: if a smooth complex cubic
fourfold is rational, then its
primitive cohomology is isomorphic -
as a rational Hodge structure - to the
shifted middle cohomology of a
projective K3 surface.

\medskip\noindent
{\bf AI-generated summary.}

The aim is to build a stronger
birational obstruction for cubic
fourfolds by mixing two kinds of data.
This is the first dimension in which
the rationality answer for cubics is
not uniform: some are rational, and
others are not.

First, classical Hodge theory says that
blow-ups preserve the Hochschild degree
$p-q$.
Thus, if a cubic fourfold were
rational, then some blow-up center in
any weak factorization to
$\mathbb{P}^4$ would have to contribute
the missing $H^{(2)}$ piece.
This already forces the center to be
surface-like, and eventually K3-like.

Second, quantum cohomology gives the
theory of atoms.
Hodge theory alone gets one only to
something like a surface with nef
canonical class.
To force the surface to be K3, one
uses Iritani's quantum blow-up formula.
The atoms are generalized eigenspaces
of quantum multiplication by the Euler
vector field.

The technical work is to make these
atoms compatible across different
blow-up factorizations.
This is why one introduces evaluation
maps into a common non-archimedean
power-series field $S$:
otherwise the completed coefficient
rings depend on the chosen blow-up
map.

The payoff is the following strategy.
If $X$ is rational, then weak
factorization together with Iritani's
theorem forces some blow-up center
$Z$ to fail a chosen property $(*)$.
One then chooses $(*)$ so that the
possible failing centers are
geometrically meaningful.
In KKPY, showing that no such center
exists gives irrationality for very
general cubic fourfolds.
Here the property is refined so that
rationality of a cubic fourfold forces
the primitive rational Hodge structure
of $X$ to come from a projective K3
surface.

In slogan form:
\[
\text{rational cubic fourfold}
\quad\Longrightarrow\quad
\text{K3 center in weak factorization}
\]
\[
\Longrightarrow
\quad
H^4(X,\mathbb{Q})_{\mathrm{prim}}
\text{ is of K3 type.}
\]

\medskip\noindent
{\bf Introduction.}

Let
\[
X_3\subset \mathbb{P}^5
\]
be a cubic fourfold.
All varieties in this talk are
smooth, projective, and over
$\mathbb{C}$.

\begin{theorem}[Weak factorization]
\label{hodge-birational-atoms-2026-theorem-let-x6}
Let $X$ and $Y$ be two birational
varieties.
Then there exists a weak
factorization
\[
X=X_0\dashrightarrow X_1
\dashrightarrow \cdots
\dashrightarrow X_{n-1}
\dashrightarrow X_n=Y,
\]
where each dashed arrow is either a
morphism or its inverse, and each
morphism is the blow-up along a smooth
center.
\end{theorem}

\begin{remark}
\label{hodge-birational-atoms-2026-remark-reference-weak}
The reference is the weak
factorization theorem of
Abramovich--Karu--Matsuki--Wlodarczyk,
also proved independently by
Wlodarczyk.
In these notes, the related blow-up
technology appeared earlier in the
discussion of Iritani's blow-up
formula.
\end{remark}

Now look at what happens for one
blow-up.

The blow-up formula says that
$D^b(Y)$ can be obtained from
$D^b(X)$ by adding or removing
semi-orthogonal pieces of the form
$D^b(Z)$, where $Z$ is the center of
the blow-up.

\begin{theorem}
\label{hodge-birational-atoms-2026-theorem-blow-up}
Let
\[
\widetilde X=\operatorname{Bl}_Z X,
\]
where $Z\hookrightarrow X$ has
codimension $r$.
Then
\[
H^{p,q}(\widetilde X)
\cong
H^{p,q}(X)
\oplus
\bigoplus_{k=1}^{r-1}
H^{p-k,q-k}(Z).
\]
\end{theorem}

\begin{remark}
\label{hodge-birational-atoms-2026-remark-not-only}
Not only is the bidegree $(p,q)$
controlled, but the difference
$p-q$ is preserved.
\end{remark}

\begin{definition}
\label{hodge-birational-atoms-2026-definition-define-h}
Define
\[
H^{(k)}(X)
=
\bigoplus_{p-q=k} H^{p,q}(X).
\]
We call $k$ the {\it Hochschild
degree}.
\end{definition}

\noindent
The point is that $H^{(n)}$ and $H^{(n-1)}$
are birational invariants, where $n=\dim X$:

\begin{example}
\label{hodge-birational-atoms-2026-example-let-n}
Let $n=\dim X=\dim Y$, and suppose
$X$ is birational to $Y$.
Then
\[
\dim H^{(n)}(X)=\dim H^{(n)}(Y).
\]
The same holds for $H^{(n-1)}$.
This follows directly from the
blow-up formula and the definition.

For example, let
\[
Q\subset \mathbb{P}^4
\]
be a quintic threefold.
Then
\[
H^{3,0}(Q)=\mathbb{C},
\qquad
H^{3,0}(\mathbb{P}^3)=0.
\]
Thus a quintic threefold is never
rational.
\end{example}

\begin{example}
\label{hodge-birational-atoms-2026-example-let-x}
Let
\[
X\subset \mathbb{P}^5
\]
be a smooth cubic fourfold.
Its Hodge diamond is
\[
\begin{array}{ccccccccc}
&&&&1&&&&\\
&&&0&&0&&&\\
&&0&&1&&0&&\\
&0&&0&&0&&0&\\
0&&1&&21&&1&&0\\
&0&&0&&0&&0&\\
&&0&&1&&0&&\\
&&&0&&0&&&\\
&&&&1&&&&
\end{array}
\]
whereas the Hodge diamond of
$\mathbb{P}^4$ is
\[
\begin{array}{ccccccccc}
&&&&1&&&&\\
&&&0&&0&&&\\
&&0&&1&&0&&\\
&0&&0&&0&&0&\\
0&&0&&1&&0&&0\\
&0&&0&&0&&0&\\
&&0&&1&&0&&\\
&&&0&&0&&&\\
&&&&1&&&&
\end{array}
\]
For the cubic, the first nonzero
off-center Hochschild degree is
$H^{(2)}$, while for $\mathbb{P}^4$
all off-center Hochschild degrees
vanish.
Thus direct comparison with
$\mathbb{P}^4$ forces nontrivial
centers in any factorization.

The point is that this invariant is obtained
by comparing the leftmost columns in the
Hodge diamond.
\end{example}

If
\[
X\dashrightarrow \mathbb{P}^4
\]
is birational, then in every weak
factorization there is a blow-up
center $Z$ such that
\[
H^{(2)}(Z)\neq 0.
\]
As a first observation, $Z$ must be a
surface.
From the classification of surfaces,
the minimal model $\Sigma$ of such a
surface has $K_\Sigma$ nef.

A variety with $K_X$ nef has only one
atom.
This shows that the nefness of $K_X$
is important for Gromov--Witten
theory.

\begin{remark}[Hassett's conjecture]
\label{hodge-birational-atoms-2026-remark-cubic-fourfold}
A cubic fourfold $X$ is rational if
and only if
\[
H^4(X,\mathbb{Z})_{\mathrm{prim}}
\overset{\sim}{\longrightarrow}
H^2(S,\mathbb{Z})(-1)
\]
as integral Hodge structures, i.e.
$\mathbb{Z}$-HS, for some
projective K3 surface $S$.
\end{remark}

\begin{remark}
\label{hodge-birational-atoms-2026-remark-point-hassett}
The point is that Hassett's conjecture
is an integral statement.
It concerns integral Hodge structures,
not only rational ones.
\end{remark}

\begin{theorem}
\label{hodge-birational-atoms-2026-theorem-direct-implication}
The direct implication of the
conjecture holds with rational Hodge
structures.
That is, if $X$ is rational, then
\[
H^4(X,\mathbb{Q})_{\mathrm{prim}}
\cong
H^2(S,\mathbb{Q})(-1)
\]
as rational Hodge structures, for some
projective K3 surface $S$.
\end{theorem}

\begin{lemma}[KKPY]
\label{hodge-birational-atoms-2026-lemma-very-general}
A very general cubic fourfold $X$ is
irrational.
\end{lemma}

The obstruction is carried by the
transcendental classes in
\[
H^4(X,\mathbb{Q})_{\mathrm{prim}}.
\]

\begin{remark}[On the proof of the theorem]
\label{hodge-birational-atoms-2026-remark-x-rational}
If $X$ is rational, then in every weak
factorization there exists a blow-up
center $Z$ which is a K3 surface.
In fact, to pass from $K_Z$ being nef
to $Z$ being a K3 surface, one needs a
little more than Hodge theory.
This leads to the next topic.
\end{remark}

\medskip\noindent
{\bf Atom construction.}

To pass from $K_\Sigma$ nef to
$\Sigma$ being a K3 surface, we use a
refined invariant composed with
$H^{(2)}$.

The blow-up formula is an isomorphism
of quantum rings.

Let us recall the basics of quantum
cohomology.

$$
\xymatrix{
\mathcal C \ar[r]^f & X
}
$$
\[
\overline{\mathcal M}_{g,n}(X,\beta)=
\left\{
f:
C\to X:
C \text{ curve of genus } g,
f_*([C])=\beta\in H_2(X),
x_1,\dots,x_n\in C
\right\}/\sim
\]
Here $\mathcal C$ is the universal
curve.

$$
\xymatrix{
\mathcal C \ar[r]^f \ar[d]& X \\
\overline{\mathcal M}_{g,n}(X,\beta)
\ar@/^/[u]^{x_1,\dots,x_n}
}
$$

There exists
\[
[\overline{\mathcal M}_{g,n}(X,\beta)]^{
\mathrm{vir}}
\]
an algebraic cycle of the expected
dimension.
The virtual dimension is
\[
\operatorname{vdim}_{\mathbb{C}}
=
(1-g)\bigl(\dim_{\mathbb{C}}X-3\bigr)
+n
+\int_\beta c_1(X).
\]

\begin{definition}
\label{hodge-birational-atoms-2026-definition-compute-gromov}
To compute the Gromov--Witten
invariants, let
\[
\alpha_1,\dots,\alpha_n\in H^\bullet(X).
\]
The {\it Gromov--Witten invariant} is
something like counting curves:
\[
\left\langle{}
\alpha_1,\dots,\alpha_n
\right\rangle{}^X_{g,n,\beta}
=
\int_{
[\overline{\mathcal M}_{g,n}(X,\beta)]^{
\mathrm{vir}}
}
\prod_i \operatorname{ev}_i^*
\bigl(\alpha_i\bigr)
\in \mathbb{Q}.
\]

Unless $g=0,\beta=0,n=2$, we have
the string equation
\[
\left\langle{}
\alpha_1,\dots,\alpha_n,\mathbb{1}_X
\right\rangle{}
=0.
\]
And the divisor equation, if
\[
\gamma\in H^1(X)\text{ or }H^2(X),
\]
then
\[
\left\langle{}
\alpha_1,\dots,\alpha_n,\gamma
\right\rangle{}
=
\int_\beta \gamma
\left\langle{}
\alpha_1,\dots,\alpha_n
\right\rangle{}.
\]
\end{definition}

\begin{definition}
\label{hodge-birational-atoms-2026-definition-let-i}
Let $(\phi_i)$ be a graded basis of
$H(X,\mathbb{Q})$.
Let $(\phi^j)$ be the dual basis, so
that
\[
\int_X \phi_i\vee \phi^j
=
\delta_{ij}.
\]
For every graded vector subspace
\[
V\subset H(X,\mathbb{Q}),
\]
choose a graded basis
$(\alpha_k)$ of $V$.
For example, one can take
\[
V=H^{\mathrm{even}},
\qquad
V=H^{(0)},
\qquad
V=H_{\mathrm{Hdg}}.
\]
Put a formal variable $T_k$ for each
basis element $\alpha_k$, and define
\[
\tau=\sum_k T_k\alpha_k.
\]
Then
\[
\phi_i\star_\tau\phi_j
=
\sum_{k,n,\beta}
\left\langle{}
\phi_i,\phi_j,\phi_k,\tau^n
\right\rangle{}^X_{0,3+n,\beta}
\frac{\phi^k}{n!}Q^\beta .
\]
\end{definition}

\noindent
Another invariant is the Euler vector field:

\begin{definition}
\label{hodge-birational-atoms-2026-definition-euler-vector}
The {\it Euler vector field} is
\[
\operatorname{Eu}_\tau
=
c_1(X)
+
\sum_k
\frac{2-\deg \alpha_k}{2}
T_k\alpha_k.
\]
\end{definition}

\noindent
One of the most important objects in theory
is:

\begin{definition}
\label{hodge-birational-atoms-2026-definition-statement2}
\[
\kappa_\tau(-)
:=
\operatorname{Eu}_\tau\star_\tau(-)
\in
\operatorname{End}
\bigl(H(X,\mathbb{Q})_R\bigr),
\]
where
\[
R=\mathbb{Q}[[Q,(T_k)]]
\]
is the base-change ring introduced
because of the formal variables above.
\end{definition}

\medskip\noindent
{\bf Properties.}

\begin{itemize}
\item
\[
\kappa_{\tau+\lambda\operatorname{id}_X}
=
\kappa_\tau+\lambda\operatorname{id}.
\]
(string)
\item $\kappa_\tau$ is a morphism of
$\mathbb{Q}$-Hodge structures.
\end{itemize}

\noindent
Iritani's formula is like the blow-up formula
we had before but upgraded to the quantum
level:

\begin{theorem}[Iritani]
\label{hodge-birational-atoms-2026-theorem-let-x7}
Let
\[
\widetilde X=\operatorname{Bl}_Z X.
\]
Then
\[
\kappa_{\widehat\tau}^{\widetilde X}
\]
is conjugate, as an endomorphism, to
\[
\kappa_{\tau_0}^X
\oplus
\bigoplus_{k=1}^{r-1}
\kappa_{\tau_k}^Z
\]
in
\[
\operatorname{End}_{\widehat R_{\mathbb{C}}}
\bigl(
H(\widetilde X,\mathbb{C})_
{\widehat R_{\mathbb{C}}}
\bigr).
\]
Here the parameters are related by a
change of variables
\[
\widehat\tau
\longleftrightarrow
(\tau_0,\tau_1,\ldots,\tau_{r-1}).
\]
\end{theorem}

\noindent
Here
\[
R_{\mathbb{C}}
=
\mathbb{C}[[Q,(T_k)]],
\qquad
\widehat R_{\mathbb{C}}
=
\mathbb{C}[q^{\pm 1}][[Q,(T_k)]],
\]
and
\[
R_{\mathbb{C}}
\hookrightarrow
\widehat R_{\mathbb{C}}.
\]

\begin{remark}
\label{hodge-birational-atoms-2026-remark-def-invariance}
$\mathbb{C}$-HS are deformation invariants,
while $\mathbb{Q}$-HS aren't.
\end{remark}

\noindent
These rings are written with respect
to $X$.
One could also ask for the
corresponding rings with respect to
$\widetilde X$, for example
\[
R_{\mathbb{C}}^{\widetilde X}
=
\mathbb{C}[[\widetilde Q,
(\widetilde T_k)]].
\]
There is a curve class
$\ell_E\subset \operatorname{Exc}$
which is a line in the exceptional
divisor.
We do not write its precise definition
here, but under the change of rings one
has
\[
\widetilde Q^{\ell_E}
\longmapsto
q^{2(r-1)}.
\]
There is also a map
\[
R_{\mathbb{C}}^Z
\longrightarrow
\widehat R_{\mathbb{C}},
\]
which need not be injective, but at
least gives a ring map.

\begin{remark}
\label{hodge-birational-atoms-2026-remark-iritani-s}
Iritani's formula comes from a
localization formula.
Do not trust the change of variables:
prove it.
\end{remark}

\begin{definition}
\label{hodge-birational-atoms-2026-definition-atom-generalized}
An {\it atom} is a generalized
eigenspace of
$\kappa_{\widetilde\tau}^{\widetilde X}$,
namely:
$$
\mathbb{E}_{\tau,\lambda}
=\Ker ((\kappa_\tau-\lambda)^m),
\qquad  \gg 1.
$$
\end{definition}

Up to this point everything is
deformation invariant, since we are in
Gromov--Witten theory.
If we want an invariant which is not
deformation invariant, we need to add
an object which is not deformation
invariant: Hodge structures.

\medskip\noindent
{\bf Refinements.}

Replace $\mathbb{C}$ by a field $K$.
Replace
\[
V=H(X,\mathbb{Q})
\]
by either
\[
H(X)_{\mathrm{Hdg}}
\qquad\text{or}\qquad
H^{(0)}(X).
\]
We now keep track of rational Hodge
structures.

The operator can be written as
\[
\kappa_\tau
=
\sum_i s_i K_i,
\]
where $(s_i)$ is a basis of
$K/\mathbb{Q}$, and the $K_i$ are
morphisms of rational Hodge
structures.

\medskip\noindent
{\bf Big issue.}

The ring $\widehat R$ inverts
\[
\widetilde Q^{\ell_E},
\]
and this depends on the morphism
\[
\widetilde X\to X.
\]
For example, in a diagram
\[
\xymatrix{
& \widetilde X \ar[dl]\ar[dr] &\\
X_1 && X_2
}
\]
the resulting rings
$\widehat R_1$ and $\widehat R_2$
need not be isomorphic.

\begin{remark}
\label{hodge-birational-atoms-2026-remark-difficulties-rings2}
One of the greatest difficulties of the
theory if to compare the rings
$\hat{R}_1$ and $\hat{R}_2$.
Here comes Levi-Civita …[missing]
\end{remark}

The solution is to look, for every
$R$ and every $X$, for a diagram
\[
\xymatrix{
R^X \ar[r]\ar@{^{(}->}[d] & S\\
\widehat R \ar@{-->}[ur] &
}
\]
where $S$ is a common ring independent
of $X$, of $\widetilde X$, and of the
chosen map $\widetilde X\to X$.

The ring $S$ must be:
\begin{itemize}
\item a ring of power series;
\item large enough to contain roots
and related elements needed by the
change of variables.
\end{itemize}

In sum, take
\[
S=\mathbb{Q}[a^{\mathbb{Q}}]
[b^{\pm 1}].
\]
The factor $b^{\pm 1}$ keeps track of
degree.
We will ignore this point for now.

Here
\[
\mathbb{Q}[a^{\mathbb{Q}}]
\ni
\sum_{n\in\mathbb{Q}} x_n a^n,
\qquad x_n\in\mathbb{Q},
\]
with the condition that, for every
$m\in\mathbb{Q}$,
\[
\#\{n:x_n\neq 0,\ n<m\}<+\infty.
\]
Thus these are generalized power
series with rational exponents and
well-ordered support toward
$-\infty$.

This is a field.
On this field one defines a
non-archimedean absolute value.
For
\[
0\neq f=\sum_{n\in\mathbb{Q}}x_n a^n,
\]
put
\[
\rho(f)=\min\{n:x_n\neq 0\}
\]
and define
\[
|f|:=2^{-\rho(f)}.
\]

\medskip\noindent
{\bf Producing the ring map.}

How do we produce the ring map
\[
R^X\to S?
\]
Since
\[
R^X=\mathbb{Q}[[Q,T_k]],
\]
one is tempted to define the map on
the generators $Q$ and $T_k$ and then
extend by linearity.
This does not always work:
there are convergence issues.
The absolute value above is useful
precisely because it gives a criterion
for when such a map on generators can
be extended.

Now let
\[
\operatorname{ev}:\widehat R_K\to S_K
\]
be an evaluation map.
From
\[
\operatorname{ev}(\kappa_\tau)
\]
we compute the generalized
eigenspaces
\[
\mathbb{E}_{\operatorname{ev},\lambda}
=
\ker
\bigl(
\operatorname{ev}(\kappa_\tau)-\lambda
\bigr)^m
\]
for $m\gg 1$.
These are the atoms.

For almost all evaluations
$\operatorname{ev}$ and for all
$\lambda$, if a property $(*)$ holds
on
$\mathbb{E}_{\operatorname{ev},\lambda}$,
then we say that $X$ satisfies $(*)$.
Thus the point is to cook up a
property $(*)$ of vector spaces.

\begin{remark}
\label{hodge-birational-atoms-2026-remark-how-to}
The upshot is that eigenspaces
are conjugation-invariant, so that
we can apply Iritani's formula.
\end{remark}

\begin{proposition}
\label{hodge-birational-atoms-2026-proposition-z-satisfies}
If $Z$ satisfies $(*)$, then $X$
satisfies $(*)$ if and only if
\[
\widetilde X=\operatorname{Bl}_Z X
\]
satisfies $(*)$.
\end{proposition}

\begin{proof}
(Uses the string equation.)
$Z$ satisfies $(*)$, $X$ not,
exists $\tau_0$, exists $\lambda$,
$\mathbb{E}_{\tau_0,\lambda}$ not $(*)$.
Choose $\tau_k$ such that 
$\mathbb{E}_{\tau_k}=0$
$\xymatrix{\ar@{~>}[r]&}$,
$\mathbb{E}_{\tilde{\tau},\lambda}$ not
$(*)$.
\end{proof}

\begin{lemma}[Corollary]
\label{hodge-birational-atoms-2026-lemma-iritani-corollary2}
Assume that
\[
\begin{cases}
X\text{ is birational to }Y,\\
X\text{ fails }(*),\\
Y\text{ satisfies }(*).
\end{cases}
\]
Then there exists a blow-up center
$Z$ which fails $(*)$.
\end{lemma}

Define
\[
\rho_{\operatorname{ev},\lambda}
=
\dim
\bigl(
\mathbb{E}_{\operatorname{ev},\lambda}
\cap H_{\mathrm{Hdg}}
\bigr),
\]
\[
\nu^{(k)}_{\operatorname{ev},\lambda}
=
\dim
\bigl(
\mathbb{E}_{\operatorname{ev},\lambda}
\cap H^{(k)}
\bigr),
\]
and
\[
\gamma_{\operatorname{ev},\lambda}
=
\operatorname{rk}
\left(
\operatorname{ev}(\kappa_\tau)-\lambda
\right)_
{\mathbb{E}_{\operatorname{ev},\lambda}}.
\]
Here
$\mathbb{E}_{\operatorname{ev},\lambda}$
is the generalized tangent space.

The property from the KKPY paper is
\[
(*)_{\mathrm{KKPY}}
\quad
\Longleftrightarrow
\quad
\nu^{(2)}=0
\ \text{or}\
\rho\geq 3.
\]
i.e., that $H^{(2)}=0$ or
at least three algebraic cycles,
because of
point, $H$ and fundamental class.

The property from the present paper is
\[
(*)_{\mathrm{new}}
\quad
\Longleftrightarrow
\quad
\nu^{(2)}=0
\ \text{or}\
\nu^{(1)}\neq 0
\ \text{or}\
\gamma\geq 2,
\]
which is a modification of the former,
such that the only surfaces that satisfy it
are K3.

Thus, by the theorem, one gets a
blow-up center $Z$ failing $(*)$.
The question becomes:
which points, curves, and surfaces can
fail this property?

For KKPY one shows that there is no
center $Z$ failing $(*)$.
This gives the contradiction used to
prove the desired non-birationality
statement.

\section{Birational Geometry of
Hyperk\"ahler Manifolds and the
Hu--Yau Conjecture}
\label{hodge-birational-atoms-2026-label-hodge-birational3}

\medskip\noindent
{\bf Speaker.}

Andrey Soldatenkov

\medskip\noindent
{\bf Title.}

Birational geometry of hyperk\"ahler
manifolds and the Hu--Yau conjecture

\medskip\noindent
{\bf Abstract.}

The talk will be based on a joint work
with Katya Amerik and Misha Verbitsky.
I will try to explain our proof of the
Hu-Yau conjecture for compact
hyperkähler manifolds. The conjecture
predicts a description ``in codimension
two'' of bimeromorphic maps between
holomorphic symplectic manifolds. It
states that any such map should be a
composition of certain elementary
transformations called Mukai flops, at
least if we ignore some ``bad locus''
that can be of codimension three or
higher. The conjecture has been proven
by Wierzba and Wiśniewski in the case
of symplectic fourfolds. Building on
their proof, and using more recent
results on the structure of the Kähler
cone of hyperkähler manifolds, we
extend the result to arbitrary
dimension.

\begin{definition}
\label{hodge-birational-atoms-2026-definition-hyperk-ahler}
A {\it hyperk\"ahler manifold} is a
compact
smooth manifold $M$ with a metric $g$
such that
\[
\pi_1(M)=1
\]
and
\[
\operatorname{Hol}(g)
=
\operatorname{Sp}(n).
\]
\end{definition}

\begin{remark}
\label{hodge-birational-atoms-2026-remark-such-m}
Such an $M$ admits a triple of complex
structures $I,J,K$ such that
\[
IJ=-JI=K
\]
and such that $g$ is K\"ahler for all
of them.
We have the K\"ahler forms
$\omega_I$, $\omega_J$, and
$\omega_K$.
\end{remark}

Consider the $2$-form
\[
\sigma_I=\omega_J+\sqrt{-1}\,\omega_K.
\]
Then $\sigma_I$ is a holomorphic
symplectic form of type $(2,0)$.
In fact, it is nondegenerate.

Let $\Omega^k_{M_I}$ be the bundle of
$I$-holomorphic $k$-forms on $M$.
Then $\sigma_I^n$ is a nowhere
vanishing section of the canonical
bundle
\[
K_{M_I}=\Omega^{2n}_{M_I}.
\]
Thus $M_I$ is an IHS.

\medskip\noindent
{\bf Turning to variational geometry.}

\begin{example}
\label{hodge-birational-atoms-2026-example-let-s2}
Let $S$ be a complex K3 surface, for
example
\[
S=
\{x_0^4+x_1^4+x_2^4+x_3^4=0\}
\subset \mathbb{P}^3.
\]
The metric is not immediate here:
it is a consequence of the
Calabi--Yau theorem.
But it does exist:
there is a hyperk\"ahler metric making
$S$ a hyperk\"ahler manifold.
\end{example}

\begin{example}
\label{hodge-birational-atoms-2026-example-let-s2-duplicate-2}
Let $S$ be a complex K3 surface.
The symmetric power $S^{(n)}$ is
singular, but admits a resolution of
singularities
\[
r:S^{[n]}\to S^{(n)},
\]
where $S^{[n]}$ is the Hilbert scheme
of length $n$ subschemes of $S$.
The manifold $S^{[n]}$ is
hyperk\"ahler.
\end{example}

\begin{example}
\label{hodge-birational-atoms-2026-example-let-t}
Let
\[
T=\mathbb{C}^2/\mathbb{Z}^4.
\]
The Albanese morphism is
\[
a:T^{[n+1]}\to T,
\qquad
(x_0,\ldots,x_n)\mapsto
\sum_i x_i.
\]
The generalized Kummer variety
\[
K_n(T)=a^{-1}(0)
\]
is hyperk\"ahler.
\end{example}

\begin{example}
\label{hodge-birational-atoms-2026-example-o-grady}
There are O'Grady's exceptional
hyperk\"ahler manifolds of dimensions
$6$ and $10$.
\end{example}

Today we will be more interested in
the Hilbert scheme example.
In fact, these are all the examples we
know.

\medskip\noindent
{\bf Mukai flops.}

Today we discuss birational maps of
these manifolds.
Birational maps of manifolds with
trivial canonical bundle are similar
to flops.
They are modeled by so-called Mukai
flops.

Mukai flops are local models of a
birational transformation between
symplectic manifolds.

Recall the cotangent space
construction, which gives a symplectic
manifold.
Look at the cotangent bundle of
$\mathbb{P}^n$.
The idea is to contract the zero
section of this manifold, and then
blow it up in a different direction.

Let
\[
\dim_{\mathbb{C}}(V)=k+1
\]
and
\[
X=
\operatorname{Tot}
\bigl(\Omega^1_{\mathbb{P}(V)}\bigr).
\]
To contract the zero section we use
the Euler sequence
\[
0\to \Omega^1_{\mathbb{P}(V)}
\to
V^*\otimes \mathcal{O}_{\mathbb{P}(V)}
(-1)
\to
\mathcal{O}_{\mathbb{P}(V)}
\to 0.
\]
Also
\[
\mathcal{O}_{\mathbb{P}(V)}(-1)
\hookrightarrow
V\otimes \mathcal{O}_{\mathbb{P}(V)}.
\]
Hence
\[
\Omega^1_{\mathbb{P}(V)}
\hookrightarrow
\operatorname{End}(V)
\otimes
\mathcal{O}_{\mathbb{P}(V)}.
\]

\begin{remark}
\label{hodge-birational-atoms-2026-remark-point-sequence}
Here is the point of the sequence
calculation.
At a point $[u]\in\mathbb{P}(V)$,
the cotangent space is
\[
\Omega^1_{\mathbb{P}(V),[u]}
\cong
\operatorname{Hom}
(V/\mathbb{C}u,\mathbb{C}u).
\]
Such a map can be viewed as an
endomorphism $A$ of $V$ by first
quotienting by $\mathbb{C}u$ and then
landing in $\mathbb{C}u$.
Thus
\[
\operatorname{im}(A)\subset \mathbb{C}u,
\qquad
A(u)=0.
\]
In particular $A^2=0$ and
$\operatorname{rk}(A)\leq 1$.
Conversely, an endomorphism satisfying
these conditions gives a cotangent
vector at $[u]$.
\end{remark}

Therefore
\[
X=
\left\{
([u],A)\in
\mathbb{P}(V)\times
\operatorname{End}(V)
\middle|
\operatorname{im}(A)\subset
\mathbb{C}u,\ A^2=0
\right\}.
\]

The zero section is
$\mathbb{P}(V)\subset X$.
The projection
\[
\pi:X\to \operatorname{End}(V)
\]
contracts the zero section to the
origin.
Let
\[
Y=\pi(X)
=
\{A\in\operatorname{End}(V)
\mid
\operatorname{rk}(A)\leq 1,\ A^2=0\}.
\]

Analogously, put
\[
X'=
\left\{
([u],B)\in
\mathbb{P}(V^*)\times
\operatorname{End}(V^*)
\middle|
\operatorname{im}(B)\subset
\mathbb{C}u,\ B^2=0
\right\}
\]
and
\[
Y'=\pi'(X')
=
\{B\in\operatorname{End}(V^*)
\mid
\operatorname{rk}(B)\leq 1,\ B^2=0\}.
\]
Then $Y\cong Y'$ via
\[
A\mapsto A^*.
\]
Hence we get a diagram
\[
\xymatrix{
X=\operatorname{Tot}
\bigl(\Omega^1_{\mathbb{P}(V)}\bigr)
\ar@{-->}[rr]^{\mu}
\ar[dr]_{\pi}
&&
X'=\operatorname{Tot}
\bigl(\Omega^1_{\mathbb{P}(V^*)}\bigr)
\ar[dl]^{\pi'}
\\
& Y\cong Y' &
}
\]
where
\[
\mu([u],A)=([\ker(A)],A^*).
\]

The upshot is this.
We blow down the zero section, which
is a Lagrangian subvariety, and then
perform the other blow-up using the
dual vector space $V^*$.

Use the above diagram to construct a
local model of a birational map.
Let
\[
U=Y\cap B_\varepsilon(0),
\]
where
\[
B_\varepsilon(0)\subset
\operatorname{End}(V)
\]
is a small ball.
Then
\[
\xymatrix{
\pi^{-1}(U)\times\Delta^d
\ar@{-->}[rr]^{\mu\times\operatorname{id}}
\ar[dr]_{\pi\times\operatorname{id}}
&&
(\pi')^{-1}(U)\times\Delta^d
\ar[dl]^{\pi'\times\operatorname{id}}
\\
& U\times\Delta^d &
}
\]
where
\[
\Delta=\{z\in\mathbb{C}: |z|<1\}.
\]

\begin{definition}
\label{hodge-birational-atoms-2026-definition-let-m3}
Let
\[
\varphi:M\dashrightarrow M'
\]
be a bimeromorphic map between
complex manifolds.
We say that $\varphi$ is
{\it Mukai's elementary transformation
in codimension $k$} if there exists a
commutative diagram
\[
\xymatrix{
M \ar@{-->}[rr]^{\varphi}
\ar[dr]_{p}
&&
M' \ar[dl]^{p'}
\\
& N &
}
\]
where $p$ and $p'$ are proper small
bimeromorphic contractions, and for
every
\[
x\in N_{\mathrm{sing}}
\]
there exists an open neighborhood of
$x$ over which the diagram is
isomorphic to the local model
\[
U\times \Delta^d.
\]
\end{definition}

\begin{remark}
\label{hodge-birational-atoms-2026-remark-prefer-describe}
We prefer to describe the situation in
terms of contractions, as in the
diagram
\[
\xymatrix{
M \ar@{-->}[rr]^{\varphi}
\ar[dr]_{p}
&&
M' \ar[dl]^{p'}
\\
& N &
}
\]
When we contract something of
codimension bigger than $1$, the
contracted space $N$ will be singular.
This is why the definition refers to
$N_{\mathrm{sing}}$.
Further, $N_{\mathrm{sing}}$ is
smooth.
So $N$ is singular, but not too badly.

rather than in terms of blow-ups, as
in a diagram
\[
\xymatrix{
& \widetilde M \ar[dl]\ar[dr] &\\
M \ar@{-->}[rr] && M'.
}
\]
\end{remark}

\medskip\noindent
{\bf Hu--Yau conjecture.}

Let
\[
\varphi:M\dashrightarrow M'
\]
be a bimeromorphic map between
hyperk\"ahler manifolds.
Hu and Yau conjectured that, up to
removing from $M$ and $M'$ some
subvarieties of codimension strictly
bigger than $2$, the map $\varphi$ is
a composition of several Mukai
elementary transformations.

\begin{definition}
\label{hodge-birational-atoms-2026-definition-let-m3-duplicate-2}
Let
\[
\psi:M\dashrightarrow M'
\]
be a bimeromorphic map between two
complex manifolds.
We say that $\psi$ is a {\it Hu--Yau
transformation} if there exist closed
subvarieties
\[
Z\subset M,
\qquad
Z'\subset M',
\]
of codimension strictly bigger than
$2$ such that
\[
\psi:M\setminus Z
\dashrightarrow
M'\setminus Z'
\]
is a Mukai elementary transformation
in codimension two.
\end{definition}

\begin{theorem}[Wierzba--Wi\'sniewski]
\label{hodge-birational-atoms-2026-theorem-assume-m}
Assume that $M$ is a symplectic
manifold of dimension $4$.
Let
\[
\pi:M\to N
\]
be a small proper bimeromorphic
contraction such that $N$ is
quasiprojective.
Then the exceptional locus of $\pi$
is a disjoint union of several copies
of $\mathbb{P}^2$, and $\pi$ is
locally analytically isomorphic to the
contraction of the zero section in
\[
\operatorname{Tot}
\bigl(\Omega^1_{\mathbb{P}^2}\bigr).
\]
\end{theorem}

This is the first step of the Mukai
flop.
It is the contraction of
$\mathbb{P}^2$, which sits in its
cotangent bundle as the zero section.
Thus this theorem is exactly the
picture explained above.

\begin{theorem}[Wierzba--Wi\'sniewski]
\label{hodge-birational-atoms-2026-theorem-hu-yau}
The Hu--Yau conjecture holds in
dimension four.
More precisely, if
\[
\varphi:M\dashrightarrow M'
\]
is a birational map between projective
hyperk\"ahler fourfolds, then
$\varphi$ is a composition of several
Mukai flops.
\end{theorem}

The approach to this conjecture is
different.
Instead of using the MMP, one uses the
wall-and-chamber structure.
Let us state the result first.

\begin{theorem}
\label{hodge-birational-atoms-2026-theorem-amerik-soldatenkov}
[Amerik--Soldatenkov--Verbitsky]
Let $M$ and $M'$ be compact
hyperk\"ahler manifolds, and let
\[
\varphi:M\dashrightarrow M'
\]
be a bimeromorphic map which is not
regular in codimension two.
Then $\varphi$ is a composition of a
finite number of Hu--Yau
transformations.
\end{theorem}

\medskip\noindent
{\bf Idea of the proof.}

\medskip\noindent
{\bf Wall crossing.}

To describe birational transformations
of hyperk\"ahler manifolds, we look at
K\"ahler cones.
Look at $H^2(M)$ with its
Beauville--Bogomolov--Fujiki form
$q$.
The restriction of $q$ to
$H^{1,1}(M)$ has signature
\[
(1,b_2-3).
\]
Let $\mathcal K_M$ be the K\"ahler
cone of $M$, and let $\mathcal C_M$
be the connected component of the
positive cone
\[
\{x\in H^{1,1}(M): q(x)>0\}
\]
which contains $\mathcal K_M$.

\begin{theorem}[Amerik--Verbitsky]
\label{hodge-birational-atoms-2026-theorem-exists-locally}
There exists a locally finite
collection of hyperplanes
\[
W_i\subset \mathcal C_M
\]
given by orthogonal complements to the
MBM classes, such that for any
bimeromorphic map
\[
\varphi:M\dashrightarrow M'
\]
with $M'$ hyperk\"ahler, the set
\[
\varphi^*\mathcal K_{M'}
\]
is a connected component of
\[
\mathcal C_M\setminus \bigcup_i W_i.
\]
The union of such chambers is the
birational cone.
\end{theorem}

\begin{definition}
\label{hodge-birational-atoms-2026-definition-bimeromorphic-map}
A {\it bimeromorphic map}
\[
\varphi:M\dashrightarrow M'
\]
between two hyperk\"ahler manifolds is
a {\it wall-crossing flop} if
\[
\mathcal K_M
\quad\text{and}\quad
\varphi^*\mathcal K_{M'}
\]
share a common wall.
\end{definition}

For any wall-crossing flop there
exists a diagram
\[
\xymatrix{
M \ar@{-->}[rr]^{\varphi}
\ar[dr]_{p}
&&
M' \ar[dl]^{p'}
\\
& N &
}
\]
where $N$ is a singular symplectic
variety.
Now $N$ may be more singular than
before.
The maps $p$ and $p'$ are proper small
bimeromorphic contractions.

Since the collection of walls in
$\mathcal C_M$ is locally finite, any
bimeromorphic map
\[
\varphi:M\dashrightarrow M'
\]
between hyperk\"ahler manifolds is a
composition of a finite number of
wall-crossing flops.

\begin{example}
\label{hodge-birational-atoms-2026-example-let-s3}
Let
\[
S\subset \mathbb{P}^3
\]
be a quartic K3 surface containing a
line $C$.
Take
\[
M=S^{[3]}.
\]
Picard rank $1$ is not very
interesting from the point of view of
wall-and-chamber structure.
The pencil of planes passing through
$C$ defines an elliptic fibration
\[
f:S\to \mathbb{P}^1
\]
with fiber $F$.
In the basis $[C],[F]$, the
intersection form on $\operatorname{Pic}
(S)$ is
\[
\begin{pmatrix}
-2 & 3\\
3 & 0
\end{pmatrix}.
\]
Moreover
\[
\operatorname{Pic}(M)
\cong
\operatorname{Pic}(S)\oplus\mathbb{Z}\delta,
\]
where
\[
\delta=[\Delta]/2,
\]
and $\Delta$ is the exceptional
divisor of the contraction
\[
S^{[3]}\to S^{(3)}.
\]
Thus $\delta$ defines a wall in the
wall-and-chamber structure.
\end{example}

There is a picture of the birational
ample cone of $S^{[3]}$, as seen in
the Poincare hyperbolic plane.

Some MBM classes on $M$ are:
\begin{itemize}
\item $\delta$, the exceptional
divisor of the contraction
$S^{[3]}\to S^{(3)}$;
\item $\alpha=[C]$, the divisor of
subschemes that intersect $C$;
\item $\epsilon=2[F]-\delta$, the
closure of the image of
\[
\operatorname{Hilb}^2(S\setminus F)
\times S
\dashrightarrow
S^{(3)};
\]
\item $\gamma=2[C]-\delta$, the set
of subschemes supported on $C$, which
is isomorphic to
\[
C^{(3)}\cong \mathbb{P}^3.
\]
\end{itemize}

We have decomposed everything into a
sequence of wall-crossing flops.
We need to show that those
wall-crossing flops are in fact
Hu--Yau transformations.

Consider a wall-crossing flop
\[
\xymatrix{
M \ar@{-->}[rr]^{\varphi}
\ar[dr]_{p}
&&
M' \ar[dl]^{p'}
\\
& N. &
}
\]

\begin{theorem}[Bakker--Lehn]
\label{hodge-birational-atoms-2026-theorem-exists-small}
There exists a small locally trivial
deformation of the above diagram which
makes $M$, $M'$, and $N$ projective.
\end{theorem}

\begin{theorem}[Kaledin]
\label{hodge-birational-atoms-2026-theorem-symplectic-variety}
The symplectic variety $N$ admits a
finite locally closed stratification
such that the strata are symplectic
manifolds.
\end{theorem}

\begin{remark}
\label{hodge-birational-atoms-2026-remark-open-stratum}
The open stratum of the Kaledin
stratification is the smooth locus of
$N$.
The symplectic forms on the strata
glue to a nice Poisson structure.
\end{remark}

Consider a wall-crossing flop
\[
\xymatrix{
M \ar@{-->}[rr]^{\varphi}
\ar[dr]_{p}
&&
M' \ar[dl]^{p'}
\\
& N &
}
\]
such that the indeterminacy locus of
$\varphi$ has a component of
codimension $2$.
Let
\[
N'\subset N_{\mathrm{sing}}
\]
be a stratum of codimension $4$, and
let $x\in N'$.
Use the Poisson structure on $N$ to
construct a slice
\[
S\subset N
\]
passing through $x$, transverse to
$N'$, and symplectic.

Then
\[
p_S:p^{-1}(S)\to S
\]
is a small contraction of symplectic
fourfolds.
Applying the theorem of
Wierzba--Wi\'sniewski, $p_S$ is the
contraction of a copy of
$\mathbb{P}^2$.
Thus over $N'$ we have a
$\mathbb{P}^2$-bundle for both $p$
and $p'$.
It follows that $\varphi$ is a
Hu--Yau transformation.
This completes the proof.

\section{Holonomy of the Obata
connection on Joyce hypercomplex
manifolds}
\label{hodge-birational-atoms-2026-section-holonomy-obata-joyce}

\medskip\noindent
{\bf Speaker.}

Beatrice Brienza

\medskip\noindent
{\bf Title.}

Holonomy of the Obata connection on
Joyce hypercomplex manifolds

\medskip\noindent
{\bf Abstract.}

The Obata connection on a hypercomplex
manifold is the unique torsion-free
connection that preserves the
hypercomplex structure. Up to taking
the product with a torus, Joyce
constructed left-invariant
hypercomplex structures on compact
semisimple Lie groups and certain
related homogeneous spaces.
Soldatenkov showed in 2011 that every
Joyce hypercomplex structure on SU(3)
has Obata holonomy equal to the
quaternionic general linear group and
it was expected that the same should
hold for all Joyce hypercomplex
manifolds. For all such group
manifolds except for SU(2n+1), we
will show that the holonomy group is
strictly contained in the quaternionic
general linear group. The case of
SU(2n+1) is more subtle, however for
every n > 1, there still exist
infinitely many Joyce hypercomplex
structures with Obata holonomy
strictly contained in the quaternionic
general linear group and at least one
with full holonomy on SU(5). This talk
is based on a joint work with Udhav
Fowdar, Giovanni Gentili and Luigi
Vezzoni.

\medskip\noindent
{\bf Upshot.} 
The talk starts recalling the analogue
of Berger's theorem for torsion-free
(not necessarily metric) connections.
The idea is then to apply this to the
Obata connection.

\medskip\noindent
{\bf Notes.}

\begin{theorem}[Berger]
\label{hodge-birational-atoms-2026-theorem-let-m2}
Let $M$ be an oriented simply
connected Riemannian manifold of
dimension $n$.
Assume that $M$ is irreducible and
not locally symmetric.
Then its Riemannian holonomy group
must be one of the following:
\[
\operatorname{SO}(n),\quad
\operatorname{U}(m),\quad
\operatorname{SU}(m),\quad
\operatorname{Sp}(m),
\]
\[
\operatorname{Sp}(m)\operatorname{Sp}(1),
\quad
G_2,\quad
\operatorname{Spin}(7).
\]
Here $n=2m$ for the unitary cases,
$n=4m$ for the symplectic cases,
$n=7$ for $G_2$, and $n=8$ for
$\operatorname{Spin}(7)$.
\end{theorem}

At the time this list was a list of
possible irreducible Riemannian
holonomies.
Examples of Riemannian manifolds with
holonomy $G_2$ and
$\operatorname{Spin}(7)$ were found
by Bryant--Salamon in 1989.

\medskip\noindent
{\bf Possible generalization.}

Consider torsion-free connections
which are not metric.
In fact, this idea was discussed in
the same paper by Berger.
The theorem actually still makes sense
when the signature of the metric is
not Riemannian.

Merkulov--Schwachh\"ofer gave in 1999
the complete list of all possible
irreducible holonomy groups for
torsion-free connections.

What about the Obata connection?
Its holonomy group fits this
scenario.

\medskip\noindent
{\bf Hypercomplex geometry in a
nutshell.}

Let $M$ be a manifold.
A {\it hypercomplex structure} on $M$
is a
triple $(I,J,K)$ of complex structures
on $M$ satisfying the quaternionic
relations:
\[
I^2=J^2=K^2=IJK=-1.
\]
Then
\[
\dim_{\mathbb{R}} M=4n.
\]
We call $(M,I,J,K)$ a {\it
hypercomplex manifold}.

\begin{example}
\label{hodge-birational-atoms-2026-example-manifold-n}
The manifold $\mathbb{H}^n$, with
the complex structures given by left
multiplication by $i$, $j$, and $k$,
is a hypercomplex manifold.
\end{example}

\medskip\noindent
{\bf Question.}

Is any hypercomplex manifold locally
isomorphic to
\[
(\mathbb{H}^n,i,j,k)?
\]

For each
\[
L\in\{I,J,K\},
\]
we have
\[
(M,L)\cong_{\mathrm{loc}}
(\mathbb{C}^{2n},i).
\]
That is, each complex structure is
locally standard by the
Newlander--Nirenberg theorem.
However, in general
\[
(M,I,J,K)\not\cong_{\mathrm{loc}}
(\mathbb{H}^n,i,j,k).
\]
Thus there is no direct
hypercomplex version of the
Newlander--Nirenberg theorem.


\begin{theorem}[Obata]
\label{hodge-birational-atoms-2026-theorem-tm-exists}
On $TM$ there exists a unique
torsion-free connection $\nabla$
which preserves $I$, $J$, and $K$.
\end{theorem}

Furthermore,
\[
(M,I,J,K)\cong_{\mathrm{loc}}
(\mathbb{H}^n,i,j,k)
\]
if and only if
\[
\operatorname{Hol}^0(\nabla)=\{e\}.
\]

That is, a hypercomplex manifold is locally
isomorphic to $\mathbb{H}^n$
with $i,j,k$ precisely when the
Obata connection is flat.

\medskip\noindent
{\bf Examples of Obata-flat
manifolds.}

\begin{enumerate}
\item
$(\mathbb{H}^n,i,j,k)$.
\item Many hypercomplex nilmanifolds.
There are many examples among the
$2$-step ones; see for instance
DF, ABB, and BD.
\item Hopf surfaces
\[
\mathbb{H}\setminus\{0\}/\langle q\rangle,
\]
where
\[
q\in\mathbb{H}^*
\quad\text{and}\quad
|q|>1.
\]
\end{enumerate}

The holonomy of the Obata connection
satisfies
\[
\operatorname{Hol}(\nabla)
\subseteq
\operatorname{GL}(n,\mathbb{H}).
\]

Only three subgroups of
$\operatorname{GL}(n,\mathbb{H})$
appear in the
Merkulov--Schwachh\"ofer list.

\begin{enumerate}
\item $\operatorname{Sp}(n)$.
We have
\[
\operatorname{Hol}(\nabla)
\subseteq
\operatorname{Sp}(n)
\]
if and only if the Obata connection is
the Levi-Civita connection of a
metric.
Equivalently, $(g,I,J,K)$ is
hyperk\"ahler. This brings us back
to the Berger case.

\item
\[
\operatorname{SL}(n,\mathbb{H})
=
[\operatorname{GL}(n,\mathbb{H}),
\operatorname{GL}(n,\mathbb{H})].
\]
Hypercomplex nilmanifolds have Obata
holonomy contained in
$\operatorname{SL}(n,\mathbb{H})$;
see BDV.

\item The most mysterious case is
$\operatorname{GL}(n,\mathbb{H})$.
Soldatenkov showed that
$\operatorname{SU}(3)$, with any
compact invariant hypercomplex
structure, has Obata holonomy
\[
\operatorname{GL}(2,\mathbb{H}).
\]
In fact, this is our first example of
a Joyce hypercomplex manifold.
Before this, it wasn't known whether
there existed a hypercomplex manifold
with full holonomy group.
\end{enumerate}

\medskip\noindent
{\bf Joyce's decomposition.}

Let $G$ be a compact semisimple Lie
group.
Let
\[
T\subseteq G
\]
be a maximal torus, and put
\[
\mathfrak g=\operatorname{Lie}(G),
\qquad
\mathfrak t=\operatorname{Lie}(T).
\]
Let $r=\operatorname{rk}(G)$, and let
$\Delta$ be a system of ordered roots
with respect to $\mathfrak t^\mathbb C$.



\medskip\noindent
{\bf Step 1.}

Pick a maximal positive root
\[
\alpha\in\Delta.
\]
Set
\[
\widetilde{\partial}_1
=
\mathfrak g_\alpha
\oplus
\mathfrak g_{-\alpha}
\oplus
[\mathfrak g_\alpha,\mathfrak g_{-\alpha}]
\cong
\mathfrak{sl}(2,\mathbb{C}),
\]
and
\[
\partial_1
=
\mathfrak g\cap\widetilde{\partial}_1
\cong
\mathfrak{su}(2).
\]
Then
\[
\mathfrak g_0=\mathfrak g
=
\partial_1
\oplus
\mathfrak c_1
\oplus
\mathfrak f_1,
\]
where $\mathfrak c_1$ is the
centralizer of $\partial_1$, and
$\mathfrak f_1$ is the orthogonal
complement of
\[
\partial_1\oplus\mathfrak c_1
\]
with respect to the Killing form $B$.

\medskip\noindent
{\bf Step 2.}

Write
\[
\mathfrak c_1
=
\mathfrak g_1\oplus\mathfrak b_1.
\]
Apply Step 1 to $\mathfrak g_1$.

\medskip\noindent
{\bf Step $m+1$.}

At the final step the centralizer is
abelian:
\[
\mathfrak c_{m+1}
=
\mathfrak b_{m+1}.
\]
Thus the inductive process stops when
the centralizer is only abelian.

We obtain
\[
\mathfrak g
=
\mathfrak b
\oplus
\bigoplus_{j=1}^m \partial_j
\oplus
\bigoplus_{j=1}^m \mathfrak f_j.
\]
Here
\[
\mathfrak b
=
\mathfrak b_1\oplus\cdots
\oplus\mathfrak b_{m+1}
\]
is an abelian subalgebra of dimension
$r-m$.
Each $\partial_j$ is a subalgebra
isomorphic to $\mathfrak{su}(2)$.
Moreover,
\[
[\mathfrak b,\partial_j]=0,
\qquad
[\partial_i,\partial_j]=0
\quad (i\ne j).
\]
The spaces $\mathfrak f_j$ have
dimension $4k_j$, and
\[
[\partial_i,\mathfrak f_j]=0
\quad (i<j),
\]
while
\[
[\partial_i,\mathfrak f_i]
\subseteq
\mathfrak f_i.
\]
Thus $\partial_i\cong\mathfrak{su}(2)$
acts on $\mathfrak f_i$ through the
standard action on $\mathbb{C}^{2k_i}$.
This is the Joyce decomposition of
$\mathfrak g$.

\medskip\noindent
{\bf Some examples.}

Consider $\operatorname{Sp}(2)$.
We write
\[
\mathfrak{sp}(2)
=
\{A\in\mathfrak{gl}(2,\mathbb{H})
\mid A+\overline{A}^t=0\}.
\]
Thus an element of $\mathfrak{sp}(2)$
has the form
\[
A=
\begin{pmatrix}
x & y\\
-\overline y & z
\end{pmatrix},
\qquad
x,z\in\operatorname{Im}(\mathbb{H}).
\]
In this case one has
\[
\mathfrak{sp}(2)
=
\partial_1\oplus\partial_2
\oplus\mathfrak f_1.
\]
Here
\[
\partial_1
=
\left\{
\begin{pmatrix}
x & 0\\
0 & 0
\end{pmatrix}
\right\},
\]
with the remaining pieces to be filled
in.

One may also consider colored Dynkin
diagrams.
In OP, 1998, the colors account for
the Joyce decomposition.
Using this other technique one obtains
the same decomposition, namely
\[
\mathfrak{sp}(k)
=
\partial_1\oplus\mathfrak c_1
\oplus\mathfrak f_1.
\]
One can also obtain
\[
\mathfrak{sp}(k)
=
\partial_1\oplus\mathfrak{sp}(k-1)
\oplus(\text{missing term}).
\]

\medskip\noindent
{\bf Hypercomplex structures.}

\noindent
Starting from any compact semisimple
Lie group, we can obtain a hypercomplex
manifold by multiplying by a torus.

Let $G$ be compact and semisimple,
with
\[
\operatorname{rk}(G)=r.
\]
Consider
\[
\mathbb{T}^{2m-r}\times G.
\]
Its Lie algebra is
\[
(2m-r)\mathfrak u(1)\oplus\mathfrak g
=
(2m-r)\mathfrak u(1)
\oplus\mathfrak b
\oplus
\bigoplus_{j=1}^m\partial_j
\oplus
\bigoplus_{j=1}^m\mathfrak f_j.
\]
The idea here is that the $\mathfrak{u}(1)$
factor is isomorphic to $\mathbb{R}^n$
by a dimension count 
(this is the torus part).
Joining
(via choosing a basis and certain
specific pairing with the $\partial_j$
and $\mathfrak{f}_j$ parts)
this with a copy of $\mathfrak{su}(2)$
we get a copy of the quaternions:
$\mathfrak{u}(1) \oplus \mathfrak{su}(2)
\cong \mathbb{H}$ with its
standard quaternionic structure.

Using the Joyce decomposition and
additional formulas, Joyce constructs
a hypercomplex structure
$(I,J,K)$ on this Lie algebra.
This is a hypercomplex structure in
the sense of Joyce, 1992.

\noindent
It turns out that any compact Lie group
with a left invariant hypercomplex
structure has to be constructed this way:

\begin{theorem}[DT, BGP]
\label{hodge-birational-atoms-2026-theorem-joyce-hypercomplex}
Joyce hypercomplex manifolds are in
one-to-one correspondence with compact
Lie groups with left-invariant
hypercomplex structures.
\end{theorem}

Let $G$ be simple.
The dimension of the torus can be
checked from a result by Spindel and
collaborators. This is a very important
result in this theory.

\[
\begin{array}{c|l}
\dim & \text{example}\\
\hline
4 &
\operatorname{SU}(2)\times S^1,
\text{ Obata-flat, simplest case}\\
8 &
\operatorname{SU}(3),
\operatorname{Hol}(\nabla)
=\operatorname{GL}(2,\mathbb{H})\\
4n &
\text{Folklore conjecture (see below)}
\end{array}
\]

\medskip\noindent
{\bf Folklore conjecture.}

Any Joyce hypercomplex manifold in the
list of Spindel and collaborators,
except for
\[
S^1\times\operatorname{SU}(2),
\]
has full Obata holonomy, that is,
\[
\operatorname{Hol}(\nabla)
=
\operatorname{GL}(n,\mathbb{H}).
\]

\begin{theorem}[BFGV]
\label{hodge-birational-atoms-2026-theorem-any-joyce}
For any Joyce hypercomplex manifold in
the list of Spindel and collaborators,
with the exception of
$\operatorname{SU}(2k+1)$, one has
\[
\operatorname{Hol}(\nabla)
\subset
\operatorname{GL}(n,\mathbb{H}).
\]
Moreover, this inclusion is strict.
\end{theorem}

\medskip\noindent
{\bf Idea of proof.}

We may assume that $\dim M\geq 8$.
Let
\[
M=\mathbb{T}^{2m-r}\times G.
\]
Then
\[
(2m-r)\mathfrak u(1)\oplus\mathfrak g
\]
decomposes as
\[
(2m-r)\mathfrak u(1)
\oplus\mathfrak b
\oplus
\bigoplus_{j=1}^m\partial_j
\oplus
\bigoplus_{j=1}^m\mathfrak f_j.
\]
There are further formulas describing
the corresponding decomposition of
the Lie algebra.

\medskip\noindent
{\bf Step 1.}

\begin{lemma}[BFGV]
\label{hodge-birational-atoms-2026-lemma-let-m2}
Let
\[
M=\mathbb{T}^{2m-r}\times G,
\]
and let $(I,J,K)$ be a Joyce
hypercomplex structure on $M$.
Equivalently, $(I,J,K)$ is a
left-invariant hypercomplex structure.
Then
\[
\nabla_X e_1
=
\begin{cases}
-X, & X\in\mathfrak h_j\oplus
\mathfrak f_j,\\
0, & \text{otherwise}.
\end{cases}
\]
\end{lemma}

More details of the formula remain to
be filled in.

\medskip\noindent
{\bf Step 2.}

By Step 1, the holonomy representation
is reducible.
Therefore the holonomy group cannot be
$\operatorname{GL}(n,\mathbb{H})$.
In fact, if $p\in M$ and $\gamma$ is
a loop based at $p$, then the parallel
transport $P_\gamma$ has the form
\[
P_\gamma=
\begin{pmatrix}
\text{missing}
\end{pmatrix}.
\]

\medskip\noindent
{\bf Step 3.}

Look at every simple Lie group and
check whether it has a trivial
$\mathfrak f_j$ summand.
All of them except
$\operatorname{SU}(2k+1)$ have at
least one such summand.

We have already seen that
\[
\mathfrak{sp}(k)
=
\partial_1\oplus\mathfrak{sp}(k-1)
\oplus\mathfrak f_1.
\]
By induction, in the case $k-1$ we
start from
\[
\mathfrak{sp}(2)
=
\partial_1\oplus\partial_2
\oplus\mathfrak f_1,
\qquad
\mathfrak f_2=\{0\}.
\]
Assume that $\mathfrak{sp}(k-1)$ has
one trivial $\mathfrak g_j$ summand.
Since $\mathfrak{sp}(k-1)$ appears in
the decomposition of
$\mathfrak{sp}(k)$, we look for the
$\mathfrak f_j$ component inside the
$\mathfrak{sp}(k-1)$ factor.

\medskip\noindent
{\bf The $\operatorname{SU}(2k+1)$
case.}

For $k=1$, this is the Soldatenkov
case.
What happens when $k>1$?

For $k=2$, the Joyce decomposition of
$\mathfrak{su}(5)$ is
(the centralizer,
i.e. abelian, coming from
the Joyce decomposition construction, 
part is $\mathfrak{su}(3)
\oplus \mathfrak{u}(1)$)
\[
\mathfrak{su}(5)
=
\partial_1
\oplus
\underbrace{\mathfrak{su}(3)}_{
\langle e_1\rangle
\oplus\partial_1\oplus\mathfrak f_2}
\oplus
\underbrace{\mathfrak u(1)}_{
\langle\delta_1\rangle}
\oplus
\mathfrak f_1.
\]
The point is that we can find a copy
of $\mathfrak{su}(3)$ inside.

\begin{theorem}[BFGV]
\label{hodge-birational-atoms-2026-theorem-joyce-hypercomplex2}
If the Joyce hypercomplex structure
has block form
\[
(I,J,K)=
\begin{pmatrix}
a & 0\\
b & d
\end{pmatrix},
\]
then
\[
\operatorname{Hol}(\nabla)
\subset
\operatorname{GL}(n,\mathbb{H}).
\]
Moreover, this inclusion is strict.
\end{theorem}

There is a decomposition of the form
\[
\mathfrak{su}(5)
=
\mathfrak b_1\oplus\mathfrak f_1
\oplus\mathfrak b_2\oplus
\mathfrak f_2
\oplus(\text{missing term}).
\]
The point is that $\mathfrak{su}(3)$
is not only a subalgebra, but a
hypercomplex subalgebra.
It is preserved by the Obata
connection.
Thus we get a reduction of the
holonomy.
The same procedure works for other
$\operatorname{SU}(2k+1)$, but not
for $2k+1<5$.

Beatrice et al. tried to check the opposite
case, when $(I,J,K) = $ 2x2 matrix
but with $c$ nonzero. Computationally hard.
They conjecture that in this case
the holonomy is full.

Florit: in that case, when $c$ is nonzero,
did you try to check that it has trivial
holonomy? Because this implies
irreducibility. Answer: yes, and miserably
failed.

\section{Bogomolov--Guan precursors
and quasi-diagonals in a product of
elliptic curves}
\label{hodge-birational-atoms-2026-section-bogomolov-guan-quasi-diagonals}

\medskip\noindent
{\bf Speaker.}

Misha Verbitsky

\medskip\noindent
{\bf Title.}

Bogomolov--Guan precursors and
quasi-diagonals in a product of
elliptic curves

\medskip\noindent
{\bf Joint work.}

Leila Abubakarova and Alexandra
Kuznetsova.

The main point of the talk is that
quasidiagonals in products of elliptic
curves control projective subvarieties
of the Bogomolov--Guan precursor.
This gives a concrete way to study the
precursor through the geometry of
elliptic curves and their products.

\begin{definition}
\label{hodge-birational-atoms-2026-definition-let-l}
Let $L$ be an ample line bundle on an
elliptic curve $E$.
Let
\[
p_i:E^n\to E
\]
denote the projection to the $i$th
component.
A {\it quasidiagonal} is a curve
\[
S\subset E^n
\]
such that, for all $i$ and $j$, the
line bundle
\[
V_{i,j}
=
\left.
p_i^*L
\right|_S
\otimes
\left(
\left.
p_j^*L
\right|_S
\right)^{-1}
\]
is torsion.
In other words, there exists $k>0$
such that
\[
V_{i,j}^{\otimes k}
=
\mathcal{O}_S.
\]
\end{definition}

\begin{remark}
\label{hodge-birational-atoms-2026-remark-n-2}
For $n=2$, a curve
\[
S\subset E\times E
\]
is a correspondence from $E$ to $E$.
Since both factors are the same
elliptic curve $E$, the two
projections restrict to maps
\[
p_1|_S,p_2|_S:S\to E.
\]
Thus both pullbacks of $L$ live on the
same curve $S$, and can be compared.
The quasidiagonal condition does not
use a pushforward; it only says that
the ratio of these two pullbacks is
torsion.
\end{remark}

The theorem says there are at most
countably many quasidiagonals for any
given pair $(E,L)$.

\begin{example}
\label{hodge-birational-atoms-2026-example-let-e}
Let
\[
\nu:E\to E
\]
be an automorphism of order $d$.
Given a line bundle $L_1$ on $E$, set
\[
L=
\bigotimes_{i=0}^{d-1}(\nu^i)^*L_1.
\]
Then $L$ is $\nu$-invariant.
Therefore the graph
\[
\Gamma_\nu\subset E^2
\]
of $\nu$ is a quasidiagonal for
$(E,L)$.
\end{example}

\medskip\noindent
{\bf Open questions.}

\begin{enumerate}
\item Can a quasidiagonal have
geometric genus $p_g(C)>1$?
There are examples which are elliptic
and singular.

\item Call a line bundle
$L\in\operatorname{Pic}^n(E)$
commensurable with
$\mathcal{O}(n[e])$ if
\[
L\otimes\mathcal{O}(n[e])^{-1}
\]
is torsion.
Is it true that there are no
restricted quasidiagonals on $(E,L)$
when $L$ is not commensurable with
$\mathcal{O}(n[e])$?

\item Consider $E^n$, where $n$ is
odd and $E$ does not admit complex
multiplication.
Are there any restricted
quasidiagonals in $E^n$, for some
ample line bundle $L$ on $E$?

\item In the case $n=3$ and
\[
E=\mathbb{C}/\mathbb{Z}[\zeta],
\]
where $\zeta$ is a primitive third
root of unity, there is an example of
a restricted quasidiagonal in $E^3$.
Does $E^3$ contain any restricted
quasidiagonal except this one?
\end{enumerate}

\begin{remark}
\label{hodge-birational-atoms-2026-remark-questions-may}
These questions may be accessible
projects for students learning
algebraic geometry.
\end{remark}

\medskip\noindent
{\bf Hilbert schemes.}

\begin{definition}
\label{hodge-birational-atoms-2026-definition-let-m4}
Let $M$ be a complex surface.
The {\it Hilbert scheme of points}
$M^{[n]}$
is a natural smooth resolution of the
$n$th symmetric power $M^{(n)}$.
\end{definition}

\begin{remark}
\label{hodge-birational-atoms-2026-remark-m-holomorphically}
If $M$ is holomorphically symplectic,
then $M^{[n]}$ is also
holomorphically symplectic.
This follows from Serre duality.
\end{remark}

\begin{remark}
\label{hodge-birational-atoms-2026-remark-hilbert-scheme}
The Hilbert scheme of a K3 surface is
simply connected and holomorphically
symplectic.
The Hilbert scheme of a torus $T$ is
not simply connected, but the fiber of
its Albanese map
\[
T^{[n]}\to T
\]
has finite fundamental group.
The universal cover of this fiber is
called a generalized Kummer variety.
In this way one obtains the two main
examples of simply connected
holomorphically symplectic manifolds.
\end{remark}

Todorov conjectured that any compact
simply connected holomorphically
symplectic manifold should be
K\"ahler.
Guan produced counterexamples:
compact simply connected
holomorphically symplectic manifolds
which are not K\"ahler.
Bogomolov then used Guan's example to
produce what are now called
Bogomolov--Guan manifolds.

\begin{definition}
\label{hodge-birational-atoms-2026-definition-let-l2}
Let $L$ be a line bundle on an
elliptic curve $E$ with
\[
c_1(L)\neq 0.
\]
Let $\widetilde S$ be the
$\mathbb{C}^*$-bundle on $E$ obtained
by removing the zero section:
\[
\widetilde S=\operatorname{Tot}(L)
\setminus 0.
\]
Fix a complex number $\lambda$ with
$|\lambda|>1$, and let
\[
h_\lambda:\widetilde S\to\widetilde S
\]
be the corresponding homothety.
The quotient
\[
\widetilde S/\langle h_\lambda\rangle
\]
is called a {\it primary Kodaira
surface}, or {\it Kodaira--Thurston
surface}, or {\it Kodaira surface}.
\end{definition}

\begin{remark}
\label{hodge-birational-atoms-2026-remark-kodaira-surface}
The Kodaira surface is a principal
elliptic fibration
\[
p:\operatorname{Kod}\to E,
\]
with fiber another elliptic curve
$E_1$.
It is holomorphically symplectic.
Therefore
\[
\operatorname{Kod}^{[n]}
\]
is also holomorphically symplectic,
but it is not simply connected.
\end{remark}

Applying $p$ to each component gives
a map
\[
\operatorname{Kod}^{(n)}\to E^{(n)}.
\]
Composing with the Hilbert--Chow
resolution
\[
\operatorname{Kod}^{[n]}
\to
\operatorname{Kod}^{(n)}
\]
and with addition on the elliptic
curve gives a holomorphic map
\[
\operatorname{Kod}^{[n]}\to E.
\]
Let
\[
\mathcal D\subset \operatorname{Kod}^{[n]}
\]
be the fiber over $0\in E$.
If $\Omega$ is the holomorphic
symplectic form on
$\operatorname{Kod}^{[n]}$, then
\[
T\mathcal D^{\perp\Omega}
\subset
T\mathcal D.
\]
The corresponding rank-one
distribution is the characteristic
foliation.

The leaf space of this foliation is
holomorphically symplectic, but may
have orbifold singularities.
Equivalently, in the precursor model
one considers
\[
D\subset \operatorname{Kod}^{n+1}
\]
defined by
\[
D=
\left\{
(z_1,\ldots,z_{n+1})
\mid
\sum_i p(z_i)=0
\right\}.
\]
The diagonal $E_1$-action on $D$ is
by leaves of the characteristic
foliation.
The quotient
\[
D/E_1
\]
is the restricted
{\it Bogomolov--Guan precursor}.

\begin{theorem}
\label{hodge-birational-atoms-2026-theorem-let-p}
Let
\[
P:X\to B
\]
be the Lagrangian fibration on the
Bogomolov--Guan precursor, and let
\[
Z\subset X
\]
be an irreducible subvariety.
Then $Z$ is projective if and only if
\[
P(Z)\subset B\subset E^{n+1}
\]
is either a point or a quasidiagonal.
\end{theorem}

\begin{definition}
\label{hodge-birational-atoms-2026-definition-nilmanifold-smooth2}
A {\it nilmanifold} is a smooth
manifold equipped with a transitive
action of a nilpotent Lie group.
\end{definition}

\begin{remark}
\label{hodge-birational-atoms-2026-remark-all-nilmanifolds2}
All nilmanifolds are obtained as
quotients
\[
M=G/\Gamma,
\]
where $G$ is a nilpotent Lie group
and $\Gamma\subset G$ is a discrete
subgroup.
The Bogomolov--Guan precursor is
particularly easy to study because it
is a complex nilmanifold.
\end{remark}

\begin{definition}
\label{hodge-birational-atoms-2026-definition-let-g4}
Let $\mathfrak g$ be a real Lie
algebra.
An {\it integrable complex structure}
on $\mathfrak g$ is a subalgebra
\[
\mathfrak g^{1,0}
\subset
\mathfrak g\otimes_{\mathbb{R}}
\mathbb{C}
\]
such that
\[
\mathfrak g^{1,0}\oplus
\overline{\mathfrak g^{1,0}}
=
\mathfrak g\otimes_{\mathbb{R}}
\mathbb{C}.
\]
\end{definition}

\begin{remark}
\label{hodge-birational-atoms-2026-remark-such-decomposition2}
Such a decomposition is exactly the
same thing as a complex structure on
the Lie algebra.
Indeed, one lets the complex structure
act by $\sqrt{-1}$ on
$\mathfrak g^{1,0}$ and by
$-\sqrt{-1}$ on
\[
\overline{\mathfrak g^{1,0}}.
\]
The condition that $\mathfrak g^{1,0}$
is a Lie subalgebra is precisely the
integrability condition.
\end{remark}

\begin{definition}
\label{hodge-birational-atoms-2026-definition-complex2}
A {\it complex nilmanifold} is a
nilmanifold
\[
M=G/\Gamma
\]
equipped with a complex structure such
that the projection
\[
(G,I)\to M
\]
is holomorphic, where $I$ is a
left-invariant complex structure on
$G$.
\end{definition}

\begin{theorem}[Bishop]
\label{hodge-birational-atoms-2026-theorem-let-m3}
Let $(M,I)$ be a compact complex
manifold, and let $\omega$ be a
Hermitian form.
Fix $a\in\mathbb{R}$, and let
\[
\mathcal Z_a
\]
be the set of $k$-cycles
\[
Z\subset M
\]
such that
\[
\int_Z \omega^k\leq a.
\]
Then $\mathcal Z_a$ is compact.
\end{theorem}

\begin{remark}
\label{hodge-birational-atoms-2026-remark-not-clear}
It is not clear at this point why this
is relevant.
\end{remark}

\begin{definition}
\label{hodge-birational-atoms-2026-definition-hermitian-form2}
A Hermitian form $\omega$ on a complex
manifold is {\it SKT} or {\it
pluriclosed} if
\[
dd^c\omega=0.
\]
\end{definition}

\begin{theorem}[Bishop-Gromov]
\label{hodge-birational-atoms-2026-theorem-let-p2}
Let $p\in M$ and $0\leq t\leq i(p)$.
Denote by $B_t(p)$ the ball of radius
$t$ centered at $p$, and by $B_{t,k}$
the ball of radius $t$ in the space of
constant curvature $k$.
If $\operatorname{Ric}\geq k$, then
\[
\frac{\operatorname{Vol}(B_t(p))}
{\operatorname{Vol}(B_{t,k})}
\]
is a nonincreasing function of $t$.
\end{theorem}

\begin{remark}
\label{hodge-birational-atoms-2026-remark-bishop-gromov}
This is the Bishop-Gromov volume
comparison theorem, not Bishop's
compactness theorem for cycles.
\end{remark}

\begin{definition}
\label{hodge-birational-atoms-2026-definition-compact-complex}
A compact complex variety $X$ is
{\it Moishezon} if it is bimeromorphic
to a projective variety.
\end{definition}

Campana's work is foundational in the
birational study of rational curves
and rational connectedness.
The guiding idea is to connect points
of a manifold using rational curves,
or more generally families of curves
of varying type.
The notes also use a theorem about a
holomorphic map
\[
M\to B
\]
with Moishezon fibers and Moishezon
base.
The precise relation between that
theorem and rational connectedness is
left open here.

\begin{definition}
\label{hodge-birational-atoms-2026-definition-hermitian-form2-duplicate-2}
A Hermitian form $\omega$ on a complex
manifold is {\it SKT} or {\it
pluriclosed} if
\[
dd^c\omega=0.
\]
\end{definition}

\begin{definition}
\label{hodge-birational-atoms-2026-definition-nilmanifold-smooth2-duplicate-2}
A {\it nilmanifold} is a smooth
manifold equipped with a transitive
action of a nilpotent Lie group.
\end{definition}

\begin{remark}
\label{hodge-birational-atoms-2026-remark-all-nilmanifolds2-duplicate-2}
All nilmanifolds are obtained as
quotients
\[
M=G/\Gamma,
\]
where $G$ is a nilpotent Lie group
and $\Gamma\subset G$ is a discrete
subgroup.
The Bogomolov--Guan precursor is
particularly easy to study because it
is a complex nilmanifold.
\end{remark}

\begin{definition}
\label{hodge-birational-atoms-2026-definition-let-g4-duplicate-2}
Let $\mathfrak g$ be a real Lie
algebra.
An {\it integrable complex structure}
on $\mathfrak g$ is a subalgebra
\[
\mathfrak g^{1,0}
\subset
\mathfrak g\otimes_{\mathbb{R}}
\mathbb{C}
\]
such that
\[
\mathfrak g^{1,0}\oplus
\overline{\mathfrak g^{1,0}}
=
\mathfrak g\otimes_{\mathbb{R}}
\mathbb{C}.
\]
\end{definition}

\begin{remark}
\label{hodge-birational-atoms-2026-remark-such-decomposition2-duplicate-2}
Such a decomposition is exactly the
same thing as a complex structure on
the Lie algebra.
Indeed, one lets the complex structure
act by $\sqrt{-1}$ on
$\mathfrak g^{1,0}$ and by
$-\sqrt{-1}$ on
$\overline{\mathfrak g^{1,0}}$.
The condition that $\mathfrak g^{1,0}$
is a Lie subalgebra is precisely the
integrability condition.
\end{remark}

\begin{definition}
\label{hodge-birational-atoms-2026-definition-complex2-duplicate-2}
A {\it complex nilmanifold} is a
nilmanifold
\[
M=G/\Gamma
\]
equipped with a complex structure such
that the projection
\[
(G,I)\to M
\]
is holomorphic, where $I$ is a
left-invariant complex structure on
$G$.
\end{definition}

\begin{theorem}
\label{hodge-birational-atoms-2026-theorem-most-countably}
There are at most countably many
quasidiagonals for any given pair
$(E,L)$.
\end{theorem}

\section{Quantum cohomology and
irrationality of
Gushel-Mukai fourfolds}
\label{hodge-birational-atoms-2026-section-quantum-cohomology-gm}

\medskip\noindent
{\bf Speaker.}

Vladimiro Benedetti

\medskip\noindent
{\bf Title.}

Quantum cohomology and irrationality
of Gushel-Mukai fourfolds

\medskip\noindent
{\bf Abstract.}

Gushel-Mukai fourfolds behave very
similarly to cubic fourfolds: they are
Fano manifolds of K3-type, one can
associate to them a (general) IHS
manifold, their rationality is
conjecturally controlled by their
cohomology. In this talk, I will
explain how one can get the
irrationality of very general
Gushel-Mukai fourfolds through the
theory of atoms. This will be done,
as for cubics, by computing the small
quantum cohomology of Gushel-Mukai
fourfolds. Moreover, thanks to a
suitable deformation of the small
quantum cohomology ring, we will also
deduce that a rational Gushel-Mukai
fourfold has the same rational
cohomology as some K3 surface. This
is a joint work with L. Manivel and
N. Perrin.

\medskip\noindent
{\bf AI-generated summary.}

The talk studies rationality of
Gushel--Mukai fourfolds through the
same circle of ideas that appears for
cubic fourfolds: Hodge theory, period
domains, associated K3 surfaces, and
atoms coming from quantum cohomology.
The Hodge-theoretic part describes the
vanishing lattice, the period map, and
Noether--Lefschetz divisors
$\mathcal D_d$; Hodge-special
fourfolds are precisely those whose
periods acquire extra algebraic
classes. This leads to the Hodge
rationality expectation: rational
Gushel--Mukai fourfolds should have
associated K3 surfaces.

The proof strategy uses small quantum
cohomology. Quantum multiplication by
the Euler vector field decomposes into
generalized eigenspaces, and these
pieces behave like atoms under
birational modifications. For a
rational variety, weak factorization
forces such atoms to appear either
from projective space or from the
centers of blow-ups. Computing the
quantum cohomology of a
Gushel--Mukai fourfold and a suitable
deformation of it shows that the
relevant Jordan blocks cannot be
accounted for unless a K3-type piece
appears. Consequently a Hodge general,
in particular very general,
Gushel--Mukai fourfold is irrational;
more precisely, rationality forces
the vanishing cohomology to be the
rational Hodge structure of a
projective K3 surface, up to the
usual shift.

\medskip\noindent
{\bf Notes.}

\begin{definition}
\label{hodge-birational-atoms-2026-definition-variety-x}
A variety $X$ is {\it rational} if
there is a birational map
\[
        X\dashrightarrow \mathbb{P}^n.
\]
Equivalently, there are Zariski open
subsets
\[
        U\subset X,
        \qquad
        V\subset \mathbb{P}^n
\]
which are isomorphic.
\end{definition}

The first question is how rationality
behaves under deformations.
We want to deform the complex
structure and ask whether rationality
is preserved.
Equivalently, one asks for the locus
of rational varieties inside the
moduli space.

\medskip\noindent
{\bf Motivation.}

Conjecture (now theorem) [KKPY].
A very general cubic fourfold is
irrational.

Let us study Fano fourfolds with
\[
H^{(2)}\neq 0.
\]
Some examples are:
\begin{itemize}
\item Birational to cubics.
\item Birational to Gushel--Mukai
fourfolds.
\item Blow-ups with K3 center.
\item Quartic bundles.
\item K\"uc\"u\"kc\"u's examples.
\end{itemize}

\medskip\noindent
{\bf Cubic fourfolds.}

For cubic threefolds the situation is
rigid in the opposite direction: all
of them are irrational.
For cubic fourfolds
\[
        C\subset \mathbb{P}^5
\]
the moduli space $\mathcal M$ has
dimension $20$.
One can ask for the rational locus
\[
        \operatorname{Rat}(C)
        \subset \mathcal M.
\]
This locus is a countable union of
divisors.

The interesting point is that the
Hodge structure controls the complex
structure.
More precisely, the Hodge
decomposition
\[
        H^k(C)=\bigoplus_{p+q=k}
        H^{p,q}(C)
\]
determines the complex structure.
Katzarkov--Kontsevich--Pantev--Yu
used this idea to prove that a very
general cubic fourfold is irrational.

\medskip\noindent
{\bf Gushel--Mukai fourfolds.}

Gushel--Mukai fourfolds are very
similar to cubic fourfolds.
They are the main objects of this
talk.

\begin{definition}
\label{hodge-birational-atoms-2026-definition-gushel-mukai}
A {\it Gushel--Mukai variety} is a
Fano $n$-fold $X$ with
\[
        \operatorname{Pic}(X)=
        \mathbb{Z}H,
        \qquad
        \operatorname{index}(X)=n-2,
        \qquad
        H^n=10.
\]
\end{definition}

In the ordinary case one can write
\[
        X=H\cap Q\cap G(2,V_5),
\]
or equivalently
\[
        X=H\cap Q
        \subset G(2,V_5).
\]

The Hodge structure decomposes as
\[
        H^\bullet(X)
        =
        H^\bullet(X)_{\rm amb}
        \oplus H^4(X)_{\rm van}.
\]
Correspondingly, the Hodge diamond of
$X$ is the sum of the ambient and
vanishing Hodge diamonds.
Note that the Hodge diamond resembles
that of a K3.

\[
\begin{array}{ccccccccc}
&&&&1&&&&\\
&&&0&&0&&&\\
&&0&&1&&0&&\\
&0&&0&&0&&0&\\
0&&1&&22&&1&&0\\
&0&&0&&0&&0&\\
&&0&&1&&0&&\\
&&&0&&0&&&\\
&&&&1&&&&
\end{array}
\]
\[
=
\begin{array}{ccccccccc}
&&&&1&&&&\\
&&&0&&0&&&\\
&&0&&1&&0&&\\
&0&&0&&0&&0&\\
0&&0&&2&&0&&0\\
&0&&0&&0&&0&\\
&&0&&1&&0&&\\
&&&0&&0&&&\\
&&&&1&&&&
\end{array}
\quad
\oplus
\quad
\begin{array}{ccccccccc}
&&&&0&&&&\\
&&&0&&0&&&\\
&&0&&0&&0&&\\
&0&&0&&0&&0&\\
0&&1&&20&&1&&0\\
&0&&0&&0&&0&\\
&&0&&0&&0&&\\
&&&0&&0&&&\\
&&&&0&&&&
\end{array}.
\]

\medskip\noindent
{\bf Periods of Gushel--Mukai
fourfolds.}

Recall the lattice
\[
        \Lambda=
        (H^4(X,\mathbb{Z})_{\rm van},
        \cup),
\]
which is deformation invariant.
The period map is
\[
        p:\mathcal M_{\rm GM}
        \longrightarrow \mathcal P,
        \qquad
        [X]\longmapsto [H^{3,1}(X)].
\]
Here
\[
        \mathcal P
        =
        \left\{
        [\omega]\in
        \mathbb{P}(\Lambda
        \otimes_{\mathbb{Z}}\mathbb{C})
        \mid
        \omega\cup\omega=0,\quad
        \omega\cup\overline{\omega}>0
        \right\}
        /\widetilde O(\Lambda)
\]
is the period domain.

\begin{remark}
\label{hodge-birational-atoms-2026-remark-have-m}
We have
\[
        \dim(\mathcal M_{\rm GM})=24,
        \qquad
        \dim(\mathcal P)=20.
\]
\end{remark}

\begin{remark}[Debarre-Kuznetzov]
\label{hodge-birational-atoms-2026-remark-debarre-kun2}
GM in fibers of $p$ are birational.
\end{remark}

\medskip\noindent
{\bf Hodge classes using period domain.}

The Hodge classes in degree $4$ are
\[
        H^4(X)_{\rm Hdg}
        =
        H^4(X)_{\rm amb}
        \oplus
        \left(
        \Lambda\cap
        (H^{3,1}(X)\oplus
        \overline{H^{3,1}(X)})^\perp
        \right),
\]
where the ambient cohomology is
always algebraic.
The idea is that for very general GM,
$$
H^\bullet(X)^{\text{Hdg}}
=H^\bullet(X)^{\text{amb} }
$$
That is, for a very general
Gushel--Mukai fourfold, the Hodge
classes are exactly the ambient
classes.

If $\alpha\in\Lambda$, then it
defines a divisor
\[
        \mathcal D_d
        =
        \{
        [\omega]\in\mathcal P
        \mid
        \omega\cup\alpha=0
        \}
        \subset \mathcal P.
\]
Here $d$ is the discriminant of
\[
        \alpha^\perp\subset\Lambda.
\]

\begin{definition}
\label{hodge-birational-atoms-2026-definition-gushel-mukai2}
A Gushel--Mukai fourfold is
{\it Hodge-special} if
\[
        p([X])\in
        \bigcup_d \mathcal D_d.
\]
\end{definition}

\noindent
It looks like the next step
is to analyse different values of $d$:

\begin{itemize}
\item For which $d$, $\mathcal{D}_d \neq  0$.
\item For which $d$, there exists
an associated K3 surface: $\alpha^\perp
\simeq H^2(K3, \mathbb{Z})_{\text{prim}}$.
\end{itemize}

\begin{remark}
\label{hodge-birational-atoms-2026-remark-rationality-hodge}
{\bf Rationality, Hodge version.}
Conjecture 1: 
Gushel--Mukai fourfold is rational
if and only if it has a Hodge
associated K3 surface.
\end{remark}

\medskip\noindent
{\bf Evidence.}

The conjecture is known in the cases
\[
        \mathcal D_{10}
        \quad\text{and}\quad
        \mathcal D_{20}.
\]
The first case is due to
Debarre--Iliev--Manivel, and the
second to Hoff and Sagli\'ano.

More evidence is the analogy with cubics:
a cubic is rational if and only if
there exists a Hodge K3 associated.

\begin{remark}
\label{hodge-birational-atoms-2026-remark-associated-k3}
If the associated K3 surface is
projective, then there exists a
non-ambient Hodge class.
Thus $X$ is not Hodge general.
\end{remark}

\medskip\noindent
{\bf Beware.}

For Gushel--Mukai fourfolds there are
two nonequivalent rationality
conjectures.

\begin{remark}
\label{hodge-birational-atoms-2026-remark-rationality}
{\bf Rationality, derived version.}
Conjecture 2: 
a Gushel--Mukai fourfold $X$ is
rational if and only if there exists
a derived associated K3 surface, that
is,
\[
        \mathcal K_X\simeq
        D^b(\operatorname{K3}),
\]
where
\[
        D^b(X)
        =
        \left\langle{}
        \mathcal K_X,
        \mathcal O_X,
        \mathcal U_X^\vee,
        \mathcal O_X(1),
        \mathcal U_X^\vee(1)
        \right\rangle{} .
\]
\end{remark}

The upshot is that the first
conjecture is a divisorial condition,
whereas the second conjecture is a
higher-codimension condition.

\begin{theorem}[Debarre--Iliev--Manivel]
\label{hodge-birational-atoms-2026-theorem-divisor-d}
The divisor $\mathcal D_d$ is nonempty
if and only if the following
conditions hold:
\begin{itemize}
\item If $d\equiv 0\pmod 3$, then
[missing].
\item [missing]
\end{itemize}
\end{theorem}

\begin{theorem}
\label{hodge-birational-atoms-2026-theorem-k3-surface2}
There exists a projective K3 surface
$S$ such that
\[
        \alpha^\perp
        \cong
        (H^2(S,\mathbb{Z})_{\rm prim},
        -\cup)
\]
if and only if either
\[
        d\equiv 2\pmod 8
        \qquad\text{or}\qquad
        d\equiv 4\pmod 8
\]
and every odd prime dividing $d$ is
congruent to $1$ modulo $4$.
\end{theorem}

\begin{remark}
\label{hodge-birational-atoms-2026-remark-k3-surface2}
There was a small additional detail at
the end of the statement which was
missed in the notes.
\end{remark}

\begin{theorem}[Benedetti--Manivel--
\label{hodge-birational-atoms-2026-theorem-perrin-hodge}
Perrin]
\label{hodge-birational-atoms-2026-theorem-bmp2}
A Hodge general Gushel--Mukai fourfold
is not rational.
More precisely, if a Gushel--Mukai
fourfold $X$ is rational, then there
exists a projective K3 surface
$\Sigma$ and an isomorphism of
rational Hodge structures
\[
        H^4(X,\mathbb{Q})_{\rm van}
        \cong
        H^2(\Sigma,\mathbb{Q})(-1).
\]
\end{theorem}

\begin{remark}
\label{hodge-birational-atoms-2026-remark-also-partial}
There are also partial results on the
birational geometry of singular
(nodal) Gushel--Mukai fourfolds due
to Benedetti--Gu\'er\'e--Manivel--
Perrin.
\end{remark}

\begin{remark}
\label{hodge-birational-atoms-2026-remark-number-field}
We can also state the theorem
by using some number field (we don't
know exactly which) instead of $\mathbb{Q}$
(and that $(-1)$ term, I suppose)
in Theorem \ref{hodge-birational-atoms-2026-theorem-bmp2}.
\end{remark}

\medskip\noindent
{\bf Strategy via atoms.}

Consider the small quantum cohomology
ring
\[
        (QH(X),\star,\operatorname{Eu}_0),
        \qquad
        \operatorname{Eu}_0=c_1(X).
\]
One computes the matrix of
\[
        \operatorname{Eu}_0\star
\]
using periods, GKM theory, and
geometry.

Let $E_\lambda$ be the generalized
eigenspace with eigenvalue $\lambda$
which contains $H^{3,1}(X)$.
Set
\[
        \rho_\lambda
        =
        \dim(E_\lambda\cap
        QH^\bullet(X)_{\rm Hdg}),
\]
\[
        \nu_\lambda
        =
        \dim(E_\lambda\cap
        QH^{(2)}(X)),
\]
\[
        \nu'_\lambda
        =
        \dim(E_\lambda\cap
        QH^{(1)}(X)),
\]
and
\[
        \gamma_\lambda
        =
        \operatorname{rank}
        \left.
        (\operatorname{Eu}_0-\lambda)
        \right|_{E_\lambda}.
\]

One can think of small quantum cohomology
just as a ring attached to the variety.
It decomposes as
\[
        QH(X)=\bigoplus_\lambda E_\lambda
\]
into generalized eigenspaces.

These generalized eigenspaces are
preserved under the blow-up and
blow-down processes which link $X$ to
$\mathbb{P}^4$. So if $X$ was rational,
we have something of the kind
\[
\xymatrix@C=1.6em@R=1.8em{
& \widetilde X_0
\ar[dl]_{\pi_0}
\ar[dr]^{\pi'_0}
&& \widetilde X_1
\ar[dl]_{\pi_1}
\ar[dr]^{\pi'_1}
&& \cdots \ar[dr]\ar[dl]
&& \widetilde X_k
\ar[dl]_{\pi_k}
\ar[dr]^{\pi'_k}
\\
X=X_0
\ar@{-->}[rr]^{\varphi_0}
&& X_1
\ar@{-->}[rr]^{\varphi_1}
&& X_2
&& \cdots
\ar@{-->}[rr]^{\varphi_k}
&& X_{k+1}=\mathbb{P}^4 .
}
\]
The generalized eigenspace appears as
an atom of some blow-up center.

Thus the idea is that these atoms
appear for rational varieties and are
preserved by rational manipulations.
More precisely, by Iritani,
$\mathbb{E}_0$ should appear in the quantum
cohomology
of $\mathbb{P}^4$ or in a smooth center.
Studying them can therefore rule out
rationality for some varieties.

\medskip\noindent
{\bf Quantum cohomology setup.}

Let $X$ be a Gushel--Mukai fourfold
with
\[
        \operatorname{Pic}(X)=
        \mathbb{Z}[\mathcal O_X(1)].
\]
We start from
\[
        H(X)=(H^\bullet(X,\mathbb Q),
        \cup)
\]
and deform the cup product to obtain
the small quantum cohomology ring
\[
        QH(X)=(H^\bullet(X)[q],\star).
\]
Here one may use small quantum
cohomology, but in general one has to
work with series.

The product has the form
\[
        \sigma\star\sigma'
        =
        \sigma\cup\sigma'
        +
        \sum_{d\geq 1}q^d
        \sum_{\tau\in H(X)}
        \langle\sigma,\sigma',\tau
        \rangle_d\,\tau^\vee .
\]

Let $M_{3,d}(X)$ be the moduli space
of stable maps
\[
        f:C\to X
\]
from a nodal genus zero curve with
three marked points
\[
        p_1,p_2,p_3,
\]
with finite automorphism group and
\[
        f_*[C]=d[\mathcal O_X(1)]^\vee.
\]
For $i=1,2,3$ there are evaluation
maps
\[
        \operatorname{ev}_i:
        M_{3,d}(X)\to X,
        \qquad
        \operatorname{ev}_i([f])=f(p_i).
\]
The Gromov--Witten invariant is
\[
        \langle\sigma,\sigma',\tau
        \rangle_d
        =
        \int_{[M_{3,d}(X)]^{\rm vir}}
        \operatorname{ev}_1^*\sigma
        \cup
        \operatorname{ev}_2^*\sigma'
        \cup
        \operatorname{ev}_3^*\tau .
\]
Roughly, this counts degree $d$
curves passing through the chosen
cycles.

\begin{remark}
\label{hodge-birational-atoms-2026-remark-missing-point}
There is a missing point here: one has
to prove that this rough counting
interpretation applies in the needed
case.
\end{remark}

\medskip\noindent
{\bf Monodromy action.}

Let
\[
        \mathcal F
        =
        \{
        \text{smooth }
        X\hookrightarrow G(2,V_5)
        \}
\]
be the family of hypersurfaces.
The monodromy action gives an action
of
\[
        \pi_1(\mathcal F)
\]
on
\[
        H^4(X,\mathbb Z)_{\rm van},
\]
and this action is irreducible.

Gromov--Witten invariants are
monodromy, hence deformation,
invariants.

\begin{lemma}
\label{hodge-birational-atoms-2026-lemma-let-1}
Let
\[
        \sigma_1,\ldots,\sigma_k
        \in H^\bullet(X)_{\rm amb}
\]
and let
\[
        \tau\in H^4(X)_{\rm van}.
\]
Then
\[
        \langle
        \sigma_1,\ldots,\sigma_k,\tau
        \rangle_d=0.
\]
\end{lemma}

As consequences, if
\[
        \alpha\in H^\bullet(X)_{\rm amb},
\]
then
\[
        \alpha\star QH^\bullet(X)_{\rm amb}
        \subset QH^\bullet(X)_{\rm amb},
\]
\[
        \alpha\star QH^4(X)_{\rm van}
        \subset QH^4(X)_{\rm van},
\]
and
\[
        \left.
        (\alpha\star)
        \right|_{QH^4(X)_{\rm van}}
        =
        \lambda\,{\rm id}.
\]
These statements hold for any ambient
deformation of $QH(X)$.
The same statements hold for the
Euler vector field.
Thus, with respect to the
decomposition into ambient and
vanishing parts, one has a block form
\[
        \operatorname{Eu}_0\star
        =
        \begin{pmatrix}
        * & 0\\
        0 & \lambda\,{\rm id}
        \end{pmatrix}.
\]
where the first column corresponds
to the ambient cohomology,
while the right column
corresponds to the vanishing cohomology.

For a deformation parameter
$t_\alpha$, the corresponding quantum
Euler vector field has the form
\[
        \operatorname{Eu}_{t_\alpha}
        =
        c_1(X)
        -
        \left(1-
        \frac{\deg(\alpha)}{2}
        \right)t_\alpha .
\]

\medskip\noindent
{\bf Enumerativity.}

For the atom calculation one computes
\[
        \operatorname{Eu}_{t_\alpha}\star.
\]
At $\alpha=0$ one has
\[
        \operatorname{Eu}_0=2h.
\]
Moreover
\[
        \left.
        \ker(\operatorname{Eu}_0\star)
        \right|_{QH^4(X)_{\rm van}}
        =
        0,
\]
because the index is at least $2$.

\begin{proposition}[Enumerativity]
\label{hodge-birational-atoms-2026-proposition-varieties-f}
The varieties $F_1(X)$ and $G_2(X)$
are smooth and irreducible.
For $d=1$ and $d=2$, after using the
$\operatorname{PGL}_5$-action, the
Gromov--Witten invariant
\[
        \langle
        \alpha,\ldots,\alpha
        \rangle_d
\]
is the degree of a suitable
enumerative cycle.
\end{proposition}

\begin{remark}
\label{hodge-birational-atoms-2026-remark-proposition-very}
This proposition is very incomplete in
these notes: the hypotheses and the
degree formula still have to be
filled in.
\end{remark}

\medskip\noindent
{\bf Computing invariants.}

A basis for $H^\bullet(X)_{\rm amb}$
is
\[
        1,\ h,\ \sigma_2,\ \sigma_{11},
        \ I=h^\vee,\ {\rm pt}.
\]
Here
\[
        \sigma_2=[\text{missing}],
        \qquad
        \sigma_{11}=[\text{missing}],
        \qquad
        I=h^\vee=[\text{missing}].
\]
The explicit expressions were missed
in the notes.

\medskip\noindent
{\bf Eigenspace decomposition.}

The point is to compute
\[
        \left.
        \operatorname{Eu}_0\star
        \right|_{H^4(X)_{\rm amb}}.
\]
This gives a matrix whose
characteristic polynomial is
\[
        P(T)
        =
        T^2
        (T^2-44qT-16q^2).
\]
Therefore
\[
        \ker(\operatorname{Eu}_0\star)
        =
        \ker(\operatorname{Eu}_0\star)
        \cap QH^\bullet(X)_{\rm amb}
        \oplus H^4(X)_{\rm van}.
\]
Equivalently,
\[
        \ker(\operatorname{Eu}_0\star)
        =
        \langle\alpha,\beta\rangle
        \oplus H^4(X)_{\rm van}.
\]

\medskip\noindent
{\bf A deformation.}

By deforming $QH(X)$, one finds that
the ambient part becomes a Jordan
block in the matrix of quantum
multiplication.

\begin{remark}
\label{hodge-birational-atoms-2026-remark-details}
The details of this deformation were
missed in the notes.
\end{remark}

\medskip\noindent
{\bf Strategy via atoms, continued.}

\begin{itemize}
\item Consider the small quantum
cohomology ring with its Euler vector
field.
\item Compute the matrix of
$\operatorname{Eu}_0\star$.
\item Study the generalized
eigenspaces.
\item By KKPY and Gu\'er\'e, following
weak factorization, $E_\lambda$
should appear as a generalized
eigenspace of either
$QH_t(\mathbb{P}^4)$ or
$QH_t$ of a smooth center.
\item By KKPY, if
$\rho_\lambda\leq 2$ and
$\nu_\lambda\neq 0$, then
$E_\lambda$ is not a generalized
eigenspace of the quantum cohomology
of a variety of dimension at most $2$,
or of $\mathbb{P}^4$.
This implies irrationality.
\item By Gu\'er\'e, rationality
implies that, if
$\nu_\lambda\neq 0$,
$\nu'_\lambda=0$, and
$\gamma_\lambda\leq 1$, then there
exists a K3 small center and a map
\[
        H^\bullet(\operatorname{K3})
        \longrightarrow H^4(X)_{\rm van}.
\]
That is,
$$
QH^\bullet_t(K3) 
\xrightarrow{\sim} 
\mathbb{E}_0
$$

\item If $\gamma_\lambda=1$, then
separating Jordan blocks gives a
Hodge-theoretic relation from
$H^2(\operatorname{K3})$ to
$H^4(C)_{\rm van}$.
Even if you deform, you get only
one eigenvalue.
\end{itemize}

The upshot is that, under deformation,
one must send a Jordan block to a
Jordan block.

\providecommand{\SU}{\operatorname{SU}}
\providecommand{\G}{\operatorname{G}}
\providecommand{\rG}{\mathrm{G}}

\section{Isometric solutions to the
heterotic G2-system}
\label{hodge-birational-atoms-2026-section-isometric-solutions-heterotic-g2-system}

\medskip\noindent
{\bf Speaker.}

Viviana del Barco

\medskip\noindent
{\bf Title.}

Isometric solutions to the heterotic
G2-system

\medskip\noindent
{\bf Abstract.}

I will report on new solutions to the
heterotic G-system with non-abelian
gauge group, both compact and
non-compact, over certain 3-Sasakian
manifolds and 2-step nilmanifolds.
Moreover, we construct families of
solutions on torus bundles over
manifolds carrying $\SU(3)$-structures
with special torsion forms. Our
methods allow us to produce solutions
to the heterotic $\G_2$-system on a
fixed Riemannian manifold endowed with
non-equivalent $\G_2$-structures. This
is achieved by modifying either the
gauge group of the corresponding
principal bundle or a torsion class of
the $\rG_2$-structure. In particular,
we obtain solutions within the same
isometry class of a $\rG_2$-structure
exhibiting both vanishing and
non-vanishing scalar torsion. This is
a joint work with Andrés J. Moreno and
Udhav Fowdar.

\medskip\noindent
{\bf AI-generated summary.}

The talk constructs new solutions to
the heterotic $G_2$-system, whose data
consist of an integrable
$G_2$-structure, a gauge connection,
and a Bianchi identity relating the
torsion of the characteristic
connection to the curvature of the
gauge field.
The first part recalls the basic
$G_2$-geometry: torsion forms,
integrability, the characteristic
connection, and the instanton and
Bianchi equations.

The main constructions come from two
ansätze.
The first uses manifolds with three
orthogonal vector fields, producing
coclosed $G_2$-structures on
nilmanifolds and on certain
$3$-Sasakian spaces, with compact and
non-compact gauge groups.
The second uses principal
$S^1$-bundles over $6$-manifolds with
$\operatorname{SU}(3)$-structures,
giving families
\[
        \varphi_t
\]
of integrable, generally non-coclosed
$G_2$-structures.
In both methods, the key point is that
one can obtain different heterotic
solutions on the same underlying
Riemannian manifold, often with the
same connection but different torsion
forms.
This explains the phrase
``isometric solutions'' in the title.

\medskip\noindent
{\bf Outline.}

\begin{enumerate}
\item Preliminaries on $G_2$-structures
and the heterotic $G_2$-system.
\item Previous solutions to the
system.
\item New solutions by two different
methods.
\end{enumerate}

\medskip\noindent
{\bf References.}

\begin{itemize}
\item {\it $G_2$-instantons on
2-step nilpotent Lie groups}.
\item {\it Isometric solutions to the
heterotic $G_2$-system}.
\end{itemize}

\medskip\noindent
{\bf Preliminaries on $G_2$-structures.}

\begin{definition}
\label{hodge-birational-atoms-2026-definition-g-2}
A {\it $G_2$-manifold} is a
$7$-dimensional manifold $M$ equipped
with a $G_2$-structure.
Such a structure is a $3$-form
\[
        \varphi\in\Omega^3(M)
\]
such that, at every point $p\in M$,
\[
        \iota_v\varphi\wedge
        \iota_v\varphi\wedge\varphi
        \neq 0
\]
for every nonzero
\[
        v\in T_pM.
\]
\end{definition}

Such a structure defines a Riemannian
metric $g_\varphi$ and an orientation
by the formula
\[
        g_\varphi(v,w)\,
        \operatorname{vol}_\varphi
        =
        \frac{1}{6}
        \iota_v\varphi\wedge
        \iota_w\varphi\wedge\varphi .
\]

\medskip\noindent
{\bf $G_2$-module decompositions.}

A $G_2$-structure gives
decompositions
\[
        \Lambda^2(M)
        =
        \Lambda^2_7
        \oplus
        \Lambda^2_{14},
\]
and
\[
        \Lambda^3(M)
        =
        \langle\varphi\rangle
        \oplus
        \Lambda^3_7
        \oplus
        \Lambda^3_{27}.
\]
Each summand has an explicit
description.

\medskip\noindent
{\bf Torsion forms.}

The decomposition allows one to define
the torsion forms
\[
        \tau_0\in\, C^\infty(M),
        \qquad
        \tau_1\in\Lambda^1(M),
\]
\[
        \tau_2\in\Lambda^2_{14},
        \qquad
        \tau_3\in\Lambda^3_{27}.
\]
They are defined by the components of
$d\varphi$ and $d\psi$, where
\[
        \psi=*_\varphi\varphi,
\]
with respect to the $G_2$-module
decompositions:
\[
        d\varphi
        =
        \tau_0\psi
        +
        3\tau_1\wedge\varphi
        +
        *_\varphi\tau_3,
\]
and
\[
        d\psi
        =
        4\tau_1\wedge\psi
        +
        \tau_2\wedge\varphi .
\]

The $G_2$-structure $\varphi$ is
called:
\begin{itemize}
\item integrable if
\[
        d\psi=4\tau_1\wedge\psi;
\]
\item coclosed if [missing];
\item purely coclosed if [missing];
\item nearly parallel if [missing];
\item parallel if [missing].
\end{itemize}

\begin{theorem}[Friedrich--Ivanov, 2002]
\label{hodge-birational-atoms-2026-theorem-being-integrable}
Being integrable is equivalent to the
existence of a metric connection on
$M$ with skew-symmetric torsion such
that
\[
        \nabla^c\varphi=0.
\]
\end{theorem}

\medskip\noindent
{\bf Heterotic $G_2$-system.}

Let $\varphi$ be an integrable
$G_2$-structure on a differentiable
manifold $M$.
Consider:
\begin{itemize}
\item a principal $G$-bundle
\[
        P\to M;
\]
\item a connection $1$-form $A$;
\item an
$\operatorname{Ad}(G)$-invariant
nondegenerate symmetric bilinear form
\[
        \langle\ ,\ \rangle_{\mathfrak g}
\]
on
\[
        \mathfrak g=\operatorname{Lie}(G).
\]
\end{itemize}

\begin{definition}
\label{hodge-birational-atoms-2026-definition-say-satisfies}
We say that $(\varphi,A)$ satisfies
the {\it heterotic $G_2$-system}
in the sense of da Silva et al.\
2025
if the following system of equations
holds:
\[
        d\psi=4\tau_1\wedge\psi,
\]
\[
        F_A\wedge\psi=0,
\]
and
\[
        dT_\varphi
        =
        \langle F_A\wedge F_A
        \rangle_{\mathfrak g}.
\]
The first equation is the integrability
condition, the second is the
instanton condition, and the third is
the heterotic Bianchi identity.
\end{definition}

Here
\[
        T_\varphi
        =
        \frac{1}{6}\tau_0\varphi
        +
        *_\varphi(\tau_1\wedge\varphi)
        -
        \tau_3
\]
is the torsion of the canonical
connection.

\begin{remark}
\label{hodge-birational-atoms-2026-remark-particular-case}
This is a particular case of a
heterotic system.
There is a more general way to define
such systems.
\end{remark}

\medskip\noindent
{\bf Previous approaches.}

Given a manifold with an integrable
$G_2$-structure $(M,\varphi)$, it is
not clear in general which bundles,
connections, and
$\operatorname{Ad}$-invariant metrics
define solutions to the heterotic
$G_2$-system.
The following list is non-exhaustive.

\begin{itemize}
\item Harland--N\"olle, 2012:
the canonical connection $\nabla^c$
is an instanton when $\varphi$ is
nearly parallel.
\item Ivanov--Ivanov, 2005:
solutions on
\[
        M=S^1\times \Gamma/N_{3,2}.
\]
\item Fern\'andez et al., 2011:
solutions on nilmanifolds
\[
        \Gamma/\mathbb H_7,
\]
where $\mathbb H_7$ is the
$7$-dimensional Heisenberg group.
\item Moroianu et al., 2025:
solutions on $2$-step nilmanifolds.
\item Lotay--Sa Earp, 2023:
approximate solutions to the system
on contact Calabi--Yau manifolds.
\item Galdeano--Stecker, 2025:
approximate solutions on
$3$-$(\alpha,\delta)$-Sasakian
manifolds, with an exact solution
occurring in a degenerate case.
\end{itemize}

\medskip\noindent
{\bf Part 2: nilmanifold solutions.}

In joint work with Clarke and Moreno,
we considered the following setup:
\begin{itemize}
\item
\[
        M=\Gamma\backslash N,
\]
where $N$ is a $2$-step nilpotent Lie
group of dimension $7$ and
$\Gamma\subset N$ is a discrete
cocompact subgroup;
\item $\varphi$ is a coclosed
left-invariant $G_2$-structure on
$N$;
\item $A=\nabla^c$ is the
characteristic connection.
\end{itemize}

First we aimed at solving the
instanton condition for $\nabla^c$.
We obtained the following theorem.

\begin{theorem}
\label{hodge-birational-atoms-2026-theorem-exists-g}
There exists a $G_2$-structure
$\varphi$ on
\[
        M=\Gamma\backslash N
\]
for which the characteristic
connection is a $G_2$-instanton if
and only if the Lie algebra
$\mathfrak n$ is one of the following:
\[
        \mathfrak n=\mathfrak h_7,
        \qquad
        \mathfrak n=\mathbb R
        \oplus\mathfrak n_{3,2},
        \qquad
        \mathfrak n=\mathfrak h_{\mathbb H}.
\]
\end{theorem}

In each case one can check that
$\nabla^c$ also solves the heterotic
Bianchi identity.
Thus each of the above gives a
solution to the heterotic
$G_2$-system.
The gauge groups are
\[
        \operatorname{SU}(3)
        \subset G_2.
\]

The Riemannian Lie groups obtained
are naturally reductive, so
\[
        \nabla^cT=0,
        \qquad
        \nabla^cR=0.
\]
The torsion is not closed.

\begin{remark}
\label{hodge-birational-atoms-2026-remark-solutions-differ}
These solutions differ from the ones
studied by Moroianu et al.\ in 2025.
\end{remark}

\begin{remark}
\label{hodge-birational-atoms-2026-remark-solution}
The solution corresponding to the
first Lie algebra was already known
from Fern\'andez et al.\ in 2011.
\end{remark}

\begin{remark}
\label{hodge-birational-atoms-2026-remark-solution2}
The solution corresponding to the
second Lie algebra is different from
the one in Ivanov--Ivanov.
\end{remark}

\medskip\noindent
{\bf Part 3: new solutions by two
methods.}

\medskip\noindent
{\bf First ansatz.}

\begin{itemize}
\item This applies to manifolds with a
triple of orthogonal vector fields.
\item It produces examples with both
\[
        \tau_0=0
        \qquad\text{and}\qquad
        \tau_0\neq 0.
\]
\item It applies to nilmanifolds and
$3$-Sasakian manifolds.
\item It gives coclosed
$G_2$-structures, so
\[
        \tau_1=\tau_2=0.
\]
\item It provides compact and
non-compact structure groups $G$.
\end{itemize}

\medskip\noindent
{\bf Second ansatz.}

\begin{itemize}
\item This applies to manifolds which
are $S^1$-bundles over manifolds with
$\operatorname{SU}(3)$-structures.
\item It gives integrable solutions,
which are non-coclosed in general:
\[
        \tau_1\neq 0,
        \qquad
        \tau_2=0.
\]
\item Again one obtains both cases
\[
        \tau_0=0
        \qquad\text{and}\qquad
        \tau_0\neq 0.
\]
\item There are $1$-parameter families
of solutions over the same Riemannian
manifold, with nonequivalent
$G_2$-structures and fixed connection.
\end{itemize}

The main feature in both cases is that
the solutions are isometric.

\medskip\noindent
{\bf First ansatz: the manifold and
the $G_2$-structure.}

Let $(M,g)$ be an oriented Riemannian
spin manifold.
Assume that $(M,g)$ admits a triple of
orthogonal vector fields
\[
        e_5,e_6,e_7.
\]
Then
\[
        \varphi
        =
        \sigma^+_1\wedge e^5
        +
        \sigma^+_2\wedge e^6
        +
        \sigma^+_3\wedge e^7
        +
        e^{567}
\]
is an $\operatorname{SO}(3)$-family
of $G_2$-structures, where
\[
        \sigma^+_1,\sigma^+_2,
        \sigma^+_3
\]
are self-dual $2$-forms on the
distribution [missing].

\medskip\noindent
{\bf First ansatz: the gauge group and
connections.}

Let $G$ be a $3$-dimensional Lie
group.
[missing]

\medskip\noindent
{\bf First ansatz: solutions on
nilmanifolds.}

Assume
\[
        M=\Gamma\backslash N,
        \qquad
        \mathfrak n=\operatorname{Lie}(N),
\]
with a global coframe
\[
        \{e^i\}_{i=1}^7
\]
satisfying the structure equation
[missing].

Any Lie algebra satisfying that
equation with $\delta\neq 0$ is
$2$-step nilpotent and isomorphic to
one of
\[
        \mathbb R\oplus\mathfrak n_{3,2},
        \qquad
        \mathfrak n_{7,3,A},
\]
\[
        \mathfrak n_{7,3,B_1},
        \qquad
        \mathfrak n_{7,3,D},
        \qquad
        \mathfrak h_{\mathbb H}.
\]

\medskip\noindent
{\bf First ansatz: results.}

\begin{theorem}
\label{hodge-birational-atoms-2026-theorem-let-m4}
Let
\[
        M=\Gamma\backslash N
\]
be a nilmanifold, where $\Gamma$ is a
cocompact lattice.
There exists a coclosed
$G_2$-structure $\varphi$ on $M$ and
a $G$-connection $A$ solving the
heterotic $G_2$-system in the
following cases:
\begin{enumerate}
\item
\[
        G=\mathbb T^3
\]
and
\[
        \mathfrak n\cong
        \mathbb R^2\oplus\mathfrak h_5,
        \quad
        \mathbb R\oplus
        \mathfrak h_3^{\mathbb C},
        \quad
        \mathfrak h_{\mathbb H}.
\]
\item
\[
        G=\operatorname{SU}(2).
\]
[missing]
\item
\[
        G=\operatorname{SL}(2,\mathbb R).
\]
[missing]
\end{enumerate}
Moreover, in the nonabelian cases,
both the underlying metric and volume
are the same, but the $G_2$-structures
$\varphi$ are distinct.
\end{theorem}

Let $(M,g)$ be a $3$-Sasakian
manifold.
We denote by $\varphi_{\rm ts}$ the
$3$-Sasakian nearly parallel
$G_2$-structure, by
$\varphi_{\rm np}$ the squashed
nearly parallel $G_2$-structure, and
by $\widehat\varphi_{\rm ts}$ a
coclosed, but not nearly parallel,
$G_2$-structure related to
$\varphi_{\rm ts}$.

\begin{theorem}
\label{hodge-birational-atoms-2026-theorem-let-m5}
Let $M$ be either $S^7$ or the
Aloff--Wallach space $N^{1,1}$.
For the $G_2$-structure
$\varphi_{\rm ts}$, there exist
$G$-connections solving the heterotic
$G_2$-system in the following cases:
\begin{enumerate}
\item
\[
        G=\operatorname{SL}(2,\mathbb R)
        \times\operatorname{SU}(3)
        \times G_2
\]
and $M=S^7$.
\item
\[
        G=\operatorname{SL}(2,\mathbb R)
        \times \operatorname{U}(1)
\]
and $M=N^{1,1}$.
\end{enumerate}
[something else missing]
\end{theorem}

\medskip\noindent
{\bf Second ansatz: the manifold and
$G_2$-structure.}

Let $M$ be a principal $S^1$-bundle
over a $6$-manifold $Q$ with an
$\operatorname{SU}(3)$-structure
\[
        (Q,g_\omega,J,\omega,\gamma).
\]
Let $\eta$ be a connection $1$-form
on $M$.
Consider the family of
$G_2$-structures on $M$
\[
        \varphi_t
        =
        \eta\wedge\omega
        +
        \operatorname{Re}
        \left(
        e^{it}(\gamma_+
        + i\gamma_-)
        \right).
\]

Analogously to the $G_2$ case, the
spaces $\Lambda^\bullet(Q)$ decompose
as $\operatorname{SU}(3)$-modules.
For example,
\[
        \Lambda^2(Q)
        =
        \langle\omega\rangle
        \oplus
        \Lambda^2_6
        \oplus
        \Lambda^2_8
        \oplus\cdots .
\]
More of this decomposition is missing
from the notes.

\begin{proposition}
\label{hodge-birational-atoms-2026-proposition-g-2}
The $G_2$-structure $\varphi_t$ is
integrable for all
\[
        t\in[0,2\pi)
\]
if and only if
\[
        \sigma_2=\pi_2=0,
        \qquad
        \pi_1=2\nu_1,
\]
and $d\eta$ is of type $(1,1)$.
In this case one has formulas for
\[
        T_{\varphi_t},
        \qquad
        \tau_0(t),
        \qquad
        \tau_1(t),
\]
which were missed in the notes.
\end{proposition}

Here $(d\eta)_0$ denotes the
$\omega$-component of $d\eta$, and
$T_\omega$ is the torsion form of the
Bismut connection of
\[
        (\omega_+,\gamma_+).
\]

\begin{remark}
\label{hodge-birational-atoms-2026-remark-more-details}
More details of this proposition are
missing.
\end{remark}

\begin{theorem}
\label{hodge-birational-atoms-2026-theorem-let-m6}
Let $M$ be a principal $S^1$-bundle
over
\[
        (Q,\omega,\gamma_+).
\]
Assume that $d\eta$ and the torsion
forms of the $\operatorname{SU}(3)$-
structure are such that
$T_{\varphi_t}$ is constant.
Assume also that
\[
        (\omega,\gamma_+,A)
\]
solves the $\operatorname{SU}(3)$
heterotic Bianchi identity, and
[missing].
Then the pullback connection
\[
        \eta\oplus A
\]
is a solution of the heterotic
$G_2$-system for every
\[
        t\in[0,2\pi).
\]
\end{theorem}

\begin{example}
\label{hodge-birational-atoms-2026-example-let-n2}
Let
\[
        \mathfrak n_{3,2}
        =
        \langle e_2,\ldots,e_7\rangle
\]
with structure equations
\[
        de^i=0
        \quad
        \text{for } i=2,3,4,
\]
and [missing].

There is a left-invariant
$\operatorname{SU}(3)$-structure
induced on $N_{3,2}$, given by
\[
        \omega=[\text{missing}],
        \qquad
        \gamma_+=[\text{missing}].
\]
On the manifold
\[
        M=S^1\times
        \Gamma\backslash N_{3,2}
\]
we take
\[
        \eta=e^1
\]
and [missing].
The theorem above gives a
$1$-parameter family of solutions to
the heterotic $G_2$-system on $M$,
with the same connection $A$ but with
varying $\tau_0$ and $\tau_1$.
In particular, these structures are
not coclosed in general.

The family $\varphi_t$ contains both a
coclosed $G_2$-structure, for $t=0$,
and integrable but not coclosed
$G_2$-structures for other values of
$t$.
\end{example}

\section{Moduli spaces of pointed
genus one curves and logarithms of
algebraic numbers}
\label{hodge-birational-atoms-2026-section-pointed-genus-one-curves-logarithms-algebraic-numbers}

\medskip\noindent
{\bf Speaker.}

Tiago Fonseca

\medskip\noindent
{\bf Title.}

Moduli spaces of pointed genus one
curves and logarithms of algebraic
numbers

\medskip\noindent
{\bf Abstract.}

This is joint work with Francis Brown.
I will explain the relation between
motives of moduli spaces of genus one
curves with marked points and the
so-called Gross-Zagier conjecture,
according to which certain special
values of automorphic Green's
functions are logarithms of algebraic
numbers.

\medskip\noindent
{\bf Goals of the talk.}

\begin{itemize}
\item Recall the Gross--Zagier
log-algebraicity conjecture.
\item Relate it to the moduli spaces
\[
        \mathcal M_{1,n}.
\]
\item Discuss geometric and motivic
aspects.
\end{itemize}

\medskip\noindent
{\bf Notation.}

\begin{itemize}
\item $\mathbb H$ denotes the upper
half-plane.
Recall that $\operatorname{SL}_2
(\mathbb Z)$ acts on $\mathbb H$ by
M\"obius transformations.
\item Klein's $J$-invariant is a map
\[
        J:\mathbb H\to\mathbb C
\]
which induces an isomorphism
\[
        \operatorname{SL}_2(\mathbb Z)
        \backslash\mathbb H
        \cong
        \mathbb C.
\]
\item The moduli space
$\mathcal M_{1,1}$ parametrizes genus
$1$ curves with one marked point, and
the moduli interpretation is
\[
        \operatorname{SL}_2(\mathbb Z)
        \backslash\mathbb H
        \cong
        \mathcal M_{1,1},
\]
via
\[
        z\longmapsto
        (E_z=\mathbb C/(\mathbb Z
        +\mathbb Zz),[0]).
\]
\end{itemize}

\medskip\noindent
{\bf AI-generated note on $J$.}

Klein's $J$-function is the classical
modular function which classifies
elliptic curves over $\mathbb C$ up
to isomorphism.
It is invariant under the action of
\[
        \operatorname{SL}_2(\mathbb Z)
\]
on the upper half-plane, so two
points $z,z'\in\mathbb H$ give the
same value of $J$ exactly when the
corresponding marked complex tori are
isomorphic.
Thus $J$ gives a coordinate on the
coarse moduli space of elliptic
curves.
In the more standard normalization,
one often writes
\[
        j(z)=q^{-1}+744+196884q+\cdots,
        \qquad
        q=e^{2\pi iz},
\]
and Klein's function is
\[
        J=j/1728.
\]

\medskip\noindent
{\bf Arithmetic of $J$.}

The function $J$ satisfies many
interesting arithmetic properties.
One of them is that this highly
transcendental function sends
quadratic numbers to algebraic
numbers.

\begin{theorem}[Kronecker]
\label{hodge-birational-atoms-2026-theorem-z-h}
If $z\in\mathbb H$ is a quadratic
number, then
\[
        J(z)\in\overline{\mathbb Q}.
\]
\end{theorem}

\begin{example}[Weber]
\label{hodge-birational-atoms-2026-example-one-has}
One has
\[
        J(\sqrt{-14})
        =
        6^{-3}
        \left(
        323+228\sqrt 2
        +(231+161\sqrt 2)
        \sqrt{2\sqrt 2-1}
        \right)^3.
\]
\end{example}

Geometrically, this says that an
elliptic curve with too many
endomorphisms, meaning
\[
        \operatorname{End}(E)
        \supset \mathbb Z,\quad
        \operatorname{End}(E)\not =
        \mathbb Z,
\]
is defined over $\overline{\mathbb Q}$.
We say that such a point
$z\in\mathbb H$ is a {\it CM point},
where CM stands for complex
multiplication.

\begin{quote}
The theory of complex multiplication,
which forms a powerful link between
number theory and analysis, is not
only the most beautiful part of
mathematics but also of all science.

\hfill D. Hilbert
\end{quote}

\medskip\noindent
{\bf Gross--Zagier Green functions.}

In the 1980s, Gross and Zagier
published
{\it Heegner points and derivatives of
$L$-series}.
This proves a key case of the
Birch--Swinnerton-Dyer conjecture.
The proof is interesting here because
it involves a long explicit
calculation with the higher Green
functions
\[
        G_s(z,w)
        =
        -s
        \sum_{\gamma\in
        \operatorname{SL}_2(\mathbb Z)}
        Q_{s-1}
        \left(
        1+
        \frac{|z-\gamma w|^2}
        {2\operatorname{Im}(z)
        \operatorname{Im}(\gamma w)}
        \right),
\]
where $z,w\in\mathbb H$ are not in
the same orbit, and $s\geq 1$ is a
parameter.

These are higher versions of
\[
        G_1(z,w)=\log|J(z)-J(w)|^2,
\]
but their geometric interpretation is
less clear.
Thus the function appears in the
proof, but it is not immediately clear
what it means geometrically.

When $s\geq 1$ is an integer, $G_s$
has remarkable arithmetic properties.

\begin{itemize}
\item The Gross--Zagier conjecture is
known for
\[
        s\in\{1,2,3,4,5,7\}.
\]
If $z$ and $w$ are CM points of
discriminants $d_z$ and $d_w$, then
\[
        G_s(z,w)
        =
        (d_zd_w)^{(s-1)/2}
        \log|\alpha|
\]
for some
\[
        \alpha\in\overline{\mathbb Q}.
\]
There is a more complicated version
for other values of $s$.
\item For
\[
        G_1(z,w)=\log|J(z)-J(w)|^2,
\]
this follows from Kronecker's theorem.
\item Mellit computed the example
\[
        G_2\left(
        \frac{1+\sqrt{-7}}{2},i
        \right)
        =
        \frac{8}{\sqrt 7}
        \log(8-3\sqrt 7).
\]
\item Previous work includes
Gross--Kohnen--Zagier 1987,
Zhang 1997, Mellit 2008,
Viazovska 2011, Li 2023, and
Bruinier--Li--Yang 2025.
\end{itemize}

The goal is to relate the
Gross--Zagier conjecture to algebraic
geometry.

\medskip\noindent
{\bf Differential forms on
$M_{1,n}$.}

The case $s=1$ is already suggestive.
Since
\[
        G_1(z,w)=\log|J(z)-J(w)|^2
\]
is a logarithm, it can be expressed
in terms of integrals. Equivalently,
with one variable fixed, one may think
of it as a primitive
\[
        \nu+\overline{\nu},
\]
where $\nu$ is a $1$-form on
$M_{1,1}$ with poles along the
corresponding value of $J$.

Let $\overline M_{1,n}$ denote the
Deligne--Mumford compactification,
whose points are stable curves of
genus one with $n$ marked points.
There is an action of $S_n$ on
$\overline M_{1,n}$ by permuting the
marked points.

\begin{theorem}[Brown--Fonseca]
\label{hodge-birational-atoms-2026-theorem-unique-real}
There are unique real-analytic
$(n-1)$-forms $g^{p,q}$ on
\[
        \overline M_{1,n}\setminus D_w
\]
such that
\[
        d g^{p,q}
        =
        \nu^{p,q}
        +
        \overline{\nu^{q,p}}.
\]
The forms $\nu^{p,q}$ define algebraic
$n$-forms on
\[
        \overline M_{1,n}\setminus D_w.
\]
Their cohomology classes are
$S_n$-alternating and satisfy
\[
        H^n(
        \overline M_{1,n}\setminus D_w
        )_{\operatorname{sgn}}
        =
        H^n(\overline M_{1,n})
        _{\operatorname{sgn}}
        \oplus
        \bigoplus_{p+q=n-1}
        \mathbb C[\nu^{p,q}].
\]
\end{theorem}

To summarize, for $n\geq 3$ one gets
an $n\times n$ matrix of Green
functions.
Its entries are coefficients of
$(n-1)$-forms $g^{p,q}$ such that
\[
        d g^{p,q}
        =
        \nu^{p,q}
        +
        \overline{\nu^{q,p}},
\]
where the $\nu^{p,q}$ are algebraic
$n$-forms.

\medskip\noindent
{\bf Motives.}

Let $K\subset\mathbb C$ be a number
field.
Many cohomology theories with constant
coefficients for algebraic varieties
over $K$ include:

\[
\begin{array}{c|c|c}
\text{Betti} &
\text{de Rham} &
\ell\text{-adic} \\
H^\bullet_{\mathrm{B}}(X) &
H^\bullet_{\mathrm{dR}}(X) &
H^\bullet_\ell(X) \\
H^\bullet_{\mathrm{sing}}
(X(\mathbb C);\mathbb Q)
&
H^\bullet_{\mathrm{Zar}}
(X;\Omega^\bullet_{X/K})
&
\varprojlim_n
H^\bullet_{\mathrm{et}}
(X,\mathbb Z/\ell^n)
\otimes \mathbb Q_\ell
\\
\mathbb Q &
K &
\mathbb Q_\ell \\
{\rm complex~conj.} &
\text{Hodge filtration} &
\text{Galois rep.}
\end{array}
\]

Grothendieck's comparison theorem
shows that the coefficients are
periods:
\[
        H^\bullet_{\mathrm{dR}}(X)
        \otimes_K \mathbb C
        \xrightarrow{\sim}
        H^\bullet_{\mathrm{B}}(X)
        \otimes_{\mathbb Q} \mathbb C,
\]
\[
        \omega \mapsto
        \left(
        [\sigma] \mapsto
        \int_\sigma \omega
        \right).
\]
These integrals are called periods.
The motive $H(X)$ is a linear algebra
object containing all cohomological
information of $X$ for all Weil
cohomology theories.

\medskip\noindent
{\bf Voevodsky motives.}

\begin{itemize}
\item $DM(K)$ is Voevodsky's
triangulated category of motives over
$K$.
Morally, it is the bounded derived
category of the true abelian category
of motives over $K$.
\item $DM(K)$ is $\mathbb Q$-linear,
triangulated, pseudo-abelian, and
rigid tensor.
\item A subvariety $Y\subset X$
determines a relative cohomology motive
\[
        H(X,Y)\in DM(K).
\]
We write $H(X,\emptyset)=H(X)$.
\item There is a distinguished triangle
\[
        H(Y)[-1]
        \to
        H(X,Y)
        \to
        H(X)
        \xrightarrow{+1}
\]
This is the relative cohomology long
exact sequence.
\item There are realization functors
\[
        R_{\mathrm{B}}:DM(K)\to
        D\mathrm{Vect}_{\mathbb Q},
        \qquad
        R_{\mathrm{dR}}:DM(K)\to
        D\mathrm{Vect}_K.
\]
For example,
\[
        R_{\mathrm{B}}(X,Y)
        =
        H^\bullet_{\mathrm{B}}(X,Y).
\]
\item Grothendieck's comparison
theorem can be rephrased as
\[
        R_{\mathrm{dR}}
        \otimes_K \mathbb C
        \cong
        R_{\mathrm{B}}
        \otimes_{\mathbb Q} \mathbb C.
\]
\end{itemize}

The basic motives are the building
blocks for all such cohomological
realizations.

\medskip\noindent
{\bf Basic motives.}

\begin{itemize}
\item The trivial motive is
\[
        \mathbb Q(0)=H(\operatorname{pt}).
\]
\item The Lefschetz motive is
\[
        \mathbb Q(-1)
        =
        H(\mathbb A^1\setminus\{0\},
        \{1\})[1].
\]
\item More generally,
\[
        \mathbb Q(n)
        =
        \mathbb Q(-1)^{\otimes -n},
        \qquad n\in\mathbb Z.
\]
\item The Tate twist is
\[
        M(n)=M\otimes \mathbb Q(n).
\]
\item For $\alpha\in K^\times$, the
Kummer motive is
\[
        \mathcal K_\alpha
        =
        H(\mathbb A^1\setminus\{0\},
        \{1,\alpha\})[1].
\]
\end{itemize}

There is a distinguished triangle
\[
        \mathbb Q(0)
        \to
        \mathcal K_\alpha
        \to
        \mathbb Q(-1)
        \xrightarrow{+1}
\]
The periods of these Kummer motives
are logarithms.
This is one way to think about
logarithms in algebraic geometry.

If you want to know whether a number
is a logarithm, you attach a motive to
it and look for the corresponding
extension of Tate motives.
In that way, the logarithm is read off
from the motive's mixed structure.

\begin{theorem}[Brown--Fonseca]
\label{hodge-birational-atoms-2026-theorem-let-n}
Let $n\geq 3$.
\begin{enumerate}
\item The functions
$G_{r,s}^{p,q}(z,w)$ are single-valued
periods of a motive
$M_{n,z,w}$.
\item The motive $F_{n,z}$ is a
subquotient of
$\operatorname{Sym}^{n-1}H^1(E_z)$.
If $z$ is a CM point, then
\[
        F_{n,z}
        =
        \mathbb Q\!\left(
        \frac{n-1}{2}
        \right)
        \oplus
        \widetilde F_{n,z}.
\]
\item If $z$ and $w$ are CM points and
\[
        H(\overline M_{1,n})
        _{\operatorname{sgn}}[1]
        =0,
\]
then one gets a subquotient
\[
        \mathbb Q\!\left(\frac{n-1}{2}\right)
        \to
        \mathcal{GZ}_{n,z,w}
        \to
        \mathbb Q\!\left(\frac{n-3}{2}\right)
        \xrightarrow{+1}
\]
\end{enumerate}
If $n=2s-1$, then
$G_{r,s}^{p,q}(z,w)$ is the
single-valued period of
$\mathcal{GZ}_{n,z,w}$.
The Gross--Zagier conjecture holds
for the corresponding CM points.
\end{theorem}

We know that
\[
        R_{\mathrm{B}}
        H(\overline M_{1,n})
        _{\operatorname{sgn}}[n]
        =0
\]
for
\[
        s\in\{1,2,3,4,5,6,7\}.
\]

\begin{remark}
\label{hodge-birational-atoms-2026-remark-functor-r}
\noindent{\bf Conservativity.}
The functor
\[
        R_{\mathrm{B}}:DM(K)\to
        D\mathrm{Vect}_{\mathbb Q}
\]
is conservative.
\end{remark}

This would imply the Gross--Zagier
conjecture.
We do not know how to prove this in
general.
The vanishing of the motives above is
not out of reach.
For instance, conservativity is known
on the subcategory of mixed Tate
motives.
So this is a completely geometric
proof of the conjecture in that case.
It works because mixed Tate motives
are conservative.

\begin{theorem}[Brown--Fonseca]
\label{hodge-birational-atoms-2026-theorem-motive-m}
The motive of $\overline M_{1,3}$ is
mixed Tate.
In particular, the above argument
proves the Gross--Zagier conjecture
for $n=3$, that is, for $G_2(z,w)$.
\end{theorem}

We also expect the mixed Tate property
for all $n\leq 10$.
Compare with Belorousski's rationality
result.

\section{Atomic Semi-Orthogonal
Decompositions for Derived Categories
of Surfaces}
\label{hodge-birational-atoms-2026-section-atomic-sod-surfaces}

\medskip\noindent
{\bf Speaker.}

Alexey Elagin

\medskip\noindent
{\bf Title.}

Atomic Semi-Orthogonal
Decompositions for Derived Categories
of Surfaces.

\medskip\noindent
{\bf Abstract.}

I will talk about our recent work
with Julia Schneider and Evgeny
Shinder.
We construct, for any surface with a
group action, a canonical "atomic"
$G$-invariant semi-orthogonal
decomposition of the derived category,
so that these decompositions are
compatible with geometric operations
such as blow-ups and conic bundles.
Non-trivial "atoms" of these
decompositions are birational
invariants, this enables a uniform
approach to birational classification
of $G$-surfaces in terms of atoms.
The construction of atomic
decompositions is based on the
$G$-minimal model program in dimension
two and on the Sarkisov link
decomposition.
Our work is inspired by the atoms by
Katzarkov-Kontsevich-Pantev-Yu,
although the approach is more
elementary and makes no use of quantum
cohomology.

\medskip\noindent
{\bf AI-generated note.}

Elagin uses Sarkisov's program in the
same structural way as the Cremona
notes do in Section
\ref{hodge-birational-atoms-2026-section-mmp-and-sarkisov}:
factor a birational map into
elementary links, then read off the
effect of each link on the derived
category.
For surfaces with group action, this
lets him build the atomic
semi-orthogonal decomposition
compatibly with blowups and conic
bundles, and then show that the
nontrivial atoms are birational
invariants.

\medskip\noindent
Let $X$ be a variety over a field
$\mathbb k$.
We write
\[
D^b(X):=D^b(\operatorname{Coh}(X)).
\]
Then $D^b(X)$ is a $k$-linear
triangulated category.

\medskip\noindent
Triangulated categories are the
derived-category version of exact
categories.
Triangles replace extensions, kernels,
and cokernels.

\medskip\noindent
If $X$ and $Y$ are isomorphic, then
$D^b(X)\simeq D^b(Y)$.
The converse is false in general.
By Bondal--Orlov, if $X$ and $Y$ are
smooth projective and $K_X$ or $-K_X$
is ample, then an equivalence
$D^b(X)\simeq D^b(Y)$ implies
$X\simeq Y$.

\begin{definition}
\label{hodge-birational-atoms-2026-definition-semi-orthogonal}
A semi-orthogonal decomposition
$D^b(X)=\langle A_1,\dots,A_n\rangle$
means that
$\operatorname{Hom}(A_j,A_i)=0$
for $j>i$, and the $A_i$ generate
$D^b(X)$.
\end{definition}

\begin{example}[Beilinson]
\label{hodge-birational-atoms-2026-example-sod-d}
There is an SOD
\[
D^b(\mathbb P^n)
=\langle
\mathcal O(-n),\dots,
\mathcal O(-1),\mathcal O
\rangle.
\]
\end{example}

\begin{definition}
\label{hodge-birational-atoms-2026-definition-object-e}
An object $E$ in $D^b(X)$ is
exceptional if
\[
\operatorname{Ext}^i(E,E)=0
\quad\text{for }i\neq 0,
\qquad
\operatorname{Hom}(E,E)=\mathbb k.
\]
A collection of objects is exceptional
if each object is exceptional and they
are semi-orthogonal.
\end{definition}

\begin{theorem}[Orlov]
\label{hodge-birational-atoms-2026-theorem-x-z}
If $\widetilde X=\operatorname{Bl}_Z X$
is the blow up of a smooth center $Z$,
then
\[
D^b(\widetilde X)
=\langle
\underbrace{D^b(Z),\dots,D^b(Z)}_
{\operatorname{codim}(Z,X)-1},
p^*D^b(X)
\rangle.
\]
\end{theorem}

\begin{theorem}[Bondal--Orlov]
\label{hodge-birational-atoms-2026-theorem-standard-flop}
For a standard flop $X\dashrightarrow Y$
there is an equivalence of derived
categories.
For a standard flip $X\dashrightarrow Y$
the category $D^b(Y)$ is a
semi-orthogonal component of
$D^b(X)$.
\end{theorem}

\medskip\noindent
The idea is to interpret the MMP in the
language of derived categories.

\begin{definition}
\label{hodge-birational-atoms-2026-definition-let-x2}
Let $X$ and $Y$ be birational
varieties.
We say that $X$ and $Y$ are
$K$-equivalent if there is a birational
diagram
\[
\xymatrix{
& Z \ar[dl]_p \ar[dr]^q \\
X \ar@{-->}[rr]^\varphi && Y
}
\]
with $p^*K_X=q^*K_Y$.
\end{definition}

\begin{remark}[Kawamata]
\label{hodge-birational-atoms-2026-remark-slogan-reducing}
The slogan is that reducing the
canonical class $K_X$ is reflected by
reducing the derived category
$D^b(X)$ by a semi-orthogonal
component.
If $X$ and $Y$ are $K$-equivalent,
then
$D^b(X)\simeq D^b(Y)$.
Conversely, if
$D^b(X)\simeq D^b(Y)$,
then $X$ and $Y$ are $K$-equivalent.
\end{remark}

\begin{example}
\label{hodge-birational-atoms-2026-example-x-p}
For $X\subset \mathbb P^5$ a smooth
cubic hypersurface, there is an SOD
\[
D^b(X)=\langle
\mathcal A_X,
\mathcal O(-2),\mathcal O(-1),
\mathcal O
\rangle.
\]
The category $\mathcal A_X$ is
triangulated and $2$-Calabi--Yau:
the Serre functor is isomorphic to
the shift functor.
\end{example}

\begin{remark}[Kuznetsov]
\label{hodge-birational-atoms-2026-remark-x-rational2}
$X$ is rational if and only if
$\mathcal A_X$ is equivalent to
$D^b(S)$ for a K3 surface $S$.
\end{remark}

\begin{remark}
\label{hodge-birational-atoms-2026-remark-natural-approach}
A natural approach is to decompose
semi-orthogonally as far as possible.
One could try to call an SOD maximal
if all its components are
semi-orthogonally indecomposable.
Then one might define atoms of $X$ as
the components of a maximal SOD of
$D^b(X)$.
A nontrivial atom would be one that
does not appear as an atom of
$D^b(Z)$ for
$\dim Z\leq \dim X-2$.
The blow-up formula would then imply
that nontrivial atoms are birational
invariants.
This would give the DK conjecture in
the globally generated case, and one
direction of Kuznetsov's conjecture.

Unfortunately this does not work.
Atoms are not well defined this way,
because SODs are not unique.
For a blow-up of $\mathbb P^2$ at
ten general points, one gets an SOD
with a phantom component.
This shows that the same category may
admit different maximal SODs.
The point of Kontsevich's idea is that
atomic decompositions should be extra
data, not a natural phenomenon.
\end{remark}

\begin{definition}
Let $X$ be a smooth projective variety.
We say that $D^b(X)$ admits an atomic
decomposition if
$$
D^b(X)=\langle A_1,\dots,A_n\rangle
$$
where the $A_i$ are atoms.
An atom is one of the components in
such a decomposition.
\end{definition}

\begin{definition}
\label{hodge-birational-atoms-2026-definition-g-atomic}
A $G$-atomic theory is given by
\begin{itemize}
\item a class of $G$-varieties and
$G$-morphisms $\mathcal V$;
\item a class of $G$-invariant atomic
SODs of $D^b(X)$ for $X\in\mathcal V$.
\end{itemize}
It is subject to:
\begin{itemize}
\item[] 1. as before;
\item[] 2. as before.
\end{itemize}
The point is that the atomic
decomposition is part of the data.
\end{definition}

\begin{remark}
\label{hodge-birational-atoms-2026-remark-prototypical-case}
A prototypical case is
\[
D^b(X)=\langle
\operatorname{Ker}Rf_*,
f^*D^b(Y)
\rangle,
\qquad
Rf_*\mathcal O_X=\mathcal O_Y.
\]
\end{remark}

\begin{definition}
\label{hodge-birational-atoms-2026-definition-sod-d}
For an SOD
\[
D^b(X)=\langle A_1,\dots,A_i,
\dots,A_n\rangle
\]
and $1\leq i\leq n-1$, the $i$th
right mutation is the unique SOD
\[
D^b(X)=\langle A_1,\dots,A_{i-1},
A_{i+1},\widetilde A_i,
A_{i+2},\dots,A_n\rangle.
\]
Left mutation is defined similarly.
\end{definition}

\begin{remark}
\label{hodge-birational-atoms-2026-remark-subcategories-i}
The subcategories $\widetilde A_i$
and $A_i$ are equivalent.
Mutations define an action of the
braid group $\operatorname{Br}_n$ on
the set of SODs.
\end{remark}

\begin{definition}
\label{hodge-birational-atoms-2026-definition-let-v}
Let
\[
\mathcal V_{k,2}^{\mathrm{Bir},G}
=\{
G\text{-surfaces over }k,\,
G\text{-equivariant birational
morphisms}
\}.
\]
\end{definition}

\begin{theorem}[Elagin--Schneider--Shinder]
\label{hodge-birational-atoms-2026-theorem-let-k}
Let $k=\overline k$ and let $G$ be a
group.
Then there exists a $G$-atomic theory
for $\mathcal V_{k,2}^{\mathrm{Bir},G}$.
\end{theorem}

\begin{remark}
\label{hodge-birational-atoms-2026-remark-let-k}
Let $k$ be perfect.
Then there is an atomic theory for
$\mathcal V_{k,2}^{\mathrm{Bir}}$.
\end{remark}

\medskip\noindent
The main operation is the blow up or
blow down of
\[
\widetilde X \to X
\]
along a finite $G$-orbit
$Z=\{p_1,\dots,p_n\}\subset X$.
The exceptional divisor is a
disjoint union
$E=E_1\cup\cdots\cup E_n$ of $(-1)$-
curves.
\[
\xymatrix{
\widetilde X \ar[d] & E=
E_1\cup\cdots\cup E_n
\ar@{_{(}->}[l] \\
X & Z=\{p_1,\dots,p_n\}
\ar@{_{(}->}[l]
}
\]

\begin{definition}
\label{hodge-birational-atoms-2026-definition-such-contraction}
If such a contraction exists, we say
that $X$ is $G$-minimal.
\end{definition}

\begin{theorem}[…,Iskovskikh-Manin,…]
\label{hodge-birational-atoms-2026-theorem-given-g}
Given a $G$-surface $Y$, there exists
a chain of blow ups
\[
Y=X_0\to X_1\to\cdots\to X_n=X
\]
where $X$ is $G$-minimal.
Such $X$ is one of:
\begin{enumerate}
\item a surface with $K_X$ nef;
\item a del Pezzo surface with
$\operatorname{Pic}^G(X)\cong\mathbb Z$;
\item a conic bundle
$X\to C$ over a curve with
$\operatorname{Pic}^G(X/C)\cong\mathbb Z$.
\end{enumerate}
\end{theorem}

\begin{remark}
\label{hodge-birational-atoms-2026-remark-x-type}
If $X$ is of type $1$, then $X$ and
the chain $Y\to\cdots\to X$ are
unique.
If $X$ is of type $2$ or $3$, they may
not be unique.
\medskip\noindent
A Mori fibre space over a point, that
is, a $G$-minimal del Pezzo surface
$X\to \operatorname{pt}$, has degree
$d=K_X^2$ in $\{1,\dots,9\}$.
A Mori fibre space over a curve,
that is, a relatively $G$-minimal
conic bundle $X\to C$, has degree
$d=K_X^2$ in $\{0,\dots,8\}$.
\end{remark}

\begin{remark}
\label{hodge-birational-atoms-2026-remark-let-f}
Let $f:\widetilde Y\to Y$ be the blow
up of a $G$-orbit.
An atomic SOD is
\[
D^b(\widetilde Y)
=\langle
\mathcal O_{E_1}(-1),\dots,
\mathcal O_{E_d}(-1),
f^*B_1,\dots,f^*B_m
\rangle,
\]
where
\[
D^b(Y)=\langle B_1,\dots,B_m\rangle
\]
is an atomic SOD.
\end{remark}

\begin{definition}
\label{hodge-birational-atoms-2026-definition-exceptional}
An exceptional collection
$E_1,\dots,E_n$ is called a block if
\[
\operatorname{Ext}^i(E_j,E_k)=0
\quad\text{for all }i\text{ and }j\neq k.
\]
\end{definition}

\begin{remark}
\label{hodge-birational-atoms-2026-remark-g-minimal}
For a $G$-minimal surface $X/B$ we
write $\mathcal B(X/B)$ for the
standard atomic decomposition.
It is defined by cases.
\begin{itemize}
\item If $K_X$ is nef, then
\[
\mathcal B(X)=\langle D^b(X)\rangle.
\]
\item If $\pi:X\to C$ is a $G$-conic
bundle and $C$ is irrational, then
\[
\mathcal B(X/C)=\langle
\ker \pi^*,\,
\pi^*D^b(C)\rangle.
\]
\item If $X$ is a $G$-minimal del
Pezzo surface of degree $d\leq 4$,
then
\[
\mathcal B(X)
=\langle
\mathcal O_X^\perp,\,
\langle\mathcal O_X\rangle
\rangle.
\]
\item If $\pi:X\to\mathbb P^1$ is a
$G$-Mori conic bundle of degree
$d\leq 4$, then
\[
\mathcal B(X/P^1)
=\langle
\ker \pi^*,\,
\langle\pi^*\mathcal O_{P^1}(-1)\rangle,
\langle\mathcal O_X\rangle
\rangle.
\]
\item If $X$ is birationally rich,
then the standard decompositions are
\[
\mathcal B(\mathbb P^2)
=\langle
\langle\mathcal O(-2H)\rangle,
\langle\mathcal O(-H)\rangle,
\langle\mathcal O\rangle
\rangle,
\]
\[
\mathcal B(\mathbb P^1\times\mathbb P^1)
=\langle
\langle\mathcal O(-h_1-h_2)\rangle,
\langle\mathcal O(-h_1),
\mathcal O(-h_2)\rangle,
\langle\mathcal O\rangle
\rangle,
\]
\[
\mathcal B(X_6)
=\langle
\langle\mathcal O(-H_1),
\mathcal O(-H_2)\rangle,
\langle\mathcal O(-h_1),\mathcal O(-h_2),
\mathcal O(-h_3)\rangle,
\langle\mathcal O\rangle
\rangle,
\]
and
\[
\mathcal B(X_5)
=\langle
\langle\mathcal E\rangle,
\langle\mathcal O(-h_1),\dots,
\mathcal O(-h_5)\rangle,
\langle\mathcal O\rangle
\rangle.
\]
\end{itemize}
\end{remark}

\begin{theorem}
\label{hodge-birational-atoms-2026-theorem-any-birational}
Any birational isomorphism between
$G$-Mori fibre spaces
$X_1/B_1\dashrightarrow X_2/B_2$
decomposes into a composition of
$G$-links, or is an isomorphism.
\end{theorem}

\begin{definition}
\label{hodge-birational-atoms-2026-definition-sarkisov-g}
A Sarkisov $G$-link is a birational
diagram of Mori fibre spaces obtained
by one of the standard Sarkisov
types.
\end{definition}

\begin{theorem}
\label{hodge-birational-atoms-2026-theorem-classification}
There is a classification of links.
\end{theorem}

\begin{example}
\label{hodge-birational-atoms-2026-example-link-type}
A link of type II with
$(\deg X_1,\deg Z,\deg X_2)=(9,7,8)$:
\[
X_1=\mathbb P^2,\qquad X_2=F_0.
\]
Here $Z$ is the blow up of $X_1$ at
an orbit of order $2$ and of $X_2$
at a fixed point.
The mutations are as follows.
\end{example}

\medskip\noindent
Step 3.
Classification of links.

\medskip\noindent
Step 4.
There is nothing to do due to the
smart definition given in Step 1.

\medskip\noindent
{\bf Plan.}

1. Define SODs for nonminimal surfaces
via minimal ones.
2. Define atomic SODs for minimal
surfaces by considering cases.
3. Check that the definitions do not
depend on the choices.
4. Check the axioms of an atomic
theory.

\medskip\noindent
{\bf Notes from now on
are from Elagin's second talk.
Some ideas may be repeated.}

These atoms are different from the
atoms discussed earlier in the notes.
They play the same role in detecting
obstructions to birational morphisms.

\begin{definition}[Atoms
in the sense
of derived categories]
(Incomplete definition…)
Let $X$ be a smooth projective variety.
We say that $D^b(X)$ admits an atomic
decomposition if
$$
D^b(X)=\langle A_1,\dots,A_n\rangle
$$
where the $A_i$ are atoms.
An atom is one of the components in
such a decomposition.
\end{definition}

The Brauer class
$$
a=[A]\in \operatorname{Br}(\mathbb k)
$$
has index $3$.

$$
\operatorname{SB}(A)\cong
\operatorname{SB}(A')
$$
if and only if
$$
A\cong A'
$$
if and only if
$$
a=a'.
$$

\begin{theorem}[Amitsur]
$$
\operatorname{SB}(A)\sim
\operatorname{SB}(A')
$$
if and only if
$$
a=\pm a'
$$
if and only if
$$
\langle a\rangle=\langle a'\rangle
$$
as subgroups in $\operatorname{Br}(\mathbb
k)$.
\end{theorem}

\begin{theorem}[Koll\'ar]
Let $C_1$ and $C_2$ be conics over
$\mathbb k$, with
$$
C_i=\operatorname{SB}(A_i),
$$
where $A_i$ are central simple
algebras of degree $2$.
Then
$$
C_1\times C_2\sim C_1'\times C_2'
$$
if and only if
$$
\langle a_1,a_2\rangle
=
\langle a_1',a_2'\rangle
$$
as subgroups in $\operatorname{Br}(\mathbb
k)$.
\end{theorem}

\begin{example}
Let $X$ be a minimal del Pezzo surface
of degree $6$.
Let
$$
\overline X
=
X\times_{\mathbb k}\overline{\mathbb k}
\cong
\operatorname{Bl}_{p_1,p_2,p_3}
(\mathbb P^2),
$$
and let
$$
G=\operatorname{Gal}(
\overline{\mathbb k}/\mathbb k).
$$

The invariant $(L_2,a_3)$:
\begin{itemize}
\item $H_1,H_2\in \operatorname{Pic}
(\overline X)$ are two classes of
contractions $\overline X\to \mathbb P^2$.
They are permuted by the Galois action.
\item $L_2/\mathbb k$ is the splitting
field of $H_1$ or $H_2$.
Then
$$
[L_2:\mathbb k]=2.
$$
\item
$$
G_2=\operatorname{Gal}(
\overline{\mathbb k}/L_2)
=\operatorname{Stab}(H_1)
=\operatorname{Stab}(H_2)\subset G.
$$
\item The $G_2$-equivariant map
$$
|H_1|:\overline X\to \mathbb P^2
$$
descends to a birational map
$$
X_{L_2}\to \operatorname{SB}(A_3)
$$
for some central simple algebra $A_3$
of degree $3$ over $L_2$.
\item
$$
a_3:=[A_3]\in \operatorname{Br}(L_2).
$$
\end{itemize}

The invariant $(L_3,a_2)$ is introduced
similarly using three classes
$h_1,h_2,h_3\in \operatorname{Pic}
(\overline X)$ that define conic bundles
$$
\overline X\to \mathbb P^1.
$$
\end{example}

\begin{example}
Atoms of a Severi--Brauer surface:
$$
X=\operatorname{SB}(A),
\qquad
a=[A]\in \operatorname{Br}(\mathbb k),
\qquad
3a=0.
$$
The atoms are
$$
(\mathbb k,2a),\quad
(\mathbb k,a),\quad
(\mathbb k,0).
$$
The subgroup is
$$
\langle a\rangle=\{0,a,2a\}.
$$
\end{example}

\begin{example}
Let $X=C_1\times C_2$, where
$$
C_1=\operatorname{SB}(A_1),
\qquad
C_2=\operatorname{SB}(A_2),
$$
and $A_1,A_2$ are central simple
algebras of degree $2$.

Over $\overline{\mathbb k}$ one has
$$
D^b(\mathbb P^1\times \mathbb P^1)
=
\langle
\mathcal O(-1,-1),
\mathcal O(-1,0),
\mathcal O(0,-1),
\mathcal O
\rangle.
$$

Over $\mathbb k$,
$$
D^b(C_1\times C_2)
=
\langle
\mathcal E_{-1,-1},
\mathcal E_{-1,0},
\mathcal E_{0,-1},
\mathcal O
\rangle.
$$

The atoms are
$$
(\mathbb k,a_1+a_2),\quad
(\mathbb k,a_1),\quad
(\mathbb k,a_2),\quad
(\mathbb k,0).
$$
\end{example}

\begin{definition}
Let $X$ be a smooth projective variety.
We say that $D^b(X)$ admits an atomic
decomposition if
$$
D^b(X)=\langle A_1,\dots,A_n\rangle
$$
where the $A_i$ are atoms.
An atom is one of the components in
such a decomposition.
\end{definition}

\begin{definition}
An atom $T$ is {\it small} if
$$
T\cong D^b(\Lambda),
$$
where $\Lambda$ is a division
$\mathbb k$-algebra.
It may be noncommutative.
\end{definition}

\begin{definition}
An atom $T$ is {\it trivial} if
$$
T\cong D^b(F),
$$
where $F$ is a field extension of
$\mathbb k$.
\end{definition}

\begin{remark}
A small atom $D^b(\Lambda)$ is given
by a pair
$$
(F,a),
$$
where
$$
F=Z(\Lambda)
$$
and
$$
a=[\Lambda]\in \operatorname{Br}(F).
$$
\end{remark}

\begin{definition}
A surface $X$ is {\it birationally
rich} if $X$ is birational to a MFS
of degree at least $5$.
\end{definition}

\begin{remark}
A MGS surface $X$ is birationally rich
if and only if
$$
\deg(X)\geq 5.
$$
\end{remark}

\begin{theorem}
Let $X$ be a surface over a perfect
field.
\begin{itemize}
\item $X$ is birationally rich if and
only if all atoms of $X$ are small.
\item $X$ is rational if and only if
all atoms of $X$ are trivial.
\end{itemize}
\end{theorem}

\begin{example}
Del Pezzo surface of degree $6$.

Over $\overline{\mathbb k}$,
$$
D^b(\operatorname{Bl}_3(\mathbb P^2))
=
\langle
\langle \mathcal O(-H_1),
\mathcal O(-H_2)\rangle,
\langle \mathcal O(-h_1),
\mathcal O(-h_2),
\mathcal O(-h_3)\rangle,
\langle \mathcal O\rangle
\rangle.
$$

Over $\mathbb k$,
$$
D^b(X)
=
\langle
\mathcal E_2,
\mathcal E_2,
\mathcal O
\rangle.
$$

The atoms are
$$
(L_2,a_3),\quad
(L_3,a_2),\quad
(\mathbb k,0).
$$

If $a_3$ or $a_2$ is not zero, then
the corresponding atom is a birational
invariant.
\end{example}

\begin{example}
For $X$ a minimal del Pezzo surface of
degree $4$, the atomic decomposition
is
$$
D^b(X)=\langle \mathcal A_X,\mathcal O
\rangle.
$$
The category $\mathcal A_X$ is called
the Kuznetsov component.
\end{example}

\begin{theorem}[ESS, Shramov--Trepalin,
2025]
Let $X$ and $X'$ be minimal del Pezzo
surfaces of degree $4$.
Then the following conditions are
equivalent:
\begin{enumerate}
\item $X$ and $X'$ are isomorphic.
\item $X$ and $X'$ are birational.
\item $\mathcal A_X\cong \mathcal A_{X'}$.
\end{enumerate}
\end{theorem}

\section{Enriched Inflection and Secant
Planes for Linear Series on Algebraic
Curves}
\label{hodge-birational-atoms-2026-section-cotterill}

\medskip\noindent
{\bf Speaker.}

Ethan Cotterill

\medskip\noindent
{\bf Title.}

Enriched Inflection and Secant
Planes for Linear Series on Algebraic
Curves.

\medskip\noindent
{\bf Abstract.}

The fundamental local invariant of a
linear series on an algebraic curve is
its inflection, as measured by the
vanishing of local derivatives in a
point.
Secant planes, which emerge from the
rank behavior of tuples of points on
the curve, may be regarded as (the
manifestations of) a global analogue
of inflection.
Here we report on joint work with
Darago, Garay López, Han, Shaska and
Zhang in which we compute the classes
of inflection points and of tuples of
points spanning secant planes in the
Grothendieck--Witt ring of an
arbitrary perfect base field of
characteristic not equal to $2$.

\medskip\noindent
{\bf Rank Loci in AG.}

Enumerative algebraic geometry
classically deals with the following
paradigm.
We want to count the number of points
in the degeneracy locus of a map of
vector bundles.

\medskip\noindent
Given a map of vector bundles
$\varphi:E\to F$ over a complex
algebraic variety $X$, we would like
to count the number of points
$x\in X$ over which $\varphi$ has
rank at most $k$, for a fixed
$k\in\mathbb N$.

\medskip\noindent
One feature of this paradigm over
$\mathbb C$ is that the number of
solutions, counted with
multiplicities, is uniquely
determined.
This is the principle of conservation
of number.

\medskip\noindent
For other choices of base field
$\mathbb F$, the number of
$\mathbb F$-rational solutions is not
typically unique.
For example, the number of lines on a
cubic over $\mathbb C$ is $27$, but
over $\mathbb R$ it depends on the
topology of the real locus.

\medskip\noindent
{\bf Enriched AG.}

Enrichments, or weights, are used to
address the loss of conservation of
number.

\medskip\noindent
{\bf 1. GW theory.}
Over $\mathbb C$, one counts
holomorphic maps
$\mathbb P^1\to X$ subject to generic
incidence and contact conditions
with linear subvarieties of $X$,
for example with a point of $X$.
Trying to replicate this over
$\mathbb C$ leads to Welschinger
invariants, that is, GW invariants
with signs.

\medskip\noindent
{\bf 2. Quadratic-enriched AG.}
Over an arbitrary field $\mathbb F$,
solutions are weighted by
isomorphism classes of symmetric
bilinear forms over $\mathbb F$.
Welschinger invariants are a special
case.
Conservation of number is achieved
in $\operatorname{GW}(\mathbb F)$,
the Grothendieck--Witt group of
$\mathbb F$.

\medskip\noindent
{\bf 3. Random AG.}
Compute expectations for
cardinalities of zero-dimensional
solution spaces with respect to
distinguished probability
distributions, for example over
$\mathbb R$ or $\mathbb Q_p$.

\medskip\noindent
{\bf Topological Degrees over
$\mathbb R$.}

Consider a polynomial $f(x)$ on the
real line.
For every zero, check the sign of the
derivative of $f$ at that point.
The topological degree of $f$ is
\[
\sum_{f(p)=0}\operatorname{sign}(f'(p)).
\]

\medskip\noindent
More generally, any continuously
differentiable map between compact
oriented $n$-manifolds has a degree,
which generalizes to the induced map
in homology
\[
f_*:H_*(M)\to H_*(N).
\]

\medskip\noindent
{\bf Local ring.}

\medskip\noindent
In the 1970s, Eisenbud--Levine and
independently Khimshiashvili
discovered how to compute local
degrees at not necessarily regular
values using algebraic geometry.

\medskip\noindent
For a germ of a map
\[
f=(f_1,\dots,f_n):(\mathbb R^n,0)
\to(\mathbb R^n,0),
\]
define the local algebra of $f$ as
\[
\mathcal Q(f)=
\mathbb R[[x_1,\dots,x_n]]/
\langle f_1,\dots,f_n\rangle.
\]
Assuming the zero is isolated, this
is finite-dimensional, that is, an
Artinian algebra.

\medskip\noindent
Much more is true.
By the early 1960s, it was noted
that $\mathcal Q(f)$ is a Gorenstein
ring, with one-dimensional socle
spanned by the residue class of
\[
\det J_f.
\]

\medskip\noindent
It also comes with a distinguished
symmetric nondegenerate bilinear
form, called the Scheja--Storch
form.
Eisenbud--Levine and independently
Khimshiashvili showed that
\[
\deg_0(f)=\operatorname{sgn}\phi.
\]

\begin{theorem}[EKL]
\label{hodge-birational-atoms-2026-theorem-signature-form}
The signature of the form is the
topological degree.
\end{theorem}

\medskip\noindent
The talk continued.


\section{Atom applications and
MMP}
\label{hodge-birational-atoms-2026-section-atom-applications-and-mmp2}

\medskip\noindent
{\bf Speaker.}

Ludmil Katzarkov

\medskip\noindent
{\bf Title.}

Atom applications and MMP.

\medskip\noindent
{\bf Abstract.}

In this talk we introduce the
theory of Atoms.
Applications and connection with
MMP will be discussed.

\medskip\noindent
These are notes of the second instance
of Ludmil's talk, which was different
than the first.

\medskip\noindent
{\bf Plan}

\begin{enumerate}
\item Example
\item Categorical approach
\item CFT to the rescue
\item Where do we go now
\end{enumerate}

\medskip\noindent
Consdier a 3-dimensional cubic

$$
x_0^3+\cdots+x_4^3.
$$

If you degenerate this cubic, create
one singular point, and, as was known
to Arquimides(?) it becomes rational.
The multiplicity-2 line intersections
give us a map to
$\mathbb{P}^3$,
which is what
makes it rational.

Let's have a look at this example from
the point of view of Hodge theory. The
hodge diamond of this cubic is:

[Put the hodge diamong]

If you introduce one singular point,
then you are actually going to reduce
by ... and after a smooth resultion,
the hodge diamond of the resolved cubic
looks as follows:

[Hodge diamond with 4's instead of
5's]

It's well known that the variety is
blow up of $\mathbb{P}^3$
[along a quadric?] and we
blow down the quadric.

Now let's discuss LG models. In this
smooth 3d cubic,

$$
W=
\frac{(x+y+z)^3}{xyz}+z
$$

So the LG model is
$(\mathbb{C}^*,W)$. By
HMS, .. is equivalent to the category
of Lagrangians, but we won't use this
in this talk.

In this picture, we have now atoms
interpretation in the cohomology of
our smooth cubic, which is split as
direct sum :

$$
H=\bigoplus H_\lambda
=H_{\lambda_0}+H_{\lambda_1}
+H_{\lambda_2}
$$

In this picture, there is a very
singular fiber... *figure denoting
fibers...*.

This is all very classical. Other
classical constructions like
degenerations and birational geometry
actually fit quite well with LG models
and the behaviours of atoms.

\medskip\noindent
{\bf Galois spectrum} 

Initially introduced by Rafael Rouqet.
Later Orlov developed this introducing
Orlov spectrum. When playing with the LG
of the cubic we noticed there are two
Lagrangian models --- we never published
this.

Let $\alpha$ be a generator of a category---
that is, you can get to any object in this
category by sums, summands, shifts and
triangles. Cf. Alexander's talk.

Associate a number to each of the generators
by forgetting about this operation. The
generation time of alpha is the minimum
number of triangles that we need to get
from $\alpha$ to any other object
in the category. Denote this number by
$t(\alpha)$.

The smallest of the $t(\alpha)$ is called
the Rouquet dimension of the category,
which was conjectured by Rouquet
to be the dimension of $X$.

Orlov introduced the notion of Orlov
spectrum. This is old generation times:
$$
\text{OSpec}(\varphi)
= t(\alpha_1),…, t(\alpha_n).
$$

\begin{example}
\label{hodge-birational-atoms-2026-example-orlov-spectrum2}
$$
\begin{aligned}
D^b(\mathbb{P}^1)&=\left\langle{}1,2\right\rangle{}\\
D^b(E)&=\left\langle{}1,2,3,4\right\rangle{}.
\end{aligned}
$$
\end{example}

\noindent
Using mirror symmetry a lot was done
using generation time of Fukaya categories
and [missing]. As a result of that, the
following conjecture was made:
there is a gap which appears, the
gap of the Orlov spectrum, then the 
conjecture says that the maximum gap
of the Orlov spectrum is bigger than
$\dim (X) -2 $, then $X$ is not rational.

These numbers are very hard to compute.

The maker of CFT was Mostow in the 80s
(confirm). Other important names are
Knizhnik-Zamwbolchicov, among others.
At the time theory of singularity was
proliferating in moscow.
Steen-Brink spectrum, as part of CFT,
who stated several conjectures,
many of which have been proved by 
Arnolds and others.

This correction was not observed in Moscow,
before mirror symmetry, but by Vata-Cecov,
leading to Ditti equations.
A zero dimensional singularity.

In a paper we wrote with M Balland, D Esvero,
and (another?) by Y. Lin.

Vafa understood the connection with the
asymptotics of the [missing name] equation.

Conifold transition. The properties of this
spectrum, the CFT spectrum, …, all this
could lead to birational invariants.
A lot of work is done when CFT is
not related to geometry.

But when geometry helps you, you can do a
lot. In a way that's what [name] will 
be doing from now on.

We will be using quantum differential
equation.

This connection appeared already in
the toric case, where the group
of birational transformations
of a toric variety was expressible
of the behvd at 00 of the correpondng
Lg model
the proptis of the crma grp where exppressed
of the 1dim 2 dim circuits
of the toric combinatorics

1 dimensional circuits corresponds
to generators of the the bir trnasformation
2 circuits to sarkisov links
and ?? dim ciurcits
correspon to relations in the cremona gp.

The new components to all this
is the differential equations ---
that's what I plan to tell you about.

I want to bring to examples here.

Consider the cubic cross $\mathbb{P}^1$,
i.e. $X_3 \times \mathbb{P}^1$.
And also a quartic multiplied by itself,
$X_4 \times X_4$.

A 5-dimensional Fano: the intersection
of quadric and cubic in $\mathbb{P}^7$,
i.e. $Q_2 \cap C_3 \subset \mathbb{P}^7$.
Here we have an interesting Hodge diamond:

\[
\begin{array}{ccccc}
&&1&&\\
&&1&&\\
1&83&&83&1\\
&&1&&\\
&&1&&
\end{array}
\]
the mirror, as observed by Candelas 
corresponds to [missing, related to
a CY].

Now for  the 7-dimensional cubic
$C_3^7 \subset \mathbb{P}^8$

$$
\xymatrix{
C_3^7 \ar[d]\\
C_{3,0}^7
}
$$

\[
\begin{array}{ccccc}
&&1&&\\
&&1&&\\
1&84&&84&1\\
&&1&&\\
&&1&&
\end{array}
\]

\noindent
Here we have a 3-dimensional CY
associated.

The idea is that if we know a lot about
$Q_2 \cap C_3$ we can study $C_3^7$.

\medskip\noindent
Consider the quantum differential equation
$$
\left(\frac{\partial}{\partial u}
-\frac{K}{u^2}+\frac{G}{u}\right)
\psi(u)=0.
$$
This equation has subexponential growth
solutions
$$
\psi(u)=u^{\delta i}\text{log}(…)
(r+… \text{exp} …)
$$

We all $u^{\delta i }$ the 
{\it asymptotics},
which appeared in Leonardo's talk.
$$
\delta = 2 \text{maximum asymptotic}
=2 \text{max} |D_i|
$$
\begin{theorem}
\label{hodge-birational-atoms-2026-theorem-asymtptics-lud2}
Let $X$ be a 
Fano hypersurface 
of degree $d$ in $\mathbb{P}^N$.
Then
\begin{enumerate}
\item $\delta = \dim X - 2 \frac{N+1-d}{d}$
\item  if $\delta> \dim X -2$ then
$X$ is not rational. 
\end{enumerate}
\end{theorem}

\noindent
In the case of $C^3_3$, the 3-dimensional
cubic, $\delta = 3-2^{\frac{5-3}{3}}=
\frac{5}{3}$. In the case of
$X_3 \times \mathbb{P}^1$,
we can use the corresponding LG model,
which for $\mathbb{P}^1$ is
$z+\frac{1}{z}=W$.
There is no 2-dimensional atom which looks
like that [figure of singular fibers].
This proves that it is not rational!
Which of course is not something new.
This would also work for many Fanos times
$\mathbb{P}^1$.

The example of $X_4 \times X_4$
was proven that if you take birationally
rigid Fanos of big dimension and big degrees
then the product is of degree as well.
For $X_4$, $\delta=3-2^{\frac{5-4}{4}}
=\frac{5}{2}$. So for $X_4 \times X_4$
we have $\delta(X_4\times X_4)=5$, 
which is bigger than $\dim X-2$, and
so according to the theorem
we get that it's not rational.

We can apply this for 4-dimensional quintics:
we compute that $\delta=4-2 \frac{6-5}{5}
=\frac{18}{5}$, so
$\delta(X_5\times X_5)=\frac{18}{5}\cdot 2
> 6$.

Now for the intersection of a quadric
and a cubic in $\mathbb{P}^7$,
we cannot use asympotics only.
Denote it by $Y$.
This is a 5-dimensional Fano.
We have four atoms:
$$
H(Y)=H_{\lambda_0} + H_{\lambda_1}
+H_{\lambda_2}+H_{\lambda_3}.
$$
The way this is done is:
$$
\delta= \dim X -2^{\frac{8-3-2}{3}} 
$$
(we substract the degrees of the
hypersurfaces
and divide by the maximal one).

\medskip\noindent
Let $Z$ be the 
minimal model of algebraic variety in
dimension 3, CY, general type, 
there exists a resolution of $Z$
which has only curves as fibers
[Corti, Kologhiros].
Then the quantum differential equation 
splits. And this way, as combination
of atoms and birational geometry
we see get non rationality.

\medskip\noindent
To combine atoms and MMP, use this roadmap
for nonrationality:
start with the splitting of atoms
$H = \bigoplus H_{\lambda_i}$
and restrict the quantum differential
equation. We have $\Delta_{-k}, …,\Delta_k$
If $k< \dim X$,
$\Delta_{-k}>\dim X -2$,
$\Delta_k = \dim X -2$.
And use your knowledge on mirror symmetry
to rule out the example.
Apply MMP and atoms to rule out this case.

\section{Degeneracy loci of Poisson
fivefolds}
\label{hodge-birational-atoms-2026-section-degeneracy-loci-poisson-fivefolds}

\medskip\noindent
{\bf Speaker.}

Jorge Vit\'orio Pereira

\medskip\noindent
{\bf Institution.}

Instituto de Matem\'atica Pura e
Aplicada

\medskip\noindent
{\bf Abstract.}

Let $X$ be a Fano manifold of dimension
five and let
$$
\sigma\in H^0(X,\wedge^2 T_X)
$$
be a holomorphic Poisson structure.
We prove that the rank zero locus
$D_0(\sigma)$ has an irreducible
component of dimension at least $1$ and
that the rank at most two locus
$D_2(\sigma)$ has an irreducible
component of dimension at least $3$.
This confirms a conjecture by Bondal in
dimension five.
The proof uses:
(i) an algebraic integrability criterion
for codimension one foliations on weak
Fano manifolds;
(ii) a lemma of Gualtieri--Pym on
degeneracy loci of ample Poisson
modules; and
(iii) a generalization of a result by
Esteves-Kleiman on the zero loci of
Pfaff fields along invariant
subvarieties.


Let $X$ be a complex manifold.

\begin{definition}
\label{hodge-birational-atoms-2026-definition-poisson-structure}
A Poisson structure is a global
bivector
$$
\sigma\in H^0(X,\wedge^2 T_X)
$$
such that
$$
[\sigma,\sigma]=0,
$$
where the bracket is the Schouten
bracket.
\end{definition}

This allows us to define a Hamiltonian
map
$$
\mathcal O_X\to T_X,
$$
given by
$$
f\mapsto \langle \sigma,df\rangle.
$$

We also get a morphism
$$
\sigma^\sharp:\Omega_X^1\to T_X,
$$
and the map
$$
j:\mathcal O_X\to \Omega_X^1.
$$

\begin{definition}
\label{hodge-birational-atoms-2026-definition-degeneracy-locus}
The $2k$-th degeneracy locus of
$\sigma$ is
$$
D_{2k}(\sigma)
=
\{x\in X:\operatorname{rk}
(\sigma^\sharp_x)\leq 2k\}.
$$
This is the vanishing of a Pfaffian
and is thus an algebraic variety.
Moreover it is a strong Poisson
subvariety.
\end{definition}

\medskip\noindent
{\bf Conjecture (Bondal).}

Suppose that $X$ is a Fano manifold,
that is, $K_X$ is antiample, and that
$\sigma\in \wedge^2T_X$ is a Poisson
structure.
Then
$$
\dim D_{2k}(\sigma)\geq 2k+1.
$$

We expect it to have dimension $2k$,
so the conjecture says it is in fact
bigger than expected.

\medskip\noindent
{\bf Known cases.}

\begin{itemize}
\item Max degeneracy:
$D_{2n-2}$ in odd dimension
(Polishchuk, 1997).
Using index and Baum--Bott theory for
foliations.
\item Generalization:
Beauville observed the same argument
works for $D_{2r-2}$ for any maximal
rank $2r$.
\item Even dimension:
$D_{2n-2}$ follows from nontriviality
of the anticanonical bundle.
\item Submaximal, even dimension:
$D_{2n-4}$ was proved by
Gualtieri--Pym in 2013.
\end{itemize}

\begin{theorem}[Druel--Pereira--Pym--
Touzet]
\label{hodge-birational-atoms-2026-theorem-dppt}
Let $X$ be a Fano manifold of
dimension $2n+1$, with $n>1$, and let
$\sigma$ have rank $2n$.
Then
$$
D_{2n-4}(\sigma)
$$
has an irreducible component of
dimension at least $2n-3$.
\end{theorem}

This gives Bondal's conjecture in
dimension $5$.

\medskip\noindent
{\bf Strategy of proof.}

\begin{enumerate}
\item Consider a saturated foliation
$\mathcal F_\sigma\subset T_X$.
\item Two main cases according to
$\operatorname{codim}(\operatorname{Sing}
(\mathcal F_\sigma))$:
\begin{itemize}
\item $\operatorname{codim}\geq 3$:
algebraic integrability implies
reduction to even dimension.
\item $\operatorname{codim}=2$:
detailed analysis via modular
residues and cohomology
(obstructions).
\end{itemize}
\end{enumerate}

\medskip\noindent
{\bf Three main techniques.}

\medskip\noindent
{\bf Technique 1: algebraic
integrability.}

\begin{theorem}[Druel--Pereira--Pym--
Touzet]
\label{hodge-birational-atoms-2026-theorem-dppt-int}
Let $X$ be weak Fano and let
$\mathcal F$ be codimension $1$ on
$X$ with
$$
\operatorname{codim}
(\operatorname{Sing}(\mathcal F))
\geq 3.
$$
Then $\mathcal F$ is defined by a
morphism
$$
f:X\to \mathbb{P}^1
$$
with irreducible fibers.
\end{theorem}

\begin{itemize}
\item A foliation $\mathcal F$ of codim
$q$ on $X$ is a saturated involutive
subsheaf
$T\mathcal F\subset T_X$ of corank $q$.
\item Normal sheaf: the dual of
$\operatorname{Ann}(T\mathcal F)
\subset \Omega_X^1$.
\item Singular locus:
$\operatorname{Sing}(\mathcal F)$ is
where $T_X/T\mathcal F$ is not locally
free.
\end{itemize}

\medskip\noindent
{\bf Key facts.}

\begin{itemize}
\item $N_{\mathcal F}\cdot C=2$
(Baum--Bott for tangent curves).
\item If $C$ is not tangent, then
$N_{\mathcal F}\cdot C\geq 2$.
\end{itemize}

\begin{definition}
A Pfaff field of degree $p$ on $X$ is
a pair $(L,\eta)$ where $L$ is a line
bundle and
$$
\eta:\Omega_X^p\to L
$$
is $\mathcal O_X$-linear.
Equivalently,
$$
\eta\in H^0(X,\wedge^p T_X\otimes L),
$$
by contraction with $\eta$.
\end{definition}

\medskip\noindent
{\bf Singular locus.} [missing]

\medskip\noindent
{\bf Invariant subvarieties.} [missing]

\medskip\noindent
{\bf Poisson structures give Pfaff
fields.}

\medskip\noindent
{\bf Proposition.}

Let $X$ be a compact K\"ahler manifold
of dimension $n$, and let
$$
\eta:\Omega_X^p\to L
$$
be a Pfaff field.
Let $S\subset Y\subset X$ be closed
subvarieties with:
\begin{enumerate}
\item $Y$ is $\eta$-invariant,
irreducible, and
$$
\dim Y=p+r.
$$
\item $\dim S\leq p-2$.
\item $Y^\circ:=Y\setminus S$ is
smooth.
\item $\eta|_T$ does not vanish in
codimension $1$, and defines a
foliation $\mathcal F^\circ$ of
dimension $p$.
\item $\mathcal F^\circ$ admits a
closed strongly directed positive
current $T^\circ$.
\end{enumerate}
Then the morphism
$$
H^\eta:H^p(X,\Omega_X^p)\to H^p(X,L)
$$
is nonzero.

\medskip\noindent
{\bf Remark.}

Condition $5$ is automatic when
$r=0$.
The current of integration along
$Y^\circ$ works.

\medskip\noindent
{\bf Connection with the work of
Esteves--Kleiman.}

\begin{itemize}
\item Work on algebraic varieties over
arbitrary fields.
\item If $Y$ is $\eta$-invariant and
$\dim \operatorname{Sing}(Y,\sigma|_Y)
\leq p-2$, then the cohomological maps
vanish.
\item Etc.
\end{itemize}

\medskip\noindent
{\bf Proposition.}

Let $X$ be a compact K\"ahler
Poisson manifold.
Suppose $Y\subset X$ is closed and
Poisson.
\begin{enumerate}
\item $\sigma|_Y$ has generic rank
$k$.
\item $\dim \operatorname{Sing}
(Y,\sigma|_Y)\leq 2k-2$.
\item The symplectic foliation
supports a strongly directed positive
current.
\end{enumerate}
Then
$$
h^{0,2k}(X)>0.
$$
In other words, we deduce the
existence of a $(0,2k)$-form.

\medskip\noindent
{\bf Modular residues.}

Gualtieri--Pym introduced modular
residues:
$$
\operatorname{Res}^{\mathrm{mod}}_k(\sigma)
\in
H^0\!\left(
D_{2k}(\sigma),
\mathfrak X_{D_{2k}(\sigma)}^{2k+1}
\right).
$$
Here $\mathfrak X_{D_{2k}(\sigma)}$ is
the singular version of the tangent
sheaf.

The local description is omitted here;
it involves a Poisson connection and a
holomorphic connection.

\medskip\noindent
{\bf Lemma (Gualtieri--Pym).}

If $Y\subset D_{2k}(\sigma)$ is
irreducible, then
$$
\operatorname{Sing}(Y,\sigma|_Y)
$$
has a component of dimension at least
$2k-1$.

\medskip\noindent
{\bf Proof of the main theorem.}

\medskip\noindent
{\bf Case 1.}

$\operatorname{codim}
(\operatorname{Sing}(\mathcal F))
\geq 3$.

\begin{itemize}
\item By the integrability theorem,
$\sigma$ defines a morphism
$$
f:X\to \mathbb{P}^1.
$$
\item A general fiber $F$ is a Fano
manifold of dimension $2n$.
\item $F$ is a Poisson subvariety.
\item Apply the known even-dimensional
result of Gualtieri--Pym:
$$
D_{2n-4}(\cdots).
$$
\end{itemize}

Done.
This is the easy case.

\medskip\noindent
{\bf Case 2.}

$\operatorname{codim}
(\operatorname{Sing}(\mathcal F))=2$.

\medskip\noindent
{\bf Setup.}

Let $Y$ be an irreducible component of
\[
\operatorname{Sing}(\mathcal F),
\]
with
$$
\dim Y=2k-1.
$$
Set
$$
S:=\operatorname{Sing}
(Y,\sigma|_Y).
$$
Then $S$ and $Y$ are strong Poisson
subvarieties of
$$
D_{2k-2}(\sigma),
$$
and
$$
\dim S\leq 2k-2.
$$

\medskip\noindent
{\bf Three subcases.}

\begin{itemize}
\item $\dim S=2k-3$.
\item $\dim S=2k-2$:
use the previous proposition and other
considerations.
\item $\dim S\leq 2k-4$:
the most delicate case.
Set
$$
\rho=
\operatorname{Res}^{\mathrm{mod}}_{k-1}
(\sigma)|_{Y^\circ}.
$$
\end{itemize}

\medskip\noindent
{\bf Subcase 2.1.}

If $\rho$ vanishes somewhere, then the
lemma gives a component.

\medskip\noindent
{\bf Subcase 2.2.}

The remaining case is omitted here.

\medskip\noindent
{\bf Summary.}

\begin{itemize}
\item Bondal's conjecture is confirmed
for the submaximal degeneracy locus of
odd-dimensional Fano manifolds.
\item This proves the conjecture in
dimension five.
\item The proof combines:
\begin{itemize}
\item algebraic integrability of
codimension $1$ foliations on weak Fano
manifolds;
\item modular residues of Poisson
structures;
\item Esteves--Kleiman type cohomological
constraints via positive currents on
compact K\"ahler manifolds.
\end{itemize}
\end{itemize}

\section{Spectral decomposition of
nc-Hodge structures and F-bundles}
\label{hodge-birational-atoms-2026-section-nc-hodge-spectral}

\medskip\noindent
{\bf Speaker.}

Shaowu Zhang

\medskip\noindent
{\bf Title.}

Spectral decomposition of
nc-Hodge structures and F-bundles.

\medskip\noindent
{\bf Abstract.}

Non-commutative Hodge structures
generalize classical Hodge
structures for non-commutative
spaces.
They consist of de Rham and Betti
data, encoded as Stokes structures.
These structures naturally arise in
mirror symmetry in the context of the
Gamma conjectures.
In this talk, I will first discuss
the behavior of Betti data under the
topological Laplace transform,
relating this to the spectral
decomposition of non-commutative
Hodge structures.
I will then talk about the formal
and non-Archimedean analogues of
nc-Hodge structures, known as
F-bundles, focusing on the spectral
decomposition theorem for maximal
F-bundles and its connection to
Hodge atoms.

\section{Blowups, Gale duality and moduli
spaces}
\label{hodge-birational-atoms-2026-section-blowups-gale-duality}

\medskip\noindent
{\bf Speaker.}

Carolina Araujo

\medskip\noindent
{\bf Title.}

Blowups, Gale duality and moduli
spaces.

\medskip\noindent
{\bf Abstract.}

We discuss the birational geometry of
blowups of projective spaces at points
in general position.
For this we explore Gale duality,
a correspondence between sets of
$n=r+s+2$ points in
$\mathbb{P}^s$ and $\mathbb{P}^r$.
For small values of $s$, this duality
has a geometric manifestation:
the blowup of $\mathbb{P}^r$ at $n$
points can be realized as a moduli
space of vector bundles on the blowup
of $\mathbb{P}^s$ at the Gale dual
points.

Joint work with Castravet, Kaur, Martinelli.

{\bf Notes} 

Setup. Consider a collection of points
in ``very general position''.
Consider $X$, the blow up of $\mathbb{P}^n$
in at $P_1,…,P_k$.
Problem: describe all the ``rational
contraction'' that variety admits.
If the number of points is small
the geometry is simple.
As the number grows, the geometry becomes
increasingly more complicated and eventually
becomes infinite.

Consider the example of $Bl_P \mathbb{P}^2$.
It's a $\mathbb{P}^1$ bundle over
$\mathbb{P}^1$. Now increase the number of
points, say,m $Bl_{P_1,P_2}\mathbb{P}^2$
now we have two exceptional divisors.
Consider the strict transform of a line
$\ell \subset \mathbb{P}^2$, $\tilde{\ell}$
which is isomorphic to $\mathbb{P}^1$
and has $\ell^2 = -1$, which
by Castelnuovo's thoerem we know
we can blow down, yielding a quadric.
These are all the rational contractions
of this blow-up.

Finally consider three points,
$Bl_{P_1,P_2,P_3}\mathbb{P}^2$
this is del Pezzo 3.
Consdier the stric transform of the
three lines passing therough this points.
Blowing down leads to the standard
Cremona transformation!
If we lift it to $X$, we obtain
an automorphism of $X$.

Denote $H=\pi^*  \mathcal{O}_{\mathbb{P}^2}
(1)$ and $E_i$ the exceptional divisor over
$P_i$.
$\text{Pic}(X)$ and $N^1=\text{Pic}(X)
\otimes_\mathbb{Z} \mathbb{R} cong
\mathbb{R}^4$. A reflection $q^*$
across a hyperplane,
$q: (x:y:z) \mapsto (? ? ? )$
which preserves $Nef(X) \subset Eff(X)
\subset N^1$.
We can construct such transformations
in similar ways, and extending to 
reflections in $\text{Pic}$L
$q_{ij l}^*:N^1(X) \to N^1(X)$ isometry.
We can do this for any collection of 3
points, and we obtain a Weyl subgroup
$$
W := \left\langle{}q_{ij l}\right\rangle{} \subset
O(N^1)
$$
whose action preserves $(-1)$-curves
and the cones $Nef(X)\subset E ff(X) \subset
N^1(X)$.

$W$ is finite $\iff$ $k \leq  8$.

\begin{theorem}
\label{hodge-birational-atoms-2026-theorem-nagata}
blowup of $\mathbb{P}^2$ ??
\end{theorem}

\begin{itemize}
\item 

for $k \leq  8$ X has finitely many (-1)
curves ($-K_X$ is ample.

\item for $k \geq 9$, $X$ has inifintely many
$(-1)$-curves

$$
\begin{cases}
k=9\qquad &-K_X \text{ is nef} \\
k>9\qquad &-K_X\text{ is not nef}
\end{cases}
$$
\item $k=9$ [Cone diagram]

\item $k\geq 10$ [Cone diagram]

\end{itemize}

\begin{theorem}[D'urso, impa!]
\label{hodge-birational-atoms-2026-theorem-durso}
$Nef(X_{K_X \leq 0})$ has a finite polyhedral
fundamental domain for the action of $W$.
\end{theorem}

\begin{example}
\label{hodge-birational-atoms-2026-example-blowup-moridream}
$X = Bl_{P_1,P_2}\mathbb{P}^n
\xrightarrow{\pi}\mathbb{P}^n$,
$n\geq  3$.

\begin{itemize}
\item $H = \pi^*\mathcal{O}_{\mathbb{P}^n}
(1)$.
\item $E_i$ exceptional divisors
\item $H_i$ strict trfm of hyperplane
\item $H_{12}$ strict transform
of hyperplane through $P_1$ and $P_2$.
$$
Nef(X) \subset Mov(X) \subset Ef f(X)
\subset N^1(X).
$$

$$
[\text{Diagram of a triangle of the
Nef cone
see cremona-transformations.tex]}
$$

$$
\xymatrix{
&Y \supset F=\mathbb{P}^1 \times
\mathbb{P}^{n-2}\ar[dl]
^{Bl_{\tilde{\ell}_{12}}}
\ar[dr]^{\text{Cont}_F}\\
\mathbb{P}^1 \cong \ell_{12}
\subset X \ar@{-->}[rr]^{f\ell_p}
\ar[dr]^{\text{contr}_{\tilde{\ell}_{12}}}&&
 X' \supset Z \cong \mathbb{P}^{n-2}
\ar[dl] \\
& W
}
$$
\end{itemize}
\end{example}

\noindent
Finally, the blow up of $\mathbb{P}^n$
at $k \geq  n+1$ very generap points.

For $I \subset \{ 1,…, k \}$
$|I|=n+1$ we obtain reflections
$(q_l)^*:N_1(X_l) \cong N^1(X_l)$.

$$
  \xymatrix{
  X_l \ar@{-->}[r]^{\cong \text{ pseudo}}
  \ar[d] & X_l \ar[d]\\
  \mathbb{P}^n \ar@{-->}[r]^{q_l} &
  \mathbb{P}^n
  }
  $$
Coble pairing on $\text{Pic}(X)$:
\begin{itemize}
\item $H^2=n-1$
\item ??aaaa
\item ??
\end{itemize}

\medskip\noindent
{\bf Mori dream spaces} 

\begin{definition}
\label{hodge-birational-atoms-2026-definition-cox-ring}
$$
\text{Cox}(X) = \sum_{D \in \text{Pic}(X)}
H^0(X,\mathcal{O}_X(D).
$$
A variety is a {\it Mori dream space} 
is the Cox ring above is finitely generated.
\end{definition}

\begin{theorem}[Batryev-Popov, 2004]
\label{hodge-birational-atoms-2026-theorem-batryev-popv04}
For $Bl$ of $\mathbb{P}^2$ at $k$ points
in very general postn
$X $ is MDS iff $k \leq  0$ iff $-K_X$
is ample.
\end{theorem}

\begin{definition}
\label{hodge-birational-atoms-2026-definition-very-general-position}
[Missed it, picture]
\end{definition}

\begin{theorem}[Hu-Keel, 2000]
\label{hodge-birational-atoms-2026-theorem-hu-keel}
Let $X$ be a Mori Dream space, then
\begin{itemize}
\item The effective cone $Ef f(X)$
\begin{itemize}
\item is generated by finitely many
prime divisors
\item decomposes into subcones
$\leftrightarrow$ rational 
contractions of $X$
\end{itemize}
\item The nef cone is generate by finitely
many
semiample divisors.
\item The movable cone decomposes into
subcones.
\end{itemize}
\end{theorem}

Upshot:
The birational geometry of MDS is finite!
And… we know exactly which MDS
are out there.

\begin{theorem}[Mukari, 2001,
Castravel-Tevelev, 2005]
\label{hodge-birational-atoms-2026-theorem-teve}
Tje blowup of $\mathbb{P}^n$ at $k$ 
points is a MDS iff **conditions**
\end{theorem}

\medskip\noindent
{\bf Problems.} 

\begin{itemize}
\item Mori chamber decomposition of $Ef f(X)$
in the MDS cases.
\item What can be said in the non-MDS cases?
(we shall focus on this one today!)
in the above theorem, $k=n+2$ and 
$n \geq 5$.
\end{itemize}

\medskip\noindent
{\bf Gale duality.} 

\noindent
Let $n, s \in \mathbb{Z}_{\geq 0}$
and $k= n+s+2$

Gale correspondence is:
$$
\left\{ \{P_1,…,P_k\} \subset \mathbb{P}^n
\right\}/\sim
\leftrightarrow
\left\{ \{ Q_1,…,Q_n\} \subset \mathbb{P}^s
\right\}/\sim
$$
where the bijection is up to
projective equivalence.

\begin{example}
\label{hodge-birational-atoms-2026-example-gale-s-equal-1}
Take $n+3$ points in $\mathbb{P}^1$
here's how I give $n+3$
points in $\mathbb{P}^n$:
there's a unique rational curve passing
through… [check, this is for $s=1$
and is very simple].
\end{example}

\noindent
More generally, we would use a matrix,
its transpose, and complete with
the identity matrix.
This is due to Gale in the 50's,
when he studied convex geometry, polytopes,
etc. But is also related to the work
of Coble (more recently).

This Gale duality is compatible with
Cremona transformations.
Let 
$$
P_1,…,P_{n+1}, P_{n+2},…, P_{n+s+2}
\leftrightarrow
Q_1,…,Q_{n+1},Q_{n+2},…,
Q_{n+s+2}
$$
$$
iP_1,…,P_{n+1},…,P_{n+1},
\hat{P}_{n+2},…,
\hat{P}_{n+s+2}
\leftrightarrow
\hat{Q}_1,…,\hat{Q}_{n+1}
,Q_{n+1},…,Q_{n+s+2}
$$
And in fact, this has related
Dynkyn diagrams. [Obviously
an explanation of why this has
to do with Cremona transformations is 
missing]

\medskip\noindent
Now, there exists a linear isomorphism
$$
\varphi:N^1(X) \to N^1(X)
$$
which
\begin{itemize}
\item sends the canonical to the canoinca
\item  and restricts to an isometry
$K^\perp_X \to K^\perp_Y$.
\end{itemize}

\noindent
This is due to the relation
if Gale duality with
the Cremona transformations
explained above.

\begin{theorem}[Mukai 2001,Baue r 1991]
\label{hodge-birational-atoms-2026-theorem}
Description of the Mori chamber
decomposition of the effective cone
when $X$ is the blowup of $\mathbb{P}^n$
at $n+3$ very general points.
\end{theorem}

\noindent
(This is the case we highlighted
above as ``the interesting for us''
based on the description
of the MDS there exist theorem.)

\begin{theorem}[Mukai 2001, Casagrande-
Codogni-Fanelli 2019]
\label{hodge-birational-atoms-2026-theorem-duplicate-2}
Description of mori chamber when $X$ is bloup
of $\mathbb{P}^4$ at $8$ v general pts.
\end{theorem}

\begin{theorem}[A-Castravet-Kaur-Matrinelli]
\label{hodge-birational-atoms-2026-theorem-duplicate-3}
PArtiall extended this description for??
\end{theorem}

\medskip\noindent
{\bf Idea of proof.} 

$P_1,…,P_{n+4} \in \mathbb{P}^n$
and dual Gale points $Q_1,…,Q_{n+4} 
\in \mathbb{P}^2$.

$X =$ blowup of $\mathbb{P}^n$ at $
P_1,…,P_{n+4}$,
$S = $ blowup of $\mathbb{P}^2$ at
$Q_1,…,Q_{n+4}$.

On $S$ we fix invariants $r=2, c_1=-K_S$
and $c_2=2$.

$L \in Amp(S)$ $\xymatrix{\ar@{~>}[r]&}$
Moduli space of Gieseker $L$-semistable
torsion free sheaves,
$\overline{\mathcal{M}}_L$.

Variation of $\overline{\mathcal{M}_L}$
as $L$ varies.
Friedman-Qin (95), Göttsche (96),
Matsuki-Wentworth (97).

\begin{itemize}
\item Chamber decomposition on 
$Amp(S)$ ($L \equiv L' \iff
\overline{\mathcal{M}}_L =
\overline{\mathcal{M}}_{L'}$.
\item (In the 90s, what people
where doing…) birational maps
between models on $K_S <0$.
\end{itemize}

\begin{theorem}[A-Castravet-Kaur-Martinelli]
\label{hodge-birational-atoms-2026-theorem-ackm}
\begin{itemize}
\item $\exists L_0 \in \text{Amp}(S)$
such that $X \cong \overline{
\mathcal{M}}_{L_0}$.
\item Mori chamber decomposition on
$Eff(X) \cap (K_X<0)$.
\end{itemize}
\end{theorem}

\section{SYZ for almost abelian
Lie groups}
\label{hodge-birational-atoms-2026-section-lino}

\medskip\noindent
{\bf Speaker.}

Lino Anderson da Silva Grama
(IMECC)

\medskip\noindent
{\bf Title.}

SYZ for almost abelian Lie groups.

\medskip\noindent
{\bf Abstract.}

In this talk, we discuss SYZ mirror
symmetry in the non-K\"ahler context,
in the sense of Lau--Tseng--Yau.
We construct SYZ dual pairs for a
family of solvmanifolds called almost
abelian, and we explore applications
related to dualities between
Bott--Chern and Tseng--Yau
cohomologies.
This is joint work in progress with
L. Cavenaghi, L. Katzarkov, and
P. Muniz.

\begin{definition}
\label{hodge-birational-atoms-2026-definition-su3-iib}
An $\operatorname{SU}(3)$ structure
$(X,\omega,\Omega)$ is supersymmetric
of type IIB if $\Omega$ defines an
honest complex structure and $\omega$
defines a balanced metric. That is,
$$
d\Omega=0
$$
and
$$
d(\omega^2)=0.
$$
For the conformal factor $F$, there
is $\rho_B$, the RR flux source
current. It is a $(4,4)$-form.
\end{definition}

\begin{remark}
\label{hodge-birational-atoms-2026-remark-su3-holonomy}
There is a system of equations.
When $\omega$ is symplectic and $F$
is constant, $X$ is a CY and
$$
\rho_B=0.
$$
In this case $X$ has
\(\operatorname{SU}(3)\) holonomy.
\end{remark}

\begin{remark}
\label{hodge-birational-atoms-2026-remark-su3-system}
We also have a system of three
equations which determine the
structure.
If $d\Omega=0$, we again get a CY
manifold.
\end{remark}

\begin{definition}
\label{hodge-birational-atoms-2026-definition-su3-iia}
An $\operatorname{SU}(3)$ structure
is supersymmetric of type IIA if
$$
d\omega=0
$$
and
$$
d(\operatorname{Re}\Omega)=0.
$$
Now $\rho_A$ is an exact $3$-form.
\end{definition}

\begin{definition}
\label{hodge-birational-atoms-2026-definition-su-n-semi-flat}
An $\operatorname{SU}(n)$ structure
$(X,\omega,\Omega)$ is semi-flat if
$X$ admits a Lagrangian torus bundle
$$
\mu:X\to B
$$
and
$$
\omega,\Omega\in\Omega_B^*(X,\mathbb{C}).
$$
That is, they are base dependent forms.
\end{definition}

\begin{definition}
\label{hodge-birational-atoms-2026-definition-polarization-switch}
The polarization switch operator
\(\mathcal{P}\) acts on
\(\Omega_B^{p,q}(X)\) locally as
$$
\alpha=\sum f_{I,J}\,dz_I\wedge
d\bar z_J
\mapsto
\mathcal{P}(\alpha)
=\sum f_{I,J}\,dz_I\wedge dz_J.
$$
\end{definition}

\medskip\noindent
{\bf Fiber product setup.}

Let
$$
X\times_B \widehat X
$$
be the fiber product of dual torus
bundles over $B$.
It is the correspondence space for
the Fourier--Mukai transform.

\medskip\noindent
{\bf Fourier--Mukai transform.}

The transform is a push-pull operator
with a correction term given by a
wedge product with an exponential.

\begin{theorem}[Lau--Tseng--Yau, 2015]
\label{hodge-birational-atoms-2026-theorem-lau-tseng-yau}
Let $\omega$ be a torus invariant real
(1,1)-form on $X$. Let
$$
\widehat\Omega
=\operatorname{FT}(e^{2\omega}).
$$
Then
\begin{itemize}
\item $(X,\Omega,\omega)$ is
\(\operatorname{SU}(n)\) if and only if
\item $(\widehat X,\widehat\omega,
\widehat\Omega)$ is the dual
\(\operatorname{SU}(n)\) structure.
\end{itemize}
\end{theorem}

\medskip\noindent
{\bf SYZ for $3$-d solvmanifolds.}

Bedulli--Vannini (2024) study SYZ
mirror symmetry in the setting of
$3$-dimensional solvmanifolds.
Their result fits the semi-flat and
Fourier--Mukai picture above.

\medskip\noindent
Up to isomorphism, there are four
simply connected unimodular solvable
Lie groups of dimension $3$.

\begin{itemize}
\item $\mathbb{R}^3$.
\item $\operatorname{Nil}^3$.
\item $\operatorname{Sol}^3$.
\item $G_{\alpha}$, the remaining
non-nilpotent unimodular solvable
case.
\end{itemize}

\begin{theorem}[Bedulli--Vannini, 2024]
\label{hodge-birational-atoms-2026-theorem-bedulli-vannini}
For dual torus bundles
$$
\widehat X\to B
\qquad\text{and}\qquad
X\to B
$$
\begin{itemize}
\item $\widehat X$ is a solvmanifold
for the semidirect product of $G$
with $(\mathbb{R}^n)^*$.
\item $X$ is a solvmanifold for the
semidirect product of $G$ with
$\mathbb{R}^n$.
\item An invariant type IIA structure
on $X$ corresponds to an invariant
type IIB structure on
\(\widehat X\).
\end{itemize}
Here
$$
\rho:G\to\operatorname{GL}(n,\mathbb{R})
$$
is the linear part of the affine
representation.
\end{theorem}

\begin{definition}
\label{hodge-birational-atoms-2026-definition-almost-abelian}
A Lie group is called almost abelian
if its Lie algebra $\mathfrak{g}$ has
an abelian ideal $\mathfrak{h}$ of
codimension $1$.
That is, choosing
$e_0\notin\mathfrak{h}$,
$\mathfrak{g}$ decomposes as
the sum of $\mathbb{R}e_0$ and
$\mathfrak{h}$.
Such a Lie algebra is also called
almost abelian.

Fix a basis $\mathcal{B}$ of
\(\mathfrak{h}\), and let $A$ be the
matrix representing
\(\operatorname{ad}_{e_0}|_{\mathfrak{h}}\)
with respect to $\mathcal{B}$.
If $\dim \mathfrak{h}=d$, then
\(\mathfrak{g}\) is determined by
$$
[e_0,v]=Av.
$$
\end{definition}

\medskip\noindent
There are two sides, the complex side
and the symplectic side.

\begin{theorem}
\label{hodge-birational-atoms-2026-theorem-two-sides}
The pairs constructed above on each
side are supersymmetric of types IIB
and IIA.
The duality is given by the
Fourier--Mukai transform.
\end{theorem}

\begin{remark}
\label{hodge-birational-atoms-2026-remark-proof-idea}
Idea of the proof.

Type IIB for $(X,\Omega,\omega)$:
\begin{itemize}
\item $d\Omega=0$ by construction.
\item To get $d(\omega^{n-1})=0$, use
\(\operatorname{tr} A=0\).
\item Compute $d\omega$.
\item Then
$$
d(\omega^{n-1})=\cdots
$$
\end{itemize}

Type IIA:
\begin{itemize}
\item Show
$$
\operatorname{FT}(e^{2\omega})=0?
$$
\item Use the duality to get the
symplectic side.
\end{itemize}
\end{remark}

\begin{theorem}[Andrada--Orfila, 2016]
\label{hodge-birational-atoms-2026-theorem-andrada-orfila}
Let
$$
G
$$
be the semidirect product of
\(\mathbb{R}\) and \(\mathbb{R}^{2n+1}\)
via $\varphi$.
Then $G$ admits a lattice if and
only if there exists $t_0\neq 0$ such
that $\varphi(t_0)$ is conjugate to an
integer matrix.
In that case, the lattice is
the group $t_0\mathbb{Z}\times \Gamma$,
with multiplication induced by
$\varphi(t_0)$, where
$\Gamma\subset\mathbb{R}^{2n+1}$
is a lattice.
\end{theorem}

\begin{theorem}
\label{hodge-birational-atoms-2026-theorem-lattice-existence}
There is another lattice existence
theorem for solvable Lie groups.
\end{theorem}

\medskip\noindent
{\bf Cohomology theories.}

\begin{definition}
\label{hodge-birational-atoms-2026-definition-real-polarization}
Let $X$ be a real $2n$-fold and let
$\omega$ be a nondegenerate $2$-form.
A real polarization with respect to
$\omega$ is an integrable Lagrangian
distribution
$$
\{\Lambda_x\subset T_xX:x\in U\}
$$
over an open dense subset $U\subset X$.

Given a real polarization
$(\Lambda,U)$ on $(X,\omega)$ and a
metric on $X$, there is a splitting
of the tangent bundle
$$
TU=\Lambda\oplus\Lambda^\perp
$$
and a splitting of forms
$$
\Omega^k(U)
=\bigoplus_{p+q=k}
\Omega_\Lambda^{(p,q)}(U).
$$
Here $\Omega_\Lambda^{(p,q)}(U)$ is
the subspace of forms of type
$(p,q)$ with respect to $\Lambda$.
\end{definition}

\begin{definition}
\label{hodge-birational-atoms-2026-definition-bott-chern}
The Bott--Chern cohomology of a
complex manifold $X$ is
$$
H^{p,q}_{BC}(X)
=\frac{
\ker d\cap\Omega^{p,q}(X)
}{
\operatorname{im}\partial\bar\partial
}.
$$
It is the complex-side analogue of a
polarized cohomology theory.
\end{definition}

\begin{definition}
\label{hodge-birational-atoms-2026-definition-tseng-yau}
Let
$(\widehat X,\widehat\omega)$ be a
Lagrangian bundle over $B$ with real
polarization $\Delta$.
The Tseng--Yau cohomology is
$$
H^{(p,q)^\Delta}_{B,TY}(\widehat X)
=\frac{
\ker(d+d^\Lambda)\cap
\Omega_\Delta^{(p,q)}(\widehat X)
}{
\operatorname{im}dd^\Lambda
}.
$$
The point is that it is similar to
Bott--Chern cohomology, but using
$d^\Delta$ on the symplectic side.
\end{definition}

\begin{theorem}[Lau--Tseng--Yau, 2015]
\label{hodge-birational-atoms-2026-theorem-ty-bc-iso}
If $\Delta$ is a real polarization on
\(\widehat X\) coming from the bundle
structure, then the Fourier--Mukai
transform gives an isomorphism
between the Bott--Chern and
Tseng--Yau cohomologies.
\end{theorem}

\medskip\noindent
The problem of computing Bott--Chern
cohomology for solvmanifolds was
studied by Angella--Kasuya.

\begin{theorem}[Angella--Kasuya]
\label{hodge-birational-atoms-2026-theorem-angella-kasuya}
Let $G$ be a connected solvable Lie
group, and let $\Gamma$ be a lattice
in $G$.
Then the Bott--Chern cohomology of
\(G/\Gamma\) can be computed from
invariant forms under suitable
assumptions.
\end{theorem}

\medskip\noindent
We use this theorem for
$$
G\times\mathbb{R}^n
\cong
\mathbb{C}\times\mathbb{C}^{n-1},
$$
where
$$
\phi(x_1+iy_1)
=\operatorname{diag}
\bigl(
e^{x_1\lambda_2},
\dots,
e^{x_1\lambda_n}
\bigr).
$$

\begin{theorem}
\label{hodge-birational-atoms-2026-theorem-bc-computation}
This theorem gives an expression for
the Bott--Chern cohomology
$$
H^{p,q}_{BC}
\bigl(
G\times\mathbb{R}^n
/
\Gamma\times\Xi
\bigr).
$$
\end{theorem}

\begin{lemma}
\label{hodge-birational-atoms-2026-lemma-bc-ty-dimensions}
For
$$
h^{p,q}_{BC}
=\dim H^{p,q}_{BC}
$$
and
$$
h^{p,q}_{TY}
=\dim H^{p,q}_{TY},
$$
we have
$$
h^{p,q}_{BC}=N(p,q)
$$
and
$$
h^{n-p,q}_{TY}=N(n-p,q).
$$
\end{lemma}

\begin{lemma}
\label{hodge-birational-atoms-2026-lemma-hodge-number-relations}
The Hodge numbers satisfy
$$
h^{p,q}_{BC}=h^{q,p}_{BC}
$$
and
$$
h^{p,q}_{BC}=h^{n-p,n-q}_{BC}.
$$
\end{lemma}

\begin{definition}
\label{hodge-birational-atoms-2026-definition-symplectic-hodge-star}
The symplectic Hodge star operator
is the operator
$$
\star_s:\Omega^k(\widehat X)\to
\Omega^{2n-k}(\widehat X)
$$
defined by
$$
\alpha\wedge(\star_s\beta)
\;=\;
\langle \alpha,\beta\rangle
\frac{\widehat\omega^n}{n!}
$$
for forms $\alpha,\beta$ of degree
$k$.
\end{definition}

\begin{theorem}[Brylinski]
\label{hodge-birational-atoms-2026-theorem-brylinski}
If $(\widehat X,\widehat\omega)$ is a
symplectic $2n$-manifold, then there
is an isomorphism
$$
H^{2n-k}_{d^\Lambda}(\widehat X)
\simeq
H^k_{dR}(\widehat X)
$$
given by the symplectic Hodge star
operator.
\end{theorem}

\begin{definition}
\label{hodge-birational-atoms-2026-definition-tseng-yau-dplus}
Tseng and Yau proposed the symplectic
cohomology
$$
H^k_{d+d^\Lambda}(\widehat X)
=\frac{
\ker(d+d^\Lambda)\cap\Omega^k(\widehat X)
}{
\operatorname{im}d\cap\Omega^k(\widehat X)
+\operatorname{im}d^\Lambda\cap
\Omega^k(\widehat X)
}.
$$
\end{definition}

\begin{definition}
\label{hodge-birational-atoms-2026-definition-tseng-yau-dd}
Tseng and Yau also proposed
$$
H^k_{dd^\Lambda}(\widehat X)
=\frac{
\ker(dd^\Lambda)\cap\Omega^k(\widehat X)
}{
\operatorname{im}d\cap\Omega^k(\widehat X)
+\operatorname{im}d^\Lambda\cap
\Omega^k(\widehat X)
}.
$$
\end{definition}

\medskip\noindent
Properties.
\begin{itemize}
\item If $X$ is compact, both
cohomologies are finite.
\item One can do Hodge theory:
harmonic forms and Hodge
decomposition.
\item They are not isomorphic to
de Rham cohomology.
\item They are invariant under
symplectomorphisms and Hamiltonian
isotopy.
\end{itemize}

\begin{theorem}[Tseng--Yau]
\label{hodge-birational-atoms-2026-theorem-tseng-yau}
If $(\widehat X,\widehat\omega)$ is a
compact symplectic manifold, then
$$
H^k_{d+d^\Lambda}(\widehat X)
\simeq
H^{2n-k}_{dd^\Lambda}(\widehat X).
$$
\end{theorem}

\begin{definition}
\label{hodge-birational-atoms-2026-definition-ddlambda-lemma}
We say that $\widehat X$ satisfies
the $dd^\Lambda$ lemma if
\begin{enumerate}
\item every form that is both
$d$-exact and $d^\Lambda$-closed is
$dd^\Lambda$-exact,
\item every form that is both
$d^\Lambda$-exact and $d$-closed is
$dd^\Lambda$-exact,
\item every form that is both
$d$-closed and $d^\Lambda$-closed is
cohomologous to a $d$-exact form
only if it is $dd^\Lambda$-exact.
\end{enumerate}
\end{definition}

\begin{theorem}
\label{hodge-birational-atoms-2026-theorem-ddlambda-dR}
The $dd^\Lambda$ lemma holds if and
only if the natural map to de Rham
cohomology is an isomorphism.
\end{theorem}

\begin{remark}
\label{hodge-birational-atoms-2026-remark-kaehler-ddlambda}
If $\widehat X$ is a K\"ahler
manifold, then it satisfies the
$dd^\Lambda$ lemma.
\end{remark}

\medskip\noindent
Questions.
\begin{itemize}
\item Is there a more natural way to
obtain the cohomologies introduced
above?
\item What is the right notion of
symplectic cohomology here?
\end{itemize}

\medskip\noindent
New proposal of symplectic cohomology.

Introduce a formal variable $u$ and
consider the complex
$$
d+ud^\Lambda.
$$
This gives four possible cohomologies:
\begin{itemize}
\item $dd^\Lambda$ cohomology.
\item Periodic cyclic cohomology.
\item The symplectic analogue of
Bott--Chern cohomology.
\item The symplectic analogue of
Aeppli cohomology.
\end{itemize}

Let $N_1=\ker(d_1)$ and $N_2=\ker(d_2)$.
If $N_1\subset N_2$, then the inclusion
induces two parallel exact sequences.

\begin{lemma}
\label{hodge-birational-atoms-2026-lemma-cokernels}
There is a canonical isomorphism
between the cokernels.
\end{lemma}

\begin{theorem}
\label{hodge-birational-atoms-2026-theorem-symplectic-bc}
Let $(\widehat X,\widehat\omega)$ be
a symplectic manifold.
Then the symplectic Bott--Chern
cohomology has odd and even parts.
These are the two graded pieces of
the new proposal.
\end{theorem}

\medskip\noindent
Let $\widehat X\to B$ and $X\to B$
be non-K\"ahler SYZ mirror partners.
Denote by $H_{B;BC,s}$ the cohomology
restricted to the basic forms.

\begin{theorem}
\label{hodge-birational-atoms-2026-theorem-fourier-nukari}
The Fourier--Mukai transform induces
an isomorphism
$$
H_{B;BC,c}(\widehat X^\bullet)
\simeq
H_{B;BC,s}(X^\bullet).
$$
\end{theorem}

\begin{lemma}
Fourier--Mukai gives an isomorphism
between the odd and even parts.
\end{lemma}

\section{Rationally inequivalent
points on complete intersections}
\label{hodge-birational-atoms-2026-section-olivier-martin}

\medskip\noindent
{\bf Speaker.}

Olivier Martin
(Institut de Mathematique Pure et
Appliquee)

\medskip\noindent
{\bf Title.}

Rationally inequivalent points on
complete intersections.

\medskip\noindent
{\bf Abstract.}

Let $X\subset \mathbb{P}^{n+k}$ be a
very general complete intersection of
multidegree $(d_1,\dots,d_k)$.
Inspired by work of Voisin, Chen--
Lewis--Shen conjectured in 2021 that
the set of points rationally
equivalent to, and different from, a
fixed point $x\in X$ has dimension at
most
$$
2n-\sum(d_i-1).
$$
The authors proved this conjecture for
hypersurfaces ($k=1$), and Riedl--
Yang proved it when
$$
\sum(d_i-1)\le 2n
$$
or
$$
\sum(d_i-1)\ge 2n+2.
$$
I will explain how to generalize the
technique of Riedl--Yang to tackle the
full conjecture and to obtain finer
information about rational
equivalence of points on complete
intersections.

Let $X$ be a smooth projective variety and
$Z_0(X)$ a free abelian group on the points
of $X$.
$\alpha,\beta \in Z_0(X)$ are {\it rationally
equivalent} if there exist
\begin{itemize}
\item $S_n$ projective curves $C_i$
\item morphisms $f_i:C_i \to X$
\item a rational functions $\phi_i \in
\mathbb{C}(t)$
\end{itemize}
such that
$$
\alpha = \beta= \sum I f(\text{div} \phi_i)
$$
which leads us to the Chow group of $X$:
$$
\text{CH}_0(X)
=Z_0(X)/\sim.
$$
\begin{example}
\label{hodge-birational-atoms-2026-example-rat-equiv}
\begin{itemize}
\item $x,y \in \mathbb{P}^1$,
$x\sim_{\text{rat}}y$,
$x, y \in X$ lie on a rational curve in $X$
then $x\sim_{\text{rat}}y$.
\item $S$ a K3 surface over $\mathbb{P}^2$
ramified over a smooth sextic $C \subset \mathbb{P}^2$.
Any two points on $\mathbb{P}^2$ are
rationally equivalent, this implies
that $2x'$ and $2y'$ are also 
rationally equivalent, it is a nontrivial
fact that actually $x'=f^{-1}(x)$ and $y'
=f^{-1}(y)$ are also rationally equivalent.
\item Rational equivalence has more to do
with
raional curves in $\text{Sym}^\vee (X)$
than on $X$ itself.
\end{itemize}
\end{example}

\noindent
Now consider
$$
\begin{aligned}
X  &\longrightarrow CH^0(X) \\
x &\longmapsto [x]
\end{aligned}
$$

\noindent
The fiber that contains $x$ is called the
{\it orbit} for rational equivalence of $x$ and
is denoted $R_{X,x}$.
This is a countable union of closed subsets.

\begin{theorem}[Mumord-Roitman 69-71]
\label{hodge-birational-atoms-2026-theorem-mumford-roitman}
If there exists a positive integer $0<p\leq
\dim X$ s.t. $H^\bullet(\Omega^p_X)\neq 0$
then for $x \in X$ very general
$\text{codim} R_{X,x} \geq p$.
In particular, if $K_X$ is effective then
$X$ is not covered by positive dimension
orbits.
\end{theorem}

\noindent
upshot: if you have forms, orbits are..
small?

\begin{theorem}[Voisin, 94-96]
\label{hodge-birational-atoms-2026-theorem-voisin9496}
If $X \subset \mathbb{P}^{n+k}$
is a very general complete intersection
of multidigree $(d_1,…,d_k)$
such that $\sum(d_i-1) \geq  2n-2$
then no two points on $X$ are rationally
equivalent to one another.
(The Chow map $X \to CH^0(X)$ is injective.)
\end{theorem}

\noindent
It was conjectured in 2001 by 
Chen-Lewiz-Shein (two from Alberta,
one from China) that
if $X \subset \mathbb{P}^{n+k}$
v.g. complete intersection of m. d.
$(d_1,…,d_k)$ then $\forall x \in X$,
$\dim R_{X,x} \leq  2n- \sum(d_i-1)$.

This would imply Theorem
\ref{hodge-birational-atoms-2026-theorem-voisin9496}.

\begin{remark}
\label{hodge-birational-atoms-2026-remark-rat-equiv-chow}
\begin{itemize}
\item CLS prove it for hypersurfaces
\item  Voisin's theorem
implies the conjecture
for $\sum(d_i-1) \geq 2n+2$.
\item Riedl-Yang 22 prove when
$\sum (d_i-1) \neq 2n$.
\end{itemize}
\end{remark}

\medskip\noindent
{\bf Expected theorem (M)}

Same setup: $X \subset \mathbb{P}^{n+k}$
very general complete intersection
of m.d. $(d_1,…,d_k)$ then 
$\dim \tilde{R}_X \leq  4n-2
\sum(d-1)$.

Theorem $\implies $ Conjecture,
$R_{X,x} \times R_{X,x} \subset R_X$.

\begin{example}
\label{hodge-birational-atoms-2026-example-very-general-ci}
If $X \subset \mathbb{P}^3$ is a 
very general quintic surface then CLS says
$R_{X,x}$ is a countable union for all x.
Theorem says that $\{x: R_{X,x}\neq \{x\}\}$
is also countable.
\end{example}

\medskip\noindent
{\bf Rieell-Yang method.}

\noindent
The idea is to use this method to 
prove the expected theorem.
This method should work for a wide
variety of other problems [list, including
hyperbolicity].

Let $B_n$ be the parameter space
for smooth complete intersections
in $\mathbb{P}^{n+k}$ of n. d.
$(d_1,…,d_k)$ and its universal family
$$
\xymatrix{
\mathfrak{X}_n\ar[d]\\
B_n.
}
$$

\begin{definition}
\label{hodge-birational-atoms-2026-definition-parametrized-lin-subsp}
A {\it parametrized $(n+k)$-linear
subspace} of $\mathbb{P}^{m+k}$
is an injective linear map
$\Lambda: \mathbb{P}^{n+k} \to
\mathbb{P}^{m+k}$.
The set of such is denoted $G_{n,m}$.
$\Lambda$ is called
{\it adapted} to a pointed
complete intersection $(X,x)$
if (1) x inn $\text{Im}(\Lambda)$
and (2) $\Lambda^2(X) \subset
\mathbb{P}^{n+k}$ is a smooth complete
intersection.
$$
\Lambda^*(X,x):=
(\Lambda^{-1}(X),\Lambda^{-1}(X) \in
\mathfrak{X}_n.
$$
\end{definition}

\noindent
The idea is to take a complete intersection
on the RHS and pull it back.

\begin{definition}
\label{hodge-birational-atoms-2026-definition-constructible-subvars}
Given $\mathcal{Z} 
\subset \mathfrak{X}_n$
a construcible subset, we let
$$
\mathcal{Z}^{+} =
\{(X,x) \in \mathfrak{X}_{n+1}:
\exists \Lambda \in G_{n,n+1}
\text{ adapted to }(X,x)
\text{ s.t. }
\Lambda^*(X,x) \in \mathcal{Z}\}.
$$
\end{definition}
\noindent
Somehow this is a natural
operation to study subsets of complete
intersections; or even
``towers'' of subvarieties.

The main property that they proved is

\begin{theorem}[RY]
\label{hodge-birational-atoms-2026-theorem-ry}
If $\mathcal{Z} \subset \mathfrak{X}_n$
is a constructible subset of 
positive codimension
then
$$
\text{codim}_{\mathfrak{X}_{n+1}}
\mathcal{Z}^+
< \text{codim}_{\mathfrak{X}_n}\mathcal{Z}.
$$
\end{theorem}

\noindent
Let's sketch how to use to prove
RY '22 theorem (see above).

\begin{theorem}[RY 22]
\label{hodge-birational-atoms-2026-theorem-ry22}
Let $X \subset \mathbb{P}$ be a very general
complete intersection w.d. $(d_1,…,d_k)$,
$\sum_i(d_i-1)<2n$ for all $x \in X$.
Then 
$$
\dim R_{X,x} \leq  2n- 
\sum (d_i-1).
$$
\end{theorem}

\begin{proof}
(Sketch.)
First can assume
$$
\sum(d_i-1) > n.
$$
Let $c:= \sum (d_i-1)-n-1$.
$K_{\mathfrak{X}_{n,t}}=\mathcal{O}(c)$,
which measures how far we are from a CY.

Suppose towards a contradiction that
the following dominates the base $B_n$
$$
\left\{(X,x): \dim R_{X,x} > 2n -
\sum_i(d_i-1)\right\}
$$
If this is the case, there exists
$\mathcal{Z}$ contained in the set above, 
irreducible with relative dimension
$> 2n - \sum_i (d_i-1)$ over $B_n$.
This implies that
$$
\begin{aligned}
\mathcal{Z}&\leq  n- \left(2n
-\sum_j(d_j-1)+1\right)\\
&=\sum_{j}(d_j-1)-n-1\\
&=c.
\end{aligned}
$$

\noindent
Then
$$
\text{codim}\mathcal{Z}^{+c}
\leq  \text{codim} \mathcal{Z} -c
=0,
$$
so $\mathcal{Z}^{+c}$ is dense,
so for $t \in B_{n+c}$ generic,
$\mathcal{Z}_t \subset \mathfrak{X}_{n+c,t}$
is dense.
$$
\mathcal{Z}_t \subset
\left\{x \in X:
\dim R_{X,p}>2n
-\sum_i (d_i-1)>0\right\}.
$$
i.e. $\mathfrak{X}_{n+c,t}$
is covered by positive dimension
orbits, contradicting
Mumford-Rotman.
\end{proof}

\medskip\noindent
{\bf A bipointed RY method.} 

$\mathcal{Z} \subset (\mathfrak{X}_x
\times_{B_n}\mathfrak{X}_n)$,
where the latter is a
{\it bipointed complete intersection}.

$$
\mathcal{Z} \subset
\left\{(X,x_1,x_2)
:x_1 \sim x_2,
x_1 \neq  x_2\right\}.
$$
The idea is that when it
``stabilizes'' the codimension decreases,
namely the naive hope is that
$\dim \mathcal{Z}^+ < \text{codim}
\mathcal{Z}$.
However, this is false. But in a controllable
way: we can understand when it fails.

\begin{example}
\label{hodge-birational-atoms-2026-example-situations-when-fails}
Suppose $\mathcal{Z}$ is the diagonal,
$$
\mathcal{Z}= \Delta=
\left\{(X,x,x)\right\} \subset
(\mathfrak{X}_n \times_{B_n}
\mathfrak{X}_n).
$$
When we stabilize the diagonal,
we get the diagonal again:
$$
\Delta^+_n = \Delta_{n+1}
$$
but the codimension on LHS is $n$,
and on the RHS is $n+1$.
\end{example}

\begin{example}
\label{hodge-birational-atoms-2026-example-another-when-fails}
Suppose
$$
\mathcal{Z} = 
\{(X,x_1,x_2)
: \text{span}(x_1,x_2)
\text{ is tangent to }X\}.
$$
In this case, the codimension
does not decrease! That 
would lead to a contradiction
using the tangent lines.
\end{example}

\begin{definition}
\label{hodge-birational-atoms-2026-definition-sp-equivalence}
$(X,x_1,x_2) \sim_{\text{}}
(X',x_1',x_2')$
if $\dim \text{span}(x_1,x_2)
= \dim \text{span}(x_1',x_2')$.
$$
\begin{aligned}
\exists\phi: 
\text{span}(x_1,x_2) &\longrightarrow 
\text{span}(x_1',x_2') \\
x_i &\longmapsto x_i
\end{aligned}
$$

$$
\phi^* :
I_{X'}|_{\ell'}
\xrightarrow{\sim}
I_X|_\ell
$$
\end{definition}

\noindent
Suppose, for the purpose of the
exposition, that the $\text{codim}
\mathcal{Z}^+< 
\text{codim} \mathcal{Z}$
does hold for all $\mathcal{Z}$.

We want to show
$$
\dim \tilde{R}_X
\leq  4n-2
\sum(d_i-1).
$$
Suppose not.
Then there exists
$$
\mathcal{Z}
\subset
\left\{(X,x_1,x_2):
x_1 \neq  x_2\right\}
$$
\begin{itemize}
\item irreducible
\item  dominates the base,
\item $\dim \mathcal{Z} >
4n-2 \sum (d_i-1)$.
\end{itemize}

\noindent
thus, $\text{codim}\mathcal{Z}< 2n-\left(
4n-2\sum(d_i-1)\right)
=2\sum(d_i-1)-2n$.

$$
\begin{aligned}
\text{codim}\mathcal{Z}^{+c}&\leq 
2 \sum(d_i-1)-1
-\sum(d_i-1)+n+1\\
&=\sum(d_i -1)-n\\
&=c-1.
\end{aligned}
$$

\noindent
If $pr_1:\mathcal{Z}^{+c}
\to \mathfrak{X}_{n+c}$
is dominant we win
because the fiber is positive dimensional.

But saturation and projections are
compatible.
$$
pr_1(\mathcal{Z}^+)
=(pr_1(\mathcal{Z}))^+
$$
So it's enough to show
$pr(\mathcal{Z}) \text{codim} \leq c$.
If it has $\text{codim} >c$,
the fixing $t \in B_n$
generic,
$pr_1(\mathcal{Z})$
has $\dim \leq  n-c$
so the fibers have
$\text{codim}$
in $\mathcal{Z}$ bigger than
…


% --- End hodge-birational-atoms-2026.tex ---

% --- Begin seminars-algebraic-geometry.tex ---
\chapter{Seminars in Algebraic Geometry}
\label{chapter-seminars-algebraic-geometry}









\section{Smoothable compactified Jacobians of
nodal curves}
\label{seminars-algebraic-geometry-section-smoot-compa-jacob-nodal}

\noindent
Nicola Pagani, University of Liverpool and
Bologna.
Seminar of Algebraic Geometry UFF.
August 20, 2025.

\medskip\noindent
{\bf Abstract.} Building from examples, we
introduce an abstract notion of a
'compactified Jacobian' of a nodal curve. We
then define a compactified Jacobian
to be 'smoothable' whenever it arises as the
limit of Jacobians of smooth
curves. We give a complete combinatorial
characterization of smoothable
compactified Jacobians in terms of some 'vine
stability conditions', which we
will also introduce. This is a joint work
with Fava and Viviani.

\medskip\noindent



Let $C$ be a smooth curve and $d \in
\mathbb{Z}$. Define
$$
J^d_C=\{L:\text{$L$ is a line bundle of
degree $d$}\}/\sim
$$
which is a smooth projective variety of
dimension $g(C)$.

If $C$ is nodal we still can consider
$J^d_C$.
\begin{enumerate}
\item One connected component. Then the
Jacobian is $\mathbb{P}^1$ minus two
points. This is not universally closed, so it
is not proper.
\item Two components intersecting at one
point.
The pullback of the normalization splits the
degree in
intintely many ways, givieng that $J^{-1}_C$
is an infinite set of points. This
is not of finite type, so it is not proper.
\item The curve has two components
intersecting at two points.
This gives $J^{-2}_C$,
 which is a mixture of the two former items.
 (Probably not proper too.)
\end{enumerate}

\medskip\noindent
Now consider
$$
\text{TF}_C^d=\{\mathcal{F}:
\text{coherent on $C$, torsion-free, rank-1
on $C$}\}/\sim
$$
This satisfies the existence oart of the valu
point of properness.

Now we consider the moduli. Now we consider
the ideal sheaf of the (singular?)
point(s?):
\begin{enumerate}
\item One component. The stack is proper!
\item Two components intersecting once. Now
we get stacky points,
$x=[\bullet/\mathbb{G}_m]$. These points have
generic stabilizer. The resulting
stack is not separated because a morphism of
a curve, say $\mathbb{P}^1$ minus a
point … there are infinitely many ways to
extend a morphism from this thing to a
line bundle. So you cannot include any of
these stacky points. Recall that a
sheaf is {\it simple} if its automorphism
group is $\mathbb{G}_m$.
\item  The ideal sheaf of both nodes
$\mathcal{I}(N_1,N_2)$ has a positive
dimensional automorphism
group. The stack is not proper.
\end{enumerate}

\begin{definition}
\label{seminars-algebraic-geometry-definition-fine-compactified-Jacobian}
A {\it fined compactified Jacobian} of $C$ is
an open connected substack of
$\text{TF}^d(C)$ that is also proper.
\end{definition}

\begin{remark}
\label{seminars-algebraic-geometry-remark-algebraic-space}
This thing is automatically an algebraic
space.
\end{remark}

\begin{definition}
\label{seminars-algebraic-geometry-definition-compactified-Jacobian}
A {\it compactified Jacobian} is an open
connected of $\text{TF}^d(C)$ that
admits a proper, good moduli space.
\end{definition}

Consider the Artin stack $\mathfrak{X}
\xrightarrow{\Gamma}X$ […]
 is a {\it good
moduli space} if
\begin{enumerate}
\item Every moduli factors
$$
\xymatrix{
\mathfrak{X} \ar[r]\ar[rd]&\mathcal{I}\text{
(ACC. space)}\\
& X
}
$$
\item
$\pi_*\mathcal{O}_{\mathfrak{X}}=\mathcal{O}
_X$.
\end{enumerate}

\medskip\noindent
We expect to find a notion of stability
condition to produce these things
[…] $[\bullet/ \mathbb{G}_m]$ would be the
polystable representative.

\begin{definition}
\label{seminars-algebraic-geometry-definition-smoot-compa-jacob}
A compactified Jacobian $\overline{J_C}$ is
{\it smoothable} if all smoothings
$\mathcal{C}\to \Delta=\{0,\eta\}$ (with
$\mathcal{C}_0=C$),
$$
J^d_{\mathcal{C}_\eta}\cup C \to
\overline{J_C}
$$
is proper.
\end{definition}

\medskip\noindent

\begin{definition}
\label{seminars-algebraic-geometry-definition-connected}
Let $X$ be a curve.
$$
\text{BCON}(X)=\{Y \subseteq X\text{s.t.
}Y,Y^c \text{ are connected}\}
$$
\end{definition}

\begin{definition}
\label{seminars-algebraic-geometry-definition-v-satbility}
A $v$-curve is a generalization of items (2)
and (3) in the lists above
[it looks like two long snakes $\sim$ that
intersect several times, and $t$ is
the number of nodes].
A {\it $v$-condition} is a pair $n=(n_1,n_2)$
such that
$$
n_1+n_2=\begin{cases}
d+1-t\qquad &\text{we say the s.c. is
nondegenerate} \\
d-t\qquad &\text{degenerate}
\end{cases}
$$
$\mathcal{F}$ on $X$ is  {\it
$n$-(semi)stable} if
$\text{deg}\mathcal{F}_{X_i}>n_i$
($\text{deg}\mathcal{F}_{X_i}\geq n_i$)
for $i=1,2$.
\end{definition}
$\mathcal{F}_{X_i}=\mathcal{F}|_{X_i}$
torsion.
$$
\text{deg}(\mathcal{F}_{X_i})+\text{deg}
(\mathcal{F}_{X_2})
=d-|\text{sing}(F)|.
$$
Then
$$
\overline{J_C}(n)=\{\mathcal{F}\text{ is
semistable}\},
$$
a smooth compact Jacobian.

\begin{definition}
\label{seminars-algebraic-geometry-definition-n-degeneracy}
A {\it degeneratation} of $v$-stab. on $X$ is
$n:\text{BCON}(X) \to \mathbb{Z}$
such that
\begin{enumerate}
\item
$$
n_Y+n_{Y^c}+|Y \cap Y^c|=
\begin{cases}
d+1\qquad &\text{ we say $Y$ is
$n$-nondegenerate} \\
d\qquad &\text{$Y$ is $n$-degenerate}
\end{cases}
$$
\item $Y_i$ no pa. common component
$n_{Y_1}+n_{Y_2}+\ldots$
$$
\xymatrix{
&Y_1\ar[dr]\ar[dl]\\
Y_2\ar[ur]\ar[rr]&&Y_3\ar[ll]\ar[ul]
}
$$
\end{enumerate}
\end{definition}

\begin{theorem}[-, et al]
\label{seminars-algebraic-geometry-theorem-bijec-stabi-condi-nodal}
(bijection between stability conditions and
nodal curves) The map
$$
\left\{ \substack{\text{sm. comp.} \\
\text{Jac of $X$}} \right\} \to
\left\{ \substack{\text{$v$-stab.} \\
\text{cond. of $X$}} \right\}
$$
$$
n\mapsto
\overline{J_X}(n)=\{n\text{-semistable
sheaves}\}
$$
is a bijection. (The arrow should be from
right to left!)
\end{theorem}

F. Viviani had proved it for fine compact
Jacobians.

\section{Equivariant spaces of matrices of
constant rank}
\label{seminars-algebraic-geometry-section-equiv-space-matri-const}

\noindent
Ada Boralevi, France.
Algebraic Geometry Seminar, IMPA.
August 27, 2025.

\medskip\noindent
{\bf Abstract.} A space of matrices of
constant rank is a vector subspace V, say
of dimension n+1, ofthe set of matrices of
size axb over a field k, such that
any nonzero element of V has fixed rank r. It
is a classical problem to look for
different ways to construct such spaces of
matrices. In this talk I will give an
introduction up to the state of the art of
the topic, and report on my latest
joint project with D. Faenzi and D. Fratila,
where we give a classification of
all spaces of matrices of constant corank one
associated to irreducible
representation of a reductive group.

\medskip\noindent

We are interested in vector spaces $U \subset
\text{Mat}_{m,n}(\mathbb{C})$,
with $m \leq n$, of {\it constant rank}, i.e.
such that for all $f \in U$,
$r:=\text{rank}f$ is the same.

Let $\ell(r,m,n):=\text{max}\dim U:U$ is of
rank $r$.

\medskip\noindent
{\bf Questions.}
\begin{enumerate}
\item $\ell(r,m,n)=$? In general not known.
\item Find relations among $\ell, r,m$ and
$n$.
\item Construction of examples and
classification.
\end{enumerate}

\begin{example}
\label{seminars-algebraic-geometry-example-what-we-know}
\begin{enumerate}
\item $\ell(1,m,n)=n$,
$$
\begin{pmatrix}
x_1&x_2&\cdots &x_n\\
\\
\\
\\
\end{pmatrix}
$$
\item $\text{rank}=2$? There are two cases
([Atkinson '83], [Eisenbud-Haus '88])
\begin{itemize}
\item Compression spaces,
$$
\begin{pmatrix}
*&*&\cdots &*\\
*&0&\cdots & 0\\
\vdots &\\
*&0&\cdots&0
\end{pmatrix}
$$
\item Skew-symmetric matrices of $3 \times
3$.
\end{itemize}
\item $\ell(r,m,n) \geq n-r+1$. Because you
can put a matrix of
$m\times n$ (with the first $r$ rows that can
have nonzero entries):
$$
\begin{pmatrix}
x_1 &  x_2 &  \cdots & & x_{n-r+1}\\
& x_1 & x_2 & \cdots &\\
\\
\\
\\
\end{pmatrix}
$$
\end{enumerate}
\end{example}

\begin{theorem}[Westwick '86]
\label{seminars-algebraic-geometry-theorem-Westwick}
\begin{enumerate}
\item $n-r+1 \leq \ell(r,m,n) \leq 2n-2r+1$.
\item If $n-r+1 \not|
\frac{(m-1)!}{(r-1)!}\implies \ell$
\end{enumerate}
\end{theorem}

We can see these spaces as (subvarieties?) of
determinantal varieties
$M_n=\{f \in
\text{Mat}_{m,n}(\mathbb{C}):\text{rank}(f)
\leq  r\}$.
[Interpretation via secant varieties inside
Segre embedding.]

\medskip\noindent
Consider a map $\varphi: U \to \Hom(V,W)$.
Then $\varphi \in U^* \otimes  V^* \otimes
W$.
We get that
$\varphi$ is of constant rank if and only if
some kernel, image and cokernel are vector
bundles.

\medskip\noindent
{\bf Focus of today.} What happens when
$U,V,W$ are irreducible
representations of a complex reductive group
$G$?

\medskip\noindent
{\bf Question.} What is the natural
equivariant morphism
$$
U \to \Hom(V,W)=V^* \otimes  W
$$
of constant rank?

\medskip\noindent
Consider the case of $G=\text{SL}_2$. All the
irreducible representations (which are
self-dual) are given by
$V(m-1)\cong
\mathbb{C}[x,y]_{\text{deg}=m-1}$.

Recall the Clebsh-Gordon decomposition ($m
\leq  n$)
$$
V(m-1)\otimes
V(n-1)=\bigoplus_{i=0}^{m-1}V(n-m+2i)
$$

\begin{theorem}[B-Faenzi-Lella '22]
\label{seminars-algebraic-geometry-theorem-BFL22}
$$
V(n+m-2)\hookrightarrow \Hom(V(m-1),V(n-1))
$$
is of constant rank (corank 1) if and only if
$$
n-m+2i|m-1
$$
\end{theorem}

\begin{theorem}[B. Faenzi, Fratila '25]
\label{seminars-algebraic-geometry-theorem-BFF25}
Let $V(\nu)$, $V(\mu)$ and $V(\lambda)$ be
irreducible representations
of a comple reductive group $G$, with
$$
\dim(V(\mu)) \leq \dim(V(\lambda))_-
$$
If there exists a morphism of representations
$$
\varphi: V (\nu) to \Hom (V (\lambda),V(\mu))
$$
then $\varphi$ is of constant corank 1 if and
only if
there exists a simple root $\alpha_i$ such
that
\begin{enumerate}
\item $\lambda=\mu+\nu-\alpha_i$,
\item $\nu$ is a multiple of $\nu$,
\item  $\nu$ is a multple of $\omega_i$
\end{enumerate}
\end{theorem}

\section{Irreducibility of the Hilbert scheme
of points and the class of 2-step
ideals}
\label{seminars-algebraic-geometry-section-irred-hilbe-schem-point}

\noindent
Michele Graffeo, UFF, September 29, 2025.
The notes, including the nested Hilbert
scheme discussion, are now in
Hilbert schemes, Section
\ref{hilbert-schemes-section-graffeo-two-step-ideals}.

\section{Normal and Restricted Tangent
Bundles of Rational
Curves in Hypersurfaces}
\label{seminars-algebraic-geometry-section-restricted-tangent-bundles}

\noindent
Lucas Mioranci, IMPA.
Algebra Seminar, IMPA.
October 8, 2025.

\medskip\noindent
{\bf Abstract.}
Let $X\subset \mathbb{P}^n$ be a degree $d$
hypersurface containing a smooth
rational curve $C$ of degree $e$. The normal
bundle $N_{C/X}$ and the restricted
tangent bundle $T_X|_C$ split as direct sums
of line bundles of the form
$\bigoplus_i
\mathcal{O}_{\mathbb{P}^1}(a_i)\qquad$ called
their splitting type.
The splitting type of $N_{C/X}$ and
${T_X}|_C\;\;$ controls the local structure
of the space of rational curves and
determines how many general points of $X$ we
can interpolate by deforming the curve $C$.
By combining explicit computations
of $T_X|_C\;\;$ and an induction argument, I
classify all triples $(e,d,n)$ such
that a general degree $d$ hypersurface
$X\subset \mathbb{P}^n$ contains a
rational curve $C$ of degree $e$ whose
restricted tangent bundle ${T_X}|_C\;\;$
is balanced. The case of quadrics is
particular, in which we show that
odd-degree rational curves do not interpolate
the expected number of points on
quadric hypersurfaces.  For the normal
bundle, I compute explicit examples of
hypersurfaces $X$ for all possible splitting
types of $N_{C/X}$ when $C$ is the
rational normal curve. Additionally, for
$d\ge 3$, we compute the dimension of
the space of hypersurfaces $X$ such that
$N_{C/X}$ has a given splitting type.



\medskip\noindent

$\mathcal{O}$ with no subindex means
$\mathcal{O}_{\mathbb{P}^1}$.
$$
\xymatrix{
&0\ar[d]&0\ar[d]\\
&T_{\mathbb{P}^1}=\mathcal{O}_(2)\ar[d]\ar[r]
^=
&T_{\mathbb{P}^1}=\mathcal{O}(2)\ar[d]\\
0\ar[r]&T_{X}|_{C}\ar[d]\ar[r]&T_{\mathbb{P}
^n}|_{C}\ar[d]\ar[r]&
N_{X/\mathbb{P}^n}|_{C}\cong
\mathcal{O}(de)\ar[d]_=\ar[r]&0\\
0\ar[r]&
N_{C/X}\ar[d]\ar[r]&N_{C/\mathbb{P}^n}\ar[d]
\ar[r]&
N_{X/\mathbb{P}^n}|_{C}\cong
\mathcal{O}(de)\ar[r]&0\\
&0&0
}
$$
Throughout this talk
$X$ is a degree $d$ hypersurface in
$\mathbb{P}^n$
and $C \subset X$ is a smooth rational curve
of degree $e$.

Since $C$ is rational we can think of it
as a map $f:\mathbb{P}^1 \to X$.
Then we can pullback any bundle over $C$
back to $\mathbb{P}^1$,
and by Birkhoff-? Theorem
we know that vector bundles over
$\mathbb{P}^1$
split as
$$
E\cong \bigoplus_{i=1}^r
\mathcal{O}_{\mathbb{P}^1}(a_i),
\qquad a_1\leq \ldots \leq a_n.
$$
This decomposition is called the
{\it splitting type} of $E$.
We say $E$ is {\it balanced} if
$|a_i-a_j|\leq 1$ for all $i,j$.

Balancedness is an open condition,
i.e. if you have a family of bundles over
$\mathbb{P}^1$
and one of them is balanced, then all of them
are,
(also you can think it's a generic
condition).

We know the normal bundle $N_{C/X}$ has rank
$n-2$
and degree  $e(n-d+1)-2$.
And the tangent bundle $T_{X/C}$
has rank $n-1$ and degree $e(n-d+1)$.

\begin{example}
\label{seminars-algebraic-geometry-example-conic}
A conic $C \subset \mathbb{P}^3$.
Since a conic in $\mathbb{P}^3$ is contained
in a plane,
this obliges to have a ``distinguished normal
direction'',
$\mathcal{O}(2)$. Indeed,
$$
\begin{aligned}
N_{C/\mathbb{P}^3}&\cong(N_{Q/\mathbb{P}^3}
\oplus
N_{\mathbb{P}^2/\mathbb{P}^3}|_{C}\\
&\cong(\mathcal{O}(2) \oplus
\mathcal{O}(1))|_{C}\\
&\cong\mathcal{O}(1)\oplus \mathcal{O}(2).
\end{aligned}
$$
\end{example}

\medskip\noindent
Now we discuss interpolation of general
points.
Let $p_1,\ldots,p_n \in \mathbb{P}^1$
and $x_1,\ldots, x_n \in X$
be general points.
Let $f:\mathbb{P}^1 \to X$ be such that
$f(p_i)=x_i$.

Consider the space of morphisms mapping
the  $p_i$ to $x_i$. It's known that its
tangent space
at $f$
is given by
$$
T_{[f]}\Mor(\mathbb{P}^1,X,p_i \mapsto x_i)
\cong H^0(\mathbb{P}^1,f^*T_X(-n)).
$$
\noindent
Deformations of $f$ (mapping $p_i \mapsto
x_i$)
dominate $X$ if $a_i-n_i \geq 0$.
Also, deformations of $f$ interpolate up
to $a_1+f$ general points in $X$.

\medskip\noindent
How many points can deformations of the curve
interpolate?
This means (I think) that the deformations of
the curve
are still such that $p_i \mapsto  x_i$.
``And for the normal bundle, interpolations
means that the curves contain the points''.

If $T_X|_{C}$ is balanced, $C$ can
interpolate up to
$$
\left\lfloor \frac{e(n+1-d)}{n_1}
\right\rfloor +1
$$
general points. The case of $N_{C/X}$ is
similar.

\medskip\noindent
Next we survey some relevant results.

\begin{theorem}[Coskun-Riedl, 2018]
\label{seminars-algebraic-geometry-theorem-CR18}
A general Fano $X$ of degree $d \geq 2$
contains a degree $c$ rational smooth curve
 $C$ with balanced $N_{C/X}$
for every $1 \leq e \leq  n$.
\end{theorem}

The idea is to observe that the space of
curves has
balanced normal bundle.

\begin{theorem}[Ran, 2021]
\label{seminars-algebraic-geometry-theorem-R21}
Extends the above for
$1 \leq  e \leq 2n-2$
and $d \geq 4$.
\end{theorem}

Idea: general hypersurface contains a curve
of balanced normal bundle.

\begin{theorem}[Ran, 2024]
\label{seminars-algebraic-geometry-theorem-R24}
$X$ general Fano hypersurface.
There exist smooth rational curve with  $C$
with balanced tangent bundle $T_{X}|_{C}$
for $e$ in some arithmetic progressions.
\end{theorem}

$e$ is in general very large.
They show there are curves of arbitrarily
large
degrees with balanced tangent bundle.

\medskip\noindent
Now we describe the new results.

\begin{theorem}
\label{seminars-algebraic-geometry-theorem-Lucas}
Let $X \subseteq \mathbb{P}^n$ be Fano
hypersurface of degree $3 \leq d \leq n$.

\begin{enumerate}
\item If $C \subset X$ is a rational smooth
curve
of degree $e$ with
$e \leq  \frac{n-1}{n+1-d}$,
then $T_X|_{C}$ is not balanced.

\item A general hypersurface $X$
contains rational curves $C$
of degree $e$ with balanced
$T_X|_{C}$ for every
$e> \frac{n-1}{n+1-d}$.
\end{enumerate}
\end{theorem}

\begin{theorem}
\label{seminars-algebraic-geometry-theorem-lucas2}
Let $X \subseteq \mathbb{P}^n$ be a smooth
quadric hypersurface.

\begin{enumerate}
\item If $e \leq 2$ is even,
then exists a curve $C$ of degree $e$
with $T_X|_{C}\cong \mathcal{O}(e)^{n-1}$.
(I.e. we can interpolate $e+2$ points
if $e$ is even.)

\item If $e \geq 1$ is odd,
then there is no curve $C$ of degree $e$
with balanced $T_X|_C$.
The best we can do is the most balanced
bundle
before balancedness, namely
$T_X|_C \cong
\mathcal{O}(e-1)\oplus\mathcal{O}(e)
\oplus\mathcal{O}(e+1)$.
(I.e. we can only interpolate $e$ points if
$e$ is odd.)
\end{enumerate}
\end{theorem}

Next we give the main ideas in the proof.
Let $C$ be the rational normal curve
of degree $e$ in $\mathbb{P}^1$,
$$
\begin{aligned}
C:\mathbb{P}^1  &\longrightarrow \mathbb{P}^n
\\
(s:t) &\longmapsto
(s^2:s^{e-1}t:\ldots:t^e:0:\ldots:0).
\end{aligned}
$$
$X=V(F)$.

Then find some examples. Then,

\begin{proposition}
\label{seminars-algebraic-geometry-proposition-proposition}
If $\mu(T_X|_{C})=\frac{e(n+1-d)}{n-1}\leq 1$
then $T_X|_{C}$ is not balanced.
\end{proposition}

\noindent
Use the proposition to show that $T_X|_C$
is not balanced if $d \leq  n$
and $e le \frac{ n-1}{n+1-d}$.

This implies Theorem \ref{seminars-algebraic-geometry-theorem-Lucas}.

\begin{lemma}
\label{seminars-algebraic-geometry-lemma-split-normal-bundle}
If $N_{C/X}\cong \bigoplus_i
\mathcal{O}(a_i)$
with $a_i < 4$ for all $i$,
then $T_{X}|_C \cong N_{C/X} \oplus
\mathcal{O}(2)$.
\end{lemma}

\noindent
This implies that
$\text{Ext}^1(N_{C/X},\mathcal{O}(2))=0$.

The most important step is the following,
used as the induction step:

\begin{proposition}
\label{seminars-algebraic-geometry-proposition-induction}
If $T_Y|_C$ is balanced for some degree $d$
hypersurface $Y \subseteq \mathbb{P}^{n-1}$,
then there exists a degree $d$ hypersurface
$C \subseteq \mathbb{P}^n$ with balanced
$T_X|_C$.
\end{proposition}

\noindent
Combining the examples, Lemma
\ref{seminars-algebraic-geometry-lemma-split-normal-bundle}
and Proposition \ref{seminars-algebraic-geometry-proposition-induction},
we establish that for $e \leq
\text{max}\{2d-2,n\}$.

The following argument is used to show that
actually what we showed so far is enough for
all degrees.

\begin{lemma}
\label{seminars-algebraic-geometry-lemma-smoothing-argument}
$C=C_1 \cup  C_2$, $C_1 \cong \mathbb{P}^1$.
$E$ vector bundle on $C$ such that
$E|_{C_1}$ is balanced,
$E_{C_2}$ is perfectly balanced
(numerical condition that works when
$e=n-1$).
Then if $E$ is the specialization of a family
of vector bundles $E'$ on $\mathbb{P}^1$,
then $E'$ is balanced.
\end{lemma}

\noindent
Basically, glue and smooth low degree curves
(i.e. $e \leq  \text{max}\{2d-2,n\}$)
to get all degrees


\medskip\noindent
Now we prove the result for quadrics,
i.e. $X$ is a quadraic.
For even $e$ compute examples and conclude
 $T_X|_C \cong \mathcal{O}(e)^{n-1}$.

For odd $e$, compute examples,
$T_X|_C\cong \mathcal{O}(c-1) \oplus
\mathcal{O}(2)^{3}\oplus \mathcal{O}(-1)$.

\begin{proposition}
\label{seminars-algebraic-geometry-proposition-interpolation2}
$n\geq 3$, $m\leq  e+1$.
There is a degree $e$ curve $\phi_e$
interpolating
$m$ general points if and only if there is a
degree $e-2$ curve $\psi_{e-2}$ interpolating
$m-2$ general points in $X$.
\end{proposition}

\noindent
The idea is to do induction.
Suppose you can interpolate $e$ points,
then show you can interpolate $e+2$.

Chose $y_1,\ldots,y_{m-2}$ on the quadric.
Choose a line in $\mathbb{P}^{m-1}$
intersecting the quadric,
and $x_1,\ldots,x_{m-1}$ points on the line.
Join the points on the line to the points on
the quadric,
Then you get somem points of intersection
with the quadric.
Those points and the two points of
intersection of the line and the quadric
let us construct a scroll, which intersects
the quadric.
The intersection of the scroll and the
quadric
(I guess the scroll contains the lines, and
so the points $y_i$)
gives the curve passing through the $y_i$.

\section{Regular but non smooth curves of
genus 3}
\label{seminars-algebraic-geometry-section-regular-but-non-smooth}

\noindent
Cesar Hilario, .
Algebra Seminar, IMPA.
October 15, 2025.

\medskip\noindent
{\bf Abstract.} Regular but non-smooth curves
are a unique feature of
geometry in positive characteristic, that
results from the fact that over an
imperfect field the notion of regularity is
weaker than the notion of
smoothness. In the setting of algebraic
geometry over an algebraically closed
field, these curves correspond to fibrations
by singular curves, which are
fibrations of relative dimension 1 whose
fibers are singular. The most famous
examples are arguably the so-called
quasi-elliptic curves and quasi-elliptic
fibrations, which play a key role in the
Bombieri-Mumford classification of
algebraic surfaces in characteristics 2 and
3. In this talk I will discuss the
case of genus 3 in characteristic 2, with an
eye towards the classification of
regular plane projective quartic curves that
become rational after base change.

\medskip\noindent
arXiv 2409.05464.

\medskip\noindent
$C$ will denote a projective geometrically
integral curve over
a not necessaily closed field $K$.
Geometricallt integral
means that $C \otimes_K \overline{K}$ is
integra.
Suppose that $p=\text{char}(K) \geq 0$
and $g=h_1(\mathcal{O}_C)$ is of genus $C$.

We define
\begin{enumerate}
\item $C$ regular iff local rings of $C$ are
regular (DVR).
\item $C$ smooth iff $C$ regular and $K
\otimes_K \overline{K}$ regular.
\end{enumerate}

\noindent
Thus it's obvious that smoothness implies
regularity.
In this talk we explain why the converse is
not true.

Note that if $K$ is perfect (e.g. $p=0$ or
$K=\overline{K}$) implies
 smooth = regular. Recall that $K$ is
 imperfect iff $p>0$
the powers of $K$, $K^p$, is a proper
subfield of $K$
(it's always a subfield but we ask it's
proper).

\medskip\noindent
Assume that $C$ is regular and $p>0$.
The genus of rthe normalization is less than
or equal to the genus of the curve,
i.e.
$$
\tilde{g}=h^1(\mathcal{O}_{\widehat{C
\otimes_K \overline{K}}})\leq g.
$$
We know that $C$ is smooth iff
$\overline{g}=g$.
So $g-\overline{g}$ is a measure of
non-smoothness.

\begin{theorem}[Tate]
\label{seminars-algebraic-geometry-theorem-tate}
$\frac{p-1}{2}$ divides $g-\overline{g}$.
\end{theorem}

As a corollary, we get that if $C$ is
non-smooth,
$g-\overline{g}>0$ and so $p \leq  2g+1$.

\begin{example}
\label{seminars-algebraic-geometry-example-attain}
There are cases in which the bound is sharp,
namely $C:y^2=x^p+t$, with $t \in K\setminus
K^p$.
Then $g=\frac{p-1}{2}$ and $\bar{g}=0$.
Over $\overline{K}$ rhere exists a singular
point,
$(x,y)=(-t^{1/p},0)$ which
{\bf is not visible over $K$}, that is we get
that
$t^{1/p}\not \in K$.
\end{example}

\begin{example}
\label{seminars-algebraic-geometry-example-elliptic}
$\{\text{elliptic curves}\}=\{g =
\overline{g}=1,\text{$\exists K$-rational
point }\}=\{\text{smooth cubic curves with
$K$-rational point}\}$.
$y^2=x^3+4g_i x+27$.
\end{example}

\begin{example}
\label{seminars-algebraic-geometry-example-quasielliptic}
$g=1$, $\overline{g}=0$, $p=2,3$,
quasielliptic curves. (Important in the
classification of surfaces
in positive characteristic.)
\end{example}

\medskip\noindent
Now we make a digression to describe
Kodaira-Enriques (over $\mathbb{C}$)
in positive characteristic. Bombieri-Mumford:
{\bf new objects} appear,
quasielliptic surfaces.

Let $k=\overline{k}$, $S$ smooth elliptic
({\it quasielliptic}) surface over $k$.
Suppose we have a fibration
where the general fiber is an elliptic curve
an elliptic curve (resp. a plane
cuspidal cubic) and the generic fiber an
elliptic (resp. quasielliptic)
curve over $k(B)$.

$$
\xymatrix{
S\ar[d]\\
B
}
$$

\medskip\noindent
Back to our main content, let's assume that
$g=3$.
$$
\xymatrix{
\{g=1,\exists K\text{-rat. pt}\}\ar[r]&\ar[l]
\left\{\substack{\text{reg. cubic curves} \\
\text{with $K$-rat pt}}\right\}
}
$$
Recall that genus-degree formula
$$
g=\frac{(d-1)(d-2)}{2}=1,
$$
So then $d=3,4$ gives  $g=3$ (right?).

No we present our main theorem, which is a
classification result:

\begin{theorem}
\label{seminars-algebraic-geometry-theorem-classification-quartics}
\begin{itemize}
\item $C$ regular non-hyperelliptic curve
over $K$.
\item  $g=3$, $\overline{g}=0$.
\item $p=2$.
\end{itemize}
(Not that non-hyperllipticity of $C$ and
$g=3$ imply that $C$ is a quartic
over $K$, and that  $\overline{g}=0$ implies
$C$ is geometrically rational.)
Then $C$ is isomorphic to one of the
following quartics.

\begin{enumerate}
\item $y^4+az^4+xz^3+bx^2z^2+cx^4=0$ where
$a,b,c \in K$, $c \not \in K^2$.
(Notice the imperfectness of $K$ is crucial.)

\item $y^4+az^4+bx^2y^2+cx^2z^2+bx^3z+dx^4=0$
for $a,b,c,d, \in K$,
$a \not \in K^2$, $b\not =0$.

\item
\item
\item
\end{enumerate}
\end{theorem}

\noindent
Now we enumerate properties of these
families.

\begin{itemize}
\item $C$ is a purely inseparable double
cover of a quasielliptic curve.
To understand this consider the induced map
from
Frobenius map $K \to K$, $a \mapsto a^p$ and
do base change:
$$
\xymatrix{
C^{(p)}\ar[r]\ar[d]\ar@{}[dr]|-{\lrcorner}&
C\ar[d]\\
\Spec K\ar[r]^{\substack{\text{induced from}
\\ \text{Frobenius}}}&\Spec K
}
$$
This gives a map $C \to C^{(p)}$ by universal
property of pullback.

$$
\xymatrix{
C \ar[rrrd]^{\text{Frobenius}}
\ar@{-->}[rrd]\ar[rrdd] & & \\
& &  C^{(p)}\ar[d] \ar[r] & C \ar[d] \\
& &  \Spec K\ar[r] & \Spec K
}
$$

Then consider the normalization of
$C^{(p)}=C_1$ which
turns out to be quasielliptic; by universal
property of normalization we can lift:
$$
\xymatrix{
& \widehat{C^{(p)}}\ar[d]\\
C\ar@{.>}[ur]\ar[r]&C^{(p)}.
}
$$
Here the lower arrow is purely inseparable of
degree 2.

\item There exists a unique non smooth point
$x \in C$.

\item Now we iterate the construction of the
first item:
$$
\xymatrix{
\underbrace{C}_{g=3}\ar[r]&\underbrace{C_1}
_{g=1}\ar[r]&
\underbrace{C_2}_{g_2-9}\ar[r]&
\underbrace{C_3}_{g_3=0}\ar[r]&
\underbrace{C_4}_{g_4=0}\ar[r]&\cdots
}
$$
Let $x$ be the non-smooth (non-rational)
point of $C$. Map it under
these maps to $x_1$( non-smooth,rational),
then  $x_2$ (smooth, rational?),
then $x_3 $ (smooth, rational.), then $x_4$
(smooth, rational), etc.

We have
$$
\begin{matrix}
&\text{$x_2$ rational}& \text{$x$ canoical
divisor} &
E=K(c_2)\\
(i)& Y & Y & Y\\
(ii) & N & Y & N\\
(iii) & Y & N & N\\
(iv) & N & N & Y\\
(v) & N & N & N
\end{matrix}
$$
Where the last column is explained next.

\item Canonical field of $C$ = subfield of
$K(C)$ generated
by the quotients of all non-zero holomorphic
differentials
=  $K(C)$ ($C$ non-hyperelliptic).

\item Pseudocanonical field of $C$ = …
non-zero \underline{exact}
holomorphic differential := $E$. So one has
$[K(c),E]=4=p^2$.

$$
\xymatrix{
&K(C)\ar[ld]_{4}\ar[rd]^{4}\\
E&&K(C_2).
}
$$
\end{itemize}

\noindent
In conclusion, we manage to characterize
these families.

\begin{example}
\label{seminars-algebraic-geometry-example-chain-of-maps}
In case 1, let $a=b=0$. We get
\begin{equation}
\label{seminars-algebraic-geometry-equation-curve-pencil}
C: y^4+xz^3+\underbrace{c}_{\in K\setminus
K^2}x^4=0
\end{equation}
$$
\underbrace{C_1}_{\text{quasiell.}}:xy^2+z^3+
cx^3=0.
$$
Further
$$
\xymatrix{
C\ar[r]&C_1\ar[r]&C_2 \cong\mathbb{P}^1\ar[r]
&C_3 \cong \mathbb{P}^1\ar[r]&C_4\cong
\mathbb{P}^1\ar[r]&\cdots.
}
$$
\end{example}

\medskip\noindent
With $C$ one can construct a pencil
(fibration over $\mathbb{P}^1$)
of quartics.
Let $k=\overline{k}$ (new base field from
what we fixed at the beginning!).
Idea: replace $c$ for $t$ Equation
\ref{seminars-algebraic-geometry-equation-curve-pencil}.
 $K=k(\mathbb{P}^1)$,
 $S=V(t_0(y^4+xz^3)+t_1x^4)
\subseteq \mathbb{P}^2_{(x:y:z)}\times
\mathbb{P}_{(t_0:t_1)}$.
This has a singular point, so we blow up:
$$
\xymatrix{
\tilde{S}\ar[d]\ar@/^1em/[dd]^f\\
S\ar[d]\\\mathbb{P}^1
}
$$
The generic fiber is the curve
$y^4+xz^3+tx^4=0$,
so that $t_1/t_0 \in k(\mathbb{P}^1)=K$.
Fibers are plane rational quartics,
and the singular fiber $f_1(0:1)$ is a
configuration of lines
given by a Dynkin diagram
$$
\xymatrix{
E_1 \ar@{-}[r]& E_2\ar@{-}[r]
& E_3 \ar@{-}[r]& E_4\ar@{-}[d]\ar@{-}[r]
& \cdots \ar@{-}[r]& E_{15}\\
& & & E
}
$$
We do the same with $C_1$: define
$S^1=V(t_0(xy^2+z^3)+tx^3) \subseteq
\mathbb{P}^2 \times \mathbb{P}^1$.
Then $S_1$ is also singular, and we blow up
to obtain a smooth
quasielliptic fibration ($\tilde{S^1}$ is
quasielliptic).
The fibers are plane cuspidal cubics.
The singular fiber an arrangement of curves
$$
\xymatrix{
F_1 \ar@{-}[r]& F_2\ar@{-}[r]&
F_3\ar@{-}[d]\ar@{-}[r]& \cdots \ar@{-}[r]&
F_8\\
& & F
}
$$

In fact, there is a correspondence between
the
singular fibers of the two fibrations.

\section{Variations on a theorem of Baum and
Bott}
\label{seminars-algebraic-geometry-section-baum-and-bott}

\noindent
Stéphane Druel, Université Claude Bernard
Lyon.
Algebra Seminar, IMPA.
November 5, 2025.

\medskip\noindent
{\bf Abstract.}
In this talk, I will discuss an analogue of
the Baum-Bott vansihing theorem for
the tangent bundle of the space of leaves of
regular foliations together with
applications to transversely projectively
flat foliations of dimension one. This
is work in progress.

\medskip\noindent



Baum-Bott 70'.

Let $E$ be a vector bundle on $X$
complex manifold is a sheaf of 1-jets of $E$:
$$
J'_X(E)= E \oplus \Omega^1_X \otimes E
$$
where equality is the category of sheavs of
abelian groups.
Where the $\mathcal{O}_X$ structure is given
by
$$
\underbrace{f}_{\substack{\text{local
function} \\ \text{on }U}}
(\underbrace{e}_{\substack{\text{local
section} \\ \text{of $E$ on  $U$}}}
, \underbrace{\alpha}_{\substack{\text{local
section} \\
\text{of $\Omega^1_X \otimes E$}}})
=(fe,f\alpha-df \otimes e)
$$
The {\it Atiyah class} of $E$ is the element
$$
\text{at}_X(E) \in H^1(X,\Omega_X^1 \otimes
\text{End}(E))
$$
corresponding to the exact sequence of
$\mathcal{O}_X$-modules
$$
\xymatrix{
0\ar[r]&\Omega_X^1\otimes
E\ar[r]&J_X^1(E)\ar[r]&E\ar[r]&0.
}
$$

\begin{remark}
\label{seminars-algebraic-geometry-remark-holomorphic-connection}
 $\text{at}(E)=0$ if and only if there exists
$\nabla:E \to \Omega^1_X \otimes E$ map of
abelian sheaves
such that $e \mapsto  (e-\nabla e)$ is
$\mathcal{O}_X$-linear
$$
-\nabla fe=-f\nabla e - df \otimes e
$$
if and only if $E$ admits a holomorphic
connextion.
\end{remark}

\begin{example}
\label{seminars-algebraic-geometry-example-line-bundl-holom-conne}
If $E$ is a line bundle, at
$(E)=(d\text{log}f_{ij})_{ij}
\in H^1(X,\Omega_X^1)$, $E=(f_{ij})$.
\end{example}

\begin{theorem}[Atiyah, '57]
\label{seminars-algebraic-geometry-theorem-atiyah-57}
Let $r =\text{rank}E$, $k$ an integer, $X$ be
compact Kähler
and $\Sigma_k$ a polynomial of degree $k$ on
$M_r$
such that  $\Sigma_k(g)$ is the $k$-th
elementary
symmetric function on the characteristic
zeros of
$g \in M_r$. $S_k=$ polynomial given by
polarization.
$$
\begin{aligned}
M_r^{\otimes k}&\to \mathbb{C}\\
\text{End}E^{\otimes k}& \to \mathcal{O}_X\\
\text{Ex., $k=1$}\qquad
\text{End}E&\xrightarrow{\text{Tr}}
\mathcal{O}_X.
\end{aligned}
$$

$$
\xymatrix{
H^1(X,\Omega_X^1 \otimes
\text{End}(E))^{\otimes k}\ar[r]^{\cup }
&H^k((\Omega^1_X)^{\otimes k}\otimes
\text{End}E^{\otimes k}\ar[r]
&H^k(X,\Omega_X^k)\ar[r]
&H^{2k}(X,\mathbb{C})\\
\text{at}(E)^{\otimes k}\ar[rrr]
&&&(2\pi i)^kc_{2k}(E).
}
$$
(Double check the term $(\Omega^1_X)^{\otimes
k}$…
I think there's a typo.)

Partial connection. $\mathcal{F}$ regular
distribution.
$\mathcal{F} \subseteq T_X$ subbundle.
An {\it $\mathcal{F}$-connection} on $E$ is
$\nabla:E \to \mathcal{F}^\vee \otimes E$
such that
$$
\nabla(fe)=f\nabla e+f_{\mathcal{F}}f \otimes
e
$$
where
$d_{\mathcal{F}}:\mathcal{O}_X\xrightarrow{d}
\Omega_X^1
\to\mathcal{F}^\vee$.
\end{theorem}

\begin{proposition}[Bott-Baum, '70]
\label{seminars-algebraic-geometry-proposition-bott-baum}
If $E$ admits an $\mathcal{F}$-connection
then all characteristic
classes of degree  $>\text{codim}\mathcal{F}$
vanish.
\end{proposition}

\begin{proof}
There exists an $\mathcal{F}$-connection if
and only if
$$
\xymatrix{
0\ar[r]&\Omega_X^1 \otimes
X\ar[r]\ar[d]&J_X^1\ar[r]&E\ar@{=}[d]\ar[r]&
0\\
0\ar[r]&\mathcal{F}^\vee\otimes E\ar[r]&
\mathcal{F}^\vee \otimes E \oplus E\ar[r]&
E\ar[r]&0
}
$$
if and only if $\text{at}(E) \in
H^1(X,\Omega_X^1\otimes\text{End}E)$
maps to 0 in
$H^1(X,\mathcal{F}^\vee\otimes\text{End}E)$
is and only if $\text{at}(E)$ comes from
$$
H^1(X,N^\vee\otimes\text{End}E)\to
H^1(X,\Omega_X^1\otimes\text{End}E)
$$
where $N=T_{X/\mathcal{F}}$.
\end{proof}

\begin{example}[Bott's partial connection]
\label{seminars-algebraic-geometry-example-bott-partial-connection}
Suppose $\mathcal{F}$ is a regular foliation
and $N=T_{X/\mathcal{F}}$ has a
$\mathcal{F}$-connection.
Take $n$ a local section of $N$,  $v$ a local
section of $\mathcal{F}$
(over the same open) and suppose exists a
section $t$ of $T_X$
that is a lift of $t$, i.e. $p(t)=n$ for
$p:T_X \to N$.
Then $\nabla_v^Bn=p([v,t])$ is a section of
$N$.
This defines an $\mathcal{F}$-connection on
$N$.
\end{example}

\begin{lemma}[Beauville]
\label{seminars-algebraic-geometry-lemma-beauville}
Let $X$ be a complex manifold and $E$ is a
direct submanifold of $T_X$.
Then $\text{at}(E) \in H^1(X,\Omega_X^1
\otimes \text{End}E)$
comes from $H^1(X,E^\vee \otimes
\text{End}(E))$.
\end{lemma}

\begin{proof}
Consider
$$
T_X= E \oplus F \xrightarrow{p}E.
$$
We claim that there exists an
$\mathcal{F}$-connection on $E$.
Let $v$ be a local section of $F$,
$e$ a local section of $E$.
Define $\nabla_ve=p([v,e])$.
One checks that this satisfies the required
properties.
\end{proof}

\noindent
Notice that no involutivity is imposed on $E$
- it's just a distribution.

\medskip\noindent
Now we will discuss some variations of
Baum-Bott's theorem
\ref{seminars-algebraic-geometry-theorem-baum-bott}.

Let $\mathcal{F}$ be  regular foliation, $E$
a vector bundle on $X$ equipped with a flat
$\mathcal{F}$-connection.
For example, $\nabla^B$ flat.

Recall that  $(E,\nabla)$ can be thought of
as
$\text{Mon}(\mathcal{F})$-equivariant vector
bundles
where $\text{Mon}(\mathcal{F})$ is the
monodromy groupoid.

Let $X/\mathcal{F}$ be the space of leaves,
and
$\mathcal{O}_{X/\mathcal{F}}=\mathcal{O}
_X^{d_{\mathcal{F}}}=
\{\text{functions $f$ such that
$d_{\mathcal{F}}f=0$}\}$.

It is a fact that $E^{\nabla}$ is a locally
free $\mathcal{O}_{X/\mathcal{F}}$
sheaf of rank the rank of $E$ and such that
$E=E^{\nabla}\otimes_{F_{X/\mathcal{F}}}
\mathcal{O}_X$.

Now one defines bundles of 1-forms and
tangent bundle as follows.
 $\Omega_{X/\mathcal{F}}^1=(N_{\mathcal{F}}^*
 )^{\nabla B}$
and
$T_{X/\mathcal{F}}=N_{\mathcal{F}}^{\nabla
B}$,
which are locally free sheaves of
$\mathcal{O}_{X/\mathcal{F}}$-modules
of rank $q=\text{codim}\mathcal{F}$.

The Atiyah class of $(E,\nabla)$ is
$$
J^1_{X/\mathcal{F}}(E,\nabla)=E^{\nabla}
\oplus
\Omega_{X/\mathcal{F}}^1
\otimes_{\mathcal{O}_{X/\mathcal{F}}}
E^{\nabla}
$$
with $\mathcal{O}_{X/\mathcal{F}}$-module
structure given by
 $$
(e,\alpha)=(f\alpha,f\alpha-\underbrace{df}
_{\in
\Omega^1_{X/\mathcal{F}}}
\otimes e.
$$
Then
$$
\text{at}_{X/\mathcal{F}}(E,\nabla)
\in
H^1(X,\Omega_{X/\mathcal{F}}
^1\otimes_{\mathcal{O}_{X/\mathcal{F}}}
\text{End}_{\mathcal{O}_{X/\mathcal{F}}}
E^{\nabla})
$$
corresponding to
$$
\xymatrix{
0\ar[r]&\Omega_{X/\mathcal{F}}
^1\otimes_{\mathcal{O}_{X/\mathcal{F}}}\ar[r]
&J_{X/\mathcal{F}}^1(E)\ar[r]&E^{\nabla}
\ar[r]&0.
}
$$

\begin{lemma}
\label{seminars-algebraic-geometry-lemma-atiya-class-maps-atiya-class}
$\text{at}_{X/\mathcal{F}}(E,\nabla)$ maps to
$\text{at}_X(E)$ under
$$
H^1(X,\Omega_{X/\mathcal{F}}
^1\otimes_{\mathcal{O}_{X/\mathcal{F}}}
\text{End}_{\mathcal{O}_{X/\mathcal{F}}}
E^{\nabla}
\to H^1(X,\Omega_X \otimes \text{End}E).
$$
\end{lemma}

\noindent
Partial connections. Let $G \subseteq
T_{X/\mathcal{F}}$
be a subbundle a {\it $G$-connection}
$(E,\nabla)$ is
$$
\begin{aligned}
D: E^{\nabla} &\longrightarrow G^\vee\otimes
E^{\nabla} \\
\text{s.t.} D(fe)&=fDe+d_Gf\otimes e
\end{aligned}
$$
where
$d_G:\mathcal{O}_{X/\mathcal{F}}\to\Omega_{X/
\mathcal{F}}^1\to
G^\vee$.

It's a fact that a $G$-connection
$(E,\nabla)$ exists if and only if
$\text{at}_{X/\mathcal{F}}(E,\nabla)$ maps to
0 in
$H^1(X,G^\vee
\otimes_{G_{X/\mathcal{F}}}\text{End}
_{\mathcal{O}_{X/\mathcal{F}}}
E^{\nabla}$.

The analogue of Beauville's lemma
\ref{seminars-algebraic-geometry-lemma-beauville} is

\begin{lemma}[Beauville]
\label{seminars-algebraic-geometry-lemma-analogue-beauville}
Suppose $N_{\mathcal{F}}=\mathcal{A} \oplus
B$ such that
$\nabla^B\mathcal{A} \subseteq
\mathcal{F}^\vee
\otimes_{\mathcal{O}_X}\mathcal{A}$
and $\nabla^B \mathcal{B} \subseteq
\mathcal{F}^\vee\otimes_{\mathcal{O}_X}
\mathcal{B}$.
Let $\nabla_{\mathcal{A}}$ and
$\nabla_{\mathcal{B}}$ be a flat
 $\mathcal{F}$-connection on $\mathcal{A}$
 and $\mathcal{B}$.

Then
$\text{at}_{X/\mathcal{F}}(\mathcal{A},
\nabla_{\mathcal{A}})
\in
H^1(X,\Omega_{X/\mathcal{F}}
^1\otimes_{\mathcal{O}_{X/\mathcal{F}}}
\text{End}_{\mathcal{O}_{X/\mathcal{F}}}
\mathcal{A}^{\nabla_{\mathcal{A}}}$
comes from
$H^1(X,(\mathcal{A}^{\nabla_{\mathcal{A}}})
^\vee
\otimes_{X/\mathcal{F}}\text{End}
_{\mathcal{O}_{X/\mathcal{F}}}\mathcal{A}
^{\nabla_{\mathcal{A}}}$.
\end{lemma}

\begin{proof}
Same proof but we need a Lie bracket on
$T_{X/\mathcal{F}}=N_{\mathcal{F}}\subseteq
N_{\mathcal{F}}\overset{p}{\leftarrow}T_X$.
Let $v_1,v_2$ be local sections on
$T_{X/\mathcal{F}}$
with $v_1=p(t_1)$ and $v_2=p(t_2)$.
Define
$$
[v_1,v_2]=p([t_1,t_2]).
$$
One can check that this a Lie bracket on
$T_{X/\mathcal{F}}$.
Suppose that $t_1 \in \mathcal{F}$. Then
$$
p([t_1,t_2])\overset{\text{def}}{=}
\nabla_{t_1}^{\mathcal{B}}v_2=0.
$$
For $v_3 \in \mathcal{F}$,
$$
\nabla_{v_3}^{\mathcal{B}}[v_1,v_2]
=p([v_3,[t_1,t_2]])=\pm\nabla_{v_3}
^{\mathcal{B}}v_1\pm
\nabla_{v_3}^{\mathcal{B}}v_2=0.
$$
\end{proof}

\begin{lemma}
\label{seminars-algebraic-geometry-lemma-vanishing-chern-class}
Suppose that $\mathcal{F}$ is a regular
foliation on $X$.
Let $N_{\mathcal{F}}=\mathcal{L}_1 \oplus
\ldots \oplus \mathcal{L}_q$
be line bundles on $X$,
$\nabla^{\mathcal{B}}\mathcal{L}_i\subseteq
\mathcal{F}^\vee\otimes\mathcal{L}
_i\rightsquigarrow
\nabla_i$.
Suppose
$(\mathcal{L}_i,\nabla_i)\cong(\mathcal{L}_j,
\nabla_j)$
for all $i,j$.
Then $c_1(\mathcal{L}_i)=0$.
\end{lemma}

\begin{proof}
By the Comparison Lemma,
$$
\begin{aligned}
H^1(X,\Omega_{X/\mathcal{F}}^1
&\longrightarrow H^1(X,\Omega_X^1) \\
\text{at}_{X/\mathcal{F}}(\mathcal{L},
\nabla_i)
&\longmapsto
\text{at}(\mathcal{L}_i),
\end{aligned}
$$

\noindent
it suffices to prove that
$\text{at}_{X/\mathcal{F}}(\mathcal{L}_i,
\nabla_i)$
vanishes.
$$
\text{at}_{X/\mathcal{F}}(\mathcal{L}_i,
\nabla_i)
\in H^1(X,\Omega_{X/\mathcal{F}}^1)
=\bigoplus_{j=1}^9
H^1(X,(\mathcal{L}_j^{\nabla_j})^\vee)
$$
comes from
$H^1(X,(\mathcal{L}_i^{\nabla_i})^\vee)$.
\end{proof}

\begin{example}
\label{seminars-algebraic-geometry-example-degree-2-line-bundle-on-e}
$A=E^1 \times B^{n-1}$, an abelian variety.
Let $ L$ be the pullback of a degree 2
divisor on $E$.
Then there exists  $T_X
\to\!\!\!\!\!\to L^{\oplus n-1}$.
\end{example}

\begin{proposition}
\label{seminars-algebraic-geometry-proposition}
$\mathcal{F}$ regular foliation of dimension
1 on $X$ complex projective
of dimension $\geq 3$.
Suppose $\mathcal{F}$ is transversally
projective flat
and that $X$ contains a rational curve.
Then $\mathcal{F}$ is a $\mathbb{P}^1$-bundle
and it is tangent
to $\mathcal{F}$.
\end{proposition}

\medskip\noindent
Now we explain what it means to be
projectively flat.
It is when the projectivization of the normal
bundle $\mathbb{P}(N)$
is induced by a representation of $\pi_1(X)
\to \text{PGL}_g$
it comes with an Ehresmann connection on
$$
\xymatrix{
\mathbb{P}(N)=\mathbb{P}\ar[d]^{f}\\
X
}
$$
$T_{\mathbb{P}}=T_f \oplus g$.

Now, it is also called transversally
projectively flat
if $\nabla^B$ is a partial connection on $N$
and there exists a foliation $\mathcal{H}$ on
$$
\xymatrix{
\mathbb{P}(N)\ar[d]\\ X
}
$$
such that $\mathcal{H}$ projects onto
$\mathcal{F}$.
Transversality means that $\mathcal{H}
\subseteq \mathcal{G}$.

Finally, transversally projectively flat
means that
it is induced on a
$\text{Mon}(\mathcal{F})$-module
by a representation of
$\pi_1(\text{Mon}(\mathcal{F}) \to
\text{PGL}_g$.

\section{Um passeio nos trabalhos de
Gerd Faltings}
\label{seminars-algebraic-geometry-section-faltings}

\noindent
Hossein Movasati, IMPA.
Pelestra especial,
May 19, 2026.

\medskip\noindent
{\bf Abstract.}
Nesta palestra vamos dar um passeio nos
trabalhos de G. Faltings, começando por sua
demonstração da conjetura de Mordell em
1983, que lhe trouxe a medalha Fields. Uma
boa parte de sua carreira foi elaborar o
lado aritmético de um dicionário heurístico
proposto por P. Deligne em 1970, que tem
origem na teoria conjetural de motivos de
Grothendieck. Isso inclui versão p-ádica de
teoria de Hodge, domínios de períodos,
correspondência de Simpson, etc. Várias
propriedades aritméticas de variedades
abelianas são outro foco da carreira dele
que vamos abordar nesta palestra. Em 2026
ele ganhou o prêmio Abel por suas
contribuições fundamentais em geometria
aritmética.

\medskip\noindent
{\bf Upshot.} 
Elliptic curves are diophantine equations.
Some are simple and fun to solve,
some are extremely complicated ---
like Fermat's last theorem.
Finding solutions to a diophantine
equation is like finding $k$-points
on an algebraic curve for suitable
choices of fields.
In this sense,
Mordell's conjecture, now Faltings'
theorem is a statement of finiteness
of solutions of diophantine equations.
Other results are adaptations
of results from complex geometry
to finite fields.

\medskip\noindent
{\bf Notes.}

Finding roots of a polynomial $f \in
\mathbb{Q}[x,y]$ becomes the geometric
problem of finding the $\mathbb{Q}$-rational
points of $\operatorname{Spec}(\mathbb{Q}[x,y]/(f))$.
Or, one can pass to the projective version
$\operatorname{Proj}(\mathbb{Q}[x,y,z]/(F))$,
where $F$ is the homogenization of $f$.
We write
$C(\mathbb{Q})=\{[x:y:z] \in
\mathbb{P}^2(\mathbb{Q}) \mid
F(x,y,z)=0\}$
and similarly $C(\mathbb{C})$ for the
analogous set.
 
It turns out that the topology
of the curve, via its genus, strongly
constrains the $\mathbb{Q}$-rational points.
The easiest cases are:
\begin{itemize}
\item ($g=0$.) Here $d=1,2$;
then the curve is projectively equivalent to
a conic, and over $\mathbb{Q}$ it has either
no rational points or infinitely many.

\item ($g=1$.) Here $d=3$; this is more difficult
but still well-studied; see the
Birch--Swinnerton-Dyer conjecture. 

\item ($g\geq 2$.) Here $d\geq 4$.
This is much more complicated!
Faltings'
theorem implies that $C(\mathbb{Q})$ is
finite.
\end{itemize}

\noindent
This is the former Mordell conjecture
from 1962.

\begin{theorem}[Faltings, 1983]
For any smooth curve of genus $\geq 2$ over
$\mathbb{Q}$, we have $\#C(\mathbb{Q})<
\infty$.
\end{theorem}

The same theorem remains true after
replacing $\mathbb{Q}$ by a number field,
that is, a finite field extension of
$\mathbb{Q}$.

\begin{theorem}[Manin 1963, Grauer 1965]
\label{seminars-algebraic-geometry-theorem-manin-grauer}
For any smooth non isotrivial curve of
genus $\geq 2$ over $k=\mathbb{C}(t)$
we have
$$
\# C(k) < \infty.
$$
\end{theorem}

\begin{proof}
Using Gauss-Manin connection ---
the name of which was invented by 
Grothendieck.
\end{proof}

\noindent
The main difference between
$\mathbb{C}(t)$
and $\mathbb{Q}$
is that in $\mathbb{C}(t)$
we have a derivation
$$
\frac{\partial }{\partial t}:
\mathbb{C}(t) \to \mathbb{C}(t)
$$
satisfying the Leibniz identity.

\medskip\noindent
The proof is somehow related
to Shafarevich and Tate conjectures.
The former states there's only
finitely many isomorphism
classes of abelian varieties of a
given dimension with
fixed polarization degree, defined
over a field with
good reduction outside a finite set
of places.

Tate conjecture was proved for finite
fields by Tate in 1966,
and by Faltings for number fields
in 1983.
It relates maps between the tangent
spaces of two abelian varieties
and maps between the varieties
themselves:
$$
\Hom_K(A,B) \otimes \mathbb{Z}_p
\cong
\Hom_{Gal(\tilde{K}/K)}
(T_pA,T_pB)
$$
where $K$ is a field of characteristic
$\neq p$, and $A,B$ are abelian over $K$.

Further, the proof uses Arakelov's
theory of arithmetic genus
and Faltings height.

\medskip\noindent
{\bf Abelian varieties.} 

Any curve (over any field)
can be embedded in some variety
called the Jacobian variety.
Over $\mathbb{C}$,
abelian varieties are just tori:
$$
\mathbb{C}^n/\Lambda
\cong_{C^\infty}
S^1\times… S^1,\qquad  2n \text{ times}.
$$

An abelian variety of dimension one
over a field $k$ is called
an {\it elliptic curve}.
They are given by diophantine
equations of degree 3 in two variables:
$$
y^2+a_1xy+a_3y
=
x^3+a_2x^2+a_4x,
\qquad a_1,…,a_6 \in k.
$$

The name was invented by A. Weil.
Before him, many mathematicians like E.
Picard,
G. Humbert, were using ``hyperelliptic
surface'' for Abelian surfaces.

\medskip\noindent
{\bf Shavarevich's conjecture.} 
What is the number of non-isomorphic elliptic
curves over $\mathbb{Q}$ and with bad
reduction only for primes $2,3$?

\medskip\noindent
A consequence of Tate conjecture
and Chebotarev density theorem is:;

\begin{theorem}
\label{seminars-algebraic-geometry-theorem-tate-chebot-consequence}
LEt $E_i$, $i=1,2$ be two
elliptic curves over $\mathbb{Q}$.
TFAE:
\begin{enumerate}
\item $E_1$ is isogeneous to $E_2$
over $\mathbb{Q}$.
\item For all except a finite number
of primes,
$$
\# E_1(\mathbb{F}_p)
=\# E_2(\mathbb{F}_p).
$$
\end{enumerate}
\end{theorem}

\begin{definition}
\label{seminars-algebraic-geometry-definition-isogeny}
{\it AI-generated.}
An isogeny between elliptic curves over
$\mathbb{Q}$ is a nonzero morphism of
elliptic curves defined over $\mathbb{Q}$
that is also a group homomorphism and has
finite kernel.
\end{definition}

Here isogenous means that there exists an
isogeny $f:E_1 \to E_2$ over $\mathbb{Q}$.

\medskip\noindent
One of the key ingredients
in Falting's proof
is the notion of {\it Faltings height},
which are integrals used to
control the numerator and denominator
or rational solutions of a (diophantine?)
equation.

For example, for $r = n/m$,
$h(r)=n^2+m^2$. For all $M \in \mathbb{R}$,
$$
\{r \in \mathbb{Q}| h(r)<M\}
$$
is finite.

\begin{proposition}[Main property
of height functions]
\label{seminars-algebraic-geometry-proposition-faltings-height}
For any positive integer $M$,
the number of elliptic curves $E/\mathbb{Q}$
with $h_{\text{Falt}}<M$ is finite.
\end{proposition}

\noindent
On the curve given by Equation
\ref{seminars-algebraic-geometry-equation}
consider the 1-form
$$
\omega=\frac{dx}{2y+a_1x+a_3}.
$$
The Falting's height is basically…
[missing! an integral involving $\omega$…]

\medskip\noindent
\begin{itemize}
\item Elliptic integrals $\leq 1900$.
\item Poincaré, Pircard, $\sim 1900$.
\item Singular, de Rham 1900-1950.
\item Weil: cohomology over the integers
(enter arithmetic).
\item Grothendieck: étale, crystalline, etc.
\end{itemize}

\noindent
Then is Deligne's dictionary,
in a 1970 paper, where $\ell$-adic
cohomology is introduced
(I think --- abstract in french, check).

\medskip\noindent
From complex to arithmetic geometry:

\begin{itemize}
\item Endilchkeitssätze
fur abelsche Varietäten über Zahlkörpern.
\item 1984: Calculus on arithmetic surfaces.
\item 1988: p-adic Hodge theory.
\item 1990: Degeneration of abelian
varieties (with Chai Ching-li).
\item 1992: Lectures on the
arithmetic Riemann-Roch theorem.
\item 1999: Does there exist an arithmetic
Kodaira-Spencer class? (!!!)
\item 2005: A p-adic Simpson corresponding.
\item 2020: covering of p-adic period
domains.
\end{itemize}

\medskip\noindent
{\bf p-adic Hodge theory.}

Grothendieck and Deligne
noticed that we can obtain
the Hodge decomposition of cohomology
using filtrations, which sit
in the realm of polynomials ---
no harmonic analysis needed.

The idea of Deligne (no, Faltings?) 
was to use
this proof for the $p$-adic case.

\section{Grothendieck--Teichm\"uller theory: from arithmetic geometry to arithmetic topology}
\label{seminars-algebraic-geometry-section-grothendieck-teichmuller-arithmetic-topology}

\noindent
Marco Boggi.
Qua 27 mai 2026, 15:30.
SALA 228.

\medskip\noindent
{\bf Abstract.}
Grothendieck--Teichm\"uller theory originated
in arithmetic geometry through the study of
Galois actions on geometric \'etale
fundamental groups. In this talk, we argue
that arithmetic topology provides a natural
setting for applications of
Grothendieck--Teichm\"uller theory. More
specifically, we explain how our theory of
automorphisms of profinite mapping class
groups may lead to new results on profinite
rigidity questions for mapping class groups.

\medskip\noindent
{\bf Warning.}
These are live notes from the talk,
written with unusually heavy AI assistance
since that day I forgot my glasses.
Some theorem
statements may be imprecise; the goal is to
record the ideas and references from the day.\medskip\noindent

\medskip\noindent
The main upshot is that GT theory allows us
to pass the problem to a free group in three
generators, and that the talk aims to
contrast arithmetic and topological
properties.

A more personally significant 
upshot, explained next,
is that the étale fundamental
group of the moduli space of curves
is the profinite completion of
the fundamental group of any curve in 
the moduli, which makes sense since the
underlying topological space of all curves
in the moduli is the same.

\medskip\noindent
Historically, this is very much in the spirit
of Grothendieck's 1984 \emph{Esquisse d'un
programme}. Drinfeld then introduced the
Grothendieck--Teichm\"uller group in 1991,
making precise a framework that was already
expected to have arithmetic applications.

\medskip\noindent
The point is that $\mathcal M_{g,n}$ is
not just an abstract moduli space: it is
built from one fixed topological surface
$S_{g,n}$. A point of $\mathcal M_{g,n}$
is essentially a complex structure on this
same underlying surface.

\medskip\noindent
Thus the topological fundamental group of
the moduli space records how a family of
curves can move around and come back to the
same curve with its underlying surface
changed by a diffeomorphism. This gives the
mapping class group:
$$
\pi_1^{\mathrm{top}}
\bigl(\mathcal M_{g,n}(\mathbb C)\bigr)
\cong
\Gamma_{g,n}.
$$
Here
$$
\Gamma_{g,n}
=
\operatorname{Diff}^+(S_{g,n})/
\operatorname{Diff}_0(S_{g,n}).
$$

\medskip\noindent
The \'etale fundamental group is the
algebraic-geometric version of the
fundamental group. Instead of remembering
all covering spaces, it remembers finite
algebraic covers. Over $\mathbb C$, this
amounts to taking the profinite completion
of the usual fundamental group. Therefore
one gets
$$
\pi_1^{\acute{e}t}
\bigl(\mathcal M_{g,n}
\otimes_{\mathbb Q}\overline{\mathbb Q}\bigr)
\cong
\widehat{\Gamma}_{g,n}.
$$

\medskip\noindent
So the profinite mapping class group is not
an artificial object: it is literally the
\'etale fundamental group of the moduli
space of curves. This is why
Grothendieck--Teichm\"uller theory naturally
acts on profinite mapping class groups.

\medskip\noindent
The starting point is the \'etale
fundamental group of the projective line with
three points removed. Over $\mathbb{C}$, the
space
$$
\mathbb{P}^1_{\mathbb{C}} \setminus \{0,1,\infty\}
$$
is a sphere with three punctures, so its
topological fundamental group is a free group
on two generators:
$$
\pi_1^{\mathrm{top}}
\bigl(\mathbb{P}^1_{\mathbb{C}} \setminus
\{0,1,\infty\}\bigr)
\cong F_2.
$$
After profinite completion one gets
$$
\pi_1^{\acute{e}t}
\bigl(\mathbb{P}^1_{\overline{\mathbb{Q}}} \setminus
\{0,1,\infty\}\bigr)
\cong \widehat{F}_2,
$$
the free profinite group on two generators.

\medskip\noindent
Now let
$$
X=\mathbb{P}^1_{\mathbb{Q}} \setminus
\{0,1,\infty\}.
$$
There is an exact sequence of profinite
groups
$$
1\to
\pi_1^{\acute{e}t}(X_{\overline{\mathbb{Q}}})
\to
\pi_1^{\acute{e}t}(X)
\to
\operatorname{Gal}(\overline{\mathbb{Q}}/\mathbb{Q})
\to 1.
$$
Because the geometric fundamental group is
\(\widehat{F}_2\), the absolute Galois group
acts by outer automorphisms on a free
profinite group:
$$
\operatorname{Gal}(\overline{\mathbb{Q}}/\mathbb{Q})
\longrightarrow
\operatorname{Out}(\widehat{F}_2).
$$
This is one of the first places where
Grothendieck's vision becomes concrete.
The theorem says that this action is
faithful.

\medskip\noindent
In particular, the absolute Galois group can
be studied through its action on the
combinatorics of finite coverings of
\(\mathbb{P}^1\setminus\{0,1,\infty\}\).
This is the origin of dessins d'enfants.
A finite cover of
\(\mathbb{P}^1\setminus\{0,1,\infty\}\)
corresponds to a finite quotient of
\(\widehat{F}_2\), and the Galois action
permutes these quotients in a very rigid way.
That rigidity is what eventually leads to the
Grothendieck--Teichm\"uller group.

\medskip\noindent
The point of the talk is that a similar
picture appears in arithmetic topology.
Mapping class groups also act on profinite
fundamental groups of surfaces, and one can
study their automorphisms in a way that
mirrors the Galois action above. This
suggests profinite rigidity questions: to
what extent does the profinite completion of
a mapping class group remember the original
group? The Grothendieck--Teichm\"uller
framework provides a natural language for
these problems, and may lead to new rigidity
results.

\medskip\noindent
In the arithmetic case there are more
automorphisms than in the topological case,
and this is what makes the theory useful for
applications. The speaker suggested that
these methods should lead to a proof of a
failure of profinite rigidity: more
precisely, if \(d(S)\) denotes the complexity
of the surface, then there should exist
finite-index subgroups \(\Gamma\) and
\(\Gamma'\) of \(\Gamma(S)\) such that
$$
\widehat{\Gamma}\cong \widehat{\Gamma'}
\qquad\text{but}\qquad
\Gamma\not\cong \Gamma'.
$$

\medskip\noindent
For \(g=0\), the speaker can prove that every
pointwise inner automorphism of
\(\widehat{\Gamma}(S)\) is actually inner.
That is, if an automorphism acts on every
single element as conjugation, then the
automorphism itself is conjugation.

\medskip\noindent
As a corollary, the Grothendieck--Teichm\"uller
group acts faithfully on the conjugacy
classes of elements of the profinite
completion of the braid group
\(\widehat{B}_n\). This is interesting from
the topological point of view because of the
Alexander theorem, which characterizes
braids via their relationship with links.

\medskip\noindent
Another application mentioned in the talk is
Furusho's work on finite-type (Vassiliev)
knot theory. The GT action can be used to
organize knot-theoretic information in a way
that is compatible with the profinite braid
group viewpoint.

\medskip\noindent
Let $B$ denote the profinite mapping class
group under consideration. We write
\(\operatorname{Aut}^{\mathbb{I}_0}(B)\) for
the subgroup of \(\operatorname{Aut}(B)\)
consisting of those automorphisms that
preserve the procyclic subgroups generated by
the elements \(C_r\) arising from Dehn
twists. The centralizer of the natural image
of the Grothendieck--Teichm\"uller action in
\(\operatorname{Aut}^{\mathbb{I}_0}(B)\) is
what we call the Grothendieck--Teichm\"uller
group \(G(c,0)\).

\medskip\noindent
For a connected oriented surface
\(S_{g,n}\) of genus \(g\) with \(n\)
punctures and \(\chi(S_{g,n})<0\), the
mapping class group is
the group of orientation-preserving
diffeomorphisms modulo isotopy to the
identity. More precisely, if
\(\Gamma_{g,n}\) denotes the mapping class
group of \(S_{g,n}\), then
\(\Gamma_{g,n}=\mathrm{P}\Gamma(S_{g,n})\)
in the pure case.

\begin{theorem}\label{seminars-algebraic-geometry-thm:mapping-class-moduli-space}
There is a natural isomorphism
$$
\Gamma_{g,n}\cong
\pi_1^{\mathrm{top}}
\bigl(\mathcal{M}_{g,n}(\mathbb{C})\bigr),
$$
where \(\mathcal{M}_{g,n}\) is the moduli
space of smooth genus \(g\) curves with
\(n\) marked points.
\end{theorem}

\begin{proof}
This is the classical identification of the
mapping class group with the orbifold
fundamental group of the moduli space. In the
case of ordered marked points, the pure
mapping class group corresponds to the
fundamental group of the pointed moduli
space.
\end{proof}

\medskip\noindent
There is also a standard exact sequence
relating the pure mapping class group to the
full mapping class group:
$$
1\to \widehat{\mathrm{P}\Gamma}_{g,n}
\to \widehat{\Gamma}_{g,n}
\to \Sigma_n
\to 1.
$$
At the discrete level this comes from the
action of the mapping class group on the set
of punctures, and the quotient records the
induced permutation of the punctures.

\begin{remark}
How can we study the Deligne--Mumford
boundary through points on a surface? This
appears to be a standard method discussed in
the talk.
\end{remark}

\medskip\noindent
The theorem implies that the \'etale
fundamental group of the moduli space is the
profinite completion of the mapping class
group:
$$
\pi_1^{\acute{e}t}
\bigl(\mathcal{M}_{g,n}\otimes_{\mathbb{Q}}\overline{\mathbb{Q}}\bigr)
\cong
\widehat{\Gamma}_{g,n}.
$$
In the pure case this is
\(\widehat{\mathrm{P}\Gamma}(S_{g,n})\).

\begin{theorem}[Hoshi--Mochizuki]\label{seminars-algebraic-geometry-thm:hoshi-mochizuki-centralizer}
Let \(\Sigma_n\) act on the profinite
mapping class group in the natural way by
permuting the punctures. Then there is a
natural injection
$$
Z_{\operatorname{Out}^{\mathbb{I}_0}
(\widehat{\Gamma}_{g,n})}(\Sigma_n)
\hookrightarrow
Z_{\operatorname{Out}^{\mathbb{I}_0}
(\widehat{\Gamma}_{g,n+1})}(\Sigma_{n+1}),
$$
and this map is an isomorphism for \(n\ge 5\).
\end{theorem}

\begin{proof}
This is the centralizer comparison used in
the study of profinite mapping class groups
and their puncture-permutation symmetries.
The low-puncture cases require separate
analysis, while for \(n\ge 5\) the natural
stabilization map becomes an isomorphism.
\end{proof}

\begin{theorem}[Hoshi--Minamide--Mochizuki]\label{seminars-algebraic-geometry-thm:hoshi-minamide-mochizuki}
There is a natural isomorphism
$$
\operatorname{Out}^{\mathbb{I}_0}
\bigl(\mathrm{P}\check\Gamma(S)\bigr)
\cong
\Sigma_n \times \widehat{\operatorname{GT}}.
$$
Equivalently, in the genus zero punctured
sphere case one may write
$$
\operatorname{Out}^{\mathbb{I}_0}
\bigl(\mathrm{P}\check\Gamma_{0,n}\bigr)
\cong
\Sigma_n \times \widehat{\operatorname{GT}}.
$$
\end{theorem}

\begin{proof}
This is the theorem that removes the extra
rigidity hypotheses and identifies the full
outer automorphism group, under the
\(\mathbb{I}_0\)-condition, as the product of
the puncture permutation group and the
profinite Grothendieck--Teichm\"uller group.
\end{proof}

\medskip\noindent
There is also a slogan in the genus-zero
setting:
$$
\operatorname{Out}\bigl(\widehat{\Gamma}_{0,n}\bigr)
\cong \widehat{\operatorname{GT}}
$$
for the appropriate completion and rigidity
condition encoded by the exponent on
\(\operatorname{Out}\). More generally, for
small genus \(g\leq 2\), the speaker proved a
corresponding isomorphism
$$
\operatorname{Out}^{\ast}\bigl(\widehat{\Gamma}_{g,n}\bigr)
\cong
\Sigma_n \times \widehat{\operatorname{GT}},
$$
where the superscript \(\ast\) indicates the
same rigidity condition as in the previous
theorems.

\section{Hodge structures in toric ambient}
\label{seminars-algebraic-geometry-section-hodge-structures-toric-ambient}

\noindent
Roberto Villaflor, USM, Chile.
GADEPS Focused Meeting June 29, 2026, IMPA.

\medskip\noindent
Joint with L. Broring, G. Martinez,
A. Viergever.

\medskip\noindent
{\bf Hodge structures over a perfect
field.} Let $X$ be a smooth
projective variety over $\mathbb{C}$.
We have the Hodge-De Rham spectral sequence
$E^{p,q}_1=H^q(X,\Omega_X^p) 
\implies \mathbb{H}^k(X,\Omega_X^\bullet)
\simeq H^k_{dR}(X/\mathbb{C})$.
It degenerates at $E_1$, giving
a non-canonical isomorphism
$$
H^\bullet_{dR}(X/\mathbb{C})
\simeq
\bigoplus_{p+q=k}H^q(X,\Omega_X^p)
$$
We may ask ourselves whether this is also
true for other fields.

\begin{remark}[Deligne-Illusie]
\label{seminars-algebraic-geometry-remark-derham-to-hodge-not-degenerate}
In general, for $X/k$ perfect field,
this may not degenerate at $E_1$.
\end{remark}

\noindent
For $X = \{F=0\} \subseteq \mathbb{P}^{n+1}$
smooth degree $d$ hypersurface over
$\mathbb{C}$,
\begin{enumerate}
\item 
$$
H^q(X,\Omega_X^p)
\simeq \begin{cases}
H^q(X,\Omega_X^p)_{\operatorname{prim}}\qquad
&p \neq q \\
H^p(X,\Omega_X^p)_{\operatorname{prim}}\oplus
\mathbb{C}\theta |_X^p\qquad &
\end{cases}
$$
where $\theta$ is the class of a hyperplane.

\item
$H^q(X,\Omega_X^p)_{\operatorname{prim}}=0$
for $p \neq q \neq  n$ 
(Lefschetz hyperplane theorem).

\item Griffith basis theorem for $p+q=n$.
\end{enumerate}

\begin{remark}
\label{seminars-algebraic-geometry-remark-lps}
For $X/k$ smooth projective, $k$ perfect
of characteristic $\neq 2$,
Levine-Pepin-Srinivas '24: (1), (2)
and (3) still hold.
\end{remark}

\medskip\noindent
{\bf Quadratic Euler characterstic.} 
Let $X$ be a CW-complex of dimension $n$.
$$
\chi_{\operatorname{top}}(X)
=\sum_{i=1}^n(-1)^i b_i(X),
$$
where
$b_i=\operatorname{rank}H_i(X,\mathbb{Z})$
is the $i$-th Betti number.

\begin{proposition}
\label{seminars-algebraic-geometry-proposition-properties-of-euler}
\begin{enumerate}
\item $\chi_{\operatorname{top}}$
is homoropy invariant.
\item $Y \subseteq X$,
$\chi_{\operatorname{top}}(X\setminus Y)
=\chi_{\operatorname{top}}(X)
-\chi_{\operatorname{top}}(Y)+1$.
\item $F \hookrightarrow E \to B$ fiber
bundle, then $\chi_{\operatorname{top}}(E) =
\chi_{\operatorname{top}}(B)
\chi_{\operatorname{top}}(F)$

\item $\hat{\chi_{\operatorname{top}}}(X
\wedge Y) =
\hat{\chi_{\operatorname{top}}}(X)
\hat{\chi_{\operatorname{top}}}(Y)$.
\end{enumerate}
\end{proposition}

\begin{definition}
\label{seminars-algebraic-geometry-definition-stably-homotopic}
$X$ is {\it stably homotopic} to $Y$
if
$$
\Sigma^kX \equiv \Sigma^kY
$$
where $\Sigma^k$ is the smash product
of $k$ 1-spheres.
\end{definition}

\noindent
By the reduced Euler characteristic formula,
if $X$ and $Y$ are stably homotopic,
then $\chi_{\operatorname{top}}(X)
=\chi_{\operatorname{top}}(Y)$.

\medskip\noindent
Morel-Boevodky '99: we have
a stable motivic homotopy category.

There exists an analogous construction
of categorical Euler characteristic
$X$ quasi-projective smooth $k$-scheme.
$\chi(X/k) \in GW(k)$
where $GW(k)$ is the Grothendieck-Witten
ring. In characteristic $\neq 2$, 
bilinear symmetric forms = group
generated by the monoid of isometry
classes of quadratic forms with $\oplus$.

The upshot is that instead
of assigning a number, now we put
a quadratic form to our scheme.

\begin{remark}
\label{seminars-algebraic-geometry-remark-real-field-euler}
If $k \subset \mathbb{R}$ is a real field,
the rank of $\chi(X/k) =
\chi_{\operatorname{top}}(X(\mathbb{C}))$,
and the signature of $\chi(X/k)$
is
$\chi_{\operatorname{top}}(X(\mathbb{R}))$.
\end{remark}

\begin{theorem}[Levine '20]
\label{seminars-algebraic-geometry-theorem-levine20}
$X/k$ smooth projective, $k$ perfect,
char $\neq 2$, then
\begin{enumerate}
\item If $\dim X$ is odd, then
$\chi(X/k) = mH$, $m \in \mathbb{Z}$,
$$
H=\begin{pmatrix}1&0\\ 0&-1\end{pmatrix}
=\langle 1\rangle+ \langle-1\rangle.
$$
\item If $\dim X$ is even, then
$\chi(X/k)=mH+Q$, $m \in \mathbb{Z}$
and $Q$ is determined by the first
page of the Hodge-to-de-Rham spectral
sequence:
$$
Q: H^{n/2}(X,\Omega^{n/2}_X)\times
H^{n/2}(X,\Omega_X^{n/2})
\xrightarrow{\cup }
H^n(X,\Omega_X^n)
\xrightarrow{\operatorname{trace}}k.
$$
\end{enumerate}
\end{theorem}

\noindent
By LPS '24, for $X = \{F=0\}
\subset \mathbb{P}^{n+1}_k$,
$Q = Q_{\operatorname{prim}} \oplus
\langle 1\rangle$ and
$$
Q_{\operatorname{prim}}:
J(X) _{(d-2)(n/2+1)}\times
J(X)_{(d-2)(n/2+1)}\xrightarrow{\cdot}
J(X)_{(d+1)(n+2)}
=\det (\text{Hess}((F)))
$$

\medskip\noindent
{\bf Non-split toric varieties.}
Recall: $X/\mathbb{C}$ is toric
if it admits an algebraic torus action
$T = (\mathbb{C}^\times)^n$, $n = \dim X$
such that there exists $x \in X$
with $T \simeq T\cdot x$ (dense open orbit).
When $X$ is normal there exists a
combinatorial object called the fan
$\Sigma = \bigcup_{k=1}^n \Sigma(k)$
consisting of cones inside $N_{\mathbb{R}}
\simeq M^4_{\mathbb{R}} \simeq \mathbb{R}^n$
where $M = \Hom(T,\mathbb{C}^\times)
\simeq \mathbb{Z}^n$
group of characters of $T$
for $\sigma \in \Sigma(k)$
if is a cone $\sigma \subseteq N_{\mathbb{R}}$
of dimension $k$, rational polyhedral
and strongly convex.

\begin{itemize}
\item for all $\sigma \in \Sigma$,
each face of $\sigma$ is in $\Sigma$,
\item for all $\sigma,\sigma' \in \Sigma$
$\sigma \cap \sigma' \in \Sigma$.
\end{itemize}

\noindent
For each $\sigma \in \Sigma$, we have
$\sigma^\vee \subseteq M_{\mathbb{R}}$
$\implies $ $\sigma^\vee \cap M \subset M$
is a semigroup, which gives
$\Spec \mathbb{C}[\sigma^\vee\cap M]
\supseteq \Spec \mathbb{C}[M]
\simeq T$.
Denote $U_\sigma = \Spec
\mathbb{C}[\sigma^\vee \cap M]$.
$\sigma \subseteq \sigma' \implies
(\sigma')^\vee \subseteq \sigma^\vee \implies 
U_\sigma \subseteq U_{\sigma'}$.

The idea is that the fan gives us the data
of how to glue the affine toric varieties
$\Spec \mathbb{C}[\sigma^\vee \cap M]$
into a toric variety.

When the rays of all cones ar $\mathbb{R}$-
linearly independent, we say
$X_{\Sigma}$ is simplicial.
\begin{itemize}
\item $X_\Sigma$ is an orbifold
(finite quotient singularity).
\item We can put homogeneous coordinates
on $X_{\Sigma}$, the {\it Cox ring}:
$$
\operatorname{Cox}(X_{\Sigma})
= \bigoplus_{\beta \in \text{Cl}(X_{\Sigma}}
H^0(\mathcal{O}_{X_{\Sigma}}(\beta)
$$
since $X_{\Sigma}\setminus T
= \bigcup_{D \in
\text{Div}_{T}(X_{\Sigma}})$.
\end{itemize}

\noindent
Each $\text{Div}_T(X_{\Sigma})
=\{D_{\rho_i}:i 1,…,r\}$,
$\Sigma(1)= \{\rho_1,…,\rho_r\}$.
$\operatorname{Cox}(X_{\Sigma})
=\mathbb{C}[X_{\rho_1},…,X_{\rho_r}]$.
This ring is naturally graded by
$\text{Cl}(X_{\Sigma})$,
$\text{deg}(X_{\rho_i})=
[D_{\rho_i}] \in \text{Cl}(X_{\Sigma})
=\bigoplus_{i=1}^r \mathbb{Z}\cdot
D_{\rho_i}/\text{div}(M)$.

A non-split torus $T$ is a $k$-linear
algebraic group such that
$T_{\overline{k}}\simeq
(\overline{k}^\times)$ is a split torus.

\begin{itemize}
\item $S^1$ is a non-split
$\mathbb{R}$-torus.
Given a $k$-torus $T$, a $T$-toric
variety is a $k$-scheme $T$ acting on $X$
such that $T_{\overline{k}}$ acts on
$X_{\overline{k}}$ becomes a usual
toric variety.
\end{itemize}

\noindent
The Galois group $\Gamma=\operatorname{Gal}
(\overline{k}/k)$ acts on $T$ and $X$.

\begin{theorem}
\label{seminars-algebraic-geometry-theorem-toric-primitive-hodge}
$Y \subseteq X$, $T$-toric,
$$
H^q(Y,\Omega_Y^p)_{\operatorname{prim}}
\simeq
(Cox(X_{\overline{k}})
/\langle \frac{\partial F}{\partial x_1}
,…,\frac{\partial F}{\partial x_r}\rangle)
^\Gamma_{[\gamma]}
$$
\end{theorem}

\section{Determinantal surfaces
and Noether-Lefschetz theory}
\label{seminars-algebraic-geometry-section-determinantal-surfaces-nl}

\noindent
Manuel Leal, IMPA.
GADEPS Focused Meeting June 30, 2026, IMPA.

\medskip\noindent
The talk studies determinantal surfaces
inside the parameter space of smooth
degree $d$ surfaces in $\mathbb{P}^3$.
Such a surface contains natural ACM curves,
which give extra line bundles besides
$\mathcal{O}_X(1)$.
Thus determinantal surfaces lie in the
Noether-Lefschetz locus. The main point is
that, for the degree type $(a,b)$, the
determinantal locus is itself a
Noether-Lefschetz component.

\medskip\noindent
Given a determinantal surface given
by the vanishing of the determinant
of a matrix $S$, we can define two curves:
$\mathcal{C}_+$ and $\mathcal{C}_-$
both given by the vanishing of maximal
rank minors of the matrix of $S$ after
removing everything except the highest/lowest
degree forms of the matrix.

These two curves are arithmetically
Cohen-McCaulay.

\begin{theorem}[L,LH,V, 24]
\label{seminars-algebraic-geometry-theorem-llhv24}
$$
\xymatrix{
&\{(\mathcal{C}_+,X):\mathcal{C}_+ \subseteq X\}
\ar[dr]\ar[dl]\\
\mathcal{H}_{a,b}\subset
\operatorname{Hilb}_{\mathbb{P}^3}
& &
\det(a,b) \subseteq U_d
}
$$
where
$U_d$ is the open subset of
$|\mathcal{O}_{\mathbb{P}^3}(d)|$
parametrizing smooth surfaces of degree $d$.
The locus $\det(a,b)\subseteq U_d$
parametrizes surfaces cut out by the
determinant of a matrix whose entries have
degree type $(a,b)$.
The scheme $\mathcal{H}_{a,b}$ is the
Hilbert scheme component parametrizing
the corresponding arithmetically
Cohen-McCaulay curves.
\end{theorem}

\noindent
This was proved using
\begin{theorem}[Ellingsrad 75]
\label{seminars-algebraic-geometry-theorem-ellin75}
$\mathcal{H}_{a,b} \subseteq
\operatorname{Hilb}_{\mathbb{P}^3}$
is an irreducible component gen. smooth
and of dimension…
\end{theorem}

\begin{proposition}
\label{seminars-algebraic-geometry-proposition-duplicate-2}
\begin{itemize}
\item $\pi^{-1}(X)
=
\coprod_{\mathcal{C}_+ \subseteq X}
|\mathcal{O}_\ell(\mathcal{C}_+)|$.
\item $H^0(X,\mathcal{C}_+)
=
\chi(X,\mathcal{C}_+)$.
\end{itemize}
\end{proposition}

\medskip\noindent
{\bf Noether-Lefschetz theorem.} 
If $d \geq  4$ and $X \in U_d$
is very general,
$$
\text{Pic}\mathbb{P}^3
\xrightarrow{\simeq}
\text{Pic}X 
=
\langle \mathcal{O}_X(1)\rangle
\simeq \mathbb{Z}
$$
where the firs map is a resolution.

\begin{definition}
\label{seminars-algebraic-geometry-definition-nl-locus}
The {\it Noether-Lefschetz locus} is
$$
NL(d)
=
\{ X \in U_d: \text{Pic}(X)\not\simeq
\mathbb{Z}\}.
$$
\end{definition}

\noindent
$$
NL(d) = \bigcup_{\operatorname{countable}}
\Sigma
$$
where $\Sigma$ are the NL-compononents.

\begin{theorem}[L,LH,V 24]
\label{seminars-algebraic-geometry-theorem-llhv24ii}
$\det(a,b)$ are NL components.
\end{theorem}

\begin{theorem}[Lopez, 91]
\label{seminars-algebraic-geometry-theorem-lopez-91}
A ver general degree $d$ surface
$X \supseteq \mathcal{C}_-$,
$$
\text{Pic}(X) = 
\langle \mathcal{O}_X(1),
\mathcal{O}_X(\mathcal{C}_-)\rangle.
$$
\end{theorem}

\noindent
This implies that $\det(a,b) \subseteq
NL(d)$,
and this implies that
$\det(a,b) \subseteq \Sigma$ for some
NL component.

Take $\{ X_s\} \subseteq \Sigma$
and $X_s = X \in \det(a,b)$.
There is a family of line bundles
$\{\mathcal{L}_s \in \text{Pic}(X)\setminus
\langle\mathcal{O}_{X_s}(\ell)\rangle\}$.
Since $\rho(X) = 2$ we can assume that
$\mathcal{L}_0 =
\mathcal{O}_X(\mathcal{C}_+)$.
$H^1(X,\mathcal{O}_X(\mathcal{C}_+))=0$
this implies that there exists
$\{\mathcal{C}_s \subseteq X_s\}$ such that
$\mathcal{C}_0 = \mathcal{C}_{+}$.
By the theorem of Elli, this implies that
$\mathcal{C}_s \in \mathcal{H}_{a,b}$ for
gen. s. This implies that
$X_s \supseteq \mathcal{C}_s$ is in 
$\det(a,b)$. This concludes the ideas
behind the proof.

\section{On a counterexample to a conjecture
of J. Harris for octic surfaces}
\label{seminars-algebraic-geometry-section-counterexample-harris-conj}

\noindent
Hossein Movasati, IMPA.
GADEPS Focused Meeting June 30, 2026, IMPA.

\medskip\noindent
Based on works by Hossein Movasati:
2021 - Special components of
Noether-Lefschetz. 2026 - On a
counterexample…

\medskip\noindent
Let $T $ be the parameter space of smooth
degree $d$ surfaces on $\mathbb{P}^3$ over
$\mathbb{C}$.
Consider
$$
x_0^d+x_1^d+x_2^d+x_3^d
+\sum_{\text{deg}x^\alpha =d}t_\alpha
x^\alpha,
$$
where there are 165 parameters $t_\alpha$.
$d=8$.

$$
NL_\alpha
=
\{t \in T|\rho(X_t)>1\}
$$
where $\rho$ is the Picard number.
As explained in Manuel's talk,
$NL_\alpha$ is a union of enumerable
sets and
$$
5=d-3 \leq  \text{codim}(\text{components} )
\leq \binom{d-1}{3}=h^{2,0}
$$
Components which attain the upper bound
are called general components, and those
for which the right-hand inequality
is strict are called special.
General components are dense in $T$.

Harris' conjecture: the number of special
components is finite.
In the 90's, Voisin found a counter-example
for large $d$.
The conjecture is known to hold
for $d=5$. For $d=6,7$ ``the conjecture
must be true''. Thus, the next interesting
case is $d=8$.

\begin{exercise}
\label{seminars-algebraic-geometry-exercise-computation-codim}
The codimension of the set of octic
surfaces containing a line is 5.
The codimension of the set of octic surfaces
containing a complete intersection
of type $(3,3)$ is 27.
Consider $x_0^8+x_1^8+x_2^8+x_3^8$
and $C_1:f_1=g_1=0$, $C_2:f_3=g_3=0$ where
the $f_i$ and $g_i$ are factors
of the Fermat equation. Then $V_{C_1}$
and $V_{C_2}$ are smooth and intersect
transversely, $\text{codim}(V_{C_1}\cap
V_{C_2})=5+27=32$.
\end{exercise}

\noindent
Let
$$
\delta_0=[C_1]+r[C_2]
\in H^2(X_0,Q) \cap H^{1,1} 
$$
where $C_1,C_2$ are as in the
exercise and $X_0$ is the Fermat,
and $r \in \mathbb{Q}$ is a parameter,
making the above a pencil of classes. 

Recall that Hodge conjecture for the case
of surfaces is called Lefschetz theorem
on $(1,1)$-classes.

Also,
$$
\delta_t \in H^2(X_t,Q),\qquad 
t \in (T,0).
$$
Local Noether-Lefschetz locus:
$$
V_{\delta_0}:=
\{t \in (T,0):
\delta_t \in H^2(X_t,\mathbb{Q})
\cap H^{1,1}(X_t)\}.
$$

In principle, this $V_{\delta_t}$
is in some component of the
Nother-Lefschetz locus.

\medskip\noindent
Conjecture: for all $r \in \mathbb{Q}$,
$r \neq  0$, except a finite number,
$V_{\delta_0}$ are distinct, smooth
varieties of codimension 31.
(Recall that $31 < 35 = h^{2,0}$.)


\section{Asymptotic behaviour of tropical
rank functions}
\label{seminars-algebraic-geometry-section-asymptotic-tropical-rank}

\noindent
Eduardo Vital, Bielefeld University.
Seminar of Algebra, IMPA.
August 26, 2026, 15:30, room 228.

\medskip\noindent
{\bf Abstract.} In this talk, we study the
asymptotic behaviour of two notions of rank
for linear series on tropical curves.
Although these two notions can differ for a
fixed divisor, we show that they have the
same asymptotic growth. This leads to a
natural notion of tropical volume.

We explain the main ideas behind this
result, including the role of uniform bounds
for the Baker--Norine rank and tropical
independence. As a consequence, both notions
of rank satisfy an asymptotic
Riemann--Roch formula. We also discuss
examples where the two ranks differ at
finite level, but nevertheless have
identical asymptotic behaviour. Finally, we
relate the tropical volume to the classical
volume of divisors through tropicalization.
This is joint work with Ana Maria Botero and
Alex K\"uronya.

\medskip\noindent
{\bf Comment.}
Here I learnt that scheme theory
for tropical geometry has not been
developed. The speaker and his collaborators
are interested in developing
Cech cohomology for higher dimensions
(and in developing scheme theory
for tropical geometry).

\section{\texorpdfstring
{$C^\infty$}{C-infinity}-Superrings and
\texorpdfstring
{$C^\infty$}{C-infinity}-Superschemes}
\label{seminars-algebraic-geometry-section-rizzo-smooth-superrings}

\noindent
Pedro Rizzo.
Seminar of Algebra, IMPA.
September 9, 2026, 15:30, room 228.

\medskip\noindent
{\bf Abstract.}
Superalgebraic and supergeometric structures
generalize classical nongraded structures to
graded ones. This extension introduces an
interplay between commuting and anticommuting
elements, enriching tools of commutative
algebra, differential geometry, and algebraic
geometry. Although originally motivated by
theoretical physics, supergeometry has
evolved into a robust mathematical framework
for studying objects such as supermanifolds
and superschemes.

In this talk, we introduce
$C^\infty$-superrings and their associated
geometric spaces, which we term
$C^\infty$-superschemes. We establish an
equivalence of categories between these
algebraic and geometric structures.
Furthermore, we present a generalization of
Batchelor’s theorem to $C^\infty$-superspaces
extending the classical result that every
smooth supermanifold is (noncanonically)
isomorphic to the sheaf of sections of a
vector bundle over its underlying smooth
manifold.

Conceptually, $C^\infty$-superrings
synthesize the features of $C^\infty$-rings
and superrings. Since smooth supermanifolds
form a core subclass of
$C^\infty$-superscheme, this framework
provides a natural scheme-theoretic
generalization of supergeometry in the smooth
setting. This is a joint work with Danilo
Olarte and Alexander Torres-Gómez (both at
Universidad de Antioquia) and based on the
preprints: arXiv:2508.09900v3 and
arXiv:2605.07169v2



\medskip\noindent
{\bf Introduction.}
$\mathbb{Z}_2$-structures
are inspired by commutative ``bosonic''
and anticommutative ``anti-bosonic''
elements. Its origins lie in theoretical
physics during the 70's and 80's, driven by
developments in supersymmetric field theory.

These notions are important in
supergeometry.
 
Synthetic differential geometry
also emerged in the 80's. Dubuc, Moerdijk,
Reyes and others constructed bridges between
differential and algebraic geometry.

Derived differential geometry attempts to
use techniques of algebraic geometry to
$C^\infty$ rings rather than commutative
rings.

Nishimura and Yetter explored this problem
using Weil superalgebras (which are a
particular case of $C^\infty$-superrings).
Liu and Yau introduced a notion of
$C^\infty$-superring to propose a work in
theoretical physics.

\medskip\noindent
{\bf Context and goal.}

Goal I: to establish structured mathematical
foundation of $C^\infty$ superspaces.

See results by Batchelor and Donagi-Witten.

Goal II: we propose the notion of Batchelor
space, which attempts to generalize
Batchelor's result.

\medskip\noindent
{\bf Category of $C^\infty$-rings}

\begin{definition}[C infinity ring]
\label{seminars-algebraic-geometry-definition-cinf-ring}

\end{definition}

\begin{itemize}
\item A morphism s…
\item $\mathfrak{C}$ is {\it finitely
generated} if…
\item An {\it ideal} is…
\item Obvious example: the
$\mathbb{R}$-algebra of smooth functions on
a differentiable manifold.
\item Every finitely-generated
$C^\infty$-ring is isomorphic to a quotient
of $C^\infty(\mathbb{R}^n)$ by some ideal.
\item An {\it $\mathbb{R}$-point} is a
$C^\infty$-ring morphism $x:\mathfrak{C}\to
\mathbb{R}$.
\item {\it Localization} is…
\item $\mathfrak{C}$ is called {\it fair} if
it is finitely generated and
germ-determined.
\end{itemize}

\medskip\noindent
{\bf $C^\infty$-ringed space.}

\begin{definition}
\label{seminars-algebraic-geometry-definition-cinf-ringed-space}

\end{definition}

\begin{itemize}
\item The stalk is localization…
\item Other expected results…
\end{itemize}

\noindent $\Spec$ and $\Gamma$ (sections)
functors are…

\begin{theorem}[Joyce]
\label{seminars-algebraic-geometry-theorem-joyce-adjoint-functors}
$\Spec$ and $\Gamma$ are adjoint.
\end{theorem}

\medskip\noindent
{\bf $C^\infty$-modules and sheaves.}

\begin{itemize}
\item A {\it module over a $C^\infty$-ring} 
is…
\item 
\end{itemize}

\medskip\noindent
{\bf Superring: basic theory.}

\begin{itemize}
\item A {\it superring} is a
$\mathbb{Z}_2$-graded unitary
supercommutative ring.
\item
$\mathfrak{R}=\mathfrak{R}_{\overline{0}}
\oplus \mathfrak{R}_{\overline{1}}$
such that $\mathfrak{R}_i,
\mathfrak{R}_j \subseteq
\mathfrak{R}_{i+j}$ (mod 2).
\item A {\it morphism} is…
\item An {\it ideal} is…
\item The quotient $\mathfrak{R}/I$
is a superring…
\item The {\it canonical superideal} is…
which gives the {\it superreduced ring}…
\end{itemize}

\noindent
Batchler's result concerns precisely when
can we ``lift'' a property from the reduced
ring to all of it.

\begin{example}
\label{seminars-algebraic-geometry-example-polynomial-superring}
Polynomial superring over the field
$\mathbb{K}$…
\end{example}

\begin{itemize}
\item We define $\Lambda(\mathfrak{R})$
as…
\item {\it Localization} is…
\end{itemize}

\medskip\noindent
{\bf $C^\infty$-superrings.}
Now we mix the two constructions so far.

\begin{definition}
\label{seminars-algebraic-geometry-definition-cinf-superring}
A {\it $C^\infty$-superring} is…
\end{definition}



\begin{proposition}
\label{seminars-algebraic-geometry-proposition-superring}
(Essentially describes the object
we shall call {\it split}.)
\end{proposition}

\noindent
[Missing content] We may thus define
super-versions of the adjoint functors
$\Spec$ and $\Gamma$…

\begin{theorem}[A]
\label{seminars-algebraic-geometry-theorem-a-superring}
$s\Spec$ and $s\Gamma$ are adjoint.
\end{theorem}

\begin{theorem}[B]
\label{seminars-algebraic-geometry-theorem-b-superring}
\end{theorem}

\medskip\noindent
{\bf Split spaces.}
{\it Split spaces} are… $\mathcal{X}$ is
{\it projected} if…

\begin{theorem}[C]
\label{seminars-algebraic-geometry-theorem-c-superring}
Every Batchelor space is split.
\end{theorem}


\medskip\noindent
{\bf Ongoing work.}
With Bilson Castro and Alexander Torres.
Based on Bruzzo et al., understand
$C^\infty$-superstacks.

\section{Quantised Picard Fuchs equations
and quantum periods}
\label{seminars-algebraic-geometry-section-garay-quantum-periods}

\noindent
Mauricio Garay.
Seminário de Geometria, Aritmética e
Equações Diferenciais de Períodos
(GADEPs), IMPA.
October 14, 2026, 10:30. Room 345.

\medskip\noindent
{\bf Abstract.}
The founders of quantum mechanics (Born,
Heisenberg, Dirac and Jordan) proposed an
interpretation of the spectrum as a
quantised version of normal forms in
classical mechanics, hence the name quantum
mechanics. The inverse of this normal form,
improperly called the quantum Birkhoff
normal form, defines a function which is a
quantum analog of the period integrals.
We show that like in the classical case,
these satisfy ordinary linear differential
equations whose degree is bounded by the
rank of homology like for ordinary
Picard-Fuchs theory. Computations for the
standard cubic oscillator show that the
coefficients of the quantum-period expansion
are integers, suggesting a new kind of
quantisation of the classical Gromov-Witten
theory of $(K_{\mathbb{P}^2})$.

\medskip\noindent
{\bf Upshot.}
Classically, integrating along a closed
curve in an energy level measures an
action, related to the area it encloses.
As the energy varies, these integrals
satisfy linear differential equations:
Picard--Fuchs equations. The talk replaces
commuting coordinates by quantum
operators and constructs a corrected
version of this area as a formal series.
Its main claim is that these quantum
periods still satisfy linear differential
equations controlled by the topology of
the classical fibres. Thus the familiar
geometry of periods survives quantization.
The observed integer coefficients suggest
a further enumerative connection; the
abstract presents this as a suggestion.

\medskip\noindent
{\bf Source.}
\href{https://impa.br/seminarios/}
{IMPA seminar announcement},
consulted October 14, 2026.
Abstract reproduced verbatim from the
text supplied by Daniel; mathematical
notation is typeset in LaTeX.

\medskip\noindent
{\bf Lecture notes (provisional).}
He first inverted the Birkhoff function.
The precise normalization is not yet
recorded. The next step studies an action
integral on a moving cycle in $f=t$:
$$
A(t)=\int_{\gamma_t}y\,dx.
$$
The Gauss--Manin derivative differentiates
its cohomology class as the fibre varies:
$$
A'(t)=\int_{\gamma_t}\eta,
\qquad
df\wedge\eta=dy\wedge dx.
$$
This explains the board notation
$(dy\wedge dx)/df$: it denotes a
Gelfand--Leray form, not ordinary division.
The wedge order fixes the sign.

He computes $A,A',A''$ to obtain a
Picard--Fuchs equation: a linear relation
among these functions with coefficients
depending on $t$. The idea is to find a
relation among the corresponding
cohomology classes, then integrate.
The exact equation and the link with
the inverted Birkhoff function remain
to be recorded. This is a reconstruction
of the reported board notes, not a
verified transcription.

\medskip\noindent
{\bf Quantum Birkhoff normal form.}
Born--Heisenberg--Jordan (1925),
as mentioned in the talk.
Now $q,p$ do not commute:
$$
[q,p]=qp-pq=h,\qquad h=i\hbar,
$$
$$
H=qp+q^3+p^3.
$$
The board uses $\mathbb{C}[[q,p,h]]$,
with multiplication subject to this
commutation relation.

The reported normal-form statement
involves a series
$$
\psi(qp)=a_0qp+a_1(qp)^2+\cdots.
$$
The idea is the noncommutative analogue
of Birkhoff normal form: transform $H$
into a function of $qp$.

{\it Formula to check.}
The transcription reads
$$
e^{-f/h}=\psi(qp).
$$
The role of $H$ and the transformation
are missing from this formula; it is
retained as an incomplete board note,
not asserted as a theorem.

\medskip\noindent
{\bf Quantum area and the classical limit.}
The next definition clarifies the role
of the inverse Birkhoff function:
$$
A(t)=\psi^{-1}(t).
$$
This is called the quantum area.
The reported expansion is
$$
A(t)=A_0(t)+hA_1(t)+h^2A_2(t)+\cdots.
$$
Our provisional understanding is that
$A_0$ recovers the classical action,
up to the normalization used in the talk,
and the higher terms are quantum
corrections. The aim is a quantum
Picard--Fuchs equation for this series.

An integrality statement was reported as
$$
V\in\mathbb{Z}[q,p]
\quad\Longrightarrow\quad
A\in\mathbb{Z}[t][[\hbar]].
$$
The definition of $V$, hypotheses,
normalization, and order of completion
still need checking. This is a board
transcription, not a verified theorem.
The distinction between $h$ and $\hbar$
also matters since $h=i\hbar$.

A possible spectrum formula was noted as
$$
\operatorname{Sp}_{BHJ}(H)=\psi(n\hbar).
$$
The first factor may be an integer $n$;
its range and the spectral convention
remain unconfirmed. This does not yet resolve
the earlier exponential formula.

\medskip\noindent
{\bf Quantum differential form
(reported theorem, unverified).}
The hypotheses include
$f\in\mathbb{C}[q,p]$ with Morse critical
points. A further hypothesis was heard
as “Infinite sing ifolita”, possibly
“isolated singularities at infinity”;
this interpretation is unconfirmed.
The reported conclusion is a formal form
$$
\alpha=\alpha(f)\in\Omega^1[[h]],
\qquad
\alpha=y\,dx+h\alpha_1+h^2\alpha_2+\cdots.
$$
The ambient space and allowed poles
have not yet been recorded.
The integral formula was transcribed as
$$
\frac{\partial A_i}{\partial t}
=\int_{\gamma_1}\alpha.
$$
The indices and derivative need checking:
it is unclear whether this concerns
$A$, its coefficients, or its derivative.
We do not identify these alternatives.
The intended picture is that quantum
data can be expressed using a formal
series of differential forms on the
classical fibres.

The reported Picard--Fuchs “degree” is
$2g(f^{-1}(t))$. Here degree presumably
means differential order. The announcement
states a bound by homology rank.
Equality, compactification, and possible
contributions at infinity remain to check.
See the background reference below.

\medskip\noindent
{\bf Deformation quantization and ordering.}
This calculation uses $pq-qp=h$,
the opposite sign to the commutator
recorded earlier. With $h$ central,
$$
qpq=q(qp+h)=q^2p+hq.
$$
Normal ordering places $q$ before $p$.
On commuting polynomial symbols, the
corresponding star product satisfies
$$
q\ast_N p=qp,\qquad p\ast_N q=qp+h,
$$
$$
q\ast_N p\ast_N q=q^2p+hq.
$$
Symmetric ordering gives the Moyal
convention
$$
q\ast_M p=qp-\frac{h}{2},
\qquad
p\ast_M q=qp+\frac{h}{2}.
$$
Both have $p\ast q-q\ast p=h$.
The terms $hq$ and $h/2$ correct
algebraic slips in the transcription.

An ordering map $Q$ encodes operator
multiplication on symbols by
$$
Q(F\ast G)=Q(F)Q(G).
$$
Compare {\it Geometric prequantization},
Sections
\ref{geometric-prequantization-%
section-dirac-rule} and
\ref{geometric-prequantization-%
section-kostant-souriau-formula}.
Prequantization supplies a commutator
representation; it is not itself the
invertible symbol map $Q$ used here.

The analogous reordering of field modes
appears in {\it Vertex algebras},
Definition
\ref{vertex-algebras-%
definition-creat-annih-opera}
and Exercise
\ref{vertex-algebras-%
exercise-most-important}:
$$
:a(z)b(z):=a(z)_+b(z)+b(z)a(z)_-.
$$
The common mechanism is reordering
noncommuting operators and retaining
commutator corrections. The field
construction involves infinitely many
modes and is not an identification of
the two products.

\medskip\noindent
{\bf Trace formula
(reported theorem, unverified).}
The final formula was transcribed as
$$
A'(t)=\operatorname{Tr}
\left(\frac{k}{f-t}\right).
$$
The meaning of $k$, the trace, the inverse
of $f-t$, and the hypotheses were not
recorded. The formula is retained as a
board transcription, without identifying
the trace with an ordinary matrix or
operator trace. No precise theorem
reference has been verified.

\medskip\noindent
{\bf Background and connections.}
See {\it Variations of Hodge Structures},
Section \ref{variations-of-hodge-structures-%
section-quantum-periods-background},
for periods, Gauss--Manin,
Picard--Fuchs equations, and links to
the quantization and Seiberg--Witten notes.

\section{Irrationality loci and
Noether--Lefschetz divisors in the moduli
space of K3s}
\label{seminars-algebraic-geometry-section-chen-irrationality-k3}

\noindent
Nathan Chen.
Algebra Seminar, IMPA.
October 14, 2026, 15:30.
Room 228.

\medskip\noindent
{\bf Abstract.}
The goal of this talk is twofold. In the
first part, we will present a family of
surfaces where the locus of points in the
base with constant degree of irrationality is
not constructible. In the second part, we
will explore the notation of gonality loci in
the moduli space of K3 surfaces. Among other
things, we prove that the polarized gonality
loci are always unions of Noether--Lefschetz
divisors.

\medskip\noindent
{\bf Source.}
\href{https://impa.br/seminarios/}
{IMPA seminar announcement},
consulted October 14, 2026.


\medskip\noindent
{\bf Setting.}
We work over $\mathbb{C}$.
The base $T$ is quasi-projective and
irreducible. Maps
$\mathcal{C}\to T$ or $\mathcal{S}\to T$
are called families, with flatness,
properness, smoothness, etc., specified
as needed. Below the curves and surfaces
are smooth and projective.

\medskip\noindent
{\bf Motivation: curves.}
Let $\mathcal{C}\to T$ be a smooth
projective family of curves.
The gonality is
$$
\operatorname{gon}(C)=\min_f\deg(f),
\qquad f:C\to\mathbb{P}^1,
$$
where $f$ ranges over nonconstant
morphisms.
For $\dim T=1$, gonality can only drop
at special points:
$$
\operatorname{gon}(C_0)
\leq\operatorname{gon}(C_t)
$$
for general $t$.
Constant-gonality loci are locally closed.
These are the motivating facts from
curve theory reported in the talk;
an exact reference remains to be added.
The special fibre belongs on the left,
correcting the transcription.

\medskip\noindent
{\bf Rationality: surfaces.}
For higher dimensions, start by asking
about rationality.
Recall the invariants
$$
q(S)=h^1(S,\mathcal{O}_S),
\qquad
P_2(S)=h^0(S,\omega_S^{\otimes 2}).
$$

\begin{theorem}[Castelnuovo]
\label{seminars-algebraic-geometry-theorem-chen-castelnuovo}
A smooth projective complex surface $S$
is rational if and only if
$$
q(S)=P_2(S)=0.
$$
\end{theorem}

\noindent
Reference: Beauville,
{\it Complex Algebraic Surfaces},
second edition, Chapter V, Theorem V.1.
This corrects the garbled equality
in the board transcription.

\medskip\noindent
{\bf Consequence reported in the talk.}
In a smooth projective family of surfaces
over a connected smooth complex base,
if $S_0$ is rational, then every $S_t$
is rational. The proof uses constancy
of $q$ and $P_2$ in the family.
The deformation-invariance reference
for this step remains to be added.

\medskip\noindent
{\bf Higher dimensions.}

\begin{theorem}[Hassett--Pirutka--Tschinkel]
\label{seminars-algebraic-geometry-theorem-chen-hpt}
There exists a smooth family
$\mathcal{X}\to T$ of projective
fourfolds such that:
\begin{enumerate}
\item A very general $X_t$ is not stably
rational, hence is not rational.
\item The locus of $t$ for which $X_t$
is rational is dense in the Euclidean
topology.
\end{enumerate}
Each fourfold is a quadric surface
bundle over $\mathbb{P}^2$.
\end{theorem}

\noindent
Reference: Hassett, Pirutka and Tschinkel,
{\it Stable rationality of quadric
surface bundles over surfaces},
Acta Mathematica 220 (2018), 341--365,
Theorem 1;
\href{https://arxiv.org/abs/1603.09262}
{author manuscript}.

\medskip\noindent
{\bf Degree of irrationality.}
If rationality behaves this way, what
happens to a numerical measure of
irrationality?

\begin{definition}
\label{seminars-algebraic-geometry-definition-chen-irrationality}
For a smooth projective variety $X$
of dimension $n$, its
{\it degree of irrationality} is
$$
\operatorname{irr}(X)=\min_f\deg(f),
\qquad f:X\dashrightarrow\mathbb{P}^n,
$$
where $f$ ranges over dominant
generically finite rational maps.
\end{definition}

\noindent
Thus $\operatorname{irr}(X)=1$ precisely
when $X$ is rational.
For a smooth family of surfaces,
can this invariant drop?

\begin{theorem}[Chen--Stapleton]
\label{seminars-algebraic-geometry-theorem-chen-stapleton-specialization}
Let $\mathcal{S}\to T$ be a smooth
projective family of surfaces over a
smooth irreducible complex curve,
with $0\in T$ and $q(S_0)=0$.
Then, for very general $t$,
$$
\operatorname{irr}(S_0)
\leq\operatorname{irr}(S_t).
$$
\end{theorem}

\noindent
Reference: Chen and Stapleton,
{\it Fano hypersurfaces with arbitrarily
large degrees of irrationality},
Forum of Mathematics, Sigma 8 (2020),
e24, Proposition C;
\href{https://arxiv.org/abs/1908.02803}
{author manuscript}.
The published statement assumes
$q(S_t)=0$ for very general $t$;
constancy of $q$ gives the formulation
above. The name transcribed as
Stephensen is Stapleton.

\medskip\noindent
{\bf Consequence of today's talk
(reported, unverified).}
Attribution recorded as B\_CD;
the full names and exact reference
remain to be checked.
There exists a smooth family of
projective surfaces $\mathcal{S}\to T$
and an integer $d\geq3$ such that:
\begin{enumerate}
\item For very general $t$,
$\operatorname{irr}(S_t)=d$.
\item There are infinitely many parameters
with $\operatorname{irr}(S_t)<d$,
and this locus is Euclidean dense.
\end{enumerate}

\medskip\noindent
{\bf Bootstrapping.}
The speaker says that bootstrapping a
certain family of K3 surfaces produces
families of surfaces of general type.
Here this means using the K3 examples
as input to a further construction,
transferring their irrationality
phenomenon to the new surfaces.
This is an interpretation of the phrase;
the actual construction has not yet
been recorded.


\medskip\noindent
The speaker also mentioned products with
$\mathbb{P}^n$ for higher-dimensional
examples. The required control of
irrationality was not recorded.

\medskip\noindent
{\bf The degree-two case.}
Chen--Esser prove that irrationality
at most two specializes in families
of integral klt surfaces over a smooth
complex curve. In a smooth projective
family with very general irrationality
two, all fibres have irrationality two:
a rational fibre would force rationality
throughout, by the earlier argument.

Reference: Chen and Esser,
{\it Group actions and irrationality
in surface families},
Selecta Mathematica 31 (2025), article 96.
Theorem 1.2 and Corollary 3.1 were
checked in the
\href{https://arxiv.org/abs/2410.18427}
{author manuscript}, version 1.

\medskip\noindent
{\bf K3s and hyperelliptic curves.}
The following equivalence was reported
in the seminar; its exact source
remains to be verified.
For a projective K3 surface $X$:
\begin{enumerate}
\item $\operatorname{irr}(X)=2$.
\item $X$ contains a smooth hyperelliptic
curve of genus at least two.
\item $X$ contains a smooth genus-two
curve, or admits a genus-one fibration
$X\to\mathbb{P}^1$ with a multisection
of degree at most two.
\end{enumerate}
The implication from (1) to (2) was
proved in the seminar; the proof was
not transcribed. In (3), $X$ replaces
the accidental $C$ in the notes.

\medskip\noindent
{\bf Lattice criterion
(board transcription, unverified).}
A further equivalence for
$\operatorname{irr}(X)=2$ used a line
bundle or one of these intersection
matrices in $\operatorname{Pic}(X)$:
$$
\begin{pmatrix}0&2\\2&0\end{pmatrix},
\qquad
\begin{pmatrix}0&1\\1&-2\end{pmatrix},
\qquad
\begin{pmatrix}0&2\\2&-2\end{pmatrix}.
$$
The other alternative was transcribed
as an ample $L$ with $L^2=0$.
This is impossible: an ample line
bundle on a surface has positive square.
A genus-two curve has square two by
adjunction, suggesting a square-two
condition. The precise positivity and
base-point-freeness hypotheses must
be checked before replacing the claim.
Embedding and primitivity conditions
for the lattices also remain to check.

\medskip\noindent
{\bf Quasi-polarized moduli.}
Write $\mathcal{F}_g$ for the moduli
space of quasi-polarized K3 surfaces
$(X,L)$, with $L$ primitive, big and nef,
and
$$
L^2=2g-2.
$$

\medskip\noindent
{\bf Corollary reported in the seminar.}
For each $g\geq3$, $\mathcal{F}_g$
contains infinitely many divisors whose
generic member has irrationality two.
The attribution and published reference
have not yet been identified.
The phrase about finitely many covers
of $\mathbb{F}_n$ was also recorded;
its intended quantifiers are unclear.

\medskip\noindent
{\bf AI-reconstructed proof outline
(incomplete, not a verified proof).}
Impose an extra divisor class which
produces a degree-two map to a rational
surface. A possible route is to seek
rank-two Picard lattices generated by
$L,H$ with intersection matrix
$$
\begin{pmatrix}
2g-2&a\\
a&2
\end{pmatrix}.
$$
Require $H$ nef and base-point-free.
Riemann--Roch and vanishing then give
$h^0(X,H)=3$; since $H^2=2$, the map
defined by $|H|$ has degree two onto
$\mathbb{P}^2$. A K3 is not rational,
so its irrationality is exactly two.

For an admissible primitive embedding,
requiring $H$ to remain algebraic gives
a Noether--Lefschetz divisor in
$\mathcal{F}_g$. To conclude, produce
infinitely many such embeddings with
$L,H$ nef in the same chamber and prove
that their divisors are distinct.
The varying determinants
$$
4g-4-a^2
$$
suggest a way to distinguish lattices,
but do not prove the embedding and
positivity claims. These lattice steps
were not transcribed and remain a gap
in this reconstruction.

A square-zero pencil $|E|$ instead
gives a genus-one fibration over
$\mathbb{P}^1$. A degree-two
multisection supplies a degree-two
linear series on its generic fibre,
leading to a rational map to a ruled
surface, birational to $\mathbb{F}_n$.
The pencil alone is not a degree-two
map of surfaces.

Relevant classical reference:
Saint-Donat, {\it Projective models
of K-3 surfaces}, American Journal
of Mathematics 96 (1974), 602--639,
Theorem 5.2 on hyperelliptic systems.
Its use is cross-checked in
Garbagnati--Sarti,
{\it Projective models of K3 surfaces
with an even set}, Section 4,
\href{https://sarti.pages.math.cnrs.fr/%
home/ottonodali1312.pdf}
{author manuscript}.
This does not verify the infinite-divisor
assertion.

\medskip\noindent
{\bf Higher gonality loci.}
We now relate irrationality to the
existence of higher-gonality curves.

The unpolarized locus
$G_{e;g}\subseteq\mathcal{F}_g$
consists of $(X,L)$ admitting a smooth
curve $C\subset X$ with
$$
\operatorname{gon}(C)=e,\qquad g(C)\geq2.
$$
The curve need not lie in a multiple
of the chosen polarization.

The polarized locus
$P^m_{e;g}\subseteq\mathcal{F}_g$
requires a smooth $C\in|mL|$ with
$$
\operatorname{gon}(C)=e.
$$
Here $m>0$. Smoothness is made explicit
to match the preceding definition;
if singular curves were intended, the
gonality convention needs checking.

\medskip\noindent
{\bf Example reported in the seminar
(unverified).}
For every
$$
g\equiv7\pmod{9},
$$
the locus $G_{3;g}$ has an irreducible
component of codimension two not
contained in $G_{2;g}$.
The congruence sign replaces the
approximation sign in the transcription.
The exact range of $g$ and the lattice
criterion were not recorded.

\medskip\noindent
{\bf Moretti example
(reported, not fully verified).}
A polarized K3 surface with
$$
\operatorname{Pic}(X)=\mathbb{Z}L,
\qquad L^2=10
$$
was said to have
$\operatorname{irr}(X)=3$ but no
smooth trigonal curve. Here $g=6$.
Thus a degree-three rational map need
not arise from a smooth trigonal curve.
Related source: Moretti,
{\it The polarized degree of
irrationality of K3 surfaces},
\href{https://arxiv.org/abs/2303.07289}
{arXiv:2303.07289}.
An exact theorem covering both parts
of this example remains to be checked.

\medskip\noindent
{\bf Final theorem
(reported, unverified).}
The loci $P^m_{e;g}$ are either empty,
all of $\mathcal{F}_g$, or finite
unions of Noether--Lefschetz divisors.
The exact hypotheses and reference
were not recorded. Check whether the
statement uses gonality equal to $e$
or at most $e$, and the loci themselves
or their closures. The announcement
confirms the union-of-divisors theme,
but not these finer distinctions.

% --- End seminars-algebraic-geometry.tex ---

% --- Begin 16th-alga-meeting-2026.tex ---
\chapter{16th ALGA meeting - Commutative Algebra and Algebraic Geometry}
\label{chapter-16th-alga-meeting-2026}


Commutative Algebra and
Algebraic Geometry}

\medskip\noindent
{\bf Conference information.}

\href{%
https://impa.br/evento/%
16th-alga-meeting-commutative-algebra-%
and-algebraic-geometry/%
}{16th ALGA meeting -
Commutative Algebra and
Algebraic Geometry}

\href{%
https://drive.google.com/file/d/%
1hMHDOJ7PgjAUEXZZZTPlzIAJmYVuu0NM/%
view?usp=sharing%
}{Program.}

August 31 to September 4, 2026.

IMPA - Ricardo Ma\~n\'e Auditorium,
Rio de Janeiro,
Rio de Janeiro, Brazil.

Information accessed August 29, 2026.

\medskip\noindent
{\bf Organizing Committee.}

Carolina Araujo (IMPA),
Juliana Coelho (UFF),
Suely Lima (IMPA),
Marco Pacini (UFF),
Eduardo Vital
(Bielefeld University, Germany).

\medskip\noindent
{\bf Scientific Committee.}

Alicia Dickenstein (UBA, Argentina),
Abramo Hefez (UFF),
Marcos Jardim (Unicamp),
Ragni Piene
(University of Oslo, Norway),
Aron Simis (UFPE),
Israel Vainsencher (UFMG).

\section{Talks of special interest}
\label{16th-alga-meeting-2026-section-talks-special-interest}

These recommendations are based on current
Heap Project interests in Hodge theory,
birational geometry, moduli of sheaves,
stability conditions, abelian varieties,
and Cremona transformations.

\begin{enumerate}

\item Mihnea Popa,
{\it Negative K-theory and Hodge theory}.
Tuesday, September 1, 2026,
4:30--5:30 p.m.
This is the closest match to the
Hodge-theoretic and derived interests.

\item Svetlana Makarova,
{\it Contractions via moduli of sheaves}.
Friday, September 4, 2026,
11:30 a.m.--12:30 p.m.
This directly connects moduli of sheaves
with birational contractions.

\item Alex Massarenti,
{\it Rationality, unirationality, and
rational connectedness}.
Tuesday, September 1, 2026,
9:00--10:00 a.m.
This belongs to the central birational
geometry direction.

\item Nivaldo Medeiros,
{\it Homaloidal hypersurfaces and a product
formula for gradient maps}.
Wednesday, September 2, 2026,
11:45 a.m.--12:45 p.m.
This is directly adjacent to Cremona
transformations and rational maps.

\item Renato Vidal Martins,
{\it Stability conditions for coherent
systems on integral curves}.
Thursday, September 3, 2026,
10:00--11:00 a.m.
This overlaps strongly with stability,
sheaves, and moduli.

\end{enumerate}

\medskip\noindent
{\bf Also highly relevant.}

\begin{itemize}

\item Anita Rojas,
{\it Abelian varieties: Subvarieties and
Decompositions},
Thursday, September 3, 2026,
11:30 a.m.--12:30 p.m.

\item Margarida Melo,
{\it Compactified spaces of roots of
bundles over curves via logarithmic
geometry},
Monday, August 31, 2026,
10:00--11:00 a.m.

\item Fillipo Viviani,
{\it On the classification of compactified
universal Jacobians},
Friday, September 4, 2026,
2:00--3:00 p.m.

\item Olivier Martin,
{\it Rationally inequivalent points on
complete intersections},
Tuesday, September 1, 2026,
2:00--3:00 p.m.

\item Alex Abreu,
{\it Computing Classes on
$\overline{\mathcal{M}}_{g,n}$},
Monday, August 31, 2026,
3:00--4:00 p.m.

\item Giusi Capobianco,
{\it The tropical theta divisor},
during the PYR 3 session on
Thursday, September 3, 2026,
2:00--3:00 p.m.

\end{itemize}

If only one complete day is possible,
Tuesday has the strongest concentration
of especially relevant talks.

\section{Subsets of $P^4$ with no four
points on a plane}
\label{16th-alga-meeting-2026-section-subsets-p4}

\medskip\noindent
{\bf Lecture.}

Felipe Voloch,
University of Canterbury, New Zealand.

\medskip\noindent
{\bf Schedule.}

Monday, August 31, 2026,
9:00--10:00 a.m.,
IMPA - Ricardo Ma\~n\'e Auditorium.

\medskip\noindent
A {\it track} is a subset
of $\mathbb{P}^4$ containing no four
distinct points on a plane.
This generalizes the notion of
{\it ?}, which are sets in $\mathbb{P}^2$
with no three points in a line.
A theorem by Segre says that over
finite field $\mathbb{F}_q$,
any such set is a conic if and only
if its size is precisely [number].

\begin{theorem}
\label{16th-alga-meeting-2026-theorem-voloch}
If 3 is not a square in $\mathbb{F}_q$,
then the set
$$
[set]
$$
is a complete track of size $2q+1$.
\end{theorem}

\section{Compactified spaces of roots of
bundles over curves via logarithmic
geometry}
\label{16th-alga-meeting-2026-section-compactified-roots-bundles}

\medskip\noindent
{\bf Lecture.}

Margarida Melo,
Roma Tre University, Italy.

\medskip\noindent
{\bf Schedule.}

Monday, August 31, 2026,
10:00--11:00 a.m.,
IMPA - Ricardo Ma\~n\'e Auditorium.

\medskip\noindent
{\bf Upshot.}
After defining a specific moduli problem,
irreducible and connected components
are described. Logarithmic geometry
is an essential ingredient.
(I think that boundary points of the
compactified moduli might be encoded
by combinatorial data, just like my SR
schemes.)
The most important objective of this work
is to describe the strata: they are 
in bijective correspondence with
connected components $S_\Gamma$,
where $\Gamma$ is a graph.
Question: will the strata
of the logarithmic compactification
match the original compactification? Yes.



\medskip\noindent
{\bf Universal spaces of roots.}

Let $r\geq 2$, let $k$ be an integer,
and let
$$
\mathbf{m}=(m_1,\ldots,m_n).
$$
Consider the universal smooth pointed
curve
$$
\pi\colon \mathcal{C}_{g,n}
\longrightarrow \mathcal{M}_{g,n}
$$
with sections $\sigma_1,\ldots,\sigma_n$.

\begin{definition}
The {\it universal bundle of type
$(k,\mathbf{m})$} is
$$
\mathcal{N}_{k,\mathbf{m}}
=
\omega_{\pi}^{\otimes k}
\otimes
\mathcal{O}_{\mathcal{C}_{g,n}}
\left(
\sum_{i=1}^{n}m_i\sigma_i
\right).
$$
\end{definition}

Thus, over a smooth pointed curve
$(C,p_1,\ldots,p_n)$, the bundle whose
roots are sought is
$$
\mathcal{N}_{C,k,\mathbf{m}}
=
\omega_C^{\otimes k}
\otimes
\mathcal{O}_C
\left(
\sum_{i=1}^{n}m_ip_i
\right).
$$

\begin{definition}
An {\it $r$th root of
$\mathcal{N}_{C,k,\mathbf{m}}$} is a line
bundle $L$ on $C$ together with an
isomorphism
$$
L^{\otimes r}
\longrightarrow
\mathcal{N}_{C,k,\mathbf{m}}.
$$
The {\it universal space of $r$th roots}
is the stack
$$
\mathcal{S}^{1/r}_{g,n}
(k,\mathbf{m})
$$
parametrizing such pointed curves,
line bundles, and isomorphisms.
\end{definition}

A necessary numerical condition is
$$
r\mid
k(2g-2)+\sum_{i=1}^{n}m_i.
$$
For an $n$th root, the notation requested
in the talk is
$$
\mathcal{S}^{1/n}_{g,n}
(k,\mathbf{m}).
$$

Over the smooth locus, forgetting the
root gives a finite cover of
$\mathcal{M}_{g,n}$.  When $r$ is
invertible, each nonempty geometric fiber
has $r^{2g}$ isomorphism classes.

\medskip\noindent
{\bf Log curves and tropicalization.}

An $n$-marked logarithmic curve of genus
$g$ is an $n$-marked nodal curve of genus
$g$, together with logarithmic data that
records a deformation parameter at each
node.  This is what
$\mathcal{M}_{g,n}^{\log}$ parametrizes.
Its tropicalization is the dual graph of
the curve, with the parameter at each node
used as the length of the corresponding
edge.

The following theorem is due to
Cavalieri, Chan, Ulirsch, and Wise.

\begin{theorem}
\label{16th-alga-meeting-2026-theorem-log-curves-tropicalization}
Let $g,n\geq 0$ and suppose that
$$
2g-2+n>0.
$$
The moduli functor
$$
a^{*}\mathcal{M}_{g,n}^{\mathrm{trop}}
$$
is represented by a logarithmic algebraic
stack that is logarithmically \'etale over
the base field.  The natural morphism
$$
\mathrm{trop}_{g,n}\colon
\mathcal{M}_{g,n}^{\log}
\longrightarrow
a^{*}\mathcal{M}_{g,n}^{\mathrm{trop}},
$$
which sends a logarithmic curve to its
family of dual tropical curves, is strict,
smooth, and surjective.
\end{theorem}

\medskip\noindent
{\bf Known compactifications.}

\begin{itemize}

\item Beauville's 1977 compactification
uses admissible double covers.  Over a
smooth curve, an unramified double cover
is equivalent to a nontrivial line bundle
$\eta$ satisfying
$$
\eta^{\otimes 2}=\mathcal{O}_C.
$$
Admissible covers supply boundary points
over stable curves and compactify the
corresponding Prym data.

\item Abreu--Pacini--Secco construct the
2022 moduli space of roots of divisors on
tropical curves.  Its combinatorics is
governed by a poset of flows, and it is
related to the skeleton of Jarvis's space
of nets of limit roots.

\end{itemize}

\medskip\noindent
{\bf Aim of the talk.}

Describe a nice compactification of this
root space over the Deligne--Mumford
compactification.  Here nice means that
the boundary strata can be described
explicitly through tropicalization.
The desired picture is
$$
\xymatrix@C=4em@R=3em{
\mathcal{S}^{1/r}_{g,n}
(k,\mathbf{m})
\ar@{^{(}->}[r]\ar[d]
&
\overline{\mathcal{S}}^{1/r}_{g,n}
(k,\mathbf{m})
\ar[d]
\\
\mathcal{M}_{g,n}
\ar@{^{(}->}[r]
&
\overline{\mathcal{M}}_{g,n}.
}
$$

The strata of
$\overline{\mathcal{M}}_{g,n}$ are
controlled by stable marked dual graphs.
In the tropical root space, these graphs
are decorated by flows recording how a
root degenerates at the nodes.  Edge
contractions describe incidence among
the strata.  The aim is to obtain an
equally explicit stratification of
$\overline{\mathcal{S}}^{1/r}_{g,n}
(k,\mathbf{m})$.

\medskip\noindent
{\bf References.}

\begin{itemize}

\item A. Beauville,
\href{%
https://doi.org/%
10.1007/BF01418373%
}{{\it Prym varieties and the Schottky
problem}},
{\it Inventiones Mathematicae} 41
(1977), 149--196.

\item D. Abramovich, L. Caporaso,
and S. Payne,
\href{%
https://arxiv.org/abs/%
1212.0373%
}{{\it The tropicalization of the moduli
space of curves}}.

\item R. Cavalieri, M. Chan,
M. Ulirsch, and J. Wise,
\href{%
https://arxiv.org/abs/%
1704.03806%
}{{\it A moduli stack of tropical
curves}}.

\item A. Abreu and M. Pacini,
\href{%
https://arxiv.org/abs/%
1806.05527%
}{{\it The universal tropical Jacobian
and the skeleton of the Esteves' universal
Jacobian}}.

\item M. Melo, S. Molcho, M. Ulirsch,
and F. Viviani,
\href{%
https://arxiv.org/abs/%
2108.04711%
}{{\it Tropicalization of the universal
Jacobian}}.

\item A. Abreu, M. Pacini,
and M. Secco,
\href{%
https://arxiv.org/abs/%
2204.13471%
}{{\it On moduli spaces of roots in
algebraic and tropical geometry}}.

\item L. Caporaso, M. Melo,
and M. Pacini,
\href{%
https://arxiv.org/abs/%
1902.07803%
}{{\it Tropicalizing the moduli space of
spin curves}}.

\end{itemize}

\section{Limit canonical series}
\label{16th-alga-meeting-2026-section-limit-canonical-series}

\medskip\noindent
{\bf Lecture.}

Omid Amini,
\'Ecole Polytechnique, France.

\medskip\noindent
{\bf Schedule.}

Monday, August 31, 2026,
11:30 a.m.--12:30 p.m.,
IMPA - Ricardo Ma\~n\'e Auditorium.

\medskip\noindent
{\bf Abstract.}

No abstract is available on the official
site as of August 29, 2026.

\section{On the analytical equivalence of
plane curves}
\label{16th-alga-meeting-2026-section-analytical-equivalence-curves}

\medskip\noindent
{\bf Lecture.}

Marcelo Escudeiro,
UEM.

\medskip\noindent
{\bf Schedule.}

Monday, August 31, 2026,
2:00--3:00 p.m.,
IMPA - Ricardo Ma\~n\'e Auditorium.

\medskip\noindent
{\bf Abstract.}

No abstract is available on the official
site as of August 29, 2026.

\section{Computing Classes on
$\overline{\mathcal{M}}_{g,n}$}
\label{16th-alga-meeting-2026-section-computing-classes-moduli}

\medskip\noindent
{\bf Lecture.}

Alex Abreu,
UFF.

\medskip\noindent
{\bf Schedule.}

Monday, August 31, 2026,
3:00--4:00 p.m.,
IMPA - Ricardo Ma\~n\'e Auditorium.

\medskip\noindent
{\bf Abstract.}
An important idea of this talk is to
provide language for
describing classes in the moduli space of
curves.

\medskip\noindent
{\bf AI-written summary (speculative).}

The talk will likely begin with the moduli space
$\overline{\mathcal{M}}_{g,n}$ of stable genus-$g$ curves with
$n$ marked points and with the problem of computing the cycle classes of
geometrically meaningful loci inside it.  Such classes live in its Chow or
cohomology ring, and in favorable cases they belong to the tautological ring,
where they can be described using $\psi$- and $\kappa$-classes together with
boundary strata indexed by stable graphs.  Thus ``computing a class'' should
mean translating a geometric construction into an explicit combination of
these standard generators, using operations such as pullback, pushforward,
and gluing along the boundary.

Given Abreu's work on compactified Jacobians, Abel maps, and degenerations of
curves, I would expect the main examples to come from divisor theory on the
universal curve or universal Jacobian---possibly Brill--Noether-type loci or
related cycles pulled back to $\overline{\mathcal{M}}_{g,n}$.  The difficult
part is what happens when a smooth curve degenerates: multidegrees and
stability conditions can change across boundary walls, producing correction
terms supported on boundary strata.  The talk may explain how the combinatorics
of dual graphs, perhaps together with tropical geometry, organizes those
corrections and leads to formulas that can actually be evaluated.

\section{Crossing the boundary, in search
of limits (Eduardo Esteves: The First
Thirty Years of a Mathematical Journey)}
\label{16th-alga-meeting-2026-section-crossing-boundary-gatto}

\medskip\noindent
{\bf Lecture.}

Letterio Gatto,
Politecnico di Torino, Italy.

\medskip\noindent
{\bf Schedule.}

Monday, August 31, 2026,
4:30--5:30 p.m.,
IMPA - Ricardo Ma\~n\'e Auditorium.

\medskip\noindent
This talk was in honor of Eduardo's work.
His main PhD contribution concerns
a moduli space of related to the Jacobian
variety. The second,
independent of the first, result in his PhD
concerns Weierstrass points. Eduardo did his
PhD in MIT under the supervision
of Kleiman.

Smooth projective curves of genus $g$
are nice. One would wish to consider
families
of them. But even a pencil of
irreducible conics can have degenerate
members. One cannot find nontrivial
families without degenerate maps.
The moduli space $\mathcal{M}_g$ of smooth
curves is not compact. Deligne and Mumford
(1969) achieved a ``minimal''
compactification
$\overline{\mathcal{M}}_g$:
$$
\overline{\mathcal{M}}_g\setminus\mathcal{M}_g
=
\{C|C \operatorname{stable}\}.
$$
The central question in Eduardo's
mathematics
is: let $C$ be a stable curve
and $\pi:\mathcal{C} \to B$ a smoothing.
How does the classically defined
geometry for smooth curves
behave when limiting to a curve in the
boundary?
$$
\lim_{\eta \to 0}
\operatorname{Jac}(\mathcal{C}_\eta)=?,\qquad 
\lim_{\eta\to 0} g_d^r(\mathcal{C})=?
$$
See Compactifying the relativa Jacobian over
families of reduced curves where Eduardo
achieves the compactification by
constructing an algebraic space
parametrizing simple, torsion-free, rank 1
sheaves, which becomes a scheme after an
étale base change. The construction …

After constructing the relative Jacobian,
there are many things to do: Abel maps,
autoduality, Brill-Noether…

\section{Rationality, unirationality, and
rational connectedness}
\label{16th-alga-meeting-2026-section-rationality-unirationality}

\medskip\noindent
{\bf Lecture.}

Alex Massarenti,
University of Ferrara, Italy.

\medskip\noindent
{\bf Schedule.}

Tuesday, September 1, 2026,
9:00--10:00 a.m.,
IMPA - Ricardo Ma\~n\'e Auditorium.

\medskip\noindent
Aim: construct a family of smooth
unirational
$X_6 \subseteq \mathbb{P}(1,1,1,2,3)$.

This is the analogue of what we know for
quartics $X_4 \subseteq \mathbb{P}^4$.

The following construction is due to Segre.
$X_4 \subseteq \mathbb{P}^4$ is Fano
and thus is rationally connected.
Let $C \subseteq X_4$ be a rational curve.
Pick a point $x \in C$ and
let $C_x$ be the conic parametrizing
lines through $x$ intersecting $X_4$
multiplicity at least 3 in $x$.

Surface conic bundle $S_C$
over $\mathbb{C}$ $\implies $
Rational.
$\varphi:S_c \dashrightarrow
X_4$. Then the image $\varphi(S_C) = S
\subseteq X_4$ is a rational surface
containing $C$.

Redo the same construction replacing $C$
with $S$.

[Missing [art]

These $S$ exist and Segre called them there
surfaces with separable assymptotics.
With them can get a family of dimension 54
in $\mathbb{P}^{69}$
of smooth unirational quartics

Recall the aim is to construct a smooth
unirational $X_6 \subseteq
\mathbb{P}(1,1,1,2,3)$.
Now construct a family of dimension 50
of unirational $X_6$.



\section{Scrollar invariants of singular
curves on toric surfaces}
\label{16th-alga-meeting-2026-section-scrollar-invariants}

\medskip\noindent
{\bf Lecture.}

Karl Christ,
University of Turin, Italy.

\medskip\noindent
{\bf Schedule.}

Tuesday, September 1, 2026,
10:00--11:00 a.m.,
IMPA - Ricardo Ma\~n\'e Auditorium.

\medskip\noindent
For the motivation of this work see
Deopark-Patel 15, a stratification of the
Harick space of degree $d$ covers
of $\mathbb{P}^1$. The question whether 
the strata are irreducible is addressed.

This talk defines the scroll invariants of a
degree-$d$ map $C \to \mathbb{P}^1$, defined
as the tuple $e_1,…,e_{d-1}$ such that
$f_*\mathcal{O}_C =
\mathcal{O}_{\mathbb{P}^1}\oplus
\mathcal{O}(e_1) \oplus … \oplus
\mathcal{O}(e_{d-1})$.

Later we consider a polarized toric surface
($S_{\Delta},L_\Delta)$ 
given by a (lattice?) polytope $\Delta$.
The scroll invariants of a curve
$C \in |L_\Delta|$ are determined by
the combinatorial data.


\section{Minimal Generators for Pinched
Veronese Varieties}
\label{16th-alga-meeting-2026-section-minimal-generators-veronese}

\medskip\noindent
{\bf Lecture.}

Aline Andrade,
UFMG.

\medskip\noindent
{\bf Schedule.}

Tuesday, September 1, 2026,
11:30 a.m.--12:30 p.m.,
IMPA - Ricardo Ma\~n\'e Auditorium.

\medskip\noindent
{\bf Abstract.}

No abstract is available on the official
site as of August 29, 2026.

\section{Rationally inequivalent points on
complete intersections}
\label{16th-alga-meeting-2026-section-inequivalent-points}

\medskip\noindent
{\bf Lecture.}

Olivier Martin,
IMPA.

\medskip\noindent
{\bf Schedule.}

Tuesday, September 1, 2026,
2:00--3:00 p.m.,
IMPA - Ricardo Ma\~n\'e Auditorium.

\medskip\noindent
{\bf Abstract.}

No abstract is available on the official
site as of August 29, 2026.

\section{Regular non-smooth quartic curves}
\label{16th-alga-meeting-2026-section-regular-nonsmooth-quartics}

\medskip\noindent
{\bf Presentation by Young Researcher
(PYR).}

Cesar Hilario,
PUC DEL PERU.

\medskip\noindent
{\bf Schedule.}

Tuesday, September 1, 2026,
PYR 1 session, 3:00--4:00 p.m.,
IMPA - Ricardo Ma\~n\'e Auditorium.

An individual presentation time within
the session is not listed on the
official site as of August 29, 2026.

\medskip\noindent
{\bf Abstract.}

No abstract is available on the official
site as of August 29, 2026.

\section{Generalized symmetry in the
vanishing of Ext}
\label{16th-alga-meeting-2026-section-generalized-symmetry-ext}

\medskip\noindent
{\bf Presentation by Young Researcher
(PYR).}

Paulo Martins,
ICMC USP.

\medskip\noindent
{\bf Schedule.}

Tuesday, September 1, 2026,
PYR 1 session, 3:00--4:00 p.m.,
IMPA - Ricardo Ma\~n\'e Auditorium.

An individual presentation time within
the session is not listed on the
official site as of August 29, 2026.

\medskip\noindent
{\bf Abstract.}

No abstract is available on the official
site as of August 29, 2026.

\section{Some Homological Conjectures over
Idealization Rings}
\label{16th-alga-meeting-2026-section-homological-conjectures-rings}

\medskip\noindent
{\bf Presentation by Young Researcher
(PYR).}

Igor Nascimento,
UFPR -- Jandaia do Sul.

\medskip\noindent
{\bf Schedule.}

Tuesday, September 1, 2026,
PYR 1 session, 3:00--4:00 p.m.,
IMPA - Ricardo Ma\~n\'e Auditorium.

An individual presentation time within
the session is not listed on the
official site as of August 29, 2026.

\medskip\noindent
{\bf Abstract.}

No abstract is available on the official
site as of August 29, 2026.

\section{On automorphisms of hypersurfaces
in Fano and Calabi Yau varieties}
\label{16th-alga-meeting-2026-section-automorphisms-hypersurfaces}

\medskip\noindent
{\bf Presentation by Young Researcher
(PYR).}

Ana Quedo,
Universitr Degli Studi Di Ferrara.

\medskip\noindent
{\bf Schedule.}

Tuesday, September 1, 2026,
PYR 1 session, 3:00--4:00 p.m.,
IMPA - Ricardo Ma\~n\'e Auditorium.

An individual presentation time within
the session is not listed on the
official site as of August 29, 2026.

\medskip\noindent
It looks like Quedo is working on a
generalization of the previous result
obtained with Paiva and Zikas,
now looking at the automorphism
group of a variety $F$ contained in 
Quedo defines an invariant called
$t_1(R(F,H))$, which plays an essential role
in her theorem:

\begin{theorem}[Quedo]
\label{16th-alga-meeting-2026-theorem-quedo}
Let $F$ be a smooth simply connected
projective variety with $\dim F\geq 3$,
and let $H$ be a normally generated line
bundle on $F$. Assume that either
\begin{enumerate}
\item $\omega_F\cong\mathcal{O}_F$; or
\item $F$ is a Fano variety and
\[
\omega_F\cong H^{\otimes r}
\]
for some integer $r<0$.
\end{enumerate}
Let
\[
X\in\lvert H^{\otimes d}\rvert
\]
be a smooth member. If
\[
d>t_1(R(F,H))
\qquad\text{and}\qquad
d\neq-r,
\]
then
\[
\operatorname{Aut}(X)
\cong\operatorname{Aut}(F;X)
\cong\operatorname{Aut}(\mathbb{P}^N;X),
\]
where
\[
i_H:F\longrightarrow\mathbb{P}^N
\]
is the embedding induced by the complete
linear system $\lvert H\rvert$.
\end{theorem}

\begin{proof}
Some ideas are based on work
by Matsumuri-Muki.
\end{proof}

As an application,

\begin{lemma}[Corollary]
\label{16th-alga-meeting-2026-lemma-quedo}
Let $F$ be a smooth Fano threefold such
that $-K_F$ is very ample, and suppose
that $F$ is not a quartic hypersurface
in $\mathbb{P}^4$. Let
\[
X\in\lvert-dK_F\rvert
\]
be a smooth member. If $d>3$, then
\[
\operatorname{Aut}(X)
\cong\operatorname{Aut}(F;X)
\cong\operatorname{Aut}(\mathbb{P}^N;X),
\]
where
\[
i:F\longrightarrow\mathbb{P}^N
\]
is the anticanonical embedding.
\end{lemma}

\begin{proof}
By \cite{CPS07},
\[
t_1(R(F,-K_F))\leq 3,
\]
so the result follows from
Theorem~\ref{16th-alga-meeting-2026-theorem-quedo}.

That is, the proof connects the theorem with
an already-known result.
\end{proof}

\section{Negative K-theory and Hodge theory}
\label{16th-alga-meeting-2026-section-negative-k-theory}

\medskip\noindent
{\bf Lecture.}

Mihnea Popa,
Harvard University, USA.

\medskip\noindent
{\bf Schedule.}

Tuesday, September 1, 2026,
4:30--5:30 p.m.,
IMPA - Ricardo Ma\~n\'e Auditorium.

\medskip\noindent
Joint with A. Burke and R. Shen.

\medskip\noindent
$X$ complex variety gives a sequence
of abelian groups indexed by the integers,
$K_i(X)$, $i \in \mathbb{Z}$.

The Grothendieck group of vector bundles on
$X$, $K_0(X)$ is an example. In an ad-hoc
way, there's also $K_1(X)$, which roughly
is the automorphisms classes of vector
bundles.

Bass defined {\it negative $K$-theory}.
$K_{-1}$ solves a lifting problem
for vector bundle classes on the
intersection
$X \times \mathbb{C}^*$. Then we define
the $K_{-i}$. The existence of certain exact
sequece is the fundamental theorem
of negative $K$-theory.

\begin{theorem}[Grothendieck]
\label{16th-alga-meeting-2026-theorem-grothendieck-k}
If $X$ is smooth, then $K_j(X) =0 $ for
$j<0$.
\end{theorem}

\noindent
Hence the negative $K$-groups are invariants
of singularities!

For example, for curves, it's a free group
generated on the branches of a singularity,
hence:
$$
K_{-1}(\operatorname{node})= \mathbb{Z},
\qquad 
K_{-1}(\operatorname{cusp}) = 0.
$$

Crucial question: are $K$-groups
homotopy invariant? e.g. If you take
product with affine space, do they remain
the same.

\begin{theorem}[Quillen]
\label{16th-alga-meeting-2026-theorem-quillen-k}
If $X$ is smooth, then $K_i(X) \simeq K_i(X
\times \mathbb{A}^1)$ for all $i$.
\end{theorem}

\noindent
Bloch-Murty in 79 show this fails in the
singular case. Much later Weibel showed that
$K_{-2}(X) \simeq K_{-2}(X \times
\mathbb{A}^1)$.

\medskip\noindent
If $K_i(X) \simeq K_i(X \times \mathbb{A}^r)$
for all $r \geq  1$ we say $X$ is
$K_i$-regular.

All we know on this is Weibel's conjecture.
[…]

\begin{theorem}[Weibel, 01]
\label{16th-alga-meeting-2026-theorem-weibel01}
LEt $X$ be quasi projective surface. There is an
exact sequence from $\operatorname{Pic}$ to
$\operatorname{Cl}$ ($X$ must be singular)
with cokernel $K_{-1}(X)$ [incomplete…]
\end{theorem}

\medskip\noindent
Weibel introduced homotopy $K$-theory.

\medskip\noindent
{\bf Understand $K$-regularity of $X$
in terms of higher singularities.}
Du Bois singularities have two main
properties:

\begin{itemize}
\item The natural map $H^i(X,\mathbb{C}) \to
H^i(X,\mathcal{O}_X)$ is a surjection for
all $i$.

\item Log canonical singularities are Du
Boid (Kollár-Kovács).
\end{itemize}

\noindent
Du Bois singularities $\implies$
is $K_{-n+1}$-regular. (It was already
$K_n$-regular, so Du Bois allows one step
further.)

Du Bois singularities connect the
commutative algebra of $X$ (the Kähler
differentials from Hartshorne)
with the Hodge theory.

\begin{definition}
\label{16th-alga-meeting-2026-definition-du-bois}
$X$ has {\it $m$-Du Boid singularities} 
if the natural map $\Omega_X^p \to
\underline{\Omega}_X^p$ is an isomorphism
for $p \leq m$.
\end{definition}

\begin{theorem}[Mustata-Olano-Popa-Witaszek,
Jung-Kim-Saito-Yoon]
\label{16th-alga-meeting-2026-theorem-mopwjksy}
[Du Bois singularity iff minimal exponent
$\tilde{\alpha}(f)$ condition.]
\end{theorem}

[Properties of minimal exponent.]

\begin{theorem}[Shen 25]
\label{16th-alga-meeting-2026-theorem-shen25}
$X$ complex variety has $m$-Du Bois
singularities and $s=\dim
X_{\operatorname{sing}}$, then $X$ is
$K_{\operatorname{number}}$-regular.
Converse if $X$ is affine LCI.
\end{theorem}

\noindent
Bloch-Murty example provides counterexample.

\medskip\noindent
{\bf Main conjecture.}
Let $i >0 $ be positive integer. If
$W_{j-i}H^j(X,\mathbb{Q})=0$ for all
$j \leq  2n-i$ [the weights are sufficiently
high] then $KH_{-i}(X)_{\mathbb{Q}}=0$
[the homotopy $K$-group is torsion].

In particular, if all $H^j(X,\mathbb{Q})$
are pure Hodge structures, then
$KH_{-i}(X)_{\mathbb{Q}}=0$ for all $i>0$.

\noindent
Example. rational homology manifolds (e.g.
quotient singularities).

\noindent
Example. $X$ threefold with rational
singularities. Then all Hodge structures of
$X$ are pure iff global
$\mathbb{Q}$-factoriality defect of $X$ is
equal to the sum of the local analytic
$\mathbb{Q}$-factoriality defects (Park-P,
25).

\begin{remark}
\label{16th-alga-meeting-2026-remark-motivic-conjecture}
The motivation of the conjecture is related
to motivic cohomology (which in turn is
related to cycles).
\end{remark}

\begin{definition}[Technical]
\label{16th-alga-meeting-2026-definition-technical}
$X$ has {\it $m$-rational singularities}
if the composition 
$$
\Omega_X^p \to \underline{\Omega}_X^p
\to \mathbb{D}(\underline{\Omega}_X^{n-p}
$$
is an isomorphism…
[The map from Kähler differential to Hodge
theoretic forms satisfies certain duality.]
\end{definition}

\noindent
Properties:
\begin{itemize}
\item 
\item 

\item $m$-rational singularities iff
$m$-Du Bois.
 \end{itemize}

\begin{theorem}[Park-P. '24]
\label{16th-alga-meeting-2026-theorem-pp24}
[Criterion for bounding weights in the Hodge
structure from below.]
\end{theorem}

\noindent
Hodge-Du Bois numbers/diamond. The speaker
and collaborators were after this, i.e.
theorem above was originally phrased in
terms of symmetries of such diamond.

\noindent
{\bf Conjecture.} Let $X$ be complex variety
with $m$-rational singularities and $m \leq
\frac{n-2}{2}$ with isolated singularities
(for simplicity).
Then 
$$
K_{-i}(X) =0 
$$
for $i>$ (involves $m$).

\medskip\noindent
Partial results joint with Burke and Shen 
(roughly):

\begin{itemize}
\item Holds for surfaces
\item Holds for certain $n$ if assuming
klt singularities.
\item Results for $\mathbb{Z}$-coefficients.
\item Holds on certain classes of varieties:
toric, varieties with quotient
singularities, cones over smooth varieties.
\item They follow from Bloch-Beilinson
conjectural filtrationon Chow groups, which
follows from the standard conjectures on
mixed motives.
\end{itemize}

\noindent
These conjectures are completely out of
reach at the moment!

\medskip\noindent
{\bf The reason for all this.}
Eigenspaces of some actions - motivic
cohomology.
In smooth case, this is the same as Chow
groups, or higher Chow groups.
Complexes of Chow groups, apply BB
machinery… but that's how Deligne
constructed the mixed Hodge structures!
You basically get to the weight filtration.
That's why the conjecture makes sense.


\medskip\noindent
{\bf Techniques.}

\begin{theorem}
\label{16th-alga-meeting-2026-theorem-chow-alpha}
$X$ threefold, $f:Y \to X$ resolution of
singularities, then the restriction map on
codimension-two cycles
$$
\alpha_{\mathbb{Q}}:
CH^2(Y)_{\mathbb{Q}}
\to
\bigoplus_{j}CH^2(S_j)_{\mathbb{Q}}
$$
is surjective.
\end{theorem}

\medskip\noindent
{\bf Extra.} The speaker mentioned
some computations that worked out well
on high-dimensional Kummer varieties
(quotients of abelian varieties).




\section{New Results on Linkage and Licci
Ideals}
\label{16th-alga-meeting-2026-section-linkage-licci-ideals}

\medskip\noindent
{\bf Lecture.}

Claudia Polini,
University of Notre Dame, USA.

\medskip\noindent
{\bf Schedule.}

Wednesday, September 2, 2026,
9:00--10:00 a.m.,
IMPA - Ricardo Ma\~n\'e Auditorium.

\medskip\noindent
{\bf History of linkage.} 
Idea: put a curve inside a complete
intersection
and compute residuals (Cayley, etc.)

\medskip\noindent
{\bf Modern formulation.}
Let $(R,\mathfrak{m})$ be a Gorenstein ring
and let $I,J \subseteq R$.

\begin{definition}
\label{16th-alga-meeting-2026-definition-direct-link}
$I$ and $J$ are {\it linked}, $I \sim J$,
if there exists a complete intersection
$\mathfrak{a}=(\alpha_1,…,\alpha_g)$
so that
$$
J= \mathfrak{a}:I,
\qquad 
I=\mathfrak{a}:J.
$$
\end{definition}

The idea is that the complete intersecion
$\mathfrak{a}$ is the union of 
$I$ and $J$ (the corresponding varieties
to these ideals, of course).

The idea is to look for an equivalence
relation - but $\sim$ is not initially a
symmetric relation. But iteration
gives the {\it linkage classes}:
$$
\mathcal{L}(I)=\{J:I \sim I_1 \sim…
\sim I_k=J\}
$$
An ideal is {\it licci} if its
linkage class contains a complete
intersection.

The hope: every link moves us toward
something simpler.

Example: we can link, after one iteration,
the twisted cubic to a line.

\begin{theorem}[Peskine-Szpiro]
\label{16th-alga-meeting-2026-theorem-ps}
If $I$ and $J$ are linked by $\mathfrak{a}$
in a Gorenstein ring,then
$R/I$ is CM iff $R/J$ is CM.
Moreove, $\omega_{R/I} \cong J/\mathfrak{a}$
and $\omega_{R/J} \cong I/\mathfrak{a}$.
\end{theorem}

\noindent
That is, morally, if two ideals are
linked then their invariants are the same.

This leads to the definition of
{\it deviation} of an ideal, which measures
how far is the ideal from a complete
intersection.

Any ideal linked to a Gorenstein is either
complete intersection or almost complete
intersection.

\medskip\noindent
The case of codimension 2 is nice:

\begin{theorem}[Apéry-Gaeta-Peskine-Szpiro]
\label{16th-alga-meeting-2026-theorem-agps}
If $R$ is Gorenstein and
$\operatorname{ht}I=2$ then $I$ is licci iff
$I$ is perfect.
\end{theorem}

\noindent
For codimension 3:

\begin{itemize}
\item Watanabe: $R=1$ (Gorenstein) or
$d \leq 1$ implies licci.

\item Guerreri-Ni-Weyman: if $\mathbb{Q}
\subset R$ then $I$ is licci when…
\end{itemize}

\noindent
Other results: Let $I$ be perfect
of height $g$, with
$r=\operatorname{type}(R/I)$ and 
$d = \mu(I) -g$.
\begin{itemize}
\item Herzog-Miller; Vasconcelos-Villareal:
if $g = 4$ and $I$ is generically a complete
intersection, then $I$ is licci when ($R/I$
is Gorenstein and $d \leq 2$) or
($r \leq 2$ …

\item 

\item 
\end{itemize}

\medskip\noindent
{\bf Objective.} Find a necessary condition
for an ideal to be licci, what properties
are preserved linkage?

Also we want sufficient conditions: what
forces an ideal to be licci?

\medskip\noindent
{\bf Licci ideals leave footprints
in their resolutions}
Morally, if $I$ is licci certain invariant
must grow rapidly:

\begin{theorem}[Huneke-Ulrich]
\label{16th-alga-meeting-2026-theorem-hu}
If $I_{\mathfrak{m}}$ is licci, then
$$
D_g > (g-1)d_1.
$$
Equivalently,
$$
\operatorname{reg}(S/I)>
(g-1)(\operatorname{indeg}I-1).
$$
\end{theorem}

\medskip\noindent
{\bf Conjecture (Huneke-Polini-Ulrich).}
If $I$ is homogeneous and licci, then 
$\mu(I) \leq D_g$. 

[Partial results explained.]

Let $S = k[x_1,…,x_n]$ and let $I \subset S$
be a square-free monomial licci ideal.
The conjecture would imply:

\begin{lemma}[Corollary of the conjecture]
\label{16th-alga-meeting-2026-lemma-coro-conj1}
$I$ satisfies $G_{\infty}$,
that is
$$
\mu(I_P)\leq  \dim S_p
$$
for every $P \supset I$.
\end{lemma}

\begin{lemma}[Corollary of the conjecture]
\label{16th-alga-meeting-2026-lemma-coro-conj2}
$I$ is of linear type, that is
its Rees algebra is naturally isomorphic
to its symmetric algebra:
$$
\mathcal{R}(I) = S[It]
\simeq \operatorname{Sym} (?)…
$$
\end{lemma}

\medskip\noindent
{\bf A sufficient criterion when $I$
contains a maximal regular sequence
of quadrics.}

\begin{theorem}[Huneke-Poini-Ulrich]
\label{16th-alga-meeting-2026-theorem-hpu}

\end{theorem}

\medskip\noindent
{\bf A new class of licci ideals.}

\begin{theorem}[Polini-Ulrich]
\label{16th-alga-meeting-2026-theorem-pu}

\end{theorem}




\section{Moduli of parabolic bundles on an
elliptic curve}
\label{16th-alga-meeting-2026-section-parabolic-bundles-elliptic}

\medskip\noindent
{\bf Presentation by Young Researcher
(PYR).}

Roberto Alvarenga,
UNESP.

\medskip\noindent
{\bf Schedule.}

Wednesday, September 2, 2026,
PYR 2 session, 10:00--11:15 a.m.,
IMPA - Ricardo Ma\~n\'e Auditorium.

An individual presentation time within
the session is not listed on the
official site as of August 29, 2026.

\medskip\noindent
{\bf Abstract.}

No abstract is available on the official
site as of August 29, 2026.

\section{Algebraic Geometry Codes on
Surfaces}
\label{16th-alga-meeting-2026-section-algebraic-geometry-codes}

\medskip\noindent
{\bf Presentation by Young Researcher
(PYR).}

Luana Costa,
IMPA.

\medskip\noindent
{\bf Schedule.}

Wednesday, September 2, 2026,
PYR 2 session, 10:00--11:15 a.m.,
IMPA - Ricardo Ma\~n\'e Auditorium.

An individual presentation time within
the session is not listed on the
official site as of August 29, 2026.

\medskip\noindent
{\bf Abstract.}

No abstract is available on the official
site as of August 29, 2026.

\section{The minimal free resolution of
full Nagata algebras}
\label{16th-alga-meeting-2026-section-minimal-resolution-nagata}

\medskip\noindent
{\bf Presentation by Young Researcher
(PYR).}

Adina Veronica Remor,
UFPR.

\medskip\noindent
{\bf Schedule.}

Wednesday, September 2, 2026,
PYR 2 session, 10:00--11:15 a.m.,
IMPA - Ricardo Ma\~n\'e Auditorium.

An individual presentation time within
the session is not listed on the
official site as of August 29, 2026.

\medskip\noindent
{\bf Abstract.}

No abstract is available on the official
site as of August 29, 2026.

\section{Some contributions to the theory
of quasi-homological dimensions}
\label{16th-alga-meeting-2026-section-quasi-homological-dimensions}

\medskip\noindent
{\bf Presentation by Young Researcher
(PYR).}

Victor Daniel Mendoza Rubio,
USP-S\~ao Carlos.

\medskip\noindent
{\bf Schedule.}

Wednesday, September 2, 2026,
PYR 2 session, 10:00--11:15 a.m.,
IMPA - Ricardo Ma\~n\'e Auditorium.

An individual presentation time within
the session is not listed on the
official site as of August 29, 2026.

\medskip\noindent
{\bf Abstract.}

No abstract is available on the official
site as of August 29, 2026.

\section{Rational Points in Weighted
Projective Spaces over Finite Fields}
\label{16th-alga-meeting-2026-section-rational-points-weighted}

\medskip\noindent
{\bf Presentation by Young Researcher
(PYR).}

Sajad Salami,
UERJ.

\medskip\noindent
{\bf Schedule.}

Wednesday, September 2, 2026,
PYR 2 session, 10:00--11:15 a.m.,
IMPA - Ricardo Ma\~n\'e Auditorium.

An individual presentation time within
the session is not listed on the
official site as of August 29, 2026.

\medskip\noindent
{\bf Abstract.}

No abstract is available on the official
site as of August 29, 2026.

\section{Homaloidal hypersurfaces and a
product formula for gradient maps}
\label{16th-alga-meeting-2026-section-homaloidal-hypersurfaces}

\medskip\noindent
{\bf Lecture.}

Nivaldo Medeiros,
UFF.

\medskip\noindent
{\bf Schedule.}

Wednesday, September 2, 2026,
11:45 a.m.--12:45 p.m.,
IMPA - Ricardo Ma\~n\'e Auditorium.

\medskip\noindent
{\bf Abstract.}

No abstract is available on the official
site as of August 29, 2026.

\section{Minimal-Degree Foliations on
Cominuscule Grassmannians}
\label{16th-alga-meeting-2026-section-minimal-degree-foliations}

\medskip\noindent
{\bf Lecture.}

Alan Muniz,
UFPE.

\medskip\noindent
{\bf Schedule.}

Thursday, September 3, 2026,
9:00--10:00 a.m.,
IMPA - Ricardo Ma\~n\'e Auditorium.

\medskip\noindent
{\bf Upshot.}
The main result is that every
minimal-degree codimension-one foliation
on a cominuscule Grassmannian is a pencil
of hyperplane sections.
The proof uses representation theory and
Bott--Borel--Weil to compute the minimal
twists $l_X(p)$ at which differential
$p$-forms appear.
In higher codimension, the authors
construct minimal-degree foliations from
flags and, on the Cayley plane, from
pencils of octonionic lines.


\medskip\noindent
\begin{definition}
\label{16th-alga-meeting-2026-definition-distribution}
Let $X$ be a smooth complex projective
variety.
A {\it codimension $q$ distribution}
$\mathcal{F}$ on $X$ is given by a
saturated subsheaf
$$
T_{\mathcal{F}}\subset TX
$$
of corank $q$, called the tangent sheaf of
$\mathcal{F}$.

$\mathcal{F}$ is a foliation if
$$
[T_{\mathcal{F}},T_{\mathcal{F}}]
\subset T_{\mathcal{F}}.
$$

$T\mathcal{F}$ is also given by the kernel
of a contraction map against an
integrable twisted $q$-form.
$$
\omega\in H^0(\Omega_X^q\otimes L)
$$
$$
T\mathcal{F}\hookrightarrow TX
\xrightarrow{\;\lrcorner\,\omega\;}
\Omega_X^{q-1}\otimes L
$$
\end{definition}

\begin{definition}
\label{16th-alga-meeting-2026-definition-space-of-foliations}
Then we define the space of foliations
$$
\operatorname{Fol}^q(X,L)
=
\left\{
\begin{array}{c}
[\omega]\in
\mathbb{P}(H^0(\Omega_X^q(L)))
\text{ integrable and not}\\
\text{vanishing in codimension one}
\end{array}
\right\}.
$$
\end{definition}

\medskip\noindent
{\bf The general problem:}
Given $X$, $q$, and $L$, determine the
space of foliations
$\operatorname{Fol}^q(X,L)$.

For $X=\mathbb{P}^n$,
$L=\mathcal{O}_{\mathbb{P}^n}(d+q+1)$,
where $d$ is the degree.
For $d=0$,
$$
\operatorname{Fol}^q
(\mathbb{P}^n,
\mathcal{O}_{\mathbb{P}^n}(q+1))
\simeq G(q+1,n+1).
$$

Every foliation is induced by a linear
projection
(Cerveau and Deserti - 2005).

For $d=1$, there are two irreducible
components of logarithmic type:
$$
\operatorname{Fol}^q
(\mathbb{P}^n,
\mathcal{O}_{\mathbb{P}^n}(q+2))
=
\operatorname{Log}^q(1^q,2)
\cup
\operatorname{Log}^q(1^{q+2}).
$$

(Jouanolou - 1979, for $q=1$, and
Loray-Pereira-Touzet - 2013, for
$q\geq 2$)

For $d=2$,
$\operatorname{Fol}^1
(\mathbb{P}^n,
\mathcal{O}_{\mathbb{P}^n}(4))$
has 6 irreducible components
(Cerveau-Lins Neto - 1996).
For $q\geq 2$ there are partial results
(Pereira-dos Santos - 2025) and
(Corrêa-M. - 2023).

\medskip\noindent
{\bf What about other $X$?}

Consider
$i\colon X\hookrightarrow\mathbb{P}^N$
an embedding.
Then we get restriction maps
$$
i_p^*\colon
H^0(\Omega_{\mathbb{P}^N}^q(l))
\longrightarrow
H^0(\Omega_X^q(l))
$$
and may ask whether the foliations on $X$
restrictions of foliations on
$\mathbb{P}^N$.

For $q=1$ and $l=0$ (i.e. $d=0$),
Araujo, Corrêa and Massarenti (2017)
obtained
$$
\operatorname{Fol}^1
(X,\mathcal{O}_X(2))
\simeq
\operatorname{Fol}^1
(\mathbb{P}^N,
\mathcal{O}_{\mathbb{P}^N}(2))
$$
for $X$ a smooth weighted complete
intersection and for $X=G(2,5)$.

\medskip\noindent
The result by Araujo-Corrêa-Massarenti is
obtained by killing cohomology, but
couldn't prove it for other Grassmanians.
Faenzi said this could be an interesting
thread.
Let, say, $X=G(k,V)$ 
which embeds equivariantly (the
Plucker embedding): $G(k,V)\hookrightarrow
\mathbb{P}(\Lambda^kV)$,
just like the map $i:X \hookrightarrow
\mathbb{P}^N$ above.
Perhaps this can be thought of as
representations; in the case of 
Grassmanian the group is
$\operatorname{GL}(n)$. Note that
the representations might have different
weights.




\begin{definition}
\label{16th-alga-meeting-2026-definition-cominuscule-grassmannian}
Consider
$$
X=G/P
$$
a rational homogeneous variety.
\begin{itemize}
\item $X$ is called {\it cominuscule} if
$T_X$ is completely reducible;

\item $X$ is called
{\it generalized Grassmannian} if
$$
\operatorname{Pic}(X)\simeq\mathbb{Z}.
$$
\end{itemize}

$X$ is called a
{\it cominuscule Grassmannian} if it is
both.
\end{definition}

They are the following:
\begin{enumerate}
\item classical Grassmannians
$$
G(k,n)=\operatorname{SL}(n)/P_k;
$$

\item Quadrics
$$
Q^n=\operatorname{SO}(n+2)/P_1,
\qquad n\geq 3;
$$

\item Spinor varieties
$$
OG(n,2n)=\operatorname{SO}(2n)/P_n;
$$

\item Lagrangian Grassmannians
$$
LG(n,2n);
$$

\item The Cayley plane
$$
E_6/P_1;
$$

\item The Freudenthal variety
$$
E_7/P_7.
$$
\end{enumerate}


The idea is: $X = G/P$ where $P$ is a
parabolic subgroup. There is a one-to-one
correspondence between the set of bundles
over $X$ and representations of $P$.
Again, the representations carry weights:
$\lambda$ weight $\to$ $E_\lambda$
irreducible vector bundle.


$p \in X$ fixed by $P$, then $P$ acts on
$T_xX = \bigoplus_\lambda V_\lambda$.

$P$ can be written as a matrix
$$
\begin{pmatrix}A&X\\ 0&B\end{pmatrix}
\in \operatorname{SL}(V).
$$
where $B \in \operatorname{SL}_{n-k}$.




\begin{theorem}[Benedetti, Faenzi,
M. - 2023]
\label{16th-alga-meeting-2026-theorem-minimal-degree-foliations}
Let
$$
X\subset\mathbb{P}(V)
$$
be one of the following (primitively
embedded) cominuscule Grassmannians:
\begin{enumerate}
\item a Grassmannian of lines, $G(2,n)$;

\item a spinor variety $OG(n,2n)$, for
$n=4,5$;

\item the Cayley plane, $E_6/P_1$;

\item one of the following:
$LG(3,6)$, $G(3,6)$, $OG(6,12)$ or
$E_7/P_7$.
\end{enumerate}

Then, the minimal-degree codimension-one
foliations ($q=1$ and
$L=\mathcal{O}_X(2)$) are pencils of
hyperplane sections.
\end{theorem}

\medskip\noindent
{\bf Question:}
What happens to the other cominuscule
varieties?



\medskip\noindent
The theorem was obtained by looking at the
representation theory, and the [?]
used in Araujo-Corrêa-Massarenti applies
here.

The idea one immediately thinks is that
using other machinery other than just
representation theory we can obtain other
results.

\medskip\noindent
{\bf General Theory:}
Fix $X$ such that
$\operatorname{Pic}(X)=\mathbb{Z}$, and
$\mathcal{F}$ given by
$$
\omega\in H^0(\Omega_X^1(2)).
$$

\begin{itemize}
\item (Araujo-Druel - 2013) If
$c_1(T\mathcal{F})>0$ then either
$T\mathcal{F}$ semistable and
$\mathcal{F}$ is algebraically integrable,
or there is a subfoliation
$\mathcal{G}\subset\mathcal{F}$ such that
$$
c_1(T\mathcal{G})
\geq c_1(T\mathcal{F}).
$$

\item (LPT - 2013) If $X$ is simply
connected, and $T\mathcal{F}$ semistable,
then $\mathcal{F}$ is a pencil of
hyperplane sections.
\end{itemize}

Then one may try to prove that there is no
such $\mathcal{G}$.
It is sufficient to show that
$$
H^0(\Omega_X^p(2))=0
$$
for $1<p<\dim X$.

So we started our journey in search of
$$
l_X(p)=
\min\left\{
l\in\mathbb{Z}
\mid
h^0(\Omega_X^p(l))>0
\right\}.
$$

\medskip\noindent
{\bf Case $G(k,n)$}

$$
G(k,n)=\operatorname{SL}_n/P_k
$$
where $P_k$ is the subgroup of matrices
$$
\begin{bmatrix}
A&*\\
0&B
\end{bmatrix}
$$
for $A\in\operatorname{SL}_k$ and
$B\in\operatorname{SL}_{n-k}$.
So
$$
\text{Vector bundles on }G(k,n)
\Longleftrightarrow
\text{Representations of }P_k
$$

Each representation of $P_k$ corresponds
to a weight vector
$$
\lambda=
(a_1,\ldots,a_k;
a_{k+1},\ldots,a_n)
$$
with
$$
a_1\geq a_2\geq\cdots\geq a_k,
\quad\text{and}\quad
a_{k+1}\geq\cdots\geq a_n.
$$

\begin{itemize}
\item Operations on vector bundles
correspond to operations on
representations, hence on weights
(plethysm).

\item Bott-Borel-Weil implies
$$
H^0(E_\lambda)\neq 0
$$
if and only if
$$
a_k\geq a_{k+1}.
$$
\end{itemize}

$$
0\longrightarrow U
\longrightarrow
\mathcal{O}_{G(k,n)}^{\oplus n}
\longrightarrow Q
\longrightarrow 0
$$

\begin{itemize}
\item $\mathcal{O}_{G(k,n)}(1)$
corresponds to
$$
(1,\ldots,1;0,\ldots,0);
$$

\item $U^\vee$ corresponds to
$$
(1,\ldots,0);
$$

\item $Q$ corresponds to
$$
(0,\ldots,0;1,0,\ldots,0);
$$

\item
$$
TX=Q\otimes U^\vee
$$
corresponds to
$$
(1,\ldots,0;0,\ldots,0,-1)
$$

\item
$$
\Omega^1_{G(k,n)}=U\otimes Q^\vee
$$
corresponds to
$$
(0,\ldots,-1;1,0,\ldots,0);
$$
\end{itemize}

$$
\Omega^p_{G(k,n)}(l)
=
\bigoplus_{\mu\vdash p}
\Gamma^\mu U\otimes
\Gamma^{\mu'}Q^\vee\otimes
\mathcal{O}_{G(k,n)}(l)
$$
$$
\rightsquigarrow
\bigoplus_{\mu\vdash p}
(l-\mu_k,\ldots,l-\mu_1;
\mu'_1,\ldots,\mu'_{n-k})
$$

Then we have
$$
l_{G(k,n)}(p)
=
\min\left\{
l\in\mathbb{Z}
\mid
h^0(\Omega^p_{G(k,n)}(l))>0
\right\}
$$
$$
=
\min\left\{
\mu_1+\mu'_1
\mid
\mu\vdash p
\right\}.
$$

That amounts to minimizing the perimeter
of a Young diagram contained in a
$(n-k)\times k$ rectangle of fixed area
$p$.
Therefore,
$$
l_{G(k,n)}(p)
=
\begin{cases}
k+\left\lceil\dfrac{p}{k}\right\rceil,
&
k^2<p<k(n-k),
\\[6pt]
\left\lceil2\sqrt{p}\right\rceil,
&
p\leq k^2.
\end{cases}
$$

We obtain that $l_X(p)\geq 3$ for
$p\geq 2$ and $X$ any cominuscule
Grassmannian.

\begin{theorem}
Let $X$ be a cominuscule Grassmannian and
let $\mathcal{F}$ be a codimension-one
foliation of minimal degree on $X$.
Then $\mathcal{F}$ is given by a pencil
of hyperplane sections.
In particular, set-theoretically,
$$
\operatorname{Fol}(X,1,2)
\simeq
G(2,H^0(\mathcal{O}_X(1))).
$$
\end{theorem}

We also obtain that $l_X(p)\geq 4$ for
$p\geq 3$ and $X$ any cominuscule
Grassmannian, except $IG(n,2n)$.

\begin{theorem}
Let $X$ be a cominuscule Grassmannian,
except $IG(n,2n)$, $n\geq 2$.
Let $\mathcal{F}$ be a codimension-one
foliation on $X$ with
$L=\mathcal{O}_X(3)$.
Then either
\begin{enumerate}
\item $\mathcal{F}$ is logarithmic of type
$(1,2)$, i.e., there exist
$$
f^2,g\in H^0(\mathcal{O}_X(2))
$$
such that $\mathcal{F}$ is given by the
fibers of
$$
(f^2:g)\colon
X\dashrightarrow\mathbb{P}^1;
$$
or

\item
$\mathcal{F}=\phi^*\mathcal{G}$ is the
pullback of a foliation $\mathcal{G}$ on a
surface $S$.
Moreover, the fibers of $\phi$ are
rationally connected and define a
minimal-degree foliation with semistable
tangent sheaf.
\end{enumerate}
\end{theorem}


For $X=IG(n,2n)$ we have:

\begin{theorem}
For every $n\geq 2$,
$$
\operatorname{Fol}^1
\bigl(
IG(n,2n),
\mathcal{O}_{IG(n,2n)}(3)
\bigr)
$$
has at least three irreducible components.
\end{theorem}

\medskip\noindent
{\bf Idea:}
$$
\operatorname{Fol}^1
\bigl(
IG(n,2n),
\mathcal{O}_{IG(n,2n)}(3)
\bigr)
\supset
\operatorname{Log}^1(1,2)
\cup
\operatorname{Log}^1(1,1,1).
$$

For $n=2$, $IG(2,4)\simeq Q^3$ a
quadric 3-fold.
By a result of Loray, Pereira, and Touzet,
$$
\operatorname{Fol}^1
\bigl(Q^3,\mathcal{O}_{Q^3}(3)\bigr)
=
\operatorname{Log}^1(1,2)
\cup
\operatorname{Log}^1(1,1,1)
\cup
\operatorname{Aff}.
$$

For $n\geq 3$ there are ``good'' rational
maps
$$
\phi\colon IG(n,2n)\dashrightarrow Q^3.
$$
We can show that for
$\mathcal{G}\in\operatorname{Aff}$,
$\phi^*\mathcal{G}$ is not contained in a
logarithmic component.

\begin{remark}
For $\mathbb{P}^n$ we have
$l(p)=p+1$.
But for $G(k,n)$, $l(p)$ can be lower.
So there are forms that do not extend to
the ambient space.
\end{remark}

Then we started looking for
higher-codimensional foliations of
minimal degree.

\begin{theorem}
Let $p=de$ for two integers $d,e$ such
that
$$
1\leq e\leq k,
\qquad
1\leq d\leq h:=n-k.
$$
Let $V$ be an $n$-dimensional complex
vector space.
Then
$$
\operatorname{Flag}(h-d,h+e,V)
\subset
\operatorname{Fol}^{de}
\bigl(G(k,V),d+e\bigr).
$$

More explicitly, given
$$
[W_{h-d}\subset W_{h+e}\subset V]
\in
\operatorname{Flag}(h-d,h+e,V)
$$
one defines a foliation $\mathcal{F}$ by
the fibers of a rational map
$$
\varphi\colon
G(k,V)\dashrightarrow
G(e,W_{h+e}/W_{h-d}).
$$
The tangent sheaf $T\mathcal{F}$ fits in
the following exact sequence
$$
0\longrightarrow
W_{h-d}\otimes U^\vee
\longrightarrow
T\mathcal{F}
\longrightarrow
(V/W_{h+e})^\vee
\otimes
\bigl(V/(U+W_{h-d})\bigr)^{\vee\vee}
\longrightarrow 0.
$$

In particular, $\mathcal{F}$ is of
minimal degree when $l(p)=d+e$.
\end{theorem}

The Cayley plane
$\mathbb{OP}^2\simeq E_6/P_1$
has the following table.

\begin{center}
\begin{tabular}{c c c c | c}
\hline
$p\backslash i$
& $1$ & $2$ & $3$ & $l(p)$\\
\hline
$1$
& $-2\lambda_1+\lambda_3$
& -- & -- & $2$\\

$2$
& $-3\lambda_1+\lambda_4$
& -- & -- & $3$\\

$3$
& $-4\lambda_1+\lambda_2+\lambda_5$
& -- & -- & $4$\\

$4$
& $-5\lambda_1+2\lambda_2+\lambda_6$
& $-5\lambda_1+2\lambda_5$
& -- & $5$\\

$5$
& $-6\lambda_1+3\lambda_2$
& $-6\lambda_1+\lambda_2+
\lambda_5+\lambda_6$
& -- & $6$\\

$6$
& $-7\lambda_1+2\lambda_2+\lambda_5$
& $-7\lambda_1+\lambda_4+2\lambda_6$
& -- & $7$\\

$7$
& $-8\lambda_1+\lambda_2+
\lambda_4+\lambda_6$
& $-8\lambda_1+\lambda_3+3\lambda_6$
& -- & $8$\\

$8$
& $-9\lambda_1+2\lambda_4$
& $-9\lambda_1+\lambda_2+
\lambda_3+6\lambda_6$
& $-8\lambda_1+4\lambda_6$
& $8$\\

$9$
& $-10\lambda_1+\lambda_3+
\lambda_4+\lambda_6$
& $-9\lambda_1+\lambda_2+3\lambda_6$
& -- & $9$\\

$10$
& $-10\lambda_1+\lambda_4+2\lambda_6$
& $-11\lambda_1+2\lambda_3+\lambda_5$
& -- & $10$\\

$11$
& $-12\lambda_1+3\lambda_3$
& $-11\lambda_1+\lambda_3+
\lambda_5+\lambda_6$
& -- & $11$\\

$12$
& $-12\lambda_1+2\lambda_3+\lambda_6$
& $-11\lambda_1+2\lambda_5$
& -- & $11$\\

$13$
& $-12\lambda_1+\lambda_3+\lambda_5$
& -- & -- & $12$\\

$14$
& $-12\lambda_1+\lambda_4$
& -- & -- & $12$\\

$15$
& $-12\lambda_1+\lambda_2$
& -- & -- & $12$\\
\hline
\end{tabular}%
\end{center}

\medskip\noindent
{\bf TABLE 2.}
Weights $\beta_{p,i}$ decomposing
$\Omega_X^p$, for $X$ the Cayley plane,
and
$$
l(p)=
\min\left\{
l\in\mathbb{Z}
\mid
H^0(\Omega_X^p(l))\neq 0
\right\}.
$$

For every $p\in\mathbb{OP}^2$, take
$\mathcal{F}_p$ the pencil of octonionic
lines $\mathbb{OP}^1$ passing through
$p$.

\begin{theorem}
The foliation $\mathcal{F}_p$ is given by
an integrable form in
$$
H^0\bigl(\Omega_{\mathbb{OP}^2}^8(8)\bigr).
$$
\end{theorem}

\section{Stability conditions for coherent
systems on integral curves}
\label{16th-alga-meeting-2026-section-stability-coherent-systems}

\medskip\noindent
{\bf Lecture.}

Renato Vidal Martins,
UFMG.

\medskip\noindent
{\bf Schedule.}

Thursday, September 3, 2026,
10:00--11:00 a.m.,
IMPA - Ricardo Ma\~n\'e Auditorium.

\medskip\noindent
{\bf Objective.} Develop new
stability conditions for coherent sheaves
(other than the standard ones).

\medskip\noindent
{\bf Outline}

\begin{enumerate}
\item Categorical Notion of Stability
\begin{itemize}
\item Stability Conditions on Abelian
Categories
\item Pre-Stability Conditions (PSC) and
(Bridgeland) Stability Conditions (SC) on
Triangulated Categories
\end{itemize}

\item Stability Conditions for Coherent
Systems
\begin{itemize}
\item Coherent Systems
\item Le Potier and Clifford Type Theorems
\item Construction of a three-parameter
family of PSC's
\item Bogomolov-Gieseker Inequality to get
SC's
\end{itemize}

\item Space of Stability Conditions
\begin{itemize}
\item Main Results
\item Applications
\end{itemize}
\end{enumerate}

joint work with Marcos Jardim and Leonardo
Roa-Leguizamon

\medskip\noindent
{\bf Stability Functions}

$\mathcal{A}:=$ abelian category with
Grothendieck group $K_0(\mathcal{A})$,

$Z:K_0(\mathcal{A})\to\mathbb{C}$ an
additive homomorphism.

$Z$ is a {\bf stability function} if, for
all $0\neq E\in\mathcal{A}$,
$$
\Im(Z(E))\geq 0
\quad\text{and}\quad
\Re(Z(E))<0
\text{ if }\Im(Z(E))=0.
$$

The slope of $E\in\mathcal{A}$ is
$$
M(E):=
\begin{cases}
-\dfrac{\Re(Z(E))}{\Im(Z(E))}
&\text{if }\Im(Z(E))\neq 0,\\[6pt]
\infty
&\text{if }\Im(Z(E))=0.
\end{cases}
$$

$0\neq E\in\mathcal{A}$ is {\bf stable}
(resp. {\bf semistable}) if, for all
proper $0\neq F\subset E$,
$$
M(F)<M(E)
\quad
\text{(resp. $\leq$).}
$$

\medskip\noindent
{\bf (Bridgeland) Stability Conditions}

$\mathcal{D}:=$ triangulated category,
$\Lambda=$ finite rank lattice,
$\nu:K_0(\mathcal{D})\to\Lambda$

A {\bf pre-stability condition} on
$\mathcal{D}$ is a pair
$(\mathcal{A},Z)$ where

$\mathcal{A}$ is the heart of a bounded
$t$-structure on $\mathcal{D}$

$Z:K_0(\mathcal{A})\to\mathbb{C}$ is a
homomorphism, which factors into
$$
\begin{array}{ccc}
K_0(\mathcal{A})
&\xrightarrow{\ Z\ }&
\mathbb{C}
\\
&\searrow_{\nu}\quad\nearrow&
\\[-6pt]
&\Lambda&
\end{array}
$$
such that $(\mathcal{A},Z)$ is a
stability condition

$(\mathcal{A},Z)$ is a
{\bf (Bridgeland) stability condition} if
the support property holds:
$$
\inf
\left\{
\frac{|Z(E)|}{\|\nu(E)\|}
\ \middle|\
0\neq E\in\mathcal{A}\text{ semistable}
\right\}
>0.
$$

\medskip\noindent
{\bf Coherent Systems on Integral Curves}

$C=$ integral and complete curve over an
a.c. field

A {\bf coherent system} of type $(n,d,k)$
on $C$ is a triple
$\mathcal{E}=(E,V,\varphi)$ where

$E=$ coherent sheaf of rank $n$ and degree
$d$ on $C$

$V=$ $k$-dimensional vector space

$\varphi:V\to H^0(E)$ morphism of vector
spaces

A morphism
$(E,V,\varphi)\to(E',V',\varphi')$ is:
$\alpha:E\to E'$, $\beta:V\to V'$,
$$
\begin{array}{ccc}
V&\xrightarrow{\ \varphi\ }&H^0(E)
\\
{\scriptstyle\beta}\downarrow
&&
\downarrow{\scriptstyle H^0(\alpha)}
\\
V'&\xrightarrow{\ \varphi'\ }&H^0(E')
\end{array}
$$

$\operatorname{CohSys}(C):=$ category of
coherent systems on $C$

\medskip\noindent
{\bf Example 2}

$\mathcal{A}:=\operatorname{CohSys}(C)$

$\mathcal{E}=$ coherent system of type
$(n,d,k)$ on $C$. For $\alpha>0$, set
$$
Z_\alpha(\mathcal{E})
=-(d+\alpha k)+\sqrt{-1}\,n
$$

The $\alpha$-slope of $\mathcal{E}$ is
$$
\mu_\alpha(\mathcal{E}):=
\begin{cases}
\dfrac{d+\alpha k}{n}
&\text{if }n>0,\\[6pt]
\infty
&\text{if }n=0.
\end{cases}
$$

$\operatorname{CohSys}(C)$ is the heart
of the $t$-structure of
$D^b(\operatorname{CohSys}(C))$
$$
\frac{|Z_\alpha(\mathcal{E})|^2}
{\|\nu(\mathcal{E})\|^2}
=
\frac{(d+\alpha k)^2+n^2}
{n^2+d^2+k^2}
\geq
\begin{cases}
\alpha^2&\text{if }0<\alpha\leq 1,\\
1&\text{if }\alpha>1.
\end{cases}
$$

(we needed $\alpha$-semistability)

\medskip\noindent
{\bf Le Potier and Clifford Type Theorems}

$E\in\operatorname{Coh}(C)$ of rank $n$
and degree $d$, and $\mu(E):=d/n$

\medskip\noindent
{\bf Le Potier:}
If $C$ is smooth, $E$ semistable and
$d\geq 0$, then
$h^0(E)\leq d+n$

\medskip\noindent
{\bf Thm 1:}
If $\mu_{\min}(E)\geq 0$, then
$h^0(E)\leq d+n$

\medskip\noindent
{\bf Clifford (G. Xiao):}

$C$ smooth, $E$ semistable bundle,
$0\leq\mu(E)\leq 2g-2$, then
$h^0(E)\leq d/2+n$

\medskip\noindent
{\bf Thm 2:}
If
$0\leq\mu_{\min}(E)\leq\mu_{\max}(E)
\leq 2g-2$, then
$h^0(E)\leq d/2+n$

\medskip\noindent
$\mathcal{E}=(E,V,\varphi)
\in\operatorname{CohSys}(C)$ of type
$(n,d,k)$, and
$\mu_\alpha(\mathcal{E})
:=(d+\alpha k)/n$

\medskip\noindent
{\bf LN}
$C$ smooth, $\mathcal{E}$
$\alpha$-semistable, $E$ bundle,
$\mu_\alpha(\mathcal{E})\geq 0$, then
$k\leq d+n$

\medskip\noindent
{\bf Thm 3:}
If $E$ is pure and
$\mu_{\alpha\mathrm{min}}(\mathcal{E})
\geq 0$, then
$k\leq d+n$

\medskip\noindent
{\bf LN}
$C$ smooth, $\mathcal{E}$
$\alpha$-semistable, $E$ bundle,
$0\leq\mu_\alpha\mathcal{E}\leq 2g-2
\Rightarrow k\leq d/2+n$

\medskip\noindent
{\bf Thm 4:}
If $E$ is pure,
$0\leq
\mu_{\alpha\mathrm{min}}(\mathcal{E})
\leq
\mu_{\alpha\mathrm{max}}(\mathcal{E})
\leq 2g-2$, then
$k\leq d/2+n$

\medskip\noindent
{\bf 3-Parameter Family of Conditions on
the Derived Category}

\medskip\noindent
{\bf Part 1:} building categories in
$D^b(\operatorname{CohSys}(C))$

For $\alpha,\beta\in\mathbb{R}$, with
$\alpha\geq 0$, set:
$$
\mathcal{T}_\alpha^\beta
:=
\left\{
\mathcal{E}\in\operatorname{CohSys}(C)
\ \middle|\
\mu_\alpha(\mathcal{G})>\beta
\text{ whenever }
\mathcal{E}\to\!\!\!\!\!\to\mathcal{G}
\right\}
$$
$$
\mathcal{F}_\alpha^\beta
:=
\left\{
\mathcal{E}\in\operatorname{CohSys}(C)
\ \middle|\
\mu_\alpha(\mathcal{F})\leq\beta
\text{ for every }
0\neq\mathcal{F}\subset\mathcal{E}
\right\}
$$
$$
\operatorname{CohSys}_\alpha^\beta(C)
:=
\left\{
\mathcal{E}\in
D^b(\operatorname{CohSys}(C))
\ \middle|\
\begin{array}{l}
\mathcal{H}^i(\mathcal{E})=0
\text{ for all }i\neq 0,-1,\\
\mathcal{H}^0(\mathcal{E})
\in\mathcal{T}_\alpha^\beta
\text{ and }
\mathcal{H}^{-1}(\mathcal{E})
\in\mathcal{F}_\alpha^\beta
\end{array}
\right\}.
$$

Note. I had to leave the talk at about
half an hour after it had started.


\section{Abelian varieties: Subvarieties and
Decompositions}
\label{16th-alga-meeting-2026-section-abelian-varieties-subvarieties-decompositions}

\medskip\noindent
{\bf Lecture.}

Anita Rojas,
Universidad de Chile, Chile.

\medskip\noindent
{\bf Schedule.}

Thursday, September 3, 2026,
11:30 a.m.--12:30 p.m.,
IMPA - Ricardo Ma\~n\'e Auditorium.


\begin{definition}
\label{16th-alga-meeting-2026-definition-abelian}
A $g$-dimensional {\it complex torus} is
$V/\Lambda$ where $V \cong \mathbb{C}^g$
and $\Lambda \simeq \mathbb{Z}^{2g}$.
\end{definition}

\noindent
A complex torus then can be embedded
in $\mathbb{P}^N$ is called an abelian
variety.

\begin{definition}
\label{16th-alga-meeting-2026-definition-polarization}
A {\it polarization} on $T=V/\Lambda$
is $H:V \times V \to \mathbb{C}$ hermitian,
$H(v,v)>0$ and
$J n H|_{\Lambda \times \Lambda}
\subseteq \mathbb{Z}$.
This is the same as $E:V \times V \to \mathbb{R}$
real alternating form such that
$E(iu,iv) = E(u,v)$,
$E(iu,v)>0$, $E|_{\Lambda \times \Lambda}
\subseteq \mathbb{Z}$.

The {\it type of the polarization}
is $(d_1,…,d_g)$ such that $d_j \in
\mathbb{Z}_{>0}$, $d_i|d_{i+1}$.
If the $d_j=1$ for all $j$, 
$(T,V/\Lambda,E_0)$ is a
{\it principally polarized abelian variety}
\end{definition}

\noindent
Taking coordinates,
$$
\begin{aligned}
\exists \beta=:  &\longrightarrow  \\
 &\longmapsto 
\end{aligned}
$$

Period matrix is (D,Z).

\medskip\noindent
Let $A=(T=V/\Lambda,E_0)$
and $f:A \to A'$ a homomorphism.
Taking coordinates
$\exists \alpha \in M(g'\circ g
\mathbb{C})$,
$R \in M(g,.?, \mathbb{Z})$.
$$
\alpha \Pi_A
=
\Pi_{A'}R
$$
If $f^*(E')=E$, then
$R \in \operatorname{Sp}^D(2g,\mathbb{Z})$.


$$
\operatorname{Sp}^D(2g,\mathbb{Z})
=\{M JM=J\},
$$
$$
J=\begin{pmatrix}0&D\\ D&0\end{pmatrix}
$$

$$
\xymatrix{
\mathbb{H}_g
\ar[d]\\
\mathcal{A}_g
}
$$



\medskip\noindent

\begin{theorem}[Poincaré complete
irreducibility]
\label{16th-alga-meeting-2026-theorem-poincare-irreduc}
$A=(T=V/\Lambda,E_0)$ $\implies $
$\exists X_1,…,X_s$ abelian
subbundle… up to conjugacy of
\begin{enumerate}
\item $\Hom(X_i,X_j) = \{0\}$.
\item 
\end{enumerate}
\end{theorem}

\begin{definition}
\label{16th-alga-meeting-2026-definition-completely-decomposable}
$A$ is completely decomposable if
$A \sim E_1\times… \times E_g$.
\end{definition}

\begin{itemize}
\item $X$ simple ? is said of {\it CM-type} 
if
$$
\xymatrix{
\operatorname{End}_R(A) \ar[d]\\
X\ar[d]^g\\
\mathbb{R}
}
$$
\item  $A$ is {\it not simple} if
it is of CM-type if its simple factos
(Poincaré decomposition) are.
\end{itemize}

\medskip\noindent
\begin{itemize}
\item '93 Ekedahl-Serre:
are there completely decomposable Jacobian
varieties in every dimension?
Is there a bound for the $g$ of such
objects? $g = 1297$.

\item $\approx$ 87 (Coleman). Conjectured
that there are only finitely many Jacobians
of $CM$-type for type $g$. This is false for
$g \leq 7$, and the proof of this fact
is done with groups actions, i.e. construct
Jacobian with given automorphism group.
\end{itemize}

\medskip\noindent
{\bf Our approach.}
$$
\varepsilon:G \to \operatorname{Aut}(A)
$$

$$
\tilde{\varepsilon}:
\mathbb{Q}[G] \to
\operatorname{End}_{\mathbb{Q}}(A)
$$

$$
\mathbb{Q}[G] = Q_1 \times … \times Q_r
$$
for $r= \#\operatorname{Im}G/$.

$Q_j$ is generated by a antal idempotent
$\implies A_j := \operatorname{Im}(e_j)$
and $A_1 \times … A_r \to A$,
isotypical.
$\Hom(A_j,A_i) \neq  \{ 0\}$,
$i \neq  j$,
$\Hom_G(A_j,A_i) = \{0\}$.

$\forall e_j \exists f_{1_j}… f_{n_j}$
$\implies$
$$
B_j^{n_1}\times … \times B_r^{n_r}
\to A
$$
in S.A.D. The $B_k$ are called {\it
primitive factors}.

\medskip\noindent
{\bf Plan.} Study
$J \in \mathcal{J}_g$ such that
$\operatorname{Aut}(J) \neq  \{\pm 1\}$.
\begin{itemize}
\item Perform G.A.D.
\item Decompose $B_j$ further.
\end{itemize}

$\operatorname{NS}_{\mathbb{Q}}(A)
=
\{H: V \times V \to \mathbb{C}:
H \operatorname{hermitian},
\operatorname{Im}H|_{\Lambda \times \Lambda}
\subseteq \mathbb{Z}\}$.
There is a correspondence between
$\operatorname{NS}_{\mathbb{Q}}$
and the endomorphisms symmetric w.r.t
Rosat involution
$$
\begin{aligned}
\varphi: \operatorname{NS}_{\mathbb{Q}}(A)
&\longrightarrow
\operatorname{End}_\mathbb{Q}^S(A) \\
H &\longmapsto 
X_{H^0}^{-1} X_H.
\end{aligned}
$$
\noindent
Recall that subvarieties of an abelian
variety are in correspondence with symmetric
idempotents.

Recall that by Lefschetz 1,1 theorem,
$$
\begin{aligned}
\mathcal{H}:
\operatorname{NS}_{\mathbb{Q}}(A)
&\longrightarrow H^{1,1}(A) \cap
H^2(X,\mathbb{Q}) \\
E &\longmapsto 
- \sum_{?j}\varepsilon?\wedge d? \wedge dg?.
\end{aligned}
$$
\noindent
If we have the period matrix of the
primitive factos we can find subvarities and
decompose them.

\medskip\noindent
We shall explain the method that allows us
to decompose full the primitive factors.

\medskip\noindent
{\bf Finding subvarieties}

Given $A=(V/\Lambda,E)$,
$\Pi=(D\ Z)$,
$D=\operatorname{diag}(d_1,\ldots,d_g)$.
If $x_1,\ldots,x_{2g}$ are the real
coordinate functions of
$\Lambda\otimes\mathbb{R}$ associated to
the basis of $\Lambda$ and $z_i$ the
complex coordinate functions with respect
to the basis of $V$. Then there is a
correspondence between:
\begin{enumerate}
\item abelian subvarieties of $A$ of
dimension $n$ and

\item differential forms
$\omega\in H^2(A,\mathbb{Q})$ such that
\begin{itemize}
\item
$$
\omega\wedge dz_1\wedge\cdots\wedge dz_g
=0
$$

\item
$$
\omega^{\wedge r}
\wedge\theta^{\wedge(g-r)}
=
\begin{cases}
\displaystyle
\left(
\prod_{j=1}^g d_j\,n!
\prod_{k=n+1}^g(k-r)
\right)\cdot\omega_0
&\text{if }1\leq r\leq n,\\[8pt]
0
&\text{if }n<r\leq g,
\end{cases}
$$
with
$$
\omega_0=(-1)^g
dx_1\wedge dx_{g+1}\wedge\cdots
\wedge dx_g\wedge dx_{2g}.
$$
\end{itemize}
\end{enumerate}

Robert Auffarth, Herbert Lange and
Anita M. Rojas.
{\it A criterion for an abelian variety
to be non-simple.}
J. Pure Appl. Algebra, 221 (2017),
no. 8, 1906--1925.

\medskip\noindent
{\bf Let $(A,E)$ be a polarized abelian
surface}

with period matrix
$\Pi_A=(D\ \tau)$.
Then, there is a quadric surface
$\mathcal{E}_{\Pi}\subset\mathbb{R}^3$
and a lifting
$$
\mathcal{L}\colon
\mathcal{E}_{\Pi}\longrightarrow
\mathbb{R}^5
$$
with
$$
\mathcal{E}_{\Pi}
=
\left\{
(x,y,z)\in\mathbb{R}^3
\ \middle|\
\begin{aligned}
Ax^2+By^2+Cz^2+Dyz+Exz+Fxy\\
{}+Gx+Hy+Jz=0
\end{aligned}
\right\}.
$$

% The displayed formulas defining
% A,\ldots,J and the two additional
% coordinates of the lifting are illegible
% in the supplied photograph.

such that
\begin{itemize}
\item the Riemann relations for $\Pi$
imply that $\mathcal{E}_{\Pi}$ is compact;

\item the elliptic curves contained in $A$
are in correspondence with the rational
points of $\mathcal{E}_{\Pi}$ with
rational lifting.
\end{itemize}

\medskip\noindent
{\bf Experimental approach}

Ekedahl--Serre, Coleman, Beauville:
construct explicit families of Jacobians
whose period matrices we can write down,
and then decompose further. Describing the
factors, their endomorphism algebras, etc.

Three ideas:
\begin{enumerate}
\item a method for $G$-invariant
subvarieties (period matrix),

\item a method for images of endomorphisms
(period matrix),

\item the criterion for finding abelian
subvarieties.
\end{enumerate}

Application: A genus 101 curve with
completely decomposable Jacobian.

R. E. Rodr\'iguez, A. M. Rojas.
{\it Period matrices for abelian
varieties: Algorithms and applications.}
Journal of Algebra 683 (2025), 364--388.

\medskip\noindent
{\bf First method: a $G$-subvariety
$A_e=\operatorname{Im}e$}

Given a symplectic $G$-action on
$A=(V/\Lambda,E)$ and a central
idempotent $e\in\mathbb{Q}[G]$, put
$h'=\dim A_e$.

\begin{itemize}
\item {\bf Input:}
$$
\rho_r\colon
G\longrightarrow
\operatorname{Sp}^D
(2g\times 2g,\mathbb{Z})
$$
and $\rho_r(e)$

\item {\bf Output:}
\begin{enumerate}
\item A basis $\beta_{A_e}$ of
$\Lambda_{A_e}$,
$\rho_r(i_{A_e})$ for the embedding
$i_{A_e}\colon A_e\to A$, and
$D_{A_e}$ the type of the induced
polarization on $A_e$.

\item The rational representation
$$
\rho_{r,A_e}\colon
G\longrightarrow
\operatorname{Sp}^{D_{A_e}}
(2h'\times 2h',\mathbb{Z})
$$
of $G$ in the basis $\beta_{A_e}$, where
$\rho_{r,A_e}(g)$ is the unique matrix in
$\operatorname{Sp}^{D_{A_e}}
(2h'\times 2h',\mathbb{Z})$ satisfying
\begin{equation}
\label{16th-alga-meeting-2026-equation-restricted-representation}
\rho_r(g)\cdot\rho_r(i_{A_e})
=
\rho_r(i_{A_e})\cdot\rho_{r,A_e}(g).
\tag{1}
\end{equation}
\end{enumerate}
\end{itemize}

\medskip\noindent
{\bf Application: $g=101$}

$$
G=\operatorname{SmallGroup}(800,980)
\simeq
(C_5\times C_5)
\operatorname{semidirectproduct}
\bigl((C_8\times
C_2)\operatorname{semidirectproduct} C_2\bigr)
$$
acts on a curve $X_{101}$ of genus 101
with signature $(0;8,8,2)$. A GAD for its
Jacobian $JX_{101}$ is nice:
\begin{equation}
\label{16th-alga-meeting-2026-equation-gad-genus-101}
E\times S\times E_0^2\times E_1^8
\times\cdots\times E_{12}^8
\longrightarrow JX_{101}.
\tag{3}
\end{equation}

We find:
\begin{itemize}
\item $S$ is a $G$-invariant abelian
surface of type $(10,10)$.

\item The representation on $\Lambda_S$
for the generators of $G$ is
$$
\rho_{r,S}(a)
=
\begin{pmatrix}
0&1&-1&1\\
0&-1&0&-1\\
1&-1&1&-1\\
1&1&0&0
\end{pmatrix}
$$
and
$$
\rho_{r,S}(b)
=
\begin{pmatrix}
-1&0&-1&0\\
1&0&1&-1\\
1&1&0&0\\
-1&1&-1&1
\end{pmatrix}.
$$
\end{itemize}

\medskip\noindent
The Riemann matrix
$Z_S\in\mathbb{H}_2$ fixed by
$\rho_{r,S}(G)$:
$$
Z_S
=
\begin{pmatrix}
\dfrac{1+i\sqrt{2}}{2}&-\dfrac12\\[6pt]
-\dfrac12&\dfrac{1+i\sqrt{2}}{2}
\end{pmatrix}.
$$

Then
$$
\mathcal{E}_{\Pi}\colon
x^2+xy+y^2-xz+z^2+\frac12x+y=0,
$$
and we find rational points with rational
lifting:
$$
\left(-\frac12,-\frac3{10},\frac1{10}\right),
\qquad
\left(-\frac12,-\frac15,-\frac35\right),
\qquad
\left(-\frac12,0,-\frac12\right),
$$
among many.

For one of them, we get the isogeny
$$
s\colon
E_{1+i\sqrt{2}}
\times
E_{\frac{1+i\sqrt{2}}{3}}
\longrightarrow S.
$$
Therefore,
$$
\boxed{
\begin{gathered}
E\times E_{1+i\sqrt{2}}
\times E_{\frac{1+i\sqrt{2}}{3}}
\times E_0^2\times E_1^8
\times\cdots\times E_{12}^8\\
\longrightarrow JX_{101}.
\end{gathered}
}
$$


































































\section{Mustafin Varieties and Linked
Projective Spaces}
\label{16th-alga-meeting-2026-section-mustafin-varieties}

\medskip\noindent
{\bf Presentation by Young Researcher
(PYR).}

Felipe de Le\'on Saenz Angel,
IMPA.

\medskip\noindent
{\bf Schedule.}

Thursday, September 3, 2026,
PYR 3 session, 2:00--3:00 p.m.,
IMPA - Ricardo Ma\~n\'e Auditorium.

An individual presentation time within
the session is not listed on the
official site as of August 29, 2026.

\medskip\noindent
{\bf From lattice configurations to
representations}

\begin{enumerate}
\item $R$ is a DVR with fraction field
$K$, residue field $k$, and uniformizer
$\pi$.

\item For a lattice $\Lambda\subseteq K^d$,
let
$$
[\Lambda]:=
\{\alpha\Lambda\mid\alpha\in K^*\}.
$$

\item Let $\Gamma$ be a finite convex
collection of lattices in $K^d$, i.e.,
$$
[\pi^a\Lambda\cap\pi^b\Lambda']
\in\Gamma
$$
for all $[\Lambda],[\Lambda']\in\Gamma$
and $a,b\in\mathbb{Z}$.
\end{enumerate}

\medskip\noindent
{\bf The adjacency quiver and its
representation}

The vertices of $Q(\Gamma)$ are the
elements of $\Gamma$. There is an arrow
$u\to v$ if $u$ and $v$ are adjacent,
i.e., there exist representatives
$\Lambda_u\in u$ and $\Lambda_v\in v$
satisfying
$$
\pi\Lambda_u\subset\Lambda_v
\subset\Lambda_u.
$$

Set
$$
V(\Gamma)_v=\Lambda_v/\pi\Lambda_v.
$$
For each arrow $u\to v$,
$$
\varphi_{uv}
=
\overline{\pi^{m_{uv}}}\colon
V(\Gamma)_u\longrightarrow V(\Gamma)_v,
\qquad
m_{uv}
=
\min\{m\mid
\pi^m\Lambda_u\subseteq\Lambda_v\}.
$$

\medskip\noindent
{\bf Mustafin varieties}

Let
$$
\Gamma=
\{[\Lambda_1],\ldots,[\Lambda_n]\}
$$
be as above. Each $\Lambda_i$ defines an
$R$-model $\mathbb{P}(\Lambda_i)$ of
$\mathbb{P}_K^{d-1}$.

\medskip\noindent
{\bf Deligne scheme and Mustafin variety}

The Deligne scheme associated to $\Gamma$
is
$$
\mathcal{M}(\Gamma)
:=
\overline{\Delta(\mathbb{P}_K^{d-1})}
\subseteq
\mathbb{P}(\Lambda_1)
\times_R\cdots\times_R
\mathbb{P}(\Lambda_n).
$$

Its special fiber $\mathcal{M}(\Gamma)_k$
is the Mustafin variety.

The scheme $\mathcal{M}(\Gamma)$ is flat
over $R$, its special fiber is reduced,
and
$$
\boxed{
\mathcal{M}(\Gamma)_k
\cong
\operatorname{Gr}_{(1,\ldots,1)}
\bigl(V(\Gamma)_k\bigr)
}.
$$

Hahn, M.A., Li, B.
{\it Mustafin varieties, moduli spaces and
tropical geometry.}
manuscripta math. 166, 167--197 (2021).

https://doi.org/10.1007/s00229-020-01237-8

\medskip\noindent
{\bf A local model for
$\operatorname{LP}(\mathfrak{V})$ and
global questions}

\medskip\noindent
{\bf Theorem (Esteves, -)}

If $\mathfrak{V}$ is a nontrivial exact
finitely generated linked net, then every
point
$$
x\in\operatorname{LP}(\mathfrak{V})
$$
has an open neighborhood isomorphic to an
open subset of
$\mathcal{M}(\Gamma_x)_k$ for some convex
lattice configuration $\Gamma_x$.

\medskip\noindent
Can these local models be globalized?

\medskip\noindent
{\bf I. LP $\longrightarrow$ Mustafin}

Does there exist a single configuration
$\Gamma$ such that
$$
\operatorname{LP}(\mathfrak{V})
\cong\mathcal{M}(\Gamma)_k\ ?
$$

Unknown in general.

\medskip\noindent
{\bf II. Mustafin $\longrightarrow$ LP}

Given $\Gamma$, does there exist an exact
linked net $\mathfrak{W}$ such that
$$
\mathcal{M}(\Gamma)_k
\cong\operatorname{LP}(\mathfrak{W})\ ?
$$

Yes --- this is our direction.

\section{The tropical theta divisor}
\label{16th-alga-meeting-2026-section-tropical-theta-divisor}

\medskip\noindent
{\bf Presentation by Young Researcher
(PYR).}

Giusi Capobianco,
Universita Degli Studi Di Roma Tor Vergata.

\medskip\noindent
{\bf Schedule.}

Thursday, September 3, 2026,
PYR 3 session, 2:00--3:00 p.m.,
IMPA - Ricardo Ma\~n\'e Auditorium.

An individual presentation time within
the session is not listed on the
official site as of August 29, 2026.

\medskip\noindent
{\bf Abstract.}

No abstract is available on the official
site as of August 29, 2026.

\section{Non-algebraicity of foliations via
reduction modulo 2}
\label{16th-alga-meeting-2026-section-nonalgebraicity-foliations}

\medskip\noindent
{\bf Presentation by Young Researcher
(PYR).}

Jo\~ao Paulo Figueredo,
UFF.

\medskip\noindent
{\bf Schedule.}

Thursday, September 3, 2026,
PYR 3 session, 2:00--3:00 p.m.,
IMPA - Ricardo Ma\~n\'e Auditorium.

An individual presentation time within
the session is not listed on the
official site as of August 29, 2026.

\medskip\noindent
{\bf Introduction.}

Fixed $d\geq 0$, one can define the space
$$
\operatorname{Fol}_d(\mathbb{P}_{\mathbb{C}}^2)
=
\left\{
[v]\in
\mathbb{P}\left(
H^0\left(
\mathbb{P}_{\mathbb{C}}^2,
T_{\mathbb{P}_{\mathbb{C}}^2}^{1}(d-1)
\right)
\right)
\mid
\operatorname{codim}(Z(v))\geq 2
\right\}
$$
which parametrizes foliations of degree
$d$ on $\mathbb{P}_{\mathbb{C}}^2$.

\medskip\noindent
{\bf Theorem (Jouanolou 1979)}

The subset of
$\operatorname{Fol}_d(\mathbb{P}_{\mathbb{C}}^2)$
parameterizing foliations of degree
$d>1$ having no invariant curves is a
countable intersection of open subsets of
$\operatorname{Fol}_d(\mathbb{P}_{\mathbb{C}}^2)$,
which is dense.

\medskip\noindent
{\bf Introduction.}

In this talk, I will present the following
result:

\medskip\noindent
{\bf Theorem (-, Mendson, 2026)}

Let $\mathcal{F}$ be a foliation on
$\mathbb{P}_{\mathbb{C}}^2$ of degree
$d>1$ such that:

\medskip\noindent
(a) $\mathcal{F}$ is defined over a number
field;

\medskip\noindent
(b) $\mathcal{F}$ is not dicritical;

\medskip\noindent
(c) $\mathcal{F}$ has good reduction
modulo 2;

\medskip\noindent
(d) The reduction modulo 2 of
$\mathcal{F}$ has irreducible 2-divisor.

\medskip\noindent
Then $\mathcal{F}$ doesn't have an
invariant curve.

\section{Multiple Tor modules and
Mayer-Vietoris spectral sequences}
\label{16th-alga-meeting-2026-section-multiple-tor-modules}

\medskip\noindent
{\bf Presentation by Young Researcher
(PYR).}

Rafael Holanda,
UFPE.

\medskip\noindent
{\bf Schedule.}

Thursday, September 3, 2026,
PYR 3 session, 2:00--3:00 p.m.,
IMPA - Ricardo Ma\~n\'e Auditorium.

An individual presentation time within
the session is not listed on the
official site as of August 29, 2026.

\medskip\noindent
{\bf Standard grading versus standard
multigrading}

\begin{center}
\small
\begin{tabular}{
|p{0.30\textwidth}|
p{0.20\textwidth}|
p{0.27\textwidth}|}
\hline
$\mathbb{P}$
&
$\mathbb{P}^{r-1}$
&
$\mathbb{P}^{n_1-1}\times\cdots\times
\mathbb{P}^{n_r-1}$
\\ \hline
Sheaf cohomology
&
{\v C}ech wrt $\mathfrak{m}$
&
{\v C}ech wrt $B$
\\ \hline
Cohomology of $\mathbb{P}$
&
only 1 (dual to $\omega_R$)
&
many (1 dual to $\omega_R$)
\\ \hline
Compute cohom.\ via free resolutions
&
direct
&
spectral sequences
\\ \hline
Grothendieck-Serre formula
&
$H_{\mathfrak{m}}$
&
$H_B$
\\ \hline
Regularity $\times$ linear resolutions
&
yes
&
quasi-linear
\\ \hline
Regularity $\times$ Betti numbers
&
yes
&
no
\\ \hline
Regularity $\times$ Tor modules
&
yes
&
yes, Tor of $R/(X_i)$
\\ \hline
\end{tabular}
\end{center}

\medskip\noindent
{\bf Why no}

Given a homogeneous ideal $I$ of $R$,
$$
C_I(M)^*
:=
\bigcup_i
\operatorname{Supp}_{\mathbb{Z}^k}
\bigl(H_I^i(M)\bigr)^*
\qquad\text{and}\qquad
T_I(M)^*
:=
\bigcup_i
\operatorname{Supp}_{\mathbb{Z}^k}
\bigl(\operatorname{Tor}_i^R(M,R/I)\bigr)^*.
$$

\medskip\noindent
{\bf Theorem (Chardin-\_ '22)}

\begin{enumerate}
\item[(1)]
$$
C_{\mathfrak{m}}(M)^*
=
T_{\mathfrak{m}}(M)^*
-
\sum_i n_i e_i.
$$

\item[(2)]
$$
C_B(M)^*
=
\bigcup_i C_{B_i}(M)^*
=
\bigcup_{j_1<\cdots<j_p}
C_{B_{j_1}+\cdots+B_{j_p}}(M)^*.
$$
\end{enumerate}

\medskip\noindent
{\bf Tor modules}

$I_1,I_2,I_3$ ideals: the matrix
$$
\begin{pmatrix}
0&1&-1\\
-1&0&1\\
1&-1&0
\end{pmatrix}
$$
induces maps
$$
R/I_2I_3\oplus R/I_1I_3\oplus R/I_1I_2
\xrightarrow{\varphi}
R/I_1\oplus R/I_2\oplus R/I_3
$$
and
$$
R/I_1\oplus R/I_2\oplus R/I_3
\xrightarrow{\psi}
R/(I_2+I_3)\oplus R/(I_1+I_3)
\oplus R/(I_1+I_2).
$$

We can prove that
$$
\operatorname{coker}(\varphi)
\simeq
\operatorname{coker}(\psi)
\simeq
R/(I_1+I_2+I_3)
=
\operatorname{Tor}_0^R
(R/I_1,R/I_2,R/I_3).
$$
The matrix
$$
\begin{pmatrix}1&1&1\end{pmatrix}
$$
induces maps
$$
\ker(\psi)
\xleftarrow{\kappa_\psi}
R/I_1I_2I_3
\xrightarrow{\kappa_\varphi}
\ker(\varphi).
$$
Assuming
$$
I_i\otimes_R I_j
\simeq
I_iI_j
=
I_i\cap I_j
\quad\text{for }i\neq j
\quad
\bigl(
\operatorname{Tor}_2^R(R/I_i,R/I_j)
=
\operatorname{Tor}_1^R(R/I_i,R/I_j)
=
0
\bigr),
$$
$$
\operatorname{coker}(\kappa_\varphi)
\simeq
\operatorname{Tor}_1^R
(R/I_1,R/I_2,R/I_3)
\simeq
\operatorname{coker}(\kappa_\psi),
$$
$$
\ker(\kappa_\psi)
\simeq
\operatorname{Tor}_2^R
(R/I_1,R/I_2,R/I_3)
\simeq
\ker(\kappa_\varphi)
\simeq
I_1\cap I_2\cap I_3/I_1I_2I_3.
$$

\medskip\noindent
{\bf Multiple Tor modules}

$I_1,\ldots,I_n\subseteq R$ ideals,
$F_\bullet^i$ canonical free resolution
of $R/I_i$,
$$
F_\bullet
=
F_\bullet^1\otimes\cdots\otimes
F_\bullet^n.
$$
$$
\operatorname{Tor}_i^R
(R/I_1,\ldots,R/I_n)
:=
H_i(F_\bullet).
$$

$$
T_{p,q}
:=
\bigoplus_{j_1<\cdots<j_p}
\operatorname{Tor}_q^R
(R/I_{j_1},\ldots,R/I_{j_p})
$$
and
$$
\mathbb{H}_{p,q}
=
\bigoplus_{j_1<\cdots<j_p}
H_{q+1}
\left(
{}^+F_\bullet^{I_{j_1},\ldots,I_{j_p}}
[p-1]
\right).
$$

\medskip\noindent
{\bf Theorem (Banerjee-Chardin-\_ '26)}

{\it There are Mayer-Vietoris-type
spectral sequences}
$$
\textnormal{(i)}
\qquad
E_{n-p,q}^1
=
T_{p,q}
\Longrightarrow
\mathbb{H}_{n,q-p}.
$$
$$
\textnormal{(ii)}
\qquad
E_{p,q}^1
=
\mathbb{H}_{p,q}
\Longrightarrow
T_{n,p+q}.
$$

\medskip\noindent
{\bf Back to multigrading setting}

$$
C_I(M)^*
:=
\bigcup_i
\operatorname{Supp}_{\mathbb{Z}^k}
\bigl(H_I^i(M)\bigr)^*
\qquad\text{and}\qquad
T_I(M)^*
:=
\bigcup_i
\operatorname{Supp}_{\mathbb{Z}^k}
\bigl(\operatorname{Tor}_i^R(M,R/I)\bigr)^*.
$$

\medskip\noindent
{\bf Theorem (Banerjee-Chardin-\_ '26)}

\begin{enumerate}
\item[(i)]
$$
\bigcup_{j_1<\cdots<j_p}
T_{B_{j_1},\ldots,B_{j_p}}(M)
=
\bigcup_{j_1<\cdots<j_p}
T_{B_{j_1}+\cdots+B_{j_p}}(M).
$$

\item[(ii)]
$$
C_B(M)^*
=
\bigcup_{j_1<\cdots<j_p}
T_{B_{j_1}+\cdots+B_{j_p}}(M)^*
=
\bigcup_{j_1<\cdots<j_p}
T_{B_{j_1},\ldots,B_{j_p}}(M)^*
\supseteq
T_B(M)^*.
$$
\end{enumerate}


\section{Configurations of Skew Lines on
Cubic Surfaces}
\label{16th-alga-meeting-2026-section-skew-lines-cubic-surfaces}

\medskip\noindent
{\bf Presentation by Young Researcher
(PYR).}

W\'allace de Sousa,
UFPB.

\medskip\noindent
{\bf Schedule.}

Thursday, September 3, 2026,
PYR 3 session, 2:00--3:00 p.m.,
IMPA - Ricardo Ma\~n\'e Auditorium.

An individual presentation time within
the session is not listed on the
official site as of August 29, 2026.

\medskip\noindent
{\bf Abstract.}

No abstract is available on the official
site as of August 29, 2026.

\section{Generalized multiplicities and
criteria for integral dependence}
\label{16th-alga-meeting-2026-section-generalized-multiplicities}

\medskip\noindent
{\bf Lecture.}

Bernd Ulrich,
Purdue University, USA.

\medskip\noindent
{\bf Schedule.}

Friday, September 4, 2026,
10:00--11:00 a.m.,
IMPA - Ricardo Ma\~n\'e Auditorium.

\medskip\noindent
Joint work with Claudia Polini and Ngo Viet
Trung.

\medskip\noindent
{\bf Notes structure.} Note for this lecture
are written like this: first my personal
notes taken during the lecture,
then the AI-transcribed photos of the
slides.

\medskip\noindent
Somehow the parameter space $T$
is the singular locus of the family $Z$.
That's interesting. The fibers have
different kinds of singularities.
Central question: under which conditions
are the fibers singular in the same way,
i.e. (Whitney) equisingular.
The catch: we want to find this condition
in terms of the fibers, not of the family
itself.

\medskip\noindent
{\bf Definition of Whitney equisingularity.}
Condition A is not enough for
equisingularity, we use condition B to fix
this: not only the tangent space converges
but also the secant lines converge
(to something inside the tangent space at
the special fiber).

Condition A: ensure convergence,
condition W, the in which it convergence is
good.

\medskip\noindent
{\bf Teissier theorem.}
Note that the ideal in the theorem
only depends on the fiber, just as
we want in the main objective of this talk.

In trying to generalize the theory beyond
hypersurfaces we note that the Jacobian
ideal of higher-codimension subvarieties
does not carry enough information. Thus we
consider submodules of the Jacobian
ideal. That's what Gaffney's observation
says. Namely, that integral depdendence
of submodules of the Jacobian ideal
is equivalent to Whitney equisingularity.

\medskip\noindent
{\bf PSID.} Principle by Tessier:
integral dependence can be detected
by a numerical condition which depends only in
the fibers.

Tessier theorem in ithis slide is designed
to describe the special fiber of the family
from the generic fiber, then describe the
whole family and finally look at the special
fiber within the family.

Note that this only works for isolated
singularities in hypersurfaces. Such
restrictions motivate a generalization of
the notion of multiplicity.
Gaffney-Gassler already remove the
isolated singularity condition; and the
other teams proved in different cases for
isolated singularities on higher
codimension.

We see that there is not a single notion
that would work in all cases.

\medskip\noindent
{\bf Integral dependence, Hilbert-Samuel
multiplicity.}
The Rees algebra is a polynomial ring
in $It$---that's an easy way of remembering
the definition.


\medskip\noindent
{\bf Generalized multiplicities.}
Achilles-Manaresi's $j$-multiplicity is
motivated by intersection theory.

\medskip\noindent
{\bf $\varepsilon$-multiplicity.}
The GGPTUV theorem is a ``Rees theorem''.

\medskip\noindent
{\bf Proof.} Relate $j$-multiplicity
with $e(A)$.

\medskip\noindent
{\bf Now we start with the transcribed
slides.}

\medskip\noindent
{\bf Families of germs of hypersurfaces}

Let
$$
f\in\mathbb{C}\{x_1,\ldots,x_n,t\}
$$
be a germ at $0$ of a holomorphic
function (irreducible), $f(0)=0$.
$$
\begin{array}{c}
\mathbb{C}^n\times\mathbb{C}
\supset Z:\quad f=0
\\
\downarrow\,\pi
\\
0\times\mathbb{C}=T
\\
\Vert
\\
\operatorname{Sing}Z
\end{array}
\qquad
\text{germ at $0$ of the $t$-axis}.
$$
Thus we have the flat family
$$
\left\{
Z_\lambda=\pi^{-1}(\lambda)
\right\}.
$$

How can one decide whether the family is
(Whitney) equisingular?

\medskip\noindent
{\bf Whitney Conditions}

\medskip\noindent
{\bf Whitney condition A}

$Z$ satisfies Whitney's condition A at
$0$ along $T$ if for every sequence of
points
$$
z_i\in Z\setminus T
$$
so that $z_i$ converge to $0$ and the
tangent spaces $T_{Z,z_i}$ converge to a
space
$$
H=\lim T_{Z,z_i},
$$
one has
$$
T\subset H.
$$

``The limiting tangent space contains the
parameter axis''

\medskip\noindent
{\bf Whitney condition B}

$Z$ satisfies Whitney's condition B at
$0$ along $T$ if for every sequence of
points
$$
z_i\in Z\setminus T
$$
and every sequence of points $t_i\in T$
so that
$$
\lim z_i=0=\lim t_i,
$$
$H=\lim T_{Z,z_i}$ exists, and the secant
lines $\overline{z_it_i}$ converge to a
line
$$
\ell=\lim\overline{z_it_i},
$$
one has
$$
\ell\subset H.
$$

``The limiting tangent space contains the
limiting secant line is contained in the
limiting tangent space''

\medskip\noindent
``The tangent spaces converge to the
parameter axis''

\medskip\noindent
{\bf Condition W}

The tangent spaces converge to the
parameter axis at least as fast as the
points do.
$$
W\Longleftrightarrow B
$$
the family is WHITNEY EQUISINGULAR.

\medskip\noindent
{\bf Theorem [Teissier]}

Condition W is equivalent to the integral
dependence of
$$
\mathcal{I}
:=
(x_1,\ldots,x_n)
\left(
\frac{\partial f}{\partial x_1},
\ldots,
\frac{\partial f}{\partial x_n}
\right)
\mathcal{O}_{Z,0}
\subset
\mathcal{J}
:=
\mathcal{I}
+
\frac{\partial f}{\partial t}
\mathcal{O}_{Z,0}.
$$

Gaffney has extended this to the case
when
$$
Z:f_1=\cdots=f_e=0
$$
is not a hypersurface; but in this case
$$
\frac{\partial f}{\partial x_1},
\ldots,
\frac{\partial f}{\partial x_n},
\frac{\partial f}{\partial t}
$$
have to be replaced by the columns of the
Jacobian matrix of $f_1,\ldots,f_e$ and
the ideals
$$
\mathcal{I}\subset\mathcal{J}
$$
are replaced by submodules
$$
\mathcal{U}\subset\mathcal{E}
$$
of $(\mathcal{O}_{Z,0})^e$.

\medskip\noindent
{\bf PSID: Principle of Specialization of
Integral Dependence}

\medskip\noindent
{\bf Goal.}
Deduce the integral dependence of
$\mathcal{I}\subset\mathcal{J}$ on the
total space $Z$ of the family, hence the
Whitney equisingularity of the family,
from numerical fiber-wise conditions that
only depend on $Z_\lambda$, the members of
the family.

\medskip\noindent
{\bf Principle [Teissier]}

$$
\operatorname{mult}
(\mathcal{I}\mathcal{O}_{Z,\lambda})
\text{ constant for }\lambda\in T
\Longrightarrow
\text{integral dependence of }
\mathcal{I}\subset\mathcal{J}
\text{ on }Z.
$$

\medskip\noindent
{\bf Theorem [Teissier]}

The PSID holds for families of
hypersurfaces with isolated
singularities, with mult the
Hilbert-Samuel multiplicity.

The Hilbert-Samuel multiplicity is only
defined in this case, any generalization
requires generalized notions of
multiplicity. They should satisfy
\begin{itemize}
\item a Rees theorem
\item upper semicontinuity in families
\end{itemize}

The PSID has been proved for
\begin{itemize}
\item families of hypersurfaces
[Gaffney, Gassler] with
$\operatorname{mult}=$the multiplicity
sequence
\item families of isolated singularities
[Gaffney, Kleiman],
[Kleiman, U, Validashti],
[Rangachev]
\end{itemize}

\medskip\noindent
{\bf Integral dependence, Hilbert-Samuel
multiplicity}

Let $R$ be a Noetherian ring and $I$ an
ideal.
$$
R[t]\supset
\mathcal{R}(I)
=
R[It]
=
\bigoplus_{n\geq 0}I^nt^n,
$$
the Rees algebra of $I$.

Let $I\subset J$ be ideals. Then
$$
J\text{ integral over }I
\Longleftrightarrow
\mathcal{R}(I)\subset\mathcal{R}(J)
\text{ integral extension of rings}
$$
$$
\Longleftrightarrow
J^n=I^{n-r}J^r
\qquad
\forall n\geq r.
$$
Hence
$$
J^n\subset I^{n-r}\subset J^{n-r}
\qquad
\forall n\geq r.
$$

From now on:
$$
(R,\mathfrak{m},k)
\text{ Noetherian local},\qquad
\widehat{R}\text{ equidimensional},
\qquad
d=\dim R>0.
$$

\medskip\noindent
{\bf The case where $I$ is
$\mathfrak{m}$-primary:}

\medskip\noindent
{\bf Theorem [Rees 1961]}

$$
I\subset J\text{ integral}
\Longleftrightarrow
e(I)=e(J).
$$

The associated graded ring is
$$
G
=
\operatorname{gr}_I(R)
=
\mathcal{R}(I)\otimes_R R/I
\simeq
\bigoplus_{n\geq 0}I^n/I^{n+1}.
$$

The Hilbert-Samuel multiplicity is
$$
e(I)
=
e(\operatorname{gr}_I(R))
=
\lim_{n\to\infty}
\frac{(d-1)!}{n^{d-1}}
\lambda(I^n/I^{n+1})
\in\mathbb{Z}_{>0}
$$
$$
=
\lim_{n\to\infty}
\frac{d!}{n^d}\lambda(R/I^n).
$$

\medskip\noindent
{\bf Generalized multiplicities}

\medskip\noindent
{\bf The case where $I\subseteq R$ is
arbitrary:}

\begin{itemize}
\item
$j$-multiplicity
[Achilles-Manaresi, 1993]:
\end{itemize}
$$
j(I)
=
\lim_{n\to\infty}
\frac{(d-1)!}{n^
{d-1}}
\lambda\left(
H_{\mathfrak{m}}^0(I^n/I^{n+1})
\right)
=
\begin{cases}
e\left(H_{\mathfrak{m}}^0(G)\right),
&
\text{if }\dim G\otimes_R k=d,
\\
0,
&
\text{otherwise},
\end{cases}
$$
$$
j(I)\in\mathbb{Z}_{\geq 0}.
$$
Here
$$
\dim G\otimes_R k=\ell(I)
$$
is the analytic spread of $I$.

\medskip\noindent
{\bf Theorem [Flenner, Manaresi 2001]}

$$
I\subset J\text{ integral}
\Longleftrightarrow
j(I_p)=j(J_p)
\qquad
\forall p\in V(I).
$$

\begin{itemize}
\item
$\varepsilon$-multiplicity:
\end{itemize}
$$
\varepsilon(I)
=
\limsup_{n\to\infty}
\frac{d!}{n^d}
\lambda\left(
H_{\mathfrak{m}}^0(R/I^n)
\right)
\in\mathbb{R}_{\geq 0}.
$$

\begin{itemize}
\item
Multiplicity sequence
[Achilles, Manaresi 1997;
Gaffney, Gassler 1999]:
\end{itemize}
$$
\operatorname{gr}_{\mathfrak{m}}
\bigl(\operatorname{gr}_I(R)\bigr)[u,v],
$$
with $u,v$ variables of degrees
$(1,0),(0,1)$, is a standard bigraded
$k$-algebra of dimension $d+2$.
Thus it has a bivariate Hilbert polynomial
of the form
$$
\sum_{i=0}^d
\frac{c_i(I)}{(d-i)!\,i!}
r^{d-i}s^i
+
\text{lower terms},
\qquad
c_i(I)\in\mathbb{Z}_{\geq 0}.
$$
The sequence
$$
c_0(I),\ldots,c_d(I)
$$
is the multiplicity sequence of $I$.

\medskip\noindent
{\bf Fact.}
$$
c_d(I)=j(I).
$$

\medskip\noindent
{\bf Theorem [Gaffney, Gassler 1999;
Polini, Trung, U, Validashti 2020]}

$$
I\subset J\text{ integral}
\Longleftrightarrow
c_i(I)=c_i(J)
\qquad
\forall\operatorname{ht}I\leq i\leq d.
$$

\begin{itemize}
\item
Mixed multiplicities
[Bhattacharya 1955]:
\end{itemize}
$$
\operatorname{gr}_{\mathfrak{m}}
\bigl(\mathcal{R}(I)\bigr)
$$
is a standard bigraded $k$-algebra of
dimension $d+1$.
The normalized leading coefficients of
its bivariate Hilbert polynomial
$$
e_0,\ldots,e_d
$$
are the mixed multiplicities of $I$.

\medskip\noindent
{\bf The graded case}

Joint with Claudia Polini and Ngo Viet
Trung.

Inspired by [Das, Roy, Trivedi 2025].

\medskip\noindent
{\bf Theorem [Simis, U, Vasconcelos 2001]}

Let
$$
A\subset B
$$
be a homogeneous inclusion of standard
graded $k$-algebras such that
$$
\dim A=\dim B,
$$
$A$ is a domain, and $B$ is
equidimensional. Then
$$
\textnormal{(a)}
\qquad
\operatorname{rank}_A(B)\cdot e(A)
\leq e(B),
$$
and
$$
\textnormal{(b)}
\qquad
\operatorname{rank}_A(B)\cdot e(A)
=
e(B)
\Longleftrightarrow
A\subset B\text{ integral}.
$$

Contrary to (a): if $I\subset J$ are
$\mathfrak{m}$-primary ideals, then
$$
e(I)\geq e(J).
$$

\medskip\noindent
{\bf Theorem [Flenner, Manaresi 2001]}

$$
j(I_p)\leq j(J_p)
\qquad
\forall p\in V(I)
\Longleftrightarrow
I\subset J\text{ integral}.
$$

From now on:
$$
(R,\mathfrak{m})
\text{ standard graded $k$-algebra},
\qquad
\text{equidimensional},
\qquad
d=\dim R>0,
$$
and
$$
I\subset J
\text{ homogeneous ideals},
\qquad
I\not\subseteq\sqrt{0}.
$$

\medskip\noindent
{\bf The case where $I$ and $J$ are
generated in degree $\delta$:}

\medskip\noindent
{\bf Theorem [Polini, Trung, U 2025]}

Suppose that $R$ is a domain and
$$
\ell(I)=d.
$$
Then
$$
\textnormal{(a)}
\qquad
j(I)\leq j(J),
$$
and
$$
\textnormal{(b)}
\qquad
j(I)\geq j(J)
\Longleftrightarrow
I\subset J\text{ integral}.
$$

\medskip\noindent
{\bf Proof.}

Let
$$
A=k[I_\delta]
\subset
B=k[J_\delta]
\subset
C=k[\mathfrak{m}_\delta],
$$
all standard graded $k$-algebras of
dimension $d$.
By [Kustin, Polini, U 2016],
$$
j(I)
=
\delta\cdot
\operatorname{rank}_A(C)\cdot e(A)
$$
$$
=
\delta\cdot
\operatorname{rank}_B(C)\cdot
\operatorname{rank}_A(B)\cdot e(A)
$$
$$
\leq
\delta\cdot
\operatorname{rank}_B(C)\cdot e(B)
\qquad
\text{by [SUV 99]}
$$
$$
=
j(J)
\qquad
\text{by [KPU 16]}.
$$
Moreover,
$$
j(I)=j(J)
\Longleftrightarrow
\operatorname{rank}_A(B)\cdot e(A)=e(B)
\Longleftrightarrow
A\subset B\text{ integral}
\Longleftrightarrow
I\subset J\text{ integral}.
$$
\hfill$\square$

\medskip\noindent
{\bf Theorem [Polini, Trung, U 2025]}

The following are equivalent:
\begin{enumerate}
\item[(1)]
$I\subset J$ is integral;

\item[(2)]
$$
e(\mathcal{R}(I))
=
e(\mathcal{R}(J));
$$

\item[(3)]
$$
\sum e_i(I)
=
\sum e_i(J);
$$

\item[(4)]
$I_p\subset J_p$ is integral for every
$p\neq\mathfrak{m}$;

\item[(5)]
$$
j(I_p)\leq j(J_p)
\qquad
\forall p\neq\mathfrak{m}.
$$
\end{enumerate}

Proof of $(5)\Rightarrow(1)$:
$$
j(I)\leq j(J)
$$
holds by the previous Theorem.


\medskip\noindent
{\bf The case where $I$ and $J$ are
arbitrary homogeneous ideals:}

Let $I$ and $J$ be generated in degrees
$<\delta$.

We use an idea from
[Das, Roy, Trivedi 2025]:
$$
S=R[x],
\qquad
\mathcal{I}=(IS)_{\geq\delta}
\subset
\mathcal{J}=(JS)_{\geq\delta}.
$$
These ideals are generated by forms of
degree $\delta$, and
$$
\ell(\mathcal{I})=\dim S.
$$

\medskip\noindent
{\bf Theorem [Polini, Trung, U 2025]}

The following are equivalent:
\begin{enumerate}
\item[(1)]
$I\subset J$ is integral;

\item[(2)]
$\mathcal{I}\subset\mathcal{J}$ is
integral;

\item[(3)]
$$
e(\mathcal{R}(\mathcal{I}))
=
e(\mathcal{R}(\mathcal{J}));
$$

\item[(4)]
$$
\sum e_i(\mathcal{I})
=
\sum e_i(\mathcal{J});
$$

\item[(5)]
$$
j(\mathcal{I})=j(\mathcal{J}).
$$
\end{enumerate}

$(1)\Longleftrightarrow(5)$ had been
proved in [Das, Roy, Trivedi 2025].

\section{Contractions via moduli of sheaves}
\label{16th-alga-meeting-2026-section-contractions-moduli-sheaves}

\medskip\noindent
{\bf Lecture.}

Svetlana Makarova,
Australian National University,
Australia.

\medskip\noindent
{\bf Schedule.}

Friday, September 4, 2026,
11:30 a.m.--12:30 p.m.,
IMPA - Ricardo Ma\~n\'e Auditorium.

\medskip\noindent
{\bf Notes structure.} Note for this lecture
are written like this: first my personal
notes taken during the lecture,
then the AI-transcribed photos of the
slides.


\medskip\noindent
{\bf What's a moduli problem?} 
Let $T$ be the base scheme.
$\mathcal{M}(T)$ is a groupoid, the families
parametrized by $T $ + automorphisms.
If $T' \to T$ is a map, then there is a
map $\mathcal{M}(T) \to \mathcal{M}(T')$.
This defines a pseudofunctor. 
Hence the definition of stack (gluing
condition missing explanation though).

Some stacks are better than other.
The better ones, where we can do geometry,
are called algebraic.

\medskip\noindent
{\bf Example 1.}
$\mathcal{H}_1=$ ideal sheaves of length-1
subschemes. E.g. $X = \mathbb{P}^1$, then
$\mathcal{H}_1 = B\mathbb{G}_m$.
E.g. $\dim X = 2$ then $\mathcal{H}_1=
X \times B \mathbb{G}_m$ (uses
Hartog's lemma).

\medskip\noindent
{\bf Existence criterion.}
To further relate good moduli spaces
to good quotients:

\begin{theorem}[Alper]
\label{16th-alga-meeting-2026-theorem-alper}
If $\mathcal{M} \xrightarrow{\varphi}M$
is a good moduli space, then
\begin{enumerate}
\item the map is universally submersive 
(here universally means that it is
submersive for any base change).

\item (Universal map to algebraic spaces.)
$\varphi$ is initial among morphisms
to algebraic spaces. This means that
for any map $\mathcal{M} \to T$ there exists
a unique $M \to T$ which factors $\varphi$.

\item For every algebraically closed field
$K$, the $K$-points of $M$ are in bijection
with the $K$-points of the stack
up to the closure equivalence, i.e.
$$
M(K) =
\mathcal{M}(K)/\operatorname{closure}\sim
$$

\end{enumerate}
\end{theorem}



\medskip\noindent
{\bf Example 3: stack
$\mathcal{U}_{\{\mathbb{P}^1\},1}$.}
$i:\mathbb{P}^1 \hookrightarrow X$
non-removable (which means that
the corresponding point in the Hilbert
scheme is isolated) such that
$\mathcal{N}_{\mathbb{P}^1/X}$
is non-positive (recall that on $\mathbb{P}^1$
this can be split as a sum of line bundles).

Figure: square divided in three, horizontal
bar on bottom, creating rectangle. The
remaining space has a diagonal. Left of
diagonal is $\mathcal{H}_1$,
right of diagonal has:
$$
F \operatorname{ s.t.}
O \to \mathcal{O}_{\mathbb{P}^1}(-1)
\to F \to \mathcal{I}_{\mathbb{P}^1}
\to 0
$$
bottom rectangle has
$$
F = i_*\mathcal{O}_{\mathbb{P}^1}(-1)
\oplus \mathcal{I}_{\mathbb{P}^1}.
$$

This square maps to $\bar{X}=
\operatorname{Hilb}^1(X)_{\{\mathbb{P}^1\}}$.

There is a new diagonal: its the previous
diagonal, which intersects the bottom
horizontal bar in the middle of the square,
reflected by a vertical line on the middle.
In the leftmost section, let $x \in
\mathbb{P}^1$. Then its ideal sheaf lies
in a short exact sequence
$$
0 \to \mathcal{I}_{\mathbb{P}^1}
\to \mathcal{I}_x
\to i_*\mathcal{O}_{\mathbb{P}^1}(-1)
\to 0.
$$
We thus have
$$
\xymatrix{
\mathcal{H}_1 \ar@{^{(}->}[r]
\ar[d]
&\operatorname{square}\ar[d]
&\mathcal{U}_{\{\mathbb{P}^1\},1}^1
\ar@{^{(}->}[l]\ar[d]\\
\mathbb{P}^1\subset X
\ar[r]
& \ast \in \bar{X} = 
\operatorname{Hilb}^1(X)_{\{\mathbb{P}^1\}}
&X^0=\operatorname{Hilb}^1(X)_{\{\mathbb{P}^1\}}
\ar@{^{(}->}[l]
}
$$




\medskip\noindent
{\bf Theorem A.}
Figure: same square divided in three.
Bottom: split torsion type. Left: ideal
type. Right: non-split torsion type.

Also:
$$
\xymatrix{
\mathcal{H}_{\ell}
\ar@{^{(}->}[r]\ar[d]^{gms}
& \mathcal{U}_{\alpha,\ell}\ar[d]^{gms}\\
\operatorname{Hilb}(X)\ar[r]
&\operatorname{Hilb}^{\ell}(X)_{\alpha}.
}
$$

$\alpha \subset \operatorname{Hilb}(X)$
bounded gen and closed, its elements are:
\begin{itemize}
\item integral, and
\item rational or divisors.
\end{itemize}

\noindent
(LR).
\begin{itemize}
\item all $Y \in \alpha$ have the same
dimension.

\item $\Hom(\mathcal{I}_Y,i_*G)=0)$,
$i:Y' \hookrightarrow X$ for any t.f.
$\operatorname{rk} = ?$ $G$ s.t.
$P_G=?$.
\end{itemize}



\medskip\noindent
{\bf $\mathbb{P}^1$-fibration
surgeries, $\ell=1$.}
$$
\xymatrix{
\alpha \subset \operatorname{Hilb}(X)
\ar@{^{(}->}[r]&
Y\subset X\ar[d]^{\mathbb{P}^1}\\
& B.
}
$$
$$
\xymatrix{
\mathcal{H}_1 \ar@{^{(}->}[r]\ar[d]
& U_{\alpha,1}\ar[d]
& \mathcal{U}_{\alpha,1}^0
\ar[d]\ar@{_{(}->}[l]\\
X \ar[r]& B \subset \bar{X}
&X^0 \ar[l]\\
Y \ar@{^{(}->}[u]\ar[r]^{\mathbb{P}^1}
&B \ar@{^{(}->}[u]
&Y^0\ar@{^{(}->}[u]\ar[l]_{\mathbb{P}^1}
}
$$
where arrows labelled $\mathbb{P}^1$
are $\mathbb{P}^1$-bundles.

\medskip\noindent
{\bf DK conjecture.}
In fact, Bridgeland's (result?)
is a particular case of ours.
Their proof of the DK conjecture uses
some result of Orlov to produce the
embedding of derived categories.








\medskip\noindent
{\bf From this point on these are the
transcribed slides.}

The main theorem was proved during
Andres' two-week visit to the MSI in
2025, supported by the ECRVP
programme.

Other ways to visit Australia:
\begin{itemize}
\item SMRI International Visitor Program

\item The MSI--CNRS collaboration
supports visits from academics
affiliated in France
\end{itemize}

\medskip\noindent
{\bf Motivating question}

$X$ -- smooth projective variety over a
field $k$

$j\colon Y\longrightarrow X$ -- closed
subvariety

$\pi\colon Y\longrightarrow B$ --
surjective morphism

\medskip\noindent
{\bf Answers}

\begin{itemize}
\item
[Artin '70] Good understanding of the
formal completion of $Y\subset X$
$\rightsquigarrow$ conditions for
existence of $\overline X$ as an
algebraic space

\item
[Koll\'ar, Mori '98] Contractions
associated to $\omega_X$-negative
extremal faces in the cone of curves using
a semiample line bundle on $X$
\end{itemize}

\medskip\noindent
{\bf Our approach}

Reinterpret $X$ as a moduli problem

\medskip\noindent
{\bf Stacks}

\medskip\noindent
{\bf Definition}

$$
\operatorname{Groupoid}
=
\text{category, all morphisms are
isomorphisms}
$$

$$
\operatorname{Gpd}
=
\text{2-category of groupoids}
$$

\medskip\noindent
{\bf Definition}

$$
\operatorname{Stack}
=
\mathcal M\colon
\operatorname{Sch}_k^{\mathrm{op}}
\longrightarrow
\operatorname{Gpd}
+
\text{gluing condition}
$$

\medskip\noindent
{\bf Definition}

$\mathcal M$ algebraic $=$

$\exists$ smooth surjection
$U\longrightarrow\mathcal M$ from some
scheme $U$,

and
$\mathcal M\longrightarrow
\mathcal M\times_k\mathcal M$
is representable

\medskip\noindent
{\bf Example 0: $B\mathbb G_m$}

$k$ a field

$$
T\in\operatorname{Sch}_k
\rightsquigarrow
\mathcal M(T)
:=
\{
\mathbb G_m\text{-torsors over }T
\}
$$

Morphisms in $\mathcal M(T)=$
isomorphisms

\medskip\noindent
{\bf Definition}

Denote this stack $\mathcal M$ by
$B\mathbb G_m$

\medskip\noindent
{\bf Example 1: ideal sheaves of one
point}

$X$ smooth projective variety over any
field $k$

$$
T\in\operatorname{Sch}_k
\rightsquigarrow
$$

$$
\mathcal H_1(T)
:=
\left\{
\mathcal I_Z\text{ ideal sheaf}
\ \middle|\
\begin{array}{l}
Z\subset X\times T
\text{ a }T\text{-flat subscheme};\\
\forall t\in T:
Z|_{X\times t}
\text{ length }1\text{ subscheme}
\end{array}
\right\}
$$

Morphisms in $\mathcal H_1(T)=$
isomorphisms of sheaves

\medskip\noindent
{\bf Important cases}

\begin{itemize}
\item If $X=\mathbb P^1$, then
$$
\mathcal H_1\simeq B\mathbb G_m
$$

\item If $\dim X>1$, then
$$
\mathcal H_1
\simeq
X\times B\mathbb G_m
$$
(Hartogs' lemma)
\end{itemize}

\medskip\noindent
{\bf Example 2: ideal sheaves of several
points}

$X$ smooth projective variety over a
field $k$

$\ell>0$ natural number

$$
T\in\operatorname{Sch}_k
\rightsquigarrow
$$

$$
\mathcal H_\ell(T)
:=
\left\{
\mathcal I_Z\text{ ideal sheaf}
\ \middle|\
\begin{array}{l}
Z\subset X\times T
\text{ a }T\text{-flat subscheme};\\
\forall t\in T:
Z|_{X\times t}
\text{ length }\ell\text{ subscheme}
\end{array}
\right\}
$$

Morphisms in $\mathcal H_\ell(T)=$
isomorphisms of sheaves

\medskip\noindent
{\bf Important cases}

\begin{itemize}
\item If $X=\mathbb P^1$, then
$$
\mathcal H_\ell\simeq B\mathbb G_m
$$

\item If $\dim X>1$, then
$$
\mathcal H_\ell
\simeq
\operatorname{Hilb}^{\ell}(X)
\times B\mathbb G_m
$$
\end{itemize}

\medskip\noindent
{\bf Good moduli spaces}

\medskip\noindent
{\bf Definition}

Good moduli space (gms) for $\mathcal M$
$=$ qcqs morphism
$$
\varphi\colon\mathcal M\longrightarrow M
$$
to an algebraic space $M$ such that
\begin{itemize}
\item
$$
\varphi_*\colon
\operatorname{QCoh}\mathcal M
\longrightarrow
\operatorname{QCoh}M
\quad\text{exact}
$$

\item
$$
\mathcal O_M
\longrightarrow
\varphi_*\mathcal O_{\mathcal M}
\quad\text{isomorphism}
$$
\end{itemize}

\medskip\noindent
{\bf Example}

$G$ (linearly reductive) acts on variety
$V$,

$$
V\longrightarrow M
\text{ good quotient}
\Longrightarrow
[V/G]\longrightarrow M
\text{ gms}
$$

\medskip\noindent
{\bf Existence criterion
[Alper--Halpern-Leistner--Heinloth
'18--'23]}

$\mathcal M$ algebraic, finite type,
affine diagonal

\medskip\noindent
{\bf Theorem}

$$
\mathcal M
\text{ has a separated gms }M
\Longleftrightarrow
$$

$$
\mathcal M
\text{ is locally linearly reductive,
$\Theta$-reductive and S-complete}
$$

\medskip\noindent
{\bf Theorem}

$$
M\text{ is proper}
\Longleftrightarrow
\mathcal M
\text{ satisfies existence of valuative
criterion for properness}
$$

\medskip\noindent
{\bf Examples}

\medskip\noindent
{\bf Example 0}

$$
B\mathbb G_m
\longrightarrow
\operatorname{Spec}k
\qquad\text{gms}
$$

\medskip\noindent
{\bf Example 1: ideal sheaves of 1 point}

$$
X\times B\mathbb G_m
\longrightarrow
X
\qquad\text{gms}
$$

\medskip\noindent
{\bf Example 2: ideal sheaves of
$\ell$ points}

$$
\operatorname{Hilb}^{\ell}(X)
\times B\mathbb G_m
\longrightarrow
\operatorname{Hilb}^{\ell}(X)
\qquad\text{gms}
$$

\medskip\noindent
{\bf Example 3: stack
$\mathcal U_{\{\mathbb P^1\},1}$}

\medskip\noindent
{\bf Setup}

\begin{itemize}
\item
$i\colon\mathbb P^1\hookrightarrow X$
rigid smooth rational curve

(i.e. $\{\mathbb P^1\}\subset
\operatorname{Hilb}(X)$ is open and
closed)

\item
$\mathcal N_{\mathbb P^1/X}$ splits into
l.b. of degree $\leq 0$
\end{itemize}

Points of
$\mathcal U_{\{\mathbb P^1\},1}$ of two
types:
\begin{itemize}
\item
$$
F\simeq\mathcal I_x
\quad\text{for }x\in X
$$

\item $F$ that fits into SES
$$
0\longrightarrow
i_*\bigl(\mathcal O_{\mathbb P^1}(-1)\bigr)
\longrightarrow F
\longrightarrow
\mathcal I_{\mathbb P^1}
\longrightarrow 0
$$
\end{itemize}

\medskip\noindent
{\bf Theorem A}

Can replace:
\begin{itemize}
\item
$$
\{\mathbb P^1\}
\rightsquigarrow
\alpha\subset\operatorname{Hilb}(X)
$$

\item
$$
1
\rightsquigarrow
\text{length }\ell\geq 1
$$

\item $+$ conditions
\end{itemize}

\medskip\noindent
{\bf Theorem A [FH, M '25]}

$$
\mathcal U_{\alpha,\ell}
\text{ has a gms, denoted }
\operatorname{Hilb}^{\ell}(X)_\alpha
$$

\medskip\noindent
{\bf Contracting rational curves}

\medskip\noindent
{\bf Known results}

\begin{itemize}
\item
[Reid '83] $\dim X=3$,
$$
\mathcal N_{\mathbb P^1/X}
\simeq
\mathcal O(-1)\oplus\mathcal O(-1)
\text{ or }
\mathcal O\oplus\mathcal O(-2)
$$
$$
\Longrightarrow
\mathbb P^1\text{ contracts}
$$

\item
[Jim\'enez '92] $C$ -- smooth curve of
genus $g$ -- contracts
$$
\Longleftrightarrow
\exists\text{ l.b. }\mathcal L
\text{ on a nbhd of }C\text{ in }X
\text{ s.t.:}
$$
$$
\deg(\mathcal L|_C)\leq g-2
\text{ and }
H^1(\widehat X,\widehat{\mathcal L})
\text{ is finite-dimensional}
$$
\end{itemize}

\medskip\noindent
{\bf Corollary of Theorems A and B}

Note:
$\mathcal N_{\mathbb P^1/X}$ splits as
$\bigoplus_i\mathcal O(n_i)$.
If all $n_i\leq 0$, then:

$\mathbb P^1$ contracts in $\overline X$
and gets replaced by
$$
\mathbb P H^1
\bigl(
\mathbb P^1,
\mathcal N_{\mathbb P^1/X}(-1)
\bigr)
$$
in $X^\circ$.

\medskip\noindent
{\bf Abel-Jacobi contractions}

$$
\dim X=2
$$

$\alpha\subset\operatorname{Hilb}(X)$
open and closed subfamily of geom.
integral curves

$$
\mathcal C\longrightarrow\alpha
\quad\text{universal curve}
$$

Abel-Jacobi morphism
$$
AJ_\ell\colon
\operatorname{Hilb}^{\ell}(\mathcal C/\alpha)
\longrightarrow
\overline{\operatorname{Pic}}
(\mathcal C/\alpha)_{\chi(\mathcal O)-\ell}
$$

$$
\xymatrix@C=24pt@R=26pt{
&
\operatorname{Hilb}^{\ell}(\mathcal C/\alpha)
\ar[rr]
\ar[dl]_{AJ_\ell}
\ar[d]
&&
\operatorname{Hilb}^{\ell}(X)
\ar[d]
\\
\overline{\operatorname{Pic}}
(\mathcal C/\alpha)_{\chi(\mathcal O)-\ell}
&
H_\ell
\ar[l]
\ar[rr]
&&
\operatorname{Hilb}^{\ell}(X)_\alpha
}
$$

\medskip\noindent
{\bf A Mukai flop}

$X\longrightarrow\mathbb P^1$ -- elliptic
K3 with two singular fibers $X_a$, $X_b$
of type $I_3$

$X_a=C_1\cup C_2\cup C_3$ -- cycle of
three $\mathbb P^1$'s

$\alpha=\{C_1,C_2,C_3\}$ -- isolated
points in $\operatorname{Hilb}(X)$

$$
\ell\geq 2
\Longrightarrow
\text{condition (LR) is satisfied}
$$

\medskip\noindent
{\bf Corollary of Theorems A and B}

\begin{itemize}
\item
$\operatorname{Hilb}^{\ell}(X)
\longrightarrow
\operatorname{Hilb}^{\ell}(X)_\alpha$
is a contraction of
$\bigsqcup\operatorname{Hilb}^{\ell}(C_i)$

\item
$\operatorname{Hilb}^{\ell}(X)^\circ_\alpha$
is a Mukai flop
\end{itemize}

$$
\operatorname{Hilb}^{2}(X)_\alpha
\text{ is not\ldots}
$$

\begin{itemize}
\item \ldots a scheme [Namikawa '01]

\item \ldots a Bridgeland moduli space
(all such admit a strictly nef l.b.)
\end{itemize}

\medskip\noindent
{\bf $\mathbb P^1$-fibration surgeries,
$\ell=1$}

$X$ -- smooth projective variety

$Y\subset X$ -- subvariety with a
$\mathbb P^1$-fibration
$Y\longrightarrow B$ s.t.

$B$ corresponds to an open and closed
subscheme
$\alpha\subset\operatorname{Hilb}(X)$

$\mathcal N_{Y/X}$ is fiberwise
non-positive

\medskip\noindent
{\bf Corollary of Theorems A and B}

We have a surgery diagram
$$
\xymatrix@C=55pt@R=28pt{
X\ar[dr]_{p_1}
&&
X^\circ\ar[dl]^{p_2}
\\
&
\overline X
&
}
$$

\medskip\noindent
{\bf DK conjecture}

$X$ -- smooth projective variety

$$
X\supset Y\longrightarrow B
\quad\text{-- $\mathbb P^1$-fibration}
$$

$$
B\longleftrightarrow
\alpha\subset\operatorname{Hilb}(X)
$$

Assume $X^\circ$ is smooth.

$\mathcal N_{Y/X}$ is fiberwise
non-positive.

$$
\xymatrix@C=55pt@R=28pt{
X\ar[dr]_{p_1}
&&
X^\circ\ar[dl]^{p_2}
\\
&
\overline X
&
}
$$

\medskip\noindent
{\bf DK conjecture
[Kawamata '18, Conj. 1.2]}

If
$$
\forall b\in B:
\quad
\omega_X\cdot Y_b\leq 0,
$$
then
$$
\exists\text{ an embedding of an
admissible subcategory }
D^b(X^\circ)\longrightarrow D^b(X).
$$

\medskip\noindent
{\bf DK conjecture}

\medskip\noindent
{\bf Theorem [Bridgeland '02]}

$X$ -- complex projective threefold with
terminal singularities,

$$
X\longrightarrow\overline X
\longleftarrow X^\circ
$$
-- crepant resolutions.

Then
$$
D^b(X^\circ)\simeq D^b(X).
$$

\medskip\noindent
{\bf Theorem [FH, M '25]}

Verified the DK conjecture in the above
setting.

Our assumptions
$$
\Longrightarrow
\omega_X\cdot Y_b=-2,-1,0.
$$

\begin{itemize}
\item Case $\omega_X\cdot Y_b=-2$:
$$
X^\circ=\emptyset
$$

\item Case $\omega_X\cdot Y_b=-1$:
$$
X^\circ\simeq\overline X
=
\text{Fujiki-Nakano contraction}
$$

\item Case $\omega_X\cdot Y_b=0$:
$$
D^b(X^\circ)\simeq D^b(X)
$$
\end{itemize}



\section{Enriched inflection for linear
series on algebraic curves}
\label{16th-alga-meeting-2026-section-enriched-inflection-linear-series}

\medskip\noindent
{\bf Poster.}

E. Cotterill
(Imecc--Unicamp),
with I. Darago, C. Han, and N. Zhang.

\texttt{ethan@unicamp.br}

\medskip\noindent
{\bf Abstract.}

The fundamental local invariant of a
linear series on an algebraic curve is its
inflection, as measured by vanishing of
local derivatives in a point.
With Darago, Han, and Zhang we computed
the classes of inflection points and of
secant planes in the Grothendieck--Witt
ring of an arbitrary perfect base field
of characteristic $\neq 2$.

\medskip\noindent
{\bf Rank loci in algebraic geometry.}

Enumerative algebraic geometry deals
traditionally with the following basic
paradigm: given a map between vector
bundles
$$
\varphi\colon E\longrightarrow F
$$
over a complex algebraic variety $X$,
we'd like to count the ``number'' of
points $x\in X$ over which $\varphi_x$
has rank $\leq k$, for fixed
$k\in\mathbb{N}$.

These rank loci satisfy the principle of
{\it conservation of number}: the number
of solutions, counted with
multiplicities, is uniquely determined.
For other choices of base field $F$, the
number of $F$-rational solutions is not
typically unique.

\medskip\noindent
{\bf 1.\ Topological degrees over
$\mathbb{R}$.}

Given any continuously differentiable
map between compact oriented
$n$-manifolds
$$
f\colon X\longrightarrow Y,
$$
we may define its {\it topological degree}
to be the integer associated with the
induced map on homology
$$
f_*\colon
\mathbb{Z}=H_n(X,\mathbb{Z})
\longrightarrow
H_n(Y,\mathbb{Z})=\mathbb{Z}.
$$
It gives a signed count of preimages,
that is conserved as the point in the
target varies.

These global topological degrees are sums
of local topological degrees:
$$
\deg f
=
\sum_{x\in f^{-1}(y)}
\deg_x f,
$$
where for simplicity we will always
assume that $x\in X$ of $f$ is isolated
in $f^{-1}(y)$, $y=f(x)$.

Whenever $x$ is a {\it regular} value,
the differential
$$
df\colon T_xX\longrightarrow T_{f(x)}Y
$$
is an isomorphism, given by a Jacobian
matrix $J_f(x)$.
In this situation
$$
\deg_x f=\operatorname{sgn}\det J_f(x);
$$
the sign detects whether $f$ is locally
orientation-preserving or reversing.

In the 1970s, Eisenbud--Levine and
independently Khimshiashvili discovered
how to compute local degrees at
not-necessarily regular values using
algebraic geometry.

To be concrete, assume
$$
f=(f_1,\ldots,f_n)\colon
(\mathbb{R}^n,0)
\longrightarrow
(\mathbb{R}^n,0)
$$
is the germ of a map, in which each of
the component functions is $C^\infty$.

We impose that $f$ be
{\it finitely-determined}, i.e., that for
some finite $k\in\mathbb{N}$, $f$ and its
order-$k$ Taylor expansion are
contact-equivalent.

We can then identify each $f_i$ with its
Taylor series expansion; the
{\it local algebra} of $f$ is
$$
Q(f)
:=
\mathbb{R}[[x_1,\ldots,x_n]]
/
(f_1,\ldots,f_n).
$$

Assume the zero of $f$ in $0$ is
isolated.
$Q(f)$ is an Artinian local ring.

$Q(f)$ comes equipped with a distinguished
symmetric nondegenerate bilinear form:
the {\it Scheja--Storch form} $\Phi$.
Eisenbud--Levine and (independently)
Khimshiashvili showed that
$$
\deg_0 f=\operatorname{sgn}(\Phi).
$$

\medskip\noindent
{\bf Thm (EKL).}
Let $\varphi$ be any real-linear
functional
$$
\phi\colon Q(f)\longrightarrow\mathbb{R},
$$
chosen so that
$$
\phi(J_f)=1.
$$
The signature of the bilinear form
$$
\langle a,b\rangle_\phi
:=
\phi(a,b)
$$
is precisely $\deg_0 f$.
In particular, the signature of the
Scheja--Storch form computes the local
topological degree.

\medskip\noindent
{\bf Finding the Scheja--Storch form.}

To motivate the Scheja--Storch form,
assume
$$
A
:=
\mathbb{R}[y_1,\ldots,y_n]
\xrightarrow{\;f\;}
\mathbb{R}[x_1,\ldots,x_n]
=:
B,
\qquad
y_i\longmapsto f_i(x_1,\ldots,x_n)
$$
is an extension of rings that makes $A$
into a finite $B$-module; and assume the
only preimage of $0$ in the
complexification $f_{\mathbb{C}}$ of $f$
is $0$; the global coordinate ring
$$
\mathbb{R}[x_1,\ldots,x_n]/I
$$
is then isomorphic to the local algebra
$$
\mathbb{R}[[x_1,\ldots,x_n]]/I.
$$

There is a flat family of algebras
$$
Q_y(f)
:=
\mathbb{R}[x_1,\ldots,x_n]
/
(f_1-y_1,\ldots,f_n-y_n),
\qquad
y=(y_1,\ldots,y_n),
$$
parameterized by $\mathbb{R}^n$; the
central fiber is
$$
Q_0(f)=R/I.
$$
When $y$ is a regular value, $Q_y(f)$ is
a product of $s$ copies of $\mathbb{R}$
(corresponding to the $\mathbb{R}$-valued
points $x_1,\ldots,x_s$ in the fiber) and
the topological degree is
$$
\sum_{i=1}^s
\operatorname{sgn}(J_f(x_i)).
$$

With our assumptions, $B$ is a global
complete intersection over $A$
(in particular Gorenstein), so
$\omega_{B/A}$ is a free $B$-module of
rank $1$; and its canonical generator is
the
$A=\mathbb{R}[y_1,\ldots,y_n]$-linear
functional
$$
T\colon
\mathbb{R}[x_1,\ldots,x_n]
\longrightarrow
\mathbb{R}[y_1,\ldots,y_n],
\qquad
P\longmapsto
\operatorname{Tr}\left(\frac{P}{J_f}\right),
$$
where $\operatorname{Tr}(\beta)$ denotes
the trace of the
$\mathbb{R}(y_1,\ldots,y_n)$-linear
endomorphism of
$\mathbb{R}(x_1,\ldots,x_n)$ given by
multiplication by $\beta$.

$T$ restricts to $\mathbb{R}$-linear
functionals $T_y$ on every $Q_y(f)$;
$\phi$ is a limit of bilinear forms on
fibers obtained from renormalizations of
these.

\medskip\noindent
{\bf Quadratically-enriched AG.}

10 years ago, Kass and Wickelgren showed
that the EKL topological local degree over
$\mathbb{R}$ are specific instances of
{\it motivic} local degrees arising in
$\mathbb{A}^1$-homotopy theory.
Signs generalize to elements of the
{\it Grothendieck--Witt} group
$\operatorname{GW}(F)$ over an arbitrary
base field $F$, $\operatorname{char}(F)
\neq 2$.
This is the group completion under
$\oplus$ of the semiring with respect to
$\oplus,\otimes$ of nondegenerate,
symmetric bilinear forms.

$\operatorname{GW}(F)$ is generated by
classes $\langle a\rangle$ associated to
the bilinear form
$$
(x,y)\longmapsto axy,
\qquad
a\in F^*,
$$
modulo the relations
$$
\langle ab^2\rangle=\langle a\rangle;
\qquad
\langle a\rangle+\langle b\rangle
=
\langle ab(a+b)\rangle,
\quad
a+b\neq 0;
$$
$$
\langle a\rangle+\langle-a\rangle
=
\langle1\rangle+\langle-1\rangle.
$$
Here the hyperbolic class
$$
\mathbb{H}
:=
\langle 1\rangle+\langle-1\rangle
$$
is distinguished.

\medskip\noindent
{\bf Instructive examples:}
$$
\operatorname{GW}(\mathbb{C})
=
\mathbb{Z}
\qquad
\text{(only invariant is the rank)};
$$
$$
\operatorname{GW}(\mathbb{R})
=
\mathbb{Z}\times\mathbb{Z}
\qquad
\text{(rank and signature)};
$$
$$
\operatorname{GW}(\mathbb{F}_q)
=
\mathbb{Z}\times\mathbb{Z}/2\mathbb{Z}
\qquad
\text{(rank and discriminant, modulo
squares)}.
$$

In order to quadratically enrich a rank
locus on a (smooth) variety $X$, two
conditions should be satisfied:
\begin{enumerate}
\item
{\it Realizability as an Euler class:}
The rank locus should represent the top
Chern class of a vector bundle $E$ of
rank $\dim(X)$.

\item
{\it Relative orientability:}
The line bundle
$$
\operatorname{Hom}
\bigl(\det(T_X),\det(E)\bigr)
$$
is isomorphic to a tensor square.
(In order to compute local Euler indices,
this isomorphism needs to be explicit.)
\end{enumerate}

Now let $Y$ be an algebraic curve.
Degree-$d$ holomorphic maps
$$
Y\longrightarrow\mathbb{P}^r
$$
are prescribed by {\it linear series}
$g_d^r$: pairs $(L,V)$ where $L$ is a
degree-$d$ line bundle on $Y$ and
$$
V\subset H^0(C,L)
$$
is an $(r+1)$-dimensional subspace of
holomorphic sections.

The fundamental {\it local} invariant of
a $g_d^r$ is its {\it inflection} in a
point $p\in Y$, which is the cumulative
deviation of the $p$-vanishing orders of
the $g_d^r$ from the generic sequence
$(0,1,\ldots,r)$.
Sum over all points of $C$ to compute the
{\it total inflection} of the $g_d^r$.

Inflection points of $(L,V)$ comprise the
non-surjectivity locus of
$$
V\otimes\mathcal{O}
\longrightarrow
J^{r+1}(L),
$$
where $J^{r+1}(L)$ is the bundle over $C$
with fibers
$$
H^0\bigl(L/L(-(r+1)p)\bigr).
$$
The {\it Pl\"ucker degree} is
$$
\gamma_C
:=
c_1\bigl(\det J^{r+1}(L)\bigr).
$$

Now assume that $Y$ is a hyperelliptic
curve ramified in $\infty_Y$ over
$\infty$.

\medskip\noindent
{\bf Thm (CDH).}
Let $F$ be a field with
$\operatorname{char}(F)\neq 2$, and let
$$
L=\mathcal{O}(2\ell\infty_Y),
$$
where $\ell\geq 1$ is a positive integer.
Associated to the complete linear series
$|L|$ on $Y$ there is a well-defined
arithmetic inflection class
$$
[\operatorname{Inf}]_{\mathbb{A}^1}
\in
\operatorname{GW}(F)
$$
given by
$$
\operatorname{Inf}_{\mathbb{A}^1}
=
\frac{\gamma_C}{2}\mathbb{H}.
$$

\section{On the classification of
compactified universal Jacobians}
\label{16th-alga-meeting-2026-section-compactified-jacobians}

\medskip\noindent
{\bf Lecture.}

Fillipo Viviani,
Roma Tre University, Italy.

\medskip\noindent
{\bf Schedule.}

Friday, September 4, 2026,
2:00--3:00 p.m.,
IMPA - Ricardo Ma\~n\'e Auditorium.

\medskip\noindent
{\bf Problem/definition.}
The universal jacobian paramaetrizes pointed
genus $g$ curves with a line bundle.
$\mathcal{M}_{g,n}$ paramterizes stable
curves.
We want to classify all modular
compactifications of the universal Jacobian
$\mathcal{J}_{g,n}$.

$$
\xymatrix{
\mathcal{J}_{g,n} \ar@{^{(}->}[r]\ar[d]&?\\
\mathcal{M}_{g,n}\ar[r]&
\overline{\mathcal{M}}_{g,n}.
}
$$

Let us fix the Euler caracteristic in this
work and use instead
$$
\mathcal{J}_{g,n}^\chi.
$$

\begin{definition}
\label{16th-alga-meeting-2026-definition-compactified-uni-jac-stk}
A {\it compactified universal Jacobian
stack}
$$
TF_{g,n}^{\chi}
=
\{(C,T):C \in \overline{M}_{g,n},
\operatorname{rk}I=1\operatorname{torsion-free-sheaf},\operatorname{characterstic}I=\chi\}
$$
such that there exists a relative proper
good moduli space relative gms:
$$
\xymatrix{
\overline{\mathcal{J}}_{g,n}^{\chi}
\ar[d]\\
\overline{J}_{g,n}^{\chi}\ar[d]\\
\overline{\mathcal{M}}_{g,n}
}
$$
where $\overline{\mathcal{J}}_{g,n}^{\chi}$
is a stack, $\overline{J}_{g,n}^{\chi}$
is a space,
and the arrow below is representable and
proper.
\end{definition}

\noindent
\begin{remark}
\label{16th-alga-meeting-2026-remark-other-ways}
There are other ways of define these
objectives, with semistable curves and line
bundles.
\end{remark}

\noindent
Eduardo has contributed in constructing
these kinds of objects.
We intend to classify them.

\medskip\noindent
\begin{definition}
\label{16th-alga-meeting-2026-definition-fine}
A compactified universal Jacobian
is called {\it fine}
if it is obtained by rigidifying
$\overline{\mathcal{J}}_{g,n}^{\chi}$,
namely
$\overline{J}_{g,n}^{\chi}=
\overline{\mathcal{J}}_{g,n}^{\chi}\underline{//}
\mathbb{G}_m$.
It is equivalent that it consists
only of simple indecomposable sheaves.
\end{definition}

\medskip\noindent
{\bf Main technical definition
used to classify these objects.}

\begin{definition}
\label{16th-alga-meeting-2026-definition-vine}
A {\it vine function of type $(g,n)$
and characterstic $\chi$} is
a function $\sigma$ from the half vine graphs of
type $(g,n)$ to the $\mathbb{Z}$ such that
\begin{enumerate}
\item For every half vine graph
$x=(e;h,A \subseteq [n])$
such that $\sigma(x)+\sigma(x^c)-\chi
\in \{0,1\}$ if =0 then
$x$ is $\sigma$-degenerate
$$
D(\sigma)=\{x \in \mathbb{D}_{g,n}:
x \sigma\operatorname{-degen}\}.
$$
\item For any triangle
$\Delta=\{x_1,x_2,x_3\}$.
If two are degenerate, say $x_i,x_j \in
D(\sigma)$ then so is the third, i.e. $x_k
\in D(\sigma)$. Also, 
$$
\sum_n
\sigma(x_i)-\chi = 
\begin{cases}
0\qquad &\operatorname{if} x_1,x_2,x_3
\in D(\sigma) \\
1\qquad &\operatorname{if only one is}\\
\{1,2\}&\operatorname{none is}.
\end{cases}
$$
\end{enumerate}
\end{definition}

Let $\Sigma_{g,n}^{\chi}$ be the set
of vine functions. We make it a poset with
$\sigma \geq \sigma_2 \iff
\sigma_1(x) \geq \sigma_2(x)$ for all
$x \in \mathbb{D}_{g,n}$.

\begin{theorem}[Main Theorem A]
\label{16th-alga-meeting-2026-theorem-main-a}
There exists an anti-isomorphism of posets
between $\Sigma_{g,n}^{\chi}$ and the
poset of compactified universal Jacobians
over $\overline{\mathcal{M}}_{g,n}$ of
characteristic $\chi$.
\end{theorem}

\noindent
The map is
$$
\sigma \mapsto 
\overline{\mathcal{J}}_{g,n}(\sigma)
=
\begin{cases}
(C,I)\in TF_{g,n}^{\chi}\\
\chi(I_Y) \geq
\sigma(\operatorname{type}(Y))\forall Y
\in \operatorname{BCon}(C)
\end{cases}
$$
\noindent
where
$$
\operatorname{BCon}(C)
=
\begin{cases}
Y \subseteq C \\
Y,Y^c=\overline{C\_Y}\qquad
&\operatorname{connected}
\end{cases}
$$
$$
Y \mapsto (Y \wedge Y^c,
p_a(Y),
\{a \in [n]:p_a \in Y\}.
$$

This implies that the Euler caracteristic is
$$
\chi(I_Y)
\subseteq
\begin{cases}
\sigma(\operatorname{type}_C(Y)
+|\{p \in Y \cap
Y^c,I\operatorname{free-at}p\}|\qquad &
\operatorname{if} \operatorname{type}_C(Y)\\
\sigma(\operatorname{type}(Y))
+|\{p\in Y \cap
Y^C,I\operatorname{free-at}p\}|-1\qquad
\operatorname{otherwise}&
\end{cases}
$$

\medskip\noindent
{\bf Sketch of proof.}

\begin{enumerate}
\item (FP-) $\overline{\mathcal{J}}_{g,n}(\sigma)$
is compactified universal Jacobian
stack. We want a locally linear reductive
$\theta$-compatible $S$-complete universally
closed map
$$
\overline{\mathcal{J}}_{g,n}(\sigma)
\to
\overline{\mathcal{M}}_{g,n}.
$$
$$
\Sigma_{g,n}^{\chi}
a\to\{\sigma_C:
\operatorname{BCon}(C)\xrightarrow{\operatorname{type}_C}
\mathbb{D}_{g,n}\xrightarrow{\sigma}
\mathbb{Z}\}_{c \in
\overline{\mathcal{M}}_{g,n}}
$$
This map becomes an isomorphism if we take
inverse limit on the codomain.
Vine stability condition.

For any
$$
\xymatrix{
\mathcal{C}\ar[d]\\ T,
}
$$
$\lim_{\leftarrow} \{\sigma_t\}_{t \in T}$

Thus, this new stability condition gives us
a new way to construct the compactified
Jacobian.

\item (FPV2) for any $C$ nodal there is a
correspondence of compactified Jacobians
for $C$ that are limits of Jacobian for
smoothings of $C$ and
$\overline{J}_C(\sigma)$.

\end{enumerate}

That is, on step one we obtain as fibers
$\overline{J}_C(\sigma)$, and then
we use those in step 2.

\medskip\noindent
{\bf Classical compactified universal
Jacobian.}

\begin{definition}
\label{16th-alga-meeting-2026-definition-universal-num-polariztns}
$\operatorname{Pol}_{g,n}^{\chi}$
is the {\it universal numerical
polarizations}, the set of all
$\phi:\mathbb{D}_{g,n} \to \mathbb{R}$ such
that
\begin{enumerate}
\item $\phi(x) + \phi(x^c) = \chi$
for all $x \in \mathbb{D}_{g,n}$.
\item for all $D=\{x_1,x_2,x_3\}$,
$\sum\phi(x_i)=\chi$.
\end{enumerate}
\end{definition}

$$
\xymatrix{
\operatorname{Pol}_{g,n}^{\chi}
\ar[d]\\
\Sigma_{g,n}^{\chi}.
}
$$
$$
\phi\mapsto \left\lceil \phi \right\rceil 
=
\sigma(\phi)(x)=
\left\lceil \phi(x) \right\rceil 
$$

It's easy to see this is a vine function.
$\operatorname{Im}\left\lceil - \right\rceil=$
classical vine functions.

\medskip\noindent
$$
\overline{\mathcal{J}}_{g,n}(\phi)
=
\overline{\mathcal{J}}_{g,n}(\sigma(\phi))
=\{(C,I) \in TF_{g,n}^{\chi},
\chi(I_Y) \geq 
\phi(\operatorname{type}(Y))\},
$$
classical compactified universal Jacobians.

\medskip\noindent
{\bf Two constructions.}

The first one is due to Eduardo, who gave
it for an arbitrary family. Margarita then
did a stacky version of that.

\begin{theorem}[Esteves-Melo]
\label{16th-alga-meeting-2026-theorem-esteves-melo}
$E$ vector bundle over the universal curve
$$
\xymatrix{
\overline{\mathcal{C}}_{g,n}\ar[d]\\
\overline{\mathcal{M}}_{g,n},
}
$$
namely $E$ vector bundle on
$\overline{\mathcal{C}}_{g,n}$
with $\mu(E)=-\chi$.
\end{theorem}

\noindent
Fix a stable curve $Y$, restrict the vector bundle
over the fiber,  thus obtain a vector bundle
over $Y$. Then take the Euler characterstic
of that, thus obtaining the slope:

$$
\mapsto  \overline{ \mathcal{J}}_{g,n}^E
=
\begin{cases}
(C,I)\in TF_{g,n}^{\chi} \\
\chi(I_Y) \geq  -\mu(E|_Y)
\end{cases}
$$
exists $\phi(E)(x) = -\mu(E|_T)$
for all $Y \in \operatorname{BCon}(C)$
such that $\operatorname{type}(Y)=x$.

A special case of this construction is
$$
\overline{\mathcal{J}}_{g,n}
=
\overline{\mathcal{J}}_{g,n}^E.
$$

\medskip\noindent
The second construction is due to
Kass-Pagani. Let $L$ be an $\mathbb{R}$-line
bundle on $\overline{\mathcal{C}}_{g,n}$
with degree $\chi$.
We map it to
$$
\overline{\mathcal{J}}_{g,n}(L)
=
\begin{cases}
(C,I) \\
\chi(I_Y) \geq \operatorname{deg}(L|_Y)
\end{cases}
$$
This is also a particular case of our
construction, since, on $x$, we look at the
degree of $L$ restricted to $Y$, so:
$$
\operatorname{Pol}_{g,n}^\chi
\ni \phi(L)(x) = \operatorname{deg}(L|_Y).
$$

\medskip\noindent
Those are the two known constructions.

\medskip\noindent
{\bf Theorem.}

\begin{theorem}[affine case:Esteves,
in general: FPV1-2]
\label{16th-alga-meeting-2026-theorem-comp-jac}
\begin{enumerate}
\item 
$$
\begin{aligned}
\operatorname{RelPic}
(\overline{\mathcal{C}}_{g,n}
/\overline{\mathcal{M}}_{g,n}
&\xrightarrow{\simeq}
\operatorname{Pol}_{g,n}^x \\
L &\longmapsto \phi(L),
\end{aligned}
$$
where $\operatorname{deg}=\chi$;
real affine vector space.
This implies,
$$
\langle w_\pi\theta(\Delta_i),\operatorname{boundaries}
\rangle,
$$
that
$$
\{\overline{\mathcal{J}}_{g,n}(\phi)\}
=
\{\overline{\mathcal{J}}_{g,n}(L)\}
=
\{\overline{\mathcal{J}}_{g,n}^E\}
$$
where these are the classical, Kass-Pagani
and
Esteves-Melo, correspondingly.

\item Consider the universal curve above the
universal compactified jacobian:
$$
\xymatrix{
\overline{\mathcal{C}_{g,n}
\times_{\overline{\mathcal{M}}_{g,n}}}
\overline{\mathcal{J}}_{g,n}^E/\mathcal{I}
\ar[d]^\pi\\
\overline{\mathcal{J}}_{g,n}^E,
}
$$
$$
\mathcal{L}_E
:= \det_\pi(\mathcal{I} \otimes
p_1 \times E)^{-1}
$$
(Esteves polarization)
descends to a relatively ample line bundle
on.
\end{enumerate}
\end{theorem}


\medskip\noindent
{\bf Final result.}
We are almost done: we classified universal
compactified Jacobian. Are there any
nonclassical u.c. Jacobian, are they all
classical? The theorem 

\begin{theorem}[Main theorem B]
\label{16th-alga-meeting-2026-theorem-main-b}
There exist non-classical compactified
universal Jacobian (almost always, we
must exclude some exceptions)
$\iff$
$$
\begin{cases}
g>0\\
n>0\\
2g+n \geq f.
\end{cases}
$$

\noindent
The second case is very interesting.
Indeed, for $n=0$ there exists a unique,
for any $\chi$, universal compactified
Jacobian space over
$\overline{\mathcal{M}}_g$,
$$
\overline{J}^\chi
\left(\frac{\chi}{2g-2}w_\pi\right)
$$
Caporaso.

\item (V) A compactified universal Jacobian
space is classical $\iff $
projective/$\overline{\mathcal{M}}_{g,n}$.
(Similarly for stacks.)

\end{theorem}






% --- End 16th-alga-meeting-2026.tex ---

% --- Begin iv-jornadas-lefschetz-2026.tex ---
\chapter{IV Brazilian Workshop on Lefschetz Properties}
\label{chapter-iv-jornadas-lefschetz-2026}


Lefschetz Properties}

\medskip\noindent
UFBA, Salvador, Bahia.
September 28--October 2, 2026.

\href{%
https://jornadaslp.github.io/%
jornadas-lp/programacao.html%
}{Official programme.}
Talk titles and times checked on
September 28, 2026.
The sections follow the programme order.
Unannounced titles are marked explicitly.

\section{Introduction}
\label{iv-jornadas-lefschetz-2026-section-introduction}

\noindent
These notes start from Daniel's notes of
the introductory discussion on
September 28, 2026.
The conventions and qualifications below
make the geometric motivation precise.
The HP starting points are
\emph{Complex Geometry},
``Lefschetz theorems'', and
\emph{Commutative Algebra},
``Artinian algebras''.

\begin{theorem}[Hard Lefschetz]
\label{iv-jornadas-lefschetz-2026-theorem-hard-lefschetz}
Let $X$ be a compact Kähler manifold of
complex dimension $n$, with Kähler form
$\omega$.
Define the Lefschetz operator by
$$
\begin{aligned}
L:H^j(X,\mathbb{R})
&\longrightarrow H^{j+2}(X,\mathbb{R}),\\
\alpha&\longmapsto[\omega]\smile\alpha.
\end{aligned}
$$
For every $0\leq j\leq n$, the map
$$
L^{n-j}:H^j(X,\mathbb{R})
\longrightarrow H^{2n-j}(X,\mathbb{R})
$$
is an isomorphism.
For a smooth complex projective variety,
one may take the first Chern class of an
ample line bundle and use rational
coefficients.
\end{theorem}

\begin{remark}[Geometric meaning]
\label{iv-jornadas-lefschetz-2026-remark-geometric-meaning}
The point of Lefschetz properties is to
control the ranks of intersection maps.
Let $X$ be a smooth complex projective
variety and $L$ an ample line bundle.
Some power $M=L^{\otimes m}$ is very
ample. Choose a basis
$s_0,\ldots,s_N$ of $H^0(X,M)$.
On a local trivialization with frame $e$,
write $s_j=f_j e$. The embedding is
$$
i:X\hookrightarrow\mathbb{P}^N,
\qquad
i(x)=[f_0(x):\cdots:f_N(x)].
$$
The sections have no common zero.
Changing the frame multiplies all the
$f_j$ by the same nowhere-zero function,
so the projective point is unchanged.

Now take a hyperplane
$$
H=\{a_0z_0+\cdots+a_Nz_N=0\}.
$$
Substituting the coordinates of $i(x)$
gives
$$
\begin{aligned}
i(x)\in H
&\Longleftrightarrow
\sum_{j=0}^N a_j f_j(x)=0\\
&\Longleftrightarrow
\sum_{j=0}^N a_j s_j(x)=0.
\end{aligned}
$$
Thus the hyperplane pulls back to the
zero divisor of the global section
$$
s=\sum_{j=0}^N a_j s_j\in H^0(X,M).
$$
This is the link between hyperplanes and
the original line bundle: the pulled-back
hyperplane divisor has associated line
bundle $M$. Hence
$$
i^*\mathcal{O}_{\mathbb{P}^N}(1)
\cong M=L^{\otimes m}.
$$
Since the first Chern class turns tensor
products into sums, the hyperplane class
on $X$ is
$$
c_1(M)=m\,c_1(L).
$$
Consequently, multiplication by $c_1(L)$
means, up to a positive scalar,
intersection with a hyperplane section.

Hard Lefschetz says that suitable powers
of this multiplication give isomorphisms
between complementary cohomology degrees.
It therefore controls intersection maps
and constrains cohomology dimensions.
An arbitrary divisor class need not
have this property.
A compact Kähler manifold need not have
an ample line bundle; if it does,
the resulting embedding makes it
projective.

The algebraic WLP and SLP ask whether
some degree-one element has analogous
maximal-rank multiplication maps,
without requiring a geometric origin.
For a degeneration with toric components,
Hard Lefschetz applies to the components
that are smooth and projective.
This helps understand their intersection
theory, but does not by itself determine
their gluing or the smoothability of the
whole special fibre.
\end{remark}

\medskip\noindent
\textbf{Motivational question.}
Quais propriedades deve satisfazer uma
álgebra para ser candidata a ser a
cohomologia de uma variedade?

For a connected smooth complex projective
variety, its rational cohomology is
finite-dimensional and graded-commutative,
with a unit and $H^0=\mathbb{Q}$.
It satisfies Poincaré duality and Hard
Lefschetz for an ample class, and carries
further Hodge-theoretic structure.
These are necessary conditions; we are not
claiming that they are sufficient.

\begin{remark}[Which grading?]
\label{iv-jornadas-lefschetz-2026-remark-grading}
Multiplication by a Kähler class raises
cohomological degree by two.
To compare with commutative algebra, put
$$
A_k=H^{2k}(X,\mathbb{Q}).
$$
The ample class now has degree one.
The even cohomology is commutative, but
need not be generated by $A_1$.
Moreover, odd cohomology can be nonzero:
the entire cohomology algebra is generally
only graded-commutative.
Thus a standard graded commutative algebra
is a more restrictive model than the
cohomology algebra of an arbitrary variety.
\end{remark}

\begin{example}[Products of projective
spaces]
\label{iv-jornadas-lefschetz-2026-example-projective-products}
Let $a_0,\ldots,a_r$ be positive integers
and put
$$
X=\prod_{i=0}^{r}\mathbb{P}^{a_i-1}.
$$
If $x_i$ denotes the pullback of the
hyperplane class from the $i$th factor,
then
$$
H^*(X,\mathbb{Q})\cong
\frac{\mathbb{Q}[x_0,\ldots,x_r]}
{(x_0^{a_0},\ldots,x_r^{a_r})}.
$$
Each $x_i$ has cohomological degree two;
after regrading, it has algebra degree one.
The class
$$
\ell=x_0+\cdots+x_r
$$
is ample, and its Lefschetz operator is
multiplication by $\ell$.
Writing $N=\sum_{i=0}^{r}(a_i-1)$, Hard
Lefschetz gives isomorphisms
$$
\times\ell^{N-2k}:A_k\longrightarrow A_{N-k}
$$
for $0\leq k\leq N/2$.
The Lefschetz decomposition then gives
maximal rank for every power between any
two graded pieces.
\end{example}

For instance, for
$X=\mathbb{P}^1\times\mathbb{P}^1$,
$$
A=\mathbb{Q}[x,y]/(x^2,y^2),
\qquad \ell=x+y.
$$
The graded dimensions are $(1,2,1)$.
Multiplication sends $1$ to $x+y$,
sends both $x$ and $y$ to $xy$, and
$$
\ell^2=2xy.
$$
Thus $A_0\longrightarrow A_1$ is
injective, $A_1\longrightarrow A_2$ is
surjective, and
$\times\ell^2:A_0\longrightarrow A_2$
is an isomorphism.

\medskip\noindent
\textbf{Realization question.}
Which graded algebras can be realized as
the rational cohomology algebra of a
smooth complex projective variety?
This is a broad realization problem,
not a criterion supplied by SLP alone.
One must specify the class of varieties,
coefficients, and grading when asking it.

\begin{definition}[Standard graded algebra]
\label{iv-jornadas-lefschetz-2026-definition-standard-graded}
Let $K$ be a field of characteristic zero.
A \emph{standard graded algebra} is a
commutative graded $K$-algebra
$$
A=\bigoplus_{k\geq0}A_k
$$
with $A_0=K$, generated by a
finite-dimensional space $A_1$.
Equivalently, $A=Q/I$ for a polynomial ring
$Q=K[X_0,\ldots,X_r]$, with variables of
degree one and $I$ homogeneous.
Here an \emph{Artinian algebra} means that
$A$ is finite-dimensional over $K$, as in
the HP notes.
For standard graded algebras this agrees
with the descending-chain definition.
\end{definition}

\begin{definition}[Maximal rank]
\label{iv-jornadas-lefschetz-2026-definition-maximal-rank}
For $k\geq0$ and $t\geq1$, say that $A$
has the \emph{maximal rank property in
degree $k$ and range $t$} if there is a
linear form $\ell\in A_1$ such that
$$
\times\ell^t:A_k\longrightarrow A_{k+t}
$$
has rank
$$
\min\{\dim_K A_k,\dim_K A_{k+t}\}.
$$
Thus the map is injective or surjective,
as dictated by the two dimensions.
\end{definition}

\begin{remark}[Understanding rank]
\label{iv-jornadas-lefschetz-2026-remark-understanding-rank}
Recall that the rank of a linear map
is the dimension of its image:
$$
\operatorname{rank}(f)
=\dim_K\operatorname{im}(f).
$$
It counts the independent directions
reached in the target. Maximal rank
means that no input directions are lost
(injectivity), or all output directions
are reached (surjectivity), whichever
the dimensions permit.

For multiplication by $\ell^t$, put
$$
m=\dim_K A_k,\qquad n=\dim_K A_{k+t}.
$$
Having maximal rank means:
\begin{itemize}
\item If $m<n$, the map is injective.
\item If $m>n$, the map is surjective.
\item If $m=n$, the map is an isomorphism.
\end{itemize}
Thus WLP requires injectivity where
the graded dimensions increase,
surjectivity where they decrease,
and isomorphisms where they agree.
SLP makes the same demand between
$A_k$ and $A_{k+t}$ for every $t\geq1$.
\end{remark}

\begin{remark}
\label{iv-jornadas-lefschetz-2026-remark-maximal-rank-terminology}
The target is $A_{k+t}$ and the linear
form belongs to $A_1$.
This is the local terminology of these
notes. In the literature, ``Maximal Rank
Property'' can instead refer to
multiplication by general forms of each
degree, which need not be powers of a
linear form.
\end{remark}

\begin{definition}[Strong Lefschetz
property]
\label{iv-jornadas-lefschetz-2026-definition-slp}
An Artinian standard graded algebra $A$
has the \emph{Strong Lefschetz Property}
(SLP) if there is a single
$\ell\in A_1$ for which
$$
\begin{aligned}
\times\ell^t:A_k&\longrightarrow A_{k+t},\\
a&\longmapsto \ell^t a.
\end{aligned}
$$
has maximal rank for every $k\geq0$ and
every $t\geq1$.
Such an $\ell$ is a \emph{strong Lefschetz
element}.
\end{definition}

\medskip\noindent
\textbf{Question.}
Can we characterize all Artinian standard
graded algebras having SLP?

\begin{remark}[Why Artinian algebras?]
\label{iv-jornadas-lefschetz-2026-remark-why-artinian}
Estudamos álgebras artinianas pensando
em anéis de cohomologia de variedades.

Daniel's question: então toda álgebra de
cohomologia é artiniana?
For a compact smooth manifold, the total
cohomology over a field is
finite-dimensional.
The same finiteness holds for complex
algebraic varieties of finite type.
So the finite-dimensional part of the
intuition is correct in these settings.
It is not true for arbitrary topological
spaces.

There is still a commutativity qualification:
the full cohomology is graded-commutative.
Its even part is a commutative
finite-dimensional algebra, hence Artinian
in the HP sense.
When odd cohomology vanishes, this
qualification disappears for the full ring.
Standard generation in degree one after
regrading is an additional condition.
\end{remark}

\begin{definition}[Weak Lefschetz property]
\label{iv-jornadas-lefschetz-2026-definition-wlp}
An Artinian standard graded algebra $A$
has the \emph{Weak Lefschetz Property}
(WLP) if there is a single
$\ell\in A_1$ such that
$$
\begin{aligned}
\times\ell:A_k&\longrightarrow A_{k+1},\\
a&\longmapsto \ell a.
\end{aligned}
$$
has maximal rank for every $k\geq0$.
Such an $\ell$ is a \emph{weak Lefschetz
element}.
\end{definition}

The same $\ell$ must work in every
degree, and for SLP in every range.
Degree one alone does not guarantee
these rank conditions.

SLP implies WLP, by taking $t=1$.
The converse does not hold in general.

\begin{lemma}[Change of variables]
\label{iv-jornadas-lefschetz-2026-lemma-change-of-variables}
SLP and WLP are invariant under graded
algebra isomorphisms, in particular under
invertible linear changes of variables.
\end{lemma}

\begin{proof}
If $\varphi:A\longrightarrow B$ is a
graded algebra isomorphism, then
$$
\varphi(\ell^t a)=\varphi(\ell)^t\varphi(a).
$$
The isomorphisms on the graded pieces
therefore identify the multiplication maps
and preserve their ranks.
\end{proof}

\begin{lemma}[Powers of independent forms]
\label{iv-jornadas-lefschetz-2026-lemma-powers-of-forms}
Let $\ell_0,\ldots,\ell_r$ be linearly
independent linear forms in
$Q=K[X_0,\ldots,X_r]$.
For positive integers $a_i$, the algebra
$$
Q/(\ell_0^{a_0},\ldots,\ell_r^{a_r})
$$
has SLP in characteristic zero.
\end{lemma}

\begin{proof}
The invertible substitution
$X_i\longmapsto\ell_i$ identifies this
algebra with the monomial complete
intersection
$$
Q/(X_0^{a_0},\ldots,X_r^{a_r}).
$$
Over $\mathbb{Q}$, its SLP follows from the
product-of-projective-spaces example.
The multiplication matrices for the element
$X_0+\cdots+X_r$ have integer entries.
Their nonzero maximal minors stay nonzero
over every field of characteristic zero.
Apply the change-of-variables lemma.
\end{proof}

Linear independence is essential for this
argument: arbitrary linear forms need not
define a change of variables or even give
an Artinian quotient.

\begin{definition}[Artinian complete
intersection]
\label{iv-jornadas-lefschetz-2026-definition-complete-intersection}
An \emph{Artinian complete intersection}
is an algebra
$$
A=K[X_0,\ldots,X_r]/(f_0,\ldots,f_r),
$$
where the homogeneous forms $f_i$ form a
regular sequence and the quotient is
finite-dimensional.
A \emph{regular sequence} means that each
$f_i$ is a non-zero-divisor modulo the
preceding forms, and the generated ideal
is proper.
The algebra is \emph{equigenerated} if all
the $f_i$ have the same degree.
If there are no linear relations, its
codimension is $r+1=\dim_K A_1$.
\end{definition}

\begin{lemma}[Generic complete
intersections]
\label{iv-jornadas-lefschetz-2026-lemma-generic-complete-intersections}
Fix positive degrees $a_0,\ldots,a_r$.
A general Artinian complete intersection
of those degrees has SLP in
characteristic zero.
\end{lemma}

\begin{proof}
The complete-intersection locus in the
parameter space of tuples of forms of
these degrees is a nonempty Zariski-open
set. Its quotients have a fixed Hilbert
function and vanish above degree
$\sum_i(a_i-1)$.
Locally in this parameter space, choose
bases for the graded pieces.
For a fixed linear form, maximal rank of
each multiplication map is an open
condition: some maximal minor is nonzero.
Only finitely many maps are relevant.

At the tuple $(X_0^{a_0},\ldots,X_r^{a_r})$,
the form $X_0+\cdots+X_r$ works for all
these maps. Their simultaneous maximal
rank therefore holds on a nonempty open
set of tuples.
\end{proof}

Here ``general'' or ``generic'' means
belonging to a suitable nonempty
Zariski-open set in the parameter space
with the degrees fixed.
This conclusion uses openness in addition
to invariance under change of variables.
It does not say that every complete
intersection is a change of variables of
a monomial one.

\medskip\noindent
\textbf{The big problem.}
Does every Artinian complete intersection
in characteristic zero have WLP?
Does every such algebra have SLP?
The conjectural answer to both is yes.
The following established cases concern
\emph{WLP}, and must not be quoted as
general SLP theorems:

\begin{enumerate}
\item Every codimension-three Artinian
complete intersection has WLP.
This is the theorem of Harima, Migliore,
Nagel and Watanabe; see
\href{%
https://arxiv.org/abs/math/0208201%
}{The Weak and Strong Lefschetz
Properties for Artinian K-Algebras}.

\item In codimension four, equigenerated
complete intersections with common
relation degree $2\leq d\leq5$ have WLP.
Boij, Migliore, Miró-Roig and Nagel prove
the cases $d=3,4,5$ and recall the known
quadratic case in their introduction;
see \href{%
https://arxiv.org/abs/2212.09890%
}{On the Weak Lefschetz Property for
height four equigenerated complete
intersections}.
Here $d$ is the degree of each defining
form, not the socle degree.
\end{enumerate}

\begin{remark}[Characteristic matters]
\label{iv-jornadas-lefschetz-2026-remark-characteristic}
The characteristic-zero hypothesis is
part of these statements.
For example, in characteristic two,
every linear form in $K[x,y]/(x^2,y^2)$
has square zero, so this algebra fails
SLP although it is a monomial complete
intersection.
\end{remark}

\section{Thiago Holleben: problem session}
\label{iv-jornadas-lefschetz-2026-section-thiago-problems}

September 28, 2026.
Morning introduction and afternoon
problem-session exercises.

\begin{definition}[Artinian monomial
quotient]
\label{iv-jornadas-lefschetz-2026-definition-monomial-quotient}
Let $\Delta$ be a simplicial complex on
$\{1,\ldots,n\}$ and let
$Q=\mathbb{K}[x_1,\ldots,x_n]$.
Its \emph{Stanley--Reisner ideal}
$I_\Delta$ is generated by squarefree
monomials $x_{i_1}\cdots x_{i_j}$ whose
supports are not faces of $\Delta$.
For positive integers $a_1,\ldots,a_n$,
consider the \emph{Artinian monomial
quotient}
$$
A(\Delta;\mathbf{a})=
\frac{\mathbb{K}[x_1,\ldots,x_n]}
{I_\Delta+(x_1^{a_1},\ldots,x_n^{a_n})}.
$$
\end{definition}

The powers of the variables make this
quotient finite-dimensional: it is spanned
by monomials $x_1^{b_1}\cdots x_n^{b_n}$
with $0\leq b_i<a_i$.
This uses the Stanley--Reisner convention
in HP's \emph{Stanley-Reisner schemes}.

\begin{definition}[Edge ideal]
\label{iv-jornadas-lefschetz-2026-definition-edge-ideal}
For a finite simple graph $G$ on
$\{1,\ldots,n\}$, its \emph{edge ideal}
is
$$
I(G)=(x_ix_j\mid\{i,j\}\in E(G))
\subseteq Q.
$$
For a fixed integer $d\geq2$, Thiago's
handout uses the notation
$$
A(G,d)=\frac{\mathbb{K}[x_1,\ldots,x_n]}
{I(G)+(x_1^d,\ldots,x_n^d)}.
$$
The shorthand $A(G)$ below presupposes a
fixed choice of the exponents.
\end{definition}

\begin{remark}[Why add the powers?]
\label{iv-jornadas-lefschetz-2026-remark-why-thiago-adds-powers}
Adding a polynomial to the ideal means
declaring that polynomial to be zero in
the quotient. For example, take a graph
with two vertices joined by one edge.
Its edge ideal is $(xy)$, so we start with
$$
R=\mathbb{K}[x,y]/(xy).
$$
Here $xy=0$: a monomial using both
vertices is forbidden. But this does not
say that $x^2=0$ or $y^2=0$.
We can repeat either vertex indefinitely.
A vector-space basis of $R$ is
$$
1,x,y,x^2,y^2,x^3,y^3,\ldots.
$$
There are infinitely many basis elements,
so $R$ is infinite-dimensional.

Now add the squares:
$$
A=\mathbb{K}[x,y]/(xy,x^2,y^2).
$$
We have imposed two new rules,
$x^2=0$ and $y^2=0$. All higher powers
also vanish, since, for example,
$x^3=x\,x^2=0$.
Every element of $A$ can now be written
uniquely as
$$
a+bx+cy,\qquad a,b,c\in\mathbb{K}.
$$
Thus $A$ has basis $1,x,y$ and dimension
three. Its only nonzero graded pieces are
$A_0$ and $A_1$.

In general, the edge relations forbid
using adjacent vertices together.
The additional relation $x_i^d=0$
forbids repeating vertex $i$ $d$ times.
Each exponent is now at most $d-1$,
so only finitely many monomials survive.

This is why Thiago adds the powers:
he wants a finite-dimensional graded
algebra, with a last nonzero degree,
on which to study Lefschetz properties.
This finiteness resembles the cohomology
algebras in the introduction.
Before adding powers, each graded piece
was already finite-dimensional; what was
missing was finiteness of the whole
algebra. Adding powers achieves that,
but does not guarantee WLP or SLP.
\end{remark}

\begin{example}[The four-vertex path]
\label{iv-jornadas-lefschetz-2026-example-path-edge-ideal}
For the graph with edges
$\{1,2\}$, $\{2,3\}$ and $\{3,4\}$,
$$
I(G)=(x_1x_2,x_2x_3,x_3x_4).
$$
Daniel recognizes here the combinatorial
example used in his poster.
To compare the ideals, keep track of
which complex is being used:
$I(G)$ records the edges as forbidden
pairs. It is the Stanley--Reisner ideal
of the \emph{independence complex} of $G$,
whose faces are vertex sets containing
no edge of $G$.
In this example its maximal faces are
$$
\{1,3\},\quad\{1,4\},\quad\{2,4\}.
$$
Thus it is not the Stanley--Reisner ideal
of the complex whose edges are those of
$G$ with the same vertex labels.
\end{example}

\begin{definition}[Nagata idealization]
\label{iv-jornadas-lefschetz-2026-definition-nagata-idealization}
For a commutative algebra $A$ and an
$A$-module $M$, the \emph{Nagata
idealization}, or \emph{trivial extension},
is the vector space $A\oplus M$ with
product
$$
(a,m)(b,n)=(ab,an+bm).
$$
In particular, the product of any two
elements of $0\oplus M$ is zero.
\end{definition}

\medskip\noindent
\textbf{From $A(G)$ to a Gorenstein
algebra.}
For a finite-dimensional algebra $A$, the
standard choice is its dual module
$$
A^\vee=\operatorname{Hom}_{\mathbb{K}}
(A,\mathbb{K}),
\qquad (a\varphi)(b)=\varphi(ab).
$$
Then
$$
B=A\oplus A^\vee
$$
is an Artinian Gorenstein algebra.
Here this means that, for our local
Artinian algebra $B$, the socle has
dimension one; the socle consists of
elements annihilated by its maximal ideal.
Indeed, an element $(a,\varphi)$ in the
socle must have $a=0$ by multiplying by
all $(0,\psi)$, and then $\varphi$ must
vanish on the maximal ideal of $A$.
These functionals form the dual of
$A/\mathfrak{m}_A=\mathbb{K}$.

This explains the construction mentioned
in the introduction. An arbitrary module
$M$ does not automatically give a
Gorenstein idealization.
For graded Lefschetz questions, one must
also specify the grading shift of the
dual module and check standard generation;
the ungraded construction alone does not
assert those properties.

\medskip\noindent
\textbf{Question from the session.}
Find interesting examples of Gorenstein
algebras arising from Nagata idealizations
of $A(G)$.

\medskip\noindent
\textbf{Graph-theoretic questions.}
For fixed exponents, find properties of
$G$ guaranteeing that $A(G)$ has WLP,
and properties guaranteeing that $A(G)$
fails WLP.

\medskip\noindent
\textbf{Combinatorial approach and
probabilistic goal.}
O problema acima pode ser abordado
combinatoriamente.
Thiago's main goal, as described in the
session, is to establish that a random
algebra $A(G)$ fails WLP with high
probability in suitable random models.
An asymptotic version would ask that the
probability of failure tends to one as
the number of vertices tends to infinity.
The graph distribution, exponent $d$,
and parameter regime must be specified;
this is not a claim about every random
model. The handout distinguishes known
asymptotic results for certain random
simplicial-complex models from the desired
graph criteria in this problem session.

\medskip\noindent
\textbf{More concrete objective.}
Find criteria for failure of WLP using
the matching and independence polynomials
of $G$ and its complement.
A matching is a set of edges with no
shared vertices; an independent vertex
set contains no edge.
Thus these polynomials encode counts of
two different kinds of configurations.
The complement graph has exactly the
edges absent from $G$, between distinct
vertices.

For example, when $d=2$, a basis of
$A(G,2)_i$ consists of squarefree monomials
supported on independent $i$-element
vertex sets. Hence its Hilbert series is
the independence polynomial.
These counts give the dimensions of the
multiplication maps; the problem is to
use combinatorics to control their ranks.

\medskip\noindent
\textbf{Representation-theoretic idea.}
``When we have an operator, we are
interested in knowing its kernel.''
Work in characteristic zero and put
$$
A=\frac{\mathbb{K}[x_1,\ldots,x_n]}
{(x_1^2,\ldots,x_n^2)},
\qquad L=x_1+\cdots+x_n.
$$
Consider multiplication
$$
U=\times L:A_i\longrightarrow A_{i+1}.
$$
Permuting variables gives an action of
$S_n$ on $A$ commuting with $U$.
Hence its kernel is an $S_n$-representation.

\begin{remark}[Attribution to check]
\label{iv-jornadas-lefschetz-2026-remark-specht-attribution}
The live notes record ``Lefschetz--Bott:
the kernel of multiplication by $L$ is
spanned by Specht polynomials''.
The name has not been verified.
The precise statement below distinguishes
the raising and lowering operators.
\end{remark}

Write $x_S=\prod_{j\in S}x_j$.
These squarefree monomials form a basis
of $A$. Define linear maps on this basis:
$$
D(x_S)=\sum_{j\in S}x_{S\setminus\{j\}},
\qquad C(x_S)=x_{\{1,\ldots,n\}\setminus S}.
$$
Thus $D:A_i\longrightarrow A_{i-1}$
lowers degree, while
$C:A_i\longrightarrow A_{n-i}$
complements the support.
The definition of $D$ uses the unique
squarefree representative; ordinary
partial differentiation does not descend
to the quotient as a derivation.
Directly on basis monomials,
$$
CU=DC,\qquad C^2=1.
$$
Consequently,
$$
\ker(U|_{A_i})=C(\ker(D|_{A_{n-i}})).
$$

\begin{definition}[Two-row Specht
polynomial]
\label{iv-jornadas-lefschetz-2026-definition-specht-polynomial}
Choose $j$ disjoint pairs of indices
from $\{1,\ldots,n\}$, written
$$
(a_1,b_1),\ldots,(a_j,b_j).
$$
The corresponding \emph{two-row Specht
polynomial} is
$$
P=\prod_{r=1}^{j}(x_{a_r}-x_{b_r}).
$$
This is the polynomial attached to a
Young tableau with row lengths $(n-j,j)$,
whose two-entry columns are these pairs.
\end{definition}

\begin{theorem}[Specht description of
kernels]
\label{iv-jornadas-lefschetz-2026-theorem-specht-kernels}
For $0\leq j\leq\lfloor n/2\rfloor$,
$\ker(D|_{A_j})$ is spanned by the
Specht polynomials just defined.
For $j>n/2$, this kernel is zero.
Therefore, if $i\geq\lceil n/2\rceil$,
$\ker(U|_{A_i})$ is spanned by their
complements $C(P)$ with $j=n-i$.
For $i<n/2$, $\ker(U|_{A_i})=0$.
\end{theorem}

The lowering-operator statement is
Theorem 14 of Morita, Wachi and Watanabe,
\href{https://hdl.handle.net/2115/69582}{%
Zero-dimensional Gorenstein algebras
with the action of the symmetric group},
in the 2005 preprint.
The statement for $U$ follows from $C$.
Thiago's handout points to Section 9.3
of \emph{The Lefschetz Properties}.

For example, when $n=3$, the polynomial
$x_1-x_2$ lies in $\ker D$, whereas
$$
L(x_1-x_2)=x_1x_3-x_2x_3\ne0.
$$
Its complement $x_2x_3-x_1x_3$ does lie
in $\ker U$.
This shows why the distinction matters.

\medskip\noindent
\textbf{Session handout.}
\href{%
https://jornadaslp.github.io/%
jornadas-lp/files/%
problemas-matchings-wlp.pdf%
}{Matchings and WLP}, September 2026.
The handout fixes the notation $A(G,d)$;
the idealization question above is
recorded from the oral introduction.
For related constructions, see
\href{%
https://arxiv.org/abs/2112.09434%
}{On the Lefschetz Property for quotients
by monomial ideals containing squares
of variables}.

\medskip\noindent
\textbf{Afternoon exercises:
squares and independent sets.}
September 28, 2026.
In these exercises the exponent is two:
$$
\begin{aligned}
R&=\mathbb K[x_1,\ldots,x_n],\\
A(G)&=R/(I(G)+(x_1^2,\ldots,x_n^2)).
\end{aligned}
$$
A surviving monomial cannot repeat a
vertex, because the squares vanish.
It cannot contain adjacent vertices,
because edge monomials vanish.
Consequently, a basis of $A(G)_r$
consists of the monomials supported on
independent sets of size $r$.

\medskip\noindent
\textbf{Example: the pentagon.}
Number the vertices cyclically.
Then
$$
\begin{aligned}
I(C_5)=(&x_1x_2,x_2x_3,x_3x_4,\\
       &x_4x_5,x_5x_1).
\end{aligned}
$$
The algebra is
$$
A(C_5)=
R/(I(C_5)+(x_1^2,\ldots,x_5^2)).
$$
Its graded bases are:
$$
\begin{array}{c|l}
0&1\\
1&x_1,x_2,x_3,x_4,x_5\\
2&x_1x_3,x_1x_4,x_2x_4,x_2x_5,x_3x_5
\end{array}
$$
All higher pieces vanish: every set
of three pentagon vertices contains
an edge.
For the hexagon, independent triples
do exist. Its degree-three basis is
$$
x_1x_3x_5,\qquad x_2x_4x_6.
$$

\medskip\noindent
\textbf{Multiplication adds a vertex.}
Put
$$
L=x_1+\cdots+x_n.
$$
For an independent vertex set $S$,
write $x_S=\prod_{i\in S}x_i$.
Let $E(S)$ denote the vertices outside
$S$ adjacent to no vertex in $S$.
Then
$$
Lx_S=\sum_{j\in E(S)}x_{S\cup\{j\}}.
$$
Thus the multiplication matrix records
containment of independent sets of
successive sizes.

For the pentagon, consider
$$
\times L:A(C_5)_1\longrightarrow A(C_5)_2.
$$
Using the bases in the order above,
its matrix, with row and column labels,
is
$$
\begin{array}{c|ccccc}
 &x_1&x_2&x_3&x_4&x_5\\ \hline
x_1x_3&1&0&1&0&0\\
x_1x_4&1&0&0&1&0\\
x_2x_4&0&1&0&1&0\\
x_2x_5&0&1&0&0&1\\
x_3x_5&0&0&1&0&1
\end{array}
$$
For example, the first column expresses
$$
Lx_1=x_1x_3+x_1x_4.
$$

\medskip\noindent
\textbf{Which incidence matrix?}
The rows correspond to nonedges of $G$.
Therefore this is the unsigned
edge--vertex incidence matrix of the
\emph{complement graph} $\overline G$,
not of $G$.
There are two entries equal to one
in each row.

Here is the precise connected-graph
criterion, over a field of
characteristic different from two.
Let $H$ be connected, with $v$ vertices
and $e$ edges, and let $M_H$ be its
unsigned edge--vertex incidence matrix.
If $H$ is bipartite, then
$$
\operatorname{rank}M_H=v-1.
$$
Otherwise,
$$
\operatorname{rank}M_H=v.
$$
Hence maximal rank holds exactly when
$H$ is a tree or is nonbipartite.
For a connected graph, being a forest
means being a tree.
Nonbipartiteness already excludes trees.

Indeed, a vector $a$ lies in the kernel
exactly when every edge satisfies
$$
a_i+a_j=0.
$$
Connectedness propagates all coordinates
from one value, with alternating signs.
A bipartition permits one free value;
an odd cycle forces that value to zero.
Comparing with $\min(e,v)$ gives the
maximal-rank criterion.

To use this for $A(G)_1\to A(G)_2$,
apply it to $H=\overline G$.
Connectedness of $G$ alone does not
imply connectedness of $\overline G$.
For the pentagon the complement is
another pentagon, so the displayed
matrix is invertible when the
characteristic is not two.

\medskip\noindent
\textbf{Aside: a rational map.}
The surviving quadratic monomials also
define a rational map of projective
spaces. For the pentagon it is
$$
\begin{aligned}
\mathbb P^4&\dashrightarrow\mathbb P^4,\\
[x_1:\cdots:x_5]&\longmapsto\\
[\,x_1x_3:x_1x_4:&x_2x_4:x_2x_5:x_3x_5\,].
\end{aligned}
$$
These are polynomials in $R$, evaluated
where they do not all vanish.
Their exponent matrix is the same
incidence matrix.
The speaker's further use of this map
was not recorded.

\medskip\noindent
\textbf{Next example: the seven-vertex
path.}
Let $G=P_7$, with edges between
consecutive vertices.
An independent set has no consecutive
indices.
The bases recorded in the session are
$$
\begin{aligned}
B_2&=\{x_ix_j:1\leq i<j\leq7,\
j-i\geq2\},\\
B_3&=\{x_ix_jx_k:1\leq i<j<k\leq7,\\
&\hspace{35pt}j-i\geq2,\ k-j\geq2\}.
\end{aligned}
$$
There are fifteen basis pairs and ten
basis triples. Explicitly, the triples
are
$$
\begin{aligned}
B_3=\{&
x_1x_3x_5,x_1x_3x_6,x_1x_3x_7,\\
&x_1x_4x_6,x_1x_4x_7,x_1x_5x_7,\\
&x_2x_4x_6,x_2x_4x_7,x_2x_5x_7,\\
&x_3x_5x_7\}.
\end{aligned}
$$
These bases prepare the next map:
$$
\times L:A(P_7)_2\longrightarrow A(P_7)_3.
$$
Its matrix has ten rows and fifteen
columns. The entry is one precisely
when the column's pair is contained
in the row's triple, and zero otherwise.
For example,
$$
L(x_1x_3)=
x_1x_3x_5+x_1x_3x_6+x_1x_3x_7.
$$
Maximal rank here means surjectivity,
or rank ten.
Unlike the previous map, each row now
has three entries equal to one.
The graph incidence theorem above
does not directly settle this map.

The point of the exercise is to turn
multiplication into combinatorics:
which independent pairs extend to which
independent triples, and do the resulting
columns span the target?
WLP requires maximal rank in every
degree, not only from degree one to two.

\medskip\noindent
\textbf{Why switch to the transpose?}
The goal is to decide WLP from the
combinatorics of $G$.
For monomial algebras over an infinite
field, it suffices to test the sum $L$
of the variables.
Every multiplication map must have
maximal rank.

Write
$$
U_{i-1}:A(G)_{i-1}\longrightarrow A(G)_i
$$
for multiplication by $L$.
In the squarefree monomial bases,
its transpose is the removal map
$$
\begin{aligned}
D_i&:A(G)_i\longrightarrow A(G)_{i-1},\\
D_i(x_S)&=\sum_{j\in S}x_{S\setminus\{j\}}.
\end{aligned}
$$
These bases identify the vector spaces
with their duals for this notation.
For example,
$$
D_2(x_1x_3)=x_1+x_3.
$$
Thus $U$ adds a vertex and $D$ removes
a vertex. The downward map in the live
transcription is a transpose, not
multiplication by a power of $L$.

A nonzero kernel alone need not mean
failure: dimensions may force one.
For the seven-vertex path, however,
$U_2$ goes from dimension fifteen
to dimension ten, so maximal rank is
equivalent to
$$
\ker D_3=0.
$$
A nonzero linear combination of allowed
triples whose removals all cancel
would therefore prove failure.

\medskip\noindent
\textbf{Why Specht polynomials appear.}
Temporarily omit the graph relations:
$$
B=R/(x_1^2,\ldots,x_n^2).
$$
Every squarefree monomial is now allowed.
The removal map has cancellations such as
$$
D_1(x_1-x_2)=1-1=0.
$$
More generally, products of differences
on disjoint pairs are killed by $D$.
In characteristic zero, for
$0\leq i\leq\lfloor n/2\rfloor$,
the kernel in degree $i$ is spanned by
$$
\prod_{r=1}^{i}(x_{a_r}-x_{b_r}),
$$
where all $2i$ indices are distinct.
These are the two-row Specht polynomials
of shape $(n-i,i)$.
Above the middle degree, the kernel is
zero. See
Theorem~\ref{iv-jornadas-lefschetz-2026-theorem-specht-kernels}
for the earlier precise statement.

To return to $A(G)$, seek combinations
in this kernel supported entirely on
independent sets of $G$.
Removing a vertex preserves independence,
so the graph removal map is the
restriction of the one for $B$.
Simply deleting forbidden terms from
one Specht polynomial need not preserve
the cancellation.

The live attribution was heard as
Lefschetz--Bott; it remains unverified.
The recorded shape and degree bound
were incomplete and are replaced above
by the characteristic-zero formulation.

\medskip\noindent
\textbf{Characteristic warning from
the session.}
The speaker mentioned failure when the
characteristic $p$ divides a degree $i$.
The precise degree and dimension
hypotheses were not captured.
Such divisibility can produce extra
cancellation, but it is not recorded
here as a blanket WLP-failure criterion.

\medskip\noindent
\textbf{A cancellation compatible with
the edge ideal.}
Take the graph on four vertices with
$$
I(G)=(x_1x_3,x_2x_4).
$$
It consists of two disjoint edges.
In the algebra
$$
A(G)=
\frac{\mathbb K[x_1,x_2,x_3,x_4]}
{(x_1x_3,x_2x_4,x_1^2,x_2^2,x_3^2,x_4^2)},
$$
consider
$$
\begin{aligned}
F&=x_1x_2-x_2x_3+x_3x_4-x_1x_4\\
 &=(x_1-x_3)(x_2-x_4).
\end{aligned}
$$
Each of its four monomials survives:
none repeats a vertex or contains
either forbidden edge.
Therefore $F\ne0$ in $A(G)_2$.
Under the removal map,
$$
\begin{aligned}
D_2(F)&=(x_1+x_2)-(x_2+x_3)\\
&\quad +(x_3+x_4)-(x_1+x_4)\\
&=0.
\end{aligned}
$$
This is a Specht polynomial whose
cancellation remains valid with the
graph relations imposed.

Both $A(G)_1$ and $A(G)_2$ have
dimension four. Hence this nonzero
transpose-kernel vector proves that
$$
\times L:A(G)_1\longrightarrow A(G)_2
$$
does not have maximal rank.
Over an infinite field, the monomial
criterion for testing $L$ therefore
shows that this algebra fails WLP.

\medskip\noindent
\textbf{Hilbert series and independent
sets.}
Let $s_r(G)$ be the number of independent
vertex sets of size $r$, including
$s_0(G)=1$.
The \emph{independence polynomial} is
$$
I_G(t)=\sum_{r\geq0}s_r(G)t^r.
$$
Here this polynomial notation is
distinct from the edge ideal $I(G)$.
The squarefree monomial basis gives
$$
\begin{aligned}
\operatorname{Hilb}(A(G),t)
&=\sum_{r\geq0}\dim_{\mathbb K}A(G)_r\,t^r\\
&=I_G(t)\\
&=1+nt+|E(\overline G)|t^2+\cdots.
\end{aligned}
$$
Thus these coefficients give the
dimensions of the source and target
of every Lefschetz multiplication map.

\begin{remark}[Hilbert series,
Hilbert polynomial and deformations]
\label{iv-jornadas-lefschetz-2026-remark-session-hilbert-reminder}
For a standard graded algebra $A$,
the \emph{Hilbert function} counts
independent coefficients of homogeneous
polynomials modulo the equations:
$$
h_A(m)=\dim_{\mathbb K}A_m.
$$
The \emph{Hilbert series} packages all
these dimensions into one formal series:
$$
H_A(t)=\sum_{m\geq0}h_A(m)t^m.
$$
The variable $t$ labels the degree.
The \emph{Hilbert polynomial} $P_A(m)$
instead satisfies
$$
P_A(m)=h_A(m)
$$
for all sufficiently large integers $m$.
It may miss the initial values.
For an Artinian graded algebra such as
$A(G)$, all sufficiently high pieces
vanish. Its Hilbert polynomial is zero,
although its Hilbert series is a useful
nonzero polynomial in $t$.

For an SR ring, surviving monomials
have face supports. Counting them gives
$$
H_{\mathbb K[\Delta]}(t)
=\sum_{F\in\Delta}
\left(\frac{t}{1-t}\right)^{|F|}.
$$
Thus face counts determine the series.
Each vertex of a support may occur
repeatedly; killing squares would
instead count only squarefree monomials.

For Daniel's sphere $M$, the face counts
and resulting series are
$$
\begin{aligned}
(f_0,f_1,f_2,f_3)&=(8,28,40,20),\\
H_{\mathbb K[M]}(t)
&=\frac{1+4t+10t^2+4t^3+t^4}{(1-t)^4}\\
&=1+8t+36t^2+104t^3+\cdots.
\end{aligned}
$$
Here $36=8+28$ counts squares and
edge monomials, while
$$
104=8+2\cdot28+40
$$
counts cubes, repeated-variable edge
monomials and triangle monomials.
The Hilbert polynomial is
$$
P(m)=\frac{10}{3}m^3+\frac{14}{3}m.
$$
It agrees with the Hilbert function
for every $m\geq1$, but not for $m=0$.

This gives geometric information directly
from the combinatorics: the projective
scheme has dimension three and degree
twenty. The degree is $3!$ times the
leading coefficient of $P$.

For deformation theory, distinguish two
flatness statements.
A flat projective family, with a fixed
relative polarization, has locally
constant Hilbert polynomial.
Thus any such smoothing has the same
dimension, degree and polynomial.
The entire Hilbert series is preserved
in a flat family of graded algebras
whose graded pieces are finite locally
free over the base: every graded
dimension is then locally constant.
Projective flatness alone does not
automatically preserve all the initial
graded dimensions of coordinate rings.

These are constraints on a possible
smoothing, not a proof that it exists
or a determination of its isomorphism
class.
\end{remark}

\medskip\noindent
\textbf{Counting matchings.}
A \emph{matching} is a set of edges
having no common vertices.
Let $m_r(H)$ count matchings of size $r$
in a graph $H$, with $m_0(H)=1$.
The convention used in the session is
the matching generating polynomial
$$
\begin{aligned}
M_H(t)&=\sum_{r\geq0}m_r(H)t^r\\
&=1+|E(H)|t+m_2(H)t^2+\cdots.
\end{aligned}
$$
This is the unsigned counting convention;
other uses of matching polynomial
include alternating signs.
The session considers $H=G$ or
$H=\overline G$.

For the four-vertex example above,
$$
\begin{aligned}
I_G(t)&=1+4t+4t^2,\\
M_G(t)&=1+2t+t^2,\\
M_{\overline G}(t)&=1+4t+2t^2.
\end{aligned}
$$
Independent sets count allowed monomials;
matchings count disjoint pairs.
The latter pairs can index factors
of a Specht polynomial.
Here the matching with edges
$\{1,3\}$ and $\{2,4\}$ indexes
$(x_1-x_3)(x_2-x_4)$.

A matching alone does not guarantee a
kernel element for $A(G)$:
the expanded polynomial must be supported
on independent sets, or one must find
a linear combination with that property.
No general equality between matching
counts and kernel dimensions was
recorded at this point in the session.

\medskip\noindent
\textbf{Final topic: induced matchings.}

\begin{definition}[Induced matching]
\label{iv-jornadas-lefschetz-2026-definition-session-induced-matching}
A set $S\subseteq E(G)$ is an
\emph{induced matching} if its edges
have disjoint endpoints and the subgraph
on all these endpoints has exactly
the edges of $S$.
Thus no additional edge joins endpoints
belonging to different chosen edges.

The \emph{induced matching number},
denoted $\operatorname{indmat}(G)$ here,
is the maximum number of edges in an
induced matching of $G$.
\end{definition}

This is the usual definition; see
\href{%
https://arxiv.org/abs/2410.11288%
}{Maximal and maximum induced matchings
in connected graphs}.

For the path on vertices $1,\ldots,7$,
the edges $\{1,2\}$ and $\{4,5\}$
form an induced matching.
The edges $\{1,2\}$ and $\{3,4\}$
form a matching but not an induced one:
the extra edge $\{2,3\}$ joins their
endpoints.
Three chosen edges would require six
endpoints and at least two intervening
vertices, so
$$
\operatorname{indmat}(P_7)=2.
$$

\begin{remark}[Neighborhood notation
to check]
\label{iv-jornadas-lefschetz-2026-remark-induced-neighborhoods}
The live transcription gives, for
distinct chosen edges,
$$
N[u,v]\cap N[a,b]=\emptyset.
$$
If $N[u,v]$ means the union of the
closed vertex neighborhoods of $u,v$,
this is stronger than the standard
induced-matching condition.
Two chosen edges may have a common
neighbor outside their endpoints.

For example, in the five-vertex path,
$\{1,2\}$ and $\{4,5\}$ form an induced
matching, but the two closed
neighborhoods both contain vertex $3$.
Using this neighborhood convention,
the standard condition is instead
$$
N[\{u,v\}]\cap\{a,b\}=\emptyset
$$
for distinct edges of the matching.
The speaker's notation should be checked.
\end{remark}

\medskip\noindent
\textbf{Connection with the cancellations.}
An induced matching with edges
$\{a_r,b_r\}$, for $1\leq r\leq s$,
gives the polynomial
$$
F=\prod_{r=1}^{s}(x_{a_r}-x_{b_r}).
$$
Every expanded term chooses one endpoint
from each edge.
Because the matching is induced, these
choices are independent sets.
The terms are distinct surviving
monomials, so $F\ne0$ in $A(G)_s$.
Removing either endpoint in each factor
gives opposite contributions, hence
$$
D_s(F)=0.
$$
This explains why induced matchings
are relevant to the preceding Specht
polynomials.
To conclude WLP failure from this
kernel, one must still compare
the dimensions of successive pieces.
In particular, it proves failure if
$$
\dim A(G)_{s-1}\geq\dim A(G)_s.
$$

\section{Cox--Gorenstein algebras:
Bruzzo's problem session}
\label{iv-jornadas-lefschetz-2026-section-cox-gorenstein-session}

September 28, 2026, morning.
Presented by a collaborator of Bruzzo;
the speaker's name was not recorded.
These are Daniel's live notes.
Incomplete statements are marked below.

\medskip\noindent
\textbf{Goal.}
Investigate structural properties of
Cox--Gorenstein algebras: Artinian
algebras graded by finitely generated
abelian groups, arising from Cox rings
of toric varieties.

\begin{definition}[Toric variety]
\label{iv-jornadas-lefschetz-2026-definition-session-toric-variety}
In the definition recorded in the talk,
a \emph{toric variety} is an irreducible
variety $X$ containing a torus
$$
T\cong(\mathbb{C}^{*})^d
$$
as a Zariski-open subset, such that
the action of $T$ on itself extends
to an action on $X$.
\end{definition}

The fan description below uses the
usual normal toric varieties.

\medskip\noindent
\textbf{Cox ring.}
For a complete toric variety $X$, the
speaker defined its Cox ring using a
direct sum. The full direct-sum formula
was not captured in the live notes.
For $X=\mathbb{P}_{\Sigma}$, the
description recorded was
$$
\operatorname{Cox}(X)
=\mathbb{K}[x_{\eta}\mid
\eta\in\Sigma(1)].
$$
There is one variable for each ray of
the fan. The grading group is the
finitely generated abelian group
$\operatorname{Cl}(X)$.
Here $\Sigma(1)$ denotes the rays.

\medskip\noindent
\textbf{Cox--Gorenstein and duality.}
Let $G$ be a finitely generated abelian
group and let
$$
A=\bigoplus_{g\in G}A_g
$$
be an Artinian $G$-graded algebra of
finite type.

\begin{definition}[As recorded]
\label{iv-jornadas-lefschetz-2026-definition-session-cox-gorenstein}
The algebra $A$ is
\emph{Cox--Gorenstein} if there is
$\rho\in G$ such that
$$
\operatorname{Soc}(A)
=(0:\mathfrak{m})=A_{\rho}
\cong\mathbb{K}.
$$
The \emph{socle} consists of elements
annihilated by $\mathfrak{m}$.
\end{definition}

The precise hypotheses specifying
$\mathfrak{m}$ in the graded setting
remain to be copied from the slides.

\medskip\noindent
\textbf{Poincaré duality.}
The talk's formulation asks for a
maximal degree $\rho\in G$ with
$A_{\rho}\cong\mathbb{K}$ and perfect
multiplication pairings
$$
A_g\times A_{\rho-g}
\longrightarrow A_{\rho}.
$$
A pairing is perfect if it identifies
each factor with the vector-space
dual of the other.
The convention for maximality in $G$
was not recorded.

\medskip\noindent
\textbf{Theorem stated in the talk.}
\par\nopagebreak\noindent
Under the speaker's hypotheses,
$A$ is Cox--Gorenstein if and only if
it has Poincaré duality.
The missing grading conventions above
must be supplied for a full statement.

\medskip\noindent
\textbf{Toric Macaulay--Matlis duality.}
\par\nopagebreak\noindent
A bijection was presented between
$G$-homogeneous ideals $I\subseteq Q$
and $G$-graded $Q$-submodules
$M\subseteq S$, where
$$
S=\mathbb{K}[x_1,\ldots,x_s],
\qquad
Q=\mathbb{K}[X_1,\ldots,X_s].
$$
The module action, grading conventions,
and full hypotheses were not captured.

\medskip\noindent
\textbf{Toric quotient and critical
degree.}
Let $X$ be a complete simplicial toric
variety. The notes record
$$
f_0,\ldots,f_n\in B(\Sigma),
\qquad \deg(f_i)=\alpha_i,
$$
with no common zero on $X$, and
$\beta$ the anticanonical class.
The formula for the critical degree
$\rho$ was only partially recorded,
as a sum. A conjecture was mentioned;
its precise statement remains to be
copied. Problem 4 below refers to the
codimension-one conjecture.

\medskip\noindent
\textbf{Toric Lefschetz properties.}
The setup recorded was
$$
A=Q/I=\bigoplus_{g\in G}A_g,
\qquad
L=\langle X_1,\ldots,X_n\rangle
\subseteq A,
\qquad L_g=L\cap A_g.
$$
Here the brackets denote the linear
span of the images of the variables.
The speaker introduced the toric strong
Lefschetz property (TSLP) and toric weak
Lefschetz property (TWLP).
Their precise definitions were not
captured. The key change is that
variables now have group-valued degrees.

\medskip\noindent
\textbf{Problem 1: Cox--Gorenstein
monomial algebras.}
The classical starting point is that
Artinian Gorenstein monomial algebras
are complete intersections of the form
$$
A=\frac{\mathbb{K}[x_1,\ldots,x_n]}
{(x_1^{a_1},\ldots,x_n^{a_n})}.
$$
With the standard grading, these have
SLP in characteristic zero.

Assume $\operatorname{char}\mathbb{K}=0$,
let $G$ be finitely generated abelian,
and give
$$
S=\mathbb{K}[x_1,\ldots,x_n]
$$
a $G$-grading.
Does
$$
A=S/(x_1^{a_1},\ldots,x_n^{a_n})
$$
have TSLP or TWLP?

\medskip\noindent
\textbf{Problem 2: codimension two.}
In characteristic zero, every standard
graded Artinian Gorenstein algebra of
codimension two satisfies SLP.
What can be said for codimension-two
Cox--Gorenstein algebras with a
$G$-grading?

\medskip\noindent
\textbf{Problem 3: higher Jacobians.}
Can a higher Jacobian-type criterion
be established in the $G$-graded
setting? Investigate the resulting
Jordan types.

\medskip\noindent
\textbf{Problem 4: artinization in the
critical degree.}
\par\nopagebreak\noindent
In the Cattani--Cox--Dickenstein
setting, let
$$
A=S/(f_0,\ldots,f_n).
$$
Assume that the codimension-one
conjecture holds for $A$.
Let $B$ be the natural artinization
obtained by truncating at the critical
degree. Is $B$ Cox--Gorenstein?

The exact truncation defining $B$ and
the conjecture's full statement remain
to be supplied from the slides.

\section{Marcos Jardim: exercise sessions}
\label{iv-jornadas-lefschetz-2026-section-jardim-problems}

September 28--October 1, 2026.

\medskip\noindent
\label{iv-jornadas-lefschetz-2026-paragraph-jardim-sections}
\textbf{Sections, Rao modules and linkage}
\par\nopagebreak\noindent
Can two curves be joined by a chain
of links? The exercises ask for more:
can they come from sections of twists
of one common rank-two bundle?

Work on $X=\mathbb{P}^3_{\mathbb{C}}$.
Write $\mathcal{O}(a)=\mathcal{O}_X(a)$,
$F(a)=F\otimes\mathcal{O}(a)$ and
$H^i_*(F)=\bigoplus_t H^i(X,F(t))$.

For a rank-two bundle $E$ with
$\det E=\mathcal{O}(c)$, a section
$s\in H^0(E(a))$ whose zero scheme is
a curve $C$ of pure codimension two gives
\begin{equation}
0\longrightarrow\mathcal{O}(-a)
\xrightarrow{s}E
\longrightarrow\mathcal{I}_C(c+a)
\longrightarrow0.
\label{iv-jornadas-lefschetz-2026-equation-jardim-section}
\end{equation}
Locally $s=(f,g)$ and $C=V(f,g)$.
Two equations of codimension two in
the smooth space $X$ form a regular
sequence: thus $C$ is locally a complete
intersection. Consequently the map
$(u,v)\mapsto fv-gu$ has image $(f,g)$
and kernel generated by $s$.
Determinants give the twist $c+a$.
The inclusion is \emph{saturated}:
its quotient is torsion-free.

\textbf{Which curves arise this way?}
Exactly the locally complete intersection
curves with $\omega_C\cong\mathcal{O}_C(m)$
for some integer $m$ (\emph{subcanonical}
curves), by Hartshorne--Serre
\hyperref[item-jardim-serre]{[1]}.
Even the smooth twisted cubic fails:
$\deg\omega_C=-2$, whereas
$\deg\mathcal{O}_C(m)=3m$.

The \emph{Rao module} of a locally
Cohen--Macaulay curve is
$$
M_C=H^1_*(\mathcal{I}_C).
$$
It records sections on $C$ that do not
extend to ambient polynomials, together
with multiplication by polynomials.
The curve is arithmetically
Cohen--Macaulay precisely when $M_C=0$.

\textbf{Rao's theorem.}
For \emph{all locally Cohen--Macaulay
curves} in $X$, isomorphism of Rao modules
up to shift is equivalent to belonging
to the same \emph{even liaison class}
\hyperref[item-jardim-rao]{[2]}.
A liaison is a chain of
complete-intersection links; ``even''
counts the links ($2,4,\ldots$), not
the rank of a bundle.
No presentation by a rank-two bundle
is required by the theorem.

\textbf{What a common bundle gives.}
Since $H^1(\mathcal{O}(m))=
H^2(\mathcal{O}(m))=0$ for every $m$,
two sequences \eqref{iv-jornadas-lefschetz-2026-equation-jardim-section}
with the same $E$ give
$$
M_{C_1}[c+a_1]\cong H^1_*(E)
\cong M_{C_2}[c+a_2],
\qquad N[e]_t=N_{t+e}.
$$
Thus a common bundle implies even
liaison. The converse fails: a line
and a smooth conic already give a
counterexample; see the
\hyperref[paragraph-jardim-line-conic]
{line--conic calculation below}.

\textbf{Why introduce sections?}
They can produce the geometry of links,
not just an existence criterion.
For sections independent at the
generic point,
put $A_i=\mathcal{O}(-a_i)$ and
$T=E/(A_1+A_2)$, identifying each $A_i$
with its image. The determinant defines
a surface $S=Z(s_1\wedge s_2)$ containing
both curves. Quotienting in either order
gives
$$
\mathcal{I}_{C_1/S}(c+a_1)
\cong T\cong
\mathcal{I}_{C_2/S}(c+a_2).
$$
If $S$ is smooth, this says
$C_2\sim C_1+(a_2-a_1)H_S$, where
$H_S$ is its hyperplane class.
Choose a general $D\sim mH_S-C_1$
for $m$ large. Then both $C_1+D$ and
$C_2+D$ are hypersurface sections of
$S$: this gives $C_1\leftrightarrow
D\leftrightarrow C_2$, two links.
For $S=V(F)$ and $C_1\cup D=V(F,G)$,
the residual ideal is
$I_D=(F,G):I_{C_1}$.

\medskip\noindent
\label{iv-jornadas-lefschetz-2026-paragraph-jardim-line-conic}
\textbf{A line and a smooth conic}
\par\nopagebreak\noindent
Let $L$ be a line and $Q$ a smooth conic.
Their ideal resolutions are
$$
\begin{aligned}
0&\to\mathcal{O}(-2)
\to\mathcal{O}(-1)^{\oplus2}
\to\mathcal{I}_L\to0,\\
0&\to\mathcal{O}(-3)
\to\mathcal{O}(-1)\oplus\mathcal{O}(-2)
\to\mathcal{I}_Q\to0.
\end{aligned}
$$
Hence both Rao modules vanish: they
are in the same even liaison class.
Individually they are zero schemes of
sections of $\mathcal{O}(1)^{\oplus2}$
and $\mathcal{O}(1)\oplus\mathcal{O}(2)$.

\textbf{A common rank-two bundle is
impossible.}
Adjunction for a regular section gives
$\omega_C\cong\mathcal{O}_C(c+2a-4)$.
On $L$ this forces $c+2a_1=2$;
on $Q$ it forces $c+2a_2=3$.
Their difference would be
$2(a_2-a_1)=1$.
The same Rao module does not enforce
the determinant needed by a common bundle.

\textbf{A common reflexive sheaf forces
coplanarity.}
Suppose instead that a rank-two
reflexive $E$ fits into
$$
\begin{aligned}
0&\to\mathcal{O}(-a_1)\xrightarrow{s_1}E
\to\mathcal{I}_L(c+a_1)\to0,\\
0&\to\mathcal{O}(-a_2)\xrightarrow{s_2}E
\to\mathcal{I}_Q(c+a_2)\to0.
\end{aligned}
$$
Put $h=c_1(\mathcal{O}(1))$.
The resolutions above give
$$
\begin{aligned}
c(\mathcal{I}_L(t))
&=1+th+h^2+(2-t)h^3,\\
c(\mathcal{I}_Q(t))
&=1+th+2h^2+(6-2t)h^3.
\end{aligned}
$$
Multiply by $1-a_i h$ and compare
the $h^2$ coefficients for $E$:
$$
1-a_1(c+a_1)=2-a_2(c+a_2),
\qquad
1=(a_2-a_1)(c+a_1+a_2).
$$
Meanwhile
$$
s_1\wedge s_2\in H^0(\mathcal{O}(d)),
\qquad d=c+a_1+a_2.
$$
\textbf{The determinant is nonzero.}
The integer equation gives $a_1\ne a_2$;
choose $i,j$ with $a_i<a_j$.
If $s_1\wedge s_2=0$, the composite
$\mathcal{O}(-a_i)\to E\to
\mathcal{I}_{C_j}(c+a_j)$, with
$C_1=L$, $C_2=Q$, vanishes generically,
hence everywhere because its target
is torsion-free. Thus $s_i$ factors
through $\mathcal{O}(-a_j)$. But
$$
\operatorname{Hom}(\mathcal{O}(-a_i),
\mathcal{O}(-a_j))
=H^0(\mathcal{O}(a_i-a_j))=0:
$$
there are no polynomials of negative
degree. This would force $s_i=0$,
a contradiction.
Thus $d\geq0$, and the integer equation
forces $d=1$ and $a_2-a_1=1$.
Composing either section with the other
quotient gives this same determinant
in $H^0(\mathcal{I}_L(1))$ and
$H^0(\mathcal{I}_Q(1))$, up to sign.
Therefore
$$
\boxed{L\cup Q\subset H=Z(s_1\wedge s_2).}
$$
The remaining Chern coefficients give
$c=-2a_1$, $c_2(E)=(1+a_1^2)h^2$ and
$c_3(E)=2h^3$. The last equality again
rules out local freeness. These are
necessary conditions, not an existence
proof for the common sheaf.

\textbf{How many extensions are available?}
For $C=L,Q$, the
\hyperref[paragraph-jardim-sheaves]
{reminder below} gives
$$
\operatorname{Ext}^1
(\mathcal{I}_C(c+a),\mathcal{O}(-a))
\cong H^1(\mathcal{O}_C(c+2a-4))^\vee.
$$
The dimensions are, respectively,
$$
\max\{0,3-c-2a_1\},\qquad
\max\{0,7-2c-4a_2\}.
$$
Here $L,Q\cong\mathbb{P}^1$, but
$\mathcal{O}_Q(1)$ has degree two.
Nonzero extension classes need not give
locally free middle sheaves.

\medskip\noindent
\label{iv-jornadas-lefschetz-2026-paragraph-jardim-sheaves}
\textbf{Reminder on sheaves and Ext}
\par\nopagebreak\noindent
All sheaves below are coherent.
Rank and torsion are taken over an
integral variety. A sheaf is
\emph{locally free} if locally it is
$\mathcal{O}^{\oplus r}$; these are
exactly sheaves of sections of vector
bundles. Rank is dimension at the
generic point, and alone says nothing
about local freeness everywhere.

\textbf{Torsion and reflexivity.}
Over a domain, $m$ is torsion if
$fm=0$ for some nonzero scalar $f$.
Torsion-free means this forces $m=0$.
For instance $R/(x)$ has torsion over
$R=\mathbb{C}[x]$, although it has no
nonzero nilpotents.
Write $F^*=\mathcal{H}om(F,\mathcal{O}_X)$.
The canonical evaluation map is
$$
F\longrightarrow F^{**},\qquad
s\longmapsto(\varphi\mapsto\varphi(s)).
$$
It is injective precisely when $F$ is
torsion-free, and an isomorphism
precisely when $F$ is \emph{reflexive}
\hyperref[item-jardim-stacks]{[4]}.
Locally free implies reflexive, which
implies torsion-free. For example,
$\mathcal{I}_L^{**}=\mathcal{O}\ne
\mathcal{I}_L$ for a line in
$\mathbb{P}^3$.

\textbf{Twisting and exactness.}
We have $F(a)(b)=F(a+b)$ and
$\operatorname{Hom}(\mathcal{O}(-a),F)
=H^0(F(a))$.
Tensoring preserves every short exact
sequence precisely for a \emph{flat}
sheaf. Locally free sheaves are flat:
if $F\cong\mathcal{O}^{\oplus r}$, then
$$
A\otimes F\cong A^{\oplus r},\qquad
a\otimes e_i\longmapsto
(0,\ldots,a,\ldots,0).
$$
Thus tensoring locally makes $r$ copies
of the sequence. In particular, twisting
by a line bundle is exact.
For arbitrary $F$, tensoring is only
right exact: the first injection may
be lost.

\textbf{Global Ext.}
$\operatorname{Ext}^1(B,A)$ classifies
sequences $0\to A\to E\to B\to0$
with end terms fixed; zero is the split
extension. Serre duality on
$\mathbb{P}^3$ gives
$$
\begin{aligned}
\operatorname{Ext}^1
(\mathcal{I}_C(c+a),\mathcal{O}(-a))
&\cong\operatorname{Ext}^2
(\mathcal{O}(-a),\mathcal{I}_C(c+a-4))^\vee\\
&\cong H^2(\mathcal{I}_C(c+2a-4))^\vee\\
&\cong H^1(\mathcal{O}_C(c+2a-4))^\vee.
\end{aligned}
$$
The second step uses the locally free
first argument; the last uses the ideal
sequence and intermediate vanishing
for $\mathcal{O}(m)$.

\textbf{Sheaf Ext: local freeness and
reflexivity.}
Unlike global Ext, sheaf Ext records
local modules:
$$
\mathcal{E}xt^i(F,\mathcal{O})_x
\cong\operatorname{Ext}^i_{\mathcal{O}_{X,x}}
(F_x,\mathcal{O}_{X,x}).
$$
For torsion-free $F$ on a smooth
projective threefold, put
$\mathcal{E}^i=\mathcal{E}xt^i
(F,\mathcal{O})$.
Then \hyperref[item-jardim-hl]
{[3, Section 1.1]}:
\begin{itemize}
\item $F$ is locally free exactly when
$\mathcal{E}^1=\mathcal{E}^2=0$.
\item $F$ is reflexive exactly when
$\mathcal{E}^2=0$ and $\mathcal{E}^1$
has finite support.
\end{itemize}
Thus for reflexive $F$, vanishing of
$\mathcal{E}^1$ tests local freeness.
For a general torsion-free sheaf one
must also check second Ext:
$\mathcal{E}xt^1
(\mathcal{I}_p,\mathcal{O})=0$
but
$\mathcal{E}xt^2(\mathcal{I}_p,\mathcal{O})
\cong\mathcal{O}_p$.

\medskip\noindent
\textbf{References.}
\begin{enumerate}
\raggedright
\item\label{iv-jornadas-lefschetz-2026-item-jardim-serre}
E. Arrondo,
\emph{A home-made Hartshorne--Serre
correspondence}, Theorem 1,
\href{https://arxiv.org/abs/math/0610015}
{arXiv:math/0610015}.
\item\label{iv-jornadas-lefschetz-2026-item-jardim-rao}
R. Hartshorne,
\emph{On Rao's Theorems and the
Lazarsfeld--Rao Property},
\href{https://arxiv.org/abs/math/0302078}
{arXiv:math/0302078}.
\item\label{iv-jornadas-lefschetz-2026-item-jardim-hl}
D. Huybrechts and M. Lehn,
\emph{The Geometry of Moduli Spaces
of Sheaves}, Section 1.1.
\item\label{iv-jornadas-lefschetz-2026-item-jardim-stacks}
The Stacks Project,
\href{%
https://stacks.math.columbia.edu/tag/0AVT%
}
{Tag 0AVT}: reflexive sheaves.
\end{enumerate}

\section{Compressed Artinian Gorenstein
algebras}
\label{iv-jornadas-lefschetz-2026-section-lenin-bezerra}

Lenin Alexandre de Almeida Bezerra (UPE).

Segunda 28/09, 11:40.

\medskip\noindent
\textbf{Notes.}
To be added during the talk.

\section{Expect the unexpected: probabilistic
failure of WLP in monomial algebras}
\label{iv-jornadas-lefschetz-2026-section-thiago-holleben}

Thiago Holleben (Max Planck Institute).

Segunda 28/09, 14:30.

\medskip\noindent
\textbf{Upshot.}
Thiago studies which combinatorial
properties determine the failure of a
Lefschetz property.

\medskip\noindent
\textbf{History and motivation.}
Stanley's work around 1980 connects
Lefschetz properties with inequalities
for Hilbert functions and face numbers.
For a standard graded Artinian algebra,
WLP implies a unimodal Hilbert function:
it increases weakly and then decreases
weakly. SLP imposes stronger conditions.

The historical list follows the talk;
dates and attributions not checked
against the slides remain provisional.
\begin{enumerate}
\item Watanabe (1987, as recorded).
The guiding mantra is
\emph{tudo tem Lefschetz}.
Lefschetz rank conditions are open in
families with fixed Hilbert function.
Failure is closed in that setting.
An open locus is dense on an irreducible
component if it is nonempty there.
Thus openness alone does not prove that
most Gorenstein algebras have SLP.

The basic characteristic-zero example is
$$
\mathbb{Q}[x_1,\ldots,x_n]/
(x_1^{a_1},\ldots,x_n^{a_n}),
$$
which has SLP.
\item Reid--Roberts--Roitman (1991,
as recorded).
Does every Artinian complete intersection
in characteristic zero have WLP or SLP?
\item Migliore--Miró-Roig (2003,
as recorded).
For a polynomial ring of dimension $n$,
how many generators can an Artinian ideal
have while SLP is still guaranteed?
The notes called the threshold $A(n)$.
The intended extremum and hypotheses
need checking: \emph{smallest} with
\emph{at most} does not express a sharp
threshold without more conditions.
\item Not recorded.
\item Not recorded.
\item For an Artinian monomial algebra,
it suffices to test
$$
L=x_1+\cdots+x_n
$$
for WLP or SLP over an infinite field.
The precise paper attribution is missing.
Diagonal changes of variables identify
linear forms with all coefficients
nonzero, and preserve a monomial ideal.
\end{enumerate}

\medskip\noindent
\textbf{Three roles for WLP.}
\begin{itemize}
\item Algebra: a tool for other questions,
including regularity.
\item Geometry: failure often accompanies
interesting phenomena, such as Togliatti
systems, Perazzo forms and unexpected
hypersurfaces.
\item Combinatorics: existence translates
into inequalities for sequences, often
unknown beforehand. Examples include
$g$-theorems and Kähler packages.
\end{itemize}

The main question of the talk is:
\emph{What is the combinatorial meaning
of failure of WLP?}
Two guiding questions:
\begin{enumerate}
\item Should failure be expected in a
combinatorial setting?
\item What unexpected combinatorial
behaviour does failure explain?
\end{enumerate}

\medskip\noindent
\textbf{Question 1: probability.}
Nagel (2023, as recorded): given a
probability space $\mathcal P$ of
monomial algebras, what is the probability
that an algebra fails WLP?
The passage quoted during the talk:
\begin{quote}
Although the WLP appears to be present
with high probability most of the time,
we propose in Conjecture 3.9 that there
is one scenario where the produced
algebras fail to have the WLP with
high probability.
\end{quote}
The precise reference and probability
model should be added from the slides.

\medskip\noindent
\textbf{Question 1: simplicial complexes.}
Let $\Delta$ have vertex set $[n]$.
Its \emph{Stanley--Reisner ideal} is
$$
I_\Delta=(x_F:F\notin\Delta),
\qquad
x_F=\prod_{i\in F}x_i.
$$
Its \emph{face numbers} are
$$
f_i=\#\{F\in\Delta:|F|=i+1\}.
$$

\begin{theorem}[Holleben, 2026]
\label{iv-jornadas-lefschetz-2026-theorem-holleben-homology-wlp}
Suppose $\dim\Delta=d>0$ and
$$
\widetilde H_d(\Delta;\mathbb K)\ne0,
\qquad f_{d-1}\geq f_d.
$$
Then the Artinian algebra
$$
A_\Delta=
\frac{\mathbb K[x_1,\ldots,x_n]}
{I_\Delta+
(x_1^{d+2},\ldots,x_n^{d+2})}
$$
fails WLP.
\end{theorem}

The live transcription used a strict
inequality. Theorem 1.2 of
\href{%
https://arxiv.org/abs/2501.12108%
}{Coinvariant stresses, Lefschetz
properties and random complexes}
allows equality.
Since $d>0$, ordinary and reduced
homology agree in degree $d$.

For the boundary of a simplex,
$n=d+2$, and the algebra is
$$
\frac{\mathbb K[x_1,\ldots,x_n]}
{(x_1\cdots x_n,x_1^n,\ldots,x_n^n)}.
$$
This recovers the known failure for
this family of monomial almost complete
intersections; it is not a claim about
every monomial almost complete
intersection.

\begin{remark}[Why add the powers?]
\label{iv-jornadas-lefschetz-2026-remark-holleben-talk-powers}
The Stanley--Reisner relations forbid
certain combinations of vertices.
They usually leave arbitrarily high
powers of each individual variable.
Adding $x_i^{d+2}=0$ puts a ceiling on
each exponent: surviving monomials have
every exponent at most $d+1$.
Only finitely many monomials remain.
The quotient is therefore Artinian,
so the usual Artinian WLP question
applies. The particular exponent $d+2$
is part of this theorem; merely making
the algebra finite does not force
failure of WLP.
\end{remark}

\medskip\noindent
\textbf{Random topology saves the day.}
The probability model is
\emph{Linial--Meshulam}:
$Y_d(n,p)$ contains the full
$(d-1)$-skeleton, and each possible
$d$-face is included independently
with probability $p$.

For every fixed $d>0$, there is
$$
c_d<d+1
$$
such that, whenever
$$
c_d<c<d+1,\qquad p=c/n,
$$
the probability that $A_\Delta$
fails WLP tends to one as $n$ tends
to infinity, for
$\Delta\in Y_d(n,c/n)$.
This is Theorem 1.4(2) of the paper above.
Write $\mathcal F_{n,d}$ for the event
of failure. Then
$$
\lim_{n\to\infty}
\mathbb P(\mathcal F_{n,d})=1.
$$

The talk reports numerical experiments
placing $c_d$ in a very narrow interval.
The numerical bounds were not recorded.

\medskip\noindent
\textbf{The everything-fails conjecture.}
The talk suggests a much broader
asymptotic predominance of failure when
the number of variables and relevant
degrees become large.
The proposed ratio has the schematic
form
$$
\frac{\#\mathcal F_{n,D}}
{\#\mathcal A_{n,D}}
\longrightarrow1.
$$
Here $\mathcal A_{n,D}$ would be a
specified finite family of Artinian
monomial algebras, and
$\mathcal F_{n,D}$ its failing members.
These symbols are placeholders:
the counting family, degree bounds and
order of limits were not recorded.
There is no uniform counting probability
on all Artinian monomial algebras
without specifying such a model.
This conjecture must be distinguished
from the proved random-complex theorem.

\medskip\noindent
\textbf{Question 2: inverse systems.}
For clarity, use characteristic zero
and the differential convention.
Let
$$
R=\mathbb K[x_1,\ldots,x_n],
\qquad S=\mathbb K[y_1,\ldots,y_n].
$$
The action is
$$
f\circ F=
f(\partial_{y_1},\ldots,\partial_{y_n})F.
$$
The \emph{inverse system} of $I\subset R$
is the $R$-submodule
$$
I^\perp=
\{F\in S:f\circ F=0
\ \forall f\in I\}.
$$
Conversely, for an $R$-submodule
$M\subset S$, its annihilator is
$$
\operatorname{Ann}_R(M)=
\{f\in R:f\circ F=0
\ \forall F\in M\}.
$$
Thus the inverse system lives in $S$;
its annihilator is an ideal of $R$.
The speaker's exact action convention
was not transcribed.

\medskip\noindent
\textbf{Unexpected hypersurfaces.}
Let $X\subset\mathbb P^N$ be a finite
set of points and let $P$ be general.
A degree-$q$ hypersurface through $X$
with multiplicity at least $m$ at $P$
is \emph{unexpected} when
$$
\begin{aligned}
\dim_{\mathbb K}(I_X\cap I_P^m)_q
&>\max\left\{0,\,
\dim_{\mathbb K}(I_X)_q
-\binom{N+m-1}{N}\right\}.
\end{aligned}
$$
The binomial coefficient is the expected
number of conditions at the fat point.
Unexpectedness means more solutions
than this count predicts; equality
would be the expected behaviour.

\medskip\noindent
\textbf{Spheres and Gorenstein rings.}
A triangulation of a sphere has a
Gorenstein Stanley--Reisner ring.
Quotienting by a homogeneous system
of parameters gives an Artinian
Gorenstein algebra.
This does not assert that adjoining
all the individual variable powers
preserves the Gorenstein property.

\medskip\noindent
\textbf{Unexpected parameters and
balanced spheres.}
A $d$-dimensional complex is
\emph{balanced} if its graph admits
a proper coloring with $d+1$ colors.

\begin{theorem}[Holleben, 2026]
\label{iv-jornadas-lefschetz-2026-theorem-holleben-balanced-sop}
Let $\Delta$ be a simplicial sphere
of dimension $d>0$, over a field
of characteristic zero.
It admits a $2$-unexpected system of
parameters if and only if it is balanced.
Every such system is linear and gives
the same ideal, together with
$I_\Delta$, as the colored system
of parameters.
\end{theorem}

For a proper coloring
$$
\rho:[n]\longrightarrow[d+1],
$$
the \emph{colored system of parameters}
is given, for $1\leq j\leq d+1$, by
$$
\theta_j=\sum_{\rho(v)=j}x_v.
$$
Uniqueness here means equality of ideals,
not equality of ordered lists of forms.
The technical definition of
$2$-unexpected remains to be transcribed.

The sphere hypothesis and the formulation
of uniqueness are supplied from
Theorem 4.5 and Corollary 4.6 of
\href{%
https://arxiv.org/abs/2510.10884%
}{From points to complexes:
A concept of unexpectedness for
simplicial complexes}.

\section{Milnor Algebras, Higher Jacobians
and
Lefschetz-Type Maps}
\label{iv-jornadas-lefschetz-2026-section-giovanna-ilardi}

Giovanna Ilardi (Università di Napoli).

Terça 29/09, 9:30.

\medskip\noindent
\textbf{Notes.}
\begin{quote}
Lefschetz-type maps are controlled by
geometry attached to the gradient system
\end{quote}

\begin{remark}[AI note: gradient system]
\label{iv-jornadas-lefschetz-2026-remark-ai-gradient-system}
For a homogeneous polynomial $f$ of
degree $d\geq2$, put
$$
R=\mathbb{C}[x_0,\ldots,x_n],
\qquad f_i=\frac{\partial f}{\partial x_i}.
$$
The \emph{gradient system} is the linear
system spanned by these first partial
derivatives: the family of hypersurfaces
$$
V\left(\sum_{i=0}^n a_i f_i\right)
\subseteq\mathbb{P}^n,
\qquad a_i\in\mathbb{C},
$$
where the sum is a nonzero polynomial,
considered up to multiplication by a
nonzero scalar. Each member has degree
$d-1$.

These same derivatives define the
\emph{gradient map}, also called the
\emph{polar map}:
$$
\begin{aligned}
\Phi_f:\mathbb{P}^n
&\dashrightarrow(\mathbb{P}^n)^*,\\
p&\longmapsto[f_0(p):\cdots:f_n(p)].
\end{aligned}
$$
It is defined wherever the derivatives
do not all vanish. For reduced $f$,
their common zero scheme is the singular
scheme of $V(f)$; at a smooth point of
$V(f)$, the image represents its tangent
hyperplane.

Algebraically, the derivatives generate
the \emph{Jacobian ideal}, and its
quotient is the \emph{Milnor algebra}:
$$
J(f)=(f_0,\ldots,f_n),
\qquad M(f)=R/J(f).
$$
Thus the gradient system gives geometric
objects attached to the same derivatives
that define the algebra in the talk.
This is background terminology, following
Dimca--Gondim--Ilardi,
\href{https://arxiv.org/abs/1902.09146}
{Higher order Jacobians, Hessians and
Milnor algebras}, introduction.
\end{remark}

\section{Homaloidal hypersurfaces and a
product
formula for gradient maps}
\label{iv-jornadas-lefschetz-2026-section-nivaldo-junior}

Nivaldo Nunes de Medeiros Júnior (UFF).

Terça 29/09, 10:30.

\medskip\noindent
\textbf{Abstract (paraphrase).}
The talk presents two constructions of
homaloidal hypersurfaces. One uses
Chern--Schwartz--MacPherson (CSM) classes
to compute gradient multidegrees
for a product of homogeneous polynomials
in disjoint variable sets. The other
associates to each Cremona map of
$\mathbb{P}^n$ a homaloidal hypersurface
of Perazzo type in $\mathbb{P}^{2n+1}$.
Applications produce codimension-one
foliations with birational Gauss maps,
extending Fassarella--Pereira's
construction. Joint work in progress with
Thiago Fassarella.
Source: the abstract linked from the
\href{%
https://jornadaslp.github.io/%
jornadas-lp/programacao.html%
}{official programme}.

$$
\xymatrix@C=1.5em@R=2em{
&\text{\shortstack{Homaloidal\\
hypersurfaces}}
\ar@{<->}[dl]\ar@{<->}[dr]&\\
\text{Cremona maps}\ar[rr]&&
\text{\shortstack{``Cremona''\\foliations}}
}
$$
This is the blackboard diagram as
recorded by Daniel: the bottom arrow has
no reverse arrow. It records constructions
discussed in the talk, with changes of
ambient dimension allowed.

\medskip\noindent
\textbf{Status of these notes.}
This is an incomplete record, based on
the six photographs in
\texttt{IMG\_4003.pdf} and Daniel's
blackboard notes. Many slides and parts
of the explanations were missed,
including the argument immediately before
the closing sentence. The material below
follows the photographs unless labelled
as an AI note or reconstruction.

\medskip\noindent
\textbf{Gradient maps.}
Let $f$ be a nonzero homogeneous
polynomial of positive degree $d$ in
$m+1$ variables over $\mathbb{C}$, and set
$$
X=V(f)\subset\mathbb{P}^m.
$$
Its \emph{gradient map}, also called the
\emph{polar map}, is
$$
\begin{aligned}
\nabla f:\mathbb{P}^m
&\dashrightarrow(\mathbb{P}^m)^\vee,\\
p&\longmapsto
\left[
\frac{\partial f}{\partial x_0}(p):
\cdots:
\frac{\partial f}{\partial x_m}(p)
\right].
\end{aligned}
$$
Coordinates identify the dual projective
space with another copy of
$\mathbb{P}^m$. A \emph{Cremona map} is
a birational map between these projective
spaces: it has a rational inverse.
The polynomial $f$, or its hypersurface,
is \emph{homaloidal} when $\nabla f$ is
birational.

For $d\geq2$, the gradient map fails to
be dominant precisely when the Hessian
determinant vanishes identically:
$$
\det\left(
\frac{\partial^2f}
{\partial x_i\partial x_j}
\right)_{0\leq i,j\leq m}=0.
$$
Here dominant means that the image is
Zariski dense. The degree restriction
separates off the linear example below.

For a reduced hypersurface, restricting
the gradient map to its smooth locus
gives its Gauss map: a point is sent to
its projective tangent hyperplane.
The dual variety $X^\vee$ is the closure
of these tangent hyperplanes.
$$
\xymatrix{
\mathbb{P}^m
\ar@{-->}[r]^{\nabla f}&
(\mathbb{P}^m)^\vee\\
X\ar@{^{(}->}[u]
\ar@{-->}[r]^{\gamma_X}&
X^\vee\ar@{^{(}->}[u]
}
$$

\medskip\noindent
\textbf{AI note: multidegrees.}
To read the next slide, let
$\phi:\mathbb{P}^m\dashrightarrow
\mathbb{P}^m$ be a rational map.
Its $i$th projective degree $d_i(\phi)$
counts, with multiplicities, points of a
general $i$-plane in the source that map
into a general codimension-$i$ linear
space in the target, away from the
indeterminacy locus. More precisely this
intersection is taken on the closure of
the graph. We use the slide convention
$$
p_\phi(t)=\sum_{i=0}^m d_i(\phi)t^i.
$$
Thus $d_0=1$, whereas $d_m$ is the
number of preimages of a general point
when the map is dominant, and is zero
otherwise. Birationality means $d_m=1$.
This topological degree differs from the
degree of the coordinate polynomials.
For a Cremona map, its \emph{type}
records the polynomial degrees of the
map and its inverse, after common factors
have been cancelled.

\medskip\noindent
\textbf{Easy examples.}
\begin{enumerate}
\item If $f$ is linear, its gradient map
is constant, and $p_{\nabla f}(t)=1$.
\item A smooth quadric has an invertible
linear gradient map. Consequently
$$
p_{\nabla f}(t)=1+t+\cdots+t^m.
$$
\item The coordinate hyperplane
arrangement is defined by
$$
f=x_0\cdots x_m.
$$
Its polar map is the standard Cremona
transformation. On the open set where
all coordinates are nonzero it is
$$
[x_0:\cdots:x_m]\longmapsto
[1/x_0:\cdots:1/x_m].
$$
It is its own inverse, has type $(m,m)$,
and its multidegrees are
$$
d_i=\binom{m}{i},
\qquad p_{\nabla f}(t)=(1+t)^m.
$$
The slide credits Gonzalez-Sprinberg--Pan,
Huh, and Fassarella--Medeiros.
\end{enumerate}

\medskip\noindent
\textbf{Guiding problems and overview.}
The two problems are to construct
homaloidal hypersurfaces and to classify
them. The slides give the following
picture over $\mathbb{C}$.
\begin{enumerate}
\item Dolgachev's classification:
a reduced homaloidal plane curve is
exactly one of a smooth conic, a smooth
conic together with a tangent line, or
three nonconcurrent lines.
\item Multiplicities of irreducible
factors do not change the multidegrees.
If $f_{\mathrm{red}}$ is obtained by
retaining each factor once, then
$$
d_i(\nabla f)=
d_i(\nabla f_{\mathrm{red}}).
$$
The slide credits work of Dimca and
Papadima, Aluffi, Fassarella and Pereira,
and Huh.
\item Huh's theorem, proving Dimca's
conjecture: a reduced hypersurface of
degree at least three in
$\mathbb{P}^m$, with $m\geq3$ and only
isolated singularities, is not homaloidal.
Thus higher-degree homaloidal examples
must have a positive-dimensional singular
locus.
\item A reduced hyperplane arrangement
in $\mathbb{P}^m$ is homaloidal exactly
when it consists of $m+1$ hyperplanes in
general position. The slide credits
Bruno and Fassarella--Medeiros.
\end{enumerate}

\medskip\noindent
\textbf{Higher-dimensional examples.}
For every $m\geq3$ and $d\geq3$, there
are irreducible reduced homaloidal
hypersurfaces of degree $d$ in
$\mathbb{P}^m$.
The slide mentions duals of certain
rational scrolls
(Ciliberto--Russo--Simis), and Huh's
approach using CSM classes, Newton
polytopes, and mixed volumes.
Its explicit example is
$$
\begin{aligned}
f={}&z_0z_m^{d-1}+z_1z_2z_m^{d-2}
+z_2^d\\
&+(z_3^2+\cdots+z_{m-1}^2)z_m^{d-2}.
\end{aligned}
$$
For $m=3$ the final sum is empty.
The photographs do not include the proof
of this example.

\medskip\noindent
\textbf{From a hypersurface to a
foliation: the photographed construction.}
Given $X=V(f)\subset\mathbb{P}^m$ of
degree $a$, introduce $z=x_{m+1}$.
The Fassarella--Pereira construction is
the foliation $\mathcal{F}$ on
$\mathbb{P}^{m+1}$ defined by
$$
\omega=z\,df-af\,dz.
$$
Its leaves are the level hypersurfaces of
$f/z^a$, equivalently the members of the
pencil generated by $f$ and $z^a$:
$$
\lambda f+\mu z^a=0,
\qquad [\lambda:\mu]\in\mathbb{P}^1.
$$
If $p(t)$ is the multidegree polynomial
of $\nabla f$ and $P(t)$ that of the
Gauss map of $\mathcal{F}$, the last
photograph states
$$
P(t)=(1+t)p(t)\cdot1.
$$
The slide identifies this as the case
$g=y_0$ of the talk's construction.
The general product formula and its proof
were not photographed.

\medskip\noindent
\textbf{AI note: why this is a foliation.}
This continues the description by
1-forms in HP's algebraic-geometry notes.
Away from its singularities, a
codimension-one foliation specifies
tangent hyperplanes which fit together
into local hypersurfaces, the leaves.
Here the identity
$$
d(f/z^a)=z^{-a-1}\omega
$$
shows that those hyperplanes are tangent
to the level sets. In particular
$\omega\wedge d\omega=0$.
As usual, common factors of its
coefficients are cancelled when defining
the singular foliation.
The Gauss map sends a regular point to
the tangent hyperplane of its leaf:
$$
G_{\mathcal{F}}([x:z])=
[z f_{x_0}:\cdots:z f_{x_m}:-af].
$$
Its topological degree equals that of
$\nabla f$. Thus a homaloidal $f$
produces a foliation with birational
Gauss map, the meaning of ``Cremona
foliation'' used here.
See Fassarella--Medeiros,
\href{https://arxiv.org/abs/1108.2676}
{On the polar degree of projective
hypersurfaces}, Remark 2.2, and
Fassarella--Pereira,
\href{https://arxiv.org/abs/0705.1867}
{On the degree of polar transformations},
Theorem 2, for the multidegree relation.

\medskip\noindent
\textbf{AI reconstruction: the missing
Perazzo-type argument.}
The official abstract confirms the
Cremona-to-Perazzo construction and the
dimensions below, but the sources located
do not give its formula.
The following is an independently checked
construction consistent with Daniel's
recollection. It is not a transcription
or a verified identification of Nivaldo's
missing slides.

For $n\geq1$, start with any Cremona map
$$
\begin{aligned}
\psi:\mathbb{P}^n&\dashrightarrow
\mathbb{P}^n,\\
[y]&\longmapsto
[q_0(y):\cdots:q_n(y)],
\end{aligned}
$$
where the $q_i$ have the same positive
degree $e$ and no common factor.
Introduce $n+1$ new variables and form
$$
F(x,y)=\sum_{i=0}^n x_iq_i(y).
$$
This has degree $e+1$ and defines a
hypersurface in $\mathbb{P}^{2n+1}$.
Writing $J_q$ for the square Jacobian
matrix, its gradient is
$$
\begin{aligned}
J_q(y)&=
\left(\frac{\partial q_i}
{\partial y_j}\right)_{0\leq i,j\leq n},\\
\nabla F([x:y])&=
[q(y):J_q(y)^{\mathsf T}x].
\end{aligned}
$$
The first block retains the original
Cremona map. Since $\psi$ is birational
over $\mathbb{C}$, $J_q$ is invertible
at a general point. Indeed its derivative
is invertible on projective tangent
spaces, while Euler's identity gives
$$
J_q(y)y=e\,q(y)
$$
in the remaining radial direction.

For a general target $[u:v]$, use the
rational inverse of $\psi$ to recover
$[y]$ from $[u]$. Choose a representative
$y$ rationally on an affine chart, and
write $q(y)=c u$. Then set
$$
x=c\,(J_q(y)^{\mathsf T})^{-1}v.
$$
The gradient of this point is
$[c u:c v]=[u:v]$. Replacing $y$ by
$s y$ replaces $c$ by $s^e c$ and $x$
by $s x$, so the recovered point
$[x:y]$ is well defined.
This gives a rational inverse to
$\nabla F$: the constructed hypersurface
is homaloidal.

For example, the standard plane Cremona
map gives the cubic in $\mathbb{P}^5$
$$
F=x_0y_1y_2+x_1y_0y_2+x_2y_0y_1.
$$
Finally, apply the photographed
construction to this $F$: in
$\mathbb{P}^{2n+2}$ take
$$
\omega=z\,dF-(e+1)F\,dz.
$$
Its Gauss map is birational. This supplies
a concrete route from an arbitrary
Cremona map to a homaloidal hypersurface
and then to a Cremona foliation.

\medskip\noindent
\textbf{AI note: which Perazzo example?}
For comparison, a standard Perazzo
example is the cubic threefold in
$\mathbb{P}^4$ given by
$$
x_0u^2+x_1uv+x_2v^2=0
$$
Its first three partial derivatives
satisfy an algebraic relation:
$$
(uv)^2=(u^2)(v^2).
$$
Its Hessian vanishes and it is not
homaloidal. The reconstruction above
uses the same pattern, linear in one
block of variables, but the coefficients
now define a Cremona map in a second
block of the same size.
This distinguishes the vanishing-Hessian
example from the construction above.
It does not identify the specific
``Perazzo surface'' Daniel recalls.
For the classical terminology, see
Miró-Roig--Pérez Díez,
\href{https://arxiv.org/abs/2402.09188}
{Perazzo hypersurfaces and the weak
Lefschetz property}, introduction.

\medskip\noindent
\textbf{Closing sentence, recorded by
Daniel.}
\begin{quote}
``This shows that the space of the
homaloidal hypersurfaces in higher
dimensions are at least as complicated
as spaces of birational maps''
\end{quote}
\textbf{AI interpretation.}
The reconstruction explains a possible
meaning: arbitrary Cremona maps can be
encoded in higher-dimensional homaloidal
equations, and recovered from their first
block of partial derivatives with the
variable splitting retained.
The precise statement about parameter
spaces made in the talk was not captured;
these notes do not assert an equivalence
of moduli spaces.

\section{Title to be announced: Kézia
Patrícia
Mestre Carvalho}
\label{iv-jornadas-lefschetz-2026-section-kezia-carvalho}

Kézia Patrícia Mestre Carvalho (UFPE).

Terça 29/09, 11:40.

The official programme has not yet
announced a title.

\medskip\noindent
\textbf{Notes.}
To be added during the talk.

\section{Title to be announced: João Paulo
Costalonga}
\label{iv-jornadas-lefschetz-2026-section-joao-costalonga}

João Paulo Costalonga (UFES).

Terça 29/09, 14:30.

The official programme has not yet
announced a title.

\medskip\noindent
\textbf{Notes.}
To be added during the talk.

\section{Bourbaki schemes of three generated
ideals}
\label{iv-jornadas-lefschetz-2026-section-marcos-jardim}

Marcos Jardim (UNICAMP).

Quarta 30/09, 9:30.

\medskip\noindent
\textbf{Notes.}

Upshot: for ideals generated by three
elements, compute what are the possible
Tjurina and Bourbaki numbers. We note there
are values which cannot be attained and
investigate why.

\medskip\noindent
\textbf{The Bourbaki exact sequence.}
Let $R=k[x_1,\ldots,x_n]$, with $n\geq3$,
and let $J=(f_1,f_2,f_3)$ be minimally
generated by forms of degree $d$, with
no nonconstant common factor. Use the
coefficient grading on the syzygy module:
$$
\operatorname{Syz}(J)
=\ker\bigl(R^3\longrightarrow J(d)\bigr).
$$
Here the map sends $(a_1,a_2,a_3)$ to
$a_1f_1+a_2f_2+a_3f_3$.
Choose a nonzero syzygy $\nu$ of the
smallest possible coefficient degree $e$.
The short exact sequence is
$$
0\longrightarrow R(-e)
\xrightarrow{\nu}\operatorname{Syz}(J)
\longrightarrow M_\nu\longrightarrow0,
$$
where the first map sends $1$ to $\nu$.
If $M_\nu$ is not free, it is a shifted
homogeneous ideal of height two:
$$
M_\nu\cong I_\nu(e-d).
$$
Thus the Bourbaki sequence becomes
$$
0\longrightarrow R(-e)
\xrightarrow{\nu}\operatorname{Syz}(J)
\longrightarrow I_\nu(e-d)
\longrightarrow0.
$$
The ideal $I_\nu$ is the Bourbaki ideal,
and its quotient defines the degree
$$
\operatorname{Bour}(J)=\deg(R/I_\nu).
$$
This degree is independent of the choice
of $\nu$ of degree $e$. In the free case,
$M_\nu\cong R(e-d)$ and the Bourbaki
degree is defined to be zero.
For a plane curve $X=V(F)$, take
$J=J_F$ and write
$\operatorname{Bour}(X)$.
The sequence and grading are checked
against Jardim--Nasrollah Nejad--Simis,
\href{https://arxiv.org/abs/2308.11467}
{The Bourbaki Degree of Plane Projective
Curves}, Section 2.1, and the extension
to three-generated ideals in
\href{https://arxiv.org/abs/2604.04252}
{Jardim--Monteiro--Nasrollah Nejad},
Section 3.

The talk starts defining the Bourbaki degree
via a well-known exact sequence associated
to the first syzygy module of an ideal.

Motivation: investigate why the construction
for cubics does not give Bourbaki degree 2
in any case. Likewise, for quartics, we do
not find any $X$ with
$$
\operatorname{Bour}(X)=5.
$$

The idea is that this is not an
``integrability'' phenomenon---something else
must be going on. Here integrability
means that the polynomials are the partial
derivatives of a polynomial of one degree
higher.

This led to an ``existence theorem'',
which I think formalizes this ``missing
Bourbaki dimension'' phenomenon.
I think this theorem somehow discusses
two numbers $(e,j)$ instead of only one
as in the motivating examples.

This leads to ``three kinds of gaps''
(photo).

\medskip\noindent
\textbf{Three kinds of gaps.}
For a fixed $d$ and a numerically admissible
pair $(e,j)$:
\begin{itemize}
\item \emph{Structural gap}: no
$d$-equigenerated triple realizes $(e,j)$
in a fixed dimension;
\item \emph{Dimensional gap}: the pair is
impossible in dimension $n_0$, but occurs
in a higher dimension;
\item \emph{Integrable gap}: in three
variables, the pair is realized by a
triple, but not by the gradient of a
reduced plane curve.
\end{itemize}
The slide's inclusion diagram reads:
pairs realized by a gradient are among
those realized by a triple, which are
among the numerically admissible pairs.

\medskip\noindent
\textbf{Gradient ideals.}
Consider the special case $n=3$:
$$
R=k[x_1,x_2,x_3],\qquad F\in R_{d+1},
\qquad f_i=\partial_iF.
$$
The corresponding three-generated ideal
is the gradient ideal
$$
J_F=(\partial_1F,\partial_2F,\partial_3F).
$$
If $X=V(F)\subset\mathbb{P}^2$, then
the syzygy module can be regarded as the
graded module of global sections of the
logarithmic tangent sheaf $T_F$:
$$
\operatorname{Syz}(J_F)\cong
\bigoplus_{t\in\mathbb{Z}}
H^0(\mathbb{P}^2,T_F(t)).
$$
Bourbaki ideals for plane curves were
studied with A. Nasrollah Nejad and
A. Simis in
\href{https://arxiv.org/abs/2308.11467}
{arXiv:2308.11467}.

\medskip\noindent
\textbf{Gradient ideals of plane curves.}
Let $X=V(F)\subset\mathbb{P}^2$ be a
reduced plane curve of degree $d+1$, not
a cone. Its gradient ideal
$$
J_F=(F_x,F_y,F_z)
$$
is generated by three forms of degree $d$.
The two numerical invariants become
geometric: the slide identifies
$$
\deg(R/J_F)=\tau(X)
$$
and describes $\operatorname{Syz}(J_F)$
as logarithmic vector fields along $X$.
Here $\tau(X)$ is the total Tjurina number.
Therefore
$$
\boxed{\operatorname{Bour}(X)+\tau(X)
=d^2-ed+e^2}.
$$
This is the formula of
Jardim--Nasrollah Nejad--Simis,
\href{https://arxiv.org/abs/2308.11467}
{Theorem 2.1}.
The realization problem now asks which
triples can be integrated to a gradient
triple.

\emph{Convention for the degree formula.}
The equality with $\deg(R/J_F)$ is for
singular $X$. For smooth $X$, the
projective singular scheme is empty and
$\tau(X)=0$; the boxed formula still holds.
The sheaf $T_F$ uses the slide's grading,
namely the sheaf associated to the
coefficient-graded syzygy module above.

See
\href{https://arxiv.org/abs/2604.04252}
{``The bourbaki degree of syzygy module of
$2\times4$ matrices''} and
\href{https://arxiv.org/abs/2609.17799}
{``Geometry of three-generated ideals in
the plane''}.

The three slide transcriptions above are
from \texttt{IMG\_4017.pdf}.

\section{Non-dominant quadratic maps and
cubics with
vanishing Hessian}
\label{iv-jornadas-lefschetz-2026-section-thiago-fassarela}

Thiago Fassarella (UFF).

Quarta 30/09, 10:30.

\medskip\noindent
\textbf{Basic idea.}
These are brief notes from the four
photographs in \texttt{IMG\_4020.pdf}.
For a cubic $f$, its partial derivatives
are quadrics. Over $\mathbb{C}$,
vanishing of the Hessian determinant
means that the polar map is not dominant.
The talk uses the geometry of quadratic
maps with lower-dimensional image to
study such cubics.

\medskip\noindent
\textbf{The four quadratic maps.}
The four examples are visible on the
blackboard in the third photograph:
$$
\begin{aligned}
g_1([x:y:z:w])&=[xz:xw:yz:yw],\\
g_2([x:y:z:w])&=[x^2:xy:y^2:xz+yw],\\
g_3([x:y:z:w])&=[x^2:xy:y^2:zw],\\
g_4([x:y:z:w])&=[x^2:xy:y^2:z^2+xw].
\end{aligned}
$$
These are the four preceding examples
referred to in the photographed theorem.

\begin{theorem}[Photographed slide]
\label{iv-jornadas-lefschetz-2026-theorem-fassarella-quadratic-maps}
Let
$$
g:\mathbb{P}^3\dashrightarrow\mathbb{P}^3
$$
be defined by four linearly independent
quadrics with no common factor.
Suppose the image closure $Y$ is a
surface. Then
$$
2\leq\deg Y\leq4.
$$
If $g$ does not factor through a linear
projection, it is equivalent to one of
$g_1,g_2,g_3,g_4$.
Equivalence permits projective coordinate
changes in both source and target.
\end{theorem}

\medskip\noindent
\textbf{AI clarification of the pictures.}
With target coordinates $[A:B:C:D]$,
the image of $g_1$ is the smooth quadric
$$
AD=BC.
$$
The other three images are the quadric
cone
$$
AC=B^2.
$$
Thus the image is not always a cone.
The general fibre closures of $g_1,g_2$
are lines; those of $g_3,g_4$ are conics.
For example, over $[1:t:t^2:s]$,
the last three maps have fibre closures
$$
\begin{array}{ll}
g_2:&y=tx,\quad z+tw=sx,\\
g_3:&y=tx,\quad zw=sx^2,\\
g_4:&y=tx,\quad z^2+xw=sx^2.
\end{array}
$$
These descriptions follow directly by
substituting into the formulas.

\medskip\noindent
\textbf{Foliating projective space by
lines.}
This was the part Daniel particularly
liked. Here a foliation by lines is a
family with exactly one member through a
general point of $\mathbb{P}^3$;
singularities are allowed.
The classical possibilities are:
\begin{enumerate}
\item \emph{A star.} All lines through
one fixed point $p$.
\item \emph{Secants of a twisted cubic.}
Join pairs of points on a twisted cubic,
including tangent lines as limits.
\item \emph{Moving pencils centred on a
line.} Fix a line $L$. For each plane
$\Pi$ through $L$, choose
$p_\Pi=\psi(\Pi)\in L$, where $\psi$ is a
nonconstant morphism from the pencil of
planes to $L$. In $\Pi$, take all lines
through $p_\Pi$.
\item \emph{A line and a rational curve.}
Take lines joining $L$ to a rational
curve $C$, where a general plane through
$L$ meets $C$ in exactly one point
outside $L$. Two skew lines give a
basic example.
\end{enumerate}
This list is reconstructed from
Fassarella,
\href{https://arxiv.org/abs/0809.2761}
{Foliations with degenerate Gauss maps
on $\mathbb{P}^4$},
Section 4.2, Theorem 1 (Kummer).
It is distinct from the four quadratic
maps above.

The photographed proof sketch says that
line fibres give such a foliation.
A star forces factorization through the
linear projection from its centre;
the non-projection line-fibre cases give
$g_1$ and $g_2$.
For $g_1$, the lines join the two skew
base-locus lines
$$
\{x=y=0\},\qquad\{z=w=0\}.
$$

\medskip\noindent
\textbf{Lower-dimensional picture.}
The first slide gives this overview of
cubics with vanishing Hessian:
\begin{itemize}
\item In $\mathbb{P}^n$, $n\leq3$:
only cones.
\item In $\mathbb{P}^4$:
Gordan--Noether (1876), Franchetta (1954);
Fassarella, Ferrer and Gondim (2021),
with the connection to
$\operatorname{Gor}(1,5,5,1)$.
\item In $\mathbb{P}^5,\mathbb{P}^6$:
Perazzo (1900), Gondim--Russo (2015);
Bezerra, Ferrer and Gondim (2025), with
$\operatorname{Gor}(1,n+1,n+1,1)$
for $n=5,6$.
\end{itemize}
The names and dates here follow the slide.

\medskip\noindent
\textbf{The final slide: cubics in
$\mathbb{P}^7$.}
In the slide's notation, an intrinsic
linear space in the singular locus gives
a presentation
$$
\begin{gathered}
\langle Z^*\rangle
\subset\operatorname{Sing}(X),\\
f=\sum_{i=0}^{k}y_i g_i(x)+h(x).
\end{gathered}
$$
The $g_i$ are quadrics and $h$ is cubic.
The quadratic-map classification supplies
the possible systems $(g_i)$:
\begin{itemize}
\item If $\dim\langle Z^*\rangle=2$:
$$
(x_0^2,x_0x_1,x_1^2).
$$
\item If $\dim\langle Z^*\rangle=3$:
a projection of the Veronese surface
from a line, or
$$
(x_0x_2,x_0x_3,x_1x_2,x_1x_3).
$$
\item If $\dim\langle Z^*\rangle=4$:
a projection of the Veronese surface
from a point.
\end{itemize}
The slide also says that projective
equivalence within each class can be
determined. Its preceding definitions
and full hypotheses were not photographed.

\medskip\noindent
\textbf{AI reminder: projection centres.}
The Veronese surface lies in
$\mathbb{P}^5$. Projection from a line
has target $\mathbb{P}^3$; projection
from a point has target $\mathbb{P}^4$.
The centre is the linear subspace where
projection is undefined.

\section{On Artinian Gorenstein Algebras with
Macaulay Dual Generator of Fixed Waring
Rank}
\label{iv-jornadas-lefschetz-2026-section-janaine-martins}

Janaíne Geralda Mesquita Martins (UFMG).

Quarta 30/09, 11:40.

\medskip\noindent
\textbf{Notes.}
This was a PhD student's project defense,
a stage before the PhD defense. She has
collaborated with Barbara, Rodrigo, Pedro
and Miró. Her project has three main
topics. The second is splitting of
syzygy bundles; another is Jordan
decomposition of algebras (Daniel does
not recall whether this was the first
or third topic).

\section{Higgs Grassmannians}
\label{iv-jornadas-lefschetz-2026-section-ugo-bruzzo}

Ugo Bruzzo (UFMG).

Quinta 01/10, 9:30.

\medskip\noindent
\textbf{Notes.}
Brief reconstruction from Daniel's
recollections, with definitions checked
against the papers linked below.
Daniel followed closely but took few
notes; the promised slides are awaited.

\medskip\noindent
\textbf{Higgs bundles.}
On a smooth complex projective variety
$X$, a Higgs bundle is a vector bundle
$E$ with an $\mathcal{O}_X$-linear map
$$
\phi:E\longrightarrow E\otimes\Omega_X^1,
\qquad \phi\wedge\phi=0.
$$
Unlike a connection, it acts pointwise:
$\phi(fs)=f\phi(s)$, with no $df$ term.

\medskip\noindent
\textbf{Quotients and Grassmannians.}
The relative Grassmannian
$\pi:G=\operatorname{Gr}_q(E)\to X$
represents rank-$q$ locally free quotients:
for $g:T\to X$, maps $T\to G$ over $X$
correspond to locally free quotients
$g^*E\twoheadrightarrow Q$, up to
isomorphism. Its universal
sequence is
$$
0\longrightarrow S\longrightarrow\pi^*E
\longrightarrow Q\longrightarrow0.
$$
The Higgs Grassmannian is the closed
subscheme cut out by vanishing of
$$
S\longrightarrow\pi^*E
\xrightarrow{\pi^*\phi}
\pi^*E\otimes\pi^*\Omega_X^1
\longrightarrow Q\otimes\pi^*\Omega_X^1.
$$
Thus the kernel is preserved by the
Higgs field, which descends to the
quotient. It can be singular or nonreduced.
See Bruzzo--Graffeo--Graña Otero,
\href{https://arxiv.org/abs/2605.29147}%
{\emph{Higgs Grassmannians}} (2026);
that paper uses the subbundle convention.
Daniel recalls this construction from an
earlier talk; see the Higgs-bundle passage
in \emph{Algebraic Geometry}.

\medskip\noindent
\textbf{Nefness in any dimension.}
A line bundle $L$ is nef when
$$
L\cdot C=\deg(L|_C)\geq0
$$
for every integral projective curve $C$.
This pairs divisors with curves; it needs
no surface intersection form.
A vector bundle $E$ is nef when
$\mathcal{O}_{\mathbb{P}(E)}(1)$ is nef,
using one-dimensional quotients.
Higgs-nefness is a different, recursive
condition using universal Higgs quotients
and the determinant.

\medskip\noindent
\textbf{Curve semistability and
discriminant.}
Curve semistability means that every
pullback to a smooth projective curve
$f:C\to X$ is Higgs-semistable:
$$
\frac{\deg F}{\operatorname{rank}F}
\leq\frac{\deg(f^*E)}{\operatorname{rank}E}
$$
for each nonzero proper Higgs-invariant
subbundle $F$. The pulled-back field uses
$f^*\Omega_X^1\to\Omega_C^1$.
On curves no polarization is needed.
For $r=\operatorname{rank}E$, set
$$
\Delta(E)=c_2(E)-\frac{r-1}{2r}c_1(E)^2
\in H^4(X,\mathbb{Q}).
$$
The original conjecture of Bruzzo and
Graña Otero equated curve semistability
with Higgs semistability for some
polarization together with $\Delta(E)=0$.
The latter implies the former; the
conjectural direction was the converse.

\medskip\noindent
\textbf{Ranks two and four.}
The conjecture holds in rank two
(the parameter is rank, not dimension):
Bruzzo--Graña Otero--Hernández Ruipérez,
\href{https://arxiv.org/abs/2208.10554}%
{\emph{On a conjecture about Higgs bundles
for rank 2 and some inequalities}}.
Huang's August 2026
\href{https://arxiv.org/abs/2608.22448}%
{counterexample preprint}
constructs a curve-semistable rank-four
Higgs bundle on a surface with
$$
c_1(E)=0,\qquad\int_X c_2(E)=10.
$$
Adding trivial Higgs summands gives
counterexamples in every rank $r\geq4$.
Daniel's question about possible AI
involvement remains unverified.

\medskip\noindent
\textbf{Hodge limits and the proposed
reformulation.}
The relevant degeneration scales the
Higgs field. In moduli, its Hodge limit
is written
$$
(E_0,\phi_0)=\lim_{t\to0}(E,t\phi).
$$
The limit is a fixed point of scaling.
When represented by a Higgs bundle, it
has a Hodge grading
$$
E_0=\bigoplus_p E_0^p,
\qquad
\phi_0(E_0^p)\subset
E_0^{p-1}\otimes\Omega_X^1.
$$
This is a degeneration of the bundle
and field on fixed $X$; the underlying
holomorphic bundle can change, so the
limit need not be $(E,0)$.
See Simpson, \emph{The ubiquity of
variations of Hodge structure}.

According to Daniel's recollection,
Bruzzo and collaborators plan to impose
an extra condition on the Hodge limit
and expect Huang's example to be excluded.
The precise condition was not recorded.
The discussion also involved ``three'';
whether this meant rank three or base
dimension three must be checked in the
slides. This paragraph records the talk,
not a verified reformulated conjecture.

No slide link was found on the official
programme when checked on October 1.

\section{Almost Cohen-Macaulay height two
ideals
equigenerated by three forms}
\label{iv-jornadas-lefschetz-2026-section-thiago-cabral}

Thiago Fiel da Costa Cabral (UFPE).

Quinta 01/10, 10:30.

\medskip\noindent
\textbf{Notes.}
For our discussion of minimal resolutions
and grading shifts, see
\emph{Commutative algebra}, Section
\ref{commutative-algebra-%
section-minimal-resolutions}.
The following section,
\ref{commutative-algebra-%
section-hilbert-burch}, compares
Hilbert--Burch for perfect height-two ideals
with the 2026 theorem of
Burity--Fiel--Ramos--Simis for
almost Cohen--Macaulay height-two ideals
equigenerated by three forms.

\section{Multiple Tor modules}
\label{iv-jornadas-lefschetz-2026-section-rafael-holanda}

Rafael Ferreira Holanda (UFPE).

Quinta 01/10, 11:40.

\medskip\noindent
\textbf{Notes.}
There are two constructions generalizing
the Mayer--Vietoris sequence for two
ideals of a commutative ring $R$:
$\mathbb{P}$ uses products and
$\mathbb{S}$ uses sums.
They are complexes; their exactness
requires hypotheses.

The statements below are checked against
Banerjee--Chardin--Holanda (2026),
\href{https://arxiv.org/abs/2604.20646}{%
\emph{Multiple Tor modules: rigidity and
Mayer--Vietoris spectral sequences}}
(BCH26). They supplement these partial
talk notes; the theorem numbering is
that of the paper.

\medskip\noindent
\textbf{From two ideals to several.}
Recall the exact sequence
$$
\begin{gathered}
0\longrightarrow R/(I_1\cap I_2)
\longrightarrow R/I_1\oplus R/I_2\\
\longrightarrow R/(I_1+I_2)
\longrightarrow0.
\end{gathered}
$$
Replacing the intersection by the product
can destroy injectivity on the left:
$$
\operatorname{Tor}_1^R(R/I_1,R/I_2)
\cong (I_1\cap I_2)/(I_1I_2).
$$
This is already a link between Tor
vanishing and Mayer--Vietoris exactness.

For $I_1,\ldots,I_n$, set
$\mathbb{S}^0=R/(I_1\cdots I_n)$ and
$\mathbb{P}_0=0$. For $1\leq p\leq n$,
$$
\begin{aligned}
\mathbb{S}^p
&=\bigoplus_{i_1<\cdots<i_p}
R/(I_{i_1}+\cdots+I_{i_p}),\\
\mathbb{P}_p
&=\bigoplus_{i_1<\cdots<i_p}
R/(I_{i_1}\cdots I_{i_p}).
\end{aligned}
$$
The differentials are alternating sums
of the natural quotient maps: increasing
$p$ for $\mathbb{S}$ and decreasing $p$
for $\mathbb{P}$.

\medskip\noindent
\textbf{Canonical free resolutions.}
Rafael uses canonical free resolutions,
rather than minimal ones, to define
multiple Tor. One standard construction
starts with
$$
F(M)=\bigoplus_{m\in M}Re_m,
\qquad e_m\longmapsto m.
$$
Take the free module on every element
of the kernel, then repeat.
No generating set needs to be chosen,
but the resulting resolution can be
enormous and impractical to compute.
This explains the existence-versus-use
point from the talk; the paper does not
spell out the canonical construction.

Choose free resolutions
$F^{(a)}_\bullet\longrightarrow M_a$.
The multiple Tor modules are
$$
\operatorname{Tor}_i^R(M_1,\ldots,M_s)
=H_i\!\left(\operatorname{Tot}
\left(\bigotimes_{a=1}^s
F^{(a)}_\bullet\right)\right).
$$
Here the total complex groups together
tensor terms whose homological degrees
sum to the same integer.
One may use any flat resolutions.
The modules are independent of this
choice and symmetric in the entries.
In degree zero they give
$$
\operatorname{Tor}_0^R(M_1,\ldots,M_s)
\cong M_1\otimes_R\cdots\otimes_R M_s.
$$
The background is EGA III, part I,
Section 6. For example, flat base change
$R\longrightarrow A$ gives
$$
\begin{aligned}
&\operatorname{Tor}_i^A
(M_1\otimes_R A,\ldots,M_s\otimes_R A)\\
&\qquad\cong
\operatorname{Tor}_i^R(M_1,\ldots,M_s)
\otimes_R A.
\end{aligned}
$$
See BCH26, Section 2, especially
Proposition 2.4, citing EGA III, 6.9.2.

\medskip\noindent
\textbf{Question about flatness.}
Wouldn't flat mean that Tor vanishes for
all degrees? Yes, for ordinary Tor:
$M$ is flat precisely when
$$
\operatorname{Tor}_i^R(M,N)=0
$$
for every $R$-module $N$ and every $i>0$.
It suffices to check $i=1$ for every $N$.
Degree zero is the tensor product.
Tensoring a resolution with a flat module
preserves exactness, which explains the
positive-degree vanishing.

For multiple Tor, one flat entry alone
does not suffice. For example,
$$
\operatorname{Tor}_i^R(R,M,N)
\cong\operatorname{Tor}_i^R(M,N),
$$
which can be nonzero. Flat terms in the
chosen resolutions do not mean that
the modules being resolved are flat.

\begin{theorem}[BCH26, Theorem 1.1]
\label{iv-jornadas-lefschetz-2026-theorem-rafael-rigidity}
Suppose $R$ is local and regular and
contains a field. Take nonzero finitely
generated modules $M_1,\ldots,M_s$ and put
$$
T_i=\operatorname{Tor}_i^R
(M_1,\ldots,M_s),\qquad
M=M_1\otimes_R\cdots\otimes_R M_s.
$$
For $i>0$, $T_i=0$ implies $T_j=0$ for
every $j\geq i$. It also implies
$$
\operatorname{Tor}_i^R
(M_1,\ldots,M_t)=0
$$
for every $t<s$, in any ordering.

Moreover, these conditions are equivalent:
\begin{enumerate}
\item $T_1=0$ and $M$ is Cohen--Macaulay.
\item The following equality holds:
$$
\operatorname{codim}_R M
=\sum_{a=1}^s\operatorname{pd}_R M_a.
$$
\item Each $M_a$ is Cohen--Macaulay and
their intersection is proper, meaning
$$
\operatorname{codim}_R M
=\sum_{a=1}^s\operatorname{codim}_R M_a.
$$
\end{enumerate}
Here codimension means $\dim R-\dim M$,
and $\operatorname{pd}_R$ denotes
projective dimension.
\end{theorem}

Unlike the flatness criterion, rigidity
tests vanishing for one fixed collection
of modules. It does not assert that
those modules are flat.

\medskip\noindent
\textbf{Tor independence.}
The ideals are called Tor-independent
when
$$
\operatorname{Tor}_i^R
(R/I_1,\ldots,R/I_n)=0
$$
for every $i>0$. Strong Tor independence
requires this for every subcollection.

\begin{proposition}[BCH26, Proposition 1.4]
\label{iv-jornadas-lefschetz-2026-proposition-rafael-spectral}
Assume strong Tor independence of
$I_1,\ldots,I_n$. For any $R$-module $N$
there are two spectral sequences, with
$1\leq p\leq n$ and $q\geq0$:
$$
\begin{aligned}
E^1_{p,q}
&=\bigoplus_{i_1<\cdots<i_p}
\operatorname{Tor}_q^R\!\left(
N,\frac{R}{I_{i_1}+\cdots+I_{i_p}}
\right)\\
&\Longrightarrow
\operatorname{Tor}_{q-p+1}^R\!\left(
N,\frac{R}{I_1\cdots I_n}\right),
\end{aligned}
$$
and
$$
\begin{aligned}
{}'E^1_{p,q}
&=\bigoplus_{i_1<\cdots<i_p}
\operatorname{Tor}_q^R\!\left(
N,\frac{R}{I_{i_1}\cdots I_{i_p}}
\right)\\
&\Longrightarrow
\operatorname{Tor}_{q+p-1}^R\!\left(
N,\frac{R}{I_1+\cdots+I_n}\right).
\end{aligned}
$$
\end{proposition}

The first passes from sums to the product;
the second passes from products to the sum.
These are Corollaries 4.7 and 4.11 in
the body of the paper.

Another result relates Tor independence
to exactness of the complexes above.
Under strong Tor independence both
$$
0\longrightarrow\mathbb{S}^0
\longrightarrow\mathbb{S}^1
\longrightarrow\cdots
\longrightarrow\mathbb{S}^n
\longrightarrow0
$$
and
$$
\begin{gathered}
0\longrightarrow\mathbb{P}_n
\longrightarrow\cdots
\longrightarrow\mathbb{P}_1\\
\longrightarrow R/(I_1+\cdots+I_n)
\longrightarrow0
\end{gathered}
$$
are exact (BCH26, Theorem 1.3).

\medskip\noindent
\textbf{Final result: Noetherian setting.}
Daniel recalls a final BCH26 result
requiring the ring to be Noetherian,
possibly the most important result
of the talk. Its precise statement
still needs to be recovered from the
slides or from Rafael.
The rigidity theorem above assumes
regularity, locality and containment of
a field; Noetherianity alone is not its
hypothesis. These notes do not identify
it with the final result from memory.

\section{Hessians, curvature and lightlike
hypersurfaces}
\label{iv-jornadas-lefschetz-2026-section-thiago-silva}

Thiago Dias de Oliveira Silva (UFRPE).

Quinta 01/10, 14:30.

\medskip\noindent
\textbf{Notes.}
To be added during the talk.

\section{Veronese-Avoiding Hypersurfaces and
Lefschetz Properties}
\label{iv-jornadas-lefschetz-2026-section-abbas-nejad}

Abbas Nasrollah Nejad (UNIFESP).

Sexta 02/10, 9:30.

\medskip\noindent
\textbf{Notes.}
To be added during the talk.

\section{Newton Duality and Macaulay Inverse
System
of Gorenstein Ideals}
\label{iv-jornadas-lefschetz-2026-section-dayane-lira}

Dayane Santos de Lira (UFERSA).

Sexta 02/10, 10:30.

\medskip\noindent
\textbf{Notes.}
To be added during the talk.

% --- End iv-jornadas-lefschetz-2026.tex ---

\part{Moduli and Deformations}
\label{part-moduli-and-deformations}

% --- Begin deformations.tex ---
\chapter{Deformation theory}
\label{chapter-deformations}





\section{Introduction}
\label{deformations-section-intro}

\noindent
\cite{Sernesi-Overview} has a great way of
motivating deformation
theory via the moduli problem.
Here let's just put the following important
meaning for
``deformation theory'' (as opposed to
studying flat families!):

\begin{quote}
``{\it Deformation theory} is the study of
infinitesimal deformaitons
as a tool to understand the local structure
of the moduli space.
The goal is to be able to describe the
restriction
of the universal family to a small
neighborhood of $m \in \mathcal{M}$,
or, more precisely, its restriction to the
germ of $M$ at $m$.

What is interesting here is that we can study
order and infinitesimal
deformations even though the functor $F$ is
not representable or simply we don’t
yet know it is. This is the most frequent
case. Such an investigation will
reveal the infinitesimal properties at $[X]$
of a yet unknown global structure
on $M$ which will be hopefully understood at
a subsequent stage of the
investigation.  In order words it turns out
to be possible and convenient to
separate the {\it global moduli problem} from
the {\it local moduli problem},
and deformation theory studies the latter,
with the purpose of constructing a
family of deformations of a given object
parametrized by the spectrum of a local
ring, and having properties as close as
possible to a universal property.''
\end{quote}


\section{A short history of deformation theory}
\label{deformations-section-history-deformation-theory}

\noindent
The following chronology explains where the
Hilbert scheme, deformation functors and
miniversal deformation spaces enter the
story.

\begin{enumerate}
\item From 1949 to 1954 Grothendieck worked
mainly in functional analysis: topological
tensor products, nuclear spaces and the
theory of Banach spaces. By the end of 1954
he had decided to leave this subject. In
1955 his interests shifted first to
homological algebra and then increasingly to
algebraic geometry. See the
\href{https://research-portal.uu.nl/ws/portalfiles/portal/254168358/Life_and_Work_of_Alexander_Grothendieck.pdf}
{biographical account of Chai and Oort}.

\item In 1957 he obtained the
Grothendieck--Riemann--Roch theorem. At the
1958 International Congress of
Mathematicians he announced his program for
the cohomology of abstract algebraic
varieties. Thus his dream of a general
cohomological framework was already present
before the construction of the Hilbert
scheme. See
\href{https://webusers.imj-prg.fr/~leila.schneps/grothendieckcircle/CohomologyVarieties.pdf}
{Grothendieck's 1958 ICM paper}.

\item In 1960 the first volume of EGA,
\emph{Le langage des sch\'emas}, appeared.
Algebraic geometry was of course already an
old subject; what Grothendieck was creating
was its modern scheme-theoretic language.

\item In May 1961 Grothendieck proved the
representability of the Hilbert functor. More
precisely, if $X\to S$ is projective and a
Hilbert polynomial $P$ is fixed, then the
functor of flat families of closed
subschemes of $X$ with Hilbert polynomial
$P$ is represented by a projective
$S$-scheme $\operatorname{Hilb}^P_{X/S}$.
See
\href{https://www.numdam.org/item/SB_1960-1961__6__249_0/}
{Grothendieck's Bourbaki expos\'e 221}.

This result belonged to his wider program of
\emph{construction theorems}: formulate a
geometric classification problem as a
functor, and then prove that the functor is
represented by a scheme. Hilbert schemes,
Quot schemes and Picard schemes are central
examples of this philosophy.

\item There was simultaneously an analytic
line of development. Kodaira and Spencer
developed the deformation theory of compact
complex manifolds in 1958, and in 1962
Kuranishi constructed locally complete, or
miniversal, analytic families.

\item In 1968 Schlessinger's
\emph{Functors of Artin Rings} gave precise
criteria for an infinitesimal deformation
functor to be prorepresentable, or at least
to possess a hull. The hull is precisely the
formal miniversal deformation space that we
shall define below. See
\href{https://doi.org/10.1090/S0002-9947-1968-0217093-3}
{Schlessinger's original paper}.

\item In 1969 Michael Artin proved his
approximation and algebraization theorems.
These results form a bridge from formal
moduli problems to actual algebraic
families. See
\href{https://www.numdam.org/articles/10.1007/BF02684596/}
{Artin's 1969 approximation paper}.
\end{enumerate}

\noindent
After 1961 Grothendieck continued developing
the EGA and SGA machinery: formal geometry,
topoi, descent, the \`etale fundamental
group and \`etale cohomology. The last of
these supplied the main cohomological
machinery for the Weil conjectures. His more
ambitious vision was the theory of
\emph{motives}: a universal geometric source
from which the different reasonable
cohomology theories should arise. This
program was not completed. In particular,
Grothendieck's standard conjectures on
algebraic cycles remain open. He left the
IH\'ES in 1970, and Deligne proved the final
Weil conjecture in 1974 by a different
route.

\begin{remark}[Algebraization in the
Gr\"unbaum paper]
Two related results must not be conflated.
Grothendieck's existence theorem, developed
in EGA III in 1961, gives effectivity and
algebraization results over a complete local
base in the proper or projective setting.
Artin's later approximation and
algebraization theorems allow one, under
suitable hypotheses, to pass from a formal
versal deformation to an algebraic family
of finite type. See
\href{https://www.numdam.org/item/?id=PMIHES_1961__11__5_0}
{EGA III} and the
\href{https://stacks.math.columbia.edu/tag/04UW}
{Stacks Project's historical summary}.

Consequently, the sentence in our current
handout saying that properness and
``Grothendieck's Existence Theorem'' turn a
formal family directly into a family over
an algebraic curve is probably too quick.
The passage from a formal family over
$\operatorname{Spf}k[[s]]$ to a family over
an algebraic curve generally involves Artin
approximation or algebraization. This step
must be stated carefully in the paper.
\end{remark}


\section{Cotangent complex and
cotangent cohomology}
\label{deformations-section-cotangent-complex-primer}

\noindent
Recall from
Section~\ref{commutative-algebra-section-kahler-differentials}
of \emph{Commutative Algebra} that
$\Omega^1_P$ is the $P$-module generated by
formal symbols $df$, for $f\in P$, modulo
linearity, the Leibniz rule and vanishing
on the base field.

The module $\Omega^1_R$ records
first-order differential information:
the symbols $df$ obey linearity and the
Leibniz rule, together with the relations
coming from the algebra $R$.
These relations hold in $\Omega^1_R$,
but the quotient does not retain how
they were imposed, or the relations
among those relations.
The cotangent complex keeps track of
this additional data.
Its extra terms record successive levels
of relations, not higher differential
forms $\Omega^2_R,\Omega^3_R$.


Why use differential forms to study
infinitesimal changes? A map
$r\mapsto r+\varepsilon D(r)$, with
$\varepsilon^2=0$, preserves multiplication
exactly when $D$ satisfies the Leibniz
rule. Such derivations are encoded by
$$
\operatorname{Der}_k(R,R)
\cong\operatorname{Hom}_R(\Omega^1_{R/k},R).
$$
Thus forms describe tangent information
by duality. Dualizing the whole cotangent
complex in the derived sense gives the
\emph{tangent complex}
$$
\operatorname{RHom}_R(L_{R/k},R).
$$
Using the free-module model below, this
means applying $\operatorname{Hom}_R(-,R)$
term by term and reversing the arrows.
Its cohomology modules are precisely
$T^i_{R/k}$. Their interpretation in terms
of deformations and obstructions is a
further theorem, discussed afterward.

Now let $I\subset P$ be an ideal and set
$$
R:=P/I.
$$
We want to understand $\Omega^1_R$. That is our starting point.

What changes when we pass from $P$ to $R$? Every $f\in I$ becomes zero in $R$. Therefore, in $\Omega^1_R$ we must have
$$
df=d(0)=0.
$$
Thus, to construct $\Omega^1_R$ from the differentials upstairs, we must kill the differentials $df$ of the equations $f\in I$.

There is one preliminary point. $\Omega^1_P$ is a $P$-module, whereas $\Omega^1_R$ must be an $R=P/I$-module. In an $R$-module every element of $I$ acts as zero. Thus we first force
$$
f\omega=0\qquad(f\in I,\ \omega\in\Omega^1_P)
$$
by passing to
$$
\Omega^1_P/I\Omega^1_P.
$$
In other words, we keep the differentials from $P$, but make their coefficients behave as coefficients in $R$. Equivalently,
$$
\Omega^1_P/I\Omega^1_P\cong\Omega^1_P\otimes_P R.
$$

Now we have the correct ambient $R$-module
$$
\Omega^1_P\otimes_P R.
$$
But it still contains the classes of $df$ for $f\in I$. In $\Omega^1_R$ these must become zero, because $f=0$ in $R$ and hence $df=d(0)=0$. Therefore we quotient them out:
$$
\Omega^1_R\cong
\frac{\Omega^1_P\otimes_P R}
{\langle df:f\in I\rangle}.
$$
So the next question is how to describe systematically the submodule $\langle df:f\in I\rangle$. This is where $I/I^2$ enters.

Why does this quotient really give $\Omega^1_R$? By definition, $\Omega^1_R$ has formal differentials
$$
d\bar p
$$
for elements
$$
\bar p\in R=P/I.
$$
Every $\bar p$ is represented by some $p\in P$, so its differential comes from $dp$. But if $p$ and $p+f$, with $f\in I$, represent the same element of $R$, then their differentials must agree:
$$
d(p+f)=dp.
$$
Equivalently, we must impose
$$
df=0
$$
for every $f\in I$. This is exactly the quotient by
$$
\langle df:f\in I\rangle
$$
above.

We therefore try to record the differentials that are being killed by the map
$$
I\longrightarrow\Omega^1_P\otimes_P R,
\qquad f\longmapsto df.
$$
Its image is precisely
$$
\langle df:f\in I\rangle.
$$

Moreover, if $f,g\in I$, then
$$
d(fg)=f\,dg+g\,df=0
$$
in $\Omega^1_P\otimes_P R$, because coefficients from $I$ act as zero there. Thus $I^2$ lies in the kernel, so the map factors through
$$
I/I^2\longrightarrow\Omega^1_P\otimes_P R,
\qquad [f]\longmapsto df.
$$
Its image is exactly the submodule that we quotient out. Hence we obtain the exact sequence
$$
I/I^2\longrightarrow
\Omega^1_P\otimes_P R\longrightarrow
\Omega^1_R\longrightarrow0.
$$
by simply putting its cokernel.

\noindent
The exact sequence above explains why
$\Omega^1_R$ is not enough for deformation
theory. By taking the cokernel we keep the
resulting differentials, but forget the map
that produced the relations among them.
The cotangent complex keeps this relation
data, and also higher relations among
relations.

\begin{remark}[Polynomial resolution]
\label{deformations-remark-polynomial-resolution-cotangent}
The general construction does not continue
the sequence
$$
I/I^2\longrightarrow
\Omega^1_P\otimes_P R
$$
arbitrarily to the left. Instead, one goes
back and first resolves the $A$-algebra
$R$ itself by polynomial $A$-algebras.
There is a standard simplicial polynomial
resolution
$$
P_\bullet\longrightarrow R,
$$
which may be pictured schematically as
$$
\cdots P_2
\rightrightarrows P_1
\rightrightarrows P_0
\longrightarrow R.
$$
Here $P_0$ is a free polynomial
$A$-algebra mapping onto $R$; the next
levels encode relations, relations among
relations, and so on. Simplicial means
that at level $n$ there are several
compatible face maps
$$
d_i:P_n\longrightarrow P_{n-1},
\qquad 0\leq i\leq n.
$$
Such a polynomial resolution exists. For
the standard construction, let $F(B)$ be
the free polynomial $A$-algebra with one
variable $[b]$ for every element $b\in B$
and iterate:
$$
P_n=F^{n+1}(R).
$$
Evaluation gives the face maps. See
\href{https://stacks.math.columbia.edu/%
tag/08PL}{Stacks, Tag 08PL}.
\end{remark}

\noindent
Now take K\"ahler differentials at every
level,
$$
\Omega^1_{P_n/A},
$$
and change coefficients from $P_n$ to $R$:
$$
\Omega^1_{P_n/A}\otimes_{P_n}R.
$$
The face maps induce maps between these
$R$-modules. Taking their alternating sums
turns them into an ordinary complex.

\begin{definition}
\label{deformations-definition-cotangent-complex-primer}
The \emph{cotangent complex} $L_{R/A}$ is
the resulting complex of $R$-modules:
$$
\cdots\longrightarrow
\Omega^1_{P_2/A}\otimes_{P_2}R
\longrightarrow
\Omega^1_{P_1/A}\otimes_{P_1}R
\longrightarrow
\Omega^1_{P_0/A}\otimes_{P_0}R
\longrightarrow0.
$$
Equivalently,
$$
L^{-n}_{R/A}=\Omega^1_{P_n/A}\otimes_{P_n}R,
\qquad n\geq0.
$$
\end{definition}

\noindent
Thus the construction is
$$
R\xleftarrow{\text{resolve}}P_\bullet
\xrightarrow{\Omega^1_{-/A}}\Omega^1_{P_\bullet/A}
\xrightarrow{\text{change coefficients to }R}L_{R/A}.
$$
In particular,
$$
H^0(L_{R/A})\cong\Omega^1_{R/A},
$$
while the negative degrees retain the
additional relation data discarded by
$\Omega^1_{R/A}$. See
\href{https://stacks.math.columbia.edu/%
tag/08QF}{Stacks, Tag 08QF}.


\medskip\noindent
The complex above is independent of the
chosen polynomial resolution up to
quasi-isomorphism. We now write $k$ for
the coefficient ring denoted by $A$ above.

\begin{definition}
\label{deformations-definition-cotangent-modules}
The \emph{cotangent cohomology modules}
of $R$ over $k$ are
$$
T^i_{R/k}
:=
\operatorname{Ext}^i_R(L_{R/k},R),
\qquad i\geq0.
$$
More generally, for an $R$-module $M$,
$$
T^i(R/k;M)
:=
\operatorname{Ext}^i_R(L_{R/k},M).
$$
\end{definition}

\noindent
Here Ext has a complex as its first
argument. Explicitly,
$$
\operatorname{Ext}^i_R(L_{R/k},M)
=
H^i\bigl(
\operatorname{RHom}_R(L_{R/k},M)
\bigr).
$$
For the polynomial model above, the
complex is bounded above and consists of
free $R$-modules. Thus one can compute
this by applying
$\operatorname{Hom}_R(-,M)$
to that complex and taking cohomology.
In particular,
$$
\begin{aligned}
T^1_{R/k}
&=\operatorname{Ext}^1_R(L_{R/k},R),\\
T^2_{R/k}
&=\operatorname{Ext}^2_R(L_{R/k},R).
\end{aligned}
$$

\medskip\noindent
The interpretation of $T^1_{R/k}$ as
first-order deformations is a theorem,
not part of its definition. The next
section develops that interpretation;
its proof is not yet complete.

\section{Why $T^1$ parametrizes
first-order deformations}
\label{deformations-section-first-order-algebra-deformations}

\noindent
We now start the algebraic deformation theory that will lead to the cotangent cohomology groups $T^1$ and $T^2$.

Let $k$ be a field and let
$$
D:=k[\varepsilon]/(\varepsilon^2)
$$
be the \emph{dual numbers}. The map
$$
D\longrightarrow k,
\qquad \varepsilon\longmapsto0,
$$
lets us recover the original object from a deformation by setting $\varepsilon=0$.

\begin{definition}
\label{deformations-definition-first-order-deformation-algebra}
Let $R$ be a $k$-algebra. A \emph{first-order deformation} of $R$ is a flat $D$-algebra $\widetilde R$ together with an isomorphism
$$
\widetilde R/\varepsilon\widetilde R\cong R.
$$
Thus reduction modulo $\varepsilon$ recovers the original algebra $R$.
\end{definition}

\noindent
Equivalently, $R$ is the special fibre of $\widetilde R$. Indeed, the $D$-algebra $\widetilde R$ corresponds to a family
$$
\operatorname{Spec}\widetilde R\longrightarrow\operatorname{Spec}D,
$$
and the fibre over $\varepsilon=0$ has coordinate ring
$$
\widetilde R\otimes_D k
\cong\widetilde R/\varepsilon\widetilde R
\cong R.
$$
Here tensoring with $k$ simply means setting
$\varepsilon=0$. Indeed,
$$
k=D/(\varepsilon),
$$
and we use the general identity
$$
M\otimes_D(D/I)\cong M/IM.
$$
This is Proposition
\ref{commutative-algebra-proposition-tensor-ideals} in
\emph{Commutative Algebra}: tensoring a module
with a quotient ring is the same as quotienting
the module by that ideal.
Taking
$$
M=\widetilde R,\qquad I=(\varepsilon)
$$
gives the displayed isomorphism.

\noindent
The deformation problem is to understand the different $D$-algebras $\widetilde R$ which reduce to $R$. In geometric language these are infinitesimal deformations over the dual numbers, but for now we work entirely with rings.

The quotient map therefore gives an exact sequence
$$
0\longrightarrow\varepsilon\widetilde R
\longrightarrow\widetilde R
\longrightarrow R\longrightarrow0.
$$
Flatness identifies the kernel with another copy of $R$. Indeed, start with
$$
0\longrightarrow(\varepsilon)\longrightarrow D
\longrightarrow k\longrightarrow0.
$$
Because $\widetilde R$ is flat over $D$, tensoring with $\widetilde R$ preserves injectivity and gives
$$
0\longrightarrow
(\varepsilon)\otimes_D\widetilde R
\longrightarrow\widetilde R
\longrightarrow k\otimes_D\widetilde R
\longrightarrow0.
$$
We already know
$$
k\otimes_D\widetilde R\cong R.
$$
Moreover,
$$
(\varepsilon)\cong k
$$
as $D$-modules, since $\varepsilon^2=0$. Hence
$$
(\varepsilon)\otimes_D\widetilde R
\cong k\otimes_D\widetilde R
\cong R.
$$
The image of the first map is precisely $\varepsilon\widetilde R$, so
$$
\varepsilon\widetilde R\cong R.
$$
Thus the kernel of the deformation sequence is isomorphic to $R$, but we keep the intrinsic notation $\varepsilon\widetilde R$ rather than replacing it by a symbolic copy of $R$. Finally, this kernel is square-zero because
$$
(\varepsilon\widetilde R)^2
\subseteq\varepsilon^2\widetilde R=0.
$$
Thus a first-order deformation gives a square-zero extension whose quotient is $R$ and whose kernel $\varepsilon\widetilde R$ is isomorphic to $R$.

\begin{exercise}
\label{deformations-exercise-module-extension-r-splits}
Suppose for a moment that we had an extension of $R$-modules
$$
0\longrightarrow R\longrightarrow E\longrightarrow R\longrightarrow0.
$$
Prove that it must split. Hint: $R$ is free of rank one as an $R$-module, hence projective. Conclude that
$$
\operatorname{Ext}^1_R(R,R)=0.
$$
Thus there is no interesting deformation information in extensions of $R$ by $R$ \emph{as $R$-modules}.
\end{exercise}

\noindent
Our actual deformation sequence
$$
0\longrightarrow\varepsilon\widetilde R
\longrightarrow\widetilde R
\longrightarrow R\longrightarrow0
$$
is not canonically an extension of $R$-modules: $\widetilde R$ is naturally a $D$-algebra, and there is in general no canonical map $R\to\widetilde R$. Nevertheless, as $k$-vector spaces the sequence splits; see Exercise
\ref{homological-algebra-exercise-split-short-exact-sequence-modules}
in \emph{Homological Algebra} for the equivalent characterizations of a split short exact sequence.

\begin{exercise}
\label{deformations-exercise-short-exact-vector-spaces-split}
Prove that every short exact sequence of $k$-vector spaces
$$
0\longrightarrow U\longrightarrow V\longrightarrow W\longrightarrow0
$$
splits. Hint: choose a basis of $W$ and lift its basis vectors to $V$.
\end{exercise}

\noindent
Therefore choose a $k$-linear section
$$
s:R\longrightarrow\widetilde R.
$$
We will not introduce notation such as $\varepsilon R$. We keep $\varepsilon\widetilde R$ as the actual kernel and use the isomorphism
$$
\theta:\varepsilon\widetilde R\xrightarrow{\sim}R
$$
obtained above whenever we want to identify an element of the kernel with an element of $R$. This lets us isolate where the genuinely new information can live: in the multiplication.

\begin{exercise}
\label{deformations-exercise-first-order-deformed-multiplication}
Choose such a $k$-linear section $s$. For $a,b\in R$, compare
$$
s(a)s(b)
$$
with
$$
s(ab).
$$
Show that
$$
s(a)s(b)-s(ab)
\in\ker(\widetilde R\to R)
=\varepsilon\widetilde R.
$$
Using the identification
$$
\varepsilon\widetilde R\cong R,
$$
this difference defines an element $\mu(a,b)\in R$, hence a map
$$
\mu:R\times R\longrightarrow R.
$$
Now use the ring axioms in $\widetilde R$ to determine the conditions on $\mu$. In particular, show from distributivity that $\mu$ is $k$-bilinear, from commutativity that
$$
\mu(a,b)=\mu(b,a),
$$
and from associativity that
$$
a\mu(b,c)-\mu(ab,c)+\mu(a,bc)-\mu(a,b)c=0.
$$
\end{exercise}

\begin{exercise}
\label{deformations-exercise-change-section-multiplication-correction}
The map $\mu$ depends on the chosen $k$-linear section
$$
s:R\longrightarrow\widetilde R.
$$
Choose another section $s'$. Show that
$$
s'(a)-s(a)\in\varepsilon\widetilde R
$$
for every $a\in R$, and hence, using $\varepsilon\widetilde R\cong R$, the change of section is encoded by a $k$-linear map
$$
h:R\longrightarrow R.
$$
If $\mu$ and $\mu'$ are the multiplication corrections obtained from $s$ and $s'$, show that
$$
\mu'(a,b)-\mu(a,b)
=
a\,h(b)-h(ab)+h(a)b,
$$
up to the corresponding sign convention for $h$.
\end{exercise}

\noindent
Thus a deformation is not represented intrinsically by one particular $\mu$. The ring axioms impose a cocycle condition on $\mu$, while changing the arbitrary section changes $\mu$ by a coboundary.

\begin{remark}[Relation with Hochschild cohomology]
The discussion above has produced the quotient
$$
\frac{
\{\mu\text{ satisfying the multiplication conditions}\}
}{
\{\text{changes coming from }h:R\to R\}
}.
$$
This is literally a cohomology group. Put
$$
C^1=\operatorname{Hom}_k(R,R),\qquad
C^2=\operatorname{Hom}_k(R\otimes_kR,R),\qquad
C^3=\operatorname{Hom}_k(R^{\otimes3},R).
$$
The relevant differentials are
$$
(dh)(a,b)=a\,h(b)-h(ab)+h(a)b
$$
and
$$
(d\mu)(a,b,c)
=a\mu(b,c)-\mu(ab,c)+\mu(a,bc)-\mu(a,b)c.
$$
Hence associativity says $d\mu=0$, while changing the section changes $\mu$ by $dh$. These are the Hochschild differentials, so for associative algebra deformations this quotient is $HH^2(R,R)$. Our algebras are commutative, so $\mu$ must also respect commutativity; the cotangent complex gives the intrinsic cohomology theory appropriate to commutative algebra deformations.
\end{remark}

\medskip\noindent
With the cotangent cohomology modules
from the preceding section defined,
we now relate $T^1_{R/k}$ to first-order
deformations. This argument is not yet
complete: it currently stops after
lifting $P_0$ to the extension.

\noindent
We now want to understand why $T^1$ has anything to do with deformations. Start with a square-zero extension of $k$-algebras
$$
0\longrightarrow M\longrightarrow\widetilde R
\xrightarrow{q}R\longrightarrow0,
\qquad M^2=0.
$$
We will imitate the argument for ordinary module extensions from \emph{Homological Algebra}. Take the polynomial resolution used to define the cotangent complex,
$$
P_\bullet\longrightarrow R.
$$
At level zero we have a polynomial $k$-algebra mapping onto $R$,
$$
P_0\longrightarrow R.
$$
Since $P_0$ is freely generated as a $k$-algebra, choose in $\widetilde R$ a lift of the image of each polynomial generator of $P_0$. This determines a $k$-algebra map
$$
\widetilde p_0:P_0\longrightarrow\widetilde R
$$
such that the diagram commutes, equivalently
$$
q\circ\widetilde p_0=p_0.
$$
This is the algebra analogue of lifting $P_0\to B$ to $P_0\to E$ in the ordinary $\operatorname{Ext}^1$ argument.

\section{Why $T^2$ contains
obstruction classes}
\label{deformations-section-obstruction-classes}

\section{Torus action and AC multidegrees}
\label{deformations-section-torus-action-ac-multidegrees}

\noindent
Our goal is to understand the multidegree
pieces of $T^2$ used by
Altmann and Christophersen (AC) in
\emph{Cotangent cohomology of
Stanley--Reisner rings} \cite{jan3}.
They compute $T^1$ and $T^2$ from the
simplicial complex. Before studying their
combinatorial formulas, we explain what
the pieces being computed mean:
they are weight spaces for the torus
that independently rescales the variables.
Thus the multidegree decomposition has a
representation-theoretic interpretation.
We will build this interpretation in
three steps.

\begin{enumerate}
\item The torus acts on the Stanley--Reisner
algebra by independently rescaling
the variables.
\item This action induces a linear
representation on $T^2$.
\item This representation decomposes into
weight spaces, which are precisely the
multidegree pieces.
\end{enumerate}

\noindent
We now explain the first item.
Put
$$
S=\mathbb{C}[x_1,\ldots,x_n],
\qquad T=(\mathbb{C}^{\times})^n.
$$
Independent rescaling means choosing
one nonzero scalar $t_i$ for each variable.
For $t=(t_1,\ldots,t_n)\in T$, define
the algebra automorphism
$$
\rho_t:S\longrightarrow S,
\qquad x_i\longmapsto t_i x_i.
$$
Its inverse is $\rho_{t^{-1}}$.
These automorphisms give a torus action:
$$
\rho_t\circ\rho_u=\rho_{tu},
\qquad \rho_{(1,\ldots,1)}=\operatorname{id}.
$$

\begin{proposition}
\label{deformations-proposition-monomial-torus-action}
For every monomial ideal $I\subset S$,
independent rescaling descends to a
torus action on the quotient algebra $S/I$.
\end{proposition}

\begin{proof}
Let $m=x_1^{a_1}\cdots x_n^{a_n}$ be
a monomial generator of $I$. Then
$$
\rho_t(m)=
(t_1^{a_1}\cdots t_n^{a_n})m.
$$
Thus $\rho_t(m)\in I$.
Every $f\in I$ is a finite sum
$$
f=\sum_j h_jm_j,
\qquad h_j\in S,
$$
where the $m_j$ are monomial generators.
Consequently,
$$
\rho_t(f)=
\sum_j\rho_t(h_j)\rho_t(m_j)\in I.
$$
This proves $\rho_t(I)\subseteq I$.
Applying the same argument to $t^{-1}$
gives the reverse inclusion, so
$$
\rho_t(I)=I.
$$

We can therefore define
$$
\overline{\rho}_t:S/I\longrightarrow S/I,
\qquad [f]\longmapsto[\rho_t(f)].
$$
This is well-defined: if $[f]=[g]$, then
$f-g\in I$, hence
$\rho_t(f)-\rho_t(g)\in I$.
It is an algebra automorphism with inverse
$\overline{\rho}_{t^{-1}}$.
The action identities on $S$ pass to
the quotient:
$$
\overline{\rho}_t\circ\overline{\rho}_u
=\overline{\rho}_{tu},
\qquad
\overline{\rho}_{(1,\ldots,1)}
=\operatorname{id}.
$$
\end{proof}

\noindent
A Stanley--Reisner ideal is monomial,
so the proposition applies to our algebra.


\medskip\noindent
For now we accept items 2 and 3:
the action on $R=S/I$ induces an
algebraic torus representation on $T^2$,
and its weight spaces are the AC
multidegree pieces. We postpone the
construction through the cotangent
complex.

The general torus decomposition is stated
in Theorem~\ref
{representation-theory-theorem-torus-weights}
of \emph{Representation Theory}.
That section also explains its relation
to Lie algebra weight spaces.
This is standard representation theory;
AC's contribution is the combinatorial
description of the individual cotangent
cohomology pieces, not the existence of
the weight decomposition.

Write $t\cdot v$ for the induced action
on $T^2$. What does a weight describe?
First keep two different objects apart.
An element
$$
t=(t_1,\ldots,t_n)\in T
$$
chooses the rescaling: multiply $x_i$
by $t_i$. A weight will instead record
how a vector responds to all these
choices.

Start with a familiar example in the
polynomial ring:
$$
\rho_t(x_1^2x_2)=t_1^2t_2x_1^2x_2.
$$
Every rescaling preserves the line
spanned by this monomial.
The scalar changes with $t$, but the
rule is always $t_1^2t_2$.
We label this rule by the exponent
vector $(2,1,0,\ldots,0)$.

Now suppose a nonzero vector $v\in T^2$
also stays on its own line under every
rescaling. Write
$$
t\cdot v=\chi(t)v.
$$
The function $\chi$ records the scalar
for every choice of $t$.
Acting by $u$ and then by $t$ is the
same as acting by $tu$, so
$$
\chi(tu)v
=t\cdot(u\cdot v)
=\chi(u)\chi(t)v.
$$
Since $v\neq0$, we obtain
$$
\chi(tu)=\chi(t)\chi(u).
$$
The scalars are nonzero because the
action is invertible. Such an algebraic
homomorphism
$\chi:T\to\mathbb{C}^{\times}$
is called a \emph{character}.
Here this means a one-dimensional
representation; it agrees with the
trace character in dimension one.

Why do integer vectors appear?
The algebraic characters of $T$ are
precisely the Laurent monomials
$$
\chi_\alpha(t)=t^\alpha
=\prod_{i=1}^n t_i^{\alpha_i},
\qquad \alpha\in\mathbb{Z}^n.
$$
This classification is a standard fact;
see \cite[Section 4g and Chapter 12,
Sections a--d]{milne-ag-2022}.
The exponents are integers, and negative
ones are allowed since each $t_i$ is
nonzero. Thus $t$ specifies a
transformation, while $\alpha$ labels
a scaling rule. They belong to
different sets.

We group together all vectors with
the same scaling rule:
$$
T^2_\alpha=
\{v\in T^2:t\cdot v=t^\alpha v
\ \forall t\in T\}.
$$
This space can have dimension greater
than one: several independent vectors
can respond by the same scalar.
The theorem we are accepting says that
every vector is a unique finite sum of
vectors belonging to these spaces:
\begin{equation}
\label{deformations-equation-ac-weight-decomposition}
T^2=\bigoplus_{\alpha\in\mathbb{Z}^n}
T^2_\alpha.
\end{equation}
This is \cite[Theorem 12.12]
{milne-ag-2022}, applied to the
algebraic torus representation on $T^2$.
These are the pieces computed by AC.

The following comparison is only a bridge
to the weight decompositions we already
know from Lie algebras. It is not needed
for the deformation argument and may
be skipped.
To recognize the Lie algebra version,
put $t_i=e^{s h_i}$ and differentiate
at $s=0$. For $v\in T^2_\alpha$,
$$
t(s)\cdot v
=e^{s\sum_i\alpha_i h_i}v
\quad\Longrightarrow\quad
\pi(h)v=\left(\sum_i\alpha_i h_i\right)v.
$$
Thus the Lie algebra weight is the
linear functional
$\lambda_\alpha(h)=\sum_i\alpha_i h_i$.
This is the simultaneous-eigenvector
condition from
Definition~\ref
{lie-algebras-definition-weight-space}.

For the graded projective problem we
use the total-degree-zero part:
$$
(T^2)_0=
\bigoplus_{\sum_i\alpha_i=0}T^2_\alpha.
$$
Indeed, setting all $t_i=c$ gives the
scalar $c^{\sum_i\alpha_i}$.
Total degree zero means that this
common rescaling acts trivially.

There is a second compatibility that
we also accept for now. Let $G$ be a
group of vertex permutations preserving
the simplicial complex; in our
Gr\"unbaum argument this is the chosen
subgroup isomorphic to $S_3$.
Use the convention $g(x_i)=x_{g(i)}$.
Then
\begin{equation}
\label{deformations-equation-ac-weight-permutation}
g(T^2_\alpha)=T^2_{g\alpha},
\qquad
(g\alpha)_i=\alpha_{g^{-1}(i)}.
\end{equation}
This says that relabelling the vertices
also relabels the entries of the weight.
For example, swapping vertices 1 and 2
sends the space of weight
$(a,b,\alpha_3,\ldots,\alpha_n)$
to the space of weight
$(b,a,\alpha_3,\ldots,\alpha_n)$.
It need not preserve each weight space
separately.

The reason is that a permutation
conjugates independent rescaling to
independent rescaling with permuted
scalars. The induced actions on $T^2$
retain this compatibility.
Consequently we can organize the AC
pieces into $G$-orbits and study one
representative at a time.
Permutations fixing a weight can still
act nontrivially inside its weight
space; understanding that action will
be essential for finding invariant
obstruction classes.

\section{Structures over the Dual Numbers}
\label{deformations-section-structures-dual-numbers}

\noindent
Other than the contents of
\cite[Chapter 1, Section 2]
{Hartshorne-deformation}, which
is most of the content of this section,
let us recall from Ariel's
lecture that he proved a proposition
with 4 equivalent conditions
for flatness over the dual numbers:
being a free module, having
$\operatorname{Tor}^1=0$, and two more
involving annihilators.

\begin{definition}
\label{deformations-definition-dual-numbers-first-order}
The very first deformation question to study
is structures over the \emph{dual numbers}
$$
D=k[t]/t^2.
$$
That is, one gives a structure (e.g., a
scheme, or a scheme with a subscheme, or a
scheme with a sheaf on it) over $k$ and one
seeks to classify extensions of this
structure over the dual numbers. These are
also called \emph{first-order deformations}.
To ensure that our structure is evenly
spread out over the base, we will always
assume that the extended structure is flat
over $D$. Flatness is the technical condition
that corresponds to the intuitive idea of a
deformation.
\end{definition}

\begin{lemma}
\label{deformations-lemma-flatness-tor}
A module $M$ over a noetherian ring $A$ is
flat if and only if for every prime ideal
$\mathfrak{p}\subseteq A$,
$$
\operatorname{Tor}^A_1
(M,A/\mathfrak{p})=0.
$$
\end{lemma}

\begin{proposition}
\label{deformations-proposition-square-zero-flatness}
Let $A'\to A$ be a surjective homomorphism
of noetherian rings whose kernel $J$ has
square zero. Then an $A'$-module $M'$ is flat
over $A'$ if and only if
\begin{enumerate}
\item
$$
M=M'\otimes_{A'}A
$$
is flat over $A$, and
\item the natural map
$$
M\otimes_AJ\longrightarrow M'
$$
is injective.
\end{enumerate}
\end{proposition}

\begin{definition}
\label{deformations-definition-embedded-first-order}
Let $X$ be a scheme over $k$ and let $Y$ be
a closed subscheme of $X$. We define a
\emph{deformation of $Y$ over $D$ in $X$} to
be a closed subscheme
$$
Y'\subseteq X'=X\times D,
$$
flat over $D$, such that
$$
Y'\times_D k=Y.
$$
\end{definition}

We consider the affine case first. Then $X$
corresponds to a $k$-algebra $B$, and $Y$ is
defined by an ideal $I\subseteq B$. We are
seeking ideals
$$
I'\subseteq B'=B[t]/t^2
$$
with $B'/I'$ flat over $D$ and such that the
image of $I'$ in $B=B'/tB'$ is just $I$.
Note that
$$
(B'/I')\otimes_Dk=B/I.
$$
Since $B$ is automatically flat over $k$,
by (2.2) the flatness of
$B'/I'$ over $D$ is equivalent to the
exactness of the sequence
$$
0\longrightarrow B/I
\xrightarrow{t}B'/I'
\longrightarrow B/I
\longrightarrow 0.
$$

\begin{proposition}
\label{deformations-proposition-affine-first-order}
In the situation above, to give
$I'\subseteq B'$ such that $B'/I'$ is flat
over $D$ and the image of $I'$ in $B$ is $I$
is equivalent to giving an element
$$
\varphi\in\operatorname{Hom}_B(I,B/I).
$$
In particular, $\varphi=0$ corresponds to the
trivial deformation given by
$$
I'=I\oplus tI
$$
inside $B'\cong B\oplus tB$.
\end{proposition}

\begin{definition}
\label{deformations-definition-normal-sheaf}
Before stating the conclusion, we will
define the \emph{normal sheaf of $Y$ in
$X$}. Note that the group
$$
\operatorname{Hom}_X
(\mathcal{I},\mathcal{O}_Y)
$$
can be regarded as
$$
H^0\bigl(X,
\mathcal{H}om_X
(\mathcal{I},\mathcal{O}_Y)\bigr).
$$
Furthermore, homomorphisms of $\mathcal{I}$
to $\mathcal{O}_Y$ factor through
$\mathcal{I}/\mathcal{I}^2$, which is a sheaf
on $Y$. So
$$
\mathcal{H}om_X
(\mathcal{I},\mathcal{O}_Y)
=
\mathcal{H}om_Y
(\mathcal{I}/\mathcal{I}^2,
\mathcal{O}_Y),
$$
and this latter sheaf is called the normal
sheaf of $Y$ in $X$, and is denoted by
$\mathcal{N}_{Y/X}$. If $X$ is nonsingular
and $Y$ is a locally complete intersection
in $X$, then $\mathcal{I}/\mathcal{I}^2$ is
locally free, so $\mathcal{N}_{Y/X}$ is
locally free also and can be called the
\emph{normal bundle of $Y$ in $X$}. This
terminology derives from the fact that if
$Y$ is also nonsingular, there is an exact
sequence
$$
0\longrightarrow\mathcal{T}_Y
\longrightarrow\mathcal{T}_X|_Y
\longrightarrow\mathcal{N}_{Y/X}
\longrightarrow 0,
$$
where $\mathcal{T}_Y$ and $\mathcal{T}_X$
denote the tangent sheaves to $Y$ and $X$,
respectively. In this case, therefore,
$\mathcal{N}_{Y/X}$ is the usual normal
bundle.
\end{definition}

\begin{theorem}
\label{deformations-theorem-embedded-first-order}
Let $X$ be a scheme over a field $k$, and let
$Y$ be a closed subscheme of $X$. Then the
deformations of $Y$ over $D$ in $X$ are in
natural one-to-one correspondence with
elements of $H^0(Y,\mathcal{N}_{Y/X})$, the
zero element corresponding to the trivial
deformation.
\end{theorem}

\begin{lemma}[Corollary]
\label{deformations-lemma-hilbert-tangent-first-order}
If $Y$ is a closed subscheme of the
projective space $X=\mathbb{P}^n_k$, then the
Zariski tangent space of the Hilbert scheme
$H$ at the point $y$ corresponding to $Y$ is
isomorphic to $H^0(Y,\mathcal{N}_{Y/X})$.
\end{lemma}

Next we consider Situation B. Let $X$ be a
scheme over $k$ and let $L$ be an invertible
sheaf on $X$. We will study the set of
isomorphism classes of invertible sheaves
$L'$ on
$$
X'=X\times D
$$
such that $L'\otimes\mathcal{O}_X\cong L$.
In this case flatness is automatic, because
$L'$ is locally free and $X'$ is flat over
$D$.

\begin{proposition}
\label{deformations-proposition-line-bundle-first-order}
Let $X$ be a scheme over $k$, and $L$ an
invertible sheaf on $X$. The set of
isomorphism classes of invertible sheaves
$L'$ on $X\times D$ such that
$L'\otimes\mathcal{O}_X\cong L$ is in natural
one-to-one correspondence with elements of
the group $H^1(X,\mathcal{O}_X)$.
\end{proposition}

\begin{definition}
\label{deformations-definition-coherent-first-order}
Let $X$ be a scheme over $k$, and let $F$ be
a coherent sheaf on $X$. We define a
\emph{deformation of $F$ over $D$} to be a
coherent sheaf $F'$ on $X'=X\times D$, flat
over $D$, together with a homomorphism
$F'\to F$ such that the induced map
$$
F'\otimes_D k\longrightarrow F
$$
is an isomorphism. We say that two such
deformations $F'_1\to F$ and $F'_2\to F$ are
\emph{equivalent} if there is an isomorphism
$$
F'_1\xrightarrow{\sim}F'_2
$$
compatible with the given maps to $F$.
\end{definition}

\begin{theorem}
\label{deformations-theorem-coherent-first-order}
Let $X$ be a scheme over $k$, and let $F$ be
a coherent sheaf on $X$. The (equivalence
classes of) deformations of $F$ over $D$ are
in natural one-to-one correspondence with
the elements of the group
$\operatorname{Ext}^1_X(F,F)$, where the
zero-element corresponds to the trivial
deformation.
\end{theorem}

\begin{remark}
\label{deformations-remark-unmarked-coherent-first-order}
Given $F$ on $X$, we can also pose a
different problem, like the one in (2.6),
namely to classify isomorphism classes of
coherent sheaves $F'$ on $X'$, flat over
$D$, such that $F'\otimes_D k$ is isomorphic
to $F$ (without specifying the isomorphism).
This set need not be the same as the set of
deformations of $F$, but we can explain
their relationship as follows. The group
$\operatorname{Aut}F$ of automorphisms of
$F$ acts on the set of deformations of $F$
by letting $\alpha\in\operatorname{Aut}F$
applied to $f:F'\to F$ be
$\alpha f:F'\to F$. Now let $f:F'\to F$ and
$g:F''\to F$ be two deformations of $F$. One
sees easily that $F'$ and $F''$ are
isomorphic as sheaves on $X'$ if and only if
there exists an
$\alpha\in\operatorname{Aut}F$ such that
$\alpha f$ and $g$ are equivalent as
deformations of $F$. Thus the set of $F'$'s
up to isomorphism as sheaves on $X'$ is the
orbit space of $\operatorname{Ext}^1_X(F,F)$
under the action of $\operatorname{Aut}F$.
This kind of subtle distinction will play an
important role in questions of
pro-representability (Chapter 3).
\end{remark}

\begin{lemma}[Corollary]
\label{deformations-lemma-vector-bundle-first-order}
If $E$ is a vector bundle over $X$, then the
deformations of $E$ over $D$ are in natural
one-to-one correspondence with the elements
of $H^1(X,\mathcal{E}nd E)$, where
$\mathcal{E}nd E=\mathcal{H}om(E,E)$ is the
sheaf of endomorphisms of $E$. The trivial
deformation corresponds to the zero element.
\end{lemma}



\section{Deformations of schemes}
\label{deformations-section-deformations-of-schemes}

\noindent
The following definition is from
[lucas-defos], which in turn comes from
\cite{Sernesi-deformations}

\begin{definition}
\label{deformations-definition-deformation}
Let $X$ be an algebraic $\mathbb{C}$-scheme.
\begin{enumerate}
\item A {\it deformation} of $X$ is a
Cartesian diagram $\xi$
$$
\xymatrix{
X\ar[r]\ar[d]&\mathcal{X}\ar[d]^{\pi}\\
\Spec \mathbb{C}\ar[r]^s&S
}
$$
where $\pi$ is a flat surjective morphism of
algebraic $\mathbb{C}$-schemes and
$S$ is connected. (Recall that flatness
accounts for ``continuity''.)
\item A {\it local deformation} of $X$ is a
deformation $\xi$ where
$S=\Spec A$ for $A$ a noetherian local
$\mathbb{C}$ algebra with residue
field $\mathbb{C}$.
\item An {\it infinitesimal deformation of
$X$} is a local deformation with $A$
an artinian local $\mathbb{C}$-algebra with
residue field $\mathbb{C}$. $X$ is
called {\it rigid} if all infinitesimal
deformations are trivial.
\item An {\it inifinitesimal deformation of
order $n$} is an infinitesimal
deformation when $S=\Spec
(\mathbb{C}[t]/(t^{n+1})$.
\end{enumerate}
\end{definition}


\section{The deformation functor
\texorpdfstring
{$\operatorname{Def}_X$}{Def X}}
\label{deformations-section-deformation-functor}

\noindent
The correct object parametrizing abstract
infinitesimal deformations is initially a
functor, not a scheme.

Fix a field $k$ and a $k$-scheme $X$. Let
$\mathbf{Art}_k$ denote the category of
local Artinian $k$-algebras whose residue
field is $k$. Its morphisms are local
$k$-algebra homomorphisms.

\begin{definition}
\label{deformations-definition-deformation-artin}
Let $A\in\mathbf{Art}_k$. A
\emph{deformation of $X$ over $A$} is a
flat morphism
$$
\pi:\mathcal{X}_A\longrightarrow
\operatorname{Spec}A
$$
together with an isomorphism, called the
marking of the special fiber,
$$
\iota:\mathcal{X}_A
\times_{\operatorname{Spec}A}
\operatorname{Spec}k\xrightarrow{\sim}X.
$$
An isomorphism between two such deformations
is an isomorphism of $A$-schemes compatible
with their markings.
\end{definition}

\begin{definition}
\label{deformations-definition-def-x}
The \emph{deformation functor of $X$} is
defined as follows. For each
$A\in\mathbf{Art}_k$, let $\mathcal{D}_X(A)$
be the set of pairs $(\pi,\iota)$ in the
previous definition. Its value at $A$ is
$$
\operatorname{Def}_X(A)
=
\mathcal{D}_X(A)/{\cong}.
$$
These sets form the functor
$$
\begin{aligned}
\operatorname{Def}_X:
\mathbf{Art}_k&\longrightarrow
\mathbf{Sets},\\
A&\longmapsto
\operatorname{Def}_X(A).
\end{aligned}
$$
If $A\to B$ is a homomorphism, the
corresponding map
$\operatorname{Def}_X(A)\to
\operatorname{Def}_X(B)$ is given by base
change:
$$
\mathcal{X}_A\longmapsto
\mathcal{X}_A\times_{\operatorname{Spec}A}
\operatorname{Spec}B.
$$
\end{definition}

\noindent
A richer version assigns to $A$ the
\emph{groupoid} of deformations and their
isomorphisms. Passing to the set of
isomorphism classes forgets the automorphisms
of each deformation. Infinitesimal
automorphisms can prevent the resulting
set-valued functor from being representable.

\medskip\noindent
\textbf{Comparison with the Hilbert
functor.}
If $X\subset Y$ is a fixed embedding, the
local Hilbert functor parametrizes flat
deformations
$$
\mathcal{X}_A\subset Y\times
\operatorname{Spec}A
$$
of $X$ \emph{inside the fixed ambient
scheme $Y$}. There is a natural
transformation
$$
\operatorname{Def}_{X\subset Y}
\longrightarrow
\operatorname{Def}_X
$$
which forgets the embedding. The functor on
the left describes embedded deformations;
the one on the right describes abstract
deformations.

In our Stanley--Reisner problem, perturbing
the homogeneous generators of
$I_X\subset k[x_0,\ldots,x_7]$ initially
studies embedded deformations of
$X\subset \mathbb{P}^7$. It belongs first
to the local Hilbert problem, although
one can subsequently forget the embedding
and obtain an abstract deformation.

\medskip
\noindent
\textbf{Universal and miniversal formal
deformations.}

\noindent
Let $R$ be a complete local Noetherian
$k$-algebra with residue field $k$. It
defines a functor. Let $\mathbf{C}_k$ denote
the category of such algebras and local
$k$-algebra homomorphisms. We regard every
object of $\mathbf{Art}_k$ as an object of
$\mathbf{C}_k$. Then
$$
h_R(A)=\operatorname{Hom}_{\mathbf{C}_k}
(R,A).
$$
Geometrically, this is the functor associated
to the formal scheme $\operatorname{Spf}R$.

If there is a natural isomorphism
$$
h_R\xrightarrow{\sim}
\operatorname{Def}_X,
$$
then $\operatorname{Def}_X$ is
\emph{prorepresentable}. Equivalently,
$\operatorname{Spf}R$ carries a universal
formal deformation: every infinitesimal
deformation is obtained by pullback along a
unique homomorphism $R\to A$.

In general this is too much to expect. One
therefore asks for a natural transformation
$$
h_R\longrightarrow\operatorname{Def}_X
$$
which is formally smooth and induces an
isomorphism on tangent spaces. Formal
smoothness means that compatible data over
the two rings in a small extension can be
lifted to $h_R$ over the larger ring.
Concretely, for every small extension
$B\longrightarrow A$, the map
$$
h_R(B)\longrightarrow
h_R(A)\times_{\operatorname{Def}_X(A)}
\operatorname{Def}_X(B)
$$
is required to be surjective.

The tangent space of the deformation functor
is
$$
T\operatorname{Def}_X=
\operatorname{Def}_X
\big(k[\varepsilon]/(\varepsilon^2)\big).
$$
Thus the minimality condition says that
$$
T_o\operatorname{Spf}R
\xrightarrow{\sim}
T\operatorname{Def}_X.
$$

\begin{definition}
\label{deformations-definition-miniversal-deformation}
A formal deformation over
$\operatorname{Spf}R$ satisfying the two
conditions above is called a
\emph{miniversal deformation}. The formal
scheme $\operatorname{Spf}R$ is called a
\emph{miniversal deformation space}, a
\emph{semiuniversal deformation space}, or
a \emph{hull} of $\operatorname{Def}_X$.
\end{definition}

\noindent
The crucial distinction is that
$\operatorname{Def}_X$ is the functor,
whereas $\operatorname{Spf}R$, together
with the natural transformation above, is
a miniversal space for it.
A miniversal family contains all
infinitesimal deformation directions and has
no redundant tangent parameters, but the map
producing a given deformation need not be
unique. If it were unique for every Artinian
base, the family would be universal and the
functor would be prorepresentable.

Schlessinger's theorem gives conditions under
which a functor on $\mathbf{Art}_k$ has such
a hull, and stronger conditions under which
it is prorepresentable. Once its lifting
conditions hold, finite-dimensionality of
the tangent space is the key finiteness
condition for the existence of a hull.
Proper schemes in standard finite-type
settings have this finiteness. Infinitesimal
automorphisms can prevent a hull from being
universal.

The cotangent complex gives the standard
tangent and obstruction spaces
$$
T^1_X=\operatorname{Ext}^1
(\mathbb{L}_{X/k},\mathcal{O}_X),
\qquad
T^2_X=\operatorname{Ext}^2
(\mathbb{L}_{X/k},\mathcal{O}_X).
$$
The second group is an obstruction space:
obstructions to lifting lie in it, although
not every element must occur as an
obstruction.

Suppose that $T^1_X$ is finite-dimensional
and that a hull exists. A minimal
presentation of its complete local ring has
the form
$$
R\cong k[[t_1,\ldots,t_r]]/J,
\qquad
J\subset (t_1,\ldots,t_r)^2,
\qquad
r=\dim_kT^1_X.
$$
The ideal $J$ records the higher-order
lifting conditions detected by the
obstruction theory. Its lack of linear terms
expresses the tangent-space minimality of the
presentation.

Finally, a formal one-parameter deformation
is only a formal arc in this miniversal
space:
$$
\operatorname{Spf}k[[s]]
\longrightarrow
\operatorname{Spf}R.
$$
Concretely, an arc through the origin is a morphism
$$
\gamma:\operatorname{Spf}k[[t]]
\longrightarrow
\operatorname{Spf}R
$$
sending the closed point $t=0$ to the point representing
$X$. Since morphisms of affine formal schemes reverse the
arrows on rings, this is equivalently a continuous local
$k$-algebra homomorphism
$$
\gamma^\#:R\longrightarrow k[[t]].
$$
If
$$
R\cong k[[x_1,\ldots,x_r]]/J,
$$
then specifying the arc amounts to choosing power series
$$
x_i\longmapsto x_i(t)\in t\,k[[t]]
$$
which satisfy the equations in $J$. Thus the arc selects a
one-dimensional formal path through the possibly much
larger deformation space. Pulling the miniversal family
back along $\gamma$ produces a formal one-parameter
deformation over $\operatorname{Spf}k[[t]]$. In particular,
the arc is a device for passing from the full deformation
space to a one-parameter deformation; it is not a section
of the total family.

In our problem the Maurer--Cartan
construction seeks a compatible deformation
along such an arc. If that formal deformation
is algebraized, one can then ask whether its
generic fiber is smooth. The construction
does not by itself produce the whole
miniversal deformation space.

\section{The Maurer--Cartan equation in complex geometry}
\label{deformations-section-maurer-cartan-complex-geometry}

\noindent
Let $X$ be a complex manifold and let $M$ be its
underlying real manifold. The complex structure is
equivalently encoded by the splitting
$$
T_{\mathbb C}M=T^{1,0}X\oplus T^{0,1}X.
$$
Thus, to deform the complex structure we may deform
the subbundle $T^{0,1}X$ inside $T_{\mathbb C}M$.

A sufficiently small deformation which stays transverse
to $T^{1,0}X$ is the graph of a bundle map
$$
\phi:T^{0,1}X\longrightarrow T^{1,0}X.
$$
Equivalently,
$$
\phi\in A^{0,1}(T^{1,0}X),
$$
since
$$
\operatorname{Hom}(T^{0,1}X,T^{1,0}X)
\cong (T^{0,1}X)^*\otimes T^{1,0}X.
$$
The deformed antiholomorphic tangent bundle is
$$
T^{0,1}_\phi
=\{v+\phi(v):v\in T^{0,1}X\}.
$$

Why does such a map really describe a deformation of the
subbundle? This is simply fiberwise linear algebra. Suppose
$$
E=V\oplus W
$$
and let $W'\subset E$ be another subspace of the same
dimension as $W$, transverse to $V$. Then the projection
$$
\pi_W:E\longrightarrow W
$$
restricts to an isomorphism
$$
\pi_W|_{W'}:W'\xrightarrow{\sim}W.
$$
Therefore for every $w\in W$ there is a unique vector of
$W'$ projecting to $w$. It has the form
$$
w+\phi(w)
$$
for a unique $\phi(w)\in V$. Thus there is a unique linear
map
$$
\phi:W\longrightarrow V
$$
such that
$$
W'=\{w+\phi(w):w\in W\}=\operatorname{graph}(\phi).
$$

Applying this at each $x\in X$ with
$$
E=T_{\mathbb C,x}M,\qquad
V=T^{1,0}_xX,\qquad
W=T^{0,1}_xX,
$$
shows that every nearby $n$-plane transverse to
$T^{1,0}_xX$ is uniquely the graph of a map
$$
\phi_x:T^{0,1}_xX\longrightarrow T^{1,0}_xX.
$$
If these planes vary smoothly with $x$, the maps $\phi_x$
assemble into the bundle map $\phi$ above. In other words,
$$
\operatorname{Hom}(T^{0,1}_xX,T^{1,0}_xX)
$$
gives local coordinates near $T^{0,1}_xX$ in the
Grassmannian of $n$-planes in $T_{\mathbb C,x}M$. Thus
$\phi$ is literally the slope of the deformed subbundle
relative to the original splitting.

Here ``sufficiently small'' just guarantees that
$T^{0,1}_\phi$ and its complex conjugate remain
transverse, so that they define an almost complex
structure. For example, after choosing a Hermitian
metric and assuming $X$ compact, it is enough to take
$\phi$ sufficiently small in the $C^0$ norm. Stronger
topologies such as $C^k$, H\"older or Sobolev topologies
are useful when studying the differential equation below.

The remaining question is whether this almost complex
structure is integrable. By the Newlander--Nirenberg
theorem, integrability is equivalent to
$$
[T^{0,1}_\phi,T^{0,1}_\phi]
\subseteq T^{0,1}_\phi.
$$

To compute this condition, choose local holomorphic
coordinates $z^1,\ldots,z^n$ for the original complex
structure and write
$$
\phi=
\phi^k_{\bar i}\,d\bar z^i\otimes
\frac{\partial}{\partial z^k}.
$$
Then $T^{0,1}_\phi$ is locally spanned by
$$
e_{\bar i}
=\frac{\partial}{\partial\bar z^i}
+\phi^k_{\bar i}\frac{\partial}{\partial z^k}.
$$
A direct computation gives
$$
[e_{\bar i},e_{\bar j}]
=
\left(
\frac{\partial\phi^k_{\bar j}}{\partial\bar z^i}
-\frac{\partial\phi^k_{\bar i}}{\partial\bar z^j}
+\phi^\ell_{\bar i}
\frac{\partial\phi^k_{\bar j}}{\partial z^\ell}
-\phi^\ell_{\bar j}
\frac{\partial\phi^k_{\bar i}}{\partial z^\ell}
\right)
\frac{\partial}{\partial z^k}.
$$
The first two terms are the components of
$\bar\partial\phi$, while the last two are the
components of $\frac12[\phi,\phi]$, where the bracket
is the natural extension of the Lie bracket of vector
fields to $T^{1,0}X$-valued forms. Hence the
integrability condition is precisely
$$
\boxed{
\bar\partial\phi+\frac12[\phi,\phi]=0.
}
$$
This is the \emph{Maurer--Cartan equation}.

Thus the differential graded Lie algebra
$$
A^{0,\bullet}(T^{1,0}X),
\qquad d=\bar\partial,
$$
with this bracket naturally appears when one studies
deformations of complex structures: its
Maurer--Cartan elements are precisely the small
integrable deformations written in these graph
coordinates.

Up to this point, $\phi$ has represented one single
nearby complex structure: its graph is the deformed
subbundle $T^{0,1}_\phi$. There has been no deformation
parameter $t$ yet.

To discuss first order, second order and lifting, we now
make an additional choice: consider a one-parameter family
of deformations passing through the original complex
structure,
$$
t\longmapsto \phi_t,
\qquad
\phi_0=0.
$$
Thus we regard the family as a curve,
that is, as vector-valued 
map of a real parameter $t$, 
with values in the
vector space
$$
A^{0,1}(T^{1,0}X)
=\Gamma\!\left(\operatorname{Hom}
(T^{0,1}X,T^{1,0}X)\right).
$$
If this family depends analytically on $t$, we may Taylor
expand it at $t=0$. Absorbing the usual factorials into the
coefficients, we write
$$
\phi(t)=t\phi_1+t^2\phi_2+t^3\phi_3+\cdots.
$$
In formal deformation theory we instead take such a formal
power series as the object from the start. The condition
$\phi(0)=0$ explains why there is no constant term: at
$t=0$ we recover the original complex structure.

For this one-parameter formal deformation, the complete
local base ring is
$$
\mathbb C[[t]]
=\varprojlim_n \mathbb C[t]/(t^{n+1}).
$$
Thus a formal deformation over $\mathbb C[[t]]$ means a
compatible system of finite-order deformations over
$$
\operatorname{Spec}\big(\mathbb C[t]/(t^{n+1})\big),
$$
or equivalently a deformation over the formal base
$\operatorname{Spf}\mathbb C[[t]]$. The series $\phi(t)$
above is therefore naturally a power series in the
coordinate $t$ of this formal base.

Here ``first order'' means the coefficient of $t^1$ after
expanding the Maurer--Cartan equation as a formal power
series in $t$. Indeed,
$$
\bar\partial\phi(t)
=t\bar\partial\phi_1+t^2\bar\partial\phi_2+\cdots,
$$
while bilinearity of the bracket gives
$$
[\phi(t),\phi(t)]
=t^2[\phi_1,\phi_1]
+2t^3[\phi_1,\phi_2]+\cdots.
$$
Thus
$$
t\bar\partial\phi_1
+t^2\left(
\bar\partial\phi_2+\frac12[\phi_1,\phi_1]
\right)
+\cdots=0.
$$
A formal power series is zero if and only if each of its
coefficients is zero. Hence
$$
\boxed{\text{order }t:\quad \bar\partial\phi_1=0}
$$
and
$$
\boxed{
\text{order }t^2:\quad
\bar\partial\phi_2+\frac12[\phi_1,\phi_1]=0.
}
$$
Equivalently, working to first order means working modulo
$t^2$. Then
$$
\phi(t)\equiv t\phi_1\pmod{t^2},
$$
and the bracket term disappears because it is already of
order $t^2$. This is why first-order deformation theory is
linear, whereas the first obstruction appears at second
order.

Thus the bracket produces the first quadratic
obstruction: {\bf the first-order direction $\phi_1$ lifts
to second order only if
$\frac12[\phi_1,\phi_1]$ is
$\bar\partial$-exact}.
The corresponding obstruction is therefore
the cohomology class
$$
\left[\frac{1}{2}[\phi_1,\phi_1]\right]
\in
H^{0,2}(X,T^{1,0}X).
$$
If that class is zero, then we have
Maurer-Cartan and we can choose $\phi_2$
to lift $\phi_1$ to second order.
But if this class is nonzero, we are stuck.


After quotienting first-order deformations by
infinitesimal changes of coordinates, their classes lie
in
$$
H^{0,1}(X,T^{1,0}X)\cong H^1(X,T_X),
$$
which is exactly the space that appears in the
Kodaira--Spencer map below.

The isomorphism above is simply the Dolbeault theorem for
the holomorphic vector bundle $T_X$. The sheaf $T_X$ of
holomorphic vector fields has the Dolbeault resolution
$$
0\longrightarrow T_X
\longrightarrow \mathcal A^{0,0}(T^{1,0}X)
\xrightarrow{\bar\partial}
\mathcal A^{0,1}(T^{1,0}X)
\xrightarrow{\bar\partial}
\mathcal A^{0,2}(T^{1,0}X)
\xrightarrow{\bar\partial}\cdots .
$$
This is an exact resolution of $T_X$, and the sheaves
$\mathcal A^{0,q}(T^{1,0}X)$ are fine, hence acyclic for
global sections. Therefore the sheaf cohomology of $T_X$
is computed by the cohomology of the complex of global
sections
$$
A^{0,0}(T^{1,0}X)
\xrightarrow{\bar\partial}
A^{0,1}(T^{1,0}X)
\xrightarrow{\bar\partial}
A^{0,2}(T^{1,0}X)
\longrightarrow\cdots .
$$
Consequently,
$$
H^q(X,T_X)\cong H^{0,q}_{\bar\partial}
(X,T^{1,0}X),
$$
and in particular
$$
H^1(X,T_X)\cong H^{0,1}_{\bar\partial}
(X,T^{1,0}X).
$$
So the same quotient that appears naturally from the
Maurer--Cartan linearization---$\bar\partial$-closed
first-order directions modulo $\bar\partial$-exact
infinitesimal coordinate changes---is precisely the usual
sheaf cohomology group $H^1(X,T_X)$.

\medskip\noindent
The point to retain is that the Maurer--Cartan equation
is not an extra formal gadget imposed on deformation
theory. Already for ordinary complex structures it is
simply the integrability condition for the deformed
object. Its linear term describes first-order
deformations, its bracket gives the quadratic
obstruction, and the full equation organizes lifting to
all orders.

\section{Kodaira-Spencer map}
\label{deformations-section-kodaira-spencer-map}

\noindent
Let $\mathcal{X} \to B$ be a
deformation of $X$.
Kodaira-Spencer map is
$$
T_pB \to H^1(X,T_X)
$$
constructed as follows.

\begin{enumerate}
\item Since $\pi$ is a submersion,
we can take submersion charts all along
the fiber $X = \pi^{-1}(p)$. This gives
a cover $\mathcal{U}_i\simeq U_i \times V_i$ 
where
$V_i$ is the part of the chart taken
using the coordinates of $B$
(this is just inverse function
theorem in its level-set chart theorem
incarnation).

\item In the intersections we have
transition functions:
$$
\xymatrix{
\mathcal{U}_i \cap \mathcal{U}_j
\ar[r]^{\varphi_i}& 
U_i \times V \ar[d]^{\varphi_{ij}}\\
\mathcal{U}_i \cap\mathcal{U}_j
\ar[r]_{\varphi_j}&
U_j \times V
}
$$
where $V:= V_i \cap V_j$ may thought
of as a neighbourhood of $B$
which must contain $p$.
This means that we can apply the
Regular Level-set theorem using the
same coordinates for the $V$-part
of the trivializations.

Then the transition functions
$\varphi_{ij}$ look like this:
$(\psi_{ij}(x,t),t)$ 
where $x$ is a variable
in $X = \pi^{-1}(p)$, 
$t$ is a variable in $B$
and $\psi_{ij}$ is a transition map from the
original atlas of $\mathcal{X}$.

\item Taking a tangent vector $v$ 
at $p \in B$ we get an 
integral curve $\gamma$.
Differentiating we obtain a vector field
on $U_i \cap U_j$:
$$
\theta_{ij}=\frac{d}{dt}\Big|_{t=0}
\psi_{ij}(x,\gamma(t)).
$$
\item Since the $\psi_{ij}$
satisfy the cocycle condition
(by the definition of the atlas
transition maps 
$\psi_{ij}=\psi_i^{-1}\circ \psi_j^{-1}$),
so do the $\theta_{ij}$, constituting
a cocycle in $H^1(T_X)$.
\end{enumerate}


\section{Kodaira-Spencer correspondence}
\label{deformations-section-kodaia-spencer-correspondence}


\medskip\noindent
After a few months of the previous text in
this section
I think I have understood the construction of
Kodaira-Spencer map for embedded
deformations,
which is \cite[Proposition
3.2.1(ii)]{Sernesi-deformations}.

But I will actually briefly outline a
particular case
of this, which is easier: \cite[Examples
3.2.4(ii)]{Sernesi-deformations}.
This the case when $X \subset Y$ is a Cartier
divisor.
Then $X$ is cut out locally by a set of
functions $f_i$
defined on an affine cover $U_i$ of $Y$.
Then we have a line bundle defined by the
transition functions $f_{ij}=f_i/f_j$.
This is the line bundle $\mathcal{O}_Y(X)$,
which
by the adjunction formula is also the normal
bundle $N_{X/Y}$
of $X$ in $Y$.

Now consider a first order deformation
$$
\xymatrix{
X\ar[r]\ar[d]
&\mathcal{X}\subset Y\times
B_\varepsilon\ar[d]^\pi\\
0\ar[r]
&B_\varepsilon
}
$$

The great thing about first order
deformations
is that the structure sheaf of $\mathcal{X}$
on an affine open $U_i$ of $Y$
looks like the tensor product
$\mathcal{O}_{U_i}\otimes
\mathbb{C}[\varepsilon]/(\varepsilon^2)$
(\href{https://stacks.math.columbia.edu/tag/%
01I4}{01I4}).
By definition of tensor product this just
says
that sections look like $f+\varepsilon g$
for $f,g \in \Gamma(\mathcal{O}_{U_i})$.

If $\mathcal{X}$ is a Cartier divisor on $Y
\times B_\varepsilon$
(it is; but I'm missing a formal proof;
the intuition is that an ``infinitesimal
thickening'' $\mathcal{X}$
of $X$ has the same underlying topological
space
but different scheme structure;
\href{https://stacks.math.columbia.edu/tag/%
01JL}{01JL}
\href{https://stacks.math.columbia.edu/tag/%
01I4}{01I4}),
it is given as the vanishing locus
of some local sections $f_i+\varepsilon g_i$.
(Perhaps the $f_i$ coincide with the
distinguished
section of $X$.)
The line bundle associated to $\mathcal{X}$
is given by the cocycle
$$
F_{ij}=f_{ij}+\varepsilon g_{ij}
$$
for $g_{ij}=\frac{g_i-f_{ij}g_j}{f_j}$.
Indeed, note that
$\frac{1}{f_j+\varepsilon
g_j}=\frac{1}{f_j}-\varepsilon
\frac{g_j}{f^2_j}$
and compute $F_{ij}=(f_i+\varepsilon
g_j)/(f_j+\varepsilon g_j)$.

Then
$$
f_i+\varepsilon g_i=(f_{ij}+\varepsilon
g_{ij})(f_j+\varepsilon g_j)
$$
gives
$$
g_i=f_{ij}g_j+g_{ij}f_j
$$
and since $f_j$ vanishes along $X$,
we are left with $g_i=f_{ij}g_j$
which is how a section of the sheaf $N_{X/Y}$
is constructed by definition of sheaf.


\begin{example}
\label{deformations-example-modul-space-nonsi-riema}
\begin{reference}
\cite[Example 3.1]{continued-fractions}
\end{reference}
Fix $g \geq 2$. The dimension of the
deformation space of a nonsingular
projective curve $X$ is $3g-3$. This is ``the
dimension of the moduli space of
curves of genus $g$''.

We can compute this number by Riemann-Roch
formula on the bundle $-K_X$.
Indeed, since $X$ is a curve,
$\Omega_X^1=K_X$ and thus $-K_X=T_X$. We get
$$
\begin{aligned}
h^0(-K)-h^0(K-(-K))&=\text{deg}(-K)-g+1\\
&=-2g+2-g+1\qquad \substack{\text{degree
additive and}\\ \text{deg}K=2g-2}\\
&=-3g+3
\end{aligned}
$$
Now by Serre duality
$h^0(K-(-K))=h^1(-K)=h^1(T_X)$, and it turns
out that a
Riemann surface of genus $g \geq 2$ has no
holomorphic vector fields, so that
$h^0(-K)=h^0(T_X)=0$.
\end{example}

\begin{exercise}
\label{deformations-exercise-deformations-of-curves-in-K3}
Let $C$ be a smooth genus $g$ curve which can
be embedded in a K3 surface  $M$,
and $X$ the family of all deformations of $C$
in $M$.
\begin{enumerate}
\item Prove that $\dim X \leq  g$.
\item Let $\mathcal{X}_g$ be the space of all
curves of genus $g$ (smooth?)
which can be possibly embedded to a K3
surface. Prove that each irreducible
component $Z$ of $\mathcal{X}_g$ satisfies
$\dim_\mathbb{C} \leq  g+19$. Deduce
that there exists a compact complex curve
which cannot be embedded in a K3
surface.
\end{enumerate}
\end{exercise}

\begin{exercise}
\label{deformations-exercise}
Let $C$ be a smooth curve embedded in a K3
surface $X$. Show that the dimension
of the deformation space of $C$ is $\leq g$.
\end{exercise}

\begin{proof}
% \begin{enumerate}
% \item
The deformation space of a variety is the
space of isomorphism classes of
deformations as explained above. It turns out
that there is a way to associate
1-cocyles of the tangent sheaf to
deformations, so that in fact the deformation
space $\text{Def}_1$ is isomorphic to
$H^{1}(X,\mathcal{T}_X)$ for any variety
$X$.

For our curve $C$ we thus know that the
dimension of the space of deformations
(deformations not necessarily contained in
$X$) is $h^1(\mathcal{T}_C)$.  The
family of deformations of $C$ that are
contained in $X$ is the Hilbert scheme of
curves with fixed Hilbert polynomial $P(t)$
after quotienting by $\mathbb{C}s$,
where $s$ is the section whose vanishing
locus is  $C$.
This says that the number we are looking for
is $h^{0}(X,\mathcal{O}(C))-1$.

Before showing that in fact
$h^0(X,\mathcal{O}(C))=1+g$,
let's explain why this Hilbert scheme
construction actually works.
Intuitively, any curve lying in the same
component
as $C$ would be deformation-equivalent to $C$
--- at least we know that it's Hilbert
polynomial is the same.
But it's not clear what's the deformation
involving these two
(this question should be addressed!).
Further, this doesn't explain Vladimiro's
suggesion of ``taking quotient by
$\mathbb{C}s$''.
To me it looks like I would just compute the
dimension
of the tangent space of $\text{Hilb}^P(M)$ at
$C$.

Now I will show that
$h^0(X,\mathcal{O}(C))=1+g$ (see
\cite[Lemma 1.2.1, Remark 1.2.2]{huk}).

Consider the ideal sheaf exact sequence
twisted by
$\mathcal{O}(C)$:
$$
\xymatrix{
0\ar[r]&\mathcal{O}_X\ar[r]&\mathcal{O}(C)
\ar[r]
&\mathcal{O}_X(C) \otimes
\mathcal{O}_C=\mathcal{O}(C)|_{C}\ar[r]&0
}
$$
By $X$ being a K3 we know that
$H^{1}(\mathcal{O}_X)=0$, so that we have the
short exact sequence in cohomology
$$
\xymatrix{
0\ar[r]&H^{0}(\mathcal{O}_X)\ar[r]&H^{0}
(\mathcal{O}(C))\ar[r]
&H^{0}(\mathcal{O}(C)|_{C})\ar[r]&0
}
$$
so that
$h^0(\mathcal{O}(C))=1+h^1(\mathcal{O}(C)
|_{C})$.
A version of adjunction formula says
$\omega_C \cong (K_X
\otimes\mathcal{O}(C)|_{C}$, and using that
$K_X=\mathcal{O}_X$ we obtain
$h^0(\mathcal{O}(C)|_{C})=h^0(\omega_C)=g$.

To

\medskip\noindent
For the record I put other thoughts I went
through in solving this exercise.

Recall from \cite[p. 146]{gri} that the
normal bundle
$\mathcal{N}$ of a hypersurface of a smooth
variety $X$ satisfies
$\mathcal{N}^\vee\cong\mathcal{O}_X(-C)|_{C}
$.
Taking duals we get that
$\mathcal{N}\cong\mathcal{O}_X(C)|_{C}$.

By adjunction formula
$2g-2=(\mathcal{O}(C),\mathcal{O}(C))$.
Applying Riemann-Roch
to $\mathcal{O}(C)$ (which is by definition
the dual of the ideal sheaf of $C$,
which is a line bundle on $X$), we obtain
that
$\chi(\mathcal{O}(C))=2+\frac{1}{2}
(\mathcal{O}(C),\mathcal{O}(C))$.
Then
$\chi(\mathcal{O}(C))=g+1$.

Now recall that
$\chi(\mathcal{O}(C))=h^0(\mathcal{O}(C)
-h^1(\mathcal{O}(C))+h^2(\mathcal{O}(C))$. By
Serre duality and $X$ being a K3
surface we see that $h^2(\mathcal{O}(C))\cong
h^{0}(\mathcal{O}(-C))$, which is
the ideal sheaf of $C$. Any section of such a
sheaf would vanish along $C$, and
since $X$ is compact we conclude there cannot
be any such section.

Now we show that also $h^1(\mathcal{O}(C))=0$
to conclude that
$h^0(\mathcal{O}(C))=g+1$.

We thus conclude that $h^0(\mathcal{O}(C))$




Since $C$ is smooth we can use the normal
exact sequence
$$
\xymatrix{
0\ar[r]&\mathcal{T}_C\ar[r]&\mathcal{T}
_X|_{C}\ar[r]&\mathcal{N}\ar[r]&0
}
$$
Taking Euler characteristic we see that
$\chi(\mathcal{T}_C)+\chi(\mathcal{N})=
\chi(\mathcal{T}_X|_{C})$.


This means that we should be done once we
compute $\chi(\mathcal{T}_X|_{C})$.
For this we can use Riemman-Roch formula for
coherent sheaves on a curve, which
tells us that
$$
\text{deg}(\mathcal{T}_X|_{C})=
\chi(\mathcal{T}_X|_{C})
-\text{rk}(\mathcal{T}_X|_{C})
\cdot\chi(\mathcal{O}_C)
$$
But of course we know that since $C$ is a
curve, by Serre duality we get that
$\chi(\mathcal{O}_C)=1-g$. (Indeed:
$h^1(\mathcal{O}_C)=h^0(\Omega^1_C):=p_a(C)
$.)
And now the question is what is
the degree. Apparently this is just the first
Chern class $c_1$ of the
restricted tangent bundle. So what is it? And
what is $h^0(\mathcal{T}_C)$ if
the genus is 0 or 1, and that's it.

I
know that $h^0(\mathcal{T}_X)=0$ and
$h^1(\mathcal{T}_X)=20$ by $X$ being a K3
surface, but I'm not sure what happens when
we restrict to $C$.

If $g \geq 2$ we know that
$h^0(\mathcal{T}_C)=0$, so that
$\boxed{\chi(\mathcal{T}_C)=h^1(\mathcal{T}
_C)}$.
\bigskip
To compute the
latter Euler characteristic, which is given
by definition by
$\chi(\mathcal{T}_X|_{C})=h^0(\mathcal{T}
_X|_{C})-h^1(\mathcal{T}_X|_{C})$,
we
first note that $h^0(\mathcal{T}_X|_{C})=0$,
because this is the dimension of
the global holomorphic vector fields on $X$
restricted to $C$, which is constant along
$C$ since $X$ is
smooth. And then this number is actually zero
by Hodge numbers of a K3 surface.

and taking cohomology long exact sequence we
obtain
$$
\xymatrix{
\cdots\ar[r]&H^{0}(T_X)\ar[r]&H^{0}
(\mathcal{N})\ar[r]&H^{1}(T_C)\ar[r]&\\
H^{1}(T_X)\ar[r]&H^{1}(\mathcal{N})\ar[r]&
\cdots
}
$$
\medskip\noindent
Or is it $h^1(\mathcal{N})$? See
\cite{HarrMorr}. If this was the case, then
we
can use adjunction formula as above to get
that
$2g-2=(\mathcal{N},\mathcal{N})=\text{deg}_C
\mathcal{N}$. Then we may find
$h^0(\mathcal{N})$ via Riemann-Roch:
$$
h^0(\mathcal{N})-h^0(K_C-\mathcal{N})=
\text{deg}\mathcal{N}-g+1
$$
Note that
$h^0(N_C-\mathcal{N})=h^1(-K_C+N+K_C)=
h^1(\mathcal{N})$
via Serre
duality, and by Riemann-Roch on a surface as
above we see that
$$
g+1=\chi(\mathcal{N})=h^0(\mathcal{N})-
h^1(\mathcal{N})
\implies
h^1(\mathcal{N})=-g-1+h^0(\mathcal{N})
$$
so that
$$
\begin{aligned}
h^0(\mathcal{N})-(-g-1+h^0(\mathcal{N}))&=
\text{deg}\mathcal{N}-g+1\\
\implies \text{oops! I lost
$h^0(\mathcal{N})$ in this operation…}
\end{aligned}
$$
Maybe if I just use normal exact sequence and
realise that
$$
h^0(\mathcal{N})=h^0(\mathcal{T}_X|_{C})-
h^0(\mathcal{T}_C)
$$
I know that $h^0(\mathcal{T}_C)=0$ for $g>1$,
so the question is how to compute
the restricted holomorphic vector fields.
% \end{enumerate}
\end{proof}

\section{Deformations of Stanley-Reisner schemes}
\label{deformations-section-sr}

\noindent
My aim in this section is to cover the
minimal
theory needed to apply results in \cite{jan2}
to my deformation problem. The main theorem
I expect to use is the following:

\begin{theorem}
\label{deformations-theorem-deformations-threefolds}
\begin{reference}
\cite[Theorem 5.7]{jan2}
\end{reference}
If $\mathcal{K}$ is a 3-dimensional manifold
then
$$
\dim T^1_{\mathbb{P}(K)}
=11d_3+5e_3+3e_4+e_{\geq 5}+c_{\geq
6}+5f_1^{(3)}+2f_1^{(4)}+h^2(\mathcal{K}).
$$
\end{theorem}

\noindent
I need to do some interpretation first.
From \cite[p. 1]{jan3}:

\begin{quotation}
  ``For each $\mathbb{C}$-algebra $A$,
there is a cohomology theory providing
modules $T^i_A$.
However, only $T^1_A$ and $T^2_A$ are
relevant for the deformation
theory of $\operatorname{Spec} A$. The module $T^1_A$
collects the infinitesimal
deformations of $A$ and $T^2_A$ contains the
obstructions
for lifting these deformations to decent
parameter spaces.
Eventually, $\text{Der}_{\mathbb{C}}(A,A)$
will be called $T_A^0$."
\end{quotation}

\cite[p. 3]{jan3}:
\begin{quotation}
``In general, for a finitely generated
$\mathbb{C}$-algebra $A$,
the modules $T_A^i$ allow the following ad
hoc definitions:
Let $P=\mathbb{C}[x_0,…,x_n]$ mapping onto
$A$ so that $A \simeq P/I$
for an ideal $I$. The $T_A^1$ is the cokernel
of the natural map
$\text{Der}_\mathbb{C}(P,P) \to \Hom_P(I,A)$.
Moreovr, […] and we obtain $T^2_A$ as the
cokernel
of the induced map
$\Hom_P(P^m,A)\to\Hom_A(R/R_0,A)$.''
\end{quotation}

\begin{remark}[dani]
\label{deformations-remark-dani1}
Let's think more about that definition of
$T^1_A$
as the cokernel of the ``natural map''
$\text{Der}_\mathbb{C}(P,P) \to \Hom_P(I,A)$.

Pick a derivation $\chi \in
\text{Der}_\mathbb{C}(P,P)$.
We want a map $I \to A=P/I$.
Pick an element $a \in I$ and map it to
$\chi(a)$.
And that's it. It might be zero or not.
If the $\chi$ we picked maps any $a \in I$ to
some element in $I$,
then this $\chi$ is in the kernel of the
natural map.
So the image consists of the deformations
coming from infinitesimal coordinate
changes. The module $T^1_A$ is the target
modulo this image, namely the cokernel.
\end{remark}

\noindent
Unfortunately this doesn't say much about the
following crucial fact:
the $T^i_A$ modules are
$\mathbb{Z}^{n+1}$-graded.
But here goes an explanation.
We use the notation that $e_p$ is a ``face
not in the complex'',
which of course corresponds to a generating
monomial $x^p$
of the SR ideal $I$.
It has a degree, of course (given by the
``support'' of $p$).
Now here's a rather sloppy step but:
since the cotangent modules are built upon
$I$,
who is given by $e_p$,
we just notice that the degree of $e_p$
along with a (very natural, too) notion of
degree among relations,
induce a grading in all the cotangent modules
$T_A^i$.
Thus, the grading is morally just the degree
of the monomials,
which is in fact the support of $p$, a vector
$\mathbf{c} \in \mathbb{Z}^{n+1}$.

\medskip\noindent
{\bf The $a-b$ grading.}
Next in line is the decomposition
$\mathbf{c}=\mathbf{a}-\mathbf{b}$
into positive minus negative part. But how
exactly…? 

Let $I_{\mathcal M}$ be a Stanley--Reisner
ideal. Its generators are monomials
corresponding to missing faces of
$\mathcal M$.

A first-order deformation replaces a
forbidden monomial by that monomial plus a
legal monomial. Schematically,
$$
x_b \quad \leadsto \quad x_b + t x_a .
$$

Here $x_b$ is the original Stanley--Reisner
generator. Thus $b$ records the missing face.

The monomial $x_a$ is the correction term.
Its support is a face of the complex, or at
least lies in a face/facet region of the
complex. Thus $a$ records what is inserted
by the deformation.

The multidegree of this deformation is
$$
\deg(x_a)-\deg(x_b)=a-b.
$$

For example, the direction $t_{48}$ sends
the missing triangle $348$ to the monomial
$x_2x_5^2$. Thus
$$
x_3x_4x_8
\quad\leadsto\quad
x_3x_4x_8+t x_2x_5^2.
$$
So in this case
$$
b=348,
\qquad
a=\deg(x_2x_5^2).
$$

In words: $b$ is the forbidden face being
deformed, and $a$ is the legal monomial
inserted as the first-order correction.




\section{Artin's criterion for flatness}
\label{deformations-section-artins-criterion}

The notes on minimal free resolutions,
syzygies, grading shifts and Macaulay2
notation are now in the commutative algebra
notebook, Section
\ref{commutative-algebra-%
section-minimal-resolutions}.
Here we use relations among the generators
of the defining ideal to apply Artin's
criterion for flatness
(\cite[Theorem 3.15]{Graffeo-Alessi}).

\noindent
The roadmap is
(I recall I couldn't get to a point
of absolute understanding when I last
dived into the details of
how to apply this theorem):
\begin{enumerate}
\item Start with a finitely generated
$R$-module $R/I$ with $I = (f_1,…,f_m)$.
Compute it's minimal resolution.

\item Choose an ideal $I_A$ inside
$R \otimes_\mathbb{C} A := R_A$
generated by elements $F_i$
such that $F_i-f_i \in \mathfrak{m}$
where $m$ is the maximal ideal of $A$.
This means that $F_i$ is a {\it lift} of
$f_i$.

\item Choose lifts $S_j$ of the syzygies
$s_i$ of $R/I$ in the same way.
Ask that the lifted syzygies $S_j$
are syzygies of the lifted complex
(this is to say:
$R^\ell_A \to R^m_A\to R_A \to R_A/I_A$
is exact.)

\item In this construction, when we tensor
by  $\mathbb{C}=A/\mathfrak{m}$
we are literally killing the new variables,
thus coming back to the fiber of
the closed point.

\item Then $A \to R_A/I_A$ is flat
with fiber over the closed point
$R \to R/I$.
\end{enumerate}

\begin{theorem}[Artin's criterion for flatness]
\label{deformations-theorem-artin-criterion}
Let $R=\mathbb{C}[x_1,…,x_n]$ be the polynomial ring
in $n$ variables and complex coefficients.
Let also $A$ be a local Artinian algebra of finite type over $\mathbb{C}$.
Let also $I \subset R$ and $I_A \subset R_A=R\otimes A$.
Suppose that there is an exact sequence
\begin{equation}
\label{deformations-equation-artin-criterion1}
\xymatrix{
R^\ell\ar[r]
&R^m\ar[r]
&R\ar[r]
&R/I\ar[r]
&0
}
\end{equation}
\noindent
and a complex
\begin{equation}
\label{deformations-equation-artin-criterion2}
\xymatrix{
R^\ell_A\ar[r]
&R^m_A\ar[r]
&R_A\ar[r]
&R_A/I_A\ar[r]
&0
}
\end{equation}

\noindent
such that the part
$$
\xymatrix{
R^m_A\ar[r]
&R_A\ar[r]
&R_A/I_A\ar[r]
&0
}
$$
is exact. Then, if
\ref{deformations-equation-artin-criterion2}$\otimes_A\mathbb{C}
=$\ref{deformations-equation-artin-criterion1},
the module $R_A/I_A$ is $A$-flat.
\end{theorem}

\section{Cotangent obstructions and
the Gr\"unbaum quadrics}
\label{deformations-section-cotangent-quadrics}

\noindent
The question is concrete: given a flat
first-order deformation of our equations,
can we add second-order terms and keep it
flat? We shall construct the space in which
the failure of this attempt lives.

\medskip\noindent
\textbf{Equations and their syzygies.}
Work over $k=\mathbb{C}$. Put
$$
S=k[x_1,\ldots,x_n],\qquad
I_0=(f_1,\ldots,f_m),\qquad A=S/I_0.
$$
The generators are homogeneous, with
$d_i=\deg f_i$. Write them as a row,
and introduce their generator module:
$$
F_0=(f_1,\ldots,f_m),\qquad
E=\bigoplus_{i=1}^m S(-d_i).
$$
If $\epsilon_i$ denotes the $i$th basis
vector of $E$, then $F_0(\epsilon_i)=f_i$.
The shift gives $\epsilon_i$ degree $d_i$,
so $F_0:E\to S$ preserves degree. Let
$$
\mathrm{Rel}=\ker(F_0:E\to S).
$$
Choose a graded free module $P$ and a matrix
$R_0:P\to E$ whose columns generate
\emph{all} of $\mathrm{Rel}$. Thus
$$
F_0R_0=0.
$$
Each column records an identity among the
generators, with polynomial coefficients.
The columns need not be independent or
minimal. In the Gr\"unbaum computation,
$m=16$ and a minimal $R_0$ has 30 columns.

Here is the precise flatness statement we
need, in the setting of
Theorem~\ref{deformations-theorem-artin-criterion}.
Let $B$ be a local Artinian $k$-algebra
with residue field $k$, and put
$S_B=S\otimes_k B$ and $E_B=E\otimes_k B$.
Choose lifts $F$ of the row $F_0$. Then
$S_B/(F)$ is flat over $B$ if and only if
every syzygy of $F_0$ lifts to a syzygy of
$F$. Equivalently, there is a matrix $R$
reducing to $R_0$ with $FR=0$ in $S_B$.

For justification, apply the
\href{https://stacks.math.columbia.edu/%
tag/051K}{Stacks Project, Lemma 10.101.6
(Tag 051K)}, with the maximal ideal of $B$.
It says that flatness is equivalent to
$$
\operatorname{Tor}^B_1(S_B/(F),k)=0.
$$
Indeed, write $J=(F)$. Tensoring
$0\to J\to S_B\to S_B/J\to0$ with $k$
identifies this Tor group with the kernel
of $J\otimes_B k\to S$. A presentation
$E_B\to J$ shows that this kernel vanishes
exactly when the relations among the $f_i$
come from relations among their lifts.
This also explains why generators of the
syzygy module suffice.

The coefficients of a syzygy are therefore
allowed to change. Requiring $FR_0=0$
would impose an extra condition: flatness
only requires some lift $R$ of $R_0$.

\medskip\noindent
\textbf{Expanding to second order.}
Put $B_r=k[t]/(t^{r+1})$, with $t$ of
degree zero in the polynomial grading.
Over $B_2$ write
$$
\begin{aligned}
F&=F_0+tF_1+t^2F_2,\\
R&=R_0+tR_1+t^2R_2.
\end{aligned}
$$
The rows $F_1,F_2:E\to S$ correct the
defining polynomials. The matrices
$R_1,R_2:P\to E$ correct the coefficients
of the syzygies. All maps preserve degree;
each $F_i$ has size $1\times m$, and each
$R_i$ has the same size as $R_0$.

Multiplication, followed by comparison of
the coefficients of $t$ and $t^2$, gives
\begin{equation}
\label{deformations-equation-syzygy-lifts}
\begin{aligned}
F_0R_1+F_1R_0&=0,\\
F_0R_2+F_1R_1+F_2R_0&=0.
\end{aligned}
\end{equation}
The constant coefficient already vanishes
because $F_0R_0=0$. These two identities
are exactly $FR=0$ modulo $t^3$.

Suppose the first identity has been solved.
Then $F_1,R_1$ give a flat deformation over
$B_1$. At the next step the known residual
is $F_1R_1$. We must find $F_2,R_2$ such
that
$$
F_1R_1=-F_0R_2-F_2R_0.
$$
Thus the obstruction asks whether this
residual can be removed by corrections.
It does not ask whether the residual
itself is zero.

For comparison with first-order theory,
put a bar over a polynomial to denote its
class in $A$. The first identity implies
$\overline{F_1}R_0=0$. Consequently the rule
$$
f_i\longmapsto\overline{(F_1)_i}
$$
defines an element of
$\operatorname{Hom}_S(I_0,A)_0$:
it respects every relation among the
$f_i$. This is the embedded tangent space
from
Proposition~%
\ref{deformations-proposition-affine-first-order}.

\medskip\noindent
\textbf{The quotient defining cotangent
cohomology.}
Some syzygies exist for any list of
polynomials, simply by commutativity. Let
$K\subset\mathrm{Rel}$ be the submodule
generated by the Koszul syzygies
$$
k_{ij}=f_i\epsilon_j-f_j\epsilon_i.
$$
The quotient $\mathrm{Rel}/K$ is an
$A$-module. To check this, take
$s=\sum_j s_j\epsilon_j\in\mathrm{Rel}$.
Then
$$
f_i s=\sum_j s_j
(f_i\epsilon_j-f_j\epsilon_i)\in K,
$$
since $\sum_j s_jf_j=0$. Thus $I_0$
annihilates $\mathrm{Rel}/K$.

Every $S$-linear map $\phi:E\to A$
vanishes on $K$, because the $f_i$ act as
zero on $A$. Restriction therefore gives
$$
\begin{aligned}
\rho:\operatorname{Hom}_S(E,A)_0
&\longrightarrow
\operatorname{Hom}_A(\mathrm{Rel}/K,A)_0,\\
\rho(\phi)(s+K)&=\phi(s).
\end{aligned}
$$
If $\phi$ is represented by a row $H$ of
polynomials, its values on the chosen
syzygy generators are $\overline{H R_0}$.

\begin{definition}
\label{deformations-definition-cotangent-two}
The degree-zero part of the
\emph{second cotangent cohomology} of
$A$ over $k$, with coefficients in $A$, is
$$
T^2_{A,0}=\operatorname{coker}(\rho).
$$
Explicitly,
\begin{equation}
\label{deformations-equation-cotangent-two-quotient}
T^2_{A,0}=
\frac{
\operatorname{Hom}_A(\mathrm{Rel}/K,A)_0
}{
\rho\bigl(\operatorname{Hom}_S(E,A)_0\bigr)
}.
\end{equation}
\end{definition}

\noindent
The subscript $0$ means that maps preserve
the usual total degree in the $x_i$.
It does not mean constant polynomial
entries. In our project $E=S(-3)^{16}$,
so a degree-zero map $E\to A$ sends each
basis vector to a cubic in $A$.
The superscript $2$ is the cohomological
index; it is unrelated to this grading.

This is the concrete construction in
\cite[Section 0.2]{Fausk}. Her generator
module is called $F$, her polynomial ring
is $R$, and her Koszul submodule is
$\mathrm{Rel}_0$. Her source for the
restriction map agrees with ours because
$K\subset I_0E$.
The construction and its independence
from the presentation are also explained
in \href{https://math.berkeley.edu/%
\string~robin/math274root.pdf}
{Hartshorne's online lectures}, Section 3,
especially Lemmas 3.2 and 3.3.
These references identify the quotient
with cotangent cohomology; we do not need
to construct a cotangent complex here.

The following theorem explains why this
quotient detects flat lifts at every order.
Write $T^2_A=T^2(A/k,A)$ for the same
cokernel without the degree restriction.

\begin{theorem}[Obstruction to a flat lift]
\label{deformations-theorem-obstruction-flat-lift}
Let $A$ be a finitely generated
$k$-algebra. Let $B'\to B$ be a surjection
of local Artinian $k$-algebras with
residue field $k$. Write $J=\ker(B'\to B)$
and assume
$$
\mathfrak{m}_{B'}J=0.
$$
Let $A_B$ be a flat, finitely generated
$B$-algebra, with a specified
identification $A_B\otimes_B k\cong A$.
There is a canonical obstruction
$$
\operatorname{ob}(A_B,B')
\in T^2_A\otimes_k J.
$$
It vanishes if and only if there exists
a flat, finitely generated $B'$-algebra
$A_{B'}$ with an identification
$$
A_{B'}\otimes_{B'}B\cong A_B.
$$
Such a lift can be realized by lifting
the generators of the ideal in a chosen
polynomial presentation of $A_B$.
\end{theorem}

\noindent
For a proof, see Hartshorne,
\href{https://math.berkeley.edu/%
\string~robin/math274root.pdf}
{\emph{Lectures on Deformation Theory}}
(2004), Theorem 10.1(a), p.~58
(PDF page 64; proof on pp.~59--60),
and Corollary 10.3, p.~61.

\medskip\noindent
For the rings $B_r$ above, take
$$
B'=B_{r+1},\qquad B=B_r,\qquad
J=(t^{r+1})/(t^{r+2})\cong k.
$$
Thus a lift extends the specified family
by one order while preserving flatness.
For graded deformations with $\deg t=0$,
the calculation takes place in
$T^2_{A,0}$: the maps preserve degree,
so corrections can be chosen in degree
zero. We now construct the obstruction
explicitly at second order.

\medskip\noindent
\textbf{Why the residual belongs to this
particular target.}
For an arbitrary $s\in\mathrm{Rel}$,
choose a first-order lift $s+tu$ satisfying
$$
F_0u=-F_1s.
$$
Such a $u\in E$ exists because the
first-order deformation is flat. Define
$$
\beta(s)=\overline{F_1u}\in A.
$$
There are three small checks.
\begin{enumerate}
\item If $u'$ is another choice, then
$u'-u\in\mathrm{Rel}$. Since $F_1$ sends
every syzygy into $I_0$, we have
$\overline{F_1(u'-u)}=0$. Thus $\beta(s)$
does not depend on the choice of lift.
Adding or multiplying choices of $u$
also proves that $\beta$ is $S$-linear.

\item For $s=k_{ij}$, take
$$
u=(F_1)_i\epsilon_j-(F_1)_j\epsilon_i.
$$
Then $F_0u=-F_1k_{ij}$, while
$$
F_1u=(F_1)_i(F_1)_j
-(F_1)_j(F_1)_i=0.
$$
Hence $\beta(K)=0$; this reflects the fact
that Koszul syzygies always lift.

\item The degrees in the construction
match. Therefore $\beta$ induces a
degree-zero $A$-linear map
$$
\beta:\mathrm{Rel}/K\longrightarrow A.
$$
\end{enumerate}
On the columns of $R_0$, this map is
represented by $\overline{F_1R_1}$.
The construction proves that these values
respect relations among the syzygies.
An arbitrary row with the same number of
entries need not have this property.

\medskip\noindent
\textbf{Which part can the corrections
remove?}
Reducing the second identity in
\ref{deformations-equation-syzygy-lifts} modulo $I_0$
removes $F_0R_2$, and leaves
$$
\beta+\rho(\overline{F_2})=0.
$$
Consequently the obstruction class is
$$
\operatorname{ob}_2(F_1)=[\beta]
\in T^2_{A,0}.
$$
If this class is zero, choose $F_2$ with
$\rho(\overline{F_2})=-\beta$. Each entry
of $F_1R_1+F_2R_0$ then belongs to $I_0$.
Expressing its negative in the generators
$f_i$ supplies the columns of $R_2$.
Thus a second-order flat lift exists.
Conversely, any such lift gives this
identity and hence a zero class.

Notice that $F_2$ is an arbitrary correction
row of the appropriate degrees. It need
not itself represent a first-order
deformation. This is why the denominator
uses all of $\operatorname{Hom}_S(E,A)_0$.

We have already checked independence of
$R_1$ and of the chosen syzygy generators.
For completeness, replacing $F_1$ by
$F_1+F_0H$ merely changes its polynomial
representatives. Taking $R_1-HR_0$ as the
new syzygy correction changes $\beta$ by
$$
-\rho(\overline{F_1H}),
$$
so its class is unchanged. Here $H:E\to E$
has degree zero. Changes of the generator
presentation give the canonical
identifications cited above. Infinitesimal
coordinate changes also preserve the class:
they extend to coordinate changes over
$k[t]/(t^3)$, and the obstruction
construction is natural under these
isomorphisms. Thus the class depends on the
intrinsic first-order deformation.

\begin{example}
\label{deformations-example-removing-residual}
Take $S=k[x,y,z]$ and $I_0=(x^2,xy)$.
All syzygies are multiples of the column
$r=(y,-x)^{\mathsf{T}}$: cancellation of
$x$ in $ax^2+bxy=0$ gives $ax+by=0$.
Choose
$$
F_0=(x^2,xy),\quad F_1=(0,yz),\quad
R_0=\begin{pmatrix}y\\-x\end{pmatrix},
\quad R_1=\begin{pmatrix}0\\z\end{pmatrix}.
$$
Then $F_0R_1=xyz$ and $F_1R_0=-xyz$.
The first-order equation holds, but
$$
F_1R_1=yz^2\neq0
$$
even in $A=S/(x^2,xy)$. Nevertheless,
$F_2=(-z^2,0)$ gives $F_2R_0=-yz^2$,
so $R_2=0$ completes the second-order lift.
Indeed the product is identically zero:
$$
(x^2-t^2z^2,xy+tyz)
\begin{pmatrix}y\\-x+tz\end{pmatrix}=0.
$$

Here $\mathrm{Rel}=Sr$ and $K=xSr$.
Giving a map $\mathrm{Rel}/K\to A$ amounts
to choosing its value on $r$, annihilated
by $x$. Since
$$
\operatorname{ann}_A(x)=(x,y),\qquad
\rho(h_1,h_2)(r)=yh_1-xh_2,
$$
restriction is surjective, also in degree
zero: every cubic in $(x,y)\subset A$ has
this form with quadratic $h_1,h_2$.
Thus $T^2_{A,0}=0$. The nonzero residual
represents zero in the quotient.
\end{example}

\medskip\noindent
\textbf{Where the quadratic polynomials
come from.}
The intrinsic tangent space is
$$
V=T^1_{A,0}=
\frac{\operatorname{Hom}_S(I_0,A)_0}
{\operatorname{im}
(\operatorname{Der}_k(S,A)_0
\to\operatorname{Hom}_S(I_0,A)_0)}.
$$
The denominator consists of infinitesimal
coordinate changes. Choose a basis
$v_1,\ldots,v_d$ of $V$ and polynomial
representatives for it. For each basis
vector solve the first-order syzygy
equation once. Linearity then gives
$$
F_1(y)=\sum_{a=1}^d y_a U_a,\qquad
R_1(y)=\sum_{a=1}^d y_a V_a.
$$
Here $U_a$ is a row representing $v_a$,
and $V_a$ is its syzygy correction matrix.
The $y_a$ are coefficients of a tangent
vector, not the original variables $x_i$.

The residual is now
$$
F_1(y)R_1(y)=
\sum_{a,b}y_ay_bU_aV_b.
$$
It is quadratic in $y$. Reduction modulo
$I_0$ and passage to the cokernel are
linear operations. Therefore they give
a quadratic map
$$
q_2:V\longrightarrow T^2_{A,0}.
$$
Choose a basis $e_1,\ldots,e_r$ of
$T^2_{A,0}$. In this basis,
\begin{equation}
\label{deformations-equation-quadratic-coordinates}
q_2(y)=\sum_{i=1}^r Q_i(y)e_i,
\qquad Q_i=e_i^*\circ q_2.
\end{equation}
The dual vector $e_i^*$ extracts the
$i$th coordinate. This is literally the
meaning of the polynomial $Q_i$: it is
one coordinate of the obstruction class.

Computationally, choose monomial bases in
the relevant graded pieces of $A$. Encode
maps on syzygies by their coefficients,
impose the relations among syzygies and
the condition that $K$ maps to zero, and
then row-reduce the restriction map $\rho$.
Apply the resulting quotient coordinates
to $\overline{F_1(y)R_1(y)}$. This is
linear elimination producing the $Q_i$.
Equivalently one eliminates the unknown
coefficients of $F_2,R_2$ from the
second-order equation. Identifying the
result with $T^2$ requires the target and
quotient just constructed; the cokernel
of an arbitrary matrix of equations is
not automatically cotangent cohomology.

Changing the basis $e_i$ takes invertible
linear combinations of the $Q_i$.
Changing tangent coordinates substitutes
an invertible linear change in $y$.
The obstruction map itself is intrinsic.
There can be fewer than $r$ independent
quadratic coordinates: the assertion that
all $r$ are independent needs checking.

\medskip\noindent
\textbf{The numbers in the Gr\"unbaum
project.}
The computational source here is
\href{https://github.com/danimalabares/%
grunbaum-zero-context-proof/%
tree/%
ee984e9541e9cb9259728e6dc7b3a1bc0fda8a9b}
{the zero-context proof repository},
at the displayed pinned revision.
Its variables $a,\ldots,h$ use a different
vertex labelling from our
Stanley--Reisner chapter. Direct comparison
of the 16 monomials identifies the ideals
under the substitution
$$
(a,b,c,d,e,f,g,h)\longmapsto
(x_3,x_1,x_7,x_4,x_2,x_6,x_5,x_8).
$$

The audit
\href{https://github.com/danimalabares/%
grunbaum-zero-context-proof/blob/%
ee984e9541e9cb9259728e6dc7b3a1bc0fda8a9b/%
enumerate/SYMMETRY_AUDIT.md}
{\texttt{enumerate/SYMMETRY\_AUDIT.md}}
points to the following exact computations.
Only their tangent and obstruction data
are needed here.
\begin{enumerate}
\item The script
\href{https://github.com/danimalabares/%
grunbaum-zero-context-proof/blob/%
ee984e9541e9cb9259728e6dc7b3a1bc0fda8a9b/%
reconstruct/coordinate_split_rank.m2}
{\texttt{coordinate\_split\_rank.m2}}
computes the normal space and the image
of the 64 linear coordinate changes.
The dimensions are
$$
\dim_k\operatorname{Hom}_S(I_0,A)_0=109,
\qquad \dim_k V=109-56=53.
$$
There are 104 standard cubic monomials in
$A$, so the initial coefficient space has
$16\cdot104=1664$ dimensions. Imposing
the syzygies leaves 109; quotienting by
coordinate changes removes 56 more.
Thus 53 counts intrinsic tangent
directions, not all embedded directions.

\item The script
\href{https://github.com/danimalabares/%
grunbaum-zero-context-proof/blob/%
ee984e9541e9cb9259728e6dc7b3a1bc0fda8a9b/%
reconstruct/audit_ct1_ct2_basis.m2}
{\texttt{audit\_ct1\_ct2\_basis.m2}}
and its
\href{https://github.com/danimalabares/%
grunbaum-zero-context-proof/blob/%
ee984e9541e9cb9259728e6dc7b3a1bc0fda8a9b/%
reconstruct/audit_ct1_ct2_basis_output.txt}
{recorded output} report 53 basis columns
for \texttt{CT\^{}1(\{0\},F0)} and 27
for \texttt{CT\^{}2(\{0\},F0)}.
The latter matrix has size $30\times27$:
its rows refer to the syzygy generators,
and its columns represent a basis of
$$
W=T^2_{A,0},\qquad \dim_k W=27.
$$
The interpretation of \texttt{CT\^{}2}
as cotangent cohomology is specified in
the \href{https://macaulay2.com/doc/%
Macaulay2/share/doc/Macaulay2/%
VersalDeformations/html/___C__T.html}
{Macaulay2 package documentation}.

\item The same Macaulay2 script computes
a fresh order-two versal deformation.
It checks that its 27 quadratic
coordinates have rank 27 and span the
same space as the stored quadratic list.
Separately,
\href{https://github.com/danimalabares/%
grunbaum-zero-context-proof/blob/%
ee984e9541e9cb9259728e6dc7b3a1bc0fda8a9b/%
enumerate/audit_symmetry_formal.py}
{\texttt{audit\_symmetry\_formal.py}}
extracts the degree-two parts of
\texttt{reconstruct/}%
\texttt{p1\_quartic\_component.m2}
and checks their rank over $\mathbb{Q}$.
For example, in that file's coordinates
(with $t_i$ renamed $y_i$), the first two
quadrics are
$$
\begin{aligned}
Q_1&=-y_{40}y_{46}+y_{43}y_{49},\\
Q_2&=y_{11}y_{43}-y_1y_{46}.
\end{aligned}
$$
These are coefficients in the obstruction
space, not additional equations in the
original variables $x_1,\ldots,x_8$.
\end{enumerate}

\noindent
\emph{Verification status (8 September
2026).} The monomial relabelling, the
dimensions 109 and 56, and the rank 27 of
the stored quadrics were recomputed with
exact rational arithmetic while preparing
this section. For the tangent calculation,
all pair syzygies of the monomial ideal
were imposed. For the quadratic rank,
each of the 27 polynomials has a monomial
absent from the other 26; these give a
diagonal coefficient minor with nonzero
entries. Macaulay2 was unavailable, so the
independent computation of
$\dim T^2_{A,0}=27$ and the identification
of the stored span with the freshly
computed primary obstruction coordinates
were checked against the cited scripts
and their recorded output, not rerun.
Those two statements remain computational
inputs to this explanation.

The calculations use rational coefficients;
extension from $\mathbb{Q}$ to
$\mathbb{C}$ preserves these dimensions
and ranks. The 27 independent quadrics
give an injective map
$$
W^*\longrightarrow
\operatorname{Sym}^2(V^*),\qquad
\lambda\longmapsto\lambda\circ q_2.
$$
Thus their span is identified with $W^*$.
Independence is an extra computational
fact, not a consequence of $\dim W=27$.
It also does not say that the quadrics
form a regular sequence or impose
codimension 27.

\medskip\noindent
\textbf{What vanishing proves.}
For a specified tangent vector $y$,
$$
Q_1(y)=\cdots=Q_{27}(y)=0
$$
is equivalent to the existence of a flat
lift from $k[t]/(t^2)$ to $k[t]/(t^3)$.
It does not by itself give an all-orders
deformation or a smooth fibre.

At the next stage, the known residual is
$F_1R_2+F_2R_1$, and the new corrections
are $F_0R_3+F_3R_0$. More generally, after
a flat lift through order $n$, the
coefficient to be removed at order $n+1$
is
$$
\sum_{i=1}^n F_iR_{n+1-i}.
$$
Theorem~\ref{deformations-theorem-obstruction-flat-lift}
assigns
its class to the same space $T^2_{A,0}$,
tensored with the one-dimensional kernel
$(t^{n+1})/(t^{n+2})$. These later classes
depend on the previously chosen lift.
The original quadrics describe only the
first of these lifting tests.

\section{Gröbner denerations}
\label{deformations-section-grobner-degenerations}

\noindent
Here's the roadmap:
\begin{enumerate}
\item Start with a smooth variety
with algebra $R/I$ that we suspect
might degenerate to the singular
variety.

\item Choose a weight vector
$\omega \in \mathbb{R}^{n+1}$
where $R = \mathbb{C}[x_0,…,x_n]$.
Define
\begin{enumerate}
\item Initial form
$\text{in}_\omega(f) = $
sum of monomials in $f= \sum c_a x^a$
such that $a\cdot \omega$
is maximal.

\item Initial ideal generated
by initial forms $f \in I$.

\item Homoegenization form
$\tilde{f} = \sum c_a x^at^{b-\omega\cdot a}$
where $b = \text{max}\{\omega \cdot a\}$.

\item Homogenization ideal $\tilde{I}$
generated by $\tilde{f}$ for $f \in I$.
\end{enumerate}
\item Apply theorem: then
$R[t]/\tilde{I}$
is flat over $\mathbb{C}[t]$
with fiber $R/\text{in}_\omega(I)$
over $(0)$ and fiber $R/I$ over 
any $(t-u)$ for $u\neq 0$.
\end{enumerate}

\begin{theorem}
\label{deformations-theorem-grobner-flat-morphism}
For any ideal $I \subset k[x_1,\ldots,x_r]$,
the $k[t]$-algebra $S[t]/\tilde{I}$ is free
---and thus flat--- as a $k[t]$ module.

[…].

Thus $S[t]/\tilde{I}$ is a flat family over
$k[t]$ of quotients of $S$ whose fiber over
$0$ is $S/\text{in}_\lambda(I)$ and whose
fiber over any $(t-u)$ for $0 \neq u \in k$
is $S/I$.

\end{theorem}

\medskip\noindent
Now I'll put the original notes on this:


\section{Miscellaneous}
\label{deformations-section-misc}

Look for the ``Perfect obstruction theory''.
This is used by Bruzzo for deformations of
(super)stacks,
work with E. Paiva and ---.


% --- End deformations.tex ---

% --- Begin derived-categories-of-sheaves.tex ---
\chapter{Derived categories of sheaves}
\label{chapter-derived-categories-of-sheaves}



\noindent
Minicourse by Lyalya Guseva (Universit\'e
Bourgogne Franche-Comt\'e, France; HSE
University, Moscow, Russia), CIMPA school
Florian\'opolis 2025.

\medskip\noindent
Notes at
\href{%
http://github.com/danimalabares/cimpa-floripa
}{%
github.com/danimalabares/cimpa-floripa
}

\bigskip\noindent

{\bf Abstract.} Derived category of coherent
sheaves is a convenient environment of
investigating algebraic geometry of a
variety. It provides useful techniques and
gives a perspective point of view. On the
level of derived categories one can see
unexpected connections which are not visible
on the classical level. I will try to give an
introduction into the techniques of derived
categories and semiorthogonal decomposition
and (hopefully) will present many examples of
applications of derived categories in
algebraic geometry.


{\bf Upshot.} The derived category of an
abelian category is the localization of its
homotopy category with respect to the class
of quasi-isomorphisms.
\begin{itemize}
\item The objects of $D(\mathcal{A})$ are
complexes in $\mathcal{A}$ and the
arrows are equivalence classes of roofs.
\item $D(\mathcal{A})$ is additive but not
abelian unless $\mathcal{A}$ is
semisimple.
\item $D(\text{Vect}_k)=$ graded vector
spaces.
\item For all $A,B\in \mathcal{A}$,
$$
\text{Hom}_{D(\mathcal{A})}(A,B[i])=
\begin{cases}
0\qquad &i<0 \\
\text{Hom}_{\mathcal{A}}(A,B)\qquad &i=0\\
\text{Ext}^i(A,B)\qquad &i>0
\end{cases}
$$
\item A sequence
$$
\xymatrix{
A\ar[r]^f&B\ar[r]^b&\text{Cone}(f)\ar[r]^a&
A[I]\ar[r]&0
}
$$
gives long exact sequences
$\text{Hom}^\bullet_{D(\mathcal{A})}(C,\cdot)
$
and
$\text{Hom}^\bullet_{D(\mathcal{A})}(\cdot,C)
$,
which are the distinguished triangles.
\end{itemize}

\section{Basic concepts}

\begin{definition}
\label{derived-categories-of-sheaves-definition-abelian-category}
An {\it abelian category} $\mathcal{A}$ is a
category such that
\begin{enumerate}
\item $\forall X,Y \in \mathcal{A}$,
$\text{Hom}_{\mathcal{A}}(X,Y)$ is an
abelian group such that
$\text{Hom}_{\mathcal{A}}(X,Y)
\times\text{Hom}
(Y,Z)\to
\text{Hom}(X,Y)$ is bilinear.
\item There is a zero object
$0_{\mathcal{A}}$ such that for every $A \in
\mathcal{A}$,
$\text{Hom}(0_{\mathcal{A}},A)=\text{Hom}(A,
0)
=0$.
\item There exist finite products.
\item For every  $f \in \text{Hom}(A,B)$
there exist kernel and cokernel such
that
$$
\xymatrix{
\Ker f\ar[r]&A\ar[r]^{f}&B\ar[r]&\Coker f
}
$$
\end{enumerate}
\end{definition}

\begin{definition}
\label{derived-categories-of-sheaves-definition-additive-category}
$\mathcal{A}$ satisfying 1-3 is called  {\it
additive category}.
\end{definition}

Examples: abelian groups, $R$-modules,
$\text{Coh}(X)$. Non-example: vector bundles
over $X$, it is an additive category but not
abelian.

\medskip\noindent

Now consider the category
$\text{Com}(\mathcal{A})$, the category of
complexes in $\mathcal{A}$, with
quasi-isomorphisms:

\begin{definition}
\label{derived-categories-of-sheaves-definition-quasi-isomorphism}
$f\in
\text{Hom}_{\text{Com}(\mathcal{A})}(A',B')$
is a {\it quasi-isomorphism} if it induces
isomorphisms $H^i(f):H^i(A^{\bullet})\cong
H^i(B^{\bullet})$ for all $i$.
\end{definition}

\begin{definition}
\label{derived-categories-of-sheaves-definition-derived-category}
Let $\mathcal{A}$ be an abelian category. The
{\it derived category} of $\mathcal{A}$ is
the localization
$\text{Com}(\mathcal{A})[QIS]^{-1}$, where
$QIS$ is the class of quasi-isomorphisms.
\end{definition}

But what is localization? Given a class of
morphisms $S$ inside the class of morphisms
of a given category, the localization will
allow us to invert these morphisms.

\begin{definition}
\label{derived-categories-of-sheaves-definition-localization-of-category}
Let $\mathcal{C}$ be a category, $S$ a class
of morphisms in $\mathcal{C}$ and
$\text{Iso}(\mathcal{C})$ the class of all
isomorphisms of $\mathcal{C}$. The {\it
localization} of $\mathcal{C}$ with respect
to $S$ is the category $\mathcal{C}[S]^{-1}$
with an isomorphism
$$
Q:\mathcal{C}\to
\mathcal{C}[S]^{i-\mathcal{A}}
$$
that satisfy
\begin{enumerate}
\item $Q(S) \subset
\text{Iso}(\mathcal{C}[S]^{-1})$
\item (Universality.) For any
$F:\mathcal{C}\to \tilde{\mathcal{C}}$ and
$Q:\mathcal{C}\to \mathcal{C}[S^{-1}]$ there
exists a unique $G:\mathcal{C}[S^{-1}]\to
\tilde{ \mathcal{C}}$ making the diagram
commute, i.e. $F=GQ$.
\end{enumerate}
\end{definition}

\begin{exercise}
\label{derived-categories-of-sheaves-exercise-uniqueness-of-localization}
Show that if $\mathcal{C}[S^{-1}]$ exists, it
must be unique up to isomorphism.
\end{exercise}

\begin{definition}
\label{derived-categories-of-sheaves-definition-bounded-derived-category}
If instead of $\text{Com}(\mathcal{A})$ we
consider
$$
\text{Com}^b(\mathcal{A}):=\{A^{\bullet} \in
\text{Com}(\mathcal{A}):
A^i=0\forall  |i|\gg 0\}
$$
Then we obtain the {\it bounded derived
category} $D^b(\mathcal{A})$.
\end{definition}

\section{Motivation}
\label{derived-categories-of-sheaves-section-motivation}

Let $X$ be a smooth projective variety over
$\mathbb{C}$. Then we can associate
$X\rightsquigarrow \text{Coh}(X)$.

\begin{theorem}[Gabriel]
\label{derived-categories-of-sheaves-theorem-Gabriel}
Let $X$ and $Y$ be smooth projective
varieties. If $\text{Coh}(X)\cong
\text{Coh}(Y)$ then $X\cong Y$.
\end{theorem}

This means we may reconstruct $X$ from its
coherent sheaf category. But it is not such
an interesting invariant. It's more
interesting the derived category. So the
motivation for studying derived categories
are:

\begin{enumerate}
\item (Beilinson 1979.)
$\text{Coh}(\mathbb{P}^n)\hookrightarrow
D^b(\mathbb{P}^n)$ gives good information
about coherent sheaves on $\mathbb{P}^n$.
\item (Mukai.) If $\mathcal{A}$ is an abelian
variety and $\hat{\mathcal{A}}$ is
the dual abelian variety,
$D^b(\mathcal{A})\cong
D^b(\hat{\mathcal{A}})$.
\item (Bondal, Orlov, 1990's.) Let $X$ be a
smooth projective variety with an
ample or anti-ample canonical bundle
$\omega_X$. Suppose that $D^b(X)\cong D^b(Y)$
for a smooth projective variety $Y$. Then $Y
\cong X$.
\end{enumerate}

\begin{remark}
\label{derived-categories-of-sheaves-dcos-label-001}
If you consider not only the bounded derived
category $D^b(X)$ but consider it with the
tensor product $D^b \times D^b(X)
\xrightarrow{\otimes} D^b(X)$ allow to
reconstruct $X$.
\end{remark}

I hope you are now convinced this is
interesting.

\section{Explicit description of localization
under
some assumption}
\label{derived-categories-of-sheaves-dcos-label-002}

If $S$ satisfies some properties called the
{\it right Ore conditions} then
$\mathcal{C}[S^{-1}]$ can be explicitly
described. The conditions are:

\begin{enumerate}
\item All identities should belong to $S$ and
if any two morphisms from $f,g$
and $f \circ g$ (whenever composition is
defined) belong to $S$, then all three belong
to $S$.
\item For any diagram
$$
\xymatrix{
X\ar[d]_{f}&W\ar@{.>}[l]_{\tilde{s}\in
S}\ar@{.>}[d]\\
Z&\ar[l]^{s\in S}Y
}
$$
can be completed to a commutative square.
\item If for $f,g\in
\text{Hom}_{\mathcal{C}}(A,B)$ and $u:B \to
U$ such athat
$u\circ f=u\circ g$ there exists $t:C \to A$
such that $ft=gt$, that is,
$$
\xymatrix{
C\ar[r]^{t}&A\ar[r]^{f,g}&B\ar[r]^{u}&U
}
$$
\end{enumerate}

The quasi-isomorphisms don't satisfy
condition 2.

\begin{proposition}
\label{derived-categories-of-sheaves-dcos-label-003}
Let $\mathcal{C}$ be a category and $S\subset
\Mor(\mathcal{C})$ satisfy right Ore
conditions. Then $\mathcal{C}[S^{-1}]$ can be
described as follows:
\begin{itemize}
\item
$\text{Ob}(\mathcal{C}[S^{-1}]=\text{Ob}
(\mathcal{C})$.
\item
$\Mor_{\mathcal{C}[S^{-1}]}(A,B)=
\{\text{equivalence
classes of
roofs}\}$, where a roof is
$$
\xymatrix{
& C\ar[dl]_{S \ni s}\ar[dr]^{f}\\
A&&B
}
$$
where $fs^{-1}\sim gt^{-1}$ if there exists a
commutative diagram
$$
\xymatrix{
& & \tilde{W} \ar[ld]^{a \in S} \ar[dr]^{b
\in S} & & \\
& C \ar[ld]^{s \in S} \ar[rr]^{f} & &
\tilde{C} \ar[rd]^{t} & \\
A & & & & B
}
$$
and composition is defined as follows:
$$
\xymatrix{
& & \tilde{W}\ar[ld]^{a \in S}\ar[dr]^{b \in
S}\\
&Z\ar[ld]^{s \in S}\ar[dr]^{f} & &
\tilde{Z}\ar[dr]\ar[dl]^{t}\\
A & &B & & C
}
$$
where
$$
(fs^{-1})\circ (gt^{-1})=(sa^{-1})\circ(gb)
$$
\end{itemize}
\end{proposition}

$\mathcal{C}[S^{-1}]$ always exists but
morphisms are difficult to describe. The
proof of the proposition is an exercise:

\begin{proof}
\begin{enumerate}
\item This is an equivalence relation.
\item Composition doesn't depend on $a$ and
$b$.
\item Composition doesn't depend on the
choice of the roof in equivalence
classes.
\item Composition is associative.
\item $\text{Id}_X \times \text{Id}_X^{-1}$
is identity morphism.
\end{enumerate}
\end{proof}

{\bf Definition of} $Q:\mathcal{C} \to
\mathcal{C}[S^{-1}]$. (Missing.)

{\bf Universal property.}

$$
\xymatrix{
\mathcal{C}\ar[r]^{F}\ar[dr]^{Q}&
\tilde{\mathcal{C}}\\
&\mathcal{C}[S^{-1}]\ar@{.>}[u]
}
$$
defines a universal property.

\begin{proposition}
\label{derived-categories-of-sheaves-proposition-Ore-conditions}
Let $\mathcal{C}$ be an additive category and
$S$ satisfy right Ore conditions. Then
$\mathcal{C}[S^{-1}]$ is also additive and
$Q:\mathcal{C} \to
\mathcal{C}[S^{-1}]$ is additive.
\end{proposition}

\begin{proof}
Some indications, then exercise.
\end{proof}

The point is that with Ore conditions you can
really define what is addition of fractions.

\section{Ore conditions for category of
complexes}
\label{derived-categories-of-sheaves-dcos-label-004}

Let $\mathcal{C}=\text{Com}(\mathcal{A})$ and
$S=\text{QIS}$. Then QIS {\bf does not}
satisfy Ore conditions: second and third
properties fail. But we can consider instead
the homotopy category
$\mathcal{H}(\mathcal{A})$. But what is that.

\begin{definition}
\label{derived-categories-of-sheaves-definition-homotopy-category}
\begin{enumerate}
\item
A morphism $f \in
\text{Hom}_{\text{Com}\mathcal{A}}
(A^{\bullet}
,B^{\bullet})$
is {\it homotopically equivalent to zero},
i.e. $f \sim_{\text{hom}}0$ if for all $i\in
\mathbb{Z}$, $h^i:A^i \to B^i$ such that
$f^i=d_B^{-1}\circ h^i \circ d_A^i$.
$$
\xymatrix{
A^{i-1}\ar[r]&
A^{i}\ar[r]\ar@{.>}[dl]^{h^i}\ar[d]^{f^i}&
A^{i+1}\ar@{.>}[dl]^{h^{i+1}}\\
B^{i-1}\ar[r]& B^i\ar[r]& B^{i+1}
}
$$
\item $f\sim_{\text{hom}}g$ are {\it
homotopically equivalent} if $f-g
\sim_{\text{hom}}0$.
\item The {\it homotopy category}
$\mathcal{H}(\mathcal{A})$ of $\mathcal{A}$
has objects
$\text{Ob}(\mathcal{H}(\mathcal{A})=
\text{Com}
(\mathcal{A})$
and morphisms
$$
\text{Hom}_{\mathcal{H}(\mathcal{A})
}(A^{\bullet},B^{\bullet})
=\text{Hom}_{\text{Com}\mathcal{A}}(A,B
)\Big/\{f\sim_{\text{hom}}0\}.
$$
\end{enumerate}
\end{definition}

\begin{exercise}
\label{derived-categories-of-sheaves-exercise-homotopy-category}
\begin{enumerate}
\item $f\sim_{\text{hom}}0$,
$H^i(f):H^i(A)\xrightarrow{0}H^i(B)$.
\item Morphisms homotopic to zero form an
ideal in $\text{Com}(\mathcal{A})$,
that is, abelian subgroups w.r.t.
compositions.
\item $\mathcal{H}(\mathcal{A})$ is an
additive category and
$\text{Com}(\mathcal{A}) \to
\mathcal{H}(\mathcal{A})$ is additive.
\end{enumerate}
\end{exercise}

{\bf Objectives.}
\begin{enumerate}
\item QIS satisfy the Ore conditions in
$\mathcal{H}(\mathcal{A})$.
\item
$D(\mathcal{A})=\text{Com}(\mathcal{A})
[\text{QIS}^{-1}]
\cong
\mathcal{H}(\mathcal{A})[\text{QIS}^{-1}]$.
This will allow to describe
$D(\mathcal{A})$ explicitly as an additive
category.
\item $\mathcal{H}(\mathcal{A})$ and
$D(\mathcal{A})$ are not abelian, but
triangulated!
\end{enumerate}

When we lose the property of an abelian
category, we lose the notion of exact
sequence.

{\bf Goal.} Now we shall prove that QIS
satisfy right Ore conditions on
$\mathfrak{X}(f)$ and $D(\mathcal{A})\cong
\mathcal{H}[\text{QIS}]^{-1}$.

\begin{definition}
\label{derived-categories-of-sheaves-definition-shift-functor}
The {\it shift functor} is given by
$$
\begin{aligned}
[1]: \text{Com}(\mathcal{A}) &\longrightarrow
\text{Com}(\mathcal{A}) \\
A^\bullet &\longmapsto (A^\bullet[1])
\end{aligned}
$$
where
$$
(A^\bullet[1])^i=A^{i+1},\qquad
d_{A[1]}=-d_A
$$
\end{definition}

\begin{definition}
\label{derived-categories-of-sheaves-definition-mapping-cone}
For a morphsm $f \in
\text{Hom}_{\text{Com}(\mathcal{A})
}(A^\bullet,B^{\bullet})$,
its {\it mapping
cone} is the complex
$$
\text{Cone}(f):=B^{\bullet}\oplus
A^{\bullet}[1]
$$
with
$$
\begin{aligned}
d_{\text{Cone}(f)}^i: A^{i+1}\oplus B^{i}
&\longrightarrow A^{i+2}\oplus B^{i+1} \\
(a^{i+1},b^i) &\longmapsto
(-d_Aa^{i+1},d_Bb^i+f(a^{i+1}))
\end{aligned}
$$
\end{definition}

\begin{exercise}
\label{derived-categories-of-sheaves-exercise-cone-is-complex}
$\text{Cone}(f) \in \text{Com}(\mathcal{A})$.
\end{exercise}

Exact sequence in $\text{Com}(\mathcal{A})$:
$$
\xymatrix{
0\ar[r]&B\ar[r]&\text{Cone}(f)\ar[r]&
A^{\bullet}[1]\ar[r]&0
}
$$
\begin{definition}
\label{derived-categories-of-sheaves-definition-contractible-complex}
$A^\bullet \in \mathcal{H}(\mathcal{A})$ is
{\it contractible} if $\mathcal{A}
\cong 0$ in $\mathcal{H}(\mathcal{A})$, that
is,
$\text{Id}_{A^{\bullet}}\sim_{\text{hom}}0$.
\end{definition}

\begin{exercise}
\label{derived-categories-of-sheaves-dcos-label-005}
$f \in \text{Hom}(A^{\bullet},B^{\bullet})$
is isomorphism in $\mathcal{H}(\mathcal{A})$
if and only if $\text{Cone}(f)$ is
contractible.
\end{exercise}

\begin{definition}
\label{derived-categories-of-sheaves-definition-Hom-of-complexes}
Let  $A^{\bullet},B^{\bullet}\in
\text{Com}(\mathcal{A})$. Then we define
$\underline{\text{Hom}}(A^{\bullet},
B^{\bullet})
\in \text{Com}(\text{Ab})$,
$$
\text{Hom}^i(A,B)=\prod_{n \in
\mathbb{Z}}\text{Hom}(A^n,B^{n+i})
$$
$$
\begin{aligned}
d^i_{\underline{\text{Hom}}}:
\underline{\text{Hom}}(A,B)^i
&\longrightarrow
\underline{\text{Hom}}(A,B)^{i+1} \\
f=\prod f^n &\longmapsto g
\end{aligned}
$$
and $g^n=d_Bf^n-(-1)^i f^{n+1}d_A$.
\end{definition}

\begin{exercise}
\label{derived-categories-of-sheaves-exercise-Hom}
\begin{enumerate}
\item Check that $\underline{\text{Hom}}$ is
a complex.
\item $\Ker
d^i_{\underline{\text{Hom}}}
=\text{Hom}_{\text{Com}(\mathcal{A})
}(A^{\bullet},B^{\bullet}[i])$.
\item
$\text{Im}d^{i-1}_{\text{Hom}}=
\{fsim_{\text{hom}}0\}$
\end{enumerate}
\end{exercise}

Notation:
$\text{Hom}^i(A^{\bullet},B^{\bullet})=
\text{Hom}(A^{\bullet},B^{\bullet}[i])$.

\begin{lemma}
\label{derived-categories-of-sheaves-lemma-cone-long-exact-sequence}
Consider the following composition
$$
\xymatrix{
A^\bullet\ar[r]_f&  B^{\bullet}\ar[r]^b&
\text{Cone}(f)\ar[r]^a &
A^{\bullet}[1]
}
$$
(which is your first {\it distinguished
triangle in life}) induces the following long
exact sequence for all $C^{\bullet \in
\text{Com}(\mathcal{A})}$:
$$
\xymatrix{
\cdots\ar[r]
&\text{Hom}^i_{\mathcal{H}(\mathcal{A})
}(C^{\bullet},A^{\bullet})\ar[r]
&\text{Hom}_{\mathcal{H}}(\mathcal{A})
^i(C^{\bullet},B^{\bullet})\ar[r]&\\
\ar[r]&\text{Hom}^i_{\mathcal{H}(\mathcal{A})
}(C,\text{Cone}(f))\ar[r]
&\text{Hom}^{i+1}_{\mathcal{H}(\mathcal{A})
}(C^{\bullet},A^{\bullet})\ar[r]
&\cdots
}
$$
and the same happens for
$\text{Hom}^i_{\mathcal{H}(\mathcal{A})
}(\cdot,C^{\bullet})$.
\end{lemma}

It is very important that we went in
$\mathcal{H}(\mathcal{A})$!

\begin{proof}[Idea of proof]
(The long exact sequence is the cohomology
exact sequence associated to a short exact
sequence. Also some exercise from the past is
hidden in the details.)

The short exact sequence
$$
\xymatrix{
0\ar[r]&B^{\bullet}\ar[r]&\text{Cone}(f)
\ar[r]
&A[1]\ar[r]&0
}
$$
splits. This gives
$$
\xymatrix{
0\ar[r]&B^i\ar[r]&B^i\oplus
A^{i+1}\ar[r]&A^{i+1}\ar[r]&0
}
$$
Which in turn gives
$$
\xymatrix{
0\ar[r]&\underline{\text{Hom}}(C^{\bullet},
B^{\bullet})
\ar[r]&\underline{\text{Hom}}(C^{\bullet},
\text{Cone}(f))
\ar[r]&\underline{\text{Hom}}(C^{\bullet},
A[1]
\ar[r]&0
}
$$
and {\it then} we take the cohomology long
exact sequence.
\end{proof}

\begin{exercise}
\label{derived-categories-of-sheaves-exercise-simple}
$bf=ab=fa[-1]=0$ in
$\mathcal{H}(\mathcal{A})$.
\end{exercise}

\begin{exercise}
\label{derived-categories-of-sheaves-dcos-label-006}
Prove that the right Ore conditions are not
satisfied in $\text{Com}(\mathcal{A})$.
\end{exercise}

\begin{proposition}
\label{derived-categories-of-sheaves-dcos-label-007}
QIS satisfy the right Ore conditions in
$\mathcal{H}(\mathcal{A})$.
\end{proposition}

\begin{proof}
In lecture we proved condition 1 and part of
condition 2 of the right Ore conditions, see
Subsection
\ref{subsection-explicit-description-of-local
ization}.
Finishing the proof was
left as exercise, including Ore condition 3.
\end{proof}

\bigskip\noindent
Now we shall show that
$$
D(\mathcal{A})\cong
\mathcal{H}(\mathcal{A})[\text{QIS}]
$$
$$
H:\text{Com}(\mathcal{A}\to
\mathcal{H}(\mathcal{A})
$$
\begin{proposition}
\label{derived-categories-of-sheaves-proposition-localization-functor}
The localization functor
$$
\xymatrix{
Q:\text{Com}(\mathcal{A})\ar[d]_H\ar[r]
&
D(\mathcal{A})\cong\mathcal{H}(\mathcal{A})
[\text{QIS}]^{-1}\\
\mathcal{H}(\mathcal{A})\ar[ur]_{Q'}
}
$$
can be be decomposed as $Q' \circ H$.
Moreover $Q' \cong
\mathcal{H}(\mathcal{A})$, where
$Q_{\mathcal{H}(\mathcal{A})}:\mathcal{H}
(\mathcal{A})
\to
\mathcal{H}(\mathcal{A})[\text{QIS}]^{-1}$.
\end{proposition}

\begin{proof}
Done in class, with some exercises for us to
complete.
\end{proof}

This shows that the derived category is
equivalent to the homotopy category with
localization. We have the following
corollary:

\begin{lemma}
\label{derived-categories-of-sheaves-lemma}
We can specifically describe
$$
D(\mathcal{A}): \text{Ob}(D(\mathcal{A}) =
\text{Com}(\mathcal{A})
$$
$$
\text{Hom}_{D(\mathcal{A})}(A^{\bullet},
B^{\bullet})
=\{\text{equivelnce classes of roofs}\}
$$
Recall that a roof is
$$
\xymatrix{
& C\ar[dl]_{\text{QIS}\ni S}\ar[dr]^{f}\\
A& & B
}
$$
\end{lemma}

\bigskip\noindent
Now we shall show that $D(\mathcal{A})$ is
additive. (Or have we already proved this?
ie. from the corollary?)

\bigskip\noindent
Why $\mathcal{H}(\mathcal{A})$ and
$D(\mathcal{A})$ are not abelian: there are
not so many injective and surjective
morphisms.

\begin{exercise}
\label{derived-categories-of-sheaves-dcos-label-008}
$f \in \text{Hom}_{\mathcal{H}(\mathcal{A})}
(A^{\bullet},B^{\bullet})$ is injective if
and only if there exists  $g: B \to
A^{\bullet} \in
\text{Hom}_{\mathcal{H}(\mathcal{A})}
(B^{\bullet},A^{\bullet})$ such that $g
\circ f=\text{id}_A$. Formulate the same for
surjections.
\end{exercise}

\section{Some properties of the derived
category}
\label{derived-categories-of-sheaves-dcos-label-009}

For every $i \in \mathbb{Z}$ define
$$
\begin{aligned}
\mathcal{A} &\longrightarrow D(\mathcal{A})
\\
A &\longmapsto A[i], \quad \{0 \to
\overset{i}{A} \to 0\}
\end{aligned}
$$

So $A[i]$ is {\it that} complex. It is
``concentrated in degree i''

{\bf Goal:} understand morphisms between
$M,N[i]$.

\begin{definition}
\label{derived-categories-of-sheaves-definition-embedding-in-Com}
Let $A \in \text{Com}(\mathcal{A})$. For all
$n$,
$$
\tau_{\leq n}(A^{\bullet})^i:=\begin{cases}
A^i\qquad &i<n \\
\Ker d^n\qquad &i=n\\
0,\quad &i>n
\end{cases}
$$
then
$$
\tau_{\geq  n}(A^{\bullet})^i=\begin{cases}
0\qquad
&i<n \\
A^n/\text{Im}d^{n-1}\qquad &i=n\\
A^i,\qquad &i>n
\end{cases}
$$
and $A^{\bullet}\to\tau_{\geq n}A^{\bullet}$.
\end{definition}

Then
$$
\tau_{\leq n}:\text{Com}(\mathcal{A}) \to
\text{Com}(\mathcal{A})
$$
defines a functor $\tau_{\leq
n}(\text{QIS})\subset \text{QIS}$ where
$$
\tau_{\leq }(f \sim _{\text{hom}}0) \subset
\sim_{\text{hom}}0
\implies
$$
$$
\begin{aligned}
\tau_{\leq n}: \mathcal{H}(\mathcal{A})
&\longrightarrow
\mathcal{H}(\mathcal{A}) \\
D(\mathcal{A}) &\longmapsto D(\mathcal{A})
\end{aligned}
$$
the same for $\tau_{\geq n}$.

\medskip
The following lemma allows us to understand
how are morphisms behaved.

\begin{lemma}
\label{derived-categories-of-sheaves-lemma-Homs}
Let $A,B \in \mathcal{A}$. Then
\begin{enumerate}
\item
$\text{Hom}_{D(\mathcal{A})}(A,B[-i])=0$ for
all $i>0$.
\item
$\text{Hom}_{\mathcal{A}}(A,B)=\text{Hom}
_{D(\mathcal{A})}(A,B)$.
\end{enumerate}
In particular, $\mathcal{A} \hookrightarrow
A(\mathcal{A})$ is fully-faithful.
\end{lemma}

Now we introduce a notion of Ext, which will
recover the usual Ext (derived functors).

\begin{definition}
\label{derived-categories-of-sheaves-definition-Ext}
For $i>0$,
$$
\text{Ext}^i_{\mathcal{A}}(A,B)=\text{Hom}
_{D(\mathcal{A})}(A,B[i]).
$$
\end{definition}
\begin{definition}[Yoneda]
\label{derived-categories-of-sheaves-definition-Yoneda}
For $i>0$ let
$$
\text{Ext}_Y^i(A,B):=
\{\substack{\text{equivalence
classes of exact sequences}\\
B\to K^{i-1}\to K^{i-2}\to \ldots \to K^0 \to
A}\}
$$
and say that two such sequences are
equivalent if there is a commutative diagram
$$
\xymatrix{
B\ar[r]&  K^{i-1}\ar[r]\ar[d]&
K^{i-2}\ar[r]\ar[d]& \cdots\ar[r]\ar[d]
&K^0\ar[r]\ar[d]&A\\
B\ar[r]&  \tilde{K}^{i-1}\ar[r]&
\tilde{K}^{i-2}\ar[r]& \cdots\ar[r]
&\tilde{K}^0\ar[r]&A
}
$$
and are equivalent if there exists a sequence
of equivalent elements
$$
K \sim K_0 \sim \ldots \sim \tilde{K}
$$
\end{definition}

\begin{proposition}
\label{derived-categories-of-sheaves-proposition-Ext}
$$
\text{Ext}_{\mathcal{A}}^i(A,B) \cong
\text{Ext}_Y^i(A,B)
$$
\end{proposition}

\begin{proof}
Sketched in class.
\end{proof}

\begin{exercise}
\label{derived-categories-of-sheaves-dcos-label-010}
$D(\mathcal{A})$ and
$\mathcal{H}(\mathcal{A})$ are abelian if and
only if $\mathcal{A}$ is semisimple, that is,
all exact sequences in $\mathcal{A}$ split.
If $\mathcal{A}$ is semisimple, then
$D(\mathcal{A})\cong \bigoplus_{i inn
\mathbb{Z}}\mathcal{A}[-i]$.
\end{exercise}

\section{Triangulated categories}
\label{derived-categories-of-sheaves-section-triangulated-categories}

\begin{definition}
\label{derived-categories-of-sheaves-definition-triangulated-category}
An additive category $\mathcal{T}$ is a {\it
triangulated category} if it has the
following structures:
\begin{itemize}
\item $[1]:\mathcal{T} \to \mathcal{T}$
autoequiv.
\item Class of sequences of the form
$$
\xymatrix{
A\ar[r]&B\ar[r]&C\ar[r]&A[1]
}
$$
which are called {\it distinguished
triangles}, and satisfy the following axioms:
\begin{itemize}
\item (Axiom 1.)
\begin{itemize}
\item $\xymatrix{
A\ar[r]^{\text{id}}&A \ar[r]& 0\ar[r]&A[1] }$
is a distinguished triangle.
\item Any triangle $A \to B \to C \to A[I]$
isomorphic to a distinguished
triangle is a distinguished triangle.
\item Any morphism $(A \xrightarrow{f}B)\in
\mathcal{T}$ can be completed to a
distinguished triangle $A\xrightarrow{f}B \to
\mathcal{K} \to A[I]$, and $\mathcal{K}$ is
called the {\it cone} of $f$.
\end{itemize}
\item (Axiom 2.) A triangle
$$
\xymatrix{
A\ar[r]^f &  B \ar[r]^g &  C\ar[r]^h & A[1]
}
$$
is distinguished if and only if
$$
\xymatrix{
B\ar[r]^g &  C \ar[r]^g &  A[1]\ar[r]^{f[1]}
& B[1]
}
$$
is distinguihsed.

\item (Axiom 3.) For any two distinguished
triangles
$$
\xymatrix{
A\ar[r]^f\ar[d]_a &  B \ar[r]\ar[d]^b &  C
\ar@{.>}[d]^{u}\ar[r]&
A[1]\ar[d]^{a[1]}\\
A'\ar[r]_{f'}& B'\ar[r]& C'\ar[r]&A'[1]
}
$$
and any $a,b$ such that $bf=f'a$ there exists
$u$ that makes the above diagram commute (it
is not unique!)

\item (Octahedral axiom.) Any two
distinguished triangles with a common vertex
of the form
$$
\xymatrix{
X\ar[r]^u& Y\ar[r]^v\ar[d]& Z\ar[r]&X[1]\\
& Y'\ar[d]\\
& W\ar[d]\\
& Y[1]
}
$$
can be completed to the following commutative
diagram:
$$
\xymatrix{
X\ar[d]_{=}\ar[r]^u& Y\ar[r]^v\ar[d]&
Z\ar[r]&  X[1]\ar[d]^=\\
X\ar[r]&  Y'\ar[d]\ar[r]&  Z'\ar[d]\ar[r]&
X[1]\ar[d]^{u[1]}\\
& W\ar[r]^=\ar[d]_h & W\ar[r]^h\ar[d]&
Y[1]\\
& Y[1]\ar[r]_{v[1]}&Z[1]
}
$$
\end{itemize}
\end{itemize}
\end{definition}

{\bf Properties.} Let
$X\xrightarrow{f}Y\xrightarrow{g}
Z\xrightarrow{h}X[1]$
be a distinguished triangle.

\begin{lemma}
\label{derived-categories-of-sheaves-dcos-label-011}
\begin{enumerate}
\item $gf=hg=f[1]h=0$.
\item It induces the following sequence
$$
\xymatrix{
\cdots\ar[r]&Z[-1]\ar[r]& X\ar[r]& Y\ar[r]&
Z\ar[r]&X[1]
\ar[r]&Y[1]\ar[r]& Z[1]\ar[r]&\cdots
}
$$
such that any sequence of length $h$
(contained in it?) is a distinguished
triangle.
\end{enumerate}
\end{lemma}

\begin{proof}
$$
\xymatrix{
Y\ar[r]^g\ar[d]_=& Z\ar[d]_=\ar[r]^h&
C[1]\ar[r]^{f[1]}\ar[d]&
Y[1]\ar[d]^{g[1]}\\
Z\ar[r]^{=}&Z\ar[r]& 0\ar[r]&Z[1]
}
$$
and by axiom 3 we get $gf=0$. By Axiom 2 we
get the second statement.
\end{proof}

\begin{lemma}
\label{derived-categories-of-sheaves-dcos-label-012}
If $X \to Y \to Z \to X[1]$ is a
distinguished triangle, then for every $U \in
\mathcal{T}$ we have the following long exact
sequences:
$$
\xymatrix{
\cdots\ar[r]&\text{Hom}(U,X[i])\ar[r]&
\text{Hom}(U,Y[i]\ar[r]&
\text{Hom}(U,Z[i])\ar[r]&\text{Hom}(U,X[i+1]
\ar[r]&\cdots
}
$$
and the same happens for
$\text{Hom}(\cdot,U)$
\end{lemma}

\begin{proof}
Done in lecture, uses previous lemma and
Axiom 3.
\end{proof}

We have the following corollary:

\begin{lemma}
\label{derived-categories-of-sheaves-dcos-label-013}
If in Axiom 3 two morphisms are isomorphisms,
then the constructed one is also an
isomorphism.
\end{lemma}

\begin{proof}
Exercise. {\bf Hint.} Use Yoneda lemma.
\end{proof}

\begin{exercise}
\label{derived-categories-of-sheaves-dcos-label-014}
If
$A\xrightarrow{f}B\xrightarrow{g}
C\xrightarrow{h}A[1]$
is a distinguished triangle, then $h=0$ if
and only if $f$ has a left inverse, if and
only if $g$ has a right inverse.
\end{exercise}

As a corollary,

\begin{lemma}
\label{derived-categories-of-sheaves-lemma-cone-of-morphism-is-unique}
The cone of a morphism is unique up to
isomorphism.
\end{lemma}

However, the cone is not functorial, which is
a problem.

\medskip\noindent
Now we want to show that the homotopy
category $\mathcal{H}(\mathcal{A})$ of an
abelian category $\mathcal{A}$ has a
triangulated structure. We need to say which
are the distinguished triangles: but we
already did this for the derived category:
they are
$$
\xymatrix{
A\ar[r]^f&B\ar[r]^b&\text{Cone}(f)\ar[r]^a&
A[I]\ar[r]&0
}
$$

\begin{theorem}
\label{derived-categories-of-sheaves-dcos-label-015}
$\mathcal{H}(\mathcal{A})$ is triangulated.
\end{theorem}

\begin{proof}
It looks like Axiom 1 should be immediate
from our construction. Also Axiom 2 (indeed,
we have a notion of $\text{Cone}(f)$ which we
obviously expect to satisfy the property of
the cone $\mathcal{K}$. Axiom 3 uses an
exercise using the notion of cylinder.
Proving the octahedral axiom is also an
exercise with a big hint.
\end{proof}

{\bf $D(\mathcal{A})$ is triangulated.} Let
$\mathcal{C}$ be triangulated, and $S$ a
class of morphisms satisfying the right Ore
conditions. We know that
$\mathcal{C}[S]^{-1}$ is an additive
category, but when is it triangulated?

\begin{definition}
\label{derived-categories-of-sheaves-dcos-label-016}
$S$ is {\it compatible} with a triangulated
structure if
\begin{enumerate}
\item $s \in S$, $S[1]\in S$,
\item In Axiom 3, if $a,b \in S$ then we can
choose $u \in S$.
\end{enumerate}
\end{definition}

\begin{proposition}
\label{derived-categories-of-sheaves-dcos-label-017}
It $\mathcal{T}$ is triangulated an $S$
satisfies the right Ore conditions and is
compatible with the triangulated structure of
$\mathcal{T}$ then $\mathcal{T}[S]^{-1}$ is
triangulated.
\end{proposition}

\begin{proof}
Axioms 1-3 proved, Octahedral axiom is
exercise with hint.
\end{proof}

We finally obtain:

\begin{lemma}
\label{derived-categories-of-sheaves-dcos-label-031}
$D(\mathcal{A})$ is triangulated.
\end{lemma}

\begin{proof}
Two or three lines, uses 5 lemma, uses long
exact sequences of $\text{Hom}$.
\end{proof}

\section{Derived functors}
\label{derived-categories-of-sheaves-section-derived-functors}

\begin{definition}
\label{derived-categories-of-sheaves-definition-additive-exact-functor}
An additive functor $F:\mathcal{T} \to
\mathcal{T}'$ between triangulated categories
is {\it exact} if
\begin{enumerate}
\item $F \circ
[1]_{\mathcal{T}}=[1]_{\mathcal{T}'}\circ F$.
\item Maps distinguished triangles to
distinguished triangles.
\end{enumerate}
\end{definition}

\medskip\noindent
Here are some bounded notions:
$$
\begin{aligned}
\text{Com}^\bullet(\mathcal{A})&\\
\text{Com}^+(\mathcal{A})&=\{A\in
\text{Com}(\mathcal{A})
|A^i=0,\quad i\ll 0\}\\
\text{Com}^-(\mathcal{A})&=\{A\in
\text{Com}(\mathcal{A})|i \gg 0\}\\
\text{Com}^b(\mathcal{A})&=\{A \in
\text{Com}(\mathcal{A})||i|\gg 0\}
\end{aligned}
$$
each of which gives a derived category $D^*$
for $*=b,+$ or $-$.

Now consider an exact functor of abelian
categories $F:\mathcal{A} \to
\mathcal{B}$. We would like to obtain a
functor
$D^*(\mathcal{A}) \to D^*(\mathcal{B})$.

In general functors are not exact. Thus our
goal now is:
\begin{itemize}
\item From a left exact functor
$F:\mathcal{A} \to \mathcal{B}$ obtain
$RF:D^+(\mathcal{A})\to D^+(\mathcal{B})$,
\item From a right exact functor
$G:\mathcal{A}\to \mathcal{B}$ obtain
$LG:D^-(\mathcal{A})\to D^-(\mathcal{B})$.
\end{itemize}
Given an additive category $\mathcal{I}
\subset \mathcal{A}$ of injective objects in
$\mathcal{A}$, we have  $\mathcal{H}^*
(\mathcal{I})\subset
\mathcal{H}^* (\mathcal{A})
\xrightarrow{G_{\mathcal{A}}}D^*(\mathcal{A})
$.

\begin{lemma}
\label{derived-categories-of-sheaves-dcos-label-018}
Suppose that $\mathcal{A}$ has enough
injective objects, $\forall A^\bullet \in
\mathcal{A}$ can be embedded $A
\hookrightarrow I \in\mathcal{I}$. Then
$i:\mathcal{H}^+(\mathcal{I}) \to
D^+(\mathcal{A})$ is an equivalence.
\end{lemma}

\begin{proof}
Not proved.
\end{proof}

\section{Semiorthogonal decompositions}
\label{derived-categories-of-sheaves-dcos-label-032}

{\bf Idea.} To split $D^*(\mathcal{A})$ in
smaller pieces that are easier.

\begin{definition}
\label{derived-categories-of-sheaves-dcos-label-019}
$\mathcal{A} \subset \mathcal{T}$ is a {\it
strictly full triangulated subcategory} if
\begin{itemize}
\item
$\text{Hom}_{\mathcal{A}}(A,B)=\text{Hom}
_{\mathcal{T}}(A,B)$
for all $A,B
\in \mathcal{A}$.
\item If  $B \in \mathcal{T}$ is isomorphic
to an object from $\mathcal{A}$,
then $B \in \mathcal{A}$.
\end{itemize}
\end{definition}

\begin{definition}
\label{derived-categories-of-sheaves-dcos-label-020}
Let $\mathcal{T}$ be a triangulated category
and $\mathcal{A}$, $\mathcal{B}$ strictly
full trinagulated subcategories in
$\mathcal{T}$. Then
$\mathcal{T}=\left\langle{}\mathcal{A},\mathcal{B}
\right\rangle{}$
is a {\it semiorthogonal decomposition} if
\begin{itemize}
\item $\text{Hom}_{\mathcal{T}}(B,A)=0$ for
all $B \in \mathcal{B}$ and $A\in
\mathcal{A}$.
\item For all $T \in \mathcal{T}$ there
exists a distinguished triangle
$$
T_B \to T \to T_{\mathcal{A}}\to
T_{\mathcal{B}}[1]
$$
where $T_B \in \mathcal{B}$ and
$T_\mathcal{A} \in \mathcal{A}$.
\end{itemize}
\end{definition}

\begin{proposition}[Properties of
semiorthogonal decomposition]
\label{derived-categories-of-sheaves-dcos-label-021}
\begin{enumerate}
\item Functoriality: for any $f \in T \to
T'$,
$$
\xymatrix{
T_{\mathcal{B}}\ar@{.>}[d]_{\exists
f_{\mathcal{B}}}\ar[r]& T\ar[d]^f\ar[r]
&
T_{\mathcal{A}}\ar@{.>}[d]^{f_{\mathcal{A}}}
\ar[r]&
T_{\mathcal{B}}[1]\\
T_{\mathcal{B}}'\ar[r]& T'\ar[r]&
T'_{\mathcal{A}}\ar[r]& T'_{\mathcal{B}}[1]
}
$$
there exist $f_{\mathcal{B}}$ and
$f_{\mathcal{A}}$ as in the diagram.
\end{enumerate}
\end{proposition}

\begin{proof}
For $f_{\mathcal{A}}$,apply
$\text{Hom}_{\mathcal{T}}(\cdot,
T_{\mathcal{A}
}')$
to the upper distinguished triangle, obtain a
long exact sequence\ldots{}
\end{proof}

Functoriality: defnote by $\alpha:\mathcal{A}
\hookrightarrow  \mathcal{T}$ the inclusion,
then we have a functor $\alpha^*:\mathcal{T}
\to \mathcal{A}$ given by $T \to
T_{\mathcal{A}}$, $f \mapsto
f_{\mathcal{A}}$. Similarly for
$\beta:\mathcal{B} \hookrightarrow
\mathcal{T}$ we have a functor $\beta
^*:\mathcal{T} \to \mathcal{B}$ by $T \to
T_{\mathcal{B}}$, $f \mapsto
f_{\mathcal{B}}$.

\begin{lemma}
\label{derived-categories-of-sheaves-lemma-alpha-and-beta-adjoint}
$\alpha^*$ is left adjoint to $\alpha$ and
$\beta^*$ is right adjoint to $\beta.$
\end{lemma}

\begin{proof}
We have for any $A \in \mathcal{A}$,
$\text{Hom}_{\mathcal{A}}(\alpha^*T,A)=
\text{Hom}_{\mathcal{T}}(T,\alpha_*A)$.
But that is equal to
$\text{Hom}_{\mathcal{A}}(\mathcal{T}
_{\mathcal{A}},A)$.
\end{proof}

\begin{exercise}
\label{derived-categories-of-sheaves-exercise-orthogonal}
If
$\mathcal{T}=\left\langle{}\mathcal{A},\mathcal{B}
\right\rangle{}$,
then
$$
\mathcal{B}=\mathcal{A}^\perp=\left\langle{}B \in
\mathcal{T} |
\text{Hom}(B,A), \forall A \in
\mathcal{A}\right\rangle{}:=\Ker \alpha^*
$$
(the $\perp$ sign should be on the left of
$\mathcal{A}$ but I don't know how to write
this) and
$$
\mathcal{A}=\mathcal{B}^\perp=\left\langle{}A \in
\mathcal{T} |
\text{Hom}(B,A)=0,\forall B \in
\mathcal{B}\right\rangle{}:=\Ker B
$$
(here the $\perp$ is correct on the right
side of $\mathcal{B}$).
\end{exercise}

This construction can be reversed.

\begin{definition}
\label{derived-categories-of-sheaves-dcos-label-022}
A strictly full subcategory $\mathcal{A}
\hookrightarrow  \mathcal{T}$ is called
\begin{itemize}
\item {\it left admissible} is there exists
$\alpha^*:\mathcal{T} \to
\mathcal{A}$ left adjoint,
\item {\it right admissible} is there exists
$\alpha^!:\mathcal{T} \to
\mathcal{A}$ right adjoint.
\end{itemize}
\end{definition}

\begin{lemma}
\label{derived-categories-of-sheaves-dcos-label-023}
If $\mathcal{A} \hookrightarrow  \mathcal{T}$
is left admissible, then
$\mathcal{T}=\left\langle{}\mathcal{A},\mathcal{A}
^\perp\right\rangle{}$
is semiorthogonal decomposition. If
$\mathcal{B} \hookrightarrow  \mathcal{T}$ is
right admissible then
$\mathcal{T}=\left\langle{}\mathcal{B}^\perp,
\mathcal{B}\right\rangle{}$ is semiorthogonal
decomposition.
\end{lemma}

\begin{proof}
Uses a remark, Yoneda lemma.
\end{proof}

\begin{proposition}[Bondall, Van der Borgh]
\label{derived-categories-of-sheaves-proposition-Bondall-Van-der-Borgh}
If $\mathcal{T}$ is smooth and proper (if
$\mathcal{T}=D^b(X)$, where $X$ is smooth and
proper, we call $\mathcal{T}$ smooth and
proper) then $\mathcal{A} \subset
\mathcal{T}$ is left admissible $\iff$
$\mathcal{A}$ is right admissible $\iff$
$\mathcal{A}$ is smooth proper. If
$\mathcal{A}$ is smooth and proper and
$\mathcal{A} \hookrightarrow  \mathcal{T}$
then $\mathcal{A}$ is left and right
admissible.
\end{proposition}

\medskip\noindent
If $D^b(C) \hookrightarrow  \mathcal{T}$ is
fully faithful, then it is left and right
admissible.

{\bf Question.} When can we construct a full
faithful embedding
$$
D^b(\text{pt})\cong D^b(\text{mod}k)
\hookrightarrow  \mathcal{T}?
$$
Let $$
\begin{aligned} D^b(\text{mod}k)
&\longrightarrow \mathcal{T} \\ k
&\longmapsto E
\end{aligned}
$$
and define
$$
\varphi_E:V^\bullet \to V^\bullet \otimes E
$$
where $V^\bullet$ is a graded vector space.

When is $\varphi_E$ full faithful? Some
computations show that $\varphi_E$ is fully
faithful if and only if $\text{Hom}(E,E)=k$.
Recall from
\texttt{categories.tex} that

\begin{definition}
\label{derived-categories-of-sheaves-definition-faithful}
Let $F : \mathcal{A} \to \mathcal{B}$ be a
functor.
\begin{enumerate}
\item We say $F$ is {\it faithful} if
for any objects $x, y \in \Ob(\mathcal{A})$
the map
$$
F : \Mor_\mathcal{A}(x, y) \to
\Mor_\mathcal{B}(F(x), F(y))
$$
is injective.
\item If these maps are all bijective then
$F$ is called
{\it fully faithful}.
\end{enumerate}
\end{definition}

\medskip\noindent
\begin{definition}
\label{derived-categories-of-sheaves-definition-exceptional-object}
$E \in \mathcal{T}$ is exceptional if
$$
\text{Hom}(E,E[i])=
\begin{cases}
k\qquad &i=0 \\
0\qquad &i \neq 0
\end{cases}
$$
\end{definition}

\begin{lemma}
\label{derived-categories-of-sheaves-dcos-label-024}
If $E$ is exceptional then
$\varphi_E:D(\text{pt}) \to \mathcal{T}$ is a
fully faithful embedding. Then
$$
\mathcal{T}=\left\langle{}\varphi_E(D^b(\text{pt})
^\perp,E\right\rangle{}=
\left\langle{}E^\perp,\varphi_E(D^b(\text{pt}))
\right\rangle{}
$$
\end{lemma}

{\bf Notation.} In this case,
$\mathcal{T}=\left\langle{}E,E^\perp\right\rangle{}=
\left\langle{}E^\perp,E\right\rangle{}$,
where on the first bracket the $\perp$ should
be on the left side of $E$.

\begin{example}
\label{derived-categories-of-sheaves-dcos-label-025}
On  $\mathbb{P}^n$, $\mathcal{O}(i)$ is
exceptional. More general, if $X$ is Fano
(i.e. smooth projective with $-K_X$ ample),
then any line bundle on  $X$ is ample.
\end{example}

\begin{proof}
$\text{Ext}^i(\mathcal{L},\mathcal{L})=H^{i}
(X,\mathcal{O}_X)$,
$H^{0}(X,\mathcal{O})=k$ and
$H^{>0}(X,\mathcal{O})=0$ by Kodaira
Vanishing.
\end{proof}

For Fano varieties we have nontrivial
semiorthogonal decomposition.

\begin{remark}
\label{derived-categories-of-sheaves-dcos-label-026}
If $K_X \cong \mathcal{O}_X$ then there is no
semiorthogonal decomposition of $D^b(X)$ (for
K3 or abelian).
\end{remark}

Back to Fano,
\begin{definition}
\label{derived-categories-of-sheaves-dcos-label-027}
Let $\mathcal{A}_1,\ldots,\mathcal{A}_n \in
\mathcal{T}$ be strictly full subcategories
and
$\mathcal{T}=\left\langle{}\mathcal{A}_1,\ldots,
\mathcal{A}_n\right\rangle{}$
a {\it semiorthogonal decomposition}
\begin{itemize}
\item
$\text{Hom}(\mathcal{A}_j,\mathcal{A}_i)=0$
for all $j>i$.
\item For all $T$ there is a chain of
morphisms
$$
0 = T_n \to T_{n-1}\to \ldots \to T_1 \to
T_0=T
$$
such that $\text{Cone}(T_i \to T_{i-1}) \in
\mathcal{A}_i$.
\end{itemize}
\end{definition}

\begin{exercise}
\label{derived-categories-of-sheaves-exercise-functoriality-and-cone}
Prove functoriality of $T \to \text{Cone}(T_i
\to T_j)$ for $i>j$.
\end{exercise}

\begin{definition}
\label{derived-categories-of-sheaves-dcos-label-028}
A collection of objects $E_1,\ldots, E_n$ is
{\it exceptional} if
\begin{enumerate}
\item For all $i$, $E_i$ is exceptional.
\item
$\text{Hom}(\left\langle{}E_i\right\rangle{},
\left\langle{}E_j\right\rangle{})
=0$
for all $i>j$, where
$\left\langle{}E_j\right\rangle{}$ is a minimal triangulated
category generated by $E_i$.
\end{enumerate}
\end{definition}

We have a corollary:

\begin{lemma}
\label{derived-categories-of-sheaves-dcos-label-029}
If $E_1,\ldots,E_n$ is an exceptional
collection, then
$$
\mathcal{T}=\left\langle{}E_1,\ldots,E_n,\left\langle{}E_1,
\ldots,E_n\right\rangle{}^\perp\right\rangle{}
=\left\langle{}\left\langle{}E_1,\ldots,E_n\right\rangle{}^\perp,E_1,
\ldots,E_n\right\rangle{}
$$
(where $\perp$ should be on the left of the
bracket in the first appearence) is a
semiorthogonal decomposition.
\end{lemma}

\begin{definition}
\label{derived-categories-of-sheaves-dcos-label-033}
If $\left\langle{}E_1,\ldots,E_n\right\rangle{}^\perp=0$
(equivalently with $\perp$ on the left).
\end{definition}

\begin{exercise}
\label{derived-categories-of-sheaves-dcos-label-030}
Let $X$ be a Fano variety with $-K_X=mH$ for
$H$ ample ($m$ is the index of $X$). Then for
every line bundle $\mathcal{L}$,
$$
(\mathcal{L}(-m+1)H,\ldots,\mathcal{L}(-H)
\mathcal{L})
$$
is an exceptional collection.
\end{exercise}

\begin{proof}
Use Serre duality and Kodaira Vanishing.
\end{proof}

\medskip\noindent
Here is the promised theorem from the
beginning of this minicourse:

\begin{theorem}[Beilinson]
\label{derived-categories-of-sheaves-theorem-Beilinson}
$D^b(\mathbb{P}^n)=\left\langle{}\mathcal{O}(-n),
\ldots,\mathcal{O}\right\rangle{}$
is a full exceptional collection.
\end{theorem}

This means that we can classify all objects
of $D^b(\mathbb{P}^n)$. Because whenever we
have a semiorthogonal decomposition we get
these chains of morphisms, i.e. for all
$D^b(\mathbb{P}^n)$ there exists a chain
$$
0=T_n \to \ldots \to T_0=T
$$
such that $\text{Cone}(T_i \to T_{i-1}) \in
\left\langle{}\mathcal{O}(-i)\right\rangle{} \cong
D^b(\text{pt})$.

Probably this is the content of
\cite{Beilinson-derived}.

\begin{proof}[Proof of Beilinson theorem]
Consider the Koszul complex
$$
0\to \mathcal{O}(-n-1)\to\ldots \to
\Lambda^{2}V^\vee \otimes \mathcal{O}(-2)
\to V^{\vee}\otimes \mathcal{O}(-1) to
\mathcal{O} \to0
$$
This implies that
$$
\mathcal{O}(n-1) \in \left\langle{}
\mathcal{O}(-n),\ldots,\mathcal{O}\right\rangle{}
$$
\ldots{} proof completed in class.
\end{proof}

\section{Kuznetsov - Video}
\label{derived-categories-of-sheaves-section-video}

\noindent
In this section I go over
\href{https://www.mathnet.ru/php/
presentation.phtml?option_lang=eng&presentid=
337}{this
video} of a talk by Kuznetsov where he
explains derived category in a very synthetic
way.

\begin{example}
\label{derived-categories-of-sheaves-example-simplest-derived-category}
Perhaps the simplest example of a derived
category is when $X=\Spec k$ is a point. In
this case,
$$
D^b(X)=D^b(k)=\left\{\bigoplus_tk^{n_t}[t]:
n_t=0\text{ for }|t|\gg 0\right\}
$$
with morphisms
$$
\Hom(k,k[t])
=\begin{cases}
k\qquad &t=0 \\
0\qquad &t\neq 0.
\end{cases}
$$
\end{example}

\noindent
Now let's define semiorthogonal
decomposition. Let $\mathcal{T}$ be a
category and
$\mathcal{A}_1,\mathcal{A}_2,\ldots,
\mathcal{A}_m
\subset \mathcal{T}$ full triangulated
subcategories. This means that each of these
subcategories is invariant under the shift
functor and under taking cones; i.e. if a
complex is in $\mathcal{A}_i$ then so is its
shift, and if two complexes are in
$\mathcal{A}_i$ then so is their cone. Then
$\mathcal{A}_1,\ldots,\mathcal{A}_m$ are a
{\it semiorthogonal decomposition} of
$\mathcal{T}$ if
$\Hom(\mathcal{A}_i,\mathcal{A}_j)=0$ for
$i>j$, and if for every object $T \in
\mathcal{T}$ there exist
$$
\xymatrix@C=1.5em{
0\ar[r]&T_m\ar[rr]&&\ar[dl]T_{m-1}\ar[r]
&\cdots\ar[r]&
T_2\ar[rr]&&\ar[dl]T_1\ar[rr]&&\ar[dl]T_0=T\\
&&A_m
\ar@{-->}[ul]&&&&A_2\ar@{-->}[ul]&&
A_1\ar@{-->}[ul]
}
$$
with $A_i \in \mathcal{A}_i$. That is, it's
like if there was a filtration of $T$ with
the components in the $\mathcal{A}_i$. In
this case we write
$\mathcal{T}=\left\langle{}\mathcal{A}_1,\ldots,
\mathcal{A}_m\right\rangle{}$.

An interesting question is: suppose that you
have a full triangulated subcategory
$\mathcal{A}_1 \subset \mathcal{T}$. Is it
possible to extend this to a semiorthogonal
decomposition? Actually, usually it is
impossible. It is a complicated problem to
construct semiorthogonal decompositions. What
we need is adjoint functors.

We say that $\mathcal{A}$ is {\it admissible}
if there exist $\alpha^*\alpha^!:\mathcal{T}
\to \mathcal{A}$ left and right adjoint. It
is easy to check that if $\mathcal{A}$ is
admissible, then
$\mathcal{T}=\left\langle{}\mathcal{A},
\leavevmode^\perp\mathcal{A}\right\rangle{}$
and
$\mathcal{T}=\left\langle{}\mathcal{A}^\perp,
\mathcal{A}\right\rangle{}$
are semiorthogonal decompositions, where
these two subcategories are the {\it left and
right orthogonal subcategories}
$$
\begin{aligned}
\leavevmode^\perp\mathcal{A}&=\{T \in
\mathcal{T}: \Hom(T,\mathcal{A})=0\}\\
\mathcal{A}^\perp&=\{T \in
\mathcal{T}:\Hom(\mathcal{A},T)=0\}.
\end{aligned}
$$
In fact for the first orthogonal
decomposition we only need the left adjoint
functor, and for the second, only the right
adjoint functor.

If $X$ is a connected Calabi-Yau manifold,
it's easy to see that $D^b(X)$ has no
semiorthogonal decomposition. This is because
since the canonical class is trivial then by
Serre duality we obtain that any
semiorthogonal decomposition is in fact
completely orthogonal, but it's easy to see
that there are no completely orthogonal
decompositions in any connected algebraic
variety.





% --- End derived-categories-of-sheaves.tex ---

% --- Begin homological-mirror-symmetry.tex ---
\chapter{Homological mirror symmetry}
\label{chapter-homological-mirror-symmetry}



\section{Starting point}
\label{homological-mirror-symmetry-section-starting-point}

\noindent
Started on 28 September 2026 after
Sergey Galkin's message about quantum
tori and matrix factorizations.
These introductory notes were drafted by
ChatGPT. Research questions below are
not established results.

The relevant background in HP is
symplectic geometry (Poisson brackets),
geometric prequantization (commutators),
homological algebra (complexes), and
derived categories of sheaves.

\section{Where this sits}
\label{homological-mirror-symmetry-section-where-this-sits}

\noindent
Homological mirror symmetry relates
categories from symplectic geometry
to categories from algebraic geometry.
On the symplectic side, Fukaya categories
involve Lagrangian submanifolds and
Floer theory; these require further study.
On the algebraic side, one encounters
derived categories of coherent sheaves
and categories of matrix factorizations.

\begin{definition}
\label{homological-mirror-symmetry-definition-lg-model}
A \textit{Landau--Ginzburg model}, in
the setting used here, is a complex
manifold or smooth algebraic variety
$Y$ equipped with a holomorphic or
regular function
$$
W:Y\longrightarrow\mathbb{C}.
$$
The function $W$ is its
\textit{superpotential}.
\end{definition}

\noindent
A Fano mirror may be a Landau--Ginzburg
model. A del Pezzo surface is a smooth
projective surface with ample
anticanonical bundle.
There are two directions to distinguish:
coherent sheaves on the Fano side relate
to a symplectic category of its mirror;
the Fukaya category of the Fano relates
to matrix factorizations on the mirror,
with suitable choices and hypotheses.
This is an orientation, not a theorem
asserted in this generality.

Quantum cohomology modifies the cup
product using counts of rational curves.
Deformation quantization modifies
multiplication of functions into a
noncommutative product. These are
different constructions, even though
mirror symmetry relates their settings.

For the classical period integrals and
their differential equations, see
\ref{variations-of-hodge-structures-%
section-quantum-periods-background}
in {\it Variations of Hodge Structures}.
That section also locates the separate
quantum-period construction.

\section{Tori and Mumford degenerations}
\label{homological-mirror-symmetry-section-tori-and-mumford}

\begin{definition}
\label{homological-mirror-symmetry-definition-algebraic-torus}
The \textit{algebraic torus} of
dimension $n$ over $\mathbb{C}$ is
$$
T=(\mathbb{C}^{*})^n.
$$
Its coordinate algebra is
$$
\mathcal{O}(T)=
\mathbb{C}[x_1^{\pm1},\ldots,x_n^{\pm1}].
$$
\end{definition}

\noindent
In toric geometry one embeds this torus
as a dense open subset of a variety,
with an extending torus action.
Fans describe normal toric varieties.
Toric geometry supplies building blocks
for Mumford degenerations of abelian
varieties.

A quantum torus instead deforms the
multiplication of the Laurent algebra.
Thus the common starting point is the
algebraic torus, but the constructions
are different. A degeneration parameter
in Mumford's construction is not
automatically a quantization parameter.

\begin{definition}
\label{homological-mirror-symmetry-definition-quantum-torus}
For $q\in\mathbb{C}^{*}$, the
two-dimensional \textit{quantum torus}
is the algebra
$$
A_q=
\mathbb{C}\langle x^{\pm1},y^{\pm1}\rangle
/(xy-qyx).
$$
\end{definition}

\noindent
At $q=1$ this is the ordinary
commutative Laurent algebra.
For general $q$ it is an algebra,
not the coordinate ring of an ordinary
commutative affine variety.

\section{First quantum calculation}
\label{homological-mirror-symmetry-section-first-quantum-calculation}

\noindent
Work formally with $q=e^\hbar$.
The relation gives
$$
xy-yx=(e^\hbar-1)yx.
$$
Dividing by $\hbar$ and reducing
modulo $\hbar$ yields
$$
\{x,y\}=xy.
$$
Here the bracket is defined by the
first-order commutator of lifts.
This convention uses the commutator
without a factor of $i$.
Compare the Dirac rule in the
prequantization notebook, where a
different normalization is used.

\begin{exercise}
\label{homological-mirror-symmetry-exercise-laurent-bracket}
Use the Leibniz rule to show that
$$
\{x^a y^b,x^c y^d\}
=(ad-bc)x^{a+c}y^{b+d}
$$
for integers $a,b,c,d$.
\end{exercise}

\section{Matrix factorizations}
\label{homological-mirror-symmetry-section-matrix-factorizations}

\begin{definition}
\label{homological-mirror-symmetry-definition-matrix-factorization}
Let $R$ be a commutative ring and
$W\in R$. A finite free
\textit{matrix factorization} of $W$
consists of finite free modules
$E^0,E^1$ and maps
$$
d_0:E^0\longrightarrow E^1,
\qquad
d_1:E^1\longrightarrow E^0
$$
such that
$$
d_1d_0=W I_{E^0},
\qquad
d_0d_1=W I_{E^1}.
$$
Equivalently, an odd endomorphism $D$
of $E^0\oplus E^1$ satisfies
$$
D^2=W I.
$$
\end{definition}

\noindent
Compare a complex, whose differential
squares to zero. Reducing the maps
modulo $W$ gives a two-periodic complex
over $R/(W)$.
Defining the morphism complexes and the
homotopy category is a next step.

\begin{example}
\label{homological-mirror-symmetry-example-xy-factorization}
Over $R=\mathbb{C}[x,y]$, take
$$
D=\begin{pmatrix}0&x\\y&0\end{pmatrix}.
$$
Then
$$
D^2=\begin{pmatrix}xy&0\\0&xy\end{pmatrix}.
$$
Over the quantum algebra, the same
matrix multiplication gives instead
$$
D^2=\begin{pmatrix}xy&0\\0&yx\end{pmatrix}.
$$
This shows why the naive construction
does not transfer unchanged.
The commutative example is on the
affine plane, not on the torus:
on the torus $xy$ is a unit and
this factorization is contractible.
\end{example}

\noindent
For a left module over a noncommutative
algebra, multiplication by $W$ need
not be module-linear. Centrality of
$W$ guarantees linearity.
Failure of centrality obstructs the
ordinary definition; it does not prove
that every generalized version is
impossible.

\section{Galkin's question}
\label{homological-mirror-symmetry-section-galkin-question}

\noindent
On 28 September 2026, Sergey proposed:
\begin{enumerate}
\item Check whether Brown--Wemyss
normal-form methods can be applied
to elements of a quantum torus.
\item Investigate a $q$-deformation
of the mirror matrix factorization
category, provisionally denoted
$\operatorname{MF}(W_q)$.
\item If such a category can be
defined, seek its interpretation
in symplectic geometry of del Pezzo
surfaces.
\end{enumerate}

\noindent
Brown--Wemyss work with completed free
algebras and cyclic derivatives.
A quantum torus already imposes
relations and invertibility, so their
results do not apply automatically.
Their paper is a possible source of
methods, not a construction of the
category requested here.

First tasks: understand one ordinary
potential and its matrix factorizations;
choose an explicit $W_q$; check its
centrality; specify what category and
classical limit are intended.

\section{People and papers to follow}
\label{homological-mirror-symmetry-section-people-and-papers}

\noindent
Status checked on 28 September 2026.
These are reading records, not claims
that the papers have been understood.

\begin{itemize}
\item Victor: Daniel identifies his
PhD topic as matrix factorizations.
The likely full-name match is
Victor Ibrahim Santos El Adji,
Galkin's student working on graph
potentials. No matching arXiv paper
was located in this check.
This does not establish that none exists.
Confirm the exact title or arXiv link.

\item Alex Junior Gomez Saltachin:
\textit{Exceptional collections for
canonical stacks of log del Pezzo
surfaces with $\frac13(1,1)$
singularities},
\href{https://arxiv.org/abs/2606.18238}
{arXiv:2606.18238}.
The arXiv page checked lists version 4,
30 August 2026. His website uses the
title \textit{Derived categories and
Brauer groups of log del Pezzo surfaces
with 1/3(1,1) singularities} for the
same identifier; do not count these as
two different papers.

\item Alex's website also lists
\textit{Categorical genus for log del
Pezzo surfaces with cyclic quotient
singularities},
\href{https://arxiv.org/abs/2609.13102}
{arXiv:2609.13102}.
The listing was verified on his
\href{https://alequisgs.github.io/}
{personal website}; direct retrieval of
this arXiv page failed in this check.
\end{itemize}


\section{Hyunbin Kim: log Calabi--Yau
surfaces}
\label{homological-mirror-symmetry-section-kim-log-cy-surfaces}

\noindent
Talk: \textit{Mirror symmetry of
log Calabi--Yau surfaces, with an eye
towards 3D mirror symmetry}.
The supplied slides identify the speaker
as Hyunbin Kim and the work as joint
with Hansol Hong. The title slide says
MIST, CUHK, 27 September 2026
(Sunday). Source: the supplied PDF
\texttt{MIST - 2026-09-27.pdf}.

\medskip\noindent
\textbf{Upshot.}
For a surface $X$ with an anticanonical
boundary $D$, the complement $X\setminus D$
is the log Calabi--Yau setting.
Counts of holomorphic disks bounded by
Lagrangian torus fibers construct a
mirror potential $W$.
Beyond the toric case, disk corrections
change how the mirror charts are glued.
The talk uses tropical and energy
methods to track critical points
through surface blowups and blowdowns.

For generic K\"ahler parameters and small
exceptional areas, the slides state
that $W$ is Morse and
$$
QH^*(X)\cong\operatorname{Jac}(W).
$$
Here Morse means that the critical points
are nondegenerate; the Jacobi algebra
is the function algebra modulo the
critical-point equations.
The geometric domain matters:
extra algebraic critical points outside
that domain are excluded.
The critical values correspond to
eigenvalues of quantum multiplication
by $c_1(X)$.

\medskip\noindent
For our discussion, this gives a concrete
connection between torus fibrations,
mirror potentials and quantum cohomology.
It does not construct Galkin's proposed
noncommutative category
$\operatorname{MF}(W_q)$.
The final slides pose further questions
about higher dimensions and symplectic
quotients; these are directions for
extension, not a completed general
higher-dimensional theorem.


\section{The tautological object}
\label{homological-mirror-symmetry-section-tautological-object}

\noindent
Source: \textit{The tautological object:
Universal families of stable objects
in a triangulated category},
dated 28 September 2026, supplied as
\texttt{tautological\_object.pdf}.
The note asks where a universal family
lives when the objects being parametrized
belong to a category rather than to a
chosen variety. Its central distinction
is between a moduli space, which records
isomorphism classes, and a moduli stack,
which also retains automorphisms and
families. In the enhanced categorical
setting of the note, the stack carries
a tautological family: pulling it back
along a map that classifies a family
recovers that very family. This makes
the construction intrinsic to the
category and the chosen moduli problem,
rather than to a geometric realization.

\medskip\noindent
The practical warning is that an ordinary
universal family on the coarse moduli
space need not exist; the obstruction
is encoded by a twisting class.
When such a family does exist, tensoring
by a line bundle from the parameter
space produces another universal family.
Thus comparing two geometric models
requires attention to normalization,
even if their moduli spaces agree.
The note illustrates this on an elliptic
curve. For HMS, the relevance is learning
how families of objects behave under
categorical equivalences; for our HP
background, this connects directly to
derived categories, stability and moduli
of sheaves. This is a reading summary
of the supplied note, not a verification
of all its statements.

\section{Initial reading}
\label{homological-mirror-symmetry-section-initial-reading}

\begin{itemize}
\item V. I. Arnold,
\textit{Local Normal Forms of Functions},
Inventiones mathematicae 35 (1976),
87--109.
\href{https://webhomes.maths.ed.ac.uk/%
~v1ranick/papers/arnold15.pdf}{PDF}.
Classical local singularity theory.

\item Gavin Brown and Michael Wemyss,
\textit{Local Normal Forms of
Noncommutative Functions},
\href{https://arxiv.org/abs/2111.05900}
{arXiv:2111.05900}.
Noncommutative normal forms.

\item Denis Auroux, Ludmil Katzarkov,
and Dmitri Orlov,
\textit{Mirror symmetry for Del Pezzo
surfaces: Vanishing cycles and
coherent sheaves}.
\par
\href{https://arxiv.org/abs/math/0506166}
{arXiv:math/0506166}.
A reference for the classical
del Pezzo mirror framework, not a
solution of Galkin's quantum question.

\item Denis Auroux,
\textit{Mirror symmetry and T-duality
in the complement of an anticanonical
divisor},
\href{https://arxiv.org/abs/0706.3207}
{arXiv:0706.3207}.
Geometric motivation for mirror
potentials.
\end{itemize}

% --- End homological-mirror-symmetry.tex ---

% --- Begin geometric-invariant-theory.tex ---
\chapter{Geometric Invariant Theory}
\label{chapter-geometric-invariant-theory}



\noindent
Minicourse by Michel Brion (Universit\'e
Grenoble Alpes, Institut Fourier, France),
CIMPA school Florian\'opolis 2025.

\medskip\noindent
Notes at
\href{%
http://github.com/danimalabares/cimpa-floripa
}{%
github.com/danimalabares/cimpa-floripa
}

\bigskip\noindent

{\bf Abstract.} Geometric Invariant Theory
(GIT) studies the actions of reductive
algebraic groups on algebraic varieties and
the constructions of quotients in algebraic
geometry. One of the most relevant aspect of
GIT is its application to moduli theory: we
can indeed approach moduli problems trying to
construct moduli spaces as GIT quotients. In
this course we will present the basic notions
of geometric invariant theory, we will
illustrate its connections to moduli and
provide concrete examples of moduli spaces
obtained as GIT quotients.

\bigskip\noindent
\bigskip\noindent

{\bf Plan.}

\begin{enumerate}
\item Quotients by finite groups.
\item Linear algebraic groups.
\item GIT (affine case).
\item GIT (general case).
\end{enumerate}

\section{Introduction}
\label{geometric-invariant-theory-section-introduction}

Let $G$ be a finite group of order $N$, $G
\subset \text{Aut}(X)$,  where $X$ is an
affine variety over a field  $k=\overline{k}$
of arbitrary characteristic $p \geq 0$. $G$
acts on
$\Gamma(X,\mathcal{O}_X)=\mathcal{O}(X)$ by
$$
(gf)(x)=f(g^{-1}x).
$$
Define
$$
\mathcal{O}(X)^G:=\{f|gf=f\;\forall g \in
G\}\subset\mathcal{O}(X)
$$

\begin{proposition}
\label{geometric-invariant-theory-git-label-018}
$\mathcal{O}(X)^G$ is finitely generated and
$\mathcal{O}(X)$ is a finite module over
$\mathcal{O}(X)^G$.
\end{proposition}

\begin{proof}
Let $A:=\mathcal{O}(X) \supset
\mathcal{O}(X)^G:=B$. Then the coordinate
ring of $X$ is a finitely presented algebra
$\mathcal{O}(X)=k[f_1,\ldots,f_n]$. For every
generator of  $\mathcal{O}(X)$ define a
polynomial $P_i(t):=\prod_{g \in G}(t-gf_i)
\in A^G[t]=B[t]$.
\end{proof}

Let $Y = \text{Spec}B$ and $\pi:X\to Y$ a
dominant, finite morphism, which implies it
is surjective.

\begin{proposition}
\label{geometric-invariant-theory-proposition-fibers-of-pi-are-orbits}
The set-theoretic fibers of $\pi$ are the
orbits $Gx$
\end{proposition}

\begin{proof}
$\pi$ is $G$-invariant, i.e. constant on
orbits.
\end{proof}

\begin{proposition}
\label{geometric-invariant-theory-proposition-2-second}
The quotient map $\pi$ is open and if
$U\subset X$ is open, affine, $G$-stable,
then the restriction $\pi|_{U}$ is the
quotient map $U\to U/G$.
\end{proposition}

\begin{proposition}
\label{geometric-invariant-theory-git-label-001}
$f:X\to Z$ $G$-invariant morphism, then there
is a unique morphism $\varphi:Y
\to Z$ such that $f=\varphi \circ\pi$. This
says that $\pi$ is the categorical
quotient.
\end{proposition}

\begin{proof}[Sketch of proof]
It is clear that $\varphi$ exists as a map of
sets. It is continuous since $\pi$ is open
surjective. $\varphi^\sharp:\mathcal{O}_Z \to
\varphi_*(\mathcal{O}_Y)$.
\end{proof}

\begin{example}
\label{geometric-invariant-theory-example-quotients}
\begin{enumerate}
\item $G=\mu_n$ the roots of unity of order
$n$ (prime to $p$). $G$ acts on
$\mathbb{A}^2$ by $g(x,y)=(gx,gy)$.
$\mathcal{O}(\mathbb{A}^2)=k[x,y]$ and
$\mathcal{O}(\mathbb{A}^2)=k[x^n,x^{n-1}y,
\ldots,xy^{n-1},y^n]$.
$\mathbb{A}^2/G$ is a singular surface.
\item $G=S_n$ symmetric group, acts on
$\mathbb{A}^n$ permuting the coordinates.
$\mathcal{O}(\mathbb{A}^n)=k[x_1,\ldots,x_n]
$,
$\mathcal{O}(\mathbb{A}^2)^G=k[e_1,\ldots,
e_n]
$
where $e_i$ is the $i$-th symmetric function.
The quotient
$\pi:\mathbb{A}^n\xrightarrow{(e_1,\ldots,
e_n)
}\mathbb{A}^n$
is finite, flat of degree $n!$.
\item (Symmetric products.) $X=Z^n$, $Z$
affine algebraic variety, $G=S_n$
acting by permuting copies of $Z$. Then
$X/G:=Z^{(n)}$, the $n$-th symmetric product
of $Z$. Points of $Z^{(n)}=$effective
$0$-cycles in $Z$ of degree $n$ =
$z_1+\ldots+z_n$, $(z_i \in Z)$ =
$n_1z_1+\ldots+n_kz_k$ where $z_i$ are
distinct and $n_i$ are positive integers with
some $n$.
\end{enumerate}
\end{example}

\section{Local structure of quotients}
\label{geometric-invariant-theory-section-local-structure-of-quotients}

See \cite{AVar}.

\begin{definition}
\label{geometric-invariant-theory-definition-free-locus}
Points with trivial stabilizer.
\end{definition}

The local structure of  the symmetric product
variety $Z^{(n)}$ at the point $x$ is the
stabilizer
$$
G_x=S_{n_1}\times\ldots\times S_{n_2}\subset
S_n=G
$$
and we get
$$
Z^{(n_1)}\times\ldots\times Z^{n_r} = X/G_x
\to Z^{(n)}=X/G
$$
which is \'etale at the image of $x$.

\begin{proposition}
\label{geometric-invariant-theory-git-label-002}
If $Z$ is smooth then $Z^{(n)}$ is smooth.
\end{proposition}

\begin{proof}
Using power series rings.
\end{proof}

\begin{proposition}
\label{geometric-invariant-theory-git-label-003}
If $\dim Z \geq 2$ then $Z^{(n)}$ is
singular.
\end{proposition}

\begin{proof}
Done in lecture.
\end{proof}

\medskip\noindent
Then we have the {\it Hilbert scheme
characteristic morphism}
$$
\begin{aligned}
\gamma: \text{Hilb}_n(Z) &\longrightarrow
Z^{(n)} \\
W\subset Z &\longmapsto \sum_{z \in Z}\dim
(\mathcal{O}_{W,z}z)
\end{aligned}
$$
where $W$ must be {\bf finite of length }$n$.
See Binger, Linear determinants (older) or J.
Bertin, The punctual Hilbert Scheme.

If $Z$ is smooth of dimension 1 then $\gamma$
is an isomorphism. If $Z$ is of dimension 2
then $\gamma$ is a resolution. In general
$\gamma$ is an isomorphism above
$Z_{f_2}^{(n)}$.

\section{Linear algebraic groups}
\label{geometric-invariant-theory-section-linear-algebraic-groups}
\begin{definition}
\label{geometric-invariant-theory-definition-algebraic-group}
An {\it algebraic group} is a variety $G$
equipped with morphisms
$$
\begin{aligned}
m: G\times G &\longrightarrow G \\
(g,h) &\longmapsto gh
\end{aligned}
$$
and
$$
\begin{aligned}
i: G &\longrightarrow G \\
g &\longmapsto g^{-1}
\end{aligned}
$$
with an element $e \in G$ which satisfies the
group axioms.
\end{definition}

(In fact this is the same as a group scheme
of finite type over $k$.)

\begin{example}
\label{geometric-invariant-theory-example-algebraic-groups}
\begin{enumerate}
\item Finite groups.
\item $\mathbb{G}_a=$additive group =
$(\mathbb{A}^1,+)$ and $\mathbb{G}_m$
=multiplicative group =
$(\mathbb{A}^1\setminus\{0\},\times)$.
\item $\text{GL}_n$ = general linear group.
This is a principal open subset
($\det \neq 0$), an affine variety.
\item Classical groups
$\text{SL}_n,\text{O}_n,\text{Sp}_{2n}$
closed in
$\text{GL}_n$. Also
$\text{PGL}_n=\text{GL}_n/\mathbb{G}
_n:I_n\hookrightarrow\text{GL}_{n^2}$.
\item $(E,0)$ elliptic curve; there is a
unique algebraic group law + with
neutral element $0$. (Abelian varieties.)
\end{enumerate}
\end{example}

\begin{proposition}
\label{geometric-invariant-theory-git-label-004}
$G$ algebraic group, $G^0$ connected
component of $e$: $G^0$ is a closed normal
subgroup of $G$ and the connected components
of $G$ are the $gG^0$ for $g \in G$.
Moreover, $G^0$ is irreducible and $G$ is
smooth (which is obvious because it is smooth
somewhere and we may just translate).
\end{proposition}
$$
\xymatrix{
1\ar[r]&G^0\ar[r]&G\ar[r]&\pi_0(G)\ar[r]&1
}
$$
\begin{definition}
\label{geometric-invariant-theory-definition-action-of-algebraic-group}
An {\it action} of an algebraic group $G$ on
a variety $X$ is a morphism $a:G
\times X \to X$ which satisfies the axioms of
an action:
$$
a(g,x)=g\cdot x,\qquad  g\cdot(h\cdot
x)=gh\cdot x,\qquad  e\cdot x=x
$$
The {\it orbit} of $x \in X$ is $G\cdot
x=\{g\cdot x|g \in G\}$. The {\it stabilizer}
$G_x$ is $\{g \in G|g\cdot x=x\}$, which is a
closed subgroup of $G$.
\end{definition}

\begin{proposition}
\label{geometric-invariant-theory-proposition-properties-of-orbit}
\begin{enumerate}
\item Each orgbit $G\cdot x$ is locally
closed, smooth of dimension
$\dim(G)-\dim(G_x)$.
\item The orbit map $G \to G\cdot x$, $g
\mapsto  g\cdot x$ is faithfully flat.
\item $\overline{G\cdot x}\setminus G\cdot x$
is a union of orbits of smaller
dimension.
\item Closed orbits exist.
\end{enumerate}
\end{proposition}

\begin{proof}
Should be easy. For example, to see it's
faithfully flat you just notice it is flat
somewhere (if it is connected, is it
connected?) and then use homogeneity again.
\end{proof}

\begin{proposition}
\label{geometric-invariant-theory-git-label-005}
For any $n \geq 1$, the set $\{x \in
X|\dim(G\cdot x)\geq n\}$ is open. That is,
$\{x \in X|\dim(G_x)\leq n\}$ is open.

In particular, the points with finite
stabilizer form an open subset.
\end{proposition}

\begin{proof}
Use
$$
\begin{aligned}
\gamma: G \times X &\longrightarrow X\times X
\\
(g,x) &\longmapsto (g\cdot x,x)
\end{aligned}
$$
\end{proof}

\begin{lemma}
\label{geometric-invariant-theory-git-label-019}
Let $f:G \to H$ be an homomorphism of
algebraic groups. Then $\Ker (f)$ is a normal
subgroup of $G$. $\text{Im}(f)$ is a closed
subgroup of $H$ of dimension $\dim(G)-\dim
\Ker f$.
\end{lemma}

\begin{example}
\label{geometric-invariant-theory-git-label-006}
$G=\text{SL}_2$ acts on $k[x,y]_n:=V_n$
(these are called the {\it binary forms of
degree $n$}) by linear change of variables.
\end{example}

\begin{definition}
\label{geometric-invariant-theory-definition-linear-algebraic-group}
An algebraic group is {\it linear} if $G
\hookrightarrow \text{GL}_n$ for some $n$.
\end{definition}

\begin{example}
\label{geometric-invariant-theory-git-label-020}
$\mathbb{G}_a=\left\{\begin{pmatrix} 1&t\\
0&q
\end{pmatrix}\right\}$ is linear.
\end{example}

\begin{proposition}
\label{geometric-invariant-theory-proposition-co-action}
$G$ linear algebraic group,  $X$ a
$G$-variety, $G$ acts on
$\mathcal{O}(X)=\Gamma(X,\mathcal{O}_X)$ via
$(g\cdot f)(x)=f(g^{-1}\cdot x)$ then
$\mathcal{O}(X)$ is a union of f.d.
$G$-stable subspaces on which $G$ acts
algebraically.
\end{proposition}

\begin{proof}
From the action $a:G \times X \to X$ we get
the {\it co-action}
$$
\begin{aligned}
a^\sharp: \mathcal{O}(X) &\longrightarrow
\mathcal{O}(G \times X) \\
f &\longmapsto \sum_{i=1}^n(\varphi_i \otimes
\psi_i)
\end{aligned}
$$
\end{proof}

\begin{definition}
\label{geometric-invariant-theory-definition-G-module}
A finite dimensional vector space $V$ is
called a {\it $G$-module} if $V$ is equipped
with an action of $G$ via a homomorphism of
algebraic groups $G \to \text{GL}(V)$.

A  {\it $G$-module} is a vector space $V$
equipped with a linear action of the group
$G$ such that $V=$ union of finite
dimensional $G$-submodules.
\end{definition}

\begin{proposition}
\label{geometric-invariant-theory-proposition-G-modules}
\begin{enumerate}
\item $X$ affine $G$-variety, $G$ linear
algebraic group, then $X
\hookrightarrow  V$ finite dimensional
$G$-module as a closed $G$-stable
subvariety.
\item Every affine algebraic group is linear.
\end{enumerate}
\end{proposition}

Let $X$ be an affine $G$-variety. Then
$$
\begin{aligned}
\mathcal{O}(X) &\longrightarrow
\mathcal{O}(G\times X)
=\mathcal{O}(X)[t,t^{-1}] \\
f &\longmapsto \sum_{n \in \mathbb{Z}}f_n t^n
\end{aligned}
$$

In particular, every $\mathbb{G}_m$-module is
semisimple: every nonzero $G$-module contains
a nonzero fixed point. Indeed, let $V$ be a
finite-dimensional $G$-module. Let $v \in V$
be a $\mathbb{G}_a$-eigenvector. Then $g\cdot
v=\gamma(g)v$ for ? $X \in
\mathcal{O}(\mathbb{G}_a)^{\times}=k^{\times}
$
taking $g=0$ yields $X=1$.

$G$ connected algebraic group, then $G$ is
{\it unipotent} if it is an iterated
extension of copies of $\mathbb{G}_a$. Then
every nonzero $G$-module contains a nonzero
fixed point.

As an example consider $G=\text{SL}_2$. (We
saw that it is not linearly reductive in the
sense described in the next section.)

\section{Quotients by linearly reductive
groups}
\label{geometric-invariant-theory-git-label-007}

\begin{definition}
\label{geometric-invariant-theory-definition-linearly-reductive-group}
A linear algebraic group $G$ is {\it linearly
reductive} if every $G$-module is
semi-simple.
\end{definition}

\begin{exercise}
\label{geometric-invariant-theory-git-label-008}
$G$ is linearly reductive if and only if the
functor $(\text{$G$-modules})\to(\text{vector
spaces})$, $V \mapsto V^G$ is exact.
\end{exercise}

\begin{theorem}
\label{geometric-invariant-theory-git-label-009}
If $G$ is linearly reductive then every
connected normal unipotent subgroup $U$ of
$G$ is trivial. The converse is true in
characteristic zero.
\end{theorem}

\begin{proof}
($\implies$) Let $V$ be a $G$-module, $V\neq
0$ and $U$ as in the statement. Then $V^U\neq
0$ is a $G$-module. Let $W$ be a $G$-module.
Then $V=V^U \oplus W$. [Missing]

($\impliedby$) We show this for $k=
\mathbb{C}$. View $G \supset K$ compact
subgroup, Zariski dense (eg. $C^* \supset
S^1$ or $\text{GL}_n(\mathbb{C})\supset
\text{U}_n$.) [Missing]
\end{proof}

\begin{definition}
\label{geometric-invariant-theory-git-label-010}
A linear algebraic group $G$ is {\it
reductive} if it has a nontrivial connected
unipotent normal subgroup.
\end{definition}

The following definition is from
StacksProject, \texttt{brauer.tex}:

\begin{definition}
\label{geometric-invariant-theory-definition-simple}
Let $A$ be a $k$-algebra. We say an
$A$-module $M$ is {\it simple} if it is
nonzero and the only $A$-submodules are $0$
and $M$. We say $A$ is {\it simple} if the
only two-sided ideals of $A$ are $0$ and $A$.
\end{definition}

\begin{lemma}
\label{geometric-invariant-theory-lemma-Reynolds-operator}
$G$ linearly reductive, $V$ a $G$-module.
Then
\begin{enumerate}
\item There is a unique projection of
$G$-modules $R_V:V \to V^G$ called the
{\it Reynolds operator}.
\item
$$
\xymatrix{
V\ar[d]_{R_V}\ar[r]^f&  W\ar[d]_{R_V}\\
V^G\ar[r]^{f^G}&W^G
}
$$
commutes.
\item $A$ a $G$-algebra, $B:=A^G$. Then
$R_A(ab)=bR_A(a)$ for all $a \in A$, $b
\in B$.
\end{enumerate}
\end{lemma}

\begin{proof}
\begin{enumerate}
\item There is a unique decomposition $V=V^G
\oplus V'$ where $V'$ is the sum of
nontrivial simple $G$-submodules.
\item Implied by the decomposition $V=V^G
\oplus V'$.
\item
$$
\xymatrix{
V\ar[d]_{R_A}\ar[r]^b&  W\ar[d]_{R_A}\\
V^G\ar[r]^{b}&W^G
}
$$
commutes.
\end{enumerate}
\end{proof}

\begin{theorem}[Hilbert,Nagata]
\label{geometric-invariant-theory-git-label-011}
If $G$ is a linearly reductive group acting
on an affine variety $X$, then
$\mathcal{O}(X)^G$ is finitely generated.
\end{theorem}

\begin{proof}
Done in class, uses $R_A$ and graded Nakayama
lemma.
\end{proof}

\medskip\noindent
$\mathcal{O}(X)^G \subset \mathcal{O}(X)$
yields $\pi:X\to Y$, a $G$-invariant morphism
of affine varieties, giving a good quotient.

\begin{proposition}
\label{geometric-invariant-theory-proposition-projection}
\begin{enumerate}
\item $\pi$ is regular.
\item If  $Z \subset X$ is closed and
$G$-stable, then $\pi|_{Z}$ if the good
quotient of $Z$.
\item If $Z,Z'\subset X$ are closed
$G$-stable and disjoint, then $\pi(Z)$,
$\pi_i(Z')$ are disjoint.
\item Every fiber of $\pi$ contains a unique
closed $\mathbb{G}$-orbit.
\end{enumerate}
But if $U\subset X$ is open affine
$G$-stable, then $\pi|_{U}$ is not the good
quotient.
\end{proposition}

\begin{example}
\label{geometric-invariant-theory-example-Hopf-surface}
$G=\mathbb{G}_n$ acts on $\mathbb{A}^n$ by
scalar, i.e.
$g(x_1,\ldots,x_n)=(gx_1,\ldots,gx_n)$. Then
$\mathcal{O}(\mathbb{A}^n)^G=k$,
$\pi:\mathbb{A}^n \to \bullet$, but
$f_nU=D(f)$, where $f$ is a nonzero linear
form,
$$
\mathcal{O}(U)^{\mathbb{G}_m}=\bigcup_{d \geq
0}\frac{k[x_1,\ldots,x_n]}{f^d}
$$
\end{example}

{\bf Notation:} $Y=X/\!/G$, the space of
closed orbits.

\begin{definition}
\label{geometric-invariant-theory-definition-stable-point}
A point $x \in X$ is called {\it stable} if
its orbit $G\cdot x \subset X$ is closed and
the stabilizer  $G_x$ is finite.
\end{definition}

This definition is different from that in
Mumford's book on GIT (our definition there
is called {\it properly} stable point.)

\begin{proposition}
\label{geometric-invariant-theory-proposition-stable-points}
The set of stable points $X^s$ is open,
$G$-stable and $\pi(X^s)=Y^s$ is open.
Moreover, $X^s=\pi^{-1}(Y^s)$ and $X^s \to
Y^s$ is a geometric quotient.
\end{proposition}

Warning: there might not be any stable points
at all! (But it looks like there are some
ways to get around this.)

\begin{proof}
The first trick very elementary; just playing
with orbits. Let $x_0 \in X$ be stable. Then
$$
G\cdot x_0 \subset \pi^{-1}\pi(x_0)=\{ x \in
X
: \overline{G\cdot x}\supset G\cdot
x\}=G\cdot x_0
$$
since $\dim (G\cdot x_0)=\dim G$.

Then, roughly, remove higher dimensional
orbits, and since the remaining orbits have
all the same dimension, then they are all
closed.

The final assertion $X^s \to Y^s$ is similar
to what we did. No deep arguments.
\end{proof}

\begin{example}
\label{geometric-invariant-theory-example-action-of-SL2-on-forms}
$G=\text{SL}_n$ acts on
$X=k[x_1,\ldots,x_n]_d$ by linear change of
variables. Let's consider some values. If
$d=1$, then the quotient is
$X/\!/G=\text{pt}$. For $d=2$, the quadractic
forms, we have the discriminant $\Delta$ as
an invariant, and in fact
$\mathcal{O}(\Delta)=k[\Delta]$ {\bf (A nice
exercise!)}. Actually it looks like
$\pi=\Delta:X \to \mathbb{A}^1$, and the
stabilizer of $f$ is $\text{SO}(f)$, no
stable points.

For $d \geq 3$, and $n \geq 5$, then
$\mathcal{O}(X)^G$ is not known! But we also
have the discriminant $\Delta \in
\mathcal{O}(X)^G$ and in fact (Jordan) if
$\Delta(f) \neq 0$ then $f$ is stable. Let us
prove this fact (we don't know how did Jordan
proved this): we show that $G_f$ is finit and
stable, i.e. $\text{Lie}(G_f)=0$. The Lie
algebra of $\text{SL}_2$, the traceless
matrices, acts on $X$ by
$$
A\cdot f= \sum_{i=1}^n \ell_i \frac{\partial
f}{\partial x_i}
$$
where $\ell_i = \sum_{j=1}^n a_{ij}x_j$. Then
$A\cdot f=0$. Now $\Delta(f) \neq 0$ iff
$\frac{\partial f}{\partial x_j}$ have no
nontrivial common zero, iff, $\frac{\partial
f}{\partial x_j}$ form a regular sequence.
This says that the ring Cohen-Macaulay. Then
$$
\ell_i \frac{\partial f}{\partial x_i}
=\sum_{i \neq  j}\ell_j \frac{\partial
f}{\partial x_j}
$$
which in turn implies that
$$
\ell_i \in \left(\frac{\partial f}{\partial
x_j}, \; j \neq  i\right)
\iff \ell_i=0 \forall  i
$$
\ldots{} which says {\it what?}
\end{example}

\medskip\noindent
Let $X$ be a projective $G$-variety, $G$
linearly reductive, $L \to X$ an ample line
bundle
$$
R(X,L)=\bigoplus_{n=0}^{\infty}
\Gamma(X,L^{\otimes n})
$$
which is a graded, finitely generated
algebra, then $X=\text{Proj}R(X,L)$. If $G$
acts on $R(X,L)$ ``compatibly'', then
$R(X,L)^G$ is finitely generated:
$$
\text{Proj}R(X,L)^G=X/\!/G
$$
and we get a birational map
$$
\xymatrix{
X\ar@{.>}[r]& X/\!/G\\
U\ar@{^{(}->}[u]\ar[ur]^{\pi}
}
$$
where $U$ is the domain of definition.

If $G$ is finite, then  $X= \mathbb{P}(V)$,
$V$ finite dimensional $G$-module. $L =
\mathcal{O}_X(1)$: $G$ acts on $L$ and hence
on $R(X,L)$ (``{\it I'll explain details
soon}'')

We have a finite ring extension $R(X,L)^G =
R(X,L)$, and we get
$$
X=\text{Proj}R(X,L) \to \text{Proj}R(X,L)^G=Y
$$

\medskip\noindent
Now let $G$ be an algebraic group, $X$ a
$G$-variety, $f:L \to X$ a line bundle.

\begin{definition}
\label{geometric-invariant-theory-definition-G-linearization}
A  {\it $G$-liearization} of $L$ is a
$G$-action on $L$ such that $f$ is
equivariant and $g:L_x \to L_{g\cdot x}$ is
linear for all $g \in G, x \in X$. That is,
the $G$-action on $L$ commutes with the
$\mathbb{G}_m$-action by multiplication on
fibers.
\end{definition}

Then we get an action on the section ring:
$G$ acts on $L^{\times}=L\setminus
\{\text{zero section}\}$ and hence
$\mathcal{O}(L^{\times})$ is a $G$-module. As
we said before,
$\mathcal{O}(L^{\times})=\bigoplus_{n \in
\mathbb{Z}}\Gamma(X,L^{\otimes n})$ is a
graded module, and the grading corresponds to
the $G$-action! That is, this is a $G$-module
structure on $\Gamma(X,L^{\otimes n})$. And
$R(X,L)$ is a graded $G$-algebra.

\begin{definition}
\label{geometric-invariant-theory-definition-semi-stable-line-bundle}
$X$ a $G$-variety, $L$ a $G$-linearized line
bundle on $X$ and $x \in X$. We say $x$ is
{\it semi-stable} w.r.t. $L$ if there exists
$n \geq 1$ and $\sigma \in
\Gamma(X,L^{\otimes n})^G$ and that $x \in
X_\sigma$ and $X_\sigma$ is affine.

$x$ is  {\it stable} if in addition $G\cdot
x$ is closed in $X_\sigma$ and $G_x$ is
finite.
\end{definition}

\begin{theorem}[Mumford]
\label{geometric-invariant-theory-theorem-Mumford}
The set of semis-stable points
$X^{ss}(L)\subset X$ is open, $G$-stable and
has a good quotient $\pi:X^{ss}(L) \to Y(L)$.
(Notice everything depends on the line
bundle.) The set of stable poins $X^s(L)$ is
open $G$-stable and equal to
$\pi^{-1}\pi(X^s(L))$ and $\pi|_{X^s(L)}$ is
a geometric quotient.
\end{theorem}

\begin{proposition}
\label{geometric-invariant-theory-proposition-on-projectivization-of-V}
Let $x \in \mathbb{P}(V)$ and $x \in [v]$.
\begin{enumerate}
\item $x \in \mathbb{P}(V)^{ss}$ $\iff$ $0
\not \in \overline{G\cdot v}$.
\item $x\in \mathbb{P}(V)^s$ $\iff$ $v \in
V^s$ i.e. $G\cdot v$ is closed in $V$
and $G_v$ is finite.
\end{enumerate}
\end{proposition}

\begin{proof}
Done in class.
\end{proof}

The following is our first example of moduli
space.

\begin{example}[Smooth hypersurfaces of
degree $d$ in $\mathbb{P}^n$]
\label{geometric-invariant-theory-example-smooth-hypersurfaces}
The isomorphism classes of such surfaces are
given by: take degree $d$ polynomials,
projectivize, take away diagonal ( this gives
an open set) and finally quoetient by
$\text{SL}_{n+1}$:
$$
\begin{aligned}
\{\text{isom.
classes}\}&=\mathbb{P}(k[x_0,\ldots,x_n]_d)
_{\Delta}
/\text{SL}_{n+1}\overset{\text{open}}
{\hookrightarrow
}\\
\mathbb{P}(\cdot)^1/\text{SL}_{n+1}&
\hookrightarrow
\mathbb{P}(\cdot)^n/\!/\text{SL}_{n+1}
\end{aligned}
$$
a projective variety, normal, very singular.
\end{example}

\begin{remark}
\label{geometric-invariant-theory-git-label-012}
$G \mathbb{y} X$ categorical quotient $\pi:X
\to Y$, then $\pi$ yields a coarse moduli
space for
$$
\xymatrix{
X\ar[r]\ar[rd]& [X/G]\ar@{.>}[d]\\
&Y
}
$$
\end{remark}

\begin{proof}
$\text{Hom}(S,[X/G])$ is a category, not a
set as usual. Its objects are
$$
\xymatrix{
P\ar[d]_{\substack{\text{principal} \\
\mathbb{G}\text{-bundle}}}
\ar[rr]^{\substack{G\text{-equivariant} \\
\text{morphism }b}}&&X\\
S
}
$$
and morphisms are
$$
\xymatrix{
&X\\
P\ar[ur]^{b'}\ar[rr]^\varphi \ar[dr]&  &
P\ar[ul]_b\ar[dl]\\
&S
}
$$
(continues\ldots{})
\end{proof}

\section{Hilbert-Mumford criterion}
\label{geometric-invariant-theory-section-Hilbert-Mumford-criterion}

\begin{definition}
\label{geometric-invariant-theory-definition-one-parameter-subgroup}
A {\it one-parameter subgroup} of  $G$ is a
homomorphism $\lambda:\mathbb{G}_m \to G$. If
$G \mathbb{y} X$ and $x \in X$, we say that
$$
\lim_{t \to 0}\lambda(t)
$$
exists and equals $y \in X$ if
$\mathbb{A}^1\setminus\{0\}=\mathbb{G}_m \to
X$, $t \mapsto \lambda(t)\cdot x$ extends to
$\mathbb{A}^1 \to X$, $0 \mapsto y$.
\end{definition}

\begin{theorem}[Hilbert-Mumford criterion]
\label{geometric-invariant-theory-theorem-Hilbert-Mumford-criterion}
$V$ finite-dimensional $G$-module, $G$
linearly reductive, $v \in V$,
$\overline{G\cdot v}\supset G\cdot v$ closed
orbit. There exists a one-paremeter subgroup
$\lambda$ such that $\lim_{t \to 0}
\lambda(t)\cdot v \in G\cdot w$.
\end{theorem}

We will not prove this theorem. Here are some
corollaries.

\begin{lemma}
\label{geometric-invariant-theory-lemma-unstable}
$v \not \in V^{ss} \iff \exists  \lambda$
such that $\lim_{t\to_0}\lambda(t)\cdot v
=0$. Then $v$ is called {\it unstable} and
$$
\{\text{unstable points}\}=\text{{\it
nilcone}} \mathcal{N}(V)=V.
$$
\end{lemma}

\begin{lemma}
\label{geometric-invariant-theory-git-label-013}
$G\cdot v$ is closed $\iff$ $\lim_{t \to 0}
\lambda(t)\cdot v$ exists, then it lies in
$G\cdot v$.
\end{lemma}

\begin{proposition}
\label{geometric-invariant-theory-git-label-014}
$G\cdot v$ is affine $\iff$ $G_v$ is linearly
reductive.
\end{proposition}

\begin{proof}
($\implies$) $G_v = H \iff G\cdot v
G/H=G/\!/H$.

($\impliedby$) We show that the functor
$$
\begin{aligned}
(\text{left $H$-modules}) &\longrightarrow
(\text{f.d. vector spces}) \\
V &\longmapsto V^H
\end{aligned}
$$
(continues\ldots{})
\end{proof}

We can characterize stable points using this
proposition.

\begin{lemma}
\label{geometric-invariant-theory-git-label-015}
$v \in V^s$ $\iff$ $G\cdot v$ is closed and
$v$ is fixed by no nontrivial one-parameter
subgroup.
\end{lemma}

\begin{proof}
The reason is that if $G_v=H$ is linearly
reductive then $H^0$ is generated by images
of a one-parameter subgroup.
\end{proof}

\begin{example}[Stable points for $SL_2$
acting on degree $d$ polynomials]
\label{geometric-invariant-theory-git-label-016}
If $G=\text{SL}_2$ in characterstic zero,
then
$$
\lambda(t)=\begin{pmatrix}
t&0\\
0&t^{-1}
\end{pmatrix}
$$
every nontrivial one-parameter subgroup of
$G$ is conjugate to $\lambda(t)^n$ for a
unique $n \geq 1$. $V=V_d=k[x,y]_d\ni x^e
y^{d-e}$.
$$
\lambda(t) x^e y^{d-e}=t^{d-2e}x^ey^{d-e}
$$
$f \in V$ is unstable $\iff$ $\exists g \in
G$ such that $gb \in \left\langle{}
x^{d/2-t}y^{d/2+1},\ldots, y^d\right\rangle{}$ if
dimension is even, and $
\left\langle{}x^{(d-1)/2}y^{(d+1)/2},\ldots,
y^d\right\rangle{}$ if dimension is odd, $\iff$
$f$ has no root of multiplicity $d/2$. $f
\not \in V^s$ $\iff$ $f$ has no root of
multiplicity $\geq d/2$.

If $d$ is odd then $V^s=V^{ss}$. If $d$ is
even then $x^{d/2}y^{d/2}$ has a closed orbit
and is fixed by $\\text{Im}(\lambda)$.
\end{example}

\section{Hyperelliptic curves}
\label{geometric-invariant-theory-section-hyperelliptic-curves}

Let $C \to \mathbb{P}^1$ be a smooth
projective curve of degree $2$, genus $g$,
branched in $2g+2$ distinct points, uniquely
determined by the branch divisor.
\begin{gather*}
\{\text{isomorphism classes of hyperellipetic
curves of genus $g$}\}\\
=\mathbb{P}(V_{2g+2})_\Delta/\text{SL}_2
\hookrightarrow
\mathbb{P}(V_{2g+2})^{ss}/\!/\text{SL}_2
\end{gather*}
projective normal of dimension $2g-3$.

\section{Local structure of quotients}
\label{geometric-invariant-theory-section-local-structure-of-quotients-duplicate-2}
$X$ affine $G$-variety, $G$ reductive, $x\in
X$ such that $G\cdot x$ is closed, $H=G_x$
(reductive).

\begin{theorem}
\label{geometric-invariant-theory-theorem}
There exists an affine $H$-stable (locally ?)
subvariety $S \subset X$ such that the
sequence
$$
\xymatrix{
(g,s)H\ar@{|->}[r]&g\cdot s\\
(G\times S)/H\ar[d]\ar[r]&  X\ar[d]\\
S/\!/H\ar[r]&  X//G\\
x\ar@{|->}[r]&\overline{x}
}
$$
$h(g,s)=(gh^{-1},hs)$ is cartesian with
\'etale
horizontal arrows. If $x \in X$ is smooth,
then there is an $H$-equivariant morphism.
$$
\begin{aligned}
S &\longrightarrow T_x(X)/T_x(G\cdot
x)=N_{G\cdot x/X,x} \\
x &\longmapsto 0
\end{aligned}
$$
\end{theorem}

This finishes the statement of the theorem.
The last map may be interpreted as the
inverse of the exponential map from
differential geometry.

A corollary:

\begin{lemma}
\label{geometric-invariant-theory-lemma}
If $G_x=e$ then $X \to X/\!/G$ is a principal
$G$-bundle in a $G$-stable neighbourhood of
$x$.
\end{lemma}

\begin{exercise}
\label{geometric-invariant-theory-git-label-017}
$G$ finite, works with $S\subset X$ open.
\end{exercise}

\begin{exercise}
\label{geometric-invariant-theory-exercise}
$G=\text{SL}_2$, $V=V_G$, $x^3y^3 \in V_G$,
stabilizer
$=\mathbb{G}_m=\text{Im}(\lambda)$. Prove
that $V_G/\!/\text{SL}_2$ is singular at the
image of $x^3y^3$. ($S=\mathbb{G}_m$-module
with weights $-6,-4,0,4,6$.) {\bf Hint.} To
check for non-smoothness you can check that
$\mathcal{O}(S)^{\mathbb{G}_m}$ is not a
polynomial ring; this is equivalent to
non-smoothness for the quotient, and it
should be easy to find: you just find that
there are too many generators.
\end{exercise}


% --- End geometric-invariant-theory.tex ---

% --- Begin moduli-spaces-of-sheaves.tex ---
\chapter{Moduli Spaces of Sheaves}
\label{chapter-moduli-spaces-of-sheaves}



Minicourse by Cristina Manolache, CIMPA
school Florian\'opolis 2025.

Notes at
\href{%
http://github.com/danimalabares/cimpa-floripa
}{%
github.com/danimalabares/cimpa-floripa
}

\bigskip\noindent

{\bf Abstract.} This course will be an
introduction to Moduli Spaces of vector
bundles. A moduli space of stable vector
bundles on a smooth, algebraic variety X is a
scheme whose points are in ``natural
bijection'' to isomorphic classes of stable
vector bundles on X. Using Geometric
Invariant Theory the moduli space can be
constructed as a quotient of certain
Quot-scheme by a natural group action. We
introduce the crucial concept of stability of
vector bundles over smooth projective
varieties and we give a cohomological
characterization of the (semi)stability. The
notion of (semi)stability is needed to ensure
that the set of vector bundles one wants to
parameterize is small enough to be
parameterized by a scheme of fniite type. We
introduce the formal defniition of moduli
functor, fnie moduli space and coarse moduli
space and we recall some generalities on
moduli spaces of vector bundles. Then we
focus on vector bundles on algebraic
surfaces. Quite a lot is known in this case
and we will review the main results, some of
which will illustrate how the geometry of the
surface is refelcted in the geometry of the
moduli space. Going beyond surfaces, we
introduce the notion of monad, which allow
the classifciation of vector bundles on, for
instance, $\mathbb{P}^3$. Monads appeared in
a wide variety of contexts within algebraic
geometry, and they are very useful when we
want to construct vector bundles with
prescribed invariants like rank, determinant,
Chern classes, etc.  Finally, we will study
moduli spaces of vector bundles on higher
dimensional varieties. As we will stress, the
situation drastically difefrs and results
like the smoothness and irreducibility of
moduli spaces of stable vector bundles on
algebraic surfaces turn to be false for
moduli spaces of stable vector bundles on
higher dimensional algebraic varieties.

\bigskip\noindent

\bigskip\noindent

\section{Introduction}
\label{moduli-spaces-of-sheaves-section-introduction}

The goal is the following. $X$ smooth
projective variety over $\mathbb{C}$. We want
to define a moduli space of sheaves up to
some kind of isomorphism. Also give some
discrete invariants. And give structure of
scheme which ``respects limits''.

{\bf Functor of points.} For a given scheme
$X$, the set
$\text{Hom}(\text{Spec}\mathbb{C},X)$ is the
set of points of $X$. This motivates defining
a functor $h_X:\text{Sch}\to \text{Sets}$ by
$h_X(S)=\text{Hom}(S,X)$, and morphisms are
just mapped via composition.

And there's more: given a morphism of schemes
$\varphi:X\to Y$ we can produce
$h_\varphi:h_X \to h_Y$.

\begin{slogan}
Schemes are functors.
\end{slogan}

\begin{definition}
\label{moduli-spaces-of-sheaves-msos-label-001}
A functor $F$ is called {\it representable
(by a scheme)} if there is a scheme $X$ such
that $F\cong h_X$ and we say $X$ is a {\it
fine moduli space for $X$}.
\end{definition}

\begin{remark}
\label{moduli-spaces-of-sheaves-msos-label-002}
If $X$ exists it is unique.
\end{remark}

\begin{example}
\label{moduli-spaces-of-sheaves-example-quotient-sheaf}
$X$ scheme,
$$
\text{Quot}_X:(\text{Sch})\to (\text{Sets})
$$
$$
\text{Quot}_X(S)=\{
\mathcal{O}_{X \times S}^{\oplus
n}\to\mathcal{F}
:\mathcal{F} \text{ is a flat sheaf}\}\Big/
\sim
$$
and morphisms are mapped by pullback.
\end{example}

Here, ``flat'' accounts for a continuously
varying family of copies of $X$ over $S$,
i.e. the product $X \times S$.

\begin{theorem}
\label{moduli-spaces-of-sheaves-theorem-Quot-is-represented}
$\text{Quot}_X$ is represented by a
projective scheme.
\end{theorem}

Almost all examples of moduli spaces are
quotients.

\medskip\noindent

The one example of a moduli space everybody
knows is\ldots{}

\begin{example}[Projective space is a
quotient]
\label{moduli-spaces-of-sheaves-msos-label-016}
Take $X=$ point, and then
$$
F(S)=\{\mathcal{O}_S^{\oplus n+1}\to
\mathcal{L}\}\cong h_{\mathbb{P}^n}
$$
where $\mathcal{L}$ is a line bundle on $S$.
\end{example}

\begin{definition}
\label{moduli-spaces-of-sheaves-definition-universal-family}
If $F$ is represented by a scheme $X$ and
$\eta:F(S)\to h_X(S)$ in part we have a
bijection
$$
\begin{aligned}
F(X) &\longrightarrow h_X(X):=\text{Hom}(X,X)
\\
U &\longmapsto \text{id}
\end{aligned}
$$
\end{definition}

\bigskip\noindent

Again, flat means continuously varying.
Consider first
$$
\mathcal{M}(S)=\{\mathcal{F}\text{a fheaf on
$X\times S$ flat over }S\}/\sim
$$
and
$$
\mathcal{M}(f:T \to S)=f^*
$$
and equivalence is given by: $\mathcal{F}\sim
\mathcal{G}$ iff there exists a line bundle
$\mathcal{L}$ on $S$ such that
$\mathcal{F}\cong \mathcal{G}\otimes\pi^*S$.

\medskip\noindent
{\bf Problem.} $X:=\mathbb{P}^1$ and rank 2
sheaves on $\mathbb{P}^1$. In this case
$[\ldots{}]$

\medskip\noindent
{\bf Fact.} Fine moduli spaces are rare.

\begin{definition}
\label{moduli-spaces-of-sheaves-definition-coarse-moduli-space}
A {\it coarse moduli space} of  $X$ is a
scheme $X$ with a natural transformation
$\eta:F \to h_X$ such that
\begin{itemize}
\item $\eta_{\eta \times \mathbb{C}}:F(hec
\mathbb{C}) \to h_X(hec \mathbb{C}$
is hij.
\item For all $Y$ and $F \to h_T$ we have [$F
\to h_Y$, $F\to h_X$, $\exists !
h_X \to h_Y$.
\end{itemize}
\end{definition}

\medskip\noindent
{\bf Summary.}
\begin{enumerate}
\item Schemes are functors.
\item The best moduli spaces are represented
functors, they are called fine.
These have universal families, and all are
pullbacks of from the universal family.
\end{enumerate}

\section{Discrete data}
\label{moduli-spaces-of-sheaves-section-discrete-data}

$X$ projective and $H$ an ample divisor on
$X$. The {\it Hilbert polynomial} is
$$
P(\mathcal{F},m):=\chi(\mathcal{F}\otimes
\mathcal{O}(m))
$$

\medskip\noindent
{\bf Fact.}
$P(\mathcal{F},m)=\sum_{i=0}^{\dim
X}\alpha_i(\mathcal{F})\frac{m^i}{i!}$.

The {\it rank of $\mathcal{F}$} is
$\text{rk}(\mathcal{F}):=
\frac{\alpha_d(\mathcal{F})
}{\alpha_d(\mathcal{O}_X)}$.
And $\alpha_d(\mathcal{O}_X)$ is the {\it
degree} of $\mathfrak{X}$ w.r.t
$\mathcal{O}(1)$.

The {\it reduced Hilbert polynomial} is
$P(\mathcal{F},m):=\frac{P(\mathcal{F},m)
}{\alpha_d(\mathcal{F})}$.
The {\it slope} of $\mathcal{F}$ is
$\mu_H(\mathcal{F}):=\frac{c_1(\mathcal{F})
H^{d-1}}{\text{rk}\mathcal{F}}$.

\begin{example}[Curve of genus $g$ using
Riemann-Roch]
\label{moduli-spaces-of-sheaves-msos-label-003}
(Missing.)
\end{example}

\medskip\noindent
Here begin the contents of the second
lecture.
\begin{definition}
\label{moduli-spaces-of-sheaves-definition-Gieseker-and-mu-stability}
$\mathcal{F}$ is {\it
Gieseker/$\mu$-semistable} w.r.t. $H$ if
\begin{itemize}
\item $P_\varepsilon(\mathfrak{m})\leq
P_{\mathcal{F}}(\mathfrak{m})$ for all $0
\neq  \varepsilon \subset \mathcal{F}$ proper
subsheaf.
\item $\mu_{M}(\varepsilon) \leq
\mu_{M}(\mathcal{F})$ of rank $0<r \leq r-1$.
\end{itemize}
If the inequality is strict then
$\mathcal{F}$ is called {\it stable}.
\end{definition}

\begin{example}
\label{moduli-spaces-of-sheaves-msos-label-004}
$X$ a curve, $d=\text{deg}\mathcal{F}$,
$r=\text{rk}\mathcal{F}$,
$p=\text{deg}Hm+\frac{d}{r}+1-g$. Then
$\mu$-stability is the same as G-stability.
\end{example}
In general:
$$
\mu\text{-stability}\implies \text{
G-stability }\implies \text{G semi-stability
}\implies \mu \text{ semi-stability}
$$
\begin{remark}
\label{moduli-spaces-of-sheaves-remark-stability-depends-on-H}
Stability depends on the choice of $H$.
\end{remark}

\begin{example}[Stability depends on choice
of ample divisor]
\label{moduli-spaces-of-sheaves-msos-label-005}
$X=\mathbb{P}^1\times \mathbb{P}^1$, then
$\text{Pic}(X)
=\mathbb{Z}[\ell]\oplus\mathbb{Z}[m]$ where
$\ell^2=m^2=0$ and $\ell m=1$. Consider
$\mathcal{F}$ given by
\begin{equation}
\label{moduli-spaces-of-sheaves-equation-extensions}
\xymatrix{
0\ar[r]&\mathcal{O}_X(\ell-3m)\ar[r]&
\mathcal{F}\ar[r]&\mathcal{O}_X(3m)
\ar[r]&0
}
\end{equation}
Then $c_1(\mathcal{F})=\ell$ and
$c_2(\mathcal{F})=(\ell-3m)(3m)=3$. Then
$L=\ell+sm$; $L'=\ell+7m$. We will show that
$\mathcal{F}$ is semistable w.r.t. $L$ but
not w.r.t. $L'$. We do have non-trivial
extensions as in the sequence
\ref{moduli-spaces-of-sheaves-equation-extensions}. These are
parametrized by
$$
\text{Ext}^1(\mathcal{O}_X(3m)),
\mathcal{O}_X(-3m)\cong
H^{1}(X,\mathcal{O}_X(\ell-6m))\cong \ldots
\cong \mathbb{C}^{10}
$$
So it is semi-stable w.r.t. $L$.

Suppose $\mathcal{O}(s) \subset \mathcal{F}$.
We want to show that $\mu(\mathcal{O}(s))
\leq  \mu(\mathcal{F})$.

We claim that we have either
$\mathcal{O}(s)\hookrightarrow
\mathcal{O}(3m)$ or
$\mathcal{O}(s)\hookrightarrow
\mathcal{O}(\ell-3m)$. Claim was proved and
further computations showed the required
inequality, so $\mathcal{F}$ is
$\mu$-semistable w.r.t. $L$. On the other
side, $\mathcal{F}$ is not semistable w.r.t.
$L'$ since $\mathcal{O}(\ell-3m)
\hookrightarrow  \mathcal{F}$ (so here I
realise that probably taking Chern class will
be monotonous with respect to this
inclusion\ldots{}) which gives
$\mu_{L'}(\mathcal{O}(\ell-3m))=(\ell-3m)
(\ell+7m)=4\not
\leq \frac{7}{2}$.
\end{example}

\begin{example}[Results on $\mu$-stability]
\label{moduli-spaces-of-sheaves-example-results-on-mu-stability}
\begin{itemize}
\item All line bundles are stable.
\item If $0\to L_0 \to \mathcal{F} \to L_1
\to 0$ is a non-trivial extension
with degree $L_0=0$, $\text{deg}(L_1)=1$,
$L_i$ line bundles, then $\mathcal{F}$ is
stable.
\item $E_i$ semi-stable shows $E_1\oplus E_2$
is semistable if and only if
$\mu(E_i)=\mu(E_2)$.
\item $E$ is $\mu$-semistable if and only if
for all $L$ line bundles $E \otimes
L$ is semistable.
\item $E$ is semistable then $E^{\vee}$ is
semi-stable.
\item $E_1,E_2$ semistable then $E_1 \otimes
E_2$ is semistable.
\end{itemize}
\end{example}

Now we explain some {\bf easy criteria.} (We
shall find a condition on cohomology that is
equivalent to stability.)

\begin{definition}
\label{moduli-spaces-of-sheaves-definition-norm-sheaf}
$\mathcal{F}$ a reflexive sheaf on
$\mathbb{P}^n$ of rank $r$. We define
$\mathcal{F}_{\text{norm}}:=\mathcal{F}(k)$
where $k$ is the unique integer such that
$c_1(\mathcal{F}(k))\in \{-r+1,\ldots,c\}$.
\end{definition}

\begin{proposition}[Cohomological
characterization of stability]
\label{moduli-spaces-of-sheaves-msos-label-006}
$\mathcal{F}$ reflexive sheaf of rank  $2$,
then $\mathcal{F}$ is $\mu$-stable if and
only if
$H^{0}(\mathbb{P}^n,\mathcal{F}_{\text{norm}}
)
=0$.
\end{proposition}

\begin{proof}
$\mathcal{F}$ vector bundle of rank $r$ on a
smooth projective variety with
$\text{Pic}=\mathbb{Z}$.
\begin{itemize}
\item If
$H^{0}(X,(\Lambda^{z}\mathcal{F})
_{\text{norm}
})=0$
for all $z \in
\{1,\ldots,r-1\}$ then $\mathcal{F}$ is
$\mu$-stable.
\item If
$H^{0}(X,\Lambda^{z}\mathcal{F})_{\text{norm}
}
(-1))=0$
for all $z \in
\{1,\ldots,r-1\}$, then $\mathcal{F}$ is
$\mu$-semi-stable.
\end{itemize}
\end{proof}

\begin{example}
\label{moduli-spaces-of-sheaves-msos-label-007}
$\Omega_{\mathbb{P}^n}$ is stable, and also
$T_{\mathbb{P}^1}$. The cohomologies vanish
using Bott.
\end{example}

\begin{example}
\label{moduli-spaces-of-sheaves-example-lots-of-Os}

\end{example}

{\bf Filtrations.} $H$ ample, $\mathcal{F}$
torsion free sheaf, then there exists a
unique filtration $0 \hookrightarrow
\mathcal{F}_0\subset
\mathcal{F}_1\subset\ldots\subset\mathcal{F}
_m=\mathcal{F}$
with
$E_i=\mathcal{F}_i/\mathcal{F}_{i+1}$
$\mu$-semi-stable w.r.t. $H$ and
$\mu(E_1)>\mu(E_2)>\ldots>\mu(E_m)$. This is
called {\it Harder-Narorimnan filtration}.

Given any Chern classes, can we guarantee the
existence of a sheaf with these Chern
classes? No.

\begin{theorem}[Bogomolov Inequality]
\label{moduli-spaces-of-sheaves-theorem-Bogomolov-inequality}
$X$ smooth, projective variety of dimension
$\geq 2$, $H$ ample. If $\mathcal{F}$ is
torsion-free, $\mu$ semi-stable sheaf w.r.t.
all
$$
\Delta(\mathcal{F}):=(2rc_2(\mathcal{F})-
(r-1)c_1(\mathcal{F})^2)H^{n-2}\geq 0
$$
where $r=\text{rk}(\mathcal{F})$ and $\dim
X=n$. We call $\Delta$ the {\it
discriminant}.
\end{theorem}

\section{Boundedness}
\label{moduli-spaces-of-sheaves-section-boundedness}

We need further preparation before we can
define moduli spaces.

\begin{definition}
\label{moduli-spaces-of-sheaves-definition-bounded-family}
A family of isomorphism classes of coherent
sheaves on $X$ is called {\it bounded} if
there exists $S$ of {\bf finite type} and an
$\mathcal{O}_{S \times X}$ sheaf $F$ such
that the family is contained in $\{
\mathcal{F}_{S \times X}|s  \text{ closed in
}S\}$.
\end{definition}

\begin{definition}
\label{moduli-spaces-of-sheaves-definition-m-regularity}
A coherent sheaf is called {\it $m$-regular}
if $H^{i}(X,\mathcal{F}(m-i))=0$ for all
$i>0$.
\end{definition}

\begin{proposition}
\label{moduli-spaces-of-sheaves-proposition-m-regularity}
If $\mathcal{F}$ is $m$-regular,
\begin{itemize}
\item $F$ is $m'$ regular for all $m' \geq
m$.
\item $F(m)$ is globally generated.
\end{itemize}
\end{proposition}

\begin{definition}
\label{moduli-spaces-of-sheaves-definition-regular-infimum}
$\text{reg}\mathcal{F}: =\text{inf}\{m \in
\mathbb{Z}:\mathcal{F}\text{ is
$m$-regular}\}$
\end{definition}

\begin{lemma}
\label{moduli-spaces-of-sheaves-lemma}
The following are equivalent:
\begin{itemize}
\item $\{\mathcal{F}_i\}_{i \in I}$ is
bounded
\item The set of Hilbert polynomials
$\{P(\mathcal{F}_i)\}_{i \in I}$ is finite
and there exists a bound $\rho$ such that
$\text{reg}\mathcal{F}_i\leq \rho$ for all $i
\in I$.
\end{itemize}
\end{lemma}

The following is a difficult result that we
shall only state and use.

\begin{theorem}
\label{moduli-spaces-of-sheaves-theorem}
$X$ projective, $H$ ample, $P$ a fixed
polynomial of degree $d$. The family of
torsion-free sheaves with Hilbert polynomial
$P$ is bounded.
\end{theorem}

\begin{proposition}
\label{moduli-spaces-of-sheaves-msos-label-008}
Let $X$ be a curve of genus $g$ with fixed
Hilbert polynomial. The semi-stable sheaves
on $X$ are bounded.
\end{proposition}

\begin{proof}
Let $E$ be a semis-stable sheaf. We want $m$
such that $H^{1}(X,E(m-1))$. We apply Serre
Duality as follows:
$$
H^{1}(X,E(m-1))\overset{\text{S.D.}}{=
}\text{Hom}(E,\omega(1-m))^\vee
$$Then $E \to \omega(m+1)$ $\implies$
$\omega^\vee(m-1)\to E^\vee$. $E$ is
semi-stable then
$-\text{deg}(\omega^\vee)\leq
\frac{\text{deg}(E^\vee)}{r}$
which implies that $-2g+2+m-1\leq
-\frac{d}{r}$ for $d=\text{deg}E$, and then
$m\leq  2g-2-\frac{d}{r}+1$. If we take
$m>2g-2\frac{d}{r}+1$.
\end{proof}

\begin{proposition}
\label{moduli-spaces-of-sheaves-proposition}
If  $E$ is a stable sheaf, then
$\text{Hom}(E,E)=\mathbb{C}$.
\end{proposition}

\begin{proof}
If $f:E \to E$ factors through $E \to F \to
F$ and is not surjective, then
$\mu(E)<\mu(F)<\mu(E)$ $\implies$ $f$ is
injective. As an exercise, prove that
$E$ is stable if and only if for all $E \to
F$ surjective we have that
$\mu(E)<\mu(F)$.
\end{proof}

\section{Construction of the moduli space of
semi-stable sheaves}
\label{moduli-spaces-of-sheaves-msos-label-009}
$$
\mathcal{M}^p(s)=\left\{
\substack{\text{$\mathcal{F}$ on $X\times S$,
pet over $S$}
\\ \text{with Hilbert polynomial $P$}}
\right\} \Big/\sim
$$
where $\mathcal{F}\sim \mathcal{G}$ $\iff$
there exists $L$ on $S$ such that
$\mathcal{F}\otimes \pi^*L \cong \mathcal{G}$
where
$X \times S\xrightarrow{\pi}\mathcal{G}$.

$\mathcal{F}$ is a semi-stable sheaf is
$\exists $ $m$ such that
$\mathcal{F}$ is $m$-regular, which implies
that $\mathcal{F}(m)$ is globally
generated.

\medskip\noindent
{\bf Warning.} Following are several formulas
I don't understand and just copied
from the board as they were.
$$
\dim H^{0}(X,\mathcal{F}(m))=P(m)
$$
$$
H^{i}(X,\mathcal{F}(m-1))=0\qquad \forall i>0
$$
$$
V=\mathcal{O}_X^{\oplus P(m)}\to
\mathcal{F}(m)\to 0
$$
$$
\implies \left[ \mathcal{O}_X^{\oplus
P(m)}(-m) \to \mathcal{F} \to 0 \right]
\in \text{Quot}(\mathcal{H},P)
$$
Let $R$ be open locus in
$\text{Quot}(\mathcal{H},P)$ such that
\begin{itemize}
\item $\mathcal{F}$ is semi-stable.
\item $H^{0}(\mathcal{H}(m))\cong
H^{0}(\mathcal{F}(m))$.
\end{itemize}
$\text{GL}(V)\mathbb{y} R$ by $g[\rho]=\pi
\circ g$ for $\rho \in R$.
$$
\xymatrix{
\mathcal{O}^{\oplus
p(m)}\ar[r]^{\rho}\ar[d]_{g,\cong}&
\mathcal{F}\\
\mathcal{O}_X^{\oplus p(m)}\ar[ur]
}
$$
We want
``$R/\text{GL}(V)$''.

\begin{theorem}
\label{moduli-spaces-of-sheaves-msos-label-010}
There exists a coarse moduli space $M_{X,P}$
of semi-stable sheaves with fixed
Hilbert polynomial $P$, $\mathcal{M}_{X,P}$
is projective and
$\mathcal{M}_{X,P}^s$ of stable sheaves with
Hilbert polynomial $P$.
\end{theorem}

\medskip\noindent
{\bf Fact.} The GIT (semi-)stability
coincides with the Gieseker
(semi-)stability (see \cite{HL}).

\medskip\noindent
{\bf Remarks on working with stacks.}

\begin{itemize}
\item Instead of constructing sheaves which
are coarse/fine moduli spaces we may
find more general geometric objects called
``stacks''.
\item Advantage: Quotients always exist:
$G \mathbb{y} X \rightsquigarrow [X/G]$.
\item Stacks have universal families.
\item Disadvantage: we need to work with more
general moduli functors. Instead
of $\mathcal{M}:(\text{Sch}) \to
(\text{Sets})$, we substitute $(\text{Sets})$
with a 2-category.
\item Disadvantage: not all algebraic stacks
have coarse moduli spaces but if
the number of automorphisms of points [is
finite?].
\item Disadvantage: may not be projective
(even if proper/compact).
\end{itemize}

In sum, we may not be able to produce a
moduli space, and even if we can, it may
not be projective.

\medskip\noindent
Even in the easiest situation we can have, we
don't have universal family:

\begin{example}[The line bundles of degree
$d$ on $\mathbb{P}^1$ do not have
universal family]
\label{moduli-spaces-of-sheaves-msos-label-011}
Line bundles on $\mathbb{P}^1$ of degree $d$
are $\mathcal{O}(d)$. Them
$\mathcal{M}=\text{Spec}\mathbb{C}$. It does
not have a universal family:
$$
\xymatrix{
\mathcal{O}_{\mathbb{P}^1}\ar[r]\ar[d]&
\mathbb{C}\ar[d]\\
\mathbb{P}^1\ar[r]&\text{Spec}\mathbb{C}
}
$$
and $\mathcal{O}(d) \neq
\mathcal{O}_{\mathbb{P}^1}$, $d\neq 0$.
\end{example}


\section{Remarks on stacks}
\label{moduli-spaces-of-sheaves-section-remarks-on-stacks}

The following was not defined properly. As a
stack,
$$
\text{Pic}_{\mathbb{P}^1,d}=\left[
\cdot/\mathbb{G}_m \right] = B\mathbb{G}_m
$$
parametrized by $\mathbb{G}_m$-bundles.
$$
\xymatrix{
\mathcal{L}\setminus\{\text{zero
section}\}\ar[r]\ar[d]&\Gamma^A[?]\ar[d]\\
\mathbb{P}^1\ar[r]_\varphi&B\mathbb{G}_m
}
$$
\medskip\noindent
{\bf Tangent space.} The tangent space of
$\mathcal{M}_{X,P}$, in $\ni E$, is
$\text{Ext}^1(E,E)$ if $\text{Ext}^2(E,E)=0
\implies  \mathcal{M}_{X,P}$ is
smooth. ($\text{Ext}$ groups are defined in
Stacks Project, Algebra; this result
also appears to be a theorem in deformation
theory.)

\begin{example}
\label{moduli-spaces-of-sheaves-msos-label-012}
Let $X$ be a curve. $E$ is a vector bundle of
rank $n$, then
$$
\text{Ext}^2(E,E)=H^{2}(X,E \otimes
E^{\vee})=0
$$
then the moduli space is smooth.
\end{example}

The dimension of the stack is $\dim
T_{\mathcal{M},E}-\dim \text{Aut}(E)$, which
is just $\dim (\text{Ext}^1(E,E) - \dim
\text{Ext}^0(E,E)$, which in turn says
that $\mathcal{X}(X,E \otimes
E^\vee)=V^2(f-g)$.

\medskip\noindent
{\bf Question 2.} Once we have a moduli space
we want so study its geometry.
What can we say about $\mathcal{M}_{X,P,H}$?

Let $X$ be a projective surface and $E$ a
vector bundle.
Let $r = \text{rk} E$, $L = \Lambda^{n}E$.
Consider
$\mathcal{M}_{X,P,H}(r,L,n)$.
$n=c_2(E)$

\begin{theorem}
\label{moduli-spaces-of-sheaves-msos-label-013}
If $\Delta=2rm-(r-1)L^2 \gg 0$ then
$\mathcal{M}_{X,H}(r,L,n)$ is a normal,
generically smooth irreducible
quasi-projective variety. (If $\Delta$ is
small
then it is not true.)
\end{theorem}

\begin{theorem}
\label{moduli-spaces-of-sheaves-theorem-birrational-moduli-spaces}
Fix $H,H'$ ample. For $\Delta \gg 0$,
$\mathcal{M}_{X,H}(r,L,n)$ and
$\mathcal{M}_{X,H}(r,L,n)$ are birrational.
\end{theorem}

\medskip\noindent
Intuitive idea of what follows: there are
ample cones, and walls between them,
and when we cross the
wall, we get to an object birational to the
one we were previously at.

\begin{definition}
\label{moduli-spaces-of-sheaves-definition-wall}
Let $C_X$ be the ample cone $\mathbb{R}
\otimes \text{Num}(X)$. For
$\xi \in\text{Num}(X)$,
$W^\xi:=C_X \cap \{x \in \text{Num}(X)
\otimes\mathbb{R}:x\cdot\xi=0\}$.
$W^\xi$ is called a {\it wall} of type
$(c_1,c_2) \iff \exists G \in\text{Pic}(X)$
with $\xi = G$ such that
\begin{itemize}
\item $G+c_1$ is divisible by $2$ in
$\text{Pic}(X)$
\item $c_1^2-4c_2 \leq G^2 <0$.
\end{itemize}
\end{definition}

\begin{remark}
\label{moduli-spaces-of-sheaves-msos-label-014}
$W^\xi \neq 0$ if there exists an ample line
bundle $L$ with $L\xi =0$,
$W(c_1,c_2)=\bigcup_{\xi}W^1$.
\end{remark}

\begin{definition}
\label{moduli-spaces-of-sheaves-definition-chamber}
A {\it chamber} is a connected component of
$C_X\setminus W(c_1,c_2)$. (Perhaps
this is the same as $C_X \setminus
\bigcup_{\xi}W\xi$.)
\end{definition}

\begin{theorem}[Qin]
\label{moduli-spaces-of-sheaves-theorem-Qin}
The moduli space $M_{X,H}(2c_1,c_2)$ of rank
2 locally free sheaves only depends
on the chamber of $H$.
\end{theorem}

\begin{remark}
\label{moduli-spaces-of-sheaves-remark-remarks}
\begin{itemize}
\item $c_1^2-4c_2<0$ is Bogomolov inequality.
\item  $G^2<0$ is Hodge index theorem
\item $F \subset E$ destabilizing line bundle
$F$ $\iff$ $H$
$c_1(F)>\frac{c_1(E)H}{2}$ $\iff$
$(\underbrace{2c_1(F)-c_1(E)}_{=G})H>0$
$\implies$ $2c_1(F)=c_1(G)+c_1(E) \implies
c_1(G)+c_1(E)$ is divisible by 2.
\item Suppose we have $0 \to F \to E \to M
\otimes |_{Z}\to 0$ with $M$ a line
bundle, $Z$ codimension 2 in $X$.
\end{itemize}
\end{remark}

[Lots of computations I missed.]

\begin{exercise}
\label{moduli-spaces-of-sheaves-exercise-P1timesP1}
$X=\mathbb{P}^1\times \mathbb{P}^1$, $E$ of
rank 2, $c_1=\ell$, $c_2=3$.
$G=a\ell+bm$, $G^2=2ab \implies ab<0$,
$\Delta=c_1^2-4c_2\leq G^2=ab\implies
-12\leq 2ab<0$, $(a+1)\ell+bm$ divisible by
$2$, then $a+1$ is even.

I understand the walls and chambers were
computed:
$L_0=(\ell+7m)\in C_0$,
$L_1=(\ell+5m) \in C_1$ and $L_2=(\ell+3m)
\in C_2$.

We have
\begin{itemize}
\item For $L$ in  $C_0$,
$\mathcal{M}(2,\ell,3)=\emptyset$. (Proved.)
\item For $L$ in $C_1$,
$\mathcal{M}(2,\ell,3)\cong \mathbb{P}^5$.
\item For $L$ in $C_2$,
$\mathcal{M}(2,\ell,3)$ is an open in
$\mathbb{P}^3$.
\end{itemize}
\end{exercise}

For more general rank:

\begin{theorem}
\label{moduli-spaces-of-sheaves-theorem-general-rank}
$F \subset E$ a subsheaf of rank $r'<r$,
$\mu_H(F)=\mu_H(E)$, then
$\xi=rc_1(F)-r'c_1(E)$ satisfies $\xi H=0$
and
$\frac{r^2}{4}\Delta \leq \xi^2<0$, with
$e_2$ $\iff$ $\xi=0$.
\end{theorem}

\begin{theorem}
\label{moduli-spaces-of-sheaves-msos-label-015}
$X$ smooth irreducible projective surface,
$H,H'$ ample and $\Delta\gg 0$. Then
$\mathcal{M}_{X,H}$ and $\mathcal{M}_{X,H'}$
are birational (any rank).
\end{theorem}

For dimension $\geq 3$, the moduli space of
sheaves is in general smooth or
irreducible.

\begin{theorem}[Ein]
\label{moduli-spaces-of-sheaves-theorem-Ein}
Let $C$ be a smooth projective 3-fold, $c_1,H
\in \text{Pic}(X)$, $H$ ample.
Assume there exists $a,b$ such that
$ac_1=bH$. $\mathcal{M}_{X,H}$ rank 2 vector
bundle on $X$, $\mu$ stable with respect to
$d:=c_2H$. $\mu(d)=F$ irreducible
component of $\mathcal{M}_{X,H}(c_1,c_2)$,
$\implies $
$\lim_{d \to \infty} m(d)=\infty$.
\end{theorem}

\begin{example}
\label{moduli-spaces-of-sheaves-example}
In DT (Donaldson-Thomas) theory, $X$ 3-fold,
$$
I(X,P)=\left\{Z \subset X: \substack{Z \text{
a curve with} \\
\text{fixed Hilbert polynomial}}\right\}
$$
$P(u)=3u+1$ twisted cubics have  $P(u)=3u+1$.
But genus 1 curves with a
point also have  $P(u)=3u+1$.
$I(\mathbb{P}^3,3n+1)$ has several
components.
\end{example}


% --- End moduli-spaces-of-sheaves.tex ---

% --- Begin bridgeland-stability-general-theory.tex ---
\chapter{Bridgeland Stability: the General Theory}
\label{chapter-bridgeland-stability-general-theory}


Theory}

Minicourse by Cristian Martinez (Universidad
de los Andes, Bogot\'a, Colombia), CIMPA
school
Florian\'opolis 2025.

Notes at
\href{%
http://github.com/danimalabares/cimpa-floripa
}{%
github.com/danimalabares/cimpa-floripa
}

\bigskip\noindent

{\bf Abstract.} Bridgeland stability is a
powerful tool for extracting geometry from
homological algebra. In particular, it gives
a framework for studying moduli spaces of
objects in a triangulated category, such as
the derived category of an algebraic variety.
The subject was born as a mathematical
interpretation of work in string theory, but
has since impacted many areas, including
classical algebraic geometry, derived
categories of coherent sheaves, enumerative
geometry, homological mirror symmetry, and
symplectic geometry. The goal of this course
is to develop the foundations of Bridgeland
stability, covering the following topics: 1)
The theory of t-structures on triangulated
categories, including examples via tilting.
2) The defniition of stability conditions and
the stability manifold, as well as
Bridgeland's deformation theorem. 3)
Constructions of stability conditions. 4)
Moduli spaces of stable objects.

\bigskip\noindent
\bigskip\noindent

\noindent
{\bf Plan.}
\begin{enumerate}
\item Lecture 1: What is a stability
condition.
\item Lecture 2: Main constructions in
dimensions 2 and 3.
\item Lecture 3: Reaching the Geiseker
chamber.
\item Lecture 4: Some applications on
structures.
\end{enumerate}

\section{Recall}
\label{bridgeland-stability-general-theory-section-recall}

Recall that for $\mu_H$-stability we have the
following property.

\begin{itemize}
\item (Schur's lemma.) If $A,B \in
\text{Coh}$ are $\mu_H$ semistable with
$\mu_H(A)>\mu_H(B)$ then $\Hom(A,B)=0$.
Exercise.
\item (Harder-Narasimhan filtrations.) $E \in
\text{Coh}(X)$ then is a
filtration
$$
\xymatrix{
0\ar[r]&E_0\ar[rr]^{i_0}&&E_1\ar[dl]\ar[r]&
\cdots\ar[r]
&E_{k-1}\ar[rr]^{i_k}&&E_k=E\ar[dl]\\
&&F_1& & & & F_k
}
$$
\iffalse$$
\xymatrix{
0\ar[r]&  E_0
\ar@{^{(}->}[r]&E_1\ar[d]\ar[r]&  \cdots
\ar@{^{(}->}[r]& E_k\ar[d]=E\\ & & F_1 & &
F_k } $$\fi with $E_0$ is the torsion
subsheaf of $ E$, $F_j$ are
$\mu_H$-semistable with
$$
\mu_H(F_1)>\mu_H(F_2)>\ldots>\mu_H(F_k)
$$
\end{itemize}

Exercise: (1) $\implies$ (2).

These two results basically say that
$\text{Coh}(X)$ is generated by
$\mu_H$-semistable sheaves (via extensions)
and torsion sheaves.
$$
0\to E_{j-1}\hookrightarrow E_j\to F_j\to_0
$$
$E_0:=F_0$.

Can we do this on $D^b(X)$?

We look for a subcategory with a fixed slope
to generate the whole category\ldots{}

\medskip

The following definition is intended to
``impose'' Schur's lemma on our objects:

\begin{definition}
\label{bridgeland-stability-general-theory-definition-slicing}

A {\it slicing} $\mathcal{P}$ of $D^b(X)$
consists of full subcategories
$\mathcal{P}(\phi)$ for every $\phi \in
\mathbb{R}$ such that
\begin{enumerate}
\item
$\mathcal{P}(\phi)[1]=\mathcal{P}(\phi+1)$ /
\item If $\phi_1>\phi_2$ then
$\Hom(\mathcal{P}(\phi_1),\mathcal{P}(\phi_2)
)
=0$.
\item (HN property.) For every $E \in D^b(X)$
there is a diagram (filtration)
$$
\phi_1>\phi_2>\cdots>\phi_k
$$
$$
\xymatrix{
0=E_0\ar[rr]^{i_0}&&E_1\ar[dl]\ar[r]&
\cdots\ar[r]
&E_{k-1}\ar[rr]^{i_k}&&E_k=E\ar[dl]\\
&F_1& & & & F_k
}
$$
where $F_j=\text{Cone}(i_j) \in
\mathcal{P}(d_j)$. {\bf Exercise:} this is
unique because of (2).

\item  $F \in Q(\phi)$ is called {\it
semistable} of phase $\phi$. The simple
objects
(those which have no subobjects) in
$\mathcal{P}(\phi)$ are called {\it stable of
phase} $\phi$.

\item  $\phi_1:=\phi^+(E)$,
$\phi_k:=\phi^-(E)$.
\end{enumerate}
\end{definition}

The following definitions is related to the
former, and to Schur's lemma:

\begin{definition}
\label{bridgeland-stability-general-theory-bsgt-label-001}
The {\it heart of a bounded $t$-structure} on
$D^b(X)$ is a full additive subcategory
$\mathcal{A} \subset D^b(X)$ such that
\begin{enumerate}
\item For $i>j$, $\Hom(A[i],B[j])=0$ for $A,B
\in \mathcal{A}$.
\item For every $E \in D^b(X)$ there are
integers $k_1>k_2>\ldots>k_m$ and a
filtration
\begin{equation}
\label{bridgeland-stability-general-theory-equation-diagram-abelian-category}
\xymatrix{
0=E_0\ar[rr]^{i_0}&&E_1\ar[dl]\ar[r]&
\cdots\ar[r]
&E_{k-1}\ar[rr]^{i_k}&&E_k=E\ar[dl]\\
&A_1[k_1]& & & & A_k[k_m]
}
\end{equation}
exact triangles, with $A_j \in \mathcal{A}$.
\end{enumerate}
\end{definition}

\begin{exercise}
\label{bridgeland-stability-general-theory-exercise-A-is-abelian}
$\mathcal{A}$ is abelian.
\end{exercise}

\begin{exercise}
\label{bridgeland-stability-general-theory-exercise}
If $\mathcal{P}$ is a slicing of $D^b(X)$
then $\mathcal{P}(0,]$ (the extension closure
of $\{\mathcal{P}(\phi):0<\phi \leq 1\}$) is
the heart of a bounded $t$-structure.
\end{exercise}

\begin{example}
\label{bridgeland-stability-general-theory-bsgt-label-002}
$\text{Coh}(X)\subset D^b(X)$ is the heart of
a bounded structure.

$$
\xymatrix{
\sigma_{\leq i}(E^\bullet):\ar[d]&
\cdots\ar[r]& E^{j-2}\ar[d]_=\ar[r]&
E^{j-1}\ar[d]_=\ar[r]^{\delta}& \Ker
\delta_j\ar[d]\ar[r]&  0 \ar[d]\ar[r]
&  \cdots\\
\sigma_{\leq j+1}(E^\bullet):\ar[d]&  \cdots
\ar[r]& E^{j-2}\ar[d]\ar[r]
& E^{j-1}\ar[d]\ar[r]^{d_j}&E^j\ar[d]\ar[r]&
\Ker(\delta_{j+1})\ar[d]\\
&& 0& 0&E_2k/\Ker \delta_j& \Ker \delta_{j+1}
}
$$
\medskip
$$
\xymatrix{
\cdots\ar[r]&\sigma_{\leq
j}(E^\bullet)\ar[rr]& &\sigma_{\leq
j+1}(E^\bullet)\ar[dl]\ar[r]&  \cdots\\
&  &  H^{i+j}(E^\bullet)[-j-1]
}
$$
\end{example}

\begin{definition}
\label{bridgeland-stability-general-theory-definition-cohomologies}
The objects $A_j$ appearing in Eq.
\ref{bridgeland-stability-general-theory-equation-diagram-abelian-category} are
called the {\it cohomologies} of $E$ with
respect to the $t$-structure $\mathcal{A}$,
$A_j:=\mathcal{H}_{\mathcal{A}}^{-k_j}(E)$.
\end{definition}

\section{Constructing new hearts}
\label{bridgeland-stability-general-theory-section-constructing-new-heats}

How can we construct other $t$-structures?
Suppose that $\mathcal{A}$ is an abelian
category. A {\it torsion pair} on
$\mathcal{A}$ consists of two full
subcategories $(\mathcal{T}, \mathcal{F})$
such that
\begin{enumerate}
\item $\Hom(T,F)=0$ for all $T \in
\mathcal{T}$ and $F \in \mathcal{F}$.
\item For every $A \in \mathcal{A}$ there is
an exact sequence
$$
\xymatrix{
0\ar[r]&T\ar[r]&A\ar[r]&F\ar[r]&0
}
$$
with $T \in \mathcal{T}$ and $F \in
\mathcal{F}$.
\end{enumerate}

Again: with the first property we may prove
that the exact sequence in the second
property is unique (Exercise).

\begin{exercise}
\label{bridgeland-stability-general-theory-exercise-torsion-pair}
$\mathcal{A}=\text{Coh}(X)$, $\mathcal{T}=$
torsion sheaves, $\mathcal{F}=$ torsion free
sheaves.
\end{exercise}

\begin{proposition}
\label{bridgeland-stability-general-theory-proposition}
If $\mathcal{A}$ is a hear of a bounded
$t$-structure on $D^b(X)$ and
$(\mathcal{T},\mathcal{F})$ is a torsion pair
on $\mathcal{A}$, then
$$
\begin{aligned}
\mathcal{A}^\sharp&=\left\langle{}\mathcal{F}[1],
\mathcal{T}\right\rangle{}\\
&=\left\{E\in
D^b(X):\substack{\mathcal{H}_{\mathcal{A}}
^j(E)=0\\
\text{for $j \neq -1,0$} \\
\mathcal{H}_{\mathcal{A}}^{-1}(E)=\mathcal{F}
\\
\mathcal{H}_{\mathcal{A}}^\bullet(E) in
\mathcal{T} }\right\}
\end{aligned}
$$
which is called the {\it tilted heart
category}.
\end{proposition}

\noindent
{\bf Note.} We can think of those as
complexes $E ^\bullet$ that fit in an exact
triangle
$$
\xymatrix{
0\ar[r]&F[1]\ar[r]&E^\bullet\ar[r]&T\ar[r]&0
}
$$
(which a prori is just an exact triangle, but
it is an exact sequence in our case).

The first proof of this fact was done by
Bridgeland [confirm]:

\begin{proof}[Idea of proof]
[Picture.]
\end{proof}

We show an example to understand how the
proof works.

\begin{example}
\label{bridgeland-stability-general-theory-example}
Suppose that the $\mathcal{A}$-filtration of
an object $E$ is of the form
$$
\xymatrix{
0\ar[r]&F_1[1]\ar[r]&E\ar[r]&F_2\ar[r]&0
}
$$

Let's comute the cohomologies:
$$
\xymatrix{
0=E_0\ar[rr]^{i_0}&&E_1\ar[dl]\ar[rr]^{i_2}&&
E_2=E\ar[dl]\\
&F_1[1]&  & F_2
}
$$

$F_1,F_2 \in \mathcal{A}$. So, the
cohomologies, by definition are
$\mathcal{H}_{\mathcal{A}}^{-1}(E)=F_1$,
$\mathcal{H}_{\mathcal{A}}^0(E)=F_2$.

Now consider the following diagram:
$$
\xymatrix{
&  F[1]\ar[d]&  &  T\ar[d]\\
0 \ar[r]&  F_1[1]\ar[d]\ar[r]&  E
\ar[r]\ar[rd]_f&  F_2\ar[d]\ar[r]& 0\\
& F[1]& & F
}
$$
so wer consider
$$
\xymatrix{
& & F[1]\ar[dl]\ar[d]\\
0\ar[r]&C=\text{Cone}(f)[-1]\ar[r]&E\ar[r]^f&
F\ar[r]&0
}
$$
\end{example}

Recall that in a triangulated category,
$$
\xymatrix{
X\ar[d]\ar[r]&  0 \ar[d]\ar[r]&
X[1]\ar@{.>}[d]\ar[r]& X[1]\ar[d]\\
Y\ar[r]& Z\ar[r]& X^1\ar[r]&Y[1]
}
$$
where the leftmost square commutes.

$$
\xymatrix{
0\ar[rr]&&F[1]\ar[dl]\ar[rr]^{g}&  &  C
\ar[dl]\ar[rr]&  &  E\ar[dl]\\
&  \underbrace{F[1]}_{\mathcal{A}^\sharp[1]}&
&
\underbrace{\text{Cone}(g)}_{\mathcal{A}
^\sharp}&
&
\underbrace{F}_{\mathcal{A}^\sharp[-1]}
}
$$

\section{Stability conditions}
\label{bridgeland-stability-general-theory-section-stability-conditions}

\begin{definition}[Stability condition,
definition 1]
\label{bridgeland-stability-general-theory-definition-stability-condition-1}
A {\it stability condition} is a pair
$\sigma=(\mathcal{P},Z)$ where $\mathcal{P}$
is a slicing of $D^b(X)$, and a morphism $Z$
called the  {\it central change} of the form
$$
\xymatrix{
Z:K_0(D^b(X))\ar[rr]\ar@{->>}[dr]_\nu&  &
\mathbb{C}\\
&  \Gamma \ar[ur]_z
}
$$
where $\Gamma$ is a finite rank lattice,
satisfying that
\begin{enumerate}
\item $Z(E) \in \mathcal{E}_{>0}e^{i\pi
\phi}$ for any nonzero object $E \in
\mathcal{P}(\phi)$.
\item (Support property.)
For a fixed norm $\|\cdot\|$ on $\Gamma
\otimes \mathbb{R}$ we have
$$
C_{\sigma}=\text{inf}\left\{
\frac{|Z(E)|}{\|\nu(E)\|}:\substack{
\text{$E$ non-zero} \\ E \in
\mathcal{P}(\phi),\phi \in \mathbb{R}}
\right\} >0
$$
\end{enumerate}
\end{definition}

\begin{definition}[Stability condition,
definition 2]
\label{bridgeland-stability-general-theory-definition-stability-condition-2}
A {\it stability condition} is a pair
$\sigma=(\mathcal{A},Z)$, where $\mathcal{A}$
is the heart of a bounded $t$-structure on
$D^b(X)$, and a morphism $Z$ of the form
$$
\xymatrix{
Z:K_0(D^b(X))\ar[rr]\ar@{->>}[dr]_\nu&  &
\mathbb{C}\\
&  \Gamma \ar[ur]_Z
}
$$
where $\Gamma$ is a finite rank lattice,
satisfying that
\begin{enumerate}
\item $Z(E) \in \mathbb{R}_{>0}e^{i\pi \phi}$
for $E \in \mathcal{A}$,
$$
\left[
\frac{\text{Im}Z(E)=0}{\text{Im}Z(E)\geq
0}\implies
\text{Re}Z(E)<0 \right]
$$
\item If
$\mu_Z=-\frac{\text{Re}Z}{\text{Im}Z}$ then
objects in $\mathcal{A}$
have the HN property w.r.t. $\mu_Z$ and we
have the support property.
\end{enumerate}
\end{definition}

\begin{exercise}
\label{bridgeland-stability-general-theory-exercise-definitions-are-equivalent}
Prove that the two definitions of stability
are equivalent.
\begin{itemize}
\item Prove Schur's Lemma for
$\mu_Z$-stability.
\item $\mathcal{P}(\phi)=\mu_Z$-semistable
objects of phase $\phi$, i.e., $Z(E)
\in \mathbb{R}_{>0}e^{i\pi\phi}$ for $\phi
\in (0,1]$ and set
$\mathcal{P}(\phi+1)=\mathcal{P}(\phi)[1]$.
\item Take $\mathcal{A}(0,1]=$generated by
$\mathcal{P}(\phi)$ with
$0<\phi\leq 1$.
\end{itemize}
\end{exercise}

\begin{exercise}
\label{bridgeland-stability-general-theory-bsgt-label-003}
If $\dim X=1$, $\mathcal{A}=\text{Coh}(X)$,
$Z= - \text{deg}+i \text{rank}$,
$v=(\text{deg},\text{rank})$.
\end{exercise}

\medskip\noindent
{\bf Rephrasing of support property.} There
exists a quadratic form $Q$ on $\Gamma
\otimes \mathbb{R}$ such that $Q(v(E))\geq 0$
for $E$ $\mu_Z$-semistable, If $Z(V)=0$ then
$Q(V)<0$ for nonzero $V$. This is like
imposing Bogomolov inequality. (See
moduli-spaces-of-sheaves.)

\medskip\noindent
Let $\text{Stab}^\Gamma(X)$ be the set of
stability conditions with central charge
factoring through
$v:K_0(\mathcal{A})\to\Gamma$.

\begin{theorem}[Bridgeland]
\label{bridgeland-stability-general-theory-theorem-Bridgeland}
$\text{Stab}^\Gamma(X)$ is a complex
manifold. Moreover,
$$
\begin{aligned}
\text{Stab}^\Gamma(X) &\longrightarrow
\Hom(\Gamma,\mathbb{C}) \\
\sigma=(\mathcal{A},Z) &\longmapsto Z
\end{aligned}
$$
is a local homomorphism.
\end{theorem}

This is known as Bridgeland deformation
result. In fact, it is a consequence of the
existence of the quadratic form [reference?].

For every object in $\mathcal{A}$ where
$\sigma=(\mathcal{A},Z)$ is a stability
condition there is a HN filtration
$$
\xymatrix{
0\ar[r]&E_0\ar[rr]^{i_0}&&E_1\ar[dl]\ar[r]&
\cdots\ar[r]
&E_{m-1}\ar[rr]^{i_m}&&E_m=E\ar[dl]\\
&&F_1& & & & F_m
}
$$
$F_1 \in \mathcal{P}(\phi_1)$,
$\phi_1:=\phi^+_\sigma(E)$, $F_m \in
\mathcal{P}(\phi_m)$,
$\phi_m:=\phi_\sigma^-(E)$.

The topology on $\text{Stab}^\Gamma(X)$ is
the coarsest topology that makes the
functions
$$
\sigma \mapsto \phi^+_\sigma,\qquad \sigma
\mapsto \phi^-_\sigma, \qquad
\sigma \mapsto Z
$$
[continuous].

\medskip\noindent
Now let $\dim X=2$. Let $H,B \in
\text{NS}(X)_\mathbb{R}$, $H$ ample.
$$
\text{ch}^B(E)=e^{-B}\text{ch}(E)=
\text{``$\text{ch}(E\otimes
(-B))$''}
$$
When computing the Chern characters we obtain
$$
\begin{aligned}
\text{ch}_0^B = \text{ch}_0,&\qquad
\text{ch}_1^B=\text{ch}_1-B\text{ch}_0\\
\text{ch}_2^B&=\text{ch}_2-B\text{ch}_1+
\frac{B^2}{2}\text{ch}_0
\end{aligned}
$$

\medskip\noindent
{\bf $B$-Twisted Mumford slope.}
$$
\mu_{B,H}=
\begin{cases}
\frac{H \text{ch}_1^B}{H^2 \text{ch}_0}\qquad
&\text{if }\text{ch}_0\neq 0 \\
+\infty\qquad &\text{if }\text{ch}_0=0
\end{cases}
$$
$$
\begin{aligned}
\mu_{B,H}(E)&=\frac{H(\text{ch}_1(E)
-B\text{ch}_0(E)}{H^2\text{ch}_0(E)}\\
&=\mu_H(E)-BH
\end{aligned}
$$
$$
\mathcal{T}_{B,H}=\{E \in
\text{Coh}(X):\mu_H\text{-HN factors have }
\mu_{B,H}>0\}
$$
$$
\mathcal{F}_{B,H}=\{E \in \text{Coh}(X):
\mu_H\text{-HN factors have
}\mu_{B,H}\leq 0\}
$$
\begin{itemize}
\item $\Hom(T,F)=0$ for $T \in
\mathcal{T}_{B,H}$, $F \in \mathcal{F}_{B,H}$
(Schur's lemma).
\item If $\mathcal{E}\in\text{Coh}(X)$, take
its HN-filtration
$$
\xymatrix{
0\ar[r]&E_0\ar[rr]^{i_0}&&E_1\ar[dl]\ar[r]&
\cdots\ar[r]
&E_{m-1}\ar[rr]^{i_m}&&E_m=\mathcal{E}\ar[dl]
\\
&&F_1& & & & F_m
}
$$
$$
\mu_{B,H}(E_0)>\mu_{B,H}(F_1)>\ldots>\mu_{B,
H}
(F_m)
$$
Let $j$ such that $\mu_{B,H}(F_j)>0$ and
$\mu_{B,H}\leq 0$.
$$
0 \to \underbrace{E_j}_{\mathcal{T}_{B,H}}\to
\mathcal{E}
\to \underbrace{\mathcal{E}/E_j}_{\in
\mathcal{F}_{B,H}}\to 0
$$
\end{itemize}

Set a tilted heart
$\text{Coh}^{B,H}(X)=\left\langle{}\mathcal{F}_{B,H}
[1],
\mathcal{T}_{B,H}\right\rangle{}$.

Generators:
\begin{itemize}
\item $T$ torsion sheaves $\to
H\text{ch}_1^B(T)\geq 0$.
\item $\mathcal{E}$ $\mu_{B,H}$-semistable
sheaves with $\mu_{B,H}>0 \to
H\text{ch}_1^B (\mathcal{E})>0$.
\item $F[1]$,  $F$ is $\mu_{B,H}$-semistable
with $\mu_{B,H}\leq 0$.
\end{itemize}
Then
$$
H\text{ch}_1^B (F)\leq 0,\qquad
H\text{ch}_1^B (F[1])\geq 0
$$
since shifting by 1 changes the sign of the
Chern character. [Comment about imaginary
part of stability condition.]

Is there a central charge $Z$ such that
$\sigma_{B,H}=(\text{Coh}^{B.H}(X),Z)$ is a
stability condition?
$$
Z_{B,H}=\text{Re}Z_{B,H}+i H\text{ch}_1^B
$$
We need: $H\text{ch}_1^B(E)=0$ for some $E
\in \text{Coh}^{B,H}(X)$ then
$\text{Re}Z_{B,H}(E)<0$.

Ingredients:
\begin{enumerate}
\item (Hodge index theorem.) $(DH)^2 \geq D^2
H^2$.
\item (Bogomolov inequality.) If $E$ is
$\mu_{B,H}$-semistable then
$$
\begin{aligned}
\Delta(E)&=[\text{ch}_1(E)]^2-2\text{ch}_0(E)
-\text{ch}_2(E)\geq
0\\
&=\text{ch}_1^B(E)^2-2\text{ch}_0^B\text{ch}
_2^B(E)\geq
0
\end{aligned}
$$
We obtain the {\it BH-discriminant}
$$
\Delta_{B,H}(E):=(H\text{ch}_1^B(E))
^2-2\text{ch}_0(E)H^2\text{ch}_2^B(C)
\geq H^2\Delta(E)\geq 0
$$
If $E$ is $\mu_H$-semistable and
$\text{Im}Z_{B,H}(E)=0$, $H\text{ch}_1^B(E)=0
\implies \text{ch}_2^B(E)\leq 0$,
$\text{ch}_2^B(E)-a\text{ch}_0(E)<0$
for any $a>0$.
\end{enumerate}

\begin{theorem}[Bridgeland, Areni-Bestum]
\label{bridgeland-stability-general-theory-theorem-Bridgeland-Areni-Bersum}
$$
Z_{B,H}=-(\text{ch}_2^B-a\text{ch}_0)+
iH\text{ch}_1^B
$$
is the central charge of a stability
condition on $\text{Coh}^{B,H}(X)$ for any
$a>0$.
\end{theorem}

\begin{theorem}
\label{bridgeland-stability-general-theory-theorem-BH-discriminant}
$\Delta_{B,H}\geq 0$ for any
$v_{B,H,a}$-semistable object, where
$v_{B,H,a}=\frac{\text{ch}_2^B-a\text{ch}_0}
{H\text{ch}_1^B}$.
\end{theorem}

\medskip\noindent
For $n:=\dim X\geq 2$,
$$
\mu_{B,H}=
\begin{cases}
\frac{H^{n-1}\text{ch}_1^B}{\text{ch}_0}
\qquad
&\text{if }\text{ch}_0 \neq 0 \\
+\infty\qquad &\text{if }\text{ch}_0=0
\end{cases}
$$
$\text{Coh}^{B,H}(X)$, $\Delta_{B,H}(E)=\dim
n$ analogy. $Z_{B,H,a}$ does not satisfy the
positivity property of a central charge since
$$
Z_{B,H,a}=-(H^{n-2}\text{ch}_2^B-a\text{ch}
_0+
iH^{n-1}\text{ch}_1^B)
$$
[skyscraper sheaf argument]

\begin{exercise}
\label{bridgeland-stability-general-theory-exercise-E-is-a-sheaf}
Let $E \in \text{Coh}^{B,H}(X)$. If
$\text{ch}_0(E)>0$ and $E$ is
$\nu_{B,H,a}$-semistable for all $a\gg 0$,
then $E$ is a sheaf. Moreover, $E$ is
$\mu_{B,H}$-semistable.
\end{exercise}

\begin{proposition}[Lo,---]
\label{bridgeland-stability-general-theory-proposition-duplicate-2}
If $a>\mu_{B,H}(E)\Delta_{B,H}(E)$, then if
$E$ is {\it twisted} semistable sheaf then
$E$ is $\nu_{B,H,a}$-semistable.
\end{proposition}

\medskip\noindent
If $\dim X=3$ consider the subcategories of
$\text{Coh}^{B,H}(X)$
$$
T_{B,H,a}=\{E \in
\text{Coh}^{B,H}(X):\text{HN factors with }
\nu_{B,H,a}>0\}
$$
$$
\mathcal{F}_{B,H,a}=\{E \in
\text{Coh}^{B,H}(X):\text{HN factors with }
\nu_{B,H,a}\leq 0\}
$$
$$
\mathcal{A}=\left\langle{}\mathcal{F}_{B,H,a}[1],
\mathcal{T}_{B,H,a}\right\rangle{}
$$
supports a stability condition.

\section{GBG inequality conjecture}
\label{bridgeland-stability-general-theory-section-GBG-inequality-conjecture}

Here's the construction so far:
$$
\xymatrix@R=1.4em{
\textit{Coh}(X)\ar[r]^{\substack{\text{first}
\\
\text{tilt}}}
&
\textit{Coh}^{B,H}(Z)\ar[r]
^{\substack{\text{second}
\\ \text{tilt}}}
&  \mathcal{A}^{B,H,\alpha}\\
\underbrace{\mu_{B,H}}_{\text{Mumford
stability}}
&
\underbrace{\nu_{B,H,\alpha}}_{\text{Tilt}}
&  \underbrace{\lambda_{B,H,\alpha,s}}
_{\substack{\text{Bridgeland} \\
\text{stability}\\\text{when possible}}}\\
\Delta \geq 0
&  \Delta_{B,H}\geq 0
&\substack{\text{Conjecture} \\
\text{generalized}\\\text{BG inequality}}
}
$$
And
$$
\begin{aligned}
Z_{\beta,H,\alpha,s}&=-\left(ch^B_3-\left(s+
\frac{1}{6}\right)\alpha^2
H^2ch_1^B)\right)\qquad s>0\\
&+i\left(Hch_2^B-\frac{\alpha^2}{2}
H^3ch_0\right)
\end{aligned}
$$

\medskip\noindent
{\bf GBG-inequality.} If $E$ is
$\nu_{B,H,\alpha}$-semistable,
\begin{equation}
\label{bridgeland-stability-general-theory-equation-GBG}
\begin{aligned}
H ch_2^B(E)-\frac{\alpha^2}{2}H^3
ch_0(E)&=0\\
\implies
ch_3^{B}-\frac{1}{6}\alpha^2H^2ch_1^B(E)&\leq
0
\end{aligned}
\end{equation}

\medskip\noindent
{\bf Bayer-Maci-Toda.} If Eq.
\ref{bridgeland-stability-general-theory-equation-GBG} is true then
$(\mathcal{A}^{B,H,\alpha},Z_{B,H,\alpha,s})$
is a strability condition.

\section{Remarks on
$\nu_{B,H,\alpha}$-stability}
\label{bridgeland-stability-general-theory-bsgt-label-004}
Keep in mind that
$$
\nu_{B,H,\alpha}=
\frac{Hch_2^B-\frac{\alpha^2}
{2}ch_0H^3}
{H^2 ch_1^B}
$$
and
$$
L \in \textit{Coh}^{B,H}(X) \qquad (L[1])
$$
as long as
$$
H^2ch_1^B(L)>0 \qquad (H^2ch_1^B(L[1])\leq 0.
$$
\medskip
If $E$ is $\nu_{B,H,\alpha}$-semistable for
$\alpha\gg 0$ then $E\in\textit{Coh}(X)$ and
is semistable with respect to the stability
given by
$$
q(t)=t \frac{H^2 ch_1}{ch_0}+ \frac{H
ch_2}{ch_0},\qquad t\gg 0
$$
and if $E$ is $q(t)$-semistable then $E$ is
$\nu_{B,H,\alpha}$-semistable for
$\frac{\alpha^2}{2}H^2>\mu_{B,H}(E)\Delta_{B,
H}(E)$.

If $X$ has Picard rank $\rho(X)=1$ then
$$
\Delta_{B,H}=H^3\Delta
$$
If $L$ is a line bundle
$\Delta_{B,H}(L)=H^3\Delta(L)=0$ (as $L[1])$,
then $L$ is $\nu_{B,H,\alpha}$-semistable as
long as $L \in \textit{Coh}^{B,H}(X)$ (as
$L[1])$

\section{Search for counterexamples}
\label{bridgeland-stability-general-theory-section-search-for-counterexamples}

{\bf Idea.} Look for counterexamples with
$ch_0=0$. When does $\mathcal{O}_D$ for
$D\geq 0$ divisor on $X$ (3-fold) produces a
counterexample?

Fix $H$ ample, $B=\beta H$. We need:
\begin{enumerate}
\item $\mathcal{O}_D$ is
$\nu_{\beta,\alpha}$-semistable.
\item $\nu_{\beta,\alpha}(\mathcal{O}_D)=0$.
\item
$ch_3^B(\mathcal{O}_D)-\frac{\alpha^2}{6}
H^2ch_1^B(\mathcal{O}_D)>0$.
\end{enumerate}

\begin{lemma}
\label{bridgeland-stability-general-theory-lemma-counterexample}
If $D \geq 0$ satisfies
$$
D^3>\frac{(DH^2)^3}{4(H^3)^2}+\frac{3}{4}
\frac{(D^3H)^2}{DH^2}
$$
then $\mathcal{O}_D$ is a counterexample for
the GBG inequality
\ref{bridgeland-stability-general-theory-equation-GBG}.
\end{lemma}

\begin{example}
\label{bridgeland-stability-general-theory-example-duplicate-2}
Let $Y$ be any smooth projective 3-fold,
$X=Bl_pY$, $D$ the exceptional divisor,
$L=\pi^*A$ with $A$ ample, $H=mL-D$ ample on
$X$.
\end{example}

[Slides]

\section{Concrete constructions}
\label{bridgeland-stability-general-theory-section-concrete-constructions}

How is a wall equation? Suppose that $\dim
X=2$, or even better that $X=
\mathbb{P}^2$. Also let
$\Gamma=\supset \oplus \supset \oplus
\frac{1}{2}\mathbb{Z}$ and
$\nu=(ch_0,ch_1,ch_2):=(r,c,d)$.

$$
\begin{aligned}
\nu_{\beta,\alpha}&
=\frac{ch_2^\beta-\frac{\alpha^2}{2}ch_0H^2}
{H
ch_1^\beta}\\
&=\frac{ch_2-\beta
ch_1+\left(\frac{\beta^2}{2}-
\frac{\alpha^2}{2}\right)ch_0}{ch_1-\beta
ch_0}
\end{aligned}
$$
A wall in the $(\beta,\alpha)$-plane is
produced by an exact sequence in
$\textit{Coh}^\beta(\mathbb{P}^2)$
$$
\xymatrix{
0\ar[r]&A\ar[r]&E\ar[r]&B\ar[r]&0
}
$$
with
$\nu_{\beta,\alpha}(A)=\nu_{\beta,\alpha}(E)$
and $ch(E)=V$.

\medskip\noindent
First: numerics: $w=(r',c',d')$ and solve
$\nu_{\beta,\alpha}(w)=\nu_{\beta,\alpha}(v)
$.

$$
\begin{aligned}
\beta^2-2\beta\left(\frac{rd'-r'd}{rc'-r'c}
\right)+\alpha^2&=
\frac{d c'-d'c}{rc'-r'c}
\end{aligned}
$$
$$
c=\frac{rd'-r'd}{rc'-r'c},\qquad
R=\sqrt{c^2+\frac{d c'-d'c}{rc'-r'c}}
$$
[Picture.] Look for wall for 1-dimensional
sheaves ($r=0$, $c>0$). $c=\frac{d}{c}$ same
center.

\begin{remark}
\label{bridgeland-stability-general-theory-bsgt-label-005}
Walls either coincide or they don't
intersect. [Picture.]
\end{remark}

\begin{lemma}[Bertan's lemma]
\label{bridgeland-stability-general-theory-lemma-Bertan}
If an object $A$ destabilizes $E$ at some
point then it destabilizes $E$ at every point
of the corresponding wall.
\end{lemma}

\medskip\noindent
Take $\beta=\frac{d}{c}$. If $A
\hookrightarrow E  \xrightarrow{surj.} B$ is
short exact in
$\textit{Coh}^\beta(\mathbb{P}^2)$.
\begin{itemize}
\item If $ch(E)=(0,c,d)$ and $E$ is
a sheaf then $A$ is a sheaf, just not a
subsheaf.
$$
\xymatrix{
0\ar[r]&\underbrace{\mathcal{H}^{-1}A}
_{\implies
0}\ar[r]
&  \underbrace{\mathcal{H}^{-1}E}_{=0}\ar[r]
&  \mathcal{H}^{-1}B\ar[r]
&  \mathcal{H}^0A\ar[r]&
\mathcal{H}^0E\ar[r] &
\mathcal{H}^0B\ar[r]&0
}
$$
$\implies$ $a= \mathcal{H}^0A$ sheaf.
\begin{enumerate}
\item $A$ produces a semi-circular wall only
when $ch_0(A)>0$.
\item $0 \leq  ch_1^\beta(A) \leq
ch_1^\beta(E)=c$.
This says that there are finitely many
possibilities for $c_1^\beta(A)$.
\item At the wall $W_{A,E}$ we have $A$ is
$\nu_{\beta,\alpha}$-semistable.
$$
\Delta_{\beta,\alpha}(A)\geq 0
$$
$$
ch_1^\beta(A)^2-2ch_0(A)ch_2^\beta(A)\geq 0
$$
$$
ch_1^\beta(A)^2\geq 2ch_0(A)ch_2^\beta(A)
$$
\item Wall equation
$$
\begin{aligned}
\nu_{\frac{1}{c}\alpha}(A)&=\nu_{\frac{1}{c},
\alpha}(E)\\
&=\frac{ch_2(E)-\beta ch_1(E)}{ch_1(E)}\\
&=\frac{d-\frac{d}{c}c}{c}=0
\end{aligned}
$$
$$
\begin{aligned}
ch_2^\beta(A)-\frac{\alpha^2}{2}ch_0^\beta(A)
&
=0\\
ch_2^\beta(A)=\frac{\alpha^2}{2}ch_2^\beta(A)
&
\geq 0
\end{aligned}
$$
$\implies$ {\bf walls are finite}.
\end{enumerate}
\end{itemize}

\medskip\noindent
Moduli of $\nu_{\beta,\alpha}$-semistable
objects are projective for $X=\mathbb{P}^2$.
[Picture of walls as nested circles in a
plane.] $\nu_{\beta \alpha}$-semistability
$\implies$ King's $\Theta$-semistability.

\begin{example}
\label{bridgeland-stability-general-theory-example-duplicate-3}
$v=(0,4,-4)$, $M_H(v)\leftarrow$ sheaves can
be found in one of the following short exact
sequences in
$\textit{Coh}^{-1}(\mathbb{P}^2)$.

[Duz\'et-Muican]

\noindent
Codimension 3:
$$
\xymatrix{
0\ar[r]&\mathcal{O}(1)\ar[r]&\mathcal{F}
\ar[r]
&\mathcal{O}(-2)\ar[r]&0
}
$$
Codimension 1:
$$
\xymatrix{
0\ar[r]&\mathcal{I}_P(1)\ar[r]&\mathcal{F}
\ar[r]
&\mathcal{I}_q^\vee(-3)[q]\ar[r]&0
}
$$
Codimension 0:
$$
\xymatrix{
0\ar[r]&2\mathcal{O}\ar[r]&\mathcal{F}\ar[r]&
\mathcal{O}(-2)[1]\ar[r]&0
}
$$
The first wall exchanges
$$
\xymatrix{
0\ar[r]&\mathcal{O}(1)\ar[r]&\mathcal{F}
\ar[r]
&\mathcal{O}(-3)[1]\ar[r]&0
}
$$
for (upon wall crossing)
$$
\xymatrix{
0\ar[r]&\mathcal{O}(-3)[1]\ar[r]&E^*\ar[r]&
\mathcal{O}(1)\ar[r]&0
}
$$
$$
\begin{aligned}
\text{Ext}^1(\mathcal{O}(1),\mathcal{O}(-3)
[1]
)
=\text{Ext}^2(\mathcal{O}(1),\mathcal{O}(-3))
\cong \Hom(\mathcal{O},\mathcal{O}(1))^\vee
\end{aligned}
$$
\end{example}

\begin{theorem}[Li-Zhao]
\label{bridgeland-stability-general-theory-theorem-Li-Zhao}
IF $v$ is primitive and $\sigma,\sigma'$ are
generic geometric with respect to $ v$ then
$\mathcal{M}_{\sigma}(v)$ and
$\mathcal{M}_{\sigma'}(v)$ are birational.
\end{theorem}

\medskip\noindent
We finish this minicourse with another
application of Schur's lemma:

\begin{theorem}[Koseki]
\label{bridgeland-stability-general-theory-theorem-Koseki}
Let $X$ be a smooth projective surface over
$k=\overline{k}$ with $\text{char}(k)>0$,
then there is a constant $C_{[X]}\geq 0$ and
dependending only on the birational
equivalence clas of $X$ such that
$$
\tilde{\Delta}(E)=\Delta(E)+C_{[X]}ch_0^2(E)
\geq
0
$$
for any $\mu_H$-semistable sheaf.
\end{theorem}

Morever
\begin{enumerate}
\item If $X$ is minimal and of general type,
$$
C_{[X]}=2+5K_X^2-\chi(\mathcal{O}_X)
$$
\item If $K(X)=1$ ($K$ is the Kodaira
dimension) and $X$ is quasi-elliptic then
$$
C_{[X]}=2-\chi(\mathcal{O}_X),\qquad
\text{[char 2+3]}
$$
\item $C_{[X]}=0$
\end{enumerate}

\medskip\noindent
{\bf Corollary.}
$$
Z_{B,H,\alpha}^{C_{[X]}}=-\left(\cong_2^B-
\left(\frac{C_{[X]}}{2H^2}+\frac{\alpha^2}{2}
\right)ch_0H^2\right)
+ich_1^B(X)
$$
is a stability condition on
$\textit{Coh}^{B,H}(X)$.

\begin{remark}
\label{bridgeland-stability-general-theory-remark-inequality}
$$
\tilde{\Delta}_{B,H}=\Delta_{B,H}+
ch_0^2C_{[X]
}\leq
2
$$
on any
$\nu_{B,H,\alpha}^{C_{[X]}}$-semistable
object.
\end{remark}

\begin{lemma}[Low key lemma]
\label{bridgeland-stability-general-theory-lemma-low-key-lemma}
\begin{itemize}
\item For $\alpha\gg 0$, $\nu$-semistable are
the Gieseker-semistable sheaves.
\item IF
$\alpha^2H^2>\mu_{B,H}(E)\tilde{\Delta}_{B,H}
(E)$
and if $E$ is
Gieseker-semistable $\implies$ $E$ is
Bridgeland semistable.
\end{itemize}
\end{lemma}

\begin{theorem}
\label{bridgeland-stability-general-theory-theorem-vanishing}
$H^2>6 C_{[X]}$ then $H^1(H \otimes K_X)=0$.
(If $\text{char}(k)=0$ we can take
$C_{[X]}=0$.)
\end{theorem}

\begin{proof}[Proof idea.]
Find a stability condition $\sigma=(Z,s)$
such that
\begin{enumerate}
\item $H$,  $\mathcal{O}[1] \in \mathcal{A}$
are $\sigma$-semistable.
\item $\mu_Z(H)>\mu_Z(\mathcal{O}[1])$
$$
\implies H^{1}(H \otimes K_X) \cong
\text{Ext}^1(H,\mathcal{O})^\vee
= \Hom(H_1\mathcal{O}[1])
$$
\end{enumerate}
\end{proof}
% --- End bridgeland-stability-general-theory.tex ---

% --- Begin geometric-stability-conditions-and-group-actions.tex ---
\chapter{Geometric Stability Conditions and Group Actions}
\label{chapter-geometric-stability-conditions-and-group-actions}


Group Actions}

\noindent
Minicourse by Hannah Dell (University of
Bonn, Germany), CIMPA school Florian\'opolis
2025.

\medskip\noindent
Hannah's notes at
\href{https://www.hannahdell.me/cimpa2025}
{hannahdell.me/cimpa2025}

\medskip\noindent
These notes at
\href{%
http://github.com/danimalabares/cimpa-floripa
}{%
github.com/danimalabares/cimpa-floripa
}

\bigskip\noindent

{\bf Abstract.} Group actions on categories
arise naturally from symmetries of varieties
and quivers, but how does this interact with
Bridgeland stability? In the frist half of
this course we will introduce equivariant
categories, which generalises the category of
equivariant sheaves. Then we will show there
is a correspondence between stability
conditions on a category with a fniite group
action, and stability conditions on the
equivariant category -- this will also play a
role in Xiaolei Zhao's course. We will use
this to produce stability conditions on
quotient varieties (and stacks).

In the second half of the course, we will
apply this to study open questions about the
geometry of the stability manifold: in
particular, we will discuss "geometric
stability conditions" -- those for which all
skyscraper sheaves of points are stable. In
practice, these are constructed using slope
stability for sheaves. Some varieties have
only geometric stability conditions, whereas
in other cases, there are more (for example
if there is an equivalence with quiver
representations). Lie Fu, Chunyi Li, and
Xiaolei Zhao were the frist to provide a
general result explaining this phenomenon. In
particular, they showed that if a variety has
a fniite map to an abelian variety, then all
stability conditions are geometric. We will
test the converse on free quotients of
abelian varieties by fniite groups, including
Beauville-type and bielliptic surfaces. This
is based on joint work with Edmund Heng and
Anthony Licata.

\bigskip\noindent
\bigskip\noindent

\section*{Motivation}
\label{geometric-stability-conditions-and-group-actions-section-motivation}

The Bridgeland stability machine:

$$
[diagram]
$$
\noindent
{\bf Question 1.} How does the geometry of
$\text{Stab}(X)$ relate to $X$?
\noindent
{\bf Question 2.} How to invariants (of $X$
or $\ddot\cup $) behave under (finite) group
actions?

\medskip\noindent
{\bf Goal of these lectures.} Answer both by
studying ``free quotients'' (and see lots of
examples along the way!)

\section{Geometric stability conditions on
surfaces}
\label{geometric-stability-conditions-and-group-actions-gscag-label-001}
\noindent
{\bf Setup.} $X$ smooth projective surface
over $\mathbb{C}$.
$$
\begin{aligned}
\text{NS}(X)&:=\text{Pic}(X)/\text{Pic}^0(X)
\qquad
\text{Neron-Severi group}\\
\text{NS}_{\mathbb{R}}(X)& :=\text{NS}(X)
\otimes \mathbb{R}\\
\text{Amp}_{\mathbb{R}}(X)&:=\text{ ample
cone, i.e. $H \in
\text{NS}_{\mathbb{R}}(X)$ such that $H^2>0$}
\end{aligned}
$$

\begin{definition}
\label{geometric-stability-conditions-and-group-actions-definition-slope-and-stability}
Let $0 \neq  E \in \text{Coh}(X)$ and $H \in
\text{Amp}_{\mathbb{R}}(X)$,
\begin{enumerate}
\item The {\it $H$-slope} of $E$ is
$$
\mu_H(E):=
\begin{cases}
+\infty\qquad &\text{ch}_0(E)=0 \\
\frac{H\cdot \text{ch}_1(E)}{H^2
\text{ch}_0(E)}\qquad &\text{otherwise}
\end{cases}
$$
\item $E$ is {\it $H$-(semi)stable} if $0
\neq  F \subset E$ implies
$$
\mu_H(F) < \mu_H(E/F)\qquad (\mu_H\leq
\mu_H(E/F)
$$
\end{enumerate}
\end{definition}

\begin{example}
\label{geometric-stability-conditions-and-group-actions-example-semistable-condition}

\end{example}

\noindent
{\bf Problem.} [missing]

\begin{definition}
\label{geometric-stability-conditions-and-group-actions-gscag-label-002}
A path $\sigma=(\mathcal{A},Z)$ is a {\it
(numerical) Bridgeland stability condition}
for $D^b(X)$ if
\begin{enumerate}
\item $\mathcal{A} \subset D^b(X)$ is the
heart of an bounded $t$-structure.
\item $Z:K_0(X) \to \mathbb{C}$ homomorphism
(the {\it central charge}) such
that
\begin{itemize}
\item $0 \neq E \in \mathcal{A} \implies
Z(E) \in \mathbb{H}$, where
$\mathbb{H}$ is the upper half plane in
$\mathbb{C}$.
\item HN filtrations exist.
\item Factors via $\text{Knum}(X)$.
\end{itemize}
\item Support property.
\end{enumerate}
\end{definition}

Let $X$ be a surface, and fix  $H \in
\text{Amp}_\mathbb{R}(X)$ and $\beta \in
\mathbb{R}$. Define
$$
\begin{aligned}
\tau_{M,\beta}&:=\{E \in \text{Coh}(X):E
\overset{\text{surj.}}{\to} Q \neq 0
\implies \mu_M(Q)>\beta\}\\
\mathcal{F}_{H,\beta}&:=\{E \in
\text{Coh}(X):0 \neq  F \subset E
\implies \mu_M(F) \leq \beta\}
\end{aligned}
$$

\begin{exercise}
\label{geometric-stability-conditions-and-group-actions-exercise-torsion-pair}
$(\tau_{M,\beta},\mathcal{F}_{H,\beta})$ is a
torsion pair in $\text{Coh}(X)$.
\end{exercise}

\begin{definition}
\label{geometric-stability-conditions-and-group-actions-definition-Coh-H-beta}
$$
\text{Coh}^{H,\beta}(X):=\left\{ E \in
D^b(X):
\substack{H^0(E) \in \tau_{H,\beta} \\
H^{-1}(E) \in \mathcal{F}_{H,\beta}\\
H^i(E)=0} \right\}
$$
\end{definition}

\begin{exercise}
\label{geometric-stability-conditions-and-group-actions-exercise-new-slope}

\end{exercise}

There are notes available about the
following:

\begin{theorem}[Bridgeland '08,
Arcara-Bertam, Macri-Schmidt]
\label{geometric-stability-conditions-and-group-actions-theorem-Bridgeland}
For $\alpha\gg 0$,
$$
\mathbb{C}\cdot \{\sigma_{H,D,\alpha,\beta}=
(\text{Coh}^{H,B}(X),Z_{H,D,\alpha,\beta})\}
$$
is a continuous family of Bridgeland
stability conditions.
\end{theorem}

\noindent
{\bf Question 3.} Is this all of
$\text{Stab}(X)$?

\begin{exercise}
\label{geometric-stability-conditions-and-group-actions-exercise-O-is-stable}
$\mathcal{O}_x$ is
$\sigma_{H,D,\alpha,\beta}$-stable for all
$x\in X$.
\end{exercise}

\begin{definition}
\label{geometric-stability-conditions-and-group-actions-gscag-label-003}
$X$ smooth projective variety, $\sigma \in
\text{Stab}(X)$ is {\it geometric} if for all
$x \in X$, $\mathcal{O}_x$ is
$\sigma$-stable. Write
$\text{Stab}^{\text{geo}}(X)=$ all geometric
stability conditions.
\end{definition}

\medskip\noindent
Now we can refine Question 3:

\medskip\noindent
{\bf Question 3.1.} Does Theorem
\ref{geometric-stability-conditions-and-group-actions-theorem-Bridgeland} describe all
geometric stability conditions?

\medskip\noindent
{\bf Question 3.2.} Do there exist
nongeometric stability conditions?

\medskip\noindent
{\bf Answer 3.2.}
\begin{itemize}
\item $\dim X=1$ no! Unless $X=\mathbb{P}^1$.
\item $\dim X \leq 3$? Yes if $X$ has a full
exceptional collection.
\item $\dim X=2$?
\begin{enumerate}
\item If $X$ is an abelian surface, no!
\item Yes for $\mathbb{P}^2$, and more
generally for rational surfaces.
\item Yes for K3 surfaces.
\item $X \supset C \cong \mathbb{P}^1$ such
that $C^2 < 0$.
\end{enumerate}
\end{itemize}

\medskip\noindent
{\bf Takeaway.} $\text{Stab}(X)$ sees
something about how (1) is different to
(2)-(4) but what is it?

\medskip\noindent
The following theorem is valid in any
dimension:

\begin{theorem}[Lie Fu-Chunyi Li-Xiaolei Zhao
'21]
\label{geometric-stability-conditions-and-group-actions-theorem-Albanese-variety}
$X$ has a finite Albanese morphism ($\iff$
there exist a finite morphism to an abelian
variety), then
$$
\text{Stab}(X) = \text{Stab}^{\text{geo}}(X)
$$
\end{theorem}

\begin{proof}[Sketch of proof]
\begin{itemize}
\item $\sigma \in \text{Stab}(X)$.
\item Let $E_1,\ldots,E_k$ be the {\it
Jordan-H\"older factors} of $\mathcal{O}_X$
w.r.t. $\sigma$ (i.e. $E_i$ $\sigma$-stable
connected component containing the identity).
\item $\mathcal{L} \in \text{Pic}^0(X)$.
\end{itemize}
By [Polishchuk '07], $\otimes \mathcal{L}$
does not change stability. Also
$\mathcal{O}_x \otimes \mathcal{L}\cong
\mathcal{O}_x$ so HN filtration and
Jordan-H\"older factors are preserved
$\implies
$ $E_i \otimes \mathcal{L} \cong E_i$.
\begin{itemize}
\item $\mathcal{L} \in
\text{Pic}^0(\text{Alb}(X))$, $E_i \otimes
\text{alb}^\vee(\mathcal{L})\cong E_i$.
\end{itemize}
$$
\begin{aligned}
\implies & R\text{alb}_*(E_i)\otimes
\mathcal{L} \cong R\text{alb}_*(E_i)\qquad
\text{(projection formula)}\\
\implies &  R\text{alb}_*(E_i)\text{ has
finite support}\qquad
\text{([Polishucuk '03])}\\
\implies &E_i\text{ has finite support}\qquad
\text{($\text{alb}_*$ finite)}\\
\overset{\text{claim}}{\implies
}&\mathcal{O}_x\text{ is $\sigma$-stable}
\end{aligned}
$$
\end{proof}

\medskip\noindent
{\bf Question 4.} Albanese morphism of $X$ is
not finite implies that there exist
nongeometric stability conditions?

\medskip\noindent
The idea is to investigate this for examples
arising as free quotients. If $G \mathbb{y}
X$ is a free action of a finite group, then
$\text{alb}_X$ is finite, and for $Y:=X/G$ we
have $\text{alb}_Y$ is not finite. In
sections 2 and 3 we will compate
$\text{Stab}(X)$ and $\text{Stab}(X/G)$.

\section{Equivariant categories}
\label{geometric-stability-conditions-and-group-actions-section-equivariant-categories}

Let $G$ be a finite group and $\mathcal{D}$ a
$k$-linear category (i.e. all $\Hom$ are
$k$-vector spaces, in particular it is
additive), where $k=\bar{k}$ and
$(|G|,\text{char}k)=1$.

The following definition is by [Deligne '97].

\begin{definition}[Group action]
\label{geometric-stability-conditions-and-group-actions-definition-group-action}
A {\it (right) action} of $G$ on
$\mathcal{D}$ is the data of
\begin{enumerate}
\item $\forall g \in G$,
$\phi_g:\mathcal{D}\xrightarrow{\sim}
\mathcal{D}$.
\item $\forall g,h\in G$, a natural
isomorphism
$\mathcal{E}_{g,h}:\phi_g\circ \phi_h
\xrightarrow{\sim}\phi_{hg}$ such that
$$
\xymatrix{
\phi_f \circ\phi_g \circ\phi_h
\ar[r]^{\mathcal{E}_{g,h}}\ar[d]_{\mathcal{E}
_{f,g}}&\phi_f
\circ \phi_{hg}
\ar[d]^{\mathcal{E}_{f \circ hg}}\\
\phi_{gf}\circ
\phi_h\ar[r]^{\mathcal{E}_{gf,h}}&\phi_{hgf}
}
$$
\end{enumerate}
\end{definition}

\begin{remark}
\label{geometric-stability-conditions-and-group-actions-remark-more-than-a-homomorphism}
This is more than a group a homomorphism $G
\to \text{Aut}(\mathcal{D})$. Indeed, given
this, $\forall E \in \mathcal{D}$,
$$
\phi_g \circ \phi_h(E)\cong \phi_{gh}(E)
$$
but may not come from a natural isomorphism
of functors
$\phi_g\circ\phi_h\cong\phi_{hg}$.
\end{remark}

\begin{example}
\label{geometric-stability-conditions-and-group-actions-example-group-actions}
\begin{enumerate}
\item $X$: scheme, $G \leq \text{Aut}(X)$.
For every $g \in G$ define
$\phi_g:=g^*
:\text{Coh}(X)\xrightarrow{\sim}\text{Coh}(X)
$.
Then for every $g,h$ there are canonical
isomorphisms
$$
\phi_g \circ\phi_h=g^*\circ
h^*\xrightarrow{\sim}(hg)^* =\phi_{hg}.
$$
$G \mathbb{y} \text{Coh}(X)$ lifts to
$G\mathbb{y}D^b(X)$.

Eg. an Enriques surface, i.e.
$Y=X/\mathbb{Z}/2\mathbb{Z}$, $X$ a K3
surface,
$\mathbb{Z}/2\mathbb{Z}=\left\langle{}i\right\rangle{}$, for
$i:X\to X$ involution, $x
\in X$, $i^*
(\mathcal{O}_x)=\mathcal{O}_{i^{-1}(x)}$.

\item $Q$: acyclic quiver, $G\leq
\text{Aut}^+(Q)=$ automorphisms of the
graph, orientation preserving. Then
$G\mathbb{y}\text{Rep}Q$. Via
$(M_i,\varphi_\alpha)\mapsto
(M_{g(i)},\varphi_{g(x)})$.

Eg. [Missing]
\end{enumerate}
\end{example}

\medskip\noindent
{\bf Equivariantization.}

\begin{definition}
\label{geometric-stability-conditions-and-group-actions-definition-G-equivariant-object}
Let $G \mathbb{y}\mathcal{D}$. A {\it
$G$-equivariant} object is a pair
$(E,\{\lambda_g\}_G$,
\begin{itemize}
\item $E \in \mathcal{D}$,
\item $\lambda_g:E
\xrightarrow{\sim}\phi_g(E)$, a {\bf choice}
of isomorphism
for all $g\in G$, such that
$$
\xymatrix{
E\ar[r]^{\lambda_g}\ar[d]_{\lambda_hg,\sim}&
\phi_g(E)\\
\phi_{hg}(E)&\phi_g(\phi_n(E))\ar[l]
_{\mathcal{E}_{gh}}
}
$$
A {\it morphism of $G$-equivariant objects}
is
$$
(E,\{\lambda_g\}_G)\to(E'_g,\{\lambda'_g\}_G)
$$
is $F \in \Hom_\mathcal{D}(E,E')$ such that
for all $g \in G$,
$$
\xymatrix{
E\ar[d]_F\ar[r]^{\lambda_g}&\phi_g(E)\ar[d]
^{\phi_g(F)}\\
E'\ar[r]^{\lambda'_g}&\phi_g(E')
}
$$
Together these form a category,
$\mathcal{D}_G$, called {\it $G$-equivariant
category}.
\end{itemize}
\end{definition}

\begin{remark}
\label{geometric-stability-conditions-and-group-actions-remark-uses-isomorphisms-Egh}
This uses the isomorphisms
$\mathcal{E}_{g,h}$.
\end{remark}

\begin{example}
\label{geometric-stability-conditions-and-group-actions-example-G-equivariant-objects}
\begin{enumerate}
\item $G\mathbb{y}X$, $G$-equivariant objects
are $G$-equivariant coherent
sheaves.
$$
(\text{Coh}(X))_G=:\text{Coh}_G(X)
$$
In fact, $\text{Coh}_G(X)\cong
\text{Coh}([X/G])$.

If $G$ acts freely, then
$$
\xymatrix{
X\ar[dr]_{\pi}\ar[rr]^g &  &
X\ar[dl]_{\pi}\\
&X/G
}
$$
and $\text{Coh}(X/G)\cong \text{Coh}_G(X)$
via $E \mapsto (\pi^* E,\{\lambda_g\}_{g \in
G}$ where
$$
\xymatrix{
\lambda_g:\pi^* E\ar[dr]_{\lambda_g}\ar[r]&
(\pi \circ g)^* E\ar[d]^=\\
&g^*(\pi^*E)
}
$$
Eg. $Y=X/(\mathbb{Z}/2\mathbb{Z})$ Enriques
surface, $y \in Y$,
$\text{supp}\pi^{-1}(y)=\{x,x'=i^{-1}(x)\}$
$$
\xymatrix@R=.2em{
&\mathcal{O}_x\ar[r]^{\text{id}}&\mathcal{O}
_x\\
\pi^*
(\mathcal{O}_x)=&\oplus&\oplus&=\text{id}
^*(\pi^*(\mathcal{O}_y))\\
&\mathcal{O}_{x'}\ar[r]^{\text{id}}&
\mathcal{O}_{x'}
}
$$
$$
\xymatrix@R=.2em{
&\mathcal{O}_x\ar[r]^{i^*}\ar[ddr]&
\mathcal{O}
_x\\
\pi^*
(\mathcal{O}_y)=&\oplus&\oplus&
=i^*(\pi^*(\mathcal{O}_y))\\
&\mathcal{O}_{x'}\ar[r]^{i^*}\ar[uur]&
\mathcal{O}_{x'}
}
$$
In general,
$$
(D^b(X))_G\cong
D^b(\text{Coh}_G(X))=:D^b_G(X)
$$
\item {\bf Theorem (Demonet '10).}
$G\mathbb{y} Q$ then $\exists Q_G$ such that
$\text{Rep}(Q_G)\cong\text{Rep}(Q)_G$. Eg.
[Missing.]

\item $G \mathbb{y} \mathcal{D}$ for $G$
abelian,
$\hat{G}:=\Hom(G,K^*)\mathbb{y}\mathcal{D}_G$
by
$$
\phi_{\underbrace{x}_{\in
G}}((E,\{\lambda_g\}_G))
:=(E,\{\lambda_g\}_G)\otimes
x=(E,\{\lambda_g\cdot x(g)\}_{g \in G})
$$
Eg. $Y$ Enriques surface \ldots{} [Missing]
\end{enumerate}
\end{example}

\begin{theorem}[Elagin '15]
\label{geometric-stability-conditions-and-group-actions-theorem-Elagin}
$G$ abelian, $G\mathbb{y}\mathcal{D}$
$$
(\mathcal{D}_G)_{\hat{G}}\cong \mathcal{D}
$$
\end{theorem}

\section{Exercise lecture}
\label{geometric-stability-conditions-and-group-actions-section-exercise-lecture}

On a surface $X$, to have a Bridgeland
stability condition we wanted a map
$$
\xymatrix{
Z:K_0(X)\ar[rr]\ar[dr]&  &  \mathbb{C}\\
&  \Lambda\ar[ur]
}
$$
For a numerical B.s.c., we ask
$$
\xymatrix{
Z:K_0(X)\ar[rr]\ar[dr]&  &  \mathbb{C}\\
&  K_{num}\ar[ur]
}
$$
Recall that the {\it Grothendieck group}
$K_0$ can be defined to be the free group
generated by isomorphism calsess of objects
in $D^b(X)$ quotient by the {\it Euler
pairing}, i.e.
$$
[F]=[E]+[G] \iff E \to F \to G \to E[i]
$$
Now let
$$
\chi(E,F)=\sum_{i}(-1)^{i}\dim
\underbrace{\text{Ext}^i(E,F)}
_{=\Hom(E,F[i])}
$$
then
$$
K_{num}=K_0(X)/\underbrace{\text{null}(\chi)}
_{=\{E:\chi(E,F)=0\forall F\in K_0(X)\}}
$$
To obtain a lattice consider
$$
\xymatrix{
K_0(X)\ar[d]\ar[r]^{ch}& \CH^* (X)\ar[d]\\
K_{num}(X)\ar[r]^{ch}&\underbrace{\CH^*_{num}
(X)}
_{\simeq \underbrace{\mathbb{Z}}_{ch_0}
\oplus \text{NS}(X) \oplus
\frac{1}{2}\underbrace{\mathbb{Z}}_{ch_2}}
}
$$
Fix $B \in \text{NS}(X)_{\mathbb{R}}$.
$$
\text{NS}(X)\ni A \mapsto \text{deg}(A\cdot
B)
\in
c\cdot \mathbb{Z}
$$
for $c \in \mathbb{R}$. For a surface,
$$
\Lambda = \mathbb{Z} \oplus
\mathbb{Z}^{\oplus \text{rk}(\text{NS}(X))}
\oplus \frac{1}{2}\mathbb{Z}
$$
$$
Z_{H,D,\alpha,\beta}=(\alpha+i \beta)H^2
ch_0+(D+iH)ch_1-ch_2.
$$

\begin{exercise}
\label{geometric-stability-conditions-and-group-actions-exercise-Ox-is-simple}
$\mathcal{O}_X$ is simple (i.e. no
subobjects) in $\text{Coh}^{H,\beta}(X)$.
\end{exercise}

% --- End geometric-stability-conditions-and-group-actions.tex ---

% --- Begin hilbert-schemes.tex ---
\chapter{Hilbert schemes}
\label{chapter-hilbert-schemes}


\author{Daniel González Casanova Azuela}

\noindent
This notebook starts with the Hilbert
scheme and its description through families.
It then treats group-invariant subschemes,
generalized Kummer varieties, and seminar
notes on nested Hilbert schemes.
The later sections preserve the material
on Nakajima's Heisenberg representation.


\section{Hilbert scheme of points}
\label{hilbert-schemes-section-hilbert-scheme}

\noindent
In this section we define
the Hilbert scheme of points
on a smooth projective complex surface $X$.
The essential property of the Hilbert scheme
is that its points (over $\mathbb{C}$)
are 0-dimensional closed subschemes of $X$
of length $n$.

\medskip\noindent
Consider a contravariant functor
$$
\mathcal{H}ilb_X:(\text{Schemes})^{\text{op}}\to (\text{Sets})
$$
which assigns to any scheme $U$ the set of
flat families
of closed subschemes in $X$ parametrized by
$U$.
Any such family is by definition a flat
morphism
$$
\xymatrix@C=.1em{
Z \ar[d]^\pi&\subseteq X \times U\\
U.
}
$$
We define $\mathcal{H}ilb$ to map a morphism
of schemes
$U \to U'$ to the map of sets (of families
over $U$ and $U'$) given
by pullback. That is, we map a flat family
$Z'\to U'$
to $Z'\times_{U'}U \to U$, which remains
flat.

The Hilbert polynomial at a fibre $Z_u$ is
defined by
$P_u(m)=\chi(\mathcal{O}_{Z_u}\otimes
\mathcal{O}_X(m))$.
By flatness of $\pi$ the Hilbert polynomial
is constant along $U$
(\cite[Part III, Theorem 9.9]{Hartshorne}).
For a fixed polynomial $P$
we define $\mathcal{H}ilb^P_X$ as the
functor which maps $U$
to the set of families over $U$ with Hilbert
polynomial $P$.
A theorem by Grothendieck says that
$\mathcal{H}ilb^P_X$ is representable
by a projective scheme $\text{Hilb}^P_X$.
This means that there is a natural
transformation
$\Hom(-,\text{Hilb}^P_X) \simeq
\mathcal{H}ilb^P_X(-)$;
that is, families over any scheme $U$ are in
bijection with
morphisms $U \to \text{Hilb}^P_X$.

As with any representable functor,
we obtain a distinguished element associated
to the identity morphism
of the object representing the functor.
In our case, this is a
``universal'' family $\mathcal{Z} \to
\text{Hilb}^P_X$
with the property that any other family $Z
\to U$ is obtained as a pullback
of $\mathcal{Z}$.
More exactly, when regarding an arbitrary
family $Z \to U$ as a morphism
$f \in \Hom(U,\text{Hilb}^P_X)$, we see that
$$
\xymatrix{
\text{id}_{\text{Hilb}^P_X}\ar@{|->}[d]
&\Hom(\text{Hilb}^P_X,\text{Hilb}^P_X)\ar[r]^-{\simeq}\ar[d]_{f^*}
&\mathcal{H}ilb^P_X(\text{Hilb}^P_X)\ar[d]^{U\times_{\text{Hilb}_X^P}-}
&\mathcal{Z}\ar@{|->}[d]\\
f&\Hom(U,\text{Hilb}^P_X)\ar[r]^-{\simeq}
 &\mathcal{H}ilb^P_X(U)&
Z=U\times_{\text{Hilb}_X^P}\mathcal{Z}.
}
$$
commutes by naturality.
This shows that studying families of
subschemes of $X$ parametrized by any
scheme $U$ can be done by studying the
universal family
over the Hilbert scheme.

We can formally view subschemes of $X$
as points of $\operatorname{Hilb}_X^P$ as
follows.
Recall that, as with any other scheme,
points $h$ of $\text{Hilb}_X^P$ correspond
to morphisms
$\Spec(\kappa(h)) \to \text{Hilb}^P_X$,
where $\kappa(h)$ is the residue field of
the point
(see Stacks Project
\href{https://stacks.math.columbia.edu/tag/01J9}{01J9}).
Then for every point $h$ of the Hilbert
scheme
we have a unique subscheme
of $X$ defined as the fibre $X_h$
of the corresponding morphism,
$$
\xymatrix{
X_h\ar[r]\ar[d]\ar@{}[dr]|<{\lrcorner}&\mathcal{Z}\ar[d]
& \subset X \times \text{Hilb}^P_X\\
\Spec(\kappa(h))\ar[r]&\text{Hilb}^P_X.
}
$$
As a rather technical observation that will
not be
addressed later, we note that
to make sure that the fibres are
$\mathbb{C}$-schemes,
we must restrict our analysis to
$\mathbb{C}$-points
of the Hilbert scheme (i.e. when the residue
field is $\mathbb{C}$).

\begin{definition}
\label{hilbert-schemes-definition-hilbert-scheme-of-points}
Let $P$ be a constant polynomial given by
$P(m)=n$ for all $m \in \mathbb{Z}$.
The {\it Hilbert scheme of $n$ points in
$X$} is $X^{[n]}:=\text{Hilb}^P_X$.
\end{definition}

\noindent
This condition forces the dimension of the
fibres of any family
over $X$ to be zero. Indeed, the leading term
of the Hilbert polynomial of a projective
variety of degree $d$
and dimension $s$ is $(dm^s)/s!$, which can
only be constant
for $s=0$ (see \cite[Exercise
1.13]{HarrMorr}).

Further, any 0-dimensional subscheme of $X$
is supported on finitely many points of $X$
(this is consequence of $X$ being Noetherian:
we couldn't have infinitely many points
since they
would give an infinite chain of prime
ideals).
Finally, the {\it length} of any such
subscheme, defined by
 $\dim\Gamma(\mathcal{O}_Z)$
must be $\sum_{x \in \text{supp}
Z}\dim\Gamma(\mathcal{O}_{Z_x})=n$.
(see Stacks Project
\href{https://stacks.math.columbia.edu/tag/0AYT}{0AYT}).

\medskip\noindent
It is now easy to see that any subscheme of
$X$ consisting of $n$ distinct
points is itself a point of
$\text{Hilb}_X^P$.
Indeed, if $Z=\{x_1,\ldots,x_n\}\subset X$,
its
structure sheaf is the direct sum of
skyscraper sheaves
$$
\mathcal{O}_Z=\bigoplus_{i=1}^n
(\text{skyscraper sheaf at $x_i$})
$$
so that $\mathcal{O}_Z \otimes
\mathcal{O}_X(m)\cong\mathcal{O}_Z$
and thus the Hilbert polynomial of $Z$ is
constant $n$.
The latter isomorphism can be verified via
the exact sequence associated to the maximal
ideal
$\mathfrak{m}_{x_i}$
of the stalk $\mathcal{O}_{X,x_i}$, namely
$$
\xymatrix{
0\ar[r]
&\mathfrak{m}_{x_i}\ar[r]
&\mathcal{O}_{X,x_i}\ar[r]
&\mathcal{O}_{X,x_i}/\mathfrak{m}_{x_i}\cong\mathbb{C}\ar[r]
&0,
}
$$
which after tensoring by
$\mathcal{O}_{Z,x_i}
\cong\mathbb{C}\cong\mathcal{O}_{X,x_i}/\mathfrak{m}_{x_i}$
kills the first term.
Twisting by $m$ does not alter this
construction.

\medskip\noindent
However, not every point in $X^{[n]}$ is
given
by a set of $n$ distinct points,
nor by a set of points with multiplicites,
as the following construction suggests.
Define the {\it symmetric power} of $X$ as
$$
S^nX=\underbrace{X \times \ldots \times
X}_{n\text{ copies}}/S_n
$$
where the symmetric group $S_n$ acts by
permutating copies. Thus, a point of $S^nX$
is
an $S_n$-orbit on the cartesian product $X
\times \ldots \times X$,
i.e., an unordered $n$-tuple of points of $X$
with allowed repetition.
In algebraic geometry these are called
``algebraic cycles''
and are denoted as formal sums $\sum_{x \in
X}n_x[x]$
where $\sum_{x \in X}n_x=n$, with $n_x \in
\mathbb{N}$,
and $[x]$ are formal symbols.

By our discussion above on the length of
0-dimensional subschemes of $X$,
we have a map
$$
\begin{aligned}
\pi : X^{[n]} &\longrightarrow S^nX \\
 Z&\longmapsto \sum_{x \in \text{supp}Z}\dim
\Gamma(\mathcal{O}_{Z,x})[x]
\end{aligned}
$$
called the {\it Hilbert-Chow morphism}.
In some cases this is an isomorphism, but
not always.

\begin{theorem}[Fogarty, \cite{Fogarty}]
\label{hilbert-schemes-theorem-fogarty}
Suppose $X$ is nonsingular and $\dim X=2$,
then the following holds.
\begin{enumerate}
\item $X^{[n]}$ is nonsingular and has
dimension $2n$.
\item The Hilbert-Chow morphism is a
resolution of singularities.
\end{enumerate}
\end{theorem}




\section{The G-Hilbert scheme}
\label{hilbert-schemes-section-g-hilb}

\noindent
Start with the Hilbert functor of
Section~\ref{hilbert-schemes-section-hilbert-scheme}:
a base $B$ parametrizes finite flat families
$Z\subset X\times B$ of fixed length.
For a finite group $G$ acting on $X$, take
length $|G|$ and require $Z$ to be
$G$-invariant, with $G$ acting trivially on
$B$. There is one extra condition: on each
fibre, the action on its functions must be
isomorphic to the action on functions on a
free $G$-orbit (an orbit of $|G|$ distinct
points). This is the regular representation
condition; invariance and length alone are
not sufficient. The resulting functor is
represented by $G$-Hilb$(X)$. Its component
containing free orbits records their flat
limits, including collisions, and maps to
$X/G$ by sending a subscheme to the orbit
supporting it. This is why a Hilbert scheme
can improve a quotient: it retains limiting
subscheme information that the quotient
forgets. Whether this map resolves the
singularities requires a theorem for the
action in question. See
Bridgeland--King--Reid,
\emph{Mukai implies McKay}, Section 1,
\href{https://arxiv.org/abs/math/9908027}
{arXiv:math/9908027}.

\section{Generalized Kummer varieties}
\label{hilbert-schemes-section-generalized-kummer}

\noindent
Let $A$ be an abelian surface. Recall from
Section
\ref{hilbert-schemes-section-hilbert-scheme}
the Hilbert-Chow morphism
$
\rho:A^{[n+1]}\longrightarrow S^{n+1}A.
$

\noindent
Since $A$ is a group, there is a natural
summation map
$$
\begin{aligned}
\sigma:S^{n+1}A&\longrightarrow A,\\
\sum_i m_i[a_i]&\longmapsto
\sum_i m_i a_i.
\end{aligned}
$$
Here the sum on the right is taken using the
group law of the abelian surface. The map is
well-defined because addition on $A$ is
commutative, so it does not depend on the
ordering of the points.

\begin{definition}
\label{hilbert-schemes-definition-generalized-kummer}
The {\it generalized Kummer variety}
associated to $A$ is
$
K_n(A)
=
(\sigma\circ\rho)^{-1}(0)
\subset A^{[n+1]}.
$
\end{definition}

\begin{theorem}[Beauville]
\label{hilbert-schemes-theorem-generalized-kummer-hk}
Let $A$ be an abelian surface and let
$n\geq1$. Then $K_n(A)$ is a smooth
projective irreducible holomorphic symplectic
variety of dimension $2n$. In particular it
is a hyperk\"ahler variety.
\end{theorem}

\begin{reference}
This is one of the two standard families
constructed in \cite{beauville-kummer}.
\end{reference}

\begin{example}
\label{hilbert-schemes-example-kummer-fourfold}
For $n=2$, consider the zero-sum locus
$
Z=
\{(a_0,a_1,a_2)\in A^3:
a_0+a_1+a_2=0\}.
$
The map
$
A^2\longrightarrow Z,
\qquad
(a,b)\longmapsto(a,b,-a-b)
$
is an isomorphism. The symmetric group
$S_3$ acts on $Z$ by permuting the three
coordinates, and the zero fibre of
$
S^3A\longrightarrow A
$
is
$
Z/S_3\cong A^2/S_3.
$
The Hilbert-Chow morphism restricts to a
resolution
$
K_2(A)\longrightarrow A^2/S_3.
$
Thus $K_2(A)$ is a hyperk\"ahler fourfold
resolving the singular zero-sum quotient.
\end{example}

\section{Irreducibility of the Hilbert scheme
of points and the class of 2-step
ideals}
\label{hilbert-schemes-section-graffeo-two-step-ideals}

\noindent
Michele Graffeo, SISSA.
Algebraic Geometry Seminar, UFF.
September 29, 2025.

\medskip\noindent
These are rough notes from the talk,
relocated from the algebraic geometry
seminar notebook. Their incomplete formulas
and theorem statements still need review.
They are preserved as a reading record,
not as verified statements.

\medskip\noindent
{\bf Abstract.}
Hilbert schemes of points on a
quasi-projective variety X are classical
objects
in algebraic geometry. Roughly speaking, they
parametrise ideals of a polynomial
ring with complex coefficients having finite
colength. Although Hilbert schemes
always have a distinguished component called
the smoothable component, their
geometry is quite pathological and one of the
main open problems around them
concerns their irreducibility. In a joint
work with Giovenzana, Giovenzana,
Lella we introduce the notion of 2-step
ideals. We show that the loci
parametrising these ideals are in general not
contained in the smoothable
component of the Hilbert scheme, thus
providing new examples of
extra-components.  In my seminar I will
discuss our class of ideals and relate
them to the compressed algebras considered by
Iarrobbino in the eighties.
Finally, I will show how to extend our result
to the nested setting.


\medskip\noindent

As usual, you start with a functor and
prove representativity.

\begin{definition}
\label{hilbert-schemes-definition-hilbert-scheme}
Let $X$ be a smooth
quasi-projective variety and a nonegative
integer $d$.
Define
$$
\begin{aligned}
\underline{Hilb}^d: \text{Sch}_\mathbb{C}
&\longrightarrow \mathsf{Set} \\
B &\longmapsto \left\{Z \hookrightarrow B
\times X:
\substack{B\text{-flat} \\ B\text{-finite}\\
\text{length }d}\right\}
\end{aligned}
$$
\end{definition}

\begin{theorem}[Grothendieck]
\label{hilbert-schemes-theorem-grothendieck}
$\underline{Hilb}^d(X)$ is represented by a
quasi-projective
scheme $Hilb^d(X)$
\end{theorem}

Idea: $\mathbb{C}$-points of $Hilb^d(X)$
are in correspondence with $Z
\overset{\subset}{\text{finite}}X$
of length.

A fat point $Z= \text{Spec}(A)$ $(A,
\mathfrak{m})$ is a
local artinian $\mathbb{C}$-algebra of finite
type.

If $X$ smooth can consider $X=\mathbb{A}^n$,
and it's the same to consider
\begin{itemize}
\item $Z \hookrightarrow \mathbb{A}^n$ of
length $d$.
\item $I \subset \mathbb{C}[x_1,\ldots,x_n]$
of
colength $d$.
\item $A= \mathbb{C}[x_1,\ldots,x_n]/I$ with
$\dim_\mathbb{C}=d$.
\end{itemize}

\begin{definition}
\label{hilbert-schemes-definition-nested-hilbert-scheme}
$X$, $0 \leq d_1 \leq \ldots
d_r=\underline{d}$.
$$
Hilb^{\underline{d}}(X)=
\{Z_1\not\hookrightarrow
\ldots
\not\hookrightarrow Z_r  \not\hookrightarrow
X:
\text{len}(Z_i)=d_i\}
$$
\end{definition}

\begin{theorem}[Hartshorne (r=0), Fogarty,
Kalpan (r>1)]
\label{hilbert-schemes-theorem-hk}
$Hilb^d(X)$ is connected.
\end{theorem}

For $r=1$, $Hilb^d(\mathbb{A}^n)$ is smooth
iff
$n \leq 2$, $d \leq 3$.
If $r>1$, $n=2$, it is smooth iff $r=2$ and
$d_1=d_2=-1$.
If $r>1$ and $n>2$, it is smooth iff $r=2$,
$(d_1,d_2)=\{(1,2),(2,3)\}$.

Iror components.
\begin{itemize}
\item $(r=1)$ $n\geq 4$,
$Hilb^d(\mathbb{A}^3)$
is irreducible iff $d\leq 7$ ($\impliedby$
[E-I], $\implies $ [Mozzola].
$n=3$ irredicuble if $d \leq  11$ (8,9 Sivic;
10 Jardim et al; 11, Jelisievv et
al.)
\item $(r=2)$, $n=2$, irreducible
[G-Rosul-Sebastian].
\item $(r \geq 3)$ $n=2$, there exist $d_1
\leq \ldots \leq d_5$
such that
$Hilb^{\underline{d}}(\mathbb{A}^i)$ is
irreducible [S-R].
\end{itemize}

There is an analogy between classical in
dimension 3
and nested in dimension 2.

Schematic structure.
\begin{itemize}
\item [I] $Hilb^{21}(\mathbb{A}^4)$ has at
least a generically
non-reduced component.
\item Problem: what about $n=3$?
\item Theorem. If $Hilb^{d_1,\ldots
d_r}(\mathbb{A}^i)$ is
irreducible $\implies $
$Hilb^{(1,d_1,\ldots,d_r)}(\mathbb{A}^n$
has a generically non reduced component.
\end{itemize}

Other results.
\begin{itemize}
\item …
\end{itemize}
Hilbert scheme function.
For a local artinian $\mathbb{C}$-algebra of
finite type there is an associated graded
$\text{Gr}_\mathfrak{m}(A)
=\bigoplus_{i \geq
0}\mathfrak{m}^1/\mathfrak{m}^{i+1}$
which allows us to define the Hilbert
function
$$
\begin{aligned}
h_A: \mathbb{Z} &\longrightarrow \mathbb{N}
\\
i &\longmapsto \dim_\mathbb{C}
\mathfrak{m}_i/\mathfrak{m}_{i+1}
\end{aligned}
$$
sistability for graded $A$-modules.
The socol is the annihilator of
the ideal which gives the algebra
e.g. $\mathbb{C}[x,y](x^2,xy,y^2)$ then
$\text{Soc}(A)=\left\langle{}x,y^2\right\rangle{}_\mathbb{C}
$.

\medskip\noindent
\begin{definition}
\label{hilbert-schemes-definition-smoothable-component}
The {\it smoothable component} is
 $$
V_{sm}=\overline{\{[(z_i)_{i=1}^n \in
Hilb^{\underline{d}}(\mathbb{A}^n):
Z_r\text{ is reduced}\}}
$$
\end{definition}

\begin{definition}
\label{hilbert-schemes-definition-elementary-componenet}
$V \subset
Hilb^{\underline{d}}(\mathbb{A}^n)$ is
an {\it elementary component} if
$\forall [Z_1,\ldots,Z_r] \in V$
then $Z_r$ is a fat point.
\end{definition}

Instead of looking for all components of the
Hilbert scheme
we look for elementary components. This may
not be some simple;
sometimes there's infinitely many. There
is no method to find them.

\begin{theorem}[Irrabaro]
\label{hilbert-schemes-theorem-i}
Every irreducible component $V \subset
Hilb^d(\mathbb{A}^n)$
is generically étale locally product of
elementary components.
\end{theorem}

\begin{definition}
\label{hilbert-schemes-definition}
Let $h:\mathbb{Z} \to \mathbb{N}$ be a
function
with finite support, $|h|=\sum_{ i \in
\mathbb{Z}}h(i)=d$. Define
\begin{itemize}
\item $H_n=\{[A] \in Hilb^d(\mathbb{A}^n):
h_A \equiv h\}$.
This is closed by semicontinuity of the
function; but can be expressed
a by a representable functor.
\item The following also has a canonical
schematic structure
$$
\begin{aligned}
\pi_n:H_n  &\longrightarrow
\mathcal{H}_n=\{[A] \in H_n: A\text{ is
graded}\} \\
A &\longmapsto A=\text{Gr}_{\mathfrak{m}}(A)
\end{aligned}
$$
Where the first is $\not\hookrightarrow $ in
the second.
\end{itemize}
\end{definition}

Recall that
$$
T_{[I]}Hilb^d(\mathbb{A}^n)=\Hom_R(I,R/I)
$$
\begin{theorem}[B-B,J,G G G L]
\label{hilbert-schemes-theorem-b-b-j-g-g-g-l}
There is a graded $[I] \in \mathcal{H}_n$.
Then
$$
\begin{aligned}
T_{[I]}\mathcal{H}_n&=\Hom_R(I,R/I)_{=0}\\
T_{[I]}\pi_n([I])&='>0\\
T_{[I]}H_n&=' \geq 0.
\end{aligned}
$$
\end{theorem}

\begin{definition}
\label{hilbert-schemes-definition-compressed}
$(A,\mathfrak{m}_A)$ is {\it compressed} is
for all  $(A',\mathfrak{m}')$
with $e(A)=e(A')$ and $len(A) \geq  len(A')
$.
\end{definition}

\begin{theorem}
\label{hilbert-schemes-theorem-e}
$E_h \subset H_n$ is $H$-compressed
$E_h \neq  \emptyset$
\end{theorem}

\section{Nakajima's representation: overview}
\label{hilbert-schemes-section-intro}
\label{hilbert-schemes-section-abs}

\noindent
The main theorem in \cite{Nakajima-Heis}
shows there is an action of the
infinite-dimensional
Heisenberg algebra
on the homology groups of the Hilbert
scheme of points on a projective surface.
We explain how this representation is
constructed
and give an idea of the proof of the theorem.

\noindent
Motivated by earlier work on moduli spaces
of
torsion-free sheaves
(see \cite{Nakajima-ALE}),
Nakajima constructed in \cite{Nakajima-Heis}
a representation the infinite-dimensional
Heisenberg
algebra on the homology groups of the
Hilbert schemes
on a nonsingular projective surface.

In this work we outline the basic ideas in
Nakajima's construction
and explain the first part of the proof of
the main theorem.
Most of this work is based in
\cite{Nakajima-lectures},
with the exception of the section on
Heisenberg algebras
where I try to identify the precise version
of the Heisenberg
algebra that is being represented.
In Section \ref{hilbert-schemes-section-hilbert-scheme},
I explain the construction of the Hilbert
scheme,
providing slightly more details than what is
really essential
to understand the main theorem,
but which help to get a better idea on what
the Hilbert scheme is
and how to work with it.

This work is strongly influenced by
\cite{videos},
a series of seminars aimed at understanding
\cite{Nakajima-lectures}
documented in videos.
This was of great help
specially in Section \ref{hilbert-schemes-section-bm}, since
the properties of Borel-Moore homology are a
rather technical
but essential part of the construction.
The remaining sections, containing the main
theorem
and part of its proof,
also pick up ideas from these videos,
but are mostly based in
\cite{Nakajima-lectures}.


\section{Heisenberg algebra}
\label{hilbert-schemes-section-heisenberg}

\noindent
In this section we review four different
variants
of the Heisenberg algebra and discuss their
characters.
We conclude by showing how the character is
used
to prove that Nakajima's representation is
of highest weight.

\medskip\noindent
The most basic definition of the Heisenberg
Lie algebra is
$\hat{\mathfrak{a}}=
\bigoplus_{n \in \mathbb{Z}} \mathbb{C}h_n
\oplus\mathbb{C}K$
with bracket given by
\begin{equation}
\label{hilbert-schemes-equation-bracket-basic}
[h_m,h_n]=m\delta_{m,-n}K,\qquad
[K,\hat{\mathfrak{a}}]=0.
\end{equation}

\noindent
In class we described its highest weight
representation
$H=\mathbb{C}[x_1,x_2,\ldots]$ with
\begin{align*}
h_n&\mapsto \begin{cases}
n \frac{\partial }{\partial x_n}\qquad
&\text{if }n>0\\
x_{-n}\qquad &\text{if }n<0\\
0\qquad &\text{if }n=0
\end{cases},
\qquad \qquad  K\mapsto \text{Id}.
\end{align*}

\noindent
We introduced a grading in
$H=\bigoplus_{\Delta=0}^\infty H_\Delta$
given by the {\it energy} of monomials,
defined by
$$
\Delta(x_1^{k_1}x_2^{k_2}\ldots
x_s^{k_s})=\sum_j k_j,
$$
which gives the character formula
\begin{equation}
\label{hilbert-schemes-equation-character-heisenberg}
\sum_{\Delta=0}^\infty q^\Delta \dim H_\Delta
=\prod_{m=1}^\infty\frac{1}{1-q^m}.
\end{equation}

\noindent
In order to generalize this construction
it's worth going over why this formula holds.
Recall that the \textit{Verma module} of
weight $\mu=0$
of $\hat{\mathfrak{a}}$ is defined by
$H=
\mathcal{U}(\hat{\mathfrak{a}})
\otimes_{\mathcal{U}(\hat{\mathfrak{a}}_+)}
\mathbb{C}|0\rangle$.
Further, by the PBW theorem, this module
is isomorphic to
$\mathcal{U}(\mathfrak{n}_-)$,
the universal enveloping algebra
of the negative roots in the triangular
decomposition
(as in the finite-dimensional case, see
\cite[Lecture 24]{KLAL} and
\cite[Video 8]{videos}).
Since the subalgebra of negative roots is
abelian,
the universal enveloping algebra becomes the
symmetric algebra
and therefore we have
\begin{equation}
\label{hilbert-schemes-equation-symmetric-rank-1}
\begin{aligned}
H&\cong\text{Sym}\left(\bigoplus_{m>0}\mathbb{C}at^{-m}\right)\\
&\cong\bigotimes_{m>0}\text{Sym}(\mathbb{C}at^{-m})\\
&\cong \bigotimes_{m>0}\mathbb{C}[at^{-m}],
\end{aligned}
\end{equation}

\noindent
since the symmetric power of a vector space
is the polynomial ring in its generators,
and the symmetric power of a direct sum
is the tensor product of the symmetric
powers of the factors.
This matches our previous formula
$H=\mathbb{C}[x_1,x_2,\ldots]$ by letting
$x_m=at^{-m}$.

Putting a dummy variably $q$ to every
1-dimensional
factor and multiplying throughout, we obtain
$$
\begin{aligned}
m=1\qquad &(1+q+q^2+q^3+\ldots)\times\\
m=2\qquad &(1+q^2+q^4+q^6+\ldots)\times\\
m=3\qquad
&(1+q^3+q^6+q^9+\ldots)\times\ldots,
\end{aligned}
$$

\noindent
so that indeed the character is given by
Equation \ref{hilbert-schemes-equation-character-heisenberg}.

\medskip\noindent
Another version of this construction
we discussed in lectures
(see \cite[Section
3.1]{Jethro-vertex-algebras})
starts with any finite-dimensional vector
space $\mathfrak{h}$
equipped with a symmetric bilinear form
$(\cdot,\cdot)$.
Define
$$
\hat{\mathfrak{h}}=\mathfrak{h}\otimes\mathbb{C}[[t^{\pm 1}]]\oplus\mathbb{C}K
$$
with bracket
\begin{equation}
\label{hilbert-schemes-equation-bracket-affine}
[at^m,bt^n]=m\delta_{m,-n}(a,b)K,\qquad
[K,\hat{\mathfrak{h}}]=0,
\end{equation}

\noindent
which is just the affine Lie algebra of
$\mathfrak{h}$ when
regarded as a commutative Lie algebra.

A similar formula to Equation
\ref{hilbert-schemes-equation-symmetric-rank-1} is
\begin{equation}
\label{hilbert-schemes-equation-symmetric}
\begin{aligned}
H&\cong\text{Sym}\left(\bigoplus_{m>0}\bigoplus_{a \in B}
\mathbb{C}at^{-n}\right)\\
&\cong\bigotimes_{m>0}\bigotimes_{a \in
B}\text{Sym}(\mathbb{C}at^{-m})\\
&\cong\bigotimes_{m>0}\bigotimes_{a \in
B}\mathbb{C}[at^{-m}]
\end{aligned}
\end{equation}
where $B$ is a basis of the underlying
vector space $\mathfrak{h}$.
If $\mathfrak{h}$ has dimension $r$, this
has character
\begin{equation}
\label{hilbert-schemes-equation-character-rank-r}
\prod_{m=0}^\infty\frac{1}{(1-q^m)^r}.
\end{equation}

\medskip\noindent
Taking this construction one step further,
let's change $\mathfrak{h}$ for a
finite-dimensional
super vector space
$\mathfrak{h}=\mathfrak{h}_0\oplus
\mathfrak{h}_1$
with basis $B^{\text{even}}\amalg
B^{\text{odd}}$.
If we equip $\mathfrak{h}$ with a graded
commutative
pairing $\left\langle{}\cdot,\cdot\right\rangle{}$, we can
apply the same construction as above
to obtain the algebra
\begin{equation}
\label{hilbert-schemes-equation-super-heisenberg}
\hat{\mathfrak{h}}
=(\mathfrak{h}_0[[t^{\pm
1}]]\oplus\mathfrak{h}_1[[t^{\pm 1}]])
\oplus\mathbb{C}K
\end{equation}

\noindent
with Lie bracket
\begin{equation}
\label{hilbert-schemes-equation-super-bracket}
[at^m,bt^n]=\left\langle{}a,b\right\rangle{}m\delta_{m,-n}K,\qquad
[K,\hat{\mathfrak{h}}]=0.
\end{equation}

\noindent
One can check that this is indeed a super
Lie bracket
(which essentially holds by the pairing
being graded commutative
and $K$ being central),
i.e. $\hat{\mathfrak{h}}$ is a Lie
superalgebra.
Its highest weight module is given by
$$
\begin{aligned}
H&\cong
\text{Sym}\left(\bigoplus_{m>0}\bigoplus_{a
\in B}
\mathbb{C}at^{-m}\right)\\
&\cong
\bigotimes_{m>0}\text{Sym}\left(\bigoplus_{a
\in B^{\text{even}}}
\mathbb{C}at^{-m}\right)
\otimes\bigotimes_{m>0}
\Lambda\left(\bigoplus_{a \in B^{\text{odd}}}
\mathbb{C}at^{-m}\right).
\end{aligned}
$$

\noindent
since the universal enveloping algebra of
the odd
part is the exterior power (see \cite[Video
8]{videos}).

Now recall that while the symmetric power of
a vector
space yields the polynomial algebra on the
basis of the vector space,
the exterior power yields the algebra
generated
by the wedge products of the basic vectors.
In the case of a 1-dimensional vector space
generated
by the element $at^{-m}$, we get
$$
\Lambda(\mathbb{C}at^{-m})=\mathbb{C} \oplus
\mathbb{C}at^{-m}
$$
and thus the character contribution in this
case is simply
$
1+q^m.
$
Multiplying throughout we obtain the
character
$$
\prod_{m=0}^\infty \frac{(1+q^m)^{\dim
\mathfrak{h}^{\text{odd}}}}{
(1-q^m)^{\dim \mathfrak{h}^{\text{even}}}}.
$$

\medskip\noindent
Let's introduce one last variant of this
construction.
This time we start from a graded vector space
$\mathfrak{h}=\bigoplus_{m \in
\mathbb{Z}}\mathfrak{h}_m$,
which is automatically also a supervector
space
by the $\mathbb{Z}_2$-grading
$\mathfrak{h}^{\text{even}}=\bigoplus_{i \in
\mathbb{Z}}\mathfrak{h}_{2i}$
and $\mathfrak{h}^{\text{odd}}=\bigoplus_{i
\in \mathbb{Z}}\mathfrak{h}_{2i+1}$.
Suppose $\mathfrak{h}_i$ is
finite-dimensional with basis $B_i$,
and that $\mathfrak{h}$ is equipped with a
graded commutative
pairing $\left\langle{}\cdot,\cdot\right\rangle{}$.
The definition of $\hat{\mathfrak{h}}$ as an
algebra
is given exactly as in Equation
\ref{hilbert-schemes-equation-super-heisenberg}
with bracket given by Equation
\ref{hilbert-schemes-equation-super-bracket}.

Now we obtain
$$
\begin{aligned}
H&\cong
\text{Sym}\left(\bigoplus_{m>0}\bigoplus
_{\substack{a \in B_i \\ i \in \mathbb{Z}}}
\mathbb{C}at^{-m}\right)\\
&\cong
\bigotimes_{m>0}\text{Sym}\left(\bigoplus
_{\substack{a \in B_{2i} \\ i \in
\mathbb{Z}}}
\mathbb{C}at^{-m}\right)
\otimes\bigotimes_{m>0}
\Lambda\left(\bigoplus_{
\substack{a \in B_{2i+1} \\ i \in
\mathbb{Z}}}
\mathbb{C}at^{-m}\right).
\end{aligned}
$$

\noindent
This gives the character
$$
\sum_{\substack{\Delta\in \mathbb{Z}_{\geq
0}  \\ i \in \mathbb{Z}}}
\dim \mathfrak{h}_{i,\Delta}t^iq^\Delta
=\prod_{m>0}\prod_j(1-(-1)^jt^jq^m)^{-(-1)^j\dim \mathfrak{h}_j}.
$$

\noindent
Finally fix $\mathfrak{h}=H_*(X)$, the
singular homology ring
of a complex projective surface $X$.
Then the character of the corresponding
Heisenberg algebra is
\begin{equation}
\label{hilbert-schemes-equation-character}
\prod_{m=1}^\infty \frac{
(1+t^{2m-1}1^m)^{b_1(X)}
(1+t^{2m+1}q^m)^{b_3(X)}}
{(1-t^{2m-2}q^m)^{b_0(X)}
(1-t^{2m}q^m)^{b_2(X)}
(1-t^{2m+2}q^m)^{b_4(X)}}.
\end{equation}

\noindent
It was proved by Göttsche using the Weil
conjectures in \cite{Gottsche1990}
(and in \cite[Section
6.2]{Nakajima-lectures} by other methods)
that Equation \ref{hilbert-schemes-equation-character}
in fact equals
$$
\sum_{n=0}^\infty q^nP_t(X^{[n]}),
$$
where the {\it Poincaré polynomials}
$P_t(Z)$ are defined by
$P_t(Z)=\sum_{m\geq 0}b_m(Z)t^m$,
$b_m(Z)=\text{rank}H_m(Z)$ for
any variety $Z$.

The representation
of the Heisenberg algebra we discuss in this
work corresponds
precisely to this version of the Heisenberg
algebra.
Göttsche's formula shows that it is a
highest weight representation,
as explained in \cite[Corollary
8.16]{Nakajima-lectures}.




\section{Borel-Moore homology}
\label{hilbert-schemes-section-bm}

\noindent
In this section we define Borel-Moore (BM)
homology
and formulate its basic properties.
Most of this section is based on \cite[Video
9]{videos}.
For further details on the construction see
\cite{Fulton-Young-Tableaux}.

\medskip\noindent
Let $M$ be a a smooth, connected,
oriented, finite-dimensional manifold over
$\mathbb{R}$.
While the following results hold for any
topological space $X \subset M$
that is a proper deformation retract
of a closed neighbourhood of a manifold,
for us it's enough to think of $X \subset M$
as a
(possibly singular) variety over
$\mathbb{C}$.

In this work we consider homology and
cohomology
groups with coefficients in $\mathbb{C}$.
Recall that singular homology is defined
as the homology of the chain complex of
finite linear combinations of $i$-simplices
in $X$.
\textit{Borel-Moore homology}
$H_\bullet^{BM}(X)$
is the homology of the complex of locally
finite
linear combinations of singular simplices in
$X$,
where ``locally finite'' means that every
point $x \in X$ is in a neighbourhood
that intersects only finitely many
simplices.

A striking
difference between singular and BM homology
is that the latter is not homotopy invariant.
To see this recall that the homology groups
of  $\mathbb{R}^2$ are
$H_0(\mathbb{R}^2)=\mathbb{C}$
and $H_i(\mathbb{R}^2)=0$ for $i \neq 0$.
However, $H_2^{BM}(\mathbb{R}^2)=\mathbb{C}$
since a tiling of triangles
is a locally finite chain of 2 simplices,
has no boundary (since the boundaries of
all triangles cancel each other),
and is clearly not the boundary of a
3-simplex;
thus this tiling is a nonzero element of
$H_2^{BM}(\mathbb{R}^2)$.
Further, notice that
$H_0^{BM}(\mathbb{R}^2)=0$
since a ray of 1-simplices
is a locally finite chain
whose boundary is a point;
thus a point is a boundary and
has trivial homology class.

\medskip\noindent
Here are three important properties of BM
homology:
\begin{enumerate}
\item If $X$ is compact, $H_i^{BM}(X)\cong
H_i(X)$.
\item If $X$ is not compact,
$H_i^{BM}(X)\cong H_i(\hat{X},\infty)$
where $\hat{X}$ is the one-point
compactification of $X$ and
$H_i(\hat{X},\infty)$ is the
relative homology group.
\item
\label{hilbert-schemes-item-isomorphism-relative-cohomology}
If the dimension of the manifold
where $X$ is embedded is $m$,
there is a canonical isomorphism
$H_i^{BM}(X) \cong H^{m-i}(M,M\setminus X)$.
In other words, there is a perfect
pairing
$$
H_i^{BM}(X) \times H_{m-i}(M,M\setminus
X)\to \mathbb{C}
$$
given by ``intersection''
(where the quotes account for some
transversal intersection condition
that we would require before
computing the intersection of
two simplicial chains).
\end{enumerate}

Consider the previous example of
the tiling of $\mathbb{R}^2$
by triangles and put
$X=M=\mathbb{R}^2$.
Take any point as a 0-chain in singular
homology
and the triangle tiling as 2-chain in BM
homology.
Since the point lies in exactly one triangle
(after suitable perturbation)
their intersection is 1.
Notice this wouldn't work with
singular homology: we need an infinite
2-chain
to make sure that the whole plane
is tiled by triangles since
otherwise we cannot guarantee that the point
is contained in some triangle.

\medskip\noindent
Now we formulate the key properties
of BM homology needed to define the operators
in Nakajima's representation of the
Heisenberg algebra.

\begin{enumerate}
\item (Pushforward.)
\label{hilbert-schemes-item-pushforward}
If $f:X \to Y$ is proper
(i.e. inverse image of any compact set is
compact)
then there is a pushforward map $f_*
:H_i^{BM}(X)\to H_i^{BM}(Y)$.
This is functorial in proper maps,
i.e. preserves compositions.
Roughly, this works because, by properness,
the inverse image of
a point in $Y$ will be mapped to a compact
set of $X$,
which can intersect only finitely many
simplices of a given chain,
ensuring local finiteness of the image under
$f$ of the chain.

\item
\label{hilbert-schemes-item-pullback}
(Pullback.) If $f:X \to Y$
is a locally trivial fibration with
fibre a manifold of dimension $m$ (over
$\mathbb{R}$),
then there is a pullback map
$f^* :H_i^{BM}(Y)\to H_{i+m}^{BM}(X)$.

\item
\label{hilbert-schemes-item-intersection}
(Intersection pairing.) If $X$ and $Y$
are both ``nicely embedded'' in the same
manifold $M$, there is a bilinear
graded-commutative map
$$
\cap_M:H_i^{BM}(X)\times H_j^{BM}(Y)
\to H_{i+j-m}^{BM}(X\cap Y).
$$
Constructing this product can be done
via the isomorphism with relative singular
cohomology
(item
(\ref{hilbert-schemes-item-isomorphism-relative-cohomology})
in the previous list),
and then using usual cap product.
Again, the term ``nicely embedded'' accounts
for a transversality condition that we will
not
specify.

\item (Fundamental classes.)
\label{hilbert-schemes-item-fundamental-class}
To any variety $V$ we can associate a
fundamental class $[V]$.
More precisely,
\begin{itemize}
\item If $V$ is a smooth irreducible variety
of dimension $d$ over $\mathbb{C}$,
its {\it fundamental class} $[V]$ is the
canonical generator of
$H_{2d}^{BM}(V)\cong H^0(V)\cong \mathbb{C}$,
where the first isomorphism
is given by the isomorphism of BM homology
and
relative singular cohomology taking $V = M$
since it is smooth.

\item If $V$ is a (possibly singular)
irreducible variety of dimension $d$
over $\mathbb{C}$, its fundamental class
$[V] \in B_{2d}^{BM}(V)$
is the unique class such that
$u^* [V]=[V_{\text{reg}}]$.
This is well-defined since
the pullback in this degree turns
out to be injective (by further properties
of BM homology not discussed here)
and thus this class is unique.

\item For $V$ not necessarily irreducible
of dimension $d$, then $H_{2d}^{BM}(V)$
has basis $[V_1],\ldots,[V_k]$
where $V_i$ are the $d$-dimensional
irreducible components of $V$.
In this case we define the fundamental class
as the sum of the fundamental classes of
those
components, i.e. $[V]=[V_1]+\ldots+[V_k]$.
\end{itemize}
\end{enumerate}

\noindent
As a final remark notice that it is
as special feature of BM homology
that we can define fundamental classes for
noncompact manifolds, as opposed to
singular homology.

\section{Correspondences}
\label{hilbert-schemes-section-correspondences}

\noindent
In this section we introduce the main tool
for constructing the operators in Nakajima's
representation
of the Heisenberg algebra.

\medskip\noindent
Let $X_1,X_2$ be smooth irreducible varieties
of dimensions $d_1,d_2$.
Let $Z$ be a $d$-dimensional
closed subvariety of $X_1\times X_2$
such that the projection onto the first
factor
is proper.

The {\it correspondence} associated to $Z
\subset X_1\times X_2$
is the map
\begin{equation}
\label{hilbert-schemes-equation-correspondence}
\begin{aligned}
c_{[Z]}:H_i^{BM}(X_2)&\longrightarrow
H_{i+2d-2d_2}^{BM}(X_1)\\
c&\longmapsto p_{1,*}(p_2^*c\cap [Z])
\end{aligned}
\end{equation}

\noindent
where $p_1$ and $p_2$ are the projections in
the following diagram:
$$
\xymatrix{
&Z\ar[ld]_{p_1}\ar[rd]^{p_2} \subset
X_1\times X_2\\
X_1&&X_2.
}
$$
\noindent
That is, the correspondence is an operator
given by pulling back a class by $p_2$,
intersecting with the fundamental class of
$[Z]$,
and pushing forward by $p_1$.
This map is well-defined according
to the key properties of BM homology stated
in the previous section;
the degree of the target homology group is
computed according
to such properties.

\begin{remark}
\label{hilbert-schemes-remark-notes}
\begin{enumerate}
\item
\label{hilbert-schemes-item-any-class}
We can define correspondences more generally
replacing the fundamental class of $Z$
by any $\alpha \in H_{2d}^{BM}(Z)$.
In fact, the operators in our representation
shall be of this more general type.

\item If $Z$ is the graph of a proper
morphism
$f:X_2 \to X_1$, then $c_{[Z]}=f_*$.
This shows that correspondences generalize
the pushforward in BM homology.

\item Since $X_1$ and $X_2$ are smooth,
we may use the isomorphism of BM and
relative singular cohomology
to identify a class in the tensor product
with an endomorphism
as we do in general for $V,W$ vector spaces
via the usual isomorphism
$V \otimes W^* \cong \Hom(W,V)$, $v \otimes
\varphi \mapsto
(w \mapsto \varphi(w)v)$.
In this way the fundamental class
$[Z] \in
\bigoplus_{i+j=\dim_\mathbb{R}Z}H_i(X_1)\otimes H_j^{BM}(X_2)$
immediately gives a map
$$
\bigoplus_{n}H_* (X^{[n]}) \to
\bigoplus_{n}H_* (X^{[n\mp i]})
$$
which, in fact, coincides with Equation
\ref{hilbert-schemes-equation-correspondence}
(see \cite[Section 8.3]{Nakajima-lectures}).

\end{enumerate}
\end{remark}

\medskip\noindent
Now we show how to compose correspondences,
as we shall do in the proof of the main
theorem.

Suppose we are given two subvarieties
$Z_{12} \subset X_1 \times X_2$ and $Z_{23}
\subset X_2 \times X_3$
such that the projections to $X_2$ and $X_3$
are proper.
We define the composition of the
corresponding correspondences by
\begin{equation}
\label{hilbert-schemes-equation-composition}
[Z_{12}]\circ[Z_{23}]:=p_{13,*}(p_{12}^*
[Z_{12}]\cap p_{23}^* [Z_{23}])
\end{equation}
according to the diagram
$$
\xymatrix{
&  X_1 \times X_2 \times
X_3\ar[dl]_{p_{12}}\ar[dr]^{p_{23}}\ar[dd]_{p_{13}}\\
 X_1\times X_2& & X_2\times X_3\\
& X_1\times X_3
}
$$
which is well-defined since, in fact,
the projection $p_{13}(p_{12}^{-1}Z_{12}
\cap p_{23}^{-1} Z_{23})\to X_3$
is proper (see \cite[Section
8.3]{Nakajima-lectures}).




\section{Main construction}
\label{hilbert-schemes-section-main-construction}

\noindent
In this section we define the
correspondences used in Nakajima's
representation
of the Heisenberg algebra and formulate
the theorem which guarantees that it is so.

Let $n\geq i>0$. Define $P[i]$
as a subset of
$\coprod_nX^{[n-i]}\times X^{[n]}$
by
$$
P[i]:=\coprod_n
\{(\mathcal{J}_1,\mathcal{J}_2,x)
\in X^{[n-i]}\times X^{[n]}:
\mathcal{J}_1 \supseteq \mathcal{J}_2,
\text{supp}(\mathcal{J}_1/\mathcal{J}_2)=\{x\}
\text{ for some }x \in X\}.
$$

\noindent
For $i<0$, we swap $\mathcal{J}_1$ and
$\mathcal{J}_2$
in this definition.
Now let
$\alpha \in H_*^{BM}(X)$, $\beta \in H_*(X)$
and suppose that $i>0$. Define the classes
\begin{equation}
\label{hilbert-schemes-equation-correspondence-p}
\begin{aligned}
P_\alpha[i]&=\varpi_*(\Pi_1^*\alpha \cap
[P[i]])\\
P_\beta[i]&=\varpi_*(\Pi_1^*\beta \cap
[P[-i]])
\end{aligned}
\end{equation}
according to the projections
$$
\xymatrix{
&  \coprod_n X^{[n-i]}\times X^{[n]}\times X
\ar[dl]_{\Pi_1}\ar[dr]^{\varpi_i}\\
X&&\coprod_n X^{[n-i]}\times X^{[n]}.
}
$$
These homology classes define the operators
of the group
$\bigoplus_{n=0}^{\infty}H_*(X^{[n]})$
that will be used in the main theorem.

\medskip\noindent
Before we continue, we formulate and
prove (with the help of other results,
as explained in \cite[Section
8.3]{Nakajima-lectures})
a lemma that will be used later.

\begin{lemma}
\label{hilbert-schemes-lemma-dimension}
$\dim_\mathbb{C}P[i]\cap (X^{[n-i]}\times
X^{[n]})=2n-i+1$.
\end{lemma}

\begin{proof}
Suppose $i>0$.
Fix $\mathcal{J}_1 \in X^{[n-i]}$.
By Theorem \ref{hilbert-schemes-theorem-fogarty} we know that
$\dim X^{[n-i]}=2(n-i)$.
To find the dimension of $P[i]$
we must add to this number the dimension of
the variety
obtained by varying  $\mathcal{J}_2 \in
X^{[n]}$
satisfying
$\mathcal{J}_1/\mathcal{J}_2=\{x\}$
for some $x \in X$.
The generic case is when
$x \in
X\setminus\text{supp}(\mathcal{O}_X/\mathcal{J}_1)$
(this may be accepted under the explanation
that since $\mathcal{O}_X/\mathcal{J}_1$ is
the structure sheaf
of a subvariety, its support has less
dimension than that of $X$;
a formal proof is provided in \cite[Lemma
8.31]{Nakajima-lectures}).
In fact, the set of such $\mathcal{J}_2$ is
$\pi^{-1}(S^i_{(i)}X)$, where $\pi$ is
the Hilbert-Chow morphism and
$S^i_{(i)}(X)=\{i[x] \in S^i(X):x \in X\}$.
This notation is used in other parts of
\cite{Nakajima-lectures};
in our case just notice that $S^i_{(i)}(X)$
is the subset of the symmetric product
$S^i(X)$
consisting of cycles of a single point with
weight $i$.
Finally, $\pi^{-1}(S^i_{(i)}(X))$ has
complex dimension $i+1$
as a subvariety of the Hilbert scheme.
\end{proof}

\medskip\noindent
Now we state the main theorem,
\cite[Theorem 8.13]{Nakajima-lectures}.
Notice that we must add the assumption that
the homology group of $X$ is concentrated in
even degrees.
The theorem can be stated in a more general
form without
this assumption, but we shall not discuss it
in this work.
\begin{theorem}[Nakajima, Grojnowski]
\label{hilbert-schemes-theorem-main}
Suppose $(\text{deg}
\alpha)(\text{deg}\beta)\equiv 0 (\text{mod}
2)$
for all $\alpha,\beta \in H_*(X)$.
Then
\begin{equation}
\label{hilbert-schemes-equation-main}
[P_\alpha[i],P_\beta[j]]=(-1)^{i-1}i
\delta_{i,-j}
\left\langle{}\alpha,\beta\right\rangle{}\text{Id}.
\end{equation}

\noindent
where $\left\langle{}\alpha,\beta\right\rangle{}=p_*(\alpha
\cap \beta)$
for $p:X \to \text{pt}$ the unique morphism
from $X$ to the point pt.
\end{theorem}

\medskip\noindent
According to our constructions in
Section \ref{hilbert-schemes-section-heisenberg},
this means that we have a representation of
the Heisenberg algebra $\hat{\mathfrak{h}}$
associated to the graded vector space
$H_*(X)$
on $\bigoplus_{n=0}^{\infty}H_{*}(X^{[n]})$.

\medskip\noindent
Notice the factor $(-1)^{i-1}$
in Equation \ref{hilbert-schemes-equation-main},
which is not present in the Lie bracket
equations of any of the
variants of the Heisenberg algebra discussed
earlier in this work.
The reason for including this term in the
main theorem
is explained in
\cite[Theorem 9.15]{Nakajima-lectures}.
Moreover, this factor does not modify the
Lie algebra structure:
we can define a Lie algebra isomorphism
from, say, the most basic version of the
Heisenberg algebra
with bracket given by Equation
\ref{hilbert-schemes-equation-bracket-basic},
to the Lie algebra
$\bigoplus_{n \in
\mathbb{Z}}\mathbb{C}\tilde{h}_n
\oplus\mathbb{C}K$
with bracket
$$
[\tilde{h}_m,\tilde{h}_m]=(-1)^{m-1}m\delta_{m,-n}K,\qquad
\text{$K$ central,}
$$
by putting
$$
h_m \mapsto
\begin{cases}
(-1)^{m-1}\tilde{h}_m\qquad &m<0 \\
\tilde{h}_m\qquad &m\geq 0.
\end{cases}
$$



\section{Worked example}
\label{hilbert-schemes-section-example}

\noindent
In this section we do a worked example
following \cite{Nakajima-lectures}
and \cite[Video 10]{videos}
to illustrate the main idea of the proof.

\medskip\noindent
Fix $x \in X$ and let $\alpha=[X]$ be the BM
class of $X$,
and $\beta=[x]$ the class of the point $x$.
Take a set of distinct $n$ points
$x_1,\ldots,x_n \in X\setminus\{x\}$
none of which is $x$.
Identify them with the corresponding
ideal $\mathcal{J}$. We refer interchangeably
to the set of points and the ideal sheaf
$\mathcal{J}$.

Let's show that
$$
[P_\alpha[1],P_\beta[-1]][\mathcal{J}]=[\mathcal{J}].
$$
We start by computing
$P_\beta[-1][\mathcal{J}]$.
For this consider
$$
\xymatrix{
&X^{[n+1]}\times
X^{[n]}\ar[ld]_{p_1}\ar[rd]^{p_2}\\
X^{[n+1]}&&X^{[n]}.
}
$$
By definition, $P_{[x]}[-1]$ is given by:
pull back $\mathcal{J}$ via $p_2$,
intersect with the class
$\varpi_*(\Pi^*\alpha \cap
[P[-1]])=P_{[x]}[-1]$
and push forward
down to $X^{[n+1]}$.
This means that we push forward the cycle
corresponding to
$$
\{(\mathcal{J}, \mathcal{J}',x)
:\mathcal{J}\subseteq
\mathcal{J}',\text{supp}(\mathcal{J}'/\mathcal{J})=\{x\}\}
\subset X^{[n]} \times X^{[n+1]} \times X.
$$
The pushforward recovers the $X^{[n+1]}$
component,
that is, the schemes $\mathcal{J}'$ in
$X^{[n+1]}$ satisfying
$\text{supp}(\mathcal{J}'/\mathcal{J})=\{x\}$.
Since $x$ is not a point in $\mathcal{J}$,
the only possible choice is that
$\mathcal{J}'=\{x\}\cup \mathcal{J}$,
that is,
$$
P_{[x]}[-1][\mathcal{J}]=[\{x\}\cup
\mathcal{J}].
$$
(see figure below).

On the other hand, the correspondence
$P_{[X]}[1]$
starts with an ideal $\mathcal{J}'\in
X^{[n+1]}$
intersects it with
$$
\{(\mathcal{J}, \mathcal{J}',x)
:\mathcal{J}\subseteq
\mathcal{J}',\text{supp}(\mathcal{J}'/\mathcal{J})=\{y\},
\text{ for some }y\in X\}
$$
and projects to $X^{[n]}$. We obtain
$\sum_{y \in
\mathcal{J}'}[\mathcal{J}'-\{y\}]$.
Then it becomes clear that
$$
[P_\alpha[1],P_\beta[-1]][\mathcal{J}]=[\mathcal{J}].
$$
The following figure is taken directly from
\cite{Nakajima-lectures}:



\section{Proof of one case of the main
theorem}
\label{hilbert-schemes-section-proof}

\noindent
The proof of Theorem \ref{hilbert-schemes-theorem-main} is
split
in two different cases:
\begin{enumerate}
\item $i,j>0$ or $i,j<0$, and
\label{hilbert-schemes-item-same}
\item $i>0$, $j<0$ or $i<0$, $j>0$.
\label{hilbert-schemes-item-different}
\end{enumerate}

\noindent
In this work we explain the proof of case
(\ref{hilbert-schemes-item-same}).

\medskip\noindent
Suppose $i,j>0$ (the case $i,j<0$ is
analogous).
Recalling our general definition for
composition of correspondences
(Equation \ref{hilbert-schemes-equation-composition})
we see that
the composition $P_\alpha[i]P_\beta[j]$ is
given by
\begin{equation}
\label{hilbert-schemes-equation-composition-pij}
P_\alpha[i]P_\beta[j]
=p_{13,*}(p_{12}^*[P[i]]\cap\Pi_1^*\alpha\cap p_{23}^*[P[j]]\cap\Pi_2^*\beta).
\end{equation}

Consider the set
\begin{equation}
\label{hilbert-schemes-equation-cycle-pij}
\begin{aligned}
P[i,j]&:=p_{12}^{-1}(P[i])\cap
p_{23}^{-1}(P[j])
\cap(X^{[n-i-j]}\times X^{[n-j]}\times
X^{[n]})\\
&=\left\{(\mathcal{J}_1,\mathcal{J}_2,\mathcal{J}_3):
\mathcal{J}_1\supset\mathcal{J}_2\supset\mathcal{J}_3,
\substack{
\text{supp}(\mathcal{J}_1/\mathcal{J}_2)=\{x\}
\text{ for some }x\in X\\
\text{supp}(\mathcal{J}_2/\mathcal{J}_3)=\{y\}
\text{ for some }y \in X }\right\}
\end{aligned}
\end{equation}

\noindent
which corresponds to the diagram
$$
\xymatrix{
&  X^{[n-i-j]} \times X^{[n-j]} \times
X^{[n]}
\ar[dl]_{p_{12}}\ar[dr]^{p_{23}}\ar[dd]_{p_{13}}\\
P[j]\subset X^{[n-j]}\times X^{[n]}& &
X^{[n-i-j]}\times X^{[n-j]}\supset P[i]\\
& X^{[n-i-j]}\times X^{[n]}.
}
$$
By definition, the correspondence $P[j]$
maps the homology of $X^{[n]}$ to the
homology of $X^{[n-j]}$. In order to compose
this with $P[i]$
we must regard $P[i]$ as a map from the
homology of $X^{[n-j]}$
to the homology of $X^{[n-j-i]}$, which is
perfectly
valid since our correspondences are defined
for all $n$.
Observe that the middle factor $X^{[n-j]}$
in the triple product will become
$X^{[n-i]}$ when
we compute the composition in the opposite
order.

\medskip\noindent
To recover the correspondence as in
Equation \ref{hilbert-schemes-equation-composition-pij}
from the set $P[i,j]$ it's natural to
consider its
fundamental class.
We will do this in the following special way.
Observe that $P[i,j]$ may be decomposed as a
disjoint union
of the set $\{x \neq y\}$
of ``generic'' cases in which $x$ and $y$ in
Equation \ref{hilbert-schemes-equation-cycle-pij}
are different, and the set $\{x=y\}$ in
which they coincide.
Let
$\overline{P}[i,j]:=\overline{P[i,j]\cap\{x\neq y\}}$.
Then
$$
p_{12}^* [P[i]]\cap p_{23}^*
[P[j]]=[\overline{P}[i,j]]+\iota_*R
$$
where $R \in H_*^{BM}(P[i,j]\cap \{x=y\})$
corresponds
to the $\{x=y\}$ part of $P[i,j]$
and $\iota$ is the inclusion.

\medskip\noindent
We can carry out the same construction
for the composition $P_\beta[j]P_\alpha[i]$,
obtaining
\begin{equation}
\label{hilbert-schemes-equation-composition-pji}
P_\alpha[j]P_\beta[i]
=p_{13,*}(p_{12}^*[P[j]]\cap\Pi_2^*\alpha\cap p_{23}^*[P[i]]\cap\Pi_1^*\beta).
\end{equation}

\noindent
and
\begin{equation}
\label{hilbert-schemes-equation-pji-r}
p_{12}^* [P[j]]\cap p_{23}^*
[P[i]]=[\overline{P}[j,i]]+\iota'_*R'.
\end{equation}

\noindent
Observe that Equations
\ref{hilbert-schemes-equation-composition-pij} and
\ref{hilbert-schemes-equation-composition-pji} differ in the
classes
$\Pi_1^*\alpha$ against $\Pi_2^*\alpha$, and
correspondingly with $\beta$.
The following lemma is all we need to obtain
the result:

\begin{lemma}
\label{hilbert-schemes-lemma-pij}
\begin{enumerate}
\item There exists an isomorphism
\begin{equation}
\label{hilbert-schemes-equation-pij-isomorphism}
P[i,j]\cap\{x \neq
y\}\xrightarrow{\cong}P[j,i]\cap \{x \neq
y\},
\end{equation}

\noindent
which exchanges $\Pi_1$ and $\Pi_2$.

\item We have
\begin{equation}
\label{hilbert-schemes-equation-r}
p_{13,*}(\Pi_1^*\alpha\cap\Pi_2^*\beta\cap\iota_*R)=0,\qquad
p_{13,*}(\Pi_1\beta\cap\Pi_2^*\alpha\cap\iota_*R')=0.
\end{equation}
\end{enumerate}
\end{lemma}

\noindent
Indeed, this lemma is all we need since
$$
\begin{aligned}
P_\alpha[i]P_\beta[j]
&=(-1)^{\text{deg}\alpha}p_{13,*}((\overline{P}[i,j]+\iota_*R)\cap\Pi_1^*\alpha\cap\Pi_2^*\beta)\\
&=(-1)^{\text{deg}\alpha\text{deg}\beta}
p_{13,*}((\overline{P}[j,i]+\iota_*R')\cap\Pi_1^*\beta\cap\Pi_2^*\alpha)\\
&=P_\beta[j]P_\alpha[i].
\end{aligned}
$$

\begin{proof}[Proof of Lemma \ref{hilbert-schemes-lemma-pij}]
\begin{enumerate}
\item In this proof I assume that
``isomorphism'' means bijection, since
$P[i,j]$ is considered as a set-theoretical
intersection.

First notice that when we say that the
isomorphism, say $\varphi$,
``exchanges $\Pi_1$ and $\Pi_2$''
we mean that $\Pi_1|_{X \times X}=\Pi_2|_{X
\times X}\circ\varphi$
according to the diagram
$$
\xymatrix{
& & X\\
P[i,j] \cap \{x\neq y\}\subset
\ar[ddd]_\varphi&
\save{
X^{[n-i-j]}\times X^{[n-j]} \times X^{[n]}
\times X \times X
}
\ar[ur]^{\Pi_1}\ar[dr]_{\Pi_2}
\restore
\\
& & X\\
& & X\\
P[j,i]\cap\{x\neq y\}\subset&
X^{[n-i-j]}\times X^{[n-i]} \times X^{[n]}
\times X \times X
\ar[ur]^{\Pi_1}\ar[dr]_{\Pi_2}\\
& & X.
}
$$

Now, the point of the isomorphism $\varphi$
is that it must
map an ideal in $X^{[n-j]}$ to an ideal in
$X^{[n-i]}$.
More precisely, we map
$(\mathcal{J}_1,\mathcal{J}_2,\mathcal{J}_3)\in P[i,j] \cap \{x \neq  y\}$
to
$(\mathcal{J}_1,\mathcal{J}_2',\mathcal{J}_3) \in P[j,i]\cap \{x \neq y\}$.
The defining relation is
\begin{equation}
\label{hilbert-schemes-equation-isomorphism}
\mathcal{J}_1/\mathcal{J}_2'=\mathcal{J}_2/\mathcal{J}_3,\qquad
\mathcal{J}_2'/\mathcal{J}_3=\mathcal{J}_1/\mathcal{J}_2.
\end{equation}

\noindent
The following figure
(taken from \cite{Nakajima-lectures})
shows how this correspondence works.

The setting is the following. We have the
ideals with contentions
$\mathcal{J}_1\supset
\mathcal{J}_2,\mathcal{J}_2'\supset
\mathcal{J}_3$.
Passing to the supports, the inclusions are
reversed, i.e
the support of $\mathcal{J}_3$ contains the
rest,
and the support of $\mathcal{J}_1$ is
contained in the rest.

To prove that this function is well-defined
we need to show that for every
$\mathcal{J}_2$ there is a unique
$\mathcal{J}_2'$ satisfying Equation
\ref{hilbert-schemes-equation-isomorphism}.
Notice that by definition,
$\text{supp}\mathcal{J}_2/\mathcal{J}_3=\{x\}$
for some $x \in X$ so that, as a scheme,
$\{x\}$ has multiplicity $j$.
We define
$\mathcal{J}_2'=\mathcal{J}_1\cup \{x\}$
where $x$ has multiplicity $j$.

\medskip\noindent
(Caveat: it looks to me that the choice of
$\mathcal{J}_2'$ is not unique.
Indeed, the space of schemes of length $j$
supported
in the point $x$ was discussed in the proof
of Lemma \ref{hilbert-schemes-lemma-dimension},
where we used the fact that it has dimension
$j-1$.)

\medskip\noindent
Defining $\mathcal{J}'_2$ this way implies
that
$$
\mathcal{J}_1/\mathcal{J}_2'=\mathcal{J}_1/(\mathcal{J}_1\cup \{x\})
=\{x\}=\mathcal{J}_2/\mathcal{J}_3.
$$
Now, we also must have
$\mathcal{J}_1/\mathcal{J}_2=\{y\}$
for some $y \in X$ with multiplicity $i$.
This implies that
$\mathcal{J}_1/\mathcal{J}_3=\{x\}\cup
\{y\}$
with their multiplicities.
Then
$$
\mathcal{J}_2'/\mathcal{J}_3=(\mathcal{J}_1\cup \{x\})/\mathcal{J}_3=\{y\}
=\mathcal{J}_1/\mathcal{J}_2
$$
as required.
An inverse function
defined similarly shows this is a bijection.

\item Consider the commutative diagram
$$
\xymatrix{
P[i,j]\cap
\{x=y\}\ar[r]^-{\Pi'}\ar[d]_{\iota}&X\ar[d]^{\Delta}\\
P[i,j]\ar[r]_{\Pi_1 \times \Pi_2}&X \times X
}
$$
where $\Delta$ is the diagonal embedding
and $\Pi'$ is any of the projections $\Pi_1$
or $\Pi_2$
(which coincide on  $P[i,j]\cap \{x=y\}$).
Then
\begin{equation}
\label{hilbert-schemes-equation-lemma2-1}
\begin{aligned}
\Pi_1^*\alpha \cap \Pi_2^*\beta \cap\iota_*R
&=\iota_*(\iota^*(\Pi_1^*\alpha\cap\Pi_2^*\beta)\cap R)\\
&=\iota_*(\iota^*(\Pi_1\times\Pi_2)^*(\alpha\times\beta)\cap R)\\
&=\iota_*({\Pi'}^*\Delta^*(\alpha\times\beta)\cap R)\\
&=\iota_*({\Pi'}^*(\alpha\cap\beta)\cap R),
\end{aligned}
\end{equation}

\noindent
where the first equality is due
to the projection formula
(see Stacks Project
\href{https://stacks.math.columbia.edu/tag/01E8}{01E8})
and we stick with Nakajima's notation
$\alpha \times \beta$
for the tensor product $\alpha \otimes
\beta$.
We also have the commutative diagram
$$
\xymatrix{
p_{13}\circ \iota:P[i,j]\cap\{x=y\}
\ar[r]^-{\Pi'}\ar[d]_{p_{13 \circ
\iota}}&X\ar@{=}[d]\\
P[i+j]=\{(\mathcal{J}_1,\mathcal{J}_3)|\mathcal{J}_1 \supset \mathcal{J}_3,
\text{supp}(\mathcal{J}_1/\mathcal{J}_3)=\{x\}\}\ar[r]_-{\Pi''}&X,
}
$$
where $\Pi''$ is the usual projection to $X$
from the
set $P[i+j]$.

By Equation \ref{hilbert-schemes-equation-lemma2-1}
and the commutativity of the second diagram,
we can express the class of interest as
intersection
of two classes: one pulled back directly
from $X$,
and the other one corresponding to $R$.
That is, we have
$$
p_{13,*}\iota_*({\Pi'}^* (\alpha\cap
\beta)\cap R)
={\Pi''}^* (\alpha \cap \beta)\cap
p_{13,*}\iota_*(R).
$$
Now, by Lemma \ref{hilbert-schemes-lemma-dimension},
we know that
$$
\dim_{\mathbb{C}}P[i+j]=2n-i-j+1.
$$
On the other hand, the ``expected dimension''
is $2n-i-j+2$. This obliges $\iota_*R$ to
vanish.

\medskip\noindent
(Caveat: I understand that the ``expected
dimension''
means in \cite{Nakajima-lectures} the
dimension of
$P[i,j]\cap \{x\neq y\}$, which is claimed
to be
$2n-i-j+2$.
It is however not clear to me why this would
imply
that $p_{13,*}\iota_*R$ must vanish.)

\end{enumerate}
\end{proof}


% --- End hilbert-schemes.tex ---

% --- Begin noncommutative-abelian-surfaces-and-kummer-type-hyperkahler-varieties.tex ---
\chapter{Noncommutative Abelian Surfaces and Kummer Type Hyperk\"ahler Varieties}
\label{chapter-noncommutative-abelian-surfaces-and-kummer-type-hyperkahler-varieties}


Kummer Type Hyperk\"ahler Varieties}

\noindent
Minicourse by Xiaolei Zhao (University of
California, Santa Barbara (CA), USA), CIMPA
school Florian\'opolis 2025.

\medskip\noindent
Notes at
\href{%
http://github.com/danimalabares/cimpa-floripa
}{%
github.com/danimalabares/cimpa-floripa
}

\bigskip\noindent

{\bf Abstract.} Examples of noncommutative K3
surfaces arise from semiorthogonal
decompositions of the bounded derived
category of certain Fano varieties. The most
interesting cases are those of cubic
fourfolds and Gushel-Mukai varieties of even
dimension. Using the deep theory of families
of stability conditions, locally complete
families of hyperk\"ahler manifolds
deformation
equivalent to Hilbert schemes of points on a
K3 surface have been constructed from moduli
spaces of stable objects in these
noncommutative K3 surfaces. On the other
hand, an explicit description of a locally
complete family of hyperk\"ahler manifolds
deformation equivalent to a generalized
Kummer variety is not available from
classical geometry. In this lecture series,
we will construct families of noncommutative
abelian surfaces as equivariant categories of
the derived category of K3 surfaces which
specialize to Kummer K3 surfaces. Then we
will explain how to induce stability
conditions on them and produce examples of
locally complete families of hyperk\"ahler
manifolds of Kummer type. Based on joint work
with Arend Bayer, Alex Perry and Laura
Pertusi.

\bigskip\noindent
\bigskip\noindent

\noindent

\section{K3 surfaces and Hyperk\"ahler
varieties}
\label{noncommutative-abelian-surfaces-and-kummer-type-hyperkahler-varieties-nasak-label-001}

\begin{definition}
\label{noncommutative-abelian-surfaces-and-kummer-type-hyperkahler-varieties-definition-K3-surface}
A {\it K3 surface} (over $\mathbb{C}$) is a
smooth projective surface $S$ such that
$\omega_S \cong \mathcal{O}_S$ and
$\pi_1(S)=1$.
\end{definition}

\begin{example}[K3 surfaces]
\label{noncommutative-abelian-surfaces-and-kummer-type-hyperkahler-varieties-example-K3-surfaces}
\begin{enumerate}
\item $S\subset \mathbb{P}^3$ of degree 4. A
straightforward computation shows
the conditions of the definition are
verified.
\item $S \to \mathbb{P}^2$ double cover
ramified along a degree 6 curve.
\end{enumerate}
\end{example}

\begin{definition}
\label{noncommutative-abelian-surfaces-and-kummer-type-hyperkahler-varieties-definition-hyperkahler-variety}
A {\it hyperk\"ahler} variety is a smooth
projective variety $X$ such that
\begin{enumerate}
\item $\pi_{1}(X)=1$
\item $H^{0}(X,\Omega_X^2)=\mathbb{C} \omega$
(i.e. the space of global
holomorphic 2-forms is one dimensional) where
$\omega$ is holomorphically symplectic.
\end{enumerate}
\end{definition}

\begin{remark}
\label{noncommutative-abelian-surfaces-and-kummer-type-hyperkahler-varieties-remark-hyperkahler}
If $X$ is hyperk\"ahler, it must be even
dimensional and $\omega_X=\mathcal{O}_X$.
\end{remark}

\begin{remark}
\label{noncommutative-abelian-surfaces-and-kummer-type-hyperkahler-varieties-remark-Beauville-Bogomolov}
We just mention the name of
Beauville-Bogomolov theorem.
\end{remark}

For a while people were looking for examples
of hyperk\"ahler varieties.

\begin{example}[Hyperk\"ahler varieties]
\label{noncommutative-abelian-surfaces-and-kummer-type-hyperkahler-varieties-example-hyperkahler}
\begin{enumerate}
\item K3 surface $S$.
\item $X=S^{[n]}$ Hilbert scheme of $n$
points on K3 surface. (Moduli space of
stable sheaves on K3 $S$ of rank 1,  $c_1=0$,
$c_2=0$.)
\item $\mathcal{M}_H(s,v)=$ moduli of
$H$-stable sheaves on $S$ of class $v$.
($v$ primitive, $H$ is $v$-generic.) (Recall
the example in Cristina's course where we
studied how the moduli changes under changes
in the polarization $H$.) Let $[E] \in
\mathcal{M}_H(s,v)$. We have the Yoneda map
$$
T_{[E]}\mathcal{M}\cong
\text{Ext}^1(E,E)\times \text{Ext}^1(E,E)
\to \text{Ext}^2(E,E)\overset{\substack{S
\text{ is}\\\text{ a
K3}}}{\cong}\Hom(E,E)^*\cong \mathbb{C}
$$
where the first factor product of
$\text{Ext}^1(E,E)\times \text{Ext}^1(E,E)$
is associated to first arrow, the second
factor to the second arrow, of the following
diagram:
$$
E \to E[1] \to E[2]
$$
This moduli space is always of Picard rank 2.
\item $Y \subset \mathbb{P}^5$ a smooth cubic
fourfold.
$$
F(Y)=\{[\ell] \in \text{Gr}(2,6)|\ell
\subseteq Y\}
$$
is a hyperk\"ahler 4-fold. For a specific
choice $Y_0$ of hyperk\"ahler 4-fold in the
moduli of smooth cubic 4-folds (which is
20-dimensional) we find $F(Y_0)$ is
deformation equivalent to $S^{[2]}$ where $S$
is a K3 surface.
\end{enumerate}
\end{example}

\medskip\noindent
Here is a general picture:
\iffalse$$
\xymatrix{
\substack{\text{cubic 4-fold $Y$} \\
D^b(Y)}\ar[dr]&  &
\substack{\text{moduli of} \\ \text{stable
objects}\\ \text{in $\text{Ku}_Y$ as}
\\ \text{HK variety}}\\
&\substack{\text{sod K3} \\ \text{cat
conditions} \\ \text{Ku}_Y }
\ar[ur]_{\substack{\text{stability} \\
\text{conditions}}}
$$\fi

\medskip\noindent
Let $X,Y$ be projective smooth varieties and
$f:X\to Y$. Consider the functors
$$
Rf_*:D^b(X) \to D^b(Y)\qquad Lf^*:D^b(Y) \to
D^b(X)
$$
for $F,G \in D^b(X)$, $F\overset{L}{\otimes}G
\in D^b(X)$. {\bf Convention:} we drop the
$R$ and $L$.

\begin{definition}
\label{noncommutative-abelian-surfaces-and-kummer-type-hyperkahler-varieties-definition}
Let $K \in D^b(X \times Y)$. The {\it
Fourier-Mulai functor} is
$$
\begin{aligned}
\Phi_K: D^b(X) &\longrightarrow D^b(Y) \\
F &\longmapsto \text{pr}_{Y,*}(\text{pr}_X^*F
\otimes K)
\end{aligned}
$$
where we are using our convention --- all
functors here are derived. Here the maps are
$$
\xymatrix{
& X \times
Y\ar[dl]_{\text{pr}_X}\ar[dr]^{\text{pr}_Y}\\
X& & Y
}
$$
\end{definition}

\begin{example}
\label{noncommutative-abelian-surfaces-and-kummer-type-hyperkahler-varieties-example-Fourier-Mukai}
$f:X \to Y$ and $\Gamma_f \equiv X \times Y$
graph. Then $\Phi_{\mathcal{O}_{p_f}}=f_*$.
\end{example}

\begin{theorem}[Orlov]
\label{noncommutative-abelian-surfaces-and-kummer-type-hyperkahler-varieties-theorem-Orlov}
If $F:D^b(X) \to D^b(Y)$ is an equivalence,
then $\exists ! K \in D^b(X \times Y)$ such
that $F \cong \Phi_K$.
\end{theorem}

\begin{theorem}
\label{noncommutative-abelian-surfaces-and-kummer-type-hyperkahler-varieties-nasak-label-002}
The following are preserved by derived
equivalence:
\begin{enumerate}
\item (Bondal-Orlov.) $\dim X$.
\item (Bondal-Orlov.) $\bigoplus_{m \geq
0}H^{0}(X,\pm m K_X)$.
\item $H^{*}(X,\mathbb{Q})$, and
$\bigoplus_{p-q=i}H^{p,q}(X)$ for any $i \in
\mathbb{Z}$.
\end{enumerate}
\end{theorem}

\medskip\noindent
{\bf Non-trivial equivalences.} $S$ a K3
surface. $H^{2}(S,\mathbb{Z})$.
\begin{itemize}
\item (Hodge decomposition.)
$H^{2}(S,\mathbb{C})\cong H^{2,0} \oplus
H^{1,1} \oplus H^{0,2}$.
\item Pairing on $H^{2}(S,\mathbb{Z})$ given
by cup product.
\end{itemize}

\begin{theorem}[Torelli for K3]
\label{noncommutative-abelian-surfaces-and-kummer-type-hyperkahler-varieties-theorem-Torelli-for-K3}
Two K3 surfaces $S,S'$ are isomorphic if and
only if there exists a Hodge isometry
$\varphi:H^{2}(S,\mathbb{Z})\to
H^{2}(S',\mathbb{Z})$ (i.e. an isomorphism
that preserves the Hodge decomposition and
the pairing).
\end{theorem}

This was the classical content.

\subsection{Mukai lattice}
\label{noncommutative-abelian-surfaces-and-kummer-type-hyperkahler-varieties-subsection-Mukai-lattice}

``The Mukai lattice is some sort of
topological invariant that behaves better
with derived category.''

$$
\tilde{H}(S,\mathbb{Z})
=H^{0}(S,\mathbb{Z})\oplus
H^{2}(S,\mathbb{Z})\oplus H^{4}(S,\mathbb{Z})
$$
\begin{itemize}
\item Weight 2 Hodge structure
$$
\tilde{H}^{2,0}=H^{2,0}\qquad
\tilde{H}^{1,1}=H^0\oplus H^{1,1}\oplus H^2
\qquad \tilde{H}^{0,2}\cong H^{0,2}.
$$
\item Pairing:
$$
\left\langle{}(a,b,c),(a',b',c')\right\rangle{}=bb'-ac'-a'c
$$
\end{itemize}

\begin{theorem}[Mukai, Orlov]
\label{noncommutative-abelian-surfaces-and-kummer-type-hyperkahler-varieties-theorem-Mukai-Orlov}
If $S,S'$ are K3 surfaces, then $D^b(S) \cong
D^b(S')\iff
\tilde{H}(S,\mathbb{Z})\underset{\varphi}
{\cong}\tilde{H}(S',\mathbb{Z})$
Hodge
isometry.
\end{theorem}

\subsection{Mukai vector}
\label{noncommutative-abelian-surfaces-and-kummer-type-hyperkahler-varieties-subsection-Mukai-vector}

Let $K_0(S)$ be the free abelian group
generated by $\Ob(D^b(S))$. Consider
$$
v:K_0(S)\to \tilde{H}(S;\mathbb{Z})
$$
Where $[F]=[E]+[G]$ if $E \to F \to G
\xrightarrow{+1}$ is a distinguished
triangle.
$$
v:[E] \to \text{ch}(E)\cdot \sqrt{\det(S)}
$$
where $\text{ch}(E)$ is the Chern character.
$$
\left\langle{}v(E),v(F)\right\rangle{}=-\chi(E,F),
$$
where
$\chi(E,F):=\sum(-1)^i\dim\text{Ext}^i(E,F):=
\Hom_{D^b(S)}(E,F[i])$.

\begin{proof}[Proof of the backward
implication of Mukai-Orlov theorem]
Let
$$
\begin{aligned}
\varphi: \tilde{H}(S,\mathbb{Z})
&\longrightarrow \tilde{H}(S',\mathbb{Z}) \\
(0,0,1) &\longmapsto v\\
(-1,0,0) &\longmapsto v'
\end{aligned}
$$
The intersection matrix of $v$ and $v'$ is
$$
U=\begin{pmatrix}
0&1\\
1&0
\end{pmatrix}
$$
For simplicity assume that $v$ is of positive
rank. ($V=(a,b,c)$, where $a\in H^0$ is
positive.)

The heart of the proof uses the following
result by Mukai. There exists a nonempty
moduli space $\mathcal{M}$ of stable sheaves
on $S'$ of class $v$ $\implies$ $S'$ is a K3
surface!
$$
0=\chi(v,v)=\underbrace{\Hom(E,E)}_{=1}
-\underbrace{\text{ext}^1(E,E)}_{=2}
+\underbrace{\text{ext}^2(E,E)}_{=1}
$$
$v\cdot v'=1\implies \mathcal{M}$ is a fine
moduli space. There exists a universal family
$D^b(S \times M)\ni \mathcal{E} \to S' \times
M$.

{\bf Claim.}
$\Phi_{\mathcal{E}}:D^b(S')
\xrightarrow{\cong}
D^b(\mathcal{M})$ (general criteria).

Now
$\tilde{H}(S,\mathbb{Z})\xrightarrow{\varphi,
\cong}
\tilde{H}(S',\mathbb{Z}$
$$
\xymatrix@R=1em{
\tilde{H}(S,\mathbb{Z})\ar[r]^{\varphi,\cong}
&
\tilde{H}(S',\mathbb{Z})
\ar[r]^{\Phi_{\mathcal{E}}}&\overset{\vee}{H}
(M,\mathbb{Z})\\
(0,0,1)\ar@{|->}[r]&  v \ar@{|->}[r]&
(0,0,1)\\
(1,0,0)\ar@{|->}[rr]& & (1,0,0)
}
$$
\end{proof}

\section{Semiorthogonal decomposition and
Calabi-Yau
categories}
\label{noncommutative-abelian-surfaces-and-kummer-type-hyperkahler-varieties-section-SOD-and-CY-categories}

\begin{situation}
\label{noncommutative-abelian-surfaces-and-kummer-type-hyperkahler-varieties-situation-SOD-and-SC}
$$
\xymatrix{
D^b(Y),\substack{Y\text{ a cubic} \\
\text{4-fold}}\ar[dr]
_{\substack{\text{semiorthogonal}
\\
\text{decomposition}}}
&  &\substack{\text{moduli}\\ \text{on K3} \\
\text{as HK}}\\
&\text{Ku}_Y\substack{\text{K3} \\
\text{category}}
\ar[ur]_{\text{stability conditions}}
}
$$
\end{situation}

\begin{definition}
\label{noncommutative-abelian-surfaces-and-kummer-type-hyperkahler-varieties-definition-SOD}
Let $X$ be a projective smooth variety over
$\mathbb{C}$. A {\it SOD}
$D^b(X)=\left\langle{}A_1,\ldots,A_n\right\rangle{}$ is a
sequence $A_i$ full triangulated subcategory
such that
\begin{enumerate}
\item $\Hom(F,G)=0$ for all $F \in A_i$, $G
\in A_j$.
\item $\forall F \in D^b(X)$ $\exists
0=F_n\to F_{n-1}\to\ldots\to F_0=F$
such that $\text{Cone}(F_i \to F_{i-1})\in
A_i$.
\end{enumerate}
\end{definition}

\begin{exercise}
\label{noncommutative-abelian-surfaces-and-kummer-type-hyperkahler-varieties-exercise}
\begin{enumerate}
\item $D^b(\text{Spec}\mathbb{C})\ni V \simeq
\bigoplus_{i}H^{i}(V)[-i]$.
\item For $E \in D^b(X)$ define
$$
\begin{aligned}
\phi_E: D^b(\text{Spec}\mathbb{C})
&\longrightarrow D^b(X) \\
V &\longmapsto V\otimes
E=\bigoplus_{i}H^i(V)\otimes E[-i]
\end{aligned}
$$
Then $\phi_E$ fully faithful $\iff$ $E$ is
exceptional
$$
\left(\text{Ext}^p(E,E)=
\begin{cases}
\mathbb{C}\qquad &p=0 \\
0\qquad &p\neq 0
\end{cases}\right)
$$
\begin{definition}
\label{noncommutative-abelian-surfaces-and-kummer-type-hyperkahler-varieties-definition-duplicate-2}
$\left\langle{}E\right\rangle{}:=\phi_E(D^b(\text{Spec}
\mathbb{C}))\subset
D^b(X)$ for exceptional $E$.

\end{definition}
\item {\bf Example.} $\mathcal{O}(i)$
exceptional on $\mathbb{P}^n$.
(Beilinson.) $D^b(\mathbb{P}^n)
=\left\langle{}\mathcal{O},\mathcal{O}(1),\ldots,
\mathcal{O}(n)\right\rangle{}$
\end{enumerate}
\end{exercise}

\begin{definition}
\label{noncommutative-abelian-surfaces-and-kummer-type-hyperkahler-varieties-definition-right-admissible}
A triangulated category $A \subset D^b(X)$ is
{\it right admissible} if $\alpha:A\to
D^b(X)$ admits a right adjoint.
$\alpha^!:D^b(X)\to A$. (I.e.,
$\Hom_{D^b(X)}(\alpha(E),F)\simeq
\Hom_A(E,\alpha^!(F))$
\end{definition}

\begin{exercise}
\label{noncommutative-abelian-surfaces-and-kummer-type-hyperkahler-varieties-exercise-right-admissible}
Let $E$ be exceptional.
$\left\langle{}E\right\rangle{}\subset D^b(X)$ is right
admissible.
$$
\alpha^!=R \Hom(E,F) \in D^b(\text{Spec}
\mathbb{C})
$$
\end{exercise}

Let
$$
A^\perp:=\{F \in D^b(X):\Hom(E,F)=0, \forall
E \in ?\}
$$
\begin{lemma}
\label{noncommutative-abelian-surfaces-and-kummer-type-hyperkahler-varieties-lemma-right-admissible}
$A \subset D^b(X)$ is right admissible
$\implies$ $D^b(X)=\left\langle{}A^\perp,A\right\rangle{}$.
\end{lemma}

\begin{proof}
\begin{enumerate}
\item Is clear.
\item $0=F_2 \to F_1 \to F_0=F$, $F_1 \in A$,
$\text{Cone}(F_1 \to F) \in
A^\perp$. We use adjunction as follows:
$$
\text{id} \in
\Hom(\alpha^!(F),\alpha^!(F))=\Hom(\alpha
\alpha^!F,F)
$$
We have the following exact triangle:
$$
\xymatrix{
\alpha \alpha^!F\ar[r]^{\text{counit}}&
F\ar[r]& B
}
$$
{\bf Claim.} $B \in A^\perp$. Let $E \in A$
$$
\xymatrix@C=1em{
\cdots\ar[r]&\Hom(\alpha(E),\alpha
\alpha^!(F))\ar[dr]^{\cong}\ar[rr]&&
\Hom(\alpha(E),F)\ar[dl]_{\cong}\ar[r]&
\Hom(\alpha(E),B)\\
& & \Hom(E,\alpha^!(F))
}
$$
which implies that $\Hom(\alpha(E),B)=0$.

\end{enumerate}
\end{proof}

\begin{definition}
\label{noncommutative-abelian-surfaces-and-kummer-type-hyperkahler-varieties-definition-left-mutation}
{\it Left mutation} $\mathbb{L}_A$ is defined
as
$$
\alpha \alpha^!F \to F \to \mathbb{L}_AF.
$$
\end{definition}

\begin{exercise}
\label{noncommutative-abelian-surfaces-and-kummer-type-hyperkahler-varieties-exercise-left-mutation}
$A=\left\langle{}E\right\rangle{}$.
$$
R \Hom(E,F)\otimes E \xrightarrow{\text{ev}}F
\to \mathbb{L}_{\left\langle{}E\right\rangle{}}F
$$
\end{exercise}

{\bf Corollary.} $E_1,\ldots,E_n \in D^b(X)$
exceptional objects with
$\text{Ext}^\bullet(E_i,E_j)=0$ for $i>j$ (we
say this is an  {\it exceptional
collection}). Then
$$
D^b(X)=\left\langle{}R_x,E_1,\ldots,E_n\right\rangle{},\qquad
R_X=\left\langle{}E_1,\ldots,E_n\right\rangle{}^\perp
$$

\begin{example}
\label{noncommutative-abelian-surfaces-and-kummer-type-hyperkahler-varieties-example-Fano}
$X$ Fano, $-K_X=rH$, $H$ ample, $r>0$,
$\mathcal{O}_X,\mathcal{O}_X(H),\ldots,
\mathcal{O}_X((r-1)H)$
an exceptional collection.

($\text{Ext}^\bullet(\mathcal{O}(iH),
\mathcal{O}(jH))=
H^{\bullet}(X,\mathcal{O}(j-i)H)=0$) if
$-r<j-i<0$. $\implies$
$$
D^b(X)=\left\langle{}R_X,\mathcal{O}
_X\ldots\mathcal{O}_X(r-i)H\right\rangle{}
$$
\end{example}

\begin{definition}
\label{noncommutative-abelian-surfaces-and-kummer-type-hyperkahler-varieties-definition-Serre-functor}
A {\it Serre functor} for a $\Delta$-category
$\mathcal{D}$ is an autoeq $S_{\mathcal{D}}$
such that
$$
\Hom_{\mathcal{D}}(E,F)^\vee \cong
\Hom(F,S_{\mathcal{D}}(E)).
$$
functorially in $E,F \in \mathcal{D}$.
\end{definition}

\begin{remark}
\label{noncommutative-abelian-surfaces-and-kummer-type-hyperkahler-varieties-remark-uniqueness-of-Serre-functor}
It is unique if it exists.
\end{remark}

The following example explains why this is
called the Serre functor --- it is a
generalization of Serre duality.

\begin{example}
\label{noncommutative-abelian-surfaces-and-kummer-type-hyperkahler-varieties-nasak-label-012}
$X$ smooth projective variety of dimension
$n$. Consider
$$
S_{D^b(X)}=(-\otimes \omega_X)[n]
$$
eg. $E$ locally free on $X$.
$$
\begin{aligned}
H^{i}(X,E)&=\Hom_{D^b(X)}(\mathcal{O}_X,E[i])
\\
& =\Hom_{D^b(X)}(E[i],\omega_X(n])^* \\
&=\Hom(\mathcal{O},E^\vee \otimes
\omega_X[n-i])\\
&=H^{n-i}(E^{\vee}\otimes \omega_X)^*
\end{aligned}
$$
\end{example}

\begin{definition}
\label{noncommutative-abelian-surfaces-and-kummer-type-hyperkahler-varieties-definition-Calabi-Yau-category}
$\mathcal{D}$ is
\begin{itemize}
\item {\it Calabi-Yau category} of dimension
$n$ if $S_{\mathcal{D}}\cong[n]$.
(I.e. the Serre functor is just shifting by
$n$.)
\item {\it Fractional Calabi-Yau category} of
dimension $p/q$ if
$S^q_{\mathcal{D}}\cong[p]$ (where the
exponent $q$ just means composing the functor
$S_{\mathcal{D}}$ $q \in \mathbb{Z}$ times).
\end{itemize}
\end{definition}

\begin{theorem}[Kuznetsov]
\label{noncommutative-abelian-surfaces-and-kummer-type-hyperkahler-varieties-theorem-Kuznetsov}
$X \subset \mathbb{P}^n$ smooth Fano
hypersurface of degree $d \leq n$.
$$
D^b(X)=\left\langle{}\text{Ku}_X,\mathcal{O}_X,
\ldots,
\mathcal{O}_X(n-d)\right\rangle{}
$$
Then $\text{Ku}_X$ has a Serre functor $S$
with
$S^{d/c}\cong\left[\frac{(n+1)(d-2)}{c}
\right]
$
where $c=\text{gcd}(d,n+1)$.
\end{theorem}

\begin{example}
\label{noncommutative-abelian-surfaces-and-kummer-type-hyperkahler-varieties-example-cubic-3-and-4-folds}
$d=3$, a cubic 3-fold,  $n=4$,  $S^3=[5]$,
cubic 4-fold $n=5$,  $S=[2]$.
\end{example}

\begin{theorem}[Kuznetsov]
\label{noncommutative-abelian-surfaces-and-kummer-type-hyperkahler-varieties-theorem-Kuznetsov-2}
$X$ cubic 4-fold,
\begin{enumerate}
\item $\exists $ special $X$ such that
$\text{Ku}_X \simeq D^b(S)$, where $S$ is
a K3 surface.
\item $X$ very general $\implies$
$\text{Ku}_X \not \simeq D^b(\text{var})$.
\end{enumerate}
\end{theorem}

We shall not prove this theorems. However, by
the fourth lecture we may give the idea of
their proofs.

\medskip\noindent
{\bf Conjecture (Kuznetsov).} Cubic 4-fold
$X$ is rational if and only if $\text{Ku}_X
\simeq D^b(\text{K3 surface})$.

\medskip\noindent
[KKPY] very general $X$ is not rational.

\medskip\noindent
Let's go back and show how to prove the
statement in Example
\ref{noncommutative-abelian-surfaces-and-kummer-type-hyperkahler-varieties-example-cubic-3-and-4-folds} about
the cubic 4-fold.

\begin{lemma}
\label{noncommutative-abelian-surfaces-and-kummer-type-hyperkahler-varieties-lemma-cubic-4-fold}
Assume that
$\mathcal{D}=\left\langle{}A^{\perp},A\right\rangle{}$,
$S_{\mathcal{D}}\exists $. Then
$S^{-1}_{A^\perp}=\mathbb{L}_A\circ
S^{-1}_{\mathcal{D}}|_{A^\perp}$.
\end{lemma}

\begin{proof}
$F,G \in A^\perp$,
$$
\begin{aligned}
\Hom_{A^\perp}(F,G)
&\cong\Hom_{\mathcal{D}}(G,S_{\mathcal{D}}(F)
)
^\vee\\
&\cong\Hom_{\mathcal{D}}(S^{-1}_{\mathcal{D}}
(G),F)^\vee\\
&\cong
\Hom_{\mathcal{D}}(\mathbb{L}
_A(\underbrace{S^{-1}_{\mathcal{D}}(G))}
_{\in A^\perp},F)^\vee
\end{aligned}
$$
\end{proof}

Now: cubic 4-fold $X \subset \mathbb{P}^5$,
$$
D^b(X)=\left\langle{}\text{Ku}_X,\mathcal{O},
\mathcal{O}(1),\mathcal{O}(2)\right\rangle{}
$$
The lemma tells us that
$$
S_{\text{Ku}_X}^{-1}
=\mathbb{L}_{\left\langle{}\mathcal{O},\mathcal{O}(1)
,
\mathcal{O}(2)\right\rangle{}}
\circ (- \otimes \mathcal{O}(3)[-4])
$$

\medskip\noindent
{\bf Key observation 1.}
$X\overset{i}{\hookrightarrow}\mathbb{P}^5$,
for $F \in D^b(X)$ there exists an exact
$\Delta$:
$$
i^*i_*F \to F \to F \otimes
\mathcal{O}_X(-3)[2]
$$
Assume $F \in \text{Ku}_X$, apply
$S^{-1}_{\text{Ku}}$.

$$
\mathbb{L}_{\left\langle{}\mathcal{O},\mathcal{O}(1),
\mathcal{O}(2)\right\rangle{}}
\circ (i^* i_* F \otimes \mathcal{O}(3)[-4])
\to S_{\text{Ku}}^{-1}(F)
\to
\underbrace{\mathbb{L}_{\left\langle{}\ldots\right\rangle{}}
\circ(F[-2])}
_{\cong F[-2]\text{, as }F\in \text{Ku}_X}
$$

\medskip\noindent
{\bf Key observation 2.} $i^*i_*F \otimes
\mathcal{O}(3) \in
\left\langle{}\mathcal{O},\mathcal{O}(1),\mathcal{O}
(2)\right\rangle{}$,
so the first term
vanishes. (It's enough to show that $i_*F
\otimes \mathcal{O}(3) \in
\left\langle{}\mathcal{O}_{\mathbb{P}^5},\mathcal{O}
_{\mathbb{P}^5}(1),
\mathcal{O}_{\mathbb{P}^5}(2)\right\rangle{}
\iff i_*F \in
\underbrace{\left\langle{}\mathcal{O}(-3),\mathcal{O}
(-2),\mathcal{O}(-1)\right\rangle{}}
_{=\left\langle{}\mathcal{O},\mathcal{O}(1),
\mathcal{O}(2)\right\rangle{}^\perp}
\subset D^b(\mathbb{P}^5)$.)

\begin{align*}
R\Hom\left(
\begin{matrix}
\mathcal{O}\\ \mathcal{O}(1)\\ \mathcal{O}(2)
\end{matrix},i_*F
\right)&=0
\\
R\Hom\left(
i^*\begin{matrix}
\mathcal{O}\\ \mathcal{O}(1)\\ \mathcal{O}(2)
\end{matrix},F
\right)&=0\qquad \text{since $F \in
\text{Ku}_X$}
\end{align*}
and the left-hand-sides on the last two
equations are $\cong$.

\begin{proof}[Proof of Key observation 1
assuming $F$ is a sheaf.]
Take $i^*i_*F \to F \to G$ and consider its
pushforward
$$
i_*i^*i_*F \to i_* F \to i_*G
$$
Note that
$$
\begin{aligned}
i_*i^*i_*F \simeq
i_*F\overset{L}{\otimes}_{\mathcal{O}
_{\mathbb{P}^5}}
&\cong \left[
\underbrace{i_* F \otimes\mathcal{O}
_{\mathbb{P}^5}(-x)}_{\text{deg}-1}
\xrightarrow{\vee}
i_*F\otimes \mathcal{O}_{\mathbb{P}^5}\right]
\\
&  \simeq i_*F \oplus i_*F\otimes
\mathcal{O}_{\mathbb{P}^5}(-x)[1]\\
&\xrightarrow{\substack{ \text{proj.
to}\\\text{$1^{\text{st}}$ factor}}}i_*F\\
&  \implies i_*G \simeq i_* F \otimes
\mathbb{P}_{\mathbb{P}^5(-x)[2]}\\
& \implies G \simeq F \otimes
\mathcal{O}_X(-x)[2]\qquad i_*\text{ f.f.}
\end{aligned}
$$

\end{proof}

\section{Inducing $t$-structure and stability
conditions}
\label{noncommutative-abelian-surfaces-and-kummer-type-hyperkahler-varieties-nasak-label-003}

\begin{example}
\label{noncommutative-abelian-surfaces-and-kummer-type-hyperkahler-varieties-example-cubic-3-fold}
$X$ cubic 3-fold,
$D^b(X)=\left\langle{}\text{Ku}_X,\mathcal{O},
\mathcal{O}(1)\right\rangle{}$.
\end{example}

\noindent
{\bf Goal.} Use $t$-structure/(weak)
stability conditions on $D^b(X)$ to induce a
$t$-structure/stability conditions on $Ku_X$.

\medskip\noindent
\begin{proposition}[Key proposition]
\label{noncommutative-abelian-surfaces-and-kummer-type-hyperkahler-varieties-proposition-key}
Let
$\mathcal{D}=\left\langle{}R,E_1,\ldots,E_n\right\rangle{}$
be a semiorthogonal decomposition. Assume
that $\mathcal{A} \subset \mathcal{D}$ is the
heart of a bonded $t$-structure such that
\begin{enumerate}
\item $\forall i,$ $E_i \in \mathcal{A}$.
\item $\forall i$, $S_{\mathcal{D}}(E_i) \in
\mathcal{A}[i]$.
\end{enumerate}
Then $\mathcal{A}:=\mathcal{A}\cap R$ is the
heart of a bounded $t$-structure on $R$.
(Concretely, this means for $F \in R$, view
it as an object in $\mathcal{D}$ and take
cohomology with respect to the $t$-structure,
we need $H_{\mathcal{A}}^q(F) \in R$ $\forall
q$.)
\end{proposition}

[BLMS] using spectral sequence.

\begin{example}
\label{noncommutative-abelian-surfaces-and-kummer-type-hyperkahler-varieties-example-cubic-3-fold-continued}
For $X$ cubic 3-fold and
$D^b(X)=\left\langle{}\text{Ku}_X,\mathcal{O},
\mathcal{O}(1)\right\rangle{}$,
$$
\begin{aligned}
\left(\textit{Coh}^{-\frac{1}{2}}(X),
Z_{\alpha,-\frac{1}{2}}
=\frac{1}{2}\alpha^2H^2 ch_0
-ch_2^{-\frac{1}{2}}+iHch_1^{-\frac{1}{2}}
\right)
\end{aligned}
$$
\end{example}

\begin{exercise}
\label{noncommutative-abelian-surfaces-and-kummer-type-hyperkahler-varieties-exercise-duplicate-2}
\begin{itemize}
\item For all $\alpha>0$, $\mathcal{O}$,
$\mathcal{O}(1)$, $\mathcal{O}(-2)[1]$,
$\mathcal{O}(-1)[1]\in
\textit{Coh}^{-\frac{1}{2}}$.
\item For all $\alpha \sim 0$,
$$
\begin{aligned}
\mu_{\alpha,-\frac{1}{2}}(\mathcal{O}(-2)[1])
&<\mu_{\alpha,-\frac{1}{2}}(\mathcal{O}(-1)
[1]
)<0\\
&<\mu_{\alpha,-\frac{1}{2}}(\mathcal{O})
<\mu_{\alpha,-\frac{1}{2}}(\mathcal{O}(1)).
\end{aligned}
$$
So tilt $\textit{Coh}^{-\frac{1}{2}}(X)$ one
more time w.r.t.
$\mu_{\alpha,-\frac{1}{2}}=0$. We can apply
Key Proposition
\ref{noncommutative-abelian-surfaces-and-kummer-type-hyperkahler-varieties-proposition-key}.
\end{itemize}
\end{exercise}

\begin{example}
\label{noncommutative-abelian-surfaces-and-kummer-type-hyperkahler-varieties-example-cubic-4-fold}
How about cubic 4-fold $Y$? Idea: embed
$Ku_X\hookrightarrow
\underbrace{D^b(\mathbb{P}^3,\mathcal{D}_0)}
_{\text{conic fibration}}$.
\end{example}

\begin{theorem}
\label{noncommutative-abelian-surfaces-and-kummer-type-hyperkahler-varieties-theorem-exists-SC-for-cubic-4-fold}
There exist stability conditions on $Ku_Y$
for a cubic 4-fold $Y$.
\end{theorem}

\medskip\noindent
Now we shall give the idea of why Key
Proposition \ref{noncommutative-abelian-surfaces-and-kummer-type-hyperkahler-varieties-proposition-key} holds.

General question modeling Key Proposition
\ref{noncommutative-abelian-surfaces-and-kummer-type-hyperkahler-varieties-proposition-key}: $F: \mathcal{C} \to
\mathcal{D}$ exact functor between
triangulated categories. Say
$\mathcal{D}^{heart}$ is the heart of a
$t$-structure on $\mathcal{D}$. Under what
condition is
$$
\mathcal{C}^{heart}=\{E \in \mathcal{C}|F(E)
\in \mathcal{D}^{heart}\}
$$
the heart of a $t$-structure on
$\mathcal{C}$?

\begin{example}
\label{noncommutative-abelian-surfaces-and-kummer-type-hyperkahler-varieties-example}
Let $X,Y$ be smooth projective and $f:X \to
Y$ a finite a morphism.
$$
\xymatrix{
D^b(X)\ar[r]_{f_*} \ar@{^{(}->}[d]
&D^b(Y)
\ar@/_{1.5pc}/[l]_{f^*}\ar@{^{(}->}[d]\\
D(Q\textit{Coh}(X))=D_{qc}(X)\ar[r]_{f_*}
&\underbrace{D_{qc}(Y)
}_{\substack{\text{admits
small} \\ \text{direct sum}}}
\ar@/_{1.5pc}/[l]_{f^*}
}
$$
\end{example}

\begin{situation}
\label{noncommutative-abelian-surfaces-and-kummer-type-hyperkahler-varieties-situation}
$\tilde{\mathcal{C}},\tilde{\mathcal{D}}$
triangulated categories admitting direct
sums. $\mathcal{C}\subset
\tilde{\mathcal{C}}$, $\mathcal{D}
\subset\tilde{\mathcal{D}}$ full essentially
small triangulated subcategories.
$F:\tilde{\mathcal{C}}\to
\tilde{\mathcal{D}}$ exact functor such that
\begin{itemize}
\item $F$ commutes with direct sums.
\item $F$ has left adjoint
$\mathcal{G}:\tilde{\mathcal{D}}\to
\tilde{\mathcal{C}}$.
\item $F(\mathcal{C}) \subseteq \mathcal{D}$
and if $E \in \tilde{\mathcal{C}}$
is such that $F(E) \in \mathcal{D}$ then $E
\in \mathcal{C}$.
\item $G(\mathcal{D}) \subseteq \mathcal{
C}$.
\item For $E \in \mathcal{C}$, if $F(E)=0
\implies E=0$.
\end{itemize}
\end{situation}

\begin{definition}
\label{noncommutative-abelian-surfaces-and-kummer-type-hyperkahler-varieties-nasak-label-004}
Let $\mathcal{D}^{heart}$ be the heart of a
$t$-structure on $\mathcal{D}$. Define
$$
\begin{aligned}
\mathcal{D}^{\leq 0}&=\{E \in
\mathcal{D}|H^{i}(E)=0, i>0\}
\mathcal{D}^{\geq 0}&=\{E \in
\mathcal{D}|H^{i}(E)=0, i<0\}
\end{aligned}
$$
\end{definition}

\begin{theorem}[Polischuk]
\label{noncommutative-abelian-surfaces-and-kummer-type-hyperkahler-varieties-theorem-Polischuk}
Let $\mathcal{D}^{heart}$ be the heart of a
bounded $t$-structure on $\mathcal{D}$. If
$FG:\mathcal{D} \to \mathcal{D}$ is right
$t$-exact (i.e. $FG(\mathcal{D}^{\leq
0})\subseteq \mathcal{D}^{\leq 0}$) then
$$
\mathcal{C}^{heart}=\{E \in \mathcal{C}|F(E)
\in \mathcal{D}^{heart}\}
$$
is the heart of a bounded $t$-structure on
$\mathcal{C}$.
\end{theorem}

\begin{example}
\label{noncommutative-abelian-surfaces-and-kummer-type-hyperkahler-varieties-example-right}
$f_*:D^b(X) \to D^b(Y)$, $f_*f^*(E)=E\otimes
f_*\mathcal{O}_x$. If this is right
$t$-exact, then can induce $t$-structures,
e.g. $f:H \hookrightarrow Y$ hypersurface.
$$
E \to E \otimes f_*\mathcal{O}_H\to E\otimes
\mathcal{O}_T(-H)[1]
$$
See recent paper by Chengi Li ``Real
reduction \ldots{}''
\end{example}

\begin{example}
\label{noncommutative-abelian-surfaces-and-kummer-type-hyperkahler-varieties-example-Kuznetsov}
$D^b(X)=\left\langle{}R,E\right\rangle{}$, $i:R \to D^b(X)$.
See [Kuznetsov]. There exists a
semiorthogonal decomposition of $D_{qc}(X)
\supset \tilde{R}$, $i^* \to i$,
$$
\mathbb{L}_{\left\langle{}E\right\rangle{}}=ii^* :D^b(X) \to
D^b(X)\qquad \text{right exact?}
$$
\end{example}

The reader is really encouraged to solve the
following exercise.

\begin{exercise}
\label{noncommutative-abelian-surfaces-and-kummer-type-hyperkahler-varieties-exercise-right-exact}
Let $E \in \mathcal{D}^{heart}$, $S(E) \in
\mathcal{D}^{heart}[1]$, then $\mathbb{L}_E$
is right exact.
\end{exercise}

\noindent
{\bf Key homological algebra ingredient.} Let
$\tilde{\mathcal{C}}^{\leq 0}$ be the
smallest full subcategory of
$\tilde{\mathcal{C}}$ that
\begin{itemize}
\item contains $G(\mathcal{D}^{\leq 0})$,
\item is closed under direct sum, extension,
positive shift $[n]$, $n>0$.
\end{itemize}
$$
\tilde{\mathcal{C}}^{\geq 0}:=\{E \in
\tilde{\mathcal{C}}|
\Hom(\tilde{\mathcal{C}}^{\leq 0}[1],E)=0\}.
$$
Then this is a $t$-structure on
$\tilde{\mathcal{C}}$.

\begin{example}
\label{noncommutative-abelian-surfaces-and-kummer-type-hyperkahler-varieties-example-curve}
Let $X$ be a smooth projective curve and $p
\in X$ a point. Then $D_{qc}(X)^{\leq 0}$
contains $\mathcal{O}_p$ and is closed under
\ldots{}.
$$
D_{qc}(X)^{\geq 0}j_*
\mathcal{O}_U/\mathcal{O}_X[-1] \to
\mathcal{O}_X \to j_*\mathcal{O}_U
$$
\end{example}

\section{Moduli of objects in derived
category}
\label{noncommutative-abelian-surfaces-and-kummer-type-hyperkahler-varieties-nasak-label-005}

\subsection{Moduli on a variety}
\label{noncommutative-abelian-surfaces-and-kummer-type-hyperkahler-varieties-subsection-moduli-on-a-variety}

Let $X$ be a scheme of finite type over
$\mathbb{C}$.

\begin{definition}
\label{noncommutative-abelian-surfaces-and-kummer-type-hyperkahler-varieties-definition-perfect-complex}
An object $E \in D^b(X)$ is a {\it perfect
complex} if locally on $X$, $E$ is isomorphic
to a bounded cx of locally free sheaves
$D_{\text{perf}}(X)\subset D^b(X)$
subcategory of perfect cx. [?]
\end{definition}

\begin{remark}
\label{noncommutative-abelian-surfaces-and-kummer-type-hyperkahler-varieties-remark-perfect}
\begin{enumerate}
\item $D_{perf}(-)$ is preserved under
derived pullback.
\item $X$ smooth $\implies$
$D_{\text{perf}}(X)=D^b(X)$.
\end{enumerate}
\end{remark}

\begin{definition}
\label{noncommutative-abelian-surfaces-and-kummer-type-hyperkahler-varieties-definition-universally-gluable}
Given $T$ scheme of finite type over
$\mathbb{C}$ and $E \in
D_{\text{perf}}(X\times T)$, say $E$ is {\it
universally gluable} over $T$ if for all $t
\in T(\mathbb{C})$, $\text{Ext}^i(E_t,E_t)=0$
if $i<0$, where  $E_t$ is the derived
pullback.
\end{definition}

\begin{example}
\label{noncommutative-abelian-surfaces-and-kummer-type-hyperkahler-varieties-example-locally-free}
Locally free on $X \times T$.
\end{example}

We will consider the following functor to
formulate the moduli problem:

\begin{definition}
\label{noncommutative-abelian-surfaces-and-kummer-type-hyperkahler-varieties-nasak-label-006}
Let $X$ be a smooth projective variety over
$\mathbb{C}$. The {\it moduli stack of
perfectly universally gluable objects} is the
functor
$$
\begin{aligned}
\mathcal{M}_{\text{proj}}(X):
(\Sch^{ft}/\mathbb{C})^{op}
&\longrightarrow \mathit{Gpds} \\
T &\longmapsto  \{E \in
D_{\text{perf}}(X\times T)\text{ uniersally
gluable}/T\}
\end{aligned}
$$

\end{definition}

\begin{theorem}[Lieblich]
\label{noncommutative-abelian-surfaces-and-kummer-type-hyperkahler-varieties-theorem-Lieblich}
$\mathcal{M}_{\text{proj}}(X)$ is an
algebraic stack locally of finite type over
$\mathbb{C}$.
\end{theorem}

\subsection{Moduli of $\sigma$-semistable
objects}
\label{noncommutative-abelian-surfaces-and-kummer-type-hyperkahler-varieties-nasak-label-007}

Let $X$ be a smooth projective variety over
$\mathbb{C}$. Fix
$\sigma=(A,Z)\in\text{Stab}_{\Lambda}(X)$,
$v:K_0(X) \to \Lambda$. Fix $w \in
\Lambda$, define
$$
\begin{aligned}
\mathcal{M}_{\sigma}(w):(\Sch^{ft}/\mathbb{C}
)
^{op}
&\longrightarrow  \mathit{Gpds} \\
T &\longmapsto \left\{
\substack{\text{perfect complex $E$ on
$X\times T$ s.t.}
\\ \text{$\forall t \in T(\mathbb{C})$, $E_t$
is in $\mathcal{A}$}\\
\text{$\sigma$-ss, $V(E_t)=w$}} \right\}
\end{aligned}
$$

\medskip\noindent
{\bf Hard problems.}
\begin{enumerate}
\item Are $\mathcal{M}_{\sigma}^s(w) \subset
\mathcal{M}_{\sigma}(w)
\subset \mathcal{M}_{\text{proj}}(X)$.
\begin{enumerate}
\item Open.
\item Finite type?
\end{enumerate}
Known for $X$ surface, $\sigma$ via tilting
(Toda).
\begin{itemize}
\item \text{[}Halper-Leistner-Robotis '25].
\item Proper SC.
\end{itemize}
\item Do there exist maps to schemes
$$
\xymatrix{
\mathcal{M}_{\sigma}(w)&M_\sigma(w)\\
\mathcal{M}_{\sigma}^s(w)\ar[r]\ar@{^{(}->}
[u]
&
M_{\sigma}^s(w)\ar@{^{(}->}[u]
}
$$
where $M_\sigma(w)$ is projective and
paramatrizes s-equivalent classes.
\end{enumerate}

\begin{definition}
\label{noncommutative-abelian-surfaces-and-kummer-type-hyperkahler-varieties-definition-good-moduli-space}
Let $\mathcal{M}$ be an algebraic stack, $M$
algebraic space. $\pi:\mathcal{M} \to M$ is a
{\it good moduli space} if
\begin{enumerate}
\item $\pi_* :\text{QCoh}(\mathcal{M}) \to
\text{QCoh}(M)$ is exact.
\item $\mathcal{O}_M \xrightarrow{\simeq
}\pi_*\mathcal{O}_{\mathcal{M}}$
\end{enumerate}
\end{definition}

\begin{example}
\label{noncommutative-abelian-surfaces-and-kummer-type-hyperkahler-varieties-example-good-moduli-space}
$X$ affine, $G$ linear reductive $G
\mathbb{y} X$, $[X/G] \to X/\!/G$ is a good
moduli space. {\bf Exercise:} think this
through!
\end{example}

\begin{theorem}
[Alper-Halpern-Leistor-Heisloth]
\label{noncommutative-abelian-surfaces-and-kummer-type-hyperkahler-varieties-theorem-AHLH}
Good moduli space exists if the stack
satisfies $\Theta$-reductivity and
S-completeness.
\end{theorem}

Where, roughly,
\begin{itemize}
\item $\Theta$-reductivity means ``HN
filtration specializes'',
\item S-completeness means: two
specializations of a family of objects differ
by
``elementary modification''.
\end{itemize}

\begin{example}[Elementary modification I]
\label{noncommutative-abelian-surfaces-and-kummer-type-hyperkahler-varieties-example-elementary-modification}
$\mathbb{P}_0=\mathbb{P}^1 \times \{0\} =
\mathbb{P}^1 \times \mathbb{A}^1$,
$$
\xymatrix{
0\ar[r]&  K \ar[r]&
\mathcal{O}_{\mathbb{P}^1\times\mathbb{A}^1}
^{\oplus
2}
\ar[rr]\ar@{->>}[rd]&  &
\mathcal{O}_{\mathbb{P}^0}(1) \ar[r]&  0\\
&  &  &  \mathcal{O}_{\mathbb{P}_0^{\oplus
2}}\ar@{->>}[ur]
}
$$
\end{example}

\begin{example}[Elementary modification II]
\label{noncommutative-abelian-surfaces-and-kummer-type-hyperkahler-varieties-example-elementary-modification-2}
$C \times \mathbb{A}^1$, $E$ rank 2 vector
bundle such that
$$
\xymatrix{
0\ar[r]&\mathcal{L}_1\ar[r]&E|_{co}\ar[r]&
\mathcal{L}\ar[r]&0
}
$$
for
$\text{deg}\mathcal{L}_1=\text{deg}
\mathcal{L}
^2$.
$$
\xymatrix{
0\ar[r]&  K \ar[r]& E\ar[rr]\ar@{->>}[rd]
&  &  \mathcal{L}_2 \ar[r]&  0\\
&  &  &  E|_{c_0}\ar@{->>}[ur]
}
$$
\end{example}

\medskip\noindent
{\bf Upshot.} $\mathcal{M}_{\sigma}(w) \to
M_{\sigma}(w)$ good moduli exist.

\begin{itemize}
\item Bayer-Macri: there exists a (numerical)
divisor class on $M_{\sigma}(w)$
that intersects every curve positively.
\end{itemize}

\subsection{Mukai's theorem for K3 category}
\label{noncommutative-abelian-surfaces-and-kummer-type-hyperkahler-varieties-nasak-label-008}

\begin{proposition}
\label{noncommutative-abelian-surfaces-and-kummer-type-hyperkahler-varieties-proposition-4-fold-Mukai-lattice}
$X$ cubic 4-fold. Then there exists a Mukai
lattice (cf Subsection
\ref{noncommutative-abelian-surfaces-and-kummer-type-hyperkahler-varieties-subsection-Mukai-lattice})
$\tilde{H}(Ku_X';\mathbb{Z})$ with
\begin{itemize}
\item Weight 2 Hodge structure,
\item Pairing compatible with $-\chi$.
\end{itemize}
\end{proposition}

\begin{theorem}[BLMNPS]
\label{noncommutative-abelian-surfaces-and-kummer-type-hyperkahler-varieties-theorem-BLMNPS}
Let $0 \neq  v \in
\tilde{H}_{\text{alg}}(Ku_X;\mathbb{Z})$
primitive, $\sigma$ generic w.r.t. $v$. Then
$\mathcal{M}_{\sigma}(v)$ has a good moduli
space $M_{\sigma}(v)\neq \emptyset$ smooth
projective hyperk\"ahler variety of dimension
$v^2+2$, deformation equivalent to
$K3^{\left[ \frac{v+1}{2 } \right] }$.
\end{theorem}

(The deepest part of this theorem is showing
that $M_{\sigma}(v)$ is not empty!)

\medskip\noindent
There exists a rank 2 lattice contained in
$\tilde{H}_{alg}(Ku_X;\mathbb{Z})$ for any
$X$ $\rightsquigarrow $ $M_\sigma(v)$ deforms
in $20$-dimensional family.

\section{Noncommutative abelian surfaces and
Hyperk\"ahler of Kummer type}
\label{noncommutative-abelian-surfaces-and-kummer-type-hyperkahler-varieties-nasak-label-009}

Joint work with Boyer, Perry, Petrusi ('25?)

\subsection{Polarized hyperk\"ahler}
\label{noncommutative-abelian-surfaces-and-kummer-type-hyperkahler-varieties-subsection-polarized-hyperkahler}

Recall Definition
\ref{noncommutative-abelian-surfaces-and-kummer-type-hyperkahler-varieties-definition-hyperkahler-variety}. In the
past lectures we have seen that the spaces
$K3^{[n]}$ are examples.

Let $A$ be an abelian surface ($A \cong
\mathbb{C}^2/\Lambda
\hookrightarrow\mathbb{P}^N$). The fact that
abelian surfaces are not abelian makes
$M_H(v)$, the moduli of $H$-semistable
sheaves of topological class $V$ on $A$, not
abelian. So we consider
$$
\xymatrix{
M_H(v)\ar[d]_{\text{alb}}\\
\hat{A}\times A
}
$$
where we must assume that $v$ is primitive,
$v^2>0$ and $\sigma/H$ is $v$-gen. Then
$M:=\text{alb}^{-1}(0)$ is a hyperk\"ahler
variety called {\it generalized Kummer}.

Now consider the even cohomology
$H^{ev}(Z,\mathbb{Z})$. Let
$$
H^{ev}(A;\mathbb{Z})\supset V^\perp
\xrightarrow{\theta,\simeq }H^{2}(M)
$$
(Torelli theorem, [\ldots{}, Verbitsky])

\begin{remark}
\label{noncommutative-abelian-surfaces-and-kummer-type-hyperkahler-varieties-nasak-label-013}
$$
\Lambda=\left\langle{}v,\theta^{-1}(H)\right\rangle{}\subset
H^{ev}_{alg}(A,\mathbb{Z})
$$
positive definite rank 2 lattice.
\end{remark}

{\bf Goal.} Fix $\Lambda$. Look for:
$\Lambda$-polarized deformation of $D^b(A)$
(i.e. $\Lambda$ remains alg.).

\subsection{Deformations via equivariant
categories}
\label{noncommutative-abelian-surfaces-and-kummer-type-hyperkahler-varieties-nasak-label-010}

Consider an involution on an abelian surface,
$\mathbb{Z}/2 \mathbb{y} A$. This induces an
action $ \mathbb{Z}/2 \mathbb{y}D^b(A)$.

There are 16 singular points in the quotient
$A/(\mathbb{Z}/2)$. The resolution
$\widetilde{A/(\mathbb{Z}/2)=S}$ is the
so-called {\it Kummer K3}. (Another way to
see this is considering the quotient stack
$A$ by  $\mathbb{Z}/2$, the result is
``basically'' equivalent to $S$.)

[BKR] $D^b(S)\cong D^b(A)^{\mathbb{Z}/2}$.

[Elagin] $D^b(A) \cong
D^b(S)^{\widehat{\mathbb{Z}/2}}$ where
$\widehat{\mathbb{Z}/2}$ is an involution
acting on the derived category of $S$, and we
explain it next: there are 16 exceptional
divisors $E_1,\ldots,E_{16}\subset S$. The
involution is
$$
(- \otimes \mathcal{L})\cdot
\prod_{i=1}^{16}ST_{\mathcal{O}_{E_i}(-1)}.
$$

\medskip\noindent
{\bf Key 1.} In order to deform $D^b(A)$, we
only need to deform $D^b(S)$ together with
its involution.

\medskip\noindent
{\bf Definition/Proposition.}
\begin{itemize}
\item The derived Kummer lattice
$\tilde{K}\subset\tilde{H}^*(S;\mathbb{Z})$
is the $(-1)$-eigenspace of
$\widehat{\mathbb{Z}/2}$ action on
$\tilde{H}(S;\mathbb{Z})$.
\item Spanned by $(0,2E_i,-1)$.
\end{itemize}

\medskip\noindent
{\bf General theory (for deforming K3
surfaces).} Due to Huybrechts, Toda,
Addington-Thomas. To deform $D^b(S)$ and its
$\widehat{\mathbb{Z}/2}$-action is equivalent
to deforming $D^b(S)$ such that $\tilde{K}$
remains algebraic.

What need to remain algebraic under
deformation?
\begin{itemize}
\item $\tilde{K}^{16}$ (so that the
involution deforms).
\item $\Lambda^2$.
\end{itemize}
[The dimension of the moduli space must then
be $22-16=4$.]

\begin{lemma}
\label{noncommutative-abelian-surfaces-and-kummer-type-hyperkahler-varieties-lemma}
$$
\left\langle{}\tilde{K},\Lambda\right\rangle{}\supset
U=\begin{pmatrix}
0&1\\
1&0
\end{pmatrix}
$$
$$
\underbrace{\left\langle{}\tilde{K},\Lambda\right\rangle{}}
_{(2,16)}
=\underbrace{U}_{(1,1)
}\oplus\underbrace{\Lambda_T}_{(1,15)}
$$
\end{lemma}

\medskip\noindent
{\bf Key 2.} There exists
$$
\begin{aligned}
\Phi: D^b(S) &\longrightarrow D^b(T) \\
\tilde{H}(S) &\qquad  \tilde{H}(T)\\
\left\langle{}\tilde{K}, U\right\rangle{}\quad U &\longmapsto
H^0 \oplus H^4\\
U \oplus \Lambda_T\quad \Lambda_T
&\longmapsto \Lambda'_T \subset H^2
\end{aligned}
$$
where $ T$ is a K3 surface and $\tilde{H}(S)
\supset \left\langle{}
\tilde{K},\Lambda\right\rangle{}=U\oplus \Lambda_T$.
Deform $T$ as a
$\Lambda_T$-polarized K3 surface.

Since $\tilde{K}$ remains algebraic,
invariant deformations w/ $D^b(T)$ take the
equiv. category.

\begin{theorem}
\label{noncommutative-abelian-surfaces-and-kummer-type-hyperkahler-varieties-nasak-label-011}
There exists a quasi-finite dominant map
$$
\begin{aligned}
U  &\longrightarrow
\text{Mod}_{\Lambda_T}^{K3}
\end{aligned}
$$
such that $u \in U$.
$D^b(T_u)^{\widehat{\mathbb{Z}/2}}=$
4-dimensional family of categories deforming.
\end{theorem}

\begin{theorem}
\label{noncommutative-abelian-surfaces-and-kummer-type-hyperkahler-varieties-theorem}
For each $u \in U$, there exists an
involution $\widehat{\mathbb{Z}/2}$ stability
condition on $D^b(T_u)$.
\end{theorem}

From these two theorems we obtain
$$
\xymatrix{
\substack{\text{Kummer} \\\text{type}\\
\text{HK}}\to
\mathcal{M}_{\sigma}(D^b(T_u)
^{\widehat{\mathbb{Z}/2}},v)\ar[d]^{alb}\\
\text{abelian 4-fold}
}
$$

\begin{definition}
\label{noncommutative-abelian-surfaces-and-kummer-type-hyperkahler-varieties-definition-Weil-type}
An abelian 4-fold $W$ is of {\it Weil type}
if $\mathbb{R}(\sqrt{-d}) \hookrightarrow
\text{End}_{\mathbb{Q}}(W)$ for
$d\in\mathbb{Z}_+$ along with some mild
condition.
\end{definition}

Let $\Lambda=\left\langle{}v,w\right\rangle{}$. Consider the
abelian category $\mathcal{A}$ nc abelian
surface. [In fact,
$\text{Aut}^0(\mathcal{A})$ is an abelian
4-fold.]
$$
\xymatrix{
M_\sigma(v)\ar[dr]_{alb_v}&
\text{Aut}^0(\mathcal{A})\ar[l]\ar[r]
\ar@/^{.5pc}/[d]^{c_w}\ar@/^{-0.5pc}/[d]
_{c_v}
& M_\sigma(w)\ar[dl]^{alb_w}\\
&\text{Alb}(\mathcal{A})
}
$$
\begin{theorem}
\label{noncommutative-abelian-surfaces-and-kummer-type-hyperkahler-varieties-theorem-duplicate-2}
$c_w^{-1} \circ c_v$ is an endomorphism of
Weil type.
\end{theorem}
% --- End noncommutative-abelian-surfaces-and-kummer-type-hyperkahler-varieties.tex ---

\part{Computing}
\label{part-computing}

% --- Begin ai.tex ---
\chapter{AI}
\label{chapter-ai}



Formal proof :=
every phrase can be computer checked.

Proof assistant :=
programming language precise enough
to encode each proof step.

Lean. Needs libraries. 
Mathlib if a library of functions and
deftinitions (formalizations)
of mathematical concepts.
Currently = 115L defs vs 232K thms.
It's a language to write proofs that
computer can check.
Proof is written as code;
trusted kernel checks it has no gaps.
If it compiles, it's 100% correct.
Created by Leonardo de Moura
(from PUC-Rio!).

Proof that addition is commutative:

\begin{verbatim}
theorem add_comm( a b : Nat) : a+b=b+a:= by
induction b with
|zero => simp
|succ n ih =>rw[Nat.add_succ,ih,Nat.succ_add]
\end{verbatim}

simp means simplifiy
n+1 is denoted succ

sorry := it's an axiom

Proving := the agentic loop converges.

Example: 14,500 lines of Lean, 25 files.
Make the verifiable proof + check examples.

Claude Code = creative
Codex = more formal

-> there is not a very big infrastructure
for Lean these days!

Idea: combine codex + Lean.
This is what an "automated prover" does,
e.g. Aritstotle: leverage LLM
to create proofs in Lean.
Also AxiomProver.


Need to search mathlib to know what is already
available! Probably there's more
number theory/combinatorics than
geometry/topology (as of today).


% --- End ai.tex ---

% --- Begin computing-systems.tex ---
\chapter{Computing systems}
\label{chapter-computing-systems}



\section{What I am trying to learn}
\label{computing-systems-section-computing-systems-goal}

The immediate goal is not to become a
computer engineer.  It is to understand
enough about computers to make good
choices about my own machines.

In particular, when buying a computer I
want to be able to answer:

\begin{itemize}
\item What work is limited by the CPU?
\item How much RAM do I actually need?
\item When does a GPU matter?
\item What is VRAM, and why does it matter
for local AI?
\item Which parts of the machine can I
replace or upgrade?
\item Which restrictions come from the
hardware, and which come from the
operating system?
\item What does it mean for a computer to
be open?
\item How much of my work can remain local
and private?
\end{itemize}

For my present situation there is one more
constraint: I move around a lot.  So the
main machine has to be a laptop.  A large
desktop workstation may make sense later,
but it is not the present problem.

\section{The basic hardware}
\label{computing-systems-section-computing-systems-hardware}

A computer has several components doing
very different jobs.

\medskip\noindent
{\bf CPU.}
The central processing unit executes most
ordinary programs.  Compiling code,
running a computer algebra system, and
many mathematical computations are
primarily CPU jobs.

A CPU has several cores.  More cores help
when a program can divide its work into
parallel tasks.  A fast individual core is
still important because many programs
cannot parallelize everything.

\medskip\noindent
{\bf RAM.}
RAM is the computer's short-term working
memory.  Programs and data being actively
used live there.

Running out of RAM does not merely mean
that no more files can be stored.  The
operating system starts moving data
between RAM and the much slower SSD.  This
can make a machine feel dramatically
slower.

For my work, RAM is particularly important
because I may simultaneously have a
browser, editor, TeX compiler, terminals,
computer algebra systems, and AI tools
open.

\medskip\noindent
{\bf Storage and SSD.}
The SSD is long-term storage.  It contains
the operating system, programs, source
files, papers, datasets, and models.

An SSD is much slower than RAM but much
faster than an old mechanical hard drive.
Capacity and replaceability are separate
questions: a large soldered SSD can still
be less useful in the long run than a
standard replaceable one.

\medskip\noindent
{\bf GPU.}
The graphics processing unit was designed
for massively parallel graphics
computations.  The same kind of
parallelism is extremely useful in machine
learning.

For ordinary LaTeX, browsing, Git, and
most programming, a powerful GPU is not
very important.  For training or running
many AI models it can become one of the
most important components.

\medskip\noindent
{\bf VRAM.}
A discrete GPU normally has its own memory,
called VRAM.  A model being run on the GPU
has to fit, at least to a large extent,
inside this memory.

Thus, for local AI, it is not enough to
ask whether a machine has a fast GPU.  The
amount of memory available to the GPU can
be the decisive limitation.

Some architectures instead use a common
pool of memory shared by the CPU and GPU.
This is called unified memory.  It changes
the comparison because a model may have
access to a much larger memory pool than
on a laptop GPU with a small fixed amount
of VRAM.

\medskip\noindent
{\bf Cooling and power.}
A fast chip produces heat.  A computer can
only sustain high performance if it can
remove that heat and supply enough power.

This is one reason laptop performance
cannot be judged from the processor name
alone.  The same chip can behave
differently in two machines with different
cooling systems and power limits.

\section{The software stack}
\label{computing-systems-section-computing-systems-software}

The hardware is only the bottom layer.

Very roughly, the stack is

$$
\begin{aligned}
&\text{applications}\\
&\text{operating system}\\
&\text{drivers and kernel}\\
&\text{firmware}\\
&\text{hardware}.
\end{aligned}
$$

The operating system manages files,
memory, processes, devices, networking,
permissions, and the interface between
applications and hardware.

The kernel is the central part of the
operating system.  Drivers teach the
kernel how to communicate with particular
hardware.

Firmware is software stored very close to
the hardware itself.  For example, a
laptop contains firmware controlling parts
of the boot process and many devices.

This distinction matters because
installing Linux changes a large part
of the software stack, but it does not
change the physical machine.  Linux cannot
make soldered RAM replaceable.

\section{What openness means}
\label{computing-systems-section-computing-systems-openness}

I should avoid using the word open as if
it meant only one thing.  There are
several different freedoms.

\begin{itemize}
\item {\bf Software freedom:} can I run,
inspect, modify, and replace the software?
\item {\bf Operating-system freedom:} can
I install another operating system?
\item {\bf Hardware freedom:} can I open
the machine and replace components?
\item {\bf Repair freedom:} can I obtain
parts and repair the machine without the
manufacturer's permission?
\item {\bf Data freedom:} are my files in
portable formats that I can move elsewhere?
\item {\bf Service independence:} does my
workflow continue to work if I stop using
a particular company's cloud services?
\end{itemize}

These freedoms are related but not
identical.  A machine can run open-source
software while having soldered hardware.
A repairable laptop can still run
proprietary software.  A Linux computer
can still depend heavily on cloud services.

\section{Local AI and privacy}
\label{computing-systems-section-computing-systems-local-ai}

One reason I care about hardware now is
local AI.

A local model runs on my own machine.
This can improve privacy because the model
does not inherently require sending the
prompt and documents to an external
server.

But a local model does not by itself prove
that the whole workflow is private.  The
surrounding application may still use
telemetry, cloud synchronization, web
search, or remote APIs.

So the privacy question should always be:

\begin{quote}
Which data leave my machine, through which
program, and for what reason?
\end{quote}

For local AI I should learn to distinguish:

\begin{itemize}
\item model size;
\item numerical precision and
quantization;
\item system RAM;
\item GPU memory or unified memory;
\item CPU inference and GPU inference;
\item local applications and cloud APIs.
\end{itemize}

A powerful laptop can run useful local
models, but a laptop is constrained by
weight, battery, heat, and power.  I should
not expect a portable machine to be a
permanent substitute for a future
workstation.

\section{The Apple question}
\label{computing-systems-section-computing-systems-apple}

My worry about Apple should also be stated
precisely.

The concern is not simply that macOS is
proprietary.  I am worried about becoming
dependent on a system in which the same
company controls the hardware, operating
system, account ecosystem, cloud services,
and an increasing number of AI features.

There are several separate questions.

\medskip\noindent
{\bf Can I install my own software?}
On macOS I can already install a great deal
of software outside Apple's App Store.
For mathematical and programming work this
includes most of the ordinary Unix-style
toolchain.

So saying that Apple simply prevents me
from installing software is too crude a
description of the problem.

\medskip\noindent
{\bf Can I replace the operating system?}
This depends on the particular Mac and on
the state of Linux support for its
hardware.  Even when Linux can run, not all
hardware features necessarily work equally
well.

Thus I should never buy a particular Mac
on the assumption that it will later
become an ideal Linux laptop without first
checking that exact model.

\medskip\noindent
{\bf Can I upgrade the machine?}
This is a much clearer limitation.
Modern Mac laptops integrate many
components very tightly.  RAM and storage
are generally not ordinary user-upgradable
modules.

That means I must buy enough memory and
storage at the beginning, and hardware
obsolescence can force replacement of more
of the machine than I would like.

\medskip\noindent
{\bf Can Apple force me to use its AI?}
No general principle forces my computing
workflow to use Apple's AI services.  I
can use third-party software and local
models.

The more serious strategic worry is
ecosystem dependence: if my files,
workflows, devices, authentication, and
services become tightly tied to one
vendor, leaving becomes expensive even if
no single restriction prevents me from
leaving.

\medskip\noindent
{\bf Privacy.}
Privacy is not the same as openness.
Closed software may have strong privacy
properties, and open-source software may
still send data to remote services if it
is configured to do so.

The robust habit is therefore not to trust
a brand-level slogan.  I should understand
where my data live, what is synchronized,
which services receive it, and which parts
of my workflow can work completely
offline.

\section{Mac versus Linux for a laptop}
\label{computing-systems-section-computing-systems-mac-linux}

The choice is not good hardware versus bad
hardware.

A Mac laptop may give an unusually strong
combination of battery life, performance,
screen, trackpad, speakers, compactness,
and quiet operation.  Those are real
advantages for a person whose entire
computer has to travel in a backpack.

A Linux laptop can give much more control
over the operating system and, depending
on the machine, much better repairability
and upgradeability.

Linux also gives direct experience with
the environment used by a huge amount of
scientific computing, servers, and AI
infrastructure.

Therefore the question for my next laptop
is not which platform is morally better.
It is which compromises matter during the
period in which this laptop has to be my
entire computing system.

\section{What I should compare before buying}
\label{computing-systems-section-computing-systems-buying}

For every serious candidate, I should make
the same checklist.

\begin{enumerate}
\item CPU performance for sustained work.
\item Amount of RAM and maximum RAM.
\item Whether RAM is replaceable.
\item SSD capacity and replaceability.
\item GPU and memory available to the GPU.
\item Quality of Linux support.
\item Battery life under real workloads.
\item Weight including the charger.
\item Cooling and fan noise.
\item Screen, keyboard, and trackpad.
\item Availability of replacement parts.
\item Ease of battery replacement.
\item Expected operating-system support.
\item Ability to run useful local AI.
\item Total price with enough RAM and
storage for several years.
\end{enumerate}

The right comparison is between complete
machines configured for my actual needs,
not between the cheapest advertised
versions.

\section{What I should learn next}
\label{computing-systems-section-computing-systems-next}

A sensible learning order is:

\begin{enumerate}
\item CPU, RAM, SSD, GPU, and VRAM.
\item Processes, filesystems, and memory.
\item The shell and the Unix filesystem.
\item Package managers.
\item Permissions and users.
\item Processes and services.
\item SSH and remote computing.
\item Partitions and filesystems.
\item Booting, UEFI, and firmware.
\item Drivers and the Linux kernel.
\item Backups and disk encryption.
\item Containers and virtual machines.
\item Local AI runtimes and model formats.
\end{enumerate}

I do not need to master all of this before
buying a computer.  The point is to stop
treating the machine as a mysterious
sealed object and gradually understand
each layer well enough to control my own
workflow.

% --- End computing-systems.tex ---

\part{Humanities and Social Sciences}
\label{part-humanities-and-social-sciences}

% --- Begin economy.tex ---
\chapter{Economy}
\label{chapter-economy}





\section{Economics Nobel Prize 2026}
\label{economy-section-nobel-economics-2026}

\noindent
AI-generated summary, 12 October 2026.

\medskip\noindent
{\bf Big picture.}
Robert E. Hall and Michael Woodford
received the 2026 economics prize for
microeconomic foundations of macroeconomic
stabilisation. The question is how the
decisions of households and firms combine
into inflation, spending and employment,
and how policy can soften economic shocks.
See the Royal Swedish Academy's
\href{https://www.kva.se/en/news/%
the-prize-in-economic-sciences-2026/}%
{announcement of 12 October 2026}.

\medskip\noindent
{\bf Hall: consumption and expectations.}
Spending depends on expected future
resources, not only current income.
A temporary bonus and a permanent pay
rise need not produce the same response.
Hall studied the implications of combining
forward-looking consumption decisions with
rational expectations. These are results
within a model, not universal descriptions
of every household.

Reference: Robert E. Hall,
\emph{Stochastic Implications of the Life
Cycle-Permanent Income Hypothesis:
Theory and Evidence},
Journal of Political Economy 86(6)
(1978), pp.~971--987,
\href{https://doi.org/10.1086/260724}%
{published article}.

\medskip\noindent
{\bf Woodford: monetary policy.}
Woodford developed a framework connecting
household choices, firms' price setting
and interest-rate policy. When prices
adjust slowly, monetary policy can affect
production and employment. Expectations
about future policy also affect decisions
today: communication is part of policy.

Reference: Michael Woodford,
\emph{Interest and Prices: Foundations
of a Theory of Monetary Policy},
Princeton University Press (2003),
Chapter 1, ``The Return of Monetary
Rules'', pp.~1--58; see also
\href{https://www.jstor.org/stable/%
j.ctv30pnvmf}{the book on JSTOR}.

\medskip\noindent
{\bf Connection with these notes.}
The Selic discussion is an example of the
policy questions this framework addresses.
Woodford also appears in our SAET notes:
see Section~\ref{% 
section-fiscal-sustainability}.
There the focus shifts to fiscal policy
and more realistic models of how people
plan and learn.

\section{What energy actually means}
\label{economy-section-energy-means}

\noindent
This is an AI-generated section.

\medskip\noindent
This section follows the historical
framework developed by Vaclav Smil in
\emph{Energy and Civilization: A History}
\cite{smil2017energy}.

\medskip\noindent
{\bf Energy and useful work.}

Energy is not merely one commodity among
others. Every society depends on converting
energy into useful work: movement, heat,
light, construction, transportation, food
production, and, more recently,
computation. What changes historically is
not the existence of energy, but the ability
of human societies to capture, convert,
store, transport, and control energy flows.

\medskip\noindent
{\bf Foraging societies.}

For early humans, almost all available
energy ultimately came from photosynthesis
[but I do think that all the energy on Earth
ultimately comes from the Sun,
because, isn't oil coming from
very old fossilized organic matter?].
Plants stored solar energy, humans and
animals consumed plants or other animals,
and muscular effort converted food energy
into mechanical work.

Fire made it possible to obtain heat outside
the human body, while tools increased the
effectiveness of muscular labour. Still,
the amount of useful work available to a
society remained closely constrained by
human metabolism and by the biological
productivity of its environment.

\medskip\noindent
{\bf Agriculture and the control of solar
energy.}

Agriculture reorganized ecosystems in order
to concentrate solar energy into crops and
domesticated animals useful to humans. It
therefore represented much more than a new
method of obtaining food.

Agricultural surpluses supported larger
populations, permanent settlements,
specialized occupations, armies, and
states. Nevertheless, traditional
agricultural civilizations remained
dependent on current or recent solar flows:
food, wood, animal feed, wind, and flowing
water.

Their energy supply was consequently
limited by land, climate, seasons, and the
rate at which biological resources could
be renewed.

\medskip\noindent
{\bf Animals, water, and wind.}

Draft animals supplied mechanical power
beyond the human body. Waterwheels and
windmills converted natural flows into
useful work without requiring continuous
human or animal effort.

These sources considerably increased the
power available for farming, milling,
transport, and manufacturing. They were,
however, geographically and temporally
restricted. Rivers existed only in certain
places, winds were variable, and animals
required large quantities of land and feed.

\medskip\noindent
{\bf Fossil fuels and industrialization.}

Coal, and later oil and natural gas, allowed
societies to use energy accumulated over
geological timescales. Humanity was no
longer restricted to the energy captured
each year by contemporary biological
processes.

Fossil fuels were particularly
transformative because they were highly
concentrated, storable, transportable, and
capable of powering machines on demand.
Steam engines converted coal into
mechanical work, while internal-combustion
engines made petroleum central to modern
transportation.

The resulting increase in available power
transformed industry, agriculture,
transport, warfare, cities, and ordinary
domestic life.

\medskip\noindent
{\bf Electricity as an energy carrier.}

Electricity is not a primary source of
energy. It must be produced from some other
source, such as coal, natural gas, moving
water, nuclear reactions, wind, or
sunlight.

Its importance lies in its flexibility.
Electricity can be transported through
networks and converted into light, heat,
motion, refrigeration, communication, or
computation. Modern societies therefore
depend not only on producing energy, but
also on maintaining large and precisely
coordinated electrical systems.

\medskip\noindent
{\bf The historical lesson.}

Smil's central perspective is that the
history of civilization can partly be read
as the history of increasingly powerful
energy conversions:
\[
  \text{capture}
  \longrightarrow
  \text{conversion}
  \longrightarrow
  \text{control}
  \longrightarrow
  \text{useful work}.
\]

Civilizations are not shaped merely by the
total amount of energy they possess. They
are shaped by the concentration, quality,
reliability, location, and convertibility
of their energy sources.

This provides the material background for
the economic questions studied later.
Markets, governments, and institutions
organize who controls energy resources,
how they are used, who receives their
benefits, and who bears the environmental
and social costs.


\section{Prices versus quantities}
\label{economy-section-prices-quantities}

\noindent
This section is AI-generated.

\noindent
This section follows Martin Weitzman's
paper \emph{Prices vs. Quantities}
\cite{weitzman1974prices}.

\medskip\noindent
{\bf The policy problem.}

Suppose that a government, regulator, or
other central authority has already decided
that some economic activity must be
regulated.

Weitzman does not study why regulation is
necessary. He takes this decision as given
and asks how the regulation should be
implemented.

The regulator has two basic options.

A quantity instrument directly specifies
the permitted level of the activity. It can
take the form of a quota, target, production
requirement, or emission standard.

A price instrument instead changes the
price faced by the agents. The regulator
announces a tax, subsidy, or transfer price
and lets producers choose their own
quantity in response.

For pollution, the comparison is between
an emission standard and a pollution tax.

\medskip\noindent
{\bf A desirable quantity.}

Weitzman formulates the problem using a
desirable commodity. Let
\[
  q
\]
denote the quantity of this commodity.

In the pollution example, it is more
convenient to let $q$ denote clean air,
environmental quality, or pollution
abatement rather than pollution itself.

Producing $q$ yields a social benefit
\[
  B(q)
\]
and incurs a cost
\[
  C(q).
\]

The regulator wants to maximize total
social surplus:
\[
  W(q)=B(q)-C(q).
\]

Under complete information, the optimal
quantity $q^*$ therefore satisfies
\[
  B'(q^*)=C'(q^*).
\]

The marginal social benefit of one more
unit must equal its marginal social cost.

\medskip\noindent
{\bf Equivalence under certainty.}

Define the optimal price
\[
  p^*=B'(q^*)=C'(q^*).
\]

The regulator could directly order the
quantity $q^*$.

Alternatively, it could announce the price
$p^*$ and let the producer maximize
\[
  p^*q-C(q).
\]

The producer would then choose $q^*$.

Thus, under complete knowledge and perfect
certainty, price and quantity instruments
are formally equivalent.

The difference between them appears only
when the regulator does not know the cost
and benefit functions perfectly.

\medskip\noindent
{\bf Uncertainty.}

Weitzman writes the cost and benefit
functions as
\[
  C(q,\theta),
  \qquad
  B(q,\eta),
\]
where $\theta$ and $\eta$ represent
unknown circumstances.

The uncertainty need not be genuine
randomness. It may simply represent an
information gap between the regulator and
the people who actually produce the good.

For example, a regulator may not know the
true cost of reducing pollution because it
does not know exactly which technologies
firms can use.

\medskip\noindent
{\bf Quantity regulation.}

Under a quantity instrument, the regulator
chooses a fixed target $\widehat q$ before
the uncertainty is resolved.

It chooses $\widehat q$ to maximize
expected social surplus:
\[
  \widehat q
  =
  \operatorname*{arg\,max}_{q}
  \mathbb{E}
  \bigl[
    B(q,\eta)-C(q,\theta)
  \bigr].
\]

The advantage is control: the quantity is
fixed.

The disadvantage is rigidity: after the
true costs become known, producers must
still produce the prescribed quantity.

\medskip\noindent
{\bf Price regulation.}

Under a price instrument, the regulator
announces a fixed price $p$.

After observing its own production costs,
the producer chooses
\[
  q=h(p,\theta)
\]
to maximize
\[
  pq-C(q,\theta).
\]

Thus
\[
  C_q(h(p,\theta),\theta)=p.
\]

The regulator chooses the price that
maximizes expected social surplus:
\[
  \widehat p
  =
  \operatorname*{arg\,max}_{p}
  \mathbb{E}
  \left[
    B(h(p,\theta),\eta)
    -
    C(h(p,\theta),\theta)
  \right].
\]

The advantage is flexibility: production
responds to the costs actually encountered
by the producer.

The disadvantage is loss of quantity
control: the resulting quantity is not
known when the price is announced.

\medskip\noindent
{\bf Comparing the policies.}

Weitzman defines the comparative advantage
of prices over quantities as
\[
  \Delta
  =
  \mathbb{E}
  \left[
    W_{\mathrm{price}}
    -
    W_{\mathrm{quantity}}
  \right].
\]

Thus,
\[
  \Delta>0
\]
means that the price instrument has higher
expected social welfare, while
\[
  \Delta<0
\]
means that the quantity instrument is
better.

Using local quadratic approximations to
the cost and benefit functions, Weitzman
obtains
\[
  \Delta
  \simeq
  \frac{\sigma^2}{2(C'')^2}
  \bigl(B''+C''\bigr),
\]
where
\[
  B''<0,
  \qquad
  C''>0,
\]
and $\sigma^2$ measures uncertainty in
marginal costs.

The sign of $\Delta$ is therefore determined
by the comparison between the curvature of
benefits and the curvature of costs.

\medskip\noindent
{\bf Weitzman's decision criterion.}

If
\[
  -B''>C'',
\]
then the benefit function is more sharply
curved than the cost function. Small errors
in quantity cause large losses in social
benefit.

In this case,
\[
  \Delta<0,
\]
and quantity regulation is preferred.

If
\[
  C''>-B'',
\]
then marginal costs rise more sharply than
marginal benefits change.

In this case,
\[
  \Delta>0,
\]
and price regulation is preferred.

The criterion can be summarized as follows:
\[
  \boxed{
  \begin{array}{c}
  \text{steep marginal social benefits}
  \\[2mm]
  \Longrightarrow
  \text{quantity control},
  \\[4mm]
  \text{relatively flat marginal benefits}
  \\[2mm]
  \Longrightarrow
  \text{price control}.
  \end{array}}
\]

\medskip\noindent
{\bf Interpretation for pollution.}

Let $q$ represent pollution abatement or
environmental quality.

If being slightly below the desired level
of environmental quality would cause very
large damages, the regulator should value
precise control of the quantity. An
emission standard or cap is then relatively
attractive.

If the environmental benefit of additional
abatement changes slowly, while the cost
of abatement is uncertain, a pollution tax
is relatively attractive. Firms can adjust
their abatement according to the costs
they actually face.

Thus, the question is not whether prices
are always better than direct regulation,
or whether direct regulation is always
safer.

The regulator should compare the expected
loss caused by choosing the wrong quantity
with the expected loss caused by choosing
the wrong price.

\medskip\noindent
{\bf Scope of the result.}

Weitzman's criterion is not a complete
algorithm for government policy.

It is derived from a simplified model and
a local quadratic approximation. Legal,
political, administrative, distributive,
monitoring, and enforcement considerations
remain outside the formal comparison.

The theory provides an economic benchmark:
it identifies when uncertainty and the
local shapes of costs and benefits favour
one instrument over the other.

\section{Scheinkman}
\label{economy-section-scheinkman}

\noindent
The upshot of Scheinkman's lecture
at SAET is:
Can economic incentives be designed so that
preserving nature wins against destroying
it?

And apparently we expect his answer to be
affirmative: Properly priced carbon capture
can make forest recovery more profitable
than low-productivity cattle ranching.



\section{Growth and Redistribution with
Risk Lovers and Risk Averters}
\label{economy-section-risk-lovers}

\medskip\noindent
Timothy Kehoe,
University of Minnesota.

\medskip\noindent
This talk was given at the 25th Annual
SAET Conference, held at IMPA in Rio de
Janeiro, Brazil, from July 13 to
July 18, 2026.

\medskip\noindent
{\bf Short communication.}

\medskip\noindent
{\bf Model.}
The economy consists of two types of
agents: risk lovers and risk
averters. Risk lovers invest in a
risky, high-return technology,
whereas risk averters prefer a safer
asset with lower returns.

\medskip\noindent
{\bf Main question.}
The model studies the trade-off
between economic growth and
redistribution. The risky investment
generates growth, but also increases
inequality.

\medskip\noindent
{\bf Welfare function.}
The objective is to maximize a CES
social welfare function. Different
choices of the CES parameter assign
different weights to inequality.

\medskip\noindent
{\bf Optimal policy.}
The model computes the optimal tax
rate. This is analogous to finding an
optimal policy parameter (such as the
optimal emission price in the forest
paper): the parameter maximizing the
chosen welfare function.

\medskip\noindent
{\bf Interpretation.}
Since the risk lovers become richer,
the policy effectively taxes the rich
and redistributes resources to the
risk averters. The interesting
question is how much redistribution
can be introduced before growth is
reduced too much.

\medskip\noindent
{\bf Remarks.}
The classification "risk lover" is a
technical one: it refers to a
preference for risky gambles, not to
a political or moral category.
Likewise, "optimal" always depends on
the chosen welfare function.


\section{Real Effects of Interest-Rate
Policy and Central-Bank Asset Purchases
with Endogenous Collateral}
\label{economy-section-interest-rate-policy}

\medskip\noindent
Susan Schommer,
University of Bonn.

\medskip\noindent
This talk was given at the 25th Annual
SAET Conference, held at IMPA in Rio de
Janeiro, Brazil, from July 13 to
July 18, 2026.

\medskip\noindent
{\bf Upshot (human-written).}


\medskip\noindent
{\bf Motivation.}
In the basic flexible-price economy,
$p_1$ and $\sigma^h$ are
predetermined. This does not capture
the effects of predictable monetary
policy.

\medskip\noindent
{\bf Household problem.}
Given $(p,q,y)$, household $h$
chooses a vector in order to maximize
its utility.

\medskip\noindent
{\bf Equilibrium allocation.}
The equilibrium allocation is the
equilibrium vector obtained by
collecting the optimal choices of all
households (and firms). Its
coordinates correspond to quantities
such as consumption, investment,
asset holdings, etc.

\medskip\noindent
{\bf Theorem (equilibrium mapping).}
Every equilibrium of the basic
price-friction economy
$\mathcal E$ maps to an equilibrium
of the auxiliary flexible-price
economy $\widehat{\mathcal E}$.

\medskip\noindent
{\bf Proposition (real effects of
interest policy).}
Consider the basic price-friction
economy $\mathcal E$ with logarithmic
utility and fixed $x^{CB}_3$.
Changes in the policy interest rate
$i\ge0$ have real effects.
Specifically, variations in $i$
affect relative prices, real rates of
return, and the equilibrium
allocation.

\medskip\noindent
{\bf Example.}
One figure plots the policy interest
rate against several coordinates of
the equilibrium allocation. The
curves are not constant, showing that
different coordinates of the
equilibrium vector vary with $i$.

\medskip\noindent
{\bf Interpretation.}
The basic flexible-price model does
not detect real effects of monetary
policy. Introducing a single price
friction restores them: the model now
detects that changing the policy
interest rate changes the equilibrium
allocation.

\medskip\noindent
{\bf Extended model.}
The extended model incorporates
predictable monetary policy. The
conclusion is that only the
unforecastable component of monetary
policy has real effects, since
predictable changes are already taken
into account by the agents.

\medskip\noindent
{\bf Conventional versus
unconventional monetary policy.}
Conventional monetary policy changes
the short-term policy interest rate.
Unconventional monetary policy acts
through changes in the central-bank
balance sheet, for instance by
purchasing assets.

\medskip\noindent
{\bf Conclusions.}
The optimal interest-rate policy is
state-contingent: the optimal policy
depends on the state of the economy,
and interest rates should be raised
in high-production states.

For unconventional monetary policy,
the effects depend on the financial
constraint. Relaxing leverage
constraints improves welfare, whereas
relaxing short-sale constraints does
not.

\section{Fiscal Sustainability and
Macroeconomic Stabilization}
\label{economy-section-fiscal-sustainability}

\medskip\noindent
Michael Woodford,
Columbia University.

\medskip\noindent
This talk was given at the 25th Annual
SAET Conference, held at IMPA in Rio de
Janeiro, Brazil, from July 13 to
July 18, 2026.

\medskip\noindent
{\bf Introduction.} Policy analysis involves
working with ``infinite-horizon general
equilibrium models'' --- somehow ideal
states of equilibrium predicted by the model
which would occur only at infinite time.

[Missing content.]

Somehow this might cause that certain
{\it transfer policies} have no
consequences.

Asumption: sustainability can be treated as
axiomatic or not. The choice might determine
what monetary policy you will choose
to stabilize marcoeconomics.

Consider monetary policy rules of
Taylor-type form:
$$
i_t=\phi_\pi \pi_t + \phi_y y_t+
e^m_t,
$$
with $\phi_\pi,\phi_y \geq  0$,
$\varepsilon_t^m$ exogenous (may respond to
other ``fundamental'' disturbances),
and all variables written in terms of
deviations from a zero-inflation steady
state.

Then in a standard RANK model,
log-linearized equations can be writen in
vector form [involves a matrix].

\medskip\noindent
{\bf Recap of REE analysis.}
(Rational expectations ?.)

Idea: no feedback to nominal interest rates
from (?) inflation.

\medskip\noindent
{\bf Adding debt dynamics.}
Suppose that the only debt instrument is a
one-period nominal bond, and that all
taxes/transfers are lump-sum.

Then previous structural equations are
unchanged; but we must add an equation
describing the dynamics of public debt,
taking into account transfer policy:
[equation…]

\medskip\noindent
{\bf }
Hence REE implies
\begin{itemize}
\item AM/PF or AF/PM $\implies $ unique
non-explosive REE.
\item AM/AF $\implies$ only explosive
solutions.
\item PM/PF $\implies$ indereminacy, sunspot
equation.
\end{itemize}

\noindent
In this analysis, it does matter
whether there is a commitment to sustainable
fiscal policy.

[We're looking at how to find a proper
macroeconomics stabilization policy.]

\medskip\noindent
{\bf Questioning this analysis.}
\begin{itemize}
\item Conceptual problem: how can one settle
whether fiscal policy guarantees
sustainability or no, when this is entirely
a question about what fiscal policy promises
to do infinitely far in the future?

\item (Missing)
\end{itemize}

\medskip\noindent
{\bf Alternative approach.}

\begin{itemize}
\item Can we model the way in which people
make forward-looking decisions, without
assuming (unrealistically) that they have
definite expectations about the far
future…

\item  The idea: 
\begin{itemize}
\item Economic agents formulate plans some
finite distance into the future, on the
basis of a model of possible transitions
over that planning horizon (don't look at
all possible final states, but a reduced
set).

\item the stages that can be reached at the
horizon at which forward planning is
truncated are evaluated using a {\it learned
value function}.
\item agents don't continue to execute plans
made in the past, but re-optimize from
whatever situation they have reached.
\item value functions are update as
additional experience is accumulated,
through reinforcement learning.
(AAAAAA!!!!)
\end{itemize}
\end{itemize}

\medskip\noindent
{\bf Learning the value functions.}
\begin{itemize}
\item As part of its planning exercise in
period $t$, each houshold $h$ computes an
estimate of the value of its objective
(expected discounted utility from $t$
onward), for any (counterfactual) level of
wealth $b_t^h = b$ that is might have
brought into the period --- call this
function $v^{est}_t(b)$.

\item Each firm $f$ similarly…
\item Value function can then be updated for
the next period's exercise using error
correction rule (difference equation).
\end{itemize}


\medskip\noindent
{\bf Rational expectations?}
\begin{itemize}
\item What relationship is ther ebetween the
dynamics implied by a computational model of
this kind and REE predictions?

\item If the value functions used by
households and firms in period $t$ are the
ones implied by a particular REE they will
all take the actions associate dwith the
REE.

\item Quesion: whether learning dynamics
converge? …?
\end{itemize}

\medskip\noindent
{\bf Log-Linear leaning dynamics.}
\begin{itemize}
\item The solution to this linear system
then determins
$$
b_{t+1}=b_{t+1|_t},\qquad 
v_t^{\operatorname{est}}=x_{t|t}-
\zeta b_t
$$
and …
if so the value functions converge…

\medskip\noindent
{\bf When do learning dynamics converge?}
\begin{itemize}
\item Conclusion when
$\bar{b}=0$:convergence if and only if AM/PF
policy configuration --- either passive MP
or active FP produces an unstable
eigenvalue.

\item Continues to be true for (modestly)
values of…
\end{itemize}

\medskip\noindent
{\bf Numerical example: $\bar{b}=1.0$ GDP, 5
year horizon.} (Figure: AM/PF remains
sufficient, and almost necessary for
stability.)

\medskip\noindent
{\bf Numerical example: finer detail.}
(Similar figure.)

\medskip\noindent
{\bf Conclusions.}
\begin{itemize}
\item One shouldn't expect transfer policy
to be irrelevant to aggregate demand
determination (and hence for price
stability)

\item According to our model with
finite-horizon planning…

\item Broader lesson: when analysing
possible stabilization policies we
shouldn't be content merely to consider
the REE consistent with a given policy
commitment we should also ask whether ar
least something close to that outcome can be
achieved with a more realistic model of
peoples thought processes --- not always the
case!
\end{itemize}


\end{itemize}

\medskip\noindent
{\bf Reference.}
Ayres et al.\ (2019) study the
monetary and fiscal history of Brazil
from 1960 to 2016, especially the
hyperinflation episodes.




\section{Inflation Targeting under
Fiscal Fragility}
\label{economy-section-inflation-targeting}

\medskip\noindent
Vitor Chagas da Costa,
IMPA.

\medskip\noindent
This talk was given at the 25th Annual
SAET Conference, held at IMPA in Rio de
Janeiro, Brazil, from July 13 to
July 18, 2026.


\medskip\noindent
{\bf Question.} How to set the inflation
targe in an economy with high government
debt burden?

\begin{itemize}
\item Some authority has an inflation target
and attempts to implement policies to
achieve it. But in practice, the target is
not attained.

\item High debt burden raises the temptation
of the policymaker to use inflation as a
partial default mechanism. However, high
inflation reduces welfare and affects
monetary policy credibility.

\item (missing.)

\item Upshot of the paper:
``The target is important in balancing
this tradeoff''
\end{itemize}

\medskip\noindent
{\bf Brazilian experience with inflation
targeting regime.} (graph: effective
inflation and inflation target 
(the latter with a neighbourhood: lower
and upper limit). The actual data is much
more random and often outside the expected
neighbourhood!)

\medskip\noindent
{\bf Model.} We model an economy with two
types of agents:
\begin{itemize}
\item An altruistic policymaker acts as a
fiscal authority and central bank, under an
inflation target regime.

\item Private agents form rational
expectations and hold debt.
\item The PM chooses the combination of
monetary and fiscal policy that maximizes
welfare, given an initial debt level
\end{itemize}

\noindent
Deviation from target solution (…)

\medskip\noindent
{\bf Results.} What the policy would do is
dependent on that limit:

\begin{itemize}
\item we obtain endogenous limits that determine
target coordination ability.

\begin{itemize}
\item fiscal dominance zone: high debt,
target irrelevant
\item no crisis zone: low debt, target
delivered.
\item in between: fiscal fragility
\end{itemize}

\item In the fiscal fragility zone (like in
Brazil's economy!):
\begin{itemize}
\item expected inflation is higher than
target.
\item delivering the target is more costly.
We need high interest rate and physical
effort to attain the target.
\end{itemize}

\item Raising inflation target:
\begin{itemize}
\item reduces the real interest rate.
\item increases lower bound of FFZ.
\item reduces the time to exit the FFZ.
\end{itemize}

\item The optimal target is the lowest
possible (inflation level) such that 
it is credible.
\end{itemize}

\medskip\noindent
{\bf References}Basic to inflation: Ser and
Wallace…

\medskip\noindent
{\bf Introducing the model: policymaker}
\begin{itemize}
\item Chooses inflation $\pi_t \geq 0$
and debt $D_{t+1}$
$$
\operatorname{max}_{\pi_t,D_{t+1}}
\mathbb{E}_0 \sum_{t=0}^\infty 
\beta^t[u(c_t,g_t)-\psi(\pi_t)]
$$
subject to $g_t+(1+r_t) D_t \leq 
D_{t+1}+\tau(e-\varepsilon_t)$.

We assume:
\begin{itemize}
\item $c_t,g_t>0$
\item $u(c,g)=\rho_c+(1-\rho)v(g)$, utility
in private consuption is risk-neutral,
and $v$ is strictly concave.
\item $\psi$ is the disutility of inflation,
such that $\psi'>0$ and $\psi'' \geq 0$.
\item Endowment $e$ taxed at fixed rate $\tau$
\item $\varepsilon_t$ is the endowment
penalty for deviating from target.
\end{itemize}
\end{itemize}

\medskip\noindent
{\bf Private agents.} (Missing, formulas and
parameters.)


\medskip\noindent
{\bf Policy maker's choices.}
The first-order condition for debt from the
private agents problem implies the following
equilibrium of real interest rate:
(formula).

\medskip\noindent
{\bf Timing.} Interest rate is defined in
equilibrium.

\begin{enumerate}
\item Sunspot variable $\zeta_t$ is
realized.
\item Policymakes chooses to follow the
target or not, given realization of
$\zeta_t$.
\item …
\item …
\end{enumerate}

\begin{itemize}
\item If expected inflation is high,
by the equation, the policymaker may decide
to inflate the debt away, so, it's costly.
\end{itemize}

\medskip\noindent
{\bf Fisical Fragility.}
As a result, we find, as we already
anticipated, there are three possible zones
(which depend on belief of the agents):

\begin{itemize}
\item (No crisis zone.) 
Agents believe the goverment will
deviatie: it's optimal for the goverment to
follow the target. This is the only solution
of the model.

\item (Fiscal dominance zone.) Even if agents believe the goverment,
the government might deviate.

\item (Fiscal fragility zone.) If agents believe believe the
goverment wants to deviate (not keep
inflation on target),
\end{itemize}

\noindent
In FEZ, estimated inflation is given by
[formula].

\medskip\noindent
{\bf Main result of paper.}
Higher target improves coordination by
reducing inflation volatility.
Therefore, the rollover cost is reduced,
lowering the temptation of deviating.

\medskip\noindent
{\bf Next result.}
The boundary between no crisis and fiscal
fragility…

\medskip\noindent
{\bf Numerical example.}
In 2002 Brazilian confidence crisis.
Inflation expectations grew by
fear that the president would change
policies.

\noindent
{\bf Optimal fiscal policy.}
Figure: in fiscal fragility, the policymaker
reduces (?) to get back to no crisis zone.

We can also plot the same functions varying
the target.

\medskip\noindent
{\bf Concluding remarks.}
\begin{itemize}
\item A higher inflation target may ssist
policymaker in economies with high debt
burden that are subject to sudden loss
of confidence on the commitment with the
target.

\item Decreasing debt levels through
austerity serves as a fiscal support to
better target coordination when the economy
is in the lower part of the FFZ.

\item Model can be relevant to the current
experience with high inflation and high debt
in emergent and advanced economies (e.g.
USA, European economies).
\end{itemize}

\medskip\noindent
{\bf Audience questions:} 

There is only
one target in the model, but the whole thing
is changing the target. How to decide?


Did you consider
price-inflation targeting, average-inflation
targeting?

\section{Recursive Contracts in
Non-Convex Environments}
\label{economy-section-recursive-contracts}

\medskip\noindent
Felix Kübler,
University of Zurich.

\medskip\noindent
This talk was given at the 25th Annual
SAET Conference, held at IMPA in Rio de
Janeiro, Brazil, from July 13 to
July 18, 2026.

\medskip\noindent
{\bf Non-convexities in Aloisio's work.}
\begin{itemize}
\item Asymmetric information, 2001.
\item Non-convex budget sets, 2002.
\item Risk-loving preferences, 2018.
\end{itemize}


\medskip\noindent
{\bf Lotteries.}
\begin{itemize}
\item Many economic models involve
non-convex optimization problems.
\item Non-convexities.
\begin{enumerate}
\item Difficult to solve optimization
problems with nonconvex constraints or
objective function that is not concave.
\item Lotteries can increase the value of
the objective function.
\end{enumerate}
\item Since allowing for randomization
across allocations might increase the value
of the objective function, there has been
substantial interest in exploring lottery
solutions to models with nonconvexities
(Weiss 1976, Myerson 1982, Prescott and
Townsend 1984, Arnott-Stiglitz 1988.)

\item Questions:
\begin{enumerate}
\item Does it make sense?
\item How can we compute optimal loteries?
\end{enumerate}
\end{itemize}

\medskip\noindent
{\bf Simple example} Maximize $f$ with
constraints $f \leq  A$, $f \leq B$.

\medskip\noindent
{\bf Ramsey problem (economic example).} 
Consider the following non-convex problem:
(maximization, constraints, $T$ representing
taxes…)
The housholds know the probability
distribution…

Do lotteries ``make sense''?
Yes - if lottery is ex post (see Citanna and
Tvede 2024 and our paper for examples).


\medskip\noindent
{\bf Dynamic problem with forward looking
constraint.}

\medskip\noindent
{\bf Problems.} MM 2019 assumes concave
objective + convex constraint set. Many
important applications violate this: e.g.
Ramsey problem. Prior results (Cole-Kubler
2012)…

\medskip\noindent
{\bf Several slides missing.}

\medskip\noindent
{\bf First period.} Involving
sub-gradient descent.

\medskip\noindent
{\bf Algorithm for computing optimal
lotteries (Algorithm 1).}

\noindent
{\bf Given:} value function $D(\cdot)$,
state $(x,s)$, initial promise
$\phi$, initial multipliers
$\lambda^1, \mu^1$.

\noindent
{\bf Subgradient Descent on Multipliers.} 
For $k=1,…,N$:
\begin{enumerate}
\item Solve $a^k \in \operatorname{arg}
\operatorname{max}_a[$ long formula].
\item Pick $\phi^k(s') \in \partial
\partial D(\lambda^k+\mu^k,x',s')$.
\item Update.
\end{enumerate}

\noindent
{\bf Output.} Stage lottery.

\medskip\noindent
{\bf Slightly more interesting example.}

\medskip\noindent
{\bf Conclusion.}
\begin{itemize}
\item Equivalence between the dual and the
lottery relaxation.
\item Different recusive formulas
reflect different versions of lottery
relaxation.
\item Subgradient descend method
to recover the policy.
\item Application to Ramsey problem.
\end{itemize}

\section{From Degradation to Recovery:
Quantifying the Impact of Emission
Prices on Tropical Forests}
\label{economy-section-tropical-forests-scheinkman}

\medskip\noindent
José Scheinkman,
Columbia University.

\medskip\noindent
This talk was given at the 25th Annual
SAET Conference, held at IMPA in Rio de
Janeiro, Brazil, from July 13 to
July 18, 2026.

\medskip\noindent
{\bf Motivation.}
Want to quantify impact of emission prices
on tropical rainforests.

\begin{itemize}
\item The 
Brazilian Amazon provided the best data set,
60% of the whole Amazon. Also, has been
substatianlly deforested.

\item Many scientists predicted parts of the
Amazon were closed to being ?.

\item There's a lot of work already on the
Amazon. Bolsonaro = lots of deforestation.
This year = least deforestation.

\item Deforestation is a dynamical process.
\end{itemize}

\medskip\noindent
{\bf Today I.}

If stays the same: emit $\approx 14.5$ GT
of $CO_2$ in 30 years.

The second part: effect of this in
biodiversity, is work in progress.

\medskip\noindent
{\bf State and control variables.}
\begin{itemize}
\item $Z$…
\item $P^a$ will be a constant for today.
\end{itemize}

\begin{remark}
\label{economy-remark-carbon-cattle}
There's much better feed for carbon
than for cattle.
\end{remark}

\medskip\noindent
{\bf Robust optimization II.}
There are ``penaliazations'' (this reminds
me of RL): penalize departure.

\medskip\noindent
{\bf Planner's shadow price I
(or the previous slide).}
Use Hamiltonian Markov chain.

We check if we can predict the
information from the past!

The only tool you have is how much you value
deforestation, so you adjust the model.

\medskip\noindent
{\bf Before "biodiversity I"}
Somehow handle income/interest, money
basically, so that deforestation is no
longer convenient (maybe the whole point,
right?).

\medskip\noindent
{\bf biodiversity I.}
Arrhenius 1921 method works great
for one or two hectares.

\medskip\noindent
{\bf biodiversity II.}
Model: how many species that have n many
individuals in x hectares? Fischer's $\alpha$
determines the average number of species.
People change the $\alpha$ as a function
of the area.

\medskip\noindent
{\bf Clustering.}
Species are not evenly distributed,
so we need clustering. Clustering makes
i
\medskip\noindent
{\bf Forest cover changes…}

\medskip\noindent
{\bf There's not a conclusion slide but…}
Lesson: there's a lot of potential
for reforesting the Amazon.
People are not being paid for recuperating
the Amazon. All we do is penalize CO2
emissions. You get your biodiversity pack…
They have proposed, negotiated, to get this
done. The basis is this paper.

In Brazil we have data at municipal level
every 5 years.

Nobody knows how to value biodiversity.
``I'm gonna restore biodiversity:
give me $x$ dollars for every… what
restored?''
Biodiversity is a very complex system
--- a very complicated dynamical system.
This happens in any rainforest.
We don't understand the dynamics of
tropical rainforests.

Who does the reforestation: the government
could compensate organizations, but
most likely not doing it itself.
The average hectare has 5000 tons of CO2, 
4000 dollars. That is: here's 4000 dollars,
reforest the land.

30 years is enough to restore the Amazon ---
in 30 years you are in the 90s.
That number $\alpha$ is/takes the convergence
rate,

Aloiso: it's a very big economic
problem: who's gonna pay for this?


S: someone says:
 I'm gonna pay for this --- then someone
else will deforest… what's the use?
There are negotiations, like the ETS,
Brazil-China to make this happen.
China communist party.

% --- End economy.tex ---

% --- Begin linguistics.tex ---
\chapter{Linguistics}
\label{chapter-linguistics}





\section{Ancient Greek orthography}
\label{linguistics-section-ancient-greek-orthography}

\subsection{Breathings}

In polytonic Ancient Greek, a word beginning with a vowel normally carries a
\emph{breathing}.  The \emph{smooth breathing} (Greek \emph{psili}) marks the
absence of an initial consonantal sound, whereas the \emph{rough breathing}
(Greek \emph{daseia}) marks an initial /h/.  Thus a form transliterated
\emph{e-} with smooth breathing begins directly with the vowel, while a form
transliterated \emph{he-} with rough breathing begins with a light expiration
/h/.  The rough breathing is therefore an \emph{expiration}, not an
inspiration.

For example, the Homeric verb \emph{helisso}, ``to turn, twist, or curve'',
begins with rough breathing and is transliterated with an initial \emph{h}.
This is the verb behind the Homeric epithet \emph{amphielissai}, used of ships
and traditionally understood as ``curved on both sides'' or ``double-curved''.

A breathing is not an accent.  On an initial vowel, a breathing and an accent
may occur simultaneously.

\subsection{Accents}

The three traditional accents are the acute, grave, and circumflex.  In the
classical pitch-accent system they concern pitch rather than nasalization or
mere stress.

The acute marks a high pitch.  The circumflex can occur only on a long
syllable and represents a high--low contour within that syllable.  The grave
replaces a final acute in certain contexts when another word follows.  The
circumflex is sometimes printed in a shape resembling a tilde or a caret, but
it has nothing to do with the Portuguese tilde: it does not nasalize the
vowel.  A long vowel bearing a circumflex is long independently of the
circumflex; the accent indicates its pitch contour.

\section{Ancient Greek verbs}
\label{linguistics-section-ancient-greek-verbs}

\subsection{Regular present active endings}

A useful first model for the present indicative active is
\begin{center}
\begin{tabular}{c|cc}
 & Singular & Plural \\
\hline
1st person & \texttt{-o} & \texttt{-omen} \\
2nd person & \texttt{-eis} & \texttt{-ete} \\
3rd person & \texttt{-ei} & \texttt{-ousi(n)}
\end{tabular}
\end{center}
Thus, for example, \emph{luo} ``I release'' gives \emph{luo, lueis, luei,
luomen, luete, luousi(n)}.

The final \emph{-n} in forms such as \emph{-ousi(n)} is the so-called movable
nu: it may appear or disappear according to phonetic and metrical context.

\subsection{The verb \emph{eimi}: to be}

Ancient Greek does not make the Portuguese lexical distinction between
\emph{ser} and \emph{estar} in this verb.  The present indicative of
\emph{eimi}, ``to be'', is irregular:
\begin{center}
\begin{tabular}{c|cc}
 & Singular & Plural \\
\hline
1st person & \emph{eimi} & \emph{esmen} \\
2nd person & \emph{ei} & \emph{este} \\
3rd person & \emph{esti(n)} & \emph{eisi(n)}
\end{tabular}
\end{center}
Hence Homeric \emph{eisi(n)} means ``they are''.  Its ending should not be
confused with the regular third-person plural active ending \emph{-ousi(n)}.

\subsection{A warning about apparent endings}

A sequence at the end of a word need not be a verbal ending.  For example,
Homeric \emph{alphestai}, found in the expression \emph{aneres alphestai}, is
a noun rather than a verb; it refers to men, especially men who earn or seek
their livelihood.  Its final \emph{-stai} should therefore not be parsed as a
verbal ending merely because it resembles one.

By contrast, verbal \emph{-tai} is a genuine third-person singular
middle/passive ending in several paradigms; an immediately preceding
\emph{s} may belong to the stem or tense formation rather than to the ending
itself.

\section{Ancient Greek numbers}
\label{linguistics-section-ancient-greek-numbers}

Some basic cardinal numbers are:
\begin{center}
\begin{tabular}{c|l}
1 & \emph{heis, mia, hen} \\
2 & \emph{duo} \\
3 & \emph{treis, tria} \\
4 & \emph{tettares/tessares} \\
5 & \emph{pente} \\
6 & \emph{hex} \\
7 & \emph{hepta} \\
8 & \emph{okto} \\
9 & \emph{ennea} \\
10 & \emph{deka} \\
100 & \emph{hekaton}
\end{tabular}
\end{center}

The word \emph{hecatomb} ultimately contains \emph{hekaton}, ``one hundred'',
and \emph{bous}, ``ox/cow'': historically a ``hundred-ox'' sacrifice, though
in Homer the term need not always imply a literal count of exactly one
hundred animals.

\section{Indo-European comparison}
\label{linguistics-section-indo-european-comparison}

The similarity of number words across Indo-European languages is often hidden
by regular sound change.  Greek \emph{hekaton} ``100'', Sanskrit \emph{satam},
Latin \emph{centum} (classically pronounced with initial /k/), and English
\emph{hundred} belong to the same inherited word-family, conventionally
reconstructed from Proto-Indo-European as approximately
\texttt{*kmtom}, with the initial consonant and syllabic nasal represented more
precisely in specialist notation.

The resemblance between Sanskrit \emph{eka}, ``one'', and part of Greek
\emph{hekaton}, ``one hundred'', is therefore accidental: Sanskrit
\emph{eka} is not the cognate of Greek \emph{hekaton}.  The Sanskrit cognate
of Greek \emph{hekaton} is \emph{satam}.  A more transparent comparison is
Greek \emph{deka}, ``ten'', with Sanskrit \emph{dasa}, ``ten''.

Polish belongs to the West Slavic branch of Indo-European:
Proto-Indo-European $\to$ Balto-Slavic $\to$ Slavic $\to$ West Slavic
$\to$ Polish.  Polish \emph{sto}, ``100'', belongs to the same inherited
family as the forms above.  English instead belongs to the Germanic branch:
Proto-Indo-European $\to$ Proto-Germanic $\to$ West Germanic $\to$
Anglo-Frisian $\to$ English.


\section{Etymological notes}
\label{linguistics-section-etymological-notes}

\subsection{\emph{Aiolos} and ``eolic''}

In the \emph{Odyssey}, Aiolos (Greek \emph{Aiolos}, Latin \emph{Aeolus}) is
the keeper of the winds: Zeus has given him power to restrain or release
them. This is his mythological office, however, not the literal etymological
meaning of his name. The name is related to the Greek adjective
\emph{aiolos}, whose range of meanings includes ideas such as quick-moving,
changeful, or shifting; its deeper etymology is not entirely certain.

The modern Romance adjective ``eolic'' (Portuguese \emph{eólico}, Spanish
\emph{eólico}, Italian \emph{eolico}, etc.) develops in the other direction:
Aiolos/Aeolus, because of his mythological association with the winds,
$\to$ ``pertaining to Aeolus'' $\to$ ``pertaining to the wind''. Thus
``eolic energy'' ultimately bears the name of Aiolos; Aiolos was not named
from a Greek word meaning ``wind''.

The Homeric association is especially vivid in \emph{Odyssey} 10.1--79. Odysseus and his companions reach the floating island of Aiolos, who entertains them for a month. When Odysseus asks to continue home, Aiolos gives him an ox-hide bag in which he has bound the courses of the troublesome winds, while he leaves the West Wind free to carry the ships toward Ithaca (10.19--27). After nine days they are already close enough to Ithaca to see the fires on shore, but Odysseus falls asleep. His companions, imagining that the bag contains gold and silver secretly given to their leader, open it; the imprisoned winds burst out and drive them all the way back to Aiolos (10.28--55). Aiolos then refuses to help Odysseus a second time, taking the disastrous return as evidence that the gods hate him (10.56--79). Thus the episode gives a concrete Homeric background for the later semantic development Aiolos/Aeolus $\to$ ``of Aeolus'' $\to$ ``of the wind''. 

The initial alpha in Greek \emph{Aiolos} carries a smooth breathing, not a
rough breathing. The smooth breathing has no consonantal phonetic
realization: there is no initial /h/. This contrasts with a rough breathing,
which does mark initial /h/. Thus the breathing sign is orthographically
present, but in this case it tells us precisely that the word begins directly
with the vowel.

\subsection{Homeric ``teeth-hedge'' and Icelandic \emph{tann-garðr}}

The following note is worth preserving verbatim, both for its Homeric
linguistic observation and for the comparison with Icelandic. It is Note 3
to p.~3, Book I.64, in S.~H. Butcher and A.~Lang, \emph{The Odyssey of Homer:
Done into English Prose} (Macmillan and Co., London, 1912; first printed
1879). The note itself credits E.~Magnússon, translator of \emph{Legends of Iceland}.

\begin{quote}
\emph{ἕρκος ὀδόντων}, or teeth-hedge.

This phrase, which regards the teeth as the hedge or fence which protects the
mouth, is a precise parallel to the Icelandic \emph{tann-garðr}, i.e.
teeth-garth (Old English \emph{geard}, enclosure). The difference, however,
between the two expressions is that the Icelandic is the common prose phrase
in use up to this day, while the Greek is archaic and poetical. The Icelandic
phrase is of too old a date in the language to be directly borrowed from the
Greek; it is the genuine metaphor of a military age, in which the teeth were
looked upon as the wall guarding the castle, that is, the mouth. A long list
might be drawn up of metaphorical expressions common to Homer and the
Icelandic sagas, but independent in origin, and pointing to similar customs
and conditions of life.

For this note we are indebted to Mr. E. Magnússon, translator of 'Legends of
Iceland,' etc.
\end{quote}

The Greek expression occurs at \emph{Odyssey} 1.64. The same note is also
cross-referenced by the editors at \emph{Odyssey} 10.328, in Circe's speech
after Odysseus drinks her \emph{pharmakon} without being enchanted: she says
that no other man has resisted the charm once it has ``passed his lips''.

\section{A first map of the Northern languages}
\label{linguistics-section-northern-languages}

Here ``Northern languages'' will initially mean the North Germanic, or Nordic,
languages, while keeping their relation to the wider Germanic family visible:
\[
\text{Indo-European}\to\text{Germanic}\to\text{North Germanic}.
\]
The principal modern North Germanic languages are Icelandic, Faroese,
Norwegian, Danish, and Swedish. Old Norse is the medieval North Germanic
language of the sagas and much of the surviving Scandinavian literary
tradition. Icelandic is especially useful for reading this tradition because
its written language remains comparatively conservative, although modern
Icelandic is not simply identical with Old Norse.

This chapter will return later to the relations among Old Norse, Icelandic,
Old English, and the other Germanic languages, and to the literary question
that motivates the present note: why Northern languages and literatures were
so important to Borges. For now the Homer--Iceland comparison above is a
first concrete object to keep in view: similar metaphors may arise
independently, and resemblance by itself does not establish borrowing.

\subsection{A Finnish parallel: the living among the dead}

Butcher and Lang point to another northern comparison in their note to
\emph{Odyssey} 11.156 (p.~177 in the 1912 Macmillan edition). Anticleia,
meeting her still-living son among the dead, asks in astonishment how he has
come beneath the darkness while still alive. The editors compare this with
the \emph{Kalevala}: the living Vainamoinen enters Tuonela, the Finnish realm
of the dead, and the daughters of Tuoni likewise ask how a man who has not
died by sickness, sorrow, iron, or steel has come there.

This comparison broadens the present Northern thread beyond the North
Germanic languages. Finnish is \emph{not} Germanic or Indo-European: it
belongs to the Uralic family. The \emph{Kalevala} is a nineteenth-century
Finnish national epic compiled by Elias Lonnrot from oral poetry. Thus it
belongs here as a comparative northern literary reference, not as evidence
that Greek and Finnish descend from a common linguistic tradition.

Together with the Icelandic \emph{tann-garðr} comparison above, this gives a
second concrete question for later study of Borges's interest in northern
languages and literatures: what did he find in Old Norse/Icelandic and in
Finnish epic traditions, and how did he distinguish linguistic kinship from
literary analogy?

\section{African languages in Brazil}
\label{linguistics-section-african-words-salvador}

These notes began from words and names that
called my attention at MUNCAB and
the Casa do Benin on September 26, 2026. The
first methodological lesson is
that an African word is not a sufficient
classification: the visit produced
Yoruba-related vocabulary, Bantu-language
material, and Brazilian diasporic
categories.

Museum spelling, Brazilian religious
spelling, and fully tone-marked African
orthography are not always the same. The
forms below are therefore preserved
first as encountered; full Yoruba tone marks
and underdots should be restored
from specialist language sources rather than
guessed.

\medskip\noindent
{\bf Yoruba, Anago, and Nago.}
\emph{Yoruba} names a language and broad
ethnolinguistic/cultural field of
West Africa, especially in present-day
Nigeria and Benin. \emph{Nago} is a
diasporic category that became historically
important in Brazil, especially
Bahia; it overlaps strongly with the Yoruba
world but is not simply the modern
African self-name. The historical relation
among Yoruba, Anago, and Brazilian
Nago belongs to the Atlantic history of
naming and ethnogenesis
\cite{lovejoyojo2015lucumi,pares2010nago}.

This distinction became visually memorable
because the MUNCAB map labelled a
Yoruba region in Africa, while Brazilian
sources on Candomble spoke of Nago
cosmology and Ketu--Nago traditions.

\medskip\noindent
{\bf Omi Tutu.}
The MUNCAB exhibition title was \emph{Omi
Tutu---Agua fresca}.
\emph{Omi} means water, while \emph{tutu}
belongs to the semantic field of
coolness or freshness. In the exhibition, the
expression also carried ideas
of balance, conciliation, vitality, and
renewal \cite{muncab2026omitutu}.

\medskip\noindent
{\bf Ori Eja and Iya Eja.}
Two labels beside Hariel Revignet's works
read \emph{Ori Eja} and
\emph{Iya Eja}. A first lexical decomposition
gives:
\begin{itemize}
\item \emph{ori}: head, with a wider
      religious field involving personhood or
destiny;
\item \emph{eja}: fish;
\item \emph{iya}: mother.
\end{itemize}
Thus expressions such as fish head or head of
fish, and mother of fish, are
useful first glosses, but they are not yet a
full grammatical translation of
the artist's titles. Yoruba tone and syntax
still need checking.

\medskip\noindent
{\bf Orun, Aiye, and Olorun.}
The film title \emph{Orun Aiye---A criacao do
mundo} made three similar forms
easy to confuse:
\begin{itemize}
\item \emph{Orun}: the spiritual or invisible
      realm;
\item \emph{Aye}/\emph{Aiye}: the earthly
      world;
\item \emph{Olorun}: a name/title for the
      supreme divinity.
\end{itemize}
The museum phrase that Mae Stella returned to
Orun after death therefore does
not mean that she became Olorum.

\medskip\noindent
{\bf Ile, axe, and sacerdotal titles.}
\emph{Ile} means house/home and appears in
names such as Ile Axe Opo Afonja
and Ile Aiye. Proper names should not be
translated by mechanically joining
dictionary glosses, but the components are
still linguistically informative.

Brazilian \emph{axe} continues Yoruba
\emph{ase}. Its field includes power,
authority, efficacy, and the capacity for
something to come to pass; reducing
it simply to the vague English word energy
loses important structure
\cite{abiodun2014yoruba}.

Likewise, Brazilian \emph{orixa} continues
Yoruba \emph{orisa}. The title
\emph{ialorixa} contains the element for
mother and \emph{orixa};
\emph{babalorixa} contains \emph{baba},
father. This lexical transparency
does not by itself explain the ritual
authority of those offices.

\medskip\noindent
{\bf Exu, Pade Ona, and Bara.}

The Sandro Aiye exhibition title was
\emph{Pade Ona---Encontrar Caminhos}.
MUNCAB explained \emph{pade} as an
inaugural gesture that creates conditions
for beginning and circulation, and
\emph{ona} as path
\cite{muncab2026padeona}.

The photographed sculpture was
\emph{Bara II}.  The exact linguistic
relation between \emph{Bara} and names or
aspects of Exu was not established by the
label.  It remains a term to investigate,
not an etymology to guess.

\medskip\noindent
{\bf Ofa.}
The iron sculpture beside Diogum's biography
prompted the word \emph{ofa},
the bow and arrow associated with Oxossi.
This is a good example of a religious
term that travelled with knowledge and was
adapted to Brazilian Portuguese
spelling.

\medskip\noindent
{\bf Iyami.}
Nadia Taquary's panel used \emph{Iyami} and
rendered the term as
\emph{feiticeiras}. The beginning of the word
is related to the Yoruba word
for mother, but the concept belongs to a
larger field of ancestral maternal
and feminine power; Portuguese
\emph{feiticeira} or English witch should not
be treated as exhaustive definitions
\cite{santos2008iyami}.

The comparison with Circe is worth preserving
as a comparative-literary
question about female power, transformation,
and fear. It does not establish
a historical relation between Greek and
Yoruba traditions.

\medskip\noindent
{\bf Funfun Dudu.}
The title of the Alberto Pitta book
\emph{Funfun Dudu} called attention to
Yoruba colour vocabulary: \emph{funfun} is
associated with white and
\emph{dudu} with black/dark. The linguistic
note should remain subordinate
to the book's larger art-historical project.

\medskip\noindent
{\bf Mironga.}
The Casa do Benin exhibition used
\emph{mironga} in the semantic field of
mystery, secret, guarded knowledge, spell,
and strategy, and used the verb
\emph{mirongar} for working an image through
what can be shown and what
remains hidden.

Caldas Aulete treats \emph{mironga} as a
variant of \emph{milonga}; for the
Afro-Brazilian religious sense, it derives
\emph{milonga} from Kimbundu
\emph{mi'lona}, meaning word
\cite{aulete-milonga}. This gives a
Bantu/Angolan
linguistic layer rather than a Yoruba one.
The exact historical morphology
and the semantic path from word to mystery,
spell, or strategic concealment
still deserve checking in a specialist
Kimbundu source.

\medskip\noindent
{\bf Zumbi and Nzambi.}
The statue of Zumbi dos Palmares prompted a
memory of \emph{Nzambi} or
\emph{Nzambi a Mpungu}, a divine name in
Kongo/Central African religious
contexts. Zumbi, the historical person, and
Nzambi, the divine name, are
different referents. A linguistic relation
should not be asserted merely
from phonetic resemblance.

\medskip\noindent
{\bf Angola as two kinds of name.}
Angola is the name of a modern African state
and also the name of a Candomble
nation in Brazil. The second use emerged
through diasporic histories of
Central African peoples and Bantu languages;
it does not mean that the modern
country was transplanted intact into Brazil.

This is one example of a general warning:
modern political borders,
ethnolinguistic regions, and Brazilian
diasporic categories are different
maps.

\medskip\noindent
{\bf Akan and museum geography.}
A MUNCAB mask label said Povos Akan---Africa
Central, while a map in the
same museum placed the Akan region around
Ghana and Cote d'Ivoire, in West
Africa. The contradiction should be preserved
rather than silently repaired.

Akan names a major West African
ethnolinguistic field; it is not another word
for Yoruba, Jeje, or a Candomble nation.

\medskip\noindent
{\bf Terms still to investigate.}
The next linguistic pass should include:
\begin{itemize}
\item reliable Yoruba tone marks and
      underdots for all terms above;
\item the exact syntax of \emph{Ori Eja} and
      \emph{Iya Eja};
\item the historical relation among Yoruba,
      Anago, Nago, and Cuban Lucumi;
\item the morphology of \emph{ase},
      \emph{orisa}, \emph{iyalorisa}, and
\emph{babalorisa};
\item a specialist Kimbundu etymology of
      \emph{mironga}/\emph{milonga};
\item the history of the names Zumbi and
      Nzambi.
\end{itemize}


% --- End linguistics.tex ---

% --- Begin religion.tex ---
\chapter{Religion}
\label{chapter-religion}





\section{Method}
\label{religion-section-religion-method}

This notebook began with visits in
Salvador on September 26, 2026.  It
separates four layers:

\begin{itemize}
\item what Daniel directly saw;
\item what a museum label or curatorial
text said;
\item what a book or scholarly source
supports;
\item what remains an interpretation or
question.
\end{itemize}

The notebook itself records the photographed
objects and panels that triggered the
questions below; no separate field-notes
notebook is assumed.

\section{Rosary and Marian titles}
\label{religion-section-rosary-marian-titles}

The Rosary is both a Catholic form of
prayer and the string of beads used to
count its repeated prayers.  Our Lady of
the Rosary is Mary under a devotional
title.  She is not a different saint.

The same distinction applies to Our Lady
of Carmel, of the Conception, of Sorrows,
of the Good Death, and other invocations.
They are titles of Mary associated with
different doctrines, places, orders,
images, and devotional histories.

Our Lady of Carmel belongs to the
Carmelite tradition, whose name derives
from Mount Carmel.  The scapular is an
important Carmelite devotional object.
This tradition also explains Carmo
place-names in Salvador.

The Immaculate Conception concerns Mary's
own conception.  Catholic doctrine holds
that she was preserved from original sin
from its first instant.  This must not be
confused with the virginal conception of
Jesus.

\section{Rosario dos Pretos}
\label{religion-section-rosario-dos-pretos}

The Church of Our Lady of the Rosary of
the Black People was found closed, so its
interior was not visited.

Its historical importance lies in the
Black Catholic brotherhood associated
with it.  The brotherhood was erected in
1685, and construction of the church began
in the early eighteenth century
\cite{ipac-rosario}.

A brotherhood was not merely a prayer
group.  Such institutions organized
worship, funerals, mutual aid, festivals,
and forms of collective belonging within
a slave society.

This gives a precise question for later
study: how did enslaved and free Black
people create institutions of their own
inside the Catholic order of colonial
Brazil?

\section{Jesuits and the Cathedral}
\label{religion-section-jesuits-cathedral}

The Jesuits are members of the Society of
Jesus, a Catholic religious order founded
by Ignatius of Loyola and his companions
and approved in 1540.

Education became one of their central
activities.  This makes the connection
between the colonial Jesuit college in
Salvador and modern Jesuit institutions
such as PUC--Rio and the Universidad
Iberoamericana in Mexico conceptually
real, although the institutions belong to
very different historical periods.

Jesuits arrived in Bahia with Tom\'e de
Sousa in 1549.  The present Cathedral
Basilica was built as the church of the
Jesuit college and later became the
cathedral.

A cathedral is the church that contains
the bishop's \emph{cathedra}, or seat.  The
word names an ecclesiastical function,
not simply a large or ornate building.

\section{Saints, archangels, and dates}
\label{religion-section-saints-archangels-dates}

Saint Michael is celebrated on September
29 together with the archangels Gabriel
and Raphael in the Roman Catholic
calendar.

During the visit, a possible local Mass
for Saint Michael near the interdicted
Church of Saint Michael was discussed.
The exact 2026 time and place were not
confirmed.  It must therefore remain an
open logistical question, not a recorded
event.

\section{What Candomble is}
\label{religion-section-what-candomble-is}

Candombl\'e is a Brazilian religion.  It
formed in Brazil through the preservation,
reorganization, and transformation of
several African religious traditions under
slavery and diaspora.

This does not mean that its divine beings
or all its concepts originated in Brazil.
Orix\'as, voduns, ancestors, divination,
ritual languages, and many other elements
have African histories.  The religious
institutions called Candombl\'e are
Afro-Brazilian formations.

Three identifications must therefore be
avoided:

\begin{itemize}
\item Africa is not one culture;
\item Yoruba religion is not identical
with Candombl\'e;
\item Candombl\'e is not identical with
Umbanda.
\end{itemize}

Umbanda is also Brazilian, but it formed
through a later and different history,
strongly associated with urban southeastern
Brazil in the twentieth century
\cite{prandi1990umbanda}.

\section{Candomble nations}
\label{religion-section-candomble-nations}

A Candombl\'e ``nation'' is not a modern
nation-state.  It is a historical-ritual
classification involving language,
lineage, liturgy, divine beings, music,
and practice.

A first map is:

\begin{itemize}
\item Ketu or Nag\^o traditions are mainly
connected with Yoruba inheritances and
orix\'as;
\item Jeje traditions are mainly connected
with Ewe--Fon inheritances and voduns;
\item Angola and Congo traditions are
mainly connected with Central African,
Bantu-language inheritances and
\emph{nkisi}, called inquices in Brazil.
\end{itemize}

This is only a first map.  Historical
labels include Ketu, Jeje, Nag\^o, Ijex\'a,
Congo, Angola, and others, and ritual
languages are one important marker among
several \cite{castro1981lingua}.

The modern country Angola and the
Candombl\'e nation Angola are therefore
related historically but are not the same
kind of object.

\section{Yoruba, Nago, and Ketu}
\label{religion-section-yoruba-nago-ketu}

Yoruba names an African language and a
broad historical and cultural field.
Nag\^o is a diasporic category that became
especially important in Bahia.  It cannot
be treated as a mathematically exact
synonym for Yoruba.

Ketu is one major tradition inside the
broader Nag\^o/Yoruba field of Brazilian
Candombl\'e.  The prominence of Nag\^o
traditions in Bahia was itself historically
constructed rather than timeless
\cite{pares2010nago}.

No single orix\'a is simply ``the orix\'a
of Ketu''.  A Ketu house cultivates many
orix\'as.  Oxal\'a is a major orix\'a
associated with creation, while a
particular house can have a patron orix\'a.

\section{Supreme divinity and orixas}
\label{religion-section-supreme-divinity-orixas}

Ol\'od\`umar\`e and Olorum are names or
titles used for the supreme divinity in
Yoruba-derived cosmologies.  The orix\'as
are not simply interchangeable names for
this supreme divinity.

The visit repeatedly produced the need to
distinguish:

\begin{itemize}
\item \emph{Orun}, the spiritual or
invisible realm;
\item \emph{Aye} or \emph{Aiye}, the
visible earthly realm;
\item Olorum, the supreme divinity;
\item the orix\'as, such as Oxal\'a,
Xang\^o, Oxum, Ox\'ossi, Ogum, and others.
\end{itemize}

The film title
\emph{Orun Aiye---A cria\c{c}\~ao do mundo}
made the first two terms memorable.  A
creation narrative centred on Obatal\'a or
Oxal\'a belongs to one important body of
Yoruba religious narrative
\cite{marins2012obatala}.

\section{Orun Aiye and Orunmila}
\label{religion-section-orun-aiye-film}

The stop-motion film
\emph{Orun Aiye---A criacao do mundo}
made the cosmological vocabulary concrete.
It was directed by Jamile Coelho and Cintia
Maria.  In the voice cast, Carlinhos Brown
plays Oxala, Joao Miguel plays Olodumare,
and Jorge Washington plays Orunmila
\cite{orun-aiye-sepromi}.

Jorge Washington is also a long-standing
actor of the Bando de Teatro Olodum.  This
places the film inside a wider Black
performing-arts history in Salvador, rather
than treating the names in the credits as
isolated celebrity appearances
\cite{bando-jorge-washington}.

\section{Candomble and Santeria}
\label{religion-section-candomble-santeria}

Cuban Santer\'\i a, or Regla de Ocha, also
cultivates divine beings called
\emph{orichas}.  This does not make it the
same religion as Candombl\'e.

The useful comparison is historical:
Yoruba religious inheritances were
reorganized in different American
societies.  Nag\^o in Brazil and Lucum\'\i
in Cuba are diasporic categories with
different histories
\cite{lovejoyojo2015lucumi}.

\section{Terreiro and leadership}
\label{religion-section-terreiro-leadership}

A terreiro is not merely a building.  It is
a sacred territory, a ritual institution,
a community, and a network of initiated
relations.

An \emph{ialorix\'a} is the female
sacerdotal leader of a terreiro.  The
masculine counterpart is a
\emph{babalorix\'a}.  An ialorix\'a leads a
house; she does not command an entire
Candombl\'e nation.

There is no single central authority that
governs all Candombl\'e houses.  Houses can
be joined by initiatory genealogies,
alliances, influence, and shared ritual
histories while retaining their own
leadership and singularity.

\section{Ile Axe Opo Afonja}
\label{religion-section-opo-afonja}

Il\^e Ax\'e Op\^o Afonj\'a is a major
Ketu/Nag\^o terreiro in Salvador.  M\~ae
Aninha founded it in 1910 in a lineage
connected with Casa Branca do Engenho
Velho.  Xang\^o Afonj\'a is the patron of
the house \cite{barata2025ancestralidade}.

Its current leadership belongs to M\~ae Ana
de Xang\^o, the sixth ialorix\'a in the
line presented by the 2025 book.

The book
\emph{Pedagogia da ancestralidade}
should be read as a deep study of one
especially influential house, not as a
universal manual of Candombl\'e
\cite{barata2025ancestralidade}.

Its central idea is a pedagogy of the
terreiro.  Knowledge is transmitted through
coexistence, observation, participation,
oral memory, hierarchy, practice, care,
and the example of more experienced
members.

The house also developed civil,
educational, and cultural institutions,
including a school, library, Casa do Alak\'a,
and Museu Il\^e Ohun Lailai.  These do not
turn the terreiro into a secular museum;
they show that ritual, education, memory,
and social organization coexist.

\section{Mystery and limits of access}
\label{religion-section-mystery-limits}

The Afonj\'a book insists on preserving the
mystery of ritual life.  Photographs and
recordings may be authorized in community
contexts while remaining forbidden during
sacred rites
\cite{barata2025ancestralidade}.

This matters for any proposed visit.  A
terreiro is not a tourist attraction whose
meaning is exhausted by access.  Permission,
conduct, photography, and the difference
between public and restricted spaces are
part of the religious reality.

\section{Mae Stella de Oxossi}
\label{religion-section-mae-stella}

M\~ae Stella de Ox\'ossi, Maria Stella de
Azevedo Santos, lived from 1925 to 2018.
She led Op\^o Afonj\'a as its fifth
ialorix\'a from 1976.

The three MUNCAB timeline panels mattered
because they made her public and
intellectual trajectory concrete:

\begin{itemize}
\item she was initiated by M\~ae Senhora in
1939 and received the name Od\'e Kayod\^e;
\item she worked for decades in nursing and
public health;
\item in 1981 she travelled to Yoruba
centres in Nigeria, including Oshogbo,
Il\'e-If\`e, and Ed\'e;
\item in 1983 she intervened publicly in
the debate on syncretism during the
international orix\'a conference in
Salvador;
\item in 1986 she participated in the next
world conference in New York;
\item in 1987 she travelled to Benin in a
Brazilian cultural delegation organized by
Pierre Verger;
\item \emph{Meu tempo e agora} was
published in the early 1990s; Bahian
institutional records date the first
edition to 1993;
\item in 1999 the terreiro was protected by
IPHAN during her leadership;
\item biographical sources record that she
was among the Brazilian women indicated in
the 2005 Nobel Peace Prize initiative;
\item in 2013 she became the first
ialorix\'a and first Black woman to occupy
a chair in the Academia de Letras da Bahia;
\item her death in 2018 was described in
the religious language encountered during
the visit as a return to Orun.
\end{itemize}

These milestones match the MUNCAB timeline
and Bahian institutional biographical
records \cite{alba-mae-stella,mae-stella-alb}.
They explain why the Op\^o Afonj\'a book is
not merely a local curiosity: it gives
access to the house in which one of the
most publicly and intellectually prominent
ialorix\'as of recent Brazilian history
worked and led.

It is still better to say ``one of the most
prominent'' than to declare M\~ae Stella
the single most famous ialorix\'a, since
that would require an artificial ranking.

\section{The 1983 manifesto}
\label{religion-section-manifesto-1983}

MUNCAB reproduced the manifesto
\emph{O p\'ublico e o povo do Candombl\'e},
published in Salvador on August 12, 1983.
The panel attributed it to Stella de
Ox\'ossi, Menininha do Gantois, Tet\^e de
Ians\~a, Olga de Alaketu, and Nicinha do
Bogum.

The manifesto is a primary source.  It
rejects descriptions of Candombl\'e as a
sect or primitive animism.  It also rejects
the idea that Candombl\'e requires Catholic
forms in order to be legitimate.

Its argument about syncretism is historical
and political.  It connects concealment
and Catholic forms with slavery and
submission, while demanding the right to
practise Candombl\'e in its own terms.

It also attacks folklorization,
commercialization, press sensationalism,
and tourist consumption.  Its reference to
the Lavagem do Bonfim shows that a practice
can be religiously and politically debated
by Candombl\'e leaders themselves.

The manifesto does not claim that
syncretism never existed.  It disputes what
that history should require in the present.


\section{Water and maternity in Omi Tutu}
\label{religion-section-omi-tutu-religion}

Hariel Revignet's MUNCAB exhibition
\emph{Omi Tutu---Agua fresca} placed
water in relation with territory, body,
displacement, continuity, and return
\cite{muncab2026omitutu}.  The curatorial
text photographed during the visit also
brought Oxum, Iemanja, and Nana into a
field of fresh water, sea, maternity,
memory, gestation, and transformation.

This is useful religious vocabulary, but
a curatorial text is not a complete
theology of the orixas.  The presence of an orixa in an artwork also does not
make the work specifically Umbandist: orixas belong to Yoruba-derived
religious worlds and are central to Candomble traditions as well.  The exhibition
made those names and associations
memorable through an artistic
interpretation; later study should compare
that interpretation with sources internal
to the relevant Candomble traditions.

\section{Ofa and Oxossi}
\label{religion-section-ofa-oxossi}

A small iron sculpture by Diogum prompted a
correct visual recognition.  Its curved
form and central shaft represented the
\emph{ofa}, the bow and arrow associated
with Ox\'ossi.

The heading ``Di Ogum'' named the artist,
Diogo Oliveira, and his relation to Ogum.
It did not mean that every symbol in his
work represented Ogum.

The case is useful because material,
artist's identity, and represented attribute
can point in different directions.

\section{Airá e Oxalufã}
\label{religion-section-aira-oxalufa}

Nota de 15 de outubro de 2026.
Airá (também escrito Ayrá) não é uma
forma de Ogum. Sua associação é com
Xangô: encontra-se o nome “Xangô Airá”,
mas há tradições que o reconhecem como
orixá distinto.\footnote{
Correio Nagô, “Ayrá: o manto branco
da justiça”, 29 de junho de 2015:
\href{https://correionago.com.br/%
ayra-o-manto-branco-da-justica/}
{artigo}; e a fonte de terreiro abaixo.
Consultas em 15 de outubro de 2026.}

Numa versão do relato, Oxalufã, o velho
Oxalá, passa sete anos injustamente
preso no reino de Xangô, sem que este
saiba. Xangô descobre a injustiça e o
liberta. Como Oxalufã está debilitado,
Xangô pede a Airá que o acompanhe na
volta para casa. Airá o carrega nas
costas, acende fogueiras para aquecê-lo
e lhe conta histórias. Tornam-se
companheiros, numa relação descrita
como a de pai e filho.\footnote{
Ilê Àsé Oni Èwé Ossayn Atin L'Oyá,
“Ayrá (Airá)”, seção
“História de Ayrá”:
\href{https://www.bruxodailha.com.br/%
orixas/ayra/}{página do terreiro}.
Consulta em 15 de outubro de 2026.}

Nesta versão, Xangô liberta; Airá
carrega e cuida. Os detalhes variam
entre narrativas e terreiros.
A mesma página atribui a Airá uma chave
ligada ao clima: fica registrada como
atribuição dessa fonte, não como
símbolo universal nem identificação
segura de uma imagem isolada.

\section{Beads and cowries}
\label{religion-section-beads-cowries}

A bead is first a material object.  Its
religious meaning comes from arrangement,
colour, preparation, use, lineage, and
context.

Strings of beads can identify religious
relations and may be ritually prepared.
One cannot infer an orix\'a mechanically
from any necklace or colour combination.

Three encounters made the material central:
Hariel Revignet's bead installation,
N\'adia Taquary's strings in the colours of
Iroko, and Sandro Aiy\^e's wooden sculpture
with beads.

Cowries appeared as actual materials in
Hariel Revignet's works.  Their visual and
religious histories require a separate
study; they should not be reduced to a
generic sign of ``Africa''.

\section{Iroko and Iyami}
\label{religion-section-iroko-iyami}

N\'adia Taquary's
\emph{Ir\^ok\^o: A \`arvore c\'osmica}
connected Iroko with time, ancestry, the
life cycle, the gameleira, and the Iyami.

The museum rendered Iyami as
``feiticeiras''.  That translation is too
narrow if taken as a definition.  The
concept belongs to a larger Yoruba field of
ancestral maternal and feminine power
\cite{santos2008iyami}.

The comparison with Circe is legitimate as
a comparative question about female power,
transformation, and fear.  It does not
establish transmission or historical
kinship between Greek and Yoruba myths.

\section{Bara, Exu, and paths}
\label{religion-section-bara-exu-paths}

Sandro Aiy\^e's sculpture
\emph{B\'ar\`a II} appeared in the MUNCAB
exhibition
\emph{Pad\^e On\~a---Encontrar Caminhos}.
The official framing presents Exu through
movement, mediation, circulation, and the
opening of paths
\cite{muncab2026padeona}.

This is better than identifying the work
only from a trident-like object.  An
attribute can call attention, but the
exhibition title and curatorial context are
stronger evidence.

\section{Saint George and protection}
\label{religion-section-saint-george}

Saint George is a Christian warrior saint.
The plant called
\emph{lan\c{c}a-de-S\~ao-Jorge} takes its
popular name from the lance-like form and
the saint's martial image.

The plant is also associated in popular
practice with protection.  Associations
between Saint George and Ogum or Ox\'ossi
vary by region and religious house.  They
must not be stated as universal equations.

\section{Bembe and public religion}
\label{religion-section-bembe-public-religion}

The book on the Bemb\'e do Mercado opens a
question distinct from the life of a single
terreiro: how does Candombl\'e organize
sacred territory and collective action in
public space?

The Bemb\'e is celebrated in Santo Amaro
and involves many terreiros, public ritual,
heritage, urban territory, and memories of
abolition \cite{ufrb2026bembe}.

It belongs in the religion notebook, but
also in the histories of the Rec\^oncavo,
public space, and cultural heritage.

\section{Reading sequence}
\label{religion-section-religion-reading}

A sensible sequence after the visit is:

\begin{enumerate}
\item read
\emph{Pedagogia da ancestralidade}
as a study of Op\^o Afonj\'a;
\item read a reliable panoramic account of
Candombl\'e nations;
\item read
\emph{Na\c{c}\~oes do Candombl\'e: Jeje}
as a deliberately particular Jeje study.
The copy encountered at MUNCAB identifies
itself as Book 1 of the larger
\emph{Na\c{c}\~oes do Candombl\'e}
project; it should not be mistaken for a
general survey of the whole religion
\cite{ufrb2026nacoes};
\item compare Ketu/Nag\^o, Jeje, and
Angola/Congo without treating one as the
whole religion;
\item later compare Candombl\'e with Umbanda
and Santer\'\i a historically.
\end{enumerate}


% --- End religion.tex ---

% --- Begin arte-moderna-brasileira.tex ---
\chapter{Arte moderna brasileira}
\label{chapter-arte-moderna-brasileira}





\section{Abstração e arte brasileira no pós-guerra}

\subsection{Um mapa não linear}

A história da arte brasileira do pós-guerra não deve ser entendida como uma
sequência linear em que figuração, abstração, concretismo e neoconcretismo se
substituem. Em particular, ``figuração'' não é o nome de um primeiro movimento:
é uma categoria ampla para obras que representam figuras, objetos ou espaços
reconhecíveis.

Nos anos 1940--1960 coexistem diferentes vertentes abstratas. Duas famílias
úteis para organizar o período são a abstração informal e a abstração
geométrica. A primeira privilegia gesto, matéria, ritmo, cor e subjetividade; a
segunda privilegia relações formais e construção geométrica. Dentro desta
última se desenvolve o concretismo brasileiro nos anos 1950 e, em diálogo e
ruptura com ele, o neoconcretismo a partir de 1959.

\subsection{Concretismo e neoconcretismo}

A arte concreta busca uma linguagem visual autônoma, fundada na ordem, na
racionalidade e na construção objetiva da forma. No Brasil, um núcleo decisivo
foi o Grupo Ruptura, em São Paulo. No Rio de Janeiro, o Grupo Frente teve Ivan
Serpa como figura central e uma prática mais heterodoxa dentro do campo
geométrico.

O neoconcretismo reage ao caráter excessivamente rígido e racionalista do
programa concreto. Conserva a investigação geométrica, mas introduz uma relação
mais sensível, fenomenológica, espacial e participativa com a obra. Lygia
Clark, Lygia Pape e Hélio Oiticica são nomes centrais desta passagem. Uma
fórmula útil é: no concretismo, geometria construída; no neoconcretismo,
geometria experimentada.

Ivan Serpa é particularmente importante como figura de ligação. Além de sua
produção, sua atividade docente no MAM Rio ajudou a formar o ambiente artístico
carioca do qual emergiram experiências neoconcretas e posteriores.

Artistas como Alfredo Volpi, Maria Leontina e Milton Dacosta dialogam fortemente
com a abstração geométrica do período sem se reduzirem simplesmente aos
programas doutrinários concreto ou neoconcreto. Em Volpi, por exemplo, a
geometria das fachadas e bandeirinhas conserva uma qualidade artesanal e
pictórica muito própria.

\subsection{Abstração informal}

Contemporaneamente ao concretismo, desenvolve-se no Brasil uma abstração
informal. Aqui a tela é entendida como campo de investigação do próprio ato de
pintar: gesto, manchas, tensões, matéria e cor ganham autonomia. Essa família
tem paralelos internacionais no expressionismo abstrato norte-americano e no
Art Informel europeu.

Tomie Ohtake é uma referência importante. Sua produção dos anos 1950--1960
explora gesto, matéria e grandes relações cromáticas, em contraste com a
precisão programática da arte concreta. Algumas obras, como sua \emph{Pintura}
de 1964, podem lembrar visualmente os campos de cor de Mark Rothko, embora a
superfície de Ohtake preserve uma presença própria do gesto e da matéria.

\subsection{Genealogia internacional da abstração geométrica}

O concretismo brasileiro pertence a uma constelação internacional mais ampla de
arte abstrata e construtiva. Uma genealogia simplificada passa pelo cubismo do
início do século XX e, mais diretamente, por Malevich, Mondrian, o
construtivismo russo, De Stijl e a Bauhaus. Nos anos 1950--1960 surgem ou se
consolidam diferentes respostas locais: concretismo e neoconcretismo no Brasil,
Op Art na Europa, hard-edge painting nos Estados Unidos e tradições
construtivas importantes na Argentina e no Uruguai.

Piet Mondrian é um ancestral particularmente direto. Depois do contato com o
cubismo em Paris, sua pintura passa gradualmente da representação de árvores e
paisagens à organização autônoma de linhas, planos e cores. No neoplasticismo,
verticais e horizontais, cores primárias e relações assimétricas de equilíbrio
constituem a própria linguagem da pintura. Sua obra mostra como o cubismo é uma
das grandes rupturas que tornam possível a abstração geométrica, embora esta
não possa ser reduzida simplesmente a uma derivação do cubismo.

\subsection{Brasil e México}

Há afinidades fortes entre a abstração geométrica brasileira e a arte mexicana
do pós-guerra, mas não existe uma correspondência simples entre concretismo e
neoconcretismo brasileiros e movimentos mexicanos de mesmo tipo. No México,
a chamada Generación de la Ruptura e artistas como Vicente Rojo, Mathias
Goeritz e Gunther Gerzso participam da abertura para linguagens abstratas
internacionais, num contexto marcado também pela reação à hegemonia da Escola
Mexicana de Pintura e do muralismo.

Vicente Rojo (1932--2021), ativo no México sobretudo a partir dos anos 1950,
tem afinidades visuais com a abstração geométrica brasileira: construção em
séries, módulos, repetição, cor e uma relação importante com o design gráfico.
Maria Leontina e Milton Dacosta oferecem comparações particularmente naturais.
A relação não deve, porém, ser confundida automaticamente com contato pessoal
entre esses artistas.

Um vínculo histórico México--Brasil mais concreto aparece na poesia concreta.
O grupo brasileiro Noigandres participa de uma circulação internacional que
chega ao México; Mathias Goeritz organiza em 1966, na UNAM, a exposição
\emph{Poesía Concreta Internacional}, na qual Vicente Rojo participa. Assim,
geometria, tipografia, poesia, livro e design constituem uma zona real de
intercâmbio entre os dois países.

\subsection{Retorno da figuração e arte experimental}

Nos anos 1960--1980, a figuração reaparece com novas funções. A Nova Figuração
brasileira incorpora cultura de massa, imprensa, televisão e referências
cotidianas, e no contexto da ditadura militar torna-se também veículo de
crítica social e política. Não se trata de retornar à figuração anterior à
abstração: a imagem retorna depois do cubismo, da abstração e da cultura de
massa.

Ao mesmo tempo, a arte experimental e conceitual radicaliza outra pergunta. Em
vez de buscar apenas uma nova linguagem pictórica, questiona o próprio objeto
artístico, seus suportes, sua circulação e as instituições que o legitimam. A
ênfase pode deslocar-se do objeto para o conceito ou procedimento. Performance,
fotografia, intervenção urbana, corpo e ready-made passam a integrar o campo da
arte. No Brasil, esse desenvolvimento adquire uma dimensão política específica
sob a ditadura.

\subsection{Depois de 1990}

A partir dos anos 1990 torna-se menos convincente narrar a arte brasileira como
uma sucessão de movimentos estilísticos dominantes. Figuração, abstração,
pintura, práticas conceituais, performance, vídeo e meios digitais coexistem.
A revisão dos cânones, a ampliação dos agentes reconhecidos pelo sistema da arte
e questões de memória, território, corpo, identidade e pertencimento tornam-se
centrais. A ideia moderna de ruptura continua presente, mas já não produz uma
sequência tão clara de novos ``ismos''.


\section{Salvador: arte, memória e museu}
\label{arte-moderna-brasileira-section-arte-salvador-2026}

Estas notas surgiram da visita ao MUNCAB e à
Casa do Benin em 26 de setembro
de 2026. A distinção metodológica principal é
entre a obra, sua legenda, o
texto curatorial e uma interpretação
posterior. Esses níveis não devem ser
fundidos.

\section{MUNCAB e a 36a Bienal}
A itinerância da 36a Bienal de São Paulo no
MUNCAB usava o título
\emph{Nem todo viandante anda estradas---Da
humanidade como prática}. O texto
curatorial perguntava como a história da arte
produz presenças, ausências,
centralidades e silêncios, e recusava a ideia
de que pertencimento exige uma
identidade homogênea \cite{muncab2026bienal}.

Esse primeiro painel não foi o ponto mais
convincente da visita; a exposição de
Hariel Revignet tornou questões semelhantes
muito mais concretas através de
água, corpo, território e memória.

\section{Hariel Revignet: Omi Tutu}
\emph{Omi Tutu---Água fresca}, de Hariel
Revignet, organiza a água como
elemento de conexão entre territórios,
corpos, deslocamento, continuidade e
retorno. A artista é brasileira-gabonesa; o
Gabão, país da África Central na
costa atlântica, entra na exposição como
território familiar concreto e não
como uma África abstrata
\cite{muncab2026omitutu}.

A imagem curatorial mais memorável foi a do
rio que entra no mar: depois da
confluência já não se distingue o rio, mas
sua água não deixa de existir.
Isso fornece uma imagem forte para pensar
continuidade diaspórica sem supor a
sobrevivência intacta de uma origem.

As legendas \emph{Ori Ejá} e \emph{Iya Ejá},
ambas de 2025, acompanhavam
pinturas em acrílica com bordado de búzios.
Na fotografia feita durante a
visita, uma figura feminina escura aparece
com ventre arredondado, mãos sobre
o corpo e cauda de peixe. Como duas legendas
estavam adjacentes, a fotografia
sozinha não permite decidir qual título
correspondia à tela fotografada; essa
incerteza deve ser preservada.

\section{Rio de Miçangas---Ziguidás}
A instalação \emph{Rio de
Miçangas---Ziguidás} suspendia fios de
miçangas e
sementes sobre recipientes de barro, cabaças
e outros objetos, produzindo
visualmente uma espécie de rio vertical.

A relação com os zigidas/sikidas é biográfica
e material: textos sobre a
trajetória de Hariel descrevem esses fios de
contas usados abaixo do umbigo
como objetos ligados a riqueza, sensualidade
e fertilidade, transmitidos no
caso da artista através de sua avó gabonesa
\cite{hariel-zigida}. A obra
transforma assim um objeto corporal em
instalação espacial, sem que seu sentido
se reduza a uma única função ritual.

\section{Nádia Taquary e Iroko}
Nádia Taquary, em \emph{Iroko: A árvore
cósmica} (2025), usa fios de contas e
uma linguagem escultórica contemporânea para
trabalhar tempo, ancestralidade,
ciclo de vida e referências a Iroko e às
Iyami.

A obra é útil para perceber que vocabulário
religioso pode entrar na arte
contemporânea sem que a obra se torne uma
ilustração etnográfica neutra. A
tradução de Iyami por feiticeiras no painel
do museu é, ela própria, uma
decisão de tradução que deve ser analisada
criticamente.

\section{Sandro Aiyê e Bará}
A escultura \emph{Bará II}, de Sandro Aiyê,
foi identificada na legenda como
madeira de demolição e contas de miçangas.
Ela integrava
\emph{Padê Onã---Encontrar Caminhos},
exposição que o MUNCAB organiza em torno
de Exu, movimento, mediação, circulação e
abertura de caminhos
\cite{muncab2026padeona}.

Isso é mais forte do que identificar a obra
apenas por uma forma semelhante a
um tridente: título, material, exposição e
contexto curatorial trabalham
juntos.

\section{Alberto Pitta e J. Cunha}
Dois livros encontrados no MUNCAB abrem uma
linha importante para a história
da arte baiana.

\emph{Alberto Pitta: Fúnfún Dúdú} reúne a
produção de Pitta em estamparia,
pintura e objetos, articulando terreiro,
carnaval e cidade
\cite{barata2025pitta}. \emph{J. Cunha e o
Carnaval Negro} examina a
trajetória estética e política de J. Cunha e
sua participação na construção da
identidade visual dos blocos afro e do
carnaval de Salvador
\cite{barata2024jcunha}.

Os dois casos mostram que tecido, figurino,
grafismo, bloco afro e carnaval não
são margens da história da arte: são campos
fundamentais de produção visual
negra em Salvador.

\section{Tincoã e Os Tincoãs}
A exposição \emph{Tincoã} colocava em relação
o pássaro tincoã e o grupo vocal
baiano Os Tincoãs. O movimento curatorial
passava do canto do pássaro ao canto
musical, e daí à memória e ancestralidade.

Essa experiência fixou uma regra para estas
notas: em Salvador, arte visual,
música, religião, artesanato e memória
aparecem frequentemente no mesmo campo
cultural; separá-los por mídia pode esconder
mais do que explicar. O próprio
projeto de Mateus Aleluia descreve Os Tincoãs
como pioneiros em levar de forma
consistente à música popular brasileira
universos poéticos ligados ao
Candomblé e à Umbanda
\cite{nacoes-candomble-projeto}.


\section{Orun Aiye e Jorge Washington}
O curta de animação \emph{Orun Aiye---A
criação do mundo}, dirigido por
Jamile Coelho e Cintia Maria, apareceu no
MUNCAB como outra ponte entre arte e
cosmologia. A produção em stop motion
trabalha narrativas de criação ligadas
ao universo iorubá. Nos créditos, Carlinhos
Brown interpreta Oxalá, Jorge
Washington interpreta Orunmilá, João Miguel
interpreta Olodumaré e Carlos
Betão interpreta Bira
\cite{orun-aiye-sepromi}.

Jorge Washington é ator do Bando de Teatro
Olodum, companhia negra criada em
Salvador em 1990
\cite{bando-jorge-washington}. O encontro do
seu nome no filme
é portanto também uma conexão entre teatro
negro baiano, cinema de animação e
repertório religioso afro-brasileiro.

\section{Casa do Benin e fluxo/refluxo}
A Casa do Benin não deve ser lida apenas como
um prédio que contém objetos
africanos. A instituição foi criada em 1988
para materializar intercâmbios
Bahia--Benin, em paralelo com a Casa do
Brasil em Ouidah
\cite{casa-benin-visite}.

Pierre Verger fornece a chave histórica do
fluxo e refluxo entre as duas
margens \cite{verger2021fluxo}; a intervenção
arquitetônica de Lina Bo Bardi
faz com que o próprio edifício participe
dessa reflexão.

\section{O navio da Liberdade}
Emo de Medeiros, em \emph{O navio da
Liberdade} (2025), reelabora o navio
atlântico como imagem de liberdade em vez de
deportação. A obra pertence ao
projeto \emph{Fatumbia}, que o próprio
artista relaciona a Pierre Fatumbi
Verger e à travessia entre Bahia e África
\cite{emomedeiros-fatumbia}.

A obra é contemporânea e inventada; não deve
ser lida como fotografia
documental de um navio histórico.

\section{Mirongar e os limites da
visibilidade}
O texto \emph{Mirongar}, na Casa do Benin,
propunha trabalhar imagens entre o
que se pode mostrar e aquilo que precisa
permanecer guardado. A formulação
produz uma questão estética e ética
importante: tornar algo visível não é o
mesmo que torná-lo disponível para consumo.

Essa questão ecoa o material religioso
encontrado no MUNCAB e no livro do Opô
Afonjá, onde determinados ritos e
conhecimentos preservam deliberadamente
limites de acesso.


% --- End arte-moderna-brasileira.tex ---

% --- Begin historia-do-brasil.tex ---
\chapter{Historia do Brasil}
\label{chapter-historia-do-brasil}





\section{Guerra do Paraguai}
\label{historia-do-brasil-section-guerra-do-paraguai}

A Guerra do Paraguai ocorreu entre 1864 e
1870. Uma maneira simples de entender seu
início é separar duas perguntas: quem começou
a crise regional e quem atacou diretamente o
Brasil primeiro.

Em 1864, o Brasil interveio militarmente no
Uruguai contra o governo dos Blancos. O
presidente paraguaio Francisco Solano López,
aliado aos Blancos, considerava a intervenção
brasileira uma ameaça ao equilíbrio de poder
na região do Prata.

O Brasil, porém, não atacou o Paraguai nesse
momento. Foi o Paraguai que iniciou as
hostilidades diretas contra o Brasil: em
novembro de 1864, o governo paraguaio
apreendeu o vapor brasileiro \emph{Marquês de
Olinda}; em dezembro, tropas paraguaias
invadiram Mato Grosso.

Portanto, a sequência inicial pode ser
resumida assim: o Brasil intervém no Uruguai;
López responde atacando o Brasil. Dizer que
López simplesmente queria conquistar o Brasil
é enganoso. Havia ambições paraguaias de
poder regional, mas a guerra surgiu de uma
crise mais ampla envolvendo Brasil, Paraguai,
Argentina e Uruguai, com disputas de
fronteiras, navegação dos rios e influência
política na Bacia do Prata.

Em 1865, López tentou enviar tropas através
do território argentino para intervir no
Uruguai. O governo argentino recusou a
passagem, e o Paraguai invadiu Corrientes.
Nesse mesmo período, a intervenção brasileira
ajudou a derrubar o governo blanco uruguaio.
Venancio Flores, dos Colorados e aliado do
Brasil, assumiu o poder.

Forma-se então a Tríplice Aliança: Brasil,
Argentina e o novo governo uruguaio contra o
Paraguai. Há aqui uma ironia importante: López
entrou no conflito em parte para defender o
governo blanco do Uruguai contra a intervenção
brasileira; depois que esse governo caiu, o
novo Uruguai passou a combater ao lado do
Brasil contra López.

Uma cronologia mínima do começo da guerra é:

\begin{enumerate}
\item Brasil intervém militarmente no Uruguai
em 1864.
\item Paraguai apreende o \emph{Marquês de
Olinda} em novembro de 1864.
\item Paraguai invade Mato Grosso em dezembro
de 1864.
\item Paraguai invade Corrientes em 1865.
\item Brasil, Argentina e Uruguai formam a
Tríplice Aliança em 1865.
\end{enumerate}

Assim, à pergunta ``quem atacou quem
primeiro?'', a resposta precisa é: o Brasil
atacou primeiro no contexto da crise regional,
mas atacando o Uruguai; entre Brasil e
Paraguai, foi o Paraguai que atacou o Brasil
primeiro.


\section{Salvador e a história colonial}
\label{historia-do-brasil-section-salvador-historia-colonial}

Estas notas nasceram durante a visita a
Salvador em 26 de setembro de 2026. A
primeira pergunta foi muito básica: em que
sentido a história colonial do
Brasil começa em 1500, e em que sentido ela
começa apenas depois?

A resposta exige separar três coisas.
Primeiro, a história indígena não começa
com a chegada dos europeus. Segundo, a
chegada portuguesa em 1500 não significa
conquista imediata do território do Brasil
atual. Terceiro, fundar uma cidade
não é o mesmo que essa cidade já ser um
grande centro econômico ou político.


\section{O Brasil no quadro das Américas}
\label{historia-do-brasil-section-brasil-quadro-americas}

A história do Brasil colonial entra num
processo atlântico mais amplo. A chegada de
Cabral em 1500 não foi o primeiro contato
entre europeus e o continente americano.

Por volta do ano 1000, navegadores nórdicos
mantiveram um assentamento em L'Anse aux
Meadows, na atual Terra Nova. É o sítio
arqueológico europeu autenticado mais antigo
conhecido nas Américas
\cite{unesco-lanse-meadows}.

Em 1492, Colombo chegou ao Caribe. Esse não
foi o primeiro contato europeu com as
Américas, mas iniciou uma conexão atlântica
contínua que rapidamente produziria
conquista, colonização, tráfico, epidemias e
transformações demográficas em escala
continental.

No caso do Brasil, há hipóteses de contatos
em 1498 e janeiro de 1500; o desembarque de
Cabral na Bahia, em abril de 1500, é o
primeiro caso documentalmente seguro dentro
da narrativa portuguesa.

No mundo maia, Colombo encontrou mercadores
maias perto de Honduras em 1502. Em 1511,
náufragos espanhóis chegaram à costa do
atual Quintana Roo; Gonzalo Guerrero e
Jerónimo de Aguilar sobreviveram e viveram
entre comunidades maias
\cite{inah-primeiros-contatos-maias}.

Em 1517, Francisco Hernández de Córdoba
realizou a primeira grande expedição
espanhola deliberada à península de Yucatán.
Em 1519, Cortés iniciou a campanha que o
levaria a Tenochtitlan.

A cronologia mínima fica então:
\begin{itemize}
\item aproximadamente 1000: presença
nórdica em Terra Nova;
\item 1492: Colombo no Caribe;
\item 1498--1500: primeiros contatos
europeus discutidos com a costa brasileira;
\item abril de 1500: Cabral na Bahia;
\item 1502: encontro europeu--maia perto de
Honduras;
\item 1511: náufragos espanhóis em
Quintana Roo;
\item 1517: Hernández de Córdoba em
Yucatán;
\item 1519: início da campanha de Cortés.
\end{itemize}

Essa cronologia é apenas de contatos
europeus. Ela não marca o começo das
histórias americanas, que já tinham
milênios de desenvolvimento independente.
\section{Primeiros contatos europeus}
\label{historia-do-brasil-section-primeiros-contatos-europeus}

A pergunta sobre o primeiro contato europeu
com o território que hoje chamamos Brasil
não tem uma resposta absolutamente segura.
É preciso separar três níveis de evidência.

Uma hipótese atribui a Duarte Pacheco
Pereira uma viagem portuguesa à costa
setentrional da América do Sul em 1498.
Há historiadores que leem seu próprio relato
como registro de uma passagem pela costa do
Brasil; outros consideram a evidência
insuficiente para localizar com segurança a
expedição \cite{pacheco1498,disney-brazil}.

Em janeiro de 1500, o navegador castelhano
Vicente Yáñez Pinzón chegou certamente à
costa nordeste da América do Sul. Uma
tradição identifica seu Cabo de Santa Maria
de la Consolación com o atual Cabo de Santo
Agostinho, em Pernambuco, ao sul do Recife
\cite{amazon-first-contact,cabo-pinzon}.

Essa localização continua discutida.
Outros estudos admitem possibilidades em
Pernambuco, Rio Grande do Norte, Ceará ou
mais ao norte; há também historiadores que
não consideram demonstrado que Pinzón tenha
pisado no território do Brasil atual
\cite{funag-pinzon,disney-brazil}.

Por isso, a formulação mais segura é:
\begin{itemize}
\item possível contato português anterior:
Duarte Pacheco Pereira, 1498, controverso;
\item primeiro contato espanhol bem
documentado com a costa nordeste da América
do Sul: Pinzón, janeiro de 1500, mas com
localização exata discutida;
\item primeiro desembarque europeu em
território brasileiro com evidência
documental considerada irrefutável por
parte importante da historiografia: a frota
de Pedro Álvares Cabral, em abril de 1500,
no litoral da Bahia \cite{disney-brazil}.
\end{itemize}

Assim, a fórmula ``Cabral descobriu o
Brasil'' pertence à narrativa convencional
da expansão portuguesa; ela não resolve,
sozinha, a pergunta sobre qual europeu viu
ou tocou primeiro essa costa.
\section{Uma primeira cronologia}
\label{historia-do-brasil-section-primeira-cronologia-brasil}

A data de 1500 marca convencionalmente a
chegada da expedição de Cabral ao litoral
da atual Bahia. Isso não significa que o
território do Brasil atual tenha sido
conquistado naquele momento.

A presença portuguesa permanente se
intensifica apenas nas décadas seguintes.
São Vicente é fundada em 1532. O sistema de
capitanias hereditárias começa na década de
1530. O Governo-Geral é criado em 1548 e
instalado em Salvador com Tomé de Sousa em
1549 \cite{arquivo-nacional-tome}.

Para as três cidades que se tornarão
centrais na história do Brasil, as datas
convencionais de fundação são:

\begin{itemize}
\item Salvador: 29 de março de 1549;
\item São Paulo: 25 de janeiro de 1554;
\item Rio de Janeiro: 1 de março de 1565.
\end{itemize}

Mas essas três datas não significam a mesma
coisa historicamente.

\section{Salvador, Rio e São Paulo}
\label{historia-do-brasil-section-salvador-rio-sao-paulo}

Salvador nasce como cidade planejada pela
Coroa para ser sede do Governo-Geral. Sua
função inicial é administrativa, militar,
religiosa e portuária. Desde o começo está
ligada ao Recôncavo, ao açúcar e ao
Atlântico. Foi sede do governo colonial até
1763 \cite{arquivo-nacional-salvador}.

São Paulo nasce de outra maneira. A data de
25 de janeiro de 1554 é a fundação do
Colégio de São Paulo de Piratininga pelos
jesuítas, no planalto. O núcleo estava
ligado à catequese e às relações com as
populações indígenas locais. Não era uma
capital colonial nem um grande porto
\cite{prefeitura-sp-fundacao,monteiro2018}.

Rio de Janeiro, por sua vez, é fundado em
1 de março de 1565 por Estácio de Sá em
pleno conflito pela Baía de Guanabara. A
cidade nasce como base militar portuguesa
contra a França Antártica e seus aliados
indígenas \cite{multirio-fundacao-rio}.

Assim, uma primeira comparação é:

\begin{itemize}
\item Salvador: governo, porto e açúcar;
\item São Paulo: colégio jesuíta e núcleo
interiorano;
\item Rio de Janeiro: defesa militar da
Guanabara.
\end{itemize}

A ordem em que foram fundadas não é a ordem
de sua importância econômica posterior.
São Paulo só se torna o principal centro
econômico brasileiro muito mais tarde,
especialmente com café, ferrovias,
imigração, urbanização e industrialização.
O Rio ganha importância crescente com o
ouro e a articulação do centro-sul, até
substituir Salvador como capital em 1763.

\section{Pernambuco: Olinda e Recife}
\label{historia-do-brasil-section-pernambuco-olinda-recife}

Pernambuco entra muito cedo como um dos
principais polos da colonização portuguesa.
A Capitania de Pernambuco foi entregue a
Duarte Coelho, que tomou posse em 1535.
Olinda tornou-se a sede política da
capitania e, graças aos engenhos de açúcar
do entorno, um dos centros mais ricos da
América portuguesa no século XVI
\cite{olinda-historia}.

Recife nasce de maneira diferente. Antes de
ser cidade importante, era o porto de
Olinda. Os arrecifes naturais protegiam o
ancoradouro e, desde o século XVI, o porto
escoava o açúcar produzido na capitania.
No início do século XVII, o porto do Recife
já estava entre os mais movimentados da
América portuguesa \cite{recife-historia}.

Assim, o par funciona inicialmente como uma
cidade alta e um porto:
\begin{itemize}
\item Olinda: sede política, residencial e
religiosa da capitania;
\item Recife: porto comercial ligado ao
Atlântico e ao açúcar.
\end{itemize}

A invasão neerlandesa muda essa relação. Em
1630, os holandeses conquistam Pernambuco.
Olinda é abandonada e incendiada em 1631,
enquanto Recife se torna o centro do
domínio holandês. Sob Maurício de Nassau,
Recife/Mauritsstad recebe pontes, palácios,
aterros e outras obras urbanas
\cite{olinda-historia,recife-historia}.

Depois da expulsão dos holandeses em 1654,
Olinda volta formalmente a ser sede da
capitania, mas a primazia econômica já se
deslocara de forma duradoura para o Recife.

Isso acrescenta um quarto polo ao mapa
inicial do Brasil colonial:
\begin{itemize}
\item Salvador: capital colonial, porto e
açúcar;
\item Olinda--Recife: capitania açucareira,
sede política e grande porto;
\item Rio de Janeiro: fundação militar e,
mais tarde, centro do eixo minerador e
capital;
\item São Paulo: núcleo jesuíta interiorano,
mais tarde ligado à expansão territorial e,
séculos depois, ao café e à
industrialização.
\end{itemize}
\section{Antes de 1500: Bahia e Brasil
indígenas}
\label{historia-do-brasil-section-brasil-antes-1500}

A fundação de Salvador não ocorreu sobre
território vazio. A Baía de Todos-os-Santos
e o Recôncavo eram habitados por diferentes
grupos indígenas, entre eles grupos
Tupinambá. Fontes coloniais e estudos
históricos descrevem sociedades
agricultoras sedentárias, com cultivo de
mandioca, milho, tubérculos e outras plantas,
além de pesca, caça e redes de circulação
\cite{oleary2024horticulture,neubauer2021}.

O mesmo cuidado vale para o Brasil inteiro.
Antes da chegada portuguesa havia enorme
diversidade linguística, política e
econômica. Povos associados às línguas
Tupi ocupavam grande parte do litoral, mas
havia muitas outras famílias linguísticas
e formas de organização.

Também é incorreto imaginar que toda a
Amazônia pré-colonial fosse formada apenas
por pequenos grupos móveis. A arqueologia
documenta paisagens intensamente
transformadas, grandes aldeias, sistemas de
terraplenagem, populações densas em algumas
regiões e sociedades politicamente
complexas \cite{gomes2018indigenous,
roosevelt1999amazonia,
iriarte2018earthworks}.

Portanto, a diferença em relação ao México
não é uma oposição entre civilização e
ausência de civilização.

\section{Por que não houve uma Tenochtitlan
no Brasil?}
\label{historia-do-brasil-section-brasil-mexico-estados}

A pergunta precisa ser formulada com
cuidado. No território do Brasil existiram
sociedades complexas, mas não conhecemos,
em 1500, um Estado centralizado e tributário
comparável ao império Mexica com capital em
Tenochtitlan.

O México central havia passado por séculos
de urbanização, formação de cidades-Estado,
agricultura intensiva e competição entre
polidades. O império Mexica, formado no
século XV, cresceu sobre essa rede já muito
densa de cidades e territórios. Tenochtitlan
era uma grande capital lacustre, integrada a
mercados, sistemas de abastecimento e uma
estrutura imperial de tributos
\cite{smith2021aztec,rojas2016tenochtitlan}.

A organização Tupinambá conhecida no litoral
brasileiro seguia outra lógica. Havia
aldeias sedentárias, agricultura, chefias,
alianças e guerras organizadas, mas a
autoridade dos chefes era muito menos
centralizada. Fontes do século XVI
registraram que cada aldeia ou conjunto de
casas tinha seus próprios chefes e que a
autoridade dependia fortemente do consenso
dos seguidores \cite{monteiro2018}.

Isso tornava possível formar alianças
militares muito amplas sem produzir
necessariamente um aparelho estatal
permanente capaz de cobrar tributos de
centenas de comunidades.

A Amazônia mostra ainda outro caminho.
Existiram grandes populações e formas
complexas de organização regional, mas
muitas parecem ter sido mais
heterárquicas, isto é, com múltiplos centros
de poder, em vez de um império único e
fortemente centralizado
\cite{roosevelt1999amazonia}.

A geografia e a ecologia entram nessa
história, mas não como destino automático.
O Vale do México concentrava população,
terras agrícolas, lagos, mercados e
cidades próximas umas das outras. Isso
favorecia redes urbanas e tributárias. O
território que se tornaria o Brasil era
muito maior e ecologicamente mais diverso,
com populações distribuídas por extensas
redes costeiras e fluviais.

Mas seria errado concluir que a floresta
tropical por si só impedia sociedades
complexas: justamente a arqueologia
amazônica moderna derrubou essa velha
hipótese determinista
\cite{gomes2018indigenous,
roosevelt1999amazonia}.

A diferença, portanto, não é racial nem
uma questão de povos mais ou menos fortes.
É uma diferença histórica de demografia,
formas de agricultura, densidade de
assentamentos, redes comerciais, tradição
urbana, guerra, autoridade política e
relações ecológicas.

\section{Jesuítas e instituições coloniais}
\label{historia-do-brasil-section-historia-jesuitas-salvador}

A Catedral Basílica visitada em Salvador foi
originalmente a igreja do colégio
dos jesuítas. O edifício atual começou a ser
construído em 1657 e foi
consagrado em 1672; o painel visto durante a
visita registrava precisamente
essa origem jesuíta.

Os jesuítas haviam chegado à Bahia com Tomé
de Sousa em 1549. Religião,
educação e colonização não eram esferas
independentes: a Companhia de Jesus
participava da catequese, da formação
intelectual e da reorganização colonial
da vida indígena \cite{jesuitas1549}.

A conexão com instituições modernas como a
PUC--Rio e a Universidad
Iberoamericana é institucional e histórica:
todas pertencem à tradição
educativa jesuíta, embora separadas por
séculos e contextos radicalmente
diferentes.

\section{Açúcar, Recôncavo e porto}
\label{historia-do-brasil-section-acucar-reconcavo-salvador}

A importância inicial de Salvador está ligada
ao Recôncavo e à economia do
açúcar. Plantações e engenhos em torno da
Baía de Todos-os-Santos produziam
para o circuito atlântico, enquanto Salvador
funcionava como centro
administrativo, comercial e portuário.

Essa economia dependeu de trabalho
coercitivo. Houve escravização indígena e,
em escala crescente, tráfico atlântico de
africanos escravizados. Uma versão
posterior destas notas deve reconstruir
quantitativamente as mudanças nas
regiões africanas de embarque e sua relação
com a Bahia.

\section{Rosário dos Pretos}
\label{historia-do-brasil-section-historia-rosario-pretos}

A Igreja de Nossa Senhora do Rosário dos
Pretos estava fechada no dia da
visita, portanto o interior não foi
conhecido. A história da instituição,
porém, é central.

A Irmandade dos Homens Pretos foi erigida em
1685. A autorização para construir
sua própria igreja foi dada em 1704, e o
edifício foi construído ao longo de
décadas por pessoas negras escravizadas e
libertas
\cite{ipac-rosario}.

Essas irmandades não eram apenas grupos de
oração. Organizavam culto, festas,
sepultamento, ajuda mútua e formas de
pertencimento coletivo dentro de uma
sociedade escravista. A pergunta histórica
que fica é: como pessoas negras
criaram instituições próprias dentro da ordem
católica colonial?

\section{África não é uma origem única}
\label{historia-do-brasil-section-africa-nao-uma-origem}

O mapa do MUNCAB tornou muito concreta uma
correção metodológica: não existe
uma única origem africana para a cultura
afro-brasileira.

O mapa distinguia, entre outras regiões:
\begin{itemize}
\item áreas iorubás entre os atuais Benin e
      Nigéria;
\item regiões Akan em torno de Gana e Costa
      do Marfim;
\item povos da região Camarões--Gabão;
\item várias sociedades da África Central nas
      zonas hoje divididas entre
Congo e Angola.
\end{itemize}

Essas regiões culturais e linguísticas não
coincidem com as fronteiras dos
Estados modernos. As categorias brasileiras
Nagô, Jeje, Congo e Angola são
produtos de histórias diaspóricas e não
simples nomes de nacionalidades
africanas contemporâneas.

Um detalhe importante da visita foi uma
contradição interna no museu: uma
legenda chamou uma máscara Akan de Africa
Central, enquanto o próprio mapa
do MUNCAB situava a região Akan em Gana/Costa
do Marfim, isto é, na África
Ocidental. A divergência deve ser preservada
como exemplo de leitura crítica
de museus.

\section{Candomblé como formação diaspórica}
\label{historia-do-brasil-section-historia-candomble-formacao}

O Candomblé é brasileiro, mas sua formação é
inseparável das histórias
africanas e do tráfico atlântico. Não surgiu
pela simples cópia de um sistema
africano intacto: populações diferentes
preservaram, reorganizaram e
transformaram divindades, línguas rituais,
música, comida, plantas, parentesco
e formas de autoridade em condições novas e
violentas.

Pesquisas históricas indicam que formações
religiosas jejes e nagôs já estavam
consolidadas em Salvador na década de 1860, e
que a posterior hegemonia nagô
foi ela própria um resultado histórico, não
uma essência eterna do Candomblé
baiano \cite{pares2010nago}.

\section{Bahia, Benin e fluxo/refluxo}
\label{historia-do-brasil-section-bahia-benin-fluxo}

A Casa do Benin existe para tornar visível
uma relação histórica específica:
Bahia--Golfo do Benin. Não é simplesmente um
museu genérico sobre a África.

A República do Benin atual aparece no mapa
como um país estreito entre Togo e
Nigéria. Kétou aparece ali como lugar
concreto, e não apenas como nome ritual
reencontrado no Candomblé brasileiro.

Pierre Verger organizou a relação entre as
duas margens através da ideia de
fluxo e refluxo \cite{verger2021fluxo}. A
pergunta deixa de ser apenas o
que veio da África para o Brasil e passa a
incluir:
\begin{itemize}
\item o que foi trazido coercivamente da
      África para a Bahia;
\item o que se transformou na Bahia;
\item o que mais tarde voltou do Brasil para
      a África através de retornos,
viagens religiosas, pesquisa, comércio e
intercâmbio cultural.
\end{itemize}

A própria Casa do Benin foi inaugurada em
1988 como espaço de intercâmbio, em
paralelo com a Casa do Brasil em Ouidah
\cite{casa-benin-visite}.

\section{Zumbi e memória da resistência}
\label{historia-do-brasil-section-historia-zumbi}

Zumbi dos Palmares é uma pessoa histórica
associada à liderança do Quilombo
dos Palmares no século XVII e se tornou um
símbolo central da resistência
negra à escravidão.

Durante a visita, seu nome evocou
\emph{Nzambi}/\emph{Nzambi a Mpungu}, nome
divino em tradições centro-africanas. Essa
semelhança sonora não autoriza
identificar os dois termos; a questão
etimológica fica no caderno de
linguística.

\section{Mãe Stella na história
contemporânea}
\label{historia-do-brasil-section-historia-mae-stella}

Mãe Stella de Oxóssi não pertence apenas à
história interna de um terreiro.
Sua trajetória cruza saúde pública, religião,
viagens entre Brasil e África,
produção intelectual, educação, patrimônio e
debate público.

Foi a quinta ialorixá do Ilê Axé Opô Afonjá,
uma das casas mais importantes da
tradição Ketu/Nagô. Trabalhou por décadas na
área da saúde, viajou a centros
iorubás na Nigéria, participou de redes
internacionais da tradição dos
orixás, publicou livros e ingressou em 2013
na Academia de Letras da Bahia,
como primeira ialorixá e primeira mulher
negra a ocupar uma cadeira na
instituição \cite{mae-stella-alb}.

O manifesto de 1983 contra a subordinação do
Candomblé ao catolicismo também
pertence à história política e intelectual do
Brasil: praticantes discutindo
como escravidão, racismo, sincretismo,
turismo e folclorização deveriam ser
interpretados no presente.

\section{Carnaval, movimento negro e memória}
\label{historia-do-brasil-section-movimento-negro-salvador}

Os livros encontrados no MUNCAB apontaram
para a continuidade dessa história
no século XX e XXI. A obra de Alberto Pitta e
J. Cunha conecta tecido,
grafismo, blocos afro, carnaval e identidade
pública. Já
\emph{Apenas um cidadão}, de Zulu Araújo,
oferece uma memória em primeira
pessoa de décadas de movimento negro, Olodum,
políticas culturais e vida
pública baiana.

Esse tipo de fonte deve ser lido como memória
situada, e não como uma história
neutra e total. Justamente por isso é
valioso: mostra como instituições negras
modernas também fizeram Salvador.

\section{Bembé do Mercado e Recôncavo}
\label{historia-do-brasil-section-historia-bembe}

O livro sobre o Bembé do Mercado abriu uma
pista para sair de Salvador e pensar
o Recôncavo. Celebrado em Santo Amaro, o
Bembé põe em relação Candomblé,
espaço público, memória da abolição,
patrimônio e organização coletiva de
terreiros \cite{ufrb2026bembe}.

\section{Perguntas históricas abertas}
\label{historia-do-brasil-section-historia-salvador-open}

Ficam abertas, para uma etapa bibliográfica
posterior:
\begin{itemize}
\item uma cronologia quantitativa do tráfico
      atlântico para a Bahia;
\item as mudanças nas regiões africanas de
      origem ao longo dos séculos;
\item a relação entre escravização indígena,
      escravização africana e economia
açucareira;
\item o crescimento comparado de Salvador,
      Rio de Janeiro, São Paulo,
Pernambuco, Amazônia e Sul;
\item a afirmação vista no MUNCAB de que a
      Bahia seria o estado mais negro
fora da África, que só deve ser repetida
depois de definir população,
estatística e conjunto de comparação.
\end{itemize}


% --- End historia-do-brasil.tex ---

% --- Begin genealogia-familiar.tex ---
\chapter{Genealogía familiar}
\label{chapter-genealogia-familiar}





\section{Mapa general}
\label{genealogia-familiar-section-genealogia-mapa-general}

Estas notas reúnen lo que voy averiguando sobre
las familias González Casanova, Henríquez Ureña
y Lombardo Toledano. La regla aquí es distinguir
entre datos documentados, recuerdos familiares y
conjeturas pendientes. En particular, un apellido
no basta para decidir un origen geográfico.

La parte del árbol que por ahora está mejor
establecida puede visualizarse así:

\begin{center}
\xymatrix@C=1em@R=2em{
\txt{Pedro\\Henríquez Ureña} \ar@{-}[dr]
  && \txt{Isabel\\Lombardo Toledano}
  \ar@{-}[dl]
\\
& \txt{Natacha\\Henríquez Lombardo}
  \ar@{-}[r]
  & \txt{Pablo González\\Casanova\\
  (1922--2023)}
}
\end{center}

Para entender de dónde vienen esos apellidos hay
que abrir el árbol hacia arriba. Dos ramas son
especialmente claras: Henríquez Ureña conduce a
República Dominicana y Curazao, y Lombardo conduce
a un inmigrante italiano del siglo XIX. La rama
Casanova está documentada en Yucatán desde el
siglo XIX, pero su origen anterior todavía está
abierto.

\section{Los González Casanova}
\label{genealogia-familiar-section-genealogia-gonzalez-casanova}

Pablo González Casanova (1922--2023), rector de
la UNAM entre 1970 y 1972, fue hijo del lingüista
y filólogo Pablo González Casanova y de Concepción
del Valle Romo. Este parentesco aparece también en
las semblanzas institucionales de la UNAM.\footnote{%
UNAM, ``Valioso legado de Don Pablo'',
\url{https://www.cch.unam.mx/comunidad/valioso-legado-de-don-pablo}.}

En la generación anterior aparece la unión de dos
líneas: González y Casanova. El padre del lingüista
Pablo González Casanova fue Pablo González Aznar y
su madre Encarnación Casanova Castillo. Por ahora,
la rama Casanova que he podido reconstruir de modo
provisional es

\begin{center}
\xymatrix@C=1em@R=1.7em{
\txt{José\\García-Casanova} \ar@{-}[dr]
  && \txt{Rosa Vázquez} \ar@{-}[dl]
\\
& \txt{Francisco\\García-Casanova Vázquez}
  \ar@{-}[d]
\\
& \txt{Francisco\\García-Casanova Zaldívar}
  \ar@{-}[d]
\\
& \txt{Encarnación\\Casanova Castillo}
  \ar@{-}[d]
\\
& \txt{Pablo González\\Casanova\\
  (1889--1936)} \ar@{-}[d]
\\
& \txt{Pablo González\\Casanova\\
  (1922--2023).}
}
\end{center}

La investigación pendiente importante es encontrar
el origen de José García-Casanova. El apellido
\emph{Casanova} existe tanto en Italia como en la
Península Ibérica; por eso no hay todavía base para
decir que esta rama sea italiana. Lo que sí sabemos
es que la familia García-Casanova estaba ya asentada
en Yucatán en el siglo XIX.

\section{Pedro Henríquez Ureña}
\label{genealogia-familiar-section-genealogia-henriquez-urena}

Pedro Henríquez Ureña (1884--1946) nació en Santo
Domingo. Fue hijo de Francisco Henríquez y Carvajal
y de la poeta Salomé Ureña. Se casó con Isabel
Lombardo Toledano en México. Tuvieron dos hijas:
Natacha, nacida en México, y Sonia, nacida en
Argentina. Este último dato está confirmado por una
entrevista con Sonia Henríquez Ureña conservada por
El Colegio de México.\footnote{%
Mary Carmen Sánchez Ambriz, ``El acervo de Pedro
Henríquez Ureña a El Colegio de México. Entrevista
con Sonia Henríquez Ureña de Hlito'',
\url{https://libros.colmex.mx/wp-content/uploads/2021/04/boled_122.pdf}.}

La línea paterna de Pedro conduce a Curazao. Su
abuelo Noel Henríquez pertenecía a una familia
curazoleña. Hay una tradición familiar y
bibliográfica que relaciona esta rama con la
diáspora sefardí del Caribe y de los Países Bajos.
Éste es un punto que merece una investigación
archivística separada: antes de prolongar el árbol
hasta Europa quiero verificar los enlaces uno por
uno en registros de Curazao y, si es posible, en
archivos de comunidades sefardíes neerlandesas.

La línea materna de Pedro es la de Salomé Ureña,
una de las figuras fundamentales de la literatura y
la educación dominicanas. De este modo, esta rama
del árbol no es sólo genealógica: está directamente
ligada a la historia intelectual de la República
Dominicana y, por Pedro, a la historia intelectual
de México.

\section{Lombardo Toledano: la rama italiana}
\label{genealogia-familiar-section-genealogia-lombardo-toledano}

Aquí sí aparece una procedencia italiana
explícitamente documentada. Vicente Lombardo
Toledano (1894--1968) fue hijo de Vicente Lombardo
Carpio e Isabel Toledano. Según la cronología
biográfica preparada por Rosa María Otero y Gama de
Lombardo, sus abuelos paternos fueron \emph{Vincenzo
Lombardo Cati}, originario de Turín, Italia, y
Marcelina Carpio. Vincenzo llegó a México en
1859.\footnote{%
Centro de Estudios Filosóficos, Políticos y Sociales
Vicente Lombardo Toledano, \emph{Escritos
autobiográficos}, cronología de Rosa María Otero y
Gama de Lombardo,
\url{https://www.centrolombardo.edu.mx/%
wp-content/uploads/formidable/14/%
Escritos_autobiogr\%C3\%A1ficos_REV.pdf}.}

La misma fuente dice que los abuelos maternos de
Vicente eran Miguel Toledano y Flora Toledano,
descendientes de una antigua familia española cuyo
fundador había llegado a México en el siglo XVII.
Esto es especialmente interesante porque impide
identificar sin más ``Lombardo Toledano'' con una
sola procedencia: la rama Lombardo es italiana,
mientras que la rama Toledano está descrita allí
como española.

Isabel Lombardo Toledano, hermana de Vicente, se
casó con Pedro Henríquez Ureña. Una versión mínima
de esta rama es

\begin{center}
\xymatrix@C=1em@R=1.8em{
\txt{Vincenzo\\Lombardo Cati} \ar@{-}[dr]
  && \txt{Marcelina\\Carpio} \ar@{-}[dl]
\\
& \txt{Vicente\\Lombardo Carpio}
  \ar@{-}[dr]
  && \txt{Isabel\\Toledano} \ar@{-}[dl]
\\
&& \txt{Isabel\\Lombardo Toledano}
  \ar@{-}[d]
\\
&& \txt{Natacha\\Henríquez Lombardo.}
}
\end{center}

El siguiente paso para esta rama es localizar el
acta de nacimiento de Vincenzo en el Piamonte y
reconstruir sus padres y abuelos. Sólo entonces se
podrá precisar más que ``originario de Turín''.

\section{Una red intelectual mexicana}
\label{genealogia-familiar-section-genealogia-red-intelectual}

La genealogía se cruza de manera llamativa con la
historia intelectual mexicana. Pedro Henríquez Ureña
formó parte del núcleo del Ateneo de la Juventud
junto con figuras como Antonio Caso, Alfonso Reyes
y José Vasconcelos. Décadas después, Vicente
Lombardo Toledano fue alumno de Antonio Caso y se
convirtió en una figura central del debate político
y universitario mexicano.

Las relaciones entre Caso, Lombardo, Vasconcelos y
Henríquez Ureña fueron además personales e
institucionales, no sólo bibliográficas. La propia
cronología histórica de la UNAM registra, por
ejemplo, los conflictos universitarios de 1923
entre Lombardo y Vasconcelos, la renuncia de
Antonio Caso y la salida de Pedro Henríquez Ureña
de sus cargos docentes.\footnote{%
UNAM, cronología histórica, ``1920'', sección sobre
1923,
\url{https://www.unam.mx/acerca-de-la-unam/unam-en-el-tiempo/cronologia-historica-de-la-unam/1920}.}

Esto sugiere que una historia familiar completa no
debe ser sólo una lista de ascendientes. Hay también
una historia de redes: parentesco, amistad,
magisterio, instituciones y disputas intelectuales.

\section{Recuerdos familiares por verificar}
\label{genealogia-familiar-section-genealogia-recuerdos-familiares}

Estas notas no se toman todavía como hechos
históricos. Se conservan porque pueden orientar la
búsqueda documental.

\begin{itemize}
\item En recuerdos familiares aparecen los nombres
Charles y Mara, Henrique y Lolita, Manuel y Helena,
y el apodo ``Titopaul'' para Pablo. Falta fijar con
precisión qué generación y qué parentesco designa
cada referencia.
\item Aparece también una ``tía Sonia''. Esto es
compatible con Sonia Henríquez Ureña, hermana de
Natacha, pero no debe identificarse automáticamente
sin contexto adicional.
\item Hay un recuerdo de una conexión familiar de
los Bracho--Carpizo--Valle con Antonio Caso. Todavía
no hay identificación documental suficiente para
dibujar esa rama.
\end{itemize}

\section{Preguntas abiertas}
\label{genealogia-familiar-section-genealogia-preguntas-abiertas}

Las tres búsquedas principales que quedan abiertas
son:

\begin{enumerate}
\item \textbf{García-Casanova.} Encontrar el origen
de José García-Casanova y decidir documentalmente
si la línea anterior conduce a España, Italia u
otro lugar.
\item \textbf{Henríquez.} Reconstruir la línea de
Curazao generación por generación y comprobar la
tradición de ascendencia sefardí con registros
primarios.
\item \textbf{Lombardo.} Localizar a Vincenzo
Lombardo Cati en registros piamonteses y extender
la rama italiana hacia atrás.
\end{enumerate}

Estas notas deben actualizarse conforme aparezcan
actas civiles, parroquiales o documentos de archivo;
los árboles genealógicos en línea sirven como pistas,
no como prueba final.

% --- End genealogia-familiar.tex ---

% --- Begin history.tex ---
\chapter{History}
\label{chapter-history}





\section{Napoleón en España y en Nueva España
la retirada de Hidalgo}
\label{history-section-retirada-de-hidalgo}

% The original title, byline, and body below
% are transcribed verbatim from the DOCX.
% Only LaTeX markup and whitespace differ.
% The attribution footnote is editorial.

L. Barjau\footnote{Luis Barjau.
Texto reproducido íntegramente del documento
\emph{La retirada de Hidalgo.docx}.}

En 1807 el ejército de Napoleón Bonaparte
ingresó a España buscando la alianza para
invadir Portugal, quien se negó a cumplir
bloqueo comercial contra Gran Bretaña.
Napoleón quería someter a los ingleses
mediante guerra financiera.

Al año siguiente Fernando VII el rey Borbón
ascendió al trono de España. Pero fue
expulsado por Napoleón, que nombró a su
hermano José I al reinado que duró cinco
años, hasta 1813. Fernando y su padre, Carlos
IV, tuvieron que abdicar y fueron hechos
prisioneros en Francia. Este hecho socavó los
cimientos de la monarquía española lo cual
repercutió en todo el mundo hispano. Las
Juntas regionales y descentralizadas se
negaron a reconocer la legitimidad de los
Bonaparte y este movimiento se extendió hasta
las colonias americanas. En 1809 hubo
rebeliones en los Andes. Se establecieron
tribunales generales tanto en España como en
América. En la península, la Junta Central
fue empujada hacia el sur, en 1809 se
estableció en Sevilla, pero de allí tuvo que
huir a Cádiz donde recibió la protección de
naves británicas. Se establecieron Cortes que
incluyeron los reinos de América y con ellas
España obtuvo mucho mayor control del
habitual. Los españoles radicados en las
colonias rechazaron el concepto de igualdad
que las cortes consideraban y los criollos
prefirieron que los \emph{indios}
permanecieran controlados por España. Este
conflicto también fue un antecedente de la
Independencia. Los españoles de América se
mantuvieron en la actitud de que indígenas y
afrodescendientes no estuvieran representados
en las cortes españolas. Las cortes eran los
órganos legislativos que representaban a los
sectores de todos los reinos ibéricos frente
al Rey. Institución nacida desde la Edad
Media.

La actitud de los españoles americanos fue un
gesto revelador de la voluntad europea
respecto del Otro, el distinto, del
“salvaje”, del extraño a sus costumbres. Que
habría de permear toda la relación política y
cultural de Europa con América. Toda la
mentalidad del hombre americano. Del criollo
enriquecido, con su resentimiento contra el
peninsular español; de las castas, el pueblo,
los pobres, la confusión identitaria.

La posible representación había estado
liderada en Quito por José Mejía Lequerica
que propuso igualdad para indios y mestizos
excluyendo a africanos.

Hasta 1812 el término “americano”, refería
solo a los blancos europeos. A partir de ese
año se aceptó tal designación también para
afrodescendientes, indígenas y mestizos, lo
que así quedó consignado en la Constitución
de Cádiz, la primera de España. Y esta
influyó aunque de manera ambigua en la
emancipación mexicana. Pues introdujo los
principios democráticos que sembraron las
ideas de la soberanía y ello estimuló las
tendencias independentistas. Surgieron los
ayuntamientos constitucionales y las
elecciones locales. Por primera vez existía
el voto y la autogestión que combatiría al
absolutismo real, de haberse ejercido al pie
de la letra. En 1820, restablecida la
Constitución en España, el clero, militares y
terratenientes novohispanos temieron la
pérdida de sus privilegios por esas nuevas
leyes de las Cortes. En la guerra de
independencia mexicana algunos personajes
realistas pactaron con Vicente Guerrero con
la intención de mantener el poder. Que
aceptaban las leyes de Cádiz pero excluían
aquellas de orden decididamente liberales.
Diputados novohispanos como Miguel Ramos
Arizpe viajaron a Cádiz y su experiencia
impulsó la Constitución mexicana de 1824.

En 1810 se formaron Juntas en Caracas, Buenos
Aires y Santiago. Hubo un vacío de poder
tanto en España como en las colonias, que
propició la creación de tales Juntas, que
tomaban el control de las diferentes
regiones, porque la resistencia en España al
mandato de Bonaparte estaba fracasando. En
Sudamérica Mariano Moreno asumió el liderazgo
que intentó una Confederación que se
inspiraba en la revolución de independencia
de EUA. Los virreinatos hispanoamericanos
trataron de evitar estas Juntas. El 20 de
julio de 1810 se promulgó la Declaración de
Independencia de Colombia. Pero desde el 1 de
febrero de 1804, gracias a los líderes
Toussaint Louverture y Jean Jacques
Dessalines se había proclamado la
independencia de Haití. Y aún antes, en 1776
Estados Unidos se había independizado. El
devenir de las paradojas: España y Francia
fueron los primeros aliados de este
movimiento de liberación en Norteamérica.

El 16 de septiembre de 1810, con el grito de
Dolores, Miguel Hidalgo y Costilla dio inicio
a la revolución de Independencia en Nueva
España, la cual, después de once años de
lucha se declararía su triunfo el 27 de
septiembre de 1821.

Una larga batalla de diez años, como el
periodo antiguo de la guerra de Troya. Sólo
que en América la guerra estuvo animada
violentamente también por el rencor indígena
porque Occidente había truncado el desarrollo
de la milenaria cultura mesoamericana. Las
batallas fueron activadas por la turbamulta
desarmada, enajenada a una rabia histórica
que prefería estrellarse contra el ejército
novohispano hasta el sacrificio y la muerte.
Retumbaba en la guerra la voluntad de
eliminar a los infiltrados españoles,
despojarlos, matarlos. Miguel Hidalgo no
calculó al Polifemo que había despertado de
su antiguo sueño en las cuevas de
Chicomoztoc.

A medida que se dirigía la revolución hacia
el centro de la Nueva España el pelotón se
engrosaba como bola de nieve en picada.
Contingente de 80,000 campesinos con armas
improvisadas. Destrozar las nuevas ciudades
hispanas, liquidar a los peninsulares que
mandaban a diestra y siniestra y se
enriquecían con los bienes naturales y con la
explotación de la mano de obra.

Hidalgo no cesó de asombrarse ante la
proporción del fenómeno que había despertado.
Y cuando llegó a las puertas de la ciudad, en
La Marquesa, Monte de las Cruces (no obstante
después de haber sido recibido en Toluca con
la misma pompa de recepción que se daba a los
virreyes) y ya al borde de la gran ciudad
capital, habiendo ganado la batalla, ¡tomó la
retirada!

Abrumadora superioridad numérica la de los
independentistas, contra el ejército
virreinal organizado y al mando de Torcuato
Trujillo y Rodríguez, que había peleado en
Alcalá la Real, de Jaén, contra Napoleón
Bonaparte. Aunque el ejército realista fuera
derrotado, el choque desigual dejó un saldo
campesino de cerca de 4000 muertos.

Sin explicar nada a nadie el cura Hidalgo
regresó por el camino que lo había conducido
hasta el centro del poder hispano. “Viva
Fernando VII”, había entremezclado un año
antes en su grito inicial por la
Independencia. “Viva nuestra señora de
Guadalupe, viva Fernando VII, muera el mal
gobierno” fue la proclama completa. Confusión
y desconcierto. A favor del rey español, a
favor del cristianismo sincretizado en la
guadalupana; pero contra del poder virreinal.
¿Qué se buscaba? ¿Qué era posible? El nombre
del rey todavía implorado contra la invasión
francesa en España.

La complejidad de cualquier proceso
histórico, en la Nueva España concitó su
extremada y densa problemática. Si la
invasión napoleónica en la Península hispana
estimuló la posibilidad de que en sus reinos
americanos se observara la opción
independiente; que los criollos vislumbraran
un nacionalismo inicial; que los indígenas
exacerbaran una caótica intervención
vengativa, caótica porque a la vez afirmaban
el derecho a un catolicismo dogmático
amparado en el mito sincrético de la virgen
guadalupana; que el cura Hidalgo encabezara
el surgimiento de un fervor revolucionario
que se gestaba desde la ideología pietista
cristiana, a la vez que en la observación de
la posibilidad de la riqueza local, pero sin
la enorme extracción de bienes y servicios
ejercida por la corona española, la tragedia
mexicana de la guerra independentista habría
de durar 11 años.

La decisión de retirada en el Monte de las
Cruces la había de pagar muy cara el párroco
de Dolores. En Acatita de Baján, Coahuila, lo
tomarían prisionero; fusilado el 30 de julio
de 1811 en el patio del antiguo hospital Real
de Loreto, hoy palacio de gobierno de
Chihuahua. Otra historia hubiera sido si
Hidalgo avanzara hasta la capital. Para bien
o para mal.

Diez largos años de guerra y todas las
traiciones y fusilamientos de
independentistas, que conocemos. Fantasmas y
leyendas. Iturbide y “el abrazo de
Acatempan”. Iturbide era el Jacob de México.

Con la independencia, el país se revolvió con
su desesperado intento de convertirse en una
nación. ¿Volver al ámbito espiritual de la
cultura indígena? ¿Construir cimientos
nacionales, con cuáles modos de producción?
¿Erguirse en la explotación del subsuelo? ¿La
agricultura? ¿El comercio? ¿Como las naciones
europeas?

En las raíces de México tuvieron que arraigar
las severas contradicciones. En extremo
complejas. Sin que nadie tuviera plena
conciencia de ello. ¿Acaso Juárez fuera el
primero en vislumbrar un camino cierto? El
indígena de Guelatao que subió hasta la silla
presidencial el 21 de enero de 1858; 37 años
después de haberse consumado la
independencia. ¡Juárez! ¿Qué pensó
íntimamente? ¡Mucho daríamos por haber
conocido dicho pensamiento! Por haber podido
discutir en aquellos términos. Juárez íntimo.
Juárez secreto. El milagro de la nación
mexicana. ¿La extensión del efluvio
prehispánico? La esencia revolucionaria de
esta nación: cambiar, cambiar a toda costa,
buscar un camino de esplendor para el viejo
dolor mexicano que se arrastra por la rivera
del rio de su historia.

% --- End history.tex ---

\bibliographystyle{amsalpha}
\bibliography{my}
\end{document}
